\documentclass[11pt,a4paper,leqno,oneside,openany]{book}

\usepackage[T1]{fontenc}
\usepackage[utf8]{inputenc}
\usepackage[provide=*,french,english]{babel}
\usepackage{lmodern}
\usepackage{amsmath,amssymb,amsfonts,mathrsfs}
\usepackage{mathtools}
\usepackage{tikz}
\usetikzlibrary{cd,arrows.meta}
\usepackage{geometry}
\geometry{left=21mm,right=21mm,top=22mm,bottom=23mm,marginparwidth=15mm,marginparsep=3mm}
\usepackage{enumitem}
\usepackage[normalem]{ulem}
\let\mathunderline\underline
\renewcommand{\underline}[1]{\ifmmode\mathunderline{#1}\else\uline{#1}\fi}
\emergencystretch=3em

\DeclareUnicodeCharacter{2192}{\ensuremath{\rightarrow}}
\DeclareUnicodeCharacter{2297}{\ensuremath{\otimes}}
\DeclareUnicodeCharacter{2295}{\ensuremath{\oplus}}
\DeclareUnicodeCharacter{2228}{\ensuremath{\vee}}
\DeclareUnicodeCharacter{223C}{\ensuremath{\sim}}
\DeclareUnicodeCharacter{2248}{\ensuremath{\approx}}
\DeclareUnicodeCharacter{00D7}{\ensuremath{\times}}
\DeclareUnicodeCharacter{03B1}{\ensuremath{\alpha}}
\DeclareUnicodeCharacter{03B2}{\ensuremath{\beta}}
\DeclareUnicodeCharacter{00A7}{\S}
\DeclareUnicodeCharacter{010C}{\v{C}}
\DeclareUnicodeCharacter{00C3}{\~{A}}
\DeclareUnicodeCharacter{02D9}{\.{}}

\providecommand{\Zl}{\mathbb{Z}_{\ell}}
\providecommand{\Ql}{\mathbb{Q}_{\ell}}
\providecommand{\Gal}{\operatorname{Gal}}
\providecommand{\Spec}{\operatorname{Spec}}
\providecommand{\Isom}{\operatorname{Isom}}
\providecommand{\Aut}{\operatorname{Aut}}
\providecommand{\Hom}{\operatorname{Hom}}
\providecommand{\Ext}{\operatorname{Ext}}
\providecommand{\Biext}{\operatorname{Biext}}
\providecommand{\Coker}{\operatorname{Coker}}
\providecommand{\Ker}{\operatorname{Ker}}
\providecommand{\Im}{\operatorname{Im}}
\providecommand{\Tor}{\operatorname{Tor}}
\providecommand{\prof}{\operatorname{prof}}
\providecommand{\codim}{\operatorname{codim}}
\providecommand{\pr}{\operatorname{pr}}
\providecommand{\id}{\mathrm{id}}
\providecommand{\sym}{\mathrm{sym}}
\providecommand{\Corr}{\operatorname{Corr}}
\providecommand{\Proj}{\operatorname{Proj}}
\providecommand{\Comp}{\operatorname{Comp}}
\providecommand{\cC}{\mathcal{C}}
\providecommand{\cY}{\mathcal{Y}}
\providecommand{\cT}{\mathcal{T}}
\providecommand{\cPsi}{\Psi}
\providecommand{\Dcat}{D}
\providecommand{\La}{\Lambda}
\providecommand{\Ghat}{\widehat{G}}
\providecommand{\uHom}{\underline{\operatorname{Hom}}}
\providecommand{\cl}{\operatorname{c\ell}}
\providecommand{\e}{\mathrm{e}}




\usepackage{etoolbox,adjustbox,marginnote,xurl,hyperref}
\tikzcdset{ampersand replacement=\&}
\tikzset{ampersand replacement=\&}
\hypersetup{hidelinks,hypertexnames=false,pdfauthor={GPT-6 Astra; source authors credited in the edition},pdftitle={SGA 7 I et II - lecteur francais}}
\BeforeBeginEnvironment{tikzcd}{\begin{adjustbox}{max width=\linewidth}}
\AfterEndEnvironment{tikzcd}{\end{adjustbox}}
\setcounter{tocdepth}{1}
\setlength{\emergencystretch}{4em}
\newcount\SGAform
\SGAform=2
\ifnum\SGAform=0
  \long\def\SGAvar#1#2#3{#1}
\else\ifnum\SGAform=1
  \long\def\SGAvar#1#2#3{\SGAvmark{#3}#2}
\else
  \long\def\SGAvar#1#2#3{#2}
\fi\fi
\newcommand{\SGAvmark}[1]{\ifstrempty{#1}{}{\leavevmode\textsuperscript{\tiny\hyperlink{app-#1.1}{[#1]}}\ignorespaces}}
\newcommand{\SGApage}[2]{\marginnote{\normalfont\scriptsize #1\,\textbf{#2}}}
\newcommand{\SGAtoc}[1]{\phantomsection\addcontentsline{toc}{chapter}{#1}}
\newcommand{\SGAtocunderline}[1]{\ifmmode\mathunderline{#1}\else #1\fi}
\newcommand{\SGAsectiontoc}[2]{\phantomsection\addcontentsline{toc}{section}{\texorpdfstring{#1}{Section #2}}}

\begin{document}
\selectlanguage{french}
\frontmatter
\begin{titlepage}\centering
{\Large Séminaire de géométrie algébrique du Bois Marie\par}\vspace{20mm}
{\Huge\bfseries SGA 7\par}\vspace{8mm}
{\LARGE Groupes de monodromie en géométrie algébrique\par}\vspace{8mm}
{\Large Volumes I et II\par}\vspace{16mm}
{\Large \ifcase\SGAform Texte diplomatique\or Texte corrigé avec apparat\or Texte de lecture\fi\par}\vfill
{\large Édition du 12 novembre 2026\par}
{\large GPT-6 Astra (OpenAI, ChatGPT Pro mode)\par}
\end{titlepage}
\section*{Présentation de cette forme}
\ifcase\SGAform Le texte suit le témoin diplomatique fourni ; les coquilles de l'imprimé sont conservées. Les réparations de transcription déjà établies restent présentes. Les repères de chaque page de l'imprimé sont indiqués en marge.\or Le texte suit le témoin corrigé fourni et les décisions documentées de cette édition. Chaque passage différent du témoin diplomatique est marqué ; chaque différence individuelle est donnée dans l'apparat de l'exposé avec les deux lectures et sa source. Les notes indécidables sans scan sont séparées des variantes effectives.\or Le texte est exactement le texte corrigé. Cette forme omet les marques de variantes, l'apparat et les notes éditoriales de lecture indécidable. Elle conserve intégralement le texte, les notes des auteurs, les bibliographies, le catalogue de l'éditeur et les seuls repères des folios imprimés en marge.\fi
Les numéros marginaux renvoient exclusivement aux folios de l'imprimé : I~1--528 et II~1--438. L'introduction générale porte les repères I~V--VII (V est inféré dans les sources). La pagination du présent volume est indépendante.
\section*{Identification bibliographique}

Alexander Grothendieck, avec M. Raynaud et D. S. Rim, Groupes de monodromie en géométrie algébrique I, Séminaire de géométrie algébrique du Bois-Marie 1967–1969 (SGA 7 I), Lecture Notes in Mathematics 288, Springer, 1972 : exposés I, II, VI–IX. Pierre Deligne et Nicholas Katz, Groupes de monodromie en géométrie algébrique II (SGA 7 II), Lecture Notes in Mathematics 340, Springer, 1973 : exposés X–XXII. Les attributions détaillées des exposés sont conservées. L’exposé VI de D. S. Rim a été publié en anglais et reste en anglais dans les deux éditions linguistiques.

\section*{Crédits et provenance}

Cette construction est attribuée à GPT-6 Astra (OpenAI, ChatGPT Pro mode). La source française est AI-transcribed ; l’anglais est AI-translated, à l’exception de l’exposé VI dont l’anglais original est conservé. L’édition comporte des AI-produced corrections. Ces désignations n’attribuent à l’intelligence artificielle aucune contribution humaine. La source anglaise est la version fournie de sga-en, v2026-11-12-english-edition, commit 4a4cbac ; les couches françaises diplomatique et courante, journaux de corrections et décisions sont les enregistrements fournis de novembre 2026. La couche française d’août sert uniquement de référence au journal des changements. Les extraits anglais ci-dessous conservent les crédits des enregistrements, y compris ceux concernant les autres volumes du projet SGA ; ces derniers ne sont pas attribués à SGA 7.

\subsection*{Contributions humaines : enregistrement original}

Sources and authorship

The original SGA volumes retain their authors, seminar contributors and editors, named in the corresponding texts. The collection includes work of Alexander Grothendieck, Michèle Raynaud, Michel Raynaud, Pierre Deligne, Nicholas Katz, Michel Demazure, Michael Artin, Jean-Louis Verdier, Luc Illusie and their collaborators; this list does not replace the detailed credits in each volume.

SGA 1 retains the French retypesetting associated with Bas Edixhoven’s volunteer project, the Société mathématique de France and Michel Raynaud’s notes, from arXiv math/0206203v2 (\url{https://arxiv.org/abs/math/0206203v2}).

SGA 2 is AI-translated from the corrected French TeX of the re-edition edited by Yves Laszlo, Cohomologie locale des faisceaux cohérents et théorèmes de Lefschetz locaux et globaux (SGA 2), Documents Mathématiques 4, Société mathématique de France, 2005 (arXiv math/0511279v1 (\url{https://arxiv.org/abs/math/0511279v1})), and follows its numbering and labels. Most of the initial French typesetting was done by volunteers of Bas Edixhoven’s project (L. Bayle, N. Borne, O. Brinon, J. Buresi, M. Chardin, F. Ducrot, P. Graftiaux, F. Han, P. Karwasz, L. Koelblen, D. Madore, S. Morel, D. Naie, B. Osserman, J. Riou and V. Sécherre); C. Sabbah adapted it to the SMF format.

The SGA 2 and SGA 3 English sources acknowledge Jacob C. Reinhold’s English translation (\url{https://github.com/jcreinhold/sga/tree/e7a259f3f8608ad3edf9bf6eead3fd504dd2d23e}), snapshot \url{e7a259f3f8608ad3edf9bf6eead3fd504dd2d23e}, used for comparison and drafting. His stated CC BY 4.0 notice applies to that contribution and is preserved; it does not license the underlying French works or every contribution here. Inherited English wording remains attributed to its human source. In SGA 3, lemmas, proofs and corrections added in the Philippe Gille–Patrick Polo re-edition are given in Translator's notes and Additions that cite it.

SGA 4 is AI-translated from the French LaTeX edition of SGA 4, Tomes 1–3, distributed by Fabrice Orgogozo (normalesup.org/\textasciitilde{}forgogozo/SGA4 (\url{https://www.normalesup.org/~forgogozo/SGA4/}); source snapshot at git commit 71766d9), and uses its TeX configuration files. Its footnotes headed “N.D.E.” are that edition's editorial notes, translated.

SGA 4½ preserves Daniel Miller’s LaTeX typesetting (\url{https://github.com/dkmiller/sga4.5}) and its original typesetter’s note. It is not attributed to AI.

Additional English translation is AI-translated; subsequent corrections, editing, reconstruction and typesetting are AI-produced. The source-specific human credits and genuine third-party notices remain applicable; attribution does not establish permission or endorsement.


\subsection*{Historique de production enregistré}

Jacob C. Reinhold's English translation was used for comparison and drafting in SGA 2 and SGA 3. Recorded AI work includes OpenAI Codex (GPT-5.6 Sol: translation; GPT-6 Astra: corrections and assembly, September–October 2026) and Claude Opus 5.5 (Anthropic: the SGA 2 and SGA 3 sources of 26 October 2026, and the translation of SGA 3 revised against the 1970 printing in November 2026). GPT-6.1 Sol (OpenAI API) proposed the edits that carry the corrections of the scan check of the French text into SGA 4 (November 2026); each edit was checked by hand. The release notes record the correction rounds: spectrum notation in SGA 1 (25 September 2026), mathematical and reference corrections in SGA 1 and SGA 7 II (3 October 2026), references in SGA 3 (24 October 2026), the scan-check corrections in SGA 4 (12 November 2026).


\subsection*{Autorité de la source de SGA 7 I}

SGA 7 I is complete as the published volume stands. The Springer volume is explicitly Exposés I, II, VI, VII, VIII, IX. Its introductory account explains that Grothendieck's oral Exposés I–V were not separately written: their substance was summarized in the written Exposé I, and there are no separately written Exposés III–V. The jump from II to VI is therefore historical, not a translation or packaging gap. The current English component includes every exposé in the published volume. Its existing release filename is retained to preserve stable asset links.


\subsection*{Témoin anglais hérité}

The inherited witness documented in \url{en/evidence/supplemental_working_decisions/sga7i/INHERITED_ENGLISH_PROVENANCE_AND_ADVERSE_USE.md} has 278,864 bytes and SHA-256 \url{CFCB4C17081E16F963DCBC3FA04A43C4D4065477CB37BD931233840D287AA2E9}. Its exact blob first appears in the public-mirror history at commit \url{cb420190bbc6311cff96a08a5a11a4caf483fdff}, dated 2026-05-29 12:24:48+02:00, under the generic identity “Manuscript Typesetting Project”; the Git blob is \url{2ac760f78b6107d2d71bd2e513e5baa4e3be418d}. The local working paths and address are omitted for privacy.

That service identity and the surviving package metadata do not identify the translating model or human. No defensible evidence presently attributes the text specifically to Claude, Codex, ChatGPT, or a human translator. Its author/model lineage is therefore recorded as unresolved.

This witness is retained unchanged for provenance, terminology comparison, and adverse-history purposes. It must not be merged continuously, treated as cumulative coverage, or cited as source-aligned. Individual phrases may be reused only after direct matching to the fresh French authority. The active English reader is source-first; Expos\textbackslash{}'e VI will be taken from its already-English frozen source after a structural check.


\subsection*{Provenance du corpus français}

SGA7 I French editions

This is an exact no-overwrite binding of two separate source-language editions through Tome-I EOF: o is the printed-diplomatic layer and c is the separately corrected scholarly layer. They are not a bilingual or interleaved edition. The English edition remains separate.

The upstream layer proves complete 528/528 source-folio coverage and stable 267-page readers. Its own audit record is intentionally narrower than a critical edition: representative pages and consequential ambiguities were checked directly, but no exhaustive high-detail page-by-page source audit is claimed. The scan remains controlling for future adjudication.


SGA7 II French editions

This is an exact no-overwrite binding of two separate source-language editions through Tome-II EOF: o is the printed-diplomatic layer and c is the separately corrected scholarly layer. They are not a bilingual or interleaved edition. The English edition remains separate.

The bound layer covers scan indices 8--445, with 438 markers / 438 unique body pages and no missing or duplicate index, across Exposes X--XXII. Gate-00 on 2026-08-20 repaired the omitted scan index 407 / printed folio 400 (the terminal Expose-XXI bibliography) in both French layers and rebuilt both readers. The immutable closure receipt is gate00\_idx407\_sourcekeeper\_closure\_20260820.json. Changed loci and terminal surfaces have direct review evidence; this binding does not silently promote that bounded review to an exhaustive high-detail audit of every unchanged printed page.


\section*{Pages liminaires imprimées et tables des matières}

Les tables des matières de ces lecteurs sont produites à partir du texte transcrit ; elles ne sont pas la transcription des tables imprimées. Les couches fournies de SGA 7 II commencent à l’exposé X : les pages de titre, de publication et de table des matières qui le précèdent ne figurent pas dans ces couches. SGA 7 I conserve son introduction générale transcrite mais ne comporte pas de table imprimée transcrite séparément ; aucune page de titre ou de publication absente n’a été reconstituée. Les exposés III–V ne sont pas des transcriptions manquantes : leur statut historique est expliqué dans l’enregistrement de l’autorité des sources reproduit ci-dessous.

\section*{Droits}

Les droits sur l’œuvre originale appartiennent à son auteur ou aux représentants ou héritiers de celui-ci. Dans la mesure où nous détenons des droits sur la transcription, la traduction, les corrections et la composition produites par l’intelligence artificielle, nous les dédions au domaine public sous CC0 1.0. Cette édition est publiée à des fins éducatives et non lucratives, et nous ne revendiquons aucun droit que nous ne détenons pas. Si vous détenez des droits sur cette œuvre et vous opposez à sa publication ici, veuillez ouvrir un ticket GitHub et nous la retirerons.

\par\noindent\url{https://github.com/KokunoYumeto/sga-en/issues}

\par\noindent\url{https://creativecommons.org/publicdomain/zero/1.0/}

\clearpage

\tableofcontents
\mainmatter
\part{SGA 7 I}
\clearpage\SGAtoc{Introduction générale}
\selectlanguage{french}
\SGAvar{\SGApage{I}{V}





}{\SGApage{I}{V}






}{}\thispagestyle{empty}
\begin{center}
\underline{INTRODUCTION}
\end{center}

Soient $S$ un schéma et $f:X\longrightarrow S$ un morphisme de schémas. Si $f$ est
propre et lisse, et que le nombre premier $\ell$ est inversible sur $S$, les groupes de
cohomologie $\ell$-adique $H^i(X_{\bar s},\mathbb Z_\ell)$ des fibres géométriques de $f$ forment un système
local $\ell$-adique sur $S$. Pour $f$ seulement supposé propre, ces groupes sont les fibres
d'un faisceau $\ell$-adique $R^i f_*\mathbb Z_\ell$ sur $S$. La \underline{théorie des cycles évanescents} met
en relation la ramification de ce faisceau sur $S$ et les singularités de $f$.

\SGAvar{Nous ne considérerons que le cas où $S$ est un trait hensélien (= spectre d'un
anneau de valuation discrète hensélien). C'est en pratique le cas essentiel. Je renvoie
à (I 2.1) pour une description heuristique de la théorie. Soient $S=\Spec(V)$, $k(\bar\eta)$
une clôture algébrique du corps des fractions $k(\eta)$ de $V$ et $k(\bar s)$ la clôture algébrique
correspondante du corps résiduel $k(s)$ ($k(\bar s)$ est le corps résiduel du normalisé
$\bar V$ de $V$ dans $k(\bar\eta)$). Dans le cas particulier où $X$ est propre et plat sur
$S$, de dimension relative $n$, et où $f$ ne présente qu'un point de non lissité
$x\in X_s$ ; on définit des $\Gal(\bar\eta/\eta)$-modules $\phi^i$, de nature purement locale au voisinage
de $x$, nuls pour $i\notin[0,n]$ (I 4.2), et on construit une suite exacte longue
de $\Gal(\bar\eta/\eta)$-modules

}{Nous ne considérerons que le cas où $S$ est un trait hensélien (= spectre d'un
anneau de valuation discrète hensélien). C'est en pratique le cas essentiel. Je renvoie
à (I 2.1) pour une description heuristique de la théorie. Soient $S=\Spec(V)$, $k(\bar\eta)$
une clôture algébrique du corps des fractions $k(\eta)$ de $V$ et $k(\bar s)$ la clôture algébrique
correspondante du corps résiduel $k(s)$ ($k(\bar s)$ est le corps résiduel du normalisé
$\bar V$ de $V$ dans $k(\bar\eta)$). Dans le cas particulier où $X$ est propre et plat sur
$S$, de dimension relative $n$, et où $f$ ne présente qu'un point de non lissité
$x\in X_s$, on définit des $\Gal(\bar\eta/\eta)$-modules $\phi^i$, de nature purement locale au voisinage
de $x$, nuls pour $i\notin[0,n]$ (I 4.2), et on construit une suite exacte longue
de $\Gal(\bar\eta/\eta)$-modules

}{F7I-INTRO-002}\[
\cdots \longrightarrow H^i(X_{\bar s},\mathbb Z_\ell)
\xrightarrow{\mathrm{sp}} H^i(X_{\bar\eta},\mathbb Z_\ell)
\longrightarrow \phi^i
\longrightarrow H^{i+1}(X_{\bar s},\mathbb Z_\ell)
\longrightarrow \cdots
\]

(le $\Gal(\bar s/s)$-module $H^i(X_{\bar s},\mathbb Z_\ell)$ est regardé comme un
$\Gal(\bar\eta/\eta)$-module avec action triviale de l'inertie ; $\mathrm{sp}$ est la flèche de spécialisation).

\SGAvar{On donne aussi des critères, valables pour l'instant seulement en caractéristique
$0$, pour que les $\phi^i$ pour $i$ petits soients nuls (par exemple, si $X$ est de plus
localement d'intersection complète, $\phi^i=0$ pour $i\ne n$ (I. 4.5)).


}{On donne aussi des critères, valables pour l'instant seulement en caractéristique
$0$, pour que les $\phi^i$ pour $i$ petits soient nuls (par exemple, si $X$ est de plus
localement d'intersection complète, $\phi^i=0$ pour $i\ne n$ (I. 4.6)).


}{F7I-INTRO-003}\SGApage{I}{VI}


Le \underline{théorème de monodromie} affirme que l'action du sous-groupe d'inertie $I$
de $\Gal(\bar\eta/\eta)$ est quasi-unipotente : pour $T\in I$, il existe des entiers $N>0$,
$M>0$ tels que l'endomorphisme $(T^M-I)^N$ de $H^i(X_{\bar\eta},\mathbb Z_\ell)$ soit nul. On donne de ce
théorème deux démonstrations. La première (I.1.2), de nature arithmétique, s'applique
dès que les groupes de cohomologie considérés ont des propriétés de finitude raisonnables.
Le point clef est que, lorsque $k(s)$ est de type fini sur son sous-corps
premier, l'action de $I$ est quasi-unipotente pour toute représentation $\ell$-adique de
$\Gal(\bar\eta/\eta)$. La seconde démonstration, plus géométrique, requiert la pureté et la résolution
des singularités : elle n'est pour l'instant valable qu'en caractéristique $0$
ou pour un $H^1$. Elle apporte de précieuses informations sur l'exposant de nilpotence
$N$ ($N\leq i+1$ pour un $H^i$). Ainsi qu'on le verra ultérieurement, elle se prête bien,
sur $\mathbb C$, à la comparaison avec la théorie transcendante.

Ces résultats, appliqués au $H^1$ des variétés abéliennes, i.e. à leur module
de Tate, permettent d'étudier la réduction mod $p$ de celles-ci. On donne ainsi deux
démonstrations du théorème de réduction stable des variétés abéliennes selon lequel,
après ramification, la fibre spéciale connexe du modèle de Néron est extension d'une
variété abélienne par un tore. La première (I. 6) reprend la méthode de la démonstration
arithmétique du théorème de monodromie. La seconde (IX 3.6) s'appuie sur une
analyse beaucoup plus fine du modèle de Néron, et de ses propriétés de polarisation.

\SGAvar{Les exposés I à V de Grothendieck n'ont pas été rédigés. Ils ont été résumés
dans un exposé I. Les résultats énoncés y sont démontrés de façon succinte, mais
essentiellement complète.

}{Les exposés I à V de Grothendieck n'ont pas été rédigés. Ils ont été résumés
dans un exposé I. Les résultats énoncés y sont démontrés de façon succincte, mais
essentiellement complète.

}{F7I-INTRO-004}L'exposé II applique la méthode des pinceaux de Lefschetz et (I 5.3) à
l'étude du groupe fondamental. Pour $S$ une surface sur un corps algébriquement clos
$k$, d'exposant caractéristique $p$, on montre que le groupe profini $\pi_1^{(p)}(S,s)$, complété
en dehors de $p$ du groupe fondamental, est de pro-$(p)$-présentation finie.

Comme expliqué plus haut, les exposés III à V n'existent pas.


\SGApage{I}{VII}


\SGAvar{L'exposé VI contient, avec quelques compléments, la théorie des déformations
de Schlessinger. On y prend soin de tenir compte des automorphismes infinitésimaux
des objets qu'on classifie,et de ne pas passer trop brutalement au foncteur des classes
d'isomorphie d'objets. On y donne aussi une nouvelle construction des
$\underline{\operatorname{Ext}}^{i}(L_{X/S},-)$ ($L_{X/S}$ complexe cotangent relatif de $X/S$). Cet exposé sert dans le
reste du séminaire surtout via l'étude qui y est faite des déformations des singularités
quadratiques ordinaires.

}{L'exposé VI contient, avec quelques compléments, la théorie des déformations
de Schlessinger. On y prend soin de tenir compte des automorphismes infinitésimaux
des objets qu'on classifie, et de ne pas passer trop brutalement au foncteur des classes
d'isomorphie d'objets. On y donne aussi une nouvelle construction des
$\underline{\operatorname{Ext}}^{i}(L_{X/S},-)$ ($L_{X/S}$ complexe cotangent relatif de $X/S$). Cet exposé sert dans le
reste du séminaire surtout via l'étude qui y est faite des déformations des singularités
quadratiques ordinaires.

}{F7I-INTRO-005}Les exposés VII et VIII sont consacrés à la théorie des biextensions. Cette
théorie joue un rôle essentiel dans l'exposé IX, pour exprimer ce qu'il advient d'une
polarisation quand on passe d'une variété abélienne à un modèle de Néron de celle-ci.
Cet exposé IX contient la démonstration du théorème de réduction stable des variétés
abéliennes, et diverses applications.

La suite de ce séminaire : SGA 7 II, par P. Deligne et N. Katz, paraîtra
ultérieurement.

\begin{flushright}
Bures sur Yvette, mai 1972,\\
P. DELIGNE
\end{flushright}
\selectlanguage{french}

\ifnum\SGAform=1
\clearpage\section*{Apparat critique}\small
Les lectures ci-dessous sont données dans le TeX exact des témoins, afin de préserver les distinctions de notation. Les indices décimaux renvoient aux différences individuelles du passage marqué.

\par\medskip\noindent\hypertarget{app-F7I-INTRO-002.1}{}\textbf{F7I-INTRO-002.1}\quad p. V; inherited corrected-layer difference.\par
\noindent Imprimé : {\ttfamily\footnotesize  ;\allowbreak{}}\par
\noindent Corrigé : {\ttfamily\footnotesize ,\allowbreak{}}\par
\noindent Source : ADOPT-F7I-INTRO-002.1. \nolinkurl{fr/sga7i/diplomatic/general_introduction_body.tex}; \nolinkurl{fr/sga7i/current/general_introduction_body.tex}\par

\par\medskip\noindent\hypertarget{app-F7I-INTRO-003.1}{}\textbf{F7I-INTRO-003.1}\quad p. V; inherited corrected-layer difference.\par
\noindent Imprimé : {\ttfamily\footnotesize soients\allowbreak{}}\par
\noindent Corrigé : {\ttfamily\footnotesize soient\allowbreak{}}\par
\noindent Source : ADOPT-F7I-INTRO-003.1. \nolinkurl{fr/sga7i/diplomatic/general_introduction_body.tex}; \nolinkurl{fr/sga7i/current/general_introduction_body.tex}\par

\par\medskip\noindent\hypertarget{app-F7I-INTRO-003.2}{}\textbf{F7I-INTRO-003.2}\quad p. V; inherited corrected-layer difference.\par
\noindent Imprimé : {\ttfamily\footnotesize 5\allowbreak{}}\par
\noindent Corrigé : {\ttfamily\footnotesize 6\allowbreak{}}\par
\noindent Source : ADOPT-F7I-INTRO-003.2. \nolinkurl{fr/sga7i/diplomatic/general_introduction_body.tex}; \nolinkurl{fr/sga7i/current/general_introduction_body.tex}\par

\par\medskip\noindent\hypertarget{app-F7I-INTRO-004.1}{}\textbf{F7I-INTRO-004.1}\quad p. VI; inherited corrected-layer difference.\par
\noindent Imprimé : {\ttfamily\footnotesize succinte\allowbreak{}}\par
\noindent Corrigé : {\ttfamily\footnotesize succincte\allowbreak{}}\par
\noindent Source : ADOPT-F7I-INTRO-004.1. \nolinkurl{fr/sga7i/diplomatic/general_introduction_body.tex}; \nolinkurl{fr/sga7i/current/general_introduction_body.tex}\par

\par\medskip\noindent\hypertarget{app-F7I-INTRO-005.1}{}\textbf{F7I-INTRO-005.1}\quad p. VII; inherited corrected-layer difference.\par
\noindent Imprimé : \emph{(vide)}\par
\noindent Corrigé : \emph{(1 espace(s), 0 retour(s) de ligne dans le témoin TeX)}\par
\noindent Source : ADOPT-F7I-INTRO-005.1. \nolinkurl{fr/sga7i/diplomatic/general_introduction_body.tex}; \nolinkurl{fr/sga7i/current/general_introduction_body.tex}\par

\normalsize\fi
\clearpage\SGAtoc{Exposé I — Résumé des premiers exposés de A. Grothendieck}
\selectlanguage{french}

\SGApage{I}{1}


SGA 7


\underline{Exposé I}


\underline{\underline{Résumé des premiers exposés de A. Grothendieck}}


\underline{rédigé par P. Deligne}


\underline{Sommaire}

0. Préliminaires

1. Démonstration arithmétique du théorème de monodromie

2. Cycles évanescents

3. Démonstration géométrique du théorème de monodromie

4. Critères de nullité pour les faisceaux de cycles évanescents

5. Action de la monodromie sur les $\pi_1$

6. Appendice par P. Deligne : démonstration arithmétique du
théorème de réduction stable


\SGAsectiontoc{0. Préliminaires}{0}
0. \underline{Préliminaires}

0.0. \underline{Terminologie}. Rappelons les définitions suivantes :

(0.0.1)    \underline{trait} : spectre d'un anneau de valuation discrète. On note souvent  $\eta$  et  $s$  les points génériques et fermés d'un trait.

(0.0.2)    \underline{strictement local} : pour un anneau : local hensélien à corps résiduel séparablement clos. Pour un schéma : spectre d'un tel anneau.

(0.0.3)    \underline{point géométrique de}  $S$  (dans cet exposé) : morphisme  $\bar{x} : \mathrm{Spec}(\bar{k}) \longrightarrow S$  d'image  $x \in S$ , $k(\bar{x}) \underset{\mathrm{dfn}}{=} \bar{k}$  étant une clôture séparable de  $k(x)$ . Pour un trait hensélien  $S = \mathrm{Spec}(V)$ , on note souvent  $\bar{\eta}$  un point géométrique générique (i.e. d'image  $\eta$ ) de  $S$ . Le corps résiduel de l'anneau de valuation  $\bar{V}$  clôture intégrale de  $V$  dans  $k(\bar{\eta})$  est une extension algébrique séparablement close de  $k(s)$  et définit un point géométrique  $\bar{s}$  de  $S$  localisé en  $s$ .


\SGApage{I}{2}

(0.0.4)    \underline{localisé strict de}  $X$  \underline{en le point géométrique}  $\bar{x}$  : SGA 4 VIII 4.4.

(0.0.5)    \underline{essentiellement de type fini sur}  $S$ : pour un $S$-schéma local (resp. local hensélien, resp. strictement local) : spectre d'un anneau local (resp. de l'hensélisé d'un anneau local, resp. localisé strict) d'un schéma de type fini sur  $S$ .

(0.0.6)    $F \longmapsto F(n)$ : twist à la Tate.

(0.0.7)    Un endomorphisme  $T$  d'un espace vectoriel de dimension finie est \underline{quasi-unipotent} si ses valeurs propres sont des racines de l'unité, i.e. s'il existe $N \geqslant 1$ , $M \geqslant 1$  tels que $(T^N - 1)^M = 0$ .


0.1.  \underline{Rappels sur la cohomologie $\ell$-adique}

Soit  $X$  un schéma de type fini sur un corps algébriquement clos  $k$  d'exposant caractéristique  $p$  et  $\ell$  un nombre premier premier à  $p$ .

Les groupes de cohomologie $\ell$-adique de  $X$  ont été définis dans SGA 4 et SGA 5.

a) $H^i(X,\mathbb{Z}/\ell^n)$  est le $i^{\text{ème}}$ groupe de cohomologie du site étale  $X_{et}$  de  $X$  à valeurs dans  $\mathbb{Z}/\ell^n$ ;

b) par définition,  $H^i(X,\mathbb{Z}_\ell) = \underset{n}{\underleftarrow{\ell\mathrm{im}}}\, H^i(X,\mathbb{Z}/\ell^n)$  ;

c) par définition,  $H^i(X,\mathbb{Q}_\ell) = H^i(X,\mathbb{Z}_\ell) \otimes_{\mathbb{Z}_\ell} \mathbb{Q}_\ell$ .

Les groupes de cohomologie $\ell$-adique à supports propres de  $X$  sont définis de même à partir des  $H^i_c(X,\mathbb{Z}/\ell^n)$  de SGA 4 XVII.

Dans chacun des cas suivants, on sait prouver que ces groupes  $H^i(X)$  ont des propriétés de finitude raisonnables :

(0.1.1)    en cohomologie à support propre ;

(0.1.2)    pour  $X$  propre sur  $k$ ;

(0.1.3)    pour  $X$  lisse sur  $k$ (grâce à  a)  et à la dualité de Poincaré SGA 4 XVIII)

\SGAvar{(0.1.4)    lorsqu'on dispose de la résolution des singularités à la Hironoka pour les schémas de type fini sur  $k$ , de dimension $\leqslant \dim(X)$ ;

}{(0.1.4)    lorsqu'on dispose de la résolution des singularités à la Hironaka pour les schémas de type fini sur  $k$ , de dimension $\leqslant \dim(X)$ ;

}{F7I-I-001}(0.1.5)    pour  $i = 0, 1$ .


\SGApage{I}{3}

Dans tous ces cas, les  $H^i(X,\mathbb{Z}/\ell^n)$  sont finis, les $H^i(X,\mathbb{Z}_\ell)$  de type fini, les $H^i(X,\mathbb{Q}_\ell)$  de dimension finie et (sauf peut-être pour (0.1.5)) on dispose de suites exactes courtes scindées

\begin{equation*}
0 \longrightarrow H^i(X,\mathbb{Z}_\ell) \otimes \mathbb{Z}/\ell^n \longrightarrow H^i(X,\mathbb{Z}/\ell^n) \longrightarrow \mathrm{Tor}_1(H^{i+1}(X,\mathbb{Z}_\ell),\mathbb{Z}/\ell^n) \longrightarrow 0 \ .
\end{equation*}
De plus, la démonstration du théorème de finitude fournit chaque fois une démonstration d'un théorème de constructibilité générique : si  $k$  est le corps des fractions d'un schéma noéthérien intègre  $S$  et que  $X$  est la fibre générique de  $f_1 : X_1 \longrightarrow S$  (de type fini), il existe un ouvert non vide  $U \subset S$  tel que, sur  $U$ , les $R^i f_* \mathbb{Z}/\ell^n$  soient localement constants de formation compatible à tout changement de base.

0.2. Ne supposons plus  $k$  algébriquement clos, et soit  $\bar{k}$  une clôture algébrique (ou séparable, cela reviendrait au même) de  $k$ . Nous considérerons les groupes de cohomologie "géométriques"  $H^i(X_{\bar{k}})$  ($X_{\bar{k}}$  est le schéma déduit de  $X$  par extension des scalaires de  $k$  à  $\bar{k}$ ). Par transport de structure, le groupe de Galois  $\mathrm{Gal}(\bar{k}/k)$  agit sur ces groupes. En  $\mathbb{Q}_\ell$-cohomologie, lorsqu'on est dans l'un des cas (0.1.1) à (0.1.5), on obtient ainsi des représentations continues de  $\mathrm{Gal}(\bar{k}/k)$  dans un groupe linéaire  $\mathrm{GL}(n,\mathbb{Q}_\ell)$ .


0.3. \underline{Rappels sur les traits} (0.0.1)

Pour les démonstrations, on renvoie à $[4]$. Soient  $S$  un trait hensélien, $\eta$ , $s$ , $V$ , $\bar{V}$ , $\bar{\eta}$ , $\bar{s}$  comme en (0.0.1) (0.0.3). On note  $p$  l'exposant caractéristique de  $k(s)$  et  $\mathrm{Gal}(\bar{\eta}/\eta)$  (resp.  $\mathrm{Gal}(\bar{s}/s)$ ) le groupe de Galois  $\mathrm{Gal}(k(\bar{\eta})/k(\eta))$  (resp. ...). Par transport de structure,  $\mathrm{Gal}(\bar{\eta}/\eta)$  agit sur  $k(\bar{s})$ , d'où une application de $\mathrm{Gal}(\bar{\eta}/\eta)$  dans (en fait sur)  $\mathrm{Gal}(\bar{s}/s)$ , de noyau le \underline{groupe d'inertie}  $I$ . Le sous-corps de  $k(\bar{\eta})$  défini par  $I$  est une extension non ramifiée maximale de  $k(\eta)$ .

Le groupe profini  $I$  admet un plus grand sous-groupe  $P$  qui soit un pro-$p$-groupe. Le quotient  $I/P$  est le \underline{groupe d'inertie modéré}. On a canoniquement

\begin{equation*}
I/P \simeq \underset{n}{\underleftarrow{\lim}}\ \mu_n(k(\bar{s})) \underset{\mathrm{dfn}}{=} \widehat{\mathbb{Z}}(1)(k(\bar{s})) \ .
\end{equation*}

\SGApage{I}{4}

Si  $\phi_n$  est l'application canonique de  $I/P$  dans  $\mu_n(\bar{k})$ , et que  $u$  est un élément de  $\bar{V}$  de valuation  $1/n$ , on a pour  $\sigma \in I$

\SGAvar{\begin{equation*}
\sigma(u) \;\equiv\; \phi_n(u).u \qquad (\text{mod. éléments de valuation} > 1/n).
\end{equation*}
En résumé, on a :

}{\begin{equation*}
\sigma(u) \;\equiv\; \phi_n(\sigma).u \qquad (\text{mod. éléments de valuation} > 1/n).
\end{equation*}
En résumé, on a :

}{F7I-I-002}\begin{equation}\tag{0.3.1}
\begin{aligned}
\operatorname{Gal}(\bar\eta/\eta) &\supset I \supset P \ , \\
\operatorname{Gal}(\bar\eta/\eta)/I &\simeq \operatorname{Gal}(\bar s/s) \ , \\
I/P &\simeq \hat{\mathbf{Z}}(1)(k(\bar s)) \ , \\
P &\quad : \text{pro-p-groupe} \ (\ \{e\} \ \text{si} \ p = 1 \ ) \ ,
\end{aligned}
\end{equation}
et l'action (par automorphismes intérieurs dans  $\operatorname{Gal}(\bar\eta/\eta)$) de  $\operatorname{Gal}(\bar\eta/\eta)/I$  sur  $I/P$  s'identifie à l'action naturelle de  $\operatorname{Gal}(\bar s/s)$  sur  $\hat{\mathbf{Z}}(1)(k(\bar s))$  .

\SGAvar{0.4.  Soit plus généralement  $S$  strictement local régulier, $D$  un diviseur régulier dans  $S$  et  $p'$  l'exposant caractéristique du point générique de  $D$ . Soient  $\eta$  un point géométrique générique de  $S$  et  $\bar s$  le point géométrique localisé au point fermé  $s$ , qui s'en déduit. D'après le lemme d'Abhyankar, le groupe fondamental  $\pi_1(S\!-\!D,\bar\eta)$  admet un dévissage analogue à (0.3.1) :

}{0.4.  Soit plus généralement  $S$  strictement local régulier, $D$  un diviseur régulier dans  $S$  et  $p'$  l'exposant caractéristique du point générique de  $D$ . Soient  $\bar\eta$  un point géométrique générique de  $S$  et  $\bar s$  le point géométrique localisé au point fermé  $s$ , qui s'en déduit. D'après le lemme d'Abhyankar, le groupe fondamental  $\pi_1(S\!-\!D,\bar\eta)$  admet un dévissage analogue à (0.3.1) :

}{F7I-I-003}\begin{equation}\tag{0.4.1}
\begin{aligned}
\pi_1(S\!-\!D,\bar\eta) &\supset I \supset P \\
\pi_1(S\!-\!D,\bar\eta)/I &\simeq \operatorname{Gal}(\bar s/s) \\
I/P &\simeq \hat{\mathbf{Z}}(1)(k(\bar s)) \\
P &\quad : \text{sans quotient d'ordre premier à } p' \ .
\end{aligned}
\end{equation}
L'action par automorphismes intérieurs de  $\pi_1(S\!-\!D,\bar\eta)/I$  sur  $I/P$  s'identifie à l'action naturelle de  $\operatorname{Gal}(\bar s/s)$  sur  $\hat{\mathbf{Z}}(1)(k(\bar s))$  .


0.5.  \underline{La désingularisation de Néron}

La méthode de désingularisation de Néron fournit le résultat suivant :

\SGAvar{\underline{Lemme} 0.5.1.  \underline{Soit}  $S \longrightarrow S_o$  \underline{un morphisme de traits henséliens,}  \underline{avec}  $S = Spec(V)$ ,  $S_o = Spec(V_o)$  \underline{et}  $\pi$  \underline{une uniformisante de}  $V_o$ .  \underline{On suppose que}  $\pi$  \underline{est une uniformisante de}  $V$ ,  \underline{et que les extensions}  $k(s)/k(s_o)$  \underline{et}  $k(\eta)/k(\eta_o)$  \underline{sont séparables.}  \underline{Alors,}  $V$  \underline{est limite inductive de}  $V_o$\underline{-algébres henséliennes lisses et essentiellement de type fini}  $B_i$  \underline{sur}  $V_o$ .


}{\underline{Lemme} 0.5.1.  \underline{Soit}  $S \longrightarrow S_o$  \underline{un morphisme de traits henséliens,}  \underline{avec}  $S = Spec(V)$ ,  $S_o = Spec(V_o)$  \underline{et}  $\pi$  \underline{une uniformisante de}  $V_o$ .  \underline{On suppose que}  $\pi$  \underline{est une uniformisante de}  $V$ ,  \underline{et que les extensions}  $k(s)/k(s_o)$  \underline{et}  $k(\eta)/k(\eta_o)$  \underline{sont séparables.}  \underline{Alors,}  $V$  \underline{est limite inductive de}  $V_o$\underline{-algèbres henséliennes lisses et essentiellement de type fini}  $B_i$  \underline{sur}  $V_o$ .


}{F7I-I-004}\SGApage{I}{5}

Pour construire les  $V_o$-algèbres voulues, on prend  $V_o \subset B \subset V$  avec  $B$  de type fini, et on désingularise  $Spec(B)$  à la Néron  ([1]  § 4).

Soit  $S = Spec(V)$  un trait hensélien.

\SGAvar{(0.5.2)  Si  $S$  est d'égale caractéristique, d'uniformisante  $\pi$ , on peut appliquer (0.5.1) pour  $S_o = \mathbf{F}_p[\pi]_{(\pi)}$ . L'extension  $s/s_o$  est séparable car  $\mathbf{F}_p$  est parfait, l'extension  $\eta/\eta_o$  l'est car  $\pi$ , étant une uniformisante, n'est pas une puissance  $p^{\text{ième}}$  (si  $p \neq 1$ ) donc fait partie d'une $p$-base. On trouve que  $S$  est limite de spectres de  $\mathbf{F}_p[\pi]_{(\pi)}$-algèbre henséliennes lisses essentiellement de type fini.

}{(0.5.2)  Si  $S$  est d'égale caractéristique, d'uniformisante  $\pi$ , on peut appliquer (0.5.1) pour  $S_o = \mathbf{F}_p[\pi]_{(\pi)}$ . L'extension  $s/s_o$  est séparable car  $\mathbf{F}_p$  est parfait, l'extension  $\eta/\eta_o$  l'est car  $\pi$ , étant une uniformisante, n'est pas une puissance  $p^{\text{ième}}$  (si  $p \neq 1$ ) donc fait partie d'une $p$-base. On trouve que  $S$  est limite de spectres de  $\mathbf{F}_p[\pi]_{(\pi)}$-algèbres henséliennes lisses essentiellement de type fini.

}{F7I-I-005}(0.5.3)  Si  $S$  est complet, d'inégale caractéristique, à corps résiduel parfait  $k$ , $V$  est une extension d'Eisenstein de l'anneau des vecteurs de Witt  ( $=$ anneau de Cohen)  $W(k)$

\SGAvar{\begin{equation*}
V \simeq W(k) \ [\pi] \ /(\pi^n + \sum_0^{\,n-1} a_i \ \pi^i) \qquad .
\end{equation*}
Si on perturbe un peu les  $a_i$ , on trouve une extension isomorphe. On peut donc supposer les  $a_i$ . dans  $W(k')$ , avec  $k'$  de type fini sur  $\mathbf{F}_p$ . Soit  $k_o$  la clôture parfaite de  $k'$  dans  $k$ et

}{\begin{equation*}
V \simeq W(k) \ [\pi] \ /(\pi^n + \sum_0^{\,n-1} a_i \ \pi^i) \qquad .
\end{equation*}
Si on perturbe un peu les  $a_i$ , on trouve une extension isomorphe. On peut donc supposer les  $a_i$  dans  $W(k')$ , avec  $k'$  de type fini sur  $\mathbf{F}_p$ . Soit  $k_o$  la clôture parfaite de  $k'$  dans  $k$ et

}{F7I-I-006}\begin{equation*}
V_o = W(k_o) \ [\pi]/(\pi^n + \sum_0^{\,n-1} a_i \ \pi^i) \qquad .
\end{equation*}
Le lemme 0.5.1 s'applique à  $V/V_o$ , et le corps résiduel de  $V_o$  est une extension radicielle d'un corps de type fini sur  $\mathbf{F}_p$  .


\SGAsectiontoc{1.  Démonstration arithmétique du théorème de monodromie}{1}
1.  \underline{Démonstration arithmétique du théorème de monodromie}

Le lecteur trouvera dans  [5], Appendice, une démonstration du résultat suivant de A. Grothendieck.

\underline{Proposition} 1.1.  \underline{Soient}  $S$  \underline{un trait hensélien,}  $\ell$  \underline{un nombre premier premier à l'exposant caractéristique}  $p$  \underline{de}  $k(s)$  \underline{et}  $\rho$  \underline{une représentation continue de}  $\operatorname{Gal}(\bar\eta/\eta)$  \underline{dans}  $GL(n,\mathbf{Q}_\ell)$ . \underline{On suppose remplie la condition suivante} :

$(*_\ell)$      \underline{Aucune extension finie de}  $k(s)$  \underline{ne contient toutes les racines de l'unité d'ordre une puissance de}  $\ell$ .


\SGApage{I}{6}

\underline{Alors, il existe un sous-groupe ouvert}  $I_1$  \underline{du groupe d'inertie}  $I$  \underline{tel que}  $\rho(\sigma)$  \underline{soit unipotent pour}  $\sigma \in I_1$ .

La condition  $(*_\ell)$  est automatiquement vérifiée pour  $k(s)$  de type fini sur le corps premier (ou radiciel sur un tel corps).

Le théorème suivant résulte aussitôt de 1.1.

\SGAvar{\underline{Théorème de monodromie} 1.2.  \underline{Avec les notations précédentes,} \underline{soit}  $X$  \underline{un schéma de type fini sur}  $\eta$  \underline{et}  $H$  \underline{un espace de cohomologie de}  $X_{\bar\eta}$ , à \underline{coefficients dans}  $\mathbf{Q}_\ell$ , \underline{de  l'un des types} (1.1.1) \underline{à} (1.1.5). \underline{Si la condition}  $(*_\ell)$  \underline{est vérifiée,} \underline{l'action d'un quelconque élément de}  $I$  \underline{sur}  $H$  \underline{est quasi-unipotente.}

Les résultats de passage à la limite 0.5 permettent de se débarasser de l'hypothèse  $(*_\ell)$.

}{\underline{Théorème de monodromie} 1.2.  \underline{Avec les notations précédentes,} \underline{soit}  $X$  \underline{un schéma de type fini sur}  $\eta$  \underline{et}  $H$  \underline{un espace de cohomologie de}  $X_{\bar\eta}$ , à \underline{coefficients dans}  $\mathbf{Q}_\ell$ , \underline{de  l'un des types} (0.1.1) \underline{à} (0.1.5). \underline{Si la condition}  $(*_\ell)$  \underline{est vérifiée,} \underline{l'action d'un quelconque élément de}  $I$  \underline{sur}  $H$  \underline{est quasi-unipotente.}

Les résultats de passage à la limite 0.5 permettent de se débarrasser de l'hypothèse  $(*_\ell)$.

}{F7I-I-007}\underline{Variante} 1.3.  \underline{La condition}  $(*_\ell)$  \underline{dans} 1.2 \underline{est superflue.}

Soit  $H'_o$  un espace de  $\mathbf{Z}_\ell$-cohomologie de  $X_{\bar\eta}$  qui donne naissance à  $H$ , et  $H_o = H'_o$/torsion. On a par hypothèse  $H = H_o \otimes_{\mathbf{Z}_\ell} \mathbf{Q}_\ell$ , et l'action de  $\operatorname{Gal}(\bar\eta/\eta)$  se déduit de

\begin{equation*}
\rho \ : \ \operatorname{Gal}(\bar\eta/\eta) \longrightarrow GL(H_o) \qquad .
\end{equation*}
Soit  $I_{\ell'}$  le plus grand sous-groupe premier à  $\ell$ de  $I$ . Son image dans  $GL(H_o)$  est un groupe de Lie $\ell$-adique premier à  $\ell$ , donc est finie.

Procédons à une extension de trait  $S' \longrightarrow S$ . Ceci ne change pas  $H_o$ (SGA 4 XVI 6.6)  et  $\rho(I')$  est d'indice fini dans  $\rho(I)$ , car  $I'/P'$  est d'indice fini dans  $I/P$  et ce qui précède. Il suffit donc de prouver le théorème sur  $S'$ : on peut supposer que

a) $\operatorname{Gal}(\bar\eta/\eta)$  agit trivialement sur  $H_o/\ell H_o$  ;

b) si  $S$  est d'inégale caractéristique,  $S$  est complet à corps résiduel parfait (EGA $\mathrm{O}_{III}$  10.3.1).

Posons  $S = Spec(V)$ , définissons  $V_o$  comme en (0.5.2) et (0.5.3) et appliquons (0.5.1). On a

\begin{equation*}
V = \varinjlim B_i
\end{equation*}

\SGApage{I}{7}

avec  $S_i = Spec(B_i)$  essentiellement lisse de type fini sur  $V_o$ . Soit  $D_i \subset S_i$  d'équation  $\pi = 0$ . Pour  $i$  assez grand, $H_o$  provient d'un  $\mathbf{Z}_\ell$-faisceau constant tordu  $\mathcal{H}_i$  sur  $S_i - D_i$  (d'après la "constructibilité générique" (0.1)) et  $\mathcal{H}_i/\ell\mathcal{H}_i$  est constant.

Les dévissages (0.3.1) et (0.4.1) donnent lieu à des diagrammes commutatifs :

\begin{equation}\tag{1.3.1}
\begin{tikzcd}[column sep=small, row sep=large]
\mathrm{Gal}(\bar{\eta}/\eta) \arrow[d] \& \supset \& I \arrow[d] \& \supset \& P \arrow[d] \\
\pi_1(S_i-D_i,\bar{\eta}) \& \supset \& I_i \& \supset \& P_i
\end{tikzcd}
\qquad\text{et}\qquad
\begin{tikzcd}[row sep=large]
I/P \arrow[r, dash, "\sim"] \arrow[d] \& \hat{\mathbf{Z}}(1) \arrow[d, equal] \\
I_i/P_i \arrow[r, dash, "\sim"] \& \hat{\mathbf{Z}}(1)
\end{tikzcd}
\end{equation}



















\SGAvar{et la représentation  $H_o$  de  $Gal(\bar{\eta}/\eta)$  se déduit d'une action de  $\pi_1(S_i-D_i,\bar{\eta})$  sur $H_o$ , triviale sur  $H_o/\ell H_o$ . Le groupe  $Ker(GL(H)$ $\longrightarrow$ $GL(H/$$\ell$$H))$  est un pro-$\ell$-groupe ; $P_i$  agit donc trivialement, et on applique la variante suivante de 1.1.

\underline{Variante} 1.4.  \underline{Reprenons les notations de} 0.4 \underline{et soit}  $\rho$  \underline{une représentation continue de} $Gal(\bar{\eta}/\eta)$  \underline{dans}  $Gal(n,\mathbf{Q}_\ell)$ . \underline{On suppose que}  $k(s)$  \underline{vérifie}  $(\ast_\ell)$  \underline{et que}  $\rho(P)$ \underline{est fini}. \underline{Alors,} \underline{l'action de}  $I$  \underline{est quasi-unipotente}.


}{et la représentation  $H_o$  de  $Gal(\bar{\eta}/\eta)$  se déduit d'une action de  $\pi_1(S_i-D_i,\bar{\eta})$  sur $H_o$ , triviale sur  $H_o/\ell H_o$ . Le groupe  $Ker(GL(H_o)$ $\longrightarrow$ $GL(H_o/$$\ell$$H_o))$  est un pro-$\ell$-groupe ; $P_i$  agit donc trivialement, et on applique la variante suivante de 1.1.

\underline{Variante} 1.4.  \underline{Reprenons les notations de} 0.4 \underline{et soit}  $\rho$  \underline{une représentation continue de} $Gal(\bar{\eta}/\eta)$  \underline{dans}  $GL(n,\mathbf{Q}_\ell)$ . \underline{On suppose que}  $k(s)$  \underline{vérifie}  $(\ast_\ell)$  \underline{et que}  $\rho(P)$ \underline{est fini}. \underline{Alors,} \underline{l'action de}  $I$  \underline{est quasi-unipotente}.


}{F7I-I-008}\SGAsectiontoc{2. Cycles évanescents}{2}
2. \underline{Cycles évanescents}

2.1.    Soit  $S$  un trait strictement local, et reprenons les notations de 0.3. Par hypothèse, on a ici  $s = \bar{s}$ . Soit  $f : X \longrightarrow S$  un schéma de type fini sur  $S$ . Le formalisme des cycles évanescents est un outil pour étudier la relation entre les cohomologies de  $X_s$  et de  $X_{\bar{\eta}}$ , et l'action du groupe d'inertie sur celle de  $X_{\bar{\eta}}$ , à partir d'informations \underline{locales} (sur  $X$ ) sur le morphisme  $f$ .

Un résultat clef de ce type est le théorème de spécialisation SGA 4 XVI 2.2, selon lequel pour  $f$  propre et lisse  $X_s$  et  $X_{\bar{\eta}}$  ont même cohomologie.

Dans la théorie transcendante, si  $f : X \longrightarrow D$  est un morphisme propre d'un espace analytique  $X$  dans le disque  $D$ , les cycles évanescents apparaissent comme suit :

a) Quitte à rapetisser  $D$ , on peut supposer que  $X$  est un fibré topologique au-dessus du disque épointé  $D^{*}$  (théorème d'isotopie de Thom). Les fibres


\SGApage{I}{8}

$X_t$ ( $t \in D^{*}$ )  sont en particulier toutes isomorphes.

b) Il n'est pas déraisonnable de se représenter la fibre spéciale  $X_s$  comme déduite de "la" fibre générale  $X_t$  par contraction de polyèdres dans  $X_t$ : on aurait $\pi : X_t \longrightarrow X_s$ . La méthode des cycles évanescents consiste alors à étudier la différence entre les cohomologies de  $X_s$  et  $X_t$  à l'aide de la suite spectrale de Leray de  $\pi$ .

\SGAvar{La théorie géométrique expliquée ci-dessous est très proche de cette théorie transcendante ; il faut remplacer le point général  $t \in D$  par  $\bar{\eta}$ , le point géométrique générique de  $S$ . Inversement, on peut calquer la théorie transcendance sur la théorie géométrique, en remplaçant  $t$  ou  $\bar{\eta}$  par le revêtement universel  $\widetilde{D}^{*}$  de  $D^{*}$ , et $X_t$  ou  $X_{\bar{\eta}}$  par  $X \times_D \widetilde{D}^{*}$ . Ce faisant, on se débarasse de nombreux  $\varepsilon$  et  $\eta$ , mais on perd parfois de l'information.

}{La théorie géométrique expliquée ci-dessous est très proche de cette théorie transcendante ; il faut remplacer le point général  $t \in D$  par  $\bar{\eta}$ , le point géométrique générique de  $S$ . Inversement, on peut calquer la théorie transcendante sur la théorie géométrique, en remplaçant  $t$  ou  $\bar{\eta}$  par le revêtement universel  $\widetilde{D}^{*}$  de  $D^{*}$ , et $X_t$  ou  $X_{\bar{\eta}}$  par  $X \times_D \widetilde{D}^{*}$ . Ce faisant, on se débarrasse de nombreux  $\varepsilon$  et  $\eta$ , mais on perd parfois de l'information.

}{F7I-I-009}2.2. Fixons un anneau de torsion  $\Lambda$  dans lequel  $p$  soit inversible (par exemple $\Lambda = \mathbf{Z}/\ell^n\mathbf{Z}$ ,  $\ell$  premier à  $p$ ). Soient les morphismes :

\begin{center}
\begin{tikzcd}[column sep=large, row sep=huge]
X_s \arrow[r, "i"] \arrow[d, "f_s"] \& X \arrow[d, "f"] \& \& X_{\bar{\eta}} \arrow[ll, "\bar{j}"] \arrow[d, "f_{\bar{\eta}}"] \\
s \arrow[r] \& S \& \eta \arrow[l] \& \bar{\eta} \arrow[l]
\end{tikzcd}
\end{center}














\SGAvar{Les faisceaux de cycles évanescents  $\psi_i$  sur  $X_s$  sont les

}{Les faisceaux de cycles évanescents  $\psi^i$  sur  $X_s$  sont les

}{F7I-I-010}\begin{equation*}\tag{2.2.1}
\psi^i \;=\; i^{*}\, R^i\, \bar{j}_{*}\, \Lambda \;.
\end{equation*}
Si  $f$  est \underline{propre}, on a (SGA 4 XII 5.5)

\SGAvar{\begin{equation*}\tag{2.2.2}
i^{*} \;:\; H^p(X,R^q\,\bar{j}_{*}\,\Lambda) \xrightarrow{\ \sim\ } H^p(X_s,i^{*}R^q\,j_{*}\,\Lambda) \qquad ,
\end{equation*}
et la suite spectrale de Leray pour  $\bar{j}$  se récrit

}{\begin{equation*}\tag{2.2.2}
i^{*} \;:\; H^p(X,R^q\,\bar{j}_{*}\,\Lambda) \xrightarrow{\ \sim\ } H^p(X_s,i^{*}R^q\,\bar{j}_{*}\,\Lambda) \qquad ,
\end{equation*}
et la suite spectrale de Leray pour  $\bar{j}$  se récrit

}{F7I-I-011}\begin{equation*}\tag{2.2.3}
H^p(X_s,\psi^q) \Longrightarrow H^{p+q}(X_{\bar{\eta}},\Lambda) \qquad .
\end{equation*}
Lorsque rien ne se passe à l'infini, on peut parfois obtenir cette suite spectrale sans hypothèse de propreté.  Qu'il faille quelque hypothèse à l'infini se voit déjà sur le cas trivial  $X_s = \emptyset$ .

\SGApage{I}{9}

On tire aussitôt des définitions :

\underline{Proposition} 2.3.  \underline{Soient}  $\bar{x}$  \underline{un point géométrique de}  $X_s$  \underline{(par exemple un point fermé)} $X_{(\bar{x})}$  \underline{l'henselisé strict de}  $X$  \underline{en}  $\bar{x}$  \underline{et}  $X_{(\bar{x})\bar{\eta}}$   \underline{la fibre géométrique générique de} \underline{la projection de}  $X_{(\bar{x})}$  \underline{sur}  $S$ .  \underline{On a}

\begin{equation*}
\psi^i \;=\; H^i(X_{(\bar{x})\bar{\eta}},\Lambda)
\end{equation*}
\underline{Corollaire} 2.4.  \underline{Là où}  $f$  \underline{est lisse,}  $\psi^i = 0$  \underline{pour}  $i > 0$  \underline{et}  $\psi^{0} = \Lambda$ .

C'est SGA 4 XV 2.1.

Le groupe  Gal($\bar{\eta}/\eta$)  (égal par hypothèse au groupe d'inertie) agit par transport de structure sur les   $\psi^i$ , et

\underline{Scholie} 2.5.  \underline{La suite spectrale} 2.2.3 \underline{est}  Gal($\bar{\eta}/\eta$)-\underline{équivariante}.

\underline{Variante} 2.6.  Soient  $\bar{S}$ $=$ Spec($\bar{V}$)  et  $\bar{X} = X \times_S \bar{S}$ .  $X_{\bar{\eta}}$  est le complément de $X_s = X_{\bar{s}}$  dans  $\bar{X}$ . Posons

\begin{equation*}
\phi^i \;=\; \underline{H}^{\,i+1}_{X_s}\,(\bar{X}\;,\Lambda\;)
\end{equation*}
(faisceaux de cohomologie à support). On a une suite exacte de faisceaux

\begin{equation*}
0 \longrightarrow \phi^{-1} \longrightarrow \Lambda \longrightarrow \psi^{0} \longrightarrow \phi^{0} \longrightarrow 0
\end{equation*}
et, pour  $i > 0$ , des isomorphismes  $\psi^i \xrightarrow{\ \sim\ } \phi^i$ . Pour  $f$  lisse, les  $\phi^i$  sont nuls.

\SGAvar{Posons  $\phi^i_{gl} = H^i_{X_s}(\bar{X}\;,\Lambda)$  (analogue global des  $\phi^i$ ). On dispose d'une suite spectrale

}{Posons  $\phi^i_{gl} = H^{i+1}_{X_s}(\bar{X}\;,\Lambda)$  (analogue global des  $\phi^i$ ). On dispose d'une suite spectrale

}{F7I-I-012}\begin{equation*}\tag{2.6.1}
H^p(X_s,\phi^q) \Longrightarrow \phi^{p+q}_{gl}
\end{equation*}
et d'une suite exacte longue de cohomologie qui, pour  $f$  propre, se récrit

\begin{equation*}\tag{2.6.2}
\dots \longrightarrow H^i(X_s\;,\Lambda\;) \longrightarrow H^i(X_{\bar{\eta}}\;,\Lambda\;) \longrightarrow \phi^i_{gl} \longrightarrow \dots
\end{equation*}
\underline{Variante} 2.7.  Soit  $K_t$  l'extension modérément ramifiée maximale de  $K$  dans  $\bar{K}$ . On pose  $X_t = X \times_S$ Spec($K_t$) . Soit  $\bar{j}_t$  la projection de  $X_t$  dans  $X$ . On définit des variantes modérées  ("tame")  des  $\psi^i$  en posant


\SGApage{I}{10}

\begin{equation*}
\psi^{i}_{t} \;=\; R^{i}\,\bar{j}_{t*}\,\Lambda
\end{equation*}
Utilisant que  $P$  est un pro-$p$-groupe, avec  $p$  inversible dans  $\Lambda$ , on trouve que

\begin{equation}\tag{2.7.1}
H^{*}(X_{t},\Lambda) \;=\; H^{*}(X_{\bar{\eta}},\Lambda)^{P} \qquad ,
\end{equation}
\begin{equation}\tag{2.7.2}
\psi^{i}_{t} \;=\; \psi^{iP} \qquad .
\end{equation}
Pour  $f$  propre, la suite spectrale de Leray de  $\bar{j}_{t}$  se déduit de 2.2.3 par passage aux  $P$-invariants :

\begin{equation}\tag{2.7.3}
H^{p}(X_{s},\psi^{q}_{t}) \;\Longrightarrow\; H^{p+q}(X_{t},\Lambda) \;=\; H^{p+q}(X_{\bar{\eta}},\Lambda)^{P}\ .
\end{equation}

\SGAsectiontoc{3. Démonstration géométrique du théorème de monodromie}{3}
3. \underline{Démonstration géométrique du théorème de monodromie}

Soit  $\Lambda$  comme en 2.2.

\SGAvar{\underline{Conjecture} 3.1 (\underline{pureté}). \underline{Soient} $A$ \underline{un anneau excellent strictement local} (0.2) \underline{régulier et} $D$ \underline{un diviseur régulier de} Spec($A$) \underline{défini par un paramètre} $x$ . \underline{Soit} $U$ $=$ Spec($A$) $-$ $D$ $=$ Spec($A\left[1/X\right]$ ). \underline{On a}

}{\underline{Conjecture} 3.1 (\underline{pureté}). \underline{Soient} $A$ \underline{un anneau excellent strictement local} (0.0.2) \underline{régulier et} $D$ \underline{un diviseur régulier de} Spec($A$) \underline{défini par un paramètre} $x$ . \underline{Soit} $U$ $=$ Spec($A$) $-$ $D$ $=$ Spec($A\left[1/x\right]$ ). \underline{On a}

}{F7I-I-013}\begin{equation}\tag{3.1.1}
H^{q}(U,\Lambda) \;=\; \left\{\begin{array}{ll} \Lambda & \text{pour } q = 0 \\[4pt] \Lambda(-1) & \text{pour } q = 1 \text{ (notation } (0.\ 0.6)) \\[4pt] 0 & \text{pour } q > 1 \end{array}\right.
\end{equation}
Voici des cas où cette conjecture est connue (SGA 4 XIX) :

pour  $q = 0$ ou $1$ , pour  $A$  de caractéristique  $0$ , pour  $A$  d'égale caractéristique  $p$  lorsqu'on dispose de la résolution des singularités en dimension $\leqslant$ dim($A$)  et en caractéristique  $p$ , ou pour  (Spec($A$),$D$)  un couple lisse (théorème de pureté relatif  SGA 4 XVI 3.7).

\SGAvar{Lorsqu'on dispose de la pureté, et que  $D = \sum\limits_{i \in B} D_{i}$  est un diviseur à croisements normaux dans  Spec $A$ , défini par une partie d'un système régulier de paramètre, la cohomologie de  $U$ $=$ Spec($A$) $-$ $D$  est donnée par


}{Lorsqu'on dispose de la pureté, et que  $D = \sum\limits_{i \in B} D_{i}$  est un diviseur à croisements normaux dans  Spec $A$ , défini par une partie d'un système régulier de paramètres, la cohomologie de  $U$ $=$ Spec($A$) $-$ $D$  est donnée par


}{F7I-I-014}\SGApage{I}{11}

\begin{equation}\tag{3.1.2}
\begin{aligned} H^{0}(U,\Lambda) &= \Lambda \\[4pt] H^{1}(U,\Lambda) &= \Lambda(-1)^{B} \\[4pt] H^{q}(U,\Lambda) &= \overset{q}{\bigwedge} H^{1}(U,\Lambda) \qquad \text{(donné par le cup-produit)} \end{aligned}
\end{equation}
Ces formules sont identiques à celles qui, sur  $\mathbb{C}$ , donnent la cohomologie de $D^{*B} \times D^{k}$ .

3.2. Soit  $S$  un trait strictement local. On garde les notations (0.0.1) (0.0.3). Soit  $f : X \longrightarrow S$  un morphisme de type fini. On suppose que  $X$  est régulier, que  $X_{\eta}$  est lisse et que  $X_{s}$  est un diviseur à croisements normaux dans  $X$ . On supposera que ce diviseur est (globalement, et non seulement localement pour la topologie étale) somme de diviseurs réguliers : $(X_{s})_{\mathrm{red}} = \sum\limits_{i \in B} D_{i}$  (le cas général peut se récupérer par localisation). Notons  $m_{i}$  leurs multiplicités : $X_{s} = \sum\limits_{i \in B} m_{i}\, D_{i}$ .

Soient  $x$  un point de  $X_{s}$ , $\bar{x}$  un point géométrique de  $X$  localisé en  $x$  (par exemple  $x$  fermé,  $x = \bar{x}$ ) ,  $C = \{i \mid x \in D_{i}\}$  et  $K$  le complexe concentré en degrés  $0$  et  $-1$

\begin{equation*}
K \;:\; \mathbf{Z}^{C} \;\xrightarrow{\;\;d\;\;}\; \mathbf{Z} \;:\; d((n_{j})) \;=\; \sum\limits_{j} m_{j}\; n_{j} \qquad .
\end{equation*}
\underline{Théorème} 3.3. (\underline{utilise la pureté}). \underline{Les groupes de cycles évanescents modérés} $\psi^{i}_{t\bar{x}}$ \underline{sont les suivants} :

(i) \underline{On a un isomorphisme canonique de} $\Lambda$-\underline{algèbres graduées}

\begin{equation*}
\psi^{*}_{t\ \bar{x}} \;\simeq\; \overset{*}{\bigwedge}(H_{1}(K) \otimes \Lambda(-1)) \otimes_{\Lambda} \; \psi^{0}_{t\bar{x}} \qquad .
\end{equation*}
(ii) $\psi^{0}_{t\ \bar{x}}$ \underline{est une} $\Lambda$-\underline{algèbre} $\Lambda^{E}$ , \underline{pour} $E$ \underline{un ensemble sur lequel le groupe d'inertie modéré} $I_{t} = I/P = \hat{\mathbf{Z}}(1)(k(s))$ \underline{agit transitivement}, \underline{de nombre d'éléments le plus grand diviseur premier à} $p$ \underline{de} $\#$ $H_{0}(K)$ .

Voici la méthode de démonstration.

\SGAvar{a) Soient  $X_{(\bar{x})}$  le localisé strict de  $X$  en  $\bar{x}$ , $U = X_{(\bar{x})} - X_{(\bar{x})s}$  (le complément d'un diviseur à croisements normaux), $t_{j}$  une équation locale de  $D_{j}$ , $X_{n} = X_{(x)}\left[(t_{j}^{1/n})\right]$ (pour  $(n,p) = 1$ ; c'est un schéma régulier) et  $U_{n}$  l'image réciproque de  $U$  dans  $X_{n}$ . D'après 3.1.2, on a  $H^{q}(U_{n},\Lambda) = \overset{q}{\bigwedge}(\Lambda(-1)^{C})$  ; pour  $m = x d$ , l'application "image réciproque" de  $H^{q}(U_{n},\Lambda)$  dans  $H^{q}(U_{m},\Lambda)$  est la multiplication


}{a) Soient  $X_{(\bar{x})}$  le localisé strict de  $X$  en  $\bar{x}$ , $U = X_{(\bar{x})} - X_{(\bar{x})s}$  (le complément d'un diviseur à croisements normaux), $t_{j}$  une équation locale de  $D_{j}$ , $X_{n} = X_{(\bar{x})}\left[(t_{j}^{1/n})\right]$ (pour  $(n,p) = 1$ ; c'est un schéma régulier) et  $U_{n}$  l'image réciproque de  $U$  dans  $X_{n}$ . D'après 3.1.2, on a  $H^{q}(U_{n},\Lambda) = \overset{q}{\bigwedge}(\Lambda(-1)^{C})$  ; pour  $m = n d$ , l'application "image réciproque" de  $H^{q}(U_{n},\Lambda)$  dans  $H^{q}(U_{m},\Lambda)$  est la multiplication


}{F7I-I-015}\SGApage{I}{12}

par  $d^{q}$ . Posant  $\widetilde{U} = \varprojlim\limits_{(n,p)=1} U_{n}$ , on a donc

\begin{equation}\tag{3.3.1}
\left\{\begin{array}{l} H^{0}(\widetilde{U},\Lambda) \;=\; \Lambda \\[8pt] H^{q}(\widetilde{U},\Lambda) \;=\; \varinjlim\limits_{n} H^{q}(U_{n},\Lambda) \;=\; 0 \ \ \text{pour}\ \ q \neq 0 \ . \end{array}\right.
\end{equation}
b) $\widetilde{U}$  est un prorevêtement de  $U$  de groupe  $\hat{\mathbf{Z}}(1)(k(s))^{C}$ . Par définition, on a  $\psi^{*}_{t\ \bar{x}} = H^{*}(U_{\eta_{t}} ,\Lambda)$ ; $\widetilde{U}$  est un prorevêtement de groupe  $H_{1}(K) \otimes \hat{\mathbf{Z}}(1)(k(s))$  de chacune des composantes connexes de  $U_{\eta_{t}}$ , et celles-ci sont permutées transitivement par le groupe d'inertie modéré. De (3.3.1) et de la suite spectrale d'Hochschild-Serre, on tire alors les formules 3.3.

En termes imagés, on peut dire que  $U_{\bar{\eta}}$  est la fibre d'un morphisme  $\delta$  d'un $K(\hat{\mathbf{Z}}(1)(k(s))^{C},1)$  cohomologique  $U$  dans le  $K(\hat{\mathbf{Z}}(1)(k(s)),1)$  cohomologique  $\eta$ , ce morphisme  $\delta$  induisant  $d = \Sigma\, m_{j}\, n_{j} \;:\; \hat{\mathbf{Z}}(1)(k(s))^{C} \longrightarrow \hat{\mathbf{Z}}(1)(k(s))$ .

\underline{Corollaire} 3.4. (\underline{utilise la pureté}). \underline{Supposons} $f$ \underline{propre}. \underline{Soit} $N$ \underline{le plus petit commun multiple des multiplicités} $m_{j}$ (j$\in$B). \underline{Soit} $q \geqslant 0$ , \underline{et} $q'$ \underline{la borne inférieure de} $q$ \underline{et du nombre maximum de diviseurs} $D_{j}$ \underline{passant par un même point}. \underline{Pour} $T$ \underline{dans le groupe d'inertie modéré} $I_{t} =I/P$ , \underline{on a}

\SGAvar{\begin{equation*}
(T^{N} - I)^{q'} \;=\; 0 \quad \underline{\text{sur}} \ \ H^{q}(X_{\bar{\eta}} \,,\Lambda\ )^{P} \quad .
\end{equation*}
Résulte aussitôt de 3.3 et de la variante 2.7 de 2.5.

}{\begin{equation*}
(T^{N} - I)^{q'+1} \;=\; 0 \quad \underline{\text{sur}} \ \ H^{q}(X_{\bar{\eta}} \,,\Lambda\ )^{P} \quad .
\end{equation*}
Résulte aussitôt de 3.3 et de la variante 2.7 de 2.5.

}{F7I-I-016}Les résultats de pureté requis sont disponibles pour les  $H^{1}$ , en caractéristique $0$ , ou pour  $S$  essentiellement de type fini sur un corps.

\underline{Théorème} 3.5. \underline{Soit} $C$ \underline{une courbe projective et lisse sur le corps des fractions} k($\eta$) \underline{d'un anneau de valuation discrète} $V$ . \underline{Le groupe d'inertie} $I$ \underline{agit sur} $H^{1}(C_{\bar{\eta}},\mathbf{Q}_{\ell})$ \underline{par automorphismes quasi-unipotents d'échelon} $2$ : \underline{il existe} $N$ \underline{tel que,} \underline{pour} $T \in I$ , \underline{on ait}

\begin{equation*}
(T^{N} - I)^{2} \;=\; 0 \quad \text{sur} \ \ H^{1}(C_{\bar{\eta}},\mathbf{Q}_{\ell}) \qquad .
\end{equation*}
On se ramène à supposer  $V$  strictement local. L'assertion est qu'il existe un sous-groupe ouvert  $I_{1}$  de  $I$  tel que  $(T - I)^{2} = 0$  sur  $H^{1}$  pour  $T \in I_{1}$ .

\SGApage{I}{13}

Il nous est loisible d'étendre les scalaires de  $k(\eta)$  à une extension finie (remplacer  $V$  par son normalisé dans cette extension). Puisque l'image du pro-$p$-groupe  $P$  dans un groupe de  Lie $\ell$-adique  $GL(n,Q_\ell)$  est finie, on peut en particulier supposer que  $P$  agit trivialement.

Soit  $C_o$  un modèle minimal de  $C$  sur  $V$  (pour l'existence de  $C_o$ , basée sur le théorème d'Abhyankar, cf. la discussion dans  [3] § 2, p. 87 ; on pourrait ici se ramener au cas  $V$  excellent en passant au complété de  $V$ : invoquer SGA 4 XVI 1.6).

Soit  $C_1$  un modèle régulier de  $C$ , déduit de  $C_o$  par éclatements, de fibre spéciale un diviseur à croisements  normaux. Les arguments 3.3, 3.4 s'appliquent à  $C_1$ , $\psi_t^0$ , $\psi_t^1$  et  $H^1(C_{\bar\eta},\mathbb{Z}/\ell^n)^P = H^1(C_{\bar\eta},\mathbb{Z}/\ell^n)$  avec  $N$  indépendant de  $n$ , et 3.5 en résulte.

\SGAvar{\underline{Remarque} 3.6.    Soit  $A$  une variété abélienne sur  $k(\eta)$ . Puisque  $A$  est quotient d'une jacobienne, on trouve encore que l'action de  $I$  sur  $H^1(A_\eta,\mathbb{Q}_\ell)$  ou  $T_\ell(A)$  est quasi-unipotente d'échelon 2.

\underline{Remarque} 3.7.   L'action de  $I$  sur  $H^1(X_{\bar\eta},\mathbb{Q})$  est en fait quasi-unipotente d'échelon 2 pour tout schéma  $X$  de type fini sur  $\eta$ .

}{\underline{Remarque} 3.6.    Soit  $A$  une variété abélienne sur  $k(\eta)$ . Puisque  $A$  est quotient d'une jacobienne, on trouve encore que l'action de  $I$  sur  $H^1(A_{\bar\eta},\mathbb{Q}_\ell)$  ou  $T_\ell(A)$  est quasi-unipotente d'échelon 2.

\underline{Remarque} 3.7.   L'action de  $I$  sur  $H^1(X_{\bar\eta},\mathbb{Q}_\ell)$  est en fait quasi-unipotente d'échelon 2 pour tout schéma  $X$  de type fini sur  $\eta$ .

}{F7I-I-017}Le théorème 3.5 (ou 3.6) est à la base de la démonstration du théorème de réduction stable pour les variétés abéliennes qui sera donnée dans l'exposé IX § 3 .

3.8.   Il est clair que lorsqu'on dispose de la pureté et de la résolution des singularités (par exemple en caractéristique 0 ) la méthode précédente démontre le théorème de monodromie, et fournit des bornes pour l'exposant de nilpotence. Ainsi, en égale caractéristique 0 , on trouve que pour  $X$  propre et lisse sur  $\eta$  et  $T$  dans un sous-groupe ouvert de  $I$ , on a

\begin{equation*}
(T - 1)^{q+1} \;=\; 0 \quad \text{sur} \quad H^q(X_{\bar\eta}\,,\, \mathbb{Z}_\ell) \qquad ( \, T \in I_1 \subset I \, ) \; .
\end{equation*}

\SGAsectiontoc{4. Critères de nullité pour les faisceaux de cycles évanescents}{4}
4. \underline{Critères de nullité pour les faisceaux de cycles évanescents}

4.1.  Soient  $S$  un trait strictement local et  $f : X \longrightarrow S$  un schéma de type fini sur  $S$ . Posons  $n = \dim(X_\eta)$ .

Si  $X'$  est l'adhérence schématique de  $X_\eta$ dans  $X$ , on a :

a) Sur  $X_s - X'_s$ ,  $\phi^{-1} = \Lambda$  et les autres  $\phi^i$  sont nuls ;


\SGApage{I}{14}

b) sur  $X'_s$ , les  $\phi^i$  et  $\psi^i$  relatifs à  $X/S$  ou à  $X'/S$  coïncident, et  $\phi^{-1} = 0$ .

Puisque  $X'$  est plat sur  $S$ , ceci permet souvent de se ramener au cas où  $X$  est plat sur  $S$ .

\underline{Théorème} 4.2.

(i) \underline{On a}  $\psi^i = 0$  \underline{pour}  $i > n$ .

(ii) \underline{Plus précisément}, \underline{pour}  $x$  \underline{un point de}  $X_s$ \underline{dont l'adhérence soit de} \underline{dimension}  $k$  ( $k = deg tr(k(x)/k(s))$ ) \underline{et pour}  $\bar{x}$  \underline{un point géométrique localisé en} $x$ , \underline{on a}

\begin{equation*}
\psi^i_{\bar{x}} \;=\; 0 \quad \underline{\text{pour}} \quad i > n-k \;\, , \;\, \text{i.e.}
\end{equation*}
\begin{equation*}
d(\psi^i) \;\leqslant\; n-i \qquad\qquad (\underline{\text{notation de}}\ \text{SGA 4 XIV 2.1}).
\end{equation*}
On a  $\psi^i_{\bar{x}} = H^i(X_{(\bar{x})\bar\eta}\, ,\, \Lambda)$ , et  $X_{(\bar{x})\bar\eta}$  est limite projective de schémas affines de dimension  $\leqslant$  $n$ sur  $k(\bar\eta)$ . L'assertion (i) résulte donc du théorème de Lefschetz affine (SGA 4 XIV 3.2) par passage à la limite.

Pour prouver (ii), on peut supposer que  $X$  est plat sur  $S$  (4.1) et qu'il existe un  $S$-morphisme quasi-fini et plat  $g : X \longrightarrow \mathbb{A}^k_S$  tel que  $g(x)$  soit le point générique de  $\mathbb{A}^k_s$ . Soit  $S'$  le trait localisé strict de  $\mathbb{A}^k_S$  en  $\bar{x}$ . On a

a) $S'_{\bar\eta}$  est le spectre d'un corps, et  $\mathrm{Gal}(\overline{\eta'}/s'_{\bar\eta})$  est un $p$-groupe  $Q$ .

b) D'après (i) pour  $g$ , on a  $H^i(X_{(x)\overline{\eta'}},\Lambda) = 0$  pour  $i > n-k$ .

c) On a  $H^i(X_{(x)\bar\eta},\Lambda) = H^i(X_{(x)\overline{\eta'}},\Lambda)^Q$  , d'où le théorème.

\underline{Corollaire} 4.3.  \underline{Si}  $f$  \underline{est propre et plat et que la dimension du lieu de non lissité} \underline{de}  $f$  \underline{dans}  $X_s$  \underline{est}  $\leqslant d$ , \underline{le morphisme de spécialisation}  $H^i(X_s) \longrightarrow H^i(X_{\bar\eta})$  \underline{est} \underline{un isomorphisme pour}  $i > n+d+1$  \underline{et un épimorphisme pour}  $i = n+d+1$ .

\SGAvar{4.4.  Pour prouver, dans certains cas, que les  $\phi^i$  pour  $i$  petit sont nuls, on s'appuyera sur la théorie de la dualité locale (directement ou via des théorèmes de Lefschetz locaux de SGA 2 XIV). Cet outil n'est disponible qu'en caractéristique 0 (ou en égale caractéristique  $p$ lorsqu'on dispose de la résolution des singularités à la Hironaka dans les dimensions considérées).


}{4.4.  Pour prouver, dans certains cas, que les  $\phi^i$  pour  $i$  petit sont nuls, on s'appuiera sur la théorie de la dualité locale (directement ou via des théorèmes de Lefschetz locaux de SGA 2 XIV). Cet outil n'est disponible qu'en caractéristique 0 (ou en égale caractéristique  $p$ lorsqu'on dispose de la résolution des singularités à la Hironaka dans les dimensions considérées).


}{F7I-I-018}\SGApage{I}{15}

\underline{Théorème} 4.5.  \underline{Supposons}  $S$  \underline{de caractéristique} 0 . \underline{Soient}  $f : X$ $\longrightarrow$ $S$  \underline{un morphisme de type fini et}  $x$  \underline{un point fermé de}  $X_s$ . \underline{On suppose que}

(a) $x$  \underline{est un point de non lissité isolé dans sa fibre}

(b) $X$  \underline{est de profondeur géométrique relative} $\geqslant$ $n$ \underline{en}  $x$ : \underline{au voisinage de} $x$ , $X$ \underline{s'identifie à un sous-schéma défini par} $k$ \underline{équations dans un schéma lisse de dimension relative} $N$ \underline{en} $x$ , \underline{avec}  $n$ $\leqslant$ $N-k$ .

\underline{Alors}  $\phi^i_x = 0$  \underline{pour}  $i < n$ .

\SGAvar{Soit  $X_x$  le localisé strict de  $X$  en  $x$  et  $V = X_x - \{x\}$  . Pour tout trait  $S'$  fini sur  $S$ , soient de même  $X'_x = X_x \times_S S'$  et  $V' = V \times_S S'$ . Pour  $S$  le normalisé de  $S$  dans  $\bar{\eta}$ , soient  $\bar{X}_x = X_x \times_S \bar{S}$ ,  $\bar{V} = V \times_S \bar{S}$ . On a

}{Soit  $X_x$  le localisé strict de  $X$  en  $x$  et  $V = X_x - \{x\}$  . Pour tout trait  $S'$  fini sur  $S$ , soient de même  $X'_x = X_x \times_S S'$  et  $V' = V \times_S S'$ . Pour  $\bar{S}$  le normalisé de  $S$  dans  $\bar{\eta}$ , soient  $\bar{X}_x = X_x \times_S \bar{S}$ ,  $\bar{V} = V \times_S \bar{S}$ . On a

}{F7I-I-019}(1) Si $\phi^i = 0$  sur  $(X_x)_s$  (i.e. en toute générisation de  $x$  dans  $X_s$ ) pour $i$ $\leqslant$ $m$ , alors

\SGAvar{\begin{equation*}
H^{i-1}_{\{x\}}(\bar{X}_x) \;=\; \widetilde{H}^{i}(\bar{V}) \;\xrightarrow{\ \sim\ }\; \widetilde{H}^{i}(\bar{V}_{\bar{\eta}}) \;=\; \phi^i \qquad\qquad \text{pour}\ \ i \leqslant m\ .
\end{equation*}
Les  $\phi^i$  sont en effet des faisceaux de cohomologie à support dans  $V_s$  de  $\bar{V}$ ; sous les hypothèses de 4.5, (a) est d'application pour  $m =$ $\infty$   (SGA 4 XV 2.1).

}{\begin{equation*}
H^{i-1}_{\{x\}}(\bar{X}_x) \;=\; \widetilde{H}^{i}(\bar{V}) \;\xrightarrow{\ \sim\ }\; \widetilde{H}^{i}(\bar{V}_{\bar{\eta}}) \;=\; \phi^i \qquad\qquad \text{pour}\ \ i \leqslant m\ .
\end{equation*}
Les  $\phi^i$  sont en effet des faisceaux de cohomologie à support dans  $V_s$  de  $\bar{V}$ ; sous les hypothèses de 4.5, (1) est d'application pour  $m =$ $\infty$   (SGA 4 XV 2.1).

}{F7I-I-020}(2) L'hypothèse (b) de 4.5 est stable par changement de base. D'après SGA 2 XIV 5.6, elle implique que

\begin{equation*}
\widetilde{H}^{i}(V') \;=\; H^{i-1}_{\{x\}}(X'_x) \;=\; 0 \qquad\qquad \text{pour}\ \ i < n\ ;
\end{equation*}
on conclut en notant que  $\widetilde{H}^{i}(\bar{V})$  est limite inductive des  $\widetilde{H}^{i}(V')$  .

D'après 4.2. et 4.5, on a

\underline{Corollaire} 4.6.  \underline{Pour}  $S$  \underline{un trait hensélien de caractéristique} 0 ,  $X/S$  \underline{une intersection complète relative de dimension relative}  $n$  \underline{et}  $x$ $\in$ $X_s$  \underline{un point de non lissité isolé dans sa fibre}, \underline{on a}  $\phi^i_x = 0$  \underline{pour}  $i$ $\neq$ $n$ .

\underline{Corollaire} 4.7.  \underline{Sous les hypothèses de} 4.6,  $\phi^n_x$  \underline{est un}  $\Lambda$-\underline{module plat}.

\SGAvar{Pour  $\Lambda'$  une $\Lambda$-algèbre finie, un general non sense (SGA 4 XVII 5.2.11) fournit en effet, pour les  $\phi^i_{\Lambda}$ , relatifs à  $\Lambda'$ ,


}{Pour  $\Lambda'$  une $\Lambda$-algèbre finie, un general non sense (SGA 4 XVII 5.2.11) fournit en effet, pour les  $\phi^i_{\Lambda'}$ , relatifs à  $\Lambda'$ ,


}{F7I-I-021}\SGApage{I}{16}

\begin{equation*}
\phi^{n-i}_{\Lambda'} \;=\; \mathrm{Tor}^{\Lambda}_{i}\ (\Lambda',\phi^n_{\Lambda})
\end{equation*}
et on applique 4.6 à  $\Lambda'$  .

\SGAvar{\underline{Variante} 4.8.  Dans 4.5, supprimons la condition (a) et soit  $Z$  le lieu de non de  $f$ . On peut encore conclure

}{\underline{Variante} 4.8.  Dans 4.5, supprimons la condition (a) et soit  $Z$  le lieu de non-lissité de  $f$ . On peut encore conclure

}{F7I-I-022}\begin{equation*}
\phi^i_x \;=\; 0 \qquad \text{pour}\ \ i < n - \dim Z_s\ .
\end{equation*}
On procède par récurrence descendante sur  dim $\overline{\{x\}}$ . L'hypothèse de récurrence permet d'appliquer 4.5 (1) avec  $m = n -$ dim $Z_s$ , et on achève la démonstration comme précédemment.

4.9.  Il est possible d'obtenir un critère de nullité utilisant seulement des informations sur les fibres géométriques de  $f$ . Dans le cas des intersections complètes, le résultat est toutefois d'une unité moins bon que 4.5.

\underline{Proposition} 4.10.  \underline{Sous les hypothèses générales de} 4.5, \underline{supposons que}

(a)  $x$  \underline{est un point de non lissité isolé dans sa fibre.}

\SGAvar{(b) \underline{Pour tout point}  $y$  \underline{de}  $X_{\bar{\eta}}$  , \underline{d'adhérence dans}  $X_{\bar{\eta}}$  \underline{de dimension} $d(y)$ , \underline{qui soit une générisation de}  $x$ , \underline{on a}  prof et$_y(X_{\bar{\eta}})$ $\geqslant$ $n - d(g)$ ,  i.e. $H^i_{\{y\}}(X_{\bar{\eta}}) = 0$ \underline{pour}  $i$ $\leqslant$ $n - d(g)$

}{(b) \underline{Pour tout point}  $y$  \underline{de}  $X_{\bar{\eta}}$  , \underline{d'adhérence dans}  $X_{\bar{\eta}}$  \underline{de dimension} $d(y)$ , \underline{qui soit une générisation de}  $x$ , \underline{on a}  prof et$_y(X_{\bar{\eta}})$ $\geqslant$ $n - d(y)$ ,  i.e. $H^i_{\{y\}}(X_{\bar{\eta}}) = 0$ \underline{pour}  $i$ $\leqslant$ $n - d(y)$

}{F7I-I-023}(c)  prof et$_x(X_s)$ $\geqslant$ $n$ ,

\underline{Alors,} $\phi^i_x = 0$ \underline{pour}  $i < n-1$ .

Gardons les notations de 4.5 ; on peut appliquer 4.5 (1) avec  $m =$ $\infty$ . D'après le théorème de Lefschetz local  (SGA 2 XIV), l'hypothèse (b) fournit

\begin{center}
\begin{tikzcd}[row sep=large, column sep=huge]
H^i(V') \arrow[r, "\sim"] \& H^i(V_s) \& {(\ i < n\text{-}1\ )} \\
H^i(V') \arrow[r, hook] \& H^i(V_s) \& {(\ i = n\text{-}1\ )}
\end{tikzcd}
\end{center}














et on conclut par (c) que  $H^i(\bar{V}) = \varinjlim H^i(V') = 0$  pour  $i$ $\leqslant$ $n-1$ , d'où 4.10.

4.11.  Comme en 4.8, on peut donner de 4.10 une variante ne supposant pas (a).

\SGApage{I}{17}


\SGAsectiontoc{5.  Action de la monodromie sur les  $\pi_1$}{5}
5.  \underline{Action de la monodromie sur les}  $\pi_1$

Dans cet § ( = exposé de Grothendieck du 19/11/1968), nous ferons usage du théorème de réduction semi-stable pour les courbes. Artin et Winters [2] ont donné une démonstration directe de ce théorème, qu'on peut aussi déduire [3] du théorème analogue pour les variétés abéliennes (IX § 3, utilisant 3.6). Nous utiliserons aussi la théorie de  Schlessinger (exposé VI de  Rim).

5.1.  Soient  $k$  un corps algébriquement clos, $C$ une courbe projective sur  $k$  et $x_1 \ldots x_n$  un ensemble fini de points fermés de  $C$ . On dit que  $X = (C ; x_1,.., x_n)$ est \underline{stable} si

(a) $C$ est réduite et ses seules singularités sont des points doubles à tangentes distinctes ;

(b) les  $x_i$  sont distincts, et des points lisses de  $C$ ;

\SGAvar{(c) si  $C_1$  est une composante irréductible de  $C$ , et que  $E$  est l'ensemble des points de la normalisée  $C_1'$  de  $C_1$  au-dessus d'un point singulier de  $C_1$  ou d'un des  $x_i$ , alors, si  $C_1'$  est de genre  $0$  (resp. $1$) ,  $E$  a au moins  $3$ (resp $1$) éléments.

}{(c) si  $C_1$  est une composante irréductible de  $C$ , et que  $E$  est l'ensemble des points de la normalisée  $C_1'$  de  $C_1$  au-dessus d'un point singulier de  $C$  ou d'un des  $x_i$ , alors, si  $C_1'$  est de genre  $0$  (resp. $1$) ,  $E$  a au moins  $3$ (resp $1$) éléments.

}{F7I-I-024}La condition (c) équivaut à chacune des suivantes :

(c')  $X$  n'a pas d'automorphismes infinitésimaux non triviaux.

(c'')  Si  $\omega$  est le faisceau inversible dualisant, le faisceau inversible $\omega' = \omega(\Sigma\ x_i)$  est ample.

Un  $S$-schéma  $C$ , muni de sections  $x_1$, ...,  $x_n$  est une \underline{courbe stable}  sur $S$  si ses fibres géométriques sont des courbes stables. Dans [3] §§ 1,2, seul le cas $n = 0$  est considéré. L'extension des résultats de loc. cit. au cas général est triviale :

(1) Pour  $X/k$  stable,  $H^1(C,\omega'^{\,\otimes n}) = 0$  pour  $n \geqslant 2$  et  $\omega'^{\,\otimes n}$  est très ample pour  $n \geqslant 3$  (cf. [3] 1.2).

(2) Le schéma formel  $M$  des modules de  $X$  est lisse, et  $C$  reste singulier au-dessus d'un diviseur à croisements normaux de  $M$ .


\SGApage{I}{18}

(3) Pour  $X$  et  $Y$  stables sur  $S$, $\mathrm{I}\underline{\mathrm{som}}_{S}(X,Y)$  est fini et non ramifié sur $S$  (isomorphismes compatibles aux points marqués).

(4) Pour  $X/k$  stable,  $\mathrm{Aut}_k(X) \longrightarrow \mathrm{Aut}_k\ \mathrm{Pic}^{\circ}(C)$  est injectif.

On a enfin

\SGAvar{\underline{Proposition} 5.1.1.  \underline{Soit}  $X_\eta = (C_\eta ; x_1, \ldots, x_n)$  \underline{stable sur le point générique}  $\eta$ \underline{d'un trait}  $S$ , \underline{avec}  $C_\eta$  \underline{lisse. Il existe une extension finie séparable}  $\eta'$  \underline{de}  $\eta$ \underline{et une courbe stable}  $X'$  \underline{sur le normalisé}  $S'$  \underline{de}  $S$  \underline{dans}  $\eta$  \underline{telle que}

\begin{equation*}
X' \otimes_{S'} \eta' = X_{\eta} \otimes_{\eta} \eta' \ .
\end{equation*}
Cela résulte du théorème de réduction semi-stable habituel qui permet de prendre  $\eta'$  tel que la fibre spéciale géométrique du modèle minimal  $C''$  de  $C_{\eta'}$ sur  $S'$  vérifie (a). On éclate ensuite, de façon itérée s'il le faut, des points \underline{lisses} de  $C''$  pour que  (b) soit vérifiée. On contracte alors les chaînes de composantes rationelles lisses de self intersection deux ne contenant aucun  $\overline{x}_i$  de la fibre spéciale et on obtient  $C'$  (cf. [3], preuve de 2.3 p. 88).

}{\underline{Proposition} 5.1.1.  \underline{Soit}  $X_\eta = (C_\eta ; x_1, \ldots, x_n)$  \underline{stable sur le point générique}  $\eta$ \underline{d'un trait}  $S$ , \underline{avec}  $C_\eta$  \underline{lisse. Il existe une extension finie séparable}  $\eta'$  \underline{de}  $\eta$ \underline{et une courbe stable}  $X'$  \underline{sur le normalisé}  $S'$  \underline{de}  $S$  \underline{dans}  $\eta'$  \underline{telle que}

\begin{equation*}
X' \otimes_{S'} \eta' = X_{\eta} \otimes_{\eta} \eta' \ .
\end{equation*}
Cela résulte du théorème de réduction semi-stable habituel qui permet de prendre  $\eta'$  tel que la fibre spéciale géométrique du modèle minimal  $C''$  de  $C_{\eta'}$ sur  $S'$  vérifie (a). On éclate ensuite, de façon itérée s'il le faut, des points \underline{lisses} de  $C''$  pour que  (b) soit vérifiée. On contracte alors les chaînes de composantes rationnelles lisses de self intersection $-2$ ne contenant aucun  $\overline{x}_i$  de la fibre spéciale et on obtient  $C'$  (cf. [3], preuve de 2.3 p. 88).

}{F7I-I-025}On peut vérifier comme dans loc. cit. que  $C'$  existe dès que la jacobienne de  $C_\eta$  a réduction stable. L'hypothèse de lissité sur  $C_\eta$  est par ailleurs superflue;

5.2.  Soit  $X_\eta$  un schéma géométriquement connexe de type fini sur le corps des fractions d'un trait strictement local  $S$ . Pour  $\xi$  un point géométrique de  $X_{\overline{\eta}}$  on a une suite exacte

\SGAvar{\begin{equation}\tag{5.2.1}
\begin{tikzcd}[column sep=large, row sep=tiny]
0 \arrow[r, dash] \& \pi_1(X_{\overline{\eta}},\xi) \arrow[r] \& \pi_1(X,\xi) \arrow[r] \& \mathrm{Gal}(\overline{\eta}/\eta) \arrow[r] \& 0
\end{tikzcd}
\end{equation}\ ,










}{\begin{equation}\tag{5.2.1}
\begin{tikzcd}[column sep=large, row sep=tiny]
0 \arrow[r] \& \pi_1(X_{\overline{\eta}},\xi) \arrow[r] \& \pi_1(X,\xi) \arrow[r] \& \mathrm{Gal}(\overline{\eta}/\eta) \arrow[r] \& 0
\end{tikzcd}
\end{equation}\ ,










}{F7I-I-026}d'où un homomorphisme de  $\mathrm{Gal}(\overline{\eta}/\eta)$  dans le groupe des automorphismes extérieurs de $\pi_1(X_{\overline{\eta}})$ . Passant au plus grand quotient premier à  $p$  $\pi_1^{(p)}(X_{\overline{\eta}})$  de  $\pi_1(X_{\overline{\eta}})$ , on trouve

\begin{equation}\tag{5.2.2}
\begin{tikzcd}[column sep=large, row sep=tiny]
\mathrm{Gal}(\overline{\eta}/\eta) \arrow[r] \& \mathrm{Aut}\ \mathrm{ext}\ (\pi_1^{(p)}\ (X_{\overline{\eta}}))
\end{tikzcd}
\end{equation}\ .

\SGAvar{Supposons que  $X_\eta$  ait un point rationnel  $x$ , et soit  $\overline{x}$  le point géométrique  $\overline{\eta} \longrightarrow \eta \xrightarrow{\ x\ } X_\eta$ . Pour  $\xi = \overline{x}$ , la suite exacte (5.2.2) splitte canoniquement, d'où


}{Supposons que  $X_\eta$  ait un point rationnel  $x$ , et soit  $\overline{x}$  le point géométrique  $\overline{\eta} \longrightarrow \eta \xrightarrow{\ x\ } X_\eta$ . Pour  $\xi = \overline{x}$ , la suite exacte (5.2.1) splitte canoniquement, d'où


}{F7I-I-027}\SGApage{I}{19}

\begin{equation}\tag{5.2.3}
[x] \ : \ \mathrm{Gal}(\bar{\eta}/\eta) \longrightarrow \mathrm{Aut}\ (\pi_1(X_{\bar{\eta}},\bar{x})) \longrightarrow \mathrm{Aut}\ (\pi_1^{(p)}(X_{\bar{\eta}},\bar{x}))\ .
\end{equation}
\underline{Théorème} 5.3. \underline{L'image de} $P$ \underline{par} 5.2.2 (\underline{ou} 5.2.3, \underline{cela revient au même}) \underline{est finie}.


5.3.1. \underline{Réduction au cas où} $X_\eta$ \underline{est géométriquement normal et intègre.}

\SGAvar{Soient $X'_\eta$ le normalisé de $X_\eta$ , $(X'_{i\eta}\,)_i$ la famille des composantes connexes de $X'_\eta$ , $X''_\eta = X'_\eta \times_{X_\eta} X'_\eta$ et $(X''_{j\eta}\,)_j$ les composantes connexes de $X''_\eta$ . Quitte à passer à une extension finie de $k(\eta)$, on peut supposer que les $X'_{i\eta}$ sont géométriquement normaux et intègres, avec un point rationnel $x_i$ , et que les $X''_{j\eta}$ sont géométriquement connexes, avec un point rationnel $y_j$ . Pour $k = 1$ ou $2$ et $p_k(y_j) \in X'_{i\eta}$ , choisissons une classe de chemins $\ell_{kij}$ de $\mathrm{pr}_k(\bar{y}_j)$ à $\bar{x}_i$ : c'est un élément d'un torseur sous $\pi_1(X'_{i\bar{\eta}},\bar{x}_i)$ . Soit $\bar{\ell}_{kij}$ le chemin "premier à $p$ " correspondant (le même modulo $\mathrm{Ker}(\pi \longrightarrow \pi^{(p)})$ . On peut prendre $\bar{\ell}_{kij}$ invariant par $P$ : l'ensemble des chemins premiers à $p$ est $\varprojlim$ d'ensembles finis d'ordre premier à $p$ sur lesquels le pro-$p$-groupe $P$ agit.

Dès lors, si un sous-groupe d'indice fini $P'$ de $P$ agit trivialement sur les $\pi_1^{(p)}(X'_{i\bar{\eta}},\bar{x}_i)$ , il résulte des formules de Van Kampen (i.e. de ce que $X'_{i\bar{\eta}} \longrightarrow X_{i\bar{\eta}}$ est de descente effective pour les revêtements étales) que $P'$ agit trivialement sur $\pi_1^{(p)}(X_{\bar{\eta}},\bar{x})$ .

5.3.2. Pour $X_{\bar{\eta}}$ normal et $U$ un ouvert, $\pi_1(U)$ s'envoie sur $\pi_1(X_{\bar{\eta}})$ ; on peut donc supposer $X_{\bar{\eta}}$ lisse et affine. Coupant par les sections hyperplanes \underline{génériques} (ceci remplace $S$ par le localisé strict de $S$ au point générique de la fibre spéciale d'un $\mathbb{A}^n_S$) et appliquant Bertini, on se ramène à supposer que $X_\eta$ est \underline{lisse, géométriquement connexe et de dimension un}. Quitte à passer à une extension finie de $k(\eta)$ , on peut supposer que $X_\eta$ est le complément dans une courbe projective et lisse $\bar{X}_\eta$ d'un ensemble fini de points rationnels $x_i$ ,(dont il nous est loisible d'augmenter le nombre. Appliquant le théorème de réduction semi-stable (5.1.1), on se ramène à supposer qu'il existe une courbe stable $X = (C ; x_1 ,\ldots, x_n)$ sur $S$ telle que $X_\eta = C_\eta - \{x_1 ,\ldots, x_n\}$ . On applique alors le résultat suivant


}{Soient $X'_\eta$ le normalisé de $X_\eta$ , $(X'_{i\eta}\,)_i$ la famille des composantes connexes de $X'_\eta$ , $X''_\eta = X'_\eta \times_{X_\eta} X'_\eta$ et $(X''_{j\eta}\,)_j$ les composantes connexes de $X''_\eta$ . Quitte à passer à une extension finie de $k(\eta)$, on peut supposer que les $X'_{i\eta}$ sont géométriquement normaux et intègres, avec un point rationnel $x_i$ , et que les $X''_{j\eta}$ sont géométriquement connexes, avec un point rationnel $y_j$ . Pour $k = 1$ ou $2$ et $p_k(y_j) \in X'_{i\eta}$ , choisissons une classe de chemins $\ell_{kij}$ de $\mathrm{pr}_k(\bar{y}_j)$ à $\bar{x}_i$ : c'est un élément d'un torseur sous $\pi_1(X'_{i\bar{\eta}},\bar{x}_i)$ . Soit $\bar{\ell}_{kij}$ le chemin "premier à $p$ " correspondant (le même modulo $\mathrm{Ker}(\pi \longrightarrow \pi^{(p)})$) . On peut prendre $\bar{\ell}_{kij}$ invariant par $P$ : l'ensemble des chemins premiers à $p$ est $\varprojlim$ d'ensembles finis d'ordre premier à $p$ sur lesquels le pro-$p$-groupe $P$ agit.

Dès lors, si un sous-groupe d'indice fini $P'$ de $P$ agit trivialement sur les $\pi_1^{(p)}(X'_{i\bar{\eta}},\bar{x}_i)$ , il résulte des formules de Van Kampen (i.e. de ce que $X'_{\bar{\eta}} \longrightarrow X_{\bar{\eta}}$ est de descente effective pour les revêtements étales) que $P'$ agit trivialement sur $\pi_1^{(p)}(X_{\bar{\eta}},\bar{x})$ .

5.3.2. Pour $X_{\bar{\eta}}$ normal et $U$ un ouvert, $\pi_1(U)$ s'envoie sur $\pi_1(X_{\bar{\eta}})$ ; on peut donc supposer $X_{\bar{\eta}}$ lisse et affine. Coupant par les sections hyperplanes \underline{génériques} (ceci remplace $S$ par le localisé strict de $S$ au point générique de la fibre spéciale d'un $\mathbb{A}^n_S$) et appliquant Bertini, on se ramène à supposer que $X_\eta$ est \underline{lisse, géométriquement connexe et de dimension un}. Quitte à passer à une extension finie de $k(\eta)$ , on peut supposer que $X_\eta$ est le complément dans une courbe projective et lisse $\bar{X}_\eta$ d'un ensemble fini de points rationnels $x_i$ ,(dont il nous est loisible d'augmenter le nombre). Appliquant le théorème de réduction semi-stable (5.1.1), on se ramène à supposer qu'il existe une courbe stable $X = (C ; x_1 ,\ldots, x_n)$ sur $S$ telle que $X_\eta = C_\eta - \{x_1 ,\ldots, x_n\}$ . On applique alors le résultat suivant


}{F7I-I-028}\SGApage{I}{20}

\underline{Théorème} 5.4. \underline{Soient} $S$ \underline{un schéma strictement local,} $f : \bar{X} \longrightarrow S$ \underline{un morphisme propre et plat, de fibre spéciale une courbe dont les seuls points de non lissité sont des points doubles ordinaires et} $Y$ \underline{un sous-schéma de} $\bar{X}$ , \underline{réunion d'un nombre fini de sections disjointes contenues dans l'ouvert de lissité}. Soit $x$ \underline{un point rationnel de} $X = \bar{X} - Y$ . \underline{Soit} $U$ \underline{l'ouvert de} $S$ \underline{au-dessus duquel} $X$ \underline{est lisse, et} $\xi$ \underline{un point géométrique de} $U$ . \underline{Alors l'image de} $\pi_1(U,\xi)$ \underline{dans} $\mathrm{Aut}(\pi_1^{(p)}(X_\xi,\ x\xi))$ \underline{est abélienne et première à} $p$ .

On prendra garde qu'on ne peut faire agir $\pi_1(U,\xi)$ sur $\pi_1^{(p)}(X_\xi,x\xi)$ que parce que les $\pi_1^{(p)}$ vérifient le théorème de spécialisation.

On se ramène au cas $S$ complet, puis au cas universel où $S$ est un schéma formel de modules pour la fibre spéciale, munie de ses sections. $S$ est alors un anneau de séries formelles sur un anneau de Cohen : $S = \mathrm{Spec}(W[[t_1 \ldots t_n]]\ )$ et $U$ est le complément d'un diviseur à croisements normaux d'équation $t_1 \ldots t_m = 0$ . Dans ce cas universel, d'après le lemme d'Abhyankar, $\pi_1(U,\xi)$ est abélien et premier à $p$ , d'où le théorème.

\underline{Remarque} 5.5. On peut démontrer un résultat analogue à 5.4 en dimension relative quelconque en considérant une section plane générique de dimension un et en appliquant Bertini.


\SGAsectiontoc{6. Appendice : démonstration arithmétique du théorème de réduction stable}{6}
6. \underline{Appendice} : \underline{démonstration arithmétique du théorème de réduction stable}

(par P. Deligne)

\SGAvar{Soient $S$ un trait et $A_\eta$ une variété abélienne sur $k(\eta)$ , de dimension $d \neq 0$ . Pour $S' \longrightarrow S$ un morphisme local de traits, on note $A_{\eta'}$ la variété abélienne sur $k(\eta')$ déduite de $A_\eta$ par extension des scalaires. Soient $\mathcal{A}_{S'}$ le modèle de Néron de $A_{\eta'}$ sur $S'$ , $\mathcal{A}_{S's'}$ sa fibre spéciale et $\mathcal{A}^{\circ}_{S's'}$ la composante neutre de $\mathcal{A}_{S's'}$ . Le théorème de réduction stable est le suivant.

}{Soient $S$ un trait et $A_\eta$ une variété abélienne sur $k(\eta)$ , de dimension $d \neq 0$ . Pour $S' \longrightarrow S$ un morphisme local de traits, on note $A_{\eta'}$ la variété abélienne sur $k(\eta')$ déduite de $A_\eta$ par extension des scalaires. Soient $\mathcal{A}_{S'}$ le modèle de Néron de $A_{\eta'}$ sur $S'$ , $\mathcal{A}_{S',s'}$ sa fibre spéciale et $\mathcal{A}^{\circ}_{S',s'}$ la composante neutre de $\mathcal{A}_{S',s'}$ . Le théorème de réduction stable est le suivant.

}{F7I-I-029}\SGApage{I}{21}

\SGAvar{\underline{Théorème} 6.1. \underline{Pour} $S'$ \underline{convenable,} $A_{\eta'}$ \underline{a réduction stable,} \underline{i.e.} $\mathcal{A}^{\circ}_{S's'}$ \underline{est extension d'une variété abélienne par un tore.}

}{\underline{Théorème} 6.1. \underline{Pour} $S'$ \underline{convenable,} $A_{\eta'}$ \underline{a réduction stable,} \underline{i.e.} $\mathcal{A}^{\circ}_{S',s'}$ \underline{est extension d'une variété abélienne par un tore.}

}{F7I-I-030}Il résulte assez facilement de 6.1. qu'on peut prendre $S'$ tel que $k(\eta')$ soit une extension séparable finie de $k(\eta)$ . Supposons tout d'abord que le corps résiduel de $S$ est de type fini sur le corps premier (le même argument marcherait pour $k(s)$ radiciel sur un tel corps).

Soient $\ell$ un nombre premier premier à l'exposant caractéristique résiduel $p$ , $T_{\ell}(A)$ le module de Tate, $V_{\ell}(A) = T_{\ell}(A) \otimes \mathbf{Q}_{\ell}$ et $\rho$ : $\mathrm{Gal}(\bar{\eta}/\eta) \longrightarrow \mathrm{GL}(T_{\ell}(A))$ la représentation naturelle. Une extension des scalaires préliminaire  nous ramène au cas où le groupe d'inertie $I$ agit de façon unipotente (1.2) donc à travers son plus grand pro-$\ell$-groupe quotient $\mathbf{Z}_{\ell}(1)$ (0.3). Soient $T$ l'image dans $\mathrm{GL}(T_{\ell}(A))$ d'un générateur de $\mathbf{Z}_{\ell}(1)$ et $r$ le plus grand entier $\geqslant 0$ tel que $(T-1)^{r} \neq 0$ . On a :

\underline{Lemme} 6.2. \underline{L'application}

\begin{equation*}
\sigma \in I \ , \ v \in V_{\ell}(A) \longmapsto (\rho(\sigma)-1)^{r} \ (v)
\end{equation*}
\underline{induit des morphisme et isomorphisme}

\SGAvar{\begin{equation}\tag{6.2.1}
\begin{tikzcd}[column sep=huge, row sep=huge]
V_{\ell}(A)_{I} \otimes \mathbf{Q}_{\ell}(r) \arrow[r, "M"] \arrow[d] \& V_{\ell}(A)^{I} \\
V_{\ell}(A)/\mathrm{Ker}(T-1)^{r} \otimes \mathbf{Q}_{\ell}(r) \arrow[r, "\sim"] \& \mathrm{Im}(T-1)^{r} \arrow[u]
\end{tikzcd}\quad .
\end{equation}














}{\begin{equation}\tag{6.2.1}
\begin{tikzcd}[column sep=huge, row sep=huge]
V_{\ell}(A)_{I} \otimes \mathbf{Q}_{\ell}(r) \arrow[r, "M"] \arrow[d] \& V_{\ell}(A)^{I} \\
\bigl(V_{\ell}(A)/\mathrm{Ker}(T-1)^{r}\bigr) \otimes \mathbf{Q}_{\ell}(r) \arrow[r, "\sim"] \& \mathrm{Im}(T-1)^{r} \arrow[u]
\end{tikzcd}\quad .
\end{equation}














}{F7I-I-031}$V_{\ell}(A)_{I}$ \underline{et} $V_{\ell}(A)^{I}$ \underline{sont des} $\mathrm{Gal}(\bar{s}/s)$\underline{-modules et} $M$ \underline{est galoisien.}

6.3. On dit qu'une représentation $\ell$-adique $\tau$ : $\mathrm{Gal}(\bar{s}/s) \longrightarrow \mathrm{GL}(V)$ est de poids $n$ ($n \in \mathbf{Z}$ ) si la condition suivante est vérifiée :

(*) Pour un modèle convenable $X$ de $k(s)$ ($X$ = une variété irréductible de type fini sur le corps premier, avec $k(s) = k(X)$), $\tau$ se factorise par $\pi_{1}(X,\bar{s})$ et pour $x$ un point fermé de $X$ , les valeurs propres de $\tau(F_{x})$ sont des nombres algébriques dont les conjugués complexes sont tous de valeur absolue $|k(x)|^{n/2}$ ($F_{x}$ désigne le Frobenius géométrique $\phi_{x}^{-1}$ ).


\SGApage{I}{22}

Si $V$ (resp. $W$ ) est de poids $n$ (resp. $m$ ) et que $n \neq m$ , on a $\mathrm{Hom}_{\mathrm{Gal}}(V,W) = 0$ .

6.4. Il résulte de la propriété universelle du modèle de Néron que $V_{\ell}(A)^{I} = V_{\ell}(\mathcal{A}^{\circ}_{S,s})$ . Si $a$ est la dimension du plus grand quotient abélien de $\mathcal{A}^{\circ}_{S,s}$ et $m$ la dimension de la partie multiplicative de $\mathcal{A}^{\circ}_{S,s}$ , d'après Weil, $V_{\ell}(A)^{I}$ est donc extension d'une représentation de dimension $2a$ et poids $-1$ par une représentation de dimension $m$ et poids $-2$ .

\SGAvar{6.5. Soit $A^{*}$ la variété abélienne duale de $A$ . Rappelons que $V_{\ell}(A^{*})$ et $V_{\ell}(A)$ (donc $V_{\ell}(A^{*})^{I}$ et $V_{\ell}(A)_{I}$ sont en dualité à valeurs dans $\mathbf{Q}_{\ell}(1)$. Si $a^{*}$ et $m^{*}$ sont définis pour $A^{*}$ comme $a$ et $m$ pour $A$ , il résulte de 6.4 que $V_{\ell}(A)_{I}$ est extension d'une représentation de dimension $m^{*}$ et de poids 0 par une représentation de dimension $2a^{*}$ et de poids $-1$ .

}{6.5. Soit $A^{*}$ la variété abélienne duale de $A$ . Rappelons que $V_{\ell}(A^{*})$ et $V_{\ell}(A)$ (donc $V_{\ell}(A^{*})^{I}$ et $V_{\ell}(A)_{I}$) sont en dualité à valeurs dans $\mathbf{Q}_{\ell}(1)$. Si $a^{*}$ et $m^{*}$ sont définis pour $A^{*}$ comme $a$ et $m$ pour $A$ , il résulte de 6.4 que $V_{\ell}(A)_{I}$ est extension d'une représentation de dimension $m^{*}$ et de poids 0 par une représentation de dimension $2a^{*}$ et de poids $-1$ .

}{F7I-I-032}6.6. Le module galoisien $\mathbf{Q}_{\ell}(r)$ est de poids $-2r$ . Puisque $\mathrm{Im}(T-1)^{r} \neq 0$ , on déduit de 6.4, 6.5 et de l'isomorphisme 6.2 que $r \leqslant 1$ . Si $r = 0$ , on a

\begin{equation*}
\left\{
\begin{array}{l}
a \ + \ m \ \leqslant \ d \\[2ex]
2a + \ m \ = \ 2d \ ,
\end{array}
\right.
\end{equation*}
donc $m = 0$ , $a = d$ et $A_{\eta}$ a bonne réduction. Supposons que $r = 1$ . On tire alors de 6.2 que $\mathrm{Im}(T-1)$ est de poids $-2$ et que $V_{\ell}(A)/V_{\ell}(A)^{I}$ est de poids 0 ; ces espaces ont même dimension $m_{1} \leqslant m$ .

On a alors

\begin{equation*}
\dim V_{\ell}(A)^{I} \ = \ 2d \ - \ m_{1} \ = \ 2a \ + \ m \ ,
\end{equation*}
\begin{equation*}
2 \dim \mathcal{A}_{S,s} \ \geqslant \ 2 \ (a + m) \ \geqslant \ 2 \ a \ + \ m \ + \ m_{1} \ = \ 2 \ d \ = \ 2 \dim \mathcal{A}_{S,s} \ ,
\end{equation*}
de sorte que $A_{\eta}$ a réduction stable ($\dim \mathcal{A}_{S,s} = a + m$) . On a obtenu en même temps que $m_{1} = m$ , i.e. que $\mathrm{Im}(T-1) \subset V_{\ell}(A)^{I} \simeq V_{\ell}(\mathcal{A}^{\circ}_{S,s})$ est le $V_{\ell}$ de la partie multiplicative de $\mathcal{A}^{\circ}_{S,s}$ .

6.7. Pour traiter le cas général, nous utiliserons (0.5) (cf. 1.3). On se ramène à supposer que


\SGApage{I}{23}

(a) les conditions de (0.5.2) ou (0.5.3)  sont vérifiées ;

(b) $\mathrm{Gal}(\bar{\eta}/\eta)$  agit trivialement sur  $A_\ell$ ;

(c) l'action de  $I$  sur  $T_\ell(A)$  est unipotente.

Sous ces hypothèses, nous prouverons que  $A_\eta$  a réduction stable. L'hypothèse (a) permet d'appliquer (0.5.1) pour avoir  $S = \varprojlim S_i$  comme en 1.3, dont nous reprenons les notations. Soit  $s_i$  le point fermé de  $S_i$ .

\SGAvar{Pour  $i$  assez grand, la composante neutre  $\mathcal{A}^\circ$  du modèle de Néron  $\mathcal{A}_S$  provient d'un schéma en groupe lisse à fibres connexes  $\mathcal{A}_i$  sur  $S_i$ , et  $(\mathcal{A}_i)_\ell$  est constant sur  $S_i - D$ . Le groupe  $\pi_1(S_i - D_i , \bar{\eta})$  agit sur  $T_\ell(A)$ , et agit trivialement sur  $T_\ell(A)/\ell\, T_\ell(A) \simeq A_\ell$ . Comme en 1.3, on voit que  $P_i$  agit trivialement. On a alors par (1.3.1) :

\begin{equation*}
T_\ell(A)_I \;=\; T_\ell(A)_{I_i} \;,\qquad T_\ell(\mathcal{A}_{i,s_i}) \;=\; T_\ell(\mathcal{A}^\circ_s) \;=\; T_\ell(A)^I \;=\; T_\ell(A^{I_i})
\end{equation*}
et (6.2.1) est un diagramme de  $\mathrm{Gal}(\bar{s}_i/s_i)$-modules auquel on peut appliquer les arguments 6.4 à 6.6.


}{Pour  $i$  assez grand, la composante neutre  $\mathcal{A}^\circ$  du modèle de Néron  $\mathcal{A}_S$  provient d'un schéma en groupe lisse à fibres connexes  $\mathcal{A}_i$  sur  $S_i$ , et  $(\mathcal{A}_i)_\ell$  est constant sur  $S_i - D_i$ . Le groupe  $\pi_1(S_i - D_i , \bar{\eta})$  agit sur  $T_\ell(A)$ , et agit trivialement sur  $T_\ell(A)/\ell\, T_\ell(A) \simeq A_\ell$ . Comme en 1.3, on voit que  $P_i$  agit trivialement. On a alors par (1.3.1) :

\begin{equation*}
T_\ell(A)_I \;=\; T_\ell(A)_{I_i} \;,\qquad T_\ell(\mathcal{A}_{i,s_i}) \;=\; T_\ell(\mathcal{A}^\circ_s) \;=\; T_\ell(A)^I \;=\; T_\ell(A)^{I_i}
\end{equation*}
et (6.2.1) est un diagramme de  $\mathrm{Gal}(\bar{s}_i/s_i)$-modules auquel on peut appliquer les arguments 6.4 à 6.6.


}{F7I-I-033}\SGApage{I}{24}


\underline{BIBLIOGRAPHIE}

[1]	M. ARTIN	Algebraic approximation of structures over complete local rings. Publ. Math. I.H.E.S. \underline{36} -1969. pp. 23-58.

[2]	M. ARTIN and G. WINTERS. Degenerate fibres and reduction of curves.

\SGAvar{[3]	P. DELIGNE and D. MUMFORD. The irreductibility of the space of curves of a given genus. Publ. Math. I.H.E.S. \underline{36} - 1969. pp. 75 - 110.

}{[3]	P. DELIGNE and D. MUMFORD. The irreducibility of the space of curves of a given genus. Publ. Math. I.H.E.S. \underline{36} - 1969. pp. 75 - 110.

}{F7I-I-034}[4]	J.P. SERRE	Corps locaux. Publ. Math. univ. Nancago-Hermann 1962.

[5]	J.P. SERRE and J. TATE. Good reduction of abelian varieties. Ann. of Math. \underline{88} 3 - 1968 - pp. 492 - 517.
\selectlanguage{french}

\ifnum\SGAform=1
\clearpage\section*{Apparat critique}\small
Les lectures ci-dessous sont données dans le TeX exact des témoins, afin de préserver les distinctions de notation. Les indices décimaux renvoient aux différences individuelles du passage marqué.

\par\medskip\noindent\hypertarget{app-F7I-I-001.1}{}\textbf{F7I-I-001.1}\quad p. 2; misprint\_pass.\par
\noindent Imprimé : {\ttfamily\footnotesize Hironoka\allowbreak{}}\par
\noindent Corrigé : {\ttfamily\footnotesize Hironaka\allowbreak{}}\par
\noindent Source : 7I:0. \nolinkurl{fixes/FIXLOG_7I_CURRENT.jsonl#607}; \nolinkurl{fixes/MISPRINTS_7I.json#0}\par

\par\medskip\noindent\hypertarget{app-F7I-I-002.1}{}\textbf{F7I-I-002.1}\quad p. 4; printed misprint.\par
\noindent Imprimé : {\ttfamily\footnotesize u\allowbreak{}}\par
\noindent Corrigé : {\ttfamily\footnotesize \textbackslash{}sigma\allowbreak{}}\par
\noindent Source : 7I:2. \nolinkurl{fixes/MISPRINTS_7I.json#2}\par

\par\medskip\noindent\hypertarget{app-F7I-I-003.1}{}\textbf{F7I-I-003.1}\quad p. 4; misprint\_pass.\par
\noindent Imprimé : \emph{(vide)}\par
\noindent Corrigé : {\ttfamily\footnotesize \textbackslash{}bar\allowbreak{}}\par
\noindent Source : 7I:3. \nolinkurl{fixes/FIXLOG_7I_CURRENT.jsonl#608}; \nolinkurl{fixes/MISPRINTS_7I.json#3}\par

\par\medskip\noindent\hypertarget{app-F7I-I-004.1}{}\textbf{F7I-I-004.1}\quad p. 4; misprint\_pass.\par
\noindent Imprimé : {\ttfamily\footnotesize algébres\allowbreak{}}\par
\noindent Corrigé : {\ttfamily\footnotesize algèbres\allowbreak{}}\par
\noindent Source : 7I:4. \nolinkurl{fixes/FIXLOG_7I_CURRENT.jsonl#609}\par

\par\medskip\noindent\hypertarget{app-F7I-I-005.1}{}\textbf{F7I-I-005.1}\quad p. 5; misprint\_pass.\par
\noindent Imprimé : {\ttfamily\footnotesize algèbre\allowbreak{}}\par
\noindent Corrigé : {\ttfamily\footnotesize algèbres\allowbreak{}}\par
\noindent Source : 7I:5. \nolinkurl{fixes/FIXLOG_7I_CURRENT.jsonl#610}\par

\par\medskip\noindent\hypertarget{app-F7I-I-006.1}{}\textbf{F7I-I-006.1}\quad p. 5; printed misprint.\par
\noindent Imprimé : {\ttfamily\footnotesize  .\allowbreak{} }\par
\noindent Corrigé : \emph{(2 espace(s), 0 retour(s) de ligne dans le témoin TeX)}\par
\noindent Source : 7I:6. \nolinkurl{fixes/MISPRINTS_7I.json#6}\par

\par\medskip\noindent\hypertarget{app-F7I-I-007.1}{}\textbf{F7I-I-007.1}\quad p. 6; inherited corrected-layer difference.\par
\noindent Imprimé : {\ttfamily\footnotesize 1\allowbreak{}}\par
\noindent Corrigé : {\ttfamily\footnotesize 0\allowbreak{}}\par
\noindent Source : ADOPT-F7I-I-007.1. \nolinkurl{fr/sga7i/diplomatic/expose_I_body.tex}; \nolinkurl{fr/sga7i/current/expose_I_body.tex}\par

\par\medskip\noindent\hypertarget{app-F7I-I-007.2}{}\textbf{F7I-I-007.2}\quad p. 6; inherited corrected-layer difference.\par
\noindent Imprimé : {\ttfamily\footnotesize 1\allowbreak{}}\par
\noindent Corrigé : {\ttfamily\footnotesize 0\allowbreak{}}\par
\noindent Source : ADOPT-F7I-I-007.2. \nolinkurl{fr/sga7i/diplomatic/expose_I_body.tex}; \nolinkurl{fr/sga7i/current/expose_I_body.tex}\par

\par\medskip\noindent\hypertarget{app-F7I-I-007.3}{}\textbf{F7I-I-007.3}\quad p. 6; misprint\_pass.\par
\noindent Imprimé : {\ttfamily\footnotesize débarasser\allowbreak{}}\par
\noindent Corrigé : {\ttfamily\footnotesize débarrasser\allowbreak{}}\par
\noindent Source : 7I:7. \nolinkurl{fixes/FIXLOG_7I_CURRENT.jsonl#611}\par

\par\medskip\noindent\hypertarget{app-F7I-I-008.1}{}\textbf{F7I-I-008.1}\quad p. 7; designated-witness attribution.\par
\noindent Imprimé : \emph{(vide)}\par
\noindent Corrigé : {\ttfamily\footnotesize \_o\allowbreak{}}\par
\noindent Source : WITNESS-F7I-I-008.1. \nolinkurl{fr/sga7i/diplomatic/expose_I_body.tex}; \nolinkurl{fr/sga7i/current/expose_I_body.tex}\par

\par\medskip\noindent\hypertarget{app-F7I-I-008.2}{}\textbf{F7I-I-008.2}\quad p. 7; designated-witness attribution.\par
\noindent Imprimé : \emph{(vide)}\par
\noindent Corrigé : {\ttfamily\footnotesize \_o\allowbreak{}}\par
\noindent Source : WITNESS-F7I-I-008.2. \nolinkurl{fr/sga7i/diplomatic/expose_I_body.tex}; \nolinkurl{fr/sga7i/current/expose_I_body.tex}\par

\par\medskip\noindent\hypertarget{app-F7I-I-008.3}{}\textbf{F7I-I-008.3}\quad p. 7; designated-witness attribution.\par
\noindent Imprimé : \emph{(vide)}\par
\noindent Corrigé : {\ttfamily\footnotesize \_o\allowbreak{}}\par
\noindent Source : WITNESS-F7I-I-008.3. \nolinkurl{fr/sga7i/diplomatic/expose_I_body.tex}; \nolinkurl{fr/sga7i/current/expose_I_body.tex}\par

\par\medskip\noindent\hypertarget{app-F7I-I-008.4}{}\textbf{F7I-I-008.4}\quad p. 7; printed misprint.\par
\noindent Imprimé : {\ttfamily\footnotesize Gal\allowbreak{}}\par
\noindent Corrigé : {\ttfamily\footnotesize GL\allowbreak{}}\par
\noindent Source : 7I:8. \nolinkurl{fixes/MISPRINTS_7I.json#8}\par

\par\medskip\noindent\hypertarget{app-F7I-I-009.1}{}\textbf{F7I-I-009.1}\quad p. 8; printed misprint.\par
\noindent Imprimé : {\ttfamily\footnotesize transcendance\allowbreak{}}\par
\noindent Corrigé : {\ttfamily\footnotesize transcendante\allowbreak{}}\par
\noindent Source : 7I:10. \nolinkurl{fixes/MISPRINTS_7I.json#10}\par

\par\medskip\noindent\hypertarget{app-F7I-I-009.2}{}\textbf{F7I-I-009.2}\quad p. 8; misprint\_pass.\par
\noindent Imprimé : {\ttfamily\footnotesize débarasse\allowbreak{}}\par
\noindent Corrigé : {\ttfamily\footnotesize débarrasse\allowbreak{}}\par
\noindent Source : 7I:11. \nolinkurl{fixes/FIXLOG_7I_CURRENT.jsonl#612}\par

\par\medskip\noindent\hypertarget{app-F7I-I-010.1}{}\textbf{F7I-I-010.1}\quad p. 8; printed misprint.\par
\noindent Imprimé : {\ttfamily\footnotesize \_\allowbreak{}}\par
\noindent Corrigé : {\ttfamily\footnotesize \textasciicircum{}\allowbreak{}}\par
\noindent Source : 7I:12. \nolinkurl{fixes/MISPRINTS_7I.json#12}\par

\par\medskip\noindent\hypertarget{app-F7I-I-011.1}{}\textbf{F7I-I-011.1}\quad p. 8; printed misprint.\par
\noindent Imprimé : \emph{(vide)}\par
\noindent Corrigé : {\ttfamily\footnotesize \textbackslash{}bar\{\allowbreak{}}\par
\noindent Source : 7I:13. \nolinkurl{fixes/MISPRINTS_7I.json#13}\par

\par\medskip\noindent\hypertarget{app-F7I-I-011.2}{}\textbf{F7I-I-011.2}\quad p. 8; printed misprint.\par
\noindent Imprimé : \emph{(vide)}\par
\noindent Corrigé : {\ttfamily\footnotesize \}\allowbreak{}}\par
\noindent Source : 7I:13. \nolinkurl{fixes/MISPRINTS_7I.json#13}\par

\par\medskip\noindent\hypertarget{app-F7I-I-012.1}{}\textbf{F7I-I-012.1}\quad p. 9; misprint\_pass.\par
\noindent Imprimé : \emph{(vide)}\par
\noindent Corrigé : {\ttfamily\footnotesize \{\allowbreak{}}\par
\noindent Source : 7I:14. \nolinkurl{fixes/FIXLOG_7I_CURRENT.jsonl#613}; \nolinkurl{fixes/MISPRINTS_7I.json#14}\par

\par\medskip\noindent\hypertarget{app-F7I-I-012.2}{}\textbf{F7I-I-012.2}\quad p. 9; misprint\_pass.\par
\noindent Imprimé : \emph{(vide)}\par
\noindent Corrigé : {\ttfamily\footnotesize +1\}\allowbreak{}}\par
\noindent Source : 7I:14. \nolinkurl{fixes/FIXLOG_7I_CURRENT.jsonl#613}; \nolinkurl{fixes/MISPRINTS_7I.json#14}\par

\par\medskip\noindent\hypertarget{app-F7I-I-013.1}{}\textbf{F7I-I-013.1}\quad p. 10; inherited corrected-layer difference.\par
\noindent Imprimé : \emph{(vide)}\par
\noindent Corrigé : {\ttfamily\footnotesize 0.\allowbreak{}}\par
\noindent Source : ADOPT-F7I-I-013.1. \nolinkurl{fr/sga7i/diplomatic/expose_I_body.tex}; \nolinkurl{fr/sga7i/current/expose_I_body.tex}\par

\par\medskip\noindent\hypertarget{app-F7I-I-013.2}{}\textbf{F7I-I-013.2}\quad p. 10; inherited corrected-layer difference.\par
\noindent Imprimé : {\ttfamily\footnotesize X\allowbreak{}}\par
\noindent Corrigé : {\ttfamily\footnotesize x\allowbreak{}}\par
\noindent Source : ADOPT-F7I-I-013.2. \nolinkurl{fr/sga7i/diplomatic/expose_I_body.tex}; \nolinkurl{fr/sga7i/current/expose_I_body.tex}\par

\par\medskip\noindent\hypertarget{app-F7I-I-014.1}{}\textbf{F7I-I-014.1}\quad p. 10; printed misprint.\par
\noindent Imprimé : {\ttfamily\footnotesize paramètre\allowbreak{}}\par
\noindent Corrigé : {\ttfamily\footnotesize paramètres\allowbreak{}}\par
\noindent Source : 7I:17. \nolinkurl{fixes/MISPRINTS_7I.json#17}\par

\par\medskip\noindent\hypertarget{app-F7I-I-015.1}{}\textbf{F7I-I-015.1}\quad p. 11; misprint\_in\_print.\par
\noindent Imprimé : \emph{(vide)}\par
\noindent Corrigé : {\ttfamily\footnotesize \textbackslash{}bar\{\allowbreak{}}\par
\noindent Source : FIX-7I-D-0006. \nolinkurl{fixes/FIXLOG_7I.jsonl#6}\par

\par\medskip\noindent\hypertarget{app-F7I-I-015.2}{}\textbf{F7I-I-015.2}\quad p. 11; misprint\_in\_print.\par
\noindent Imprimé : \emph{(vide)}\par
\noindent Corrigé : {\ttfamily\footnotesize \}\allowbreak{}}\par
\noindent Source : FIX-7I-D-0006. \nolinkurl{fixes/FIXLOG_7I.jsonl#6}\par

\par\medskip\noindent\hypertarget{app-F7I-I-015.3}{}\textbf{F7I-I-015.3}\quad p. 11; printed misprint.\par
\noindent Imprimé : {\ttfamily\footnotesize x\allowbreak{}}\par
\noindent Corrigé : {\ttfamily\footnotesize n\allowbreak{}}\par
\noindent Source : 7I:19. \nolinkurl{fixes/MISPRINTS_7I.json#19}\par

\par\medskip\noindent\hypertarget{app-F7I-I-016.1}{}\textbf{F7I-I-016.1}\quad p. 12; inherited corrected-layer difference.\par
\noindent Imprimé : \emph{(vide)}\par
\noindent Corrigé : {\ttfamily\footnotesize +1\allowbreak{}}\par
\noindent Source : ADOPT-F7I-I-016.1. \nolinkurl{fr/sga7i/diplomatic/expose_I_body.tex}; \nolinkurl{fr/sga7i/current/expose_I_body.tex}\par

\par\medskip\noindent\hypertarget{app-F7I-I-017.1}{}\textbf{F7I-I-017.1}\quad p. 13; printed misprint.\par
\noindent Imprimé : \emph{(vide)}\par
\noindent Corrigé : {\ttfamily\footnotesize \{\textbackslash{}bar\allowbreak{}}\par
\noindent Source : 7I:21. \nolinkurl{fixes/MISPRINTS_7I.json#21}\par

\par\medskip\noindent\hypertarget{app-F7I-I-017.2}{}\textbf{F7I-I-017.2}\quad p. 13; printed misprint.\par
\noindent Imprimé : \emph{(vide)}\par
\noindent Corrigé : {\ttfamily\footnotesize \}\allowbreak{}}\par
\noindent Source : 7I:21. \nolinkurl{fixes/MISPRINTS_7I.json#21}\par

\par\medskip\noindent\hypertarget{app-F7I-I-017.3}{}\textbf{F7I-I-017.3}\quad p. 13; printed misprint.\par
\noindent Imprimé : \emph{(vide)}\par
\noindent Corrigé : {\ttfamily\footnotesize \_\textbackslash{}ell\allowbreak{}}\par
\noindent Source : 7I:22. \nolinkurl{fixes/MISPRINTS_7I.json#22}\par

\par\medskip\noindent\hypertarget{app-F7I-I-018.1}{}\textbf{F7I-I-018.1}\quad p. 14; misprint\_pass.\par
\noindent Imprimé : {\ttfamily\footnotesize appuyera\allowbreak{}}\par
\noindent Corrigé : {\ttfamily\footnotesize appuiera\allowbreak{}}\par
\noindent Source : 7I:24. \nolinkurl{fixes/FIXLOG_7I_CURRENT.jsonl#614}\par

\par\medskip\noindent\hypertarget{app-F7I-I-019.1}{}\textbf{F7I-I-019.1}\quad p. 15; misprint\_pass.\par
\noindent Imprimé : \emph{(vide)}\par
\noindent Corrigé : {\ttfamily\footnotesize \textbackslash{}bar\{\allowbreak{}}\par
\noindent Source : 7I:25. \nolinkurl{fixes/FIXLOG_7I_CURRENT.jsonl#615}\par

\par\medskip\noindent\hypertarget{app-F7I-I-019.2}{}\textbf{F7I-I-019.2}\quad p. 15; misprint\_pass.\par
\noindent Imprimé : \emph{(vide)}\par
\noindent Corrigé : {\ttfamily\footnotesize \}\allowbreak{}}\par
\noindent Source : 7I:25. \nolinkurl{fixes/FIXLOG_7I_CURRENT.jsonl#615}\par

\par\medskip\noindent\hypertarget{app-F7I-I-020.1}{}\textbf{F7I-I-020.1}\quad p. 15; misprint\_pass.\par
\noindent Imprimé : {\ttfamily\footnotesize a\allowbreak{}}\par
\noindent Corrigé : {\ttfamily\footnotesize 1\allowbreak{}}\par
\noindent Source : 7I:26. \nolinkurl{fixes/FIXLOG_7I_CURRENT.jsonl#616}\par

\par\medskip\noindent\hypertarget{app-F7I-I-021.1}{}\textbf{F7I-I-021.1}\quad p. 15; misprint\_pass.\par
\noindent Imprimé : \emph{(vide)}\par
\noindent Corrigé : {\ttfamily\footnotesize '\allowbreak{}}\par
\noindent Source : 7I:27. \nolinkurl{fixes/FIXLOG_7I_CURRENT.jsonl#617}\par

\par\medskip\noindent\hypertarget{app-F7I-I-022.1}{}\textbf{F7I-I-022.1}\quad p. 16; designated-witness attribution.\par
\noindent Imprimé : \emph{(vide)}\par
\noindent Corrigé : {\ttfamily\footnotesize -lissité\allowbreak{}}\par
\noindent Source : WITNESS-F7I-I-022.1. \nolinkurl{fr/sga7i/diplomatic/expose_I_body.tex}; \nolinkurl{fr/sga7i/current/expose_I_body.tex}\par

\par\medskip\noindent\hypertarget{app-F7I-I-023.1}{}\textbf{F7I-I-023.1}\quad p. 16; printed misprint.\par
\noindent Imprimé : {\ttfamily\footnotesize g\allowbreak{}}\par
\noindent Corrigé : {\ttfamily\footnotesize y\allowbreak{}}\par
\noindent Source : 7I:29. \nolinkurl{fixes/MISPRINTS_7I.json#29}\par

\par\medskip\noindent\hypertarget{app-F7I-I-023.2}{}\textbf{F7I-I-023.2}\quad p. 16; printed misprint.\par
\noindent Imprimé : {\ttfamily\footnotesize g\allowbreak{}}\par
\noindent Corrigé : {\ttfamily\footnotesize y\allowbreak{}}\par
\noindent Source : 7I:30. \nolinkurl{fixes/MISPRINTS_7I.json#30}\par

\par\medskip\noindent\hypertarget{app-F7I-I-024.1}{}\textbf{F7I-I-024.1}\quad p. 17; inherited corrected-layer difference.\par
\noindent Imprimé : {\ttfamily\footnotesize \_1\allowbreak{}}\par
\noindent Corrigé : \emph{(vide)}\par
\noindent Source : ADOPT-F7I-I-024.1. \nolinkurl{fr/sga7i/diplomatic/expose_I_body.tex}; \nolinkurl{fr/sga7i/current/expose_I_body.tex}\par

\par\medskip\noindent\hypertarget{app-F7I-I-025.1}{}\textbf{F7I-I-025.1}\quad p. 18; printed misprint.\par
\noindent Imprimé : \emph{(vide)}\par
\noindent Corrigé : {\ttfamily\footnotesize '\allowbreak{}}\par
\noindent Source : 7I:31. \nolinkurl{fixes/MISPRINTS_7I.json#31}\par

\par\medskip\noindent\hypertarget{app-F7I-I-025.2}{}\textbf{F7I-I-025.2}\quad p. 18; misprint\_pass.\par
\noindent Imprimé : {\ttfamily\footnotesize rationelles\allowbreak{}}\par
\noindent Corrigé : {\ttfamily\footnotesize rationnelles\allowbreak{}}\par
\noindent Source : 7I:32. \nolinkurl{fixes/FIXLOG_7I_CURRENT.jsonl#618}; \nolinkurl{fixes/MISPRINTS_7I.json#32}\par

\par\medskip\noindent\hypertarget{app-F7I-I-025.3}{}\textbf{F7I-I-025.3}\quad p. 18; designated-witness attribution.\par
\noindent Imprimé : {\ttfamily\footnotesize deux\allowbreak{}}\par
\noindent Corrigé : {\ttfamily\footnotesize \$-2\$\allowbreak{}}\par
\noindent Source : WITNESS-F7I-I-025.3. \nolinkurl{fr/sga7i/diplomatic/expose_I_body.tex}; \nolinkurl{fr/sga7i/current/expose_I_body.tex}\par

\par\medskip\noindent\hypertarget{app-F7I-I-026.1}{}\textbf{F7I-I-026.1}\quad p. 18; printed misprint.\par
\noindent Imprimé : {\ttfamily\footnotesize ,\allowbreak{} dash\allowbreak{}}\par
\noindent Corrigé : \emph{(vide)}\par
\noindent Source : 7I:33. \nolinkurl{fixes/MISPRINTS_7I.json#33}\par

\par\medskip\noindent\hypertarget{app-F7I-I-027.1}{}\textbf{F7I-I-027.1}\quad p. 18; printed misprint.\par
\noindent Imprimé : {\ttfamily\footnotesize 2\allowbreak{}}\par
\noindent Corrigé : {\ttfamily\footnotesize 1\allowbreak{}}\par
\noindent Source : 7I:35. \nolinkurl{fixes/MISPRINTS_7I.json#35}\par

\par\medskip\noindent\hypertarget{app-F7I-I-028.1}{}\textbf{F7I-I-028.1}\quad p. 19; printed misprint.\par
\noindent Imprimé : \emph{(vide)}\par
\noindent Corrigé : {\ttfamily\footnotesize )\allowbreak{}}\par
\noindent Source : 7I:37. \nolinkurl{fixes/MISPRINTS_7I.json#37}\par

\par\medskip\noindent\hypertarget{app-F7I-I-028.2}{}\textbf{F7I-I-028.2}\quad p. 19; designated-witness attribution.\par
\noindent Imprimé : {\ttfamily\footnotesize i\allowbreak{}}\par
\noindent Corrigé : \emph{(vide)}\par
\noindent Source : WITNESS-F7I-I-028.2. \nolinkurl{fr/sga7i/diplomatic/expose_I_body.tex}; \nolinkurl{fr/sga7i/current/expose_I_body.tex}\par

\par\medskip\noindent\hypertarget{app-F7I-I-028.3}{}\textbf{F7I-I-028.3}\quad p. 19; printed misprint.\par
\noindent Imprimé : {\ttfamily\footnotesize i\allowbreak{}}\par
\noindent Corrigé : \emph{(vide)}\par
\noindent Source : 7I:38. \nolinkurl{fixes/MISPRINTS_7I.json#38}\par

\par\medskip\noindent\hypertarget{app-F7I-I-028.4}{}\textbf{F7I-I-028.4}\quad p. 19; printed misprint.\par
\noindent Imprimé : \emph{(vide)}\par
\noindent Corrigé : {\ttfamily\footnotesize )\allowbreak{}}\par
\noindent Source : 7I:39. \nolinkurl{fixes/MISPRINTS_7I.json#39}\par

\par\medskip\noindent\hypertarget{app-F7I-I-029.1}{}\textbf{F7I-I-029.1}\quad p. 20; inherited corrected-layer difference.\par
\noindent Imprimé : \emph{(vide)}\par
\noindent Corrigé : {\ttfamily\footnotesize ,\allowbreak{}}\par
\noindent Source : ADOPT-F7I-I-029.1. \nolinkurl{fr/sga7i/diplomatic/expose_I_body.tex}; \nolinkurl{fr/sga7i/current/expose_I_body.tex}\par

\par\medskip\noindent\hypertarget{app-F7I-I-029.2}{}\textbf{F7I-I-029.2}\quad p. 20; inherited corrected-layer difference.\par
\noindent Imprimé : \emph{(vide)}\par
\noindent Corrigé : {\ttfamily\footnotesize ,\allowbreak{}}\par
\noindent Source : ADOPT-F7I-I-029.2. \nolinkurl{fr/sga7i/diplomatic/expose_I_body.tex}; \nolinkurl{fr/sga7i/current/expose_I_body.tex}\par

\par\medskip\noindent\hypertarget{app-F7I-I-029.3}{}\textbf{F7I-I-029.3}\quad p. 20; inherited corrected-layer difference.\par
\noindent Imprimé : \emph{(vide)}\par
\noindent Corrigé : {\ttfamily\footnotesize ,\allowbreak{}}\par
\noindent Source : ADOPT-F7I-I-029.3. \nolinkurl{fr/sga7i/diplomatic/expose_I_body.tex}; \nolinkurl{fr/sga7i/current/expose_I_body.tex}\par

\par\medskip\noindent\hypertarget{app-F7I-I-030.1}{}\textbf{F7I-I-030.1}\quad p. 21; designated-witness attribution.\par
\noindent Imprimé : \emph{(vide)}\par
\noindent Corrigé : {\ttfamily\footnotesize ,\allowbreak{}}\par
\noindent Source : WITNESS-F7I-I-030.1. \nolinkurl{fr/sga7i/diplomatic/expose_I_body.tex}; \nolinkurl{fr/sga7i/current/expose_I_body.tex}\par

\par\medskip\noindent\hypertarget{app-F7I-I-031.1}{}\textbf{F7I-I-031.1}\quad p. 21; inherited corrected-layer difference.\par
\noindent Imprimé : \emph{(vide)}\par
\noindent Corrigé : {\ttfamily\footnotesize \textbackslash{}bigl(\allowbreak{}}\par
\noindent Source : ADOPT-F7I-I-031.1. \nolinkurl{fr/sga7i/diplomatic/expose_I_body.tex}; \nolinkurl{fr/sga7i/current/expose_I_body.tex}\par

\par\medskip\noindent\hypertarget{app-F7I-I-031.2}{}\textbf{F7I-I-031.2}\quad p. 21; designated-witness attribution.\par
\noindent Imprimé : \emph{(vide)}\par
\noindent Corrigé : {\ttfamily\footnotesize \textbackslash{}bigr)\allowbreak{}}\par
\noindent Source : WITNESS-F7I-I-031.2. \nolinkurl{fr/sga7i/diplomatic/expose_I_body.tex}; \nolinkurl{fr/sga7i/current/expose_I_body.tex}\par

\par\medskip\noindent\hypertarget{app-F7I-I-032.1}{}\textbf{F7I-I-032.1}\quad p. 22; misprint\_pass.\par
\noindent Imprimé : \emph{(vide)}\par
\noindent Corrigé : {\ttfamily\footnotesize )\allowbreak{}}\par
\noindent Source : 7I:40. \nolinkurl{fixes/FIXLOG_7I_CURRENT.jsonl#619}; \nolinkurl{fixes/MISPRINTS_7I.json#40}\par

\par\medskip\noindent\hypertarget{app-F7I-I-033.1}{}\textbf{F7I-I-033.1}\quad p. 23; printed misprint.\par
\noindent Imprimé : \emph{(vide)}\par
\noindent Corrigé : {\ttfamily\footnotesize \_i\allowbreak{}}\par
\noindent Source : 7I:41. \nolinkurl{fixes/MISPRINTS_7I.json#41}\par

\par\medskip\noindent\hypertarget{app-F7I-I-033.2}{}\textbf{F7I-I-033.2}\quad p. 23; printed misprint.\par
\noindent Imprimé : \emph{(vide)}\par
\noindent Corrigé : {\ttfamily\footnotesize )\allowbreak{}}\par
\noindent Source : 7I:42. \nolinkurl{fixes/MISPRINTS_7I.json#42}\par

\par\medskip\noindent\hypertarget{app-F7I-I-033.3}{}\textbf{F7I-I-033.3}\quad p. 23; printed misprint.\par
\noindent Imprimé : {\ttfamily\footnotesize )\allowbreak{}}\par
\noindent Corrigé : \emph{(vide)}\par
\noindent Source : 7I:42. \nolinkurl{fixes/MISPRINTS_7I.json#42}\par

\par\medskip\noindent\hypertarget{app-F7I-I-034.1}{}\textbf{F7I-I-034.1}\quad p. 24; printed misprint.\par
\noindent Imprimé : {\ttfamily\footnotesize irreductibility\allowbreak{}}\par
\noindent Corrigé : {\ttfamily\footnotesize irreducibility\allowbreak{}}\par
\noindent Source : 7I:43. \nolinkurl{fixes/MISPRINTS_7I.json#43}\par

\normalsize\fi
\clearpage\SGAtoc{Exposé II — Propriétés de finitude du groupe fondamental}
\selectlanguage{french}

\SGApage{I}{25}


\underline{Exposé II}


\underline{\underline{Propriétés de finitude du groupe fondamental}}


\underline{par Michèle RAYNAUD}

Soient  $k$  un corps séparablement clos de caractéristique  $p \geqslant 0$ , $X$  un schéma connexe de type fini sur  $k$  et  $\Pi_1^{(p)}(X)$  le plus grand quotient "premier à $p$ " de  $\Pi_1(X)$ . On montre, dans cet exposé, que, si  $X$  est une surface, $\Pi_1^{(p)}(X)$  est topologiquement de présentation finie et qu'il en est de même pour  $X$  quelconque si l'on admet la résolution des singularités.

Dans le cas d'une surface, la méthode de démonstration, due à J.P. Murre, consiste à se ramener au cas où l'on a une fibration au-dessus de  $P^1$  ayant les propriétés (i), (ii) et (iii) de la proposition ci-dessous.

\underline{Proposition} 2.1.  \underline{Soient}  $k$  \underline{un corps algébriquement clos},  $S$  \underline{une surface projective}, \underline{irréductible}, \underline{lisse sur}  $k$ . \underline{Alors on peut trouver un éclatement}

\[ g : S' \longrightarrow S \quad , \]
\underline{où}  $S'$  \underline{est une surface régulière}, \underline{munie d'une fibration}

\[ f : S' \longrightarrow P^1 \quad , \]
\underline{tels que}  $f$  \underline{ait une section}  $\sigma$  , \underline{et que les conditions suivantes soient satisfaites} :

(i)   $f$  \underline{est lisse aux points de}  $\sigma(P^1)$ .

(ii)  \underline{les fibres de}  $f$  \underline{sont des courbes géométriquement intègres n'ayant que des singularités ordinaires.}

(iii) \underline{il existe un ouvert non vide}  $U$  \underline{de}  $P^1$  \underline{tel que les fibres de}  $f$  \underline{aux points de}  $U$  \underline{soient lisses.}

La proposition résulte de l'existence d'un plongement projectif  $S \longrightarrow P^r$  tel qu'on puisse trouver un pinceau de Lefschetz par rapport à ce plongement (XVII 2.5). Soit en effet  $D = \{H_t\}_{t \in P^1}$  un tel pinceau d'hyperplans ; rappelons que l'on
\SGApage{I}{26}
a les propriétés suivantes (loc. cit. 2.2) :

a) l'axe  $\Delta$  de  $D$  coupe transversalement  $S$ ; le sous-schéma fermé  $\Delta \cap S$  est donc réduit et formé d'un ensemble fini de points fermés de  $S$ , $x_1,...,x_n$ .

b) il existe un ouvert non vide  $U$  de  $P^1$  tel que, pour tout point  $t \in U$ , $H_t$  coupe transversalement  $S$ .

c) pour tout point  $t \in P^1-U$ , $H_t$  coupe transversalement  $S$  sauf en un point qui est un point singulier quadratique ordinaire.

Par chaque point  $x$  de  $S$  n'appartenant pas à  $\Delta$  passe un unique hyperplan  $H_t$ , et l'application qui, à  $x$ , associe  $t \in P^1$ , est une application rationnelle  $a$  de  $S$  dans  $P^1$  de domaine de définition  $S-(\Delta \cap S)$ .

On considère l'éclatement  $g : S' \longrightarrow S$  du sous-schéma fermé  $S \cap \Delta$  de  $S$ . Le schéma  $S'$  est régulier et la fibre  $S'_{x_i}$  au-dessus d'un point  $x_i$  s'identifie au fibré projectif défini par l'espace tangent à  $S$  en  $x_i$ , donc est isomorphe à  $P^1$  (EGA IV 19.4.2, 19.4.4). De plus l'application rationnelle  $f = ag$  est partout définie (XVIII 3.1). Enfin, on peut associer à chaque point  $x_i$  la section de  $f$  qui, à un point  $t \in P^1$ , fait correspondre le vecteur tangent en  $x_i$  à  $H_t \cap S$ .

Vérifions les conditions (i), (ii) et (iii). On note que, pour chaque  $t \in P^1$ , les courbes  $S'_t$  et  $S_t = H_t \cap S$  sont isomorphes ; le morphisme  $S'_t \longrightarrow S_t$  est en effet un isomorphisme sauf peut-être en les points  $x_i$ , points où  $S_t$  est régulière ; comme de plus  $S'_t$  est intersection complète dans  $S'$ , $S'_t$  est intègre, donc isomorphe à  $S_t$ .

\SGAvar{Les conditions (i) et (ii) ne sont autres que les conditions a) et b). Enfin  $f$  est lisse aux points de  $\sigma$  car  $f$  est plat et car les fibres  $S'_t$  sont géométriquement régulières aux points  $\sigma(t)$ .


}{Les conditions (ii) et (iii) ne sont autres que les conditions b) et c). Enfin  $f$  est lisse aux points de  $\sigma$  car  $f$  est plat et car les fibres  $S'_t$  sont géométriquement régulières aux points  $\sigma(t)$ .


}{F7I-II-001}\SGApage{I}{27}

\underline{Corollaire} 2.2. \underline{Les notations étant celles de} 2.1, \underline{soit} $Y$ \underline{un diviseur à croisements normaux dans} $S$ , $Y' = Y \times_S S'$ . \underline{Alors on peut choisir} $f$ , $g$ , $U$ \underline{tels que} $Y'\times_{P^1}U$ \underline{soit un diviseur à croisements normaux relativement à} $U$ .

\SGAvar{On sait que, s'il existe un pinceau de Lefschetz pour un plongement $S \longrightarrow P^r$ , les pinceaux de Lefschetz forment un ouvert de la variété des droites de $P^r$ (XVII 3.2.1). Il en est de même de l'ensemble des pinceaux de Lefschetz dont l'axe ne rencontre pas $Y$ et tels qu'il existe un ouvert non vide $U_1$ de $P^1$ tel que l'intersection de $Y$ et $H_t$ soit lisse si $t \in U_1$ . Il suffit alors de construire $f$ et $g$ à partir d'un pinceau satisfaisant à ces deux conditions et de remplacer par $U \cap U_1$ .

}{On sait que, s'il existe un pinceau de Lefschetz pour un plongement $S \longrightarrow P^r$ , les pinceaux de Lefschetz forment un ouvert de la variété des droites de $P^r$ (XVII 3.2.1). Il en est de même de l'ensemble des pinceaux de Lefschetz dont l'axe ne rencontre pas $Y$ et tels qu'il existe un ouvert non vide $U_1$ de $P^1$ tel que l'intersection de $Y$ et $H_t$ soit lisse si $t \in U_1$ . Il suffit alors de construire $f$ et $g$ à partir d'un pinceau satisfaisant à ces deux conditions et de remplacer $U$ par $U \cap U_1$ .

}{F7I-II-002}2.3.0. Soit $G$ un groupe profini. Rappelons que $G$ est dit topologiquement de génération finie s'il existe un entier $n$ tel que $G$ soit isomorphe à un quotient du pro-groupe libre à $n$ générateurs $F(n)$ ; soit $K = Ker(F(n) \longrightarrow G)$ ; on dit que $G$ est topologiquement de présentation finie si, de plus, on peut trouver un nombre fini d'éléments $k_1,...,k_r$ de $K$ tels que le plus petit sous-groupe invariant fermé de $K$ , contenant $k_1,...,k_r$ , soit égal à $K$ .

\SGAvar{\underline{Théorème} 2.3.1. \underline{Soient} $k$ \underline{un corps séparablement clos de caractéristique} $p \geqslant 0$ , $X$ \underline{un schéma connexe de type fini sur} $k$ . \underline{On suppose que les schémas de type fini et de dimension} $\leqslant dim(X)$ \underline{sur une clôture algébrique de} $k$ \underline{sont fortement désingularisables} (SGA 5 I 3.1.5) (\underline{condition réalisée si} $dim X = 2$ \underline{d'après} [1] \underline{ou si} $k = 0$ \underline{d'après} [2]). \underline{Alors le groupe fondamental} $\Pi_1^{(p)}(X)$ \underline{est topologiquement de présentation finie}.

}{\underline{Théorème} 2.3.1. \underline{Soient} $k$ \underline{un corps séparablement clos de caractéristique} $p \geqslant 0$ , $X$ \underline{un schéma connexe de type fini sur} $k$ . \underline{On suppose que les schémas de type fini et de dimension} $\leqslant dim(X)$ \underline{sur une clôture algébrique de} $k$ \underline{sont fortement désingularisables} (SGA 5 I 3.1.5) (\underline{condition réalisée si} $dim X = 2$ \underline{d'après} [1] \underline{ou si} $\operatorname{char}k = 0$ \underline{d'après} [2]). \underline{Alors le groupe fondamental} $\Pi_1^{(p)}(X)$ \underline{est topologiquement de présentation finie}.

}{F7I-II-003}On peut supposer $k$ algébriquement clos d'après SGA 1 IX 4.10.


1) \underline{Réduction au cas où} $X$ \underline{est une surface régulière}.

\SGAvar{Si $f : X' \longrightarrow X$ est un morphisme de type fini, on pose $X'' = X'\times_X S'$ et on note $(X'_a)_{a \in A}$ et $(X''_b)_{b \in B}$ l'ensemble des composantes connexes de $X'$ et $X''$ respectivement. Lorsque $f$ est un morphisme de descente effective pour la catégorie des revêtements étales, le groupe fondamental $\Pi_1(X)$ peut se calculer explicitement
\SGApage{I}{28}
en terme des groupes fondamentaux $\Pi_1(X'_a)$ et $\Pi_1(X''_b)$ ; il en résulte (SGA 1 IX 5.2, 5.3) que, si les $\Pi_1(X'_a)$ sont topologiquement de génération finie, il en est de même de $\Pi_1(X)$ ; si les $\Pi_1(X'_a)$ sont topologiquement de présentation finie et les $\Pi_1(X''_b)$ topologiquement de génération finie, $\Pi_1(X)$ est topologiquement de présentation finie.

}{Si $f : X' \longrightarrow X$ est un morphisme de type fini, on pose $X'' = X'\times_X X'$ et on note $(X'_a)_{a \in A}$ et $(X''_b)_{b \in B}$ l'ensemble des composantes connexes de $X'$ et $X''$ respectivement. Lorsque $f$ est un morphisme de descente effective pour la catégorie des revêtements étales, le groupe fondamental $\Pi_1(X)$ peut se calculer explicitement
\SGApage{I}{28}
en terme des groupes fondamentaux $\Pi_1(X'_a)$ et $\Pi_1(X''_b)$ ; il en résulte (SGA 1 IX 5.2, 5.3) que, si les $\Pi_1(X'_a)$ sont topologiquement de génération finie, il en est de même de $\Pi_1(X)$ ; si les $\Pi_1(X'_a)$ sont topologiquement de présentation finie et les $\Pi_1(X''_b)$ topologiquement de génération finie, $\Pi_1(X)$ est topologiquement de présentation finie.

}{F7I-II-004}Comme $X$ est désingularisable, on peut trouver un schéma régulier $X'$ et un morphisme propre birationnel $f : X' \longrightarrow X$ ; comme $f$ est un morphisme de descente effective pour la catégorie des revêtements étales, il suffit de prouver le théorème lorsqu'on fait l'hypothèse supplémentaire que $X$ est régulier. On en déduit en effet que $\Pi_1(X)$ est topologiquement de génération finie ; sachant que $\Pi_1(X')$ est topologiquement de présentation finie et les $\Pi_1(X''_b)$ topologiquement de génération finie, on en déduit que $\Pi_1(X)$ est topologiquement de présentation finie.

Soit maintenant $(U_i)_{i \in I}$ un recouvrement fini de $X$ par des ouverts affines et soit $f : \coprod U_i \longrightarrow X$ le morphisme canonique. Comme $f$ est un morphisme de descente effective pour la catégorie des revêtements étales, on voit qu'il suffit de prouver le théorème pour les $U_i$ ; on est donc ramené au cas où $X$ est affine et lisse.

On suppose désormais $X$ affine, lisse sur $k$ , et $dim X > 2$ . On peut alors appliquer le théorème de Lefschetz pour les schémas quasi-projectifs ([3] V 7.4) ; si $Y$ est une section hyperplane "assez générale" de $X$ , on a un isomorphisme

\begin{equation*}
\Pi_1(Y) \simeq \Pi_1(X) .
\end{equation*}
Il suffit donc de prouver le théorème pour $Y$ , ce qui nous ramène au cas où $dim X = 2$ .


2) \underline{Cas où} $X$ \underline{est une surface lisse sur} $k$

D'après le théorème de désingularisation forte des surfaces, on peut trouver une surface $S$ lisse, intègre, projective et une immersion ouverte $i : X \longrightarrow S$ telle que le diviseur $Y = S - X$ soit à croisements normaux. On applique le corollaire 2.2 dont on reprend les notations.

\SGApage{I}{29}

Montrons qu'il suffit de prouver que le groupe fondamental de $X' = X\times_S S'$ est topologiquement de présentation finie. Comme $X$ et $X'$ sont réguliers et comme on a

\SGAvar{\begin{equation*}
\mathrm{codim}(X- \bigcup_{1\leqslant i\leqslant n} \{x_i\},X) \geqslant 2 \qquad ,
\end{equation*}
tout revêtement étale de $X'$ induit sur $X- \bigcup_{1\leqslant i\leqslant n}\{x_i\}$ un revêtement étale qui se prolonge de façon unique en un revêtement étale de $X$ (SGA 2 1.11) ; ceci montre que l'on a un isomorphisme

}{\begin{equation*}
\mathrm{codim}\left(\bigcup_{1\leqslant i\leqslant n} \{x_i\},X\right) \geqslant 2 \qquad ,
\end{equation*}
tout revêtement étale de $X'$ induit sur $X- \bigcup_{1\leqslant i\leqslant n}\{x_i\}$ un revêtement étale qui se prolonge de façon unique en un revêtement étale de $X$ (SGA 2 1.11) ; ceci montre que l'on a un isomorphisme

}{F7I-II-005}\begin{equation*}
\Pi_1(X') \simeq \Pi_1(X) \qquad .
\end{equation*}
On considère le diagramme cartésien

\begin{tikzcd}
X' \arrow[d, "h"'] \& X'_U \arrow[l] \arrow[d, "h_U"] \\
P^1 \& U \arrow[l]
\end{tikzcd}
\qquad ,











\SGAvar{où $h = f|X'$ . On note $L$ l'ensemble des nombres premiers distincts de $p$ et $\bar{\eta}$ un point géométrique au-dessus du point générique $\eta$ de $P^1$ . Vérifions que les hypothèses de SGA 1 XIII 4.8 sont satisfaites. Le morphisme $h$ a une section $\sigma$ car il en est ainsi de $f$ . Prouvons la condition a) de SGA 1 XIII 4.7 ; les fibres de $h$ étant géométriquement intègres, le morphisme $h$ est 0-acyclique et localement 0-acyclique (SGA 4 XV 4.1, 1.16) ; il reste à montrer que, pour tout faisceau d'ensembles localement constant $F$ sur $X'$ , $(F,h)$ est cohomologiquement propre en dimension $\leqslant 0$ . Si $\xi$ est un point fermé de $P^1$ et si $Z$ désigne la fermeture intégrale de Spec $\mathcal{O}_{P^1,\xi}$ dans $\bar{\eta}$ le schéma $X'_Z = X'\times_{P^1} Z$ vérifie la condition $(S_2)$ aux points fermés de $X'$ et est régulier en tout autre point, donc est normal ; par suite un revêtement étale de $X'_Z$ dont la restriction à $X'_{\bar{\eta}}$ est connexe est nécessairement connexe. Ceci démontre notre assertion, d'où a). Comme $X'_U$ est le complémentaire dans $S'_U$ d'un diviseur à croisements normaux relativement à $U$ , $h_U$ est localement 1-asphérique pour $L$ (SGA 4 XV 2.1) et, pour tout faisceau de $L$-groupes constant fini $F$ sur $X'_U$ , $(F,h_U)$ est cohomologiquement propre en dimension $\leqslant 1$ (SGA 1 XIII 2.4), d'où la condition b). Comme $P^1$ est normal et $T = P^1 - U$ de
\SGApage{I}{30}
codimension 1 , on a

\begin{equation*}
\text{prof ét}_T(S) \geqslant 2 \qquad ,
\end{equation*}
d'où la condition c). La condition

}{où $h = f|X'$ . On note $L$ l'ensemble des nombres premiers distincts de $p$ et $\bar{\eta}$ un point géométrique au-dessus du point générique $\eta$ de $P^1$ . Vérifions que les hypothèses de SGA 1 XIII 4.8 sont satisfaites. Le morphisme $h$ a une section $\sigma$ car il en est ainsi de $f$ . Prouvons la condition a) de SGA 1 XIII 4.7 ; les fibres de $h$ étant géométriquement intègres, le morphisme $h$ est 0-acyclique et localement 0-acyclique (SGA 4 XV 4.1, 1.16) ; il reste à montrer que, pour tout faisceau d'ensembles localement constant $F$ sur $X'$ , $(F,h)$ est cohomologiquement propre en dimension $\leqslant 0$ . Si $\xi$ est un point fermé de $P^1$ et si $Z$ désigne la fermeture intégrale de Spec $\mathcal{O}_{P^1,\xi}$ dans $\bar{\eta}$ le schéma $X'_Z = X'\times_{P^1} Z$ vérifie la condition $(S_2)$ aux points fermés de $X'_Z$ et est régulier en tout autre point, donc est normal ; par suite un revêtement étale de $X'_Z$ dont la restriction à $X'_{\bar{\eta}}$ est connexe est nécessairement connexe. Ceci démontre notre assertion, d'où a). Comme $X'_U$ est le complémentaire dans $S'_U$ d'un diviseur à croisements normaux relativement à $U$ , $h_U$ est localement 1-asphérique pour $L$ (SGA 4 XV 2.1) et, pour tout faisceau de $L$-groupes constant fini $F$ sur $X'_U$ , $(F,h_U)$ est cohomologiquement propre en dimension $\leqslant 1$ (SGA 1 XIII 2.4), d'où la condition b). Comme $P^1$ est normal et $T = P^1 - U$ de
\SGApage{I}{30}
codimension 1 , on a

\begin{equation*}
\text{prof ét}_T(P^1) \geqslant 2 \qquad ,
\end{equation*}
d'où la condition c). La condition

}{F7I-II-006}\begin{equation*}
\text{prof hop}_x(X') \geqslant 3
\end{equation*}
est valable en tout point fermé de $X'$ d'après le théorème de pureté, d'où la condition e). Enfin, pour tout point $t \in P^1$ , $h$ est lisse en $\sigma(t)$ , donc est localement 1-asphérique pour $L$ en $\sigma(t)$ , d'où la condition d).

D'après SGA 1 XIII 4.7, si $a$ est un point géométrique de $X'_{\bar{\eta}}$ , on a une suite exacte

\SGAvar{\begin{equation*}
1 \longrightarrow \Pi_1^{(p)}(X'_{\bar{\eta}},a) \longrightarrow \Pi'_1(X'_U,a) \longrightarrow \Pi_1(U,a) \longrightarrow 1 \qquad ,
\end{equation*}
et le choix de la section $\sigma$ définit un scindage de cette suite exacte. Notons $K$ l'image de $\Pi_1(U,a)$ dans $\mathrm{Aut}(\Pi_1^{(p)}(X'_{\bar{\eta}},a))$ ; le groupe des coinvariants sous $K$ , $\Pi_1^{(p)}(X'_{\bar{\eta}},a)_K$ , est le quotient de $\Pi_1^{(p)}(X'_{\bar{\eta}},a)$ par le sous-groupe invariant fermé engendré par les éléments du type $x^{-1}qx$ , où $x \in \Pi_1^{(p)}(X'_{\bar{\eta}},a)$ , $q \in K$ . D'après loc. cit. (4.7.4), on a un isomorphisme

\begin{equation*}
\Pi_1^{(p)}(X',a) \simeq \Pi_1^{(p)}(X'_{\bar{\eta}},a) \qquad .
\end{equation*}
Soit $P^1 - U = \coprod_{1\leqslant i\leqslant n} \{t_i\}$ et, pour chaque $i$ , soit $\bar{T}_i$ le localisé strict de $P^1$ en un point géométrique $\bar{t}_i$ au-dessus de $t_i$ et $s_i$ le point générique de $\bar{T}_i$ . Un revêtement étale de $U$ , dont la restriction à chaque $s_i$ est triviale, se prolonge en un revêtement étale de $P^1$ , donc est trivial. Pour chaque $i$ , soit $\bar{k}(s_i)$ la clôture algébrique de $k(s_i)$ ; l'assertion ci-dessus revient à dire que $\Pi_1(U,a)$ est égal au sous-groupe invariant fermé engendré par les images des groupes de Galois, $\mathrm{Gal}(\bar{k}(s_i),k(s_i))$ , dans $\Pi_1(U,a)$ . D'après I 5.3, l'image $K_i$ de $\mathrm{Gal}(\bar{k}(s_i),k(s_i))$ dans $\mathrm{Aut}(\Pi_1^{(p)}(X'_{\bar{\eta}},a))$ est finie. Comme $K$ est le plus petit sous-groupe invariant fermé contenant tous les $K_i$ , et comme $\Pi_1^{(p)}(X'_{\bar{\eta}},a)$ est topologiquement de présentation finie (SGA 1 XIII 3.2), ceci prouve que $\Pi_1^{(p)}(X'_{\bar{\eta}},a)_K$ , donc aussi $\Pi_1^{(p)}(X)$ qui lui est isomorphe, est topologiquement de présentation finie.


}{\begin{equation*}
1 \longrightarrow \Pi_1^{(p)}(X'_{\bar{\eta}},a) \longrightarrow \Pi'_1(X'_U,a) \longrightarrow \Pi_1(U,a) \longrightarrow 1 \qquad ,
\end{equation*}
et le choix de la section $\sigma$ définit un scindage de cette suite exacte. Notons $K$ l'image de $\Pi_1(U,a)$ dans $\mathrm{Aut}(\Pi_1^{(p)}(X'_{\bar{\eta}},a))$ ; le groupe des coinvariants sous $K$ , $\Pi_1^{(p)}(X'_{\bar{\eta}},a)_K$ , est le quotient de $\Pi_1^{(p)}(X'_{\bar{\eta}},a)$ par le sous-groupe invariant fermé engendré par les éléments du type $x^{-1}q(x)$ , où $x \in \Pi_1^{(p)}(X'_{\bar{\eta}},a)$ , $q \in K$ . D'après loc. cit. (4.7.4), on a un isomorphisme

\begin{equation*}
\Pi_1^{(p)}(X',a) \simeq \Pi_1^{(p)}(X'_{\bar{\eta}},a)_K \qquad .
\end{equation*}
Soit $P^1 - U = \coprod_{1\leqslant i\leqslant n} \{t_i\}$ et, pour chaque $i$ , soit $\bar{T}_i$ le localisé strict de $P^1$ en un point géométrique $\bar{t}_i$ au-dessus de $t_i$ et $s_i$ le point générique de $\bar{T}_i$ . Un revêtement étale de $U$ , dont la restriction à chaque $s_i$ est triviale, se prolonge en un revêtement étale de $P^1$ , donc est trivial. Pour chaque $i$ , soit $\bar{k}(s_i)$ la clôture algébrique de $k(s_i)$ ; l'assertion ci-dessus revient à dire que $\Pi_1(U,a)$ est égal au sous-groupe invariant fermé engendré par les images des groupes de Galois, $\mathrm{Gal}(\bar{k}(s_i),k(s_i))$ , dans $\Pi_1(U,a)$ . D'après I 5.3, l'image $K_i$ de $\mathrm{Gal}(\bar{k}(s_i),k(s_i))$ dans $\mathrm{Aut}(\Pi_1^{(p)}(X'_{\bar{\eta}},a))$ est finie. Comme $K$ est le plus petit sous-groupe invariant fermé contenant tous les $K_i$ , et comme $\Pi_1^{(p)}(X'_{\bar{\eta}},a)$ est topologiquement de présentation finie (SGA 1 XIII 3.2), ceci prouve que $\Pi_1^{(p)}(X'_{\bar{\eta}},a)_K$ , donc aussi $\Pi_1^{(p)}(X)$ qui lui est isomorphe, est topologiquement de présentation finie.


}{F7I-II-007}\SGApage{I}{31}

\underline{Remarque} 2.3.2.  La résolution des singularités n'est pas nécessaire pour démontrer le théorème 2.3.1 dans le cas où  $X$  est l'ouvert complémentaire d'un diviseur à croisements normaux d'un schéma projectif, lisse sur  $k$ .

\begin{center}\rule{1.5cm}{0.4pt}\end{center}

\underline{BIBLIOGRAPHIE}

\SGAvar{[1]\quad S. Abhyankar, Local uniformization of algebraic surfaces over ground field of characteristic  $p = 0$ , Amer. J. of Math., t.63, 1956.

}{[1]\quad S. Abhyankar, Local uniformization on algebraic surfaces over ground fields of characteristic  $p \neq 0$ , Ann. of Math., t. 63, no. 3, 1956, p. 491--526.

}{F7I-II-008}[2]\quad H. Hironaka, Resolution of singularities of an algebraic variety over a field of characteristic zero, Ann. of Math., t. 79, 1964.

[3]\quad M. Raynaud, Théorème de Lefschetz en cohomologie des faisceaux cohérents et en cohomologie étale, thèse, à paraître.
\selectlanguage{french}

\ifnum\SGAform=1
\clearpage\section*{Apparat critique}\small
Les lectures ci-dessous sont données dans le TeX exact des témoins, afin de préserver les distinctions de notation. Les indices décimaux renvoient aux différences individuelles du passage marqué.

\par\medskip\noindent\hypertarget{app-F7I-II-001.1}{}\textbf{F7I-II-001.1}\quad p. 26; printed misprint.\par
\noindent Imprimé : {\ttfamily\footnotesize i\allowbreak{}}\par
\noindent Corrigé : {\ttfamily\footnotesize ii\allowbreak{}}\par
\noindent Source : 7I:44. \nolinkurl{fixes/MISPRINTS_7I.json#44}\par

\par\medskip\noindent\hypertarget{app-F7I-II-001.2}{}\textbf{F7I-II-001.2}\quad p. 26; printed misprint.\par
\noindent Imprimé : {\ttfamily\footnotesize ii\allowbreak{}}\par
\noindent Corrigé : {\ttfamily\footnotesize iii\allowbreak{}}\par
\noindent Source : 7I:44. \nolinkurl{fixes/MISPRINTS_7I.json#44}\par

\par\medskip\noindent\hypertarget{app-F7I-II-001.3}{}\textbf{F7I-II-001.3}\quad p. 26; designated-witness attribution.\par
\noindent Imprimé : {\ttfamily\footnotesize a\allowbreak{}}\par
\noindent Corrigé : {\ttfamily\footnotesize b\allowbreak{}}\par
\noindent Source : WITNESS-F7I-II-001.3. \nolinkurl{fr/sga7i/diplomatic/expose_II_body.tex}; \nolinkurl{fr/sga7i/current/expose_II_body.tex}\par

\par\medskip\noindent\hypertarget{app-F7I-II-001.4}{}\textbf{F7I-II-001.4}\quad p. 26; designated-witness attribution.\par
\noindent Imprimé : {\ttfamily\footnotesize b\allowbreak{}}\par
\noindent Corrigé : {\ttfamily\footnotesize c\allowbreak{}}\par
\noindent Source : WITNESS-F7I-II-001.4. \nolinkurl{fr/sga7i/diplomatic/expose_II_body.tex}; \nolinkurl{fr/sga7i/current/expose_II_body.tex}\par

\par\medskip\noindent\hypertarget{app-F7I-II-002.1}{}\textbf{F7I-II-002.1}\quad p. 27; misprint\_pass.\par
\noindent Imprimé : \emph{(vide)}\par
\noindent Corrigé : {\ttfamily\footnotesize \$U\$\allowbreak{} }\par
\noindent Source : 7I:45. \nolinkurl{fixes/FIXLOG_7I_CURRENT.jsonl#620}; \nolinkurl{fixes/MISPRINTS_7I.json#45}\par

\par\medskip\noindent\hypertarget{app-F7I-II-003.1}{}\textbf{F7I-II-003.1}\quad p. 27; printed misprint.\par
\noindent Imprimé : \emph{(vide)}\par
\noindent Corrigé : {\ttfamily\footnotesize \textbackslash{}operatorname\{char\}\allowbreak{}}\par
\noindent Source : 7I:46. \nolinkurl{fixes/MISPRINTS_7I.json#46}\par

\par\medskip\noindent\hypertarget{app-F7I-II-004.1}{}\textbf{F7I-II-004.1}\quad p. 27; printed misprint.\par
\noindent Imprimé : {\ttfamily\footnotesize S\allowbreak{}}\par
\noindent Corrigé : {\ttfamily\footnotesize X\allowbreak{}}\par
\noindent Source : 7I:47. \nolinkurl{fixes/MISPRINTS_7I.json#47}\par

\par\medskip\noindent\hypertarget{app-F7I-II-005.1}{}\textbf{F7I-II-005.1}\quad p. 29; designated-witness attribution.\par
\noindent Imprimé : \emph{(vide)}\par
\noindent Corrigé : {\ttfamily\footnotesize \textbackslash{}left\allowbreak{}}\par
\noindent Source : WITNESS-F7I-II-005.1. \nolinkurl{fr/sga7i/diplomatic/expose_II_body.tex}; \nolinkurl{fr/sga7i/current/expose_II_body.tex}\par

\par\medskip\noindent\hypertarget{app-F7I-II-005.2}{}\textbf{F7I-II-005.2}\quad p. 29; printed misprint.\par
\noindent Imprimé : {\ttfamily\footnotesize X-\allowbreak{} }\par
\noindent Corrigé : \emph{(vide)}\par
\noindent Source : 7I:48. \nolinkurl{fixes/MISPRINTS_7I.json#48}\par

\par\medskip\noindent\hypertarget{app-F7I-II-005.3}{}\textbf{F7I-II-005.3}\quad p. 29; designated-witness attribution.\par
\noindent Imprimé : \emph{(vide)}\par
\noindent Corrigé : {\ttfamily\footnotesize \textbackslash{}right\allowbreak{}}\par
\noindent Source : WITNESS-F7I-II-005.3. \nolinkurl{fr/sga7i/diplomatic/expose_II_body.tex}; \nolinkurl{fr/sga7i/current/expose_II_body.tex}\par

\par\medskip\noindent\hypertarget{app-F7I-II-006.1}{}\textbf{F7I-II-006.1}\quad p. 29; printed misprint.\par
\noindent Imprimé : \emph{(vide)}\par
\noindent Corrigé : {\ttfamily\footnotesize \_Z\allowbreak{}}\par
\noindent Source : 7I:49. \nolinkurl{fixes/MISPRINTS_7I.json#49}\par

\par\medskip\noindent\hypertarget{app-F7I-II-006.2}{}\textbf{F7I-II-006.2}\quad p. 30; printed misprint.\par
\noindent Imprimé : {\ttfamily\footnotesize S\allowbreak{}}\par
\noindent Corrigé : {\ttfamily\footnotesize P\textasciicircum{}1\allowbreak{}}\par
\noindent Source : 7I:50. \nolinkurl{fixes/MISPRINTS_7I.json#50}\par

\par\medskip\noindent\hypertarget{app-F7I-II-007.1}{}\textbf{F7I-II-007.1}\quad p. 30; inherited corrected-layer difference.\par
\noindent Imprimé : {\ttfamily\footnotesize qx\allowbreak{}}\par
\noindent Corrigé : {\ttfamily\footnotesize q(x)\allowbreak{}}\par
\noindent Source : ADOPT-F7I-II-007.1. \nolinkurl{fr/sga7i/diplomatic/expose_II_body.tex}; \nolinkurl{fr/sga7i/current/expose_II_body.tex}\par

\par\medskip\noindent\hypertarget{app-F7I-II-007.2}{}\textbf{F7I-II-007.2}\quad p. 30; printed misprint.\par
\noindent Imprimé : \emph{(vide)}\par
\noindent Corrigé : {\ttfamily\footnotesize \_K\allowbreak{}}\par
\noindent Source : 7I:51. \nolinkurl{fixes/MISPRINTS_7I.json#51}\par

\par\medskip\noindent\hypertarget{app-F7I-II-008.1}{}\textbf{F7I-II-008.1}\quad p. 31; printed misprint.\par
\noindent Imprimé : {\ttfamily\footnotesize of\allowbreak{}}\par
\noindent Corrigé : {\ttfamily\footnotesize on\allowbreak{}}\par
\noindent Source : 7I:52. \nolinkurl{fixes/MISPRINTS_7I.json#52}\par

\par\medskip\noindent\hypertarget{app-F7I-II-008.2}{}\textbf{F7I-II-008.2}\quad p. 31; printed misprint.\par
\noindent Imprimé : {\ttfamily\footnotesize field\allowbreak{}}\par
\noindent Corrigé : {\ttfamily\footnotesize fields\allowbreak{}}\par
\noindent Source : 7I:52. \nolinkurl{fixes/MISPRINTS_7I.json#52}\par

\par\medskip\noindent\hypertarget{app-F7I-II-008.3}{}\textbf{F7I-II-008.3}\quad p. 31; printed misprint.\par
\noindent Imprimé : {\ttfamily\footnotesize =\allowbreak{}}\par
\noindent Corrigé : {\ttfamily\footnotesize \textbackslash{}neq\allowbreak{}}\par
\noindent Source : 7I:52. \nolinkurl{fixes/MISPRINTS_7I.json#52}\par

\par\medskip\noindent\hypertarget{app-F7I-II-008.4}{}\textbf{F7I-II-008.4}\quad p. 31; printed misprint.\par
\noindent Imprimé : {\ttfamily\footnotesize Amer.\allowbreak{} J\allowbreak{}}\par
\noindent Corrigé : {\ttfamily\footnotesize Ann\allowbreak{}}\par
\noindent Source : 7I:52. \nolinkurl{fixes/MISPRINTS_7I.json#52}\par

\par\medskip\noindent\hypertarget{app-F7I-II-008.5}{}\textbf{F7I-II-008.5}\quad p. 31; inherited corrected-layer difference.\par
\noindent Imprimé : \emph{(vide)}\par
\noindent Corrigé : \emph{(1 espace(s), 0 retour(s) de ligne dans le témoin TeX)}\par
\noindent Source : ADOPT-F7I-II-008.5. \nolinkurl{fr/sga7i/diplomatic/expose_II_body.tex}; \nolinkurl{fr/sga7i/current/expose_II_body.tex}\par

\par\medskip\noindent\hypertarget{app-F7I-II-008.6}{}\textbf{F7I-II-008.6}\quad p. 31; inherited corrected-layer difference.\par
\noindent Imprimé : \emph{(vide)}\par
\noindent Corrigé : {\ttfamily\footnotesize no.\allowbreak{} 3,\allowbreak{} }\par
\noindent Source : ADOPT-F7I-II-008.6. \nolinkurl{fr/sga7i/diplomatic/expose_II_body.tex}; \nolinkurl{fr/sga7i/current/expose_II_body.tex}\par

\par\medskip\noindent\hypertarget{app-F7I-II-008.7}{}\textbf{F7I-II-008.7}\quad p. 31; inherited corrected-layer difference.\par
\noindent Imprimé : \emph{(vide)}\par
\noindent Corrigé : {\ttfamily\footnotesize ,\allowbreak{} p.\allowbreak{} 491--526\allowbreak{}}\par
\noindent Source : ADOPT-F7I-II-008.7. \nolinkurl{fr/sga7i/diplomatic/expose_II_body.tex}; \nolinkurl{fr/sga7i/current/expose_II_body.tex}\par

\normalsize\fi
\clearpage\SGAtoc{Exposé VI — FORMAL DEFORMATION THEORY}
\selectlanguage{english}

\SGApage{I}{32}


S G A  7


\underline{Exposé VI}


\underline{\underline{FORMAL DEFORMATION THEORY}}


\underline{by  D. S. RIM}


\SGAsectiontoc{1. Formal existence theorem}{1}
1. \underline{Formal existence theorem}

We recall the notion of cofibered category (S.G.A.1.IV). Let  $\underline{C}$  be a fixed category. A cofibered category  $\underline{A}$  over  $\underline{C}$  is, by definition, a category  $\underline{A}$  together with a functor  $P : \underline{A} \to \underline{C}$  such that

\underline{Ax.1.}  For any map  $f$  in  $\underline{C}$  and any object  $a$  in  $\underline{A}$  with  $P(a) =$ the source of  $f$  , there exists an  $f$-morphism  $\alpha : a \to b$  which is cocartesian (i.e. for any  $f$-morphism  $\alpha' : a \to b'$  there exists a unique  $T$-morphism  $\tau : b \to b'$  in  $\underline{A}$  such that  $\alpha' = \tau\cdot\alpha$  where  $T =$ the target of  $f$ ).

\SGAvar{\underline{Ax.2.}  A composit of cocartesian morphism in  $\underline{A}$  is again cocartesian.

}{\underline{Ax.2.}  A composite of cocartesian morphisms in  $\underline{A}$  is again cocartesian.

}{F7I-VI-001}The following properties of a cofibered category  $P : \underline{A} \to \underline{C}$  follow immediately from the definitions:

(cancellation)  Let  $a \overset{\alpha}{\underset{\beta}{\rightrightarrows}} a'$  be any two cocartesian maps such that  $P(\alpha) = P(\beta)$  . If there exists a cocartesian map  $b \overset{\gamma}{\to} a$  such that  $\alpha\gamma = \beta\gamma$ , then  $\alpha = \beta$  .

(factorization)  Given cocartesian maps  $a \overset{\alpha}{\to} a'$  ,  $a \overset{\beta}{\to} a''$  , there exists a cocartesian map  $a' \overset{\gamma}{\to} a''$  such that  $\gamma\alpha = \beta$  if and only if there exists  $f : P(a') \to P(a'')$  in  $\underline{C}$  such that  $f\cdot P(\alpha) = P(\beta)$  . Furthermore  $\gamma$  is unique.

Let  $\underline{A}$  be a cofibered category over  $\underline{C}$  . For each object  $S$  in  $\underline{C}$  we denote by  $\underline{A}(S)$  the subcategory of  $\underline{A}$  whose objects are the objects  $a$  in  $\underline{A}$  such that  $P(a) = S$  , and  $\mathrm{Hom}_{\underline{A}(S)}(a,a') = \{\alpha \in \mathrm{Hom}_{\underline{A}}(a,a') \mid P(\alpha) = \mathrm{id}_S\}$  . A cofibered groupoid over  $\underline{C}$  is a cofibered category  $\underline{A}$  over  $\underline{C}$  such that  $\underline{A}(S)$
\SGApage{I}{33}
is a groupoid for every object  $S$  in  $\underline{C}$ . We recall that a groupoid is a category in which every morphism is an isomorphism.  Thus, a cofibered groupoid over $\underline{C}$  is a cofibered category  $P : \underline{A} \to \underline{C}$  such that every map in  $\underline{A}$  is cocartesian. Furthermore, for every map  $\alpha \in F\ell(\underline{A})$  ,  $\alpha$  is an isomorphism in  $\underline{A}$  if and only if  $P(\alpha)$  is an isomorphism.

\SGAvar{Throughout the rest of this section, we fix a complete noetherian local ring  $\Lambda$  with the residue field  $k$  , and we set  $\underline{C}_{\Lambda}$  to be the category of artinian local  $\Lambda$-algebras with the same residue field  $k$  , where the maps are local homomorphisms over  $\Lambda$  . From 1.2 on, a cofibered groupoid  $\underline{A}$  over  $\underline{C}_{\Lambda}$ is always assumed to have the property that  $\underline{A}(k) = \{e\}$  , where  $\{e\}$  denotes the category in which  $\mathrm{Hom}_{\{e\}}(a,\ b)$  consists of one and only one element for any  two objects  $a,\ b$  in  $\{e\}$  . A map  $a' \xrightarrow{\ \alpha\ } a$  in  $\underline{A}$  will be called simply a $\underline{C}_{\Lambda}$ -map if  $p(a') = p(a)$  and  $p(\alpha) = 1_{p(a)}$  . A map  $a' \xrightarrow{\ \alpha\ } a$  on  $\underline{A}$  will be called surjective, by abuse of language, if  $p(\alpha) : p(a') \longrightarrow p(a)$  is surjective in  $\underline{C}_{\Lambda}$  .

}{Throughout the rest of this section, we fix a complete noetherian local ring  $\Lambda$  with the residue field  $k$  , and we set  $\underline{C}_{\Lambda}$  to be the category of artinian local  $\Lambda$-algebras with the same residue field  $k$  , where the maps are local homomorphisms over  $\Lambda$  . From 1.2 on, a cofibered groupoid  $\underline{A}$  over  $\underline{C}_{\Lambda}$ is always assumed to have the property that  $\underline{A}(k) = \{e\}$  , where  $\{e\}$  denotes the category in which  $\mathrm{Hom}_{\{e\}}(a,\ b)$  consists of one and only one element for any  two objects  $a,\ b$  in  $\{e\}$  . A map  $a' \xrightarrow{\ \alpha\ } a$  in  $\underline{A}$  will be called simply a $\underline{C}_{\Lambda}$ -map if  $p(a') = p(a)$  and  $p(\alpha) = 1_{p(a)}$  . A map  $a' \xrightarrow{\ \alpha\ } a$  in  $\underline{A}$  will be called surjective, by abuse of language, if  $p(\alpha) : p(a') \longrightarrow p(a)$  is surjective in  $\underline{C}_{\Lambda}$  .

}{F7I-VI-002}\underline{Ex.1.1.(a)}  Let  $F : \underline{C} \longrightarrow$(categories) be a functor. Define the category $\underline{F}$  as follows:  an object in  $\underline{F}$  is a pair  $(S,a)$  where  $S \in \mathrm{ob}(\underline{C})$  and $a \in \mathrm{ob}F(S)$  , and a map  $(S,a) \to (S',a')$  is a pair  $(f,u')$  where  $f : S \to S'$  is a map in  $\underline{C}$  and  $u' : F(f)a \to a'$  is a map in  $F(S')$  . Composition of two maps $(S,a) \xrightarrow{\ (f,u')\ } (S',a') \xrightarrow{\ (g,u'')\ } (S'',a'')$  is given by $(gf,\ u''\cdot(F(g)u')) : (S,a) \to (S'',a'')$  . With respect to the functor  $P : \underline{F} \to \underline{C}$ given by  $(S,a) \mapsto S$  ,  $\underline{F}$  becomes a cofibered category over  $\underline{C}$  . If  $F(S)$  is a groupoid for every  $S \in \mathrm{ob}(\underline{C})$  , then  $\underline{F}$  is a cofibered groupoid over  $\underline{C}$  . In particular, a functor  $F : \underline{C} \to$ (groups) yields a cofibered group, and a functor $F : \underline{C} \to$ (Ens)  yields a cofibered discrete groupoid. We also note that a natural transformation  $\varphi : F \to G$  , where  $F,G : \underline{C} \to$ (categories) are functors, induces a cocartesian functor  $\varphi : \underline{F} \to \underline{G}$  over  $\underline{C}$  .

\underline{Ex.1.1.(b)}  Let  $X$  ,  $Y$  be schemes over  $\Lambda$  together with a fixed
\SGApage{I}{34}
isomorphism  $X \underset{\Lambda}{\otimes} k \xleftrightarrow{\ \tau\ } Y \underset{\Lambda}{\otimes} k$  . Then the functor  $\mathrm{Isom}_{X,Y} : \underline{C}_{\Lambda} \to$ (groupoids) given by  $\mathrm{Isom}_{X,Y}(S) =$ the set of isomorphisms  $X \underset{\Lambda}{\otimes} S \xleftrightarrow{\ \sigma\ } Y \underset{\Lambda}{\otimes} S$  such that $\sigma \underset{\Lambda}{\otimes} k = \tau$  yields a cofibered groupoid  $P : \underline{\mathrm{Isom}}_{X,Y} \to \underline{C}_{\Lambda}$  .

\underline{Ex.1.1.(c)}  Let  $X_0$  be a fixed scheme over  $k$  , and let  $\underline{M}_{X_0}$  be the category consisting of deformations of  $X_0$  over  $\underline{C}_{\Lambda}$  (see §4). Thus an object in $\underline{M}_{X_0}$  is a pair  $(R,X)$  where  $R$  is in  $\underline{C}_{\Lambda}$  and  $X$  is a deformation of  $X_0$  to $R$  . Then the functor  $P : \underline{M}_{X_0} \to \underline{C}_{\Lambda}$  given by  $(R,X) \mathrel{\mapstochar\rightsquigarrow} R$  defines a cofibered groupoid over  $\underline{C}_{\Lambda}$  .

\underline{Ex.1.1.(d)}  For any groupoid  $X$  we denote by  $[X]$  the discrete groupoid consisting of the isomorphism classes of objects in  $X$  . If  $P : \underline{F} \to \underline{C}_{\Lambda}$  is a cofibered groupoid, we obtain a functor  $[\underline{F}] : \underline{C}_{\Lambda} \to$ (Ens)  given by $[\underline{F}](S) = [\underline{F}(S)]$  , and in turn a cofibered discrete groupoid, which is, by abuse of notation, denoted by  $[\underline{F}]$  .

\underline{Definition} 1.2.  \underline{A cofibered groupoid}  $\underline{A}$  \underline{over}  $\underline{C}_{\Lambda}$  \underline{is called semi-homogeneous if the following properties hold}:

(1) \underline{every diagram}

\begin{center}
\begin{tikzcd}
a'' \arrow[r, dashed] \arrow[d, dotted] \& a' \arrow[d] \\
a_1 \arrow[r] \& a
\end{tikzcd}
\end{center}

\underline{of solid arrows in}  $\underline{A}$  \underline{where}  $a' \to a$  \underline{is surjective can be completed into a commutative diagram including dotted arrows}.

(2) \underline{Let}

\begin{center}
\begin{tikzcd}[column sep=large]
\& R \\
R \underset{k}{\times} k[\varepsilon] \arrow[ur, "\pi_1"] \arrow[dr, "\pi_2"'] \& \\
\& k[\varepsilon]
\end{tikzcd}
\end{center}

\underline{be the projection maps, where}  $k[\varepsilon]$  \underline{is the ring of dual numbers over}  $k$  , \underline{and let}

\begin{center}
\begin{tikzcd}[column sep=huge]
a' \arrow[dr, "\alpha"] \& \\
\& a \\
b' \arrow[ur, "\beta"'] \&
\end{tikzcd}
\end{center}

\underline{be}  $\pi_1$\underline{-maps in}  $\underline{A}$  . \underline{Then there exists a}  $\underline{C}_{\Lambda}$  \underline{-map}  $\gamma : a' \longrightarrow b'$  \underline{with}  $\alpha = \beta\cdot\gamma$ \underline{if and only if}  $(\pi_2)_{*}(a') \approx (\pi_2)_{*}(b')$  .

\SGAvar{\underline{Remark} 1.3.(a)  A reformulation of the notion in terms of  2-categories will be given later in 2.6(b). We note that if  $\underline{A}$  is semi-homogeneous. Then so is  $[\underline{A}]$  , and in particular 1.2(2) entails the canonical map
\SGApage{I}{35}
$[\underline{A}(R \times_{k} k[\epsilon])] \longrightarrow [\underline{A}(R)] \times [\underline{A}(k[\epsilon])]$
to be bijective.

}{\underline{Remark} 1.3.(a)  A reformulation of the notion in terms of  2-categories will be given later in 2.6(b). We note that if  $\underline{A}$  is semi-homogeneous, then so is  $[\underline{A}]$  , and in particular 1.2(2) entails the canonical map
\SGApage{I}{35}
$[\underline{A}(R \times_{k} k[\epsilon])] \longrightarrow [\underline{A}(R)] \times [\underline{A}(k[\epsilon])]$
to be bijective.

}{F7I-VI-003}\underline{Remark} 1.3.(b)    Consider a commutative diagram

\begin{center}
\begin{tikzcd}[row sep=large, column sep=large]
 \& a' \arrow[dl, "v_{1}"'] \arrow[dr, "v_{2}"] \& \\
a_{1} \arrow[dr, "u_{1}"'] \& \& a_{2} \arrow[dl, "u_{2}"] \\
 \& a \&
\end{tikzcd}
\end{center}









in a cofibered category  $\underline{A}$  over  $\underline{C}_{\Lambda}$ . Set  $R = P(a)$ , $R' = P(a')$  and
$R_{i} = P(a_{i})$  for  $i = 1, 2$ . If we set  $a''$  to be the  "direct image"  under the
map  $(P(v_{1}),P(v_{2})) : R' \to R_{1} \times_{R} R_{2}$ , we obtain a commutative diagram

\begin{center}
\begin{tikzcd}[row sep=large, column sep=large]
 \& a' \arrow[d] \arrow[ddl, "v_{1}"'] \arrow[ddr, "v_{2}"] \& \\
 \& a'' \arrow[dl] \arrow[dr] \& \\
a_{1} \arrow[dr, "u_{1}"'] \& \& a_{2} \arrow[dl, "u_{2}"] \\
 \& a \&
\end{tikzcd}
\end{center}













Thus 1.2(1) can be replaced by

(1)':  every diagram

\begin{tikzcd}[cramped, column sep=small, row sep=small]
a_{1} \arrow[dr] \& \& a_{2} \arrow[dl] \\
 \& a \&
\end{tikzcd}  in  $\underline{A}$  above  \begin{tikzcd}[cramped, column sep=small, row sep=small]
R_{1} \arrow[dr] \& \& R_{2} \arrow[dl] \\
 \& R \&
\end{tikzcd}  in  $\underline{C}_{\Lambda}$ , where









$R_{2} \to R$  is surjective can be completed to a commutative diagram

\begin{center}
\begin{tikzcd}[row sep=large, column sep=large]
 \& a' \arrow[dl] \arrow[dr] \& \\
a_{1} \arrow[dr] \& \& a_{2} \arrow[dl] \\
 \& a \&
\end{tikzcd}  above  \begin{tikzcd}[row sep=large, column sep=large]
 \& R_{1} \times_{R} R_{2} \arrow[dl, "\pi_{1}"'] \arrow[dr, "\pi_{2}"] \& \\
R_{1} \arrow[dr] \& \& R_{2} \arrow[dl] \\
 \& R \&
\end{tikzcd}
\end{center}












where  $\pi_{i}$  ($i = 1,2$)  are projections.

\SGAvar{\underline{Remark} 1.3.(c)  If  $\underline{A}$  is a semi-homogeneous cofibered groupoid over  $\underline{C}_{\Lambda}$ , then
 $\underline{A}(k [\epsilon])]$ is canonically provided with a vector space structure over  $k$ , where
$k [\epsilon]$ is the ring of dual numbers over  $k$ . Indeed, let  $\underline{\underline{V}}$  be the category of
\SGApage{I}{36}
finite-dimensional vector spaces over  $k$ . Then any functor  $H : \underline{\underline{V}} \to (Ens)$
preserving the product is canonically factored through

}{\underline{Remark} 1.3.(c)  If  $\underline{A}$  is a semi-homogeneous cofibered groupoid over  $\underline{C}_{\Lambda}$ , then
 $[\underline{A}(k [\epsilon])]$ is canonically provided with a vector space structure over  $k$ , where
$k [\epsilon]$ is the ring of dual numbers over  $k$ . Indeed, let  $\underline{\underline{V}}$  be the category of
\SGApage{I}{36}
finite-dimensional vector spaces over  $k$ . Then any functor  $H : \underline{\underline{V}} \to (Ens)$
preserving the product is canonically factored through

}{F7I-VI-004}\begin{center}
\begin{tikzcd}[row sep=large, column sep=large]
\underline{\underline{V}} \arrow[rr, "H"] \arrow[dr, "\overline{H}"'] \& \& (Ens) \\
 \& \underline{\underline{V}} \arrow[ur] \&
\end{tikzcd}
\end{center}

where  $\underline{\underline{V}} \to (Ens)$  is the forgetful functor. On the other hand, the functor
$G : \underline{\underline{V}} \longrightarrow \underline{C}_{\wedge}$  given by  $W \longmapsto k[W] = k \oplus W$ , with  $W^{2} = 0$ , transforms
products in  $\underline{\underline{V}}$  into fibre products over  $k$  in  $\underline{C}_{\wedge}$ . If  $\underline{A}$  is a semi-homogeneous cofibered groupoid over  $\underline{C}_{\wedge}$ , then  $[\underline{A}(k[V] \times_{k} k[W])] \longrightarrow [\underline{A}(k[V])] \times [\underline{A}(k[W])]$
is bijective by 1.3(a), and the functor  $V \longmapsto [\underline{A}(k[V])]$  commutes with products,
hence the result.

\underline{Definition} 1.4. (1) \underline{Let}  $R \in ob(\underline{C}_{\wedge})$ . \underline{A small extension of}  $R$  \underline{in}  $\underline{C}_{\wedge}$  \underline{is a}
\underline{surjective map}  $R' \xrightarrow{\ f\ } R$  \underline{in}  $\underline{C}_{\wedge}$  \underline{such that}  $(\mathrm{Ker}\ f)^{2} = 0$  \underline{and}  $\mathrm{Ker}\ f \simeq k$ .
\underline{In other words,}  $R' \xrightarrow{\ f\ } R$  \underline{in}  $\underline{C}_{\wedge}$  \underline{is a small extension of}  $R$  \underline{if}

$0 \longrightarrow k \xrightarrow{\ t\ } R' \xrightarrow{\ f\ } R \longrightarrow 0$
\underline{is exact for some}  $t \in R'$  \underline{with}  $t^{2} = 0$ .

        (2)  \underline{Let}  $\underline{A}$  \underline{be a cofibered groupoid over}  $\underline{C}_{\wedge}$ , \underline{and let}  $a$  \underline{be an object}
\underline{in}  $\underline{A}$ . \underline{A small prolongation of}  $a$  \underline{is an arrow}  $a' \xrightarrow{\ \alpha\ } a$  \underline{in}  $\underline{A}$  \underline{such that}
$P(\alpha)$  \underline{is a small extension in}  $\underline{C}_{\wedge}$ . \underline{An arrow}  $a' \xrightarrow{\ \alpha\ } a$  \underline{in}  $\underline{A}$  \underline{will be called}
"\underline{versal for small prolongations of}  $a$  "  \underline{if, for every small prolongation}
$a'_{1} \longrightarrow a$  \underline{in}  $\underline{A}$ , \underline{there exists a commutative diagram}

\begin{center}
\begin{tikzcd}[row sep=large, column sep=large]
a' \arrow[rr] \arrow[dr, "\alpha"'] \& \& a'_{1} \arrow[dl] \\
 \& a \&
\end{tikzcd}
\end{center}

\underline{Proposition} 1.5.    \underline{Let}  $\underline{A}$  \underline{be a semi-homogeneous cofibered groupoid over}  $\underline{C}_{\wedge}$ .

(1) \underline{Let}  \begin{tikzcd}[cramped, row sep=tiny, column sep=small]
 \& R \\
R \times_{k} k[\epsilon] \arrow[ur, "\pi_{1}"] \arrow[dr, "\pi_{2}"'] \& \\
 \& k[\epsilon]
\end{tikzcd}  \underline{be projection maps in}  $\underline{C}_{\wedge}$ , \underline{and let}  $a' \xrightarrow{\ \alpha\ } a$








\SGApage{I}{37}
\underline{be a}  $\pi_1$-\underline{map.} \underline{Then there exists a map}  $a \xrightarrow{\ \beta\ } a'$  \underline{such that}  $\alpha\,\beta = 1_a$ \underline{if and only if there exists a map}  $a \longrightarrow (\pi_2)_* a'$  .

(2)  \underline{Let}  $R' \xrightarrow{\ f\ } R$  \underline{be a small extension in}  $\underline{C}_{\Lambda}$  , \underline{and consider a commutative} \underline{diagram}

\SGAvar{\begin{center}
\begin{tikzcd}[row sep=large, column sep=small]
 \& a'' \arrow[dl, "\alpha'"'] \arrow[dr, "\alpha'_1"] \& \& \& \& \& R'\underset{R}{\times}R' \arrow[dl] \arrow[dr] \& \\
a' \arrow[dr, "\alpha"'] \& \& a' \arrow[dl, "\alpha_1"] \& \& \underline{above} \& R' \arrow[dr] \& \& R' \arrow[dl] \\
 \& a \& \& \& \& \& R \&
\end{tikzcd}
\end{center}
















\underline{Define the canonical map}  $\tilde{f} : R' \underset{R}{\times} R' \longrightarrow k[\mathrm{Ker}\ f]$  by  $\tilde{f}(r'_1,\ r'_2) = r'_1(0)+(r'_1-r'_2)$ . \underline{If there exists a map}  $a' \longrightarrow (\tilde{f})_*(a'')$  , \underline{then there exists a map}  $\gamma: a' \longrightarrow a$ \underline{such that}  $\alpha = \alpha_1 \cdot \gamma$  .

\underline{Proof}  (1) The part  "only if"  is clear. Now assume that there exists a map $\varphi : a \longrightarrow (\pi_2)_* a'$  . If we set  $h : R \longrightarrow R \underset{k}{\times} k[\epsilon]$  by  $h = 1 \times P(\varphi)$  , then the composite map  $R \xrightarrow{\ h\ } R \underset{h}{\times} k[\epsilon] \xrightarrow{\ \pi_1\ } R$  is the identity, and therefore we get maps  $a \overset{\gamma}{\underset{\alpha_1}{\rightleftarrows}} f_* a$  such that  $\alpha_1 \cdot \gamma = 1_a$  where  $P(\gamma) = h$  and  $P(\alpha_1) = \pi_1$  . On the other hand,  $\pi_2 \cdot h = P(\varphi)$  and hence there exists a  $\pi_2$ -map $h_* a \longrightarrow (\pi_2)_* a'$  . This means that  $(\pi_2)_*(h_* a) = (\pi_2)_* a'$  , and therefore by 1.2(2) we have a  $\underline{C}_{\Lambda}$ -map  $u : h_* a \longrightarrow a'$  such that  $\alpha \cdot u = \alpha_1$  . Define $\beta : a \longrightarrow a'$  by  $\beta = u \cdot \gamma$  . Then  $\alpha \cdot \beta = \alpha \cdot u \cdot \gamma = \alpha_1 \cdot \gamma = 1_a$  .

(2) The map  $1 \times \tilde{f} : R' \underset{R}{\times} h' \longrightarrow R' \underset{k}{\times} k[\mathrm{Ker}\ f]$  is an isomorphism compatible with the projections on the first factor, and  $\pi_2 \cdot (1 \times \tilde{f}) = \tilde{f}$  where $\pi_2 : R' \underset{k}{\times} k[\mathrm{Ker}\ f] \longrightarrow k[\mathrm{Ker}\ f]$  is the projection map. We note that $k[\mathrm{Ker}\ f] \simeq k[\epsilon]$  since  $R' \xrightarrow{\ f\ } R$  is a small extension. Therefore, if there exists a map  $a' \longrightarrow (\tilde{f})_* a'' = (\pi_2)_* [(1\ \tilde{f})_* a'']$  , then by (1) above there exists a map  $\gamma': a' \longrightarrow a''$  such that  $\alpha' \cdot \gamma' = 1_{a'}$  . Define  $\gamma : a' \longrightarrow a$ by  $\gamma = \alpha'_1 \cdot \gamma'$  . We then have  $\alpha_1 \cdot \gamma = \alpha$  .


}{\begin{center}
\begin{tikzcd}[row sep=large, column sep=small]
 \& a'' \arrow[dl, "\alpha'"'] \arrow[dr, "\alpha'_1"] \& \& \& \& \& R'\underset{R}{\times}R' \arrow[dl] \arrow[dr] \& \\
a' \arrow[dr, "\alpha"'] \& \& a'_1 \arrow[dl, "\alpha_1"] \& \& \underline{above} \& R' \arrow[dr] \& \& R' \arrow[dl] \\
 \& a \& \& \& \& \& R \&
\end{tikzcd}
\end{center}
















\underline{Define the canonical map}  $\tilde{f} : R' \underset{R}{\times} R' \longrightarrow k[\mathrm{Ker}\ f]$  by  $\tilde{f}(r'_1,\ r'_2) = r'_1(0)+(r'_1-r'_2)$ . \underline{If there exists a map}  $a' \longrightarrow (\tilde{f})_*(a'')$  , \underline{then there exists a map}  $\gamma: a' \longrightarrow a'_1$ \underline{such that}  $\alpha = \alpha_1 \cdot \gamma$  .

\underline{Proof}  (1) The part  "only if"  is clear. Now assume that there exists a map $\varphi : a \longrightarrow (\pi_2)_* a'$  . If we set  $h : R \longrightarrow R \underset{k}{\times} k[\epsilon]$  by  $h = 1 \times P(\varphi)$  , then the composite map  $R \xrightarrow{\ h\ } R \underset{k}{\times} k[\epsilon] \xrightarrow{\ \pi_1\ } R$  is the identity, and therefore we get maps  $a \overset{\gamma}{\underset{\alpha_1}{\rightleftarrows}} h_* a$  such that  $\alpha_1 \cdot \gamma = 1_a$  where  $P(\gamma) = h$  and  $P(\alpha_1) = \pi_1$  . On the other hand,  $\pi_2 \cdot h = P(\varphi)$  and hence there exists a  $\pi_2$ -map $h_* a \longrightarrow (\pi_2)_* a'$  . This means that  $(\pi_2)_*(h_* a) = (\pi_2)_* a'$  , and therefore by 1.2(2) we have a  $\underline{C}_{\Lambda}$ -map  $u : h_* a \longrightarrow a'$  such that  $\alpha \cdot u = \alpha_1$  . Define $\beta : a \longrightarrow a'$  by  $\beta = u \cdot \gamma$  . Then  $\alpha \cdot \beta = \alpha \cdot u \cdot \gamma = \alpha_1 \cdot \gamma = 1_a$  .

(2) The map  $1 \times \tilde{f} : R' \underset{R}{\times} R' \longrightarrow R' \underset{k}{\times} k[\mathrm{Ker}\ f]$  is an isomorphism compatible with the projections on the first factor, and  $\pi_2 \cdot (1 \times \tilde{f}) = \tilde{f}$  where $\pi_2 : R' \underset{k}{\times} k[\mathrm{Ker}\ f] \longrightarrow k[\mathrm{Ker}\ f]$  is the projection map. We note that $k[\mathrm{Ker}\ f] \simeq k[\epsilon]$  since  $R' \xrightarrow{\ f\ } R$  is a small extension. Therefore, if there exists a map  $a' \longrightarrow (\tilde{f})_* a'' = (\pi_2)_* [(1 \times \tilde{f})_* a'']$  , then by (1) above there exists a map  $\gamma': a' \longrightarrow a''$  such that  $\alpha' \cdot \gamma' = 1_{a'}$  . Define  $\gamma : a' \longrightarrow a'_1$ by  $\gamma = \alpha'_1 \cdot \gamma'$  . We then have  $\alpha_1 \cdot \gamma = \alpha$  .


}{F7I-VI-005}\SGApage{I}{38}

\underline{Corollary} 1.6.   \underline{Let}  $\underline{A}$  \underline{be a semi-homogeneous cofibered groupoid over}  $\underline{C}_{\Lambda}$ \underline{and consider a diagram}

\begin{center}
\begin{tikzcd}[column sep=small]
a' \arrow[dr, "\alpha"'] \& \& a_1 \arrow[dl, "\alpha_1"] \\
 \& a \&
\end{tikzcd}
\end{center}

\SGAvar{\underline{in}  $\underline{A}$  \underline{where}  $\alpha_1$  \underline{is a small prolongation.} \underline{Assume that}  $a$  \underline{is versal for small} \underline{prolongations of}  $e$  . \underline{Then there exists a map}  $a' \overset{\ \rho\ }{\relbar\joinrel\relbar\joinrel\relbar} a_1$  \underline{such that}  $\alpha = \alpha_1 \rho$ \underline{in}  $\underline{A}$  \underline{if and only if there exists a map}  $f : P(a') \rightarrow P(a_1)$  \underline{such that} $P(\alpha) = P(\alpha_1)\cdot f$  \underline{in}  $\underline{C}_{\Lambda}$  .

}{\underline{in}  $\underline{A}$  \underline{where}  $\alpha_1$  \underline{is a small prolongation.} \underline{Assume that}  $a$  \underline{is versal for small} \underline{prolongations of}  $e$  . \underline{Then there exists a map}  $a' \xrightarrow{\ \rho\ } a_1$  \underline{such that}  $\alpha = \alpha_1 \rho$ \underline{in}  $\underline{A}$  \underline{if and only if there exists a map}  $f : P(a') \rightarrow P(a_1)$  \underline{such that} $P(\alpha) = P(\alpha_1)\cdot f$  \underline{in}  $\underline{C}_{\Lambda}$  .

}{F7I-VI-006}\underline{Proof}:  Assume that  $P(\alpha) = P(\alpha_1)\cdot f$  where  $f : P(a') \rightarrow P(a_1)$  . Replacing  $a'$ by its direct image under  $f$  , we may and shall assume that  $P(\alpha) = P(\alpha_1)$  . Set  $R' = P(a_1)$  ,  $R = P(a)$  . Since  $\underline{A}$  is semi-homogeneous, we get a commutative diagram

\begin{center}
\begin{tikzcd}[row sep=large, column sep=small]
 \& a'' \arrow[dl, "\beta"'] \arrow[dr, "\beta_1"] \& \& \& \& \& R'\underset{R}{\times}R' \arrow[dl] \arrow[dr] \& \\
a' \arrow[dr, "\alpha"'] \& \& a_1 \arrow[dl, "\alpha_1"] \& \& \text{above} \& R' \arrow[dr] \& \& R' \arrow[dl] \\
 \& a \& \& \& \& \& R \&
\end{tikzcd}
\end{center}
















Since  $a$  is versal for small prolongations of  $e$  , so is  $a'$ and therefore our statement follows from 1.5(2).

\SGAvar{\underline{Definition} 1.7.   \underline{For each object}  $R$  \underline{in}  $\underline{C}_{\Lambda}$ \underline{we put}  $\overline{\Omega}_R = M/(m_{\Lambda}R+M^2)$  \underline{where}  $\underline{m}_{\Lambda}$ $M$  \underline{are the maximal ideals of}  $\Lambda$ ,  $R$  \underline{respectively.} \underline{It is a finite-dimensional vector space over}  $k$  , \underline{and}  $R \longmapsto \overline{\Omega}_R$  \underline{defines a functor}  $\overline{\Omega} : \underline{C}_{\Lambda} \rightarrow \underline{\underline{V}}$  \underline{where} $\underline{\underline{V}}$  \underline{is the category of finite-dimensional vector spaces over}  $k$  . \underline{If}  $\underline{A}$  \underline{is a} \underline{cofibered category over}  $\underline{C}_{\Lambda}$  , \underline{then the composite functor}  $\underline{A} \xrightarrow{\ P\ } \underline{C}_{\Lambda} \rightarrow \underline{\underline{V}}$  \underline{is, by} \underline{abuse of notation denoted by}  $\overline{\Omega}$  \underline{again.}

}{\underline{Definition} 1.7.   \underline{For each object}  $R$  \underline{in}  $\underline{C}_{\Lambda}$ \underline{we put}  $\overline{\Omega}_R = M/(m_{\Lambda}R+M^2)$  \underline{where}  $\underline{m}_{\Lambda}$ , $M$  \underline{are the maximal ideals of}  $\Lambda$ ,  $R$  \underline{respectively.} \underline{It is a finite-dimensional vector space over}  $k$  , \underline{and}  $R \longmapsto \overline{\Omega}_R$  \underline{defines a functor}  $\overline{\Omega} : \underline{C}_{\Lambda} \rightarrow \underline{\underline{V}}$  \underline{where} $\underline{\underline{V}}$  \underline{is the category of finite-dimensional vector spaces over}  $k$  . \underline{If}  $\underline{A}$  \underline{is a} \underline{cofibered category over}  $\underline{C}_{\Lambda}$  , \underline{then the composite functor}  $\underline{A} \xrightarrow{\ P\ } \underline{C}_{\Lambda} \rightarrow \underline{\underline{V}}$  \underline{is, by} \underline{abuse of notation denoted by}  $\overline{\Omega}$  \underline{again.}

}{F7I-VI-007}\underline{Proposition} 1.8.   \underline{Let}  $\underline{A}$  \underline{be a semi-homogeneous cofibered groupoid over}  $\underline{C}_{\Lambda}$  . \underline{Then, given an object}  $a$  \underline{in}  $\underline{A}$  \underline{versal for small prolongations of}  $e$  , \underline{there} \underline{exists}  $a' \rightarrow a$  \underline{such that}
\SGApage{I}{39}
(i)  $a'$  \underline{is versal for small prolongations of}  $a$

(ii)  $p(a')$ $\longrightarrow$ $p(a)$  \underline{is surjective; and the kernel is annihilated by the maximal ideal of}  $p(a)$  .

(iii)  $\overline{\Omega}_{a'} \longrightarrow \overline{\Omega}_{a}$  \underline{is an isomorphism}.

\underline{Proof}:  Set  $R = P(a)$  and consider an exact sequence  $0 \to J \to S \to R \to 0$  where $S$  is a formal power-series ring over  $\wedge$  in a finite number of variables, and $J \subset M^2$  where  $M$  is the maximal ideal of  $S$  . Let  $\Sigma$  be the set of ideals  $J_1$ in  $S$  such that (i)  $J \supset J_1 \supset MJ$  and (ii) there exists  $\pi_1$-morphism  $a_1 \to a$  in  $\underline{A}$  where  $\pi_1 : S/J_1 \to R$  is the canonical map. The property 1.2(1) of  $\underline{A}$  together with the fact that  $J/MJ$  is a finite dimensional vector space over  $k$  entails that there exists a smallest ideal in the family  $\Sigma$  . Let it be  $J'$  and choose a  $\pi'$-morphism  $a' \overset{g}{\to} a$  where  $\pi' : S/J' \to R$  is the canonical map. Note that $\overline{\Omega}_{a'} \to \overline{\Omega}_{a}$  is an isomorphism since  $J' \subset J \subset M^2$  . We claim that  $a' \overset{g}{\to} a$  is versal for small prolongations of  $a$  . Indeed, let  $a_1 \to a$  be a small prolongation. If we put  $R_1 = P(a_1)$  , we obtain an exact commutative diagram

\begin{center}
\begin{tikzcd}[column sep=small]
0 \arrow[r] \& J/MJ \arrow[r] \arrow[d, "h"] \& S/MJ \arrow[r] \arrow[d, "h_1"] \& R \arrow[r] \arrow[d, equal] \& 0 \\
0 \arrow[r] \& k \arrow[r] \& R_1 \arrow[r] \& R \arrow[r] \& 0
\end{tikzcd}
\end{center}









\SGAvar{If  $h = 0$  , then  $\pi_1 : R_1 \to R$  admits a cross-section so that  $h_1$  induces $f : S/J'$ $\to R_1$  . If  $h \neq 0$  , then  $h_1 : S/MJ \to R_1$  is surjective and hence  $h_1$ induces  $f : S/J'$ $\to R_1$  by the definiton of the ideal  $J'$  . Therefore in either cases we obtain a commutative diagram

\begin{center}
\begin{tikzcd}
S/J' \arrow[rr, "f", no head] \arrow[dr, "\pi'"'] \& \& R_1 \arrow[dl, "\pi_1"] \\
\& R \&
\end{tikzcd}
\end{center}

}{If  $h = 0$  , then  $\pi_1 : R_1 \to R$  admits a cross-section so that  $h_1$  induces $f : S/J'$ $\to R_1$  . If  $h \neq 0$  , then  $h_1 : S/MJ \to R_1$  is surjective and hence  $h_1$ induces  $f : S/J'$ $\to R_1$  by the definition of the ideal  $J'$  . Therefore in either case we obtain a commutative diagram

\begin{center}
\begin{tikzcd}
S/J' \arrow[rr, "f"] \arrow[dr, "\pi'"'] \& \& R_1 \arrow[dl, "\pi_1"] \\
\& R \&
\end{tikzcd}
\end{center}

}{F7I-VI-008}Since  $a$  is versal for small prolongations of  $e$  , it follows from 1.6 that we obtain a commutative diagram
\SGApage{I}{40}
\begin{tikzcd}
a' \arrow[rr, "\rho"] \arrow[dr] \& \& a_1 \arrow[dl] \\
\& a \&
\end{tikzcd}

Let  $P : \underline{A} \to \underline{C}_\wedge$  be a cofibered category (in groupoids). It is then clear that we obtain a cofibered category  $\mathrm{pro}(P) : \mathrm{pro}(\underline{A}) \to \mathrm{pro}(\underline{C}_\wedge)$  (in groupoids) (pour la définition de la catégorie de pro-objets  $\mathrm{pro}(\underline{A})$  , voir A. Grothendieck, Sém. Bourbaki 195).

\begin{center}
\begin{tikzcd}
\underline{A} \arrow[r] \arrow[d, "P"'] \& \mathrm{pro}(\underline{A}) \arrow[d, "\mathbf{pro}(P)"] \\
\underline{C}_\wedge \arrow[r] \& \mathrm{pro}(\underline{C}_\wedge)
\end{tikzcd}
\end{center}









Now let  $\widehat{\underline{C}}_\wedge$  be the full subcategory of  pro-$\underline{C}_\wedge$  in which $\mathrm{ob}(\widehat{\underline{C}}_\wedge) = \{(R_i)_{i \in \underline{N}}$  ,  $R_i \in \mathrm{ob}(\underline{C}_\wedge) \mid \varphi_i : R_i \to R_{i-1}$  is surjective for all  $i$  and induces isomorphisms  $\underline{m}_i/\underline{m}_\wedge + \underline{m}_i^2 \to \underline{m}_{i-1}/\underline{m}_\wedge + \underline{m}_{i-1}^2$  for all sufficiently large  $i$  $\}$ where  $\underline{N}$  stands for the ordered category of natural numbers, and  $\underline{m}_i$  is the maximal ideal of  $R_i$  . We set  $\widehat{\underline{A}}$  to be the full subcategory of pro-$\underline{A}$  in which

\begin{equation*}
\mathrm{ob}(\widehat{\underline{A}}) = \{a \in \mathrm{ob}(\mathrm{pro}(\underline{A})) \mid \mathrm{pro}(P)\ (a) \in \mathrm{ob}(\widehat{\underline{C}}_\wedge)\}
\end{equation*}
We then obtain a cofibered category  $\widehat{P} : \widehat{\underline{A}} \to \widehat{\underline{C}}_\wedge$  (in groupoids) with a commutative diagram

\begin{center}
\begin{tikzcd}
\underline{A} \arrow[r, hook] \arrow[d, "P"'] \& \widehat{\underline{A}} \arrow[d, "\widehat{P}"] \\
\underline{C}_\wedge \arrow[r, hook] \& \widehat{\underline{C}}_\wedge
\end{tikzcd}
\end{center}










We note that the category  $\widehat{\underline{C}}_\wedge$  is canonically equivalent to the category of complete local $\wedge$-algebras of formally finite type over  $\wedge$  with the fixed residue field  $k$  . Henceforth we shall make such identification.

\underline{Definition} 1.9.   \underline{Let}  $\underline{A}$  \underline{be a cofibered groupoid over}  $\underline{C}_\wedge$  . \underline{An object}  $\xi$  \underline{in} $\widehat{\underline{A}}$  \underline{is called a quasi-initial object of}  $\widehat{\underline{A}}$  \underline{if for every object}  $a$  \underline{in}  $\widehat{\underline{A}}$  \underline{there exists a map}  $\xi \to a$  \underline{in}  $\widehat{\underline{A}}$  . \underline{A quasi-initial object}  $\xi$  \underline{of}  $\widehat{\underline{A}}$  \underline{is called minimal}
\SGApage{I}{41}
\underline{if the morphism}  $\xi : h_R \to [\hat{\underline{A}}]$  \underline{induces a bijection} $\xi(k[\epsilon])$: $\mathrm{Hom}_{\wedge\text{-alg}}(R,k[\epsilon]) \to [\underline{A}(k[\epsilon])]$ , \underline{where} $R = P(\xi)$ . \underline{An object} $\xi$ \underline{in} $\hat{\underline{A}}$ \underline{is called a versal formal object of} $\underline{A}$ \underline{if every surjective map} $\alpha : a' \to a$ \underline{in} $\underline{A}$ \underline{induces a surjective map} $\mathrm{Hom}_{\hat{\underline{A}}}(\xi,a') \to \mathrm{Hom}_{\hat{\underline{A}}}(\xi,a)$ .

\underline{A versal formal object} $\xi$ \underline{of} $\underline{A}$ \underline{is called minimal if the morphism} $\xi : h_R \to [\hat{\underline{A}}]$ \underline{induces a bijection} $\xi(k[\epsilon]) : \mathrm{Hom}_{\wedge\text{-alg}}(R,k[\epsilon]) \to [\underline{A}(k[\epsilon])]$ \underline{where} $R = P(\xi)$ .

\underline{Proposition} 1.10.  \underline{Let} $\underline{A}$ \underline{be a cofibered groupoid over} $\underline{C}_{\wedge}$ .

  (i)  \underline{every (minimally) versal formal object of} $\underline{A}$ \underline{is a quasi-initial (minimal) object of} $\hat{\underline{A}}$ .

\SGAvar{  (ii) \underline{If} $\xi$ \underline{is an object in} $\hat{A}$ \underline{such that} $\widetilde{\xi}(k[\epsilon])$: $\mathrm{Hom}_{\wedge\text{-alg}}(R,k[\epsilon])\to [\underline{A}(k[\epsilon])]$ \underline{is injective, where} $P(\xi) = R$ , \underline{then every endomorphism of} $\xi$ \underline{in} $\hat{A}$ \underline{is an automorphism. In particular a quasi-initial minimal object of} $\hat{\underline{A}}$ , \underline{if it exists, is unique up to (non-canonical) isomorphism.}

\underline{Proof}: (i) let $\xi$ be a versal formal object of $\underline{A}$ . We must show that, for any object $a$ in $\hat{\underline{A}}$ , we have a map $\xi \to a$ . By the definition of $\hat{\underline{A}}$ , there exists a sequence of surjective maps in $\underline{A}$ : $--- \to a_n \xrightarrow{\ \alpha_n\ } a_{n-1} \to --- \to a_1 \overset{\alpha_1}{\relbar\joinrel\relbar} a_0$ such that $a = \varprojlim_n a_n$ . Assume that there exists maps $\varphi_i : \xi \to a_i$ such that $\alpha_i\varphi_i = \varphi_{i-1}$ for all $i < n$ . Since $\xi$ is a versal formal object of $\underline{A}$ and $\alpha_n : a_n \to a_{n-1}$ is surjective in $\underline{A}$ , we obtain a map $\xi \xrightarrow{\ \varphi_n\ } a_n$ such that $\alpha_n\varphi_n = \varphi_{n-1}$ . If we set $\varphi = \varprojlim_n \varphi_n$ , we have $\varphi : \xi \to a$ .

}{  (ii) \underline{If} $\xi$ \underline{is an object in} $\hat{\underline{A}}$ \underline{such that} $\widetilde{\xi}(k[\epsilon])$: $\mathrm{Hom}_{\wedge\text{-alg}}(R,k[\epsilon])\to [\underline{A}(k[\epsilon])]$ \underline{is injective, where} $P(\xi) = R$ , \underline{then every endomorphism of} $\xi$ \underline{in} $\hat{\underline{A}}$ \underline{is an automorphism. In particular a quasi-initial minimal object of} $\hat{\underline{A}}$ , \underline{if it exists, is unique up to (non-canonical) isomorphism.}

\underline{Proof}: (i) let $\xi$ be a versal formal object of $\underline{A}$ . We must show that, for any object $a$ in $\hat{\underline{A}}$ , we have a map $\xi \to a$ . By the definition of $\hat{\underline{A}}$ , there exists a sequence of surjective maps in $\underline{A}$ : $--- \to a_n \xrightarrow{\ \alpha_n\ } a_{n-1} \to --- \to a_1 \overset{\alpha_1}{\relbar\joinrel\relbar} a_0$ such that $a = \varprojlim_n a_n$ . Assume that there exist maps $\varphi_i : \xi \to a_i$ such that $\alpha_i\varphi_i = \varphi_{i-1}$ for all $i < n$ . Since $\xi$ is a versal formal object of $\underline{A}$ and $\alpha_n : a_n \to a_{n-1}$ is surjective in $\underline{A}$ , we obtain a map $\xi \xrightarrow{\ \varphi_n\ } a_n$ such that $\alpha_n\varphi_n = \varphi_{n-1}$ . If we set $\varphi = \varprojlim_n \varphi_n$ , we have $\varphi : \xi \to a$ .

}{F7I-VI-009}(ii) Let $\xi \xrightarrow{\ \alpha\ } \xi$ be a map in $\hat{A}$ . Then for every $\lambda \in \mathrm{Hom}_{\wedge\text{-alg}}(R,k[\epsilon])$ we have $\lambda_*\xi = \lambda_*\alpha\xi \simeq (\lambda\cdot P(\alpha))_*\xi$ and hence by the hypothesis on $\xi$ we have that $\lambda = \lambda\cdot P(\alpha)$ for all $\lambda \in \mathrm{Hom}_{\wedge\text{-alg}}(R,k[\epsilon])$ , and therefore $P(\alpha) : R \to R$ is surjective. However it is trivial to see that every surjective endomorphism of a noetherian ring is always an automorphism and hence $P(\alpha)$ is an automorphism of $R$ and therefore $\alpha$ is an automorphism of $\xi$ .

\underline{Theorem} 1.11. (Schlessinger) \underline{Let} $\underline{A}$ \underline{be a cofibered groupoid over} $\underline{C}_{\Lambda}$ . \underline{Then} $\underline{A}$
\SGApage{I}{42}
\underline{admits a minimally-versal formal object if and only if}

  (i)  $\underline{A}$  \underline{is semi-homogeneous}

  (ii)  $[\underline{A}(k[\epsilon])]$  \underline{is a finite-dimensional vector space over}  $k$  (cf. 1.3(c)).

\underline{Proof}; (\underline{Necessity}) Let $\xi$ in $\hat{\underline{A}}$ be a minimally-versal formal object of $\underline{A}$ , and consider a diagram

\begin{center}
\begin{tikzcd}[column sep=small, row sep=large]
 a_1 \arrow[dr] \& \& a_2 \arrow[dl, "\alpha"] \\
 \& a \&
\end{tikzcd}
\end{center}

\SGAvar{in $\underline{A}$ where $\alpha$ is surjective. Since $\xi$ is a quasi-initial object of $\hat{\underline{A}}$ , there exists a map $\xi \to a_1$ . Since $\alpha$ is surjective, the composit map $\xi \to a_1 \to a$ can be lifted to $a_2$ i.e., we obtain a commutative diagram

}{in $\underline{A}$ where $\alpha$ is surjective. Since $\xi$ is a quasi-initial object of $\hat{\underline{A}}$ , there exists a map $\xi \to a_1$ . Since $\alpha$ is surjective, the composite map $\xi \to a_1 \to a$ can be lifted to $a_2$ i.e., we obtain a commutative diagram

}{F7I-VI-010}\begin{center}
\begin{tikzcd}[column sep=small, row sep=large]
 \& \xi \arrow[dl] \arrow[dr] \& \\
 a_1 \arrow[dr] \& \& a_2 \arrow[dl, "\alpha"] \\
 \& a \&
\end{tikzcd}
\end{center}

which proves the condition 1.2(1). We now prove the condition 1.2(2). Let

\SGAvar{\begin{center}
\begin{tikzcd}[column sep=large]
 \& S \\
 S{\times}k[\epsilon] \arrow[ur, "\pi_1"] \arrow[dr, "\pi_2"'] \& \\
 \& k[\epsilon]
\end{tikzcd}
\end{center}

}{\begin{center}
\begin{tikzcd}[column sep=large]
 \& S \\
 S{\times}_k k[\epsilon] \arrow[ur, "\pi_1"] \arrow[dr, "\pi_2"'] \& \\
 \& k[\epsilon]
\end{tikzcd}
\end{center}

}{F7I-VI-011}be the projections,

and let

\begin{center}
\begin{tikzcd}[column sep=small, row sep=large]
 a' \arrow[dr, "\alpha"'] \& \& b' \arrow[dl, "\beta"] \\
 \& a \&
\end{tikzcd}
\end{center}

be $\pi_1$-maps in $\underline{A}$ . By the definition of $\xi$ , we obtain a commutative diagram

\begin{center}
\begin{tikzcd}[column sep=small, row sep=large]
 \& \xi \arrow[dl, "u"'] \arrow[dr, "v"] \& \\
 a' \arrow[dr, "\alpha"'] \& \& b' \arrow[dl, "\beta"] \\
 \& a \&
\end{tikzcd}
\end{center}

\SGAvar{and in particular we have $\pi_1\cdot P(u) = \pi_1\cdot P(v)$ . Assume now that $(\pi_2)_*a' \approx (\pi_2)_*b'$ . We then have $\pi_2\cdot P(u) = \pi_2\cdot P(v)$ because of the minimality of
\SGApage{I}{43}
$\xi$ , and hence  $P(u) = \pi_1 P(u)\times \pi_2 P(u) = \pi_1 P(v)\times \pi_2 P(v) = P(v)$  . Since  $\underline{A}$  is a cofibered groupoid, we obtain a unique isomorphism  $a' \overset{\sigma}{\to} b'$  in $\underline{A}(S\times k[\epsilon])$  such that  $\sigma u = v$  . Then  $\beta\sigma u = \alpha u$  entails that  $\beta\sigma= \alpha$  since  $\underline{A}$ is a cofibered category.

(\underline{Sufficiency}) Let  $c_1,c_2,\dots,c_n$  be objects in  $\underline{A}(k[\epsilon])$  which form a $k$-basis of the vector space  $[\underline{A}(k[\epsilon])]$  , and choose  $a_1$  in  $\underline{A}(k[\epsilon]\times_k\dots\times_k k[\epsilon])$ such that  $(\pi_j)_* a_1 \approx c_j$  . Then  $\tilde{a}(k[\epsilon]) : \mathrm{Hom}_{\Lambda\text{-alg}}(R_1,k[\epsilon]) \to [\underline{A}(k[\epsilon])]$  is bijective, where  $R_1 = P(a_1)$  . In particular  $a_1 \to e$  is versal for small prolongations of  $e$  . Choose  $a_2 \to a_1$  which is versal for small prolongations of  $a_1$ such that  $\overline{\Omega}_{a_2} \to \overline{\Omega}_{a_1}$  is bijective. We then choose  $a_3 \to a_2$  , versal for small prolongations of  $a_2$  such that  $\overline{\Omega}_{a_3} \to \overline{\Omega}_{a_2}$  , is bijective etc., and set $\xi = \underset{n}{\underleftarrow{\mathrm{Lim}}}\, a_n$  . We claim that  $\xi$  is a minimally-versal formal object of  $\underline{A}$  . Indeed, since  $\overline{\Omega}_{a_n} \to \overline{\Omega}_{a_{n-1}}$  is bijective, for all  $n$  ,  $\overline{\Omega}_{\xi} \to \overline{\Omega}_{a_1}$  is bijective so that $\xi(k[\epsilon]) : \mathrm{Hom}_{\Lambda\text{-alg}}(R,k[\epsilon] \to [\underline{A}(k[\epsilon])]$  is bijective. Now consider a diagram

}{and in particular we have $\pi_1\cdot P(u) = \pi_1\cdot P(v)$ . Assume now that $(\pi_2)_*a' \approx (\pi_2)_*b'$ . We then have $\pi_2\cdot P(u) = \pi_2\cdot P(v)$ because of the minimality of
\SGApage{I}{43}
$\xi$ , and hence  $P(u) = \pi_1 P(u)\times \pi_2 P(u) = \pi_1 P(v)\times \pi_2 P(v) = P(v)$  . Since  $\underline{A}$  is a cofibered groupoid, we obtain a unique isomorphism  $a' \overset{\sigma}{\to} b'$  in $\underline{A}(S\times_k k[\epsilon])$  such that  $\sigma u = v$  . Then  $\beta\sigma u = \alpha u$  entails that  $\beta\sigma= \alpha$  since  $\underline{A}$ is a cofibered category.

(\underline{Sufficiency}) Let  $c_1,c_2,\dots,c_n$  be objects in  $\underline{A}(k[\epsilon])$  which form a $k$-basis of the vector space  $[\underline{A}(k[\epsilon])]$  , and choose  $a_1$  in  $\underline{A}(k[\epsilon]\times_k\dots\times_k k[\epsilon])$ such that  $(\pi_j)_* a_1 \approx c_j$  . Then  $\tilde{a}_1(k[\epsilon]) : \mathrm{Hom}_{\Lambda\text{-alg}}(R_1,k[\epsilon]) \to [\underline{A}(k[\epsilon])]$  is bijective, where  $R_1 = P(a_1)$  . In particular  $a_1 \to e$  is versal for small prolongations of  $e$  . Choose  $a_2 \to a_1$  which is versal for small prolongations of  $a_1$ such that  $\overline{\Omega}_{a_2} \to \overline{\Omega}_{a_1}$  is bijective. We then choose  $a_3 \to a_2$  , versal for small prolongations of  $a_2$  such that  $\overline{\Omega}_{a_3} \to \overline{\Omega}_{a_2}$  is bijective etc., and set $\xi = \underset{n}{\underleftarrow{\mathrm{Lim}}}\, a_n$  . We claim that  $\xi$  is a minimally-versal formal object of  $\underline{A}$  . Indeed, since  $\overline{\Omega}_{a_n} \to \overline{\Omega}_{a_{n-1}}$  is bijective, for all  $n$  ,  $\overline{\Omega}_{\xi} \to \overline{\Omega}_{a_1}$  is bijective so that $\xi(k[\epsilon]) : \mathrm{Hom}_{\Lambda\text{-alg}}(R,k[\epsilon]) \to [\underline{A}(k[\epsilon])]$  is bijective. Now consider a diagram

}{F7I-VI-012}\begin{center}
\begin{tikzcd}
 \& a' \arrow[d, "\alpha"] \\
\xi \arrow[r, "h"] \& a
\end{tikzcd}
\end{center}

\SGAvar{where  $\alpha$  is surjective in  $\underline{A}$  . We must show that  $h$  can be lifted to  $a'$  . For this, one may assume, by induction, that  $a'$  is a small prolongation of  $a$  . Since  $P(a)$  is an artinian ring,  $h : \xi \to a$  is factored as a composit map $\xi \to a_q \longrightarrow a$  for some  $q$  . Since  $\underline{A}$  is semi-homogeneous, we get a diagram

}{where  $\alpha$  is surjective in  $\underline{A}$  . We must show that  $h$  can be lifted to  $a'$  . For this, one may assume, by induction, that  $a'$  is a small prolongation of  $a$  . Since  $P(a)$  is an artinian ring,  $h : \xi \to a$  is factored as a composite map $\xi \to a_q \longrightarrow a$  for some  $q$  . Since  $\underline{A}$  is semi-homogeneous, we get a diagram

}{F7I-VI-013}\begin{center}
\begin{tikzcd}
 \& a'' \arrow[r] \arrow[d, "\alpha'"] \& a' \arrow[d, "\alpha"] \\
\xi \arrow[r] \& a_q \arrow[r] \& a
\end{tikzcd}
\end{center}

where  $\alpha'$  is also a small prolongation. We then obtain a commutative diagram

\begin{center}
\begin{tikzcd}
 \& \& a'' \arrow[d, "\alpha'"] \\
\xi \arrow[r] \& a_{q+1} \arrow[r] \arrow[ur] \& a_q
\end{tikzcd}
\end{center}

and therefore we are done.

\underline{Remark} 1.12.  Let
\SGApage{I}{44}
\begin{tikzcd}
\dots\dots \arrow[r] \& a_n \arrow[r] \& a_{n-1} \arrow[r] \& \dots \arrow[r] \& a_2 \arrow[r] \& a_1 \arrow[r] \& a_o = e
\end{tikzcd}

\SGAvar{be an infinite sequences of arrows in  $\underline{A}$  such that

}{be an infinite sequence of arrows in  $\underline{A}$  such that

}{F7I-VI-014}\begin{equation*}
\left\{\begin{array}{l}
a_i \ \text{ is versal for small prolongations of } \ a_{i-1} \\[2pt]
P(a_i) \longrightarrow P(a_{i-1}) \ \text{ is surjective} \\[2pt]
\overline{\Omega}_{a_i} \dashrightarrow \overline{\Omega}_{a_{i-1}} \ \text{ is bijective for } \ i >> 0
\end{array}\right.
\end{equation*}
Then the above proof shows that  $\underleftarrow{\mathrm{Lim}}\, a_i$  is a versal formal object of  $\underline{A}$  .

\underline{Remark} 1.13.  Let  $\underline{A}$  be a cofibered groupoid over  $\underline{C}_\wedge$  . If  $\underline{A}$  has the property 1.11(i) and (ii), then so does  $[\underline{A}]$  , and  $\xi \in \mathrm{ob}(\hat{\underline{A}})$  is a minimally-versal formal object of  $\underline{A}$  if and only if it is so for  $[\underline{A}]$  by 1.10(ii).

\underline{Remark} 1.14.  Let  $\underline{A}$  be a cofibered groupoid over  $\underline{C}_\wedge$  , admitting a minimally-versal formal object  $\xi$  . Then an invariant of  $R = P(\xi)$  , which is a complete local  $\wedge$-algebra of formally finite type over  $\wedge$  , is automatically an invariant of  $\underline{A}$  over  $\underline{C}_\wedge$  . Therefore, for example, we shall say that  $\underline{A}$  is (formally) smooth over  $\underline{C}_\wedge$  if  $R$  is (formally) smooth over  $\wedge$  i.e., a formal power-series algebra over  $\wedge$  , or we say that  $\underline{A}$  is irreducible over  $\underline{C}_\wedge$  if  $\mathrm{Spec}(R)$  is irreducible etc. In the next section we introduce the notion of  "smoothness"  of a functor  $\underline{A} \to \underline{B}$  over  $\underline{C}_\wedge$  , where  $\underline{A}, \underline{B}$  are cofibered categories in groupoids over  $\underline{C}_\wedge$  . It turns out (cf. 2.4) that  "$\underline{A}$  is (formally) smooth over  $\underline{C}_\wedge$"  is equivalent to  " $P : \underline{A} \to \underline{C}_\wedge$  is a smooth functor"  , and therefore no confusion arises. We also define the dimension of  $\underline{A}$  over  $\underline{C}_\wedge$  by  $\dim_{\underline{C}_\wedge} \underline{A} = \dim\, k \otimes_\wedge R$  . For this we note the following fact:  Let  $\underline{C}_k$  be the category of artinian local $k$-algebras with the same residue field  $k$  . Then  $\underline{C}_k \to \underline{C}_\wedge$  is a fully-faithful imbedding, whose (left) adjoint functor  $\underline{C}_\wedge \to \underline{C}_k$  is given by  $R \to k \otimes_\wedge R$  . Let $\underline{A}|\underline{C}_k$  be the pull-back of  $P : \underline{A} \to \underline{C}_\wedge$  under  $\underline{C}_k \to \underline{C}_\wedge$  . Then it is trivial to see that if  $\xi$  is a minimally-versal formal object for  $\underline{A}$  over  $\underline{C}_\wedge$  , then  $k \otimes_\wedge \xi$ is a minimally-versal formal object for  $\underline{A}|\underline{C}_k$  over  $\underline{C}_k$  , where  $\xi \to k \otimes_\wedge \xi$  is a cocartesian arrow above  $R \to k \otimes_\wedge R$  . It follows that  $\dim_{\underline{C}_\wedge} \underline{A} = \dim_{\underline{C}_k} (\underline{A}|\underline{C}_k)$  .


\SGApage{I}{45}

\SGAvar{1.15.  In the remaining of this paragraph, we biefly explain how to modify the preceeding theory in order to allow for residue field extensions.

}{1.15.  In the remaining of this paragraph, we briefly explain how to modify the preceding theory in order to allow for residue field extensions.

}{F7I-VI-015}Let $\Lambda$ and $K$ be complete noetherian local rings, with $K$ a $\Lambda$-algebra. We assume that

(a)  the natural map from $\Lambda$ to $K$ is a local homomorphism;

(b)  the residue field $K'$ of $K$ is an extension of finite type of the residue field $k$ of $\Lambda$ .

The most interesting case is when $K = K'$ and $K$ is a finite inseparable extension of $k$ .

We denote by $\hat{C}_{\Lambda,K}$ the category whose objects are the complete noetherian local rings $R$ , provided with a structure of $\Lambda$-algebra and with a $\Lambda$-epimorphism $s : R \longrightarrow K$ .

We denote by $C_{\Lambda,K}$ the subcategory of $\hat{C}_{\Lambda,K}$ consisting of the $(R, s)$ such that $\mathrm{Ker}(s)$ is a $R$-module of finite length. The category $\hat{C}_{\Lambda,K}$ can be identified with a full subcategory of $\mathrm{Pro}(C_{\Lambda,K})$ .

If $K$ is artinian, so is any object in $C_{\Lambda,K}$ .

For any finite dimensional vector space $V$ over $K'$ we denote by $D_V(K)$ the algebra of dual numbers over $K$ defined by $V$ ; its elements are of the form $k + v\,\varepsilon$ (with $k \in K$ , $v \in V$ and $\varepsilon^2 = 0$ : $(k + v\,\varepsilon)\,(k' + v'\,\varepsilon) = (k\,k') + (k\,v' + k'\,v)\varepsilon$  ). The map $s : D_V(K) \to K$ : $s(k + v\,\varepsilon) = k$ , turns $D_V(K)$ into an object of $C_{\Lambda,K}$ . The role played before by $k[\varepsilon]$ will be played here by $D_{K'}(K)$ .

\SGAvar{As explained in 1.8, any cofibered groupoid $\underline{A}$ over $C_{\Lambda,K}$ extends in a natural way to a cofibered category $\hat{\underline{A}}$ over $\hat{C}_{\Lambda,K}$ . A map in $\underline{A}$ or $\hat{\underline{A}}$ will be called \underline{surjective} if the corresponding homomorphism of algebras in $C_{\Lambda,k}$ or in $\hat{C}_{\Lambda,K}$ is.

}{As explained in 1.8, any cofibered groupoid $\underline{A}$ over $C_{\Lambda,K}$ extends in a natural way to a cofibered category $\hat{\underline{A}}$ over $\hat{C}_{\Lambda,K}$ . A map in $\underline{A}$ or $\hat{\underline{A}}$ will be called \underline{surjective} if the corresponding homomorphism of algebras in $C_{\Lambda,K}$ or in $\hat{C}_{\Lambda,K}$ is.

}{F7I-VI-016}\SGApage{I}{46}

\underline{We assume that} $\underline{A}(K) = \{e\}$ .

\underline{Definition} 1.16.  \underline{A cofibered groupoïd} $\underline{A}$ \underline{over} $C_{\Lambda,K}$ \underline{is called} \underline{semi-homogeneous} \underline{if the following properties hold:}

(I)  \underline{every diagram}

\begin{center}
\begin{tikzcd}[row sep=large, column sep=large]
a'' \arrow[r, dashed] \arrow[d, dashed] \& a' \arrow[d] \\
a_1 \arrow[r] \& a
\end{tikzcd}
\end{center}











\underline{of solid arrows in} $\underline{A}$ , \underline{where} $a' \longrightarrow a$ \underline{is surjective, can be completed into} \underline{a commutative diagram including dotted arrows;}

(II)  \underline{For any} $R$ , \underline{the natural map}

\begin{equation*}
[\underline{A}(R \times_K D_{K'}(K))] \to [\underline{A}(R) \times \underline{A}\,(D_{K'}(K))] \simeq [\underline{A}(R)] \times [\underline{A}(D_{K'}(K))]
\end{equation*}
\underline{is bijective.}

If $\underline{A}$ is semi-homogeneous, so is $[\underline{A}]$ .

1.17.  Let us first consider the restriction of $\underline{A}$ to the full subcategory of $C_{\Lambda,K}$ whose objects are the $D_V(K)$ . Let $\Omega$ be the vector space $\Omega^1_{K/\Lambda} \otimes_K K'$ , and $d : K \to \Omega$ the derivative. A map $f : D_V(K) \to D_W(K)$ can be written:

\SGAvar{\begin{equation*}
f(k + v\,\varepsilon) = k + (u(v) + D\,d\,k)
\end{equation*}
with $u \in \mathrm{Hom}_{K'}(V, W)$ and $D \in \mathrm{Hom}_{K'}(\Omega, W)$ . If we identify $f$ with $(u, D)$ , the composition law is

}{\begin{equation*}
f(k + v\,\varepsilon) = k + (u(v) + D\,d\,k)\,\varepsilon
\end{equation*}
with $u \in \mathrm{Hom}_{K'}(V, W)$ and $D \in \mathrm{Hom}_{K'}(\Omega, W)$ . If we identify $f$ with $(u, D)$ , the composition law is

}{F7I-VI-017}\begin{equation*}
(u, D) \circ (u', D') = (uu', D + u\,D')\ .
\end{equation*}
The condition (II) entails

(II')  the functor from $K'$-vector spaces to sets

\begin{equation*}
V \longmapsto [\underline{A}(D_V(K))]
\end{equation*}
commutes with products.

\underline{Lemma} 1.18.  \underline{Assume condition} (II') \underline{is satisfied by} $\underline{A}$ . \underline{Then}
\SGApage{I}{47}
(1)  $[\underline{A}(D_{K'}(K))]$  \underline{carries a unique vector space structure, such that}  $[\underline{A}(D_{V}(K))] \simeq V \otimes_{K'} [\underline{A}(D_{K'}(K))]$  \underline{as a functor in}  $V$ .

(2)   \underline{There is a linear map}  $\gamma : \operatorname{Hom}(\Omega,K') \rightarrow [\underline{A}(D_{K'}(K))]$  \underline{such that, for any}  $D \in \operatorname{Hom}(\Omega, K')$  ,

\SGAvar{\begin{equation*}
[\underline{A}((k, D))] : [\underline{A}(D_{K'}(K))] \rightarrow [\underline{A}(D_{K'}(K))] \quad \text{is} \quad x \longmapsto x + \gamma(D) \quad .
\end{equation*}
(3)   \underline{For any}  $(u, D) : D_{V}(K) \longrightarrow D_{W}(K)$  , \underline{the map}

}{\begin{equation*}
[\underline{A}((1, D))] : [\underline{A}(D_{K'}(K))] \rightarrow [\underline{A}(D_{K'}(K))] \quad \text{is} \quad x \longmapsto x + \gamma(D) \quad .
\end{equation*}
(3)   \underline{For any}  $(u, D) : D_{V}(K) \longrightarrow D_{W}(K)$  , \underline{the map}

}{F7I-VI-018}\begin{equation*}
(u, D) : [\underline{A}(D_{V}(K))] \rightarrow [\underline{A}(D_{W}(K))] \quad \text{i.e.}
\end{equation*}
\begin{equation*}
(u, D) : V \otimes [\underline{A}(D_{K'}(K))] \rightarrow W \otimes [\underline{A}(D_{K'}(K))] \quad \underline{\text{is}}
\end{equation*}
\begin{equation*}
x \longmapsto u(x) + \gamma(D) \quad .
\end{equation*}
Let us introduce the notations

\begin{equation*}
\begin{aligned}
F(V) \quad &= \quad [\underline{A}(D_{V}(K))] \\
A \quad &= \quad F(K') \\
F(u,D) : F(V) \rightarrow F(W) \quad &= \ \text{the map induced by} \ \ (u,D) : D_{V}(K) \rightarrow D_{W}(K) \\
F(D) \quad : F(V) \rightarrow F(V) \quad &: \quad F(D) = F(1, D) \quad .
\end{aligned}
\end{equation*}
By (II') and the fact that  $K'$  is a  $K'$-vector space object in the category of $K'$-vector spaces,  $F(K')$  is a  $K'$-vector space in (Ens), hence a  $K'$-vector space, and (1) is clear. Let us now identify  $F(V)$  and  $V \otimes A$  .

We deduce from the identity

\begin{equation*}
(u, D) = (1, D) \circ (u, 0) \quad \text{that}
\end{equation*}
\begin{equation*}
F(u, D) = F(D) \circ (u \otimes 1_{A}) \quad .
\end{equation*}
The functoriality of  $F$  amounts to

\begin{equation*}
\begin{aligned}
F(0) &= \mathrm{Id} \\
F(D) \circ F(D') &= F(D + D') \\
(u \otimes 1) \circ F(D) \ &= F(u\,D) \circ (u \otimes 1_{A})
\end{aligned}
\end{equation*}
For any  $D : \Omega \rightarrow V$  and any  $n \in N$  , if  $i$  and  $p$  are the maps

\SGAvar{\begin{equation*}
i \ ; \ V \lhook\joinrel\longrightarrow V \oplus K'^{\,n} \qquad\qquad p : V \oplus K'^{\,n} \longrightarrow K^{n} \quad , 
\end{equation*}

}{\begin{equation*}
i \ ; \ V \lhook\joinrel\longrightarrow V \oplus K'^{\,n} \qquad\qquad p : V \oplus K'^{\,n} \longrightarrow K'^{\,n} \quad , 
\end{equation*}

}{F7I-VI-019}\SGApage{I}{48}

\SGAvar{ then  $p i D = 0$  and the following diagramm is commutative

\begin{center}
\begin{tikzcd}[column sep=huge, row sep=large]
V \otimes A \arrow[r] \arrow[d, "F(D)"] \& (V \oplus K'^{\,n}) \otimes A \arrow[r] \arrow[d, "{F(i,D)}"] \& K'^{\,n}\!\otimes A \arrow[d, equal] \\
V \otimes A \arrow[r] \& (V \oplus K'^{\,n}) \otimes A \arrow[r] \& K'^{\,n}\!\otimes A
\end{tikzcd}
\end{center}













Let  $x : K^{n} \rightarrow V$  be a map,  $x_{i} = x(e_{i})$  , and let  $a_{i}$  be a family of elements of  $A$  . We have  $(1_{V} + x) \circ i \circ D = D$  . Hence, if

\begin{equation*}
F(i\,D)\ (\Sigma\ e_{i}\ a_{i}) = \Sigma\ e_{i}\ a_{i} + a \qquad (a \in V \otimes H) \quad , \ \text{then}
\end{equation*}
\begin{equation*}
F(D)\ (\Sigma\ x_{i}\ a_{i}) \ = \Sigma\ x_{i}\ a_{i} + a
\end{equation*}
and  $a$  does not depend on the  $x_{i}$  . Hence,

\begin{equation*}
\begin{aligned}
F(D)\ (y) \ &= y + G(D) &&\qquad , \ \text{with} \\
G(D + D') &= G(D) + G(D') & \\
u\ G(D) \quad &= F(u\,D) &&\qquad (\text{for any} \ \ u : V \rightarrow W) \quad .
\end{aligned}
\end{equation*}
One easily checks that any such  $G$  is defined, for some  $\gamma \in \Omega \otimes A$  , by the formula

}{ then  $p i D = 0$  and the following diagram is commutative

\begin{center}
\begin{tikzcd}[column sep=huge, row sep=large]
V \otimes A \arrow[r] \arrow[d, "F(D)"] \& (V \oplus K'^{\,n}) \otimes A \arrow[r] \arrow[d, "{F(i\,D)}"] \& K'^{\,n}\!\otimes A \arrow[d, equal] \\
V \otimes A \arrow[r] \& (V \oplus K'^{\,n}) \otimes A \arrow[r] \& K'^{\,n}\!\otimes A
\end{tikzcd}
\end{center}













Let  $x : K'^{\,n} \rightarrow V$  be a map,  $x_{i} = x(e_{i})$  , and let  $a_{i}$  be a family of elements of  $A$  . We have  $(1_{V} + x) \circ i \circ D = D$  . Hence, if

\begin{equation*}
F(i\,D)\ (\Sigma\ e_{i}\ a_{i}) = \Sigma\ e_{i}\ a_{i} + a \qquad (a \in V \otimes A) \quad , \ \text{then}
\end{equation*}
\begin{equation*}
F(D)\ (\Sigma\ x_{i}\ a_{i}) \ = \Sigma\ x_{i}\ a_{i} + a
\end{equation*}
and  $a$  does not depend on the  $x_{i}$  . Hence,

\begin{equation*}
\begin{aligned}
F(D)\ (y) \ &= y + G(D) &&\qquad , \ \text{with} \\
G(D + D') &= G(D) + G(D') & \\
u\ G(D) \quad &= G(u\,D) &&\qquad (\text{for any} \ \ u : V \rightarrow W) \quad .
\end{aligned}
\end{equation*}
One easily checks that any such  $G$  is defined, for some  $\gamma \in \Omega \otimes A$  , by the formula

}{F7I-VI-020}\begin{equation*}
G(D) = D \otimes 1_{A}\ (\gamma) \quad ,
\end{equation*}
as promised in (2) (3).

\underline{Definition} 1.19.   (1) \underline{An object}  $\xi$  \underline{in}  $\widehat{\underline{A}}$  \underline{is called}   \underline{versal formal object if} \underline{every surjective map}  $\alpha : a' \longrightarrow a$  \underline{in}  $\underline{A}$  \underline{induces a surjective map}

\SGAvar{\begin{equation*}
\operatorname{Hom}_{\widehat{\underline{A}}}(\xi, a') \longrightarrow \operatorname{Hom}_{\widehat{\underline{A}}}\ (\xi, a) \quad .
\end{equation*}
(2)   \underline{A formal versal object}  $\xi$  \underline{in}  $\widehat{A}$  , \underline{with image}  $(R, s)$  \underline{in}  $\widehat{C}_{\Lambda,K}$  , \underline{is} \underline{called  minimal  if for any two maps}  $\varphi_{1}$ , $\varphi_{2} : R \rightrightarrows D_{K'}(K)$  , \underline{if}  $\varphi_{1}(\xi)$  \underline{is iso-} \underline{morphic to}  $\varphi_{2}(\xi)$  , \underline{then there exists a derivation}  $D$  \underline{of}  $K$  \underline{into}  $K'$  \underline{such that}  $(1, D)\ (\varphi_{1}) = \varphi_{2}$  , i.e.  $\varphi_{2}\text{-}\varphi_{2} = D\,s$  .

}{\begin{equation*}
\operatorname{Hom}_{\widehat{\underline{A}}}(\xi, a') \longrightarrow \operatorname{Hom}_{\widehat{\underline{A}}}\ (\xi, a) \quad .
\end{equation*}
(2)   \underline{A formal versal object}  $\xi$  \underline{in}  $\widehat{A}$  , \underline{with image}  $(R, s)$  \underline{in}  $\widehat{C}_{\Lambda,K}$  , \underline{is} \underline{called  minimal  if for any two maps}  $\varphi_{1}$ , $\varphi_{2} : R \rightrightarrows D_{K'}(K)$  , \underline{if}  $\varphi_{1}(\xi)$  \underline{is iso-} \underline{morphic to}  $\varphi_{2}(\xi)$  , \underline{then there exists a derivation}  $D$  \underline{of}  $K$  \underline{into}  $K'$  \underline{such that}  $(1, D)\ (\varphi_{1}) = \varphi_{2}$  , i.e.  $\varphi_{2}\text{-}\varphi_{1} = D\,s$  .

}{F7I-VI-021}\underline{Theorem} 1.20.   \underline{If}  $\underline{A}$  \underline{is semi-homogeneous and if}  ,

\begin{equation*}
\dim_{K'}\ ([A(D_{K'}(K))]) < \infty \quad ,
\end{equation*}
\underline{then}  $\underline{A}$  \underline{admits a minimally versal formal object, and any two of them are}
\SGApage{I}{49}
\underline{isomorphic}.

With the notations of 1.18, let  $H^*$  be a subspace of  $[\underline{A}(D_{K'}(K))]$ , supplementary to the image of  $\gamma$ , and let  $H$  be the dual of  $H^*$ . We define $R_o = D_H(K)$ . Let  $\xi_o$  be an object in  $\underline{A}(R_o)$ , whose image in $[\underline{A}(R_o)] \simeq H \otimes [\underline{A}(D_{K'}(K))] \simeq \operatorname{Hom}(H^*,[\underline{A}(D_{K'}(K))])$  is the inclusion of  $H^*$  in $[\underline{A}(D_{K'}(K))]$ . Then

(1)  The maps  $\operatorname{Hom}_{C_{\wedge,K}}(D_H(K),\ D_V(K)) \longrightarrow [\underline{A}(D_V(K))]$  are surjective .

(2)  The minimality condition 1.19(2)  is valid for  $(R_o,\ \xi_o)$  .

Choose now some  $S \in \hat{C}_{\wedge,K}$ , formally smooth over  $\wedge$ , and an epimorphism $t : S \longrightarrow R_o$  such that

\begin{equation*}
\Omega_{R_o/\wedge} \otimes K' \overset{\sim}{\longrightarrow} \Omega_{S/\wedge} \otimes K' \quad .
\end{equation*}
If  $M$  is the maximal ideal of  $S$ , we inductively define, for  $n \geq 1$ , $R_n = S/I_n$  and  $\xi_n \in \operatorname{Ob}(\underline{A}(R_n))$  by the properties that

(a)  $I_{n+1}$  is the intersection of the ideals  $I : I_n M \subset I \subset I_n$ , such that  $\xi_n$ extends to an object  $\xi$  in  $\underline{A}(S/I)$  .

(b)  $\xi_{n+1}$  is an extension of  $\xi_n$  .

If  $R = \varprojlim R_n$  and  $\xi = \varprojlim \xi_n$ , then one checks that  $(R,\ \xi)$  is a minimally-versal formal object, and the only one up to isomorphism.


\SGAsectiontoc{2. Prorepresentable cofibered groupoids.}{2}
2. \underline{Prorepresentable cofibered groupoids.}

\SGAvar{Let  $\underline{A}$  be a cofibered groupoid over  $\underline{C}_\wedge$ , admitting a minimally-versal formal object  $\xi$ . In general we can not recover  $\underline{A}$  (up to  $\underline{C}_\wedge$-equivalence) out of  $\xi$  alone..One has to know when two maps  $R \overset{f}{\underset{g}{\rightrightarrows}} S$  yield a  $\underline{C}_\wedge$-isomorphism $\xi(f) \longrightarrow \xi(g)$ , where  $R = P(\xi)$ . As an application of 1.11  and 1.12 we now consider the question of  "recovering  $\underline{A}$ "  up to  $\underline{C}_\wedge$-equivalence. Indeed, we
\SGApage{I}{50}
shall have a quite satisfactory answer applicable to a wide class of cofibered groupoids.

}{Let  $\underline{A}$  be a cofibered groupoid over  $\underline{C}_\wedge$ , admitting a minimally-versal formal object  $\xi$ . In general we can not recover  $\underline{A}$  (up to  $\underline{C}_\wedge$-equivalence) out of  $\xi$  alone. One has to know when two maps  $R \overset{f}{\underset{g}{\rightrightarrows}} S$  yield a  $\underline{C}_\wedge$-isomorphism $\xi(f) \longrightarrow \xi(g)$ , where  $R = P(\xi)$ . As an application of 1.11  and 1.12 we now consider the question of  "recovering  $\underline{A}$ "  up to  $\underline{C}_\wedge$-equivalence. Indeed, we
\SGApage{I}{50}
shall have a quite satisfactory answer applicable to a wide class of cofibered groupoids.

}{F7I-VI-022}A number of cofibered groupoids which usually occur in practice enjoy a slightly stronger property than those of semi-homogeneous ones, as it will be defined explicitly in 2.5. To investigate them it is convenient to rephrase what we have done so far in terms of  "smooth functors" . All the notations and terminologies in the previous section remain in force throughout in this section.

\underline{Definition} 2.1.  \underline{Let}  $\varphi$ : $\underline{A} \longrightarrow \underline{B}$  \underline{be a functor of cofibered groupoids over}  $\underline{C}_\wedge$ . \underline{We say that}  $\varphi$  \underline{is smooth if every surjective map}  $\beta : b' \longrightarrow \varphi(a)$  \underline{in}  $\underline{B}$  \underline{can be completed to a commutative diagram}

\begin{equation}\tag{*}
\begin{tikzcd}[row sep=large, column sep=huge]
\varphi(a') \arrow[r] \arrow[dr, "\varphi(\alpha)"'] \& b' \arrow[d, "\beta"] \\
 \& \varphi(a)
\end{tikzcd}
\end{equation}










\underline{for some}  $\alpha : a' \longrightarrow a$  \underline{in}  $\underline{A}$ . \underline{We say that}  $\varphi$  \underline{is minimally smooth if it is smooth and}  $\varphi(k[\epsilon]) : [\underline{A}(k[\epsilon])] \longrightarrow [\underline{B}(k[\epsilon])]$  \underline{is bijective}.

\underline{Remark} 2.2(a)  Let  $\varphi$ : $\underline{A} \longrightarrow \underline{B}$  be a functor of cofibered groupoids over  $\underline{C}_\wedge$ . We note that  " $\varphi$  is smooth" $\Leftrightarrow$ "every surjective map  $\beta : b' \longrightarrow \varphi(a)$  in  $\underline{B}$ can be completed into a commutative diagram (*) such that  $\varphi(a') \longrightarrow b'$  is a $\underline{C}_\wedge$-isomorphism"  $\Leftrightarrow$  "every small prolongation  $\beta : b' \longrightarrow \varphi(a)$  in  $\underline{B}$  can be completed into a commutative diagram (*) ".

\underline{Remark} 2.2(b)  For any cofibered groupoid  $\underline{A}$  over  $\underline{C}_\wedge$ , the natural functor $\underline{A} \longrightarrow [\underline{A}]$  is minimally smooth.

\underline{Remark} 2.2(c)  For each object  $R$  in  $\hat{\underline{C}}_\wedge$ , we set  $h_R : \underline{C}_\wedge \longrightarrow$ (Ens)  to be the functor defined by  $h_R(S) = \operatorname{Hom}_{\wedge\text{-alg}}(R,\ S)$  for all  $S$  in  $\underline{C}_\wedge$ . It associates a cofibered discrete groupoid  $\underline{h}_R$  over  $\underline{C}_\wedge$ , and a map  $R \longrightarrow S$  in  $\underline{C}_\wedge$
\SGApage{I}{51}
induces a functor  $\underline{h}_S \longrightarrow \underline{h}_R$  of cofibered groupoids. The functor  $\underline{h}_S \longrightarrow \underline{h}_R$  induced by a map  $R \longrightarrow S$  in  $\underline{C}_\wedge$  is smooth if and only if  $S$  is a formal power-series ring over  $R$ .

\SGAvar{\underline{Remark} 2.2(d)  Given an object  $R$  in  $\widehat{\underline{C}}_\wedge$ , let us denote by  $(R,\ \underline{C}_\wedge)$  the category of arrows in  $\widehat{\underline{C}}_\wedge$  with the source  $R$  and a target in  $\underline{C}_\wedge$ . Let  $\underline{A}$  be a cofibered groupoid over  $\underline{C}_\wedge$ . Given  $\xi \in \mathrm{ob}\ \widehat{\underline{A}}$  with  $P(\xi) = R$ , let  $\underline{A}' \longrightarrow (R,\ \underline{C}_\wedge)$  be the pull-back of  $\underline{A} \longrightarrow \underline{C}_\wedge$  under the canonical functor  $(R,\ \underline{C}_\wedge) \longrightarrow \underline{C}_\wedge$  and choose a section  $\xi^{\#}$  such that  $\xi^{\#}(1_R) = \xi$ . (This simply means that we choose, for each map  $R \xrightarrow{\ f\ } S$ , a $f$-map  $\xi \longrightarrow \xi(f)$ ). We then obtain a functor  $\xi^{\#} : \underline{h}_R \longrightarrow \underline{A}$  over  $\underline{C}_\wedge$  by sitting  $(R \xrightarrow{\ f\ } S) \longmapsto \xi(f)$ . Of course this functor does not depend on the choice of  $\xi^{\#}$  up to natural equivalence of functors. Now  " $\xi^{\#} : \underline{h}_R \longrightarrow \underline{A}$  is (minimally) smooth"  means that every surjective map  $a' \xrightarrow{\ \alpha\ } \xi(f)$  in  $\underline{A}$  can be completed into a commutative diagram

}{\underline{Remark} 2.2(d)  Given an object  $R$  in  $\widehat{\underline{C}}_\wedge$ , let us denote by  $(R,\ \underline{C}_\wedge)$  the category of arrows in  $\widehat{\underline{C}}_\wedge$  with the source  $R$  and a target in  $\underline{C}_\wedge$ . Let  $\underline{A}$  be a cofibered groupoid over  $\underline{C}_\wedge$ . Given  $\xi \in \mathrm{ob}\ \widehat{\underline{A}}$  with  $P(\xi) = R$ , let  $\underline{A}' \longrightarrow (R,\ \underline{C}_\wedge)$  be the pull-back of  $\underline{A} \longrightarrow \underline{C}_\wedge$  under the canonical functor  $(R,\ \underline{C}_\wedge) \longrightarrow \underline{C}_\wedge$  and choose a section  $\xi^{\#}$  such that  $\xi^{\#}(1_R) = \xi$ . (This simply means that we choose, for each map  $R \xrightarrow{\ f\ } S$ , a $f$-map  $\xi \longrightarrow \xi(f)$ ). We then obtain a functor  $\xi^{\#} : \underline{h}_R \longrightarrow \underline{A}$  over  $\underline{C}_\wedge$  by setting  $(R \xrightarrow{\ f\ } S) \longmapsto \xi(f)$ . Of course this functor does not depend on the choice of  $\xi^{\#}$  up to natural equivalence of functors. Now  " $\xi^{\#} : \underline{h}_R \longrightarrow \underline{A}$  is (minimally) smooth"  means that every surjective map  $a' \xrightarrow{\ \alpha\ } \xi(f)$  in  $\underline{A}$  can be completed into a commutative diagram

}{F7I-VI-023}\begin{center}
\begin{tikzcd}[row sep=large, column sep=huge]
\xi(f') \arrow[r] \arrow[dr, "\xi(\ \ )"'] \& a' \arrow[d, "\alpha"] \\
 \& \xi(f)
\end{tikzcd}
\end{center}











(and  $\underline{h}_R(k[\epsilon]) \longrightarrow \underline{A}(k[\epsilon])$   is bijective), which means precisely that  $\xi$  is a (minimally) versal formal object of  $\underline{A}$  over  $\underline{C}_\wedge$ . Thus we conclude that the existence of a (minimally) versal formal object of  $\underline{A}$  is equivalent to the existence of a (minimally) smooth functor  $\underline{h}_R \longrightarrow \underline{A}$  over  $\underline{C}_\wedge$  for some  $R \in \mathrm{ob}(\widehat{\underline{C}}_\wedge)$ .

The following general properties on smoothness are trivial to verify.

\SGAvar{\underline{Proposition} 2.3.  \underline{Let}  $\underline{A} \xrightarrow{\ \varphi\ } \underline{B} \xrightarrow{\ \varphi\ } \underline{C}$  \underline{be functors of cofibered groupoids over}  $\underline{C}_\wedge$ .

}{\underline{Proposition} 2.3.  \underline{Let}  $\underline{A} \xrightarrow{\ \varphi\ } \underline{B} \xrightarrow{\ \psi\ } \underline{C}$  \underline{be functors of cofibered groupoids over}  $\underline{C}_\wedge$ .

}{F7I-VI-024}(a)  \underline{If}  $\varphi, \psi$  \underline{are both smooth, then so is}  $\psi.\varphi$ .

(b)  \underline{If}  $\psi.\varphi$  \underline{is smooth and}  $\varphi$  \underline{is surjective on isomorphism classes of objects,}
\SGApage{I}{52}
\underline{then}  $\psi$  \underline{is smooth.}

\SGAvar{\underline{Remark} 2.4.  Let  $\underline{A}$  be a cofibered groupoids over  $\underline{C}_\wedge$ , admitting a versal formal object  $\xi$ . This means, by 2.2(d), that  $\tilde{\xi} : \underline{h}_R \to \underline{A}$  is smooth, where  $R = P(\xi)$ . Therefore we have,  " $\underline{A}$  is (formally) smooth over  $\underline{C}_\wedge$  (c.f. 1.13) i.e.,  $R$  is a formal power-series algebra over  $\wedge$ "  $\Leftrightarrow$  "the composite functor  $\underline{h}_R \xrightarrow{\ \tilde{\xi}\ } \underline{A} \xrightarrow{\ P\ } \underline{C}_\wedge$   $(= \underline{h}_\wedge)$   is smooth"  $\Leftrightarrow$  " $\underline{A} \xrightarrow{\ P\ } \underline{C}_\wedge$  is a smooth functor"  by 2.3(b).

}{\underline{Remark} 2.4.  Let  $\underline{A}$  be a cofibered groupoid over  $\underline{C}_\wedge$ , admitting a versal formal object  $\xi$ . This means, by 2.2(d), that  $\tilde{\xi} : \underline{h}_R \to \underline{A}$  is smooth, where  $R = P(\xi)$ . Therefore we have,  " $\underline{A}$  is (formally) smooth over  $\underline{C}_\wedge$  (c.f. 1.14) i.e.,  $R$  is a formal power-series algebra over  $\wedge$ "  $\Leftrightarrow$  "the composite functor  $\underline{h}_R \xrightarrow{\ \tilde{\xi}\ } \underline{A} \xrightarrow{\ P\ } \underline{C}_\wedge$   $(= \underline{h}_\wedge)$   is smooth"  $\Leftrightarrow$  " $\underline{A} \xrightarrow{\ P\ } \underline{C}_\wedge$  is a smooth functor"  by 2.3(b).

}{F7I-VI-025}We now consider a stronger property than being semi-homogeneous.

\underline{Definition} 2.5.  \underline{Let}  $\underline{A}$  \underline{be a cofibered groupoid over}  $\underline{C}_\wedge$ . \underline{It is called homogeneous if every diagram}

\begin{center}
\begin{tikzcd}[row sep=large, column sep=huge]
 \& a' \arrow[d] \\
b' \arrow[r] \& a
\end{tikzcd}
\end{center}

\underline{where}  $a' \to a$  \underline{is surjective admits a fibre-product}  $a' \underset{a}{\times} b'$  \underline{in}  $\underline{A}$  \underline{such that}

\begin{equation*}
P(a' \times b') = P(a') \underset{P(a)}{\times} P(b') \ .
\end{equation*}
\underline{Remark} 2.6(a)  Let  $\underline{A}$  be a cofibered groupoid over  $\underline{C}_\wedge$ , and consider a diagram

(*) \qquad
\begin{tikzcd}[row sep=small, column sep=small]
 \& a' \arrow[dl] \arrow[dr] \& \\
a_1 \arrow[dr] \& \& a_2 \arrow[dl] \\
 \& a \&
\end{tikzcd}
\qquad \text{in } $\underline{A}$ \text{ above the diagram} \qquad
\begin{tikzcd}[row sep=small, column sep=small]
 \& R_1 \underset{R}{\times} R_2 \arrow[dl, "\pi_1"'] \arrow[dr, "\pi_2"] \& \\
R_1 \arrow[dr] \& \& R_2 \arrow[dl] \\
 \& R \&
\end{tikzcd}
















in  $\underline{C}_\wedge$  where  $\pi_i$  are the projections. If a fibre-product  $a_1 \underset{a}{\times} a_2$  exists in  $\underline{A}$ , we obtain a unique map  $\varphi : a' \to a_1 \underset{a}{\times} a_2$  compatible with the projections. If  $P(a_1 \underset{a}{\times} a_2) = R_1 \underset{R}{\times} R_2$ , then  $P(\varphi)$  is the identity map on  $R_1 \underset{R}{\times} R_2$  and therefore  $\varphi$  is a  $\underline{C}_\wedge$-isomorphism i.e.,  $a' = a_1 \underset{a}{\times} a_2$ . In other words, if  $\underline{A}$  is a homogeneous cofibered groupoid over  $\underline{C}_\wedge$ , every  diagram (*) defines a fibre-product  $a_1 \underset{a}{\times} a_2$  provided  $R_1 \to R$  is surjective in  $\underline{C}_\wedge$ . This fact entails immediately
\SGApage{I}{53}
that a homogeneous cofibered groupoid is, a fortiori, semi-homogeneous.


\underline{Remark} 2.6(b)   Consider a diagram

\begin{center}
\begin{tikzcd}[row sep=large, column sep=large]
X \arrow[dr, "\varphi"'] \& \& Y \arrow[dl, "\psi"] \\
 \& Z \&
\end{tikzcd}
\end{center}

of categories and functors. The ordinary fibre-product  $X \times_Z Y$  is a category in which an object is a pair  $(x,y)$  such that  $\varphi(x) = \psi(y)$ , and a map $(x,y) \to (x', y')$  consists of a pair  $\alpha : x \to x'$ , $\beta : y \to y'$  such that $\varphi(\alpha) = \psi(\beta)$ . The fibre-product  $X \times^2_Z Y$  in the sense of  2-categories is a category in which an object is a triple  $(x,y; \gamma)$ , where  $\varphi(x) \xrightarrow{\gamma} \psi(y)$  is an isomorphism, and a map  $(x,y; \gamma) \to (x',y'; \gamma')$  consists of a pair  $x \xrightarrow{\alpha} x'$   $y \xrightarrow{\beta} y'$  with a commutative diagram

\begin{center}
\begin{tikzcd}[row sep=huge, column sep=huge]
\varphi(x) \arrow[r, "\varphi(\alpha)"] \arrow[d, "\gamma"'] \& \varphi(x') \arrow[d, "\gamma'"] \\
\psi(y) \arrow[r, "\psi(\beta)"'] \& \psi(y')
\end{tikzcd}
\end{center}









One can rephrase the notions of (semi)homogeneous cofibered groupoid or smooth functor in terms of fibre-product of categories (in the sense of 2-category). Now let  $P : \underline{A} \to \underline{C}_\wedge$  be a cofibered category, and consider an object $R_1 \times_R R_2$  in  $\underline{C}_\wedge$ . If we choose a direct image functor  $(\pi_i)_*: \underline{A}(R_1 \times_R R_2) \to \underline{A}(R_i)$ , where  $\pi_i : R_1 \times_R R_2 \to R_i$  is the projection for  $i = 1,2$,  we obtain a functor $\underline{A}(R_1 \times_R R_2) \to \underline{A}(R_1) \times^2_{\underline{A}(R)} \underline{A}(R_2)$ . It is clear from the definition of cofibered category that this functor, up to canonical natural equivalence, does not depend on the choice of direct image functors. It is now clear from the definition that a cofibered groupoid  $\underline{A}$  over  $\underline{C}_\wedge$  is homogeneous over  $\underline{C}_\wedge$  if and only if $\underline{A}(R_1 \times_R R_2) \to \underline{A}(R_1) \times^2_{\underline{A}(R)} \underline{A}(R_2)$  is an equivalence of categories whenever  $R_1 \to R$  is surjective in  $\underline{C}_\wedge$ , and  $\underline{A}$  is semi-homogeneous if and only if

\SGAvar{\begin{equation*}
[\underline{A}(R \times_R R_2)] \longrightarrow [\underline{A}(R_1) \times^2_{\underline{A}(R)} \underline{A}(R_2)] 
\end{equation*}

}{\begin{equation*}
[\underline{A}(R_1 \times_R R_2)] \longrightarrow [\underline{A}(R_1) \times^2_{\underline{A}(R)} \underline{A}(R_2)] 
\end{equation*}

}{F7I-VI-026}\SGApage{I}{54}

 is surjective whenever  $R_1 \longrightarrow R$  is surjective in  $\underline{C}_\wedge$ , and is bijective whenever  $R = k$  and  $R_1 = k[\epsilon]$ . Now let  $\varphi : \underline{A} \to \underline{B}$  be a functor of cofibered groupoids over  $\underline{C}_\wedge$ . For each map  $R' \to R$  in  $\underline{C}$  , $\varphi$  induces a functor $\underline{A}(R') \to \underline{A}(R) \times^2_{\underline{B}(R)} \underline{B}(R')$  uniquely determined up to a (canonical) natural equivalence. It is again trivial to see that  $\varphi$  is smooth if and only if the functor $A(R') \to \underline{A}(R) \times^2_{\underline{B}(R)} \underline{B}(R')$  is surjective on isomorphism classes of objects, whenever $R' \to R$  is surjective.


\underline{Remark} 2.6(c). If  $\underline{A}$  is a homogeneous cofibered groupoid over  $\underline{C}_\wedge$ , the functor $A(k[V \times W]) \to \underline{A}(k[V])\times \underline{A}(k[W])$         is an equivalence of categories for any finite-dimensional vector spaces  $V, W$  over  $k$  (We recall our convention $\underline{A}(k) = \{e\}$  ). In particular, we have

\begin{equation*}
\mathrm{Aut}_{\underline{A}(k[V\times W])}(1_{V\times W}) \xrightarrow{\ \sim\ } \mathrm{Aut}_{\underline{A}(k[V])}(1_V)\times\mathrm{Aut}_{\underline{A}(k[W])}(1_W)
\end{equation*}
where  $1_V$  is the direct image of  $e$  under the unique map  $k \to k[V]$  in  $\underline{C}_\wedge$  . It follows that if  $\underline{A}$  is homogeneous over  $\underline{C}_\wedge$ , $\mathrm{Aut}(1_\epsilon)$  as well as  $[\underline{A}(k[\epsilon])]$ are canonically provided with vector space structures over  $k$  .


\underline{Remark} 2.6(d). For any object  $R$  in  $\hat{\underline{C}}_\wedge$ , the cofibered (discrete) groupoid $\underline{h}_R$  over  $\underline{C}_\wedge$  is homogeneous.

If  $\underline{A}$  is semi-homogeneous, then so is  $[\underline{A}]$ . However, the homogeneity of  $\underline{A}$  does not imply in general the homogeneity of  $[\underline{A}]$ . Indeed we have


\underline{Proposition} 2.7.  \underline{Let}  $\underline{A}$  \underline{be a homogeneous cofibered groupoid over}  $\underline{C}_\wedge$ . \underline{The following statements are equivalent.}

(i)  $[\underline{A}]$  \underline{is homogeneous i.e.}  $[\underline{A}(R' \times_R S)] \longrightarrow [\underline{A}(R')] \times_{[\underline{A}(R)]} [\underline{A}(S)]$  \underline{is bijective} \underline{for every small extension}  $R' \to R$  .

(ii) \underline{For every small prolongation}  $a' \to a$  \underline{in}  $\underline{A}$ , \underline{the map}

\begin{equation*}
\mathrm{Aut}_{\underline{A}(R')}(a') \longrightarrow \mathrm{Aut}_{\underline{A}(R)}(a) 
\end{equation*}

\SGApage{I}{55}

 \underline{is surjective, where}  $R' = P(a')$  ,  $R = P(a)$  .

\underline{Proof}.  (i) $\Rightarrow$ (ii)  Let  $R' \xrightarrow{\ f\ } R$  be a small extension in  $\underline{C}_{\wedge}$  and let  $a' \to a$  be a f-map in  $\underline{A}$ . Given  $\tau \in \mathrm{Aut}_{\underline{A}(R)}(a)$  , we consider  $a' \times_{a} a'$  and  $a' \times_{a} a'_{\tau}$  where  $a'_{\tau} \longrightarrow a$  denotes the composite map  $a' \longrightarrow a \xrightarrow{\ \tau\ } a$  . Since  $[A(R' \times_{R} R')] \longrightarrow [\underline{A}(R')] \times_{[\underline{A}(R)]} [\underline{A}(R')]$  is bijective, we conclude the existence of an isomorphism  $a' \times_{a} a' \xrightarrow{\ u\ } a' \times_{a} a'_{\tau}$  , and  $u$  induces, via the projection on the second factor, an isomorphism  $a' \xrightarrow{\ \tau'\ } a'$  which is a lifting of  $\tau$  .

(ii) $\Rightarrow$ (i)  Let  $a''_{i}$  ($i=1,2$)  be objects in  $\underline{A}(R' \times_{R} S)$  which have the same image under  $[\underline{A}(R' \times_{R} S)] \to [\underline{A}(R')] \times_{[\underline{A}(R)]} [\underline{A}(S)]$  . We then obtain a diagram of solid arrows

\begin{center}
\begin{tikzcd}[column sep=small, row sep=small]
\&\& a''_{2} \arrow[ddll] \arrow[ddrr] \&\& \\
\&\& a''_{1} \arrow[u, dashed] \arrow[dl] \arrow[dr] \&\& \\
a'_{2} \arrow[ddrr] \& a'_{1} \arrow[l, "\sim"] \arrow[dr] \&\& b_{1} \arrow[r, "\tau"] \arrow[dl] \& b_{2} \arrow[ddll] \\
\&\& a_{1} \arrow[d, "\bar{\tau}"] \&\& \\
\&\& a_{2} \&\&
\end{tikzcd}
\qquad \text{above} \qquad
\begin{tikzcd}[column sep=small, row sep=small]
\& R' \times_{R} S \arrow[dl] \arrow[dr] \& \\
R' \arrow[dr] \&\& S \arrow[dl] \\
\& R \&
\end{tikzcd}
\end{center}



































where  $\bar{\tau}$ ; $a_{1} \to a_{2}$  is the induced isomorphism from  $\tau$  so that the square box on the lower right corner is commutative. If the condition (ii) is verified, one can always replace the given isomorphism  $a'_{1} \overset{\sim}{\to} a'_{2}$  so that the square box on the left corner is also commutative. Since  $a''_{i} = a'_{i} \times_{a_{i}} b_{i}$  for  $i = 1, 2$  by  2.6(a), we obtain the isomorphism  $a''_{1} \overset{\sim}{\to} a''_{2}$  . This proves that  $[\underline{A}(R' \times_{R} S)] \longrightarrow [\underline{A}(R')] \times_{[\underline{A}(R)]} [\underline{A}(S)]$  is injective, and hence is bijective since  $\underline{A}$  is homogeneous by hypothesis.

Now let  $\underline{A}$  be a homogeneous cofibered groupoid over  $\underline{C}_{\Lambda}$  , and let  $R' \xrightarrow{\ f\ } R$  be a small extension in  $\underline{C}_{\Lambda}$  . We have a canonical map
\SGApage{I}{56}
$\tilde{f} : R' \times_{R} R' \longrightarrow k[\mathrm{Ker}\ f]$  given by  $(r'_{1}\ ,\ r'_{2}) \longrightarrow r'_{1}(0) + (r'_{1} - r'_{2})$  , and  $1 \times \tilde{f} : R' \times_{R} R' \longrightarrow R' \times_{k} k[\mathrm{Ker}\ f]$  is an isomorphism. Given any two f-maps  $a_{i} \to a$  ($i=1,2$) in  $\underline{A}$  ,  $a_{1} \times_{a} a_{2}$  is an object over  $R' \times_{R} R'$  and in turn we may consider  $\tau(a_{1},\ a_{2}) = (\tilde{f})_{*}(a_{1} \times_{a} a_{2})$  which is an object in  $\underline{A}(k[\mathrm{Ker}\ f])$  . Since  $1 \times \tilde{f} : R' \times_{R} R' \longrightarrow R' \times_{k} k[\mathrm{Ker}\ f]$  is an isomorphism, we have the isomorphism  $a_{1} \times_{a} a_{2} \simeq a_{1} \times \tau(a_{1},\ a_{2})$  compatible with the projections on the first factor. The statement 1.5(2) takes a more definite form for a homogeneous groupoid.

\underline{Lemma} 2.8.  \underline{Let}  $\underline{A}$  \underline{be a homogeneous cofibered groupoid over}  $\underline{C}_{\Lambda}$  , \underline{and}  $a_{i} \xrightarrow{\ \alpha_{i}\ } a$  ($i = 1, 2$)  f-\underline{maps in}  $\underline{A}$  . \underline{Then}

(1)  \underline{There exists}  $\rho$  : $a_{1} \to a_{2}$  \underline{in}  $\underline{A}$  \underline{with}  $\alpha_{2} \cdot \rho = \alpha_{1}$  \underline{if and only if} \underline{there exists an arrow}  $a_{1} \to \tau(a_{1},\ a_{2})$  \underline{in}  $\underline{A}$  .

(2)  \underline{There exists}  $\rho : a_{1} \to a_{2}$  \underline{with}  $\alpha_{2} \cdot \rho = \alpha_{1}$  \underline{such that}  $P(\rho) = \mathrm{id}_{R'}$ \underline{if and only if}  $\tau(a_{1},\ a_{2}) = 0$  \underline{in}  $[\underline{A}(k[\mathrm{Ker}\ f])]$  .

\SGAvar{\underline{Proof}:  We prove both at the same time. There exists  $\rho : a_{1} \to a_{2}$  in  $\underline{A}$  with  $\alpha_{2} \cdot \rho = \alpha_{1}$  (such that  $P(\rho) = \mathrm{id}_{R'}$ )  $\Leftrightarrow$  there exists  $h : a_{1} \to a_{1} \times_{a} a_{2}$  in  $\underline{A}$  with  $\mathrm{pr}_{1} \circ h = \mathrm{id}_{a_{1}}$  (such that  $P(h)$ = diagonal map)  $\Leftrightarrow$ there exists  $g : a_{1} \to a_{1} \times \tau(a_{1},\ a_{2})$  in  $\underline{A}$  with  $\mathrm{pr}_{1} \circ g = \mathrm{id}_{a_{1}}$  (such that  $P(g) : R' \to R' \times_{k} k[\mathrm{Ker}\ f]$  is given by  $r' \to (r',\ r'(0))$  $\Leftrightarrow$ there exists an arrow  $\gamma : a_{1} \to \tau(a_{1},\ a_{2})$  in  $\underline{A}$  (such that  $P(\gamma) : R' \to k[\mathrm{Ker}\ f]$  is given by  $r' \to r'(0))$  . However, there exists an arrow  $\gamma : a_{1} \to \tau(a_{1},\ a_{2})$  in  $\underline{A}$  such that  $P(\gamma) : R' \to k[\mathrm{Ker}\ f]$  is factored through  $R' \to k \to k[\mathrm{Ker}\ f]$  if and only if  $\tau(a_{1},\ a_{2}) = 0$  .

}{\underline{Proof}:  We prove both at the same time. There exists  $\rho : a_{1} \to a_{2}$  in  $\underline{A}$  with  $\alpha_{2} \cdot \rho = \alpha_{1}$  (such that  $P(\rho) = \mathrm{id}_{R'}$ )  $\Leftrightarrow$  there exists  $h : a_{1} \to a_{1} \times_{a} a_{2}$  in  $\underline{A}$  with  $\mathrm{pr}_{1} \circ h = \mathrm{id}_{a_{1}}$  (such that  $P(h)$ = diagonal map)  $\Leftrightarrow$ there exists  $g : a_{1} \to a_{1} \times \tau(a_{1},\ a_{2})$  in  $\underline{A}$  with  $\mathrm{pr}_{1} \circ g = \mathrm{id}_{a_{1}}$  (such that  $P(g) : R' \to R' \times_{k} k[\mathrm{Ker}\ f]$  is given by  $r' \to (r',\ r'(0)))$  $\Leftrightarrow$ there exists an arrow  $\gamma : a_{1} \to \tau(a_{1},\ a_{2})$  in  $\underline{A}$  (such that  $P(\gamma) : R' \to k[\mathrm{Ker}\ f]$  is given by  $r' \to r'(0))$  . However, there exists an arrow  $\gamma : a_{1} \to \tau(a_{1},\ a_{2})$  in  $\underline{A}$  such that  $P(\gamma) : R' \to k[\mathrm{Ker}\ f]$  is factored through  $R' \to k \to k[\mathrm{Ker}\ f]$  if and only if  $\tau(a_{1},\ a_{2}) = 0$  .

}{F7I-VI-027}\underline{Proposition} 2.9.  (a)  \underline{Let}  $\underline{A} \xrightarrow{\ \varphi\ } \underline{B} \xrightarrow{\ \psi\ } \underline{C}$  \underline{be functors of cofibered groupoids} \underline{over}  $\underline{C}_{\Lambda}$  . \underline{Assume that}

(i)  $\underline{B}$  \underline{is semi-homogeneous and}  $\underline{C}$  \underline{is homogeneous}

(ii)  $\psi(k[\epsilon]) : [\underline{B}(k[\epsilon])] \longrightarrow [\underline{C}(k[\epsilon])]$  \underline{is injective.}

\underline{Then the composite functor}  $\psi \circ \varphi$  \underline{is smooth if and only if}  $\varphi$  \underline{and}  $\psi$  \underline{are both}
\SGApage{I}{57}
\underline{smooth.}

(b)  \underline{Let} $\varphi : \underline{A} \longrightarrow \underline{B}$ \underline{be a smooth functor of cofibered groupoids over} $\underline{C}_{\Lambda}$ . \underline{If} $\underline{A}$ \underline{is semi-homogeneous and} $\underline{B}$ \underline{is homogeneous, then} $\widehat{\varphi} : \widehat{\underline{A}} \to \widehat{\underline{B}}$ \underline{is} \underline{smooth} (i.e., \underline{for any surjective map} $b' \xrightarrow{\ \beta\ } \widehat{\varphi}(a)$ \underline{in} $\widehat{\underline{B}}$ , \underline{there exists} $\alpha : a' \to a$ \underline{in} $\widehat{\underline{A}}$ \underline{and a} $\widehat{\underline{C}}_{\Lambda}$\underline{-isomorphism} $\rho : \widehat{\varphi}(a') \xrightarrow{\ \sim\ } b'$ \underline{such that} $\beta \cdot \rho = \varphi(\alpha))$.

\SGAvar{\underline{Proof}: (a)  Assume that the composite functor $\psi \circ \varphi$ is smooth. If $\varphi$ is smooth, then so is $\psi$ by 2.3(b), and therefore it suffices to show that $\varphi$ is smooth. Let a small prolongation $b' \xrightarrow{\ \beta\ } \varphi(a)$ in $\underline{B}$ be given. Since the composite functor $\varphi \circ \psi$ is smooth, there exists an arrow $a' \xrightarrow{\ \alpha\ } a$ in $\underline{A}$ and a $\underline{C}_{\Lambda}$-isomorphism $\rho : \psi(b') \longrightarrow \psi\ \varphi(a')$ with a commutative diagram

}{\underline{Proof}: (a)  Assume that the composite functor $\psi \circ \varphi$ is smooth. If $\varphi$ is smooth, then so is $\psi$ by 2.3(b), and therefore it suffices to show that $\varphi$ is smooth. Let a small prolongation $b' \xrightarrow{\ \beta\ } \varphi(a)$ in $\underline{B}$ be given. Since the composite functor $\psi \circ \varphi$ is smooth, there exists an arrow $a' \xrightarrow{\ \alpha\ } a$ in $\underline{A}$ and a $\underline{C}_{\Lambda}$-isomorphism $\rho : \psi(b') \longrightarrow \psi\ \varphi(a')$ with a commutative diagram

}{F7I-VI-028}\begin{center}
\begin{tikzcd}[column sep=huge, row sep=huge]
\psi(b') \ar[rr, "\rho"] \ar[dr, "\psi(\beta)"'] \& \& \psi\ \varphi(a') \ar[dl, "\psi\ \varphi(\alpha)"] \\
 \& \psi\ \varphi(a) \&
\end{tikzcd}
\end{center}

Since $\underline{B}$ is semi-homogeneous, we can choose a commutative diagram

\begin{center}
\begin{tikzcd}[column sep=small, row sep=small]
 \& b'' \ar[dl] \ar[dr] \& \& \& \& R' \underset{R}{\times} R' \ar[dl] \ar[dr] \& \\
b' \ar[dr, "\beta"'] \& \& \varphi(a') \ar[dl, "\varphi(\alpha)"] \& \text{above} \& R' \ar[dr] \& \& R' \ar[dl] \\
 \& \varphi(a) \& \& \& \& R \&
\end{tikzcd}
\end{center}













\SGAvar{where  $R = P(a)$ , $R' = P(b') = P(a')$ . Since $\underline{C}$ is homogeneous, $\psi(b'') = \psi(b') \underset{\psi\ \varphi(a)}{\times} \psi\varphi(a')$ by 2.6(a), and it follows from 2.8(2) that $0 = \tau(\psi\ \varphi(a'), \psi(b')) = (\widetilde{f})_{*}\psi(b'')$ , and hence $(\widetilde{f})_{*}(b'') = 0$ be the hypothesis (ii). It follows from 1.5(2) that there exists a map $\gamma : b' \longrightarrow \varphi(a')$ such that $\beta = \varphi(\alpha)\cdot\gamma$ .

}{where  $R = P(a)$ , $R' = P(b') = P(a')$ . Since $\underline{C}$ is homogeneous, $\psi(b'') = \psi(b') \underset{\psi\ \varphi(a)}{\times} \psi\varphi(a')$ by 2.6(a), and it follows from 2.8(2) that $0 = \tau(\psi\ \varphi(a'), \psi(b')) = (\widetilde{f})_{*}\psi(b'')$ , and hence $(\widetilde{f})_{*}(b'') = 0$ by the hypothesis (ii). It follows from 1.5(2) that there exists a map $\gamma : b' \longrightarrow \varphi(a')$ such that $\beta = \varphi(\alpha)\cdot\gamma$ .

}{F7I-VI-029}(b)  Let a surjective map $\beta : b' \to \varphi(a)$ in $\widehat{\underline{B}}$ be given. This means that we are given a sequence of commutative diagrams
\SGApage{I}{58}
\begin{tikzcd}[column sep=large]
{} \ar[r, dashed] \& b_3' \ar[r, "\beta_2'"] \ar[d, "\beta_3"] \& b_2' \ar[r, "\beta_1'"] \ar[d, "\beta_2"] \& b_1' \ar[d, "\beta_1"] \\
{} \ar[r, dashed] \& \varphi(a_3) \ar[r, "\varphi(\alpha_2)"'] \& \varphi(a_2) \ar[r, "\varphi(\alpha_1)"'] \& \varphi(a_1)
\end{tikzcd}















in $\underline{B}$ , where $\alpha_{v}$ , $\beta_{v}$ , $\beta_{v}'$ are surjective for all $v$ , such that $\beta = \varprojlim_{v} \beta_{v}$ . To show that $\widehat{\varphi}$ is smooth, it suffices to prove that a commutative diagram below of solid arrows in which $\alpha_{n}$ , $\beta_{n}$ , $\beta_{n+1}$ , $\beta_{n}'$ are surjective can be completed into a commutative diagram including dotted arrows.

\begin{center}
\begin{tikzcd}[column sep=large, row sep=large]
\varphi(a_{n+1}') \ar[rrrr, dashed, "\varphi(\ \ )"] \ar[dr, dashed, leftrightarrow] \ar[ddr, dashed, "\varphi(\ \ )"'] \& \& \& \& \varphi(a_n') \ar[dl, leftrightarrow] \ar[ddl, "\varphi(\alpha_n')"] \\
 \& b_{n+1}' \ar[rr, "\beta_n'"] \ar[d, "\beta_{n+1}"] \& \& b_n' \ar[d, "\beta_n"'] \& \\
 \& \varphi(a_{n+1}) \ar[rr, "\varphi(\alpha_n)"'] \& \& \varphi(a_n) \&
\end{tikzcd}
\end{center}


















\SGAvar{Now $\underline{A}$ is, by hypotheses, semi-homogeneous and $\alpha_{n} : a_{n}' \to a_{n}$ is surjective, and therefore there exists a commutative diagram

}{Now $\underline{A}$ is, by hypotheses, semi-homogeneous and $\alpha'_{n} : a_{n}' \to a_{n}$ is surjective, and therefore there exists a commutative diagram

}{F7I-VI-030}\begin{center}
\begin{tikzcd}[column sep=small, row sep=small]
 \& a' \ar[dl, "\gamma"'] \ar[dr, "\gamma'"] \& \& \& \& R_{n+1} \underset{R_n}{\times} R_n' \ar[dl] \ar[dr] \& \& \\
a_{n+1} \ar[dr, "\alpha_n"'] \& \& a_n' \ar[dl, "\alpha_n'"] \& \text{in } \underline{A} \text{ over} \& R_{n+1} \ar[dr] \& \& R_n' \ar[dl] \& \text{in } \underline{C}_{\Lambda} . \\
 \& a_n \& \& \& \& R_n \& \&
\end{tikzcd}
\end{center}













The hypothesis that $\underline{B}$ is homogeneous entails that

\begin{tikzcd}[column sep=huge, row sep=huge]
 \& \varphi(a') \ar[dl, "\varphi(\gamma)"'] \ar[dr, "\varphi(\gamma')"] \& \\
\varphi(a_{n+1}) \ar[dr, "\varphi(\alpha_n)"'] \& \& \varphi(a_n') \ar[dl, "\varphi(\alpha_n')"] \\
 \& \varphi(a_n) \&
\end{tikzcd}









\SGApage{I}{59}
is a diagram of a fibre-product in $\underline{B}$ , and therefore we obtain a unique map $\delta : b'_{n+1} \to \varphi(a')$ in $\underline{B}$ with an appropriate commutativity. Since $\beta_{n+1}$ , $\beta'_n$ are both surjective, so is $\delta$ and therefore, since $\varphi$ is smooth, there exists $a'_{n+1} \xrightarrow{\ \alpha'\ } a'$ in $\underline{A}$ and a $\underline{C}_{\Lambda}$-isomorphism $\varphi(a'_{n+1}) \xrightarrow[\sim]{\rho} b'_{n+1}$ such that $\delta \cdot \rho = \varphi(\alpha')$ . We then obtain the commutative diagram

\SGAvar{\begin{center}
\begin{tikzcd}[row sep=huge, column sep=huge]
\varphi(a'_{n+1})
  \arrow[rrr, "\varphi(\gamma'\cdot\alpha')"]
  \arrow[dr, "\rho"]
  \arrow[ddr, "\varphi(\gamma\cdot\alpha')"']
\& \& \& \varphi(a'_n)
  \arrow[dl, "\sim"]
  \arrow[ddl, "\varphi(\alpha'_n)"] \\
\& b'_{n+1} \arrow[r] \arrow[d] \& b'_n \arrow[d] \& \\
\& \varphi(a_{n+1} \arrow[r] \& \varphi(a_n) \&
\end{tikzcd}
\end{center}


















}{\begin{center}
\begin{tikzcd}[row sep=huge, column sep=huge]
\varphi(a'_{n+1})
  \arrow[rrr, "\varphi(\gamma'\cdot\alpha')"]
  \arrow[dr, "\rho"]
  \arrow[ddr, "\varphi(\gamma\cdot\alpha')"']
\& \& \& \varphi(a'_n)
  \arrow[dl, "\sim"]
  \arrow[ddl, "\varphi(\alpha'_n)"] \\
\& b'_{n+1} \arrow[r] \arrow[d] \& b'_n \arrow[d] \& \\
\& \varphi(a_{n+1}) \arrow[r] \& \varphi(a_n) \&
\end{tikzcd}
\end{center}


















}{F7I-VI-031}\underline{Corollary} 2.10.  (a) \underline{Let} $\varphi : \underline{A} \to \underline{B}$ \underline{be a smooth functor of cofibered groupoids over} $\underline{C}_{\Lambda}$ . \underline{Let} $\xi$ \underline{be a versal formal object of} $\underline{A}$ , \underline{and} $\eta$ \underline{a minimally-versal formal object of} $\underline{B}$ . \underline{If} $\underline{B}$ \underline{is homogeneous, then for any map} $\beta : \eta \to \varphi(\xi)$ \underline{in} $\underline{\hat{B}}$ , $P(\xi)$ \underline{is a formal power-series ring over} $P(\eta)$ \underline{through} $P(\beta) : P(\eta) \to P(\varphi(\xi)) = P(\xi)$ .

(b)  \underline{Let} $\underline{A}$ \underline{be a homogeneous cofibered groupoid over} $\underline{C}_{\Lambda}$ , \underline{and let} $\xi$ , $\eta$ \underline{be versal formal objects of} $\underline{A}$ . \underline{If} $\eta$ \underline{is minimal, then for any map} $\alpha : \eta \to \xi$ \underline{in} $\hat{A}$ , $P(\xi)$ \underline{is a formal power-series ring over} $P(\eta)$ \underline{through} $P(\alpha) : P(\eta) \to P(\xi)$ .

\underline{Proof}:  (a)  A map $\beta : \eta \to \varphi(\xi)$ induces a commutative diagram

\begin{center}
\begin{tikzcd}[row sep=huge, column sep=huge]
\underline{h}_{P(\xi)} \arrow[r, "\widetilde{\beta}"] \arrow[d, "\widetilde{\xi}"'] \& \underline{h}_{P(\eta)} \arrow[d, "\widetilde{\eta}"] \\
\underline{A} \arrow[r, "\varphi"'] \& \underline{B}
\end{tikzcd}
\end{center}









Since $\widetilde{\xi}$ , $\varphi$ are smooth, $\widetilde{\eta}\cdot\widetilde{\beta} = \varphi\cdot\widetilde{\xi}$ is smooth. On the other hand,
\SGApage{I}{60}
$\widetilde{\eta}(k[\epsilon]) : \underline{h}_{P(\eta)}(k[\epsilon]) \to [\underline{B}(k[\epsilon])]$ is an isomorphism and therefore $\widetilde{\beta}$ is smooth by 2.9(a), which means that $P(\xi)$ is a formal power-series ring over $P(\eta)$ .

(b)  follows from (a) by setting $\underline{A} = \underline{B}$ .

\SGAvar{Finally we consider the question of  "representability"  of a cofibered groupoid. Let $\underline{C}$ be a fixed category. Each object $R$ in $\underline{C}$ defines a functor $h_R : \underline{C} \to$ (Ens)  given by  $h_R(S) = \mathrm{Hom}_{\underline{C}}(R,S)$ , and in turn we obtain a functor $h : \underline{C}^0 \to \underline{\mathrm{Hom}}(\underline{C},(\mathrm{Ens}))$ which is fully-faithful. A  "cocartegory in $\underline{C}$ "  is a diagram

}{Finally we consider the question of  "representability"  of a cofibered groupoid. Let $\underline{C}$ be a fixed category. Each object $R$ in $\underline{C}$ defines a functor $h_R : \underline{C} \to$ (Ens)  given by  $h_R(S) = \mathrm{Hom}_{\underline{C}}(R,S)$ , and in turn we obtain a functor $h : \underline{C}^0 \to \underline{\mathrm{Hom}}(\underline{C},(\mathrm{Ens}))$ which is fully-faithful. A  "cocategory in $\underline{C}$ "  is a diagram

}{F7I-VI-032}\begin{tikzcd}[column sep=huge]
h_{R_1} \times_{h_{R_0}} h_{R_1} \arrow[r] \& h_{R_1} \arrow[r, shift left=3, "s_1"] \arrow[r, shift right=3, "s_2"'] \& h_{R_0} \arrow[l, "\epsilon" description]
\end{tikzcd}
$\qquad\text{with } s_i \cdot \epsilon = \text{identity for } i = 1,2$










\begin{equation*}
(\text{where } \ h_{R_1} \times_{h_{R_0}} h_{R_1} \ = \ (h_{R_1})_{s_1} \times_{h_{R_0}} {}_{s_2}(h_{R_1})) \ \text{ such that}
\end{equation*}
(i)  $c$  is associative and is compatible with  $s_i$  for  $i = 1,2$

\begin{multline*}
(\mathrm{ii}) \quad h_{R_1} \xrightarrow{\ \sim\ } h_{R_0} \times_{h_{R_0}} h_{R_1} \xrightarrow{\ \epsilon \times 1\ } h_{R_1} \times_{h_{R_0}} h_{R_1} \xrightarrow{\ c\ } h_{R_1}
\\[2ex]
\phantom{(\mathrm{ii})} \quad h_{R_1} \xrightarrow{\ \sim\ } h_{R_1} \times_{h_{R_0}} h_{R_0} \xrightarrow{\ 1 \times \epsilon\ } h_{R_1} \times_{h_{R_0}} h_{R_1} \xrightarrow{\ c\ } h_{R_1}
\\[1ex]
\text{are both identities.}
\end{multline*}
If  $(R_0, R_1, s_1, s_2, \epsilon, c)$  is a cocategory in $\underline{C}$ , then for an object $S$ in $\underline{C}$ , we can define a category  $h_{(R_0,\ R_1,\ \ldots)}(S)$  by setting  $\mathrm{ob}\ h_{(R_0,\ R_1,\ldots)}(S) = \mathrm{Hom}_{\underline{C}}(R_0,\ S)$ , and  $\mathrm{Fl}\ h_{(R_0,\ R_1,\ldots)}(S) = \mathrm{Hom}_{\underline{C}}(R_1,S)$  in which the composition rule of arrows is given by  $c : h_{R_1} \times_{h_{R_0}} h_{R_1} \to h_{R_1}$ . Then $h_{(R_0,\ R_1,\ldots)} : \underline{C} \to$ (categories)  is a functor and in turn defines a cofibered category $\underline{h}_{(R_0,\ R_1,\ldots)}$ over $\underline{C}$ . For instance, any object $R$ in $\underline{C}$ defines a discrete cocategory $\underline{h}_R = \underline{h}_{(R,\ R,\ s_1,\ s_2,\ \epsilon,\ c)}$ by setting $s_1$, $s_2$, $\epsilon$, $c$ to be all equal to the identity map. A homomorphism  $(R_0,\ R_1,\ldots) \to (R_0{}',\ R_1{}',\ldots)$
\SGApage{I}{61}
of cocategories in $\underline{C}$ is a pair of maps $h_{R_0} \to h_{R'_0}$ , $h_{R_1} \to h_{R'_1}$ in $\underline{C}$ such that the induced map

\begin{equation*}
h_{(R_0,\ R_1,\ldots)}(S) \to h_{(R_0',\ R_1',\ldots)}(S)
\end{equation*}
is a functor for every object $S$ in $\underline{C}$ . For instance, each cocategory $(R_0,\ R_1,\ldots)$ is provided with a homomorphism $\epsilon : \underline{h}_{R_0} \to \underline{h}_{(R_0,\ R_1,\ldots)}$ of cocategories. A cogroupoid in $\underline{C}$ is by definition a cocategory $(R_0,\ R_1,\ldots)$ in $\underline{C}$ such that $h_{(R_0,\ R_1,\ldots)}(S)$ is a groupoid for every object $S$ in $\underline{C}$ . In other words, a cogroupoid in $\underline{C}$ is a cocategory $(R_0,\ R_1,\ldots)$ in $\underline{C}$ such that $\underline{h}_{(R_0,\ R_1,\ldots)}$ is a cofibered category in groupoids over $\underline{C}$ .

In practice, a cocategory or cogroupoid occur in a following fashion: Let $P : \underline{F} \to \underline{C}$ be a cofibered category, and let $\xi$ be a fixed object in $\underline{F}$ . Set $R_0 = P(\xi)$ and let $(R_0,\ \underline{C})$ be the category of arrows in $\underline{C}$ with the source $R_0$ . Let $P : \underline{F}' \to (R_0,\ \underline{C})$ be the pull-back of $\underline{F}$ under the canonical functor $(R_0,\ \underline{C}) \to \underline{C}$ , and choose a cocartesian section $\xi^{\#}$ of $\underline{F}'$ over $(R_0,\ \underline{C})$ such that $\xi^{\#}(\mathrm{id}_{R_0}) = \xi$ . (This simply means that we choose, for each arrow $R_0 \xrightarrow{f} S$ in $\underline{C}$ , a cocartesian $f$-map $\xi \to \xi^{\#}(f)$ ). Define the functor

\SGAvar{\begin{equation*}
\mathrm{Hom}_{\xi^{\#}} : (R_0 \amalg R_0,\ \underline{C}) \to (\mathrm{Ens}) \qquad \text{given by}
\end{equation*}
\begin{equation*}
X = (f_1,\ f_2) \longmapsto \mathrm{Hom}_{\underline{F}(S)}(\xi^{\#}(f_1),\ \xi^{\#}(f_2)) \qquad \text{where}
\end{equation*}
$S$ $=$ the target of $f_1$ ( $=$ the target of $f_2$ ). We note that this functor does not depend on the choice of $\xi^{\#}$ up to (canonical) natural equivalence. Henceforth, by abuse of notation, we shall simply write $\xi$ for $\xi^{\#}$ . If this functor is representable by an object $X_0 = (R_0 \ \overset{s_1}{\underset{s_2}{=\!=}}\ R_1)$ in $(R_0 \qquad R_0,\ \underline{C})$ i.e. if there exists $R_0$-map $\varphi : \xi(s_1) \to \xi(s_2)$ in $\underline{F}$ such that the induced map $\widetilde{\varphi} : \mathrm{Hom}_{(R_0 \amalg R_0,\ \underline{C})}(X_0,\ X) \to \mathrm{Hom}_{\xi}(X)$ is bijective for all $X$ in $(R_0 \amalg R_0,\ \underline{C})$ , then the object $X_0$ defines a cocategory $(R_0,\ R_1,\ldots)$ in $\underline{C}$ in a natural way. Indeed, a triple of arrows
\SGApage{I}{62}
\begin{tikzcd}[column sep=huge]
R_0 \arrow[r, shift left=2.5, "u"] \arrow[r, "v"] \arrow[r, shift right=2.5, "w"'] \& S
\end{tikzcd}

}{\begin{equation*}
\mathrm{Hom}_{\xi^{\#}} : (R_0 \amalg R_0,\ \underline{C}) \to (\mathrm{Ens}) \qquad \text{given by}
\end{equation*}
\begin{equation*}
X = (f_1,\ f_2) \longmapsto \mathrm{Hom}_{\underline{F}(S)}(\xi^{\#}(f_1),\ \xi^{\#}(f_2)) \qquad \text{where}
\end{equation*}
$S$ $=$ the target of $f_1$ ( $=$ the target of $f_2$ ). We note that this functor does not depend on the choice of $\xi^{\#}$ up to (canonical) natural equivalence. Henceforth, by abuse of notation, we shall simply write $\xi$ for $\xi^{\#}$ . If this functor is representable by an object $X_0 = (R_0 \ \overset{s_1}{\underset{s_2}{\rightrightarrows}}\ R_1)$ in $(R_0 \amalg R_0,\ \underline{C})$ i.e. if there exists $R_0$-map $\varphi : \xi(s_1) \to \xi(s_2)$ in $\underline{F}$ such that the induced map $\widetilde{\varphi} : \mathrm{Hom}_{(R_0 \amalg R_0,\ \underline{C})}(X_0,\ X) \to \mathrm{Hom}_{\xi}(X)$ is bijective for all $X$ in $(R_0 \amalg R_0,\ \underline{C})$ , then the object $X_0$ defines a cocategory $(R_0,\ R_1,\ldots)$ in $\underline{C}$ in a natural way. Indeed, a triple of arrows
\SGApage{I}{62}
\begin{tikzcd}[column sep=huge]
R_0 \arrow[r, shift left=2.5, "u"] \arrow[r, "v"] \arrow[r, shift right=2.5, "w"'] \& S
\end{tikzcd}

}{F7I-VI-033}in $\underline{C}$ defines a composition

\begin{equation*}
\mathrm{Hom}_{\underline{F}(S)}(\xi(v),\ \xi(w)) \times \mathrm{Hom}_{\underline{F}(S)}(\xi(u),\ \xi(v)) \to \mathrm{Hom}_{\underline{F}(S)}(\xi(u),\ \xi(w))
\end{equation*}
and in turn yields the composition rule

\begin{equation*}
h_{R_1}(S) \underset{h_{R_0}(S)}{\times} h_{R_1}(S) \to h_{R_1}(S)
\end{equation*}
depending functorially on $S$ . Consequently the functor $\mathrm{Hom}_{\xi}$ defines the cocategory $(R_0,\ R_1,\ldots)$ in $\underline{C}$ and in turn defines a cofibered category $\underline{h}_{(R_0,\ R_1,\ldots)}$ in $\underline{C}$ . By definition, ob $\underline{h}_{(R_0,\ R_1,\ldots)} = \coprod_{S} \mathrm{Hom}_{\underline{C}}(R_0,\ S)$ (disjoint union for $S \in \mathrm{ob}(\underline{C})$ ), and an arrow $(R_0 \xrightarrow{f} S) \to (R_0 \xrightarrow{f'} S')$ is a pair $S \xrightarrow{h} S'$ and $R_1 \xrightarrow{u} S'$ in $\underline{C}$ such that $u \cdot s_1 = h \cdot f$ and $u \cdot s_2 = f'$ , and the structural functor $P : \underline{h}_{(R_0,\ R_1,\ldots)} \to \underline{C}$ is given by $(R_0 \xrightarrow{f} S) \longmapsto S$ . Now the pair $(\xi,\ \varphi)$ defines the functor

\SGAvar{\begin{equation*}
(\xi,\ \varphi) : \underline{h}_{(R_0,\ R_1,\ldots)} \to \underline{F} \quad \text{over} \quad \underline{C} \quad \text{by}
\end{equation*}
\begin{equation*}
(R_0 \xrightarrow{f} S) \ \longmapsto\ \xi(f)
\end{equation*}
\begin{equation*}
[(R_0 \xrightarrow{f} S) \xrightarrow{\ (h,u)\ } (R_0 \to S')] \longmapsto \ \text{the composit map}\ \ \xi(f) \xrightarrow{\ h_{*}\ } \xi(hf) \xrightarrow{\ u(\varphi)\ } \xi(f') \ .
\end{equation*}
We note that, for each object $S$ in $\underline{C}$ , the functor $(\xi,\varphi)(S)$   $\underline{h}_{(R_0,R_1,\ldots)}(S) \to \underline{F}(S)$ is fully-faithful by virtue of the definition of $(R_0 \rightrightarrows R_1)$ coming from the representability of the functor $\mathrm{Hom}_{\xi}$ . Consequently the functor $(\xi,\varphi) : \underline{h}_{(R_0,R_1,\ldots)} \to \underline{F}$ is an $\underline{C}$-equivalence if and only if $\xi$ is a quasi-initial object of $\underline{F}$ over $\underline{C}$ i.e., for every object $a$ in $\underline{F}$ there exists a cocartesian arrow $\xi \to a$ . Instead of this functor $\mathrm{Hom}_{\xi}$ one may also consider the functor

}{\begin{equation*}
(\xi,\ \varphi) : \underline{h}_{(R_0,\ R_1,\ldots)} \to \underline{F} \quad \text{over} \quad \underline{C} \quad \text{by}
\end{equation*}
\begin{equation*}
(R_0 \xrightarrow{f} S) \ \longmapsto\ \xi(f)
\end{equation*}
\begin{equation*}
[(R_0 \xrightarrow{f} S) \xrightarrow{\ (h,u)\ } (R_0 \to S')] \longmapsto \ \text{the composite map}\ \ \xi(f) \xrightarrow{\ h_{*}\ } \xi(hf) \xrightarrow{\ u(\varphi)\ } \xi(f') \ .
\end{equation*}
We note that, for each object $S$ in $\underline{C}$ , the functor $(\xi,\varphi)(S) :$   $\underline{h}_{(R_0,R_1,\ldots)}(S) \to \underline{F}(S)$ is fully-faithful by virtue of the definition of $(R_0 \rightrightarrows R_1)$ coming from the representability of the functor $\mathrm{Hom}_{\xi}$ . Consequently the functor $(\xi,\varphi) : \underline{h}_{(R_0,R_1,\ldots)} \to \underline{F}$ is a $\underline{C}$-equivalence if and only if $\xi$ is a quasi-initial object of $\underline{F}$ over $\underline{C}$ i.e., for every object $a$ in $\underline{F}$ there exists a cocartesian arrow $\xi \to a$ . Instead of this functor $\mathrm{Hom}_{\xi}$ one may also consider the functor

}{F7I-VI-034}\begin{equation*}
\mathrm{Isom}_{\xi} : (R_0 \amalg R_0,\ \underline{C}) \to (\mathrm{Ens}) 
\end{equation*}

\SGApage{I}{63}

 given by  $X = (f_1,f_2) \longmapsto \mathrm{Isom}_{\underline{F}(S)}(\xi(f_1),\xi(f_2))$  where  $S$  is the target of  $f_1$ ( $=$ the target of  $f_2$ ). If this functor is representable by an object  $(R_0 \rightrightarrows R_1)$ in  $(R_0 \amalg R_0, \underline{C})$  , then it defines  a cogroupoid  $(R_0,R_1,\ldots)$  in  $\underline{C}$  together with the functor  $(\xi,\varphi) : \underline{h}_{(R_0,R_1,\ldots)} \to \underline{F}$ , which is, when  $\underline{F}$  is a cofibered groupoid, a  $\underline{C}$-equivalence if and only if  $\xi$  is a quasi-initial object of  $\underline{F}$  over  $\underline{C}$ .


\underline{Definition and proposition} 2.11.  \underline{A cogroupoid}  $(R_0,R_1,\ldots)$  \underline{in}  $\hat{\underline{C}}_\Lambda$  \underline{is called smooth if any one of the following equivalent conditions holds:}

 (i)  $R_1$  \underline{is a formal power-series ring over}  $R_0$  \underline{via} $s_1$  (\underline{and}  $s_2$ )

 (ii)  $\epsilon : \underline{h}_{R_0} \longrightarrow \underline{h}_{(R_0,R_1,\ldots)}$  \underline{is smooth over}  $\underline{C}_\Lambda$ .

\underline{A minimally smooth cogroupoid or a normalized smooth cogroupoid in}  $\hat{\underline{C}}_\Lambda$  \underline{is a smooth cogroupoid}  $(R_0,R_1,\ldots)$  \underline{in}  $\hat{\underline{C}}_\Lambda$  \underline{such that}  $\underline{h}_{(R_0,R_1,\ldots)}(k[\epsilon])$  \underline{is totally disconnected} i.e.,  $s_1(k[\epsilon]) = s_2(k[\epsilon])$  .

\SGAvar{\underline{Proof}:  $\epsilon : \underline{h}_{R_0} \to \underline{h}_{(R_0,R_1,\ldots)}$  is smooth over  $\underline{C}_\Lambda \Leftrightarrow$ given any surjective map  $(S',f') \xrightarrow{\ (h,u)\ } (S,\ f)$  in  $\underline{h}_{(R_0,R_1,\ldots)}$  , there exists a lifting  $f''$  of  $f$ to  $S'$  together with an isomorphism  $(S',f'') \xrightarrow{\ (1,u')\ } (S',f') \Leftrightarrow$ given  $u \in h_{R_1}(S)$ and a lifting  $f'$  or  $s_1(u)$  to  $S'$  , there exists a lifting  $u'$  of  $u$  to  $S'$ such that  $s_1(u') = f'$  , where  $S' \to S$  is an arbitrary surjective map in $\underline{C}_\Lambda \Leftrightarrow s_1 : \underline{h}_{R_1} \to \underline{h}_{R_0}$  is smooth over  $\underline{C}_\Lambda$ .


}{\underline{Proof}:  $\epsilon : \underline{h}_{R_0} \to \underline{h}_{(R_0,R_1,\ldots)}$  is smooth over  $\underline{C}_\Lambda \Leftrightarrow$ given any surjective map  $(S',f') \xrightarrow{\ (h,u)\ } (S,\ f)$  in  $\underline{h}_{(R_0,R_1,\ldots)}$  , there exists a lifting  $f''$  of  $f$ to  $S'$  together with an isomorphism  $(S',f'') \xrightarrow{\ (1,u')\ } (S',f') \Leftrightarrow$ given  $u \in h_{R_1}(S)$ and a lifting  $f'$  of  $s_1(u)$  to  $S'$  , there exists a lifting  $u'$  of  $u$  to  $S'$ such that  $s_1(u') = f'$  , where  $S' \to S$  is an arbitrary surjective map in $\underline{C}_\Lambda \Leftrightarrow s_1 : \underline{h}_{R_1} \to \underline{h}_{R_0}$  is smooth over  $\underline{C}_\Lambda$ .


}{F7I-VI-035}\underline{Proposition} 2.12.  \underline{Let}  $(R_0,R_1,\ldots)$  \underline{be a cogroupoid in}  $\hat{\underline{C}}_\Lambda$ .

(a)  \underline{If}  $(R_0,R_1,\ldots)$  \underline{is smooth, then}  $\underline{h}_{(R_0,R_1,\ldots)}$  \underline{is homogeneous.}

(b)  \underline{If}  $s_1(k[\epsilon]) = s_2(k[\epsilon])$  , \underline{then the following statements are equivalent}

          (i)  $(R_0,R_1,\ldots)$  \underline{is smooth}

         (ii)  $\underline{h}_{(R_0,R_1,\ldots)}$  \underline{is homogeneous}

        (iii)  $\underline{h}_{(R_0,R_1,\ldots)}$  \underline{is semi-homogeneous}


\SGApage{I}{64}

(c)  \underline{If}  $\underline{h}_{(R_0,R_1,\ldots)}$  \underline{is semi-homogeneous, then it is equivalent to}  $\underline{h}_{(S_0,S_1,\ldots)}$  \underline{for some minimally smooth cogroupoid}  $S_*$  .

\underline{Proof}:   (a)  Consider a diagram

\begin{center}
\begin{tikzcd}[column sep=huge, row sep=large]
(S_1,f_1) \arrow[dr, "{(h_1,u_1)}"'] \& \& (S_2,f_2) \arrow[dl, "{(h_2,u_2)}"] \\
\& (S,f) \&
\end{tikzcd}
\end{center}

in  $\underline{h}_{(R_0,R_1,\ldots)}$  where  $h_2 : S_2 \longrightarrow S$  is surjective in  $\underline{C}_\Lambda$ . Since $s_1 : \underline{h}_{R_1} \longrightarrow \underline{h}_{R_0}$  is smooth by hypothesis,  $u_1^{-1}\, u_2$  can be lifted to some $v \in h_{R_1}(S_2)$  such that  $s_1(v) = f_2$ . Since  $h_2 \cdot s_2(v) = h_1 \cdot f_1$ , $f_1 \times s_2(v) \in h_{R_0}(S_1 \underset{S}{\times} S_2)$  is defined and we obtain the commutative diagram

\begin{center}
\begin{tikzcd}[column sep=huge, row sep=large]
\& (S_1 \underset{S}{\times} S_2,\ f_1 \times s_2(v)) \arrow[dl, "\pi_1"'] \arrow[dr, "{(\pi_2,v^{-1})}"] \& \\
(S_1,f_1) \arrow[dr, "{(h_1,u_1)}"'] \& \& (S_2,f_2) \arrow[dl, "{(h_2,u_2)}"] \\
\& (S,f) \&
\end{tikzcd}
\end{center}










We claim that the above diagram defines a fibre-product  $(S_1,f_1) \underset{(S,f)}{\times} (S_2,f_2)$  in $\underline{h}_{(R_0,R_1,\ldots)}$ . Indeed, for any object  $(T,\ g)$  in  $\underline{h}_{(R_0,R_1,\ldots)}$  we have

{\small \begin{multline*}
\mathrm{Hom}((T,g),(S_1\underset{S}{\times}S_2,f_1\times s_2(v)))=\{(h_1'\times h_2',w_1\times w_2)\in h_T(S_1\underset{S}{\times}S_2)\times h_{R_1}(S_1\underset{S}{\times}S_2)\,|\,s_1(w_1\times w_2)=(h_1'\times h_2')\ g\ ,\\
s_2(w_1\times w_2)=f_1\times s_2(v)\} = \{(h_1'\times h_2',w_1,w_2) \in h_T(S_1\underset{S}{\times}S_2)\ \times h_{R_1}(S_1)\times h_{R_1}(S_2)\,|\,h_1w_1 = h_2w_2,\\
s_1(w_1) = h_1'\ g,\ s_1(w_2) = h_2'\ g,\ s_2(w_1) = f_1,\ s_2(w_2) = s_2(v)\} =\\
\{(h_1'\times h_2',w_1,w_2) \in h_T(S_1\underset{S}{\times}S_2)\times h_{R_1}(S_1)\times h_{R_1}(S_2)\,|\ u_1^{-1}\ u_2\ h_2(v^{-1}w_2)=h_1\ w_1\ ,\ s_1(w_1) = h_1'\ \ g, 
\end{multline*}}

\SGApage{I}{65}

\SGAvar{{\footnotesize \begin{multline*}s_2(w_1)=f_1, \quad s_1(v^{-1}w_2)=h_2'g, \quad s_2(v^{-1}w_2)=f_2\} \;\overset{\sim}{\longrightarrow}\; \{(h_1'\times h_2', w_1, w_2') \in h_T(S_1\times_S S_2)\times h_{R_1}(S_1)\times h_{R_1}(S_2)\;\big| \\ u_2\, h_2\, w_2' \;=\; u_1\, h_1\, w_1, \quad s_1(w_1) \;=\; h_1'g, \quad s_2(w_1)=\; f_1, \quad s_1(w_2') \;=\; h_2'g, \quad s_2(w_2') \;=\; f_2\} \;= \\ \mathrm{Hom}((T,\ g),\ (S_1,\ f_1) \;\times_{\mathrm{Hom}((T,g),\ (S,f))}\; \mathrm{Hom}((T,\ g),\ (S_2,\ f_2))\ .\end{multline*}}
(b)  It suffices to show (iii) $\Rightarrow$ (i) . Since $\underline{h}_{(R_0,R_1,\ldots)}$ is semi-homogeneous and $\dim[\underline{h}_{(R_0,R_1,\ldots)}(k[\epsilon])] = \dim h_{R_0}(k[\epsilon]) < \infty$ , $\underline{h}_{(R_0,R_1,\ldots)}$ admits a minimally versal formal object. On the other hand, $(R_0,\ 1_{R_0})$ is a quasi-initial minimal object of $\widehat{\underline{h}}_{(R_0,R_1,\ldots)}$ since $\underline{h}_{R_0}(k[\epsilon]) \to [\underline{h}_{(R_0,R_1,\ldots,)}(k[\epsilon])]$ is bijective, and therefore $(R_0,\ 1_{R_0})$ is a minimally versal formal object of $\underline{h}_{(R_0,R_1,\ldots)}$ by 1.10(ii), which means that $\epsilon : \underline{h}_{R_0} \longrightarrow \underline{h}_{(R_0,R_1,\ldots)}$ is minimally smooth.

(c)  Let $(S_0,\ \xi_0)$ be a minimally versal formal object for $\underline{h}_{(R_0,R_1,\ldots)}$ (1.11). For any object $\xi : R_0 \to T$ in $\underline{h}_{(R_0,\ R_1,\ldots)}(T)$ , the functor $\mathrm{Isom}(\xi,\xi)$ from $C_\Lambda/T \otimes_\Lambda T$ to (sets) which to any $R$ associates the sets of isomorphisms between the two images of $\xi$ on $R : R_0 \overset{\xi}{\longrightarrow} T \rightrightarrows T \otimes_\Lambda T \longrightarrow R$ , is represented by $R_1 \otimes_{R_0 \hat{\otimes}_\Lambda R_0} (T \otimes_\Lambda T)$ . By a passage to limit, one get that $\mathrm{Isom}(\xi_0,\ \xi_0)$ is pro-represented by some $S_1$ , and $\underline{h}_{(R_0,R_1,\ldots)}$ is equivalent to $\underline{h}_{(S_0,S_1,\ldots)}$ . By (b), $S_*$ is minimally smooth.

}{{\footnotesize \begin{multline*}s_2(w_1)=f_1, \quad s_1(v^{-1}w_2)=h_2'g, \quad s_2(v^{-1}w_2)=f_2\} \;\overset{\sim}{\longrightarrow}\; \{(h_1'\times h_2', w_1, w_2') \in h_T(S_1\times_S S_2)\times h_{R_1}(S_1)\times h_{R_1}(S_2)\;\big| \\ u_2\, h_2\, w_2' \;=\; u_1\, h_1\, w_1, \quad s_1(w_1) \;=\; h_1'g, \quad s_2(w_1)=\; f_1, \quad s_1(w_2') \;=\; h_2'g, \quad s_2(w_2') \;=\; f_2\} \;= \\ \mathrm{Hom}((T,\ g),\ (S_1,\ f_1)) \;\times_{\mathrm{Hom}((T,g),\ (S,f))}\; \mathrm{Hom}((T,\ g),\ (S_2,\ f_2))\ .\end{multline*}}
(b)  It suffices to show (iii) $\Rightarrow$ (i) . Since $\underline{h}_{(R_0,R_1,\ldots)}$ is semi-homogeneous and $\dim[\underline{h}_{(R_0,R_1,\ldots)}(k[\epsilon])] = \dim h_{R_0}(k[\epsilon]) < \infty$ , $\underline{h}_{(R_0,R_1,\ldots)}$ admits a minimally versal formal object. On the other hand, $(R_0,\ 1_{R_0})$ is a quasi-initial minimal object of $\widehat{\underline{h}}_{(R_0,R_1,\ldots)}$ since $\underline{h}_{R_0}(k[\epsilon]) \to [\underline{h}_{(R_0,R_1,\ldots)}(k[\epsilon])]$ is bijective, and therefore $(R_0,\ 1_{R_0})$ is a minimally versal formal object of $\underline{h}_{(R_0,R_1,\ldots)}$ by 1.10(ii), which means that $\epsilon : \underline{h}_{R_0} \longrightarrow \underline{h}_{(R_0,R_1,\ldots)}$ is minimally smooth.

(c)  Let $(S_0,\ \xi_0)$ be a minimally versal formal object for $\underline{h}_{(R_0,R_1,\ldots)}$ (1.11). For any object $\xi : R_0 \to T$ in $\underline{h}_{(R_0,\ R_1,\ldots)}(T)$ , the functor $\mathrm{Isom}(\xi,\xi)$ from $\underline{C}_\Lambda/T \hat{\otimes}_\Lambda T$ to (sets) which to any $R$ associates the set of isomorphisms between the two images of $\xi$ on $R : R_0 \overset{\xi}{\longrightarrow} T \rightrightarrows T \hat{\otimes}_\Lambda T \longrightarrow R$ , is represented by $R_1 \otimes_{R_0 \hat{\otimes}_\Lambda R_0} (T \otimes_\Lambda T)$ . By a passage to limit, one gets that $\mathrm{Isom}(\xi_0,\ \xi_0)$ is pro-represented by some $S_1$ , and $\underline{h}_{(R_0,R_1,\ldots)}$ is equivalent to $\underline{h}_{(S_0,S_1,\ldots)}$ . By (b), $S_*$ is minimally smooth.

}{F7I-VI-036}\underline{Definition} 2.13.  \underline{A cofibered groupoid} $\underline{A}$ \underline{over} $\underline{C}_\Lambda$ \underline{is called (minimally) pro-representable if it is} $\underline{C}_\Lambda$\underline{-equivalent to} $\underline{h}_{(R_0,R_1,\ldots)}$ \underline{for some (normalized) cogroupoid} $(R_0,R_1,\ldots)$ \underline{in} $\widehat{\underline{C}}_\Lambda$ .

\underline{Proposition} 2.14.  \underline{If} $(R_0,\ R_1,\ldots)$ \underline{is a normalized cogroupoid in} $\widehat{\underline{C}}_\Lambda$ , \underline{then} \underline{every homomorphism} $(\alpha_0,\alpha_1) : (R_0,R_1,\ldots) \to (R_0,R_1,\ldots)$ \underline{which induces} $\underline{C}_\Lambda$\underline{-equivalence} $\underline{h}_{(\alpha_0,\alpha_1)} : \underline{h}_{(R_0,R_1,\ldots)} \longrightarrow \underline{h}_{(R_0,\ R_1,\ldots)}$ \underline{is an isomorphism. In}
\SGApage{I}{66}
\underline{particular, for any cofibered groupoid} $\underline{A}$ \underline{over} $\underline{C}_\Lambda$ , \underline{a normalized prorepresentation, if it exists, is unique up to an isomorphism.}

\underline{Proof}:  Consider the commutative diagram

\begin{center}
\begin{tikzcd}[column sep=huge, row sep=huge]
\underline{h}_{(R_0,R_1,\ldots)} \arrow[r, "{\underline{h}_{(\alpha_0,\alpha_1)}}"] \& \underline{h}_{(R_0,R_1,\ldots)} \\
\underline{h}_{R_0} \arrow[u, "\epsilon"] \arrow[r, "\underline{h}_{\alpha_0}"'] \& \underline{h}_{R_0} \arrow[u, "\epsilon"']
\end{tikzcd}
\end{center}













\SGAvar{where $\underline{h}_{(\alpha_0,\alpha_1)}$ is a $\underline{C}_\Lambda$-equivalence. Since $\underline{h}_{(R_0,R_1,\ldots)}(k[\epsilon])$ is totally disconnected, $[\underline{h}_{(\alpha_0,\alpha_1)}(k[\epsilon])$ and $[\epsilon(k[\epsilon])]$ are both bijections, and hence so is $\underline{h}_{\alpha_0}(k[\epsilon])$ .

}{where $\underline{h}_{(\alpha_0,\alpha_1)}$ is a $\underline{C}_\Lambda$-equivalence. Since $\underline{h}_{(R_0,R_1,\ldots)}(k[\epsilon])$ is totally disconnected, $[\underline{h}_{(\alpha_0,\alpha_1)}(k[\epsilon])]$ and $[\epsilon(k[\epsilon])]$ are both bijections, and hence so is $\underline{h}_{\alpha_0}(k[\epsilon])$ .

}{F7I-VI-037}Consequently, $\alpha_0 : R_0 \to R_0$ is a surjective endomorphism and hence is an automorphism. It then follows that $\alpha_1$ is also an automorphism of $R_1$ .

\underline{Lemma} 2.15.  \underline{Let} $A,\ B$ \underline{be discrete groupoids over} $\underline{C}_\wedge$ \underline{and} $\varphi : A \to B$ \underline{a minimally smooth functor}. \underline{If} $A,\ B$ \underline{are homogeneous}, \underline{then} $\varphi$ \underline{is a} $\underline{C}_\wedge$\underline{-equivalence.}

\underline{Proof}:  We must show that $\varphi(S) : A(S) \longrightarrow B(S)$ is bijective for all $S \in \mathrm{ob}(\underline{C}_\wedge)$ . We note that $\varphi(S)$ is surjective for all $S \in \mathrm{ob}(\underline{C}_\wedge)$ and bijective for $S = k[\epsilon]$ . We argue by induction on the length of $S$ . Let $S' \overset{f}{\longrightarrow} S$ be a small extension in $\underline{C}_\wedge$ . Since $A,\ B$ are homogeneous discrete groupoids, the canonical isomorphism $1 \times \tilde{f} : S' \times_S S' \longrightarrow S' \times_k k[\mathrm{Ker}\ f]$ induces the commutative diagram
\SGApage{I}{67}
\begin{tikzcd}[column sep=huge, row sep=huge]
A(S') \underset{A(S)}{\times} A(S') \arrow[r, "\sim"] \arrow[d] \& A(S') \times A(k[\mathrm{Ker}\ f]) \arrow[d] \\
B(S') \underset{B(S)}{\times} B(S') \arrow[r, "\sim"] \& B(S') \times B(k[\mathrm{Ker}\ f])
\end{tikzcd}













Therefore  $A(S')$ , $B(S')$  are principal homogeneous sets over  $A(S)$ , $B(S)$  with the structural group  $A(k[\mathrm{Ker}\ f]) \xrightarrow{\ \sim\ } B(k[\mathrm{Ker}\ f])$  . Consequently if  $\varphi(S)$  is bijective, then  $\varphi(S')$  is injective and hence is bijective.

\SGAvar{\underline{Corollary} 2.16.  \underline{Let}  $\underline{A}$  \underline{be a homogeneous groupoid over}  $\underline{C}_{\wedge}$  \underline{such that} $\mathrm{Aut}_{\underline{A}(k[\epsilon])}(1_{\epsilon})$  \underline{is of finite-dimensional. Then for any}  $\xi \in \mathrm{ob}(\hat{\underline{A}})$  , \underline{the functor} $\mathrm{Isom}_{\xi} : \underline{C}_{R \hat{\otimes}_{\wedge} R} \longrightarrow (\mathrm{Ens})$  \underline{is prorepresentable, where}  $R = P(\xi)$  .

}{\underline{Corollary} 2.16.  \underline{Let}  $\underline{A}$  \underline{be a homogeneous groupoid over}  $\underline{C}_{\wedge}$  \underline{such that} $\mathrm{Aut}_{\underline{A}(k[\epsilon])}(1_{\epsilon})$  \underline{is finite-dimensional. Then for any}  $\xi \in \mathrm{ob}(\hat{\underline{A}})$  , \underline{the functor} $\mathrm{Isom}_{\xi} : \underline{C}_{R \hat{\otimes}_{\wedge} R} \longrightarrow (\mathrm{Ens})$  \underline{is prorepresentable, where}  $R = P(\xi)$  .

}{F7I-VI-038}\underline{Proof}:  Assume that  $\mathrm{Isom}_{\xi}$  is homogeneous. We then have, by 1.11, a minimally smooth functor  $\Phi : \underline{h}_{R_{1}} \to \underline{\mathrm{Isom}}_{\xi}$  , which is a  $\underline{C}_{R \hat{\otimes}_{\wedge} R}$-equivalence by 2.15 i.e., $\mathrm{Isom}_{\xi}$  is prorepresentable. Therefore it suffices to show that  $\mathrm{Isom}_{\xi}$  is homogeneous. Consider a diagram

\begin{center}
\begin{tikzcd}[column sep=large, row sep=large]
T_1 \arrow[dr] \& \& T_2 \arrow[dl] \\
 \& T \&
\end{tikzcd}
\end{center}

in  $\underline{C}_{R \hat{\otimes}_{\wedge} R}$  where  $T_{2} \longrightarrow T$  is surjective. Since  $\underline{A}$  is homogeneous, $\xi_{i}(T_{1}) \underset{\xi_{i}(T)}{\times} \xi_{i}(T_{2})$  exists in  $\underline{A}$  , and we have a canonical isomorphism

$\xi_{i}(T_{1} \underset{T}{\times} T_{2}) \simeq \xi_{i}(T_{1}) \underset{\xi_{i}(T)}{\times} \xi_{i}(T_{2})$  for  $i = 1,2$  (cf. 2.6(a)). It follows that

$\mathrm{Isom}_{\xi}(T_{1} \underset{T}{\times} T_{2}) \longrightarrow \mathrm{Isom}_{\xi}(T_{1}) \underset{\mathrm{Isom}_{\xi}(T)}{\times} \mathrm{Isom}_{\xi}(T_{2})$  is an isomorphism.

\underline{Theorem} 2.17.  \underline{Let}  $\underline{A}$  \underline{be a cofibered groupoid over}  $\underline{C}_{\wedge}$  . \underline{Then it is prorepresentable by a (normalized) smooth cogroupoid if and only if}
\SGApage{I}{68}
(1)  $\underline{A}$  \underline{is homogeneous}

(2)  $[A(k[\epsilon])]$  \underline{and}  $[\mathrm{Aut}(1_{\epsilon})]$  \underline{are both finite-dimensional vector spaces over}  $k$  .

\underline{Proof}:  Assume that  $\underline{A}$  is prorepresentable by a smooth cogroupoid i.e.,  $\underline{A}$  is  $\underline{C}_{\wedge}$-equivalent to  $\underline{h} = \underline{h}_{(R_{0},R_{1},\ldots)}$  for some smooth cogroupoid  $(R_{0},R_{1},\ldots)$  in  $\hat{\underline{C}}_{\wedge}$  . Since  $\underline{h}$  is homogeneous by 2.12(a), and

\SGAvar{\begin{equation*}
\begin{aligned}
[\underline{h}(k[\epsilon])] &= \mathrm{Coker}\ (\mathrm{Hom}_{\wedge\text{-alg}}(R_{1},k[\epsilon] \ \underset{s_{2}}{\overset{s_{1}}{\rightrightarrows}}\ \mathrm{Hom}_{\wedge\text{-alg}}(R_{0},k[\epsilon])) \qquad \text{as well as} \\[2ex]
\mathrm{Aut}_{\underline{h}}(1_{\epsilon}) &= \mathrm{Ker}(\mathrm{Hom}_{\wedge\text{-alg}}(R_{1},k[\epsilon]) \xrightarrow{\ s_{1}\ } \mathrm{Hom}_{\wedge\text{-alg}}(R_{0},k[\epsilon])) \\[2ex]
&= \mathrm{Ker}(\mathrm{Hom}_{\wedge\text{-alg}}(R_{1},k[\epsilon]) \xrightarrow{\ s_{2}\ } \mathrm{Hom}_{\wedge\text{-alg}}(R_{0},k[\epsilon])) \qquad \text{are both finite-dimensional}
\end{aligned}
\end{equation*}
vector spaces over  $k$  , we must have the same properties on  $\underline{A}$  . Conversely assume (1) and (2) on  $\underline{A}$  . We claim that  $\underline{A}$  is prorepresentable by a normalized smooth cogroupoid. Since  $\underline{A}$  is homogeneous and  $[\underline{A}(k[\epsilon])]$  is a finite-dimensional vector space over  $k$  ,  $\underline{A}$  admits a minimally versal formal object  $\xi$  . Set  $P(\xi) = R_{0}$  and consider the functor  $\mathrm{Isom}_{\xi} : \underline{C}_{R_{0} \hat{\otimes}_{\wedge} R_{0}} \to (\mathrm{Ens})$  given by

\begin{equation*}
T \longmapsto \mathrm{Isom}_{\underline{A}(T)}(\xi_{1}(T),\xi_{2}(T))
\end{equation*}
where  $\xi_{i}(T) = \xi(f_{i})$  and  $f_{i} : R_{0} \to T$  denotes the composit maps

}{\begin{equation*}
\begin{aligned}
[\underline{h}(k[\epsilon])] &= \mathrm{Coker}\ (\mathrm{Hom}_{\wedge\text{-alg}}(R_{1},k[\epsilon]) \ \underset{s_{2}}{\overset{s_{1}}{\rightrightarrows}}\ \mathrm{Hom}_{\wedge\text{-alg}}(R_{0},k[\epsilon])) \qquad \text{as well as} \\[2ex]
\mathrm{Aut}_{\underline{h}}(1_{\epsilon}) &= \mathrm{Ker}(\mathrm{Hom}_{\wedge\text{-alg}}(R_{1},k[\epsilon]) \xrightarrow{\ s_{1}\ } \mathrm{Hom}_{\wedge\text{-alg}}(R_{0},k[\epsilon])) \\[2ex]
&= \mathrm{Ker}(\mathrm{Hom}_{\wedge\text{-alg}}(R_{1},k[\epsilon]) \xrightarrow{\ s_{2}\ } \mathrm{Hom}_{\wedge\text{-alg}}(R_{0},k[\epsilon])) \qquad \text{are both finite-dimensional}
\end{aligned}
\end{equation*}
vector spaces over  $k$  , we must have the same properties on  $\underline{A}$  . Conversely assume (1) and (2) on  $\underline{A}$  . We claim that  $\underline{A}$  is prorepresentable by a normalized smooth cogroupoid. Since  $\underline{A}$  is homogeneous and  $[\underline{A}(k[\epsilon])]$  is a finite-dimensional vector space over  $k$  ,  $\underline{A}$  admits a minimally versal formal object  $\xi$  . Set  $P(\xi) = R_{0}$  and consider the functor  $\mathrm{Isom}_{\xi} : \underline{C}_{R_{0} \hat{\otimes}_{\wedge} R_{0}} \to (\mathrm{Ens})$  given by

\begin{equation*}
T \longmapsto \mathrm{Isom}_{\underline{A}(T)}(\xi_{1}(T),\xi_{2}(T))
\end{equation*}
where  $\xi_{i}(T) = \xi(f_{i})$  and  $f_{i} : R_{0} \to T$  denotes the composite maps

}{F7I-VI-039}\begin{center}
\begin{tikzcd}[column sep=large]
R_0 \arrow[r, shift left=1.2ex, "s_1"] \arrow[r, shift right=1.2ex, "s_2"'] \& R_0 \otimes_{\wedge} R_0 \arrow[r] \& T \ .
\end{tikzcd}
\end{center}









\SGAvar{This functor is pro-representable by 2.16, i.e. there exists  $R_{1} \in \mathrm{ob}(\hat{\underline{C}}_{R_{0} \hat{\otimes}_{\wedge} R_{0}})$  and an isomorphism  $\xi_{1}(R_{1}) \xrightarrow{\ \varphi\ } \xi_{2}(R_{1})$  such that  $\Phi : h_{R_{1}} \to \mathrm{Isom}_{\xi}$  is an isomorphism over  $\underline{C}_{R_{0} \hat{\otimes}_{\wedge} R_{0}}$  . Then  $(R_{0},R_{1},\ldots)$  defines a cogroupoid in  $\hat{\underline{C}}_{\wedge}$  and the functor $(\xi,\varphi) : \underline{h}_{(R_{0},R_{1},\ldots)} \to \underline{A}$  is a  $\underline{C}_{\wedge}$-equivalence. Since  $\xi$  is a minimally versal formal object of  $\underline{A}$  , the composit functor  $\underline{h}_{(R_{0},R_{0},\ldots)} \xrightarrow{\ \epsilon\ } \underline{h}_{(R_{0},R_{1},\ldots)} \xrightarrow{\ (\xi,\varphi)\ } \underline{A}$ is, by definition, minimally smooth and therefore so is  $\epsilon$  since  $(\xi,\varphi)$  is a $\underline{C}_{\wedge}$-equivalence i.e.,  $(R_{0},R_{1},\ldots)$  is a normalized smooth cogroupoid in  $\hat{\underline{C}}_{\wedge}$  .


}{This functor is pro-representable by 2.16, i.e. there exists  $R_{1} \in \mathrm{ob}(\hat{\underline{C}}_{R_{0} \hat{\otimes}_{\wedge} R_{0}})$  and an isomorphism  $\xi_{1}(R_{1}) \xrightarrow{\ \varphi\ } \xi_{2}(R_{1})$  such that  $\Phi : h_{R_{1}} \to \mathrm{Isom}_{\xi}$  is an isomorphism over  $\underline{C}_{R_{0} \hat{\otimes}_{\wedge} R_{0}}$  . Then  $(R_{0},R_{1},\ldots)$  defines a cogroupoid in  $\hat{\underline{C}}_{\wedge}$  and the functor $(\xi,\varphi) : \underline{h}_{(R_{0},R_{1},\ldots)} \to \underline{A}$  is a  $\underline{C}_{\wedge}$-equivalence. Since  $\xi$  is a minimally versal formal object of  $\underline{A}$  , the composite functor  $\underline{h}_{(R_{0},R_{0},\ldots)} \xrightarrow{\ \epsilon\ } \underline{h}_{(R_{0},R_{1},\ldots)} \xrightarrow{\ (\xi,\varphi)\ } \underline{A}$ is, by definition, minimally smooth and therefore so is  $\epsilon$  since  $(\xi,\varphi)$  is a $\underline{C}_{\wedge}$-equivalence i.e.,  $(R_{0},R_{1},\ldots)$  is a normalized smooth cogroupoid in  $\hat{\underline{C}}_{\wedge}$  .


}{F7I-VI-040}\SGApage{I}{69}

\underline{Corollary} 2.18.  \underline{A functor}  $F : \underline{C}_\Lambda \longrightarrow$ (Ens)  \underline{is pro-representable if and only if it is homogeneous and}  $\dim F(k[\epsilon]) < \infty$  .

\underline{Corollary} 2.19.  \underline{Every smooth cogroupoid in}  $\hat{\underline{C}}_\Lambda$  \underline{can be normalized up to isomorphism.}

Given a cogroupoid  $(R_0,R_1,\dots)$  in  $\hat{\underline{C}}_\Lambda$ , we set  $\bar{R}_1 = R_1/J$  where  $J$  is the ideal in  $R_1$  generated by  $\{s_1(r) - s_2(r) \mid r \in R_0\}$ . Then  $(R_0,\bar{R}_1,\dots)$  defines a formal group scheme over  $R_0$  with respect to the induced composition rule. Now let  $\underline{A}$  be a cofibered groupoid over  $\underline{C}_\Lambda$  prorepresented by a normalized smooth cogroupoid  $(R_0,R_1,\dots)$  via  $(\xi,\varphi)$ . Let  $\underline{C}_k \hookrightarrow \underline{C}_{R_0} \hookrightarrow \underline{C}_{R_0 \hat{\otimes}_\Lambda R_0}$  be the embeddings corresponding to the surjective maps  $R_0 \hat{\otimes}_\Lambda R_0 \longrightarrow R_0 \longrightarrow k$ , and let  $\mathrm{Aut}_\xi$ , $\mathrm{Aut}^o$  be the restriction functors of  $\mathrm{Isom}_\xi$  on  $\underline{C}_{R_0}$ , $\underline{C}_k$  respectively. In other words,  $\mathrm{Aut}_\xi(S)$ $=$ the automorphism group of  $\xi(S)$  for every  $S \in \mathrm{ob}(\underline{C}_{R_0})$ , and  $\mathrm{Aut}^o(S)$ $=$ the automorphism group of  $1_S$  for every  $S \in \mathrm{ob}(\underline{C}_k)$ , where  $1_S$ $=$ the direct image of  $1 \in \underline{A}(k)$  under the map  $k \to S$ . Since  $R_1$  prorepresents the functor  $\mathrm{Isom}_\xi$  on  $\underline{C}_{R_0 \hat{\otimes}_\Lambda R_0}$ , it follows that  $\bar{R}_1$  prorepresents the functor  $\mathrm{Aut}_\xi$  on  $\underline{C}_{R_0}$ , and  $k \otimes_{R_0} \bar{R}_1$  prorepresents  $\mathrm{Aut}^o$  on  $\underline{C}_k$ . Though  $\underline{A}$  is prorepresented by a smooth cogroupoid  $(R_0,R_1,\dots)$  , neither  $(R_0,\bar{R}_1,\dots)$  nor  $(k,\ k \otimes_{R_0} \bar{R}_1,\dots)$  are in general smooth. Indeed we have

\SGAvar{\underline{Corollary} 2.20.  \underline{Let}  $\underline{A}$  \underline{be a homogeneous cogroupoid over}  $\underline{C}_\Lambda$  \underline{prorepresented by a normalized smooth cogroupoid}  $(R_0,R_1,\dots)$  \underline{via}  $(\xi,\varphi)$  .

}{\underline{Corollary} 2.20.  \underline{Let}  $\underline{A}$  \underline{be a homogeneous cofibered groupoid over}  $\underline{C}_\Lambda$  \underline{prorepresented by a normalized smooth cogroupoid}  $(R_0,R_1,\dots)$  \underline{via}  $(\xi,\varphi)$  .

}{F7I-VI-041}(a)  $(R_0,\bar{R}_1,\dots)$  \underline{and}  $(k,\ k \otimes_{R_0} \bar{R}_1,\dots)$  \underline{prorepresent the functors}  $\mathrm{Aut}_\xi$  \underline{and}  $\mathrm{Aut}^o$  \underline{respectively.}

(b)  \underline{The following statements are equivalent}

(i)   \underline{the formal group scheme}  $\underline{h}_{(R_0,\bar{R}_1,\dots)}$  \underline{is formally smooth over}  $R_0$  ;

(ii)  \underline{for every surjective arrow}  $a' \to a$  \underline{in}  $\underline{A}$  ,
$\mathrm{Aut}_{P(a')}(a') \longrightarrow \mathrm{Aut}_{P(a)}(a)$  \underline{is surjective},

(iii)  $R_0$  \underline{prorepresents the functor}  $[\underline{A}]$  \underline{over}  $\underline{C}_\Lambda$  .


\SGApage{I}{70}

\underline{Proof}:  We have already noted (a). As for (b),  (i) $\Leftrightarrow$ (ii)  is clear whereas (ii) $\Leftrightarrow$ (iii)  follows from 2.7 and 2.15, since  $\tilde{\xi} : h_{R_0} \longrightarrow [\underline{A}]$  is minimally smooth.

\SGAvar{\underline{Remark} 2.21.  Let  $\underline{A}$  be a cofibered groupoid over  $\underline{C}_\Lambda$  , and let  $\underline{A}|\underline{C}_k$  be the pull-back of  $p : \underline{A} \longrightarrow \underline{C}_\Lambda$  under  $\underline{C}_k \hookrightarrow \underline{C}_\Lambda$ . Then it is trivial to see that if  $\underline{A}$  is prorepresented by  $(R_0,R_1,\dots)$  via  $(\xi,\varphi)$  , then  $\underline{A}|\underline{C}_k$  is prorepresented by  $(k \otimes_\Lambda R_0\ ,\ k \otimes_\Lambda R_1,\dots)$  via  $(k \otimes_\Lambda \xi\ ,\ k \otimes_\Lambda \varphi)$  where  $k \otimes_\Lambda \xi$ $=$ its direct image of  $\xi$  under the map  $R_0 \longrightarrow k \otimes_\Lambda R_0$ ,  $k \otimes_\Lambda \varphi$ $=$ the image of  $\varphi$  under  $\mathrm{Isom}_{\underline{A}(R_0)}(\xi(s_1),\ \xi(s_2)) \longrightarrow \mathrm{Isom}_{\underline{A}(k \otimes_\Lambda R_1)}(k \otimes_\Lambda \xi(s_1),\ k \otimes_\Lambda \xi(s_2))$ . If  $(R_0,R_1,\dots)$  is normalized, then so is  $(k \otimes_\Lambda R_0,\ k \otimes_\Lambda R_1,\dots)$  , and the relative dimension of  $R_1$  over  $R_0$  is equal to the relative dimension of  $k \otimes_\Lambda R_1$  over  $k \otimes_\Lambda R_0$  since  $R_1$  is smooth over  $R_0$  .

\underline{Remark} 2.22.  Let  $\underline{A}$  be prorepresented by smooth cogroupoid  $(R_0,R_1,\dots)$  and  $(R'_0,\ R'_1,\dots)$  where we assume that  $(R_0,R_1,\dots)$  is normalized. We then obtain a homomorphism  $(i_0,i_1) : (R_0,R_1,\dots) \longrightarrow (R'_0,\ R'_1,\dots)$  . Now  $R'_0$  is a formal power-series ring over  $R_0$  by 2.10(b), and the homomorphism  $(i_0,i_1)$  induces an isomorphism  $R'_0 \otimes_{R_0} (R_0,R_1,\dots) \overset{\sim}{\longrightarrow} (R'_0,\ R'_1,\dots)$  .


}{\underline{Remark} 2.21.  Let  $\underline{A}$  be a cofibered groupoid over  $\underline{C}_\Lambda$  , and let  $\underline{A}|\underline{C}_k$  be the pull-back of  $p : \underline{A} \longrightarrow \underline{C}_\Lambda$  under  $\underline{C}_k \hookrightarrow \underline{C}_\Lambda$ . Then it is trivial to see that if  $\underline{A}$  is prorepresented by  $(R_0,R_1,\dots)$  via  $(\xi,\varphi)$  , then  $\underline{A}|\underline{C}_k$  is prorepresented by  $(k \otimes_\Lambda R_0\ ,\ k \otimes_\Lambda R_1,\dots)$  via  $(k \otimes_\Lambda \xi\ ,\ k \otimes_\Lambda \varphi)$  where  $k \otimes_\Lambda \xi$ $=$ the direct image of  $\xi$  under the map  $R_0 \longrightarrow k \otimes_\Lambda R_0$ ,  $k \otimes_\Lambda \varphi$ $=$ the image of  $\varphi$  under  $\mathrm{Isom}_{\underline{A}(R_1)}(\xi(s_1),\ \xi(s_2)) \longrightarrow \mathrm{Isom}_{\underline{A}(k \otimes_\Lambda R_1)}(k \otimes_\Lambda \xi(s_1),\ k \otimes_\Lambda \xi(s_2))$ . If  $(R_0,R_1,\dots)$  is normalized, then so is  $(k \otimes_\Lambda R_0,\ k \otimes_\Lambda R_1,\dots)$  , and the relative dimension of  $R_1$  over  $R_0$  is equal to the relative dimension of  $k \otimes_\Lambda R_1$  over  $k \otimes_\Lambda R_0$  since  $R_1$  is smooth over  $R_0$  .

\underline{Remark} 2.22.  Let  $\underline{A}$  be prorepresented by smooth cogroupoids  $(R_0,R_1,\dots)$  and  $(R'_0,\ R'_1,\dots)$  where we assume that  $(R_0,R_1,\dots)$  is normalized. We then obtain a homomorphism  $(i_0,i_1) : (R_0,R_1,\dots) \longrightarrow (R'_0,\ R'_1,\dots)$  . Now  $R'_0$  is a formal power-series ring over  $R_0$  by 2.10(b), and the homomorphism  $(i_0,i_1)$  induces an isomorphism  $R'_0 \otimes_{R_0} (R_0,R_1,\dots) \overset{\sim}{\longrightarrow} (R'_0,\ R'_1,\dots)$  .


}{F7I-VI-042}\SGAsectiontoc{3.  The coherent sheaves  $\protect\underline{D}^i_X$}{3}
3.  \underline{The coherent sheaves}  $\underline{D}^i_X$

One of the nice functors which is usually not pro-representable but nevertheless admits a formal versal object is the deformation functor. To this end we recall a few basic facts on commutative ring extensions and its obstruction theory. This section contains nothing new (possibly except in its systematic approach via Grothendieck cohomology on site). A purpose of this section is to make this expose as much self-contained. (Indeed we prove here more than we need later). The basic facts are stated in theorem 3.12. Throughout rings are meant to be commutative rings with unit.


\SGApage{I}{71}

Given a ring  $R$  we denote by  $\underline{C}_R$  the category of $R$-algebras. For each object  $A$  in  $\underline{C}_R$  we set

Cov($A$) $=$ the set of all surjective arrows  $\varphi : A' \to A$  in  $\underline{C}_R$  with target  $A$ .

It clearly defines a pretopology and in turn we obtain a site  $\underline{C}_R$ . For each object  $A$  in  $\underline{C}_R$ , we shall denote by  $\underline{C}_{A|R}$  the category of arrows in  $\underline{C}_R$  with target  $A$ . We may and shall provide the category  $\underline{C}_{A|R}$  with the topology induced from the forgetful functor  $j_A : \underline{C}_{A|R} \to \underline{C}_R$ . We note that the category  $\underline{C}_{A|R}$  has a final object  $(1_A : A \to A)$  and an initial object  $(i_A : R \to A)$ , and admits a fibre-product as well as coproduct. Henceforth we shall denote the object  $(B \to A)$  in  $\underline{C}_{A|R}$ , by abuse of notation, simply by  $B$  in case when there is no possible confusion. Given an abelian category  $\underline{C}$ , a  $\underline{C}$-valued sheaf on  $\underline{C}_{A|R}$  is, by definition, a functor  $F : \underline{C}^{O}_{A|R} \to \underline{C}$  such that, for any surjective arrow  $B' \to B$  in  $\underline{C}_{A|R}$ ,

\begin{equation*}
0 \to F(B) \to F(B') \rightrightarrows F(B' \underset{B}{\times} B')
\end{equation*}
is exact. The  \underline{additive}  presheaves are those which preserve coproduct i.e.,  $F(B_1 \underset{R}{\otimes} B_2) = F(B_1) \oplus F(B_2)$ . For example, an $A$-module  $M$  defines the functor

\begin{equation*}
\mathrm{Der}_R(-,M): \underline{C}^{O}_{A|R} \to (A\text{-mod})
\end{equation*}
given by  $B \to \mathrm{Der}_R(B,M) =$ $R$-linear derivations of  $B$  with values in  $M$  regarded as $B$-module via the structural map  $B \to A$ , and it is clearly an additive sheaf on  $\underline{C}^{O}_{A|R}$ , where ($A$-mod) denotes the category of $A$-modules. A presheaf  $F$  on  $\underline{C}_{A|R}$  with values in  ($A$-mod)  is called "finite type" if  $F(B)$  is an $A$-module of finite type whenever  $B$  is an $R$-algebra of finite type. For example, if  $A$  is noetherian and  $M$  is an $A$-module of finite type, then  $\mathrm{Der}_R(-,M)$  is an additive sheaf of finite type on the site  $\underline{C}_{A|R}$ .

A few more words on notations and basic facts in homological algebra:

3.1(a)  Throughout the rest of this section, we shall denote by  $\widetilde{\underline{C}}_{A|R}$  (or  $\overset{\vee}{\underline{C}}_{A|R}$ )  the category of sheaves (or presheaves) with values in ($A$-mod). We have the (left exact) injection functor  $i_{A|R}: \widetilde{\underline{C}}_{A|R} \to \overset{\vee}{\underline{C}}_{A|R}$ , and the associated-sheaf (exact)
\SGApage{I}{72}
functor  $G_{A|R}: \overset{\vee}{\underline{C}}_{A|R} \to \widetilde{\underline{C}}_{A|R}$  which is adjoint to the functor  $i_{A|R}$ . An object  $B \in ob(\underline{C}_{A|R})$  induces a continuous functor  $j : \underline{C}_{B|R} \to \underline{C}_{A|R}$  and in turn a commutative diagram

\begin{equation}\tag{*}
\begin{tikzcd}[column sep=huge, row sep=huge]
\widetilde{\underline{C}}_{A|R} \arrow[r, "\widetilde{j}"] \arrow[d, "i_{A|R}"'] \& \widetilde{\underline{C}}_{B|R} \arrow[d] \\
\overset{\vee}{\underline{C}}_{A|R} \arrow[r, "j^{\vee}"'] \& \overset{\vee}{\underline{C}}_{B|R}
\end{tikzcd}
\end{equation}
















where the restriction functor  $\widetilde{j}$  (as well as  $j^{\vee}$ )  is an exact functor, since the site  $\underline{C}_{B|R}$  is nothing but the category of arrows in  $\underline{C}_{A|R}$  having target  $B$ , with the induced topology. A ring homomorphism  $S \to R$  induces a continuous functor  $\beta : \underline{C}_{A|R} \to \underline{C}_{A|S}$ , and in turn a commutative diagram

\begin{equation}\tag{**}
\begin{tikzcd}[column sep=huge, row sep=huge]
\widetilde{\underline{C}}_{A|S} \arrow[r, "\widetilde{\beta}"] \arrow[d, "i_{A|S}"'] \& \widetilde{\underline{C}}_{A|R} \arrow[d] \\
\overset{\vee}{\underline{C}}_{A|S} \arrow[r, "\beta^{\vee}"'] \& \overset{\vee}{\underline{C}}_{A|R}
\end{tikzcd}
\end{equation}















Note that  $\widetilde{\beta}$  is not necessarily an exact functor.

3.1(b)  The right derived functor of the injection functor  $i_{A|R}: \widetilde{\underline{C}}_{A|R} \to \overset{\vee}{\underline{C}}_{A|R}$  will be denoted by  $\underline{H}^{\circ}_{A|R}$ . For each sheaf  $F$  on  $\underline{C}_{A|R}$ , we set

\SGAvar{\begin{equation*}
H^{\circ}(A|R,F) = [H^{\circ}_{A|R}(F)](A)
\end{equation*}
For any  $B' \to B$  in  Cov($B$) , let us define the functor  $H^{\circ}((B' \to B),-) : \overset{\vee}{\underline{C}}_{A|R} \to (A\text{-mod})$  by  $H^{\circ}((B'\to B),F) = Ker(F(B') \rightrightarrows F(B'\underset{B}{\times} B'))$ . The right derived functors of  $H^{\circ}((B' \to B),-)$  will be denoted by  $H^{\circ}((B'\to B),-)$ . One knows that the  $H^{\circ}((B' \to B),F)$  are the cohomology groups of the (Čech)
\SGApage{I}{73}
semi-simplicial module

}{\begin{equation*}
H^{\circ}(A|R,F) = [H^{\circ}_{A|R}(F)](A)
\end{equation*}
For any  $B' \to B$  in  Cov($B$) , let us define the functor  $H^{0}((B' \to B),-) : \overset{\vee}{\underline{C}}_{A|R} \to (A\text{-mod})$  by  $H^{0}((B'\to B),F) = Ker(F(B') \rightrightarrows F(B'\underset{B}{\times} B'))$ . The right derived functors of  $H^{0}((B' \to B),-)$  will be denoted by  $H^{\circ}((B'\to B),-)$ . One knows that the  $H^{\circ}((B' \to B),F)$  are the cohomology groups of the (Čech)
\SGApage{I}{73}
semi-simplicial module

}{F7I-VI-043}\begin{center}
\begin{tikzcd}[column sep=large]
F(B') \arrow[r, shift left=2] \arrow[r, shift right=2] \& F(B'\times_B B') \arrow[r, shift left=3] \arrow[r] \arrow[r, shift right=3] \& F(B'\ \times_B B'\ \times_B B') \arrow[r, shift left=4.5] \arrow[r, shift left=1.5] \arrow[r, shift right=1.5] \arrow[r, shift right=4.5] \& \cdots
\end{tikzcd}
\end{center}








\SGAvar{The right derived functor of  $\check{\underline{H}}^{0}_{A|R} : \underline{C}_{A|R} \to \underline{C}_{A|R}$  will be denoted by  $\check{\underline{H}}^{\bullet}_{A|R}$ ,

}{The right derived functor of  $\check{\underline{H}}^{0}_{A|R} : \overset{\vee}{\underline{C}}_{A|R} \to \overset{\vee}{\underline{C}}_{A|R}$  will be denoted by  $\check{\underline{H}}^{\bullet}_{A|R}$ ,

}{F7I-VI-044}where

\begin{equation*}
[\check{\underline{H}}^{0}_{A|R}(F)](B) \;=\; \lim_{\{B' \to B\}\, \in\, \mathrm{Cov}(B)} H^{0}(\{B' \to B\},F) \ .
\end{equation*}
For each presheaf  $F$  on  $\underline{C}_{A|R}$  , we set

\begin{equation*}
\check{H}^{\bullet}(A|R,F) \;=\; [\check{H}^{\bullet}_{A|R}(F)](A) \;=\; \lim_{\{A' \to A\}\, \in\, \mathrm{Cov}(A)} H^{\bullet}(\{A' \to A\},F)
\end{equation*}
for any sheaf  $F$  on  $\underline{C}_{A|R}$  , we have the usual spectral sequence

\begin{equation*}
\check{H}^{p}(A|R,\underline{H}^{q}_{A|R}(F)) \;\Rightarrow\; H^{p+q}(A|R,F)
\end{equation*}
where we note that  $\check{H}^{0}(A|R,\underline{H}^{q}(F)) = 0$  for all  $q > 0$ .

In particular

\SGAvar{\begin{equation*}
\left\{
\begin{array}{l}
H^{1}(A|R,F) \;=\; H^{1}(A|R,F) \\[2ex]
0 \to \check{H}^{2}(A|R,F) \to H^{2}(A|R,F) \to \check{H}^{1}(A|R,\underline{H}^{1}_{A\ R}(F)) \to \check{H}^{3}(A|R,F) \quad \text{is exact.}
\end{array}
\right.
\end{equation*}
3.1(c)   If  $B$  is an object in  $\underline{C}_{A|R}$  , the commutative diagram (*) in 2.1(a)

}{\begin{equation*}
\left\{
\begin{array}{l}
\check{H}^{1}(A|R,F) \;=\; H^{1}(A|R,F) \\[2ex]
0 \to \check{H}^{2}(A|R,F) \to H^{2}(A|R,F) \to \check{H}^{1}(A|R,\underline{H}^{1}_{A|R}(F)) \to \check{H}^{3}(A|R,F) \quad \text{is exact.}
\end{array}
\right.
\end{equation*}
3.1(c)   If  $B$  is an object in  $\underline{C}_{A|R}$  , the commutative diagram (*) in 3.1(a)

}{F7I-VI-045}together with the fact that  $\tilde{j}$  as well as  $j^{\vee}$  are exact functors yields the

canonical isomorphism

\begin{equation*}
H^{\bullet}(B|R,\tilde{j}F) \;=\; [\underline{H}^{\bullet}_{A|R}(F)](B)
\end{equation*}
On the other hand, a ring homomorphism  $S → R$  induces the commutative diagram (**)

\SGAvar{in 2.1(a). Now  $\tilde{\beta}$  is not necessarily an exact functor. However,  $β$  preserves a

}{in 3.1(a). Now  $\tilde{\beta}$  is not necessarily an exact functor. However,  $β$  preserves a

}{F7I-VI-046}fibre-product, and thus we obtain a spectral sequence

\begin{equation*}
\underline{H}^{p}_{A|R}(R^{q}\tilde{\beta}F) \;\Rightarrow\; \check{\beta}\underline{H}^{p+q}_{A|S}(F) \qquad \text{i.e.,} \qquad H^{p}(A|R,R^{q}\tilde{\beta}F) \;\Rightarrow\; H^{p+q}(A|S,F)
\end{equation*}
Since we have the canonical isomorphism  $\tilde{\beta} = G_{A|R}\circ\check{\beta}\circ i_{A|S}$  and  $\beta^{\#} = G_{A|R}\circ\check{\beta}$  is an
\SGApage{I}{74}
exact functor, we have  $R^{\bullet}\tilde{\beta}F = \beta^{\#}\underline{H}^{\bullet}_{A|S}(F)$  . Thus the spectral sequence above can

now be written as

\begin{equation*}
H^{p}(A|R,\beta^{\#}\underline{H}^{q}_{A|S}(F)) \;\Rightarrow\; H^{p+q}(A|S,F) \ .
\end{equation*}
\underline{Lemma} 3.2.  \underline{Given}  $B \in ob(\underline{C}_{A|R})$  , \underline{choose a surjective arrow}  $P → B$  in  $\underline{C}_{B|R}$  ,

\underline{where}  $P$  \underline{is a polynomial algebra over}  $R$  . \underline{Then for any}  $F \in ob(\check{\underline{C}}_{A|R})$  , \underline{the}

\underline{natural map}  $H^{\bullet}(\{P \to B\},F) \to \check{\underline{H}}^{\bullet}(F)\ (B)$  \underline{is an isomorphism.}

\SGAvar{\underline{Proof}:  It suffice to see that the natural map

}{\underline{Proof}:  It suffices to see that the natural map

}{F7I-VI-047}\begin{equation*}
H^{0}(\{P \to B\},F) \to \check{\underline{H}}^{0}(F)\ (B) \;=\; \varinjlim_{\{B' \to B\}\, \in\, \mathrm{Cov}(B)} H^{0}(\{B' \to B\},F)
\end{equation*}
is an isomorphism, which is clear since, for any  $\{B' \to B\} \in \mathrm{Cov}(B)$  , there exists

an arrow  $\{P \to B\} \to \{B' \to B\}$  .


\underline{Corollary} 3.3.

(a)  \underline{Assume that}  $A$  \underline{is noetherian}. \underline{If}  $F \in ob(\widetilde{\underline{C}}_{A|R})$  \underline{is of finite type,}

\underline{then}  $\underline{H}^{n}(F)$  \underline{is of finite type for all}  $n$  .

(b)  \underline{For any polynomial algebra}  $P$  \underline{over}  $R$  , \underline{we have}  $H^{n}(P|R,F) = 0$

\underline{for all}  $n > 0$  \underline{and for all}  $F \in ob(\widetilde{\underline{C}}_{P|R})$  .

(c)  \underline{A diagram}  $0 \to F' \to F \to F'' \to 0$  \underline{in}  $\widetilde{\underline{C}}_{A|R}$  \underline{is exact if and only if}

$0 \to F'(P) \to F(P) \to F''(P) \to 0$  \underline{is exact for every}  $P \in ob(\underline{C}_{A|R})$  \underline{which is a polynomial algebra over}  $R$  .

\underline{Proof}:

(a)  In view of the spectral sequence

\begin{equation*}
\check{H}^{p}(B|R,\underline{H}^{q}(F)) \;\Rightarrow\; H^{p+q}(B|R,F)
\end{equation*}
together with the fact that  $\check{H}^{0}(\underline{H}^{q}(F)) = 0$  for all  $q > 0$  , it suffices to

\SGAvar{show that  $\underline{H}^{n}(G)$  is of finite type for all  $n$  and for all presheaf  $G$  of finite

}{show that  $\underline{H}^{n}(G)$  is of finite type for all  $n$  and for all presheaves  $G$  of finite

}{F7I-VI-048}type. However  $\check{H}^{\bullet}(B|R,G) = H^{\bullet}(\{P \to B\},G)$  , where  $P$  is a polynomial algebra over

$R$  and  $\{P \to B\} \in \mathrm{Cov}(B)$  , and furthermore  $H^{\bullet}(\{P \to B\},G)$  is simply the
\SGApage{I}{75}
cohomology gotten from the complex

\begin{multline*}
\dots\dots\dots\dots\dots G(P \underset{B}{\times} P \underset{B}{\times} P) \;\substack{\longleftarrow\\ \longleftarrow\\ \longleftarrow}\; G(P \underset{B}{\times} P) \substack{\longleftarrow\\ \longleftarrow}\; G(P) \ .
\end{multline*}
If  $B$  is an R-algebra of finite type, then we may take  $P$  to be of finite type over  $R$  so that  $P\underset{B}{\times}\ ....\ \underset{B}{\times}P$  are all R-algebras of finite type. If  $G$  is of finite type, then  $G(P\underset{B}{\times}\ ....\ \underset{B}{\times}P)$  are all A-modules of finite type, so that  $H^{\bullet}(\{P \to B\},G)$  is an A-module of finite type since  $A$  is noetherian.

\SGAvar{(b)  Let  $P$  be a polynomial algebra over  $R$ . If  $B \to P$  is any surjective R-algebra map, then it amits a section and hence  $H^{i}(\{B \to P\},G) = 0$  for all  $i > 0$ , so that  $\check{H}^{i}(P|R,\ G) = 0$  for all  $i > 0$  and for all presheaves  $G$  on  $\underline{C}_{P|R}$ . It follows that the spectral sequence  $\check{H}^{p}(P|R,\ \underline{H}^{q}_{P|R}(F)) \Rightarrow H^{p+q}(P|R,\ F)$  degenerates and hence  $\check{H}^{0}(P|R,\ \underline{H}^{n}(F)) \simeq H^{n}(P|R,\ F)$  for all  $n$ . However  $\check{H}^{0}(P|R,\ \underline{H}^{n}(F)) = 0$  for all  $n > 0$  whenever  $F$  is a sheaf, and hence  $H^{n}(P|R,\ F) = 0$  for all  $n > 0$  and for all  $F \in ob(\widetilde{\underline{C}}_{P|R})$ .

}{(b)  Let  $P$  be a polynomial algebra over  $R$ . If  $B \to P$  is any surjective R-algebra map, then it admits a section and hence  $H^{i}(\{B \to P\},G) = 0$  for all  $i > 0$ , so that  $\check{H}^{i}(P|R,\ G) = 0$  for all  $i > 0$  and for all presheaves  $G$  on  $\underline{C}_{P|R}$ . It follows that the spectral sequence  $\check{H}^{p}(P|R,\ \underline{H}^{q}_{P|R}(F)) \Rightarrow H^{p+q}(P|R,\ F)$  degenerates and hence  $\check{H}^{0}(P|R,\ \underline{H}^{n}(F)) \simeq H^{n}(P|R,\ F)$  for all  $n$ . However  $\check{H}^{0}(P|R,\ \underline{H}^{n}(F)) = 0$  for all  $n > 0$  whenever  $F$  is a sheaf, and hence  $H^{n}(P|R,\ F) = 0$  for all  $n > 0$  and for all  $F \in ob(\widetilde{\underline{C}}_{P|R})$ .

}{F7I-VI-049}(c)  If  $0 \to F' \to F \to F'' \to 0$  is an exact sequence in  $\widetilde{\underline{C}}_{A|R}$ , then

\SGAvar{\begin{equation*}
0 \to F'(B) \to F(B) \to F''(B) \to H'(B|R,\ F') \to H'(B|R,\ F) \to \ ....
\end{equation*}
is exact. If  $P \in ob(\underline{C}_{A|R})$  is a polynomial algebra over  $R$ , then  $H^{1}(P|R,\ F') = 0$  by 3.3(b) and hence

\begin{equation*}
0 \to F' \to F \to F'' \to 0
\end{equation*}
is exact. As for the other direction, let  $\underline{P}_{A|R}$  be the full subcategory of  $\underline{C}_{A|R}$  in which

}{\begin{equation*}
0 \to F'(B) \to F(B) \to F''(B) \to H^{1}(B|R,\ F') \to H^{1}(B|R,\ F) \to \ ....
\end{equation*}
is exact. If  $P \in ob(\underline{C}_{A|R})$  is a polynomial algebra over  $R$ , then  $H^{1}(P|R,\ F') = 0$  by 3.3(b) and hence

\begin{equation*}
0 \to F'(P) \to F(P) \to F''(P) \to 0
\end{equation*}
is exact. As for the other direction, let  $\underline{P}_{A|R}$  be the full subcategory of  $\underline{C}_{A|R}$  in which

}{F7I-VI-050}\begin{equation*}
ob(\underline{P}_{A|R}) = \{P \in ob(\underline{C}_{A|R})|\ \ P \text{ is a polynomial algebra over } R\} \ ,
\end{equation*}
provided with the induced topology, and let  $p : \underline{P}_{A|R} \to \underline{C}_{A|R}$  be the inclusion functor. Since every object in  $\underline{C}_{A|R}$  can be covered by an object in  $\underline{P}_{A|R}$ , it follows from a comparison lemma (          ) that  $\widetilde{p} : \widetilde{\underline{C}}_{A|R} \to \widetilde{\underline{P}}_{A|R}$  is an equivalence of categories. Consequently, if  $0 \to F'(P) \to F(P) \to F''(P) \to 0$  is exact for every  $P \in ob(\underline{P}_{A|R})$ , then  $0 \to \widetilde{p}F' \to \widetilde{p}F \to \widetilde{p}F'' \to 0$  is exact in  $\widetilde{\underline{P}}_{A|R}$  and hence  $0 \to F' \to F \to F'' \to 0$  is exact in  $\widetilde{\underline{C}}_{A|R}$  .


\SGApage{I}{76}

\SGAvar{\underline{Remark} 3.4.  Every surjective arrow in  $\underline{P}_{A|R}$  admits a section. It follows immediately that every presheaf on  $\underline{P}_{A|R}$  is actually a sheaf i.e.,  $\widetilde{\underline{P}}_{A|R} = \underline{P}_{A|R}$ . Consequently we find that the restriction functor  $\widetilde{\underline{C}}_{A|R} \to \check{\underline{P}}_{A|R}$  is an equivalence of categories.

}{\underline{Remark} 3.4.  Every surjective arrow in  $\underline{P}_{A|R}$  admits a section. It follows immediately that every presheaf on  $\underline{P}_{A|R}$  is actually a sheaf i.e.,  $\widetilde{\underline{P}}_{A|R} = \check{\underline{P}}_{A|R}$ . Consequently we find that the restriction functor  $\widetilde{\underline{C}}_{A|R} \to \check{\underline{P}}_{A|R}$  is an equivalence of categories.

}{F7I-VI-051}\underline{Lemma} 3.5.  \underline{Consider a diagram of R-algebras}

\begin{center}
\begin{tikzcd}
 \& C \arrow[d] \\
B' \arrow[r, "\text{surjective}"'] \& B
\end{tikzcd}
\end{center}

\underline{and let}  $S$  \underline{be an R-algebra.}

(i)  \underline{If}  $\mathrm{Tor}^{R}_{1}(B,S) = 0$  , \underline{then the canonical map}

\begin{equation*}
(B'\underset{B}{\times} C) \underset{R}{\otimes} S \to (B'\underset{R}{\otimes} S) \underset{(B \underset{R}{\otimes} S)}{\times} (C \underset{R}{\otimes} S)
\end{equation*}
\underline{is an isomorphism.}

\SGAvar{(ii)  \underline{Assume that}  $B'$  \underline{is R-flat.} \underline{If}  $\mathrm{Tor}^{R}_{i}(C,S) = 0$  \underline{for all}  $0 < i < n$  \underline{and}  $\mathrm{Tor}^{R}_{i}(B,S) = 0$  \underline{for all}  $0 < i < n$  , \underline{then}  $\mathrm{Tor}^{R}_{i}(B'\underset{B}{\times}C,\ S) = 0$  \underline{for all}  $0 < i < n$  .

}{(ii)  \underline{Assume that}  $B'$  \underline{is R-flat.} \underline{If}  $\mathrm{Tor}^{R}_{i}(C,S) = 0$  \underline{for all}  $0 < i < n$  \underline{and}  $\mathrm{Tor}^{R}_{i}(B,S) = 0$  \underline{for all}  $0 < i \leq n$  , \underline{then}  $\mathrm{Tor}^{R}_{i}(B'\underset{B}{\times}C,\ S) = 0$  \underline{for all}  $0 < i < n$  .

}{F7I-VI-052}\underline{Proof}:  Let  $0 \to J \to B' \to B \to 0$  be exact.

\SGAvar{(i)  If  $\mathrm{Tor}^{R}_{i}(B,\ S) = 0$  , then

}{(i)  If  $\mathrm{Tor}^{R}_{1}(B,\ S) = 0$  , then

}{F7I-VI-053}\begin{equation*}
0 \to J\underset{R}{\otimes}S \to B'\underset{R}{\otimes}S \to B\underset{R}{\otimes}S \to 0
\end{equation*}
is exact so that

\begin{equation*}
0 \to J\underset{R}{\otimes}S \to (B'\underset{R}{\otimes}S) \underset{(B\underset{R}{\otimes}S)}{\times} (C\underset{R}{\otimes}S) \to C\underset{R}{\otimes}S \to 0
\end{equation*}
is exact. We thus obtain an exact commutative diagram

\begin{tikzcd}
 \& J \underset{R}{\otimes} S \arrow[r] \arrow[d, equal] \& (B'\underset{B}{\times}C)\underset{R}{\otimes}S \arrow[r] \arrow[d] \& C \underset{R}{\otimes} S \arrow[r] \arrow[d, equal] \& 0 \\
0 \arrow[r] \& J \underset{R}{\otimes} S \arrow[r] \& (B'\underset{R}{\otimes}S) \underset{(B\underset{R}{\otimes}S)}{\times} (C \underset{R}{\otimes}S) \arrow[r] \& C \underset{R}{\otimes} S \arrow[r] \& 0
\end{tikzcd}















\SGAvar{\SGApage{I}{77}
and therefore $(B'\otimes_B C)\otimes_R S \to (B'\otimes_R S)\ \times_{(B\otimes_R S)}\ (C \otimes_R S)$  is an isomorphism.

}{\SGApage{I}{77}
and therefore $(B'\times_B C)\otimes_R S \to (B'\otimes_R S)\ \times_{(B\otimes_R S)}\ (C \otimes_R S)$  is an isomorphism.

}{F7I-VI-054}(ii)  Since  $B'$  is R-flat and  $\mathrm{Tor}^R_i(B,\ S) = 0$  for all  $0 < i \leq n$   by hypothesis, we have  $\mathrm{Tor}^R_i(J,\ S) = 0$  for all  $0 < i < n$  . Then the exact sequence  $0 \to J \to B'\times_B C \to C \to 0$  together with the hypothesis that  $\mathrm{Tor}^R_i(C,\ S) = 0$  for all  $0 < i < n$  entails that  $\mathrm{Tor}^R_i(B'\times_B C,\ S) = 0$  for all  $0 < i < n$  .


\underline{Corollary} 3.6.

\SGAvar{(a)  \underline{Let}  $A, S$  \underline{be} R-\underline{algebras. If}  $\mathrm{Tor}^R_1(A,\ S) = 0$  , \underline{then the canonical map}  $\check{H}^{\cdot}(A\otimes_R S|S,\ G) \to \check{H}^{\cdot}(A|R,\ SG)$  \underline{is an isomorphism for any}  $G \in \mathrm{ob}(\check{\underline{C}}_{A\otimes_R S|S})$  \underline{where the functor}  $S : \underline{C}_{A|R} \to \underline{C}_{A\otimes_R S|S}$  \underline{is given by}  $X \longmapsto X\otimes_R S$  . \underline{In particular if}  $S$  \underline{is} R-\underline{flat, then the canonical map}  $\check{S}\check{H}^{\cdot}(G) \to \check{H}(\check{S}G)$  \underline{is an isomorphism for every}  $G \in \mathrm{ob}(\check{\underline{C}}_{A\otimes_R S|S})$  .

}{(a)  \underline{Let}  $A, S$  \underline{be} R-\underline{algebras. If}  $\mathrm{Tor}^R_1(A,\ S) = 0$  , \underline{then the canonical map}  $\check{H}^{\cdot}(A\otimes_R S|S,\ G) \to \check{H}^{\cdot}(A|R,\ \check{S}G)$  \underline{is an isomorphism for any}  $G \in \mathrm{ob}(\check{\underline{C}}_{A\otimes_R S|S})$  \underline{where the functor}  $S : \underline{C}_{A|R} \to \underline{C}_{A\otimes_R S|S}$  \underline{is given by}  $X \longmapsto X\otimes_R S$  . \underline{In particular if}  $S$  \underline{is} R-\underline{flat, then the canonical map}  $\check{S}\underline{\check{H}}^{\cdot}(G) \to \underline{\check{H}}^{\cdot}(\check{S}G)$  \underline{is an isomorphism for every}  $G \in \mathrm{ob}(\check{\underline{C}}_{A\otimes_R S|S})$  .

}{F7I-VI-055}(b)  \underline{Let}  $P \to B$  \underline{be a surjective} R-\underline{algebra map, where}  $P$  \underline{is a polynomial algebra over}  $R$  . \underline{If}  $\mathrm{Tor}^R_i(B,\ S) = 0$  \underline{for all}  $0 < i \leq n$  , \underline{then} $\mathrm{Tor}^R_i(P\times_B P\times_B \ldots \times_B P,\ S) = 0$  \underline{for all}  $0 < i < n$  .


\underline{Proof}:

(a)  Given  $B \in \mathrm{ob}(\underline{C}_{A|R})$  , choose  $\{P \to B\} \in \mathrm{Cov}(B)$  where  $P$  is a polynomial algebra over  $R$  . If  $\mathrm{Tor}^R_1(B,\ S) = 0$  , the natural map

\begin{equation*}
(\overbrace{P\times_B \ \ldots\ \times_B P}^{\,n\,})\otimes_R S \to (P\otimes_R S)\times_{(B\otimes_R S)} \ldots \times_{(B\otimes_R S)} (P\otimes_R S)
\end{equation*}
is an isomorphism for all  $n$  by 3.5(i), and hence  $\check{H}^{\cdot}(B|R,\ \check{S}G) = H^{\cdot}(\{P \to B\},\ \check{S}G) = H^{\cdot}(\{P\otimes_R S \to B\otimes_R S\},\ G) = \check{H}^{\cdot}(B\otimes_R S|S,\ G)$  . If  $S$  is R-flat, then  $\mathrm{Tor}^R_1(X,\ S) = 0$  for every  $X \in \mathrm{ob}(\underline{C}_{A|R})$  and hence $\check{H}^{\cdot}(X\otimes_R S|S,\ G) \to \check{H}^{\cdot}(X|R,\ \check{S}G)$  is an isomorphism for every  $X \in \mathrm{ob}(\underline{C}_{A|R})$  i.e. $\check{S}\underline{\check{H}}^{\cdot}(G) \to \underline{\check{H}}^{\cdot}(\check{S}G)$  is an isomorphism.

(b)  Set  $P^{(m)} = \overbrace{P\times_B \ \ldots\ \times_B P}^{\,m+1\,}$  . If  $m = 0$  there is nothing to prove since  $P$  is R-flat. Assume that  $\mathrm{Tor}^R_i(P^{(m-1)},\ S) = 0$  for all  $0 < i < n$  . Since  $P \to B$  is surjective and  $P^{(m)} = P\times_B P^{(m-1)}$  , it follows from 3.5(ii) that $\mathrm{Tor}^R_i(P^{(m)},\ S) = 0$  for all  $0 < i < n$  .


\SGApage{I}{78}


\underline{Proposition} 3.7.

(a)  \underline{Let}  $A, S$  \underline{be} R-\underline{algebras. If}  $\mathrm{Tor}^R_i(A,\ S) = 0$  \underline{for all}  $0 < i \leq n$  , \underline{then for any sheaf}  $F$  \underline{on}  $\underline{C}_{A\otimes_R S|S}$  , \underline{the canonical map}  $H^i(A\otimes_R S|S,\ F) \to H^i(A|R,\ \widetilde{S}F)$ \underline{is an isomorphism for all}  $i \leq n$  , \underline{where}  $S : \underline{C}_{A|R} \to \underline{C}_{A\otimes_R S|S}$  \underline{is given by} $X \longmapsto X\otimes_R S$  .

(b)  \underline{Let}  $P$  \underline{be a polynomial algebra over}  $R$  . \underline{Then for any additive sheaf}  $F$  \underline{on}  $\underline{C}_{P\otimes_R A|R}$  , \underline{the canonical map}  $H^n(P\otimes_R A|R,\ F) \to H^n(A|R,\ F)$  \underline{is an isomorphism for all}  $n > 0$  .

(c)  \underline{Let}  $S \to R \to A$  \underline{be ring homomorphisms. Then for any additive sheaf}  $F$  \underline{on}  $\underline{C}_{A|S}$  , \underline{we have an exact sequence}

{\small \begin{multline*}
0 \to H^0(A|R,\ F_R) \to H^0(A|S,\ F) \to H^0(R|S,\ F) \to H^1(A|R,\ F_R) \to H^1(A|S,\ F) \to H^1(R|S,\ F) \to \\ \ldots\ldots\ldots \to H^n(A|R,\ F_R) \to H^n(A|S,\ F) \to H^n(R|S,\ F) \to H^{n+1}(A|R,\ F_R) \to \ldots.
\end{multline*}}
\underline{where}  $F_R$  \underline{is a sheaf on}  $\underline{C}_{A|R}$ \underline{defined by}

\begin{equation*}
F_R(X) = \mathrm{Ker}(F(X) \to F(R))
\end{equation*}
\underline{for all}  $X \in \mathrm{ob}(\underline{C}_{A|R})$  \underline{via the forgetful functor}  $\underline{C}_{A|R} \to \underline{C}_{A|S}$  .


\underline{Proof}:

(a)  Consider the continuous functor

\SGAvar{\begin{equation*}
S : \underline{C}_{A|R} \to \underline{C}_{A\otimes_R S|S}
\end{equation*}
where  $S(X) = X\otimes_R S$  . Since the functor  $S$  sends a polynomial algebras over  $R$  to polynomial algebras over  $S$  ,  $\widetilde{S} : \widetilde{\underline{C}}_{A\otimes_R S|S} \to \widetilde{\underline{C}}_{A|R}$  is an exact functor by 3.3(c) and in turn we obtain the canonical map  $\check{S}\underline{H}^{\cdot}_{A\otimes_R S|S}(F) \to \underline{H}^{\cdot}_{A|R}(\widetilde{S}F)$  , and in turn we obtain the commutative diagram of spectral sequences:

}{\begin{equation*}
S : \underline{C}_{A|R} \to \underline{C}_{A\otimes_R S|S}
\end{equation*}
where  $S(X) = X\otimes_R S$  . Since the functor  $S$  sends polynomial algebras over  $R$  to polynomial algebras over  $S$  ,  $\widetilde{S} : \widetilde{\underline{C}}_{A\otimes_R S|S} \to \widetilde{\underline{C}}_{A|R}$  is an exact functor by 3.3(c) and in turn we obtain the canonical map  $\check{S}\underline{H}^{\cdot}_{A\otimes_R S|S}(F) \to \underline{H}^{\cdot}_{A|R}(\widetilde{S}F)$  , and in turn we obtain the commutative diagram of spectral sequences:

}{F7I-VI-056}\begin{center}
\begin{tikzcd}[column sep=huge, row sep=large]
E^{p,q}_S(S(X)) = \check{H}^p(X\otimes_R S|S,\ \underline{H}^q_{A\otimes_R S|S}(F)) \arrow[r, Rightarrow] \arrow[d] \& H^{p+q}(X\otimes_R S|S,\ F) \arrow[dd] \\
\check{H}^p(X|R,\ \check{S}\underline{H}^q_{A\otimes_R S|S}(F)) \arrow[d] \& \\
E^{p,q}_R(X) =\ \ \check{H}^p\ (X|R,\ \underline{H}^q_{A|R}(\widetilde{S}F)) \arrow[r, Rightarrow] \& H^{p+q}(X|R,\ \widetilde{S}F)
\end{tikzcd}
\end{center}




















\SGApage{I}{79}

\SGAvar{We note that  $E^{p,q}_{S}(S(X)) \to \check{H}^{p}(X|R,\ \check{S}\underline{H}^{q}_{A\underset{R}{\otimes}S|S}(F))$  is an isomorphism whenever $\operatorname{Tor}^{R}_{1}(X,\ S) = 0$  by 3.6(a). Assume now that $[\check{S}\underline{H}^{q}_{A\underset{R}{\otimes}S|S}(F)](X) \to [\underline{H}^{q}_{A|R}(\widetilde{S}F)](X)$  is an isomorphism for all  $q < n$  and for all $X \in \mathrm{ob}(\underline{C}_{A|R})$  such that  $\operatorname{Tor}^{R}_{i}(X,\ S) = 0$  for all  $0\ \ \ i\ \ \ n$ . It follows immediately from 3.6(b) that  $E^{p,q}_{S}(S(X)) \to E^{p,q}_{R}(X)$  is an isomorphism for all  $q < n$ and for all  $p$ , provided  $\operatorname{Tor}^{R}_{i}(X,\ S) = 0$  for all  $0 < i \leq n$ . Since $E^{0,n}_{S}(S(X)) \to E^{0,n}_{R}(X)$  is an isomorphism (since both sides are zero), we get that $H^{n}(X\underset{R}{\phantom{\otimes}}S|S,\ F) \to H^{n}(X|R,\ \widetilde{S}F)$  is an isomorphism whenever  $\operatorname{Tor}^{R}_{i}(X,\ S) = 0$  for all $0 < i \leq n$  i.e.,  $[\check{S}\underline{H}^{i}_{A\underset{R}{\otimes}S|S}(F)](X) \to [\underline{H}^{i}_{A|R}(\widetilde{S}F)](X)$  is an isomorphism for all $i \leq n$  whenever  $\operatorname{Tor}^{R}_{i}(X,\ S) = 0$  for all  $0 < i \leq n$ .

(b)  Let  $P$  be a polynomial algebra over  $R$  and  let $f : \underline{C}_{A|R} \to \underline{C}_{A\underset{R}{\otimes}P|R}$  be the continuous functor given by  $f(X) = X\underset{R}{\otimes}P$ . Let  $P' \to X$ be a surjective arrow in  $\underline{C}_{A|R}$  where  $P'$  is a polynomial algebra over  $R$ . Then $P'\underset{R}{\otimes}P \to X\underset{R}{\otimes}P$  is a surjective arrow in  $\underline{C}_{A\underset{R}{\otimes}P|R}$ , and  $P'\underset{R}{\otimes}P$  is again a polynomial algebra over  $R$ . Since  $P$  is flat over  $R$ , it follows from 3.5(i) that $\check{H}^{\cdot}(X|R,\ \check{f}G) = H^{\cdot}(\{P' \to X\},\ \check{f}G) = H\ (\{P'\underset{R}{\otimes}P \to X\underset{R}{\otimes}P\},\ G) = \check{H}^{\cdot}(X\underset{R}{\otimes}P|R,\ G)$ for every  $G \in \mathrm{ob}(\underline{\check{C}}_{A\underset{R}{\otimes}P|R})$ . It follows that  $\widetilde{f} : \underline{\widetilde{C}}_{A\underset{R}{\otimes}P|R} \to \underline{C}_{A|R}$  sends acyclic ones to acyclic ones. On the other hand,  $f$  sends polynomial algebras over  $R$  to polynomial algebras over  $R$ , and hence  $\widetilde{f} : \underline{\widetilde{C}}_{A\underset{R}{\otimes}P|R} \to\!\underline{\widetilde{C}}_{A|R}$  is an exact functor by 3.2(c). It follows that  $H^{\cdot}(A\underset{R}{\otimes}P|R,\ F) \to H^{\cdot}(A|R,\ \widetilde{f}F)$  is an isomorphism for every $F \in \mathrm{ob}(\underline{\widetilde{C}}_{A\underset{R}{\otimes}P|R})$ . Now for each  $X \in \mathrm{ob}(\underline{C}_{A|R})$ , we have canonical maps

}{We note that  $E^{p,q}_{S}(S(X)) \to \check{H}^{p}(X|R,\ \check{S}\underline{H}^{q}_{A\underset{R}{\otimes}S|S}(F))$  is an isomorphism whenever $\operatorname{Tor}^{R}_{1}(X,\ S) = 0$  by 3.6(a). Assume now that $[\check{S}\underline{H}^{q}_{A\underset{R}{\otimes}S|S}(F)](X) \to [\underline{H}^{q}_{A|R}(\widetilde{S}F)](X)$  is an isomorphism for all  $q < n$  and for all $X \in \mathrm{ob}(\underline{C}_{A|R})$  such that  $\operatorname{Tor}^{R}_{i}(X,\ S) = 0$  for all  $0 < i < n$ . It follows immediately from 3.6(b) that  $E^{p,q}_{S}(S(X)) \to E^{p,q}_{R}(X)$  is an isomorphism for all  $q < n$ and for all  $p$ , provided  $\operatorname{Tor}^{R}_{i}(X,\ S) = 0$  for all  $0 < i \leq n$ . Since $E^{0,n}_{S}(S(X)) \to E^{0,n}_{R}(X)$  is an isomorphism (since both sides are zero), we get that $H^{n}(X\underset{R}{\otimes}S|S,\ F) \to H^{n}(X|R,\ \widetilde{S}F)$  is an isomorphism whenever  $\operatorname{Tor}^{R}_{i}(X,\ S) = 0$  for all $0 < i \leq n$  i.e.,  $[\check{S}\underline{H}^{i}_{A\underset{R}{\otimes}S|S}(F)](X) \to [\underline{H}^{i}_{A|R}(\widetilde{S}F)](X)$  is an isomorphism for all $i \leq n$  whenever  $\operatorname{Tor}^{R}_{i}(X,\ S) = 0$  for all  $0 < i \leq n$ .

(b)  Let  $P$  be a polynomial algebra over  $R$  and  let $f : \underline{C}_{A|R} \to \underline{C}_{A\underset{R}{\otimes}P|R}$  be the continuous functor given by  $f(X) = X\underset{R}{\otimes}P$ . Let  $P' \to X$ be a surjective arrow in  $\underline{C}_{A|R}$  where  $P'$  is a polynomial algebra over  $R$ . Then $P'\underset{R}{\otimes}P \to X\underset{R}{\otimes}P$  is a surjective arrow in  $\underline{C}_{A\underset{R}{\otimes}P|R}$ , and  $P'\underset{R}{\otimes}P$  is again a polynomial algebra over  $R$ . Since  $P$  is flat over  $R$ , it follows from 3.5(i) that $\check{H}^{\cdot}(X|R,\ \check{f}G) = H^{\cdot}(\{P' \to X\},\ \check{f}G) = H^{\cdot}\ (\{P'\underset{R}{\otimes}P \to X\underset{R}{\otimes}P\},\ G) = \check{H}^{\cdot}(X\underset{R}{\otimes}P|R,\ G)$ for every  $G \in \mathrm{ob}(\underline{\check{C}}_{A\underset{R}{\otimes}P|R})$ . It follows that  $\widetilde{f} : \underline{\widetilde{C}}_{A\underset{R}{\otimes}P|R} \to \underline{\widetilde{C}}_{A|R}$  sends acyclic ones to acyclic ones. On the other hand,  $f$  sends polynomial algebras over  $R$  to polynomial algebras over  $R$ , and hence  $\widetilde{f} : \underline{\widetilde{C}}_{A\underset{R}{\otimes}P|R} \to\!\underline{\widetilde{C}}_{A|R}$  is an exact functor by 3.3(c). It follows that  $H^{\cdot}(A\underset{R}{\otimes}P|R,\ F) \to H^{\cdot}(A|R,\ \widetilde{f}F)$  is an isomorphism for every $F \in \mathrm{ob}(\underline{\widetilde{C}}_{A\underset{R}{\otimes}P|R})$ . Now for each  $X \in \mathrm{ob}(\underline{C}_{A|R})$ , we have canonical maps

}{F7I-VI-057}\begin{equation*}
(\widetilde{f}F)\ (X) = F(X\underset{R}{\otimes}P) \to F(X)\oplus F(P)\ \ ,
\end{equation*}
and hence a canonical map  $\widetilde{f}F \to F\oplus F(P)$  where  $F(P)$  denotes a constant sheaf (and  $F$  on the right hand side denotes, strictly speaking, the restriction of  $F$ on  $\underline{C}_{A|R}$ ) . If  $F$  is an additive sheaf, then  $(\widetilde{f}F)\ (X) = F(X\underset{R}{\otimes}P) \overset{\sim}{\to} F(X)\oplus F(P)$ is an isomorphism for every  $X \in \mathrm{ob}(\underline{C}_{A|R})$  and hence  $\widetilde{f}F \overset{\sim}{\to} F\oplus F(P)$  is an isomorphism. Therefore 


\SGApage{I}{80}

\begin{equation*}
 H^{n}(A|R,\ \widetilde{f}F) \overset{\sim}{\to} H^{n}(A|R,\ F\oplus F(P)) = H^{n}(A|R,\ F)
\end{equation*}
for all  $n > 0$  since  $F(P)$ , being a constant sheaf, is flask. It follows that $H^{n}(A\underset{R}{\otimes}P|R,\ F) \to H^{n}(A|R,\ F)$  is an isomorphism for all  $n > 0$ , provided  $F$  is an additive sheaf on  $\underline{C}_{A\underset{R}{\otimes}P|R}$ .

(c)  (communicated from Rinehart) Consider the forgetful functor $\beta : \underline{C}_{A|R} \to \underline{C}_{A|S}$ . We have a spectral sequence

\SGAvar{\begin{equation*}
E^{p,q} = H^{p}(A|R,\ \beta^{\#}\underline{H}^{q}_{A|S}(F))\ \Longrightarrow\ H^{p+q}(A|S,\ F)
\end{equation*}
where  $\beta^{\#} = G_{A|R}\circ\check{\beta}$ ,  $G_{A|R} : \underline{\check{C}}_{A|R} \to \underline{\widetilde{C}}_{A|R}$  is "associated-sheaf"  functor. Now let  $F$  be an additive sheaf. We claim that  $\beta^{\#}\underline{H}^{q}_{A|S}(F)$  is a constant sheaf for all  $q > 0$ . Indeed, if  $P$  is a polynomial algebra over  $R$ , then  $P = R\underset{S}{\otimes}P'$ for a polynomial algebra  $P'$  over  $S$  and hence

}{\begin{equation*}
E^{p,q} = H^{p}(A|R,\ \beta^{\#}\underline{H}^{q}_{A|S}(F))\ \Longrightarrow\ H^{p+q}(A|S,\ F)
\end{equation*}
where  $\beta^{\#} = G_{A|R}\circ\check{\beta}$ ,  $G_{A|R} : \underline{\check{C}}_{A|R} \to \underline{\widetilde{C}}_{A|R}$  is the "associated-sheaf"  functor. Now let  $F$  be an additive sheaf. We claim that  $\beta^{\#}\underline{H}^{q}_{A|S}(F)$  is a constant sheaf for all  $q > 0$ . Indeed, if  $P$  is a polynomial algebra over  $R$ , then  $P = R\underset{S}{\otimes}P'$ for a polynomial algebra  $P'$  over  $S$  and hence

}{F7I-VI-058}\begin{equation*}
[\beta^{\#}\underline{H}^{q}_{A|S}(F)](P) = H^{q}(P|S,\ F) = H^{q}(R\underset{S}{\otimes}P'|S,\ F) \to H^{q}(R|S,\ F)
\end{equation*}
is an isomorphism for all  $q > 0$  by 3.7(b), and surjective for  $q = 0$ . It follows from 3.3(c) that the canonical map  $\beta^{\#}\underline{H}^{q}_{A|S}(F) \to H^{q}(R|S,\ F) =$ constant sheaf is an isomorphism for all  $q > 0$ , and is surjective for  $q = 0$ . It follows that $E^{p,q} = H^{p}(A|R,\ \beta^{\#}\underline{H}^{q}_{A|S}(F)) = 0$  for all  $p,\ q > 0$ , and hence the above spectral sequence yields the exact sequence

\begin{equation*}
(*)\quad 0 \to E^{1,0} \to H^{1} \to E^{0,1}\!\to\ E^{2,0} \to H^{2} \to E^{0,2} \to E^{3,0}\!\to\ H^{3} \to\ ....
\end{equation*}
where we note that  $E^{0,q} = [\beta^{\#}\underline{H}^{q}_{A|S}(F)](A) \overset{\sim}{\to} H^{q}(R|S,\ F)$  for all  $q > 0$  and $E^{p,0} = H^{p}(A|R,\ \widetilde{\beta}F)$ . On the other hand, the exact sequence  $0 \to F_{R} \to \widetilde{\beta}F \to F(R) \to 0$ gives us the exact sequence

{\footnotesize (**)\quad\begin{tikzcd}[column sep=tiny, row sep=large]
0 \arrow[r] \& H^{0}(A|R,F_{R}) \arrow[r] \& H^{0}(A|R,\widetilde{\beta}F) \arrow[r] \arrow[d, equal] \& H^{0}(A|R,F(R)) \arrow[r] \arrow[d, equal] \& H^{1}(A|R,F_{R}) \arrow[r] \& H^{1}(A|R,\widetilde{\beta}F) \arrow[r] \& 0 \\
 \& \& H^{0}(A|S,F) \& H^{0}(R|S,F) \& \& \&
\end{tikzcd}}










\SGAvar{The exact sequences (*) and (**) combined gives the desired exact sequence.

}{The exact sequences (*) and (**) combined give the desired exact sequence.

}{F7I-VI-059}\SGApage{I}{81}


\underline{Corollary} 3.8.

(a)  \underline{If}  $A, B$  \underline{are} $R$-\underline{algebras such that}  $\mathrm{Tor}^{R}_{i}(A, B) = 0$  \underline{for all} $0 < i \leq n$  , \underline{then for any additive sheaf}  $F \in \mathrm{ob}(\widetilde{\underline{C}}_{A\underset{R}{\otimes}B|R})$  , \underline{the canonical map}

\begin{equation*}
H^{i}(A\underset{R}{\otimes}B|R,\ F) \rightarrow H^{i}(A|R,\ F) \oplus H^{i}(B|R,\ F)
\end{equation*}
\underline{is an isomorphism for all}  $i \leq n$  .

(b)  \underline{Let}  $A$  \underline{be an} $R$-\underline{algebra such that}  $\mathrm{Tor}^{R}_{i}(A, A) = 0$  \underline{for all}  $i > 0$  . \underline{If}  $A\underset{R}{\otimes}A \rightarrow A$  \underline{is an isomorphism, then for any additive sheaf}  $F \in \mathrm{ob}(\widetilde{\underline{C}}_{A|R})$  , \underline{we have}  $H^{i}(A|R,\ F) = 0$  \underline{for all}  $i$  .


\underline{Proof:}

\SGAvar{(a)  Let  $F \in \mathrm{ob}(\widetilde{\underline{C}}_{A\underset{R}{\otimes}B|R})$  be additive. By 3.7(c), the ring homomorphisms $R \rightarrow A \rightarrow A\underset{R}{\otimes}B$  gives us the exact sequence

\begin{equation}\tag{*}
... \rightarrow H^{i}(A\underset{R}{\otimes}B\ A,\ F_{A}) \rightarrow H^{i}(A\underset{R}{\otimes}B\ R,\ F) \rightarrow H^{i}(A|R,\ F) \rightarrow ....
\end{equation}
consider now the commutative diagram

}{(a)  Let  $F \in \mathrm{ob}(\widetilde{\underline{C}}_{A\underset{R}{\otimes}B|R})$  be additive. By 3.7(c), the ring homomorphisms $R \rightarrow A \rightarrow A\underset{R}{\otimes}B$  give us the exact sequence

\begin{equation}\tag{*}
... \rightarrow H^{i}(A\underset{R}{\otimes}B|A,\ F_{A}) \rightarrow H^{i}(A\underset{R}{\otimes}B|R,\ F) \rightarrow H^{i}(A|R,\ F) \rightarrow ....
\end{equation}
consider now the commutative diagram

}{F7I-VI-060}\begin{center}
\begin{tikzcd}[column sep=huge, row sep=huge]
H^{i}(A\underset{R}{\otimes}B|A,\ F_{A}) \arrow[r] \arrow[d] \& H^{i}(A\underset{R}{\otimes}B|R,\ F) \arrow[d] \\
H^{i}(B|R,\ \widetilde{A}F_{A}) \arrow[r, "\approx"] \& H^{i}(B|R,\ F)
\end{tikzcd}
\end{center}









where the functor  $A : \underline{C}_{B|R} \rightarrow \underline{C}_{A\underset{R}{\otimes}B|A}$  is given by  $X \longmapsto A\underset{R}{\otimes}X$  , and the map $H^{i}(B|R,\ \widetilde{A}F_{A}) \rightarrow H^{i}(B|R,\ F)$  is induced from the map  $\widetilde{A}F_{A} \rightarrow F$  , which is an isomorphism since  $F$  is additive.  Note that the commutativity of the above diagram, because of its functorial nature, is reduced to the case  $i = 0$  .

If  $\mathrm{Tor}^{R}_{i}(A, B) = 0$  for all  $0 < i \leq n$  , then  $H^{i}(A\underset{R}{\otimes}B|A,\ F_{A}) \rightarrow H^{i}(B|R,\ \widetilde{A}F_{A})$  is an isomorphism for all  $i \leq n$  by 3.7(a), and hence the composite map

\begin{equation*}
H^{i}(A\underset{R}{\otimes}B|A,\ F_{A}) \rightarrow H^{i}(A\underset{R}{\otimes}B|R,\ F) \rightarrow H^{i}(B|R,\ F)
\end{equation*}
and (because of the symmetry) 


\SGApage{I}{82}

\begin{equation*}
 H^{i}(A\underset{R}{\otimes}B|B,\ F_{B}) \rightarrow H^{i}(A\underset{R}{\otimes}B|R,\ F) \rightarrow H^{i}(A|R,\ F)
\end{equation*}
are isomorphisms for all  $i \leq n$  . This implies that the exact sequence (*) breaks up to short split-exact sequences

\begin{equation*}
0 \rightarrow H^{i}(A\underset{R}{\otimes}B|A,\ F_{A}) \rightarrow H^{i}(A\underset{R}{\otimes}B|R,\ F) \rightarrow H^{i}(A|R,\ F) \rightarrow 0
\end{equation*}
for all  $i \leq n$  , and that

\begin{equation*}
H^{i}(A\underset{R}{\otimes}B|R,\ F) \rightarrow H^{i}(A|R,\ F) \oplus H^{i}(B|R,\ F)
\end{equation*}
is an isomorphism for all  $i \leq n$  .

(b)  Since  $\mathrm{Tor}^{R}_{i}(A,A) = 0$  for all  $i > 0$  , it follows from 3.8(a) that

\begin{equation*}
H^{i}(A\underset{R}{\otimes}A|R,\ F) \rightarrow H^{i}(A|R,\ F) \oplus H^{i}(A|R,\ F)
\end{equation*}
is an isomorphism for all  $i$  . On the other hand,  $A\underset{R}{\otimes}A \rightarrow A$  is an isomorphism by hypothesis and hence  $H^{i}(A|R,\ F) \rightarrow H^{i}(A\underset{R}{\otimes}A|R,\ F)$  is an isomorphism for all  $i$  . This means that the diagonal map  $H^{i}(A|R,\ F) \rightarrow H^{i}(A|R,\ F) \oplus H^{i}(A|R,\ F)$  is an isomorphism for all  $i$  i.e.,  $H^{i}(A|R,\ F) = 0$  for all  $i$  .

\underline{Corollary} 3.9:  Let  $S$  be a multiplicative subset of an $R$-algebra  $A$  . If  $F \in \mathrm{ob}\!\left(\widetilde{\underline{C}}_{S^{-1}A|R}\right)$  is additive such that  $F_{A}$  is also additive, then the canonical map

\begin{equation*}
H^{i}(S^{-1}A|R,\ F) \rightarrow H^{i}(A|R,\ F)
\end{equation*}
is an isomorphism for all  $i$  , where  $F_{A} \in \mathrm{ob}\!\left(\widetilde{\underline{C}}_{S^{-1}A|A}\right)$  is defined by $F_{A}(X) = \mathrm{Ker}(F(X) \rightarrow F(A))$  for all  $X \in \mathrm{ob}\!\left(\underline{C}_{S^{-1}A|A}\right)$  .

\SGAvar{\underline{Proof}:  Since  $F$  is additive, the ring homomorphisms  $R \rightarrow A \rightarrow S^{-1}A$  gives us the exact sequence

}{\underline{Proof}:  Since  $F$  is additive, the ring homomorphisms  $R \rightarrow A \rightarrow S^{-1}A$  give us the exact sequence

}{F7I-VI-061}\begin{equation*}
.... \rightarrow H^{n}(S^{-1}A|A,\ F_{A}) \rightarrow H^{n}(S^{-1}A|R,\ F) \rightarrow H^{n}(A|R,\ F) \rightarrow ....
\end{equation*}
On the other hand, since  $S^{-1}A\underset{A}{\otimes}S^{-1}A \rightarrow S^{-1}A$  is an isomorphism and  $F_{A}$  is additive by hypothesis, it follows from 3.8(b) that  $H^{n}(S^{-1}A|A,\ F_{A}) = 0$  for all  $n$  , and hence  $H^{n}(S^{-1}A|R,\ F) \rightarrow H^{n}(A|R,\ F)$  is an isomorphism for all  $n$  .


\SGApage{I}{83}

Now let  $A$  be an $R$-algebra. An $A$-module  $M$  defines an additive sheaf $\mathrm{Der}_R(-,\ M) \in \mathrm{ob}(\widetilde{\underline{C}}_{A|R})$  by  $X \to \mathrm{Der}_R(X,\ M)$  . We now make the following definition:


\underline{Definition} 3.10:  \underline{Let}  $A$  \underline{be an} $R$-\underline{algebra. For each} $A$-\underline{module}  $M$  \underline{we set}

\begin{equation*}
D^i(A|R,\ M) = H^i(A|R,\ \mathrm{Der}_R(-,\ M)) \ \ .
\end{equation*}
We want to make an explicit computation of low-dimensional ones i.e., $D^i(A|R,\ M)$  for  $i = 1,\ 2$  : Choose a surjective arrow  $P \to A$  in  $\underline{C}_{A|R}$  where  $P$  is a polynomial algebra over  $R$  , and let  $0 \to I \to P \to A \to 0$  be exact. Then


\underline{Lemma} 3.10:

\begin{equation*}
\begin{aligned}
D^1(A|R,\ M) &= \mathrm{Coker}(\mathrm{Der}_R(P,\ M) \to \mathrm{Hom}_P(I,\ M)) \\
D^2(A|R,\ M) &= H^1(A|R,\ \underline{H}^1_{A|R}(\mathrm{Der}_R(-,\ M)))
\end{aligned}
\end{equation*}
\underline{Proof}:  Consider the simplicial algebra

\begin{multline}\tag{*}
........\ \ P_A^{(2)} \ \mathrel{\substack{\to\\ \to\\ \to}} \ P_A^{(1)} \ \mathrel{\substack{\to\\ \to}} \ P_A^{(0)}
\end{multline}
where  $P_A^{(n)} = \overbrace{P\times_A \ldots \times_A P}^{\,n+1}$  . For each  $n$  we have the exact sequence

\begin{equation*}
0 \to I^{(n)} \to P_A^{(n)} \to A \to 0
\end{equation*}
where  $I^{(n)} = \overbrace{I\times I\times \ldots \times I}^{\,n+1}$  , and it is easy to see that

(a):

\begin{equation*}
0 \to \mathrm{Hom}_A(\overline{I}^{(n)}/\Delta(\overline{I}),\ M) \to \mathrm{Der}_R(P_A^{(n)},\ M) \xrightarrow{\ \mathrm{Der}(\Delta,\ M)\ } \mathrm{Der}_R(P,\ M) \to 0
\end{equation*}
is exact for all  $n$  , where  $\overline{I} = I/I^2$  and  $\Delta$  stands for the diagonal map. We also note the exact sequence

(b):

\SGAvar{\begin{equation*}
0 \to \mathrm{Hom}_A(\overline{I}^{(n)}/\Delta(\overline{I}),\ M) \to \mathrm{Hom}_A(\overline{I}^{(n)},\ M) \xrightarrow{\ \mathrm{Hom}(\Delta,\ M)\ } \mathrm{Hom}_A(I,\ M) \to 0
\end{equation*}
The exact sequences (a) and (b) yield the respective exact sequences of simplicial $A$-modules. Since  $\{I \to 0\}$ , $\{I \to I\}$  admits sections, we have

}{\begin{equation*}
0 \to \mathrm{Hom}_A(\overline{I}^{(n)}/\Delta(\overline{I}),\ M) \to \mathrm{Hom}_A(\overline{I}^{(n)},\ M) \xrightarrow{\ \mathrm{Hom}(\Delta,\ M)\ } \mathrm{Hom}_A(\overline{I},\ M) \to 0
\end{equation*}
The exact sequences (a) and (b) yield the respective exact sequences of simplicial $A$-modules. Since  $\{I \to 0\}$ , $\{I \to I\}$  admit sections, we have

}{F7I-VI-062}\begin{equation*}
H^i(\{I \to 0\},\ \mathrm{Hom}_A(-,\ M)) = H^i(\{I \to I\},\ \mathrm{Hom}_A(-,\ M)) = 0 
\end{equation*}

\SGApage{I}{84}

 for all  $i > 0$  , and it follows readily that

\begin{equation*}
\left\{\begin{array}{l}
0 \to \mathrm{Der}_R(A,\ M) \to \mathrm{Der}_R(P,\ M) \to \mathrm{Hom}_A(\overline{I},\ M) \to H^1(\{P \to A\},\ \mathrm{Der}_R(-,\ M)) \to 0 \\[4pt]
\hspace{9cm}\text{is exact} \\[4pt]
H^i(\{P \to A\},\ \mathrm{Der}_R(-,\ M)) = 0 \ \ \text{for all} \ \ i > 1
\end{array}\right.
\end{equation*}
It follows from 3.1(b) that

\begin{equation*}
\begin{aligned}
D^1(A|R,\ M) &= \check{H}^1(A|R,\ \mathrm{Der}_R(-,\ M)) = H^1(\{P \to A\},\ \mathrm{Der}_R(-,\ M)) \\
&= \mathrm{Coker}(\mathrm{Der}_R(P,\ M) \to \mathrm{Hom}_A(\overline{I},\ M)) \\
&= \mathrm{Coker}(\mathrm{Der}_R(P,\ M) \to \mathrm{Hom}_P(I,\ M)) \ \ , \\
D^2(A|R,\ M) &= H^1(A|R,\ \underline{H}^1_{A|R}(\mathrm{Der}_R(-,\ M))) \ \ .
\end{aligned}
\end{equation*}
Now choose a free $P$-module  $F$  and a $P$-linear map  $F \xrightarrow{\ \gamma\ } P$  such that  $\gamma(F) = I$  . Associated to $P$-linear map  $\gamma : F \to P$  , there is a Koszul complex, which is simply denoted, by abuse of notation, by  $K_I$  .


\SGAvar{\underline{Lemma} 3.11:   $D^2(A|R,\ M) \simeq \mathrm{Coker}(H^1(K_I,\ M) \to \mathrm{Hom}_A(H_1(K^I),\ M))$ .

}{\underline{Lemma} 3.11:   $D^2(A|R,\ M) \simeq \mathrm{Coker}(H^1(K_I,\ M) \to \mathrm{Hom}_A(H_1(K_I),\ M))$ .

}{F7I-VI-063}\underline{Proof}:  We have

\begin{equation*}
\begin{aligned}
D^2(A|R,\ M) &= H^1(A|R,\ H^1_{A|R}(\mathrm{Der}_R(-,\ M))) \\
&= \mathrm{Ker}(H^1(P_A^{(1)}|R,\ M) \ \mathrel{\substack{\to\\ \to\\ \to}} \ H^1(P_A^{(2)}|R,\ M))
\end{aligned}
\end{equation*}
since  $[\underline{H}^1_{A|R}(\mathrm{Der}_R(-,\ M))](P_A^0) = [\underline{H}^1_{A|R}(\mathrm{Der}_R(-,\ M))](P) = 0$  .

Let  $0 \to K \to F \xrightarrow{\ \gamma\ } I \to 0$  be exact. If we set  $P[F]$  to be the symmetric algebra over  $P$  generated by  $F$  , then  $P[F]$  and  $P[F]\otimes P[F]$  are polynomial algebras over  $R$  and we obtain the surjective arrows

\begin{equation*}
\begin{aligned}
P[F] \to P_A^{(1)} &= P\times_A P &&\text{given by} && X \to (0,\ \gamma(X)) \\
P[F]\otimes_P P[F] \to P_A^{(2)} &= P\times_A P\times_A P && && X\otimes 1 \to (0,\ \gamma(X),\ 0),\ 1\otimes X \to (0,\ 0,\ \gamma(X))
\end{aligned}
\end{equation*}
where  $X \in F$  . If we set  $J = \mathrm{Ker}(P[F] \to P_A^{(1)})$ and
$J' = \mathrm{Ker}(P[F]\otimes_P P[F] \to P_A^{(2)})$  , then 


\SGApage{I}{85}

 
\begin{align*}
J &= (F) \cap (F_\gamma) = (K) + (F\cdot F_\gamma)\\
J' &= (F\otimes 1 + 1\otimes F) \cap (F_\gamma\otimes 1 + 1\otimes F) \cap (F\otimes 1 + 1\otimes F_\gamma)\\
&= (K\otimes 1 + 1\otimes K) + (F\cdot F_\gamma\otimes 1 + 1\otimes F\cdot F_\gamma) + (F\otimes F)
\end{align*}

where  $F_\gamma$ $=$ the ideal generated by  $\{X - \gamma(X)|X \in F\}$  . Now  $\pi_i : P_A^{(2)} \to P_A^{(1)}$  $i = 0, 1, 2$  can be lifted to  $\hat{\pi}_i : P[F]\underset{P}{\otimes}P[F] \to P[F]$  where, for each  $X \in F$  ,

$$\pi_0 = \begin{cases} X\otimes 1 \to \gamma(X) - X\\[3ex] 1\otimes X \to X\end{cases}\qquad\qquad \pi_1 = \begin{cases} X\otimes 1 \to 0\\[3ex] 1\otimes X \to X\end{cases}\qquad\qquad \pi_2 = \begin{cases} X\otimes 1 \to X\\[3ex] 1\otimes X \to 0\end{cases}\qquad .$$
A trivial computation shows that  $\pi_0 - \pi_1 + \pi_2$  sends  $K\otimes 1 + 1\otimes K + 1\otimes F\cdot F_\gamma$  to zero, whereas it sends  $(F\cdot F_\gamma\otimes 1 + F\otimes F)$  onto  $(F\cdot F_\gamma)$  . Now for any $D \in \mathrm{Der}_R(P[F]\underset{P}{\otimes}P[F], M)$  ,  $D(F\cdot F_\gamma\otimes 1 + F\otimes F) = 0$  since  $I\cdot M = 0$  , and consequently

\begin{align*}
D^2(A|R, M) &= \mathrm{Ker}(D^1(P_A^{(1)}|R, M) \mathrel{\substack{\to\\ \to\\ \to}} D^1(P_A^{(2)}|R, M))\\
&= \{[f] \in D^1(P_A^{(1)}|R, M)|f(F\cdot F_\gamma) = 0\}\\
&= \mathrm{Coker}(\mathrm{Der}_R(P[F], M) \to \mathrm{Hom}_{P[F]}((F) \cap (F_\gamma)/(F\cdot F_\gamma), M))\\
&= \mathrm{Coker}(\mathrm{Der}_P(P[F], M) \to \mathrm{Hom}_{P[F]}((F) \cap (F_\gamma)/(F)\cdot(F_\gamma), M))\quad .
\end{align*}

\SGAvar{Now
$$\begin{aligned}
0 \to (F) \to P[F] &\overset{\;\;P[0]\;\;}{\relbar\joinrel\relbar\joinrel\relbar} P \to 0\\
0 \to (F_\gamma) \to P[F] &\overset{\;\;P[\gamma]\;\;}{\relbar\joinrel\relbar\joinrel\relbar} P \to 0 \quad .
\end{aligned}$$
are exact and hence  $(F) \cap (F_\gamma)/(F)\cdot(F_\gamma) = \mathrm{Tor}_1^{P[F]}(P, P_\gamma) = H_1(K_I)$  where  $K_I$  is the Koszul complex associated  to the $P$-linear map  $F \xrightarrow{\gamma} I(\subset P)$  . We thus have

\begin{align*}
D^2(A|R, M) &= \mathrm{Coker}(\mathrm{Hom}_P(F, M) \to \mathrm{Hom}_P(H_1(K_I), M))\\
&= \mathrm{Coker}(H^1(K_I, M) \to \mathrm{Hom}_A(H^1(K_I), M))\quad .
\end{align*}

}{Now
$$\begin{aligned}
0 \to (F) \to P[F] &\xrightarrow{\;\;P[0]\;\;} P \to 0\\
0 \to (F_\gamma) \to P[F] &\xrightarrow{\;\;P[\gamma]\;\;} P \to 0 \quad .
\end{aligned}$$
are exact and hence  $(F) \cap (F_\gamma)/(F)\cdot(F_\gamma) = \mathrm{Tor}_1^{P[F]}(P, P_\gamma) = H_1(K_I)$  where  $K_I$  is the Koszul complex associated  to the $P$-linear map  $F \xrightarrow{\gamma} I(\subset P)$  . We thus have

\begin{align*}
D^2(A|R, M) &= \mathrm{Coker}(\mathrm{Hom}_P(F, M) \to \mathrm{Hom}_P(H_1(K_I), M))\\
&= \mathrm{Coker}(H^1(K_I, M) \to \mathrm{Hom}_A(H_1(K_I), M))\quad .
\end{align*}

}{F7I-VI-064}We can now state all the basic facts needed for our purpose, summarizing the various results we have already established.


\SGApage{I}{86}


\underline{Theorem} 3.12.

(1)  \underline{If}  $0 \to M' \to M \to M'' \to 0$  \underline{is an exact sequence of} $A$-\underline{modules, then}

\begin{gather*}
0 \to D^0(A|R, M') \to D^0(A|R, M) \to D^0(A|R, M'') \to D^1(A|R, M') \to\\
D^1(A|R, M)\to\;\; D^1(A|R, M'') \to\; ........ \;\to D^n(A|R, M') \to\\
D^n(A|R, M) \to D^n(A|R, M'') \to \;...
\end{gather*}

\underline{is exact}.

(2)  \underline{Let}  $S \to R \to A$  \underline{be ring homomorphisms. Then for any $A$-module}  $M$  \underline{we have the exact sequence}

\SGAvar{\begin{gather*}
0 \to D^0(A\;R, M) \to D^0(A\;S, M) \to D^0(R\;S, M) \to D^1(A\;R, M) \to\\
D^1(A\;S, M) \to D^1(R\;S, M) \to\; ........ \;\to D^n(A\;R, M) \to\\
D^n(A\;S, M) \to D^n(R\;S, M) \to \;....
\end{gather*}

}{\begin{gather*}
0 \to D^0(A|R, M) \to D^0(A|S, M) \to D^0(R|S, M) \to D^1(A|R, M) \to\\
D^1(A|S, M) \to D^1(R|S, M) \to\; ........ \;\to D^n(A|R, M) \to\\
D^n(A|S, M) \to D^n(R|S, M) \to \;....
\end{gather*}

}{F7I-VI-065}(3)  \underline{Let}  $A, S$  \underline{be $R$-algebras such that}  $\mathrm{Tor}_i^R(A, S) = 0$  \underline{for all}  $i > 0$  . \underline{Then for any}  $A\underset{R}{\otimes}S$-\underline{module}  $M$  , \underline{the canonical map}  $D^i(A\underset{R}{\otimes}S|S, M) \to D^i(A|R, M)$  \underline{is an} \underline{isomorphism for all}  $i$  . \underline{In particular, if}  $A$  \underline{is} $R$-\underline{flat we have}

$$D^i(A\otimes S|S, M) \xrightarrow{\sim} D^i(A|R, M)$$
\underline{for all}  $i$  \underline{and for all} $R$-\underline{algebras}  $S$  .

(4)  \underline{Let}  $S$  \underline{be a multiplicative subset of}  $A$  . \underline{Then for any}  $S^{-1}A$-\underline{module} $N$  , \underline{the canonical map}

$$D^i(S^{-1}A|R, N) \to D^i(A|R, N)$$
\underline{is an isomorphism for all}  $i$  . \underline{Consequently, for any} $A$-\underline{module}  $M$  \underline{we have the} \underline{canonical isomorphism}

$$D^i(S^{-1}A|R, S^{-1}M) \xrightarrow{\sim} S^{-1}D^i(A|R, M)$$
\underline{for all}  $i$  .

(5)  \underline{If}  $A, B$  \underline{are} $R$-\underline{algebras such that}  $\mathrm{Tor}_i^R(A, B) = 0$  \underline{for all} $0 < i \leq n$  , \underline{then for any}  $A\underset{R}{\otimes}B$-\underline{module}  $M$  , \underline{the canonical map}

$$D^i(A\underset{R}{\otimes}B|R, M) \to D^i(A|R, M) \oplus D^i(B|R, M)$$
\underline{is an isomorphism for all}  $i \leq n$  . \underline{In particular, if}  $\mathrm{Tor}_i^R(A, A) = 0$  \underline{for all}
\SGApage{I}{87}
$i > 0$  \underline{and}

\begin{equation*}
\Delta \ : \ \mathrm{Spec}(A) \ \to \ \mathrm{Spec}(A) \underset{\mathrm{Spec}(R)}{\times} \mathrm{Spec}(A)
\end{equation*}
\underline{is an open immersion} (\underline{for instance}  $A/R$  \underline{etale}), \underline{then}  $D^i(A|R, M) = 0$  \underline{for all} $i > 0$  \underline{and for all} $A$-\underline{modules}  $M$ .

(6) \underline{Let}  $A$  \underline{be an} $R$-\underline{algebra of essentially finite type over}  $R$ , \underline{where} $R$  \underline{is noetherian, and let}  $E$  \underline{be a flat} $A$-\underline{module. Then we have a canonical isomorphism}

\begin{equation*}
E \underset{A}{\otimes} D^i(A|R, \ M) \ \overset{\sim}{\longrightarrow} \ D^i(A|R, \ E \underset{A}{\otimes} M)
\end{equation*}
\underline{for all}  $i \geq 0$  \underline{and for all} $A$-\underline{modules}  $M$ .

\SGAvar{(7) \underline{If}  $R$  \underline{is noetherian and}  $A$  \underline{is an} $R$-\underline{algebra of essentially finite} \underline{type, then}  $D^i(A|R, M)$  \underline{is} $A$-\underline{module of finite type for all}  $i$ , \underline{provided}  $M$  \underline{is} \underline{an} $A$-\underline{module of finite type}.

}{(7) \underline{If}  $R$  \underline{is noetherian and}  $A$  \underline{is an} $R$-\underline{algebra of essentially finite} \underline{type, then}  $D^i(A|R, M)$  \underline{is an} $A$-\underline{module of finite type for all}  $i$ , \underline{provided}  $M$  \underline{is} \underline{an} $A$-\underline{module of finite type}.

}{F7I-VI-066}(8) \underline{Let}  $P$  \underline{be a polynomial algebra over}  $R$  \underline{and let}  $0 \to I \to P \to A \to 0$ \underline{be exact. Then}  $D^0(A|R, M) = \mathrm{Der}_R(A, M)$  ,

\SGAvar{\begin{equation*}
\begin{aligned}
D^1(A|R, \ M) \ &= \ \mathrm{Coker}(\mathrm{Der}_R(P, \ M) \ \to \ \mathrm{Hom}_A(I/I^2, \ M)\\
&= \ \underline{\text{the set of isomorphism classes of}}\\
&\qquad\ \ \text{R-}\underline{\text{algebra extensions of}} \ \ A \ \ \underline{\text{by}} \ \ M \ ,\\
D^2(A|R, \ M) \ &= \ \mathrm{Coker}(H^1(K_I, \ M) \ \to \ \mathrm{Hom}_A(H_1(K_I), \ M)) \ .
\end{aligned}
\end{equation*}

}{\begin{equation*}
\begin{aligned}
D^1(A|R, \ M) \ &= \ \mathrm{Coker}(\mathrm{Der}_R(P, \ M) \ \to \ \mathrm{Hom}_A(I/I^2, \ M))\\
&= \ \underline{\text{the set of isomorphism classes of}}\\
&\qquad\ \ \text{R-}\underline{\text{algebra extensions of}} \ \ A \ \ \underline{\text{by}} \ \ M \ ,\\
D^2(A|R, \ M) \ &= \ \mathrm{Coker}(H^1(K_I, \ M) \ \to \ \mathrm{Hom}_A(H_1(K_I), \ M)) \ .
\end{aligned}
\end{equation*}

}{F7I-VI-067}\underline{Proof}:

(1) If  $0 \to M' \to M \to M'' \to 0$  is an exact sequence of $A$-modules, then for any polynomial algebra  $P$  over  $R$  ,  $0 \to \mathrm{Der}_R(P, M') \to \mathrm{Der}_R(P, M) \to \mathrm{Der}_R(P, M'') \to 0$ is exact and hence  $0 \to \mathrm{Der}_R(-, M') \to \mathrm{Der}_R(-, M) \to \mathrm{Der}_R(-, M'') \to 0$  is an exact sequence of sheaves by 3.3(c).

(2) $\mathrm{Der}_S(-, M)$  is an additive sheaf on  $\underline{C}_{A|S}$  , and $\mathrm{Der}_R(B, M) = \mathrm{Ker}(\mathrm{Der}_S(B, M) \to \mathrm{Der}_S(R, M))$  for every  $B \in \mathrm{ob}(\underline{C}_{A|R})$  . Therefore our assertion follows from 3.7(c).

(3) Let  $S : \underline{C}_{A|R} \to \underline{C}_{A \underset{R}{\otimes} S|S}$  be the functor given by  $X \longmapsto X \underset{R}{\otimes} S$  . Then $\widetilde{S}\mathrm{Der}_S(-, M) = \mathrm{Der}_R(-, M)$  and hence our assertion follows from 3.7(a).

\SGApage{I}{88}

(4)  Follows from 3.9. Indeed,  $F = \mathrm{Der}_R(-, N)$  is an additive sheaf on $\underline{C}_{S^{-1}A|R}$  , and  $F_A \in \mathrm{ob}\left(\underline{C}^{\sim}_{S^{-1}A|A}\right)$  is also additive since $F_A(X) = \mathrm{Ker}(F(X) \to F(A)) = \mathrm{Ker}(\mathrm{Der}_R(X, N) \to \mathrm{Der}_R(A, N)) = \mathrm{Der}_A(X, N)$  i.e., $F_A = \mathrm{Der}_A(-, N)$  .

(5)  The first assertion follows from 3.8(a). Now assume that $\mathrm{Tor}^R_i(A, A) = 0$  for all  $i > 0$  and that

\begin{equation*}
\Delta \ : \ \mathrm{Spec}(A) \ \to \ \mathrm{Spec}(A) \underset{\mathrm{Spec}(R)}{\times} \mathrm{Spec}(A)
\end{equation*}
is an open immersion i.e.,  $(A \underset{R}{\otimes} A)_{\varphi^{-1}(\underline{m})}$  $\to A_{\underline{m}}$  is an isomorphism for all maximal ideals  $\underline{m}$  of  $A$  , where  $\varphi : A \underset{R}{\otimes} A \to A$  is the multiplication map. Then the first assertion combined with 3.12(4) entails that the diagonal map

\begin{equation*}
\Delta_{\underline{m}} \ : \ D^i(A|R, \ M)_{\underline{m}} \ \to \ D^i(A|R, \ M)_{\underline{m}} \ \oplus \ D^i(A|R, \ M)_{\underline{m}}
\end{equation*}
is an isomorphism for every maximal ideal  $\underline{m}$  of  $A$  and hence

\begin{equation*}
\Delta \ : \ D^i(A|R, \ M) \ \to \ D^i(A|R, \ M) \ \oplus \ D^i(A|R, \ M)
\end{equation*}
is an isomorphism, which entails that  $D^i(A|R, M) = 0$  for all  $i$  .

(6)  In view of 3.12(4), we may assume that  $A$  is an R-algebra of finite type. Now consider the map

\begin{equation*}
E \underset{A}{\otimes} \mathrm{Der}_R(-, \ M) \ \to \ \mathrm{Der}_R(-, \ E \underset{A}{\otimes} M)
\end{equation*}
in  $\underline{C}^{\sim}_{A|R}$  . Since  $E$  is A-flat, we have an isomorphism

\begin{equation*}
E \underset{A}{\otimes} \mathrm{Der}_R(B, \ M) \ \overset{\sim}{\longrightarrow} \ \mathrm{Der}_R(B, \ E \underset{A}{\otimes} M)
\end{equation*}
whenever  $B \in \mathrm{ob}(\underline{C}_{A|R})$  is of finite type over  $R$  . It follows from 3.3(b) that

\begin{equation*}
D^i(A|R, \ E \underset{A}{\otimes} \mathrm{Der}_R(-, \ M)) \ \to \ D^i(A|R, \ E \underset{A}{\otimes} M)
\end{equation*}
is an isomorphism. On the other hand, the functor  $E \underset{A}{\otimes} -$  :  $(A\text{-mod}) \to (A\text{-mod})$  is exact, and therefore we have

\begin{equation*}
E \underset{A}{\otimes} D^i(A|R, \ M) \ = \ D^i(A|R, \ E \underset{A}{\otimes} \mathrm{Der}_R(-, \ M)) \ \ ,
\end{equation*}
and hence our assertion.

(7)  follows from 3.3(a).


\SGApage{I}{89}

(8)  In view of 3.10 and 3.11, it suffices to show that $\mathrm{Coker}(\mathrm{Der}_R(P, M) \to \mathrm{Hom}_A(I/I^2, M)) = \mathrm{Exalcomm}(A|R, M)$ , where  $\mathrm{Exalcomm}(A|R, M)$ denotes the set of isomorphism classes of R-algebra extensions of  $A$  by  $M$ . Let $0 \to M \to A' \to A \to 0$   be an R-algebra extension of  $A$  by  $M$ . Since  $P$  is a polynomial algebra over  $R$ , we obtain a commutative diagram

\begin{center}
\begin{tikzcd}[column sep=large]
0 \arrow[r] \& I \arrow[r] \arrow[d] \& P \arrow[r] \arrow[d] \& A \arrow[r] \arrow[d, equal] \& 0 \\
0 \arrow[r] \& M \arrow[r] \& A' \arrow[r] \& A \arrow[r] \& 0
\end{tikzcd}
\end{center}











\SGAvar{where  $I \to M$  is uniquely determined by the extensions  $A'$  up to  $\mathrm{Der}_R(P, M)$  . We thus obtain a map

}{where  $I \to M$  is uniquely determined by the extension  $A'$  up to  $\mathrm{Der}_R(P, M)$  . We thus obtain a map

}{F7I-VI-068}\begin{equation*}
\Phi \ : \ \mathrm{Exalcomm}(A|R, M) \to \mathrm{Coker}(\mathrm{Der}_R(P, M) \to \mathrm{Hom}_A(I/I^2, M))  \ .
\end{equation*}
Conversely a P-linear map  $I \to M$  yields an extension by push-out

\begin{center}
\begin{tikzcd}[column sep=large]
0 \arrow[r] \& I \arrow[r] \arrow[d] \& P \arrow[r] \arrow[d] \& A \arrow[r] \arrow[d, equal] \& 0 \\
0 \arrow[r] \& M \arrow[r] \& A' \arrow[r] \& A \arrow[r] \& 0
\end{tikzcd}
\end{center}










\SGAvar{Two P-linear maps  $I \rightrightarrows M$  which differ by  $\mathrm{Der}_R(P, M)$  is trivially seen to yield isomorphic extensions and therefore we obtain a map

\begin{equation*}
\Psi \ : \ \mathrm{Coker}(\mathrm{Der}_R(P, M) \to \mathrm{Hom}_A(I/I^2, M)) \to \mathrm{Exalcomm}(A|R, M)  \ .
\end{equation*}
It is straightforward to verify that  $\Phi$, $\Psi$  are inverse process to each other.

\underline{Corollary} 3.13.  (a)  R-\underline{algebra}  $A$  \underline{is formally smooth over}  $R$  \underline{if and only if} $D^1(A|R, M) = 0$  \underline{for every} A-\underline{module}  $M$  .

}{Two P-linear maps  $I \rightrightarrows M$  which differ by  $\mathrm{Der}_R(P, M)$  are trivially seen to yield isomorphic extensions and therefore we obtain a map

\begin{equation*}
\Psi \ : \ \mathrm{Coker}(\mathrm{Der}_R(P, M) \to \mathrm{Hom}_A(I/I^2, M)) \to \mathrm{Exalcomm}(A|R, M)  \ .
\end{equation*}
It is straightforward to verify that  $\Phi$, $\Psi$  are inverse processes to each other.

\underline{Corollary} 3.13.  (a)  The R-\underline{algebra}  $A$  \underline{is formally smooth over}  $R$  \underline{if and only if} $D^1(A|R, M) = 0$  \underline{for every} A-\underline{module}  $M$  .

}{F7I-VI-069}(b)  \underline{Let}  $R$  \underline{be noetherian and}  $A$  \underline{an} R-\underline{algebra of finite type.} \underline{Then}  $A$ \underline{is a relative complete-intersection over}  $R$  \underline{if and only if}  $D^2(A|R, M) = 0$  \underline{for every} A-\underline{module}  $M$  .

(c)  \underline{Let}  $R$  \underline{be noetherian and}  $A$  \underline{an} R-\underline{algebra of finite type.} \underline{Then for} \underline{any multiplicative subset}  $S$  \underline{in}  $A$  , \underline{we have a canonical isomorphism}

\begin{equation*}
D^i(S^{-1}A|R, \ S^{-1}M) \xrightarrow{\ \sim\ } S^{-1}D^i(A|R, M)  \ \underline{\text{for all}} \ \ i \geq 0  \ .
\end{equation*}

\SGApage{I}{90}

(d)  \underline{Let}  $R$  \underline{be noetherian and}  $A$  \underline{a local} R-\underline{algebra of essentially finite} \underline{type. We denote by}  $\hat{A}$  \underline{the} \underline{m}-\underline{adic completion of}  $A$  \underline{where}  $\underline{m}$ = \underline{the maximal ideal of} $A$  .  \underline{Then for any} A-\underline{module}  $E$  \underline{of finite type we have a canonical isomorphism}

\begin{equation*}
D^1(\hat{A}|R, \ \hat{A} \underset{A}{\otimes} E) \xrightarrow{\ \sim\ } \hat{A} \underset{A}{\otimes} D^1(A|R, E)  \ .
\end{equation*}
\underline{Proof}:  (a)  follows immediately from  $D^1(A|R, M) = \mathrm{Exalcomm}(A|R, M)$ .

(b)  Let  $P$  be a polynomial algebra over  $R$  in a finite number of variables and let  $0 \to I \to P \to A \to 0$  be exact. Localizing if necessary, we may assume that  $I$  is generated by a P-regular sequence which would imply that  $H_1(K_I) = 0$ and therefore  $D^2(A|R, M) = \mathrm{Coker}(H^1(K_I, M) \to \mathrm{Hom}_A(H_1(K_I), M)) = 0$  . Conversely assume that  $D^2(A|R, M) = 0$  for every A-module  $M$  .

Let  $0 \to I \to P \to A \to 0$  be exact where  $P$  is a polynomial algebra over $R$  in a finite number of variables. We have

\SGAvar{\begin{equation*}
0 = D^2(A|R, M) = \mathrm{Coker}(H^1(K_I, M) \to \mathrm{Hom}_A(H_1(K_I), M))
\end{equation*}
i.e.,  $H^1(K_I, M) \to \mathrm{Hom}_A(H_1(K_I), M)$  is surjective for all A-module  $M$  . Choose an exact sequence  $0 \to K \to F \to I \to 0$  where  $F$  is a free P-module of finite type. We note that  $H_1(K_I) = K/K_0$  where  $K_0 = \mathrm{Im}(\overset{2}{\wedge} F \to \overset{1}{\wedge} F)$  . The fact that $D^2(A|R, H_1(K_I)) = 0$  i.e.,

\begin{equation*}
H^1(K_I, \ H_1(K_I)) \to \mathrm{Hom}_A(H_1(K_I), \ H_1(K_I))  \ \text{ is surjective}
\end{equation*}
entails that  $0 \to K/K_0 \to F/IF \to I/I^2 \to 0$  is exact and splits as A-modules and in particular  $I/I^2$  is projective as A-module. Therefore, localizing  $A$  if necessary, we may assume that  $I/I^2$  is  A-free, and that  $F/IF \to I/I^2$  is an isomorphism. Then the fact that  $0 \to K/K_0 \to F/IF \to I/I^2 \to 0$  is exact and splits as A-modules entails that  $H_1(K_I) = K/K_0 = 0$  . This means that  $I$  is generated by a P-regular sequence i.e.,  $A$  is a relative complete intersection over  $R$  .

}{\begin{equation*}
0 = D^2(A|R, M) = \mathrm{Coker}(H^1(K_I, M) \to \mathrm{Hom}_A(H_1(K_I), M))
\end{equation*}
i.e.,  $H^1(K_I, M) \to \mathrm{Hom}_A(H_1(K_I), M)$  is surjective for every A-module  $M$  . Choose an exact sequence  $0 \to K \to F \to I \to 0$  where  $F$  is a free P-module of finite type. We note that  $H_1(K_I) = K/K_0$  where  $K_0 = \mathrm{Im}(\overset{2}{\wedge} F \to \overset{1}{\wedge} F)$  . The fact that $D^2(A|R, H_1(K_I)) = 0$  i.e.,

\begin{equation*}
H^1(K_I, \ H_1(K_I)) \to \mathrm{Hom}_A(H_1(K_I), \ H_1(K_I))  \ \text{ is surjective}
\end{equation*}
entails that  $0 \to K/K_0 \to F/IF \to I/I^2 \to 0$  is exact and splits as A-modules and in particular  $I/I^2$  is projective as an A-module. Therefore, localizing  $A$  if necessary, we may assume that  $I/I^2$  is  A-free, and that  $F/IF \to I/I^2$  is an isomorphism. Then the fact that  $0 \to K/K_0 \to F/IF \to I/I^2 \to 0$  is exact and splits as A-modules entails that  $H_1(K_I) = K/K_0 = 0$  . This means that  $I$  is generated by a P-regular sequence i.e.,  $A$  is a relative complete intersection over  $R$  .

}{F7I-VI-070}(c)  follows immediately from 3.12(4) and (6).

(d)  Firstly it follows from 3.12(6) that

\SGAvar{\begin{equation*}
D^i(A\;R, \ \hat{A}\underset{A}{\otimes}E) \xrightarrow{\ \sim\ } \hat{A}\underset{A}{\otimes}D^i(A|R, E)
\end{equation*}
is an isomorphism for all  $i$  , and therefore it suffices to show that
\SGApage{I}{91}
$D^1(\hat{A}|R,\ \hat{A}\otimes_A E) \to D^1(A|R,\ \hat{A}\otimes_A E)$  is an isomorphism. Our proof is based on the description of  $D^1$  in terms of algebra extensions. Let  $0 \to \hat{A}\otimes_A E \to A' \xrightarrow{\ \varphi\ } A \to 0$  be an R-algebra extension. The fact that  $\hat{A}\otimes_A E$  is an $\hat{A}$-module of finite type entails readily that  $A'$  is a local R-algebra, complete with respect to $\underline{m}'$-adic topology where  $\underline{m}'$ = the maximal ideal of  $A'$ . Therefore if there is an R-algebra map  $\psi : A \to A'$  such that  $\varphi \circ \psi = \mathrm{id}_A$ , then  $\psi$  extends uniquely to  $\hat{\psi} : \hat{A} \to A'$  such that  $\varphi \circ \hat{\psi} = \mathrm{id}_{\hat{A}}$ . This proves that

\begin{equation*}
D^1(\hat{A}|R,\ \hat{A}\otimes_A E) \to D^1(A|R,\ \hat{A}\otimes_A E)
\end{equation*}
is injective. Now let  $0 \to \hat{A}\otimes_A E \to B \xrightarrow{\ \varphi\ } A \to 0$  be an R-algebra extension. Firstly we claim that the topology on  $\hat{E} = \hat{A}\otimes_A E$  induced from  $\underline{m}_B$-adic topology on  $B$  coincides with the topology on  $\hat{E}$  as $\hat{A}$-module. Indeed, since  $A$  is essentially of finite type over  $R$ , one can find a subalgebra  $B_0$  of  $B$  such that  $B_0$  is a local R-algebra of essentially finite type over  $R$  and that  $\varphi(B_0) = A$ . Set  $E_0 = E \cap B_0$  and denote by  $\underline{m}_B$ , $\underline{m}_0$ , $\underline{m}$  the maximal ideals of  $B$ , $B_0$ , $A$  respectively. Then  $\underline{m}_B = \underline{m}_0 + \hat{E}$  and hence  $\underline{m}_B^{n+1} = \underline{m}_0^{n+1} + \underline{m}^n\hat{E}$  for all  $n$ . Since  $B_0$  is noetherian, there exists an integer  $r > 0$  by Artin-Rees Theorem such that

\begin{equation*}
\begin{aligned}
\hat{E} \cap \underline{m}_B^{n+1} &= \underline{m}^n\hat{E} + \hat{E} \cap \underline{m}_0^{n+1} = \underline{m}^n\hat{E} + E_0 \cap \underline{m}_0^{n+1}\\
&= \underline{m}^n\hat{E} + \underline{m}^{n+1-r}(E_0 \cap \underline{m}_0^{\,r}) \subset \underline{m}^{n+1-r}\hat{E}
\end{aligned}
\end{equation*}
for all  $n \geq r$ , which proves our assertion about the two topologies on  $E$  . In particular  $\hat{E} = \hat{A}\otimes_A E$  is complete with respect to the topology induced from $\underline{m}_B$-adic topology on  $B$ , and hence we obtain the exact commutative diagram

}{\begin{equation*}
D^i(A|R, \ \hat{A}\underset{A}{\otimes}E) \xrightarrow{\ \sim\ } \hat{A}\underset{A}{\otimes}D^i(A|R, E)
\end{equation*}
is an isomorphism for all  $i$  , and therefore it suffices to show that
\SGApage{I}{91}
$D^1(\hat{A}|R,\ \hat{A}\otimes_A E) \to D^1(A|R,\ \hat{A}\otimes_A E)$  is an isomorphism. Our proof is based on the description of  $D^1$  in terms of algebra extensions. Let  $0 \to \hat{A}\otimes_A E \to A' \xrightarrow{\ \varphi\ } \hat{A} \to 0$  be an R-algebra extension whose restriction to $A$ splits. The fact that  $\hat{A}\otimes_A E$  is an $\hat{A}$-module of finite type entails readily that  $A'$  is a local R-algebra, complete with respect to $\underline{m}'$-adic topology where  $\underline{m}'$ = the maximal ideal of  $A'$ . Therefore if there is an R-algebra map  $\psi : A \to A'$  lifting the canonical map $A \to \hat{A}$ , then  $\psi$  extends uniquely to  $\hat{\psi} : \hat{A} \to A'$  such that  $\varphi \circ \hat{\psi} = \mathrm{id}_{\hat{A}}$ . This proves that

\begin{equation*}
D^1(\hat{A}|R,\ \hat{A}\otimes_A E) \to D^1(A|R,\ \hat{A}\otimes_A E)
\end{equation*}
is injective. Now let  $0 \to \hat{A}\otimes_A E \to B \xrightarrow{\ \varphi\ } A \to 0$  be an R-algebra extension. Firstly we claim that the topology on  $\hat{E} = \hat{A}\otimes_A E$  induced from  $\underline{m}_B$-adic topology on  $B$  coincides with the topology on  $\hat{E}$  as $\hat{A}$-module. Indeed, since  $A$  is essentially of finite type over  $R$ , one can find a subalgebra  $B_0$  of  $B$  such that  $B_0$  is a local R-algebra of essentially finite type over  $R$  and that  $\varphi(B_0) = A$ . Set  $E_0 = \hat{E} \cap B_0$  and denote by  $\underline{m}_B$ , $\underline{m}_0$ , $\underline{m}$  the maximal ideals of  $B$ , $B_0$ , $A$  respectively. Then  $\underline{m}_B = \underline{m}_0 + \hat{E}$  and hence  $\underline{m}_B^{n+1} = \underline{m}_0^{n+1} + \underline{m}^n\hat{E}$  for all  $n$ . Since  $B_0$  is noetherian, there exists an integer  $r > 0$  by Artin-Rees Theorem such that

\begin{equation*}
\begin{aligned}
\hat{E} \cap \underline{m}_B^{n+1} &= \underline{m}^n\hat{E} + \hat{E} \cap \underline{m}_0^{n+1} = \underline{m}^n\hat{E} + E_0 \cap \underline{m}_0^{n+1}\\
&= \underline{m}^n\hat{E} + \underline{m}^{n+1-r}(E_0 \cap \underline{m}_0^{\,r}) \subset \underline{m}^{n+1-r}\hat{E}
\end{aligned}
\end{equation*}
for all  $n \geq r$ , which proves our assertion about the two topologies on  $\hat{E}$  . In particular  $\hat{E} = \hat{A}\otimes_A E$  is complete with respect to the topology induced from $\underline{m}_B$-adic topology on  $B$ , and hence we obtain the exact commutative diagram

}{F7I-VI-071}\begin{center}
\begin{tikzcd}
0 \arrow[r] \& \hat{E} \arrow[r] \arrow[d, equal] \& B \arrow[r] \arrow[d] \& A \arrow[r] \arrow[d] \& 0 \\
0 \arrow[r] \& \hat{E} \arrow[r] \& \hat{B} \arrow[r] \& \hat{A} \arrow[r] \& 0
\end{tikzcd}
\end{center}












\SGAvar{where  $\hat{B}$  stands for  $\underline{m}_B$-adic completion of  $B$  . This proves that  $D^1$
$D^1(\hat{A}|R,\ \hat{E}) \to D^1(A|R,\ \hat{E})$  is surjective and hence we are done.


}{where  $\hat{B}$  stands for  $\underline{m}_B$-adic completion of  $B$  . This proves that
$D^1(\hat{A}|R,\ \hat{E}) \to D^1(A|R,\ \hat{E})$  is surjective and hence we are done.


}{F7I-VI-072}\SGApage{I}{92}

\underline{Corollary} 3.14. \underline{Let}  $R$  \underline{be noetherian and}  $A$  \underline{an R-algebra of finite type, smooth everywhere except at one closed point}  $x \in \mathrm{Spec}(A)$  . \underline{We then have isomorphisms}

\begin{equation*}
D^1(A|R,\ A) \simeq D^1(A_x|R,\ A_x) \simeq D^1(\hat{A}_x|R,\ \hat{A}_x)\ .
\end{equation*}
\underline{Proof}: Since  $A$  is smooth over  $R$  except at  $x \in \mathrm{Spec}(A)$  ,  $D^1(A|R,\ A)$  is an A-module of finite type with its support on the set  $\{x\}$  , and therefore we have

\begin{equation*}
D^1(A|R,\ A) = D^1(A|R,\ A)_x = \widehat{D^1(A|R,\ A)}_x\ .
\end{equation*}
Therefore our assertion follows from 3.13.

\SGAvar{\underline{Remark} 3.15.  Let  $R$  be topological ring and  $A$  a topological R-algebra. We then define, for any A-module  $M$  ,

\begin{equation*}
D^i_{\mathrm{top}}(A|R,\ E) = \varinjlim_{\alpha} D^i(A_\alpha|R_\alpha,\ A_\alpha \otimes_\alpha E)
\end{equation*}
where  $A_\alpha = A/J_\alpha$ , $R_\alpha = R/I_\alpha$  , and  $J_\alpha$ , $I_\alpha$  are both open in  $A, R$  respectively. Then  $A$  is formally smooth topological R-algebra means that  $D^1_{\mathrm{top}}(A|R,\ E) = 0$  for every A-module  $E$  .

}{\underline{Remark} 3.15.  Let  $R$  be a topological ring and  $A$  a topological R-algebra. We then define, for any A-module  $E$  ,

\begin{equation*}
D^i_{\mathrm{top}}(A|R,\ E) = \varinjlim_{\alpha} D^i(A_\alpha|R_\alpha,\ A_\alpha \otimes_A E)
\end{equation*}
where  $A_\alpha = A/J_\alpha$ , $R_\alpha = R/I_\alpha$  , and  $J_\alpha$ , $I_\alpha$  are both open in  $A, R$  respectively. Then to say that  $A$  is formally smooth as a topological R-algebra means that  $D^1_{\mathrm{top}}(A|R,\ E) = 0$  for every A-module  $E$  .

}{F7I-VI-073}        Now let  $X$  be a scheme over a ring  $R$  and  $\underline{E}$  a quasi-coherent  X-module. We then define the sheaves  $\underline{D}^i_{X|R}(\underline{E}) = \underline{D}^i(X|R,\ \underline{E})$  on  $X$  via presheaves on  $X$  given by

\SGAvar{\begin{equation*}
[\underline{D}^i_{X|R}(\underline{E})](U) = D^i(\Gamma(U,\ \underline{O}_X)|R,\ \Gamma(U,\ \underline{E}))
\end{equation*}
for every affine open set  $U \subset X$  . If  $R$  is noetherian and  $X$  is of finite type over  $R$  , then the property 3.13(a) shows that  $\underline{D}^i_{X|R}(\underline{E})$  can actually be constructed canonically, as a quasi-coherent X-module.

        In case when there is no possible confusion we shall abreviate by setting  $\underline{D}^i_{X|R} = \underline{D}^i_{X|R}(\underline{O}_X)$  . All the statements of 3.13 can now be translated in terms of these quasi-coherent sheaves  $\underline{D}^i_{X|R}(\underline{E})$  , which we leave to the readers. We list below a few properties which is most relevant to our purpose.

}{\begin{equation*}
[\underline{D}^i_{X|R}(\underline{E})](U) = D^i(\Gamma(U,\ \underline{O}_X)|R,\ \Gamma(U,\ \underline{E}))
\end{equation*}
for every affine open set  $U \subset X$  . If  $R$  is noetherian and  $X$  is of finite type over  $R$  , then the property 3.13(c) shows that  $\underline{D}^i_{X|R}(\underline{E})$  can actually be constructed canonically, as a quasi-coherent X-module.

        In case when there is no possible confusion we shall abbreviate by setting  $\underline{D}^i_{X|R} = \underline{D}^i_{X|R}(\underline{O}_X)$  . All the statements of 3.13 can now be translated in terms of these quasi-coherent sheaves  $\underline{D}^i_{X|R}(\underline{E})$  , which we leave to the readers. We list below a few properties which are most relevant to our purpose.

}{F7I-VI-074}\underline{Corollary} 3.15. \underline{Let}  $R$  \underline{be noetherian and}  $X$  \underline{an} R-\underline{scheme of finite type. Then}
\SGApage{I}{93}
(a)  \underline{For any coherent} X-\underline{module}  $\underline{E}$ , $\underline{D}^i_{X|R}(\underline{E})$  \underline{is again a coherent} X-\underline{module} \underline{for all}  $i$ , \underline{and furthermore}  $\operatorname{Supp} \underline{D}^i_{X|R}(\underline{E}) \subset X^{\mathrm{Sing}}$  \underline{for all}  $i > 0$  \underline{where} $X^{\mathrm{Sing}} = \{x \in X| X$ \underline{is non-smooth over}  $R$  \underline{at the point}  $x \in X\}$ .

(b)  \underline{If}  $X$  \underline{is flat over}  $R$ , \underline{then for any} R-\underline{algebra}  $S$ , \underline{we have the} \underline{canonical isomorphism}

\begin{equation*}
\underline{D}^i_{X\underset{R}{\otimes}S\,|\,S}(\underline{E}) \overset{\sim}{\to} \underline{D}^i_{X|R}(\underline{E})
\end{equation*}
\underline{for all}  $i$ .

(c)  \underline{If}  $X$  \underline{is a relative complete intersection over}  $R$ , \underline{then for every} \underline{surjective ring homomorphism}  $S \to R$  \underline{we have the exact sequence}

\begin{equation*}
0 \to \underline{D}^1_{X|R}(\underline{E}) \to \underline{D}^1_{X|S}(\underline{E}) \to \underline{\mathrm{Hom}}_{\underline{O}_X}(I/I^2\underset{R}{\otimes}\underline{O}_X, \underline{E}) \to 0
\end{equation*}
\underline{where}  $I = \mathrm{Ker}(S \to R)$ . \underline{If, furthermore,}  $R$  \underline{is local with the residue field}  $k$ \underline{and}  $X$  \underline{is} R-\underline{flat,} \underline{then we have the exact sequence}

\SGAvar{\begin{equation*}
0 \to \underline{D}^1_{X_0|h}(\underline{O}_{X_0}) \to \underline{D}^1_{X|S}(\underline{O}_{X_0}) \to D^1(R|S,\ k)\underset{k}{\otimes}\underline{O}_{X_0} \to 0
\end{equation*}
\underline{where}  $X_0 = X\underset{R}{\otimes}k$ .

\underline{Proof}:  (a)  Follows from 3.12(6) and (8) whereas (b) follows from 3.12(3).

(c)  Since  $X$  is a relative complete-intersection over  $R$ , we have $\underline{D}^i_{X|R}(\underline{E}) = 0$  for all  $i \geq 2$  by 3.12(10) and hence

}{\begin{equation*}
0 \to \underline{D}^1_{X_0|k}(\underline{O}_{X_0}) \to \underline{D}^1_{X|S}(\underline{O}_{X_0}) \to D^1(R|S,\ k)\underset{k}{\otimes}\underline{O}_{X_0} \to 0
\end{equation*}
\underline{where}  $X_0 = X\underset{R}{\otimes}k$ .

\underline{Proof}:  (a)  Follows from 3.12(7) and the vanishing for smooth algebras (cf. 3.13(a)), whereas (b) follows from 3.12(3).

(c)  Since  $X$  is a relative complete-intersection over  $R$ , we have $\underline{D}^i_{X|R}(\underline{E}) = 0$  for all  $i \geq 2$  by the complete-intersection vanishing (cf. 3.13(b)) and hence

}{F7I-VI-075}\begin{equation*}
0 \to \underline{D}^0_{X|R}(\underline{E}) \to \underline{D}^0_{X|S}(\underline{E}) \to \underline{D}^0_{R|S}(\underline{E}) \to \underline{D}^1_{X|R}(\underline{E}) \to \underline{D}^1_{X|S}(\underline{E}) \to \underline{D}^1_{R|S}(\underline{E}) \to 0
\end{equation*}
is exact by 3.12(2), where  $\underline{D}^i_{R|S}(\underline{E})$  are quasi-coherent sheaves of  $X$  defined by

\begin{equation*}
[\underline{D}^i_{R|S}(\underline{E})](U) = D^i(R|S,\ \underline{E}(U))
\end{equation*}
for every affine open subset  $U \subset X$ . In particular,  $\underline{D}^0_{R|S}(\underline{E}) = 0$  since  $S \to R$ is surjective, and

\SGAvar{\begin{equation*}
\underline{D}^1_{R|S}(\underline{E}) = \underline{\mathrm{Hom}}_{\underline{O}_X}(I/I^2\underset{R}{\otimes}\underline{O}_X, \underline{E})
\end{equation*}
by 3.12(7), and hence the first exact sequence. If  $X$  is R-flat, we have
\SGApage{I}{94}
$\underline{D}^i_{X_0|k}(\underline{E}) \overset{\sim}{\to} \underline{D}^i_{X|R}(\underline{E})$  for all  $i$  and for all quasi-coherent $X_0$-module  $\underline{E}$  and in particular

}{\begin{equation*}
\underline{D}^1_{R|S}(\underline{E}) = \underline{\mathrm{Hom}}_{\underline{O}_X}(I/I^2\underset{R}{\otimes}\underline{O}_X, \underline{E})
\end{equation*}
by 3.12(8), and hence the first exact sequence. If  $X$  is R-flat, we have
\SGApage{I}{94}
$\underline{D}^i_{X_0|k}(\underline{E}) \overset{\sim}{\to} \underline{D}^i_{X|R}(\underline{E})$  for all  $i$  and for every quasi-coherent $X_0$-module  $\underline{E}$  and in particular

}{F7I-VI-076}\begin{equation*}
\underline{D}^1_{X_0|k}(\underline{O}_{X_0}) \overset{\sim}{\to} \underline{D}^1_{X|R}(\underline{O}_{X_0})
\end{equation*}
from which the second exact sequence follows.

\SGAvar{\underline{Remark} 3.16:  The quasi-coherent sheaves  $\underline{D}^i_{X|R}(\underline{E})$  we have defined should appear as  $E_2^{0,i}$  of a spectral sequence relating the local ones and the global ones. L. Illusie has obtained such a theory (Lecture Notes in Mathematics 239). His method is via  "homotopical algebra"  generalizing the methods of M. Andre and D. Quillen rather than  "cohomological algebra" . It seems that the approach we have taken here can be globalized via Grothendieck cohomology on the category of schemes with an appropriate topology generalizing the one on the affine category.


}{\underline{Remark} 3.16:  The quasi-coherent sheaves  $\underline{D}^i_{X|R}(\underline{E})$  we have defined should appear as  $E_2^{0,i}$  of a spectral sequence relating the local ones and the global ones. L. Illusie has obtained such a theory (Lecture Notes in Mathematics 239). His method is via  "homotopical algebra"  generalizing the methods of M. André and D. Quillen rather than  "cohomological algebra" . It seems that the approach we have taken here can be globalized via Grothendieck cohomology on the category of schemes with an appropriate topology generalizing the one on the affine category.


}{F7I-VI-077}\SGAsectiontoc{4.  Formal moduli of deformations}{4}
4.  \underline{Formal moduli of deformations}

One amongst the nice cofibered categories in which the isomorphism classes of objects are usually not prorepresentable, but nevertheless admit a formal versal object is the cofibered category of deformations. We now come back to the situations and notations of paragraph 1. Thus, throughout the section,  $\Lambda$  is a fixed complete local noetherian ring with the residue field  $k$ , and  $\underline{C}_{\Lambda}$  is the category of artinian local $\Lambda$-algebras with the same residue field  $k$ .

Let  $X_0$  be a fixed scheme over the field  $k$ . By an (infinitesimal) deformation of  $X_0$ , we mean a pair  $(R, X)$  where  $R \in \mathrm{ob}(\underline{C}_{\Lambda})$  and  $X$  is a flat R-scheme together with a commutative diagram

\begin{tikzcd}[column sep=huge, row sep=huge]
X_0 \arrow[r, "i_X"] \arrow[d] \& X \arrow[d, "\mathrm{flat}"] \\
\mathrm{Spec}(k) \arrow[r] \& \mathrm{Spec}(R)
\end{tikzcd}












\SGAvar{\SGApage{I}{95}
such that the induced morphism $X_0 \to X \times_{\mathrm{spec}(R)} \mathrm{Spec}(k)$ is an isomorphism. We note that, since $R$ is an artinian local ring, $i_X : X_0 \to X$ is necessarily a closed immersion which is a homeomorphism on underlying topological spaces. If $(R, X)$ , $(R', X')$ are deformations of $X_0$ to $R, R'$ respectively, we define an arrow $(R', X') \to (R, X)$ to consist of a map $R' \xrightarrow{\ \varphi\ } R$ in $\underline{C}_\Lambda$ and a morphism $\Phi : X \to X'$ of schemes with a commutative diagram

}{\SGApage{I}{95}
such that the induced morphism $X_0 \to X \times_{\mathrm{Spec}(R)} \mathrm{Spec}(k)$ is an isomorphism. We note that, since $R$ is an artinian local ring, $i_X : X_0 \to X$ is necessarily a closed immersion which is a homeomorphism on underlying topological spaces. If $(R, X)$ , $(R', X')$ are deformations of $X_0$ to $R, R'$ respectively, we define an arrow $(R', X') \to (R, X)$ to consist of a map $R' \xrightarrow{\ \varphi\ } R$ in $\underline{C}_\Lambda$ and a morphism $\Phi : X \to X'$ of schemes with a commutative diagram

}{F7I-VI-078}\begin{center}
\begin{tikzcd}[column sep=huge, row sep=large]
X \arrow[rr, "\Phi"] \arrow[ddd] \& \& X' \arrow[ddd] \\
 \& X_0 \arrow[ul, "i_X"'] \arrow[ur, "i_{X'}"] \arrow[d] \& \\
 \& \mathrm{Spec}(k) \arrow[dl] \arrow[dr] \& \\
\mathrm{Spec}(R) \arrow[rr, "\mathrm{Spec}(\varphi)"'] \& \& \mathrm{Spec}(R')
\end{tikzcd}
\end{center}












\SGAvar{We thus obtain the category $\underline{M}_{X_0}$ of (infinitesimal) deformations of $X_0$ in which the objects and the arrows are defined as above. The category $\underline{M}_{X_0}$ is provided with a functor $p : \underline{M}_{X_0} \to \underline{C}_\Lambda$ defined by $(R, X) \to R$ . Now let $(R, X) \in \mathrm{ob}(\underline{M}_{X_0})$ and a map $R \xrightarrow{\ \varphi\ } S$ in $\underline{C}_\Lambda$ be given. We then obtain a commutative diagram

}{We thus obtain the category $\underline{M}_{X_0}$ of (infinitesimal) deformations of $X_0$ in which the objects and the arrows are defined as above. The category $\underline{M}_{X_0}$ is provided with a functor $P : \underline{M}_{X_0} \to \underline{C}_\Lambda$ defined by $(R, X) \to R$ . Now let $(R, X) \in \mathrm{ob}(\underline{M}_{X_0})$ and a map $R \xrightarrow{\ \varphi\ } S$ in $\underline{C}_\Lambda$ be given. We then obtain a commutative diagram

}{F7I-VI-079}\begin{center}
\begin{tikzcd}[column sep=huge, row sep=huge]
X \times_{\mathrm{Spec}(R)} \mathrm{Spec}(S) \arrow[r, "p_1"] \arrow[d, "\mathrm{flat}"'] \& X \arrow[d, "\mathrm{flat}"] \\
\mathrm{Spec}(S) \arrow[r, "\mathrm{Spec}(\varphi)"'] \& \mathrm{Spec}(R)
\end{tikzcd}
\end{center}









and $\varphi$-arrow $(\varphi, p_1) : (R, X) \to (S, X \times_{\mathrm{Spec}(R)} \mathrm{Spec}(S))$ is obviously cocartesian in $\underline{M}_{X_0}$ . It follows that the category $\underline{M}_{X_0}$ is a cofibered category over $\underline{C}_\Lambda$ via $P : \underline{M}_{X_0} \to \underline{C}_\Lambda$ . We observe the following facts:


\SGApage{I}{96}

4.1(a): Let $(1_R, \Phi) : (R, X_1) \to (R, X_2)$ be an arrow in $\underline{M}_{X_0}$ i.e., $\Phi : X_2 \to X_1$ is a morphism over $R$ such that $X_2 \otimes_R k \to X_1 \otimes_R k$ is an isomorphism. Since $R$ is artinian and $X_i$ is $R$-flat, it follows immediately that $\Phi$ is an isomorphism so that $(1_R, \Phi) : (R, X_1) \to (R, X_2)$ is an isomorphism in $\underline{M}_{X_0}$ . Consequently, $\underline{M}_{X_0}$ is a cofibered category in groupoids over $\underline{C}_\Lambda$ .

4.1(b): Suppose that $X_0$ is affine, say $X_0 = \mathrm{Spec}(A_0)$ . Then for any object $(R, X) \in \mathrm{ob}(\underline{M}_{X_0})$ , $X$ is also affine. Consequently, for any $R \in \mathrm{ob}(\underline{C}_\Lambda)$ , the category $\underline{M}_{X_0}(R)$ is canonically equivalent, via the sections, to the category in which an object is a flat algebra $A$ over $R$ together with an $R$-algebra map $j_A : A \to A_0$ inducing an isomorphism $A \otimes_R k \xrightarrow{\ \sim\ } A_0$ , and a map $\Phi : A_1 \to A_2$ is an $R$-algebra map with a commutative diagram

\begin{center}
\begin{tikzcd}[column sep=huge, row sep=huge]
A_1 \arrow[r, "\Phi"] \arrow[d, "j_{A_1}"'] \& A_2 \arrow[d, "j_{A_2}"] \\
A_0 \arrow[r, equal] \& A_0
\end{tikzcd}
\end{center}

4.1(c): Consider the cofibered category $\widehat{\underline{M}}_{X_0}$ in groupoids over $\widehat{\underline{C}}_\Lambda$ . By definition, the category $\widehat{\underline{M}}_{X_0}(R)$ for any $R \in \mathrm{ob}(\widehat{\underline{C}}_\Lambda)$ is canonically equivalent to the category of formal schemes $\mathfrak{X}$ , adic over the formal spectrum $\mathrm{Spf}(R)$ , such that $\mathfrak{X}_n = \mathfrak{X} \otimes_R R/\underline{m}^{n+1}$ is flat over $R/\underline{m}^{n+1}$ for every $n \geq 0$ (or what amounts to the same, if $X_0$ is noetherian, such that $\mathfrak{X}$ is flat over $\mathrm{Spf}(R)$) . Consequently, if $X_0$ is noetherian, $\widehat{\underline{M}}_{X_0}$ is canonically $\widehat{\underline{C}}_\Lambda$-equivalent to the category of deformations of $X_0$ over the objects in $\widehat{\underline{C}}_\Lambda$ , in the sense of formal schemes.

4.1(d): Let $(R, X) \in \mathrm{ob}(\underline{M}_{X_0})$ i.e., we have a commutative diagram
\SGApage{I}{97}
\[
\begin{tikzcd}
X_0 \arrow[r, "i_X"] \arrow[d] \& X \arrow[d, "\text{flat}"] \\
\mathrm{Spec}(k) \arrow[r] \& \mathrm{Spec}(R)
\end{tikzcd}
\]

such that  $X_0 \to X \times_{\mathrm{Spec}(R)} \mathrm{Spec}(k)$  is an isomorphism. Now let  $U_0 \subset X_0$  be a fixed open subscheme of  $X_0$ . Since  $i_X : X_0 \to X$  is a closed immersion which is a homeomorphism on the underlying topological spaces, we obtain a commutative diagram

\begin{center}
\begin{tikzcd}
U_0 \arrow[r, "i_X|U_0"] \arrow[d] \& X|U_0 \arrow[d, "\text{flat}"] \\
\mathrm{Spec}(k) \arrow[r] \& \mathrm{Spec}(R)
\end{tikzcd}
\end{center}









such that  $U_0 \to (X|U_0) \times_{\mathrm{Spec}(R)} \mathrm{Spec}(k) = X \otimes_R k|U_0$  is an isomorphism, so that $(R, X|U_0)$  is a deformation of  $U_0$  to  $R$ . Furthermore, for any arrow $(\varphi, \Phi) : (R, X) \to (R', X')$  in  $\underline{M}_{X_0}$ ,  $(\varphi, \Phi|U_0) : (R, X|U_0) \to (R', X'|U_0)$  is an arrow in  $\underline{M}_{U_0}$ . Consequently we have the restriction functor  $\underline{M}_{X_0} \to \underline{M}_{U_0}$  over $\underline{C}_\Lambda$ . One may also consider a localization as well as a formalization at a point. Namely let  $x_0$  be a fixed point in  $X_0$ . If  $(R, X) \in \mathrm{ob}(\underline{M}_{X_0})$ , then $\mathrm{Spec}(\underline{O}_{X,x_0})$  is a deformation of  $\mathrm{Spec}(\underline{O}_{X_0,x_0})$  to  $R$  and  $\mathrm{Spec}(\hat{\underline{O}}_{X,x_0})$  is a deformation of  $\mathrm{Spec}(\hat{\underline{O}}_{X_0,x_0})$  over  $R$ , where  $\hat{\ }$  stands for the completion with respect to the linear topology given by the powers of the maximal ideal. We thus obtain functors  $\underline{M}_{X_0} \to \underline{M}_{\mathrm{Spec}(\underline{O}_{X_0,x_0})} \to \underline{M}_{\mathrm{Spec}(\hat{\underline{O}}_{X_0,x_0})}$ .

Firstly we establish that  $\underline{M}_{X_0}$  is a homogeneous cofibered category over $\underline{C}_\Lambda$ .

\underline{Lemma} 4.2.  \underline{Consider a diagram}

\begin{center}
\begin{tikzcd}
 \& R' \arrow[d] \\
R \arrow[r, "\text{surjective}"'] \& \bar{R}
\end{tikzcd}
\end{center}

\underline{in}  $\underline{C}_\Lambda$ , \underline{and let}  $E, E'$  \underline{be flat modules over}  $R, R'$  \underline{respectively. Then for any} $R'$-\underline{homomorphism}  $E' \to \bar{E} = E \otimes_R \bar{R}$ , \underline{the fibre-product}  $E \times_{\bar{E}} E'$  \underline{is flat as} $R \times_{\bar{R}} R'$-\underline{module}


\SGApage{I}{98}

\underline{Proof}:  Let  $0 \to I \to R \to \bar{R} \to 0$  be exact so that  $0 \to I \to R \times_{\bar{R}} R' \to R' \to 0$  is exact. Since  $E$  is $R$-flat,  $0 \to I \otimes_R E \to E \to \bar{E} \to 0$  is exact and hence we obtain the exact sequence

\SGAvar{\begin{equation}\tag{*}
0 \to I \otimes_R E \to E \times_{\bar{E}} E' \to E' \to 0 \ .
\end{equation}
Set  $S = R \times_{\bar{R}} R$ . Since  $I$  is nilpotent,  $E \times_{\bar{E}} E'$  is flat as an $S$-module if and only if  $(E \times_{\bar{E}} E') \otimes_S R'$  is flat as an $R'$-module and  $\mathrm{Tor}_1^S (E \times_{\bar{E}} E', R') = 0$ . However, we note that  $I \otimes_R E = I \otimes_S (E \times_{\bar{E}} E')$  and  $E' = R' \otimes_S (E \times_{\bar{E}} E')$  and therefore the exact sequence (*) can be rewritten as

}{\begin{equation}\tag{*}
0 \to I \otimes_R E \to E \times_{\bar{E}} E' \to E' \to 0 \ .
\end{equation}
Set  $S = R \times_{\bar{R}} R'$ . Since  $I$  is nilpotent,  $E \times_{\bar{E}} E'$  is flat as an $S$-module if and only if  $(E \times_{\bar{E}} E') \otimes_S R'$  is flat as an $R'$-module and  $\mathrm{Tor}_1^S (E \times_{\bar{E}} E', R') = 0$ . However, we note that  $I \otimes_R E = I \otimes_S (E \times_{\bar{E}} E')$  and  $E' = R' \otimes_S (E \times_{\bar{E}} E')$  and therefore the exact sequence (*) can be rewritten as

}{F7I-VI-080}\begin{equation*}
0 \ \to \ I \otimes_S (E \times_{\bar{E}} E') \ \to \ E \times_{\bar{E}} E' \ \to \ R' \otimes_S (E \times_{\bar{E}} E') \ \to \ 0 \ .
\end{equation*}
This means that  $\mathrm{Tor}_1^S (R', E \times_{\bar{E}} E') = 0$  and  $R' \otimes_S (E \times_{\bar{E}} E') = E'$  is flat as $R'$-module, and therefore  $E \times_{\bar{E}} E'$  is $S$-flat.

\underline{Corollary} 4.3.  \underline{Given a diagram}

\begin{center}
\begin{tikzcd}[column sep=large]
(R,\ X) \arrow[dr, "{(\varphi,\ \Phi)}"'] \& \& (R',\ X') \arrow[dl, "{(\varphi',\ \Phi')}"] \\
 \& (\bar{R},\ \bar{X}) \&
\end{tikzcd}
\end{center}








\underline{in}  $\underline{M}_{X_0}$ , \underline{the amalgamated sum}  $X \underset{\bar{X}}{\amalg} X'$  \underline{(in the sense of schemes) is a flat scheme over}  $R \times_{\bar{R}} R'$ , \underline{provided}  $\varphi : R \to \bar{R}$  \underline{is surjective. Consequently,}  $\underline{M}_{X_0}$ \underline{is a homogeneous cofibered category in groupoids over}  $\underline{C}_\Lambda$ .

\SGAvar{\underline{Proof}:  For each affine open  $U_0 \subset X_0$ , we have by definition that $X \underset{\bar{X}}{\amalg} X'|U_0 \simeq \mathrm{Spec}(\Gamma(U_0,\ X) \times_{\Gamma(U_0,\bar{X})} \Gamma(U_0,\ X'))$ . Consequently, if  $\varphi : R \to \bar{R}$ is surjective,  $X \underset{\bar{X}}{\amalg} X'$  is a flat scheme over  $R \times_{\bar{R}} R'$  by 4.2, and the morphism $i_X \underset{i_{\bar{X}}}{\amalg} i_{X'} : X_0 \to X \underset{\bar{X}}{\phantom{\amalg}} X'$  induces an isomorphism  $X_0 \to (X \underset{\bar{X}}{\amalg} X') \otimes_S k$ , where $S = R \times_{\bar{R}} R'$ . Therefore  $(R \times_{\bar{R}} R' ,\ X \underset{\bar{X}}{\amalg} X') \in \mathrm{ob}(\underline{M}_{X_0})$  and we obviously have that  $(R \times_{\bar{R}} R',\ X \underset{\bar{X}}{\amalg} X') = (R,\ X) \times_{(\bar{R},\bar{X})} (R',\ X')$  in the category  $\underline{M}_{X_0}$ , which
\SGApage{I}{99}
means that $\underline{M}_{X_0}$ is homogeneous over $\underline{C}_\Lambda$ .

}{\underline{Proof}:  For each affine open  $U_0 \subset X_0$ , we have by definition that $X \underset{\bar{X}}{\amalg} X'|U_0 \simeq \mathrm{Spec}(\Gamma(U_0,\ X) \times_{\Gamma(U_0,\bar{X})} \Gamma(U_0,\ X'))$ . Consequently, if  $\varphi : R \to \bar{R}$ is surjective,  $X \underset{\bar{X}}{\amalg} X'$  is a flat scheme over  $R \times_{\bar{R}} R'$  by 4.2, and the morphism $i_X \underset{i_{\bar{X}}}{\amalg} i_{X'} : X_0 \to X \underset{\bar{X}}{\amalg} X'$  induces an isomorphism  $X_0 \to (X \underset{\bar{X}}{\amalg} X') \otimes_S k$ , where $S = R \times_{\bar{R}} R'$ . Therefore  $(R \times_{\bar{R}} R' ,\ X \underset{\bar{X}}{\amalg} X') \in \mathrm{ob}(\underline{M}_{X_0})$  and we obviously have that  $(R \times_{\bar{R}} R',\ X \underset{\bar{X}}{\amalg} X') = (R,\ X) \times_{(\bar{R},\bar{X})} (R',\ X')$  in the category  $\underline{M}_{X_0}$ , which
\SGApage{I}{99}
means that $\underline{M}_{X_0}$ is homogeneous over $\underline{C}_\Lambda$ .

}{F7I-VI-081}\underline{Lemma} 4.4.  \underline{For any scheme} $X_0$ \underline{over} $k$ , \underline{we have a canonical exact sequence}

\begin{equation*}
0 \;\to\; H^{1}(X_{0},\ \underline{D}^{0}_{X_{0}}) \;\to\; [\underline{M}_{X_{0}}(k[\varepsilon])] \;\to\; H^{0}(X_{0},\ \underline{D}^{1}_{X_{0}})
\end{equation*}
\underline{where} $\underline{D}^{i}_{X_{0}}$ \underline{are as defined in paragraph 3.}

\underline{Proof}:  (Special case)  Assume that $X_{0}$ is affine and let $\Gamma(X_{0},\ \underline{O}_{X_{0}}) = A_{0}$ . Then as it was noted in 4.1(b), the category $\underline{M}_{X_{0}}(k[\varepsilon])$ is canonically equivalent to the category $\underline{E}$ , where

\SGAvar{\begin{equation*}
\mathrm{ob}(\underline{E}) : \text{object in } \underline{E} \text{ is a flat algebra } A \text{ over } k[\varepsilon]
\end{equation*}
together with a  $k[\varepsilon]$ -algebra map  $j_{A} : A \to A_{0}$  inducing the isomorphism $k \otimes_{k[\varepsilon]} j_{A} : k \otimes_{k[\varepsilon]} A \overset{\sim}{\relbar\joinrel\relbar\joinrel\relbar} A_{0}$ . Alternatively, an object in $\underline{E}$ is a $k$-algebra $A$  together with an exact sequence  $0 \to A_{0} \overset{\varepsilon}{\relbar\joinrel\relbar\joinrel\relbar} A \overset{j_{A}}{\relbar\joinrel\relbar\joinrel\relbar} A_{0} \to 0$  in which $\varepsilon(A_{0})^{2} = 0$  as an ideal in  $A$ .

}{\begin{equation*}
\mathrm{ob}(\underline{E}) : \text{object in } \underline{E} \text{ is a flat algebra } A \text{ over } k[\varepsilon]
\end{equation*}
together with a  $k[\varepsilon]$ -algebra map  $j_{A} : A \to A_{0}$  inducing the isomorphism $k \otimes_{k[\varepsilon]} j_{A} : k \otimes_{k[\varepsilon]} A \overset{\sim}{\longrightarrow} A_{0}$ . Alternatively, an object in $\underline{E}$ is a $k$-algebra $A$  together with an exact sequence  $0 \to A_{0} \overset{\varepsilon}{\longrightarrow} A \overset{j_{A}}{\longrightarrow} A_{0} \to 0$  in which $\varepsilon(A_{0})^{2} = 0$  as an ideal in  $A$ .

}{F7I-VI-082}F$\ell$($\underline{E}$) :  An arrow  $A_{1} \to A_{2}$  in  $\underline{E}$  with respect to first description of the object in  $\underline{E}$  is a  $k[\varepsilon]$-algebra map  $\Phi : A_{1} \to A_{2}$  such that  $j_{A_{1}} = j_{A_{2}} \circ \Phi$ . Therefore, an arrow  $A_{1} \to A_{2}$  in  $\underline{E}$  with respect to the alternative description of objects in  $\underline{E}$  is a $k$-algebra map  $\Phi : A_{1} \to A_{2}$  yielding a commutative diagram

\begin{center}
\begin{tikzcd}[column sep=large]
0 \arrow[r] \& A_{0} \arrow[r, "\varepsilon"] \arrow[d, equal] \& A_{1} \arrow[r, "j_{A_{1}}"] \arrow[d, "\Phi"] \& A_{0} \arrow[r] \arrow[d, equal] \& 0 \\
0 \arrow[r] \& A_{0} \arrow[r] \& A_{2} \arrow[r, "j_{A_{2}}"] \& A_{0} \arrow[r] \& 0
\end{tikzcd}\quad .
\end{center}
















\SGAvar{It follows that  $[\underline{M}_{X_{0}}(k[\varepsilon])]$  is canonically isomorphic, via  $\Gamma(X_{0},\ -)$ , to the set of isomorphism classes of $k$-algebra extensions of  $A_{0}$  by  $A_{0}$ . In other words, we have a canonical bijection  $[\underline{M}_{X_{0}}(k[\varepsilon])] \overset{\sim}{\relbar\joinrel\relbar\joinrel\relbar} H^{0}(X_{0},\ \underline{D}^{1}_{X_{0}})$  via $\Gamma(X_{0},\ -)$, which is trivially seen to respect the vector space structures on both sides.


}{It follows that  $[\underline{M}_{X_{0}}(k[\varepsilon])]$  is canonically isomorphic, via  $\Gamma(X_{0},\ -)$ , to the set of isomorphism classes of $k$-algebra extensions of  $A_{0}$  by  $A_{0}$ . In other words, we have a canonical bijection  $[\underline{M}_{X_{0}}(k[\varepsilon])] \overset{\sim}{\longrightarrow} H^{0}(X_{0},\ \underline{D}^{1}_{X_{0}})$  via $\Gamma(X_{0},\ -)$, which is trivially seen to respect the vector space structures on both sides.


}{F7I-VI-083}\SGApage{I}{100}

\SGAvar{(General case)  For each affine open  $U_{0} \subset X_{0}$ , we get a canonical map $[\underline{M}_{X_{0}}(k[\varepsilon])] \to H^{0}(U_{0},\ \underline{D}^{1}_{X_{0}})$  defined as the composite map $[\underline{M}_{X_{0}}(k[\varepsilon])] \to [\underline{M}_{U_{0}}(k[\varepsilon])] \overset{\sim}{\relbar\joinrel\relbar\joinrel\relbar} H^{0}(U_{0},\ \underline{D}^{1}_{X_{0}})$  and hence we obtain a canonical map $[\underline{M}_{X_{0}}(k[\varepsilon])] \to H^{0}(X_{0},\ \underline{D}^{1}_{X_{0}})$ . Its kernel is the set of isomorphism classes of locally-trivial deformations of  $X_{0}$  to  $k[\varepsilon]$ , i.e., the set of isomorphism classes of deformations of  $X_{0}$  to  $k[\varepsilon]$ , locally (on $X$ ) isomorphic to $X_{0} \otimes_{k} k[\varepsilon]$ . Since  $\underline{D}^{0}_{X_{0}} = \underline{\mathrm{Der}}_{k}\,(\underline{O}_{X_{0}},\ \underline{O}_{X_{0}}) = \underline{\mathrm{Aut}}(X_{0} \otimes_{k} k[\varepsilon]|X_{0})$  $=$ the sheaf of germs of automorphisms of  $X_{0} \otimes_{k} k[\varepsilon]$  inducing the identity on the subscheme $X_{0}$ , it is canonically isomorphic to  $H^{1}(X_{0},\ \underline{D}^{0}_{X_{0}})$ , so that we obtain the desired exact sequence.

}{(General case)  For each affine open  $U_{0} \subset X_{0}$ , we get a canonical map $[\underline{M}_{X_{0}}(k[\varepsilon])] \to H^{0}(U_{0},\ \underline{D}^{1}_{X_{0}})$  defined as the composite map $[\underline{M}_{X_{0}}(k[\varepsilon])] \to [\underline{M}_{U_{0}}(k[\varepsilon])] \overset{\sim}{\longrightarrow} H^{0}(U_{0},\ \underline{D}^{1}_{X_{0}})$  and hence we obtain a canonical map $[\underline{M}_{X_{0}}(k[\varepsilon])] \to H^{0}(X_{0},\ \underline{D}^{1}_{X_{0}})$ . Its kernel is the set of isomorphism classes of locally-trivial deformations of  $X_{0}$  to  $k[\varepsilon]$ , i.e., the set of isomorphism classes of deformations of  $X_{0}$  to  $k[\varepsilon]$ , locally (on $X$ ) isomorphic to $X_{0} \otimes_{k} k[\varepsilon]$ . Since  $\underline{D}^{0}_{X_{0}} = \underline{\mathrm{Der}}_{k}\,(\underline{O}_{X_{0}},\ \underline{O}_{X_{0}}) = \underline{\mathrm{Aut}}(X_{0} \otimes_{k} k[\varepsilon]|X_{0})$  $=$ the sheaf of germs of automorphisms of  $X_{0} \otimes_{k} k[\varepsilon]$  inducing the identity on the subscheme $X_{0}$ , it is canonically isomorphic to  $H^{1}(X_{0},\ \underline{D}^{0}_{X_{0}})$ , so that we obtain the desired exact sequence.

}{F7I-VI-084}\underline{Theorem} 4.5.  \underline{Let} $X_{0}$ \underline{be a scheme of finite type over} $k$ \underline{for which either} (i) $X_{0}$ \underline{is proper over} $k$ \underline{or} (ii) $X_{0}$ \underline{is affine but} $X_{0}^{\mathrm{sing}} = \{x \in X_{0}|X_{0}$ \underline{is non-smooth over} $k$ \underline{at the point} $x \}$ \underline{is a finite set. Then}

(a)  $\underline{M}_{X_{0}}$  \underline{admits a formal versal object} (\underline{called a formal versal deformation of}  $X_{0}$ ) \underline{over}  $\underline{C}_{\Lambda}$ . \underline{If}  $(R,\ \mathfrak{X})$  \underline{is a formal versal deformation of}  $X_{0}$ , \underline{then for any} $(S,\ \mathfrak{Y}) \in \mathrm{ob}(\widehat{\underline{M}}_{X_{0}})$ , \underline{there is a map}  $(\varphi,\Phi) : (R,\mathfrak{X}) \to (S,\mathfrak{Y})$  \underline{in}  $\widehat{\underline{M}}_{X_{0}}$  \underline{such that the diagram below is cartesian}

\begin{center}
\begin{tikzcd}[column sep=huge, row sep=large]
\mathfrak{Y} \arrow[r, "\Phi"] \arrow[d] \& \mathfrak{X} \arrow[d] \\
\mathrm{Spf}(S) \arrow[r, "\mathrm{Spf}(\varphi)"'] \& \mathrm{Spf}(R)
\end{tikzcd}\quad .
\end{center}









\underline{Furthermore a formal versal deformation}  $(R,\mathfrak{X})$  of  $X_{0}$  \underline{prorepresents the functor} $[\underline{M}_{X_{0}}]$  \underline{if and only if, for every surjective map}  $R' \to R$  \underline{in}  $\underline{C}_{\Lambda}$  \underline{and a deformation} $(R',\ X')$  \underline{of}  $X_{0}$  \underline{to}  $R'$ , \underline{the homomorphism}

\begin{equation*}
\mathrm{Aut}_{R'}(X'|X_{0}) \to \mathrm{Aut}_{R}(X' \underset{R'}{\otimes} R|X_{0})
\end{equation*}
\underline{is surjective.}

(b)  \underline{If}  $X_{0}$  \underline{is proper over}  $k$ , \underline{then}  $\underline{M}_{X_{0}}$  \underline{is smoothly pseudo-representable}
\SGApage{I}{101}
\underline{over} $\underline{C}_{\Lambda}$ .

\SGAvar{\underline{Proof}:  (a)  Since  $\underline{M}_{X_{0}}$  is a homogeneous cofibered category in groupoids over $\underline{C}_{\Lambda}$ by 4.3, it suffices to see that  $[\underline{M}(k[\epsilon])]$  is a finite-dimensional vector space over  $k$  by 1.11 and 2.9(b).

}{\underline{Proof}:  (a)  Since  $\underline{M}_{X_{0}}$  is a homogeneous cofibered category in groupoids over $\underline{C}_{\Lambda}$ by 4.3, it suffices to see that  $[\underline{M}_{X_{0}}(k[\epsilon])]$  is a finite-dimensional vector space over  $k$  by 1.11 and 2.9(b).

}{F7I-VI-085}Now consider the exact sequence

\SGAvar{\begin{equation}\tag{*}
0 \to H^{1}(X_{0}, \underline{D}^{0}_{X_{0}}) \to [\underline{M}_{X_{0}}(k[\epsilon])] \to H^{0}(X_{0}, \underline{D}^{1}_{X_{0}}) \ .
\end{equation}
If  $X_{0}$  is proper over  $k$  , then  $H^{1}(X_{0}, \underline{D}^{0}_{X_{0}})$  and  $H^{0}(X_{0}, \underline{D}^{1}_{X_{0}})$  are both finite dimensional since  $\underline{D}^{i}_{X_{0}}$  are coherent for all  $i \geq 0$  by 2.13(a). Now assume that $X_{0}$  is affine such that  $X_{0}^{\mathrm{sing}}$  is a finite set. Since  $\underline{D}^{1}_{X_{0}}$  is a coherent sheaf on  $X_{0}$  such that  $\mathrm{Supp}\ \underline{D}^{1}_{X_{0}} \subset X_{0}^{\mathrm{sing}}$  by 2.13(a), we find that  $H^{0}(X_{0}, \underline{D}^{1}_{X_{0}})$  is a finite-dimensional vector space over  $k$  , whereas  $H^{1}(X_{0}, \underline{D}^{0}_{X_{0}}) = 0$  since  $X_{0}$ is affine. Therefore, in either case, it follows from the exact sequence (*) that $[\underline{M}_{X_{0}}(k[\epsilon])]$  is a finite dimensional vector space over  $k$  . The second statement of (a) follows from 2.7.

}{\begin{equation}\tag{*}
0 \to H^{1}(X_{0}, \underline{D}^{0}_{X_{0}}) \to [\underline{M}_{X_{0}}(k[\epsilon])] \to H^{0}(X_{0}, \underline{D}^{1}_{X_{0}}) \ .
\end{equation}
If  $X_{0}$  is proper over  $k$  , then  $H^{1}(X_{0}, \underline{D}^{0}_{X_{0}})$  and  $H^{0}(X_{0}, \underline{D}^{1}_{X_{0}})$  are both finite dimensional since  $\underline{D}^{i}_{X_{0}}$  are coherent for all  $i \geq 0$  by 3.15(a). Now assume that $X_{0}$  is affine such that  $X_{0}^{\mathrm{sing}}$  is a finite set. Since  $\underline{D}^{1}_{X_{0}}$  is a coherent sheaf on  $X_{0}$  such that  $\mathrm{Supp}\ \underline{D}^{1}_{X_{0}} \subset X_{0}^{\mathrm{sing}}$  by 3.15(a), we find that  $H^{0}(X_{0}, \underline{D}^{1}_{X_{0}})$  is a finite-dimensional vector space over  $k$  , whereas  $H^{1}(X_{0}, \underline{D}^{0}_{X_{0}}) = 0$  since  $X_{0}$ is affine. Therefore, in either case, it follows from the exact sequence (*) that $[\underline{M}_{X_{0}}(k[\epsilon])]$  is a finite dimensional vector space over  $k$  . The second statement of (a) follows from 2.7.

}{F7I-VI-086}          (b)  If  $X_{0}$  is proper,  $\mathrm{Aut}\ (X_{0} \otimes_{k} k[\epsilon]|X_{0}) = H^{0}(X_{0}, \underline{D}^{0}_{X_{0}})$  is a finite dimensional vector space and hence our statement follows from 2.13.

\underline{Remark 4.6(a)}:  Let  $X_{0}$  be affine. Then a formal versal deformation of  $X_{0}$ exists if and only if  $\underline{D}^{1}_{X_{0}}$  has its support on a finite number of points. However this does not imply that non-smooth points of  $X_{0}$  are all isolated. There are affine schemes  $X_{0}$  over $k$  having a non-isolated non-smooth point but  $\mathrm{supp}\ \underline{D}^{1}_{X_{0}}$ is a finite set.

\underline{Remark 4.6(b)}:  Henceforth a deformation of  $X_{0}$  is meant by an object in  $\underline{M}_{X_{0}}$  . In case when there is a possible confusion, an object in  $\underline{M}_{X_{0}}$  will be referred as an infinitesimal deformation of  $X_{0}$  .

\underline{Remark 4.6(c)}:  Let  $X_{0}$  be a scheme over  $k$  . We have fixed a complete local noetherian ring  $\Lambda$  with the residue field  $k$  , and considered the deformations
\SGApage{I}{102}
of  $X_{0}$  over the category  $\underline{C}_{\Lambda}$  . (In practice one usually takes  $\Lambda = k$  or  $W(k)$ if  $k$  is a perfect field of positive characteristic). Now  $\underline{C}_{k} \subset \underline{C}_{\Lambda}$  is a fully faithful imbedding, and that  $\underline{M}_{X_{0}}|\underline{C}_{k}$  is simply the category of deformations of $X_{0}$  over  $S$  where  $S \in \mathrm{ob}(\underline{C}_{k})$  . Therefore if  $(R, \mathfrak{X})$  is a formal versal deformation of  $X_{0}$  over  $\underline{C}_{\Lambda}$  , then  $(k \otimes_{\Lambda} R, k \otimes_{\Lambda} \mathfrak{X})$  is a formal versal deformation of $X_{0}$  over  $\underline{C}_{k}$  . It follows, for example, that  $\dim_{\underline{C}_{\Lambda}} \underline{M}_{X_{0}} = \dim k \otimes_{\Lambda} R$  does not depend on the choice of  $\Lambda$  but depends only on  $X_{0}$  .

          One can generalize the above theorem in many ways. Let us for instance consider a pair  $Y_{0} \to X_{0}$  where  $X_{0}$  is a scheme of finite type over  $k$  and  $Y_{0}$ a fixed closed subscheme of  $X_{0}$  . For each  $R \in \mathrm{ob}(\underline{C}_{\Lambda})$  , a deformation of $(X_{0}, Y_{0})$  to  $R$  is meant to be a commutative diagram

\SGAvar{\begin{tikzcd}[row sep=large, column sep=huge]
Y_0 \arrow[rrr] \arrow[dr, hook] \arrow[ddr] \& \& \& Y \arrow[dl, hook] \arrow[ddl, "\text{flat}"] \\
 \& X_0 \arrow[r] \arrow[d] \& X \arrow[d, no head, "\text{flat}"] \& \\
 \& \mathrm{Spec}(k) \arrow[r] \& \mathrm{Spec}(R) \&
\end{tikzcd}


















}{\begin{tikzcd}[row sep=large, column sep=huge]
Y_0 \arrow[rrr] \arrow[dr, hook] \arrow[ddr] \& \& \& Y \arrow[dl, hook] \arrow[ddl, "\text{flat}"] \\
 \& X_0 \arrow[r] \arrow[d] \& X \arrow[d, "\text{flat}"] \& \\
 \& \mathrm{Spec}(k) \arrow[r] \& \mathrm{Spec}(R) \&
\end{tikzcd}


















}{F7I-VI-087}such that  $(X_{0}, Y_{0}) \to (X \times_{\mathrm{Spec}(R)} \mathrm{Spec}(k), Y \times_{\mathrm{Spec}(R)} \mathrm{Spec}(k))$  is an isomorphism. Defining a morphism  $(R'; X', Y') \to (R; X, Y)$  of deformations of  $(X_{0}, Y_{0})$  to be the obvious ones, we obtain the cofibered category  $\underline{M}_{X_{0},Y_{0}}$  in groupoids over $\underline{C}_{\Lambda}$  . (This is a process often used in rigidification of automorphisms in moduli problems of deformations.) One sees immediately, by the same argument as in 4.3 based on 4.2, that  $\underline{M}_{X_{0},Y_{0}}$  is a homogeneous cofibered category in groupoids over $\underline{C}_{\Lambda}$  .

\underline{Lemma} 4.7.  $0 \to H^{0}(Y_{0}, \underline{\mathrm{Hom}}_{\underline{O}_{Y_{0}}} (\underline{I}/\underline{I}^{2}, \underline{O}_{Y_{0}})) \to [\underline{M}_{X_{0},Y_{0}}(k[\epsilon])] \to [\underline{M}_{X_{0}}(k[\epsilon])]$ \underline{is exact, where}  $\underline{I}$  \underline{is the ideal sheaf of}  $\underline{O}_{X_{0}}$  \underline{defining}  $Y_{0}$  i.e. 


\SGApage{I}{103}

\SGAvar{\begin{equation*}
 0 \to \underline{I} \to \underline{O}_{X_0} \to \underline{O}_{Y_0} \to 0 \qquad \underline{\text{is exact,}}
\end{equation*}
\underline{Proof}: Let $\underline{B}$ be the full subcategory of $\underline{M}_{X_0,Y_0}(k[\epsilon])$ whose objects are of the form $(k[\epsilon]; X \times_{\mathrm{Spec}(k)} \mathrm{Spec}(k[\epsilon]), Y)$ . We then have $[B] = \mathrm{Ker}\,([\underline{M}_{X_0,Y_0}(k[\epsilon])] \to [\underline{M}_{X_0}(k[\epsilon])]\,)$ . The category $\underline{B}$ is the groupoid of deformations of $Y_0$ to $k[\epsilon]$ "inside $X_0 \times_{\mathrm{Spec}(k)} \mathrm{Spec}(k[\epsilon])$ ", i.e., an object in $\underline{B}$ is a commutative diagram

}{\begin{equation*}
 0 \to \underline{I} \to \underline{O}_{X_0} \to \underline{O}_{Y_0} \to 0 \qquad \underline{\text{is exact,}}
\end{equation*}
\underline{Proof}: Let $\underline{B}$ be the full subcategory of $\underline{M}_{X_0,Y_0}(k[\epsilon])$ whose objects are of the form $(k[\epsilon]; X_0 \times_{\mathrm{Spec}(k)} \mathrm{Spec}(k[\epsilon]), Y)$ . We then have $[B] = \mathrm{Ker}\,([\underline{M}_{X_0,Y_0}(k[\epsilon])] \to [\underline{M}_{X_0}(k[\epsilon])]\,)$ . The category $\underline{B}$ is the groupoid of deformations of $Y_0$ to $k[\epsilon]$ "inside $X_0 \times_{\mathrm{Spec}(k)} \mathrm{Spec}(k[\epsilon])$ ", i.e., an object in $\underline{B}$ is a commutative diagram

}{F7I-VI-088}\begin{center}
\begin{tikzcd}[column sep=large, row sep=large]
Y \arrow[rr, "i_Y"] \arrow[dr, "\text{flat}"'] \& \& X_0 \times_{\mathrm{Spec}(k)} \mathrm{Spec}(k[\epsilon]) \arrow[dl] \\
 \& \mathrm{Spec}(k[\epsilon]) \&
\end{tikzcd}
\end{center}










such that $\mathrm{Spec}(k) \times_{\mathrm{Spec}(k[\epsilon])} i_Y = i_{Y_0}$ , and an arrow $(Y_1, i_{Y_1}) \to (Y_2, i_{Y_2})$ is a $k[\epsilon]$-morphism $\Phi : Y_1 \to Y_2$ (necessarily an isomorphism) such that $\mathrm{Spec}(k) \times_{\mathrm{Spec}(k[\epsilon])} \Phi = \mathrm{id}_{Y_0}$ and $i_{Y_2} \circ \Phi = i_{Y_1}$ . Alternatively, an object in $\underline{B}$ is a pair $(Y, j_Y)$ , where $Y$ is a deformation of $Y_0$ to $k[\epsilon]$ and $j_Y : Y \to X_0$ is a morphism over $k$ such that the composite morphism $Y_0 \to Y \xrightarrow{\;j_Y\;} X_0$ is $i_{Y_0}$ , and an arrow $(Y_1, j_{Y_1}) \to (Y_2, j_{Y_2})$ is a morphism $\Phi : Y_1 \to Y_2$ such that $k \otimes_{k[\epsilon]} \Phi = \mathrm{id}_{Y_0}$ and $j_{Y_2} \circ \Phi = j_{Y_1}$ . Therefore the association

\SGAvar{$(Y, j_Y) \;\longmapsto\;$
\begin{tikzcd}[column sep=large]
0 \arrow[r] \& \underline{I} \arrow[r] \arrow[d] \& \underline{O}_{X_0} \arrow[r] \arrow[d, "\mathfrak{J}_Y"] \& \underline{O}_{Y_0} \arrow[r] \arrow[d, equal] \& 0 \\
0 \arrow[r] \& \underline{O}_{Y_0} \arrow[r, "\epsilon"] \& \underline{O}_Y \arrow[r] \& \underline{O}_{Y_0} \arrow[r] \& 0
\end{tikzcd}















}{$(Y, j_Y) \;\longmapsto\;$
\begin{tikzcd}[column sep=large]
0 \arrow[r] \& \underline{I} \arrow[r] \arrow[d] \& \underline{O}_{X_0} \arrow[r] \arrow[d, "j_Y"] \& \underline{O}_{Y_0} \arrow[r] \arrow[d, equal] \& 0 \\
0 \arrow[r] \& \underline{O}_{Y_0} \arrow[r, "\epsilon"] \& \underline{O}_Y \arrow[r] \& \underline{O}_{Y_0} \arrow[r] \& 0
\end{tikzcd}















}{F7I-VI-089}gives rise to an isomorphism

\begin{equation*}
[B] \;\xrightarrow{\;j|\underline{I}\;}\; \mathrm{Hom}_{\underline{O}_{X_0}}(\underline{I},\, \underline{O}_{Y_0}) = \mathrm{Hom}_{\underline{O}_{Y_0}}(\underline{I}/\underline{I}^2,\, \underline{O}_{Y_0}) = H^0(Y_0,\, \underline{\mathrm{Hom}}_{\underline{O}_{Y_0}}(\underline{I}/\underline{I}^2,\, \underline{O}_{Y_0}))\;.
\end{equation*}

\SGApage{I}{104}

\underline{Remark} 4.8.  For a more general formula implying the above lemma, see SGA3 III 4.4. One may also note the following fact. One can define for any $S$-scheme $X$ and a quasi-coherent $X$-module $\underline{E}$ a quasi-coherent $X$-module $\underline{D}^i_{X|S}(\underline{E})$ for all $i \geq 0$ , generalizing the case when $S$ is affine. As usual we simply write $\underline{D}^i_{X|S}(\underline{O}_X) = \underline{D}^i_{X|S}$ . We then have

\SGAvar{\begin{equation*}
H^0(Y_0,\, \underline{\mathrm{Hom}}_{Y_0}(\underline{I}/\underline{I}^2,\, \underline{O}_{Y_0})) = H^0(Y_0,\, \underline{D}^1_{Y_0|X})
\end{equation*}
where $0 \to \underline{I} \to \underline{O}_{X_0} \to \underline{O}_{Y_0} \to 0$ is exact.

}{\begin{equation*}
H^0(Y_0,\, \underline{\mathrm{Hom}}_{Y_0}(\underline{I}/\underline{I}^2,\, \underline{O}_{Y_0})) = H^0(Y_0,\, \underline{D}^1_{Y_0|X_0})
\end{equation*}
where $0 \to \underline{I} \to \underline{O}_{X_0} \to \underline{O}_{Y_0} \to 0$ is exact.

}{F7I-VI-090}\underline{Theorem} 4.9.  \underline{Let} $X_0$ \underline{be} $k$-\underline{scheme of finite type and} $Y_0$ \underline{a closed subscheme of} $X_0$ . \underline{Assume that either} (i) $X_0$ \underline{is proper over} $k$ \underline{or} (ii) $X_0$ \underline{is affine and smooth over} $k$ \underline{outside a finite number of points and} $Y_0$ \underline{is proper over} $k$ . \underline{We then have}

(a) $\underline{M}_{X_0,\,Y_0}$ \underline{admits a formal versal deformation over} $\underline{C}_\wedge$ , \underline{and it pro-represents the functor} $[\underline{M}_{X_0,\,Y_0}]$ \underline{if and only if, for every surjective map} $R' \to R$ \underline{in} $\underline{C}_\wedge$ \underline{and a deformation} $(R'; X', Y')$ ,

\SGAvar{\begin{equation*}
\mathrm{Aut}_{R'}((X', Y')|X_0) \to \mathrm{Aut}_R(X' \otimes_{R'} R,\, Y',\, \otimes_{R'} R)|X_0)
\end{equation*}
\underline{is surjective.}

}{\begin{equation*}
\mathrm{Aut}_{R'}((X', Y')|X_0) \to \mathrm{Aut}_R((X' \otimes_{R'} R,\, Y' \otimes_{R'} R)|X_0)
\end{equation*}
\underline{is surjective.}

}{F7I-VI-091}(b) \underline{If furthermore} $H^0(X_0,\, \underline{D}^0_{Y_0|k} \times_{\underline{D}^0_{X_0|k}(\underline{O}_{Y_0})} \underline{D}^0_{X_0|k})$ \underline{is finite dimensional,} \underline{then} $\underline{M}_{X_0,\,Y_0}$ \underline{is smoothly pseudo-representable.}

\underline{Proof}: It is obvious from 4.4 and 4.6 together with 1.11, 2.7, and 2.13.

\SGAvar{We now consider the passage to localization and completion in the deformation theory     of schemes. Firstly on a notation: Let $X_0$ be a $k$-scheme of finite type. Given a fixed point $x_0 \in X_0$ , we have natural functors

}{We now consider the passage to localization and completion in the deformation theory     of schemes. First, some notation: Let $X_0$ be a $k$-scheme of finite type. Given a fixed point $x_0 \in X_0$ , we have natural functors

}{F7I-VI-092}\begin{equation*}
\underline{M}_{X_0} \to \underline{M}_{\mathrm{Spec}(\underline{O}_{X_0,\, x_0})} \to \underline{M}_{\mathrm{Spec}(\widehat{\underline{O}}_{X_0,\, x_0})} 
\end{equation*}

\SGApage{I}{105}

 and in turn we obtain the functors

\SGAvar{\begin{equation*}
\widehat{\underline{M}}_{X_0} \to \widehat{\underline{M}}_{\mathrm{Spec}(\underline{O}_{X_0,\ x_0})} \to \widehat{\underline{M}}_{\mathrm{Spec}(\widehat{\underline{O}}_{X_0,\ x_0})}\ .
\end{equation*}
As it was remarked in 4.1(c), $\underline{M}_{X_0}$  is canonically equivalent to the category in which an object is a pair  $(R,\ \mathfrak{X})$  where  $R \in \mathrm{ob}(\underline{C}_{\Lambda})$  and  $\mathfrak{X}$  is a flat $R$-adic formal scheme together with a fixed isomorphism  $\mathfrak{X}_0 \longleftrightarrow X_0$ . Henceforth we shall write, when there is no possible confusion, the images of  $(R,\ \mathfrak{X})$  under $\underline{M}_{X_0} \to \underline{M}_{\mathrm{Spec}(\underline{O}_{X_0,\ x_0})}$  and  $\underline{M}_{X_0} \to \underline{M}_{\mathrm{Spec}(\underline{O}_{X_0,\ x_0})}$  by  $(R,\ \mathfrak{X}_{(x_0)})$  and  $(R,\ \mathfrak{X}_{(x_0)})$ respectively. Therefore if  $\mathfrak{X} = \varinjlim\limits_{\nu} \mathfrak{X}_{\nu}$  where  $\mathfrak{X}_{\nu} = \mathfrak{X} \underset{\mathrm{Spf}(R)}{\times} \mathrm{Spec}(R/\underline{m}^{\nu+1})$    then

\begin{equation*}
\mathfrak{X}_{(x_0)} = \varinjlim\limits_{\nu} \mathrm{Spec}(\underline{O}_{\mathfrak{X}_{\nu},\ x_0})
\end{equation*}
\begin{equation*}
\mathfrak{X}_{(\hat{X}_0)} = \varinjlim\limits_{\nu} \mathrm{Spec}(\widehat{\underline{O}}_{\mathfrak{X}_{\nu},\ x_0})\ .
\end{equation*}
\underline{Theorem} 4.10. \underline{Let}  $X_0$  \underline{be an affine scheme of finite type over}  $k$  .

}{\begin{equation*}
\widehat{\underline{M}}_{X_0} \to \widehat{\underline{M}}_{\mathrm{Spec}(\underline{O}_{X_0,\ x_0})} \to \widehat{\underline{M}}_{\mathrm{Spec}(\widehat{\underline{O}}_{X_0,\ x_0})}\ .
\end{equation*}
As was remarked in 4.1(c), $\widehat{\underline{M}}_{X_0}$  is canonically equivalent to the category in which an object is a pair  $(R,\ \mathfrak{X})$  where  $R \in \mathrm{ob}(\widehat{\underline{C}}_{\Lambda})$  and  $\mathfrak{X}$  is a flat $R$-adic formal scheme together with a fixed isomorphism  $\mathfrak{X}_0 \longleftrightarrow X_0$ . Henceforth we shall denote, when there is no possible confusion, the images of  $(R,\ \mathfrak{X})$  under $\widehat{\underline{M}}_{X_0} \to \widehat{\underline{M}}_{\mathrm{Spec}(\underline{O}_{X_0,\ x_0})}$  and  $\widehat{\underline{M}}_{X_0} \to \widehat{\underline{M}}_{\mathrm{Spec}(\widehat{\underline{O}}_{X_0,\ x_0})}$  by  $(R,\ \mathfrak{X}_{(x_0)})$  and  $(R,\ \mathfrak{X}_{(\hat{x}_0)})$ respectively. Therefore if  $\mathfrak{X} = \varinjlim\limits_{\nu} \mathfrak{X}_{\nu}$  where  $\mathfrak{X}_{\nu} = \mathfrak{X} \underset{\mathrm{Spf}(R)}{\times} \mathrm{Spec}(R/\underline{m}^{\nu+1})$    then

\begin{equation*}
\mathfrak{X}_{(x_0)} = \varinjlim\limits_{\nu} \mathrm{Spec}(\underline{O}_{\mathfrak{X}_{\nu},\ x_0})
\end{equation*}
\begin{equation*}
\mathfrak{X}_{(\hat{x}_0)} = \varinjlim\limits_{\nu} \mathrm{Spec}(\widehat{\underline{O}}_{\mathfrak{X}_{\nu},\ x_0})\ .
\end{equation*}
\underline{Theorem} 4.10. \underline{Let}  $X_0$  \underline{be an affine scheme of finite type over}  $k$  .

}{F7I-VI-093}(a)  \underline{If}  $X_0$  \underline{is locally a complete-intersection, then}  $\underline{M}_{X_0}$  \underline{is smooth over}  $\underline{C}_{\Lambda}$  .

(b)  \underline{Assume that, for a fixed closed point}  $x \in X_0$  ,  $X_0 - \{x\}$  \underline{is smooth over}  $k$  . \underline{Then}

\begin{equation*}
\underline{M}_{X_0} \to \underline{M}_{\mathrm{Spec}(\underline{O}_{X_0,\ x})} \to \underline{M}_{\mathrm{Spec}(\widehat{\underline{O}}_{X_0,\ x})}
\end{equation*}
\underline{are both minimally smooth functors. In particular, if}  $(R,\ \mathfrak{X})$  \underline{is a formal versal deformation of}  $X_0$  , \underline{then}  $(R,\ \mathfrak{X}_{(x)})$  \underline{and}  $(R,\ \mathfrak{X}_{(\hat{x})})$  \underline{are formal versal deformations of}  $\mathrm{Spec}(\underline{O}_{X_0,\ x})$  \underline{and}  $\mathrm{Spec}(\widehat{\underline{O}}_{X_0,\ x})$  \underline{respectively}.

\SGAvar{\underline{Proof}:  In this proof we go back to the notations of the cohomology of commutative algebras defined in the paragraph 3. We set  $X_0 = \mathrm{Spec}(A_0)$


}{\underline{Proof}:  In this proof we go back to the notations of the cohomology of commutative algebras defined in paragraph 3. We set  $X_0 = \mathrm{Spec}(A_0)$


}{F7I-VI-094}\SGApage{I}{106}

(a)  It suffices to show that  $P : [\underline{M}_{X_0}] \to \underline{C}_{\Lambda}$  is smooth.

\SGAvar{Let  $(R,\ \mathrm{Spec}(A))$  be a deformation of  $\mathrm{Spec}(A_0)$  to  $R$  , where  $R \in \mathrm{ob}(\underline{C}_{\Lambda})$  . Given a small extension  $R' \to R$  in  $\underline{C}_{\Lambda}$  (i.e.,  $0 \to k \to R' \to R \to 0$  is exact), we must show that  $(R,\ \mathrm{Spec}(A))$  can be lifted to  $R'$  . The exact sequence of  $D^i$  with coefficient  $A_0$  applied to the ring homomorphisms  $R' \to R \to A$  gives us the exact sequence

\begin{equation*}
\ldots \to D^1(A|R,\ A_0) \to D^1(A|R',\ A_0) \to D^1(R|R',\ k) \underset{k}{\otimes} A_0 \to D^2(A|R,\ A_0) \to \ldots
\end{equation*}
Since  $R' \to R$  is a small extension, $D^1(R'|R,\ k)$  is 1-dimensional vector space over  $k$  generated by the element (denoted by  $[R']$) represented by the extension $0 \to k \to R' \to R \to 0$  . On the other hand,  $A$  is $R$-flat and hence $D^i(A|R,\ A_0) = D^i(A_0|k,A_0)$  for all  $i > 0$  and in particular  $D^2(A|R,\ A_0) = 0$ by 3.12(11) since  $A_0$  is locally a complete-intersection. Consequently we get the exact sequence

}{Let  $(R,\ \mathrm{Spec}(A))$  be a deformation of  $\mathrm{Spec}(A_0)$  to  $R$  , where  $R \in \mathrm{ob}(\underline{C}_{\Lambda})$  . Given a small extension  $R' \to R$  in  $\underline{C}_{\Lambda}$  (i.e.,  $0 \to k \to R' \to R \to 0$  is exact), we must show that  $(R,\ \mathrm{Spec}(A))$  can be lifted to  $R'$  . The exact sequence for  $D^i$  with coefficients in  $A_0$  applied to the ring homomorphisms  $R' \to R \to A$  gives us the exact sequence

\begin{equation*}
\ldots \to D^1(A|R,\ A_0) \to D^1(A|R',\ A_0) \to D^1(R|R',\ k) \underset{k}{\otimes} A_0 \to D^2(A|R,\ A_0) \to \ldots
\end{equation*}
Since  $R' \to R$  is a small extension, $D^1(R|R',\ k)$  is a 1-dimensional vector space over  $k$  generated by the element (denoted by  $[R']$) represented by the extension $0 \to k \to R' \to R \to 0$  . On the other hand,  $A$  is $R$-flat and hence $D^i(A|R,\ A_0) = D^i(A_0|k,A_0)$  for all  $i > 0$  and in particular  $D^2(A|R,\ A_0) = 0$ by 3.12(8) since  $A_0$  is locally a complete-intersection. Consequently we get the exact sequence

}{F7I-VI-095}\begin{equation*}
\to D^1(A_0|k,\ A_0) \to D^1(A|R',\ A_0) \to D^1(R|R',\ k) \underset{k}{\otimes} A_0 \to 0
\end{equation*}
This means that there exists an $R'$-algebra extension of  $A$  by  $A_0$  , denoted by $A'$  , such that  $[A'] \to [R'] \otimes 1$  under the map  $D^1(A|R',\ A_0) \to D^1(R|R',\ k) \underset{k}{\otimes} A_0$  , which means that the deformation  $(R,\ \mathrm{Spec}(A))$  can be lifted to a deformation $(R',\ \mathrm{Spec}(A'))$  .

(b)  By 3.14 we have the isomorphisms

\SGAvar{\begin{equation*}
D^1(A_0|k,\ A_0) \simeq D^1((A_0)_x|k,\ (A_0)_x) \simeq D^1(\widehat{(A_0)}_x|k,\ \widehat{(A_0)}_x)
\end{equation*}
i.e.,   $[\underline{M}_{X_0}(k[\epsilon])] \simeq [\underline{M}_{\mathrm{Spec}(\underline{O}_{X_0,\ x})}(k[\epsilon])] \simeq [\underline{M}_{\mathrm{Spec}(\widehat{\underline{O}}_{X_0,\ x})}(k[\epsilon])]$ .

}{\begin{equation*}
D^1(A_0|k,\ A_0) \simeq D^1((A_0)_x|k,\ (A_0)_x) \simeq D^1(\widehat{(A_0)_x}|k,\ \widehat{(A_0)_x})
\end{equation*}
i.e.,   $[\underline{M}_{X_0}(k[\epsilon])] \simeq [\underline{M}_{\mathrm{Spec}(\underline{O}_{X_0,\ x})}(k[\epsilon])] \simeq [\underline{M}_{\mathrm{Spec}(\widehat{\underline{O}}_{X_0,\ x})}(k[\epsilon])]$ .

}{F7I-VI-096}It suffices to show, by 2.7(a), that

\SGAvar{\begin{equation*}
[\underline{M}_{X_0}] \to [\underline{M}_{\mathrm{Spec}(\underline{O}_{X_0,\ x})}] \to [\underline{M}_{\mathrm{Spec}(\widehat{\underline{O}}_{X_0,\ x})}]
\end{equation*}
are both smooth functors. Let  $(R,\ \mathrm{Spec}(A))$  be a deformation of  $X_0 = \mathrm{Spec}(A_0)$  , and let  $0 \to k \to R' \to R \to 0$  be a small extension in  $\underline{C}_{\Lambda}$  . As above, the ring homomorphisms  $R' \to R \to A$  gives us the exact commutative diagram
\SGApage{I}{107}
{\small
\[
\begin{tikzcd}[column sep=small, row sep=large]
{} \arrow[r] \& D^1(A_0|k,A_0) \arrow[r] \arrow[d, "\approx"'] \& D^1(A|R',A_0) \arrow[r, "p"] \arrow[d] \& D^1(R|R',k)\underset{k}{\otimes}A_0 \arrow[r] \arrow[d] \& D^2(A_0|k,A_0) \arrow[r] \arrow[d, "\approx"] \& \dots \\
{} \arrow[r] \& D^1(A_0|k,(A_0)_x) \arrow[r] \& D^1(A|R',(A_0)_x) \arrow[r] \& D^1(R|R',k)\underset{k}{\otimes}(A_0)_x \arrow[r] \& D^2(A_0|k,(A_0)_x) \arrow[r] \& \dots \\
{} \arrow[r] \& D^1((A_0)_x|k,(A_0)_x) \arrow[r] \arrow[u] \& D^1(A_x|R',(A_0)_x) \arrow[r, "p_x"] \arrow[u] \& D^1(R|R',k)\underset{k}{\otimes}(A_0)_x \arrow[r] \arrow[u] \& D^2((A_0)_x|k,(A_0)_x) \arrow[r] \arrow[u] \& \dots
\end{tikzcd}
\]}























where the indicated isomorphism comes from the fact that  $A_0$  is smooth over  $k$  except at one closed point  $x$  so that

$$D^i(A_0|k,\ (A_0)_x) = D^i(A_0|k,\ A_0)_x = D^i(A_0|k,\ A_0)$$
for all  $i > 0$ . Now suppose that we are given an arrow $(R',\ \mathrm{Spec}(B)) \to (R,\ \mathrm{Spec}(A_x))$  in  $[\underline{M}_{\mathrm{Spec}((A_0)_x)}]$ . It means that we are given a class $[B] \in D^1(A_x|R',\ (A_0)_x)$  such that  $p_x([B]) = [R']$ . Then a trivial diagram chasing shows that there exists $[A'] \in D^1(A|R',\ A_0)$  such that $p([A']) = [R']$  and  $A'_x = B$  i.e.,  $(R',\ \mathrm{Spec}(A'))$  is a deformation of  $\mathrm{Spec}(A_0)$ to  $R'$  such that  $\mathrm{Spec}(A')_x = \mathrm{Spec}(B)$ . This proves that $\underline{M}_{X_0} \to \underline{M}_{\mathrm{Spec}(\underline{O}_{X_0,\ x})}$ is smooth. The smoothness of the functor  $\underline{M}_{\mathrm{Spec}((A_0)_x)} \to \underline{M}_{\mathrm{Spec}(\widehat{(A_0)_x})}$ comes from the similar diagram together with the fact that

$$D^i((A_0)_x|k,\ \widehat{(A_0)_x}) = \widehat{D^i}((A_0)_x|k,\ (A_0)_x) = D^i((A_0)_x|k,\ (A_0)_x)$$
for all  $i > 0$  since  $x$  is an only possible non-smooth point.

}{\begin{equation*}
[\underline{M}_{X_0}] \to [\underline{M}_{\mathrm{Spec}(\underline{O}_{X_0,\ x})}] \to [\underline{M}_{\mathrm{Spec}(\widehat{\underline{O}}_{X_0,\ x})}]
\end{equation*}
are both smooth functors. Let  $(R,\ \mathrm{Spec}(A))$  be a deformation of  $X_0 = \mathrm{Spec}(A_0)$  , and let  $0 \to k \to R' \to R \to 0$  be a small extension in  $\underline{C}_{\Lambda}$  . As above, the ring homomorphisms  $R' \to R \to A$  give us the exact commutative diagram
\SGApage{I}{107}
{\small
\[
\begin{tikzcd}[column sep=small, row sep=large]
{} \arrow[r] \& D^1(A_0|k,A_0) \arrow[r] \arrow[d, "\approx"'] \& D^1(A|R',A_0) \arrow[r, "p"] \arrow[d] \& D^1(R|R',k)\underset{k}{\otimes}A_0 \arrow[r] \arrow[d] \& D^2(A_0|k,A_0) \arrow[r] \arrow[d, "\approx"] \& \dots \\
{} \arrow[r] \& D^1(A_0|k,(A_0)_x) \arrow[r] \& D^1(A|R',(A_0)_x) \arrow[r] \& D^1(R|R',k)\underset{k}{\otimes}(A_0)_x \arrow[r] \& D^2(A_0|k,(A_0)_x) \arrow[r] \& \dots \\
{} \arrow[r] \& D^1((A_0)_x|k,(A_0)_x) \arrow[r] \arrow[u] \& D^1(A_x|R',(A_0)_x) \arrow[r, "p_x"] \arrow[u] \& D^1(R|R',k)\underset{k}{\otimes}(A_0)_x \arrow[r] \arrow[u] \& D^2((A_0)_x|k,(A_0)_x) \arrow[r] \arrow[u] \& \dots
\end{tikzcd}
\]}























where the indicated isomorphisms come from the fact that  $A_0$  is smooth over  $k$  except at one closed point  $x$  so that

$$D^i(A_0|k,\ (A_0)_x) = D^i(A_0|k,\ A_0)_x = D^i(A_0|k,\ A_0)$$
for all  $i > 0$ . Now suppose that we are given an arrow $(R',\ \mathrm{Spec}(B)) \to (R,\ \mathrm{Spec}(A_x))$  in  $[\underline{M}_{\mathrm{Spec}((A_0)_x)}]$ . It means that we are given a class $[B] \in D^1(A_x|R',\ (A_0)_x)$  such that  $p_x([B]) = [R']$ . Then a diagram chase shows that there exists $[A'] \in D^1(A|R',\ A_0)$  such that $p([A']) = [R']$  and  $A'_x = B$  i.e.,  $(R',\ \mathrm{Spec}(A'))$  is a deformation of  $\mathrm{Spec}(A_0)$ to  $R'$  such that  $\mathrm{Spec}(A')_x = \mathrm{Spec}(B)$ . This proves that $\underline{M}_{X_0} \to \underline{M}_{\mathrm{Spec}(\underline{O}_{X_0,\ x})}$ is smooth. The smoothness of the functor  $\underline{M}_{\mathrm{Spec}((A_0)_x)} \to \underline{M}_{\mathrm{Spec}(\widehat{(A_0)_x})}$ comes from the similar diagram together with the fact that

$$D^i((A_0)_x|k,\ \widehat{(A_0)_x}) = \widehat{D^i((A_0)_x|k,\ (A_0)_x)} = D^i((A_0)_x|k,\ (A_0)_x)$$
for all  $i > 0$  since  $x$  is the only possible non-smooth point.

}{F7I-VI-097}\underline{Theorem} 4.11.  \underline{Let}  $X_0$  \underline{be a} (\underline{separated}) $k$-\underline{scheme of finite type, smooth outside a finite closed set}  $X_0^{\mathrm{Sing}} = \{x_1,\dots,x_r\}$  \underline{of}  $X_0$ . \underline{Let}  $U_0^1,\dots,U_0^r$  \underline{be affine open neighbourhoods of}  $x_1,\dots,x_r$  \underline{such that}  $U_0^i \cap X_0^{\mathrm{Sing}} = \{x_i\}$ . \underline{If} $H^2(X_0,\ \underline{D}^0_{X_0}) = 0$ , \underline{then the natural functor}  $\underline{M}_{X_0} \to \underline{M}_{U_0^1} \times\dots\times \underline{M}_{U_0^r}$   \underline{is smooth.}

\underline{Proof}: We can complete  $U_0^1, \dots ,U_0^r$  into an affine open covering of  $X_0$  by choosing $U_0^{r+1},\dots,U_0^N$  such that   $U_0^j \cap X_0^{\mathrm{Sing}} = \emptyset$  for all  $j > r$ . Now let $R' \to R$  be a small extension in  $\underline{C}_{\Lambda}$ , and let  $(R',\ U'^1),\dots,(R',\ U'^r)$  be deformations of  $U_0^1,\dots,U_0^r$  to  $R'$  such that  $U'^i \underset{R'}{\otimes} R = X|U_0^i$ ,  $1 \le i \le r$ , for
\SGApage{I}{108}
some deformation  $(R,\ X)$  of  $X_0$  to  $R$ . We must show that there exists a deformation  $(R',\ X')$  of  $X_0$  to  $R'$  such that  $X' \underset{R'}{\otimes} R = X$  and  $X'|U_0^i = U'^i$  for $1 \le i \le r$ . Since  $U_0^j$  for  $j > r$  is smooth over  $k$ ,  $X|U_0^j$  can be lifted, by 4.9(a), to a deformation  $(R',\ U'^j)$  i.e., there exist deformations  $(R',\ U'^j)$ of  $U_0^j$  to  $R'$  such that  $U'^j \underset{R'}{\otimes} R = X|U_0^j$  for all  $j > r$ . Since, for any pair $i \ne j$ ,  $U_0^i \cap U_0^j$  is affine and smooth over  $k$ , we obtain an isomorphism

\SGAvar{$$\tau_{ij} : U'^i|U_0^i \cap U_0^j \xrightarrow{\ \sim\ } U'^j|U_0^i \cap U_0^j$$
such that  $\tau_{ij} \underset{R'}{\otimes} R = \mathrm{id}_{X|U_0^i \cap U_0^j}$ , and in turn defines a cohomology class in $H^2(X_0\ \underline{D}^0_{X_0})$  given by  $U_0^i \cap U_0^j \cap U_0^k \to \tau_{ij}\tau_{jk}\tau_{ik}^{\ -1}$ . Since  $H^2(X_0,\ \underline{D}^0_{X_0}) = 0$  by hypothesis, the obstruction for glueing these to obtain a global deformation vanishes, and hence we get a deformation  $(R',\ X')$  of  $X_0$  to  $R'$  such that $X' \underset{R'}{\otimes} R = X$  and  $X'|U_0^j \simeq U'^j$ . This completes our proof.

\underline{Remark} 4.12.  The previous proof assumes  $X_0$  to be separated, which is unnecessary. In fact, for a fixed deformation  $(R,\ X)$  of  $X_0$  and deformations  $(R',\ U'^j)$  of $U_0^i$ , the deformations  $(R',\ V')$  of open set  $V$  of  $X$  to  $R'$  compatible with the deformation  $(R',\ U'^j)$  already given, form a  "gerbe"  (in Giraud's terminology) whose structure sheaf is  $\underline{D}^0_{X_0}$ . Hence the obstruction to the existence of a section of this  "gerbe"  (i.e., a global deformation  $(R',\ X')$  of  $X$  compatible with the given datum  $(R',\ U'^j))$  is in $H^2(X_0,\ \underline{D}^0_{X_0})$  which is zero by assumption.

}{$$\tau_{ij} : U'^i|U_0^i \cap U_0^j \xrightarrow{\ \sim\ } U'^j|U_0^i \cap U_0^j$$
such that  $\tau_{ij} \underset{R'}{\otimes} R = \mathrm{id}_{X|U_0^i \cap U_0^j}$ , and in turn defines a cohomology class in $H^2(X_0,\ \underline{D}^0_{X_0})$  given by  $U_0^i \cap U_0^j \cap U_0^k \to \tau_{ij}\tau_{jk}\tau_{ik}^{\ -1}$ . Since  $H^2(X_0,\ \underline{D}^0_{X_0}) = 0$  by hypothesis, the obstruction for glueing these to obtain a global deformation vanishes, and hence we get a deformation  $(R',\ X')$  of  $X_0$  to  $R'$  such that $X' \underset{R'}{\otimes} R = X$  and  $X'|U_0^j \simeq U'^j$ . This completes our proof.

\underline{Remark} 4.12.  The previous proof assumes  $X_0$  to be separated, which is unnecessary. In fact, for a fixed deformation  $(R,\ X)$  of  $X_0$  and deformations  $(R',\ U'^j)$  of $U_0^j$ , the deformations  $(R',\ V')$  of open set  $V$  of  $X$  to  $R'$  compatible with the deformation  $(R',\ U'^j)$  already given, form a  "gerbe"  (in Giraud's terminology) whose structure sheaf is  $\underline{D}^0_{X_0}$ . Hence the obstruction to the existence of a section of this  "gerbe"  (i.e., a global deformation  $(R',\ X')$  of  $X$  compatible with the given datum  $(R',\ U'^j))$  is in $H^2(X_0,\ \underline{D}^0_{X_0})$  which is zero by assumption.

}{F7I-VI-098}\underline{Corollary} 4.13.  \underline{Let}  $X_0$  \underline{be a} $k$-\underline{scheme of finite type, smooth outside a finite subset}  $X_0^{\mathrm{Sing}} = \{x_1,\dots,x_r\}$  of  $X_0$  \underline{such that}  $H^2(X_0,\ \underline{D}^0_{X_0}) = 0$ . \underline{Then}

\SGAvar{$$\text{(a)}\quad \underline{M}_{X_0} \to \underline{M}_{\mathrm{Spec}(\widehat{\underline{O}}_{X_0,x_1})} \times\dots\times \underline{M}_{\mathrm{Spec}(\widehat{\underline{O}}_{X_0,x_r})}\quad \underline{\text{is smooth.}}$$
$$\underline{\text{In particular,}}\ \ \dim \underline{M}_{X_0} = \dim H^1(X_0,\ \underline{D}^0_{X_0}) + \sum_{c=1}^{r} \ \dim \underline{M}_{\mathrm{Spec}(\widehat{\underline{O}}_{X_0,x_0})}\ .$$
(b)  \underline{If furthermore}  $X_0$  \underline{is locally a complete-intersection, then}  $\underline{M}_{X_0}$ \underline{is smooth.}

}{$$\text{(a)}\quad \underline{M}_{X_0} \to \underline{M}_{\mathrm{Spec}(\widehat{\underline{O}}_{X_0,x_1})} \times\dots\times \underline{M}_{\mathrm{Spec}(\widehat{\underline{O}}_{X_0,x_r})}\quad \underline{\text{is smooth.}}$$
$$\underline{\text{In particular,}}\ \ \dim \underline{M}_{X_0} = \dim H^1(X_0,\ \underline{D}^0_{X_0}) + \sum_{c=1}^{r} \ \dim \underline{M}_{\mathrm{Spec}(\widehat{\underline{O}}_{X_0,x_c})}\ .$$
(b)  \underline{If furthermore}  $X_0$  \underline{is locally a complete-intersection, then}  $\underline{M}_{X_0}$ \underline{is smooth.}

}{F7I-VI-099}\SGApage{I}{109}

\underline{Proof}:

(a)  The first statement follows from 4.10 and 4.9(b). Now let  $(R, \mathfrak{X})$  and  $(R_1, \mathfrak{X}_1),...,(R_r, \mathfrak{X}_r)$  be formal versal deformations of  $X_0$  and  $\mathrm{Spec}(\underline{O}_{X_0,x_1}),...,\mathrm{Spec}(\underline{O}_{X_0,x_r})$  respectively. We then obtain a commutative diagram

\begin{center}
\begin{tikzcd}[column sep=huge, row sep=huge]
\underline{h}_R \arrow[r] \arrow[d] \& \underline{h}_{R_1 \widehat{\otimes}_\Lambda ... \widehat{\otimes}_\Lambda R_r} \arrow[d] \\
\underline{M}_{X_0} \arrow[r] \& \underline{M}_{\mathrm{Spec}(\widehat{\underline{O}}_{X_0,x_1})} \times ... \times \underline{M}_{\mathrm{Spec}(\widehat{\underline{O}}_{X_0,x_r})}
\end{tikzcd}
\end{center}















Since the composite functor

\begin{equation*}
\underline{h}_R \to \underline{M}_{X_0} \to \underline{M}_{\mathrm{Spec}(\widehat{\underline{O}}_{X_0,x_1})} \times ... \times \underline{M}_{\mathrm{Spec}(\widehat{\underline{O}}_{X_0,x_r})}
\end{equation*}
is smooth whereas

\begin{equation*}
\underline{h}_{R_1 \widehat{\otimes}_\Lambda ... \widehat{\otimes}_\Lambda R_r} \to \underline{M}_{\mathrm{Spec}(\widehat{\underline{O}}_{X_0,x_1})} \times .... \times \underline{M}_{\mathrm{Spec}(\widehat{\underline{O}}_{X_0,x_r})}
\end{equation*}
is minimally smooth, it follows from 2.7(a) that

$\underline{h}_R \to \underline{h}_{R_1 \widehat{\otimes}_\Lambda ... \widehat{\otimes}_\Lambda R_r}$   is smooth i.e.,  $R$  is a formal power-series ring over  $R_1 \widehat{\otimes}_\Lambda ... \widehat{\otimes}_\Lambda R_r$ . Since

\SGAvar{\begin{equation*}
0 \to H^1(X_0,\underline{D}^0_{X_0}) \to [\underline{M}_{X_0}(k[\epsilon])] \to [\underline{M}_{\mathrm{Spec}(\widehat{\underline{O}}_{X_0,x_1})}(k[\epsilon])] ... [\underline{M}_{\mathrm{Spec}(\widehat{\underline{O}}_{X_0,x_r})}(k[\epsilon])]
\end{equation*}
is exact by 4.4 i.e.,

}{\begin{equation*}
0 \to H^1(X_0,\underline{D}^0_{X_0}) \to [\underline{M}_{X_0}(k[\epsilon])] \to [\underline{M}_{\mathrm{Spec}(\widehat{\underline{O}}_{X_0,x_1})}(k[\epsilon])] \times ... \times [\underline{M}_{\mathrm{Spec}(\widehat{\underline{O}}_{X_0,x_r})}(k[\epsilon])]
\end{equation*}
is exact by 4.4 i.e.,

}{F7I-VI-100}\begin{equation*}
0 \to H^1(X_0, \underline{D}^0_{X_0}) \to \mathrm{Hom}_{\Lambda-\mathrm{alg}}(R, k[\epsilon]) \to \mathrm{Hom}_{\Lambda-\mathrm{alg}}(R_1 \widehat{\otimes}_\Lambda ... \widehat{\otimes}_\Lambda R_r, k[\epsilon])
\end{equation*}
is exact, it follows that  $R$  is a formal power-series ring over    $R_1 \widehat{\otimes}_\Lambda ... \widehat{\otimes}_\Lambda R_r$ in $d$-variables where  $d = \dim_k H^1(X_0, \underline{D}^0_{X_0})$  .

Therefore

\begin{equation*}
\begin{aligned}
\dim \underline{M}_{X_0} &= \dim k \otimes_\Lambda R = d + \dim k \otimes_\Lambda (R_1 \widehat{\otimes}_\Lambda ... \widehat{\otimes}_\Lambda R_r) \\
&= d + \dim(k \otimes_\Lambda R_1) \widehat{\otimes}_k ... \widehat{\otimes}_k (k \otimes_\Lambda R_r) \\
&= d + \sum_{i=1}^{r} \dim k \otimes_\Lambda R_i = d + \sum_{i=1}^{r} \dim \underline{M}_{\mathrm{Spec}(\widehat{\underline{O}}_{X_0,x_i})}
\end{aligned}
\end{equation*}

\SGApage{I}{110}

(b)  follows from 4.10 and 4.9(a).

\SGAvar{4.14.    We close this paragraph with an explicit construction of a formal versal deformation in a simple situation. In fact, if  $X_0$  is affine and is a complete-intersection, it is trivial to write down explicitely. Thus let

\begin{equation*}
A_0 = k[x_1,...,x_n]/(f_1,...,f_r) \qquad S/J
\end{equation*}
where  $S = k[x_1,...,x_n]$  ,  $J = (f_1,...,f_r)$  , and  $f_1,...,f_r$  is an $S$-regular sequence. We then obtain the exact sequence

\begin{equation*}
... \to \mathrm{Der}_k(S, A_0) \to \mathrm{Hom}_{A_0}(J/J^2, A_0) \to D^1(A_0|k, A_0) \to 0
\end{equation*}
We note that  $J/J^2$  is  $A_0$-free module of rank  $r$  generated by $f_1(\mathrm{mod}\ J^2),...,f_r(\mathrm{mod}\ J^2)$  , since  $f_1,...,f_r$  is an $S$-regular sequence. We now assume that  $D^1(A_0|k, A_0)$  is a finite-dimensional vector space over  $k$  , say  $A$ has isolated non-smooth points only. Let  $d = \dim_k D^1(A_0|k, A_0)$  and choose $p_j : J/J^2 \to A_0$   $(j = 1, 2,...,d)$  which form a $k$-basis of  $D^1(A_0|k, A_0)$  . Now choose  $P_{jk} \in \Lambda[x_1,...,x_n]$  such that  $p_j(f_i) = P_{ji}$  (mod $\underline{m}_\Lambda + J$)  where  $\underline{m}_\Lambda$  is the maximal ideal of  $\Lambda$  . We then consider

\begin{equation*}
\begin{aligned}
R &= \Lambda[[t_1,...,t_d]] \\
B &= R[x_1,...,x_n]/(G_1,...,G_r)
\end{aligned}
\end{equation*}
where  $G_j = F_j - \sum_{i=1}^{d} t_i P_{ji}$  and  $F_j$  is a polynomial with coefficients in  $\Lambda$ gotten from  $f_j$  by lifting the coefficients from  $k$  to  $\Lambda$  . Since $R[x_1,...,x_n]$  is $R$-flat and  $G_1,...,G_r$  modulo  $\underline{m}_R$  (which becomes  $f_1,...,f_r$) is a regular sequence in  $k \otimes_R R[x_1,...,x_n]$  , it follows that  $B$  is $R$-flat, and $k \otimes_R B = A_0$  . Furthermore,

}{4.14.    We close this paragraph with an explicit construction of a formal versal deformation in a simple situation. In fact, if  $X_0$  is affine and is a complete-intersection, it is trivial to write down explicitly. Thus let

\begin{equation*}
A_0 = k[x_1,...,x_n]/(f_1,...,f_r) = S/J
\end{equation*}
where  $S = k[x_1,...,x_n]$  ,  $J = (f_1,...,f_r)$  , and  $f_1,...,f_r$  is an $S$-regular sequence. We then obtain the exact sequence

\begin{equation*}
... \to \mathrm{Der}_k(S, A_0) \to \mathrm{Hom}_{A_0}(J/J^2, A_0) \to D^1(A_0|k, A_0) \to 0
\end{equation*}
We note that  $J/J^2$  is  $A_0$-free module of rank  $r$  generated by $f_1(\mathrm{mod}\ J^2),...,f_r(\mathrm{mod}\ J^2)$  , since  $f_1,...,f_r$  is an $S$-regular sequence. We now assume that  $D^1(A_0|k, A_0)$  is a finite-dimensional vector space over  $k$  , say  $A_0$ has isolated non-smooth points only. Let  $d = \dim_k D^1(A_0|k, A_0)$  and choose $p_j : J/J^2 \to A_0$   $(j = 1, 2,...,d)$  which form a $k$-basis of  $D^1(A_0|k, A_0)$  . Now choose  $P_{jk} \in \Lambda[x_1,...,x_n]$  such that  $p_j(f_i) = P_{ji}$  (mod $\underline{m}_\Lambda + J$)  where  $\underline{m}_\Lambda$  is the maximal ideal of  $\Lambda$  . We then consider

\begin{equation*}
\begin{aligned}
R &= \Lambda[[t_1,...,t_d]] \\
B &= R[x_1,...,x_n]/(G_1,...,G_r)
\end{aligned}
\end{equation*}
where  $G_j = F_j - \sum_{i=1}^{d} t_i P_{ij}$  and  $F_j$  is a polynomial with coefficients in  $\Lambda$ gotten from  $f_j$  by lifting the coefficients from  $k$  to  $\Lambda$  . Since $R[x_1,...,x_n]$  is $R$-flat and  $G_1,...,G_r$  modulo  $\underline{m}_R$  (which becomes  $f_1,...,f_r$) is a regular sequence in  $k \otimes_R R[x_1,...,x_n]$  , it follows that  $B$  is $R$-flat, and $k \otimes_R B = A_0$  . Furthermore,

}{F7I-VI-101}\begin{equation*}
\mathrm{Hom}_{\Lambda-\mathrm{alg}}(R, k[\epsilon]) \to D^1(A_0|k, A_0) = [\underline{M}_{\mathrm{Spec}(A_0)}(k[\epsilon])]
\end{equation*}
is bijective by our construction of  $B$  . Therefore, if we set 


\SGApage{I}{111}

\begin{equation*}
 \mathrm{Spf}(B) \;=\; \varinjlim_{v}\; \mathrm{Spec}(R/\underline{m}^{v+1}\ \underset{R}{\otimes}\ B)
\end{equation*}
then  $(R,\ \mathrm{Spf}(B))$  is an object in  $\widehat{\underline{M}}_{\mathrm{Spec}(A_0)}$  and the functor  $\underline{h}_R \to \underline{M}_{\mathrm{Spec}(A_0)}$  induced by  $(R,\ \mathrm{Spf}(B))$  is smooth by 1.6, since  $R$  is a formal power-series ring over  $\Lambda$  . In other words,  $(R,\ \mathrm{Spf}(B))$  is a formal versal object of  $\underline{M}_{\mathrm{Spec}(A_0)}$  .

We note, by 4.10(b), that if we set

\begin{equation*}
B = R[[X_1,\ldots,X_n]]/(G_1,\ldots,G_r)\ \ ,
\end{equation*}
then  $(R,\ \mathrm{Spf}(B))$  is a formal versal deformation of  $\mathrm{Spec}(A_0)$  where  $A_0 = k[[X_1,\ldots,X_n]]/(f_1,\ldots,f_r)$  .

\underline{Remark} 4.15.  Let  $X_0$  be a $k$-scheme of finite type with isolated non-smooth points only, and let  $(R,\ \mathfrak{X})$  be a formal versal deformation of  $X_0$  . The above construction shows that if  $X_0$  is affine and locally a complete-intersection, then  $\mathfrak{X}$  is a formalization of an $R$-scheme of finite type. This fact can easily be generalized to the case when  $X_0$  is not necessarily affine but  $H^2(X_0,\ \underline{D}^0_{X_0}) = 0$  .

\SGAvar{\underline{Remark} 4.16.  The \underline{global}  theory of the relative cotangent complex (the latter theory being reduced to the truncated complex of length 1 , which is sufficient however for many purposes of deformation theory) allows to generalize and simplify certain results, such as 4.10. Namely let  $S$  be a scheme,  $S_0 \subset S' \subset S$  two subschemes of  $S$  defined by ideals  $J \supset I$  such that  $JI = 0$  ,  $X'$  a scheme flat over  $S'$  ,  $L. = L^{X_0/S_0}$  the relative cotangent complex of  $X_0$  over  $S_0$  , which is an element of the derived category  $D(X_0,\ \underline{O}_X)$  . We consider the category  $\underline{F}$  of all flat schemes  $X$  over  $S$  which extend  $X'/S'$  , i.e. together with an  $S'$-isomorphism  $X \underset{S}{\times} S' \simeq X'$  . Then

(a)  For an object  $X'$  of  $\underline{F}$  , the group of automorphisms is canonically isomorphic to  $\mathrm{Ext}^0(X_0;\ L.,I_{X_0})$  (where  $I_{X_0}$  is the inverse image of  $I$  , viewed as a module on  $S_0$  , on  $X_0)$  .


}{\underline{Remark} 4.16.  The \underline{global}  theory of the relative cotangent complex (the latter theory being reduced to the truncated complex of length 1 , which is sufficient however for many purposes of deformation theory) allows to generalize and simplify certain results, such as 4.10. Namely let  $S$  be a scheme,  $S_0 \subset S' \subset S$  two subschemes of  $S$  defined by ideals  $J \supset I$  such that  $JI = 0$  ,  $X'$  a scheme flat over  $S'$  ,  $L. = L^{X_0/S_0}$  the relative cotangent complex of  $X_0$  over  $S_0$  , which is an element of the derived category  $D(X_0,\ \underline{O}_{X_0})$  . We consider the category  $\underline{F}$  of all flat schemes  $X$  over  $S$  which extend  $X'/S'$  , i.e. together with an  $S'$-isomorphism  $X \underset{S}{\times} S' \simeq X'$  . Then

(a)  For an object  $X$  of  $\underline{F}$  , the group of automorphisms is canonically isomorphic to  $\mathrm{Ext}^0(X_0;\ L.,I_{X_0})$  (where  $I_{X_0}$  is the inverse image of  $I$  , viewed as a module on  $S_0$  , on  $X_0)$  .


}{F7I-VI-102}\SGApage{I}{112}

(b)  The set of isomorphism classes of objects of  $\underline{F}$  is empty or a torsor (= a principal homogeneous space) under the group  $\mathrm{Ext}^1(X_0;\ L.,I_{X_0})$  .

(c)  One can define a canonical element  $c \in \mathrm{Ext}^2(X_0;\ L.,I_{X_0})$  whose vanishing is necessary and sufficient for  $\underline{F}$  to be non empty.

In order to study the  $\mathrm{Ext}^i(X_0;\ L.,I_{X_0})$  , it is often convenient to write

\SGAvar{\begin{equation*}
\mathrm{Ext}^i(X_0;L.,I_{X_0}) = H^i(X_0,\ \underline{RHom}(L.,I_{X_0}))\ \ .
\end{equation*}
and to notice that the complex  $\underline{RHom}(L.,I_{X_0})$  has as cohomology sheaves the sheaves  $\underline{D}^i_{X_0}(I_{X_0})$  introduced in §3. For  $i = 0, 1, 2$,  these sheaves may be viewed respectively as the sheaf of automorphisms of any object  $X$  of  $\underline{F}$  (understood:  automorphisms inducing the identity on  $X')$  , the sheaf of indeterminacy of isomorphism classes of local solutions to the flat deformation problem, and the sheaf of obstructions to the existence of local solutions. On the other hand, the standart spectral sequence gives us

\begin{equation*}
\mathrm{Ext}^*(X_0 \quad L.,I_{X_0}) \leftarrow E_2^{p,q} = H^p(X_0,\ \underline{D}^q_{X_0}(I_{X_0}))\ \ ,
\end{equation*}
which yields the well-known isomorphism

}{\begin{equation*}
\mathrm{Ext}^i(X_0;L.,I_{X_0}) = H^i(X_0,\ \underline{RHom}(L.,I_{X_0}))\ \ .
\end{equation*}
and to notice that the complex  $\underline{RHom}(L.,I_{X_0})$  has as cohomology sheaves the sheaves  $\underline{D}^i_{X_0}(I_{X_0})$  introduced in §3. For  $i = 0, 1, 2$,  these sheaves may be viewed respectively as the sheaf of automorphisms of any object  $X$  of  $\underline{F}$  (understood:  automorphisms inducing the identity on  $X')$  , the sheaf of indeterminacy of isomorphism classes of local solutions to the flat deformation problem, and the sheaf of obstructions to the existence of local solutions. On the other hand, the standard spectral sequence gives us

\begin{equation*}
\mathrm{Ext}^*(X_0;\ L.,I_{X_0}) \leftarrow E_2^{p,q} = H^p(X_0,\ \underline{D}^q_{X_0}(I_{X_0}))\ \ ,
\end{equation*}
which yields the well-known isomorphism

}{F7I-VI-103}\begin{equation*}
\mathrm{Aut}(X/X') \simeq H^0(X_0,\underline{D}^0_{X_0}(I_{X_0})) \simeq H^0(X_0,\underline{Hom}(\underline{\Omega}^1_{X_0/S_0},I_{X_0})) = \mathrm{Hom}(\underline{\Omega}^1_{X_0/S_0},I_{X_0})\ \ ,
\end{equation*}
and the exact sequence in low degrees

\SGAvar{\begin{equation*}
0 - H^1(X_0,\underline{D}^0_{X_0}) - \mathrm{Ext}^1 - H^0(X_0,\underline{D}^1_{X_0}) - H^2(X_0,\underline{D}^0_{X_0}) - \mathrm{Ext}^2
\end{equation*}
This shows in particular that if  $H^2(X_0,\underline{D}^0_{X_0}) = 0$  , and if there exists any local flat deformation  $X$  of  $X'/S'$  , then there exists a flat global deformation which corresponds to given local deformations. More precisely and generally, for a given local deformations (i.e. a section of the  "sheaf of isomorphism classes of local flat deformations" ), the argument 4.10 shows that the obstruction to finding a global flat deformation compatible with these is in  $H^2(X_0,\ \underline{D}^0_{X_0})$  , hence such a global flat deformation exists if  $H^2(X_0,\ \underline{D}^0_{X_0}) = 0$  : this is essentially 4.10.


}{\begin{equation*}
0 - H^1(X_0,\underline{D}^0_{X_0}) - \mathrm{Ext}^1 - H^0(X_0,\underline{D}^1_{X_0}) - H^2(X_0,\underline{D}^0_{X_0}) - \mathrm{Ext}^2
\end{equation*}
This shows in particular that if  $H^2(X_0,\underline{D}^0_{X_0}) = 0$  , and if there exists any local flat deformation  $X$  of  $X'/S'$  , then there exists a flat global deformation which corresponds to given local deformations. More precisely and generally, for given local deformations (i.e. a section of the  "sheaf of isomorphism classes of local flat deformations" ), the argument 4.10 shows that the obstruction to finding a global flat deformation compatible with these is in  $H^2(X_0,\ \underline{D}^0_{X_0})$  , hence such a global flat deformation exists if  $H^2(X_0,\ \underline{D}^0_{X_0}) = 0$  : this is essentially 4.10.


}{F7I-VI-104}\SGApage{I}{113}

On the other hand, the spectral sequence shows that the  $\mathrm{Ext}^2$  is zero whenever we have the relations

\SGAvar{\begin{equation}\tag{*}
H^0(X_0,\ \underline{D}^2_{X_0}\ ) = 0 \quad ,\quad H^1(X_0,\ \underline{D}^1_{X_0}\ ) = 0 \quad ,\quad H^2(X_0,\ \underline{D}^0_{X_0}\ ) = 0 \quad .
\end{equation}
The first relation is always satisfied if  $X_0$  is a relative complete intersection over  $S_0$  (as then  $\underline{D}^2_{X_0} = 0$  for any module  $M$  on  $X_0$ ), the second is satisfied provided the support of  $\underline{D}^1_{X_0}$  is discrete, as is the case if  $X_0$  is smooth over  $S_0$  except on a discrete subset of  $X_0$  over  $S_0$ . When the previous three relations $(*)$ are all satisfied, then by (c) above,  $\underline{F}$  is non empty, which generalizes 4.12(b).

More geometrically, and without using the global theory of the relative cotangentcomplex (but only the globalizing of the functors  $\underline{D}^i_{X_0}$  ), one can get a simple geometric interpretation of the three successive obstructions to deforming flatly  $X'$  to  $X$ , lying in the three left hand terms of $(*)$. Namely by the local theory, the first obstruction  $c_1$  in  $H^0(X_0,\ \underline{D}^2_{X_0})$  can be defined as the obstruction to finding local solutions, when this obstruction vanishes, the sheaf  $\underline{F}$  of isomorphism classes of local solutions of the problem is (by the local theory) in a canonical way a torsor under  $\underline{D}^1_{X_0}$ , and the second obstruction  $c_2$ , lying in  $H^1(X_0,\ \underline{D}^1_{X_0})$ , is the obstruction to finding a section of this torsor, i.e. of finding a  "coherent system of isomorphism classes of local solutions" ; if this second obstruction also vanishes, then, choosing arbitrarily a section of  $\underline{F}$ , one gets, as noted above, an obstruction  $c_3$  in  $H^2(X_0,\ \underline{D}^0_{X_0})$  to finding a global solution corresponding to the given system of (isomorphism classes of) local solutions. Using the fact that the section of  $F$  is determined up to adding a section of  $\underline{D}^1_{X_0}$ , we get a map  $d : H^0(X_0,\ \underline{D}^1_{X_0}) \to H^2(X_0,\ \underline{D}^0_{X_0})$ , and a solution to the deformation problem exists if and only if  $c_1 = 0$ ,  $c_2 = 0$ , and  $c_3$  is zero as an element of Coker $d$  (each  $c_i$  being defined when the previous ones are zero). Of course, the groups  $H^0(X_0,\ \underline{D}^2_{X_0})$ ,  $H^1(X_0,\ \underline{D}^1_{X_0})$  and
\SGApage{I}{114}
$H^2(X_0,\ \underline{D}^0_{X_0})/\mathrm{Im}\ d$  are just the initial terms  $E_2^{0,2}$ ,  $E_2^{1,1}$ , and  $E_3^{2,0}$  of the spectral sequence above,  $d$  being the differential (up to sign ?)

}{\begin{equation}\tag{*}
H^0(X_0,\ \underline{D}^2_{X_0}\ ) = 0 \quad ,\quad H^1(X_0,\ \underline{D}^1_{X_0}\ ) = 0 \quad ,\quad H^2(X_0,\ \underline{D}^0_{X_0}\ ) = 0 \quad .
\end{equation}
The first relation is always satisfied if  $X_0$  is a relative complete intersection over  $S_0$  (as then  $\underline{D}^2_{X_0} = 0$  for any module  $M$  on  $X_0$ ), the second is satisfied provided the support of  $\underline{D}^1_{X_0}$  is discrete, as is the case if  $X_0$  is smooth over  $S_0$  except on a discrete subset of  $X_0$  over  $S_0$ . When the previous three relations $(*)$ are all satisfied, then by (c) above,  $\underline{F}$  is non empty, which generalizes 4.13(b).

More geometrically, and without using the global theory of the relative cotangent complex (but only the globalizing of the functors  $\underline{D}^i_{X_0}$  ), one can get a simple geometric interpretation of the three successive obstructions to deforming flatly  $X'$  to  $X$ , lying in the three left hand terms of $(*)$. Namely by the local theory, the first obstruction  $c_1$  in  $H^0(X_0,\ \underline{D}^2_{X_0})$  can be defined as the obstruction to finding local solutions, when this obstruction vanishes, the sheaf  $\underline{F}$  of isomorphism classes of local solutions of the problem is (by the local theory) in a canonical way a torsor under  $\underline{D}^1_{X_0}$ , and the second obstruction  $c_2$ , lying in  $H^1(X_0,\ \underline{D}^1_{X_0})$ , is the obstruction to finding a section of this torsor, i.e. of finding a  "coherent system of isomorphism classes of local solutions" ; if this second obstruction also vanishes, then, choosing arbitrarily a section of  $\underline{F}$ , one gets, as noted above, an obstruction  $c_3$  in  $H^2(X_0,\ \underline{D}^0_{X_0})$  to finding a global solution corresponding to the given system of (isomorphism classes of) local solutions. Using the fact that the section of  $\underline{F}$  is determined up to adding a section of  $\underline{D}^1_{X_0}$ , we get a map  $d : H^0(X_0,\ \underline{D}^1_{X_0}) \to H^2(X_0,\ \underline{D}^0_{X_0})$ , and a solution to the deformation problem exists if and only if  $c_1 = 0$ ,  $c_2 = 0$ , and  $c_3$  is zero as an element of Coker $d$  (each  $c_i$  being defined when the previous ones are zero). Of course, the groups  $H^0(X_0,\ \underline{D}^2_{X_0})$ ,  $H^1(X_0,\ \underline{D}^1_{X_0})$  and
\SGApage{I}{114}
$H^2(X_0,\ \underline{D}^0_{X_0})/\mathrm{Im}\ d$  are just the initial terms  $E_2^{0,2}$ ,  $E_2^{1,1}$ , and  $E_3^{2,0}$  of the spectral sequence above,  $d$  being the differential (up to sign ?)

}{F7I-VI-105}\begin{equation*}
d_2 : E_2^{0,1} \to E_2^{2,0} \quad .
\end{equation*}

\SGAsectiontoc{5.  Formal Jacobian subschemes}{5}
5.  \underline{Formal Jacobian subschemes}

Let  $X_0$  be a scheme of finite type over a field  $k$ , and let  $(R,\ \mathfrak{X})$  be a formal versal deformation of  $X_0$ . Then  $\mathfrak{X}$  is not an ordinary scheme over  $R$  but an $R$-adic formal scheme. The purpose of this section is to clarify  "the non-smooth locus of  $\mathfrak{X}$  over  $R$ ".

\SGAvar{We fistly recall the definition of Fitting ideals associated to a module of finite type: Let  $A$  be a commutative ring and  $E$  an $A$-module of finite type. Choose a presentation

}{We first recall the definition of Fitting ideals associated to a module of finite type: Let  $A$  be a commutative ring and  $E$  an $A$-module of finite type. Choose a presentation

}{F7I-VI-106}\begin{equation*}
K \to F \to E \to 0
\end{equation*}
where  $F$  is locally-free of constant rank  $n$ . We define

\SGAvar{\begin{equation*}
I^P_A(E) = \mathrm{Im}(\overset{n-p}{\wedge} K \otimes \overset{n-p}{\wedge} F^* \to A) \qquad \text{for}\quad p = 0,\ 1,\ 2,\ \ldots
\end{equation*}
The ideals  $I^P_A(E)$  thus defined does not depend on the choice of a presentation and therefore are invariants of $A$-module  $E$ , called the Fitting ideals of $A$-module  $E$ . The following properties readily follow from the definition.

}{\begin{equation*}
I^p_A(E) = \mathrm{Im}(\overset{n-p}{\wedge} K \otimes \overset{n-p}{\wedge} F^* \to A) \qquad \text{for}\quad p = 0,\ 1,\ 2,\ \ldots
\end{equation*}
The ideals  $I^p_A(E)$  thus defined do not depend on the choice of a presentation and therefore are invariants of the $A$-module  $E$ , called the Fitting ideals of the $A$-module  $E$ . The following properties readily follow from the definition.

}{F7I-VI-107}5.1.    (a) For any ring homomorphism  $A \to B$ , the canonical map

\SGAvar{\begin{equation*}
B \underset{A}{\otimes} A/I^p_A(E) \overset{\sim}{\to} B/I^p_B(B \underset{A}{\otimes} E)
\end{equation*}
is an isomorphism for all  $P$ . In particular we have

}{\begin{equation*}
B \underset{A}{\otimes} A/I^p_A(E) \overset{\sim}{\to} B/I^p_B(B \underset{A}{\otimes} E)
\end{equation*}
is an isomorphism for all  $p$ . In particular we have

}{F7I-VI-108}\begin{equation*}
I^p_{S^{-1}A}(S^{-1}E) = S^{-1}I^p_A(E)
\end{equation*}
for any multiplicative subset  $S$  in  $A$ .


\SGApage{I}{115}

(b)  We have  $I_A^0(E) \subset I_A^1(E) \subset I_A^2(E) \subset \ \dots$ , and, for each point $x \in \operatorname{Spec}(A)$ , $I_A^p(E)_x = A_x$  for all  $p \geq \dim_{\varkappa(x)} E(x)$ , where  $E(x) = \varkappa(x) \underset{A}{\otimes} E$ .

(c)  $x \in \operatorname{Spec}(A) - \operatorname{Supp} A/I_A^p(E)$  if and only if, for any (or for some) presentation

\begin{equation*}
K \to F \to E \to 0
\end{equation*}
where  $F \approx A^n$ , $K_x$  contains  $A_x^{n-p}$  which is isomorphic to a direct summand of $F_x$ . In particular, we have that

\begin{equation*}
\operatorname{Supp} A/I_A^0(E) = \operatorname{Supp} E \ .
\end{equation*}
(d)

\begin{equation*}
I_A^r(E_1 \oplus E_2) = \sum_{p+q=r} \ I_A^p(E_1) \cdot I_A^q(E_2)
\end{equation*}
(e)  If  $A$  is a noetherian adic ring, then

\SGAvar{\begin{equation*}
A/I_A^p(E) = \varprojlim_{v} A_v \big/ I_{A_v}^p (E_v)
\end{equation*}
where  $A = \varprojlim_{v} A_v$  and  $E_v = E \underset{A}{\otimes} A_v$ .

}{\begin{equation*}
A/I_A^p(E) = \varprojlim_{\nu} A_\nu \big/ I_{A_\nu}^p (E_\nu)
\end{equation*}
where  $A = \varprojlim_{\nu} A_\nu$  and  $E_\nu = E \underset{A}{\otimes} A_\nu$ .

}{F7I-VI-109}In view of the property 5.1(a), our definition of Fitting ideals is readily globalized. Thus let  $X$  be any scheme, and  $\underline{E}$  a quasi-coherent $X$-module of finite presentation. We then obtain a sequence of quasi-coherent ideal sheaves $\underline{I}_X^p(\underline{E})$  for  $p \geq 0$  defined by

\SGAvar{\begin{equation*}
\Gamma(U, \ \underline{I}_X^p(\underline{E})) \xrightarrow{\sim} I^p_{\Gamma(U, \ \underline{O}_x)}(\Gamma(U, \ \underline{E}))
\end{equation*}
for every affine open subset  $U \subset X$ .

Now let  $X \to Y$  be a morphism of schemes. The relative Kähler-differentials $\underline{\Omega}_{X|Y}(= \underline{\Omega}^1_{X|Y})$  is a quasi-coherent $X$-module. If the morphism  $X \to Y$  is of essentially finite presentation, then  $\underline{\Omega}_{X|Y}$  is a quasi-coherent $X$-module of finite presentation, and in turn we may consider the ideal sheaves  $\underline{I}_X^p(\underline{\Omega}_{X|Y})$  ..

}{\begin{equation*}
\Gamma(U, \ \underline{I}_X^p(\underline{E})) \xrightarrow{\sim} I^p_{\Gamma(U, \ \underline{O}_X)}(\Gamma(U, \ \underline{E}))
\end{equation*}
for every affine open subset  $U \subset X$ .

Now let  $X \to Y$  be a morphism of schemes. The relative module of Kähler differentials $\underline{\Omega}_{X|Y}(= \underline{\Omega}^1_{X|Y})$  is a quasi-coherent $X$-module. If the morphism  $X \to Y$  is of essentially finite presentation, then  $\underline{\Omega}_{X|Y}$  is a quasi-coherent $X$-module of finite presentation, and in turn we may consider the ideal sheaves  $\underline{I}_X^p(\underline{\Omega}_{X|Y})$  .

}{F7I-VI-110}\underline{Definition} 5.2.  \underline{Let}  $X \to Y$  \underline{be a morphism of essentially finite presentation.} \underline{For each integer}  $p \geq 0$ , \underline{we define}  $J_X^p(X|Y)$  \underline{to be the closed subscheme of}  $X$
\SGApage{I}{116}
\underline{defined by the ideal sheaf}  $\underline{I}_X^p(X|Y) = \underline{I}_X^p(\underline{\Omega}_{X|Y})$ , \underline{called the Jacobian subschemes} \underline{of}  $X$  \underline{over}  $Y$ . \underline{For each point}  $x \in X$ , \underline{we set}  $X_x = \operatorname{Spec}(\underline{O}_{X,x})$ .

\underline{Theorem} 5.3.  \underline{Let}  $X \to Y$  \underline{be a morphism of essentially finite presentation. Then}

(1)  \underline{For any base change}  $Y \leftarrow Y'$ , \underline{the canonical map} $J^p_{X'}(X'|Y') \xrightarrow{\sim} J^p_X(X|Y) \underset{Y}{\times} Y'$  \underline{is an isomorphism for all}  $p$  \underline{where}  $X' = X \underset{Y}{\times} Y'$ .

(2)  \underline{For each point}  $x \in X$ , \underline{we have the canonical isomorphism} $J^p_{X_x}(X_x|Y) \xrightarrow{\sim} J^p_X(X|Y)_x$ .

(3)  $J^0_X(X|Y) \supset J^1_X(X|Y) \supset J^2_X(X|Y) \supset \dots,$  \underline{and}  $x \notin J^p_X(X|Y)$  \underline{for all} $p \geq \dim_{\varkappa(x)} \Omega_{X|Y}(x)$ .

(4)  $x \notin J^p_X(X|Y)$  \underline{if and only if there exists an open neighbourhood}  $U$ \underline{of}  $x$  \underline{which admits a closed immersion into a}  $Y$-\underline{scheme}  $U'$  \underline{such that}  $U'$  \underline{is} \underline{smooth over}  $Y$  at  $x'$  \underline{and}  $\dim_{x'} U'_{f(x)} = p$  \underline{where}  $x'$  \underline{is the image of}  $x$  \underline{under} $U \to U'$ .

\begin{equation*}
(5) \quad \underline{I}^r_{X_1 \underset{Y}{\times} X_2}(X_1 \underset{Y}{\times} X_2 | Y) = \sum_{p+q=r} \ \pi_1^* \underline{I}^p_{X_1}(X_1|Y) \cdot \pi_2^* \underline{I}^q_{X_2}(X_2|Y)
\end{equation*}
\underline{where}  $\pi_i : X_1 \underset{Y}{\times} X_2 \to X_i$  \underline{is the projection}.

\SGAvar{\underline{Proof}:  (3) follows from 5.1(b) whereas (1) follows from 5.1(a) together with the fact that  $P^*\underline{\Omega}_{X|Y} \xrightarrow{\sim} \underline{\Omega}_{X'|Y'}$  is an isomorphism where  $P : X \underset{Y}{\times} Y' \to X$  is the projection. (2) also follows from 5.1(a) since  $\underline{O}_{X,x} \underset{\underline{O}_X}{\otimes} \underline{\Omega}_{X|Y} \xrightarrow{\sim} \underline{\Omega}_{X_x|Y}$  is an isomorphism (5) follows from 5.1(d) and the isomorphism

}{\underline{Proof}:  (3) follows from 5.1(b) whereas (1) follows from 5.1(a) together with the fact that  $P^*\underline{\Omega}_{X|Y} \xrightarrow{\sim} \underline{\Omega}_{X'|Y'}$  is an isomorphism where  $P : X \underset{Y}{\times} Y' \to X$  is the projection. (2) also follows from 5.1(a) since  $\underline{O}_{X,x} \underset{\underline{O}_X}{\otimes} \underline{\Omega}_{X|Y} \xrightarrow{\sim} \underline{\Omega}_{X_x|Y}$  is an isomorphism. (5) follows from 5.1(d) and the isomorphism

}{F7I-VI-111}\begin{equation*}
\pi_1^*\underline{\Omega}_{X_1|Y} \oplus \pi_2^*\underline{\Omega}_{X_2|Y} \xrightarrow{\sim} \underline{\Omega}_{X_1 \underset{Y}{\times} X_2|Y}
\end{equation*}
(4)  Assume the existence of a diagram

\SGAvar{\begin{tikzcd}[column sep=huge, row sep=huge]
U \arrow[d, "f"'] \& U' \arrow[l, "i"] \arrow[dl] \\
Y \&
\end{tikzcd}











}{\begin{tikzcd}[column sep=huge, row sep=huge]
U \arrow[r, "i"] \arrow[d, "f"'] \& U' \arrow[dl] \\
Y \&
\end{tikzcd}











}{F7I-VI-112}\SGApage{I}{117}
where $U'$ is smooth over $Y$ at $x' = i(x)$ and $\dim_{x'}U'_{f(x)} = p$ . Then $\dim_{\varkappa(x')}\underline{\Omega}_{U'|Y}(x') = p$ and $i^{*}\underline{\Omega}_{U'|Y} \to \underline{\Omega}_{U|Y} \to 0$ is exact, so that $\dim_{\varkappa(x)}\underline{\Omega}_{U|Y}(x) \leq p$ which implies that $x \notin J^{p}_{U}(U|Y)$ by 5.3(3). Since $J^{p}_{U}(U|Y) = J^{p}_{X}(X|Y)|U$ , this means that $x \notin J^{p}_{X}(X|Y)$ . Now assume conversely that $x \notin J^{p}_{X}(X|Y)$ .

The question being local, we may assume that $X, Y$ are both affine, say $X = \mathrm{Spec}(B)$ , $Y = \mathrm{Spec}(A)$ . Choose a presentation over $A$

\SGAvar{\begin{equation*}
0 \to J \to P \to B \to 0
\end{equation*}
where $P = A[x_1,\ldots,x_n]$ . Then $J/J^{2} \overset{d}{\to} B \underset{P}{\otimes} \Omega_{P|A} \to \Omega_{B|A} \to 0$ is exact and $B \underset{P}{\otimes} \Omega_{P|A}$ is $B$-free of rank $n$ , and therefore $I^{p}_{B}(\Omega_{B|A}) = \mathrm{Im}(\overset{n-p}{\wedge}J/J^{2} \underset{B}{\otimes} (B \underset{P}{\otimes} \overset{n-p}{\wedge}\Omega_{P|A})^{*} \to B)$ . Therefore, $x \notin J^{p}_{X}(X|Y)$ $\Leftrightarrow$ there exists $f_1,\ldots,f_{n-p}$ in $J$ such that $\left[\det\left(\dfrac{f_i}{x_j}\right)\right](x) \neq 0$ for some choice of $n-p$ variables among $x_1, x_2,\ldots,x_n$   there exists $f_1, f_2,\ldots,f_{n-p}$ in $J$ such that $\mathrm{Spec}(P/(f_1,\ldots,f_{n-p})) \to \mathrm{Spec}(A)$ is smooth at $x'$ where $x'$ is the image of $x$ under $\mathrm{Spec}(B) \to \mathrm{Spec}(P/(f_1,\ldots,f_{n-p}))$ . Thus it suffices to take $U' = \mathrm{Spec}(P/(f_1,\ldots,f_{n-p}))$ .

}{\begin{equation*}
0 \to J \to P \to B \to 0
\end{equation*}
where $P = A[x_1,\ldots,x_n]$ . Then $J/J^{2} \overset{d}{\to} B \underset{P}{\otimes} \Omega_{P|A} \to \Omega_{B|A} \to 0$ is exact and $B \underset{P}{\otimes} \Omega_{P|A}$ is $B$-free of rank $n$ , and therefore $I^{p}_{B}(\Omega_{B|A}) = \mathrm{Im}(\overset{n-p}{\wedge}J/J^{2} \underset{B}{\otimes} (B \underset{P}{\otimes} \overset{n-p}{\wedge}\Omega_{P|A})^{*} \to B)$ . Therefore, $x \notin J^{p}_{X}(X|Y)$ $\Leftrightarrow$ there exists $f_1,\ldots,f_{n-p}$ in $J$ such that $\left[\det\left(\dfrac{\partial f_i}{\partial x_j}\right)\right](x) \neq 0$ for some choice of $n-p$ variables among $x_1, x_2,\ldots,x_n$   there exists $f_1, f_2,\ldots,f_{n-p}$ in $J$ such that $\mathrm{Spec}(P/(f_1,\ldots,f_{n-p})) \to \mathrm{Spec}(A)$ is smooth at $x'$ where $x'$ is the image of $x$ under $\mathrm{Spec}(B) \to \mathrm{Spec}(P/(f_1,\ldots,f_{n-p}))$ . Thus it suffices to take $U' = \mathrm{Spec}(P/(f_1,\ldots,f_{n-p}))$ .

}{F7I-VI-113}The following is a generalization of the usual Jacobian criteria for relative situation.

\underline{Theorem} 5.4.   \underline{Let} $X \overset{f}{\to} Y$ \underline{be a morphism of essentially finite presentation. Assume that there exists a fixed integer} $n \geq 0$ \underline{with either conditions:}

(1)  $f$  \underline{is flat and} $\dim_{x}X_{f(x)} = n$ \underline{for every point} $x \in X$ .

(2)  $f$  \underline{is a relative complete-intersection with virtual dimension} $n$ (SGA6 VIII 1.9).

\underline{Then} $x \notin J^{n}_{X}(X|Y)$ \underline{if and only if} $f$ \underline{is smooth at} $x$ .

\underline{Proof:}  We have, by 5.3(4) that $x \notin J^{n}_{X}(X|Y)$ if and only if there exists a diagram
\SGApage{I}{118}
\begin{tikzcd}
U \arrow[r, hook, "i"] \arrow[d, "f"'] \& U' \arrow[dl, "f'"] \\
Y \&
\end{tikzcd}








such that $f'$ is smooth at $x'$ and $\dim_{x'}U'_{f(x)} = n$ , where $i$ is a closed immersion and $x' = i(x)$ . Thus it suffices to show that either conditions (1) or (2) entails the isomorphism of $i$ at the point $x$ . Set $A = \underline{O}_{Y,f(x)}$ , $B' = O_{U',x'}$ , $B = O_{U,x}$ and let

\begin{equation*}
0 \to I \to B' \to B \to 0
\end{equation*}
be exact. We must show that $I = 0$ . Assume (1). Since $B$ is $A$-flat,

\begin{equation*}
0 \to I/\underline{m}I \to B'/\underline{m}B' \to B/\underline{m}B \to 0
\end{equation*}
is exact where $\underline{m}$ is the maximal ideal of $A$ . By hypothesis, $\dim_{x}U_{f(x)} = \dim_{x'}U'_{f(x)}$ so that $\dim B'/\underline{m}B' = \dim B/\underline{m}B$ . However $B'/\underline{m}B'$ is a regular local ring and therefore it entails that $I/\underline{m}I = 0$ so that $I = 0$ . Now assume (2), and consider the exact sequence

\begin{equation*}
I/I^{2} \overset{d}{\to} B \underset{B'}{\otimes} \Omega_{B'|A} \to \Omega_{B|A} \to 0
\end{equation*}
Since $X$ is a relative complete-intersection over $Y$ and $U' \to Y$ is smooth at $x$ , it follows that $d : I/I^{2} \to B \underset{B'}{\otimes} \Omega_{B'|A}$ is injective at $x$ and $I/I^{2}$ and $B \underset{B'}{\otimes} \Omega_{B'|A}$ are both free at $x$ . Therefore by the definition of virtual dimension of relative complete-intersection we must have that

\begin{equation*}
n = \dim_{\varkappa(x)}\Omega_{B'|A}(x') - \dim_{\varkappa(x)}I/I^{2}(x) = n - \dim_{\varkappa(x)}I/I^{2}(x)
\end{equation*}
i.e., $I/I^{2} = 0$ i.e., $I = 0$ .

\underline{Definition} 5.5.   \underline{Let} $X \overset{f}{\to} Y$ \underline{be a morphism of essentially finite presentation.} \underline{In case when there exists a fixed integer} $n$ \underline{with the property} (1) \underline{or} (2) \underline{in} 5.4,
\SGApage{I}{119}
\underline{we set}  $\underline{I}_X(X|Y) = \underline{I}^n_X(X|Y)$  \underline{and}  $J_X(X|Y) = J^n_X(X|Y)$ .

We can extend this notion to formal schemes in an obvious manner.

\SGAvar{Thus let $\mathfrak{X}$ be a noetherian formal scheme and $\underline{E}$ a coherent $\mathfrak{X}$-module. Then we can define the ideal sheaf $\underline{I}^P_{\mathfrak{X}}(\underline{E})$ by setting

\begin{equation*}
\underline{I}^P_{\mathfrak{X}}(\underline{E})(U) \;=\; I^P_{\mathfrak{X}(U)}(\underline{E}(U))
\end{equation*}
for every affine open $U \subset \mathfrak{X}$ . If $\mathfrak{X} = \varinjlim_{v} X_v$ , then $\underline{E} = \varinjlim_{v} \underline{E}_v$ where $\underline{E}_v = \underline{E} \otimes_{\underline{O}_{\mathfrak{X}}} \underline{O}_{X_v}$ and it follows from the fact that $\mathfrak{X}$ is noetherian that

\begin{equation*}
\underline{I}^P_{\mathfrak{X}}(\underline{E}) \;=\; \varprojlim_{v} \underline{I}^P_{X_v}(\underline{E}_v)
\end{equation*}
Now let $S$ be a noetherian formal scheme and $\mathfrak{X}$ an $S$-adic formal scheme of essentially finite type. Let $\underline{I}_{\mathfrak{X}}$ be the ideal sheaf defining the closed immersion $\Delta : \mathfrak{X} \to \mathfrak{X} \underset{S}{\times} \mathfrak{X}$ (= the fibre-product in the category of formal schemes) i.e.,

}{Thus let $\mathfrak{X}$ be a noetherian formal scheme and $\underline{E}$ a coherent $\mathfrak{X}$-module. Then we can define the ideal sheaf $\underline{I}^p_{\mathfrak{X}}(\underline{E})$ by setting

\begin{equation*}
\underline{I}^p_{\mathfrak{X}}(\underline{E})(U) \;=\; I^p_{\mathfrak{X}(U)}(\underline{E}(U))
\end{equation*}
for every affine open $U \subset \mathfrak{X}$ . If $\mathfrak{X} = \varinjlim_{v} X_v$ , then $\underline{E} = \varprojlim_{v} \underline{E}_v$ where $\underline{E}_v = \underline{E} \otimes_{\underline{O}_{\mathfrak{X}}} \underline{O}_{X_v}$ and it follows from the fact that $\mathfrak{X}$ is noetherian that

\begin{equation*}
\underline{I}^p_{\mathfrak{X}}(\underline{E}) \;=\; \varprojlim_{v} \underline{I}^p_{X_v}(\underline{E}_v)
\end{equation*}
Now let $S$ be a noetherian formal scheme and $\mathfrak{X}$ an $S$-adic formal scheme of essentially finite type. Let $\underline{I}_{\mathfrak{X}}$ be the ideal sheaf defining the closed immersion $\Delta : \mathfrak{X} \to \mathfrak{X} \underset{S}{\times} \mathfrak{X}$ (= the fibre-product in the category of formal schemes) i.e.,

}{F7I-VI-114}\begin{equation*}
0 \to \underline{I}_{\mathfrak{X}} \to \underline{O}_{\mathfrak{X} \underset{S}{\times} \mathfrak{X}} \to \underline{O}_{\mathfrak{X}} \to 0
\end{equation*}
is exact. We define $\underline{\Omega}_{\mathfrak{X}|S} = \underline{I}_{\mathfrak{X}}/I^2_{\mathfrak{X}}$ . Since $\mathfrak{X}$ is an $S$-adic formal scheme of essentially finite type, it follows that $\mathfrak{X} = \varinjlim_{v} X_v$ where $X_v = \mathfrak{X} \underset{S}{\times} S_v$ , and hence $\underline{\Omega}_{\mathfrak{X}|S} \simeq \varprojlim_{v} \underline{\Omega}_{X_v|S_v}$ is a coherent $\mathfrak{X}$-module. We note that for any formal $S$-scheme $S'$ of finite type, we have a canonical isomorphism

\begin{equation*}
\underline{\Omega}_{\mathfrak{X} \underset{S}{\times} S'|S'} \;\simeq\; \underline{O}_{S'} \otimes_{\underline{O}_S} \underline{\Omega}_{\mathfrak{X}|S} \quad .
\end{equation*}
\underline{Definition} 5.6. \underline{Let} $S$ \underline{be a noetherian formal scheme and} $\mathfrak{X}$ \underline{an} $S$-\underline{adic formal scheme of essentially finite type. We define the ideal sheaf} $\underline{I}^p_{\mathfrak{X}}(\mathfrak{X}|S)$ \underline{by}

\begin{equation*}
\underline{I}^p_{\mathfrak{X}}(\mathfrak{X}|S) \;=\; \underline{I}^p_{\mathfrak{X}}(\underline{\Omega}_{\mathfrak{X}|S})
\end{equation*}
\underline{and we set} $J^p_{\mathfrak{X}}(\mathfrak{X}|S)$ \underline{the formal closed subscheme of} $\mathfrak{X}$ \underline{defined by} $\underline{I}^p_{\mathfrak{X}}(\mathfrak{X}|S)$ .

Since $\Omega_{\mathfrak{X}|S} = \varprojlim_{n} \Omega_{X_n|S_n}$ we have 


\SGApage{I}{120}

\begin{equation*}
 I^p_{\mathfrak{X}}(\mathfrak{X}|S) \;=\; \varprojlim_{n} I^p_{X_n}(X_n|S_n)
\end{equation*}
\begin{equation*}
J^p_{\mathfrak{X}}(\mathfrak{X}|S) \;=\; \varinjlim_{n} J^p_{X_n}(X_n|S_n) \quad .
\end{equation*}
$\mathfrak{X}_{(x)}$ is an $S$-adic formal scheme of essentially finite type and we have a commutative diagram

\begin{center}
\begin{tikzcd}[column sep=huge, row sep=large]
\mathfrak{X}_{(\widehat{x})} \arrow[rr, "i_{(\widehat{x})}"] \arrow[dr] \& \& \mathfrak{X} \\
\& \mathfrak{X}_{(x)} \arrow[ur, "i_{(x)}"'] \&
\end{tikzcd}
\end{center}












We note that if $x$ is an isolated point of $\mathfrak{X}$ , then $i_{(x)}$ , $i_{(\widehat{x})}$ are both isomorphisms to the connected component of $\mathfrak{X}$ containing the point $x$ .

In particular, if $|\mathfrak{X}| = \{x_1,\dots,x_m\}$ is a finite discrete space, then we have a canonical identification $\mathfrak{X} = \coprod_{1 \leq i \leq m} \mathfrak{X}_{(x_i)}$ (disjoint union).

\underline{Corollary} 5.7. \underline{Let} $S$ \underline{be a noetherian formal scheme, and} $\mathfrak{X}$ \underline{an} $S$-\underline{adic formal scheme of essentially finite type. We then have}

\SGAvar{(1) $J^p_{\mathfrak{X}}(\mathfrak{X}|S)$ \underline{commutes with base change} i.e., $J^p_{\mathfrak{X}'}(\mathfrak{X}'|S') = J^p_{\mathfrak{X}}(\mathfrak{X}|S) \underset{S}{\times} S'$ \underline{where} $\mathfrak{X} = \mathfrak{X} \underset{S}{\times} S'$ .

(2) \underline{The connected components of} $J^p_{\mathfrak{X}}(\mathfrak{X}|S)$ \underline{is in bijective correspondence with that of} $J^p_{X_0}(X_0|S_0)$ . \underline{For each point} $x \in |J^p_{X_0}(X_0|S_0)|$ \underline{we have a canonical isomorphism} $J^p_{\mathfrak{X}}(\mathfrak{X}|S)_{(x)} \simeq J^p_{\mathfrak{X}_{(x)}}(\mathfrak{X}_{(x)}|S)$ . \underline{In particular if}

}{(1) $J^p_{\mathfrak{X}}(\mathfrak{X}|S)$ \underline{commutes with base change} i.e., $J^p_{\mathfrak{X}'}(\mathfrak{X}'|S') = J^p_{\mathfrak{X}}(\mathfrak{X}|S) \underset{S}{\times} S'$ \underline{where} $\mathfrak{X}' = \mathfrak{X} \underset{S}{\times} S'$ .

(2) \underline{The connected components of} $J^p_{\mathfrak{X}}(\mathfrak{X}|S)$ \underline{are in bijective correspondence with those of} $J^p_{X_0}(X_0|S_0)$ . \underline{For each point} $x \in |J^p_{X_0}(X_0|S_0)|$ \underline{we have a canonical isomorphism} $J^p_{\mathfrak{X}}(\mathfrak{X}|S)_{(x)} \simeq J^p_{\mathfrak{X}_{(x)}}(\mathfrak{X}_{(x)}|S)$ . \underline{In particular if}

}{F7I-VI-115}\begin{equation*}
|J^p_{X_0}(X_0|S_0)| \;=\; \{x_1,\dots,x_m\}
\end{equation*}
\underline{is discrete, then}

\SGAvar{\begin{equation*}
J^p_{\mathfrak{X}}(\mathfrak{X}|S) \;-\; \coprod_{1 \leq i \leq m} J^p_{\mathfrak{X}_{(x_i)}}(\mathfrak{X}_{(x_i)}|S)
\end{equation*}
\underline{(disjoint union)}.


}{\begin{equation*}
J^p_{\mathfrak{X}}(\mathfrak{X}|S) \;=\; \coprod_{1 \leq i \leq m} J^p_{\mathfrak{X}_{(x_i)}}(\mathfrak{X}_{(x_i)}|S)
\end{equation*}
\underline{(disjoint union)}.


}{F7I-VI-116}\SGApage{I}{121}

(3) \underline{Let} $S$ \underline{be an ordinary noetherian scheme and} $X$ \underline{an} $S$-\underline{scheme of essentially finite type. Given a closed subscheme} $S_0 \subset S$ , \underline{let} $\widehat{S}$ \underline{be the formal completion of} $S$ \underline{along} $S_0$ .

\underline{If we set}

\begin{equation*}
\widehat{X} \;=\; X \underset{S}{\times} \widehat{S} \;=\; \underline{\text{the formal completion of}}\;\; X \;\; \underline{\text{along}} \;\; X_0 \;=\; X \underset{S}{\times} S_0 \;\;,\;\; \underline{\text{then}}
\end{equation*}
\begin{equation*}
J^{p}_{\widehat{X}}(\widehat{X}|\widehat{S}) \;=\; J^{p}_{X}(X|S) \underset{S}{\times} \widehat{S} \;=\; \underline{\text{the formal completion of}} \;\; J^{p}_{X}(X|S) \;\; \underline{\text{along}} \;\; J^{p}_{X_0}(X_0|S_0) \;\;.
\end{equation*}
(4) \underline{If there exists a fixed integer} $n \geq 0$ \underline{with the property} 5.4(1) \underline{or} (2) \underline{for} $X_\nu \to S_\nu$ \underline{for all} $\nu$ , \underline{then} $x \notin J^{n}_{\mathfrak{X}}(\mathfrak{X}|S)$ \underline{if and only if} $\mathfrak{X}$ \underline{is formally smooth over} $S$ \underline{at} $x$ .

\underline{Proof}: (1) follows from $\underline{\Omega}_{\mathfrak{X}'|S'} \simeq \underline{O}_{S'} \widehat{\otimes}_{\underline{O}_S} \underline{\Omega}_{\mathfrak{X}|S}$ whereas (2) follows from

\SGAvar{\begin{equation*}
\underline{O}_{\mathfrak{X}_{(x)}} \;\widehat{\otimes}_{\underline{O}_{\mathfrak{X}}}\; \underline{\Omega}_{\mathfrak{X}|S} \;\simeq\; \underline{\Omega}_{\mathfrak{X}_{(x)}|S} \;\;.
\end{equation*}
\begin{equation*}
(3)\quad J^{p}_{\widehat{X}}(\widehat{X}|\widehat{S}) \;=\; \varinjlim_{\nu} \; J^{p}_{X \underset{S}{\times} S_\nu}\!\left(X \underset{S}{\times} S_\nu \,\middle|\, S_\nu\right) \;=\; \varinjlim_{\nu} \; J^{p}_{X}(X|S) \underset{S}{\times} S_\nu \;=\; J^{p}_{X}(X|S) \underset{S}{\times} \widehat{S} \;\;.
\end{equation*}
\begin{multline*}
(4)\quad \text{Since } \mathfrak{X} \text{ is an S-adic formal scheme, } \mathfrak{X} \to S \text{ is formally smooth at}\\ x \Leftrightarrow X_\nu \to S_\nu \text{ is formally smooth at } x \text{ for all } \nu \;\Leftrightarrow\; x \notin J^{n}_{X_\nu}(X_\nu|S_\nu) \text{ for all } \nu\\ \Leftrightarrow\; x \notin \varinjlim_{\nu} J^{n}_{X_\nu}(X_\nu|S_\nu) \;=\; J^{n}_{\mathfrak{X}}(\mathfrak{X}|S) \;\;.
\end{multline*}
\underline{Remark} 5.8. Let $\mathfrak{X}$ be an $S$-adic formal scheme of essentially finite type. Given a point $x \in \mathfrak{X}$ we also consider the $S$-adic formal scheme $\mathfrak{X}_{(\hat{X})} = \varinjlim_{\nu} \mathrm{Spec}\big(\widehat{\underline{O}_{X_\nu,x}}\big)$ . We thus have the natural morphisms $\mathfrak{X}_{(\hat{X})} \to \mathfrak{X}_{(x)} \to \mathfrak{X}$ . Now $\widehat{\underline{O}_{X_0,x}}$ is in general not essentially finite type over $S_{(y)}$ where $y$ = the image of $x$ under $\mathfrak{X} \to S$ . For simplicity, let us set $A_\nu = \underline{O}_{X_\nu,x}$ , $R_\nu = \underline{O}_{S_\nu,y}$ . Now the natural map $\Omega_{A_\nu|R_\nu} \to \Omega_{\widehat{A}_\nu|R_\nu}$ induces an isomorphism $\widehat{\Omega_{A_\nu|R_\nu}} \to \widehat{\Omega_{A_\nu|R_\nu}}$ , and therefore $\varprojlim_{\nu} \widehat{\Omega_{A_\nu|R_\nu}} \;\xrightarrow{\;\sim\;}\; \lim_{\nu} \widehat{\Omega_{\widehat{A}_\nu|R_\nu}}$ where $\widehat{\phantom{xx}}$ stands for the separable completion of a module with respect to the topology induced from the ring. Therefore if we define $\widehat{\Omega_{\mathfrak{X}_{(\hat{x})}|S}} = \varprojlim_{\nu} \widehat{\Omega_{\widehat{\underline{O}}_{X_\nu,x}|\underline{O}_{S_\nu}}}$ , then $\widehat{\Omega_{\mathfrak{X}_{(\hat{x})}|S}}$ is a coherent $\mathfrak{X}_{(\hat{x})}$-module. Thus if
\SGApage{I}{122}
we set $J^{p}_{\mathfrak{X}_{(\hat{x})}}(\mathfrak{X}_{(\hat{x})}|S)$ to be the formal closed subscheme of $\mathfrak{X}_{(\hat{x})}$ defined by the ideal sheaf $I^{p}_{\mathfrak{X}_{(\hat{x})}}(\mathfrak{X}_{(\hat{x})}|S) = I^{p}_{\mathfrak{X}_{(\hat{x})}}(\Omega_{\mathfrak{X}_{(\hat{x})}|S})$ , then we have

}{\begin{equation*}
\underline{O}_{\mathfrak{X}_{(x)}} \;\widehat{\otimes}_{\underline{O}_{\mathfrak{X}}}\; \underline{\Omega}_{\mathfrak{X}|S} \;\simeq\; \underline{\Omega}_{\mathfrak{X}_{(x)}|S} \;\;.
\end{equation*}
\begin{equation*}
(3)\quad J^{p}_{\widehat{X}}(\widehat{X}|\widehat{S}) \;=\; \varinjlim_{\nu} \; J^{p}_{X \underset{S}{\times} S_\nu}\!\left(X \underset{S}{\times} S_\nu \,\middle|\, S_\nu\right) \;=\; \varinjlim_{\nu} \; J^{p}_{X}(X|S) \underset{S}{\times} S_\nu \;=\; J^{p}_{X}(X|S) \underset{S}{\times} \widehat{S} \;\;.
\end{equation*}
\begin{multline*}
(4)\quad \text{Since } \mathfrak{X} \text{ is an S-adic formal scheme, } \mathfrak{X} \to S \text{ is formally smooth at}\\ x \Leftrightarrow X_\nu \to S_\nu \text{ is formally smooth at } x \text{ for all } \nu \;\Leftrightarrow\; x \notin J^{n}_{X_\nu}(X_\nu|S_\nu) \text{ for all } \nu\\ \Leftrightarrow\; x \notin \varinjlim_{\nu} J^{n}_{X_\nu}(X_\nu|S_\nu) \;=\; J^{n}_{\mathfrak{X}}(\mathfrak{X}|S) \;\;.
\end{multline*}
\underline{Remark} 5.8. Let $\mathfrak{X}$ be an $S$-adic formal scheme of essentially finite type. Given a point $x \in \mathfrak{X}$ we also consider the $S$-adic formal scheme $\mathfrak{X}_{(\hat{x})} = \varinjlim_{\nu} \mathrm{Spec}\big(\widehat{\underline{O}_{X_\nu,x}}\big)$ . We thus have the natural morphisms $\mathfrak{X}_{(\hat{x})} \to \mathfrak{X}_{(x)} \to \mathfrak{X}$ . Now $\widehat{\underline{O}_{X_0,x}}$ is in general not essentially finite type over $S_{(y)}$ where $y$ = the image of $x$ under $\mathfrak{X} \to S$ . For simplicity, let us set $A_\nu = \underline{O}_{X_\nu,x}$ , $R_\nu = \underline{O}_{S_\nu,y}$ . Now the natural map $\Omega_{A_\nu|R_\nu} \to \Omega_{\widehat{A}_\nu|R_\nu}$ induces an isomorphism $\widehat{\Omega_{A_\nu|R_\nu}} \to \widehat{\Omega_{\widehat{A}_\nu|R_\nu}}$ , and therefore $\varprojlim_{\nu} \widehat{\Omega_{A_\nu|R_\nu}} \;\xrightarrow{\;\sim\;}\; \varprojlim_{\nu} \widehat{\Omega_{\widehat{A}_\nu|R_\nu}}$ where $\widehat{\phantom{xx}}$ stands for the separated completion of a module with respect to the topology induced from the ring. Therefore if we define $\widehat{\Omega_{\mathfrak{X}_{(\hat{x})}|S}} = \varprojlim_{\nu} \widehat{\Omega_{\widehat{\underline{O}}_{X_\nu,x}|\underline{O}_{S_\nu,y}}}$ , then $\widehat{\Omega_{\mathfrak{X}_{(\hat{x})}|S}}$ is a coherent $\mathfrak{X}_{(\hat{x})}$-module. Thus if
\SGApage{I}{122}
we set $J^{p}_{\mathfrak{X}_{(\hat{x})}}(\mathfrak{X}_{(\hat{x})}|S)$ to be the formal closed subscheme of $\mathfrak{X}_{(\hat{x})}$ defined by the ideal sheaf $\underline{I}^{p}_{\mathfrak{X}_{(\hat{x})}}(\mathfrak{X}_{(\hat{x})}|S) = \underline{I}^{p}_{\mathfrak{X}_{(\hat{x})}}(\widehat{\Omega}_{\mathfrak{X}_{(\hat{x})}|S})$ , then we have

}{F7I-VI-117}\begin{equation*}
J^{p}_{\mathfrak{X}_{(x)}}(\mathfrak{X}_{(x)}|S)_{(\hat{x})} \;\simeq\; J^{p}_{\mathfrak{X}_{(\hat{x})}}(\mathfrak{X}_{(\hat{x})}|S) \;\;.
\end{equation*}
We note that if $x$ is an isolated point of $J^{p}_{\mathfrak{X}}(\mathfrak{X}|S)$ , then $J^{p}_{\mathfrak{X}_{(\hat{x})}}(\mathfrak{X}_{(\hat{x})}|S) \to J^{p}_{\mathfrak{X}_{(x)}}(\mathfrak{X}_{(x)}|S)$ is an isomorphism.

\SGAvar{To consider the "image" of $J^{p}_{\mathfrak{X}}(\mathfrak{X}|S)$ under the morphism $\mathfrak{X} \xrightarrow{\,f\,} S$ , one may simply take the formal subscheme of $S$ defined by the ideal sheaves $\mathrm{Ker}(\underline{O}_S \to \underline{O}_{\mathfrak{X}}/I^{p}_{\mathfrak{X}}(\mathfrak{X}|S))$ . However, it is more convenient for our purpose to use the Fitting ideal again.

}{To consider the "image" of $J^{p}_{\mathfrak{X}}(\mathfrak{X}|S)$ under the morphism $\mathfrak{X} \xrightarrow{\,f\,} S$ , one may simply take the formal subscheme of $S$ defined by the ideal sheaf $\mathrm{Ker}(\underline{O}_S \to \underline{O}_{\mathfrak{X}}/\underline{I}^{p}_{\mathfrak{X}}(\mathfrak{X}|S))$ . However, it is more convenient for our purpose to use the Fitting ideal again.

}{F7I-VI-118}\underline{Definition} 5.9. \underline{Let} $S$ \underline{be a noetherian formal scheme and} $f : \mathfrak{X} \to S$ \underline{an} $S$-\underline{adic formal scheme of essentially finite type. Assume that} $J^{p}_{\mathfrak{X}}(\mathfrak{X}|S)$ \underline{is proper over} $S$ , \underline{so that} $f_{*}\underline{O}_{J^{p}_{\mathfrak{X}}(\mathfrak{X}|S)}$ \underline{is a coherent} $\underline{O}_S$-\underline{module. We then define the ideal sheaf} $\underline{I}^{p}_{S}(\mathfrak{X}|S)$ \underline{of} $\underline{O}_S$ \underline{by}

\begin{equation*}
\underline{I}^{p}_{S}(\mathfrak{X}|S) \;=\; \underline{I}^{0}_{S}\big(f_{*}\underline{O}_{J^{p}_{\mathfrak{X}}(\mathfrak{X}|S)}\big)
\end{equation*}
\underline{and set} $J^{p}_{S}(\mathfrak{X}|S)$ \underline{to be the closed formal subscheme of} $S$ \underline{defined by the ideal sheaf} $\underline{I}^{p}_{S}(\mathfrak{X}|S)$ . \underline{In case when there exists a fixed integer} $n$ \underline{having the property} (1) \underline{or} (2) \underline{in} 5.4 \underline{for} $X_\nu \to S_\nu$ \underline{for all} $\nu$ , \underline{we set}

\begin{equation*}
\underline{I}_{\mathfrak{X}}(\mathfrak{X}|S) \;=\; \underline{I}^{n}_{\mathfrak{X}}(\mathfrak{X}|S) \;\;,\qquad J_{\mathfrak{X}}(\mathfrak{X}|S) \;=\; J^{n}_{\mathfrak{X}}(\mathfrak{X}|S)
\end{equation*}
\begin{equation*}
\underline{I}_{S}(\mathfrak{X}|S) \;=\; \underline{I}^{n}_{S}(\mathfrak{X}|S) \;\;,\qquad J_{S}(\mathfrak{X}|S) \;=\; J^{n}_{S}(\mathfrak{X}|S)
\end{equation*}
\underline{Corollary} 5.10. \underline{Let} $S$ \underline{be a noetherian formal scheme and} $\mathfrak{X} \xrightarrow{\,f\,} S$ \underline{an} $S$-\underline{adic formal scheme, essentially of finite type such that} $J^{p}_{\mathfrak{X}}(\mathfrak{X}|S)$ \underline{is proper over} $S$ . \underline{Then}

\begin{equation*}
(1)\quad \big|J^{p}_{S}(\mathfrak{X}|S)\big| \;=\; \underline{\text{the image of}} \;\; \big|J^{p}_{\mathfrak{X}}(\mathfrak{X}|S)\big| \;\; \underline{\text{under}} \;\; f : \mathfrak{X} \to S \;\;.
\end{equation*}

\SGApage{I}{123}

\SGAvar{(2)  $J^p_S(\mathfrak{X}|S)$  \underline{commute with base change}  i.e.,  $J^p_{S'}(\mathfrak{X}'|S') = J^p_S(\mathfrak{X}|S) \times_S S'$  .

}{(2)  $J^p_S(\mathfrak{X}|S)$  \underline{commutes with base change}  i.e.,  $J^p_{S'}(\mathfrak{X}'|S') = J^p_S(\mathfrak{X}|S) \times_S S'$  .

}{F7I-VI-119}(3)  If  $J^p_{\mathfrak{X}}(\mathfrak{X}|S) = \coprod_{1\leq i\leq m} \mathfrak{Y}^{(i)}$  \underline{(disjoint union), then}

$\underline{I}^p_S(\mathfrak{X}|S) = \prod_{1\leq i\leq m} \underline{I}^0_S(f_*\underline{O}_{\mathfrak{Y}(i)})$  . \underline{In particular if}

\begin{equation*}
|J^p_{X_0}(X_0|S_0)| = \{x_1,\ldots,x_m\}
\end{equation*}
\underline{is discrete, then}

\begin{equation*}
\underline{I}^p_S(\mathfrak{X}|S) = \prod_{1\leq i\leq m} \underline{I}^p_S(\mathfrak{X}_{(x_i)}|S) = \prod_{1\leq i\leq m} \underline{I}^p_S(\mathfrak{X}_{(\hat{x}_i)}|S)
\end{equation*}
(4)  \underline{Let}  $S$  \underline{be an ordinary noetherian scheme and}  $X$  \underline{an} $S$-\underline{scheme of essentially finite type. Given a closed subscheme}  $S_0 \to S$  , \underline{let}  $\hat{S}$  \underline{be the formal completion of}  $S$  \underline{along}  $S_0$  . \underline{If we set}  $\hat{X} = X \times_S \hat{S}$ $=$ \underline{the formal completion of}  $X$  \underline{along}  $X_0 = X \times_S S_0$  , \underline{then}  $J^p_{\hat{S}}(\hat{X}|\hat{S}) = J^p_S(X|S) \times_S \hat{S}$ $=$ \underline{the formal completion of}  $J^p_S(X|S)$  \underline{along}  $J^p_{S_0}(X_0|S_0)$  .

\underline{Proof}:  Supp $f_*\underline{O}_{J^p_{\mathfrak{X}}(\mathfrak{X}|S)}$ $=$ $|J^p_S(\mathfrak{X}|S)|$  by 5.1(c), and hence (1) follows. (2) and (4) are clear from 5.1(a), and (3) follows from 5.1(d), and the isomorphism $J^p_{\mathfrak{X}_{(\hat{x})}}(\mathfrak{X}_{(\hat{x})}|S) \xrightarrow{\sim} J^p_{\mathfrak{X}_{(x)}}(\mathfrak{X}_{(x)}|S)$  for an isolated point  $x$  in  $J^p_{\mathfrak{X}}(\mathfrak{X}|S)$  .

\SGAvar{\underline{Remark} 5.11.  Our definition of Jacobian ideal, adopted from the classical one over base field, commutes with base change and gives a smoothness criteria under a mild condition of 5.4. However one can define, probably in many different ways, a canonical subscheme which gives the non-smooth locus in a general setting. We shall outline one here:  Given an $A$-algebra  $B$  , define  $I_{B|A} = \bigcap_E \mathrm{Ann}\ D^1(B|A,E)$ where  $E$  runs through all  $B$-modules of finite type. This ideal  $I_{B|A}$  commutes with localizations by 3.13(c), and therefore, given a morphism  $X \xrightarrow{\ f\ } Y$  , we obtain the ideal sheaf  $\underline{I}_{X|Y}$  . If  $f$  is a morphism of essentially finite type, then  $f$  is smooth at  $x$  if and only if  $x \notin \mathrm{Supp}\ \underline{O}_X/\underline{O}_{X|Y}$  by 3.13(a). This notion is in an obvious way carried over to adic formal schemes of essentially
\SGApage{I}{124}
finite type and thus it is equally suitable for deformation theory.


}{\underline{Remark} 5.11.  Our definition of Jacobian ideal, adopted from the classical one over base field, commutes with base change and gives a smoothness criteria under a mild condition of 5.4. However one can define, probably in many different ways, a canonical subscheme which gives the non-smooth locus in a general setting. We shall outline one here:  Given an $A$-algebra  $B$  , define  $I_{B|A} = \bigcap_E \mathrm{Ann}\ D^1(B|A,E)$ where  $E$  runs through all  $B$-modules of finite type. This ideal  $I_{B|A}$  commutes with localizations by 3.13(c), and therefore, given a morphism  $X \xrightarrow{\ f\ } Y$  , we obtain the ideal sheaf  $\underline{I}_{X|Y}$  . If  $f$  is a morphism of essentially finite type, then  $f$  is smooth at  $x$  if and only if  $x \notin \mathrm{Supp}\ \underline{O}_X/\underline{I}_{X|Y}$  by 3.13(a). This notion is in an obvious way carried over to adic formal schemes of essentially
\SGApage{I}{124}
finite type and thus it is equally suitable for deformation theory.


}{F7I-VI-120}\SGAsectiontoc{6.  Non-degenerate quadratic singularities}{6}
6.  \underline{Non-degenerate quadratic singularities}

The objective in this section is to investigate the deformations of a $k$-scheme having non-degenerate quadratic singularities only. In view of 4.13, it becomes purely a local problem in an appropriate situation.

We go back to the notations of  §4. Thus  $\Lambda$  is a fixed complete noetherian local ring with the residue field  $k$  , and  $\underline{C}_\Lambda$  stands for the category of local artinian $\Lambda$-algebras with the same residue field  $k$  .

\SGAvar{Let  $X$  be a $k$-scheme of finite type, and let  $x$  be a closed point of $X$  . We recall that the point  $x$  is said to be a non-degenerate quadratic singularity rational over  $k$  if  $\hat{O}_{X,x}$ $-$ $k$ $[[x_1,\ldots,x_n]]/(f)$  with $(\partial f/\partial x_1,\ldots,\partial f/\partial x_n) = (x_1,\ldots,x_n)$  . One  can always assume  $f$  to be a quadratic form. Indeed we have

}{Let  $X$  be a $k$-scheme of finite type, and let  $x$  be a closed point of $X$  . We recall that the point  $x$  is said to be a non-degenerate quadratic singularity rational over  $k$  if  $\hat{O}_{X,x}$ $\simeq$ $k$ $[[x_1,\ldots,x_n]]/(f)$  with $(\partial f/\partial x_1,\ldots,\partial f/\partial x_n) = (x_1,\ldots,x_n)$  . One  can always assume  $f$  to be a quadratic form. Indeed we have

}{F7I-VI-121}\underline{Lemma} 6.1.  \underline{Let}  $f \in k[[x_1,\ldots,x_n]]$  . \underline{If}  $(\partial f/\partial x_1,\ldots,\partial f/\partial x_n) = M$  $(= (x_1,\ldots,x_n))$, \underline{then one can find}  $x_1',\ldots,x_n'$  \underline{with}  $x_i = x_i'$  (mod $M^2$)  \underline{such that} $f = q(x_1',\ldots,x_n')$  \underline{where}  $q$  \underline{is a quadratic form.}

\SGAvar{\underline{Proof}:  We do this by a successive approximation. Write  $f = f_2 + h_3$  where  $f_2$ is a quadratic form and  $h_3 \in M^3$  . Since  $(\partial f/\partial x_1,\ldots,\partial f/\partial x_n) = M$  , there exists $\epsilon_1^{(1)},\ldots,\epsilon_n^{(1)}$  in  $M^2$  such that  $\Sigma\ \epsilon_i^{(1)}\ \partial f/\partial x_i \equiv h_3$  (mod $M^4$)  . If we set $X_i^{(1)} = X_i - \epsilon_i^{(1)}$  , then  $f(x_1^{(1)},\ldots,x_n^{(1)}) =$ $f(x_1,\ldots,x_n) - \Sigma\ \epsilon_i^{(1)}\ \partial f/\partial x_i + \ldots \equiv f_2(x_1,\ldots,x_n)$  (mod $M^4$)  i.e., $f(x_1^{(1)},\ldots,x_n^{(1)}) \equiv f_2(x_1,\ldots,x_n) + h_4$  with  $h_4 \in M^4$  . Now $(\partial f/\partial x_1^{(1)},\ldots,\partial f/\partial x_n^{(1)}) = (\partial f/\partial x_1,\ldots,\partial f/\partial x_n) = M$  and therefore we can choose
\SGApage{I}{125}
$\varepsilon_1^{(2)},\dots,\varepsilon_n^{(2)}$  in  $M^3$  such that  $\Sigma\ \varepsilon_i^{(2)}\ \partial f/\partial x_u^{(1)} \equiv h_4\ (\mathrm{mod}\ M^5)$ . Then $f(x_1^{(1)} - \varepsilon_1^{(2)},\dots,x_n^{(1)} - \varepsilon_n^{(2)}) = f(x_1^{(1)},\dots,x_n^{(1)}) - \Sigma\ \varepsilon_i^{(2)}\partial f/\partial x_i^{(1)} + \dots \equiv f_2(x_1,\dots,x_n)\ (\mathrm{mod}\ M^5)$ . Continuing this process we can find $\varepsilon_i^{(1)},\ \varepsilon_i^{(2)},\ \varepsilon_i^{(3)},\dots$  $(i = 1,2,\dots,n)$  with  $\varepsilon_i^{(v)} \in M^{v+1}$  such that if we set $x_i^{(v)} = x_i - (\varepsilon_i^{(1)} + \dots + \varepsilon_i^{(v)})$  then  $f(x_1^{(v)},\dots,x_n^{(v)}) \equiv f_2(x_1,\dots,x_n)$ $(\mathrm{mod}\ M^{v+3})$ . Thus if we set  $\varepsilon_i = \sum_{v} \varepsilon_i^{(v)}$  then  $f(x_1-\varepsilon_1,\dots,x_n-\varepsilon_n) =$ $= f_2(x_1,\dots,x_n)$  i.e.,  $f(x_1,\dots,x_n) = f_2(x_1+\varepsilon_1,\dots,x_n+\varepsilon_n)$ .

}{\underline{Proof}:  We do this by a successive approximation. Write  $f = f_2 + h_3$  where  $f_2$ is a quadratic form and  $h_3 \in M^3$  . Since  $(\partial f/\partial x_1,\ldots,\partial f/\partial x_n) = M$  , there exists $\epsilon_1^{(1)},\ldots,\epsilon_n^{(1)}$  in  $M^2$  such that  $\Sigma\ \epsilon_i^{(1)}\ \partial f/\partial x_i \equiv h_3$  (mod $M^4$)  . If we set $x_i^{(1)} = x_i - \epsilon_i^{(1)}$  , then  $f(x_1^{(1)},\ldots,x_n^{(1)}) =$ $f(x_1,\ldots,x_n) - \Sigma\ \epsilon_i^{(1)}\ \partial f/\partial x_i + \ldots \equiv f_2(x_1,\ldots,x_n)$  (mod $M^4$)  i.e., $f(x_1^{(1)},\ldots,x_n^{(1)}) \equiv f_2(x_1,\ldots,x_n) + h_4$  with  $h_4 \in M^4$  . Now $(\partial f/\partial x_1^{(1)},\ldots,\partial f/\partial x_n^{(1)}) = (\partial f/\partial x_1,\ldots,\partial f/\partial x_n) = M$  and therefore we can choose
\SGApage{I}{125}
$\varepsilon_1^{(2)},\dots,\varepsilon_n^{(2)}$  in  $M^3$  such that  $\Sigma\ \varepsilon_i^{(2)}\ \partial f/\partial x_i^{(1)} \equiv h_4\ (\mathrm{mod}\ M^5)$ . Then $f(x_1^{(1)} - \varepsilon_1^{(2)},\dots,x_n^{(1)} - \varepsilon_n^{(2)}) = f(x_1^{(1)},\dots,x_n^{(1)}) - \Sigma\ \varepsilon_i^{(2)}\partial f/\partial x_i^{(1)} + \dots \equiv f_2(x_1,\dots,x_n)\ (\mathrm{mod}\ M^5)$ . Continuing this process we can find $\varepsilon_i^{(1)},\ \varepsilon_i^{(2)},\ \varepsilon_i^{(3)},\dots$  $(i = 1,2,\dots,n)$  with  $\varepsilon_i^{(v)} \in M^{v+1}$  such that if we set $x_i^{(v)} = x_i - (\varepsilon_i^{(1)} + \dots + \varepsilon_i^{(v)})$  then  $f(x_1^{(v)},\dots,x_n^{(v)}) \equiv f_2(x_1,\dots,x_n)$ $(\mathrm{mod}\ M^{v+3})$ . Thus if we set  $\varepsilon_i = \sum_{v} \varepsilon_i^{(v)}$  then  $f(x_1-\varepsilon_1,\dots,x_n-\varepsilon_n) =$ $= f_2(x_1,\dots,x_n)$  i.e.,  $f(x_1,\dots,x_n) = f_2(x_1+\varepsilon_1,\dots,x_n+\varepsilon_n)$ .

}{F7I-VI-122}\underline{Lemma} 6.2.   \underline{Let}  $A_0$  \underline{be a local} $k$-\underline{algebra, essentially of finite type such that}

\begin{equation*}
\hat{A}_0 \simeq k[[x_1,\dots,x_n]]/(f)
\end{equation*}
\underline{where}  $f$  \underline{is a non-degenerate quadratic form, and let}  $(R,\ \mathfrak{X})$  \underline{be a formal versal deformation of}  $X_0 = \mathrm{Spec}(A_0)$  . \underline{Then}

(1)  $J_{\mathfrak{X}}(\mathfrak{X}|R) \rightarrow J_R(\mathfrak{X}|R) \rightarrow \mathrm{Spec}(\Lambda)$  \underline{are both isomorphisms.}

\SGAvar{(2)  $J_R(\mathfrak{X}|R)$  \underline{is a principal ideal in}  $R$  \underline{generated by a free formal generator of}  $R$  \underline{over}  $\Lambda$  i.e.  $I_R(\mathfrak{X}|R) = (t)$  \underline{and}  $R = \Lambda[[t]]$  .

}{(2)  $I_R(\mathfrak{X}|R)$  \underline{is a principal ideal in}  $R$  \underline{generated by a free formal generator of}  $R$  \underline{over}  $\Lambda$  i.e.  $I_R(\mathfrak{X}|R) = (t)$  \underline{and}  $R = \Lambda[[t]]$  .

}{F7I-VI-123}(3)  $\mathfrak{X} \rightarrow \mathrm{Spec}(\Lambda)$  \underline{is formally smooth.}

\SGAvar{\underline{Proof}:  We set  $\mathfrak{X} = \mathrm{Spf}(A)$  . By 4.10(b),  $(R,\ \mathrm{Spf}(\hat{A}))$  is a formal versal deformation of  $\hat{X}_0 = \mathrm{Spec}(\hat{A}_0)$  where  $\hat{A}_0 = \varprojlim_{v} A \otimes_R R_v$  . Now  $\hat{A}_0 \simeq k[[x_1,\dots,x_n]]/(f)$  and  $(\partial f/\partial x_1,\dots,\partial f/\partial x_n) = (x_1,\dots,x_n)$  in $k[[x_1,\dots,x_n]]$  . It follows readily that  $[\underline{M}_{\hat{X}_0}(k[\varepsilon])] \simeq k$  and, by 4.14, we get

}{\underline{Proof}:  We set  $\mathfrak{X} = \mathrm{Spf}(A)$  . By 4.10(b),  $(R,\ \mathrm{Spf}(\hat{A}))$  is a formal versal deformation of  $\hat{X}_0 = \mathrm{Spec}(\hat{A}_0)$  where  $\hat{A} = \varprojlim_{v} A \otimes_R R_v$  . Now  $\hat{A}_0 \simeq k[[x_1,\dots,x_n]]/(f)$  and  $(\partial f/\partial x_1,\dots,\partial f/\partial x_n) = (x_1,\dots,x_n)$  in $k[[x_1,\dots,x_n]]$  . It follows readily that  $[\underline{M}_{\hat{X}_0}(k[\varepsilon])] \simeq k$  and, by 4.14, we get

}{F7I-VI-124}\begin{equation*}
R \simeq \Lambda[[t]] \quad\text{and}\quad \hat{A} \simeq R[[x_1,\dots,x_n]]/(F - t)
\end{equation*}
where  $F$  is a quadratic polynomial in  $\Lambda[x_1,\dots,x_n]$  $(\subset R[[x_1,\dots,x_n]])$ obtained from  $f$  by lifting the coefficients to  $\Lambda$  . Now

\SGAvar{\begin{equation*}
0 \rightarrow \hat{A} \ \overset{df}{\relbar\joinrel\relbar\joinrel\relbar}\ \hat{A}dx_1 \oplus \dots \oplus \hat{A}dx_n \ \relbar\joinrel\relbar\joinrel\relbar\joinrel\relbar\joinrel\relbar\ \hat{\Omega}_{\hat{A}|R} \rightarrow 0 
\end{equation*}

}{\begin{equation*}
0 \rightarrow \hat{A} \ \overset{dF}{\longrightarrow}\ \hat{A}dx_1 \oplus \dots \oplus \hat{A}dx_n \ \xrightarrow{\hspace{3em}}\ \hat{\Omega}_{\hat{A}|R} \rightarrow 0 
\end{equation*}

}{F7I-VI-125}\SGApage{I}{126}

\SGAvar{ is exact where  $df : \hat{A} \rightarrow \hat{A}dx_1 \oplus \dots \oplus \hat{A}dx_n$  is given by $1 \rightarrow (\partial f/\partial x_1)dx_1 + \dots + (\partial f/\partial x_n)dx_n$  , and therefore $I_{\hat{A}}(\hat{A}|R) = (\partial F/\partial x_1,\dots,\partial F/\partial x_n) = (x_1,\dots,x_n)$  , and hence $\hat{A}/I_{\hat{A}}(\hat{A}|R) = \hat{A}/(x_1,\dots,x_n) = R[[x_1,\dots,x_n]]/(x_1,\dots,x_n,F - t) \simeq \Lambda$ i.e., the canonical map  $\Lambda \rightarrow \hat{A}/I_{\hat{A}}(\hat{A}|R)$  is an isomorphism. It follows that $\hat{A} \rightarrow A/I_A(A|R)$  is an isomorphism, and consequently  $\Lambda \rightarrow R/I_R(A|R)$  is also an isomorphism. This proves (1) and (2). On the other hand, (3) follows from $\hat{A} \simeq R[[x_1,\dots,x_n]]/(F - t) \simeq \Lambda[[t,x_1,\dots,x_n]]/(F - t) \simeq \Lambda[[x_1,\dots,x_n]]$   as $\Lambda$-algebras.

}{ is exact where  $dF : \hat{A} \rightarrow \hat{A}dx_1 \oplus \dots \oplus \hat{A}dx_n$  is given by $1 \rightarrow (\partial F/\partial x_1)dx_1 + \dots + (\partial F/\partial x_n)dx_n$  , and therefore $I_{\hat{A}}(\hat{A}|R) = (\partial F/\partial x_1,\dots,\partial F/\partial x_n) = (x_1,\dots,x_n)$  , and hence $\hat{A}/I_{\hat{A}}(\hat{A}|R) = \hat{A}/(x_1,\dots,x_n) = R[[x_1,\dots,x_n]]/(x_1,\dots,x_n,F - t) \simeq \Lambda$ i.e., the canonical map  $\Lambda \rightarrow \hat{A}/I_{\hat{A}}(\hat{A}|R)$  is an isomorphism. It follows that $\hat{A}/I_{\hat{A}}(\hat{A}|R) \rightarrow A/I_A(A|R)$  is an isomorphism, and consequently  $\Lambda \rightarrow R/I_R(A|R)$  is also an isomorphism. This proves (1) and (2). On the other hand, (3) follows from $\hat{A} \simeq R[[x_1,\dots,x_n]]/(F - t) \simeq \Lambda[[t,x_1,\dots,x_n]]/(F - t) \simeq \Lambda[[x_1,\dots,x_n]]$   as $\Lambda$-algebras.

}{F7I-VI-126}\underline{Corollary} 6.3.  \underline{Let}  $X_0$  \underline{be a k-scheme of finite type, and let}  $x$  \underline{in}  $X_0$  \underline{be a non-degenerate quadratic singularity rational over}  $k$  .

(1)  \underline{Let}  $(R,\ \mathfrak{X})$ ,  $(R,\ \mathfrak{Y})$  \underline{be any two deformations of}  $X_0$  \underline{where} $R \in \mathrm{ob}(\hat{\underline{C}}_{\Lambda})$  . \underline{Assume that}  $k = k^2$  . \underline{Then}  $\mathfrak{X}_{(\widehat{x})} \simeq \mathfrak{Y}_{(\widehat{x})}$  \underline{as} $R$-\underline{adic formal schemes if and only if}  $I_R(\mathfrak{X}_{(\mathrm{x})}|R) = I_R(\mathfrak{Y}_{(\mathrm{x})}|R)$  .

(2)  \underline{Let}  $R \in \mathrm{ob}(\hat{\underline{C}}_{\Lambda})$  \underline{be a regular local ring, and let}  $(R,\ \mathfrak{X})$  \underline{be a deformation of}  $X_0$  . \underline{Then}  $O_{\mathfrak{X},x}$  \underline{is a regular local ring if and only if} $I_R(\mathfrak{X}_{(\mathrm{x})}|R) \not\subset \underline{m}_R^{\,2}$  i.e., \underline{if and only if}  $I_R(\mathfrak{X}_{(\mathrm{x})}|R)$  \underline{is a regular divisor in} $\mathrm{Spf}(R)$  .

\underline{Proof:}  Let  $(R_0,\ Z)$  be a formal versal deformation of  $\mathrm{Spec}(\hat{\underline{O}}_{X_0,x})$  . Since  $x$ in  $X_0$  is a non-degenerate quadratic singularity rational over  $k$  , we may assume that  $R_0 \simeq \Lambda[[t]]$  and  $Z \simeq \mathrm{Spf}(R_0[[x_1,\dots,x_n]]/(F - t))$  as in 6.2.

(1)  $(R,\ \mathfrak{X}_{(\widehat{x})})$ ,  $(R,\ \mathfrak{Y}_{(\widehat{x})})$  are deformations of  $\mathrm{Spec}(\hat{\underline{O}}_{X_0,x})$  , and therefore they are induced from  $(R_0,\ Z)$  by continuous local $\Lambda$-algebra maps $\varphi,\ \psi : R_0 \rightarrow R$  respectively, i.e., if we set  $a = \varphi(t)$  ,  $b = \psi(t)$  , then

\SGAvar{\begin{equation*}
\mathfrak{X}_{(\widehat{x})} \simeq \mathrm{Spf}(R[[X_1,\dots,X_n]]/(R - a)) \quad ,
\end{equation*}
\begin{equation*}
\mathfrak{Y}_{(\widehat{x})} \simeq \mathrm{Spf}(R[[X_1,\dots,X_n]]/(F - b))
\end{equation*}
as $R$-adic formal schemes. Assume that  $I_R(\mathfrak{X}_{(\mathrm{x})}|R) = I_R(\mathfrak{Y}_{(\mathrm{x})}|R)$  i.e.,
\SGApage{I}{127}
$(a) = (b)$ in $R$ . Then $a = ab$ for some unit $u$ in $R$ . Since $k = k^2$ and $R$ is complete, we can write $u = v^2$ for some unit $v$ in $R$ , i.e., $a = v^2b$ . Then $(F - b) = (v^2F - a)$ , so that $X_i \to vX_i$ establishes an $R$-algebra isomorphism $R[[X_1,\ldots,X_n]]/(F - a) \;\overset{\sim}{\relbar\relbar\relbar}\; R[[X_1,\ldots,X_n]]/(F - b)$ since $F$ is a quadratic form. The other direction is trivial.

}{\begin{equation*}
\mathfrak{X}_{(\widehat{x})} \simeq \mathrm{Spf}(R[[X_1,\dots,X_n]]/(F - a)) \quad ,
\end{equation*}
\begin{equation*}
\mathfrak{Y}_{(\widehat{x})} \simeq \mathrm{Spf}(R[[X_1,\dots,X_n]]/(F - b))
\end{equation*}
as $R$-adic formal schemes. Assume that  $I_R(\mathfrak{X}_{(\mathrm{x})}|R) = I_R(\mathfrak{Y}_{(\mathrm{x})}|R)$  i.e.,
\SGApage{I}{127}
$(a) = (b)$ in $R$ . Then $a = ub$ for some unit $u$ in $R$ . We construct a matrix $T\in\operatorname{GL}_n(R)$ such that $F(TX)=uF(X)$.\footnote{Editorial proof repair: the printed inference that every unit of $R$ has a square root need not hold in characteristic $2$. The following argument retains the original hypotheses.} Write $\mathfrak m=\mathfrak m_R$, and interpret polynomial congruences coefficientwise. Since $k=k^2$, choose $c\in R^\times$ whose residue satisfies $\bar c^2=\bar u$, and put $T_1=cI_n$. Then $F(T_1X)\equiv uF(X)\pmod{\mathfrak m}$. By the nondegeneracy condition in Lemma 6.2, the coefficient matrix of the linear forms $\partial F/\partial X_i$ is invertible modulo $\mathfrak m$, hence over $R$.  Suppose $T_\nu\in\operatorname{GL}_n(R)$ satisfies $F(T_\nu X)\equiv uF(X)\pmod{\mathfrak m^\nu}$, with $\nu\geq1$. The forms $(\partial F/\partial X_i)(T_\nu X)$ form an $R$-basis of the linear forms. Consequently, there are linear forms $L_i$ with coefficients in $\mathfrak m^\nu$ such that \[ E_\nu(X):=uF(X)-F(T_\nu X)  =\sum_i L_i(X)\frac{\partial F}{\partial X_i}(T_\nu X). \] Indeed, every quadratic monomial is a product of two linear forms, and expressing one factor in this basis preserves its coefficient ideal. Define $H_\nu$ by $H_\nu X=(L_i(X))_i$, and set $T_{\nu+1}=T_\nu+H_\nu$. The quadratic identity gives \[ F(T_{\nu+1}X)=F(T_\nu X)+E_\nu(X)+F(H_\nu X)  \equiv uF(X)\pmod{\mathfrak m^{\nu+1}}, \] since the coefficients of $F(H_\nu X)$ lie in $\mathfrak m^{2\nu}\subseteq\mathfrak m^{\nu+1}$. Moreover, $T_{\nu+1}\equiv T_\nu\pmod{\mathfrak m^\nu}$. Completeness and separatedness therefore give a limit $T\equiv cI_n\pmod{\mathfrak m}$ with $F(TX)=uF(X)$; in particular, $T$ is invertible. The substitution $X\mapsto TX$ sends $F-a$ to $u(F-b)$ and induces the required isomorphism \[ R[[X_1,\ldots,X_n]]/(F-a)  \xrightarrow{\ \sim\ } R[[X_1,\ldots,X_n]]/(F-b). \] The converse follows from invariance of $I_R$ under $R$-isomorphisms.

}{F7I-VI-127}(2)  There exists a continuous local $\Lambda$-algebra map $\varphi : R_0 = \Lambda[[t]] \to R$ such that $\mathfrak{X}_{(\widehat{x})} \simeq \mathrm{Spf}(R[[X_1,\ldots,X_n]]/(F - \varphi(t)))$ . We note that $I_R(\mathfrak{X}_{(x)}|R) = I_R(\mathfrak{X}_{(\widehat{x})}|R) = (\varphi(t))$ , and that $R[[X_1,\ldots,X_n]]/(F - \varphi(t)) = \hat{\underline{O}}_{\mathfrak{X},x}$ is regular if and only if $\underline{O}_{\mathfrak{X},x}$ is regular. Now $R$ is regular and hence $R[[X_1,\ldots,X_n]]$ is also a regular local ring. Thus if we set $M$ to be the maximal ideal of $R[[X_1,\ldots,X_n]]$ , then $R[[X_1,\ldots,X_n]]/(F - \varphi(t))$ is regular if and only if $F - \varphi(t)$ is not in $M^2$ , i.e. if and only if $\varphi(t) \notin \underline{m}_R^{\,2}$ , i.e., if and only if $I_R(\mathfrak{X}_{(x)}|R) \not\subset \underline{m}_R^{\,2}$ .

\underline{Theorem} 6.4.  \underline{Let} $X_0$ \underline{be a} $k$-\underline{scheme of finite type with isolated non-smooth points only, and let} $(R, \mathfrak{X})$ \underline{be a formal versal deformation of} $X_0$ . \underline{Denote by} $\{z_1,\ldots,z_m\}$ \underline{the set of non-smooth points in} $X_0$ . \underline{Assume that} $H^2(X_0, \underline{D}^0_{X_0}) = 0$ , \underline{and all} $z_i$'s \underline{are non-degenerate quadratic singularity rational over} $k$ . \underline{Then}

(1)  $R$ \underline{is formally smooth over} $\Lambda$ i.e., $R$ \underline{is a formal power-series ring over} $\Lambda$ ;

(2)  $J_{\mathfrak{X}}(\mathfrak{X}|R) = \coprod_{1\leq i\leq m} J_{\mathfrak{X}_{(z_i)}}(\mathfrak{X}_{(z_i)}|R)$  \underline{(disjoint union), and}

$J_{\mathfrak{X}_{(z_i)}}(\mathfrak{X}_{(z_i)}|R) \;\overset{\sim}{\relbar\relbar\relbar}\; J_R(\mathfrak{X}_{(z_i)}|R)$  \underline{is an isomorphism};

(3)  $I_R(\mathfrak{X}_{(z_i)}|R)$ \underline{is a principal ideal and in turn so is}

$I_R(\mathfrak{X}|R) = \prod_{1\leq\ i\leq\ m} I_R(\mathfrak{X}_{(z_i)}|R)$ . \underline{Furthermore if we set} $I_R(\mathfrak{X}_{(z_i)}|R) = (t_i)$ , \underline{then} $t_1,\ldots,t_m$ \underline{can be extended to a system of free formal generators of} $R$ \underline{over} $\Lambda$ (\underline{as formal power-series ring over} $\Lambda$ ).

(4)  $\mathfrak{X} \to \mathrm{Spec}(\Lambda)$ \underline{is smooth}, i.e. $J_{\mathfrak{X}}(\mathfrak{X}|\Lambda) = \varphi$ .


\SGApage{I}{128}

(5)  \underline{If} $X_0$ \underline{is a complete irreducible curve, then} $\dim \underline{M}_{X_0} = 3g - 3 + a$ \underline{where} $g = \dim_k H^1(X_0, \underline{O}_{X_0})$ \underline{and} $a = \dim_k H^0(X_0, \underline{D}^0_{X_0})$ .

\SGAvar{\underline{Proof:}  (1) follows from (2) and (3): The first statement of (2) follows from 5.7(2). Now let $(R_i, \mathfrak{Y}_i)$ be a formal versal deformation of $\mathrm{Spec}(\hat{\underline{O}}_{X_0,z_i})$ . Since $(R, \mathfrak{X}_{(\widehat{z}_i)})$ is a deformation of $\mathrm{Spec}(\hat{\underline{O}}_{X_0,z_i})$ , there exists a continuous local $\Lambda$-algebra maps $\varphi_i = R_i \to R$ with a cartesian diagram

}{\underline{Proof:}  (1) follows from (2) and (3): The first statement of (2) follows from 5.7(2). Now let $(R_i, \mathfrak{Y}_i)$ be a formal versal deformation of $\mathrm{Spec}(\hat{\underline{O}}_{X_0,z_i})$ . Since $(R, \mathfrak{X}_{(\widehat{z}_i)})$ is a deformation of $\mathrm{Spec}(\hat{\underline{O}}_{X_0,z_i})$ , there exists a continuous local $\Lambda$-algebra maps $\varphi_i : R_i \to R$ with a cartesian diagram

}{F7I-VI-128}\begin{center}
\begin{tikzcd}[column sep=huge, row sep=huge]
\mathfrak{X}_{(\widehat{z}_i)} \arrow[r] \arrow[d] \& \mathfrak{Y}_i \arrow[d] \\
\mathrm{Spf}(R) \arrow[r, "{\mathrm{Spf}(\varphi_i)}"'] \& \mathrm{Spf}(R_i)
\end{tikzcd}
\end{center}













Therefore we have

\begin{equation*}
J_{\mathfrak{X}_{(z_i)}}(\mathfrak{X}_{(z_i)}|R) = J_{\mathfrak{X}_{(\widehat{z}_i)}}(\mathfrak{X}_{(\widehat{z}_i)}|R) \simeq J_{\mathfrak{Y}_i}(\mathfrak{Y}_i|R_i) \underset{\mathrm{Spf}(R_i)}{\times} \mathrm{Spf}(R)
\end{equation*}
\begin{equation*}
J_R(\mathfrak{X}_{(z_i)}|R) = J_R(\mathfrak{X}_{(\widehat{z}_i)}|R) \simeq J_{R_i}(\mathfrak{Y}_i|R_i) \underset{\mathrm{Spf}(R_i)}{\times} \mathrm{Spf}(R) \quad .
\end{equation*}
Since $J_{\mathfrak{Y}_i}(\mathfrak{Y}_i|R_i) \to J_{R_i}(\mathfrak{Y}_i|R_i)$ is an isomorphism by 6.2(1), we have that the canonical morphism $J_{\mathfrak{X}_{(z_i)}}(\mathfrak{X}_{(z_i)}|R) \to J_R(\mathfrak{X}_{(z_i)}|R)$ is an isomorphism. Now $I_{R_i}(\mathfrak{Y}_i|R_i) = (u_i) = (u_i)$ with $R_i = \Lambda[[u_i]]$ by 6.2(2), and therefore

\SGAvar{\begin{equation*}
I_R(\mathfrak{X}_{(z_i)}|R) = I_R(\mathfrak{X}_{(\widehat{z}_i)}|R) = (t_i)
\end{equation*}
where $t_i = \varphi_i(u_i)$ , and therefore $I_R(\mathfrak{X}|R) = \prod_{1\leq\ i\leq\ m} I_R(\mathfrak{X}_{(z_i)}|R) = (\ \prod_{i=1}^{n} t_i\ )$ .

}{\begin{equation*}
I_R(\mathfrak{X}_{(z_i)}|R) = I_R(\mathfrak{X}_{(\widehat{z}_i)}|R) = (t_i)
\end{equation*}
where $t_i = \varphi_i(u_i)$ , and therefore $I_R(\mathfrak{X}|R) = \prod_{1\leq\ i\leq\ m} I_R(\mathfrak{X}_{(z_i)}|R) = (\ \prod_{i=1}^{m} t_i\ )$ .

}{F7I-VI-129}Now $\underline{M}_{X_0} \to \underline{M}_{X_0,(\widehat{z}_1)} \times \ldots \times \underline{M}_{X_0,(\widehat{z}_m)}$ is smooth by and therefore $R_1 \hat{\otimes}_{\Lambda} \ldots \hat{\otimes}_{\Lambda} R_m \to R$ is formally smooth by 4.13, which implies that $t_1,\ldots,t_m$ is extended to a system of free formal generators of $R$ over $\Lambda$ .


\SGApage{I}{129}

\SGAvar{(4)  The projection morphism  $\mathfrak{X}_{(\widehat{z}_i)} \overset{\sim}{\rightarrow} \mathfrak{Y}_i \times_{\mathrm{Spf}(R_i)} \mathrm{Spf}(R) \rightarrow \mathfrak{Y}_i$  is formally smooth since  $\mathrm{Spf}(R) \overset{\ \mathrm{Spf}(\varphi_i)\ }{\relbar\joinrel\relbar\joinrel\relbar\joinrel\relbar} \mathrm{Spf}(R_i)$  is formally smooth. On the other hand, $\mathfrak{Y}_i \rightarrow \mathrm{Spec}(\Lambda)$  is formally smooth by 6.2(3)  and therefore  $\mathfrak{X}_{(\widehat{z}_i)} \rightarrow \mathrm{Spec}(\Lambda)$  is formally smooth. Consequently  $\mathfrak{X} \rightarrow \mathrm{Spec}(\Lambda)$  is formally smooth since  $\{z_1,\dots,z_m\}$  is the set of non-smooth points of  $X_0$ .

}{(4)  The projection morphism  $\mathfrak{X}_{(\widehat{z}_i)} \overset{\sim}{\rightarrow} \mathfrak{Y}_i \times_{\mathrm{Spf}(R_i)} \mathrm{Spf}(R) \rightarrow \mathfrak{Y}_i$  is formally smooth since  $\mathrm{Spf}(R) \xrightarrow{\ \mathrm{Spf}(\varphi_i)\ } \mathrm{Spf}(R_i)$  is formally smooth. On the other hand, $\mathfrak{Y}_i \rightarrow \mathrm{Spec}(\Lambda)$  is formally smooth by 6.2(3)  and therefore  $\mathfrak{X}_{(\widehat{z}_i)} \rightarrow \mathrm{Spec}(\Lambda)$  is formally smooth. Consequently  $\mathfrak{X} \rightarrow \mathrm{Spec}(\Lambda)$  is formally smooth since  $\{z_1,\dots,z_m\}$  is the set of non-smooth points of  $X_0$ .

}{F7I-VI-130}(5)  Since  $R$  is formally smooth over  $\Lambda$  and

\begin{equation*}
0 \rightarrow H^1(X_0,\ \underline{D}^0_{X_0})\ \rightarrow [\underline{M}_{X_0}(k[\varepsilon])] \rightarrow H^0(X_0,\ \underline{D}^1_{X_0}) \rightarrow 0
\end{equation*}
is exact by 4.4 and 4.11, we have

\begin{equation*}
\begin{aligned}
\dim \underline{M}_{X_0} \ &= \ \dim k \otimes R \ = \ \dim_k H^1(X_0,\ \underline{D}^0_{X_0}) + \dim H^0(X_0,\ \underline{D}^1_{X_0})\\
&= \ \dim_k H^1(X_0,\ \underline{D}^0_{X_0}) + m \ .
\end{aligned}
\end{equation*}
Then it suffices to show that

\begin{equation*}
\dim_k H^1(X_0,\ \underline{D}^0_{X_0}) + m = 3g - 3 + \dim_k H^0(X_0,\ \underline{D}^0_{X_0})
\end{equation*}
i.e.,  $\chi(\underline{D}^0_{X_0}) = 3 - 3g + m$

\SGAvar{where we set  $\chi(\underline{E}) = \dim_k H^0(X_0,\ \underline{E}) - \dim_k H^1(X_0,\ \underline{E})$  for any coherent  $X_0$-module  $\underline{E}$ . Since  $\lambda(\underline{D}^0_{X_0})$  does not change when one replaces the base field  $k$  by its algebraic closure, we may assume that  $k$  is algebraically closed. Set  $A = \underline{O}_{X_0,z}$ , where  $z$  is a non-smooth point of  $X_0$ , and let  $A'$  be its normalization. By hypothesis, we have  $A \simeq k[[x,y]]/(f)$  where  $f$  is a non-degenerate quadratic form. Since  $k$  is algebraically closed, we may assume that  $f = xy$ , i.e., $A \simeq k[[x,y]]/(xy)$ . From this we conclude immediately that  $\ell(A'/A) = 1$  and  $\mathrm{Der}_k(A,\ \underline{m}_A) = \mathrm{Der}_k(A,\ A)$ . Let  $C$  be the conductor of  $A$  in  $A'$ . Since $A$ is a complete-intersection we must have that  $\ell(A/C) = \ell(A'/A) = 1$  so that  $C = \underline{m}_A$ . We thus obtain that  $A' = \{a' \in K |\ a'\underline{m}_A \subset \underline{m}_A\}$  where  $K$  is the fraction field of  $A$ . We now claim that  $\mathrm{Der}_k(A,\ A) \subset \mathrm{Der}_k(A',\ A')$  as $A$-modules contained in  $\mathrm{Der}_k(K,\ K)$ . Indeed, let  $X \in \mathrm{Der}_k(A,\ A)$ . Then for any  $a' \in A'$  and
\SGApage{I}{130}
$x \in \underline{m}_A$  we have  $xa' \in \underline{m}_A$  so that  $X(xa') \in \underline{m}_A$  and  $a'X(x) \in \underline{m}_A$  since  $\mathrm{Der}_k(A,\ \underline{m}_A) = \mathrm{Der}_k(A,\ A)$ . Consequently  $X(xa') = a'X(x) + xX(a')$  entails that  $xX(a') \in \underline{m}_A$  for all  $x \in \underline{m}_A$ , and therefore  $X(a') \in A'$ . Furthermore, the exact sequence

}{where we set  $\chi(\underline{E}) = \dim_k H^0(X_0,\ \underline{E}) - \dim_k H^1(X_0,\ \underline{E})$  for any coherent  $X_0$-module  $\underline{E}$ . Since  $\chi(\underline{D}^0_{X_0})$  does not change when one replaces the base field  $k$  by its algebraic closure, we may assume that  $k$  is algebraically closed. Set  $A = \underline{O}_{X_0,z}$ , where  $z$  is a non-smooth point of  $X_0$ , and let  $A'$  be its normalization. By hypothesis, we have  $A \simeq k[[x,y]]/(f)$  where  $f$  is a non-degenerate quadratic form. Since  $k$  is algebraically closed, we may assume that  $f = xy$ , i.e., $A \simeq k[[x,y]]/(xy)$ . From this we conclude immediately that  $\ell(A'/A) = 1$  and  $\mathrm{Der}_k(A,\ \underline{m}_A) = \mathrm{Der}_k(A,\ A)$ . Let  $C$  be the conductor of  $A$  in  $A'$ . Since $A$ is a complete-intersection we must have that  $\ell(A/C) = \ell(A'/A) = 1$  so that  $C = \underline{m}_A$ . We thus obtain that  $A' = \{a' \in K |\ a'\underline{m}_A \subset \underline{m}_A\}$  where  $K$  is the fraction field of  $A$ . We now claim that  $\mathrm{Der}_k(A,\ A) \subset \mathrm{Der}_k(A',\ A')$  as $A$-modules contained in  $\mathrm{Der}_k(K,\ K)$ . Indeed, let  $X \in \mathrm{Der}_k(A,\ A)$ . Then for any  $a' \in A'$  and
\SGApage{I}{130}
$x \in \underline{m}_A$  we have  $xa' \in \underline{m}_A$  so that  $X(xa') \in \underline{m}_A$  and  $a'X(x) \in \underline{m}_A$  since  $\mathrm{Der}_k(A,\ \underline{m}_A) = \mathrm{Der}_k(A,\ A)$ . Consequently  $X(xa') = a'X(x) + xX(a')$  entails that  $xX(a') \in \underline{m}_A$  for all  $x \in \underline{m}_A$ , and therefore  $X(a') \in A'$ . Furthermore, the exact sequence

}{F7I-VI-131}\begin{equation*}
0 \ \rightarrow\ \mathrm{Der}_k(A,\ A) \rightarrow \mathrm{Der}_k(A',\ A') \rightarrow A'/\underline{m}_A A' \rightarrow\ 0
\end{equation*}
entails the exact sequence

\begin{equation*}
0 \ \rightarrow\ \mathrm{Der}_k(A,\ A) \rightarrow \mathrm{Der}_k(A',\ A') \rightarrow A'/C \rightarrow 0
\end{equation*}
Therefore if we set  $X'_0$  to be the normalization of  $X_0$ , then

\begin{equation*}
0 \rightarrow \underline{D}^0_{X_0} \rightarrow \underline{D}^0_{X'_0}\ \rightarrow \underline{O}_{X'_0}/C \rightarrow 0
\end{equation*}
is exact. Since  $X_0$  is locally a complete-intersection, we have  $\ell(\underline{O}_{X'_0}/C) = 2\ell(\underline{O}_{X'_0}/\underline{O}_{X_0}) = 2m$ , whereas  $\chi(\underline{D}^0_{X'_0}) = 3 - 3(g-m)$  by Riemann-Roch theorem. Consequently

\SGAvar{\begin{equation*}
\chi(\underline{D}^0_{X^0_0}) = \chi(\underline{D}^0_{X'_0}) - \ell(\underline{O}_{X'_0}/C) = 3 - 3(g-m) - 2m = 3 - 3g + m \ .
\end{equation*}
This completes our proof.

}{\begin{equation*}
\chi(\underline{D}^0_{X_0}) = \chi(\underline{D}^0_{X'_0}) - \ell(\underline{O}_{X'_0}/C) = 3 - 3(g-m) - 2m = 3 - 3g + m \ .
\end{equation*}
This completes our proof.

}{F7I-VI-132}\underline{Remark} 6.5.  In the formula 6.4(5), the number  $a = H^0(X_0,\ \underline{D}^0_{X_0})$  can be easily computed in terms of Riemann-Roch theorem. Thus

\begin{equation*}
a = \left\{ \begin{array}{lll} 0 & \text{if} & g \geq 2 \\[4pt] 1 & \text{if} & g = 1 \\[4pt] 3 & \text{if} & g = 0 \end{array} \right.
\end{equation*}
Finally we add a remark on characteristic 2. If the characteristic of the base field  $k$  is  $2$ , there is no non-degenerate quadratic form on  $n$  variables if  $n$  is odd. In other words there does not exist a non-degenerate quadratic singularity of even dimension if  $\mathrm{char}\, k = 2$ . In this case one may consider a quadratic form of maximum non-degeneracy: 


\SGApage{I}{131}

\underline{Definition} 6.6.  \underline{Let}  $X$  \underline{be a} k-\underline{scheme of finite type.} \underline{A closed point}  $z \in X$  \underline{is called an ordinary quadratic singularity rational over}  $k$  \underline{if the following holds:}

\SGAvar{(i)  \underline{If either}  $char k \neq 2$  \underline{or}  $char k = 2$  \underline{and}  $dim O_{X,z}$  \underline{is odd, then} $a$  \underline{is a non-degenerate quadratic singularity rational over}  $k$  .

}{(i)  \underline{If either}  $char k \neq 2$  \underline{or}  $char k = 2$  \underline{and}  $dim O_{X,z}$  \underline{is odd, then} $z$  \underline{is a non-degenerate quadratic singularity rational over}  $k$  .

}{F7I-VI-133}(ii)  \underline{If}  $char k = 2$  \underline{and}  $dim O_{X,z}$  \underline{is even, then} $\bar{k} \underset{k}{\otimes} \hat{O}_{X,z} \simeq \bar{k}[[X_0, X_1, \ldots, X_{2n}]]/(f)$  \underline{with}

\begin{equation*}
f = X_0^2 + X_1X_2 + X_3X_4 + \ldots + X_{2n-1}X_{2n}
\end{equation*}
\underline{where}  $\bar{k}$  \underline{is the algebraic closure of}  $k$  .

\underline{Lemma} 6.7.  \underline{Let}  $k$  \underline{be a field of characteristic}  $2$  \underline{and}  \underline{set} $A_0 = k[[X_0, X_1, \ldots ,X_{2n}]]/(f)$  \underline{where}

\begin{equation*}
f = X_0^2 + X_1X_2 + X_3X_4 + \ldots + X_{2n-1}X_{2n} \quad ,
\end{equation*}
\underline{and let}  $(R, \mathfrak{X})$  \underline{be a formal versal deformation of}  $X_0 = \mathrm{Spec}(A_0)$  . \underline{Then} $R = \Lambda[[t, t_1]]$  , \underline{where}  $t, t_1$  \underline{are free formal generators of}  $R$  \underline{over}  $\Lambda$  , \underline{such} \underline{that}  $I_R(\mathfrak{X}|R) = (t^2)$  , \underline{and}  $\mathfrak{X} \to \mathrm{Spec}(\Lambda)$  \underline{is formally smooth.}

\underline{Proof:}  We set  $\mathfrak{X} = \mathrm{Spf}(A)$  . Since

\begin{equation*}
(f, \frac{\partial f}{\partial X_0} , \ldots , \frac{\partial f}{\partial X_{2n}} ) = (X_0^2, X_1, \ldots , X_{2n}) \quad ,
\end{equation*}
it follows that  $D^1(A_0|k, A_0) \simeq  k[X_0]/(X_0^2)$  and therefore  $[\underline{M}_{X_0}(k[\epsilon])]$  is 2-dimensional vector space over  $k$  . Furthermore, it follows from 4.14 that

\begin{equation*}
R \simeq \Lambda[[t, t_1]] \quad \text{and} \quad A \simeq R[[X_0, X_1, \ldots ,X_{2n}]]/(F)
\end{equation*}
where  $F = f + t_1 + tX_0$  . Now the exact sequence

\begin{center}
\begin{tikzcd}[column sep=large]
0 \arrow[r] \& A \arrow[r, "dF"] \& AdX_0 \oplus \ldots \oplus AdX_{2n} \arrow[r] \& \hat{\Omega}_{A|R} \arrow[r] \& 0
\end{tikzcd}
\end{center}










where

\SGAvar{\begin{equation*}
dF = 1 \to \sum_{0 \leq i \leq 2n} \frac{\partial f}{\partial X_i} dX_i + dX_0 \quad , 
\end{equation*}

}{\begin{equation*}
dF = 1 \to \sum_{0 \leq i \leq 2n} \frac{\partial f}{\partial X_i} dX_i + t\,dX_0 \quad , 
\end{equation*}

}{F7I-VI-134}\SGApage{I}{132}

 entails that

\begin{equation*}
I_A(A|R) = (t, x_1, x_2, \ldots, x_{2n}) = (x_0^2 + t_1, t, x_1, \ldots, x_{2n})
\end{equation*}
so that

\begin{equation*}
A/I_A(A|R) \simeq A/(X_0^2 + t_1, t, x_1, \ldots, x_{2n}) \simeq \Lambda[[t_1, X_0]]/(X_0^2 - t_1)  \ .
\end{equation*}
Since  $\Lambda[[t_1, X_0]]/(X_0^2 - t_1) \simeq R/(t) \oplus R/(t)$  as R-modules, it follows that $I_R(A|R) = (t^2)$  , which proves our first assertion. On the other hand,

\begin{equation*}
\begin{aligned} A &\simeq \Lambda[[t, t_1, x_o, x_1, \ldots, x_{2n}]]/(X_0^2 + X_1X_2 + \ldots + X_{2n-1}X_{2n} + tX_0 + t_1) \\ &\simeq \Lambda[[t, X_0, X_1,\ldots,X_{2n}]] \end{aligned}
\end{equation*}
and therefore  $A$  is formally smooth over  $\Lambda$  .

The above statement 6.7 yields a similar result as 6.4, proof of which we leave to the readers.

\underline{Corollary} 6.8.  \underline{Let}  $X_0$  \underline{be a} k-\underline{scheme of finite type with isolated non-smooth points only such that}  $H^2(X_0, \underline{D}^0_{X_0}) = 0$  , \underline{and let}  $(R, \mathfrak{X})$  \underline{be a formal versal deformation of}  $X_0$  . \underline{Denote by}  $\{z_1, \ldots, z_m\}$  \underline{the set of non-smooth points in} $X_0$  , \underline{and assume that}  $char k = 2$  , $dim X_0$  \underline{is even, and all}  $z_i$'s  \underline{are ordinary quadratic singularities. Then}

(1)  $R$  \underline{is formally smooth over}  $\Lambda$

(2)  $I_R(\mathfrak{X}_{(z_i)}|R)$  \underline{is a principal ideal and in turn so is}

\SGAvar{$I_R(\mathfrak{X}|R) = \mathop{\hphantom{\bigcap}}_{1 \leq i \leq m} I_R(\mathfrak{X}_{(z_i)}|R)$  . \underline{Furthermore,}  $I_R(\mathfrak{X}_{(z_i)}|R) = (t_i^2)$  \underline{and} $t_1, \ldots , t_m$  \underline{can be extended to a system of free formal generators of}  $R$  \underline{over} $\Lambda$  .

}{$I_R(\mathfrak{X}|R) = \prod_{1 \leq i \leq m} I_R(\mathfrak{X}_{(z_i)}|R)$  . \underline{Furthermore,}  $I_R(\mathfrak{X}_{(z_i)}|R) = (t_i^2)$  \underline{and} $t_1, \ldots , t_m$  \underline{can be extended to a system of free formal generators of}  $R$  \underline{over} $\Lambda$  .

}{F7I-VI-135}(3)  $\mathfrak{X} \to \mathrm{Spec}(\Lambda)$  \underline{is formally smooth}

(4)  \underline{If}  $k$  \underline{is algebraically closed and}  $X_0$  \underline{is a complete irreducible curve, then}  $dim \underline{M}_{X_0} = 3g - 3 + a$  \underline{where}  $g = dim_k H^1(X_0, \underline{O}_{X_0})$  \underline{and} $a = dim_k H^0(X_0, \underline{D}^0_{X_0})$  .
\selectlanguage{french}

\ifnum\SGAform=1
\clearpage\section*{Apparat critique}\small
Les lectures ci-dessous sont données dans le TeX exact des témoins, afin de préserver les distinctions de notation. Les indices décimaux renvoient aux différences individuelles du passage marqué.

\par\medskip\noindent\hypertarget{app-F7I-VI-001.1}{}\textbf{F7I-VI-001.1}\quad p. 32; misprint\_pass.\par
\noindent Imprimé : {\ttfamily\footnotesize composit\allowbreak{}}\par
\noindent Corrigé : {\ttfamily\footnotesize composite\allowbreak{}}\par
\noindent Source : 7I:53. \nolinkurl{fixes/FIXLOG_7I_CURRENT.jsonl#621}\par

\par\medskip\noindent\hypertarget{app-F7I-VI-001.2}{}\textbf{F7I-VI-001.2}\quad p. 32; misprint\_pass.\par
\noindent Imprimé : {\ttfamily\footnotesize morphism\allowbreak{}}\par
\noindent Corrigé : {\ttfamily\footnotesize morphisms\allowbreak{}}\par
\noindent Source : 7I:53. \nolinkurl{fixes/FIXLOG_7I_CURRENT.jsonl#621}\par

\par\medskip\noindent\hypertarget{app-F7I-VI-002.1}{}\textbf{F7I-VI-002.1}\quad p. 33; misprint\_pass.\par
\noindent Imprimé : {\ttfamily\footnotesize on\allowbreak{}}\par
\noindent Corrigé : {\ttfamily\footnotesize in\allowbreak{}}\par
\noindent Source : 7I:55. \nolinkurl{fixes/FIXLOG_7I_CURRENT.jsonl#622}; \nolinkurl{fixes/MISPRINTS_7I.json#55}\par

\par\medskip\noindent\hypertarget{app-F7I-VI-003.1}{}\textbf{F7I-VI-003.1}\quad p. 34; misprint\_pass.\par
\noindent Imprimé : {\ttfamily\footnotesize .\allowbreak{}}\par
\noindent Corrigé : {\ttfamily\footnotesize ,\allowbreak{}}\par
\noindent Source : 7I:56. \nolinkurl{fixes/FIXLOG_7I_CURRENT.jsonl#623}; \nolinkurl{fixes/MISPRINTS_7I.json#56}\par

\par\medskip\noindent\hypertarget{app-F7I-VI-003.2}{}\textbf{F7I-VI-003.2}\quad p. 34; misprint\_pass.\par
\noindent Imprimé : {\ttfamily\footnotesize Then\allowbreak{}}\par
\noindent Corrigé : {\ttfamily\footnotesize then\allowbreak{}}\par
\noindent Source : 7I:56. \nolinkurl{fixes/FIXLOG_7I_CURRENT.jsonl#623}; \nolinkurl{fixes/MISPRINTS_7I.json#56}\par

\par\medskip\noindent\hypertarget{app-F7I-VI-004.1}{}\textbf{F7I-VI-004.1}\quad p. 35; misprint\_pass.\par
\noindent Imprimé : \emph{(vide)}\par
\noindent Corrigé : {\ttfamily\footnotesize [\allowbreak{}}\par
\noindent Source : 7I:57. \nolinkurl{fixes/FIXLOG_7I_CURRENT.jsonl#624}; \nolinkurl{fixes/MISPRINTS_7I.json#57}\par

\par\medskip\noindent\hypertarget{app-F7I-VI-005.1}{}\textbf{F7I-VI-005.1}\quad p. 37; misprint\_pass.\par
\noindent Imprimé : \emph{(vide)}\par
\noindent Corrigé : {\ttfamily\footnotesize \_1\allowbreak{}}\par
\noindent Source : 7I:60. \nolinkurl{fixes/FIXLOG_7I_CURRENT.jsonl#626}; \nolinkurl{fixes/MISPRINTS_7I.json#60}\par

\par\medskip\noindent\hypertarget{app-F7I-VI-005.2}{}\textbf{F7I-VI-005.2}\quad p. 37; misprint\_pass.\par
\noindent Imprimé : \emph{(vide)}\par
\noindent Corrigé : {\ttfamily\footnotesize '\_1\allowbreak{}}\par
\noindent Source : 7I:61. \nolinkurl{fixes/FIXLOG_7I_CURRENT.jsonl#627}; \nolinkurl{fixes/MISPRINTS_7I.json#61}\par

\par\medskip\noindent\hypertarget{app-F7I-VI-005.3}{}\textbf{F7I-VI-005.3}\quad p. 37; misprint\_in\_print.\par
\noindent Imprimé : {\ttfamily\footnotesize h\allowbreak{}}\par
\noindent Corrigé : {\ttfamily\footnotesize k\allowbreak{}}\par
\noindent Source : FIX-7I-D-0025. \nolinkurl{fixes/FIXLOG_7I.jsonl#25}\par

\par\medskip\noindent\hypertarget{app-F7I-VI-005.4}{}\textbf{F7I-VI-005.4}\quad p. 37; misprint\_pass.\par
\noindent Imprimé : {\ttfamily\footnotesize f\allowbreak{}}\par
\noindent Corrigé : {\ttfamily\footnotesize h\allowbreak{}}\par
\noindent Source : 7I:59. \nolinkurl{fixes/FIXLOG_7I_CURRENT.jsonl#625}; \nolinkurl{fixes/MISPRINTS_7I.json#59}\par

\par\medskip\noindent\hypertarget{app-F7I-VI-005.5}{}\textbf{F7I-VI-005.5}\quad p. 37; misprint\_pass.\par
\noindent Imprimé : {\ttfamily\footnotesize h\allowbreak{}}\par
\noindent Corrigé : {\ttfamily\footnotesize R\allowbreak{}}\par
\noindent Source : 7I:62. \nolinkurl{fixes/FIXLOG_7I_CURRENT.jsonl#628}; \nolinkurl{fixes/MISPRINTS_7I.json#62}\par

\par\medskip\noindent\hypertarget{app-F7I-VI-005.6}{}\textbf{F7I-VI-005.6}\quad p. 37; misprint\_pass.\par
\noindent Imprimé : {\ttfamily\footnotesize \textbackslash{}\allowbreak{} }\par
\noindent Corrigé : {\ttfamily\footnotesize  \textbackslash{}times\allowbreak{} }\par
\noindent Source : 7I:63. \nolinkurl{fixes/FIXLOG_7I_CURRENT.jsonl#629}; \nolinkurl{fixes/MISPRINTS_7I.json#63}\par

\par\medskip\noindent\hypertarget{app-F7I-VI-005.7}{}\textbf{F7I-VI-005.7}\quad p. 37; misprint\_pass.\par
\noindent Imprimé : \emph{(vide)}\par
\noindent Corrigé : {\ttfamily\footnotesize '\_1\allowbreak{}}\par
\noindent Source : 7I:64. \nolinkurl{fixes/FIXLOG_7I_CURRENT.jsonl#630}; \nolinkurl{fixes/MISPRINTS_7I.json#64}\par

\par\medskip\noindent\hypertarget{app-F7I-VI-006.1}{}\textbf{F7I-VI-006.1}\quad p. 38; misprint\_in\_print.\par
\noindent Imprimé : {\ttfamily\footnotesize \textbackslash{}overset\allowbreak{}}\par
\noindent Corrigé : {\ttfamily\footnotesize \textbackslash{}xrightarrow\allowbreak{}}\par
\noindent Source : FIX-7I-D-0027. \nolinkurl{fixes/FIXLOG_7I.jsonl#27}\par

\par\medskip\noindent\hypertarget{app-F7I-VI-006.2}{}\textbf{F7I-VI-006.2}\quad p. 38; designated-witness attribution.\par
\noindent Imprimé : {\ttfamily\footnotesize \}\{\textbackslash{}relbar\textbackslash{}joinrel\textbackslash{}re\allowbreak{}lbar\textbackslash{}joinrel\textbackslash{}relbar\allowbreak{}}\par
\noindent Corrigé : \emph{(vide)}\par
\noindent Source : WITNESS-F7I-VI-006.2. \nolinkurl{fr/sga7i/diplomatic/expose_VI_body.tex}; \nolinkurl{fr/sga7i/current/expose_VI_body.tex}\par

\par\medskip\noindent\hypertarget{app-F7I-VI-007.1}{}\textbf{F7I-VI-007.1}\quad p. 38; misprint\_pass.\par
\noindent Imprimé : \emph{(vide)}\par
\noindent Corrigé : {\ttfamily\footnotesize  ,\allowbreak{}}\par
\noindent Source : 7I:66. \nolinkurl{fixes/FIXLOG_7I_CURRENT.jsonl#631}; \nolinkurl{fixes/MISPRINTS_7I.json#66}\par

\par\medskip\noindent\hypertarget{app-F7I-VI-008.1}{}\textbf{F7I-VI-008.1}\quad p. 39; misprint\_pass.\par
\noindent Imprimé : {\ttfamily\footnotesize definiton\allowbreak{}}\par
\noindent Corrigé : {\ttfamily\footnotesize definition\allowbreak{}}\par
\noindent Source : 7I:68. \nolinkurl{fixes/FIXLOG_7I_CURRENT.jsonl#632}; \nolinkurl{fixes/MISPRINTS_7I.json#68}\par

\par\medskip\noindent\hypertarget{app-F7I-VI-008.2}{}\textbf{F7I-VI-008.2}\quad p. 39; misprint\_pass.\par
\noindent Imprimé : {\ttfamily\footnotesize cases\allowbreak{}}\par
\noindent Corrigé : {\ttfamily\footnotesize case\allowbreak{}}\par
\noindent Source : 7I:69. \nolinkurl{fixes/FIXLOG_7I_CURRENT.jsonl#633}; \nolinkurl{fixes/MISPRINTS_7I.json#69}\par

\par\medskip\noindent\hypertarget{app-F7I-VI-008.3}{}\textbf{F7I-VI-008.3}\quad p. 39; misprint\_in\_print.\par
\noindent Imprimé : {\ttfamily\footnotesize ,\allowbreak{} no\allowbreak{} head\allowbreak{}}\par
\noindent Corrigé : \emph{(vide)}\par
\noindent Source : FIX-7I-D-0029. \nolinkurl{fixes/FIXLOG_7I.jsonl#29}\par

\par\medskip\noindent\hypertarget{app-F7I-VI-009.1}{}\textbf{F7I-VI-009.1}\quad p. 41; misprint\_in\_print.\par
\noindent Imprimé : \emph{(vide)}\par
\noindent Corrigé : {\ttfamily\footnotesize \textbackslash{}underline\{\allowbreak{}}\par
\noindent Source : FIX-7I-D-0031. \nolinkurl{fixes/FIXLOG_7I.jsonl#31}\par

\par\medskip\noindent\hypertarget{app-F7I-VI-009.2}{}\textbf{F7I-VI-009.2}\quad p. 41; misprint\_in\_print.\par
\noindent Imprimé : \emph{(vide)}\par
\noindent Corrigé : {\ttfamily\footnotesize \}\allowbreak{}}\par
\noindent Source : FIX-7I-D-0031. \nolinkurl{fixes/FIXLOG_7I.jsonl#31}\par

\par\medskip\noindent\hypertarget{app-F7I-VI-009.3}{}\textbf{F7I-VI-009.3}\quad p. 41; misprint\_in\_print.\par
\noindent Imprimé : \emph{(vide)}\par
\noindent Corrigé : {\ttfamily\footnotesize \textbackslash{}underline\{\allowbreak{}}\par
\noindent Source : FIX-7I-D-0032. \nolinkurl{fixes/FIXLOG_7I.jsonl#32}\par

\par\medskip\noindent\hypertarget{app-F7I-VI-009.4}{}\textbf{F7I-VI-009.4}\quad p. 41; misprint\_in\_print.\par
\noindent Imprimé : \emph{(vide)}\par
\noindent Corrigé : {\ttfamily\footnotesize \}\allowbreak{}}\par
\noindent Source : FIX-7I-D-0032. \nolinkurl{fixes/FIXLOG_7I.jsonl#32}\par

\par\medskip\noindent\hypertarget{app-F7I-VI-009.5}{}\textbf{F7I-VI-009.5}\quad p. 41; misprint\_pass.\par
\noindent Imprimé : {\ttfamily\footnotesize exists\allowbreak{}}\par
\noindent Corrigé : {\ttfamily\footnotesize exist\allowbreak{}}\par
\noindent Source : 7I:75. \nolinkurl{fixes/FIXLOG_7I_CURRENT.jsonl#634}; \nolinkurl{fixes/MISPRINTS_7I.json#75}\par

\par\medskip\noindent\hypertarget{app-F7I-VI-010.1}{}\textbf{F7I-VI-010.1}\quad p. 42; printed misprint.\par
\noindent Imprimé : {\ttfamily\footnotesize composit\allowbreak{}}\par
\noindent Corrigé : {\ttfamily\footnotesize composite\allowbreak{}}\par
\noindent Source : 7I:77. \nolinkurl{fixes/MISPRINTS_7I.json#77}\par

\par\medskip\noindent\hypertarget{app-F7I-VI-011.1}{}\textbf{F7I-VI-011.1}\quad p. 42; inherited corrected-layer difference.\par
\noindent Imprimé : \emph{(vide)}\par
\noindent Corrigé : {\ttfamily\footnotesize \_k\allowbreak{} }\par
\noindent Source : ADOPT-F7I-VI-011.1. \nolinkurl{fr/sga7i/diplomatic/expose_VI_body.tex}; \nolinkurl{fr/sga7i/current/expose_VI_body.tex}\par

\par\medskip\noindent\hypertarget{app-F7I-VI-012.1}{}\textbf{F7I-VI-012.1}\quad p. 43; inherited corrected-layer difference.\par
\noindent Imprimé : \emph{(vide)}\par
\noindent Corrigé : {\ttfamily\footnotesize \_k\allowbreak{}}\par
\noindent Source : ADOPT-F7I-VI-012.1. \nolinkurl{fr/sga7i/diplomatic/expose_VI_body.tex}; \nolinkurl{fr/sga7i/current/expose_VI_body.tex}\par

\par\medskip\noindent\hypertarget{app-F7I-VI-012.2}{}\textbf{F7I-VI-012.2}\quad p. 43; printed misprint.\par
\noindent Imprimé : \emph{(vide)}\par
\noindent Corrigé : {\ttfamily\footnotesize \_1\allowbreak{}}\par
\noindent Source : 7I:78. \nolinkurl{fixes/MISPRINTS_7I.json#78}\par

\par\medskip\noindent\hypertarget{app-F7I-VI-012.3}{}\textbf{F7I-VI-012.3}\quad p. 43; printed misprint.\par
\noindent Imprimé : {\ttfamily\footnotesize ,\allowbreak{} }\par
\noindent Corrigé : \emph{(vide)}\par
\noindent Source : 7I:79. \nolinkurl{fixes/MISPRINTS_7I.json#79}\par

\par\medskip\noindent\hypertarget{app-F7I-VI-012.4}{}\textbf{F7I-VI-012.4}\quad p. 43; printed misprint.\par
\noindent Imprimé : \emph{(vide)}\par
\noindent Corrigé : {\ttfamily\footnotesize )\allowbreak{}}\par
\noindent Source : 7I:80. \nolinkurl{fixes/MISPRINTS_7I.json#80}\par

\par\medskip\noindent\hypertarget{app-F7I-VI-013.1}{}\textbf{F7I-VI-013.1}\quad p. 43; printed misprint.\par
\noindent Imprimé : {\ttfamily\footnotesize composit\allowbreak{}}\par
\noindent Corrigé : {\ttfamily\footnotesize composite\allowbreak{}}\par
\noindent Source : 7I:81. \nolinkurl{fixes/MISPRINTS_7I.json#81}\par

\par\medskip\noindent\hypertarget{app-F7I-VI-014.1}{}\textbf{F7I-VI-014.1}\quad p. 44; misprint\_pass.\par
\noindent Imprimé : {\ttfamily\footnotesize sequences\allowbreak{}}\par
\noindent Corrigé : {\ttfamily\footnotesize sequence\allowbreak{}}\par
\noindent Source : 7I:82. \nolinkurl{fixes/FIXLOG_7I_CURRENT.jsonl#635}; \nolinkurl{fixes/MISPRINTS_7I.json#82}\par

\par\medskip\noindent\hypertarget{app-F7I-VI-015.1}{}\textbf{F7I-VI-015.1}\quad p. 45; misprint\_pass.\par
\noindent Imprimé : {\ttfamily\footnotesize biefly\allowbreak{}}\par
\noindent Corrigé : {\ttfamily\footnotesize briefly\allowbreak{}}\par
\noindent Source : 7I:83. \nolinkurl{fixes/FIXLOG_7I_CURRENT.jsonl#636}; \nolinkurl{fixes/MISPRINTS_7I.json#83}\par

\par\medskip\noindent\hypertarget{app-F7I-VI-015.2}{}\textbf{F7I-VI-015.2}\quad p. 45; misprint\_pass.\par
\noindent Imprimé : {\ttfamily\footnotesize preceeding\allowbreak{}}\par
\noindent Corrigé : {\ttfamily\footnotesize preceding\allowbreak{}}\par
\noindent Source : 7I:84. \nolinkurl{fixes/FIXLOG_7I_CURRENT.jsonl#637}; \nolinkurl{fixes/MISPRINTS_7I.json#84}\par

\par\medskip\noindent\hypertarget{app-F7I-VI-016.1}{}\textbf{F7I-VI-016.1}\quad p. 45; misprint\_pass.\par
\noindent Imprimé : {\ttfamily\footnotesize k\allowbreak{}}\par
\noindent Corrigé : {\ttfamily\footnotesize K\allowbreak{}}\par
\noindent Source : 7I:85. \nolinkurl{fixes/FIXLOG_7I_CURRENT.jsonl#638}; \nolinkurl{fixes/MISPRINTS_7I.json#85}\par

\par\medskip\noindent\hypertarget{app-F7I-VI-017.1}{}\textbf{F7I-VI-017.1}\quad p. 46; misprint\_pass.\par
\noindent Imprimé : \emph{(vide)}\par
\noindent Corrigé : {\ttfamily\footnotesize \textbackslash{},\textbackslash{}varepsilon\allowbreak{}}\par
\noindent Source : 7I:86. \nolinkurl{fixes/FIXLOG_7I_CURRENT.jsonl#639}; \nolinkurl{fixes/MISPRINTS_7I.json#86}\par

\par\medskip\noindent\hypertarget{app-F7I-VI-018.1}{}\textbf{F7I-VI-018.1}\quad p. 47; misprint\_pass.\par
\noindent Imprimé : {\ttfamily\footnotesize k\allowbreak{}}\par
\noindent Corrigé : {\ttfamily\footnotesize 1\allowbreak{}}\par
\noindent Source : 7I:87. \nolinkurl{fixes/FIXLOG_7I_CURRENT.jsonl#640}; \nolinkurl{fixes/MISPRINTS_7I.json#87}\par

\par\medskip\noindent\hypertarget{app-F7I-VI-019.1}{}\textbf{F7I-VI-019.1}\quad p. 47; printed misprint.\par
\noindent Imprimé : \emph{(vide)}\par
\noindent Corrigé : {\ttfamily\footnotesize '\allowbreak{}}\par
\noindent Source : 7I:89. \nolinkurl{fixes/MISPRINTS_7I.json#89}\par

\par\medskip\noindent\hypertarget{app-F7I-VI-019.2}{}\textbf{F7I-VI-019.2}\quad p. 47; printed misprint.\par
\noindent Imprimé : \emph{(vide)}\par
\noindent Corrigé : {\ttfamily\footnotesize \textbackslash{},\allowbreak{}}\par
\noindent Source : 7I:89. \nolinkurl{fixes/MISPRINTS_7I.json#89}\par

\par\medskip\noindent\hypertarget{app-F7I-VI-020.1}{}\textbf{F7I-VI-020.1}\quad p. 48; misprint\_pass.\par
\noindent Imprimé : {\ttfamily\footnotesize diagramm\allowbreak{}}\par
\noindent Corrigé : {\ttfamily\footnotesize diagram\allowbreak{}}\par
\noindent Source : 7I:90. \nolinkurl{fixes/FIXLOG_7I_CURRENT.jsonl#641}; \nolinkurl{fixes/MISPRINTS_7I.json#90}\par

\par\medskip\noindent\hypertarget{app-F7I-VI-020.2}{}\textbf{F7I-VI-020.2}\quad p. 48; misprint\_pass.\par
\noindent Imprimé : {\ttfamily\footnotesize ,\allowbreak{}}\par
\noindent Corrigé : {\ttfamily\footnotesize \textbackslash{},\allowbreak{}}\par
\noindent Source : 7I:91. \nolinkurl{fixes/FIXLOG_7I_CURRENT.jsonl#642}\par

\par\medskip\noindent\hypertarget{app-F7I-VI-020.3}{}\textbf{F7I-VI-020.3}\quad p. 48; misprint\_pass.\par
\noindent Imprimé : \emph{(vide)}\par
\noindent Corrigé : {\ttfamily\footnotesize '\allowbreak{}}\par
\noindent Source : 7I:92. \nolinkurl{fixes/FIXLOG_7I_CURRENT.jsonl#643}; \nolinkurl{fixes/MISPRINTS_7I.json#92}\par

\par\medskip\noindent\hypertarget{app-F7I-VI-020.4}{}\textbf{F7I-VI-020.4}\quad p. 48; misprint\_pass.\par
\noindent Imprimé : \emph{(vide)}\par
\noindent Corrigé : {\ttfamily\footnotesize \textbackslash{},\allowbreak{}}\par
\noindent Source : 7I:92. \nolinkurl{fixes/FIXLOG_7I_CURRENT.jsonl#643}; \nolinkurl{fixes/MISPRINTS_7I.json#92}\par

\par\medskip\noindent\hypertarget{app-F7I-VI-020.5}{}\textbf{F7I-VI-020.5}\quad p. 48; misprint\_pass.\par
\noindent Imprimé : {\ttfamily\footnotesize H\allowbreak{}}\par
\noindent Corrigé : {\ttfamily\footnotesize A\allowbreak{}}\par
\noindent Source : 7I:93. \nolinkurl{fixes/FIXLOG_7I_CURRENT.jsonl#644}; \nolinkurl{fixes/MISPRINTS_7I.json#93}\par

\par\medskip\noindent\hypertarget{app-F7I-VI-020.6}{}\textbf{F7I-VI-020.6}\quad p. 48; misprint\_pass.\par
\noindent Imprimé : {\ttfamily\footnotesize F\allowbreak{}}\par
\noindent Corrigé : {\ttfamily\footnotesize G\allowbreak{}}\par
\noindent Source : 7I:94. \nolinkurl{fixes/FIXLOG_7I_CURRENT.jsonl#645}; \nolinkurl{fixes/MISPRINTS_7I.json#94}\par

\par\medskip\noindent\hypertarget{app-F7I-VI-021.1}{}\textbf{F7I-VI-021.1}\quad p. 48; misprint\_pass.\par
\noindent Imprimé : {\ttfamily\footnotesize 2\allowbreak{}}\par
\noindent Corrigé : {\ttfamily\footnotesize 1\allowbreak{}}\par
\noindent Source : 7I:96. \nolinkurl{fixes/FIXLOG_7I_CURRENT.jsonl#646}; \nolinkurl{fixes/MISPRINTS_7I.json#96}\par

\par\medskip\noindent\hypertarget{app-F7I-VI-022.1}{}\textbf{F7I-VI-022.1}\quad p. 49; misprint\_pass.\par
\noindent Imprimé : {\ttfamily\footnotesize .\allowbreak{}}\par
\noindent Corrigé : \emph{(1 espace(s), 0 retour(s) de ligne dans le témoin TeX)}\par
\noindent Source : 7I:98. \nolinkurl{fixes/FIXLOG_7I_CURRENT.jsonl#647}; \nolinkurl{fixes/MISPRINTS_7I.json#98}\par

\par\medskip\noindent\hypertarget{app-F7I-VI-023.1}{}\textbf{F7I-VI-023.1}\quad p. 51; misprint\_pass.\par
\noindent Imprimé : {\ttfamily\footnotesize sitting\allowbreak{}}\par
\noindent Corrigé : {\ttfamily\footnotesize setting\allowbreak{}}\par
\noindent Source : 7I:99. \nolinkurl{fixes/FIXLOG_7I_CURRENT.jsonl#648}; \nolinkurl{fixes/MISPRINTS_7I.json#99}\par

\par\medskip\noindent\hypertarget{app-F7I-VI-024.1}{}\textbf{F7I-VI-024.1}\quad p. 51; misprint\_pass.\par
\noindent Imprimé : {\ttfamily\footnotesize \textbackslash{}varphi\allowbreak{}}\par
\noindent Corrigé : {\ttfamily\footnotesize \textbackslash{}psi\allowbreak{}}\par
\noindent Source : 7I:100. \nolinkurl{fixes/FIXLOG_7I_CURRENT.jsonl#649}; \nolinkurl{fixes/MISPRINTS_7I.json#100}\par

\par\medskip\noindent\hypertarget{app-F7I-VI-025.1}{}\textbf{F7I-VI-025.1}\quad p. 52; misprint\_pass.\par
\noindent Imprimé : {\ttfamily\footnotesize groupoids\allowbreak{}}\par
\noindent Corrigé : {\ttfamily\footnotesize groupoid\allowbreak{}}\par
\noindent Source : 7I:101. \nolinkurl{fixes/FIXLOG_7I_CURRENT.jsonl#650}; \nolinkurl{fixes/MISPRINTS_7I.json#101}\par

\par\medskip\noindent\hypertarget{app-F7I-VI-025.2}{}\textbf{F7I-VI-025.2}\quad p. 52; misprint\_pass.\par
\noindent Imprimé : {\ttfamily\footnotesize 13\allowbreak{}}\par
\noindent Corrigé : {\ttfamily\footnotesize 14\allowbreak{}}\par
\noindent Source : 7I:102. \nolinkurl{fixes/FIXLOG_7I_CURRENT.jsonl#651}; \nolinkurl{fixes/MISPRINTS_7I.json#102}\par

\par\medskip\noindent\hypertarget{app-F7I-VI-026.1}{}\textbf{F7I-VI-026.1}\quad p. 53; misprint\_pass.\par
\noindent Imprimé : \emph{(vide)}\par
\noindent Corrigé : {\ttfamily\footnotesize \_1\allowbreak{}}\par
\noindent Source : 7I:104. \nolinkurl{fixes/FIXLOG_7I_CURRENT.jsonl#652}; \nolinkurl{fixes/MISPRINTS_7I.json#104}\par

\par\medskip\noindent\hypertarget{app-F7I-VI-027.1}{}\textbf{F7I-VI-027.1}\quad p. 56; misprint\_pass.\par
\noindent Imprimé : \emph{(vide)}\par
\noindent Corrigé : {\ttfamily\footnotesize )\allowbreak{}}\par
\noindent Source : 7I:110. \nolinkurl{fixes/FIXLOG_7I_CURRENT.jsonl#653}; \nolinkurl{fixes/MISPRINTS_7I.json#110}\par

\par\medskip\noindent\hypertarget{app-F7I-VI-028.1}{}\textbf{F7I-VI-028.1}\quad p. 57; printed misprint.\par
\noindent Imprimé : {\ttfamily\footnotesize \textbackslash{}varphi\allowbreak{}}\par
\noindent Corrigé : {\ttfamily\footnotesize \textbackslash{}psi\allowbreak{}}\par
\noindent Source : 7I:111. \nolinkurl{fixes/MISPRINTS_7I.json#111}\par

\par\medskip\noindent\hypertarget{app-F7I-VI-028.2}{}\textbf{F7I-VI-028.2}\quad p. 57; printed misprint.\par
\noindent Imprimé : {\ttfamily\footnotesize \textbackslash{}psi\allowbreak{}}\par
\noindent Corrigé : {\ttfamily\footnotesize \textbackslash{}varphi\allowbreak{}}\par
\noindent Source : 7I:111. \nolinkurl{fixes/MISPRINTS_7I.json#111}\par

\par\medskip\noindent\hypertarget{app-F7I-VI-029.1}{}\textbf{F7I-VI-029.1}\quad p. 57; printed misprint.\par
\noindent Imprimé : {\ttfamily\footnotesize be\allowbreak{}}\par
\noindent Corrigé : {\ttfamily\footnotesize by\allowbreak{}}\par
\noindent Source : 7I:112. \nolinkurl{fixes/MISPRINTS_7I.json#112}\par

\par\medskip\noindent\hypertarget{app-F7I-VI-030.1}{}\textbf{F7I-VI-030.1}\quad p. 58; designated-witness attribution.\par
\noindent Imprimé : \emph{(vide)}\par
\noindent Corrigé : {\ttfamily\footnotesize '\allowbreak{}}\par
\noindent Source : WITNESS-F7I-VI-030.1. \nolinkurl{fr/sga7i/diplomatic/expose_VI_body.tex}; \nolinkurl{fr/sga7i/current/expose_VI_body.tex}\par

\par\medskip\noindent\hypertarget{app-F7I-VI-031.1}{}\textbf{F7I-VI-031.1}\quad p. 59; printed misprint.\par
\noindent Imprimé : \emph{(vide)}\par
\noindent Corrigé : {\ttfamily\footnotesize )\allowbreak{}}\par
\noindent Source : 7I:114. \nolinkurl{fixes/MISPRINTS_7I.json#114}\par

\par\medskip\noindent\hypertarget{app-F7I-VI-032.1}{}\textbf{F7I-VI-032.1}\quad p. 60; misprint\_pass.\par
\noindent Imprimé : {\ttfamily\footnotesize cocartegory\allowbreak{}}\par
\noindent Corrigé : {\ttfamily\footnotesize cocategory\allowbreak{}}\par
\noindent Source : 7I:116. \nolinkurl{fixes/FIXLOG_7I_CURRENT.jsonl#654}; \nolinkurl{fixes/MISPRINTS_7I.json#116}\par

\par\medskip\noindent\hypertarget{app-F7I-VI-033.1}{}\textbf{F7I-VI-033.1}\quad p. 61; misprint\_in\_print.\par
\noindent Imprimé : {\ttfamily\footnotesize =\textbackslash{}!=\allowbreak{}}\par
\noindent Corrigé : {\ttfamily\footnotesize \textbackslash{}rightrightarrows\allowbreak{}}\par
\noindent Source : FIX-7I-D-0041. \nolinkurl{fixes/FIXLOG_7I.jsonl#41}\par

\par\medskip\noindent\hypertarget{app-F7I-VI-033.2}{}\textbf{F7I-VI-033.2}\quad p. 61; misprint\_pass.\par
\noindent Imprimé : {\ttfamily\footnotesize \textbackslash{}qquad\allowbreak{}}\par
\noindent Corrigé : {\ttfamily\footnotesize \textbackslash{}amalg\allowbreak{}}\par
\noindent Source : 7I:118. \nolinkurl{fixes/FIXLOG_7I_CURRENT.jsonl#655}; \nolinkurl{fixes/MISPRINTS_7I.json#118}\par

\par\medskip\noindent\hypertarget{app-F7I-VI-034.1}{}\textbf{F7I-VI-034.1}\quad p. 62; misprint\_pass.\par
\noindent Imprimé : {\ttfamily\footnotesize composit\allowbreak{}}\par
\noindent Corrigé : {\ttfamily\footnotesize composite\allowbreak{}}\par
\noindent Source : 7I:120. \nolinkurl{fixes/FIXLOG_7I_CURRENT.jsonl#656}; \nolinkurl{fixes/MISPRINTS_7I.json#120}\par

\par\medskip\noindent\hypertarget{app-F7I-VI-034.2}{}\textbf{F7I-VI-034.2}\quad p. 62; misprint\_pass.\par
\noindent Imprimé : \emph{(vide)}\par
\noindent Corrigé : {\ttfamily\footnotesize  :\allowbreak{}}\par
\noindent Source : 7I:121. \nolinkurl{fixes/FIXLOG_7I_CURRENT.jsonl#657}; \nolinkurl{fixes/MISPRINTS_7I.json#121}\par

\par\medskip\noindent\hypertarget{app-F7I-VI-034.3}{}\textbf{F7I-VI-034.3}\quad p. 62; misprint\_pass.\par
\noindent Imprimé : {\ttfamily\footnotesize an\allowbreak{}}\par
\noindent Corrigé : {\ttfamily\footnotesize a\allowbreak{}}\par
\noindent Source : 7I:122. \nolinkurl{fixes/FIXLOG_7I_CURRENT.jsonl#658}; \nolinkurl{fixes/MISPRINTS_7I.json#122}\par

\par\medskip\noindent\hypertarget{app-F7I-VI-035.1}{}\textbf{F7I-VI-035.1}\quad p. 63; misprint\_pass.\par
\noindent Imprimé : {\ttfamily\footnotesize or\allowbreak{}}\par
\noindent Corrigé : {\ttfamily\footnotesize of\allowbreak{}}\par
\noindent Source : 7I:123. \nolinkurl{fixes/FIXLOG_7I_CURRENT.jsonl#659}; \nolinkurl{fixes/MISPRINTS_7I.json#123}\par

\par\medskip\noindent\hypertarget{app-F7I-VI-036.1}{}\textbf{F7I-VI-036.1}\quad p. 65; misprint\_pass.\par
\noindent Imprimé : \emph{(vide)}\par
\noindent Corrigé : {\ttfamily\footnotesize )\allowbreak{}}\par
\noindent Source : 7I:124. \nolinkurl{fixes/FIXLOG_7I_CURRENT.jsonl#660}\par

\par\medskip\noindent\hypertarget{app-F7I-VI-036.2}{}\textbf{F7I-VI-036.2}\quad p. 65; misprint\_pass.\par
\noindent Imprimé : {\ttfamily\footnotesize ,\allowbreak{}}\par
\noindent Corrigé : \emph{(vide)}\par
\noindent Source : 7I:125. \nolinkurl{fixes/FIXLOG_7I_CURRENT.jsonl#661}; \nolinkurl{fixes/MISPRINTS_7I.json#125}\par

\par\medskip\noindent\hypertarget{app-F7I-VI-036.3}{}\textbf{F7I-VI-036.3}\quad p. 65; printed misprint.\par
\noindent Imprimé : \emph{(vide)}\par
\noindent Corrigé : {\ttfamily\footnotesize \textbackslash{}underline\{\allowbreak{}}\par
\noindent Source : 7I:126. \nolinkurl{fixes/MISPRINTS_7I.json#126}\par

\par\medskip\noindent\hypertarget{app-F7I-VI-036.4}{}\textbf{F7I-VI-036.4}\quad p. 65; designated-witness attribution.\par
\noindent Imprimé : \emph{(vide)}\par
\noindent Corrigé : {\ttfamily\footnotesize \}\allowbreak{}}\par
\noindent Source : WITNESS-F7I-VI-036.4. \nolinkurl{fr/sga7i/diplomatic/expose_VI_body.tex}; \nolinkurl{fr/sga7i/current/expose_VI_body.tex}\par

\par\medskip\noindent\hypertarget{app-F7I-VI-036.5}{}\textbf{F7I-VI-036.5}\quad p. 65; designated-witness attribution.\par
\noindent Imprimé : \emph{(vide)}\par
\noindent Corrigé : {\ttfamily\footnotesize \textbackslash{}hat\{\allowbreak{}}\par
\noindent Source : WITNESS-F7I-VI-036.5. \nolinkurl{fr/sga7i/diplomatic/expose_VI_body.tex}; \nolinkurl{fr/sga7i/current/expose_VI_body.tex}\par

\par\medskip\noindent\hypertarget{app-F7I-VI-036.6}{}\textbf{F7I-VI-036.6}\quad p. 65; designated-witness attribution.\par
\noindent Imprimé : \emph{(vide)}\par
\noindent Corrigé : {\ttfamily\footnotesize \}\allowbreak{}}\par
\noindent Source : WITNESS-F7I-VI-036.6. \nolinkurl{fr/sga7i/diplomatic/expose_VI_body.tex}; \nolinkurl{fr/sga7i/current/expose_VI_body.tex}\par

\par\medskip\noindent\hypertarget{app-F7I-VI-036.7}{}\textbf{F7I-VI-036.7}\quad p. 65; designated-witness attribution.\par
\noindent Imprimé : {\ttfamily\footnotesize sets\allowbreak{}}\par
\noindent Corrigé : {\ttfamily\footnotesize set\allowbreak{}}\par
\noindent Source : WITNESS-F7I-VI-036.7. \nolinkurl{fr/sga7i/diplomatic/expose_VI_body.tex}; \nolinkurl{fr/sga7i/current/expose_VI_body.tex}\par

\par\medskip\noindent\hypertarget{app-F7I-VI-036.8}{}\textbf{F7I-VI-036.8}\quad p. 65; inherited corrected-layer difference.\par
\noindent Imprimé : \emph{(vide)}\par
\noindent Corrigé : {\ttfamily\footnotesize \textbackslash{}hat\{\allowbreak{}}\par
\noindent Source : ADOPT-F7I-VI-036.8. \nolinkurl{fr/sga7i/diplomatic/expose_VI_body.tex}; \nolinkurl{fr/sga7i/current/expose_VI_body.tex}\par

\par\medskip\noindent\hypertarget{app-F7I-VI-036.9}{}\textbf{F7I-VI-036.9}\quad p. 65; inherited corrected-layer difference.\par
\noindent Imprimé : \emph{(vide)}\par
\noindent Corrigé : {\ttfamily\footnotesize \}\allowbreak{}}\par
\noindent Source : ADOPT-F7I-VI-036.9. \nolinkurl{fr/sga7i/diplomatic/expose_VI_body.tex}; \nolinkurl{fr/sga7i/current/expose_VI_body.tex}\par

\par\medskip\noindent\hypertarget{app-F7I-VI-036.10}{}\textbf{F7I-VI-036.10}\quad p. 65; printed misprint.\par
\noindent Imprimé : {\ttfamily\footnotesize get\allowbreak{}}\par
\noindent Corrigé : {\ttfamily\footnotesize gets\allowbreak{}}\par
\noindent Source : 7I:127. \nolinkurl{fixes/MISPRINTS_7I.json#127}\par

\par\medskip\noindent\hypertarget{app-F7I-VI-037.1}{}\textbf{F7I-VI-037.1}\quad p. 66; printed misprint.\par
\noindent Imprimé : \emph{(vide)}\par
\noindent Corrigé : {\ttfamily\footnotesize ]\allowbreak{}}\par
\noindent Source : 7I:128. \nolinkurl{fixes/MISPRINTS_7I.json#128}\par

\par\medskip\noindent\hypertarget{app-F7I-VI-038.1}{}\textbf{F7I-VI-038.1}\quad p. 67; misprint\_pass.\par
\noindent Imprimé : {\ttfamily\footnotesize  of\allowbreak{}}\par
\noindent Corrigé : \emph{(vide)}\par
\noindent Source : 7I:129. \nolinkurl{fixes/FIXLOG_7I_CURRENT.jsonl#662}\par

\par\medskip\noindent\hypertarget{app-F7I-VI-039.1}{}\textbf{F7I-VI-039.1}\quad p. 68; misprint\_pass.\par
\noindent Imprimé : \emph{(vide)}\par
\noindent Corrigé : {\ttfamily\footnotesize )\allowbreak{}}\par
\noindent Source : 7I:131. \nolinkurl{fixes/FIXLOG_7I_CURRENT.jsonl#663}; \nolinkurl{fixes/MISPRINTS_7I.json#131}\par

\par\medskip\noindent\hypertarget{app-F7I-VI-039.2}{}\textbf{F7I-VI-039.2}\quad p. 68; misprint\_pass.\par
\noindent Imprimé : {\ttfamily\footnotesize composit\allowbreak{}}\par
\noindent Corrigé : {\ttfamily\footnotesize composite\allowbreak{}}\par
\noindent Source : 7I:132. \nolinkurl{fixes/FIXLOG_7I_CURRENT.jsonl#664}; \nolinkurl{fixes/MISPRINTS_7I.json#132}\par

\par\medskip\noindent\hypertarget{app-F7I-VI-040.1}{}\textbf{F7I-VI-040.1}\quad p. 68; misprint\_pass.\par
\noindent Imprimé : {\ttfamily\footnotesize composit\allowbreak{}}\par
\noindent Corrigé : {\ttfamily\footnotesize composite\allowbreak{}}\par
\noindent Source : 7I:133. \nolinkurl{fixes/FIXLOG_7I_CURRENT.jsonl#665}; \nolinkurl{fixes/MISPRINTS_7I.json#133}\par

\par\medskip\noindent\hypertarget{app-F7I-VI-041.1}{}\textbf{F7I-VI-041.1}\quad p. 69; designated-witness attribution.\par
\noindent Imprimé : {\ttfamily\footnotesize cogroupoid\allowbreak{}}\par
\noindent Corrigé : {\ttfamily\footnotesize cofibered\allowbreak{} groupoid\allowbreak{}}\par
\noindent Source : WITNESS-F7I-VI-041.1. \nolinkurl{fr/sga7i/diplomatic/expose_VI_body.tex}; \nolinkurl{fr/sga7i/current/expose_VI_body.tex}\par

\par\medskip\noindent\hypertarget{app-F7I-VI-042.1}{}\textbf{F7I-VI-042.1}\quad p. 70; printed misprint.\par
\noindent Imprimé : {\ttfamily\footnotesize its\allowbreak{}}\par
\noindent Corrigé : {\ttfamily\footnotesize the\allowbreak{}}\par
\noindent Source : 7I:136. \nolinkurl{fixes/MISPRINTS_7I.json#136}\par

\par\medskip\noindent\hypertarget{app-F7I-VI-042.2}{}\textbf{F7I-VI-042.2}\quad p. 70; misprint\_pass.\par
\noindent Imprimé : {\ttfamily\footnotesize 0\allowbreak{}}\par
\noindent Corrigé : {\ttfamily\footnotesize 1\allowbreak{}}\par
\noindent Source : 7I:135. \nolinkurl{fixes/FIXLOG_7I_CURRENT.jsonl#666}; \nolinkurl{fixes/MISPRINTS_7I.json#135}\par

\par\medskip\noindent\hypertarget{app-F7I-VI-042.3}{}\textbf{F7I-VI-042.3}\quad p. 70; inherited corrected-layer difference.\par
\noindent Imprimé : {\ttfamily\footnotesize cogroupoid\allowbreak{}}\par
\noindent Corrigé : {\ttfamily\footnotesize cogroupoids\allowbreak{}}\par
\noindent Source : ADOPT-F7I-VI-042.3. \nolinkurl{fr/sga7i/diplomatic/expose_VI_body.tex}; \nolinkurl{fr/sga7i/current/expose_VI_body.tex}\par

\par\medskip\noindent\hypertarget{app-F7I-VI-043.1}{}\textbf{F7I-VI-043.1}\quad p. 72; misprint\_pass.\par
\noindent Imprimé : {\ttfamily\footnotesize \textbackslash{}circ\allowbreak{}}\par
\noindent Corrigé : {\ttfamily\footnotesize 0\allowbreak{}}\par
\noindent Source : 7I:137. \nolinkurl{fixes/FIXLOG_7I_CURRENT.jsonl#667}; \nolinkurl{fixes/MISPRINTS_7I.json#137}\par

\par\medskip\noindent\hypertarget{app-F7I-VI-043.2}{}\textbf{F7I-VI-043.2}\quad p. 72; misprint\_pass.\par
\noindent Imprimé : {\ttfamily\footnotesize \textbackslash{}circ\allowbreak{}}\par
\noindent Corrigé : {\ttfamily\footnotesize 0\allowbreak{}}\par
\noindent Source : 7I:138. \nolinkurl{fixes/FIXLOG_7I_CURRENT.jsonl#668}; \nolinkurl{fixes/MISPRINTS_7I.json#138}\par

\par\medskip\noindent\hypertarget{app-F7I-VI-043.3}{}\textbf{F7I-VI-043.3}\quad p. 72; misprint\_pass.\par
\noindent Imprimé : {\ttfamily\footnotesize \textbackslash{}circ\allowbreak{}}\par
\noindent Corrigé : {\ttfamily\footnotesize 0\allowbreak{}}\par
\noindent Source : 7I:139. \nolinkurl{fixes/FIXLOG_7I_CURRENT.jsonl#669}; \nolinkurl{fixes/MISPRINTS_7I.json#139}\par

\par\medskip\noindent\hypertarget{app-F7I-VI-044.1}{}\textbf{F7I-VI-044.1}\quad p. 73; misprint\_pass.\par
\noindent Imprimé : \emph{(vide)}\par
\noindent Corrigé : {\ttfamily\footnotesize \textbackslash{}overset\{\textbackslash{}vee\}\{\allowbreak{}}\par
\noindent Source : 7I:140. \nolinkurl{fixes/FIXLOG_7I_CURRENT.jsonl#670}; \nolinkurl{fixes/MISPRINTS_7I.json#140}\par

\par\medskip\noindent\hypertarget{app-F7I-VI-044.2}{}\textbf{F7I-VI-044.2}\quad p. 73; misprint\_pass.\par
\noindent Imprimé : \emph{(vide)}\par
\noindent Corrigé : {\ttfamily\footnotesize \}\allowbreak{}}\par
\noindent Source : 7I:140. \nolinkurl{fixes/FIXLOG_7I_CURRENT.jsonl#670}; \nolinkurl{fixes/MISPRINTS_7I.json#140}\par

\par\medskip\noindent\hypertarget{app-F7I-VI-044.3}{}\textbf{F7I-VI-044.3}\quad p. 73; misprint\_pass.\par
\noindent Imprimé : \emph{(vide)}\par
\noindent Corrigé : {\ttfamily\footnotesize \textbackslash{}overset\{\textbackslash{}vee\}\{\allowbreak{}}\par
\noindent Source : 7I:140. \nolinkurl{fixes/FIXLOG_7I_CURRENT.jsonl#670}; \nolinkurl{fixes/MISPRINTS_7I.json#140}\par

\par\medskip\noindent\hypertarget{app-F7I-VI-044.4}{}\textbf{F7I-VI-044.4}\quad p. 73; misprint\_pass.\par
\noindent Imprimé : \emph{(vide)}\par
\noindent Corrigé : {\ttfamily\footnotesize \}\allowbreak{}}\par
\noindent Source : 7I:140. \nolinkurl{fixes/FIXLOG_7I_CURRENT.jsonl#670}; \nolinkurl{fixes/MISPRINTS_7I.json#140}\par

\par\medskip\noindent\hypertarget{app-F7I-VI-045.1}{}\textbf{F7I-VI-045.1}\quad p. 73; misprint\_pass.\par
\noindent Imprimé : \emph{(vide)}\par
\noindent Corrigé : {\ttfamily\footnotesize \textbackslash{}check\{\allowbreak{}}\par
\noindent Source : 7I:141. \nolinkurl{fixes/FIXLOG_7I_CURRENT.jsonl#671}; \nolinkurl{fixes/MISPRINTS_7I.json#141}\par

\par\medskip\noindent\hypertarget{app-F7I-VI-045.2}{}\textbf{F7I-VI-045.2}\quad p. 73; misprint\_pass.\par
\noindent Imprimé : \emph{(vide)}\par
\noindent Corrigé : {\ttfamily\footnotesize \}\allowbreak{}}\par
\noindent Source : 7I:141. \nolinkurl{fixes/FIXLOG_7I_CURRENT.jsonl#671}; \nolinkurl{fixes/MISPRINTS_7I.json#141}\par

\par\medskip\noindent\hypertarget{app-F7I-VI-045.3}{}\textbf{F7I-VI-045.3}\quad p. 73; misprint\_pass.\par
\noindent Imprimé : {\ttfamily\footnotesize \textbackslash{}\allowbreak{} }\par
\noindent Corrigé : {\ttfamily\footnotesize |\allowbreak{}}\par
\noindent Source : 7I:142. \nolinkurl{fixes/FIXLOG_7I_CURRENT.jsonl#672}; \nolinkurl{fixes/MISPRINTS_7I.json#142}\par

\par\medskip\noindent\hypertarget{app-F7I-VI-045.4}{}\textbf{F7I-VI-045.4}\quad p. 73; misprint\_pass.\par
\noindent Imprimé : {\ttfamily\footnotesize 2\allowbreak{}}\par
\noindent Corrigé : {\ttfamily\footnotesize 3\allowbreak{}}\par
\noindent Source : 7I:143. \nolinkurl{fixes/FIXLOG_7I_CURRENT.jsonl#673}; \nolinkurl{fixes/MISPRINTS_7I.json#143}\par

\par\medskip\noindent\hypertarget{app-F7I-VI-046.1}{}\textbf{F7I-VI-046.1}\quad p. 73; misprint\_pass.\par
\noindent Imprimé : {\ttfamily\footnotesize 2\allowbreak{}}\par
\noindent Corrigé : {\ttfamily\footnotesize 3\allowbreak{}}\par
\noindent Source : 7I:144. \nolinkurl{fixes/FIXLOG_7I_CURRENT.jsonl#674}; \nolinkurl{fixes/MISPRINTS_7I.json#144}\par

\par\medskip\noindent\hypertarget{app-F7I-VI-047.1}{}\textbf{F7I-VI-047.1}\quad p. 74; misprint\_pass.\par
\noindent Imprimé : {\ttfamily\footnotesize suffice\allowbreak{}}\par
\noindent Corrigé : {\ttfamily\footnotesize suffices\allowbreak{}}\par
\noindent Source : 7I:145. \nolinkurl{fixes/FIXLOG_7I_CURRENT.jsonl#675}; \nolinkurl{fixes/MISPRINTS_7I.json#145}\par

\par\medskip\noindent\hypertarget{app-F7I-VI-048.1}{}\textbf{F7I-VI-048.1}\quad p. 74; misprint\_pass.\par
\noindent Imprimé : {\ttfamily\footnotesize presheaf\allowbreak{}}\par
\noindent Corrigé : {\ttfamily\footnotesize presheaves\allowbreak{}}\par
\noindent Source : 7I:146. \nolinkurl{fixes/FIXLOG_7I_CURRENT.jsonl#676}; \nolinkurl{fixes/MISPRINTS_7I.json#146}\par

\par\medskip\noindent\hypertarget{app-F7I-VI-049.1}{}\textbf{F7I-VI-049.1}\quad p. 75; misprint\_pass.\par
\noindent Imprimé : {\ttfamily\footnotesize amits\allowbreak{}}\par
\noindent Corrigé : {\ttfamily\footnotesize admits\allowbreak{}}\par
\noindent Source : 7I:147. \nolinkurl{fixes/FIXLOG_7I_CURRENT.jsonl#677}; \nolinkurl{fixes/MISPRINTS_7I.json#147}\par

\par\medskip\noindent\hypertarget{app-F7I-VI-050.1}{}\textbf{F7I-VI-050.1}\quad p. 75; misprint\_pass.\par
\noindent Imprimé : {\ttfamily\footnotesize '\allowbreak{}}\par
\noindent Corrigé : {\ttfamily\footnotesize \textasciicircum{}\{1\}\allowbreak{}}\par
\noindent Source : 7I:148. \nolinkurl{fixes/FIXLOG_7I_CURRENT.jsonl#678}; \nolinkurl{fixes/MISPRINTS_7I.json#148}\par

\par\medskip\noindent\hypertarget{app-F7I-VI-050.2}{}\textbf{F7I-VI-050.2}\quad p. 75; misprint\_pass.\par
\noindent Imprimé : {\ttfamily\footnotesize '\allowbreak{}}\par
\noindent Corrigé : {\ttfamily\footnotesize \textasciicircum{}\{1\}\allowbreak{}}\par
\noindent Source : 7I:148. \nolinkurl{fixes/FIXLOG_7I_CURRENT.jsonl#678}; \nolinkurl{fixes/MISPRINTS_7I.json#148}\par

\par\medskip\noindent\hypertarget{app-F7I-VI-050.3}{}\textbf{F7I-VI-050.3}\quad p. 75; misprint\_pass.\par
\noindent Imprimé : \emph{(vide)}\par
\noindent Corrigé : {\ttfamily\footnotesize (P)\allowbreak{}}\par
\noindent Source : 7I:149. \nolinkurl{fixes/FIXLOG_7I_CURRENT.jsonl#679}; \nolinkurl{fixes/MISPRINTS_7I.json#149}\par

\par\medskip\noindent\hypertarget{app-F7I-VI-050.4}{}\textbf{F7I-VI-050.4}\quad p. 75; misprint\_pass.\par
\noindent Imprimé : \emph{(vide)}\par
\noindent Corrigé : {\ttfamily\footnotesize (P)\allowbreak{}}\par
\noindent Source : 7I:149. \nolinkurl{fixes/FIXLOG_7I_CURRENT.jsonl#679}; \nolinkurl{fixes/MISPRINTS_7I.json#149}\par

\par\medskip\noindent\hypertarget{app-F7I-VI-050.5}{}\textbf{F7I-VI-050.5}\quad p. 75; misprint\_pass.\par
\noindent Imprimé : \emph{(vide)}\par
\noindent Corrigé : {\ttfamily\footnotesize (P)\allowbreak{}}\par
\noindent Source : 7I:149. \nolinkurl{fixes/FIXLOG_7I_CURRENT.jsonl#679}; \nolinkurl{fixes/MISPRINTS_7I.json#149}\par

\par\medskip\noindent\hypertarget{app-F7I-VI-051.1}{}\textbf{F7I-VI-051.1}\quad p. 76; misprint\_pass.\par
\noindent Imprimé : \emph{(vide)}\par
\noindent Corrigé : {\ttfamily\footnotesize \textbackslash{}check\{\allowbreak{}}\par
\noindent Source : 7I:150. \nolinkurl{fixes/FIXLOG_7I_CURRENT.jsonl#680}; \nolinkurl{fixes/MISPRINTS_7I.json#150}\par

\par\medskip\noindent\hypertarget{app-F7I-VI-051.2}{}\textbf{F7I-VI-051.2}\quad p. 76; misprint\_pass.\par
\noindent Imprimé : \emph{(vide)}\par
\noindent Corrigé : {\ttfamily\footnotesize \}\allowbreak{}}\par
\noindent Source : 7I:150. \nolinkurl{fixes/FIXLOG_7I_CURRENT.jsonl#680}; \nolinkurl{fixes/MISPRINTS_7I.json#150}\par

\par\medskip\noindent\hypertarget{app-F7I-VI-052.1}{}\textbf{F7I-VI-052.1}\quad p. 76; printed misprint.\par
\noindent Imprimé : {\ttfamily\footnotesize <\allowbreak{}}\par
\noindent Corrigé : {\ttfamily\footnotesize \textbackslash{}leq\allowbreak{}}\par
\noindent Source : 7I:151. \nolinkurl{fixes/MISPRINTS_7I.json#151}\par

\par\medskip\noindent\hypertarget{app-F7I-VI-053.1}{}\textbf{F7I-VI-053.1}\quad p. 76; misprint\_in\_print.\par
\noindent Imprimé : {\ttfamily\footnotesize i\allowbreak{}}\par
\noindent Corrigé : {\ttfamily\footnotesize 1\allowbreak{}}\par
\noindent Source : FIX-7I-D-0066. \nolinkurl{fixes/FIXLOG_7I.jsonl#66}\par

\par\medskip\noindent\hypertarget{app-F7I-VI-054.1}{}\textbf{F7I-VI-054.1}\quad p. 77; misprint\_in\_print.\par
\noindent Imprimé : {\ttfamily\footnotesize \textbackslash{}otimes\allowbreak{}}\par
\noindent Corrigé : {\ttfamily\footnotesize \textbackslash{}times\allowbreak{}}\par
\noindent Source : FIX-7I-D-0067. \nolinkurl{fixes/FIXLOG_7I.jsonl#67}\par

\par\medskip\noindent\hypertarget{app-F7I-VI-055.1}{}\textbf{F7I-VI-055.1}\quad p. 77; printed misprint.\par
\noindent Imprimé : {\ttfamily\footnotesize SG\allowbreak{}}\par
\noindent Corrigé : {\ttfamily\footnotesize \textbackslash{}check\{S\}G\allowbreak{}}\par
\noindent Source : 7I:154. \nolinkurl{fixes/MISPRINTS_7I.json#154}\par

\par\medskip\noindent\hypertarget{app-F7I-VI-055.2}{}\textbf{F7I-VI-055.2}\quad p. 77; printed misprint.\par
\noindent Imprimé : \emph{(vide)}\par
\noindent Corrigé : {\ttfamily\footnotesize \textbackslash{}underline\{\allowbreak{}}\par
\noindent Source : 7I:155. \nolinkurl{fixes/MISPRINTS_7I.json#155}\par

\par\medskip\noindent\hypertarget{app-F7I-VI-055.3}{}\textbf{F7I-VI-055.3}\quad p. 77; printed misprint.\par
\noindent Imprimé : \emph{(vide)}\par
\noindent Corrigé : {\ttfamily\footnotesize \}\allowbreak{}}\par
\noindent Source : 7I:155. \nolinkurl{fixes/MISPRINTS_7I.json#155}\par

\par\medskip\noindent\hypertarget{app-F7I-VI-055.4}{}\textbf{F7I-VI-055.4}\quad p. 77; printed misprint.\par
\noindent Imprimé : \emph{(vide)}\par
\noindent Corrigé : {\ttfamily\footnotesize \textbackslash{}underline\{\allowbreak{}}\par
\noindent Source : 7I:155. \nolinkurl{fixes/MISPRINTS_7I.json#155}\par

\par\medskip\noindent\hypertarget{app-F7I-VI-055.5}{}\textbf{F7I-VI-055.5}\quad p. 77; printed misprint.\par
\noindent Imprimé : \emph{(vide)}\par
\noindent Corrigé : {\ttfamily\footnotesize \}\}\textasciicircum{}\{\textbackslash{}cdot\allowbreak{}}\par
\noindent Source : 7I:155. \nolinkurl{fixes/MISPRINTS_7I.json#155}\par

\par\medskip\noindent\hypertarget{app-F7I-VI-056.1}{}\textbf{F7I-VI-056.1}\quad p. 78; printed misprint.\par
\noindent Imprimé : {\ttfamily\footnotesize  a\allowbreak{}}\par
\noindent Corrigé : \emph{(vide)}\par
\noindent Source : 7I:156. \nolinkurl{fixes/MISPRINTS_7I.json#156}\par

\par\medskip\noindent\hypertarget{app-F7I-VI-057.1}{}\textbf{F7I-VI-057.1}\quad p. 79; printed misprint.\par
\noindent Imprimé : {\ttfamily\footnotesize \textbackslash{}\allowbreak{} \textbackslash{}\allowbreak{} \textbackslash{}\allowbreak{} }\par
\noindent Corrigé : {\ttfamily\footnotesize  <\allowbreak{} }\par
\noindent Source : 7I:157. \nolinkurl{fixes/MISPRINTS_7I.json#157}\par

\par\medskip\noindent\hypertarget{app-F7I-VI-057.2}{}\textbf{F7I-VI-057.2}\quad p. 79; printed misprint.\par
\noindent Imprimé : {\ttfamily\footnotesize \textbackslash{}\allowbreak{} \textbackslash{}\allowbreak{} \textbackslash{}\allowbreak{} }\par
\noindent Corrigé : {\ttfamily\footnotesize  <\allowbreak{} }\par
\noindent Source : 7I:157. \nolinkurl{fixes/MISPRINTS_7I.json#157}\par

\par\medskip\noindent\hypertarget{app-F7I-VI-057.3}{}\textbf{F7I-VI-057.3}\quad p. 79; printed misprint.\par
\noindent Imprimé : {\ttfamily\footnotesize \textbackslash{}phantom\{\allowbreak{}}\par
\noindent Corrigé : \emph{(vide)}\par
\noindent Source : 7I:158. \nolinkurl{fixes/MISPRINTS_7I.json#158}\par

\par\medskip\noindent\hypertarget{app-F7I-VI-057.4}{}\textbf{F7I-VI-057.4}\quad p. 79; designated-witness attribution.\par
\noindent Imprimé : {\ttfamily\footnotesize \}\allowbreak{}}\par
\noindent Corrigé : \emph{(vide)}\par
\noindent Source : WITNESS-F7I-VI-057.4. \nolinkurl{fr/sga7i/diplomatic/expose_VI_body.tex}; \nolinkurl{fr/sga7i/current/expose_VI_body.tex}\par

\par\medskip\noindent\hypertarget{app-F7I-VI-057.5}{}\textbf{F7I-VI-057.5}\quad p. 79; designated-witness attribution.\par
\noindent Imprimé : \emph{(vide)}\par
\noindent Corrigé : {\ttfamily\footnotesize \textasciicircum{}\{\textbackslash{}cdot\}\allowbreak{}}\par
\noindent Source : WITNESS-F7I-VI-057.5. \nolinkurl{fr/sga7i/diplomatic/expose_VI_body.tex}; \nolinkurl{fr/sga7i/current/expose_VI_body.tex}\par

\par\medskip\noindent\hypertarget{app-F7I-VI-057.6}{}\textbf{F7I-VI-057.6}\quad p. 79; printed misprint.\par
\noindent Imprimé : \emph{(vide)}\par
\noindent Corrigé : {\ttfamily\footnotesize \textbackslash{}widetilde\{\allowbreak{}}\par
\noindent Source : 7I:160. \nolinkurl{fixes/MISPRINTS_7I.json#160}\par

\par\medskip\noindent\hypertarget{app-F7I-VI-057.7}{}\textbf{F7I-VI-057.7}\quad p. 79; printed misprint.\par
\noindent Imprimé : \emph{(vide)}\par
\noindent Corrigé : {\ttfamily\footnotesize \}\allowbreak{}}\par
\noindent Source : 7I:160. \nolinkurl{fixes/MISPRINTS_7I.json#160}\par

\par\medskip\noindent\hypertarget{app-F7I-VI-057.8}{}\textbf{F7I-VI-057.8}\quad p. 79; printed misprint.\par
\noindent Imprimé : {\ttfamily\footnotesize 2\allowbreak{}}\par
\noindent Corrigé : {\ttfamily\footnotesize 3\allowbreak{}}\par
\noindent Source : 7I:161. \nolinkurl{fixes/MISPRINTS_7I.json#161}\par

\par\medskip\noindent\hypertarget{app-F7I-VI-058.1}{}\textbf{F7I-VI-058.1}\quad p. 80; misprint\_pass.\par
\noindent Imprimé : \emph{(vide)}\par
\noindent Corrigé : {\ttfamily\footnotesize  the\allowbreak{}}\par
\noindent Source : 7I:162. \nolinkurl{fixes/FIXLOG_7I_CURRENT.jsonl#682}; \nolinkurl{fixes/MISPRINTS_7I.json#162}\par

\par\medskip\noindent\hypertarget{app-F7I-VI-059.1}{}\textbf{F7I-VI-059.1}\quad p. 80; printed misprint.\par
\noindent Imprimé : {\ttfamily\footnotesize gives\allowbreak{}}\par
\noindent Corrigé : {\ttfamily\footnotesize give\allowbreak{}}\par
\noindent Source : 7I:163. \nolinkurl{fixes/MISPRINTS_7I.json#163}\par

\par\medskip\noindent\hypertarget{app-F7I-VI-060.1}{}\textbf{F7I-VI-060.1}\quad p. 81; misprint\_pass.\par
\noindent Imprimé : {\ttfamily\footnotesize gives\allowbreak{}}\par
\noindent Corrigé : {\ttfamily\footnotesize give\allowbreak{}}\par
\noindent Source : 7I:165. \nolinkurl{fixes/FIXLOG_7I_CURRENT.jsonl#684}; \nolinkurl{fixes/MISPRINTS_7I.json#165}\par

\par\medskip\noindent\hypertarget{app-F7I-VI-060.2}{}\textbf{F7I-VI-060.2}\quad p. 81; misprint\_pass.\par
\noindent Imprimé : {\ttfamily\footnotesize \textbackslash{}\allowbreak{} }\par
\noindent Corrigé : {\ttfamily\footnotesize |\allowbreak{}}\par
\noindent Source : 7I:164. \nolinkurl{fixes/FIXLOG_7I_CURRENT.jsonl#683}; \nolinkurl{fixes/MISPRINTS_7I.json#164}\par

\par\medskip\noindent\hypertarget{app-F7I-VI-060.3}{}\textbf{F7I-VI-060.3}\quad p. 81; misprint\_pass.\par
\noindent Imprimé : {\ttfamily\footnotesize \textbackslash{}\allowbreak{} }\par
\noindent Corrigé : {\ttfamily\footnotesize |\allowbreak{}}\par
\noindent Source : 7I:164. \nolinkurl{fixes/FIXLOG_7I_CURRENT.jsonl#683}; \nolinkurl{fixes/MISPRINTS_7I.json#164}\par

\par\medskip\noindent\hypertarget{app-F7I-VI-061.1}{}\textbf{F7I-VI-061.1}\quad p. 82; misprint\_pass.\par
\noindent Imprimé : {\ttfamily\footnotesize gives\allowbreak{}}\par
\noindent Corrigé : {\ttfamily\footnotesize give\allowbreak{}}\par
\noindent Source : 7I:166. \nolinkurl{fixes/FIXLOG_7I_CURRENT.jsonl#685}; \nolinkurl{fixes/MISPRINTS_7I.json#166}\par

\par\medskip\noindent\hypertarget{app-F7I-VI-062.1}{}\textbf{F7I-VI-062.1}\quad p. 83; misprint\_pass.\par
\noindent Imprimé : \emph{(vide)}\par
\noindent Corrigé : {\ttfamily\footnotesize \textbackslash{}overline\{\allowbreak{}}\par
\noindent Source : 7I:167. \nolinkurl{fixes/FIXLOG_7I_CURRENT.jsonl#686}; \nolinkurl{fixes/MISPRINTS_7I.json#167}\par

\par\medskip\noindent\hypertarget{app-F7I-VI-062.2}{}\textbf{F7I-VI-062.2}\quad p. 83; misprint\_pass.\par
\noindent Imprimé : \emph{(vide)}\par
\noindent Corrigé : {\ttfamily\footnotesize \}\allowbreak{}}\par
\noindent Source : 7I:167. \nolinkurl{fixes/FIXLOG_7I_CURRENT.jsonl#686}; \nolinkurl{fixes/MISPRINTS_7I.json#167}\par

\par\medskip\noindent\hypertarget{app-F7I-VI-062.3}{}\textbf{F7I-VI-062.3}\quad p. 83; misprint\_pass.\par
\noindent Imprimé : {\ttfamily\footnotesize admits\allowbreak{}}\par
\noindent Corrigé : {\ttfamily\footnotesize admit\allowbreak{}}\par
\noindent Source : 7I:168. \nolinkurl{fixes/FIXLOG_7I_CURRENT.jsonl#687}; \nolinkurl{fixes/MISPRINTS_7I.json#168}\par

\par\medskip\noindent\hypertarget{app-F7I-VI-063.1}{}\textbf{F7I-VI-063.1}\quad p. 84; misprint\_pass.\par
\noindent Imprimé : {\ttfamily\footnotesize \textasciicircum{}\allowbreak{}}\par
\noindent Corrigé : {\ttfamily\footnotesize \_\allowbreak{}}\par
\noindent Source : 7I:169. \nolinkurl{fixes/FIXLOG_7I_CURRENT.jsonl#688}; \nolinkurl{fixes/MISPRINTS_7I.json#169}\par

\par\medskip\noindent\hypertarget{app-F7I-VI-064.1}{}\textbf{F7I-VI-064.1}\quad p. 85; misprint\_in\_print.\par
\noindent Imprimé : {\ttfamily\footnotesize \textbackslash{}overset\allowbreak{}}\par
\noindent Corrigé : {\ttfamily\footnotesize \textbackslash{}xrightarrow\allowbreak{}}\par
\noindent Source : FIX-7I-D-0078. \nolinkurl{fixes/FIXLOG_7I.jsonl#78}\par

\par\medskip\noindent\hypertarget{app-F7I-VI-064.2}{}\textbf{F7I-VI-064.2}\quad p. 85; designated-witness attribution.\par
\noindent Imprimé : {\ttfamily\footnotesize \}\{\textbackslash{}relbar\textbackslash{}joinrel\textbackslash{}re\allowbreak{}lbar\textbackslash{}joinrel\textbackslash{}relbar\allowbreak{}}\par
\noindent Corrigé : \emph{(vide)}\par
\noindent Source : WITNESS-F7I-VI-064.2. \nolinkurl{fr/sga7i/diplomatic/expose_VI_body.tex}; \nolinkurl{fr/sga7i/current/expose_VI_body.tex}\par

\par\medskip\noindent\hypertarget{app-F7I-VI-064.3}{}\textbf{F7I-VI-064.3}\quad p. 85; misprint\_in\_print.\par
\noindent Imprimé : {\ttfamily\footnotesize \textbackslash{}overset\allowbreak{}}\par
\noindent Corrigé : {\ttfamily\footnotesize \textbackslash{}xrightarrow\allowbreak{}}\par
\noindent Source : FIX-7I-D-0079. \nolinkurl{fixes/FIXLOG_7I.jsonl#79}\par

\par\medskip\noindent\hypertarget{app-F7I-VI-064.4}{}\textbf{F7I-VI-064.4}\quad p. 85; designated-witness attribution.\par
\noindent Imprimé : {\ttfamily\footnotesize \}\{\textbackslash{}relbar\textbackslash{}joinrel\textbackslash{}re\allowbreak{}lbar\textbackslash{}joinrel\textbackslash{}relbar\allowbreak{}}\par
\noindent Corrigé : \emph{(vide)}\par
\noindent Source : WITNESS-F7I-VI-064.4. \nolinkurl{fr/sga7i/diplomatic/expose_VI_body.tex}; \nolinkurl{fr/sga7i/current/expose_VI_body.tex}\par

\par\medskip\noindent\hypertarget{app-F7I-VI-064.5}{}\textbf{F7I-VI-064.5}\quad p. 85; misprint\_pass.\par
\noindent Imprimé : {\ttfamily\footnotesize \textasciicircum{}\allowbreak{}}\par
\noindent Corrigé : {\ttfamily\footnotesize \_\allowbreak{}}\par
\noindent Source : 7I:174. \nolinkurl{fixes/FIXLOG_7I_CURRENT.jsonl#689}; \nolinkurl{fixes/MISPRINTS_7I.json#174}\par

\par\medskip\noindent\hypertarget{app-F7I-VI-065.1}{}\textbf{F7I-VI-065.1}\quad p. 86; misprint\_pass.\par
\noindent Imprimé : {\ttfamily\footnotesize \textbackslash{};\allowbreak{}}\par
\noindent Corrigé : {\ttfamily\footnotesize |\allowbreak{}}\par
\noindent Source : 7I:175. \nolinkurl{fixes/FIXLOG_7I_CURRENT.jsonl#690}; \nolinkurl{fixes/MISPRINTS_7I.json#175}\par

\par\medskip\noindent\hypertarget{app-F7I-VI-065.2}{}\textbf{F7I-VI-065.2}\quad p. 86; misprint\_pass.\par
\noindent Imprimé : {\ttfamily\footnotesize \textbackslash{};\allowbreak{}}\par
\noindent Corrigé : {\ttfamily\footnotesize |\allowbreak{}}\par
\noindent Source : 7I:175. \nolinkurl{fixes/FIXLOG_7I_CURRENT.jsonl#690}; \nolinkurl{fixes/MISPRINTS_7I.json#175}\par

\par\medskip\noindent\hypertarget{app-F7I-VI-065.3}{}\textbf{F7I-VI-065.3}\quad p. 86; misprint\_pass.\par
\noindent Imprimé : {\ttfamily\footnotesize \textbackslash{};\allowbreak{}}\par
\noindent Corrigé : {\ttfamily\footnotesize |\allowbreak{}}\par
\noindent Source : 7I:175. \nolinkurl{fixes/FIXLOG_7I_CURRENT.jsonl#690}; \nolinkurl{fixes/MISPRINTS_7I.json#175}\par

\par\medskip\noindent\hypertarget{app-F7I-VI-065.4}{}\textbf{F7I-VI-065.4}\quad p. 86; misprint\_pass.\par
\noindent Imprimé : {\ttfamily\footnotesize \textbackslash{};\allowbreak{}}\par
\noindent Corrigé : {\ttfamily\footnotesize |\allowbreak{}}\par
\noindent Source : 7I:175. \nolinkurl{fixes/FIXLOG_7I_CURRENT.jsonl#690}; \nolinkurl{fixes/MISPRINTS_7I.json#175}\par

\par\medskip\noindent\hypertarget{app-F7I-VI-065.5}{}\textbf{F7I-VI-065.5}\quad p. 86; misprint\_pass.\par
\noindent Imprimé : {\ttfamily\footnotesize \textbackslash{};\allowbreak{}}\par
\noindent Corrigé : {\ttfamily\footnotesize |\allowbreak{}}\par
\noindent Source : 7I:175. \nolinkurl{fixes/FIXLOG_7I_CURRENT.jsonl#690}; \nolinkurl{fixes/MISPRINTS_7I.json#175}\par

\par\medskip\noindent\hypertarget{app-F7I-VI-065.6}{}\textbf{F7I-VI-065.6}\quad p. 86; misprint\_pass.\par
\noindent Imprimé : {\ttfamily\footnotesize \textbackslash{};\allowbreak{}}\par
\noindent Corrigé : {\ttfamily\footnotesize |\allowbreak{}}\par
\noindent Source : 7I:175. \nolinkurl{fixes/FIXLOG_7I_CURRENT.jsonl#690}; \nolinkurl{fixes/MISPRINTS_7I.json#175}\par

\par\medskip\noindent\hypertarget{app-F7I-VI-065.7}{}\textbf{F7I-VI-065.7}\quad p. 86; misprint\_pass.\par
\noindent Imprimé : {\ttfamily\footnotesize \textbackslash{};\allowbreak{}}\par
\noindent Corrigé : {\ttfamily\footnotesize |\allowbreak{}}\par
\noindent Source : 7I:175. \nolinkurl{fixes/FIXLOG_7I_CURRENT.jsonl#690}; \nolinkurl{fixes/MISPRINTS_7I.json#175}\par

\par\medskip\noindent\hypertarget{app-F7I-VI-065.8}{}\textbf{F7I-VI-065.8}\quad p. 86; misprint\_pass.\par
\noindent Imprimé : {\ttfamily\footnotesize \textbackslash{};\allowbreak{}}\par
\noindent Corrigé : {\ttfamily\footnotesize |\allowbreak{}}\par
\noindent Source : 7I:175. \nolinkurl{fixes/FIXLOG_7I_CURRENT.jsonl#690}; \nolinkurl{fixes/MISPRINTS_7I.json#175}\par

\par\medskip\noindent\hypertarget{app-F7I-VI-065.9}{}\textbf{F7I-VI-065.9}\quad p. 86; misprint\_pass.\par
\noindent Imprimé : {\ttfamily\footnotesize \textbackslash{};\allowbreak{}}\par
\noindent Corrigé : {\ttfamily\footnotesize |\allowbreak{}}\par
\noindent Source : 7I:175. \nolinkurl{fixes/FIXLOG_7I_CURRENT.jsonl#690}; \nolinkurl{fixes/MISPRINTS_7I.json#175}\par

\par\medskip\noindent\hypertarget{app-F7I-VI-066.1}{}\textbf{F7I-VI-066.1}\quad p. 87; misprint\_pass.\par
\noindent Imprimé : \emph{(vide)}\par
\noindent Corrigé : {\ttfamily\footnotesize  an\allowbreak{}}\par
\noindent Source : 7I:176. \nolinkurl{fixes/FIXLOG_7I_CURRENT.jsonl#691}; \nolinkurl{fixes/MISPRINTS_7I.json#176}\par

\par\medskip\noindent\hypertarget{app-F7I-VI-067.1}{}\textbf{F7I-VI-067.1}\quad p. 87; misprint\_in\_print.\par
\noindent Imprimé : \emph{(vide)}\par
\noindent Corrigé : {\ttfamily\footnotesize )\allowbreak{}}\par
\noindent Source : FIX-7I-D-0081. \nolinkurl{fixes/FIXLOG_7I.jsonl#81}\par

\par\medskip\noindent\hypertarget{app-F7I-VI-068.1}{}\textbf{F7I-VI-068.1}\quad p. 89; misprint\_pass.\par
\noindent Imprimé : {\ttfamily\footnotesize extensions\allowbreak{}}\par
\noindent Corrigé : {\ttfamily\footnotesize extension\allowbreak{}}\par
\noindent Source : 7I:178. \nolinkurl{fixes/FIXLOG_7I_CURRENT.jsonl#692}; \nolinkurl{fixes/MISPRINTS_7I.json#178}\par

\par\medskip\noindent\hypertarget{app-F7I-VI-069.1}{}\textbf{F7I-VI-069.1}\quad p. 89; misprint\_pass.\par
\noindent Imprimé : {\ttfamily\footnotesize is\allowbreak{}}\par
\noindent Corrigé : {\ttfamily\footnotesize are\allowbreak{}}\par
\noindent Source : 7I:179. \nolinkurl{fixes/FIXLOG_7I_CURRENT.jsonl#693}; \nolinkurl{fixes/MISPRINTS_7I.json#179}\par

\par\medskip\noindent\hypertarget{app-F7I-VI-069.2}{}\textbf{F7I-VI-069.2}\quad p. 89; misprint\_pass.\par
\noindent Imprimé : {\ttfamily\footnotesize process\allowbreak{}}\par
\noindent Corrigé : {\ttfamily\footnotesize processes\allowbreak{}}\par
\noindent Source : 7I:180. \nolinkurl{fixes/FIXLOG_7I_CURRENT.jsonl#694}; \nolinkurl{fixes/MISPRINTS_7I.json#180}\par

\par\medskip\noindent\hypertarget{app-F7I-VI-069.3}{}\textbf{F7I-VI-069.3}\quad p. 89; designated-witness attribution.\par
\noindent Imprimé : \emph{(vide)}\par
\noindent Corrigé : {\ttfamily\footnotesize The\allowbreak{} }\par
\noindent Source : WITNESS-F7I-VI-069.3. \nolinkurl{fr/sga7i/diplomatic/expose_VI_body.tex}; \nolinkurl{fr/sga7i/current/expose_VI_body.tex}\par

\par\medskip\noindent\hypertarget{app-F7I-VI-070.1}{}\textbf{F7I-VI-070.1}\quad p. 90; designated-witness attribution.\par
\noindent Imprimé : {\ttfamily\footnotesize all\allowbreak{}}\par
\noindent Corrigé : {\ttfamily\footnotesize every\allowbreak{}}\par
\noindent Source : WITNESS-F7I-VI-070.1. \nolinkurl{fr/sga7i/diplomatic/expose_VI_body.tex}; \nolinkurl{fr/sga7i/current/expose_VI_body.tex}\par

\par\medskip\noindent\hypertarget{app-F7I-VI-070.2}{}\textbf{F7I-VI-070.2}\quad p. 90; inherited corrected-layer difference.\par
\noindent Imprimé : \emph{(vide)}\par
\noindent Corrigé : {\ttfamily\footnotesize an\allowbreak{} }\par
\noindent Source : ADOPT-F7I-VI-070.2. \nolinkurl{fr/sga7i/diplomatic/expose_VI_body.tex}; \nolinkurl{fr/sga7i/current/expose_VI_body.tex}\par

\par\medskip\noindent\hypertarget{app-F7I-VI-071.1}{}\textbf{F7I-VI-071.1}\quad p. 90; printed misprint.\par
\noindent Imprimé : {\ttfamily\footnotesize \textbackslash{};\allowbreak{}}\par
\noindent Corrigé : {\ttfamily\footnotesize |\allowbreak{}}\par
\noindent Source : 7I:183. \nolinkurl{fixes/MISPRINTS_7I.json#183}\par

\par\medskip\noindent\hypertarget{app-F7I-VI-071.2}{}\textbf{F7I-VI-071.2}\quad p. 91; inherited corrected-layer difference.\par
\noindent Imprimé : \emph{(vide)}\par
\noindent Corrigé : {\ttfamily\footnotesize \textbackslash{}hat\{\allowbreak{}}\par
\noindent Source : ADOPT-F7I-VI-071.2. \nolinkurl{fr/sga7i/diplomatic/expose_VI_body.tex}; \nolinkurl{fr/sga7i/current/expose_VI_body.tex}\par

\par\medskip\noindent\hypertarget{app-F7I-VI-071.3}{}\textbf{F7I-VI-071.3}\quad p. 91; inherited corrected-layer difference.\par
\noindent Imprimé : \emph{(vide)}\par
\noindent Corrigé : {\ttfamily\footnotesize \}\allowbreak{}}\par
\noindent Source : ADOPT-F7I-VI-071.3. \nolinkurl{fr/sga7i/diplomatic/expose_VI_body.tex}; \nolinkurl{fr/sga7i/current/expose_VI_body.tex}\par

\par\medskip\noindent\hypertarget{app-F7I-VI-071.4}{}\textbf{F7I-VI-071.4}\quad p. 91; inherited corrected-layer difference.\par
\noindent Imprimé : \emph{(vide)}\par
\noindent Corrigé : {\ttfamily\footnotesize  whose\allowbreak{} restriction\allowbreak{} to\allowbreak{} \$A\$\allowbreak{} splits\allowbreak{}}\par
\noindent Source : ADOPT-F7I-VI-071.4. \nolinkurl{fr/sga7i/diplomatic/expose_VI_body.tex}; \nolinkurl{fr/sga7i/current/expose_VI_body.tex}\par

\par\medskip\noindent\hypertarget{app-F7I-VI-071.5}{}\textbf{F7I-VI-071.5}\quad p. 91; inherited corrected-layer difference.\par
\noindent Imprimé : {\ttfamily\footnotesize such\allowbreak{}}\par
\noindent Corrigé : {\ttfamily\footnotesize lifting\allowbreak{}}\par
\noindent Source : ADOPT-F7I-VI-071.5. \nolinkurl{fr/sga7i/diplomatic/expose_VI_body.tex}; \nolinkurl{fr/sga7i/current/expose_VI_body.tex}\par

\par\medskip\noindent\hypertarget{app-F7I-VI-071.6}{}\textbf{F7I-VI-071.6}\quad p. 91; inherited corrected-layer difference.\par
\noindent Imprimé : {\ttfamily\footnotesize that\allowbreak{} }\par
\noindent Corrigé : {\ttfamily\footnotesize the\allowbreak{} canonical\allowbreak{} map\allowbreak{} }\par
\noindent Source : ADOPT-F7I-VI-071.6. \nolinkurl{fr/sga7i/diplomatic/expose_VI_body.tex}; \nolinkurl{fr/sga7i/current/expose_VI_body.tex}\par

\par\medskip\noindent\hypertarget{app-F7I-VI-071.7}{}\textbf{F7I-VI-071.7}\quad p. 91; inherited corrected-layer difference.\par
\noindent Imprimé : {\ttfamily\footnotesize \textbackslash{}varphi\allowbreak{}}\par
\noindent Corrigé : {\ttfamily\footnotesize A\allowbreak{}}\par
\noindent Source : ADOPT-F7I-VI-071.7. \nolinkurl{fr/sga7i/diplomatic/expose_VI_body.tex}; \nolinkurl{fr/sga7i/current/expose_VI_body.tex}\par

\par\medskip\noindent\hypertarget{app-F7I-VI-071.8}{}\textbf{F7I-VI-071.8}\quad p. 91; inherited corrected-layer difference.\par
\noindent Imprimé : {\ttfamily\footnotesize \textbackslash{}circ\allowbreak{}}\par
\noindent Corrigé : {\ttfamily\footnotesize \textbackslash{}to\allowbreak{}}\par
\noindent Source : ADOPT-F7I-VI-071.8. \nolinkurl{fr/sga7i/diplomatic/expose_VI_body.tex}; \nolinkurl{fr/sga7i/current/expose_VI_body.tex}\par

\par\medskip\noindent\hypertarget{app-F7I-VI-071.9}{}\textbf{F7I-VI-071.9}\quad p. 91; inherited corrected-layer difference.\par
\noindent Imprimé : {\ttfamily\footnotesize \textbackslash{}psi\allowbreak{} =\allowbreak{} \textbackslash{}mathrm\allowbreak{}}\par
\noindent Corrigé : {\ttfamily\footnotesize \textbackslash{}hat\allowbreak{}}\par
\noindent Source : ADOPT-F7I-VI-071.9. \nolinkurl{fr/sga7i/diplomatic/expose_VI_body.tex}; \nolinkurl{fr/sga7i/current/expose_VI_body.tex}\par

\par\medskip\noindent\hypertarget{app-F7I-VI-071.10}{}\textbf{F7I-VI-071.10}\quad p. 91; inherited corrected-layer difference.\par
\noindent Imprimé : {\ttfamily\footnotesize id\allowbreak{}}\par
\noindent Corrigé : {\ttfamily\footnotesize A\allowbreak{}}\par
\noindent Source : ADOPT-F7I-VI-071.10. \nolinkurl{fr/sga7i/diplomatic/expose_VI_body.tex}; \nolinkurl{fr/sga7i/current/expose_VI_body.tex}\par

\par\medskip\noindent\hypertarget{app-F7I-VI-071.11}{}\textbf{F7I-VI-071.11}\quad p. 91; inherited corrected-layer difference.\par
\noindent Imprimé : {\ttfamily\footnotesize \_A\allowbreak{}}\par
\noindent Corrigé : \emph{(vide)}\par
\noindent Source : ADOPT-F7I-VI-071.11. \nolinkurl{fr/sga7i/diplomatic/expose_VI_body.tex}; \nolinkurl{fr/sga7i/current/expose_VI_body.tex}\par

\par\medskip\noindent\hypertarget{app-F7I-VI-071.12}{}\textbf{F7I-VI-071.12}\quad p. 91; printed misprint.\par
\noindent Imprimé : \emph{(vide)}\par
\noindent Corrigé : {\ttfamily\footnotesize \textbackslash{}hat\{\allowbreak{}}\par
\noindent Source : 7I:190. \nolinkurl{fixes/MISPRINTS_7I.json#190}\par

\par\medskip\noindent\hypertarget{app-F7I-VI-071.13}{}\textbf{F7I-VI-071.13}\quad p. 91; printed misprint.\par
\noindent Imprimé : \emph{(vide)}\par
\noindent Corrigé : {\ttfamily\footnotesize \}\allowbreak{}}\par
\noindent Source : 7I:190. \nolinkurl{fixes/MISPRINTS_7I.json#190}\par

\par\medskip\noindent\hypertarget{app-F7I-VI-071.14}{}\textbf{F7I-VI-071.14}\quad p. 91; printed misprint.\par
\noindent Imprimé : \emph{(vide)}\par
\noindent Corrigé : {\ttfamily\footnotesize \textbackslash{}hat\{\allowbreak{}}\par
\noindent Source : 7I:191. \nolinkurl{fixes/MISPRINTS_7I.json#191}\par

\par\medskip\noindent\hypertarget{app-F7I-VI-071.15}{}\textbf{F7I-VI-071.15}\quad p. 91; printed misprint.\par
\noindent Imprimé : \emph{(vide)}\par
\noindent Corrigé : {\ttfamily\footnotesize \}\allowbreak{}}\par
\noindent Source : 7I:191. \nolinkurl{fixes/MISPRINTS_7I.json#191}\par

\par\medskip\noindent\hypertarget{app-F7I-VI-072.1}{}\textbf{F7I-VI-072.1}\quad p. 91; designated-witness attribution.\par
\noindent Imprimé : {\ttfamily\footnotesize  \$D\textasciicircum{}1\$\allowbreak{}}\par
\noindent Corrigé : \emph{(vide)}\par
\noindent Source : WITNESS-F7I-VI-072.1. \nolinkurl{fr/sga7i/diplomatic/expose_VI_body.tex}; \nolinkurl{fr/sga7i/current/expose_VI_body.tex}\par

\par\medskip\noindent\hypertarget{app-F7I-VI-073.1}{}\textbf{F7I-VI-073.1}\quad p. 92; designated-witness attribution.\par
\noindent Imprimé : \emph{(vide)}\par
\noindent Corrigé : {\ttfamily\footnotesize  a\allowbreak{}}\par
\noindent Source : WITNESS-F7I-VI-073.1. \nolinkurl{fr/sga7i/diplomatic/expose_VI_body.tex}; \nolinkurl{fr/sga7i/current/expose_VI_body.tex}\par

\par\medskip\noindent\hypertarget{app-F7I-VI-073.2}{}\textbf{F7I-VI-073.2}\quad p. 92; printed misprint.\par
\noindent Imprimé : {\ttfamily\footnotesize M\allowbreak{}}\par
\noindent Corrigé : {\ttfamily\footnotesize E\allowbreak{}}\par
\noindent Source : 7I:193. \nolinkurl{fixes/MISPRINTS_7I.json#193}\par

\par\medskip\noindent\hypertarget{app-F7I-VI-073.3}{}\textbf{F7I-VI-073.3}\quad p. 92; printed misprint.\par
\noindent Imprimé : {\ttfamily\footnotesize \textbackslash{}alpha\allowbreak{}}\par
\noindent Corrigé : {\ttfamily\footnotesize A\allowbreak{}}\par
\noindent Source : 7I:194. \nolinkurl{fixes/MISPRINTS_7I.json#194}\par

\par\medskip\noindent\hypertarget{app-F7I-VI-073.4}{}\textbf{F7I-VI-073.4}\quad p. 92; inherited corrected-layer difference.\par
\noindent Imprimé : \emph{(vide)}\par
\noindent Corrigé : {\ttfamily\footnotesize  to\allowbreak{} say\allowbreak{} that\allowbreak{}}\par
\noindent Source : ADOPT-F7I-VI-073.4. \nolinkurl{fr/sga7i/diplomatic/expose_VI_body.tex}; \nolinkurl{fr/sga7i/current/expose_VI_body.tex}\par

\par\medskip\noindent\hypertarget{app-F7I-VI-073.5}{}\textbf{F7I-VI-073.5}\quad p. 92; inherited corrected-layer difference.\par
\noindent Imprimé : \emph{(vide)}\par
\noindent Corrigé : {\ttfamily\footnotesize  as\allowbreak{} a\allowbreak{}}\par
\noindent Source : ADOPT-F7I-VI-073.5. \nolinkurl{fr/sga7i/diplomatic/expose_VI_body.tex}; \nolinkurl{fr/sga7i/current/expose_VI_body.tex}\par

\par\medskip\noindent\hypertarget{app-F7I-VI-074.1}{}\textbf{F7I-VI-074.1}\quad p. 92; printed misprint.\par
\noindent Imprimé : {\ttfamily\footnotesize a\allowbreak{}}\par
\noindent Corrigé : {\ttfamily\footnotesize c\allowbreak{}}\par
\noindent Source : 7I:195. \nolinkurl{fixes/MISPRINTS_7I.json#195}\par

\par\medskip\noindent\hypertarget{app-F7I-VI-074.2}{}\textbf{F7I-VI-074.2}\quad p. 92; printed misprint.\par
\noindent Imprimé : {\ttfamily\footnotesize abreviate\allowbreak{}}\par
\noindent Corrigé : {\ttfamily\footnotesize abbreviate\allowbreak{}}\par
\noindent Source : 7I:196. \nolinkurl{fixes/MISPRINTS_7I.json#196}\par

\par\medskip\noindent\hypertarget{app-F7I-VI-074.3}{}\textbf{F7I-VI-074.3}\quad p. 92; printed misprint.\par
\noindent Imprimé : {\ttfamily\footnotesize is\allowbreak{}}\par
\noindent Corrigé : {\ttfamily\footnotesize are\allowbreak{}}\par
\noindent Source : 7I:197. \nolinkurl{fixes/MISPRINTS_7I.json#197}\par

\par\medskip\noindent\hypertarget{app-F7I-VI-075.1}{}\textbf{F7I-VI-075.1}\quad p. 93; printed misprint.\par
\noindent Imprimé : {\ttfamily\footnotesize h\allowbreak{}}\par
\noindent Corrigé : {\ttfamily\footnotesize k\allowbreak{}}\par
\noindent Source : 7I:198. \nolinkurl{fixes/MISPRINTS_7I.json#198}\par

\par\medskip\noindent\hypertarget{app-F7I-VI-075.2}{}\textbf{F7I-VI-075.2}\quad p. 93; designated-witness attribution.\par
\noindent Imprimé : {\ttfamily\footnotesize 6\allowbreak{}}\par
\noindent Corrigé : {\ttfamily\footnotesize 7\allowbreak{}}\par
\noindent Source : WITNESS-F7I-VI-075.2. \nolinkurl{fr/sga7i/diplomatic/expose_VI_body.tex}; \nolinkurl{fr/sga7i/current/expose_VI_body.tex}\par

\par\medskip\noindent\hypertarget{app-F7I-VI-075.3}{}\textbf{F7I-VI-075.3}\quad p. 93; designated-witness attribution.\par
\noindent Imprimé : \emph{(vide)}\par
\noindent Corrigé : {\ttfamily\footnotesize the\allowbreak{} vanishing\allowbreak{} for\allowbreak{} smooth\allowbreak{} algebras\allowbreak{} }\par
\noindent Source : WITNESS-F7I-VI-075.3. \nolinkurl{fr/sga7i/diplomatic/expose_VI_body.tex}; \nolinkurl{fr/sga7i/current/expose_VI_body.tex}\par

\par\medskip\noindent\hypertarget{app-F7I-VI-075.4}{}\textbf{F7I-VI-075.4}\quad p. 93; designated-witness attribution.\par
\noindent Imprimé : {\ttfamily\footnotesize 8\allowbreak{}}\par
\noindent Corrigé : {\ttfamily\footnotesize cf.\allowbreak{} 3.13(a\allowbreak{}}\par
\noindent Source : WITNESS-F7I-VI-075.4. \nolinkurl{fr/sga7i/diplomatic/expose_VI_body.tex}; \nolinkurl{fr/sga7i/current/expose_VI_body.tex}\par

\par\medskip\noindent\hypertarget{app-F7I-VI-075.5}{}\textbf{F7I-VI-075.5}\quad p. 93; designated-witness attribution.\par
\noindent Imprimé : \emph{(vide)}\par
\noindent Corrigé : {\ttfamily\footnotesize ),\allowbreak{}}\par
\noindent Source : WITNESS-F7I-VI-075.5. \nolinkurl{fr/sga7i/diplomatic/expose_VI_body.tex}; \nolinkurl{fr/sga7i/current/expose_VI_body.tex}\par

\par\medskip\noindent\hypertarget{app-F7I-VI-075.6}{}\textbf{F7I-VI-075.6}\quad p. 93; designated-witness attribution.\par
\noindent Imprimé : \emph{(vide)}\par
\noindent Corrigé : {\ttfamily\footnotesize the\allowbreak{} complete-intersectio\allowbreak{}n\allowbreak{} vanishing\allowbreak{} (cf.\allowbreak{} }\par
\noindent Source : WITNESS-F7I-VI-075.6. \nolinkurl{fr/sga7i/diplomatic/expose_VI_body.tex}; \nolinkurl{fr/sga7i/current/expose_VI_body.tex}\par

\par\medskip\noindent\hypertarget{app-F7I-VI-075.7}{}\textbf{F7I-VI-075.7}\quad p. 93; printed misprint.\par
\noindent Imprimé : {\ttfamily\footnotesize 12\allowbreak{}}\par
\noindent Corrigé : {\ttfamily\footnotesize 13\allowbreak{}}\par
\noindent Source : 7I:200. \nolinkurl{fixes/MISPRINTS_7I.json#200}\par

\par\medskip\noindent\hypertarget{app-F7I-VI-075.8}{}\textbf{F7I-VI-075.8}\quad p. 93; designated-witness attribution.\par
\noindent Imprimé : {\ttfamily\footnotesize 10\allowbreak{}}\par
\noindent Corrigé : {\ttfamily\footnotesize b)\allowbreak{}}\par
\noindent Source : WITNESS-F7I-VI-075.8. \nolinkurl{fr/sga7i/diplomatic/expose_VI_body.tex}; \nolinkurl{fr/sga7i/current/expose_VI_body.tex}\par

\par\medskip\noindent\hypertarget{app-F7I-VI-076.1}{}\textbf{F7I-VI-076.1}\quad p. 93; printed misprint.\par
\noindent Imprimé : {\ttfamily\footnotesize 7\allowbreak{}}\par
\noindent Corrigé : {\ttfamily\footnotesize 8\allowbreak{}}\par
\noindent Source : 7I:201. \nolinkurl{fixes/MISPRINTS_7I.json#201}\par

\par\medskip\noindent\hypertarget{app-F7I-VI-076.2}{}\textbf{F7I-VI-076.2}\quad p. 94; designated-witness attribution.\par
\noindent Imprimé : {\ttfamily\footnotesize all\allowbreak{}}\par
\noindent Corrigé : {\ttfamily\footnotesize every\allowbreak{}}\par
\noindent Source : WITNESS-F7I-VI-076.2. \nolinkurl{fr/sga7i/diplomatic/expose_VI_body.tex}; \nolinkurl{fr/sga7i/current/expose_VI_body.tex}\par

\par\medskip\noindent\hypertarget{app-F7I-VI-077.1}{}\textbf{F7I-VI-077.1}\quad p. 94; inherited corrected-layer difference.\par
\noindent Imprimé : {\ttfamily\footnotesize Andre\allowbreak{}}\par
\noindent Corrigé : {\ttfamily\footnotesize André\allowbreak{}}\par
\noindent Source : ADOPT-F7I-VI-077.1. \nolinkurl{fr/sga7i/diplomatic/expose_VI_body.tex}; \nolinkurl{fr/sga7i/current/expose_VI_body.tex}\par

\par\medskip\noindent\hypertarget{app-F7I-VI-078.1}{}\textbf{F7I-VI-078.1}\quad p. 95; misprint\_in\_print.\par
\noindent Imprimé : {\ttfamily\footnotesize spec\allowbreak{}}\par
\noindent Corrigé : {\ttfamily\footnotesize Spec\allowbreak{}}\par
\noindent Source : FIX-7I-D-0103. \nolinkurl{fixes/FIXLOG_7I.jsonl#103}\par

\par\medskip\noindent\hypertarget{app-F7I-VI-079.1}{}\textbf{F7I-VI-079.1}\quad p. 95; inherited corrected-layer difference.\par
\noindent Imprimé : {\ttfamily\footnotesize p\allowbreak{}}\par
\noindent Corrigé : {\ttfamily\footnotesize P\allowbreak{}}\par
\noindent Source : ADOPT-F7I-VI-079.1. \nolinkurl{fr/sga7i/diplomatic/expose_VI_body.tex}; \nolinkurl{fr/sga7i/current/expose_VI_body.tex}\par

\par\medskip\noindent\hypertarget{app-F7I-VI-080.1}{}\textbf{F7I-VI-080.1}\quad p. 98; printed misprint.\par
\noindent Imprimé : \emph{(vide)}\par
\noindent Corrigé : {\ttfamily\footnotesize '\allowbreak{}}\par
\noindent Source : 7I:184. \nolinkurl{fixes/MISPRINTS_7I.json#184}\par

\par\medskip\noindent\hypertarget{app-F7I-VI-081.1}{}\textbf{F7I-VI-081.1}\quad p. 98; printed misprint.\par
\noindent Imprimé : {\ttfamily\footnotesize \textbackslash{}phantom\{\allowbreak{}}\par
\noindent Corrigé : \emph{(vide)}\par
\noindent Source : 7I:185. \nolinkurl{fixes/MISPRINTS_7I.json#185}\par

\par\medskip\noindent\hypertarget{app-F7I-VI-081.2}{}\textbf{F7I-VI-081.2}\quad p. 98; printed misprint.\par
\noindent Imprimé : {\ttfamily\footnotesize \}\allowbreak{}}\par
\noindent Corrigé : \emph{(vide)}\par
\noindent Source : 7I:185. \nolinkurl{fixes/MISPRINTS_7I.json#185}\par

\par\medskip\noindent\hypertarget{app-F7I-VI-082.1}{}\textbf{F7I-VI-082.1}\quad p. 99; misprint\_in\_print.\par
\noindent Imprimé : {\ttfamily\footnotesize \textbackslash{}relbar\textbackslash{}joinrel\textbackslash{}relb\allowbreak{}ar\textbackslash{}joinrel\textbackslash{}relbar\allowbreak{}}\par
\noindent Corrigé : {\ttfamily\footnotesize \textbackslash{}longrightarrow\allowbreak{}}\par
\noindent Source : FIX-7I-D-0084. \nolinkurl{fixes/FIXLOG_7I.jsonl#84}\par

\par\medskip\noindent\hypertarget{app-F7I-VI-082.2}{}\textbf{F7I-VI-082.2}\quad p. 99; misprint\_in\_print.\par
\noindent Imprimé : {\ttfamily\footnotesize \textbackslash{}relbar\textbackslash{}joinrel\textbackslash{}relb\allowbreak{}ar\textbackslash{}joinrel\textbackslash{}relbar\allowbreak{}}\par
\noindent Corrigé : {\ttfamily\footnotesize \textbackslash{}longrightarrow\allowbreak{}}\par
\noindent Source : FIX-7I-D-0085. \nolinkurl{fixes/FIXLOG_7I.jsonl#85}\par

\par\medskip\noindent\hypertarget{app-F7I-VI-082.3}{}\textbf{F7I-VI-082.3}\quad p. 99; misprint\_in\_print.\par
\noindent Imprimé : {\ttfamily\footnotesize \textbackslash{}relbar\textbackslash{}joinrel\textbackslash{}relb\allowbreak{}ar\textbackslash{}joinrel\textbackslash{}relbar\allowbreak{}}\par
\noindent Corrigé : {\ttfamily\footnotesize \textbackslash{}longrightarrow\allowbreak{}}\par
\noindent Source : FIX-7I-D-0085. \nolinkurl{fixes/FIXLOG_7I.jsonl#85}\par

\par\medskip\noindent\hypertarget{app-F7I-VI-083.1}{}\textbf{F7I-VI-083.1}\quad p. 99; misprint\_in\_print.\par
\noindent Imprimé : {\ttfamily\footnotesize \textbackslash{}relbar\textbackslash{}joinrel\textbackslash{}relb\allowbreak{}ar\textbackslash{}joinrel\textbackslash{}relbar\allowbreak{}}\par
\noindent Corrigé : {\ttfamily\footnotesize \textbackslash{}longrightarrow\allowbreak{}}\par
\noindent Source : FIX-7I-D-0087. \nolinkurl{fixes/FIXLOG_7I.jsonl#87}\par

\par\medskip\noindent\hypertarget{app-F7I-VI-084.1}{}\textbf{F7I-VI-084.1}\quad p. 100; misprint\_in\_print.\par
\noindent Imprimé : {\ttfamily\footnotesize \textbackslash{}relbar\textbackslash{}joinrel\textbackslash{}relb\allowbreak{}ar\textbackslash{}joinrel\textbackslash{}relbar\allowbreak{}}\par
\noindent Corrigé : {\ttfamily\footnotesize \textbackslash{}longrightarrow\allowbreak{}}\par
\noindent Source : FIX-7I-D-0088. \nolinkurl{fixes/FIXLOG_7I.jsonl#88}\par

\par\medskip\noindent\hypertarget{app-F7I-VI-085.1}{}\textbf{F7I-VI-085.1}\quad p. 101; printed misprint.\par
\noindent Imprimé : \emph{(vide)}\par
\noindent Corrigé : {\ttfamily\footnotesize \_\{X\_\{0\}\}\allowbreak{}}\par
\noindent Source : 7I:204. \nolinkurl{fixes/MISPRINTS_7I.json#204}\par

\par\medskip\noindent\hypertarget{app-F7I-VI-086.1}{}\textbf{F7I-VI-086.1}\quad p. 101; inherited corrected-layer difference.\par
\noindent Imprimé : {\ttfamily\footnotesize 2\allowbreak{}}\par
\noindent Corrigé : {\ttfamily\footnotesize 3\allowbreak{}}\par
\noindent Source : ADOPT-F7I-VI-086.1. \nolinkurl{fr/sga7i/diplomatic/expose_VI_body.tex}; \nolinkurl{fr/sga7i/current/expose_VI_body.tex}\par

\par\medskip\noindent\hypertarget{app-F7I-VI-086.2}{}\textbf{F7I-VI-086.2}\quad p. 101; inherited corrected-layer difference.\par
\noindent Imprimé : {\ttfamily\footnotesize 13\allowbreak{}}\par
\noindent Corrigé : {\ttfamily\footnotesize 15\allowbreak{}}\par
\noindent Source : ADOPT-F7I-VI-086.2. \nolinkurl{fr/sga7i/diplomatic/expose_VI_body.tex}; \nolinkurl{fr/sga7i/current/expose_VI_body.tex}\par

\par\medskip\noindent\hypertarget{app-F7I-VI-086.3}{}\textbf{F7I-VI-086.3}\quad p. 101; inherited corrected-layer difference.\par
\noindent Imprimé : {\ttfamily\footnotesize 2\allowbreak{}}\par
\noindent Corrigé : {\ttfamily\footnotesize 3\allowbreak{}}\par
\noindent Source : ADOPT-F7I-VI-086.3. \nolinkurl{fr/sga7i/diplomatic/expose_VI_body.tex}; \nolinkurl{fr/sga7i/current/expose_VI_body.tex}\par

\par\medskip\noindent\hypertarget{app-F7I-VI-086.4}{}\textbf{F7I-VI-086.4}\quad p. 101; inherited corrected-layer difference.\par
\noindent Imprimé : {\ttfamily\footnotesize 13\allowbreak{}}\par
\noindent Corrigé : {\ttfamily\footnotesize 15\allowbreak{}}\par
\noindent Source : ADOPT-F7I-VI-086.4. \nolinkurl{fr/sga7i/diplomatic/expose_VI_body.tex}; \nolinkurl{fr/sga7i/current/expose_VI_body.tex}\par

\par\medskip\noindent\hypertarget{app-F7I-VI-087.1}{}\textbf{F7I-VI-087.1}\quad p. 102; designated-witness attribution.\par
\noindent Imprimé : {\ttfamily\footnotesize no\allowbreak{} head,\allowbreak{} }\par
\noindent Corrigé : \emph{(vide)}\par
\noindent Source : WITNESS-F7I-VI-087.1. \nolinkurl{fr/sga7i/diplomatic/expose_VI_body.tex}; \nolinkurl{fr/sga7i/current/expose_VI_body.tex}\par

\par\medskip\noindent\hypertarget{app-F7I-VI-088.1}{}\textbf{F7I-VI-088.1}\quad p. 103; misprint\_pass.\par
\noindent Imprimé : \emph{(vide)}\par
\noindent Corrigé : {\ttfamily\footnotesize \_0\allowbreak{}}\par
\noindent Source : 7I:206. \nolinkurl{fixes/FIXLOG_7I_CURRENT.jsonl#695}; \nolinkurl{fixes/MISPRINTS_7I.json#206}\par

\par\medskip\noindent\hypertarget{app-F7I-VI-089.1}{}\textbf{F7I-VI-089.1}\quad p. 103; misprint\_in\_print.\par
\noindent Imprimé : {\ttfamily\footnotesize \textbackslash{}mathfrak\{J\}\allowbreak{}}\par
\noindent Corrigé : {\ttfamily\footnotesize j\allowbreak{}}\par
\noindent Source : FIX-7I-D-0107. \nolinkurl{fixes/FIXLOG_7I.jsonl#107}\par

\par\medskip\noindent\hypertarget{app-F7I-VI-090.1}{}\textbf{F7I-VI-090.1}\quad p. 104; misprint\_pass.\par
\noindent Imprimé : \emph{(vide)}\par
\noindent Corrigé : {\ttfamily\footnotesize \_0\allowbreak{}}\par
\noindent Source : 7I:208. \nolinkurl{fixes/FIXLOG_7I_CURRENT.jsonl#696}; \nolinkurl{fixes/MISPRINTS_7I.json#208}\par

\par\medskip\noindent\hypertarget{app-F7I-VI-091.1}{}\textbf{F7I-VI-091.1}\quad p. 104; designated-witness attribution.\par
\noindent Imprimé : \emph{(vide)}\par
\noindent Corrigé : {\ttfamily\footnotesize (\allowbreak{}}\par
\noindent Source : WITNESS-F7I-VI-091.1. \nolinkurl{fr/sga7i/diplomatic/expose_VI_body.tex}; \nolinkurl{fr/sga7i/current/expose_VI_body.tex}\par

\par\medskip\noindent\hypertarget{app-F7I-VI-091.2}{}\textbf{F7I-VI-091.2}\quad p. 104; misprint\_pass.\par
\noindent Imprimé : {\ttfamily\footnotesize ,\textbackslash{},\allowbreak{}}\par
\noindent Corrigé : \emph{(vide)}\par
\noindent Source : 7I:209. \nolinkurl{fixes/FIXLOG_7I_CURRENT.jsonl#697}; \nolinkurl{fixes/MISPRINTS_7I.json#209}\par

\par\medskip\noindent\hypertarget{app-F7I-VI-092.1}{}\textbf{F7I-VI-092.1}\quad p. 104; inherited corrected-layer difference.\par
\noindent Imprimé : {\ttfamily\footnotesize Firstly\allowbreak{}}\par
\noindent Corrigé : {\ttfamily\footnotesize First,\allowbreak{}}\par
\noindent Source : ADOPT-F7I-VI-092.1. \nolinkurl{fr/sga7i/diplomatic/expose_VI_body.tex}; \nolinkurl{fr/sga7i/current/expose_VI_body.tex}\par

\par\medskip\noindent\hypertarget{app-F7I-VI-092.2}{}\textbf{F7I-VI-092.2}\quad p. 104; inherited corrected-layer difference.\par
\noindent Imprimé : {\ttfamily\footnotesize on\allowbreak{} a\allowbreak{}}\par
\noindent Corrigé : {\ttfamily\footnotesize some\allowbreak{}}\par
\noindent Source : ADOPT-F7I-VI-092.2. \nolinkurl{fr/sga7i/diplomatic/expose_VI_body.tex}; \nolinkurl{fr/sga7i/current/expose_VI_body.tex}\par

\par\medskip\noindent\hypertarget{app-F7I-VI-093.1}{}\textbf{F7I-VI-093.1}\quad p. 105; inherited corrected-layer difference.\par
\noindent Imprimé : {\ttfamily\footnotesize it\allowbreak{} }\par
\noindent Corrigé : \emph{(vide)}\par
\noindent Source : ADOPT-F7I-VI-093.1. \nolinkurl{fr/sga7i/diplomatic/expose_VI_body.tex}; \nolinkurl{fr/sga7i/current/expose_VI_body.tex}\par

\par\medskip\noindent\hypertarget{app-F7I-VI-093.2}{}\textbf{F7I-VI-093.2}\quad p. 105; inherited corrected-layer difference.\par
\noindent Imprimé : \emph{(vide)}\par
\noindent Corrigé : {\ttfamily\footnotesize \textbackslash{}widehat\{\allowbreak{}}\par
\noindent Source : ADOPT-F7I-VI-093.2. \nolinkurl{fr/sga7i/diplomatic/expose_VI_body.tex}; \nolinkurl{fr/sga7i/current/expose_VI_body.tex}\par

\par\medskip\noindent\hypertarget{app-F7I-VI-093.3}{}\textbf{F7I-VI-093.3}\quad p. 105; inherited corrected-layer difference.\par
\noindent Imprimé : \emph{(vide)}\par
\noindent Corrigé : {\ttfamily\footnotesize \}\allowbreak{}}\par
\noindent Source : ADOPT-F7I-VI-093.3. \nolinkurl{fr/sga7i/diplomatic/expose_VI_body.tex}; \nolinkurl{fr/sga7i/current/expose_VI_body.tex}\par

\par\medskip\noindent\hypertarget{app-F7I-VI-093.4}{}\textbf{F7I-VI-093.4}\quad p. 105; inherited corrected-layer difference.\par
\noindent Imprimé : \emph{(vide)}\par
\noindent Corrigé : {\ttfamily\footnotesize \textbackslash{}widehat\{\allowbreak{}}\par
\noindent Source : ADOPT-F7I-VI-093.4. \nolinkurl{fr/sga7i/diplomatic/expose_VI_body.tex}; \nolinkurl{fr/sga7i/current/expose_VI_body.tex}\par

\par\medskip\noindent\hypertarget{app-F7I-VI-093.5}{}\textbf{F7I-VI-093.5}\quad p. 105; inherited corrected-layer difference.\par
\noindent Imprimé : \emph{(vide)}\par
\noindent Corrigé : {\ttfamily\footnotesize \}\allowbreak{}}\par
\noindent Source : ADOPT-F7I-VI-093.5. \nolinkurl{fr/sga7i/diplomatic/expose_VI_body.tex}; \nolinkurl{fr/sga7i/current/expose_VI_body.tex}\par

\par\medskip\noindent\hypertarget{app-F7I-VI-093.6}{}\textbf{F7I-VI-093.6}\quad p. 105; inherited corrected-layer difference.\par
\noindent Imprimé : {\ttfamily\footnotesize write\allowbreak{}}\par
\noindent Corrigé : {\ttfamily\footnotesize denote\allowbreak{}}\par
\noindent Source : ADOPT-F7I-VI-093.6. \nolinkurl{fr/sga7i/diplomatic/expose_VI_body.tex}; \nolinkurl{fr/sga7i/current/expose_VI_body.tex}\par

\par\medskip\noindent\hypertarget{app-F7I-VI-093.7}{}\textbf{F7I-VI-093.7}\quad p. 105; inherited corrected-layer difference.\par
\noindent Imprimé : \emph{(vide)}\par
\noindent Corrigé : {\ttfamily\footnotesize \textbackslash{}widehat\{\allowbreak{}}\par
\noindent Source : ADOPT-F7I-VI-093.7. \nolinkurl{fr/sga7i/diplomatic/expose_VI_body.tex}; \nolinkurl{fr/sga7i/current/expose_VI_body.tex}\par

\par\medskip\noindent\hypertarget{app-F7I-VI-093.8}{}\textbf{F7I-VI-093.8}\quad p. 105; inherited corrected-layer difference.\par
\noindent Imprimé : \emph{(vide)}\par
\noindent Corrigé : {\ttfamily\footnotesize \}\allowbreak{}}\par
\noindent Source : ADOPT-F7I-VI-093.8. \nolinkurl{fr/sga7i/diplomatic/expose_VI_body.tex}; \nolinkurl{fr/sga7i/current/expose_VI_body.tex}\par

\par\medskip\noindent\hypertarget{app-F7I-VI-093.9}{}\textbf{F7I-VI-093.9}\quad p. 105; inherited corrected-layer difference.\par
\noindent Imprimé : \emph{(vide)}\par
\noindent Corrigé : {\ttfamily\footnotesize \textbackslash{}widehat\{\allowbreak{}}\par
\noindent Source : ADOPT-F7I-VI-093.9. \nolinkurl{fr/sga7i/diplomatic/expose_VI_body.tex}; \nolinkurl{fr/sga7i/current/expose_VI_body.tex}\par

\par\medskip\noindent\hypertarget{app-F7I-VI-093.10}{}\textbf{F7I-VI-093.10}\quad p. 105; inherited corrected-layer difference.\par
\noindent Imprimé : \emph{(vide)}\par
\noindent Corrigé : {\ttfamily\footnotesize \}\allowbreak{}}\par
\noindent Source : ADOPT-F7I-VI-093.10. \nolinkurl{fr/sga7i/diplomatic/expose_VI_body.tex}; \nolinkurl{fr/sga7i/current/expose_VI_body.tex}\par

\par\medskip\noindent\hypertarget{app-F7I-VI-093.11}{}\textbf{F7I-VI-093.11}\quad p. 105; designated-witness attribution.\par
\noindent Imprimé : \emph{(vide)}\par
\noindent Corrigé : {\ttfamily\footnotesize \textbackslash{}widehat\{\allowbreak{}}\par
\noindent Source : WITNESS-F7I-VI-093.11. \nolinkurl{fr/sga7i/diplomatic/expose_VI_body.tex}; \nolinkurl{fr/sga7i/current/expose_VI_body.tex}\par

\par\medskip\noindent\hypertarget{app-F7I-VI-093.12}{}\textbf{F7I-VI-093.12}\quad p. 105; designated-witness attribution.\par
\noindent Imprimé : \emph{(vide)}\par
\noindent Corrigé : {\ttfamily\footnotesize \}\allowbreak{}}\par
\noindent Source : WITNESS-F7I-VI-093.12. \nolinkurl{fr/sga7i/diplomatic/expose_VI_body.tex}; \nolinkurl{fr/sga7i/current/expose_VI_body.tex}\par

\par\medskip\noindent\hypertarget{app-F7I-VI-093.13}{}\textbf{F7I-VI-093.13}\quad p. 105; designated-witness attribution.\par
\noindent Imprimé : \emph{(vide)}\par
\noindent Corrigé : {\ttfamily\footnotesize \textbackslash{}widehat\{\allowbreak{}}\par
\noindent Source : WITNESS-F7I-VI-093.13. \nolinkurl{fr/sga7i/diplomatic/expose_VI_body.tex}; \nolinkurl{fr/sga7i/current/expose_VI_body.tex}\par

\par\medskip\noindent\hypertarget{app-F7I-VI-093.14}{}\textbf{F7I-VI-093.14}\quad p. 105; designated-witness attribution.\par
\noindent Imprimé : \emph{(vide)}\par
\noindent Corrigé : {\ttfamily\footnotesize \}\allowbreak{}}\par
\noindent Source : WITNESS-F7I-VI-093.14. \nolinkurl{fr/sga7i/diplomatic/expose_VI_body.tex}; \nolinkurl{fr/sga7i/current/expose_VI_body.tex}\par

\par\medskip\noindent\hypertarget{app-F7I-VI-093.15}{}\textbf{F7I-VI-093.15}\quad p. 105; designated-witness attribution.\par
\noindent Imprimé : \emph{(vide)}\par
\noindent Corrigé : {\ttfamily\footnotesize \textbackslash{}widehat\{\allowbreak{}}\par
\noindent Source : WITNESS-F7I-VI-093.15. \nolinkurl{fr/sga7i/diplomatic/expose_VI_body.tex}; \nolinkurl{fr/sga7i/current/expose_VI_body.tex}\par

\par\medskip\noindent\hypertarget{app-F7I-VI-093.16}{}\textbf{F7I-VI-093.16}\quad p. 105; designated-witness attribution.\par
\noindent Imprimé : \emph{(vide)}\par
\noindent Corrigé : {\ttfamily\footnotesize \}\allowbreak{}}\par
\noindent Source : WITNESS-F7I-VI-093.16. \nolinkurl{fr/sga7i/diplomatic/expose_VI_body.tex}; \nolinkurl{fr/sga7i/current/expose_VI_body.tex}\par

\par\medskip\noindent\hypertarget{app-F7I-VI-093.17}{}\textbf{F7I-VI-093.17}\quad p. 105; printed misprint.\par
\noindent Imprimé : \emph{(vide)}\par
\noindent Corrigé : {\ttfamily\footnotesize \textbackslash{}hat\{\allowbreak{}}\par
\noindent Source : 7I:211. \nolinkurl{fixes/MISPRINTS_7I.json#211}\par

\par\medskip\noindent\hypertarget{app-F7I-VI-093.18}{}\textbf{F7I-VI-093.18}\quad p. 105; printed misprint.\par
\noindent Imprimé : \emph{(vide)}\par
\noindent Corrigé : {\ttfamily\footnotesize \}\allowbreak{}}\par
\noindent Source : 7I:211. \nolinkurl{fixes/MISPRINTS_7I.json#211}\par

\par\medskip\noindent\hypertarget{app-F7I-VI-093.19}{}\textbf{F7I-VI-093.19}\quad p. 105; misprint\_in\_print.\par
\noindent Imprimé : {\ttfamily\footnotesize X\allowbreak{}}\par
\noindent Corrigé : {\ttfamily\footnotesize x\allowbreak{}}\par
\noindent Source : FIX-7I-D-0112. \nolinkurl{fixes/FIXLOG_7I.jsonl#112}\par

\par\medskip\noindent\hypertarget{app-F7I-VI-094.1}{}\textbf{F7I-VI-094.1}\quad p. 105; inherited corrected-layer difference.\par
\noindent Imprimé : {\ttfamily\footnotesize the\allowbreak{} }\par
\noindent Corrigé : \emph{(vide)}\par
\noindent Source : ADOPT-F7I-VI-094.1. \nolinkurl{fr/sga7i/diplomatic/expose_VI_body.tex}; \nolinkurl{fr/sga7i/current/expose_VI_body.tex}\par

\par\medskip\noindent\hypertarget{app-F7I-VI-095.1}{}\textbf{F7I-VI-095.1}\quad p. 106; inherited corrected-layer difference.\par
\noindent Imprimé : {\ttfamily\footnotesize of\allowbreak{}}\par
\noindent Corrigé : {\ttfamily\footnotesize for\allowbreak{}}\par
\noindent Source : ADOPT-F7I-VI-095.1. \nolinkurl{fr/sga7i/diplomatic/expose_VI_body.tex}; \nolinkurl{fr/sga7i/current/expose_VI_body.tex}\par

\par\medskip\noindent\hypertarget{app-F7I-VI-095.2}{}\textbf{F7I-VI-095.2}\quad p. 106; inherited corrected-layer difference.\par
\noindent Imprimé : {\ttfamily\footnotesize coefficient\allowbreak{}}\par
\noindent Corrigé : {\ttfamily\footnotesize coefficients\allowbreak{} in\allowbreak{}}\par
\noindent Source : ADOPT-F7I-VI-095.2. \nolinkurl{fr/sga7i/diplomatic/expose_VI_body.tex}; \nolinkurl{fr/sga7i/current/expose_VI_body.tex}\par

\par\medskip\noindent\hypertarget{app-F7I-VI-095.3}{}\textbf{F7I-VI-095.3}\quad p. 106; printed misprint.\par
\noindent Imprimé : {\ttfamily\footnotesize '\allowbreak{}}\par
\noindent Corrigé : \emph{(vide)}\par
\noindent Source : 7I:229. \nolinkurl{fixes/MISPRINTS_7I.json#229}\par

\par\medskip\noindent\hypertarget{app-F7I-VI-095.4}{}\textbf{F7I-VI-095.4}\quad p. 106; printed misprint.\par
\noindent Imprimé : \emph{(vide)}\par
\noindent Corrigé : {\ttfamily\footnotesize '\allowbreak{}}\par
\noindent Source : 7I:229. \nolinkurl{fixes/MISPRINTS_7I.json#229}\par

\par\medskip\noindent\hypertarget{app-F7I-VI-095.5}{}\textbf{F7I-VI-095.5}\quad p. 106; designated-witness attribution.\par
\noindent Imprimé : \emph{(vide)}\par
\noindent Corrigé : {\ttfamily\footnotesize  a\allowbreak{}}\par
\noindent Source : WITNESS-F7I-VI-095.5. \nolinkurl{fr/sga7i/diplomatic/expose_VI_body.tex}; \nolinkurl{fr/sga7i/current/expose_VI_body.tex}\par

\par\medskip\noindent\hypertarget{app-F7I-VI-095.6}{}\textbf{F7I-VI-095.6}\quad p. 106; inherited corrected-layer difference.\par
\noindent Imprimé : {\ttfamily\footnotesize 11\allowbreak{}}\par
\noindent Corrigé : {\ttfamily\footnotesize 8\allowbreak{}}\par
\noindent Source : ADOPT-F7I-VI-095.6. \nolinkurl{fr/sga7i/diplomatic/expose_VI_body.tex}; \nolinkurl{fr/sga7i/current/expose_VI_body.tex}\par

\par\medskip\noindent\hypertarget{app-F7I-VI-096.1}{}\textbf{F7I-VI-096.1}\quad p. 106; inherited corrected-layer difference.\par
\noindent Imprimé : {\ttfamily\footnotesize \}\allowbreak{}}\par
\noindent Corrigé : \emph{(vide)}\par
\noindent Source : ADOPT-F7I-VI-096.1. \nolinkurl{fr/sga7i/diplomatic/expose_VI_body.tex}; \nolinkurl{fr/sga7i/current/expose_VI_body.tex}\par

\par\medskip\noindent\hypertarget{app-F7I-VI-096.2}{}\textbf{F7I-VI-096.2}\quad p. 106; inherited corrected-layer difference.\par
\noindent Imprimé : \emph{(vide)}\par
\noindent Corrigé : {\ttfamily\footnotesize \}\allowbreak{}}\par
\noindent Source : ADOPT-F7I-VI-096.2. \nolinkurl{fr/sga7i/diplomatic/expose_VI_body.tex}; \nolinkurl{fr/sga7i/current/expose_VI_body.tex}\par

\par\medskip\noindent\hypertarget{app-F7I-VI-096.3}{}\textbf{F7I-VI-096.3}\quad p. 106; inherited corrected-layer difference.\par
\noindent Imprimé : {\ttfamily\footnotesize \}\allowbreak{}}\par
\noindent Corrigé : \emph{(vide)}\par
\noindent Source : ADOPT-F7I-VI-096.3. \nolinkurl{fr/sga7i/diplomatic/expose_VI_body.tex}; \nolinkurl{fr/sga7i/current/expose_VI_body.tex}\par

\par\medskip\noindent\hypertarget{app-F7I-VI-096.4}{}\textbf{F7I-VI-096.4}\quad p. 106; inherited corrected-layer difference.\par
\noindent Imprimé : \emph{(vide)}\par
\noindent Corrigé : {\ttfamily\footnotesize \}\allowbreak{}}\par
\noindent Source : ADOPT-F7I-VI-096.4. \nolinkurl{fr/sga7i/diplomatic/expose_VI_body.tex}; \nolinkurl{fr/sga7i/current/expose_VI_body.tex}\par

\par\medskip\noindent\hypertarget{app-F7I-VI-097.1}{}\textbf{F7I-VI-097.1}\quad p. 106; inherited corrected-layer difference.\par
\noindent Imprimé : {\ttfamily\footnotesize gives\allowbreak{}}\par
\noindent Corrigé : {\ttfamily\footnotesize give\allowbreak{}}\par
\noindent Source : ADOPT-F7I-VI-097.1. \nolinkurl{fr/sga7i/diplomatic/expose_VI_body.tex}; \nolinkurl{fr/sga7i/current/expose_VI_body.tex}\par

\par\medskip\noindent\hypertarget{app-F7I-VI-097.2}{}\textbf{F7I-VI-097.2}\quad p. 107; inherited corrected-layer difference.\par
\noindent Imprimé : {\ttfamily\footnotesize isomorphism\allowbreak{}}\par
\noindent Corrigé : {\ttfamily\footnotesize isomorphisms\allowbreak{}}\par
\noindent Source : ADOPT-F7I-VI-097.2. \nolinkurl{fr/sga7i/diplomatic/expose_VI_body.tex}; \nolinkurl{fr/sga7i/current/expose_VI_body.tex}\par

\par\medskip\noindent\hypertarget{app-F7I-VI-097.3}{}\textbf{F7I-VI-097.3}\quad p. 107; inherited corrected-layer difference.\par
\noindent Imprimé : {\ttfamily\footnotesize comes\allowbreak{}}\par
\noindent Corrigé : {\ttfamily\footnotesize come\allowbreak{}}\par
\noindent Source : ADOPT-F7I-VI-097.3. \nolinkurl{fr/sga7i/diplomatic/expose_VI_body.tex}; \nolinkurl{fr/sga7i/current/expose_VI_body.tex}\par

\par\medskip\noindent\hypertarget{app-F7I-VI-097.4}{}\textbf{F7I-VI-097.4}\quad p. 107; inherited corrected-layer difference.\par
\noindent Imprimé : {\ttfamily\footnotesize trivial\allowbreak{} }\par
\noindent Corrigé : \emph{(vide)}\par
\noindent Source : ADOPT-F7I-VI-097.4. \nolinkurl{fr/sga7i/diplomatic/expose_VI_body.tex}; \nolinkurl{fr/sga7i/current/expose_VI_body.tex}\par

\par\medskip\noindent\hypertarget{app-F7I-VI-097.5}{}\textbf{F7I-VI-097.5}\quad p. 107; inherited corrected-layer difference.\par
\noindent Imprimé : {\ttfamily\footnotesize chasing\allowbreak{}}\par
\noindent Corrigé : {\ttfamily\footnotesize chase\allowbreak{}}\par
\noindent Source : ADOPT-F7I-VI-097.5. \nolinkurl{fr/sga7i/diplomatic/expose_VI_body.tex}; \nolinkurl{fr/sga7i/current/expose_VI_body.tex}\par

\par\medskip\noindent\hypertarget{app-F7I-VI-097.6}{}\textbf{F7I-VI-097.6}\quad p. 107; inherited corrected-layer difference.\par
\noindent Imprimé : {\ttfamily\footnotesize \}\allowbreak{}}\par
\noindent Corrigé : \emph{(vide)}\par
\noindent Source : ADOPT-F7I-VI-097.6. \nolinkurl{fr/sga7i/diplomatic/expose_VI_body.tex}; \nolinkurl{fr/sga7i/current/expose_VI_body.tex}\par

\par\medskip\noindent\hypertarget{app-F7I-VI-097.7}{}\textbf{F7I-VI-097.7}\quad p. 107; inherited corrected-layer difference.\par
\noindent Imprimé : \emph{(vide)}\par
\noindent Corrigé : {\ttfamily\footnotesize \}\allowbreak{}}\par
\noindent Source : ADOPT-F7I-VI-097.7. \nolinkurl{fr/sga7i/diplomatic/expose_VI_body.tex}; \nolinkurl{fr/sga7i/current/expose_VI_body.tex}\par

\par\medskip\noindent\hypertarget{app-F7I-VI-097.8}{}\textbf{F7I-VI-097.8}\quad p. 107; inherited corrected-layer difference.\par
\noindent Imprimé : {\ttfamily\footnotesize an\allowbreak{}}\par
\noindent Corrigé : {\ttfamily\footnotesize the\allowbreak{}}\par
\noindent Source : ADOPT-F7I-VI-097.8. \nolinkurl{fr/sga7i/diplomatic/expose_VI_body.tex}; \nolinkurl{fr/sga7i/current/expose_VI_body.tex}\par

\par\medskip\noindent\hypertarget{app-F7I-VI-098.1}{}\textbf{F7I-VI-098.1}\quad p. 108; printed misprint.\par
\noindent Imprimé : \emph{(vide)}\par
\noindent Corrigé : {\ttfamily\footnotesize ,\allowbreak{}}\par
\noindent Source : 7I:230. \nolinkurl{fixes/MISPRINTS_7I.json#230}\par

\par\medskip\noindent\hypertarget{app-F7I-VI-098.2}{}\textbf{F7I-VI-098.2}\quad p. 108; printed misprint.\par
\noindent Imprimé : {\ttfamily\footnotesize i\allowbreak{}}\par
\noindent Corrigé : {\ttfamily\footnotesize j\allowbreak{}}\par
\noindent Source : 7I:231. \nolinkurl{fixes/MISPRINTS_7I.json#231}\par

\par\medskip\noindent\hypertarget{app-F7I-VI-099.1}{}\textbf{F7I-VI-099.1}\quad p. 108; misprint\_pass.\par
\noindent Imprimé : {\ttfamily\footnotesize 0\allowbreak{}}\par
\noindent Corrigé : {\ttfamily\footnotesize c\allowbreak{}}\par
\noindent Source : 7I:232. \nolinkurl{fixes/FIXLOG_7I_CURRENT.jsonl#698}; \nolinkurl{fixes/MISPRINTS_7I.json#232}\par

\par\medskip\noindent\hypertarget{app-F7I-VI-100.1}{}\textbf{F7I-VI-100.1}\quad p. 109; misprint\_pass.\par
\noindent Imprimé : \emph{(vide)}\par
\noindent Corrigé : {\ttfamily\footnotesize \textbackslash{}times\allowbreak{} }\par
\noindent Source : 7I:233. \nolinkurl{fixes/FIXLOG_7I_CURRENT.jsonl#699}; \nolinkurl{fixes/MISPRINTS_7I.json#233}\par

\par\medskip\noindent\hypertarget{app-F7I-VI-100.2}{}\textbf{F7I-VI-100.2}\quad p. 109; misprint\_pass.\par
\noindent Imprimé : \emph{(vide)}\par
\noindent Corrigé : {\ttfamily\footnotesize  \textbackslash{}times\allowbreak{}}\par
\noindent Source : 7I:233. \nolinkurl{fixes/FIXLOG_7I_CURRENT.jsonl#699}; \nolinkurl{fixes/MISPRINTS_7I.json#233}\par

\par\medskip\noindent\hypertarget{app-F7I-VI-101.1}{}\textbf{F7I-VI-101.1}\quad p. 110; misprint\_pass.\par
\noindent Imprimé : {\ttfamily\footnotesize explicitely\allowbreak{}}\par
\noindent Corrigé : {\ttfamily\footnotesize explicitly\allowbreak{}}\par
\noindent Source : 7I:234. \nolinkurl{fixes/FIXLOG_7I_CURRENT.jsonl#700}; \nolinkurl{fixes/MISPRINTS_7I.json#234}\par

\par\medskip\noindent\hypertarget{app-F7I-VI-101.2}{}\textbf{F7I-VI-101.2}\quad p. 110; misprint\_pass.\par
\noindent Imprimé : {\ttfamily\footnotesize \textbackslash{}qquad\allowbreak{}}\par
\noindent Corrigé : {\ttfamily\footnotesize =\allowbreak{}}\par
\noindent Source : 7I:235. \nolinkurl{fixes/FIXLOG_7I_CURRENT.jsonl#701}; \nolinkurl{fixes/MISPRINTS_7I.json#235}\par

\par\medskip\noindent\hypertarget{app-F7I-VI-101.3}{}\textbf{F7I-VI-101.3}\quad p. 110; misprint\_pass.\par
\noindent Imprimé : \emph{(vide)}\par
\noindent Corrigé : {\ttfamily\footnotesize \_0\allowbreak{}}\par
\noindent Source : 7I:236. \nolinkurl{fixes/FIXLOG_7I_CURRENT.jsonl#702}; \nolinkurl{fixes/MISPRINTS_7I.json#236}\par

\par\medskip\noindent\hypertarget{app-F7I-VI-101.4}{}\textbf{F7I-VI-101.4}\quad p. 110; misprint\_pass.\par
\noindent Imprimé : {\ttfamily\footnotesize ji\allowbreak{}}\par
\noindent Corrigé : {\ttfamily\footnotesize ij\allowbreak{}}\par
\noindent Source : 7I:237. \nolinkurl{fixes/FIXLOG_7I_CURRENT.jsonl#703}; \nolinkurl{fixes/MISPRINTS_7I.json#237}\par

\par\medskip\noindent\hypertarget{app-F7I-VI-102.1}{}\textbf{F7I-VI-102.1}\quad p. 111; misprint\_pass.\par
\noindent Imprimé : \emph{(vide)}\par
\noindent Corrigé : {\ttfamily\footnotesize \{\allowbreak{}}\par
\noindent Source : 7I:213. \nolinkurl{fixes/FIXLOG_7I_CURRENT.jsonl#704}; \nolinkurl{fixes/MISPRINTS_7I.json#213}\par

\par\medskip\noindent\hypertarget{app-F7I-VI-102.2}{}\textbf{F7I-VI-102.2}\quad p. 111; misprint\_pass.\par
\noindent Imprimé : \emph{(vide)}\par
\noindent Corrigé : {\ttfamily\footnotesize \_0\}\allowbreak{}}\par
\noindent Source : 7I:213. \nolinkurl{fixes/FIXLOG_7I_CURRENT.jsonl#704}; \nolinkurl{fixes/MISPRINTS_7I.json#213}\par

\par\medskip\noindent\hypertarget{app-F7I-VI-102.3}{}\textbf{F7I-VI-102.3}\quad p. 111; misprint\_pass.\par
\noindent Imprimé : {\ttfamily\footnotesize '\allowbreak{}}\par
\noindent Corrigé : \emph{(vide)}\par
\noindent Source : 7I:214. \nolinkurl{fixes/FIXLOG_7I_CURRENT.jsonl#705}; \nolinkurl{fixes/MISPRINTS_7I.json#214}\par

\par\medskip\noindent\hypertarget{app-F7I-VI-103.1}{}\textbf{F7I-VI-103.1}\quad p. 112; misprint\_pass.\par
\noindent Imprimé : {\ttfamily\footnotesize standart\allowbreak{}}\par
\noindent Corrigé : {\ttfamily\footnotesize standard\allowbreak{}}\par
\noindent Source : 7I:216. \nolinkurl{fixes/FIXLOG_7I_CURRENT.jsonl#707}; \nolinkurl{fixes/MISPRINTS_7I.json#216}\par

\par\medskip\noindent\hypertarget{app-F7I-VI-103.2}{}\textbf{F7I-VI-103.2}\quad p. 112; misprint\_pass.\par
\noindent Imprimé : {\ttfamily\footnotesize  \textbackslash{}quad\allowbreak{} }\par
\noindent Corrigé : {\ttfamily\footnotesize ;\textbackslash{}\allowbreak{} }\par
\noindent Source : 7I:215. \nolinkurl{fixes/FIXLOG_7I_CURRENT.jsonl#706}; \nolinkurl{fixes/MISPRINTS_7I.json#215}\par

\par\medskip\noindent\hypertarget{app-F7I-VI-104.1}{}\textbf{F7I-VI-104.1}\quad p. 112; misprint\_pass.\par
\noindent Imprimé : {\ttfamily\footnotesize a\allowbreak{} }\par
\noindent Corrigé : \emph{(vide)}\par
\noindent Source : 7I:218. \nolinkurl{fixes/FIXLOG_7I_CURRENT.jsonl#708}; \nolinkurl{fixes/MISPRINTS_7I.json#218}\par

\par\medskip\noindent\hypertarget{app-F7I-VI-105.1}{}\textbf{F7I-VI-105.1}\quad p. 113; designated-witness attribution.\par
\noindent Imprimé : {\ttfamily\footnotesize 12\allowbreak{}}\par
\noindent Corrigé : {\ttfamily\footnotesize 13\allowbreak{}}\par
\noindent Source : WITNESS-F7I-VI-105.1. \nolinkurl{fr/sga7i/diplomatic/expose_VI_body.tex}; \nolinkurl{fr/sga7i/current/expose_VI_body.tex}\par

\par\medskip\noindent\hypertarget{app-F7I-VI-105.2}{}\textbf{F7I-VI-105.2}\quad p. 113; printed misprint.\par
\noindent Imprimé : {\ttfamily\footnotesize cotangentcomplex\allowbreak{}}\par
\noindent Corrigé : {\ttfamily\footnotesize cotangent\allowbreak{} complex\allowbreak{}}\par
\noindent Source : 7I:219. \nolinkurl{fixes/MISPRINTS_7I.json#219}\par

\par\medskip\noindent\hypertarget{app-F7I-VI-105.3}{}\textbf{F7I-VI-105.3}\quad p. 113; printed misprint.\par
\noindent Imprimé : \emph{(vide)}\par
\noindent Corrigé : {\ttfamily\footnotesize \textbackslash{}underline\{\allowbreak{}}\par
\noindent Source : 7I:220. \nolinkurl{fixes/MISPRINTS_7I.json#220}\par

\par\medskip\noindent\hypertarget{app-F7I-VI-105.4}{}\textbf{F7I-VI-105.4}\quad p. 113; printed misprint.\par
\noindent Imprimé : \emph{(vide)}\par
\noindent Corrigé : {\ttfamily\footnotesize \}\allowbreak{}}\par
\noindent Source : 7I:220. \nolinkurl{fixes/MISPRINTS_7I.json#220}\par

\par\medskip\noindent\hypertarget{app-F7I-VI-106.1}{}\textbf{F7I-VI-106.1}\quad p. 114; designated-witness attribution.\par
\noindent Imprimé : {\ttfamily\footnotesize fistly\allowbreak{}}\par
\noindent Corrigé : {\ttfamily\footnotesize first\allowbreak{}}\par
\noindent Source : WITNESS-F7I-VI-106.1. \nolinkurl{fr/sga7i/diplomatic/expose_VI_body.tex}; \nolinkurl{fr/sga7i/current/expose_VI_body.tex}\par

\par\medskip\noindent\hypertarget{app-F7I-VI-107.1}{}\textbf{F7I-VI-107.1}\quad p. 114; printed misprint.\par
\noindent Imprimé : {\ttfamily\footnotesize P\allowbreak{}}\par
\noindent Corrigé : {\ttfamily\footnotesize p\allowbreak{}}\par
\noindent Source : 7I:222. \nolinkurl{fixes/MISPRINTS_7I.json#222}\par

\par\medskip\noindent\hypertarget{app-F7I-VI-107.2}{}\textbf{F7I-VI-107.2}\quad p. 114; printed misprint.\par
\noindent Imprimé : {\ttfamily\footnotesize P\allowbreak{}}\par
\noindent Corrigé : {\ttfamily\footnotesize p\allowbreak{}}\par
\noindent Source : 7I:223. \nolinkurl{fixes/MISPRINTS_7I.json#223}\par

\par\medskip\noindent\hypertarget{app-F7I-VI-107.3}{}\textbf{F7I-VI-107.3}\quad p. 114; printed misprint.\par
\noindent Imprimé : {\ttfamily\footnotesize does\allowbreak{}}\par
\noindent Corrigé : {\ttfamily\footnotesize do\allowbreak{}}\par
\noindent Source : 7I:224. \nolinkurl{fixes/MISPRINTS_7I.json#224}\par

\par\medskip\noindent\hypertarget{app-F7I-VI-107.4}{}\textbf{F7I-VI-107.4}\quad p. 114; designated-witness attribution.\par
\noindent Imprimé : \emph{(vide)}\par
\noindent Corrigé : {\ttfamily\footnotesize the\allowbreak{} }\par
\noindent Source : WITNESS-F7I-VI-107.4. \nolinkurl{fr/sga7i/diplomatic/expose_VI_body.tex}; \nolinkurl{fr/sga7i/current/expose_VI_body.tex}\par

\par\medskip\noindent\hypertarget{app-F7I-VI-107.5}{}\textbf{F7I-VI-107.5}\quad p. 114; designated-witness attribution.\par
\noindent Imprimé : \emph{(vide)}\par
\noindent Corrigé : {\ttfamily\footnotesize  the\allowbreak{}}\par
\noindent Source : WITNESS-F7I-VI-107.5. \nolinkurl{fr/sga7i/diplomatic/expose_VI_body.tex}; \nolinkurl{fr/sga7i/current/expose_VI_body.tex}\par

\par\medskip\noindent\hypertarget{app-F7I-VI-108.1}{}\textbf{F7I-VI-108.1}\quad p. 114; printed misprint.\par
\noindent Imprimé : {\ttfamily\footnotesize P\allowbreak{}}\par
\noindent Corrigé : {\ttfamily\footnotesize p\allowbreak{}}\par
\noindent Source : 7I:226. \nolinkurl{fixes/MISPRINTS_7I.json#226}\par

\par\medskip\noindent\hypertarget{app-F7I-VI-109.1}{}\textbf{F7I-VI-109.1}\quad p. 115; inherited corrected-layer difference.\par
\noindent Imprimé : {\ttfamily\footnotesize v\allowbreak{}}\par
\noindent Corrigé : {\ttfamily\footnotesize \textbackslash{}nu\allowbreak{}}\par
\noindent Source : ADOPT-F7I-VI-109.1. \nolinkurl{fr/sga7i/diplomatic/expose_VI_body.tex}; \nolinkurl{fr/sga7i/current/expose_VI_body.tex}\par

\par\medskip\noindent\hypertarget{app-F7I-VI-109.2}{}\textbf{F7I-VI-109.2}\quad p. 115; inherited corrected-layer difference.\par
\noindent Imprimé : {\ttfamily\footnotesize v\allowbreak{}}\par
\noindent Corrigé : {\ttfamily\footnotesize \textbackslash{}nu\allowbreak{}}\par
\noindent Source : ADOPT-F7I-VI-109.2. \nolinkurl{fr/sga7i/diplomatic/expose_VI_body.tex}; \nolinkurl{fr/sga7i/current/expose_VI_body.tex}\par

\par\medskip\noindent\hypertarget{app-F7I-VI-109.3}{}\textbf{F7I-VI-109.3}\quad p. 115; inherited corrected-layer difference.\par
\noindent Imprimé : {\ttfamily\footnotesize v\allowbreak{}}\par
\noindent Corrigé : {\ttfamily\footnotesize \textbackslash{}nu\allowbreak{}}\par
\noindent Source : ADOPT-F7I-VI-109.3. \nolinkurl{fr/sga7i/diplomatic/expose_VI_body.tex}; \nolinkurl{fr/sga7i/current/expose_VI_body.tex}\par

\par\medskip\noindent\hypertarget{app-F7I-VI-109.4}{}\textbf{F7I-VI-109.4}\quad p. 115; inherited corrected-layer difference.\par
\noindent Imprimé : {\ttfamily\footnotesize v\allowbreak{}}\par
\noindent Corrigé : {\ttfamily\footnotesize \textbackslash{}nu\allowbreak{}}\par
\noindent Source : ADOPT-F7I-VI-109.4. \nolinkurl{fr/sga7i/diplomatic/expose_VI_body.tex}; \nolinkurl{fr/sga7i/current/expose_VI_body.tex}\par

\par\medskip\noindent\hypertarget{app-F7I-VI-109.5}{}\textbf{F7I-VI-109.5}\quad p. 115; inherited corrected-layer difference.\par
\noindent Imprimé : {\ttfamily\footnotesize v\allowbreak{}}\par
\noindent Corrigé : {\ttfamily\footnotesize \textbackslash{}nu\allowbreak{}}\par
\noindent Source : ADOPT-F7I-VI-109.5. \nolinkurl{fr/sga7i/diplomatic/expose_VI_body.tex}; \nolinkurl{fr/sga7i/current/expose_VI_body.tex}\par

\par\medskip\noindent\hypertarget{app-F7I-VI-109.6}{}\textbf{F7I-VI-109.6}\quad p. 115; inherited corrected-layer difference.\par
\noindent Imprimé : {\ttfamily\footnotesize v\allowbreak{}}\par
\noindent Corrigé : {\ttfamily\footnotesize \textbackslash{}nu\allowbreak{}}\par
\noindent Source : ADOPT-F7I-VI-109.6. \nolinkurl{fr/sga7i/diplomatic/expose_VI_body.tex}; \nolinkurl{fr/sga7i/current/expose_VI_body.tex}\par

\par\medskip\noindent\hypertarget{app-F7I-VI-109.7}{}\textbf{F7I-VI-109.7}\quad p. 115; inherited corrected-layer difference.\par
\noindent Imprimé : {\ttfamily\footnotesize v\allowbreak{}}\par
\noindent Corrigé : {\ttfamily\footnotesize \textbackslash{}nu\allowbreak{}}\par
\noindent Source : ADOPT-F7I-VI-109.7. \nolinkurl{fr/sga7i/diplomatic/expose_VI_body.tex}; \nolinkurl{fr/sga7i/current/expose_VI_body.tex}\par

\par\medskip\noindent\hypertarget{app-F7I-VI-109.8}{}\textbf{F7I-VI-109.8}\quad p. 115; inherited corrected-layer difference.\par
\noindent Imprimé : {\ttfamily\footnotesize v\allowbreak{}}\par
\noindent Corrigé : {\ttfamily\footnotesize \textbackslash{}nu\allowbreak{}}\par
\noindent Source : ADOPT-F7I-VI-109.8. \nolinkurl{fr/sga7i/diplomatic/expose_VI_body.tex}; \nolinkurl{fr/sga7i/current/expose_VI_body.tex}\par

\par\medskip\noindent\hypertarget{app-F7I-VI-110.1}{}\textbf{F7I-VI-110.1}\quad p. 115; printed misprint.\par
\noindent Imprimé : {\ttfamily\footnotesize x\allowbreak{}}\par
\noindent Corrigé : {\ttfamily\footnotesize X\allowbreak{}}\par
\noindent Source : 7I:227. \nolinkurl{fixes/MISPRINTS_7I.json#227}\par

\par\medskip\noindent\hypertarget{app-F7I-VI-110.2}{}\textbf{F7I-VI-110.2}\quad p. 115; inherited corrected-layer difference.\par
\noindent Imprimé : \emph{(vide)}\par
\noindent Corrigé : {\ttfamily\footnotesize module\allowbreak{} of\allowbreak{} }\par
\noindent Source : ADOPT-F7I-VI-110.2. \nolinkurl{fr/sga7i/diplomatic/expose_VI_body.tex}; \nolinkurl{fr/sga7i/current/expose_VI_body.tex}\par

\par\medskip\noindent\hypertarget{app-F7I-VI-110.3}{}\textbf{F7I-VI-110.3}\quad p. 115; inherited corrected-layer difference.\par
\noindent Imprimé : {\ttfamily\footnotesize -\allowbreak{}}\par
\noindent Corrigé : \emph{(1 espace(s), 0 retour(s) de ligne dans le témoin TeX)}\par
\noindent Source : ADOPT-F7I-VI-110.3. \nolinkurl{fr/sga7i/diplomatic/expose_VI_body.tex}; \nolinkurl{fr/sga7i/current/expose_VI_body.tex}\par

\par\medskip\noindent\hypertarget{app-F7I-VI-110.4}{}\textbf{F7I-VI-110.4}\quad p. 115; printed misprint.\par
\noindent Imprimé : {\ttfamily\footnotesize .\allowbreak{}}\par
\noindent Corrigé : \emph{(vide)}\par
\noindent Source : 7I:228. \nolinkurl{fixes/MISPRINTS_7I.json#228}\par

\par\medskip\noindent\hypertarget{app-F7I-VI-111.1}{}\textbf{F7I-VI-111.1}\quad p. 116; misprint\_pass.\par
\noindent Imprimé : \emph{(vide)}\par
\noindent Corrigé : {\ttfamily\footnotesize .\allowbreak{}}\par
\noindent Source : 7I:239. \nolinkurl{fixes/FIXLOG_7I_CURRENT.jsonl#710}; \nolinkurl{fixes/MISPRINTS_7I.json#239}\par

\par\medskip\noindent\hypertarget{app-F7I-VI-112.1}{}\textbf{F7I-VI-112.1}\quad p. 116; designated-witness attribution.\par
\noindent Imprimé : \emph{(vide)}\par
\noindent Corrigé : {\ttfamily\footnotesize r,\allowbreak{} "i"]\allowbreak{} \textbackslash{}arrow[\allowbreak{}}\par
\noindent Source : WITNESS-F7I-VI-112.1. \nolinkurl{fr/sga7i/diplomatic/expose_VI_body.tex}; \nolinkurl{fr/sga7i/current/expose_VI_body.tex}\par

\par\medskip\noindent\hypertarget{app-F7I-VI-112.2}{}\textbf{F7I-VI-112.2}\quad p. 116; designated-witness attribution.\par
\noindent Imprimé : {\ttfamily\footnotesize  \textbackslash{}arrow[l,\allowbreak{} "i"]\allowbreak{}}\par
\noindent Corrigé : \emph{(vide)}\par
\noindent Source : WITNESS-F7I-VI-112.2. \nolinkurl{fr/sga7i/diplomatic/expose_VI_body.tex}; \nolinkurl{fr/sga7i/current/expose_VI_body.tex}\par

\par\medskip\noindent\hypertarget{app-F7I-VI-113.1}{}\textbf{F7I-VI-113.1}\quad p. 117; misprint\_pass.\par
\noindent Imprimé : \emph{(vide)}\par
\noindent Corrigé : {\ttfamily\footnotesize \textbackslash{}partial\allowbreak{} }\par
\noindent Source : 7I:240. \nolinkurl{fixes/FIXLOG_7I_CURRENT.jsonl#711}; \nolinkurl{fixes/MISPRINTS_7I.json#240}\par

\par\medskip\noindent\hypertarget{app-F7I-VI-113.2}{}\textbf{F7I-VI-113.2}\quad p. 117; misprint\_pass.\par
\noindent Imprimé : \emph{(vide)}\par
\noindent Corrigé : {\ttfamily\footnotesize \textbackslash{}partial\allowbreak{} }\par
\noindent Source : 7I:240. \nolinkurl{fixes/FIXLOG_7I_CURRENT.jsonl#711}; \nolinkurl{fixes/MISPRINTS_7I.json#240}\par

\par\medskip\noindent\hypertarget{app-F7I-VI-114.1}{}\textbf{F7I-VI-114.1}\quad p. 119; misprint\_in\_print.\par
\noindent Imprimé : {\ttfamily\footnotesize P\allowbreak{}}\par
\noindent Corrigé : {\ttfamily\footnotesize p\allowbreak{}}\par
\noindent Source : FIX-7I-D-0131. \nolinkurl{fixes/FIXLOG_7I.jsonl#131}\par

\par\medskip\noindent\hypertarget{app-F7I-VI-114.2}{}\textbf{F7I-VI-114.2}\quad p. 119; misprint\_in\_print.\par
\noindent Imprimé : {\ttfamily\footnotesize P\allowbreak{}}\par
\noindent Corrigé : {\ttfamily\footnotesize p\allowbreak{}}\par
\noindent Source : FIX-7I-D-0132. \nolinkurl{fixes/FIXLOG_7I.jsonl#132}\par

\par\medskip\noindent\hypertarget{app-F7I-VI-114.3}{}\textbf{F7I-VI-114.3}\quad p. 119; misprint\_in\_print.\par
\noindent Imprimé : {\ttfamily\footnotesize P\allowbreak{}}\par
\noindent Corrigé : {\ttfamily\footnotesize p\allowbreak{}}\par
\noindent Source : FIX-7I-D-0132. \nolinkurl{fixes/FIXLOG_7I.jsonl#132}\par

\par\medskip\noindent\hypertarget{app-F7I-VI-114.4}{}\textbf{F7I-VI-114.4}\quad p. 119; misprint\_pass.\par
\noindent Imprimé : {\ttfamily\footnotesize \textbackslash{}varinjlim\allowbreak{}}\par
\noindent Corrigé : {\ttfamily\footnotesize \textbackslash{}varprojlim\allowbreak{}}\par
\noindent Source : 7I:245. \nolinkurl{fixes/FIXLOG_7I_CURRENT.jsonl#712}; \nolinkurl{fixes/MISPRINTS_7I.json#245}\par

\par\medskip\noindent\hypertarget{app-F7I-VI-114.5}{}\textbf{F7I-VI-114.5}\quad p. 119; misprint\_in\_print.\par
\noindent Imprimé : {\ttfamily\footnotesize P\allowbreak{}}\par
\noindent Corrigé : {\ttfamily\footnotesize p\allowbreak{}}\par
\noindent Source : FIX-7I-D-0133. \nolinkurl{fixes/FIXLOG_7I.jsonl#133}\par

\par\medskip\noindent\hypertarget{app-F7I-VI-114.6}{}\textbf{F7I-VI-114.6}\quad p. 119; misprint\_in\_print.\par
\noindent Imprimé : {\ttfamily\footnotesize P\allowbreak{}}\par
\noindent Corrigé : {\ttfamily\footnotesize p\allowbreak{}}\par
\noindent Source : FIX-7I-D-0133. \nolinkurl{fixes/FIXLOG_7I.jsonl#133}\par

\par\medskip\noindent\hypertarget{app-F7I-VI-115.1}{}\textbf{F7I-VI-115.1}\quad p. 120; printed misprint.\par
\noindent Imprimé : \emph{(vide)}\par
\noindent Corrigé : {\ttfamily\footnotesize '\allowbreak{}}\par
\noindent Source : 7I:246. \nolinkurl{fixes/MISPRINTS_7I.json#246}\par

\par\medskip\noindent\hypertarget{app-F7I-VI-115.2}{}\textbf{F7I-VI-115.2}\quad p. 120; inherited corrected-layer difference.\par
\noindent Imprimé : {\ttfamily\footnotesize is\allowbreak{}}\par
\noindent Corrigé : {\ttfamily\footnotesize are\allowbreak{}}\par
\noindent Source : ADOPT-F7I-VI-115.2. \nolinkurl{fr/sga7i/diplomatic/expose_VI_body.tex}; \nolinkurl{fr/sga7i/current/expose_VI_body.tex}\par

\par\medskip\noindent\hypertarget{app-F7I-VI-115.3}{}\textbf{F7I-VI-115.3}\quad p. 120; inherited corrected-layer difference.\par
\noindent Imprimé : {\ttfamily\footnotesize that\allowbreak{}}\par
\noindent Corrigé : {\ttfamily\footnotesize those\allowbreak{}}\par
\noindent Source : ADOPT-F7I-VI-115.3. \nolinkurl{fr/sga7i/diplomatic/expose_VI_body.tex}; \nolinkurl{fr/sga7i/current/expose_VI_body.tex}\par

\par\medskip\noindent\hypertarget{app-F7I-VI-116.1}{}\textbf{F7I-VI-116.1}\quad p. 120; printed misprint.\par
\noindent Imprimé : {\ttfamily\footnotesize -\allowbreak{}}\par
\noindent Corrigé : {\ttfamily\footnotesize =\allowbreak{}}\par
\noindent Source : 7I:247. \nolinkurl{fixes/MISPRINTS_7I.json#247}\par

\par\medskip\noindent\hypertarget{app-F7I-VI-117.1}{}\textbf{F7I-VI-117.1}\quad p. 121; misprint\_in\_print.\par
\noindent Imprimé : {\ttfamily\footnotesize X\allowbreak{}}\par
\noindent Corrigé : {\ttfamily\footnotesize x\allowbreak{}}\par
\noindent Source : FIX-7I-D-0142. \nolinkurl{fixes/FIXLOG_7I.jsonl#142}\par

\par\medskip\noindent\hypertarget{app-F7I-VI-117.2}{}\textbf{F7I-VI-117.2}\quad p. 121; misprint\_in\_print.\par
\noindent Imprimé : {\ttfamily\footnotesize X\allowbreak{}}\par
\noindent Corrigé : {\ttfamily\footnotesize x\allowbreak{}}\par
\noindent Source : FIX-7I-D-0143. \nolinkurl{fixes/FIXLOG_7I.jsonl#143}\par

\par\medskip\noindent\hypertarget{app-F7I-VI-117.3}{}\textbf{F7I-VI-117.3}\quad p. 121; printed misprint.\par
\noindent Imprimé : \emph{(vide)}\par
\noindent Corrigé : {\ttfamily\footnotesize \textbackslash{}widehat\{\allowbreak{}}\par
\noindent Source : 7I:250. \nolinkurl{fixes/MISPRINTS_7I.json#250}\par

\par\medskip\noindent\hypertarget{app-F7I-VI-117.4}{}\textbf{F7I-VI-117.4}\quad p. 121; printed misprint.\par
\noindent Imprimé : \emph{(vide)}\par
\noindent Corrigé : {\ttfamily\footnotesize \}\allowbreak{}}\par
\noindent Source : 7I:250. \nolinkurl{fixes/MISPRINTS_7I.json#250}\par

\par\medskip\noindent\hypertarget{app-F7I-VI-117.5}{}\textbf{F7I-VI-117.5}\quad p. 121; printed misprint.\par
\noindent Imprimé : {\ttfamily\footnotesize \textbackslash{}lim\allowbreak{}}\par
\noindent Corrigé : {\ttfamily\footnotesize \textbackslash{}varprojlim\allowbreak{}}\par
\noindent Source : 7I:251. \nolinkurl{fixes/MISPRINTS_7I.json#251}\par

\par\medskip\noindent\hypertarget{app-F7I-VI-117.6}{}\textbf{F7I-VI-117.6}\quad p. 121; designated-witness attribution.\par
\noindent Imprimé : {\ttfamily\footnotesize separable\allowbreak{}}\par
\noindent Corrigé : {\ttfamily\footnotesize separated\allowbreak{}}\par
\noindent Source : WITNESS-F7I-VI-117.6. \nolinkurl{fr/sga7i/diplomatic/expose_VI_body.tex}; \nolinkurl{fr/sga7i/current/expose_VI_body.tex}\par

\par\medskip\noindent\hypertarget{app-F7I-VI-117.7}{}\textbf{F7I-VI-117.7}\quad p. 121; designated-witness attribution.\par
\noindent Imprimé : \emph{(vide)}\par
\noindent Corrigé : {\ttfamily\footnotesize ,y\allowbreak{}}\par
\noindent Source : WITNESS-F7I-VI-117.7. \nolinkurl{fr/sga7i/diplomatic/expose_VI_body.tex}; \nolinkurl{fr/sga7i/current/expose_VI_body.tex}\par

\par\medskip\noindent\hypertarget{app-F7I-VI-117.8}{}\textbf{F7I-VI-117.8}\quad p. 122; inherited corrected-layer difference.\par
\noindent Imprimé : \emph{(vide)}\par
\noindent Corrigé : {\ttfamily\footnotesize \textbackslash{}underline\{\allowbreak{}}\par
\noindent Source : ADOPT-F7I-VI-117.8. \nolinkurl{fr/sga7i/diplomatic/expose_VI_body.tex}; \nolinkurl{fr/sga7i/current/expose_VI_body.tex}\par

\par\medskip\noindent\hypertarget{app-F7I-VI-117.9}{}\textbf{F7I-VI-117.9}\quad p. 122; inherited corrected-layer difference.\par
\noindent Imprimé : \emph{(vide)}\par
\noindent Corrigé : {\ttfamily\footnotesize \}\allowbreak{}}\par
\noindent Source : ADOPT-F7I-VI-117.9. \nolinkurl{fr/sga7i/diplomatic/expose_VI_body.tex}; \nolinkurl{fr/sga7i/current/expose_VI_body.tex}\par

\par\medskip\noindent\hypertarget{app-F7I-VI-117.10}{}\textbf{F7I-VI-117.10}\quad p. 122; designated-witness attribution.\par
\noindent Imprimé : \emph{(vide)}\par
\noindent Corrigé : {\ttfamily\footnotesize \textbackslash{}underline\{\allowbreak{}}\par
\noindent Source : WITNESS-F7I-VI-117.10. \nolinkurl{fr/sga7i/diplomatic/expose_VI_body.tex}; \nolinkurl{fr/sga7i/current/expose_VI_body.tex}\par

\par\medskip\noindent\hypertarget{app-F7I-VI-117.11}{}\textbf{F7I-VI-117.11}\quad p. 122; designated-witness attribution.\par
\noindent Imprimé : \emph{(vide)}\par
\noindent Corrigé : {\ttfamily\footnotesize \}\allowbreak{}}\par
\noindent Source : WITNESS-F7I-VI-117.11. \nolinkurl{fr/sga7i/diplomatic/expose_VI_body.tex}; \nolinkurl{fr/sga7i/current/expose_VI_body.tex}\par

\par\medskip\noindent\hypertarget{app-F7I-VI-117.12}{}\textbf{F7I-VI-117.12}\quad p. 122; printed misprint.\par
\noindent Imprimé : \emph{(vide)}\par
\noindent Corrigé : {\ttfamily\footnotesize \textbackslash{}widehat\{\allowbreak{}}\par
\noindent Source : 7I:253. \nolinkurl{fixes/MISPRINTS_7I.json#253}\par

\par\medskip\noindent\hypertarget{app-F7I-VI-117.13}{}\textbf{F7I-VI-117.13}\quad p. 122; designated-witness attribution.\par
\noindent Imprimé : \emph{(vide)}\par
\noindent Corrigé : {\ttfamily\footnotesize \}\allowbreak{}}\par
\noindent Source : WITNESS-F7I-VI-117.13. \nolinkurl{fr/sga7i/diplomatic/expose_VI_body.tex}; \nolinkurl{fr/sga7i/current/expose_VI_body.tex}\par

\par\medskip\noindent\hypertarget{app-F7I-VI-118.1}{}\textbf{F7I-VI-118.1}\quad p. 122; inherited corrected-layer difference.\par
\noindent Imprimé : {\ttfamily\footnotesize sheaves\allowbreak{}}\par
\noindent Corrigé : {\ttfamily\footnotesize sheaf\allowbreak{}}\par
\noindent Source : ADOPT-F7I-VI-118.1. \nolinkurl{fr/sga7i/diplomatic/expose_VI_body.tex}; \nolinkurl{fr/sga7i/current/expose_VI_body.tex}\par

\par\medskip\noindent\hypertarget{app-F7I-VI-118.2}{}\textbf{F7I-VI-118.2}\quad p. 122; inherited corrected-layer difference.\par
\noindent Imprimé : \emph{(vide)}\par
\noindent Corrigé : {\ttfamily\footnotesize \textbackslash{}underline\{\allowbreak{}}\par
\noindent Source : ADOPT-F7I-VI-118.2. \nolinkurl{fr/sga7i/diplomatic/expose_VI_body.tex}; \nolinkurl{fr/sga7i/current/expose_VI_body.tex}\par

\par\medskip\noindent\hypertarget{app-F7I-VI-118.3}{}\textbf{F7I-VI-118.3}\quad p. 122; inherited corrected-layer difference.\par
\noindent Imprimé : \emph{(vide)}\par
\noindent Corrigé : {\ttfamily\footnotesize \}\allowbreak{}}\par
\noindent Source : ADOPT-F7I-VI-118.3. \nolinkurl{fr/sga7i/diplomatic/expose_VI_body.tex}; \nolinkurl{fr/sga7i/current/expose_VI_body.tex}\par

\par\medskip\noindent\hypertarget{app-F7I-VI-119.1}{}\textbf{F7I-VI-119.1}\quad p. 123; misprint\_pass.\par
\noindent Imprimé : {\ttfamily\footnotesize commute\allowbreak{}}\par
\noindent Corrigé : {\ttfamily\footnotesize commutes\allowbreak{}}\par
\noindent Source : 7I:254. \nolinkurl{fixes/FIXLOG_7I_CURRENT.jsonl#714}; \nolinkurl{fixes/MISPRINTS_7I.json#254}\par

\par\medskip\noindent\hypertarget{app-F7I-VI-120.1}{}\textbf{F7I-VI-120.1}\quad p. 123; misprint\_pass.\par
\noindent Imprimé : {\ttfamily\footnotesize O\allowbreak{}}\par
\noindent Corrigé : {\ttfamily\footnotesize I\allowbreak{}}\par
\noindent Source : 7I:255. \nolinkurl{fixes/FIXLOG_7I_CURRENT.jsonl#715}; \nolinkurl{fixes/MISPRINTS_7I.json#255}\par

\par\medskip\noindent\hypertarget{app-F7I-VI-121.1}{}\textbf{F7I-VI-121.1}\quad p. 124; misprint\_pass.\par
\noindent Imprimé : {\ttfamily\footnotesize -\allowbreak{}}\par
\noindent Corrigé : {\ttfamily\footnotesize \textbackslash{}simeq\allowbreak{}}\par
\noindent Source : 7I:256. \nolinkurl{fixes/FIXLOG_7I_CURRENT.jsonl#716}; \nolinkurl{fixes/MISPRINTS_7I.json#256}\par

\par\medskip\noindent\hypertarget{app-F7I-VI-122.1}{}\textbf{F7I-VI-122.1}\quad p. 124; misprint\_pass.\par
\noindent Imprimé : {\ttfamily\footnotesize X\allowbreak{}}\par
\noindent Corrigé : {\ttfamily\footnotesize x\allowbreak{}}\par
\noindent Source : 7I:257. \nolinkurl{fixes/FIXLOG_7I_CURRENT.jsonl#717}; \nolinkurl{fixes/MISPRINTS_7I.json#257}\par

\par\medskip\noindent\hypertarget{app-F7I-VI-122.2}{}\textbf{F7I-VI-122.2}\quad p. 124; misprint\_pass.\par
\noindent Imprimé : {\ttfamily\footnotesize X\allowbreak{}}\par
\noindent Corrigé : {\ttfamily\footnotesize x\allowbreak{}}\par
\noindent Source : 7I:257. \nolinkurl{fixes/FIXLOG_7I_CURRENT.jsonl#717}; \nolinkurl{fixes/MISPRINTS_7I.json#257}\par

\par\medskip\noindent\hypertarget{app-F7I-VI-122.3}{}\textbf{F7I-VI-122.3}\quad p. 125; misprint\_pass.\par
\noindent Imprimé : {\ttfamily\footnotesize u\allowbreak{}}\par
\noindent Corrigé : {\ttfamily\footnotesize i\allowbreak{}}\par
\noindent Source : 7I:258. \nolinkurl{fixes/FIXLOG_7I_CURRENT.jsonl#718}; \nolinkurl{fixes/MISPRINTS_7I.json#258}\par

\par\medskip\noindent\hypertarget{app-F7I-VI-123.1}{}\textbf{F7I-VI-123.1}\quad p. 125; printed misprint.\par
\noindent Imprimé : {\ttfamily\footnotesize J\allowbreak{}}\par
\noindent Corrigé : {\ttfamily\footnotesize I\allowbreak{}}\par
\noindent Source : 7I:259. \nolinkurl{fixes/MISPRINTS_7I.json#259}\par

\par\medskip\noindent\hypertarget{app-F7I-VI-124.1}{}\textbf{F7I-VI-124.1}\quad p. 125; inherited corrected-layer difference.\par
\noindent Imprimé : {\ttfamily\footnotesize \_0\allowbreak{}}\par
\noindent Corrigé : \emph{(vide)}\par
\noindent Source : ADOPT-F7I-VI-124.1. \nolinkurl{fr/sga7i/diplomatic/expose_VI_body.tex}; \nolinkurl{fr/sga7i/current/expose_VI_body.tex}\par

\par\medskip\noindent\hypertarget{app-F7I-VI-125.1}{}\textbf{F7I-VI-125.1}\quad p. 125; designated-witness attribution.\par
\noindent Imprimé : {\ttfamily\footnotesize df\allowbreak{}}\par
\noindent Corrigé : {\ttfamily\footnotesize dF\allowbreak{}}\par
\noindent Source : WITNESS-F7I-VI-125.1. \nolinkurl{fr/sga7i/diplomatic/expose_VI_body.tex}; \nolinkurl{fr/sga7i/current/expose_VI_body.tex}\par

\par\medskip\noindent\hypertarget{app-F7I-VI-125.2}{}\textbf{F7I-VI-125.2}\quad p. 125; misprint\_in\_print.\par
\noindent Imprimé : {\ttfamily\footnotesize \textbackslash{}relbar\textbackslash{}joinrel\textbackslash{}relb\allowbreak{}ar\textbackslash{}joinrel\textbackslash{}relbar\allowbreak{}}\par
\noindent Corrigé : {\ttfamily\footnotesize \textbackslash{}longrightarrow\allowbreak{}}\par
\noindent Source : FIX-7I-D-0163. \nolinkurl{fixes/FIXLOG_7I.jsonl#163}\par

\par\medskip\noindent\hypertarget{app-F7I-VI-125.3}{}\textbf{F7I-VI-125.3}\quad p. 125; misprint\_in\_print.\par
\noindent Imprimé : {\ttfamily\footnotesize \textbackslash{}relbar\textbackslash{}joinrel\textbackslash{}relb\allowbreak{}ar\textbackslash{}joinrel\textbackslash{}relbar\textbackslash{}jo\allowbreak{}inrel\textbackslash{}relbar\textbackslash{}joinrel\allowbreak{}\textbackslash{}relbar\allowbreak{}}\par
\noindent Corrigé : {\ttfamily\footnotesize \textbackslash{}xrightarrow\{\textbackslash{}hspace\allowbreak{}\{3em\}\}\allowbreak{}}\par
\noindent Source : FIX-7I-D-0164. \nolinkurl{fixes/FIXLOG_7I.jsonl#164}\par

\par\medskip\noindent\hypertarget{app-F7I-VI-126.1}{}\textbf{F7I-VI-126.1}\quad p. 126; inherited corrected-layer difference.\par
\noindent Imprimé : {\ttfamily\footnotesize df\allowbreak{}}\par
\noindent Corrigé : {\ttfamily\footnotesize dF\allowbreak{}}\par
\noindent Source : ADOPT-F7I-VI-126.1. \nolinkurl{fr/sga7i/diplomatic/expose_VI_body.tex}; \nolinkurl{fr/sga7i/current/expose_VI_body.tex}\par

\par\medskip\noindent\hypertarget{app-F7I-VI-126.2}{}\textbf{F7I-VI-126.2}\quad p. 126; inherited corrected-layer difference.\par
\noindent Imprimé : {\ttfamily\footnotesize f\allowbreak{}}\par
\noindent Corrigé : {\ttfamily\footnotesize F\allowbreak{}}\par
\noindent Source : ADOPT-F7I-VI-126.2. \nolinkurl{fr/sga7i/diplomatic/expose_VI_body.tex}; \nolinkurl{fr/sga7i/current/expose_VI_body.tex}\par

\par\medskip\noindent\hypertarget{app-F7I-VI-126.3}{}\textbf{F7I-VI-126.3}\quad p. 126; inherited corrected-layer difference.\par
\noindent Imprimé : {\ttfamily\footnotesize f\allowbreak{}}\par
\noindent Corrigé : {\ttfamily\footnotesize F\allowbreak{}}\par
\noindent Source : ADOPT-F7I-VI-126.3. \nolinkurl{fr/sga7i/diplomatic/expose_VI_body.tex}; \nolinkurl{fr/sga7i/current/expose_VI_body.tex}\par

\par\medskip\noindent\hypertarget{app-F7I-VI-126.4}{}\textbf{F7I-VI-126.4}\quad p. 126; designated-witness attribution.\par
\noindent Imprimé : \emph{(vide)}\par
\noindent Corrigé : {\ttfamily\footnotesize /I\_\{\textbackslash{}hat\{A\}\}(\textbackslash{}hat\{A\}\allowbreak{}|R)\allowbreak{}}\par
\noindent Source : WITNESS-F7I-VI-126.4. \nolinkurl{fr/sga7i/diplomatic/expose_VI_body.tex}; \nolinkurl{fr/sga7i/current/expose_VI_body.tex}\par

\par\medskip\noindent\hypertarget{app-F7I-VI-127.1}{}\textbf{F7I-VI-127.1}\quad p. 126; printed misprint.\par
\noindent Imprimé : {\ttfamily\footnotesize R\allowbreak{}}\par
\noindent Corrigé : {\ttfamily\footnotesize F\allowbreak{}}\par
\noindent Source : 7I:268. \nolinkurl{fixes/MISPRINTS_7I.json#268}\par

\par\medskip\noindent\hypertarget{app-F7I-VI-127.2}{}\textbf{F7I-VI-127.2}\quad p. 127; printed misprint.\par
\noindent Imprimé : {\ttfamily\footnotesize ab\allowbreak{}}\par
\noindent Corrigé : {\ttfamily\footnotesize ub\allowbreak{}}\par
\noindent Source : 7I:269. \nolinkurl{fixes/MISPRINTS_7I.json#269}\par

\par\medskip\noindent\hypertarget{app-F7I-VI-127.3}{}\textbf{F7I-VI-127.3}\quad p. 127; designated-witness attribution.\par
\noindent Imprimé : {\ttfamily\footnotesize Since\allowbreak{}}\par
\noindent Corrigé : {\ttfamily\footnotesize We\allowbreak{} construct\allowbreak{} a\allowbreak{} matrix\allowbreak{}}\par
\noindent Source : WITNESS-F7I-VI-127.3. \nolinkurl{fr/sga7i/diplomatic/expose_VI_body.tex}; \nolinkurl{fr/sga7i/current/expose_VI_body.tex}\par

\par\medskip\noindent\hypertarget{app-F7I-VI-127.4}{}\textbf{F7I-VI-127.4}\quad p. 127; designated-witness attribution.\par
\noindent Imprimé : {\ttfamily\footnotesize k\allowbreak{} =\allowbreak{} k\textasciicircum{}2\allowbreak{}}\par
\noindent Corrigé : {\ttfamily\footnotesize T\textbackslash{}in\textbackslash{}operatorname\{GL\allowbreak{}\}\_n(R)\allowbreak{}}\par
\noindent Source : WITNESS-F7I-VI-127.4. \nolinkurl{fr/sga7i/diplomatic/expose_VI_body.tex}; \nolinkurl{fr/sga7i/current/expose_VI_body.tex}\par

\par\medskip\noindent\hypertarget{app-F7I-VI-127.5}{}\textbf{F7I-VI-127.5}\quad p. 127; designated-witness attribution.\par
\noindent Imprimé : {\ttfamily\footnotesize and\allowbreak{}}\par
\noindent Corrigé : {\ttfamily\footnotesize such\allowbreak{} that\allowbreak{} \$F(TX)=uF(X)\$.\textbackslash{}footn\allowbreak{}ote\{Editorial\allowbreak{} proof\allowbreak{} repair:\allowbreak{} the\allowbreak{} printed\allowbreak{} inference\allowbreak{} that\allowbreak{} every\allowbreak{} unit\allowbreak{} of\allowbreak{}}\par
\noindent Source : WITNESS-F7I-VI-127.5. \nolinkurl{fr/sga7i/diplomatic/expose_VI_body.tex}; \nolinkurl{fr/sga7i/current/expose_VI_body.tex}\par

\par\medskip\noindent\hypertarget{app-F7I-VI-127.6}{}\textbf{F7I-VI-127.6}\quad p. 127; inherited corrected-layer difference.\par
\noindent Imprimé : {\ttfamily\footnotesize is\allowbreak{}}\par
\noindent Corrigé : {\ttfamily\footnotesize has\allowbreak{}}\par
\noindent Source : ADOPT-F7I-VI-127.6. \nolinkurl{fr/sga7i/diplomatic/expose_VI_body.tex}; \nolinkurl{fr/sga7i/current/expose_VI_body.tex}\par

\par\medskip\noindent\hypertarget{app-F7I-VI-127.7}{}\textbf{F7I-VI-127.7}\quad p. 127; inherited corrected-layer difference.\par
\noindent Imprimé : {\ttfamily\footnotesize complete\allowbreak{}}\par
\noindent Corrigé : {\ttfamily\footnotesize a\allowbreak{} square\allowbreak{} root\allowbreak{} need\allowbreak{} not\allowbreak{} hold\allowbreak{} in\allowbreak{} characteristic\allowbreak{} \$2\$.\allowbreak{} The\allowbreak{} following\allowbreak{} argument\allowbreak{} retains\allowbreak{} the\allowbreak{} original\allowbreak{} hypotheses.\}\allowbreak{} Write\allowbreak{} \$\textbackslash{}mathfrak\allowbreak{} m=\textbackslash{}mathfrak\allowbreak{} m\_R\$\allowbreak{}}\par
\noindent Source : ADOPT-F7I-VI-127.7. \nolinkurl{fr/sga7i/diplomatic/expose_VI_body.tex}; \nolinkurl{fr/sga7i/current/expose_VI_body.tex}\par

\par\medskip\noindent\hypertarget{app-F7I-VI-127.8}{}\textbf{F7I-VI-127.8}\quad p. 127; inherited corrected-layer difference.\par
\noindent Imprimé : {\ttfamily\footnotesize we\allowbreak{}}\par
\noindent Corrigé : {\ttfamily\footnotesize and\allowbreak{}}\par
\noindent Source : ADOPT-F7I-VI-127.8. \nolinkurl{fr/sga7i/diplomatic/expose_VI_body.tex}; \nolinkurl{fr/sga7i/current/expose_VI_body.tex}\par

\par\medskip\noindent\hypertarget{app-F7I-VI-127.9}{}\textbf{F7I-VI-127.9}\quad p. 127; inherited corrected-layer difference.\par
\noindent Imprimé : {\ttfamily\footnotesize can\allowbreak{}}\par
\noindent Corrigé : {\ttfamily\footnotesize interpret\allowbreak{}}\par
\noindent Source : ADOPT-F7I-VI-127.9. \nolinkurl{fr/sga7i/diplomatic/expose_VI_body.tex}; \nolinkurl{fr/sga7i/current/expose_VI_body.tex}\par

\par\medskip\noindent\hypertarget{app-F7I-VI-127.10}{}\textbf{F7I-VI-127.10}\quad p. 127; inherited corrected-layer difference.\par
\noindent Imprimé : {\ttfamily\footnotesize write\allowbreak{}}\par
\noindent Corrigé : {\ttfamily\footnotesize polynomial\allowbreak{} congruences\allowbreak{} coefficientwise.\allowbreak{} Since\allowbreak{}}\par
\noindent Source : ADOPT-F7I-VI-127.10. \nolinkurl{fr/sga7i/diplomatic/expose_VI_body.tex}; \nolinkurl{fr/sga7i/current/expose_VI_body.tex}\par

\par\medskip\noindent\hypertarget{app-F7I-VI-127.11}{}\textbf{F7I-VI-127.11}\quad p. 127; designated-witness attribution.\par
\noindent Imprimé : {\ttfamily\footnotesize u\allowbreak{} }\par
\noindent Corrigé : {\ttfamily\footnotesize k\allowbreak{}}\par
\noindent Source : WITNESS-F7I-VI-127.11. \nolinkurl{fr/sga7i/diplomatic/expose_VI_body.tex}; \nolinkurl{fr/sga7i/current/expose_VI_body.tex}\par

\par\medskip\noindent\hypertarget{app-F7I-VI-127.12}{}\textbf{F7I-VI-127.12}\quad p. 127; designated-witness attribution.\par
\noindent Imprimé : {\ttfamily\footnotesize  v\allowbreak{}}\par
\noindent Corrigé : {\ttfamily\footnotesize k\allowbreak{}}\par
\noindent Source : WITNESS-F7I-VI-127.12. \nolinkurl{fr/sga7i/diplomatic/expose_VI_body.tex}; \nolinkurl{fr/sga7i/current/expose_VI_body.tex}\par

\par\medskip\noindent\hypertarget{app-F7I-VI-127.13}{}\textbf{F7I-VI-127.13}\quad p. 127; inherited corrected-layer difference.\par
\noindent Imprimé : \emph{(vide)}\par
\noindent Corrigé : {\ttfamily\footnotesize ,\allowbreak{}}\par
\noindent Source : ADOPT-F7I-VI-127.13. \nolinkurl{fr/sga7i/diplomatic/expose_VI_body.tex}; \nolinkurl{fr/sga7i/current/expose_VI_body.tex}\par

\par\medskip\noindent\hypertarget{app-F7I-VI-127.14}{}\textbf{F7I-VI-127.14}\quad p. 127; inherited corrected-layer difference.\par
\noindent Imprimé : {\ttfamily\footnotesize for\allowbreak{} some\allowbreak{} unit\allowbreak{}}\par
\noindent Corrigé : {\ttfamily\footnotesize choose\allowbreak{}}\par
\noindent Source : ADOPT-F7I-VI-127.14. \nolinkurl{fr/sga7i/diplomatic/expose_VI_body.tex}; \nolinkurl{fr/sga7i/current/expose_VI_body.tex}\par

\par\medskip\noindent\hypertarget{app-F7I-VI-127.15}{}\textbf{F7I-VI-127.15}\quad p. 127; inherited corrected-layer difference.\par
\noindent Imprimé : {\ttfamily\footnotesize v\allowbreak{}}\par
\noindent Corrigé : {\ttfamily\footnotesize c\textbackslash{}in\allowbreak{} R\textasciicircum{}\textbackslash{}times\allowbreak{}}\par
\noindent Source : ADOPT-F7I-VI-127.15. \nolinkurl{fr/sga7i/diplomatic/expose_VI_body.tex}; \nolinkurl{fr/sga7i/current/expose_VI_body.tex}\par

\par\medskip\noindent\hypertarget{app-F7I-VI-127.16}{}\textbf{F7I-VI-127.16}\quad p. 127; inherited corrected-layer difference.\par
\noindent Imprimé : {\ttfamily\footnotesize in\allowbreak{}}\par
\noindent Corrigé : {\ttfamily\footnotesize whose\allowbreak{} residue\allowbreak{} satisfies\allowbreak{}}\par
\noindent Source : ADOPT-F7I-VI-127.16. \nolinkurl{fr/sga7i/diplomatic/expose_VI_body.tex}; \nolinkurl{fr/sga7i/current/expose_VI_body.tex}\par

\par\medskip\noindent\hypertarget{app-F7I-VI-127.17}{}\textbf{F7I-VI-127.17}\quad p. 127; inherited corrected-layer difference.\par
\noindent Imprimé : {\ttfamily\footnotesize R\allowbreak{}}\par
\noindent Corrigé : {\ttfamily\footnotesize \textbackslash{}bar\allowbreak{} c\textasciicircum{}2=\textbackslash{}bar\allowbreak{} u\allowbreak{}}\par
\noindent Source : ADOPT-F7I-VI-127.17. \nolinkurl{fr/sga7i/diplomatic/expose_VI_body.tex}; \nolinkurl{fr/sga7i/current/expose_VI_body.tex}\par

\par\medskip\noindent\hypertarget{app-F7I-VI-127.18}{}\textbf{F7I-VI-127.18}\quad p. 127; inherited corrected-layer difference.\par
\noindent Imprimé : \emph{(1 espace(s), 0 retour(s) de ligne dans le témoin TeX)}\par
\noindent Corrigé : \emph{(vide)}\par
\noindent Source : ADOPT-F7I-VI-127.18. \nolinkurl{fr/sga7i/diplomatic/expose_VI_body.tex}; \nolinkurl{fr/sga7i/current/expose_VI_body.tex}\par

\par\medskip\noindent\hypertarget{app-F7I-VI-127.19}{}\textbf{F7I-VI-127.19}\quad p. 127; inherited corrected-layer difference.\par
\noindent Imprimé : {\ttfamily\footnotesize i.e.,\allowbreak{}}\par
\noindent Corrigé : {\ttfamily\footnotesize and\allowbreak{} put\allowbreak{}}\par
\noindent Source : ADOPT-F7I-VI-127.19. \nolinkurl{fr/sga7i/diplomatic/expose_VI_body.tex}; \nolinkurl{fr/sga7i/current/expose_VI_body.tex}\par

\par\medskip\noindent\hypertarget{app-F7I-VI-127.20}{}\textbf{F7I-VI-127.20}\quad p. 127; inherited corrected-layer difference.\par
\noindent Imprimé : {\ttfamily\footnotesize a\allowbreak{} }\par
\noindent Corrigé : {\ttfamily\footnotesize T\_1\allowbreak{}}\par
\noindent Source : ADOPT-F7I-VI-127.20. \nolinkurl{fr/sga7i/diplomatic/expose_VI_body.tex}; \nolinkurl{fr/sga7i/current/expose_VI_body.tex}\par

\par\medskip\noindent\hypertarget{app-F7I-VI-127.21}{}\textbf{F7I-VI-127.21}\quad p. 127; inherited corrected-layer difference.\par
\noindent Imprimé : {\ttfamily\footnotesize  v\textasciicircum{}2b\allowbreak{}}\par
\noindent Corrigé : {\ttfamily\footnotesize cI\_n\allowbreak{}}\par
\noindent Source : ADOPT-F7I-VI-127.21. \nolinkurl{fr/sga7i/diplomatic/expose_VI_body.tex}; \nolinkurl{fr/sga7i/current/expose_VI_body.tex}\par

\par\medskip\noindent\hypertarget{app-F7I-VI-127.22}{}\textbf{F7I-VI-127.22}\quad p. 127; inherited corrected-layer difference.\par
\noindent Imprimé : \emph{(1 espace(s), 0 retour(s) de ligne dans le témoin TeX)}\par
\noindent Corrigé : \emph{(vide)}\par
\noindent Source : ADOPT-F7I-VI-127.22. \nolinkurl{fr/sga7i/diplomatic/expose_VI_body.tex}; \nolinkurl{fr/sga7i/current/expose_VI_body.tex}\par

\par\medskip\noindent\hypertarget{app-F7I-VI-127.23}{}\textbf{F7I-VI-127.23}\quad p. 127; inherited corrected-layer difference.\par
\noindent Imprimé : \emph{(vide)}\par
\noindent Corrigé : {\ttfamily\footnotesize F\allowbreak{}}\par
\noindent Source : ADOPT-F7I-VI-127.23. \nolinkurl{fr/sga7i/diplomatic/expose_VI_body.tex}; \nolinkurl{fr/sga7i/current/expose_VI_body.tex}\par

\par\medskip\noindent\hypertarget{app-F7I-VI-127.24}{}\textbf{F7I-VI-127.24}\quad p. 127; inherited corrected-layer difference.\par
\noindent Imprimé : {\ttfamily\footnotesize F\allowbreak{}}\par
\noindent Corrigé : {\ttfamily\footnotesize T\_1X)\textbackslash{}equiv\allowbreak{}}\par
\noindent Source : ADOPT-F7I-VI-127.24. \nolinkurl{fr/sga7i/diplomatic/expose_VI_body.tex}; \nolinkurl{fr/sga7i/current/expose_VI_body.tex}\par

\par\medskip\noindent\hypertarget{app-F7I-VI-127.25}{}\textbf{F7I-VI-127.25}\quad p. 127; inherited corrected-layer difference.\par
\noindent Imprimé : {\ttfamily\footnotesize -\allowbreak{}}\par
\noindent Corrigé : {\ttfamily\footnotesize uF(X)\textbackslash{}pmod\{\textbackslash{}mathfrak\allowbreak{}}\par
\noindent Source : ADOPT-F7I-VI-127.25. \nolinkurl{fr/sga7i/diplomatic/expose_VI_body.tex}; \nolinkurl{fr/sga7i/current/expose_VI_body.tex}\par

\par\medskip\noindent\hypertarget{app-F7I-VI-127.26}{}\textbf{F7I-VI-127.26}\quad p. 127; inherited corrected-layer difference.\par
\noindent Imprimé : {\ttfamily\footnotesize b)\allowbreak{}}\par
\noindent Corrigé : {\ttfamily\footnotesize m\}\$.\allowbreak{}}\par
\noindent Source : ADOPT-F7I-VI-127.26. \nolinkurl{fr/sga7i/diplomatic/expose_VI_body.tex}; \nolinkurl{fr/sga7i/current/expose_VI_body.tex}\par

\par\medskip\noindent\hypertarget{app-F7I-VI-127.27}{}\textbf{F7I-VI-127.27}\quad p. 127; inherited corrected-layer difference.\par
\noindent Imprimé : {\ttfamily\footnotesize =\allowbreak{}}\par
\noindent Corrigé : {\ttfamily\footnotesize By\allowbreak{}}\par
\noindent Source : ADOPT-F7I-VI-127.27. \nolinkurl{fr/sga7i/diplomatic/expose_VI_body.tex}; \nolinkurl{fr/sga7i/current/expose_VI_body.tex}\par

\par\medskip\noindent\hypertarget{app-F7I-VI-127.28}{}\textbf{F7I-VI-127.28}\quad p. 127; inherited corrected-layer difference.\par
\noindent Imprimé : {\ttfamily\footnotesize (v\textasciicircum{}2F\allowbreak{}}\par
\noindent Corrigé : {\ttfamily\footnotesize the\allowbreak{}}\par
\noindent Source : ADOPT-F7I-VI-127.28. \nolinkurl{fr/sga7i/diplomatic/expose_VI_body.tex}; \nolinkurl{fr/sga7i/current/expose_VI_body.tex}\par

\par\medskip\noindent\hypertarget{app-F7I-VI-127.29}{}\textbf{F7I-VI-127.29}\quad p. 127; inherited corrected-layer difference.\par
\noindent Imprimé : {\ttfamily\footnotesize -\allowbreak{}}\par
\noindent Corrigé : {\ttfamily\footnotesize nondegeneracy\allowbreak{}}\par
\noindent Source : ADOPT-F7I-VI-127.29. \nolinkurl{fr/sga7i/diplomatic/expose_VI_body.tex}; \nolinkurl{fr/sga7i/current/expose_VI_body.tex}\par

\par\medskip\noindent\hypertarget{app-F7I-VI-127.30}{}\textbf{F7I-VI-127.30}\quad p. 127; inherited corrected-layer difference.\par
\noindent Imprimé : {\ttfamily\footnotesize a)\$\allowbreak{}}\par
\noindent Corrigé : {\ttfamily\footnotesize condition\allowbreak{}}\par
\noindent Source : ADOPT-F7I-VI-127.30. \nolinkurl{fr/sga7i/diplomatic/expose_VI_body.tex}; \nolinkurl{fr/sga7i/current/expose_VI_body.tex}\par

\par\medskip\noindent\hypertarget{app-F7I-VI-127.31}{}\textbf{F7I-VI-127.31}\quad p. 127; inherited corrected-layer difference.\par
\noindent Imprimé : \emph{(vide)}\par
\noindent Corrigé : {\ttfamily\footnotesize in\allowbreak{} Lemma\allowbreak{} 6.2\allowbreak{}}\par
\noindent Source : ADOPT-F7I-VI-127.31. \nolinkurl{fr/sga7i/diplomatic/expose_VI_body.tex}; \nolinkurl{fr/sga7i/current/expose_VI_body.tex}\par

\par\medskip\noindent\hypertarget{app-F7I-VI-127.32}{}\textbf{F7I-VI-127.32}\quad p. 127; inherited corrected-layer difference.\par
\noindent Imprimé : {\ttfamily\footnotesize so\allowbreak{}}\par
\noindent Corrigé : {\ttfamily\footnotesize the\allowbreak{}}\par
\noindent Source : ADOPT-F7I-VI-127.32. \nolinkurl{fr/sga7i/diplomatic/expose_VI_body.tex}; \nolinkurl{fr/sga7i/current/expose_VI_body.tex}\par

\par\medskip\noindent\hypertarget{app-F7I-VI-127.33}{}\textbf{F7I-VI-127.33}\quad p. 127; inherited corrected-layer difference.\par
\noindent Imprimé : {\ttfamily\footnotesize that\allowbreak{}}\par
\noindent Corrigé : {\ttfamily\footnotesize coefficient\allowbreak{} matrix\allowbreak{} of\allowbreak{} the\allowbreak{} linear\allowbreak{} forms\allowbreak{}}\par
\noindent Source : ADOPT-F7I-VI-127.33. \nolinkurl{fr/sga7i/diplomatic/expose_VI_body.tex}; \nolinkurl{fr/sga7i/current/expose_VI_body.tex}\par

\par\medskip\noindent\hypertarget{app-F7I-VI-127.34}{}\textbf{F7I-VI-127.34}\quad p. 127; inherited corrected-layer difference.\par
\noindent Imprimé : \emph{(vide)}\par
\noindent Corrigé : {\ttfamily\footnotesize \textbackslash{}partial\allowbreak{} F/\textbackslash{}partial\allowbreak{} }\par
\noindent Source : ADOPT-F7I-VI-127.34. \nolinkurl{fr/sga7i/diplomatic/expose_VI_body.tex}; \nolinkurl{fr/sga7i/current/expose_VI_body.tex}\par

\par\medskip\noindent\hypertarget{app-F7I-VI-127.35}{}\textbf{F7I-VI-127.35}\quad p. 127; inherited corrected-layer difference.\par
\noindent Imprimé : {\ttfamily\footnotesize \_i\allowbreak{} \textbackslash{}to\allowbreak{} vX\allowbreak{}}\par
\noindent Corrigé : \emph{(vide)}\par
\noindent Source : ADOPT-F7I-VI-127.35. \nolinkurl{fr/sga7i/diplomatic/expose_VI_body.tex}; \nolinkurl{fr/sga7i/current/expose_VI_body.tex}\par

\par\medskip\noindent\hypertarget{app-F7I-VI-127.36}{}\textbf{F7I-VI-127.36}\quad p. 127; inherited corrected-layer difference.\par
\noindent Imprimé : {\ttfamily\footnotesize establishes\allowbreak{}}\par
\noindent Corrigé : {\ttfamily\footnotesize is\allowbreak{} invertible\allowbreak{} modulo\allowbreak{} \$\textbackslash{}mathfrak\allowbreak{} m\$,\allowbreak{} hence\allowbreak{} over\allowbreak{} \$R\$.\allowbreak{} Suppose\allowbreak{} \$T\_\textbackslash{}nu\textbackslash{}in\textbackslash{}operatorna\allowbreak{}me\{GL\}\_n(R)\$\allowbreak{} satisfies\allowbreak{} \$F(T\_\textbackslash{}nu\allowbreak{} X)\textbackslash{}equiv\allowbreak{} uF(X)\textbackslash{}pmod\{\textbackslash{}mathfrak\allowbreak{} m\textasciicircum{}\textbackslash{}nu\}\$,\allowbreak{} with\allowbreak{} \$\textbackslash{}nu\textbackslash{}geq1\$.\allowbreak{} The\allowbreak{} forms\allowbreak{} \$(\textbackslash{}partial\allowbreak{} F/\textbackslash{}partial\allowbreak{} X\_i)(T\_\textbackslash{}nu\allowbreak{} X)\$\allowbreak{} form\allowbreak{}}\par
\noindent Source : ADOPT-F7I-VI-127.36. \nolinkurl{fr/sga7i/diplomatic/expose_VI_body.tex}; \nolinkurl{fr/sga7i/current/expose_VI_body.tex}\par

\par\medskip\noindent\hypertarget{app-F7I-VI-127.37}{}\textbf{F7I-VI-127.37}\quad p. 127; designated-witness attribution.\par
\noindent Imprimé : {\ttfamily\footnotesize algebra\allowbreak{}}\par
\noindent Corrigé : {\ttfamily\footnotesize basis\allowbreak{} of\allowbreak{} the\allowbreak{} linear\allowbreak{} forms.\allowbreak{} Consequently,\allowbreak{} there\allowbreak{} are\allowbreak{} linear\allowbreak{} forms\allowbreak{} \$L\_i\$\allowbreak{} with\allowbreak{} coefficients\allowbreak{} in\allowbreak{} \$\textbackslash{}mathfrak\allowbreak{} m\textasciicircum{}\textbackslash{}nu\$\allowbreak{} such\allowbreak{} that\allowbreak{} \textbackslash{}[\allowbreak{} E\_\textbackslash{}nu(X):=uF(X)-F(T\_\allowbreak{}\textbackslash{}nu\allowbreak{} X)\allowbreak{} =\textbackslash{}sum\_i\allowbreak{} L\_i(X)\textbackslash{}frac\{\textbackslash{}partial\allowbreak{} F\}\{\textbackslash{}partial\allowbreak{} X\_i\}(T\_\textbackslash{}nu\allowbreak{} X).\allowbreak{} \textbackslash{}]\allowbreak{} Indeed,\allowbreak{} every\allowbreak{} quadratic\allowbreak{} monomial\allowbreak{} is\allowbreak{} a\allowbreak{} product\allowbreak{} of\allowbreak{} two\allowbreak{} linear\allowbreak{} forms,\allowbreak{} and\allowbreak{} expressing\allowbreak{} one\allowbreak{} factor\allowbreak{} in\allowbreak{} this\allowbreak{} basis\allowbreak{} preserves\allowbreak{} its\allowbreak{} coefficient\allowbreak{} ideal.\allowbreak{} Define\allowbreak{} \$H\_\textbackslash{}nu\$\allowbreak{} by\allowbreak{} \$H\_\textbackslash{}nu\allowbreak{} X=(L\_i(X))\_i\$,\allowbreak{} and\allowbreak{} set\allowbreak{} \$T\_\{\textbackslash{}nu+1\}=T\_\textbackslash{}nu+H\_\textbackslash{}\allowbreak{}nu\$.\allowbreak{} The\allowbreak{} quadratic\allowbreak{} identity\allowbreak{} gives\allowbreak{} \textbackslash{}[\allowbreak{} F(T\_\{\textbackslash{}nu+1\}X)=F(T\_\textbackslash{}n\allowbreak{}u\allowbreak{} X)+E\_\textbackslash{}nu(X)+F(H\_\textbackslash{}nu\allowbreak{} X)\allowbreak{} \textbackslash{}equiv\allowbreak{} uF(X)\textbackslash{}pmod\{\textbackslash{}mathfrak\allowbreak{} m\textasciicircum{}\{\textbackslash{}nu+1\}\},\allowbreak{} \textbackslash{}]\allowbreak{} since\allowbreak{} the\allowbreak{} coefficients\allowbreak{} of\allowbreak{} \$F(H\_\textbackslash{}nu\allowbreak{} X)\$\allowbreak{} lie\allowbreak{} in\allowbreak{} \$\textbackslash{}mathfrak\allowbreak{} m\textasciicircum{}\{2\textbackslash{}nu\}\textbackslash{}subseteq\textbackslash{}ma\allowbreak{}thfrak\allowbreak{} m\textasciicircum{}\{\textbackslash{}nu+1\}\$.\allowbreak{} Moreover,\allowbreak{} \$T\_\{\textbackslash{}nu+1\}\textbackslash{}equiv\allowbreak{} T\_\textbackslash{}nu\textbackslash{}pmod\{\textbackslash{}mathfrak\allowbreak{} m\textasciicircum{}\textbackslash{}nu\}\$.\allowbreak{} Completeness\allowbreak{} and\allowbreak{} separatedness\allowbreak{} therefore\allowbreak{} give\allowbreak{} a\allowbreak{} limit\allowbreak{} \$T\textbackslash{}equiv\allowbreak{} cI\_n\textbackslash{}pmod\{\textbackslash{}mathfrak\allowbreak{} m\}\$\allowbreak{} with\allowbreak{} \$F(TX)=uF(X)\$;\allowbreak{} in\allowbreak{} particular,\allowbreak{} \$T\$\allowbreak{} is\allowbreak{} invertible.\allowbreak{} The\allowbreak{} substitution\allowbreak{} \$X\textbackslash{}mapsto\allowbreak{} TX\$\allowbreak{} sends\allowbreak{} \$F-a\$\allowbreak{} to\allowbreak{} \$u(F-b)\$\allowbreak{} and\allowbreak{} induces\allowbreak{} the\allowbreak{} required\allowbreak{}}\par
\noindent Source : WITNESS-F7I-VI-127.37. \nolinkurl{fr/sga7i/diplomatic/expose_VI_body.tex}; \nolinkurl{fr/sga7i/current/expose_VI_body.tex}\par

\par\medskip\noindent\hypertarget{app-F7I-VI-127.38}{}\textbf{F7I-VI-127.38}\quad p. 127; designated-witness attribution.\par
\noindent Imprimé : {\ttfamily\footnotesize \$R[[X\_1,\textbackslash{}ldots,X\_n]]\allowbreak{}/(F\allowbreak{} -\allowbreak{} a)\allowbreak{} \textbackslash{};\textbackslash{}overset\{\textbackslash{}sim\}\{\textbackslash{}re\allowbreak{}lbar\textbackslash{}relbar\textbackslash{}relbar\}\textbackslash{}\allowbreak{};\allowbreak{}}\par
\noindent Corrigé : {\ttfamily\footnotesize \textbackslash{}[\allowbreak{}}\par
\noindent Source : WITNESS-F7I-VI-127.38. \nolinkurl{fr/sga7i/diplomatic/expose_VI_body.tex}; \nolinkurl{fr/sga7i/current/expose_VI_body.tex}\par

\par\medskip\noindent\hypertarget{app-F7I-VI-127.39}{}\textbf{F7I-VI-127.39}\quad p. 127; designated-witness attribution.\par
\noindent Imprimé : \emph{(vide)}\par
\noindent Corrigé : {\ttfamily\footnotesize -a)\allowbreak{} \textbackslash{}xrightarrow\{\textbackslash{}\allowbreak{} \textbackslash{}sim\textbackslash{}\allowbreak{} \}\allowbreak{}}\par
\noindent Source : WITNESS-F7I-VI-127.39. \nolinkurl{fr/sga7i/diplomatic/expose_VI_body.tex}; \nolinkurl{fr/sga7i/current/expose_VI_body.tex}\par

\par\medskip\noindent\hypertarget{app-F7I-VI-127.40}{}\textbf{F7I-VI-127.40}\quad p. 127; designated-witness attribution.\par
\noindent Imprimé : \emph{(vide)}\par
\noindent Corrigé : {\ttfamily\footnotesize R[[X\_1,\textbackslash{}ldots,X\_n]]/\allowbreak{}(F\allowbreak{}}\par
\noindent Source : WITNESS-F7I-VI-127.40. \nolinkurl{fr/sga7i/diplomatic/expose_VI_body.tex}; \nolinkurl{fr/sga7i/current/expose_VI_body.tex}\par

\par\medskip\noindent\hypertarget{app-F7I-VI-127.41}{}\textbf{F7I-VI-127.41}\quad p. 127; designated-witness attribution.\par
\noindent Imprimé : {\ttfamily\footnotesize \$\allowbreak{}}\par
\noindent Corrigé : {\ttfamily\footnotesize .\allowbreak{}}\par
\noindent Source : WITNESS-F7I-VI-127.41. \nolinkurl{fr/sga7i/diplomatic/expose_VI_body.tex}; \nolinkurl{fr/sga7i/current/expose_VI_body.tex}\par

\par\medskip\noindent\hypertarget{app-F7I-VI-127.42}{}\textbf{F7I-VI-127.42}\quad p. 127; designated-witness attribution.\par
\noindent Imprimé : {\ttfamily\footnotesize since\allowbreak{} \$F\$\allowbreak{} is\allowbreak{} a\allowbreak{} quadratic\allowbreak{} form.\allowbreak{}}\par
\noindent Corrigé : {\ttfamily\footnotesize \textbackslash{}]\allowbreak{}}\par
\noindent Source : WITNESS-F7I-VI-127.42. \nolinkurl{fr/sga7i/diplomatic/expose_VI_body.tex}; \nolinkurl{fr/sga7i/current/expose_VI_body.tex}\par

\par\medskip\noindent\hypertarget{app-F7I-VI-127.43}{}\textbf{F7I-VI-127.43}\quad p. 127; inherited corrected-layer difference.\par
\noindent Imprimé : {\ttfamily\footnotesize other\allowbreak{}}\par
\noindent Corrigé : {\ttfamily\footnotesize converse\allowbreak{}}\par
\noindent Source : ADOPT-F7I-VI-127.43. \nolinkurl{fr/sga7i/diplomatic/expose_VI_body.tex}; \nolinkurl{fr/sga7i/current/expose_VI_body.tex}\par

\par\medskip\noindent\hypertarget{app-F7I-VI-127.44}{}\textbf{F7I-VI-127.44}\quad p. 127; inherited corrected-layer difference.\par
\noindent Imprimé : {\ttfamily\footnotesize direction\allowbreak{}}\par
\noindent Corrigé : {\ttfamily\footnotesize follows\allowbreak{}}\par
\noindent Source : ADOPT-F7I-VI-127.44. \nolinkurl{fr/sga7i/diplomatic/expose_VI_body.tex}; \nolinkurl{fr/sga7i/current/expose_VI_body.tex}\par

\par\medskip\noindent\hypertarget{app-F7I-VI-127.45}{}\textbf{F7I-VI-127.45}\quad p. 127; inherited corrected-layer difference.\par
\noindent Imprimé : {\ttfamily\footnotesize is\allowbreak{}}\par
\noindent Corrigé : {\ttfamily\footnotesize from\allowbreak{}}\par
\noindent Source : ADOPT-F7I-VI-127.45. \nolinkurl{fr/sga7i/diplomatic/expose_VI_body.tex}; \nolinkurl{fr/sga7i/current/expose_VI_body.tex}\par

\par\medskip\noindent\hypertarget{app-F7I-VI-127.46}{}\textbf{F7I-VI-127.46}\quad p. 127; inherited corrected-layer difference.\par
\noindent Imprimé : {\ttfamily\footnotesize trivial\allowbreak{}}\par
\noindent Corrigé : {\ttfamily\footnotesize invariance\allowbreak{} of\allowbreak{} \$I\_R\$\allowbreak{} under\allowbreak{} \$R\$-isomorphisms\allowbreak{}}\par
\noindent Source : ADOPT-F7I-VI-127.46. \nolinkurl{fr/sga7i/diplomatic/expose_VI_body.tex}; \nolinkurl{fr/sga7i/current/expose_VI_body.tex}\par

\par\medskip\noindent\hypertarget{app-F7I-VI-128.1}{}\textbf{F7I-VI-128.1}\quad p. 128; misprint\_pass.\par
\noindent Imprimé : {\ttfamily\footnotesize =\allowbreak{}}\par
\noindent Corrigé : {\ttfamily\footnotesize :\allowbreak{}}\par
\noindent Source : 7I:273. \nolinkurl{fixes/FIXLOG_7I_CURRENT.jsonl#719}; \nolinkurl{fixes/MISPRINTS_7I.json#273}\par

\par\medskip\noindent\hypertarget{app-F7I-VI-129.1}{}\textbf{F7I-VI-129.1}\quad p. 128; misprint\_pass.\par
\noindent Imprimé : {\ttfamily\footnotesize n\allowbreak{}}\par
\noindent Corrigé : {\ttfamily\footnotesize m\allowbreak{}}\par
\noindent Source : 7I:274. \nolinkurl{fixes/FIXLOG_7I_CURRENT.jsonl#720}; \nolinkurl{fixes/MISPRINTS_7I.json#274}\par

\par\medskip\noindent\hypertarget{app-F7I-VI-130.1}{}\textbf{F7I-VI-130.1}\quad p. 129; misprint\_in\_print.\par
\noindent Imprimé : {\ttfamily\footnotesize \textbackslash{}overset\allowbreak{}}\par
\noindent Corrigé : {\ttfamily\footnotesize \textbackslash{}xrightarrow\allowbreak{}}\par
\noindent Source : FIX-7I-D-0195. \nolinkurl{fixes/FIXLOG_7I.jsonl#195}\par

\par\medskip\noindent\hypertarget{app-F7I-VI-130.2}{}\textbf{F7I-VI-130.2}\quad p. 129; designated-witness attribution.\par
\noindent Imprimé : {\ttfamily\footnotesize \}\{\textbackslash{}relbar\textbackslash{}joinrel\textbackslash{}re\allowbreak{}lbar\textbackslash{}joinrel\textbackslash{}relbar\textbackslash{}\allowbreak{}joinrel\textbackslash{}relbar\allowbreak{}}\par
\noindent Corrigé : \emph{(vide)}\par
\noindent Source : WITNESS-F7I-VI-130.2. \nolinkurl{fr/sga7i/diplomatic/expose_VI_body.tex}; \nolinkurl{fr/sga7i/current/expose_VI_body.tex}\par

\par\medskip\noindent\hypertarget{app-F7I-VI-131.1}{}\textbf{F7I-VI-131.1}\quad p. 129; misprint\_pass.\par
\noindent Imprimé : {\ttfamily\footnotesize \textbackslash{}lambda\allowbreak{}}\par
\noindent Corrigé : {\ttfamily\footnotesize \textbackslash{}chi\allowbreak{}}\par
\noindent Source : 7I:277. \nolinkurl{fixes/FIXLOG_7I_CURRENT.jsonl#721}; \nolinkurl{fixes/MISPRINTS_7I.json#277}\par

\par\medskip\noindent\hypertarget{app-F7I-VI-132.1}{}\textbf{F7I-VI-132.1}\quad p. 130; misprint\_in\_print.\par
\noindent Imprimé : {\ttfamily\footnotesize \textasciicircum{}0\allowbreak{}}\par
\noindent Corrigé : \emph{(vide)}\par
\noindent Source : FIX-7I-D-0200. \nolinkurl{fixes/FIXLOG_7I.jsonl#200}\par

\par\medskip\noindent\hypertarget{app-F7I-VI-133.1}{}\textbf{F7I-VI-133.1}\quad p. 131; misprint\_pass.\par
\noindent Imprimé : {\ttfamily\footnotesize a\allowbreak{}}\par
\noindent Corrigé : {\ttfamily\footnotesize z\allowbreak{}}\par
\noindent Source : 7I:262. \nolinkurl{fixes/FIXLOG_7I_CURRENT.jsonl#722}\par

\par\medskip\noindent\hypertarget{app-F7I-VI-134.1}{}\textbf{F7I-VI-134.1}\quad p. 131; inherited corrected-layer difference.\par
\noindent Imprimé : \emph{(vide)}\par
\noindent Corrigé : {\ttfamily\footnotesize t\textbackslash{},\allowbreak{}}\par
\noindent Source : ADOPT-F7I-VI-134.1. \nolinkurl{fr/sga7i/diplomatic/expose_VI_body.tex}; \nolinkurl{fr/sga7i/current/expose_VI_body.tex}\par

\par\medskip\noindent\hypertarget{app-F7I-VI-135.1}{}\textbf{F7I-VI-135.1}\quad p. 132; misprint\_pass.\par
\noindent Imprimé : {\ttfamily\footnotesize \textbackslash{}mathop\{\textbackslash{}hphantom\{\textbackslash{}b\allowbreak{}igcap\}\}\allowbreak{}}\par
\noindent Corrigé : {\ttfamily\footnotesize \textbackslash{}prod\allowbreak{}}\par
\noindent Source : 7I:264. \nolinkurl{fixes/FIXLOG_7I_CURRENT.jsonl#723}; \nolinkurl{fixes/MISPRINTS_7I.json#264}\par

\normalsize\fi
\clearpage\SGAtoc{Exposé VII — BIEXTENSIONS DE FAISCEAUX DE GROUPES}
\selectlanguage{french}

\SGApage{I}{133}


S G A 7


\underline{E X P O S E \ VII}


\underline{\underline{BIEXTENSIONS DE FAISCEAUX DE GROUPES}}


\underline{par A. Grothendieck}


\SGAsectiontoc{0. Introduction}{0}
0. \underline{Introduction}

La notion de \underline{biextension} a été introduite par D. MUMFORD [6], dans le contexte des groupes formels, pour exprimer de façon commode les relations entre les groupes formels associés à un schéma abélien et le schéma abélien dual. Dans le présent exposé, nous faisons une étude systématique de cette notion dans le contexte général des faisceaux en groupes sur un site donné, et montrons notamment que cette notion s'insère de façon remarquablement simple dans le formalisme cohomologique classique des $\overset{L}{\otimes}$ et $\mathrm{R}\underline{\mathrm{Hom}}$ (3.6.4 et 3.6.5). Dans l'exposé suivant, nous expliciterons les cas particuliers les plus intéressants, et notamment le cas des biextensions de schémas en groupes par le groupe multiplicatif $\underline{G}_{m}$ (en travaillant avec des sites tels que le site zariskien, ou le site fpqc, ou quelque site intermédiaire) ; on   verra en particulier que si $A$ et $B$ sont deux schémas abéliens sur un schéma $S$, les biextensions de $(A,B)$ par $\underline{G}_{m}$ ont une interprétation remarquable comme étant, essentiellement, les classiques correspondances divisorielles sur $A,B$ . Ainsi se trouve explicité de façon fort commode la signification cohomologique de la notion de correspondance divisorielle (implicite déjà dans le théorème de dualité de Tate des variétés abéliennes sur des corps $p$-adiques). Cette interprétation des correspondances divisorielles est également commode pour "suivre" le destin d'une telle correspondance, quand on fait "dégénérer"
\SGApage{I}{134}
$A$ et $B$ (supposés définis sur le corps des fractions d'un anneau de valuation discrète) à l'aide de la théorie de Néron : les résultats de VIII nous montreront qu'on trouve par exemple une biextension du couple des composantes neutres des modèles de Néron, d'où une biextension des fibres spéciales connexes des modèles de Néron, jouant le rôle d'une spécialisation de la correspondance divisorielle de départ. Ces résultats seront utilisés dans IX, en conjonction avec le théorème de monodromie (I 3) , pour prouver certains résultats sur les modèles de Néron des variétés abéliennes  (le plus important étant le "théorème de réduction semi-stable" pour le modèle de Néron).

Dans tout ce qui suit, on travaille (sauf mention expresse du contraire) dans un $\underline{U}$-topos fixé, noté $\underline{T}$, i.e. (SGA 4 IV) $\underline{T}$ est une catégorie équivalente à la catégorie des faisceaux sur un $\underline{U}$-site convenable (qu'on peut prendre, si on veut, égal à $\underline{T}$ , muni de sa topologie canonique). On identifiera souvent, dans la terminologie, les objets de $\underline{T}$ avec les faisceaux (d'ensembles) sur $\underline{T}$. Les Groupes, extensions de Groupes etc. dont il s'agira seront tous relatifs à ce topos $\underline{T}$.


\SGAsectiontoc{1. Compléments sur les extensions de Groupes}{1}
1. \underline{Compléments sur les extensions de Groupes}

1.1.   Soient $P$ et $G$ deux Groupes de $\underline{T}$, et

\begin{equation}\tag{1.1.1}
\begin{tikzcd}[column sep=large]
1 \arrow[r] \& G \arrow[r, "i"] \& E \arrow[r, "j"] \& P \arrow[r] \& 1
\end{tikzcd}
\end{equation}









une extension de $P$ par $G$ ; donc $E$ est un Groupe, $j : E \longrightarrow P$ est un épimorphisme de Groupes, et $i$ induit un isomorphisme de $G$ avec Ker $j$. Nous allons réinterpréter la donnée d'une telle extension en termes différents, qui nous seront commodes dans la définition et l'étude des biextensions.


\SGApage{I}{135}

1.1.2.  Tout d'abord, l'homomorphisme $i$ permet de faire opérer $G$ de deux façons sur $E$, à gauche et à droite, en posant

\SGAvar{\begin{equation*}
gxg' = i(g)xi(g')
\end{equation*}
pour $g,g' \in G(S)$, $x \in E(S)$ ($S$ un objet quelconque de $\underline{T}$). Si on considère $E$ grâce à $j$ comme un objet "au-dessus de $P$" i.e. comme un objet du topos induit $E_{/P}$ , alors il revient au même de dire que le groupe induit $G_P = G\times P$ par $G$ dans $\underline{T}_{/P}$ opère sur $E$ à gauche et à droite. De plus, il est bien connu que l'une et l'autre opération fait de $E$ un \underline{torseur} (à gauche resp. à droite) sur $P$, de groupe $G_P$ , et que l'on obtient même, plus précisément, un \underline{bitorseur} sous $(G_P,G_P)$. Rappelons [2] que cela signifie que $E$ est un torseur à droite $E_d$ par l'action droite de $G_P$ , et que l'action gauche induit un isomorphisme de $G_P$ avec le Groupe des automorphismes de $E_d$ (= faisceau en groupes des automorphismes de $E$ qui commutent à l'action droite de $G_P$ sur $E$). Cette condition est équivalente à la condition symétrique, savoir que l'action gauche de $G_P$ sur $E$ fait de $E$ un torseur à gauche $E_s$, et que l'action droite induit un isomorphisme du Groupe opposé $G_P^{\;o}$ de $G_P$ avec le Groupe des automorphismes de $E_s$.

1.1.2.1.  Lorsque $\underline{T}$ est le "topos ponctuel" i.e. est la catégorie des ensembles (équivalente à la catégorie des faisceaux sur l'espace topologique ponctuel), donc que $P$ et $G$ sont des groupes ordinaires, alors la donnée d'un bitorseur $E$ sous  $(G_P,G_P)$  équivaut manifestement à la donnée, pour tout $p \in P$, d'un bitorseur $E_p$ sous $(G,G)$ (savoir la fibre de $E$ en $p$). Une interprétation analogue peut-être donnée dans le cas général, par la remarque que la donnée d'un bitorseur $E$ sous $(G_P,G_P)$ équivaut à la donnée, pour tout "point" $p \in P(S)$ ($S$ un objet quelconque de $\underline{T}$), d'un bitorseur $E_p$ sous
\SGApage{I}{136}
le couple de Groupes $(G_S,G_S)$ sur $S$, de façon que la formation de $E_p$ soit "compatible avec les changements de base pour tout morphisme $S' \longrightarrow S$ " ; i.e. si $p' : S' \longrightarrow P$ est le morphisme composé, on suppose donné de plus un isomorphisme de bitorseurs de $P_p$, avec l'image inverse de $P_p$ par $S' \longrightarrow S$, - ces isomorphismes étant soumis à la condition habituelle de transitivité pour le cas d'un composé de deux morphismes $S'' \longrightarrow S' \longrightarrow S$. Quand on explicite cet énoncé en langage canonique (SGA 1 VI    ), il se réduit à la tautologie suivante : la donnée d'un bitorseur sous le couple de Groupes $(G_P,G_P)$ sur $P$ équivaut à une section cartésienne, sur la catégorie $\underline{T}_{/P}$ , de la catégorie fibrée des bitorseurs de groupe $(G_S,G_S)$ sur une base variable $S \in \mathrm{ob}\ \underline{T}_{/P}$ ; plus précisément encore, le foncteur "valeur en $P$" induit une équivalence de la catégorie des sections cartésiennes en question, avec la catégorie des bitorseurs sous $(G_P,G_P)$. Cet énoncé résulte bien entendu trivialement du fait que $P$ est objet final de la catégorie base $\underline{T}_{/P}$ de notre catégorie fibrée ; donc pour \underline{toute} catégorie fibrée sur $\underline{T}_{/P}$ le foncteur "valeur en $S$" est une équivalence entre la catégorie des sections cartésiennes de ladite catégorie fibrée, et la catégorie fibre en $P$. Néanmoins, malgré la nature tautologique, cette interprétation nous sera utile pour exprimer certaines compatibilités dans un langage proche du cas "ensembliste", en utilisant des points à valeurs dans un objet $S$ au lieu des points ordinaires.

}{\begin{equation*}
gxg' = i(g)xi(g')
\end{equation*}
pour $g,g' \in G(S)$, $x \in E(S)$ ($S$ un objet quelconque de $\underline{T}$). Si on considère $E$ grâce à $j$ comme un objet "au-dessus de $P$" i.e. comme un objet du topos induit $\underline{T}_{/P}$ , alors il revient au même de dire que le groupe induit $G_P = G \times P$ par $G$ dans $\underline{T}_{/P}$ opère sur $E$ à gauche et à droite. De plus, il est bien connu que l'une et l'autre opération fait de $E$ un \underline{torseur} (à gauche resp. à droite) sur $P$, de groupe $G_P$ , et que l'on obtient même, plus précisément, un \underline{bitorseur} sous $(G_P,G_P)$. Rappelons [2] que cela signifie que $E$ est un torseur à droite $E_d$ par l'action droite de $G_P$ , et que l'action gauche induit un isomorphisme de $G_P$ avec le Groupe des automorphismes de $E_d$ (= faisceau en groupes des automorphismes de $E$ qui commutent à l'action droite de $G_P$ sur $E$). Cette condition est équivalente à la condition symétrique, savoir que l'action gauche de $G_P$ sur $E$ fait de $E$ un torseur à gauche $E_s$, et que l'action droite induit un isomorphisme du Groupe opposé $G_P^{\;o}$ de $G_P$ avec le Groupe des automorphismes de $E_s$.

1.1.2.1.  Lorsque $\underline{T}$ est le "topos ponctuel" i.e. est la catégorie des ensembles (équivalente à la catégorie des faisceaux sur l'espace topologique ponctuel), donc que $P$ et $G$ sont des groupes ordinaires, alors la donnée d'un bitorseur $E$ sous  $(G_P,G_P)$  équivaut manifestement à la donnée, pour tout $p \in P$, d'un bitorseur $E_p$ sous $(G,G)$ (savoir la fibre de $E$ en $p$). Une interprétation analogue peut être donnée dans le cas général, par la remarque que la donnée d'un bitorseur $E$ sous $(G_P,G_P)$ équivaut à la donnée, pour tout "point" $p \in P(S)$ ($S$ un objet quelconque de $\underline{T}$), d'un bitorseur $E_p$ sous
\SGApage{I}{136}
le couple de Groupes $(G_S,G_S)$ sur $S$, de façon que la formation de $E_p$ soit "compatible avec les changements de base pour tout morphisme $S' \longrightarrow S$ " ; i.e. si $p' : S' \longrightarrow P$ est le morphisme composé, on suppose donné de plus un isomorphisme de bitorseurs de $E_{p'}$, avec l'image inverse de $E_p$ par $S' \longrightarrow S$, - ces isomorphismes étant soumis à la condition habituelle de transitivité pour le cas d'un composé de deux morphismes $S'' \longrightarrow S' \longrightarrow S$. Quand on explicite cet énoncé en langage canonique (SGA 1 VI    ), il se réduit à la tautologie suivante : la donnée d'un bitorseur sous le couple de Groupes $(G_P,G_P)$ sur $P$ équivaut à une section cartésienne, sur la catégorie $\underline{T}_{/P}$ , de la catégorie fibrée des bitorseurs de groupe $(G_S,G_S)$ sur une base variable $S \in \mathrm{ob}\ \underline{T}_{/P}$ ; plus précisément encore, le foncteur "valeur en $P$" induit une équivalence de la catégorie des sections cartésiennes en question, avec la catégorie des bitorseurs sous $(G_P,G_P)$. Cet énoncé résulte bien entendu trivialement du fait que $P$ est objet final de la catégorie base $\underline{T}_{/P}$ de notre catégorie fibrée ; donc pour \underline{toute} catégorie fibrée sur $\underline{T}_{/P}$ le foncteur "valeur en $P$" est une équivalence entre la catégorie des sections cartésiennes de ladite catégorie fibrée, et la catégorie fibre en $P$. Néanmoins, malgré la nature tautologique, cette interprétation nous sera utile pour exprimer certaines compatibilités dans un langage proche du cas "ensembliste", en utilisant des points à valeurs dans un objet $S$ au lieu des points ordinaires.

}{F7I-VII-001}1.1.3.  La structure de bitorseur 1.1.2 sous $(G_P,G_P)$ sur $E$, en termes de la structure d'extension (1.1.1), n'utilisait qu'une partie de la structure multiplicative de $E$, savoir la multiplication de couples de "points" dont l'un au moins provient d'un point de $G$ (via $i$). Nous allons maintenant interpréter la structure multiplicative de $E$ comme une structure supplémen
\SGApage{I}{137}
taire sur le bitorseur envisagé. Pour ceci, considérons d'abord deux sections $p$, $p'$ de $P$, d'où, avec les notations de 1.1.2.1 , deux bitorseurs $E_p$ , $E_{p'}$ sous $(G,G)$, savoir les images inverses dans $E$ desdites sections $p,p' : e \longrightarrow P$ (NB $e$ désigne l'objet final de $\underline{T}$). Il est trivial alors que la multiplication de $E$ induit un morphisme d'objets de $\underline{T}$

\SGAvar{\begin{equation*}
\varphi : \quad P_p \times P_{p'} \ \longrightarrow\ P_{pp'}
\end{equation*}
qui satisfait, en termes des structures gauches et droites des trois bitorseurs en jeu (et en prenant comme d'habitude des points à valeurs dans un objet $S \in \mathrm{ob}\ \underline{T}$ arbitraire)

\begin{equation}\tag{1.1.3.0}
\left\{\ \begin{array}{l} \varphi(gx,x') = g\ \varphi(x,x')\ ,\ \ \varphi(x,x'g') = \varphi(x,x')g' \\[3ex] \qquad\qquad\qquad \varphi(xg,x') = \varphi(x,gx') \end{array}\right.
\end{equation}
$(g,g' \in G(S)$, $x \in P_p(S)$, $x' \in P_{p'}(S))$. La troisième de ces relations s'interprète, en termes du produit fibré contracté [2]

\begin{equation*}
P_p \overset{G}{\times} P_{p'} \qquad ,
\end{equation*}
en disant que le morphisme $\varphi$ se factorise en

\begin{equation*}
P_p \times P_{p'} \ \longrightarrow\ P_p \overset{G}{\times} P_{p'} \ \xrightarrow{\ \varphi_{p,p'}\ }\ P_{pp'} \qquad ,
\end{equation*}
où la première flèche est le morphisme canonique de passage au quotient. D'autre part, les deux premières relations signifient que le morphisme

\begin{equation}\tag{1.1.3.1}
\varphi_{p,p'} \ :\ P_p \overset{G}{\times} P_{p'} \ \longrightarrow\ P_{pp'}
\end{equation}
qu'on vient d'obtenir est un homomorphisme de \underline{bitorseurs}, compte tenu de la structure de bitorseur sur le premier membre pour laquelle les opérations
\SGApage{I}{138}
gauches (resp. droites) de $G$ proviennent des opérations gauches (resp. droites) de $G$ sur le premier facteur $P_p$ (resp. sur le deuxième facteur $P_{p'}$). Notons que (comme tout homomorphisme de bitorseurs sous un couple de Groupes donné) l'homomorphisme $\varphi_{p,p'}$ est nécessairement un \underline{isomorphisme}.

}{\begin{equation*}
\varphi : \quad E_p \times E_{p'} \ \longrightarrow\ E_{pp'}
\end{equation*}
qui satisfait, en termes des structures gauches et droites des trois bitorseurs en jeu (et en prenant comme d'habitude des points à valeurs dans un objet $S \in \mathrm{ob}\ \underline{T}$ arbitraire)

\begin{equation}\tag{1.1.3.0}
\left\{\ \begin{array}{l} \varphi(gx,x') = g\ \varphi(x,x')\ ,\ \ \varphi(x,x'g') = \varphi(x,x')g' \\[3ex] \qquad\qquad\qquad \varphi(xg,x') = \varphi(x,gx') \end{array}\right.
\end{equation}
$(g,g' \in G(S)$, $x \in E_p(S)$, $x' \in E_{p'}(S))$. La troisième de ces relations s'interprète, en termes du produit fibré contracté [2]

\begin{equation*}
E_p \overset{G}{\times} E_{p'} \qquad ,
\end{equation*}
en disant que le morphisme $\varphi$ se factorise en

\begin{equation*}
E_p \times E_{p'} \ \longrightarrow\ E_p \overset{G}{\times} E_{p'} \ \xrightarrow{\ \varphi_{p,p'}\ }\ E_{pp'} \qquad ,
\end{equation*}
où la première flèche est le morphisme canonique de passage au quotient. D'autre part, les deux premières relations signifient que le morphisme

\begin{equation}\tag{1.1.3.1}
\varphi_{p,p'} \ :\ E_p \overset{G}{\times} E_{p'} \ \longrightarrow\ E_{pp'}
\end{equation}
qu'on vient d'obtenir est un homomorphisme de \underline{bitorseurs}, compte tenu de la structure de bitorseur sur le premier membre pour laquelle les opérations
\SGApage{I}{138}
gauches (resp. droites) de $G$ proviennent des opérations gauches (resp. droites) de $G$ sur le premier facteur $E_p$ (resp. sur le deuxième facteur $E_{p'}$). Notons que (comme tout homomorphisme de bitorseurs sous un couple de Groupes donné) l'homomorphisme $\varphi_{p,p'}$ est nécessairement un \underline{isomorphisme}.

}{F7I-VII-002}La définition de l'isomorphisme précédent (1.1.3.1) s'étend aussitôt au cas où $p,p'$, au lieu d'être des sections de $P$, sont des points de $P$ à valeurs dans un objet quelconque $S$ de $\underline{T}$ : il suffit d'appliquer ce qui précède au topos induit $\underline{T}_{/S}$ et à l'extension de Groupes sur $S$ déduite de (1.1.1) par changement de base . On constate alors que la formation des isomorphismes $\varphi_{p,p'}$ précédents est "compatible avec tout changement de base" $S' \longrightarrow S$ (i.e. correspond à un homomorphisme de sections cartésiennes de la catégorie fibrée qu'on pense sur la catégorie $\underline{T}_{/P\times P}$ des doubles flèches $p,p' : S \rightrightarrows P$ de but $P$). La donnée du système des $\varphi_{p,p'}$ satisfaisant cette condition revient encore, bien entendu, à la donnée de $\varphi_{p,p'}$ dans le cas "universel", correspondant au cas où $S = P\times P$ et où $p,p'$ sont respectivement les deux projections $\mathrm{pr}_i$ ($i=1,2$) de $P\times P$ dans $P$. Il s'interprète alors comme un isomorphisme de bitorseurs sous $(G_{P\times P},G_{P\times P})$ :

\begin{equation}\tag{1.1.3.2}
\varphi \ :\ \mathrm{pr}_1^*(E)\ \ \mathrm{pr}_2^*(E)\ \overset{\sim}{\longrightarrow}\ \pi^*(E) \qquad ,
\end{equation}
où

\begin{equation*}
\pi \ :\ P\times P \ \longrightarrow\ P
\end{equation*}
est le morphisme d'objets de $\underline{T}$ définissant la structure multiplicative du Groupe $P$. C'est cependant sous la forme des données (1.1.3.1) que les propriétés de compatibilité qui vont suivre s'énoncent le plus commodément, en évitant le recours à des diagrammes moins proches de l'intuition.


\SGApage{I}{139}

\SGAvar{1.1.4.  En écrivant les relations (1.1.3.0), qui nous ont permis d'exprimer la loi de composition donnée sur $E$ en termes d'isomorphismes de bitorseurs (1.1.3.1), on a utilisé l'associativité de cette loi de composition seulement dans le cas du produit de trois "points" de $E$ dont un au moins provient d'un point de $G$. Si on veut exprimer l'associativité complète en termes des $\varphi_{p,p'}$ , on est amené à écrire, pour trois points $p,p',p"$ de $P$ (à valeurs dans un objet $S$ de $\underline{T}$), que le diagramme suivant est commutatif :

}{1.1.4.  En écrivant les relations (1.1.3.0), qui nous ont permis d'exprimer la loi de composition donnée sur $E$ en termes d'isomorphismes de bitorseurs (1.1.3.1), on a utilisé l'associativité de cette loi de composition seulement dans le cas du produit de trois "points" de $E$ dont un au moins provient d'un point de $G$. Si on veut exprimer l'associativité complète en termes des $\varphi_{p,p'}$ , on est amené à écrire, pour trois points $p,p',p''$ de $P$ (à valeurs dans un objet $S$ de $\underline{T}$), que le diagramme suivant est commutatif :

}{F7I-VII-003}\begin{equation}\tag{1.1.4.1}
\begin{tikzcd}[column sep=huge, row sep=large]
 \& E_{pp'}E_{p''} \arrow[dr, "{\varphi_{pp',p''}}"] \& \\
E_{p}E_{p'}E_{p''} \arrow[ur, "{\varphi_{p,p'}\wedge id}"] \arrow[dr, "{id\wedge\varphi_{p',p''}}"'] \& \& E_{pp'p''} \\
 \& E_{p}E_{p'p''} \arrow[ur, "{\varphi_{p,p'p''}}"'] \&
\end{tikzcd}\quad .
\end{equation}


















\SGAvar{Dans ce diagramme, nous avons écrit le produit contracté de bitorseurs sans signes $\overset{G_S}{\times}$ , étant entendu que la multiplication contractéede deux bitorseurs $Q.R$ se fait en utilisant (pour le passage au quotient dans $Q$$\times$$R$) l'opération du groupe structural $G$ \underline{à droite sur} $Q$ \underline{et à gauche sur} $R$. On sait par ailleurs que cette opération de produit contracté estassociative à un isomorphisme canonique près [2], et nous faisons l'identification $(E_{p}E_{p'})E_{p''} \simeq E_{p}(E_{p'}E_{p''})$, en écrivant simplement les deux membres comme $E_{p}E_{p'}E_{p''}$. Les conventions évidentes de cette nature seront utilisées dans la suite de cet exposé, sans être expressément signalées à chaque fois.

1.1.4.2.  De même que la donnée des isomorphismes $\varphi_{p,p'}$ (1.1.3.1) se ramène au cas "universel" $S=P$$\times$$P$, $p=\mathrm{pr}_1,p' = \mathrm{pr}_2$ envisagé dans (1.1.3.2), de même
\SGApage{I}{140}
pour la vérification des conditions d'associativité (1.1.4.1), on peut encore se borner au cas "universel" où $S=P\times P\times P$, $p=\mathrm{pr}_1$,$p'=\mathrm{pr}_2$,$p"=\mathrm{pr}_3$ . Nous laissons au lecteur le soin d'écrire le diagramme en question avec les notations de son choix.

}{Dans ce diagramme, nous avons écrit le produit contracté de bitorseurs sans signes $\overset{G_S}{\times}$ , étant entendu que la multiplication contractée de deux bitorseurs $Q.R$ se fait en utilisant (pour le passage au quotient dans $Q$$\times$$R$) l'opération du groupe structural $G$ \underline{à droite sur} $Q$ \underline{et à gauche sur} $R$. On sait par ailleurs que cette opération de produit contracté est associative à un isomorphisme canonique près [2], et nous faisons l'identification $(E_{p}E_{p'})E_{p''} \simeq E_{p}(E_{p'}E_{p''})$, en écrivant simplement les deux membres comme $E_{p}E_{p'}E_{p''}$. Les conventions évidentes de cette nature seront utilisées dans la suite de cet exposé, sans être expressément signalées à chaque fois.

1.1.4.2.  De même que la donnée des isomorphismes $\varphi_{p,p'}$ (1.1.3.1) se ramène au cas "universel" $S=P$$\times$$P$, $p=\mathrm{pr}_1,p' = \mathrm{pr}_2$ envisagé dans (1.1.3.2), de même
\SGApage{I}{140}
pour la vérification des conditions d'associativité (1.1.4.1), on peut encore se borner au cas "universel" où $S=P\times P\times P$, $p=\mathrm{pr}_1$,$p'=\mathrm{pr}_2$,$p''=\mathrm{pr}_3$ . Nous laissons au lecteur le soin d'écrire le diagramme en question avec les notations de son choix.

}{F7I-VII-004}1.1.5.  Réciproquement, les Groupes $P$ et $G$ de $\underline{T}$ étant supposés donnés, supposons donné de plus un bitorseur

\SGAvar{\begin{equation}\tag{1.1.5.1}
j \; : \; E \longrightarrow P
\end{equation}
sous $(G_{P},G_{P})$ (NB les deux actions de $G_{P}$ sur $E$, ou ce qui revient au même, de $G$ sur $P$, sont sous-entendues dans les notations), et d'autre part un système d'isomorphismes de bitorseurs sur une base $S$ variable

}{\begin{equation}\tag{1.1.5.1}
j \; : \; E \longrightarrow P
\end{equation}
sous $(G_{P},G_{P})$ (NB les deux actions de $G_{P}$ sur $E$, ou ce qui revient au même, de $G$ sur $E$, sont sous-entendues dans les notations), et d'autre part un système d'isomorphismes de bitorseurs sur une base $S$ variable

}{F7I-VII-005}\begin{equation}\tag{1.1.5.2}
\varphi_{p,p'} \; : \; E_{p}E_{p'} \;\overset{\sim}{\longrightarrow}\; E_{pp'} \qquad (p,p' \in P(S)),
\end{equation}
compatibles avec changement de base $S' \longrightarrow S$ ; en d'autres termes, un isomorphisme de bitorseurs (1.1.3.2). Supposons que les diagrammes correspondants (1.1.4.1) sont commutatifs. Je dis qu'\underline{il existe sur} $E$ \underline{une unique structure d'extension de} $P$ \underline{par} $G$, \underline{telle que les données précédentes se déduisent de cette structure par les constructions de} 1.1.2 \underline{et} 1.1.3.

Notons d'abord que lorsqu'on prend dans (1.1.5.2) pour $p,p'$ la section unité $1$ de $P$, on trouve un isomorphisme

\begin{equation}\tag{1.1.5.3}
\varphi_{1,1} \; : \; E_{1} \;\overset{\sim}{\longleftarrow}\; E_{1} \overset{G}{\times} E_{1} \qquad ,
\end{equation}
d'où, en prenant le produit contracté (à gauche, disons) des deux membres avec le bitorseur inverse $E_{1}^{-1}$ de $E_{1}$ [2], un isomorphisme canonique

\begin{equation}\tag{1.1.5.4}
E_{1} \;\simeq\; \text{bitorseur unité } {}_{s}G_{d}
\end{equation}
( $= G$ opérant sur lui-même par translations gauche et droite), compte tenu
\SGApage{I}{141}
de l'isomorphisme canonique

\begin{equation*}
E_1^{-1}.E_1 \;\; \simeq \;\; \text{bitorseur unité } {}_sG_d
\end{equation*}
de loc.cit (valable pour tout bitorseur). Il revient au même de dire qu'on a une section $1_E$ de $E_1$ donc de $E$ :

\begin{equation}\tag{1.1.5.5}
1_E \in \Gamma(E_1) = \mathrm{Hom}(e,E_1) \subset \Gamma(E)
\end{equation}
(savoir l'image par (1.1.5.4) de la section marquée de ${}_sG_d$), satisfaisant la relation

\begin{equation*}
g.1_E = 1_E.g \qquad (g \in G(S)) \quad ,
\end{equation*}
et par construction cette section est caractérisée par la condition

\begin{equation}\tag{1.1.5.6}
\varphi_{1,1}(1_E,1_E) = 1_E \qquad .
\end{equation}
D'autre part, la donnée d'un morphisme de torseurs $\varphi$ (1.1.3.2) équivaut manifestement à la donnée d'un morphisme

\begin{equation}\tag{1.1.5.7}
\pi_E : E\times E \longrightarrow E \quad , \qquad \text{noté aussi } (x,y) \longmapsto xy
\end{equation}
satisfaisant les conditions

\begin{equation}\tag{1.1.5.8}
j(xy) = j(x)j(y) \qquad (x,y \in E(S))
\end{equation}
(i.e. $j$ est compatible avec les structures multiplicatives $\pi_E$ resp. $\pi$ sur $E$ resp. $P$), et les trois relations (1.1.3.0) (en commençant de préférence par la dernière de celles-ci). Enfin, on définit un morphisme d'objets de \underline{$T$}

\begin{equation*}
i : G \longrightarrow E
\end{equation*}
par la formule

\begin{equation}\tag{1.1.5.9}
i(g) = g.1_E = 1_E.g 
\end{equation}

\SGApage{I}{142}

\SGAvar{ (l'identité des deux derniers membres dans cette formule ayant déjà été noté avec la construction de $1_E$). Ceci posé, je dis que \underline{la loi de multiplication} $\pi_E$ \underline{sur} $E$ \underline{fait de} $E$ \underline{un} Groupe, \underline{admettant} $1_E$ \underline{comme section unité}, \underline{que} $j : E \rightarrow P$ \underline{et} $i : G \rightarrow E$ \underline{sont des homomorphismes de Groupes qui font du Groupe} $E$ \underline{une extension de} $P$ \underline{par} $G$.

}{ (l'identité des deux derniers membres dans cette formule ayant déjà été notée avec la construction de $1_E$). Ceci posé, je dis que \underline{la loi de multiplication} $\pi_E$ \underline{sur} $E$ \underline{fait de} $E$ \underline{un} Groupe, \underline{admettant} $1_E$ \underline{comme section unité}, \underline{que} $j : E \rightarrow P$ \underline{et} $i : G \rightarrow E$ \underline{sont des homomorphismes de Groupes qui font du Groupe} $E$ \underline{une extension de} $P$ \underline{par} $G$.

}{F7I-VII-006}En effet, l'associativité de la loi de composition de $E$ n'est autre que la commutativité des diagrammes (1.1.4.1). La construction de $\pi_E$ via les isomorphismes $\varphi_{p,p'}$ (1.1.3.1) montre d'autre part que pour tout $x \in E(S)$, les translations à gauche et à droite par $x$ dans $E(S)$ sont injectives. Comme $1_E1_E = 1_E$ par (1.1.5.6), on en conclut aussitôt que $1_E$ est une unité bilatère de $E$. Nous laissons au lecteur la vérification de l'existence d'un inverse d'un $x \in E(S)$ sur $p \in P(S)$, utilisant un isomorphisme canonique déduit de $\varphi_{p,p^{-1}}$ :

\SGAvar{\begin{equation}\tag{1.1.5.10}
P_{p^{-1}} \;\simeq\; (P_p)^{o} \qquad .
\end{equation}
Donc $E$ est bien un \underline{Groupe}. On a déjà signalé que $j$ est compatible avec les lois de composition (1.1.5.8), et il en est de même de $i$, car on aura, pour $g,g' \in G(S)$ :

}{\begin{equation}\tag{1.1.5.10}
E_{p^{-1}} \;\simeq\; (E_p)^{o} \qquad .
\end{equation}
Donc $E$ est bien un \underline{Groupe}. On a déjà signalé que $j$ est compatible avec les lois de composition (1.1.5.8), et il en est de même de $i$, car on aura, pour $g,g' \in G(S)$ :

}{F7I-VII-007}\begin{multline*}
i(g)i(g') = (g.1_E)(g'.1_E) = (g.1_E.g')\ 1_E \quad (\text{dernière relation} \\ \text{de } (1.1.3.0)) = \ g.1_E.g' \quad (\text{car } 1_E \text{ est unité}) = g.(1_E.g') = g.(g'.1_E) \\ (1.1.5.9) = (gg').1_E = i(gg'). \qquad .
\end{multline*}
Comme $j$ est la projection structurale d'un torseur, c'est un épimorphisme d'objets de \underline{$T$}, il reste à voir que $i$ induit un isomorphisme de $G$ avec $Ker j = E_1$, ce qui est clair par la définition (1.1.5.9) et la construction de $1_E$ .


\SGApage{I}{143}

1.1.6.  Nous admettrons que les raisonnements précédents prouvent que pour un topos fixé \underline{$T$}, les constructions précédentes fournissent une \underline{équivalence entre la catégorie des extensions de Groupes de} \underline{$T$} (1.1.1) (où $P$ et $G$ comme $E$ sont supposés variables), \underline{et la catégorie des systèmes} ($P,G,E,\varphi$), \underline{où} $P$ \underline{et} $G$ \underline{sont des Groupes de} \underline{$T$}, $E$ \underline{un bitorseur sous} $(G_P,G_P)$ \underline{sur} $P$, \underline{enfin} $\varphi$ \underline{un isomorphisme de bitorseurs} (1.1.3.2) (ou si on préfère, un système d'isomorphismes de bitorseurs (1.1.3.1) compatibles aux changements de base), \underline{satisfaisant la condition d'associativité exprimée par la commutativité du diagramme} (1.1.4.2) (ou si on préfère, des diagrammes (1.1.4.1)).

Nous laissons au lecteur le soin de faire l'explicitation évidente, en termes du dictionnaire précédent, de la notion d'image inverse (resp. directe) d'une extension (1.1.1) par un morphisme de Groupes

\begin{equation*}
P' \longrightarrow P \qquad (\text{resp. } G \longrightarrow G'\ ) \qquad ,
\end{equation*}
ou par ces deux opérations simultanément. Nous lui laissons également le soin d'expliciter une compatibilité de l'équivalence de catégories précédente avec le foncteur image inverse $f^*$ par un morphisme de topos

\begin{equation*}
f : \underline{T}' \longrightarrow \underline{T} \qquad ;
\end{equation*}
pour ceci, il y a lieu d'utiliser la donnée $\varphi$ sous la forme (1.1.3.2) (de préférence aux $\varphi_{p,p'}$ de (1.1.3.1)), qui définit alors une donnée correspondante $f^*(\varphi)$ sur le topos $\underline{T}'$.

\SGAvar{1.2. \underline{Cas commutatif}. Dans toute la suite à partir de 1.4, nous ne nous intéresserons qu'aux extensions \underline{commutatives} de Groupes. Nous supposerons donc $P$ et $G$ commutatifs, et nous proposons ici d'expliciter, en termes
\SGApage{I}{144}
de la description ($P,G,E,\varphi$) d'une extension en termes de bitorseurs, le fait que $E$ est un Groupe commutatif. Une première condition évidemment nécessaire (et qui exprime en fait celle que $i(G)$ soit un sous-Groupe \underline{central} de $E$) est que les bitorseurs $E_p$ ($p \in P(S)$) soient des torseurs ordinaires, i.e. que les opérations gauches et droites de $G$ sur $P_p$ coîncident. Il revient évidemment au même de dire que $E$ sur $P$ est un torseur ordinaire sous $G_P$ . (Il sera donc inutile dorénavant de parler de bitorseurs, et on ne regardera plus que des torseurs ordinaires.) Cette condition préliminaire étant supposée remplie, la condition de commutativité de $E$ s'exprime alors par la commutativité des diagrammes suivants (où $p,p' \in P(S)$) :

}{1.2. \underline{Cas commutatif}. Dans toute la suite à partir de 1.4, nous ne nous intéresserons qu'aux extensions \underline{commutatives} de Groupes. Nous supposerons donc $P$ et $G$ commutatifs, et nous proposons ici d'expliciter, en termes
\SGApage{I}{144}
de la description ($P,G,E,\varphi$) d'une extension en termes de bitorseurs, le fait que $E$ est un Groupe commutatif. Une première condition évidemment nécessaire (et qui exprime en fait celle que $i(G)$ soit un sous-Groupe \underline{central} de $E$) est que les bitorseurs $E_p$ ($p \in P(S)$) soient des torseurs ordinaires, i.e. que les opérations gauches et droites de $G$ sur $E_p$ coïncident. Il revient évidemment au même de dire que $E$ sur $P$ est un torseur ordinaire sous $G_P$ . (Il sera donc inutile dorénavant de parler de bitorseurs, et on ne regardera plus que des torseurs ordinaires.) Cette condition préliminaire étant supposée remplie, la condition de commutativité de $E$ s'exprime alors par la commutativité des diagrammes suivants (où $p,p' \in P(S)$) :

}{F7I-VII-008}\begin{equation}\tag{1.2.1}
\begin{tikzcd}[column sep=large, row sep=large]
E_p\ E_{p'} \arrow[r, "{\varphi_{p,p'}}"] \arrow[d, "\mathrm{sym}"'] \& E_{pp'} \arrow[d, "\mathrm{id}"] \\
E_{p'}E_p \arrow[r, "{\varphi_{p',p}}"] \& E_{p'p}
\end{tikzcd}\quad ,
\end{equation}














où la première flèche verticale est l'isomorphisme canonique de commutativité (défini grâce au fait que $G$ est maintenant commutatif ; on a de plus $pp' = p'p$ puisque $P$ est commutatif). Cette condition est évidemment équivalente à la condition correspondante dans la situation universelle S=P$\times$P, $p=pr_1$,$p' = pr_2$, que nous nous dispensons de réécrire ici.


1.3.  \underline{Relations avec les torseurs rigidifiés ou birigidifiés}

1.3.1.0.  Soient $P$ un objet de $T$ muni d'une section $e_P$, $G$ un Groupe de \underline{$T$}, $E$ un torseur sous $G_P$ . On appelle \underline{rigidification de} $E$ (relativement à la section donnée $e_P$ de $P$) une trivialisation du torseur $E_1 = e_P^*(E)$ sous $G$,
\SGApage{I}{145}
i.e. une section de ce torseur. La terminologie adoptée est inspirée par le cas où $G$ est commutatif et où l'on a

\begin{equation}\tag{1.3.1.1}
\mathrm{Hom}_{pt}(P,G) \;\simeq\; (e) \qquad ,
\end{equation}
où le premier membre désigne l'ensemble des morphismes transformant $e_P$ en la section unité $e_G$ de $G$, -relation qui équivaut au fait que l'homomorphisme naturel $\Gamma(G) \longrightarrow \mathrm{Hom}(P,G)$, associant à une section $g$ de $G$ le morphisme constant $P \longrightarrow G$ de valeur $g$, est un isomorphisme :

\begin{equation}\tag{1.3.1.2}
\Gamma(G) \;\xrightarrow{\;\sim\;}\; \mathrm{Hom}(P,G) \qquad .
\end{equation}
On voit en effet immédiatement que les relations équivalentes précédentes équivalent encore à dire que tout automorphisme du torseur $E$ qui respecte la rigidification donnée est l'automorphisme identique (puisque les automorphismes en question correspondent précisément aux éléments du premier membre de (1.3.1.1)) ; en particulier, moyennant la condition (1.3.1.1), la catégorie des torseurs sous $G_P$ rigidifiés relativement à $e_P$ est "rigide".

\SGAvar{1.3.2.  Soit de plus $Q$ un objet de \underline{$T$} muni d'une section $e_Q$ , et désignons maintenant par $E$ un torseur sous $G_{P\times Q}$ . Comme $e_Q$ définit une section $e_{P\times Q/P}$ de $P\times Q$ sur $P$, on peut considérer sur $E$ les rigidifications $\alpha$ relativement à cette section (en se plaçant dans le topos induit $\underline{T}_{/P}$), qui s'explicitent comme les sections du torseur induit par $E$ sur $P\times e_Q$. Echangeant le rôle de $P$ et $Q$, on peut considérer les rigidifications $\beta$ sur $E$ relatives à la section $e_{P\times Q/Q}$ de $P\times Q$ sur $Q$ définie par $e_P$. Deux telles rigidifications partielles sont dites \underline{compatibles} si les sections $\alpha$ et $\beta$ coîncident sur $e_P\times e_Q$ , i.e. définissent une même section du torseur $E_{1,1}$ sous $G$, image
\SGApage{I}{146}
inverse de $E$ par la section $(e_P,e_Q)$ de $P\times Q$. On appelle aussi \underline{birigidification} du torseur $E$ relativement au couple de sections $(e_P,e_Q)$, un couple de deux rigidifications partielles compatibles. Si $G$ est commutatif, on voit tout de suite que pour que le torseur $E$ admette une birigidification $(\alpha,\beta)$, il faut et il suffit qu'il admette des rigidifications partielles $\alpha$ et $\beta$, et on peut trouver alors une birigidification dans laquelle $\alpha$ (ou $\beta$) est choisie à l'avance ; en effet, il suffit alors de corriger $\beta$ par le morphisme $Q \longrightarrow G$ composé $Q \longrightarrow e \xrightarrow{\lambda} G$, où la deuxième flèche $\lambda$ est définie par la propriété que $\beta e_Q = (\alpha e_P)\lambda$. Toujours dans le cas commutatif, notons que si $E,E'$ sont deux torseurs birigidifiés sous $G_{P\times Q}$, il revient au même de dire qu'ils sont isomorphes comme $G_{P\times Q}$-torseurs tout court, ou comme $G_{P\times Q}$ - torseurs birigidifiés ; en d'autres termes, s'il existe un isomorphisme $u : E \longrightarrow E'$ de $G_{P\times Q}$-torseurs, on peut trouver un isomorphisme $v : E \longrightarrow E'$ de $G_{P\times Q}$-torseurs qui respecte les birigidifications. En effet, regardant l'écart de $u$ pour les birigidifications données, on trouve des morphismes $f:P \longrightarrow G$ et $g : Q \longrightarrow G$ tels que $fe_P=ge_Q$, soit $a \in G(e)$ ; il existe alors un morphisme $P\times Q \longrightarrow G$ dont les composés avec $e_{P\times Q/P}$ et $e_{P\times Q/Q}$ soient respectivement $f$ et $g$ (savoir $h(p,q) = f(p)g(q)a^{-1}$), et on peut prendre $v = uh$ .

1.3.3.  La terminologie de 1.3.1 s'étend aussitôt au cas où on a un \underline{bitorseur} $E$ sous $(G_P,G_P)$, étant entendu qu'une rigidification du bitorseur $E$ doit être un isomorphisme de $e_P(E)$ avec le bitorseur trivial $_sG_d$ , i.e. est défini par une section $e_E$ de $e_P(E)$ qui est \underline{centrale}, i.e. qui satisfait à la condition $g.e_E = e_E.g$ ($g \in G(S)$. On étend de même la terminologie des
\SGApage{I}{147}
birigidifications au cas des bitorseurs.

}{1.3.2.  Soit de plus $Q$ un objet de \underline{$T$} muni d'une section $e_Q$ , et désignons maintenant par $E$ un torseur sous $G_{P\times Q}$ . Comme $e_Q$ définit une section $e_{P\times Q/P}$ de $P\times Q$ sur $P$, on peut considérer sur $E$ les rigidifications $\alpha$ relativement à cette section (en se plaçant dans le topos induit $\underline{T}_{/P}$), qui s'explicitent comme les sections du torseur induit par $E$ sur $P\times e_Q$. Echangeant le rôle de $P$ et $Q$, on peut considérer les rigidifications $\beta$ sur $E$ relatives à la section $e_{P\times Q/Q}$ de $P\times Q$ sur $Q$ définie par $e_P$. Deux telles rigidifications partielles sont dites \underline{compatibles} si les sections $\alpha$ et $\beta$ coïncident sur $e_P\times e_Q$ , i.e. définissent une même section du torseur $E_{1,1}$ sous $G$, image
\SGApage{I}{146}
inverse de $E$ par la section $(e_P,e_Q)$ de $P\times Q$. On appelle aussi \underline{birigidification} du torseur $E$ relativement au couple de sections $(e_P,e_Q)$, un couple de deux rigidifications partielles compatibles. Si $G$ est commutatif, on voit tout de suite que pour que le torseur $E$ admette une birigidification $(\alpha,\beta)$, il faut et il suffit qu'il admette des rigidifications partielles $\alpha$ et $\beta$, et on peut trouver alors une birigidification dans laquelle $\alpha$ (ou $\beta$) est choisie à l'avance ; en effet, il suffit alors de corriger $\beta$ par le morphisme $Q \longrightarrow G$ composé $Q \longrightarrow e \xrightarrow{\lambda} G$, où la deuxième flèche $\lambda$ est définie par la propriété que $\beta e_Q = (\alpha e_P)\lambda$. Toujours dans le cas commutatif, notons que si $E,E'$ sont deux torseurs birigidifiés sous $G_{P\times Q}$, il revient au même de dire qu'ils sont isomorphes comme $G_{P\times Q}$-torseurs tout court, ou comme $G_{P\times Q}$ - torseurs birigidifiés ; en d'autres termes, s'il existe un isomorphisme $u : E \longrightarrow E'$ de $G_{P\times Q}$-torseurs, on peut trouver un isomorphisme $v : E \longrightarrow E'$ de $G_{P\times Q}$-torseurs qui respecte les birigidifications. En effet, regardant l'écart de $u$ pour les birigidifications données, on trouve des morphismes $f:P \longrightarrow G$ et $g : Q \longrightarrow G$ tels que $fe_P=ge_Q$, soit $a \in G(e)$ ; il existe alors un morphisme $P\times Q \longrightarrow G$ dont les composés avec $e_{P\times Q/P}$ et $e_{P\times Q/Q}$ soient respectivement $f$ et $g$ (savoir $h(p,q) = f(p)g(q)a^{-1}$), et on peut prendre $v = uh$ .

1.3.3.  La terminologie de 1.3.1 s'étend aussitôt au cas où on a un \underline{bitorseur} $E$ sous $(G_P,G_P)$, étant entendu qu'une rigidification du bitorseur $E$ doit être un isomorphisme de $e_P^*(E)$ avec le bitorseur trivial $_sG_d$ , i.e. est défini par une section $e_E$ de $e_P^*(E)$ qui est \underline{centrale}, i.e. qui satisfait à la condition $g.e_E = e_E.g$ ($g \in G(S)$). On étend de même la terminologie des
\SGApage{I}{147}
birigidifications au cas des bitorseurs.

}{F7I-VII-009}1.3.4.  Revenons au cas où on se donne deux Groupes $P$ et $G$  et choisissons la section unité $e_P$ pour définir des rigidifications de bitorseurs sur $P$, et des birigidifications de bitorseurs sur $P\times P$. On a des foncteurs naturels

\begin{equation}\tag{1.3.4.1}
\mathbf{EXT}(P;\ G) \xrightarrow{\ i\ } \mathrm{BITORSRIG}(P,G) \xrightarrow{\ \delta\ } \mathbf{BITORSBIRIG}(P,P;G) \quad ,
\end{equation}
où les trois expressions écrites désignent respectivement la catégorie des extensions de Groupes de $P$ par $G$, la catégorie des bitorseurs sous $(G_P,G_P)$ rigidifiés pour $e_P$, et la catégorie des bitorseurs sous $(G_{P\times P},G_{P\times P})$ birigidifiés pour $(e_P,e_P)$. On définit le foncteur $i$ en associant à l'extension $E$ de $P$ par $G$ le bitorseur correspondant sur $P$ (1.1.6), rigidifié par la section unité de $E$. On définit le foncteur $\delta$ par la formule

\begin{equation}\tag{1.3.4.2}
\delta(E) = \pi^*(E)\mathrm{pr}_1^*(E)^{-1}\mathrm{pr}_2^*(E)^{-1} \quad ,
\end{equation}
où $\pi$ : $P\times P \longrightarrow P$ est la loi de composition de $P$, $\mathrm{pr}_1$ et $\mathrm{pr}_2$ sont les deux projections. On a une \underline{birigidification canonique sur} $\delta(E)$. En effet, la restriction de $\delta(E)$ à $P\times e_P$ est canoniquement isomorphe à $\mathrm{id}_P^*(E)\mathrm{id}_P^*(E)^{-1}e^*(E)^{-1}$, où $e = \mathrm{pr}_2\mathrm{inj}_1$ est le morphisme constant $P \longrightarrow P$ de valeur $e_P$ , de sorte que grâce à la rigidification donnée de $E$, on a un isomorphisme $e^*(E) \xrightarrow{\ \sim\ }$ bitorseur trivial, d'où un isomorphisme de la restriction de $\delta(E)$ à $P\times e_P$ avec le bitorseur trivial. On définit de façon symétrique une rigidification partielle relativement à $e_P\times P$, et on vérifie trivialement que ces deux rigidifications partielles sont compatibles.


\SGApage{I}{148}

\SGAvar{1.3.4.3.  On observera que le composé $\delta i$ des foncteurs $\delta$ et $i$ de (1.3.4.1) est canoniquement isomorphe au foncteur trivial, i.e. au foncteur constant de valeur la bitorseur birigidifié trivial sur $P\times P$. En effet, pour toute extension $E$ de $P$ par $G$, l'isomorphisme $\varphi$ de (1.1.3.2) définit évidemment une trivialisation de $\delta i(E)$ compatible avec sa birigidification canonique.

\underline{Proposition} 1.3.5. \underline{Soient} $P$ \underline{et} $Q$ \underline{deux Groupes de} $\underline{T}$ .

a)  \underline{Supposons que tout morphisme de faisceaux d'ensembles de} $P\times P$ \underline{dans} $G$ \underline{qui est trivial sur} $P\times e_P$ \underline{et sur} $e_P\times P$ \underline{soit trivial}. \underline{Alors le foncteur} $i$ \underline{de} (1.3.4.1) \underline{est pleinement fidèle}. \underline{En particulier}, \underline{sur un bitorseur} $E$ \underline{sous} $(G_P,G_P)$, \underline{biridifié relativement à} $e_P$ \underline{i.e.} \underline{muni d'une section centrale} $e_E$ \underline{de} $E_1 = e_P^*(E)$, \underline{il y a au plus une structure d'extension de} $P$ \underline{par} $G$, \underline{compatible avec sa structure de bitorseur et admettant} $e_E$ \underline{comme section unité}.

b)  \underline{Supposons de plus que tout morphisme de faisceaux d'ensembles de} $P\times P\times P$ \underline{dans} $G$ \underline{qui est trivial sur les produits partiels} $P\times P\times e_P$ , $P\times e_P\times P$ \underline{et} $e_P\times P\times P$ \underline{soit trivial}. \underline{Soit} $E$ \underline{un torseur sous} $G_P$ . \underline{Pour qu'il existe sur} $E$ \underline{une structure d'extension de} $P$ \underline{par} $G$ \underline{compatible avec sa structure de bitorseur}, \underline{il faut et il suffit que le bitorseur sous} $(G_{P\times P},G_{P\times P})$ \underline{sous-jacent à} $\delta(E) = \pi^*(E)\mathrm{pr}_1^*(E)^{-1}\mathrm{pr}_2^*(E)^{-1}$ \underline{soit trivial}. \underline{Lorsqu'il en est ainsi}, \underline{pour toute section centrale} $e_E$ de $E_1 = e_P^*(E)$, \underline{i.e.} \underline{toute rigidification de} $E$ \underline{relativement à} $e_P$, \underline{il existe une structure d'extension unique sur} $E$, \underline{compatible avec sa structure de bitorseur et admettant} $e_E$ \underline{comme section unité}, \underline{i.e.} $E$ , \underline{considéré comme bitorseur rigidifié relativement à} $e_P$, \underline{est dans l'image essentielle de} $i$ . \underline{D'autre part}, \underline{il existe une et une seule trivialisation du bitorseur birigidifié} $\delta(E)$.


}{1.3.4.3.  On observera que le composé $\delta i$ des foncteurs $\delta$ et $i$ de (1.3.4.1) est canoniquement isomorphe au foncteur trivial, i.e. au foncteur constant de valeur le bitorseur birigidifié trivial sur $P\times P$. En effet, pour toute extension $E$ de $P$ par $G$, l'isomorphisme $\varphi$ de (1.1.3.2) définit évidemment une trivialisation de $\delta i(E)$ compatible avec sa birigidification canonique.

\underline{Proposition} 1.3.5. \underline{Soient} $P$ \underline{et} $G$ \underline{deux Groupes de} $\underline{T}$ .

a)  \underline{Supposons que tout morphisme de faisceaux d'ensembles de} $P\times P$ \underline{dans} $G$ \underline{qui est trivial sur} $P\times e_P$ \underline{et sur} $e_P\times P$ \underline{soit trivial}. \underline{Alors le foncteur} $i$ \underline{de} (1.3.4.1) \underline{est pleinement fidèle}. \underline{En particulier}, \underline{sur un bitorseur} $E$ \underline{sous} $(G_P,G_P)$, \underline{rigidifié relativement à} $e_P$ \underline{i.e.} \underline{muni d'une section centrale} $e_E$ \underline{de} $E_1 = e_P^*(E)$, \underline{il y a au plus une structure d'extension de} $P$ \underline{par} $G$, \underline{compatible avec sa structure de bitorseur et admettant} $e_E$ \underline{comme section unité}.

b)  \underline{Supposons de plus que tout morphisme de faisceaux d'ensembles de} $P\times P\times P$ \underline{dans} $G$ \underline{qui est trivial sur les produits partiels} $P\times P\times e_P$ , $P\times e_P\times P$ \underline{et} $e_P\times P\times P$ \underline{soit trivial}. \underline{Soit} $E$ \underline{un bitorseur sous} $(G_P,G_P)$ . \underline{Pour qu'il existe sur} $E$ \underline{une structure d'extension de} $P$ \underline{par} $G$ \underline{compatible avec sa structure de bitorseur}, \underline{il faut et il suffit que le bitorseur sous} $(G_{P\times P},G_{P\times P})$ \underline{sous-jacent à} $\delta(E) = \pi^*(E)\mathrm{pr}_1^*(E)^{-1}\mathrm{pr}_2^*(E)^{-1}$ \underline{soit trivial}. \underline{Lorsqu'il en est ainsi}, \underline{pour toute section centrale} $e_E$ de $E_1 = e_P^*(E)$, \underline{i.e.} \underline{toute rigidification de} $E$ \underline{relativement à} $e_P$, \underline{il existe une structure d'extension unique sur} $E$, \underline{compatible avec sa structure de bitorseur et admettant} $e_E$ \underline{comme section unité}, \underline{i.e.} $E$ , \underline{considéré comme bitorseur rigidifié relativement à} $e_P$, \underline{est dans l'image essentielle de} $i$ . \underline{D'autre part}, \underline{il existe une et une seule trivialisation du bitorseur birigidifié} $\delta(E)$.


}{F7I-VII-010}\SGApage{I}{149}

En termes du diagramme (1.3.4.1), on peut donc dire que l'image essentielle du foncteur pleinement fidèle $i$ est "le noyau" du foncteur $\delta$, ou encore, de façon imagée, que le diagramme (1.3.4.1) est une "suite exacte de catégories ponctuées", - et même de catégories à structures de pseudo-groupes (dans la terminologie introduite dans 2.5.1 ci-dessous).

\SGAvar{Démontrons 1.3.5.  a)  Soient $E,E'$ deux extensions de $P$ par $G$, $f : E \longrightarrow E'$ un morphisme des bitorseurs rigidifiés sous-jacents, montrons que $f$ est un morphisme d'extensions, ou ce qui revient manifestement au même, un morphisme de Groupes. Il faut prouver que l'on a $f(xy) = f(x)f(y)$, or (désignant par $j' : E' \longrightarrow P$ le morphisme structural du torseur $E'$ sur $P$) on aura $j'f(xy) = j'(f(x)f(y))$, car les deux membres sont respectivement égaux à $j(f(xy))$ et à $j(f(x)f(y))$, qui sont égaux puisque $j$ est multiplicatif. On a donc

}{Démontrons 1.3.5.  a)  Soient $E,E'$ deux extensions de $P$ par $G$, $f : E \longrightarrow E'$ un morphisme des bitorseurs rigidifiés sous-jacents, montrons que $f$ est un morphisme d'extensions, ou ce qui revient manifestement au même, un morphisme de Groupes. Il faut prouver que l'on a $f(xy) = f(x)f(y)$, or (désignant par $j' : E' \longrightarrow P$ le morphisme structural du torseur $E'$ sur $P$) on aura $j'f(xy) = j'(f(x)f(y))$, car les deux membres sont respectivement égaux à $j(xy)$ et à $j(x)j(y)$, qui sont égaux puisque $j$ est multiplicatif. On a donc

}{F7I-VII-011}\begin{equation*}
f(xy) = u(j(x),j(y))f(x)f(y) \qquad ,
\end{equation*}
où

\begin{equation*}
u : P \times P \longrightarrow G
\end{equation*}
est un morphisme bien déterminé. Faisant $x=1$ puis $y=1$ dans la relation précédente, on trouve que l'on a (compte tenu de $f(1)=1$)

\begin{equation*}
u(p,1) = u(1,p) = 1 \qquad ,
\end{equation*}
i.e. $u$ est trivial sur $P \times e_P$ et sur $e_P \times P$, d'où il résulte par hypothèse que $u$ est trivial, donc $f(xy) = f(x)f(y)$, cqfd.

\SGAvar{b)  La nécessité de la condition envisagée sur $E$ pour qu'il admette une structure d'extension compatible avec sa structure de bitorseur étant triviale (1.3.4.3), il reste à prouver sa suffisance pour prouver la première assertion de b). Or, utilisant une trivialisation $\varphi$
\SGApage{I}{150}
de $\delta(E)$, on en déduit par image inverse par la section unité de $P \times P$ une trivialisation de $E_1 = e_P(E)$ (comparer 1.5.1), i.e. une rigidification de $E$ relativement à $e_P$ . Supposons-la fixée. Une trivialisation donnée $\varphi$ de $\delta(E)$ n'est pas nécessairement compatible avec cette birigidification ; ces restrictions à $P \times e_P$ et à $e_P \times P$ s'expriment, en termes des rigidifications partielles naturelles de $\delta(E)$, par des morphismes \underline{centraux}

}{b)  La nécessité de la condition envisagée sur $E$ pour qu'il admette une structure d'extension compatible avec sa structure de bitorseur étant triviale (1.3.4.3), il reste à prouver sa suffisance pour prouver la première assertion de b). Or, utilisant une trivialisation $\varphi$
\SGApage{I}{150}
de $\delta(E)$, on en déduit par image inverse par la section unité de $P \times P$ une trivialisation de $E_1 = e_P^*(E)$ (comparer 1.5.1), i.e. une rigidification de $E$ relativement à $e_P$ . Supposons-la fixée. Une trivialisation donnée $\varphi$ de $\delta(E)$ n'est pas nécessairement compatible avec cette birigidification ; ces restrictions à $P \times e_P$ et à $e_P \times P$ s'expriment, en termes des rigidifications partielles naturelles de $\delta(E)$, par des morphismes \underline{centraux}

}{F7I-VII-012}\begin{equation*}
f : P \longrightarrow G \ , \ g : P \longrightarrow G \ ,
\end{equation*}
ayant mêmes restrictions à la section unité de $P$ : $f(1) = g(1) = a \in \Gamma G$ . Alors le morphisme \underline{central}

\SGAvar{\begin{equation*}
h : P \times P \longrightarrow G \quad , \quad h(p,q) = f(p)g(q)a^{-1}
\end{equation*}
a comme restrictions à $P \times e_P$ et $e_P \times P$ les morphismes $f$ et $g$. Par suite, corrigeant $\varphi$ par $h$ (i.e. le remplaçant par $\varphi(p,q)h(p,q)^{-1}$), on trouve une trivialisation de $\delta(E)$ compatible avec sa birigidification. On notera d'ailleurs qu'une telle trivialisation est déterminée modulo un morphisme central $P \times P \longrightarrow G$ qui est trivial sur $P \times e_P$ et sur $e_P \times P$, donc qui est trivial, donc cette $\varphi$ est unique, ce qui est la dernière assertion de 1.3.5 b). La démonstration sera terminée si nous prouvons que cette $\varphi$, considérée comme un isomorphisme de la forme (1.1.3.2), satisfait à la condition d'associativité (1.1.4.1), compte tenu que l'on sait déjà par le choix de $\varphi$ que $e_E$ est une unité bilatère pour la loi de composition définie par $\varphi$ . Or, les deux composés dans (1.1.4.1) sont des isomorphismes de bitorseurs sous $(G_{P \times P \times P} , G_{P \times P \times P})$, donc différent par un morphisme central

}{\begin{equation*}
h : P \times P \longrightarrow G \quad , \quad h(p,q) = f(p)g(q)a^{-1}
\end{equation*}
a comme restrictions à $P \times e_P$ et $e_P \times P$ les morphismes $f$ et $g$. Par suite, corrigeant $\varphi$ par $h$ (i.e. le remplaçant par $\varphi(p,q)h(p,q)^{-1}$), on trouve une trivialisation de $\delta(E)$ compatible avec sa birigidification. On notera d'ailleurs qu'une telle trivialisation est déterminée modulo un morphisme central $P \times P \longrightarrow G$ qui est trivial sur $P \times e_P$ et sur $e_P \times P$, donc qui est trivial, donc cette $\varphi$ est unique, ce qui est la dernière assertion de 1.3.5 b). La démonstration sera terminée si nous prouvons que cette $\varphi$, considérée comme un isomorphisme de la forme (1.1.3.2), satisfait à la condition d'associativité (1.1.4.1), compte tenu que l'on sait déjà par le choix de $\varphi$ que $e_E$ est une unité bilatère pour la loi de composition définie par $\varphi$ . Or, les deux composés dans (1.1.4.1) sont des isomorphismes de bitorseurs sous $(G_{P \times P \times P} , G_{P \times P \times P})$, donc diffèrent par un morphisme central

}{F7I-VII-013}\begin{equation*}
u : P \times P \times P \longrightarrow G \qquad ,
\end{equation*}
grâce à la formule

\begin{equation*}
(xy)z = u(jx,jy,jz) \ x(yz) \qquad .
\end{equation*}

\SGApage{I}{151}

\SGAvar{Faisant  dans cette formule successivement $y=z=1$, puis $z=x=1$, puis $x=y=1$, et utilisant le fait que $1$ est unité bilatère pour la loi de composition envisagée, on trouve que $u$ est trivial sur les trois morceaux $P\times e_P\times e_P$ , $e_P\times P\times e_P$ , $e_P\times e_P\times P$, ce qui implique par hypothèse qu'il est trivial, donc que $(xy)z = x(yz)$, ce qui achève la démonstration.

}{Faisant dans cette formule successivement $z=1$, puis $x=1$, puis $y=1$, et utilisant le fait que $1$ est unité bilatère pour la loi de composition envisagée, on trouve que $u$ est trivial sur les trois morceaux $P\times P\times e_P$ , $e_P\times P\times P$ , $P\times e_P\times P$, ce qui implique par hypothèse qu'il est trivial, donc que $(xy)z = x(yz)$, ce qui achève la démonstration.

}{F7I-VII-014}\underline{Corollaire} 1.3.6. \underline{Supposons que tout morphisme de faisceaux d'ensembles} $P\times G \longrightarrow G$ \underline{qui est trivial sur} $P\times e_G$ \underline{et sur} $e_P\times G$ \underline{soit trivial}. \underline{Alors}, \underline{si} $G$ \underline{est commutatif}, $G$ \underline{est central dans toute extension de Groupes de} $P$ \underline{par} $G$. \underline{Si} $P$ \underline{est également commutatif}, \underline{et si l'hypothèse de} 1.3.5 a) \underline{est vérifiée}, \underline{alors toute extension de Groupes} $E$ \underline{de} $P$ \underline{par} $G$ \underline{est commutative}.

\SGAvar{Pour la première assertion, il faut prouver que l'opération de $P$ sur $G$ définie par l'extension $G$ est triviale. Or cette opération s'exprime par un morphisme $u : P\times G \longrightarrow G$, tel que les restrictions à $P\times e_G$ et à $e_P\times G$ du morphisme correspondant $(p,g) \longmapsto u(p,g)g^{-1}$ sont évidemment triviales, donc ce dernier morphisme est trivial i.e. on a $u(p,g) = g$, ce qu'on voulait établir. Notons maintenant que si $E$ est une extension de $P$ par $G$, alors le Groupe opposé $E^o$ à $E$, qui est une extension de $P^o$ par $G^o$, a comme structure de bitorseur sous-jacente $i(E^o)$ celle déduite du bitorseur $i(E)$ en y échangeant les opérations gauches et droites. Lorsque $E$ est une extension centrale, on a donc $i(E)=i(E^o)$, et si de plus $P$ est commutatif i.e. $P=P^o$, on trouve par l'assertion d'unicité 1.3.5 a) que $E=E^o$ i.e. que $E$ est commutatif. Cela établit le corollaire.


}{Pour la première assertion, il faut prouver que l'opération de $P$ sur $G$ définie par l'extension $E$ est triviale. Or cette opération s'exprime par un morphisme $u : P\times G \longrightarrow G$, tel que les restrictions à $P\times e_G$ et à $e_P\times G$ du morphisme correspondant $(p,g) \longmapsto u(p,g)g^{-1}$ sont évidemment triviales, donc ce dernier morphisme est trivial i.e. on a $u(p,g) = g$, ce qu'on voulait établir. Notons maintenant que si $E$ est une extension de $P$ par $G$, alors le Groupe opposé $E^o$ à $E$, qui est une extension de $P^o$ par $G^o$, a comme structure de bitorseur sous-jacente $i(E^o)$ celle déduite du bitorseur $i(E)$ en y échangeant les opérations gauches et droites. Lorsque $E$ est une extension centrale, on a donc $i(E)=i(E^o)$, et si de plus $P$ est commutatif i.e. $P=P^o$, on trouve par l'assertion d'unicité 1.3.5 a) que $E=E^o$ i.e. que $E$ est commutatif. Cela établit le corollaire.


}{F7I-VII-015}\SGApage{I}{152}

1.3.7. Un cas particulier intéressant où toutes les hypothèses envisagées dans 1.3.5 et 1.3.6 sont satisfaites est celui où on a

\begin{equation}\tag{1.3.7.1}
\underline{\mathrm{Hom}}_{pt}(P,G) \;=\; e \ \text{(faisceau final)} \qquad ,
\end{equation}
où le premier membre désigne le faisceau des morphismes de faisceaux d'ensembles ponctués de $P$ dans $G$. Cette condition s'écrit encore, indépendamment des sections marquées, sous la forme (analogue à (1.3.1.2))

\begin{equation}\tag{1.3.7.2}
G \;\simeq\; \underline{\mathrm{Hom}}_{ens}(P,G) \qquad ,
\end{equation}
où le deuxième membre désigne le faisceau des morphismes pour les faisceaux d'ensembles sous-jacents à $P$ et $G$. Ces conditions signifient aussi que pour tout objet $Q$ de $\underline{T}$, l'application $u \longmapsto u\circ pr_2$

\begin{equation}\tag{1.3.7.3}
\mathrm{Hom}_{ens}(Q,G) \;\xrightarrow{\ \sim\ }\; \mathrm{Hom}_{ens}(P\times Q,G)
\end{equation}
est bijective. On voit que cette propriété de $P$ (pour $G$ fixé) est stable par produits cartésiens, donc en particulier qu'elle implique la même propriété pour les produits cartésiens $P\times P$, $P\times P\times P$ etc. Il est immédiat dès lors qu'elle implique les hypothèses faites dans 1.3.5 et dans 1.3.6 ; de plus, elle implique que les catégories qui interviennent dans (1.3.4.1) sont toutes \underline{rigides}, i.e. que tout automorphisme d'un objet d'une de ces catégories est l'identité. De plus, comme la condition (1.3.7.1) sur $P,G$ est stable par tout changement de base $S \longrightarrow e$, il s'ensuit aussitôt par descente, du résultat d'unicité 1.3.5 a), que pour qu'un bitorseur rigidifié $E$ sur $P$ corresponde bien à une extension de $P$ par $G$, il suffit qu'il en soit ainsi \underline{localement sur} $e$, ce qui (en vertu de 1.3.5 b)) s'exprime aussi par la nullité de la section

\begin{equation*}
\underline{\delta}(E) \in \Gamma\, R^1 h_*(G_{P\times P}) 
\end{equation*}

\SGApage{I}{153}



 définie par $\delta(E)$, où

\SGAvar{\begin{equation*}
h : P\times P \longrightarrow e
\end{equation*}
est le morphisme structural.

}{\begin{equation*}
h : P \times P \longrightarrow e
\end{equation*}
est le morphisme structural.

}{}1.3.8.  Soit $p : P \longrightarrow S$ un morphisme de schémas, tel que $p$ soit \underline{cohérent} (i.e. quasi-compact et quasi-séparé) \underline{et plat}, et tel qu'on ait

\begin{equation}\tag{1.3.8.1}
p_*(\underline{O}_P) \xleftarrow{\ \sim\ } \underline{O}_S \quad .
\end{equation}
Utilisant le fait que (grâce à la cohérence de $p$) la formation de $p_*(\underline{O}_P)$ commute à tout changement de base plat, on constate aussitôt que les conditions envisagées sur le schéma relatif $P/S$ sont stables par produits cartésiens, et sont vérifiées en particulier pour les puissances cartésiennes $P\times P$, $P\times P\times P$, ... Elles sont également stables par tout changement de base plat.

\SGAvar{Si maintenant $G$ est un schéma \underline{affine} sur $S$, il résulte aussitôt de (1.3.8.1) que la relation (1.3.7.2) est vérifiée, puisque l'on a $\mathrm{Hom}_S(P,G) = \mathrm{Hom}_{alg}(\underline{A},p_*(\underline{O}_P)) \simeq \mathrm{Hom}_{alg}(\underline{A},\underline{O}_S)$, si $\underline{A}$ est l'Algèbre quasi-cohérente sur $S$ telle que $G = Spec(A)$. D'après ce qu'on vient de dire, la même conclusion s'applique aux produits $P\times P$, $P\times P\times P$, et de même, si $G$ est plat sur $S$, la même conclusion s'applique au schéma relatif $P_G$ sur $G$ vis-à-vis du schéma relatif $G_G$ sur $G$.

}{Si maintenant $G$ est un schéma \underline{affine} sur $S$, il résulte aussitôt de (1.3.8.1) que la relation (1.3.7.2) est vérifiée, puisque l'on a $\mathrm{Hom}_S(P,G) = \mathrm{Hom}_{alg}(\underline{A},p_*(\underline{O}_P)) \simeq \mathrm{Hom}_{alg}(\underline{A},\underline{O}_S)$, si $\underline{A}$ est l'Algèbre quasi-cohérente sur $S$ telle que $G = \operatorname{Spec}(\underline{A})$. D'après ce qu'on vient de dire, la même conclusion s'applique aux produits $P\times P$, $P\times P\times P$, et de même, si $G$ est plat sur $S$, la même conclusion s'applique au schéma relatif $P_G$ sur $G$ vis-à-vis du schéma relatif $G_G$ sur $G$.

}{F7I-VII-017}Lorsque $P$ et $G$ sont des schémas en groupes sur $S$, et qu'on travaille avec une topologie sur $(\mathrm{Sch})_{/S}$ moins fine que la topologie canonique, ce qui permet d'identifier $P$ et $G$ à des faisceaux sur le site précédent, on voit que tout morphisme de $P$, $P\times P$, $P\times P\times P$ ... dans $G$, compatible avec les sections marquées, est trivial ; de même si $G$ est plat, 

\SGApage{I}{154}



 tout $S$-morphisme $P\times G$ $\longrightarrow$ $G$ compatible avec les sections marquées est trivial, comme on voit en l'interprétant comme un $G$-morphisme de $P_G$ dans $G_G$ . Cela montre donc que l'on est sous les conditions d'application de 1.3.5, et également de 1.3.6 lorsque $G$ est plat. De plus, les trois catégories intervenant dans (1.3.4.1) sont rigides.

\SGAvar{1.3.8.2.  Le cas le plus intéressant où la condition (1.3.8.1) est vérifiée, et reste vérifiée après toute extension de la base, est celui où $P$ est un \underline{schéma abélien} sur $S$ ; on est donc même sous les conditions de 1.3.7. Dans ce cas, les conclusions de 3.5 et de 3.6 sont donc valables pour tout Groupe $G$ affine sur $S$. C'est le cas initialement envisagé par J.P. SERRE, dont les développements précédents sont visiblement inspirés. Un cas intéressant où (1.3.8.1) est vérifié sans l'être universellement est celui où $P$ est le \underline{modèle de Néron} d'un schéma abélien sur le corps des fractions d'un anneau de valuation discrète de spectre $S$.


}{1.3.8.2.  Le cas le plus intéressant où la condition (1.3.8.1) est vérifiée, et reste vérifiée après toute extension de la base, est celui où $P$ est un \underline{schéma abélien} sur $S$ ; on est donc même sous les conditions de 1.3.7. Dans ce cas, les conclusions de 1.3.5 et de 1.3.6 sont donc valables pour tout Groupe $G$ affine sur $S$. C'est le cas initialement envisagé par J.P. SERRE, dont les développements précédents sont visiblement inspirés. Un cas intéressant où (1.3.8.1) est vérifié sans l'être universellement est celui où $P$ est le \underline{modèle de Néron} d'un schéma abélien sur le corps des fractions d'un anneau de valuation discrète de spectre $S$.


}{F7I-VII-018}1.4.  \underline{Extensions d'un Groupe commutatif libre}

Soit $I$ un objet de $\underline{T}$, et prenons

\begin{equation}\tag{1.4.1}
P = \mathbb{Z}[I] \quad ,
\end{equation}
le $\mathbb{Z}$-Module libre engendré par $I$ (SGA 4 IV 11 ). A toute extension $E$ de $P$ par un $\mathbb{Z}$-Module donné $G$ (NB il ne sera plus par la suite question que d'extensions commutatives de groupes), associons le torseur correspondant sous $G_P$, puis sa restriction à $I$ par le morphisme canonique

\begin{equation*}
I \longrightarrow P = \mathbb{Z}[I] \quad .
\end{equation*}

\SGApage{I}{155}

On trouve ainsi un foncteur canonique

\begin{equation}\tag{1.4.2}
\mathrm{EXT}(P,G) \ \longrightarrow\ \mathrm{TORS}(I,G_I) \qquad .
\end{equation}
Je dis que ce foncteur est une \underline{équivalence de catégories}. Il revient évidemment au même de dire que pour tout $G$, les applications

\begin{equation}\tag{1.4.3}
\mathrm{Hom}_{\mathbb{Z}}(P,G) \ \longrightarrow\ H^0(I,G_I) = G(I)
\end{equation}
\begin{equation}\tag{1.4.4}
\mathrm{Ext}^1_{\mathbb{Z}}(P,G) \ \longrightarrow\ H^1(I,G_I)
\end{equation}
sont des isomorphismes. En fait, plus généralement, lorsque $A$ est un Anneau (pas nécessairement commutatif) de \underline{$T$} et $P = \underline{A}[I]$ le $A$-Module libre engendré par $I$, on trouve des isomorphismes canoniques fonctoriels en le $A$-Module variable $G$

\SGAvar{\begin{equation}\tag{1.4.5}
\mathrm{Ext}^i_A(A[I],\ G) \ \xrightarrow{\ \sim\ }\ H^i(I,G_I) \qquad ,
\end{equation}
qui pour $A=\mathbb{Z}$ et $i=0,1$ coîncident avec (1.4.3) et (1.4.4). Pour définir les isomorphismes (1.4.5), on se rappelle que le premier membre est le foncteur dérivé du foncteur $G \longmapsto \mathrm{Hom}_A(A[I],G)$, qui par définition de $A[I]$ est isomorphe au foncteur $G \longmapsto G(I) = H^0(I,G_I)$. Comme $G \longmapsto G_I$ transforme Modules injectifs en Modules injectifs de $\underline{T}_{/I}$ (SGA 4 V     ), on en déduit la conclusion voulue. On laisse au lecteur le soin de vérifier la compatibilité annoncée avec les homomorphismes (1.4.3) et (1.4.4).

}{\begin{equation}\tag{1.4.5}
\mathrm{Ext}^i_A(A[I],\ G) \ \xrightarrow{\ \sim\ }\ H^i(I,G_I) \qquad ,
\end{equation}
qui pour $A=\mathbb{Z}$ et $i=0,1$ coïncident avec (1.4.3) et (1.4.4). Pour définir les isomorphismes (1.4.5), on se rappelle que le premier membre est le foncteur dérivé du foncteur $G \longmapsto \mathrm{Hom}_A(A[I],G)$, qui par définition de $A[I]$ est isomorphe au foncteur $G \longmapsto G(I) = H^0(I,G_I)$. Comme $G \longmapsto G_I$ transforme Modules injectifs en Modules injectifs de $\underline{T}_{/I}$ (SGA 4 V     ), on en déduit la conclusion voulue. On laisse au lecteur le soin de vérifier la compatibilité annoncée avec les homomorphismes (1.4.3) et (1.4.4).

}{F7I-VII-019}On notera que le foncteur 1.4.2 est compatible aux images inverses par un morphisme de topos

\begin{equation*}
f : \underline{T}' \ \longrightarrow\ \underline{T} \qquad .
\end{equation*}

\SGApage{I}{156}


\SGAsectiontoc{2. La notion de biextension  de faisceaux abéliens}{2}
2. \underline{La notion de biextension  de faisceaux abéliens}

2.0.  Dans toute la suite du présent numéro, $P,Q,G$ désignent des Groupes commutatifs du topos \underline{$T$} . Nous aurons à utiliser fréquemment le diagramme cartésien

(2.0.1) \begin{tikzcd}
P \arrow[d] \& P\times Q \arrow[l] \arrow[d] \\
e \& Q \arrow[l]
\end{tikzcd} \qquad .













\SGAvar{et les Groupes $G_P$ (resp. $G_Q$ resp. $G_{P\times Q}$) sur $P$ (resp. $Q$, resp. $P\times Q$) images inverses du Groupe $G$ sur $e$. On identifiera $G_{P\times Q}$ indifféremment à l'image inverse $(G_P)_{P\times Q}$ de $G_P$ par $\mathrm{pr}_1$ ou $(G_Q)_{P\times Q}$ de $G_Q$ par $\mathrm{pr}_2$ . Dans ces notations, les structures de Groupes de $P,Q,P\times Q$ n'interviennent pas. Nous n'aurons d'ailleurs partiquement jamais à utiliser la structure de groupe produit sur $P\times Q$. Par contre, il nous sera utile de considérer sur $P\times Q$ la structure de $P$-Groupe image inverse de celle de $Q$ sur $e$,

}{et les Groupes $G_P$ (resp. $G_Q$ resp. $G_{P\times Q}$) sur $P$ (resp. $Q$, resp. $P\times Q$) images inverses du Groupe $G$ sur $e$. On identifiera $G_{P\times Q}$ indifféremment à l'image inverse $(G_P)_{P\times Q}$ de $G_P$ par $\mathrm{pr}_1$ ou $(G_Q)_{P\times Q}$ de $G_Q$ par $\mathrm{pr}_2$ . Dans ces notations, les structures de Groupes de $P,Q,P\times Q$ n'interviennent pas. Nous n'aurons d'ailleurs pratiquement jamais à utiliser la structure de groupe produit sur $P\times Q$. Par contre, il nous sera utile de considérer sur $P\times Q$ la structure de $P$-Groupe image inverse de celle de $Q$ sur $e$,

}{F7I-VII-020}\begin{equation}\tag{2.0.2}
P\times Q = Q_P \qquad ,
\end{equation}
et de même la structure de $Q$-Groupe image inverse de celle de $P$ sur $e$,

\SGAvar{\begin{equation}\tag{2.0.3}
P\times Q = P_Q \qquad .
\end{equation}
Ainsi, sur $P$ on  aura à considérer les deux Groupes $G_P$ et $Q_P$ , et sur $Q$ les deux Groupes $G_Q$ et $P_Q$ . Si $E$ est un torseur sur $P\times Q$ de Groupe $G_{P\times Q}$ , il peut être envisagé indifféremment comme un torseur sur $Q_P$ de groupe $G_{P\times Q} = (G_P)_{Q_P}$ , ou comme un torseur sur $P_Q$ de groupe $G_{P\times Q}=(G_Q)_{P_Q}$ . Cela a donc un sens de considérer sur $E$ une structure d'extension (commutative) de $Q_P$ par $G_P$ compatible avec sa structure de torseur, ou (symétriquement)
\SGApage{I}{157}
une structure d'extension de $P_Q$ par $G_Q$ compatible avec sa structure de torseur. De plus, 1.6, appliqué aux topos induits $T_{/P}$ resp. $T_{/Q}$, nous fournit une façon commode d'expliciter de telles structures. Utilisant les isomorphismes canoniques $T_{/P\times Q} \simeq (T_{/P})_{/P\times Q}$ et $T_{/P\times Q} \simeq (T_{/Q})_{/P\times Q}$ , on trouve les descriptions suivantes. Une structure d'extension de $Q_P$ par $G_P$ sur le torseur $E$ équivaut à la donnée, pour tout objet $S$ de $\underline{T}$ et des éléments $p,p' \in P(S)$, $q \in Q(S)$, d'un isomorphisme de torseurs sous $G_S$

}{\begin{equation}\tag{2.0.3}
P\times Q = P_Q \qquad .
\end{equation}
Ainsi, sur $P$ on  aura à considérer les deux Groupes $G_P$ et $Q_P$ , et sur $Q$ les deux Groupes $G_Q$ et $P_Q$ . Si $E$ est un torseur sur $P\times Q$ de Groupe $G_{P\times Q}$ , il peut être envisagé indifféremment comme un torseur sur $Q_P$ de groupe $G_{P\times Q} = (G_P)_{Q_P}$ , ou comme un torseur sur $P_Q$ de groupe $G_{P\times Q}=(G_Q)_{P_Q}$ . Cela a donc un sens de considérer sur $E$ une structure d'extension (commutative) de $Q_P$ par $G_P$ compatible avec sa structure de torseur, ou (symétriquement)
\SGApage{I}{157}
une structure d'extension de $P_Q$ par $G_Q$ compatible avec sa structure de torseur. De plus, 1.6, appliqué aux topos induits $T_{/P}$ resp. $T_{/Q}$, nous fournit une façon commode d'expliciter de telles structures. Utilisant les isomorphismes canoniques $T_{/P\times Q} \simeq (T_{/P})_{/P\times Q}$ et $T_{/P\times Q} \simeq (T_{/Q})_{/P\times Q}$ , on trouve les descriptions suivantes. Une structure d'extension de $P_Q$ par $G_Q$ sur le torseur $E$ équivaut à la donnée, pour tout objet $S$ de $\underline{T}$ et des éléments $p,p' \in P(S)$, $q \in Q(S)$, d'un isomorphisme de torseurs sous $G_S$

}{F7I-VII-021}\begin{equation}\tag{2.0.4}
\varphi_{p,p';q} \quad : \quad E_{p,q}\,E_{p',q} \;\overset{\sim}{\longrightarrow}\; E_{pp',q} \qquad ,
\end{equation}
satisfaisant aux conditions qu'on va rappeler ; comme d'habitude (cf. 1.1.2.1), pour tout point $(p,q) \in (P\times Q)(S) = P(S)\times Q(S)$, on désigne par $E_{p,q}$ le torseur sous $G_S$ image inverse du torseur $E$ sous $G_{P\times Q}$ par $(p,q) : S \longrightarrow P\times Q$. Les conditions à imposer aux données (2.0.4) sont la \underline{fonctorialité en} $S$ (i.e. la compatibilité aux images inverses par des morphismes $S' \longrightarrow S$ dans $\underline{T}$), l'\underline{associativité} qui s'exprime par la commutativité des diagrammes de la forme (1.1.4.1)

\SGAvar{\begin{equation}\tag{2.0.5}
\begin{tikzcd}[column sep=huge, row sep=large]
 \& E_{pp',q}E_{p'',q} \arrow[rd, "{\varphi_{pp',p'';q}}"] \& \\
E_{p,q}E_{p',q}E_{p'',q} \arrow[ru, "{\varphi_{p,p;q}\wedge \mathrm{id}}"] \arrow[rd, "{\mathrm{id}\wedge\varphi_{p',p'';q}}"'] \& \& E_{pp'p'',q} \\
 \& E_{p,q}E_{p'p'',q} \arrow[ru, "{\varphi_{p,p'p'';q}}"'] \&
\end{tikzcd}
\end{equation}

















}{\begin{equation}\tag{2.0.5}
\begin{tikzcd}[column sep=huge, row sep=large]
 \& E_{pp',q}E_{p'',q} \arrow[rd, "{\varphi_{pp',p'';q}}"] \& \\
E_{p,q}E_{p',q}E_{p'',q} \arrow[ru, "{\varphi_{p,p';q}\wedge \mathrm{id}}"] \arrow[rd, "{\mathrm{id}\wedge\varphi_{p',p'';q}}"'] \& \& E_{pp'p'',q} \\
 \& E_{p,q}E_{p'p'',q} \arrow[ru, "{\varphi_{p,p'p'';q}}"'] \&
\end{tikzcd}
\end{equation}

















}{F7I-VII-022}(où $p,p',p'' \in P(S)$, $q \in Q(S)$), enfin la \underline{commutativité} qui s'exprime par la commutativité des diagrammes de la forme (1.2.1) 


\SGApage{I}{158}

(2.0.6) 
\begin{tikzcd}[column sep=huge, row sep=huge]
E_{p,q}E_{p',q} \arrow[r, "{\varphi_{p,p';q}}"] \arrow[d, "\mathrm{sym}"'] \& E_{pp',q} \arrow[d, "\mathrm{id}"] \\
E_{p',q}E_{p,q} \arrow[r, "{\varphi_{p',p;q}}"] \& E_{p'p,q}
\end{tikzcd}

\SGAvar{(où $p,p' \in P(S)$, $q \in Q(S)$). On explicite de même une structure d'extension de $P_Q$ par $G_Q$ sur $E$ à l'aide d'isomorphismes de torseurs sous $G_S$

\begin{equation}\tag{2.0.7}
\psi_{p;\ q,q'}: \quad E_{p,q}E_{p,q'} \;\overset{\sim}{\longrightarrow}\; E_{pp',q} \qquad ,
\end{equation}
pour $p \in P(S)$ et $q,q' \in Q(S)$, les $\psi$ étant fonctoriels en $S$, et satisfaisant les conditions d'associativité et de commutativité symétriques de (2.0.5) et (2.0.6), s'exprimant par la commutativité des diagrammes

\begin{equation}\tag{2.0.8}
\begin{tikzcd}[column sep=huge, row sep=large]
 \& E_{p,qq'}E_{p,q''} \arrow[rd, "{\psi_{p;qq',q''}}"] \& \\
E_{p,q}E_{p,q'}E_{p,q''} \arrow[ru, "{\psi_{p;q,q'}\wedge \mathrm{id}}"] \arrow[rd, "{\mathrm{id}\wedge\psi_{p;q',q''}}"'] \& \& E_{p,qq'q''} \\
 \& E_{p,q}E_{p,q'q''} \arrow[ru, "\psi_{p;q;q'q''}"'] \&
\end{tikzcd}
\end{equation}

















}{(où $p,p' \in P(S)$, $q \in Q(S)$). On explicite de même une structure d'extension de $Q_P$ par $G_P$ sur $E$ à l'aide d'isomorphismes de torseurs sous $G_S$

\begin{equation}\tag{2.0.7}
\psi_{p;\ q,q'}: \quad E_{p,q}E_{p,q'} \;\overset{\sim}{\longrightarrow}\; E_{p,qq'} \qquad ,
\end{equation}
pour $p \in P(S)$ et $q,q' \in Q(S)$, les $\psi$ étant fonctoriels en $S$, et satisfaisant les conditions d'associativité et de commutativité symétriques de (2.0.5) et (2.0.6), s'exprimant par la commutativité des diagrammes

\begin{equation}\tag{2.0.8}
\begin{tikzcd}[column sep=huge, row sep=large]
 \& E_{p,qq'}E_{p,q''} \arrow[rd, "{\psi_{p;qq',q''}}"] \& \\
E_{p,q}E_{p,q'}E_{p,q''} \arrow[ru, "{\psi_{p;q,q'}\wedge \mathrm{id}}"] \arrow[rd, "{\mathrm{id}\wedge\psi_{p;q',q''}}"'] \& \& E_{p,qq'q''} \\
 \& E_{p,q}E_{p,q'q''} \arrow[ru, "\psi_{p;q,q'q''}"'] \&
\end{tikzcd}
\end{equation}

















}{F7I-VII-023}($p \in P(S)$, $q,q',q'' \in Q(S)$),

\begin{equation}\tag{2.0.9}
\begin{tikzcd}[column sep=huge, row sep=huge]
E_{p,q}E_{p,q'} \arrow[r, "{\psi_{p;q,q'}}"] \arrow[d, "\mathrm{sym}"'] \& E_{p,qq'} \arrow[d, "\mathrm{id}"] \\
E_{p,q'}E_{p,q} \arrow[r, "{\psi_{p;q',q}}"] \& E_{p,q'q}
\end{tikzcd}
\end{equation}











($p \in P(S)$, $q,q' \in Q(S)$).


\SGApage{I}{159}

\SGAvar{\underline{Définition} 2.1. \underline{Soient} $P,Q,G$ \underline{des Groupes abéliens de} \underline{$T$}, $E$ \underline{un torseur sous}
$G_{P\times Q}$, \underline{muni de plus d'une structure d'extension de} $Q_P$ \underline{par} $G_P$ \underline{défini par}
\underline{des isomorphismes} (2.0.4) (\underline{fonctoriels en} $S$, \underline{soumis aux conditions s'exprimant par la commutativité des diagrammes} (2.0.5) \underline{et} (2.0.6)), \underline{et d'une}
\underline{structure d'extension de} $P_Q$ \underline{par} $G_Q$, \underline{défini par des isomorphismes} (2.0.7)
(\underline{fonctoriels en} $S$, \underline{soumis aux conditions exprimées par la commutativité}
\underline{des diagrammes} (2.0.8) \underline{et} (2.0.9)). \underline{On dit que les deux structures précédentes sont} compatibles \underline{si pour tout objet} $S$ \underline{de} \underline{$T$} \underline{et pour}
$p,p' \in P(S)$, $q,q' \in Q(S)$, \underline{on a commutativité dans le diagramme}

}{\underline{Définition} 2.1. \underline{Soient} $P,Q,G$ \underline{des Groupes abéliens de} \underline{$T$}, $E$ \underline{un torseur sous}
$G_{P\times Q}$, \underline{muni de plus d'une structure d'extension de} $P_Q$ \underline{par} $G_Q$ \underline{définie par}
\underline{des isomorphismes} (2.0.4) (\underline{fonctoriels en} $S$, \underline{soumis aux conditions s'exprimant par la commutativité des diagrammes} (2.0.5) \underline{et} (2.0.6)), \underline{et d'une}
\underline{structure d'extension de} $Q_P$ \underline{par} $G_P$, \underline{définie par des isomorphismes} (2.0.7)
(\underline{fonctoriels en} $S$, \underline{soumis aux conditions exprimées par la commutativité}
\underline{des diagrammes} (2.0.8) \underline{et} (2.0.9)). \underline{On dit que les deux structures précédentes sont} compatibles \underline{si pour tout objet} $S$ \underline{de} \underline{$T$} \underline{et pour}
$p,p' \in P(S)$, $q,q' \in Q(S)$, \underline{on a commutativité dans le diagramme}

}{F7I-VII-024}\begin{equation}\tag{2.1.1}
\begin{tikzcd}[column sep=1.2em, row sep=2.2em]
 \& \& E_{p,qq'}E_{p',qq'} \arrow[rrdd, "{\varphi_{p,p';qq'}}"] \& \& \\
E_{p,q}E_{p,q'}E_{p',q}E_{p',q'} \arrow[rru, "{\psi_{p;q,q'}\wedge\psi_{p';q,q'}}"] \arrow[d] \& \& \& \& \\
E_{p,q}E_{p',q}E_{p,q'}E_{p',q'} \arrow[rrd, "{\varphi_{p,p';q}\wedge\varphi_{p,p';q'}}"'] \& \& \& \& E_{pp',qq'} \\
 \& \& E_{pp',q}E_{pp',q'} \arrow[rru, "{\psi_{pp';q,q'}}"'] \& \&
\end{tikzcd}
\qquad ,
\end{equation}






















\SGAvar{\underline{où la flèche verticale est déduite de l'isomorphisme de symétrie}
$E_{p,q'}E_{p',q} \overset{\sim}{\longrightarrow} E_{p',q}E_{p,q'}$ . \underline{On appelle} biextension de $(P,Q)$ par $G$ \underline{un}
\underline{torseur} $E$ \underline{sous} $G_{P\times Q}$ , \underline{muni d'une structure d'extension de} $Q_P$ \underline{par} $G_P$ \underline{et}
\underline{d'une structure d'extension de} $G_Q$ \underline{par} $P_Q$, \underline{ces deux structures étant}
\underline{compatibles au sens précédent.}


}{\underline{où la flèche verticale est déduite de l'isomorphisme de symétrie}
$E_{p,q'}E_{p',q} \overset{\sim}{\longrightarrow} E_{p',q}E_{p,q'}$ . \underline{On appelle} biextension de $(P,Q)$ par $G$ \underline{un}
\underline{torseur} $E$ \underline{sous} $G_{P\times Q}$ , \underline{muni d'une structure d'extension de} $P_Q$ \underline{par} $G_Q$ \underline{et}
\underline{d'une structure d'extension de} $Q_P$ \underline{par} $G_P$, \underline{ces deux structures étant}
\underline{compatibles au sens précédent.}


}{F7I-VII-025}\SGApage{I}{160}

2.1.1.  On peut dire encore qu'une structure de biextension sur le
$G_{P\times Q}$-torseur $E$ est définie par deux lois de composition non partout définies$\varphi(x,y)$ et $\psi(x,y)$ sur $E$, la première étant définie lorsque
$\mathrm{pr}_2 j(x) = \mathrm{pr}_2 j(y)$, la deuxième lorsque $\mathrm{pr}_1 j(x) = \mathrm{pr}_1 j(y)$ (où $j : E \to P\times Q$
est le morphisme structural du torseur). Les axiomes, en plus des associativités du type

\SGAvar{\begin{equation}\tag{2.1.1.1}
\varphi(gx,y) = \varphi(x,gy) = g\varphi(x,y) \qquad (x,y \in E(S), \; g \in G(S))
\end{equation}
et des relations analogues en $\psi$, s'exprimant alors par l'associativité
et la commutativité de ces lois de compositions partielles, et par la
relation

}{\begin{equation}\tag{2.1.1.1}
\varphi(gx,y) = \varphi(x,gy) = g\varphi(x,y) \qquad (x,y \in E(S), \; g \in G(S))
\end{equation}
et des relations analogues en $\psi$, s'expriment alors par l'associativité
et la commutativité de ces lois de compositions partielles, et par la
relation

}{F7I-VII-026}\begin{equation}\tag{2.1.1.2}
\varphi(\psi(x,y),\psi(z,t)) \;=\; \psi(\varphi(x,z),\varphi(y,t))
\end{equation}
pour $x,y,z,t$ respectivement dans $E_{p,q}(S),E_{p,q'}(S),E_{p',q}(S),E_{p',q'}(S)$.

2.2.  \underline{Sections unité}, Il y a lieu de considérer la section unité de $E$,
considéré comme extension de $Q_P$ par $G_P$ ,

\SGAvar{\begin{equation}\tag{2.2.1}
e_{E/P} \in \Gamma(E/P) = \mathrm{Hom}_P(P,E) \qquad ,
\end{equation}
et de façon symétrique la section unité de $E$, considéré comme extension
de $P_G$ par $G_Q$ ,

\begin{equation}\tag{2.2.2}
e_{E/Q} \in \Gamma(E/Q) = \mathrm{Hom}_Q(Q,E) \qquad .
\end{equation}
En particulier, \underline{le torseur} $E$ \underline{au-dessus de} $P\times Q$ (de Groupe $G_{P\times Q}$) \underline{est canoniquement trivialisé au-dessus de} $1_P\times Q$ \underline{par} $e_{E/P}$, \underline{et au-dessus de} $P \times 1_Q$ \underline{par}
$e_{E/Q}$ . De plus, \underline{ces trivialisations coïncident au-dessus de la section}
\underline{unité} $1_P\times 1_Q$ \underline{de} $P\times Q$, et définissent donc une même section 

}{\begin{equation}\tag{2.2.1}
e_{E/P} \in \Gamma(E/P) = \mathrm{Hom}_P(P,E) \qquad ,
\end{equation}
et de façon symétrique la section unité de $E$, considéré comme extension
de $P_Q$ par $G_Q$ ,

\begin{equation}\tag{2.2.2}
e_{E/Q} \in \Gamma(E/Q) = \mathrm{Hom}_Q(Q,E) \qquad .
\end{equation}
En particulier, \underline{le torseur} $E$ \underline{au-dessus de} $P\times Q$ (de Groupe $G_{P\times Q}$) \underline{est canoniquement trivialisé au-dessus de} $P\times 1_Q$ \underline{par} $e_{E/P}$, \underline{et au-dessus de} $1_P \times Q$ \underline{par}
$e_{E/Q}$ . De plus, \underline{ces trivialisations coïncident au-dessus de la section}
\underline{unité} $1_P\times 1_Q$ \underline{de} $P\times Q$, et définissent donc une même section 

}{F7I-VII-027}\SGApage{I}{161}

\SGAvar{\begin{equation}\tag{2.2.3}
 e_{E} \in \Gamma(E) \qquad ,
\end{equation}
qu'on appelera par abus de langage la \underline{section unité} (ou \underline{section neutre}) de $E$. Pour prouver l'égalité annoncée des deux sections $e,e'$ de $E$ déduites respectivement de $e_{E/P}$ et $e_{E/Q}$ , on note qu'elles sont caractérisées respectivement comme les sections unités dans l'image inverse $E_{1_{P},1_{Q}}$ dans $E$ de la section unité de $P\times Q$, pour les deux lois de compositions $\varphi$ et $\psi$. Appliquant alors (2.1.1.2) pour les sections $e,e',e',e$ de $E_{1,1}$ , on trouve que des deux membres sont égaux respectivement à $e$ et $e'$, d'où $e=e'$ comme annoncé. On en conclut de plus que \underline{les deux lois de composition} $\varphi$ \underline{et} $\psi$ \underline{sur} $E_{1,1}$ \underline{sont identiques}. En effet, il est immédiat qu'elles correspondent à la loi de Groupe de $G$ par transport de structure, à l'aide des morphismes $g \longmapsto g.e$ et $g \longmapsto g.e'$ respectivement, donc elles sont identiques puisque $e=e'$.

Notons également que l'extension $E$ de $Q_{P}$ par $G_{P}$ sur la base $P$ donne, par restriction au-dessus de la section $1_{P}$ de $P$, une structure d'extension de $Q$ par $G$ sur l'image inverse de $1_{P}\times Q$ dans $G$. On constate aussitôt que l'homomorphisme $e_{E/Q}$ est un homomorphisme multiplicatif pour cette structure d'extension, qui est donc canoniquement scindée. De façon symétrique, l'extension de $P$ par $G$ obtenue par restriction de l'extension $E$ de $P_{Q}$ par $G_{Q}$ sur $Q$ au-dessus de la section $1_{Q}$ de $Q$, est canoniquement scindée par $e_{E/P}$ .


}{\begin{equation}\tag{2.2.3}
 e_{E} \in \Gamma(E) \qquad ,
\end{equation}
qu'on appellera par abus de langage la \underline{section unité} (ou \underline{section neutre}) de $E$. Pour prouver l'égalité annoncée des deux sections $e,e'$ de $E$ déduites respectivement de $e_{E/P}$ et $e_{E/Q}$ , on note qu'elles sont caractérisées respectivement comme les sections unités dans l'image inverse $E_{1_{P},1_{Q}}$ dans $E$ de la section unité de $P\times Q$, pour les deux lois de composition $\varphi$ et $\psi$. Appliquant alors (2.1.1.2) pour les sections $e,e',e',e$ de $E_{1,1}$ , on trouve que les deux membres sont égaux respectivement à $e$ et $e'$, d'où $e=e'$ comme annoncé. On en conclut de plus que \underline{les deux lois de composition} $\varphi$ \underline{et} $\psi$ \underline{sur} $E_{1,1}$ \underline{sont identiques}. En effet, il est immédiat qu'elles correspondent à la loi de Groupe de $G$ par transport de structure, à l'aide des morphismes $g \longmapsto g.e$ et $g \longmapsto g.e'$ respectivement, donc elles sont identiques puisque $e=e'$.

Notons également que l'extension $E$ de $Q_{P}$ par $G_{P}$ sur la base $P$ donne, par restriction au-dessus de la section $1_{P}$ de $P$, une structure d'extension de $Q$ par $G$ sur l'image inverse de $1_{P}\times Q$ dans $E$. On constate aussitôt que l'homomorphisme $e_{E/Q}$ est un homomorphisme multiplicatif pour cette structure d'extension, qui est donc canoniquement scindée. De façon symétrique, l'extension de $P$ par $G$ obtenue par restriction de l'extension $E$ de $P_{Q}$ par $G_{Q}$ sur $Q$ au-dessus de la section $1_{Q}$ de $Q$, est canoniquement scindée par $e_{E/P}$ .


}{F7I-VII-028}\SGApage{I}{162}

2.3. \underline{Morphismes de biextensions}. Soient $E$ une biextension de $(P,Q)$ par $G$, et de même $E'$ une biextension de $(P',Q')$ par $G'$, où $P',Q',G'$ sont encore des Groupes commutatifs de \underline{$T$}. On appelle \underline{morphisme de biextensions} de $E$ dans $E'$ un système $(u,v,w,f)$, où

\begin{equation}\tag{2.3.1}
u:P \longrightarrow P' \ , \ v:Q \longrightarrow Q' \ , \ w:G \longrightarrow G'
\end{equation}
sont des morphismes de Groupes, et où

\begin{equation}\tag{2.3.2}
f \, : \, E \longrightarrow E'
\end{equation}
est un morphisme des faisceaux d'ensembles sous-jacents à $E,E'$, de telle façon que $f$ soit à la fois un homomorphisme d'extensions de Groupes sur $P$ resp. $P'$ associé aux homomorphismes de Groupes relatifs

\begin{equation*}
u \times v \ : \ Q_{P} \longrightarrow Q'_{P'} \qquad , \qquad u \times w : G_{P} \longrightarrow G'_{P'} \qquad ,
\end{equation*}
et un homomorphisme d'extensions de Groupes sur $Q$ resp. $Q'$ associé aux homomorphismes de Groupes relatifs

\begin{equation*}
u \times v \ : \ P_{Q} \longrightarrow P'_{Q'} \qquad , \qquad v \times w : G_{Q} \longrightarrow G'_{Q'} \qquad .
\end{equation*}
On notera d'ailleurs que $u,v,w$ sont connus par la seule connaissance de $f$, car à partir de $f$ on récupère $u \times v$ par passage au quotient, donc aussi $u$ et $v$ par restriction à $P \times 1_{Q}$ resp. $1_{P} \times Q$, et on récupère $w$ par restriction de $f$ à $G.e_{E}$ . Aussi il n'y a pas d'inconvénient à identifier le morphisme $(u,v,w,f)$ à $f$.

Moyennant la précaution habituelle de rajouter aux morphismes leur source et leur but (génuflexion canonique), on voit donc que \underline{les biextensions du topos} \underline{$T$} \underline{sont les objets d'une catégorie} dont les objets et les morphismes sont ceux définis en 2.1 et 2.3, et où la composition
\SGApage{I}{163}
des morphismes est définie de façon évidente. De plus, par $(P,Q,G,E) \longmapsto (P,Q,G)$ on obtient un foncteur canonique

\begin{equation}\tag{2.3.3}
\mathrm{BIEXT}(\underline{T}) \longrightarrow \mathrm{GRAB}(\underline{T})\times\mathrm{GRAB}(\underline{T})\times\mathrm{GRAB}(\underline{T}) \quad ,
\end{equation}
où $\mathrm{BIEXT}(\underline{T})$ désigne la catégorie des biextensions sur $\underline{T}$, et $\mathrm{GRAB}(\underline{T})$ celle des Groupes abéliens sur $T$. Il y a lieu d'ailleurs de décomposer le foncteur canonique précédent en deux, savoir $E \longmapsto (P,Q)$ :

\begin{equation}\tag{2.3.4}
\mathrm{BIEXT}(\underline{T}) \longrightarrow \mathrm{GRAB}(\underline{T}) \times \underline{\mathrm{GRAB}}(T) \quad ,
\end{equation}
et $E \longmapsto G$ :

\begin{equation}\tag{2.3.5}
\mathrm{BIEXT}(\underline{T}) \longrightarrow \mathrm{GRAB}(\underline{T}) \quad .
\end{equation}

2.4. \underline{La catégorie} $\mathrm{BIEXT}(P,Q,G)$  \underline{pour} $P,Q,G$ \underline{variables}.

\SGAvar{Je dis que \underline{le foncteur canonique} (2.3.4) \underline{en fibrant}, et le \underline{foncteur canonique} (2.3.5) \underline{est cofibrant}. La première assertion signifie donc notamment que pour une biextension $E$ de $(P,Q)$ par $G$, et pour des morphismes de Groupes abéliens

}{Je dis que \underline{le foncteur canonique} (2.3.4) \underline{est fibrant}, et le \underline{foncteur canonique} (2.3.5) \underline{est cofibrant}. La première assertion signifie donc notamment que pour une biextension $E$ de $(P,Q)$ par $G$, et pour des morphismes de Groupes abéliens

}{F7I-VII-029}\begin{equation*}
u : P' \longrightarrow P \ , \ v : Q' \longrightarrow Q \ ,
\end{equation*}
on peut trouver une biextension $E'$ de $(P',Q')$ par $G$, et un morphisme de biextensions

\SGAvar{\begin{equation*}
f : E' \longrightarrow E
\end{equation*}
compatible avec $u,v,\mathrm{id}_G$ , qui jouisse de la propriété universelle habituelle pour les morphismes de biextensions dans $E$ compatibles avec $u,v,\mathrm{id}_G$ . La seconde assertion s'explicite de même, et signifie notamment que pour une biextension $E$ de $(P,G)$ par $G$, et un morphisme de Groupes abéliens

}{\begin{equation*}
f : E' \longrightarrow E
\end{equation*}
compatible avec $u,v,\mathrm{id}_G$ , qui jouisse de la propriété universelle habituelle pour les morphismes de biextensions dans $E$ compatibles avec $u,v,\mathrm{id}_G$ . La seconde assertion s'explicite de même, et signifie notamment que pour une biextension $E$ de $(P,Q)$ par $G$, et un morphisme de Groupes abéliens

}{F7I-VII-030}\begin{equation*}
w : G \longrightarrow G' \quad , 
\end{equation*}

\SGApage{I}{164}

il existe une biextension $E'$ de  $(P,Q)$  par $G'$, et un morphisme de biextensions

\begin{equation*}
g : E \longrightarrow E'
\end{equation*}
compatible avec $\mathrm{id}_P,\mathrm{id}_Q$ et $w$, qui jouisse de la propriété universelle habituelle pour les morphismes de biextensions de $E$ compatibles avec $\mathrm{id}_P,\mathrm{id}_Q$ et $w$. Nous laissons au lecteur le soin d'expliciter les deux constructions dont nous affirmons ici l'existence, constructions qui se font aisément en paraphrasant simplement les constructions analogues bien connues dans le cas des extensions de Groupes ordinaires.

Ainsi, désignant, pour trois Groupes abéliens $P,Q,G$ de $\underline{T}$, par

\SGAvar{\begin{equation}\tag{2.4.1}
\mathrm{BIEXT}(P,Q;G)
\end{equation}
la catégorie des biextensions de $(P;Q)$ par $G$, définie comme la catégorie-fibre de (2.3.3) en $(P,Q,G)$ (donc en y prenant comme morphismes les morphismes de biextensions qui sont compatibles avec les morphismes identiques de $P,Q,G$), on voit que \underline{cette catégorie} "\underline{dépend de façon covariante de} $G$, \underline{et de façon contravariante de} $P,Q$ ". De plus, pour les variances précédentes, \underline{la catégorie} $\mathrm{BIEXT}(P,Q;G)$ \underline{dépend} additivement \underline{de chacun des trois arguments} $P,Q,G$. Explicitons la signification précise de cette assertion, pour l'argument $G$ par exemple :

}{\begin{equation}\tag{2.4.1}
\mathrm{BIEXT}(P,Q;G)
\end{equation}
la catégorie des biextensions de $(P,Q)$ par $G$, définie comme la catégorie-fibre de (2.3.3) en $(P,Q,G)$ (donc en y prenant comme morphismes les morphismes de biextensions qui sont compatibles avec les morphismes identiques de $P,Q,G$), on voit que \underline{cette catégorie} "\underline{dépend de façon covariante de} $G$, \underline{et de façon contravariante de} $P,Q$ ". De plus, pour les variances précédentes, \underline{la catégorie} $\mathrm{BIEXT}(P,Q;G)$ \underline{dépend} additivement \underline{de chacun des trois arguments} $P,Q,G$. Explicitons la signification précise de cette assertion, pour l'argument $G$ par exemple :

}{F7I-VII-031}a)   $\mathrm{BIEXT}(P,Q;0)$ est équivalent à la catégorie ponctuelle, et

b)   pour deux Groupes abéliens $G,G'$ de $\underline{T}$, le foncteur

\SGAvar{\begin{equation}\tag{2.4.2}
\mathrm{BIEXT}(P,G;G\times G') \longrightarrow \mathrm{BIEXT}(P,Q;G) \times \mathrm{BIEXT}(P,Q;G')
\end{equation}
déduit des deux projections $G \times G' \longrightarrow G$ et $G\times G' \longrightarrow G'$ est une \underline{équivalence de catégories}. Dans le langage de [4], on dit aussi que pour $P,Q$ fixés, la catégorie des biextensions de $(P,Q)$ par des Groupes $G$ variables est,
\SGApage{I}{165}
via le foncteur $E \longmapsto G$ induit par (2.3.5), une \underline{catégorie cofibrée additive} sur la catégorie GRAB(\underline{$T$}).

\underline{Remarque} 2.4.3. Avec la terminologie de loc.cit., la catégorie cofibrée précédente est d'ailleurs non seulement additive, mais même "\underline{exacte à gauche}" [4, 2.2]. Nous ne vérifierons pas ici cette propriété, que nous retrouverons plus bas (3.6) comme conséquence de faits plus précis. Signalons seulement que la théorie générale de loc. cit. nous permettrait alors d'expliciter la structure de cette catégorie cofibrée en termes d'un complexe de chaînes $L$. dans GRAB(\underline{$T$})(ne dépendant que de $P$ et de $Q$), à équivalence de catégories fibrées près. Nous trouverons également plus bas une construction directe de ce complexe $L$. comme étant le complexe $P \overset{\mathbb{L}}{\otimes} Q$ (produit tensoriel pris dans une catégorie dérivée)..


}{\begin{equation}\tag{2.4.2}
\mathrm{BIEXT}(P,Q;G\times G') \longrightarrow \mathrm{BIEXT}(P,Q;G) \times \mathrm{BIEXT}(P,Q;G')
\end{equation}
déduit des deux projections $G \times G' \longrightarrow G$ et $G\times G' \longrightarrow G'$ est une \underline{équivalence de catégories}. Dans le langage de [4], on dit aussi que pour $P,Q$ fixés, la catégorie des biextensions de $(P,Q)$ par des Groupes $G$ variables est,
\SGApage{I}{165}
via le foncteur $E \longmapsto G$ induit par (2.3.5), une \underline{catégorie cofibrée additive} sur la catégorie GRAB(\underline{$T$}).

\underline{Remarque} 2.4.3. Avec la terminologie de loc.cit., la catégorie cofibrée précédente est d'ailleurs non seulement additive, mais même "\underline{exacte à gauche}" [4, 2.2]. Nous ne vérifierons pas ici cette propriété, que nous retrouverons plus bas (3.6) comme conséquence de faits plus précis. Signalons seulement que la théorie générale de loc. cit. nous permettrait alors d'expliciter la structure de cette catégorie cofibrée en termes d'un complexe de chaînes $L$. dans GRAB(\underline{$T$})(ne dépendant que de $P$ et de $Q$), à équivalence de catégories fibrées près. Nous trouverons également plus bas une construction directe de ce complexe $L$. comme étant le complexe $P \overset{\mathbb{L}}{\otimes} Q$ (produit tensoriel pris dans une catégorie dérivée).


}{F7I-VII-032}2.5. \underline{Structure de la catégorie} BIEXT$(P,Q;G)$

Les considérations générales formelles de [4, n° 1] permettent de tirer quelques conséquences utiles du fait que les  catégories BIEXT$(P,Q;G)$  dépendent additivement de $G$.

\SGAvar{2.5.1. Tout d'abord, pour $G$ fixé, la catégorie BIEXT$(P,Q;G)$ est munie d'une structure de composition, permettant à deux biextensions $E,E'$ de $(P,Q )$ par $G$ de faire correspondre fonctoriellement une \underline{biextension produit}, notée $E.E'$. On peut considérer que cette construction ne fait que paraphraser la construction bien connue du produit de deux extensions d'un Groupe commutatif par un autre ; on notera de plus qu'elle est bien compatible avec cette dernière notion, quand on interprète $E$ et $E'$, soient comme des extensions de Groupes de $Q_P$ par $G_P$, soit comme des extensions de Groupes de $P_Q$ par $G_Q$. En particulier, $EE'$ est, en tant que torseur
\SGApage{I}{166}
sous $G_{P \times Q}$, le produit des torseurs sous $G_{P \times Q}$ définis par $E$ et $E'$. La formation de $E.E'$ est de plus \underline{associative et commutative} à isomorphismes canoniques près, et laloi de composition envisagée dans BIEXT$(P,Q;G)$ admet un "objet neutre bilatère", appelée la \underline{biextension triviale de} $(P,Q)$ \underline{par} $G$. Celle-ci peut être construite (loc. cit.) comme la biextension déduite de l'unique biextension (à isomorphisme près) $E_o$ de $(P,Q)$ par le Groupe nul $O$, grâce à l'homomorphisme  $O \longrightarrow G$. On vérifie aussitôt sur cette description que la biextension triviale $E$ peut être aussi décrite comme le torseur trivial

}{2.5.1. Tout d'abord, pour $G$ fixé, la catégorie BIEXT$(P,Q;G)$ est munie d'une structure de composition, permettant à deux biextensions $E,E'$ de $(P,Q )$ par $G$ de faire correspondre fonctoriellement une \underline{biextension produit}, notée $E.E'$. On peut considérer que cette construction ne fait que paraphraser la construction bien connue du produit de deux extensions d'un Groupe commutatif par un autre ; on notera de plus qu'elle est bien compatible avec cette dernière notion, quand on interprète $E$ et $E'$, soit comme des extensions de Groupes de $Q_P$ par $G_P$, soit comme des extensions de Groupes de $P_Q$ par $G_Q$. En particulier, $EE'$ est, en tant que torseur
\SGApage{I}{166}
sous $G_{P \times Q}$, le produit des torseurs sous $G_{P \times Q}$ définis par $E$ et $E'$. La formation de $E.E'$ est de plus \underline{associative et commutative} à isomorphismes canoniques près, et la loi de composition envisagée dans BIEXT$(P,Q;G)$ admet un "objet neutre bilatère", appelé la \underline{biextension triviale de} $(P,Q)$ \underline{par} $G$. Celle-ci peut être construite (loc. cit.) comme la biextension déduite de l'unique biextension (à isomorphisme près) $E_o$ de $(P,Q)$ par le Groupe nul $O$, grâce à l'homomorphisme  $O \longrightarrow G$. On vérifie aussitôt sur cette description que la biextension triviale $E$ peut être aussi décrite comme le torseur trivial

}{F7I-VII-033}\begin{equation*}
E \;=\; P \times Q \times G \qquad ,
\end{equation*}
ses deux structures de Groupes sur $P$ resp  sur $Q$ étant déduites des structures de Groupes produits de $Q \times G$ resp. de $P \times G$ par changement de base $P \longrightarrow e$ resp. $Q \longrightarrow e$.

Enfin, pour toute biextension $E$ de $(P,Q)$ par $G$, on peut définir \underline{la biextension inverse}, notée $E^{-1}$ , munie d'un isomorphisme de $EE^{-1}$ avec la biextension triviale, ce qui détermine $E^{-1}$ à isomorphisme unique près, et permet de définir la dépendance fonctorielle par rapport à $E$.

En somme, la catégorie BIEXT$(P,Q;G)$ apparaît comme l'analogue d'un objet-groupe de la catégorie (Cat) des catégories, à cela près que les associativités, commutativités, relations d'unité et d'inverses ne sont valables que modulo des isomorphismes donnés près. On pourra dire que cette catégorie est un \underline{pseudo-groupe} (de la 2-catégorie (Cat)).


\SGApage{I}{167}

2.5.2. Pour $P, Q, G$ fixés, il résulte des propriétés précédentes les faits suivants. Pour toute biextension $E$ de $(P,Q)$ par $G$, le monoïde des endomorphismes de $E$ est canoniquement isomorphe au monoïde des endomorphismes de la biextension triviale \underline{$1$}, en associant à tout endomorphisme $u$ de celle-ci l'endomorphisme $u$.id$_{E}$ de \underline{$1$}.$E$ $\simeq$ $E$ . D'autre part, le monoïde des endomorphismes de \underline{$1$} est un \underline{groupe commutatif}, que nous noterons

\SGAvar{\begin{equation}\tag{2.5.2.1}
\mathrm{Biext}^{0}(P,Q;G) \quad ;
\end{equation}
sa loi de composition est aussi donnéepar $(u,v)$ $\longmapsto$ $u.v$ $\in$ End(\underline{$1$}.\underline{$1$}) $\simeq$ End(\underline{$1$}). Enfin, la loi de composition $E,E'$ $\longmapsto$ $EE'$ définit dans l'ensemble

\begin{equation}\tag{2.5.2.2}
\mathrm{Biext}^{1}(P,Q;G)
\end{equation}
des classes d'isomorphismesde biextensions de $(P,Q)$ par $G$ une loi de \underline{groupe commutatif}, l'inverse de la classe de $E$ étant celle du bitorseur inverse $E^{-1}$.

}{\begin{equation}\tag{2.5.2.1}
\mathrm{Biext}^{0}(P,Q;G) \quad ;
\end{equation}
sa loi de composition est aussi donnée par $(u,v)$ $\longmapsto$ $u.v$ $\in$ End(\underline{$1$}.\underline{$1$}) $\simeq$ End(\underline{$1$}). Enfin, la loi de composition $E,E'$ $\longmapsto$ $EE'$ définit dans l'ensemble

\begin{equation}\tag{2.5.2.2}
\mathrm{Biext}^{1}(P,Q;G)
\end{equation}
des classes d'isomorphismes de biextensions de $(P,Q)$ par $G$ une loi de \underline{groupe commutatif}, l'inverse de la classe de $E$ étant celle du bitorseur inverse $E^{-1}$.

}{F7I-VII-034}2.5.3. Lorsqu'on se donne un homomorphisme de Groupes $G$ $\longrightarrow$ $G'$, le foncteur correspondant

\begin{equation}\tag{2.5.3.1}
\mathrm{BIEXT}(P,Q;G) \longrightarrow \mathrm{BIEXT}(P,Q;G')
\end{equation}
est compatible avec les lois $E.E'$ et $E^{-1}$ envisagées dans 2.5.1, à des isomorphismes canoniques près. Il s'ensuit en particulier que les deux applications évidentes

\begin{equation}\tag{2.5.3.2}
\mathrm{Biext}^{i}(P,Q;G) \longrightarrow \mathrm{Biext}^{i}(P,Q;G') \qquad (i=0,1)
\end{equation}
sont des \underline{homomorphismes de groupes}. De plus, si $E$ est une biextension de $(P,Q)$ par $G$, et $E'$ la biextension correspondante de $(P,Q)$ par $G'$, l'homomorphisme naturel

\begin{equation}\tag{2.5.3.3}
\mathrm{Hom}(E,E) = \mathrm{Aut}(E) \longrightarrow \mathrm{Hom}(E',E') = \mathrm{Aut}(E')
\end{equation}
n'est autre que (2.5.3.2) pour $i = 0$.

\SGApage{I}{168}

\SGAvar{2.5.4. On peut d'ailleurs expliciter aisément le groupe Biext$^{0}(P,Q;G)$. En effet, si $E$ est une biextension de $(P,Q)$ par $G$, les automorphismes $u$ de $E$ sont d'abord des automorphismes du torseur sous $G_{P\times Q}$ sous-jacent à $E$, et à ce titre est donnés par les morphismes

}{2.5.4. On peut d'ailleurs expliciter aisément le groupe Biext$^{0}(P,Q;G)$. En effet, si $E$ est une biextension de $(P,Q)$ par $G$, les automorphismes $u$ de $E$ sont d'abord des automorphismes du torseur sous $G_{P\times Q}$ sous-jacent à $E$, et à ce titre sont donnés par les morphismes

}{F7I-VII-035}\begin{equation*}
u_{0} : P\times Q \longrightarrow G \qquad ,
\end{equation*}
par la formule

\begin{equation*}
u(x) = u_{0}(j(x))x \qquad .
\end{equation*}
Ceci posé, pour que $u$ soit compatible avec la première (resp. la seconde) structure d'extension sur $E$, on voit tout de suite qu'il faut et il suffit que $u_{0}$ soit additif par rapport au second (resp. par rapport au premier) argument. Donc \underline{pour que} $u$ \underline{soit un automorphisme de la structure de biextension de} $E$, \underline{il faut et il suffit que} $u_{0}$ \underline{soit bilinéaire}. Il en résulte donc une \underline{bijection canonique}

\begin{equation}\tag{2.5.4.1}
\mathrm{Biext}^{0}(P,Q;G) \simeq \mathrm{Hom}(P\otimes Q,G) \qquad .
\end{equation}
On vérifie facilement que celle-ci est indépendante de $E$, qu'elle respecte les structures de groupes, et qu'elle est fonctorielle en $G$ (pour la loi fonctorielle du premier membre explicité dans (2.5.3.2)).

Nous déterminerons plus bas le groupe Biext$^{1}(P,Q;G)$, et trouverons qu'il est canoniquement, et fonctoriellement en $G$,isomorphe à Ext$^{1}(P \overset{\mathbb{L}}{\otimes} Q, G)$.

2.6. Nous avons dans 2.5 fait varier uniquement $G$, en gardant $P$ et $Q$ fixes. Il convient maintenant de dire quelques mots concernant la variation de $P,Q$. Le fait essentiel, dont les autres découleront formellement,est que la variance contravariante de BIEXT$(P,Q;G)$ en $P,Q$ \underline{commute} à sa variance covariante en $G$, de sorte que les catégories envisagées peuvent être considérées
\SGApage{I}{169}
comme les catégories fibres d'une \underline{catégorie cofibrée}

\SGAvar{\begin{equation}\tag{2.6.1}
\mathrm{BIEXT}'(\underline{T}) \longrightarrow \mathrm{GRAB}(\underline{T})^{o}\times\mathrm{GRAB}(\underline{T})^{o}\times \mathrm{GRAB}(\underline{T}) \quad ,
\end{equation}
où les exposants $^{o}$ indiquent les catégories opposées, (le foncteur structural envisagé étant $E \longmapsto (P,Q,G)$). A ce titre, la catégorie cofibrée envisagée est additive en chacun de ces trois arguments, ou encore, comme nous dirons, elle est \underline{tri-additive}. En paraphrasant des arguments standards que nous nous dispensons de reproduire ici, on en conclut que la structure de composition $E.E'$ sur la catégorie fibre BIEXT$(P,P';G)$ peut se décrire indifféremment en utilisant la structure additive de l'un quelconque des trois arguments (i.e. les trois structures de pseudogroupe obtenues sont canoniquement isomorphes). Il en résulte à son tour que les groupes abéliens

\begin{equation}\tag{2.6.2}
\mathrm{Biext}^{i}(P,Q;G) \quad , \quad i=0,1 \quad ,
\end{equation}
peuvent être regardées comme deux foncteurs tri-linéaires en les arguments $P,Q,G$ dans GRAB($\underline{T}$)$^{o}$, GRAB($\underline{T}$)$^{o}$ et GRAB$(T)$, ou si on préfère comme des \underline{foncteurs trilinéaires} en $P,Q,G$ variables dans GRAB($\underline{T}$), contravariants en $P,Q$, covariants en $G$, à valeurs dans la catégorie GRAB des groupes abéliens. On vérifie de plus que pour cette variance ainsi précisée, l'isomorphisme (2.5.4.1) est fonctoriel en $P,Q$ et $G$, et non seulement en $G$. Il en sera d'ailleurs de même plus bas pour la détermination de Biext$^{1}$ en termes d'un Ext$^{1}$ .


}{\begin{equation}\tag{2.6.1}
\mathrm{BIEXT}'(\underline{T}) \longrightarrow \mathrm{GRAB}(\underline{T})^{o}\times\mathrm{GRAB}(\underline{T})^{o}\times \mathrm{GRAB}(\underline{T}) \quad ,
\end{equation}
où les exposants $^{o}$ indiquent les catégories opposées, (le foncteur structural envisagé étant $E \longmapsto (P,Q,G)$). A ce titre, la catégorie cofibrée envisagée est additive en chacun de ces trois arguments, ou encore, comme nous dirons, elle est \underline{tri-additive}. En paraphrasant des arguments standards que nous nous dispensons de reproduire ici, on en conclut que la structure de composition $E.E'$ sur la catégorie fibre BIEXT$(P,Q;G)$ peut se décrire indifféremment en utilisant la structure additive de l'un quelconque des trois arguments (i.e. les trois structures de pseudogroupe obtenues sont canoniquement isomorphes). Il en résulte à son tour que les groupes abéliens

\begin{equation}\tag{2.6.2}
\mathrm{Biext}^{i}(P,Q;G) \quad , \quad i=0,1 \quad ,
\end{equation}
peuvent être regardés comme deux foncteurs tri-linéaires en les arguments $P,Q,G$ dans GRAB($\underline{T}$)$^{o}$, GRAB($\underline{T}$)$^{o}$ et GRAB$(T)$, ou si on préfère comme des \underline{foncteurs trilinéaires} en $P,Q,G$ variables dans GRAB($\underline{T}$), contravariants en $P,Q$, covariants en $G$, à valeurs dans la catégorie GRAB des groupes abéliens. On vérifie de plus que pour cette variance ainsi précisée, l'isomorphisme (2.5.4.1) est fonctoriel en $P,Q$ et $G$, et non seulement en $G$. Il en sera d'ailleurs de même plus bas pour la détermination de Biext$^{1}$ en termes d'un Ext$^{1}$ .


}{F7I-VII-036}2.7.\underline{Biextension symétrique d'une biextension donnée.}

\SGAvar{Si $E$ est une biextension de $(P,Q)$ par $G$, il y a lieu de définir la \underline{biextension symétrique} $^{s}E$ de $E$, qui est une biextension de $(Q,P)$ par $G$. Le faisceau d'ensembles sous-jacent à $^{s}E$ est celui sous-jacent à $E$, les opérations de $G$ sur $E$ et sur $^{s}E$ sont les mêmes, le morphisme structural $^{s}E$ $\longrightarrow$ P$\times$Q est le composé du morphisme structural $E$ $\longrightarrow$ P$\times$Q avec la symétrie P$\times$Q $\longrightarrow$ Q$\times$P, enfin les structures
\SGApage{I}{170}
d'extensions gauches et droites de $^{s}E$ sont respectivement les structures d'extension droite et gauche de $E$. Il est évident que l'opération $E$ $\longmapsto$ $^{s}E$ est un automorphisme de la catégorie BIEXT($\underline{T}$) des biextensions, et en fait un \underline{automorphisme involutif} :

}{Si $E$ est une biextension de $(P,Q)$ par $G$, il y a lieu de définir la \underline{biextension symétrique} $^{s}E$ de $E$, qui est une biextension de $(Q,P)$ par $G$. Le faisceau d'ensembles sous-jacent à $^{s}E$ est celui sous-jacent à $E$, les opérations de $G$ sur $E$ et sur $^{s}E$ sont les mêmes, le morphisme structural $^{s}E$ $\longrightarrow$ Q$\times$P est le composé du morphisme structural $E$ $\longrightarrow$ P$\times$Q avec la symétrie P$\times$Q $\longrightarrow$ Q$\times$P, enfin les structures
\SGApage{I}{170}
d'extensions gauches et droites de $^{s}E$ sont respectivement les structures d'extension droite et gauche de $E$. Il est évident que l'opération $E$ $\longmapsto$ $^{s}E$ est un automorphisme de la catégorie BIEXT($\underline{T}$) des biextensions, et en fait un \underline{automorphisme involutif} :

}{F7I-VII-037}\begin{equation*}
{}^{s}({}^{s}E) = E \quad .
\end{equation*}
Tout énoncé concernant des bitorseurs $E$, ... permet de déduire un énoncé symétrique, en appliquant l'énoncé donné à $^{s}E$, .... On pourra l'appeler l'énoncé déduit "en échangeant la gauche et la droite". Bien entendu, on se contentera en général d'expliciter un seul des deux énoncés symétriques, en choisissant entre sa gauche et sa droite, en suivant par exemple ses préférences politiques.


2.8.  \underline{Changement de topos.} \underline{Localisation.}

Supposons donné un morphisme de topos

\begin{equation*}
f \;:\; \underline{T}' \longrightarrow \underline{T} \quad .
\end{equation*}
Alors le foncteur "image inverse" $f^{*}$ transforme une biextension de $(P,Q)$ par $G$ en une biextension $f^{*}(E)$ de $(f^{*}(P),f^{*}(Q))$ par $f^{*}(G)$, comme chaque fois qu'on a une espèce de structure qui peut s'expliciter "en termes de limites projectives finies et de limites inductives quelconques". Pour montrer la compatibilité des deux structures d'extensions sur $f^{*}(E)$ déduites de celles de $E$, on exprimera cette compatibilité par la commutativité du diagramme (2.1.1) dans le cas "universel", i.e. pour S=P$\times$P$\times$Q$\times$Q, $p,p',q,q'$ étant les quatre projections de $S$ sur ces quatre facteurs. Comme de juste, on obtiendra par $E \longmapsto f^{*}(E)$ un foncteur canonique

\begin{equation*}
f^{*} \;:\; \mathrm{BIEXT}(\underline{T}) \longrightarrow \mathrm{BIEXT}(\underline{T}') \quad .
\end{equation*}

\SGApage{I}{171}



Nous utiliserons par la suite sans autre mention que toutes les constructions que nous pourrons être amenés à effectuer au moyen de biextensions sont compatibles aux images inverses par morphismes de topos. Il en est en particulier ainsi des opérations de produit $EE'$ et de passage à l'inverse $E^{-1}$, considérées dans 2.5.

On pourra prendre en particulier pour $\underline{T}'$ un \underline{topos induit} $T_{/S'}$ , et pour $f$ le morphisme canonique (dit "de localisation" (SGA 4 IV 5))

\SGAvar{\begin{equation*}
f = i_{S'} \; : \; \underline{T}_{/S'} \longrightarrow \underline{T} \qquad .
\end{equation*}
On appelle alors $f^{*}(E)$ la \underline{biextension induite} par $E$ sur $S'$, et on la notera aussi, comme d'habitude, $E|S'$ ou $E_{S'}$ . Comme l'opération de restriction possède évidemment les propriétés habituelles de transitivité, on en conclut que les catégories BIEXT$(\underline{T}_{/S})$, pour un objet variable $S$ de $\underline{T}$, forment les fibres d'une \underline{catégorie fibrée} sur $\underline{T}$ (SGA 1 VI 8). En fait, on vérifie trivialement que cette catégorie est même un \underline{champ} sur $\underline{T}$ au sens de Giraud [2], i.e. que "les morphismes et les objets des fibres de cette catégorie fibrée peuvent se recoller". Par exemple, pour se donner une biextension à isomorphisme canonique près sur $S$, il suffit de se donner une famille couvrante (i.e. épimorphique) $S_i \longrightarrow S$, pour tout $i \in I$ une biextension $E_i$ sur $S_i$ (i.e. dans $\underline{T}_{/S_i}$), pour tout couple $(i,j)$ un isomorphisme de $E_{iS_j}$ avec $E_{jS_i}$, ces isomorphismes étant soumis à la condition de cocycles habituelle. On appelera \underline{champ des biextensions sur} $\underline{T}$ le champ qu'on vient de définir. Le foncteur (2.3.3), donc aussi les foncteurs déduits (2.3.4) et (2.3.5), lorsqu'on y remplace $\underline{T}$ par $\underline{T}_{/S}$ avec $S$ variable, correspondent à des morphismes de champs sur $\underline{T}$. De même les opérations $E.E'$ et $E^{-1}$ (2.5.1) dans les catégories BIEXT$(P_S,Q_S;G_S)$ pour $S$ variable, $P,Q,G$ étant des Groupes commutatifs donnés sur $\underline{T}$.


}{\begin{equation*}
f = i_{S'} \; : \; \underline{T}_{/S'} \longrightarrow \underline{T} \qquad .
\end{equation*}
On appelle alors $f^{*}(E)$ la \underline{biextension induite} par $E$ sur $S'$, et on la notera aussi, comme d'habitude, $E|S'$ ou $E_{S'}$ . Comme l'opération de restriction possède évidemment les propriétés habituelles de transitivité, on en conclut que les catégories BIEXT$(\underline{T}_{/S})$, pour un objet variable $S$ de $\underline{T}$, forment les fibres d'une \underline{catégorie fibrée} sur $\underline{T}$ (SGA 1 VI 8). En fait, on vérifie trivialement que cette catégorie est même un \underline{champ} sur $\underline{T}$ au sens de Giraud [2], i.e. que "les morphismes et les objets des fibres de cette catégorie fibrée peuvent se recoller". Par exemple, pour se donner une biextension à isomorphisme canonique près sur $S$, il suffit de se donner une famille couvrante (i.e. épimorphique) $S_i \longrightarrow S$, pour tout $i \in I$ une biextension $E_i$ sur $S_i$ (i.e. dans $\underline{T}_{/S_i}$), pour tout couple $(i,j)$ un isomorphisme de $E_{iS_j}$ avec $E_{jS_i}$, ces isomorphismes étant soumis à la condition de cocycles habituelle. On appellera \underline{champ des biextensions sur} $\underline{T}$ le champ qu'on vient de définir. Le foncteur (2.3.3), donc aussi les foncteurs déduits (2.3.4) et (2.3.5), lorsqu'on y remplace $\underline{T}$ par $\underline{T}_{/S}$ avec $S$ variable, correspondent à des morphismes de champs sur $\underline{T}$. De même les opérations $E.E'$ et $E^{-1}$ (2.5.1) dans les catégories BIEXT$(P_S,Q_S;G_S)$ pour $S$ variable, $P,Q,G$ étant des Groupes commutatifs donnés sur $\underline{T}$.


}{F7I-VII-038}\SGApage{I}{172}




2.9. \underline{Relations avec les torseurs birigidifiés}

Avec la terminologie de 1.3.2, nous avons vu dans 2.2 que pour toute biextension $E$ de $(P,Q)$ par $G$, la structure de torseur sous $G_{P\times Q}$ sous-jacente à $E$ est canoniquement birigidifiée relativement au couple des sections unités de $P$ et de $Q$. On trouve ainsi un foncteur canonique

\begin{equation}\tag{2.9.1}
\mathrm{BIEXT}(P,Q;G) \xrightarrow{\;i\;} \mathrm{TORSBIRIG}(P,Q;G) \qquad ,
\end{equation}
où le deuxième membre désigne la catégorie des torseurs birigidifiés sur $P\times Q$ (sous $G_{P\times Q}$).

Utilisant les sections unités de $P$ et de $Q$, on trouve, pour tout produit de $n$ objets isomorphes à $P$ ou $Q$, un morphisme canonique $inj_i$ du $i$.ème facteur dans ledit produit. Si on se donne un morphisme de ce produit dans $G$, on pourra ainsi parler de la restriction de ce morphisme au $i$.ème facteur. Ceci posé, nous considérons les produits

\begin{equation}\tag{2.9.2}
\left\{ \begin{array}{l} P\times Q \\[2mm] P\times Q\times Q, \; P\times Q\times Q\times Q \\[2mm] P\times P\times Q, \; P\times P\times P\times Q \\[2mm] P\times P\times Q\times Q \end{array} \right. \qquad ,
\end{equation}
\underline{Proposition} 2.9.3. \underline{Supposons que tout morphisme d'un des produits} (2.9.2) \underline{dans} $G$ \underline{qui est trivial sur chacun des facteurs de ce produit} (\underline{plongés dans ledit grâce aux sections unité comme on vient de l'expliquer}) \underline{soit trivial}. \underline{Alors} :

a)  \underline{Le foncteur canonique} (2.9.1) \underline{est pleinement fidèle}, \underline{et les deux membres de} 2.9.1 \underline{sont des catégories rigides}. \underline{En particulier}, \underline{si} $E$ \underline{est un torseur sous} $G_{P\times Q}$ \underline{birigidifié pour} $(e_P,e_Q)$, \underline{il existe sur} $E$ \underline{au plus une structure de biextension compatible avec sa structure de torseur birigidifié}.


\SGApage{I}{173}

b)\quad \underline{Pour qu'il existe une telle structure de biextension sur} $E$, \underline{il faut et il suffit qu'il existe sur} $E$ \underline{une structure d'extension de} $Q_P$ \underline{par} $G_P$ \underline{compatible avec sa structure de torseur sous} $G_P$, \underline{et qu'il existe sur} $E$ \underline{une structure d'extension de} $P_Q$ \underline{par} $G_Q$ \underline{compatible avec sa structure de torseur sous} $G_Q$ (cf. 1.3.5 b)). \underline{Il existe dans ce cas sur} $E$ \underline{une unique structure d'extension de} $Q_P$ \underline{par} $G_P$ \underline{compatible avec les structures de torseur sous} $G_P$ \underline{et la rigidification donnée sur} $P\times e_Q$, \underline{et une unique structure d'extension de} $P_Q$ \underline{par} $G_Q$ \underline{compatible avec la structure de torseur sous} $G_Q$ \underline{et la rigidification  donnée sur} $e_P\times Q$, \underline{et ces deux structures d'extension sont compatibles} (2.1).

\SGAvar{\underline{Démonstration}. a)\quad Si $E$ est un torseur birigidifié sur $P\times Q$, le groupe de ses automorphismes s'identifie au groupe des morphismes $P\times Q \longrightarrow G$ qui sont triviaux sur les deux facteurs partiels, donc d'après l'hypothèse appliqué au produit $P\times Q$, c'est l'automorphisme identique. Cela prouve l'assertion de rigidité. Pour montrer que si $E,E'$ sont des biextensions de $(P,Q)$ par $G$, tout morphisme de bitorseurs birigidifiés $f : E \longrightarrow E'$ est un morphisme de biextensions, il suffit d'appliquer 1.3.5 a) aux deux structures d'extension définies par $E,E'$ sur $P$ resp. sur $Q$ (ce qui amène à utiliser l'hypothèse pour les deux produits $P\times Q\times Q$ et $P\times P\times Q$ de (2.9.2)).

b)\quad La nécessité de la condition énoncée sur le torseur birigidifié $E$ pour qu'il provienne d'une biextension étant claire, prouvons la suffisance. Il résulte d'abord de 1.3.5 b) que la condition énoncée implique l'existence de structures d'extensions uniques de $Q_P$ par $G_P$ et de $P_Q$ par $G_P$ , compatibles avec la structure de torseur et les rigidifications. (On utilise l'hypothèse pour les produits $P\times Q\times Q\times Q$ et $P\times P\times P\times Q$
\SGApage{I}{174}
dans (2.9.2)). Il reste à prouver que ces deux structures d'extension sont compatibles, i.e. qu'on a commutativité dans le diagramme (2.1.1). Or les deux composés de ce diagramme sont des isomorphismes de torseurs sur $P\times P\times Q\times Q$, donc ils diffèrent par un morphisme

}{\underline{Démonstration}. a)\quad Si $E$ est un torseur birigidifié sur $P\times Q$, le groupe de ses automorphismes s'identifie au groupe des morphismes $P\times Q \longrightarrow G$ qui sont triviaux sur les deux facteurs partiels, donc d'après l'hypothèse appliquée au produit $P\times Q$, c'est l'automorphisme identique. Cela prouve l'assertion de rigidité. Pour montrer que si $E,E'$ sont des biextensions de $(P,Q)$ par $G$, tout morphisme de bitorseurs birigidifiés $f : E \longrightarrow E'$ est un morphisme de biextensions, il suffit d'appliquer 1.3.5 a) aux deux structures d'extension définies par $E,E'$ sur $P$ resp. sur $Q$ (ce qui amène à utiliser l'hypothèse pour les deux produits $P\times Q\times Q$ et $P\times P\times Q$ de (2.9.2)).

b)\quad La nécessité de la condition énoncée sur le torseur birigidifié $E$ pour qu'il provienne d'une biextension étant claire, prouvons la suffisance. Il résulte d'abord de 1.3.5 b) que la condition énoncée implique l'existence de structures d'extensions uniques de $Q_P$ par $G_P$ et de $P_Q$ par $G_Q$ , compatibles avec la structure de torseur et les rigidifications. (On utilise l'hypothèse pour les produits $P\times Q\times Q\times Q$ et $P\times P\times P\times Q$
\SGApage{I}{174}
dans (2.9.2)). Il reste à prouver que ces deux structures d'extension sont compatibles, i.e. qu'on a commutativité dans le diagramme (2.1.1). Or les deux composés de ce diagramme sont des isomorphismes de torseurs sur $P\times P\times Q\times Q$, donc ils diffèrent par un morphisme

}{F7I-VII-039}\begin{equation*}
u : P\times P\times Q\times Q \longrightarrow G \qquad .
\end{equation*}
Utilisant le fait que les trivialisations partielles de $E$ sont des sections unités pour les lois $\varphi$ et $\psi$ respectivement, on en conclut facilement que $u$ est trivial sur chacun de ces quatre facteurs, donc par l'hypothèse appliquée au produit $P\times P\times Q\times Q$ de (2.9.2), qu'il est trivial, ce qui achève la démonstration.

\underline{Corollaire} 2.9.4. \underline{Soient} $P$, $Q$, $G$ \underline{des Groupes commutatifs satisfaisant les relations}

\SGAvar{\begin{equation}\tag{2.9.4.1}
\underline{\mathrm{Hom}}_{pt}(P,G) = e \quad , \quad \underline{\mathrm{Hom}}_{pt}(Q,G) = e
\end{equation}
(cf. 1.3.7). \underline{Alors les conditions préliminaires de} 2.9.2 \underline{sont satisfaites, donc les conclusions} a) \underline{et} b) \underline{de cette proposition sont valables}, \underline{De plus, si} $E$ \underline{est un torseur sous} $G_{P\times Q}$ \underline{birigidifié pour} $(e_P,e_Q)$, \underline{pour qu'il admette une structure de biextension de} $(P,Q)$ \underline{par} $G$ \underline{compatible avec ses structures données}, \underline{il faut et il suffit que les deux morphismes de faisceaux d'ensembles ponctués qu'il définit},

}{\begin{equation}\tag{2.9.4.1}
\underline{\mathrm{Hom}}_{pt}(P,G) = e \quad , \quad \underline{\mathrm{Hom}}_{pt}(Q,G) = e
\end{equation}
(cf. 1.3.7). \underline{Alors les hypothèses de la Proposition} 2.9.3 \underline{sont satisfaites, donc les conclusions} a) \underline{et} b) \underline{de cette proposition sont valables}. \underline{De plus, si} $E$ \underline{est un torseur sous} $G_{P\times Q}$ \underline{birigidifié pour} $(e_P,e_Q)$, \underline{pour qu'il admette une structure de biextension de} $(P,Q)$ \underline{par} $G$ \underline{compatible avec ses structures données}, \underline{il faut et il suffit que les deux morphismes de faisceaux d'ensembles ponctués qu'il définit},

}{F7I-VII-040}\begin{equation}\tag{2.9.4.2}
Q \longrightarrow R^1 p_*(G_P) \quad , \quad P \longrightarrow R^1 q_*(G_Q)
\end{equation}
($p,q$ \underline{les morphismes structuraux} $P \to e$ et $Q \to e$), \underline{soient des morphismes de Groupes}.

\SGAvar{La première assertion de 2.9.4 a été vue dans 1.3.7. Pour la deuxième, la nécessité de la condition envisagée étant claire, prouvons
\SGApage{I}{175}
la suffisance, en appliquant 2.9.2 b) et 1.3.5 b). Notons que la multiplicativité du premier morphisme de (2.9.4.2) signifie que le torseur rigidifié $E$ sous $G_P$ satisfait à la condition de 1.3.5 b) \underline{localement sur} $P$. Nous avons déjà signalé dans 1.3.7 que cela implique que cette condition est vraie globalement sur $P$, donc que sur $E$ existe une structure d'extension de $Q_P$ par $G_P$ compatible avec sa structure de torseur rigidifié sous $G_P$. On procède de même pour voir que sur $E$ existe une structure d'extension de $P_Q$ par $G_Q$ compatible avec sa structure de torseur rigidifié sous $G_Q$. La conclusion résulte donc de 2.9.2 b).

\underline{Exemple} 2.9.5. Soient $S$ un schéma, $P$ et $Q$ deux schémas abéliens sur $S$, et $G=\underline{G}_{mS}$ le schéma en groupes multiplicatif sur $S$. On a signalé dans 1.3.8 qu'alors les conditions (2.9.4.1) sont vérifiées. D'autre part, la donnée d'un torseur birigidifié $E$ sur $P\times Q$ revient alors à la donnée d'un Module inversible birigidifié sur $P\times Q$. Rappelons qu'on appelle \underline{correspondance divisorielle sur} $P,Q$ la classe à isomorphisme près d'un tel Module inversible birigidifié (cf. par exemple [1]). Travaillant avec la topologie fppf, les morphismes (2.9.4.2) sont ici les morphismes

}{La première assertion de 2.9.4 a été vue dans 1.3.7. Pour la deuxième, la nécessité de la condition envisagée étant claire, prouvons
\SGApage{I}{175}
la suffisance, en appliquant 2.9.3 b) et 1.3.5 b). Notons que la multiplicativité du premier morphisme de (2.9.4.2) signifie que le torseur rigidifié $E$ sous $G_P$ satisfait à la condition de 1.3.5 b) \underline{localement sur} $P$. Nous avons déjà signalé dans 1.3.7 que cela implique que cette condition est vraie globalement sur $P$, donc que sur $E$ existe une structure d'extension de $Q_P$ par $G_P$ compatible avec sa structure de torseur rigidifié sous $G_P$. On procède de même pour voir que sur $E$ existe une structure d'extension de $P_Q$ par $G_Q$ compatible avec sa structure de torseur rigidifié sous $G_Q$. La conclusion résulte donc de 2.9.3 b).

\underline{Exemple} 2.9.5. Soient $S$ un schéma, $P$ et $Q$ deux schémas abéliens sur $S$, et $G=\underline{G}_{m,S}$ le schéma en groupes multiplicatif sur $S$. On a signalé dans 1.3.8 qu'alors les conditions (2.9.4.1) sont vérifiées. D'autre part, la donnée d'un torseur birigidifié $E$ sur $P\times Q$ revient alors à la donnée d'un Module inversible birigidifié sur $P\times Q$. Rappelons qu'on appelle \underline{correspondance divisorielle sur} $P,Q$ la classe à isomorphisme près d'un tel Module inversible birigidifié (cf. par exemple [1]). Travaillant avec la topologie fppf, les morphismes (2.9.4.2) sont ici les morphismes

}{F7I-VII-041}\begin{equation}\tag{2.9.5.1}
Q \longrightarrow \underline{Pic}_{P/S} \qquad , \qquad P \longrightarrow \underline{Pic}_{Q/S} \qquad ,
\end{equation}
définis par la correspondance divisorielle, et il est bien connu par le "théorème du carré" (loc. cit.) que ces morphismes sont nécessairement compatibles avec les structures de Groupes. Donc 2.9.4 nous apprend que \underline{la catégorie des biextensions de} $(P,Q)$ \underline{par} $G=G_{m\ S}$ \underline{est équivalente à la catégorie des Modules inversibles birigidifiés sur} $P\times Q$, et que c'est donc une catégorie rigide dont l'ensemble des classes à isomorphisme près s'identifie à l'ensemble des correspondances divisorielles sur $(P,Q)$.


\SGApage{I}{176}

Cette identification respecte d'ailleurs, manifestement, les lois de groupes naturelles sur ces deux ensembles (la première étant définie dans (2.5.2.2)).

\underline{Remarque} 2.9.6. On peut vérifier facilement que l'additivité du premier morphisme (2.9.4.2) équivaut à dire que le deuxième morphisme (2.9.4.2) se factorise par le sous-faisceau $\underline{Ext}^1(Q,G)$ (noter qu'en vertu de 1.3.5, le morphisme canonique

\begin{equation*}
\underline{Ext}^1(Q,G) \longrightarrow R^1q_*(G_Q)
\end{equation*}
est injectif). De même l'additivité du deuxième morphisme (2.9.4.2) signifie que le premier morphisme se factorise par $\underline{Ext}^1(P,G)$. De plus, on obtient de cette façon des isomorphismes de groupes

\SGAvar{\begin{equation}\tag{2.9.6.1}
\mathrm{Biext}(P,Q;G) \simeq \mathrm{Hom}(Q,\underline{Ext}^1(P,G)) \simeq \mathrm{Hom}(P,\underline{Ext}^1(Q,G)) \qquad .
\end{equation}
Dans le cas particulier envisagé dans 2.9.5, ces isomorphismes se réduisent aux isomorphismes bien connus

}{\begin{equation}\tag{2.9.6.1}
\mathrm{Biext}^{1}(P,Q;G) \simeq \mathrm{Hom}(Q,\underline{Ext}^1(P,G)) \simeq \mathrm{Hom}(P,\underline{Ext}^1(Q,G)) \qquad .
\end{equation}
Dans le cas particulier envisagé dans 2.9.5, ces isomorphismes se réduisent aux isomorphismes bien connus

}{F7I-VII-042}\begin{equation}\tag{2.9.6.2}
\mathrm{Corr}(P,Q) \simeq \mathrm{Hom}(P,Q') \simeq \mathrm{Hom}(Q,P') \qquad ,
\end{equation}
où $P' = \underline{Ext}^1(P,\underline{G}_m)$et $Q' = \underline{Ext}^1(Q,\underline{G}_m)$ sont les schémas abéliens duaux de $P$ et de $Q$, qui s'identifient respectivement à $\underline{Pic}^o_{P/S}$ et $\underline{Pic}^o_{Q/S}$. Dans les deux formules précédentes, les Hom désignent bien entendu les groupes d'homomorphisme \underline{de Groupes}.

\SGAvar{2.10.  \underline{Généralisations diverses de la notion de biextension}. Nous les signalons par acquit de conscience, car nous n'autons pas à les utiliser.

2.10.1. On peut définir la notion de biextension  de $(P,Q)$ par $G$ sans supposer $P$ ni $Q$ commutatifs, $G$ seul étant supposé commutatif : c'est un
\SGApage{I}{177}
torseur sous $G_{P\times Q}$, muni de structures d'entension centrale de $Q_P$ par $G_P$ et de $P_Q$ par $G_Q$, ces deux structures étant compatibles à la structure de torseur, et compatibles entre elles i.e. rendent commuatifs les diagrammes (2.1.1). La plupart des développements qui précèdent sont valables mutatis mutandis pour cette notion plus général. On fera attention que lorsque $P$ et $Q$ sont commutatifs également,la notion précédente est    plus générale que celle de 2.1, qui correspond au cas où les deux structures d'extension sous-jacentes à $E$ sont des structures de Groupes relatifs \underline{commutatifs}. On fera attention également que dans le cas non commutatif, la catégorie des biextensions de $(P,Q)$ par $G$ dépend encore additivement de $G$, mais non pas de $P$ et de $Q$.

}{2.10.  \underline{Généralisations diverses de la notion de biextension}. Nous les signalons par acquit de conscience, car nous n'aurons pas à les utiliser.

2.10.1. On peut définir la notion de biextension  de $(P,Q)$ par $G$ sans supposer $P$ ni $Q$ commutatifs, $G$ seul étant supposé commutatif : c'est un
\SGApage{I}{177}
torseur sous $G_{P\times Q}$, muni de structures d'extension centrale de $Q_P$ par $G_P$ et de $P_Q$ par $G_Q$, ces deux structures étant compatibles à la structure de torseur, et compatibles entre elles i.e. rendent commutatifs les diagrammes (2.1.1). La plupart des développements qui précèdent sont valables mutatis mutandis pour cette notion plus générale. On fera attention que lorsque $P$ et $Q$ sont commutatifs également,la notion précédente est    plus générale que celle de 2.1, qui correspond au cas où les deux structures d'extension sous-jacentes à $E$ sont des structures de Groupes relatifs \underline{commutatifs}. On fera attention également que dans le cas non commutatif, la catégorie des biextensions de $(P,Q)$ par $G$ dépend encore additivement de $G$, mais non pas de $P$ et de $Q$.

}{F7I-VII-043}2.10.2.  En paraphrasant 2.1, on définit (pour un entier donné $n \geq 1$) la notion de $n$-\underline{extension d'une suite de} $n$ \underline{Groupes commutatifs} $(P_1,P_2,...,P_n)$ par un Groupe commutatif $G$. Elle correspond à la donnée, sur un torseur $E$ sous $G_{P_1\times..\times P_n}$ , de $n$ structures d'extensions partielles sur les produits partiels de $n-1$ facteurs, avec une relation de compatibilité du type (2.1.1) pour chaque couple d'indices distincts entre 1 et $n$.

\SGAvar{          Lorsque $n \geq 3$ et que les $P_i$ sont des schémas abéliens sur un schéma $S$, $G$ étant le groupe multiplicatif sur $\underline{G}_m$ (comparer 2.9.2), on trouve que toute $n$-extension de $(P_1,...,P_n)$ par $G$ est \underline{triviale}.

2.10.3.  La notion de biextension 2.1 peut-être considérée comme une variante de la notion d'extension commutative de Groupes commutatifs, i.e. de la notion d'extension de $\mathbb{Z}$-Modules. Lorsqu'on se donne un Anneau (commutatif pour simplifier) $A$ de $\underline{T}$, et trois $A$-Modules $P,Q,G$, on peut
\SGApage{I}{178}
définir plus généralement la notion de \underline{biextension de} $A$-\underline{Modules} de $(P,Q)$ par $G$, qui est une variante de celle d'extension de $A$-Modules. Une telle biextension est définie comme une biextension au sens de 2.1, mais les structures partielles d'extension sous-jacentes étant enrichies en des structures d'extensions de $A_P$-Modules resp. de $A_Q$-Modules (compatibles avec les structures de Modules relatifs qu'on a déjà sur $Q_P,G_P$ et $P_Q,G_Q$). Les deux structures de Modules relatifs sur $E$ peuvent s'interpréter en termes de deux opérations du monoïde multiplicatif sous-jacent à $A$ sur le faisceau d'ensembles sous-jacent à $E$, et on suppose que ces deux actions commutent l'une à l'autre; cette compatibilité s'ajoute donc à celle exprimée par la commutativité du diagramme (2.1.1), et s'exprime par la commutativité d'un diagramme analogue, que nous laissons au lecteur le soin d'écrire. Enfin, on suppose que la première opération de $A$ sur $E$ soit distributive, non seulement par rapport à la première loi de composition de $E$, mais aussi pour la seconde loi de composition de $E$ ; et de même pour la deuxième opération de $A$ sur $E$.

}{          Lorsque $n \geq 3$ et que les $P_i$ sont des schémas abéliens sur un schéma $S$, $G=\underline{G}_{m,S}$ étant le groupe multiplicatif sur $S$ (comparer 2.9.3), on trouve que toute $n$-extension de $(P_1,...,P_n)$ par $G$ est \underline{triviale}.

2.10.3.  La notion de biextension 2.1 peut être considérée comme une variante de la notion d'extension commutative de Groupes commutatifs, i.e. de la notion d'extension de $\mathbb{Z}$-Modules. Lorsqu'on se donne un Anneau (commutatif pour simplifier) $A$ de $\underline{T}$, et trois $A$-Modules $P,Q,G$, on peut
\SGApage{I}{178}
définir plus généralement la notion de \underline{biextension de} $A$-\underline{Modules} de $(P,Q)$ par $G$, qui est une variante de celle d'extension de $A$-Modules. Une telle biextension est définie comme une biextension au sens de 2.1, mais les structures partielles d'extension sous-jacentes étant enrichies en des structures d'extensions de $A_P$-Modules resp. de $A_Q$-Modules (compatibles avec les structures de Modules relatifs qu'on a déjà sur $Q_P,G_P$ et $P_Q,G_Q$). Les deux structures de Modules relatifs sur $E$ peuvent s'interpréter en termes de deux opérations du monoïde multiplicatif sous-jacent à $A$ sur le faisceau d'ensembles sous-jacent à $E$, et on suppose que ces deux actions commutent l'une à l'autre; cette compatibilité s'ajoute donc à celle exprimée par la commutativité du diagramme (2.1.1), et s'exprime par la commutativité d'un diagramme analogue, que nous laissons au lecteur le soin d'écrire. Enfin, on suppose que la première opération de $A$ sur $E$ soit distributive, non seulement par rapport à la première loi de composition de $E$, mais aussi pour la seconde loi de composition de $E$ ; et de même pour la deuxième opération de $A$ sur $E$.

}{F7I-VII-044}          On observera que le dictionnaire explicité dans 1.1.6 admet aussi une variante $A$-linéaire, qui peut s'exprimer en disant que la donnée d'une extension du $A$-Module $P$ par un autre $G$ revient à la donnée d'une famille de torseurs $E_p$ sous les $G_p$, paramétrée par les "points" $p \in P(S)$ ($S \in \mathrm{ob}\ \underline{T}$ étant un objet variable), $E_p$ étant assujetti à varier de façon $A$-\underline{linéaire} (et non seulement additive comme dans 1.1.6) en $p$. Ceci s'explicite par la donnée, en plus des isomorphismes $\varphi_{p,p'}$ (1.1.3.1), par la donnée d'isomorphismes

\SGAvar{\begin{equation*}
\mu_{a,p} \ :\ {}^a(P_p) \overset{\sim}{\longrightarrow} P_{ap} \qquad\quad (p \in P(S),\ a \in A(S)) \qquad ,
\end{equation*}
où le premier  membre désigne l'image du $G_S$-torseur $P_p$ par l'extension
\SGApage{I}{179}
des scalaires $g$ $\longmapsto$ $ag$ : $G_S \longrightarrow G_S$ . On laisse au lecteur le soin d'expliciter les conditions à imposer à ces données $\varphi$ , $\mu$ pour qu'elles correspondent bien à une extension de Modules $E$ de $P$ par $G$  (cf 3.5.4 b)).


}{\begin{equation*}
\mu_{a,p} \ :\ {}^a(E_p) \overset{\sim}{\longrightarrow} E_{ap} \qquad\quad (p \in P(S),\ a \in A(S)) \qquad ,
\end{equation*}
où le premier  membre désigne l'image du $G_S$-torseur $E_p$ par l'extension
\SGApage{I}{179}
des scalaires $g$ $\longmapsto$ $ag$ : $G_S \longrightarrow G_S$ . On laisse au lecteur le soin d'expliciter les conditions à imposer à ces données $\varphi$ , $\mu$ pour qu'elles correspondent bien à une extension de Modules $E$ de $P$ par $G$  (cf 3.5.4 b)).


}{F7I-VII-045}\SGAsectiontoc{3. Compléments d'algèbre homologique}{3}
3. \underline{Compléments d'algèbre homologique}

Le but principal du présent numéro est la définition de l'isomorphisme

\begin{equation*}
\mathrm{Biext}^1(P,Q;G) \;\simeq\; \mathrm{Ext}^1(P \overset{\mathbb{L}}{\otimes} Q,G)
\end{equation*}
annoncé plus haut. Nous aurons besoin pour ceci de certains développements d'algèbre homologique, ayant un intérêt indépendant, qui occuperont les sections 3.1  à 3.4. Les développements 3.1 et 3.3 peuvent être considérés comme des compléments à [4], et notamment à la section 6.9 de loc. cit.


3.1.  \underline{Construction d'une catégorie cofibrée additive} $\underline{\Psi}_{L.}$ \underline{en termes d'un complexe de chaînes} $L.$ \underline{dans une catégorie abélienne} $\underline{C}$.

Soit $\underline{C}$ une catégorie abélienne, $L.$ un complexe de chaînes de $\underline{C}$,

\SGAvar{\begin{equation}\tag{3.1.1}
L_2 \longrightarrow L_1 \longrightarrow L_o \longrightarrow 0 \qquad ,
\end{equation}
que nous écrirons avec la convention des degrés inférieurs (opérateur différentiel de degré $-1$). On a donc $L_i = 0$ pour $i < 0$. Nous allons définir une catégorie cofibré $\underline{\Psi} = \underline{\Psi}_{L.}$ sur $\underline{C}$ (qui,à équivalence canonique de catégories cofibrées près,sera la catégorie de même nom de loc. cit. 6.9). Les objets de $\underline{\Psi}_{L.}$ sont des couples

}{\begin{equation}\tag{3.1.1}
L_2 \longrightarrow L_1 \longrightarrow L_o \longrightarrow 0 \qquad ,
\end{equation}
que nous écrirons avec la convention des degrés inférieurs (opérateur différentiel de degré $-1$). On a donc $L_i = 0$ pour $i < 0$. Nous allons définir une catégorie cofibrée $\underline{\Psi} = \underline{\Psi}_{L.}$ sur $\underline{C}$ (qui,à équivalence canonique de catégories cofibrées près,sera la catégorie de même nom de loc. cit. 6.9). Les objets de $\underline{\Psi}_{L.}$ sont des couples

}{F7I-VII-046}\begin{equation*}
(X , \alpha) \qquad , 
\end{equation*}

\SGApage{I}{180}

 où $X$ est une extension de $L_o$ par un objet $A$ de $\underline{C}$ (dépendant de $X$)

\begin{equation}\tag{3.1.2}
0 \longrightarrow A \overset{i}{\longrightarrow} X \overset{j}{\longrightarrow} L_o \longrightarrow 0 \qquad ,
\end{equation}
et où $\alpha$ est une trivialisation de l'image inverse $d^*_o(X)$ de cette extension par $d_o:L_1 \longrightarrow L_o$, $\alpha$ étant assujetti à la condition que la trivialisation correspondante de l'image inverse $d^*_1(d^*_o(X))$ par $d_1 : L_2 \longrightarrow L_1$ est la trivialisation évidente, résultant de l'isomorphisme de transitivité

\begin{equation*}
d^*_1(d^*_o(E)) \;\simeq\; (d_o d_1)^*(E) \qquad ,
\end{equation*}
et de la relation $d_od_1=0$. Une autre façon d'exprimer la donnée de $\alpha$, soumis à la condition précédente,est de dire qu'on se donne un relèvement de $d_o : L_1 \longrightarrow L_o$ en $\alpha : L_1 \longrightarrow X$, soumis à la condition $\alpha d_1 = 0$ :

(3.1.3) \begin{tikzcd}[row sep=large, column sep=large]
 \& \& 0 \arrow[d] \\
 \& \& A \arrow[d, "i"] \\
 \& \& X \arrow[d, "j"] \\
L_2 \arrow[r, "d_1"'] \& L_1 \arrow[r, "d_o"'] \arrow[ru, "\alpha"] \& L_o \arrow[d] \\
 \& \& 0
\end{tikzcd}\qquad , \quad $\alpha\, d_1 = 0 \quad . \;(*)$
















Sous cette forme, on voit que la description des objets $(X,\alpha)$ ne dépend que du complexe de chaînes $[L.]$ \underline{tronqué à l'ordre} 1, déduit de $L.$ en "tuant les $H_i(L.)$ pour $i \geq 2$", dont les composantes sont nulles en dimension $\geq 2$, coïncident avec celles de $L.$ en dimension $\leq 0$, et dont la composante en dimension 1 est le conoyau de $d_1: L_2 \longrightarrow L_1$. On définit les homomorphismes de $(X,\alpha)$ dans $(X',\alpha')$ comme les homomorphismes d'extensions $u: X \longrightarrow X'$, compatibles avec l'identité sur $L_o$, tels que $u \alpha = \alpha'$ ; on 

(*)  Une façon plus simple de décrire  $\psi_{L.}$ $(A)$  (signalée par DELIGNE) est de le considérer comme catégorie des extensions de complexes de  $L.$  par  $A$ (regardé comme complexe réduit au degré zéro).


\SGApage{I}{181}

 compose les homomorphismes en composant les homomorphismes d'extensions correspondants. On trouve ainsi la catégorie annoncée $\underline{\Psi}_{L.}$ . Elle ne dépend, à isomorphisme canonique près, que du tronqué à l'ordre 1 $[L.]$ de $L.$ .

On a un foncteur évident, $(X,\alpha) \longmapsto A$ :

\begin{equation}\tag{3.1.4}
\pi : \underline{\Psi}_{L.} \longrightarrow \underline{C} \qquad .
\end{equation}
\underline{Ce foncteur est cofibrant}. En effet, soit

\begin{equation*}
u: A \longrightarrow B
\end{equation*}
un morphisme de $\underline{C}$, et $(X,\alpha)$ un objet de la catégorie $\underline{\Psi}_{L.}$ $(A)$ fibre en $A$. Donc $X$ est une extension de $L_o$ par $A$, et on peut considérer l'extension $u_*(X)$ de $L_o$ par $B$, image directe de $X$ par $u$, et la munir de la trivialisation partielle composée $\beta = f\alpha$, où $f : X \longrightarrow u_*(X)$ est le morphisme canonique. On vérifie alors immédiatement que le morphisme $f : (X,\alpha) \longrightarrow (u_*(X),\beta)$ est un morphisme \underline{cocartésien} pour le foncteur (3.1.4), i.e. satisfait à la propriété universelle habituelle pour les $u$-morphismes dans $\underline{\Psi}_{L.}$ de source $(X,\alpha)$ ; il suffit pour ceci d'utiliser la propriété universelle analogue pour le morphisme $f : X \longrightarrow u_*(X)$, qui est bien connue et triviale.

Enfin, la catégorie cofibrée $\underline{\Psi}_{L.}$ est \underline{additive}, i.e. elle satisfait les deux conditions :

\SGAvar{\begin{equation}\tag{3.1.5}
\underline{\Psi}_{L.} (O) \ \text{est équivalente à la catégorie ponctuelle} \qquad ,
\end{equation}
\begin{equation}\tag{3.1.6}
\underline{\Psi}_{L.} (A \times B) \longrightarrow \underline{\Psi}_{L}(A) \times \underline{\Psi}_{L.} (B) \ \text{est une équivalence} \qquad ,
\end{equation}
pour deux objets quelconques $A$ et $B$ de $\underline{C}$. La vérification est triviale, en termes des propriétés correspondantes pour la catégorie cofibrée des extensions de $L_o$ par des objets variables de $\underline{C}$. Nous utiliserons librement
\SGApage{I}{182}
les conséquences de l'additivité de $\underline{\Psi}_{L.}$, établies dans (loc. cit. 1) et rappelées plus haut (2.5).

}{\begin{equation}\tag{3.1.5}
\underline{\Psi}_{L.} (O) \ \text{est équivalente à la catégorie ponctuelle} \qquad ,
\end{equation}
\begin{equation}\tag{3.1.6}
\underline{\Psi}_{L.} (A \times B) \longrightarrow \underline{\Psi}_{L.}(A) \times \underline{\Psi}_{L.} (B) \ \text{est une équivalence} \qquad ,
\end{equation}
pour deux objets quelconques $A$ et $B$ de $\underline{C}$. La vérification est triviale, en termes des propriétés correspondantes pour la catégorie cofibrée des extensions de $L_o$ par des objets variables de $\underline{C}$. Nous utiliserons librement
\SGApage{I}{182}
les conséquences de l'additivité de $\underline{\Psi}_{L.}$, établies dans (loc. cit. 1) et rappelées plus haut (2.5).

}{F7I-VII-047}3.1.7.  Parmi ces conséquences, mentionnons le fait que pour tout objet $A$ de $\underline{C}$, il y a dans $\underline{\Psi}_{L.}$ $(A)$ un \underline{objet nul}, ou \underline{trivial} (2.5.1) ; on vérifie qu'il est défini par l'extension triviale $X = L_o + A$ de $L_o$ par $A$, munie de la trivialisation partielle définie par $\alpha : L_1 \longrightarrow X$ de composantes $d_o$ et $O$. D'autre part (2.5.2), pour tout objet $(X,\alpha)$ de $\underline{\Psi}_{L.}$ $(A)$, le groupe des automorphismes de cet objet est canoniquement isomorphe au groupe des automorphismes de l'objet trivial de $\underline{\Psi}_{L.}$ $(A)$, soit $\Psi^{o}_{L.}$ $(A)$, qui est un groupe commutatif. Bien entendu, on retrouve trivialement ici ce fait sur la description explicite de $\underline{\Psi}_{L.}$ $(A)$, et on trouve même un isomorphisme canonique

\SGAvar{\begin{multline}\tag{3.1.8}
\Psi^{o}_{L.} (A) = \mathrm{Ker}(\mathrm{Hom}(L_o,A) \to \mathrm{Hom}(L_1,A)) \simeq \mathrm{Hom}(H_o(L.),A) \\ \simeq \mathrm{Ext}^{o}(L.,A) \qquad ,
\end{multline}
induit par l'isomorphisme habituel du groupe des automorphismes d'extension de $X$ avec le groupe $\mathrm{Hom}(L_o,A)$. L'isomorphisme (3.1.7.1) est évidemment fonctoriel en $A$.

3.1.9.  Rappelons enfin (2.5.2) que grâce à la structure multiplicative naturelle dans la catégorie fibre $\underline{\Psi}_{L.}$ $(A)$, l'ensemble des classes d'isomorphie de celle-ci est munie de façon naturelle d'une structure de groupe, que nous noterons $\Psi^{1}_{L.}$ $(A)$, et que ce dernier est un foncteur additif par rapport à $A$.


}{\begin{multline}\tag{3.1.8}
\Psi^{o}_{L.} (A) = \mathrm{Ker}(\mathrm{Hom}(L_o,A) \to \mathrm{Hom}(L_1,A)) \simeq \mathrm{Hom}(H_o(L.),A) \\ \simeq \mathrm{Ext}^{o}(L.,A) \qquad ,
\end{multline}
induit par l'isomorphisme habituel du groupe des automorphismes d'extension de $X$ avec le groupe $\mathrm{Hom}(L_o,A)$. L'isomorphisme (3.1.8) est évidemment fonctoriel en $A$.

3.1.9.  Rappelons enfin (2.5.2) que grâce à la structure multiplicative naturelle dans la catégorie fibre $\underline{\Psi}_{L.}$ $(A)$, l'ensemble des classes d'isomorphie de celle-ci est muni de façon naturelle d'une structure de groupe, que nous noterons $\Psi^{1}_{L.}$ $(A)$, et que ce dernier est un foncteur additif par rapport à $A$.


}{F7I-VII-048}\SGApage{I}{183}


3.2. \underline{Détermination de} $\Psi^1_{L.} (A)$

Considérons un objet $(X,\alpha)$ (3.1.3) de $\underline{\Psi}_{L.} (A)$. Considérons le complexe de cochaînes de longueur $1$

\SGAvar{\begin{equation}\tag{3.2.1}
[X \to L_o] : \; 0 \longrightarrow X \longrightarrow L_o \longrightarrow 0 \ldots
\end{equation}
défini par le morphisme $j : X \longrightarrow L_o$, la composante de degré $0$ étant $X$. C'est une résolution de $A$, donc peut-être identifié à $A$ en tant qu'objet de la catégorie dérivée $D(\underline{C})$ [5]. Ceci posé, le diagramme (3.1.3) nous fournit un homomorphisme de degré $1$ de complexes

}{\begin{equation}\tag{3.2.1}
[X \to L_o] : \; 0 \longrightarrow X \longrightarrow L_o \longrightarrow 0 \ldots
\end{equation}
défini par le morphisme $j : X \longrightarrow L_o$, la composante de degré $0$ étant $X$. C'est une résolution de $A$, donc peut être identifié à $A$ en tant qu'objet de la catégorie dérivée $D(\underline{C})$ [5]. Ceci posé, le diagramme (3.1.3) nous fournit un homomorphisme de degré $1$ de complexes

}{F7I-VII-049}\begin{equation}\tag{3.2.2}
L. \longrightarrow [X \longrightarrow L_o] \quad ,
\end{equation}
explicité dans le diagramme à carrés anticommutatifs

\begin{equation}\tag{3.2.3}
\begin{tikzcd}[column sep=large, row sep=huge]
\cdots \& L_2 \arrow[r, "d_1"] \& L_1 \arrow[r, "d_o"] \arrow[d, "u^{-1} = -\alpha"'] \& L_o \arrow[r] \arrow[d, "u^o = \mathrm{id}"'] \& 0 \& \cdots \\
\cdots \& 0 \arrow[r] \& X \arrow[r, "j"] \& L_o \arrow[r] \& 0 \& \cdots
\end{tikzcd}\quad .
\end{equation}















(On utilisera la convention habituelle, associant au complexe de chaînes $L.$ un complexe à opérateur différentiel de degré $1$, dont la composante de degré $i$ est $L_{-i}$.). On trouve donc ainsi un homomorphisme de degré $1$ dans la catégorie dérivée

\begin{equation}\tag{3.2.3}
c(X,\alpha) : L. \longrightarrow A \qquad (\text{hom. de degré 1 dans } D(\underline{C})) \quad ,
\end{equation}
qui ne dépend manifestement que de la classe à isomorphie près de $(X,\alpha)$. On a défini ainsi une application canonique

\begin{equation}\tag{3.2.4}
\Psi^1_{L.} (A) \longrightarrow \mathrm{Ext}^1(L., A) \overset{\mathrm{dfn}}{=} \mathrm{Hom}_{D(\underline{C})}(L., A[1]) \quad .
\end{equation}

\SGApage{I}{184}

\underline{Théorème} 3.2.5. \underline{L'application} (3.2.4) \underline{est un isomorphisme de groupes,} \underline{fonctoriel en l'objet} $A$ \underline{de} $\underline{C}$.

\SGAvar{Nous laisserons au lecteur le soin d'établir l'additivité de l'application (3.2.4), et sa fonctorialité pour $A$, et nous bornerons à prouver qu'elle est bijective. Pour l'injectivité, montrons d'abord que si l'homomorphisme de complexes $u$ de (3.2.3) est nul dans $D(\underline{C})$, alors il est déjà nul dans la catégorie $K(\underline{C})$ des complexes à homotopie près, i.e. il existe un morphisme

}{Nous laisserons au lecteur le soin d'établir l'additivité de l'application (3.2.4), et sa fonctorialité pour $A$, et nous nous bornerons à prouver qu'elle est bijective. Pour l'injectivité, montrons d'abord que si l'homomorphisme de complexes $u$ de (3.2.3) est nul dans $D(\underline{C})$, alors il est déjà nul dans la catégorie $K(\underline{C})$ des complexes à homotopie près, i.e. il existe un morphisme

}{F7I-VII-050}\begin{equation*}
s : L_o \longrightarrow X
\end{equation*}
tel que l'on ait

\begin{equation}\tag{3.2.5.1}
js = u^o = \mathrm{id}_{L_o} \ , \ sd_o = - u^{-1} = \alpha \quad .
\end{equation}
En effet, l'hypothèse signifie que la conclusion devient vraie après avoir composé le morphisme de complexes envisagé avec un quasi-isomorphisme convenable $v$

\begin{equation}\tag{*}
\begin{tikzcd}[column sep=large, row sep=huge]
0 \arrow[r] \& X \arrow[r, "j"] \arrow[d] \& L_o \arrow[r] \arrow[d] \& 0 \\
0 \arrow[r] \& C^o \arrow[r, "j'"] \& C^1 \arrow[r] \& C^2 \arrow[r] \& \cdots
\end{tikzcd}\quad ,
\end{equation}














où on peut d'ailleurs (quitte à remplacer $C^i$ par $0$ pour $i \geq 2$, par $\mathrm{Im}(C^o \longrightarrow C^1) = \mathrm{Ker}(C^1 \longrightarrow C^2)$ pour $i = 1$) supposer que $C^i = 0$ pour $i \geq 2$. Mais alors, le morphisme (*) induisant l'identité sur $A$, on voit qu'il identifie $X$ au produit fibré de $L_o$ et de $C^o$ sur $C^1$, et que la donnée d'un morphisme $s' : L_o \longrightarrow C^o$ définissant une homotopie $vu \sim 0$, donc satisfaisant $j's' = v^1u^o$, se factorise par un morphisme $s : L_o \longrightarrow X$ satisfaisant $js = u^o = \mathrm{id}_{L_o}$. Mais alors $s$ définit un splittage de l'extension $X$ de $L_o$ par $A$, donc le complexe $X \longrightarrow L_o$ est isomorphe \underline{dans} $K(\underline{C})$ au complexe $K^{\cdot}$
\SGApage{I}{185}
réduit à $A$. Or pour ce dernier, il est évident que l'homomorphisme

\begin{equation*}
\mathrm{Hom}_{K(\underline{C})}(L.,K^{\cdot}) \longrightarrow \mathrm{Hom}_{D(\underline{C})}(L.,K^{\cdot})
\end{equation*}
est un isomorphisme, d'où notre assertion.

\SGAvar{Ceci dit, la donnée d'un homomorphisme $j : L_o \longrightarrow X$ satisfaisant les conditions (3.2.5.1) équivaut précisément à une trivialisation de l'extension $X$ de $L_o$ par $A$, compatible avec la trivialisation partielle $\alpha$. Il en résulte bien l'injectivité de (3.2.4).

}{Ceci dit, la donnée d'un homomorphisme $s : L_o \longrightarrow X$ satisfaisant les conditions (3.2.5.1) équivaut précisément à une trivialisation de l'extension $X$ de $L_o$ par $A$, compatible avec la trivialisation partielle $\alpha$. Il en résulte bien l'injectivité de (3.2.4).

}{F7I-VII-051}Pour la surjectivité, partons d'un élément du deuxième membre, qui peut donc être décrit, par définition, à l'aide d'un homomorphisme de degré $1$ de complexes

\begin{equation}\tag{3.2.5.2}
\begin{tikzcd}
\ldots \& L_2 \arrow[r, "d_1"] \& L_1 \arrow[r, "d_o"] \arrow[d, "u^{-1}"'] \& L_o \arrow[r] \arrow[d, "u^o"'] \& 0 \\
 \& 0 \arrow[r] \& C^o \arrow[r] \& C^1 \arrow[r] \& \ldots
\end{tikzcd}
\end{equation}
















(carrés anticommutatifs), où $C^{\cdot}$ est une résolution de $A$, et où on peut supposer comme ci-dessus que $C^i=0$ pour $i \geqslant 2$. Ainsi $C^o$ apparaît comme une extension de $C^1$ par $A$, et on peut considérer l'extension $X$ image inverse de celle-ci par $L_o \to C^1$. L'anticommutativité du carré de (3.2.5.2) montre que $-u^{-1}$ se factorise par un morphisme $\alpha : L_1 \longrightarrow X$, qui est un splittage partiel i.e. satisfait $j\alpha = d_o$ ($j$ défini dans (3.1.2)). D'ailleurs la relation $u^{-1}d_1 = 0$ dans (3.2.5.2) implique que $\alpha d_1 = 0$. Donc $(X,\alpha)$ est bien un élément de $\underline{\Psi}_{L.}(A)$. De plus, par construction même, l'homomorphisme (3.2.5.2) est le composé de l'homomorphisme (3.2.3) relatif à $(X,\alpha)$, et de l'homomorphisme

\begin{equation*}
[X \longrightarrow L_o] \longrightarrow C^{\cdot} = [C^o \longrightarrow C^1] 
\end{equation*}

\SGApage{I}{186}

 de composantes$(f,u^o)$, où $f: X \longrightarrow C^o$ est l'homomorphisme canonique déduit de la définition de l'extension $X$ comme image inverse de l'extension $C^o$. Comme l'homomorphisme précédent est un homomorphisme de résolutions de $A$ (induisant l'identité sur $A$), on en conclut, par définition de $\mathrm{Ext}^1 = \mathrm{Hom}^1_{D(\underline{C})}$, que l'élément donné de $\mathrm{Ext}^1(L.,A)$, décrit par (3.2.5.2), n'est autre que $c(X,\alpha)$, ce qui établit la surjectivité de (3.2.4) et achève la démonstration de 3.2.5.


3.3. \underline{Propriétés de la catégorie cofibrée} $\underline{\Psi}_{L.}$

Dans le présent numéro, qui n'est pas essentiel pour la compréhension de la suite, nous utilisons librement les notions de [4], notre but étant de comparer les résultats précédents à ceux de loc. cit.

\SGAvar{\underline{Théorème} 3.3.1. \underline{La catégorie cofibrée additive} $\underline{\Psi}_{L.}$ \underline{sur} $\underline{C}$ est exacte à gauche (loc. cit. 2.2) \underline{elle admet un objet quasi-maximal} (loc. cit. 4.4), \underline{et pour tout objet} $E = (X,\alpha)$ \underline{de cette catégorie}, \underline{il existe un complexe typique} $L^E_{.}$ \underline{pour} $E$ (loc. cit. 3.1). \underline{De plus}, \underline{le pro-complexe typique de} $\underline{\Psi}_{L.}$ (loc. cit. 3.7 \underline{et} 4.4) \underline{est canoniquement isomorphe dans} $D(\underline{C})$ \underline{au complexe tronqué} $[L.]$ \underline{envisagé dans} 3.1.

}{\underline{Théorème} 3.3.1. \underline{La catégorie cofibrée additive} $\underline{\Psi}_{L.}$ \underline{sur} $\underline{C}$ est exacte à gauche (loc. cit. 2.2), \underline{elle admet un objet quasi-maximal} (loc. cit. 4.4), \underline{et pour tout objet} $E = (X,\alpha)$ \underline{de cette catégorie}, \underline{il existe un complexe typique} $L^E_{.}$ \underline{pour} $E$ (loc. cit. 3.1). \underline{De plus}, \underline{le pro-complexe typique de} $\underline{\Psi}_{L.}$ (loc. cit. 3.7 \underline{et} 4.4) \underline{est canoniquement isomorphe dans} $D(\underline{C})$ \underline{au complexe tronqué} $[L.]$ \underline{envisagé dans} 3.1.

}{F7I-VII-052}Nous explicitons ci-dessous la signification des diverses assertions de 3.3.1.

\SGAvar{3.3.2. Dire la catégorie cofibrée $\underline{\Psi}_{L.}$ est \underline{exacte à gauche} signifie qu'elle est additive et satisfait de plus la condition :

}{3.3.2. Dire que la catégorie cofibrée $\underline{\Psi}_{L.}$ est \underline{exacte à gauche} signifie qu'elle est additive et satisfait de plus la condition :

}{F7I-VII-053}(3.3.2.1) Pour toute suite exacte

\begin{equation*}
0 \longrightarrow A' \stackrel{u}{\longrightarrow} A \stackrel{v}{\longrightarrow} A'' \longrightarrow 0 
\end{equation*}

\SGApage{I}{187}

 de $\underline{C}$, le foncteur naturel (notations de loc. cit. 1.8)

\SGAvar{\begin{equation*}
\underline{\Psi}_{L.}\,(A) \longrightarrow \underline{\Psi}_{L.}\,([A \longrightarrow A''])
\end{equation*}
est une équivalence de catégories.

}{\begin{equation*}
\underline{\Psi}_{L.}\,(A') \longrightarrow \underline{\Psi}_{L.}\,([A \longrightarrow A''])
\end{equation*}
est une équivalence de catégories.

}{F7I-VII-054}Rappelons que, pour une catégorie cofibrée $\underline{\Psi}$ sur $\underline{C}$, et un morphisme $v : A \longrightarrow A''$ de $\underline{C}$, $\underline{\Psi}([A \longrightarrow A''])$ désigne la catégorie formée des couples $(E,\rho)$, où $E$ est un objet de $\underline{\Psi}(A)$ et où $\rho$ est une ``trivialisation'' de l'objet correspondant $v_*(E)$ de $\underline{\Psi}(A'')$, i.e. un isomorphisme de cet objet avec l'objet nul (2.5.1) de $\underline{\Psi}(A'')$. La définition du foncteur

\begin{equation*}
\underline{\Psi}(A') \longrightarrow \underline{\Psi}([A \longrightarrow A''])
\end{equation*}
de (3.3.2.1) est évidente. Dans le cas $\underline{\Psi} = \underline{\Psi}_{L.}$ , les objets de $\underline{\Psi}_{L.}([A \longrightarrow A''])$ s'explicitent comme les triples $(X,\alpha,\rho)$, où $(X,\alpha)$ est comme dans (3.1.3), et où $\rho : X \longrightarrow A''$ est un morphisme dans $\underline{C}$ satisfaisant

\SGAvar{\begin{equation}\tag{*}
\rho i = v \quad , \quad \rho\, \alpha = 0 \quad .
\end{equation}
Lorsque $v: A \longrightarrow A''$ s'insère dans une suite exacte courte comme dans 3.2.3, alors la donnée d'un morphisme satisfaisant la première des relations (*) équivaut, comme il est bien connu (exactitude à gauche de la catégorie cofibrée des extensions de $L_o$ par des objets variables de $\underline{C}$) à la donnée d'une extension $X'$ de $L_o$ par $A' \simeq \mathrm{Ker}\, v$ qui soit une sous-extension de $X$ (on prend $X' = \mathrm{Ker}\, \rho$), et la deuxième relation (*) exprime que $\alpha$ se factorise en un $\alpha' : L_1 \longrightarrow X'$ . Ceci établit, modulo des compatibilités dont la vérification est évidente, l'exactitude à gauche explicitée dans (3.2.3).


}{\begin{equation}\tag{*}
\rho i = v \quad , \quad \rho\, \alpha = 0 \quad .
\end{equation}
Lorsque $v: A \longrightarrow A''$ s'insère dans une suite exacte courte comme dans 3.3.2.1, alors la donnée d'un morphisme satisfaisant la première des relations (*) équivaut, comme il est bien connu (exactitude à gauche de la catégorie cofibrée des extensions de $L_o$ par des objets variables de $\underline{C}$) à la donnée d'une extension $X'$ de $L_o$ par $A' \simeq \mathrm{Ker}\, v$ qui soit une sous-extension de $X$ (on prend $X' = \mathrm{Ker}\, \rho$), et la deuxième relation (*) exprime que $\alpha$ se factorise en un $\alpha' : L_1 \longrightarrow X'$ . Ceci établit, modulo des compatibilités dont la vérification est évidente, l'exactitude à gauche explicitée dans (3.3.2.1).


}{F7I-VII-055}\SGApage{I}{188}

\SGAvar{3.3.3. Rappelons qu'un objet $E_o$ d'une catégorie cofibrée $\underline{\Psi}$ sur $\underline{C}$ est dit \underline{quasi-maximal} si pour tout objet $E$ de $\underline{\Psi}$, au-dessus d'un objet $A$ de $\underline{C}$, on peut trouver un morphisme $E \longrightarrow E'$ dans $\underline{\Psi}$, induisant un \underline{monomorphisme} $A \longrightarrow A'$ dans $\underline{C}$, de telle sorte que $E$ soit \underline{majoré par} $E_o$ i.e. qu'il existe un morphisme, $E_o \longrightarrow E'$. Lorsque dans $\underline{C}$ il existe suffisamment d'objets injectifs, il revient au même de dire que tout objet de $\underline{\Psi}$ au-dessus d'un objet injectif de $\underline{C}$ est majoré par $E_o$. Dans le cas de $\underline{\Psi} = \underline{\Psi}_{L.}$ , on a l'objet quasi-maximal canonique $E_o$ suivant. On pose

\begin{equation*}
L_1' = \mathrm{Coker}\ (d_1 : L_2 \longrightarrow L_1)
\end{equation*}
et on désigne par $E_o$ l'objet $(X_o,\alpha_o)$ de $\underline{\Psi}_{L.}(L_1')$, où $X_o = L_o + L_1'$ est l'extension triviale de $L_o$ par $L_1'$, et où $\alpha_o$ a comme composantes $(d_1',\mathrm{id}_{L_1'})$ , $d_1' : L_1 \longrightarrow L_o$ étant déduit de $d_1 : L_1 \longrightarrow L_o$ par passage au quotient.

Pour voir que $E_o$ est bien quasi-maximal, notons que pour tout objet $(X,\alpha)$ de $\underline{\Psi}_{L.}$ , on peut trouver un morphisme de celui-ci dans un $(X',\alpha')$, induisant un \underline{monomorphisme} $A \longrightarrow A'$ dans $\underline{C}$, tel que l'extension $X'$ de $L_o$ par $A'$ soit \underline{scindée} : il suffit en effet de prendre $A'=X$, et $(X',\alpha') = i_*(X,\alpha)$, les notations étant celles de (3.1.2). D'autre part, il est immédiat qu'un objet $(X',\alpha')$ de $\underline{\Psi}$ tel que l'extension $X'$ soit scindée, est majorée par $E_o$.

}{3.3.3. Rappelons qu'un objet $E_o$ d'une catégorie cofibrée $\underline{\Psi}$ sur $\underline{C}$ est dit \underline{quasi-maximal} si pour tout objet $E$ de $\underline{\Psi}$, au-dessus d'un objet $A$ de $\underline{C}$, on peut trouver un morphisme $E \longrightarrow E'$ dans $\underline{\Psi}$, induisant un \underline{monomorphisme} $A \longrightarrow A'$ dans $\underline{C}$, de telle sorte que $E'$ soit \underline{majoré par} $E_o$ i.e. qu'il existe un morphisme, $E_o \longrightarrow E'$. Lorsque dans $\underline{C}$ il existe suffisamment d'objets injectifs, il revient au même de dire que tout objet de $\underline{\Psi}$ au-dessus d'un objet injectif de $\underline{C}$ est majoré par $E_o$. Dans le cas de $\underline{\Psi} = \underline{\Psi}_{L.}$ , on a l'objet quasi-maximal canonique $E_o$ suivant. On pose

\begin{equation*}
L_1' = \mathrm{Coker}\ (d_1 : L_2 \longrightarrow L_1)
\end{equation*}
et on désigne par $\bar d_o:L'_1\longrightarrow L_o$ le morphisme induit par $d_o:L_1\longrightarrow L_o$. On désigne par $E_o$ l'objet $(X_o,\alpha_o)$ de $\underline{\Psi}_{L.}(L_1')$, où $X_o = L_o + L_1'$ est l'extension triviale de $L_o$ par $L_1'$, et où $\alpha_o$ est le composé de $L_1\twoheadrightarrow L'_1$ avec $(\bar d_o,\mathrm{id}_{L_1'}):L'_1\longrightarrow L_o+L'_1$.

Pour voir que $E_o$ est bien quasi-maximal, notons que pour tout objet $(X,\alpha)$ de $\underline{\Psi}_{L.}$ , on peut trouver un morphisme de celui-ci dans un $(X',\alpha')$, induisant un \underline{monomorphisme} $A \longrightarrow A'$ dans $\underline{C}$, tel que l'extension $X'$ de $L_o$ par $A'$ soit \underline{scindée} : il suffit en effet de prendre $A'=X$, et $(X',\alpha') = i_*(X,\alpha)$, les notations étant celles de (3.1.2). D'autre part, il est immédiat qu'un objet $(X',\alpha')$ de $\underline{\Psi}$ tel que l'extension $X'$ soit scindée, est majoré par $E_o$.

}{F7I-VII-056}3.3.4. Soit $E$ un objet d'une catégorie cofibrée additive $\underline{\Psi}$ sur $\underline{C}$, au-dessus d'un objet $A$ de $\underline{C}$, et considérons le foncteur en l'objet variable $B$ de $\underline{C}$ :

\begin{equation}\tag{3.3.4.1}
B \longmapsto \mathrm{Hom}(E,\theta_B) \quad , 
\end{equation}
\SGApage{I}{189}

 où $\theta_B$ désigne l'objet trivial de $\underline{\Psi}(B)$. On dit que $E$ \underline{admet un complexe typique} si le foncteur précédent est représentable par un objet $\Omega_E$ de $C$. Comme le foncteur envisagé s'envoie canoniquement dans le foncteur $B \longmapsto \mathrm{Hom}(A,B)$ représenté par $A$, on conclut alors un homomorphisme canonique

\begin{equation}\tag{3.3.4.2}
A \longrightarrow \Omega_E
\end{equation}
et le complexe \underline{de chaînes} correspondant $A \longrightarrow \Omega_E$ de $\underline{C}$ de longueur 1 est appelé le \underline{complexe typique de} $E$.

\SGAvar{Dans le cas $\underline{\Psi} = \underline{\Psi}_{L.}$ , $E = (X,\alpha)$, on vérifie immédiatement que le foncteur (3.3.4.1) est représentable par $\mathrm{Coker}\ \alpha$, l'homomorphisme canonique (3.3.4.2) n'étant autre que le composé $A \xrightarrow{\ j\ } X \xrightarrow{\ \mathrm{can}\ } \mathrm{Coker}\ \alpha$ .

3.3.5. Lorsque une catégorie cofibrée additive $\underline{\Psi}$ sur $\underline{C}$ admet un objet quasi-maximal (3.3.3), et que les complexes typiques de ses objets existent, on prouve que le complexe typique $T.$ d'un objet quasi-maximal $E_o$ de $\underline{\Psi}$, considéré comme un objet de la \underline{catégorie dérivée} $D(\underline{C})$, ne dépend pas, à isomorphisme canonique près, du choix de $E_o$ (loc. cit. 4.4). On l'appelle le \underline{complexe typique de la catégorie cofibrée additive} $\underline{\Psi}$ . Lorsqu'on suppose enfin que $\underline{\Psi}$ est de plus exacte à gauche (3.3.2), on montre que $\underline{\Psi}$ se réconstitue, à équivalence de catégories fibrées sur $\underline{C}$ près (unique à isomorphisme unique près), en termes du complexe typique $T.$ (loc. cit. 6.10), et qu'on a des isomorphismes canoniques fonctoriels en l'objet $A$ de $\underline{C}$ (loc. cit. 5.4) :

}{Dans le cas $\underline{\Psi} = \underline{\Psi}_{L.}$ , $E = (X,\alpha)$, on vérifie immédiatement que le foncteur (3.3.4.1) est représentable par $\mathrm{Coker}\ \alpha$, l'homomorphisme canonique (3.3.4.2) n'étant autre que le composé $A \xrightarrow{\ i\ } X \xrightarrow{\ \mathrm{can}\ } \mathrm{Coker}\ \alpha$ .

3.3.5. Lorsque une catégorie cofibrée additive $\underline{\Psi}$ sur $\underline{C}$ admet un objet quasi-maximal (3.3.3), et que les complexes typiques de ses objets existent, on prouve que le complexe typique $T.$ d'un objet quasi-maximal $E_o$ de $\underline{\Psi}$, considéré comme un objet de la \underline{catégorie dérivée} $D(\underline{C})$, ne dépend pas, à isomorphisme canonique près, du choix de $E_o$ (loc. cit. 4.4). On l'appelle le \underline{complexe typique de la catégorie cofibrée additive} $\underline{\Psi}$ . Lorsqu'on suppose enfin que $\underline{\Psi}$ est de plus exacte à gauche (3.3.2), on montre que $\underline{\Psi}$ se reconstitue, à équivalence de catégories fibrées sur $\underline{C}$ près (unique à isomorphisme unique près), en termes du complexe typique $T.$ (loc. cit. 6.10), et qu'on a des isomorphismes canoniques fonctoriels en l'objet $A$ de $\underline{C}$ (loc. cit. 5.4) :

}{F7I-VII-057}\begin{equation}\tag{3.3.5.1}
\Psi^o(A) \xrightarrow{\ \sim\ } \mathrm{Ext}^o(T.,A) \ , \quad \Psi^1(A) \xrightarrow{\ \sim\ } \mathrm{Ext}^1(T.,A) \qquad .
\end{equation}

\SGApage{I}{190}

Pour achever de faire le lien avec loc. cit., nous allons prouver que dans le cas $\underline{\Psi} = \underline{\Psi}_{L.}$ , le complexe typique de $\underline{\Psi}_{L.}$ est le tronqué $[L.]$. Cela redonnera une définition des isomorphismes (3.1.8) et (3.2.4) en termes de (3.3.5.1), dont on vérifie qu'elle est compatible avec celle donnée dans 3.1 et dans 3.2.

\SGAvar{Pour prouver notre assertion, on peut évidemment se borner au cas où $L.$ est déjà tronqué à l'ordre 1 (car on a signalé dans 3.1 que le tronquage ne modifie pas, à isomorphisme de catégories cofibrées près, la catégorie $\underline{\Psi}_{L.}$ ). Mais alors il est évident que le complexe typique de l'objet quasi-maximal décrit dans 3.3.3) est canoniquement isomorphe à $L.$, compte tenu de la construction explicite des complexes typiques donnée dans 3.3.4.

}{Pour prouver notre assertion, on peut évidemment se borner au cas où $L.$ est déjà tronqué à l'ordre 1 (car on a signalé dans 3.1 que le tronquage ne modifie pas, à isomorphisme de catégories cofibrées près, la catégorie $\underline{\Psi}_{L.}$ ). Mais alors il est évident que le complexe typique de l'objet quasi-maximal décrit dans 3.3.3 est canoniquement isomorphe à $L.$, compte tenu de la construction explicite des complexes typiques donnée dans 3.3.4.

}{F7I-VII-058}\underline{Remarque} 3.3.6. Soit

\begin{equation*}
0 \longrightarrow A' \longrightarrow A \longrightarrow A'' \longrightarrow 0
\end{equation*}
une suite exacte dans $\underline{C}$, d'où une suite exacte pour les $\mathrm{Ext}^i$ :

\begin{equation*}
0 \longrightarrow \mathrm{Ext}^o(L.,A') \longrightarrow \mathrm{Ext}^o(L.,A) \rightarrow \ldots \rightarrow \mathrm{Ext}^1(L.,A) \longrightarrow \mathrm{Ext}^1(L.,A''),
\end{equation*}
(qu'on pourrait d'ailleurs continuer sur la droite en introduisant les $\mathrm{Ext}^i$ ($i \geq 2$)). Moyennant les isomorphismes (3.1.8) et (3.2.4) on en conclut une suite exacte

\SGAvar{\begin{equation}\tag{3.3.6.1}
0 \longrightarrow \Psi^o_{L.}(A') \longrightarrow \Psi^o_{L.}(A) \longrightarrow \Psi^o_{L.}(A'') \xrightarrow{\ o\ } \Psi^1_{L.}(A') \rightarrow \Psi^1_{L.}(A) \rightarrow \Psi^1_{L.}(A'').
\end{equation}
dont les flèches distinctes de la flèche notée $\partial$ sont définies simplement par fonctorialité de $\Psi^o_{L.}$ et de $\Psi^1_{L.}$ . D'autre part, dans loc. cit. (2.5.2), on définit une suite exacte (3.3.6.1), via un homomorphisme 


}{\begin{equation}\tag{3.3.6.1}
0 \longrightarrow \Psi^o_{L.}(A') \longrightarrow \Psi^o_{L.}(A) \longrightarrow \Psi^o_{L.}(A'') \xrightarrow{\ \partial\ } \Psi^1_{L.}(A') \rightarrow \Psi^1_{L.}(A) \rightarrow \Psi^1_{L.}(A'').
\end{equation}
dont les flèches distinctes de la flèche notée $\partial$ sont définies simplement par fonctorialité de $\Psi^o_{L.}$ et de $\Psi^1_{L.}$ . D'autre part, dans loc. cit. (2.5.2), on définit une suite exacte (3.3.6.1), via un homomorphisme 


}{F7I-VII-059}\SGApage{I}{191}

\begin{equation}\tag{3.3.6.2}
 \partial : \Psi^{o}(A'') \longrightarrow \Psi^{1}(A') \qquad ,
\end{equation}
pour toute catégorie cofibrée exacte à gauche $\underline{\Psi}$ sur $\underline{C}$, en associant à un élément $u$ du premier membre la classe dans le deuxième membre de l'élément de $\underline{\Psi}(A')$, déterminé à isomorphisme canonique près, dont l'image dans $\underline{\Psi}([A \longrightarrow A''])$ (notations de 3.3.2) est le couple $(X_o,u)$, où $X_o$ est l'objet trivial de $\underline{\Psi}(A)$. Ceci posé, je dis que dans le cas qui nous intéresse $\underline{\Psi} = \underline{\Psi}_{L.}$ , \underline{l'homomorphisme} (3.3.6.2) \underline{défini par voie} \underline{"géométrique"} \underline{n'est autre que l'homomorphisme bord} $\partial$ \underline{de} (3.3.6.1) \underline{changé de signe}, donc que cette dernière suite exacte n'est autre (avec les identifications faites et en changeant le signe du $\partial$) que celle de (loc. cit. (2.5.2)). Ceci est en effet contenu dans la dernière assertion de (loc. cit. 5.4), compte tenu des compatibilités signalées dans 3.3.5.

\SGAvar{\underline{Remarque} 3.3.6.3. Profitons de cette occasion pour préciser un point peu clair de [4, n$^{\circ}$ 3.11], lorsqu'on y définit les isomorphismes fondamentaux (3.11.3) et (3.11.4) (qui interviennent de façon essentielle dans la démonstration de (loc. cit. 5.4)). Le diagramme de loc. cit. p.35 est commutatif, et non anticommutatif, donc ne définit pas sans plus un homomorphisme $u : [L_1^X \longrightarrow L_o^X] \longrightarrow K^{\cdot}[1]$. Il convient de le rendre auticommutatif en y remplaçant $u^1$ par $-u^1$. (Il aurait été possible de suivre la convention opposée, en remplaçant $u^o$ par $-u^o$ ; mais c'est la convention indiquée ($u^1$ devenant $-u^1$) qui donne un isomorphisme (3.11.3) généralisant l'isomorphisme (3.5.3) de loc. cit., comme annoncé dans loc. cit.). Cette définition étant ainsi précisée, il convient de corriger l'assertion faite dans loc. cit. 5.4, les deux homomorphismes cobords qui
\SGApage{I}{192}
y sont envisagés \underline{différent par le signe.}

3.3.7. Dans (loc. cit. 6.13) on a précisé aussi la variance de la catégorie cofibrée $\underline{\Psi}_{L.}$ sur $\underline{C}$, pour $L.$ variable, en réconstituant notamment la catégorie des morphismes cocartésiens de $\underline{\Psi}_{L.}$ dans $\underline{\Psi}_{L'.}$ en termes des complexes de chaînes $L.$ et $L'.$ . On trouve en particulier que l'ensemble des classes à isomorphismes près de foncteurs cocartésiens $\underline{\Psi}_{L.} \longrightarrow \underline{\Psi}_{L'.}$ est en correspondance biunivoque canonique avec l'ensemble des isomorphismes $\mathrm{Hom}_{D(\underline{C})}(L'. ,L.)$. Nous n'utiliserons qu'une petite partie de ces renseignements, que nous développons dans la section qui suit, en utilisant la description simple de $\underline{\Psi}_{L.}$ donnée dans 3.1.


}{\underline{Remarque} 3.3.6.3. Profitons de cette occasion pour préciser un point peu clair de [4, n$^{\circ}$ 3.11], lorsqu'on y définit les isomorphismes fondamentaux (3.11.3) et (3.11.4) (qui interviennent de façon essentielle dans la démonstration de (loc. cit. 5.4)). Le diagramme de loc. cit. p.35 est commutatif, et non anticommutatif, donc ne définit pas sans plus un homomorphisme $u : [L_1^X \longrightarrow L_o^X] \longrightarrow K^{\cdot}[1]$. Il convient de le rendre anticommutatif en y remplaçant $u^1$ par $-u^1$. (Il aurait été possible de suivre la convention opposée, en remplaçant $u^o$ par $-u^o$ ; mais c'est la convention indiquée ($u^1$ devenant $-u^1$) qui donne un isomorphisme (3.11.3) généralisant l'isomorphisme (3.5.3) de loc. cit., comme annoncé dans loc. cit.). Cette définition étant ainsi précisée, il convient de corriger l'assertion faite dans loc. cit. 5.4, les deux homomorphismes cobords qui
\SGApage{I}{192}
y sont envisagés \underline{diffèrent par le signe.}

3.3.7. Dans (loc. cit. 6.13) on a précisé aussi la variance de la catégorie cofibrée $\underline{\Psi}_{L.}$ sur $\underline{C}$, pour $L.$ variable, en reconstituant notamment la catégorie des morphismes cocartésiens de $\underline{\Psi}_{L.}$ dans $\underline{\Psi}_{L'.}$ en termes des complexes de chaînes $L.$ et $L'.$ . On trouve en particulier que l'ensemble des classes à isomorphismes près de foncteurs cocartésiens $\underline{\Psi}_{L.} \longrightarrow \underline{\Psi}_{L'.}$ est en correspondance biunivoque canonique avec l'ensemble des homomorphismes $\mathrm{Hom}_{D(\underline{C})}(L'. ,L.)$. Nous n'utiliserons qu'une petite partie de ces renseignements, que nous développons dans la section qui suit, en utilisant la description simple de $\underline{\Psi}_{L.}$ donnée dans 3.1.


}{F7I-VII-060}3.4. \underline{Variance de} $\underline{\Psi}_{L.}$ \underline{avec} $L.$ .

Considérons donc un morphisme de complexes

\begin{equation}\tag{3.4.1}
f : L'. \longrightarrow L. \qquad ,
\end{equation}
donné par un diagramme commutatif

\begin{equation}\tag{3.4.2}
\begin{tikzcd}[column sep=large, row sep=large]
{} \arrow[r] \& L'_2 \arrow[r] \arrow[d, "f_2"'] \& L'_1 \arrow[r] \arrow[d, "f_1"'] \& L'_o \arrow[r] \arrow[d, "f_o"'] \& 0 \\
{} \arrow[r] \& L_2 \arrow[r] \& L_1 \arrow[r] \& L_o \arrow[r] \& 0 \& .
\end{tikzcd}
\end{equation}


















On va lui associer un foncteur canonique

\SGAvar{\begin{equation}\tag{3.4.3}
f^* : \underline{\Psi}_{L.} \longrightarrow \underline{\Psi}_{L'.} \qquad ,
\end{equation}
qui sera en fait un foncteur cocartésien de catégories cofibrées sur $\underline{C}$ (i.e. il est compatible avec les foncteurs structuraux à valeurs dans $\underline{C}$, et transforme morphismes cocartésiens en morphismes cocartésiens, cf. SGA 1 VI). Si $(X,\alpha)$ est un objet de $\underline{\Psi}_L$ . on désigne par $f^*(X,\alpha)$ l'objet
\SGApage{I}{193}
$(X',\alpha')$, où $X'$ est l'extension de $L'_o$ par $A$ image inverse de l'extension donnée $X$ de $L_o$ par $A$ à l'aide de $f_o$, et où $\alpha'$ est la trivialisation partielle déduite de la trivialisation partielle $\alpha$ à l'aide de la transitivité dans le deuxième carré de (3.4.2). (La condition de compatibilité de $\alpha'$ à la relation $d'_o d'_1 = 0$ résulte évidemment de la condition analogue pour $\alpha$.) Il est trivial par construction que le foncteur (3.4.3) commute aux foncteurs structuraux à valeurs dans $\underline{C}$. Le fait qu'il soit cocartésien résulte aussitôt du fait que dans la catégorie des extensions d'objets de $\underline{C}$, les opérations ``image directe d'extensions'' commutent aux opérations ``image inverse d'extensions''.

}{\begin{equation}\tag{3.4.3}
f^* : \underline{\Psi}_{L.} \longrightarrow \underline{\Psi}_{L'.} \qquad ,
\end{equation}
qui sera en fait un foncteur cocartésien de catégories cofibrées sur $\underline{C}$ (i.e. il est compatible avec les foncteurs structuraux à valeurs dans $\underline{C}$, et transforme morphismes cocartésiens en morphismes cocartésiens, cf. SGA 1 VI). Si $(X,\alpha)$ est un objet de $\underline{\Psi}_{L.}$, on désigne par $f^*(X,\alpha)$ l'objet
\SGApage{I}{193}
$(X',\alpha')$, où $X'$ est l'extension de $L'_o$ par $A$ image inverse de l'extension donnée $X$ de $L_o$ par $A$ à l'aide de $f_o$, et où $\alpha'$ est la trivialisation partielle déduite de la trivialisation partielle $\alpha$ à l'aide de la transitivité dans le deuxième carré de (3.4.2). (La condition de compatibilité de $\alpha'$ à la relation $d'_o d'_1 = 0$ résulte évidemment de la condition analogue pour $\alpha$.) Il est trivial par construction que le foncteur (3.4.3) commute aux foncteurs structuraux à valeurs dans $\underline{C}$. Le fait qu'il soit cocartésien résulte aussitôt du fait que dans la catégorie des extensions d'objets de $\underline{C}$, les opérations ``image directe d'extensions'' commutent aux opérations ``image inverse d'extensions''.

}{F7I-VII-061}Notons que le foncteur (3.4.3) définit des homomorphismes de foncteurs

\begin{equation}\tag{3.4.4}
f^i : \Psi^i_{L.} \longrightarrow \Psi^i_{L'.} \qquad , \qquad i=0,1
\end{equation}
(comparer 2.6). D'autre part, l'homomorphisme (3.4.1) définit des homomorphismes

\SGAvar{\begin{equation}\tag{3.4.5}
f^i : \mathrm{Ext}^i(L.,-) \qquad \mathrm{Ext}^i(L'.,-) \quad , \quad i \in \mathbb{Z} \quad .
\end{equation}
Ceci posé, je dis que pour $i=0,1$, les homomorphismes (3.4.4) et (3.4.5) sont compatibles avec les isomorphismes de la forme (3.1.8) et (3.2.4), i.e. qu'on a commutativité dans les diagrammes suivants ($A \in \mathrm{ob}\ \underline{C}$, $i=0,1$) :

}{\begin{equation}\tag{3.4.5}
f^i : \mathrm{Ext}^i(L.,-) \longrightarrow \mathrm{Ext}^i(L'.,-) \quad , \quad i \in \mathbb{Z} \quad .
\end{equation}
Ceci posé, je dis que pour $i=0,1$, les homomorphismes (3.4.4) et (3.4.5) sont compatibles avec les isomorphismes de la forme (3.1.8) et (3.2.4), i.e. qu'on a commutativité dans les diagrammes suivants ($A \in \mathrm{ob}\ \underline{C}$, $i=0,1$) :

}{F7I-VII-062}\begin{equation}\tag{3.4.6}
\begin{tikzcd}
\underline{\Psi}^i_{L.}(A) \arrow[r, "\sim"] \arrow[d] \& \mathrm{Ext}^i(L.,A) \arrow[d] \\
\underline{\Psi}^i_{L'.}(A) \arrow[r, "\sim"] \& \mathrm{Ext}^i(L'.,A) \& .
\end{tikzcd}
\end{equation}













\SGApage{I}{194}

\SGAvar{La vérification de cette compatibilité est essentiellement triviale, et laissé au lecteur.

}{La vérification de cette compatibilité est essentiellement triviale, et laissée au lecteur.

}{F7I-VII-063}\underline{Proposition} 3.4.7. \underline{Pour que le foncteur} (3.4.3) \underline{soit fidèle} (\underline{resp.} \underline{pleinement fidèle}, \underline{resp.} \underline{une équivalence de catégories}), \underline{il faut et il suffit} \underline{que} $H_o(f) : H_o(L'.) \longrightarrow H_o(L.)$ \underline{soit un épimorphisme} (\underline{resp.} \underline{que} $H_o(f)$ \underline{soit un isomorphisme et} $H_1(f)$ \underline{soit un épimorphisme}, \underline{resp.} \underline{que} $H_o(f)$ \underline{et} $H_1(f)$ \underline{soient des isomorphismes}).

\SGAvar{En effet, il faut exprimer que, $f^i$ étant l'homomorphisme fonctoriel de (3.4.4) ($i=0,1$), $f^o$ est un monomorphisme (resp. $f^o$ est un isomorphisme et $f$ un monomorphisme, resp. $f^o$ et $f^1$ sont des isomorphismes). En vertu de la commutativité de (3.4.6), ces conditions s'expriment en termes des conditions analogues sur les foncteurs $\mathrm{Ext}^i$, $i=0,1$. Dans ce cas, l'assertion n'est qu'un énoncé assez trivial d'algèbre homologique, que nous laissons au lecteur à titre d'exercice.

}{En effet, il faut exprimer que, $f^i$ étant l'homomorphisme fonctoriel de (3.4.4) ($i=0,1$), $f^o$ est un monomorphisme (resp. $f^o$ est un isomorphisme et $f^1$ un monomorphisme, resp. $f^o$ et $f^1$ sont des isomorphismes). En vertu de la commutativité de (3.4.6), ces conditions s'expriment en termes des conditions analogues sur les foncteurs $\mathrm{Ext}^i$, $i=0,1$. Dans ce cas, l'assertion n'est qu'un énoncé assez trivial d'algèbre homologique, que nous laissons au lecteur à titre d'exercice.

}{F7I-VII-064}3.4.8. Considérons maintenant une suite exacte de complexes de chaînes

\begin{equation}\tag{3.4.8.1}
0 \longrightarrow L. \xrightarrow{f} L'. \xrightarrow{g} L''. \longrightarrow 0 \qquad ,
\end{equation}
d'où des foncteurs du type (3.4.3)

\begin{center}
\begin{tikzcd}
\underline{\Psi}_{L''.} \arrow[r, "g^*"] \& \Psi_{L'.} \arrow[r, "f^*"] \& \underline{\Psi}_{L.} \& .
\end{tikzcd}
\end{center}









La propriété de transitivité qu'on devine sur les foncteurs du type (3.4.3) fournit ici un isomorphisme de foncteurs

\begin{equation*}
f^* \, g^* \quad \simeq \quad (gf)^* \quad = \quad (0_{L. \,\to\, L''.})^* \qquad , \qquad 
\end{equation*}

\SGApage{I}{195}

 et comme le dernier membre est évidemment canoniquement isomorphe au foncteur cocartésien ``nul'' de $\underline{\Psi}_{L''.}$ dans $\underline{\Psi}_{L.}$ , l'isomorphisme précédent permet de définir un foncteur canonique

\SGAvar{\begin{equation}\tag{3.4.8.2}
\underline{\Psi}_{L''.} \longrightarrow \underline{\Psi}_{[L. \longrightarrow L'.]} \quad ,
\end{equation}
où $\underline{\Psi}_{[L. \longrightarrow L'.]}$ désigne la catégorie des couples $(X',t)$, où $X'$ est un objet de $\underline{\Psi}_{L'.}$ et $t$ une ``trivialisation'' de son image inverse $f^*(X')$ i.e. un isomorphisme $f^*(X') \simeq 0_{L.,A'}$, $0_{L.,A}$ désignant l'objet nul de $\underline{\Psi}_{L.}(A)$, et $A$ étant l'objet de $C$ sous $X'$. Ceci posé, je dis que \underline{le foncteur} (3.4.8.2) \underline{est une équivalence de catégories} (ce qui complète, par une propriété d'exactitude en $L.$ de l'expression $\underline{\Psi}_{L.}(A)$, la propriété d'exactitude à gauche en $A$ établie dans 3.3.1).

Esquissons la démonstration de l'assertion précédente. La donnée d'un objet du deuxième membre de (3.4.8.2) revient à la donnée d'une extension $0 \longrightarrow A \stackrel{i'}{\longrightarrow} X' \stackrel{j'}{\longrightarrow} L'_o \longrightarrow 0$ de $L'_o$ par un objet $A$ de $C$, plus un morphisme $\alpha' : L'_1 \longrightarrow X'$ satisfaisant $j\alpha' = d'_o$ et $\alpha'd'_1 = 0$, plus un morphisme $t : L_o \longrightarrow X'$ satisfaisant $jt = f_o$ et $td_o = \alpha'f_1$, cf. diagramme ci-dessous :

\begin{center}
\begin{tikzcd}[column sep=large, row sep=large]
\& \& \& \& \& 0 \arrow[dl] \& \& 0 \arrow[dl, dashed] \\
\& \& \& \& A \arrow[dl, "i'"] \& \& A \arrow[dl, dashed, "i''"] \& \\
\& \& \& X' \arrow[dl, leftrightarrow, "j'"'] \& \& X'' \arrow[dll, dashed, "j''"] \& \& \\
0 \arrow[r] \& L_o \arrow[r, "f_o"] \arrow[rru, bend left, "t"] \& L'_o \arrow[r, "g_o"] \& L''_o \arrow[r] \& 0 \& \& \& \\
0 \arrow[r] \& L_1 \arrow[r, "f_1"] \arrow[u, "d_o"] \& L'_1 \arrow[r, "g_1"] \arrow[u, "d'_o"] \arrow[ruu, "\alpha'"] \& L''_1 \arrow[r] \arrow[u, "d''_o"] \arrow[rruu, dashed, bend right, "\alpha''"] \& 0 \& \& \& \\
0 \arrow[r] \& L_2 \arrow[r] \arrow[u] \& L'_2 \arrow[r] \arrow[u, "d'_1"] \& L''_2 \arrow[r] \arrow[u, "d''_1"] \& 0 \& \& \&
\end{tikzcd}
\end{center}

































}{\begin{equation}\tag{3.4.8.2}
\underline{\Psi}_{L''.} \longrightarrow \underline{\Psi}_{[L. \longrightarrow L'.]} \quad ,
\end{equation}
où $\underline{\Psi}_{[L. \longrightarrow L'.]}$ désigne la catégorie des couples $(X',t)$, où $X'$ est un objet de $\underline{\Psi}_{L'.}$ et $t$ une ``trivialisation'' de son image inverse $f^*(X')$ i.e. un isomorphisme $f^*(X') \simeq 0_{L.,A}$, $0_{L.,A}$ désignant l'objet nul de $\underline{\Psi}_{L.}(A)$, et $A$ étant l'objet de $C$ sous $X'$. Ceci posé, je dis que \underline{le foncteur} (3.4.8.2) \underline{est une équivalence de catégories} (ce qui complète, par une propriété d'exactitude en $L.$ de l'expression $\underline{\Psi}_{L.}(A)$, la propriété d'exactitude à gauche en $A$ établie dans 3.3.1).

Esquissons la démonstration de l'assertion précédente. La donnée d'un objet du deuxième membre de (3.4.8.2) revient à la donnée d'une extension $0 \longrightarrow A \stackrel{i'}{\longrightarrow} X' \stackrel{j'}{\longrightarrow} L'_o \longrightarrow 0$ de $L'_o$ par un objet $A$ de $C$, plus un morphisme $\alpha' : L'_1 \longrightarrow X'$ satisfaisant $j'\alpha' = d'_o$ et $\alpha'd'_1 = 0$, plus un morphisme $t : L_o \longrightarrow X'$ satisfaisant $j't = f_o$ et $td_o = \alpha'f_1$, cf. diagramme ci-dessous :

\begin{center}
\begin{tikzcd}[column sep=large, row sep=large]
\& \& \& \& \& 0 \arrow[dl] \& \& 0 \arrow[dl, dashed] \\
\& \& \& \& A \arrow[dl, "i'"] \& \& A \arrow[dl, dashed, "i''"] \& \\
\& \& \& X' \arrow[dl, "j'"'] \& \& X'' \arrow[dll, dashed, "j''"] \& \& \\
0 \arrow[r] \& L_o \arrow[r, "f_o"] \arrow[rru, bend left, "t"] \& L'_o \arrow[r, "g_o"] \& L''_o \arrow[r] \& 0 \& \& \& \\
0 \arrow[r] \& L_1 \arrow[r, "f_1"] \arrow[u, "d_o"] \& L'_1 \arrow[r, "g_1"] \arrow[u, "d'_o"] \arrow[ruu, "\alpha'"] \& L''_1 \arrow[r] \arrow[u, "d''_o"] \arrow[rruu, dashed, bend right, "\alpha''"] \& 0 \& \& \& \\
0 \arrow[r] \& L_2 \arrow[r] \arrow[u] \& L'_2 \arrow[r] \arrow[u, "d'_1"] \& L''_2 \arrow[r] \arrow[u, "d''_1"] \& 0 \& \& \&
\end{tikzcd}
\end{center}

































}{F7I-VII-065}\begin{equation*}
\left\{ \begin{array}{rcl} j'\alpha' & = & d'_o \\ \alpha'd'_1 & = & 0 \\ j't & = & f_o \\ td_o & = & \alpha'f_1 \end{array} \right.
\end{equation*}
\SGApage{I}{196}

\SGAvar{Une propriété bien connue d'exactitude à gauche en le premier argument des extensions d'un objet $L$ par un autre $A$, appliquée à la première ligne du diagramme, montre que la donnée de $t$ satisfaisant $jt = f_o$ permet d'interpréter l'extension $X'$ comme image inverse d'une extension $X''$ de $L''_o$ par $A$, de telle façon que $t$ corresponde à la trivialisation naturelle de l'image inverse de cette extension par $g_o f_o = 0$. La donnée de $\alpha'$ satisfaisant les relations $j\alpha' = d'_o$ et $\alpha'f_1 = td_o$ peut s'interpréter par essentiellement le même dictionnaire, mais appliqué à la seconde ligne du diagramme, comme provenant d'un morphisme $\alpha'' : L''_1 \longrightarrow X''$ satisfaisant $j''\alpha'' = d''_o$ (en interprétant $\alpha'$ et $\alpha''$ comme des trivialisations des images inverses des extensions $X'$ resp. $X''$ par $d'_o$ resp. $d''_o$). Enfin la relation $\alpha'd'_1 = 0$ équivaut alors à la relation $\alpha''d''_1 = 0$. Bref, la donnée de $(X',\alpha',t)$, objet de $\underline{\Psi}_{[L. \longrightarrow L'.]}$, est équivalente à celle de $(X'',\alpha'')$, objet de $\underline{\Psi}_{L.}$ , ce qui prouve (essentiellement) l'assertion que (3.4.8.2) est bien une équivalence de catégories.

}{Une propriété bien connue d'exactitude à gauche en le premier argument des extensions d'un objet $L$ par un autre $A$, appliquée à la première ligne du diagramme, montre que la donnée de $t$ satisfaisant $j't = f_o$ permet d'interpréter l'extension $X'$ comme image inverse d'une extension $X''$ de $L''_o$ par $A$, de telle façon que $t$ corresponde à la trivialisation naturelle de l'image inverse de cette extension par $g_o f_o = 0$. La donnée de $\alpha'$ satisfaisant les relations $j'\alpha' = d'_o$ et $\alpha'f_1 = td_o$ peut s'interpréter par essentiellement le même dictionnaire, mais appliqué à la seconde ligne du diagramme, comme provenant d'un morphisme $\alpha'' : L''_1 \longrightarrow X''$ satisfaisant $j''\alpha'' = d''_o$ (en interprétant $\alpha'$ et $\alpha''$ comme des trivialisations des images inverses des extensions $X'$ resp. $X''$ par $d'_o$ resp. $d''_o$). Enfin la relation $\alpha'd'_1 = 0$ équivaut alors à la relation $\alpha''d''_1 = 0$. Bref, la donnée de $(X',\alpha',t)$, objet de $\underline{\Psi}_{[L. \longrightarrow L'.]}$, est équivalente à celle de $(X'',\alpha'')$, objet de $\underline{\Psi}_{L''.}$ , ce qui prouve (essentiellement) l'assertion que (3.4.8.2) est bien une équivalence de catégories.

}{F7I-VII-066}3.4.9. Considérons toujours la suite exacte (3.4.8.1), et considérons la suite exacte des $\mathrm{Ext}^i$ qu'elle définit :

\SGAvar{{\footnotesize \begin{equation}\tag{3.4.9.1}
0 \longrightarrow \mathrm{Ext}^o(L''.,A) \stackrel{g^o}{\longrightarrow} \mathrm{Ext}^o(L'.A) \stackrel{f^o}{\longrightarrow} \mathrm{Ext}^o(L.,A) \stackrel{\partial^o}{\longrightarrow} \mathrm{Ext}^1(L''.,A) \stackrel{g^1}{\longrightarrow} \mathrm{Ext}^1(L'.,A) \stackrel{f^1}{\longrightarrow} \mathrm{Ext}^1(L.,A) \longrightarrow \quad .
\end{equation}}
Les isomorphismes canoniques du type (3.1.8) et (3.2.4) nous permettent d'interpréter ``géométriquement'' les six premiers termes de cette suite exacte en termes des catégories $\underline{\Psi}_{L.}$, $\underline{\Psi}_{L'.}$ et $\underline{\Psi}_{L''.}$ , et la compatibilité (3.4.6) nous permet d'interpréter de même géométriquement les flèches
\SGApage{I}{197}
$f^o,g^o,f^1,g^1$ de ce diagramme. Nous allons donner de même une interprétation géométrique de la flèche $\partial^o$ de ce diagramme, en utilisant le résultat qu'on vient de prouver, savoir que (3.4.8.2) est une équivalence de catégories. Soit en effet

\begin{equation*}
u \in \mathrm{Ext}^o(L.,A), \qquad \mathrm{i.e.} \ u : L_o \longrightarrow A, \ \mathrm{avec} \ ud_o = 0 \ .
\end{equation*}
Considérons l'objet $(0,\overline{u})$ de $\underline{\Psi}_{[L. \to L']}$, où $0$ est l'objet nul de $\underline{\Psi}_{L'.}$ , dont l'image inverse par $f_o$ s'identifie donc à l'objet nul de $\underline{\Psi}_{L.}$ , et où $\overline{u}$ est l'automorphisme de ce dernier objet défini par $u$, grâce à l'isomorphisme canonique (3.1.8). Grâce au fait que (3.4.8.2) est une équivalence, l'objet précédent $(0,\overline{u})$ provient d'un objet de $\underline{\Psi}_{L''.}(A)$, défini à isomorphisme unique près. La classe à isomorphisme près de cet élément est un élément de $\Psi^1_{L''.}(A) \simeq \mathrm{Ext}^1(L''.,A)$. \underline{Je dis que cet élément n'est autre que} $\partial^o(u)$.

Nous allons prouver cette compatibilité en la ramenant à la compatibilité analogue, énoncée dans [4, cor. 5.5] (où il convient de rectifier le signe comme indiqué dans 3.3.6.3) ; notre argument donnera en même temps une autre démonstration de l'équivalence (3.4.8.2), en la ramenant à l'énoncé analogue [4, prop. 2.7]. Pour ceci, considérons la catégorie opposée $\underline{C}^o$ de $\underline{C}$, et interprétons les constructions de $\underline{\Psi}_{L.}(A)$, pour $A$ fixé, et celles de 3.4.8 et de la présente section 3.4.9, en termes de la catégorie $\underline{C}^o$, et de la catégorie cofibrée $\underline{E}$ sur $C^o$ des extensions de l'objet fixé $A$ de $\underline{C}^o$ par un objet variable $L$ de $\underline{C}^o$. C'est une catégorie cofibrée additive sur $\underline{C}^o$, et même exacte à gauche sur $\underline{C}^o$ (au sens de loc. cit. 2.2). Un complexe de \underline{chaînes} $L.$ de $\underline{C}$ peut-être interprété  comme
\SGApage{I}{198}
un complexe de \underline{cochaînes} de $\underline{C}^o$, et avec cette interprétation, la construction donnée de $\underline{\Psi}_{L.}(A)$ coïncide avec celle de la catégorie opposée à $\underline{\hat{E}}(L.)$ au sens de loc. cit. 1.8. Le foncteur (3.4.8.2) s'identifie alors à l'opposé du foncteur du type de loc. cit. (2.1.2), où $E$ est remplacé par son extension $\underline{\hat{E}}$ à la catégorie des complexes de cochaînes, et le fait que le foncteur envisagé soit une équivalence est donc un cas particulier de loc. cit. 2.7 a) $\Longrightarrow$ d). D'autre part, la suite exacte (3.4.9.1) s'identifie à la suite exacte des $\mathrm{Ext}^i(A,L''.)$, $\mathrm{Ext}^i(A,L'.)$ et $\mathrm{Ext}^i(A,L.)$, modulo un signe $(-1)^i$ pour le $\partial^i$ (Verdier dixit). De même, les isomorphismes (3.1.8) et (3.2.4) s'identifient aux isomorphismes analogues (3.11.5) et (3.11.4). Donc la compatibilité énoncée est un cas particulier de loc. cit. 5.5 (rectifié comme indiqué plus haut).

}{{\footnotesize \begin{equation}\tag{3.4.9.1}
0 \longrightarrow \mathrm{Ext}^o(L''.,A) \stackrel{g^o}{\longrightarrow} \mathrm{Ext}^o(L'.,A) \stackrel{f^o}{\longrightarrow} \mathrm{Ext}^o(L.,A) \stackrel{\partial^o}{\longrightarrow} \mathrm{Ext}^1(L''.,A) \stackrel{g^1}{\longrightarrow} \mathrm{Ext}^1(L'.,A) \stackrel{f^1}{\longrightarrow} \mathrm{Ext}^1(L.,A) \longrightarrow \quad .
\end{equation}}
Les isomorphismes canoniques du type (3.1.8) et (3.2.4) nous permettent d'interpréter ``géométriquement'' les six premiers termes de cette suite exacte en termes des catégories $\underline{\Psi}_{L.}$, $\underline{\Psi}_{L'.}$ et $\underline{\Psi}_{L''.}$ , et la compatibilité (3.4.6) nous permet d'interpréter de même géométriquement les flèches
\SGApage{I}{197}
$f^o,g^o,f^1,g^1$ de ce diagramme. Nous allons donner de même une interprétation géométrique de la flèche $\partial^o$ de ce diagramme, en utilisant le résultat qu'on vient de prouver, savoir que (3.4.8.2) est une équivalence de catégories. Soit en effet

\begin{equation*}
u \in \mathrm{Ext}^o(L.,A), \qquad \mathrm{i.e.} \ u : L_o \longrightarrow A, \ \mathrm{avec} \ ud_o = 0 \ .
\end{equation*}
Considérons l'objet $(0,\overline{u})$ de $\underline{\Psi}_{[L. \to L'.]}$, où $0$ est l'objet nul de $\underline{\Psi}_{L'.}$ , dont l'image inverse par $f_o$ s'identifie donc à l'objet nul de $\underline{\Psi}_{L.}$ , et où $\overline{u}$ est l'automorphisme de ce dernier objet défini par $u$, grâce à l'isomorphisme canonique (3.1.8). Grâce au fait que (3.4.8.2) est une équivalence, l'objet précédent $(0,\overline{u})$ provient d'un objet de $\underline{\Psi}_{L''.}(A)$, défini à isomorphisme unique près. La classe à isomorphisme près de cet élément est un élément de $\Psi^1_{L''.}(A) \simeq \mathrm{Ext}^1(L''.,A)$. \underline{Je dis que cet élément n'est autre que} $\partial^o(u)$.

Nous allons prouver cette compatibilité en la ramenant à la compatibilité analogue, énoncée dans [4, cor. 5.5] (où il convient de rectifier le signe comme indiqué dans 3.3.6.3) ; notre argument donnera en même temps une autre démonstration de l'équivalence (3.4.8.2), en la ramenant à l'énoncé analogue [4, prop. 2.7]. Pour ceci, considérons la catégorie opposée $\underline{C}^o$ de $\underline{C}$, et interprétons les constructions de $\underline{\Psi}_{L.}(A)$, pour $A$ fixé, et celles de 3.4.8 et de la présente section 3.4.9, en termes de la catégorie $\underline{C}^o$, et de la catégorie cofibrée $\underline{E}$ sur $C^o$ des extensions de l'objet fixé $A$ de $\underline{C}^o$ par un objet variable $L$ de $\underline{C}^o$. C'est une catégorie cofibrée additive sur $\underline{C}^o$, et même exacte à gauche sur $\underline{C}^o$ (au sens de loc. cit. 2.2). Un complexe de \underline{chaînes} $L.$ de $\underline{C}$ peut être interprété  comme
\SGApage{I}{198}
un complexe de \underline{cochaînes} de $\underline{C}^o$, et avec cette interprétation, la construction donnée de $\underline{\Psi}_{L.}(A)$ coïncide avec celle de la catégorie opposée à $\underline{\hat{E}}(L.)$ au sens de loc. cit. 1.8. Le foncteur (3.4.8.2) s'identifie alors à l'opposé du foncteur du type de loc. cit. (2.1.2), où $E$ est remplacé par son extension $\underline{\hat{E}}$ à la catégorie des complexes de cochaînes, et le fait que le foncteur envisagé soit une équivalence est donc un cas particulier de loc. cit. 2.7 a) $\Longrightarrow$ d). D'autre part, la suite exacte (3.4.9.1) s'identifie à la suite exacte des $\mathrm{Ext}^i(A,L''.)$, $\mathrm{Ext}^i(A,L'.)$ et $\mathrm{Ext}^i(A,L.)$, modulo un signe $(-1)^i$ pour le $\partial^i$ (Verdier dixit). De même, les isomorphismes (3.1.8) et (3.2.4) s'identifient aux isomorphismes analogues (3.11.5) et (3.11.4). Donc la compatibilité énoncée est un cas particulier de loc. cit. 5.5 (rectifié comme indiqué plus haut).

}{F7I-VII-067}\underline{Remarque} 3.4.9.2. Le rédacteur, qui a très peur des signes, avoue qu'il n'a pas entièrement confiance dans la démonstration qui précède, craignant d'avoir oublié des signes dans le passage de $C$ à $C^o$. Il recommande au lecteur courageux de faire une vérification directe de la compatibilité annoncée (resp. de la compatibilité opposée). Notons que ces questions de signes ne joueront (heureusement) aucun rôle dans les exposés suivants.


3.5. \underline{Une résolution partielle plate canonique pour un $\mathbb{Z}$-Module quelconque}

Nous allons appliquer les résultats précédents au cas où $\underline{C}$ est la catégorie des Groupes commutatifs du topos $\underline{T}$ i.e. des $\mathbb{Z}$-Modules de $\underline{T}$. Soit $P$ un objet de $\underline{C}$, i.e. un groupe commutatif du topos $\underline{T}$. Nous allons définir un complexe de chaînes $L.(P)$ augmenté vers $P$, dépendant fonctoriellement de $P$ :


\SGApage{I}{199}

\begin{equation}\tag{3.5.1}
\begin{tikzcd}[column sep=large, row sep=large]
0 \arrow[r] \& \overbrace{\mathbb{Z}[P\times P]\times\mathbb{Z}[P\times P\times P]}^{L_2} \arrow[r, "d_1"] \& \overbrace{\mathbb{Z}[P\times P]}^{L_1} \arrow[r, "d_o"] \& \overbrace{\mathbb{Z}[P]}^{L_o} \arrow[r] \arrow[d, "\varepsilon"] \& 0 \\
\& \& \& P \&
\end{tikzcd}\quad ,
\end{equation}













 satisfaisant $L_i(P)=0$ pour $i \geq 3$. Les composantes de ce complexe sont explicitées dans (3.5.1) ci-dessus, en utilisant les notations de 1.4. Pour expliciter les opérateurs différentiels $d_o,d_1$ et l'augmentation $\varepsilon$, on utilise la notation $[p]$ pour désigner le point de $\mathbb{Z}[P](S)$ défini par le point $p \in P(S)$ de $P$, et de même pour les notations $[p,q]$, $[p,q,r]$, lorsque $p,q,r \in P(S)$, pour désigner des éléments de $\mathbb{Z}[P\times P](S)$ resp. de $\mathbb{Z}[P\times P\times P](S)$. Avec ces notations, on définira $\varepsilon$, $d_o$ et $d_1$ par les formules

\begin{equation}\tag{3.5.2}
\left\{
\begin{array}{l}
\varepsilon[p] = p \\[1.5ex]
d_o[p,q]= [p+q] - [p] - [q] \\[1.5ex]
d_1[p,q]= [p,q] - [q,p] \\[1.5ex]
d_1[p,q,r] = [p+q,r] - [p,q+r] + [p,q] - [q,r] \qquad .
\end{array}
\right.
\end{equation}
On vérifie facilement que les homomorphismes de Groupes commutatifs ainsi définis définissent effectivement un complexe $L.(P)$ augmenté vers $P$, i.e. qu'on a les relations

\SGAvar{\begin{equation*}
\varepsilon d_o = 0 \quad , \quad d_1 d_o = 0 \qquad .
\end{equation*}
En fait, la première relation est triviale, et la deuxième se décompose en deux, suivant les deux facteurs de $L_2(P)$, les deux relations en questions exprimant respectivement la \underline{commutativité} et l'\underline{associativité} de la
\SGApage{I}{200}
loi de composition donnée sur $P$.

}{\begin{equation*}
\varepsilon d_o = 0 \quad , \quad d_o d_1 = 0 \qquad .
\end{equation*}
En fait, la première relation est triviale, et la deuxième se décompose en deux, suivant les deux facteurs de $L_2(P)$, les deux relations en question exprimant respectivement la \underline{commutativité} et l'\underline{associativité} de la
\SGApage{I}{200}
loi de composition donnée sur $P$.

}{F7I-VII-068}Il est immédiat que le complexe augmenté ainsi défini dépend fonctoriellement de $P$, un morphisme $u : P \longrightarrow P'$ donnant naissance à un diagramme commutatif

\begin{center}
\begin{tikzcd}[column sep=huge, row sep=huge]
L.(P) \arrow[r, "L.(u)"] \arrow[d] \& L.(P') \arrow[d] \\
P \arrow[r, "u"] \& P'
\end{tikzcd}\quad .
\end{center}

Notons aussi que les composantes du complexe $L.$ sont \underline{plates}, comme tout $\mathbb{Z}$-Module libre.

\underline{Proposition} 3.5.3. \underline{On a} $H_1(L.(P))=0$, \underline{et l'augmentation} $\varepsilon$ \underline{induit un isomorphisme} $H_o(L.(P)) \simeq P$.

Pour démontrer 3.5.3, on peut se réduire par le procédé standard au cas où \underline{$T$} est la catégorie des ensembles, et où $P$ est donc un $\mathbb{Z}$-Module ordinaire, où (3.5.1) a des chances d'être ``bien connu''. Nous allons donner une démonstration directe, n'utilisant pas cette réduction, en appliquant 3.4.7 au morphisme d'augmentation

\begin{equation*}
\varepsilon : L.(P) \longrightarrow P \qquad ,
\end{equation*}
où $P$ est comme d'habitude considéré comme un complexe de chaînes réduit à la dimension zéro. On est donc ramené à prouver que pour tout objet $G$ de \underline{$C$}, le foncteur

\begin{equation}\tag{*}
\varepsilon^* : \underline{\Psi}_P(G) \longrightarrow \underline{\Psi}_{L.(P)}(G)
\end{equation}
est une équivalence de catégories. Or le premier membre de (*) est évidem
\SGApage{I}{201}
ment isomorphe à la catégorie des extensions de $\mathbb{Z}$ -Modules de $P$ par $G$. Pour expliciter le deuxième membre, on interprète les catégories d'extensions de $L_o$, $L_1$, $L_2$ par $G$, comme la catégorie des torseurs sous $G_P$, resp. la catégorie des torseurs sous $G_{P \times P}$, resp. la catégorie des couples d'un torseur sous $G_{P \times P}$ et d'un torseur sous $G_{P \times P \times P}$, grâce à 1.4. Donc un objet du second membre de (*) s'identifie à un couple $(E,\varphi)$, où $E$ est un torseur sous $G_P$ et $\varphi$ une trivialisation du torseur $E'$ sur $P \times P$ qui correspond à l'image inverse par $d_o$ de l'extension de $\mathbb{Z}(P)$ par $G$ définie par $E$, - $\varphi$ étant soumise à une condition de compatibilité avec $d_o d_1 = 0$ à préciser. Or la définition de $d_o$ (3.5.2) donne aussitôt

\begin{equation*}
E' = \pi^*(E)\mathrm{pr}_1^*(E)^{-1}\mathrm{pr}_2^*(E)^{-1} \qquad ,
\end{equation*}
et une trivialisation de ce torseur peut s'interpréter comme un isomorphisme (également noté $\varphi$)

\SGAvar{\begin{equation*}
\pi^*(E) \xleftarrow{\sim} \mathrm{pr}_1^*(E)\mathrm{pr}_2^*(E) \qquad .
\end{equation*}
D'autre part, en explicitant la condition de compatibilité de $\varphi$ avec $d_1 d_o = 0$, on trouve qu'elle s'exprime par la \underline{condition de commutativité} de (1.2.1), et par la \underline{condition d'associativité} de (1.1.4.1), ces deux conditions correspondant respectivement aux deux facteurs de $L_2$ . Le fait que (*) est une équivalence de catégories n'est alors autre que le dictionnaire de 1.1.6. Cela achève la démonstration de 3.5.3.

}{\begin{equation*}
\pi^*(E) \xleftarrow{\sim} \mathrm{pr}_1^*(E)\mathrm{pr}_2^*(E) \qquad .
\end{equation*}
D'autre part, en explicitant la condition de compatibilité de $\varphi$ avec $d_o d_1 = 0$, on trouve qu'elle s'exprime par la \underline{condition de commutativité} de (1.2.1), et par la \underline{condition d'associativité} de (1.1.4.1), ces deux conditions correspondant respectivement aux deux facteurs de $L_2$ . Le fait que (*) est une équivalence de catégories n'est alors autre que le dictionnaire de 1.1.6. Cela achève la démonstration de 3.5.3.

}{F7I-VII-069}\underline{Remarque} 3.5.4. a)  Il serait intéressant de renforcer 3.5.3 par la construction d'une \underline{résolution} $L.(P)$ de $P$, dépendant fonctoriellement de $P$, et dont les composantes seraient isomorphes (fonctoriellement en $P$) à des sommes de $\mathbb{Z}$ -modules libres engendrés par des puissances cartésiennes
\SGApage{I}{202}
de $P$. Compte tenu des isomorphismes de la forme (1.4.5), cela donnerait en effet une méthode de calcul partiel des $\mathrm{Ext}^i(P,A)$ comme aboutissement d'une suite spectrale dont le terme  initial s'expliciterait en termes des $H^j(P \times P .. \times P, A)$. (Ici nous ne trouvons que des renseignements très partiels sur les $\mathrm{Ext}^1$, à l'exclusion des $\mathrm{Ext}^i$ supérieurs). Il y a en tous cas une façon naturelle de continuer la résolution partielle $L.(P)$ d'\underline{un} cran.(*)

b) On peut donner une variante de la construction (3.5.1) et de 3.5.3 dans la catégorie des $A$-Modules, où $A$ est un Anneau fixé de $\underline{T}$. La forme du complexe $L.(P)$ devient plus compliquée, à cause de la nécessité de compliquer 1.1.6 (cf. 2.10.3). On trouve

\begin{equation}\tag{3.5.5}
\left\{
\begin{array}{lcl}
L_o & = & A[P] \\[4pt]
L_1 & = & A[P \times P] + A[A \times P] \\[4pt]
L_2 & = & A[P \times P \times P] + A[P \times P] + A[A \times A \times P] + A[A \times P \times P] \quad ,
\end{array}
\right.
\end{equation}
la structure de complexe augmenté vers $P$ étant donnée par les formules (3.5.2), complétées par les formules

\begin{equation}\tag{3.5.6}
\left\{
\begin{array}{l}
d_o[\lambda,p] = [\lambda p] - \lambda[p] \\[4pt]
d_1[\lambda,\mu,p] = [\lambda,\mu p] - [\lambda\mu ,p] + \lambda[\mu,p] \\[4pt]
d_1[\lambda,p,q] = [\lambda,p+q] - [\lambda,p] - [\lambda,q] + \lambda[p,q] - [\lambda p,\lambda q].
\end{array}
\right.
\end{equation}
La relation $d_o d_1 = 0$ résulte de la relation analogue pour la restriction des opérateurs $d_o,d_1$ aux composantes envisagées dans (3.5.2) (qui, on l'a vu, était équivalente à l'associativité et à la commutativité de la loi d'addition donnée dans $P$), et des relations $d_o d_1(\lambda,\mu,p) = 0$ et $d_o d_1(\lambda,p,q) = 0$, qui, comme on le vérifie immédiatement, équivalent 

(*)  L. BREEN a construit une telle résolution (Note du correcteur).


\SGApage{I}{203}

 respectivement à la condition que l'opération donnée de $A$ sur $P$ est associative (i.e. $\lambda(\mu p) = (\lambda\mu)p$) et distributive par rapport à l'addition (i.e. $\lambda(p+q) = \lambda p + \lambda q$). Pour un torseur $E$ sous $G_P$, la donnée d'une trivialisation de $d_o^*(E)$ revient à la donnée d'une loi de composition interne sur $E$ (satisfaisant les conditions de compatibilité déjà explicitées par ailleurs avec la loi de composition interne de $E$ et la structure de torseur), \underline{et} d'une loi de composition externe de $A$ sur $E$, compatible avec l'opération de $A$ sur $P$ et sur $G$ ; la compatibilité de cette trivialisation avec la trivialisation évidente de $d_1^* d_o^*(E)$ équivaut alors aux conditions déjà explicitées plus haut (exprimant que la loi de composition interne de $E$ est une loi de Groupe commutatif, extension de $P$ par $G$), \underline{plus} la \underline{condition d'associativité de la loi externe}, et la \underline{condition de distributivité} de celle-ci par rapport à la loi interne ; conditions qui équivalent à dire que $E$ est (pour les deux lois de composition envisagées, et sa structure de torseur) un $A$-Module extension des $A$-Modules $P$ et $G$. De ceci, on déduit que $L.$ est une résolution partielle du Module $P$(i.e. que $H_1(L.) = 0$) par essentiellement la même méthode que 3.5.3. Il serait évidemment encore intéressant de prolonger cette résolution partielle en une résolution infinie de type analogue, les composantes étant des Modules de la forme $A(A \times A \ldots \times P \times P \ldots)$ ; comparer la question analogue a) dans le cas des Groupes commutatifs.


\SGApage{I}{204}


3.6.  \underline{Application à la description de la catégorie cofibrée des biextensions d'un couple} $(P,Q)$ \underline{donné de Groupes commutatifs}

\SGAvar{Soient $P$, $Q$ deux $\mathbb{Z}$-Modules du topos $\underline{T}$. Choisissons arbitrairement une résolution plate $L'.(P)$ de $P$, coïncident en dimension $\leq 2$ avec le complexe augmenté $L.(P)$ de 3.5, - ce qui est possible grâce à 3.5.3 et grâce au fait que tout $\mathbb{Z}$-Module sur $\underline{T}$ est quotient d'un $\mathbb{Z}$-Module libre, donc plat. Choisissons de même une résolution plate $L'.(Q)$ de $Q$ coïncidant avec $L.(Q)$ en dimensions $\leq 2$. On a donc des homomorphismes canoniques

}{Soient $P$, $Q$ deux $\mathbb{Z}$-Modules du topos $\underline{T}$. Choisissons arbitrairement une résolution plate $L'.(P)$ de $P$, coïncidant en dimension $\leq 2$ avec le complexe augmenté $L.(P)$ de 3.5, - ce qui est possible grâce à 3.5.3 et grâce au fait que tout $\mathbb{Z}$-Module sur $\underline{T}$ est quotient d'un $\mathbb{Z}$-Module libre, donc plat. Choisissons de même une résolution plate $L'.(Q)$ de $Q$ coïncidant avec $L.(Q)$ en dimensions $\leq 2$. On a donc des homomorphismes canoniques

}{F7I-VII-070}\begin{equation}\tag{3.6.1}
L.(P) \longrightarrow L'.(P) \qquad , \qquad L.(Q) \longrightarrow L'.(Q) \qquad ,
\end{equation}
induisant un homomorphisme canonique

\begin{equation}\tag{3.6.2}
L.(P,Q) \overset{\mathrm{dfn}}{=} L.(P) \otimes L.(Q) \longrightarrow L'.(P,Q) \overset{\mathrm{dfn}}{=} L'.(P) \otimes L'.(Q) \simeq P \overset{\mathbb{L}}{\otimes} Q \quad ,
\end{equation}
qui est un isomorphisme sur les composantes de dimension $\leq 2$. Par suite on a, en vertu des définitions,

\SGAvar{\begin{equation}\tag{3.6.3}
\underline{\Psi}_{L.(P,Q)} = \underline{\Psi}_{L.(P',Q')} \qquad .
\end{equation}
\underline{Théorème} 3.6.4. \underline{La catégorie cofibrée additive}, \underline{sur la catégorie} $\underline{C} = \mathrm{GRAB}(\underline{T})$ \underline{des Groupes commutatifs de} $\underline{T}$, \underline{formée des biextensions de} $(P,Q)$ \underline{par un Groupe abélien} $G$ \underline{variable} (2.4), \underline{est canoniquement équivalente à la catégorie cofibrée additive} $\underline{\Psi}_{L.(P,Q)} = \underline{\Psi}_{L'.(P,Q)}$, \underline{définie en termes de} $L.(P,Q) = L.(P) \otimes L.(Q)$ \underline{comme en} 3.1.

}{\begin{equation}\tag{3.6.3}
\underline{\Psi}_{L.(P,Q)} = \underline{\Psi}_{L'.(P,Q)} \qquad .
\end{equation}
\underline{Théorème} 3.6.4. \underline{La catégorie cofibrée additive}, \underline{sur la catégorie} $\underline{C} = \mathrm{GRAB}(\underline{T})$ \underline{des Groupes commutatifs de} $\underline{T}$, \underline{formée des biextensions de} $(P,Q)$ \underline{par un Groupe abélien} $G$ \underline{variable} (2.4), \underline{est canoniquement équivalente à la catégorie cofibrée additive} $\underline{\Psi}_{L.(P,Q)} = \underline{\Psi}_{L'.(P,Q)}$, \underline{définie en termes de} $L.(P,Q) = L.(P) \otimes L.(Q)$ \underline{comme en} 3.1.

}{F7I-VII-071}De ceci, et de 3.2.5, on déduit le résultat principal du présent numéro, annoncé plus haut :


\SGApage{I}{205}

\underline{Corollaire} 3.6.5. \underline{Pour} $P$ \underline{et} $Q$ \underline{fixés}, \underline{on a un isomorphisme canonique,} \underline{fonctoriel en} $G$ :

\begin{equation}\tag{3.6.5.1}
\mathrm{Biext}^1(P,Q;G) \;\simeq\; \mathrm{Ext}^1(P \overset{\mathbb{L}}{\otimes} Q,\, G) \quad .
\end{equation}
Démontrons 3.6.5, et pour ceci explicitons d'abord les composantes de dimension $0,1$ et $2$ de $L.(P,Q)$ :

\SGAvar{\begin{equation}\tag{3.6.6}
\left\{
\begin{aligned}
&L_o(P,Q) = L_o(P) \otimes L_o(Q) = \mathbb{Z}[P\times Q] \\
&L_1(P,Q) = L_o(P) \otimes L_1(Q) + L_1(P) \otimes L_o(Q) = \mathbb{Z}[P\times Q\times Q] + \mathbb{Z}[P\times P\times Q] \\
&L_2(P,Q) = L_o(P) \otimes L_2(Q) + L_1(P) \otimes L_1(Q) + L_2(P) \otimes L_o(Q) \\
&\qquad = \mathbb{Z}[P\times Q\times Q] + \mathbb{Z}[P\times Q\times Q\times Q] \\
&\qquad\quad + \mathbb{Z}[P\times P\times Q\times Q] \\
&\qquad\quad + \mathbb{Z}[P\times P\times Q] + \mathbb{Z}[P\times P\times P\times Q] \quad .
\end{aligned}
\right.
\end{equation}
Pour expliciter la catégorie cofibrée $\underline{\Psi}_{L.(P,Q)}$ via ses catégories fibres $\underline{\Psi}_{L.(P,Q)}(G)$, on utilise 1.4 pour expliciter les extensions des $L_i(P,Q)$ (pour $i=0,1,2$) par $G$, en termes de torseurs. Ainsi, une extension de $L_o(P\times Q)$ par $G$ s'identifie à un torseur sous $G_{P\times Q}$, une extension de $L_1(P,Q)$ par $G$ s'identifie à un couple d'un torseur sous $G_{P\times Q\times Q}$ et d'un autre sous $G_{P\times P\times Q}$ , et de même une extension de $L_2(P,Q)$ par $G$ s'identifie à un système de cinq torseurs, sous les groupes déduits de $G$ par changement de base sur les bases $P\times Q\times Q, P\times Q\times Q\times Q, P\times P\times Q\times Q, P\times P\times Q, P\times P\times P\times Q$ respectivement. Utilisant ce dictionnaire, on peut interpréter un couple $(X,\alpha)$ d'une extension de $L_o(P,Q)$ par $G$, et d'une trivialisation de $d_o^*(X)$ sur $L_1(P,Q)$, comme un torseur $E$ sur $P\times Q$, muni d'un couple $(\varphi,\psi)$ de trivialisations du couple de torseurs $(E_1,E_2)$ sur $P\times Q\times Q$ resp. $P\times P\times Q$, déduit de $E$ par image inverse
\SGApage{I}{206}
à l'aide de $d_o$ : $L_1(P,Q) \longrightarrow L_o(P,Q)$. Explicitant l'opérateur différentiel $d_o$, on trouve comme dans la démonstration de 3.5.3 que $\varphi$ et $\psi$ peuvent s'interpréter comme desisomorphismes

\begin{equation*}
\varphi : \pi_Q^*(E) \overset{\sim}{\longleftarrow} q_1^*(E)q_2^*(E) \;,\quad \psi : \pi_P^*(E) \overset{\sim}{\longleftarrow} p_1^*(E)p_2^*(E) \qquad ,
\end{equation*}
où $\pi_Q,q_1q_2$ sont les trois morphismes $P\times Q\times Q \longrightarrow Q$ déduits respectivement des trois morphismes $Q\times Q \longrightarrow Q$ loi de composition, première et deuxième projection, et où $\pi_P,p_1,p_2: P\times P\times Q \longrightarrow P\times Q$ sont définis de façon symétrique. Donc $\varphi$ et $\psi$ peuvent s'interpréter comme des données du type (2.0.4) et (2.0.7) qui figurent dans 2.1. Enfin, la condition que le couple $(X,\alpha)$ soit bien un objet de $\underline{\Psi}_{L.(P,Q)}(G)$, i.e. que la trivialisation partielle $\alpha$ soit compatible avec la relation $d_od_1=0$, s'exprime en termes de la structure $(E,\varphi,\psi)$ par cinq relations de compatibilité, en termes de torseurs sur $P\times Q\times Q$, $P\times Q\times Q\times Q$, $P\times P\times Q\times Q$, $P\times P\times Q$, $P\times P\times P\times Q$ respectivement, ces cinq relations correspondant aux cinq termes intervenant dans l'expression (3.6.6) de $L_2(P,Q)$. On vérifie alors comme dans la démonstration de 3.5.3 que les deux premières relations ne sont autres que les relations (2.0.6) de commutativité et (2.0.5) d'associativité, dont la conjonction exprime que $\varphi$ définit sur $E$ une structure d'extension commutative de $Q_P$ par $G_P$ . De façon symétrique, les deux dernières conditions expriment que $\psi$ définit sur $E$ une structure d'extension commutative de $P_Q$ par $G_P$ . Enfin, la condition médiane, portant sur des torseurs sur $P\times P\times Q\times Q$, n'est autre que la commutativité du diagramme (2.1.1), qui exprime que $\varphi$ et $\psi$ sont compatibles, i.e. (compte tenu des autres conditions) qu'ils définissent sur $E$ une structure de biextension de $(P,Q)$ par $G$. La vérification de ces
\SGApage{I}{207}
interprétations est essentiellement triviale, et laissée au lecteur. Le lecteur admettra sans doute avec nous que les discours précédents définissent bien une équivalence de catégories cofibrées sur $\underline{C} = \mathrm{GRAB}(\underline{T})$, qui est celle annoncée dans 3.6.4.

}{\begin{equation}\tag{3.6.6}
\left\{
\begin{aligned}
&L_o(P,Q) = L_o(P) \otimes L_o(Q) = \mathbb{Z}[P\times Q] \\
&L_1(P,Q) = L_o(P) \otimes L_1(Q) + L_1(P) \otimes L_o(Q) = \mathbb{Z}[P\times Q\times Q] + \mathbb{Z}[P\times P\times Q] \\
&L_2(P,Q) = L_o(P) \otimes L_2(Q) + L_1(P) \otimes L_1(Q) + L_2(P) \otimes L_o(Q) \\
&\qquad = \mathbb{Z}[P\times Q\times Q] + \mathbb{Z}[P\times Q\times Q\times Q] \\
&\qquad\quad + \mathbb{Z}[P\times P\times Q\times Q] \\
&\qquad\quad + \mathbb{Z}[P\times P\times Q] + \mathbb{Z}[P\times P\times P\times Q] \quad .
\end{aligned}
\right.
\end{equation}
Pour expliciter la catégorie cofibrée $\underline{\Psi}_{L.(P,Q)}$ via ses catégories fibres $\underline{\Psi}_{L.(P,Q)}(G)$, on utilise 1.4 pour expliciter les extensions des $L_i(P,Q)$ (pour $i=0,1,2$) par $G$, en termes de torseurs. Ainsi, une extension de $L_o(P,Q)$ par $G$ s'identifie à un torseur sous $G_{P\times Q}$, une extension de $L_1(P,Q)$ par $G$ s'identifie à un couple d'un torseur sous $G_{P\times Q\times Q}$ et d'un autre sous $G_{P\times P\times Q}$ , et de même une extension de $L_2(P,Q)$ par $G$ s'identifie à un système de cinq torseurs, sous les groupes déduits de $G$ par changement de base sur les bases $P\times Q\times Q, P\times Q\times Q\times Q, P\times P\times Q\times Q, P\times P\times Q, P\times P\times P\times Q$ respectivement. Utilisant ce dictionnaire, on peut interpréter un couple $(X,\alpha)$ d'une extension de $L_o(P,Q)$ par $G$, et d'une trivialisation de $d_o^*(X)$ sur $L_1(P,Q)$, comme un torseur $E$ sur $P\times Q$, muni d'un couple $(\varphi,\psi)$ de trivialisations du couple de torseurs $(E_1,E_2)$ sur $P\times Q\times Q$ resp. $P\times P\times Q$, déduit de $E$ par image inverse
\SGApage{I}{206}
à l'aide de $d_o$ : $L_1(P,Q) \longrightarrow L_o(P,Q)$. Explicitant l'opérateur différentiel $d_o$, on trouve comme dans la démonstration de 3.5.3 que $\varphi$ et $\psi$ peuvent s'interpréter comme des isomorphismes

\begin{equation*}
\varphi : \pi_Q^*(E) \overset{\sim}{\longleftarrow} q_1^*(E)q_2^*(E) \;,\quad \psi : \pi_P^*(E) \overset{\sim}{\longleftarrow} p_1^*(E)p_2^*(E) \qquad ,
\end{equation*}
où $\pi_Q,q_1,q_2$ sont les trois morphismes $P\times Q\times Q \longrightarrow P\times Q$ déduits respectivement des trois morphismes $Q\times Q \longrightarrow Q$ loi de composition, première et deuxième projection, et où $\pi_P,p_1,p_2: P\times P\times Q \longrightarrow P\times Q$ sont définis de façon symétrique. Donc $\varphi$ et $\psi$ peuvent s'interpréter comme des données du type (2.0.4) et (2.0.7) qui figurent dans 2.1. Enfin, la condition que le couple $(X,\alpha)$ soit bien un objet de $\underline{\Psi}_{L.(P,Q)}(G)$, i.e. que la trivialisation partielle $\alpha$ soit compatible avec la relation $d_od_1=0$, s'exprime en termes de la structure $(E,\varphi,\psi)$ par cinq relations de compatibilité, en termes de torseurs sur $P\times Q\times Q$, $P\times Q\times Q\times Q$, $P\times P\times Q\times Q$, $P\times P\times Q$, $P\times P\times P\times Q$ respectivement, ces cinq relations correspondant aux cinq termes intervenant dans l'expression (3.6.6) de $L_2(P,Q)$. On vérifie alors comme dans la démonstration de 3.5.3 que les deux premières relations ne sont autres que les relations (2.0.6) de commutativité et (2.0.5) d'associativité, dont la conjonction exprime que $\varphi$ définit sur $E$ une structure d'extension commutative de $Q_P$ par $G_P$ . De façon symétrique, les deux dernières conditions expriment que $\psi$ définit sur $E$ une structure d'extension commutative de $P_Q$ par $G_Q$ . Enfin, la condition médiane, portant sur des torseurs sur $P\times P\times Q\times Q$, n'est autre que la commutativité du diagramme (2.1.1), qui exprime que $\varphi$ et $\psi$ sont compatibles, i.e. (compte tenu des autres conditions) qu'ils définissent sur $E$ une structure de biextension de $(P,Q)$ par $G$. La vérification de ces
\SGApage{I}{207}
interprétations est essentiellement triviale, et laissée au lecteur. Le lecteur admettra sans doute avec nous que les discours précédents définissent bien une équivalence de catégories cofibrées sur $\underline{C} = \mathrm{GRAB}(\underline{T})$, qui est celle annoncée dans 3.6.4.

}{F7I-VII-072}\underline{Remarque} 3.6.7. Chacune des trois généralisations signalées dans 2.10 de la notion de biextension permet de donner une variante correspondante de 3.6.4 (donc aussi de 3.6.5). La généralisation au cas d'un nombre quelconque de facteurs (cf. 2.10.2) est triviale, et la généralisation au cas de biextensions de $A$-Modules (cf. 2.10.3) est immédiate, une fois qu'on a obtenu la variante de la construction de $L.(P)$ dans ce cas, signalée dans 3.5.4 b). Plus intéressante du point de vue de l'algèbre homologique semblerait la généralisation au cas où $P$ et $Q$ ne sont pas supposés commutatifs (cf. 2.10.1). Il convient alors d'associer à un Groupe $P$ (pas nécessairement commutatif) de $\underline{T}$ un complexe tronqué à l'ordre 2, soit $K.(P)$, qui se décrit de façon analogue que $L.(P)$ dans 3.5, avec la différence que l'on prend $K_2(P) = \mathbb{Z}[P \times P \times P]$ ( en omettant donc le premier facteur qui figurait dans $L.(P)$), $d_o$ et $d_1$ étant définis comme dans (3.5.2). Ici $H_o(K.(P))$ s'identifie au Groupe $P/[P,P]$ déduit de $P$ en le rendant abélien, qu'on peut aussi identifier au faisceau $\underline{H}_1(P,\mathbb{Z})$ des groupes d'homologie de degré 1 de $P$ à coefficients dans $\mathbb{Z}$, et $H_1(K.(P))$ est de même isomorphe au faisceau $\underline{H}_2(P,\mathbb{Z} )$ ; en fait, $K.(P)$ n'est autre que le \underline{tronqué à l'ordre} 2 \underline{du complexe des chaînes de} $P$ \underline{à coefficients dans} $\mathbb{Z}$, soit $C.(P,\mathbb{Z} )$. Ceci dit, l'argument établissant 3.5.3 établit de même que la catégorie cofibrée $\underline{\Psi}_{K.(P)}$ sur $\underline{C} = \mathrm{GRAB}(\underline{T})$ est
\SGApage{I}{208}
canoniquement équivalente à la catégorie des \underline{extensions centrales} de $P$ par des Groupes commutatifs $G$ variables. On trouve de même, pour deux Groupes $P,Q$, une équivalence de la catégorie cofibrée des biextensions de $(P,Q)$ au sens de 2.10.1, avec la catégorie cofibrée $\underline{\Psi}_{K.(P,Q)}$, où $K.(P,Q) = K.(P) \overset{\mathbb{L}}{\otimes} K.(Q)$.

Ces réflexions suggèrent qu'on trouvera peut-être, pour $P$ commutatif, une résolution canonique de $P$ comme souhaitée dans 3.5.4 a), par une modification judicieuse du complexe de chaînes $C.(P,\mathbb{Z} )$.


3.7. \underline{Compatibilités diverses}

\SGAvar{Dans la présente section , nous examinons le comportement de l'équivalence de catégorie cofibrées 3.6.4, pour $P$ et $Q$ variables. La variance de la catégorie cofibrée des biextensions de $(P,Q)$ par un $G$ variable, par rapport à $P,Q$ variables, est définie en principe dans 2.4, et peut aussi se formuler par l'assertion que l'on a un foncteur cofibrant (2.6.1). On notera que la variance en $P,Q$ est \underline{contravariante}. La variance de la catégorie cofibrée $\underline{\Psi}_{L.(P,Q)}$ par rapport à $P,Q$ variables est également claire, puisque $L.(P,Q) = L.(P) \otimes L.(Q)$ dépend fonctoriellement du couple $(P,Q)$ et que $\underline{\Psi}_{L.}$ dépend de façon contravariante de $L.$ (3.4). Ceci posé, une première compatibilité s'exprime, pour tout triple de morphismes

}{Dans la présente section , nous examinons le comportement de l'équivalence de catégories cofibrées 3.6.4, pour $P$ et $Q$ variables. La variance de la catégorie cofibrée des biextensions de $(P,Q)$ par un $G$ variable, par rapport à $P,Q$ variables, est définie en principe dans 2.4, et peut aussi se formuler par l'assertion que l'on a un foncteur cofibrant (2.6.1). On notera que la variance en $P,Q$ est \underline{contravariante}. La variance de la catégorie cofibrée $\underline{\Psi}_{L.(P,Q)}$ par rapport à $P,Q$ variables est également claire, puisque $L.(P,Q) = L.(P) \otimes L.(Q)$ dépend fonctoriellement du couple $(P,Q)$ et que $\underline{\Psi}_{L.}$ dépend de façon contravariante de $L.$ (3.4). Ceci posé, une première compatibilité s'exprime, pour tout triple de morphismes

}{F7I-VII-073}\begin{equation}\tag{3.7.1}
P' \longrightarrow P \;\; , \quad Q' \longrightarrow Q \;\; , \quad G \longrightarrow G' \quad ,
\end{equation}
par le fait que le diagramme de foncteurs

\begin{equation}\tag{3.7.2}
\begin{tikzcd}
\underline{\Psi}_{L.(P,Q)}(G) \arrow[r] \arrow[d] \& \underline{\Psi}_{L.(P',Q')}(G') \arrow[d] \\
\mathrm{BIEXT}(P,Q;G) \arrow[r] \& \mathrm{BIEXT}(P',Q';G')
\end{tikzcd}
\end{equation}












\SGApage{I}{209}
est commutatif à un isomorphisme canonique près. Nous admettrons cette compatibilité, l'explicitation de l'isomorphisme de commutativité ne pouvant être qu'aussi évidente que fastidieuse (tout comme, du reste, l'explicitation complète qu'elle suppose de l'équivalence de catégories cofibrées 3.6.4).

3.7.3. De la compatibilité précédente résulte en particulier que l'isomorphisme (3.6.5.1) est fonctoriel non seulement en $G$, mais aussi en $P$ et $Q$. On peut énoncer de même, et plus trivialement, une fonctorialité analogue pour l'isomorphisme défini via (3.1.8) comme (3.6.5.1) l'était via (3.2.4) :

\begin{equation}\tag{3.7.4}
\mathrm{Biext}^{\mathrm{o}}(P,Q;G) \xrightarrow{\ \sim\ } \mathrm{Ext}^{\mathrm{o}}(P \overset{\mathbb{L}}{\otimes} Q,G) \simeq \mathrm{Hom}(P \otimes Q,G) \qquad .
\end{equation}
Mais on notera qu'on avait déjà obtenu un tel isomorphisme dans (2.5.4.1), et noté (2.6) qu'il était fonctoriel en ses trois arguments. Or il est immédiat de vérifier que les deux isomorphismes (3.7.4) et (2.5.4.1) sont égaux.

3.7.5. L'intérêt de l'isomorphisme de (3.6.5) est qu'il exprime le groupe remarquable de nature géométrique, $\mathrm{Biext}^{1}(P,Q;G)$, en termes d'un invariant d'algèbre homologique très simple $\mathrm{Ext}^{1}(P \overset{\mathbb{L}}{\otimes} Q,G)$, qui s'insère dans la suite des invariants $\mathrm{Ext}^{i}(P \overset{\mathbb{L}}{\otimes} Q,G)$, à laquelle s'applique le formalisme habituel d'algèbre homologique. En particulier, la suite de ces $\mathrm{Ext}^{i}$ est un \underline{foncteur cohomologique} au sens de [3] (ou \underline{suite liée de foncteurs} dans la terminologie de [5]) en chacun des trois arguments $P,Q,G$. En particulier, une suite exacte

\begin{equation}\tag{3.7.5.1}
0 \longrightarrow G' \longrightarrow G \longrightarrow G'' \longrightarrow 0 \qquad 
\end{equation}
\SGApage{I}{210}

 (resp.

\begin{equation}\tag{3.7.5.2}
0 \longrightarrow P' \longrightarrow P \longrightarrow P'' \longrightarrow 0,\ \text{ou}\ 0 \longrightarrow Q' \longrightarrow Q \longrightarrow Q'' \longrightarrow 0)
\end{equation}
donne naissance à une suite exacte infinie sur les $\mathrm{Ext}^{i}$, dont les premiers six termes (ne faisant intervenir que des $\mathrm{Ext}^{\mathrm{o}}$ et $\mathrm{Ext}^{1}$) s'interprètent en termes de biextensions grâce à 3.6.5 et 3.7.4. On a déjà signalé une interprétation de la suite exacte déduite de (3.7.5.1), et notamment de l'opérateur cobord

\begin{equation}\tag{3.7.5.3}
\partial : \mathrm{Biext}^{\mathrm{o}}(P,Q;G'') \longrightarrow \mathrm{Biext}^{1}(P,Q;G'),
\end{equation}
en termes géométriques, dans le cadre de la notion de \underline{catégorie cofibrée exacte à gauche} (3.3.6). On peut se proposer de donner une interprétation analogue de la suite exacte à six termes associée à l'une des suites exactes courtes (3.7.5.2), la première disons, ou ce qui revient au même compte tenu de 3.7.3, de donner une interprétation géométrique de l'homomorphisme cobord

\begin{equation}\tag{3.7.5.4}
\partial : \mathrm{Biext}^{\mathrm{o}}(P',Q;G) \longrightarrow \mathrm{Biext}^{1}(P'',Q;G) \qquad .
\end{equation}
On a effectivement :

\SGAvar{\underline{Proposition} 3.7.6. \underline{La catégorie cofibrée sur} $\mathrm{GRAB}(\underline{T})^{\mathrm{o}} \times \mathrm{GRAB}(\underline{T})^{\mathrm{o}} \times \mathrm{GRAB}(\underline{T})$ \underline{des biextensions de} $(P,Q)$ \underline{par} $G$, \underline{avec} $P,Q,G$ \underline{variables} (cf. (2.6.1)) \underline{est exacte à gauche non seulement en l'argument} $G$, \underline{mais aussi en chacun des arguments} $P$ \underline{et} $Q$. \underline{Ceci dit,} \underline{l'homomorphisme cobord} (3.7.5.4) \underline{déduit de la suite exacte des} $\mathrm{Ext}^{i}$ \underline{associée à} (3.7.5.2) (\underline{grâce aux isomorphismes} 3.6.5 \underline{et} 3.7.4) \underline{est l'homomorphisme cobord de la théorie des catégories}
\SGApage{I}{211}
\underline{cofibrées exactes à gauches sur} $\underline{C} = \mathrm{GRAB}(\underline{T})^{o}$ \underline{rappelé dans} 3.3.6, \underline{appliqué} \underline{à la catégorie cofibrée des biextensions de} $(P,Q)$ \underline{par} $G$, \underline{où} $Q$ \underline{et} $G$ \underline{sont} \underline{fixés}, \underline{et} $P$ \underline{est variable} (\underline{sous réserve d'erreur de signe dans} 3.4.9, (cf. 3.4.9.2)).

}{\underline{Proposition} 3.7.6. \underline{La catégorie cofibrée sur} $\mathrm{GRAB}(\underline{T})^{\mathrm{o}} \times \mathrm{GRAB}(\underline{T})^{\mathrm{o}} \times \mathrm{GRAB}(\underline{T})$ \underline{des biextensions de} $(P,Q)$ \underline{par} $G$, \underline{avec} $P,Q,G$ \underline{variables} (cf. (2.6.1)) \underline{est exacte à gauche non seulement en l'argument} $G$, \underline{mais aussi en chacun des arguments} $P$ \underline{et} $Q$. \underline{Ceci dit,} \underline{l'homomorphisme cobord} (3.7.5.4) \underline{déduit de la suite exacte des} $\mathrm{Ext}^{i}$ \underline{associée à} (3.7.5.2) (\underline{grâce aux isomorphismes} 3.6.5 \underline{et} 3.7.4) \underline{est l'homomorphisme cobord de la théorie des catégories}
\SGApage{I}{211}
\underline{cofibrées exactes à gauche sur} $\underline{C} = \mathrm{GRAB}(\underline{T})^{o}$ \underline{rappelé dans} 3.3.6, \underline{appliqué} \underline{à la catégorie cofibrée des biextensions de} $(P,Q)$ \underline{par} $G$, \underline{où} $Q$ \underline{et} $G$ \underline{sont} \underline{fixés}, \underline{et} $P$ \underline{est variable} (\underline{sous réserve d'erreur de signe dans} 3.4.9, (cf. 3.4.9.2)).

}{F7I-VII-074}Esquissons la démonstration de 3.7.6. Grâce à 3.6.4 et la compatibilité (3.7.2), nous interprétons les biextensions de $(P,Q)$ par $G$ comme les objets de $\underline{\Psi}_{L.(P,Q)}(G)$. Considérons l'homomorphisme

\SGAvar{\begin{equation*}
L.(P,Q)=L.(P) \otimes L.(Q) \longrightarrow P \otimes L.(Q)
\end{equation*}
défini par l'augmentation $L.(P) \longrightarrow P$ ; il résulte de 3.5.3 et de la platitude de $L.(Q)$ que cet homomorphisme induit un isomorphisme sur les objets $H_o$ et $H_1$, donc (3.4.7) une équivalence $\underline{\Psi}_{P \otimes L.(Q)} \longrightarrow \underline{\Psi}_{L.(P,Q)}$. Ceci nous permet d'interpréter les biextensions de $(P,Q)$ par $G$ comme des objets de $\underline{\Psi}_{P \otimes L.(Q)}(G)$. Ceci posé, la suite exacte (3.7.5.2) donne une suite exacte de complexe de chaînes

}{\begin{equation*}
L.(P,Q)=L.(P) \otimes L.(Q) \longrightarrow P \otimes L.(Q)
\end{equation*}
défini par l'augmentation $L.(P) \longrightarrow P$ ; il résulte de 3.5.3 et de la platitude de $L.(Q)$ que cet homomorphisme induit un isomorphisme sur les objets $H_o$ et $H_1$, donc (3.4.7) une équivalence $\underline{\Psi}_{P \otimes L.(Q)} \longrightarrow \underline{\Psi}_{L.(P,Q)}$. Ceci nous permet d'interpréter les biextensions de $(P,Q)$ par $G$ comme des objets de $\underline{\Psi}_{P \otimes L.(Q)}(G)$. Ceci posé, la suite exacte (3.7.5.2) donne une suite exacte de complexes de chaînes

}{F7I-VII-075}\begin{equation}\tag{3.7.6.1}
0 \longrightarrow P' \otimes L.(Q) \longrightarrow P \otimes L.(Q) \longrightarrow P'' \otimes L.(Q) \longrightarrow 0 \quad ,
\end{equation}
L'exactitude à gauche annoncée résulte alors de 3.4.8 appliqué à cette suite exacte de complexes.(La vérification directe serait d'ailleurs immédiate.) D'autre part, l'homomorphisme (3.7.5.4) n'est autre que l'homomorphisme cobord $\mathrm{Ext}^{o} \longrightarrow \mathrm{Ext}^{1}$ déduit de cette suite exacte, en prenant les $\mathrm{Ext}^{i}$ à valeurs dans $G$, et la dernière assertion de 3.7.6 est donc un cas particulier de 3.4.9.


\SGApage{I}{212}


3.8. \underline{Biextensions de} $(P,Q)$ \underline{trivialisées sur des extensions de} $P$ \underline{et} $Q$

Considérons d'abord deux suites exactes de Groupes abéliens

\begin{equation}\tag{3.8.1}
0 \longrightarrow L_1 \xrightarrow{\ i\ } L_o \xrightarrow{\ u\ } P \longrightarrow 0 \ , \ 0 \longrightarrow M_1 \xrightarrow{\ j\ } M_o \xrightarrow{\ v\ } Q \longrightarrow 0
\end{equation}
Nous nous proposons de décrire la catégorie des biextensions $E$ de $(P,Q)$ par un Groupe abélien $G$, telles que l'image inverse $(u,v)^{*}\ (E)$ de $E$ par $(u,v)$ soit triviale, i.e. telles qu'il existe un homomorphisme $s$ rendant commutatif le triangle suivant

\SGAvar{\begin{equation}\tag{3.8.2}
\begin{tikzcd}[column sep=huge, row sep=huge]
E \arrow[d] \arrow[dr, dash, "s"] \& \\
P \times Q \& L_o \times M_o \arrow[l, "u \times v"]
\end{tikzcd}\quad ,
\end{equation}
















}{\begin{equation}\tag{3.8.2}
\begin{tikzcd}[column sep=huge, row sep=huge]
E \arrow[d] \& \\
P \times Q \& L_o \times M_o \arrow[l, "u \times v"] \arrow[ul, "s"]
\end{tikzcd}\quad ,
\end{equation}














}{F7I-VII-076}et satisfaisant la condition d'additivité qu'il faut en chacun de ses deux arguments. La donnée d'un tel $s$ permet donc d'identifier $(u,v)^{*}(E)$ à la biextension triviale $E_o = L_o \times M_o \times G$ de $(L_o,M_o)$ par $G$. La propriété d'exactitude à gauche en $P,Q$ de la catégorie cofibrée des biextensions de $(P,Q)$ par $G$ fixé (3.7.6) montre alors que la donnée de la biextension $E$ de $(P,Q)$, munie de la trivialisation $s$ de son image inverse, équivaut essentiellement à la donnée de deux trivialisations partielles de la biextension triviale $E_o$ de $(L_o,M_o)$ par $G$, sur $(L_1,M_o)$ et $(L_o,M_1)$ respectivement, qui coïncident sur $(L_1,M_1)$. Comme les biextensions induites par $E_o$ sur $(L_1,M_o),(L_o,M_1)$ et $(L_1,M_1)$ s'identifient aux biextensions triviales, une trivialisation partielle sur $(L_1,M_o)$ (resp. sur $(L_o,M_1)$) équivaut à la donnée d'un homomorphisme (ou accouplement)

\begin{equation}\tag{3.8.3}
f : L_1 \otimes M_o \longrightarrow G \quad (\text{resp. } g: L_o \otimes M_1 \longrightarrow G) \quad ,
\end{equation}

\SGApage{I}{213}
et la condition de compatibilité envisagée signifie simplement que les deux homomorphismes en question coïncident sur $L_1 \otimes M_1$, i.e qu'on a

\begin{equation}\tag{3.8.4}
f(x_1 \otimes y_1) = g(x_1 \otimes y_1) \qquad .
\end{equation}
pour des sections $x_1,y_1$ de $L_1$ , $M_1$ sur un objet de $\underline{T}$. En termes de l'homomorphisme $s$ de (3.8.2), les formes $f$ et $g$ sont caractérisées  par les conditions

\begin{equation}\tag{3.8.5}
s(x_1,y_o) = f(x_1 \otimes y_o) e_{E/Q}(y_o) \ , \ s(x_o,y_1) = g(x_o \otimes y_1) e_{E/P}(x_o) \qquad ,
\end{equation}
qui mettent en évidence à nouveau la nécessité de la condition (3.8.4).

La trivialisation $s$ étant choisie, les autres trivialisations sont les morphismes $s'$: $L_o \times M_o \longrightarrow E$ de la forme

\begin{equation}\tag{3.8.6}
s'(x_o,y_o) = s(x_o,y_o) + h(x_o \otimes y_o) \qquad ,
\end{equation}
avec

\begin{equation*}
h \ : \ L_o \otimes M_o \longrightarrow G
\end{equation*}
une forme quelconque ; il est alors clair sur (3.8.5) que les formes $f,g$ définies par $s$ sont changées en les formes

\begin{equation}\tag{3.8.7}
\begin{aligned}
f'(x_1 \otimes y_o) &= f(x_1 \otimes y_o) + h(x_1 \otimes y_o) \\
g'(x_o \otimes y_1) &= g(x_o \otimes y_1) + h(x_o \otimes y_1) \qquad .
\end{aligned}
\end{equation}
On trouve ainsi :

\SGAvar{\underline{Proposition} 3.8.8. \underline{Le sous-groupe de} $\mathrm{Biext}^1(P,Q;G)$ \underline{formée des classes de biextensions dont l'image inverse par} $(u,v)$ \underline{est une biextension} \underline{triviale} \underline{de} $(L_o,M_o)$ (\underline{i}.\underline{e}. \underline{le noyau de} $\mathrm{Biext}^1(P,Q;G) \longrightarrow \mathrm{Biext}^1(L_o,M_o;G)$) \underline{est canoniquement isomorphe, grâce à la description qui précède, au groupe quotient}
\SGApage{I}{214}
\underline{du groupe des couples de formes} $(f,g)$ (3.8.3), \underline{satisfaisant la condition de compatibilité} (3.8.4), \underline{modulo les couples de formes qui proviennent par restriction d'une forme} $h$: $L_o \otimes M_o \longrightarrow G$.

}{\underline{Proposition} 3.8.8. \underline{Le sous-groupe de} $\mathrm{Biext}^1(P,Q;G)$ \underline{formé des classes de biextensions dont l'image inverse par} $(u,v)$ \underline{est une biextension} \underline{triviale} \underline{de} $(L_o,M_o)$ (\underline{i}.\underline{e}. \underline{le noyau de} $\mathrm{Biext}^1(P,Q;G) \longrightarrow \mathrm{Biext}^1(L_o,M_o;G)$) \underline{est canoniquement isomorphe, grâce à la description qui précède, au groupe quotient}
\SGApage{I}{214}
\underline{du groupe des couples de formes} $(f,g)$ (3.8.3), \underline{satisfaisant la condition de compatibilité} (3.8.4), \underline{modulo les couples de formes qui proviennent par restriction d'une forme} $h$: $L_o \otimes M_o \longrightarrow G$.

}{F7I-VII-077}Nous voulons maintenant interpréter cet isomorphisme en termes d'algèbre homologique, lorsqu'on tient compte de l'isomorphisme canonique $\mathrm{Biext}^1(P,Q;G) \simeq \mathrm{Ext}^1(P \overset{\mathbb{L}}{\otimes} Q,G)$ de 3.6.5. Pour ceci, interprétons les données (3.8.1) comme étant deux complexes résolutions de $P$ resp. de $Q$,

\begin{equation}\tag{3.8.9}
L. \ : \ 0 \longrightarrow L_1 \longrightarrow L_o \longrightarrow 0 \ , \ M. \ : \ 0 \longrightarrow M_1 \longrightarrow M_o \longrightarrow 0 \qquad ,
\end{equation}
et considérons le complexe produit tensoriel $L. \otimes M.$

\begin{equation*}
0 \longrightarrow L_1 \otimes M_1 \overset{d_2}{\longrightarrow} L_1 \otimes M_o + L_o \otimes M_1 \overset{d_1}{\longrightarrow} L_o \otimes M_o \longrightarrow 0 \qquad .
\end{equation*}
On sait d'ailleurs qu'on a un homomorphisme canonique dans la catégorie dérivée

\begin{equation*}
P \overset{\mathbb{L}}{\otimes} Q \longrightarrow L. \otimes M. \qquad ,
\end{equation*}
d'où un homomorphisme non moins canonique

\begin{equation}\tag{3.8.10}
H^1(\mathrm{Hom}^{.}(L. \otimes M.,G)) \longrightarrow \mathrm{Ext}^1(P \overset{\mathbb{L}}{\otimes} Q,G) \qquad .
\end{equation}
On notera d'ailleurs que le groupe des 1-cocycles du complexe $\mathrm{Hom}^{.}(L. \otimes M.,G)$ s'identifie au groupe des couples $(f,g)$ comme en (3.8.3), satisfaisant à la condition de compatibilité (3.8.4), en associant à un tel couple l'homomorphisme $F \ : \ L_1 \otimes M_o + L_o \otimes M_1 \longrightarrow G$ de composantes $(f,g)$ :

\begin{equation}\tag{3.8.11}
F(x_1 \otimes y_o + x_o \otimes y_1) = f(x_1 \otimes y_o) + g(x_o \otimes y_1) \qquad .
\end{equation}
D'ailleurs le cobord de la 0-cochaîne $h \ : \ L_o \otimes M_o \longrightarrow G$ n'est autre que le couple de ses restrictions à $L_1 \otimes M_o$ et $L_o \otimes M_1$, de sorte que le groupe des
\SGApage{I}{215}
couples $(f,g)$ modulo les couples $(h,h)$ décrit dans 3.8.8 est canoniquement isomorphe au premier membre de (3.8.10). Ceci posé, on a la compatibilité suivante :

\underline{Proposition} 3.8.12. \underline{Si} $\xi \in \mathrm{Ker}(\mathrm{Biext}^1(P,Q;G) \longrightarrow \mathrm{Biext}^1(L_o,M_o;G))$ \underline{est} \underline{décrit par le couple} $(f,g)$ \underline{comme dit dans} 3.8.8, \underline{alors l'image de} $\xi$ \underline{dans} $\mathrm{Ext}^1(P \overset{\mathbb{L}}{\otimes} Q;G)$ \underline{par l'isomorphisme} (3.6.5.1) \underline{est l'image par} (3.8.10) \underline{de} \underline{l'élément du premier membre défini par le cocycle} (3.8.11).

La vérification de cette compatibilité est laissée au lecteur.

3.8.13. On va appliquer ceci en particulier dans le cas où on prend

\begin{equation*}
L_o = \mathbb{Z}[P] \quad , \quad M_o = \mathbb{Z}[Q] \quad ,
\end{equation*}
$L_1$ et $M_1$ étant respectivement les noyaux des homomorphismes canoniques $\mathbb{Z}[P] \longrightarrow P$, $\mathbb{Z}[Q] \longrightarrow Q$. Notons que l'on a des isomorphismes canoniques

\begin{equation*}
L_o \overset{\mathbb{L}}{\otimes} M_o \quad \simeq \quad L_o \otimes M_o \quad \simeq \quad \mathbb{Z}[P\times Q] \quad ,
\end{equation*}
$L_o$ et $M_o$ étant plats, d'où on déduit grâce à (1.4.5) des isomorphismes canoniques

\begin{equation*}
\mathrm{Ext}^i(L_o \overset{\mathbb{L}}{\otimes} M_o, G) \quad \simeq \quad H^i(P\times Q, G_{P\times Q}) \quad .
\end{equation*}
De ceci et de 3.6.5, ou directement grâce aux définitions et à 1.4, on conclut que le foncteur qui à toute biextension de $(L_o,M_o)$ par $G$ associe le torseur sous $G_{P\times Q}$ sur $P\times Q$ qu'elle induit, est une équivalence de catégories. (Ceci reste valable sans utiliser de structure additive sur $P,Q$ .) Par suite, si $E$ est une biextension de $(P,Q)$ par $G$, il revient au même de se donner une trivialisation de son image inverse sur $(L_o,M_o)$, ou de se donner une section de $E$ considéré seulement comme objet au-dessus de $P\times Q$, i.e. une trivialisation de $E$ considéré comme torseur sous $G_{P\times Q}$.


\SGApage{I}{216}

Les biextensions munies de cette structure supplémentaire sont donc décrites par les couples $(f,g)$ (3.8.3) satisfaisant la compatibilité (3.8.4), et la description à isomorphisme près est celle donnée dans 3.8.8.

Notons d'ailleurs qu'avec les notations de 3.5, on peut considérer $L_1$ comme le conoyau de $L_2(P) \longrightarrow L_1(P)$, et de même pour $M_1$, ce qui permet d'expliciter la donnée de $(f,g)$ en termes d'un système d'homomorphismes

\begin{equation*}
\left\{
\begin{aligned}
f' &: L_1(P) \otimes L_o(Q) = \mathbb{Z}(P\times P\times Q) \longrightarrow G\\
g' &: L_o(P) \otimes L_1(Q) = \mathbb{Z}(P\times Q\times Q) \longrightarrow G \quad ,
\end{aligned}
\right.
\end{equation*}
i.e. d'homomorphismes de faisceaux d'\underline{ensembles}

\begin{equation*}
P\times P\times Q \longrightarrow G \quad , \quad P\times Q\times Q \longrightarrow G \quad ,
\end{equation*}
satisfaisant cinq conditions de compatibilité (exprimant que $f'$ et $g'$ peuvent se factoriser en $f$ et $g$ - ce qui fait $2\times 2$ équations à écrire, compte tenu de la forme de $L_2(P)$ et $L_2(Q)$ comme somme de deux termes - et la condition (3.8.4)). On constate alors que $f'$ et $g'$ ne sont autres que les données de multiplication partielle définissant une biextension, et les cinq conditions en question ne sont autres que les cinq axiomes des biextensions envisagés dans 2.1 (cf. démonstration de 3.6.4).


\SGApage{I}{217}


\underline{BIBLIOGRAPHIE}

{[1]}\quad M. Demazure, J. Giraud et M. Raynaud, Schémas abéliens, Séminaire \\
\hspace*{4em}Fac. Sc. Orsay 1967/68 (à paraître).

{[2]}\quad J. Giraud, Cohomologie non abélienne, thèse Paris, 1967, (à paraître).

\SGAvar{{[3]}\quad A. Grothendieck, Sur quelques points d'Algèbre homologique ; Toh. \\
\hspace*{4em}Math. Journal, vol. 9 (1956), p.119-221

{[4]}\quad A. Grothendieck, Catégories cofibrées additives et Complexe cotangent \\
\hspace*{4em}relatif, Lectures Notes in Math. n$^{\circ}$79 1968 (Springer).

}{{[3]}\quad A. Grothendieck, Sur quelques points d'Algèbre homologique ; Toh. \\
\hspace*{4em}Math. Journal, vol. 9 (1957), p.119-221

{[4]}\quad A. Grothendieck, Catégories cofibrées additives et Complexe cotangent \\
\hspace*{4em}relatif, Lecture Notes in Mathematics n$^{\circ}$79 1968 (Springer).

}{F7I-VII-078}{[5]}\quad J.L. Verdier, Catégories dérivées des catégories abéliennes \\
\hspace*{4em}(à paraître).

{[6]}\quad D. Mumford, Biextensions of formal groups, Tata Institute of \\
\hspace*{4em}Fundamental Research 1968 (à paraître).
\selectlanguage{french}

\ifnum\SGAform=1
\clearpage\section*{Apparat critique}\small
Les lectures ci-dessous sont données dans le TeX exact des témoins, afin de préserver les distinctions de notation. Les indices décimaux renvoient aux différences individuelles du passage marqué.

\par\medskip\noindent\hypertarget{app-F7I-VII-001.1}{}\textbf{F7I-VII-001.1}\quad p. 135; printed misprint.\par
\noindent Imprimé : {\ttfamily\footnotesize E\allowbreak{}}\par
\noindent Corrigé : {\ttfamily\footnotesize \textbackslash{}underline\{T\}\allowbreak{}}\par
\noindent Source : 7I:265. \nolinkurl{fixes/MISPRINTS_7I.json#265}\par

\par\medskip\noindent\hypertarget{app-F7I-VII-001.2}{}\textbf{F7I-VII-001.2}\quad p. 135; misprint\_pass.\par
\noindent Imprimé : {\ttfamily\footnotesize -\allowbreak{}}\par
\noindent Corrigé : \emph{(1 espace(s), 0 retour(s) de ligne dans le témoin TeX)}\par
\noindent Source : 7I:266. \nolinkurl{fixes/FIXLOG_7I_CURRENT.jsonl#724}\par

\par\medskip\noindent\hypertarget{app-F7I-VII-001.3}{}\textbf{F7I-VII-001.3}\quad p. 136; printed misprint.\par
\noindent Imprimé : {\ttfamily\footnotesize P\allowbreak{}}\par
\noindent Corrigé : {\ttfamily\footnotesize E\allowbreak{}}\par
\noindent Source : 7I:285. \nolinkurl{fixes/MISPRINTS_7I.json#285}\par

\par\medskip\noindent\hypertarget{app-F7I-VII-001.4}{}\textbf{F7I-VII-001.4}\quad p. 136; printed misprint.\par
\noindent Imprimé : \emph{(vide)}\par
\noindent Corrigé : {\ttfamily\footnotesize \{\allowbreak{}}\par
\noindent Source : 7I:285. \nolinkurl{fixes/MISPRINTS_7I.json#285}\par

\par\medskip\noindent\hypertarget{app-F7I-VII-001.5}{}\textbf{F7I-VII-001.5}\quad p. 136; printed misprint.\par
\noindent Imprimé : \emph{(vide)}\par
\noindent Corrigé : {\ttfamily\footnotesize '\}\allowbreak{}}\par
\noindent Source : 7I:285. \nolinkurl{fixes/MISPRINTS_7I.json#285}\par

\par\medskip\noindent\hypertarget{app-F7I-VII-001.6}{}\textbf{F7I-VII-001.6}\quad p. 136; printed misprint.\par
\noindent Imprimé : {\ttfamily\footnotesize P\allowbreak{}}\par
\noindent Corrigé : {\ttfamily\footnotesize E\allowbreak{}}\par
\noindent Source : 7I:285. \nolinkurl{fixes/MISPRINTS_7I.json#285}\par

\par\medskip\noindent\hypertarget{app-F7I-VII-001.7}{}\textbf{F7I-VII-001.7}\quad p. 136; misprint\_pass.\par
\noindent Imprimé : {\ttfamily\footnotesize S\allowbreak{}}\par
\noindent Corrigé : {\ttfamily\footnotesize P\allowbreak{}}\par
\noindent Source : 7I:286. \nolinkurl{fixes/FIXLOG_7I_CURRENT.jsonl#725}; \nolinkurl{fixes/MISPRINTS_7I.json#286}\par

\par\medskip\noindent\hypertarget{app-F7I-VII-002.1}{}\textbf{F7I-VII-002.1}\quad p. 137; printed misprint.\par
\noindent Imprimé : {\ttfamily\footnotesize P\allowbreak{}}\par
\noindent Corrigé : {\ttfamily\footnotesize E\allowbreak{}}\par
\noindent Source : 7I:287. \nolinkurl{fixes/MISPRINTS_7I.json#287}\par

\par\medskip\noindent\hypertarget{app-F7I-VII-002.2}{}\textbf{F7I-VII-002.2}\quad p. 137; printed misprint.\par
\noindent Imprimé : {\ttfamily\footnotesize P\allowbreak{}}\par
\noindent Corrigé : {\ttfamily\footnotesize E\allowbreak{}}\par
\noindent Source : 7I:287. \nolinkurl{fixes/MISPRINTS_7I.json#287}\par

\par\medskip\noindent\hypertarget{app-F7I-VII-002.3}{}\textbf{F7I-VII-002.3}\quad p. 137; printed misprint.\par
\noindent Imprimé : {\ttfamily\footnotesize P\allowbreak{}}\par
\noindent Corrigé : {\ttfamily\footnotesize E\allowbreak{}}\par
\noindent Source : 7I:287. \nolinkurl{fixes/MISPRINTS_7I.json#287}\par

\par\medskip\noindent\hypertarget{app-F7I-VII-002.4}{}\textbf{F7I-VII-002.4}\quad p. 137; printed misprint.\par
\noindent Imprimé : {\ttfamily\footnotesize P\allowbreak{}}\par
\noindent Corrigé : {\ttfamily\footnotesize E\allowbreak{}}\par
\noindent Source : 7I:288. \nolinkurl{fixes/MISPRINTS_7I.json#288}\par

\par\medskip\noindent\hypertarget{app-F7I-VII-002.5}{}\textbf{F7I-VII-002.5}\quad p. 137; printed misprint.\par
\noindent Imprimé : {\ttfamily\footnotesize P\allowbreak{}}\par
\noindent Corrigé : {\ttfamily\footnotesize E\allowbreak{}}\par
\noindent Source : 7I:288. \nolinkurl{fixes/MISPRINTS_7I.json#288}\par

\par\medskip\noindent\hypertarget{app-F7I-VII-002.6}{}\textbf{F7I-VII-002.6}\quad p. 137; printed misprint.\par
\noindent Imprimé : {\ttfamily\footnotesize P\allowbreak{}}\par
\noindent Corrigé : {\ttfamily\footnotesize E\allowbreak{}}\par
\noindent Source : 7I:289. \nolinkurl{fixes/MISPRINTS_7I.json#289}\par

\par\medskip\noindent\hypertarget{app-F7I-VII-002.7}{}\textbf{F7I-VII-002.7}\quad p. 137; printed misprint.\par
\noindent Imprimé : {\ttfamily\footnotesize P\allowbreak{}}\par
\noindent Corrigé : {\ttfamily\footnotesize E\allowbreak{}}\par
\noindent Source : 7I:289. \nolinkurl{fixes/MISPRINTS_7I.json#289}\par

\par\medskip\noindent\hypertarget{app-F7I-VII-002.8}{}\textbf{F7I-VII-002.8}\quad p. 137; printed misprint.\par
\noindent Imprimé : {\ttfamily\footnotesize P\allowbreak{}}\par
\noindent Corrigé : {\ttfamily\footnotesize E\allowbreak{}}\par
\noindent Source : 7I:290. \nolinkurl{fixes/MISPRINTS_7I.json#290}\par

\par\medskip\noindent\hypertarget{app-F7I-VII-002.9}{}\textbf{F7I-VII-002.9}\quad p. 137; printed misprint.\par
\noindent Imprimé : {\ttfamily\footnotesize P\allowbreak{}}\par
\noindent Corrigé : {\ttfamily\footnotesize E\allowbreak{}}\par
\noindent Source : 7I:290. \nolinkurl{fixes/MISPRINTS_7I.json#290}\par

\par\medskip\noindent\hypertarget{app-F7I-VII-002.10}{}\textbf{F7I-VII-002.10}\quad p. 137; printed misprint.\par
\noindent Imprimé : {\ttfamily\footnotesize P\allowbreak{}}\par
\noindent Corrigé : {\ttfamily\footnotesize E\allowbreak{}}\par
\noindent Source : 7I:290. \nolinkurl{fixes/MISPRINTS_7I.json#290}\par

\par\medskip\noindent\hypertarget{app-F7I-VII-002.11}{}\textbf{F7I-VII-002.11}\quad p. 137; printed misprint.\par
\noindent Imprimé : {\ttfamily\footnotesize P\allowbreak{}}\par
\noindent Corrigé : {\ttfamily\footnotesize E\allowbreak{}}\par
\noindent Source : 7I:290. \nolinkurl{fixes/MISPRINTS_7I.json#290}\par

\par\medskip\noindent\hypertarget{app-F7I-VII-002.12}{}\textbf{F7I-VII-002.12}\quad p. 137; printed misprint.\par
\noindent Imprimé : {\ttfamily\footnotesize P\allowbreak{}}\par
\noindent Corrigé : {\ttfamily\footnotesize E\allowbreak{}}\par
\noindent Source : 7I:290. \nolinkurl{fixes/MISPRINTS_7I.json#290}\par

\par\medskip\noindent\hypertarget{app-F7I-VII-002.13}{}\textbf{F7I-VII-002.13}\quad p. 137; printed misprint.\par
\noindent Imprimé : {\ttfamily\footnotesize P\allowbreak{}}\par
\noindent Corrigé : {\ttfamily\footnotesize E\allowbreak{}}\par
\noindent Source : 7I:291. \nolinkurl{fixes/MISPRINTS_7I.json#291}\par

\par\medskip\noindent\hypertarget{app-F7I-VII-002.14}{}\textbf{F7I-VII-002.14}\quad p. 137; printed misprint.\par
\noindent Imprimé : {\ttfamily\footnotesize P\allowbreak{}}\par
\noindent Corrigé : {\ttfamily\footnotesize E\allowbreak{}}\par
\noindent Source : 7I:291. \nolinkurl{fixes/MISPRINTS_7I.json#291}\par

\par\medskip\noindent\hypertarget{app-F7I-VII-002.15}{}\textbf{F7I-VII-002.15}\quad p. 137; printed misprint.\par
\noindent Imprimé : {\ttfamily\footnotesize P\allowbreak{}}\par
\noindent Corrigé : {\ttfamily\footnotesize E\allowbreak{}}\par
\noindent Source : 7I:291. \nolinkurl{fixes/MISPRINTS_7I.json#291}\par

\par\medskip\noindent\hypertarget{app-F7I-VII-002.16}{}\textbf{F7I-VII-002.16}\quad p. 138; printed misprint.\par
\noindent Imprimé : {\ttfamily\footnotesize P\allowbreak{}}\par
\noindent Corrigé : {\ttfamily\footnotesize E\allowbreak{}}\par
\noindent Source : 7I:292. \nolinkurl{fixes/MISPRINTS_7I.json#292}\par

\par\medskip\noindent\hypertarget{app-F7I-VII-002.17}{}\textbf{F7I-VII-002.17}\quad p. 138; printed misprint.\par
\noindent Imprimé : {\ttfamily\footnotesize P\allowbreak{}}\par
\noindent Corrigé : {\ttfamily\footnotesize E\allowbreak{}}\par
\noindent Source : 7I:292. \nolinkurl{fixes/MISPRINTS_7I.json#292}\par

\par\medskip\noindent\hypertarget{app-F7I-VII-003.1}{}\textbf{F7I-VII-003.1}\quad p. 139; inherited corrected-layer difference.\par
\noindent Imprimé : {\ttfamily\footnotesize "\allowbreak{}}\par
\noindent Corrigé : {\ttfamily\footnotesize ''\allowbreak{}}\par
\noindent Source : ADOPT-F7I-VII-003.1. \nolinkurl{fr/sga7i/diplomatic/expose_VII_body.tex}; \nolinkurl{fr/sga7i/current/expose_VII_body.tex}\par

\par\medskip\noindent\hypertarget{app-F7I-VII-004.1}{}\textbf{F7I-VII-004.1}\quad p. 139; misprint\_pass.\par
\noindent Imprimé : {\ttfamily\footnotesize contractéede\allowbreak{}}\par
\noindent Corrigé : {\ttfamily\footnotesize contractée\allowbreak{} de\allowbreak{}}\par
\noindent Source : 7I:293. \nolinkurl{fixes/FIXLOG_7I_CURRENT.jsonl#726}; \nolinkurl{fixes/MISPRINTS_7I.json#293}\par

\par\medskip\noindent\hypertarget{app-F7I-VII-004.2}{}\textbf{F7I-VII-004.2}\quad p. 139; misprint\_pass.\par
\noindent Imprimé : {\ttfamily\footnotesize estassociative\allowbreak{}}\par
\noindent Corrigé : {\ttfamily\footnotesize est\allowbreak{} associative\allowbreak{}}\par
\noindent Source : 7I:294. \nolinkurl{fixes/FIXLOG_7I_CURRENT.jsonl#727}; \nolinkurl{fixes/MISPRINTS_7I.json#294}\par

\par\medskip\noindent\hypertarget{app-F7I-VII-004.3}{}\textbf{F7I-VII-004.3}\quad p. 140; designated-witness attribution.\par
\noindent Imprimé : {\ttfamily\footnotesize "\allowbreak{}}\par
\noindent Corrigé : {\ttfamily\footnotesize ''\allowbreak{}}\par
\noindent Source : WITNESS-F7I-VII-004.3. \nolinkurl{fr/sga7i/diplomatic/expose_VII_body.tex}; \nolinkurl{fr/sga7i/current/expose_VII_body.tex}\par

\par\medskip\noindent\hypertarget{app-F7I-VII-005.1}{}\textbf{F7I-VII-005.1}\quad p. 140; printed misprint.\par
\noindent Imprimé : {\ttfamily\footnotesize P\allowbreak{}}\par
\noindent Corrigé : {\ttfamily\footnotesize E\allowbreak{}}\par
\noindent Source : 7I:295. \nolinkurl{fixes/MISPRINTS_7I.json#295}\par

\par\medskip\noindent\hypertarget{app-F7I-VII-006.1}{}\textbf{F7I-VII-006.1}\quad p. 142; misprint\_pass.\par
\noindent Imprimé : {\ttfamily\footnotesize noté\allowbreak{}}\par
\noindent Corrigé : {\ttfamily\footnotesize notée\allowbreak{}}\par
\noindent Source : 7I:279. \nolinkurl{fixes/FIXLOG_7I_CURRENT.jsonl#728}; \nolinkurl{fixes/MISPRINTS_7I.json#279}\par

\par\medskip\noindent\hypertarget{app-F7I-VII-007.1}{}\textbf{F7I-VII-007.1}\quad p. 142; printed misprint.\par
\noindent Imprimé : {\ttfamily\footnotesize P\allowbreak{}}\par
\noindent Corrigé : {\ttfamily\footnotesize E\allowbreak{}}\par
\noindent Source : 7I:280. \nolinkurl{fixes/MISPRINTS_7I.json#280}\par

\par\medskip\noindent\hypertarget{app-F7I-VII-007.2}{}\textbf{F7I-VII-007.2}\quad p. 142; printed misprint.\par
\noindent Imprimé : {\ttfamily\footnotesize P\allowbreak{}}\par
\noindent Corrigé : {\ttfamily\footnotesize E\allowbreak{}}\par
\noindent Source : 7I:280. \nolinkurl{fixes/MISPRINTS_7I.json#280}\par

\par\medskip\noindent\hypertarget{app-F7I-VII-008.1}{}\textbf{F7I-VII-008.1}\quad p. 144; printed misprint.\par
\noindent Imprimé : {\ttfamily\footnotesize P\allowbreak{}}\par
\noindent Corrigé : {\ttfamily\footnotesize E\allowbreak{}}\par
\noindent Source : 7I:281. \nolinkurl{fixes/MISPRINTS_7I.json#281}\par

\par\medskip\noindent\hypertarget{app-F7I-VII-008.2}{}\textbf{F7I-VII-008.2}\quad p. 144; misprint\_pass.\par
\noindent Imprimé : {\ttfamily\footnotesize coîncident\allowbreak{}}\par
\noindent Corrigé : {\ttfamily\footnotesize coïncident\allowbreak{}}\par
\noindent Source : 7I:282. \nolinkurl{fixes/FIXLOG_7I_CURRENT.jsonl#729}; \nolinkurl{fixes/MISPRINTS_7I.json#282}\par

\par\medskip\noindent\hypertarget{app-F7I-VII-009.1}{}\textbf{F7I-VII-009.1}\quad p. 145; misprint\_in\_print.\par
\noindent Imprimé : {\ttfamily\footnotesize coîncident\allowbreak{}}\par
\noindent Corrigé : {\ttfamily\footnotesize coïncident\allowbreak{}}\par
\noindent Source : FIX-7I-D-0204. \nolinkurl{fixes/FIXLOG_7I.jsonl#204}\par

\par\medskip\noindent\hypertarget{app-F7I-VII-009.2}{}\textbf{F7I-VII-009.2}\quad p. 146; inherited corrected-layer difference.\par
\noindent Imprimé : \emph{(vide)}\par
\noindent Corrigé : {\ttfamily\footnotesize \textasciicircum{}*\allowbreak{}}\par
\noindent Source : ADOPT-F7I-VII-009.2. \nolinkurl{fr/sga7i/diplomatic/expose_VII_body.tex}; \nolinkurl{fr/sga7i/current/expose_VII_body.tex}\par

\par\medskip\noindent\hypertarget{app-F7I-VII-009.3}{}\textbf{F7I-VII-009.3}\quad p. 146; designated-witness attribution.\par
\noindent Imprimé : \emph{(vide)}\par
\noindent Corrigé : {\ttfamily\footnotesize \textasciicircum{}*\allowbreak{}}\par
\noindent Source : WITNESS-F7I-VII-009.3. \nolinkurl{fr/sga7i/diplomatic/expose_VII_body.tex}; \nolinkurl{fr/sga7i/current/expose_VII_body.tex}\par

\par\medskip\noindent\hypertarget{app-F7I-VII-009.4}{}\textbf{F7I-VII-009.4}\quad p. 146; printed misprint.\par
\noindent Imprimé : \emph{(vide)}\par
\noindent Corrigé : {\ttfamily\footnotesize )\allowbreak{}}\par
\noindent Source : 7I:296. \nolinkurl{fixes/MISPRINTS_7I.json#296}\par

\par\medskip\noindent\hypertarget{app-F7I-VII-010.1}{}\textbf{F7I-VII-010.1}\quad p. 148; misprint\_pass.\par
\noindent Imprimé : {\ttfamily\footnotesize la\allowbreak{}}\par
\noindent Corrigé : {\ttfamily\footnotesize le\allowbreak{}}\par
\noindent Source : 7I:297. \nolinkurl{fixes/FIXLOG_7I_CURRENT.jsonl#730}; \nolinkurl{fixes/MISPRINTS_7I.json#297}\par

\par\medskip\noindent\hypertarget{app-F7I-VII-010.2}{}\textbf{F7I-VII-010.2}\quad p. 148; printed misprint.\par
\noindent Imprimé : {\ttfamily\footnotesize Q\allowbreak{}}\par
\noindent Corrigé : {\ttfamily\footnotesize G\allowbreak{}}\par
\noindent Source : 7I:298. \nolinkurl{fixes/MISPRINTS_7I.json#298}\par

\par\medskip\noindent\hypertarget{app-F7I-VII-010.3}{}\textbf{F7I-VII-010.3}\quad p. 148; misprint\_pass.\par
\noindent Imprimé : {\ttfamily\footnotesize biridifié\allowbreak{}}\par
\noindent Corrigé : {\ttfamily\footnotesize rigidifié\allowbreak{}}\par
\noindent Source : 7I:299. \nolinkurl{fixes/FIXLOG_7I_CURRENT.jsonl#731}; \nolinkurl{fixes/MISPRINTS_7I.json#299}\par

\par\medskip\noindent\hypertarget{app-F7I-VII-010.4}{}\textbf{F7I-VII-010.4}\quad p. 148; printed misprint.\par
\noindent Imprimé : {\ttfamily\footnotesize torseur\allowbreak{}}\par
\noindent Corrigé : {\ttfamily\footnotesize bitorseur\allowbreak{}}\par
\noindent Source : 7I:300. \nolinkurl{fixes/MISPRINTS_7I.json#300}\par

\par\medskip\noindent\hypertarget{app-F7I-VII-010.5}{}\textbf{F7I-VII-010.5}\quad p. 148; printed misprint.\par
\noindent Imprimé : \emph{(vide)}\par
\noindent Corrigé : {\ttfamily\footnotesize (\allowbreak{}}\par
\noindent Source : 7I:300. \nolinkurl{fixes/MISPRINTS_7I.json#300}\par

\par\medskip\noindent\hypertarget{app-F7I-VII-010.6}{}\textbf{F7I-VII-010.6}\quad p. 148; printed misprint.\par
\noindent Imprimé : \emph{(vide)}\par
\noindent Corrigé : {\ttfamily\footnotesize ,G\_P)\allowbreak{}}\par
\noindent Source : 7I:300. \nolinkurl{fixes/MISPRINTS_7I.json#300}\par

\par\medskip\noindent\hypertarget{app-F7I-VII-011.1}{}\textbf{F7I-VII-011.1}\quad p. 149; misprint\_pass.\par
\noindent Imprimé : {\ttfamily\footnotesize f(\allowbreak{}}\par
\noindent Corrigé : \emph{(vide)}\par
\noindent Source : 7I:301. \nolinkurl{fixes/FIXLOG_7I_CURRENT.jsonl#732}; \nolinkurl{fixes/MISPRINTS_7I.json#301}\par

\par\medskip\noindent\hypertarget{app-F7I-VII-011.2}{}\textbf{F7I-VII-011.2}\quad p. 149; misprint\_pass.\par
\noindent Imprimé : {\ttfamily\footnotesize )\allowbreak{}}\par
\noindent Corrigé : \emph{(vide)}\par
\noindent Source : 7I:301. \nolinkurl{fixes/FIXLOG_7I_CURRENT.jsonl#732}; \nolinkurl{fixes/MISPRINTS_7I.json#301}\par

\par\medskip\noindent\hypertarget{app-F7I-VII-011.3}{}\textbf{F7I-VII-011.3}\quad p. 149; misprint\_pass.\par
\noindent Imprimé : {\ttfamily\footnotesize f(\allowbreak{}}\par
\noindent Corrigé : \emph{(vide)}\par
\noindent Source : 7I:301. \nolinkurl{fixes/FIXLOG_7I_CURRENT.jsonl#732}; \nolinkurl{fixes/MISPRINTS_7I.json#301}\par

\par\medskip\noindent\hypertarget{app-F7I-VII-011.4}{}\textbf{F7I-VII-011.4}\quad p. 149; misprint\_pass.\par
\noindent Imprimé : {\ttfamily\footnotesize f\allowbreak{}}\par
\noindent Corrigé : {\ttfamily\footnotesize j\allowbreak{}}\par
\noindent Source : 7I:301. \nolinkurl{fixes/FIXLOG_7I_CURRENT.jsonl#732}; \nolinkurl{fixes/MISPRINTS_7I.json#301}\par

\par\medskip\noindent\hypertarget{app-F7I-VII-011.5}{}\textbf{F7I-VII-011.5}\quad p. 149; designated-witness attribution.\par
\noindent Imprimé : {\ttfamily\footnotesize )\allowbreak{}}\par
\noindent Corrigé : \emph{(vide)}\par
\noindent Source : WITNESS-F7I-VII-011.5. \nolinkurl{fr/sga7i/diplomatic/expose_VII_body.tex}; \nolinkurl{fr/sga7i/current/expose_VII_body.tex}\par

\par\medskip\noindent\hypertarget{app-F7I-VII-012.1}{}\textbf{F7I-VII-012.1}\quad p. 150; inherited corrected-layer difference.\par
\noindent Imprimé : \emph{(vide)}\par
\noindent Corrigé : {\ttfamily\footnotesize \textasciicircum{}*\allowbreak{}}\par
\noindent Source : ADOPT-F7I-VII-012.1. \nolinkurl{fr/sga7i/diplomatic/expose_VII_body.tex}; \nolinkurl{fr/sga7i/current/expose_VII_body.tex}\par

\par\medskip\noindent\hypertarget{app-F7I-VII-013.1}{}\textbf{F7I-VII-013.1}\quad p. 150; misprint\_pass.\par
\noindent Imprimé : {\ttfamily\footnotesize différent\allowbreak{}}\par
\noindent Corrigé : {\ttfamily\footnotesize diffèrent\allowbreak{}}\par
\noindent Source : 7I:303. \nolinkurl{fixes/FIXLOG_7I_CURRENT.jsonl#733}; \nolinkurl{fixes/MISPRINTS_7I.json#303}\par

\par\medskip\noindent\hypertarget{app-F7I-VII-014.1}{}\textbf{F7I-VII-014.1}\quad p. 151; inherited corrected-layer difference.\par
\noindent Imprimé : {\ttfamily\footnotesize y=\allowbreak{}}\par
\noindent Corrigé : \emph{(vide)}\par
\noindent Source : ADOPT-F7I-VII-014.1. \nolinkurl{fr/sga7i/diplomatic/expose_VII_body.tex}; \nolinkurl{fr/sga7i/current/expose_VII_body.tex}\par

\par\medskip\noindent\hypertarget{app-F7I-VII-014.2}{}\textbf{F7I-VII-014.2}\quad p. 151; inherited corrected-layer difference.\par
\noindent Imprimé : {\ttfamily\footnotesize =1\$,\allowbreak{} puis\allowbreak{} \$z=x\allowbreak{}}\par
\noindent Corrigé : \emph{(vide)}\par
\noindent Source : ADOPT-F7I-VII-014.2. \nolinkurl{fr/sga7i/diplomatic/expose_VII_body.tex}; \nolinkurl{fr/sga7i/current/expose_VII_body.tex}\par

\par\medskip\noindent\hypertarget{app-F7I-VII-014.3}{}\textbf{F7I-VII-014.3}\quad p. 151; inherited corrected-layer difference.\par
\noindent Imprimé : \emph{(vide)}\par
\noindent Corrigé : {\ttfamily\footnotesize 1\$,\allowbreak{} puis\allowbreak{} \$\allowbreak{}}\par
\noindent Source : ADOPT-F7I-VII-014.3. \nolinkurl{fr/sga7i/diplomatic/expose_VII_body.tex}; \nolinkurl{fr/sga7i/current/expose_VII_body.tex}\par

\par\medskip\noindent\hypertarget{app-F7I-VII-014.4}{}\textbf{F7I-VII-014.4}\quad p. 151; inherited corrected-layer difference.\par
\noindent Imprimé : {\ttfamily\footnotesize e\_\allowbreak{}}\par
\noindent Corrigé : \emph{(vide)}\par
\noindent Source : ADOPT-F7I-VII-014.4. \nolinkurl{fr/sga7i/diplomatic/expose_VII_body.tex}; \nolinkurl{fr/sga7i/current/expose_VII_body.tex}\par

\par\medskip\noindent\hypertarget{app-F7I-VII-014.5}{}\textbf{F7I-VII-014.5}\quad p. 151; inherited corrected-layer difference.\par
\noindent Imprimé : {\ttfamily\footnotesize e\_\allowbreak{}}\par
\noindent Corrigé : \emph{(vide)}\par
\noindent Source : ADOPT-F7I-VII-014.5. \nolinkurl{fr/sga7i/diplomatic/expose_VII_body.tex}; \nolinkurl{fr/sga7i/current/expose_VII_body.tex}\par

\par\medskip\noindent\hypertarget{app-F7I-VII-014.6}{}\textbf{F7I-VII-014.6}\quad p. 151; inherited corrected-layer difference.\par
\noindent Imprimé : {\ttfamily\footnotesize e\_\allowbreak{}}\par
\noindent Corrigé : \emph{(vide)}\par
\noindent Source : ADOPT-F7I-VII-014.6. \nolinkurl{fr/sga7i/diplomatic/expose_VII_body.tex}; \nolinkurl{fr/sga7i/current/expose_VII_body.tex}\par

\par\medskip\noindent\hypertarget{app-F7I-VII-015.1}{}\textbf{F7I-VII-015.1}\quad p. 151; printed misprint.\par
\noindent Imprimé : {\ttfamily\footnotesize G\allowbreak{}}\par
\noindent Corrigé : {\ttfamily\footnotesize E\allowbreak{}}\par
\noindent Source : 7I:304. \nolinkurl{fixes/MISPRINTS_7I.json#304}\par

\par\medskip\noindent\hypertarget{app-F7I-VII-017.1}{}\textbf{F7I-VII-017.1}\quad p. 153; inherited corrected-layer difference.\par
\noindent Imprimé : \emph{(vide)}\par
\noindent Corrigé : {\ttfamily\footnotesize \textbackslash{}operatorname\{\allowbreak{}}\par
\noindent Source : ADOPT-F7I-VII-017.1. \nolinkurl{fr/sga7i/diplomatic/expose_VII_body.tex}; \nolinkurl{fr/sga7i/current/expose_VII_body.tex}\par

\par\medskip\noindent\hypertarget{app-F7I-VII-017.2}{}\textbf{F7I-VII-017.2}\quad p. 153; inherited corrected-layer difference.\par
\noindent Imprimé : \emph{(vide)}\par
\noindent Corrigé : {\ttfamily\footnotesize \}\allowbreak{}}\par
\noindent Source : ADOPT-F7I-VII-017.2. \nolinkurl{fr/sga7i/diplomatic/expose_VII_body.tex}; \nolinkurl{fr/sga7i/current/expose_VII_body.tex}\par

\par\medskip\noindent\hypertarget{app-F7I-VII-017.3}{}\textbf{F7I-VII-017.3}\quad p. 153; inherited corrected-layer difference.\par
\noindent Imprimé : \emph{(vide)}\par
\noindent Corrigé : {\ttfamily\footnotesize \textbackslash{}underline\{\allowbreak{}}\par
\noindent Source : ADOPT-F7I-VII-017.3. \nolinkurl{fr/sga7i/diplomatic/expose_VII_body.tex}; \nolinkurl{fr/sga7i/current/expose_VII_body.tex}\par

\par\medskip\noindent\hypertarget{app-F7I-VII-017.4}{}\textbf{F7I-VII-017.4}\quad p. 153; inherited corrected-layer difference.\par
\noindent Imprimé : \emph{(vide)}\par
\noindent Corrigé : {\ttfamily\footnotesize \}\allowbreak{}}\par
\noindent Source : ADOPT-F7I-VII-017.4. \nolinkurl{fr/sga7i/diplomatic/expose_VII_body.tex}; \nolinkurl{fr/sga7i/current/expose_VII_body.tex}\par

\par\medskip\noindent\hypertarget{app-F7I-VII-018.1}{}\textbf{F7I-VII-018.1}\quad p. 154; printed misprint.\par
\noindent Imprimé : \emph{(vide)}\par
\noindent Corrigé : {\ttfamily\footnotesize 1.\allowbreak{}}\par
\noindent Source : 7I:305. \nolinkurl{fixes/MISPRINTS_7I.json#305}\par

\par\medskip\noindent\hypertarget{app-F7I-VII-018.2}{}\textbf{F7I-VII-018.2}\quad p. 154; printed misprint.\par
\noindent Imprimé : \emph{(vide)}\par
\noindent Corrigé : {\ttfamily\footnotesize 1.\allowbreak{}}\par
\noindent Source : 7I:305. \nolinkurl{fixes/MISPRINTS_7I.json#305}\par

\par\medskip\noindent\hypertarget{app-F7I-VII-019.1}{}\textbf{F7I-VII-019.1}\quad p. 155; misprint\_in\_print.\par
\noindent Imprimé : {\ttfamily\footnotesize coîncident\allowbreak{}}\par
\noindent Corrigé : {\ttfamily\footnotesize coïncident\allowbreak{}}\par
\noindent Source : FIX-7I-D-0224. \nolinkurl{fixes/FIXLOG_7I.jsonl#224}\par

\par\medskip\noindent\hypertarget{app-F7I-VII-020.1}{}\textbf{F7I-VII-020.1}\quad p. 156; misprint\_pass.\par
\noindent Imprimé : {\ttfamily\footnotesize partiquement\allowbreak{}}\par
\noindent Corrigé : {\ttfamily\footnotesize pratiquement\allowbreak{}}\par
\noindent Source : 7I:328. \nolinkurl{fixes/FIXLOG_7I_CURRENT.jsonl#734}; \nolinkurl{fixes/MISPRINTS_7I.json#328}\par

\par\medskip\noindent\hypertarget{app-F7I-VII-021.1}{}\textbf{F7I-VII-021.1}\quad p. 157; inherited corrected-layer difference.\par
\noindent Imprimé : \emph{(vide)}\par
\noindent Corrigé : {\ttfamily\footnotesize P\_\allowbreak{}}\par
\noindent Source : ADOPT-F7I-VII-021.1. \nolinkurl{fr/sga7i/diplomatic/expose_VII_body.tex}; \nolinkurl{fr/sga7i/current/expose_VII_body.tex}\par

\par\medskip\noindent\hypertarget{app-F7I-VII-021.2}{}\textbf{F7I-VII-021.2}\quad p. 157; inherited corrected-layer difference.\par
\noindent Imprimé : {\ttfamily\footnotesize \_P\allowbreak{}}\par
\noindent Corrigé : \emph{(vide)}\par
\noindent Source : ADOPT-F7I-VII-021.2. \nolinkurl{fr/sga7i/diplomatic/expose_VII_body.tex}; \nolinkurl{fr/sga7i/current/expose_VII_body.tex}\par

\par\medskip\noindent\hypertarget{app-F7I-VII-021.3}{}\textbf{F7I-VII-021.3}\quad p. 157; inherited corrected-layer difference.\par
\noindent Imprimé : {\ttfamily\footnotesize P\allowbreak{}}\par
\noindent Corrigé : {\ttfamily\footnotesize Q\allowbreak{}}\par
\noindent Source : ADOPT-F7I-VII-021.3. \nolinkurl{fr/sga7i/diplomatic/expose_VII_body.tex}; \nolinkurl{fr/sga7i/current/expose_VII_body.tex}\par

\par\medskip\noindent\hypertarget{app-F7I-VII-022.1}{}\textbf{F7I-VII-022.1}\quad p. 157; misprint\_in\_print.\par
\noindent Imprimé : \emph{(vide)}\par
\noindent Corrigé : {\ttfamily\footnotesize '\allowbreak{}}\par
\noindent Source : FIX-7I-D-0230. \nolinkurl{fixes/FIXLOG_7I.jsonl#230}\par

\par\medskip\noindent\hypertarget{app-F7I-VII-023.1}{}\textbf{F7I-VII-023.1}\quad p. 158; inherited corrected-layer difference.\par
\noindent Imprimé : \emph{(vide)}\par
\noindent Corrigé : {\ttfamily\footnotesize Q\_\allowbreak{}}\par
\noindent Source : ADOPT-F7I-VII-023.1. \nolinkurl{fr/sga7i/diplomatic/expose_VII_body.tex}; \nolinkurl{fr/sga7i/current/expose_VII_body.tex}\par

\par\medskip\noindent\hypertarget{app-F7I-VII-023.2}{}\textbf{F7I-VII-023.2}\quad p. 158; inherited corrected-layer difference.\par
\noindent Imprimé : {\ttfamily\footnotesize \_Q\allowbreak{}}\par
\noindent Corrigé : \emph{(vide)}\par
\noindent Source : ADOPT-F7I-VII-023.2. \nolinkurl{fr/sga7i/diplomatic/expose_VII_body.tex}; \nolinkurl{fr/sga7i/current/expose_VII_body.tex}\par

\par\medskip\noindent\hypertarget{app-F7I-VII-023.3}{}\textbf{F7I-VII-023.3}\quad p. 158; inherited corrected-layer difference.\par
\noindent Imprimé : {\ttfamily\footnotesize Q\allowbreak{}}\par
\noindent Corrigé : {\ttfamily\footnotesize P\allowbreak{}}\par
\noindent Source : ADOPT-F7I-VII-023.3. \nolinkurl{fr/sga7i/diplomatic/expose_VII_body.tex}; \nolinkurl{fr/sga7i/current/expose_VII_body.tex}\par

\par\medskip\noindent\hypertarget{app-F7I-VII-023.4}{}\textbf{F7I-VII-023.4}\quad p. 158; printed misprint.\par
\noindent Imprimé : {\ttfamily\footnotesize pp\allowbreak{}}\par
\noindent Corrigé : {\ttfamily\footnotesize p,qq\allowbreak{}}\par
\noindent Source : 7I:330. \nolinkurl{fixes/MISPRINTS_7I.json#330}\par

\par\medskip\noindent\hypertarget{app-F7I-VII-023.5}{}\textbf{F7I-VII-023.5}\quad p. 158; printed misprint.\par
\noindent Imprimé : {\ttfamily\footnotesize ,q\allowbreak{}}\par
\noindent Corrigé : \emph{(vide)}\par
\noindent Source : 7I:330. \nolinkurl{fixes/MISPRINTS_7I.json#330}\par

\par\medskip\noindent\hypertarget{app-F7I-VII-023.6}{}\textbf{F7I-VII-023.6}\quad p. 158; printed misprint.\par
\noindent Imprimé : {\ttfamily\footnotesize ;\allowbreak{}}\par
\noindent Corrigé : {\ttfamily\footnotesize ,\allowbreak{}}\par
\noindent Source : 7I:331. \nolinkurl{fixes/MISPRINTS_7I.json#331}\par

\par\medskip\noindent\hypertarget{app-F7I-VII-024.1}{}\textbf{F7I-VII-024.1}\quad p. 159; designated-witness attribution.\par
\noindent Imprimé : \emph{(vide)}\par
\noindent Corrigé : {\ttfamily\footnotesize P\_\allowbreak{}}\par
\noindent Source : WITNESS-F7I-VII-024.1. \nolinkurl{fr/sga7i/diplomatic/expose_VII_body.tex}; \nolinkurl{fr/sga7i/current/expose_VII_body.tex}\par

\par\medskip\noindent\hypertarget{app-F7I-VII-024.2}{}\textbf{F7I-VII-024.2}\quad p. 159; designated-witness attribution.\par
\noindent Imprimé : {\ttfamily\footnotesize \_P\allowbreak{}}\par
\noindent Corrigé : \emph{(vide)}\par
\noindent Source : WITNESS-F7I-VII-024.2. \nolinkurl{fr/sga7i/diplomatic/expose_VII_body.tex}; \nolinkurl{fr/sga7i/current/expose_VII_body.tex}\par

\par\medskip\noindent\hypertarget{app-F7I-VII-024.3}{}\textbf{F7I-VII-024.3}\quad p. 159; designated-witness attribution.\par
\noindent Imprimé : {\ttfamily\footnotesize P\allowbreak{}}\par
\noindent Corrigé : {\ttfamily\footnotesize Q\allowbreak{}}\par
\noindent Source : WITNESS-F7I-VII-024.3. \nolinkurl{fr/sga7i/diplomatic/expose_VII_body.tex}; \nolinkurl{fr/sga7i/current/expose_VII_body.tex}\par

\par\medskip\noindent\hypertarget{app-F7I-VII-024.4}{}\textbf{F7I-VII-024.4}\quad p. 159; misprint\_pass.\par
\noindent Imprimé : {\ttfamily\footnotesize défini\allowbreak{}}\par
\noindent Corrigé : {\ttfamily\footnotesize définie\allowbreak{}}\par
\noindent Source : 7I:332. \nolinkurl{fixes/FIXLOG_7I_CURRENT.jsonl#735}; \nolinkurl{fixes/MISPRINTS_7I.json#332}\par

\par\medskip\noindent\hypertarget{app-F7I-VII-024.5}{}\textbf{F7I-VII-024.5}\quad p. 159; designated-witness attribution.\par
\noindent Imprimé : \emph{(vide)}\par
\noindent Corrigé : {\ttfamily\footnotesize Q\_\allowbreak{}}\par
\noindent Source : WITNESS-F7I-VII-024.5. \nolinkurl{fr/sga7i/diplomatic/expose_VII_body.tex}; \nolinkurl{fr/sga7i/current/expose_VII_body.tex}\par

\par\medskip\noindent\hypertarget{app-F7I-VII-024.6}{}\textbf{F7I-VII-024.6}\quad p. 159; designated-witness attribution.\par
\noindent Imprimé : {\ttfamily\footnotesize \_Q\allowbreak{}}\par
\noindent Corrigé : \emph{(vide)}\par
\noindent Source : WITNESS-F7I-VII-024.6. \nolinkurl{fr/sga7i/diplomatic/expose_VII_body.tex}; \nolinkurl{fr/sga7i/current/expose_VII_body.tex}\par

\par\medskip\noindent\hypertarget{app-F7I-VII-024.7}{}\textbf{F7I-VII-024.7}\quad p. 159; designated-witness attribution.\par
\noindent Imprimé : {\ttfamily\footnotesize Q\allowbreak{}}\par
\noindent Corrigé : {\ttfamily\footnotesize P\allowbreak{}}\par
\noindent Source : WITNESS-F7I-VII-024.7. \nolinkurl{fr/sga7i/diplomatic/expose_VII_body.tex}; \nolinkurl{fr/sga7i/current/expose_VII_body.tex}\par

\par\medskip\noindent\hypertarget{app-F7I-VII-024.8}{}\textbf{F7I-VII-024.8}\quad p. 159; misprint\_pass.\par
\noindent Imprimé : {\ttfamily\footnotesize défini\allowbreak{}}\par
\noindent Corrigé : {\ttfamily\footnotesize définie\allowbreak{}}\par
\noindent Source : 7I:333. \nolinkurl{fixes/FIXLOG_7I_CURRENT.jsonl#736}; \nolinkurl{fixes/MISPRINTS_7I.json#333}\par

\par\medskip\noindent\hypertarget{app-F7I-VII-025.1}{}\textbf{F7I-VII-025.1}\quad p. 159; designated-witness attribution.\par
\noindent Imprimé : \emph{(vide)}\par
\noindent Corrigé : {\ttfamily\footnotesize P\_\allowbreak{}}\par
\noindent Source : WITNESS-F7I-VII-025.1. \nolinkurl{fr/sga7i/diplomatic/expose_VII_body.tex}; \nolinkurl{fr/sga7i/current/expose_VII_body.tex}\par

\par\medskip\noindent\hypertarget{app-F7I-VII-025.2}{}\textbf{F7I-VII-025.2}\quad p. 159; designated-witness attribution.\par
\noindent Imprimé : {\ttfamily\footnotesize \_P\allowbreak{}}\par
\noindent Corrigé : \emph{(vide)}\par
\noindent Source : WITNESS-F7I-VII-025.2. \nolinkurl{fr/sga7i/diplomatic/expose_VII_body.tex}; \nolinkurl{fr/sga7i/current/expose_VII_body.tex}\par

\par\medskip\noindent\hypertarget{app-F7I-VII-025.3}{}\textbf{F7I-VII-025.3}\quad p. 159; designated-witness attribution.\par
\noindent Imprimé : {\ttfamily\footnotesize P\allowbreak{}}\par
\noindent Corrigé : {\ttfamily\footnotesize Q\allowbreak{}}\par
\noindent Source : WITNESS-F7I-VII-025.3. \nolinkurl{fr/sga7i/diplomatic/expose_VII_body.tex}; \nolinkurl{fr/sga7i/current/expose_VII_body.tex}\par

\par\medskip\noindent\hypertarget{app-F7I-VII-025.4}{}\textbf{F7I-VII-025.4}\quad p. 159; designated-witness attribution.\par
\noindent Imprimé : {\ttfamily\footnotesize G\allowbreak{}}\par
\noindent Corrigé : {\ttfamily\footnotesize Q\allowbreak{}}\par
\noindent Source : WITNESS-F7I-VII-025.4. \nolinkurl{fr/sga7i/diplomatic/expose_VII_body.tex}; \nolinkurl{fr/sga7i/current/expose_VII_body.tex}\par

\par\medskip\noindent\hypertarget{app-F7I-VII-025.5}{}\textbf{F7I-VII-025.5}\quad p. 159; designated-witness attribution.\par
\noindent Imprimé : {\ttfamily\footnotesize Q\allowbreak{}}\par
\noindent Corrigé : {\ttfamily\footnotesize P\allowbreak{}}\par
\noindent Source : WITNESS-F7I-VII-025.5. \nolinkurl{fr/sga7i/diplomatic/expose_VII_body.tex}; \nolinkurl{fr/sga7i/current/expose_VII_body.tex}\par

\par\medskip\noindent\hypertarget{app-F7I-VII-025.6}{}\textbf{F7I-VII-025.6}\quad p. 159; designated-witness attribution.\par
\noindent Imprimé : \emph{(vide)}\par
\noindent Corrigé : {\ttfamily\footnotesize G\_\allowbreak{}}\par
\noindent Source : WITNESS-F7I-VII-025.6. \nolinkurl{fr/sga7i/diplomatic/expose_VII_body.tex}; \nolinkurl{fr/sga7i/current/expose_VII_body.tex}\par

\par\medskip\noindent\hypertarget{app-F7I-VII-025.7}{}\textbf{F7I-VII-025.7}\quad p. 159; designated-witness attribution.\par
\noindent Imprimé : {\ttfamily\footnotesize \_Q\allowbreak{}}\par
\noindent Corrigé : \emph{(vide)}\par
\noindent Source : WITNESS-F7I-VII-025.7. \nolinkurl{fr/sga7i/diplomatic/expose_VII_body.tex}; \nolinkurl{fr/sga7i/current/expose_VII_body.tex}\par

\par\medskip\noindent\hypertarget{app-F7I-VII-026.1}{}\textbf{F7I-VII-026.1}\quad p. 160; misprint\_pass.\par
\noindent Imprimé : {\ttfamily\footnotesize exprimant\allowbreak{}}\par
\noindent Corrigé : {\ttfamily\footnotesize expriment\allowbreak{}}\par
\noindent Source : 7I:335. \nolinkurl{fixes/FIXLOG_7I_CURRENT.jsonl#737}; \nolinkurl{fixes/MISPRINTS_7I.json#335}\par

\par\medskip\noindent\hypertarget{app-F7I-VII-027.1}{}\textbf{F7I-VII-027.1}\quad p. 160; printed misprint.\par
\noindent Imprimé : {\ttfamily\footnotesize G\allowbreak{}}\par
\noindent Corrigé : {\ttfamily\footnotesize Q\allowbreak{}}\par
\noindent Source : 7I:337. \nolinkurl{fixes/MISPRINTS_7I.json#337}\par

\par\medskip\noindent\hypertarget{app-F7I-VII-027.2}{}\textbf{F7I-VII-027.2}\quad p. 160; inherited corrected-layer difference.\par
\noindent Imprimé : {\ttfamily\footnotesize 1\_\allowbreak{}}\par
\noindent Corrigé : \emph{(vide)}\par
\noindent Source : ADOPT-F7I-VII-027.2. \nolinkurl{fr/sga7i/diplomatic/expose_VII_body.tex}; \nolinkurl{fr/sga7i/current/expose_VII_body.tex}\par

\par\medskip\noindent\hypertarget{app-F7I-VII-027.3}{}\textbf{F7I-VII-027.3}\quad p. 160; inherited corrected-layer difference.\par
\noindent Imprimé : \emph{(vide)}\par
\noindent Corrigé : {\ttfamily\footnotesize 1\_\allowbreak{}}\par
\noindent Source : ADOPT-F7I-VII-027.3. \nolinkurl{fr/sga7i/diplomatic/expose_VII_body.tex}; \nolinkurl{fr/sga7i/current/expose_VII_body.tex}\par

\par\medskip\noindent\hypertarget{app-F7I-VII-027.4}{}\textbf{F7I-VII-027.4}\quad p. 160; designated-witness attribution.\par
\noindent Imprimé : \emph{(vide)}\par
\noindent Corrigé : {\ttfamily\footnotesize 1\_\allowbreak{}}\par
\noindent Source : WITNESS-F7I-VII-027.4. \nolinkurl{fr/sga7i/diplomatic/expose_VII_body.tex}; \nolinkurl{fr/sga7i/current/expose_VII_body.tex}\par

\par\medskip\noindent\hypertarget{app-F7I-VII-027.5}{}\textbf{F7I-VII-027.5}\quad p. 160; designated-witness attribution.\par
\noindent Imprimé : {\ttfamily\footnotesize 1\_\allowbreak{}}\par
\noindent Corrigé : \emph{(vide)}\par
\noindent Source : WITNESS-F7I-VII-027.5. \nolinkurl{fr/sga7i/diplomatic/expose_VII_body.tex}; \nolinkurl{fr/sga7i/current/expose_VII_body.tex}\par

\par\medskip\noindent\hypertarget{app-F7I-VII-028.1}{}\textbf{F7I-VII-028.1}\quad p. 161; misprint\_pass.\par
\noindent Imprimé : {\ttfamily\footnotesize appelera\allowbreak{}}\par
\noindent Corrigé : {\ttfamily\footnotesize appellera\allowbreak{}}\par
\noindent Source : 7I:307. \nolinkurl{fixes/FIXLOG_7I_CURRENT.jsonl#738}; \nolinkurl{fixes/MISPRINTS_7I.json#307}\par

\par\medskip\noindent\hypertarget{app-F7I-VII-028.2}{}\textbf{F7I-VII-028.2}\quad p. 161; misprint\_pass.\par
\noindent Imprimé : {\ttfamily\footnotesize compositions\allowbreak{}}\par
\noindent Corrigé : {\ttfamily\footnotesize composition\allowbreak{}}\par
\noindent Source : 7I:308. \nolinkurl{fixes/FIXLOG_7I_CURRENT.jsonl#739}\par

\par\medskip\noindent\hypertarget{app-F7I-VII-028.3}{}\textbf{F7I-VII-028.3}\quad p. 161; misprint\_pass.\par
\noindent Imprimé : {\ttfamily\footnotesize des\allowbreak{}}\par
\noindent Corrigé : {\ttfamily\footnotesize les\allowbreak{}}\par
\noindent Source : 7I:309. \nolinkurl{fixes/FIXLOG_7I_CURRENT.jsonl#740}; \nolinkurl{fixes/MISPRINTS_7I.json#309}\par

\par\medskip\noindent\hypertarget{app-F7I-VII-028.4}{}\textbf{F7I-VII-028.4}\quad p. 161; printed misprint.\par
\noindent Imprimé : {\ttfamily\footnotesize G\allowbreak{}}\par
\noindent Corrigé : {\ttfamily\footnotesize E\allowbreak{}}\par
\noindent Source : 7I:310. \nolinkurl{fixes/MISPRINTS_7I.json#310}\par

\par\medskip\noindent\hypertarget{app-F7I-VII-029.1}{}\textbf{F7I-VII-029.1}\quad p. 163; misprint\_pass.\par
\noindent Imprimé : {\ttfamily\footnotesize en\allowbreak{}}\par
\noindent Corrigé : {\ttfamily\footnotesize est\allowbreak{}}\par
\noindent Source : 7I:313. \nolinkurl{fixes/FIXLOG_7I_CURRENT.jsonl#741}; \nolinkurl{fixes/MISPRINTS_7I.json#313}\par

\par\medskip\noindent\hypertarget{app-F7I-VII-030.1}{}\textbf{F7I-VII-030.1}\quad p. 163; printed misprint.\par
\noindent Imprimé : {\ttfamily\footnotesize G\allowbreak{}}\par
\noindent Corrigé : {\ttfamily\footnotesize Q\allowbreak{}}\par
\noindent Source : 7I:314. \nolinkurl{fixes/MISPRINTS_7I.json#314}\par

\par\medskip\noindent\hypertarget{app-F7I-VII-031.1}{}\textbf{F7I-VII-031.1}\quad p. 164; printed misprint.\par
\noindent Imprimé : {\ttfamily\footnotesize ;\allowbreak{}}\par
\noindent Corrigé : {\ttfamily\footnotesize ,\allowbreak{}}\par
\noindent Source : 7I:315. \nolinkurl{fixes/MISPRINTS_7I.json#315}\par

\par\medskip\noindent\hypertarget{app-F7I-VII-032.1}{}\textbf{F7I-VII-032.1}\quad p. 164; printed misprint.\par
\noindent Imprimé : {\ttfamily\footnotesize G\allowbreak{}}\par
\noindent Corrigé : {\ttfamily\footnotesize Q\allowbreak{}}\par
\noindent Source : 7I:316. \nolinkurl{fixes/MISPRINTS_7I.json#316}\par

\par\medskip\noindent\hypertarget{app-F7I-VII-032.2}{}\textbf{F7I-VII-032.2}\quad p. 165; misprint\_pass.\par
\noindent Imprimé : {\ttfamily\footnotesize .\allowbreak{}}\par
\noindent Corrigé : \emph{(vide)}\par
\noindent Source : 7I:317. \nolinkurl{fixes/FIXLOG_7I_CURRENT.jsonl#742}; \nolinkurl{fixes/MISPRINTS_7I.json#317}\par

\par\medskip\noindent\hypertarget{app-F7I-VII-033.1}{}\textbf{F7I-VII-033.1}\quad p. 165; misprint\_pass.\par
\noindent Imprimé : {\ttfamily\footnotesize soient\allowbreak{}}\par
\noindent Corrigé : {\ttfamily\footnotesize soit\allowbreak{}}\par
\noindent Source : 7I:318. \nolinkurl{fixes/FIXLOG_7I_CURRENT.jsonl#743}; \nolinkurl{fixes/MISPRINTS_7I.json#318}\par

\par\medskip\noindent\hypertarget{app-F7I-VII-033.2}{}\textbf{F7I-VII-033.2}\quad p. 166; misprint\_pass.\par
\noindent Imprimé : {\ttfamily\footnotesize laloi\allowbreak{}}\par
\noindent Corrigé : {\ttfamily\footnotesize la\allowbreak{} loi\allowbreak{}}\par
\noindent Source : 7I:319. \nolinkurl{fixes/FIXLOG_7I_CURRENT.jsonl#744}; \nolinkurl{fixes/MISPRINTS_7I.json#319}\par

\par\medskip\noindent\hypertarget{app-F7I-VII-033.3}{}\textbf{F7I-VII-033.3}\quad p. 166; misprint\_pass.\par
\noindent Imprimé : {\ttfamily\footnotesize appelée\allowbreak{}}\par
\noindent Corrigé : {\ttfamily\footnotesize appelé\allowbreak{}}\par
\noindent Source : 7I:320. \nolinkurl{fixes/FIXLOG_7I_CURRENT.jsonl#745}; \nolinkurl{fixes/MISPRINTS_7I.json#320}\par

\par\medskip\noindent\hypertarget{app-F7I-VII-034.1}{}\textbf{F7I-VII-034.1}\quad p. 167; misprint\_pass.\par
\noindent Imprimé : {\ttfamily\footnotesize donnéepar\allowbreak{}}\par
\noindent Corrigé : {\ttfamily\footnotesize donnée\allowbreak{} par\allowbreak{}}\par
\noindent Source : 7I:322. \nolinkurl{fixes/FIXLOG_7I_CURRENT.jsonl#746}; \nolinkurl{fixes/MISPRINTS_7I.json#322}\par

\par\medskip\noindent\hypertarget{app-F7I-VII-034.2}{}\textbf{F7I-VII-034.2}\quad p. 167; misprint\_pass.\par
\noindent Imprimé : {\ttfamily\footnotesize isomorphismesde\allowbreak{}}\par
\noindent Corrigé : {\ttfamily\footnotesize isomorphismes\allowbreak{} de\allowbreak{}}\par
\noindent Source : 7I:323. \nolinkurl{fixes/FIXLOG_7I_CURRENT.jsonl#747}; \nolinkurl{fixes/MISPRINTS_7I.json#323}\par

\par\medskip\noindent\hypertarget{app-F7I-VII-035.1}{}\textbf{F7I-VII-035.1}\quad p. 168; misprint\_pass.\par
\noindent Imprimé : {\ttfamily\footnotesize est\allowbreak{}}\par
\noindent Corrigé : {\ttfamily\footnotesize sont\allowbreak{}}\par
\noindent Source : 7I:324. \nolinkurl{fixes/FIXLOG_7I_CURRENT.jsonl#748}; \nolinkurl{fixes/MISPRINTS_7I.json#324}\par

\par\medskip\noindent\hypertarget{app-F7I-VII-036.1}{}\textbf{F7I-VII-036.1}\quad p. 169; printed misprint.\par
\noindent Imprimé : {\ttfamily\footnotesize P'\allowbreak{}}\par
\noindent Corrigé : {\ttfamily\footnotesize Q\allowbreak{}}\par
\noindent Source : 7I:325. \nolinkurl{fixes/MISPRINTS_7I.json#325}\par

\par\medskip\noindent\hypertarget{app-F7I-VII-036.2}{}\textbf{F7I-VII-036.2}\quad p. 169; misprint\_pass.\par
\noindent Imprimé : {\ttfamily\footnotesize regardées\allowbreak{}}\par
\noindent Corrigé : {\ttfamily\footnotesize regardés\allowbreak{}}\par
\noindent Source : 7I:326. \nolinkurl{fixes/FIXLOG_7I_CURRENT.jsonl#749}\par

\par\medskip\noindent\hypertarget{app-F7I-VII-037.1}{}\textbf{F7I-VII-037.1}\quad p. 169; designated-witness attribution.\par
\noindent Imprimé : {\ttfamily\footnotesize P\allowbreak{}}\par
\noindent Corrigé : {\ttfamily\footnotesize Q\allowbreak{}}\par
\noindent Source : WITNESS-F7I-VII-037.1. \nolinkurl{fr/sga7i/diplomatic/expose_VII_body.tex}; \nolinkurl{fr/sga7i/current/expose_VII_body.tex}\par

\par\medskip\noindent\hypertarget{app-F7I-VII-037.2}{}\textbf{F7I-VII-037.2}\quad p. 169; inherited corrected-layer difference.\par
\noindent Imprimé : {\ttfamily\footnotesize Q\allowbreak{}}\par
\noindent Corrigé : {\ttfamily\footnotesize P\allowbreak{}}\par
\noindent Source : ADOPT-F7I-VII-037.2. \nolinkurl{fr/sga7i/diplomatic/expose_VII_body.tex}; \nolinkurl{fr/sga7i/current/expose_VII_body.tex}\par

\par\medskip\noindent\hypertarget{app-F7I-VII-038.1}{}\textbf{F7I-VII-038.1}\quad p. 171; misprint\_pass.\par
\noindent Imprimé : {\ttfamily\footnotesize appelera\allowbreak{}}\par
\noindent Corrigé : {\ttfamily\footnotesize appellera\allowbreak{}}\par
\noindent Source : 7I:340. \nolinkurl{fixes/FIXLOG_7I_CURRENT.jsonl#750}; \nolinkurl{fixes/MISPRINTS_7I.json#340}\par

\par\medskip\noindent\hypertarget{app-F7I-VII-039.1}{}\textbf{F7I-VII-039.1}\quad p. 173; misprint\_pass.\par
\noindent Imprimé : {\ttfamily\footnotesize appliqué\allowbreak{}}\par
\noindent Corrigé : {\ttfamily\footnotesize appliquée\allowbreak{}}\par
\noindent Source : 7I:341. \nolinkurl{fixes/FIXLOG_7I_CURRENT.jsonl#751}; \nolinkurl{fixes/MISPRINTS_7I.json#341}\par

\par\medskip\noindent\hypertarget{app-F7I-VII-039.2}{}\textbf{F7I-VII-039.2}\quad p. 173; printed misprint.\par
\noindent Imprimé : {\ttfamily\footnotesize P\allowbreak{}}\par
\noindent Corrigé : {\ttfamily\footnotesize Q\allowbreak{}}\par
\noindent Source : 7I:342. \nolinkurl{fixes/MISPRINTS_7I.json#342}\par

\par\medskip\noindent\hypertarget{app-F7I-VII-040.1}{}\textbf{F7I-VII-040.1}\quad p. 174; designated-witness attribution.\par
\noindent Imprimé : {\ttfamily\footnotesize conditions\allowbreak{} préliminaires\allowbreak{}}\par
\noindent Corrigé : {\ttfamily\footnotesize hypothèses\allowbreak{}}\par
\noindent Source : WITNESS-F7I-VII-040.1. \nolinkurl{fr/sga7i/diplomatic/expose_VII_body.tex}; \nolinkurl{fr/sga7i/current/expose_VII_body.tex}\par

\par\medskip\noindent\hypertarget{app-F7I-VII-040.2}{}\textbf{F7I-VII-040.2}\quad p. 174; designated-witness attribution.\par
\noindent Imprimé : \emph{(vide)}\par
\noindent Corrigé : {\ttfamily\footnotesize  la\allowbreak{} Proposition\allowbreak{}}\par
\noindent Source : WITNESS-F7I-VII-040.2. \nolinkurl{fr/sga7i/diplomatic/expose_VII_body.tex}; \nolinkurl{fr/sga7i/current/expose_VII_body.tex}\par

\par\medskip\noindent\hypertarget{app-F7I-VII-040.3}{}\textbf{F7I-VII-040.3}\quad p. 174; printed misprint.\par
\noindent Imprimé : {\ttfamily\footnotesize 2\allowbreak{}}\par
\noindent Corrigé : {\ttfamily\footnotesize 3\allowbreak{}}\par
\noindent Source : 7I:343. \nolinkurl{fixes/MISPRINTS_7I.json#343}\par

\par\medskip\noindent\hypertarget{app-F7I-VII-040.4}{}\textbf{F7I-VII-040.4}\quad p. 174; misprint\_in\_print.\par
\noindent Imprimé : {\ttfamily\footnotesize ,\allowbreak{}}\par
\noindent Corrigé : {\ttfamily\footnotesize .\allowbreak{}}\par
\noindent Source : FIX-7I-D-0236. \nolinkurl{fixes/FIXLOG_7I.jsonl#236}\par

\par\medskip\noindent\hypertarget{app-F7I-VII-041.1}{}\textbf{F7I-VII-041.1}\quad p. 175; printed misprint.\par
\noindent Imprimé : {\ttfamily\footnotesize 2\allowbreak{}}\par
\noindent Corrigé : {\ttfamily\footnotesize 3\allowbreak{}}\par
\noindent Source : 7I:345. \nolinkurl{fixes/MISPRINTS_7I.json#345}\par

\par\medskip\noindent\hypertarget{app-F7I-VII-041.2}{}\textbf{F7I-VII-041.2}\quad p. 175; printed misprint.\par
\noindent Imprimé : {\ttfamily\footnotesize 2\allowbreak{}}\par
\noindent Corrigé : {\ttfamily\footnotesize 3\allowbreak{}}\par
\noindent Source : 7I:346. \nolinkurl{fixes/MISPRINTS_7I.json#346}\par

\par\medskip\noindent\hypertarget{app-F7I-VII-041.3}{}\textbf{F7I-VII-041.3}\quad p. 175; inherited corrected-layer difference.\par
\noindent Imprimé : {\ttfamily\footnotesize mS\allowbreak{}}\par
\noindent Corrigé : {\ttfamily\footnotesize m,S\allowbreak{}}\par
\noindent Source : ADOPT-F7I-VII-041.3. \nolinkurl{fr/sga7i/diplomatic/expose_VII_body.tex}; \nolinkurl{fr/sga7i/current/expose_VII_body.tex}\par

\par\medskip\noindent\hypertarget{app-F7I-VII-042.1}{}\textbf{F7I-VII-042.1}\quad p. 176; inherited corrected-layer difference.\par
\noindent Imprimé : \emph{(vide)}\par
\noindent Corrigé : {\ttfamily\footnotesize \}\textasciicircum{}\{1\allowbreak{}}\par
\noindent Source : ADOPT-F7I-VII-042.1. \nolinkurl{fr/sga7i/diplomatic/expose_VII_body.tex}; \nolinkurl{fr/sga7i/current/expose_VII_body.tex}\par

\par\medskip\noindent\hypertarget{app-F7I-VII-043.1}{}\textbf{F7I-VII-043.1}\quad p. 176; misprint\_pass.\par
\noindent Imprimé : {\ttfamily\footnotesize autons\allowbreak{}}\par
\noindent Corrigé : {\ttfamily\footnotesize aurons\allowbreak{}}\par
\noindent Source : 7I:348. \nolinkurl{fixes/FIXLOG_7I_CURRENT.jsonl#752}; \nolinkurl{fixes/MISPRINTS_7I.json#348}\par

\par\medskip\noindent\hypertarget{app-F7I-VII-043.2}{}\textbf{F7I-VII-043.2}\quad p. 177; misprint\_pass.\par
\noindent Imprimé : {\ttfamily\footnotesize entension\allowbreak{}}\par
\noindent Corrigé : {\ttfamily\footnotesize extension\allowbreak{}}\par
\noindent Source : 7I:349. \nolinkurl{fixes/FIXLOG_7I_CURRENT.jsonl#753}; \nolinkurl{fixes/MISPRINTS_7I.json#349}\par

\par\medskip\noindent\hypertarget{app-F7I-VII-043.3}{}\textbf{F7I-VII-043.3}\quad p. 177; misprint\_pass.\par
\noindent Imprimé : {\ttfamily\footnotesize commuatifs\allowbreak{}}\par
\noindent Corrigé : {\ttfamily\footnotesize commutatifs\allowbreak{}}\par
\noindent Source : 7I:350. \nolinkurl{fixes/FIXLOG_7I_CURRENT.jsonl#754}; \nolinkurl{fixes/MISPRINTS_7I.json#350}\par

\par\medskip\noindent\hypertarget{app-F7I-VII-043.4}{}\textbf{F7I-VII-043.4}\quad p. 177; misprint\_pass.\par
\noindent Imprimé : {\ttfamily\footnotesize général\allowbreak{}}\par
\noindent Corrigé : {\ttfamily\footnotesize générale\allowbreak{}}\par
\noindent Source : 7I:351. \nolinkurl{fixes/FIXLOG_7I_CURRENT.jsonl#755}; \nolinkurl{fixes/MISPRINTS_7I.json#351}\par

\par\medskip\noindent\hypertarget{app-F7I-VII-044.1}{}\textbf{F7I-VII-044.1}\quad p. 177; designated-witness attribution.\par
\noindent Imprimé : \emph{(vide)}\par
\noindent Corrigé : {\ttfamily\footnotesize =\textbackslash{}underline\{G\}\_\{m,S\}\allowbreak{}}\par
\noindent Source : WITNESS-F7I-VII-044.1. \nolinkurl{fr/sga7i/diplomatic/expose_VII_body.tex}; \nolinkurl{fr/sga7i/current/expose_VII_body.tex}\par

\par\medskip\noindent\hypertarget{app-F7I-VII-044.2}{}\textbf{F7I-VII-044.2}\quad p. 177; designated-witness attribution.\par
\noindent Imprimé : {\ttfamily\footnotesize \textbackslash{}underline\{G\}\_m\allowbreak{}}\par
\noindent Corrigé : {\ttfamily\footnotesize S\allowbreak{}}\par
\noindent Source : WITNESS-F7I-VII-044.2. \nolinkurl{fr/sga7i/diplomatic/expose_VII_body.tex}; \nolinkurl{fr/sga7i/current/expose_VII_body.tex}\par

\par\medskip\noindent\hypertarget{app-F7I-VII-044.3}{}\textbf{F7I-VII-044.3}\quad p. 177; designated-witness attribution.\par
\noindent Imprimé : {\ttfamily\footnotesize 2\allowbreak{}}\par
\noindent Corrigé : {\ttfamily\footnotesize 3\allowbreak{}}\par
\noindent Source : WITNESS-F7I-VII-044.3. \nolinkurl{fr/sga7i/diplomatic/expose_VII_body.tex}; \nolinkurl{fr/sga7i/current/expose_VII_body.tex}\par

\par\medskip\noindent\hypertarget{app-F7I-VII-044.4}{}\textbf{F7I-VII-044.4}\quad p. 177; misprint\_pass.\par
\noindent Imprimé : {\ttfamily\footnotesize -\allowbreak{}}\par
\noindent Corrigé : \emph{(1 espace(s), 0 retour(s) de ligne dans le témoin TeX)}\par
\noindent Source : 7I:353. \nolinkurl{fixes/FIXLOG_7I_CURRENT.jsonl#756}; \nolinkurl{fixes/MISPRINTS_7I.json#353}\par

\par\medskip\noindent\hypertarget{app-F7I-VII-045.1}{}\textbf{F7I-VII-045.1}\quad p. 178; printed misprint.\par
\noindent Imprimé : {\ttfamily\footnotesize P\allowbreak{}}\par
\noindent Corrigé : {\ttfamily\footnotesize E\allowbreak{}}\par
\noindent Source : 7I:354. \nolinkurl{fixes/MISPRINTS_7I.json#354}\par

\par\medskip\noindent\hypertarget{app-F7I-VII-045.2}{}\textbf{F7I-VII-045.2}\quad p. 178; printed misprint.\par
\noindent Imprimé : {\ttfamily\footnotesize P\allowbreak{}}\par
\noindent Corrigé : {\ttfamily\footnotesize E\allowbreak{}}\par
\noindent Source : 7I:354. \nolinkurl{fixes/MISPRINTS_7I.json#354}\par

\par\medskip\noindent\hypertarget{app-F7I-VII-045.3}{}\textbf{F7I-VII-045.3}\quad p. 178; printed misprint.\par
\noindent Imprimé : {\ttfamily\footnotesize P\allowbreak{}}\par
\noindent Corrigé : {\ttfamily\footnotesize E\allowbreak{}}\par
\noindent Source : 7I:355. \nolinkurl{fixes/MISPRINTS_7I.json#355}\par

\par\medskip\noindent\hypertarget{app-F7I-VII-046.1}{}\textbf{F7I-VII-046.1}\quad p. 179; editorial decision.\par
\noindent Imprimé : {\ttfamily\footnotesize cofibré\allowbreak{}}\par
\noindent Corrigé : {\ttfamily\footnotesize cofibrée\allowbreak{}}\par
\noindent Source : ED-7I-356. \nolinkurl{OPEN_ITEMS_DECIDED.tsv}\par

\par\medskip\noindent\hypertarget{app-F7I-VII-047.1}{}\textbf{F7I-VII-047.1}\quad p. 181; misprint\_pass.\par
\noindent Imprimé : \emph{(vide)}\par
\noindent Corrigé : {\ttfamily\footnotesize .\allowbreak{}}\par
\noindent Source : 7I:359. \nolinkurl{fixes/FIXLOG_7I_CURRENT.jsonl#757}; \nolinkurl{fixes/MISPRINTS_7I.json#359}\par

\par\medskip\noindent\hypertarget{app-F7I-VII-048.1}{}\textbf{F7I-VII-048.1}\quad p. 182; printed misprint.\par
\noindent Imprimé : {\ttfamily\footnotesize 7.1\allowbreak{}}\par
\noindent Corrigé : {\ttfamily\footnotesize 8\allowbreak{}}\par
\noindent Source : 7I:360. \nolinkurl{fixes/MISPRINTS_7I.json#360}\par

\par\medskip\noindent\hypertarget{app-F7I-VII-048.2}{}\textbf{F7I-VII-048.2}\quad p. 182; misprint\_pass.\par
\noindent Imprimé : {\ttfamily\footnotesize munie\allowbreak{}}\par
\noindent Corrigé : {\ttfamily\footnotesize muni\allowbreak{}}\par
\noindent Source : 7I:361. \nolinkurl{fixes/FIXLOG_7I_CURRENT.jsonl#758}; \nolinkurl{fixes/MISPRINTS_7I.json#361}\par

\par\medskip\noindent\hypertarget{app-F7I-VII-049.1}{}\textbf{F7I-VII-049.1}\quad p. 183; misprint\_pass.\par
\noindent Imprimé : {\ttfamily\footnotesize -\allowbreak{}}\par
\noindent Corrigé : \emph{(1 espace(s), 0 retour(s) de ligne dans le témoin TeX)}\par
\noindent Source : 7I:362. \nolinkurl{fixes/FIXLOG_7I_CURRENT.jsonl#759}\par

\par\medskip\noindent\hypertarget{app-F7I-VII-050.1}{}\textbf{F7I-VII-050.1}\quad p. 184; misprint\_pass.\par
\noindent Imprimé : \emph{(vide)}\par
\noindent Corrigé : {\ttfamily\footnotesize  nous\allowbreak{}}\par
\noindent Source : 7I:364. \nolinkurl{fixes/FIXLOG_7I_CURRENT.jsonl#760}; \nolinkurl{fixes/MISPRINTS_7I.json#364}\par

\par\medskip\noindent\hypertarget{app-F7I-VII-051.1}{}\textbf{F7I-VII-051.1}\quad p. 185; printed misprint.\par
\noindent Imprimé : {\ttfamily\footnotesize j\allowbreak{}}\par
\noindent Corrigé : {\ttfamily\footnotesize s\allowbreak{}}\par
\noindent Source : 7I:365. \nolinkurl{fixes/MISPRINTS_7I.json#365}\par

\par\medskip\noindent\hypertarget{app-F7I-VII-052.1}{}\textbf{F7I-VII-052.1}\quad p. 186; misprint\_pass.\par
\noindent Imprimé : \emph{(vide)}\par
\noindent Corrigé : {\ttfamily\footnotesize ,\allowbreak{}}\par
\noindent Source : 7I:366. \nolinkurl{fixes/FIXLOG_7I_CURRENT.jsonl#761}; \nolinkurl{fixes/MISPRINTS_7I.json#366}\par

\par\medskip\noindent\hypertarget{app-F7I-VII-053.1}{}\textbf{F7I-VII-053.1}\quad p. 186; misprint\_pass.\par
\noindent Imprimé : \emph{(vide)}\par
\noindent Corrigé : {\ttfamily\footnotesize  que\allowbreak{}}\par
\noindent Source : 7I:367. \nolinkurl{fixes/FIXLOG_7I_CURRENT.jsonl#762}\par

\par\medskip\noindent\hypertarget{app-F7I-VII-054.1}{}\textbf{F7I-VII-054.1}\quad p. 187; misprint\_pass.\par
\noindent Imprimé : \emph{(vide)}\par
\noindent Corrigé : {\ttfamily\footnotesize '\allowbreak{}}\par
\noindent Source : 7I:368. \nolinkurl{fixes/FIXLOG_7I_CURRENT.jsonl#763}; \nolinkurl{fixes/MISPRINTS_7I.json#368}\par

\par\medskip\noindent\hypertarget{app-F7I-VII-055.1}{}\textbf{F7I-VII-055.1}\quad p. 187; printed misprint.\par
\noindent Imprimé : \emph{(vide)}\par
\noindent Corrigé : {\ttfamily\footnotesize 3.\allowbreak{}}\par
\noindent Source : 7I:369. \nolinkurl{fixes/MISPRINTS_7I.json#369}\par

\par\medskip\noindent\hypertarget{app-F7I-VII-055.2}{}\textbf{F7I-VII-055.2}\quad p. 187; designated-witness attribution.\par
\noindent Imprimé : {\ttfamily\footnotesize 3\allowbreak{}}\par
\noindent Corrigé : {\ttfamily\footnotesize 1\allowbreak{}}\par
\noindent Source : WITNESS-F7I-VII-055.2. \nolinkurl{fr/sga7i/diplomatic/expose_VII_body.tex}; \nolinkurl{fr/sga7i/current/expose_VII_body.tex}\par

\par\medskip\noindent\hypertarget{app-F7I-VII-055.3}{}\textbf{F7I-VII-055.3}\quad p. 187; designated-witness attribution.\par
\noindent Imprimé : \emph{(vide)}\par
\noindent Corrigé : {\ttfamily\footnotesize 3.\allowbreak{}}\par
\noindent Source : WITNESS-F7I-VII-055.3. \nolinkurl{fr/sga7i/diplomatic/expose_VII_body.tex}; \nolinkurl{fr/sga7i/current/expose_VII_body.tex}\par

\par\medskip\noindent\hypertarget{app-F7I-VII-055.4}{}\textbf{F7I-VII-055.4}\quad p. 187; designated-witness attribution.\par
\noindent Imprimé : {\ttfamily\footnotesize 3\allowbreak{}}\par
\noindent Corrigé : {\ttfamily\footnotesize 1\allowbreak{}}\par
\noindent Source : WITNESS-F7I-VII-055.4. \nolinkurl{fr/sga7i/diplomatic/expose_VII_body.tex}; \nolinkurl{fr/sga7i/current/expose_VII_body.tex}\par

\par\medskip\noindent\hypertarget{app-F7I-VII-056.1}{}\textbf{F7I-VII-056.1}\quad p. 188; misprint\_pass.\par
\noindent Imprimé : \emph{(vide)}\par
\noindent Corrigé : {\ttfamily\footnotesize '\allowbreak{}}\par
\noindent Source : 7I:371. \nolinkurl{fixes/FIXLOG_7I_CURRENT.jsonl#764}; \nolinkurl{fixes/MISPRINTS_7I.json#371}\par

\par\medskip\noindent\hypertarget{app-F7I-VII-056.2}{}\textbf{F7I-VII-056.2}\quad p. 188; inherited corrected-layer difference.\par
\noindent Imprimé : \emph{(vide)}\par
\noindent Corrigé : {\ttfamily\footnotesize \textbackslash{}bar\allowbreak{} d\_o:L'\_1\textbackslash{}longrightar\allowbreak{}row\allowbreak{} L\_o\$\allowbreak{} le\allowbreak{} morphisme\allowbreak{} induit\allowbreak{} par\allowbreak{} \$d\_o:L\_1\textbackslash{}longrightar\allowbreak{}row\allowbreak{} L\_o\$.\allowbreak{} On\allowbreak{} désigne\allowbreak{} par\allowbreak{} \$\allowbreak{}}\par
\noindent Source : ADOPT-F7I-VII-056.2. \nolinkurl{fr/sga7i/diplomatic/expose_VII_body.tex}; \nolinkurl{fr/sga7i/current/expose_VII_body.tex}\par

\par\medskip\noindent\hypertarget{app-F7I-VII-056.3}{}\textbf{F7I-VII-056.3}\quad p. 188; designated-witness attribution.\par
\noindent Imprimé : {\ttfamily\footnotesize a\allowbreak{}}\par
\noindent Corrigé : {\ttfamily\footnotesize est\allowbreak{}}\par
\noindent Source : WITNESS-F7I-VII-056.3. \nolinkurl{fr/sga7i/diplomatic/expose_VII_body.tex}; \nolinkurl{fr/sga7i/current/expose_VII_body.tex}\par

\par\medskip\noindent\hypertarget{app-F7I-VII-056.4}{}\textbf{F7I-VII-056.4}\quad p. 188; designated-witness attribution.\par
\noindent Imprimé : {\ttfamily\footnotesize comme\allowbreak{}}\par
\noindent Corrigé : {\ttfamily\footnotesize le\allowbreak{}}\par
\noindent Source : WITNESS-F7I-VII-056.4. \nolinkurl{fr/sga7i/diplomatic/expose_VII_body.tex}; \nolinkurl{fr/sga7i/current/expose_VII_body.tex}\par

\par\medskip\noindent\hypertarget{app-F7I-VII-056.5}{}\textbf{F7I-VII-056.5}\quad p. 188; designated-witness attribution.\par
\noindent Imprimé : {\ttfamily\footnotesize composantes\allowbreak{}}\par
\noindent Corrigé : {\ttfamily\footnotesize composé\allowbreak{} de\allowbreak{} \$L\_1\textbackslash{}twoheadrightarr\allowbreak{}ow\allowbreak{} L'\_1\$\allowbreak{} avec\allowbreak{}}\par
\noindent Source : WITNESS-F7I-VII-056.5. \nolinkurl{fr/sga7i/diplomatic/expose_VII_body.tex}; \nolinkurl{fr/sga7i/current/expose_VII_body.tex}\par

\par\medskip\noindent\hypertarget{app-F7I-VII-056.6}{}\textbf{F7I-VII-056.6}\quad p. 188; designated-witness attribution.\par
\noindent Imprimé : \emph{(vide)}\par
\noindent Corrigé : {\ttfamily\footnotesize \textbackslash{}bar\allowbreak{} }\par
\noindent Source : WITNESS-F7I-VII-056.6. \nolinkurl{fr/sga7i/diplomatic/expose_VII_body.tex}; \nolinkurl{fr/sga7i/current/expose_VII_body.tex}\par

\par\medskip\noindent\hypertarget{app-F7I-VII-056.7}{}\textbf{F7I-VII-056.7}\quad p. 188; designated-witness attribution.\par
\noindent Imprimé : {\ttfamily\footnotesize 1'\allowbreak{}}\par
\noindent Corrigé : {\ttfamily\footnotesize o\allowbreak{}}\par
\noindent Source : WITNESS-F7I-VII-056.7. \nolinkurl{fr/sga7i/diplomatic/expose_VII_body.tex}; \nolinkurl{fr/sga7i/current/expose_VII_body.tex}\par

\par\medskip\noindent\hypertarget{app-F7I-VII-056.8}{}\textbf{F7I-VII-056.8}\quad p. 188; designated-witness attribution.\par
\noindent Imprimé : {\ttfamily\footnotesize \$\allowbreak{} ,\allowbreak{} \$d\allowbreak{}}\par
\noindent Corrigé : {\ttfamily\footnotesize :L'\allowbreak{}}\par
\noindent Source : WITNESS-F7I-VII-056.8. \nolinkurl{fr/sga7i/diplomatic/expose_VII_body.tex}; \nolinkurl{fr/sga7i/current/expose_VII_body.tex}\par

\par\medskip\noindent\hypertarget{app-F7I-VII-056.9}{}\textbf{F7I-VII-056.9}\quad p. 188; designated-witness attribution.\par
\noindent Imprimé : {\ttfamily\footnotesize '\allowbreak{} :\allowbreak{} L\_1\allowbreak{} }\par
\noindent Corrigé : \emph{(vide)}\par
\noindent Source : WITNESS-F7I-VII-056.9. \nolinkurl{fr/sga7i/diplomatic/expose_VII_body.tex}; \nolinkurl{fr/sga7i/current/expose_VII_body.tex}\par

\par\medskip\noindent\hypertarget{app-F7I-VII-056.10}{}\textbf{F7I-VII-056.10}\quad p. 188; designated-witness attribution.\par
\noindent Imprimé : {\ttfamily\footnotesize \$\allowbreak{} étant\allowbreak{} déduit\allowbreak{} de\allowbreak{} \$d\allowbreak{}}\par
\noindent Corrigé : {\ttfamily\footnotesize +L'\allowbreak{}}\par
\noindent Source : WITNESS-F7I-VII-056.10. \nolinkurl{fr/sga7i/diplomatic/expose_VII_body.tex}; \nolinkurl{fr/sga7i/current/expose_VII_body.tex}\par

\par\medskip\noindent\hypertarget{app-F7I-VII-056.11}{}\textbf{F7I-VII-056.11}\quad p. 188; designated-witness attribution.\par
\noindent Imprimé : {\ttfamily\footnotesize  :\allowbreak{} L\_1\allowbreak{} \textbackslash{}longrightarrow\allowbreak{} L\_o\allowbreak{}}\par
\noindent Corrigé : \emph{(vide)}\par
\noindent Source : WITNESS-F7I-VII-056.11. \nolinkurl{fr/sga7i/diplomatic/expose_VII_body.tex}; \nolinkurl{fr/sga7i/current/expose_VII_body.tex}\par

\par\medskip\noindent\hypertarget{app-F7I-VII-056.12}{}\textbf{F7I-VII-056.12}\quad p. 188; designated-witness attribution.\par
\noindent Imprimé : {\ttfamily\footnotesize  par\allowbreak{} passage\allowbreak{} au\allowbreak{} quotient\allowbreak{}}\par
\noindent Corrigé : \emph{(vide)}\par
\noindent Source : WITNESS-F7I-VII-056.12. \nolinkurl{fr/sga7i/diplomatic/expose_VII_body.tex}; \nolinkurl{fr/sga7i/current/expose_VII_body.tex}\par

\par\medskip\noindent\hypertarget{app-F7I-VII-056.13}{}\textbf{F7I-VII-056.13}\quad p. 188; misprint\_pass.\par
\noindent Imprimé : {\ttfamily\footnotesize majorée\allowbreak{}}\par
\noindent Corrigé : {\ttfamily\footnotesize majoré\allowbreak{}}\par
\noindent Source : 7I:375. \nolinkurl{fixes/FIXLOG_7I_CURRENT.jsonl#765}; \nolinkurl{fixes/MISPRINTS_7I.json#375}\par

\par\medskip\noindent\hypertarget{app-F7I-VII-057.1}{}\textbf{F7I-VII-057.1}\quad p. 189; printed misprint.\par
\noindent Imprimé : {\ttfamily\footnotesize j\allowbreak{}}\par
\noindent Corrigé : {\ttfamily\footnotesize i\allowbreak{}}\par
\noindent Source : 7I:377. \nolinkurl{fixes/MISPRINTS_7I.json#377}\par

\par\medskip\noindent\hypertarget{app-F7I-VII-057.2}{}\textbf{F7I-VII-057.2}\quad p. 189; misprint\_pass.\par
\noindent Imprimé : {\ttfamily\footnotesize réconstitue\allowbreak{}}\par
\noindent Corrigé : {\ttfamily\footnotesize reconstitue\allowbreak{}}\par
\noindent Source : 7I:378. \nolinkurl{fixes/FIXLOG_7I_CURRENT.jsonl#766}; \nolinkurl{fixes/MISPRINTS_7I.json#378}\par

\par\medskip\noindent\hypertarget{app-F7I-VII-058.1}{}\textbf{F7I-VII-058.1}\quad p. 190; misprint\_pass.\par
\noindent Imprimé : {\ttfamily\footnotesize )\allowbreak{}}\par
\noindent Corrigé : \emph{(vide)}\par
\noindent Source : 7I:379. \nolinkurl{fixes/FIXLOG_7I_CURRENT.jsonl#767}\par

\par\medskip\noindent\hypertarget{app-F7I-VII-059.1}{}\textbf{F7I-VII-059.1}\quad p. 190; printed misprint.\par
\noindent Imprimé : {\ttfamily\footnotesize o\allowbreak{}}\par
\noindent Corrigé : {\ttfamily\footnotesize \textbackslash{}partial\allowbreak{}}\par
\noindent Source : 7I:380. \nolinkurl{fixes/MISPRINTS_7I.json#380}\par

\par\medskip\noindent\hypertarget{app-F7I-VII-060.1}{}\textbf{F7I-VII-060.1}\quad p. 191; printed misprint.\par
\noindent Imprimé : {\ttfamily\footnotesize auticommutatif\allowbreak{}}\par
\noindent Corrigé : {\ttfamily\footnotesize anticommutatif\allowbreak{}}\par
\noindent Source : 7I:381. \nolinkurl{fixes/MISPRINTS_7I.json#381}\par

\par\medskip\noindent\hypertarget{app-F7I-VII-060.2}{}\textbf{F7I-VII-060.2}\quad p. 192; misprint\_pass.\par
\noindent Imprimé : {\ttfamily\footnotesize différent\allowbreak{}}\par
\noindent Corrigé : {\ttfamily\footnotesize diffèrent\allowbreak{}}\par
\noindent Source : 7I:382. \nolinkurl{fixes/FIXLOG_7I_CURRENT.jsonl#768}\par

\par\medskip\noindent\hypertarget{app-F7I-VII-060.3}{}\textbf{F7I-VII-060.3}\quad p. 192; misprint\_pass.\par
\noindent Imprimé : {\ttfamily\footnotesize réconstituant\allowbreak{}}\par
\noindent Corrigé : {\ttfamily\footnotesize reconstituant\allowbreak{}}\par
\noindent Source : 7I:383. \nolinkurl{fixes/FIXLOG_7I_CURRENT.jsonl#769}\par

\par\medskip\noindent\hypertarget{app-F7I-VII-060.4}{}\textbf{F7I-VII-060.4}\quad p. 192; misprint\_pass.\par
\noindent Imprimé : {\ttfamily\footnotesize isomorphismes\allowbreak{}}\par
\noindent Corrigé : {\ttfamily\footnotesize homomorphismes\allowbreak{}}\par
\noindent Source : 7I:384. \nolinkurl{fixes/FIXLOG_7I_CURRENT.jsonl#770}\par

\par\medskip\noindent\hypertarget{app-F7I-VII-061.1}{}\textbf{F7I-VII-061.1}\quad p. 192; misprint\_pass.\par
\noindent Imprimé : \emph{(vide)}\par
\noindent Corrigé : {\ttfamily\footnotesize \{\allowbreak{}}\par
\noindent Source : 7I:385. \nolinkurl{fixes/FIXLOG_7I_CURRENT.jsonl#771}; \nolinkurl{fixes/MISPRINTS_7I.json#385}\par

\par\medskip\noindent\hypertarget{app-F7I-VII-061.2}{}\textbf{F7I-VII-061.2}\quad p. 192; misprint\_pass.\par
\noindent Imprimé : \emph{(vide)}\par
\noindent Corrigé : {\ttfamily\footnotesize .\}\allowbreak{}}\par
\noindent Source : 7I:385. \nolinkurl{fixes/FIXLOG_7I_CURRENT.jsonl#771}; \nolinkurl{fixes/MISPRINTS_7I.json#385}\par

\par\medskip\noindent\hypertarget{app-F7I-VII-061.3}{}\textbf{F7I-VII-061.3}\quad p. 192; misprint\_pass.\par
\noindent Imprimé : {\ttfamily\footnotesize  .\allowbreak{}}\par
\noindent Corrigé : {\ttfamily\footnotesize ,\allowbreak{}}\par
\noindent Source : 7I:385. \nolinkurl{fixes/FIXLOG_7I_CURRENT.jsonl#771}; \nolinkurl{fixes/MISPRINTS_7I.json#385}\par

\par\medskip\noindent\hypertarget{app-F7I-VII-062.1}{}\textbf{F7I-VII-062.1}\quad p. 193; misprint\_in\_print.\par
\noindent Imprimé : {\ttfamily\footnotesize \textbackslash{}qquad\allowbreak{}}\par
\noindent Corrigé : {\ttfamily\footnotesize \textbackslash{}longrightarrow\allowbreak{}}\par
\noindent Source : FIX-7I-D-0244. \nolinkurl{fixes/FIXLOG_7I.jsonl#244}\par

\par\medskip\noindent\hypertarget{app-F7I-VII-063.1}{}\textbf{F7I-VII-063.1}\quad p. 194; misprint\_pass.\par
\noindent Imprimé : {\ttfamily\footnotesize laissé\allowbreak{}}\par
\noindent Corrigé : {\ttfamily\footnotesize laissée\allowbreak{}}\par
\noindent Source : 7I:387. \nolinkurl{fixes/FIXLOG_7I_CURRENT.jsonl#772}\par

\par\medskip\noindent\hypertarget{app-F7I-VII-064.1}{}\textbf{F7I-VII-064.1}\quad p. 194; printed misprint.\par
\noindent Imprimé : \emph{(vide)}\par
\noindent Corrigé : {\ttfamily\footnotesize \textasciicircum{}1\allowbreak{}}\par
\noindent Source : 7I:388. \nolinkurl{fixes/MISPRINTS_7I.json#388}\par

\par\medskip\noindent\hypertarget{app-F7I-VII-065.1}{}\textbf{F7I-VII-065.1}\quad p. 195; printed misprint.\par
\noindent Imprimé : {\ttfamily\footnotesize '\allowbreak{}}\par
\noindent Corrigé : \emph{(vide)}\par
\noindent Source : 7I:390. \nolinkurl{fixes/MISPRINTS_7I.json#390}\par

\par\medskip\noindent\hypertarget{app-F7I-VII-065.2}{}\textbf{F7I-VII-065.2}\quad p. 195; printed misprint.\par
\noindent Imprimé : \emph{(vide)}\par
\noindent Corrigé : {\ttfamily\footnotesize '\allowbreak{}}\par
\noindent Source : 7I:391. \nolinkurl{fixes/MISPRINTS_7I.json#391}\par

\par\medskip\noindent\hypertarget{app-F7I-VII-065.3}{}\textbf{F7I-VII-065.3}\quad p. 195; printed misprint.\par
\noindent Imprimé : {\ttfamily\footnotesize jt\allowbreak{}}\par
\noindent Corrigé : {\ttfamily\footnotesize j't\allowbreak{}}\par
\noindent Source : 7I:392. \nolinkurl{fixes/MISPRINTS_7I.json#392}\par

\par\medskip\noindent\hypertarget{app-F7I-VII-065.4}{}\textbf{F7I-VII-065.4}\quad p. 195; misprint\_in\_print.\par
\noindent Imprimé : {\ttfamily\footnotesize ,\allowbreak{} leftrightarrow\allowbreak{}}\par
\noindent Corrigé : \emph{(vide)}\par
\noindent Source : FIX-7I-D-0247. \nolinkurl{fixes/FIXLOG_7I.jsonl#247}\par

\par\medskip\noindent\hypertarget{app-F7I-VII-066.1}{}\textbf{F7I-VII-066.1}\quad p. 196; inherited corrected-layer difference.\par
\noindent Imprimé : {\ttfamily\footnotesize jt\allowbreak{}}\par
\noindent Corrigé : {\ttfamily\footnotesize j't\allowbreak{}}\par
\noindent Source : ADOPT-F7I-VII-066.1. \nolinkurl{fr/sga7i/diplomatic/expose_VII_body.tex}; \nolinkurl{fr/sga7i/current/expose_VII_body.tex}\par

\par\medskip\noindent\hypertarget{app-F7I-VII-066.2}{}\textbf{F7I-VII-066.2}\quad p. 196; inherited corrected-layer difference.\par
\noindent Imprimé : \emph{(vide)}\par
\noindent Corrigé : {\ttfamily\footnotesize '\allowbreak{}}\par
\noindent Source : ADOPT-F7I-VII-066.2. \nolinkurl{fr/sga7i/diplomatic/expose_VII_body.tex}; \nolinkurl{fr/sga7i/current/expose_VII_body.tex}\par

\par\medskip\noindent\hypertarget{app-F7I-VII-066.3}{}\textbf{F7I-VII-066.3}\quad p. 196; printed misprint.\par
\noindent Imprimé : \emph{(vide)}\par
\noindent Corrigé : {\ttfamily\footnotesize ''\allowbreak{}}\par
\noindent Source : 7I:398. \nolinkurl{fixes/MISPRINTS_7I.json#398}\par

\par\medskip\noindent\hypertarget{app-F7I-VII-067.1}{}\textbf{F7I-VII-067.1}\quad p. 196; printed misprint.\par
\noindent Imprimé : \emph{(vide)}\par
\noindent Corrigé : {\ttfamily\footnotesize ,\allowbreak{}}\par
\noindent Source : 7I:397. \nolinkurl{fixes/MISPRINTS_7I.json#397}\par

\par\medskip\noindent\hypertarget{app-F7I-VII-067.2}{}\textbf{F7I-VII-067.2}\quad p. 197; misprint\_pass.\par
\noindent Imprimé : \emph{(vide)}\par
\noindent Corrigé : {\ttfamily\footnotesize .\allowbreak{}}\par
\noindent Source : 7I:399. \nolinkurl{fixes/FIXLOG_7I_CURRENT.jsonl#773}; \nolinkurl{fixes/MISPRINTS_7I.json#399}\par

\par\medskip\noindent\hypertarget{app-F7I-VII-067.3}{}\textbf{F7I-VII-067.3}\quad p. 197; misprint\_pass.\par
\noindent Imprimé : {\ttfamily\footnotesize -\allowbreak{}}\par
\noindent Corrigé : \emph{(1 espace(s), 0 retour(s) de ligne dans le témoin TeX)}\par
\noindent Source : 7I:400. \nolinkurl{fixes/FIXLOG_7I_CURRENT.jsonl#774}; \nolinkurl{fixes/MISPRINTS_7I.json#400}\par

\par\medskip\noindent\hypertarget{app-F7I-VII-068.1}{}\textbf{F7I-VII-068.1}\quad p. 199; designated-witness attribution.\par
\noindent Imprimé : {\ttfamily\footnotesize 1\allowbreak{}}\par
\noindent Corrigé : {\ttfamily\footnotesize o\allowbreak{}}\par
\noindent Source : WITNESS-F7I-VII-068.1. \nolinkurl{fr/sga7i/diplomatic/expose_VII_body.tex}; \nolinkurl{fr/sga7i/current/expose_VII_body.tex}\par

\par\medskip\noindent\hypertarget{app-F7I-VII-068.2}{}\textbf{F7I-VII-068.2}\quad p. 199; designated-witness attribution.\par
\noindent Imprimé : {\ttfamily\footnotesize o\allowbreak{}}\par
\noindent Corrigé : {\ttfamily\footnotesize 1\allowbreak{}}\par
\noindent Source : WITNESS-F7I-VII-068.2. \nolinkurl{fr/sga7i/diplomatic/expose_VII_body.tex}; \nolinkurl{fr/sga7i/current/expose_VII_body.tex}\par

\par\medskip\noindent\hypertarget{app-F7I-VII-068.3}{}\textbf{F7I-VII-068.3}\quad p. 199; misprint\_pass.\par
\noindent Imprimé : {\ttfamily\footnotesize questions\allowbreak{}}\par
\noindent Corrigé : {\ttfamily\footnotesize question\allowbreak{}}\par
\noindent Source : 7I:402. \nolinkurl{fixes/FIXLOG_7I_CURRENT.jsonl#775}\par

\par\medskip\noindent\hypertarget{app-F7I-VII-069.1}{}\textbf{F7I-VII-069.1}\quad p. 201; inherited corrected-layer difference.\par
\noindent Imprimé : {\ttfamily\footnotesize 1\allowbreak{}}\par
\noindent Corrigé : {\ttfamily\footnotesize o\allowbreak{}}\par
\noindent Source : ADOPT-F7I-VII-069.1. \nolinkurl{fr/sga7i/diplomatic/expose_VII_body.tex}; \nolinkurl{fr/sga7i/current/expose_VII_body.tex}\par

\par\medskip\noindent\hypertarget{app-F7I-VII-069.2}{}\textbf{F7I-VII-069.2}\quad p. 201; inherited corrected-layer difference.\par
\noindent Imprimé : {\ttfamily\footnotesize o\allowbreak{}}\par
\noindent Corrigé : {\ttfamily\footnotesize 1\allowbreak{}}\par
\noindent Source : ADOPT-F7I-VII-069.2. \nolinkurl{fr/sga7i/diplomatic/expose_VII_body.tex}; \nolinkurl{fr/sga7i/current/expose_VII_body.tex}\par

\par\medskip\noindent\hypertarget{app-F7I-VII-070.1}{}\textbf{F7I-VII-070.1}\quad p. 204; misprint\_pass.\par
\noindent Imprimé : {\ttfamily\footnotesize coïncident\allowbreak{}}\par
\noindent Corrigé : {\ttfamily\footnotesize coïncidant\allowbreak{}}\par
\noindent Source : 7I:404. \nolinkurl{fixes/FIXLOG_7I_CURRENT.jsonl#776}\par

\par\medskip\noindent\hypertarget{app-F7I-VII-071.1}{}\textbf{F7I-VII-071.1}\quad p. 204; printed misprint.\par
\noindent Imprimé : \emph{(vide)}\par
\noindent Corrigé : {\ttfamily\footnotesize '\allowbreak{}}\par
\noindent Source : 7I:403. \nolinkurl{fixes/MISPRINTS_7I.json#403}\par

\par\medskip\noindent\hypertarget{app-F7I-VII-071.2}{}\textbf{F7I-VII-071.2}\quad p. 204; printed misprint.\par
\noindent Imprimé : {\ttfamily\footnotesize '\allowbreak{}}\par
\noindent Corrigé : \emph{(vide)}\par
\noindent Source : 7I:403. \nolinkurl{fixes/MISPRINTS_7I.json#403}\par

\par\medskip\noindent\hypertarget{app-F7I-VII-071.3}{}\textbf{F7I-VII-071.3}\quad p. 204; printed misprint.\par
\noindent Imprimé : {\ttfamily\footnotesize '\allowbreak{}}\par
\noindent Corrigé : \emph{(vide)}\par
\noindent Source : 7I:403. \nolinkurl{fixes/MISPRINTS_7I.json#403}\par

\par\medskip\noindent\hypertarget{app-F7I-VII-072.1}{}\textbf{F7I-VII-072.1}\quad p. 205; printed misprint.\par
\noindent Imprimé : {\ttfamily\footnotesize \textbackslash{}times\allowbreak{} }\par
\noindent Corrigé : {\ttfamily\footnotesize ,\allowbreak{}}\par
\noindent Source : 7I:405. \nolinkurl{fixes/MISPRINTS_7I.json#405}\par

\par\medskip\noindent\hypertarget{app-F7I-VII-072.2}{}\textbf{F7I-VII-072.2}\quad p. 206; misprint\_pass.\par
\noindent Imprimé : {\ttfamily\footnotesize desisomorphismes\allowbreak{}}\par
\noindent Corrigé : {\ttfamily\footnotesize des\allowbreak{} isomorphismes\allowbreak{}}\par
\noindent Source : 7I:406. \nolinkurl{fixes/FIXLOG_7I_CURRENT.jsonl#777}\par

\par\medskip\noindent\hypertarget{app-F7I-VII-072.3}{}\textbf{F7I-VII-072.3}\quad p. 206; printed misprint.\par
\noindent Imprimé : {\ttfamily\footnotesize 1q\allowbreak{}}\par
\noindent Corrigé : {\ttfamily\footnotesize 1,q\allowbreak{}}\par
\noindent Source : 7I:407. \nolinkurl{fixes/MISPRINTS_7I.json#407}\par

\par\medskip\noindent\hypertarget{app-F7I-VII-072.4}{}\textbf{F7I-VII-072.4}\quad p. 206; printed misprint.\par
\noindent Imprimé : \emph{(vide)}\par
\noindent Corrigé : {\ttfamily\footnotesize  P\textbackslash{}times\allowbreak{}}\par
\noindent Source : 7I:408. \nolinkurl{fixes/MISPRINTS_7I.json#408}\par

\par\medskip\noindent\hypertarget{app-F7I-VII-072.5}{}\textbf{F7I-VII-072.5}\quad p. 206; printed misprint.\par
\noindent Imprimé : {\ttfamily\footnotesize P\allowbreak{}}\par
\noindent Corrigé : {\ttfamily\footnotesize Q\allowbreak{}}\par
\noindent Source : 7I:409. \nolinkurl{fixes/MISPRINTS_7I.json#409}\par

\par\medskip\noindent\hypertarget{app-F7I-VII-073.1}{}\textbf{F7I-VII-073.1}\quad p. 208; misprint\_pass.\par
\noindent Imprimé : {\ttfamily\footnotesize catégorie\allowbreak{}}\par
\noindent Corrigé : {\ttfamily\footnotesize catégories\allowbreak{}}\par
\noindent Source : 7I:410. \nolinkurl{fixes/FIXLOG_7I_CURRENT.jsonl#778}\par

\par\medskip\noindent\hypertarget{app-F7I-VII-074.1}{}\textbf{F7I-VII-074.1}\quad p. 211; printed misprint.\par
\noindent Imprimé : {\ttfamily\footnotesize gauches\allowbreak{}}\par
\noindent Corrigé : {\ttfamily\footnotesize gauche\allowbreak{}}\par
\noindent Source : 7I:411. \nolinkurl{fixes/MISPRINTS_7I.json#411}\par

\par\medskip\noindent\hypertarget{app-F7I-VII-075.1}{}\textbf{F7I-VII-075.1}\quad p. 211; misprint\_pass.\par
\noindent Imprimé : {\ttfamily\footnotesize complexe\allowbreak{}}\par
\noindent Corrigé : {\ttfamily\footnotesize complexes\allowbreak{}}\par
\noindent Source : 7I:412. \nolinkurl{fixes/FIXLOG_7I_CURRENT.jsonl#779}; \nolinkurl{fixes/MISPRINTS_7I.json#412}\par

\par\medskip\noindent\hypertarget{app-F7I-VII-076.1}{}\textbf{F7I-VII-076.1}\quad p. 212; designated-witness attribution.\par
\noindent Imprimé : {\ttfamily\footnotesize ]\allowbreak{} \textbackslash{}arrow[dr,\allowbreak{} dash,\allowbreak{} "s"\allowbreak{}}\par
\noindent Corrigé : \emph{(vide)}\par
\noindent Source : WITNESS-F7I-VII-076.1. \nolinkurl{fr/sga7i/diplomatic/expose_VII_body.tex}; \nolinkurl{fr/sga7i/current/expose_VII_body.tex}\par

\par\medskip\noindent\hypertarget{app-F7I-VII-076.2}{}\textbf{F7I-VII-076.2}\quad p. 212; designated-witness attribution.\par
\noindent Imprimé : \emph{(vide)}\par
\noindent Corrigé : {\ttfamily\footnotesize  \textbackslash{}arrow[ul,\allowbreak{} "s"]\allowbreak{}}\par
\noindent Source : WITNESS-F7I-VII-076.2. \nolinkurl{fr/sga7i/diplomatic/expose_VII_body.tex}; \nolinkurl{fr/sga7i/current/expose_VII_body.tex}\par

\par\medskip\noindent\hypertarget{app-F7I-VII-077.1}{}\textbf{F7I-VII-077.1}\quad p. 213; printed misprint.\par
\noindent Imprimé : {\ttfamily\footnotesize formée\allowbreak{}}\par
\noindent Corrigé : {\ttfamily\footnotesize formé\allowbreak{}}\par
\noindent Source : 7I:417. \nolinkurl{fixes/MISPRINTS_7I.json#417}\par

\par\medskip\noindent\hypertarget{app-F7I-VII-078.1}{}\textbf{F7I-VII-078.1}\quad p. 217; misprint\_pass.\par
\noindent Imprimé : {\ttfamily\footnotesize 1956\allowbreak{}}\par
\noindent Corrigé : {\ttfamily\footnotesize 1957\allowbreak{}}\par
\noindent Source : 7I:418. \nolinkurl{fixes/FIXLOG_7I_CURRENT.jsonl#780}; \nolinkurl{fixes/MISPRINTS_7I.json#418}\par

\par\medskip\noindent\hypertarget{app-F7I-VII-078.2}{}\textbf{F7I-VII-078.2}\quad p. 217; printed misprint.\par
\noindent Imprimé : {\ttfamily\footnotesize Lectures\allowbreak{}}\par
\noindent Corrigé : {\ttfamily\footnotesize Lecture\allowbreak{}}\par
\noindent Source : 7I:419. \nolinkurl{fixes/MISPRINTS_7I.json#419}\par

\par\medskip\noindent\hypertarget{app-F7I-VII-078.3}{}\textbf{F7I-VII-078.3}\quad p. 217; designated-witness attribution.\par
\noindent Imprimé : {\ttfamily\footnotesize Math.\allowbreak{}}\par
\noindent Corrigé : {\ttfamily\footnotesize Mathematics\allowbreak{}}\par
\noindent Source : WITNESS-F7I-VII-078.3. \nolinkurl{fr/sga7i/diplomatic/expose_VII_body.tex}; \nolinkurl{fr/sga7i/current/expose_VII_body.tex}\par

\normalsize\fi
\clearpage\SGAtoc{Exposé VIII — COMPLEMENTS SUR LES BIEXTENSIONS. PROPRIETES GENERALES DES BIEXTENSIONS DES SCHEMAS EN GROUPES}
\selectlanguage{french}

\SGApage{I}{218}


S G A 7


\underline{E X P O S E  VIII}


\underline{\underline{COMPLEMENTS SUR LES BIEXTENSIONS. PROPRIETES GENERALES}}


\underline{\underline{DES BIEXTENSIONS DES SCHEMAS EN GROUPES}}

\underline{par A. Grothendieck}


\underline{S O M M A I R E}

0. Introduction

1. Cas particuliers divers sur un topos quelconque

2. Accouplements définis par une biextension

3. Biextensions de schémas en groupes $(P,Q)$ par $\underline{G}_m$ : généralités

4. Extensions et biextensions des schémas en groupes lisses et connexes sur un corps

5. Extensions et biextensions par des schémas en groupes constants tronqués sans torsion

6. Extensions et biextensions par le modèle de Néron de $\underline{G}_m$

7. Prolongements canoniques d'extensions et de biextensions par $\underline{G}_m$


\SGAsectiontoc{0. Introduction}{0}
0. \underline{Introduction}

0.1.  Nous continuons ici l'étude générale des biextensions, commencée dans l'exposé précédent. Tout d'abord nous donnons quelques compléments dans le cas d'un topos de base général, en développant notamment, avec les détails qu'il mérite, le formalisme des accouplements associés aux biextensions (n$^{\circ}$2), dont l'importance dans la théorie des schémas abéliens est bien connue. D'autre part, nous abordons l'étude des biextensions de schémas en groupes généraux, le cas le plus intéressant pour nous étant celui des biextensions par le groupe multiplicatif $\underline{G}_m$ . Nous serons cependant obligés, à titre d'intermédiaire technique, d'étudier aussi les biextensions par d'autres types de schémas en groupes (n$^{\circ}$s 5 et 6).


\SGApage{I}{219}

Le résultat le plus intéressant pour la suite est donné dans 7.1, qui donne des conditions moyennant lesquelles, pour $P$ et $Q$ lisses sur un trait de point générique $\eta$, une biextension de $(P_\eta,Q_\eta)$ par $\underline{G}_{m_\eta}$ peut se prolonger en une biextension de $(P,Q)$ par $\underline{G}_m$ , déterminée à isomorphisme unique près. Ce résultat sera fort commode dans l'exposé suivant, consacré à la théorie du modèle de Néron des variétés abéliennes. Il est probable qu'on y pourrait s'en tirer à meilleur compte, sans l'imposant poids-papier des exposés VII et VIII (qui ne serviront plus dans la suite du Séminaire) ; il nous semble cependant probable que les développements donnés dans ces exposés sont appelés à être utiles encore ailleurs.

0.2.  Lorsque nous travaillons avec des schémas en groupes sur un schéma $S$, il sera sous-entendu que nous travaillons avec le topos fppf de $S$, qui est le topos associé au $\underline{U}$-site des schémas localement de présentation finie sur $S$ (et $\in \underline{U}$), muni de la topologie fppf (SGA 3 IV 6.3). Le lecteur vérifiera d'ailleurs sans peine que la plupart des énoncés que nous donnerons seraient également valables (et essentiellement équivalents aux énoncés fppf) pour d'autres topologies sur $S$ (ayant comme catégorie sous-jacente au site une catégorie contenant au moins les schémas localement de présentation finie sur $S$), telles les topologies de Zariski, la topologie étale, la topologie fpqc. Une des raisons en est que, grâce à la théorie de la descente, la catégorie des extensions d'un schéma en groupes $P$ localement de présentation finie sur $S$ par le schéma en groupes $\underline{G}_m$ ne dépend pas, à équivalence près, de la topologie choisie, et le même énoncé vaut pour la notion de biextension. La raison technique pour laquelle il y a lieu de préférer la topologie fppf (ou une topologie plus fine, comme fpqc) aux topologies moins fines comme celle de Zariski et la topologie étale, c'est que pour tout entier $n > 0$, l'homomorphisme $n.\mathrm{id}$ de $\underline{G}_m$ dans lui-même est un \underline{épimorphisme} fppf (car il est fidèlement plat de présentation finie), mais en général pas un épimorphisme (Zar) ou (ét). (NB c'est un épimorphisme (ét) si et seulement si $n$ est premier aux caractéristiques résiduelles de $S$ et un épimorphisme (Zar) si et seulement si $S$ est vide.)


\SGApage{I}{220}

0.3.  Signalons enfin que la plupart des énoncés concernant les extensions ou les biextensions par $\underline{G}_m$ resteraient valables en remplaçant ce groupe par un \underline{tore} $G$. Comme un tel $G$ est localement isomorphe (pour fppf, et même pour (ét)) à un groupe $\underline{G}_m^r$ , il s'agit là d'une généralisation essentiellement triviale.


\SGAsectiontoc{1. Cas particuliers divers sur un topos quelconque}{1}
1. \underline{Cas particuliers divers sur un topos quelconque}

1.1.  Dans le présent numéro et le suivant, comme dans VII, $T$ désigne un topos, $P,Q$ et $G$ dénotent trois Groupes abéliens de $T$, de sorte qu'on a un isomorphisme canonique (VII 3.6.5) :

\begin{equation}\tag{1.1.1}
\mathrm{Biext}^1(P,Q;G) \;\simeq\; \mathrm{Ext}^1(P\overset{\mathbb{L}}{\otimes}Q\ ;G) \qquad\qquad .
\end{equation}
Rappelons qu'on a des isomorphismes canoniques (dits ``de Cartan'')

\begin{equation}\tag{1.1.2}
\mathrm{Ext}^1(P\overset{\mathbb{L}}{\otimes}Q,\ G) \;\simeq\; \mathrm{Ext}^1(P,R\underline{\mathrm{Hom}}(Q,G)) \;\simeq\; \mathrm{Ext}^1(Q,R\underline{\mathrm{Hom}}(P,G)) \qquad\qquad .
\end{equation}
Tous ces isomorphismes sont fonctoriels en les trois arguments $P,Q,G$.

\SGAvar{On déduit de (1.1.1) la suite exacte infinie suivante, compte tenu que les objets d'homologie de $P\overset{\mathbb{L}}{\otimes}Q$ sont les $\underline{\mathrm{Tor}}_i(P,Q)$, qui sont nuls pour $i \geqslant 1$ :

}{On déduit de (1.1.1), par la troncation standard de $P\overset{\mathbb{L}}{\otimes}Q$, dont les objets d'homologie sont les $\underline{\mathrm{Tor}}_i(P,Q)$, la suite exacte infinie suivante :

}{F7I-VIII-001}\begin{equation}\tag{1.1.3}
\begin{aligned}
&0 \longrightarrow \mathrm{Ext}^1(P\otimes Q,G) \longrightarrow \mathrm{Biext}^1(P,Q;G) \longrightarrow \mathrm{Hom}(\underline{\mathrm{Tor}}_1(P,Q),G) \longrightarrow \\
&\qquad\qquad \mathrm{Ext}^2(P\otimes Q,G) \longrightarrow \mathrm{Ext}^2(P\overset{\mathbb{L}}{\otimes} Q,G) \longrightarrow \quad \ldots. \qquad\qquad\qquad .
\end{aligned}
\end{equation}
D'autre part, les objets de cohomologie de $R\underline{\mathrm{Hom}}(Q,G)$ étant les $\underline{\mathrm{Ext}}^i(Q,G)$, la première égalité (1.1.2) donne lieu à la suite exacte à 5 termes,

\SGAvar{\begin{equation}\tag{1.1.4}
\begin{aligned}
&0 \longrightarrow \mathrm{Ext}^1(P,\underline{\mathrm{Hom}}(Q,G)) \longrightarrow \mathrm{Biext}^1(P,Q;G) \longrightarrow \mathrm{Hom}(P,\underline{\mathrm{Ext}}^1(Q,G)) \\
&\qquad\qquad \mathrm{Ext}^2(P,\underline{\mathrm{Hom}}(Q,P)) \longrightarrow \mathrm{Ext}^2(P,R\underline{\mathrm{Hom}}(Q,G)) \qquad\qquad ,
\end{aligned}
\end{equation}
\SGApage{I}{221}
associée à la suite spectrale habituelle

}{\begin{equation}\tag{1.1.4}
\begin{aligned}
&0 \longrightarrow \mathrm{Ext}^1(P,\underline{\mathrm{Hom}}(Q,G)) \longrightarrow \mathrm{Biext}^1(P,Q;G) \longrightarrow \mathrm{Hom}(P,\underline{\mathrm{Ext}}^1(Q,G)) \longrightarrow \\
&\qquad\qquad \mathrm{Ext}^2(P,\underline{\mathrm{Hom}}(Q,G)) \longrightarrow \mathrm{Ext}^2(P,R\underline{\mathrm{Hom}}(Q,G)) \qquad\qquad ,
\end{aligned}
\end{equation}
\SGApage{I}{221}
associée à la suite spectrale habituelle

}{F7I-VIII-002}\begin{equation*}
\mathrm{Ext}^*(P,R\underline{\mathrm{Hom}}(Q,G)) \;\Longleftarrow\; \mathrm{Ext}^p(P,\underline{\mathrm{Ext}}^q(Q,G)) \qquad\qquad .
\end{equation*}
Nous laissons au lecteur le soin de préciser, s'il en ressent le besoin, la première flèche de (1.1.3) (resp. de (1.1.4)) à l'aide d'un \underline{foncteur naturel pleinement fidèle}

\begin{equation}\tag{1.1.5}
\mathrm{EXT}(P\otimes Q,G) \;\longrightarrow\; \mathrm{BIEXT}(P,Q;G)
\end{equation}
(resp.

\begin{equation}\tag{1.1.6}
\mathrm{EXT}(P,\underline{\mathrm{Hom}}(Q,G)) \;\longrightarrow\; \mathrm{BIEXT}(P,Q;G) \qquad\qquad\qquad ),
\end{equation}
où EXT et BIEXT désignent les catégories d'extensions et de biextensions qu'on devine.

Bien entendu, l'isomorphisme $\mathrm{Biext}^1(P,Q;G) \simeq \mathrm{Ext}^1(Q,R\underline{\mathrm{Hom}}(P,G))$ donne naissance à une suite exacte analogue à (1.1.3), avec les rôles de $P$ et $Q$ intervertis, dont la première flèche peut se préciser à l'aide d'un foncteur pleinement fidèle

\begin{equation}\tag{1.1.7}
\mathrm{EXT}(Q,\underline{\mathrm{Hom}}(P,G)) \;\longrightarrow\; \mathrm{BIEXT}(P,Q;G)
\end{equation}
analogue à (1.1.6).

1.2.  Supposons qu'on ait

\begin{equation}\tag{1.2.1}
\underline{\mathrm{Tor}}_1(P,Q) \;=\; 0 \qquad\qquad ,
\end{equation}
alors (1.1.3) fournit un isomorphisme canonique

\begin{equation}\tag{1.2.2}
\mathrm{Biext}^1(P,Q;G) \;\simeq\; \mathrm{Ext}^1(P\otimes Q,G) \qquad\qquad .
\end{equation}
Il en est ainsi par exemple si $P$ ou $Q$ est un $\mathbb{Z}$-Module \underline{plat}. On peut évidemment préciser cet énoncé en disant que sous la condition (1.2.1), le foncteur
\SGApage{I}{222}
(1.1.5) est une équivalence de catégories.

1.3.  Supposons qu'on ait

\begin{equation}\tag{1.3.1}
P \otimes Q \;=\; 0 \qquad ,
\end{equation}
alors (1.1.3) fournit un isomorphisme canonique

\begin{equation}\tag{1.3.2}
\mathrm{Biext}^1(P,Q;G) \;\simeq\; \mathrm{Hom}(\underline{\mathrm{Tor}}_1(P,Q),G) \qquad .
\end{equation}
Le cas le plus intéressant pour nous où (1.3.1) est vérifié est celui où on dispose d'un entier $n > 0$ tel que l'on ait

\begin{equation}\tag{1.3.3}
nP=P \qquad , \quad Q \;=\; \varinjlim_i {}_{n^i}Q \qquad ,
\end{equation}
i.e. $P$ est divisible par $n$ et $Q$ est limite de sous-Groupes annulés par des puissances de $n$. Posant, suivant l'usage

\begin{equation*}
T_n(P) \;=\; ({}_{n^i}P)_{i \geq 0} \quad (\text{système projectif}), \qquad {}_{n^\infty}P \;=\; \varinjlim_i {}_{n^i}P \qquad ,
\end{equation*}
on conclut de (1.3.3) des formules plus précises que (1.3.1) :

\begin{equation}\tag{1.3.4}
P \otimes Q \;=\; 0 \quad , \quad \underline{\mathrm{Tor}}_1(P,Q) \;\simeq\; \varinjlim_i {}_{n^i}P \otimes {}_{n^i}Q \;\simeq\; T_n(P) \otimes Q \;\simeq\; {}_{n^\infty}P \otimes T_n(Q) \quad ,
\end{equation}
de sorte que l'isomorphisme (1.3.2) prend ici la forme plus concrète :

\SGAvar{\begin{equation}\tag{1.3.5}
\mathrm{Biext}^1(P,Q;G) \;\simeq\; \mathrm{Hom}(T_n(P),D(Q)) \;\simeq\; \mathrm{Hom}(Q,D(T_n(P)) \;\simeq\; \mathrm{Hom}(T_n(P) \otimes T_n(Q),T_n(G)).
\end{equation}
Précisons la signification de ces formules :

}{\begin{equation}\tag{1.3.5}
\mathrm{Biext}^1(P,Q;G) \;\simeq\; \mathrm{Hom}(T_n(P),D(Q)) \;\simeq\; \mathrm{Hom}(Q,D(T_n(P))) \;\simeq\; \mathrm{Hom}(T_n(P) \otimes T_n(Q),T_n(G)).
\end{equation}
Précisons la signification de ces formules :

}{F7I-VIII-003}a)  Dans le système inductif qui intervient dans (1.3.4), l'homomorphisme de transition

\begin{equation*}
{}_{n^i}P \otimes {}_{n^i}Q \;\longrightarrow\; {}_{n^j}P \otimes {}_{n^j}Q \qquad (i \leq j) \qquad 
\end{equation*}
\SGApage{I}{223}
 est décrit en identifiant le premier membre à ${}_{n^j}P \otimes {}_{n^i}Q$ grâce à $n^{j-i} \otimes \mathrm{id}$, et en utilisant l'homomorphisme

\begin{equation*}
\mathrm{id} \otimes \mathrm{inj} \;:\; {}_{n^j}P \otimes {}_{n^i}Q \;\longrightarrow\; {}_{n^j}P \otimes {}_{n^j}Q \qquad ,
\end{equation*}
où $\mathrm{inj} : {}_{n^i}Q \longrightarrow {}_{n^j}Q$ désigne l'inclusion.

b)  On pose, par définition

\begin{equation*}
T_n(P) \otimes Q \;\overset{\mathrm{dfn}}{=}\; \varinjlim T_n(P) \otimes {}_{n^i}Q \qquad ,
\end{equation*}
où

\SGAvar{\begin{equation*}
T_n(P) \otimes M \;\overset{\mathrm{dfn}}{=}\; {}_{n^i}P \otimes M \qquad \text{si} \qquad n^i M = 0 \qquad .
\end{equation*}
On précise de même le sens du troisième terme écrit dans (1.3.4), lorsqu'on suppose que $nQ = Q$ (de sorte que $T_n(Q)$ est, comme $T_n(P)$, un système projectif strict). Le dernier terme de (1.3.4) est un $\mathrm{Hom}$ de systèmes projectifs de groupes.

}{\begin{equation*}
T_n(P) \otimes M \;\overset{\mathrm{dfn}}{=}\; {}_{n^i}P \otimes M \qquad \text{si} \qquad n^i M = 0 \qquad .
\end{equation*}
On précise de même le sens du troisième terme écrit dans (1.3.4), lorsqu'on suppose que $nQ = Q$ (de sorte que $T_n(Q)$ est, comme $T_n(P)$, un système projectif strict). Le dernier terme de (1.3.5) est un $\mathrm{Hom}$ de systèmes projectifs de groupes.

}{F7I-VIII-004}c)  La lettre $D$ désigne $\underline{\mathrm{Hom}}(-,G)$. L'objet $D(Q)$ est défini comme le système projectif de Groupes

\SGAvar{\begin{equation*}
D(Q) \;=\; (D({}_{n^i}Q))_{i \geq 0} \;=\; (\mathrm{Hom}({}_{n^i}Q,G))_{i \geq 0} \qquad ,
\end{equation*}
le deuxième membre de (1.3.5) étant le groupe d'homomorphismes de système projectifs de groupes. Le dernier membre s'interprète plus simplement en posant

\begin{equation*}
D(T_n(P)) \;=\; \varinjlim_i D({}_{n^i}P) \qquad ,
\end{equation*}
les morphismes de transition de ce sytème inductif étant évidemment ceux qui se déduisent par transposition des morphismes de transition
\SGApage{I}{224}
$n^{j-1}$ : ${}_{n^j}P \longrightarrow {}_{n^i}P$ de $T_n(P)$.

}{\begin{equation*}
D(Q) \;=\; (D({}_{n^i}Q))_{i \geq 0} \;=\; (\mathrm{Hom}({}_{n^i}Q,G))_{i \geq 0} \qquad ,
\end{equation*}
le deuxième membre de (1.3.5) étant le groupe d'homomorphismes de systèmes projectifs de groupes. Le membre suivant s'interprète plus simplement en posant

\begin{equation*}
D(T_n(P)) \;=\; \varinjlim_i D({}_{n^i}P) \qquad ,
\end{equation*}
les morphismes de transition de ce système inductif étant évidemment ceux qui se déduisent par transposition des morphismes de transition
\SGApage{I}{224}
$n^{j-i}$ : ${}_{n^j}P \longrightarrow {}_{n^i}P$ de $T_n(P)$.

}{F7I-VIII-005}Nous laissons au lecteur le soin d'expliciter le calcul de $\underline{\mathrm{Tor}}_1(P,Q)$ indiqué dans (1.3.4), et la conséquence (1.3.5) de ce calcul. Notons seulement qu'on utilise les relations déduites de la deuxième hypothèse (1.3.3) :

\begin{equation*}
\underline{\mathrm{Tor}}_p(P,Q) \;\simeq\; \varinjlim_{i} \underline{\mathrm{Tor}}_p(P, {}_{n^i}Q) \qquad ,
\end{equation*}
qui pour $p=0$ donne immédiatement $P \otimes Q = 0$, et qui ramène le calcul du $\underline{\mathrm{Tor}}_1(P,Q)$ à celui des $\underline{\mathrm{Tor}}_1(P, {}_{n^i}Q)$ et de leurs morphismes de transition. Il reste donc, lorsque $Q$ est annulé par un entier $m = n^i$ tel que $mP = P$, à expliciter un isomorphisme

\SGAvar{\begin{equation}\tag{1.3.6}
\underline{\mathrm{Tor}}_1(P,Q) \;\stackrel{\sim}{\longleftarrow}\; {}_mP \otimes {}_mQ \qquad (mP=P \;\; , \;\; mQ = 0) \qquad .
\end{equation}
Il y a diverses façons de la faire, l'une d'elles consistant à utiliser la suite exacte des $\underline{\mathrm{Tor}}$ associée à la suite exacte

}{\begin{equation}\tag{1.3.6}
\underline{\mathrm{Tor}}_1(P,Q) \;\stackrel{\sim}{\longleftarrow}\; {}_mP \otimes {}_mQ \qquad (mP=P \;\; , \;\; mQ = 0) \qquad .
\end{equation}
Il y a diverses façons de le faire, l'une d'elles consistant à utiliser la suite exacte des $\underline{\mathrm{Tor}}$ associée à la suite exacte

}{F7I-VIII-006}\begin{equation*}
0 \longrightarrow {}_mP \longrightarrow P \;\stackrel{m}{\longrightarrow}\; P \longrightarrow 0 \qquad .
\end{equation*}
On peut aussi utiliser l'homomorphisme (2.1.1) ci-dessous ; le fait que les morphismes de transition sont ceux précisés dans a) n'est alors autre que 2.1.11. Signalons que de toutes façons, le choix de (1.3.6) revient au choix d'un signe, cf. 2.1.10.

1.4.  Supposons que l'on ait

\begin{equation}\tag{1.4.1}
\underline{\mathrm{Hom}}(Q,G) = 0,
\end{equation}
alors (1.1.4) fournit un isomorphisme canonique

\SGAvar{\begin{equation}\tag{1.4.2}
\mathrm{Biext}^1(P,Q;G) \;\stackrel{\sim}{\longrightarrow}\; \mathrm{Hom}(P, \underline{\mathrm{Ext}}^1(Q,G)) \qquad .
\end{equation}
\SGApage{I}{225}
Ce cas se présente en particulier si $Q$ est un schéma abélien et si $G = \underline{G}_m$ , cf. VII1.3.8.Donc le foncteur en $P$, $P \longmapsto \mathrm{Biext}^1(P,Q;G)$ est \underline{représentable} par $Q^* = \underline{\mathrm{Ext}}^1(Q,G)$, et on trouve en particulier une biextension canonique $E^Q$ de $(Q^*,Q)$ par $G$, défini à isomorphisme unique près (puisque le groupe des automorphismes de $E^Q$ est $\mathrm{Hom}(P,\underline{\mathrm{Hom}}(Q,G)) = 0$). On pourra appeler cette biextension, ou la biextension symétrique ${}^sE^Q$ de $(Q,Q^*)$ par $G$, la \underline{biextension de Weil} de $(Q^*,Q)$ ou $(Q,Q^*)$ par $G$.

}{\begin{equation}\tag{1.4.2}
\mathrm{Biext}^1(P,Q;G) \;\stackrel{\sim}{\longrightarrow}\; \mathrm{Hom}(P, \underline{\mathrm{Ext}}^1(Q,G)) \qquad .
\end{equation}
\SGApage{I}{225}
Ce cas se présente en particulier si $Q$ est un schéma abélien et si $G = \underline{G}_m$ , cf. VII 1.3.8. Donc le foncteur en $P$, $P \longmapsto \mathrm{Biext}^1(P,Q;G)$ est \underline{représentable} par $Q^* = \underline{\mathrm{Ext}}^1(Q,G)$, et on trouve en particulier une biextension canonique $E^Q$ de $(Q^*,Q)$ par $G$, définie à isomorphisme unique près (puisque le groupe des automorphismes de $E^Q$ est $\mathrm{Hom}(Q^*,\underline{\mathrm{Hom}}(Q,G)) = 0$). On pourra appeler cette biextension, ou la biextension symétrique ${}^sE^Q$ de $(Q,Q^*)$ par $G$, la \underline{biextension de Weil} de $(Q^*,Q)$ ou $(Q,Q^*)$ par $G$.

}{F7I-VIII-007}1.5.  Supposons que l'on ait

\begin{equation}\tag{1.5.1}
\mathrm{Hom}(P, \underline{\mathrm{Ext}}^1(Q,G)) = 0 \qquad ,
\end{equation}
ce qui est le cas en particulier si $\underline{\mathrm{Ext}}^1(Q,G) = 0$ (par exemple si $Q$ est un schéma en groupes fini localement libre ou un groupe de type multiplicatif (SGA 3 IX), et $G = \underline{G}_m$ , cf. 3.4 ci-dessous). Alors (1.1.4) fournit un isomorphisme canonique

\begin{equation}\tag{1.5.2}
\mathrm{Biext}^1(P,Q; G) \;\stackrel{\sim}{\longleftarrow}\; \mathrm{Ext}^1(P,\underline{\mathrm{Hom}}(Q,G)) \qquad ,
\end{equation}
où $\underline{\mathrm{Hom}}(Q,G)$ joue le rôle d'un ``groupe dual'' de $Q$ (relativement à $G$). Plus précisément, on déduit de 1.5.2 que le foncteur pleinement fidèle (1.1.6) est une \underline{équivalence de catégories}.

\SGAvar{1.6.  Suppossons que l'on ait simultanément les relations

\begin{equation}\tag{1.6.1}
\underline{\mathrm{Hom}}(P,G) = 0, \;\; \mathrm{Hom}(P,\underline{\mathrm{Ext}}^1(Q,G)) = 0 \qquad .
\end{equation}
En vertu de 1.4, ces deux relations impliquent respectivement qu'on a des isomorphismes

}{1.6.  Supposons que l'on ait simultanément les relations

\begin{equation}\tag{1.6.1}
\underline{\mathrm{Hom}}(P,G) = 0, \;\; \mathrm{Hom}(P,\underline{\mathrm{Ext}}^1(Q,G)) = 0 \qquad .
\end{equation}
En vertu de 1.4 et 1.5, ces deux relations impliquent respectivement qu'on a des isomorphismes

}{F7I-VIII-008}\begin{equation*}
\mathrm{Biext}^1(P,Q;G) \;\stackrel{\sim}{\longrightarrow}\; \mathrm{Hom}(Q,\underline{\mathrm{Ext}}^1(P,G)) \qquad ,
\end{equation*}
\begin{equation*}
\mathrm{Biext}^1(P,Q;G) \;\stackrel{\sim}{\longleftarrow}\; \mathrm{Ext}^1(P,\underline{\mathrm{Hom}}(Q,G)) \qquad , \qquad 
\end{equation*}
\SGApage{I}{226}
 d'où en composant un isomorphisme canonique

\SGAvar{\begin{equation}\tag{1.6.2}
\mathrm{Ext}^1(P,\underline{\mathrm{Hom}}(Q,G)) \xrightarrow{\ \sim\ } \mathrm{Hom}(Q,\underline{\mathrm{Ext}}^1(P,G)) \ ( \simeq \mathrm{Biext}^1(P,G;G)) \ ,
\end{equation}
jouant le rôle d'une formule de dualité. Notons d'ailleurs que la première relation (1.6.1) implique évidemment que $\mathrm{Hom}(P \otimes Q,G) = 0$, donc que la catégorie $\mathrm{EXT}(P,\underline{\mathrm{Hom}}(Q,G))$ des extensions (commutatives) de $P$ par $\underline{\mathrm{Hom}}(Q,G)$ est rigide (tout automorphisme d'un objet est l'identité), et qu'il en est de même de $\mathrm{BIEXT}(P,Q;G)$. La structure de ces catégories est donc connue quand on connaît le groupe des classes d'isomorphie d'objets, qui en vertu de (1.6.2) est canoniquement isomorphe à $\mathrm{Hom}(Q,\underline{\mathrm{Ext}}^1(P,Q))$.

}{\begin{equation}\tag{1.6.2}
\mathrm{Ext}^1(P,\underline{\mathrm{Hom}}(Q,G)) \xrightarrow{\ \sim\ } \mathrm{Hom}(Q,\underline{\mathrm{Ext}}^1(P,G)) \ ( \simeq \mathrm{Biext}^1(P,Q;G)) \ ,
\end{equation}
jouant le rôle d'une formule de dualité. Notons d'ailleurs que la première relation (1.6.1) implique évidemment que $\mathrm{Hom}(P \otimes Q,G) = 0$, donc que la catégorie $\mathrm{EXT}(P,\underline{\mathrm{Hom}}(Q,G))$ des extensions (commutatives) de $P$ par $\underline{\mathrm{Hom}}(Q,G)$ est rigide (tout automorphisme d'un objet est l'identité), et qu'il en est de même de $\mathrm{BIEXT}(P,Q;G)$. La structure de ces catégories est donc connue quand on connaît le groupe des classes d'isomorphie d'objets, qui en vertu de (1.6.2) est canoniquement isomorphe à $\mathrm{Hom}(Q,\underline{\mathrm{Ext}}^1(P,G))$.

}{F7I-VIII-009}Notons que l'application (1.6.2) est définie en tous cas (sans supposer (1.6.1)), et qu'elle peut se décrire aussi directement, sans passer par $\mathrm{Biext}^1(P,Q;G)$, en termes de l'accouplement naturel

\begin{equation*}
Q \times Q' \longrightarrow G
\end{equation*}
(où $Q' = \underline{\mathrm{Hom}}(Q,G)$), en l'interprétant comme un homomorphisme $Q \to \underline{\mathrm{Hom}}(Q',G)$, d'où un composé

\begin{equation}\tag{1.6.3}
\mathrm{Ext}^1(P,Q') \times Q \to \mathrm{Ext}^1(P,Q') \times \mathrm{Hom}(Q',G) \longrightarrow \mathrm{Ext}^1(P,G),
\end{equation}
qui redonne (1.6.2).

Pour un cas utile où les considérations de 1.6 s'appliquent, voir 3.7 ci-dessous.


\SGAsectiontoc{2. Accouplements définis par une biextension}{2}
2. \underline{Accouplements définis par une biextension}

2.1.  Soient $P$ et $Q$ deux Groupes abéliens sur $T$, et $n > 0$ un entier. Rappelons qu'on désigne par ${}_nP$ le noyau de la multiplication par $n \ \mathrm{id}_P$. Nous allons définir un homomorphisme canonique 


\SGApage{I}{227}

\begin{equation}\tag{2.1.1}
 \alpha = \alpha^n_{P,Q} \ : \ {}_nP \otimes {}_nQ \longrightarrow \underline{\mathrm{Tor}}_1(P,Q) \quad ,
\end{equation}
qui sera le composé

\SGAvar{\begin{equation}\tag{2.1.2}
{}_nP \otimes {}_nQ \longrightarrow \underline{\mathrm{Tor}}_1({}_nP,{}_nQ) \longrightarrow \underline{\mathrm{Tor}}_1(P,Q) \quad ,
\end{equation}
où la première flèche est $\alpha^n_{{}_nP,{}_nQ}$ , et la deuxième est celle déduite des inclusions ${}_nP \longrightarrow P$, ${}_nQ \longrightarrow Q$. On est donc ramené à définir (2.1.1) dans le cas où on a $nP = 0$, $nQ = 0$, i.e. où $P$ et $Q$ sont des $\mathbb{Z}/n\mathbb{Z}$-Modules. Supposcns seulement $nQ = 0$. Prenant une résolution plate $L.$ de $P$ (en tant que $\mathbb{Z}$-Module), on trouve des isomorphismes canoniques

}{\begin{equation}\tag{2.1.2}
{}_nP \otimes {}_nQ \longrightarrow \underline{\mathrm{Tor}}_1({}_nP,{}_nQ) \longrightarrow \underline{\mathrm{Tor}}_1(P,Q) \quad ,
\end{equation}
où la première flèche est $\alpha^n_{{}_nP,{}_nQ}$ , et la deuxième est celle déduite des inclusions ${}_nP \longrightarrow P$, ${}_nQ \longrightarrow Q$. On est donc ramené à définir (2.1.1) dans le cas où on a $nP = 0$, $nQ = 0$, i.e. où $P$ et $Q$ sont des $\mathbb{Z}/n\mathbb{Z}$-Modules. Supposons seulement $nQ = 0$. Prenant une résolution plate $L.$ de $P$ (en tant que $\mathbb{Z}$-Module), on trouve des isomorphismes canoniques

}{F7I-VIII-010}\begin{equation*}
P \overset{\mathbb{L}}{\otimes}_{\mathbb{Z}} Q \ \simeq \ L. \otimes_{\mathbb{Z}} Q \ \simeq \ (L. \otimes_{\mathbb{Z}} \mathbb{Z}/n\mathbb{Z}) \otimes_{\mathbb{Z}/n\mathbb{Z}} Q \quad ,
\end{equation*}
qu'on peut écrire comme un isomorphisme dans la catégorie dérivée

\begin{equation}\tag{2.1.3}
P \overset{\mathbb{L}}{\otimes}_{\mathbb{Z}} Q \simeq (P \overset{\mathbb{L}}{\otimes}_{\mathbb{Z}} \mathbb{Z}/n\mathbb{Z}) \overset{\mathbb{L}}{\otimes}_{\mathbb{Z}/n\mathbb{Z}} Q \quad ,
\end{equation}
donnant naissance à la suite spectrale de type homologique ``de Kunneth'' (cf. SGA 1 IV 5.2) :

\begin{equation}\tag{2.1.4}
\underline{\mathrm{Tor}}^{\mathbb{Z}}_*(P,Q) \ \Longleftarrow \ E^2_{p,q} = \underline{\mathrm{Tor}}^{\mathbb{Z}/n\mathbb{Z}}_p(\underline{\mathrm{Tor}}^{\mathbb{Z}}_q(P,\mathbb{Z}/n\mathbb{Z}),Q)
\end{equation}
Rappelons d'ailleurs qu'on a des isomorphismes canoniques

\begin{equation*}
\underline{\mathrm{Tor}}^{\mathbb{Z}}_o(P,\mathbb{Z}/n\mathbb{Z}) = P_n \ \overset{\mathrm{dfn}}{=} \ P/nP \ , \quad \underline{\mathrm{Tor}}^{\mathbb{Z}}_1(P,\mathbb{Z}/n\mathbb{Z}) \simeq {}_nP \ ,
\end{equation*}
\begin{equation*}
\underline{\mathrm{Tor}}^{\mathbb{Z}}_i(P,\mathbb{Z}/n\mathbb{Z}) = 0 \quad \text{si } i \geq 2 \ ,
\end{equation*}
déduits du calcul des $\underline{\mathrm{Tor}}^{\mathbb{Z}}_i(P,\mathbb{Z}/n\mathbb{Z})$ à l'aide de la résolution plate canonique suivante de $(\mathbb{Z}/n\mathbb{Z})_T$ :

\begin{center}
\begin{tikzcd}[column sep=large, row sep=large]
0 \arrow[r] \& \mathbb{Z}_T \arrow[r, "n"] \arrow[d, equal] \& \mathbb{Z}_T \arrow[r] \arrow[d, equal] \& 0 \\
\& \mathbb{L}_1 \& \mathbb{L}_o \& {,}
\end{tikzcd}
\end{center}














\SGApage{I}{228}

 l'augmentation $\mathbb{L}_o \longrightarrow (\mathbb{Z}/n\mathbb{Z})_T$ étant définie en envoyant $1$ sur $1$. Par suite, la suite spectrale (2.1.4) se réduit essentiellement à une suite exacte

\begin{equation}\tag{2.1.5}
0 \to \underline{\mathrm{Tor}}_2^{\mathbb{Z}/n\mathbb{Z}}(P_n,Q) \to {}_nP \otimes Q \stackrel{\alpha}{\longrightarrow} \underline{\mathrm{Tor}}_1^{\mathbb{Z}}(P,Q) \longrightarrow \underline{\mathrm{Tor}}_1^{\mathbb{Z}/n\mathbb{Z}}(P_n,Q) \to 0,
\end{equation}
où on a posé $P_n = P\otimes_{\mathbb{Z}} \mathbb{Z}/n\mathbb{Z} = P/nP$. L'homomorphisme $\alpha$ cherché est l'homomorphisme-coin $E^2_{0,1} \longrightarrow \underline{\mathrm{Tor}}_1^{\mathbb{Z}}$ qui figure dans la dernière suite exacte. Cette dernière étant évidemment fonctorielle en $P$ et en $Q$, la fonctorialité en $P$ montre que l'homomorphisme envisagé est bien le composé

\SGAvar{\begin{equation*}
{}_nP \otimes Q \longrightarrow \underline{\mathrm{Tor}}_1({}_nP,Q) \longrightarrow \underline{\mathrm{Tor}}_1(P,Q) \qquad ,
\end{equation*}
où le premier homomorphisme est l'homomorphisme $\alpha$ défini comme précédemment, mais avec $(P,Q)$ remplacé par $({}_nP,Q)$, et où le deu ième homomorphisme provient de l'inclusion ${}_nP \longrightarrow Q$.

2.1.6. Ceci définit donc bien, pour $P$ et $Q$ quelconque, un homomorphisme (2.1.1) satisfaisant à la condition de factorisation (2.1.2). De plus, pour $nQ=0$, la suite exacte (2.1.5) montre que $\alpha$ est un isomorphisme dès que $P_n$ est \underline{plat} sur $\mathbb{Z}/n\mathbb{Z}$ (ce qui est le cas en particulier si $P_n = 0$ i.e. $nP=P$ i.e. $P$ est $n$-divisible, ou si $P$ lui-même est un $\mathbb{Z}/n\mathbb{Z}$-Module plat) ou si $Q$ est un $\mathbb{Z}/n\mathbb{Z}$-Module plat.

2.1.7. Explicitons le calcul de l'homomorphisme $\alpha$ de (2.1.5). Pour le connaître, il suffit de connaître, pour toute section $q$ de $Q$ (sur un objet $S$ et $T$, disons l'objet final pour simplifier) l'homomorphisme ${}_nP \longrightarrow \underline{\mathrm{Tor}}_1^{\mathbb{Z}}(P,Q)$ correspondant. Or on peut considérer $q$ comme l'image de la section $1$ de $(\mathbb{Z}/n\mathbb{Z})_{\underline{T}}$ par un homomorphisme $u_q : (\mathbb{Z}/n\mathbb{Z})_{\underline{T}} \longrightarrow Q$ bien déterminé, et la fonctorialité de l'homomorphisme $\alpha$ en $Q$ nous
\SGApage{I}{229}
ramène donc au cas où $Q$ est le faisceau constant $(\mathbb{Z}/n\mathbb{Z})_{\underline{T}}$ , et où $q$ est la section unité de $(\mathbb{Z}/n\mathbb{Z})_{\underline{T}}$ . Dans ce cas, identifiant $\underline{\mathrm{Tor}}_1^{\mathbb{Z}}(P,Q) = \underline{\mathrm{Tor}}_1^{\mathbb{Z}}(P,\mathbb{Z}/n\mathbb{Z})$ à ${}_nP$ comme expliqué après (2.1.4) ci-dessus, l'homomorphisme cherché est l'identité, en d'autres termes, l'homomorphisme coin $E^2_{0,q} = \underline{\mathrm{Tor}}_q^{\mathbb{Z}}(P,\mathbb{Z}/n\mathbb{Z}) \otimes_{\mathbb{Z}/n\mathbb{Z}} \mathbb{Z}/n\mathbb{Z} \longrightarrow \underline{\mathrm{Tor}}_q^{\mathbb{Z}}(P,\mathbb{Z}/n\mathbb{Z})$ est l'identité. Cela résulte en effet immédiatement des définitions (et l'assertion analogue reste valable en remplaçant $\mathbb{Z}$ et son quotient $\mathbb{Z}/n\mathbb{Z}$ par n'importe quel Anneau $A$ sur $\underline{T}$ et n'importe quel Anneau-quotient $A/I$ de $A$).

}{\begin{equation*}
{}_nP \otimes Q \longrightarrow \underline{\mathrm{Tor}}_1({}_nP,Q) \longrightarrow \underline{\mathrm{Tor}}_1(P,Q) \qquad ,
\end{equation*}
où le premier homomorphisme est l'homomorphisme $\alpha$ défini comme précédemment, mais avec $(P,Q)$ remplacé par $({}_nP,Q)$, et où le deuxième homomorphisme provient de l'inclusion ${}_nP \longrightarrow P$.

2.1.6. Ceci définit donc bien, pour $P$ et $Q$ quelconques, un homomorphisme (2.1.1) satisfaisant à la condition de factorisation (2.1.2). De plus, pour $nQ=0$, la suite exacte (2.1.5) montre que $\alpha$ est un isomorphisme dès que $P_n$ est \underline{plat} sur $\mathbb{Z}/n\mathbb{Z}$ (ce qui est le cas en particulier si $P_n = 0$ i.e. $nP=P$ i.e. $P$ est $n$-divisible, ou si $P$ lui-même est un $\mathbb{Z}/n\mathbb{Z}$-Module plat) ou si $Q$ est un $\mathbb{Z}/n\mathbb{Z}$-Module plat.

2.1.7. Explicitons le calcul de l'homomorphisme $\alpha$ de (2.1.5). Pour le connaître, il suffit de connaître, pour toute section $q$ de $Q$ (sur un objet $S$ de $T$, disons l'objet final pour simplifier) l'homomorphisme ${}_nP \longrightarrow \underline{\mathrm{Tor}}_1^{\mathbb{Z}}(P,Q)$ correspondant. Or on peut considérer $q$ comme l'image de la section $1$ de $(\mathbb{Z}/n\mathbb{Z})_{\underline{T}}$ par un homomorphisme $u_q : (\mathbb{Z}/n\mathbb{Z})_{\underline{T}} \longrightarrow Q$ bien déterminé, et la fonctorialité de l'homomorphisme $\alpha$ en $Q$ nous
\SGApage{I}{229}
ramène donc au cas où $Q$ est le faisceau constant $(\mathbb{Z}/n\mathbb{Z})_{\underline{T}}$ , et où $q$ est la section unité de $(\mathbb{Z}/n\mathbb{Z})_{\underline{T}}$ . Dans ce cas, identifiant $\underline{\mathrm{Tor}}_1^{\mathbb{Z}}(P,Q) = \underline{\mathrm{Tor}}_1^{\mathbb{Z}}(P,\mathbb{Z}/n\mathbb{Z})$ à ${}_nP$ comme expliqué après (2.1.4) ci-dessus, l'homomorphisme cherché est l'identité, en d'autres termes, l'homomorphisme coin $E^2_{0,q} = \underline{\mathrm{Tor}}_q^{\mathbb{Z}}(P,\mathbb{Z}/n\mathbb{Z}) \otimes_{\mathbb{Z}/n\mathbb{Z}} \mathbb{Z}/n\mathbb{Z} \longrightarrow \underline{\mathrm{Tor}}_q^{\mathbb{Z}}(P,\mathbb{Z}/n\mathbb{Z})$ est l'identité. Cela résulte en effet immédiatement des définitions (et l'assertion analogue reste valable en remplaçant $\mathbb{Z}$ et son quotient $\mathbb{Z}/n\mathbb{Z}$ par n'importe quel Anneau $A$ sur $\underline{T}$ et n'importe quel Anneau-quotient $A/I$ de $A$).

}{F7I-VIII-011}2.1.8. On peut aussi déterminer l'homomorphisme $\alpha$ de (2.1.1) en décrivant, pour tout couple de sections $p$ de ${}_nP$ et $q$ de ${}_nQ$ (sur un objet $S$ de $\underline{T}$, que nous prendrons égal à l'objet final pour simplifier les notations), la section $(p,q)$ de $\underline{\mathrm{Tor}}_1^{\mathbb{Z}}(P,Q)$. Utilisant comme ci-dessus la fonctorialité de $\alpha$ en $P$ et en $Q$, on est donc ramené au cas où on a

\begin{equation*}
P = Q = (\mathbb{Z}/n\mathbb{Z})_{\underline{T}} \qquad ,
\end{equation*}
$p$ et $q$ étant la section unité de $(\mathbb{Z}/n\mathbb{Z})_{\underline{T}}$ , et on est donc ramené à déterminer l'homomorphisme

\SGAvar{\begin{equation*}
\alpha : (\mathbb{Z}/n\mathbb{Z})_{\underline{T}} \otimes (\mathbb{Z}/n\mathbb{Z})_{\underline{T}} \longrightarrow \underline{\mathrm{Tor}}_1^{\mathbb{Z}} (\mathbb{Z}/n\mathbb{Z})_{\underline{T}} , (\mathbb{Z}/n\mathbb{Z})_T) \qquad .
\end{equation*}
La description 2.1.7 nous montre que c'est l'isomorphisme

\begin{equation}\tag{2.1.8.1}
(\mathbb{Z}/n\mathbb{Z})_{\underline{T}} \simeq \underline{\mathrm{Tor}}_1^{\mathbb{Z}} (\mathbb{Z}/ n\mathbb{Z})_{\underline{T}} , (\mathbb{Z}/n\mathbb{Z})_T)
\end{equation}
déduit de la résolution plate canonique

}{\begin{equation*}
\alpha : (\mathbb{Z}/n\mathbb{Z})_{\underline{T}} \otimes (\mathbb{Z}/n\mathbb{Z})_{\underline{T}} \longrightarrow \underline{\mathrm{Tor}}_1^{\mathbb{Z}}((\mathbb{Z}/n\mathbb{Z})_{\underline{T}},(\mathbb{Z}/n\mathbb{Z})_T) \qquad .
\end{equation*}
La description 2.1.7 nous montre que c'est l'isomorphisme

\begin{equation}\tag{2.1.8.1}
(\mathbb{Z}/n\mathbb{Z})_{\underline{T}} \simeq \underline{\mathrm{Tor}}_1^{\mathbb{Z}}((\mathbb{Z}/ n\mathbb{Z})_{\underline{T}},(\mathbb{Z}/n\mathbb{Z})_T)
\end{equation}
déduit de la résolution plate canonique

}{F7I-VIII-012}\begin{equation}\tag{2.1.8.2}
0 \longrightarrow \mathbb{Z}_T \stackrel{n}{\longrightarrow} \mathbb{Z}_T \longrightarrow 0
\end{equation}

\SGApage{I}{230}

 du \underline{premier} facteur, $P = (\mathbb{Z}/n\mathbb{Z})_{\underline{T}}$ , qui figure dans le $\underline{Tor}_1$ .

2.1.9. La définition donnée de (2.1.1) n'est pas symétrique, et on peut en échangeant les rôles de $P$ et de $Q$ dans la construction précédente, définir un homomorphisme analogue

\begin{equation}\tag{2.1.9.1}
\beta^n_{P,Q} : {}_nP \otimes {}_nQ \longrightarrow \underline{Tor}_1(P,Q) \quad ,
\end{equation}
en se ramenant d'abord au cas $nP=0$, et utilisant une résolution plate \underline{de} $Q$. On peut aussi définir $\beta$ en termes de $\alpha$ , de façon équivalente, par la commutativité du diagramme

\begin{equation}\tag{2.1.9.2}
\begin{tikzcd}[column sep=huge, row sep=huge]
{}_nP \otimes {}_nQ \arrow[r, "\sim"] \arrow[d, "{\beta^n_{P,Q}}"'] \& {}_nQ \otimes {}_nP \arrow[d, "{\alpha^n_{Q,P}}"] \\
\underline{Tor}_1(P,Q) \arrow[r, "\sim"] \& \underline{Tor}_1(Q,P)
\end{tikzcd}\quad ,
\end{equation}















où les flèches horizontales sont les isomorphismes de symétrie. Ceci posé, les relations entre $\alpha$ et $\beta$ sont données par la formule :

\begin{equation}\tag{2.1.9.3}
\beta^n_{P,Q} \;=\; - \,\alpha^n_{P,Q} \quad ,
\end{equation}
en d'autres termes :

\underline{Lemme} 2.1.10. \underline{Soient} $P$ \underline{et} $Q$ \underline{deux Groupes commutatifs,} $n > 0$ \underline{un entier}. \underline{Alors le diagramme suivant est} anticommutatif, \underline{où} $\alpha$ \underline{est défini dans} (2.1.1) :

\begin{equation}\tag{2.1.10.1}
\begin{tikzcd}[column sep=huge, row sep=huge]
{}_nP \otimes {}_nQ \arrow[rr, "\sim"] \arrow[dd, "{\alpha^n_{P,Q}}"'] \& \& {}_nQ \otimes {}_nP \arrow[dd, "{\alpha^n_{Q,P}}"] \\
 \& - \& \\
\underline{Tor}_1(P,Q) \arrow[rr, "\sim"] \& \& \underline{Tor}_1(Q,P)
\end{tikzcd}\quad ,
\end{equation}

















\underline{les flèches horizontales étant les isomorphismes de symétrie.}


\SGApage{I}{231}

\underline{Démonstration.} Le procédé de 2.1.8 nous ramène aussitôt au cas où on a $P=Q = (\mathbb{Z}/n\mathbb{Z})_{\underline{T}}$ , et à prouver alors que les deux isomorphismes (2.1.8.1) obtenus en résolvant, par (2.1.8.2), le premier et le deuxième argument respectivement intervenant dans le $\underline{Tor}_1$ , on trouve des isomorphismes qui diffèrent par le signe (ou encore, que l'isomorphisme de symétrie sur $\underline{Tor}_1^{\mathbb{Z}}(\mathbb{Z}/n\mathbb{Z} , \mathbb{Z}/n\mathbb{Z})$ est $-\mathrm{id}$). On peut évidemment supposer pour ceci que $\underline{T}$ est le topos ponctuel, i.e qu'on travaille dans la catégorie des groupes commutatifs ordinaires. Or soit $L.$ le complexe

\begin{equation*}
L. \;=\; \qquad 0 \longrightarrow \mathbb{Z} \;\xrightarrow{\;n\;}\; \mathbb{Z} \longrightarrow 0 \quad ,
\end{equation*}
considéré comme résolution plate de $P = Q = \mathbb{Z}/n\mathbb{Z}$ . Utilisons $L.\otimes L.$ pour le calcul de $Tor_1(P,Q) = Tor_1(\mathbb{Z}/n\mathbb{Z} , \mathbb{Z}/n\mathbb{Z})$, alors les isomorphismes avec le $H_1$ de $L.\otimes\mathbb{Z}/n\mathbb{Z}$ (qui permet de définir $\alpha$) et avec le $H_1$ de $\mathbb{Z}/n\mathbb{Z}\otimes L.$ (qui permet de définir $\beta$) sont définis respectivement par l'augmentation $L. \longrightarrow \mathbb{Z}/n\mathbb{Z}$ du deuxième et du premier facteur de $L.\otimes L.$ . Il s'ensuit aussitôt que $\alpha(1,1)$ est représenté par le cycle $- 1_o\otimes 1_1 + 1_1\otimes 1_o \in (L.\otimes L.)_1 = L_o\otimes L_1 + L_1\otimes L_o$(NB on désigne par $1_o$ et $1_1$ respectivement l'élément $1$ de $L_o$ et de $L_1$), tandis que $\beta(1,1)$ est représenté par le cycle $1_o\otimes 1_1 - 1_1\otimes 1_o$ . Cela prouve donc 2.1.10.

Pour terminer ces généralités sur les accouplements $\alpha$, il faut encore préciser les relations entre les divers homomorphismes $\alpha^n_{P,Q}$ , pour $n$ variable. Le résultat principal est donné dans le

\underline{Lemme} 2.1.11. \underline{Soient} $P,Q$ \underline{deux Groupes commutatifs de} $\underline{T}$, $n$ \underline{et} $n'$ \underline{deux entiers} $> 0$ \underline{tels que} $n'$ \underline{soit un multiple de} $n$, \underline{soit} $n' = mn$. \underline{Alors on a commutativité dans le diagramme suivant}
\SGApage{I}{232}

\begin{equation}\tag{2.1.11.1}
\begin{tikzcd}[column sep=large, row sep=large]
 \& {{}_{n'}P \otimes {}_{n}Q} \arrow[ld, "P_m \otimes \mathrm{id}_{{}_{n}Q}"'] \arrow[rd, "\mathrm{id}_{{}_{n'}P} \otimes i_m"] \& \\
{{}_{n}P \otimes {}_{n}Q} \arrow[rd, "{\alpha^{n}_{P,Q}}"'] \& \& {{}_{n'}P \otimes {}_{n'}Q} \arrow[ld, "{\alpha^{n'}_{P,Q}}"] \\
 \& {\underline{\mathrm{Tor}}_1(P,Q)} \&
\end{tikzcd}\quad ,
\end{equation}















\SGAvar{\underline{où} $P_m : {}_{n'}P \longrightarrow {}_{n}P$ \underline{désigne l'homomorphisme induit par} $m.\mathrm{id}_P$, \underline{et} $i_m : {}_{n}Q \longrightarrow {}_{n'}Q$ \underline{désigne l'inclusion.} \underline{On a de même commutativité dans le} \underline{diagramme analogue} (2.11.2), \underline{symétrique de} (2.11.1), \underline{dont les sommets} \underline{sont les mêmes que ceux de} (2.11.1), \underline{sauf le sommet supérieur} ${}_{n'}P \otimes {}_{n}Q$, \underline{qui est remplacé par} ${}_{n}P \otimes {}_{n'}Q$ .

}{\underline{où} $P_m : {}_{n'}P \longrightarrow {}_{n}P$ \underline{désigne l'homomorphisme induit par} $m.\mathrm{id}_P$, \underline{et} $i_m : {}_{n}Q \longrightarrow {}_{n'}Q$ \underline{désigne l'inclusion.} \underline{On a de même commutativité dans le} \underline{diagramme analogue} (2.1.11.2), \underline{symétrique de} (2.1.11.1), \underline{dont les sommets} \underline{sont les mêmes que ceux de} (2.1.11.1), \underline{sauf le sommet supérieur} ${}_{n'}P \otimes {}_{n}Q$, \underline{qui est remplacé par} ${}_{n}P \otimes {}_{n'}Q$ .

}{F7I-VIII-013}Il suffit en effet de prouver le premier énoncé, le deuxième s'en déduit par symétrie, compte tenu de 2.1.10. On est ramené, par le procédé de 2.1.8, au cas où $\underline{T}$ est le topos ponctuel, i.e. où $P$ et $Q$ sont des groupes commutatifs ordinaires, et où on a $P = \mathbb{Z}/n'\mathbb{Z}$ , $Q = \mathbb{Z}/n'\mathbb{Z}$ . En fait, on utilisera pour $P$ la seule hypothèse $n'P=0$. Nous identifierons $\mathrm{Tor}_1(P,Q)$ à ${}_{n'}P = P$, grâce à la résolution $0 \longrightarrow \mathbb{Z} \overset{n'}{\longrightarrow} \mathbb{Z} \longrightarrow 0$ de $Q$, et de même $\mathrm{Tor}_1({}_{n}P,{}_{n}Q)$ à ${}_{n}P$ grâce à la résolution $0 \longrightarrow \mathbb{Z} \overset{n}{\longrightarrow} \mathbb{Z} \longrightarrow 0$ de ${}_{n}Q \simeq \mathbb{Z}/n\mathbb{Z}$ , dont l'augmentation est donnée en envoyant $1_o$ sur $m.1_Q$ ($1_Q$ désignant l'élément $1$ de $Q = \mathbb{Z}/n'\mathbb{Z}$ ). L'homomorphisme

\begin{center}
\begin{tikzcd}[column sep=large, row sep=large]
{\mathrm{Tor}_1({}_{n}P,{}_{n}Q)} \arrow[r] \arrow[d, no head, "\wr"'] \& {\mathrm{Tor}_1(P,Q)} \arrow[d, no head, "\wr"'] \\
{{}_{n}P} \& {{}_{n'}P = P}
\end{tikzcd}
\end{center}










\SGAvar{\SGApage{I}{233}
déduit des inclusions ${}_{n}P \hookrightarrow P$ et ${}_{n}Q \hookrightarrow Q$ n'est alors autre que l'inclusion naturelle ${}_{n}P \hookrightarrow {}_{n'}P$. Par suite, identifiant le sommetsupérieur de (2.1.11.1) à $P/nP$, grâce au générateur $m.1_Q$ de ${}_{n}Q$, permettant d'identifier ${}_{n}Q$ à $\mathbb{Z}/n\mathbb{Z}$ , le diagramme envisagé devient

}{\SGApage{I}{233}
déduit des inclusions ${}_{n}P \hookrightarrow P$ et ${}_{n}Q \hookrightarrow Q$ n'est alors autre que l'inclusion naturelle ${}_{n}P \hookrightarrow {}_{n'}P$. Par suite, identifiant le sommet supérieur de (2.1.11.1) à $P/nP$, grâce au générateur $m.1_Q$ de ${}_{n}Q$, permettant d'identifier ${}_{n}Q$ à $\mathbb{Z}/n\mathbb{Z}$ , le diagramme envisagé devient

}{F7I-VIII-014}\begin{center}
\begin{tikzcd}[column sep=large, row sep=large]
 \& {P/nP} \arrow[ld, "x \longmapsto mx"'] \arrow[rd, "x \longmapsto mx"] \& \\
{{}_{n}P} \arrow[rd, "x \longmapsto x"'] \& \& {P} \arrow[ld, "x \longmapsto x"] \\
 \& {{}_{n'}P} \&
\end{tikzcd}
\end{center}









qui est évidemment commutatif.

\underline{Corollaire} 2.1.12. \underline{Avec les notations de} 2.1.11, \underline{on a commutativité dans} \underline{le diagramme suivant}

\begin{equation}\tag{2.1.12.1}
\begin{tikzcd}[column sep=huge, row sep=huge]
{{}_{n'}P \otimes {}_{n'}Q} \arrow[r, "m \otimes m"] \arrow[d, "{\alpha^{n'}_{P,Q}}"'] \& {{}_{n}P \otimes {}_{n}Q} \arrow[d, "{\alpha^{n}_{P,Q}}"] \\
{{}_{n'}\underline{\mathrm{Tor}}_1(P,Q)} \arrow[r, "m"] \& {{}_{n}\underline{\mathrm{Tor}}_1(P,Q)}
\end{tikzcd}\quad ,
\end{equation}












\underline{où}, \underline{dans la description des flèches horizontales}, $m$ \underline{désigne les homomorphismes} ${}_{n'}E \longrightarrow {}_{n}E$ \underline{induits par} $m.\mathrm{id}_E$ (\underline{pour} $E$ \underline{égal respectivement à} $P,Q$ \underline{et à} $\underline{\mathrm{Tor}}_1(P,Q))$.

En effet, en vertu de 2.1.11 on aura, pour $x \in {}_{n'}P(S)$, $y \in {}_{n'}Q(S)$ : $\alpha^n(mx \otimes my) = \alpha^{n'}(x \otimes my)$, or le deuxième membre est égal à $m\alpha^{n'}(x \otimes y)$ puisque $x \otimes my = m(x \otimes y)$, d'où la commutativité de (2.1.12.1).


\SGApage{I}{234}

2.1.13. Soit $\ell$ un nombre premier. Pour tout Groupe commutatif $E$, désignons par $T_\ell(E)$ le système projectif

\begin{equation}\tag{2.1.13.1}
T_\ell(E) = ({}_{\ell^\nu}E)_{\nu \geq 0} \qquad ,
\end{equation}
qui dépend évidemment fonctoriellement de $E$. Alors les compatibilités 2.1.12, pour $n,n'$ de la forme $\ell^\nu,\ell^{\nu'}$ avec $\nu' \geq \nu$, nous montrent que les $\alpha^{\ell^\nu}_{P,Q}$ pour $\nu$ variable définissent un homomorphisme de systèmes projectifs, fonctoriel en $P,Q$ :

\begin{equation}\tag{2.1.13.2}
\alpha^{(\ell)}_{P,Q} : T_\ell(P) \otimes T_\ell(Q) \longrightarrow T_\ell(\underline{\mathrm{Tor}}_1(P,Q)) \qquad .
\end{equation}
La connaissance de ces homomorphismes, pour tout nombre premier $\ell$, équivaut évidemment à la connaissance des $\alpha^n_{P,Q}$ pour tout entier $n$. Les relations d'antisymétrie 2.1.10 se traduisent alors par les relations d'antisymétrie analogues $\ell$-adiques, pour tout $\ell$, exprimées par l'anticommutativité des diagrammes

\begin{equation}\tag{2.1.13.3}
\begin{tikzcd}[column sep=huge, row sep=huge]
T_\ell(P) \otimes T_\ell(Q) \arrow[r] \arrow[d, "{\alpha^{(\ell)}_{P,Q}}"'] \arrow[rd, phantom, "-" description] \& T_\ell(Q) \otimes T_\ell(P) \arrow[d, "{\alpha^{(\ell)}_{Q,P}}"] \\
T_\ell(\underline{\mathrm{Tor}}_1(P,Q)) \arrow[r] \& T_\ell(\underline{\mathrm{Tor}}_1(Q,P))
\end{tikzcd}\quad ,
\end{equation}











où les flèches horizontales sont les flèches de symétrie.

\SGAvar{\underline{Remarque} 2.1.14. On peut généraliser la construction de l'homomorphisme (2.1.1) et les résultats 2.1.10 et 2.1.11 le concernant au cas où on se donne un Anneau commutatif $\underline{A}$ sur $T$, et deux Idéaux $\underline{a}$ et $\underline{b}$ de $\underline{A}$ (qui dans le cas traités précédemment se réduisent au même Idéal $n\mathbb{Z}_T$ de
\SGApage{I}{235}
l'Anneau constant $\mathbb{Z}_T$ sur $T$). On définit alors un homomorphisme fonctoriel en les $\underline{A}$-Modules $P,Q$ :

}{\underline{Remarque} 2.1.14. On peut généraliser la construction de l'homomorphisme (2.1.1) et les résultats 2.1.10 et 2.1.11 le concernant au cas où on se donne un Anneau commutatif $\underline{A}$ sur $T$, et deux Idéaux $\underline{a}$ et $\underline{b}$ de $\underline{A}$ (qui dans le cas traité précédemment se réduisent au même Idéal $n\mathbb{Z}_T$ de
\SGApage{I}{235}
l'Anneau constant $\mathbb{Z}_T$ sur $T$). On définit alors un homomorphisme fonctoriel en les $\underline{A}$-Modules $P,Q$ :

}{F7I-VIII-015}\begin{equation}\tag{2.1.14.1}
{}_{\underline{a}}P \otimes {}_{\underline{b}}Q \ \overset{\alpha}{\longrightarrow}\ \underline{\mathrm{Hom}}_A(\underline{a} \cap \underline{b}/\underline{a}.\underline{b}\ ,\ \underline{\mathrm{Tor}}^A_1(P,Q)) \qquad ,
\end{equation}
comme composé des homomorphismes

\begin{equation*}
{}_{\underline{a}}P \otimes {}_{\underline{b}}Q \to \underline{\mathrm{Hom}}(\underline{A}/\underline{a},P) \otimes \underline{\mathrm{Hom}}(\underline{A}/\underline{b},Q) \longrightarrow \underline{\mathrm{Hom}}(\underline{\mathrm{Tor}}_1(\underline{A}/\underline{a},\underline{A}/\underline{b}),\underline{\mathrm{Tor}}_1(P,Q)) \qquad ,
\end{equation*}
compte tenu de l'isomorphisme canonique bien connu

\begin{equation}\tag{2.1.14.2}
\underline{\mathrm{Tor}}_1(\underline{A}/\underline{a},\underline{A}/\underline{b}) \ \simeq\ \underline{a} \cap \underline{b}/\underline{a}.\underline{b} \qquad ;
\end{equation}
en fait, c'est dans le choix de cet isomorphisme, suivant qu'on travaille avec la résolution $0 \longrightarrow \underline{a} \longrightarrow \underline{A} \longrightarrow 0$ de $\underline{A}/\underline{a}$ ou la résolution $0 \longrightarrow \underline{b} \to \underline{A} \to 0$ de $\underline{A}/\underline{b}$, qu'est ``caché'' le signe de 2.1.10, qui se généralise ici en l'anticommutativité du diagramme

\begin{equation}\tag{2.1.14.3}
\begin{tikzcd}[column sep=huge, row sep=huge]
{{}_{\underline{a}}P \otimes {}_{\underline{b}}Q} \arrow[r] \arrow[d] \arrow[rd, phantom, "-1" description] \& {{}_{\underline{b}}Q \otimes {}_{\underline{a}}P} \arrow[d] \\
{\underline{\mathrm{Hom}}(\underline{a} \cap \underline{b}/\underline{a}.\underline{b},\underline{\mathrm{Tor}}_1(P,Q))} \arrow[r] \& {\underline{\mathrm{Hom}}(\underline{a} \cap \underline{b}/\underline{a}.\underline{b},\underline{\mathrm{Tor}}_1(Q,P))}
\end{tikzcd}\quad ,
\end{equation}












où les flèches horizontales sont les flèches de symétrie. Lorsque $\underline{a}=\underline{b}$, cette anticommutativité peut aussi s'exprimer par le fait que l'automorphisme de symétrie dans $\underline{\mathrm{Tor}}_1(\underline{A}/\underline{a},\underline{A}/\underline{a})$ est égal à $-\mathrm{id}$ . Dans le cas particulier où $\underline{a} = \underline{b} = n\underline{A}$, $n$ étant une section de $\underline{A}$, alors la multiplication par $n$ induit un épimorphisme

\SGAvar{\begin{equation*}
\underline{A}/\underline{a} \longrightarrow \underline{a} \cap \underline{a}/\underline{a}.\underline{a} = \underline{a}/\underline{a}^2 \qquad ,
\end{equation*}
d'où un isomorphisme injectif fonctoriel en le $\underline{A}$-Module $T$ 


}{\begin{equation*}
\underline{A}/\underline{a} \longrightarrow \underline{a} \cap \underline{a}/\underline{a}.\underline{a} = \underline{a}/\underline{a}^2 \qquad ,
\end{equation*}
d'où un homomorphisme injectif fonctoriel en le $\underline{A}$-Module $T$ 


}{F7I-VIII-016}\SGApage{I}{236}

\begin{equation*}
 \underline{\mathrm{Hom}}(\underline{a}/\underline{a}^2 , T) \longrightarrow {}_{\underline{a}}T \quad ,
\end{equation*}
de sorte que dans ce cas l'homomorphisme (2.1.14.1) s'interprète comme un homomorphisme

\begin{equation}\tag{2.1.14.4}
{}_nP \otimes {}_nQ \longrightarrow {}_n\underline{\mathrm{Tor}}_1(P,Q) \quad ,
\end{equation}
plus proche de la forme (2.1.1). Supposant toujours $\underline{a} = \underline{b} = n\underline{A}$, et $n$ une section \underline{régulière} de $\underline{A}$, la suite spectrale (2.1.4) garde un sens dans le cas présent, et donne lieu à une suite exacte de termes de bas degré qui généralise (2.1.5) ($Q$ étant maintenant un $\underline{A}$-Module \underline{annulé par} $\underline{a}$) :

\begin{equation}\tag{2.1.14.5}
\underline{\mathrm{Tor}}_2^{\underline{A}}(P,Q) \longrightarrow \underline{\mathrm{Tor}}_2^{\underline{A}/n\underline{A}}(P_n,Q) \longrightarrow {}_nP \otimes Q \xrightarrow{\ \alpha\ } \underline{\mathrm{Tor}}_1^{\underline{A}}(P,Q) \longrightarrow \underline{\mathrm{Tor}}_1^{\underline{A}/n\underline{A}}(P_n,Q) \longrightarrow 0.
\end{equation}
Revenons au cas où $\underline{a}$ et $\underline{b}$ sont quelconques, et où on se donne deux sections $p$ et $q$ de $\underline{A}$, on trouve alors commutativité dans le diagramme suivant (dont (2.1.11.1) est essentiellement le cas particulier obtenu en prenant $q=1$) :

\begin{equation}\tag{2.1.14.6}
{\footnotesize \begin{tikzcd}[row sep=large, column sep=tiny]
\& {}_{p\underline{a}}P \otimes {}_{q\underline{b}}Q \arrow[dl, "p{\otimes}q"'] \arrow[dr, "i{\otimes}j"] \& \\
{}_{\underline{a}}P \otimes {}_{\underline{b}}Q \arrow[d, "\alpha"] \& \& {}_{pq\underline{a}}P \otimes {}_{pq\underline{b}}Q \arrow[d, "\alpha"] \\
\underline{\mathrm{Hom}}(\underline{a}\cap\underline{b}/\underline{a}.\underline{b},\underline{\mathrm{Tor}}_1(P,Q)) \arrow[dr, hook, "\mathrm{incl.}"'] \& \& \underline{\mathrm{Hom}}(pq\underline{a}\cap pq\underline{b} \,/\, pq\underline{a}.pq\underline{b},\ \underline{\mathrm{Tor}}_1(P,Q)) \arrow[dl, hook, "(pq)^{*}"] \\
\& \underline{\mathrm{Hom}}(\underline{a}\cap\underline{b}/pq\underline{a}.\underline{b},\underline{\mathrm{Tor}}_1(P,Q)) \&
\end{tikzcd}}\quad ,
\end{equation}



















\SGAvar{\SGApage{I}{237}
où $i : {}_{p\underline{a}}P \longrightarrow {}_{pq\underline{a}}P$ et $j : {}_{q\underline{b}}Q \longrightarrow {}_{pq\underline{b}}Q$ sont les inclusions canoniques, et $(pq)^{*}$ désigne l'homomorphisme injectif déduit par passage aux \underline{Hom} de l'homomorphisme surjectif $\underline{a}\cap\underline{b}/pq\ \underline{a}.\underline{b} \longrightarrow pq\ \underline{a}\cap\underline{b}/pq\underline{a}.pq\underline{b}$ induit par la multiplication par $pq$ . Lorsque $\underline{a} = \underline{b} = n\underline{A}$, ce diagramme se réécrit sous une forme plus simple (donnée dans (2.1.11.1) pour $q = 1$),

}{\SGApage{I}{237}
où $i : {}_{p\underline{a}}P \longrightarrow {}_{pq\underline{a}}P$ et $j : {}_{q\underline{b}}Q \longrightarrow {}_{pq\underline{b}}Q$ sont les inclusions canoniques, et $(pq)^{*}$ désigne l'homomorphisme injectif déduit par passage aux \underline{Hom} de l'homomorphisme surjectif $\underline{a}\cap\underline{b}/pq\ \underline{a}.\underline{b} \longrightarrow pq\ \underline{a}\cap pq\ \underline{b}/pq\underline{a}.pq\underline{b}$ induit par la multiplication par $pq$ . Lorsque $\underline{a} = \underline{b} = n\underline{A}$, ce diagramme se réécrit sous une forme plus simple (donnée dans (2.1.11.1) pour $q = 1$),

}{F7I-VIII-017}2.2. Considérons maintenant, en plus de $P$ et $Q$, un Groupe commutatif $G$. Nous avons alors défini dans (1.1.3) un homomorphisme canonique, fonctoriel en $P,Q$ et $G$ :

\begin{equation}\tag{2.2.1}
\mathrm{Biext}^1(P,Q;G) \longrightarrow \mathrm{Hom}(\mathrm{Tor}_1(P,Q),G) \quad .
\end{equation}
D'autre part, si $n$ est un entier $> 0$ (resp. si $\ell$ est un nombre premier) alors l'homomorphisme $\alpha^n_{P,Q}$ de (2.1.1) (resp. l'homomorphisme $\alpha^{(\ell)}_{P,Q}$ de (2.1.13.2)) donne un homomorphisme

\begin{equation*}
\mathrm{Hom}(\mathrm{Tor}_1(P,Q),G) \longrightarrow \mathrm{Hom}({}_nP \otimes {}_nQ \, , \, G)
\end{equation*}
(resp. un homomorphisme

\begin{equation*}
\mathrm{Hom}(\mathrm{Tor}_1(P,Q),G) \longrightarrow \mathrm{Hom}(T_{\ell}(P) \otimes T_{\ell}(Q), \ T_{\ell}(G)) \ ).
\end{equation*}
Ces derniers homomorphismes sont également fonctoriels en $P,Q,G$. Les composant avec (2.2.1), on en déduit des homomorphismes canoniques, fonctoriels en $P,Q,G$ :

\begin{equation}\tag{2.2.2}
\varphi^n = \varphi^n_{P,Q} : \mathrm{Biext}^1(P,Q;G) \longrightarrow \mathrm{Hom}({}_nP{\otimes}{}_nQ, \ G) \quad ,
\end{equation}
\begin{equation}\tag{2.2.3}
\varphi^{(\ell)} = \varphi^{(\ell)}_{P,Q} : \mathrm{Biext}^1(P,Q;G) \longrightarrow \mathrm{Hom}(T_{\ell}(P){\otimes}T_{\ell}(Q), \ T_{\ell}(G)) \quad .
\end{equation}
Pour

\begin{equation}\tag{2.2.4}
\xi \in \mathrm{Biext}^1(P,Q;G) \quad ,
\end{equation}
\SGApage{I}{238}
on désigne par $\varphi_\xi^n$ (resp. $\varphi_\xi^{(\ell)}$) son image par (2.2.2) (resp. (2.2.3)) :

\begin{equation}\tag{2.2.5}
\varphi_\xi^n \; : \; {}_nP \otimes {}_nQ \longrightarrow G \;\; (\text{ou } {}_nG, \text{ au choix}) \qquad ,
\end{equation}
\begin{equation}\tag{2.2.6}
\varphi_\xi^{(\ell)} : \; T_\ell(P) \otimes T_\ell(Q) \longrightarrow T_\ell(G) \qquad .
\end{equation}
2.2.7.  Evidemment la connaissance de $\varphi_\xi^{(\ell)}$ équivaut à celle des $\varphi_\xi^{(n)}$ pour tout $n$ de la forme $\ell^\nu$ , avec $\nu \geq 0$, et pour $n$ fixé, la connaissance de $\varphi_\xi^{(n)}$ équivaut à celle de l'homomorphisme

\SGAvar{\begin{equation}\tag{2.2.8}
\underline{\mathrm{Tor}}_1(P,Q) \longrightarrow G \qquad ,
\end{equation}
image de $\xi$ par (2.2.1), sur l'image de $\varphi^n_{P,Q} : {}_nP \otimes {}_nQ \longrightarrow \underline{\mathrm{Tor}}_1(P,Q;G)$. On notera que lorsque $P=nP$ et $Q=nQ$, alors l'image précédente est égale à ${}_n\underline{\mathrm{Tor}}_1(P,Q)$, comme il résulte aisément de la considération de la suite exacte (1.3.5) et de la suite exacte analogue pour $Q$ ; d'autre part, si $P$ et $Q$ sont divisibles i.e. si $nP=P,nQ=Q$ pour tout entier $n > 0$, alors on vérifie aisément que $\underline{\mathrm{Tor}}_1(P,Q)$ est de torsion i.e. $\underline{\mathrm{Tor}}_1(P,Q) = \underset{n}{\varinjlim}\; {}_n\underline{\mathrm{Tor}}_1(P,Q)$, de sorte que dans ce cas la connaissance des $\varphi_\xi^n$ pour tout $n$ équivaut à celle de (2.2.8), tandis que la connaissance de $\varphi_\xi^{(\ell)}$, pour un nombre premier donné, équivaut à celle de la restriction de (2.2.8) à la composante $\ell$-primaire de $\mathrm{Tor}\; {}_1(P,Q)$. Nous allons expliciter maintenant les propriétés de symétrie des accouplements (2.2.4), (2.2.6) associés aux biextensions. Notons d'abord qu'on a un diagramme \underline{commutatif}

}{\begin{equation}\tag{2.2.8}
\underline{\mathrm{Tor}}_1(P,Q) \longrightarrow G \qquad ,
\end{equation}
image de $\xi$ par (2.2.1), sur l'image de $\alpha^n_{P,Q} : {}_nP \otimes {}_nQ \longrightarrow \underline{\mathrm{Tor}}_1(P,Q)$. On notera que lorsque $P=nP$ et $Q=nQ$, alors l'image précédente est égale à ${}_n\underline{\mathrm{Tor}}_1(P,Q)$, comme il résulte aisément de la considération de la suite exacte (1.3.5) et de la suite exacte analogue pour $Q$ ; d'autre part, si $P$ et $Q$ sont divisibles i.e. si $nP=P,nQ=Q$ pour tout entier $n > 0$, alors on vérifie aisément que $\underline{\mathrm{Tor}}_1(P,Q)$ est de torsion i.e. $\underline{\mathrm{Tor}}_1(P,Q) = \underset{n}{\varinjlim}\; {}_n\underline{\mathrm{Tor}}_1(P,Q)$, de sorte que dans ce cas la connaissance des $\varphi_\xi^n$ pour tout $n$ équivaut à celle de (2.2.8), tandis que la connaissance de $\varphi_\xi^{(\ell)}$, pour un nombre premier donné, équivaut à celle de la restriction de (2.2.8) à la composante $\ell$-primaire de $\mathrm{Tor}\; {}_1(P,Q)$. Nous allons expliciter maintenant les propriétés de symétrie des accouplements (2.2.5), (2.2.6) associés aux biextensions. Notons d'abord qu'on a un diagramme \underline{commutatif}

}{F7I-VIII-018}\begin{equation}\tag{2.2.9}
\begin{tikzcd}[column sep=large, row sep=huge]
\mathrm{Biext}^1(P,Q;G) \arrow[r, "\sim"] \arrow[d] \& \mathrm{Biext}^1(Q,P;G) \arrow[d] \\
\mathrm{Hom}(\underline{\mathrm{Tor}}_1(P,Q),G) \arrow[r, "\sim"] \& \mathrm{Hom}(\underline{\mathrm{Tor}}_1(Q,P),G) \& ,
\end{tikzcd}
\end{equation}


















\SGApage{I}{239}

 où les flèches horizontales sont les flèches de symétrie (la première de ces flèches étant définie dans VII 2.7 ), et les flèches verticales sont (2.2.1). Cette commutativité résulte en effet de la commutativité des deux carrés composants du diagramme plus complet

\begin{equation}\tag{2.2.10}
\begin{tikzcd}[column sep=large, row sep=large]
\mathrm{Biext}^1(P,Q;G) \arrow[r, "\sim"] \arrow[d] \& \mathrm{Biext}^1(Q,P;G) \arrow[d] \\
\mathrm{Ext}^1(P \overset{\mathbb{L}}{\otimes} Q,G) \arrow[r, "\sim"] \arrow[d] \& \mathrm{Ext}^1(Q \overset{\mathbb{L}}{\otimes} P,G) \arrow[d] \\
\mathrm{Hom}(\underline{\mathrm{Tor}}_1(P,Q),G) \arrow[r, "\sim"] \& \mathrm{Hom}(\underline{\mathrm{Tor}}_1(Q,P),G) \& ,
\end{tikzcd}
\end{equation}

















où la première flèche verticale de chaque colonne est l'isomorphisme canonique de VII 3.6.5 , et la deuxième flèche est l'homomorphisme évident. La commutativité du deuxième carré de (2.2.10) est triviale, compte tenu que l'isomorphisme de symétrie $P \overset{\mathbb{L}}{\otimes} Q \xrightarrow{\;\sim\;} Q \overset{\mathbb{L}}{\otimes} P$ induit par définition l'isomorphisme de symétrie $\underline{\mathrm{Tor}}_1(P,Q) \simeq \underline{\mathrm{Tor}}_1(Q,P)$, et que pour un complexe variable $L.$ , l'homomorphisme canonique $\mathrm{Ext}^1(L.,G) \simeq \mathrm{Hom}(H_1(L.),G)$ est fonctoriel en $L.$ . Quant à la commutativité du premier carré de (2.2.10), elle résulte de l'explicitation simultanée des isomorphismes de symétrie et des isomorphismes verticaux définis dans VII 3.6.5 , et nous laissons le détail de la vérification au lecteur ; on utilisera le fait que l'isomorphisme de symétrie $L.(P) \otimes L.(Q) \simeq L.(Q) \otimes L.(P)$ induit en degré 1 l'isomorphisme

\begin{equation*}
L_o(P) \otimes L_1(Q) + L_1(P) \otimes L_o(Q) \simeq L_o(Q) \otimes L_1(P) + L_1(Q) \otimes L_o(P)
\end{equation*}
qui est donné sur les sommandes correspondants par l'isomorphisme de
\SGApage{I}{240}
symétrie \underline{ordinaire} des modules non gradués sous-jacents (donc sans intervention d'un signe).

Compte tenu de la commutativité de (2.2.9), l'anticommutativité de (2.1.10.1) se traduit, en termes des $\varphi^n$ ou des $\varphi^{(\ell)}$, par la relation suivante :

\underline{Proposition} 2.2.11. \underline{On a anticommutativité dans les diagrammes suivants} :

\begin{equation}\tag{2.2.11.1}
\begin{tikzcd}[row sep=large, column sep=huge]
\mathrm{Biext}^1(P,Q;G) \arrow[rr, "\sim"] \arrow[dd, "{\varphi^n_{P,Q}}"'] \& \& \mathrm{Biext}^1(Q,P;G) \arrow[dd, "{\varphi^n_{Q,P}}"] \\
\& - \& \\
\mathrm{Hom}({}_nP \otimes {}_nQ,G) \arrow[rr, "\sim"] \& \& \mathrm{Hom}({}_nQ \otimes {}_nP\ ,\ G) \quad ,
\end{tikzcd}
\end{equation}












\SGAvar{\begin{equation}\tag{2.2.11.2}
\begin{tikzcd}[row sep=large, column sep=huge]
\mathrm{Biext}^1(P,Q;G) \arrow[rr, "\sim"] \arrow[dd, "{\varphi^{(\ell)}_{P,Q}}"'] \& \& \mathrm{Biext}^1(Q,P;G) \arrow[dd, "{\varphi^{(\ell)}_{Q,P}}"] \\
\& - \& \\
\mathrm{Hom}(T_\ell(P)\otimes T_\ell(Q),T_\ell(G) \arrow[rr, "\sim"] \& \& \mathrm{Hom}(T_\ell(Q)\otimes T_\ell(P),T_\ell(G)) \quad ,
\end{tikzcd}
\end{equation}













}{\begin{equation}\tag{2.2.11.2}
\begin{tikzcd}[row sep=large, column sep=huge]
\mathrm{Biext}^1(P,Q;G) \arrow[rr, "\sim"] \arrow[dd, "{\varphi^{(\ell)}_{P,Q}}"'] \& \& \mathrm{Biext}^1(Q,P;G) \arrow[dd, "{\varphi^{(\ell)}_{Q,P}}"] \\
\& - \& \\
\mathrm{Hom}(T_\ell(P)\otimes T_\ell(Q),T_\ell(G)) \arrow[rr, "\sim"] \& \& \mathrm{Hom}(T_\ell(Q)\otimes T_\ell(P),T_\ell(G)) \quad ,
\end{tikzcd}
\end{equation}













}{F7I-VIII-019}\underline{où les flèches horizontales sont les flèches de symétrie}.

Signalons le corollaire important :

\underline{Corollaire} 2.2.12. \underline{Supposons} $P=Q$, \underline{et supposons que} $\xi$ \underline{soit la classe} \underline{d'une biextension} symétrique (\underline{resp.} \underline{antisymétrique}) \underline{de} $(P,Q)$ \underline{par} $G$, \underline{i}.\underline{e}. \underline{qu'on ait} ${}^s\xi = \xi$ (\underline{resp.} ${}^s\xi = -\xi$), \underline{où} ${}^s\xi$ \underline{désigne l'image de} $\xi$ \underline{par symétrie}. \underline{Alors les applications bilinéaires} $\varphi^n_\xi$(\underline{sur} ${}_nP\times{}_nP$) \underline{et} $\varphi^{(\ell)}_\xi$ (\underline{sur} $T_\ell(P)\times T_\ell(P)$) \underline{sont antisymétriques} (\underline{resp.} \underline{symétriques}).


\SGApage{I}{241}

\underline{Remarque} 2.2.13. Supposons $\xi$ \underline{symétrique}. Si $n$ (resp. $\ell$) est impair, alors il résulte de 2.2.12 que $\varphi^n_\xi$ (resp. $\varphi^{(\ell)}_\xi$)est même une forme alternée. Cette conclusion n'est pas valable en général sans la restriction $n$ resp. $\ell$ impair, comme on voit déjà dans le cas où $P$ est le groupe ordinaire $\mathbb{Z}/2\mathbb{Z}$ , et $G = \mathbb{Z}/2\mathbb{Z}$. Supposons cependant que l'on ait $2P = P$, (ce qui est le cas en particulier si $P$ est un schéma abélien, en travaillant comme d'habitude avec la topologie fppf ou fpqc), alors il est encore vrai que les formes précédentes sont alternées. Il suffit de voir que $\varphi^{2^\nu}(x,x) = 0$, or, comme $2P=P$, on peut supposer que l'élément donné $x \in {}_{2^\nu}P(S)$ est de la forme $2y$, avec $y \in {}_{2^{\nu+1}}P(S)$, et on a alors $\varphi^{2^\nu}(x,x) = 2\ \varphi^{2^{\nu+1}}(y,y) = 0$, puisque $\varphi^{2^{\nu+1}}$ est antisymétrique.

2.3. \underline{Lien avec la} "\underline{dualité de Cartier}". La description donnée dans 2.2 des accouplements définis par une biextension, via le théorème d'isomorphisme VII 3.6.5 et l'homomorphisme (2.2.1), n'est évidemment pas la définition qui avait été utilisée précédemment dans le cas des schémas abéliens (à un moment où la notion même de biextension n'existait pas encore!). Celle-ci procède via la dualité de Cartier, et nous allons donner maintenant la relation avec cette dernière.

Reprenons l'homomorphisme

\begin{equation}\tag{2.3.1}
\mathrm{Biext}^1(P,Q;G) \longrightarrow \mathrm{Hom}(P,\underline{\mathrm{Ext}}^1(Q,G))
\end{equation}
provenant de la suite exacte (1.4), qui est d'ailleurs un isomorphisme dans le cas où l'on a

\begin{equation}\tag{2.3.2}
\underline{\mathrm{Hom}}(Q,G) = 0 \quad , \quad 
\end{equation}

\SGApage{I}{242}

 condition vérifiée en particulier lorsque $Q$ est un schéma abélien sur un schéma $S$, et $G$ est le groupe multiplicatif $\underline{G}_m$ (VII 1.3.8 ). On a d'autre part un homomorphisme évident ``restriction à $n\,Q$'', $\underline{\mathrm{Ext}}^1(Q,G) \to \underline{\mathrm{Ext}}^1(nQ,G)$, d'où un homomorphisme composé

\begin{equation}\tag{2.3.4}
\mathrm{Hom}(P,\underline{\mathrm{Ext}}^1(Q,G)) \to \mathrm{Hom}(P,\underline{\mathrm{Ext}}^1(nQ,G)) \to \mathrm{Hom}({}_nP,\,{}_n\underline{\mathrm{Ext}}^1(nQ,G)) \qquad ,
\end{equation}
où la deuxième flèche est l'homomorphisme de restriction à ${}_nP$ . Composant avec (2.3.1), on trouve un homomorphisme

\begin{equation}\tag{2.3.5}
\mathrm{Biext}^1(P,Q;G) \longrightarrow \mathrm{Hom}({}_nP,\,{}_n\underline{\mathrm{Ext}}^1(nQ,G)) \qquad .
\end{equation}
D'autre part, la suite exacte

\begin{equation}\tag{2.3.6}
0 \longrightarrow {}_nQ \longrightarrow Q \xrightarrow{\;n\;} nQ \longrightarrow 0
\end{equation}
donne lieu à une suite exacte des $\underline{\mathrm{Ext}}^i(-,G)$, d'où une suite exacte courte

\begin{equation}\tag{2.3.7}
0 \longrightarrow \underline{\mathrm{Hom}}(Q,G)_n \longrightarrow \underline{\mathrm{Hom}}({}_nQ,G) \xrightarrow{\;\partial\;} {}_n\underline{\mathrm{Ext}}^1(nQ,G) \longrightarrow 0 \qquad ,
\end{equation}
où $E_n$ désigne, comme d'habitude, $E/nE$ . On en retiendra l'homomorphisme canonique surjectif

\begin{equation}\tag{2.3.8}
\underline{\mathrm{Hom}}({}_nQ,G) \xrightarrow{\hspace{3em}\partial\hspace{3em}} {}_n\underline{\mathrm{Ext}}^1(nQ,G) \qquad ,
\end{equation}
qui est d'ailleurs un isomorphisme si $\underline{\mathrm{Hom}}(Q,G)_n = 0$, a fortiori si on a (2.3.2). De (2.3.8) on déduit un homomorphisme

\begin{equation}\tag{2.3.9}
\mathrm{Hom}({}_nP,\underline{\mathrm{Hom}}({}_nQ,G)) \simeq \mathrm{Hom}({}_nP \otimes {}_nQ,G) \longrightarrow \mathrm{Hom}({}_nP,\,{}_n\underline{\mathrm{Ext}}^1(nQ,G)) \qquad ,
\end{equation}
qui est un isomorphisme si on a (2.3.2).

Tous les homomorphismes envisagés sont fonctoriels en $P,Q,G$.

Les homomorphismes (2.3.5), (2.3.9) et (2.2.2) donnent lieu à un diagramme
\SGApage{I}{243}

\SGAvar{\begin{equation}\tag{2.3.10}
\begin{tikzcd}[column sep=large, row sep=large]
 \& \mathrm{Biext}^1(P,Q;G) \arrow[ddl, "\varphi^n"'] \arrow[ddr, "(2.3.5)"] \& \\
 \& - \& \\
\mathrm{Hom}({}_nP\otimes{}_nQ,G) \arrow[rr, "(2.3.9)"] \& \& \mathrm{Hom}({}_nP,{}_n\underline{\mathrm{Ext}}^1(nQnG))
\end{tikzcd}\quad ,
\end{equation}























}{\begin{equation}\tag{2.3.10}
\begin{tikzcd}[column sep=large, row sep=large]
 \& \mathrm{Biext}^1(P,Q;G) \arrow[ddl, "\varphi^n"'] \arrow[ddr, "(2.3.5)"] \& \\
 \& - \& \\
\mathrm{Hom}({}_nP\otimes{}_nQ,G) \arrow[rr, "(2.3.9)"] \& \& \mathrm{Hom}({}_nP,{}_n\underline{\mathrm{Ext}}^1(nQ,G))
\end{tikzcd}\quad ,
\end{equation}























}{F7I-VIII-020}où la flèche horizontale (2.3.9) est un isomorphisme si on a (2.3.2).

\underline{Proposition} 2.3.11. \underline{Le diagramme} (2.3.10) \underline{est} anticommutatif.

Dans le cas où $\underline{\mathrm{Hom}}(Q,G) = 0$, cet énoncé redonne donc une caractérisation de l'homomorphisme $\varphi^n$ de (2.2.2) en termes des homomorphismes (2.3.5) et (2.3.9), qui est celle utilisée précédemment en théorie des schémas abéliens (où d'ailleurs on a $nQ = Q$).

\SGAvar{\underline{Démonstration de} 2.5.11 (*). Soient $\xi$ un élément de $\mathrm{Biext}^1(P,Q;G)$, $\varphi_\xi : {}_nP \otimes {}_nQ \longrightarrow G$ l'accouplement associé, et $u_\xi : {}_nP \longrightarrow \underline{\mathrm{Ext}}^1(nQ,G)$ l'homomorphisme image de $\xi$ par (2.3.5). Il faut prouver que $-u_\xi$ se déduit de $\varphi_\xi$ par la flèche (2.3.9), i.e. que c'est le composé

}{\underline{Démonstration de} 2.3.11 (*). Soient $\xi$ un élément de $\mathrm{Biext}^1(P,Q;G)$, $\varphi_\xi : {}_nP \otimes {}_nQ \longrightarrow G$ l'accouplement associé, et $u_\xi : {}_nP \longrightarrow \underline{\mathrm{Ext}}^1(nQ,G)$ l'homomorphisme image de $\xi$ par (2.3.5). Il faut prouver que $-u_\xi$ se déduit de $\varphi_\xi$ par la flèche (2.3.9), i.e. que c'est le composé

}{F7I-VIII-021}\begin{equation}\tag{*}
{}_nP \xrightarrow{\;\varphi_\xi\;} \underline{\mathrm{Hom}}({}_nQ,G) \xrightarrow{\;(2.3.8)\;} {}_n\underline{\mathrm{Ext}}^1(nQ,G) \qquad .
\end{equation}
Pour prouver ceci, on peut supposer donnée une section $x$ de ${}_nP$ et prouver que son image $-u_\xi$ est égale à son image par le composé (*). (Ceci nous permettrait, si nous le désirions, de nous ramener au cas où $P = (\mathbb{Z}/n\mathbb{Z})_T$ , mais cette réduction ne nous sera pas nécessaire.)

a) Choisissons des résolutions plates $L.$ et $M.$ ,

\begin{equation*}
0 \longrightarrow L_1 \longrightarrow L_o \longrightarrow 0 \quad , \quad 0 \longrightarrow M_1 \longrightarrow M_o \longrightarrow 0
\end{equation*}
de $P$ et $Q$ respectivement, une résolution plate de $P$ (par exemple) étant 

(*)\quad Pour une démonstration plus élégante, cf 2.3.19 ci-dessous.


\SGApage{I}{244}

 obtenue en écrivant $P$ comme quotient d'un $\mathbb{Z}_T$-Module plat $L_o$, et en désignant par $L_1$ le noyau de $L_o \longrightarrow P$ (qui est plat, car sans torsion). On en déduit un complexe $L.\otimes M.$

\begin{equation*}
0 \longrightarrow L_1 \otimes M_1 \overset{d_2}{\longrightarrow} L_1 \otimes M_o + L_o \otimes M_1 \overset{d_1}{\longrightarrow} L_o \otimes M_o \longrightarrow 0 \qquad ,
\end{equation*}
dont le faisceau $\underline{H}_1$ nous servira de $\underline{\mathrm{Tor}}_1(P,Q)$. Désignons par $L'_o$ resp. $M'_o$ le sous-faisceau de $L_o$ resp. $M_o$ image inverse de ${}_nP$ resp. ${}_nQ$ , de sorte que le complexe $0 \longrightarrow L_1 \longrightarrow L'_o \longrightarrow 0$ (resp. le complexe $0 \longrightarrow M_1 \longrightarrow M'_o \longrightarrow 0$) est une résolution plate de ${}_nP$ (resp. de ${}_nQ$). On peut alors expliciter l'homomorphisme

\begin{equation*}
{}_nP \otimes {}_nQ \longrightarrow \underline{\mathrm{Tor}}_1(P,Q) = \underline{H}_1(L.\otimes M.)
\end{equation*}
par passage au quotient à partir de l'homomorphisme

\begin{equation*}
L'_o \otimes M'_o \longrightarrow \underline{H}_1(L.\otimes M.)
\end{equation*}
donné par

\begin{equation*}
x_o \otimes y_o \longmapsto \text{classe du 1-cycle } [nx_o] \otimes y_o - x_o \otimes [ny_o] \qquad ,
\end{equation*}
où les crochets $[\ ]$ indiquent que l'élément envisagé est considéré comme section du sous-faisceau $L_1$ de $L_o$ resp. $M_1$ de $M_o$ . L'explicitation précédente résulte en effet du calcul fait dans la démonstration de 2.1.10.

b)\quad Prenons une résolution injective

\begin{equation*}
0 \longrightarrow C^o \overset{d^o}{\longrightarrow} C^1 \longrightarrow \quad \ldots
\end{equation*}
de $G$, et identifions $\mathrm{Biext}^1(P,Q;G)$ à $\mathrm{Ext}^1(P \overset{\mathbb{L}}{\otimes} Q,G) = H^1(\mathrm{Hom}^{\cdot}(L.\otimes M.,C^{\cdot}))$. Donc l'élément $\xi$ peut s'interpréter comme un couple $f_1,f_o$ d'homomorphismes, donnant lieu à un diagramme anticommutatif 


\SGApage{I}{245}

\begin{equation}\tag{2.3.12}
\begin{tikzcd}[column sep=huge, row sep=huge]
L_1\otimes M_o + L_o\otimes M_1 \arrow[r, "d_1"] \arrow[d, "f_1"'] \arrow[rd, phantom, "{-}" description] \& L_o\otimes M_o \arrow[d, "f_o"'] \\
C^o \arrow[r, "d^o"] \& C^1
\end{tikzcd}\quad .
\end{equation}














Par définition, l'accouplement

\begin{equation*}
\varphi_{\xi} : {}_nP \otimes {}_nQ \longrightarrow G
\end{equation*}
se déduit alors par passage au quotient de l'homomorphisme

\begin{equation*}
\varphi : L'_o \otimes M'_o \longrightarrow G
\end{equation*}
donné par

\begin{equation}\tag{2.3.13}
\varphi(x_o \otimes y_o) = f_1([nx_o] \otimes y_o - x_o \otimes [ny_o]) \in G = \mathrm{Ker}(C^o \longrightarrow C^1) \ .
\end{equation}
Nous pouvons supposer que la section donnée $x$ de ${}_nP$ se relève en une section $x_o$ de $L'_o$ (car l'assertion (2.3.11) est locale sur $T$), et la formule précédente, où $y_o$ est considéré comme variable, définit donc un homomorphisme

\SGAvar{\begin{equation}\tag{2.3.14}
\varphi' : M'_o \longrightarrow G \quad , \qquad y_o \longmapsto f_1([nx_o] \otimes y_o \to x_o \otimes [ny_o]) \quad ,
\end{equation}
qui passe au quotient pour nous donner l'homomorphisme

}{\begin{equation}\tag{2.3.14}
\varphi' : M'_o \longrightarrow G \quad , \qquad y_o \longmapsto f_1([nx_o] \otimes y_o - x_o \otimes [ny_o]) \quad ,
\end{equation}
qui passe au quotient pour nous donner l'homomorphisme

}{F7I-VIII-022}\begin{equation*}
{}_nQ \longrightarrow G \quad , \qquad y \longmapsto \varphi_{\xi}(x,y)
\end{equation*}
associé à $\varphi_{\xi}$ et à $x$ .

c)\quad Il faut prendre l'image de cet homomorphisme par l'opérateur $\partial$ de (2.3.7), associé à la suite exacte (2.3.6). Or $nQ \simeq Q/{}_nQ$ est évidemment aussi isomorphe à $M_o/M'_o$ , de sorte qu'on a un homomorphisme de suites exactes 

\SGApage{I}{246}

\begin{center}
\begin{tikzcd}[column sep=large, row sep=large]
0 \arrow[r] \& M'_o \arrow[r] \arrow[d] \& M_o \arrow[r] \arrow[d] \& M_o/M'_o \arrow[r] \arrow[d, "\wr"'] \& 0 \\
0 \arrow[r] \& {}_{n}Q \arrow[r] \& Q \arrow[r] \& nQ \arrow[r] \& 0 \& {,}
\end{tikzcd}
\end{center}














\SGAvar{ où la dernière flèche verticale est un isomorphisme. Ceci montre que la section cherchée de $\underline{Ext}^1(nQ,G)$ peut aussi être calculéecomme l'image de $\varphi'$ (2.3.14) par l'opérateur cobord

\begin{equation*}
\underline{Hom}(M'_o,G) \overset{\partial}{\longrightarrow} \underline{Ext}^1(M_o/M'_o,G) \qquad ,
\end{equation*}
ou encore la section du deuxième membre défini par l'élément du $Ext^1$ global image de $\varphi'$ par l'opérateur cobord

}{ où la dernière flèche verticale est un isomorphisme. Ceci montre que la section cherchée de $\underline{Ext}^1(nQ,G)$ peut aussi être calculée comme l'image de $\varphi'$ (2.3.14) par l'opérateur cobord

\begin{equation*}
\underline{Hom}(M'_o,G) \overset{\partial}{\longrightarrow} \underline{Ext}^1(M_o/M'_o,G) \qquad ,
\end{equation*}
ou encore la section du deuxième membre définie par l'élément du $Ext^1$ global image de $\varphi'$ par l'opérateur cobord

}{F7I-VIII-023}\begin{equation*}
Hom(M'_o\ ,G) \longrightarrow Ext^1(M_o/M'_o,G) \qquad .
\end{equation*}
La construction habituelle de l'opérateur cobord nous donne alors la recette suivante : on prolonge le composé (2.3.13)

\begin{equation*}
M'_o \overset{\varphi'}{\longrightarrow} G \longrightarrow C^o
\end{equation*}
en un homomorphisme

\begin{equation*}
\psi \, : \, M_o \longrightarrow C^o \qquad ,
\end{equation*}
ce qui est possible puisque $C^o$ est injectif, et on considère l'unique homomorphisme

\begin{equation}\tag{2.3.15}
\lambda \, : \, M_o/M'_o \longrightarrow C^1
\end{equation}
rendant commutatif le diagramme

\begin{center}
\begin{tikzcd}[column sep=large, row sep=large]
0 \arrow[r] \& M'_o \arrow[r] \arrow[d, "\varphi'"] \& M_o \arrow[r] \arrow[d, "\psi"] \& M_o/M'_o \arrow[r] \arrow[d, "\lambda"] \& 0 \\
0 \arrow[r] \& G \arrow[r] \& C^o \arrow[r] \& C^1 \& \& {.}
\end{tikzcd}
\end{center}













\SGApage{I}{247}

Il est clair que (2.3.15) est à valeurs dans les 1-cocycles, i.e. est un 1-cocycle, et que ce dernier représente l'élément cherché de $Ext^1(nQ,G) = Ext^1(M_o/M'_o,G)$ .

d) Il reste à expliciter la section de $\underline{Ext}^1(nQ,G)$ déduite de $\xi$ et de $x$ grâce à la flèche (2.3.5) du diagramme (2.3.10). Pour ceci, nous pouvons ignorer le fait que la section $x$ de $P$ est annulée par $n$, et on trouve la description suivante, en termes de relèvement choisi $x_o$ de $x$ en une section de $L_o$ : cette dernière définit un homomorphisme d'augmentation

\SGAvar{\begin{equation*}
u \, : \, M. \longrightarrow L. \otimes M. \qquad , \qquad y_o \mapsto x_o \otimes y_o \ , \ y_1 \mapsto x_o \otimes y_1 \qquad ,
\end{equation*}
et la section de $Ext^1(Q,G)$ déduite de (2.3.1) est celle qui se déduit de l'élément $\xi$ de $Ext^1(L. \otimes M.,G)$ en composant avec l'homomorphisme précédent, donc est donné par l'homomorphisme de degré 1

\begin{center}
\begin{tikzcd}[column sep=large, row sep=huge]
0 \arrow[r] \& M_1 \arrow[r, "d_1"] \arrow[d, "{y_1 \mapsto f_1(x_o \otimes y_1)}"', "{-}"] \& M_o \arrow[r] \arrow[d, "{y_o \mapsto f_1(x_o \otimes y_o)}"] \& 0 \\
0 \arrow[r] \& C^o \arrow[r, "d^o"] \& C^1 \& \ldots \& {.}
\end{tikzcd}
\end{center}















}{\begin{equation*}
u \, : \, M. \longrightarrow L. \otimes M. \qquad , \qquad y_o \mapsto x_o \otimes y_o \ , \ y_1 \mapsto x_o \otimes y_1 \qquad ,
\end{equation*}
et la section de $Ext^1(Q,G)$ déduite de (2.3.1) est celle qui se déduit de l'élément $\xi$ de $Ext^1(L. \otimes M.,G)$ en composant avec l'homomorphisme précédent, donc est donnée par l'homomorphisme de degré 1

\begin{center}
\begin{tikzcd}[column sep=large, row sep=huge]
0 \arrow[r] \& M_1 \arrow[r, "d_1"] \arrow[d, "{y_1 \mapsto f_1(x_o \otimes y_1)}"', "{-}"] \& M_o \arrow[r] \arrow[d, "{y_o \mapsto f_o(x_o \otimes y_o)}"] \& 0 \\
0 \arrow[r] \& C^o \arrow[r, "d^o"] \& C^1 \& \ldots \& {.}
\end{tikzcd}
\end{center}















}{F7I-VIII-024}Pour prendre la restriction de cette classe à $nQ \hookrightarrow Q$, utilisons l'isomorphisme $nQ \simeq M_o/M'_o$ , qui fournit une résolution plate

\begin{equation*}
0 \longrightarrow M'_o \longrightarrow M_o \longrightarrow 0
\end{equation*}
de $nQ$, l'inclusion $nQ \hookrightarrow Q$ étant remontée en un homomorphisme de résolutions

\begin{center}
\begin{tikzcd}[column sep=large, row sep=huge]
0 \arrow[r] \& M'_o \arrow[r] \arrow[d, "{y'_o \mapsto [ny'_o]}"', "{+}"] \& M_o \arrow[r] \arrow[d, "{y_o \mapsto ny_o}"] \& 0 \\
0 \arrow[r] \& M_1 \arrow[r] \& M_o \arrow[r] \& 0 \& {.}
\end{tikzcd}
\end{center}














\SGApage{I}{248}

Par composition, on voit donc que la section cherchée de $\underline{\mathrm{Ext}}^1(nQ,G)$ est associée à l'élément de $\mathrm{Ext}^1(nQ,G)$ représenté par l'homomorphisme de complexes de degré 1

\begin{center}
\begin{tikzcd}[row sep=large, column sep=large]
0 \arrow[r, dash] \& M'_o \arrow[r, dash, "d_1"] \arrow[d, "g_1"'] \& M_o \arrow[r, dash] \arrow[d, "g_o"] \& 0 \& {} \\
0 \arrow[r, dash] \& C^o \arrow[r, dash] \& C^1 \arrow[r, dash] \& \dots \& {,}
\end{tikzcd}
\end{center}














avec

\begin{equation*}
g_1(y'_o) = f_1(x_o \otimes [ny'_o]) \quad , \quad g_o(y_o) = nf_o(x_o \otimes y_o) \quad .
\end{equation*}
e) \quad Utilisant maintenant le fait que $x$ est une section de ${}_nP$, montrons que l'élément précédent de $\mathrm{Ext}^1(nQ,G)$ est opposé de celui défini par (2.3.15). Pour ceci, notons que pour tout homomorphisme

\begin{equation*}
h : M_o \longrightarrow C^o \quad ,
\end{equation*}
l'élément du $\mathrm{Ext}^1$ envisagé dans d) ne change pas si on change $g_1,g_o$ respectivement en

\begin{equation*}
g'_1 = g_1 + hd_1 \quad , \quad g'_o = g_o - d^o h \quad .
\end{equation*}
Or on dispose d'un homomorphisme $\psi : M_o \longrightarrow C^o$ construit dans c), qui par définition satisfait la condition

\begin{equation*}
\psi(y'_o) = f_1([nx_o] \otimes y'_o - x_o \otimes [ny'_o])
\end{equation*}
pour $y'_o$ section de $M'_o$, ce qui s'écrit encore

\begin{equation*}
f_1(x_o \otimes [ny'_o]) \ (= g_1(y'_o)) = f_1([nx_o] \otimes y'_o \ ) - \psi(y'_o) \quad ;
\end{equation*}
définissons donc $h : M_o \longrightarrow C^o$ par 


\SGApage{I}{249}

\begin{equation*}
 h(y_o) = -f_1([nx_o] \otimes y_o) + \psi(y_o) \quad ;
\end{equation*}
donc par construction on aura

\begin{equation*}
g'_1 = 0 \ , \ g'_o(y_o) = nf_o(x_o \otimes y_o) + d^o f_1([nx_o] \otimes y_o) - d^o \psi(y_o) \quad ,
\end{equation*}
or on a $d^o f_1 = -f_o d_1$ donc $d^o f_1([nx_o] \otimes y_o) = - f_o(nx_o \otimes y_o)$, et il reste

\begin{equation*}
g'_1 = 0 \quad , \quad g'_o = -d^o \psi \quad ,
\end{equation*}
de sorte que l'élément envisagé est celui défini par l'homomorphisme $M_o/M'_o \simeq nQ \longrightarrow C^1$ déduit de $- d^o \psi$ par passage au quotient. On trouve donc l'opposé de l'élément défini par (2.3.15).

\SGAvar{\underline{Remarques} 2.3.16. \quad La démonstration précédente donne en fait une compatibilité légèrement plus précise que celle indiquée dans 2.3.11, obtenue en choisissant une section $x$ de ${}_nP$. Cette section définit alors un homomorphisme $y \longmapsto \varphi_\xi(x,y)$ de ${}_nQ$ dans $G$, d'où par l'homomorphisme cobord relatif à (2.3.6) un élément de $\mathrm{Ext}^1(nQ,G)$. D'autre part, oubliant que $x$ est annulé par $n$, on considère l'homomorphisme $\mathbb{Z}_T \longrightarrow P$ défini par $X$, d'où par image inverse une biextension de $(\mathbb{Z}_T,Q)$ par $G$, dont la classe peut s'interpréter comme un élément de $\mathrm{Ext}^1(Q,G)$, et induit donc un élément de $\mathrm{Ext}^1(nQ,G)$. Ceci posé, les deux éléments obtenus dans $\mathrm{Ext}^1(nQ,G)$ sont \underline{opposés}.

2.3.17. \quad Profitons de l'explicitation en termes de cocycles de l'accouplement $\varphi_\xi^n : {}_nP \otimes {}_nQ \longrightarrow G$, donnée dans (2.3.12), pour donner une autre expression de cet accouplement. Supposons d'abord que la biextension $E$ de classe $\xi$ admette une section $s$ sur $P \times Q$ . Prenons, comme
\SGApage{I}{250}
dans VII 3.8.13, comme résolution plate de $P$ celle donnée par $L_o = \mathbb{Z}[P]$, $L_1$ étant le noyau de l'homomorphisme canonique $\mathbb{Z}[P] \longrightarrow P$, et définissons de même $M.$ . Alors il existe une unique trivialisation de biextension de l'image inverse de $E$ sur $(L_o,M_o)$, dont la restriction à $P \times Q$ soit donnée par $s$, appelons la encore $s$ (cf. VII 3.8.2). En vertu de VII 3.8.12, la classe de $\xi$ dans $\mathrm{Ext}^1(P \stackrel{\mathbb{L}}{\otimes} Q, G)$ est donnée parle cocycle de composantes $(f,g)$, caractérisées par

}{\underline{Remarques} 2.3.16. \quad La démonstration précédente donne en fait une compatibilité légèrement plus précise que celle indiquée dans 2.3.11, obtenue en choisissant une section $x$ de ${}_nP$. Cette section définit alors un homomorphisme $y \longmapsto \varphi_\xi(x,y)$ de ${}_nQ$ dans $G$, d'où par l'homomorphisme cobord relatif à (2.3.6) un élément de $\mathrm{Ext}^1(nQ,G)$. D'autre part, oubliant que $x$ est annulé par $n$, on considère l'homomorphisme $\mathbb{Z}_T \longrightarrow P$ défini par $x$, d'où par image inverse une biextension de $(\mathbb{Z}_T,Q)$ par $G$, dont la classe peut s'interpréter comme un élément de $\mathrm{Ext}^1(Q,G)$, et induit donc un élément de $\mathrm{Ext}^1(nQ,G)$. Ceci posé, les deux éléments obtenus dans $\mathrm{Ext}^1(nQ,G)$ sont \underline{opposés}.

2.3.17. \quad Profitons de l'explicitation en termes de cocycles de l'accouplement $\varphi_\xi^n : {}_nP \otimes {}_nQ \longrightarrow G$, donnée dans (2.3.12), pour donner une autre expression de cet accouplement. Supposons d'abord que la biextension $E$ de classe $\xi$ admette une section $s$ sur $P \times Q$ . Prenons, comme
\SGApage{I}{250}
dans VII 3.8.13, comme résolution plate de $P$ celle donnée par $L_o = \mathbb{Z}[P]$, $L_1$ étant le noyau de l'homomorphisme canonique $\mathbb{Z}[P] \longrightarrow P$, et définissons de même $M.$ . Alors il existe une unique trivialisation de biextension de l'image inverse de $E$ sur $(L_o,M_o)$, dont la restriction à $P \times Q$ soit donnée par $s$, appelons-la encore $s$ (cf. VII 3.8.2). En vertu de VII 3.8.12, la classe de $\xi$ dans $\mathrm{Ext}^1(P \stackrel{\mathbb{L}}{\otimes} Q, G)$ est donnée par le cocycle de composantes $(f,g)$, caractérisées par

}{F7I-VIII-025}\begin{equation*}
s(x_1,y_o) = f(x_1 \otimes y_o) e_{E/Q}(y_o) \quad , \quad s(x_o,y_1) = g(x_o \otimes y_1) e_{E/P}(x_o) \quad .
\end{equation*}
Si alors $x$ , $y$ sont des sections de ${}_nP$ et ${}_nQ$ , on les remonte en les sections correspondantes $x_o = [x]$, $y_o = [y]$ de $L_o$ et $M_o$, et on applique la formule (2.3.13), qui nous donne, posant $z = s(x,y)$ :

\begin{equation}\tag{2.3.18}
\varphi_{\xi}^{n}(x,y) = {}^{n}z/e_{E/Q}(y) \quad - \quad z^{n}/e_{E/P}(x) \quad ,
\end{equation}
où l'exposant $n$ à gauche resp. à droite de $z$ indique la puissance $n$.ième pour la structure multiplicative gauche resp. droite sur $E$, de sorte qu'on trouve une section de $E$ au-dessus de $(0,y)$ resp. $(x,0)$ ; la signification du ``quotient'' d'un tel élément par $e_{E/Q}(y)$ resp. $e_{E/P}(x)$ est claire.

Notons maintenant que la formule précédente (2.3.18) garde un sens, dès que $z$ est une section globale de $E$ qui relève les sections $x,y$ de ${}_nP, {}_nQ$ . Je dis que la formule (2.3.18) est valable pour toute telle section, sans supposer qu'il existe une section de $E$ sur $P \times Q$. Comme la section $(x,y)$ de $P \times Q$ se relève en tous cas localement en une section de $E$ (puisque $E \longrightarrow P \times Q$ est un épimorphisme), ceci permet donc de calculer $\varphi_{\xi}^{n}(x,y)$ dans tous les cas. Noter d'ailleurs que si on
\SGApage{I}{251}
remplace $z$ par une autre section $z'$ de $E$ relevant $(x,y)$, on aura $z' = z + g$ , où $g$ est une section de $G$, et les deux termes du deuxième membre de (2.3.18) correspondants à $z'$ sont ceux correspondants à $z$, augmentés de $ng$, de sorte que la différence est bien indépendante de $z$ , ce qui nous montre que la formule (2.3.18) définit bien un homomorphisme bien déterminé ${}_nP \otimes {}_nQ \longrightarrow G$.

\SGAvar{Reste à montrer que ce dernier est bien $\varphi^n$. Or pour ceci on est ramené aussitôt au cas où $P=Q = (\mathbb{Z}/n\mathbb{Z})_T$, de sorte que $P \times Q$ est de la forme $I_T$ , où $I = \mathbb{Z}/n\mathbb{Z} \times \mathbb{Z}/n\mathbb{Z}$ est un ensemble fini. Dans ce cas, le torseur $E$ sous $G_{P \times Q}$ (qui s'identifie à une famille de torseurs $E_i$ sous $G$, $i \in I$) devient évidemment trivial localement sur $T$, de sorte que nous sommes ramenés au cas où il existe une section globale de $E$ sur $P \times Q$, cas déjà traité.

}{Reste à montrer que ce dernier est bien $\varphi_\xi^n$. Or pour ceci on est ramené aussitôt au cas où $P=Q = (\mathbb{Z}/n\mathbb{Z})_T$, de sorte que $P \times Q$ est de la forme $I_T$ , où $I = \mathbb{Z}/n\mathbb{Z} \times \mathbb{Z}/n\mathbb{Z}$ est un ensemble fini. Dans ce cas, le torseur $E$ sous $G_{P \times Q}$ (qui s'identifie à une famille de torseurs $E_i$ sous $G$, $i \in I$) devient évidemment trivial localement sur $T$, de sorte que nous sommes ramenés au cas où il existe une section globale de $E$ sur $P \times Q$, cas déjà traité.

}{F7I-VIII-026}2.3.19.  \underline{Autre démonstration de} 2.3.11 (de DELIGNE).

a)  Réduction à $P = \mathbb{Z}/n\mathbb{Z}$ , et à $T = (\mathrm{Ens})$

b)  Soit $\partial$ la flèche suivante de $nQ$ dans ${}_nQ[1]$ :

\begin{equation*}
nQ \quad \xleftarrow[\mbox{``n''}]{\sim} \quad [\,{}_nQ \longrightarrow Q\,] \longrightarrow {}_nQ[1] \quad .
\end{equation*}
Le diagramme suivant est anticommutatif :

\SGAvar{\begin{center}
\begin{tikzcd}[row sep=huge, column sep=huge]
{nQ[o]} \arrow[r, "1"] \arrow[d, "\partial"'] \& {[Q \stackrel{n}{\longrightarrow} Q]} \\
{{}_nQ[1]} \arrow[ru, "1"'] \& {.}
\end{tikzcd}
\end{center}

















}{\begin{center}
\begin{tikzcd}[row sep=huge, column sep=huge]
{nQ[0]} \arrow[r, "1"] \arrow[d, "\partial"'] \& {[Q \stackrel{n}{\longrightarrow} Q]} \\
{{}_nQ[1]} \arrow[ru, "1"'] \& {.}
\end{tikzcd}
\end{center}

















}{F7I-VIII-027}En effet, les flèches
\SGApage{I}{252}

\begin{center}
\begin{tikzcd}[column sep=small, row sep=large]
{[{}_nQ} \arrow[r] \arrow[dd, "0"'] \& {Q]} \arrow[dd, "n"] \& \& {[{}_nQ} \arrow[r] \arrow[dd, "1"'] \& {Q]} \arrow[dd, "0"'] \& \text{sont opposées} \& {[{}_nQ} \arrow[r] \& {Q]} \arrow[ddl, "1"] \\
 \& \& \text{et} \& \& \& \text{grâce à} \quad H \, : \& \& \\
{[Q} \arrow[r, "n"] \& {Q]} \& \& {[Q} \arrow[r, "n"] \& {Q]} \& \& {[Q} \arrow[r, "n"] \& {Q]}
\end{tikzcd}
\end{center}




















c)   2.3.10 est le diagramme

\begin{center}
\begin{tikzcd}[column sep=small, row sep=small]
 \& \& \& H^1\mathrm{RHom}(\mathbb{Z} \overset{\mathbb{L}}{\otimes} Q,G) \arrow[ddd] \\
 \& \& \& + \\
H^1\mathrm{RHom}(\mathbb{Z}/n\mathbb{Z} \overset{\mathbb{L}}{\otimes} Q,G) \arrow[rr] \arrow[dd] \arrow[uurrr] \& \& H^1\mathrm{RHom}(\mathbb{Z}/n\mathbb{Z} , \mathrm{RHom}(Q,G)) \arrow[dd] \& \\
 \& \& \& \mathrm{Ext}^1(Q,G) \arrow[ddl] \\
\mathrm{Hom}(\mathrm{Tor}_1(\mathbb{Z}/n\mathbb{Z} ,Q),G) \arrow[d] \& |[draw, circle, inner sep=1pt]| {-\ ?} \& \mathrm{Hom}(\mathbb{Z}/n\mathbb{Z} ,\mathrm{Ext}^1(Q,G)) \arrow[d] \& \\
\mathrm{Hom}({}_nQ,G) \arrow[rr] \& \& \mathrm{Ext}^1(nQ,G) \&
\end{tikzcd}
\end{center}















et

c 1)   $\mathbb{Z}/n\mathbb{Z} \overset{\mathbb{L}}{\otimes} Q \simeq [Q \xrightarrow{\,n\,} Q]$, et, modulo cet isomorphisme, la 1ère flèche composée verticale est induite par

\begin{equation*}
{}_nQ[1] \longrightarrow [Q \xrightarrow{\,n\,} Q] \quad ;
\end{equation*}
c 2)   la 2\textsuperscript{ème} flèche horizontale est transposée de la flèche $\partial$ indiquée plus haut dans b);

c 3)   si on parcourt le diagramme par l' $\curvearrowright$ , il est clair qu'on tombe sur la flèche induite par $nQ[0] \longrightarrow [Q \xrightarrow{\,n\,} Q]$. D'où la conclusion 2.3.11.

\SGApage{I}{253}

2.4.   Soit

\begin{equation*}
f : \underline{T}' \longrightarrow \underline{T}
\end{equation*}
un morphisme de topos. Alors on a un homomorphisme canonique (VII 2.8)

\begin{equation*}
\mathrm{Biext}^1(P,Q;G) \longrightarrow \mathrm{Biext}^1(P',Q';G') \quad ,
\end{equation*}
où les ' dénotent les Groupes images inverses. Comme, pour tout Groupe commutatif $E$ de $\underline{T}$, on a évidemment un isomorphisme canonique

\begin{equation*}
({}_nE)' \;\simeq\; {}_n(E') \quad ,
\end{equation*}
les homomorphismes (2.2.2) sur $\underline{T}$ et sur $\underline{T}'$ permettent de définir un diagramme canonique

\begin{equation}\tag{2.4.1}
\begin{tikzcd}[column sep=huge, row sep=huge]
\mathrm{Biext}^1(P,Q;G) \arrow[r] \arrow[d, "{\varphi^n_{P,Q}}"'] \& \mathrm{Biext}^1(P',Q';G') \arrow[d, "{\varphi^n_{P',Q'}}"] \& \\
\mathrm{Hom}({}_nP \otimes {}_nQ,G) \arrow[r] \& \mathrm{Hom}({}_nP' \otimes {}_nQ',G') \& \quad .
\end{tikzcd}
\end{equation}











De même, les homomorphismes (2.2.3) donnent naissance à un diagramme analogue en termes des $T_\ell$ . Ceci posé, je dis que les diagrammes envisagés sont \underline{commutatifs}. La vérification est essentiellement triviale en termes des définitions, et laissée au lecteur. On notera d'ailleurs,plus généralement, que \underline{tous} les homomorphismes canoniques envisagés dans le présent numéro et le précédent sont \underline{compatibles} \underline{au changement de topos}, dans un sens évident. Nous utiliserons par la suite sans autre mention les compatibilités de cette nature.


\SGApage{I}{254}


\SGAsectiontoc{3. Biextensions de schémas en groupes $(P,Q)$ par $\protect\underline{G}_m$ : généralités}{3}
3. \underline{Biextensions de schémas en groupes} $(P,Q)$ \underline{par} $\underline{G}_m$ : \underline{généralités}

Dans le présent numéro, nous fixons un schéma de base $S$, et, comme annoncé au n° 0, nous travaillons dans le topos fppf de $S$.

3.1. Pour mémoire, rappelons d'abord l'exemple le plus important d'un groupe $\mathrm{Biext}^1(P,Q;\underline{G}_m)$ : celui où $P$ et $Q$ sont deux schémas abéliens sur $S$, et où le groupe $\mathrm{Biext}^1$ s'interprète comme le groupe des correspondances divisorielles sur $(P,Q)$ (VII 2.9.4). Les accouplements $\ell$-adiques du numéro précédent sont alors les accouplements bien connus [1] dans la théorie des schémas abéliens.

3.2. Plus généralement, supposons seulement $Q$ un schéma abélien, alors il est bien connu (VII 1.3.8) qu'on a

\begin{equation}\tag{3.2.1}
\underline{\mathrm{Hom}}(Q,\underline{G}_m) = 0 \qquad ,
\end{equation}
donc (1.4) on a un isomorphisme canonique

\begin{equation}\tag{3.2.2}
\mathrm{Biext}^1(P,Q;\underline{G}_m) \;\simeq\; \mathrm{Hom}(P,Q^*) \qquad ,
\end{equation}
où

\begin{equation}\tag{3.2.3}
Q^* = \underline{\mathrm{Ext}}^1(Q,\underline{G}_m)
\end{equation}
est le schéma abélien dual de $Q$. (Pour la représentabilité de $Q^*$, due à M. RAYNAUD dans le cas général, cf. [1].)

Lorsque $P$ est également un schéma abélien, on a, de façon symétrique à (3.2.2), un isomorphisme canonique

\begin{equation}\tag{3.2.4}
\mathrm{Biext}^1(P,Q;\underline{G}_m) \;\simeq\; \mathrm{Hom}(Q,P^*) \qquad .
\end{equation}

\SGApage{I}{255}

Ainsi la donnée d'une biextension $\xi$ de $(P,Q)$ par $\underline{G}_m$ peut s'interpréter, indifféremment, par la donnée d'un homomorphisme

\begin{equation*}
u : P \longrightarrow Q^*
\end{equation*}
ou d'un homomorphisme

\begin{equation*}
v : Q \longrightarrow P^* \qquad .
\end{equation*}
D'ailleurs, on sait qu'un homomorphisme de schémas abéliens $w : A \longrightarrow B$ est un isomorphisme si et seulement si pour tout entier $n > 0$, l'homomorphisme induit ${}_n w : {}_n A \longrightarrow {}_n B$ est un isomorphisme. Or dans le cas de $u$ et de $v$ ci-dessus, ${}_n u : {}_n P \longrightarrow {}_n Q^*$ s'identifie à l'homomorphisme ${}_n P \longrightarrow D({}_n Q)$ défini par l'accouplement $\varphi_\xi^n : {}_n P \times {}_n Q \longrightarrow \underline{G}_m$ (cf. 2.3.11 ), tandis que ${}_n v : {}_n Q \longrightarrow {}_n P^*$ s'identifie à l'homomorphisme ${}_n Q \longrightarrow D({}_n P)$ défini par l'accouplement $\psi_\xi^n : {}_n P \times {}_n Q \longrightarrow \underline{G}_m$ . Or on a $\varphi_\xi^n = - \psi_\xi^n$ (2.1.10), donc au signe ${}_n u$ et ${}_n v$ sont transposés l'un de l'autre. Tenant compte du fait bien connu que les ${}_n P$ et les ${}_n Q$ sont finis et localement libres sur $S$, donc satisfont à la dualité de Cartier, on trouve donc que ${}_n u$ est un isomorphisme si et seulement si ${}_n v$ est un isomorphisme, donc que $u$ est un isomorphisme si et seulement si $v$ est un isomorphisme. Une façon équivalente (un peu moins symétrique) d'énoncer ce résultat est la suivante : considérons la biextension canonique (biextension de Weil) de $(P,P^*)$ par $\underline{G}_m$ (correspondante à $v : P^* \longrightarrow P^*$ l'identité), alors l'homomorphisme canonique correspondant $u : P \longrightarrow P^{**}$ est un \underline{isomorphisme}. C'est le \underline{théorème de bidualité} de BARSOTTI et CARTIER.


\SGApage{I}{256}

3.3. Nous passons au cas où $Q$ est fini localement libre sur $S$, ou est un schéma de type multiplicatif sur $S$. Le lemme-clef dans ce cas est le résultat (bien connu) suivant :

\underline{Proposition} 3.3.1. \underline{Soit} $Q$ \underline{un schéma en groupes commutatifs sur} $S$, \underline{et supposons que} $Q$ \underline{soit fini localement libre sur} $S$, \underline{ou que} $Q$ \underline{soit de type fini et de type multiplicatif sur} $S$ (SGA 3 IX). \underline{Alors on a}

\begin{equation*}
\underline{\mathrm{Ext}}^1(Q,\underline{G}_m) = 0 \qquad .
\end{equation*}
\underline{Démonstration}. Distinguons les deux cas envisagés.

a) $Q$ est fini localement libre sur $S$. Il faut prouver que toute extension $E$ de $Q$ par $\underline{G}_m$ splitte localement pour la topologie fppf. On peut supposer évidemment $S$ affine ; il existe alors un entier $n > 0$ tel qu'on ait $nQ = 0$ (SGA 3 IX 6.1). D'autre part, on sait que $n.\mathrm{id}$ est un épimorphisme dans $\underline{G}_m$ pour fppf, donc la suite exacte ``de Kummer''

\begin{equation*}
0 \longrightarrow \mu_n \longrightarrow \underline{G}_m \xrightarrow{\ n\ } \underline{G}_m \longrightarrow 0
\end{equation*}
donne naissance à une suite exacte des $\mathrm{Ext}^i(Q, - )$, qui nous donne un \underline{épimorphisme}

\begin{equation*}
\mathrm{Ext}^1(Q,\mu_n) \longrightarrow {}_n\mathrm{Ext}^1(Q,\underline{G}_m) = \mathrm{Ext}^1(Q,\underline{G}_m) \qquad ,
\end{equation*}
la dernière égalité provenant du fait que la multiplication par $n$ étant nulle dans $Q$, est également nulle dans $\mathrm{Ext}^1(Q,\underline{G}_m)$. Cela montre que l'extension $E$ provient d'une extension $F$ de $Q$ par $\mu_n$ ; $F$ est représen
\SGApage{I}{257}
table par un schéma fini localement libre sur $S$ (théorie de la descente fpqc, SGA 1 VII). A cause des propriétés connues de la dualité de Cartier des groupes finis localement libres (SGA 3 VII 3), le dual $D(F) = \underline{\mathrm{Hom}}(F,\underline{G}_m)$ est une extension de $D(\mu_n)$ par $D(Q)$. (Il faut voir que si $A \longrightarrow B$ est un monomorphisme de groupes commutatifs finis localement libres sur $S$, alors $D(B) \longrightarrow D(A)$ est un morphisme fidèlement plat. Comme $D(B)$ et $D(A)$ sont également finis localement libres sur $S$, le critère de fidèle platitude par fibres s'applique et nous ramène au cas où $S$ est le spectre d'un corps $k$. Dans ce cas, notre assertion signifie simplement que l'homomorphisme sur les anneaux affines $k(D(A)) \longrightarrow k(D(B))$ est injectif (SGA 3 $\mathrm{VI}_B$ 11.14), or cet homomorphisme est transposé de l'homomorphisme $k(B) \longrightarrow k(A)$, qui est surjectif puisque $A \longrightarrow B$ est une immersion fermée (SGA 3 $\mathrm{VI}_B$ 1.4.2). Donc $k(D(A)) \longrightarrow k(D(B))$ est bien injectif, cqfd.) Or l'inclusion $\mu_n \to \underline{G}_m$ s'identifie à une section de $D(\mu_n) = \underline{\mathrm{Hom}}(\mu_n,\underline{G}_m)$, et tout revient à prouver que cette inclusion se prolonge, localement fppf sur $S$, en un homomorphisme $F \longrightarrow \underline{G}_m$, i.e. que la section envisagée de $D(\mu_n)$ se relève localement fppf en une section de $D(F)$. Or ceci est clair, puisque $D(F) \longrightarrow D(\mu_n)$ est fidèlement plat et de présentation finie.

\SGAvar{b) En vertu de (SGA 3 X 4.5), quitte à se localiser sur $S$ pour la topologie étale, on peut supposer que $Q$ est diagonalisable, i.e. de la forme $D(M)$, où $M$ est un $\mathbb{Z}$-module de type fini. Soit $T$ le sous-module de torsion de $M$, et $L = M/T$ le quotient, qui est isomorphisme à $\mathbb{Z}^r$ .


}{b) En vertu de (SGA 3 X 4.5), quitte à se localiser sur $S$ pour la topologie étale, on peut supposer que $Q$ est diagonalisable, i.e. de la forme $D(M)$, où $M$ est un $\mathbb{Z}$-module de type fini. Soit $T$ le sous-module de torsion de $M$, et $L = M/T$ le quotient, qui est isomorphe à $\mathbb{Z}^r$ .


}{F7I-VIII-028}\SGApage{I}{258}

\SGAvar{Alors la suite exacte $0 \longrightarrow T \longrightarrow M \longrightarrow L \longrightarrow 0$ donne une suite exacte $0 \longrightarrow D(L) \longrightarrow D(M) \longrightarrow D(L) \longrightarrow 0$ (SGA 3 IX 3.1), donc on est ramené à prouver séparément que le $\underline{\mathrm{Ext}}^1$ de $D(L)$ et de $D(T)$ avec $\underline{G}_m$ est nul. Le cas de $D(T)$ est justiciable de a), le cas de $D(L)$ se ramène à celui où $L = \mathbb{Z}$ , de sorte qu'on est réduit à vérifier que l'on a

}{Alors la suite exacte $0 \longrightarrow T \longrightarrow M \longrightarrow L \longrightarrow 0$ donne une suite exacte $0 \longrightarrow D(L) \longrightarrow D(M) \longrightarrow D(T) \longrightarrow 0$ (SGA 3 IX 3.1), donc on est ramené à prouver séparément que le $\underline{\mathrm{Ext}}^1$ de $D(L)$ et de $D(T)$ avec $\underline{G}_m$ est nul. Le cas de $D(T)$ est justiciable de a), le cas de $D(L)$ se ramène à celui où $L = \mathbb{Z}$ , de sorte qu'on est réduit à vérifier que l'on a

}{F7I-VIII-029}\begin{equation*}
\underline{\mathrm{Ext}}^1(\underline{G}_m,\underline{G}_m) = 0 \quad ,
\end{equation*}
i.e. que toute extension $E$ de $\underline{G}_m$ par $\underline{G}_m$ est localement triviale fppf. Or une telle extension est représentable par un schéma en groupes affine et plat sur $S$ (théorie de la descente fpqc SGA 1 VII), dont les fibres sont des tores, c'est donc un tore (SGA 3 X 4.9 et \qquad ). En vertu de SGA 3 X 4.5, on peut supposer que $E$ est un tore trivial $D(R) \simeq \underline{G}_m^{\,2}$. De plus, quitte à se localiser au sens de Zariski, on peut supposer que le monomorphisme $\underline{G}_m \hookrightarrow E = D(R)$ provient d'un morphisme constant $R \longrightarrow \mathbb{Z}_S$ (SGA 3 VIII 1.4), et celui-ci est nécessairement un épimorphisme (SGA 3 VIII 3.2 b)). Choisissant un $r \in R$ au-dessus de l'élément $1$ de $\mathbb{Z}$ , on trouve donc un homomorphisme $E \longrightarrow \underline{G}_m$ qui scinde l'extension $E$, ce qui prouve l'assertion voulue.

\underline{Corollaire} 3.3.2. \underline{Soit} $Q$ \underline{comme dans} 3.3.1. \underline{Alors on a}, \underline{pour tout} $P$, \underline{un isomorphisme canonique}

\begin{equation}\tag{3.3.2.1}
\mathrm{Biext}^1(P,Q;\underline{G}_m) \;\xleftarrow{\ \sim\ }\; \mathrm{Ext}^1(P,D(Q)) \quad ,
\end{equation}
\underline{où} $D(Q)$ \underline{désigne le dual de Cartier} $\underline{\mathrm{Hom}}(Q,\underline{G}_m)$ \underline{de} $Q$. \underline{Plus précisément}, \underline{le foncteur canonique} (1.1.6) \underline{est une équivalence de la catégorie des extensions de} $P$ \underline{par} $D(Q)$ \underline{avec la catégorie des biextensions de} $(P,Q)$ \underline{par} $\underline{G}_m$ .


\SGApage{I}{259}

C'est en effet un cas particulier de 1.5, compte tenu de 3.3.1.

\underline{Remarques} 3.3.3. Sous les conditions de 3.3.1 , $D(Q)$ est représentable par un schéma en groupes séparé, localement de présentation finie et localement quasi-fini sur $S$. Plus précisément, si $Q$ est fini localement libre sur $S$, il en est de même de $D(Q)$ comme il est bien connu (SGA 3 VII 3), et si $Q$ est de type multiplicatif, alors $D(Q)$ est un "\underline{groupe constant tordu à engendrement fini}" (SGA 3 X 5.1), comme il a été vu dans (SGA 3 IX 4.5 et X 5.7). Si $P$ est un schéma en groupes, toute extension de $P$ par $D(Q)$ est donc représentable, d'après un lemme connu de la théorie de la descente fpqc (SGA 3 X 5.4).

3.3.4 \qquad Il est faux \quad que pour tout schéma en groupes $Q$ affine, plat et de présentation finie sur $S$, on ait $\underline{\mathrm{Ext}}^1(Q,\underline{G}_m) = 0$, déjà pour $Q = \underline{G}_{a_S}$ , $S = \mathrm{Spec}\ k[t]/(t^2)$, $k$ étant un corps de car. $p > 0$ .

\underline{Proposition} 3.4. \underline{Soient} $P$ \underline{un schéma en groupes lisse de présentation finie sur} $S$ \underline{et à fibres connexes}, $Q$ \underline{un tore sur} $S$. \underline{Alors la catégorie des biextensions de} $(P,Q)$ \underline{par} $\underline{G}_m$ \underline{est équivalente à la catégorie ponctuelle}, \underline{i.e.} \underline{on a}

\begin{equation}\tag{3.4.1}
\mathrm{Biext}^{o}(P,Q;\underline{G}_m) = \mathrm{Biext}^1(P,Q;\ \underline{G}_m) = 0
\end{equation}
Soit en effet $M = D(Q)$, qui est un groupe constant tordu sur $S$ à fibres isomorphes à des $\mathbb{Z}^{\,r}$ (3.3.3). En vertu de 3.3.2, la catégorie envisagée est équivalente à celle des extensions de $P$ par $M$. Cette catégorie est \underline{rigide}, i.e. le groupe des automorphismes d'un objet
\SGApage{I}{260}
est réduit au groupe unité, ce qui exprime la première relation (3.4.1), qui s'écrit également $\mathrm{Hom}(P,M) = 0$. Or cette relation résulte aussitôt du fait que $M$ est étale sur $S$ et $P$ à fibres connexes. Il reste donc à prouver la relation

\SGAvar{\begin{equation*}
\mathrm{Ext}^1(P,M) = 0 \quad ,
\end{equation*}
i.e. que toute extension $E$ de $P$ par $M$ est triviale. Utilisant l'unicité d'un scindage de cette extension (s'il en existe), on est ramené à prouver que $E$ est localement triviale, ce qui nous ramène au cas où $M = \mathbb{Z}^{r}_{S}$ , et même au cas où $M = \mathbb{Z}_{S}$ . Lorsque $S$ est le spectre d'un corps, la conclusion résulte de 5.5 (i) ci-dessous. (NB nous n'utiliserons pas 3.4 avant IX 6, et en particulier il n'y aura pas de cercle vicieux !). Montrons maintenant comment réduire le cas général à ce cas particulier. Notons qu'il suffit de trouver une section $f$ de $P$ sur $P$ telle que $f(0) = 0$, car une telle section sera nécessairement multiplicative, i.e. satisfera $f(x+y)-f(x)-f(y) = 0$ : en effet, le premier membre de cette relation désigne un morphisme $P \times_{S} P \longrightarrow \mathbb{Z}_{S}$ qui est nul sur la section nulle de $P \times_{S} P$ , donc est nulle puisque $P \times_{S} P$ est à fibres connexes ($P$ étant à fibres \underline{géométriques} connexes). La conclusion voulue résulte alors du

}{\begin{equation*}
\mathrm{Ext}^1(P,M) = 0 \quad ,
\end{equation*}
i.e. que toute extension $E$ de $P$ par $M$ est triviale. Utilisant l'unicité d'un scindage de cette extension (s'il en existe), on est ramené à prouver que $E$ est localement triviale, ce qui nous ramène au cas où $M = \mathbb{Z}^{r}_{S}$ , et même au cas où $M = \mathbb{Z}_{S}$ . Lorsque $S$ est le spectre d'un corps, la conclusion résulte de 5.5 (i) ci-dessous. (NB nous n'utiliserons pas 3.4 avant IX 6, et en particulier il n'y aura pas de cercle vicieux !). Montrons maintenant comment réduire le cas général à ce cas particulier. Notons qu'il suffit de trouver une section $f$ de $E$ sur $P$ telle que $f(0) = 0$, car une telle section sera nécessairement multiplicative, i.e. satisfera $f(x+y)-f(x)-f(y) = 0$ : en effet, le premier membre de cette relation désigne un morphisme $P \times_{S} P \longrightarrow \mathbb{Z}_{S}$ qui est nul sur la section nulle de $P \times_{S} P$ , donc est nul puisque $P \times_{S} P$ est à fibres connexes ($P$ étant à fibres \underline{géométriques} connexes). La conclusion voulue résulte alors du

}{F7I-VIII-030}\underline{Lemme} 3.4.2. \underline{Soient} $S$ \underline{un schéma,} $P$ \underline{un} \underline{$S$-schéma lisse de présentation finie,} \underline{à fibres géométriques connexes,} $Z$ \underline{un groupe,} $E$ \underline{un} $Z_{P}$\underline{-torseur} (\underline{pour la topologie} fppf, \underline{la topologie étale ou la topologie} fpqc, \underline{cela revient au même,} cf. SGA 3 X 5.4 ...), $e$ \underline{une section de} $P$ \underline{sur} $S$, $e'$ \underline{une section de} $E$ \underline{sur} $S$ \underline{au-dessus de} $e$ (\underline{i}.\underline{e}. \underline{une section de} $e^{-1}(E)$ \underline{sur} $S$).


\SGApage{I}{261}

\underline{Pour que} $E$ \underline{soit trivial,} \underline{il} (\underline{faut et il}) \underline{suffit que pour tout} $s \in S$, \underline{le torseur induit} $E_{s}$ \underline{de} $P_{s}$ \underline{le soit}, \underline{et alors il existe une unique section} $f$ \underline{de} $E$ \underline{sur} $P$ \underline{telle que} $f \circ e = e'$ .

\SGAvar{L'unicité d'une telle section est claire, car deux telles sections $f$ , $g$ ``différent'' par un morphisme $h : P \longrightarrow Z_{S}$ tel que $h \circ e = 1$, et un tel morphisme est constant de valeur $1$ puisque $P$ est à fibres connexes. D'ailleurs, s'il existe une section $f$ de $E$ sur $P$, on peut toujours la corriger par un morphisme $P \longrightarrow Z_{S}$ (provenant d'un morphisme $z : S \longrightarrow Z_{S}$) pour obtenir une section $f'$ telle que $f' \circ e = e'$ : il suffit de choisir $z$ tel que $f(e(t)) = e'(t)z(t)$. Appliquant cette dernière observation aux fibres au-dessus des $s \in S$, on trouve pour tout point $s \in S$ une section bien déterminée $f_{s}$ de $E_{s}$ sur $P_{s}$, transformant $e_{s}$ en $e'_{s}$ . Tout revient à prouver que ces $f_{s}$ proviennent d'une section $f$ de $E$ sur $P$, ou ce qui revient au même (EGA IV 17.9.3), que l'ensemble $U \subset E$ réunion des $f_{s}(P_{s})$ est une partie ouverte de $E$. Supposons d'abord $S$ noethérien ; alors on vérifie facilement que $U$ est constructible, en utilisant le fait que pour tout $s \in S$, il existe un voisinage ouvert $T'$ de $s$ dans $T = {}$ adhérence de $s$ muni de la structure induite réduite, et une section de $E_{T'}$ sur $P_{T'}$ prolongeant $f_{s}$ et compatible avec les sections marquées $e$ et $e'$ (comme il résulte trivialement de EGA IV 8.8.2). En vertu du critère EGA IV 1.10.1, il reste alors à prouver que $U$ est stable par générisation. Pour ceci, on est ramené, par la méthode habituelle utilisant EGA II 7.1.9, au cas où $S$ est le spectre d'un anneau de valuation discrète. Comme $P$ est lisse sur $S$,
\SGApage{I}{262}
$P$ est donc régulier donc normal, et comme il est à fibres connexes, il est connexe. Mais alors la section donnée de $E$ au-dessus de la fibre générique de $P$ sur $S$ se prolonge de façon unique en une section $f$ de $E$ sur $P$ (EGA IV 18.10.19 et 18.10.20), évidemment compatible avec $e$ et $e'$. Donc dans ce cas $U = f(P)$ est ouvert. Cela prouve 3.4 dans le cas où $S$ est noethérien. Le cas général se ramène aisément à ce cas, par les procédés habituels de passage à la limite, utilisant EGA IV 8,9, dont nous laissons le détail au lecteur.

}{L'unicité d'une telle section est claire, car deux telles sections $f$ , $g$ ``diffèrent'' par un morphisme $h : P \longrightarrow Z_{S}$ tel que $h \circ e = 1$, et un tel morphisme est constant de valeur $1$ puisque $P$ est à fibres connexes. D'ailleurs, s'il existe une section $f$ de $E$ sur $P$, on peut toujours la corriger par un morphisme $P \longrightarrow Z_{S}$ (provenant d'un morphisme $z : S \longrightarrow Z_{S}$) pour obtenir une section $f'$ telle que $f' \circ e = e'$ : il suffit de choisir $z$ tel que $f(e(t)) = e'(t)z(t)$. Appliquant cette dernière observation aux fibres au-dessus des $s \in S$, on trouve pour tout point $s \in S$ une section bien déterminée $f_{s}$ de $E_{s}$ sur $P_{s}$, transformant $e_{s}$ en $e'_{s}$ . Tout revient à prouver que ces $f_{s}$ proviennent d'une section $f$ de $E$ sur $P$, ou ce qui revient au même (EGA IV 17.9.3), que l'ensemble $U \subset E$ réunion des $f_{s}(P_{s})$ est une partie ouverte de $E$. Supposons d'abord $S$ noethérien ; alors on vérifie facilement que $U$ est constructible, en utilisant le fait que pour tout $s \in S$, il existe un voisinage ouvert $T'$ de $s$ dans $T = {}$ adhérence de $s$ muni de la structure induite réduite, et une section de $E_{T'}$ sur $P_{T'}$ prolongeant $f_{s}$ et compatible avec les sections marquées $e$ et $e'$ (comme il résulte trivialement de EGA IV 8.8.2). En vertu du critère EGA IV 1.10.1, il reste alors à prouver que $U$ est stable par générisation. Pour ceci, on est ramené, par la méthode habituelle utilisant EGA II 7.1.9, au cas où $S$ est le spectre d'un anneau de valuation discrète. Comme $P$ est lisse sur $S$,
\SGApage{I}{262}
$P$ est donc régulier donc normal, et comme il est à fibres connexes, il est connexe. Mais alors la section donnée de $E$ au-dessus de la fibre générique de $P$ sur $S$ se prolonge de façon unique en une section $f$ de $E$ sur $P$ (EGA IV 18.10.19 et 18.10.20), évidemment compatible avec $e$ et $e'$. Donc dans ce cas $U = f(P)$ est ouvert. Cela prouve 3.4 dans le cas où $S$ est noethérien. Le cas général se ramène aisément à ce cas, par les procédés habituels de passage à la limite, utilisant EGA IV 8,9, dont nous laissons le détail au lecteur.

}{F7I-VIII-031}\underline{Remarque} 3.4.3. On peut"écraser" le lemme 3.4.2 par des références savantes, savoir (SGA 4 XV 1.15, cas $n=1$, et 4.1).

\underline{Corollaire} 3.5. \underline{Soient} $P$ \underline{un schéma en groupes lisse de présentation finie sur} $S$, \underline{à fibres connexes}, $0 \longrightarrow Q' \longrightarrow Q \longrightarrow Q'' \longrightarrow 0$ \underline{une suite exacte de Groupes commutatifs sur} $S$, \underline{avec} $Q'$ \underline{un tore}. \underline{Alors le foncteur} "\underline{image inverse de biextensions}"

\begin{equation}\tag{3.5.1}
\mathrm{BIEXT}(P,Q''; \underline{G}_m) \stackrel{\approx}{\longrightarrow} \mathrm{BIEXT}(P,Q;\underline{G}_m)
\end{equation}
\underline{est une équivalence de catégories}, \underline{en particulier les homomorphismes naturels suivants sont des isomorphismes}

\begin{equation}\tag{3.5.2}
\mathrm{Biext}^i(P,Q''; \underline{G}_m) \stackrel{\sim}{\longrightarrow} \mathrm{Biext}^i(P,Q; \underline{G}_m) \quad , \quad i = 0,1 \quad .
\end{equation}
En effet, $\mathrm{BIEXT}(P,Q'';\underline{G}_m)$ s'identifie, à équivalence près, à la catégorie des biextensions $E$ de $(P,Q)$ par $\underline{G}_m$ \underline{munies} d'une trivialisation de la biextension induite $E'$ de $(P,Q')$ par $\underline{G}_m$ (VII 3.7.6\quad ). Or pour une biextension donnée, il y a en vertu de 3.4 exactement une trivialisation de $E'$. Bien entendu, on peut également prouver le fait
\SGApage{I}{263}
que $\mathrm{Biext}^i(P,Q'';\underline{G}_m) \longrightarrow \mathrm{Biext}^i(P,Q;\underline{G}_m)$ est un isomorphisme pour $i = 0,1$ (qui équivaut à la conclusion de 3.5) en utilisant la suite exacte (VII 3.7.5\quad ) des $\mathrm{Biext}^i$ associée à la suite exacte $0 \to Q' \to Q \to Q'' \to 0$.

\underline{Remarque} 3.6. Bien entendu, l'énoncé symétrique déduit de 3.5, obtenu en échangeant la gauche et la droite, est également valable, et se déduit d'ailleurs formellement de 3.5 par l'opération de passage aux biextensions symétriques (VII 2.7\quad ). Lorsqu'on suppose, \underline{en plus} des données et hypothèses de 3.5, qu'on a les données et hypothèses symétriques, i.e. $Q$ est à fibres connexes (ou encore $Q''$ à fibres connexes, cela revient au même puisque $Q'$ est à fibres connexes), et qu'on a une suite exacte $0 \longrightarrow P' \longrightarrow P \longrightarrow P'' \longrightarrow 0$\quad avec $P'$ un tore, alors l'application successive de 3.5, et de l'énoncé symétrique appliqué à la précédente suite exacte et à $Q''$ , montrent que le foncteur

\begin{equation}\tag{3.6.1}
\mathrm{BIEXT}(P'',Q''; \underline{G}_m) \stackrel{\approx}{\longrightarrow} \mathrm{BIEXT}(P,Q;\underline{G}_m)
\end{equation}
est une équivalence de catégories, et en particulier l'homomorphisme suivant est un isomorphisme

\SGAvar{\begin{equation}\tag{3.6.2}
\mathrm{Biext}^i(P'',Q'';\underline{G}_m) \stackrel{\sim}{\longrightarrow} \mathrm{Biext}^1(P,Q;\underline{G}_m) \quad , \quad i = 0, 1 \quad .
\end{equation}
\underline{Proposition} 3.7. \underline{Soient} $A$ \underline{un schéma abélien sur} $S$, $T$ \underline{un schéma en groupes de type multiplicatif} (SGA 3 IX) \underline{et de type fini sur} $S$, $\underline{M}=\underline{\mathrm{Hom}}(T,\underline{G}_m)$ \underline{le dual de Cartier de} $T$ (\underline{qui est un groupe constant tordu à engendrement fini sur} $S$ (3.3.3)). \underline{Alors les catégories} $\mathrm{EXT}(A,T)$ \underline{et} $\mathrm{BIEXT}(A,\underline{M}; \underline{G}_m)$, \underline{formées respectivement des extensions commutatives de} $A$ \underline{par} $T$ \underline{et des biextensions de} $(A,\underline{M})$ \underline{par} $\underline{G}_m$ , \underline{sont rigides}, \underline{et canoniquement équivalentes},
\SGApage{I}{264}
\underline{et les groupes des classes d'objets de ces catégories sont canoniquement} \underline{isomorphes à} $\mathrm{Hom}(M,A^*)$, \underline{où} $A^* = \underline{\mathrm{Ext}}^1(A,\underline{G}_m)$ \underline{est le schéma abélien dual de} $A$ (3.2). \underline{En particulier, on a un isomorphisme canonique}

}{\begin{equation}\tag{3.6.2}
\mathrm{Biext}^i(P'',Q'';\underline{G}_m) \stackrel{\sim}{\longrightarrow} \mathrm{Biext}^i(P,Q;\underline{G}_m) \quad , \quad i = 0, 1 \quad .
\end{equation}
\underline{Proposition} 3.7. \underline{Soient} $A$ \underline{un schéma abélien sur} $S$, $T$ \underline{un schéma en groupes de type multiplicatif} (SGA 3 IX) \underline{et de type fini sur} $S$, $\underline{M}=\underline{\mathrm{Hom}}(T,\underline{G}_m)$ \underline{le dual de Cartier de} $T$ (\underline{qui est un groupe constant tordu à engendrement fini sur} $S$ (3.3.3)). \underline{Alors les catégories} $\mathrm{EXT}(A,T)$ \underline{et} $\mathrm{BIEXT}(A,\underline{M}; \underline{G}_m)$, \underline{formées respectivement des extensions commutatives de} $A$ \underline{par} $T$ \underline{et des biextensions de} $(A,\underline{M})$ \underline{par} $\underline{G}_m$ , \underline{sont rigides}, \underline{et canoniquement équivalentes},
\SGApage{I}{264}
\underline{et les groupes des classes d'objets de ces catégories sont canoniquement} \underline{isomorphes à} $\mathrm{Hom}(M,A^*)$, \underline{où} $A^* = \underline{\mathrm{Ext}}^1(A,\underline{G}_m)$ \underline{est le schéma abélien dual de} $A$ (3.2). \underline{En particulier, on a un isomorphisme canonique}

}{F7I-VIII-032}\begin{equation}\tag{3.7.1}
\mathrm{Ext}^1(A,T) \stackrel{\sim}{\longrightarrow} \mathrm{Hom}(\underline{M},A^*) \qquad .
\end{equation}
C'est en effet le cas particulier de 1.6 obtenu en faisant $P = A$, $Q = M$ ; en vertu du théorème de bidualité (SGA 3 X 5.7), on a bien $T \simeq \underline{\mathrm{Hom}}(M,\underline{G}_m)$, et il suffit de vérifier les conditions (1.6.1), qui s'écrivent ici

\SGAvar{\begin{equation*}
\underline{\mathrm{Hom}}(A,\underline{G}_m) = 0 \qquad , \qquad \mathrm{Hom}(A,\underline{\mathrm{Ext}}^1(\underline{M},\underline{G}_m)) = 0 \qquad .
\end{equation*}
La première relation a été déjà été rappelée dans (3.2.1), la deuxième peut se préciser en

\begin{equation*}
\underline{\mathrm{Ext}}^1(\underline{M},\underline{G}_m) = 0 \qquad .
\end{equation*}
Pour vérifier cette dernière, la question étant locale sur $S$, on peut supposer $\underline{M}$ constant $= M_S$, et en décomposant $M$ en somme de $\mathbb{Z}$-modules monogènes, on est ramené au cas où $M = \mathbb{Z}$, qui est trivial, et au cas $\mathbb{Z} = \mathbb{Z}/n\mathbb{Z}$ , qui provient du fait que $\underline{G}_m$ est $n$-divisible (ou, plus savamment, de 3.3.1).

}{\begin{equation*}
\underline{\mathrm{Hom}}(A,\underline{G}_m) = 0 \qquad , \qquad \mathrm{Hom}(A,\underline{\mathrm{Ext}}^1(\underline{M},\underline{G}_m)) = 0 \qquad .
\end{equation*}
La première relation a déjà été rappelée dans (3.2.1), la deuxième peut se préciser en

\begin{equation*}
\underline{\mathrm{Ext}}^1(\underline{M},\underline{G}_m) = 0 \qquad .
\end{equation*}
Pour vérifier cette dernière, la question étant locale sur $S$, on peut supposer $\underline{M}$ constant $= M_S$, et en décomposant $M$ en somme de $\mathbb{Z}$-modules monogènes, on est ramené au cas où $M = \mathbb{Z}$, qui est trivial, et au cas $M = \mathbb{Z}/n\mathbb{Z}$ , qui provient du fait que $\underline{G}_m$ est $n$-divisible (ou, plus savamment, de 3.3.1).

}{F7I-VIII-033}On notera que lorsque $T = \underline{G}_m$ donc $\underline{M} = \mathbb{Z}_S$ , (3.7.1) se réduit à l'isomorphisme identique (3.2.3). Bien entendu, la démonstration directe de (3.7.1) serait essentiellement triviale (sans faire intervenir le groupe $\mathrm{Ext}^1(A \stackrel{\mathbb{L}}{\otimes} \underline{M},\underline{G}_m) =\mathrm{Biext}^1(A,\underline{M},\underline{G}_m))$.


\SGApage{I}{265}


\SGAsectiontoc{4. Extensions et biextensions des schémas en groupes lisses et connexes sur un corps}{4}
4. \underline{Extensions et biextensions des schémas en groupes lisses et connexes sur un corps}

Rappelons le lemme suivant dû à Rosenlicht [3 bis, Prop.3] :

\underline{Lemme} 4.1. \underline{Soient} $k$ \underline{un corps}, $X$ \underline{et} $Y$ \underline{deux $k$-schémas de type fini, séparables et géométriquement connexes, munis de points} $e_X$ \underline{resp.} $e_Y$ \underline{rationnels sur} $k$. \underline{Alors tout morphisme} $X \times Y \longrightarrow \underline{G}_m$ \underline{qui est trivial sur les deux sous-schémas} $X \times e_Y$ \underline{et} $e_X \times Y$ \underline{est trivial}.

Nous n'utiliserons ce résultat que lorsque $X$ et $Y$ sont des schémas en groupes sur $k$, auquel cas il résultera des hypothèses faites que $X$ et $Y$ sont même lisses sur $k$. Notons que par récurrence, le résultat 4.1 implique le résultat analogue pour un nombre quelconque de facteurs. Ceci nous permet donc d'appliquer les résultats de VII 1.3 et de VII 2.9. On obtient ainsi les deux propositions suivantes :

\SGAvar{\underline{Proposition} 4.2. \underline{Soient} $P$ \underline{un groupe lisse et connexe sur le corps} $k$. \underline{Alors le foncteur naturel de la catégorie des extensions de Groupes de} $P$ \underline{par} $\underline{G}_m$ \underline{dans la catégorie des Modules inversibles sur} $P$ \underline{rigidifiés en l'élément neutre} $e_P$ \underline{de} $P$ \underline{est pleinement fidèle}. \underline{Son image essentielle est formée des Modules inversibles rigidifiés} $\underline{L}$ \underline{sur} $P$ \underline{tels que l'on ait} $\pi^*(P) \simeq \mathrm{pr}_1^*(P)\mathrm{pr}_2^*(P)$, \underline{où} $\pi : P \times P \longrightarrow P$ \underline{est la loi de composition de} $P$, \underline{et} $\mathrm{pr}_1$ \underline{et} $\mathrm{pr}_2$ \underline{sont les deux projections du produit}.

}{\underline{Proposition} 4.2. \underline{Soit} $P$ \underline{un groupe lisse et connexe sur le corps} $k$. \underline{Alors le foncteur naturel de la catégorie des extensions de Groupes de} $P$ \underline{par} $\underline{G}_m$ \underline{dans la catégorie des Modules inversibles sur} $P$ \underline{rigidifiés en l'élément neutre} $e_P$ \underline{de} $P$ \underline{est pleinement fidèle}. \underline{Son image essentielle est formée des Modules inversibles rigidifiés} $\underline{L}$ \underline{sur} $P$ \underline{tels que l'on ait} $\pi^*(L) \simeq \mathrm{pr}_1^*(L)\otimes\mathrm{pr}_2^*(L)$, \underline{où} $\pi : P \times P \longrightarrow P$ \underline{est la loi de composition de} $P$, \underline{et} $\mathrm{pr}_1$ \underline{et} $\mathrm{pr}_2$ \underline{sont les deux projections du produit}.

}{F7I-VIII-034}\underline{Proposition} 4.3. \underline{Soient} $P$ et $Q$ \underline{deux Groupes algébriques lisses et connexes sur le corps} $k$. \underline{Alors le foncteur naturel de la catégorie des biextensions de} $(P,Q)$ \underline{par} $\underline{G}_m$, \underline{dans la catégorie des Modules inversibles sur} $P \times Q$ \underline{birigidifiés par rapport aux sections unités} $e_P,e_Q$ \underline{de} $P$ \underline{resp.} $Q$,
\SGApage{I}{266}
\underline{est pleinement fidèle}, \underline{les deux catégories envisagées étant par ailleurs rigides}. \underline{Pour qu'un Module inversible birigidifié} $\mathbb{L}$ \underline{sur} $P \times Q$ \underline{provienne d'une biextension}, \underline{il faut et il suffit que ses images inverses sur} $P \times P \times Q$ \underline{et} $P \times Q \times Q$ \underline{par} $\pi_P \times \mathrm{id}_Q$ \underline{et} $\mathrm{id}_P \times \pi_Q$ \underline{satisfassent au} ``\underline{théorème du cube}'' ([1] \underline{ou} [2]).

Cette dernière condition n'est en effet que traduction des conditions envisagées dans VII 2.9.3 b), dans l'explicitation qui en est donnée dans VII 3.5 b).

4.4. J'ignore si la condition du cube envisagée dans 4.3 est automatiquement vérifiée. On sait tout au moins [2, IV 2.6] que pour $\mathbb{L}$ donné, il existe toujours un entier $n \geqslant 1$, puissance de l'exposant caractéristique de $k$, tel que $\mathbb{L}^{\otimes n}$ satisfasse à la condition envisagée, donc provienne d'une biextension de $(P,Q)$ par $\underline{G}_m$ ; et que l'on peut toujours trouver une extension radicielle $k'$ de $k$, telle que la condition voulue devienne vraie après l'extension de la base $k \longrightarrow k'$, de sorte qu'après cette extension $\mathbb{L}$ donne un Module inversible birigidifié qui provienne d'une biextension de $(P,Q)$ par $\underline{G}_m$ . En particulier, si $k$ est parfait, la condition envisagée dans 4.3 est toujours vérifiée. Ce dernier fait, qui implique formellement les précédents par un argument de trace, peut se voir aussi en appliquant 4.7 plus bas, cf. 4.9.

Rappelons sous forme de lemme le résultat suivant de EGA Err$_{\mathrm{IV}}$ 53, 21.4.13 :

\SGApage{I}{267}

\SGAvar{\underline{Lemme} 4.5. \underline{Soient} $f : X \longrightarrow Y$ \underline{un morphisme lisse surjectif à fibres intègres}, \underline{avec} $Y$ \underline{normal intègre de point générique} $y$, $\mathbb{L}$ \underline{un Module inversible sur} $X$ \underline{dont la restriction à} $X_y$ \underline{soit trivial}. \underline{Alors il existe un Module inversible} $\underline{M}$ \underline{sur} $Y$ \underline{tel que} $\mathbb{L}$ \underline{soit isomorphe à} $f^*(\underline{M})$.

}{\underline{Lemme} 4.5. \underline{Soient} $f : X \longrightarrow Y$ \underline{un morphisme lisse surjectif à fibres intègres}, \underline{avec} $Y$ \underline{normal intègre de point générique} $y$, $\mathbb{L}$ \underline{un Module inversible sur} $X$ \underline{dont la restriction à} $X_y$ \underline{soit triviale}. \underline{Alors il existe un Module inversible} $\underline{M}$ \underline{sur} $Y$ \underline{tel que} $\mathbb{L}$ \underline{soit isomorphe à} $f^*(\underline{M})$.

}{F7I-VIII-035}\underline{Proposition} 4.6. \underline{Soient} $P,Q$ \underline{deux groupes algébriques lisses et connexes sur le corps} $k$, \underline{avec} $P$ affine. \underline{On suppose de plus que} $P$ \underline{admet une suite de composition dont chacun des facteurs est un tore}, \underline{ou est isomorphe au groupe additif} $\underline{G}_a$ (\underline{condition vérifiée en particulier si} $k$ \underline{est parfait} [SGA 3, XVII 7.2.1, 4.1.1 (i) $\Longleftrightarrow$ (ii) \underline{et} 4.1.5]), \underline{ou que} $Q$ \underline{soit extension d'une variété abélienne} $B$ \underline{par un groupe algébrique lisse, connexe et affine satisfaisant à la condition précédente}. \underline{Sous ces conditions}, \underline{tout Module inversible birigidifié} $\mathbb{L}$ \underline{sur} $P \times Q$ \underline{est trivial}, \underline{à fortiori} (4.3) \underline{toute biextension de} $(P,Q)$ \underline{par} $\underline{G}_m$ \underline{est triviale}.

Signalons d'abord que comme tout automorphisme du Module birigidifié est l'identité (4.3), la théorie de la descente nous montre qu'il suffit de prouver que $\mathbb{L}$ devient trivial après extension séparable finie de $k$, ce qui nous ramène aussitôt, par passage à la limite, au cas où $k$ est séparablement clos. Comme dans ce cas tout tore sur $k$ est déployé, on voit que si $P$ admet une suite de composition comme indiquée dans 4.6, il admet aussi une suite de composition à quotients isomorphes à $\underline{G}_a$ ou $\underline{G}_m$, ce qui implique, par application répétée de 4.5, compte tenu de $\mathrm{Pic}\ (\underline{G}_{a_K}) = \mathrm{Pic}(\underline{G}_{m_K}) = 0$ pour tout corps $K$, que l'on a

\begin{equation*}
\mathrm{Pic}(P) \;=\; 0 \qquad .
\end{equation*}

\SGApage{I}{268}

Comme cette relation restera encore vraie par tout changement du corps de base $k$, on voit, en appliquant 4.5 à $\mathrm{pr}_2 : P \times Q \longrightarrow Q$ et au Module $\mathbb{L}$ , que $\mathbb{L}$ est isomorphe à un Module de la forme $\mathrm{pr}_2^*(\underline{M})$, où $\underline{M}$ est un Module inversible sur $Q$. Comme la restriction de $\mathbb{L}$ à $e_P \times Q$ est triviale, on en conclut aussitôt que $\underline{M}$ est trivial, donc $\mathbb{L}$ est trivial, ce qui établit la conclusion dans le premier cas envisagé. Dans le deuxième, où $Q$ est une extension d'un schéma abélien $B$ par un groupe $M$ admettant une suite de composition à quotients isomorphes à $\underline{G}_m$ ou à $\underline{G}_a$ , on applique 4.5 au morphisme

\SGAvar{\begin{equation*}
f : X = P \times Q \longrightarrow Y = P \times B \qquad ,
\end{equation*}
dont la fibre au point générique $y$ de $P \times B$ est un torseur sous $M_y$, nécessairement trivial en vertu de la structure particulière de $M$ et des relations classiques $H^1(K,\underline{G}_a) = 0$, $H^1(K,\underline{G}_m) = 0$, valables pour tout corps $K$ (qu'on prendra ici égal à $k(y)$). Comme on aura $\mathrm{Pic}(M) = 0$, donc $\mathrm{Pic}(X_y) = 0$, 4.5 est bien applicable et montre que $\mathbb{L}$ est isomorphe à l'image inverse d'un Module inversible $\underline{M}$ sur $Y = P \times B$. Quitte à remplacer $\underline{M}$ par $\underline{M}\,\mathrm{pr}_1^*(\underline{M}_1)^{-1}\mathrm{pr}_2^*(M_2)^{-1}$, où $\mathrm{pr}_1$ et $\mathrm{pr}_2$ sont les deux projections de $Y=P \times A$, et où $\underline{M}_1$ et $\underline{M}_2$ sont les deux restrictions de $\underline{M}$ à $P \times e_B$ et à $e_P \times B$ respectivement, on peut supposer $\underline{M}$ birigidifié, et que $\mathbb{L}$ est son image inverse birigidifiée. Cela nous ramène au cas où $Q$ est lui-même un schéma abélien. Mais alors $\mathbb{L}$ définit un morphisme canonique

\begin{equation}\tag{*}
P \longrightarrow \underline{\mathrm{Pic}}_{Q/k} \qquad ,
\end{equation}
transformant unité en unité, et comme $P$ est connexe, ce morphisme se
\SGApage{I}{269}
factorise par le schéma abélien dual $Q' = \underline{\mathrm{Pic}}^{\mathrm{o}}_{Q/k}$ . Or il est bien connu qu'un morphisme de schémas d'un groupe algébrique affine lisse et connexe dans un schéma abélien est constant. (C'est une question ``géométrique'' i.e. on peut supposer $k$ algébriquement clos, et dans ce cas on est ramené aussitôt par dévissage au cas où $P$ est $\underline{G}_a$ ou $\underline{G}_m$ ; or il est bien connu que toute application rationnelle de $\underline{P}^1_k$ dans un schéma abélien sur $k$ est constante.) Le fait que (*) soit constant, donc nul, s'interprète par la condition que $\mathbb{L}$ soit isomorphe à l'image inverse d'un Module inversible $\underline{M}$ sur $P$, et comme $\mathbb{L}$ est trivial sur $P \times e_A$ , $\underline{M}$ est trivial, donc $\mathbb{L}$ l'est également, cqfd.

\underline{Proposition} 4.7. \underline{Soient} $P,Q$ \underline{deux groupes algébriques lisses et connexes} \underline{sur le corps} $k$. \underline{Supposons que} $P$ \underline{soit extension d'un schéma abélien} $A$ \underline{par un groupe algébrique lisse}, \underline{connexe et affine} $L$ \underline{admettant une suite} \underline{de composition dont chaque facteur soit isomorphe à un tore ou à} $\underline{G}_a$ (\underline{condition toujours remplie si} $k$ \underline{est parfait}). \underline{Alors le foncteur envisagé dans} 1.3 \underline{est une équivalence de catégories}, \underline{i.e.} \underline{tout Module inversible birigidifié} $\mathbb{L}$ \underline{sur} $P \times Q$ \underline{provient d'une biextension de} $(P,Q)$. \underline{De plus}, \underline{le foncteur de catégories rigides}

}{\begin{equation*}
f : X = P \times Q \longrightarrow Y = P \times B \qquad ,
\end{equation*}
dont la fibre au point générique $y$ de $P \times B$ est un torseur sous $M_y$, nécessairement trivial en vertu de la structure particulière de $M$ et des relations classiques $H^1(K,\underline{G}_a) = 0$, $H^1(K,\underline{G}_m) = 0$, valables pour tout corps $K$ (qu'on prendra ici égal à $k(y)$). Comme on aura $\mathrm{Pic}(M) = 0$, donc $\mathrm{Pic}(X_y) = 0$, 4.5 est bien applicable et montre que $\mathbb{L}$ est isomorphe à l'image inverse d'un Module inversible $\underline{M}$ sur $Y = P \times B$. Quitte à remplacer $\underline{M}$ par $\underline{M}\,\mathrm{pr}_1^*(\underline{M}_1)^{-1}\mathrm{pr}_2^*(M_2)^{-1}$, où $\mathrm{pr}_1$ et $\mathrm{pr}_2$ sont les deux projections de $Y=P \times B$, et où $\underline{M}_1$ et $\underline{M}_2$ sont les deux restrictions de $\underline{M}$ à $P \times e_B$ et à $e_P \times B$ respectivement, on peut supposer $\underline{M}$ birigidifié, et que $\mathbb{L}$ est son image inverse birigidifiée. Cela nous ramène au cas où $Q$ est lui-même un schéma abélien. Mais alors $\mathbb{L}$ définit un morphisme canonique

\begin{equation}\tag{*}
P \longrightarrow \underline{\mathrm{Pic}}_{Q/k} \qquad ,
\end{equation}
transformant unité en unité, et comme $P$ est connexe, ce morphisme se
\SGApage{I}{269}
factorise par le schéma abélien dual $Q' = \underline{\mathrm{Pic}}^{\mathrm{o}}_{Q/k}$ . Or il est bien connu qu'un morphisme de schémas d'un groupe algébrique affine lisse et connexe dans un schéma abélien est constant. (C'est une question ``géométrique'' i.e. on peut supposer $k$ algébriquement clos, et dans ce cas on est ramené aussitôt par dévissage au cas où $P$ est $\underline{G}_a$ ou $\underline{G}_m$ ; or il est bien connu que toute application rationnelle de $\underline{P}^1_k$ dans un schéma abélien sur $k$ est constante.) Le fait que (*) soit constant, donc nul, s'interprète par la condition que $\mathbb{L}$ soit isomorphe à l'image inverse d'un Module inversible $\underline{M}$ sur $P$, et comme $\mathbb{L}$ est trivial sur $P \times e_Q$ , $\underline{M}$ est trivial, donc $\mathbb{L}$ l'est également, cqfd.

\underline{Proposition} 4.7. \underline{Soient} $P,Q$ \underline{deux groupes algébriques lisses et connexes} \underline{sur le corps} $k$. \underline{Supposons que} $P$ \underline{soit extension d'un schéma abélien} $A$ \underline{par un groupe algébrique lisse}, \underline{connexe et affine} $L$ \underline{admettant une suite} \underline{de composition dont chaque facteur soit isomorphe à un tore ou à} $\underline{G}_a$ (\underline{condition toujours remplie si} $k$ \underline{est parfait}). \underline{Alors le foncteur envisagé dans} 4.3 \underline{est une équivalence de catégories}, \underline{i.e.} \underline{tout Module inversible birigidifié} $\mathbb{L}$ \underline{sur} $P \times Q$ \underline{provient d'une biextension de} $(P,Q)$. \underline{De plus}, \underline{le foncteur de catégories rigides}

}{F7I-VIII-036}\begin{equation*}
\mathrm{BIEXT}(A,Q;\underline{G}_m) \longrightarrow \mathrm{BIEXT}(P,Q;\underline{G}_m)
\end{equation*}
\underline{est une équivalence de catégories}, \underline{en d'autres termes}, \underline{il donne naissance} \underline{à un isomorphisme de groupes} :

\begin{equation}\tag{4.7.1}
\mathrm{Biext}^1(A,Q;\underline{G}_m) \stackrel{\sim}{\longrightarrow} \mathrm{Biext}^1(P,Q;\underline{G}_m) \qquad .
\end{equation}

\SGApage{I}{270}

\SGAvar{Pour la première assertion, on se ramène encore, par descente et en utilisant 4.3, au cas où $k$ est séparablement clos. Procédant comme dans la démonstration de 4.6, on déduit de 4.5 et de l'hypothèses sur $L$ que $\mathbb{L}$ est isomorphe à l'image inverse d'un Module birigidifié $\underline{M}$ sur $A \times Q$, et pour la première assertion de 4.7, on est donc ramené au cas $A=P$. Or $\mathbb{L}$ définit un morphisme de schémas ponctués

}{Pour la première assertion, on se ramène encore, par descente et en utilisant 4.3, au cas où $k$ est séparablement clos. Procédant comme dans la démonstration de 4.6, on déduit de 4.5 et de l'hypothèse sur $L$ que $\mathbb{L}$ est isomorphe à l'image inverse d'un Module birigidifié $\underline{M}$ sur $A \times Q$, et pour la première assertion de 4.7, on est donc ramené au cas $A=P$. Or $\mathbb{L}$ définit un morphisme de schémas ponctués

}{F7I-VIII-037}\begin{equation*}
f:Q \longrightarrow \underline{\mathrm{Pic}}_{P/k} \quad ,
\end{equation*}
qui se factorise nécessairement par le schéma abélien dual $P'$ de $P$. Soit $\mathbb{L}'$ l'image inverse du faisceau de Weil sur $P \times P'$ par

\begin{equation*}
\mathrm{id}_P \times f \; : \; P \times Q \longrightarrow P \times P' \quad ,
\end{equation*}
c'est un Module birigidifié qui, par définition de $f$, ne diffère de $\mathbb{L}$ que modulo l'image inverse d'un Module inversible $\underline{M}$ sur $P$. Comme $\mathbb{L}'$ et $\mathbb{L}$ ont une restriction triviale à $P \times e_Q$, on en conclut que $\underline{M}$ est trivial, donc que $\mathbb{L}$ est isomorphe à $\mathbb{L}'$. Or nous avons signalé dans VII 2.9.4 que sur un produit de deux schémas abéliens, les correspondances divisorielles provenaient de biextensions. Il en résulte que $\mathbb{L}'$, ou ce qui revient au même, $\mathbb{L}$ , provient d'une biextension.

Revenant au cas général ($P$ extension de $A$ par $L$), il reste à prouver que (4.7.1) est bien un isomorphisme. Mais cela résulte de la suite exacte des $\mathrm{Biext}^i$ ($i=0,1$) associée à la suite exacte

\begin{equation*}
0 \longrightarrow L \longrightarrow P \longrightarrow A \longrightarrow 0 \quad ,
\end{equation*}
et des relations

\begin{equation*}
\mathrm{Biext}^0(L,Q;\underline{G}_m) = 0 \quad , \quad \mathrm{Biext}^1(L,Q;\underline{G}_m) = 0 \quad , \quad 
\end{equation*}

\SGApage{I}{271}

 dont la première est un cas particulier de l'assertion de rigidité faite dans 4.3, et la deuxième n'est autre que le premier cas de 4.6.

\underline{Corollaire} 4.8. \underline{Sous les conditions de} 4.7, \underline{supposons que} $Q$ \underline{soit extension} \underline{d'un groupe algébrique} $B$ (\underline{nécessairement lisse et connexe}) \underline{par un groupe} \underline{algébrique} $M$ \underline{qui est lisse}, \underline{connexe et} affine. \underline{Alors le foncteur de} \underline{catégories rigides}

\SGAvar{\begin{equation*}
\mathrm{BIEXT}(A,B,\underline{G}_m) \longrightarrow \mathrm{BIEXT}(P,Q;\underline{G}_m)
\end{equation*}
\underline{est une équivalence de catégories,} \underline{en d'autres termes,} \underline{elle donne naissanee} \underline{à un isomorphisme de groupes}

}{\begin{equation*}
\mathrm{BIEXT}(A,B,\underline{G}_m) \longrightarrow \mathrm{BIEXT}(P,Q;\underline{G}_m)
\end{equation*}
\underline{est une équivalence de catégories,} \underline{en d'autres termes,} \underline{elle donne naissance} \underline{à un isomorphisme de groupes}

}{F7I-VIII-038}\begin{equation}\tag{4.8.1}
\mathrm{Biext}^1(A,B;\underline{G}_m) \stackrel{\sim}{\longrightarrow} \mathrm{Biext}^1(P,Q;\underline{G}_m) \quad .
\end{equation}

Compte tenu de (4.7.1), il reste à prouver que l'on a un isomorphisme

\SGAvar{\begin{equation*}
\mathrm{Biext}^1(A,B;\underline{G}_m) \stackrel{\sim}{\longrightarrow} \mathrm{Biext}^1(A,Q;\underline{G}_m) \quad .
\end{equation*}
CeIa se déduit comme (4.7.1) de la suite exacte des $\mathrm{Biext}^i$ déduite de la suite exacte

}{\begin{equation*}
\mathrm{Biext}^1(A,B;\underline{G}_m) \stackrel{\sim}{\longrightarrow} \mathrm{Biext}^1(A,Q;\underline{G}_m) \quad .
\end{equation*}
Cela se déduit comme (4.7.1) de la suite exacte des $\mathrm{Biext}^i$ déduite de la suite exacte

}{F7I-VIII-039}\begin{equation*}
0 \longrightarrow M \longrightarrow Q \longrightarrow B \longrightarrow 0 \quad ,
\end{equation*}
et des relations

\begin{equation*}
\mathrm{Biext}^0(P,M;\underline{G}_m) = 0 \quad , \quad \mathrm{Biext}^1(P,M;\underline{G}_m) = 0 \quad ,
\end{equation*}
dont la deuxième est un cas particulier du deuxième cas envisagé dans 4.6 (où on aurait échangé le rôle des deux facteurs envisagés).


\SGApage{I}{272}

\underline{Remarque} 4.9. Soient $P,Q$ deux groupes algébriques lisses et connexes sur un corps \underline{parfait} $k$. On sait alors [3] que chacun d'eux est extension d'un schéma abélien par un groupe lisse, connexe et affine, ce dernier étant produit d'un tore par une extension multiple de groupes $\underline{G}_a$ (cf. référence donnée dans 4.6) :

\begin{equation*}
0 \longrightarrow L \longrightarrow P \longrightarrow A \longrightarrow 0 \quad , \quad 0 \longrightarrow M \longrightarrow Q \longrightarrow B \longrightarrow 0 \quad .
\end{equation*}
Donc 4.8 s'applique, et nous montre que la théorie des biextensions de $(P,Q)$ par $\underline{G}_m$ se ramène à la théorie analogue (et combien classique !) pour le couple $(A,B)$ de variétés abéliennes.

\underline{Corollaire} 4.10. \underline{Soient} $P$ \underline{et} $Q$ \underline{deux groupes algébriques lisses et connexes sur} $k$, $L$ \underline{un sous-groupe algébrique} affine \underline{lisse et connexe de} $P$, $E$ \underline{une biextension de} $(P,Q)$ \underline{par} $\underline{G}_m$, \underline{d'où par} 2.2 \underline{un accouplement déduit de} $E$

\begin{equation*}
T_\ell(P) \times T_\ell(Q) \longrightarrow T_\ell(\underline{G}_m) \quad .
\end{equation*}
\underline{Alors l'accouplement induit}

\begin{equation*}
T_\ell(L) \times T_\ell(Q) \longrightarrow T_\ell(\underline{G}_m)
\end{equation*}
\underline{est nul}.

En effet, on peut supposer évidemment $k$ algébriquement clos, auquel cas on est sous les conditions de 4.6, qui impliquent que la restriction de $E$ à $(L,Q)$ est la biextension triviale, d'où la conclusion, puisque l'accouplement induit envisagé est celui défini par cette restriction.


\SGApage{I}{273}

4.11.  Soient $P,Q$ deux groupes algébriques et lisses sur $k$, $E$ une biextension de $(P,Q)$ par $\underline{G}_m$ , et $T$ le plus grand sous-tore de $P$. On se propose de déterminer l'annulateur gauche de l'accouplement

\begin{equation}\tag{4.11.1}
T_\ell(P) \times T_\ell(Q) \longrightarrow T_\ell(\underline{G}_{m_k}) \quad ,
\end{equation}
$\ell$ étant encore un nombre premier fixé.  En vertu de 4.10, cet annulateur contient en tous cas le sous-progroupe $T_\ell(T)$ de $T_\ell(P)$. Considérons une extension $k'$ de $k$ au-dessus de laquelle $P$ et $Q$ s'écrivent comme extension d'un schéma abélien par un schéma en groupes affine, extension multiple de tores et de groupes $\underline{G}_a$ ; lorsque $P$ et $Q$ sont de rangs unipotents nuls, on peut prendre $k'=k$ ; en tous cas, on peut prendre pour $k'$ la clôture parfaite (ou une extension finie radicielle) de $k$. Soient $A'$ resp. $B'$ la partie abélienne de $P_{k'}$ resp. $Q_{k'}$ . En vertu de 4.7, la biextension $E_{k'}$ de $(P_{k'},Q_{k'})$ provient d'une biextension $F'$ de $(A',B')$, bien déterminée à isomorphisme unique près. Donc l'accouplement

\begin{equation}\tag{4.11.2}
T_\ell(P_{k'}) \times T_\ell(Q_{k'}) \longrightarrow T_\ell(\underline{G}_{mk'})
\end{equation}
déduit de (4.11.1) par changement du corps de base $k \longrightarrow k'$, et qui est aussi l'accouplement défini par la biextension $E_{k'}$ de $(P_{k'},Q_{k'})$, s'obtient à l'aide de l'accouplement

\begin{equation}\tag{4.11.3}
T_\ell(A') \times T_\ell(B') \longrightarrow T_\ell(\underline{G}_{mk'})
\end{equation}
défini par la biextension $F'$, en composant avec les homomorphismes canoniques (qui sont en fait des \underline{épimorphismes} si $P$ et $Q$ sont $\ell$-divisibles).

\begin{equation}\tag{4.11.4}
T_\ell(P_{k'}) \longrightarrow T_\ell(A') \quad , \quad T_\ell(Q_{k'}) \longrightarrow T_\ell(B') \quad .
\end{equation}

\SGApage{I}{274}

Par suite, l'annulateur gauche de l'accouplement (4.11.2) (qui est déduit de l'annulateur gauche de (4.11.1) par changement de base) n'est autre que l'image inverse par le premier homomorphisme (4.11.4) de l'annulateur gauche de l'accouplement (4.11.3). Si $C'$ est le plus grand sous-schéma abélien de $A'$ tel que la restriction de $F'$ à $(C',B')$ soit la biextension triviale (donc $C'$ est aussi la composante neutre réduite de l'homomorphisme

\SGAvar{\begin{equation}\tag{4.11.5}
A' \longrightarrow \text{schéma abélien dual de } B'
\end{equation}
défini par $F'$), l'annulateur à gauche de (4.11.3) s'identifie à $T_{\ell}(C')$ ; donc l'annulateur à gauche de l'accouplement (1.11.1) n'est autre que l'image inverse de $T_{\ell}(C') \hookrightarrow T_{\ell}(A')$ par le premier morphisme (4.11.4).

En particulier, dire que l'annulateur à gauche de l'accouplement (4.11.1) est réduit à $T_{\ell}(T)$, i.e. que l'annulateur à gauche de (4.11.2) est réduit au noyau $T_{\ell}(T')$ du premier morphisme (4.11.4), revient à dire que l'accouplement (4.11.3) est séparant à gauche i.e, que $C' = 0$, ou encore que le morphisme (4.11.5) défini par $F'$ est à noyau fini. On notera sous cette dernière forme que cette condition ne dépend pas de $\ell$, et il est clair (sous n'importe quelle forme) qu'elle ne dépend pas du choix de l'extension $k'$. Lorsque cette condition est remplie, nous dirons simplement que la biextension donnée $E$ de $(P,Q)$ par $\underline{G}_m$ est \underline{séparante à gauche}. On notera par dualité que cela signifie aussi que le morphisme

}{\begin{equation}\tag{4.11.5}
A' \longrightarrow \text{schéma abélien dual de } B'
\end{equation}
défini par $F'$), l'annulateur à gauche de (4.11.3) s'identifie à $T_{\ell}(C')$ ; donc l'annulateur à gauche de l'accouplement (4.11.1) n'est autre que l'image inverse de $T_{\ell}(C') \hookrightarrow T_{\ell}(A')$ par le premier morphisme (4.11.4).

En particulier, dire que l'annulateur à gauche de l'accouplement (4.11.1) est réduit à $T_{\ell}(T)$, i.e. que l'annulateur à gauche de (4.11.2) est réduit au noyau $T_{\ell}(T')$ du premier morphisme (4.11.4), revient à dire que l'accouplement (4.11.3) est séparant à gauche i.e. que $C' = 0$, ou encore que le morphisme (4.11.5) défini par $F'$ est à noyau fini. On notera sous cette dernière forme que cette condition ne dépend pas de $\ell$, et il est clair (sous n'importe quelle forme) qu'elle ne dépend pas du choix de l'extension $k'$. Lorsque cette condition est remplie, nous dirons simplement que la biextension donnée $E$ de $(P,Q)$ par $\underline{G}_m$ est \underline{séparante à gauche}. On notera par dualité que cela signifie aussi que le morphisme

}{F7I-VIII-040}\begin{equation}\tag{4.11.6}
B' \longrightarrow \text{schéma abélien dual de } A'
\end{equation}
défini par $F'$ est un épimorphisme (i.e. surjectif).


\SGApage{I}{275}

On définit de façon symétrique la notion de \underline{biextension séparante à droite}, qui correspond au cas où (4.11.6) est à noyau fini, i.e. (4.11.5) un épimorphisme. On dit que la biextension $E$ de $(P,Q)$ par $\underline{G}_{mk}$ est \underline{séparante} si elle est séparante à gauche et à droite, ce qui revient à dire que (4.11.5) est une isogénie, ou ce qui revient au même, que (4.11.6) est une isogénie.

Il est évident par définition que ces notions sont invariantes par toute extension du corps de base $k$.

4.12. Soient $P$ un schéma en groupes lisse et connexe, $\mathbb{L}$ un Module inversible sur $P$, et posons

\begin{equation}\tag{4.12.1}
\delta(\mathbb{L}) = \pi^{*}(\mathbb{L})\ \mathrm{pr}_{1}^{*}(\mathbb{L})^{-1}\mathrm{pr}_{2}^{*}(\mathbb{L})^{-1} \qquad ,
\end{equation}
les notations étant celles de 4.2. Comme on a déjà signalé dans VII 1.3.4, on trouve ainsi un Module inversible birigidifié sur $PXP$. En vertu de 4.3, si $\delta(\mathbb{L})$ provient d'une biextension, celle-ci est déterminée à isomorphisme unique près. Cette biextension est définie en particulier,en vertu de 4.7,si $P$ est une extension d'un schéma abélien $A$ par une extension multiple de tores et de groupes $\underline{G}_a$. En tous cas, si elle est définie, il est clair que cette biextension est symétrique. Dans le cas favorable précédent, et lorsqu'on suppose le tore maximal de $P$ déployé, on peut ramener la construction de $\delta(\mathbb{L})$ au cas où $P$ est un schéma abélien, en notant qu'elle est évidemment fonctorielle en $P$, et que d'autre part l'homomorphisme

\begin{equation}\tag{4.12.2}
\mathrm{Pic}(A) \longrightarrow \mathrm{Pic}(P) 
\end{equation}

\SGApage{I}{276}

\SGAvar{ déduit de l'homomorphisme canonique $f : A \longrightarrow P$ est surjectif grâce à 4.5 (cf. début démonstration de 4.6), de sorte que l'on peut trouver un isomorphisme

\begin{equation*}
\mathbb{L}^{*} \simeq f^{*}(\underline{M}) \qquad ,
\end{equation*}
où $\underline{M}$ est un Module inversible sur $\underline{M}$. Alors $\delta(\mathbb{L})$ ``est'' l'image inverse de la biextension $\delta(\underline{M})$ de $(A,A)$ par $\underline{G}_{mk}$ .

}{ déduit de l'homomorphisme canonique $f : P \longrightarrow A$ est surjectif grâce à 4.5 (cf. début démonstration de 4.6), de sorte que l'on peut trouver un isomorphisme

\begin{equation*}
\mathbb{L} \simeq f^{*}(\underline{M}) \qquad ,
\end{equation*}
où $\underline{M}$ est un Module inversible sur $A$. Alors $\delta(\mathbb{L})$ ``est'' l'image inverse de la biextension $\delta(\underline{M})$ de $(A,A)$ par $\underline{G}_{mk}$ .

}{F7I-VIII-041}\underline{Proposition} 4.13. \underline{Sous les conditions de} 4.12, \underline{supposons que} $\mathbb{L}$ \underline{soit un Module inversible} ample \underline{sur} $P$. \underline{Alors si la biextension} $\delta(\mathbb{L})$ \underline{de} $(P,P)$ \underline{par} $\underline{G}_{mk}$ \underline{est définie,} \underline{elle est} séparante (4.11) ; \underline{si} $P$ \underline{est extension d'un schéma abélien} $A$ \underline{par un groupe affine lisse connexe} $L$ \underline{satisfaisant la condition de} 4.7, \underline{alors la biextension de} $(A,A)$ \underline{dont} $\delta(\mathbb{L})$ \underline{provient} (\underline{en vertu de} 4.8) \underline{est de la forme} $\delta(\underline{M})$, \underline{où} $\underline{M}$ \underline{est un Module inversible} ample \underline{sur} $A$.

On est ramené au cas où $k$ est algébriquement clos. Par définition, il suffit alors de prouver, avec les notations de 4.12, que la biextension $\delta(\underline{M})$ de $(A,A)$ est séparante. Or il est prouvé dans la thèse de M. RAYNAUD [2, XI 1.11] que si $\mathbb{L} = f^{*}(\underline{M})$ est ample, il en est de même de $\underline{M}$ . Or il est bien connu par la théorie de A. WEIL que cela implique que $\delta(\underline{M})$ est séparante [1] [2], i.e. que l'homomorphisme de $A$ dans son schéma abélien dual est une isogénie.

\underline{Remarque} 4.14. Plaçons nous dans le cas favorable où $P$ est extension d'un schéma abélien $A$ par un groupe $L$ extension multiple de tores et de groupes $\underline{G}_{a}$, et considérons l'homomorphisme composé canonique 


\SGApage{I}{277}

\SGAvar{\begin{equation}\tag{4.14.1}
 \mathrm{NS}(A) \xrightarrow{\ \delta\ } \mathrm{Biext}^{1}(A,A;\underline{G}_{m}) \xrightarrow{\ \sim\ } \mathrm{Biext}^{1}(P,P;\underline{G}_{m}) \qquad .
\end{equation}
On sait par 4.8 que la deuxième flèche est bijective, d'autre part il est bien connu en théorie des variétés abéliennes que la première flèche induit un isomorphisme du groupe de Néron-Séveri $\mathrm{NS}(A)$ de $A$ sur le sous-groupe de $\mathrm{Biext}^{1}(A,A,\underline{G}_{m}) = \mathrm{Corr}(A,A)$ formé des (classes de) biextensions ( $\simeq$ correspondances divisorielles) \underline{symétriques}. De ceci résulte que le composé (4.14.1) induit un isomorphisme de $\mathrm{NS}(A)$ avec le groupe des biextensions \underline{symétriques} de $(P,P)$ par $\underline{G}_{m}$ (qui, en vertu de 4.7, peut d'ailleurs se décrire également comme le groupe des (classes de) correspondances divisorielles sur $(P,P)$ qui sont symétriques. Il s'impose d'appeler une correspondance divisorielle sur $(P,P)$ \underline{correspondance divisorielle ample}, et d'appeler encore \underline{biextension ample} de $(P,P)$ par $\underline{G}_{m}$ la biextension correspondante, lorsqu'elle provient d'un élément \underline{ample} du groupe $\mathrm{NS}(A)$. Cette condition étant manifestement invariante par changement du corps de base, on voit que la notion précédente garde un sens sans la condition restrictive mise plus haut sur $P$, en supposant seulement $P$ lisse et connexe.

}{\begin{equation}\tag{4.14.1}
 \mathrm{NS}(A) \xrightarrow{\ \delta\ } \mathrm{Biext}^{1}(A,A;\underline{G}_{m}) \xrightarrow{\ \sim\ } \mathrm{Biext}^{1}(P,P;\underline{G}_{m}) \qquad .
\end{equation}
On sait par 4.8 que la deuxième flèche est bijective, d'autre part il est bien connu en théorie des variétés abéliennes que la première flèche induit un isomorphisme du groupe de Néron-Séveri $\mathrm{NS}(A)$ de $A$ sur le sous-groupe de $\mathrm{Biext}^{1}(A,A,\underline{G}_{m}) = \mathrm{Corr}(A,A)$ formé des (classes de) biextensions ( $\simeq$ correspondances divisorielles) \underline{symétriques}. De ceci résulte que le composé (4.14.1) induit un isomorphisme de $\mathrm{NS}(A)$ avec le groupe des biextensions \underline{symétriques} de $(P,P)$ par $\underline{G}_{m}$ (qui, en vertu de 4.7, peut d'ailleurs se décrire également comme le groupe des (classes de) correspondances divisorielles sur $(P,P)$ qui sont symétriques). Il s'impose d'appeler une correspondance divisorielle sur $(P,P)$ \underline{correspondance divisorielle ample}, et d'appeler encore \underline{biextension ample} de $(P,P)$ par $\underline{G}_{m}$ la biextension correspondante, lorsqu'elle provient d'un élément \underline{ample} du groupe $\mathrm{NS}(A)$. Cette condition étant manifestement invariante par changement du corps de base, on voit que la notion précédente garde un sens sans la condition restrictive mise plus haut sur $P$, en supposant seulement $P$ lisse et connexe.

}{F7I-VIII-042}Notons, lorsqu'on suppose même $L$ ``résoluble sur $k$'' i.e. qu'il se dévisse en groupes $\underline{G}_{a}$ et $\underline{G}_{m}$, qu'il résulte de l'observation précédente, et du fait que l'application (4.12.2) est surjective, que l'image de

\begin{equation}\tag{4.14.2}
\mathrm{Pic}(P) \xrightarrow{\ \delta\ } \mathrm{Biext}^{1}(P,P;\underline{G}_{m})
\end{equation}
est formée des biextensions symétriques, et que son noyau est le sous-groupe $\mathrm{Pic}^{o}(P)$ de $\mathrm{Pic}(P)$ image de $\mathrm{Pic}^{o}(A)$, de sorte que, posant

\begin{equation*}
\mathrm{NS}(P) = \mathrm{Pic}(P)/\mathrm{Pic}^{o}(P) \ ( \simeq \mathrm{NS}(A)) \qquad , 
\end{equation*}

\SGApage{I}{278}

 on peut interpréter également le composé (4.14.1) comme un isomorphisme canonique

\SGAvar{\begin{equation}\tag{4.14.3}
\mathrm{NS}(P) \stackrel{\sim}{\longrightarrow} \mathrm{Biext}^1(P,P;\underline{G}_m)^{\mathrm{sym}}
\end{equation}
Notons enfin que le théorème de RAYNAUD cité pour 4.13 s'énonce en disant qu'un élément $\xi$ de $\mathrm{Pic}(P)$ est ample si et seulement si son image $\delta(\xi)$ par 4.14.1 est une classe de biextensions amples.


}{\begin{equation}\tag{4.14.3}
\mathrm{NS}(P) \stackrel{\sim}{\longrightarrow} \mathrm{Biext}^1(P,P;\underline{G}_m)^{\mathrm{sym}}
\end{equation}
Notons enfin que le théorème de RAYNAUD cité pour 4.13 s'énonce en disant qu'un élément $\xi$ de $\mathrm{Pic}(P)$ est ample si et seulement si son image $\delta(\xi)$ par 4.14.2 est une classe de biextensions amples.


}{F7I-VIII-043}\SGAsectiontoc{5. Extensions et biextensions par des schémas en groupes constants tronqués sans torsion}{5}
5. \underline{Extensions et biextensions par des schémas en groupes constants tronqués} \underline{sans torsion}

Nous donnons,dans ce numéro et le suivant, quelques résultats auxiliaires, qui seront utilisés au numéro 7 dans le cas $Z = \mathbb{Z}$ .

\underline{Proposition} 5.1. \underline{Soient} $S$ \underline{un schéma géométriquement unibranche} \underline{et irréductible de point générique} $\eta$ , $Z$ \underline{un groupe sans torsion,} \underline{d'où un schéma en groupes constant} $Z_S$ \underline{sur} $S$. \underline{Alors tout torseur} $E$ \underline{sous le schéma en groupes} $Z_S$ \underline{est trivial,} \underline{et le foncteur} $E \longmapsto E(\eta)$ \underline{de la catégorie des torseurs sous} $Z_S$ \underline{dans la catégorie des torseurs sous} $Z = Z_S(\eta)$ \underline{est une équivalence de catégories.}

\SGAvar{Un lemme connu de théorie de la descente (SGA 3 X 5.4) nous montre que tout torseur $E$ sous $G_S$ est représentable (ce résultat étant valable pour tout schéma en groupes sur $S$ qui est séparé, localement de présentation finie et localement quasi-fini sur $S$). Comme $Z_S$ est étale sur $S$, il en est de même par descente fpqc de tout torseur $E$ sous $Z_S$, qui est donc localement trivial pour la \underline{topologie étale} (et non seulement
\SGApage{I}{279}
la topologie fppf). Pour prouver que $E$ est trivial, il suffit donc de prouver qu'on a

\begin{equation*}
H^1(S_{\acute{e}t},Z_S) = e \quad ,
\end{equation*}
et de même l'assertion de pleine fidèlité de 5.1 se réduit à la relation

}{Un lemme connu de théorie de la descente (SGA 3 X 5.4) nous montre que tout torseur $E$ sous $Z_S$ est représentable (ce résultat étant valable pour tout schéma en groupes sur $S$ qui est séparé, localement de présentation finie et localement quasi-fini sur $S$). Comme $Z_S$ est étale sur $S$, il en est de même par descente fpqc de tout torseur $E$ sous $Z_S$, qui est donc localement trivial pour la \underline{topologie étale} (et non seulement
\SGApage{I}{279}
la topologie fppf). Pour prouver que $E$ est trivial, il suffit donc de prouver qu'on a

\begin{equation*}
H^1(S_{\acute{e}t},Z_S) = e \quad ,
\end{equation*}
et de même l'assertion de pleine fidélité de 5.1 se réduit à la relation

}{F7I-VIII-044}\begin{equation*}
H^o(S_{\acute{e}t},Z_S) \stackrel{\sim}{\longleftarrow} Z \quad .
\end{equation*}
Cette dernière étant triviale, prouvons la première. Pour ceci, considérons l'inclusion

\begin{equation*}
i : \eta \longrightarrow S \quad ,
\end{equation*}
on a alors, grâce au fait que $S$ est géométriquement unibranche

\begin{equation*}
i_*(Z_\eta) \stackrel{\sim}{\longleftarrow} Z_S
\end{equation*}
(SGA 4 IX 2.14.1), d'où par la suite exacte non commutative de bas degrés tenant lieu de suite spectrale de Leray pour $i : \eta \longrightarrow S$ (SGA 4 XII 3.2) :

\SGAvar{\begin{equation*}
e \longrightarrow H^1(S,i_*(Z_\eta)) \longrightarrow H^1(\eta,Z_\eta) \quad .... \quad ,
\end{equation*}
le fait que $H^1(S,Z_S) \longrightarrow H^1(\eta,Z_\eta))$ est injectif. On est donc ramené au cas où $S$ est le spectre d'un corps $k$. Si $\pi$ est le groupe de Galois d'une clôture séparable de $k$, on aura donc

\begin{equation*}
H^1(k,Z) = \mathrm{Hom}\ \mathrm{cont}(\pi,Z)/Z = e \quad ,
\end{equation*}
puisque $E$ est sans torsion. Cela achève la démonstration.

\underline{Corollaire} 5.2. \underline{Avec les notations de} 5.1, \underline{soit} $P$ \underline{un schéma lisse sur} $S$. \underline{Alors tout torseur} $E$ \underline{sous} $Z_P$ \underline{est trivial,} \underline{et le foncteur} $E \longmapsto E_\eta$ \underline{de la} \underline{catégorie des torseurs sur} $P_\eta$ \underline{de groupe} $Z_{P_\eta}$ \underline{dans la catégorie des torseurs}
\SGApage{I}{280}
\underline{sur} $P_\eta$ \underline{de groupe} $Z_{P_\eta}$ \underline{est une équivalence de catégories}.

En effet, $P$ est géométriquement unibranche (EGA IV 17.5.7), et comme chaque composante irréductible de $P$ domine $S$, $P$ étant plat sur $S$, l'ensemble des composantes irréductibles de $P$ est localement fini. Donc $P$ est un préschéma somme de préschémas irréductibles géométriquement unibranches $P_i$ . Ceci dit, 5.2 résulte aussitôt de 5.1 appliqué aux divers $P_i$ et aux $P_{i\eta}$ , compte tenu que le point générique $p_i$ de $P_i$ est aussi point générique de $P_i$ .

}{\begin{equation*}
e \longrightarrow H^1(S,i_*(Z_\eta)) \longrightarrow H^1(\eta,Z_\eta) \quad .... \quad ,
\end{equation*}
le fait que $H^1(S,Z_S) \longrightarrow H^1(\eta,Z_\eta)$ est injectif. On est donc ramené au cas où $S$ est le spectre d'un corps $k$. Si $\pi$ est le groupe de Galois d'une clôture séparable de $k$, on aura donc

\begin{equation*}
H^1(k,Z) = \mathrm{Hom}\ \mathrm{cont}(\pi,Z)/Z = e \quad ,
\end{equation*}
puisque $Z$ est sans torsion. Cela achève la démonstration.

\underline{Corollaire} 5.2. \underline{Avec les notations de} 5.1, \underline{soit} $P$ \underline{un schéma lisse sur} $S$. \underline{Alors tout torseur} $E$ \underline{sous} $Z_P$ \underline{est trivial,} \underline{et le foncteur} $E \longmapsto E_\eta$ \underline{de la} \underline{catégorie des torseurs sur} $P$ \underline{de groupe} $Z_P$ \underline{dans la catégorie des torseurs}
\SGApage{I}{280}
\underline{sur} $P_\eta$ \underline{de groupe} $Z_{P_\eta}$ \underline{est une équivalence de catégories}.

En effet, $P$ est géométriquement unibranche (EGA IV 17.5.7), et comme chaque composante irréductible de $P$ domine $S$, $P$ étant plat sur $S$, l'ensemble des composantes irréductibles de $P$ est localement fini. Donc $P$ est un préschéma somme de préschémas irréductibles géométriquement unibranches $P_i$ . Ceci dit, 5.2 résulte aussitôt de 5.1 appliqué aux divers $P_i$ et aux $P_{i\eta}$ , compte tenu que le point générique $p_i$ de $P_i$ est aussi point générique de $P_{i\eta}$ .

}{F7I-VIII-045}5.2.1. Du corollaire précédent résulte aussitôt l'énoncé analogue pour les bitorseurs sur $P$ sous $(Z_P,Z_P)$ : le foncteur $E \longmapsto E_\eta$ de la catégorie de ceux-ci, dans la catégorie des bitorseurs sur $P_\eta$ sous $(Z_{P_\eta} ,Z_{P_\eta} )$ est une équivalence de catégories. En effet, il suffit d'expliciter la structure de bitorseur de $E$ en termes de celle de torseur à droite sous $Z_P$ , muni d'un isomorphisme de $Z_P$ dans le commutant de $Z_P$ pour son opération sur $E$, en utilisant le fait que $E$ est nécessairement trivial donc que le commutant en question est isomorphe à $Z_P$ lui-même. On est donc ramené à la remarque triviale suivante pour les composantes irréductibles $X$ de $P$ : si $X$ est un schéma connexe non vide, le foncteur $Z \longmapsto Z_X$ de la catégorie des ensembles dans celle des $X$-schémas est pleinement fidèle.

\underline{Proposition} 5.3. \underline{Soient} $S$ \underline{et} $Z$ \underline{comme dans} 5.1, \underline{et} $P$ \underline{un schéma en groupes} \underline{lisse sur} $S$. \underline{Considérons les extensions} $E$ \underline{de schémas en groupes de} $P$ \underline{par} $Z_S$ . \underline{Alors le foncteur} $E \longmapsto E_\eta$ , \underline{allant de la catégorie de ces extensions} \underline{dans celle des extensions de schémas en groupes de} $P_\eta$ \underline{par} $Z_\eta$ , \underline{est une}
\SGApage{I}{281}
\underline{équivalence de catégories}. \underline{Si} $P$ \underline{et} $Z$ \underline{sont commutatifs,} \underline{ce foncteur induit} \underline{aussi une équivalence entre les catégories d'extensions commutatives} \underline{correspondantes}.

Notons d'abord que toute extension de faisceaux de groupes de $P$ par $Z_S$ (ou de $P_\eta$ par $Z_\eta$) donne naissance par 5.2 à un torseur trivial, et a fortiori représentable. Donc les catégories envisagées dans 5.3 peuvent aussi s'interpréter comme des catégories d'extensions de faisceaux de groupes. Celles-ci sont justiciables de la description donnée dans VII 1.1.6 et 1.2 en termes de bitorseurs sur $P$, $PXP$, $PXPXP$ (resp. sur $P_\eta$ , $P_\eta X P_\eta$ , $P_\eta X P_\eta X P_\eta$ ), sous les paires de groupes déduites de $(Z,Z)$ par changement de base. Appliquant 5.2 dans le cas des extensions commutatives, et 5.2.1 dans le cas général, aux schémas (lisses sur $S$) $P$, $PXP$ et $PXPXP$, on obtient la conclusion annoncée. On trouve de façon toute analogue :

\underline{Corollaire} 5.4. \underline{Soient} $S,Z$ \underline{comme dans} 5.1, \underline{et} $P,Q$ \underline{deux schémas en groupes} \underline{lisses sur} $S$. \underline{On suppose} $P,Q,Z$ \underline{commutatifs,} \underline{et on considère les biextensions} $E$ \underline{de} $(P,Q)$ \underline{par} $Z_S$ . \underline{Alors le foncteur} $E \longmapsto E_\eta$ \underline{qui va de la catégorie} \underline{de ces biextensions dans la catégorie des biextensions de} $(P_\eta,Q_\eta)$ \underline{par} $Z_\eta$ , \underline{est une équivalence de catégories.}

Les résultats précédents nous amènent donc à déterminer les catégories d'extensions ou de biextensions par un groupe constant sans torsion, dans le cas d'un schéma de base réduit au spectre d'un corps. On trouve alors ceci :


\SGApage{I}{282}

\underline{Proposition} 5.5. \underline{Soient} $k$ \underline{un corps}, $Z$ \underline{un groupe sans torsion}, $P$ \underline{un} $k$\underline{-schéma en groupes localement de type fini}, $P^{o}$ \underline{la composante neutre}, $\Phi = P/P^{o}$ \underline{le groupe des composantes connexes} (SGA 3 $\mathrm{VI}_{A}$ 4). \underline{Alors}

(i) \underline{Le foncteur naturel allant de la catégorie des schémas en groupes, extensions (resp. extensions commutatives) de} $\Phi$ \underline{par} $Z_{k}$ , \underline{dans la catégorie des schémas en groupes extensions (resp. extensions commutatives) de} $P$ \underline{par} $Z_{k}$ , \underline{est une équivalence de catégories. En particulier, si} $P$ \underline{est connexe, toute extension de} $P$ \underline{par} $Z_{k}$ \underline{est triviale, et la catégorie de ces extensions est rigide.}

\SGAvar{(ii) \underline{Supposons que} $\Phi$ \underline{soit un groupe de torsion (p.ex.} $P$ \underline{de type fini), et} $P$ \underline{et} $Z$ \underline{commutatifs. Alors la catégorie des extensions de} $P$ \underline{par} $Z_{k}$ \underline{est rigide, i.e. tout automorphisme d'un objet de cette catégorie est trivial. Si de plus il existe} $n > 0$ \underline{tel que} $n.\mathrm{id}_{\Phi} = 0$, \underline{par exemple si} $\Phi$ \underline{est fini} p.ex. $Z$ \underline{de type fini, alors le groupe des classes d'extensions commutatives de} $P$ \underline{par} $G_{k}$ \underline{est donné par un isomorphisme canonique}

}{(ii) \underline{Supposons que} $\Phi$ \underline{soit un groupe de torsion (p.ex.} $P$ \underline{de type fini), et} $P$ \underline{et} $Z$ \underline{commutatifs. Alors la catégorie des extensions de} $P$ \underline{par} $Z_{k}$ \underline{est rigide, i.e. tout automorphisme d'un objet de cette catégorie est trivial. Si de plus il existe} $n > 0$ \underline{tel que} $n.\mathrm{id}_{\Phi} = 0$, \underline{par exemple si} $\Phi$ \underline{est fini} p.ex. $P$ \underline{de type fini, alors le groupe des classes d'extensions commutatives de} $P$ \underline{par} $Z_{k}$ \underline{est donné par un isomorphisme canonique}

}{F7I-VIII-046}\begin{equation}\tag{5.5.1}
\mathrm{Ext}^{1}(P,Z_{k}) \quad \xleftarrow{\;\sim\;} \quad \mathrm{Hom}(\Phi\ ,\ (Z\otimes_{\mathbb{Z}}\mathbb{Q})/Z\ ) \qquad ,
\end{equation}
\underline{donné par l'opérateur cobord de la suite exacte des} $Ext$ \underline{déduite de la suite exacte}

\begin{equation}\tag{5.5.2}
0 \longrightarrow Z \longrightarrow Z\otimes_{\mathbb{Z}}\mathbb{Q} \longrightarrow (Z\otimes_{\mathbb{Z}}\mathbb{Q})/Z \longrightarrow 0 \qquad .
\end{equation}
Notons que le morphisme canonique

\begin{equation*}
P \longrightarrow \Phi
\end{equation*}
induit évidemment un isomorphisme sur les ensembles de composantes irréductibles, car c'est un morphisme plat à fibres irréductibles.


\SGApage{I}{283}

\SGAvar{Donc en vertu de 5.1, le foncteur induit de la catégorie des torseurs sous $Z_{\Phi}$ dans la catégorie des torseurs sous $Z_{P}$ est une équivalence de catégories. (On se rappelera qu'un schéma en groupes localement de type fini sur $k$ est géométriquement unibranche). La même conclusion s'applique aux catégories de bitorseurs, et aux morphismes

\begin{equation*}
P \times P \longrightarrow \Phi \times \Phi \quad , \qquad P \times P \times P \longrightarrow \Phi \times \Phi \times \Phi \qquad .
\end{equation*}
La description VII 1.1.6 des extensions de $P$ par $Z_{k}$ resp. de $\Phi$ par $Z_{k}$ (et VII 1.3 dans le cas d'extensions commutatives) implique alors les assertions de (i). Bornons nous maintenant au cas des extensions commutatives. Alors le groupe des automorphismes d'une extension de $P$ par $G_{k}$ est le groupe

\begin{equation*}
\mathrm{Hom}(P,G_{k}) \quad \simeq \quad \mathrm{Hom}(\Phi,G_{k}) \qquad ,
\end{equation*}
qui est évidemment nul si $\Phi$ est un groupe de torsion. On a de même, si $V = G\otimes_{\mathbb{Z}}\mathbb{Q}$ , et lorsque $\Phi$ est fini,

}{Donc en vertu de 5.1, le foncteur induit de la catégorie des torseurs sous $Z_{\Phi}$ dans la catégorie des torseurs sous $Z_{P}$ est une équivalence de catégories. (On se rappellera qu'un schéma en groupes localement de type fini sur $k$ est géométriquement unibranche). La même conclusion s'applique aux catégories de bitorseurs, et aux morphismes

\begin{equation*}
P \times P \longrightarrow \Phi \times \Phi \quad , \qquad P \times P \times P \longrightarrow \Phi \times \Phi \times \Phi \qquad .
\end{equation*}
La description VII 1.1.6 des extensions de $P$ par $Z_{k}$ resp. de $\Phi$ par $Z_{k}$ (et VII 1.2 dans le cas d'extensions commutatives) implique alors les assertions de (i). Bornons-nous maintenant au cas des extensions commutatives. Alors le groupe des automorphismes d'une extension de $P$ par $Z_{k}$ est le groupe

\begin{equation*}
\mathrm{Hom}(P,Z_{k}) \quad \simeq \quad \mathrm{Hom}(\Phi,Z_{k}) \qquad ,
\end{equation*}
qui est évidemment nul si $\Phi$ est un groupe de torsion. On a de même, si $V = Z\otimes_{\mathbb{Z}}\mathbb{Q}$ , et lorsque $\Phi$ est fini,

}{F7I-VIII-047}\begin{equation*}
\mathrm{Ext}^{i}(\Phi,V_{k}) = 0 \qquad \text{pour tout } i \in \mathbb{Z} \qquad ,
\end{equation*}
car les groupes envisagés sont des vectoriels sur $\mathbb{Q}$ annulés par l'entier $n > 0$. L'isomorphisme (5.5.1) en résulte aussitôt.

On prouve de la même façon essentiellement :

\SGAvar{\underline{Corollaire} 5.6. \underline{Soient} $k$ \underline{un corps}, $Z$ \underline{un groupe commutatif sans torsion}, $P$ \underline{et} $Q$ \underline{deux schémasen groupes commutatifs localement de type fini sur} $k$ , $\Phi = P/P^{o}$ \underline{et} $\Psi = Q/Q^{o}$ \underline{les groupes des composantes neutres.} \underline{Alors} :

}{\underline{Corollaire} 5.6. \underline{Soient} $k$ \underline{un corps}, $Z$ \underline{un groupe commutatif sans torsion}, $P$ \underline{et} $Q$ \underline{deux schémas en groupes commutatifs localement de type fini sur} $k$ , $\Phi = P/P^{o}$ \underline{et} $\Psi = Q/Q^{o}$ \underline{les groupes des composantes connexes.} \underline{Alors} :

}{F7I-VIII-048}(i) \underline{Le foncteur naturel, allant de la catégorie des biextensions}
\SGApage{I}{284}
\underline{de} $(\Phi,\Psi)$\underline{par} $Z_k$ \underline{dans celle des biextensions de} $(P,Q)$ \underline{par} $Z_k$ \underline{est une équivalence de catégories. Si en particulier} $P$ \underline{ou} $Q$ \underline{est connexe,} \underline{la catégorie des biextensions de} $(P,Q)$ \underline{par} $Z_k$ \underline{est équivalente à la catégorie ponctuelle.}

(ii) \underline{Supposons que} $\Phi$ \underline{ou} $\Psi$ \underline{soit un groupe de torsion.} \underline{Alors la catégorie des biextensions de} $(P,Q)$ \underline{par} $Z_k$ \underline{est rigide.} \underline{Si de plus il existe un entier} $n > 0$ \underline{tel que} $n.\mathrm{id}_{\Phi} = 0$ \underline{ou} $n.\mathrm{id}_{\Psi} = 0$, \underline{alors le groupe des classes de biextensions de} $(P,Q)$ \underline{par} $Z_k$ \underline{est donné par un isomorphisme canonique}

\SGAvar{\begin{equation}\tag{5.6.1}
\mathrm{Biext}^1(P,Q;\ G_k) \xleftarrow{\ \sim\ } \mathrm{Hom}(\Phi \otimes \Psi\ ,\ (Z \otimes_{\mathbb{Z}} \mathbb{Q})/Z) \quad ,
\end{equation}
\underline{donné par l'opérateur cobord de la suite exacte des} $\mathrm{Biext}^i$ $(i=0,1)$ \underline{déduite de la suite exacte} (5.5.2).

}{\begin{equation}\tag{5.6.1}
\mathrm{Biext}^1(P,Q;\ Z_k) \xleftarrow{\ \sim\ } \mathrm{Hom}(\Phi \otimes \Psi\ ,\ (Z \otimes_{\mathbb{Z}} \mathbb{Q})/Z) \quad ,
\end{equation}
\underline{donné par l'opérateur cobord de la suite exacte des} $\mathrm{Biext}^i$ $(i=0,1)$ \underline{déduite de la suite exacte} (5.5.2).

}{F7I-VIII-049}5.7. Soient $S$ un schéma, $T$ un sous-schéma fermé, $Z$ un ensemble. Appelons \underline{schéma constant sur} $S$ \underline{tronqué sur} $T$, \underline{de valeur} $Z$, et désignons par $Z_{S,T}$ , le schéma obtenu par recollement en prenant la somme de $Z$ copies $S_z$ ($z \in Z$) de $S$, qu'on recolle suivant l'ouvert $U = S-T$ complémentaire de $T$. Donc on a

\begin{equation}\tag{5.7.1}
Z_{S,T} \times_S U \simeq U \quad , \quad Z_{S,T} \times_S T \simeq Z_T \quad ,
\end{equation}
et $Z_{S,T}$ est un schéma \underline{étale sur} $S$. Pour un $S$-schéma variable $S'$, on a un isomorphisme fonctoriel en $S'$ :

\begin{equation}\tag{5.7.2}
\mathrm{Hom}_S(S',Z_{S,T}) \simeq \mathrm{Hom}_T(S'_T,Z_T) \quad ,
\end{equation}
la flèche étant définie par "restriction à $S'_T$ ",compte tenu du deuxième isomorphisme (5.7.1) ; nous laissons au lecteur le soin de vérifier que cette flèche est bijective. Désignant par 


\SGApage{I}{285}

\begin{equation*}
 i : T \longrightarrow S
\end{equation*}
l'inclusion, l'isomorphisme fonctoriel (5.7.2) peut aussi s'écrire comme un isomorphisme de faisceaux

\SGAvar{\begin{equation}\tag{5.7.3}
Z_{S,T} \simeq i_*(Z_T) \quad .
\end{equation}
On voit sur cette formule, ou directement sur la définition, que pour $Z$ variable, $Z_{S,T}$ dépend fonctoriellement de $Z$, ledit foncteur commutant aux limites projectives finies. Il transforme en particulier groupes en groupes, groupes commutatifs en groupes commutatifs, ce qui est également clair sur la forme (2.7.2) du foncteur représenté par $Z_{S,T}$ .

}{\begin{equation}\tag{5.7.3}
Z_{S,T} \simeq i_*(Z_T) \quad .
\end{equation}
On voit sur cette formule, ou directement sur la définition, que pour $Z$ variable, $Z_{S,T}$ dépend fonctoriellement de $Z$, ledit foncteur commutant aux limites projectives finies. Il transforme en particulier groupes en groupes, groupes commutatifs en groupes commutatifs, ce qui est également clair sur la forme (5.7.2) du foncteur représenté par $Z_{S,T}$ .

}{F7I-VIII-050}Signalons aussi qu'il est trivial sur la définition, ou sur (5.7.3), que la formation de $Z_{S,T}$ commute à tout changement de base $f : S' \longrightarrow S$, qui donne naissance à un isomorphisme de foncteurs en $Z$

\begin{equation}\tag{5.7.4}
Z_{S,T} \times_S S' \simeq Z_{S',T'} \quad , \quad \text{où } T' = T \times_S S' = T_{S'} \quad .
\end{equation}
5.8. Nous nous proposons d'étudier les torseurs $E$ sous un groupe de la forme $Z_{S,T}$ , où $Z$ est un groupe, qu'on suppose ici commutatif pour simplifier. On a un foncteur naturel

\SGAvar{\begin{equation}\tag{5.8.0}
E \longmapsto E_T
\end{equation}
de la catégorie des torseurs sous $Z_{S,T}$ dans la catégorie des torseurs sous $Z_T$ . Je dis que ce foncteur est une \underline{équivalence de catégories}. En effet, le fait qu'il est pleinement fidèle, i.e. que pour deux torseurs, $E,F$ sous $G_{S,T}$ , l'application $\mathrm{Hom}(E,F) \longrightarrow \mathrm{Hom}(E_T,F_T)$ est bijective, se ramène immédiatement par descente fpqc au cas où $E$ et $F$
\SGApage{I}{286}
sont triviaux, donc au cas $E = F = (G_{S,T})$, où ce n'est autre que la bijectivité de (5.7.2) pour $S'=S$. Prouvons que (5.8.0) est essentiellement surjectif. Soit $E_T$ un torseur sous $Z_T$ . On a déjà noté dans la démonstration de 5.1 que $E_T$ est représentable et localement trivial pour la topologie étale. Il suffit donc de prouver que

}{\begin{equation}\tag{5.8.0}
E \longmapsto E_T
\end{equation}
de la catégorie des torseurs sous $Z_{S,T}$ dans la catégorie des torseurs sous $Z_T$ . Je dis que ce foncteur est une \underline{équivalence de catégories}. En effet, le fait qu'il est pleinement fidèle, i.e. que pour deux torseurs, $E,F$ sous $Z_{S,T}$ , l'application $\mathrm{Hom}(E,F) \longrightarrow \mathrm{Hom}(E_T,F_T)$ est bijective, se ramène immédiatement par descente fpqc au cas où $E$ et $F$
\SGApage{I}{286}
sont triviaux, donc au cas $E = F = (Z_{S,T})$, où ce n'est autre que la bijectivité de (5.7.2) pour $S'=S$. Prouvons que (5.8.0) est essentiellement surjectif. Soit $E_T$ un torseur sous $Z_T$ . On a déjà noté dans la démonstration de 5.1 que $E_T$ est représentable et localement trivial pour la topologie étale. Il suffit donc de prouver que

}{F7I-VIII-051}\begin{equation*}
H^1(S_{\text{ét}}, Z_{S,T}) \longrightarrow H^1(T_{\text{ét}}, Z_T)
\end{equation*}
est bijectif, ce qui résulte de (5.7.3) et de la relation

\begin{equation*}
R^1 i_{\text{ét}\,*}(Z_T) = 0 \quad ,
\end{equation*}
cas particulier de SGA 4 VIII 5.6.

5.8.1. Soit maintenant $P$ un schéma en groupes commutatif sur $S$, et proposons-nous de déterminer la catégorie des extensions de faisceaux en groupes commutatifs de $P$ par $Z_{S,T}$ . Pour ceci, on utilise la description (VII 1.1.6, 1.3) en termes de torseurs sur $P$, $P \times P$, $P \times P \times P$, sous les groupes images inverses de $Z_{S,T}$ . Or ces images inverses sont de la même forme que $Z_{S,T}$ en vertu de (5.7.4), et on peut donc appliquer à la description de ces torseurs le résultat précédent, qui nous apprend qu'ils correspondent aux torseurs sur $P_T$, $P_T \times_T P_T$ , $P_T \times_T P_T \times_T P_T$ sous les groupes images inverses de $Z_T$ . On trouve ainsi :

\SGAvar{\underline{Proposition} 5.9. \underline{Soient} $S$ \underline{un schéma}, $T$ \underline{un sous-schéma fermé}, $G$ \underline{un groupe commutatif}, $G_{S,T}$ \underline{le schéma en groupes constant tronqué correspondant sur} $S$ \underline{défini dans} 5.7, $P$ \underline{un schéma en groupes commutatifs sur} $S$. \underline{Alors le foncteur de restriction à} $T$ \underline{induit une équivalence entre la catégorie}
\SGApage{I}{287}
\underline{des faisceaux en groupes commutatifs extensions de} $P$ \underline{par} $Z_{S,T}$ , \underline{et la catégorie des faisceaux en groupes commutatifs extensions de} $P_T$ \underline{par} $Z_T$ .

}{\underline{Proposition} 5.9. \underline{Soient} $S$ \underline{un schéma}, $T$ \underline{un sous-schéma fermé}, $Z$ \underline{un groupe commutatif}, $Z_{S,T}$ \underline{le schéma en groupes constant tronqué correspondant sur} $S$ \underline{défini dans} 5.7, $P$ \underline{un schéma en groupes commutatifs sur} $S$. \underline{Alors le foncteur de restriction à} $T$ \underline{induit une équivalence entre la catégorie}
\SGApage{I}{287}
\underline{des faisceaux en groupes commutatifs extensions de} $P$ \underline{par} $Z_{S,T}$ , \underline{et la catégorie des faisceaux en groupes commutatifs extensions de} $P_T$ \underline{par} $Z_T$ .

}{F7I-VIII-052}On trouve, par une démonstration essentiellement identique :

\underline{Corollaire} 5.10. \underline{Avec les notations de} 5.9, \underline{soit de plus} $Q$ \underline{un schéma en groupes commutatifs sur} $S$. \underline{Alors le foncteur restriction à} $T$ \underline{induit une équivalence entre la catégorie des biextensions de} $(P,Q)$ \underline{par} $Z_{S,T}$ , \underline{et la catégorie des biextensions de} $(P_T,Q_T)$ \underline{par} $Z_T$ .

\SGAsectiontoc{6. Extensions et biextensions par le modèle de Néron de $\protect\underline{G}_m$ .}{6}
6. \underline{Extensions et biextensions par le modèle de Néron de} $\underline{G}_m$ .

6.1. Soient $S$ un schéma, $T$ un sous-schéma fermé, $Z$ un ensemble, d'où un schéma constant tronqué $Z_{S,T}$ sur $S$ (5.7). On a un homomorphisme canonique, fonctoriel en $Z$, de schémas étales sur $S$ :

\begin{equation}\tag{6.1.1}
Z_S \longrightarrow Z_{S,T} \quad ,
\end{equation}
qui est un morphisme \underline{surjectif}, comme il résulte aussitôt de (5.7.1). De sa fonctorialité en $Z$ résulte que lorsque $Z$ est un groupe, (6.1.1) est un homomorphisme de schémas en groupes étales sur $S$. Son noyau est donc un schéma en groupes étale sur $S$, noté $Z_{S,U}$ , et on a une suite exacte canonique, fonctorielle en $Z$ :

\SGAvar{\begin{equation}\tag{6.1.2}
0 \longrightarrow Z_{S,U} \longrightarrow Z_S \longrightarrow Z_{S,T} \longrightarrow 0 \quad .
\end{equation}
On constate aussitôt, gra\^{c}e à (5.7.1), que $Z_{S,U}$ est la somme d'une famille de $S$-schémas $S_n$ , $n \in Z$, avec $S_n = S$ pour $n = 0$, et $S_n = U$ pour $n \neq 0$. De cette description de $Z_{S,U}$ résulte aisément que pour tout
\SGApage{I}{288}
Groupe $F$ sur $S$, l'application de restriction à $U$ induit une bijection :

}{\begin{equation}\tag{6.1.2}
0 \longrightarrow Z_{S,U} \longrightarrow Z_S \longrightarrow Z_{S,T} \longrightarrow 0 \quad .
\end{equation}
On constate aussitôt, grâce à (5.7.1), que $Z_{S,U}$ est la somme d'une famille de $S$-schémas $S_n$ , $n \in Z$, avec $S_n = S$ pour $n = 0$, et $S_n = U$ pour $n \neq 0$. De cette description de $Z_{S,U}$ résulte aisément que pour tout
\SGApage{I}{288}
Groupe $F$ sur $S$, l'application de restriction à $U$ induit une bijection :

}{F7I-VIII-053}\begin{equation}\tag{6.1.3}
\mathrm{Hom}_{gr}(Z_{S,U},F) \;\xrightarrow{\,\sim\,}\; \mathrm{Hom}_{gr}(Z_U,F|U) \qquad .
\end{equation}
6.1.4. Nous supposons par la suite $Z$ commutatif, et nous nous proposons d'examiner les extensions de Groupes commutatives de $Z_{S,T}$ par un Groupe commutatif donné $G$ sur $S$. Par la suite exacte (6.1.2), la catégorie de ces extensions est équivalente à la catégorie des extensions $F$ de $Z_S$ par $G$, munies d'une trivialisation de l'image inverse de cette extension par $Z_{S,U} \longrightarrow Z_S$ . En vertu de (6.1.3), une telle trivialisation revient à la donnée d'une trivialisation de l'extension $F|U$ de $Z_U$ par $G|U$ .

\SGAvar{Pour la suite, nous prendrons pour $Z$ le groupe $\mathbb{Z}$ des entiers, donc $\mathbb{Z}_S$ est le Groupe libre à un générateur, et VII 1.4 nous fournit une équivalence entre la catégorie des extensions de $\mathbb{Z}_Z$ par $G$, et la catégorie des torseurs sous $G$. On trouve donc, en appliquant le même résultat à $Z_U$ et $G|U$ :

\underline{Lemme} 6.2. \underline{Soient} $S$ \underline{un schéma}, $T$ \underline{un sous-schéma fermé}, $U = S-T$, $G$ \underline{un Groupe commutatif sur} $S$. \underline{La catégorie des Groupes commutatifs sur} $S$ \underline{extensions de} $\mathbb{Z}_{S,T}$ (\underline{défini dans} (5.7)) \underline{par} $G$ \underline{est équivalente à la catégorie des couples} $(F,s)$ \underline{formés d'un torseur} $F$ \underline{sous} $G$ \underline{et d'une section} $s$ \underline{de} $F|U$. \underline{Si} $\alpha : S \longrightarrow \mathbb{Z}_{S,T}$ \underline{est la section de} $\mathbb{Z}_{S,T}$ \underline{image de la section} $1$ \underline{de} $\mathbb{Z}_S$ \underline{sur} ${}_{\overline{G}}S$, \underline{à l'extension} $\overline{G}$ \underline{de} $\mathbb{Z}_{S,T}$ \underline{par} $G$ \underline{correspond le torseur} $F = \alpha^*(\overline{G})$ \underline{sous} $G$, \underline{muni de la section de} $F|U$ \underline{donnée par la restriction à} $U$ \underline{de la section unité de} $\overline{G}$ .


}{Pour la suite, nous prendrons pour $Z$ le groupe $\mathbb{Z}$ des entiers, donc $\mathbb{Z}_S$ est le Groupe libre à un générateur, et VII 1.4 nous fournit une équivalence entre la catégorie des extensions de $\mathbb{Z}_S$ par $G$, et la catégorie des torseurs sous $G$. On trouve donc, en appliquant le même résultat à $Z_U$ et $G|U$ :

\underline{Lemme} 6.2. \underline{Soient} $S$ \underline{un schéma}, $T$ \underline{un sous-schéma fermé}, $U = S-T$, $G$ \underline{un Groupe commutatif sur} $S$. \underline{La catégorie des Groupes commutatifs sur} $S$ \underline{extensions de} $\mathbb{Z}_{S,T}$ (\underline{défini dans} (5.7)) \underline{par} $G$ \underline{est équivalente à la catégorie des couples} $(F,s)$ \underline{formés d'un torseur} $F$ \underline{sous} $G$ \underline{et d'une section} $s$ \underline{de} $F|U$. \underline{Si} $\alpha : S \longrightarrow \mathbb{Z}_{S,T}$ \underline{est la section de} $\mathbb{Z}_{S,T}$ \underline{image de la section} $1$ \underline{de} $\mathbb{Z}_S$ \underline{sur} $S$, \underline{à l'extension} $\overline{G}$ \underline{de} $\mathbb{Z}_{S,T}$ \underline{par} $G$ \underline{correspond le torseur} $F = \alpha^*(\overline{G})$ \underline{sous} $G$, \underline{muni de la section de} $F|U$ \underline{donnée par la restriction à} $U$ \underline{de la section unité de} $\overline{G}$ .


}{F7I-VIII-054}\SGApage{I}{289}

\SGAvar{6.2.1. Si $\overline{G}$ correspond au couple $(F,s)$, l'image inverse $f^*_n(\overline{G})$ dans $\overline{G}$ de la section $f_n$ de $\mathbb{Z}_{S,T}$, définie par la section constante de valeur $n \in \mathbb{Z}$ de $\mathbb{Z}_S$ , s'identifie canoniquement, comme torseur sous $G$, à la puissance $n$.ième $F^{(n)}$ du $G$-torseur $F$ ; l'isomorphisme canonique de $G_U$-torseurs $f^*_n(\overline{G})|U \simeq G_U$ , définie par la première relation (5.7.1), est définie par la section $s^{(n)}$ de $F^{(n)}_U$ puissance $n$.ème de la section donnée $s$ de $F_U$ . De cette explicitation résulte la description suivante des sections de $\overline{G}$ sur $S$ : une section $\overline{g}$ de $\overline{G}$ sur $S$ est définie par

}{6.2.1. Si $\overline{G}$ correspond au couple $(F,s)$, l'image inverse $f^*_n(\overline{G})$ dans $\overline{G}$ de la section $f_n$ de $\mathbb{Z}_{S,T}$, définie par la section constante de valeur $n \in \mathbb{Z}$ de $\mathbb{Z}_S$ , s'identifie canoniquement, comme torseur sous $G$, à la puissance $n$.ième $F^{(n)}$ du $G$-torseur $F$ ; l'isomorphisme canonique de $G_U$-torseurs $f^*_n(\overline{G})|U \simeq G_U$ , défini par la première relation (5.7.1), est défini par la section $s^{(n)}$ de $F^{(n)}_U$ puissance $n$.ème de la section donnée $s$ de $F_U$ . De cette explicitation résulte la description suivante des sections de $\overline{G}$ sur $S$ : une section $\overline{g}$ de $\overline{G}$ sur $S$ est définie par

}{F7I-VIII-055}a) une section $\underline{n}$ de $\mathbb{Z}_{S,T}$ sur $S$, qu'on peut interpréter comme une section de $\mathbb{Z}_T$ i.e. une fonction localement constante sur $T$ à valeurs dans $\mathbb{Z}$ , ou si on préfère, une partition $T = \coprod\limits_{i \in \mathbb{Z}} T_i$ de $T$ en somme disjointe de parties ouvertes indexées par les $i \in \mathbb{Z}$ ;

b) une famille de sections

\SGAvar{\begin{multline*}
\overline{g}_i \in \Gamma\,(F^{(i)}/S'_i) \;\;,\;\; i \in \mathbb{Z},\;\; \text{où } S'_i = S - T'_i \;\;,\;\; T'_i = \bigcup_{\substack{j \in \mathbb{Z} \\ j \neq i}} T_j \;\;,
\end{multline*}
les sections $\overline{g}_i$ pour $i \in \mathbb{Z}$ étant liées par la condition de compatibilité suivante : pour $i,j \in \mathbb{Z}$, $g_i$ et $g_j$ coïncident sur $S'_i \cap S'_j = U$, quand on identifie $F^{(i)}|U$ et $F^{(j)}|U$ à $G_U$ , grâce à la section $s^{(i)}$ resp. $s^{(j)}$ de ces torseurs. En d'autres termes, il doit exister une section $h_{ij} \in \Gamma(G/U)$, évidemment uniquement déterminée, telle qu'on ait

\begin{equation*}
\overline{g}_i = h_{ij}s^{(i)} \;\;,\;\; \overline{g}_j = h_{ij}s^{(j)} \qquad .
\end{equation*}
Comme la description 3.2 d'une extension $\overline{G}$ de $\mathbb{Z}_{S,T}$ par $G$ est évidemment compatible avec tout changement de base $S' \longrightarrow S$, ce qui précède donne une description complète du faisceau $\overline{G}$ comme foncteur
\SGApage{I}{290}
en l'argument $S' \in \mathrm{Ob}\ \mathrm{Sch}_{/S}$.

6.2.2  Lorsque $T$ est un schéma somme de schémas connexes non vides $T_i$ , $i \in I$, alors la description 6.2.1 des sections de $\overline{G}$ sur $S$ peut se simplifier : une telle section est définie par une famille d'entiers $(n_i)_{i \in I}$ indexée par le même ensemble $I$, une famille des sections $\overline{g}_i \in \Gamma\,(F^{(n_i)}/S_i')$, pour $i \in I$, où $S_i' = S-T_i'$ , $T_i' = \coprod_{\substack{j \in I \\ j \neq i}} T_j$, les $\overline{g}_i$ satisfaisant à la condition de compatibilité suivante : pour tout couple d'léments $i,j \in I$, les restrictions de $\overline{g}_i$ et $\overline{g}_j$ à $S_i' \cap S_j' = U$ s'identifient, modulo l'identification de $F^{(n_i)}|U$ et de $F^{(n_j)}|U$, déduites des sections $s^{(n_i)}$ resp. $s^{(n_j)}$ de ces torseurs.

}{\begin{multline*}
\overline{g}_i \in \Gamma\,(F^{(i)}/S'_i) \;\;,\;\; i \in \mathbb{Z},\;\; \text{où } S'_i = S - T'_i \;\;,\;\; T'_i = \bigcup_{\substack{j \in \mathbb{Z} \\ j \neq i}} T_j \;\;,
\end{multline*}
les sections $\overline{g}_i$ pour $i \in \mathbb{Z}$ étant liées par la condition de compatibilité suivante : pour $i,j \in \mathbb{Z}$, $\overline{g}_i$ et $\overline{g}_j$ coïncident sur $S'_i \cap S'_j = U$, quand on identifie $F^{(i)}|U$ et $F^{(j)}|U$ à $G_U$ , grâce à la section $s^{(i)}$ resp. $s^{(j)}$ de ces torseurs. En d'autres termes, il doit exister une section $h_{ij} \in \Gamma(G/U)$, évidemment uniquement déterminée, telle qu'on ait

\begin{equation*}
\overline{g}_i = h_{ij}s^{(i)} \;\;,\;\; \overline{g}_j = h_{ij}s^{(j)} \qquad .
\end{equation*}
Comme la description 6.2 d'une extension $\overline{G}$ de $\mathbb{Z}_{S,T}$ par $G$ est évidemment compatible avec tout changement de base $S' \longrightarrow S$, ce qui précède donne une description complète du faisceau $\overline{G}$ comme foncteur
\SGApage{I}{290}
en l'argument $S' \in \mathrm{Ob}\ \mathrm{Sch}_{/S}$.

6.2.2  Lorsque $T$ est un schéma somme de schémas connexes non vides $T_i$ , $i \in I$, alors la description 6.2.1 des sections de $\overline{G}$ sur $S$ peut se simplifier : une telle section est définie par une famille d'entiers $(n_i)_{i \in I}$ indexée par le même ensemble $I$, une famille des sections $\overline{g}_i \in \Gamma\,(F^{(n_i)}/S_i')$, pour $i \in I$, où $S_i' = S-T_i'$ , $T_i' = \coprod_{\substack{j \in I \\ j \neq i}} T_j$, les $\overline{g}_i$ satisfaisant à la condition de compatibilité suivante : pour tout couple d'éléments $i,j \in I$, les restrictions de $\overline{g}_i$ et $\overline{g}_j$ à $S_i' \cap S_j' = U$ s'identifient, modulo l'identification de $F^{(n_i)}|U$ et de $F^{(n_j)}|U$, déduites des sections $s^{(n_i)}$ resp. $s^{(n_j)}$ de ces torseurs.

}{F7I-VIII-056}6.2.3.  Lorsqu'on suppose enfin $T$ connexe non vide, la donnée d'une section $\overline{g}$ de $\overline{G}$ sur $S$ revient à la donnée d'un entier $n$, et d'une section $\overline{g}$ de $F^{(n)}$ sur $S$. Il n'y a plus de condition de compatibilité à imposer, de sorte que la section $s$ de $F|U$ est absente de la description du groupe des sections de $\overline{G}$ sur $S$. (Mais $s$ apparaît de nouveau, bien entendu, dès qu'on se propose de donner une description des points de $\overline{G}$ à valeurs dans un $S$-schéma général $S'$.)

6.2.4.  Comme $\overline{G}$ est réunion des ``ouverts'' $f^*_n(\overline{G})$ ($n \in \mathbb{Z}$) décrits dans 6.2.1, on en conclut que $\overline{G}$ est représentable si et seulement si les torseurs $F^{(n)}$ le sont. Il suffit par exemple, pour ceci, que $G$ soit représentable et soit un schéma en groupes affine sur $S$, grâce à la théorie de la descente SGA 1 VIII 2.1.


\SGApage{I}{291}

\SGAvar{6.3.  Dans la suite, nous supposerons que $G = \underline{G}_{mS}$ , de sorte que toute extension $\underline{G}$ de $\mathbb{Z}_{S,T}$ par $G = \underline{G}_{mS}$ sera bien représentable (6.2.4), et bien entendu, sera lisse sur $S$ comme extension d'un schéma en groupes étale $\mathbb{Z}_{S,T}$ par un schéma en groupes lisse $\underline{G}_{mS}$ . De plus, la donnée d'un torseur $F$ sous $\underline{G}_m$ équivaut à la donnée d'un Module inversible $\mathbb{L}$ sur $S$ (en associant à ce dernier le torseur des sections inversibles de $\mathbb{L}$ ). Donc $\overline{G}$ est décrit par un couple $(\mathbb{L},s)$, où $\mathbb{L}$ est un Module inversible sur $S$, et $s$ une trivialisation de $\mathbb{L}$ sur $U=S-T$, i.e. une section inversible de $\mathbb{L}|U$. Lorsque $U$ est schématiquement dense dans $S$, et que $S$ est localement noethérien, ou réduit et l'ensemble de ses composantes irréductibles est localement fini, la donnée d'un tel couple équivaut à celle d'un \underline{diviseur} $D$ \underline{sur} $S$ \underline{tel que}

}{6.3.  Dans la suite, nous supposerons que $G = \underline{G}_{mS}$ , de sorte que toute extension $\overline{G}$ de $\mathbb{Z}_{S,T}$ par $G = \underline{G}_{mS}$ sera bien représentable (6.2.4), et bien entendu, sera lisse sur $S$ comme extension d'un schéma en groupes étale $\mathbb{Z}_{S,T}$ par un schéma en groupes lisse $\underline{G}_{mS}$ . De plus, la donnée d'un torseur $F$ sous $\underline{G}_m$ équivaut à la donnée d'un Module inversible $\mathbb{L}$ sur $S$ (en associant à ce dernier le torseur des sections inversibles de $\mathbb{L}$ ). Donc $\overline{G}$ est décrit par un couple $(\mathbb{L},s)$, où $\mathbb{L}$ est un Module inversible sur $S$, et $s$ une trivialisation de $\mathbb{L}$ sur $U=S-T$, i.e. une section inversible de $\mathbb{L}|U$. Lorsque $U$ est schématiquement dense dans $S$, et que $S$ est localement noethérien, ou réduit et l'ensemble de ses composantes irréductibles est localement fini, la donnée d'un tel couple équivaut à celle d'un \underline{diviseur} $D$ \underline{sur} $S$ \underline{tel que}

}{F7I-VIII-057}\begin{equation}\tag{6.3.1}
\mathrm{supp}\ D \subset T \quad ,
\end{equation}
en faisant correspondre à un tel diviseur le couple

\SGAvar{\begin{equation}\tag{6.3.2}
(\underline{O}_S(-D),s_D|U)
\end{equation}
du Module inversible $\underline{O}_S(-D)$ associé à $-D$ (EGA IV 21.2.8.1), et de la restriction à $U$ de la section rationnelle canonique $s_{-D}$ de celui-ci (EGA IV 21.2.11 et 20.2.11 (ii)). En tous cas, sans hypothèse sur $S$, la donnée d'un diviseur $D$ sur $S$ satisfaisant (6.3.1) définit un couple $(\mathbb{L},s)$ par (6.3.2), donc une extension $G$ de $\mathbb{Z}_{S,T}$ par $\underline{G}_m$ , que nous appelerons le \underline{modèle de Néron-Raynaud de} $\underline{G}_m$ , associé au diviseur $D$ et au fermé $T$ de $S$ satisfaisant (6.3.1). Si $T = \mathrm{supp}\ D$, on omet $T$ dans la terminologie.


}{\begin{equation}\tag{6.3.2}
(\underline{O}_S(-D),s_{-D}|U)
\end{equation}
du Module inversible $\underline{O}_S(-D)$ associé à $-D$ (EGA IV 21.2.8.1), et de la restriction à $U$ de la section rationnelle canonique $s_{-D}$ de celui-ci (EGA IV 21.2.11 et 20.2.11 (ii)). En tous cas, sans hypothèse sur $S$, la donnée d'un diviseur $D$ sur $S$ satisfaisant (6.3.1) définit un couple $(\mathbb{L},s)$ par (6.3.2), donc une extension $\overline{G}$ de $\mathbb{Z}_{S,T}$ par $\underline{G}_m$ , que nous appellerons le \underline{modèle de Néron-Raynaud de} $\underline{G}_m$ , associé au diviseur $D$ et au fermé $T$ de $S$ satisfaisant (6.3.1). Si $T = \mathrm{supp}\ D$, on omet $T$ dans la terminologie.


}{F7I-VIII-058}\SGApage{I}{292}

Il est immédiat que la formation de cette extension $\overline{G}$ commute à tout changement de base $f : S' \longrightarrow S$ tel que le diviseur image inverse $f^*(D)$ soit défini (EGA IV 21.4.2), en particulier à tout changement de base \underline{plat}.

\SGAvar{6.3.3.  On voit aussitôt, grâce à la description générale 6.2.1 des sections de $\overline{G}$, que celles-ci peuvent s'identifier aux sections $f$ du faisceau $\mathcal{M}^*_S$ des sections méromorphes régulières sur $S$ (EGA IV 20.1.8) dont le diviseur $\mathrm{div}(f)$ est localement (au sens de la topologie de Zariski) un multiple entier de $D$ . Si $V$ est un ouvert de $S$ tel que la restriction de $\mathrm{div}\, f$ à $V$ est égal à celle de $nD$, où $n \in \mathbb{Z}$ , alors la section de $\mathbb{Z}_{S,T}$ définie par $f$ n'est autre que l'image de la section constante de valeur $n$ de $\mathbb{Z}_S$ dans (6.1.2). La description précédente du groupe des sections de $\overline{G}$ sur $S$ reste valable au-dessus de tout ouvert de $S$, plus généralement pour le groupe de points de $\overline{G}$ à valeurs dans n'importe quel $S$-schéma plat $S'$. Comme $\overline{G}$ lui-même est plat (et même lisse) sur $S$, cette description redonne donc une caractérisation complète du modèle de Néron-Raynaud de $\underline{G}_m$ défini par $D$ et $T \supset \mathrm{supp}\, D$.

}{6.3.3.  On voit aussitôt, grâce à la description générale 6.2.1 des sections de $\overline{G}$, que celles-ci peuvent s'identifier aux sections $f$ du faisceau $\mathcal{M}^*_S$ des sections méromorphes régulières sur $S$ (EGA IV 20.1.8) dont le diviseur $\mathrm{div}(f)$ est localement (au sens de la topologie de Zariski) un multiple entier de $D$ . Si $V$ est un ouvert de $S$ tel que la restriction de $\mathrm{div}\, f$ à $V$ est égale à celle de $nD$, où $n \in \mathbb{Z}$ , alors la section de $\mathbb{Z}_{S,T}$ définie par $f$ n'est autre que l'image de la section constante de valeur $n$ de $\mathbb{Z}_S$ dans (6.1.2). La description précédente du groupe des sections de $\overline{G}$ sur $S$ reste valable au-dessus de tout ouvert de $S$, plus généralement pour le groupe de points de $\overline{G}$ à valeurs dans n'importe quel $S$-schéma plat $S'$. Comme $\overline{G}$ lui-même est plat (et même lisse) sur $S$, cette description redonne donc une caractérisation complète du modèle de Néron-Raynaud de $\underline{G}_m$ défini par $D$ et $T \supset \mathrm{supp}\, D$.

}{F7I-VIII-059}6.4.  Supposons maintenant que $S$ soit \underline{normal}, et l'ensemble de ses composantes irréductibles localement fini (donc $S$ est un schéma somme de schémas normaux intègres), et supposons que le diviseur $D$ soit positif, à multiplicités 1 (EGA IV 21.6), que $T=\mathrm{supp}\, D$, et que $T$ soit localement irréductible i.e. que ses composantes irréductibles soient disjointes. Alors pour toute fonction rationnelle $f$ sur $S$, régulière sur $U = S-T$, $\mathrm{div}\, f$ est localement un multiple de $D$, donc on a dans ce cas un isomor
\SGApage{I}{293}
phisme de groupes canonique :

\begin{equation}\tag{6.4.1}
\Gamma(S,\overline{G}) \overset{\sim}{\longrightarrow} \Gamma(U,\underline{G}_{mU}) \quad .
\end{equation}
Supposons même $T$ géométriquement unibranche et localement intègre, hypothèse qui est stable par changement de base lisse (EGA IV 17.5.7). Alors, toutes les hypothèses étant stables par changement de base, on déduit de (6.4.1) que \underline{pour} $S'$ \underline{lisse sur} $S$, \underline{on a un isomorphisme de groupes canonique}, \underline{fonctoriel en} $S'$ :

\begin{equation}\tag{6.4.2}
\mathrm{Hom}_S(S', \overline{G}) \overset{\sim}{\longrightarrow} \mathrm{Hom}_U(S'_U, \underline{G}_{mU}) \simeq \Gamma(S'_U , \underline{O}^*_{S'_U}) \quad .
\end{equation}
Comme $\overline{G}$ est lui-même lisse sur $S$, cet isomorphisme caractérise le schéma en groupes $\overline{G}$ sur $S$ à isomorphisme unique près. D'autre part, il montre que $\overline{G}$ joue le rôle d'un ``modèle de Néron'' pour le schéma en groupes $\underline{G}_{mU}$ sur l'ouvert $U$ de $S$, ce qui (joint au fait qu'il a été introduit et utilisé pour la première fois par M. RAYNAUD) justifie la terminologie adoptée dans 6.3 pour désigner $\overline{G}$ .

\underline{Théorème} 6.5. \underline{Soient} $S$ \underline{un schéma normal localement intègre}, $D$ \underline{un diviseur positif sur} $S$, \underline{sans composantes multiples}, \underline{tel que le support de} $D$ \underline{soit géométriquement unibranche et localement intègre}. \underline{Soit} $\overline{G}$ \underline{le modèle de Néron-Raynaud de} $\underline{G}_m$ \underline{sur} $S$ \underline{relatif à} $D$ (6.3), \underline{et considérons la catégorie des torseurs} $E$ \underline{sur} $S$ \underline{sous le schéma en groupes} $\overline{G}$ . \underline{Le foncteur restriction}

\SGAvar{\begin{equation*}
E \longmapsto E_U = E|U \quad ,
\end{equation*}
\underline{qui va de cette catégorie dans la catégorie des torseurs sur} $U$ \underline{sous le schéma en groupes} $\underline{G}_{mU}$ , \underline{est pleinement fidèle}. \underline{Si} $S$ \underline{est un schéma régulier} (\underline{plus généralement}, \underline{si ces anneaux locaux sont factoriels}), \underline{alors ce foncteur est même une équivalence de catégories}.


}{\begin{equation*}
E \longmapsto E_U = E|U \quad ,
\end{equation*}
\underline{qui va de cette catégorie dans la catégorie des torseurs sur} $U$ \underline{sous le schéma en groupes} $\underline{G}_{mU}$ , \underline{est pleinement fidèle}. \underline{Si} $S$ \underline{est un schéma régulier} (\underline{plus généralement}, \underline{si ses anneaux locaux sont factoriels}), \underline{alors ce foncteur est même une équivalence de catégories}.


}{F7I-VIII-060}\SGApage{I}{294}

Montrons que le foncteur envisagé est pleinement fidèle, i.e. que pour deux torseurs $E,F$ sous $\overline{G}$, l'application de restriction

\begin{equation}\tag{*}
\mathrm{Hom}(E,F) \longrightarrow \mathrm{Hom}(E_U,F_U)
\end{equation}
est bijective, où les Hom désignent les ensembles d'homomorphismes de \underline{torseurs}. Notons d'abord que cette question est manifestement \underline{locale} \underline{sur} $S$ \underline{pour la topologie de Zariski} : si $S$ est réunion d'ouverts $S_i$ tels que pour tout ouvert $S'$ de $S$ contenu dans un $S_i$ (donc en particulier pour $S'$ de la forme $S_i$ ou $S_i \cap S_j$), l'application

\begin{equation*}
\mathrm{Hom}(E_{S'},F_{S'}) \longrightarrow \mathrm{Hom}(E_{S' \cap U},F_{S' \cap U})
\end{equation*}
est bijective, alors l'application (*) est bijective. D'autre part, il résulte aussitôt de 5.8 appliqué à $S,T,\mathbb{Z}$ , et de 5.1 appliqué à $T,\mathbb{Z}$ , que tout torseur sous $\overline{G}$ est localement trivial pour la topologie de Zariski. Ces deux observations nous ramènent à prouver la bijectivité de l'application (*) dans le cas où $E=F=\overline{G}$ . Mais alors cette application n'est autre que (6.4.1), dont la bijectivité a déjà été établie.

\SGAvar{Montrons maintenant que le foncteur $E \longmapsto E_U$ envisagé est essentiellement surjectif lorsque les anneaux locaux de $X$ sont factoriels. Comme $\underline{G}_{m\,S}$ s'identifie à un sous-schéma en groupes ouvert de $\overline{G}$, il suffit de prouver que tout torseur sous $\underline{G}_{m\,U}$ est isomorphe à la restriction à $U$ d'un torseur sous $\underline{G}_{m\,S}$ . Il revient au même de dire que tout Module inversible sur $U$ est isomorphe à la restriction d'un Module inversible sur $S$, ce qui résulte en effet de l'hypothèse faite sur $S$ (EGA IV 21.6.11, où on se borne au cas localement noethérien, qui est le seul cas qui nous servira par la suite). Cela prouve 6.5.

}{Montrons maintenant que le foncteur $E \longmapsto E_U$ envisagé est essentiellement surjectif lorsque les anneaux locaux de $S$ sont factoriels. Comme $\underline{G}_{m\,S}$ s'identifie à un sous-schéma en groupes ouvert de $\overline{G}$, il suffit de prouver que tout torseur sous $\underline{G}_{m\,U}$ est isomorphe à la restriction à $U$ d'un torseur sous $\underline{G}_{m\,S}$ . Il revient au même de dire que tout Module inversible sur $U$ est isomorphe à la restriction d'un Module inversible sur $S$, ce qui résulte en effet de l'hypothèse faite sur $S$ (EGA IV 21.6.11, où on se borne au cas localement noethérien, qui est le seul cas qui nous servira par la suite). Cela prouve 6.5.

}{F7I-VIII-061}\SGApage{I}{295}

\SGAvar{6.5.1. Plaçons nous maintenant dans le cas où on suppose de plus $S$ localement noethérien et \underline{régulier}. Comme $S$ est alors localement factoriel, la donnée de $D$ équivaut à celle d'un sous-schéma fermé (que nous notons encore $D$) qui soit purement 1-codimensionnel, réduit et géométriquement unibranche. Notons aussi que les hypothèses sur $S$ et $D$ sont stables par tout changement de base lisse $S' \longrightarrow S$, de sorte que la conclusion de 6.5 s'appliquera à $S'$ muni du sous-schéma $D' = D \times_S S'$ de $S'$. Ceci dit, la description VII 1.1.6 des Groupes extensions nous donne aussitôt la conséquence suivante :

}{6.5.1. Plaçons-nous maintenant dans le cas où on suppose de plus $S$ localement noethérien et \underline{régulier}. Comme $S$ est alors localement factoriel, la donnée de $D$ équivaut à celle d'un sous-schéma fermé (que nous notons encore $D$) qui soit purement 1-codimensionnel, réduit et géométriquement unibranche. Notons aussi que les hypothèses sur $S$ et $D$ sont stables par tout changement de base lisse $S' \longrightarrow S$, de sorte que la conclusion de 6.5 s'appliquera à $S'$ muni du sous-schéma $D' = D \times_S S'$ de $S'$. Ceci dit, la description VII 1.1.6 des Groupes extensions nous donne aussitôt la conséquence suivante :

}{F7I-VIII-062}\underline{Corollaire} 6.6. \underline{Soient} $S$ \underline{un schéma localement noethérien et régulier}, $D$ \underline{un sous-schéma fermé réduit géométriquement unibranche et purement} 1-\underline{codimensionnel} (\underline{définissant donc un diviseur sur} $S$), $\overline{G}$ \underline{le modèle de Néron-Raynaud de} $\underline{G}_m$ \underline{relatif à} $D$, $P$ \underline{un schéma en groupes commutatifs lisse sur} $S$. \underline{Considérons la catégorie des extensions commutatives de faisceaux en Groupes de} $P$ \underline{par} $\overline{G}$. \underline{Alors le foncteur restriction à} $U$, $E \longrightarrow E|U$, \underline{induit une équivalence de cette catégorie avec la catégorie des extensions de} $P_U$ \underline{par} $\underline{G}_{m\,U}$ .

6.6.1. On obtient par la même méthode l'énoncé analogue, relatif au cas d'extensions pas nécessairement commutatives ($P$ n'étant plus supposé commutatif). Nous laissons au lecteur le détail de la démonstration, légèrement plus compliquée du fait qu'il faut considérer ici des bitorseurs au lieu de torseurs.

Par essentiellement la même démonstration que 6.6, on obtient de même :


\SGApage{I}{296}

\underline{Corollaire} 6.7. \underline{Soient} $S,D,\overline{G}$ \underline{comme dans} 6.6, \underline{et} $P,Q$ \underline{deux schémas en groupes commutatifs et lisses sur} $S$. \underline{On considère la catégorie des biextensions de} $(P,Q)$ \underline{par} $\overline{G}$ . \underline{Alors le foncteur restriction à} $U$, $E \longmapsto E|U$, \underline{qui va de cette catégorie dans la catégorie des biextensions de} $(P_U,Q_U)$ \underline{par} $\underline{G}_{m\,U}$, \underline{est une équivalence de catégories.}

\underline{Remarques} 6.8.\quad a)\quad Nous appliquerons 6.6 et 6.7 lorsque $S$ est un trait, et $D$ le diviseur défini par l'inclusion du point fermé $s$ de $S$ dans $S$. Dans ce cas, $U$ est le spectre du corps des fractions $K$ de $S$, et $\overline{G}$ est le modèle de Néron habituel du groupe multiplicatif sur $K$, introduit par Raynaud. On notera que, même en se bornant à ce cas, la démonstration de 6.6 et de 6.7 utilise 6.5 sur des bases plus générales que des traits (telles que $P$ et ses puissances cartésiennes ...). Néanmoins, il suffit dans ce cas de connaître l'\underline{existence} du modèle de Néron $\overline{G}$ dans le cas particulier envisagé, pour en conclure formellement 6.5 pour des schémas de base lisses sur $S$, donc pour obtenir 6.6 et 6.7. Pour nous, l'introduction de $\overline{G}$ sera surtout un intermédiaire commode pour étudier au numéro suivant le problème des prolongements d'extensions et de biextensions par le groupe multiplicatif.

b) On peut définir un modèle de Néron-Raynaud $\overline{G}$ de $\underline{G}_{m\,U}$ , satisfaisant à la condition fondamentale (6.4.2) pour tout $S'$ lisse sur $S$, sous des conditions nettement plus générales que celles envisagées plus haut, $U$ étant simplement un ouvert schématiquement dense d'un schéma localement noethérien \underline{normal} $S$. Dans ce cas, $\overline{G}$ sera un \underline{faisceau} fppf, pas nécessairement représentable. Désignant par $\underline{D}$ le faisceau, sur le petit site étale de $S$, des diviseurs sur $S$ de support contenu
\SGApage{I}{297}
dans $T=S-U$, et (par abus de langage) par la même lettre l'extension canonique de $\underline{D}$ en un faisceau sur $(\mathrm{Sch})_{/S}$ (dont les sections sur le schéma relatif $f:S' \longrightarrow S$ sont les sections de $f^{*}_{\text{ét}}(\underline{D})$ sur $S'$), le modèle cherché est une extension de Groupes

\begin{equation}\tag{6.8.1}
0 \longrightarrow \underline{G}_{mS} \longrightarrow \overline{G} \longrightarrow \underline{D} \longrightarrow 0 \quad .
\end{equation}
La construction d'une extension canonique (6.8.1) (sans hypothèse de normalité sur $S$) ne présente pas de difficultés, en utilisant la description générale VII 1.1.6, et on obtient en même temps (6.4.2) pour $S'$ \underline{étale} sur $S$. Pour avoir la même relation pour tout $S'$ lisse sur $S$, il revient au même de vérifier que la formation du faisceau $\underline{D}$ commute à tout changement de base lisse $f: S' \longrightarrow S$, ce qui pour $S$ normal résulte de EGA $\mathrm{Err}_{\mathrm{IV}}$ 53, 21.4.9.


6.9.\quad ``\underline{Autocritique}'' (observations de P. DELIGNE)(*).

Les n\textsuperscript{o}s 5 et 6, sont vraiment peu inspirants. On a périodiquement l'impression qu'on refait ``à la main'' le formalisme $i_*,i^*$, sans voir ce qui est géométrique ou ``general non-sense''.

Je propose de dégager les énoncés ``géométriques'' a)b)c) :

\SGAvar{a) (cf. 6.8 b)\quad Soit $j:V \hookrightarrow S$ un ouvert schématiquement dense d'un schéma normal noethérien $S$. Le faisceau étale $\underline{G}_{mS}$ s'identifie alors à un sous-faisceau du faisceau étale $j_{*}\underline{G}_{mU}$ . Soit $\underline{D}$ le faisceau étale quotient

\begin{equation*}
0 \longrightarrow \underline{G}_{mS} \longrightarrow j_{*} \, \underline{G}_{mS} \longrightarrow \underline{D} \longrightarrow 0 \quad .
\end{equation*}
(*)\quad Le rédacteur du présent exposé déclarant forfait, le lecteur est invité à récrire, en cas de besoin, les n\textsuperscript{o}s 5 et 6 suivant les directives jointes.


}{a) (cf. 6.8 b)\quad Soit $j:U \hookrightarrow S$ un ouvert schématiquement dense d'un schéma normal noethérien $S$. Le faisceau étale $\underline{G}_{mS}$ s'identifie alors à un sous-faisceau du faisceau étale $j_{*}\underline{G}_{mU}$ . Soit $\underline{D}$ le faisceau étale quotient

\begin{equation*}
0 \longrightarrow \underline{G}_{mS} \longrightarrow j_{*} \, \underline{G}_{mU} \longrightarrow \underline{D} \longrightarrow 0 \quad .
\end{equation*}
(*)\quad Le rédacteur du présent exposé déclarant forfait, le lecteur est invité à récrire, en cas de besoin, les n\textsuperscript{o}s 5 et 6 suivant les directives jointes.


}{F7I-VIII-063}\SGApage{I}{298}

\underline{Alors} la formation de $\underline{D}$ commute à tout changement de base lisse (EGA Err$_{\mathrm{IV}}$ 53, 21.4.9) ; \underline{de plus}, $R^1 j_* \underline{G}_m = 0$ (Th. 90) si les anneaux strictement locaux de $S$ sont factoriels.

b) Si $U = S-Y$, avec $Y$ un diviseur géométriquement unibranche, et si $i$ est l'inclusion de $Y$ dans $S$, alors

\SGAvar{\begin{equation*}
\underline{D} = i_* \mathbb{Z}_Y \ .
\end{equation*}
c) Si $Y$ est sous-schéma fermé d'un schéma $X$, et $E$ un ensemble, alors le faisceau étale $i_* \underline{E}_Y$ est représentable par un schéma étale sur $X$, de formation compatible à tout changement de base.

}{\begin{equation*}
\underline{D} = i_* \mathbb{Z}_Y \ .
\end{equation*}

c) Si $Y$ est sous-schéma fermé d'un schéma $X$, si $i:Y\hookrightarrow X$ désigne l'immersion fermée, et si $E$ est un ensemble, alors le faisceau étale $i_* \underline{E}_Y$ est représentable par un schéma étale sur $X$, de formation compatible à tout changement de base.

}{F7I-VIII-064}\underline{Corollaire}. \underline{Soient} $S$ \underline{un schéma noethérien normal}, $Y$ \underline{un diviseur géométriquement unibranche et} $U = S-Y$ :

\begin{center}
\begin{tikzcd}[column sep=huge]
U \arrow[r, hook, "j"'] \& S \& Y \arrow[l, "i"]
\end{tikzcd}\quad .
\end{center}









\underline{Il existe un unique schéma en groupes lisse} $\mathcal{G}_m$ \underline{sur} $S$, \underline{tel qu'après tout changement de base lisse}, \underline{on ait} $\mathcal{G}_m \mid S_{\text{\'et}} = j_*(\underline{G}_m)_U$ .

Les exigences précédentes déterminent en effet $\mathcal{G}_m$ en tant que faisceau (pour la topologie étale) sur la catégorie des schémas lisses sur $S$. Par a), b), la suite

\begin{equation*}
0 \longrightarrow \underline{G}_m \longrightarrow \mathcal{G}_m \xrightarrow{\ q\ } i_* \mathbb{Z}_Y \longrightarrow 0
\end{equation*}
est exacte. Par c), $i_* \mathbb{Z}_Y$ est représentable. Enfin le morphisme $q$ est représentable.

L'usage du``petit site lisse''précédent devrait permettre de déduire formellement les énoncés du n° 6 de a).


\SGApage{I}{299}


\SGAsectiontoc{7. Prolongements canoniques d'extensions et de biextensions par $\protect\underline{G}_m$}{7}
7. \underline{Prolongements canoniques d'extensions et de biextensions par} $\underline{G}_m$

\SGAvar{7.0. Nous récoltons maintenant le fruit de nos préparatifs des n°s 5 et 6. Plaçons nous d'abord dans le cas d'un schéma de base $S$ qui soit un \underline{trait}, i.e. le spectre d'un anneau de valuation discrète, et soit $\eta$ son point générique, $s$ sont point fermé, $k = k(s)$. Soit $P$ un schéma en groupes commutatifs lisse sur $S$, on posera $P_o = P_s$, fibre spéciale de $P$. Dans le présent numéro, nous ne considérons que des groupes commutatifs, et en particulier les extensions de groupes envisagées sont commutatives.

}{7.0. Nous récoltons maintenant le fruit de nos préparatifs des n°s 5 et 6. Plaçons nous d'abord dans le cas d'un schéma de base $S$ qui soit un \underline{trait}, i.e. le spectre d'un anneau de valuation discrète, et soit $\eta$ son point générique, $s$ son point fermé, $k = k(s)$. Soit $P$ un schéma en groupes commutatifs lisse sur $S$, on posera $P_o = P_s$, fibre spéciale de $P$. Dans le présent numéro, nous ne considérons que des groupes commutatifs, et en particulier les extensions de groupes envisagées sont commutatives.

}{F7I-VIII-065}\underline{Théorème} 7.1. \underline{Les notations sont celles rappelées dans} 7.0.

a) \underline{Supposons que le groupe des composantes connexes de la fibre spéciale de} $P$,

\begin{equation*}
\Phi = P_o/P_o^{\,o}
\end{equation*}
\underline{soit un groupe de torsion}, - \underline{ce qui est le cas en particulier si} $P$ \underline{est de type fini sur} $S$ ; \underline{alors le foncteur} $E \longmapsto E_\eta$ \underline{de la catégorie des extensions de} $P$ \underline{par} $\underline{G}_{mS}$ \underline{dans la catégorie des extensions de} $P_\eta$ \underline{par} $\underline{G}_{m\eta}$ \underline{est pleinement fidèle}. \underline{Supposons qu'il existe un entier} $n > 0$ \underline{tel que} $n\Phi = 0$, \underline{alors si} $E_\eta$ \underline{est une extension de} $P_\eta$ \underline{par} $\underline{G}_{m\eta}$ , \underline{il lui est associé un homomorphisme canonique}

\begin{equation}\tag{7.1.1}
d(E_\eta) : \Phi \longrightarrow (\mathbb{Q}/\mathbb{Z})_k \qquad ,
\end{equation}
\underline{dont l'annulation est nécessaire et suffisante pour que l'extension} $E_\eta$ \underline{soit isomorphe à la fibre générique d'une extension de} $P$ \underline{par} $\underline{G}_{mS}$ . En \underline{particulier}, \underline{si} $P_o$ \underline{est connexe}, \underline{le foncteur précédent} $E \longmapsto E_\eta$ \underline{est une équivalence de catégories}.


\SGApage{I}{300}

b)  \underline{Soient} $P$ \underline{et} $Q$ \underline{deux schémas en groupes lisses sur} $S$, \underline{posons}

\SGAvar{\begin{equation*}
\Phi = P_o/P_o^{\,o} \qquad , \qquad \Psi = Q_o/Q_o^{\,o} \qquad ,
\end{equation*}
\underline{et supposons que} $\Phi$ \underline{ou} $\Psi$ \underline{soit un groupe de torsion} (\underline{par exemple} $P$ \underline{ou} $Q$ \underline{de type fini sur} $S$) ; \underline{alors le foncteur} $E \longmapsto E_\eta$ \underline{dans la catégorie des biextensions de} $(P,Q)$ \underline{par} $\underline{G}_{mS}$ \underline{dans la catégorie des biextensions de} $(P_\eta,Q_\eta)$ \underline{par} $\underline{G}_{m\eta}$ \underline{est pleinement fidèle.} \underline{Supposons qu'il existe un entier} $n > 0$ \underline{tel que} $n\Phi=0$ \underline{ou} $n\Psi=0$, \underline{alors si} $E_\eta$ \underline{est une biextension de} $(P_\eta,Q_\eta)$ \underline{par} $\underline{G}_{m\eta}$ , \underline{il lui est associé un accouplement canonique}

\begin{equation}\tag{7.1.2}
d(E_\eta) : \Phi \otimes_{\mathbb{Z}} \Psi \longrightarrow (\mathbb{Q}/\mathbb{Z})_k \qquad ,
\end{equation}
\underline{dont l'annulation est nécessaire et suffisante pour que} $E$ \underline{soit isomorphe à la fibre générique d'une biextension} $E$ \underline{de} $(P,Q)$ \underline{par} $\underline{G}_{m\,S}$. \underline{En particulier, si} $P_o$ \underline{ou} $Q_o$ \underline{est connexe, le foncteur précédent} $E \longmapsto E_\eta$ \underline{est une équivalence de catégories.}

\underline{Démonstration.}  a)  Soit $\overline{G}$ le modèle de Néron-Raynaud de $\underline{G}_m$ . Le foncteur envisagé est le composé du foncteur "extension du noyau $\underline{G}_{mS} \longrightarrow \overline{G}$ " avec le foncteur "restriction à $\eta$ " défini sur la catégorie des biextensions par $\overline{G}$ . Ce dernier foncteur étant une équivalence de catégories en vertu de 6.6, on est ramené à étudier le premier foncteur. Pour ceci, on utilise la suite exacte de schémas en groupes

}{\begin{equation*}
\Phi = P_o/P_o^{\,o} \qquad , \qquad \Psi = Q_o/Q_o^{\,o} \qquad ,
\end{equation*}
\underline{et supposons que} $\Phi$ \underline{ou} $\Psi$ \underline{soit un groupe de torsion} (\underline{par exemple} $P$ \underline{ou} $Q$ \underline{de type fini sur} $S$) ; \underline{alors le foncteur} $E \longmapsto E_\eta$ \underline{de la catégorie des biextensions de} $(P,Q)$ \underline{par} $\underline{G}_{mS}$ \underline{dans la catégorie des biextensions de} $(P_\eta,Q_\eta)$ \underline{par} $\underline{G}_{m\eta}$ \underline{est pleinement fidèle.} \underline{Supposons qu'il existe un entier} $n > 0$ \underline{tel que} $n\Phi=0$ \underline{ou} $n\Psi=0$, \underline{alors si} $E_\eta$ \underline{est une biextension de} $(P_\eta,Q_\eta)$ \underline{par} $\underline{G}_{m\eta}$ , \underline{il lui est associé un accouplement canonique}

\begin{equation}\tag{7.1.2}
d(E_\eta) : \Phi \otimes_{\mathbb{Z}} \Psi \longrightarrow (\mathbb{Q}/\mathbb{Z})_k \qquad ,
\end{equation}
\underline{dont l'annulation est nécessaire et suffisante pour que} $E_\eta$ \underline{soit isomorphe à la fibre générique d'une biextension} $E$ \underline{de} $(P,Q)$ \underline{par} $\underline{G}_{mS}$. \underline{En particulier, si} $P_o$ \underline{ou} $Q_o$ \underline{est connexe, le foncteur précédent} $E \longmapsto E_\eta$ \underline{est une équivalence de catégories.}

\underline{Démonstration.}  a)  Soit $\overline{G}$ le modèle de Néron-Raynaud de $\underline{G}_m$ . Le foncteur envisagé est le composé du foncteur "extension du noyau $\underline{G}_{mS} \longrightarrow \overline{G}$ " avec le foncteur "restriction à $\eta$ " défini sur la catégorie des extensions par $\overline{G}$ . Ce dernier foncteur étant une équivalence de catégories en vertu de 6.6, on est ramené à étudier le premier foncteur. Pour ceci, on utilise la suite exacte de schémas en groupes

}{F7I-VIII-066}\begin{equation}\tag{7.1.3}
0 \longrightarrow \underline{G}_{mS} \longrightarrow \overline{G} \longrightarrow \mathbb{Z}_{S,s} \longrightarrow 0 \qquad ,
\end{equation}
la notation pour $\mathbb{Z}_{S,s}$ étant celle de 5.7. La suite exacte des $\mathrm{Ext}^i$ ($i=0,1$) associée à cette suite exacte nous montre que pour que le foncteur envisagé soit pleinement fidèle, il faut et il suffit que l'on ait 


\SGApage{I}{301}

\begin{equation}\tag{7.1.4}
 \mathrm{Hom}(P,\mathbb{Z}_{S,s}) = 0 \qquad ,
\end{equation}
et que l'obstruction, pour une extension $\overline{E}$ de $P$ par $\overline{G}$ , de provenir d'une extension $E$ de $P$ par $\underline{G}_{mS}$ se trouve dans

\begin{equation}\tag{7.1.5}
\mathrm{Ext}^1(P,\mathbb{Z}_{S,s}) \qquad .
\end{equation}
Or le premier groupe (7.1.4) est isomorphe à $\mathrm{Hom}(P_o,\mathbb{Z}_k)$ en vertu de (5.7.2), qui est isomorphe lui-même à $\mathrm{Hom}(\Phi,\ \mathbb{Z}_k)$, de sorte que (7.1.4) équivaut à la relation

\begin{equation}\tag{7.1.4 bis}
\mathrm{Hom}(\Phi,\mathbb{Z}_k) = 0 \qquad .
\end{equation}
Cette relation est certainement vérifiée si $\Phi$ est un groupe de torsion. D'autre part, si on a même $n\Phi = 0$, alors en vertu de 5.9 et 5.5 on a des isomorphismes canoniques

\begin{equation}\tag{7.1.6}
\mathrm{Ext}^1(P,\mathbb{Z}_{S,s}) \simeq \mathrm{Ext}^1(P_o,\mathbb{Z}_k) \simeq \mathrm{Hom}(\Phi,\ (\mathbb{Q}/\mathbb{Z})_k) \qquad ,
\end{equation}
ce qui achève de prouver 7.1 a).

\SGAvar{b)  La démonstration est essentiellement la même. On est ramené, grâce à 6.7, à étudier le foncteur "extension du noyau $\underline{G}_{mS} \longrightarrow \overline{G}$ " de la catégorie des biextensions de $(P,Q)$ par $\underline{G}_{mS}$ dans celle des biextensions de $(P,Q)$ par $\overline{G}$ . Cette étude se fait en utilisant la suite exacte (7.1.3), et la suite exacte des $\mathrm{Biext}^i$ ($i=0,1$) associée (VII 3.7.5). La pleine fidèlité du foncteur envisagé équivaut à la relation

\begin{equation}\tag{7.1.7}
\mathrm{Hom}(P \otimes Q,\ \mathbb{Z}_{S,s}) = 0 \qquad ,
\end{equation}
or le premier membre est isomorphe à $\mathrm{Hom}(P_o \otimes Q_o\ ,\ \mathbb{Z}_k)$ en vertu de (5.7.2),
\SGApage{I}{302}
donc à $\mathrm{Hom}(\Phi \otimes \Psi, \mathbb{Z}_k)$, groupe qui est bien nul si $\Phi$ ou $\Psi$ est un groupe de torsion. D'autre part, pour une biextension donnée $\overline{E}$ de $(P,Q)$ par $\overline{G}$, l'obstruction à ce qu'elle provienne d'une biextension $E$ de $(P,Q)$ par $\underline{G}_{mS}$ se trouve dans $\mathrm{Biext}^1(P,Q;\mathbb{Z}_{S,s})$, qui est isomorphe en vertu de 5.10 à $\mathrm{Biext}^1(P_o,Q_o, \mathbb{Z}_k)$, lui-même isomorphe (si $n\Phi = 0$ ou $n\Psi = 0$) à $\mathrm{Hom}(\Phi \otimes \Psi, Q/\mathbb{Z})$ en vertu de 5.6. Cela achève la démonstration de 7.1.

}{b)  La démonstration est essentiellement la même. On est ramené, grâce à 6.7, à étudier le foncteur "extension du noyau $\underline{G}_{mS} \longrightarrow \overline{G}$ " de la catégorie des biextensions de $(P,Q)$ par $\underline{G}_{mS}$ dans celle des biextensions de $(P,Q)$ par $\overline{G}$ . Cette étude se fait en utilisant la suite exacte (7.1.3), et la suite exacte des $\mathrm{Biext}^i$ ($i=0,1$) associée (VII 3.7.5). La pleine fidélité du foncteur envisagé équivaut à la relation

\begin{equation}\tag{7.1.7}
\mathrm{Hom}(P \otimes Q,\ \mathbb{Z}_{S,s}) = 0 \qquad ,
\end{equation}
or le premier membre est isomorphe à $\mathrm{Hom}(P_o \otimes Q_o\ ,\ \mathbb{Z}_k)$ en vertu de (5.7.2),
\SGApage{I}{302}
donc à $\mathrm{Hom}(\Phi \otimes \Psi, \mathbb{Z}_k)$, groupe qui est bien nul si $\Phi$ ou $\Psi$ est un groupe de torsion. D'autre part, pour une biextension donnée $\overline{E}$ de $(P,Q)$ par $\overline{G}$, l'obstruction à ce qu'elle provienne d'une biextension $E$ de $(P,Q)$ par $\underline{G}_{mS}$ se trouve dans $\mathrm{Biext}^1(P,Q;\mathbb{Z}_{S,s})$, qui est isomorphe en vertu de 5.10 à $\mathrm{Biext}^1(P_o,Q_o, \mathbb{Z}_k)$, lui-même isomorphe (si $n\Phi = 0$ ou $n\Psi = 0$) à $\mathrm{Hom}(\Phi \otimes \Psi, (\mathbb{Q}/\mathbb{Z})_k)$ en vertu de 5.6. Cela achève la démonstration de 7.1.

}{F7I-VIII-067}\underline{Remarque} 7.2. On peut se proposer, étant donné un schéma en groupes lisse $P$ sur un schéma localement noethérien régulier connexe $S$ de point générique $\eta$ , d'étudier le foncteur $E \longmapsto E_\eta$ de la catégorie des extensions de $P$ par $\underline{G}_{mS}$ dans celle des extensions de $P_\eta$ par $\underline{G}_{m\eta}$ . Ce problème se ramène en fait facilement à celui traité dans 7.1 a). Supposons pour simplifier $P$ de type fini sur $S$. Alors il résulte facilement de 7.1 que le foncteur envisagé est pleinement fidèle (*). On en déduit aisément que si $E_\eta$ est une extension de $P_\eta$ par $\underline{G}_{m\eta}$ , alors pour que celle-ci soit isomorphe à la fibre générique d'une extension $E$ de $P$ par $\underline{G}_{mS}$ , il faut et il suffit que pour tout $s \in S$ tel que $\dim \mathrm{Spec}(\underline{O}_{S,s}) = 1$, le problème de prolongement analogue obtenu par le changement de base $\mathrm{Spec}(\underline{O}_{S,s}) \to S$ ait une solution, - ce qui, en vertu de 7.1 a), s'exprime par la nullité d'un certain homomorphisme

\begin{equation}\tag{7.2.1}
d(E_\eta ,s) \; : \; P_s/P_s^{\,o} \longrightarrow (\mathbb{Q}/\mathbb{Z})_{k(s)} \qquad .
\end{equation}
L'assertion précédente se réduit de façon évidente à la suivante : soit $U$ un ouvert de $S$, de complémentaire $T$ tel que

\begin{equation*}
\mathrm{codim}(T,S) \geqslant 2 \qquad , \qquad 
\end{equation*}
(*) Il suffit pour ceci que $S$ soit normal (au lieu de régulier).


\SGApage{I}{303}

\SGAvar{ alors le foncteur $E \longmapsto E|U$ de la catégorie des extensions de $P$ par $\underline{G}_{mS}$ dans celle des extensions de $P_U$ par $\underline{G}_{m\,U}$ est une équivalence de catégories. Or cette assertion est essentiellement triviale en termes de la description VII 1.1.6, en utilisant le fait que pourtout $S$-schéma lisse $X$ (tel $P, P\times P, P\times P\times P$ ...), $X$ est localement factoriel et $X_T = X - X_U$ est de codimension $\geqslant 2$ dans $X$, d'où résulte que le foncteur $F \longmapsto F|X_U$ de la catégorie des torseurs sous $\underline{G}_{mX}$ dans celle des torseurs sous $\underline{G}_{mX_U}$ est une équivalence de catégories.

}{ alors le foncteur $E \longmapsto E|U$ de la catégorie des extensions de $P$ par $\underline{G}_{mS}$ dans celle des extensions de $P_U$ par $\underline{G}_{m\,U}$ est une équivalence de catégories. Or cette assertion est essentiellement triviale en termes de la description VII 1.1.6, en utilisant le fait que pour tout $S$-schéma lisse $X$ (tel $P, P\times P, P\times P\times P$ ...), $X$ est localement factoriel et $X_T = X - X_U$ est de codimension $\geqslant 2$ dans $X$, d'où résulte que le foncteur $F \longmapsto F|X_U$ de la catégorie des torseurs sous $\underline{G}_{mX}$ dans celle des torseurs sous $\underline{G}_{mX_U}$ est une équivalence de catégories.

}{F7I-VIII-068}On a des résultats tout analogues pour la question des prolongements de biextensions, qu'on laisse au lecteur le soin de formuler.

7.3. Nous allons expliciter quelques propriétés fonctorielles des obstructions définies dans (7.1.1) et (7.1.2) pour $P,Q,S$ variables, et de leurs relations entre elles. On supposera vérifiées par la suite les conditions qui permettent de définir ces obstructions. Comme les vérifications n'offrent pas de difficulté, nous les laissons aux soins du lecteur.

7.3.1. Soit

\SGAvar{\begin{equation*}
u \; : \; P' \longrightarrow P
\end{equation*}
un morphisme de schémas en groupes lisses sur $S$, $E_\eta$ une extension de $P_\eta$ par $\underline{G}_{m\eta}$ , d'où par image inverse par $u_\eta : P'_\eta \longrightarrow P_\eta$ une extension $E'_\eta$ de $P'_\eta$ par $\underline{G}_m$ . Ceci posé, on a commutativité dans le diagramme 


}{\begin{equation*}
u \; : \; P' \longrightarrow P
\end{equation*}
un morphisme de schémas en groupes lisses sur $S$, $E_\eta$ une extension de $P_\eta$ par $\underline{G}_{m\eta}$ , d'où par image inverse par $u_\eta : P'_\eta \longrightarrow P_\eta$ une extension $E'_\eta$ de $P'_\eta$ par $\underline{G}_{m\eta}$ . Ceci posé, on a commutativité dans le diagramme 


}{F7I-VIII-069}\SGApage{I}{304}

\begin{equation}\tag{7.3.1.1}
\begin{tikzcd}[row sep=large, column sep=huge]
{P'_o/{P'_o}^{o}} \arrow[dd, "u_o"'] \arrow[rd, "d(E'_\eta)"] \& \\
 \& {(\mathbb{Q}/\mathbb{Z})_k} \\
{P_o/{P_o}^{o}} \arrow[ru, "d(E_\eta)"'] \&
\end{tikzcd}\quad .
\end{equation}










\SGAvar{De même, si on des morphismes de schémas en groupes lisses sur $S$

}{De même, si on a des morphismes de schémas en groupes lisses sur $S$

}{F7I-VIII-070}\begin{equation*}
u : P' \longrightarrow P \quad , \quad v : Q' \longrightarrow Q
\end{equation*}
et une biextension $E_\eta$ de $(P_\eta ,Q_\eta)$ par $\underline{G}_{m\eta}$ , d'où par image inverse par $(u_\eta ,v_\eta)$ une biextension $E'_\eta$ de $(P'_\eta ,Q'_\eta)$ par $\underline{G}_{m\eta}$ , on a commutativité dans le diagramme

\begin{equation}\tag{7.3.1.2}
\begin{tikzcd}[row sep=large, column sep=huge]
{P'_o/{P'_o}^{o} \otimes Q'_o/{Q'_o}^{o}} \arrow[dd, "u_o \otimes v_o"'] \arrow[rd, "d(E'_\eta)"] \& \\
 \& {(\mathbb{Q}/\mathbb{Z})_k} \\
{P_o/{P_o}^{o} \otimes Q_o/{Q_o}^{o}} \arrow[ru, "d(E_\eta)"'] \&
\end{tikzcd}
\end{equation}









La vérification, essentiellement mécanique, de ces compatibilités est laissée au lecteur.

7.3.2. De la compatibilité (7.3.1.1) résulte que si $E_\eta$ est une extension de $P_\eta$ par $\underline{G}_m$ , il existe un plus grand sous-schéma en groupes $V$ de $P$ tel que la restriction de $E_\eta$ à $V_\eta$ s'étende en une extension de $V$ par
\SGApage{I}{305}
$\underline{G}_{mS}$ : c'est le sous-groupe ouvert $V$ de $P$ contenant $P_\eta$ défini par la condition

\begin{equation}\tag{7.3.2.1}
V_o/{V_o}^{o} = \mathrm{Ker}\ d(E_\eta) \quad .
\end{equation}
Si alors $u:P' \longrightarrow P$ est un morphisme de schémas en groupes lisses sur $S$, l'extension $E'_\eta$ , image inverse de $E_\eta$ par $u_\eta$ , se prolonge en une extension de $P'$ par $\underline{G}_{mS}$ si et seulement si on a

\SGAvar{\begin{equation}\tag{7.3.2.2}
u(P') \subset V \quad .
\end{equation}
Dans le cas d'une biextension $E$ de $(P_\eta , Q_\eta)$ par $\underline{G}_{m\eta}$ , il existe de même un plus grand sous-schéma en groupes $V$ de $P$ tel que la restriction de $E$ à $(V_\eta , Q_\eta)$ se prolonge en une biextension de $(V,Q)$ par $\underline{G}_{mS}$ : c'est le sous-groupe ouvert de $P$, contenant la fibre générique $P_\eta$ , défini par la condition que $V_o/{V_o}^{o}$ soit le sous-groupe de $\Phi$ annulateur gauche de l'accouplement (7.1.2), i.e. noyau de l'homomorphisme correspondant de $\Phi$ dans $\underline{\mathrm{Hom}}(\Psi, (\mathbb{Q}/\mathbb{Z})_k)$ :

\begin{equation}\tag{7.3.2.3}
V_o/{V_o}^{o} = \mathrm{Ker}\ (\Phi \xrightarrow{\ d(E_\eta)\ } \underline{\mathrm{Hom}}(\Psi ,\mathbb{Q}/\mathbb{Z} )_k)) \quad .
\end{equation}
Si alors on se donne $u: P' \longrightarrow P$ de schémas en groupes lisses sur $S$, la biextension de $(P'_\eta , Q_\eta)$ par $\underline{G}_{m\eta}$ image inverse de $E_\eta$ par $u_\eta$ se prolonge en une biextension de $(P',Q)$ par $\underline{G}_{mS}$ si et seulement si on a la relation (7.3.2.2). On a des résultats symétriques en échangeant le rôle de $P$ et de $Q$. Mais bien entendu on ne pourra en général trouver un plus grand \underline{couple} de sous-groupes $P' \longrightarrow P$, $Q' \longrightarrow Q$ tels que la restriction $E'_\eta$ de $E_\eta$ en une biextension de $(P'_\eta , Q'_\eta)$ par $\underline{G}_{m\eta}$ se prolonge en une biextension de $(P',Q')$ par $\underline{G}_{mS}$ ; on peut seulement dire que pour un couple
\SGApage{I}{306}
$(P',Q')$ donné, plus généralement pour un couple de morphismes donné $P' \longrightarrow P$, $Q' \longrightarrow Q$ de schémas en groupes lisses, le prolongement cherché existe si et seulement si la restriction de la forme $d(E_\eta)$ à $P'_o/P'^o_o \otimes Q'_o/Q'^o_o$ est nulle.

}{\begin{equation}\tag{7.3.2.2}
u(P') \subset V \quad .
\end{equation}
Dans le cas d'une biextension $E_\eta$ de $(P_\eta , Q_\eta)$ par $\underline{G}_{m\eta}$ , il existe de même un plus grand sous-schéma en groupes $V$ de $P$ tel que la restriction de $E_\eta$ à $(V_\eta , Q_\eta)$ se prolonge en une biextension de $(V,Q)$ par $\underline{G}_{mS}$ : c'est le sous-groupe ouvert de $P$, contenant la fibre générique $P_\eta$ , défini par la condition que $V_o/{V_o}^{o}$ soit le sous-groupe de $\Phi$ annulateur gauche de l'accouplement (7.1.2), i.e. noyau de l'homomorphisme correspondant de $\Phi$ dans $\underline{\mathrm{Hom}}(\Psi, (\mathbb{Q}/\mathbb{Z})_k)$ :

\begin{equation}\tag{7.3.2.3}
V_o/{V_o}^{o} = \mathrm{Ker}\ (\Phi \xrightarrow{\ d(E_\eta)\ } \underline{\mathrm{Hom}}(\Psi ,(\mathbb{Q}/\mathbb{Z} )_k)) \quad .
\end{equation}
Si alors on se donne $u: P' \longrightarrow P$ de schémas en groupes lisses sur $S$, la biextension de $(P'_\eta , Q_\eta)$ par $\underline{G}_{m\eta}$ image inverse de $E_\eta$ par $u_\eta$ se prolonge en une biextension de $(P',Q)$ par $\underline{G}_{mS}$ si et seulement si on a la relation (7.3.2.2). On a des résultats symétriques en échangeant le rôle de $P$ et de $Q$. Mais bien entendu on ne pourra en général trouver un plus grand \underline{couple} de sous-groupes $P' \longrightarrow P$, $Q' \longrightarrow Q$ tels que la restriction $E'_\eta$ de $E_\eta$ en une biextension de $(P'_\eta , Q'_\eta)$ par $\underline{G}_{m\eta}$ se prolonge en une biextension de $(P',Q')$ par $\underline{G}_{mS}$ ; on peut seulement dire que pour un couple
\SGApage{I}{306}
$(P',Q')$ donné, plus généralement pour un couple de morphismes donné $P' \longrightarrow P$, $Q' \longrightarrow Q$ de schémas en groupes lisses, le prolongement cherché existe si et seulement si la restriction de la forme $d(E_\eta)$ à $P'_o/P'^o_o \otimes Q'_o/Q'^o_o$ est nulle.

}{F7I-VIII-071}7.3.3. \underline{Symétrie}. Soit $E_\eta$ une biextension de $(P_\eta , Q_\eta)$ par $\underline{G}_{m\eta}$ , et considérons la biextension symétrique ${}^{s}E_\eta$ (VII 2.7), qui est une biextension de $(Q_\eta , P_\eta)$ par $\underline{G}_{m\eta}$ . On a alors commutativité dans le diagramme

\begin{equation}\tag{7.3.3.1}
\begin{tikzcd}[row sep=large, column sep=huge]
\Phi \otimes \Psi \arrow[dd, "\mathrm{sym}"'] \arrow[rd, "d(E_\eta)"] \& \\
 \& (\mathbb{Q}/\mathbb{Z})_k \\
\Psi \otimes \Phi \arrow[ru, "d({}^{s}E_\eta)"'] \&
\end{tikzcd}\quad ,
\end{equation}














où la flèche verticale est la symétrie canonique. (On comparera cet énoncé avec l'énoncé d'\underline{antisymétrie} de 2.2.10.) On en déduit en particulier que si $P=Q$, et si $E_\eta$ est une biextension \underline{symétrique} de $(P_\eta , Q_\eta)$ par $\underline{G}_{m\eta}$ , i.e. si $E_\eta$ est isomorphe à ${}^{s}E_\eta$ , alors la forme bilinéaire $d(E_\eta)$ sur $\Phi \times \Phi$ à valeurs dans $(\mathbb{Q}/\mathbb{Z})_k$ est \underline{symétrique} (et non alternée comme dans 2.2.12 !). On notera que si elle est nulle, i.e. si $E$ se prolonge en une biextension de $(P,P)$ par $\underline{G}_{mS}$ , cette dernière est nécessairement symétrique également, puisqu'en vertu de 7.1 a) tout isomorphisme ${}^{s}E_\eta \longrightarrow E_\eta$ se prolonge de façon unique en un isomorphisme ${}^{s}E \longrightarrow E$ .


\SGApage{I}{307}


7.3.4.  \underline{Réduction des obstructions du type} (7.1.2) \underline{à celles du type} (7.1.1).

Supposons pour fixer les idées qu'il existe un entier $n > 0$ tel que $n \Phi = 0$, et soit

\begin{equation*}
q \in Q(S)
\end{equation*}
une section de $Q$. On désigne, par abus de notation, par

\SGAvar{\begin{equation*}
q_o \in \Psi(k)
\end{equation*}
l'image dans $\Psi(k)$ de la valeur de $q$ en $s$. Considérant $E_\eta$ comme une extension de $P_{\eta Q_\eta}$ par $\underline{G}_{mQ_\eta}$ , on trouve une extension

}{\begin{equation*}
q_o \in \Psi(k)
\end{equation*}
l'image dans $\Psi(k)$ de la valeur de $q$ en $s$. Considérant $E_\eta$ comme une extension de $P_\eta \times_\eta Q_\eta$ par $\underline{G}_{mQ_\eta}$ , on trouve une extension

}{F7I-VIII-072}\begin{equation*}
q^*_\eta(E_\eta) \ , \ \text{extension de } P_\eta \ \text{par} \ \underline{G}_{m\eta} \ ,
\end{equation*}
d'où une obstruction

\SGAvar{\begin{equation}\tag{7.3.4.1}
d_q \ \overset{\mathrm{dfn}}{=} \ d(q^*_\eta(E_*)) : \Phi \longrightarrow (\mathbb{Q}/\mathbb{Z})_k \ .
\end{equation}
Ceci posé, on a la relation (où on a posé $d = d(E_\eta)$)

}{\begin{equation}\tag{7.3.4.1}
d_q \ \overset{\mathrm{dfn}}{=} \ d(q^*_\eta(E_\eta)) : \Phi \longrightarrow (\mathbb{Q}/\mathbb{Z})_k \ .
\end{equation}
Ceci posé, on a la relation (où on a posé $d = d(E_\eta)$)

}{F7I-VIII-073}\begin{equation}\tag{7.3.4.2}
d_q(x) = d(x,q_o) \ ,
\end{equation}
valable pour tout point $x$ de $\Phi$ à valeurs dans un $k$-schéma $T$ .

\SGAvar{7.3.4.3.  Cette formule, qui pour tout $q \in Q(S)$ détermine $d_q$ en termes de l'accouplement $d$ , redonne évidemment ce dernier en termes des invariants $d_q$ , du moins lorsque tout point de $\Psi$ provient d'une section de $Q$. Cette condition est d'ailleurs vérifiée automatiquement lorsque $S$ est strictement\_local, en vertu du lemme de Hensel. Comme l'accouplement $d$ est connu quand on le connaît sur la clôture séparable de $k$, et que sa formation commute à la localisation étale sur $S$ en vertu de 7.3.5.3
\SGApage{I}{308}
ci-dessous, on voit que, quitte à passer au localisé strict de $S$, la détermination de l'accouplement $d$ se ramène à celle des invariants $d_q$ . En particulier, lorsque tout point de $\Psi$ provient d'une section $q$ de $Q$, on voit que $d$ est nul (i.e. $E_\eta$ se prolonge en une biextension $E$ de $(P,Q)$ par $\underline{G}_m$) si et seulement si pour tout section $q$ de $Q$, $d_q$ est nul i.e. $q^*_\eta(E_\eta)$ se prolonge en une extension de $P$ par $\underline{G}_m$ . Lorsqu'on ne fait pas d'hypothèse sur $\Psi$, le critère analogue est valable, à condition de prendre pour $q$ des points à valeurs dans des $S'$ locaux étales sur $S$.

}{7.3.4.3.  Cette formule, qui pour tout $q \in Q(S)$ détermine $d_q$ en termes de l'accouplement $d$ , redonne évidemment ce dernier en termes des invariants $d_q$ , du moins lorsque tout point de $\Psi$ provient d'une section de $Q$. Cette condition est d'ailleurs vérifiée automatiquement lorsque $S$ est strictement\_local, en vertu du lemme de Hensel. Comme l'accouplement $d$ est connu quand on le connaît sur la clôture séparable de $k$, et que sa formation commute à la localisation étale sur $S$ en vertu de 7.3.5.3
\SGApage{I}{308}
ci-dessous, on voit que, quitte à passer au localisé strict de $S$, la détermination de l'accouplement $d$ se ramène à celle des invariants $d_q$ . En particulier, lorsque tout point de $\Psi$ provient d'une section $q$ de $Q$, on voit que $d$ est nul (i.e. $E_\eta$ se prolonge en une biextension $E$ de $(P,Q)$ par $\underline{G}_m$) si et seulement si pour toute section $q$ de $Q$, $d_q$ est nul i.e. $q^*_\eta(E_\eta)$ se prolonge en une extension de $P$ par $\underline{G}_m$ . Lorsqu'on ne fait pas d'hypothèse sur $\Psi$, le critère analogue est valable, à condition de prendre pour $q$ des points à valeurs dans des $S'$ locaux étales sur $S$.

}{F7I-VIII-074}7.3.4.4. Inversement, partant de la situation de 7.1 a), l'obstruction (7.1.1) peut s'interpréter comme une obstruction du type (7.1.2), en prenant $Q = \mathbb{Z}_S$ et la biextension de $(P,\mathbb{Z}_S)$ déduite de $E$. Pour le voir, il suffit d'appliquer ce qui précède en prenant pour $q$ la section 1 de $\mathbb{Z}_S$ .

7.3.5. \underline{Effet d'un changement de base} $S' \longrightarrow S$. On suppose donné un morphisme local dominant de traits

\begin{equation*}
f : S' \longrightarrow S \quad ,
\end{equation*}
et on désigne par $\eta'$ le point générique de $S'$, par $s'$ son point fermé, et on pose $k'=k(s')$. Si $S=\mathrm{Spec}(V)$, $S' = \mathrm{Spec}(V')$, $V'$ est donc un anneau de valuation discrète dominant l'anneau de valuation discrète $V$. Posons

\SGAvar{\begin{equation}\tag{7.3.5.1}
e(S'/S) = e(V'/V) = \mathrm{long}(V'/\underline{m}\,V) \quad ,
\end{equation}
où $\underline{m}$ est l'idéal maximal de $V$. Soient maintenant $P$ dans 7.1 a), $E_\eta$ une extension de $P_\eta$ par $\underline{G}_{m\eta}$ , $P' = P\times_S S'$, $E'_{\eta'}$ l'extension de $P'_{\eta'}$ par $\underline{G}_{m\eta'}$ image inverse de $E_\eta$ . On a alors un isomorphisme canonique 


}{\begin{equation}\tag{7.3.5.1}
e(S'/S) = e(V'/V) = \mathrm{long}(V'/\underline{m}\,V') \quad ,
\end{equation}
où $\underline{m}$ est l'idéal maximal de $V$. Soient maintenant $P$ dans 7.1 a), $E_\eta$ une extension de $P_\eta$ par $\underline{G}_{m\eta}$ , $P' = P\times_S S'$, $E'_{\eta'}$ l'extension de $P'_{\eta'}$ par $\underline{G}_{m\eta'}$ image inverse de $E_\eta$ . On a alors un isomorphisme canonique 


}{F7I-VIII-075}\SGApage{I}{309}

\begin{equation*}
 \Phi' \overset{\mathrm{dfn}}{=} P'_o/{P'}^{o}_{o} \simeq \Phi \otimes_k k' \simeq \Phi_{k'} \quad ,
\end{equation*}
avec les notations de 4.1 a), et on peut par suite considérer

\begin{equation*}
d(E_\eta)_{k'} : \Phi' \longrightarrow (\mathbb{Q}/\mathbb{Z})_{k'} \quad .
\end{equation*}
Ceci posé, on a la relation

\begin{equation}\tag{7.3.5.2}
d(E'_{\eta'}) = e(S'/S)\, d(E_\eta)_{k'} \quad .
\end{equation}
De même, sous les conditions de 7.1 b), et avec les notations évidentes, on a la relation (7.3.5.2), où $d(E_\eta)$ désigne maintenant l'accouplement (7.1.2), $d(E_\eta)_{k'}$ l'accouplement

\SGAvar{\begin{equation*}
\Phi' \otimes \Psi' \longrightarrow (\mathbb{Q}/\mathbb{Z})_{k'}
\end{equation*}
qui s'en déduit par le changement de base $k \longrightarrow k'$, et où $E'_{\eta'}$ est la biextension de $(P'_{\eta'}, Q'_{\eta'})$ par $\underline{G}_{m\eta'}$ déduite de la biextension $E_\eta$ de $(P_\eta,Q_\eta)$ par $\underline{G}_{m\eta}$ à l'aise du changement de base $\eta' \longrightarrow \eta$ .

7.3.5.3. Un cas intéressant est celui où on a $e(S'/S) = 1$ , i.e. où

}{\begin{equation*}
\Phi' \otimes \Psi' \longrightarrow (\mathbb{Q}/\mathbb{Z})_{k'}
\end{equation*}
qui s'en déduit par le changement de base $k \longrightarrow k'$, et où $E'_{\eta'}$ est la biextension de $(P'_{\eta'}, Q'_{\eta'})$ par $\underline{G}_{m\eta'}$ déduite de la biextension $E_\eta$ de $(P_\eta,Q_\eta)$ par $\underline{G}_{m\eta}$ à l'aide du changement de base $\eta' \longrightarrow \eta$ .

7.3.5.3. Un cas intéressant est celui où on a $e(S'/S) = 1$, i.e. où

}{F7I-VIII-076}\begin{equation*}
\underline{m}' = \underline{m}\, V' \quad ,
\end{equation*}
où $\underline{m}'$ désigne l'idéal maximal de $V'$. C'est le cas en particulier si $V'$ est étale sur $V$, ou si c'est le hensélisé, ou le hensélisé strict, ou le complété de $V$. Dans ce cas la relation (7.3.5.2) s'exprime simplement en disant que \underline{la formation de l'obstruction} (7.1.1) \underline{resp}. (7.1.2) \underline{commute au changement de base envisagé} $S' \longrightarrow S$.

7.3.5.4. D'autre part, la relation (7.3.5.2) dans le cas général montre que pour toute extension $E_\eta$ comme dans 7.1 a) (resp. toute biextension $E_\eta$ comme dans 7.1 b)), il existe un morphisme de traits \underline{fini} $S' \longrightarrow S$,
\SGApage{I}{310}
tel que $k(\eta')$ soit une extension séparable de $k(\eta)$, et que l'extension (resp. la biextension) $E'_{\eta'}$ déduite de $E_{\eta}$ par changement de base $\eta' \longrightarrow \eta$ se prolonge en une extension (resp. biextension) $E'$ sur $S'$ tout entier. En effet, il suffit de choisir $S'/S$ tel que $e(S'/S)$ soit un multiple de l'entier $n > 0$ figurant dans 7.1, et on sait qu'on peut toujours trouver une telle extension, avec $k(\eta')$ séparable sur $k(\eta)$.


\underline{Proposition} 7.4. \underline{Soit} $S$ \underline{un schéma réduit.}

a) \underline{Soit} $P$ \underline{un schéma en groupes lisse} \underline{sur} $S$ \underline{dont les fibres aux points maximaux de} $S$ \underline{sont connexes.} \underline{Alors le foncteur naturel}

VII 1.3.4.1

\begin{equation}\tag{7.4.1}
\mathrm{EXT}(P,\underline{G}_{mS}) \longrightarrow \mathrm{TORSRIG}(P,e_P;\underline{G}_{mS})
\end{equation}
\underline{de la catégorie des extensions de} $P$ \underline{par} $\underline{G}_{mS}$ \underline{dans la catégorie des torseurs sous} $\underline{G}_{mP}$ \underline{rigidifiés par rapport à la section unité de} $P$, \underline{est} pleinement fidèle. \underline{Si} $S$ \underline{est normal}, \underline{et} $P$ \underline{à fibres connexes}, \underline{alors un objet} $E$ \underline{du deuxième membre de} (7.4.1) \underline{appartient à l'image essentielle de ce foncteur si et seulement si pour tout point maximal} $\eta$ \underline{de} $S$, \underline{le torseur rigidifié} $E_{\eta}$ \underline{sur} $P_{\eta}$ \underline{provient d'une extension de} $P_{\eta}$ \underline{par} $\underline{G}_{m\eta}$ .

b) \underline{Soient} $P,Q$ \underline{deux schémas en groupes lisses sur} $S$ \underline{dont les fibres aux points maximaux de} $S$ \underline{sont connexes.} \underline{Alors le foncteur naturel}

VII 2.9.1

\begin{equation}\tag{7.4.2}
\mathrm{BIEXT}(P,Q;\underline{G}_{mS}) \longrightarrow \mathrm{TORSBIRIG}(P,Q;\underline{G}_{mS})
\end{equation}
\underline{de la catégorie des biextensions de} $(P,Q)$ \underline{par} $\underline{G}_{mS}$ \underline{dans la catégorie des torseurs sous} $\underline{G}_{m\,P\times_S Q}$ \underline{birigidifiés relativement au couple des}
\SGApage{I}{311}
\underline{sections unités de} $P$ \underline{et de} $Q$, \underline{est} pleinement fidèle. \underline{Si} $S$ \underline{est normal}, \underline{et} $P$ \underline{ou} $Q$ \underline{à fibres connexes}, \underline{un objet} $E$ \underline{du second membre de} (7.4.2) \underline{appartient à l'image essentielle du foncteur si et seulement si pour tout point maximal} $\eta$ \underline{de} $S$, $E_{\eta}$ \underline{provient d'une biextension de} $(P_{\eta},Q_{\eta})$ \underline{par} $\underline{G}_{m\eta}$ .

Nous prouverons a), la démonstration de b), essentiellement identique, étant laissée au lecteur. Pour l'assertion de pleine fidélité, on peut noter que 4.1 reste valable si $X$ et $Y$ sont des schémas plats sur un schéma réduit $S$, dont les fibres aux points maximaux satisfont aux conditions énoncées dans 4.1 : en effet, la famille des fibres maximales de $X\times_S Y$ (qui est plat sur $S$ réduit) est schématiquement dense dans $X\times_S Y$, de sorte qu'un morphisme de $X\times_S Y$ dans le $S$-schéma séparé $\underline{G}_{mS}$ est trivial dès qu'il l'est sur les fibres maximales. Ceci dit, la pleine fidélité de (7.4.1) se démontre comme 4.2, et la caractérisation qui y est donnée de l'image essentielle reste également valable. Il reste à prouver que si le torseur rigidifié $E$ sur $P$ est tel que les $E_{\eta}$ proviennent d'extensions, il en est de même de $E$, lorsqu'on suppose $S$ normal et $P$ à fibres connexes. Cette dernière condition implique que $P$ est de présentation finie sur $S$ (SGA 3 $\mathrm{VI}_B$ 5.5) ; la pleine fidélité nous montre d'autre part que la question est locale sur $S$, qu'on peut donc supposer local affine, et le passage à la limite habituel (utilisant le théorème de NAGATA de finitude de la clôture intégrale d'une $\mathbb{Z}$ -algèbre réduite de type fini EGA IV 7.8) nous ramène au cas où $S$ est de plus supposé noethérien. Alors 7.1 a) nous montre qu'il existe un ouvert $U \subset S$, dont le complémentaire est de codimension $\geq 2$, tel que $E_U$ provienne d'une
\SGApage{I}{312}
extension de $P_U$ par $\underline{G}_{mU}$ . En d'autres termes, l'isomorphisme cherché

\begin{equation*}
\pi^*(E) \simeq \mathrm{pr}_1^*(E)\mathrm{pr}_2^*(E)
\end{equation*}
peut être construit au-dessus de $(P \times_S P)_U$ . Or c'est là un ouvert du schéma normal $P \times_S P$, dont le complémentaire est de codimension $\geq 2$. Donc l'isomorphisme en question sur $(P \times_S P)_U$ se prolonge à $P \times_S P$ tout entier, ce qui achève la démonstration.

\underline{Corollaire} 7.5. \underline{Soient} $S$ \underline{un schéma normal}, $P$ \underline{et} $Q$ \underline{deux schémas en groupes lisses sur} $S$, \underline{à fibres maximales connexes}, $P$ \underline{ou} $Q$ \underline{étant à fibres connexes}. \underline{Supposons de plus que pour tout point maximal} $\eta$ \underline{de} $S$, $P_\eta$ \underline{ou} $Q_\eta$ \underline{soit une extension d'un schéma abélien par un tore} (\underline{plus généralement qu'il admette une suite de composition dont les facteurs soient des schémas abéliens}, \underline{des tores ou des groupes} $\underline{G}_a$). \underline{Alors le foncteur} (7.4.2) \underline{est une équivalence de catégories}.

C'est en effet une conséquence immédiate de 7.4 b) et de 4.7.


\underline{BIBLIOGRAPHIE}

\SGAvar{[1] \quad M. Demazure, J.Giraud, M.Raynaud, Schémas abéliens, Séminaire Orsay 1967/68, à parâitre.

}{[1] \quad M. Demazure, J.Giraud, M.Raynaud, Schémas abéliens, Séminaire Orsay 1967/68, à paraître.

}{F7I-VIII-077}[2] \quad M. Raynaud, Faisceaux amples sur les schémas en groupes et les espaces homogènes, thèse Paris 1968 (à paraître dans Lecture Notes, Springer).

[3] \quad M. Rosenlicht, Some basic theorems on algebraic groups, Amer. Jour. Math. vol 78 n°2 (1956), p.401-443.

[3bis] M. Rosenlicht, Some rationality questions on algebraic groups, Annali di Matematica, Série 4, n° 43 (1957), p.25-50.

[4] \quad J.P. Serre, corps locaux, Act. Sci. Ind, 1296 (1962), Hermann Paris.
\selectlanguage{french}

\ifnum\SGAform=1
\clearpage\section*{Apparat critique}\small
Les lectures ci-dessous sont données dans le TeX exact des témoins, afin de préserver les distinctions de notation. Les indices décimaux renvoient aux différences individuelles du passage marqué.

\par\medskip\noindent\hypertarget{app-F7I-VIII-001.1}{}\textbf{F7I-VIII-001.1}\quad p. 220; inherited corrected-layer difference.\par
\noindent Imprimé : \emph{(vide)}\par
\noindent Corrigé : {\ttfamily\footnotesize ,\allowbreak{} par\allowbreak{}}\par
\noindent Source : ADOPT-F7I-VIII-001.1. \nolinkurl{fr/sga7i/diplomatic/expose_VIII_body.tex}; \nolinkurl{fr/sga7i/current/expose_VIII_body.tex}\par

\par\medskip\noindent\hypertarget{app-F7I-VIII-001.2}{}\textbf{F7I-VIII-001.2}\quad p. 220; inherited corrected-layer difference.\par
\noindent Imprimé : {\ttfamily\footnotesize suite\allowbreak{}}\par
\noindent Corrigé : {\ttfamily\footnotesize troncation\allowbreak{}}\par
\noindent Source : ADOPT-F7I-VIII-001.2. \nolinkurl{fr/sga7i/diplomatic/expose_VIII_body.tex}; \nolinkurl{fr/sga7i/current/expose_VIII_body.tex}\par

\par\medskip\noindent\hypertarget{app-F7I-VIII-001.3}{}\textbf{F7I-VIII-001.3}\quad p. 220; inherited corrected-layer difference.\par
\noindent Imprimé : {\ttfamily\footnotesize exacte\allowbreak{} infinie\allowbreak{} suivante,\allowbreak{} compte\allowbreak{} tenu\allowbreak{} que\allowbreak{} les\allowbreak{} objets\allowbreak{} d'homologie\allowbreak{}}\par
\noindent Corrigé : {\ttfamily\footnotesize standard\allowbreak{}}\par
\noindent Source : ADOPT-F7I-VIII-001.3. \nolinkurl{fr/sga7i/diplomatic/expose_VIII_body.tex}; \nolinkurl{fr/sga7i/current/expose_VIII_body.tex}\par

\par\medskip\noindent\hypertarget{app-F7I-VIII-001.4}{}\textbf{F7I-VIII-001.4}\quad p. 220; designated-witness attribution.\par
\noindent Imprimé : \emph{(vide)}\par
\noindent Corrigé : {\ttfamily\footnotesize ,\allowbreak{} dont\allowbreak{} les\allowbreak{} objets\allowbreak{} d'homologie\allowbreak{}}\par
\noindent Source : WITNESS-F7I-VIII-001.4. \nolinkurl{fr/sga7i/diplomatic/expose_VIII_body.tex}; \nolinkurl{fr/sga7i/current/expose_VIII_body.tex}\par

\par\medskip\noindent\hypertarget{app-F7I-VIII-001.5}{}\textbf{F7I-VIII-001.5}\quad p. 220; designated-witness attribution.\par
\noindent Imprimé : {\ttfamily\footnotesize qui\allowbreak{}}\par
\noindent Corrigé : {\ttfamily\footnotesize la\allowbreak{}}\par
\noindent Source : WITNESS-F7I-VIII-001.5. \nolinkurl{fr/sga7i/diplomatic/expose_VIII_body.tex}; \nolinkurl{fr/sga7i/current/expose_VIII_body.tex}\par

\par\medskip\noindent\hypertarget{app-F7I-VIII-001.6}{}\textbf{F7I-VIII-001.6}\quad p. 220; designated-witness attribution.\par
\noindent Imprimé : {\ttfamily\footnotesize sont\allowbreak{}}\par
\noindent Corrigé : {\ttfamily\footnotesize suite\allowbreak{}}\par
\noindent Source : WITNESS-F7I-VIII-001.6. \nolinkurl{fr/sga7i/diplomatic/expose_VIII_body.tex}; \nolinkurl{fr/sga7i/current/expose_VIII_body.tex}\par

\par\medskip\noindent\hypertarget{app-F7I-VIII-001.7}{}\textbf{F7I-VIII-001.7}\quad p. 220; designated-witness attribution.\par
\noindent Imprimé : {\ttfamily\footnotesize nuls\allowbreak{}}\par
\noindent Corrigé : {\ttfamily\footnotesize exacte\allowbreak{}}\par
\noindent Source : WITNESS-F7I-VIII-001.7. \nolinkurl{fr/sga7i/diplomatic/expose_VIII_body.tex}; \nolinkurl{fr/sga7i/current/expose_VIII_body.tex}\par

\par\medskip\noindent\hypertarget{app-F7I-VIII-001.8}{}\textbf{F7I-VIII-001.8}\quad p. 220; designated-witness attribution.\par
\noindent Imprimé : {\ttfamily\footnotesize pour\allowbreak{}}\par
\noindent Corrigé : {\ttfamily\footnotesize infinie\allowbreak{}}\par
\noindent Source : WITNESS-F7I-VIII-001.8. \nolinkurl{fr/sga7i/diplomatic/expose_VIII_body.tex}; \nolinkurl{fr/sga7i/current/expose_VIII_body.tex}\par

\par\medskip\noindent\hypertarget{app-F7I-VIII-001.9}{}\textbf{F7I-VIII-001.9}\quad p. 220; designated-witness attribution.\par
\noindent Imprimé : {\ttfamily\footnotesize \$i\allowbreak{} \textbackslash{}geqslant\allowbreak{} 1\$\allowbreak{}}\par
\noindent Corrigé : {\ttfamily\footnotesize suivante\allowbreak{}}\par
\noindent Source : WITNESS-F7I-VIII-001.9. \nolinkurl{fr/sga7i/diplomatic/expose_VIII_body.tex}; \nolinkurl{fr/sga7i/current/expose_VIII_body.tex}\par

\par\medskip\noindent\hypertarget{app-F7I-VIII-002.1}{}\textbf{F7I-VIII-002.1}\quad p. 220; misprint\_pass.\par
\noindent Imprimé : \emph{(vide)}\par
\noindent Corrigé : {\ttfamily\footnotesize \textbackslash{}longrightarrow\allowbreak{} }\par
\noindent Source : 7I:421. \nolinkurl{fixes/FIXLOG_7I_CURRENT.jsonl#781}; \nolinkurl{fixes/MISPRINTS_7I.json#421}\par

\par\medskip\noindent\hypertarget{app-F7I-VIII-002.2}{}\textbf{F7I-VIII-002.2}\quad p. 220; printed misprint.\par
\noindent Imprimé : {\ttfamily\footnotesize P\allowbreak{}}\par
\noindent Corrigé : {\ttfamily\footnotesize G\allowbreak{}}\par
\noindent Source : 7I:422. \nolinkurl{fixes/MISPRINTS_7I.json#422}\par

\par\medskip\noindent\hypertarget{app-F7I-VIII-003.1}{}\textbf{F7I-VIII-003.1}\quad p. 222; printed misprint.\par
\noindent Imprimé : \emph{(vide)}\par
\noindent Corrigé : {\ttfamily\footnotesize )\allowbreak{}}\par
\noindent Source : 7I:423. \nolinkurl{fixes/MISPRINTS_7I.json#423}\par

\par\medskip\noindent\hypertarget{app-F7I-VIII-004.1}{}\textbf{F7I-VIII-004.1}\quad p. 223; printed misprint.\par
\noindent Imprimé : {\ttfamily\footnotesize 4\allowbreak{}}\par
\noindent Corrigé : {\ttfamily\footnotesize 5\allowbreak{}}\par
\noindent Source : 7I:424. \nolinkurl{fixes/MISPRINTS_7I.json#424}\par

\par\medskip\noindent\hypertarget{app-F7I-VIII-005.1}{}\textbf{F7I-VIII-005.1}\quad p. 223; printed misprint.\par
\noindent Imprimé : {\ttfamily\footnotesize système\allowbreak{}}\par
\noindent Corrigé : {\ttfamily\footnotesize systèmes\allowbreak{}}\par
\noindent Source : 7I:425. \nolinkurl{fixes/MISPRINTS_7I.json#425}\par

\par\medskip\noindent\hypertarget{app-F7I-VIII-005.2}{}\textbf{F7I-VIII-005.2}\quad p. 223; printed misprint.\par
\noindent Imprimé : {\ttfamily\footnotesize dernier\allowbreak{}}\par
\noindent Corrigé : {\ttfamily\footnotesize membre\allowbreak{}}\par
\noindent Source : 7I:426. \nolinkurl{fixes/MISPRINTS_7I.json#426}\par

\par\medskip\noindent\hypertarget{app-F7I-VIII-005.3}{}\textbf{F7I-VIII-005.3}\quad p. 223; printed misprint.\par
\noindent Imprimé : {\ttfamily\footnotesize membre\allowbreak{}}\par
\noindent Corrigé : {\ttfamily\footnotesize suivant\allowbreak{}}\par
\noindent Source : 7I:426. \nolinkurl{fixes/MISPRINTS_7I.json#426}\par

\par\medskip\noindent\hypertarget{app-F7I-VIII-005.4}{}\textbf{F7I-VIII-005.4}\quad p. 223; printed misprint.\par
\noindent Imprimé : {\ttfamily\footnotesize sytème\allowbreak{}}\par
\noindent Corrigé : {\ttfamily\footnotesize système\allowbreak{}}\par
\noindent Source : 7I:427. \nolinkurl{fixes/MISPRINTS_7I.json#427}\par

\par\medskip\noindent\hypertarget{app-F7I-VIII-005.5}{}\textbf{F7I-VIII-005.5}\quad p. 224; printed misprint.\par
\noindent Imprimé : {\ttfamily\footnotesize 1\allowbreak{}}\par
\noindent Corrigé : {\ttfamily\footnotesize i\allowbreak{}}\par
\noindent Source : 7I:428. \nolinkurl{fixes/MISPRINTS_7I.json#428}\par

\par\medskip\noindent\hypertarget{app-F7I-VIII-006.1}{}\textbf{F7I-VIII-006.1}\quad p. 224; misprint\_pass.\par
\noindent Imprimé : {\ttfamily\footnotesize la\allowbreak{}}\par
\noindent Corrigé : {\ttfamily\footnotesize le\allowbreak{}}\par
\noindent Source : 7I:429. \nolinkurl{fixes/FIXLOG_7I_CURRENT.jsonl#782}\par

\par\medskip\noindent\hypertarget{app-F7I-VIII-007.1}{}\textbf{F7I-VIII-007.1}\quad p. 225; designated-witness attribution.\par
\noindent Imprimé : {\ttfamily\footnotesize VII1\allowbreak{}}\par
\noindent Corrigé : {\ttfamily\footnotesize VII\allowbreak{} 1\allowbreak{}}\par
\noindent Source : WITNESS-F7I-VIII-007.1. \nolinkurl{fr/sga7i/diplomatic/expose_VIII_body.tex}; \nolinkurl{fr/sga7i/current/expose_VIII_body.tex}\par

\par\medskip\noindent\hypertarget{app-F7I-VIII-007.2}{}\textbf{F7I-VIII-007.2}\quad p. 225; inherited corrected-layer difference.\par
\noindent Imprimé : \emph{(vide)}\par
\noindent Corrigé : \emph{(1 espace(s), 0 retour(s) de ligne dans le témoin TeX)}\par
\noindent Source : ADOPT-F7I-VIII-007.2. \nolinkurl{fr/sga7i/diplomatic/expose_VIII_body.tex}; \nolinkurl{fr/sga7i/current/expose_VIII_body.tex}\par

\par\medskip\noindent\hypertarget{app-F7I-VIII-007.3}{}\textbf{F7I-VIII-007.3}\quad p. 225; printed misprint.\par
\noindent Imprimé : {\ttfamily\footnotesize défini\allowbreak{}}\par
\noindent Corrigé : {\ttfamily\footnotesize définie\allowbreak{}}\par
\noindent Source : 7I:430. \nolinkurl{fixes/MISPRINTS_7I.json#430}\par

\par\medskip\noindent\hypertarget{app-F7I-VIII-007.4}{}\textbf{F7I-VIII-007.4}\quad p. 225; printed misprint.\par
\noindent Imprimé : {\ttfamily\footnotesize P\allowbreak{}}\par
\noindent Corrigé : {\ttfamily\footnotesize Q\textasciicircum{}*\allowbreak{}}\par
\noindent Source : 7I:431. \nolinkurl{fixes/MISPRINTS_7I.json#431}\par

\par\medskip\noindent\hypertarget{app-F7I-VIII-008.1}{}\textbf{F7I-VIII-008.1}\quad p. 225; misprint\_pass.\par
\noindent Imprimé : {\ttfamily\footnotesize Suppossons\allowbreak{}}\par
\noindent Corrigé : {\ttfamily\footnotesize Supposons\allowbreak{}}\par
\noindent Source : 7I:432. \nolinkurl{fixes/FIXLOG_7I_CURRENT.jsonl#783}\par

\par\medskip\noindent\hypertarget{app-F7I-VIII-008.2}{}\textbf{F7I-VIII-008.2}\quad p. 225; misprint\_pass.\par
\noindent Imprimé : \emph{(vide)}\par
\noindent Corrigé : {\ttfamily\footnotesize  et\allowbreak{} 1.5\allowbreak{}}\par
\noindent Source : 7I:433. \nolinkurl{fixes/FIXLOG_7I_CURRENT.jsonl#784}\par

\par\medskip\noindent\hypertarget{app-F7I-VIII-009.1}{}\textbf{F7I-VIII-009.1}\quad p. 226; printed misprint.\par
\noindent Imprimé : {\ttfamily\footnotesize G\allowbreak{}}\par
\noindent Corrigé : {\ttfamily\footnotesize Q\allowbreak{}}\par
\noindent Source : 7I:434. \nolinkurl{fixes/MISPRINTS_7I.json#434}\par

\par\medskip\noindent\hypertarget{app-F7I-VIII-009.2}{}\textbf{F7I-VIII-009.2}\quad p. 226; printed misprint.\par
\noindent Imprimé : {\ttfamily\footnotesize Q\allowbreak{}}\par
\noindent Corrigé : {\ttfamily\footnotesize G\allowbreak{}}\par
\noindent Source : 7I:435. \nolinkurl{fixes/MISPRINTS_7I.json#435}\par

\par\medskip\noindent\hypertarget{app-F7I-VIII-010.1}{}\textbf{F7I-VIII-010.1}\quad p. 227; misprint\_in\_print.\par
\noindent Imprimé : {\ttfamily\footnotesize Supposcns\allowbreak{}}\par
\noindent Corrigé : {\ttfamily\footnotesize Supposons\allowbreak{}}\par
\noindent Source : FIX-7I-D-0265. \nolinkurl{fixes/FIXLOG_7I.jsonl#265}\par

\par\medskip\noindent\hypertarget{app-F7I-VIII-011.1}{}\textbf{F7I-VIII-011.1}\quad p. 228; printed misprint.\par
\noindent Imprimé : {\ttfamily\footnotesize deu\allowbreak{} ième\allowbreak{}}\par
\noindent Corrigé : {\ttfamily\footnotesize deuxième\allowbreak{}}\par
\noindent Source : 7I:437. \nolinkurl{fixes/MISPRINTS_7I.json#437}\par

\par\medskip\noindent\hypertarget{app-F7I-VIII-011.2}{}\textbf{F7I-VIII-011.2}\quad p. 228; printed misprint.\par
\noindent Imprimé : {\ttfamily\footnotesize Q\allowbreak{}}\par
\noindent Corrigé : {\ttfamily\footnotesize P\allowbreak{}}\par
\noindent Source : 7I:438. \nolinkurl{fixes/MISPRINTS_7I.json#438}\par

\par\medskip\noindent\hypertarget{app-F7I-VIII-011.3}{}\textbf{F7I-VIII-011.3}\quad p. 228; misprint\_pass.\par
\noindent Imprimé : {\ttfamily\footnotesize quelconque\allowbreak{}}\par
\noindent Corrigé : {\ttfamily\footnotesize quelconques\allowbreak{}}\par
\noindent Source : 7I:439. \nolinkurl{fixes/FIXLOG_7I_CURRENT.jsonl#785}; \nolinkurl{fixes/MISPRINTS_7I.json#439}\par

\par\medskip\noindent\hypertarget{app-F7I-VIII-011.4}{}\textbf{F7I-VIII-011.4}\quad p. 228; misprint\_pass.\par
\noindent Imprimé : {\ttfamily\footnotesize et\allowbreak{}}\par
\noindent Corrigé : {\ttfamily\footnotesize de\allowbreak{}}\par
\noindent Source : 7I:440. \nolinkurl{fixes/FIXLOG_7I_CURRENT.jsonl#786}\par

\par\medskip\noindent\hypertarget{app-F7I-VIII-012.1}{}\textbf{F7I-VIII-012.1}\quad p. 229; designated-witness attribution.\par
\noindent Imprimé : \emph{(1 espace(s), 0 retour(s) de ligne dans le témoin TeX)}\par
\noindent Corrigé : {\ttfamily\footnotesize (\allowbreak{}}\par
\noindent Source : WITNESS-F7I-VIII-012.1. \nolinkurl{fr/sga7i/diplomatic/expose_VIII_body.tex}; \nolinkurl{fr/sga7i/current/expose_VIII_body.tex}\par

\par\medskip\noindent\hypertarget{app-F7I-VIII-012.2}{}\textbf{F7I-VIII-012.2}\quad p. 229; designated-witness attribution.\par
\noindent Imprimé : \emph{(1 espace(s), 0 retour(s) de ligne dans le témoin TeX)}\par
\noindent Corrigé : {\ttfamily\footnotesize (\allowbreak{}}\par
\noindent Source : WITNESS-F7I-VIII-012.2. \nolinkurl{fr/sga7i/diplomatic/expose_VIII_body.tex}; \nolinkurl{fr/sga7i/current/expose_VIII_body.tex}\par

\par\medskip\noindent\hypertarget{app-F7I-VIII-013.1}{}\textbf{F7I-VIII-013.1}\quad p. 232; printed misprint.\par
\noindent Imprimé : \emph{(vide)}\par
\noindent Corrigé : {\ttfamily\footnotesize 1.\allowbreak{}}\par
\noindent Source : 7I:443. \nolinkurl{fixes/MISPRINTS_7I.json#443}\par

\par\medskip\noindent\hypertarget{app-F7I-VIII-013.2}{}\textbf{F7I-VIII-013.2}\quad p. 232; printed misprint.\par
\noindent Imprimé : \emph{(vide)}\par
\noindent Corrigé : {\ttfamily\footnotesize .1\allowbreak{}}\par
\noindent Source : 7I:444. \nolinkurl{fixes/MISPRINTS_7I.json#444}\par

\par\medskip\noindent\hypertarget{app-F7I-VIII-013.3}{}\textbf{F7I-VIII-013.3}\quad p. 232; printed misprint.\par
\noindent Imprimé : \emph{(vide)}\par
\noindent Corrigé : {\ttfamily\footnotesize .1\allowbreak{}}\par
\noindent Source : 7I:445. \nolinkurl{fixes/MISPRINTS_7I.json#445}\par

\par\medskip\noindent\hypertarget{app-F7I-VIII-014.1}{}\textbf{F7I-VIII-014.1}\quad p. 233; misprint\_pass.\par
\noindent Imprimé : {\ttfamily\footnotesize sommetsupérieur\allowbreak{}}\par
\noindent Corrigé : {\ttfamily\footnotesize sommet\allowbreak{} supérieur\allowbreak{}}\par
\noindent Source : 7I:446. \nolinkurl{fixes/FIXLOG_7I_CURRENT.jsonl#787}\par

\par\medskip\noindent\hypertarget{app-F7I-VIII-015.1}{}\textbf{F7I-VIII-015.1}\quad p. 234; misprint\_pass.\par
\noindent Imprimé : {\ttfamily\footnotesize traités\allowbreak{}}\par
\noindent Corrigé : {\ttfamily\footnotesize traité\allowbreak{}}\par
\noindent Source : 7I:447. \nolinkurl{fixes/FIXLOG_7I_CURRENT.jsonl#788}\par

\par\medskip\noindent\hypertarget{app-F7I-VIII-016.1}{}\textbf{F7I-VIII-016.1}\quad p. 235; printed misprint.\par
\noindent Imprimé : {\ttfamily\footnotesize isomorphisme\allowbreak{}}\par
\noindent Corrigé : {\ttfamily\footnotesize homomorphisme\allowbreak{}}\par
\noindent Source : 7I:448. \nolinkurl{fixes/MISPRINTS_7I.json#448}\par

\par\medskip\noindent\hypertarget{app-F7I-VIII-017.1}{}\textbf{F7I-VIII-017.1}\quad p. 237; inherited corrected-layer difference.\par
\noindent Imprimé : \emph{(vide)}\par
\noindent Corrigé : {\ttfamily\footnotesize  pq\textbackslash{}\allowbreak{} }\par
\noindent Source : ADOPT-F7I-VIII-017.1. \nolinkurl{fr/sga7i/diplomatic/expose_VIII_body.tex}; \nolinkurl{fr/sga7i/current/expose_VIII_body.tex}\par

\par\medskip\noindent\hypertarget{app-F7I-VIII-018.1}{}\textbf{F7I-VIII-018.1}\quad p. 238; printed misprint.\par
\noindent Imprimé : {\ttfamily\footnotesize \textbackslash{}varphi\allowbreak{}}\par
\noindent Corrigé : {\ttfamily\footnotesize \textbackslash{}alpha\allowbreak{}}\par
\noindent Source : 7I:450. \nolinkurl{fixes/MISPRINTS_7I.json#450}\par

\par\medskip\noindent\hypertarget{app-F7I-VIII-018.2}{}\textbf{F7I-VIII-018.2}\quad p. 238; printed misprint.\par
\noindent Imprimé : {\ttfamily\footnotesize ;G\allowbreak{}}\par
\noindent Corrigé : \emph{(vide)}\par
\noindent Source : 7I:451. \nolinkurl{fixes/MISPRINTS_7I.json#451}\par

\par\medskip\noindent\hypertarget{app-F7I-VIII-018.3}{}\textbf{F7I-VIII-018.3}\quad p. 238; printed misprint.\par
\noindent Imprimé : {\ttfamily\footnotesize 4\allowbreak{}}\par
\noindent Corrigé : {\ttfamily\footnotesize 5\allowbreak{}}\par
\noindent Source : 7I:452. \nolinkurl{fixes/MISPRINTS_7I.json#452}\par

\par\medskip\noindent\hypertarget{app-F7I-VIII-019.1}{}\textbf{F7I-VIII-019.1}\quad p. 240; printed misprint.\par
\noindent Imprimé : \emph{(vide)}\par
\noindent Corrigé : {\ttfamily\footnotesize )\allowbreak{}}\par
\noindent Source : 7I:453. \nolinkurl{fixes/MISPRINTS_7I.json#453}\par

\par\medskip\noindent\hypertarget{app-F7I-VIII-020.1}{}\textbf{F7I-VIII-020.1}\quad p. 243; printed misprint.\par
\noindent Imprimé : {\ttfamily\footnotesize nQnG\allowbreak{}}\par
\noindent Corrigé : {\ttfamily\footnotesize nQ,G\allowbreak{}}\par
\noindent Source : 7I:454. \nolinkurl{fixes/MISPRINTS_7I.json#454}\par

\par\medskip\noindent\hypertarget{app-F7I-VIII-021.1}{}\textbf{F7I-VIII-021.1}\quad p. 243; printed misprint.\par
\noindent Imprimé : {\ttfamily\footnotesize 5\allowbreak{}}\par
\noindent Corrigé : {\ttfamily\footnotesize 3\allowbreak{}}\par
\noindent Source : 7I:455. \nolinkurl{fixes/MISPRINTS_7I.json#455}\par

\par\medskip\noindent\hypertarget{app-F7I-VIII-022.1}{}\textbf{F7I-VIII-022.1}\quad p. 245; printed misprint.\par
\noindent Imprimé : {\ttfamily\footnotesize \textbackslash{}to\allowbreak{}}\par
\noindent Corrigé : {\ttfamily\footnotesize -\allowbreak{}}\par
\noindent Source : 7I:456. \nolinkurl{fixes/MISPRINTS_7I.json#456}\par

\par\medskip\noindent\hypertarget{app-F7I-VIII-023.1}{}\textbf{F7I-VIII-023.1}\quad p. 246; misprint\_in\_print.\par
\noindent Imprimé : {\ttfamily\footnotesize calculéecomme\allowbreak{}}\par
\noindent Corrigé : {\ttfamily\footnotesize calculée\allowbreak{} comme\allowbreak{}}\par
\noindent Source : FIX-7I-D-0281. \nolinkurl{fixes/FIXLOG_7I.jsonl#281}\par

\par\medskip\noindent\hypertarget{app-F7I-VIII-023.2}{}\textbf{F7I-VIII-023.2}\quad p. 246; misprint\_pass.\par
\noindent Imprimé : {\ttfamily\footnotesize défini\allowbreak{}}\par
\noindent Corrigé : {\ttfamily\footnotesize définie\allowbreak{}}\par
\noindent Source : 7I:458. \nolinkurl{fixes/FIXLOG_7I_CURRENT.jsonl#789}\par

\par\medskip\noindent\hypertarget{app-F7I-VIII-024.1}{}\textbf{F7I-VIII-024.1}\quad p. 247; misprint\_pass.\par
\noindent Imprimé : {\ttfamily\footnotesize donné\allowbreak{}}\par
\noindent Corrigé : {\ttfamily\footnotesize donnée\allowbreak{}}\par
\noindent Source : 7I:460. \nolinkurl{fixes/FIXLOG_7I_CURRENT.jsonl#790}\par

\par\medskip\noindent\hypertarget{app-F7I-VIII-024.2}{}\textbf{F7I-VIII-024.2}\quad p. 247; printed misprint.\par
\noindent Imprimé : {\ttfamily\footnotesize 1\allowbreak{}}\par
\noindent Corrigé : {\ttfamily\footnotesize o\allowbreak{}}\par
\noindent Source : 7I:459. \nolinkurl{fixes/MISPRINTS_7I.json#459}\par

\par\medskip\noindent\hypertarget{app-F7I-VIII-025.1}{}\textbf{F7I-VIII-025.1}\quad p. 249; printed misprint.\par
\noindent Imprimé : {\ttfamily\footnotesize X\allowbreak{}}\par
\noindent Corrigé : {\ttfamily\footnotesize x\allowbreak{}}\par
\noindent Source : 7I:461. \nolinkurl{fixes/MISPRINTS_7I.json#461}\par

\par\medskip\noindent\hypertarget{app-F7I-VIII-025.2}{}\textbf{F7I-VIII-025.2}\quad p. 250; misprint\_pass.\par
\noindent Imprimé : \emph{(1 espace(s), 0 retour(s) de ligne dans le témoin TeX)}\par
\noindent Corrigé : {\ttfamily\footnotesize -\allowbreak{}}\par
\noindent Source : 7I:463. \nolinkurl{fixes/FIXLOG_7I_CURRENT.jsonl#792}\par

\par\medskip\noindent\hypertarget{app-F7I-VIII-025.3}{}\textbf{F7I-VIII-025.3}\quad p. 250; misprint\_pass.\par
\noindent Imprimé : {\ttfamily\footnotesize parle\allowbreak{}}\par
\noindent Corrigé : {\ttfamily\footnotesize par\allowbreak{} le\allowbreak{}}\par
\noindent Source : 7I:462. \nolinkurl{fixes/FIXLOG_7I_CURRENT.jsonl#791}\par

\par\medskip\noindent\hypertarget{app-F7I-VIII-026.1}{}\textbf{F7I-VIII-026.1}\quad p. 251; inherited corrected-layer difference.\par
\noindent Imprimé : \emph{(vide)}\par
\noindent Corrigé : {\ttfamily\footnotesize \_\textbackslash{}xi\allowbreak{}}\par
\noindent Source : ADOPT-F7I-VIII-026.1. \nolinkurl{fr/sga7i/diplomatic/expose_VIII_body.tex}; \nolinkurl{fr/sga7i/current/expose_VIII_body.tex}\par

\par\medskip\noindent\hypertarget{app-F7I-VIII-027.1}{}\textbf{F7I-VIII-027.1}\quad p. 251; inherited corrected-layer difference.\par
\noindent Imprimé : {\ttfamily\footnotesize o\allowbreak{}}\par
\noindent Corrigé : {\ttfamily\footnotesize 0\allowbreak{}}\par
\noindent Source : ADOPT-F7I-VIII-027.1. \nolinkurl{fr/sga7i/diplomatic/expose_VIII_body.tex}; \nolinkurl{fr/sga7i/current/expose_VIII_body.tex}\par

\par\medskip\noindent\hypertarget{app-F7I-VIII-028.1}{}\textbf{F7I-VIII-028.1}\quad p. 257; misprint\_pass.\par
\noindent Imprimé : {\ttfamily\footnotesize isomorphisme\allowbreak{}}\par
\noindent Corrigé : {\ttfamily\footnotesize isomorphe\allowbreak{}}\par
\noindent Source : 7I:464. \nolinkurl{fixes/FIXLOG_7I_CURRENT.jsonl#793}; \nolinkurl{fixes/MISPRINTS_7I.json#464}\par

\par\medskip\noindent\hypertarget{app-F7I-VIII-029.1}{}\textbf{F7I-VIII-029.1}\quad p. 258; printed misprint.\par
\noindent Imprimé : {\ttfamily\footnotesize L\allowbreak{}}\par
\noindent Corrigé : {\ttfamily\footnotesize T\allowbreak{}}\par
\noindent Source : 7I:465. \nolinkurl{fixes/MISPRINTS_7I.json#465}\par

\par\medskip\noindent\hypertarget{app-F7I-VIII-030.1}{}\textbf{F7I-VIII-030.1}\quad p. 260; printed misprint.\par
\noindent Imprimé : {\ttfamily\footnotesize P\allowbreak{}}\par
\noindent Corrigé : {\ttfamily\footnotesize E\allowbreak{}}\par
\noindent Source : 7I:466; 7I:803. \nolinkurl{fixes/MISPRINTS_7I.json#466}; \nolinkurl{fixes/MISPRINTS_7I.json#803}\par

\par\medskip\noindent\hypertarget{app-F7I-VIII-030.2}{}\textbf{F7I-VIII-030.2}\quad p. 260; misprint\_pass.\par
\noindent Imprimé : {\ttfamily\footnotesize nulle\allowbreak{}}\par
\noindent Corrigé : {\ttfamily\footnotesize nul\allowbreak{}}\par
\noindent Source : 7I:467; 7I:804. \nolinkurl{fixes/FIXLOG_7I_CURRENT.jsonl#794}; \nolinkurl{fixes/MISPRINTS_7I.json#467}; \nolinkurl{fixes/MISPRINTS_7I.json#804}\par

\par\medskip\noindent\hypertarget{app-F7I-VIII-031.1}{}\textbf{F7I-VIII-031.1}\quad p. 261; misprint\_pass.\par
\noindent Imprimé : {\ttfamily\footnotesize différent\allowbreak{}}\par
\noindent Corrigé : {\ttfamily\footnotesize diffèrent\allowbreak{}}\par
\noindent Source : 7I:468; 7I:805. \nolinkurl{fixes/FIXLOG_7I_CURRENT.jsonl#795}; \nolinkurl{fixes/MISPRINTS_7I.json#468}; \nolinkurl{fixes/MISPRINTS_7I.json#805}\par

\par\medskip\noindent\hypertarget{app-F7I-VIII-032.1}{}\textbf{F7I-VIII-032.1}\quad p. 263; printed misprint.\par
\noindent Imprimé : {\ttfamily\footnotesize 1\allowbreak{}}\par
\noindent Corrigé : {\ttfamily\footnotesize i\allowbreak{}}\par
\noindent Source : 7I:469; 7I:806. \nolinkurl{fixes/MISPRINTS_7I.json#469}; \nolinkurl{fixes/MISPRINTS_7I.json#806}\par

\par\medskip\noindent\hypertarget{app-F7I-VIII-033.1}{}\textbf{F7I-VIII-033.1}\quad p. 264; misprint\_pass.\par
\noindent Imprimé : {\ttfamily\footnotesize  été\allowbreak{}}\par
\noindent Corrigé : \emph{(vide)}\par
\noindent Source : 7I:470; 7I:807. \nolinkurl{fixes/FIXLOG_7I_CURRENT.jsonl#796}; \nolinkurl{fixes/MISPRINTS_7I.json#470}; \nolinkurl{fixes/MISPRINTS_7I.json#807}\par

\par\medskip\noindent\hypertarget{app-F7I-VIII-033.2}{}\textbf{F7I-VIII-033.2}\quad p. 264; printed misprint.\par
\noindent Imprimé : {\ttfamily\footnotesize \textbackslash{}mathbb\{Z\}\allowbreak{}}\par
\noindent Corrigé : {\ttfamily\footnotesize M\allowbreak{}}\par
\noindent Source : 7I:808. \nolinkurl{fixes/MISPRINTS_7I.json#808}\par

\par\medskip\noindent\hypertarget{app-F7I-VIII-034.1}{}\textbf{F7I-VIII-034.1}\quad p. 265; misprint\_pass.\par
\noindent Imprimé : {\ttfamily\footnotesize Soient\allowbreak{}}\par
\noindent Corrigé : {\ttfamily\footnotesize Soit\allowbreak{}}\par
\noindent Source : 7I:473; 7I:811. \nolinkurl{fixes/FIXLOG_7I_CURRENT.jsonl#797}; \nolinkurl{fixes/MISPRINTS_7I.json#473}; \nolinkurl{fixes/MISPRINTS_7I.json#811}\par

\par\medskip\noindent\hypertarget{app-F7I-VIII-034.2}{}\textbf{F7I-VIII-034.2}\quad p. 265; designated-witness attribution.\par
\noindent Imprimé : {\ttfamily\footnotesize P\allowbreak{}}\par
\noindent Corrigé : {\ttfamily\footnotesize L\allowbreak{}}\par
\noindent Source : WITNESS-F7I-VIII-034.2. \nolinkurl{fr/sga7i/diplomatic/expose_VIII_body.tex}; \nolinkurl{fr/sga7i/current/expose_VIII_body.tex}\par

\par\medskip\noindent\hypertarget{app-F7I-VIII-034.3}{}\textbf{F7I-VIII-034.3}\quad p. 265; designated-witness attribution.\par
\noindent Imprimé : {\ttfamily\footnotesize P\allowbreak{}}\par
\noindent Corrigé : {\ttfamily\footnotesize L\allowbreak{}}\par
\noindent Source : WITNESS-F7I-VIII-034.3. \nolinkurl{fr/sga7i/diplomatic/expose_VIII_body.tex}; \nolinkurl{fr/sga7i/current/expose_VIII_body.tex}\par

\par\medskip\noindent\hypertarget{app-F7I-VIII-034.4}{}\textbf{F7I-VIII-034.4}\quad p. 265; designated-witness attribution.\par
\noindent Imprimé : \emph{(vide)}\par
\noindent Corrigé : {\ttfamily\footnotesize \textbackslash{}otimes\allowbreak{}}\par
\noindent Source : WITNESS-F7I-VIII-034.4. \nolinkurl{fr/sga7i/diplomatic/expose_VIII_body.tex}; \nolinkurl{fr/sga7i/current/expose_VIII_body.tex}\par

\par\medskip\noindent\hypertarget{app-F7I-VIII-034.5}{}\textbf{F7I-VIII-034.5}\quad p. 265; designated-witness attribution.\par
\noindent Imprimé : {\ttfamily\footnotesize P\allowbreak{}}\par
\noindent Corrigé : {\ttfamily\footnotesize L\allowbreak{}}\par
\noindent Source : WITNESS-F7I-VIII-034.5. \nolinkurl{fr/sga7i/diplomatic/expose_VIII_body.tex}; \nolinkurl{fr/sga7i/current/expose_VIII_body.tex}\par

\par\medskip\noindent\hypertarget{app-F7I-VIII-035.1}{}\textbf{F7I-VIII-035.1}\quad p. 267; misprint\_pass.\par
\noindent Imprimé : {\ttfamily\footnotesize trivial\allowbreak{}}\par
\noindent Corrigé : {\ttfamily\footnotesize triviale\allowbreak{}}\par
\noindent Source : 7I:474; 7I:812. \nolinkurl{fixes/FIXLOG_7I_CURRENT.jsonl#798}; \nolinkurl{fixes/MISPRINTS_7I.json#474}; \nolinkurl{fixes/MISPRINTS_7I.json#812}\par

\par\medskip\noindent\hypertarget{app-F7I-VIII-036.1}{}\textbf{F7I-VIII-036.1}\quad p. 268; printed misprint.\par
\noindent Imprimé : {\ttfamily\footnotesize A\allowbreak{}}\par
\noindent Corrigé : {\ttfamily\footnotesize B\allowbreak{}}\par
\noindent Source : 7I:476; 7I:814. \nolinkurl{fixes/MISPRINTS_7I.json#476}; \nolinkurl{fixes/MISPRINTS_7I.json#814}\par

\par\medskip\noindent\hypertarget{app-F7I-VIII-036.2}{}\textbf{F7I-VIII-036.2}\quad p. 269; printed misprint.\par
\noindent Imprimé : {\ttfamily\footnotesize A\allowbreak{}}\par
\noindent Corrigé : {\ttfamily\footnotesize Q\allowbreak{}}\par
\noindent Source : 7I:477; 7I:815. \nolinkurl{fixes/MISPRINTS_7I.json#477}; \nolinkurl{fixes/MISPRINTS_7I.json#815}\par

\par\medskip\noindent\hypertarget{app-F7I-VIII-036.3}{}\textbf{F7I-VIII-036.3}\quad p. 269; misprint\_pass.\par
\noindent Imprimé : {\ttfamily\footnotesize 1\allowbreak{}}\par
\noindent Corrigé : {\ttfamily\footnotesize 4\allowbreak{}}\par
\noindent Source : 7I:817. \nolinkurl{fixes/FIXLOG_7I_CURRENT.jsonl#799}\par

\par\medskip\noindent\hypertarget{app-F7I-VIII-037.1}{}\textbf{F7I-VIII-037.1}\quad p. 270; misprint\_pass.\par
\noindent Imprimé : {\ttfamily\footnotesize hypothèses\allowbreak{}}\par
\noindent Corrigé : {\ttfamily\footnotesize hypothèse\allowbreak{}}\par
\noindent Source : 7I:478. \nolinkurl{fixes/FIXLOG_7I_CURRENT.jsonl#800}\par

\par\medskip\noindent\hypertarget{app-F7I-VIII-038.1}{}\textbf{F7I-VIII-038.1}\quad p. 271; misprint\_pass.\par
\noindent Imprimé : {\ttfamily\footnotesize naissanee\allowbreak{}}\par
\noindent Corrigé : {\ttfamily\footnotesize naissance\allowbreak{}}\par
\noindent Source : 7I:481. \nolinkurl{fixes/FIXLOG_7I_CURRENT.jsonl#801}\par

\par\medskip\noindent\hypertarget{app-F7I-VIII-039.1}{}\textbf{F7I-VIII-039.1}\quad p. 271; misprint\_pass.\par
\noindent Imprimé : {\ttfamily\footnotesize CeIa\allowbreak{}}\par
\noindent Corrigé : {\ttfamily\footnotesize Cela\allowbreak{}}\par
\noindent Source : 7I:482. \nolinkurl{fixes/FIXLOG_7I_CURRENT.jsonl#802}\par

\par\medskip\noindent\hypertarget{app-F7I-VIII-040.1}{}\textbf{F7I-VIII-040.1}\quad p. 274; printed misprint.\par
\noindent Imprimé : {\ttfamily\footnotesize 1\allowbreak{}}\par
\noindent Corrigé : {\ttfamily\footnotesize 4\allowbreak{}}\par
\noindent Source : 7I:483. \nolinkurl{fixes/MISPRINTS_7I.json#483}\par

\par\medskip\noindent\hypertarget{app-F7I-VIII-040.2}{}\textbf{F7I-VIII-040.2}\quad p. 274; inherited corrected-layer difference.\par
\noindent Imprimé : {\ttfamily\footnotesize ,\allowbreak{}}\par
\noindent Corrigé : {\ttfamily\footnotesize .\allowbreak{}}\par
\noindent Source : ADOPT-F7I-VIII-040.2. \nolinkurl{fr/sga7i/diplomatic/expose_VIII_body.tex}; \nolinkurl{fr/sga7i/current/expose_VIII_body.tex}\par

\par\medskip\noindent\hypertarget{app-F7I-VIII-041.1}{}\textbf{F7I-VIII-041.1}\quad p. 276; printed misprint.\par
\noindent Imprimé : {\ttfamily\footnotesize A\allowbreak{}}\par
\noindent Corrigé : {\ttfamily\footnotesize P\allowbreak{}}\par
\noindent Source : 7I:486; 7I:1045. \nolinkurl{fixes/MISPRINTS_7I.json#486}; \nolinkurl{fixes/MISPRINTS_7I.json#1045}\par

\par\medskip\noindent\hypertarget{app-F7I-VIII-041.2}{}\textbf{F7I-VIII-041.2}\quad p. 276; printed misprint.\par
\noindent Imprimé : {\ttfamily\footnotesize P\allowbreak{}}\par
\noindent Corrigé : {\ttfamily\footnotesize A\allowbreak{}}\par
\noindent Source : 7I:486; 7I:1045. \nolinkurl{fixes/MISPRINTS_7I.json#486}; \nolinkurl{fixes/MISPRINTS_7I.json#1045}\par

\par\medskip\noindent\hypertarget{app-F7I-VIII-041.3}{}\textbf{F7I-VIII-041.3}\quad p. 276; printed misprint.\par
\noindent Imprimé : {\ttfamily\footnotesize \textasciicircum{}\{*\}\allowbreak{}}\par
\noindent Corrigé : \emph{(vide)}\par
\noindent Source : 7I:1047. \nolinkurl{fixes/MISPRINTS_7I.json#1047}\par

\par\medskip\noindent\hypertarget{app-F7I-VIII-041.4}{}\textbf{F7I-VIII-041.4}\quad p. 276; printed misprint.\par
\noindent Imprimé : {\ttfamily\footnotesize \textbackslash{}underline\{M\}\allowbreak{}}\par
\noindent Corrigé : {\ttfamily\footnotesize A\allowbreak{}}\par
\noindent Source : 7I:1046. \nolinkurl{fixes/MISPRINTS_7I.json#1046}\par

\par\medskip\noindent\hypertarget{app-F7I-VIII-042.1}{}\textbf{F7I-VIII-042.1}\quad p. 277; misprint\_pass.\par
\noindent Imprimé : \emph{(vide)}\par
\noindent Corrigé : {\ttfamily\footnotesize )\allowbreak{}}\par
\noindent Source : 7I:487; 7I:1048. \nolinkurl{fixes/FIXLOG_7I_CURRENT.jsonl#803}; \nolinkurl{fixes/MISPRINTS_7I.json#487}; \nolinkurl{fixes/MISPRINTS_7I.json#1048}\par

\par\medskip\noindent\hypertarget{app-F7I-VIII-043.1}{}\textbf{F7I-VIII-043.1}\quad p. 278; inherited corrected-layer difference.\par
\noindent Imprimé : {\ttfamily\footnotesize 1\allowbreak{}}\par
\noindent Corrigé : {\ttfamily\footnotesize 2\allowbreak{}}\par
\noindent Source : ADOPT-F7I-VIII-043.1. \nolinkurl{fr/sga7i/diplomatic/expose_VIII_body.tex}; \nolinkurl{fr/sga7i/current/expose_VIII_body.tex}\par

\par\medskip\noindent\hypertarget{app-F7I-VIII-044.1}{}\textbf{F7I-VIII-044.1}\quad p. 278; printed misprint.\par
\noindent Imprimé : {\ttfamily\footnotesize G\allowbreak{}}\par
\noindent Corrigé : {\ttfamily\footnotesize Z\allowbreak{}}\par
\noindent Source : 7I:488. \nolinkurl{fixes/MISPRINTS_7I.json#488}\par

\par\medskip\noindent\hypertarget{app-F7I-VIII-044.2}{}\textbf{F7I-VIII-044.2}\quad p. 279; misprint\_pass.\par
\noindent Imprimé : {\ttfamily\footnotesize fidèlité\allowbreak{}}\par
\noindent Corrigé : {\ttfamily\footnotesize fidélité\allowbreak{}}\par
\noindent Source : 7I:490. \nolinkurl{fixes/FIXLOG_7I_CURRENT.jsonl#804}; \nolinkurl{fixes/MISPRINTS_7I.json#490}\par

\par\medskip\noindent\hypertarget{app-F7I-VIII-045.1}{}\textbf{F7I-VIII-045.1}\quad p. 279; misprint\_pass.\par
\noindent Imprimé : {\ttfamily\footnotesize )\allowbreak{}}\par
\noindent Corrigé : \emph{(vide)}\par
\noindent Source : 7I:491; 7I:1051. \nolinkurl{fixes/FIXLOG_7I_CURRENT.jsonl#805}; \nolinkurl{fixes/MISPRINTS_7I.json#491}; \nolinkurl{fixes/MISPRINTS_7I.json#1051}\par

\par\medskip\noindent\hypertarget{app-F7I-VIII-045.2}{}\textbf{F7I-VIII-045.2}\quad p. 279; printed misprint.\par
\noindent Imprimé : {\ttfamily\footnotesize E\allowbreak{}}\par
\noindent Corrigé : {\ttfamily\footnotesize Z\allowbreak{}}\par
\noindent Source : 7I:492; 7I:1052. \nolinkurl{fixes/MISPRINTS_7I.json#492}; \nolinkurl{fixes/MISPRINTS_7I.json#1052}\par

\par\medskip\noindent\hypertarget{app-F7I-VIII-045.3}{}\textbf{F7I-VIII-045.3}\quad p. 279; printed misprint.\par
\noindent Imprimé : {\ttfamily\footnotesize \_\textbackslash{}eta\allowbreak{}}\par
\noindent Corrigé : \emph{(vide)}\par
\noindent Source : 7I:493. \nolinkurl{fixes/MISPRINTS_7I.json#493}\par

\par\medskip\noindent\hypertarget{app-F7I-VIII-045.4}{}\textbf{F7I-VIII-045.4}\quad p. 279; printed misprint.\par
\noindent Imprimé : {\ttfamily\footnotesize \{\allowbreak{}}\par
\noindent Corrigé : \emph{(vide)}\par
\noindent Source : 7I:493. \nolinkurl{fixes/MISPRINTS_7I.json#493}\par

\par\medskip\noindent\hypertarget{app-F7I-VIII-045.5}{}\textbf{F7I-VIII-045.5}\quad p. 279; printed misprint.\par
\noindent Imprimé : {\ttfamily\footnotesize \_\textbackslash{}eta\}\allowbreak{}}\par
\noindent Corrigé : \emph{(vide)}\par
\noindent Source : 7I:493. \nolinkurl{fixes/MISPRINTS_7I.json#493}\par

\par\medskip\noindent\hypertarget{app-F7I-VIII-045.6}{}\textbf{F7I-VIII-045.6}\quad p. 280; printed misprint.\par
\noindent Imprimé : \emph{(vide)}\par
\noindent Corrigé : {\ttfamily\footnotesize \{\allowbreak{}}\par
\noindent Source : 7I:494. \nolinkurl{fixes/MISPRINTS_7I.json#494}\par

\par\medskip\noindent\hypertarget{app-F7I-VIII-045.7}{}\textbf{F7I-VIII-045.7}\quad p. 280; printed misprint.\par
\noindent Imprimé : \emph{(vide)}\par
\noindent Corrigé : {\ttfamily\footnotesize \textbackslash{}eta\}\allowbreak{}}\par
\noindent Source : 7I:494. \nolinkurl{fixes/MISPRINTS_7I.json#494}\par

\par\medskip\noindent\hypertarget{app-F7I-VIII-046.1}{}\textbf{F7I-VIII-046.1}\quad p. 282; printed misprint.\par
\noindent Imprimé : {\ttfamily\footnotesize Z\allowbreak{}}\par
\noindent Corrigé : {\ttfamily\footnotesize P\allowbreak{}}\par
\noindent Source : 7I:496; 7I:1064. \nolinkurl{fixes/MISPRINTS_7I.json#496}; \nolinkurl{fixes/MISPRINTS_7I.json#1064}\par

\par\medskip\noindent\hypertarget{app-F7I-VIII-046.2}{}\textbf{F7I-VIII-046.2}\quad p. 282; printed misprint.\par
\noindent Imprimé : {\ttfamily\footnotesize G\allowbreak{}}\par
\noindent Corrigé : {\ttfamily\footnotesize Z\allowbreak{}}\par
\noindent Source : 7I:495; 7I:1063. \nolinkurl{fixes/MISPRINTS_7I.json#495}; \nolinkurl{fixes/MISPRINTS_7I.json#1063}\par

\par\medskip\noindent\hypertarget{app-F7I-VIII-047.1}{}\textbf{F7I-VIII-047.1}\quad p. 283; misprint\_pass.\par
\noindent Imprimé : {\ttfamily\footnotesize rappelera\allowbreak{}}\par
\noindent Corrigé : {\ttfamily\footnotesize rappellera\allowbreak{}}\par
\noindent Source : 7I:497. \nolinkurl{fixes/FIXLOG_7I_CURRENT.jsonl#806}\par

\par\medskip\noindent\hypertarget{app-F7I-VIII-047.2}{}\textbf{F7I-VIII-047.2}\quad p. 283; misprint\_pass.\par
\noindent Imprimé : {\ttfamily\footnotesize 3\allowbreak{}}\par
\noindent Corrigé : {\ttfamily\footnotesize 2\allowbreak{}}\par
\noindent Source : 7I:1066. \nolinkurl{fixes/FIXLOG_7I_CURRENT.jsonl#1027}\par

\par\medskip\noindent\hypertarget{app-F7I-VIII-047.3}{}\textbf{F7I-VIII-047.3}\quad p. 283; misprint\_pass.\par
\noindent Imprimé : \emph{(1 espace(s), 0 retour(s) de ligne dans le témoin TeX)}\par
\noindent Corrigé : {\ttfamily\footnotesize -\allowbreak{}}\par
\noindent Source : 7I:1067. \nolinkurl{fixes/FIXLOG_7I_CURRENT.jsonl#1028}; \nolinkurl{fixes/MISPRINTS_7I.json#1067}\par

\par\medskip\noindent\hypertarget{app-F7I-VIII-047.4}{}\textbf{F7I-VIII-047.4}\quad p. 283; printed misprint.\par
\noindent Imprimé : {\ttfamily\footnotesize G\allowbreak{}}\par
\noindent Corrigé : {\ttfamily\footnotesize Z\allowbreak{}}\par
\noindent Source : 7I:498; 7I:1068. \nolinkurl{fixes/MISPRINTS_7I.json#498}; \nolinkurl{fixes/MISPRINTS_7I.json#1068}\par

\par\medskip\noindent\hypertarget{app-F7I-VIII-047.5}{}\textbf{F7I-VIII-047.5}\quad p. 283; printed misprint.\par
\noindent Imprimé : {\ttfamily\footnotesize G\allowbreak{}}\par
\noindent Corrigé : {\ttfamily\footnotesize Z\allowbreak{}}\par
\noindent Source : 7I:499; 7I:1069. \nolinkurl{fixes/MISPRINTS_7I.json#499}; \nolinkurl{fixes/MISPRINTS_7I.json#1069}\par

\par\medskip\noindent\hypertarget{app-F7I-VIII-047.6}{}\textbf{F7I-VIII-047.6}\quad p. 283; printed misprint.\par
\noindent Imprimé : {\ttfamily\footnotesize G\allowbreak{}}\par
\noindent Corrigé : {\ttfamily\footnotesize Z\allowbreak{}}\par
\noindent Source : 7I:499; 7I:1069. \nolinkurl{fixes/MISPRINTS_7I.json#499}; \nolinkurl{fixes/MISPRINTS_7I.json#1069}\par

\par\medskip\noindent\hypertarget{app-F7I-VIII-047.7}{}\textbf{F7I-VIII-047.7}\quad p. 283; printed misprint.\par
\noindent Imprimé : {\ttfamily\footnotesize G\allowbreak{}}\par
\noindent Corrigé : {\ttfamily\footnotesize Z\allowbreak{}}\par
\noindent Source : 7I:500; 7I:1070. \nolinkurl{fixes/MISPRINTS_7I.json#500}; \nolinkurl{fixes/MISPRINTS_7I.json#1070}\par

\par\medskip\noindent\hypertarget{app-F7I-VIII-048.1}{}\textbf{F7I-VIII-048.1}\quad p. 283; printed misprint.\par
\noindent Imprimé : {\ttfamily\footnotesize schémasen\allowbreak{}}\par
\noindent Corrigé : {\ttfamily\footnotesize schémas\allowbreak{} en\allowbreak{}}\par
\noindent Source : 7I:1071. \nolinkurl{fixes/MISPRINTS_7I.json#1071}\par

\par\medskip\noindent\hypertarget{app-F7I-VIII-048.2}{}\textbf{F7I-VIII-048.2}\quad p. 283; printed misprint.\par
\noindent Imprimé : {\ttfamily\footnotesize neutres\allowbreak{}}\par
\noindent Corrigé : {\ttfamily\footnotesize connexes\allowbreak{}}\par
\noindent Source : 7I:502; 7I:1072. \nolinkurl{fixes/MISPRINTS_7I.json#502}; \nolinkurl{fixes/MISPRINTS_7I.json#1072}\par

\par\medskip\noindent\hypertarget{app-F7I-VIII-049.1}{}\textbf{F7I-VIII-049.1}\quad p. 284; printed misprint.\par
\noindent Imprimé : {\ttfamily\footnotesize G\allowbreak{}}\par
\noindent Corrigé : {\ttfamily\footnotesize Z\allowbreak{}}\par
\noindent Source : 7I:503; 7I:1073. \nolinkurl{fixes/MISPRINTS_7I.json#503}; \nolinkurl{fixes/MISPRINTS_7I.json#1073}\par

\par\medskip\noindent\hypertarget{app-F7I-VIII-050.1}{}\textbf{F7I-VIII-050.1}\quad p. 285; printed misprint.\par
\noindent Imprimé : {\ttfamily\footnotesize 2\allowbreak{}}\par
\noindent Corrigé : {\ttfamily\footnotesize 5\allowbreak{}}\par
\noindent Source : 7I:1074. \nolinkurl{fixes/MISPRINTS_7I.json#1074}\par

\par\medskip\noindent\hypertarget{app-F7I-VIII-051.1}{}\textbf{F7I-VIII-051.1}\quad p. 285; printed misprint.\par
\noindent Imprimé : {\ttfamily\footnotesize G\allowbreak{}}\par
\noindent Corrigé : {\ttfamily\footnotesize Z\allowbreak{}}\par
\noindent Source : 7I:1075. \nolinkurl{fixes/MISPRINTS_7I.json#1075}\par

\par\medskip\noindent\hypertarget{app-F7I-VIII-051.2}{}\textbf{F7I-VIII-051.2}\quad p. 286; inherited corrected-layer difference.\par
\noindent Imprimé : {\ttfamily\footnotesize G\allowbreak{}}\par
\noindent Corrigé : {\ttfamily\footnotesize Z\allowbreak{}}\par
\noindent Source : ADOPT-F7I-VIII-051.2. \nolinkurl{fr/sga7i/diplomatic/expose_VIII_body.tex}; \nolinkurl{fr/sga7i/current/expose_VIII_body.tex}\par

\par\medskip\noindent\hypertarget{app-F7I-VIII-052.1}{}\textbf{F7I-VIII-052.1}\quad p. 286; inherited corrected-layer difference.\par
\noindent Imprimé : {\ttfamily\footnotesize G\allowbreak{}}\par
\noindent Corrigé : {\ttfamily\footnotesize Z\allowbreak{}}\par
\noindent Source : ADOPT-F7I-VIII-052.1. \nolinkurl{fr/sga7i/diplomatic/expose_VIII_body.tex}; \nolinkurl{fr/sga7i/current/expose_VIII_body.tex}\par

\par\medskip\noindent\hypertarget{app-F7I-VIII-052.2}{}\textbf{F7I-VIII-052.2}\quad p. 286; inherited corrected-layer difference.\par
\noindent Imprimé : {\ttfamily\footnotesize G\allowbreak{}}\par
\noindent Corrigé : {\ttfamily\footnotesize Z\allowbreak{}}\par
\noindent Source : ADOPT-F7I-VIII-052.2. \nolinkurl{fr/sga7i/diplomatic/expose_VIII_body.tex}; \nolinkurl{fr/sga7i/current/expose_VIII_body.tex}\par

\par\medskip\noindent\hypertarget{app-F7I-VIII-053.1}{}\textbf{F7I-VIII-053.1}\quad p. 287; misprint\_pass.\par
\noindent Imprimé : {\ttfamily\footnotesize gra\textbackslash{}\textasciicircum{}\{c\}e\allowbreak{}}\par
\noindent Corrigé : {\ttfamily\footnotesize grâce\allowbreak{}}\par
\noindent Source : 7I:506. \nolinkurl{fixes/FIXLOG_7I_CURRENT.jsonl#807}\par

\par\medskip\noindent\hypertarget{app-F7I-VIII-054.1}{}\textbf{F7I-VIII-054.1}\quad p. 288; printed misprint.\par
\noindent Imprimé : {\ttfamily\footnotesize Z\allowbreak{}}\par
\noindent Corrigé : {\ttfamily\footnotesize S\allowbreak{}}\par
\noindent Source : 7I:507; 7I:1056. \nolinkurl{fixes/MISPRINTS_7I.json#507}; \nolinkurl{fixes/MISPRINTS_7I.json#1056}\par

\par\medskip\noindent\hypertarget{app-F7I-VIII-054.2}{}\textbf{F7I-VIII-054.2}\quad p. 288; printed misprint.\par
\noindent Imprimé : {\ttfamily\footnotesize \{\}\_\{\textbackslash{}overline\{G\}\}\allowbreak{}}\par
\noindent Corrigé : \emph{(vide)}\par
\noindent Source : 7I:1057. \nolinkurl{fixes/MISPRINTS_7I.json#1057}\par

\par\medskip\noindent\hypertarget{app-F7I-VIII-055.1}{}\textbf{F7I-VIII-055.1}\quad p. 289; misprint\_pass.\par
\noindent Imprimé : {\ttfamily\footnotesize définie\allowbreak{}}\par
\noindent Corrigé : {\ttfamily\footnotesize défini\allowbreak{}}\par
\noindent Source : 7I:1058. \nolinkurl{fixes/FIXLOG_7I_CURRENT.jsonl#1029}\par

\par\medskip\noindent\hypertarget{app-F7I-VIII-055.2}{}\textbf{F7I-VIII-055.2}\quad p. 289; misprint\_pass.\par
\noindent Imprimé : {\ttfamily\footnotesize définie\allowbreak{}}\par
\noindent Corrigé : {\ttfamily\footnotesize défini\allowbreak{}}\par
\noindent Source : 7I:1059. \nolinkurl{fixes/FIXLOG_7I_CURRENT.jsonl#1030}; \nolinkurl{fixes/MISPRINTS_7I.json#1059}\par

\par\medskip\noindent\hypertarget{app-F7I-VIII-056.1}{}\textbf{F7I-VIII-056.1}\quad p. 289; misprint\_pass.\par
\noindent Imprimé : \emph{(vide)}\par
\noindent Corrigé : {\ttfamily\footnotesize \textbackslash{}overline\{\allowbreak{}}\par
\noindent Source : 7I:508. \nolinkurl{fixes/FIXLOG_7I_CURRENT.jsonl#808}; \nolinkurl{fixes/MISPRINTS_7I.json#508}\par

\par\medskip\noindent\hypertarget{app-F7I-VIII-056.2}{}\textbf{F7I-VIII-056.2}\quad p. 289; misprint\_pass.\par
\noindent Imprimé : \emph{(vide)}\par
\noindent Corrigé : {\ttfamily\footnotesize \}\allowbreak{}}\par
\noindent Source : 7I:508. \nolinkurl{fixes/FIXLOG_7I_CURRENT.jsonl#808}; \nolinkurl{fixes/MISPRINTS_7I.json#508}\par

\par\medskip\noindent\hypertarget{app-F7I-VIII-056.3}{}\textbf{F7I-VIII-056.3}\quad p. 289; misprint\_pass.\par
\noindent Imprimé : \emph{(vide)}\par
\noindent Corrigé : {\ttfamily\footnotesize \textbackslash{}overline\{\allowbreak{}}\par
\noindent Source : 7I:508. \nolinkurl{fixes/FIXLOG_7I_CURRENT.jsonl#808}; \nolinkurl{fixes/MISPRINTS_7I.json#508}\par

\par\medskip\noindent\hypertarget{app-F7I-VIII-056.4}{}\textbf{F7I-VIII-056.4}\quad p. 289; misprint\_pass.\par
\noindent Imprimé : \emph{(vide)}\par
\noindent Corrigé : {\ttfamily\footnotesize \}\allowbreak{}}\par
\noindent Source : 7I:508. \nolinkurl{fixes/FIXLOG_7I_CURRENT.jsonl#808}; \nolinkurl{fixes/MISPRINTS_7I.json#508}\par

\par\medskip\noindent\hypertarget{app-F7I-VIII-056.5}{}\textbf{F7I-VIII-056.5}\quad p. 289; printed misprint.\par
\noindent Imprimé : {\ttfamily\footnotesize 3\allowbreak{}}\par
\noindent Corrigé : {\ttfamily\footnotesize 6\allowbreak{}}\par
\noindent Source : 7I:1061. \nolinkurl{fixes/MISPRINTS_7I.json#1061}\par

\par\medskip\noindent\hypertarget{app-F7I-VIII-056.6}{}\textbf{F7I-VIII-056.6}\quad p. 290; misprint\_pass.\par
\noindent Imprimé : {\ttfamily\footnotesize léments\allowbreak{}}\par
\noindent Corrigé : {\ttfamily\footnotesize éléments\allowbreak{}}\par
\noindent Source : 7I:510. \nolinkurl{fixes/FIXLOG_7I_CURRENT.jsonl#809}\par

\par\medskip\noindent\hypertarget{app-F7I-VIII-057.1}{}\textbf{F7I-VIII-057.1}\quad p. 291; printed misprint.\par
\noindent Imprimé : {\ttfamily\footnotesize \textbackslash{}underline\allowbreak{}}\par
\noindent Corrigé : {\ttfamily\footnotesize \textbackslash{}overline\allowbreak{}}\par
\noindent Source : 7I:1076. \nolinkurl{fixes/MISPRINTS_7I.json#1076}\par

\par\medskip\noindent\hypertarget{app-F7I-VIII-058.1}{}\textbf{F7I-VIII-058.1}\quad p. 291; printed misprint.\par
\noindent Imprimé : \emph{(vide)}\par
\noindent Corrigé : {\ttfamily\footnotesize \{-\allowbreak{}}\par
\noindent Source : 7I:1077. \nolinkurl{fixes/MISPRINTS_7I.json#1077}\par

\par\medskip\noindent\hypertarget{app-F7I-VIII-058.2}{}\textbf{F7I-VIII-058.2}\quad p. 291; printed misprint.\par
\noindent Imprimé : \emph{(vide)}\par
\noindent Corrigé : {\ttfamily\footnotesize \}\allowbreak{}}\par
\noindent Source : 7I:1077. \nolinkurl{fixes/MISPRINTS_7I.json#1077}\par

\par\medskip\noindent\hypertarget{app-F7I-VIII-058.3}{}\textbf{F7I-VIII-058.3}\quad p. 291; printed misprint.\par
\noindent Imprimé : \emph{(vide)}\par
\noindent Corrigé : {\ttfamily\footnotesize \textbackslash{}overline\{\allowbreak{}}\par
\noindent Source : 7I:1078. \nolinkurl{fixes/MISPRINTS_7I.json#1078}\par

\par\medskip\noindent\hypertarget{app-F7I-VIII-058.4}{}\textbf{F7I-VIII-058.4}\quad p. 291; printed misprint.\par
\noindent Imprimé : \emph{(vide)}\par
\noindent Corrigé : {\ttfamily\footnotesize \}\allowbreak{}}\par
\noindent Source : 7I:1078. \nolinkurl{fixes/MISPRINTS_7I.json#1078}\par

\par\medskip\noindent\hypertarget{app-F7I-VIII-058.5}{}\textbf{F7I-VIII-058.5}\quad p. 291; misprint\_pass.\par
\noindent Imprimé : {\ttfamily\footnotesize appelerons\allowbreak{}}\par
\noindent Corrigé : {\ttfamily\footnotesize appellerons\allowbreak{}}\par
\noindent Source : 7I:513; 7I:1079. \nolinkurl{fixes/FIXLOG_7I_CURRENT.jsonl#810}; \nolinkurl{fixes/MISPRINTS_7I.json#513}; \nolinkurl{fixes/MISPRINTS_7I.json#1079}\par

\par\medskip\noindent\hypertarget{app-F7I-VIII-059.1}{}\textbf{F7I-VIII-059.1}\quad p. 292; misprint\_pass.\par
\noindent Imprimé : {\ttfamily\footnotesize égal\allowbreak{}}\par
\noindent Corrigé : {\ttfamily\footnotesize égale\allowbreak{}}\par
\noindent Source : 7I:516; 7I:1080. \nolinkurl{fixes/FIXLOG_7I_CURRENT.jsonl#811}; \nolinkurl{fixes/MISPRINTS_7I.json#1080}\par

\par\medskip\noindent\hypertarget{app-F7I-VIII-060.1}{}\textbf{F7I-VIII-060.1}\quad p. 293; misprint\_pass.\par
\noindent Imprimé : {\ttfamily\footnotesize ces\allowbreak{}}\par
\noindent Corrigé : {\ttfamily\footnotesize ses\allowbreak{}}\par
\noindent Source : 7I:517; 7I:1081. \nolinkurl{fixes/FIXLOG_7I_CURRENT.jsonl#812}; \nolinkurl{fixes/MISPRINTS_7I.json#1081}\par

\par\medskip\noindent\hypertarget{app-F7I-VIII-061.1}{}\textbf{F7I-VIII-061.1}\quad p. 294; printed misprint.\par
\noindent Imprimé : {\ttfamily\footnotesize X\allowbreak{}}\par
\noindent Corrigé : {\ttfamily\footnotesize S\allowbreak{}}\par
\noindent Source : 7I:1082. \nolinkurl{fixes/MISPRINTS_7I.json#1082}\par

\par\medskip\noindent\hypertarget{app-F7I-VIII-062.1}{}\textbf{F7I-VIII-062.1}\quad p. 295; misprint\_pass.\par
\noindent Imprimé : \emph{(1 espace(s), 0 retour(s) de ligne dans le témoin TeX)}\par
\noindent Corrigé : {\ttfamily\footnotesize -\allowbreak{}}\par
\noindent Source : 7I:1083. \nolinkurl{fixes/FIXLOG_7I_CURRENT.jsonl#1031}\par

\par\medskip\noindent\hypertarget{app-F7I-VIII-063.1}{}\textbf{F7I-VIII-063.1}\quad p. 297; printed misprint.\par
\noindent Imprimé : {\ttfamily\footnotesize V\allowbreak{}}\par
\noindent Corrigé : {\ttfamily\footnotesize U\allowbreak{}}\par
\noindent Source : 7I:519; 7I:1089. \nolinkurl{fixes/MISPRINTS_7I.json#519}; \nolinkurl{fixes/MISPRINTS_7I.json#1089}\par

\par\medskip\noindent\hypertarget{app-F7I-VIII-063.2}{}\textbf{F7I-VIII-063.2}\quad p. 297; printed misprint.\par
\noindent Imprimé : {\ttfamily\footnotesize mS\allowbreak{}}\par
\noindent Corrigé : {\ttfamily\footnotesize mU\allowbreak{}}\par
\noindent Source : 7I:520; 7I:1090. \nolinkurl{fixes/MISPRINTS_7I.json#520}; \nolinkurl{fixes/MISPRINTS_7I.json#1090}\par

\par\medskip\noindent\hypertarget{app-F7I-VIII-064.1}{}\textbf{F7I-VIII-064.1}\quad p. 298; inherited corrected-layer difference.\par
\noindent Imprimé : \emph{(0 espace(s), 1 retour(s) de ligne dans le témoin TeX)}\par
\noindent Corrigé : \emph{(0 espace(s), 2 retour(s) de ligne dans le témoin TeX)}\par
\noindent Source : ADOPT-F7I-VIII-064.1. \nolinkurl{fr/sga7i/diplomatic/expose_VIII_body.tex}; \nolinkurl{fr/sga7i/current/expose_VIII_body.tex}\par

\par\medskip\noindent\hypertarget{app-F7I-VIII-064.2}{}\textbf{F7I-VIII-064.2}\quad p. 298; inherited corrected-layer difference.\par
\noindent Imprimé : \emph{(vide)}\par
\noindent Corrigé : {\ttfamily\footnotesize si\allowbreak{} \$i:Y\textbackslash{}hookrightarrow\allowbreak{} X\$\allowbreak{} désigne\allowbreak{} l'immersion\allowbreak{} fermée,\allowbreak{} }\par
\noindent Source : ADOPT-F7I-VIII-064.2. \nolinkurl{fr/sga7i/diplomatic/expose_VIII_body.tex}; \nolinkurl{fr/sga7i/current/expose_VIII_body.tex}\par

\par\medskip\noindent\hypertarget{app-F7I-VIII-064.3}{}\textbf{F7I-VIII-064.3}\quad p. 298; inherited corrected-layer difference.\par
\noindent Imprimé : \emph{(vide)}\par
\noindent Corrigé : {\ttfamily\footnotesize  si\allowbreak{}}\par
\noindent Source : ADOPT-F7I-VIII-064.3. \nolinkurl{fr/sga7i/diplomatic/expose_VIII_body.tex}; \nolinkurl{fr/sga7i/current/expose_VIII_body.tex}\par

\par\medskip\noindent\hypertarget{app-F7I-VIII-064.4}{}\textbf{F7I-VIII-064.4}\quad p. 298; inherited corrected-layer difference.\par
\noindent Imprimé : \emph{(vide)}\par
\noindent Corrigé : {\ttfamily\footnotesize  est\allowbreak{}}\par
\noindent Source : ADOPT-F7I-VIII-064.4. \nolinkurl{fr/sga7i/diplomatic/expose_VIII_body.tex}; \nolinkurl{fr/sga7i/current/expose_VIII_body.tex}\par

\par\medskip\noindent\hypertarget{app-F7I-VIII-065.1}{}\textbf{F7I-VIII-065.1}\quad p. 299; misprint\_pass.\par
\noindent Imprimé : {\ttfamily\footnotesize sont\allowbreak{}}\par
\noindent Corrigé : {\ttfamily\footnotesize son\allowbreak{}}\par
\noindent Source : 7I:521; 7I:1091. \nolinkurl{fixes/FIXLOG_7I_CURRENT.jsonl#813}; \nolinkurl{fixes/MISPRINTS_7I.json#1091}\par

\par\medskip\noindent\hypertarget{app-F7I-VIII-066.1}{}\textbf{F7I-VIII-066.1}\quad p. 300; misprint\_pass.\par
\noindent Imprimé : {\ttfamily\footnotesize dans\allowbreak{}}\par
\noindent Corrigé : {\ttfamily\footnotesize de\allowbreak{}}\par
\noindent Source : 7I:522; 7I:1092. \nolinkurl{fixes/FIXLOG_7I_CURRENT.jsonl#814}; \nolinkurl{fixes/MISPRINTS_7I.json#522}; \nolinkurl{fixes/MISPRINTS_7I.json#1092}\par

\par\medskip\noindent\hypertarget{app-F7I-VIII-066.2}{}\textbf{F7I-VIII-066.2}\quad p. 300; printed misprint.\par
\noindent Imprimé : \emph{(vide)}\par
\noindent Corrigé : {\ttfamily\footnotesize \_\textbackslash{}eta\allowbreak{}}\par
\noindent Source : 7I:1093. \nolinkurl{fixes/MISPRINTS_7I.json#1093}\par

\par\medskip\noindent\hypertarget{app-F7I-VIII-066.3}{}\textbf{F7I-VIII-066.3}\quad p. 300; inherited corrected-layer difference.\par
\noindent Imprimé : {\ttfamily\footnotesize m\textbackslash{},S\allowbreak{}}\par
\noindent Corrigé : {\ttfamily\footnotesize mS\allowbreak{}}\par
\noindent Source : ADOPT-F7I-VIII-066.3. \nolinkurl{fr/sga7i/diplomatic/expose_VIII_body.tex}; \nolinkurl{fr/sga7i/current/expose_VIII_body.tex}\par

\par\medskip\noindent\hypertarget{app-F7I-VIII-066.4}{}\textbf{F7I-VIII-066.4}\quad p. 300; printed misprint.\par
\noindent Imprimé : {\ttfamily\footnotesize biextensions\allowbreak{}}\par
\noindent Corrigé : {\ttfamily\footnotesize extensions\allowbreak{}}\par
\noindent Source : 7I:524; 7I:1094. \nolinkurl{fixes/MISPRINTS_7I.json#524}; \nolinkurl{fixes/MISPRINTS_7I.json#1094}\par

\par\medskip\noindent\hypertarget{app-F7I-VIII-067.1}{}\textbf{F7I-VIII-067.1}\quad p. 301; misprint\_pass.\par
\noindent Imprimé : {\ttfamily\footnotesize fidèlité\allowbreak{}}\par
\noindent Corrigé : {\ttfamily\footnotesize fidélité\allowbreak{}}\par
\noindent Source : 7I:525. \nolinkurl{fixes/FIXLOG_7I_CURRENT.jsonl#815}\par

\par\medskip\noindent\hypertarget{app-F7I-VIII-067.2}{}\textbf{F7I-VIII-067.2}\quad p. 302; designated-witness attribution.\par
\noindent Imprimé : \emph{(vide)}\par
\noindent Corrigé : {\ttfamily\footnotesize (\textbackslash{}mathbb\{\allowbreak{}}\par
\noindent Source : WITNESS-F7I-VIII-067.2. \nolinkurl{fr/sga7i/diplomatic/expose_VIII_body.tex}; \nolinkurl{fr/sga7i/current/expose_VIII_body.tex}\par

\par\medskip\noindent\hypertarget{app-F7I-VIII-067.3}{}\textbf{F7I-VIII-067.3}\quad p. 302; printed misprint.\par
\noindent Imprimé : \emph{(vide)}\par
\noindent Corrigé : {\ttfamily\footnotesize \}\allowbreak{}}\par
\noindent Source : 7I:1085. \nolinkurl{fixes/MISPRINTS_7I.json#1085}\par

\par\medskip\noindent\hypertarget{app-F7I-VIII-067.4}{}\textbf{F7I-VIII-067.4}\quad p. 302; designated-witness attribution.\par
\noindent Imprimé : \emph{(vide)}\par
\noindent Corrigé : {\ttfamily\footnotesize )\_k\allowbreak{}}\par
\noindent Source : WITNESS-F7I-VIII-067.4. \nolinkurl{fr/sga7i/diplomatic/expose_VIII_body.tex}; \nolinkurl{fr/sga7i/current/expose_VIII_body.tex}\par

\par\medskip\noindent\hypertarget{app-F7I-VIII-068.1}{}\textbf{F7I-VIII-068.1}\quad p. 303; misprint\_pass.\par
\noindent Imprimé : {\ttfamily\footnotesize pourtout\allowbreak{}}\par
\noindent Corrigé : {\ttfamily\footnotesize pour\allowbreak{} tout\allowbreak{}}\par
\noindent Source : 7I:527. \nolinkurl{fixes/FIXLOG_7I_CURRENT.jsonl#816}\par

\par\medskip\noindent\hypertarget{app-F7I-VIII-069.1}{}\textbf{F7I-VIII-069.1}\quad p. 303; inherited corrected-layer difference.\par
\noindent Imprimé : \emph{(vide)}\par
\noindent Corrigé : {\ttfamily\footnotesize \{\allowbreak{}}\par
\noindent Source : ADOPT-F7I-VIII-069.1. \nolinkurl{fr/sga7i/diplomatic/expose_VIII_body.tex}; \nolinkurl{fr/sga7i/current/expose_VIII_body.tex}\par

\par\medskip\noindent\hypertarget{app-F7I-VIII-069.2}{}\textbf{F7I-VIII-069.2}\quad p. 303; inherited corrected-layer difference.\par
\noindent Imprimé : \emph{(vide)}\par
\noindent Corrigé : {\ttfamily\footnotesize \textbackslash{}eta\}\allowbreak{}}\par
\noindent Source : ADOPT-F7I-VIII-069.2. \nolinkurl{fr/sga7i/diplomatic/expose_VIII_body.tex}; \nolinkurl{fr/sga7i/current/expose_VIII_body.tex}\par

\par\medskip\noindent\hypertarget{app-F7I-VIII-070.1}{}\textbf{F7I-VIII-070.1}\quad p. 304; misprint\_pass.\par
\noindent Imprimé : \emph{(vide)}\par
\noindent Corrigé : {\ttfamily\footnotesize  a\allowbreak{}}\par
\noindent Source : 7I:528. \nolinkurl{fixes/FIXLOG_7I_CURRENT.jsonl#817}\par

\par\medskip\noindent\hypertarget{app-F7I-VIII-071.1}{}\textbf{F7I-VIII-071.1}\quad p. 305; inherited corrected-layer difference.\par
\noindent Imprimé : \emph{(vide)}\par
\noindent Corrigé : {\ttfamily\footnotesize \_\textbackslash{}eta\allowbreak{}}\par
\noindent Source : ADOPT-F7I-VIII-071.1. \nolinkurl{fr/sga7i/diplomatic/expose_VIII_body.tex}; \nolinkurl{fr/sga7i/current/expose_VIII_body.tex}\par

\par\medskip\noindent\hypertarget{app-F7I-VIII-071.2}{}\textbf{F7I-VIII-071.2}\quad p. 305; inherited corrected-layer difference.\par
\noindent Imprimé : \emph{(vide)}\par
\noindent Corrigé : {\ttfamily\footnotesize \_\textbackslash{}eta\allowbreak{}}\par
\noindent Source : ADOPT-F7I-VIII-071.2. \nolinkurl{fr/sga7i/diplomatic/expose_VIII_body.tex}; \nolinkurl{fr/sga7i/current/expose_VIII_body.tex}\par

\par\medskip\noindent\hypertarget{app-F7I-VIII-071.3}{}\textbf{F7I-VIII-071.3}\quad p. 305; misprint\_pass.\par
\noindent Imprimé : \emph{(vide)}\par
\noindent Corrigé : {\ttfamily\footnotesize (\allowbreak{}}\par
\noindent Source : 7I:529. \nolinkurl{fixes/FIXLOG_7I_CURRENT.jsonl#818}; \nolinkurl{fixes/MISPRINTS_7I.json#529}\par

\par\medskip\noindent\hypertarget{app-F7I-VIII-072.1}{}\textbf{F7I-VIII-072.1}\quad p. 307; inherited corrected-layer difference.\par
\noindent Imprimé : {\ttfamily\footnotesize \{\allowbreak{}}\par
\noindent Corrigé : {\ttfamily\footnotesize \textbackslash{}eta\allowbreak{} \textbackslash{}times\_\allowbreak{}}\par
\noindent Source : ADOPT-F7I-VIII-072.1. \nolinkurl{fr/sga7i/diplomatic/expose_VIII_body.tex}; \nolinkurl{fr/sga7i/current/expose_VIII_body.tex}\par

\par\medskip\noindent\hypertarget{app-F7I-VIII-072.2}{}\textbf{F7I-VIII-072.2}\quad p. 307; inherited corrected-layer difference.\par
\noindent Imprimé : {\ttfamily\footnotesize \}\allowbreak{}}\par
\noindent Corrigé : \emph{(vide)}\par
\noindent Source : ADOPT-F7I-VIII-072.2. \nolinkurl{fr/sga7i/diplomatic/expose_VIII_body.tex}; \nolinkurl{fr/sga7i/current/expose_VIII_body.tex}\par

\par\medskip\noindent\hypertarget{app-F7I-VIII-073.1}{}\textbf{F7I-VIII-073.1}\quad p. 307; printed misprint.\par
\noindent Imprimé : {\ttfamily\footnotesize *\allowbreak{}}\par
\noindent Corrigé : {\ttfamily\footnotesize \textbackslash{}eta\allowbreak{}}\par
\noindent Source : 7I:1095. \nolinkurl{fixes/MISPRINTS_7I.json#1095}\par

\par\medskip\noindent\hypertarget{app-F7I-VIII-074.1}{}\textbf{F7I-VIII-074.1}\quad p. 308; misprint\_pass.\par
\noindent Imprimé : {\ttfamily\footnotesize tout\allowbreak{}}\par
\noindent Corrigé : {\ttfamily\footnotesize toute\allowbreak{}}\par
\noindent Source : 7I:531. \nolinkurl{fixes/FIXLOG_7I_CURRENT.jsonl#819}\par

\par\medskip\noindent\hypertarget{app-F7I-VIII-075.1}{}\textbf{F7I-VIII-075.1}\quad p. 308; misprint\_pass.\par
\noindent Imprimé : \emph{(vide)}\par
\noindent Corrigé : {\ttfamily\footnotesize '\allowbreak{}}\par
\noindent Source : 7I:1098. \nolinkurl{fixes/FIXLOG_7I_CURRENT.jsonl#1032}; \nolinkurl{fixes/MISPRINTS_7I.json#1098}\par

\par\medskip\noindent\hypertarget{app-F7I-VIII-076.1}{}\textbf{F7I-VIII-076.1}\quad p. 309; misprint\_pass.\par
\noindent Imprimé : {\ttfamily\footnotesize aise\allowbreak{}}\par
\noindent Corrigé : {\ttfamily\footnotesize aide\allowbreak{}}\par
\noindent Source : 7I:532; 7I:1099. \nolinkurl{fixes/FIXLOG_7I_CURRENT.jsonl#820}; \nolinkurl{fixes/MISPRINTS_7I.json#532}; \nolinkurl{fixes/MISPRINTS_7I.json#1099}\par

\par\medskip\noindent\hypertarget{app-F7I-VIII-076.2}{}\textbf{F7I-VIII-076.2}\quad p. 309; designated-witness attribution.\par
\noindent Imprimé : \emph{(1 espace(s), 0 retour(s) de ligne dans le témoin TeX)}\par
\noindent Corrigé : \emph{(vide)}\par
\noindent Source : WITNESS-F7I-VIII-076.2. \nolinkurl{fr/sga7i/diplomatic/expose_VIII_body.tex}; \nolinkurl{fr/sga7i/current/expose_VIII_body.tex}\par

\par\medskip\noindent\hypertarget{app-F7I-VIII-077.1}{}\textbf{F7I-VIII-077.1}\quad p. 312; inherited corrected-layer difference.\par
\noindent Imprimé : {\ttfamily\footnotesize parâitre\allowbreak{}}\par
\noindent Corrigé : {\ttfamily\footnotesize paraître\allowbreak{}}\par
\noindent Source : ADOPT-F7I-VIII-077.1. \nolinkurl{fr/sga7i/diplomatic/expose_VIII_body.tex}; \nolinkurl{fr/sga7i/current/expose_VIII_body.tex}\par

\normalsize\fi
\clearpage\SGAtoc{Exposé IX — MODELES DE NERON ET MONODROMIE}
\selectlanguage{french}

\SGApage{I}{313}


S G A 7


\underline{E X P O S E \ \ IX}


\underline{\underline{MODELES DE NERON ET MONODROMIE}}


\underline{par A. Grothendieck}

(avec un appendice par M. Raynaud)


\underline{S O M M A I R E}

0. Introduction

1. Modèle de Néron des schémas abéliens : notations ; l'accouplement canonique $\Phi_o \times \Phi'_o \longrightarrow (Q/\mathbb{Z})_k$ , et la biextension canonique $W^o$ de $(A^o,{A'}^o)$

2. Partie fixe et partie torique de $T_\ell(A_K)$. Critères de bonne réduction. Théorème d'orthogonalité pour $\ell \neq p$

3. Cas de la réduction semi-stable. Le théorème de réduction semi-stable

4. Application à une conjecture de Serre-Tate et au conducteur

5. Le théorème d'orthogonalité dans le cas $\ell = p$ (cas semi-stable) Dualité des schémas abéliens $B_o$ et $B'_o$ . Caractérisation de la partie fixe. Critères de bonne réduction, de bonne réduction essentielle et de réduction semi-stable

6. Remarques sur la construction de la partie torique et la partie fixe de $T_p(A)$ dans le cas de réduction non semi-stable

7. L'extension de Raynaud $G^\natural$ sur $S$ attachée au schéma abélien $A_K$ à réduction semi-stable. Dualité des schémas abéliens $B,B'$ sur $S$

8. L'extension de Raynaud $A_K^{\natural\,o}$ dans le cas semi-stable, et le ind-groupe $A_K^{\natural}$

9. Définition de l'accouplement de monodromie $\underline{M}_\ell \otimes \underline{M}'_\ell \longrightarrow \mathbb{Z}_{\ell S}$

10. Propriétés de l'accouplement de monodromie. Théorème d'intégrité et de positivité

\SGAvar{11. Composantes connexes du modèle de Néron : relation de dualité, comportement assymptotique


}{11. Composantes connexes du modèle de Néron : relation de dualité, comportement asymptotique


}{F7I-IX-001}\SGApage{I}{314}

12. Modèle de Néron d'une jacobienne et formule de Picard-Lefschetz

13. Liens avec la théorie transcendante : cas analytique complexe

13.1. Rappels sur les groupes analytiques sur $S$

13.2. Complétion formelle le long d'une fibre

13.3. L'extension de RAYNAUD analytique

13.4. Interprétation en termes de ``structures de Hodge mixtes'' de DELIGNE

13.5. Cas d'une situation relativement algébrique

14. Liens avec la théorie transcendante : cas rigide-analytique L'homomorphisme canonique $\underline{M}_K \longrightarrow A_K^{\natural\,o}$ .


\SGAsectiontoc{0. Introduction}{0}
0. \underline{Introduction}

0.1. Dans le présent exposé, nous appliquons les résultats des exposés précédents à l'étude du \underline{modèle de Néron} (1.1) $A$ d'un schéma abélien $A_K$ défini sur le corps des fonctions $K$ du trait hensélien $S$. Nous verrons en effet que les propriétés les plus importantes du modèle de Néron peuvent s'interpréter en termes de l'action du groupe de monodromie locale $I$ sur le module de Tate $T_\ell(A_K(\overline{K}))$ (où $\ell$ est un nombre premier $\neq p$). Comme le dual de ce dernier s'interprète comme le groupe de cohomologie $\ell$-adique $H^1(A_{\overline{K}},\mathbb{Z}_\ell)$, et que pour tout schéma projectif et lisse $X$ sur $K$, $H^1(X_{\overline{K}},\mathbb{Z}_\ell)$ s'identifie à $H^1(A_{\overline{K}},\mathbb{Z}_\ell)$, où $A_{\overline{K}}$ est la variété d'Albanese de $X_{\overline{K}}$, on voit donc que le présent exposé peut aussi être considéré comme une étude détaillée des \underline{phénomènes de monodromie locale pour les} $H^1$ $\ell$-\underline{adiques} (ou mieux encore, pour les $H^1$ ``motiviques'') \underline{des variétés projectives et lisses sur} $K$. Dans cette optique, il semble clair que les principaux résultats du présent exposé sont destinés à être englobés dans une ``théorie de Néron'' pour des motifs de poids quelconque, i.e. pour des $H^i$ ($\ell$-adiques, ou de De Rham, ou de Hodge etc) avec $i$ quelconque, qu'on ne commence qu'à entrevoir à l'heure actuelle. (Cf. à ce sujet [9], et plus particulièrement les conjectures de DELIGNE 9.8 à 9.13 du rapport cité.)


\SGApage{I}{315}

0.2.  Le résultat le plus important du présent exposé est le "\underline{théorème de réduction semi-stable}" 3.6, affirmant qu'après extension séparable finie $K'$ de $K$, le modèle de Néron de $A_{K'}$ a une fibre spéciale dont la composante neutre est extension d'un schéma abélien par un tore. Ce résultat avait été conjecturé par J.P. SERRE en 1964, et le principe de la démonstration que nous en donnons ici (*) remonte à la même année. Il apparaît (3.5) que cette propriété de "réduction semi-stable" équivaut au fait que l'action du groupe de monodromie $I$ sur $T_\ell(A_K(\overline{K}))$ est "essentiellement unipotente d'échelon 2" ($\ell$ étant un nombre premier fixé distinct de la caractéristique résiduelle $p$ de $S$), et sous cette forme a été prouvée en termes de $H^1$ $\ell$-adique dans l'exposé III comme cas particulier du "\underline{théorème de monodromie}", en utilisant la résolution des singularités des schémas excellents de dimension 2 (due à S. ABHYANKAR). Pour établir l'équivalence des deux propriétés, nous aurons besoin d'autre part d'un "théorème d'orthogonalité" 2.4, dû à IGUSA [11] dans le cas particulier où $A_K$ est la jacobienne d'une courbe lisse $X_K$ sur $K$, fibre générique d'un schéma projectif régulier et connexe $X$ sur $S$ dont la fibre spéciale géométrique n'a comme seuls points singuliers que des points singuliers quadratiques ordinaires. C'est d'ailleurs là un cas de réduction semi-stable, qui ne suffirait pas pour notre propos ; c'est surtout pour établir le théorème d'orthogonalité sous sa forme générale que nous utilisons l'outil commode que fournit la théorie des biextensions, développée dans les deux exposés précédents. Signalons d'ailleurs qu'une démonstration très différente du théorème de réduction semi-stable, n'utilisant pas la résolution des singularités,

(*) Il est développé dans une lettre du conférencier à J.P. SERRE du 30.10.64.


\SGApage{I}{316}

a été trouvée par D. MUMFORD, en utilisant sa théorie des fonctions $\theta$ (non publié) ; malheureusement, il est obligé de se limiter au cas d'une caractéristique résiduelle $\neq 2$.

\SGAvar{0.3.  Le théorème de réduction semi-stable est l'outil principal utilisé dans les paragraphes ultérieurs pour l'étude des modèles de Néron, que nous entreprenons ici du \underline{point de vue} "\underline{algébrique}". Il semble indispensable également dans l'étude des variétés abéliennes sur $K$ du \underline{point de vue} "\underline{rigide-analytique}" ou des "fonctions $\theta$ " (ces termes étant synonymes, parait-il, lorsqu'il est question de variétés abéliennes sur des corps valués complets ...), telle qu'elle a été développée récemment par D. MUMFORD et M. RAYNAUD (à la suite de travaux de TATE et MORIKAWA). (Nous dirons quelques mots aux \S{} 13, 14 sur les liens entre les points de vue algébrique et transcendant.) Enfin, ce même théorème est un ingrédient-clef de la démonstration du "\underline{théorème de réduction semi-stable des courbes algébriques}" de ARTIN-DELIGNE-MUMFORD [3] (*), dont nous avons déjà eû à faire usage dans l'étude de l'action de monodromie sur le groupe fondamental ("théorème d'action essentiellement modérée" V \hspace{2em}), et qui est utilisé à nouveau dans la démonstration du théorème d'intégrité et de positivité 10.4 dans le présent exposé (**).

}{0.3.  Le théorème de réduction semi-stable est l'outil principal utilisé dans les paragraphes ultérieurs pour l'étude des modèles de Néron, que nous entreprenons ici du \underline{point de vue} "\underline{algébrique}". Il semble indispensable également dans l'étude des variétés abéliennes sur $K$ du \underline{point de vue} "\underline{rigide-analytique}" ou des "fonctions $\theta$ " (ces termes étant synonymes, paraît-il, lorsqu'il est question de variétés abéliennes sur des corps valués complets ...), telle qu'elle a été développée récemment par D. MUMFORD et M. RAYNAUD (à la suite de travaux de TATE et MORIKAWA). (Nous dirons quelques mots aux \S{} 13, 14 sur les liens entre les points de vue algébrique et transcendant.) Enfin, ce même théorème est un ingrédient-clef de la démonstration du "\underline{théorème de réduction semi-stable des courbes algébriques}" de ARTIN-DELIGNE-MUMFORD [3] (*), dont nous avons déjà eu à faire usage dans l'étude de l'action de monodromie sur le groupe fondamental ("théorème d'action essentiellement modérée" V \hspace{2em}), et qui est utilisé à nouveau dans la démonstration du théorème d'intégrité et de positivité 10.4 dans le présent exposé (**).

}{F7I-IX-002}0.4.  Le présent exposé se distingue des autres exposés du Séminaire par le fait que dans l'étude du module de TATE $T_\ell(A_K)$, nous avons systématiquement tenu compte aussi du cas où $\ell$ est égal à la caractéristique

(*)  La démonstration de ce théorème avait été donnée par P. DELIGNE dans le Séminaire oral, mais il n'a pas été jugé utile de la reproduire dans les textes de Séminaire.

(**) Il existe une autre démonstration de ce dernier théorème, via la théorie rigide-analytique, n'utilisant pas le théorème de réduction semi-stable.


\SGApage{I}{317}

\SGAvar{résiduelle $p$. Il se trouve que le langage des \underline{groupes de Barsotti-Tate} [33] (*) permet d'étendre à ce cas tous les phénomènes les plus importants établis pour le cas $\ell \neq p$, sous la réserve toutesfois, parfois, de disposer du théorème fondamental de TATE [33 th. 4], établi pour l'instant seulement dans le cas où $K$ est de caractéristique nulle. Bien entendu, dans ce dernier cas le pro-$p$-groupe de Barsotti-Tate $T_p(A_K)$ s'interprète encore en termes de l'action du groupe de monodromie $I$ sur le module $p$-adique $T_p(A_K(\overline{K}))$, mais on fera attention que la plupart des énoncés que pourrait suggérer cette analogie superficielle avec le cas $\ell \neq p$ sont grossièrement faux. (Ainsi, il n'est plus vrai que l'action du groupe de monodromie soit essentiellement unipotente pour $\ell = p$.) Mais utilisant des traductions plus judicieuses (telle par exemple 5.13 plus bas), il est clair dès maintenant que l'étude des pro-$p$-groupes de Barsotti-Tate définis par voie géométrique sur le corps des fonctions d'un trait à caractéristique résiduelle $p > 0$, qu'il soit ou non d'inégales caractéristiques,doit être considéré comme faisant partie de la théorie de la "monodromie locale", où ces groupes jouent le rôle de "\underline{systèmes locaux} $p$-\underline{adiques}" sur $\mathrm{Spec}(K)$. En fait, ce sont essentiellement les "systèmes locaux" qui s'introduisent dans l'\underline{étude}, \underline{du point de vue} "$p$-\underline{adique }", \underline{du} $H^1$ \underline{des variétés projectives et lisses sur} $K$. Pour l'étude des $H^i$ supérieurs du point de vue $p$-adique, qui reste à faire, il apparaît qu'il faudra élargir convenablement la catégorie des pro-$p$-groupes de Barsotti-Tate sur $K$ resp. sur $S$, en s'inspirant des points de vue fournis par la "cohomologie cristalline".

}{résiduelle $p$. Il se trouve que le langage des \underline{groupes de Barsotti-Tate} [33] (*) permet d'étendre à ce cas tous les phénomènes les plus importants établis pour le cas $\ell \neq p$, sous la réserve toutefois, parfois, de disposer du théorème fondamental de TATE [33 th. 4], établi pour l'instant seulement dans le cas où $K$ est de caractéristique nulle. Bien entendu, dans ce dernier cas le pro-$p$-groupe de Barsotti-Tate $T_p(A_K)$ s'interprète encore en termes de l'action du groupe de monodromie $I$ sur le module $p$-adique $T_p(A_K(\overline{K}))$, mais on fera attention que la plupart des énoncés que pourrait suggérer cette analogie superficielle avec le cas $\ell \neq p$ sont grossièrement faux. (Ainsi, il n'est plus vrai que l'action du groupe de monodromie soit essentiellement unipotente pour $\ell = p$.) Mais utilisant des traductions plus judicieuses (telle par exemple 5.13 plus bas), il est clair dès maintenant que l'étude des pro-$p$-groupes de Barsotti-Tate définis par voie géométrique sur le corps des fonctions d'un trait à caractéristique résiduelle $p > 0$, qu'il soit ou non d'inégales caractéristiques,doit être considérée comme faisant partie de la théorie de la "monodromie locale", où ces groupes jouent le rôle de "\underline{systèmes locaux} $p$-\underline{adiques}" sur $\mathrm{Spec}(K)$. En fait, ce sont essentiellement les "systèmes locaux" qui s'introduisent dans l'\underline{étude}, \underline{du point de vue} "$p$-\underline{adique }", \underline{du} $H^1$ \underline{des variétés projectives et lisses sur} $K$. Pour l'étude des $H^i$ supérieurs du point de vue $p$-adique, qui reste à faire, il apparaît qu'il faudra élargir convenablement la catégorie des pro-$p$-groupes de Barsotti-Tate sur $K$ resp. sur $S$, en s'inspirant des points de vue fournis par la "cohomologie cristalline".

}{F7I-IX-003}(*) ou "groupes $p$-divisibles" dans la terminologie de TATE.


\SGApage{I}{318}

0.5. L'exposé donné ici, très proche des exposés oraux jusqu'au \S 4, est nettement plus détaillé et complet à partir des paragraphes suivants, qui avaient été seulement partiellement esquissés. Il ne sera pas utilisé dans les exposés ultérieurs. Par contre, certains résultats du présent exposé (10.4 et 12.5) utilisent la \underline{formule de Picard-Lefschetz}, qui sera établie seulement ultérieurement.


\underline{LEITFADEN}

\begin{center}
\begin{tikzcd}[row sep=large, column sep=small]
 \& 1 \arrow[d, dash] \& \& \& \& \\
 \& 2 \arrow[d, dash] \& \& \& \& \\
 \& 3 \arrow[dl, dash] \arrow[dd, dash] \& \& \& \& \\
4 \& \& \& \& \& \\
 \& 5 \arrow[dl, dash] \arrow[ddd, dash] \arrow[ddr, dash] \arrow[drrr, dash] \& \& \& \& \\
6 \& \& \& \& 9 \arrow[dd, dash] \& \\
 \& \& 7 \arrow[dl, dash] \arrow[dddl, dash] \arrow[ddd, dash] \& \& \& \\
 \& 8 \arrow[dd, dash] \arrow[ddr, dash] \& \& \& 10 \arrow[dl, dash] \arrow[dr, dash] \arrow[ddlll, dash] \arrow[ddll, dash] \& \\
 \& \& \& 11 \& \& 12 \\
 \& 13 \& 14 \& \& \&
\end{tikzcd}
\end{center}























\SGApage{I}{319}


\SGAvar{\SGAsectiontoc{1. Modèle de Néron des schémas abéliens : notations ; l'accouplement canonique $\Phi_o \times \Phi'_o \longrightarrow (\mathbb{Q}/\mathbb{Z})_k$ , et la biextension cancnique $W^o$ de $(A^o,{A'}^{o})$.}{1}
1. \underline{Modèle de Néron des schémas abéliens} : \underline{notations} ; \underline{l'accouplement canonique} $\Phi_o \times \Phi'_o \longrightarrow (\mathbb{Q}/\mathbb{Z})_k$ , \underline{et la biextension cancnique} $W^o$ \underline{de} $(A^o,{A'}^{o})$.

1.0.  Dans toute la suite du présent exposé, on suppose fixé un trait $S$, pour lequel nous utilisons nos notations habituelles $s$, $\eta$, $k$, $X$ etc. de I 1 . Nous supposons donné un schéma abélien $A_K$ sur $K$, dont le schéma abélien dual [5] sera noté $A'_K$ . La dualité entre $A_K$ et $A'_K$ est donnée par une correspondance divisorielle $W_K$ sur $(A_K,A'_K)$, i.e. un Module inversible sur $A_K \times A'_K$ , birigidifié relativement aux sections unité de $A_K,A'_K$ , lequel est déterminé à isomorphisme unique près par sa classe (cf. VIII 3.1)

}{\SGAsectiontoc{1. Modèle de Néron des schémas abéliens : notations ; l'accouplement canonique $\Phi_o \times \Phi'_o \longrightarrow (\mathbb{Q}/\mathbb{Z})_k$ , et la biextension canonique $W^o$ de $(A^o,{A'}^{o})$.}{1}
1. \underline{Modèle de Néron des schémas abéliens} : \underline{notations} ; \underline{l'accouplement canonique} $\Phi_o \times \Phi'_o \longrightarrow (\mathbb{Q}/\mathbb{Z})_k$ , \underline{et la biextension canonique} $W^o$ \underline{de} $(A^o,{A'}^{o})$.

1.0.  Dans toute la suite du présent exposé, on suppose fixé un trait $S$, pour lequel nous utilisons nos notations habituelles $s$, $\eta$, $k$, $K$ etc. de I 1 . Nous supposons donné un schéma abélien $A_K$ sur $K$, dont le schéma abélien dual [5] sera noté $A'_K$ . La dualité entre $A_K$ et $A'_K$ est donnée par une correspondance divisorielle $W_K$ sur $(A_K,A'_K)$, i.e. un Module inversible sur $A_K \times A'_K$ , birigidifié relativement aux sections unité de $A_K,A'_K$ , lequel est déterminé à isomorphisme unique près par sa classe (cf. VIII 3.1)

}{F7I-IX-004}\begin{equation}\tag{1.0.1}
w_K \in \mathrm{Biext}^1(A_K,A'_K;\ \mathbb{G}_{m\,K}) \qquad .
\end{equation}
Un rôle important sera joué par l'accouplement correspondant (VIII 2.2)

\begin{equation}\tag{1.0.2}
\varphi : T_\ell(A_K) \times T_\ell(A'_K) \longrightarrow \mathbb{Z}_\ell(1)_K \overset{\mathrm{dfn}}{=} T_\ell(\mathbb{G}_{mK}) \qquad ,
\end{equation}
où $\ell$ est un nombre premier quelconque. Lorsque $\ell \neq p = \mathrm{car}.k$, la donnée de $\varphi$ équivaut à celle d'un accouplement de $\mathbb{Z}_\ell$-modules libres de type fini

\begin{equation}\tag{1.0.3}
\varphi : T_\ell(A_K(\overline{K})) \times T_\ell(A'_K(\overline{K})) \longrightarrow T_\ell(\overline{K}*) \qquad ,
\end{equation}
compatible avec l'action de

\begin{equation*}
\pi = \mathrm{Gal}(\overline{K}/K)
\end{equation*}
sur les trois modules en jeu, $\overline{K}$ étant une clôture séparable de $K$. Il est connu (VIII 3.2) que l'accouplement en question est une dualité parfaite, i.e.  qu'il  identifie chacun des modules $T_\ell(A_K(\overline{K}))$, $T_\ell(A'_K(\overline{K}))$ au ``dual'' de l'autre à valeurs dans le module inversible $\mathbb{Z}_\ell(1) = T_\ell(\overline{K}*)$.


\SGApage{I}{320}

On supposera parfois choisie une polarisation de $A_K$ , qui est donc un élément

\begin{equation}\tag{1.0.4}
\xi \in \mathrm{NS}(A_K/K) = \underline{\mathrm{NS}}_{A_K/K}(K) \qquad ,
\end{equation}
où

\begin{equation*}
\underline{\mathrm{NS}}_{A_K/K} \overset{\mathrm{dfn}}{=} \underline{\mathrm{Pic}}_{A_K/K} / (\underline{\mathrm{Pic}}{}^{o}_{A_K/K} = A'_K) \qquad .
\end{equation*}
Elle définit de la façon habituelle un homomorphisme

\begin{equation}\tag{1.0.5}
\psi_\xi : A_K \longrightarrow A'_K \qquad ,
\end{equation}
d'où un homomorphisme

\begin{equation}\tag{1.0.6}
T_\ell(\psi_\xi) \ \text{ou} \ \psi_\xi : T_\ell(A_K) \longrightarrow T_\ell(A'_K) \qquad ,
\end{equation}
et une forme bilinéaire $\varphi_\xi$

\begin{equation}\tag{1.0.7}
\varphi_\xi(x,y) = \varphi(x,\psi_\xi(y)) : T_\ell(A_K) \times T_\ell(A_K) \longrightarrow \mathbb{Z}_\ell(1)_K \qquad ,
\end{equation}
qu'on peut aussi interpréter pour $\ell \neq p$ comme accouplement de $\mathbb{Z}_\ell$-modules libres de type fini à opérateurs :

\begin{equation}\tag{1.0.8}
\varphi_\xi : T_\ell(A_K(\overline{K})) \times T_\ell(A_K(\overline{K})) \longrightarrow T_\ell(\overline{K}*) \qquad .
\end{equation}
Ces constructions ont un sens pour tout élément du groupe de Néron-Séveri de $A_K$, et fournissent une forme $\varphi_\xi$ (1.0.7) ou (1.0.8) \underline{alternée} (VIII 2.2.11). Le fait que $\xi$ soit une \underline{polarisation} implique que $\psi_\xi$ dans (1.0.5) est une isogénie, ou ce qui revient au même puisque $\dim A_K = \dim A'_K$, un épimorphisme de Groupes, ou encore que $T_\ell(\psi_\xi)$ dans (1.0.6) est injectif (ou encore, a un conoyau fini), ou enfin que la forme $\varphi_\xi$ de (1.0.7) est \underline{séparante}.


\SGApage{I}{321}

1.1. Rappelons qu'il existe sur $S$ un \underline{schéma lisse} $A$, appelé \underline{modèle de Néron de} $A_K$ , qui représente le foncteur

\begin{equation}\tag{1.1.1}
S' \longmapsto \mathrm{Hom}\ (S'_\eta,A_K)
\end{equation}
sur la catégorie des $S$-schémas lisses. On a donc par définition un isomorphisme de foncteurs en $S'$, schéma relatif lisse sur $S$ :

\begin{equation}\tag{1.1.2}
\mathrm{Hom}_S(S',A) \quad \simeq \quad \mathrm{Hom}_K(S'_\eta,A_K) \qquad .
\end{equation}
La construction de $A$ est donnée par NERON [19] lorsque le corps résiduel $k$ de $S$ est parfait, et a été étendue par RAYNAUD au cas général au cours d'un séminaire non rédigé qu'il a donné à l'IHES en 1966/67 (cf.[21] pour l'énoncé des résultats principaux obtenus par RAYNAUD, complétant les résultats de NERON de loc. cit.). Nous admettrons ici ce résultat d'existence, ainsi que le fait que $A$ est un \underline{schéma} quasi-projectif sur $S$. Appliquant (1.1.2) pour $S'$ lisse \underline{sur} $\eta$, on constate de plus que le morphisme canonique suivant est un isomorphisme

\begin{equation}\tag{1.1.3}
A \times_S \mathrm{Spec}(K) \xrightarrow{\ \sim\ } A_K \qquad ,
\end{equation}
ce qui justifie le couple de notations $A,A_K$ . On écrira aussi, le cas échéant, $A_\eta$ au lieu de $A_K$ . Utilisant (1.1.3) comme isomorphisme d'identification, la bijection canonique (1.1.2) n'est autre que l'application de \underline{restriction aux fibres génériques.}

\SGAvar{Comme le foncteur (1.1.1) a évidemment une structure de Groupe commutatif, il s'ensuit que le modèle de Néron $A$ qui le représente a lui-même une structure de \underline{schéma en groupes commutatif lisse sur} $S$. La suite de cet exposé est consacré à l'étude de ce schéma en groupes.


}{Comme le foncteur (1.1.1) a évidemment une structure de Groupe commutatif, il s'ensuit que le modèle de Néron $A$ qui le représente a lui-même une structure de \underline{schéma en groupes commutatif lisse sur} $S$. La suite de cet exposé est consacrée à l'étude de ce schéma en groupes.


}{F7I-IX-005}\SGApage{I}{322}

Un invariant important, déduit de la considération du modèle de Néron, est évidemment le schéma en groupes commutatifs, lisse et de type fini sur $k$:

\begin{equation}\tag{1.1.4}
A_o = A \times_S s = A_k \qquad .
\end{equation}
Suivant les notations générales, on désigne par $A_o^{\ o}$ sa composante neutre, et on considère également le schéma en groupes

\begin{equation}\tag{1.1.5}
\Phi_o = A_o/A_o^{\ o}
\end{equation}
des composantes connexes de $A_o$, qui est un schéma en groupes fini et étale sur $k$. En termes de la clôture séparable $\bar{k}$ de $k$, il est donc déterminé par le groupe fini ordinaire

\begin{equation}\tag{1.1.6}
\Phi_o(\bar{k}) \qquad ,
\end{equation}
avec les opérations naturelles de

\begin{equation*}
\pi_o = \mathrm{Gal}(\bar{k}/k)
\end{equation*}
sur ce groupe.

En même temps que le modèle de Néron $A$ de $A_K$ , nous étudierons

\begin{equation}\tag{1.1.7}
A' = \text{modèle de Néron de } A'_K \qquad ,
\end{equation}
qui est également un schéma en groupes commutatifs, lisse et de type fini sur $S$. On désignera par

\SGAvar{\begin{equation}\tag{1.1.8}
\Phi'_o = A'_o/{A'}_o^{\ o}
\end{equation}
le groupe des composantes connexes de sa fibre spéciale $A_o$ .


}{\begin{equation}\tag{1.1.8}
\Phi'_o = A'_o/{A'}_o^{\ o}
\end{equation}
le groupe des composantes connexes de sa fibre spéciale $A'_o$ .


}{F7I-IX-006}\SGApage{I}{323}

1.2. Appliquons VIII 7.1 b) au couple $(A,A')$ et à la biextension $W_K$ de (1.0.1) $(A_K,A'_K)$ par $G_{mK}$ . On trouve donc un \underline{accouplement canonique}

\begin{equation}\tag{1.2.1}
\Phi_o \times \Phi'_o \longrightarrow (\mathbb{Q}/\mathbb{Z})_k \qquad ,
\end{equation}
i.e. un accouplement de $\pi_o$-modules

\begin{equation}\tag{1.2.2}
\Phi_o(\overline{k}) \times \Phi'_o(\overline{k}) \longrightarrow \mathbb{Q}/\mathbb{Z} \qquad .
\end{equation}
Cet accouplement apparaît ici comme obstruction au prolongement de $W_K$ en une biextension $W$ de $(A,A')$ par $\mathbb{G}_{mS}$ , et la signification de ce point de vue est explicitée dans VIII 7.3.2. Signalons tout de suite à propos de cet accouplement la

\underline{Conjecture} 1.3. \underline{L'accouplement} (1.2.1) \underline{est une dualité parfaite, i.e.} \underline{identifie chacun des groupes de composantes connexes} $\Phi_o$ \underline{et} $\Phi'_o$ \underline{au dual} \underline{à valeurs dans} $(\mathbb{Q}/\mathbb{Z})_k$ \underline{de l'autre.}

\SGAvar{1.3.1. Signalons qu'il résulte de la commutation de la formation de (1.2.1) aux changements de base formellements nets (VIII 7.3.5.3), et de la propriété analogue pour la formation des modèles de Néron, qu'il suffirait de prouver 1.3 dans le cas où $S$ est complet à corps résiduel séparablement clos. La conjecture a été prouvée par ARTIN et MAZUR dans le cas où $A_K$ est la jacobienne d'une courbe propre et lisse sur $K$, en utilisant leur théorie générale (non publiée) de l'autodualité de la jacobienne relative (au sens des catégories dérivées) d'un schéma en courbes propre sur $S$ ; ils ont prouvé la même conjecture 1.3 dans le cas où le corps résiduel est \underline{fini}, en utilisant la théorie de dualité de TATE [32]. Nous verrons également dans 11.4 et 11.5(modulo des


}{1.3.1. Signalons qu'il résulte de la commutation de la formation de (1.2.1) aux changements de base formellement nets (VIII 7.3.5.3), et de la propriété analogue pour la formation des modèles de Néron, qu'il suffirait de prouver 1.3 dans le cas où $S$ est complet à corps résiduel séparablement clos. La conjecture a été prouvée par ARTIN et MAZUR dans le cas où $A_K$ est la jacobienne d'une courbe propre et lisse sur $K$, en utilisant leur théorie générale (non publiée) de l'autodualité de la jacobienne relative (au sens des catégories dérivées) d'un schéma en courbes propre sur $S$ ; ils ont prouvé la même conjecture 1.3 dans le cas où le corps résiduel est \underline{fini}, en utilisant la théorie de dualité de TATE [32]. Nous verrons également dans 11.4 et 11.5(modulo des


}{F7I-IX-007}\SGApage{I}{324}

\SGAvar{vérifications de compatibilité) que 1.3 est vrai pour les composantes $\ell$-primaires pour $\ell \neq p$, et que 1.3 est vrai dans le cas de la ``réduction semi-stable''. Ces résultats rendent 1.3 extrêmement plausible. D'ailleurs 1.3 peut également être considéré comme un sous-produit d'une théorie de dualité générale pour les schémas en groupes sur les corps de fractions d'anneaux de valuation discrète complets, formellement analogue à celle de SGA 5 I, qui reste heuristique à l'heure actuelle, et qui contiendrait également la théorie du corps de classe local sous la forme que lui a donné SERRE [27].

}{vérifications de compatibilité) que 1.3 est vrai pour les composantes $\ell$-primaires pour $\ell \neq p$, et que 1.3 est vrai dans le cas de la ``réduction semi-stable''. Ces résultats rendent 1.3 extrêmement plausible. D'ailleurs 1.3 peut également être considéré comme un sous-produit d'une théorie de dualité générale pour les schémas en groupes sur les corps de fractions d'anneaux de valuation discrète complets, formellement analogue à celle de SGA 5 I, qui reste heuristique à l'heure actuelle, et qui contiendrait également la théorie du corps de classe local sous la forme que lui a donnée SERRE [27].

}{F7I-IX-008}1.3.2. La validité de 1.3 pour un schéma abélien donné $A_K$ revient à dire que l'accouplement (1.2.1) est \underline{séparant} à gauche et à droite. D'ailleurs le fait qu'il soit séparant à gauche signifie, en vertu de VIII 7.3.2, que le seul sous-schéma en groupes ouvert $V$ de $A$ tel que $W_K$ puisse se prolonger en une biextension de $(V,A')$ par $G_{mS}$ est la composante neutre $A^o$ de $A$. On a une interprétation symétrique pour la condition que (1.2.1) soit séparant à droite.

1.4. De façon générale si

\begin{equation}\tag{1.4.1}
U \subset A \ \mbox{et} \ U' \subset A'
\end{equation}
sont des sous-schémas en groupes ouverts tels que $W_K$ se prolonge en une biextension de $(U,U')$ par $\mathbb{G}_{mS}$ (ce qui, en vertu de VIII 7.3.2, s'explicite par la condition que les sous-groupes

\SGAvar{\begin{equation}\tag{1.4.2}
\Phi'_o = U_o/U_o^{\,o} \subset \Phi_o \quad , \quad \Psi'_o = {U'}_o/{U'}_o^{\,o} \subset \Psi_o
\end{equation}
de $\Phi_o$ et de $\Psi_o$ définis par $U$ et $U'$sont orthogonaux l'un à l'autre pour


}{\begin{equation}\tag{1.4.2}
F_o = U_o/U_o^{\,o} \subset \Phi_o \quad , \quad F'_o = {U'}_o/{U'}_o^{\,o} \subset \Phi'_o
\end{equation}
de $\Phi_o$ et de $\Phi'_o$ définis par $U$ et $U'$ sont orthogonaux l'un à l'autre pour


}{F7I-IX-009}\SGApage{I}{325}

\SGAvar{l'accouplement (1.2.1)), on désigne par $W$ ce prolongement (unique à isomorphe unique près) de $W_K$ . Il est déterminé à isomorphisme unique près par sa classe

}{l'accouplement (1.2.1)), on désigne par $W$ ce prolongement (unique à isomorphisme unique près) de $W_K$ . Il est déterminé à isomorphisme unique près par sa classe

}{F7I-IX-010}\begin{equation}\tag{1.4.3}
w \in \mathrm{Biext}^1(U,U' \; ; \; \mathbb{G}_{mS}) \quad ,
\end{equation}
elle-même caractérisée par la condition que son image dans $\mathrm{Biext}^1(A_K,A'_K;\mathbb{G}_{m\,K})$ soit $\mathbf{w}_K$ de (1.0.1).

On peut prendre par exemple pour $U$ et $U'$ les composantes neutres $A^o$ et $A'{}^o$ de $A$ et $A'$. On trouve ainsi une \underline{biextension canonique} $W^o$ \underline{ou} \textbf{W de} $(A^o,A'{}^o)$ \underline{par} $\mathbb{G}_{mS}$, donnant lieu à une classe

\begin{equation}\tag{1.4.4}
w = w^o \in \mathrm{Biext}^1(A^o,A'{}^o;\mathbb{G}_{mS}) \quad ,
\end{equation}
avec laquelle nous travaillerons principalement. Evidemment les classes $w$ (1.4.3) relatives à divers couples $(U,U')$ ont toutes $w^o$ (1.4.4) pour restriction à $(A^o,A'{}^o)$.

\underline{Proposition} 1.5. \underline{Considérons la fibre spéciale} $W_s$ \underline{de la biextension canonique} (1.4.4) $W$ \underline{de} $(A^o,A'{}^o)$ \underline{par} $\mathbb{G}_{mS}$, \underline{de sorte que} $W_s$ \underline{est une biextension de} $(A^o_o,A'{}^o_o)$ \underline{par} $\mathbb{G}_{mk}$ . \underline{Cette dernière est séparante} (VIII 4.11).

Choisissons en effet un Module inversible ample $\mathbb{L}$ sur $A^o$. Alors $\mathbb{L}_\eta$ définit un homomorphisme (cf. 1.0.5)

\begin{equation*}
A_K \longrightarrow A'_K \quad ,
\end{equation*}
qui se prolonge en un homomorphisme $A \longrightarrow A'$, d'où

\begin{equation*}
\psi : A^o \longrightarrow A'{}^o \quad ,
\end{equation*}
et on peut considérer la biextension $W_{\mathbb{L}}$ de $(A^o,A^o)$ par $\mathbb{G}_{mS}$, image inverse


\SGApage{I}{326}

de $W$ par $(\mathrm{id}_{A^o} , \psi)$. La restriction de $W_{\mathbb{L}}$ aux fibres génériques n'est alors autre que la biextension $\delta(\mathbb{L}_\eta)$ de (VIII 4.12), donc en vertu de VIII 7.4, $W_{\mathbb{L}}$ est associé au Module birigidifié $\delta(\mathbb{L})$, et par suite $(W_{\mathbb{L}})_s$ est défini par $\delta(\mathbb{L}_s)$. D'autre part, il peut également être décrit comme l'image inverse de la biextension $W_s$ par les morphismes

\begin{equation*}
\mathrm{id}_{A^o_o} : A^o_o \longrightarrow A^o_o \quad , \quad \psi_o : A^o_o \longrightarrow A'{}^o_o \quad .
\end{equation*}
Comme $\mathbb{L}$ donc $\mathbb{L}_s$ est ample, il résulte de VIII 4.13 que $W_{\mathbb{L}_s}$ est séparant, et en particulier séparant à gauche, et \underline{à fortiori} $W_s$ est donc séparant à gauche. Par symétrie, $W_s$ est également séparant à droite, ce qui prouve 1.5.

Notons que nous préciserons considérablement 1.5 dans le cas où $A^o_o$ est de rang unipotent nul (5.4).


\SGAsectiontoc{2. Partie fixe et partie torique de $T_\ell(A_K)$. Critères de bonne réduction. Théorème d'orthogonalité pour $\ell \neq p$}{2}
2. \underline{Partie fixe et partie torique de} $T_\ell(A_K)$. \underline{Critères de bonne réduction.} \underline{Théorème d'orthogonalité pour} $\ell \neq p$

2.1.  Comme tout groupe algébrique commutatif sur $k$, $A_o$ admet un plus grand sous-tore, dont la formation commute à toute extension du corps de base [SGA 3 XVII 7.2.1, XII 1.12] :

\begin{equation}\tag{2.1.1}
T_o \subset A_o \quad .
\end{equation}
Lorsque $k$ est parfait, $T_o$ est caractérisé par la condition que $A^o_o/T_o$ soit extension d'un schéma abélien $B_o$ sur $k$ par un groupe unipotent (lisse et connexe) [25], qui admet nécessairement une suite de composition à quotients tous isomorphes à $G_{ak}$ (SGA 3 XVII 4.1.3)


\SGApage{I}{327}

\begin{equation}\tag{2.1.2}
0 \longrightarrow U_o \longrightarrow A_o^o/T_o \longrightarrow B_o \longrightarrow 0 \quad .
\end{equation}
En tous cas, $k$ parfait ou non, \underline{si} $A_o^o$ est extension d'un schéma abélien $B_o$ par un schéma affine lisse et connexe $L_o$ :

\begin{equation}\tag{2.1.3}
0 \longrightarrow L_o \longrightarrow A_o^o \longrightarrow B_o \longrightarrow 0 \quad ,
\end{equation}
on sait que cette structure d'extension est unique [25], et que $T_o$ est alors $\subset L_o$, donc $T_o$ est le plus grand sous-tore de $L_o$, de sorte qu'on a une suite exacte (2.1.2), avec

\begin{equation}\tag{2.1.4}
U_o = L_o/T_o \quad .
\end{equation}
2.1.5. Signalons que $A_o^o$ est de \underline{rang unipotent nul} (i.e. que, sur une clôture algébrique de $k$, il n'admet aucun sous-groupe isomorphe à $G_a$) si et seulement si $A_o^o/T_o$ est un schéma abélien $B_o$, de sorte qu'on sera alors sous les conditions précédentes, avec $U_o = 0$.

Si $\ell$ est un nombre premier distinct de la caractéristique de $k$, de sorte que $A_o^o$ est $\ell$-divisible, l'inclusion (2.1.1) donne lieu à une suite exacte de pro-schémas en groupes :

\begin{equation}\tag{2.1.6}
\begin{tikzcd}[row sep=large, column sep=large]
0 \arrow[r] \& T_\ell(T_o) \arrow[r] \& T_\ell(A_o^o) \arrow[r] \arrow[d, "\wr"] \& T_\ell(A_o^o/T_o) \arrow[r] \arrow[d, "\wr"] \& 0 \\
\& \& T_\ell(A_o) \& T_\ell(A_o/T_o) \& {,}
\end{tikzcd}
\end{equation}












dont les termes peuvent aussi s'interpréter comme des faisceaux $\ell$-adiques libres constants tordus ; lorsque l'on dispose d'une suite exacte (2.1.2), par exemple $k$ parfait ou $A_o^o$ de rang unipotent nul, on aura

\begin{equation}\tag{2.1.7}
T_\ell(A_o^o/T_o) \overset{\sim}{\longrightarrow} T_\ell(B_o) \quad ,
\end{equation}

\SGApage{I}{328}

\SGAvar{de sorte que la suite exacte (2.1.5) s'écrit alors

}{de sorte que la suite exacte (2.1.6) s'écrit alors

}{F7I-IX-011}\begin{equation}\tag{2.1.8}
0 \longrightarrow T_\ell(T_o) \longrightarrow T_\ell(A_o^o) \longrightarrow T_\ell(B_o) \longrightarrow 0 \quad .
\end{equation}
Posons

\begin{equation}\tag{2.1.9}
n = \dim A_K = \dim A_o \ , \ \mu = \dim T_o \ , \ \lambda = \dim U_o \ , \ \alpha = \dim B_o \ ,
\end{equation}
de sorte qu'on a

\begin{equation}\tag{2.1.10}
n = \mu + \lambda + \alpha \quad ,
\end{equation}
\begin{equation}\tag{2.1.11}
\left\{
\begin{array}{l}
\mathrm{rang}\ T_\ell(T_o) = \mu \ , \ \mathrm{rang}\ T_\ell(A_o^o/T_o) = \mathrm{rang}\ T_\ell(B_o) = 2\alpha \quad , \\[1ex]
\qquad\qquad \mathrm{rang}\ T_\ell(A_o^o) = \mu + 2\alpha \quad ,
\end{array}
\right.
\end{equation}
La deuxième relation étant un théorème bien connu de Weil [5] [15].

\SGAvar{2.1.12. Les entiers $n$, $\mu$, $\lambda$, $\alpha$ sont définis manifestement sans même supposer qu'on dispose d'une suite exacte (2.1.2), en passant à une clôture algébrique de $k$ ; ils sont connus respectivement sous le nom de dimension de $A_o$, \underline{rang réductif de} $A_o$ $=$ dimension de $T_o$ , \underline{rang unipotent de} $A_o$, \underline{rang abélien de} $A_o$ . Il sont encore reliés par les relations (2.1.10) et (2.1.11) (en omettant dans ces dernières le terme rang $T_\ell(B_o)$...).

}{2.1.12. Les entiers $n$, $\mu$, $\lambda$, $\alpha$ sont définis manifestement sans même supposer qu'on dispose d'une suite exacte (2.1.2), en passant à une clôture algébrique de $k$ ; ils sont connus respectivement sous le nom de dimension de $A_o$, \underline{rang réductif de} $A_o$ $=$ dimension de $T_o$ , \underline{rang unipotent de} $A_o$, \underline{rang abélien de} $A_o$ . Ils sont encore reliés par les relations (2.1.10) et (2.1.11) (en omettant dans ces dernières le terme rang $T_\ell(B_o)$...).

}{F7I-IX-012}2.1.13. Lorsque $k$ est de caractéristique $p > 0$, et que $A_o^o$ est $p$-divisible, ou ce qui revient au même, de rang unipotent nul (cf. 2.1.5), alors il y a lieu également de considérer le pro-groupe $T_p(A_o^o)$, qui sera un \underline{groupe de Barsotti-Tate} (ou \underline{groupe} $p$-\underline{divisible} dans la terminologie de [33]). Alors la suite exacte (2.1.8) reste valable pour $\ell = p$ et c'est une suite exacte de groupes de Barsotti-Tate. Ici on aura $\lambda = 0$, donc

\begin{equation}\tag{2.1.13.1}
n = \mu + \alpha \quad ,
\end{equation}

\SGApage{I}{329}

\SGAvar{et les formules (2.1.11) restent valables, en y interprétant le rang comme le rang (ou ``hauteur'') des groupes-de Barsotti-Tate envisagés.

}{et les formules (2.1.11) restent valables, en y interprétant le rang comme le rang (ou ``hauteur'') des groupes de Barsotti-Tate envisagés.

}{F7I-IX-013}2.1.14.  On introduit également pour ${A'}^o_o$ des invariants

\begin{equation*}
{T'}_o \subset {A'}^o_o \qquad , \qquad {U'}_o \qquad , \qquad {B'}_o \qquad ,
\end{equation*}
où ${T'}_o$ est la partie torique de ${A'}^o_o$, ${B'}_o$ sa partie abélienne (définie notamment si $k$ est parfait ou si ${A'}^o_o$ est de rang unipotent nul), ${U'}_o$ sa partie unipotente (définie par exemple si $k$ est parfait, ou, trivialement si ${A'}^o_o$ est de rang unipotent nul). On trouve une suite exacte correspondant à (2.1.6) ou (2.1.8), que nous ne réécrirons pas. On pourrait introduire pour ${A'}^o_o$ des entiers $\mu'$, $\lambda'$, $\alpha'$ comme dans (2.1.9), mais on verra plus bas (2.2.7) que ce sont les mêmes que pour $A^o_o$, ce qui nous dispense d'introduire une nouvelle notation pour eux.

2.2.  Nous nous proposons de préciser les relations entre $T_\ell(A_K)$ et $T_\ell(A_o) = T_\ell(A^o_o)$, en faisant intervenir $T_\ell(A^o)$. Rappelons d'abord :

\underline{Lemme} 2.2.1. \underline{Soit} $m > 0$ \underline{un entier}. \underline{Alors les conditions suivantes sont équivalentes} :

(i)   $m \,\mathrm{id}_{A^o}$ \underline{est plat}.

(ii)  $m \,\mathrm{id}_{A^o}$ \underline{est surjectif}.

(iii) $m \,\mathrm{id}_{A^o}$ \underline{est quasi-fini}.

\SGAvar{(iv)  $m$ \underline{est premier à la caractéristique de} $k$, ou $A^o_o$ \underline{est de rang unipotent nul} (cf. 1.5.5).

}{(iv)  $m$ \underline{est premier à la caractéristique de} $k$, ou $A^o_o$ \underline{est de rang unipotent nul} (cf. 2.1.5).

}{F7I-IX-014}\underline{De plus}, \underline{si les conditions envisagées sont vérifiées pour} $m$, \underline{elles le sont pour} $m^\nu$ \underline{pour tout} $\nu > 0$, \underline{les groupes}


\SGApage{I}{330}

\begin{equation*}
{}_{m^\nu}A^o \overset{\mathrm{dfn}}{=} \mathrm{Ker}(m^\nu \,\mathrm{id}_{A^o})
\end{equation*}
\underline{sont quasi-finis}, \underline{séparés et plats sur} $S$, \underline{et ils forment pour} $\nu$ \underline{variable} \underline{un système projectif de schémas en groupes}, \underline{à morphismes de transition fidèlement plats} $x \longmapsto m^{\nu'-\nu}x$ : ${}_{m^{\nu'}}A^o \longrightarrow {}_{m^\nu}A^o$ .

En effet, le critère de platitude par fibres (EGA IV 11.3.5) nous montre que (i) tout comme (ii), (iii) sont des conditions \underline{sur les fibres}, et elles sont équivalentes par (SGA 3 $\mathrm{VI}_A$ 5.6), parce que les fibres envisagées sont lisses et connexes. D'ailleurs on voit sur la forme (ii) et par un théorème classique de Weil [5] [15] qu'elles sont vérifiées sur la fibre générique, qui est un schéma abélien, il suffit donc de l'exprimer sur la fibre spéciale. L'équivalence avec (iv) est alors également bien connue, grâce à la théorie de structure. Sous n'importe quelle forme, la stabilité de la condition envisagée par $m \longmapsto m^\nu$ est triviale, d'où le fait que ces conditions impliquent que les ${}_{m^\nu}A^o$ sont des schémas en groupes plats quasi-finis sur $S$, évidemment séparés sur $S$ puisque $A^o$ l'est. La dernière assertion sur la fidèle platitude des morphismes de transition envisagés, qui revient à dire que ce sont des épimorphismes pour la topologie fppf, résulte formellement du fait que $m \,\mathrm{id}_{A^o}$ est fidèlement plat, donc un épimorphisme pour la topologie envisagée.

2.2.2.  Par suite, nous prendrons pour $m$ un nombre premier $\ell$ satisfaisant les conditions de 2.2.1. On désigne alors par

\begin{equation}\tag{2.2.2.1}
T_\ell(A^o) = ({}_{\ell^\nu}A^o)_{\nu > 0}
\end{equation}


\SGApage{I}{331}

le système projectif décrit dans 2.2.1. On notera que par les deux changements de base $\eta \longrightarrow S$ et $s \longrightarrow S$, il redonne les systèmes projectifs $T_{\ell}(A_{K})$ et $T_{\ell}(A_{o}^{\,o}) = T_{\ell}(A_{o})$ envisagés dans 1.0 et 2.1 respectivement.

2.2.3. Supposons maintenant $S$ hensélien. On sait alors que tout schéma quasi-fini et séparé $X$ sur $S$ se décompose canoniquement en une somme

\begin{equation}\tag{2.2.3.1}
X = X^{f} \coprod X' \qquad ,
\end{equation}
où $X^{f}$ est \underline{fini sur} $S$, et où $X'_{o} = \emptyset$ i.e. $X'$ est réduit à sa fibre générique $X'_{\eta}$ ; donc $X^{f}$ a même fibre spéciale que $X$. Cette décomposition est évidemment fonctorielle en $X$. On en conclut en particulier que si $X$ est un schéma en groupes sur $S$, alors $X^{f}$ en est un sous-schéma en groupes. Si $X$ est plat resp. étale sur $S$, il en est évidemment de même de $X^{f}$.

Nous appliquerons ceci aux composants ${}_{\ell^{\nu}}A^{o}$ de $T_{\ell}(A^{o})$, d'où un sous-système projectif de groupes

\SGAvar{\begin{equation}\tag{2.2.3.2}
(({}_{\ell^{\nu}}A^{o})^{f})_{\nu > 0} \quad \stackrel{\mathrm{dfn}}{=} \quad T_{\ell}(A^{o})^{f} \subset T_{\ell}(A^{o}) \qquad ,
\end{equation}
qu'on appelera la \underline{partie fixe} du module de Tate $T_{\ell}(A^{o})$. On a évidemment, par construction, des isomorphismes canoniques, fonctoriels en $A$ :

}{\begin{equation}\tag{2.2.3.2}
(({}_{\ell^{\nu}}A^{o})^{f})_{\nu > 0} \quad \stackrel{\mathrm{dfn}}{=} \quad T_{\ell}(A^{o})^{f} \subset T_{\ell}(A^{o}) \qquad ,
\end{equation}
qu'on appellera la \underline{partie fixe} du module de Tate $T_{\ell}(A^{o})$. On a évidemment, par construction, des isomorphismes canoniques, fonctoriels en $A$ :

}{F7I-IX-015}\begin{equation}\tag{2.2.3.3}
T_{\ell}(A^{o})^{f} \times_{S} s \quad \simeq \quad T_{\ell}(A^{o}) \times_{S} s \quad \simeq \quad T_{\ell}(A_{o}^{\,o}) \qquad .
\end{equation}
D'autre part, la fibre générique de $T_{\ell}(A^{o})^{f}$ définit un sous-pro-groupe canonique de $T_{\ell}(A_{K})$, appelée encore ``\underline{partie fixe}'' de ce module de Tate, soit

\begin{equation}\tag{2.2.3.4}
T_{\ell}(A_{K})^{f} \subset T_{\ell}(A_{K}) \qquad .
\end{equation}

\SGApage{I}{332}

2.2.4. Lorsque $\ell \neq p = \mathrm{car}.k$, alors $T_{\ell}(A^{o})^{f}$ s'interprète comme un faisceau $\ell$-adique constant tordu sur $S$ (les ${}_{\ell^{\nu}}A^{o}$ étant finis \underline{étales} sur $S$). Comme le foncteur

\begin{equation*}
X \longmapsto X_{o} = X \times_{S} s
\end{equation*}
de la catégorie des schémas finis étales sur $S$ dans celle des schémas finis étales sur $k$ est une équivalence de catégories, $S$ étant hensélien, et qu'on a donc une équivalence analogue pour les catégories de faisceaux $\ell$-adiques constants tordus sur $S$ resp. sur $k$, la relation (2.2.3.3) donnant la fibre spéciale de $T_{\ell}(A^{o})^{f}$ caractérise donc cette partie fixe à isomorphisme unique près. D'autre part, en tant que sous-faisceau $\ell$-adique de $T_{\ell}(A^{o})$, elle est évidemment caractérisée par sa fibre générique (2.2.3.4), ou ce qui revient au même, par le sous $\mathbb{Z}_{\ell}$-module correspondant, stable par $\pi = \mathrm{Gal}(\overline{K}/K)$ :

\begin{equation}\tag{2.2.4.1}
T_{\ell}(A_{K})^{f}(\overline{K}) \quad \stackrel{\mathrm{dfn}}{=} \quad T_{\ell}(A(\overline{K}))^{f} \subset T_{\ell}(A(\overline{K})) \qquad .
\end{equation}
En fait, on a le résultat suivant :

\underline{Proposition} 2.2.5. \underline{Supposons comme ci-dessus} $\ell \neq p$, \underline{et} $S$ \underline{hensélien.} \underline{Alors le sous-module} (2.2.4.1) \underline{de} $T_{\ell}(A(\overline{K})) = T_{\ell}(A_{K})(\overline{K})$ \underline{correspondant à la partie fixe} $T_{\ell}(A_{K})^{f}$ (2.2.3.4) \underline{de} $T_{\ell}(A_{K})$ \underline{est le sous-module} $T_{\ell}(A(\overline{K}))^{I}$ \underline{des invariants sous le groupe d'inertie} $I \subset \pi$.

Ceci résultera, par passage à la limite projective, de la relation analogue, valable pour tout entier $m > 0$ premier à $p = \mathrm{car}.k$ :

\begin{equation}\tag{2.2.5.1}
({}_{m}A)^{f}(\overline{K}) = ({}_{m}A(\overline{K}))^{I} \qquad .
\end{equation}

\SGApage{I}{333}


\SGAvar{Pour prouver cette relation, introduisons le localisé strict $\tilde{S}$ de $S$, de corps des fractions $\tilde{K}$ la sous-extension non ramifiée maximale de $\overline{K}$ sur $K$. Comme $({}_{m}A)^{f}$ est fini étale sur $S$, ses points à valeurs dans $\overline{K}$ sont déjà à valeurs dans $\tilde{K}$, donc le premier membre de (2.2.5.1) n'est autre que $\mathrm{Hom}_{S}(\tilde{S} , ({}_{m}A)^{f})$, tandis que le deuxième est, par définition de $\tilde{K}$,égal à ${}_{m}A(\tilde{K})$. Or le premier membre, par définition de $({}_{m}A)^{f}$, compte tenu que $\tilde{S}$ est une limite projective de schémas $S_{i}$ étales finis sur $S$, est aussi égal à $\mathrm{Hom}_{S}(\tilde{S} , {}_{m}A) = {}_{m}\mathrm{Hom}_{S}(\tilde{S},A)$, et par définition du modèle de Néron on a $\mathrm{Hom}_{S}(\tilde{S},A) = A(\tilde{K})$, ce qui donne bien l'égalité annoncée (6.2.5.1).

}{Pour prouver cette relation, introduisons le localisé strict $\tilde{S}$ de $S$, de corps des fractions $\tilde{K}$ la sous-extension non ramifiée maximale de $\overline{K}$ sur $K$. Comme $({}_{m}A)^{f}$ est fini étale sur $S$, ses points à valeurs dans $\overline{K}$ sont déjà à valeurs dans $\tilde{K}$, donc le premier membre de (2.2.5.1) n'est autre que $\mathrm{Hom}_{S}(\tilde{S} , ({}_{m}A)^{f})$, tandis que le deuxième est, par définition de $\tilde{K}$,égal à ${}_{m}A(\tilde{K})$. Or le premier membre, par définition de $({}_{m}A)^{f}$, compte tenu que $\tilde{S}$ est une limite projective de schémas $S_{i}$ étales finis sur $S$, est aussi égal à $\mathrm{Hom}_{S}(\tilde{S} , {}_{m}A) = {}_{m}\mathrm{Hom}_{S}(\tilde{S},A)$, et par définition du modèle de Néron on a $\mathrm{Hom}_{S}(\tilde{S},A) = A(\tilde{K})$, ce qui donne bien l'égalité annoncée (2.2.5.1).

}{F7I-IX-016}Pour en déduire 2.2.5, on note qu'on déduit de (2.2.5.1) que

\begin{equation}\tag{*}
T_{\ell}(A(\overline{K}))^{I} \;=\; \varprojlim_{\nu} ({}_{\ell^{\nu}}A)^{f}(\overline{K}) \supset \lim\ ({}_{\ell^{\nu}}A^{o})^{f}(\overline{K}) \quad ;
\end{equation}
mais si $(x_{\nu})$ est un élément du deuxième membre, alors pour tout entier $\nu \geq 0$, $x_{\nu}$ est une section de $A$ qui est infiniment divisible par $\ell$, donc sa valeur en $s$ est un élément de $A_{o}(k)$ infiniment divisible par $\ell$, donc il est dans $A_{o}^{\,o}(k)$, ce qui prouve l'égalité des deux derniers membres de (*) et achève de prouver 2.2.5.

\underline{Proposition} 2.2.6. \underline{Soient} $C_{K}$ \underline{un deuxième schéma abélien sur} $K$, \underline{de modèle de Néron} $C$, \underline{et considérons un homomorphisme} $u_{K}:C_{K} \longrightarrow A_{K}$, \underline{se prolongeant donc en un homomorphisme} $u: C \longrightarrow A$. \underline{Supposons que} $u_{K}$ \underline{soit une} isogénie. \underline{Alors pour tout nombre premier} $\ell,T_{\ell}(u_{K}):T_{\ell}(C_{K}) \longrightarrow T_{\ell}(A_{K})$ \underline{induit une isogénie des parties fixes}

\begin{equation}\tag{2.2.6.1}
T_{\ell}(C_{K})^{f} \longrightarrow T_{\ell}(A_{K})^{f} \quad ,
\end{equation}

\SGApage{I}{334}


\underline{et de même le morphisme induit par} $u_{o} : C_{o} \longrightarrow A_{o}$

\begin{equation}\tag{2.2.6.2}
T_{\ell}(u_{o}^{\,o}) : T_{\ell}(C_{o}^{\,o}) \longrightarrow T_{\ell}(A_{o}^{\,o})
\end{equation}
\underline{est une isogénie.} \underline{Enfin,} $u_{o} : C_{o} \longrightarrow A_{o}$ \underline{induit une isogénie du plus grand sous-tore de} $C_{o}$ \underline{dans le plus grand sous-tore} $T_{o}$ \underline{de} $A_{o}$, \underline{et au-dessus d'une extension} $k'$ \underline{de} $k$ \underline{sur laquelle les} ``\underline{parties abéliennes}'' \underline{de} $C_{o}$ \underline{et de} $A_{o}$ \underline{sont définies} (\underline{par exemple sur la clôture parfaite de} $k$), $u_{o}$ \underline{induit une isogénie entre les parties abéliennes de} $C_{ok'}$ \underline{et de} $A_{ok'}$ .

Cet énoncé est conséquence formelle du fait que $T_{\ell}(A_{K})^{f}$ , $T_{\ell}(A_{o}^{\,o})$, le plus grand sous-tore de $A_{o}$, enfin la ``partie abélienne'' de $A_{ok'}$ , sont pour $A_{K}$ variable des \underline{foncteurs additifs} en $A_{K}$, et que le fait que $u_{K}$ soit une isogénie s'exprime par la condition qu'il existe un morphisme $v_{K}:A_{K} \longrightarrow C_{K}$ tel que l'on ait $u_{K}v_{K} = n.\mathrm{id}$, $v_{K}u_{K} = n.\mathrm{id}$, avec $n$ un entier $> 0$.

Notons en particulier :

\SGAvar{\underline{Corollaire} 2.2.7. \underline{Si} $A_{K}$ \underline{et} $C_{K}$ \underline{sont des schémas abéliens isogénes sur} $K$, \underline{alors les invariants correspondants} $\mu$, $\lambda$, $\alpha$ (2.1.11) \underline{sont les mêmes pour} $A_{K}$ \underline{et pour} $C_{K}$ . \underline{En particulier, ils sont les mêmes pour} $A_{K}$ \underline{et son dual} $A'_{K}$ .

}{\underline{Corollaire} 2.2.7. \underline{Si} $A_{K}$ \underline{et} $C_{K}$ \underline{sont des schémas abéliens isogènes sur} $K$, \underline{alors les invariants correspondants} $\mu$, $\lambda$, $\alpha$ (2.1.11) \underline{sont les mêmes pour} $A_{K}$ \underline{et pour} $C_{K}$ . \underline{En particulier, ils sont les mêmes pour} $A_{K}$ \underline{et son dual} $A'_{K}$ .

}{F7I-IX-017}\underline{Remarques} 2.2.8. Dans les deux énoncés précédents 2.2.6 et 2.2.7, on a supposé $S$ hensélien. Il est d'autre part immédiat que 2.2.7, ainsi que les assertions de 2.2.6 concernant les propriétés de l'homomorphisme induit $u_{o} : C_{o} \longrightarrow A_{o}$ , sont valables sans cette hypothèse. On se ramène en effet au cas traité par le changement de base $S^{h} \longrightarrow S$, où $S^{h}$ est le hensélisé de $S$, compte tenu du fait (immédiat en vertu des définitions) que la formation des modèles de Néron commute à ce changement de base.


\SGApage{I}{335}

\underline{Corollaire} 2.2.9. (Cas de \underline{bonne réduction.}) \underline{Soient} $S$ \underline{un trait}, $A_K$ \underline{un schéma abélien sur le corps des fractions} $K$ \underline{de} $S$, $A$ \underline{le modèle de Néron de} $A_K$, $\ell$ \underline{un nombre premier} $\neq$ car.k . \underline{Les conditions suivantes sont équivalentes.}

(i) \underline{Il existe un schéma abélien sur} $S$ \underline{dont la fibre générique est isomorphe à} $A_K$ .

(ii) $A$ \underline{est un schéma abélien.}

(ii$^o$) $A^o$ \underline{est un schéma abélien.}

(iii) $A_o$ \underline{est un schéma abélien.}

\SGAvar{(iii$^o$) $A_o^{\,o}$ \underline{est un schéma abélien}, \underline{i.e.} (\underline{avec les notations de} (2.1.11) \underline{on a} $\lambda = \mu = 0$ .

}{(iii$^o$) $A_o^{\,o}$ \underline{est un schéma abélien}, \underline{i.e.} (\underline{avec les notations de} (2.1.11)) \underline{on a} $\lambda = \mu = 0$ .

}{F7I-IX-018}(iv) $A$ \underline{est propre sur} $S$.

(iv$^o$) $A^o$ \underline{est propre sur} $S$.

(v) $T_\ell(A_K)$ \underline{est} "\underline{non ramifié sur} $S$" \underline{i.e.} \underline{se prolonge en un faisceau} $\ell$\underline{-adique constant tordu sur} $S$, \underline{i.e.} \underline{le groupe d'inertie} $I$ \underline{opère trivialement sur} $T_\ell(A_K(\overline{K}))$.

(v bis) \underline{Pour tout entier} $\nu \geq 0$, ${}_{\ell^\nu}A_K$ \underline{est} "\underline{non ramifié sur} $S$" \underline{i.e. se prolonge en un revêtement étale sur} $S$, \underline{i.e.} \underline{le groupe d'inertie} $I$ \underline{opère trivialement sur} ${}_{\ell^\nu}A_K(\overline{K})$.

\underline{Enfin}, \underline{un schéma abélien sur} $S$ \underline{comme dans} (i) \underline{est un modèle de Néron de} $A_K$ , \underline{et en particulier est déterminé à isomorphisme unique près.}

2.2.9.1. On dit que $A_K$ a \underline{bonne réduction} sur $S$ s'il satisfait la condition (i) ci-dessus. Le critère (v) ou (v bis) est connu sous le nom de \underline{critère de} NERON-OGG-CHAFAREVITCH de bonne réduction.


\SGApage{I}{336}

Prouvons 2.2.9. Le fait qu'un schéma abélien sur $S$ satisfait à la propriété universelle d'un modèle de Néron, qui se réduit au fait qu'une section rationnelle d'un schéma abélien sur une base régulière est partout définie, est bien connu (cf. p.ex. [5]). Cela prouve l'équivalence (i) $\Longleftrightarrow$ (ii) $\Longleftrightarrow$ (ii$^o$), ainsi que la dernière assertion de 2.2.9. Les conditions envisagées impliquent évidemment (iii), (iv) (iv$^o$), qui impliquent chacune (iii$^o$), prouvons que (iii$^o$) $\Longrightarrow$ (ii$^o$) (ce qui prouvera l'équivalence des conditions (i) à (iv$^o$)). Or cela résulte du fait qu'un schéma en groupes lisse sur un schéma $S$, à fibres des schémas abéliens, est propre sur $S$ (donc est un schéma abélien) en vertu de (EGA IV 15.7.10). D'autre part, compte tenu de 2.2.5 et du fait que $T_\ell(A_K(\overline{K}))/T_\ell(A_K(\overline{K}))^I$ est manifestement un $\mathbb{Z}_\ell$-module sans torsion, la condition (v) s'exprime par l'égalité des rangs de $T_\ell(A_K(\overline{K}))$ et de sa "partie fixe", i.e. (par (2.1.11) et (2.2.3.3)) par la condition

\begin{equation*}
\mu + 2\alpha = 2(\lambda + \mu + \alpha) \qquad ,
\end{equation*}
i.e. $2\lambda + \mu = 0$, i.e. $\lambda = \mu = 0$, qui n'est autre que (iii$^o$). Enfin l'équivalence de (v) et (v bis) est claire par définition. Cela achève la démonstration de 2.2.9.

2.3. Nous supposons à nouveau $S$ hensélien, et $\ell \neq p =$ car. k. Utilisons maintenant l'inclusion

\begin{equation}\tag{2.3.1}
T_\ell(T_o) \hookrightarrow T_\ell(A_o^{\,o})
\end{equation}
de (2.1.6), et l'équivalence de catégories rappelée dans 2.2.4. On trouve un sous-faisceau $\ell$-adique constant tordu canonique de $T_\ell(A^{\,o})^f$,


\SGApage{I}{337}

\SGAvar{appelée la \underline{partie torique de} $T_\ell(A^o)$,

}{appelé la \underline{partie torique de} $T_\ell(A^o)$,

}{F7I-IX-019}\begin{equation}\tag{2.3.2}
T_\ell(A^o)^t \hookrightarrow T_\ell(A^o)^f \hookrightarrow T_\ell(A^o) \quad ,
\end{equation}
caractérisé par la condition que sa fibre spéciale soit donnée par l'inclusion (2.3.1) :

\begin{equation}\tag{2.3.3}
T_\ell(A^o)^t \times_S s = T_\ell(T_o) \hookrightarrow T_\ell(A_o^{\,o}) \quad .
\end{equation}
Nous donnerons d'ailleurs plus loin une interprétation de la partie torique de $T_\ell(A^o)$ comme le module de Tate $T_\ell(T)$ d'un tore $T$ sur $S$ qui relève le tore $T_o$ sur $k$, interprétation qui permettra d'étendre la définition de la partie torique au cas $\ell = p$. Pour l'instant, et pour fixer les idées, nous nous bornons à examiner le cas $\ell \neq p$, auquel cas la partie torique de $T_\ell(A^o)$ peut donc se décrire par le sous-$\mathbb{Z}_\ell$-module correspondant de $T_\ell(A(\overline{K}))$ :

\begin{equation}\tag{2.3.4}
T_\ell(A(\overline{K}))^t \hookrightarrow T_\ell(A(\overline{K}))^f \hookrightarrow T_\ell(A(\overline{K})) \quad ,
\end{equation}
stable par l'action de $\pi$, où le premier terme est défini comme $T_\ell(A^o)^t(\overline{K})$, et est appelé encore comme de juste le \underline{sous-module torique} du module de Tate $T_\ell(A(\overline{K}))$. Ce dernier apparaît donc comme muni d'une filtration canonique à trois crans (2.3.4). Appliquant ces réflexions au schéma abélien dual $A'_K$ , on trouve de même une filtration à trois crans

\SGAvar{\begin{equation}\tag{2.3.5}
T_\ell(A'(\overline{K}))^t \hookrightarrow T_\ell(A'(\overline{K}))^f \hookrightarrow T_\ell(A'(\overline{K})) \quad .
\end{equation}
Les termes médians dans (2.3.4) et (2.3.5) (``parties fixes'')étant décrits en termes des structures de $\pi$-modules des modules de Tate de $A(\overline{K})$ , $(A'(\overline{K})$ comme les modules d'invariants sous le groupe d'inertie $I$ (1.6.5), les parties toriques sont complètement décrites en termes des parties fixes


}{\begin{equation}\tag{2.3.5}
T_\ell(A'(\overline{K}))^t \hookrightarrow T_\ell(A'(\overline{K}))^f \hookrightarrow T_\ell(A'(\overline{K})) \quad .
\end{equation}
Les termes médians dans (2.3.4) et (2.3.5) (``parties fixes'')étant décrits en termes des structures de $\pi$-modules des modules de Tate de $A(\overline{K})$ , $A'(\overline{K})$ comme les modules d'invariants sous le groupe d'inertie $I$ (2.2.5), les parties toriques sont complètement décrites en termes des parties fixes


}{F7I-IX-020}\SGApage{I}{338}

\SGAvar{et de l'accouplement $\varphi$ de (1.0.2), gra\^ce au théorème suivant (dont un cas particulier se trouve déjà dans [11]).

\underline{Thoérème} 2.4. (\underline{Théorème d'orthogonalité}). \underline{Soient} $S$ \underline{un trait hensélien,} $A_K$ \underline{un schéma abélien sur son corps des fractions} $K$, $A'_K$ \underline{le schéma abélien dual,} $\ell$ \underline{un nombre premier distinct de la caractéristique résiduelle de} $S$, \underline{d'où sur les modules de Tate} $T_\ell(A\ (\overline{K}))$ \underline{et} $T_\ell(A'(\overline{K}))$ \underline{la filtration canonique} (2.3.4) \underline{resp.} (2.3.5) \underline{par la} ``\underline{partie fixe}'' \underline{et par la} ``\underline{partie torique}''. \underline{Alors la partie torique} $T_\ell(A(\overline{K}))^t$ \underline{est l'intersection de la partie fixe} $T_\ell(A(\overline{K}))^f$ \underline{avec l'orthogonal de la partie fixe} $T_\ell(A'(\overline{K}))^f = T_\ell(A'(\overline{K}))^I$ \underline{pour l'accouplement canonique} $\varphi$ (1.0.3), \underline{donc symétriquement}, \underline{la partie torique} $T_\ell(A'(\overline{K}))^t$ \underline{est l'intersection de} $T_\ell(A'(\overline{K}))^f$ \underline{avec l'orthogonal de la partie fixe} $T_\ell(A(\overline{K}))^f = T_\ell(A(\overline{K}))^I$. (NB \underline{l'exposant} $I$ \underline{dénote le sous-module des invariants sous l'action du groupe d'inertie} $I$.)

\underline{Démonstration} \quad a) \quad Prouvons que $T_\ell(A(\overline{K}))^t$ est orthogonal à $T_\ell(A'(\overline{K}))^f$. Pour ceci, nous allons interpréter la restriction de la forme $\varphi$ aux parties fixes des modules de Tate, à l'aide de la biextension canonique $W^o = W$ de (1.4.4), prolongement canonique de la biextension $W_K$ de $(A_K ; A'_K)$ par $\mathbb{G}_{mK}$ en une biextension de $(A^o , {A'}^o)$ par $\mathbb{G}_{mS}$ . En vertu de VIII 2.2, à $W$ est associé un accouplement

\begin{equation}\tag{2.4.1}
\varphi_w : T_\ell(A^o) \times T_\ell({A'}^o) \longrightarrow T_\ell(\mathbb{G}_{mS}) \quad ,
\end{equation}
dont l'accouplement $\varphi$ de (1.0.3) est déduit en passant aux points à valeurs dans $\overline{K}$ . Il en résulte que la restriction de $\varphi$ aux parties fixes envisagées dans 2.4 s'obtient à partir de l'accouplement

}{et de l'accouplement $\varphi$ de (1.0.2), grâce au théorème suivant (dont un cas particulier se trouve déjà dans [11]).

\underline{Théorème} 2.4. (\underline{Théorème d'orthogonalité}). \underline{Soient} $S$ \underline{un trait hensélien,} $A_K$ \underline{un schéma abélien sur son corps des fractions} $K$, $A'_K$ \underline{le schéma abélien dual,} $\ell$ \underline{un nombre premier distinct de la caractéristique résiduelle de} $S$, \underline{d'où sur les modules de Tate} $T_\ell(A\ (\overline{K}))$ \underline{et} $T_\ell(A'(\overline{K}))$ \underline{la filtration canonique} (2.3.4) \underline{resp.} (2.3.5) \underline{par la} ``\underline{partie fixe}'' \underline{et par la} ``\underline{partie torique}''. \underline{Alors la partie torique} $T_\ell(A(\overline{K}))^t$ \underline{est l'intersection de la partie fixe} $T_\ell(A(\overline{K}))^f$ \underline{avec l'orthogonal de la partie fixe} $T_\ell(A'(\overline{K}))^f = T_\ell(A'(\overline{K}))^I$ \underline{pour l'accouplement canonique} $\varphi$ (1.0.3), \underline{donc symétriquement}, \underline{la partie torique} $T_\ell(A'(\overline{K}))^t$ \underline{est l'intersection de} $T_\ell(A'(\overline{K}))^f$ \underline{avec l'orthogonal de la partie fixe} $T_\ell(A(\overline{K}))^f = T_\ell(A(\overline{K}))^I$. (NB \underline{l'exposant} $I$ \underline{dénote le sous-module des invariants sous l'action du groupe d'inertie} $I$.)

\underline{Démonstration} \quad a) \quad Prouvons que $T_\ell(A(\overline{K}))^t$ est orthogonal à $T_\ell(A'(\overline{K}))^f$. Pour ceci, nous allons interpréter la restriction de la forme $\varphi$ aux parties fixes des modules de Tate, à l'aide de la biextension canonique $W^o = W$ de (1.4.4), prolongement canonique de la biextension $W_K$ de $(A_K , A'_K)$ par $\mathbb{G}_{mK}$ en une biextension de $(A^o , {A'}^o)$ par $\mathbb{G}_{mS}$ . En vertu de VIII 2.2, à $W$ est associé un accouplement

\begin{equation}\tag{2.4.1}
\varphi_W : T_\ell(A^o) \times T_\ell({A'}^o) \longrightarrow T_\ell(\mathbb{G}_{mS}) \quad ,
\end{equation}
dont l'accouplement $\varphi$ de (1.0.3) est déduit en passant aux points à valeurs dans $\overline{K}$ . Il en résulte que la restriction de $\varphi$ aux parties fixes envisagées dans 2.4 s'obtient à partir de l'accouplement

}{F7I-IX-021}\begin{equation}\tag{2.4.2}
T_\ell(A^o)^f \times T_\ell({A'}^o)^f \longrightarrow T_\ell(\mathbb{G}_{mS})
\end{equation}

\SGApage{I}{339}

restriction de (2.4.1), en passant aux points à valeurs dans $\overline{K}$ . Comme les deux termes qui interviennent dans le premier membre de (2.4.2) sont maintenant des faisceaux $\ell$-adiques constants tordus, pour étudier ce dernier accouplement, on peut se borner à étudier l'accouplement induit sur les fibres spéciales. Or les fibres des deux termes envisagés ne sont autres que les modules de Tate $T_{\ell}(A_{o}^{\,o}),T_{\ell}({A'}^{o}_{o})$ (2.2.3.3), et l'accouplement déduit de (2.4.2) par passage aux fibres spéciales n'est autre que l'accouplement

\begin{equation}\tag{2.4.3}
T_{\ell}(A_{o}^{\,o}) \times T_{\ell}({A'}^{o}_{o}) \longrightarrow T_{\ell}(\mathbb{G}_{mk})
\end{equation}
déduit de la biextension $W_{o}$ de $(A_{o}^{\,o},{A'}^{o}_{o})$ par $\mathbb{G}_{mk}$ , restriction de $W$ aux fibres spéciales.

Rappelons maintenant (2.3) que la partie torique de $T_{\ell}(A(\overline{K}))$ s'obtient en passant aux points à valeurs dans $\overline{K}$ dans le sous-faisceau $\ell$-adique constant tordu $T_{\ell}(A^{o})^{t}$ de $T_{\ell}(A^{o})^{f}$, donc la relation d'orthogonalité à prouver s'exprime en disant que l'accouplement induit par (2.4.2)

\begin{equation}\tag{2.4.4}
T_{\ell}(A^{o})^{t} \times T_{\ell}({A'}^{o})^{f} \longrightarrow T_{\ell}(\mathbb{G}_{mS})
\end{equation}
est identiquement nul. Or il suffit de le voir pour l'accouplement correspondant sur les fibres spéciales, qui par la définition de la partie torique n'est autre que l'accouplement

\SGAvar{\begin{equation}\tag{2.4.5}
T_{\ell}(T_{o}) \times T_{\ell}(A_{o}^{\,o}) \longrightarrow T_{\ell}(\mathbb{G}_{mk})
\end{equation}
induit par (2.4.3). Or en vertu de VIII 4.10 cet accouplement est nécessairement nul, ce qui prouve l'orthogonalité annoncée.


}{\begin{equation}\tag{2.4.5}
T_{\ell}(T_{o}) \times T_{\ell}(A_{o}^{\prime\,o}) \longrightarrow T_{\ell}(\mathbb{G}_{mk})
\end{equation}
induit par (2.4.3). Or en vertu de VIII 4.10 cet accouplement est nécessairement nul, ce qui prouve l'orthogonalité annoncée.


}{F7I-IX-022}\SGApage{I}{340}

b)  Prouvons que l'inclusion

\SGAvar{\begin{equation}\tag{2.4.6}
T_{\ell}(A(\overline{K}))^{t} \subset T_{\ell}(A(\overline{K}))^{f} \cap (T_{\ell}(A'(\overline{K}))^{f})^{\perp}
\end{equation}
est une égalité, où le signe $\perp$ désigne l'orthogonal. Il revient au même de dire que dans l'accouplement (2.4.2), l'annulateur du deuxième facteur, qui contient la partie torique $T_{\ell}(A^{o})^{t}$ comme on vient de voir, est \underline{égale} à cette dernière. Comme il s'agit encore de relations entre faisceaux $\ell$-adiques constants tordus, il suffit de vérifier cette assertion sur les fibres spéciales, i.e. que l'inclusion déjà obtenue (1.10)

}{\begin{equation}\tag{2.4.6}
T_{\ell}(A(\overline{K}))^{t} \subset T_{\ell}(A(\overline{K}))^{f} \cap (T_{\ell}(A'(\overline{K}))^{f})^{\perp}
\end{equation}
est une égalité, où le signe $\perp$ désigne l'orthogonal. Il revient au même de dire que dans l'accouplement (2.4.2), l'annulateur du deuxième facteur, qui contient la partie torique $T_{\ell}(A^{o})^{t}$ comme on vient de voir, est \underline{égal} à cette dernière. Comme il s'agit encore de relations entre faisceaux $\ell$-adiques constants tordus, il suffit de vérifier cette assertion sur les fibres spéciales, i.e. que l'inclusion déjà obtenue (2.4.7)

}{F7I-IX-023}\begin{equation}\tag{2.4.7}
T_{\ell}(T_{o}) \subset T_{\ell}({A'}^{o}_{o})^{\perp}
\end{equation}
est une égalité ; en d'autres termes, il faut vérifier que la biextension $W_{o}$ de $(A_{o}^{\,o},{A'}^{o}_{o})$ est \underline{non dégénérée à gauche} (VIII 4.11). Or c'est ce que nous avons déjà noté dans 1.5, comme conséquence d'un résultat de la thèse de RAYNAUD.

2.5.  On peut donner une forme équivalente au théorème d'orthogonalité, en utilisant une polarisation $\xi$ de $A_{K}$ , et la forme bilinéaire alternée (1.0.8), compatible aux opérations de $\pi$ ,

\begin{equation}\tag{2.5.1}
\varphi_{\xi} : T_{\ell}(A(\overline{K})) \times T_{\ell}(A(\overline{K})) \longrightarrow T_{\ell}(\overline{K}*) \quad ,
\end{equation}
déduite de la forme $\varphi$ (1.0.3) par transport de structure sur le deuxième argument à l'aide de l'isogénie de $\mathbb{Z}_{\ell}$-modules

\begin{equation*}
T_{\ell}(A(\overline{K})) \longrightarrow T_{\ell}(A'(\overline{K}))
\end{equation*}
induite par l'isogénie $\psi_{\xi} : A_{K} \longrightarrow A'_{K}$ définie par $\xi$ . L'égalité dans (2.4.6) prend alors la forme simple


\SGApage{I}{341}

\begin{equation}\tag{2.5.2}
W = V \cap V^{\perp} \quad ,
\end{equation}
où on a posé

\begin{equation}\tag{2.5.3}
\left\{
\begin{array}{l}
U = T_{\ell}(A(\overline{K})) \quad , \quad V = U^{I} = \text{"partie fixe" de } T_{\ell}(A(\overline{K})) \\[3mm]
\hspace{4em} W = \text{partie torique de } T_{\ell}(A(\overline{K})) \quad ,
\end{array}
\right.
\end{equation}
et où $\perp$ désigne l'orthogonal par rapport à la forme de polarisation (2.5.1). Les modules précédents donnent alors lieu au diagramme d'inclusions

\begin{equation}\tag{2.5.4}
\begin{tikzcd}[column sep=large, row sep=large]
\& \& V \arrow[dl, hook', "2\lambda"'] \& \& \\
U \& W^{\perp} = V + V^{\perp} \arrow[l, hook', "\mu"'] \& \& W = V \cap V^{\perp} \arrow[ul, hook', "2\alpha"'] \arrow[dl, hook, "2\lambda"] \& 0 \arrow[l, hook', "\mu"'] \\
\& \& V^{\perp} \arrow[ul, hook, "2\alpha"] \& \&
\end{tikzcd} \quad ,
\end{equation}



















\SGAvar{où les entiers $2\alpha$, $2\lambda$, $\mu$ désignent les rangs du quotients de deux modules consécutifs dans le diagramme, les notations étant celles de (2.1.9). (Le fait que $W$ soit de rang $\mu$ est la définition de $\mu$, que $V/W$ soit de rang $2\alpha$ est la deuxième relation (2.1.11) jointe à la suite exacte (2.1.6), d'où par orthogonalité rang $U/W^{\perp} = \mu$ et rang $W^{\perp}/V^{\perp} = 2\alpha$ ; enfin les formules

\begin{equation}\tag{2.5.5}
\mathrm{rang}\ V/V \cap V^{\perp} \;=\; \mathrm{rang}\ V/W = 2\lambda \quad ,
\end{equation}
et la formule duale, résultent des précédentes, compte tenu de (2.1.10) et de rang $U = n$.)

}{où les entiers $2\alpha$, $2\lambda$, $\mu$ désignent les rangs des quotients de deux modules consécutifs dans le diagramme, les notations étant celles de (2.1.9). (Le fait que $W$ soit de rang $\mu$ est la définition de $\mu$, que $V/W$ soit de rang $2\alpha$ est la deuxième relation (2.1.11) jointe à la suite exacte (2.1.6), d'où par orthogonalité rang $U/W^{\perp} = \mu$ et rang $W^{\perp}/V^{\perp} = 2\alpha$ ; enfin les formules

\begin{equation}\tag{2.5.5}
\mathrm{rang}\ V^{\perp}/V \cap V^{\perp} \;=\; \mathrm{rang}\ V^{\perp}/W = 2\lambda \quad ,
\end{equation}
et la formule duale, résultent des précédentes, compte tenu de (2.1.10) et de rang $U = 2n$.)

}{F7I-IX-024}\underline{Remarques} 2.6.  a) Nous avons obtenu ici le théorème d'orthogonalité 2.4 comme conséquence facile de la théorie des biextensions développée dans les numéros précédents. On peut démontrer l'\underline{inclusion} (2.4.6) sans utiliser la théorie des biextensions, par un argument arithmétique dans


\SGApage{I}{342}

le style de la démonstration du "théorème de monodromie" donnée dans I, en se ramenant à une situation où $S$ est de corps résiduel fini, et où le résultat d'orthogonalité résulte d'une simple considération de \underline{poids} pour l'opération de frobenius (utilisant le résultat de A. WEIL sur les valeurs propres dudit frobenius). (Il est inutile d'invoquer ici la résolution des singularités, en énonçant le résultat voulu pour des schémas en groupes $A$ sur des schémas $S$, sous des conditions suffisamment générales.) Il ne semble pas malheureusement qu'on puisse obtenir par cette méthode l'égalité dans l'inclusion précédente, qui sera utilisée, de façon essentielle semble-t-il, dans la démonstration du "théorème de réduction semi-stable" au n° suivant.

b) Le théorème d'orthogonalité 2.4 permet d'expliciter les propriétés les plus importantes du modèle de Néron d'une variété abélienne $A_{K}$ en termes de la représentation du groupe d'inertie $I$ sur le module de Tate $T_{\ell}(A(\overline{K}))$, $\ell$ étant un nombre premier fixé distinct de la caractéristique résiduelle. (Le cas $\ell = p$ sera examiné à part plus loin (n° 5)). Comme premier énoncé de ce type, mais n'utilisant pas le théorème d'orthogonalité, rappelons le critère de bonne réduction 2.2.9. Nous allons donner de nombreux autres exemples par la suite, en commençant au n° suivant (3.5).


\SGApage{I}{343}


\SGAsectiontoc{3. Cas de la réduction semi-stable. Le théorème de réduction semi-stable.}{3}
3. \underline{Cas de la réduction semi-stable. Le théorème de réduction semi-stable.}

Donnons d'abord un résultat auxiliaire dans le style de SGA 3 :

\SGAvar{\underline{Proposition} 3.1. \underline{Soient} $S$ \underline{un schéma}, $G$ \underline{un} $S$-\underline{schéma en groupes commutatifs}, \underline{plat}, \underline{séparé}, \underline{de présentation finie sur} $S$, $H$ \underline{un schéma en groupes quasi-séparé et localement de type fini sur} $S$, $u: G \longrightarrow H$ \underline{un homomorphisme de} $S$-\underline{schémas en groupes}, $K = \mathrm{Ker}\ u$ . \underline{Pour tout} $s \in S$, \underline{désignant par} $\bar{s}$ \underline{le spectre d'une clôture algébrique de} $k(s)$, \underline{on note} $\mu(s)$ (\underline{resp.} $\alpha(s)$) \underline{le rang réductif} (\underline{resp. le rang abélien}) \underline{de} $K_{s}$ , \underline{i.e.} \underline{la dimension du plus grand sous-tore} (\underline{resp. de la plus grande variété abélienne quotient}) \underline{de} $(K_{\bar{s}})^{o}_{\mathrm{red}}$. \underline{Soient} $s \in S$, $t \in S$ \underline{une générisation de} $S$, $p$ \underline{l'exposant caractéristique de} $k(s)$.

}{\underline{Proposition} 3.1. \underline{Soient} $S$ \underline{un schéma}, $G$ \underline{un} $S$-\underline{schéma en groupes commutatifs}, \underline{plat}, \underline{séparé}, \underline{de présentation finie sur} $S$, $H$ \underline{un schéma en groupes quasi-séparé et localement de type fini sur} $S$, $u: G \longrightarrow H$ \underline{un homomorphisme de} $S$-\underline{schémas en groupes}, $K = \mathrm{Ker}\ u$ . \underline{Pour tout} $s \in S$, \underline{désignant par} $\bar{s}$ \underline{le spectre d'une clôture algébrique de} $k(s)$, \underline{on note} $\mu(s)$ (\underline{resp.} $\alpha(s)$) \underline{le rang réductif} (\underline{resp. le rang abélien}) \underline{de} $K_{s}$ , \underline{i.e.} \underline{la dimension du plus grand sous-tore} (\underline{resp. de la plus grande variété abélienne quotient}) \underline{de} $(K_{\bar{s}})^{o}_{\mathrm{red}}$. \underline{Soient} $s \in S$, $t \in S$ \underline{une générisation de} $s$, $p$ \underline{l'exposant caractéristique de} $k(s)$.

}{F7I-IX-025}a) \underline{Pour tout entier} $n > 0$ \underline{premier à} $p$, ${}_{n}K$ \underline{est étale de type fini et séparé sur} $S$ \underline{au voisinage de} $s$, \underline{et on a la divisibilité} :

\begin{equation}\tag{3.1.1}
\mathrm{rang}({}_{n}K_{s}) \mid \mathrm{rang}({}_{n}K_{t}) \quad .
\end{equation}
b) \underline{On a}

\begin{equation}\tag{3.1.2}
\mu(s) + 2\alpha(s) \leq \mu(t) + 2\alpha(t) \quad .
\end{equation}
c) \underline{Si} $(K_{\bar{t}})^{o}_{\mathrm{red}}$ \underline{est unipotent}, \underline{il en est de même de} $(K_{\bar{s}})^{o}_{\mathrm{red}}$ .

d) \underline{Supposons} $G_{s}$ \underline{de rang unipotent nul}, $H$ \underline{lisse sur} $S$, \underline{et} $u_{t}$ \underline{une isogénie}. \underline{Alors} $u$ \underline{est plat et quasi-fini au-dessus d'un voisinage ouvert} $U$ \underline{de} $s$, \underline{a fortiori} $K|U$ \underline{est plat et quasi-fini.}

e) \underline{Sous les conditions de} d), \underline{si} $u_{t}$ \underline{est de plus un monomorphisme} (\underline{donc un isomorphisme de} $G_{t}$ \underline{sur un sous-groupe ouvert de} $H_{t}$), \underline{alors} $u_{s}$ \underline{est un isomorphisme de} $G_{s}$ \underline{sur un sous-groupe ouvert de} $H_{s}$,


\SGApage{I}{344}

\underline{et si de plus} $S$ \underline{est irréductible de point générique} $t$, \underline{alors il existe un voisinage ouvert} $U$ \underline{de} $s$ \underline{tel que} $u$ \underline{induise un isomorphisme de} $G|U$ \underline{sur un sous-groupe ouvert de} $H|U$.

\underline{Démonstration.}

\SGAvar{a) On peut supposer que $n$ est inversible sur $S$, et alors $n.\mathrm{id}_{G}$ est étale fibre par fibre (SGA 3 XV 1.3), donc, $G$ étant plat de présentation finie sur $S$, il résulte du critère d'étalité par fibres EGA IV 17.8.2 que $n.\mathrm{id}_{G}$ est étale, donc son noyau ${}_{n}G$ est étale ; il est clair d'ailleurs qu'il est séparé et de présentation finie sur $S$, puisque $G$ l'est. On voit de même que ${}_{n}H$ est localement de type fini, quasi-séparé et \underline{formellement net} sur $S$ (*) donc sa section unité est une immersion ouverte (EGA IV 17.4.2 b')), donc le noyau de ${}_{n}u: {}_{n}G \longrightarrow {}_{n}H$ est un sous-schéma ouvert de ${}_{n}G$, donc étale sur $S$, qui est retrocompact dans ${}_{n}G$ donc de présentation finie sur $S$, et séparé sur $S$ puisque $G$ l'est. Or ce noyau n'est évidemment autre que ${}_{n}K$ . La formule (3.1.1) résulte des propriétés précédentes de ${}_{n}K$, comme on voit en se ramenant par changement de base au cas où $S$ est local hensélien de point fermé $s$, et en utilisant alors la décomposition de ${}_{n}K$ en ${}_{n}K^{f}$ et un schéma à fibre spéciale vide (cf. 2.2.3). En fait, cet argument montre, plus précisément, qu'il existe un \underline{monomorphisme} pour les fibres géométriques

}{a) On peut supposer que $n$ est inversible sur $S$, et alors $n.\mathrm{id}_{G}$ est étale fibre par fibre (SGA 3 XV 1.3), donc, $G$ étant plat de présentation finie sur $S$, il résulte du critère d'étalité par fibres EGA IV 17.8.2 que $n.\mathrm{id}_{G}$ est étale, donc son noyau ${}_{n}G$ est étale ; il est clair d'ailleurs qu'il est séparé et de présentation finie sur $S$, puisque $G$ l'est. On voit de même que ${}_{n}H$ est localement de type fini, quasi-séparé et \underline{formellement net} sur $S$ (*) donc sa section unité est une immersion ouverte (EGA IV 17.4.2 b')), donc le noyau de ${}_{n}u: {}_{n}G \longrightarrow {}_{n}H$ est un sous-schéma ouvert de ${}_{n}G$, donc étale sur $S$, qui est rétrocompact dans ${}_{n}G$ donc de présentation finie sur $S$, et séparé sur $S$ puisque $G$ l'est. Or ce noyau n'est évidemment autre que ${}_{n}K$ . La formule (3.1.1) résulte des propriétés précédentes de ${}_{n}K$, comme on voit en se ramenant par changement de base au cas où $S$ est local hensélien de point fermé $s$, et en utilisant alors la décomposition de ${}_{n}K$ en ${}_{n}K^{f}$ et un schéma à fibre spéciale vide (cf. 2.2.3). En fait, cet argument montre, plus précisément, qu'il existe un \underline{monomorphisme} pour les fibres géométriques

}{F7I-IX-026}\begin{equation*}
{}_{n}K(\bar{s}) \lhook\joinrel\longrightarrow {}_{n}K(\bar{t}) \quad .
\end{equation*}
(*) Du moins si $H$ est commutatif ; dans le cas général, il restera vrai que ${}_{n}H$ est net sur $S$ le long de sa section unité, ce qui suffit pour notre propos.


\SGApage{I}{345}

\SGAvar{b)   Soit $n > 0$ premier à l'exposant caractéristique de $k(s)$ et aux rangs des quotients des fibres géométriques de $K$ en $s$ et $t$ par leurs composantes neutres réduites, alors les deux membres de (3.1.1) s'obtiennent en élévant $n$ respectivement à la puissance d'exposant l'un et l'autre membre de (3.1.2), comme il résulte de la structure connue des groupes algébriques commutatifs. Donc b) résulte de a).

c)   Comme l'hypothèse $(K_{\overline{t}})^{o}_{s}$ unipotent équivaut aux relations $\mu(t) = \alpha(t) = 0$, et de même pour $K_{s}$, l'assertion c) résulte formellement de b).

}{b)   Soit $n > 0$ premier à l'exposant caractéristique de $k(s)$ et aux rangs des quotients des fibres géométriques de $K$ en $s$ et $t$ par leurs composantes neutres réduites, alors les deux membres de (3.1.1) s'obtiennent en élevant $n$ respectivement à la puissance d'exposant l'un et l'autre membre de (3.1.2), comme il résulte de la structure connue des groupes algébriques commutatifs. Donc b) résulte de a).

c)   Comme l'hypothèse $(K_{\overline{t}})^{o}_{\mathrm{red}}$ unipotent équivaut aux relations $\mu(t) = \alpha(t) = 0$, et de même pour $K_{s}$, l'assertion c) résulte formellement de b).

}{F7I-IX-027}d)   L'hypothèse que $u_{t}$ soit une isogénie implique que $K_{t}$ est fini, donc on est sous les conditions de c), et par suite $(K_{\overline{s}})^{o}_{\mathrm{red}}$ est unipotent. Comme $G_{s}$ est de rang unipotent nul, on en déduit que $K_{s}$ est fini. D'ailleurs le fait que $u_{t}$ est une isogénie implique $\dim G_{t} = \dim H_{t}$, et le fait que $G$ et $H$ sont plats localement de présentation finie implique $\dim G_{s} = \dim G_{t}$, $\dim H_{s} = \dim H_{t}$ (SGA 3 $VI_{B}$ 4.3), d'où $\dim G_{s} = \dim H_{s}$ . Comme $\dim \mathrm{Ker}\ u_{s} = 0$, on en conclut que $\dim H_{s} = \dim u_{s}(G_{s})$ (SGA 3 $VI_{B}$ 1.2.), ce qui équivaut à $u_{s}$ plat grâce à l'hypothèse $H_{s}$ lisse (SGA 3 $VI_{A}$ 5.6). On en conclut qu'il existe un voisinage ouvert $U$ de $s$ tel que $u$ soit plat et quasi-fini en les points de la section unité sur $U$ (SGA 3 $VI_{B}$ 2.2 (i) et 2.5 (ii)). Il en résulte, comme pour $u_{s}$, que $u_{s'}$ est plat pour tout $s' \in U$, donc $u|G_{U}$ est plat et quasi-fini, l'étant sur chaque fibre.

\SGAvar{e)   Pour la premier assertion, on se ramène aussitôt par changement de base au cas où $S$ est intègre de point générique $t$, donc à prouver la deuxième assertion de e). Celle-ci résulte aussitôt de d) et du


}{e)   Pour la première assertion, on se ramène aussitôt par changement de base au cas où $S$ est intègre de point générique $t$, donc à prouver la deuxième assertion de e). Celle-ci résulte aussitôt de d) et du


}{F7I-IX-028}\SGApage{I}{346}

\SGAvar{\underline{Lemme} 3.1.3. \underline{Soient} $S$ \underline{un schéma}, $K$ \underline{un} \underline{S-schéma quasi-fini}, \underline{de présentation finie}, \underline{séparé et plat sur} $S$, $s \in S$, $t \in T$ \underline{une générisation de} $s$. \underline{Alors on a}

}{\underline{Lemme} 3.1.3. \underline{Soient} $S$ \underline{un schéma}, $K$ \underline{un} \underline{S-schéma quasi-fini}, \underline{de présentation finie}, \underline{séparé et plat sur} $S$, $s \in S$, $t \in S$ \underline{une générisation de} $s$. \underline{Alors on a}

}{F7I-IX-029}\begin{equation*}
\mathrm{rang}\ K_{s} \;\leq\; \mathrm{rang}\ K_{t} \quad .
\end{equation*}
\underline{Si rang} $K_{\eta} \leq 1$ \underline{pour tout point maximal de} $S$, \underline{alors} $K \longrightarrow S$ \underline{est une immersion ouverte} (\underline{donc un isomorphisme s'il admet une section}).

Pour la première assertion, on se ramène par localisation stricte au cas où $S$ est strictement local. Décomposant alors $K$ en somme d'un schéma fini $K^{f}$ et d'un schéma à fibre spéciale vide, on aura

\begin{equation*}
\mathrm{rang}\ K_{s} = \mathrm{rang}\ K^{f}_{s} = \mathrm{rang}\ K^{f}_{t} \leq \mathrm{rang}\ K_{t} \quad ,
\end{equation*}
la deuxième égalité provenant du fait que $K^{f}$ est plat, fini et de présentation finie sur $S$, donc localement libre sur $S$.

Pour la deuxième assertion, ce qui précède montre déjà que pour tout $s \in S$ on a rang $K_{s} \leq 1$, donc $K_{s}$ est vide ou isomorphe à $\mathrm{Spec}\ k(s)$. Par suite $K \longrightarrow S$ est un monomorphisme (SGA 3 VI$_{B}$ 2.11), et étant plat et de présentation finie c'est une immersion ouverte (EGA IV 17.9.1), ce qui achève la démonstration.

\underline{Remarque} 3.1.4. Dans 3.1 a) et la suite de la proposition 3.1, on ne donne aucun renseignement sur ${}_{n}K_{s}$ lorsque $n$ est une puissance de $p$ (lorsque $p > 1$). Lorsque $G$ est lisse (et $S$ localement noethérien ou $H$ de présentation finie sur $S$) tout au moins, on peut donner un tel renseignement, par essentiellement la même méthode que dans la


\SGApage{I}{347}

démonstration de a) (en se ramenant au cas $S$ local noethérien complet, et utilisant SGA 3 XV 1.6 c), X 3.2, IX 3.6, IX 5.2), savoir :

\SGAvar{\begin{equation}\tag{3.1.4.1}
\mathrm{rang\ mult}\ ({}_n K_s) \mid \mathrm{rang\ mult}({}_n K_t) \quad ,
\end{equation}
où rang mult désigne le rang du plus grand sous-groupe de type multiplicatif du groupe algébrique commutatif envisagé (SGA 3 XVII 7.2.1) ; on notera qu'un groupe algébrique de type multiplicatif de torsion est fini. La formule (3.1.4.1) est valable pour tout entier $n \geq 1$. On conclut, grâce à SGA XVII 4.6.1 (vi), que dans c), si $K_t$ est unipotent (SGA 3 VII 1.1, il en est de même de $K_s$.

}{\begin{equation}\tag{3.1.4.1}
\mathrm{rang\ mult}\ ({}_n K_s) \mid \mathrm{rang\ mult}({}_n K_t) \quad ,
\end{equation}
où rang mult désigne le rang du plus grand sous-groupe de type multiplicatif du groupe algébrique commutatif envisagé (SGA 3 XVII 7.2.1) ; on notera qu'un groupe algébrique de type multiplicatif de torsion est fini. La formule (3.1.4.1) est valable pour tout entier $n \geq 1$. On conclut, grâce à SGA 3 XVII 4.6.1 (vi), que dans c), si $K_t$ est unipotent (SGA 3 XVII 1.1), il en est de même de $K_s$.

}{F7I-IX-030}\underline{Proposition} 3.2. \underline{Soient} $S$ \underline{un schéma localement noethérien régulier intègre de dimension} 1, $K$ \underline{son corps des fonctions rationnelles,} $A_K$ \underline{un schéma abélien sur} $K$, $A$ \underline{son modèle de Néron sur} $S$, \underline{qui est un schéma en groupes lisse et séparé sur} $S$ (\underline{et quasi-projectif sur} $S$ \underline{si} $S$ \underline{est noethérien} [ 23 , VIII 2, 3°]). \underline{Les conditions suivantes sont équivalentes} :

(i) \underline{Pour tout} $s \in S$, $A_s^o$ \underline{est de rang unipotent nul,} \underline{i.e.} (2.1.5) \underline{est une extension d'un schéma abélien par un tore.}

(ii) \underline{Il existe un schéma en groupes} $G$ \underline{lisse,} \underline{séparé de type fini sur} $S$, \underline{prolongeant} $A_K$, \underline{et dont les fibres sont de rang unipotent nul.}

\underline{De plus,} \underline{si} $G$ \underline{est comme dans} (ii), \underline{alors l'unique morphisme} $G \longrightarrow A$ \underline{induisant l'isomorphisme donné sur les fibres génériques est un isomorphisme de} $G$ \underline{sur un sous-groupe ouvert de} $A$, \underline{donc induit un isomorphisme des composantes neutres} $G^o \xrightarrow{\sim} A^o$ .


\SGApage{I}{348}

3.2.1. Notons tout de suite que la définition du modèle de Néron $A$ de $A_K$ est la même que dans le cas local, rappelée dans 1.1. L'existence de ce modèle de Néron, et le fait qu'il est séparé de type fini sur $S$, est bien connue, et se ramène en fait immédiatement au cas local. Bien entendu, pour tout point \underline{fermé} $s \in S$, $S' = \mathrm{Spec}(\underline{O}_{S,s})$ est un trait, et $A \times_S S'$ n'est autre que le modèle de Néron de $A_K$ sur $S'$, dans le sens envisagé dans 1.1.

Démontrons 3.2. Il est trivial que i) implique ii), car il suffit de prendre $G = A$ . Il reste donc à prouver que si $G$ est comme dans (ii), alors $G \longrightarrow A$ est une immersion ouverte. Or c'est en effet un cas particulier de 3.1 e).

\underline{Corollaire} 3.3. \underline{Avec les notations de} 3.2, \underline{supposons les composantes neutres des fibres du modèle de Néron de rang unipotent nul}. \underline{Soit} $S'$ \underline{un schéma satisfaisant aux mêmes hypothèses que} $S$, $K'$ \underline{son corps des fonctions rationnelles}, $f: S' \longrightarrow S$ \underline{un morphisme dominant}, $A_{K'} \simeq A_K \otimes_K K' \simeq$ \underline{fibre générique de} $A \times_S S'$, $A'$ \underline{le modèle de Néron de} $A_{K'}$, \underline{et considérons le morphisme unique}

\begin{equation*}
u \; : \; A \times_S S' \longrightarrow A'
\end{equation*}
\underline{induisant l'identité sur les fibres génériques}. \underline{Alors} $u$ \underline{est une immersion ouverte,} \underline{et induit donc un isomorphisme sur les composantes neutres}

\begin{equation}\tag{3.3.1}
A^o \times_S S' \xrightarrow{\ \sim\ } {A'}^{o} \quad .
\end{equation}
Il suffit en effet d'appliquer 3.2 dans le cas de $S'$, en prenant $G \simeq A \times_S S'$ .


\SGApage{I}{349}

3.3.2. On voit donc que les propriétés équivalentes de 3.2 impliquent une propriété de "stabilité" partielle pour les modèles de Néron ; on fera attention que lorsque $A_K$ n'a pas bonne réduction i.e. que $A$ n'est pas un schéma abélien, il n'est \underline{pas} vrai que pour tout $S'$ comme dans 3.3 (et même pour $S'$ fini sur $S$, et $K'/K$ séparable ...), le morphisme $A \times_S S' \longrightarrow A'$ lui-même soit un isomorphisme, i.e. induise des isomorphismes sur les groupes de composantes connexes des fibres (cf. 11.10 plus bas pour le comportement de ces groupes de composantes connexes par extension de la base). Ces observations, renforcées par 3.9 plus bas, justifient la terminologie suivante, introduite par MUMFORD (et qui s'harmonise aussi, paraît-il, avec les notions de même nom étudiées dans son livre [17]):

\underline{Définition} 3.4. \underline{Les notations étant celles de} 3.2, \underline{on dit que le schéma abélien} $A_K$ \underline{sur} $K$ \underline{a} réduction semi-stable sur $S$ \underline{si les conditions équivalentes} (i) \underline{et} (ii) \underline{de} 3.2 \underline{sont vérifiées.}

\SGAvar{3.4.0. On notera que sous la forme (ii), cette condition ne fait pas appel à la notion de modèle de Néron. D'autre part, elle se ramène immédiatement par localisation au cas où $S$ est un trait, comme il résulte par exemple de 3.2.1. D'ailleurs, comme la formation du modèle de Néron commute au changement de base de $S$ vers son hensélisé, ou son hensélisé strict, la notion se ramène même au cas $S$ où $S$ est un trait strictement local.


}{3.4.0. On notera que sous la forme (ii), cette condition ne fait pas appel à la notion de modèle de Néron. D'autre part, elle se ramène immédiatement par localisation au cas où $S$ est un trait, comme il résulte par exemple de 3.2.1. D'ailleurs, comme la formation du modèle de Néron commute au changement de base de $S$ vers son hensélisé, ou son hensélisé strict, la notion se ramène même au cas où $S$ est un trait strictement local.


}{F7I-IX-031}\SGApage{I}{350}

\underline{Proposition} 3.5. (\underline{critère galoisien de réduction semi-stable}). \underline{On suppose} $S$ \underline{un trait.} \underline{Les conditions suivantes sur le schéma abélien} $A_K$ \underline{sont équivalentes} (\underline{les notations étant celles de} (2.5.3), \underline{appliquées à la situation déduite de} $(A_K,S)$ \underline{par passage au hensélisé de} $S$) :

(i) $A_K$ \underline{a réduction semi-stable sur} $S$ (3.4).

(ii) \underline{On a,} \underline{avec les notations} (2.5.3), (2.5.4) :

\SGAvar{\begin{equation*}
V \subset V^{\perp} \qquad .
\end{equation*}
(iii) \underline{On a pour la partie torique} $W$ \underline{du module de Tate} :

}{\begin{equation*}
V^{\perp} \subset V \qquad .
\end{equation*}
(iii) \underline{On a pour la partie torique} $W$ \underline{du module de Tate} :

}{F7I-IX-032}\begin{equation*}
W = V^{\perp} \qquad .
\end{equation*}
(iv) \underline{L'action du groupe d'inertie} $I$ \underline{sur} $U = T_{\ell}(A(\overline{K}))$ \underline{est unipotente d'échelon} 2 (\underline{i.e.} \underline{il existe un sous-module} $U'$ \underline{de} $U$ \underline{stable par} $I$ \underline{tel que} $I$ \underline{opère trivialement sur} $U'$ \underline{et sur} $U/U'$).

On peut supposer $S$ hensélien (3.4.0). Avec les notations de 2.1.9, (i) signifie que l'on a

\SGAvar{\begin{equation*}
\lambda = 0 \qquad ,
\end{equation*}
ce qui en vertu de (2.5.5) équivaut à $V = V \cap V^{\perp}$ i.e. à (ii). Comme $W=V\cap V^{\perp}$ par le théorème d'orthogonalité 2.4, on a (ii) $\Longleftrightarrow$ (iii). D'ailleurs (iv) signifie aussi que $I$ opère trivialement sur $U^{I} = V$, ce qui par dualité signifie qu'il opère trivialement sur $V^{\perp}$ i.e. que $V^{\perp} \subset U^{I}=V$, ce qui est (ii). Cela achève la démonstration.

}{\begin{equation*}
\lambda = 0 \qquad ,
\end{equation*}
ce qui en vertu de (2.5.5) équivaut à $V^{\perp} = V \cap V^{\perp}$ i.e. à (ii). Comme $W=V\cap V^{\perp}$ par le théorème d'orthogonalité 2.4, on a (ii) $\Longleftrightarrow$ (iii). D'ailleurs (iv) signifie aussi que $I$ opère trivialement sur $U/U^{I}=U/V$, ce qui par dualité signifie qu'il opère trivialement sur $V^{\perp}$ i.e. que $V^{\perp} \subset U^{I}=V$, ce qui est (ii). Cela achève la démonstration.

}{F7I-IX-033}\underline{Remarques} 3.5.1. Les conditions (ii) et (iv) ne dépendent visiblement que de l'action de $I$ sur le $\mathbb{Q}_{\ell}$-vectoriel $U \otimes_{\mathbb{Z}_{\ell}} \mathbb{Q}_{\ell}$ , et auraient pu être formulées en termes de cette action. Signalons aussi que ces conditions sont manifestement stables par tout changement de base par un morphisme


\SGApage{I}{351}

de traits $S' \longrightarrow S$. Enfin, si (iv) est satisfaite pour $A_K$, elle l'est aussi manifestement pour tout sous-schéma abélien et tout schéma abélien quotient de $A_K$ , et en particulier pour tout schéma abélien \underline{isogène} à $A$ (par exemple pour $A'_K$).

\SGAvar{\underline{Corollaire} 3.5.2. \underline{Soit} $g$ \underline{un générateur topologique de} $\mathbb{Z}_\ell(1) = T_\ell(\overline{k}^*)$. \underline{Les conditions équivalentes de} 3.5 \underline{équivalent aussi à la condition suivante} :

(v) \underline{L'action du groupe d'inertie} $I$ \underline{sur} $U = T_\ell(A(\overline{K}))$ \underline{se fait à travers le plus grand quotient pro-$\ell$-fini} $I(\ell)$ \underline{de} $I$, \underline{qui est canoniquement isomorphe à} $\mathbb{Z}_\ell(1)$ $(I \qquad )$, \underline{et l'opération} $g_U$ \underline{du générateur} $g$ \underline{sur} $U$ \underline{satisfait la condition}

}{\underline{Corollaire} 3.5.2. \underline{Soit} $g$ \underline{un générateur topologique de} $\mathbb{Z}_\ell(1) = T_\ell(\overline{k}^*)$. \underline{Les conditions équivalentes de} 3.5 \underline{équivalent aussi à la condition suivante} :


(v) \underline{L'action du groupe d'inertie} $I$ \underline{sur} $U = T_\ell(A(\overline{K}))$ \underline{se fait à travers le plus grand quotient pro-$\ell$-fini} $I(\ell)$ \underline{de} $I$, \underline{qui est canoniquement isomorphe à} $\mathbb{Z}_\ell(1)$, \underline{et l'opération} $g_U$ \underline{du générateur} $g$ \underline{sur} $U$ \underline{satisfait la condition}

}{F7I-IX-034}\begin{equation}\tag{3.5.3}
(1 - g_U)^2 \ = \ 0 \quad .
\end{equation}
En effet, pour toute action unipotente continue d'un groupe profini $I$ sur un espace vectoriel $\ell$-adique, on voit aussitôt que l'image de $I$ est un pro-$\ell$-groupe. Il s'ensuit que (v) équivaut à (iv).

Nous pouvons maintenant prouver le résultat central du présent exposé :

\underline{Théorème} 3.6. (\underline{théorème de réduction semi-stable}). \underline{Soient} $S$ \underline{un schéma noethérien régulier connexe de dimension} 1, $K$ \underline{le corps des fonctions rationnelles de} $S$, $A_K$ \underline{un schéma abélien sur} $K$. \underline{Alors il existe une extension finie galoisienne} $K'$ \underline{de} $K$, \underline{telle que le schéma abélien} $A_{K'} = A_K \otimes_K K'$ \underline{ait réduction semi-stable} (3.4) \underline{sur le normalisé} $S'$ \underline{de} $S$ \underline{dans} $K'$.


\SGApage{I}{352}

\SGAvar{On observera que $S'$ est noethérien régulier connexe de dimension 1, ce qui donne un sens à la conclusion énoncée. Pour prouver 3.6, choisissons une clôture séparable $\overline{K}$ de $K$, et écrivons l'extension $\overline{K}$ comme limite inductive de ses sous-extensions galoisiennes finies $K_i$ . Soit $S_i$ le normalisé de $S$ dans $K_i$ , et soit $U$ le plus grand ouvert de $S$ sur lequel $A_K$ ait bonne réduction. Comme $S$ est noethérien connexe de dimension 1, et $U \neq \emptyset$, on voit que $T = S-U$ est un ensemble fini de point fermés de $S$, et que pour tout $t \in T$, $\underline{O}_{S,t}$ est un anneau de valuation discrète. Il suffit de trouver, pour tout $t \in T$, un indice $i = i(t)$ tel que $A_{K_i}$ ait réduction semi-stable en les points de $S_i$ au-dessus de $t$ ; en effet, il suffira de prendre alors pour $K'$ le composé des extensions $K_{i(t)}$ pour $t \in T$. Cela nous ramène donc au cas où $S$ est un trait. D'autre part, comme de groupe $\mathrm{Gal}(K'/K)$ permute entre eux les points de $S_i$ au-dessus du point fermé $s$ de $S$, il suffit de trouver $K'$ de telle façon que $A_{K'}$ ait réduction semi-stable en \underline{un} point du normalisé $S'$ ; d'ailleurs il est inutile ici que $K'$ soit galoisien sur $K$, il suffit de trouver $K'$ extension finie étale de $K$ satisfaisant la condition précédente, car la clôture galoisienne de cette extension satisfera à la même condition. Utilisant le critère 3.5 (iv) on voit donc que 3.6 revient à l'assertion suivante :

\underline{Corollaire} 3.7. \underline{On suppose à nouveau} $S$ \underline{un trait}. \underline{Alors il existe un sous-groupe ouvert} $\pi'$ \underline{de} $\pi$ \underline{tel que l'action} $I' = I \cap \pi'$ \underline{sur} $U = T_\ell(A(\overline{K}))$ \underline{soit unipotente d'échelon} 2, \underline{ou encore} : \underline{il existe un sous-groupe ouvert} $I'$ \underline{de} $I$ \underline{dont l'action sur} $U$ \underline{soit unipotente d'échelon} 2.


}{On observera que $S'$ est noethérien régulier connexe de dimension 1, ce qui donne un sens à la conclusion énoncée. Pour prouver 3.6, choisissons une clôture séparable $\overline{K}$ de $K$, et écrivons l'extension $\overline{K}$ comme limite inductive de ses sous-extensions galoisiennes finies $K_i$ . Soit $S_i$ le normalisé de $S$ dans $K_i$ , et soit $U$ le plus grand ouvert de $S$ sur lequel $A_K$ ait bonne réduction. Comme $S$ est noethérien connexe de dimension 1, et $U \neq \emptyset$, on voit que $T = S-U$ est un ensemble fini de points fermés de $S$, et que pour tout $t \in T$, $\underline{O}_{S,t}$ est un anneau de valuation discrète. Il suffit de trouver, pour tout $t \in T$, un indice $i = i(t)$ tel que $A_{K_i}$ ait réduction semi-stable en les points de $S_i$ au-dessus de $t$ ; en effet, il suffira de prendre alors pour $K'$ le composé des extensions $K_{i(t)}$ pour $t \in T$. Cela nous ramène donc au cas où $S$ est un trait. D'autre part, comme le groupe $\mathrm{Gal}(K'/K)$ permute entre eux les points de $S_i$ au-dessus du point fermé $s$ de $S$, il suffit de trouver $K'$ de telle façon que $A_{K'}$ ait réduction semi-stable en \underline{un} point du normalisé $S'$ ; d'ailleurs il est inutile ici que $K'$ soit galoisien sur $K$, il suffit de trouver $K'$ extension finie étale de $K$ satisfaisant la condition précédente, car la clôture galoisienne de cette extension satisfera à la même condition. Utilisant le critère 3.5 (iv) on voit donc que 3.6 revient à l'assertion suivante :

\underline{Corollaire} 3.7. \underline{On suppose à nouveau} $S$ \underline{un trait}. \underline{Alors il existe un sous-groupe ouvert} $\pi'$ \underline{de} $\pi$ \underline{tel que l'action de} $I' = I \cap \pi'$ \underline{sur} $U = T_\ell(A(\overline{K}))$ \underline{soit unipotente d'échelon} 2, \underline{ou encore} : \underline{il existe un sous-groupe ouvert} $I'$ \underline{de} $I$ \underline{dont l'action sur} $U$ \underline{soit unipotente d'échelon} 2.


}{F7I-IX-035}\SGApage{I}{353}

On a un isomorphisme canonique (compatible aux opérations de $\pi$) :

\begin{equation}\tag{*}
H^1(A_{\overline{K}}, \mathbb{Z}_\ell) \simeq \mathrm{Hom}(T_\ell(A(\overline{K})), \mathbb{Z}_\ell) \quad ,
\end{equation}
qui montre que l'assertion 3.7 équivaut à la même assertion pour les opérations de $I$ sur le premier membre de (*). Or cette assertion a été prouvée dans III \hspace{2em} (en utilisant la résolution des singularités des schémas excellents de dimension 2).

\underline{Corollaire} 3.8. \underline{Les conditions équivalentes} (i) \underline{à} (iv) \underline{de} 3.5 \underline{équivalent aussi à la condition suivante} :

(vi) \quad \underline{L'action de} $I$ \underline{sur} $U = T_\ell(A(\overline{K}))$ \underline{est unipotente}.

\SGAvar{En effet, pour prouver que (v) implique (vi), on note que si un groupe $I$ opérant de façon unipotente sur un vectoriel $E$ de dimension finie sur un corps de caractéristique nulle (en l'occurrence $\mathbb{Q}_\ell$) admet un sous-groupe d'indice fini $I'$ opérant avec l'échelon $m$ (en l'occurrence $m=2$), alors $I$ lui-même opère avec l'échelon $m$. En effet, $I'$ étant d'indice fini dans $I$, le groupe algébrique $G' \subset \underline{\mathrm{Gl}}(E)$ qu'il définit est d'indice fini dans celui, $G$, défini par $I$, or $G$ étant unipotent est connexe (caractéristique nulle !), donc $G=G'$, donc $I$ et $I'$ ont même ``échelon d'unipotence''.

}{En effet, pour prouver que (vi) implique (iv), on note que si un groupe $I$ opérant de façon unipotente sur un vectoriel $E$ de dimension finie sur un corps de caractéristique nulle (en l'occurrence $\mathbb{Q}_\ell$) admet un sous-groupe d'indice fini $I'$ opérant avec l'échelon $m$ (en l'occurrence $m=2$), alors $I$ lui-même opère avec l'échelon $m$. En effet, $I'$ étant d'indice fini dans $I$, le groupe algébrique $G' \subset \underline{\mathrm{Gl}}(E)$ qu'il définit est d'indice fini dans celui, $G$, défini par $I$, or $G$ étant unipotent est connexe (caractéristique nulle !), donc $G=G'$, donc $I$ et $I'$ ont même ``échelon d'unipotence''.

}{F7I-IX-036}On obtient également une réciproque à 3.3 :

\underline{Corollaire} 3.9. \underline{Les notations sont celles de} 3.2. \underline{Les conditions} (i) \underline{et} (ii) \underline{de} 3.2 \underline{équivalent aussi à la condition suivante} :

(iii) \quad \underline{Pour toute extension étale} $K'$ \underline{de} $K$, \underline{désignant par} $S'$ \underline{le normalisé de} $S$ \underline{dans} $K'$, \underline{par} $A'$ \underline{le modèle de Néron de} $A_{K'} = A\otimes_S K'=A_K\otimes_K K'$, \underline{le morphisme canonique}


\SGApage{I}{354}

\begin{equation*}
A\times_S S' \longrightarrow A'
\end{equation*}
\underline{induit un isomorphisme sur les composantes neutres}.

La nécessité de la condition a été vue en effet, sous une forme plus forte, dans 3.3. Pour voir la suffisance, on choisit $K'$ comme dans 3.6. Comme les fibres de ${A'}^{\mathrm{o}}$ sont de rang unipotent nul, il en est de même de celles de $(A\times_S S')^{\mathrm{o}} = A^{\mathrm{o}}\times_S S'$, donc aussi de celles de $A^{\mathrm{o}}$ dont les précédentes se déduisent par extension du corps de base.


\SGAsectiontoc{4. Application à une conjecture de SERRE-TATE et au conducteur}{4}
4. \underline{Application à une conjecture de SERRE-TATE et au conducteur}

4.1. \quad Nous supposons $S$ hensélien, et gardons les notations de (2.5.3), avec $\ell \neq p$ (*). Pour étudier l'action de $\pi = \mathrm{Gal}(\overline{K}/K)$ et de son sous-groupe d'inertie $I$ sur $U = T_\ell(A(\overline{K}))$, nous allons utiliser le fait qu'il existe un sous-groupe ouvert $I'$ de $I$, que nous prendrons même invariant dans $\pi$, tel que l'action de $I'$ sur $U$ soit unipotente d'échelon 2, i.e. tel que l'on ait

\begin{equation*}
U^{I'} \supset (U^{I'})^{\perp} \quad .
\end{equation*}
Il est clair que $U^{I'}$ ne dépend pas alors du choix particulier de $I'$, on pourrait l'appeler \underline{le sous-module essentiellement fixe} de $U$, et nous le noterons $U^{\mathrm{ef}}$, son orthogonal prenant alors le nom de \underline{sous-module essentiellement torique} de $U$, qui sera noté $U^{\mathrm{et}}$. On a donc une filtration en trois crans

\begin{equation}\tag{4.1.1}
U \supset U^{\mathrm{ef}} \supset U^{\mathrm{et}} \supset 0 \quad ,
\end{equation}
évidemment invariante sous l'action de $\pi$ . Notons par la même occasion

\SGAvar{(*) Le cas $\ell \neq p$ sera étudié dans les \S\S 5 et 6 ci-dessous.


}{(*) Le cas $\ell = p$ sera étudié dans les \S\S 5 et 6 ci-dessous.


}{F7I-IX-037}\SGApage{I}{355}

\SGAvar{qu'il existe un plus grand sous-groupe de $I$ qui opère de façon unipotente (ou ce qui revient au même, unipotente d'échelon deux) sur $U$ : c'est le sous-groupe $I'$ des éléments qui opèrent de façon unipotente, qui est aussi noyau de la représentation de $I$ dans le gradué associé à la filtration (4.1.1) de $U$, ou ce qui revient ici au même, le noyau de la représentation induite de $I$ sur $U^{ef}$. En vertu du critère galoisien de réduction semi-stable (3.5), ce sous-groupe de $I'$ correspond à \underline{la plus petite extension galoisienne} $\tilde{K}'$ de $\tilde{K}$ (l'extension non ramifiée maximale de $K$) \underline{telle que} $A_{\tilde{K}'}$ \underline{ait une réduction semi-stable}. Cette description montre en même temps que ce sous-groupe $I'$ de $I$ ne dépend pas du choix du nombre premier $\ell \neq p$.

}{qu'il existe un plus grand sous-groupe de $I$ qui opère de façon unipotente (ou ce qui revient au même, unipotente d'échelon deux) sur $U$ : c'est le sous-groupe $I'$ des éléments qui opèrent de façon unipotente, qui est aussi noyau de la représentation de $I$ dans le gradué associé à la filtration (4.1.1) de $U$, ou ce qui revient ici au même, le noyau de la représentation induite de $I$ sur $U^{ef}$. En vertu du critère galoisien de réduction semi-stable (3.5), ce sous-groupe de $I$ correspond à \underline{la plus petite extension galoisienne} $\tilde{K}'$ de $\tilde{K}$ (l'extension non ramifiée maximale de $K$) \underline{telle que} $A_{\tilde{K}'}$ \underline{ait une réduction semi-stable}. Cette description montre en même temps que ce sous-groupe $I'$ de $I$ ne dépend pas du choix du nombre premier $\ell \neq p$.

}{F7I-IX-038}Notons que la représentation de $\pi$ sur $U/U^{ef}$ est connue à isogénie près quand on connaît la représentation sur $U^{ef}$, donc sur $U^{et}$, grâce à l'isogénie induite par l'accouplement de polarisation (2.5.1) :

\begin{equation}\tag{4.1.2}
U/U^{ef} \longrightarrow (U^{et})^{\vee}(1) \qquad \mbox{une isogénie,}
\end{equation}
où le (1) désigne le twist de Tate par $\mathbb{Z}_\ell(1) = T_\ell(\bar{K}^*)$ ; l'homomorphisme précédent est en effet, évidemment, compatible avec les opérations de $\pi$ . (NB. On a un \underline{isomorphisme} canonique $U/U^{ef} \simeq ({U'}^{et})^{\vee} (1)$, où $U' = T_\ell(A'(\bar{K})))$. Il s'impose donc d'étudier de plus près l'action de $\pi$ sur $U^{ef}$.

Or par construction, celle-ci est triviale sur le sous-groupe invariant $I'$ de $I$, elle se factorise donc en une représentation du groupe quotient

\begin{equation}\tag{4.1.3}
\pi' = \pi/I' \qquad ,
\end{equation}

\SGApage{I}{356}

action dont nous allons donner maintenant une description ``indépendante de $\ell$ ''.

4.2. Soit donc $K'$ une extension galoisienne de $K$, correspondant à un sous-groupe invariant ouvert $V$ de $\pi$ tel que $V \cap I = I'$, le groupe de Galois $G$ de $K'/K$ apparaissant donc comme une extension

\begin{equation}\tag{4.2.1}
1 \longrightarrow \Gamma \longrightarrow G \longrightarrow G_o \longrightarrow 1 \qquad ,
\end{equation}
où

\begin{equation}\tag{4.2.2}
\Gamma = I/I'
\end{equation}
est le sous-groupe d'inertie, et $G_o$ le groupe de Galois de la plus grande sous-extension non ramifiée $K''$ de $K'/K$ . La suite exacte (4.2.1) est donc image de la suite exacte analogue

\begin{equation}\tag{4.2.3}
1 \longrightarrow \Gamma \longrightarrow \pi' \longrightarrow \pi_o \longrightarrow 1 \qquad ,
\end{equation}
où $\pi'$ est le groupe (4.1.3) opérant sur $U^{ef}$. On a donc un diagramme d'inclusions de corps

\begin{equation}\tag{4.2.4}
\begin{tikzcd}[column sep=large, row sep=large]
K \arrow[r, "G_o"] \arrow[rd, dash, "G"'] \arrow[rr, bend left=25, dash, "\pi_o"] \& K''=\tilde{K}\cap K' \arrow[r] \arrow[d, "\Gamma"] \& \tilde{K} \arrow[d, "\Gamma"] \arrow[dd, bend left=40, dash, "I"] \\
\& K' \arrow[r] \& \tilde{K}' \arrow[d, "I'"] \\
\& \& \bar{K}
\end{tikzcd}
\end{equation}

















\SGAvar{Soient $S'$ le normalisé de $S$ dans $K'$, $A^*$ le modèle de Néron de $A_{K'}$. Alors par le critère galoisien de réduction semi-stable (4.5), $A_{K'}$ \underline{a réduction semi-stable sur} $S'$ \underline{i.e.} $A^{*o}_{o}$ \underline{est extension d'un schéma}


}{Soient $S'$ le normalisé de $S$ dans $K'$, $A^*$ le modèle de Néron de $A_{K'}$. Alors par le critère galoisien de réduction semi-stable (3.5), $A_{K'}$ \underline{a réduction semi-stable sur} $S'$ \underline{i.e.} $A^{*o}_{o}$ \underline{est extension d'un schéma}


}{F7I-IX-039}\SGApage{I}{357}

\SGAvar{\underline{abélien par un tore} ; \underline{d'autre part,} \underline{par transport de structure}, $G$ \underline{opère} \underline{sur} $A^*$ \underline{donc sur} $A^{*o}$ \underline{par automorphisme de façon compatible avec l'action donnée de} $G$ \underline{sur le schéma de base} $S'$ \underline{de} $A^*$. Il est clair alors que l'\underline{action de} $\pi'$ \underline{sur} $U^{ef} = T_{\ell}(A^{*o})^{f}(\overline{K})$ \underline{est déduit par transport de structure de l'action de} $\pi'$ \underline{sur} $(A^{*o}/K',\widetilde{K}')$, $\pi'$ \underline{équivaut sur} $A^*$ \underline{via} $G$. Tenant compte de l'isomorphisme canonique

\begin{equation}\tag{4.2.6}
U^{ef} \simeq T_{\ell}(A^{*o})^{f}(\overline{K}) \simeq T_{\ell}(A^{*o}_{o})(\overline{k})
\end{equation}
(valable pour tout faisceau $\ell$-adique constant tordu sur $S'$, cf. I \hspace{2.5em}), on trouve alors la description suivante de l'action de $\pi'$ sur $U^{ef}$ :

}{\underline{abélien par un tore} ; \underline{d'autre part,} \underline{par transport de structure}, $G$ \underline{opère} \underline{sur} $A^*$ \underline{donc sur} $A^{*o}$ \underline{par automorphisme de façon compatible avec l'action donnée de} $G$ \underline{sur le schéma de base} $S'$ \underline{de} $A^*$. Il est clair alors que l'\underline{action de} $\pi'$ \underline{sur} $U^{ef} = T_{\ell}(A^{*o})^{f}(\overline{K})$ \underline{est déduite par transport de structure de l'action de} $\pi'$ \underline{sur} $(A^{*o}/K',\widetilde{K}')$, $\pi'$ \underline{équivaut sur} $A^*$ \underline{via} $G$. Tenant compte de l'isomorphisme canonique

\begin{equation}\tag{4.2.6}
U^{ef} \simeq T_{\ell}(A^{*o})^{f}(\overline{K}) \simeq T_{\ell}(A^{*o}_{o})(\overline{k})
\end{equation}

(valable pour tout faisceau $\ell$-adique constant tordu sur $S'$, cf. Exposé I [référence interne omise dans la source]), on trouve alors la description suivante de l'action de $\pi'$ sur $U^{ef}$ :

}{F7I-IX-040}On considère le schéma en groupes lisse connexe $A^{*o}_{o}$ sur le corps résiduel $k' \hookrightarrow \overline{k}$ de $S'$, qui est une extension quasi-galoisienne de $k$ de groupe de Galois $G_{o}$. Le groupe quotient $G$ de $\pi'$ opère sur $A^{*o}_{o}$ de façon compatible avec son opération sur $k'$ via $G_{o}$. D'ailleurs $\pi'$ opère sur $\overline{k}$ via son quotient $\pi_{o}$, donc il opère sur la situation $(A^{*o}_{o}/k',\overline{k})$. Donc \underline{par transport de structure} $\pi'$ opère sur

\begin{equation*}
U^{ef} = T_{\ell}(A^{*o}_{o})(\overline{k}) \quad ,
\end{equation*}
et c'est l'opération cherchée de $\pi'$. En particulier, le sous-groupe $\Gamma$ de $G$ qui opère trivialement sur $k'$ opère sur $A^{*o}_{o}$ par $k'$-automorphismes, et la représentation correspondante de $\Gamma = I/I'$ sur $T_{\ell}(A^{*o}_{o})(\overline{k})$ est la représentation déduite de la représentation naturelle du groupe d'inertie $I$ sur $U^{ef}$.

D'autre part, la partie essentiellement torique $U^{et}$ de $U$ correspond simplement au sous-module

\begin{equation}\tag{4.2.7}
U^{et} \simeq T_{\ell}(T^{*}_{o})(\overline{k}) \subset U^{ef} \simeq T_{\ell}(A^{*o}_{o})(\overline{k}) \quad ,
\end{equation}

\SGApage{I}{358}

où $T^{*}_{o}$ est le plus grand sous-tore de $A^{*o}_{o}$. Posant de même

\begin{equation}\tag{4.2.8}
B^{*}_{o} = A^{*o}_{o}/T^{*}_{o}
\end{equation}
pour la partie abélienne de $A^{*o}_{o}$, l'opération de $G$ sur $A^{*o}_{o}$ induit séparément des opérations de $G$ sur $T^{*}_{o}$ et sur $B^{*}_{o}$, et l'opération de $\pi'$ sur $U^{et}$ et sur $U^{ef}/U^{et}$ se déduit des opérations précédentes à l'aide de l'isomorphisme (4.2.7) et de l'isomorphisme correspondant

\begin{equation}\tag{4.2.9}
U^{ef}/U^{et} \simeq T_{\ell}(B^{*}_{o})(\overline{k}) \quad .
\end{equation}
De ceci, et de l'isogénie (4.1.2), on déduit, pour le caractère $g \longmapsto \mathrm{Tr}\, g_{U}$ de l'action de $\pi'$ (ou de $\pi$) sur $U$, la formule

\begin{equation}\tag{4.2.10}
\mathrm{Tr}\, g_{U_{\ell}} = \mathrm{Tr}\, g_{B^{*}_{o},\ell} + \mathrm{Tr}\, g_{T^{*}_{o},\ell} + \chi_{\ell}(g)\mathrm{Tr}\, g^{-1}_{T^{*}_{o},\ell} \quad ,
\end{equation}
où au premier membre on a écrit $U_{\ell}$ pour $U$, pour rappeler le nombre premier choisi $\ell \neq p$, où $g_{B^{*}_{o},\ell}$ resp. $g_{T^{*}_{o},\ell}$ désigne l'action de $g$ sur le module de Tate $T_{\ell}(B^{*}_{o})(\overline{k})$ resp. $T_{\ell}(T^{*}_{o})(\overline{k})$, déduit de son action sur $B^{*}_{o}$ resp. $T^{*}_{o}$ qu'on vient d'expliciter, et où

\SGAvar{\begin{equation}\tag{4.2.11}
\chi_{\ell} : \pi' \longrightarrow \pi_{o} \longrightarrow \mathbb{Z}^{*}_{\ell}
\end{equation}
est le caractère de $\pi'$ (ou de $\pi_{o}$) déduit de l'action sur $\mathbb{Z}_{\ell}(1) = T_{\ell}(\overline{k}^{*})$, - caractère qui est trivial sur $\Gamma^{o}$ .

}{\begin{equation}\tag{4.2.11}
\chi_{\ell} : \pi' \longrightarrow \pi_{o} \longrightarrow \mathbb{Z}^{*}_{\ell}
\end{equation}
est le caractère de $\pi'$ (ou de $\pi_{o}$) déduit de l'action sur $\mathbb{Z}_{\ell}(1) = T_{\ell}(\overline{k}^{*})$, - caractère qui est trivial sur $\Gamma$ .

}{F7I-IX-041}Nous sommes maintenant en mesure de prouver un cas particulier d'une conjecture de SERRE-TATE [29, Appendice] (*).

(*) C'est essentiellement le cas $i=1$ desdites conjectures pour les $H^{i}$ .


\SGApage{I}{359}

\underline{Théorème} 4.3. \underline{Soient} $S$ \underline{un trait hensélien de corps résiduel} $k$, \underline{corps des fonctions} $K$, $A_K$ \underline{un schéma abélien sur} $K$, \underline{considérons la représentation de} $\pi = \mathrm{Gal}(\overline{K}/K)$ \underline{sur le module de Tate} $U_\ell = T_\ell(A(\overline{K}))$, \underline{où} $\ell$ \underline{est un nombre premier distinct de la caractéristique résiduelle.}

a) \underline{Il existe un sous-groupe ouvert} $I'$ \underline{du groupe d'inertie qui opère de façon unipotente sur} $U_\ell$, \underline{de sorte que le caractère de la représentation de} $\pi$ \underline{sur} $U_\ell$ \underline{est en fait une fonction sur} $\pi/I'$ . \underline{De plus, sur} $\Gamma = I/I'$ \underline{cette fonction est à valeurs entières, et est indépendante de} $\ell$ . \underline{Plus généralement, pour tout} $g \in I$, \underline{les coefficients du polynôme caractéristique de} $g_{U_\ell}$ \underline{sont des entiers rationnels indépendants de} $\ell$ .

\SGAvar{b) \underline{Supposons que} $k$ \underline{soit un corps fini à} $q$ \underline{éléments, de sorte que} $\pi_o = \pi/I \simeq \mathrm{Gal}(\overline{k}/k)$ \underline{s'identifie au complété} $\widehat{\mathbb{Z}}$ \underline{de} $\mathbb{Z}$, \underline{au générateur} $1$ \underline{de} $\mathbb{Z}$ \underline{correspondant l'automorphisme de frobenius} $f_q : x \longmapsto x^q$ \underline{de} $\overline{k}/k$. \underline{Alors pour tout} $g \in \pi$ \underline{dont l'image dans} $\pi_o$ \underline{est dans} $\mathbb{Z}$, \underline{la trace} (\underline{et plus généralement tout coefficients du polynôme caractéristique}) \underline{de} $g_{U_\ell}$ \underline{est un nombre rationnel indépendant de} $\ell$, \underline{qui est même un entier si l'image de} $g$ \underline{dans} $\pi_o$ \underline{est de la forme} $f_q^{\,n}$, \underline{avec} $n \geq 0$.

}{b) \underline{Supposons que} $k$ \underline{soit un corps fini à} $q$ \underline{éléments, de sorte que} $\pi_o = \pi/I \simeq \mathrm{Gal}(\overline{k}/k)$ \underline{s'identifie au complété} $\widehat{\mathbb{Z}}$ \underline{de} $\mathbb{Z}$, \underline{au générateur} $1$ \underline{de} $\mathbb{Z}$ \underline{correspondant l'automorphisme de frobenius} $f_q : x \longmapsto x^q$ \underline{de} $\overline{k}/k$. \underline{Alors pour tout} $g \in \pi$ \underline{dont l'image dans} $\pi_o$ \underline{est dans} $\mathbb{Z}$, \underline{la trace} (\underline{et plus généralement tout coefficient du polynôme caractéristique}) \underline{de} $g_{U_\ell}$ \underline{est un nombre rationnel indépendant de} $\ell$, \underline{qui est même un entier si l'image de} $g$ \underline{dans} $\pi_o$ \underline{est de la forme} $f_q^{\,n}$, \underline{avec} $n \geq 0$.

}{F7I-IX-042}Prouvons par exemple les assertions sur les traces. Les assertions sur les coefficients du polynôme caractéristique se prouveraient de la même façon, mais on peut noter aussi qu'elles résultent formellement des assertions faites sur les traces, grâce au calcul de ces coefficients en termes des traces des itérés (formules introduisant cependant des dénominateurs, donc ne donnant pas l'intégralité annoncée), et à un lemme de DWORK-FATOU (dont une jolie démonstration


\SGApage{I}{360}

est donnée par KLEIMAN dans [14 2.8]).

a) On utilise la formule (4.2.10), qui devient ici

\begin{equation}\tag{4.3.0}
\mathrm{Tr}\, g_{U_\ell} \;=\; \mathrm{Tr}\, g_{B^{*}_{o},\ell} \;+\; \mathrm{Tr}\, g_{T^{*}_{o},\ell} \;+\; \mathrm{Tr}\, g^{-1}_{T^{*}_{o},\ell} \qquad ,
\end{equation}
où maintenant $g_{B^{*}_{o}}$ et $g_{T^{*}_{o}}$ sont des $k'$-automorphismes de $B^{*}_{o}$ resp. de $T^{*}_{o}$ . Il suffit donc de montrer que les deux traces correspondantes sont des entiers rationnels indépendants de $\ell$ ce qui est un énoncé ``géométrique'' général sur les groupes algébriques. Pour $g_{B^{*}_{o}}$ c'est un résultat bien connu de WEIL sur la théorie des variétés abéliennes [5\ \ ][15\ ] ; pour $g_{T^{*}_{o}}$ c'est pratiquement trivial, en explicitant la représentation de $\Gamma$ sur $T = T^{*}_{o}$ à l'aide de la représentation correspondante sur le groupe dual $M'$, isomorphe à $\mathbb{Z}^{\mu}$ : la trace cherchée est contragrédiente de celle de cette dernière représentation, grâce à l'isomorphisme canonique

\begin{equation*}
T_\ell(T) \;\simeq\; \check{M}_\ell(1) \qquad \stackrel{\mathrm{dfn}}{=} \qquad \check{M}' \otimes_{\mathbb{Z}} \mathbb{Z}_\ell(1)
\end{equation*}
b) Si l'image de $g$ dans $\pi_o$ est $f_q^{\,n}$, la formule (4.2.10) devient maintenant

\begin{equation}\tag{4.3.1}
\mathrm{Tr}\, g_{U,\ell} \;=\; \mathrm{Tr}\, g_{B^{*}_{o},\ell} \;+\; \mathrm{Tr}\, g_{T^{*}_{o},\ell} \;+\; q^{n}\mathrm{Tr}\, g^{-1}_{T^{*}_{o},\ell} \qquad .
\end{equation}
Nous sommes ainsi ramenés à prouver le lemme suivant (où on fera $k \longmapsto \overline{k}'$, $A \longmapsto A^{*\,o}_{o} \otimes_{k'} \overline{k}$) :

\underline{Lemme} 4.3.2. \underline{Soient} $k$ \underline{la clôture algébrique d'un corps fini,} $A$ \underline{un schéma en groupes commutatif de type fini sur} $k$, $g$ \underline{un automorphisme de ce schéma en groupes, compatible avec l'automorphisme} $f$ \underline{de} $\mathrm{Spec}(k)$ \underline{déduit par} transport de structure \underline{d'un automorphisme de frobenius itéré} $x \longmapsto x^q$ (\underline{où} $q = p^{r}$ \underline{est une puissance de la caractéristique}, \underline{avec} $r \geq 0$),


\SGApage{I}{361}

-\underline{donc on a pour l'action de} $g$ \underline{sur des fonctions} $\varphi$ :

\begin{equation*}
g^*(\lambda\varphi) = f^*(\lambda)g^*(\varphi) = \lambda^{q^{-1}} g^*(\varphi) = \lambda^{p^{-r}} g^*(\varphi) \qquad (\lambda \in k) \quad .
\end{equation*}
\underline{Considérons l'automorphisme} $g_\ell$ \underline{induit par} $g$ \underline{sur} $T_\ell(A(k))$ \underline{par} "\underline{transport de structure}", \underline{où} $\ell$ \underline{est un nombre premier} $\neq p$. \underline{Alors la trace de} $g_\ell$ (\underline{plus généralement}, \underline{tout coefficient de son polynôme caractéristique}) \underline{est un entier rationnel}, \underline{indépendant de} $\ell$, \underline{et ses valeurs propres} \underline{sont des entiers algébriques de valeur absolue égale à} $q^{1/2}$ \underline{ou} $q$. \underline{Si} $A$ \underline{est un schéma abélien}, \underline{les valeurs absolues sont} $q^{1/2}$ ; \underline{si} $A$ \underline{est affine}, \underline{les valeurs propres sont de la forme} $q\varepsilon_i$, \underline{où les} $\varepsilon_i$ \underline{sont} \underline{des racines de l'unité}.

En effet, considérons l'endomorphisme $f_q^A$ "puissance $q$.ième" du schéma absolu sous-jacent à $A$, qui est défini comme l'identité sur les points et l'élévation à la puissance $q$.ième sur le faisceau structural. C'est un endomorphisme du schéma en groupes $A$ compatible avec l'endomorphisme $f_q^{\mathrm{Spec}(k)} = f^{-1}$ sur la base $\mathrm{Spec}(k)$, donc le composé

\begin{equation}\tag{4.3.3}
F = f_q^A\, g
\end{equation}
est un endomorphisme du schéma en groupes $A$, compatible avec l'identité sur $\mathrm{Spec}(k)$, i.e. c'est un $k$-endomorphisme du schéma en groupes $A$. Je dis que l'automorphisme $\Gamma_g$ de $A(k)$ défini par $g$ par transport de structure, explicité par la formule

\SGAvar{\begin{equation}\tag{4.3.4}
\Gamma_g(u) = g \circ u \circ f^{-1} \qquad \mbox{pour } u \in A(k) \quad ,
\end{equation}
peut aussi de décrire par

}{\begin{equation}\tag{4.3.4}
\Gamma_g(u) = g \circ u \circ f^{-1} \qquad \mbox{pour } u \in A(k) \quad ,
\end{equation}
peut aussi se décrire par

}{F7I-IX-043}\begin{equation}\tag{4.3.5}
\Gamma_g(u) = F(u) \qquad \mbox{pour } u \in A(k) \quad .
\end{equation}

\SGApage{I}{362}

\SGAvar{Pour le vérifier, il suffit de voir que l'application ensembliste sous-jacente aux morphismes $\mathrm{Spec}(k) \longrightarrow A$ représentés par les deux membres sont les mêmes, ce qui nous réduit à prouver l'égalité déduite de (4.3.5) en y composant le premier membre à gauche avec $f_q^A$, et le deuxième membre à droite avec $f^{-1}$ (puisque $f_q^A$ et $f^{-1}$ induisent l'identité sur les ensembles sous-jacents aux schémas source). Mais la formule en question est alors conséquence immédiate de (4.3.3) et de (4.3.4).

}{Pour le vérifier, il suffit de voir que les applications ensemblistes sous-jacentes aux morphismes $\mathrm{Spec}(k) \longrightarrow A$ représentés par les deux membres sont les mêmes, ce qui nous réduit à prouver l'égalité déduite de (4.3.5) en y composant le premier membre à gauche avec $f_q^A$, et le deuxième membre à droite avec $f^{-1}$ (puisque $f_q^A$ et $f^{-1}$ induisent l'identité sur les ensembles sous-jacents aux schémas source). Mais la formule en question est alors conséquence immédiate de (4.3.3) et de (4.3.4).

}{F7I-IX-044}Appliquant (4.3.5) aux $u \in {}_{\ell^\nu}A(k)$, et passant à la limite sur $\nu$, on trouve que $g_\ell$ peut aussi s'interpréter comme l'endomorphisme $T_\ell(F)$ de $T_\ell(A(k))$ associé à l'endomorphisme $F$ de $A$. On sait d'autre part que $T_\ell(F)$ a un polynôme caractéristique qui est à coefficients entiers indépendants de $\ell$, en se ramenant par dévissage au cas où $A$ est, soit un schéma abélien, soit un tore, où le résultat en question était déjà utilisé dans la partie a) de notre démonstration. De plus, cette démonstration prouve également les assertions faites sur les valeurs propres, compte tenu, dans le cas d'un schéma abélien, des classiques résultats de WEIL [15].

Cela prouve 4.3.2 et achève la démonstration de 4.3, et prouve en même temps :

\underline{Corollaire} 4.4. \underline{Plaçons nous sous les conditions de} 4.3 b). \underline{Alors} \underline{pour tout élément} $g$ \underline{dans} $\pi'$ \underline{dont l'image dans} $\pi_o$ \underline{est de la forme} $f_q^n$ \underline{avec} $n \geq 0$, \underline{les valeurs propres des automorphismes induits de} $U^{et}$, $U^{ef}/U^{et}$, $U/U^{ef}$ \underline{sont des entiers algébriques dont les valeurs absolues} \underline{sont respectivement} $q^n$, $q^{n/2}$ \underline{et} 1. \underline{En ce qui concerne les deux morceaux} \underline{extrêmes}, \underline{on peut dire}, \underline{plus précisément}, \underline{que les valeurs propres}


\SGApage{I}{363}

\underline{correspondantes sont respectivement} $(q^n\varepsilon_i)_{1\leq i\leq\mu}$ \underline{et} $(\varepsilon_i)_{1\leq i\leq\mu}$, \underline{où les} $\varepsilon_i(1\leq i\leq\mu)$ \underline{sont des} racines de l'unité.

\SGAvar{En fait, la formule (4.1.2) donne pour le dernier morceau les valeurs propres $q^n/(q^n\varepsilon_i) = \varepsilon_i^{-1}$, mais comme $\varepsilon_i$ est une racine de l'unité, on a $\varepsilon_i^{-1} = \overline{\varepsilon_i}$ (conjugué complexe), et comme la famille des $\varepsilon_i$ est stable par conjugaison complexe (comme famille des solutions d'une équation à coefficients entiers), la famille des $\varepsilon_i^{-1}$ est aussi celle des $\varepsilon_i$, ce qui prouve la derniere assertion de 4.4.

}{En fait, la formule (4.1.2) donne pour le dernier morceau les valeurs propres $q^n/(q^n\varepsilon_i) = \varepsilon_i^{-1}$, mais comme $\varepsilon_i$ est une racine de l'unité, on a $\varepsilon_i^{-1} = \overline{\varepsilon_i}$ (conjugué complexe), et comme la famille des $\varepsilon_i$ est stable par conjugaison complexe (comme famille des solutions d'une équation à coefficients entiers), la famille des $\varepsilon_i^{-1}$ est aussi celle des $\varepsilon_i$, ce qui prouve la dernière assertion de 4.4.

}{F7I-IX-045}4.5. Nous allons donner une application de 4.3 a) à la théorie du conducteur local des schémas abéliens. Supposons $k$ algébriquement clos (*). Rappelons (SGA 5 X 6.2 ou [20]) que pour tout schéma en $\ell$-groupes commutatifs étale $F_K$ sur $K$, ($\ell \neq p$), on définit un entier

\begin{equation}\tag{4.5.1}
\alpha_s(F_K) = \alpha_s(i_*(F_K))
\end{equation}
($i : \eta \longrightarrow S$ désignant l'inclusion), de la façon suivante : on choisit une extension galoisienne finie $K'$ de $K$, de groupe $\Gamma$, qui déploie $F_K$, de sorte que $F_K$ soit défini par un $\ell$-groupe fini $M$ sur lequel $\Gamma$ opère. D'autre part, on associe à l'extension $K'$ de $K$ un certain $\mathbb{Z}_\ell[\Gamma]$-\underline{module projectif} (défini à isomorphisme près)

\begin{equation}\tag{4.5.2}
\mathrm{Sw} = \mathrm{Sw}(K'/K,\mathbb{Z}_\ell) \qquad ,
\end{equation}
et on pose

\begin{equation}\tag{4.5.3}
\alpha_s(F_K) = \mathrm{long}\ \mathrm{Hom}_{\mathbb{Z}_\ell(\Gamma)}(\mathrm{Sw},M) \qquad .
\end{equation}
On vérifie (loc. cit) que cette expression ne dépend pas du choix de l'extension $K'/K$. D'ailleurs, le fait que $\mathrm{Sw}$ soit un $\mathbb{Z}_\ell[\Gamma]$-module

(*) Il suffirait en fait que $k$ soit \underline{parfait}, mais par passage au hensélisé strict de $S$, on se ramène aussitôt au cas envisagé.


\SGApage{I}{364}

projectif implique que $\alpha_s(F_K)$ se comporte additivement pour les suites exactes.

Nous allons appliquer cette définition au cas du faisceau ${}_\ell A_K$, où $A_K$ est notre schéma abélien sur $K$. On a donc, dans ce cas, un entier

\begin{equation}\tag{4.5.4}
\alpha_s(A_K;\ell) = \mathrm{long}\ \mathrm{Hom}_{\mathbb{Z}_\ell(\Gamma)}(\mathrm{Sw},U_\ell \otimes_{\mathbb{Z}_\ell} \mathbb{F}_\ell) \ , \ U_\ell = T_\ell(A(\overline{K})) \ ,
\end{equation}
$\Gamma$ étant choisi comme il a été dit (choix dépendant à priori de $\ell$). L'entier (4.5.4) s'appelle l'\underline{exposant conducteur immodéré} (ou "\underline{sauvage}") (*) du schéma abélien $A_K$ par rapport à l'anneau de valuation discrète $V \subset K$ ; par passage au hensélisé, il est défini bien sûr sans supposer $V$ hensélien, et il se ramène aussitôt au cas où $S$ est \underline{strictement} local. Nous nous proposons de montrer le

\underline{Corollaire} 4.6. \underline{L'exposant conducteur immodéré} (4.5.4) \underline{ne dépend pas du choix du nombre premier} $\ell \neq p$.

\SGAvar{Nous pouvons supposer $S$ strictement local. Nous allons donner le calcul de (4.5.4), en remplaçant $\mathrm{Sw}$ par n'importe quel autre $\mathbb{Z}_\ell(\Gamma)$-module projectif $L$. Notons que la filtration (4.1.1) de $U_\ell$ se fait par des sous-$\mathbb{Z}_\ell$-modules qui sont manisfestement des facteurs directs, donc

}{Nous pouvons supposer $S$ strictement local. Nous allons donner le calcul de (4.5.4), en remplaçant $\mathrm{Sw}$ par n'importe quel autre $\mathbb{Z}_\ell(\Gamma)$-module projectif $L$. Notons que la filtration (4.1.1) de $U_\ell$ se fait par des sous-$\mathbb{Z}_\ell$-modules qui sont manifestement des facteurs directs, donc

}{F7I-IX-046}(*) Il y a lieu de définir aussi l'\underline{exposant conducteur modéré} de $A_K$ par rapport à $V$ comme l'entier différence des longueurs sur $\mathbb{Z}/\ell\mathbb{Z}$ des fibres géométriques générique et spéciale de ${}_\ell A^o$, qui est indépendant de $\ell$ et égal à $2\lambda + \mu$ avec les notations de (1.1.11) ; \underline{l'exposant conducteur} (total) de $A_K$ est défini comme la somme des deux entiers précédents; c'est le terme local qui intervient dans la formule d'Euler-Poincaré pour ${}_\ell A^o$ de SGA 5 X 6.2, [20].


\SGApage{I}{365}

cette filtration induit une filtration analogue de $U(\ell) = U_\ell \otimes_{\mathbb{Z}_\ell} \mathbb{F}_\ell$ , dont les facteurs successifs sont déduits des facteurs $V_i(i=1,2,3)$

\begin{equation*}
U^{et} \; , \; U^{ef}/U^{et} \; , \; U/U^{ef}
\end{equation*}
intervenant dans la filtration (4.1.1) par tensorisation par $\mathbb{F}_\ell$ . Bien entendu, cette filtration est stable sous $\Gamma$, et donne par additivité une décomposition de long $\mathrm{Hom}_{\mathbb{Z}_\ell(\Gamma)}(L,U(\ell))$ en somme de trois termes, longueurs des groupes

\begin{equation}\tag{*}
\mathrm{Hom}_{\mathbb{Z}_\ell(\Gamma)}(L,V_i \otimes_{\mathbb{Z}_\ell} \mathbb{F}_\ell) \;\simeq\; (\mathrm{Hom}_{\mathbb{Z}_\ell(\Gamma)}(L,V_i)) \otimes_{\mathbb{Z}_\ell} \mathbb{F}_\ell \qquad ,
\end{equation}
l'isomorphisme écrit dans (*) résultant encore de l'hypothèse de projectivité faite sur $L$ ; on suppose seulement $K'$ choisi assez grand pour que $\Gamma$ opère sur les $V_i$ . Or $\mathrm{Hom}_{\mathbb{Z}_\ell(\Gamma)}(L,V_i)$ est évidemment un $\mathbb{Z}_\ell$-module sans torsion, donc son \underline{rang} est égal au deuxième membre de (*). Si $\varphi_i$ est le caractère de la représentation de $\Gamma$ dans $V_i$, et $\psi$ le caractère de la représentation de $I$ dans $\check{L}$ , on trouve donc

\begin{equation}\tag{4.6.1}
\mathrm{long}\ \mathrm{Hom}_{\mathbb{Z}_\ell(\Gamma)}(L,U(\ell)) = \sum_{1 \leq i \leq 3} N^{-1} \sum_{g \in \Gamma} \psi(g)\varphi_i(g)
\end{equation}
\begin{equation*}
= N^{-1} \sum_{g \in \Gamma} \psi(g)\varphi(g), \quad \mbox{avec } N = \mathrm{card}\ \Gamma,
\end{equation*}
où $\varphi(g) = \Sigma\, \varphi_i(g)$ est aussi la fonction sur $\Gamma$ déduite par passage au quotient de la fonction sur le groupe d'inertie $I$ caractère de la représentation de monodromie de $I$ sur $U$. Or en vertu de 4.3 a) ce caractère est une fonction à \underline{valeurs entières rationnelles indépendante} de $\ell$ . Prenant $L = \mathrm{Sw}$, il en est de même du caractère $\psi(g)$ de $L$, qui est égal à $a-N\delta+1$, où $\delta$ est la ``\underline{fonction de Dirac}'' de $\Gamma$, et $a$ le ``\underline{conducteur}


\SGApage{I}{366}

\underline{d'Artin}'' de $K'/K$ [26]. Cela prouve donc 4.6, en mettant en évidence en même temps l'expression suivante, indépendante de $\ell$, de l'exposant conducteur immodéré :

\begin{equation}\tag{4.6.2}
\alpha_s(A_K) = N^{-1} \sum_{g \in \Gamma} (a(g)-N\delta(g)+1)\, \varphi(g) \qquad ,
\end{equation}
où $\Gamma$ est le groupe de Galois d'une extension galoisienne finie $K'$ de $K$ qui splitte $U^{ef}$ (par exemple la plus petite telle extension, cf. 4.1), $a(g)$ le conducteur d'Artin sur $\Gamma$, et $\varphi(g)$ la fonction trace sur $\Gamma$ explicitée dans (4.3.0). Dans le cas de réduction semi-stable, on peut prendre $K'=K$, donc $\Gamma = e$, et on trouve

\begin{equation}\tag{4.6.3}
\alpha_s(A_K) = 0 \qquad (\mbox{cas de réduction semi-stable}).
\end{equation}
Pour terminer, explicitons un critère de réduction semi-stable qui résulte des réflexions générales de 4.2 :

\SGAvar{\underline{Proposition} 4.7. (\underline{Critère de M.} RAYNAUD \underline{de réduction semi-stable}). \underline{Soient} $S$ \underline{un trait,} $A_K$ \underline{un schéma abélien défini sur son corps de fonctions} $K$, $n$ \underline{en entier} $> 1$ \underline{premier à la caractéristique résiduelle} $p$ (*), \underline{et tel que} ${}_nA_K$ \underline{soit non ramifié par rapport à} $S$ \underline{i}.\underline{e}. \underline{tel que le groupe d'inertie opère trivialement sur} ${}_nA(\overline{K})$. \underline{Alors,} \underline{ou bien} $A_K$ \underline{a une réduction semi-stable sur} $S$, \underline{ou bien} $n = 2$, \underline{et si} $S$ \underline{est strictement local}, $A_K$ \underline{acquiert une réduction semi-stable sur une extension finie galoisienne} $K'$ \underline{de} $K$ \underline{à groupe de Galois isomorphe à} $(\mathbb{Z}/2\mathbb{Z})^r$.

}{\underline{Proposition} 4.7. (\underline{Critère de M.} RAYNAUD \underline{de réduction semi-stable}). \underline{Soient} $S$ \underline{un trait,} $A_K$ \underline{un schéma abélien défini sur son corps de fonctions} $K$, $n$ \underline{un entier} $> 1$ \underline{premier à la caractéristique résiduelle} $p$ (*), \underline{et tel que} ${}_nA_K$ \underline{soit non ramifié par rapport à} $S$ \underline{i}.\underline{e}. \underline{tel que le groupe d'inertie opère trivialement sur} ${}_nA(\overline{K})$. \underline{Alors,} \underline{ou bien} $A_K$ \underline{a une réduction semi-stable sur} $S$, \underline{ou bien} $n = 2$, \underline{et si} $S$ \underline{est strictement local}, $A_K$ \underline{acquiert une réduction semi-stable sur une extension finie galoisienne} $K'$ \underline{de} $K$ \underline{à groupe de Galois isomorphe à} $(\mathbb{Z}/2\mathbb{Z})^r$.

}{F7I-IX-047}(*) Comme il résulte d'un exemple qui m'a été communiqué par J. TATE, 4.7 ne reste pas valable quand $n$ est égal à la caractéristique résiduelle $p$, même si $K$ est de caractéristique nulle et si $A_K$ est une courbe elliptique à multiplication complexe.


\SGApage{I}{367}

On peut évidemment supposer $S$ strictement local. Avec les notations de 4.1 et de 4.2, prenons $K'$ le plus petit possible, i.e. supposons que $\Gamma$ opère \underline{fidèlement} sur $A^{*o}_o$. Par hypothèse, $\Gamma$ opère trivialement sur ${}_nA(\overline{K})$, et à fortiori sur ${}_nA(\overline{K})^{\mathrm{ef}} \simeq {}_nA^*_o(k) \supset {}_nA^{*o}_o(k)$. On conclut alors par les deux lemmes suivants :

\underline{Lemme} 4.7.1. (J.P. SERRE). \underline{Soient} $G$ \underline{un schéma en groupes commutatif sur un corps} $k$, \underline{extension d'un schéma abélien par un tore,} $n$ \underline{un nombre entier} $\geq 2$. \underline{Alors tout automorphisme d'ordre fini de} $G$ \underline{induisant l'identité sur} ${}_nG$ \underline{est réduit à l'identité si} $n \neq 2$, \underline{et est l'identité ou d'ordre deux si} $n = 2$.

\underline{Lemme} 4.7.2. \underline{Un groupe fini dont tout élément} $\neq 1$ \underline{est d'ordre} $2$ \underline{est isomorphe à} $(\mathbb{Z}/2\mathbb{Z})^r$, \underline{pour un} $r \geq 0$ \underline{convenable}.

\SGAvar{4.7.3. Rappelons la démonstration de 4.7.1, qui est bien connue. On note que grâce au fait que $n\,\mathrm{id}_G$ est un épimorphisme, l'hypothèse que $u$ induise l'identité sur ${}_nG$ i.e. que $1 - u$ induise $0$ sur ${}_nG$, s'exprime par la condition que l'on ait $1 - u = nv$ dans l'anneau des endomorphismes de $G$. Pour prouver la conclusion annoncée, on est ramené (en décomposant le groupe cyclique fini engendré par $u$ en ses composantes $p$-primaires, pour les divers nombres premiers $p$ divisant l'ordre de $u$) au cas où $u$ est d'ordre $p^N$, puissance d'un nombre premier. Soit $E$ le sous-anneau de $\mathrm{End}(G)$ engendré par $v$, qui est un quotient de $\mathbb{Z}[T]/(T^{p^N}-1)$, donc de type fini sur $\mathbb{Z}$, et qui de plus est sans torsion (ceci est conséquence formelle du fait que $G$ est divisible). Il se plonge donc dans $E_{\mathbb{Q}} = E\otimes_{\mathbb{Z}} \mathbb{Q}$, qui


}{4.7.3. Rappelons la démonstration de 4.7.1, qui est bien connue. On note que grâce au fait que $n\,\mathrm{id}_G$ est un épimorphisme, l'hypothèse que $u$ induise l'identité sur ${}_nG$ i.e. que $1 - u$ induise $0$ sur ${}_nG$, s'exprime par la condition que l'on ait $1 - u = nv$ dans l'anneau des endomorphismes de $G$. Pour prouver la conclusion annoncée, on est ramené (en décomposant le groupe cyclique fini engendré par $u$ en ses composantes $p$-primaires, pour les divers nombres premiers $p$ divisant l'ordre de $u$) au cas où $u$ est d'ordre $p^N$, puissance d'un nombre premier. Soit $E$ le sous-anneau de $\mathrm{End}(G)$ engendré par $v$, qui est un quotient de $\mathbb{Z}[T]/((1-nT)^{p^N}-1)$, donc de type fini sur $\mathbb{Z}$, et qui de plus est sans torsion (ceci est conséquence formelle du fait que $G$ est divisible). Il se plonge donc dans $E_{\mathbb{Q}} = E\otimes_{\mathbb{Z}} \mathbb{Q}$, qui


}{F7I-IX-048}\SGApage{I}{368}

\SGAvar{apparaît comme un quotient de $\mathbb{Q}[T]/(T^{p^N}-1)$, qui est un composé direct de corps extensions finies de $\mathbb{Q}$. Il en est donc de même de $E_{\mathbb{Q}}$. Considérant les images de $E,u,v$ dans les corps composants $K_i$ de $E_{\mathbb{Q}}$, on est donc ramené à voir que si $K = \mathbb{Q}(u)$ est un corps engendré par une racine primitive $p^n$-ème de l'unité $u$, et si $n$ est un entier $\geq 2$ tel que $u-1 \equiv 0$ mod $n$ (dans l'anneau des entiers de $K$), alors on a $u=1$, ou $u=-1$ et $n=2$. Or ceci résulte aussitôt du fait bien connu que si $u \neq 1$ i.e. $m \geq 1$, l'idéal $(u-1)$ est premier et l'idéal $(p)$ en est une puissance, dont l'exposant est $> 1$ sauf si $p=2,m=1$ i.e. si $u=-1$. Du fait que $(u-1)$ est premier et divise $p$ résulte en effet que $n \geq 2$ ne peut diviser $u-1$ que si $n=p$, et $p$ ne divise $u-1$ que si $(u-1) = (p)$ i.e. $u = -1$, donc $n = 2$.

}{apparaît comme un quotient de $\mathbb{Q}[T]/(T^{p^N}-1)$, qui est un composé direct de corps extensions finies de $\mathbb{Q}$. Il en est donc de même de $E_{\mathbb{Q}}$. Considérant les images de $E,u,v$ dans les corps composants $K_i$ de $E_{\mathbb{Q}}$, on est donc ramené à voir que si $K = \mathbb{Q}(u)$ est un corps engendré par une racine primitive $p^m$-ème de l'unité $u$, et si $n$ est un entier $\geq 2$ tel que $u-1 \equiv 0$ mod $n$ (dans l'anneau des entiers de $K$), alors on a $u=1$, ou $u=-1$ et $n=2$. Or ceci résulte aussitôt du fait bien connu que si $u \neq 1$ i.e. $m \geq 1$, l'idéal $(u-1)$ est premier et l'idéal $(p)$ en est une puissance, dont l'exposant est $> 1$ sauf si $p=2,m=1$ i.e. si $u=-1$. Du fait que $(u-1)$ est premier et divise $p$ résulte en effet que $n \geq 2$ ne peut diviser $u-1$ que si $n=p$, et $p$ ne divise $u-1$ que si $(u-1) = (p)$ i.e. $u = -1$, donc $n = 2$.

}{F7I-IX-049}4.7.4. Démontrons 4.7.2. Il suffit de prouver que $\Gamma$ est commutatif. Or si $x, y \in \Gamma$, on aura $(xy)(xy) = 1$, ce qui peut s'écrire aussi, puisque $x^2 = 1$, $y^2 = 1$ , $xy = (xy)^{-1} = y^{-1}x^{-1} = yx$, cqfd.

\SGAvar{\underline{Remarque} 4.7.5. La démonstration donnée 4.7.3, jointe à 4.7.2, nous montrent que si $G$ est un Groupe commutatif divisible d'un Topos, $n \geq 2$ un entier, et $\Gamma$ un groupe fini, opèrant fidèlement sur $G$ en induisant l'identité sur ${}_nG$, alors $\Gamma$ est réduit à l'identité si $n \geq 3$, et est isomorphe à un groupe $(\mathbb{Z}/2\mathbb{Z})^r$ si $n=2$.


}{\underline{Remarque} 4.7.5. La démonstration donnée dans 4.7.3, jointe à 4.7.2, nous montre que si $G$ est un Groupe commutatif divisible d'un Topos, $n \geq 2$ un entier, et $\Gamma$ un groupe fini, opérant fidèlement sur $G$ en induisant l'identité sur ${}_nG$, alors $\Gamma$ est réduit à l'identité si $n \geq 3$, et est isomorphe à un groupe $(\mathbb{Z}/2\mathbb{Z})^r$ si $n=2$.


}{F7I-IX-050}\SGApage{I}{369}


\SGAsectiontoc{5. Le théorème d'orthogonalité dans le cas $\ell = p$ (cas semi-stable). Dualité des schémas abéliens $B_{\mathrm{o}}$ et $B'_{\mathrm{o}}$. Caractérisation de la partie fixe. Critères de bonne réduction, de bonne réduction essentielle et semi-stable.}{5}
5. \underline{Le théorème d'orthogonalité dans le cas} $\ell = p$ (\underline{cas semi-stable}). \underline{Dualité des schémas abéliens} $B_{\mathrm{o}}$ \underline{et} $B'_{\mathrm{o}}$. \underline{Caractérisation de la partie fixe}. \underline{Critères de bonne réduction, de bonne réduction essentielle et semi-stable}.

5.1. Plaçons nous dans le cas 2.2.3 où $S$ est hensélien et $A^{\mathrm{o}}$ est $\ell$-divisible, de sorte que la \underline{partie fixe}

\begin{equation*}
T_\ell(A^{\mathrm{o}})^f \hookrightarrow T_\ell(A)
\end{equation*}
est définie comme le pro-$\ell$-groupe de Barsotti-Tate formé des $({}_{\ell^\nu}A^{\mathrm{o}})^f$ . Nous allons maintenant donner une définition de la partie torique (2.3.2) de $T_\ell(A^{\mathrm{o}})$, indépendante de l'hypothèse $\ell \neq p$ faite dans 2.3, mais supposant en revanche $S$ \underline{complet} (*). Pour ceci, introduisons pour tout entier $n \geq 0$ le schéma artinien local

\begin{equation}\tag{5.1.1}
S_n = \mathrm{Spec}(V/\underline{m}^{n+1}) \qquad ,
\end{equation}
n.ème voisinage infinitésimal de $s$ dans $S$, et pour tout $S$-schéma $X$ soit $X_n = X\times_S S_n$. Comme $(A^{\mathrm{o}})_n = (A_n)^{\mathrm{o}}$ est lisse sur $S_n$, il est connu (SGA 3 IX 3.6 bis) que $T_{\mathrm{o}}$ est la fibre spéciale d'un unique sous-tore

\SGAvar{\begin{equation}\tag{5.1.2}
T_n \subset A_n^{\mathrm{o}} \qquad .
\end{equation}
D'après l'unicité des $T_n$, ceux-ci se déduisent les uns des autres par les changements de base $S_{n'} \longrightarrow S_n$ . En d'autres termes, il définissent un \underline{sous-tore formel} $\hat{T}$ du groupe $\hat{A}^{\mathrm{o}}$ (dans la catégorie des schémas formels sur $S$) complété formel de $A^{\mathrm{o}}$ le long de la fibre spéciale :

}{\begin{equation}\tag{5.1.2}
T_n \subset A_n^{\mathrm{o}} \qquad .
\end{equation}
D'après l'unicité des $T_n$, ceux-ci se déduisent les uns des autres par les changements de base $S_{n'} \longrightarrow S_n$ . En d'autres termes, ils définissent un \underline{sous-tore formel} $\hat{T}$ du groupe $\hat{A}^{\mathrm{o}}$ (dans la catégorie des schémas formels sur $S$) complété formel de $A^{\mathrm{o}}$ le long de la fibre spéciale :

}{F7I-IX-051}\begin{equation}\tag{5.1.3}
\hat{T} \subset \hat{A}^{\mathrm{o}} \qquad .
\end{equation}
(*) Cette hypothèse n'est d'ailleurs pas essentielle, comme on montre dans 6.1 ci-dessous.


\SGApage{I}{370}

Comme le foncteur $X \longmapsto \hat{X}$ (complétion formelle le long de la fibre spéciale) induit une équivalence entre la catégorie des schémas finis sur $S$, et celle des schémas formels finis sur $S$ (EGA III 4.8), nous allons identifier simplement un schéma fini sur $S$ avec le schéma formel correspondant. Avec cette identification, la définition des "parties fixes" nous donne aussitôt pour tout $\nu \geq 0$ :

\begin{equation}\tag{5.1.4}
({}_{\ell^\nu}A^{\mathrm{o}})^f \ = \ {}_{\ell^\nu}(\hat{A}^{\mathrm{o}}) \qquad ,
\end{equation}
ce qu'on peut écrire aussi comme une identité de systèmes projectifs

\begin{equation}\tag{5.1.5}
T_\ell(A^{\mathrm{o}})^f \ = \ T_\ell(\hat{A}^{\mathrm{o}}) \qquad .
\end{equation}
On obtient donc un sous-pro-groupe de Barsotti-Tate de la partie fixe de $T_\ell(A^{\mathrm{o}})$ en prenant $T_\ell(\hat{T}) \subset T_\ell(\hat{A}^{\mathrm{o}})$. On définira alors la \underline{partie torique} de $T_\ell(A^{\mathrm{o}})$ comme

\begin{equation}\tag{5.1.6}
T_\ell(A^{\mathrm{o}})^t \stackrel{\mathrm{dfn}}{=} T_\ell(\hat{T}) \subset T_\ell(\hat{A}^{\mathrm{o}}) = T_\ell(A^{\mathrm{o}})^f \qquad .
\end{equation}
Donc par définition, on aura

\begin{equation}\tag{5.1.7}
T_\ell(A^{\mathrm{o}})^t \ = \ ({}_{\ell^\nu}\hat{T})_{\nu \geq 0} \qquad ,
\end{equation}
le $S$-groupe fini et plat ${}_{\ell^\nu}\hat{T}$ étant lui-même décrit comme la limite inductive des $S_n$-groupes finis et plats

\begin{equation}\tag{5.1.8}
{}_{\ell^\nu}(T_n) \ , \ n \ \mbox{variable} \qquad ,
\end{equation}
i.e. son algèbre affine est décrite comme la limite projective des algèbres affines des groupes (5.1.8). Bien entendu, on aura

\begin{equation}\tag{5.1.9}
(T_\ell(A^{\mathrm{o}})^t)_n \ = \ T_\ell(T_n) \qquad ,
\end{equation}

\SGApage{I}{371}

ce qui, pour $n=0$, n'est autre que la formule (2.3.3) par laquelle on avait caractérisé la partie torique (comme groupe pro-fini-étale) dans le cas $\ell \neq p$. Ceci montre que la définition (5.1.6) est compatible avec celle donnée dans 2.3 lorsque $\ell \neq p$.

Passant aux fibres génériques, les parties fixe et torique de $T_\ell(A^o)$ nous définissent des parties fixes et toriques dans $T_\ell(A_K)$ :

\begin{equation}\tag{5.1.10}
T_\ell(A_K)^t \subset T_\ell(A_K)^f \subset T_\ell(A_K) \quad .
\end{equation}
Ce sont des inclusions de pro-$\ell$-groupes de Barsotti-Tate. Considérons la filtration analogue, correspondante au schéma abélien dual $A'_K$ :

\begin{equation}\tag{5.1.11}
T_\ell(A'_K)^t \subset T_\ell(A'_K)^f \subset T_\ell(A'_K) \quad .
\end{equation}
On a alors le résultat suivant, qui généralise 2.4 (du moins dans le cas $S$ complet, auquel on pourrait facilement se réduire dans 2.4) :

\SGAvar{\underline{Théorème} 5.2. \underline{Supposons que} $S$ \underline{soit un trait} complet, \underline{et soient} $A_K$ \underline{un schéma abélien sur} $K$, $A'_K$ \underline{le schéma abélien dual}, $\ell$ \underline{un nombre premier tel que le modèle de Néron connexe} $A^o$ \underline{soit} $\ell$\underline{-divisible} (cf. 2.2.1), \underline{de sorte qu'on dispose alors des filtrations} (5.1.10) \underline{et} (5.1.11) \underline{des pro-groupes} $T_\ell(A_K)$ \underline{et} $T_\ell(A_{K'})$. \underline{Ceci posé}, \underline{les parties toriques se reconstituent en termes des parties fixes par la formule}

\begin{equation}\tag{5.2.1}
T_\ell(A_K)^t = T_\ell(A_K)^f \ \cap \ (T_\ell(A'_K)^f)^\perp
\end{equation}
\underline{et la formule analogue obtenu en échangeant les rôles de} $A_K$ \underline{et de} $A'_K$ , \underline{où le signe} $\perp$ \underline{désigne l'orthogonal par rapport à l'accouplement canonique} (1.0.2).


}{\underline{Théorème} 5.2. \underline{Supposons que} $S$ \underline{soit un trait} complet, \underline{et soient} $A_K$ \underline{un schéma abélien sur} $K$, $A'_K$ \underline{le schéma abélien dual}, $\ell$ \underline{un nombre premier tel que le modèle de Néron connexe} $A^o$ \underline{soit} $\ell$\underline{-divisible} (cf. 2.2.1), \underline{de sorte qu'on dispose alors des filtrations} (5.1.10) \underline{et} (5.1.11) \underline{des pro-groupes} $T_\ell(A_K)$ \underline{et} $T_\ell(A'_K)$. \underline{Ceci posé}, \underline{les parties toriques se reconstituent en termes des parties fixes par la formule}

\begin{equation}\tag{5.2.1}
T_\ell(A_K)^t = T_\ell(A_K)^f \ \cap \ (T_\ell(A'_K)^f)^\perp
\end{equation}
\underline{et la formule analogue obtenue en échangeant les rôles de} $A_K$ \underline{et de} $A'_K$ , \underline{où le signe} $\perp$ \underline{désigne l'orthogonal par rapport à l'accouplement canonique} (1.0.2).


}{F7I-IX-052}\SGApage{I}{372}

\underline{Enfin}, \underline{lorsque} $A_K$ \underline{est à réduction semi-stable sur} $S$ (\underline{ce qui résulte des autres hypothèses dans le cas} $\ell = p$), \underline{alors on a même la formule}

\SGAvar{\begin{equation}\tag{5.2.2}
T_\ell(A_K)^t = (T_\ell(A'_K)^f)^\perp
\end{equation}
5.2.3. La signification de la formule (5.2.1) étant bien claire si $\ell$ est premier à la caractéristique de $K$ (où elle peut s'interpréter en passant aux $\mathbb{Z}_\ell$-modules ordinaires, à action de $\pi = \mathrm{Gal}(\overline{K}/K)$, formés des $\overline{K}$-points des pro-groupes en jeu, comme dans 2.4), explicitons-là dans le cas général, en utilisant la notion de dual

}{\begin{equation}\tag{5.2.2}
T_\ell(A_K)^t = (T_\ell(A'_K)^f)^\perp
\end{equation}
5.2.3. La signification de la formule (5.2.1) étant bien claire si $\ell$ est premier à la caractéristique de $K$ (où elle peut s'interpréter en passant aux $\mathbb{Z}_\ell$-modules ordinaires, à action de $\pi = \mathrm{Gal}(\overline{K}/K)$, formés des $\overline{K}$-points des pro-groupes en jeu, comme dans 2.4), explicitons-la dans le cas général, en utilisant la notion de dual

}{F7I-IX-053}\begin{equation}\tag{5.2.3.1}
D(G) = \mathrm{Hom}(G,T_\ell(\mathbb{G}_m))
\end{equation}
d'un pro-$\ell$-groupe de Barsotti-Tate $G = (G(\nu))$, qui par définition est le système projectif des duals de Cartier $D(G(\nu))$ (SGA 3 VII 3.3.1 et VIII 1), où les morphismes de transition se définissent de façon évidente (par transposition à partir des morphismes d'\underline{inclusion} entre les $G(\nu)$, par exemple). La théorie de dualité des variétés abéliennes nous apprend que l'accouplement mentionné dans 5.2 définit en fait un \underline{isomorphisme}

\begin{equation}\tag{5.2.3.2}
T_\ell(A'_K) \ \simeq \ D(T_\ell(A_K)) \quad ,
\end{equation}
d'où par composition un morphisme

\begin{equation}\tag{5.2.3.3}
T_\ell(A_K)^f \longrightarrow D(T_\ell(A'_K)^f) \quad ,
\end{equation}
et la formule (5.2.1) affirme que le noyau de l'homomorphisme précédent (qui à priori est un sous-pro-groupe du premier terme) n'est autre que la partie torique $T_\ell(A_K)^t$ . On notera \underline{qu'il s'agit d'un isomorphisme de pro-objets}, \underline{et non de systèmes projectifs} : nous n'affirmons pas (et il


\SGApage{I}{373}

\SGAvar{serait faux en général) que pour tout entier $\nu$, $T_\ell(A_K)^t(\nu)$ soit égal au noyau de $T_\ell(A_K)^f(\nu) \longrightarrow D(T_\ell(A'_K)^f)(\nu)$ ; ce dernier pourra être plus grand à priori. (Nous verrons cependant plus bas qu'il ne l'est pas si $A_K$ a réduction semi-stable, ce qui est le cas en particulier, d'après les hypohtèses faites dans 5.2, lorsque $\ell = p$).

}{serait faux en général) que pour tout entier $\nu$, $T_\ell(A_K)^t(\nu)$ soit égal au noyau de $T_\ell(A_K)^f(\nu) \longrightarrow D(T_\ell(A'_K)^f)(\nu)$ ; ce dernier pourra être plus grand à priori. (Nous verrons cependant plus bas qu'il ne l'est pas si $A_K$ a réduction semi-stable, ce qui est le cas en particulier, d'après les hypothèses faites dans 5.2, lorsque $\ell = p$).

}{F7I-IX-054}5.2.4.  La démonstration de (5.2.1) ne fait que répéter, pratiquement sans changement, celle de 2.4. On note que l'accouplement induit

\begin{equation}\tag{5.2.4.1}
T_\ell(A_K)^f \quad \times \quad T_\ell(A'_K)^f \longrightarrow T_\ell(\mathbb{G}_{mK})
\end{equation}
donnant naissance à (5.2.3.3) est induit par un accouplement

\begin{equation}\tag{5.2.4.2}
T_\ell(A^o)^f \quad \times \quad T_\ell({A'}^o)^f \longrightarrow T_\ell(\mathbb{G}_{mS})
\end{equation}
de pro-groupes de Barsotti-Tate sur $S$, lui-même induit par l'accouplement de pro-groupes

\SGAvar{\begin{equation}\tag{5.2.4.3}
T_\ell(A^o) \quad \times \quad T_\ell({A'}^o) \longrightarrow T_\ell(\mathbb{G}_{mS})
\end{equation}
déduite de la biextension $W$ de $(A^o,{A'}^o)$ par $\mathbb{G}_{mS}$ . Par suite, l'accouplement

}{\begin{equation}\tag{5.2.4.3}
T_\ell(A^o) \quad \times \quad T_\ell({A'}^o) \longrightarrow T_\ell(\mathbb{G}_{mS})
\end{equation}
déduit de la biextension $W$ de $(A^o,{A'}^o)$ par $\mathbb{G}_{mS}$ . Par suite, l'accouplement

}{F7I-IX-055}\begin{equation}\tag{5.2.4.4}
T_\ell(A_o^{o}) \quad \times \quad T_\ell({A'}_o^{o}) \longrightarrow T_\ell(\mathbb{G}_{mk}) \qquad ,
\end{equation}
déduit de (5.2.4.2) par passage aux fibres spéciales, n'est autre que celui déduit de la biextension $W_o$ de $(A_o^{o},{A'}_o^{o})$ par $\mathbb{G}_{mk}$ . En vertu de VIII 4.11 et de 1.5, l'annulateur gauche de l'accouplement (5.2.4.4) est donc égal à $T_\ell(T_o)$, qui n'est autre que la fibre spéciale de la partie torique $T_\ell(A^o)^t$. On conclut donc, en revenant à la définition des annulateurs gauches explicitée dans 5.2.3, à l'aide du


\SGApage{I}{374}

\underline{Lemme} 5.2.5. \underline{Soit} $u : G \longrightarrow H$ \underline{un homomorphisme de pro}-$\ell$-\underline{groupes de Barsotti-Tate sur} $S$. \underline{Si} $u_o : G_o \longrightarrow H_o$ \underline{est nul}, \underline{alors} $u$ \underline{est nul}. \underline{Si} $u_o : G_o \longrightarrow H_o$ \underline{est une isogénie} (resp. \underline{un isomorphisme}), $u$ \underline{est une isogénie} (resp. \underline{un isomorphisme}).

Nous appliquerons en effet la première assertion de 5.2.5 au morphisme composé

\begin{equation*}
T_\ell(A^o)^t \longrightarrow T_\ell(A^o)^f \longrightarrow D(T_\ell({A'}^o)^f) \qquad ,
\end{equation*}
où le deuxième morphisme est celui déduit de l'accouplement (5.2.4.2), on trouve que ce composé est nul, donc que l'accouplement envisagé s'annule sur la partie torique à gauche. Par symétrie, il s'annule donc sur la partie torique à droite, donc il provient d'un accouplement sur les quotients, ou ``morceaux abéliens''

\begin{equation*}
T_\ell(A^o)^{ab} \quad \times \quad T_\ell({A'}^o)^{ab} \longrightarrow T_\ell(\mathbb{G}_m) \qquad ,
\end{equation*}
lesquels quotients sont encore des pro-groupes de Barsotti-Tate, comme on vérifie aisément. Il reste à prouver que le morphisme correspondant à cet accouplement

\begin{equation*}
T_\ell(A^o)^{ab} \longrightarrow D(T_\ell({A'}^o)^{ab})
\end{equation*}
est une isogénie, ce qui résulte encore de 5.2.5.

Nous laisserons au lecteur la démonstration (facile) de 5.2.5, qui est un résultat valable sur toute base noethérienne connexe. Nous trouverons d'ailleurs plus bas (7.4.7), par une analyse un peu plus précise des biextensions $W_n$ des $(A_n^{o},{A'}_n^{o})$, une démonstration directe assez évidente des deux cas particuliers où nous avons eu à utiliser 5.2.5.


\SGApage{I}{375}

5.2.6. Supposons maintenant que $A_K$ ait une réduction semi-stable sur $S$, de sorte que nous pouvons prendre pour $\ell$ n'importe quel nombre premier (distinct ou non de $p$). D'ailleurs, en vertu du critère galoisien de réduction semi-stable 3.5, on a pour $\ell \neq p$

\begin{equation}\tag{5.2.7}
(T_\ell(A'_K)^f)^\perp \quad (= T_\ell(A_K)^t) \subset T_\ell(A_K)^f \quad .
\end{equation}
Il reste à prouver, pour terminer la preuve de 5.2, que cette inclusion est valable sans restriction sur $\ell$.

En effet, le membre de gauche est canoniquement isomorphe à $D(T_\ell(A'_K)/T_\ell(A'_K)^f)$, et est donc de ``rang'' (ou ``hauteur'') égal à $2n - 2\alpha - \mu = \mu$ (notations de (2.1.9), où par hypothèse $\lambda = 0$), qui est aussi celui de $T_\ell(A_K)^t$. Comme l'inclusion

\SGAvar{\begin{equation}\tag{*}
T_\ell(A_K)^t \subset \ T_\ell(A_K)
\end{equation}
est \underline{stricte}, dans le sens qu'elle induit déjà un monomorphisme pour tous les termes du systèmes projectifs envisagés, il s'ensuit aisément que le premier membre de $(*)$ est maximal parmi les sous-pro-groupes de Barsotti-Tate de $T_\ell(A_K)$ ayant même rang que lui, de sorte qu'on a bien l'égalité de (5.2.7), qui équivaut à l'inclusion entre les termes extrêmes de 5.2.7 par le théorème d'orthogonalité.

}{\begin{equation}\tag{*}
T_\ell(A_K)^t \subset \ T_\ell(A_K)
\end{equation}
est \underline{stricte}, dans le sens qu'elle induit déjà un monomorphisme pour tous les termes des systèmes projectifs envisagés, il s'ensuit aisément que le premier membre de $(*)$ est maximal parmi les sous-pro-groupes de Barsotti-Tate de $T_\ell(A_K)$ ayant même rang que lui, de sorte qu'on a bien l'égalité de (5.2.7), qui équivaut à l'inclusion entre les termes extrêmes de 5.2.7 par le théorème d'orthogonalité.

}{F7I-IX-056}5.3. Utilisant la correspondance par orthogonalité entre sous-pro-groupes de Barsotti-Tate stricts (au sens précisé plus haut) des pro-groupes duals $T_\ell(A_K)$ et $T_\ell(A'_K)$, on déduit formellement de (5.2.2) que l'accouplement (1.0.2) entre ces pro-groupes induit un accouplement

\begin{equation}\tag{5.3.2}
T_\ell(A_K)^{ab} \ \times \ T_\ell(A'_K)^{ab} \longrightarrow T_\ell(\mathbb{G}_{mK})
\end{equation}

\SGApage{I}{376}

qui est une \underline{dualité parfaite}, i.e. identifiant chacun des deux facteurs du premier membre au dual de l'autre (au sens de 5.2.2). (NB les exposants ab désignent les ``parties abéliennes'', quotient de la partie fixe par la partie torique). Il est clair d'autre part que (5.3.2) est induit par l'accouplement

\begin{equation}\tag{5.3.3}
T_\ell(A^o)^{ab} \ \times \ T_\ell(A^{\prime o})^{ab} \longrightarrow T_\ell(\mathbb{G}_{mS})
\end{equation}
déduit par passages aux quotients de l'accouplement (5.2.4.2). Celui-ci étant une dualité parfaite sur les fibres génériques, i.e. l'homomorphisme correspondant

\begin{equation}\tag{5.3.4}
T_\ell(A^o)^{ab} \longrightarrow D(T_\ell(A^{\prime o})^{ab})
\end{equation}
étant un isomorphisme sur les fibres génériques, nous avons évidemment envie d'en déduire qu'il en est de même sur $S$ tout entier, ou ce qui revient au même, sur les fibres spéciales. Cette déduction est valable, trivialement, si $\ell \neq p$ ; elle est valable également, du moins dans le cas d'inégales caractéristiques, si $\ell = p$, grâce à un théorème profond de TATE [33, th.4]. Nous allons cependant pouvoir l'établir directement, sans l'aide du théorème de Tate, plus bas (7.4.6).

Pour préciser la signification géométrique du résultat que nous désirons établir, notons que l'accouplement déduit de (5.3.3) par passage aux fibres spéciales n'est autre que l'accouplement

\begin{equation}\tag{5.3.5}
T_\ell(B_o) \ \times \ T_\ell(B'_o) \longrightarrow T_\ell(\mathbb{G}_{mk})
\end{equation}
déduit de la biextension $W_o$ de $(A_o^{\,o}, A_o^{\prime\,o})$ par passage au quotient, $B_o$ et $B'_o$


\SGApage{I}{377}

désignant respectivement les parties abéliennes de $A_o^o$ et $A_o^{\prime o}$. Comme, en vertu de VIII 4.8, la biextension $W_o$ provient d'une biextension $W_o^{ab}$ de $(B_o,B'_o)$ par $\mathbb{G}_{mk}$, déterminée à isomorphisme unique près, l'accouplement (5.3.5) peut aussi s'interpréter comme déduit de la correspondance divisorielle $W_o^{ab}$, et l'homomorphisme correspondant

\begin{equation}\tag{5.3.6}
T_\ell(B_o) \longrightarrow D(T_\ell(B'_o)) \simeq T_\ell(D(B'_o))
\end{equation}
(où dans le dernier terme, la lettre $D$ désigne le schéma abélien dual) n'est autre que celui déduit par fonctorialité de $T_\ell$ de l'homomorphisme

\begin{equation}\tag{5.3.7}
B_o \longrightarrow D(B'_o)
\end{equation}
défini par $W_o^{ab}$. Donc le fait que (5.3.5) soit une dualité parfaite signifie aussi, par le dictionnaire bien connu (cf. [ 5]) que le noyau de (5.3.7), qui est un groupe fini sur $k$ en vertu de 1.5, a une composante $\ell$-primaire nulle. Dire qu'il en est ainsi pour tout $\ell$ signifie donc que l'isogénie (5.3.7) est un isomorphisme, i.e. que la correspondance divisorielle $W_o^{ab}$ est une dualité parfaite entre $B_o$ et $B'_o$ . Les considérations qui précédaient montrent seulement que le noyau de (5.3.7) est un $p$-groupe. Enonçons donc le résultat plus complet :

\SGAvar{\underline{Théorème} 5.4. \underline{Soit} $S$ \underline{un trait, et} $A_K$ \underline{un schéma abélien sur} $S$ \underline{ayant une réduction semi-stable sur} $S$. \underline{Désignons par} $B_o$, $B'_o$ \underline{les parties abéliennes des modèles de Néron de} $A_K$ \underline{et du schéma abélien dual} $A'_K$ \underline{respectivement}. \underline{Alors la correspondance divisorielle canonique} $W_o^{ab}$ \underline{sur} $(B_o,B'_o)$, \underline{déduite par} VIII 4.8 \underline{de la biextension canonique} $W_o$ \underline{de} $(A_o^o,A_o^{\prime o})$ (\underline{fibre spéciale de la biextension} $W$ \underline{de} 1.4), \underline{est une dualité parfaite entre} $B_o$ \underline{et} $B'_o$,


}{\underline{Théorème} 5.4. \underline{Soit} $S$ \underline{un trait, et} $A_K$ \underline{un schéma abélien sur} $K$ \underline{ayant une réduction semi-stable sur} $S$. \underline{Désignons par} $B_o$, $B'_o$ \underline{les parties abéliennes des modèles de Néron de} $A_K$ \underline{et du schéma abélien dual} $A'_K$ \underline{respectivement}. \underline{Alors la correspondance divisorielle canonique} $W_o^{ab}$ \underline{sur} $(B_o,B'_o)$, \underline{déduite par} VIII 4.8 \underline{de la biextension canonique} $W_o$ \underline{de} $(A_o^o,A_o^{\prime o})$ (\underline{fibre spéciale de la biextension} $W$ \underline{de} 1.4), \underline{est une dualité parfaite entre} $B_o$ \underline{et} $B'_o$,


}{F7I-IX-057}\SGApage{I}{378}

\underline{i.e.} \underline{définit un isomorphisme de chacun de ces schémas abéliens avec le dual de l'autre}.

La démonstration de 5.4 dans le cas général va être donnée au n° 7, comme conséquence d'une construction fort naturelle et importante due à RAYNAUD. Notons seulement ici que pour prouver 5.4, on peut supposer $S$ complet, compte tenu de 3.3.

5.5. Pour exprimer commodément les facteurs extrêmes dans les filtrations (5.1.10) et (5.1.11), dans le cas de la réduction semi-stable, introduisons les schémas en groupes constants tordus sur $S$

\begin{equation}\tag{5.5.1}
\underline{M}'_o = D(T_o) \quad , \quad \underline{M}_o = D(T'_o) \quad ,
\end{equation}
duals des tores maximaux $T_o$, $T'_o$ des fibres spéciales de $A,A'$ (SGA 3 X 5.7). Comme $S$ est local complet, ces schémas en groupes peuvent être considérés comme les fibres spéciales de schémas constants tordus

\begin{equation}\tag{5.5.2}
\underline{M}' \quad , \quad \underline{M} \ \text{sur} \ S \qquad ;
\end{equation}
ainsi, si le tore $T_o$ est déployé (par exemple si $k$ est séparablement clos) i.e. si $\underline{M}'_o$ est le groupe constant défini par le groupe ordinaire

\begin{equation}\tag{5.5.3}
M' = \mathrm{Hom}_{k-gr}(T_o,\mathbb{G}_{mk}) \quad ,
\end{equation}
alors $\underline{M}'$ est le groupe constant défini par $M'$ :

\begin{equation*}
\underline{M}' = M'_S \quad .
\end{equation*}
On peut d'ailleurs noter que, puisque $T_o$ et $T'_o$ sont isogènes (2.2.6), $T'_o$ est déployé si $T_o$ l'est, de sorte qu'il y a lieu d'introduire également


\SGApage{I}{379}

dans ce cas

\begin{equation}\tag{5.5.4}
M = \mathrm{Hom}_{k\text{-}gr}(T'_o, \mathbb{G}_{mk}) \quad ,
\end{equation}
et on aura alors

\begin{equation*}
\underline{M} \;=\; M_S \quad .
\end{equation*}
Ceci posé, on aura évidemment, pour tout entier $n$, des isomorphismes

\begin{equation}\tag{5.5.5}
\underline{M}'_n \simeq D(T_n) \quad , \quad \underline{M}_n \simeq D(T'_n) \quad ,
\end{equation}
où $T'_n$ est défini en termes de $A'$ comme $T_n$ en termes de $A$ (5.1.2), d'où on déduit des isomorphismes canoniques (où on pose $\mathbb{Z}_\ell(1) = T_\ell(\mathbb{G}_{mS})$) :

\begin{equation}\tag{5.5.6}
T_\ell(T_n) \simeq \check{\underline{M}}'_n \otimes_{\mathbb{Z}} \mathbb{Z}_\ell(1) \quad , \quad T_\ell(T'_n) \simeq \check{\underline{M}}_n \otimes_{\mathbb{Z}} \mathbb{Z}_\ell(1) \quad ,
\end{equation}
et par passage à la limite sur $n$ des isomorphismes canoniques :

\begin{equation}\tag{5.5.7}
\left\{\begin{aligned} T_\ell(A^o)^t &= T_\ell(\hat{T}) \simeq \check{\underline{M}}' \otimes \mathbb{Z}_\ell(1)\\[2ex] T_\ell({A'}^o)^t &= T_\ell(\hat{T}') \simeq \check{\underline{M}} \otimes \mathbb{Z}_\ell(1) \quad . \end{aligned}\right.
\end{equation}
Dans le cas de la réduction semi-stable, on en conclut par dualité, en utilisant le théorème d'orthogonalité 5.2 et la relation (5.2.2), des isomorphismes canoniques :

\SGAvar{\begin{equation}\tag{5.5.8}
\left\{\begin{aligned} T_\ell(A_K)/T_\ell(A_K)^f \quad &\simeq \quad \underline{M}_K \otimes_{\mathbb{Z}} \mathbb{Z}_\ell\\[2ex] T_\ell(A')/T_\ell(A'_K)^f \quad &\simeq \quad \underline{M}'_K \otimes_{\mathbb{Z}} \mathbb{Z}_\ell \quad . \end{aligned}\right.
\end{equation}
On trouve de même, grâce à la même formule (5.2.2), des isomorphismes canoniques :

}{\begin{equation}\tag{5.5.8}
\left\{\begin{aligned} T_\ell(A_K)/T_\ell(A_K)^f \quad &\simeq \quad \underline{M}_K \otimes_{\mathbb{Z}} \mathbb{Z}_\ell\\[2ex] T_\ell(A'_K)/T_\ell(A'_K)^f \quad &\simeq \quad \underline{M}'_K \otimes_{\mathbb{Z}} \mathbb{Z}_\ell \quad . \end{aligned}\right.
\end{equation}
On trouve de même, grâce à la même formule (5.2.2), des isomorphismes canoniques :

}{F7I-IX-058}\begin{equation}\tag{5.5.9}
\left\{\begin{aligned} T_\ell(A_K)/T_\ell(A_K)^t \quad &\simeq \quad D(T_\ell(A'_K)^f)\\[2ex] T_\ell(A'_K)/T_\ell(A'_K)^t \quad &\simeq \quad D(T_\ell(A_K)^f) \quad , \end{aligned}\right.
\end{equation}

\SGApage{I}{380}

\SGAvar{où on se rappelera que les groupes $T_\ell(A_K)^f$ et $T_\ell(A'_K)^f$ qui figurent au second membre sont les restrictions à $\eta$ de pro-$\ell$-groupes de Barsotti-Tate définis sur $S$, $T_\ell(A^o)^f$ et $T_\ell({A'}^o)^f$ respectivement ; la lettre $D$ désigne ici les duals au sens des pro-groupes de Barsotti-Tate, savoir $\underline{\mathrm{Hom}}( - , \mathbb{Z}_\ell(1))$. Ainsi :

}{où on se rappellera que les groupes $T_\ell(A_K)^f$ et $T_\ell(A'_K)^f$ qui figurent au second membre sont les restrictions à $\eta$ de pro-$\ell$-groupes de Barsotti-Tate définis sur $S$, $T_\ell(A^o)^f$ et $T_\ell({A'}^o)^f$ respectivement ; la lettre $D$ désigne ici les duals au sens des pro-groupes de Barsotti-Tate, savoir $\underline{\mathrm{Hom}}( - , \mathbb{Z}_\ell(1))$. Ainsi :

}{F7I-IX-059}\underline{Proposition} 5.6. \underline{Supposons} $S$ \underline{complet et que} $A_K$ \underline{ait réduction semi-stable sur} $S$. \underline{Alors pour tout nombre premier} $\ell$ (\underline{qui peut être égal à la caractéristique résiduelle de} $S$), \underline{le pro}-$\ell$-\underline{groupe de Barsotti-Tate} $T_\ell(A_K)^t$ \underline{est la fibre générique d'un pro-groupe de Barsotti-Tate sur} $S$, "\underline{de type torique}", \underline{donné par} (5.5.7), \underline{tandis que} $T_\ell(A_K)/T_\ell(A_K)^f$ \underline{est la fibre générique d'un pro-groupe de Barsotti-Tate sur} $S$ \underline{de type pro-étale}, \underline{donné par} (5.5.8). \underline{Enfin les groupes} $T_\ell(A_K)^f$ \underline{et} $T_\ell(A_K)/T_\ell(A_K)^t$ \underline{sont les fibres génériques de pro-groupes de Barsotti-Tate définis sur} $S$, \underline{égaux respectivement à} $T_\ell(A^o)^f$ \underline{et à} $D(T_\ell({A'}^o)^f) = \underline{\mathrm{Hom}}\,(T_\ell({A'}^o)^f, \mathbb{Z}_\ell(1))$.

\SGAvar{5.6.1. On notera qu'il revient au même pour le revêtement étale $S'$ de déployer $T_\ell(A_K)^t$ ou de déployer $T_\ell(A_K)/T_\ell(A_K)^f$, car cela signifie respectivement que $S'$ déploye les groupes constants tordus $\underline{M}$ et $\underline{M}'$, or ceux-ci sont isogènes, puisque $T_o$ et $T'_o$ le sont. On voit en même temps qu'il existe une plus petite extension étale $S'$ de $S$ qui fait l'affaire.

}{5.6.1. On notera qu'il revient au même pour le revêtement étale $S'$ de déployer $T_\ell(A_K)^t$ ou de déployer $T_\ell(A_K)/T_\ell(A_K)^f$, car cela signifie respectivement que $S'$ déploie les groupes constants tordus $\underline{M}'$ et $\underline{M}$, or ceux-ci sont isogènes, puisque $T_o$ et $T'_o$ le sont. On voit en même temps qu'il existe une plus petite extension étale $S'$ de $S$ qui fait l'affaire.

}{F7I-IX-060}Signalons aussi la conséquence facile suivante des définitions de 2.2.3 :


\SGApage{I}{381}

\underline{Proposition} 5.7.  \underline{Supposons} $A^o$ \underline{$\ell$-divisible}, $S$ \underline{hensélien}. \underline{Soit} $F_K$ \underline{un} \underline{faisceau $\ell$-adique constructible sur} $K$ \underline{qui soit} \underline{``non ramifié''} \underline{i.e.} \underline{qui} \underline{se prolonge en un faisceau $\ell$-adique constant tordu} $F$ \underline{sur} $S$, \underline{et soit} $u: F_K \longrightarrow T_\ell(A_K)$ \underline{un homomorphisme de pro-groupes}, \underline{alors} $u$ \underline{se factorise} \underline{en un homomorphisme de} $F$ \underline{dans} $T_\ell(A_K)^f$.

En effet, on aura

\begin{equation}\tag{5.7.1}
u = (u(\nu))_{\nu \geqslant 0}, \ \mbox{avec} \ u(\nu): F_K(\nu) \longrightarrow T_\ell(A_K)(\nu) = {}_{\ell^\nu}A_K \quad .
\end{equation}
Par la propriété universelle de $A$, comme $F(\nu)$ est étale donc lisse sur $S$, le morphisme $u(\nu)$ provient d'un homomorphisme

\begin{equation*}
v(\nu) \ : \ F(\nu) \longrightarrow {}_{\ell^\nu}A \quad ,
\end{equation*}
et par définition de la partie fixe, cet homomorphisme se factorise en

\begin{equation*}
v(\nu) \ : \ F(\nu) \longrightarrow ({}_{\ell^\nu}A)^f \quad .
\end{equation*}
Comme l'image du morphisme induit sur les fibres spéciales tombe dans la partie infiniment $\ell$-divisible de $A_o$, on conclut que $v(\nu)$ se factorise même par $({}_{\ell^\nu}A^o)^f$. Par passage à la limite sur $\nu$, on trouve un homomorphisme

\begin{equation*}
v: F \longrightarrow T_\ell(A^o)^f \quad ,
\end{equation*}
et le morphisme induit sur les fibres génériques donne évidemment la factorisation cherchée.

\underline{Remarques} 5.7.2. Lorsque $\ell \neq p$, il est clair que la propriété 5.7 du sous-pro-groupe $T_\ell(A_K)^f$ de $T_\ell(A_K)$ caractérise celui-ci de façon unique, comme le plus grand sous-faisceau $\ell$-adique de $T_\ell(A_K)$ qui soit ``non ramifié''. Lorsque $\ell = p$, l'énoncé 5.7 cesse d'être bien raisonnable, et on va donner


\SGApage{I}{382}

une propriété nettement plus forte que celle énoncée dans 5.7, qui pourra alors servir de caractérisation intrinsèque de la partie fixe de $T_p(A_K)$ (en termes de la seule structure de pro-groupe de Barsotti-Tate de $T_p(A_K)$), - du moins dans le cas d'inégales caractéristiques :

\underline{Théorème} 5.8. (M. RAYNAUD). \underline{Soit} $A_K$ \underline{un schéma abélien à réduction semi-stable sur le corps des fonctions} $K$ \underline{d'un trait hensélien} $S$. \underline{Soit} $\ell$ \underline{un} \underline{nombre premier (le cas intéressant étant celui où} $\ell = p$), \underline{et considérons} \underline{le pro-$\ell$-groupe de Barsotti-Tate} $X_\eta = T_\ell(A_K)$ \underline{et sa} \underline{``partie fixe''}$X_\eta^f$ , \underline{qui est la fibre générique du pro-$\ell$-groupe de Barsotti-Tate} $X^f = T_\ell(A^o)^f$. \underline{Supposons, si} $\ell = p$, \underline{qu'on ait} car. $K = 0$\underline{(de sorte que nous disposons} \underline{sur} $S$ \underline{du théorème de TATE} [33, th.4]).\underline{Soit} $F$ \underline{un pro-$\ell$-groupe de} \underline{Barsotti-Tate sur} $S$. \underline{Alors tout homomorphisme} $F_\eta \longrightarrow X_\eta$ \underline{se factorise par} $X_\eta^f$

Passant aux ind-$\ell$-groupes associés aux pro-$\ell$-groupes de Barsotti-Tate envisagés, l'énoncé précédent devient équivalent au suivant, grâce au théorème cité de TATE :

\underline{Corollaire} 5.9. \underline{Si} $F$ \underline{est un} ind-$\ell$-\underline{groupe de Barsotti-Tate,} \underline{alors tout} \underline{homomorphisme} $u_\eta : F_\eta \longrightarrow A_\eta$ \underline{se prolonge (de façon unique) en un homomorphisme} $u: F \longrightarrow A$.

En d'autres termes, la propriété ``néronienne'' caractéristique (1.1.2) du modèle de Néron reste valable pour des homomorphismes de ind-schémas en groupes de Barsotti-Tate, sans hypothèse de lissité sur celui-ci. (Seul le cas des groupes ind-étales résulterait directement de la définition.)


\SGApage{I}{383}

\underline{Remarque} 5.9.1.  En fait, M. RAYNAUD démontre 5.9 sans supposer $A_K$ à réduction semi-stable, et même pour tout schéma en groupes néronien $A$ sur $S$ (dont la fibre générique n'a pas besoin d'être un schéma abélien) ; bien sûr, il est également inutile dans la formulation 5.9 que $S$ soit hensélien. La démonstration de RAYNAUD de ce cas plus général procède par réduction au cas particulier envisagé dans 5.8, en utilisant une propriété inédite dans la construction des modèles de Néron. Nous nous bornerons ici à donner la démonstration de RAYNAUD de 5.9.

\underline{Démonstration de} 5.9. On sait [33] que $F$ est une extension d'un groupe de Barsotti-Tate $F'$ qui est ind-étale par un groupe de Barsotti-Tate $F''$ à fibre spéciale connexe, de sorte que pour tout entier $n$, on a une suite exacte

\SGAvar{\begin{equation*}
0 \longrightarrow F''(n) \longrightarrow F(n) \longrightarrow F'(n) \longrightarrow 0 \qquad ,
\end{equation*}
avec $F'(n)$ connexe et $F''(n)_o$ purement infinitésimal. L'homomorphisme donné $u_\eta$ se définit par une suite d'homomorphismes

}{\begin{equation*}
0 \longrightarrow F''(n) \longrightarrow F(n) \longrightarrow F'(n) \longrightarrow 0 \qquad ,
\end{equation*}
avec $F'(n)$ étale et $F''(n)_o$ purement infinitésimal. L'homomorphisme donné $u_\eta$ se définit par une suite d'homomorphismes

}{F7I-IX-061}\begin{equation*}
u_\eta(n) : F(n)_\eta \longrightarrow A_\eta \qquad ,
\end{equation*}
et tout revient à prolonger séparément ceux-ci en des $u(n) : F(n) \longrightarrow A$. Le lemme 5.9.2 ci-dessous nous ramène alors à prouver que l'homomorphisme $u''_\eta : F''_\eta \longrightarrow A_\eta$ induit par $u_\eta$ se prolonge en $u'' : F'' \longrightarrow A$. Par suite, \underline{on peut supposer} $F$ \underline{à fibre spéciale connexe}.

Nous pouvons considérer $u_\eta$ comme un homomorphisme de $F_\eta$ dans le ind-$\ell$-groupe de Barsotti-Tate $Y_\eta$ associé à $X_\eta = T_\ell(A_K)$, et considérer alors le composé


\SGApage{I}{384}

\begin{equation}\tag{*}
v_\eta : F_\eta \longrightarrow Y_\eta \longrightarrow Y_\eta/Y_\eta^f \qquad .
\end{equation}
En vertu de 5.6 ce dernier groupe se prolonge en un groupe de Barsotti-Tate ind-étale $Z$ sur $S$ (savoir $\underline{M} \otimes_{\mathbb{Z}} (Q_\ell/\mathbb{Z}_\ell)$). En vertu du théorème de TATE cité dans 5.8, l'homomorphisme (*) se prolonge alors en un homomorphisme

\begin{equation*}
v : F \longrightarrow Z \qquad .
\end{equation*}
Comme $F$ est à fibre spéciale connexe et que $Z$ est ind-étale, ce morphisme est \underline{nul}. Donc il en est de même de $v_\eta$, donc $u_\eta$ se factorise par $Y_\eta^f$ . C'est ce qu'il fallait démontrer !

Reste à prouver le

\underline{Lemme} 5.9.2. \underline{Soient} $S$ \underline{un trait}, $A$ \underline{un schéma en groupes néronien sur} $S$ (1.1.2),

\begin{equation*}
1 \longrightarrow G' \longrightarrow G \longrightarrow G'' \longrightarrow 1
\end{equation*}
\underline{une extension de schémas en groupes sur} $S$, \underline{avec} $G''$ \underline{lisse,}

\begin{equation*}
u : G' \longrightarrow A
\end{equation*}
\underline{un homomorphisme de groupes}, \underline{et}

\begin{equation*}
v_\eta : G_\eta \longrightarrow A_\eta
\end{equation*}
\underline{un homomorphisme de groupes prolongeant} $u_\eta$. \underline{Supposons que l'extension de} $G''$ \underline{par} $A$ \underline{image de l'extension} $G$ \underline{par} $u$ \underline{soit représentable} (*), - \underline{ce qui est le cas par exemple si} $G \longrightarrow G''$ \underline{admet une section localement pour la topologie} \underline{``finie localement libre''} (cf. SGA 3 IV 6.3) \underline{et si} $A$ \underline{est quasi-projectif sur} $S$. \underline{Alors il existe un unique morphisme de groupes}

\begin{equation*}
v : G \longrightarrow A
\end{equation*}
\underline{prolongeant} $v_\eta$ \underline{et} $u$.

(*) Des résultats inédits de M. RAYNAUD impliquent d'ailleurs que cette condition est toujours satisfaite.


\SGApage{I}{385}

En effet, la dernière hypothèse faite nous ramène au cas où $G'=A$, $u$ = identité, et où $v_\eta : G_\eta \longrightarrow A_\eta$ est donc une projection, qu'on se propose de prolonger en une projection $v: G \longrightarrow A$. Or $G$ est \underline{lisse} comme extension de deux schémas en groupes lisses, de sorte que la conclusion résulte de la définition de la propriété néronienne de $A$.

\underline{Corollaire} 5.10. \underline{Soient} $A_K$ \underline{un schéma abélien sur le corps des fonctions} $K$ \underline{d'un trait} $S$, $\ell$ \underline{un nombre premier.} \underline{Lorsque} $\ell = p$, \underline{on suppose} car.$K = 0$, \underline{pour pouvoir disposer du théorème de TATE déjà invoqué dans} 5.8. \underline{Pour que} $A_K$ \underline{ait bonne réduction sur} $S$, \underline{il faut et il suffit que le pro-$\ell$-groupe de Barsotti-Tate} $T_\ell(A_K)$ ``\underline{ait bonne réduction sur} $S$'', \underline{i.e.} \underline{se prolonge en un groupe de Barsotti-Tate sur} $S$.

(C'est le résultat conjecturé dans [8, 4.9].)

\SGAvar{La nécessité étant claire, il reste à prouver la suffisance ; notons d'ailleurs que le cas $\ell \neq p$ a déjà été traité (2.2.9). On peut supposer $S$ hensélien (2.2.9). Supposons d'abord que $A_K$ ait réduction semi-stable sur $S$. Alors 5.8 appliqué à $F = X$ nous montre qu'on a $T_\ell(A_K) = T_\ell(A_K)^f$, ce qui signifie aussi que le tore maximal de $A_o$ est nul, i.e. que $A_o^o$ est un schéma abélien, de sorte que la conclusion résulte de (2.2.9 (iii${}^o$)). Traitons maintenant le cas général. En vertu de 2.2.9 on peut supposer évidemment $S$ strictement local. Soir $K'$ une extension galoisienne finie telle que $A_{K'}$ ait réduction semi-stable sur le normalisé $S'$ de $S$ dans $K'$ (3.6). Si $A^*$ désigne le modèle de Néron correspondant, $G = \mathrm{Gal}(K'/K)$ opère sur la situation $(A^*,S')$, donc sur


}{La nécessité étant claire, il reste à prouver la suffisance ; notons d'ailleurs que le cas $\ell \neq p$ a déjà été traité (2.2.9). On peut supposer $S$ hensélien (2.2.9). Supposons d'abord que $A_K$ ait réduction semi-stable sur $S$. Alors 5.8 appliqué à $F = X$ nous montre qu'on a $T_\ell(A_K) = T_\ell(A_K)^f$, ce qui signifie aussi que le tore maximal de $A_o$ est nul, i.e. que $A_o^o$ est un schéma abélien, de sorte que la conclusion résulte de (2.2.9 (iii${}^o$)). Traitons maintenant le cas général. En vertu de 2.2.9 on peut supposer évidemment $S$ strictement local. Soit $K'$ une extension galoisienne finie telle que $A_{K'}$ ait réduction semi-stable sur le normalisé $S'$ de $S$ dans $K'$ (3.6). Si $A^*$ désigne le modèle de Néron correspondant, $G = \mathrm{Gal}(K'/K)$ opère sur la situation $(A^*,S')$, donc sur


}{F7I-IX-062}\SGApage{I}{386}

$A^{*o}_o$, et d'après ce qu'on a vu dans 4.1, il suffit de prouver que $G$ opère trivialement sur ce dernier, ce qui impliquera en effet que $A_K$ a réduction semi-stable, i.e. qu'on est dans le cas déjà traité. Or soit $X$ un groupe de Barsotti-Tate sur $S$ qui prolonge $X_\eta = T_\ell(A_K)$. En vertu du théorème d'unicité de TATE déjà invoqué, il existe un (unique) isomorphisme

\SGAvar{\begin{equation*}
X \times_S S' \quad \simeq \quad T_\ell(A^{*o})
\end{equation*}
qui induit l'identité sur les fibres génériques. Par transport de structure, cet isomorphisme est compatible avec les opérations de $G$. Or il est clair que $X$ opère trivialement sur la fibre spéciale de $X \times_S S'$, donc il opère trivialement sur $T_\ell(A^{*o}_o)$, d'où il résulte évidemment qu'il opère trivialement sur $A^{*o}_o$, ce qui achève la démonstration.

}{\begin{equation*}
X \times_S S' \quad \simeq \quad T_\ell(A^{*o})
\end{equation*}
qui induit l'identité sur les fibres génériques. Par transport de structure, cet isomorphisme est compatible avec les opérations de $G$. Or il est clair que $G$ opère trivialement sur la fibre spéciale de $X \times_S S'$, donc il opère trivialement sur $T_\ell(A^{*o}_o)$, d'où il résulte évidemment qu'il opère trivialement sur $A^{*o}_o$, ce qui achève la démonstration.

}{F7I-IX-063}\underline{Remarque} 5.11. Une fois le théorème de TATE prouvé sans hypothèse restrictive sur car.$K$, les résultats précédents 5.8, 5.9, 5.10, ainsi que 5.13 plus bas, seraient évidemment établis sans une telle restriction.

Posons une définition, qui pourra être encore utile ailleurs :

\underline{Définition} 5.12. \underline{Soient} $S$ \underline{un trait hensélien,} $\ell$ \underline{un nombre premier}, $X_\eta$ \underline{un pro}-$\ell$\underline{-groupe de Barsotti-Tate sur le point générique} $\eta$ \underline{de} $S$. \underline{On dit que} $X$ a bonne réduction sur $S$ \underline{si} $X$ \underline{se prolonge en un groupe de Barsotti-Tate sur} $S$ (cf. 5.10), \underline{qu'il a} bonne réduction virtuelle sur $S$, \underline{s'il existe une filtration} $(\mathrm{Fil}^i(X_\eta))_{i \geq 0}$ \underline{de} $X$ \underline{par des sous-groupes de Barsotti-Tate stricts} (*) \underline{telle que les groupes de Barsotti-Tate} $\mathrm{Gr}^i(X_\eta) = \mathrm{Fil}^i(X_\eta)/\mathrm{Fil}^{i+1}(X_\eta)$ \underline{aient bonne réduction sur} $S$ ; \underline{si on peut choisir la filtration de sorte qu'elle soit de longueur} $\leq n$ (\underline{i.e.} \underline{que} $n$

(*) i.e. qui sont des sous-objets comme systèmes projectifs, i.e. terme à terme.


\SGApage{I}{387}

\underline{au plus parmi les} $\mathrm{Gr}^i$ \underline{sont} $\neq 0$), \underline{on dit que} $X$ \underline{a} bonne réduction virtuelle d'échelon $n$ sur $S$. \underline{Enfin on dit que} $X$ a bonne réduction essentielle sur $S$ (resp. bonne réduction virtuelle essentielle, resp. bonne réduction virtuelle essentielle d'échelon $n$) \underline{s'il existe une extension finie} $K'$ \underline{telle que} $X_{K'}$ \underline{ait bonne réduction} (resp. \underline{bonne réduction virtuelle,} resp. \underline{bonne réduction virtuelle d'échelon} $n$) \underline{sur le normalisé} $S'$ \underline{dans} $K'$.

5.12.1.  Lorsque $\ell \neq p$, de sorte que la catégorie des pro-$\ell$-groupes de Barsotti-Tate sur $\eta$ resp. sur $S$ s'interprète en termes de représentations $\ell$-adiques de $\pi = \pi_1(\eta)$ resp. de $\pi_o = \pi_1(S)$, on voit aussitôt que les six notions introduites s'expriment en termes de la représentation naturelle du groupe d'inertie $I$ sur $X_\eta(\overline{K})$, par la propriété respectivement que celle-ci soit triviale, resp. unipotente, resp. unipotente d'échelon $n$, resp. essentiellement triviale, resp. essentiellement unipotente, resp. essentiellement unipotente d'échelon $n$.

5.12.2.  Supposons qu'on prenne $X_\eta = T_\ell(A_K)$, où $A_K$ est un schéma abélien sur $K$. On a vu dans 5.6 que si $A_K$ a réduction semi-stable sur $S$, alors $X_\eta$ a bonne réduction virtuelle d'échelon 2 (et on verra une réciproque dans un instant) ; donc en vertu du théorème de réduction semi-stable 3.6, $X_\eta$ \underline{a bonne réduction essentielle virtuelle d'échelon} 2. Il convient de regarder cet énoncé comme donnant la ``bonne'' formulation du théorème de monodromie géométrique III\hspace{2em}, dans le cas du $H^1$ d'un schéma propre et lisse sur $K$, \underline{sans exclure le cas} $\ell = p$ (dont il n'avait pas été question dans loc. cit.). Il conviendrait de reformuler également le théorème


\SGApage{I}{388}

(ou plutôt, pour l'instant, la conjecture) de monodromie géométrique pour les $H^i$ supérieurs, donné dans III\hspace{2em}, de façon à tenir compte des coefficients ``$p$-adiques'', comme il a été dit dans 0.4.

Donnons un énoncé récapitulatif des relations entre les propriétés de réduction des schémas abéliens sur $K$, et des groupes de Barsotti-Tate associés :

\underline{Proposition} 5.13.  \underline{Soient} $S$ \underline{un trait hensélien}, $A_K$ \underline{un schéma abélien sur son corps des fractions}, $\ell$ \underline{un nombre premier}. \underline{Lorsque} $\ell = p$, \underline{on suppose que} $K$ \underline{est de car}. \underline{nulle}(cf. 5.11).

\SGAvar{a) \underline{Pour que} $A$ \underline{ait} bonne réduction sur $S$ (2.2.9), \underline{il faut et il suffit que} $T_\ell(A_K)$ \underline{ait bonne réduction sur} $S$ (5.12).

}{a) \underline{Pour que} $A_K$ \underline{ait} bonne réduction sur $S$ (2.2.9), \underline{il faut et il suffit que} $T_\ell(A_K)$ \underline{ait bonne réduction sur} $S$ (5.12).

}{F7I-IX-064}b) \underline{Pour que} $A_K$ \underline{ait} bonne réduction essentielle sur $S$, \underline{i.e. qu'il existe une extension finie} $K'$ \underline{de} $K$ \underline{telle que} $A_{K'}$ \underline{ait bonne réduction sur le normalisé} $S'$ \underline{de} $S$ \underline{dans} $K'$, \underline{il faut et il suffit que} $T_\ell(A_K)$ \underline{ait bonne réduction essentielle sur} $S$ (5.12). \underline{On peut alors prendre pour} $K'$ \underline{une extension finie} séparable \underline{de} $K$.

c) \underline{Les conditions suivantes sont équivalentes} :

\SGAvar{1) $A$ \underline{a} réduction semi-stable sur $S$ (3.4).

}{1) $A_K$ \underline{a} réduction semi-stable sur $S$ (3.4).

}{F7I-IX-065}2) $T_\ell(A_K)$ \underline{a bonne réduction virtuelle d'échelon} 2 \underline{sur} $S$ (5.12).

3) $T_\ell(A_K)$ \underline{a bonne réduction virtuelle sur} $S$.

\underline{Démonstration}. Le cas a) n'est autre que 5.10 et a été mis pour mémoire.

b) La première assertion est une conséquence triviale de a) et des définitions. Pour prouver la dernière assertion dans b), on peut supposer que $\ell \neq p$, mais alors grâce à a) il suffit de prendre


\SGApage{I}{389}

\SGAvar{pour $K'$ une extension finie séparable de $K$ correspondant à un sous-groupe ouvert $\pi'$ de $\pi = \mathrm{Gal}(\overline{K}/\overline{K})$ dont l'intersection avec le groupe d'inertie $I$ est contenue dans le noyau de la représentation de $I$ sur $T_\ell(A_K(\overline{K}))$, lequel est ouvert par hypothèse (cf. 5.12.1).

}{pour $K'$ une extension finie séparable de $K$ correspondant à un sous-groupe ouvert $\pi'$ de $\pi = \mathrm{Gal}(\overline{K}/K)$ dont l'intersection avec le groupe d'inertie $I$ est contenue dans le noyau de la représentation de $I$ sur $T_\ell(A_K(\overline{K}))$, lequel est ouvert par hypothèse (cf. 5.12.1).

}{F7I-IX-066}c) On a déjà signalé que 1)$\Longrightarrow$ 2) (5.12.2), il est trivial que 2) $\Longrightarrow$ 3), reste à prouver que 3) $\Longrightarrow$ 1). Pour $\ell \neq p$, ce n'est autre que le critère galoisien de réduction semi-stable 3.5. Pour $\ell = p$, on peut procéder ainsi. Choisissons un nombre premier $\ell' \neq p$, et appliquons à celui-ci les considérations de 4.1. Supposons $S$ strictement local, ce qui est loisible (3.4.0), prenons une extension galoisienne finie $K'$ de $K$, de groupe $G$, telle que $A_{K'}$ ait réduction semi-stable. Il faut prouver que la représentation de $G$ sur les trois quotients intervenant dans (4.1.1) est triviale, ou ce qui revient au même, que la représentation somme est triviale, ou encore, que la trace de cette dernière est ``triviale'', i.e. constante. Or nous avons vu dans 4.3 a) que cette fonction trace est indépendante du choix de $\ell' \neq p$, et dans 6.4 plus bas nous indiquons une méthode de calcul de cette fonction en termes de $T_p(A_K)$ également. Ce calcul donne trivialement une fonction constante dans le cas où $T_p(A_K)$ a bonne réduction virtuelle, et cela achève la démonstration

\underline{Remarque} 5.13.1. Il semble plausible que si $X_\eta$ est un pro-$\ell$-groupe de Barsotti-Tate sur $\eta$, qui a bonne réduction virtuelle sur $S$, et bonne réduction virtuelle essentielle d'échelon $n$, alors il a bonne réduction virtuelle d'échelon $n$ . (C'est vrai pour $\ell \neq p$.)


\SGApage{I}{390}

Le résultat suivant, conséquence facile de 5.7 et du théorème de réduction semi-stable 3.6, m'a été signalé par J.P. SERRE :

\underline{Proposition} 5.14. (\underline{Critère de bonne réduction essentielle.}) \underline{Soient} $S$ \underline{un trait de corps de fonctions} $K$, $A_K$ \underline{un schéma abélien sur} $K$, $\ell$ \underline{un nombre premier distinct de la caractéristique de} $K$.

a) \underline{Supposons} $\ell \neq p = \mathrm{car}.k.$ \underline{Alors les conditions suivantes sont équivalentes} :

(i) $A_K$ \underline{a} essentiellement bonne réduction \underline{sur} $S$, \underline{i.e.} \underline{il existe une extension finie séparable} $K'$ \underline{de} $K$ \underline{telle que} $A_{K'}$ \underline{ait bonne réduction relativement au normalisé} $S'$ \underline{de} $S$ \underline{dans} $K'$.

(ii) \underline{L'action du groupe d'inertie} $I$ \underline{sur} $U_\ell = T_\ell(A(\overline{K}))$ \underline{est essentiellement triviale}, \underline{i.e.} \underline{est triviale sur un sous-groupe ouvert de} $I$, \underline{i.e.} \underline{se fait par l'intermédiaire d'un groupe quotient fini} $\Gamma$ \underline{de} $I$.

(iii)\underline{L'action du groupe d'inertie} $I$ \underline{sur} $U_\ell \otimes_{\mathbb{Z}_\ell} \mathbb{Q}_\ell$ \underline{est semi-simple}.

\SGAvar{b) \underline{Supposons} $\ell = p$, \underline{alors}, \underline{si} $n = \dim A_K \neq 0$, \underline{l'image de} $I$ \underline{dans le groupe des automorphismes de} $U_\ell$ \underline{n'est pas finie}, \underline{plus précisément son image dans le groupe} $\mathbb{Z}_\ell$ \underline{des automorphismes de} $\stackrel{2n}{\wedge} U_\ell$ \underline{est d'indice fini}. \underline{Pour que} $A_K$ \underline{ait essentiellement bonne réduction}, \underline{il suffit} (\underline{mais il n'est nullement nécessaire}) \underline{que l'action de} $I$ \underline{sur} $U_\ell \otimes_{\mathbb{Z}_\ell} \mathbb{Q}_\ell$ \underline{soit semi-simple}.

}{b) \underline{Supposons} $\ell = p$, \underline{alors}, \underline{si} $n = \dim A_K \neq 0$, \underline{l'image de} $I$ \underline{dans le groupe des automorphismes de} $U_\ell$ \underline{n'est pas finie}, \underline{plus précisément son image dans le groupe} $\mathbb{Z}_\ell^{\times}$ \underline{des automorphismes de} $\stackrel{2n}{\wedge} U_\ell$ \underline{est d'indice fini}. \underline{Pour que} $A_K$ \underline{ait essentiellement bonne réduction}, \underline{il suffit} (\underline{mais il n'est nullement nécessaire}) \underline{que l'action de} $I$ \underline{sur} $U_\ell \otimes_{\mathbb{Z}_\ell} \mathbb{Q}_\ell$ \underline{soit semi-simple}.

}{F7I-IX-067}\underline{Démonstration}. Dans le cas $\ell \neq p$, l'équivalence de (i) et (ii) est une conséquence triviale du critère de bonne réduction 2.2.9, et (ii) implique (iii) puisque $\mathbb{Q}_\ell$ est de caractéristique nulle. L'homomorphisme $\pi = \mathrm{Gal}(\overline{K}/K) \longrightarrow \mathbb{Z}^*_\ell \simeq \mathrm{Aut}(\stackrel{2n}{\wedge} U_\ell)$ est égal à $\chi^n$, où $\chi$ est induit par


\SGApage{I}{391}

le caractère

\SGAvar{\begin{equation*}
\mathbb{Z}_\ell^{*} = \mathrm{Aut}(\mathbb{Z}_\ell(1))
\end{equation*}
provenant de l'action de $I$ sur les racines $\ell^\nu$-èmes de l'unité, et on sait que si $\ell = p$, l'image de $I$ par ce caractère est bien d'indice fini [33] ; le fait que l'action de $I$ sur $U_p \otimes \mathbb{Q}_p$ puisse ne pas être semi-simple, même si $A_K$ est une courbe elliptique avec bonne réduction, se voit dans [28, IV A 2.4]. Il reste à prouver que (quel que soit $\ell$) (iii) implique (i). Pour ceci, nous utilisons le théorème de réduction semi-stable 3.6, qui nous permet de nous ramener aussitôt au cas où $A_K$ a une réduction semi-stable sur $S$ (puisque l'hypothèse (iii) est conservée en remplaçant $K$ par une extension séparable finie $K'$, i.e. $I$ par un sous-groupe ouvert $I'$ (\quad\quad ). Lorsque $\ell \neq p$, l'action de $I$ sur $U_\ell \otimes \mathbb{Q}_\ell$ est unipotente (3.5 (v)) et semi-simple, donc triviale, et on gagne par le critère de bonne réduction 2.2.9. Donnons une démonstration valable sans restriction sur $\ell$, et qui donne un résultat légèrement plus précis :

}{\begin{equation*}
\mathbb{Z}_\ell^{*} = \mathrm{Aut}(\mathbb{Z}_\ell(1))
\end{equation*}
provenant de l'action de $I$ sur les racines $\ell^\nu$-èmes de l'unité, et on sait que si $\ell = p$, l'image de $I$ par ce caractère est bien d'indice fini [33] ; le fait que l'action de $I$ sur $U_p \otimes \mathbb{Q}_p$ puisse ne pas être semi-simple, même si $A_K$ est une courbe elliptique avec bonne réduction, se voit dans [28, IV A 2.4]. Il reste à prouver que (quel que soit $\ell$) (iii) implique (i). Pour ceci, nous utilisons le théorème de réduction semi-stable 3.6, qui nous permet de nous ramener aussitôt au cas où $A_K$ a une réduction semi-stable sur $S$ (puisque l'hypothèse (iii) est conservée en remplaçant $K$ par une extension séparable finie $K'$, i.e. $I$ par un sous-groupe ouvert $I'$ (\quad\quad )). Lorsque $\ell \neq p$, l'action de $I$ sur $U_\ell \otimes \mathbb{Q}_\ell$ est unipotente (3.5 (v)) et semi-simple, donc triviale, et on gagne par le critère de bonne réduction 2.2.9. Donnons une démonstration valable sans restriction sur $\ell$, et qui donne un résultat légèrement plus précis :

}{F7I-IX-068}\underline{Corollaire} 5.15. \underline{Les notations étant celles de} 5.14, \underline{supposons que} $A_K$ \underline{ait une réduction semi-stable sur} $S$, \underline{et que la partie fixe} $(U_{\ell\mathbb{Q}_\ell})^f \stackrel{\mathrm{dfn}}{=} (U_\ell^f)_{\mathbb{Q}_\ell}$ \underline{admette un supplémentaire dans} $U_{\ell\mathbb{Q}_\ell}$, \underline{stable sous} $I$. \underline{Alors} $A_K$ \underline{a bonne réduction.}

Il suffit de prouver que sous ces conditions, on a $U_\ell = U_\ell^f$, ce qui par 5.6 signifie en effet que la partie torique $T_o$ de $A_o^o$ est nulle, donc que $A_o^o$ est un schéma abélien, ce qui implique bien que $A$ est un schéma abélien (2.2.9). On peut évidemment supposer $k$ séparablement clos. Alors l'intersection du supplémentaire postulé dans 5.15 avec $U_\ell$


\SGApage{I}{392}

définit un sous-pro-$\ell$-groupe $F_K$ de $T_\ell(A_K)$, isogène à la partie cotorique $T_\ell(A_K)/T_\ell(A_K)^f$. Comme celle-ci est proétale non ramifiée en vertu de 5.6, il en est de même de $F$, et on conclut par 5.7 que $F_K$ est contenu dans $T_\ell(A_K)^f$, donc que $F_K = 0$ ce qui prouve bien que $U_\ell$ est égal à sa partie fixe.


\SGAsectiontoc{6. Remarques sur la construction de la partie torique et la partie fixe de $T_p(A)$ dans le cas de réduction non semi-stable.}{6}
6. \underline{Remarques sur la construction de la partie torique et la partie fixe de} $T_p(A)$ \underline{dans le cas de réduction non semi-stable.}

\SGAvar{6.1. La construction de la partie torique de $T_\ell(A_K)$ donnée dans 5.1 garde un sens dès que $S$ est complet, sans supposer $A^o$ $\ell$-divisible i.e. (dans le cas $p = \ell$) sans supposer qu'on soit dans le cas de la réduction semi-stable (hypothèse qui avait servi seulement pour pouvoir formuler le théorème d'orthogonalité 5.2). On peut d'ailleurs construire $T_p(A_K)^t$ en supposant seulement $S$ hensélien, en montrant que la partie torique $T_p(A_{\hat{K}})^t$ sur le complété, construit par la méthode de 5.1, provient d'un sous-pro-groupe de Barsotti-Tate de $T_p(A_K)$ (évidemment déterminé de façon unique). Plus précisément, on arrive à construire un $p$-pro-groupe de Barsotti-Tate $T_p(A)^t$ \underline{sur} $S$ et un homomorphisme $T_p(A)^t \longrightarrow T_p(A)$ de systèmes projectifs, qui soit un monomorphisme terme à terme, et qui par image inverse sur $\hat{S}$ donne les objets analogues construits par la méthode de 5.1. Ceci se voit en paraphrasant la construction de 5.1, mais au lieu de relever le tore maximal de $T_o$ de $A_o$ tout entier, on se borne, pour tout entier $\nu \geq 0$, à relever le sous-groupe ${}_{p^\nu}T_o$ en un


}{6.1. La construction de la partie torique de $T_\ell(A_K)$ donnée dans 5.1 garde un sens dès que $S$ est complet, sans supposer $A^o$ $\ell$-divisible i.e. (dans le cas $p = \ell$) sans supposer qu'on soit dans le cas de la réduction semi-stable (hypothèse qui avait servi seulement pour pouvoir formuler le théorème d'orthogonalité 5.2). On peut d'ailleurs construire $T_p(A_K)^t$ en supposant seulement $S$ hensélien, en montrant que la partie torique $T_p(A_{\hat{K}})^t$ sur le complété, construit par la méthode de 5.1, provient d'un sous-pro-groupe de Barsotti-Tate de $T_p(A_K)$ (évidemment déterminé de façon unique). Plus précisément, on arrive à construire un $p$-pro-groupe de Barsotti-Tate $T_p(A)^t$ \underline{sur} $S$ et un homomorphisme $T_p(A)^t \longrightarrow T_p(A)$ de systèmes projectifs, qui soit un monomorphisme terme à terme, et qui par image inverse sur $\hat{S}$ donne les objets analogues construits par la méthode de 5.1. Ceci se voit en paraphrasant la construction de 5.1, mais au lieu de relever le tore maximal $T_o$ de $A_o$ tout entier, on se borne, pour tout entier $\nu \geq 0$, à relever le sous-groupe ${}_{p^\nu}T_o$ en un


}{F7I-IX-069}\SGApage{I}{393}

sous-groupe de type multiplicatif $({}_{p^\nu}A^o)^t$ de ${}_{p^\nu}A^o$. Pour l'existence et l'unicité d'un tel relèvement, on note que le foncteur $(\mathrm{Sch})^o_{/S} \longrightarrow \mathrm{Ens}$, $S' \longmapsto$ ensemble des sous-groupes finis de type multiplicatif de $A^o_{S'}$ , est représentable par un schéma $X$ de type fini sur $S$ (cf. démonstration de SGA 3 XI 3.12 a)) qui est étale sur $S$ en vertu de SGA 3 IX 3.6 bis ; il y a par suite une section et une seule de $X$ sur $S$ qui prolonge une section donnée de $X_s$ sur $s$, ce qui prouve notre assertion. Ce raisonnement montre d'ailleurs, plus généralement, que pour tout sous-groupe \underline{fini et de type multiplicatif} $H_o$ de $A_o$, il y a un unique sous-groupe fini et de type multiplicatif $H$ de $A$ qui le prolonge ; par suite, pour tout sous-$p$-groupe de Barsotti-Tate $H_o$ de $A_o$ qui soit \underline{toroïdal} (i.e. dont les composantes sont de type multiplicatif), il existe un unique sous-$p$-groupe de Barsotti-Tate $H$ de $A$ qui le prolonge, et qui fournit par suite un sous-$p$-groupe de Barsotti-Tate de ${}_{p^\infty}A_K$, ou, si on préfère, un sous-$p$-pro-groupe de Barsotti-Tate de $T_p(A_K)$. Ce procédé s'applique notamment au cas où $H_o$ est le plus grand sous-$p$-groupe de Barsotti-Tate toroïdal de $A_o$, lequel s'obtient en prenant le plus grand sous-groupe $p$-divisible $\overline{A}_o$ de $A_o^o$($= p^N A_o^o$ pour $N$ grand), et la ``partie toroïdale'' de ${}_{p^\infty}\overline{A}_o$ (définie par exemple grâce à SGA 3 XVII 7.2.1). On obtient ainsi un sous-groupe de Barsotti-Tate de ${}_{p^\infty}A_K$, qu'on pourrait appeler la \underline{partie toroïdale} de ce dernier.

6.2. Reprenons le plus grand sous-schéma en groupes $p$-divisible $\overline{A}_o$ de $A_o^o$, et considérons le $p$-groupe de Barsotti-Tate $F_o = {}_{p^\infty}(\overline{A}_o)$ ; il est immédiat qu'il peut être caractérisé comme le plus grand sous-$p$-groupe de Barsotti-Tate de $A_o$ (ou de $A_o^o$). Il est naturel de se poser la question


\SGApage{I}{394}

\SGAvar{de prolonger ce groupe de Barsotti-Tate en un sous-groupe de Barsotti-Tate $F$ de $A$, qui jouerait alors le rôle joué par la \underline{partie fixe} $T_\ell(A)^f$ de (2.2.3.2) dans le cas où $A^o$ est $\ell$-divisible. Dans le cas semi-stable, où on a $\overline{A}_o = A_o^o$,la question se résoud par l'affirmative par la méthode de 2.2. Dans le cas général, lorsque le problème envisagé a une solution, on voit facilement qu'elle est unique, et que les sous-pro-groupes correspondants de $T_p(A_K)$ et de $T_p(A'_K)$ (qui méritent alors encore le nom de \underline{partie fixe} des pro-groupes envisagés) satisfont encore au théorème d'orthogonalité 5.2, enfin que la partie fixe $T_p(A_K)^f$ peut encore se caractériser comme indiqué dans 5.8.

}{de prolonger ce groupe de Barsotti-Tate en un sous-groupe de Barsotti-Tate $F$ de $A$, qui jouerait alors le rôle joué par la \underline{partie fixe} $T_\ell(A)^f$ de (2.2.3.2) dans le cas où $A^o$ est $\ell$-divisible. Dans le cas semi-stable, où on a $\overline{A}_o = A_o^o$,la question se résout par l'affirmative par la méthode de 2.2. Dans le cas général, lorsque le problème envisagé a une solution, on voit facilement qu'elle est unique, et que les sous-pro-groupes correspondants de $T_p(A_K)$ et de $T_p(A'_K)$ (qui méritent alors encore le nom de \underline{partie fixe} des pro-groupes envisagés) satisfont encore au théorème d'orthogonalité 5.2, enfin que la partie fixe $T_p(A_K)^f$ peut encore se caractériser comme indiqué dans 5.8.

}{F7I-IX-070}6.2.1. Montrons cependant qu'en général le problème envisagé n'a pas une solution affirmative, i.e. qu'\underline{il n'est pas possible en général de définir la partie fixe} $T_p(A)^f$ (cas de la réduction non semi-stable). En effet, en vertu d'un théorème de relèvement de SERRE-TATE [30], lorsque $S$ est complet, l'existence de la partie fixe équivaut à celle de relèvements infinitésimaux $\overline{A}_n$ de $\overline{A}_o$ dans les $A_n$, ou encore à celle d'un schéma en groupes lisse $\overline{A}$ sur $S$, extension d'un schéma abélien par un tore $T$, de fibre spéciale $\overline{A}_o$, et d'un homomorphisme formel $\widehat{\overline{A}} \longrightarrow \widehat{A}$ prolongeant l'inclusion sur les fibres spéciales. Donc, dans le cas où $\overline{A}_o$ est un schéma abélien i.e. où $T_o = 0$, cela équivaut à l'existence d'un schéma abélien $\overline{A}$ prolongeant $\overline{A}_o$ et d'un homomorphisme $\overline{A} \longrightarrow A$ prolongeant l'inclusion des fibres spéciales. Il suffira donc de donner un cas où $T_o = 0$, où $A_o^o$ est non unipotent et non propre, et où $A_K$ est simple sur $K$.


\SGApage{I}{395}

\SGAvar{6.2.2. Voici un exemple, dû à M. RAYNAUD, basé sur l'observation que si $S$ est un trait, $S'$ le normalisé de $S$ dans une extension finie séparable $K'$ de $K$, et $A'$ un schéma abélien sur $S'$, alors le schéma en groupes $A = \prod_{S'/S} A'/S'$ sur $S$ est néronien, i.e. s'identifie au modèle de Néron de sa fibre générique (qui est bien un schéma abélien). La vérification de ce fait est immédiate en termes des définitions et du fait que $A'$ lui-même est néronien sur $S'$. Alors $A_o = \prod_{S'_s/s}(A' \times_{S'} S'_s)/S'_s$ est une extension d'un schéma abélien, non nul si $A' \neq 0$, par un schéma unipotent connexe, non nul si $A' \neq 0$ et si $S'$ est ramifié sur $S$. Il suffira donc de donner un exemple où $A' \neq 0$, $S'$ ramifié sur $S$ (trait local complet), et $A_K = \prod_{K'/K} A'_{K'}/K'$ est simple sur $K'$. Or prenons pour $K'$ une extension quadratique ramifiée de $S$ (il en existe toujours), pour $A'$ une courbe elliptique sur $S'$. Alors $A$ est de dimension relative 2, et dire que $A_K$ est simple sur $K$ signifie que l'on ne peut trouver un homomorphisme non nul $B_K \longrightarrow A_K$, où $B_K$ est un schéma abélien de dimension 1 sur $K$. Or la donnée d'un tel homomorphisme revient à la donnée d'un homomorphisme non nul $B_{K'} \longrightarrow A'_{K'}$ , qui sera donc nécessairement une isogénie, i.e, la condition que $A_K$ soit simple sur $K$ signifie que $A'_{K'}$ ne soit pas isogène à une courbe elliptique provenant de $K$. Mais dans le cas contraire, l'invariant $j$ de $A'$ serait algébriquement dépendant (sur le corps premier) du conjugué $\bar{j}$ de $j$ par le $K$-automorphisme non trivial de $K'$, et si $j$ est choisi tel que $a = j + \bar{j} = \mathrm{Tr}_{K'/K} j$ soit transcendant sur le corps premier, il revient au même de dire que $j$ est algébriquement dépendant de $a$. Or, pour des raisons de cardinalité, il existe un élément


}{6.2.2. Voici un exemple, dû à M. RAYNAUD, basé sur l'observation que si $S$ est un trait, $S'$ le normalisé de $S$ dans une extension finie séparable $K'$ de $K$, et $A'$ un schéma abélien sur $S'$, alors le schéma en groupes $A = \prod_{S'/S} A'/S'$ sur $S$ est néronien, i.e. s'identifie au modèle de Néron de sa fibre générique (qui est bien un schéma abélien). La vérification de ce fait est immédiate en termes des définitions et du fait que $A'$ lui-même est néronien sur $S'$. Alors $A_o = \prod_{S'_s/s}(A' \times_{S'} S'_s)/S'_s$ est une extension d'un schéma abélien, non nul si $A' \neq 0$, par un schéma unipotent connexe, non nul si $A' \neq 0$ et si $S'$ est ramifié sur $S$. Il suffira donc de donner un exemple où $A' \neq 0$, $S'$ ramifié sur $S$ (trait local complet), et $A_K = \prod_{K'/K} A'_{K'}/K'$ est simple sur $K$. Or prenons pour $K'$ une extension quadratique ramifiée de $K$ (il en existe toujours), pour $A'$ une courbe elliptique sur $S'$. Alors $A$ est de dimension relative 2, et dire que $A_K$ est simple sur $K$ signifie que l'on ne peut trouver un homomorphisme non nul $B_K \longrightarrow A_K$, où $B_K$ est un schéma abélien de dimension 1 sur $K$. Or la donnée d'un tel homomorphisme revient à la donnée d'un homomorphisme non nul $B_{K'} \longrightarrow A'_{K'}$ , qui sera donc nécessairement une isogénie, i.e, la condition que $A_K$ soit simple sur $K$ signifie que $A'_{K'}$ ne soit pas isogène à une courbe elliptique provenant de $K$. Mais dans le cas contraire, l'invariant $j$ de $A'$ serait algébriquement dépendant (sur le corps premier) du conjugué $\bar{j}$ de $j$ par le $K$-automorphisme non trivial de $K'$, et si $j$ est choisi tel que $a = j + \bar{j} = \mathrm{Tr}_{K'/K} j$ soit transcendant sur le corps premier, il revient au même de dire que $j$ est algébriquement dépendant de $a$. Or, pour des raisons de cardinalité, il existe un élément


}{F7I-IX-071}\SGApage{I}{396}

$a \in V$ qui soit transcendant sur le sous-corps premier $P$ de $K$, puis un élément $b$ de $V'$ (normalisé de $V$ dans $K'$) qui soit transcendant sur $P(a)$ et tel que $\mathrm{Tr}_{V'/V} b = 0$, et il suffit de prendre $j = a + b \in V'$. On peut d'ailleurs supposer que $b \in \underline{r}(V')$, et que l'image de $a$ dans $k = V/\underline{m}$ est un élément donné de $k$, qu'on choisira distinct des valeurs exceptionnelles $j_o = 0$ et $j_o = 12^3$ (NB c'est possible même si $k$ n'a que deux éléments, car alors les deux valeurs envisagées de $j_o$ coïncident). Donc l'image $j_o$ de $j$ dans $k(s')$ est $\neq 0,12^3$, et on sait alors qu'il existe bien une courbe elliptique sur $S'$ d'invariant $j$ .

6.3. Revenons au cas où on suppose seulement $S$ hensélien sans plus. En analogie avec des définitions de 4.1 dans le cas $\ell \neq p$, on peut définir alors une filtration de $T_p(A_K)$,

\SGAvar{\begin{equation}\tag{6.3.1}
T_p(A_K) \supset T_p(A_K)^{ef} \supset T_p(A_K)^{et} \supset 0
\end{equation}
par des sous-pro-groupes (stricts) de Barsotit-Tate, appelés respectivement la partie \underline{essentiellement fixe} et la partie \underline{essentiellement torique} de $T_p(A_K)$. Pour ceci, on choisit une extension galoisienne finie $K'$ de $K$ telle que $A_{K'}$ ait une réduction semi-stable sur le normalisé $S'$ de $S$ dans $K'$, de sorte que la partie fixe $T_p(A_{K'})^f$ de $T_p(A_{K'})=(T_p(A_K))_{K'}$ est définie (2.2.3), ainsi que la partie torique (6.1), qui n'est d'ailleurs autre que l'orthogonale de la partie fixe $T_p(A'_K)^f$ de $T_p(A'_K)$ (5.2). Il est clair que l'on a $T_p(A_{K'})^t \subset T_p(A_{K'})^f$ (résulte de 5.3, ou directement de la construction 6.1 de la partie torique), et que le groupe de Galois de $K'/K$ laisse invariants ces deux sous-pro-groupes stricts de $(T_p(A_K))_{K'}$ .


}{\begin{equation}\tag{6.3.1}
T_p(A_K) \supset T_p(A_K)^{ef} \supset T_p(A_K)^{et} \supset 0
\end{equation}
par des sous-pro-groupes (stricts) de Barsotti-Tate, appelés respectivement la partie \underline{essentiellement fixe} et la partie \underline{essentiellement torique} de $T_p(A_K)$. Pour ceci, on choisit une extension galoisienne finie $K'$ de $K$ telle que $A_{K'}$ ait une réduction semi-stable sur le normalisé $S'$ de $S$ dans $K'$, de sorte que la partie fixe $T_p(A_{K'})^f$ de $T_p(A_{K'})=(T_p(A_K))_{K'}$ est définie (2.2.3), ainsi que la partie torique (6.1), qui n'est d'ailleurs autre que l'orthogonale de la partie fixe $T_p(A'_{K'})^f$ de $T_p(A'_{K'})$ (5.2). Il est clair que l'on a $T_p(A_{K'})^t \subset T_p(A_{K'})^f$ (résulte de 5.3, ou directement de la construction 6.1 de la partie torique), et que le groupe de Galois de $K'/K$ laisse invariants ces deux sous-pro-groupes stricts de $(T_p(A_K))_{K'}$ .


}{F7I-IX-072}\SGApage{I}{397}

Donc par descente on en déduit la filtration (6.3.1) voulue. Elle ne dépend pas du choix de $K'$, comme il résulte de la propriété 3.3 de la réduction semi-stable. De plus, on aura évidemment par descente la propriété d'orthogonalité

\SGAvar{\begin{equation}\tag{6.3.2}
T_p(A_K)^{et} = (T_p(A'_k)^{ef})^{\perp}, \; T_p(A_K)^{ef} = (T_p(A'_k)^{et})^{\perp} \qquad .
\end{equation}
Il serait possible de donner une caractérisation intrinsèque de la filtration (6.3.1), en termes du pro-groupe de Barsotti-Tate $X = T_p(A_K)$ sur le corps des fractions $K$ de $S$, dans l'esprit de 5.8, - du moins, pour l'instant, dans le cas d'inégales caractéristiques : la partie essentiellement fixe est le plus grand sous-pro-groupe de Barsotti-Tate qui ait ``essentiellement bonne réduction sur $S$'', i.e. qui, après extension finie séparable $K'$ de $K$, donne un pro-groupe de Barsotti-Tate $X_{K'}$ ayant bonne réduction sur le normalisé $S'$ de $S$ dans $K$, i.e. se prolonge en un pro-groupe de Barsotti-Tate sur $S'$ (nécessairement unique à isomorphisme unique près, par le théorème de TATE [33, th. 4] déjà cité). La partie essentiellement torique se définit alors en termes de la partie essentiellement fixe par orthogonalité (6.3.2).

Signalons aussi la conséquence immédiate suivante de 5.8 et des définitions,qui pour $\ell = p$ se réduit d'ailleurs à 5.12 a) :

\underline{Proposition} 6.3.3. \underline{Soit} $\ell$ \underline{un nombre premier}. \underline{Si} $\ell = p$, \underline{on suppose} car.$K = 0$ (*). \underline{Pour que} $A_K$ \underline{ait essentiellement bonne réduction sur} $S$ (5.12 a)), \underline{il faut et il suffit que la partie essentiellement fixe} $T_\ell(A_K)^{ef}$ (4.1.1 \underline{et} 6.3.1) \underline{soit égal à} $T_\ell(A_K)$ \underline{tout entier}.

}{\begin{equation}\tag{6.3.2}
T_p(A_K)^{et} = (T_p(A'_K)^{ef})^{\perp}, \; T_p(A_K)^{ef} = (T_p(A'_K)^{et})^{\perp} \qquad .
\end{equation}
Il serait possible de donner une caractérisation intrinsèque de la filtration (6.3.1), en termes du pro-groupe de Barsotti-Tate $X = T_p(A_K)$ sur le corps des fractions $K$ de $S$, dans l'esprit de 5.8, - du moins, pour l'instant, dans le cas d'inégales caractéristiques : la partie essentiellement fixe est le plus grand sous-pro-groupe de Barsotti-Tate qui ait ``essentiellement bonne réduction sur $S$'', i.e. qui, après extension finie séparable $K'$ de $K$, donne un pro-groupe de Barsotti-Tate $X_{K'}$ ayant bonne réduction sur le normalisé $S'$ de $S$ dans $K'$, i.e. se prolonge en un pro-groupe de Barsotti-Tate sur $S'$ (nécessairement unique à isomorphisme unique près, par le théorème de TATE [33, th. 4] déjà cité). La partie essentiellement torique se définit alors en termes de la partie essentiellement fixe par orthogonalité (6.3.2).

Signalons aussi la conséquence immédiate suivante de 5.8 et des définitions,qui pour $\ell = p$ se réduit d'ailleurs à 5.13 b) :

\underline{Proposition} 6.3.3. \underline{Soit} $\ell$ \underline{un nombre premier}. \underline{Si} $\ell = p$, \underline{on suppose} car.$K = 0$ (*). \underline{Pour que} $A_K$ \underline{ait essentiellement bonne réduction sur} $S$ (5.13 b)), \underline{il faut et il suffit que la partie essentiellement fixe} $T_\ell(A_K)^{ef}$ (4.1.1 \underline{et} 6.3.1) \underline{soit égale à} $T_\ell(A_K)$ \underline{tout entier}.

}{F7I-IX-073}(*) Hypothèse sans doute inutile, cf. 5.11.


\SGApage{I}{398}

\SGAvar{6.4. Notons que lorsque $\ell \neq p$, nous avons mis en évidence dans 4.3 a) une fonction à valeurs entières remarquable sur le groupe d'inertie $I$ de $S$, d'ailleurs indépendante du choix de $\ell$, et dont la connaissance, dans le cas d'un corps résiduel parfait, suffit à reconstruire l'exposant conducteur de $A$ (4.5), en prenant le produit scalaire, sur un quotient assez grand $\Gamma = I/I'$ de $I$, de la fonction précédente et du conducteur d'Artin sur $\Gamma$ (cf. (4.6.2)). Il peut être intéressant d'observer que cette fonction à valeurs entières sur $I$ peut se réconstituer également par la connaissance de $T_\ell(A_K)$ pour $\ell = p$ - du moins (pour le moment) en inégales caractéristiques. Lorsque le corps résiduel $k$ est fini, une observation analogue s'applique à la fonction à valeurs dans $\mathbb{Q}$, définie sur le sous-groupe $\pi'$ de $\pi$ formé des éléments dont l'image dans $\pi_o$ est de la forme $f_q^{\,n}$ $(n \in \mathbb{Z}\,)$, considérée dans 4.3 b). Indiquons le principe de la vérification de ces assertions. On utilise la filtration (6.3.1), et on est réduit à attacher une fonction convenable (sur $I$ resp. sur $\pi'$) à un $p$-pro-groupe de Barsoti-Tate $Y$ sur $K$ qui a ``essentiellement bonne réduction sur $S$ '', au sens précisé dans 6.5. ($Y$ sera un des trois quotients provenant de la filtration (6.3.1).) Dans le cas de 4.3 a), on se ramène au cas d'un corps résiduel parfait (ou même algébriquement clos, si on veut) par extension résiduelle convenable dans $S$ (EGA $0_{\mathrm{III}}$ 10.3.1). Prenons alors une extension galoisienne finie $K'$ de $K$ telle que $Y_{K'}$ ait bonne réduction sur $S'$, d'où un groupe de Barsotti-Tate $Y_{S'}$ sur $S'$ et une opération de $G$ sur $Y_{S'}$ compatible avec son action sur $S'$. Si $k'$ est le corps résiduel de $S'$, $G$


}{6.4. Notons que lorsque $\ell \neq p$, nous avons mis en évidence dans 4.3 a) une fonction à valeurs entières remarquable sur le groupe d'inertie $I$ de $S$, d'ailleurs indépendante du choix de $\ell$, et dont la connaissance, dans le cas d'un corps résiduel parfait, suffit à reconstruire l'exposant conducteur de $A$ (4.5), en prenant le produit scalaire, sur un quotient assez grand $\Gamma = I/I'$ de $I$, de la fonction précédente et du conducteur d'Artin sur $\Gamma$ (cf. (4.6.2)). Il peut être intéressant d'observer que cette fonction à valeurs entières sur $I$ peut se réconstituer également par la connaissance de $T_\ell(A_K)$ pour $\ell = p$ - du moins (pour le moment) en inégales caractéristiques. Lorsque le corps résiduel $k$ est fini, une observation analogue s'applique à la fonction à valeurs dans $\mathbb{Q}$, définie sur le sous-groupe $\pi'$ de $\pi$ formé des éléments dont l'image dans $\pi_o$ est de la forme $f_q^{\,n}$ $(n \in \mathbb{Z}\,)$, considérée dans 4.3 b). Indiquons le principe de la vérification de ces assertions. On utilise la filtration (6.3.1), et on est réduit à attacher une fonction convenable (sur $I$ resp. sur $\pi'$) à un $p$-pro-groupe de Barsotti-Tate $Y$ sur $K$ qui a ``essentiellement bonne réduction sur $S$ '', au sens précisé dans 5.12. ($Y$ sera un des trois quotients provenant de la filtration (6.3.1).) Dans le cas de 4.3 a), on se ramène au cas d'un corps résiduel parfait (ou même algébriquement clos, si on veut) par extension résiduelle convenable dans $S$ (EGA $0_{\mathrm{III}}$ 10.3.1). Prenons alors une extension galoisienne finie $K'$ de $K$ telle que $Y_{K'}$ ait bonne réduction sur $S'$, d'où un groupe de Barsotti-Tate $Y_{S'}$ sur $S'$ et une opération de $G$ sur $Y_{S'}$ compatible avec son action sur $S'$. Si $k'$ est le corps résiduel de $S'$, $G$


}{F7I-IX-074}\SGApage{I}{399}

\SGAvar{opère (de façon non nécessairement fidèle) sur $k'$ (le corps des invariants étant $k$), et sur $Y_{k'}$ da façon compatible avec son action sur $k'$. Or, d'après la théorie de DIEUDONNÉ [16 bis], au $p$-pro-groupe de Barsotti-Tate $Y_{k'}$ est canoniquement associé  un module $M$ libre de type fini (de rang égal au rang de $Y_{k'}$) sur l'anneau $W(k')$ des vecteurs de Witt sur $k'$, et $G$ opère évidemment sur $M$ de façon compatible avec ses opérations sur $W(k')$ (via ses opérations sur $k'$). On en déduit une opération $W(k')$-linéaire du sous-groupe d'inertie $\Gamma$ de $G$ sur $M$ (qui est un quotient du groupe d'inertie $I$), d'où une fonction trace sur $\Gamma$ donc sur $I$, qui est la fonction cherchée.Elle est a priori à valeurs dans $W(k')$ (*), mais coïncide avec la fonction analogue définie en termes des $T_\ell$ pour $\ell \neq p$, lorsque l'opération de $G$ sur $Y_{S'}$ s'obtient à partir d'une opération de $G$ sur un schéma en groupes lisse et $p$-divisible $C'$ sur $S'$, en prenant $Y = T_p(C')$ (cf. [25 bis, th.7 p.737]) ; elle est donc dans ce cas à valeurs entières. On procède de façon analogue dans le cas de 4.3 b) pour définir une fonction trace sur $\pi'$, en utilisant le procédé de la démonstration de 4.3.2 pour ``linéariser'' l'opération dans $M$ définie par un élément $g \in \pi$ qui se trouve au-dessus d'un $f_q^n$ avec $n \geqslant 0$.


}{opère (de façon non nécessairement fidèle) sur $k'$ (le corps des invariants étant $k$), et sur $Y_{k'}$ de façon compatible avec son action sur $k'$. Or, d'après la théorie de DIEUDONNÉ [16 bis], au $p$-pro-groupe de Barsotti-Tate $Y_{k'}$ est canoniquement associé  un module $M$ libre de type fini (de rang égal au rang de $Y_{k'}$) sur l'anneau $W(k')$ des vecteurs de Witt sur $k'$, et $G$ opère évidemment sur $M$ de façon compatible avec ses opérations sur $W(k')$ (via ses opérations sur $k'$). On en déduit une opération $W(k')$-linéaire du sous-groupe d'inertie $\Gamma$ de $G$ sur $M$ (qui est un quotient du groupe d'inertie $I$), d'où une fonction trace sur $\Gamma$ donc sur $I$, qui est la fonction cherchée.Elle est a priori à valeurs dans $W(k')$ (*), mais coïncide avec la fonction analogue définie en termes des $T_\ell$ pour $\ell \neq p$, lorsque l'opération de $G$ sur $Y_{S'}$ s'obtient à partir d'une opération de $G$ sur un schéma en groupes lisse et $p$-divisible $C'$ sur $S'$, en prenant $Y = T_p(C')$ (cf. [25 bis, th.7 p.737]) ; elle est donc dans ce cas à valeurs entières. On procède de façon analogue dans le cas de 4.3 b) pour définir une fonction trace sur $\pi'$, en utilisant le procédé de la démonstration de 4.3.2 pour ``linéariser'' l'opération dans $M$ définie par un élément $g \in \pi$ qui se trouve au-dessus d'un $f_q^n$ avec $n \geqslant 0$.


}{F7I-IX-075}\SGAsectiontoc{7. L'extension de Raynaud $G^\natural$ sur $S$ attachée au schéma abélien $A_K$ à réduction semi-stable. Dualité des schémas abéliens $B,B'$ sur $S$.}{7}
7. \underline{L'extension de Raynaud} $G^\natural$ \underline{sur} $S$ \underline{attachée au schéma abélien} $A_K$ \underline{à réduction semi-stable.} \underline{Dualité des schémas abéliens} $B,B'$ \underline{sur} $S$.

7.1.  Soit $S$ un schéma noethérien spectre d'un anneau $V$, séparé et complet pour la topologie définie par un idéal $\underline{m}$ de $V$. Soit $A$ un schéma en groupes lisse sur $S$, $A^o$ sa composante neutre, et supposons que $A_o/A_o^{\,o}$ soit fini (donc fini étale) sur $S_o$, et que $A_o^{\,o}$ soit une extension d'un schéma abélien

(*) En fait, il est immédiat que cette fonction est toujours à valeurs dans $W(k)$, étant invariante par $\mathrm{Gal}(k'/k)$.


\SGApage{I}{400}

par un tore. Nous avons posé comme d'habitude $A_n = A\times_S S_n$, où $S_n = \mathrm{Spec}(V/\underline{m}^{n+1})$. Nous nous intéresserons surtout au cas où $S$ est un trait, et où la fibre générique $A_K$ est un schéma abélien à réduction semi-stable sur $S$, $A$ étant le modèle de Néron de $A_K$ .

Soit $T_o$ le tore maximal de $A_o$, de sorte que $B_o = A_o^{\,o}/T_o$ est un schéma abélien. La construction de 5.1 s'applique alors pour fournir un sous-tore formel $\hat{T}$ du complété formel $\hat{A}$. On sait que pour tout $n$, le quotient

\begin{equation}\tag{7.1.1}
B_n = A_n^{\,o}/T_n
\end{equation}
est représentable (SGA 3 $\mathrm{VI}_A$ 3.2), plus précisément c'est un schéma plat et de type fini sur $S_n$, dont l'image inverse sur $S_o$ est le schéma abélien $B_o$. Par suite $B_n$ est un \underline{schéma abélien} sur $S_n$, et pour $n$ variable, les $B_n$ se déduisent les uns des autres par changement de base, donc ils définissent un \underline{schéma abélien formel}

\begin{equation}\tag{7.1.2}
\hat{B} = (B_n)_{n \geqslant 0} \qquad .
\end{equation}
Nous allons donner des conditions impliquant que ce schéma abélien est \underline{algébrisable} (EGA III 5.4.1), plus précisément que c'est le complété formel d'un schéma abélien projectif sur $S$, que nous noterons alors $B$. En vertu de loc.cit. $B$ est déterminé à isomorphisme unique près (comme $S$-schéma ; la structure de groupe de $\hat{B}$ provient alors de façon unique d'une structure de groupe de $B$, en vertu de loc. cit. également). Ce schéma abélien $B$ sera alors appelé \underline{la partie abélienne} du schéma $A$, et parfois noté $A^{ab}$.


\SGApage{I}{401}

\SGAvar{7.1.3. La projectivité de $\hat{B}$ signifie, en vertu de EGA III 5.4.5, que l'on peut trouver un Module inversible $\underline{\hat{M}}$ sur $\hat{B}$, i.e. une suite de Modules inversibles $\underline{M}_n$ sur les $B_n$ qui se ``recollent'' pour $n$ variable, qui soit \underline{ample} relativement à $S$ i.e. qui soit tel que $\underline{M}_o$ soit ample. L'idée naturelle serait de partir d'un Module inversible ample $\underline{M}$ de $G$, et de montrer que pour tout $n$, la restriction $\underline{L}_n$ de $\underline{L}$ à $G_n$ provient par image inverse d'un Module inversible $M_n$ sur $B_n$. J'ignore si de tels $M_n$ existent nécessairement (et en doute fort). Nous allons suivre une méthode un peu différente, en construisant une \underline{polarisation} de $B$, i.e. un système projectif de \underline{polarisations} $\rho_n$ des $B_n$ qui se recollent. On sait alors [5] que les doubles de ces polarisations définissent canoniquement des Modules inversibles $\underline{M}_n$ (symétriques) sur les $B_n$, qui par leur nature canonique se recolleront nécessairement, et fourniront le $\underline{M}$ cherché.

}{7.1.3. La projectivité de $\hat{B}$ signifie, en vertu de EGA III 5.4.5, que l'on peut trouver un Module inversible $\underline{\hat{M}}$ sur $\hat{B}$, i.e. une suite de Modules inversibles $\underline{M}_n$ sur les $B_n$ qui se ``recollent'' pour $n$ variable, qui soit \underline{ample} relativement à $S$ i.e. qui soit tel que $\underline{M}_o$ soit ample. L'idée naturelle serait de partir d'un Module inversible ample $\underline{L}$ sur $A$, et de montrer que pour tout $n$, la restriction $\underline{L}_n$ de $\underline{L}$ à $A_n$ provient par image inverse d'un Module inversible $M_n$ sur $B_n$. J'ignore si de tels $M_n$ existent nécessairement (et en doute fort). Nous allons suivre une méthode un peu différente, en construisant une \underline{polarisation} de $\hat{B}$, i.e. un système projectif de \underline{polarisations} $\rho_n$ des $B_n$ qui se recollent. On sait alors [5] que les doubles de ces polarisations définissent canoniquement des Modules inversibles $\underline{M}_n$ (symétriques) sur les $B_n$, qui par leur nature canonique se recolleront nécessairement, et fourniront le $\underline{\hat{M}}$ cherché.

}{F7I-IX-076}Nous allons interpréter un élément de $\mathrm{NS}(B_n/S_n)$ comme une correspondance divisorielle symétrique sur $(B_n,B_n)$ [5], i.e. comme une biextension de $(B_n,B_n)$ par $\mathbb{G}_{m\,S_n}$. En vertu de VIII 3.5, ces biextensions correspondent biunivoquement aux biextensions de $(A_n,A_n)$ par $\mathbb{G}_{mS}$. Pour obtenir un système cohérent de telles biextensions $W_n$ il suffit de partir d'une

\begin{equation}\tag{7.1.4}
W, \text{ biextension symétrique de } (A,A) \text{ par } \mathbb{G}_{mS} ,
\end{equation}
en prenant les $W_n$ se déduisant de $W$ par les changements de base $S_n \longrightarrow S$. Nous devons de plus supposer que l'élément de $\mathrm{NS}(B_o)$ défini par $W_o$ est une polarisation de $B_o$ pour obtenir ainsi une polarisation de $\hat{B}$.


\SGApage{I}{402}

\underline{Proposition} 7.1.5. \underline{Avec les notations précédentes pour} $S,A$, $\hat{B}$, \underline{supposons que} $S$ \underline{soit} normal. \underline{Alors} $\hat{B}$ \underline{est algébrisable}, \underline{et le schéma abélien} $B$ \underline{sur} $S$ \underline{dont il provient est projectif sur} $S$.

On peut évidemment supposer $A=A^o$, et $S$ intègre. D'après M. RAYNAUD [23 XI 1.13], l'hypothèse que $S$ est normal implique que $A$ est quasi-projectif sur $S$. Soit donc $\underline{L}$ un Module inversible sur $A$ qui est relativement ample, et formons le Module birigidifié $\delta(\underline{L}) = \pi^*(\underline{L})\mathrm{pr}_1^*(\underline{L})^{-1}\mathrm{pr}_2^*(\underline{L})^{-1}$ sur $A\times_S A$ comme dans VIII 4.12. Comme les fibres de $A_o$ sont de rang unipotent nul, il en est de même de celles de $A$ (résulte facilement de 3.1), donc en vertu de VIII 7.5, $\delta(\underline{L})$ définit une biextension $W$ de $(A,A)$, évidemment symétrique, qui en vertu de VIII 4.13 (i.e. essentiellement en vertu de RAYNAUD [23 XI 1.11]) satisfait à la condition d'amplitude voulue en les points de $S_o$.

\underline{Remarque} 7.1.6. On peut prouver, sans hypothèse de normalité sur $S$, que $\hat{B}$ est toujours algébrisable, en se ramenant au cas normal comme dans 13.5.3 a).

7.2. Notons maintenant le

\underline{Lemme} 7.2.1. \underline{Soient} $S$ \underline{le spectre d'un anneau noethérien} $V$, \underline{séparé et complet pour la topologie définie par un idéal} $\underline{m}$ \underline{de} $V$, $C$ \underline{la catégorie des schémas en groupes} $\Lambda$ \underline{sur} $S$, \underline{lisses de type fini sur} $S$, \underline{tels que} $\Lambda/\Lambda^o$ \underline{soit un schéma en groupes étale fini sur} $S$ \underline{et que} $\Lambda^o$ \underline{soit une extension d'un schéma abélien} $B$ \underline{par un tore isotrivial} $T$ (SGA 3 IX 1.1),


\SGApage{I}{403}

\SGAvar{$C'$ \underline{la catégorie des groupes dans la catégorie des schémas formels de type fini} (EGA I 10.3.3 ) \underline{sur} $\mathrm{Specf}(A)$, \underline{et} $A \longmapsto \hat{A}$ \underline{le foncteur de} $C$ \underline{dans} $C'$ "\underline{complété formel le long de} $A_o = A \times_S S_o$ ", \underline{où} $S_o = \mathrm{Spec}(V/\underline{m})$. \underline{Alors} :

}{$C'$ \underline{la catégorie des groupes dans la catégorie des schémas formels de type fini} (EGA I 10.13.3 ) \underline{sur} $\mathrm{Specf}(V)$, \underline{et} $A \longmapsto \hat{A}$ \underline{le foncteur de} $C$ \underline{dans} $C'$ "\underline{complété formel le long de} $A_o = A \times_S S_o$ ", \underline{où} $S_o = \mathrm{Spec}(V/\underline{m})$. \underline{Alors} :

}{F7I-IX-077}a) \underline{Le foncteur précédent est pleinement fidèle.}

\SGAvar{b) \underline{Soit} $\mathcal{A}$ \underline{un objet de} $C'$, \underline{lisse sur} $\hat{S} = \mathrm{Specf}(A)$ (\underline{i.e.} \underline{tel que les schémas ordinaires} $A_n = A \times_{\hat{S}} S_n$ \underline{sur les} $S_n = \mathrm{Spec}(A/\underline{m}^{n+1})$ \underline{soient lisses}). \underline{Supposons que} $A_o/A_o^{\,o}$ \underline{soit un schéma en groupes étale fini sur} $S_o$, \underline{et que} $A_o^{\,o}$ \underline{soit extension d'un schéma abélien} $B_o$ \underline{par un tore isotrivial} $T_o$ . \underline{Soient} $\mathcal{A}^o = (A_n^{\,o})_{n \geq 0}$, $\mathcal{A}/\mathcal{A}^o = (A_n/A_n^{\,o})$ $n \geq 0$ . \underline{Alors} $\mathcal{A}/\mathcal{A}^o$ \underline{est étale fini sur} $\hat{S}$, \underline{et} $\mathcal{A}^o$ \underline{est de façon essentiellement unique, et fonctorielle en} $\mathcal{A}$, \underline{extension d'un schéma formel abélien} $\mathcal{B}$ \underline{sur} $S$ \underline{par un tore formel} $\mathcal{T}$ \underline{sur} $S$, \underline{se réduisant respectivement suivant} $B_o$ \underline{et} $T_o$. \underline{Pour que} $\mathcal{A}$ \underline{appartienne à l'image essentielle du foncteur précédent} $C \longrightarrow C'$ (\underline{ou}, \underline{comme nous dirons aussi}, \underline{pour qu'il soit} "algébrisable"), \underline{il faut et il suffit qu'il en soit ainsi de} $\mathcal{B}$ .

}{b) \underline{Soit} $\mathcal{A}$ \underline{un objet de} $C'$, \underline{lisse sur} $\hat{S} = \mathrm{Specf}(V)$ (\underline{i.e.} \underline{tel que les schémas ordinaires} $A_n = \mathcal{A} \times_{\hat{S}} S_n$ \underline{sur les} $S_n = \mathrm{Spec}(V/\underline{m}^{n+1})$ \underline{soient lisses}). \underline{Supposons que} $A_o/A_o^{\,o}$ \underline{soit un schéma en groupes étale fini sur} $S_o$, \underline{et que} $A_o^{\,o}$ \underline{soit extension d'un schéma abélien} $B_o$ \underline{par un tore isotrivial} $T_o$ . \underline{Soient} $\mathcal{A}^o = (A_n^{\,o})_{n \geq 0}$, $\mathcal{A}/\mathcal{A}^o = (A_n/A_n^{\,o})$ $n \geq 0$ . \underline{Alors} $\mathcal{A}/\mathcal{A}^o$ \underline{est étale fini sur} $\hat{S}$, \underline{et} $\mathcal{A}^o$ \underline{est de façon essentiellement unique, et fonctorielle en} $\mathcal{A}$, \underline{extension d'un schéma formel abélien} $\mathcal{B}$ \underline{sur} $S$ \underline{par un tore formel} $\mathcal{T}$ \underline{sur} $S$, \underline{se réduisant respectivement suivant} $B_o$ \underline{et} $T_o$. \underline{Pour que} $\mathcal{A}$ \underline{appartienne à l'image essentielle du foncteur précédent} $C \longrightarrow C'$ (\underline{ou}, \underline{comme nous dirons aussi}, \underline{pour qu'il soit} "algébrisable"), \underline{il faut et il suffit qu'il en soit ainsi de} $\mathcal{B}$ .

}{F7I-IX-078}7.2.1.1. Notons d'abord, sur un schéma quelconque $S$, que si $A$ est un schéma en groupes lisse tel que $A^o$ soit extension d'un schéma abélien par un tore, cette structure d'extension est essentiellement unique, et est fonctorielle en $A$, dans un sens évident. En effet, notre assertion se ramène aussitôt à celle-ci : tout morphisme d'un tore $T$ dans un schéma abélien $B$ est nul. Or il est bien connu qu'il en est ainsi sur les fibres géométriques [15], et on conclut grâce à SGA 3 IX 5.3. Appliquant ce résultat aux systèmes inductifs de schémas en groupes lisses $A_n$ sur


\SGApage{I}{404}

les $S_n$ intervenant dans 7.2.1, on en conclut un résultat analogue pour les schémas formels en groupes sur $S$. D'ailleurs, on a déjà signalé dans 7.1 que si $A_o^{\,o}$ est une extension d'un schéma abélien par un tore, il en est de même de $A_n^{\,o}$. Ceci prouve donc la première assertion de b).

7.2.1.2. Supposons d'autre part que le schéma formel en groupes $\mathcal{A}$ sur $S$ soit de la forme $\hat{A}$, avec $A \in \mathrm{ob}\ C$. Si $A^o$ est extension d'un schéma abélien $B$ par un tore $T$, on voit immédiatement qu'on a $\mathcal{B} \simeq \hat{B}$, ce qui prouve que $\mathcal{A}$ algébrisable implique $\mathcal{B}$ algébrisable. Réciproquement, supposons $\mathcal{B}$ algébrisable, i.e. qu'il provient d'un schéma abélien $B$, et prouvons que $\mathcal{A}$ l'est. Prouvons d'abord que $\mathcal{A}^o = (A_n^{\,o})$ est algébrisable. Pour ceci, notons que les extensions $A_n^{\,o}$ de $B_n$ par $T_n$ peuvent s'interpréter, en vertu de VIII 3.7, comme des homomorphismes

\SGAvar{\begin{equation}\tag{*}
\underline{M}_n \longrightarrow B'_n \qquad ,
\end{equation}
où $\underline{M}_n$ est le groupe constant tordu sur $S_n$ dual de $T_n$ (SGA 3 X 5.7), et $B'_n$ le schéma abélien dual de $B_n$. Comme le foncteur restriction à $S_o$, ou à un $S_n$ quelconque, induit une équivalence de la catégorie des tores isotriviaux sur $S$ dans celle des tores isotriviaux sur $S_o$ (SGA 3 X 3.2), les $T_n$ proviennent d'un tore $T$ sur $S$, défini à isomorphisme unique près, dont nous désignerons par $M$ le groupe dual;sa restriction à $S_n$ est donc canoniquement isomorphe à $\underline{M}_n$. La propriété d'isotrivialité de $T$ s'interprète par le fait que $\underline{M}$ est un schéma somme de schémas $\underline{M}^i$ \underline{finis} sur $S$ (SGA 3 X 5.11). Il en résulte aussitôt que les homomorphismes (*) proviennent d'un unique homomorphisme

}{\begin{equation}\tag{*}
\underline{M}_n \longrightarrow B'_n \qquad ,
\end{equation}
où $\underline{M}_n$ est le groupe constant tordu sur $S_n$ dual de $T_n$ (SGA 3 X 5.7), et $B'_n$ le schéma abélien dual de $B_n$. Comme le foncteur restriction à $S_o$, ou à un $S_n$ quelconque, induit une équivalence de la catégorie des tores isotriviaux sur $S$ dans celle des tores isotriviaux sur $S_o$ (SGA 3 X 3.2), les $T_n$ proviennent d'un tore $T$ sur $S$, défini à isomorphisme unique près, dont nous désignerons par $\underline{M}$ le groupe dual;sa restriction à $S_n$ est donc canoniquement isomorphe à $\underline{M}_n$. La propriété d'isotrivialité de $T$ s'interprète par le fait que $\underline{M}$ est un schéma somme de schémas $\underline{M}^i$ \underline{finis} sur $S$ (SGA 3 X 5.11). Il en résulte aussitôt que les homomorphismes (*) proviennent d'un unique homomorphisme

}{F7I-IX-079}\begin{equation}\tag{**}
\underline{M} \longrightarrow B' \qquad ,
\end{equation}

\SGApage{I}{405}

où $B'$ est le schéma abélien dual de $B$. Cet homomorphisme à son tour s'interprète comme une extension de $B$ par $T$, soit $A^{o}$, qui fournit manifestement une algébrisation de $\mathcal{A}^{o}$.

\SGAvar{7.2.1.3.\quad D'autre part, comme par hypothèse $A_{o}/A_{o}^{\;o}$ est un schéma en groupes fini $\Phi_{o}$ sur $S$, les $A_{n}/A_{n}^{\;o}$ sont également des schémas en groupes $\Phi_{n}$ finis sur $S$ (leur représentabilité résultant de \hspace{3cm}), qui sont d'ailleurs nécessairement étales sur $S$. Il existe alors un schéma en groupes étale $\Phi$ sur $S$, défini à isomorphisme unique près, donnant naissance au schéma formel en groupes $\mathcal{A}/\mathcal{A}^{o}$ défini par les $\Phi_{n}$. Or, $G = A^{o}$ étant divisible, on voit aussitôt que la catégorie des extensions de faisceaux fppf d'un groupe fini localement constant sur $S$, tel que $\Phi$, par le schéma en groupes $A^{o}$, est canoniquement équivalente à la catégorie des torseurs sous le faisceau

}{7.2.1.3.\quad D'autre part, comme par hypothèse $A_{o}/A_{o}^{\;o}$ est un schéma en groupes fini $\Phi_{o}$ sur $S_o$, les $A_{n}/A_{n}^{\;o}$ sont également des schémas en groupes $\Phi_{n}$ finis sur $S_n$ (leur représentabilité résultant de \hspace{3cm}), qui sont d'ailleurs nécessairement étales sur $S_n$. Il existe alors un schéma en groupes étale $\Phi$ sur $S$, défini à isomorphisme unique près, donnant naissance au schéma formel en groupes $\mathcal{A}/\mathcal{A}^{o}$ défini par les $\Phi_{n}$. Or, $G = A^{o}$ étant divisible, on voit aussitôt que la catégorie des extensions de faisceaux fppf d'un groupe fini localement constant sur $S$, tel que $\Phi$, par le schéma en groupes $A^{o}$, est canoniquement équivalente à la catégorie des torseurs sous le faisceau

}{F7I-IX-080}\begin{equation*}
H = \underline{\mathrm{Hom}}_{\mathrm{gr}}(\Phi \; , \; G) \qquad ,
\end{equation*}
l'équivalence en question étant compatible avec tout changement de base. Appliquant la même description des extensions de $\Phi_{n}$ par $G_{n}$ en termes de torseurs sous $H_{n} \simeq \underline{\mathrm{Hom}}_{\mathrm{gr}}(\Phi_{n},G_{n})$, les $A_{n}$ nous définissent un système cohérent de torseurs $P_{n}$ sous les $H_{n}$. Or, $G$ étant une extension d'un schéma abélien par un tore, on en conclut aisément que $H$ est un schéma en groupes \underline{fini et localement libre sur} $S$, d'où on conclut facilement que le foncteur $P \longmapsto \hat{P} = (P_{n})$ définit une équivalence entre la catégorie des torseurs sur $S$ sous $H$, avec la catégorie des torseurs ``formels'' $\hat{P}$ sous $\hat{H}$, i.e. des systèmes cohérents de torseurs $P_{n}$ sous les $H_{n}$. Donc


\SGApage{I}{406}

\SGAvar{les $P_{n}$ définis par les extensions $A_{n}$ proviennent d'un même torseur $P$ sur $S$ sous $H$. On en conclut une extension $A$ de $\Phi$ par $G = A^{o}$ qui ``algébrise'' $\mathcal{A}$ , à cela près qu'il fait encore vérifier qu'elle est représentable. Or on peut regarder $A$ comme un torseur sur $\Phi$, de groupe $A_{\Phi}^{o}$ , et il est immédiat que ce torseur devient trivial après l'extension de la base par le morphisme fini surjectif localement libre $P \times_{S} \Phi \longrightarrow \Phi$ (puisque le changement de base $P \longrightarrow S$ scindé l'extension envisagée $A$ de $\Phi$ par $A^{o}$). On peut alors conclure par \hspace{1.5cm} un des nombreux critères d'effectivité de descente de M. RAYNAUD [23 XI 3.1 6)].

}{les $P_{n}$ définis par les extensions $A_{n}$ proviennent d'un même torseur $P$ sur $S$ sous $H$. On en conclut une extension $A$ de $\Phi$ par $G = A^{o}$ qui ``algébrise'' $\mathcal{A}$ , à cela près qu'il faut encore vérifier qu'elle est représentable. Or on peut regarder $A$ comme un torseur sur $\Phi$, de groupe $A_{\Phi}^{o}$ , et il est immédiat que ce torseur devient trivial après l'extension de la base par le morphisme fini surjectif localement libre $P \times_{S} \Phi \longrightarrow \Phi$ (puisque le changement de base $P \longrightarrow S$ scinde l'extension envisagée $A$ de $\Phi$ par $A^{o}$). On peut alors conclure par \hspace{1.5cm} un des nombreux critères d'effectivité de descente de M. RAYNAUD [23 XI 3.1 6)].

}{F7I-IX-081}7.2.1.4.\quad Cela achève de prouver b), et il reste seulement à prouver l'assertion a) de pleine fidélité. Nous laissons au lecteur le soin de vérifier que celle-ci est contenue essentiellement dans la méthode de démonstration de b), ou, s'il le préfère, d'en chercher une démonstration directe indépendante.

\SGAvar{7.2.1.5.\quad On notera que la référence au délicat résultat de descente de RAYNAUD devient inutile si on suppose $A^{o}$ quasi-projectif sur $S$, ou ce qui revient au même, $B$ projectif sur $S$. Cette condition est satisfait en particulier quand on part du schéma formel en groupes $\hat{A}$ sur $S$, où $A$ est comme dans 7.1.5, - ce qui est le seul cas où nous allons utiliser 7.2.1.

}{7.2.1.5.\quad On notera que la référence au délicat résultat de descente de RAYNAUD devient inutile si on suppose $A^{o}$ quasi-projectif sur $S$, ou ce qui revient au même, $B$ projectif sur $S$. Cette condition est satisfaite en particulier quand on part du schéma formel en groupes $\hat{A}$ sur $S$, où $A$ est comme dans 7.1.5, - ce qui est le seul cas où nous allons utiliser 7.2.1.

}{F7I-IX-082}7.2.2.\quad Nous allons appliquer 7.2.1 au schéma formel en groupes $\hat{A}$ envisagé dans 7.1, dans le cas favorable envisagé dans 7.1 où le schéma abélien $\hat{B}$ est algébrisable (cf. 7.1.5), et où on suppose de plus que le tore $T_{o}$ sur $S_{o}$ est isotrivial. On trouve donc un schéma en groupes $A^{\natural}$ sur $S$, appelé le \underline{Groupe de Raynaud} de $A$ sur $S$, caractérisé à isomor
\SGApage{I}{407}
phisme près par un isomorphisme canonique de schémas formels en groupes sur $\mathrm{Specf}(V) = \hat{S}$ :

\begin{equation}\tag{7.2.3}
(A^{\natural})^{\wedge} \simeq \hat{A} \quad .
\end{equation}
Il résulte de 7.2.1 que

\begin{equation}\tag{7.2.4}
\Phi = A^{\natural}/A^{\natural o}
\end{equation}
est un groupe fini étale sur $S$, et que la composante neutre

\begin{equation*}
A^{\natural o} \quad = \quad A^{o\natural}
\end{equation*}
est munie d'une structure d'extension

\begin{equation}\tag{7.2.5}
0 \longrightarrow T \longrightarrow A^{\natural o} \longrightarrow B \longrightarrow 0 \quad ,
\end{equation}
où $B$ est le schéma abélien sur $S$ dont le complété formel est le schéma formel abélien $\hat{B}$ construit dans 7.1, et $T$ est caractérisé (à isomorphisme unique près) comme le tore isotrivial sur $S$ qui relève $T_{o}$. Cette structure d'extension (7.2.5) donne, par restriction aux $S_{n}$ , les structures d'extension sur les $A_{n}^{\natural o} = A_{n}^{o}$ définies par (7.1.1).

7.2.6. Il résulte d'autre part immédiatement de 7.2.1 que le Groupe de Raynaud $A^{\natural}$ , ainsi que la suite exacte (7.2.5), dépendent fonctoriellement de $A$ (satisfaisant aux conditions envisagées), et que leur formation commute à tout changement de base $S' \longrightarrow S$ associé à un homomorphisme d'anneaux topologiques adiques $V \longrightarrow V'$.

7.3. Considérons maintenant, pour un nombre premier fixé $\ell$, le pro-$\ell$-groupe $T_{\ell}(A^{o})$. Identifiant un schéma fini sur $S$ avec le schéma formel correspondant, et définissant la ``\underline{partie fixe}'' $T_{\ell}(A^{o})^{f}$ comme


\SGApage{I}{408}

d'habitude, on trouve un isomorphisme comme (5.1.5)

\begin{equation}\tag{7.3.1}
T_{\ell}(A^{o})^{f} \simeq T_{\ell}(\hat{A}^{o}) \quad ,
\end{equation}
et définissant la ``\underline{partie torique}'' comme dans (5.1.6)

\begin{equation}\tag{7.3.2}
T_{\ell}(A^{o})^{t} = T_{\ell}(\hat{T}) \simeq T_{\ell}(T) \subset T_{\ell}(A^{o})^{f} \simeq T_{\ell}(\hat{A}^{o}) \quad ,
\end{equation}
on trouve aussitôt, pour la ``\underline{partie abélienne}'', un isomorphisme

\begin{equation}\tag{7.3.3}
T_{\ell}(A^{o})^{ab} \stackrel{\mathrm{dfn}}{=} T_{\ell}(A^{o})^{f}/T_{\ell}(A^{o})^{t} \simeq T_{\ell}(\hat{B}) \quad .
\end{equation}
On n'a pas encore ici utilisé l'hypothèse que $\hat{B}$ est algébrisable, hypothèse qui nous a permis de construire le schéma abélien $B$ et l'extension de Raynaud de $B$ par $T$. Comme on a évidemment un isomorphisme canonique (cas particulier de (7.3.1) appliqué au schéma abélien $B$ !) $T_{\ell}(\hat{B}) \simeq T_{\ell}(B)$, l'isomorphisme (7.3.3) nous donne aussi :

\begin{equation}\tag{7.3.4}
T_{\ell}(A^{o})^{ab} \simeq T_{\ell}(B) \quad .
\end{equation}
Notons aussi, en passant, que (7.2.3) et (7.3.1) impliquent aussi un isomorphisme canonique

\begin{equation}\tag{7.3.5}
T_{\ell}(A^{o})^{f} \simeq T_{\ell}(A^{\natural o}) \quad ,
\end{equation}
compte tenu que le deuxième membre est déjà un groupe pro-fini, i.e. identique à sa partie fixe. La suite exacte canonique

\SGAvar{\begin{equation}\tag{7.3.6}
0 \longrightarrow T_{\ell}(A^{o})^{t} \longrightarrow T_{\ell}(A^{o})^{f} \longrightarrow T_{\ell}(B) \longrightarrow 0 \quad ,
\end{equation}
qui peut en tout état de cause (sans hypothèse d'algébrisabilité sur $\hat{B}$) s'interprèter comme la suite exacte déduite par application du foncteur $T_{\ell}$ de la suite exacte

}{\begin{equation}\tag{7.3.6}
0 \longrightarrow T_{\ell}(A^{o})^{t} \longrightarrow T_{\ell}(A^{o})^{f} \longrightarrow T_{\ell}(B) \longrightarrow 0 \quad ,
\end{equation}
qui peut en tout état de cause (sans hypothèse d'algébrisabilité sur $\hat{B}$) s'interpréter comme la suite exacte déduite par application du foncteur $T_{\ell}$ de la suite exacte

}{F7I-IX-083}\begin{equation*}
0 \longrightarrow \hat{T} \longrightarrow \hat{A}^{o} \longrightarrow \hat{B} \longrightarrow 0 \quad ,
\end{equation*}

\SGApage{I}{409}

peut aussi être considérée comme déduite par application du foncteur $T_{\ell}$ de la suite exacte canonique de schémas en groupes ordinaires

\begin{equation*}
0 \longrightarrow T \longrightarrow A^{\natural o} \longrightarrow B \longrightarrow 0 \qquad .
\end{equation*}
7.4. Supposons maintenant donné un deuxième schéma en groupes $A'$ sur $S$, satisfaisant aux mêmes hypothèses que $G$, et supposons $A$, $A'$ à fibres connexes. Soit $W$ une biextension de $(A,A')$ par $\mathbb{G}_{mS}$. On en déduit donc un accouplement (VIII 2.2)

\begin{equation}\tag{7.4.1}
T_{\ell}(A) \times T_{\ell}(A') \longrightarrow T_{\ell}(\mathbb{G}_{mS}) = \mathbb{Z}_{\ell}(1)
\end{equation}
et par restriction aux parties fixes un accouplement

\SGAvar{\begin{equation}\tag{7.4.2}
T_{\ell}(A)^{f} \times T_{\ell}(A')^{f} = T_{\ell}(\hat{G}) \times T_{\ell}(\hat{G}') \longrightarrow \mathbb{Z}_{\ell}(1) \qquad ,
\end{equation}
qu'on peut aussi interprèter comme l'accouplement défini par la biextension $\hat{W}$ de $(\hat{A},\hat{A}')$ par $\hat{\mathbb{G}}_{mS}$ déduite de $W$ par complétion (en définissant ici, pour simplifier, une biextension de schémas en groupes formels comme une famille cohérente de biextensions $W_{n}$ des $(A_{n},A_{n}')$ par $\mathbb{G}_{mS_{n}}$ ). Or on sait (VIII 3.5) que $W_{n}$ provient d'une biextension $W_{n}^{ab}$ de $(B_{n},B_{n}')$ (où $B_{n}'$ est la partie abélienne de $G_{n}'$), de sorte que (7.4.2) se factorise en

}{\begin{equation}\tag{7.4.2}
T_{\ell}(A)^{f} \times T_{\ell}(A')^{f} = T_{\ell}(\hat{G}) \times T_{\ell}(\hat{G}') \longrightarrow \mathbb{Z}_{\ell}(1) \qquad ,
\end{equation}
qu'on peut aussi interpréter comme l'accouplement défini par la biextension $\hat{W}$ de $(\hat{A},\hat{A}')$ par $\hat{\mathbb{G}}_{mS}$ déduite de $W$ par complétion (en définissant ici, pour simplifier, une biextension de schémas en groupes formels comme une famille cohérente de biextensions $W_{n}$ des $(A_{n},A_{n}')$ par $\mathbb{G}_{mS_{n}}$ ). Or on sait (VIII 3.5) que $W_{n}$ provient d'une biextension $W_{n}^{ab}$ de $(B_{n},B_{n}')$ (où $B_{n}'$ est la partie abélienne de $G_{n}'$), de sorte que (7.4.2) se factorise en

}{F7I-IX-084}\begin{equation}\tag{7.4.3}
T_{\ell}(A)^{ab} \times T_{\ell}(A')^{ab} = T_{\ell}(\hat{B}) \times T_{\ell}(\hat{B}') \longrightarrow \mathbb{Z}_{\ell}(1) \qquad ,
\end{equation}
qui est aussi l'accouplement défini par la biextension

\begin{equation*}
\hat{W}^{ab} \;=\; (W_{n}^{ab})
\end{equation*}
de $(\hat{B},\hat{B}')$ par $\hat{\mathbb{G}}_{mS}$ . Lorsque $\hat{B}$, $\hat{B}'$ sont algébrisables, de sorte qu'on


\SGApage{I}{410}

\SGAvar{dispose de schémas abéliens $B,B'$ sur $S$, alors en vertu de EGA IV 5.1.4 $\hat{W}^{ab}$ provient d'une biextension, définie à isomorphisme unique près,

}{dispose de schémas abéliens $B,B'$ sur $S$, alors en vertu de EGA III 5.1.4 $\hat{W}^{ab}$ provient d'une biextension, définie à isomorphisme unique près,

}{F7I-IX-085}\begin{equation*}
W^{ab} \quad \text{de} \quad (B,B') \;\; \text{par} \;\; \mathbb{G}_{mS} \qquad ,
\end{equation*}
et l'accouplement correspondant

\begin{equation}\tag{7.4.4}
T_{\ell}(B) \times T_{\ell}(B') \longrightarrow T_{\ell}(\mathbb{G}_{mS})
\end{equation}
est canoniquement isomorphe à (7.4.3).

Désignant par $DB'$ le schéma abélien dual de $B'$, $W^{ab}$ définit un homomorphisme

\begin{equation}\tag{7.4.5}
u : B \longrightarrow D(B') \qquad ,
\end{equation}
dont les propriétés reflètent celles des accouplements (7.4.4) i.e. (7.4.3). On voit par exemple que l'accouplement (7.4.3) est une dualité parfaite en un point $t \in S$ si et seulement si $u_{t} : B_{t} \longrightarrow D(B')_{t}$ est une isogénie dont le noyau est de rang premier à $\ell$ . Comme ceci implique que $u$ lui-même est une isogénie ($S$ étant local donc connexe), et que $\mathrm{Ker}\,u$ est plat et fini sur $S$, donc que $\mathrm{Ker}\,u$ a même rang en tout point, on voit que la condition qu'on vient d'envisager est indépendante du point choisi $t$.

7.4.6. Ceci achève de prouver en particulier 5.4, puisque sous les conditions de ce théorème (où on peut supposer $S$ complet), appliquant ce qui précède à $A^{o}$ et ${A'}^{o}$, on a établi que l'accouplement (7.4.3) est pour tout $\ell$ une dualité parfaite au point générique $\eta$, de sorte qu'il en est de même au point $s$.


\SGApage{I}{411}

\SGAvar{7.4.7. On voit de même par les considérations précédentes que le théorème d'orthogonalité 5.2 devient clair, sans recours au lemme 5.2.4. En effet, nous avons établi en cours de route dans 7.4 que l'accouplement $(5.2.3.1)=(7.4.2)$ s'annule sur les parties toriques de part et d'autre, donc se factorise en un accouplement (7.4.3), dont il reste à établir qu'il est séparé \underline{sur les fibres génériques} des pro-groupes envisagés, sachant qu'il l'est sur les fibres spéciales. Cela provient en effet du fait que si (7.4.5) induit une isogénie sur les fibres spéciales, c'est une isogénie.

}{7.4.7. On voit de même par les considérations précédentes que le théorème d'orthogonalité 5.2 devient clair, sans recours au lemme 5.2.5. En effet, nous avons établi en cours de route dans 7.4 que l'accouplement $(5.2.4.2)=(7.4.2)$ s'annule sur les parties toriques de part et d'autre, donc se factorise en un accouplement (7.4.3), dont il reste à établir qu'il est séparé \underline{sur les fibres génériques} des pro-groupes envisagés, sachant qu'il l'est sur les fibres spéciales. Cela provient en effet du fait que si (7.4.5) induit une isogénie sur les fibres spéciales, c'est une isogénie.

}{F7I-IX-086}7.4.8. Revenons aux conditions générales de 7.4, et supposons donné un ouvert $U \subset S-S_o$, tel que $A_U$ et $A'_U$ soient des \underline{schémas abéliens}, duals l'un de l'autre par la biextension $W_U$ . Considérons les restrictions à $U$ de la suite exacte (7.3.6) et de la suite exacte analogue relative à $A'$, de sorte qu'on trouve sur les pro-$\ell$-groupes de Barsotti-Tate $T_\ell(A_U)$ et $T_\ell(A'_U)$ des filtrations canoniques

\begin{equation}\tag{7.4.9}
\left\{\begin{aligned} T_\ell(A_U) &\supset T_\ell(A_U)^f \supset T_\ell(A_U)^t \supset 0 \\[4ex] T_\ell(A'_U) &\supset T_\ell(A'_U)^f \supset T_\ell(A'_U)^t \supset 0 \qquad . \end{aligned}\right.
\end{equation}
Ceci posé, je dis que \underline{ces deux filtrations sont duales l'une de l'autre}, relativement à l'accouplement (7.4.1) (qui est ici une dualité parfaite (VIII 3.2)), i.e. qu'on a encore la relation (5.2.2) et la relation symétrique (avec $A_K,A'_K$ remplacés par $A_U$ et $A'_U$). On a vu d'ailleurs ci-dessus que cet énoncé d'orthogonalité revient aussi à dire que


\SGApage{I}{412}

la biextension $W^{ab}$ de $(B,B')$ par $\mathbb{G}_{mS}$ est une dualité parfaite sur $U$ (ou ce qui revient au même si $U$ est dense dans $S$, sur $S$). C'est une propriété qu'il suffit de vérifier en tout point $\eta$ de $U$, et alors l'argument habituel nous ramène au cas où $S$ est un trait de point générique $\eta$ , qui est déjà connu.

La relation d'orthogonalité nous permet d'autre part d'étendre au cas présent les formules (5.5.7), (5.5.8) et (5.5.9), où on remplace l'opération de restriction au point générique $\eta$ par celle de restriction à l'ouvert $U$ ; par suite, 5.6 s'étend également au cas présent. Nous nous dispensons de réécrire ici les formules et l'énoncé précédents dans le contexte présent.

Résumons, du point de vue des modèles de Néron, les résultats obtenus concernant les parties abéliennes et leurs accouplements :

\underline{Proposition} 7.5. \underline{Soit} $S$ \underline{un trait complet}, \underline{de corps de fonctions} $K$. A \underline{tout schéma abélien} $A_K$ \underline{sur} $K$ \underline{ayant une réduction semi-stable sur} $S$, \underline{le procédé de} 7.1,\underline{appliqué à la composante neutre du modèle de Néron} $A$, \underline{de} $A_K$ \underline{permet de définir la} partie abélienne $B = A^{ab}$ \underline{de celui-ci}, \underline{qui est un schéma abélien sur} $S$, \underline{dépendant fonctoriellement de} $A_K$ . \underline{Pour tout autre schéma abélien} $A'_K$ \underline{à réduction semi-stable}, \underline{on peut à toute correspondance divisorielle} $W_K$ \underline{sur} $(A_K,A'_K)$, \underline{associer une correspondance divisorielle unique} $W^{ab}$ \underline{sur} $(A^{ab},{A'}^{ab})$, \underline{telle que pour tout entier} $n \geq 0$, $W_n^{ab}$ \underline{ait comme image inverse par} $A_n^o \longrightarrow B_n = A_n^{oab}$ , ${A'}^{o}_{n} \longrightarrow B'_n = {A'}^{oab}_{n}$ \underline{la biextension} $W_n$ \underline{de} $(A_n^o,{A'}^{o}_{n})$ \underline{déduite de} $W$ \underline{par réduction}, \underline{où} $W$ \underline{est le prolongement canonique} (1.4) \underline{de} $W_K$ \underline{à} $(A^o,{A'}^{o})$.


\SGApage{I}{413}

\underline{La formation de} $W^{ab}$ \underline{est compatible avec la formation d'images inverses de correspondances par des morphismes de schémas abéliens} $\overline{A}_K \longrightarrow A_K$ , $\overline{A'}_K \longrightarrow A'_K$ . \underline{Si} $W_K$ \underline{est une dualité parfaite}, \underline{il en est de même de} $W^{ab}$. \underline{Si} $W_K$ \underline{est la correspondance divisorielle de} $(A_K,A_K)$ \underline{associée à une polarisation de} $A_K$, \underline{alors} $W^{ab}$ \underline{est associé à une polarisation de} $A^{ab}$.

\SGAvar{L'assertion de compatibilité de la formation des $W^{ab}$ aves les images inverses est triviale, grâce à la caractérisation indiquée des $W^{ab}$. Le complément relatif aux polarisations a été vu dans 7.1.5 en même temps que la démonstration de l'algébrisabilité de $\hat{B}$.

}{L'assertion de compatibilité de la formation des $W^{ab}$ avec les images inverses est triviale, grâce à la caractérisation indiquée des $W^{ab}$. Le complément relatif aux polarisations a été vu dans 7.1.5 en même temps que la démonstration de l'algébrisabilité de $\hat{B}$.

}{F7I-IX-087}7.6. Supposons toujours que $S$ est un trait complet, $A$ étant le modèle de Néron d'un schéma abélien $A_K$ sur $K$. On pose alors

\begin{equation}\tag{7.6.1}
A_K^{\natural} \;=\; (A^{\natural})_K \quad ,
\end{equation}
où $A^{\natural}$ est le schéma de Raynaud (7.2.2), extension successive du Groupe fini étale $\Phi$ par le schéma abélien $B$ par le tore $T$. Donc $A_K^{\natural}$ est lui-même extension de $\Phi_K$ par $A_K^{\natural\,o} = A^{\natural o}_K = A^{o\natural}_K$, qui est extension du schéma abélien $B_K$ par le tore $T_K$ . Pour tout nombre premier $\ell$ , on a en vertu de (7.3.5) et (7.3.2) :

\begin{equation}\tag{7.6.2}
T_{\ell}(A_K^{\natural}) \;\simeq\; T_{\ell}(A_K^{\natural\,o}) \;\simeq\; T_{\ell}(A_K)^{f} \; , \; T_{\ell}(T_K) \;\simeq\; T_{\ell}(A_K)^{t} \quad .
\end{equation}
7.6.3. On notera qu'en vertu de 3.3, la formation de $A_K^{\natural\,o}$ commute à tout changement de base pour un morphisme de traits complets $S' \longrightarrow S$.


\SGApage{I}{414}


\SGAsectiontoc{8. L'extension de Raynaud $A_K^{\natural\,o}$ dans le cas non semi-stable, et le ind-groupe $A_K^{\flat}$}{8}
8. \underline{L'extension de Raynaud} $A_K^{\natural\,o}$ \underline{dans le cas non semi-stable}, \underline{et le ind-groupe} $A_K^{\flat}$

8.1. Supposons toujours le trait $S$ complet, mais ne supposons plus nécessairement que $A_K$ ait une réduction semi-stable. Soit comme d'habitude $K'$ une extension galoisienne finie de $K$ telle que $A_{K'}$ ait une réduction semi-stable par rapport au normalisé $S'$ de $S$ dans $K'$. Considérons l'extension de Raynaud $A_{K'}^{\natural}$ associée à $A_{K'}$ (7.6.1). Il est clair que le groupe de Galois de $K'/K$ opère par transport de structure sur ce schéma en groupes, d'où par descente galoisienne un schéma en groupes sur $K$, soit $A_{K,K'}^{\natural}$ , qui dépend du choix de l'extension $K'$. Cependant sa composante neutre n'en dépend pas à isomorphisme canonique près grâce à 7.6.3, et on la notera simplement $A_K^{\natural\,o}$. On voit par descente que $A_K^{\natural\,o}$ est une extension

\SGAvar{\begin{equation}\tag{8.1.1}
0 \longrightarrow T_K \longrightarrow A_K^{\natural\,o} \longrightarrow B_K \longrightarrow 0 \quad ,
\end{equation}
où $T_K$ est un tore et $B_K$ un schéma abélien sur $K$. De même, en vertu de la définition des parties essentiellement fixe et essentiellement torique de $T_{\ell}(A_K)$ (4.1 et 6.5), les isomorphismes (7.6.2) donnent des isomorphismes canoniques

}{\begin{equation}\tag{8.1.1}
0 \longrightarrow T_K \longrightarrow A_K^{\natural\,o} \longrightarrow B_K \longrightarrow 0 \quad ,
\end{equation}
où $T_K$ est un tore et $B_K$ un schéma abélien sur $K$. De même, en vertu de la définition des parties essentiellement fixe et essentiellement torique de $T_{\ell}(A_K)$ (4.1 et 6.3), les isomorphismes (7.6.2) donnent des isomorphismes canoniques

}{F7I-IX-088}\begin{equation}\tag{8.1.2}
T_{\ell}(A_K^{\natural}) \;\simeq\; T_{\ell}(A_K)^{ef} \; , \; T_{\ell}(T_K) \;\simeq\; T_{\ell}(A_K)^{et} \quad ,
\end{equation}
d'où par passage au quotient

\begin{equation}\tag{8.1.3}
T_{\ell}(B_K) \;\simeq\; T_{\ell}(A_K)^{ef}/T_{\ell}(A_K)^{et} \quad .
\end{equation}

\SGApage{I}{415}

Ceci détermine donc deux des quotients intervenant dans la filtration

\SGAvar{\begin{equation}\tag{8.1.4}
T_\ell(A_K) \supset T_\ell(A_K)^{ef} \supset T_\ell(A_K)^{et} \supset 0
\end{equation}
en termes de l'extension de schémas en groupes (8.1.1). Pour déterminer le troisième quotient $T_\ell(A_K)/T_\ell(A_K)^{ef}$, on note qu'en vertu de (6.5.2) on a un isomorphisme canonique

}{\begin{equation}\tag{8.1.4}
T_\ell(A_K) \supset T_\ell(A_K)^{ef} \supset T_\ell(A_K)^{et} \supset 0
\end{equation}
en termes de l'extension de schémas en groupes (8.1.1). Pour déterminer le troisième quotient $T_\ell(A_K)/T_\ell(A_K)^{ef}$, on note qu'en vertu de (6.3.2) on a un isomorphisme canonique

}{F7I-IX-089}\begin{equation}\tag{8.1.5}
T_\ell(A_K)/T_\ell(A_K)^{ef} \;\simeq\; D(T_\ell(A'_K)^{et}) \quad ,
\end{equation}
où $D = \underline{Hom}( - , T_\ell(\mathbb{G}_m))$ désigne le dual de Cartier (au sens des pro-groupes de Barsotti-Tate), et où $A'_K$ désigne le schéma abélien dual de $A_K$ . Considérons alors la suite exacte du type (8.1.1)

\begin{equation}\tag{8.1.6}
0 \longrightarrow T'_K \longrightarrow {A'_K}^{\natural\,o} \longrightarrow B'_K \longrightarrow 0
\end{equation}
associée au schéma abélien $A'_K$ sur $K$, et les isomorphismes du type (8.1.2) correspondants, la formule (8.1.5) nous donne donc

\begin{equation}\tag{8.1.7}
T_\ell(A_K)/T_\ell(A_K)^{ef} \;\simeq\; D(T_\ell(T'_K)) \;\simeq\; \underline{M}_K \otimes_{\mathbb{Z}} \mathbb{Z}_\ell \quad .
\end{equation}
où on a posé (en généralisant les notations introduites dans (5.5.2)) :

\begin{equation}\tag{8.1.8}
\underline{M}_K = D(T'_K) \quad , \qquad \underline{M}'_K = D(T_K) \quad ,
\end{equation}
et où maintenant $D$ désigne le dual de Cartier $\underline{Hom}( - , \mathbb{G}_{mK})$ au sens des schémas en groupes ordinaires. Ainsi $\underline{M}_K$ et $\underline{M}'_K$ sont des Groupes constants tordus sur $K$, qui sur $K'$ deviennent constants et isomorphes à des groupes de la forme $\mathbb{Z}^{\,r}$ ($r=\dim T = \dim T'$).


\SGApage{I}{416}

\SGAvar{8.2. Les conditions sont celles de 8.1. Il est clair que l'extension (8.1.1) dépend fonctoriellement du schéma abélien $A_K$, et que sa formation commute à tout morphisme de trait $\tilde{S} \longrightarrow S$. Choisissant une polarisation de $A_K$, d'où une isogénie $A_K \longrightarrow A_{K'}$, on en déduit un homomorphisme de (8.1.1) dans (8.1.6), dont on voit aussitôt (en regardant les $T_\ell$) que c'est une isogénie terme à terme. De même, utilisant 7.5 et la descente, on voit que toute biextension par $\mathbb{G}_{mK}$ de schémas abéliens $A_K$, $\bar{A}_K$ sur $K$ définit une biextension de $(B_K,\bar{B}_K)$ par $\mathbb{G}_{mK}$ , où $B_K$ , $\bar{B}_K$ désignent les quotients abéliens des extensions de Raynaud qui correspondent à $A_K$ et à $\bar{A}_K$ respectivement ; lorsque $W$ est une dualité parfaite, il en est de même de la biextension correspondante de $(B_K,\bar{B}_K)$. En particulier, avec les notations de 7.6, les schémas abéliens $B_K$ et $B'_K$ sont canoniquement duals l'un de l'autre. La correspondance entre biextensions sur $(A_K,\bar{A}_K)$ et biextensions sur $(B_K, \bar{B}_K)$ a des propriétés de fonctorialité évidentes pour $A_K,\bar{A}_K$ variables, qui se déduisent par descente des propriétés analogues énoncées dans 7.5, et dont nous laissons l'énoncé au lecteur. Nous laissons également au lecteur la compatibilité de l'application précédente avec tout changement de traits $\tilde{S} \longrightarrow S$.

8.3. Revenons au groupe $A_{K,K'}^{\natural}$ , et étudions sa variance en $K'$. Il est clair que deux extensions isomorphes $K'$, $K'_1$ de $K'$ donnent des groupes \underline{canoniquement} isomorphes, de sorte qu'on peut se borner aux sous-extensions $K'$ de la clôture séparable $\bar{K}$ de $K$. Si alors $K' \hookrightarrow K''$ sont deux telles sous-extensions galoisiennes finies de $\bar{K}$, il résulte aussitôt de 3.3


}{8.2. Les conditions sont celles de 8.1. Il est clair que l'extension (8.1.1) dépend fonctoriellement du schéma abélien $A_K$, et que sa formation commute à tout morphisme de trait $\tilde{S} \longrightarrow S$. Choisissant une polarisation de $A_K$, d'où une isogénie $A_K \longrightarrow A'_K$, on en déduit un homomorphisme de (8.1.1) dans (8.1.6), dont on voit aussitôt (en regardant les $T_\ell$) que c'est une isogénie terme à terme. De même, utilisant 7.5 et la descente, on voit que toute biextension par $\mathbb{G}_{mK}$ de schémas abéliens $A_K$, $\bar{A}_K$ sur $K$ définit une biextension de $(B_K,\bar{B}_K)$ par $\mathbb{G}_{mK}$ , où $B_K$ , $\bar{B}_K$ désignent les quotients abéliens des extensions de Raynaud qui correspondent à $A_K$ et à $\bar{A}_K$ respectivement ; lorsque $W$ est une dualité parfaite, il en est de même de la biextension correspondante de $(B_K,\bar{B}_K)$. En particulier, avec les notations de 7.6, les schémas abéliens $B_K$ et $B'_K$ sont canoniquement duals l'un de l'autre. La correspondance entre biextensions sur $(A_K,\bar{A}_K)$ et biextensions sur $(B_K, \bar{B}_K)$ a des propriétés de fonctorialité évidentes pour $A_K,\bar{A}_K$ variables, qui se déduisent par descente des propriétés analogues énoncées dans 7.5, et dont nous laissons l'énoncé au lecteur. Nous laissons également au lecteur la compatibilité de l'application précédente avec tout changement de traits $\tilde{S} \longrightarrow S$.

8.3. Revenons au groupe $A_{K,K'}^{\natural}$ , et étudions sa variance en $K'$. Il est clair que deux extensions isomorphes $K'$, $K'_1$ de $K$ donnent des groupes \underline{canoniquement} isomorphes, de sorte qu'on peut se borner aux sous-extensions $K'$ de la clôture séparable $\bar{K}$ de $K$. Si alors $K' \hookrightarrow K''$ sont deux telles sous-extensions galoisiennes finies de $\bar{K}$, il résulte aussitôt de 3.3


}{F7I-IX-090}\SGApage{I}{417}

que l'on a une \underline{immersion ouverte} canonique

\SGAvar{\begin{equation}\tag{8.3.1}
A^{\natural}_{K,K'} \hookrightarrow A^{\natural}_{K,K''} \qquad ,
\end{equation}
permettant d'identifier le premier groupe à un sous-groupe ouvert du deuxième. Pour $K'$ variable, on voit que les $A_{K,K'}$ forment ainsi un système inductif de groupes lisses de type fini sur $K$, à morphismes de transition des immersions ouvertes. Si on pose

}{\begin{equation}\tag{8.3.1}
A^{\natural}_{K,K'} \hookrightarrow A^{\natural}_{K,K''} \qquad ,
\end{equation}
permettant d'identifier le premier groupe à un sous-groupe ouvert du deuxième. Pour $K'$ variable, on voit que les $A^{\natural}_{K,K'}$ forment ainsi un système inductif de groupes lisses de type fini sur $K$, à morphismes de transition des immersions ouvertes. Si on pose

}{F7I-IX-091}\begin{equation}\tag{8.3.2}
\Phi_{K,K'} \;=\; A^{\natural}_{K,K'}/A^{\natural}_{K,K'}{}^{o} \qquad ,
\end{equation}
les $\Phi_{K,K'}$ eux-mêmes forment donc un système inductif de groupes finis étales sur $K$, qui sera étudié plus en détail plus bas.

Nous poserons

\SGAvar{\begin{equation}\tag{8.3.3}
A^{\flat}_{K} \;=\; \varinjlim_{K'} \; A^{\natural}_{K,K'} \qquad ,
\end{equation}
la limite étant prise dans la catégorie des faisceaux fppf sur $K$. Comme les morphismes de transition dans le système inductif envisagé dans des immersions ouvertes, il est d'ailleurs clair que cette limite est en fait représentable par un schéma en groupes (en général non quasi-compact) sur $K$, que nous désignerons comme de juste par le même symbole $A^{\flat}_{K}$ . On notera qu'on a, pour ce schéma en groupes

}{\begin{equation}\tag{8.3.3}
A^{\flat}_{K} \;=\; \varinjlim_{K'} \; A^{\natural}_{K,K'} \qquad ,
\end{equation}
la limite étant prise dans la catégorie des faisceaux fppf sur $K$. Comme les morphismes de transition dans le système inductif envisagé sont des immersions ouvertes, il est d'ailleurs clair que cette limite est en fait représentable par un schéma en groupes (en général non quasi-compact) sur $K$, que nous désignerons comme de juste par le même symbole $A^{\flat}_{K}$ . On notera qu'on a, pour ce schéma en groupes

}{F7I-IX-092}\begin{equation}\tag{8.3.4}
A_{K}^{\flat\,o} \;=\; A_{K}^{\natural\,o} \quad , \quad A^{\flat}_{K} \,/\, A_{K}^{\flat\,o} \;\simeq\; \varinjlim_{K'} \Phi_{K,K'} \qquad .
\end{equation}
Ce dernier groupe sera déterminé dans 11.10 : il est canoniquement isomorphe à $\underline{M}_{K}\otimes\mathbb{Q}/\mathbb{Z}$ (avec la notation (8.1.8)).


\SGApage{I}{418}


\SGAsectiontoc{9. Définition de l'accouplement de monodromie $\protect\underline{M}_{\ell}\otimes\protect\underline{M}'_{\ell} \longrightarrow \mathbb{Z}_{\ell S}$ .}{9}
9. \underline{Définition de l'accouplement de monodromie} $\underline{M}_{\ell}\otimes\underline{M}'_{\ell} \longrightarrow \mathbb{Z}_{\ell S}$ .

\SGAvar{9.1. Dans le présent numéro on suppose $S$ hensélien, et que $A_{K}$ ait réduction semi-stable sur $S$, ou ce qui revient au même, que $A'_{K}$ ait réduction semi-stable sur $S$ (3.5.1). On dispose alors sur $S$ de deux $\mathbb{Z}$-Modules constants tordus $\underline{M}$ et $M'$ (5.5.2), qui sur un revêtement fini étale convenable de $S$ (donc sur $S$ lui-même si $S$ est strictement local i.e. si $k$ est séparablement clos) se déploient en des groupes constants de valeur $M$ resp. $M'$ isomorphes (non canoniquement) à $\mathbb{Z}^{r}$, $r$ étant la dimension commune des tores maximaux de $A_{o}$ et de $A'_{o}$. Nous allons considérer, pour tout nombre premier $\ell$, les faisceaux $\ell$-adiques correspondants

}{9.1. Dans le présent numéro on suppose $S$ hensélien, et que $A_{K}$ ait réduction semi-stable sur $S$, ou ce qui revient au même, que $A'_{K}$ ait réduction semi-stable sur $S$ (3.5.1). On dispose alors sur $S$ de deux $\mathbb{Z}$-Modules constants tordus $\underline{M}$ et $\underline{M}'$ (5.5.2), qui sur un revêtement fini étale convenable de $S$ (donc sur $S$ lui-même si $S$ est strictement local i.e. si $k$ est séparablement clos) se déploient en des groupes constants de valeur $M$ resp. $M'$ isomorphes (non canoniquement) à $\mathbb{Z}^{r}$, $r$ étant la dimension commune des tores maximaux de $A_{o}$ et de $A'_{o}$. Nous allons considérer, pour tout nombre premier $\ell$, les faisceaux $\ell$-adiques correspondants

}{F7I-IX-093}\begin{equation}\tag{9.1.1}
\underline{M}_{\ell} \;=\; \underline{M}\otimes_{\mathbb{Z}}\mathbb{Z}_{\ell} \quad , \quad \underline{M}'_{\ell} \;=\; \underline{M}'\otimes_{\mathbb{Z}}\mathbb{Z}_{\ell} \qquad ,
\end{equation}
dont la signification géométrique est explicitée par les formules (5.5.7) et (5.5.8). Dans le présent numéro, nous allons définir un accouplement canonique (dit ``\underline{accouplement de monodromie}'')

\begin{equation}\tag{9.1.2}
u_{\ell} \;:\; \underline{M}_{\ell} \otimes \underline{M}'_{\ell} \longrightarrow \mathbb{Z}_{\ell} \qquad ,
\end{equation}
qui sera défini en termes des pro-$\ell$-groupes de Barsotti-Tate $T_{\ell}(A^{o})$ et $T_{\ell}(A'{}^{o})$, et même de la fibre générique $T_{\ell}(A_{K})$ du premier de ces groupes (sous réserve, lorsque $\ell=p$, de disposer du théorème de Tate [33, th.4], i.e. pour l'instant qu'on soit alors dans le cas d'inégales caractéristiques). Lorsque $\ell \neq p$, il revient au même de dire que l'accouplement (9.1.2) est décrit en termes du $\pi$-module galoisien

\begin{equation}\tag{9.1.3}
U \;=\; T_{\ell}(A_{K})(\overline{K}) \;=\; T_{\ell}(A_{\overline{K}})
\end{equation}

\SGApage{I}{419}

envisagé dans $(2.5.3)$, et plus précisément, nous verrons que sa connaissance équivaut essentiellement à celle de l'action du groupe de monodromie $I$ sur $U$.

\SGAvar{9.2. \underline{Cas} $\ell \neq p$ . Plaçons nous donc d'abord dans le cas, plus simple, où $\ell \neq p$, en reprenant alors les notations introduites dans $(2.5.3)$. Notons qu'à torsion près, la relation $V = U^I$ signifie aussi que $V^{\perp}$ est le module engendré par les éléments $gx - x$ ($u \in U$, $g \in I$), et qu'on a en tous cas

}{9.2. \underline{Cas} $\ell \neq p$ . Plaçons nous donc d'abord dans le cas, plus simple, où $\ell \neq p$, en reprenant alors les notations introduites dans $(2.5.3)$. Notons qu'à torsion près, la relation $V = U^I$ signifie aussi que $V^{\perp}$ est le module engendré par les éléments $gx - x$ ($x \in U$, $g \in I$), et qu'on a en tous cas

}{F7I-IX-094}\begin{equation}\tag{9.2.1}
gx - x \in V^{\perp} \ \ \mathrm{pour}\ x \in U,\ g \in I \qquad .
\end{equation}
On a donc, pour $g \in I$ :

\SGAvar{\begin{equation}\tag{9.2.2}
gx = x + u(g).x \qquad ,
\end{equation}
où $g \longmapsto u(g)$ est un homomorphisme de $U$ dans $V^{\perp}$ qui s'annule sur $V= U^I$, ou ce qui revient au même

}{\begin{equation}\tag{9.2.2}
gx = x + u(g).x \qquad ,
\end{equation}
où $u(g)$ est un homomorphisme de $U$ dans $V^{\perp}$ qui s'annule sur $V= U^I$, ou ce qui revient au même

}{F7I-IX-095}\begin{equation}\tag{9.2.3}
u \in \mathrm{Hom}(U/V\ ,\ V^{\perp}) \ \simeq\ (U/V)^{\vee} \otimes V^{\perp} \qquad .
\end{equation}
Dans le cas envisagé de la réduction semi-stable, i.e. quand on a $V^{\perp} \subset V$, la formule $(9.2.2)$ exprime que l'opération $g$ dans $U$ est une \underline{transvection} parallèle à $V^{\perp} = W$, et la composition de ces transvections n'est autre que l'addition des homomorphismes $u$ correspondants $(9.2.3)$. Donc l'opération du groupe d'inertie $I$ sur $U$ s'explicite entièrement alors par un homomorphisme de groupes

\begin{equation}\tag{9.2.4}
I \longrightarrow (U/V)^{\vee} \otimes V^{\perp} \qquad .
\end{equation}

\SGApage{I}{420}

\SGAvar{Or par le théorème d'orthogonalité appliqué a la fois pour $A_K$ et $A'_K$ on a

\begin{equation*}
V^{\perp} = W \quad ,\quad (U/V)^{\vee} \ \simeq\ W'(-1) \qquad ,
\end{equation*}
où $W'$ est la partie torique de $U'=T_{\ell}(A'(\overline{K}))$, et où le $(-1)$ désigne le twist de Tate. D'autre part, l'homomorphisme $(9.2.4)$ se factorise par $I(\ell) \simeq \mathbb{Z}_{\ell}(1)$, de sorte que l'homomorphisme $(9.2.4)$ peut s'interprèter comme un homomorphisme canonique

}{Or par le théorème d'orthogonalité appliqué à la fois pour $A_K$ et $A'_K$ on a

\begin{equation*}
V^{\perp} = W \quad ,\quad (U/V)^{\vee} \ \simeq\ W'(-1) \qquad ,
\end{equation*}
où $W'$ est la partie torique de $U'=T_{\ell}(A'(\overline{K}))$, et où le $(-1)$ désigne le twist de Tate. D'autre part, l'homomorphisme $(9.2.4)$ se factorise par $I(\ell) \simeq \mathbb{Z}_{\ell}(1)$, de sorte que l'homomorphisme $(9.2.4)$ peut s'interpréter comme un homomorphisme canonique

}{F7I-IX-096}\begin{equation}\tag{9.2.5}
u: \mathbb{Z}_{\ell}(1) \longrightarrow W'\otimes W(-1) \qquad ,
\end{equation}
ou encore comme un élément

\begin{equation}\tag{9.2.6}
u \in W' \otimes W\ (-2) \qquad .
\end{equation}
Se rappelant de la définition de $W$ et de $W'$, comme les groupes de points à valeurs dans $\overline{K}$ des deux faisceaux $\ell$-adiques constants tordus $T_{\ell}(A^o)^t$ et $T_{\ell}({A'}^{o})^t$, dont les fibres spéciales sont respectivement $T_{\ell}(T_o)$ et $T_{\ell}(T'_o)$, on trouve les isomorphismes canoniques

\SGAvar{\begin{equation}\tag{9.2.7}
\left\{ \begin{array}{l} W \ \simeq\ T_{\ell}(T_o)(\overline{k})) \ \ \simeq\ \check{M} \otimes \mathbb{Z}_{\ell}(1) \\[3ex] W' \ \simeq\ T_{\ell}(T_o(\overline{k})) \ \ \simeq\ M' \otimes \mathbb{Z}_{\ell}(1) \qquad , \end{array} \right.
\end{equation}
$M$ et $M'$ désignant les groupes des caractères relativement à $\overline{k}$ de $T_o$ et de $T'_o$ respectivement :

}{\begin{equation}\tag{9.2.7}
\left\{ \begin{array}{l} W \ \simeq\ T_{\ell}(T_o)(\overline{k}) \ \ \simeq\ \check{M} \otimes \mathbb{Z}_{\ell}(1) \\[3ex] W' \ \simeq\ T_{\ell}(T'_o(\overline{k})) \ \ \simeq\ \check{M}' \otimes \mathbb{Z}_{\ell}(1) \qquad , \end{array} \right.
\end{equation}
$M$ et $M'$ désignant les groupes des caractères relativement à $\overline{k}$ de $T_o$ et de $T'_o$ respectivement :

}{F7I-IX-097}\begin{equation}\tag{9.2.8}
T_o = D(M) \quad ,\quad T'_o = D(M') \qquad .
\end{equation}
Par suite, l'élément $u$ de $(9.2.6)$, qui  exprime l'opération du groupe d'inertie $I$ sur le module de Tate $U$ de $A_K$ , s'interprète comme un élément


\SGApage{I}{421}

canonique de $(\check{M}'\otimes_{\mathbb{Z}}\check{M})\otimes_{\mathbb{Z}}\mathbb{Z}_\ell$ , d'ailleurs invariant par $\pi_o=\pi/I=\mathrm{Gal}(\bar{k}/k)$, puisque par définition il est invariant par transport de structure :

\SGAvar{\begin{equation}\tag{9.2.9}
u \in (\check{M}'\otimes_{\mathbb{Z}}\check{M})\otimes_{\mathbb{Z}}\mathbb{Z}_\ell)^I \;\simeq\; \mathrm{Hom}_{\mathbb{Z}}(\underline{M}\otimes\underline{M}',\mathbb{Z}_\ell) \;\simeq\; \mathrm{Hom}_{\mathbb{Z}_\ell}(\underline{M}_\ell\otimes\underline{M}'_\ell,\mathbb{Z}_\ell) \quad .
\end{equation}
C'est (à la notation $u$ au lieu de $u_\ell$ près) l'élément (9.1.2) que nous voulions définir.

}{\begin{equation}\tag{9.2.9}
u \in ((\check{M}'\otimes_{\mathbb{Z}}\check{M})\otimes_{\mathbb{Z}}\mathbb{Z}_\ell)^{\pi_o} \;\simeq\; \mathrm{Hom}_{\mathbb{Z}}(\underline{M}\otimes\underline{M}',\mathbb{Z}_\ell) \;\simeq\; \mathrm{Hom}_{\mathbb{Z}_\ell}(\underline{M}_\ell\otimes\underline{M}'_\ell,\mathbb{Z}_\ell) \quad .
\end{equation}
C'est (à la notation $u$ au lieu de $u_\ell$ près) l'élément (9.1.2) que nous voulions définir.

}{F7I-IX-098}Pour définir (9.1.2) dans le cas général, nous aurons besoin de quelques préliminaires, que nous allons développer dans les deux sections suivantes .


9.3. \underline{Extensions panachées.}

Soient $\underline{C}$ une catégorie abélienne, et

\SGAvar{\begin{equation}\tag{9.3.1}
P,\ R,\ Q
\end{equation}
trois objets de $\underline{C}$. On s'intéresse à la classification des objets $X$ de $\underline{C}$, munie d'une filtration en trois crans

}{\begin{equation}\tag{9.3.1}
P,\ R,\ Q
\end{equation}
trois objets de $\underline{C}$. On s'intéresse à la classification des objets $X$ de $\underline{C}$, munis d'une filtration en trois crans

}{F7I-IX-099}\begin{equation}\tag{*}
X^o = X \supset X^1 \supset X^2 \supset X^3 = 0 \quad ,
\end{equation}
avec des isomorphismes donnés

\SGAvar{\begin{equation*}
X^o/X^1 \xrightarrow{\ \sim\ } P \quad , \quad X^1/X^2 \xrightarrow{\ \sim\ } R, \quad X^2/X^3 = X^2 \xrightarrow{\ \sim\ } Q \quad .
\end{equation*}
Posant

}{\begin{equation*}
X^o/X^1 \xrightarrow{\ \sim\ } Q \quad , \quad X^1/X^2 \xrightarrow{\ \sim\ } R, \quad X^2/X^3 = X^2 \xrightarrow{\ \sim\ } P \quad .
\end{equation*}
Posant

}{F7I-IX-100}\begin{equation*}
E = X^o/X^2 \quad , \quad F = X^1/X^3 = X^1 \quad ,
\end{equation*}
on voit que $X$ donne naissance à deux extensions ordinaires

\begin{equation}\tag{9.3.2}
0 \longrightarrow P \longrightarrow F \longrightarrow R \longrightarrow 0 \quad , \quad 0 \longrightarrow R \longrightarrow E \longrightarrow Q \longrightarrow 0 \quad ,
\end{equation}
de $R$ par $P$ et de $Q$ par $R$ respectivement. Notre point de vue sera ici


\SGApage{I}{422}

\SGAvar{de partir, en plus des objets (9.3.1), de structures d'extensions (9.3.2) supposées également données, et de considérer les structures filtrées à trois crans $X$ correspondantes. De façon précise, nous appelerons \underline{extension panachée} de l'extension $E$ par l'extension $F$ un objet $X$ de $\underline{C}$, muni d'une filtration à trois crans (*), et d'isomorphismes donnés

}{de partir, en plus des objets (9.3.1), de structures d'extensions (9.3.2) supposées également données, et de considérer les structures filtrées à trois crans $X$ correspondantes. De façon précise, nous appellerons \underline{extension panachée} de l'extension $E$ par l'extension $F$ un objet $X$ de $\underline{C}$, muni d'une filtration à trois crans (*), et d'isomorphismes donnés

}{F7I-IX-101}\begin{equation}\tag{9.3.3}
E \simeq X^o/X^2 \quad , \quad F \simeq X^1/X^3 = X^1 \quad ,
\end{equation}
compatibles aux filtrations i.e. transformant respectivement $X^1/X^2$ en $\mathrm{Im}(R \longrightarrow E)$ et $X^2$ en $\mathrm{Im}(P,F)$,et tels enfin que l'isomorphisme $R \simeq F/P$ ($\simeq X^1/X^2$) déduit de (9.3.3) soit égal à celui déduit de la première suite exacte (9.3.2). Pour $P,Q,R$ et les extensions $E$ et $F$ fixées, les extensions panachées de $E$ par $F$ forment une catégorie de façon évidente, qui est en fait un groupoïde (toute flèche y est inversible). Le groupe des automorphismes d'une extension panachée $X$ de $E$ par $F$ est isomorphe canoniquement au groupe $\mathrm{Hom}(Q,P)$,grâce à l'isomorphisme

\begin{equation}\tag{9.3.4}
\mathrm{Hom}(Q,P) \xrightarrow{\ \sim\ } \mathrm{Aut}_{\mathrm{extpan}}(X) \quad , \quad u \longmapsto \mathrm{id}_X + \bar{u} \quad ,
\end{equation}
où $\bar{u}$ est l'endomorphisme de l'objet de $\underline{C}$ sous-jacent à $X$ composé de $X \longrightarrow Q \xrightarrow{\ u\ } P \longrightarrow X$, où les flèches extrêmes sont l'épimorphisme canonique resp. le monomorphisme canonique provenant de la structure d'extension panachée donnée sur $X$.

Il y aurait lieu de préciser les structures de la catégorie $\mathrm{EXTPAN}(E,F)$ des extensions panachées de $E$ par $F$, en considérant celle-ci comme un``torseur'' sous la ``catégorie-groupe'' (VII 2.5) $\mathrm{EXT}(Q,P)$ des


\SGApage{I}{423}

extensions de $Q$ par $P$, à l'aide d'un foncteur naturel (analogue à la composition de Brauer des extensions ordinaires)

\begin{equation}\tag{9.3.5}
\omega : \mathrm{EXT}(Q,P) \ \times \ \mathrm{EXTPAN}(E,F) \longrightarrow \mathrm{EXTPAN}(E,F), \quad \omega(G,X) = G \wedge X \ ,
\end{equation}
qui est tel que pour tout objet $X$ de $\mathrm{EXTPAN}(E,F)$, le foncteur partiel $G \longrightarrow G \wedge X$ est une équivalence de $\mathrm{EXT}(Q,P)$ dans $\mathrm{EXTPAN}(E,F)$. Contentons-nous de donner la description sommaire du foncteur $\omega$ . Pour ceci, regardons une extension panachée $X$ de $E$ par $F$ comme définissant une extension ordinaire (également notée $X$) de $Q$ par $F$ (considéré comme objet de $\underline{C}$), et pour toute extension $G$ de $Q$ par $P$, désignons par $\overline{G}$ l'extension de $Q$ par $F$ qu'elle définit, grâce à l'inclusion $P \longrightarrow F$. On pose alors

\SGAvar{\begin{equation}\tag{9.3.6}
G \wedge X = \omega(G,X) \;\overset{\mathrm{dfn}}{=}\; \overline{G} \wedge X \ ,
\end{equation}
où le signe $\wedge$ désigne la composition de Brauer dans la catégorie des extensions de $Q$ par $F$. La filtration en trois crans du deuxième membre $X'$ de (9.3.6) est claire, ainsi que le deuxième des isomorphismes du type (9.3.3) associé à $X'$ ; pour définir le premier isomorphisme du type (9.3.3) pour $X'$, on note que le foncteur $Y \longmapsto \widetilde{Y}$, $\mathrm{EXT}(Q,F) \longrightarrow \mathrm{EXT}(Q,R)$ déduit du morphisme donné $F \longrightarrow R$ (9.3.2) est compatible avec la composition de Brauer, donc transforme $X'$ en un $\widetilde{X}'$ canoniquement isomorphe à $\widetilde{\overline{G}} \wedge \widetilde{\overline{X}}$ . Or le premier isomorphisme (9.3.3) nous donne un isomorphisme $\widetilde{\overline{X}} \simeq F$, tandis qu'il est clair qu'on a un isomorphisme canonique de $\widetilde{\overline{G}} =$ (extension déduite de $G$ par le composé $P \longrightarrow F \longrightarrow R$) avec l'extension triviale, puisque le composé envisagé $P \longrightarrow R$ est nul. Il en résulte un isomorphisme canonique $\widetilde{X}' \simeq F$, ce qui précise la structure d'extension


}{\begin{equation}\tag{9.3.6}
G \wedge X = \omega(G,X) \;\overset{\mathrm{dfn}}{=}\; \overline{G} \wedge X \ ,
\end{equation}
où le signe $\wedge$ désigne la composition de Brauer dans la catégorie des extensions de $Q$ par $F$. La filtration en trois crans du deuxième membre $X'$ de (9.3.6) est claire, ainsi que le deuxième des isomorphismes du type (9.3.3) associé à $X'$ ; pour définir le premier isomorphisme du type (9.3.3) pour $X'$, on note que le foncteur $Y \longmapsto \widetilde{Y}$, $\mathrm{EXT}(Q,F) \longrightarrow \mathrm{EXT}(Q,R)$ déduit du morphisme donné $F \longrightarrow R$ (9.3.2) est compatible avec la composition de Brauer, donc transforme $X'$ en un $\widetilde{X}'$ canoniquement isomorphe à $\widetilde{\overline{G}} \wedge \widetilde{\overline{X}}$ . Or le premier isomorphisme (9.3.3) nous donne un isomorphisme $\widetilde{\overline{X}} \simeq E$, tandis qu'il est clair qu'on a un isomorphisme canonique de $\widetilde{\overline{G}} =$ (extension déduite de $G$ par le composé $P \longrightarrow F \longrightarrow R$) avec l'extension triviale, puisque le composé envisagé $P \longrightarrow R$ est nul. Il en résulte un isomorphisme canonique $\widetilde{X}' \simeq E$, ce qui précise la structure d'extension


}{F7I-IX-102}\SGApage{I}{424}

panachée de $E$ par $F$ sur le deuxième membre de (9.3.6). Nous laissons au lecteur le soin de définir $\omega$ sur les flèches de $\mathrm{EXT}(Q,P) \times \mathrm{EXTPAN}(E,F)$, de vérifier que le foncteur (9.3.5) a bien la propriété annoncée, savoir que pour toute extension panachée \underline{fixée} $X$, le foncteur $G \longmapsto G \wedge X$ est une équivalence de $\mathrm{EXT}(Q,P)$ avec $\mathrm{EXTPAN}(E,F)$, et de prouver la formule d'associativité

\begin{equation}\tag{9.3.7}
G_1 \wedge (G_2 \wedge X) \ \simeq \ (G_1 \wedge G_2) \wedge X \ ,
\end{equation}
où $G_1 \wedge G_2$ désigne la composition de Brauer dans $\mathrm{EXT}(Q,P)$.

De ce qui précède, on déduit aussitôt les parties a) et b) de la

\underline{Proposition} 9.3.8. \underline{Etant donné deux extensions} $E$ \underline{et} $F$ (9.3.2) \underline{dans la catégorie abélienne} $\mathbf{C}$, \underline{on considère la catégorie} $\mathrm{EXTPAN}(E,F)$ \underline{des extensions panachées de} $E$ \underline{par} $F$. \underline{On a alors ce qui suit :}

a) $\mathrm{EXTPAN}(E,F)$ \underline{est un groupoïde. Pour tout objet de celui-ci, le groupe des automorphismes est canoniquement isomorphe à} $\mathrm{Ext}^{0}(Q,P) = \mathrm{Hom}(Q,P)$.

b) \underline{L'ensemble} $\mathrm{Extpan}(E,F)$ \underline{des classes à isomorphisme près d'objets de} $\mathrm{EXTPAN}(E,F)$ \underline{est vide, ou est de façon naturelle un torseur sous le groupe} $\mathrm{Ext}^{1}(Q,P)$.

c) \underline{On peut de façon canonique définir une classe}

\begin{equation}\tag{9.3.9}
c(E,F) \ \in \ \mathrm{Ext}^{2}(Q,P) \ ,
\end{equation}
\underline{dont l'annulation est nécessaire et suffisante pour que l'on ait} $\mathrm{Extpan}(E,F) \neq \emptyset$ , \underline{i.e. pour qu'il existe une extension panachée de} $E$ \underline{par} $F$.


\SGApage{I}{425}

\SGAvar{Il reste à prouver d). Pour ceci on note qu'une extension panachée de $E$ par $F$ peut aussi être décrit comme une extension $X$ de $E$ par $P$, munie d'un isomorphisme d'extensions de $R$ par $P$

\begin{equation*}
X^1 \;\simeq\; F \qquad ,
\end{equation*}
où $X^1$ désigne l'image de l'extension $X$ par le morphisme $R \longrightarrow E$ donné dans (9.3.2). Il existe donc une telle extension panachée si et seulement si il existe un élément de $\mathrm{Ext}^1(E,P)$ donc l'image dans $\mathrm{Ext}^1(R,P)$ est la classe $\xi$ de l'extension $F$. Considérant alors la suite exacte des $\mathrm{Ext}^i( - ,P)$ associée à la deuxième suite exacte (9.3.2) :

\begin{equation*}
.... \longrightarrow \mathrm{Ext}^1(E,P) \longrightarrow \mathrm{Ext}^1(R,P) \stackrel{\partial}{\longrightarrow} \mathrm{Ext}^2(R,Q) \ ... \qquad ,
\end{equation*}
il suffit de définir $c(E,F)$ par la formule

}{Il reste à prouver c). Pour ceci on note qu'une extension panachée de $E$ par $F$ peut aussi être décrite comme une extension $X$ de $E$ par $P$, munie d'un isomorphisme d'extensions de $R$ par $P$

\begin{equation*}
X^1 \;\simeq\; F \qquad ,
\end{equation*}
où $X^1$ désigne l'image de l'extension $X$ par le morphisme $R \longrightarrow E$ donné dans (9.3.2). Il existe donc une telle extension panachée si et seulement si il existe un élément de $\mathrm{Ext}^1(E,P)$ dont l'image dans $\mathrm{Ext}^1(R,P)$ est la classe $\xi$ de l'extension $F$. Considérant alors la suite exacte des $\mathrm{Ext}^i( - ,P)$ associée à la deuxième suite exacte (9.3.2) :

\begin{equation*}
.... \longrightarrow \mathrm{Ext}^1(E,P) \longrightarrow \mathrm{Ext}^1(R,P) \stackrel{\partial}{\longrightarrow} \mathrm{Ext}^2(Q,P) \ ... \qquad ,
\end{equation*}
il suffit de définir $c(E,F)$ par la formule

}{F7I-IX-103}\begin{equation}\tag{9.3.10}
c(E,F) = \partial \xi \quad , \quad \mbox{où } \xi = c\ell(F) \in \mathrm{Ext}^1(R,P) \qquad .
\end{equation}
Supposons maintenant que nous ayons un foncteur exact

\begin{equation}\tag{9.3.11}
r \,:\, \underline{C} \longrightarrow \underline{C}{}' \qquad ,
\end{equation}
où $\underline{C}{}'$ est une deuxième catégorie abélienne. Désignons par $P',Q',R',E',F'$ les objets de $\underline{C}{}'$ déduits de $P,Q,R,E,F$ en appliquant le foncteur $r$, de sorte que $E'$ est une extension de $Q'$ par $R'$, et $F'$ une extension de $R'$ par $P'$. Il est alors clair qu'on a commutativité à isomorphisme canonique près dans le diagramme de foncteurs

\SGAvar{\begin{equation}\tag{9.3.12}
\begin{tikzcd}
\mathrm{EXT}(Q,P) \times \mathrm{EXTPAN}(E,F) \arrow[r, "\omega"] \arrow[d] \& \mathrm{EXTPAN}(E,F) \arrow[d] \\
\mathrm{EXT}(Q',P') \times \mathrm{EXTPAN}(E',F') \arrow[r, "\omega"] \& \mathrm{EXTPAN}(E',F') \quad ,
\end{tikzcd}
\end{equation}














}{\begin{equation}\tag{9.3.12}
\begin{tikzcd}
\mathrm{EXT}(Q,P) \times \mathrm{EXTPAN}(E,F) \arrow[r, "\omega"] \arrow[d] \& \mathrm{EXTPAN}(E,F) \arrow[d] \\
\mathrm{EXT}(Q',P') \times \mathrm{EXTPAN}(E',F') \arrow[r, "\omega'"] \& \mathrm{EXTPAN}(E',F') \quad ,
\end{tikzcd}
\end{equation}














}{F7I-IX-104}\SGApage{I}{426}

où les flèches verticales désignent les foncteurs évidents induits par le foncteur exact $r$, et où $\omega'$ est la flèche analogue à $\omega$ (9.3.5), avec $C$ remplacé par $C'$. On en déduit des compatibilités analogues des structures envisagées dans 9.3.8 avec le foncteur $r$, que nous laissons au soin du lecteur à expliciter.

Nous utiliserons la conséquence suivante de cette commutativité. Supposons que $\mathrm{EXTPAN}(E,F)$ ne soit pas vide, et choisissons alors un objet $X$ dans $\mathrm{EXTPAN}(E,F)$. Alors grâce à $G' \longmapsto G' \wedge X'$ (où $X'=r(X)$), la catégorie $\mathrm{EXTPAN}(E',F')$ est équivalente à la catégorie $\mathrm{EXT}(Q',P')$, donc un élément $Y'$ de $\mathrm{EXTPAN}(E',F')$ est connu à isomorphisme près quand on connaît la classe d'isomorphie de l'élément correspondant de $\mathrm{EXT}(Q',P')$, i.e. quand on connaît l'élément correspondant de $\mathrm{Ext}^1(Q' ,P')$, soit

\begin{equation*}
c_X(Y') \in \mathrm{Ext}^1(Q',P') \qquad .
\end{equation*}
Lorsqu'on remplace $X$ par un autre objet $X_1 = H \wedge X$, où $H$ est une extension de $Q$ par $P$, on voit aussitôt que l'on obtient

\SGAvar{\begin{equation*}
c_{X_1}(Y') \;=\; c_X(Y) - \mathrm{cl}(H') \qquad ,
\end{equation*}
grâce à la formule d'associativité (9.3.7). Par suite, pour $Y'$ fixé, et $X$ variable, les $c_X(Y')$ parcourent exactement une classe

}{\begin{equation*}
c_{X_1}(Y') \;=\; c_X(Y') - \mathrm{cl}(H') \qquad ,
\end{equation*}
grâce à la formule d'associativité (9.3.7). Par suite, pour $Y'$ fixé, et $X$ variable, les $c_X(Y')$ parcourent exactement une classe

}{F7I-IX-105}\begin{equation}\tag{9.3.13}
c(Y') \in \mathrm{Ext}^1(Q',P')/\ \mathrm{Im}\ \mathrm{Ext}^1(Q,P) \qquad ,
\end{equation}
qui est donc canoniquement associée à l'extension panachée $Y'$ de $E'$ par $F'$. Par construction, sa nullité est nécessaire et suffisante pour que l'extension panachée $Y'$ soit isomorphe à la transformée par $r$ d'une extension panachée $X$ de $E$ par $F$.


\SGApage{I}{427}

9.4. Soient $S$ un trait hensélien, de point générique $\eta$, $n$ un entier $> 0$, $\Lambda = \mathbb{Z}/n\mathbb{Z}$ , $\underline{C}$ la catégorie des $\Lambda$-Modules sur $S$ (pour le site fppf sur $S$, cf. VIII 0.2), $\underline{C}'$ la catégorie des $\Lambda$-Modules sur $\eta$ . Soient $P,R,Q$ des objets de $\underline{C}$, et $E,F$ des extensions de $Q$ par $R$ et de $R$ par $P$ respectivement (9.3.2). On suppose que $Q$ est un Module localement libre de type fini sur $\Lambda_S$, et que $P$ soit représentable par un groupe de type multiplicatif fini sur $S$ (SGA 3 IX), donc qu'il soit de la forme

\begin{equation*}
P = D(Q^*) \quad ,
\end{equation*}
où $D = \underline{\mathrm{Hom}}( -, \mathbb{G}_{mS})$ désigne la dualité de Cartier. On considère les restrictions de $P,R,Q,E,F$ à $\eta$, soient $P_\eta$, $R_\eta$ ...., et on suppose donnée une extension panachée (9.3) $X_\eta$ de $E_\eta$ par $F_\eta$. Sous ces conditions, nous allons définir canoniquement un accouplement

\begin{equation}\tag{9.4.1}
c(X_\eta) : Q \otimes_\Lambda Q^* \longrightarrow \Lambda_S \quad (\Lambda = \mathbb{Z}/n\mathbb{Z}) \quad .
\end{equation}
Pour ceci, \underline{supposons d'abord} $S$ \underline{strictement local}, de sorte que $Q$ et $Q^*$ sont constants, soient $Q = Q_{oS}$ , $Q^* = Q^*_{oS}$, et la donnée de (9.4.1) revient à celle d'un élément

\begin{equation}\tag{9.4.2}
c(X_\eta) \in \check{Q}_o \otimes \check{Q}^*_o \quad ,
\end{equation}
où le signe $\vee$ désigne le dual ordinaire $\mathrm{Hom}( - , \Lambda)$. Pour ceci, nous appliquons les considérations de9.3, en prenant pour $r: \underline{C} \longrightarrow \underline{C}'$ le foncteur restriction. Notons d'abord qu'il existe une extension panachée de $E$ par $F$ : en effet, l'obstruction à l'existence d'une telle extension se trouve dans $\mathrm{Ext}^2(Q,P)$ (9.3.8 c)), et il suffit de prouver qu'on a

\begin{equation*}
\mathrm{Ext}^2(Q,P) = 0 \quad .
\end{equation*}

\SGApage{I}{428}

Ceci est un cas particulier du

\underline{Lemme} 9.4.3.  \underline{Sous les conditions précédentes, avec} $S$ \underline{strictement local,} \underline{on a des isomorphismes}

\begin{equation*}\tag{a)}
\mathrm{Ext}^i_\Lambda(Q,P) \simeq H^i(S,\check{Q}\otimes P) \ , \quad \mathrm{Ext}^i_\Lambda(Q_\eta,P_\eta) \simeq H^i(\eta,\check{Q} \otimes P) \quad ,
\end{equation*}
\underline{d'autre part on a un isomorphisme canonique}

\begin{equation}\tag{b)}
\check{Q}\otimes P \simeq D(Q\otimes Q^*) \quad ,
\end{equation}
\underline{enfin,} \underline{pour tout groupe constant fini} $H = H_{oS}$ \underline{sur} $S$, \underline{on a}

\SGAvar{\begin{equation}\tag{c)}
H^i(S,D(H) = H^i(\eta,D(H)_\eta) = 0 \quad \underline{\text{pour}}\ i \geq 2 \quad ,
\end{equation}
\underline{et des isomorphismes canoniques}

\begin{equation}\tag{d)}
H^1(S,D(H)) \simeq G_m(S)\otimes \check{H}_o = V^*\otimes \check{H}_o \ , \quad H^1(\eta,D(H_\eta)) \simeq G_m(\eta) \otimes H_o = K^*\otimes \check{H}_o \quad ,
\end{equation}
\underline{où le signe} $\vee$ \underline{désigne ici} $\mathrm{Hom}( - , Q/\mathbb{Z} )$.

Les formules a) et b) résultent trivialement du fait que $Q$ est un $\Lambda_S$-Module localement libre de type fini. Pour prouver c), on est évidemment ramené au cas où $H_o = \mathbb{Z}/m\mathbb{Z}$ i.e. $D(H) = \mu_m$ , et alors la auite exacte de cohomologie associée à la suite exacte de Kummer

}{\begin{equation}\tag{c)}
H^i(S,D(H)) = H^i(\eta,D(H)_\eta) = 0 \quad \underline{\text{pour}}\ i \geq 2 \quad ,
\end{equation}
\underline{et des isomorphismes canoniques}

\begin{equation}\tag{d)}
H^1(S,D(H)) \simeq G_m(S)\otimes \check{H}_o = V^*\otimes \check{H}_o \ , \quad H^1(\eta,D(H_\eta)) \simeq G_m(\eta) \otimes \check{H}_o = K^*\otimes \check{H}_o \quad ,
\end{equation}
\underline{où le signe} $\vee$ \underline{désigne ici} $\mathrm{Hom}( - , Q/\mathbb{Z} )$.

Les formules a) et b) résultent trivialement du fait que $Q$ est un $\Lambda_S$-Module localement libre de type fini. Pour prouver c), on est évidemment ramené au cas où $H_o = \mathbb{Z}/m\mathbb{Z}$ i.e. $D(H) = \mu_m$ , et alors la suite exacte de cohomologie associée à la suite exacte de Kummer

}{F7I-IX-106}\begin{equation*}
0 \longrightarrow \mu_m \longrightarrow \mathbb{G}_m \overset{m.\mathrm{id}}{\longrightarrow} \mathbb{G}_m \longrightarrow 0
\end{equation*}
nous ramène à prouver les formules

\begin{equation*}
H^i(S,\mathbb{G}_m) = H^i(\eta,\mathbb{G}_m) = 0 \quad \text{pour}\ i \geq 1 \quad .
\end{equation*}
Les $H^i$ sont égaux aux $H^i_{\text{ét}}$ calculés pour le site étale [ 6 , th. 11.7], ils sont donc nuls pour $S$ puisque $S$ est strictement local. Le fait qu'ils sont nuls aussi pour $\eta$ est un fait bien connu [ 6 , th. 1.1].


\SGApage{I}{429}

Reste à décrire des isomorphismes d). En fait, sur une base quelconque $U$, on a un homomorphisme canonique

\begin{equation}\tag{9.4.3.1}
\mathbb{G}_m(U) \otimes \check{H}_o \longrightarrow H^1(U,D(H_{oU})) \quad ,
\end{equation}
qu'on décrit, lorsque $H_o$ est annulé par l'entier $m \geqslant 0$, à l'aide de l'homomorphisme cobord $\mathbb{G}_m(U) \stackrel{\partial}{\longrightarrow} H^1(U,\mu_m)$ associé à la suite exacte de Kummer, en prenant le composé

\SGAvar{\begin{equation*}
\mathbb{G}_m(U) \otimes \check{H}_o \longrightarrow H^1(U ,\mu_m) \otimes \check{H}^o \longrightarrow H^1(U,\mu_m \otimes \check{H}_o) \longrightarrow H^1(U, D(H_{oS})) \quad ,
\end{equation*}
où la deuxième flèche se déduit de l'homomorphisme évident $\check{H}_o \longrightarrow \mathrm{Hom}(\mu_m, \mu_m \otimes \check{H}_o)$, et la troisième est un cas particulier de b) (où on remplace $n$ par $m$). Nous laissons au lecteur le soin de vérifier que (9.4.3.1) ne dépend pas du choix de l'entier $m$ tel que $mH_o = 0$ . D'autre part je dis que l'homomorphisme (9.4.3.1) est toujours injectif, et que son conoyau est nul si $\mathrm{Pic}(U) = 0$. (En fait, le conoyau est toujours isomorphe à $\mathrm{Hom}(H_o,\mathrm{Pic}(U))$, mais peu importe ici.) Pour voir ceci, on est ramené en effet au cas où $H_o$ est de la forme $\mathbb{Z}/n\mathbb{Z}$ , et alors l'assertion résulte aussitôt de la suite exacte de cohomologie associée à la suite exacte de Kummer (*) ci-dessus. Ceci prouve la formule d). On en conclut aussitôt :

}{\begin{equation*}
\mathbb{G}_m(U) \otimes \check{H}_o \longrightarrow H^1(U ,\mu_m) \otimes \check{H}_o \longrightarrow H^1(U,\mu_m \otimes \check{H}_o) \longrightarrow H^1(U, D(H_{oU})) \quad ,
\end{equation*}
où la deuxième flèche se déduit de l'homomorphisme évident $\check{H}_o \longrightarrow \mathrm{Hom}(\mu_m, \mu_m \otimes \check{H}_o)$, et la troisième est un cas particulier de b) (où on remplace $n$ par $m$). Nous laissons au lecteur le soin de vérifier que (9.4.3.1) ne dépend pas du choix de l'entier $m$ tel que $mH_o = 0$ . D'autre part je dis que l'homomorphisme (9.4.3.1) est toujours injectif, et que son conoyau est nul si $\mathrm{Pic}(U) = 0$. (En fait, le conoyau est toujours isomorphe à $\mathrm{Hom}(H_o,\mathrm{Pic}(U))$, mais peu importe ici.) Pour voir ceci, on est ramené en effet au cas où $H_o$ est de la forme $\mathbb{Z}/n\mathbb{Z}$ , et alors l'assertion résulte aussitôt de la suite exacte de cohomologie associée à la suite exacte de Kummer (*) ci-dessus. Ceci prouve la formule d). On en conclut aussitôt :

}{F7I-IX-107}\underline{Corollaire} 9.4.4. \underline{Avec les notations de} 9.4.3 \underline{on a un isomorphisme canonique}

\SGAvar{\begin{equation*}\tag{e)}
H^1(\eta,D(H)_\eta)/ \mathrm{Im}\, H^1(S,D(H) \quad \simeq \quad \check{H}_o \quad ,
\end{equation*}
\underline{d'où} \underline{(en combinant avec} a) \underline{et} b)) \underline{un isomorphisme canonique}

}{\begin{equation*}\tag{e)}
H^1(\eta,D(H)_\eta)/ \mathrm{Im}\, H^1(S,D(H)) \quad \simeq \quad \check{H}_o \quad ,
\end{equation*}
\underline{d'où} \underline{(en combinant avec} a) \underline{et} b)) \underline{un isomorphisme canonique}

}{F7I-IX-108}\begin{equation}\tag{9.4.5}
\mathrm{Ext}^1(Q_\eta,P_\eta)/\mathrm{Im}\, \mathrm{Ext}^1(Q,P) \simeq \check{Q}_o \otimes \check{Q}_o^* \quad .
\end{equation}

\SGApage{I}{430}

Revenant à la définition de (9.4.2) lorsque $S$ est strictement local, on note que la relation $\mathrm{Ext}^2(Q,P) = 0$ implique la possibilité de définir un élément canonique (9.3.13), où on prend pour $Y'$ l'extension panachée $X_\eta$ donnée de $E_\eta$ par $F_\eta$ . La formule (9.4.5) permet d'identifier cet élément $c(X_\eta)$ à un élément de $\check{Q}_o \otimes \check{Q}_o^*$, ce qui précise la définition de (9.4.2). On trouve de plus, en vertu des considérations générales de 9.3, que l'élément en question s'interprète comme l'obstruction au prolongement de $X_\eta$ en une extension panachée $X$ de $E$ par $F$.

9.4.6. La définition de (9.4.1) dans le cas $S$ hensélien se ramène maintenant immédiatement au cas strictement local, par descente galoisienne à partir du hensélisé strict de $S$.

\underline{Proposition} 9.4.7. \underline{Les données} $P,Q,R,E,F$ \underline{sont comme ci-dessus, sauf qu'il est inutile de supposer le trait} $S$ \underline{hensélien. On considère le foncteur}

\begin{equation}\tag{9.4.7.1}
X \longmapsto X_\eta
\end{equation}
\underline{de la catégorie des extensions panachées de} $E$ \underline{par} $F$ \underline{dans la catégorie des extensions panachées de} $E_\eta$ \underline{par} $F_\eta$. \underline{Pour toute extension panachée} $X_\eta$ \underline{de} $E_\eta$ \underline{par} $F_\eta$, \underline{on considère l'accouplement}

\begin{equation}\tag{9.4.7.2}
c_o(X_\eta) : Q_o \otimes Q_o^* \longrightarrow \Lambda_s
\end{equation}
\underline{induit par l'accouplement du type} (9.4.1) \underline{défini sur le} hensélisé $S^h$ \underline{de} $S$. \underline{Ceci posé, le foncteur} (9.4.7.1) \underline{est} pleinement fidèle, \underline{et pour une extension panachée} $X_\eta$ \underline{de} $E_\eta$ \underline{par} $F_\eta$ , $X_\eta$ \underline{appartient à l'image essentielle de ce foncteur si et seulement si l'accouplement} (9.4.7.2) \underline{est nul.}


\SGApage{I}{431}

9.4.7.3.  Cet énoncé garde un sens et reste valable lorsqu'on suppose seulement que $S$ est un schéma noethérien régulier connexe de dimension 1 : au lieu d'un seul accouplement (9.4.7.2), il faut alors considérer un accouplement $Q_s \otimes Q_s^* \longrightarrow \Lambda_s$ pour chaque point fermé $s$ de $S$. C'est sous cette forme plus générale qu'il nous sera le plus commode de prouver 9.4.7.

a)  Pour toute extension panachée $X$, l'application

\SGAvar{\begin{equation*}
\mathrm{End}(X) \longrightarrow \mathrm{End}(X_\eta)
\end{equation*}
est bijective. Comme les deux membres sont respectivement les sections de $\underline{\mathrm{Hom}}(P,Q) = D(Q \otimes Q^*)$ sur $S$ et sur $\eta$ , l'assertion provient du fait que $S$ est normal et que $D(Q \otimes Q^*)$ est fini sur $S$. De ceci résulte aussitôt que si $X$ et $Y$ sont deux extensions panachées, et si on sait déjà qu'elles sont isomorphes, alors

}{\begin{equation*}
\mathrm{End}(X) \longrightarrow \mathrm{End}(X_\eta)
\end{equation*}
est bijective. Comme les deux membres sont respectivement les sections de $\underline{\mathrm{Hom}}(Q,P) = D(Q \otimes Q^*)$ sur $S$ et sur $\eta$ , l'assertion provient du fait que $S$ est normal et que $D(Q \otimes Q^*)$ est fini sur $S$. De ceci résulte aussitôt que si $X$ et $Y$ sont deux extensions panachées, et si on sait déjà qu'elles sont isomorphes, alors

}{F7I-IX-109}\begin{equation*}
\underline{\mathrm{Hom}}(X,Y) \longrightarrow \underline{\mathrm{Hom}}(X_\eta,\ Y_\eta)
\end{equation*}
est bijectif.

b)  Si $X$ et $Y$ sont deux extensions panachées, tout isomorphisme $X_\eta \xrightarrow{\sim} Y_\eta$ provient d'un isomorphisme $X \longrightarrow Y$. Grâce à a) la question est locale pour la topologie étale, ce qui nous permet de nous ramener au cas où $S$ est strictement local. Mais dans ce cas 9.4.3 nous montre que l'application

\SGAvar{\begin{equation*}
\mathrm{Ext}^1_\Lambda(P,Q) \longrightarrow \mathrm{Ext}^1_\Lambda(P_\eta\ ,\ Q_\eta)
\end{equation*}
s'identifie à l'application

}{\begin{equation*}
\mathrm{Ext}^1_\Lambda(Q,P) \longrightarrow \mathrm{Ext}^1_\Lambda(Q_\eta\ ,\ P_\eta)
\end{equation*}
s'identifie à l'application

}{F7I-IX-110}\begin{equation*}
V^* \otimes \check{H}_o \longrightarrow K^* \otimes \check{H}_o
\end{equation*}


\SGApage{I}{432}

qui est injective (puisque la suite exacte $0 \longrightarrow V^* \longrightarrow K^* \longrightarrow \mathbb{Z} \longrightarrow 0$ est scindée), donc elle est injective, ce qui implique (grâce au sorite 9.3) que si $X_\eta$ et $Y_\eta$ sont isomorphes, il en est de même de $X$ et de $Y$. On gagne donc par a).

Evidemment a) et b) impliquent que le foncteur $X \longmapsto X_\eta$ est \underline{pleinement} fidèle.

\SGAvar{c)  Soit $X_\eta$ une extension panachée sur $\eta$, prouvons qu'elle se prolonge à $S$ si et seulement si les accouplements (9.4.6.2) sont nuls. Ici encore, l'unicité déjà prouvée montre que la question est locale pour la topologie étale, ce qui nous ramène encore au cas $S$ strictement local. Alors la conclusion est vraie par construction, comme on avait signalé juste avant 9.4.6.

}{c)  Soit $X_\eta$ une extension panachée sur $\eta$, prouvons qu'elle se prolonge à $S$ si et seulement si les accouplements (9.4.7.2) sont nuls. Ici encore, l'unicité déjà prouvée montre que la question est locale pour la topologie étale, ce qui nous ramène encore au cas $S$ strictement local. Alors la conclusion est vraie par construction, comme on avait signalé juste avant 9.4.6.

}{F7I-IX-111}Ainsi la démonstration de 9.4.7 est achevée.

9.5.  Supposons maintenant fixé un nombre premier $\ell$, des pro-$\ell$-groupes de Barsotti-Tate sur $S$,

\begin{equation}\tag{9.5.1}
P = (P(n))_{n \geq 0} \ , \quad R = (R(n))_{n \geq 0} \ , \quad Q = (Q(n))_{n \geq 0} \quad ,
\end{equation}
et des extensions de pro-$\ell$-groupes de Barsotti-Tate

\begin{equation}\tag{9.5.2}
0 \longrightarrow P \longrightarrow F \longrightarrow R \longrightarrow 0 \ , \quad 0 \longrightarrow R \longrightarrow E \longrightarrow Q \longrightarrow 0 \quad ,
\end{equation}
définissant donc pour tout entier $n \geq 0$ des extensions de $\Lambda_n$-Modules où $\Lambda_n = \mathbb{Z}/\ell^{n+1}\mathbb{Z}$ :

\begin{equation}\tag{9.5.3}
0 \longrightarrow P(n) \longrightarrow F(n) \longrightarrow R(n) \longrightarrow 0, \quad 0 \longrightarrow R(n) \longrightarrow E(n) \longrightarrow Q(n) \longrightarrow 0 \ .
\end{equation}
Nous supposerons $Q$ de type étale et $P$ de type multiplicatif. Supposons


\SGApage{I}{433}

\SGAvar{enfin donnée sur $\eta$ un pro-$\ell$-groupe de Barsotti-Tate $X_\eta$ qui soit \underline{extension panachée} de $E_\eta$ par $F_\eta$, par quoi nous entendons qu'il est muni d'une filtration en trois crans (en tant que pro-$\ell$-groupe de Barsotti-Tate), donnant lieu à des extensions partielles munies d'isomorphismes avec $E_\eta$ et $F_\eta$ respectivement. Il s'ensuit que pour tout $n$, $X_\eta(n)$ est muni d'une structure d'extension panachée de $E(n)_\eta$ par $F(n)_\eta$. Nous pouvons lui faire correspondre alors un accouplement du type (9.4.1) :

}{enfin donné sur $\eta$ un pro-$\ell$-groupe de Barsotti-Tate $X_\eta$ qui soit \underline{extension panachée} de $E_\eta$ par $F_\eta$, par quoi nous entendons qu'il est muni d'une filtration en trois crans (en tant que pro-$\ell$-groupe de Barsotti-Tate), donnant lieu à des extensions partielles munies d'isomorphismes avec $E_\eta$ et $F_\eta$ respectivement. Il s'ensuit que pour tout $n$, $X_\eta(n)$ est muni d'une structure d'extension panachée de $E(n)_\eta$ par $F(n)_\eta$. Nous pouvons lui faire correspondre alors un accouplement du type (9.4.1) :

}{F7I-IX-112}\begin{equation}\tag{9.5.4}
c(X(n)_\eta) : Q(n) \otimes Q^*(n) \longrightarrow \Lambda_n = \mathbb{Z}/\ell^{n+1}\mathbb{Z} \quad ,
\end{equation}
où les $Q^*(n)$ sont les groupes finis étales sur $S$ définis par

\begin{equation*}
P(n) \simeq D(Q^*(n)) \; , \quad \mathrm{i.e.} \quad Q^*(n) = D(P(n)) \quad ,
\end{equation*}
$D$ désignant encore la dualité de Cartier. Les $Q^*(n)$ forment un pro-$\ell$-groupe de façon évidente, soit $Q^*$. Nous laissons au lecteur le soin de vérifier que pour $n$ variable, les homomorphismes (9.5.4) définissent un homomorphisme de systèmes projectifs, i.e. un homomorphisme de faisceaux $\ell$-adiques constants tordus sur $S$ :

\begin{equation}\tag{9.5.5}
c(X_\eta) : Q \otimes Q^* \longrightarrow \mathbb{Z}_{\ell S} \quad .
\end{equation}
Signalons en passant le résultat suivant (qui ne sera pas utilisé dans la suite) :

\underline{Proposition} 9.5.6. \underline{Les données} $P,Q,R,E,F$ \underline{sont les mêmes que ci-dessus, sauf qu'il est inutile de supposer le trait} $S$ \underline{hensélien. On considère le foncteur}

\begin{equation}\tag{9.5.6.1}
X \longmapsto X_\eta
\end{equation}
\underline{de la catégorie des extensions panachées de Barsotti-Tate de} $E$ \underline{par} $F$


\SGApage{I}{434}

\underline{dans la catégorie des extensions panachées de} $E_\eta$ \underline{par} $F_\eta$ . \underline{Pour tout groupe de Barsotti-Tate extension panachée} $X_\eta$ \underline{de} $E_\eta$ \underline{par} $F_\eta$ , \underline{on considère l'accouplement correspondant}

\begin{equation}\tag{9.5.6.2}
Q_o \otimes Q^*_o \longrightarrow \mathbb{Z}_{\ell s} \quad ,
\end{equation}
\underline{induit par l'accouplement du type} (9.5.5) \underline{défini sur le} hensélisé \underline{de} $S$. \underline{Ceci posé, le foncteur} (9.5.6.1) \underline{est pleinement fidèle, et son image essentielle est formée des groupes de Barsotti-Tate} $X_\eta$ \underline{extensions panachées de} $E_\eta$ \underline{par} $F_\eta$ \underline{telles que l'accouplement} (9.5.6.2) \underline{correspondant soit nul}. \underline{De plus, lorsque} car.$K = 0$ (\underline{de sorte qu'on peut appliquer le théorème de TATE} [33 th. 4]), \underline{alors} $X_\eta$ \underline{appartient à l'image essentielle si et seulement si} $X_\eta$ \underline{se prolonge en un groupe de Barsotti-Tate sur} $S$.

9.5.6.3. Nous laissons au lecteur le soin d'étendre cet énoncé au cas où on suppose seulement que $S$ est un schéma noethérien régulier connexe de dimension 1.

\SGAvar{La première assertion de 9.5.6 est une conséquence immédiate de 9.4.6, appliqué aux composants des groupes de Barsotti-Tate envisagés. Il reste à prouver que si un groupe de Barsotti-Tate $X_\eta$ , extension panachée de $E_\eta$ par $F_\eta$ , se prolonge en un groupe de Barsotti-Tate $X_\eta$ sur $S$, il se prolonge aussi comme extension panachée. Or en vertu du théorème de TATE [33], l'homomorphisme $X_\eta \longrightarrow Q_\eta$ se prolonge de façon unique en un homomorphisme

}{La première assertion de 9.5.6 est une conséquence immédiate de 9.4.7, appliqué aux composants des groupes de Barsotti-Tate envisagés. Il reste à prouver que si un groupe de Barsotti-Tate $X_\eta$ , extension panachée de $E_\eta$ par $F_\eta$ , se prolonge en un groupe de Barsotti-Tate $X$ sur $S$, il se prolonge aussi comme extension panachée. Or en vertu du théorème de TATE [33], l'homomorphisme $X_\eta \longrightarrow Q_\eta$ se prolonge de façon unique en un homomorphisme

}{F7I-IX-113}\begin{equation*}
X \longrightarrow Q \quad .
\end{equation*}

\SGApage{I}{435}

\SGAvar{Pour tout entier $n$, l'homomorphisme correspondant $X(n) \longrightarrow Q(n)$, étant surjectif sur les fibres génériques, est surjectif ; comme $Q(n)$ est étale sur $n$, donc $X(n) \longrightarrow Q(n)$ est plat, il s'ensuit que c'est un épimorphisme. Donc $X \longrightarrow Q$ \underline{est un épimorphisme composant par composant}, Si donc $\overline{F}$ désigne le noyau de $X \longrightarrow Q$, il s'ensuit que $\overline{F}$ est un groupe de Barsotti-Tate et $X$ est une extension de $Q$ par $\overline{F}$ . D'ailleurs l'isomorphisme $F_\eta \xrightarrow{\sim} \overline{F}_\eta$ se prolonge, par le théorème de TATE, en un isomorphisme $F \simeq \overline{F}$ . Par suite, $X$ apparaît comme une extension de $Q$ par $F$. Considérons alors le quotient $\overline{E} = X/P$, on voit encore, par application du théorème de TATE, qu'on a un isomorphisme $E \simeq \overline{E}$ . Il s'ensuit que $X$ est une extension panachée de $E$ par $F$ prolongeant l'extension panachée $X_\eta$ , cqfd.

}{Pour tout entier $n$, l'homomorphisme correspondant $X(n) \longrightarrow Q(n)$, étant surjectif sur les fibres génériques, est surjectif ; comme $Q(n)$ est étale sur $S$, donc $X(n) \longrightarrow Q(n)$ est plat, il s'ensuit que c'est un épimorphisme. Donc $X \longrightarrow Q$ \underline{est un épimorphisme composant par composant}. Si donc $\overline{F}$ désigne le noyau de $X \longrightarrow Q$, il s'ensuit que $\overline{F}$ est un groupe de Barsotti-Tate et $X$ est une extension de $Q$ par $\overline{F}$ . D'ailleurs l'isomorphisme $F_\eta \xrightarrow{\sim} \overline{F}_\eta$ se prolonge, par le théorème de TATE, en un isomorphisme $F \simeq \overline{F}$ . Par suite, $X$ apparaît comme une extension de $Q$ par $F$. Considérons alors le quotient $\overline{E} = X/P$, on voit encore, par application du théorème de TATE, qu'on a un isomorphisme $E \simeq \overline{E}$ . Il s'ensuit que $X$ est une extension panachée de $E$ par $F$ prolongeant l'extension panachée $X_\eta$ , cqfd.

}{F7I-IX-114}\underline{Remarques} 9.5.6.4.  a)  Bien entendu, la dernière assertion de 9.5.6 devrait être valable sans restriction sur $K$, puisqu'il devrait en être ainsi du théorème de Tate invoqué.

\SGAvar{b)  On peut donner une autre démonstration du critère de bonne réduction 5.10, en se ramenant comme dans 5.10 au cas où $A_K$ est à réduction semi-stable, et en appliquant alors 9.5.6 à $T_\ell(A_K)$ considérée comme extension panachée (cf. 9.6), qui implique $u_\ell = 0$, et utilisant 11.6.2 a) plus bas qui implique que l'on a $T_\ell(A_K)^t = 0$.

}{b)  On peut donner une autre démonstration du critère de bonne réduction 5.10, en se ramenant comme dans 5.10 au cas où $A_K$ est à réduction semi-stable, et en appliquant alors 9.5.6 à $T_\ell(A_K)$ considérée comme extension panachée (cf. 9.6), qui implique $u_\ell = 0$, et utilisant 11.5.2 a) plus bas qui implique que l'on a $T_\ell(A_K)^t = 0$.

}{F7I-IX-115}9.5.8.  Lorsque $\ell \neq p$, les groupes de Barsotti-Tate envisagés sur $S$ resp. sur $K$ s'interprètent simplement en termes de faisceaux $\ell$-adiques constants tordus sur $S$, ou de façon équivalente, en termes de représentations de $\pi_o = \mathrm{Gal}(\overline{k}/k)$ resp. de $\pi = \mathrm{Gal}(\overline{K}/K)$ sur des $\mathbb{Z}_\ell$-modules libres


\SGApage{I}{436}

\SGAvar{de type fini. On constate alors que l'homomorphisme (9.5.6) est l'homomorphisme canonique qui exprime l'action du sous-groupe d'inertie $I \subset \pi$ sur $X_\eta(\overline{K})$, défini par la méthode de 9.2. Pour le vérifier, on se ramène immédiatement à la situation de 9.4, en y supposant  que $n$ est égal à une puissance de $\ell$, et que $R$ est un $\Lambda$-Module localement constant (tout comme $Q$ et $P$), de sorte qu'il en est également de même de $E$, $F$ et de $X_\eta$ . La vérification de l'identité des deux définitions est alors laissée au lecteur.

}{de type fini. On constate alors que l'homomorphisme (9.5.5) est l'homomorphisme canonique qui exprime l'action du sous-groupe d'inertie $I \subset \pi$ sur $X_\eta(\overline{K})$, défini par la méthode de 9.2. Pour le vérifier, on se ramène immédiatement à la situation de 9.4, en y supposant  que $n$ est égal à une puissance de $\ell$, et que $R$ est un $\Lambda$-Module localement constant (tout comme $Q$ et $P$), de sorte qu'il en est également de même de $E$, $F$ et de $X_\eta$ . La vérification de l'identité des deux définitions est alors laissée au lecteur.

}{F7I-IX-116}9.6.  Nous sommes maintenant en mesure de définir l'accouplement canonique annoncé (9.1.2), sans restriction sur $\ell$, de telle façon que pour $\ell \neq p$ cette définition coïncide avec celle donnée dans 9.2. Pour ceci, on applique les considérations de 9.5 à la situation suivante :

\begin{equation}\tag{9.6.1}
\left\{
\begin{aligned}
&P = T_\ell(A^o)^t = \check{\underline{M}}{}'_{\ell}(1) \ , \quad Q = M_\ell \ , \quad R = T_\ell(A^o)^f / T_\ell(A^o)^t \qquad , \\[4ex]
&\qquad\quad F = T_\ell(A^o)^f \ , \quad E = D(F') = D(T_\ell({A'}^o)^f) \qquad ,
\end{aligned}
\right.
\end{equation}
où dans la dernière formule $D$ désigne le dual de Cartier au sens des pro-$\ell$-groupes de Barsotti-Tate, i.e $D(Y)$ est le système \underline{projectif} des $D(Y(n))$, pour les morphismes de transition habituels. La structure de $F$ comme extension de $R$ par $P$ est claire. On a de même sur $F' = T_\ell({A'}^o)^f$ une structure d'extension de $R'$ par $P' = T_\ell({A'}^o)^t = \check{\underline{M}}_\ell(1)$. Appliquant la dualité de Cartier, on en conclut que $E = D(F')$ est muni d'une structure d'extension de $D(P') = M_\ell$ par $D(R')$. Or $D(R')$ est canoniquement isomorphe à $R$ en vertu de 5.5 (prouvé dans 7.4.6), ce qui précise la structure de $E$ comme extension de $Q$ par $R$.


\SGApage{I}{437}

Comme extension panachée $X_\eta$ de $E_\eta$ par $F_\eta$ , on prendra

\begin{equation}\tag{9.6.2}
X_\eta = T_\ell(A_K) = X^0_\eta \supset X^1_\eta = T_\ell(A_K)^f \supset X^2_\eta = T_\ell(A_K)^t \supset 0 \quad .
\end{equation}
L'isomorphisme d'extensions de $R_\eta$ par $P_\eta$

\begin{equation*}
X^1_\eta \simeq F_\eta
\end{equation*}
est clair. L'isomorphisme d'extensions de $Q$ par $R$ :

\SGAvar{\begin{equation*}
X_\eta/X^2_\eta \simeq F_\eta = D(F'_\eta)
\end{equation*}
est clair grâce au théorème d'orthogonalité 5.2.

Nous disposons donc des données requises pour définir un accouplement (9.5.6), qui ici (comme $Q^* = \underline{M}'_\ell$) prend la forme (9.1.2). Le fait que cette définition soit compatible avec celle de 9.2 résulte de 9.5.8. Les considérations de 5.8 montrent d'autre part, sous réserve de la validité du théorème de TATE [33] (donc en tous cas dans le cas où car $K = 0$) que l'accouplement en question peut se décrire en termes du seul $\ell$-pro-groupe de Barsotti-Tate $X_\eta = T_\ell(A_K)$ sur $S$ (puisque sa filtration en trois crans, et les extensions $F$ et $E$ qui prolongent $X^1_\eta$ et $X_\eta/X^2_\eta$ , peuvent se définir canoniquement en termes du seul pro-$\ell$-groupe $X = T_\ell(A_K)$.

}{\begin{equation*}
X_\eta/X^2_\eta \simeq E_\eta = D(F'_\eta)
\end{equation*}
est clair grâce au théorème d'orthogonalité 5.2.

Nous disposons donc des données requises pour définir un accouplement (9.5.5), qui ici (comme $Q^* = \underline{M}'_\ell$) prend la forme (9.1.2). Le fait que cette définition soit compatible avec celle de 9.2 résulte de 9.5.8. Les considérations de 5.8 montrent d'autre part, sous réserve de la validité du théorème de TATE [33] (donc en tous cas dans le cas où car $K = 0$) que l'accouplement en question peut se décrire en termes du seul $\ell$-pro-groupe de Barsotti-Tate $X_\eta = T_\ell(A_K)$ sur $\eta$ (puisque sa filtration en trois crans, et les extensions $F$ et $E$ qui prolongent $X^1_\eta$ et $X_\eta/X^2_\eta$ , peuvent se définir canoniquement en termes du seul pro-$\ell$-groupe $X = T_\ell(A_K)$).

}{F7I-IX-117}\underline{Remarque} 9.7. Ne supposons plus nécessairement que $A_K$ ait réduction semi-stable sur le trait hensélien $S$, et soit $K'$ une extension galoisienne finie de $K$ telle que $A_{K'}$ ait réduction semi-stable sur le normalisé $S'$ de $S$ dans $K'$ (3.6). Avec les notations de (8.1.8), on définit alors, par descente galoisienne à partir de l'accouplement de monodromie pour $A_{K'}$ ,


\SGApage{I}{438}

un accouplement

\begin{equation}\tag{*}
v_\ell^{K'} : \underline{M}_K \otimes \underline{M}'_K \longrightarrow \mathbb{Z}_{\ell K} \quad .
\end{equation}
Ce dernier à priori dépend du choix de $K'$ , mais il résulte de (10.4.3) ci-dessous que l'accouplement

\begin{equation}\tag{9.7.1}
u_\ell : \underline{M}_K \otimes \underline{M}'_K \longrightarrow \mathbb{Q}_{\ell K}
\end{equation}
donné par

\begin{equation*}
u_\ell = (1/e(S'/S))\; v_\ell^{K'}
\end{equation*}
n'en dépend pas, où $e(S'/S) = \mathrm{long}(V'/\underline{m}V')$ . On pourra appeler encore (9.7.1) l'\underline{accouplement de monodromie} pour le schéma abélien $A_K$ sur $K$ et le nombre premier $\ell$ envisagé, puisque dans le cas où $A_K$ lui-même a réduction semi-stable, il coïncide évidemment avec la restriction à $\eta$ de l'accouplement de monodromie (9.1.2). Il résulte de 10.4 plus bas que cet accouplement est en fait toujours à valeurs dans $\mathbb{Q}_K$, et est indépendant de $\ell$ .

\SGAvar{On peut se demander s'il est à valeurs dans $\mathbb{Z}_K$ , ce qui revient à dire que (9.7.1) est à valeurs dans $\mathbb{Z}_{\ell K}$ pour tout nombre premier $\ell$ . C'est le cas du moins si $A_K$ est une courbe elliptique, le seul cas qui demande une vérification étant alors celui où $\underline{M}_K \simeq \underline{M}'_K$ est de rang 1, de sorte que $\underline{M}_K \otimes \underline{M}'_K$ est canoniquement isomorphe à $\mathbb{Z}_K$ , et que l'accouplement de monodromie (9.7.1) peut s'interprèter simplement comme un nombre rationnel ordinaire. On trouve grâce aux résultats de NERON [19] que celui-ci est égal à $v(j)$ , $j$ étant l'invariant absolu de $A_K$ et $v$ la valuation normalisée de $K$ .


}{On peut se demander s'il est à valeurs dans $\mathbb{Z}_K$ , ce qui revient à dire que (9.7.1) est à valeurs dans $\mathbb{Z}_{\ell K}$ pour tout nombre premier $\ell$ . C'est le cas du moins si $A_K$ est une courbe elliptique, le seul cas qui demande une vérification étant alors celui où $\underline{M}_K \simeq \underline{M}'_K$ est de rang 1, de sorte que $\underline{M}_K \otimes \underline{M}'_K$ est canoniquement isomorphe à $\mathbb{Z}_K$ , et que l'accouplement de monodromie (9.7.1) peut s'interpréter simplement comme un nombre rationnel ordinaire. On trouve grâce aux résultats de NERON [19] que celui-ci est égal à $v(j)$ , $j$ étant l'invariant absolu de $A_K$ et $v$ la valuation normalisée de $K$ .


}{F7I-IX-118}\SGApage{I}{439}


\SGAsectiontoc{10. Propriétés de l'accouplement de monodromie. Théorème d'intégrité et de positivité}{10}
10. \underline{Propriétés de l'accouplement de monodromie}. \underline{Théorème d'intégrité et de positivité}

Dans le présent paragraphe, nous allons donner tout d'abord certaines propriétés formelles essentiellement triviales des accouplements de monodromie, le détail (un peu fastidieux) des vérifications étant laissé au lecteur. D'autre part, nous indiquons un théorème plus profond (10.4) concernant ces accouplements, et indiquons comment on peut se réduire pour sa démonstration à un cas spécial de jacobiennes, qui sera traité (par voie transcendante) au par. 12 (qui est d'ailleurs logiquement indépendant du par. 11).

10.1. \underline{Transport de structure}. Il est clair, sur la définition donnée de l'accouplement de monodromie (9.1.2), que cet accouplement est compatible avec le transport de structure en la situation $(S,A_K,A_K')$.

\SGAvar{10.2. \underline{Symétrie}. Si nous partons du schéma dual $A_K'$ de $A_K$ et lui appliquons la construction du par. 9, nous devons introduire le schéma dual $(A_K')'$ de $A_K$ , qui est canoniquement isomorphe à $A_K$ par le théorème de bidualité ([5] ou VIII 3.2). Désignons par $u_\ell^{A_K}$ l'accouplement de monodromie (9.1.2) relativement à $A_K$ , de sorte qu'on trouve de même un accouplement

}{10.2. \underline{Symétrie}. Si nous partons du schéma dual $A_K'$ de $A_K$ et lui appliquons la construction du par. 9, nous devons introduire le schéma dual $(A_K')'$ de $A_K'$ , qui est canoniquement isomorphe à $A_K$ par le théorème de bidualité ([5] ou VIII 3.2). Désignons par $u_\ell^{A_K}$ l'accouplement de monodromie (9.1.2) relativement à $A_K$ , de sorte qu'on trouve de même un accouplement

}{F7I-IX-119}\begin{equation}\tag{10.2.1}
u_\ell^{A_K'} : \underline{M}'_\ell \otimes \underline{M}_\ell \longrightarrow \mathbb{Z}_\ell \qquad .
\end{equation}
Ceci posé, je dis que ce dernier se déduit de $u_\ell^{A_K}$ par la symétrie

\begin{equation*}
s : \underline{M}_\ell \otimes \underline{M}'_\ell \;\xleftarrow{\;\sim\;}\; \underline{M}'_\ell \otimes \underline{M}_\ell \qquad ,
\end{equation*}
i.e. on a la formule


\SGApage{I}{440}

\begin{equation}\tag{10.2.2}
u_\ell^{A_K'} \;=\; u_\ell^{A_K}\, s \qquad .
\end{equation}
Compte tenu de la propriété connue d'anticommutativité des accouplements $\varphi_\xi$ ([5] ou VIII 2.2.11), la formule de symétrie (10.2.2) peut se reformuler de la façon suivante, en termes d'une polarisation fixée $\xi$ de $A_K$ , qui définit un homomorphisme $\widetilde{\xi} : A_K \longrightarrow A_K'$ , d'où un homomorphisme correspondant ($t$ étant l'initiale de "torique")

\begin{equation}\tag{10.2.3}
\xi_t : \underline{M} \longrightarrow \underline{M}' \qquad ,
\end{equation}
permettant de former l'accouplement

\begin{equation}\tag{10.2.4}
\varphi^t_{\xi\ell} : \underline{M}_\ell \otimes \underline{M}_\ell \longrightarrow \mathbb{Z}_\ell \; , \; \varphi^t_{\xi,\ell}(x,y) = u_\ell^{A_K}(x,\xi_t(y)) \qquad :
\end{equation}
la forme précédente $\varphi^t_{\xi,\ell}$ est \underline{symétrique} :

\begin{equation}\tag{10.2.5}
\varphi^t_{\xi\ell}(x,y) = \varphi^t_{\xi\ell}(y,x) \qquad .
\end{equation}
En fait, cet énoncé est valable pour tout élément $\xi$ du groupe de Néron-Séveri de $A_K$ (pas nécessairement une polarisation) ; mais ce n'est que lorsque $\xi$ est non dégénéré i.e. définit une \underline{isogénie} $A_K \to A_K'$ donc aussi une isogénie (10.2.3), que la relation de symétrie (10.2.5) est équivalente à la relation de symétrie (10.2.2). D'ailleurs dans le cas où $\xi$ est une polarisation, on va préciser considérablement (10.2.5) dans 10.4 b) plus bas.

Signalons enfin une troisième façon d'interpréter la relation de symétrie. Pour ceci, considérons une biextension ($=$ correspondance divisorielle) $\xi$ sur un couple $(A_K,\overline{A}_K)$ de schémas abéliens à réduction semi-stable sur $K$, qui définit un homomorphisme

\begin{equation*}
A_K \longrightarrow \overline{A}_K' \qquad ,
\end{equation*}

\SGApage{I}{441}

\SGAvar{et par passage aux modèles de Néron définit donc un homomorphisme sur les tores maximaux des fibres epéciales, donc en sens inverse un homomorphisme sur les Groupes de caractères, d'où un homomorphisme

}{et par passage aux modèles de Néron définit donc un homomorphisme sur les tores maximaux des fibres spéciales, donc en sens inverse un homomorphisme sur les Groupes de caractères, d'où un homomorphisme

}{F7I-IX-120}\begin{equation}\tag{10.2.6}
\xi_t \; : \; \overline{M} \longrightarrow \underline{M}' \quad ,
\end{equation}
où $\overline{M}$ est défini en termes de $\overline{A}_K$ comme $\underline{M}$ en termes de $A_K$ . On peut donc former l'accouplement

\SGAvar{\begin{equation}\tag{10.2.7}
\varphi^t_\xi \; : \; \underline{M}_\ell \otimes \overline{\underline{M}}_\ell \longrightarrow \mathbb{Z}_\ell \quad , \quad \varphi^t(x,y) = \varphi^t(x,\xi_t(y)) \quad .
\end{equation}
Notons en passant que pour un schéma abélien $A_K$ a réduction semi-stable variable, le Groupe constant tordu $\underline{M}$ est un foncteur covariant en $A_K$ , et qu'alors la formation des accouplements (10.2.7) associés aux biextensions de schémas abéliens à réduction semi-stable sur $(K,S)$ est compatible, en un sens évident, avec la formation des images inverses de biextensions. Bien entendu, lorsque $\xi$ est la biextension canonique ( ``de Weil'' ) de $(A_K , A'_K)$, l'accouplement (10.2.7) n'est autre que l'accouplement de monodromie (9.2.1). - Ceci posé, la relation de symétrie peut s'exprimer de la façon suivante : considérons la symétrique ${}^{s}\xi$ de $\xi$ , qui est une biextension de $(\overline{A}_K,A_K)$, d'où un accouplement

}{\begin{equation}\tag{10.2.7}
\varphi^t_\xi \; : \; \underline{M}_\ell \otimes \overline{\underline{M}}_\ell \longrightarrow \mathbb{Z}_\ell \quad , \quad \varphi^t(x,y) = u_\ell(x,\xi_t(y)) \quad .
\end{equation}
Notons en passant que pour un schéma abélien $A_K$ à réduction semi-stable variable, le Groupe constant tordu $\underline{M}$ est un foncteur covariant en $A_K$ , et qu'alors la formation des accouplements (10.2.7) associés aux biextensions de schémas abéliens à réduction semi-stable sur $(K,S)$ est compatible, en un sens évident, avec la formation des images inverses de biextensions. Bien entendu, lorsque $\xi$ est la biextension canonique ( ``de Weil'' ) de $(A_K , A'_K)$, l'accouplement (10.2.7) n'est autre que l'accouplement de monodromie (9.1.2). - Ceci posé, la relation de symétrie peut s'exprimer de la façon suivante : considérons la symétrique ${}^{s}\xi$ de $\xi$ , qui est une biextension de $(\overline{A}_K,A_K)$, d'où un accouplement

}{F7I-IX-121}\begin{equation*}
\varphi^t_{s_\xi} \; : \; \overline{\underline{M}}_\ell \otimes \underline{M}_\ell \longrightarrow \mathbb{Z}_\ell \quad ;
\end{equation*}
on a alors la relation de symétrie

\begin{equation}\tag{10.2.8}
\varphi^t_{s_\xi} = \varphi^t_\xi \circ s \quad , \quad \text{où} \quad s: \overline{\underline{M}}_\ell \otimes \underline{M}_\ell \xrightarrow{\ \sim\ } \underline{M}_\ell \otimes \overline{\underline{M}}_\ell
\end{equation}
est encore l'isomorphisme de symétrie. On comparera cette relation à la relation d'antisymétrie de VIII 2.2.11.


\SGApage{I}{442}


10.3. \underline{Changement de trait.} Soit maintenant

\begin{equation}\tag{10.3.1}
g: S' \longrightarrow S
\end{equation}
un morphisme dominant de traits henséliens, associé à un homomorphisme local injectif d'anneaux de valuation discrète

\begin{equation*}
V \longrightarrow V' \quad ;
\end{equation*}
posons comme d'habitude

\begin{equation}\tag{10.3.2}
e(S'/S) = \mathrm{long}_V V'/\underline{m}V' \quad ,
\end{equation}
$\underline{m}$ étant l'idéal maximal de $V$. On sait alors que l'homomorphisme induit par (10.3.1) sur les parties ``modérées'' des groupes d'inertie

\begin{equation}\tag{10.3.3}
g_* \; : \; I'_t \simeq \prod_{\ell \neq p} \mathbb{Z}_\ell(1)(k') \longrightarrow I_t = \prod_{\ell \neq p} \mathbb{Z}_\ell(1)(k)
\end{equation}
est donné par

\begin{equation}\tag{10.3.4}
g_*(x) = e(S'/S)\, x \quad ,
\end{equation}
en identifiant les deux membres de (10.3.3) à l'aide des isomorphismes évidents $\mathbb{Z}_\ell(1)(k) \xrightarrow{\ \sim\ } \mathbb{Z}_\ell(1)(k')$ (I \qquad ).

\SGAvar{Notons d'autre part que, pour un schéma abélien $A_K$ sur $K$ à réduction semi-stable sur $S$, la formation des réseaux tordus $\underline{M}$ , $\underline{M}'$ commute à tout changement de base $S' \longrightarrow S$ comme dessus, comme il résulte aussitôt des définitions (notamment 3.4). On peut donc considérer l'accouplement $u_\ell^{A_K}$ défini par $A_{K'}$ ($K'$ = corps des fractions de $V'$) comme un accouplement

}{Notons d'autre part que, pour un schéma abélien $A_K$ sur $K$ à réduction semi-stable sur $S$, la formation des réseaux tordus $\underline{M}$ , $\underline{M}'$ commute à tout changement de base $S' \longrightarrow S$ comme dessus, comme il résulte aussitôt des définitions (notamment 3.4). On peut donc considérer l'accouplement $u_\ell^{A_{K'}}$ défini par $A_{K'}$ ($K'$ = corps des fractions de $V'$) comme un accouplement

}{F7I-IX-122}\begin{equation*}
u_\ell^{A_{K'}} \; : \; \underline{M}_{\ell S'} \otimes \underline{M}'_{\ell S'} \longrightarrow \mathbb{Z}_{\ell S'} \quad ,
\end{equation*}

\SGApage{I}{443}

\SGAvar{et le comparer à l'accouplement $u_{\ell S'}^{A_K}$, déduit de $u_\ell^{A_K}$ par changement de base. Lorsque $\ell \neq p$, la construction des accouplements $u_\ell$ donnée dans 9.2, jointe à la formule (10.2.3), montre inmédiatement qu'on a alors

}{et le comparer à l'accouplement $u_{\ell S'}^{A_K}$, déduit de $u_\ell^{A_K}$ par changement de base. Lorsque $\ell \neq p$, la construction des accouplements $u_\ell$ donnée dans 9.2, jointe à la formule (10.3.4), montre immédiatement qu'on a alors

}{F7I-IX-123}\begin{equation}\tag{10.3.5}
u_\ell^{A_{K'}} = e(S'/S) \; u_\ell^{A_K} \quad .
\end{equation}
Cette formule reste en fait valable également pour $\ell = p$, comme on voit en revenant à la construction générale donnée dans 9.3 à 9.6, en se ramenant à un énoncé de compatibilité analogue à (10.3.5) pour l'isomorphisme 9.4.4 e), pour un changement de traits strictement locaux $S' \longrightarrow S$. (On notera que l'isomorphisme en question est déduit de l'isomorphisme $K^*/V^* \simeq \mathbb{Z}$ déduit de la \underline{norme}, donnant donc lieu au diagramme commutatif

\begin{center}
\begin{tikzcd}
K^*/V^* \arrow[r, "\sim"] \arrow[d] \& \mathbb{Z} \arrow[d, "e(S'/S)"] \\
K^{\prime *}/V^{\prime *} \arrow[r, "\sim"] \& \mathbb{Z} \& .)
\end{tikzcd}
\end{center}











Voici un résultat plus profond concernant les accouplements de monodromie, dont on donnera une démonstration par voie transcendante au par. 12, en nous appuyant sur des résultats qui seront démontrés ultérieurement dans le Séminaire. (Nous n'utiliserons d'ailleurs pas 10.4 dans la suite du Séminaire, pas plus que les autres résultats du présent exposé.)


\SGApage{I}{444}

\underline{Théorème} 10.4. \underline{Soient} $S$ \underline{un trait hensélien,} $A_K$ \underline{un schéma abélien sur} $K$ \underline{à réduction semi-stable sur} $S$. \underline{Considérons les réseaux} $\underline{M}$ \underline{et} $\underline{M}'$ \underline{associés comme dans} 9.1 (cf. 5.5.2). \underline{Alors on a ce qui suit} :

a) \underline{Il existe un unique accouplement}

\begin{equation}\tag{10.4.1}
u: \underline{M} \otimes \underline{M}' \longrightarrow \mathbb{Z}_S
\end{equation}
\underline{tel que pour tout nombre premier} $\ell$, $u \otimes_{\mathbb{Z}} \mathbb{Z}_\ell$ \underline{soit l'accouplement de monodromie} $u_\ell$ \underline{de} (9.1.2). \underline{Cet accouplement est séparant.}

b) \underline{Soit} $\xi$ \underline{une polarisation de} $A_K$ , \underline{définissant donc un homomorphisme} (10.2.3)

\begin{equation*}
\xi_t : \underline{M} \longrightarrow \underline{M}' \quad ,
\end{equation*}
\underline{et considérons la forme}

\begin{equation}\tag{10.4.2}
\varphi_\xi^t : \underline{M} \otimes \underline{M} \longrightarrow \mathbb{Z}_S \quad , \quad \varphi^t(x,y) = u(x,\xi_t(y)) \quad .
\end{equation}
\underline{Cette forme est} symétrique, positive \underline{et} non dégénérée.

Bien entendu, en termes de la fibre géométrique $M_o$ de $\underline{M}$ en un point géométrique donné sur $S$,l'assertion précédente signifie que la forme correspondante

\SGAvar{\begin{equation*}
M_o \otimes M_o \longrightarrow \mathbb{Z}
\end{equation*}
est symétrique, positive et non dégénérée. Quand à l'accouplement (10.4.1), on l'appelera encore l'\underline{accouplement de monodromie}.


}{\begin{equation*}
M_o \otimes M_o \longrightarrow \mathbb{Z}
\end{equation*}
est symétrique, positive et non dégénérée. Quant à l'accouplement (10.4.1), on l'appellera encore l'\underline{accouplement de monodromie}.


}{F7I-IX-124}\SGApage{I}{445}


10.5. \underline{Réduction de} 10.4 \underline{à un cas particulier}

10.5.1. L'unicité d'un $u$ satisfaisant la condition énoncée (même pour un seul $\ell$) est claire. D'autre part, l'existence signifie simplement que chaque $u_\ell$ applique en fait $\underline{M} \otimes \underline{M}'$ dans $\mathbb{Z}_S$ , et que l'homomorphisme ainsi obtenu $\underline{M} \otimes \underline{M}' \longrightarrow \mathbb{Z}_S$ est indépendant de $\ell$ . En fait, il suffit même de prouver l'assertion analogue pour les $u_\ell \otimes_{\mathbb{Z}_\ell} \mathbb{Q}_\ell$ et un $u_{\mathbb{Q}} : \underline{M}_{\mathbb{Q}} \otimes \underline{M}'_{\mathbb{Q}} \longrightarrow \mathbb{Q}_S$ .

Supposons alors qu'on ait établi l'existence après changement de base $S' \longrightarrow S$ comme dans 9.3. Alors il résulte immédiatement de (10.3.5) qu'il existe un homomorphisme

\SGAvar{\begin{equation*}
u : \underline{M} \otimes \underline{M}' \longrightarrow \frac{1}{e(S'/S)} \mathbb{Z}_S \subset Q_S
\end{equation*}
compatible avec les $u_\ell \otimes_{\mathbb{Z}_\ell} \mathbb{Q}_\ell$ , d'où a). Quant à b), il est trivial alors que s'il est vrai pour $A_{K'}$ , il l'est aussi pour $A_K$ .

}{\begin{equation*}
u : \underline{M} \otimes \underline{M}' \longrightarrow \frac{1}{e(S'/S)} \mathbb{Z}_S \subset \mathbb{Q}_S
\end{equation*}
compatible avec les $u_\ell \otimes_{\mathbb{Z}_\ell} \mathbb{Q}_\ell$ , d'où a). Quant à b), il est trivial alors que s'il est vrai pour $A_{K'}$ , il l'est aussi pour $A_K$ .

}{F7I-IX-125}Par suite, nous pouvons, pour prouver 10.4, nous permettre n'importe quel changement de traits $S' \longrightarrow S$ dominant. Ceci nous permet en particulier de supposer $S$ complet à corps résiduel algébriquement clos. Dans ce cas, $\underline{M}$ et $\underline{M}'$ sont des schémas en groupes constants, de valeur $M,M'$ des modules libres de type fini sur $\mathbb{Z}$ , et a) et b) se reformulent avantageusement en termes de ces derniers.

\SGAvar{10.5.2. Supposons maintenant que $A_K$ soit isomorphe à un facteur direct d'un schéma abélien $\bar{A}_K$ sur $k$ à réduction semi-stable, et que 10.4 soit prouvé pour ce dernier. Alors il en résulte aussitôt que 10.4 est valable également pour $A_K$ , compte tenu du comportement fonctoriel évident des réseaux $\underline{M},\underline{M}'$ , et des accouplements de monodromie, vis-à-vis des sommes


}{10.5.2. Supposons maintenant que $A_K$ soit isomorphe à un facteur direct d'un schéma abélien $\bar{A}_K$ sur $K$ à réduction semi-stable, et que 10.4 soit prouvé pour ce dernier. Alors il en résulte aussitôt que 10.4 est valable également pour $A_K$ , compte tenu du comportement fonctoriel évident des réseaux $\underline{M},\underline{M}'$ , et des accouplements de monodromie, vis-à-vis des sommes


}{F7I-IX-126}\SGApage{I}{446}

directes de schémas abéliens à réduction semi-stable.

Il est également immédiat, grâce à 10.5.1, que la validité de 10.4 ne dépend que de la classe à isogénie près de $A_K$ . Joint à la remarque précédente, et du ``théorème de réductibilité de Poincaré'' [15 \qquad ], ceci nous montre que si $A_K$ est quotient d'un schéma abélien $\bar{A}_K$ pour lequel 10.4 est valable, il en est de même pour $A_K$ .

\SGAvar{10.5.3. Utilisons maintenant le fait bien connu [31, th. 11 (p.20)] que, quitte à faire une extension finie sur $K$ (ce qui est loisible, grâce à 10.5.1), tout schéma abélien $A_K$ est quotient de la jacobienne ( $= \underline{\mathrm{Pic}}^{0}_{X_K/K}$ ) d'une courbe propre, lisse et géométrique connexe $X_K$ sur $X$. Compte tenu de 10.5.2, on est donc ramené à prouver 10.4 dans le cas d'une telle jacobienne.

}{10.5.3. Utilisons maintenant le fait bien connu [31, th. 11 (p.20)] que, quitte à faire une extension finie sur $K$ (ce qui est loisible, grâce à 10.5.1), tout schéma abélien $A_K$ est quotient de la jacobienne ( $= \underline{\mathrm{Pic}}^{0}_{X_K/K}$ ) d'une courbe propre, lisse et géométriquement connexe $X_K$ sur $K$. Compte tenu de 10.5.2, on est donc ramené à prouver 10.4 dans le cas d'une telle jacobienne.

}{F7I-IX-127}10.5.4. D'après un théorème de ARTIN-MUMFORD [3] (déjà cité et utilisé dans V \qquad ), quitte à faire encore une extension finie $K'$ de $K$, on peut supposer que $X_K$ se prolonge en un schéma projectif et plat sur $S$, régulier, et tel que la fibre spéciale géométrique $X_{\bar{s}}$ n'ait comme seuls points singuliers que des points singuliers quadratiques ordinaires (i.e. ici, des croisements normaux de deux branches). Dans ce cas, nous prouverons 10.4 a) dans 12.5 plus bas, en explicitant entièrement l'accouplement de monodromie (10.4.1) en termes des données géométriques de la situation, et nous vérifierons en même temps 10.4 b) dans le cas de la polarisation canonique de la jacobienne.


\SGApage{I}{447}

10.5.5.  Il reste à voir que si pour un $A_K$ donné à réduction semi-stable, 10.4 a) est vrai, et 10.4 b) est vrai pour \underline{une} polarisation donnée $\xi_o$ de $A_K$ , alors cette conclusion reste vraie pour \underline{toute} polarisation $\xi$ de $A_K$ . En effet, si $\widetilde{\xi}_o$ , $\widetilde{\xi}$ désignent les homomorphismes

\begin{equation*}
\widetilde{\xi}_o \;,\; \widetilde{\xi} \;:\; A_K \longrightarrow A'_K
\end{equation*}
définis par $\xi_o$ , $\xi$ , qui sont donc des isogénies, on aura

\begin{equation*}
\widetilde{\xi} = \widetilde{\xi}_o \, \alpha = {}^t\alpha^* \, \widetilde{\xi}_o \;,\; \text{avec } \alpha \in \mathrm{End}(A_K) \otimes_{\mathbb{Z}} \mathbb{Q} \qquad ,
\end{equation*}
où $\alpha \longmapsto \alpha^*$ désigne l'involution canonique de l'algèbre $\mathrm{End}(A_K) \otimes_{\mathbb{Z}} \mathbb{Q}$ définie par la polarisation $\xi_o$ , et où l'exposant $t$ dans ${}^t(\alpha^*)$ désigne le transposé (associant à un endomorphisme de $A_K$ un endomorphisme de $A'_K$). D'autre part, il est connu [15, Chap V th.3] que l'on a

\begin{equation*}
\alpha = \sum_i \beta_i^* \, \beta_i \qquad , \qquad \text{avec } \beta_i \in \mathrm{End}(A_K) \otimes_{\mathbb{Z}} \mathbb{Q} \qquad ,
\end{equation*}
d'où

\begin{equation*}
\widetilde{\xi} = \sum_i {}^t\beta_i \, \widetilde{\xi}_o \, \beta_i \qquad ,
\end{equation*}
ou encore, en termes des biextensions (= correspondances divisorielles) associées aux polarisations envisagées (VIII 4.14) :

\begin{equation}\tag{10.5.5.1}
\delta(\xi) = \sum_i \delta(\xi_o) \circ (\beta_i , \beta_i) \qquad .
\end{equation}
D'après ce qui a été dit après (10.2.7), cette formule implique (avec les notations de (10.4.2))

\begin{equation}\tag{10.5.5.2}
\varphi^t_{\xi} = \sum_i \varphi^t_{\xi_o} \, (\beta_i , \beta_i) \qquad \in \mathrm{Hom}(M_{\mathbb{Q}} \otimes_{\mathbb{Q}} M_{\mathbb{Q}} , \mathbb{Q}) \qquad .
\end{equation}

\SGApage{I}{448}

Elle implique bien que si $\varphi^t_{\xi_o}$ est symétrique et positive, il en est de même de $\varphi^t_{\xi}$ .

Enfin, le fait que $\varphi^t_{\xi_o}$ soit non dégénérée équivaut évidemment à l'assertion que l'accouplement (10.4.1) est séparant (compte tenu que $\xi_{ot} : M \otimes_{\mathbb{Z}} \mathbb{Q} \longrightarrow M' \otimes_{\mathbb{Z}} \mathbb{Q}$ est un isomorphisme, $\widetilde{\xi}_o$ étant une isogénie), et est donc également indépendant de la polarisation choisie. On peut d'ailleurs noter que pour vérifier que (10.4.1) est séparant, il suffit de vérifier la même assertion pour l'accouplement de monodromie $u_\ell$ (9.1.2) pour \underline{un} nombre premier $\ell$ . Prenant $\ell \neq p$, on peut utiliser la description 9.2 de $u_\ell$ ; le fait que $u_\ell$ soit séparant à gauche (donc des deux côtés, pour une raison de rangs) est alors trivial.


\SGAvar{\SGAsectiontoc{11. Composantes connexes du modèle de Néron : relation de dualité, comportement assymptotique}{11}
11. \underline{Composantes connexes du modèle de Néron} : \underline{relation de dualité, comportement assymptotique}

}{\SGAsectiontoc{11. Composantes connexes du modèle de Néron : relation de dualité, comportement asymptotique}{11}
11. \underline{Composantes connexes du modèle de Néron} : \underline{relation de dualité, comportement asymptotique}

}{F7I-IX-128}11.0.  Le but du présent paragraphe est l'étude du $k$-groupe fini étale

\begin{equation}\tag{11.0.1}
\Phi_o = A_o / A_o^{\,o}
\end{equation}
déjà considéré dans 1. Pour tout nombre premier $\ell$, nous désignons par

\SGAvar{\begin{equation*}
\Phi_o(\ell)
\end{equation*}
la composante $\ell$-primaire de ce dernier, de sorte que la connaissance de $\Phi_o$ équivaut à celle des $\Phi_o(1)$, grâce à la formule

}{\begin{equation*}
\Phi_o(\ell)
\end{equation*}
la composante $\ell$-primaire de ce dernier, de sorte que la connaissance de $\Phi_o$ équivaut à celle des $\Phi_o(\ell)$, grâce à la formule

}{F7I-IX-129}\begin{equation}\tag{11.0.2}
\Phi_o = \coprod_{\ell} \Phi_o(\ell) \qquad .
\end{equation}

\SGApage{I}{449}

Notons que les groupes $\Phi_{o}$ , $\Phi_{o}(\ell)$ sont déterminés respectivement par la connaissance des groupes finis ordinaires $\Phi_{o}(\bar{k})$, $\Phi_{o}(\ell)(\bar{k})$ et de l'opération naturelle de $\pi_{o} = \mathrm{Gal}(\bar{k}/k)$ sur ces derniers. D'autre part, il est clair que $\Phi_{o}(\bar{k})$ est aussi le groupe des composantes connexes du modèle de Néron de $A_{\widetilde{K}}$ , où $\widetilde{K}$ est le corps des fractions de l'hensélisé strict de $S$ relativement à la clôture séparable choisie $\bar{k}$ de $k$. Ceci nous permet donc, pour l'étude de $\Phi_{o}$ , de nous ramener au cas où $S$ est strictement local, - les opérations de $\pi_{o} \simeq \mathrm{Gal}(\widetilde{K}/K)$ sur $\Phi_{o}(\bar{k})$ se récupérant automatiquement par transport de structure. On pourrait de même se ramener au cas où $S$ est de plus complet, grâce au changement de base $\widehat{S} \longrightarrow S$, grâce au fait que la formation du modèle de Néron y commute [21, th. 3.4 b)].

11.1.  Supposons donc $S$ strictement local, de sorte que $\Phi_{o}$ peut s'identifier à un groupe fini ordinaire $\Phi_{o}(k)$, et on trouve

\begin{equation}\tag{11.1.1}
\Phi_{o}(k) \simeq A_{o}(k)/A_{o}^{o}(k) \simeq A(S)/A^{o}(S) \simeq A(K)/A(K)^{o} \quad ,
\end{equation}
où la première égalité provient du fait que $A_{o}$ est lisse sur $k$, la deuxième du ``lemme de Hensel'' compte tenu du fait que $A$ est lisse sur $S$ et $S$ hensélien, enfin la troisième provient de la relation $A(S) \stackrel{\sim}{\longrightarrow} A(K)$ et du fait que nous posons

\begin{equation}\tag{11.1.2}
A(K)^{o} \stackrel{\mathrm{dfn}}{=} A^{o}(S) \quad ,
\end{equation}
de sorte que $A(K)^{o}$ apparaît comme un sous-groupe d'indice fini naturel dans $A(K)$ (à ne pas confondre avec $A^{o}(K) = A(K)$ !). Les isomorphismes


\SGApage{I}{450}

(11.1.1) sont évidemment fonctoriels en $A$, et compatibles avec le transport de structure en $(S,A_{K})$, et en particulier avec l'action d'éventuels groupes de Galois.

Supposons maintenant que $\ell$ soit un nombre premier $\neq p$, de sorte que $A^{o}$ est $\ell$-divisible, et plus précisément $\ell\ \mathrm{id}_{A^{o}}$ est un morphisme étale et surjectif, de sorte que ce morphisme induit un endomorphisme surjectif de $A^{o}(S)=A(K)^{o}$, qui est donc un \underline{groupe} $\ell$\underline{-divisible}. Il en résulte aussitôt que pour tout entier $n \geq 0$ premier à $p = \mathrm{car}(k)$, on a la relation

\begin{equation}\tag{11.1.3}
{}_{n}\Phi_{o}(k) \simeq {}_{n}A(K)/{}_{n}A(K)^{o} \quad , \quad \mathrm{pour}\ (n,p) = 1 \quad .
\end{equation}
Or par définition (2.2.3) de la partie fixe de ${}_{n}A_{K}$ , on a

\begin{equation}\tag{11.1.4}
{}_{n}A(K)^{o} = ({}_{n}A_{K})^{f}(\bar{K}) \simeq T_{n}(A_{K})^{f}(\bar{K}) \otimes_{\mathbb{Z}_{n}} \mathbb{Z}/n\mathbb{Z} \quad .
\end{equation}
Prenant pour $n$ une puissance de $\ell$, et posant comme dans (2.5.3)

\begin{equation}\tag{11.1.5}
U = T_{\ell}(A(\bar{K})) \quad , \quad \mbox{d'où} \quad T_{\ell}(A(\bar{K}))^{f} = U^{I} \quad ,
\end{equation}
où $U^{I}$ désigne le sous-module de $U$ formé des invariants sous le groupe d'inertie $I$, on peut écrire (11.1.3) sous la forme

\begin{equation}\tag{11.1.6}
{}_{n}\Phi_{o}(k) \simeq (U_{n})^{I}/(U^{I})_{n} \quad , \quad n = \ell^{r+1} \quad ,
\end{equation}
où l'indice $n$ désigne la réduction mod. $n$ i.e. le produit tensoriel sur $\mathbb{Z}_{\ell}$ avec $\mathbb{Z}_{\ell}/n\mathbb{Z}_{\ell}$ . Les morphismes de transition entre les groupes intervenant dans les deux membres de (11.1.6) étant évidents, on déduit de (11.1.6), par passage à la limite inductive sur les puissances de $\ell$, la


\SGApage{I}{451}

\underline{Proposition} 11.2. \underline{Supposons} $S$ \underline{strictement local}, \underline{et soit} $\ell$ \underline{un nombre premier} $\neq p$. \underline{Alors on a un isomorphisme canonique}

\begin{equation}\tag{11.2.1}
\Phi_o(\ell) \;\simeq\; (U\otimes D_\ell)^I/U^I\otimes D_\ell \quad ,
\end{equation}
\underline{où} $U$ \underline{est défini dans} (11.1.5) \underline{et où on pose} $D_\ell = \mathbb{Q}_\ell/\mathbb{Z}_\ell$ .

Ainsi, le groupe $\Phi_o(\ell)$ peut s'expliciter en termes de la seule connaissance du module de Tate $U = T_\ell(A(\overline{K}))$ comme module galoisien sous l'action du groupe d'inertie $I$.

11.3. On va utiliser l'expression explicite (11.2.1) pour esquisser la définition d'un accouplement canonique

\begin{equation}\tag{11.3.1}
\Phi_o(\ell) \;\otimes\; \Phi'_o(\ell) \longrightarrow \mathbb{Q}/\mathbb{Z} \quad ,
\end{equation}
qui est une dualité parfaite, - $\Phi'_o$ désignant comme au par. 1 le groupe de composantes connexes associé au modèle de Néron du \underline{dual} $A'_K$ de $A_K$ . Il n'y a guère de doute (mais le rédacteur s'est dispensé de faire la vérification) que cet accouplement n'est autre, éventuellement au signe près, que celui défini dans (1.2.1), et le lecteur est invité à titre d'exercice d'établir la compatibilité voulue (avec le signe correct !).

Posons

\begin{equation*}
U' = T_\ell(A'(\overline{K})) \quad , \quad T = T_\ell(1)(\overline{K}) = T_\ell(\overline{K}^*) \quad .
\end{equation*}
On a \underline{un accouplement parfait} de $\mathbb{Z}_\ell$-modules libres de type fini, compatible aux actions de $I$ :

\begin{equation}\tag{11.3.2}
\varphi : U \otimes U' \longrightarrow T, \quad (T \simeq \mathbb{Z}_\ell \ \text{non canoniquement}) \quad ,
\end{equation}

\SGApage{I}{452}

et, en vertu de la structure connue de $I$ ($I$\hspace{2.5em}), une suite exacte

\begin{equation}\tag{11.3.3}
1 \longrightarrow R \longrightarrow I \longrightarrow T \longrightarrow 1 \quad ,
\end{equation}
avec $R$ un groupe pro-fini \underline{premier à} $\ell$ . Si maintenant nous \underline{partons} de données (11.3.2) et (11.3.3), et introduisons les groupes

\SGAvar{\begin{equation}\tag{11,3.4}
\Phi = (U\otimes D_\ell)^I/U^I\otimes D_\ell \;,\; \Phi' = (U'\otimes D_\ell)^I/{U'}^{I}\otimes D_\ell \;,\; (D_\ell = \mathbb{Q}_\ell/\mathbb{Z}_\ell)
\end{equation}
nous allons définir une dualité parfaite

}{\begin{equation}\tag{11.3.4}
\Phi = (U\otimes D_\ell)^I/U^I\otimes D_\ell \;,\; \Phi' = (U'\otimes D_\ell)^I/{U'}^{I}\otimes D_\ell \;,\; (D_\ell = \mathbb{Q}_\ell/\mathbb{Z}_\ell)
\end{equation}
nous allons définir une dualité parfaite

}{F7I-IX-130}\begin{equation}\tag{11.3.5}
\Phi \otimes \Phi' \longrightarrow \mathbb{Q}/\mathbb{Z} \quad ,
\end{equation}
ce qui nous donnera en particulier l'accouplement annoncé (11.3.1). Considérons la suite exacte de $I$-modules

\begin{equation}\tag{11.3.6}
0 \longrightarrow U \xrightarrow{\ i\ } U\otimes_{\mathbb{Z}_\ell}\mathbb{Q}_\ell \xrightarrow{\ j\ } U\otimes D_\ell \longrightarrow 0 \quad ,
\end{equation}
et la suite exacte de cohomologie associée, en prenant la cohomologie de $I$ définie par les \underline{cochaînes continues} à coefficients dans les groupes topologiques envisagés :

\SGAvar{\begin{equation}\tag{11.3.7}
\begin{tikzcd}[column sep=small, row sep=large]
0 \arrow[r] \& H^o(I,U) \arrow[r, "i^o"] \arrow[d, dash, "\wr"] \& H^o(I,U\otimes\mathbb{Q}_\ell) \arrow[r, "j^o"] \arrow[d, dash, "\wr"] \& H^o(I,U\otimes D_\ell) \arrow[r, "\partial"] \arrow[d, dash, "\wr"] \& H^1(I,U) \arrow[r, "i^1"] \& H^1(I,U\otimes\mathbb{Q}_\ell) \arrow[d, dash, "\wr"] \\
{} \& U^I \& (U\otimes Q_\ell)^I \simeq U^I\otimes_{\mathbb{Z}}\mathbb{Q} \& (U\otimes D_\ell)^I \& {} \& H^1(I,U)\otimes_{\mathbb{Z}}\mathbb{Q} \;.
\end{tikzcd}
\end{equation}












}{\begin{equation}\tag{11.3.7}
\begin{tikzcd}[column sep=small, row sep=large]
0 \arrow[r] \& H^o(I,U) \arrow[r, "i^o"] \arrow[d, dash, "\wr"] \& H^o(I,U\otimes\mathbb{Q}_\ell) \arrow[r, "j^o"] \arrow[d, dash, "\wr"] \& H^o(I,U\otimes D_\ell) \arrow[r, "\partial"] \arrow[d, dash, "\wr"] \& H^1(I,U) \arrow[r, "i^1"] \& H^1(I,U\otimes\mathbb{Q}_\ell) \arrow[d, dash, "\wr"] \\
{} \& U^I \& (U\otimes\mathbb{Q}_\ell)^I \simeq U^I\otimes_{\mathbb{Z}}\mathbb{Q} \& (U\otimes D_\ell)^I \& {} \& H^1(I,U)\otimes_{\mathbb{Z}}\mathbb{Q} \;.
\end{tikzcd}
\end{equation}












}{F7I-IX-131}Il est immédiat que $\Phi$ défini par (11.3.4) n'est autre que coker $j^o$ , donc est canoniquement isomorphe à Ker $i^1$, i.e. on trouve une suite exacte canonique

\begin{equation}\tag{11.3.8}
0 \longrightarrow \Phi \longrightarrow H^1(I,U) \longrightarrow H^1(I,U)\otimes_{\mathbb{Z}}\mathbb{Q} \quad ,
\end{equation}

\SGApage{I}{453}

qui montre que $\Phi$ \underline{est canoniquement isomorphe au sous-groupe de torsion de} $H^1(I,U)$.

D'autre part, de la suite exacte (11.3.3) on déduit un isomorphisme, pour tout $n = \ell^r$ :

\begin{equation}\tag{11.3.9}
H^1(I,T_n) \simeq \mathrm{Hom}(I,T_n) \xleftarrow{\sim} \mathrm{Hom}(T,T_n) \xleftarrow{\sim} \mathbb{Z}/n\mathbb{Z}
\end{equation}
(NB on fait opérer $I$ trivialement sur $T$, donc sur $T_n = T/nT$ ) ; d'où, pour tout $I$-$\mathbb{Z}_\ell$-Module fini $M$, si

\begin{equation}\tag{11.3.10}
M' = \mathrm{Hom}(M, T \otimes D_\ell)
\end{equation}
désigne son ``dual de Cartier'', un accouplement canonique

\begin{equation}\tag{11.3.11}
H^i(I,M) \otimes H^{1-i}(I,M') \longrightarrow H^1(I,T \otimes D_\ell) \simeq D_\ell \quad ,
\end{equation}
où le dernier isomorphisme est déduit des isomorphismes (11.3.9) par passage à la limite inductive sur les puissances $n$ de $\ell$. Il est bien connu que ces accouplements sont des \underline{dualités parfaites} (cf. par exemple SGA 5 I 5). Prenant dans ces accouplements $M$ de la forme $U_n$ , donc $M' = U'_n$ , on trouve par passage à la limite un accouplement parfait

\begin{equation}\tag{11.3.12}
H^o(I,\, U \otimes D_\ell) \times H^1(I,U') \longrightarrow D_\ell \quad ,
\end{equation}
où le deuxième facteur du premier membre est un $\mathbb{Z}_\ell$-module de type fini. Un nouveau passage à la limite dans les accouplements (11.3.12) (pour des morphismes de transition de la forme $\ell^r.\mathrm{id}$) nous donne un deuxième accouplement parfait

\begin{equation}\tag{11.3.13}
H^o(I,U \otimes \mathbb{Q}_\ell) \times H^1(I,U' \otimes \mathbb{Q}_\ell) \longrightarrow D_\ell
\end{equation}
de vectoriels de dimension finie sur $\mathbb{Q}_\ell$ . Désignant par $i'$, $j'$ les


\SGApage{I}{454}

morphismes définis comme dans (11.3.6), avec $U$ remplacé par $U'$, on trouve alors aisément que les homomorphismes suivants

\begin{center}
\begin{tikzcd}[column sep=6em, row sep=6em]
H^o(I,U\otimes\mathbb{Q}_\ell) \arrow[r, "j^o"] \& H^o(I,U\otimes D_\ell) \\
H^1(I,U'\otimes\mathbb{Q}_\ell) \& H^1(I,U') \arrow[l, "{i'}^{1}"']
\end{tikzcd}
\end{center}









sont \underline{transposés} l'un de l'autre, via les accouplements (11.3.12) et (11.3.13). Il en résulte aussitôt que coker $j^o$ et Ker ${i'}^{1}$ sont duals l'un de l'autre via un accouplement canonique à valeurs dans $D_\ell$ , qui est l'accouplement annoncé (11.3.5).

11.4.  La description précédente d'un accouplement (11.3.1), via l'expression (11.2.1) pour la composante $\ell$-primaire du groupe des composantes connexes, n'est valable que pour $\ell \neq p$. Supposant maintenant que $A_K$ a \underline{réduction semi-stable} sur $S$, nous allons encore définir des accouplements qui sont des dualités parfaites

\SGAvar{\begin{equation}\tag{11.4.1}
\Phi_o(\ell) \otimes \Phi'_o(\ell) \longrightarrow \mathbb{Q}/\mathbb{Z}
\end{equation}
sans restriction sur $\ell$. Comme dans 11.4, il y aurait lieu d'établir qu'on retrouve (à un signe près peut-être qu'il conviendrait de préciser) l'accouplement défini dans (1.2.1), ce qui établirait donc la conjecture 1.3 dans le cas de la réduction semi-stable, comme dans le cas où l'on se restreint aux composantes $\ell$-primaires pour $\ell \neq p$.

}{\begin{equation}\tag{11.4.1}
\Phi_o(\ell) \otimes \Phi'_o(\ell) \longrightarrow \mathbb{Q}/\mathbb{Z}
\end{equation}
sans restriction sur $\ell$. Comme dans 11.3, il y aurait lieu d'établir qu'on retrouve (à un signe près peut-être qu'il conviendrait de préciser) l'accouplement défini dans (1.2.1), ce qui établirait donc la conjecture 1.3 dans le cas de la réduction semi-stable, comme dans le cas où l'on se restreint aux composantes $\ell$-primaires pour $\ell \neq p$.

}{F7I-IX-132}Pour définir (11.4.1), nous allons donner, essentiellement, une expression de $\Phi_o(\ell)$ en termes du groupe de Barsotti-Tate $T_\ell(A_K)$


\SGApage{I}{455}

sur $K$, qui jouera dans le cas présent le même rôle que 11.2 dans le cas traité précédemment :

\underline{Théorème} 11.5. \underline{Supposons} $S$ \underline{hensélien, et que} $A_K$ \underline{ait réduction semi-stable sur} $S$. \underline{Soit} $\ell$ \underline{un nombre premier, et considérons l'accouplement de monodromie} (9.1.2), \underline{d'où un homomorphisme canonique}

\SGAvar{\begin{equation*}
\widetilde{u}_{\ell} \; ; \; \underline{M}_{\ell} \longrightarrow \underline{\check{M}}'_{\ell} \quad .
\end{equation*}
\underline{Soit} $\Phi(\ell) = \mathrm{coker}\ \widetilde{u}_{\ell}$ ; \underline{alors on a un isomorphisme canonique}

}{\begin{equation*}
\widetilde{u}_{\ell} \; : \; \underline{M}_{\ell} \longrightarrow \underline{\check{M}}'_{\ell} \quad .
\end{equation*}
\underline{Soit} $\Phi(\ell) = \mathrm{coker}\ \widetilde{u}_{\ell}$ ; \underline{alors on a un isomorphisme canonique}

}{F7I-IX-133}\begin{equation}\tag{11.5.1}
\Phi_o(\ell) \simeq \Phi(\ell)\times_S s \quad .
\end{equation}
De cet énoncé résulte immédiatement la définition d'un accouplement parfait (11.4.1), puisque on trouve, en appliquant 11.5 à $A'_K$ au lieu de $A_K$ et utilisant (10.2.2), qu'on a aussi un isomorphisme canonique

\begin{equation*}
\Phi'_o(\ell) \simeq \Phi'(\ell)\times_S s \quad , \quad \text{où } \Phi'(\ell) = \mathrm{coker}({}^t\widetilde{u}_{\ell} \; : \; \underline{M}'_{\ell} \longrightarrow \underline{\check{M}}_{\ell}) \quad .
\end{equation*}
\underline{Remarques} 11.5.2. a) Le théorème 11.5 implique notamment que $\widetilde{u}_{\ell}$ est une isogénie i.e. que $u_{\ell}$ est séparant, y compris dans le cas $\ell = p$, où cette assertion n'est pas triviale sur la définition de 9.6. On notera que nous n'utilisons pas ici le théorème (nettement plus profond) 10.4, qui sera prouvé au par. suivant.

b) Utilisant 10.4, qui nous fournit un accouplement de monodromie $u: \underline{M} \otimes \underline{M}' \longrightarrow \mathbb{Z}_S$ , et désignant par $\Phi$ le conoyau de l'homomorphisme correspondant

\begin{equation*}
\widetilde{u} \; : \; \underline{M} \longrightarrow \underline{\check{M}}' \quad ,
\end{equation*}

\SGApage{I}{456}

on trouve une formulation équivalente de 11.5 indépendante du choix d'un $\ell$, par l'isomorphisme canonique

\begin{equation}\tag{11.5.3}
\Phi_o \simeq \Phi \times_S s \quad .
\end{equation}
11.6. Pour définir l'isomorphisme (11.5.1), nous pouvons nous borner au cas où $S$ est strictement local (11.0), de sorte que $\underline{M}$ et $\underline{M}'$ sont constants, de valeurs

\begin{equation}\tag{11.6.1}
M = D(T'_o) \quad , \quad M' = D(T_o) \quad ,
\end{equation}
où $T_o$ et $T'_o$ sont les tores maximaux de $A_o$ et de $A'_o$ respectivement. Identifiant de même $\Phi_o$ à un groupe fini ordinaire, il faut définir un isomorphisme canonique entre sa composante $\Phi_o(\ell)$ et

\begin{equation*}
\Phi(\ell) = \mathrm{coker}(\widetilde{u}_{\ell} : M_{\ell} \longrightarrow \check{M}'_{\ell}) \quad .
\end{equation*}
En fait, nous allons définir directement un isomorphisme canonique

\SGAvar{\begin{equation}\tag{11.6.2}
\Phi_o(\ell) \simeq \ker(\widetilde{u}_{\ell}\otimes D_{\ell} \; : \; M_{\ell}\otimes D_{\ell} \longrightarrow \check{M}'_{\ell}\otimes D_{\ell}) \quad ;
\end{equation}
comme $\Phi_o(\ell)$, est fini, et $M_{\ell}$ et $\check{M}'_{\ell}$ ont même rang, il en résultera d'abord que $\widetilde{u}_{\ell} : M_{\ell} \longrightarrow \check{M}'_{\ell}$ est une isogénie, d'où il résulte ensuite que le deuxième membre de (11.6.2) est canoniquement isomorphe à $\Phi(\ell)$ :

}{\begin{equation}\tag{11.6.2}
\Phi_o(\ell) \simeq \ker(\widetilde{u}_{\ell}\otimes D_{\ell} \; : \; M_{\ell}\otimes D_{\ell} \longrightarrow \check{M}'_{\ell}\otimes D_{\ell}) \quad ;
\end{equation}
comme $\Phi_o(\ell)$ est fini, et $M_{\ell}$ et $\check{M}'_{\ell}$ ont même rang, il en résultera d'abord que $\widetilde{u}_{\ell} : M_{\ell} \longrightarrow \check{M}'_{\ell}$ est une isogénie, d'où il résulte ensuite que le deuxième membre de (11.6.2) est canoniquement isomorphe à $\Phi(\ell)$ :

}{F7I-IX-134}\begin{equation}\tag{11.6.3}
\Phi(\ell) = \mathrm{coker}\ \widetilde{u}_{\ell} \simeq \ker \widetilde{u}_{\ell}\otimes D_{\ell} \quad ,
\end{equation}
puisqu'on peut écrire, identifiant $M_{\ell}$ par $\widetilde{u}_{\ell}$ à un sous-module d'indice fini dans $\check{M}'_{\ell}$ :

\SGAvar{\begin{equation*}
\Phi(\ell) = V/M_{\ell} \quad , \quad \Phi'(\ell) = V/\check{M}'_{\ell} \quad ,
\end{equation*}
où $V = M_{\ell}\otimes Q_{\ell} \simeq \check{M}'_{\ell}\otimes Q_{\ell}$ . De (11.6.2) et (11.6.3) l'isomorphisme à définir (11.5.1) résulte aussitôt par composition, de sorte que tout revient à


}{\begin{equation*}
M_{\ell}\otimes D_{\ell}=V/M_{\ell} \quad , \quad \check{M}'_{\ell}\otimes D_{\ell}=V/\check{M}'_{\ell} \quad ,
\end{equation*}
où $V = M_{\ell}\otimes Q_{\ell} \simeq \check{M}'_{\ell}\otimes Q_{\ell}$ . De (11.6.2) et (11.6.3) l'isomorphisme à définir (11.5.1) résulte aussitôt par composition, de sorte que tout revient à


}{F7I-IX-135}\SGApage{I}{457}

\SGAvar{définir l'isomorphisme canonique (11.6,2).

Ecrivant $M_\ell \otimes D_\ell$ (resp. $\check{M}_\ell \otimes D_\ell$) comme limite inductive des $M_n$ (resp. des $\check{M}'_n$), où $n$ parcourt les puissances de $\ell$, il revient au même de définir (11.6.2), où des isomorphismes canoniques

\begin{equation}\tag{11.6.4}
{}_n\Phi_o \simeq \mathrm{Ker}\,(\tilde{u}_n : M_n \longrightarrow \check{M}'_n) \quad , \quad n = \ell^{r+1} \quad ,
\end{equation}
comptatibles avec les morphismes de transition évidents, pour $n$ variable (où $\tilde{u}_n$ désigne maintenant, par abus de notations, l'homomorphisme déduit de $\tilde{u}_\ell : M_\ell \longrightarrow \check{M}'_\ell$ par réduction mod. $n$).

}{définir l'isomorphisme canonique (11.6.2).

Ecrivant $M_\ell \otimes D_\ell$ (resp. $\check{M}'_\ell \otimes D_\ell$) comme limite inductive des $M_n$ (resp. des $\check{M}'_n$), où $n$ parcourt les puissances de $\ell$, il revient au même de définir (11.6.2), ou des isomorphismes canoniques

\begin{equation}\tag{11.6.4}
{}_n\Phi_o \simeq \mathrm{Ker}\,(\tilde{u}_n : M_n \longrightarrow \check{M}'_n) \quad , \quad n = \ell^{r+1} \quad ,
\end{equation}
compatibles avec les morphismes de transition évidents, pour $n$ variable (où $\tilde{u}_n$ désigne maintenant, par abus de notations, l'homomorphisme déduit de $\tilde{u}_\ell : M_\ell \longrightarrow \check{M}'_\ell$ par réduction mod. $n$).

}{F7I-IX-136}Pour ceci, considérons le schéma en groupes ${}_nA$ sur $S$, qui est plat et quasi-fini sur $S$ (2.2.1), grâce à l'hypothèse de réduction semi-stable. On a une filtration canonique

\begin{equation}\tag{11.6.5}
{}_nA \supset ({}_nA^o)^f \supset ({}_nA^o)^t \supset 0 \quad ,
\end{equation}
où l'exposant $f$ (resp. l'exposant $t$) désigne la ``partie fixe'' (2.2.3) (resp. la ``partie torique'' (5.1)) du schéma en groupe qu'il affecte ; pour la définition de la partie torique dans 5.1, on notera qu'on avait supposé $S$ complet, ce qui est loisible ici (11.0). (On a d'ailleurs vu dans 6.1 que la construction se généralise au cas $S$ hensélien, si on y tient ...) Soit

\begin{equation}\tag{11.6.6}
{}_n\Psi = {}_nA/({}_nA^o)^f \quad ,
\end{equation}
c'est un schéma en groupes plat, séparé et quasi-fini sur $S$ (SGA 3 $\mathrm{VI}_A$ 3.2 et $\mathrm{VI}_B$ 9.2) de fibre générique

\begin{equation*}
{}_n\Psi_K = ({}_nA_K)/({}_nA_K)^f \quad ,
\end{equation*}

\SGApage{I}{458}

où par définition $({}_nA_K)^f = ({}_nA^o)^f_K$ désigne le r.ème composant $(n=\ell^{r+1})$ du pro-groupe de Barsotti-Tate $T_\ell(A_K)^f$ . La première des relations (5.5.8) nous fournit alors un isomorphisme canonique

\begin{equation}\tag{11.6.7}
{}_n\Psi_K \;\simeq\; \underline{M}_{nK} \qquad \mbox{où} \qquad \underline{M}_n \;\stackrel{\mathrm{dfn}}{=}\; \underline{M} \otimes_{\mathbb{Z}} (\mathbb{Z}/n\mathbb{Z}) \qquad .
\end{equation}
En fait, on a :

\underline{Lemme} 11.6.8. \underline{Le schéma en groupes} ${}_n\Psi$ (11.6.6) \underline{sur} $S$ \underline{est étale sur} $S$, \underline{et l'isomorphisme} (11.6.7) \underline{se prolonge de façon unique en un morphisme}

\begin{equation}\tag{11.6.8.1}
{}_n\Psi \lhook\joinrel\longrightarrow \underline{M}_n
\end{equation}
\underline{qui est une immersion ouverte}.

Pour voir que ${}_n\Psi$ est étale, il reste à vérifier en vertu de (11.6.7) qu'il l'est en les points de sa fibre spéciale, ou encore que cette dernière est étale, ce qui est trivial, la fibre en question étant ${}_nA_o/{}_nA_o^o \subset A_o/A_o^o = \Phi_o$ . Il s'ensuit que ${}_n\Psi$ est un schéma normal, et comme $\underline{M}_n$ est fini sur $S$, l'existence d'un prolongement (11.6.8.1) en résulte aussitôt. C'est un morphisme de schémas en groupes étales sur $S$, donc son noyau est un groupe étale sur $S$, et comme ${}_n\Psi$ donc ce noyau est séparé, et sa fibre générique est le groupe unité par (11.6.7), c'est le groupe unité (3.1.3). Donc (11.6.8.1) est un monomorphisme étale, donc une immersion ouverte (EGA IV 17.9.1), ce qui prouve le lemme.

Posons maintenant

\begin{equation}\tag{11.6.9}
{}_n\Phi = ({}_nA)^f/({}_nA^o)^f \subset {}_n\Psi \quad ,
\end{equation}
de sorte que ${}_n\Phi$ est un sous-schéma fini ouvert de ${}_n\Psi$ , puisque $({}_nA)^f$


\SGApage{I}{459}

est un sous-schéma en groupes fini ouvert de ${}_nA$ ; plus précisément, on voit immédiatement qu'on a

\begin{equation}\tag{11.6.10}
{}_n\Phi \simeq ({}_n\Psi)^f \quad .
\end{equation}
Composant (11.6.9) et (11.6.8.1), on trouve une immersion ouverte

\begin{equation}\tag{11.6.11}
{}_n\Phi \hookrightarrow \underline{M}_n \quad .
\end{equation}
D'autre part, on a

\begin{equation*}
{}_n\Phi \times_S s = {}_n\Psi \times_S s = {}_nA_o/{}_nA_o^o \quad ,
\end{equation*}
et comme $A_o^o$ est $n$-divisible (2.2.1), on voit immédiatement que le dernier membre s'identifie à ${}_n\Phi_o$ ; donc on trouve un isomorphisme

\begin{equation}\tag{11.6.12}
{}_n\Phi \times_S s \simeq {}_n\Phi_o \quad ,
\end{equation}
qui montre que ${}_n\Phi$ est défini (à isomorphisme unique près) comme le schéma fini étale sur $S$ dont la fibre spéciale est ${}_n\Phi_o$ . Ceci dit, pour définir l'isomorphisme cherché (11.6.4) via l'immersion ouverte (11.6.11), il suffit de prouver le

\underline{Lemme} 11.6.13. \underline{L'image de l'homomorphisme canonique} (11.6.11) \underline{est égale au noyau de} $\tilde{u}_n : \underline{M}_n \longrightarrow \check{\underline{M}}'_n$ .

La compatibilité voulue des homomorphismes (11.6.11) avec les homomorphismes de transition pour $n$ variable est d'ailleurs triviale sur les définitions, de sorte qu'avec 11.6.13, la démonstration de 11.5 sera achevée.


\SGApage{I}{460}

Pour prouver le lemme, notons que

\begin{equation*}
K_n = \ker \tilde{u}_n
\end{equation*}
est aussi l'annulateur à gauche de l'accouplement canonique

\begin{equation*}
u_n : \underline{M}_n \times \underline{M}'_n \longrightarrow (\mathbb{Z}/n\mathbb{Z})_S \quad ,
\end{equation*}
dont la définition explicite est donnée dans 9.4 et 9.6, en termes des extensions

\begin{equation}\tag{11.6.14}
\begin{tikzcd}[column sep=small, row sep=large]
0 \arrow[r] \& \check{\underline{M}}'_n(1) \arrow[r] \& F_n \arrow[r] \arrow[d, equal] \& R_n \arrow[r] \& 0 \& , \& 0 \arrow[r] \& \check{\underline{M}}_n(1) \arrow[r] \& F'_n \arrow[r] \arrow[d, equal] \& R'_n \arrow[r] \& 0 \\
\& \& ({}_nA^o)^f \& \& \& \& \& \& ({}_nA'^o)^f \& \&
\end{tikzcd}
\end{equation}
















\SGAvar{(où $R_n$ et $R'_n$ sont duals de Cartier l'un de l'autre), et de l'extension panachée $X_{n\eta}$ de $E_{n\eta}$ par $F_{n\eta}$ , fibre générique de (11.7.5), où on pose

\begin{equation}\tag{11.6.15}
\begin{tikzcd}[column sep=large, row sep=large]
E_n = D(F_n) \; , \& 0 \arrow[r] \& R_n \arrow[r] \arrow[d, dash, "\wr"'] \& E_n \arrow[r] \& \underline{M}_n \arrow[r] \& 0 \& . \\
\& \& D(R'_n) \& \& \& \&
\end{tikzcd}
\end{equation}











Il résulte alors immédiatement de 9.4.6 que $K_n$ est le plus grand des sous-groupes finis étales $K'_n$ de $\underline{M}_n$ tels que l'image inverse $\overline{X}_{n\eta}$ de $K'_n$ dans l'extension panachée $X_{n\eta}$, regardée comme extension panachée (de $\overline{E}_{n\eta}$ par $F_{n\eta}$ , où $\overline{E}_n$ est le sous-groupe extension de $K_n$ par $R_n$ qu'on devine), se prolonge en une extension panachée de $\overline{E}_n$ par $F_n$. Or cette condition est évidemment remplie pour le sous-groupe ${}_n\Phi$ (11.6.11) de $\underline{M}_n$ , puisqu'on dispose de l'extension panachée $({}_nA)^f$ correspondante sur $S$. On a donc ${}_n\Phi \subset K_n$ . Pour prouver l'inclusion opposée, considérons l'ex
\SGApage{I}{461}
tension panachée $\overline{X}_n$ de $\overline{E}_n$ par $F_n$ prolongeant l'extension panachée $\overline{X}_{n\eta}$ , définie comme ci-dessus en y prenant $K'_n = K_n$ . C'est donc une extension

}{(où $R_n$ et $R'_n$ sont duals de Cartier l'un de l'autre), et de l'extension panachée $X_{n\eta}$ de $E_{n\eta}$ par $F_{n\eta}$ , fibre générique de (11.6.5), où on pose

\begin{equation}\tag{11.6.15}
\begin{tikzcd}[column sep=large, row sep=large]
E_n = D(F'_n) \; , \& 0 \arrow[r] \& R_n \arrow[r] \arrow[d, dash, "\wr"'] \& E_n \arrow[r] \& \underline{M}_n \arrow[r] \& 0 \& . \\
\& \& D(R'_n) \& \& \& \&
\end{tikzcd}
\end{equation}











Il résulte alors immédiatement de 9.4.7 que $K_n$ est le plus grand des sous-groupes finis étales $K'_n$ de $\underline{M}_n$ tels que l'image inverse $\overline{X}_{n\eta}$ de $K'_n$ dans l'extension panachée $X_{n\eta}$, regardée comme extension panachée (de $\overline{E}_{n\eta}$ par $F_{n\eta}$ , où $\overline{E}_n$ est le sous-groupe extension de $K'_n$ par $R_n$ qu'on devine), se prolonge en une extension panachée de $\overline{E}_n$ par $F_n$. Or cette condition est évidemment remplie pour le sous-groupe ${}_n\Phi$ (11.6.11) de $\underline{M}_n$ , puisqu'on dispose de l'extension panachée $({}_nA)^f$ correspondante sur $S$. On a donc ${}_n\Phi \subset K_n$ . Pour prouver l'inclusion opposée, considérons l'ex
\SGApage{I}{461}
tension panachée $\overline{X}_n$ de $\overline{E}_n$ par $F_n$ prolongeant l'extension panachée $\overline{X}_{n\eta}$ , définie comme ci-dessus en y prenant $K'_n = K_n$ . C'est donc une extension

}{F7I-IX-137}\begin{equation*}
0 \longrightarrow F_n \longrightarrow \overline{X}_n \longrightarrow K_n \longrightarrow 0 \quad ,
\end{equation*}
et on a un homomorphisme canonique

\begin{equation}\tag{*}
F_n = ({}_nA^o)^f \hookrightarrow A \quad ,
\end{equation}
qui sur la fibre générique est induit par l'homomorphisme canonique

\begin{equation}\tag{**}
\overline{X}_{n\eta} \hookrightarrow X_{n\eta} = {}_nA_K \hookrightarrow A_K \quad .
\end{equation}
Il est alors immédiat (5.9.2) que l'homomorphisme (**) se prolonge en un homomorphisme (unique)

\begin{equation}\tag{***}
\overline{X}_n \longrightarrow A
\end{equation}
induisant (*). Comme ce dernier applique la fibre générique dans ${}_nA$ , qui est un sous-schéma fermé de $A$, on conclut qu'il se factorise en $\overline{X}_n \longrightarrow {}_nA$, et comme $\overline{X}_n$ est fini, il se factorise par $({}_nA)^f$. Regardant l'homomorphisme induit sur les fibres génériques, on en conclut que $\overline{X}_{n\eta} \subset ({}_nA)^f_\eta$ i.e. $K_{n\eta} \subset {}_n\Phi_\eta$ i.e. $K_n \subset {}_n\Phi$ , cqfd.

\SGAvar{\underline{Remarques} 11.7. La construction faite dans 11.7 nous montre qu'on peut reconstruire à isomorphisme unique près l'extension panachée (11.6.5), et par suite aussi le sous-groupe $({}_nA)^f \supset ({}_nA^o)^f$ de ${}_nA$, par la connaissance des extensions (11.6.14) sur $S$ et de l'extension panachée $X_{n\eta} = ({}_nA_K)$ de $E_n$ par $F_n$ . En effet, en vertu de 9.4 on en conclut l'accouplement $u_n : \underline{M}_n \times \underline{M}'_n \longrightarrow (\mathbb{Z}/n\mathbb{Z})_S$, et la construction qui


}{\underline{Remarques} 11.7. La construction faite dans 11.6 nous montre qu'on peut reconstruire à isomorphisme unique près l'extension panachée (11.6.5), et par suite aussi le sous-groupe $({}_nA)^f \supset ({}_nA^o)^f$ de ${}_nA$, par la connaissance des extensions (11.6.14) sur $S$ et de l'extension panachée $X_{n\eta} = ({}_nA_K)$ de $E_n$ par $F_n$ . En effet, en vertu de 9.4 on en conclut l'accouplement $u_n : \underline{M}_n \times \underline{M}'_n \longrightarrow (\mathbb{Z}/n\mathbb{Z})_S$, et la construction qui


}{F7I-IX-138}\SGApage{I}{462}

\SGAvar{précède nous donne $({}_nA)^f$ comme l'extension panachée $\overline{X}_n$ de $\overline{E}_n$ ($=$ image inverse de $\mathrm{Ker}\,\widetilde{u}_n$ par $E_n \longrightarrow \underline{M}_n$) par $F_n$ qui prolonge l'extension panachée $\overline{X}_{n\eta}$ induite par $X_{n\eta}$ ; on notera (9.4.8) que cette extension panachée $\overline{X}_n$ se trouve déterminée ainsi \underline{à isomorphisme unique près}. On en déduit alors ${}_nA$ en recollant $\overline{X}_n$ et $X_{n\eta}$ suivant l'ouvert commun $\overline{X}_{n\eta}$ .

}{précède nous donne $({}_nA)^f$ comme l'extension panachée $\overline{X}_n$ de $\overline{E}_n$ ($=$ image inverse de $\mathrm{Ker}\,\widetilde{u}_n$ par $E_n \longrightarrow \underline{M}_n$) par $F_n$ qui prolonge l'extension panachée $\overline{X}_{n\eta}$ induite par $X_{n\eta}$ ; on notera (9.4.7) que cette extension panachée $\overline{X}_n$ se trouve déterminée ainsi \underline{à isomorphisme unique près}. On en déduit alors ${}_nA$ en recollant $\overline{X}_n$ et $X_{n\eta}$ suivant l'ouvert commun $\overline{X}_{n\eta}$ .

}{F7I-IX-139}Si on connaît donc les éléments de structure (9.6.1) et l'extension panachée $X_\eta = T_\ell(A_K)$ de $E_\eta$ par $F_\eta$ (ce qui revient à la connaissance des extensions (11.6.14) et de $X_{n\eta}$ pour tout $n$ de la forme $\ell^{r+1}$, $r \geq 0$), alors on sait reconstruire le système inductif (ou projectif, au choix) des extensions panachées (11.6.5).

11.8. Considérons, comme dans 8.3, la clôture séparable $\overline{K}$ de $K$ comme limite inductive de ses sous-extensions finies $K'$. Pour tout tel $K'$, nous considérons le modèle de Néron

\begin{equation*}
A[K']
\end{equation*}
de $A_{K'} = A_K \otimes_K K'$, d'où un groupe fini étale

\SGAvar{\begin{equation}\tag{11.8.1}
\Phi_o[K'] = A[K']_o, A[K']_o^o
\end{equation}
de composantes connexes, défini sur le corps résiduel du normalisé $S[K']$ de $S$ dans $K'$. Supposons pour simplifier que le corps résiduel de $S$ soit séparablement clos, de sorte qu'il en est de même de ceux des $S[K']$. Les groupes (11.8.1) s'identifient alors à des groupes finis ordinaires $\Phi[K']$, et il est immédiat que pour $K'$ variable, ceux-ci forment un système inductif, que nous nous proposons d'étudier. Nous nous intéresserons notamment au groupe


}{\begin{equation}\tag{11.8.1}
\Phi_o[K'] = A[K']_o/A[K']_o^o
\end{equation}
de composantes connexes, défini sur le corps résiduel du normalisé $S[K']$ de $S$ dans $K'$. Supposons pour simplifier que le corps résiduel de $S$ soit séparablement clos, de sorte qu'il en est de même de ceux des $S[K']$. Les groupes (11.8.1) s'identifient alors à des groupes finis ordinaires $\Phi[K']$, et il est immédiat que pour $K'$ variable, ceux-ci forment un système inductif, que nous nous proposons d'étudier. Nous nous intéresserons notamment au groupe


}{F7I-IX-140}\SGApage{I}{463}

\begin{equation}\tag{11.8.2}
\Psi \;=\; \varinjlim_{K'} \Phi[K'] \qquad .
\end{equation}
Nous savons (3.6) que pour $K'$ assez grand, disons $K' \supset K'_o$ , $A[K']$ est à réduction semi-stable ; donc pour les $K' \supset K'_o$ , les morphismes de transition sont injectifs (3.3). D'ailleurs, avec les notations de 5.5 reprises dans le présent numéro, on aura pour $K' \supset K'_o$ un isomorphisme canonique

\SGAvar{\begin{equation}\tag{11.8.3}
\Psi[K'](\ell) \;\simeq\; \ker(\widetilde{u}[K']_\ell \otimes D_\ell \;:\; M_\ell \otimes D_\ell \longrightarrow \check{M}'_\ell \otimes D_\ell) \qquad ,
\end{equation}
provenant de (11.6.2) appliqué à $A[K']$. On vérifie de plus immédiatement que les immersions correspondantes

}{\begin{equation}\tag{11.8.3}
\Phi[K'](\ell) \;\simeq\; \ker(\widetilde{u}[K']_\ell \otimes D_\ell \;:\; M_\ell \otimes D_\ell \longrightarrow \check{M}'_\ell \otimes D_\ell) \qquad ,
\end{equation}
provenant de (11.6.2) appliqué à $A[K']$. On vérifie de plus immédiatement que les immersions correspondantes

}{F7I-IX-141}\begin{equation*}
\Phi[K'](\ell) \hookrightarrow M_\ell \otimes D_\ell
\end{equation*}
sont compatibles avec les morphismes de transition pour $K'$ variable, d'où par passage à la limite un monomorphisme canonique

\begin{equation}\tag{11.8.4}
\Psi(\ell) \;\overset{\mathrm{dfn}}{=}\; \varinjlim_{K'} \Phi[K'](\ell) \hookrightarrow M_\ell \otimes D_\ell
\end{equation}
de la composante $\ell$-primaire de $\Psi$ ; d'où, en faisant la somme sur tous les $\ell$ :

\begin{equation}\tag{11.8.5}
\Psi \;\overset{\mathrm{dfn}}{=}\; \varinjlim_{K'} \Phi[K'] \longrightarrow M \otimes \mathbb{Q}/\mathbb{Z} \qquad .
\end{equation}
\underline{Théorème} 11.9. \underline{L'homomorphisme canonique} (11.8.5) \underline{est un isomorphisme}.

On sait déjà que c'est un monomorphisme, il reste à prouver que c'est un isomorphisme, et pour ceci il nous faudra utiliser la caractérisation (11.8.3) de l'image de $\Phi[K'](\ell)$, et le comportement de $\widetilde{u}(K')$ pour $K'$ variable. Ce comportement est donné par (10.3.4).


\SGAvar{\SGApage{I}{464}

Se ramenant au cas où $S$ est strictement local, ce qui est loisible, on est donc ramené, pour prouver 11.9, à l'énoncé bien connu suivant (qu'on appliquera à $S(K'_o)$) : si $S$ est un trait strictement local de corps des fractions $K$, alors pour tout entier $n \geqslant 1$, il existe une extension séparable $K'$ de $K$ de degré multiple de $n$ (cf. I\hspace{3em}).

}{\SGApage{I}{464}


Se ramenant au cas où $S$ est strictement local, ce qui est loisible, on est donc ramené, pour prouver 11.9, à l'énoncé bien connu suivant (qu'on appliquera à $S(K'_o)$) : si $S$ est un trait strictement local de corps des fractions $K$, alors pour tout entier $n \geqslant 1$, il existe une extension séparable $K'$ de $K$ de degré multiple de $n$.

}{F7I-IX-142}\underline{Corollaire} 11.10. \underline{Sous les conditions} ($S$ \underline{complet}) \underline{et avec les notations de} 8.3, \underline{le schéma en groupes} $A^{\flat}_K/A^{\flat\,o}_K$ \underline{ind-fini sur} $K$ \underline{des composantes connexes du schéma en groupes} $A^{\flat}_K$ \underline{est canoniquement isomorphe à} $\underline{M}_K \otimes \mathbb{Q}/\mathbb{Z}$ , \underline{où} $\underline{M}_K$ \underline{est le groupe constant tordu sur} $K$ \underline{défini dans} (8.1.8).

\SGAvar{C'est une conséquence immédiate de (8.2.4) et de 11.9, compte tenu de la définition de $\underline{M}_K$ .

}{C'est une conséquence immédiate de (8.3.4) et de 11.9, compte tenu de la définition de $\underline{M}_K$ .

}{F7I-IX-143}\underline{Remarques} 11.11.  a) Lorsqu'on se borne aux composantes $\ell$-primaires pour $\ell \neq p$, le théorème 11.9 est immédiat en termes de 11.2, et ne demande pas le recours aux considérations plus délicates de 11.6 (via 9.4). De même, lorsqu'on suppose $K$ \underline{de caractéristique nulle}, 11.9 se démontre sans utiliser (11.8.3), car il est immédiat que l'image de ${}_n\Phi(K')(\ell)$ dans $M_n = {}_n(M_\ell \otimes D_\ell)$ contient l'image du groupe ${}_nA(K')$ des points de ${}_nA$ à valeurs dans $K'$, qui est égal à ${}_nA(\overline{K})$ donc se surjecte sur $M_n$ pour $K'$ assez grand. Lorsque car.$K = p$, il semble par contre qu'il n'y ait pas de démonstration ``triviale'' (i.e. n'utilisant pas 9.4)de 11.10 pour la composante $p$-primaire.

b) Le fait qu'on ait un isomorphisme $\Psi \simeq (\mathbb{Q}/\mathbb{Z}\,)^r$, et la valeur correcte de $r$, avaient été conjecturés en 1964 par J.P. SERRE.


\SGApage{I}{465}


\SGAsectiontoc{12. Modèle de Néron d'une jacobienne et formule de Picard Lefschetz}{12}
12. \underline{Modèle de Néron d'une jacobienne et formule de Picard Lefschetz}

Rappelons d'abord les résultats généraux suivants de M. RAYNAUD [22, th.5 et th.7 , 22 bis] :

\underline{Théorème} 12.1 (M.RAYNAUD). \underline{Soient} $S$ \underline{un trait strictement local}, $f:X \to S$ \underline{un schéma propre et plat sur} $S$, \underline{à fibres de dimension} 1, \underline{tel que l'on ait}

\begin{equation}\tag{12.1.1}
f_*(\underline{O}_X) \simeq \underline{O}_S
\end{equation}
\underline{et que les anneaux locaux de} $X$ \underline{soient factoriels}. \underline{Pour tout point maximal} $x_i$ \underline{de} $X_o$, \underline{soient}

\SGAvar{\begin{equation}\tag{12.1.2}
d_i = \mathrm{long}\ \underline{O}_{X_o,x} \ , \ d_i^t = \mu_i \ d_i \quad ,
\end{equation}
\underline{où} $\mu_i$ \underline{est la multiplicité radicielle de} $k(x_i)$ \underline{sur} $k$ (EGA IV 4.7.4) (\underline{qui est une puissance de l'exposant caractéristique de} $k$). \underline{Soit} $d$ (resp. $d^t$) \underline{le pgcd des entiers} $d_i$ (resp. $d_i^t$) ; \underline{donc on a} $d = d^t$ \underline{si} $k$ \underline{est parfait}. \underline{Posons}

}{\begin{equation}\tag{12.1.2}
d_i = \mathrm{long}\ \underline{O}_{X_o,x_i} \ , \ d_i^t = \mu_i \ d_i \quad ,
\end{equation}
\underline{où} $\mu_i$ \underline{est la multiplicité radicielle de} $k(x_i)$ \underline{sur} $k$ (EGA IV 4.7.4) (\underline{qui est une puissance de l'exposant caractéristique de} $k$). \underline{Soit} $d$ (resp. $d^t$) \underline{le pgcd des entiers} $d_i$ (resp. $d_i^t$) ; \underline{donc on a} $d = d^t$ \underline{si} $k$ \underline{est parfait}. \underline{Posons}

}{F7I-IX-144}\begin{equation}\tag{12.1.3}
P = \underline{\mathrm{Pic}}_{X/S} : (\mathrm{Sch})_{/S}^{\,o} \longrightarrow (\mathrm{Ab}) \quad .
\end{equation}
a) \underline{Considérons les conditions suivantes} :

1) $d^t=1$.

2) $X$ \underline{possède un} $0$-\underline{cycle de degré} 1.

3) \underline{Il existe un Module inversible} $L$ \underline{sur} $X \times P^o$ \underline{tel que l'homomorphisme correspondant} $P^o_\eta \longrightarrow \underline{\mathrm{Pic}}_{X_\eta/\eta} = P_\eta$ \underline{soit l'inclusion canonique}. (NB \underline{on sait} (J.P. MURRE) \underline{que} $P_\eta$ \underline{est représentable par un schéma en groupes localement de type fini sur} $\eta$ , cf. p.ex. SGA 6 XII 1.5).


\SGApage{I}{466}

4) $X$ \underline{est cohomologiquement plat sur} $S$ (EGA III 7.8.1) (\underline{ou encore} : \underline{cohomologiquement plat en dimension} $0$, \underline{ou encore} : \underline{on a} $k \stackrel{\sim}{\longrightarrow} H^o(X_o,\underline{O}_{X_o})$), \underline{et} $d=1$.

5) $P$ \underline{est représentable} (\underline{ou encore} : \underline{représentable par un schéma en groupes lisse sur} $S$).

6) $P^o$ \underline{est représentable} (\underline{ou encore} : \underline{représentable par un schéma en groupes lisse sur} $S$). (NB \underline{Pour la définition de} $P^o$, cf. SGA 3 $\mathrm{VI}_B$ 3.1.)

7) $P^o$ \underline{est séparé sur} $S$, \underline{i.e.} \underline{satisfait au critère valuatif de séparation} EGA II 7.2.3 (\underline{ou encore} : \underline{toute section} $s$ \underline{de} $P^o$ \underline{sur} $S$ \underline{telle que} $s(\eta) = 0$ \underline{est nulle}).

8) $d=1$ .

\underline{On a alors les implications}

\begin{equation}\tag{12.1.4}
1) \Longleftrightarrow 2) \Longrightarrow 3) \Longrightarrow 4) \Longleftrightarrow 5) \Longleftrightarrow 6) \Longleftrightarrow 7) \Longrightarrow 8) \qquad .
\end{equation}
\underline{En particulier}, \underline{si} $d=d^t$, \underline{par exemple si} $k$ \underline{est parfait}, \underline{les conditions} 1) \underline{à} 8) \underline{sont équivalentes}.

b) \underline{Supposons} $d^t=1$, \underline{de sorte que} $P$ \underline{est représentable par un schéma en groupes lisse sur} $S$. \underline{Alors l'adhérence schématique de la section unité de} $P$ \underline{est un sous-schéma en groupes} $E$ \underline{étale sur} $S$, \underline{et si} $I$ \underline{est l'ensemble des composantes irréductibles réduites} $X_o^i$ \underline{de} $X_o$, \underline{on obtient un isomorphisme}

\begin{equation}\tag{12.1.5}
\mathbb{Z}^I \simeq \Gamma(S,E)
\end{equation}

\SGApage{I}{467}

\SGAvar{\underline{en associant à tout} $i \in I$ \underline{la section} $S_i$ \underline{de} $P$ \underline{sur} $S$ (\underline{nulle en} $\eta$) \underline{définie par le faisceau inversible} $\underline{O}_X(X_o^i)$ \underline{sur} $S$. (NB \underline{Comme les anneaux locaux de} $X$ \underline{sont factoriels en hypothèse, les cycles} $X_o^i$ \underline{sont bien des diviseurs, ce qui donne un sens à} $\underline{O}_X(X_i^o)$ \underline{comme le faisceau inversible associé au diviseur en question} (EGA IV 21.2.8.1).) \underline{On a}

}{\underline{en associant à tout} $i \in I$ \underline{la section} $S_i$ \underline{de} $P$ \underline{sur} $S$ (\underline{nulle en} $\eta$) \underline{définie par le faisceau inversible} $\underline{O}_X(X_o^i)$ \underline{sur} $X$. (NB \underline{Comme les anneaux locaux de} $X$ \underline{sont factoriels en hypothèse, les cycles} $X_o^i$ \underline{sont bien des diviseurs, ce qui donne un sens à} $\underline{O}_X(X_o^i)$ \underline{comme le faisceau inversible associé au diviseur en question} (EGA IV 21.2.8.1).) \underline{On a}

}{F7I-IX-145}\begin{equation}\tag{12.1.6}
P^o \cap E = \underline{\mbox{section nulle de}}\ P
\quad .
\end{equation}
c) \underline{Sous les conditions de} b), \underline{le quotient}

\SGAvar{\begin{equation}\tag{12.1.7}
Q = P/E
\end{equation}
\underline{est représentable par un schéma en groupes lisse et séparé sur} $S$, \underline{et ce dernier est} néronien, \underline{i.e. satisfait la propriété universelle} (1.1.2) \underline{en termes de sa fibre spéciale}.

}{\begin{equation}\tag{12.1.7}
Q = P/E
\end{equation}
\underline{est représentable par un schéma en groupes lisse et séparé sur} $S$, \underline{et ce dernier est} néronien, \underline{i.e. satisfait la propriété universelle} (1.1.2) \underline{en termes de sa fibre générique}.

}{F7I-IX-146}d) \underline{Sous les conditions de} c), \underline{supposons de plus que la fibre géométrique générique de} $X$ \underline{soit} intègre. \underline{Alors on a une suite exacte canonique}

\begin{equation}\tag{12.1.8}
0 \longrightarrow A \longrightarrow Q \xrightarrow{\ \deg\ } \mathbb{Z}_S \longrightarrow 0 \quad ,
\end{equation}
\underline{où} $\deg$ \underline{est induit par l'homomorphisme ``degré''} $P \longrightarrow \mathbb{Z}_S$, \underline{et où} $A$ \underline{est un schéma en groupes lisse séparé et} de type fini \underline{sur} $S$, \underline{qui est également néronien i.e. qui est un modèle de Néron de}

\begin{equation}\tag{12.1.9}
A_\eta = P_\eta^o = \underline{\mathrm{Pic}}^o_{X_\eta/\eta} \quad .
\end{equation}
\underline{D'autre part, l'homomorphisme composé} $P^o \longrightarrow P \to Q$ \underline{induit un isomorphisme}

\begin{equation}\tag{12.1.10}
P^o \xrightarrow{\ \sim\ } A^o \quad .
\end{equation}

\SGApage{I}{468}

12.1.11. Le cas qui nous intéressera est celui où $X_\eta$ est lisse sur $\eta$ , de sorte que (compte tenu de 12.1.1) la fibre géométrique générique de $X$ est lisse et connexe (donc intègre), et où $X_\eta$ possède un $0$-cycle de degré 1. Alors toutes les conditions envisagées dans 12.1 (et notamment les conditions 1) à 8) de a)) sont vérifiées, d'autre part $A_\eta = P_\eta^o$ est alors un schéma abélien, et d) nous donne une construction très explicite de son modèle de Néron $A$. En particulier (12.1.10) nous fournit un isomorphisme

\begin{equation}\tag{12.1.12}
A_o^{\,o} \simeq \underline{\mathrm{Pic}}^o_{X_o/k} \quad ,
\end{equation}
qui montre en particulier que $A$ \underline{a réduction semi-stable sur} $S$ \underline{si et seulement si} $\underline{\mathrm{Pic}}^o_{X_o/k}$ \underline{a un rang unipotent nul}. Lorsque $X$ est régulier et que $X_o$ \underline{est séparable sur} $k$, on sait qu'il en est ainsi si et seulement si la fibre géométrique spéciale $\overline{X}_o$ n'a comme seuls points singuliers que des points quadratiques ordinaires (i.e. ayant même anneau local complété que la réunion des deux axes de coordonnées de l'espace affine $\mathbf{E}^2$ en l'origine). La vérification de ce fait, allant avec la description ``bien connue'' du schéma de Picard d'une courbe propre séparable sur un corps $k$, en utilisant EGA IV 21.8.5, est laissée au lecteur à titre d'exercice. On notera d'ailleurs que le fait que ``$\overline{X}_o$ à singularités quadratiques ordinaires'' implique ``réduction semi-stable de $A_\eta = \underline{\mathrm{Pic}}^o_{X_\eta/\eta}$'' est aussi une conséquence immédiate du ``théorème de monodromie géométrique'' III \hspace{3.5em} , compte tenu du critère galoisien de réduction semi-stable 3.5 ; mais cet argument ne donne pas l'important isomorphisme (12.1.10) et sa conséquence (12.1.12).


\SGApage{I}{469}

\SGAvar{12.2. Pour la démonstration de 12.1, nous renvoyons à [22 bis] ; nous n'aurons besoin d'ailleurs que du cas $X_\eta$ lisse sur $\eta$. Notons d'ailleurs que dans [22bis] sont démontrées des relations entre schémas de Picard et modèles de Néron sous les conditions nettement plus générales que celles de 12.1 c) et d) ; comme leur énoncé est plus technique, nous avons renoncé à les inclure dans notre rappel. Signalons seulement que b), c), d) sont des conséquences faciles de a), et que dans a) les seuls points délicats sont les implications 3) $\Longrightarrow$ 4) $\Longrightarrow$ 5) et 7) $\Longrightarrow$ 4), le plus difficile étant le premier. Ce point est trivial par contre si on suppose que $X_o$ est séparable sur $k$ (qui implique que $X$ admet une section sur $S$ puisque $S$ est strictement local, donc que la condition 1), $d^t=1$, est vérifiée). Nous conseillons au lecteur de faire la démonstration dans ce cas particulièrement simple, en supposant de plus, au besoin, que $X_\eta$ est lisse donc $P^o_\eta$ un schéma abélien, ce qui permet alors de disposer du modèle de Néron $A$ de ce dernier, et de $A^o$ pour représenter $\underline{\mathrm{Pic}}^o_{X/S}$ . Nous n'utiliserons d'ailleurs 12.1 que dans ce cas particulier, et lorsque l'on suppose de plus que $\overline{X}_o$ n'a d'autres points singuliers que des points quadratiques ordinaires, i.e. (12.1.11) que $A_\eta$ est à réduction semi-stable : c'est le cas envisagé dans 12.5.4, auquel nous nous bornons à partir de maintenant.

}{12.2. Pour la démonstration de 12.1, nous renvoyons à [22 bis] ; nous n'aurons besoin d'ailleurs que du cas $X_\eta$ lisse sur $\eta$. Notons d'ailleurs que dans [22bis] sont démontrées des relations entre schémas de Picard et modèles de Néron sous les conditions nettement plus générales que celles de 12.1 c) et d) ; comme leur énoncé est plus technique, nous avons renoncé à les inclure dans notre rappel. Signalons seulement que b), c), d) sont des conséquences faciles de a), et que dans a) les seuls points délicats sont les implications 3) $\Longrightarrow$ 4) $\Longrightarrow$ 5) et 7) $\Longrightarrow$ 4), le plus difficile étant le premier. Ce point est trivial par contre si on suppose que $X_o$ est séparable sur $k$ (qui implique que $X$ admet une section sur $S$ puisque $S$ est strictement local, donc que la condition 1), $d^t=1$, est vérifiée). Nous conseillons au lecteur de faire la démonstration dans ce cas particulièrement simple, en supposant de plus, au besoin, que $X_\eta$ est lisse donc $P^o_\eta$ un schéma abélien, ce qui permet alors de disposer du modèle de Néron $A$ de ce dernier, et de $A^o$ pour représenter $\underline{\mathrm{Pic}}^o_{X/S}$ . Nous n'utiliserons d'ailleurs 12.1 que dans ce cas particulier, et lorsque l'on suppose de plus que $\overline{X}_o$ n'a d'autres points singuliers que des points quadratiques ordinaires, i.e. (12.1.11) que $A_\eta$ est à réduction semi-stable : c'est le cas envisagé dans 10.5.4, auquel nous nous bornons à partir de maintenant.

}{F7I-IX-147}12.3. Soit $C$ une courbe séparable et propre sur un corps $k$, que nous supposons pour simplifier séparablement clos (ou même algébriquement clos, si le lecteur préfère). Si $\widetilde{C}$ est la normalisée de $C$, on sait que


\SGApage{I}{470}

$\underline{\mathrm{Pic}}^o_{\widetilde{C}/k}$ est de ``rang réductif'' nul (i.e. son tore maximal est nul), de sorte que le tore maximal de $\underline{\mathrm{Pic}}^o_{C/k}$ est contenu dans

\SGAvar{\begin{equation*}
N = \mathrm{Ker}(\underline{\mathrm{Pic}}_{C/k} \longrightarrow \underline{\mathrm{Pic}}_{\widetilde{C}/k}) \quad ,
\end{equation*}
donc est isomorphe canoniquement au tore maximal de ce dernier groupe. Or la structure de $N$ s'explicite aisément à l'aide de EGA IV 21.8.5 : on trouve que $N$ est canoniquement une extension d'un tore $T$ dont on va expliciter la description, par un groupe unipotent connexe $U$ (qui se dévisse de même, sur $k$, avec une suite de composition à quotients isomorphes au groupe additif $G_a$). D'ailleurs, si $T'$ est le tore maximal de $N$, alors $T' \cap U$ est à la fois unipotent et de type multiplicatif, donc nul (SGA 3 XVII 2.4), et l'image de $T'$ dans $T$ est identique à $T$ (SGA 3 XII 7.1 e)), de sorte que $T' \longrightarrow T$ est un isomorphisme. Donc $T$ est canoniquement isomorphe au tore maximal de $N$, donc de $\underline{\mathrm{Pic}}^o_{X/k}$ .

}{\begin{equation*}
N = \mathrm{Ker}(\underline{\mathrm{Pic}}_{C/k} \longrightarrow \underline{\mathrm{Pic}}_{\widetilde{C}/k}) \quad ,
\end{equation*}
donc est isomorphe canoniquement au tore maximal de ce dernier groupe. Or la structure de $N$ s'explicite aisément à l'aide de EGA IV 21.8.5 : on trouve que $N$ est canoniquement une extension d'un tore $T$ dont on va expliciter la description, par un groupe unipotent connexe $U$ (qui se dévisse de même, sur $k$, avec une suite de composition à quotients isomorphes au groupe additif $G_a$). D'ailleurs, si $T'$ est le tore maximal de $N$, alors $T' \cap U$ est à la fois unipotent et de type multiplicatif, donc nul (SGA 3 XVII 2.4), et l'image de $T'$ dans $T$ est identique à $T$ (SGA 3 XII 7.1 e)), de sorte que $T' \longrightarrow T$ est un isomorphisme. Donc $T$ est canoniquement isomorphe au tore maximal de $N$, donc de $\underline{\mathrm{Pic}}^o_{C/k}$ .

}{F7I-IX-148}C'est pour expliciter la structure de $N$ qu'il sera commode de supposer $k$ séparablement clos, de sorte que les composantes irréductibles $C_i$ de $C$ sont géométriquement irréductibles, et que les points non normaux de $C$ sont radiciels sur $k$. Ceci dit, soit $J$ l'ensemble des composantes irréductibles de $\widetilde{C}$, et pour tout $x \in C$, soit $J(x)$ l'ensemble des ``branches'' de $C$ passant par $x$ (i.e. des points de $\widetilde{C}$ au-dessus de $x$) ; soient enfin

\begin{equation}\tag{12.3.1}
R(x) = \mathbb{Z}^{J(x)}/\mathrm{Im}\,\mathbb{Z} \quad ,
\end{equation}
où $\mathbb{Z} \longrightarrow \mathbb{Z}^{J(x)}$ est l'application diagonale, et

\begin{equation}\tag{12.3.2}
R = (\mathop{\oplus}\limits_{x} R_x)/\mathrm{Im}\,\mathbb{Z}^J \quad ,
\end{equation}

\SGApage{I}{471}

où il suffit évidemment d'étendre la somme à l'ensemble fini des points de C qui appartiennent à au moins deux branches de C. Donc $R$ est un $\mathbb{Z}$-module de type fini. On a alors un isomorphisme canonique

\SGAvar{\begin{equation}\tag{12.2.3}
T\ (\ \simeq\ \mbox{tore maximal de }\ \underline{\mathrm{Pic}}^{o}_{C/k})\ \simeq\ R \otimes \mathbb{G}_{mk}\qquad .
\end{equation}
En d'autres termes, le groupe des caractères $M$ de $T$ est donné par

\begin{equation}\tag{12.2.4}
M\ \stackrel{\mathrm{dfn}}{=}\ D(T) = \check{R} \simeq \mathrm{Ker}\ (\mathop{\oplus}\limits_{x} \check{R}(x) \longrightarrow \mathbb{Z}^{J})\ ,
\end{equation}
où $\check{R}(x)$ peut s'interpréter aussi comme l'ensemble des fonctions à valeurs entières sur $J(x)$ dont la somme des valeurs est nulle.

12.3.5. Nous aurons à utiliser seulement le cas où par un point $x$ de $C$ passent au plus deux branches (ce qui est le cas en particulier pour les points quadratiques ordinaires !). Pour un point tel que card $J(x) = 2$, $R(x)$ est donc isomorphe à $\mathbb{Z}$, mais l'isomorphisme dépendant du choix d'un ordre sur l'ensemble des deux branches $(C',C'')$ passant par $x$, ou ce qui revient au même, le choix de l'une de ses branches $C'$ (en associant à cellè-ci l'élément de $R(x)$ image canonique de $0_{C'} + 1_{C''}$ ). On trouve donc alors un isomorphisme

}{\begin{equation}\tag{12.3.3}
T\ (\ \simeq\ \mbox{tore maximal de }\ \underline{\mathrm{Pic}}^{o}_{C/k})\ \simeq\ R \otimes \mathbb{G}_{mk}\qquad .
\end{equation}
En d'autres termes, le groupe des caractères $M$ de $T$ est donné par

\begin{equation}\tag{12.3.4}
M\ \stackrel{\mathrm{dfn}}{=}\ D(T) = \check{R} \simeq \mathrm{Ker}\ (\mathop{\oplus}\limits_{x} \check{R}(x) \longrightarrow \mathbb{Z}^{J})\ ,
\end{equation}
où $\check{R}(x)$ peut s'interpréter aussi comme l'ensemble des fonctions à valeurs entières sur $J(x)$ dont la somme des valeurs est nulle.

12.3.5. Nous aurons à utiliser seulement le cas où par un point $x$ de $C$ passent au plus deux branches (ce qui est le cas en particulier pour les points quadratiques ordinaires !). Pour un point tel que card $J(x) = 2$, $R(x)$ est donc isomorphe à $\mathbb{Z}$, mais l'isomorphisme dépendant du choix d'un ordre sur l'ensemble des deux branches $(C',C'')$ passant par $x$, ou ce qui revient au même, le choix de l'une de ses branches $C'$ (en associant à celle-ci l'élément de $R(x)$ image canonique de $0_{C'} + 1_{C''}$ ). On trouve donc alors un isomorphisme

}{F7I-IX-149}\begin{equation}\tag{12.3.6}
R = \mathbb{Z}^{I}/\mathrm{Im}\ \mathbb{Z}^{J},\ \mbox{donc}\ T \simeq \mathbb{G}_{m}^{I}/\mathrm{Im}\ \mathbb{G}_{m}^{J}\ ,\ M = \mathrm{Ker}\ (\mathbb{Z}^{I} \longrightarrow \mathbb{Z}^{J})\ ,
\end{equation}
où $I$ désigne l'ensemble des $x \in C$ appartenant à deux branches de $C$, et moyennant un choix de signes comme on vient de l'expliciter.

\underline{Remarques} 12.3.7. Une façon parfois commode d'expliciter le groupe des caractères $M$ du tore maximal $T$ de $\underline{\mathrm{Pic}}^{o}_{C/k}$ est la suivante : on considère le "graphe (non orienté) de dimension 1 " (N. BOURBAKI, Groupes et Algèbres


\SGApage{I}{472}

\SGAvar{de Lie, Chap IV, Annexe) $\Gamma(C)$ dont les sommets sont les composantes irréductibles de $\widetilde{C}$, deux sommets $\widetilde{C}_i$ et $\widetilde{C}_j$ étant reliés par un ensemble d'arrêtes en correspondance bijective avec l'ensemble $\widetilde{C}_i \times_C \widetilde{C}_j$ des ensembles de deux branches \underline{distinctes} de $C$ contenues dans $\widetilde{C}_i$ resp. $\widetilde{C}_j$ passant par le même point de $C$, ensemble égal à $\widetilde{C}_i \times_C \widetilde{C}_j$ si $i \neq j$, à $\widetilde{C}_i \times_C \widetilde{C}_i - \Delta(\widetilde{C}_i)$ modulo symétrie ($\Delta : \widetilde{C} \longrightarrow \widetilde{C} \times_C \widetilde{C}$ étant la diagonale) si $i=j$ (donc vide si $C_i$ est géométriquement unibranche). Ceci posé, on a des isomorphismes canoniques

}{de Lie, Chap IV, Annexe) $\Gamma(C)$ dont les sommets sont les composantes irréductibles de $\widetilde{C}$, deux sommets $\widetilde{C}_i$ et $\widetilde{C}_j$ étant reliés par un ensemble d'arêtes en correspondance bijective avec l'ensemble $\widetilde{C}_i \times_C \widetilde{C}_j$ des ensembles de deux branches \underline{distinctes} de $C$ contenues dans $\widetilde{C}_i$ resp. $\widetilde{C}_j$ passant par le même point de $C$, ensemble égal à $\widetilde{C}_i \times_C \widetilde{C}_j$ si $i \neq j$, à $\widetilde{C}_i \times_C \widetilde{C}_i - \Delta(\widetilde{C}_i)$ modulo symétrie ($\Delta : \widetilde{C} \longrightarrow \widetilde{C} \times_C \widetilde{C}$ étant la diagonale) si $i=j$ (donc vide si $C_i$ est géométriquement unibranche). Ceci posé, on a des isomorphismes canoniques

}{F7I-IX-150}\begin{equation}\tag{12.3.7}
M = D(T) \stackrel{\sim}{\longleftarrow} H_1(\Gamma(C),\mathbb{Z}\ )\ ,\ \mbox{donc}\ \check{M} = R \simeq H^{1}(\Gamma(C),\mathbb{Z}\ )\qquad .
\end{equation}
Nous n'aurons pas besoin de cette interprétation combinatoire de $M$, et allons nous borner à définir l'homomorphisme canonique

\SGAvar{\begin{equation*}
H_1(\Gamma(C),\mathbb{Z}\ ) \longrightarrow M \qquad ,
\end{equation*}
laissant au lecteur le soin de prouver que c'est un isomorphisme. Soit donc $I$ l'ensemble des points de $C$ appartenant a au moins deux branches de $C$, $J$(comme ci-dessus) l'ensemble des composantes irréductibles de $\widetilde{C}$ (ou de $C$, cela revient au même), $K$ l'ensemble des "branches" de $C$ en les points de $I$, i.e. l'image inverse de $I$ dans $\widetilde{C}$ . On a donc une application canonique

}{\begin{equation}\tag{12.3.8}
H_1(\Gamma(C),\mathbb{Z}\ ) \longrightarrow M \qquad ,
\end{equation}
laissant au lecteur le soin de prouver que c'est un isomorphisme. Soit donc $I$ l'ensemble des points de $C$ appartenant à au moins deux branches de $C$, $J$(comme ci-dessus) l'ensemble des composantes irréductibles de $\widetilde{C}$ (ou de $C$, cela revient au même), $K$ l'ensemble des "branches" de $C$ en les points de $I$, i.e. l'image inverse de $I$ dans $\widetilde{C}$ . On a donc une application canonique

}{F7I-IX-151}\begin{equation}\tag{12.3.9}
h : K \longrightarrow I \times J\ ,\ \ h = (f,g)\qquad ,
\end{equation}
et la description (12.3.4) de $M$ s'exprime en termes de $h$ comme

\SGAvar{\begin{equation}\tag{12.3.10}
M \simeq \mathrm{Ker}\ (\mathbb{Z}^{(h)} : \mathbb{Z}^{(K)} \longrightarrow \mathbb{Z}^{(I \times J)})\qquad ,
\end{equation}
car l'expression (12.3.3) de $R = \check{M}$ peut s'interprèter manifestement


}{\begin{equation}\tag{12.3.10}
M \simeq \mathrm{Ker}\ (\mathbb{Z}^{(h)} : \mathbb{Z}^{(K)} \longrightarrow \mathbb{Z}^{(I \times J)})\qquad ,
\end{equation}
car l'expression (12.3.2) de $R = \check{M}$ peut s'interpréter manifestement


}{F7I-IX-152}\SGApage{I}{473}

\SGAvar{comme le conoyau de l'homomorphisme $\mathbb{Z}^{I \times J} \longrightarrow \mathbb{Z}^{K}$ induit par $h$, lequel est transposé de $\mathbb{Z}^{(h)}$. Notons d'ailleurs que le graphe $\Gamma(C)$ s'explicite aussi de façon évidente en termes de $h$, par $\Gamma_o = J$, $\Gamma_1 = (K \times_I K - \Delta(K))/$symétrie. D'autre part, choisissant une orientation sur chacune des arrêtes de $\Gamma(C)$ , et désignant par $\Gamma_1$ l'ensemble de ces arrêtes $(x_i,C',C'')$, on a

}{comme le conoyau de l'homomorphisme $\mathbb{Z}^{I \times J} \longrightarrow \mathbb{Z}^{K}$ induit par $h$, lequel est transposé de $\mathbb{Z}^{(h)}$. Notons d'ailleurs que le graphe $\Gamma(C)$ s'explicite aussi de façon évidente en termes de $h$, par $\Gamma_o = J$, $\Gamma_1 = (K \times_I K - \Delta(K))/$symétrie. D'autre part, choisissant une orientation sur chacune des arêtes de $\Gamma(C)$ , et désignant par $\Gamma_1$ l'ensemble de ces arêtes $(x_i,C',C'')$, on a

}{F7I-IX-153}\begin{equation}\tag{12.3.11}
H_1(\Gamma(C),\mathbb{Z}\,) = \mathrm{Ker}\ (d : \mathbb{Z}^{(\Gamma_1)} \longrightarrow \mathbb{Z}^{(J)})
\end{equation}
où l'opérateur différentiel est donné par

\begin{equation}\tag{12.3.12}
d(x,C',C'') = g(C'') - g(C') \quad .
\end{equation}
Ceci posé, l'homomorphisme (12.3.8) est induit par l'homomorphisme $v$ de complexes de chaînes :

\SGAvar{\begin{equation}\tag{12.3.13}
v = (v_o,v_1) \ : \ C.(\Gamma(C)) = \left[\ \mathbb{Z}^{(\Gamma)} \longrightarrow \mathbb{Z}^{(J)}\ \right] \to \left[\ \mathbb{Z}^{(K)} \to \mathbb{Z}^{(I \times J)}\ \right]
\end{equation}
donné par

\begin{equation}\tag{12.3.14}
v_1(x,C',C'') = (x,C'') - (x,C'), \quad v_o(\widetilde{C}_j) = (0,\ \widetilde{C}_j) \quad .
\end{equation}
12.4. Nous nous plaçons maintenant dans le cas envisagé à la fin de 12.2. En vertu de (12.1.12) et de (12.3.6) appliqué à $C = X_o$, le groupe des caractères $M$ du tore maximal $T_o$ de la fibre spéciale $A_o$ du modèle de Néron $A$ de $A = \underline{\mathrm{Pic}}^{o}_{X_\eta/\eta}$ est donné par

}{\begin{equation}\tag{12.3.13}
v = (v_o,v_1) \ : \ C.(\Gamma(C)) = \left[\ \mathbb{Z}^{(\Gamma_1)} \longrightarrow \mathbb{Z}^{(J)}\ \right] \to \left[\ \mathbb{Z}^{(K)} \to \mathbb{Z}^{(I \times J)}\ \right]
\end{equation}
donné par

\begin{equation}\tag{12.3.14}
v_1(x,C',C'') = (x,C'') - (x,C'), \quad v_o(\widetilde{C}_j) = (0,\ \widetilde{C}_j) \quad .
\end{equation}
12.4. Nous nous plaçons maintenant dans le cas envisagé à la fin de 12.2. En vertu de (12.1.12) et de (12.3.6) appliqué à $C = X_o$, le groupe des caractères $M$ du tore maximal $T_o$ de la fibre spéciale $A_o$ du modèle de Néron $A$ de $A_\eta = \underline{\mathrm{Pic}}^{o}_{X_\eta/\eta}$ est donné par

}{F7I-IX-154}\begin{equation}\tag{12.4.1}
M = \mathrm{Ker}\ (\mathbb{Z}^{I} \xrightarrow{\ d\ } \mathbb{Z}^{J}) \quad ,
\end{equation}
où $I$ est l'ensemble des points singuliers de $X_o$, et $J$ l'ensemble de ses composantes irréductibles, l'homomorphisme $\mathbb{Z}^{I} \longrightarrow \mathbb{Z}^{J}$ étant donné par


\SGApage{I}{474}

\begin{equation}\tag{12.4.2}
d(x_i) = C'_i - C''_i \quad ,
\end{equation}
$C'_i$ et $C''_i$ étant les deux composantes irréductibles de $\widetilde{C}$ qui correspondent aux branches passant par le point double $x_i$. L'opérateur $d$ dépend du choix de l'ordre $(C'_i,C''_i)$ qu'on se donne entre ces deux branches, et l'isomorphisme (12.4.1) dépend évidemment de ce système de choix. Cependant, si pour $i \in I$ on désigne par

\SGAvar{\begin{equation}\tag{12.4.3}
\varphi_{x_i} \ : \ M \longrightarrow \mathbb{Z}
\end{equation}
le composé $M \longrightarrow \mathbb{Z}^{I} \xrightarrow{\ \mathrm{pr}_i\ } \mathbb{Z}$ , où le premier homomorphisme provient de (12.4.1), on constate aussitôt (en revenant à l'expression (12.3.4)) que $\varphi_{x_i}$ est bien déterminé au signe près ; de façon plus précise, il dépend seulement de l'ordre pris entre les deux branches $C'_i,C''_i$ passant par $i$, et change de signe quand on change cet ordre. Il s'ensuit notamment que la forme billinéaire

\begin{equation}\tag{12.4.4}
\varphi_{x_i} \otimes \varphi_{x_i} \in \check{M} \otimes_{\mathbb{Z}} \check{M} \simeq \mathrm{Hom}(M \otimes M, \mathbb{Z}\,)
\end{equation}
est déterminée canoniquement en termes du point singulier $x_i$ , sans ambiguité de signe. On peut donc poser

\begin{equation}\tag{12.4.5}
u = \sum_{i} \varphi_{x_i} \otimes \varphi_{x_i} \ : \ M \otimes_{\mathbb{Z}} M \longrightarrow \mathbb{Z} \quad ,
\end{equation}
et on trouve ainsi un accouplement bien déterminé i.e. indépendant de tout choix de signes. Il est clair qu'il est \underline{défini positif}, car c'est la restriction à $M \hookrightarrow \mathbb{Z}^{I}$ de la forme quadratique standart $\Sigma\, X_i^2$ sur $\mathbb{Z}^{I}$ .


}{\begin{equation}\tag{12.4.3}
\varphi_{x_i} \ : \ M \longrightarrow \mathbb{Z}
\end{equation}
le composé $M \longrightarrow \mathbb{Z}^{I} \xrightarrow{\ \mathrm{pr}_i\ } \mathbb{Z}$ , où le premier homomorphisme provient de (12.4.1), on constate aussitôt (en revenant à l'expression (12.3.4)) que $\varphi_{x_i}$ est bien déterminé au signe près ; de façon plus précise, il dépend seulement de l'ordre pris entre les deux branches $C'_i,C''_i$ passant par $i$, et change de signe quand on change cet ordre. Il s'ensuit notamment que la forme bilinéaire

\begin{equation}\tag{12.4.4}
\varphi_{x_i} \otimes \varphi_{x_i} \in \check{M} \otimes_{\mathbb{Z}} \check{M} \simeq \mathrm{Hom}(M \otimes M, \mathbb{Z}\,)
\end{equation}
est déterminée canoniquement en termes du point singulier $x_i$ , sans ambiguïté de signe. On peut donc poser

\begin{equation}\tag{12.4.5}
u = \sum_{i} \varphi_{x_i} \otimes \varphi_{x_i} \ : \ M \otimes_{\mathbb{Z}} M \longrightarrow \mathbb{Z} \quad ,
\end{equation}
et on trouve ainsi un accouplement bien déterminé i.e. indépendant de tout choix de signes. Il est clair qu'il est \underline{défini positif}, car c'est la restriction à $M \hookrightarrow \mathbb{Z}^{I}$ de la forme quadratique standard $\Sigma\, X_i^2$ sur $\mathbb{Z}^{I}$ .


}{F7I-IX-155}\SGApage{I}{475}

La notation $M$ pour le groupe des caractères de $T_o$ n'est pas en accord a priori avec celle des paragraphes précédents, où on écrivait $M'$, $M$ étant le groupe analogue relatif au schéma abélien $A'_K$ dual de $A_K$ . Mais comme $A_K$ est une jacobienne, il est muni d'une polarisation $\xi$ canonique [5\hspace{2.5em}], induisant un isomorphisme canonique

\begin{equation}\tag{12.4.6}
\psi_{\xi} : A_K \overset{\sim}{\longrightarrow} A'_K \quad ,
\end{equation}
moyennant lequel nous identifierons $A_K$ et $A'_K$, de sorte qu'il n'y a plus lieu de distinguer $M$ et $M'$. En particulier, pour tout nombre premier $\ell$ , l'accouplement de monodromie (9.1.2) peut être considéré comme un accouplement

\begin{equation}\tag{12.4.7}
u_{\ell} : M_{\ell} \otimes_{\mathbb{Z}_{\ell}} M_{\ell} \longrightarrow \mathbb{Z}_{\ell} \quad .
\end{equation}
Ceci posé, nous pouvons énoncer le résultat principal du présent paragraphe :

\SGAvar{\underline{Théorème} 12,5.("\underline{Formule de PICARD-LEFSCHETZ}"). \underline{Soient} $S$ \underline{un trait strictement local,} $X$ \underline{un schéma régulier,} \underline{projectif et plat sur} $S$, \underline{à fibres de dimension} $1$, \underline{dont la fibre géométrique générique est lisse et connexe}, \underline{et la fibre géométrique spéciale n'admet comme seuls points singuliers que des points quadratiques ordinaires.} \underline{Avec les notations introduites dans} 12.4, \underline{pour tout nombre premier} $\ell$, \underline{l'accouplement de monodromie} (12.4.7) \underline{est égal à l'accouplement déduit de} (12.4.5) \underline{par extension des scalaires de} $\mathbb{Z}$ \underline{à} $\mathbb{Z}_{\ell}$ .


}{\underline{Théorème} 12.5.("\underline{Formule de PICARD-LEFSCHETZ}"). \underline{Soient} $S$ \underline{un trait strictement local,} $X$ \underline{un schéma régulier,} \underline{projectif et plat sur} $S$, \underline{à fibres de dimension} $1$, \underline{dont la fibre géométrique générique est lisse et connexe}, \underline{et la fibre géométrique spéciale n'admet comme seuls points singuliers que des points quadratiques ordinaires.} \underline{Avec les notations introduites dans} 12.4, \underline{pour tout nombre premier} $\ell$, \underline{l'accouplement de monodromie} (12.4.7) \underline{est égal à l'accouplement déduit de} (12.4.5) \underline{par extension des scalaires de} $\mathbb{Z}$ \underline{à} $\mathbb{Z}_{\ell}$ .


}{F7I-IX-156}\SGApage{I}{476}

12.6.  On notera que cet énoncé prouve 10.4 dans le cas particulier envisagé, lorsque dans 10.4 b) on prend pour $\xi$ la polarisation canonique envisagée dans 10.4. Comme nous l'avons montré dans 10.5, ceci suffit à entraîner 10.4 dans le cas général.

12.7.  \underline{Démonstration de} 12.5. Cas $\ell \neq p$. Elle s'appuie de façon essentielle sur la \underline{formule de Picard-Lefschetz cohomologique} dans la situation envisagée dans 12.5, formule donnant l'action du groupe de monodromie $I$ sur

\SGAvar{\begin{equation}\tag{12.7.1}
H^1(X_{\overline{\eta}} \, , \, \mathbb{Z}_\ell) \;\simeq\; H^1(A_{\overline{\eta}} \, , \, \mathbb{Z}_\ell) \;\simeq\; \mathrm{Hom}(T_\ell(A_{\overline{\eta}}), \mathbb{Z}_\ell)
\end{equation}
dans le cas $\ell \neq p$. Cette formule, qui sera démontrée (dans un cas d'ailleurs nettement plus général) dans un exposé ultérieur (Exp XV), est en effet essentiellement identique à l'assertion 12.5 pour le $\ell$ envisagé, compte tenu de la définition 9.2 de $u_\ell$ pour $\ell \neq p$,. Il reste alors à traiter le cas $\ell = p$. Nous le ferons par une méthode de réduction au cas de la caractéristique nulle (et où de plus $X_o$ n'a qu'un seul point singulier), que nous allons maintenant exposer.

\underline{Remarques} 12.7.2. On notera d'ailleurs que, une fois la réduction annoncée faite, les arguments habituels nous permettraient de nous ramener au cas où on a comme corps de base le corps $\mathbb{C}$ des complexes, et à invoquer alors la formule de Picard-Lefschetz sour la forme classique, - à cela près qu'il ne semble pas y avoir de référence satisfaisante de cette formule. La démonstration de la formule de Picard-Lefschetz donnée dans notre Séminaire procédera de façon analogue, par réduction au cas transcendant, justiciable de méthodes proprement transcendantes. On


}{\begin{equation}\tag{12.7.1}
H^1(X_{\overline{\eta}} \, , \, \mathbb{Z}_\ell) \;\simeq\; H^1(A_{\overline{\eta}} \, , \, \mathbb{Z}_\ell) \;\simeq\; \mathrm{Hom}(T_\ell(A_{\overline{\eta}}), \mathbb{Z}_\ell)
\end{equation}
dans le cas $\ell \neq p$. Cette formule, qui sera démontrée (dans un cas d'ailleurs nettement plus général) dans un exposé ultérieur (Exp XV), est en effet essentiellement identique à l'assertion 12.5 pour le $\ell$ envisagé, compte tenu de la définition 9.2 de $u_\ell$ pour $\ell \neq p$. Il reste alors à traiter le cas $\ell = p$. Nous le ferons par une méthode de réduction au cas de la caractéristique nulle (et où de plus $X_o$ n'a qu'un seul point singulier), que nous allons maintenant exposer.

\underline{Remarques} 12.7.2. On notera d'ailleurs que, une fois la réduction annoncée faite, les arguments habituels nous permettraient de nous ramener au cas où on a comme corps de base le corps $\mathbb{C}$ des complexes, et à invoquer alors la formule de Picard-Lefschetz sous la forme classique, - à cela près qu'il ne semble pas y avoir de référence satisfaisante de cette formule. La démonstration de la formule de Picard-Lefschetz donnée dans notre Séminaire procédera de façon analogue, par réduction au cas transcendant, justiciable de méthodes proprement transcendantes. On


}{F7I-IX-157}\SGApage{I}{477}

\SGAvar{peut donc dire qu'en tout état de cause, la démonstration de 12.5 (donc de (10.4) que nous proposons ici est de nature transcendante. Signalons d'ailleurs que la seule autre démonstration connue de 10.4 provient de la théorie rigide-analytique [24], et à ce titre doit être regardée également comme étant essentiellement ``transcendante''. Il serait intéressant de trouver des démonstrations purement algébriques de ces résultats (mais il n'est nullement évident que cela soit possible !).

}{peut donc dire qu'en tout état de cause, la démonstration de 12.5 (donc de (10.4)) que nous proposons ici est de nature transcendante. Signalons d'ailleurs que la seule autre démonstration connue de 10.4 provient de la théorie rigide-analytique [24], et à ce titre doit être regardée également comme étant essentiellement ``transcendante''. Il serait intéressant de trouver des démonstrations purement algébriques de ces résultats (mais il n'est nullement évident que cela soit possible !).

}{F7I-IX-158}12.8. \underline{Démonstration de} 12.5 : \underline{réduction au cas où} $S$ \underline{est de caractéristique nulle.} Nous allons utiliser un raisonnement s'appuyant sur le même principe que la démonstration du ``théorème d'action essentiellement modérée de la monodromie sur le $\pi_1$ '' dans V \hspace{4em}. Notons qu'on peut supposer $S$ complet grâce à 10.3. Soit $W$ un $p$-anneau de Cohen de corps résiduel $k$, et choisissons un homomorphisme $W \longrightarrow V$ induisant l'identité sur les corps résiduels (EGA $0_{\mathrm{IV}}$ 19.8.6). Soit $L = (x_\lambda)_{\lambda \in L}$ l'ensemble des points singuliers de $X_o$, et considérons le schéma formel des modules $\underline{S}$ de $X_o$ sur $W$, et la ``courbe universelle'' $\underline{f} : \underline{X} \longrightarrow \underline{S}$ de fibre spéciale $X_o$ (VI \hspace{3em}). On a vu dans loc. cit. que $\underline{S}$ est isomorphe au spectre d'un anneau de séries formelles $W[[T_1,\ldots,T_N]]$, et qu'on peut trouver dans $\underline{S}$ un diviseur à croisements normaux

\begin{equation}\tag{12.8.1}
D = \sum_{\lambda \in L} D_\lambda \quad ,
\end{equation}
dont les composantes irréductibles $D_\lambda$ sont indexées par l'ensemble $L$ précédent, chaque $D_\lambda$ étant décrit par une équation de la forme $T_{i_\lambda} = 0$ (où on peut supposer, si on veut, $i_\lambda = \lambda$, si on indexe $L$ par un ensemble


\SGApage{I}{478}

d'entiers $[1,r]$), $D_\lambda$ étant l'image par $\underline{f} : \underline{X} \longrightarrow \underline{S}$ de la composante $D'_\lambda$ passant par le point singulier $x_\lambda$ de $X_o$ du lieu de non lissité $D' = \bigcup_{\lambda \in L} D'_\lambda$ du morphisme $\underline{f}$. Par construction même, $f : X \longrightarrow S$ est isomorphe à l'image inverse de $\underline{f} : \underline{X} \longrightarrow \underline{S}$ par un $W$-morphisme convenable $g : S \longrightarrow \underline{S}$ :

\begin{equation}\tag{12.8.2}
\begin{tikzcd}
\underline{X} \arrow[d, "\underline{f}"'] \& X \arrow[l] \arrow[d, "f"] \\
\underline{S} \& S \arrow[l, "g"'] \& .
\end{tikzcd}
\end{equation}















Posons

\SGAvar{\begin{equation}\tag{12.8.3}
P = \underline{\mathrm{Pic}}^{o}_{\underline{X}/\underline{S}} \quad ;
\end{equation}
comme la fibre spéciale de $\underline{f}$ est géométriquement réduite de dimension 1, il s'ensuit que $\underline{F}$ est cohomologiquement plat en dimension 0 (EGA III 7.8.1) et que $P$ est formellement lisse sur $S$ (ce dernier fait résultant de ce que $H^2(X_o,\underline{O}_{X_o}) = 0$, $X_o$ étant de dimension 1). Pour simplifier, nous allons admettre que $P$ \underline{est représentable par un schéma en groupes sur} $S$ (qui sera nécessairement un schéma en groupes lisse à fibres connexes) ; c'est là un cas particulier d'un résultat inédit de MUMFORD cité dans [22](considérablement généralisé par RAYNAUD dans la note citée, dans le cas où la base est un trait). Comme aucune démonstration des résultats de MUMFORD ni de RAYNAUD n'a été publiée, le lecteur préfèrera peut-être utiliser les résultats généraux de M. ARTIN [1], qui donnent la représentabilité de $P$ (sous la seule hypothèse de partir d'un $\underline{f}$ propre, plat et cohomologiquement plat en dimension 0) par un ``espace algébrique''


}{\begin{equation}\tag{12.8.3}
P = \underline{\mathrm{Pic}}^{o}_{\underline{X}/\underline{S}} \quad ;
\end{equation}
comme la fibre spéciale de $\underline{f}$ est géométriquement réduite de dimension 1, il s'ensuit que $\underline{f}$ est cohomologiquement plat en dimension 0 (EGA III 7.8.1) et que $P$ est formellement lisse sur $\underline{S}$ (ce dernier fait résultant de ce que $H^2(X_o,\underline{O}_{X_o}) = 0$, $X_o$ étant de dimension 1). Pour simplifier, nous allons admettre que $P$ \underline{est représentable par un schéma en groupes sur} $\underline{S}$ (qui sera nécessairement un schéma en groupes lisse à fibres connexes) ; c'est là un cas particulier d'un résultat inédit de MUMFORD cité dans [22](considérablement généralisé par RAYNAUD dans la note citée, dans le cas où la base est un trait). Comme aucune démonstration des résultats de MUMFORD ni de RAYNAUD n'a été publiée, le lecteur préfèrera peut-être utiliser les résultats généraux de M. ARTIN [1], qui donnent la représentabilité de $P$ (sous la seule hypothèse de partir d'un $\underline{f}$ propre, plat et cohomologiquement plat en dimension 0) par un ``espace algébrique''


}{F7I-IX-159}\SGApage{I}{479}

au sens de ARTIN, ce qui suffirait pour notre propos. Posant

\begin{equation}\tag{12.8.4}
U = \underline{S} - \mathrm{supp}\, D \quad ,
\end{equation}
de sorte que $U$ est le plus grand ouvert de $\underline{S}$ au-dessus duquel $\underline{X}$ soit lisse, on voit donc que

\begin{equation*}
P\vert U = \underline{\mathrm{Pic}}^{o}_{\underline{X}_U/U}
\end{equation*}
est un schéma abélien sur $U$, dont l'image inverse sur $\eta$ par $g : \eta \longrightarrow U$ est le schéma abélien $A_\eta$ sur $\eta$ qu'on se propose d'étudier. Plus précisément, on a un isomorphisme canonique sur $S$ (12.1.10) :

\begin{equation}\tag{12.8.5}
P \times_{\underline{S}} S \simeq A^{o} \quad .
\end{equation}
Comme la fibre spéciale $P_o = A_o^{o}$ est $\ell$-divisible pour tout $\ell$, il en est de même de $P$ lui-même, dont toutes les fibres sont donc des extensions d'une variété abélienne par un tore. Considérons le pro-$\ell$-groupe $T_\ell(P)$ sur $\underline{S}$, dont l'image inverse par $g$ n'est autre que $T_\ell(A^{o})$. Appliquant les constructions 7.3,on définit une filtration de ce pro-groupe par une``partie fixe'' et une ``partie torique''

\begin{equation}\tag{12.8.6}
T_\ell(P) \supset T_\ell(P)^f = T_\ell(\hat{P}) \supset T_\ell(P)^t = T_\ell(T) \supset 0 \quad ,
\end{equation}
dont l'image inverse par $g$ est la filtration analogue de $T_\ell(A^{o})$, étudiée dans les par. 2. et 5.

Considérons la restriction de cette filtration au-dessus de l'ouvert $U$ de $\underline{S}$, sur lequel on dispose de l'autodualité canonique de $P_U$ comme jacobienne d'un schéma en courbes lisse sur $U$, d'où un accouplement correspondant


\SGApage{I}{480}

\SGAvar{\begin{equation}\tag{12.8.7}
\varphi : T_\ell(P_U) \otimes T_\ell(P_U) \longrightarrow T_\ell(\mathbb{G}_{mU}) \quad ,
\end{equation}
qui est une \underline{autodualité} (au sens des groupes de Barsotti-Tate) de $T_\ell(B_U)$. On a vu dans 7.4.8 \quad que pour cet accouplement, $T_\ell(P_U)^f$ \underline{et} $T_\ell(P_U)^t$ \underline{sont les orthogonaux l'un de l'autre}. De ceci résulte que

\begin{equation*}
E_U = T_\ell(P_U)/T_\ell(P_U)^t
\end{equation*}
est canoniquement isomorphe au pro-groupe de Barsotti-Tate dual de $T_\ell(P_U)^f = F_U$, donc se prolonge canoniquement en un groupe de Barsotti-Tate $E$ sur $S$, savoir

}{\begin{equation}\tag{12.8.7}
\varphi : T_\ell(P_U) \otimes T_\ell(P_U) \longrightarrow T_\ell(\mathbb{G}_{mU}) \quad ,
\end{equation}
qui est une \underline{autodualité} (au sens des groupes de Barsotti-Tate) de $T_\ell(P_U)$. On a vu dans 7.4.8 \quad que pour cet accouplement, $T_\ell(P_U)^f$ \underline{et} $T_\ell(P_U)^t$ \underline{sont les orthogonaux l'un de l'autre}. De ceci résulte que

\begin{equation*}
E_U = T_\ell(P_U)/T_\ell(P_U)^t
\end{equation*}
est canoniquement isomorphe au pro-groupe de Barsotti-Tate dual de $T_\ell(P_U)^f = F_U$, donc se prolonge canoniquement en un groupe de Barsotti-Tate $E$ sur $\underline{S}$, savoir

}{F7I-IX-160}\begin{equation}\tag{12.8.8}
E = D(F) \ , \ \text{où} \ F = T_\ell(P)^f \quad .
\end{equation}
Nous sommes alors dans une situation typique d'extensions panachées (9.3), et nous pouvons reprendre les constructions de 9.4 à 9.6 (que nous nous excusons d'avoir rédigé dans le contexte éculé et pisseux d'un schéma de base réduit à un trait !) : $T_\ell(P_U)$ est une extension panachée de $E_U$ par $F_U$, où $E$ et $F$ sont des groupes de Barsotti-Tate sur $\underline{S}$ qui sont des extensions, duales l'une de l'autre :

\SGAvar{\begin{equation}\tag{12.8.9}
\begin{tikzcd}[row sep=small, column sep=small]
0 \arrow[r] \& P \arrow[rr] \arrow[d, dash, "\wr"] \& \& F \arrow[r] \arrow[d, dash, "\wr"] \& R \arrow[r] \& 0 \& \\
\& T_\ell(P)^t \simeq \underline{\mathrm{Hom}}(M,T_\ell(\underline{\mathbb{G}}_{m\underline{S}}) \& \& T_\ell(P)^f \& \& \& \\
\& \& \& \& \& \& \\
0 \arrow[r] \& R \arrow[rr] \arrow[d, dash, "\wr"] \& \& E \arrow[r] \& Q \arrow[r] \arrow[d, dash, "\wr"] \& 0 \& \\
\& D(R) \& \& \& M \otimes \mathbb{Z}_{\ell\,\underline{S}} \& \& .
\end{tikzcd}
\end{equation}

























}{\begin{equation}\tag{12.8.9}
\begin{tikzcd}[row sep=small, column sep=small]
0 \arrow[r] \& P \arrow[rr] \arrow[d, dash, "\wr"] \& \& F \arrow[r] \arrow[d, dash, "\wr"] \& R \arrow[r] \& 0 \& \\
\& T_\ell(P)^t \simeq \underline{\mathrm{Hom}}(M,T_\ell(\mathbb{G}_{m\underline{S}})) \& \& T_\ell(P)^f \& \& \& \\
\& \& \& \& \& \& \\
0 \arrow[r] \& R \arrow[rr] \arrow[d, dash, "\wr"] \& \& E \arrow[r] \& Q \arrow[r] \arrow[d, dash, "\wr"] \& 0 \& \\
\& D(R) \& \& \& M \otimes \mathbb{Z}_{\ell\,\underline{S}} \& \& .
\end{tikzcd}
\end{equation}

























}{F7I-IX-161}\SGApage{I}{481}

(NB $R$ est canoniquement autodual, comme il résulte de l'autodualité (12.8.7) et de l'énoncé d'orthogonalité signalé plus haut.) La seule différence avec loc. cit. est que l'accouplement qu'on trouve de $M_\ell \times M_\ell$ est à valeurs dans un groupe

\SGAvar{\begin{equation*}
(\mathbb{G}_m(U)/\mathbb{G}_m(\underline{S}) \otimes_{\mathbb{Z}} \mathbb{Z}_\ell
\end{equation*}
qui n'est plus canoniquement isomorphe à $\mathbb{Z}_\ell$, mais à $\mathbb{Z}_\ell^{\,L}$, grâce à l'isomorphisme

}{\begin{equation*}
(\mathbb{G}_m(U)/\mathbb{G}_m(\underline{S})) \otimes_{\mathbb{Z}} \mathbb{Z}_\ell
\end{equation*}
qui n'est plus canoniquement isomorphe à $\mathbb{Z}_\ell$, mais à $\mathbb{Z}_\ell^{\,L}$, grâce à l'isomorphisme

}{F7I-IX-162}\begin{equation}\tag{12.8.10}
\mathbb{Z}^L \simeq \mathbb{G}_m(U)/\mathbb{G}_m(\underline{S})
\end{equation}
qui applique l'élément d'indice $\lambda$ de la base canonique de $\mathbb{Z}^L$ en la classe de $T_{i_\lambda}$ dans le deuxième membre, où $T_{i_\lambda}$ est une équation du diviseur $D_\lambda$ . On trouve ainsi un accouplement canonique

\SGAvar{\begin{equation}\tag{12.8.11}
\underline{u}_\ell : M_\ell \otimes M_\ell \longrightarrow \mathbb{Z}_\ell^{\,L} \qquad .
\end{equation}
On vérifie de façon essentiellement triviale que l'accouplement de monodromie $u_\ell$ (12.4.7) sur $S$ que nous proposons d'étudier est le composé de $u_\ell$ avec l'homomorphisme $\mathbb{Z}_\ell^{\,L} \longrightarrow \mathbb{Z}_\ell$ provenant, par tensorisation avec $Z_\ell$ sur $\mathbb{Z}$ , de l'homomorphisme

}{\begin{equation}\tag{12.8.11}
\underline{u}_\ell : M_\ell \otimes M_\ell \longrightarrow \mathbb{Z}_\ell^{\,L} \qquad .
\end{equation}
On vérifie de façon essentiellement triviale que l'accouplement de monodromie $u_\ell$ (12.4.7) sur $S$ que nous proposons d'étudier est le composé de $\underline{u}_\ell$ avec l'homomorphisme $\mathbb{Z}_\ell^{\,L} \longrightarrow \mathbb{Z}_\ell$ provenant, par tensorisation avec $Z_\ell$ sur $\mathbb{Z}$ , de l'homomorphisme

}{F7I-IX-163}\begin{center}
\begin{tikzcd}
\mathbb{G}_m(U)/\mathbb{G}_m(\underline{S}) \arrow[r] \arrow[d, dash, "\wr"'] \& \mathbb{G}_m(\eta)/\mathbb{G}_m(S) \arrow[d, dash, "\wr"'] \\
\mathbb{Z}^L \& \mathbb{Z}
\end{tikzcd}
\end{center}











induit par $g : S \longrightarrow \underline{S}$ (envoyant $\eta$ dans $U$). Cet homomorphisme est donné par un élément


\SGApage{I}{482}

\SGAvar{\begin{equation}\tag{12.8.12}
(d_\lambda)_{\lambda \in L} \in \mathbb{Z}^L \ , \ (d_\lambda \geqslant 1 \ \mbox{pour tout} \ \lambda \in L)
\end{equation}
où pour tout $\lambda \in L$, $d_\lambda$ est la multiplicité de l'image inverse dans $S$ du diviseur $D_\lambda$ sur $\underline{S}$, qui ne dépend d'ailleurs que du complété $\widehat{\underline{O}}_{X,x_\lambda}$ en tant que $V$-algèbre (VI\hspace{2.5em}). On sait de plus (VI\hspace{2.5em}) que pour un $\lambda$ donné, on a $d_\lambda = 1$ si et seulement si $x$ est \underline{régulier} au point $x_\lambda$ . Comme dans 12.5 on a supposé $X$ régulier, i.e. les $d_\lambda$ égaux à 1, on voit donc que 12.5 sera une conséquence du lemme suivant, où $X,S,g$ ne figurent plus :

}{\begin{equation}\tag{12.8.12}
(d_\lambda)_{\lambda \in L} \in \mathbb{Z}^L \ , \ (d_\lambda \geqslant 1 \ \mbox{pour tout} \ \lambda \in L)
\end{equation}
où pour tout $\lambda \in L$, $d_\lambda$ est la multiplicité de l'image inverse dans $S$ du diviseur $D_\lambda$ sur $\underline{S}$, qui ne dépend d'ailleurs que du complété $\widehat{\underline{O}}_{X,x_\lambda}$ en tant que $V$-algèbre (VI\hspace{2.5em}). On sait de plus (VI\hspace{2.5em}) que pour un $\lambda$ donné, on a $d_\lambda = 1$ si et seulement si $X$ est \underline{régulier} au point $x_\lambda$ . Comme dans 12.5 on a supposé $X$ régulier, i.e. les $d_\lambda$ égaux à 1, on voit donc que 12.5 sera une conséquence du lemme suivant, où $X,S,g$ ne figurent plus :

}{F7I-IX-164}\underline{Lemme} 12.8.13. \underline{Si} $\lambda \in L$, \underline{le composé de} (12.8.11) \underline{avec} $\mathrm{pr}_\lambda$ \underline{est égal à la forme} $\varphi_{x_\lambda} \otimes \varphi_{x_\lambda}$ (12.4.4).

\SGAvar{Pour prouver ce dernier énoncé, considérons le localisé strict complété $\underline{S}^\lambda$ de $\underline{S}$ en le point maximal de $D$, et $\underline{X}^\lambda = X \times_{S} \underline{S}^\lambda$. Alors la fibre spéciale $X_o^\lambda$ n'a qu'un seul point singulier $y_\lambda$, lequel est encore quadratique ordinaire. Dans l'accouplement de monodromie correspondant

}{Pour prouver ce dernier énoncé, considérons le localisé strict complété $\underline{S}^\lambda$ de $\underline{S}$ en le point maximal de $D_\lambda$, et $\underline{X}^\lambda = \underline{X} \times_{\underline{S}} \underline{S}^\lambda$. Alors la fibre spéciale $\underline{X}_o^\lambda$ n'a qu'un seul point singulier $y_\lambda$, lequel est encore quadratique ordinaire. Dans l'accouplement de monodromie correspondant

}{F7I-IX-165}\begin{equation}\tag{12.8.14}
u_\ell^\lambda : M_\ell^\lambda \otimes M_\ell^\lambda \longrightarrow \mathbb{Z}_\ell \qquad ,
\end{equation}
$M^\lambda$ est donc de rang 1 ou zéro. Nous laissons au lecteur le soin de vérifier les faits faciles suivants, qui explicitent, dans le cas particulier où nous en avons besoin, le comportement des constructions précédentes par localisation :

a) Il existe un unique sous-tore $T_\lambda$ de $P \times_{\underline{S}} D_\lambda$ , dont la restriction au point maximal $s_\lambda$ de $D_\lambda$ est le tore maximal de $P_{s_\lambda}$ (lequel est de dimension 0 ou 1).


\SGApage{I}{483}

b)  Considérons l'homomorphisme transposé $M \longrightarrow M^{\lambda}$, sur les groupes des caractères, de l'inclusion $T_{\lambda_{o}} \hookrightarrow T_{o} =$ tore maximal de $P_{o}$, et la forme $\varphi_{y_{\lambda}} : M^{\lambda} \longrightarrow \mathbb{Z}$, définie au signe près, définie comme dans 12.4.3 (avec $X$ remplacé par $\underline{X}^{\lambda}$). On a alors commutativité, au signe près, dans

\SGAvar{\begin{equation}\tag{12.8.15}
\begin{tikzcd}[column sep=huge, row sep=large]
M \arrow[rr] \arrow[dr, "\varphi_{x_{\lambda}}"'] \& \& M^{\lambda} \arrow[dl, "\varphi_{\lambda_{\lambda}}"] \\
 \& \mathbb{Z} \&
\end{tikzcd} \qquad .
\end{equation}











}{\begin{equation}\tag{12.8.15}
\begin{tikzcd}[column sep=huge, row sep=large]
M \arrow[rr] \arrow[dr, "\varphi_{x_{\lambda}}"'] \& \& M^{\lambda} \arrow[dl, "\varphi_{y_{\lambda}}"] \\
 \& \mathbb{Z} \&
\end{tikzcd} \qquad .
\end{equation}











}{F7I-IX-166}(En fait, on constate que le choix d'un ordre sur l'ensemble des deux branches de $X_{o}$ passant par $x_{\lambda}$ permet de faire un choix analogue en le point $y_{\lambda}$ de $\underline{X}^{\lambda}_{o}$ , et pour deux choix qui se correspondent ainsi, le diagramme précédent est commutatif.)

c)  On a commutativité dans le diagramme suivant

\begin{center}
\begin{tikzcd}[column sep=huge, row sep=large]
M_{\ell} \otimes M_{\ell} \arrow[rr] \arrow[dr, "\mathrm{pr}_{\lambda}\underline{u}_{\ell}"'] \& \& M^{\lambda}_{\ell} \otimes M^{\lambda}_{\ell} \arrow[dl, "u^{\lambda}_{\ell}"] \\
 \& \mathbb{Z}_{\ell} \&
\end{tikzcd} \qquad ,
\end{center}














où les notations sont celles de 12.8.13 et (12.8.14).

De ces relations résulte aussitôt que pour prouver 12.8.13, il suffit de prouver la relation analogue pour $\underline{X}^{\lambda}$ :

\begin{equation*}
u^{\lambda}_{\ell} = (\varphi_{y_{\lambda}} \otimes \varphi_{y_{\lambda}}) \otimes_{\mathbb{Z}} \mathbb{Z}_{\ell} \quad ,
\end{equation*}

\SGApage{I}{484}

\SGAvar{qui n'est autre que 12.5 dans le cas de $\underline{X}^{\lambda}$ sur $\underline{S}^{\lambda}$. Or $\underline{S}$ est un trait de \underline{caractéristique nulle}. Cela achève donc la réduction annoncée.

\underline{Remarques} 12.9. On notera que cette réduction nous réduit en même temps au cas où $X_{o}$ n'a qu'un seul point singulier, qui est celui qui était considéré classiquement dans la formule de Picard-Lefschetz. Notons d'ailleurs que dans ce cas, $X_{o}$ a une ou deux composantes irréductibles, ces cas correspondants respectivement à rang $M = 1$ et rang $M = 0$ (ce dernier cas étant celui où la jacobienne de la fibre générique a bonne réduction).

12.10.  La démonstration donnée dans 12.8 nous fournit un résultat plus général que 12.5, valable sans supposer que $X$ soit régulier, en introduisant les entiers $d_{\lambda} \geqslant 1$ de (12.8.12) : on trouve alors que les accouplements de monodormie se déduisent tous de l'accouplement

}{qui n'est autre que 12.5 dans le cas de $\underline{X}^{\lambda}$ sur $\underline{S}^{\lambda}$. Or $\underline{S}^\lambda$ est un trait de \underline{caractéristique nulle}. Cela achève donc la réduction annoncée.

\underline{Remarques} 12.9. On notera que cette réduction nous réduit en même temps au cas où $X_{o}$ n'a qu'un seul point singulier, qui est celui qui était considéré classiquement dans la formule de Picard-Lefschetz. Notons d'ailleurs que dans ce cas, $X_{o}$ a une ou deux composantes irréductibles, ces cas correspondant respectivement à rang $M = 1$ et rang $M = 0$ (ce dernier cas étant celui où la jacobienne de la fibre générique a bonne réduction).

12.10.  La démonstration donnée dans 12.8 nous fournit un résultat plus général que 12.5, valable sans supposer que $X$ soit régulier, en introduisant les entiers $d_{\lambda} \geqslant 1$ de (12.8.12) : on trouve alors que les accouplements de monodromie se déduisent tous de l'accouplement

}{F7I-IX-167}\begin{equation}\tag{12.10.1}
u = \sum_{\lambda \in L} d_{\lambda} \, \varphi_{x_{\lambda}} \otimes \varphi_{x_{\lambda}} \; : \; M \otimes M \longrightarrow \mathbb{Z} \quad .
\end{equation}

\SGAsectiontoc{13.  Liens avec la théorie transcendante : cas analytique complexe.}{13}
13.  \underline{Liens avec la théorie transcendante} : \underline{cas analytique complexe}.

Pour les notions générales relatives aux espaces analytiques, nous renvoyons à [2].

13.1.  \underline{Rappels sur les groupes analytiques sur} $S$.  Soit $S$ un espace analytique, $A$ un espace analytique au-dessus de $S$, muni d'une structure de groupe commutatif au-dessus de $S$ (i.e. $A$ est un groupe commutatif de la catégorie des espaces analytiques au-dessus de $S$). On suppose


\SGApage{I}{485}

$A$ \underline{lisse} sur $S$, et on désigne par

\begin{equation}\tag{13.1.1}
\underline{E} = \underline{\mathrm{Lie}}(A/S)
\end{equation}
le module localement libre sur $S$, image inverse par la section nulle du faisceau tangent relatif de $A$ sur $S$. Nous admettrons ici la définition de l'homomorphisme exponentiel

\begin{equation}\tag{13.1.2}
\exp: E \longrightarrow A \qquad ,
\end{equation}
où $E$ est le fibré vectoriel analytique sur $S$ dont le Module des sections holomorphes est $\underline{E}$ . Cet homomorphisme se définit, au voisinage d'un point $s \in S$, d'abord au voisinage de la section nulle de $E$ par les méthodes bien connues en théorie des groupes formels en car. nulle, puis se prolonge de façon évidente par multiplicativité à $E$ tout entier. C'est un homomorphisme de groupes analytiques relatifs sur $S$, qui est \underline{étale} et a comme image le sous-espace ouvert $A^{o}$ de $A$ réunion des composantes neutres des fibres de $A$ sur $S$. Si $R$ est le sous-espace analytique noyau de $\exp$, on a donc une suite exacte de groupes analytiques relatifs :

\begin{equation}\tag{13.1.3}
0 \longrightarrow R \longrightarrow E \longrightarrow A^{o} \longrightarrow 0 \qquad ,
\end{equation}
où le morphisme $E \longrightarrow A^{o}$ est surjectif et étale, donc $R$ est un groupe analytique relatif \underline{étale} sur $S$. C'est un sous-espace analytique de $E$, qui est fermé dans $E$ si et seulement si $A^{o}$ est séparé sur $S$, i.e. la section nulle de $A^{o}$ est fermée dans $A^{o}$.


\SGApage{I}{486}

13.1.4. Bien entendu, la suite exacte (13.1.3) est \underline{fonctorielle en} $A$, et commute à tout changement de base $S' \longrightarrow S$. De cette façon on trouve une équivalence de catégories, compatible aux changements de base, entre la catégorie des groupes analytiques relatifs sur $S$ lisses et à fibres connexes, et la catégorie des couples $(E,R)$ d'un fibré vectoriel localement libre $E$ sur $S$ et d'un sous-groupe analytique $R$ de $E$ étale sur $S$.

13.1.5. Lorsque $A$ est un groupe ``\underline{abéloide}'' sur $S$, par quoi nous entendons qu'il est lisse, propre et séparé sur $S$ et à fibres connexes, de sorte qu'on a en particulier $A=A^{o}$, alors $R$ est un groupe localement constant sur $S$, dont le rang en chaque point $s$ (en tant que $\mathbb{Z}$ -module) est égal à $2\ \mathrm{rang}_s(E)$, et réciproquement. De cette façon on trouve l'interprétation bien connue d'un espace abéloide sur $S$ en termes d'un ``\underline{système local}'' $R$ sur $S$ et d'une structure de fibré vectoriel analytique \underline{complexe} sur le fibré analytique réel $R \otimes_{\mathbb{Z}} \mathbb{R}$ défini par $R$.

\SGAvar{13.2. \underline{Complétion formelle le long d'une fibre}. Nous supposons donné désormais un point $s \in S$, et allons travailler localement au voisinage de $s$. Pour tout entier $n \geq 0$, désignons par $S_n$ le $n$.ième voisinage infinitésimal de $s$ dans $S$, et pour tout espace analytique $X$ sur $S$, soit de même $X_n=X\times_S S_n$ ; pour $n$ variable, le système des $X_n$ forme alors un système inductif d'espaces analytiques sur $S$, qu'on dénotera par $\hat{X}$ et appelera le \underline{complété formel de} $X$ \underline{le long de} $X_s$ . Il dépend évidemment fonctoriellement de $X$, et sa formation commute aux limites projectives


}{13.2. \underline{Complétion formelle le long d'une fibre}. Nous supposons donné désormais un point $s \in S$, et allons travailler localement au voisinage de $s$. Pour tout entier $n \geq 0$, désignons par $S_n$ le $n$.ième voisinage infinitésimal de $s$ dans $S$, et pour tout espace analytique $X$ sur $S$, soit de même $X_n=X\times_S S_n$ ; pour $n$ variable, le système des $X_n$ forme alors un système inductif d'espaces analytiques sur $S$, qu'on dénotera par $\hat{X}$ et appellera le \underline{complété formel de} $X$ \underline{le long de} $X_s$ . Il dépend évidemment fonctoriellement de $X$, et sa formation commute aux limites projectives


}{F7I-IX-168}\SGApage{I}{487}

\SGAvar{finies ; donc $X \longmapsto \hat{X}$ transforme groupes en groupes. Notons que $\hat{X}$ ne dépend en fait que de la restriction de $X$ à un voisinage ouvert arbitraire $U$ de $S$.

}{finies ; donc $X \longmapsto \hat{X}$ transforme groupes en groupes. Notons que $\hat{X}$ ne dépend en fait que de la restriction de $X$ à un voisinage ouvert arbitraire $U$ de $s$.

}{F7I-IX-169}\underline{Lemme} 13.2.1. \underline{Les données sont celles de} 13.1, \underline{et on suppose que les conditions suivantes sont vérifiées} : a) $A = A^{o}$,  b)  $R$ \underline{est constant au voisinage de} $s$, \underline{et} c) $R_s$ \underline{engendre l'espace vectoriel complexe} $E_s$ . \underline{Soit} $A'$ \underline{un deuxième groupe analytique lisse sur} $S$, \underline{et considérons l'application}

\begin{equation*}
u \longmapsto \hat{u} : \mathrm{Hom}_{S\text{-}\mathrm{gr}}(A,A') \longrightarrow \mathrm{Hom}_{\hat{S}\text{-}\mathrm{gr}}(\hat{A},\hat{A}') \quad ,
\end{equation*}
\underline{où les notations sont celles de} 13.2, \underline{et l'application analogue obtenue en se restreignant au-dessus d'un voisinage ouvert} $U$ \underline{variable de} $s$, \underline{et passant à la limite inductive} :

\begin{equation}\tag{13.2.1.1}
u \longmapsto \hat{u} : \mathrm{Hom}_{(S,s)\text{-}\mathrm{gr}}(A,A') \longrightarrow \mathrm{Hom}_{\hat{S}\text{-}\mathrm{gr}}(\hat{A},\hat{A}') \quad ,
\end{equation}
(\underline{où} $(S,s)$ \underline{désigne le} germe \underline{d'espace analytique de} $S$ \underline{en} $s$). \underline{Alors l'application précédente est} bijective.

Considérons en effet la suite exacte

\SGAvar{\begin{equation*}
0 \longrightarrow R' \longrightarrow E' \longrightarrow {A'}^{o} \longrightarrow C
\end{equation*}
analogue à (13.1.3) définie par $A'$. Compte tenu du fait que les homomorphismes $A \longrightarrow A'$ se factorisent nécessairement par ${A'}^{o}$, et de même après les changements de base $U \longrightarrow S$ et $S_n \longrightarrow S$, on peut supposer $A' = {A'}^{o}$. Compte tenu du dictionnaire 13.1.4, il faut prouver que la donnée d'un


}{\begin{equation*}
0 \longrightarrow R' \longrightarrow E' \longrightarrow {A'}^{o} \longrightarrow 0
\end{equation*}
analogue à (13.1.3) définie par $A'$. Compte tenu du fait que les homomorphismes $A \longrightarrow A'$ se factorisent nécessairement par ${A'}^{o}$, et de même après les changements de base $U \longrightarrow S$ et $S_n \longrightarrow S$, on peut supposer $A' = {A'}^{o}$. Compte tenu du dictionnaire 13.1.4, il faut prouver que la donnée d'un


}{F7I-IX-170}\SGApage{I}{488}

germe d'homomorphisme de fibrés vectoriels $E \longrightarrow E'$, appliquant $R$ dans $R'$, revient à la donnée d'un système cohérent d'homomorphismes $E_n \longrightarrow E'_n$ appliquant $R_n$ dans $R'_n$ . Or si $\underline{O}_s$ désigne l'anneau local de $S$ en $s$, la donnée d'un système cohérent d'homomorphismes $E_n \longrightarrow E'_n$ revient à la donnée d'un homomorphisme de $\hat{\underline{O}}_s$-modules

\SGAvar{\begin{equation*}
\hat{F} : \hat{\underline{E}}_s \longrightarrow \hat{\underline{E}}'_s
\end{equation*}
(où $\underline{E}_s$ et $E'_s$ sont les $\underline{O}_s$-modules libres fibres en $s$ des $\underline{O}_s$-Modules $\underline{E}$, $\underline{E}'$ correspondant à $E,E'$), tandis qu'un germe d'homomorphisme de fibrés vectoriels $E \longrightarrow E'$ revient à la donnée d'un homomorphisme de $\underline{O}_s$-modules

\begin{equation*}
F : \underline{E}_s \longrightarrow \underline{E}'_s \quad .
\end{equation*}
Cela montre déjà que l'application (13.2.1.1) est injective, et pour la surjectivité, il suffit de prouver que si on a un homomorphisme $\hat{F}$ comme dessus, tel que pour tout $n$ l'homomorphisme induit $F_n : E_n \longrightarrow E'_n$ applique $R_n$ dans $R'_n$, alors $\hat{F}$ provient d'un $F$ i.e. applique $\underline{E}_s$ dans $\underline{E}'_s$ , et le germe correspondant d'homomorphismes $E \longrightarrow E'$ applique $R$ dans $R'$. Or, comme par hypothèse c), $R_s = R_o$ engendre $E_s$, il s'ensuit par Nakayama que le groupe $R_o$ (identifié au groupe des sections de $\hat{R}$ sur $\hat{S}$, ou,au choix,au groupe des germes de sections de $R$ sur $S$) engendre $\hat{E}_s$ ; donc, pour prouver que $\hat{F}(\underline{E}_s) \subset \underline{E}'_s$ ,il suffit de prouver que $\hat{F}(R_o) \subset \underline{E}'_s$ . Or on a $\hat{F}(R_o) \subset R'_o$, où $R'_o$ est lui-même interprété comme le groupe des germes de sections de $R' \subset E'$ au voisinage de $s$, d'où l'inclusion annoncée, et par suite l'existence de $F$ . De plus on voit que $F(R_o) \subset R'_o$ , où $R_o$ et $R'_o$ sont considérés comme des groupes de germes de sections de $E$ resp. $E'$ . En vertu de b), $R_o$


}{\begin{equation*}
\hat{F} : \hat{\underline{E}}_s \longrightarrow \hat{\underline{E}}'_s
\end{equation*}
(où $\underline{E}_s$ et $\underline{E}'_s$ sont les $\underline{O}_s$-modules libres fibres en $s$ des $\underline{O}_S$-Modules $\underline{E}$, $\underline{E}'$ correspondant à $E,E'$), tandis qu'un germe d'homomorphisme de fibrés vectoriels $E \longrightarrow E'$ revient à la donnée d'un homomorphisme de $\underline{O}_s$-modules

\begin{equation*}
F : \underline{E}_s \longrightarrow \underline{E}'_s \quad .
\end{equation*}
Cela montre déjà que l'application (13.2.1.1) est injective, et pour la surjectivité, il suffit de prouver que si on a un homomorphisme $\hat{F}$ comme dessus, tel que pour tout $n$ l'homomorphisme induit $F_n : E_n \longrightarrow E'_n$ applique $R_n$ dans $R'_n$, alors $\hat{F}$ provient d'un $F$ i.e. applique $\underline{E}_s$ dans $\underline{E}'_s$ , et le germe correspondant d'homomorphismes $E \longrightarrow E'$ applique $R$ dans $R'$. Or, comme par hypothèse c), $R_s = R_o$ engendre $E_s$, il s'ensuit par Nakayama que le groupe $R_o$ (identifié au groupe des sections de $\hat{R}$ sur $\hat{S}$, ou,au choix,au groupe des germes de sections de $R$ sur $S$) engendre $\hat{\underline{E}}_s$ ; donc, pour prouver que $\hat{F}(\underline{E}_s) \subset \underline{E}'_s$ ,il suffit de prouver que $\hat{F}(R_o) \subset \underline{E}'_s$ . Or on a $\hat{F}(R_o) \subset R'_o$, où $R'_o$ est lui-même interprété comme le groupe des germes de sections de $R' \subset E'$ au voisinage de $s$, d'où l'inclusion annoncée, et par suite l'existence de $F$ . De plus on voit que $F(R_o) \subset R'_o$ , où $R_o$ et $R'_o$ sont considérés comme des groupes de germes de sections de $E$ resp. $E'$ . En vertu de b), $R_o$


}{F7I-IX-171}\SGApage{I}{489}

engendre $R$ au voisinage de $s$, d'où il résulte que $F(R) \subset R'$ au voisinage de $s$, ce qui achève la démonstration.

\SGAvar{\underline{Corollaire} 13.2.2. \underline{Si} $A$ \underline{satisfait les conditions} a), b), c) \underline{de} 13.2.2, \underline{il est déterminé à isomorphisme unique près au voisinage de} $s$ \underline{par la connaissance du} $\hat{S}$\underline{-groupe} $\hat{A}$.


}{\underline{Corollaire} 13.2.2. \underline{Si} $A$ \underline{satisfait les conditions} a), b), c) \underline{de} 13.2.1, \underline{il est déterminé à isomorphisme unique près au voisinage de} $s$ \underline{par la connaissance du} $\hat{S}$\underline{-groupe} $\hat{A}$.


}{F7I-IX-172}13.3. \underline{L'extension de Raynaud analytique}.

\SGAvar{Nous sommes maintenant en mesure de paraphraser, en Géométrie Analytique, la construction de l'extension de Raynaud donnée dans 7 . Supposons donné $A$ comme dans 13.1, avec $A = A_o$, et considérons $R_o = R_3$, qui est un $\mathbb{Z}$-module libre de type fini. Quitte à restreindre $S$ à un voisinage ouvert assez petit de $s$, interprétant $R_o$ comme le groupe des germes de sections de $R$ au voisinage de $s$, on trouve un homomorphisme

\begin{equation}\tag{13.3.1}
R_{oS} \longrightarrow R \subset E
\end{equation}
où $R_{oS}$ désigne le groupe analytique constant sur $S$ de valeur $R_o$ . Choisissant le voisinage de $s$ assez petit pour que toute section de $\underline{O}_S$ dont le germe en $s$ est nul soit nul, on voit que l'homomorphisme précédent est injectif. Son image est donc un sous-groupe analytique constant de $R$, qui est évidemment le plus grand sous-groupe constant de $R$, et qui s'appelle la partie \underline{fixe} de $R$, et est notée $R^f$ :

}{Nous sommes maintenant en mesure de paraphraser, en Géométrie Analytique, la construction de l'extension de Raynaud donnée dans 7 . Supposons donné $A$ comme dans 13.1, avec $A = A^o$, et considérons $R_o = R_s$, qui est un $\mathbb{Z}$-module libre de type fini. Quitte à restreindre $S$ à un voisinage ouvert assez petit de $s$, interprétant $R_o$ comme le groupe des germes de sections de $R$ au voisinage de $s$, on trouve un homomorphisme

\begin{equation}\tag{13.3.1}
R_{oS} \longrightarrow R \subset E
\end{equation}
où $R_{oS}$ désigne le groupe analytique constant sur $S$ de valeur $R_o$ . Choisissant le voisinage de $s$ assez petit pour que toute section de $\underline{O}_S$ dont le germe en $s$ est nul soit nulle, on voit que l'homomorphisme précédent est injectif. Son image est donc un sous-groupe analytique constant de $R$, qui est évidemment le plus grand sous-groupe constant de $R$, et qui s'appelle la partie \underline{fixe} de $R$, et est notée $R^f$ :

}{F7I-IX-173}\begin{equation}\tag{13.3.2}
R_{oS} \simeq R^f \hookrightarrow R \ , \ \mbox{induit un isomorphisme } R^f_s \simeq R_s \ .
\end{equation}
On pourra alors considérer le groupe quotient

\begin{equation}\tag{13.3.3}
A^\natural = E/R^f
\end{equation}

\SGApage{I}{490}

et on trouve une suite exacte canonique, compte tenu de (13.1.3) :

\begin{equation}\tag{13.3.4}
0 \longrightarrow R/R^f \stackrel{i}{\longrightarrow} A^\natural \longrightarrow A \longrightarrow 0 \ .
\end{equation}
On notera que $R^f$ est toujours fermé dans $E$, à condition de restreindre $S$ à un voisinage de $s$, et que par suite $A^\natural$ est séparé sur $S$. On notera aussi que l'homomorphisme $A^\natural \longrightarrow A$ est étale, donc induit un isomorphisme

\begin{equation}\tag{13.3.5}
\underline{\mathrm{Lie}}(A^\natural) \stackrel{\sim}{\longrightarrow} \underline{\mathrm{Lie}}(A) = E \ .
\end{equation}
D'autre part, pour tout $n \geqslant 0$, l'homomorphisme canonique $A^\natural \longrightarrow A$ induit un isomorphisme $A^\natural_n \simeq A_n$ (puisque $R^f_n \stackrel{\sim}{\longrightarrow} R_n$), donc il induit un isomorphisme

\begin{equation}\tag{13.3.6}
\widehat{A^\natural} \stackrel{\sim}{\longrightarrow} \hat{A} \ .
\end{equation}
Ce dernier fait montre que $A^\natural$ joue bien le rôle du groupe, encore noté $A^\natural$ , que nous avions envisagé dans le par. 7 . Il a été construit ici sans faire d'hypothèse sur la fibre $A_s = A_o$, du type ``extension d'un groupe abéloide par un tore'', ce qui sera nécessaire cependant si nous voulons récupérer l'importante structure d'extension (7.2.3) sur $A^\natural$. En revanche, l'homomorphisme $A^\natural \longrightarrow A$ de (13.3.4) est de nature essentiellement transcendante et n'a pas de contre-partie dans la théorie purement algébrique des \S~7. Nous reviendrons d'autre part au \S~14 sur la question d'une interprétation purement algébrique de l'homomorphisme $i : R/R^f \longrightarrow A^\natural$ de (13.3.4).

Supposons d'abord qu'on se donne un sous-tore complexe

\begin{equation}\tag{13.3.7}
T_o \subset A_o \ ,
\end{equation}
où par ``\underline{tore complexe}'' on entend un groupe analytique isomorphe à une


\SGApage{I}{491}

puissance cartésienne de $\mathbb{C}^*$ . Soit

\begin{equation}\tag{13.3.8}
T = T_o \times S
\end{equation}
le tore constant sur $S$ de valeur $T_o$ . Procédant comme dans 5.1, on trouve un unique homomorphisme (et même une immersion) de systèmes inductifs de groupes analytiques

\begin{equation*}
\hat{T} \longrightarrow \hat{A}^{\natural} \quad ,
\end{equation*}
qui au-dessus de $S_o$ se réduise à (13.3.7). En vertu de 13.2.2, cet homomorphisme provient d'un unique germe d'homomorphismes

\begin{equation}\tag{13.3.9}
T \longrightarrow A^{\natural} \quad ,
\end{equation}
et prenant $S$ assez petit autour de $s$, on peut supposer cet homomorphisme défini au-dessus de $S$ tout entier. L'homomorphisme correspondant

\SGAvar{\begin{equation*}
\underline{Lie}(T) \longrightarrow \underline{Lie}(A^{\natural}) \simeq E
\end{equation*}
est injectif sur les $S_n$ donc sur $\hat{S}$, donc sur $\underline{O}_s$ par descente plate, donc il existe un voisinage de $S$, qu'on peut supposer égal à $S$, en lequel $\underline{Lie}(T)$ s'identifie à un sous-Module localement facteur direct de $\underline{E}$, que nous noterons $\underline{E}^t$ ; le sous-fibré correspondant de $E$ est noté $E^t$:

}{\begin{equation*}
\underline{Lie}(T) \longrightarrow \underline{Lie}(A^{\natural}) \simeq E
\end{equation*}
est injectif sur les $S_n$ donc sur $\hat{S}$, donc sur $\underline{O}_s$ par descente plate, donc il existe un voisinage de $s$, qu'on peut supposer égal à $S$, en lequel $\underline{Lie}(T)$ s'identifie à un sous-Module localement facteur direct de $\underline{E}$, que nous noterons $\underline{E}^t$ ; le sous-fibré correspondant de $E$ est noté $E^t$:

}{F7I-IX-174}\begin{equation}\tag{13.3.10}
\underline{Lie}(T) \xrightarrow{\ \sim\ } \underline{E}^{t} \subset \underline{E} \simeq \underline{Lie}(A^{\natural}) \simeq \underline{Lie}(A) \quad , \quad E^{t} \subset E \quad .
\end{equation}
De même, désignons par

\begin{equation}\tag{13.3.11}
R_o^{t} \subset R_o \quad \mathrm{et} \quad R^{t} = R_{oS}^{t} \subset R_{oS} \simeq R^{f}
\end{equation}
le sous-réseau de $R_o$ défini par le sous-tore $T_o$ de $A_o$ , et le groupe analytique constant correspondant sur $S$, de sorte que l'homomorphisme


\SGApage{I}{492}

(13.3.9) correspond simplement à l'homomorphisme

\begin{equation*}
E^{t}/R_{oS}^{t} = E^{t}/R^{t} \longrightarrow E/R_{oS} = E/R^{f}
\end{equation*}
induit par les inclusions

\begin{equation*}
E^{t} \hookrightarrow E, \quad R^{t} \hookrightarrow R^{f} \quad .
\end{equation*}
Utilisant le fait que (13.3.9) induit une immersion sur les fibres en $s$, on en conclut que l'homomorphisme

\begin{equation*}
R_o/R_o^{t} \longrightarrow E_o/E_o^{t} = (E/E^{t})_o
\end{equation*}
est injectif, d'où il résulte aussitôt que

\begin{equation*}
R_o/R_o^{t} \longrightarrow (E/E^{t})_{s'}
\end{equation*}
est encore injectif pour $s'$ voisin de $s$, ce qui implique que, quitte à rapetisser $S$ autour de $s$, on peut supposer que (13.3.9) \underline{est une immersion fermée}.

Il s'impose alors d'introduire

\begin{equation}\tag{13.3.12}
B = A^{\natural}/T \simeq (E/E^{t})/(R^{f}/R^{t}) \quad ,
\end{equation}
de sorte qu'on obtient sur $A^{\natural}$ une structure d'extension

\SGAvar{\begin{equation}\tag{13.3.13}
0 \longrightarrow T \longrightarrow A^{\natural} \longrightarrow B \longrightarrow 0 \quad ,
\end{equation}
$B$ étant un groupe analytique relatif sur $S$, lisse et séparé sur $S$ et à fibres connexes. Par contruction même, on a donc

}{\begin{equation}\tag{13.3.13}
0 \longrightarrow T \longrightarrow A^{\natural} \longrightarrow B \longrightarrow 0 \quad ,
\end{equation}
$B$ étant un groupe analytique relatif sur $S$, lisse et séparé sur $S$ et à fibres connexes. Par construction même, on a donc

}{F7I-IX-175}\begin{equation}\tag{13.3.14}
\underline{Lie}(B) \simeq \underline{E}/\underline{E}^{t} \quad ,
\end{equation}
et la suite exacte exponentielle pour $B$


\SGApage{I}{493}

\begin{equation}\tag{13.3.15}
0 \longrightarrow R^f/R^t \longrightarrow E/E^t \longrightarrow B \longrightarrow 0 \quad .
\end{equation}
On notera que

\begin{equation}\tag{13.3.16}
R^f/R^t = (R_o/R_o^t)_S
\end{equation}
est un groupe analytique constant sur $S$. Il résulte donc de 13.1.5 que $B$ est un groupe abéloide sur $S$ au voisinage de $s$ si et seulement si sa fibre en $s$

\begin{equation*}
B_o = A_o/T_o
\end{equation*}
est un groupe analytique abéloide, de sorte que $A_o$ apparaît alors comme une extension d'un tel groupe $B_o$ par un tore $T_o$ . \underline{Toute façon d'écrire} $A_o$ \underline{sous forme d'une extension d'un groupe abéloide par un tore complexe} (*) \underline{donne donc naissance sur} $A$ , \underline{au voisinage de} $s$, \underline{à une unique structure d'extension} (13.3.13) \underline{d'un groupe abéloide} $B$ \underline{par un tore} $T$, \underline{induisant la structure d'extension donnée en} $s$.


13.4. \underline{Interprétation en termes de} ``\underline{structures de Hodge mixtes}'' \underline{de DELIGNE}.

\SGAvar{Reprenons d'abord le dictionnaire 13.1.4, que nous allons reformuler d'une façon légèrement différente, mieux adaptée à l'introduction des ``structures de Hodge'' pertinentes à la situation. Pour ceci, soit $\mathbb{R}$ le faisceau des sections de $R$ sur $S$, et introduisons le $\underline{O}_S$-Module

}{Reprenons d'abord le dictionnaire 13.1.4, que nous allons reformuler d'une façon légèrement différente, mieux adaptée à l'introduction des ``structures de Hodge'' pertinentes à la situation. Pour ceci, soit $\underline{R}$ le faisceau des sections de $R$ sur $S$, et introduisons le $\underline{O}_S$-Module

}{F7I-IX-176}\begin{equation}\tag{13.4.1}
\mathcal{E}(A) = \mathcal{E} = \underline{R} \otimes_{\mathbb{Z}} \underline{O}_S \quad ,
\end{equation}
qui est localement libre de type fini lorsque $R$ est localement constant. (NB Le faisceau $\mathcal{E}$ n'est pas nécessairement cohérent si $R$ n'est pas localement constant.) L'homomorphisme (13.1.3) $\underline{R} \longrightarrow \underline{E}$ définit alors


\SGApage{I}{494}

un homomorphisme canonique

\begin{equation}\tag{13.4.2}
\mathcal{E} = \underline{R} \otimes_{\mathbb{Z}} \underline{O}_S \longrightarrow \underline{E} \quad ,
\end{equation}
dont le noyau sera noté $\underline{F}(A)$ ou simplement $\underline{F}$, de sorte qu'on a une suite exacte

\begin{equation}\tag{13.4.3}
0 \longrightarrow \underline{F} \longrightarrow \mathcal{E} \longrightarrow \underline{E} \quad .
\end{equation}
Le cas le plus intéressant est celui où l'homomorphisme (13.4.2) est un épimorphisme, i.e. où pour tout $s \in S$, $R_s$ engendre le vectoriel complexe $E_s$ . Dans ce cas on trouve donc une suite exacte

\begin{equation}\tag{13.4.4}
0 \longrightarrow \underline{F} \longrightarrow \mathcal{E} \longrightarrow \underline{E} \longrightarrow 0 \quad ,
\end{equation}
et si $R$ est localement constant, c'est là une suite exacte courte de Modules localement libres sur $S$. La connaissance du couple $(E,R)$ est alors équivalente à celle du couple $(R,\underline{F})$, où $R$ est un groupe analytique localement constant sur $S$ à fibres des $\mathbb{Z}$ -Modules libres de type fini, et $\underline{F}$ est un sous-Module localement facteur direct de $\mathcal{E} = \underline{R} \otimes_{\mathbb{Z}} \underline{O}_S$ , tel que l'on ait de plus la relation

\begin{equation}\tag{13.4.5}
\underline{F} \cap \underline{R} = 0 \quad .
\end{equation}
Du point de vue des structures de Hodge, il y a lieu de regarder $\mathcal{E}$ comme filtré par la ``\underline{filtration de Hodge}''

\SGAvar{\begin{equation}\tag{13.4.6}
\mathrm{Fil}^{1}(\mathcal{E}) = \mathcal{E} \ , \quad \mathrm{Fil}^{0}(\mathcal{E}) = \underline{F} \ , \quad \mathrm{Fil}^{-1}(\mathcal{E}) = 0 \quad .
\end{equation}
Notons que la suite exacte (13.4.3) dépend encore fonctoriellement du groupe analytique $A$ sur $S$, et qu'elle commute aux changements de base quelconques.


}{\begin{equation}\tag{13.4.6}
\mathrm{Fil}^{-1}(\mathcal{E}) = \mathcal{E} \ , \quad \mathrm{Fil}^{0}(\mathcal{E}) = \underline{F} \ , \quad \mathrm{Fil}^{1}(\mathcal{E}) = 0 \quad .
\end{equation}
Notons que la suite exacte (13.4.3) dépend encore fonctoriellement du groupe analytique $A$ sur $S$, et qu'elle commute aux changements de base quelconques.


}{F7I-IX-177}\SGApage{I}{495}

Revenant alors aux conditions de 13.3, nous posons

\begin{equation}\tag{13.4.7}
\mathcal{E}^f = \underline{R}^f \otimes_{\mathbb{Z}} \underline{O}_S \quad , \quad \mathcal{E}^t = \underline{R}^t \otimes_{\mathbb{Z}} \underline{O}_S \quad ,
\end{equation}
ces Modules étant définis dans un voisinage convenable de $s$ (qu'on peut supposer égal à $S$), et définissent donc une filtration de $\mathcal{E}$ (appelée la \underline{filtration par le poids})

\SGAvar{\begin{equation}\tag{13.4.8}
\mathcal{E} \supset \mathcal{E}^f \supset \mathcal{E}^t \supset 0 \quad ,
\end{equation}
où il y a lieu de regarder $\mathcal{E}^t$, $\mathcal{E}^f$ et $\mathcal{E}$ respectivement comme de poids $-2,-1$ et $0$ (il s'agit donc d'une filtration croissante). On notera que comme $R^t$, $\mathcal{E}^t$ dépend du choix d'un sous-tore $T_o$ de $A_o = \Lambda_s$ . On notera que comme $R^f$ et $R^t$ sont localement constants, $\mathcal{E}^f$ \underline{et} $\mathcal{E}^t$ \underline{sont localement libres}, et que l'on peut aussi les interpréter comme

}{\begin{equation}\tag{13.4.8}
\mathcal{E} \supset \mathcal{E}^f \supset \mathcal{E}^t \supset 0 \quad ,
\end{equation}
où il y a lieu de regarder $\mathcal{E}^t$, $\mathcal{E}^f$ et $\mathcal{E}$ respectivement comme de poids $-2,-1$ et $0$ (il s'agit donc d'une filtration croissante). On notera que comme $R^t$, $\mathcal{E}^t$ dépend du choix d'un sous-tore $T_o$ de $A_o = A_s$ . On notera que comme $R^f$ et $R^t$ sont localement constants, $\mathcal{E}^f$ \underline{et} $\mathcal{E}^t$ \underline{sont localement libres}, et que l'on peut aussi les interpréter comme

}{F7I-IX-178}\begin{equation}\tag{13.4.9}
\mathcal{E}^f = \mathcal{E}(A^\natural) \quad , \quad \mathcal{E}^t = \mathcal{E}(T) \quad .
\end{equation}
La fonctorialité de la suite exacte (13.4.3) en $A$ variable, appliquée aux homomorphismes

\begin{equation*}
T \longrightarrow A^\natural \longrightarrow A \quad ,
\end{equation*}
nous montre alors aussitôt que $\underline{F}(A^\natural)$ et $\underline{F}(T)$ s'identifient respectivement à

\begin{equation}\tag{13.4.10}
\underline{F}(A^\natural) \simeq \underline{F}^f \stackrel{\mathrm{dfn}}{=} \underline{F} \cap \mathcal{E}^f \quad , \quad \underline{F}(T) = \underline{F}^t = \underline{F} \cap \mathcal{E}^t = 0 \quad ,
\end{equation}
le $=0$ provenant du fait immédiat que pour un tore $T$, on a $\underline{F}(T) = 0$, la suite exacte du type (13.4.3) devenant ici $0 \to 0 \to \mathcal{E}(T) \to \underline{Lie}(T) \to 0$. On voit donc en même temps que l'homomorphisme composé

\begin{equation*}
\mathcal{E}^t \longrightarrow \mathcal{E} \longrightarrow \underline{E}
\end{equation*}
induit un isomorphisme


\SGApage{I}{496}

\begin{equation}\tag{13.4.11}
\mathcal{E}^t \stackrel{\sim}{\longrightarrow} \underline{E}^t \ ( \simeq \underline{Lie}(T)) \quad .
\end{equation}
\underline{Supposons maintenant que} $R_s$ \underline{engendre l'espace vectoriel complexe} $E_s$ , ce qui est le cas en particulier dans le cas favorable 13.3.16 où $A_s$ est une extension d'un groupe abéloide par un tore complexe $T_o$ . Alors, quitte à restreindre $S$ à un voisinage convenable de $s$, on trouve que l'homomorphisme canonique du type (13.4.2) (avec $A$ remplacé par $A^\natural$)

\begin{equation*}
\mathcal{E}^f = \mathcal{E}(A^\natural) \longrightarrow \underline{Lie}(A^\natural) \simeq \underline{E}
\end{equation*}
est un épimorphisme, et a fortiori (13.4.2) est un épimorphisme. On a donc des suites exactes (13.4.4) et

\begin{equation}\tag{13.4.12}
0 \longrightarrow \underline{F}^f \longrightarrow \mathcal{E}^f \longrightarrow \underline{E} \longrightarrow 0 \quad ,
\end{equation}
ce qui montre en particulier que l'on a

\SGAvar{\begin{equation}\tag{13.4.13}
\underline{F} + \mathcal{E}^f = \mathcal{E} \quad ,
\end{equation}
puisque $\mathcal{E}/\underline{F} \simeq \underline{E}$ . De plus, comme $B = A^\natural/T$ est un quotient de $A$, il satisfait à la condition de surjectivité analogue, i.e. donne naissance à une suite exacte du type (13.4.4) :

\begin{equation}\tag{13.4.14}
0 \longrightarrow \underline{F}(B) \longrightarrow (B) \longrightarrow \underline{Lie}(B) \longrightarrow 0 \quad ,
\end{equation}
dont l'interprétation est immédiate en termes des Modules déjà introduits. On a en effet par (13.3.12)

}{\begin{equation}\tag{13.4.13}
\underline{F} + \mathcal{E}^f = \mathcal{E} \quad ,
\end{equation}
puisque $\mathcal{E}/\underline{F} \simeq \underline{E}$ . De plus, comme $B = A^\natural/T$ est un quotient de $A^\natural$, il satisfait à la condition de surjectivité analogue, i.e. donne naissance à une suite exacte du type (13.4.4) :

\begin{equation}\tag{13.4.14}
0 \longrightarrow \underline{F}(B) \longrightarrow \mathcal{E}(B) \longrightarrow \underline{Lie}(B) \longrightarrow 0 \quad ,
\end{equation}
dont l'interprétation est immédiate en termes des Modules déjà introduits. On a en effet par (13.3.12)

}{F7I-IX-179}\begin{equation}\tag{13.4.15}
\left\{ \begin{array}{l} \mathcal{E}(B) \simeq (\underline{R}^f / \underline{R}^t) \otimes_{\mathbb{Z}} \underline{O}_S \simeq \mathcal{E}^f / \mathcal{E}^t \quad , \\[3ex] \underline{Lie}(B) \simeq \underline{E} / \underline{E}^t \qquad\qquad\qquad , \end{array} \right.
\end{equation}

\SGApage{I}{497}

d'où on conclut, grâce à (13.4.11) et à la relation $\underline{F} \cap \mathcal{E}^t = 0$ (13.4.10), que $\underline{F}(B) = \mathrm{Ker}\ (\mathcal{E}^f/\mathcal{E}^t \longrightarrow \underline{E}/\underline{E}^t)$ est isomorphe à $\underline{F}^f$ par l'homomorphisme

\SGAvar{\begin{equation}\tag{13.4.16}
\underline{F}^f \simeq \underline{F}(A^\natural) \overset{\sim}{\longrightarrow} \underline{F}(B) \qquad .
\end{equation}
Récapitulons alors les propriétés obtenues de la filtration "par le poids" (13.4.8), et ses relations avec la "filtration de Hodge" (13.4.6) de $\mathcal{E}$ défini par le sous-Module $\underline{F}$ de $\mathcal{E}$, dans le cas favorable 13.3.16 où $A_s$ est une extension d'un groupe analytique abéloide $B_o$ par un tore $T_o$ (qu'on utilisera pour définir $R^t$, de sorte que $B$ est un groupe analytique abéloide sur $S$) :

}{\begin{equation}\tag{13.4.16}
\underline{F}^f \simeq \underline{F}(A^\natural) \overset{\sim}{\longrightarrow} \underline{F}(B) \qquad .
\end{equation}
Récapitulons alors les propriétés obtenues de la filtration "par le poids" (13.4.8), et ses relations avec la "filtration de Hodge" (13.4.6) de $\mathcal{E}$ définie par le sous-Module $\underline{F}$ de $\mathcal{E}$, dans le cas favorable 13.3.16 où $A_s$ est une extension d'un groupe analytique abéloide $B_o$ par un tore $T_o$ (qu'on utilisera pour définir $R^t$, de sorte que $B$ est un groupe analytique abéloide sur $S$) :

}{F7I-IX-180}\underline{Proposition} 13.4.17. \underline{Sous les conditions précédentes, les filtrations} \underline{sur} $\mathcal{E} = \underline{R} \otimes_{\mathbb{Z}} \underline{O}_S$ "\underline{de Hodge}" \underline{et} "\underline{par le poids}" :

\begin{equation}\tag{13.4.17.1}
\mathcal{E} \supset \underline{F} \quad , \quad \mathcal{E} \supset \mathcal{E}^f \supset \mathcal{E}^t \supset 0
\end{equation}
\underline{satisfont aux conditions suivantes} :

a) \underline{La filtration par le poids provient,} \underline{par tensorisation} $\otimes_{\mathbb{Z}} \underline{O}_S$ , \underline{d'une filtration analogue sur} $\underline{R}$

\begin{equation*}
\underline{R} \ \supset \underline{R}^f \supset \underline{R}^t \supset 0 \qquad ,
\end{equation*}
\underline{où} $\underline{R}^f$ \underline{et} $\underline{R}^t$ \underline{sont des groupes analytiques localement constants sur} $S$, \underline{de valeur des} $\mathbb{Z}$ \underline{-modules libres de type fini} (\underline{savoir} $R_o^f = R_s = \pi_1(A_s)$ \underline{et} $R_o^t = \pi_1(T_o)$).

b) \underline{On a les relations}

\begin{equation}\tag{*}
\underline{F} \cap \mathcal{E}^t = 0 \ , \ \underline{F} + \mathcal{E}^f = \mathcal{E} \qquad .
\end{equation}

\SGApage{I}{498}

\underline{En d'autres termes,} \underline{la filtration de} $\mathcal{E}^t$ \underline{induite par la filtration de} \underline{Hodge} (\underline{avec les conventions de degrés} (13.4.6)) \underline{est concentrée en} \underline{degré} $-1$, \underline{celle de} $\mathcal{E}/\mathcal{E}^f$ \underline{induite par la filtration de Hodge est concentrée} \underline{en degré} $0$.

c) \underline{Considérons la filtration de}

\begin{equation*}
\mathcal{E}^{ab} \overset{\mathrm{dfn}}{=} \mathcal{E}^f/\mathcal{E}^t \simeq (\underline{R}^{ab}) \otimes_{\mathbb{Z}} \underline{O}_S \quad , \quad \text{où } \underline{R}^{ab} = \underline{R}^f/\underline{R}^t \ \ (\text{cf. a})) \qquad ,
\end{equation*}
\underline{induite par la filtration de Hodge de} $\mathcal{E}$,\underline{donc} \underline{définie par le sous-Module} $\underline{F}^{ab}$ \underline{de} $\mathcal{E}^{ab}$ \underline{image de}

\begin{equation*}
\underline{F}^f \overset{\mathrm{dfn}}{=} \underline{F} \cap \mathcal{E}^f \longrightarrow \mathcal{E}^{ab} \qquad ,
\end{equation*}
(\underline{la flèche étant un monomorphisme grâce à la première relation} $(*)$). \underline{Alors le couple} ($\underline{R}^{ab}$, $\underline{F}^{ab} \subset \underline{R}^{ab} \otimes_{\mathbb{Z}} \underline{O}_S = \mathcal{E}^{ab}$) \underline{est isomorphe au couple} ($\underline{R}(B)$, $\underline{F}(B)$) \underline{associé à un} \textbf{groupe abéloide} \underline{sur} $S$, \underline{comme expliqué après} (13.4.4) (\underline{où donc} $B \simeq \mathcal{E}^{ab}/(F^{ab}+ R^{ab})$, \underline{les symboles non soulignés} \underline{désignant les groupes analytiques sur} $S$ \underline{correspondants}).

\SGAvar{\underline{Remarque} 13.4.18. Dans la terminologie de DELIGNE [4], la conclusion de la proposition précédente, lorsqu'on se restreint au-dessus d'un ouvert $U$ de $S$ sur lequel $R$ est localement constant, exprime que $R$, muni des deux filtrations (13.4.17.1) de $\mathcal{E} = \underline{R} \otimes_{\mathbb{Z}} \underline{O}_S$ , est une "famille analytique de structures de Hodge mixtes sur $U$ ". La situation qu'on vient d'expliciter, a servie de modèle aux conjectures de DELIGNE sur la théorie de Néron en dimensions supérieures, énoncées dans [9, 9.8 à 9.13]. Il y a cependant un trait impcrtant de la théorie algébrique


}{\underline{Remarque} 13.4.18. Dans la terminologie de DELIGNE [4], la conclusion de la proposition précédente, lorsqu'on se restreint au-dessus d'un ouvert $U$ de $S$ sur lequel $R$ est localement constant, exprime que $R$, muni des deux filtrations (13.4.17.1) de $\mathcal{E} = \underline{R} \otimes_{\mathbb{Z}} \underline{O}_S$ , est une "famille analytique de structures de Hodge mixtes sur $U$ ". La situation qu'on vient d'expliciter, a servi de modèle aux conjectures de DELIGNE sur la théorie de Néron en dimensions supérieures, énoncées dans [9, 9.8 à 9.13]. Il y a cependant un trait important de la théorie algébrique


}{F7I-IX-181}\SGApage{I}{499}

étudiée au par. 7 qui manque dans la description 13.4.17, lorsqu'on suppose que $A$ est un groupe abéloide au-dessus d'un ouvert $U$ complémentaire d'une partie analytique $Z$ rare de $S$ (mettons $Z = \{s\}$, lorsque $\dim S = 1$), de sorte que (13.1.5) $R|U$ est un système local sur $U$ : c'est que le système local $R/R^t$ sur $U$ devrait être \underline{constant}, et plus précisément, qu'il devrait être dual au système local ${R'}^{f}$, où $R'$ est le groupe analytique étale sur $S$ qui correspondrait (comme $R$ à $A$) à un groupe analytique convenable $A'$ sur $S$, tel que $A'|U$ soit le groupe abéloide ``dual'' de $A|U$. (L'analogue algébrique de cet énoncé, avec les notations du par. 5, est la dualité de $T_{\ell}(A_K)/T_{\ell}(A_K)^t$ avec $T_{\ell}(A'_K)^f$. (5.5.9), résultant du \underline{théorème de dualité} 5.2.) J'ignore si un tel énoncé est valable, sous les seules hypothèses de 13.4.17, ni même si l'action de $\pi_1(U,t)$ sur $R_t$ est unipotente ($t \in U$) ; et il ne semble pas que les méthodes naïves utilisées dans le présent paragraphe permettent de démontrer un tel résultat, même sous l'hypothèse supplémentaire où on suppose que $A|U$ soit munie d'une ``polarisation'' i.e. d'un faisceau inversible ample sur chaque fibre. C'est ce que savent démontrer cependant les ``semi-simple group people'', semble-t-il, du moins pour $S$ non singulier de dimension 1, cf. le rapport [9]. Dans le numéro qui suit, nous allons montrer que tout marche bien en faisant des hypothèses d'algébricité relative sur $A$, et faire en même temps le lien avec la construction du par. 7 .


\SGApage{I}{500}

\SGAvar{13.5. \underline{Cas d'une situation relativement algébrique}. Nous supposons maintenant que $A$ provient d'un schéma en groupes relatif sur $S$, lisse, séparé et à fibres des extensions de schémas abéliens par des tores. Nous désignerons ce schéma relatif par la même lettre $A$. (Nous nous autorisons cet abus de notation grâce au fait que le foncteur $G \longmapsto G^{an}$, qui va de la catégorie des schémas en groupes relatifs du type envisagé vers la catégorie des groupes analytiques sur $S$, est pleinement fidèle, comme il serait facile à établir grâce aux résultats du type GAGA relatifs [10].) Pour tout ce qui concerne la notion de schéma relatif, et les relations entre schémas relatifs sur $S$ et espaces analytiques sur $S$, nous renvoyons à [10] (cf. aussi [2, Exp. 15]). Si $\underline{O}$ désigne l'anneau local de $S$ en $s$, le schéma relatif $A$ est connu au voisinage de $s$ par la connaissance du schéma ordinaire $A_{\underline{O}}$ sur $\mathrm{Spec}(\underline{O})$, déduit de $A$ par localisation en $s$. Il est clair alors que le complété formel de $A_{\underline{O}}$ le long de sa fibre spéciale (isomorphe à $A_s = A_o$) est isomorphe au complété formel du groupe analytique $A$ le long de $A_s$ (13.2), en identifiant les $A_n = A \times_S S_n$ à des schémas algébriques sur $S$, par l'abus de notations introduit ci-dessus. Dans le cas favorable où $\hat{A}$, regardé comme schéma formel en groupes sur $\hat{\underline{O}}$, est algébrisable (7.2.1), ce qui est le cas en particulier si $S$ est normal en $s$ et $A$ quasi-projectif sur $S$ (7.1.5) on obtient donc un schéma en groupes $A^\natural_{\hat{\underline{O}}}$ sur $\mathrm{Spec}(\hat{\underline{O}})$, extension d'un schéma abélien $B_{\hat{\underline{O}}}$ par un tore. Pour simplifier, nous supposons que $B_{\hat{\underline{O}}}$ est projectif sur $\hat{\underline{O}}$ (ce qui est bien le cas, on le sait, si $\underline{O}$ donc $\hat{\underline{O}}$ est normal). On va alors prouver :


}{13.5. \underline{Cas d'une situation relativement algébrique}. Nous supposons maintenant que $A$ provient d'un schéma en groupes relatif sur $S$, lisse, séparé et à fibres des extensions de schémas abéliens par des tores. Nous désignerons ce schéma relatif par la même lettre $A$. (Nous nous autorisons cet abus de notation grâce au fait que le foncteur $G \longmapsto G^{an}$, qui va de la catégorie des schémas en groupes relatifs du type envisagé vers la catégorie des groupes analytiques sur $S$, est pleinement fidèle, comme il serait facile à établir grâce aux résultats du type GAGA relatifs [10].) Pour tout ce qui concerne la notion de schéma relatif, et les relations entre schémas relatifs sur $S$ et espaces analytiques sur $S$, nous renvoyons à [10] (cf. aussi [2, Exp. 15]). Si $\underline{O}$ désigne l'anneau local de $S$ en $s$, le schéma relatif $A$ est connu au voisinage de $s$ par la connaissance du schéma ordinaire $A_{\underline{O}}$ sur $\mathrm{Spec}(\underline{O})$, déduit de $A$ par localisation en $s$. Il est clair alors que le complété formel de $A_{\underline{O}}$ le long de sa fibre spéciale (isomorphe à $A_s = A_o$) est isomorphe au complété formel du groupe analytique $A$ le long de $A_s$ (13.2), en identifiant les $A_n = A \times_S S_n$ à des schémas algébriques sur $S$, par l'abus de notations introduit ci-dessus. Dans le cas favorable où $\hat{A}$, regardé comme schéma formel en groupes sur $\hat{\underline{O}}$, est algébrisable (7.2.1), ce qui est le cas en particulier si $S$ est normal en $s$ et $A$ quasi-projectif sur $S$ (7.1.5), on obtient donc un schéma en groupes $A^{\natural}_{\hat{\underline{O}}}$ sur $\mathrm{Spec}(\hat{\underline{O}})$, extension d'un schéma abélien $B_{\hat{\underline{O}}}$ par un tore. Pour simplifier, nous supposons que $B_{\hat{\underline{O}}}$ est projectif sur $\hat{\underline{O}}$ (ce qui est bien le cas, on le sait, si $\underline{O}$ donc $\hat{\underline{O}}$ est normal). On va alors prouver :


}{F7I-IX-182}\SGApage{I}{501}

\underline{Proposition} 13.5.1. \underline{Sous les conditions précédentes}, \underline{le groupe analytique} $A^\natural$ \underline{de} (13.3.3) \underline{provient}, \underline{au voisinage de} $s$, \underline{d'un schéma en groupes relatif déterminé à isomorphisme unique près}, \underline{extension d'un schéma abélien relatif projectif sur} $S$ \underline{par un tore, et si} $A^\natural_{\underline{O}}$ \underline{est le schéma en groupes qu'il définit sur} $\underline{O} = \underline{O}_{S,s}$ , \underline{alors le schéma en groupes} $A^\natural_{\underline{O}} \otimes_{\underline{O}} \hat{\underline{O}}$ \underline{est canoniquement isomorphe avec sa structure d'extension à l'extension de Raynaud} $A^\natural_{\hat{\underline{O}}}$ \underline{introduite précédemment}.

Ce dernier fait résultera immédiatement de la première assertion, compte tenu de 7.2.1 , puisque $A^\natural_{\underline{O}} \otimes_{\underline{O}} \hat{\underline{O}}$ a comme complété formel le long de la fibre spéciale $\hat{A}_{\underline{O}}$, en vertu de (13.3.6). Pour prouver l'algébricité de $A^\natural$ au voisinage de $s$, procédant comme dans 7.2.1, on est ramené aussitôt à prouver celle du quotient abélien $B$ de $A^\natural$, ce qui va résulter du

\SGAvar{\underline{Lemme} 13.5.2. \underline{Soient} $B$ \underline{un groupe analytique abéloide sur l'espace analytique} $S$, $s$ \underline{un point de} $S$ \underline{tel que le complété} $\hat{B}$ \underline{le long de la fibre} $B_s$ \underline{soit algébrisable}, \underline{i.e.} \underline{tel que les} $B_n = B \times_S S_n$ \underline{soient algébrisables}, \underline{et que} $\hat{B}$, \underline{regardé maintenant comme} schéma formel \underline{sur} $\hat{\underline{C}}_{S,s}$ , \underline{soit algébrisable}, \underline{i.e.} \underline{provienne d'un schéma abélien} $B_{\hat{\underline{O}}}$ \underline{sur} $\hat{\underline{C}}$. \underline{Supposons de plus que} $B_{\hat{\underline{O}}}$ \underline{soit projectif sur} $\hat{\underline{O}}$ . \underline{Alors} $B$ \underline{est projectif sur} $S$ \underline{au voisinage de} $s$.

}{\underline{Lemme} 13.5.2. \underline{Soient} $B$ \underline{un groupe analytique abéloide sur l'espace analytique} $S$, $s$ \underline{un point de} $S$ \underline{tel que le complété} $\hat{B}$ \underline{le long de la fibre} $B_s$ \underline{soit algébrisable}, \underline{i.e.} \underline{tel que les} $B_n = B \times_S S_n$ \underline{soient algébrisables}, \underline{et que} $\hat{B}$, \underline{regardé maintenant comme} schéma formel \underline{sur} $\hat{\underline{O}}_{S,s}$ , \underline{soit algébrisable}, \underline{i.e.} \underline{provienne d'un schéma abélien} $B_{\hat{\underline{O}}}$ \underline{sur} $\hat{\underline{O}}$. \underline{Supposons de plus que} $B_{\hat{\underline{O}}}$ \underline{soit projectif sur} $\hat{\underline{O}}$ . \underline{Alors} $B$ \underline{est projectif sur} $S$ \underline{au voisinage de} $s$.

}{F7I-IX-183}L'hypothèse signifie qu'il existe une polarisation de l'espace analytique formel $\hat{B}$, i.e. une suite de polarisations $p_n$ des $B_n$ qui se ``recollent''. La conclusion signifie qu'il existe une polarisation $p$ de $B$ au-dessus d'un voisinage de $s$. Or soit $P$ l'espace analytique sur $S$


\SGApage{I}{502}

des polarisations relatives de $B$ sur $S$, qui représente le foncteur

\SGAvar{\begin{equation*}
S' \longmapsto \text{ensemble des polarisations de } B_{S'} \text{ sur } S'
\end{equation*}
sur la catégorie des espaces analytiques $S'$ sur $S$, lequel est un sous-foncteur ouvert du foncteur $\underline{NS}_{B/S}$ de Néron-Sévéri de $S$. Des arguments standarts des plus simples, que nous nous dispensons d'expliciter ici, montrent que $\underline{NS}_{B/S}$, et par suite aussi $P$, est représentable par un espace analytique \underline{net} (i.e. non ramifié) sur $S$. L'hypothèse s'exprime alors en disant qu'il existe une section \underline{formelle} $p$ de $P$ sur $S$ en $s$. Comme $P$ est net donc localement quasi-fini sur $S$, il s'ensuit [2, Exp. 18, n$^{\circ}$1] que toute section formelle de $P$ sur $S$ en $s$ provient d'une section analytique au voisinage de $s$, i.e. d'une polarisation de $B$ au-dessus d'un voisinage de $s$, d'où la conclusion annoncée.

\underline{Remarques} 13.5.3. a) Si on omet dans 13.5.2 l'hypothèse que $B_{\hat{\underline{O}}}$ est projectif sur $S$, on peut encore conclure que $B$ est algébrisable au voisinage de $s$. Pour ceci on utilise un résultat inédit (dû indépendamment à M. ARTIN et P. DELIGNE, en utilisant la théorie des faisceaux $\underline{T}^i$ de VI 2), affirmant que si $X$ propre sur $S$ est tel que $X_{\text{réd}}$ est algébrisable, il en est de même de $X$. Ceci nous ramène au cas où $S$ est réduit. D'autre part, lorsque $S$ est normal, donc $\hat{\underline{O}}$ est normal, on sait que tout schéma abélien sur $\hat{\underline{O}}$ est projectif sur $\hat{\underline{O}}$ (cf. par exemple [23]), et 13.5.2 s'applique. Pour obtenir le cas $S$ réduit quelconque, on procède par descente non plate à partir du normalisé $S'$ de $S$, grâce au résultat (délicat) suivant, que nous ne prouverons pas ici : un morphisme propre de schémas noethériens $f : S' \longrightarrow S$ qui est un épimor
\SGApage{I}{503}
phisme (i.e. tel que $\underline{O}_S \longrightarrow f_*(\underline{O}_{S'})$ soit injectif) est un \underline{morphisme de descente effective} pour la catégorie fibrée des schémas abéliens relatifs.

}{\begin{equation*}
S' \longmapsto \text{ensemble des polarisations de } B_{S'} \text{ sur } S'
\end{equation*}
sur la catégorie des espaces analytiques $S'$ sur $S$, lequel est un sous-foncteur ouvert du foncteur $\underline{NS}_{B/S}$ de Néron-Severi de $S$. Des arguments standards des plus simples, que nous nous dispensons d'expliciter ici, montrent que $\underline{NS}_{B/S}$, et par suite aussi $P$, est représentable par un espace analytique \underline{net} (i.e. non ramifié) sur $S$. L'hypothèse s'exprime alors en disant qu'il existe une section \underline{formelle} $p$ de $P$ sur $S$ en $s$. Comme $P$ est net donc localement quasi-fini sur $S$, il s'ensuit [2, Exp. 18, n$^{\circ}$1] que toute section formelle de $P$ sur $S$ en $s$ provient d'une section analytique au voisinage de $s$, i.e. d'une polarisation de $B$ au-dessus d'un voisinage de $s$, d'où la conclusion annoncée.

\underline{Remarques} 13.5.3. a) Si on omet dans 13.5.2 l'hypothèse que $B_{\hat{\underline{O}}}$ est projectif sur $\hat{\underline{O}}$, on peut encore conclure que $B$ est algébrisable au voisinage de $s$. Pour ceci on utilise un résultat inédit (dû indépendamment à M. ARTIN et P. DELIGNE, en utilisant la théorie des faisceaux $\underline{T}^i$ de VI 2), affirmant que si $X$ propre sur $S$ est tel que $X_{\text{réd}}$ est algébrisable, il en est de même de $X$. Ceci nous ramène au cas où $S$ est réduit. D'autre part, lorsque $S$ est normal, donc $\hat{\underline{O}}$ est normal, on sait que tout schéma abélien sur $\hat{\underline{O}}$ est projectif sur $\hat{\underline{O}}$ (cf. par exemple [23]), et 13.5.2 s'applique. Pour obtenir le cas $S$ réduit quelconque, on procède par descente non plate à partir du normalisé $S'$ de $S$, grâce au résultat (délicat) suivant, que nous ne prouverons pas ici : un morphisme propre de schémas noethériens $f : S' \longrightarrow S$ qui est un épimor
\SGApage{I}{503}
phisme (i.e. tel que $\underline{O}_S \longrightarrow f_*(\underline{O}_{S'})$ soit injectif) est un \underline{morphisme de descente effective} pour la catégorie fibrée des schémas abéliens relatifs.

}{F7I-IX-184}b) Il semble que la méthode de démonstration de 13.5.2 doive permettre aussi de prouver le résultat analogue pour tout espace analytique $B$ propre sur $S$. Plus délicate sans doute est la question de savoir si, pour un espace analytique $X$ propre sur $S$, tel que le complété formel $\hat{X}$ le long de $X_s$ soit algébrisable, $X$ est algébrisable sur $S$ au voisinage de $s$. Signalons seulement qu'il y a lieu ici de prendre la notion d'algébrisabilité relative de $X$ sur $S$ dans un sens plus large que dans [10], en entendant par là qu'au voisinage de $s$, $X$ provient d'un \underline{espace algébrique} sur $\underline{O}_{S,s}$ au sens de M. ARTIN [1].

\SGAvar{13.5.4. Plaçons-nous dans les conditions de 13.5.1. Choisissons par ailleurs, pour tout entiers $n > C$, l'isomorphisme standart de groupes sur $\mathbb{C}$

}{13.5.4. Plaçons-nous dans les conditions de 13.5.1. Choisissons par ailleurs, pour tout entier $n > 0$, l'isomorphisme standard de groupes sur $\mathbb{C}$

}{F7I-IX-185}\begin{equation*}
\mu_{n\mathbb{C}} \;\simeq\; (\mathbb{Z}/n\mathbb{Z})_{C} \;,\qquad m \longmapsto \exp(2i\pi m/n) \qquad ,
\end{equation*}
ce qui nous donne, par passage à la limite sur les puissances d'un nombre fixé $\ell$, un isomorphisme canonique de faisceaux $\ell$-adiques sur $\mathbb{C}$ :

\begin{equation*}
\mathbb{Z}_{\ell}(1)_{\mathbb{C}} \;\simeq\; \mathbb{Z}_{\ell\mathbb{C}} \qquad ,
\end{equation*}
que nous utiliserons pour identifier les deux membres.

De la suite exacte (13.1.3), on déduit, en appliquant la suite exacte des $\underline{\mathrm{Ext}}^{i}_{\mathbb{Z}}(\mathbb{Z}/n\mathbb{Z}\,,\ -)$, un isomorphisme canonique


\SGApage{I}{504}

\begin{equation}\tag{13.5.5}
{}_{n}A^{\mathrm{o}} \;\simeq\; \underline{R}\otimes_{\mathbb{Z}}\mathbb{Z}/n\mathbb{Z} \qquad ,
\end{equation}
qui donne, pour $n$ parcourant les puissances d'un nombre premier $\ell$ donné, un isomorphisme de systèmes projectifs de groupes sur $S$

\SGAvar{\begin{equation}\tag{13.5.6}
T_{\ell}(A^{\mathrm{o}}) \;\simeq\; \underline{R}\otimes_{\mathbb{Z}}\mathbb{Z}_{\ell} \qquad ,
\end{equation}
où $\mathbb{Z}_{\ell}$ est considéré comme le système projectif des $\mathbb{Z}/\ell^{n+1}\mathbb{Z}$ . Dans le cas envisagé ici où $A$ provient d'un schéma en groupes séparé sur $S$, on voit alors immédiatement que les $\underline{R}\otimes\mathbb{Z}/n\mathbb{Z}$ peuvent être considérés comme des schémas relatifs, ce qui donne un sens au symbole $(\underline{R}\otimes_{\mathbb{Z}}\mathbb{Z}_{\ell})_{\underline{O}}$ (alors que $R_{\underline{O}}$ n'en aurait pas en général !), et on obtient un isomorphisme canonique

}{\begin{equation}\tag{13.5.6}
T_{\ell}(A^{\mathrm{o}}) \;\simeq\; \underline{R}\otimes_{\mathbb{Z}}\mathbb{Z}_{\ell} \qquad ,
\end{equation}
où $\mathbb{Z}_{\ell}$ est considéré comme le système projectif des $\mathbb{Z}/\ell^{n+1}\mathbb{Z}$ . Dans le cas envisagé ici où $A$ provient d'un schéma en groupes séparé sur $S$, on voit alors immédiatement que les $\underline{R}\otimes\mathbb{Z}/n\mathbb{Z}$ peuvent être considérés comme des schémas relatifs, ce qui donne un sens au symbole $(\underline{R}\otimes_{\mathbb{Z}}\mathbb{Z}_{\ell})_{\underline{O}}$ (alors que $\underline{R}_{\underline{O}}$ n'en aurait pas en général !), et on obtient un isomorphisme canonique

}{F7I-IX-186}\begin{equation}\tag{13.5.7}
T_{\ell}(A_{\underline{O}}) \;\simeq\; (\underline{R}\otimes_{\mathbb{Z}}\mathbb{Z}_{\ell})_{\underline{O}} \qquad .
\end{equation}
Il résulte alors immédiatement des constructions du par. 2 et de 13.3 que l'isomorphisme précédent transforme partie fixe en partie fixe, partie torique en partie torique :

\begin{equation}\tag{13.5.8}
T_{\ell}(A_{\underline{O}})^{f} \;\simeq\; (\underline{R}^{f}\otimes_{\mathbb{Z}}\mathbb{Z}_{\ell})_{\underline{O}} \;,\quad T_{\ell}(A_{\underline{O}})^{t} \;\simeq\; (\underline{R}^{t}\otimes_{\mathbb{Z}}\mathbb{Z}_{\ell})_{\underline{O}} \qquad .
\end{equation}
Si $A'$ est un groupe analytique relatif sur $S$ satisfaisant aux mêmes hypothèses que $A$, et si $W_{\underline{O}}$ est une biextension de $(A_{\underline{O}},A'_{\underline{O}})$ par $\mathbb{G}_{m\underline{O}}$ , définissant donc au-dessus d'un voisinage de $s$ une biextension de $(A,A')$ par $\mathbb{G}_{mS}$ (en travaillant maintenant avec la notion au sens analytique complexe, obtenue en travaillant dans le site des espaces analytiques au-dessus de $S$, pour une topologie convenable ...), alors d'après VIII 2.2 $W$ définit des accouplements


\SGApage{I}{505}

\begin{equation}\tag{13.5.9}
(\underline{R}\otimes\mathbb{Z}_\ell) \times (\underline{R}'\otimes_{\mathbb{Z}}\mathbb{Z}_\ell) \simeq T_\ell(A) \times T_\ell(A') \to \mathbb{Z}_\ell(1)_S \simeq \mathbb{Z}_{\ell S} \quad ,
\end{equation}
et il est encore clair que l'accouplement correspondant

\begin{equation}\tag{13.5.10}
T_\ell(A_{\underline{O}}) \times T_\ell(A'_{\underline{O}}) \longrightarrow \mathbb{Z}_\ell(1)_S
\end{equation}
n'est autre que celui associé, par loc. cit., à la biextension $W_{\underline{O}}$ de $(A_{\underline{O}},A'_{\underline{O}})$. Il n'est d'ailleurs pas difficile de voir que, à condition de se restreindre à un voisinage de $s$, les accouplements (13.5.9) proviennent d'un accouplement unique (ne dépendant plus de $\ell$)

\SGAvar{\begin{equation}\tag{13.5.11}
\underline{R} \otimes \underline{R}' \longrightarrow \mathbb{Z} \quad .
\end{equation}
Tout ceci posé, les résultats algébriques concernant les $T_\ell(A_{\underline{O}})$ peuvent alors être utilisés pour obtenir des conclusions sur le faisceau $\underline{R}$ . On trouve ainsi que l'accouplement (13.5.11) est nul sur $\underline{R}^{t} \times \underline{R}'{}^{f}$ et sur $\underline{R}^{f} \times \underline{R}'{}^{t}$ (7.4.8), et que dans le cas où $S$ est non singulier de dimension 1, et que, dans $U = S-\{s\}$, $A$ et $A'$ sont deux schémas abéliens relatifs duals l'un de l'autre par $W$, $\underline{R}^{t}_{U}$ et $\underline{R}'{}^{f}_{U}$ ainsi que $\underline{R}^{f}_{U}$ et $R'{}^{t}_{U}$ sont exactement l'annulateur l'un de l'autre sur un voisinage assez petit de $s$ (7.4.8). Il en résulte alors, compte tenu que (13.5.11) est une dualité parfaite sur $U$, que l'on a une dualité sur $U$ entre $\underline{R}/\underline{R}^{t}$ et $\underline{R}'{}^{f}$, ce qui implique en particulier que $\underline{R}/\underline{R}^{t}$ est constant, et a fortiori que l'action de $\pi_1(U,t)$ sur $\underline{R}_t$ est unipotente d'échelon 2. On voit de même qu'on a sur $U$ une dualité parfaite entre $\underline{R}/\underline{R}^{f}$ et $\underline{R}'{}^{t}$ ; comme ce dernier est également isomorphe au tore dual du réseau $M_S$ sur $S$, où $M$ est le groupe défini dans (5.5.4), on trouve un isomorphisme canonique


}{\begin{equation}\tag{13.5.11}
\underline{R} \otimes \underline{R}' \longrightarrow \mathbb{Z} \quad .
\end{equation}
Tout ceci posé, les résultats algébriques concernant les $T_\ell(A_{\underline{O}})$ peuvent alors être utilisés pour obtenir des conclusions sur le faisceau $\underline{R}$ . On trouve ainsi que l'accouplement (13.5.11) est nul sur $\underline{R}^{t} \times \underline{R}'{}^{f}$ et sur $\underline{R}^{f} \times \underline{R}'{}^{t}$ (7.4.8), et que dans le cas où $S$ est non singulier de dimension 1, et que, dans $U = S-\{s\}$, $A$ et $A'$ sont deux schémas abéliens relatifs duals l'un de l'autre par $W$, $\underline{R}^{t}_{U}$ et $\underline{R}'{}^{f}_{U}$ ainsi que $\underline{R}^{f}_{U}$ et $\underline{R}'{}^{t}_{U}$ sont exactement l'annulateur l'un de l'autre sur un voisinage assez petit de $s$ (7.4.8). Il en résulte alors, compte tenu que (13.5.11) est une dualité parfaite sur $U$, que l'on a une dualité sur $U$ entre $\underline{R}/\underline{R}^{t}$ et $\underline{R}'{}^{f}$, ce qui implique en particulier que $\underline{R}/\underline{R}^{t}$ est constant, et a fortiori que l'action de $\pi_1(U,t)$ sur $\underline{R}_t$ est unipotente d'échelon 2. On voit de même qu'on a sur $U$ une dualité parfaite entre $\underline{R}/\underline{R}^{f}$ et $\underline{R}'{}^{t}$ ; comme ce dernier est également isomorphe au tore dual du réseau $M_S$ sur $S$, où $M$ est le groupe défini dans (5.5.4), on trouve un isomorphisme canonique


}{F7I-IX-187}\SGApage{I}{506}

\begin{equation}\tag{13.5.12}
\underline{R}_U/\underline{R}^{f}_{U} \simeq M_U \quad ,
\end{equation}
et l'homomorphisme $i$ de (13.3.4) peut s'écrire comme un homomorphisme

\begin{equation}\tag{13.5.13}
i : M_U \longrightarrow A^{\natural} \quad ,
\end{equation}
qu'on retrouvera dans un contexte un peu différent au par. suivant.


\SGAsectiontoc{14. Liens avec la théorie transcendante : cas rigide-analytique. L'homomorphisme canonique $\protect\underline{M}_K \longrightarrow A_K^{\natural o}$.}{14}
14. \underline{Liens avec la théorie transcendante} : \underline{cas rigide-analytique}. \underline{L'homomorphisme canonique} $\underline{M}_K \longrightarrow A_K^{\natural o}$.

\SGAvar{14.1. Nous supposons connu ici les fondements de la théorie des espaces rigide  analytiques, pour lesquels on pourra consulter [12] [13] [34].

}{14.1. Nous supposons connus ici les fondements de la théorie des espaces rigide-analytiques, pour lesquels on pourra consulter [12] [13] [34].

}{F7I-IX-188}Soit $S$ un trait complet, $A_K$ un schéma abélien sur son corps des fractions. On a construit alors dans 8.1, par descente à partir de l'extension de Raynaud dans le cas semi-stable, une extension $G_K = A_K^{\natural o}$ (8.1.1)

\begin{equation}\tag{14.1.1}
0 \longrightarrow T_K \longrightarrow G_K \longrightarrow B_K \longrightarrow 0
\end{equation}
d'un schéma abélien $B_K$ par un tore $T_K$ . Considérant le schéma abélien dual $A'_K$ , on a de même une extension $G'_K = A'{}^{\natural o}_{K}$

\begin{equation}\tag{14.1.2}
0 \longrightarrow T'_K \longrightarrow G'_K \longrightarrow B'_K \longrightarrow 0 \quad ,
\end{equation}
où $B'_K$ s'identifie d'ailleurs au schéma abélien dual de $B_K$ (8.2). On désigne, comme dans (8.1.8), par $\underline{M}_K$ et $\underline{M}'_K$ les réseaux sur $K$ duals de $T'_K$ et de $T_K$ respectivement :

\begin{equation}\tag{14.1.3}
\underline{M}_K = D(T'_K) \quad , \quad \underline{M}'_K = D(T_K) \quad .
\end{equation}

\SGApage{I}{507}

Donc les structures d'extension (14.1.2) et (14.1.1) peuvent s'interpréter respectivement (VIII 3.7) comme équivalant à la donnée d'homomorphisme

\begin{equation}\tag{14.1.4}
\underline{M}_K \xrightarrow{\ p_K\ } B_K \quad , \quad \underline{M}'_K \xrightarrow{\ p'_K\ } B'_K \qquad .
\end{equation}
Soient $W_K^{\natural}$ la biextension de $(B_K,B'_K)$ par $\mathbb{G}_{mK}$ qui définit la dualité canonique entre ces schémas abéliens (8.2), et $E_K$ la biextension de $(\underline{M}_K,\underline{M}'_K)$ par $\mathbb{G}_{mK}$ image inverse de $W_K^{\natural}$ par les homomorphismes (14.1.4). Notons qu'on vérifie immédiatement que les trivialisations de $E_K$ correspondent biunivoquement, de façon canonique, aux homomorphismes

\begin{equation}\tag{14.1.5}
i \ : \ \underline{M}_K \longrightarrow G_K
\end{equation}
qui relèvent $p_K : \underline{M}_K \longrightarrow B_K$, ou, symétriquement, aux homomorphismes

\SGAvar{\begin{equation}\tag{14.1.6}
i' \ : \ \underline{M}'_K \longrightarrow G'_K
\end{equation}
qui relèvent $P'_K : \underline{M}'_K \longrightarrow B'_K$ .

}{\begin{equation}\tag{14.1.6}
i' \ : \ \underline{M}'_K \longrightarrow G'_K
\end{equation}
qui relèvent $p'_K : \underline{M}'_K \longrightarrow B'_K$ .

}{F7I-IX-189}Ceci posé, désignant par un exposant ${}^{an}$ les groupes rigide-analytiques sur $K$ associés aux groupes algébriques envisagés [13], on définit dans [24] une suite exacte canonique

\begin{equation}\tag{14.1.7}
0 \longrightarrow \underline{M}_K^{an} \xrightarrow{\ i^{an}\ } G_K^{an} \xrightarrow{\ e\ } A_K^{an} \longrightarrow 0 \qquad ,
\end{equation}
dans laquelle l'homomorphisme $i^{an}$, qui peut aussi s'interpréter comme un homomorphisme $i: \underline{M}_K \longrightarrow G_K$ pour des schémas en groupes ordinaires, \underline{relève l'homomorphisme} $p$ (14.1.4), i.e. est de la forme envisagée dans (14.1.5). Cette suite exacte est l'analogue rigide-analytique de la suite exacte (13.3.4). Elle dépend de façon fonctorielle de $A_K$ , et


\SGApage{I}{508}

sa formation est compatible avec tout changement de traits $S' \longrightarrow S$. Ce dernier fait montre d'ailleurs que, pour définir et étudier (14.1.7), on est ramené par descente finie au cas où $A_K$ a une réduction semi-stable. Dans ce cas, les schémas en groupes $T_K,G_K,B_K,\underline{M}_K$, ... sont les fibres génériques de schémas en groupes $T,G,B,M$ ... sur $S$, $G$ étant une extension du schéma abélien $B$ par le tore $T$ (par. 7). On peut alors considérer le complété formel $\hat{G}$ de $G$ le long de sa fibre spéciale, et par construction on a un isomorphisme canonique (7.2.6)

\SGAvar{\begin{equation}\tag{14.1.8}
\hat{G} \xrightarrow{\ \sim\ } \widehat{A^{o}} \qquad ,
\end{equation}
d'où par passage aux ``fibres génériques'' (qui sont des groupes ridige-analytiquessur $K$) un isomorphisme canonique

}{\begin{equation}\tag{14.1.8}
\hat{G} \xrightarrow{\ \sim\ } \widehat{A^{o}} \qquad ,
\end{equation}
d'où par passage aux ``fibres génériques'' (qui sont des groupes rigide-analytiques sur $K$) un isomorphisme canonique

}{F7I-IX-190}\begin{equation}\tag{14.1.9}
\hat{G}^{an} \xrightarrow{\ \sim\ } \widehat{A^{o}}\,{}^{an} \qquad .
\end{equation}
On a d'autre part des homomorphismes canoniques (qui sont des immersions ouvertes rigide-analytiques, mais pas en général des isomorphismes)

\SGAvar{\begin{equation}\tag{14.1.10}
\hat{G}^{an} \lhook\joinrel\longrightarrow G_K^{\ an} \quad , \quad \widehat{A^{o}}\,{}^{an} \lhook\joinrel\longrightarrow A_K^{\ an} \qquad .
\end{equation}
Ceci dit, l'homomorphisme $e$ de (14.1.7) est caractérisé dé façon unique par la commutativité du diagramme

}{\begin{equation}\tag{14.1.10}
\hat{G}^{an} \lhook\joinrel\longrightarrow G_K^{\ an} \quad , \quad \widehat{A^{o}}\,{}^{an} \lhook\joinrel\longrightarrow A_K^{\ an} \qquad .
\end{equation}
Ceci dit, l'homomorphisme $e$ de (14.1.7) est caractérisé de façon unique par la commutativité du diagramme

}{F7I-IX-191}\begin{equation}\tag{14.1.11}
\begin{tikzcd}
\hat{G}^{an} \arrow[r] \arrow[d] \& \widehat{A^{o}}\,{}^{an} \arrow[d] \\
G_K^{an} \arrow[r] \& A_K^{an}
\end{tikzcd} \quad ,
\end{equation}










où les flèches non nommées sont celles de (14.1.9) et (14.1.10).


\SGApage{I}{509}

\underline{Remarques} 14.1.12. a) On notera que l'unicité de $e$ rendant commutatif le diagramme précédent est immédiate, l'espace rigide-analytique $G_K^{\mathrm{an}}$ étant connexe. L'existence de $e$, par contre, n'est nullement triviale, et nous nous contentons à ce sujet de renvoyer à la note citée de M. RAYNAUD. D'autre part, la définition donnée dans loc. cit. de l'homomorphisme $i^{\mathrm{an}}$ est également de nature transcendante, bien qu'il puisse être interprété comme étant un objet de nature essentiellement algébrique (14.1.5) (une fois la construction de $G$ lui-même, donc de $G_K$, acquise (ce qui se fait par les constructions de \underline{géométrie formelle} du par. 7)). Je n'en connais aucune description purement algébrico-géométrique, mais proposerai plus bas (14.4) une caractérisation ``arithmétique'', qu'on obtiendra en se ramenant à une situation de type fini sur l'anneau $\mathbb{Z}$ . Il semble donc que l'homomorphisme (14.1.5) soit, du point de vue de la géométrie algébrique, un élément de structure plus caché que ceux rencontrés au cours des paragraphes précédents.

b) Comme expliqué plus haut, l'homomorphisme $i$ est canoniquement associé à une trivialisation bien déterminée de la biextension $E_K$ de $(\underline{M}_K,\underline{M}'_K)$, ou encore à un homomorphisme bien déterminé (14.1.6) relevant $p'_K$ . On montre alors que ce dernier n'est autre que l'homomorphisme qui intervient dans la suite exacte

\begin{equation}\tag{14.1.13}
0 \longrightarrow \underline{M}'^{\mathrm{an}}_K \xrightarrow{\;i'^{\mathrm{an}}\;} G'^{\mathrm{an}}_K \longrightarrow A'^{\mathrm{an}}_K \longrightarrow 0
\end{equation}
analogue à (14.1.7), associée au schéma abélien dual $A'_K$ de $A_K$ (à moins que ce ne soit l'opposé ???).


\SGApage{I}{510}


14.2. \underline{Relations de l'homomorphisme} $i:\underline{M}_K \longrightarrow G_K$ \underline{avec la monodromie.}

Soit $n$ un entier $> 0$. On déduit alors de la suite exacte (14.1.7), via la suite exacte des $\underline{\mathrm{Ext}}^i(\mathbb{Z}/n\mathbb{Z}\ , -)$ sur le site des espaces rigide-analytiques au-dessus de $K$, une suite exacte canonique (compte tenu du fait que $\underline{M}_K^{\mathrm{an}}$ est sans torsion et $G_K^{\mathrm{an}}$ est $n$-divisible) :

\begin{equation*}
0 \longrightarrow {}_nG_K^{\mathrm{an}} \longrightarrow {}_nA_K^{\mathrm{an}} \longrightarrow \underline{M}_{K,n}^{\mathrm{an}} \longrightarrow 0 \qquad ,
\end{equation*}
où $\underline{M}_{K,n} \overset{\mathrm{dfn}}{=} \underline{M}_K/n\underline{M}_K$ . Comme tous les termes qui interviennent dans cette suite exacte proviennent de schémas en groupes finis sur $K$, cette suite exacte peut se réécrire comme provenant d'une suite exacte

\begin{equation}\tag{14.2.1}
0 \longrightarrow {}_nG_K \longrightarrow {}_nA_K \longrightarrow \underline{M}_{K,n} \longrightarrow 0 \qquad .
\end{equation}
Pour $n$ variable, ces suites exactes se comportent comme il convient, et on trouve, par passage à la limite sur les puissances d'un nombre premier $\ell$, une suite exacte analogue

\begin{equation}\tag{14.2.2}
0 \longrightarrow T_\ell(G_K) \longrightarrow T_\ell(A_K) \longrightarrow \underline{M}_K \otimes_{\mathbb{Z}} \mathbb{Z}_\ell \longrightarrow 0 \qquad .
\end{equation}
Or on a déjà une telle suite exacte par voie purement algébrique, grâce à la première formule (8.1.2) et à (8.1.7), qui donnent

\begin{equation}\tag{14.2.3}
T_\ell(G_K) \overset{\sim}{\longrightarrow} T_\ell(A_K)^{\mathrm{ef}} \subset T_\ell(A_K) \ , \ T_\ell(A_K)/T_\ell(A_K)^{\mathrm{ef}} \simeq \underline{M}_K \otimes_{\mathbb{Z}} \mathbb{Z}_\ell \qquad .
\end{equation}
Ceci posé, je dis que \underline{les homomorphismes} (14.2.3) \underline{provenant de la théorie algébrique ne sont autres que ceux déduits de la suite exacte} (14.2.2) \underline{provenant de la théorie transcendante}. Pour le premier homomorphisme, correspondant au premier homomorphisme de (14.2.2), notre assertion résulte aussitôt, par réduction à la situation semi-stable, de la


\SGApage{I}{511}

\SGAvar{commutativité de (14.1.11), compte tenu que les flèches verticales y définissent des isomorphismes sur les $T_\ell$, et que $T_\ell$ appliqué aux flèches horizontales donne respectivement la flèche envisagée dans (14.2.3) et dans (14.2.2). Nous admettrons la compatibilité annoncée, concernant les deux isomorphismes obtenus entre $T_\ell(A_K)/T_\ell(A_K)^{ef}$ avec $\underline{M}_K \otimes_{\mathbb{Z}} \mathbb{Z}_\ell$ ; notons d'ailleurs que \underline{cette compatibilité} (pour un $\ell$ donné) \underline{suffit déjà à caractériser de façon unique}, \underline{par voie transcendante}, \underline{l'homomorphisme} (14.1.5) \underline{intervenant dans} (14.1.7), interprété comme un isomorphisme de $\underline{M}_K^{\mathrm{an}}$ avec $\mathrm{Ker}\, e$. On peut dire aussi que l'homomorphisme $e$ de (14.1.7) permet d'identifier les $T_\ell(A_K)/T_\ell(A_K)^{ef}$ aux $\ell$-adifiés $\underline{R}_K \otimes \mathbb{Z}_\ell$ d'un même réseau $\underline{R}_K = \mathrm{Ker}\, e$ sur $K$, et que cette ``stucture entière'' sur le système des pro-$\ell$-groupes précédents est compatible, dans un sens évident, avec la structure entière provenant des formules (8.1.7) (théorème d'orthogonalité 5.2), ce qui définit alors de façon canonique un isomorphisme

\begin{equation*}
i \;:\; \underline{M}_K \;\simeq\; \mathrm{Ker}\, e \qquad .
\end{equation*}
Comme l'homomorphisme $i$ (14.1.5) permet de réconstituer la structure d'extension habituelle de $T_\ell(A_K)$, qui est (14.2.2), donc aussi sa structure d'\underline{extension panachée} décrite dans 9.6, grâce à la suite exacte

}{commutativité de (14.1.11), compte tenu que les flèches verticales y définissent des isomorphismes sur les $T_\ell$, et que $T_\ell$ appliqué aux flèches horizontales donne respectivement la flèche envisagée dans (14.2.3) et dans (14.2.2). Nous admettrons la compatibilité annoncée, concernant les deux isomorphismes obtenus entre $T_\ell(A_K)/T_\ell(A_K)^{ef}$ avec $\underline{M}_K \otimes_{\mathbb{Z}} \mathbb{Z}_\ell$ ; notons d'ailleurs que \underline{cette compatibilité} (pour un $\ell$ donné) \underline{suffit déjà à caractériser de façon unique}, \underline{par voie transcendante}, \underline{l'homomorphisme} (14.1.5) \underline{intervenant dans} (14.1.7), interprété comme un isomorphisme de $\underline{M}_K^{\mathrm{an}}$ avec $\mathrm{Ker}\, e$. On peut dire aussi que l'homomorphisme $e$ de (14.1.7) permet d'identifier les $T_\ell(A_K)/T_\ell(A_K)^{ef}$ aux $\ell$-adifiés $\underline{R}_K \otimes \mathbb{Z}_\ell$ d'un même réseau $\underline{R}_K = \mathrm{Ker}\, e$ sur $K$, et que cette ``structure entière'' sur le système des pro-$\ell$-groupes précédents est compatible, dans un sens évident, avec la structure entière provenant des formules (8.1.7) (théorème d'orthogonalité 5.2), ce qui définit alors de façon canonique un isomorphisme

\begin{equation*}
i \;:\; \underline{M}_K \;\simeq\; \mathrm{Ker}\, e \qquad .
\end{equation*}
Comme l'homomorphisme $i$ (14.1.5) permet de reconstituer la structure d'extension habituelle de $T_\ell(A_K)$, qui est (14.2.2), donc aussi sa structure d'\underline{extension panachée} décrite dans 9.6, grâce à la suite exacte

}{F7I-IX-192}\begin{equation*}
\begin{array}{ccccccccc}
0 & \longrightarrow & T_\ell(T_K) & \longrightarrow & T_\ell(G_K) & \longrightarrow & T_\ell(B_K) & \longrightarrow & 0 \\[5ex]
 & & \underline{M}'_K \otimes \mathbb{Z}_\ell(1) & & & & & &
\end{array}
\end{equation*}

\SGApage{I}{512}

\SGAvar{déduite de (14.1.1), il permet également dans le cas semi-stable de réconstituer les \underline{accouplements de monodromie} (9.1.2)

}{déduite de (14.1.1), il permet également dans le cas semi-stable de reconstituer les \underline{accouplements de monodromie} (9.1.2)

}{F7I-IX-193}\begin{equation}\tag{14.2.4}
u_\ell \;:\; \underline{M} \otimes \underline{M}' \longrightarrow \mathbb{Z}_\ell \qquad ,
\end{equation}
ceux-ci ayant été définis précisément en termes d'extensions panachées dans 9.6. Contentons-nous d'indiquer le résultat de la traduction de la construction de 9.6. Pour ceci, notons que dans le cas semi-stable, $p_K$ et $p'_K$ proviennent d'homomorphismes

\SGAvar{\begin{equation*}
p \;:\; \underline{M} \longrightarrow B \quad , \quad \underline{M}' \longrightarrow B' \qquad ,
\end{equation*}
où $B,B'$ sont deux schémas abéliens sur $S$ duals l'un de l'autre, d'où par image inverse par $(p,p')$ une biextension $E$ de $(\underline{M},M')$ par $\mathbb{G}_{mS}$ , dont la fibre générique est la biextension $E_K$ déjà envisagée dans 4.1. Ceci dit, supposant pour simplifier $\underline{M},\underline{M}'$ constants de valeurs $M,M'$. Pour tout couple $(m,m') \in M \times M'$, $E_{m,m'}$ est un torseur sous $\mathbb{G}_{mS}$ , donc correspond à un Module inversible $L_{m,m'}$ sur l'anneau $V$ de $S$, qui admet donc une base $u_o$ définie modulo multiplication par une unité de $V$. Par suite on trouve une application canonique de ``valuation''

}{\begin{equation*}
p \;:\; \underline{M} \longrightarrow B \quad , \quad p' \;:\; \underline{M}' \longrightarrow B' \qquad ,
\end{equation*}
où $B,B'$ sont deux schémas abéliens sur $S$ duals l'un de l'autre, d'où par image inverse par $(p,p')$ une biextension $E$ de $(\underline{M},\underline{M}')$ par $\mathbb{G}_{mS}$ , dont la fibre générique est la biextension $E_K$ déjà envisagée dans 14.1. Ceci dit, supposons pour simplifier $\underline{M},\underline{M}'$ constants de valeurs $M,M'$. Pour tout couple $(m,m') \in M \times M'$, $E_{m,m'}$ est un torseur sous $\mathbb{G}_{mS}$ , donc correspond à un Module inversible $L_{m,m'}$ sur l'anneau $V$ de $S$, qui admet donc une base $u_o$ définie modulo multiplication par une unité de $V$. Par suite on trouve une application canonique de ``valuation''

}{F7I-IX-194}\begin{equation*}
\underline{v} \;:\; E_{m,m'}(K) \longrightarrow \mathbb{Z}
\end{equation*}
en associant à toute base $u$ de $L_{m,m'} \otimes_V K$ , donné par $u = \lambda u_o$ ($\lambda \in K^*$) l'élément

\begin{equation*}
\underline{v}(u) \;=\; \underline{v}(\lambda) \in \mathbb{Z} \qquad ,
\end{equation*}
où $\underline{v}$ dans le deuxième membre désigne la valuation associée à $V$ (normalisée de la façon habituelle, de sorte que son groupe de valeurs soit $\mathbb{Z}$ ) ;


\SGApage{I}{513}

l'entier obtenu ne dépend évidemment pas du choix de $u_o$ . Si alors $\varphi$ est une trivialisation de $E$, on peut considérer l'application correspondante

\begin{equation}\tag{14.2.5}
(m,m') \longmapsto \underline{v}(\varphi(m,m')) \; : \; M \times M' \longrightarrow \mathbb{Z} \quad .
\end{equation}
Ceci posé, on trouve que l'\underline{homomorphisme de monodromie} (14.2.4) \underline{est à valeurs dans} $\mathbb{Z}$ , \underline{et n'est autre que l'homomorphisme} (14.2.5) \underline{où} $\varphi$ \underline{est le splittage de} $E_K$ \underline{associé à l'homomorphisme} $i$ (14.1.5).

\underline{Remarques} 14.2.6. On retrouve en particulier (via la théorie transcendante rigide-analytique), le fait que l'accouplement de monodromie est entier et indépendant de $\ell$ (10.4 a)), prouvé par une autre méthode (également transcendante) dans 12.5. D'après M. RAYNAUD, la théorie rigide-analytique permet également de retrouver l'assertion de positivité 10.4 b) (qui est également contenue dans 12.5).

\underline{Remarques} 14.2.7. Considérons $\underline{M}_K \otimes_{\mathbb{Z}} \mathbb{Q}$ comme une extension de $\underline{M}_K \otimes_{\mathbb{Z}} \mathbb{Q}/\mathbb{Z}$ par $\underline{M}_K$ , et utilisons l'homomorphisme (14.1.5) $i : \underline{M}_K \longrightarrow G_K$ , on en déduit une extension $A_K^{\flat}$

\begin{equation}\tag{14.2.8}
0 \longrightarrow G_K \longrightarrow A_K^{\flat} \longrightarrow \underline{M}_K \otimes \mathbb{Q}/\mathbb{Z} \longrightarrow 0 \quad .
\end{equation}
Notons d'ailleurs que la donnée d'une telle extension, grâce au fait que $G_K$ est un groupe divisible, équivaut à la donnée d'une famille d'extensions de la forme (14.2.2) (savoir celles déduites de (14.2.8) par application des $T_\ell$), et $i$ étant un monomorphisme, les extensions en question définissant (14.2.8) ne sont autres que les extensions


\SGApage{I}{514}

(14.2.2) provenant de la suite exacte (14.2.1). Comme ces dernières extensions sont définies par voie purement algébrique, il en est de même de l'extension (14.2.8) ; en fait, la caractérisation qu'on vient d'en donner montre aussitôt que cette extension est canoniquement isomorphe au groupe, noté également $A_K^{\flat}$ , introduit dans (8.2.3), compte tenu de l'isomorphisme 11.11 du Groupe des composantes connexes de ce dernier avec $\underline{M}_K \otimes \mathbb{Q}/\mathbb{Z}$ .


14.3. \underline{Application à une caractérisation arithmétique de} $i: \underline{M}_K \longrightarrow G_K$ .

La caractérisation que nous avons en vue repose sur l'observation que pour tout nombre premier $\ell$, l'extension (14.2.2) de $\underline{M}_K \otimes \mathbb{Z}_\ell$ par $T_\ell(G_K)$ déduite de $i$ est connue par voie purement algébrique, en particulier on connaît l'image de $i$ par l'homomorphisme canonique

\begin{equation}\tag{14.3.1}
\mathrm{Hom}(\underline{M}_K , G_K) \longrightarrow \prod_\ell \mathrm{Ext}^1(\underline{M}_K \otimes \mathbb{Z}_\ell \, , \, T_\ell(G_K)) \quad .
\end{equation}
Dans le cas (auquel on peut toujours se ramener par descente) où $\underline{M}_K$ est constant, de valeur $M$, l'homomorphisme précédent peut s'interpréter comme l'homomorphisme

\SGAvar{\begin{equation}\tag{14.3.2}
\mathrm{Hom}(M, G_K(K)) \longrightarrow \prod_\ell \mathrm{Hom}(M, H^1(K, T_\ell(G_K))
\end{equation}
déduit de l'homomorphisme canonique

}{\begin{equation}\tag{14.3.2}
\mathrm{Hom}(M, G_K(K)) \longrightarrow \prod_\ell \mathrm{Hom}(M, H^1(K, T_\ell(G_K)))
\end{equation}
déduit de l'homomorphisme canonique

}{F7I-IX-195}\begin{equation}\tag{14.3.3}
G_K(K) \longrightarrow \prod_\ell H^1(K, T_\ell(G_K)) \quad ,
\end{equation}
provenant des homomorphismes cobords dans les suites exactes de Kummer

\begin{equation*}
0 \longrightarrow {}_n G_K \longrightarrow G_K \xrightarrow{\; n \;} G_K \longrightarrow 0 \quad .
\end{equation*}

\SGApage{I}{515}

\SGAvar{(L'homomorphisme (14.3.1) pourrait d'ailleurs s'interpréter en tous cas comme l'homomorphisme cobord analogue à (14.3.3), relatif au groupe tordu $\check{\underline{M}}_K \otimes G_K = \underline{\mathrm{Hom}}(\underline{M}_K , G_K)$.) Le noyau de (14.3.3) est manifestement formé des éléments de $G_K(K)$ qui sont \underline{divisibles par} $n$ \underline{pour tout entier} $n \geq 1$, ou, comme nous dirons, qui sont \underline{infiniment divisibles}. Ceci détermine donc déjà $i$ modulo un élément $\mathrm{Hom}(M,g)$, où $g$ est le sous-groupe de $G_K(K)$ formé des éléments infiniment divisibles de $G_K(K)$ :

}{(L'homomorphisme (14.3.1) pourrait d'ailleurs s'interpréter en tous cas comme l'homomorphisme cobord analogue à (14.3.3), relatif au groupe tordu $\check{\underline{M}}_K \otimes G_K = \underline{\mathrm{Hom}}(\underline{M}_K , G_K)$.) Le noyau de (14.3.3) est manifestement formé des éléments de $G_K(K)$ qui sont \underline{divisibles par} $n$ \underline{pour tout entier} $n \geq 1$, ou, comme nous dirons, qui sont \underline{infiniment divisibles}. Ceci détermine donc déjà $i$ modulo un élément de $\mathrm{Hom}(M,g)$, où $g$ est le sous-groupe de $G_K(K)$ formé des éléments infiniment divisibles de $G_K(K)$ :

}{F7I-IX-196}\begin{equation*}
g = \bigcap_{n \geq 1} n\, G_K(K) \qquad .
\end{equation*}
Or on a le résultat suivant :

\underline{Lemme} 14.3.4. \underline{Soient} $V$ \underline{un anneau de valuation discrète dont le corps résiduel} $k$ \underline{est de} car. $p > 0$ \underline{et de type fini sur le corps premier}, $G_K$ \underline{un schéma en groupes de type fini sur} $K$, \underline{de rang unipotent nul si} car $K = 0$. \underline{Alors le groupe} $g$ \underline{des éléments infiniment divisibles de} $G_K(K)$ \underline{est réduit à zéro}, \underline{i.e.} (\underline{si} $G_K$ \underline{est divisible}) \underline{l'homomorphisme} (14.3.3) \underline{est injectif}.

\SGAvar{Quitte à faire une extension finie $K'$ sur $K$ et de remplacer $V$ par son normalisé dans $K'$, on peut supposer que $G_K$ admet une suite de composition dans les quotients successifs sont de l'un des types suivants : groupe fini constant, variété abélienne, $\mathbb{G}_{mK}$ , $\mathbb{G}_{aK}$ . On est donc réduit aussitôt au cas où $\mathbb{G}_K$ lui-même est de l'un des types précédents. Le cas d'un groupe constant fini, ou du groupe $\mathbb{G}_{aK}$ , est trivial. Dans le cas d'un schéma abélien, il peut être considéré comme la fibre générique de son modèle de Néron $G$, de sorte que dans ce cas,


}{Quitte à faire une extension finie $K'$ sur $K$ et de remplacer $V$ par son normalisé dans $K'$, on peut supposer que $G_K$ admet une suite de composition dont les quotients successifs sont de l'un des types suivants : groupe fini constant, variété abélienne, $\mathbb{G}_{mK}$ , $\mathbb{G}_{aK}$ . On est donc réduit aussitôt au cas où $G_K$ lui-même est de l'un des types précédents. Le cas d'un groupe constant fini, ou du groupe $\mathbb{G}_{aK}$ , est trivial. Dans le cas d'un schéma abélien, il peut être considéré comme la fibre générique de son modèle de Néron $G$, de sorte que dans ce cas,


}{F7I-IX-197}\SGApage{I}{516}

\SGAvar{on est ramené à montrer que si $G$ est un schéma de type fini sur $S$, alors le sous-groupe de $G(S)$ formé des points infiniments divisibles est réduit à zéro. Or il en est ainsi pour le groupe $G(k)$, comme il résulte aussitôt du dévissage habituel de $G_o$ et du \underline{théorème de MORDELL-WEIL} [16] (disant que le groupe $G_o(k)$ est un $\mathbb{Z}$ -module de type fini si $G_o$ est un schéma abélien sur le corps de type fini $k$). Donc $g$ est contenu dans le noyau $G(S)^{(1)}$ de $G(S) \longrightarrow G(k)$. Comme le noyau de $G(S)^{(1)} \longrightarrow G(V/\underline{m}^{n+1})$ se dévisse en groupes isomorphes à des vectoriels sur $k$, (en supposant $S$ complet, ce qui est loisible) qui sont donc annulés par $p$, on voit que $g$ est contenu dans le noyau de l'homomorphisme précédent quel que soit $n$, donc qu'il est nul. Le cas où $G_K = \mathbb{G}_{mK}$ se ramène au cas précédent, en observant que dans ce cas $G_K(K) \longrightarrow K^*$ est une extension de $\mathbb{Z}$ (qui n'a pas d'éléments infiniment divisibles sauf zéro) par $V^* = \mathbb{G}_m(S)$, qui satisfait la même condition d'après ce qu'on vient de voir.

}{on est ramené à montrer que si $G$ est un schéma de type fini sur $S$, alors le sous-groupe de $G(S)$ formé des points infiniment divisibles est réduit à zéro. Or il en est ainsi pour le groupe $G(k)$, comme il résulte aussitôt du dévissage habituel de $G_o$ et du \underline{théorème de MORDELL-WEIL} [16] (disant que le groupe $G_o(k)$ est un $\mathbb{Z}$ -module de type fini si $G_o$ est un schéma abélien sur le corps de type fini $k$). Donc $g$ est contenu dans le noyau $G(S)^{(1)}$ de $G(S) \longrightarrow G(k)$. Comme le noyau de $G(S)^{(1)} \longrightarrow G(V/\underline{m}^{n+1})$ se dévisse en groupes isomorphes à des vectoriels sur $k$, (en supposant $S$ complet, ce qui est loisible) qui sont donc annulés par $p$, on voit que $g$ est contenu dans le noyau de l'homomorphisme précédent quel que soit $n$, donc qu'il est nul. Le cas où $G_K = \mathbb{G}_{mK}$ se ramène au cas précédent, en observant que dans ce cas $G_K(K) = K^*$ est une extension de $\mathbb{Z}$ (qui n'a pas d'éléments infiniment divisibles sauf zéro) par $V^* = \mathbb{G}_m(S)$, qui satisfait la même condition d'après ce qu'on vient de voir.

}{F7I-IX-198}\underline{Remarques} 14.3.4.1. On notera qu'il est essentiel dans 14.3.4 de faire intervenir la divisibilité par \underline{tous} les entiers $\geq 1$, et non seulement les puissances d'un nombre premier fixé $\ell$ ; ou encore qu'il faut regarder le noyau de l'homomorphisme (14.3.3), sans se borner au second membre à un seul facteur $H^1(K, T_\ell(G_K))$. En effet, si $G_K$ provient du schéma en groupes lisse $G$ sur $S$ et si $S$ est hensélien, tout élément du noyau de $G(S) \longrightarrow G(k)$ est infiniment $\ell$-divisible pour tout $\ell \neq p$ ! l'argument établissant 14.3.4 montre que le sous-groupe de $G_K(K)$ formé


\SGApage{I}{517}

des éléments infiniment $p$-divisibles (i.e. le noyau de $G_K(K) \to H^1(K,T_p(G_K))$ si $G_K$ est $p$-divisible) est un groupe \underline{fini}.

\SGAvar{14.3.5. Les réflexions qui précèdent nous montrent donc que lorsque le corps résiduel $k$ est un corps de car. $p > 0$, de type fini au sens absolu, alors l'homomorphisme $i$ (14.1.5) est connu entièrement par la connaissance de son image par (14.3.1), i.e. par la connaissance des extensions (14.2.2). Pour en déduire un principe de caractérisation de $i$ dans le cas général où on ne fait pas de restriction sur $k$, notons qu'au par. 7 nous avons pu construire le schéma $G$ sur $S$ sous des conditions nettement plus générales que celles où $S$ est un trait complet et qu'on dispose d'une variété abélienne à réduction semi-stable sur son corps des fractions. Notamment, désignons plus généralement par $S$ le spectre d'un anneau noethérien normal $V$, séparé et complet pour une topologie $\underline{m}$-adique, et considérons un schéma en groupes lisse et quasi-projectif $A$ sur $S$, à fibres connexes. Nous supposons que $A_o = A \times_S S_o$ ($S_o = \mathrm{Spec}(V/\underline{m})$) est une extension d'un schéma abélien $B_o$ par un tore isotrivial $T_o$, et que $A|U$ est un schéma abélien, où $U \subset S - S_o$ est un ouvert non vide. On a alors dans 7.22 construit un schéma $G$ sur $S$, extension d'un schéma abélien $B$ par un tore $T$, $G$ étant caractérisé à isomorphisme unique près par l'isomorphisme (14.1.8)

}{14.3.5. Les réflexions qui précèdent nous montrent donc que lorsque le corps résiduel $k$ est un corps de car. $p > 0$, de type fini au sens absolu, alors l'homomorphisme $i$ (14.1.5) est connu entièrement par la connaissance de son image par (14.3.1), i.e. par la connaissance des extensions (14.2.2). Pour en déduire un principe de caractérisation de $i$ dans le cas général où on ne fait pas de restriction sur $k$, notons qu'au par. 7 nous avons pu construire le schéma $G$ sur $S$ sous des conditions nettement plus générales que celles où $S$ est un trait complet et qu'on dispose d'une variété abélienne à réduction semi-stable sur son corps des fractions. Notamment, désignons plus généralement par $S$ le spectre d'un anneau noethérien normal $V$, séparé et complet pour une topologie $\underline{m}$-adique, et considérons un schéma en groupes lisse et quasi-projectif $A$ sur $S$, à fibres connexes. Nous supposons que $A_o = A \times_S S_o$ ($S_o = \mathrm{Spec}(V/\underline{m})$) est une extension d'un schéma abélien $B_o$ par un tore isotrivial $T_o$, et que $A|U$ est un schéma abélien, où $U \subset S - S_o$ est un ouvert non vide. On a alors dans 7.2.2 construit un schéma $G$ sur $S$, extension d'un schéma abélien $B$ par un tore $T$, $G$ étant caractérisé à isomorphisme unique près par l'isomorphisme (14.1.8)

}{F7I-IX-199}\begin{equation}\tag{14.3.5.1}
\hat{G} \xrightarrow{\ \sim\ } \hat{A}
\end{equation}
sur les complétés formels le long de $G_o$ et de $A_o$ . On en déduit, pour tout nombre premier $\ell$, un isomorphisme (7.3.5)


\SGApage{I}{518}

\begin{equation}\tag{14.3.5.2}
T_\ell(G) \xrightarrow{\ \sim\ } T_\ell(A)^f \quad ,
\end{equation}
d'où par restriction à $U$ une structure d'extensions de pro-$\ell$-groupes de Barsotti-Tate

\begin{equation}\tag{14.3.5.3}
\begin{tikzcd}[row sep=large, column sep=large]
0 \arrow[r] \& T_\ell(G_U) \arrow[r] \arrow[d, dash, "\wr"'] \& T_\ell(A_U) \arrow[r] \& T_\ell(A_U)/T_\ell(A_U)^f \arrow[r] \& 0 \quad . \\
\& T_\ell(A_U)^f
\end{tikzcd}
\end{equation}










\SGAvar{Supposons donné un deuxième groupe $G'$ sur $S$, satisfaisant aux mêmes hypothèses que $A$, et une biextension $W$ de $(A,A')$ par $\mathbb{G}_m$ qui, sur $U$, induise une dualité entre les schémas abéliens $A_U$ et $A'_U$ . (En fait, en vertu de VIII 7.1, la donnée de $W$ équivaut à celle d'une telle dualité entre $A_U$ et $A'_U$). Alors $W$ définit des accouplements

\begin{equation}\tag{14.3.5.4}
T_\ell(A_U) \times T_\ell(A'_U) \longrightarrow \mathbb{Z}_\ell(1) \quad ,
\end{equation}
donnant lieu au théorème d'orthogonalité (7.4.8) : $T_\ell(A_U)^f$ est l'orthogonal de $T_\ell(A'_U)^t = T_\ell(T'_U)$, où $T$ est la partie torique de l'extension de Raynaud $G' = {A'}^{\natural}$ associé à $A'$, - de sorte qu'on trouve encore un isomorphisme

\begin{equation}\tag{14.3.5.5}
T_\ell(A_U)/T_\ell(A_U)^f \simeq \underline{M}_U \otimes \mathbb{Z}_\ell \quad ,
\end{equation}
où $\underline{M} = D(T')$ est le réseau sur $S$ dual du tore $T'$ . Donc (14.3.8) et (14.3.10) nous donnent une structure d'extension

}{Supposons donné un deuxième groupe $A'$ sur $S$, satisfaisant aux mêmes hypothèses que $A$, et une biextension $W$ de $(A,A')$ par $\mathbb{G}_m$ qui, sur $U$, induise une dualité entre les schémas abéliens $A_U$ et $A'_U$ . (En fait, en vertu de VIII 7.1, la donnée de $W$ équivaut à celle d'une telle dualité entre $A_U$ et $A'_U$). Alors $W$ définit des accouplements

\begin{equation}\tag{14.3.5.4}
T_\ell(A_U) \times T_\ell(A'_U) \longrightarrow \mathbb{Z}_\ell(1) \quad ,
\end{equation}
donnant lieu au théorème d'orthogonalité (7.4.8) : $T_\ell(A_U)^f$ est l'orthogonal de $T_\ell(A'_U)^t = T_\ell(T'_U)$, où $T'$ est la partie torique de l'extension de Raynaud $G' = {A'}^{\natural}$ associé à $A'$, - de sorte qu'on trouve encore un isomorphisme

\begin{equation}\tag{14.3.5.5}
T_\ell(A_U)/T_\ell(A_U)^f \simeq \underline{M}_U \otimes \mathbb{Z}_\ell \quad ,
\end{equation}
où $\underline{M} = D(T')$ est le réseau sur $S$ dual du tore $T'$ . Donc (14.3.5.3) et (14.3.5.5) nous donnent une structure d'extension

}{F7I-IX-200}\begin{equation}\tag{14.3.5.6}
0 \longrightarrow T_\ell(G_U) \longrightarrow T_\ell(A_U) \longrightarrow \underline{M}_U \otimes \mathbb{Z}_\ell \longrightarrow 0 \quad .
\end{equation}
Ceci posé, nous pouvons énoncer la


\SGApage{I}{519}

\underline{Conjecture} 14.4. \underline{Les conditions étant celles qu'on vient d'expliciter dans} 14.3.5, \underline{il existe un homomorphisme}

\SGAvar{\begin{equation}\tag{14.4.1}
i : \underline{M}_U \longrightarrow G_U
\end{equation}
\underline{tel que, pour tout nombre premier} $\ell$, \underline{l'extension} (14.3.5.6) \underline{soit isomorphe à l'extension déduite de} $i$ \underline{comme expliqué pour} (14.2.2). \underline{De façon plus précise, il existe une et une seul façon d'associer, à toute situation du type envisagé dans} 14.3.5, \underline{un homomorphisme} (14.4.1) \underline{satisfaisant à la condition précédente, et de telle façon que la formation de} $i$ \underline{soit compatible avec tout changement de base} $S' \longrightarrow S$, \underline{associé à un homomorphisme continu} $V \longrightarrow V'$ \underline{d'anneaux noethériens adiques.}

14.4.2. Il semble que l'existence de $i$, compatible avec les changements de base, résulte de la théorie rigide-analytique de M. RAYNAUD, compte tenu du fait qu'il est en train de développer celle-ci dans un cadre plus général que celui énoncé dans 14.1, en se plaçant comme ici sur des anneaux noethériens adiques généraux. D'ailleurs, pour prouver l'existence et l'unicité énoncée dans 14.4, on est ramené aussitôt au cas où on se borne à des $S$ tels que $S_o$ soit un schéma de type fini sur $\mathbb{Z}$ : en effet, $V$ est limite inductive filtrante sous-anneaux $V_i$ de type fini sur $\mathbb{Z}$ ; qu'on peut supposer normaux grâce à NAGATA (EGA IV 7.8), et pour $i$ grand la situation $(A,A',W)$ provient d'une situation analogue sur un $V_i$, ou encore sur le séparé complété de $V_i$ pour la topologie définie par $m_i = \underline{m} V_i$ , ce qui fournit la réduction annoncée.


}{\begin{equation}\tag{14.4.1}
i : \underline{M}_U \longrightarrow G_U
\end{equation}
\underline{tel que, pour tout nombre premier} $\ell$, \underline{l'extension} (14.3.5.6) \underline{soit isomorphe à l'extension déduite de} $i$ \underline{comme expliqué pour} (14.2.2). \underline{De façon plus précise, il existe une et une seule façon d'associer, à toute situation du type envisagé dans} 14.3.5, \underline{un homomorphisme} (14.4.1) \underline{satisfaisant à la condition précédente, et de telle façon que la formation de} $i$ \underline{soit compatible avec tout changement de base} $S' \longrightarrow S$, \underline{associé à un homomorphisme continu} $V \longrightarrow V'$ \underline{d'anneaux noethériens adiques.}

14.4.2. Il semble que l'existence de $i$, compatible avec les changements de base, résulte de la théorie rigide-analytique de M. RAYNAUD, compte tenu du fait qu'il est en train de développer celle-ci dans un cadre plus général que celui énoncé dans 14.1, en se plaçant comme ici sur des anneaux noethériens adiques généraux. D'ailleurs, pour prouver l'existence et l'unicité énoncée dans 14.4, on est ramené aussitôt au cas où on se borne à des $S$ tels que $S_o$ soit un schéma de type fini sur $\mathbb{Z}$ : en effet, $V$ est limite inductive filtrante de sous-anneaux $V_i$ de type fini sur $\mathbb{Z}$ ; qu'on peut supposer normaux grâce à NAGATA (EGA IV 7.8), et pour $i$ grand la situation $(A,A',W)$ provient d'une situation analogue sur un $V_i$, ou encore sur le séparé complété de $V_i$ pour la topologie définie par $m_i = \underline{m} V_i$ , ce qui fournit la réduction annoncée.


}{F7I-IX-201}\SGApage{I}{520}

\SGAvar{14.4.3. Notons maintenant que $S-S_o$ est un schéma de Jacobson (EGA IV 10.5.8), donc pour voir que deux solutions $i$, $i'$ du problème posé, pour une situation $(A,A',W)$ donnée, sont égales, il suffit de prouver que ces homomorphismes sont égaux en tout point \underline{fermé} $\eta$ de $U$. Or introduisant le normalisé $S'$ de l'adhérence de $\eta$ dans $S$, $S'$ est un trait complet, et on peut appliquer, à l'image inverse sur $S'$ de la situation envisagée sur $S$, le résultat d'unicité de 4.3, compte tenu que le corps résiduel de $S'$ est \underline{fini}. L'unicité étant ainsi établie, elle implique la compatibilité de la formation de $i$ au changement de base, sous réserve d'existence.

}{14.4.3. Notons maintenant que $S-S_o$ est un schéma de Jacobson (EGA IV 10.5.8), donc pour voir que deux solutions $i$, $i'$ du problème posé, pour une situation $(A,A',W)$ donnée, sont égales, il suffit de prouver que ces homomorphismes sont égaux en tout point \underline{fermé} $\eta$ de $U$. Or introduisant le normalisé $S'$ de l'adhérence de $\eta$ dans $S$, $S'$ est un trait complet, et on peut appliquer, à l'image inverse sur $S'$ de la situation envisagée sur $S$, le résultat d'unicité de 14.3.4, compte tenu que le corps résiduel de $S'$ est \underline{fini}. L'unicité étant ainsi établie, elle implique la compatibilité de la formation de $i$ au changement de base, sous réserve d'existence.

}{F7I-IX-202}\underline{Remarques} 14.5. Il resterait à formuler avec la généralité souhaitable, et à établir, une compatibilité entre l'homomorphisme $i$ introduit dans 14.4 par une méthode rigide-analytique et l'homomorphisme analogue (13.5.13) de la théorie analytique complexe ; ou mieux, d'en trouver une généralisation commune !


\underline{BIBLIOGRAPHIE}

[1]\quad M. Artin, Algebraization of formal moduli I (à paraître).

[2]\quad H. Cartan, Familles d'espaces complexes et fondements de la géométrie analytique, Séminaire ENS 1960/61.

[3]\quad P. Deligne, D. Mumford, The irreducibility of the space of curves of given genus. Pub. Math. n°36 (1969).


\SGApage{I}{521}

{[}4] P. Deligne, travail en préparation sur théorie de Hodge (à paraître dans Pub. Math. 1970).

{[}5] M. Demazure, J. Giraud, M. Raynaud, Schémas abéliens, Séminaire Orsay 1967/68, à paraître.

{[}6] A. Grothendieck, Le Groupe de Brauer III, dans Dix Exposés sur la cohomologie des schémas, North Holland Pub. Cie.

\SGAvar{{[}7] A. Grothendieck, Crystals and the De Rham cohomology of schémes, (notesby I. Coates and O. Jussila), in Dix exposés sur la cohomologie des schémas, Noth Holland Pub. Cie.

}{{[}7] A. Grothendieck, Crystals and the De Rham cohomology of schemes, (notes by I. Coates and O. Jussila), in Dix exposés sur la cohomologie des schémas, North Holland Pub. Cie.

}{F7I-IX-203}{[}8] A. Grothendieck, Un théorème sur les homomorphismes de schémas abéliens, Inventions math. 2, p.59-78 (1966).

{[}9] P.A. Griffiths, Report on variation of Hodge structures (à paraître).

{[}10] M. Hakim, Topos annelés et schémas relatifs, à paraître dans Grundlehren der Math, Springer, (1970).

{[}11] J.I. Igusa, Abstract vanishing cycle theory, Proceedings of the Japan Academy, Vol 34 n$^\circ$9 (1958) p.589-593.

\SGAvar{{[}12] R. Kiehl, Der Endlichkeitssatzfür eigentliche Abbildnngen ... Inventions Math. 2, p.191-214 (1967).

}{{[}12] R. Kiehl, Der Endlichkeitssatz für eigentliche Abbildungen ... Inventions Math. 2, p.191-214 (1967).

}{F7I-IX-204}{[}13] R. Kiehl


\SGApage{I}{522}

{[}14] S.L. Kleiman, Algebraic cycles and the Weil conjectures, in Dix exposés sur la cohomologie des schémas, North Holland Pub. Cie.

{[}15] S. Lang, Abelian varieties, Interscience Publishers n$^\circ$ 7 (1959).

{[}16] S. Lang et A. Néron, Rational points of abelian varieties over function fields, Amer. Jour. Math. vol. 81, n$^\circ$1, p.95-118 (1959).

{[}16bis] U.I. Manin, Théorie des groupes formels commutatifs, Uspechi Mat. Naouk, t. XVIII, p.4-90, (1963).

{[}17] D. Mumford, Geometric invariant theory, Ergebnisse der Mathematik ... 34, Springer (1965).

{[}18] D. Mumford, Abstract Theta functions, Advanced Science Seminar in Alg. Geom. Bowdoin College 1967 (notes by H. Pittie).

{[}19] A. Néron, Modèles minimaux des variétés abéliennes sur les corps locaux et globaux, Pub. Math. 21 (1964), p5-128.

\SGAvar{{[}20] M. Raynaud, Caractéristique  d'Euler-Poincaré d'un faisceau et et cohomologie des variétés abéliennes, in Dix exposés sur la cohomologie des schémas, North Holland Pub. Cie.

}{{[}20] M. Raynaud, Caractéristique  d'Euler-Poincaré d'un faisceau et cohomologie des variétés abéliennes, in Dix exposés sur la cohomologie des schémas, North Holland Pub. Cie.

}{F7I-IX-205}{[}21] M. Raynaud, Modèles de Néron, C.R. Acad. Sci. Paris, t.262, p.413-416 (février 1966).

{[}22] M. Raynaud, Spécialisation du foncteur de Picard, C.R. Ac. Sc. t.264, p941-943 et p.1001-1004.

{[}22bis] M. Raynaud, Spécialisation du foncteur de Picard, à paraître aux Pub. Math.


\SGApage{I}{523}

{[23]} M. Raynaud, Faisceaux amples sur les schémas en groupes et les espaces homogènes, thèse, Paris 1968 (à paraître dans Lecture Notes, Springer 1970).

{[24]} M. Raynaud

{[25]} M. Rosenlicht, Some basic theorems on algebraic groups, Amer. Jour. Math., vol. 78, n$^{\circ}$2, p.401-443 (1956).

{[25bis]} J.P. Serre, Quelques propriétés des variétés abéliennes en car. $p$, Amer. J. Math. vol.80, 1958, p.715-739.

\SGAvar{{[26]} J.P. Serre, Corps locaux, Act. Sci. Ind. 1296 (1962) Herman Paris.

}{{[26]} J.P. Serre, Corps locaux, Act. Sci. Ind. 1296 (1962) Hermann Paris.

}{F7I-IX-206}{[27]} J.P. Serre, Sur les corps locaux à corps résiduel algébriquement clos, Bull. Soc. Math. France,vol.89, 1961,p.105-154.

{[28]} J.P. Serre, Abelian $\ell$-adic representations and elliptic curves, Benjamin Inc (1968).

{[29]} J.P. Serre, J. Tate, Good reduction of abelian varieties, Annals of Math. vol 88, n$^{\circ}$3 (1968), p.492-517.

{[30]} J.P. Serre, J. Tate, Elliptic curves and formal groups (notes du Summer Institute on Alg. Geom., Woods Hole 1964).

{[31]} J.P. Serre, Morphismes universels et variété d'Albanese, in Sém. Chevalley 1958/59, Exp.10.

{[32]} J. Tate, WC-groups over $p$-adic fields, Séminaire Bourbaki n$^{\circ}$156,1957.

{[33]} J. Tate, $p$-divisible groups, Proceedings of a conference on local fields, Driebergen (1966), Springer.

\SGAvar{{[34]} J. Tate, Rigid analytic spaces (míméographíé par l'IHES).


}{{[34]} J. Tate, Rigid analytic spaces (miméographié par l'IHES).


}{F7I-IX-207}\SGApage{I}{524}












\SGApage{I}{525}


Lecture Notes in Mathematics


Comprehensive leaflet on request

Vol. 111: K. H. Mayer, Relationen zwischen charakteristischen Zahlen. III, 99 Seiten. 1969. DM 16,--

Vol. 112: Colloquium on Methods of Optimization. Edited by N. N. Moiseev. IV, 293 pages. 1970. DM 18,--

Vol. 113: R. Wille, Kongruenzklassengeometrien. III, 99 Seiten. 1970. DM 16,--

Vol. 114: H. Jacquet and R. P. Langlands, Automorphic Forms on GL (2). VII, 548 pages. 1970. DM 24,--

Vol. 115: K. H. Roggenkamp and V. Huber-Dyson, Lattices over Orders I. XIX, 290 pages. 1970. DM 18,--

\SGAvar{Vol. 116: Sèminaire Pierre Lelong (Analyse) Année 1969. IV, 195 pages. 1970. DM 16,--

}{Vol. 116: Séminaire Pierre Lelong (Analyse) Année 1969. IV, 195 pages. 1970. DM 16,--

}{F7I-IX-208}Vol. 117 : Y. Meyer, Nombres de Pisot, Nombres de Salem et Analyse Harmonique. 63 pages. 1970. DM 16,--

Vol. 118: Proceedings of the 15th Scandinavian Congress, Oslo 1968. Edited by K. E. Aubert and W. Ljunggren. IV, 162 pages. 1970. DM 16,--

Vol. 119: M. Raynaud, Faisceaux amples sur les schémas en groupes et les espaces homogènes. III, 219 pages. 1970. DM 16,--

Vol. 120: D. Siefkes, Büchi's Monadic Second Order Successor Arithmetic. XII, 130 Seiten. 1970. DM 16,--

Vol. 121: H. S. Bear, Lectures on Gleason Parts. III, 47 pages. 1970. DM 16,--

Vol. 122: H. Zieschang, E. Vogt und H.-D. Coldewey, Flächen und ebene diskontinuierliche Gruppen. VIII, 203 Seiten. 1970. DM 16,--

Vol. 123: A. V. Jategaonkar, Left Principal Ideal Rings.VI, 145 pages. 1970. DM 16,--

\SGAvar{Vol. 124: Séminare de Probabilités IV. Edited by P. A. Meyer. IV, 282 pages. 1970. DM 20,--

}{Vol. 124: Séminaire de Probabilités IV. Edited by P. A. Meyer. IV, 282 pages. 1970. DM 20,--

}{F7I-IX-209}Vol. 125: Symposium on Automatic Demonstration. V, 310 pages.1970. DM 20,--

Vol. 126: P. Schapira, Théorie des Hyperfonctions. XI,157 pages.1970. DM 16,--

Vol. 127: I. Stewart, Lie Algebras. IV, 97 pages. 1970. DM 16,--

Vol. 128: M. Takesaki, Tomita's Theory of Modular Hilbert Algebras and its Applications. II, 123 pages. 1970. DM 16,--

Vol. 129: K. H. Hofmann, The Duality of Compact Semigroups and C$^*$- Bigebras. XII, 142 pages. 1970. DM 16,--

Vol. 130: F. Lorenz, Quadratische Formen über Körpern. II, 77 Seiten. 1970. DM 16,--

Vol. 131: A Borel et al., Seminar on Algebraic Groups and Related Finite Groups. VII, 321 pages. 1970. DM 22,--

Vol. 132: Symposium on Optimization. III, 348 pages. 1970. DM 22,--

Vol. 133: F. Topsøe, Topology and Measure. XIV, 79 pages. 1970. DM 16,--

Vol. 134: L. Smith, Lectures on the Eilenberg-Moore Spectral Sequence. VII, 142 pages. 1970. DM 16,--

Vol. 135: W. Stoll, Value Distribution of Holomorphic Maps into Compact Complex Manifolds. II, 267 pages. 1970. DM 18,--

\SGAvar{Vol. 136 : M. Karoubi et al., Séminaire Heidelberg-Saarbrücken-Strasbuorg sur la K-Théorie. IV, 264 pages. 1970. DM 18,--

}{Vol. 136 : M. Karoubi et al., Séminaire Heidelberg-Saarbrücken-Strasbourg sur la K-Théorie. IV, 264 pages. 1970. DM 18,--

}{F7I-IX-210}Vol. 137 : Reports of the Midwest Category Seminar IV. Edited by S. MacLane. III, 139 pages. 1970. DM 16,--

Vol. 138: D. Foata et M. Schützenberger, Théorie Géométrique des Polynômes Eulériens. V, 94 pages. 1970. DM 16,--

Vol. 139: A. Badrikian, Séminaire sur les Fonctions Aléatoires Linéaires et les Mesures Cylindriques. VII, 221 pages. 1970. DM 18,--

Vol. 140: Lectures in Modern Analysis and Applications II. Edited by C. T. Taam. VI, 119 pages. 1970. DM 16,--

Vol. 141: G. Jameson, Ordered Linear Spaces. XV, 194 pages. 1970. DM 16,--

Vol. 142: K. W. Roggenkamp, Lattices over Orders II. V, 388 pages. 1970. DM 22,--

Vol. 143: K. W. Gruenberg, Cohomological Topics in Group Theory. XIV, 275 pages. 1970. DM 20,--

Vol. 144: Seminar on Differential Equations and Dynamical Systems, II. Edited by J. A. Yorke. VIII, 268 pages. 1970. DM 20,--

Vol. 145: E. J. Dubuc, Kan Extensions in Enriched Category Theory. XVI, 173 pages. 1970. DM 16,--

Vol. 146: A. B. Altman and S. Kleiman, Introduction to Grothendieck Duality Theory. II, 192 pages. 1970. DM 18,--

Vol. 147: D. E. Dobbs, Cech Cohomological Dimensions for Commutative Rings. VI, 176 pages. 1970. DM 16,--

Vol. 148: R. Azencott, Espaces de Poisson des Groupes Localement Compacts. IX, 141 pages. 1970. DM 16,--

Vol. 149: R. G. Swan and E. G. Evans, K-Theory of Finite Groups and Orders. IV, 237 pages. 1970. DM 20,--

Vol. 150: Heyer, Dualität lokalkompakter Gruppen. XIII, 372 Seiten. 1970. DM 20,--

Vol. 151: M. Demazure et A. Grothendieck, Schémas en Groupes I. (SGA 3). XV, 562 pages. 1970. DM 24,--

Vol. 152: M. Demazure et A. Grothendieck, Schémas en Groupes II. (SGA 3). IX, 654 pages. 1970. DM 24,--

Vol. 153: M. Demazure et A. Grothendieck, Schémas en Groupes III. (SGA 3). VIII, 529 pages. 1970. DM 24,--

Vol. 154: A. Lascoux et M. Berger, Variétés Kähleriennes Compactes. VII, 83 pages. 1970. DM 16,--

Vol. 155: Several Complex Variables I, Maryland 1970. Edited by J. Horváth. IV, 214 pages. 1970. DM 18,--

Vol. 156: R. Hartshorne, Ample Subvarieties of Algebraic Varieties. XIV, 256 pages. 1970. DM 20,--

Vol. 157: T. tom Dieck, K. H. Kamps und D. Puppe, Homotopietheorie. VI, 265 Seiten. 1970. DM 20,--

Vol. 158: T. G. Ostrom, Finite Translation Planes. IV. 112 pages. 1970. DM 16,--

Vol. 159: R. Ansorge und R. Hass. Konvergenz von Differenzenverfahren für lineare und nichtlineare Anfangswertaufgaben. VIII, 145 Seiten. 1970. DM 16,--

\SGAvar{Vol. 160: L. Sucheston, Constributions to Ergodic Theory and Probability. VII, 277 pages. 1970. DM 20,--

}{Vol. 160: L. Sucheston, Contributions to Ergodic Theory and Probability. VII, 277 pages. 1970. DM 20,--

}{F7I-IX-211}Vol. 161: J. Stasheff, H-Spaces from a Homotopy Point of View. VI, 95 pages. 1970. DM 16,--

Vol. 162: Harish-Chandra and van Dijk, Harmonic Analysis on Reductive p-adic Groups. IV, 125 pages. 1970. DM 16,--

\SGAvar{Vol. 163: P. Deligne, Equations Différentielles à Points Singuliers Reguliers. III, 133 pages. 1970. DM 16,--

Vol. 164: J. P. Ferrier, Seminaire sur les Algebres Complètes. II, 69 pages. 1970. DM 16,--

}{Vol. 163: P. Deligne, Equations Différentielles à Points Singuliers Réguliers. III, 133 pages. 1970. DM 16,--

Vol. 164: J. P. Ferrier, Séminaire sur les Algèbres Complètes. II, 69 pages. 1970. DM 16,--

}{F7I-IX-212}Vol. 165: J. M. Cohen, Stable Homotopy. V, 194 pages. 1970. DM 16,--

Vol. 166: A. J. Silberger, PGL$_2$ over the p-adics: its Representations, Spherical Functions, and Fourier Analysis. VII, 202 pages. 1970. DM 18,--

Vol. 167: Lavrentiev, Romanov and Vasiliev, Multidimensional Inverse Problems for Differential Equations. V, 59 pages. 1970. DM 16,--

Vol. 168: F. P. Peterson, The Steenrod Algebra and its Applications: A conference to Celebrate N. E. Steenrod's Sixtieth Birthday. VII, 317 pages. 1970. DM 22,--

Vol. 169: M. Raynaud, Anneaux Locaux Henséliens. V, 129 pages. 1970. DM 16,--

Vol. 170: Lectures in Modern Analysis and Applications III. Edited by C. T. Taam. VI, 213 pages. 1970. DM 18,--

Vol. 171: Set-Valued Mappings, Selections and Topological Properties of $2^X$. Edited by W. M. Fleischman. X, 110 pages. 1970. DM 16,--

Vol. 172: Y.-T. Siu and G. Trautmann, Gap-Sheaves and Extension of Coherent Analytic Subsheaves. V, 172 pages. 1971. DM 16,--

Vol. 173: J. N. Mordeson and B. Vinograde, Structure of Arbitrary Purely Inseparable Extension Fields. IV, 138 pages. 1970. DM 16,--

Vol. 174: B. Iversen, Linear Determinants with Applications to the Picard Scheme of a Family of Algebraic Curves. VI, 69 pages. 1970. DM 16,--

Vol. 175: M. Brelot, On Topologies and Boundaries in Potential Theory. VI, 176 pages. 1971. DM 18,--

Vol. 176: H. Popp, Fundamentalgruppen algebraischer Mannigfaltigkeiten. IV, 154 Seiten. 1970. DM 16,--

Vol. 177: J. Lambek, Torsion Theories, Additive Semantics and Rings of Quotients. VI, 94 pages. 1971. DM 16,--

Please turn over


\SGApage{I}{526}

\noindent\rule{\linewidth}{0.4pt}

Vol. 178: Th. Bröcker und T. tom Dieck, Kobordismentheorie. XVI, 191 Seiten. 1970. DM 18,--

\SGAvar{Vol. 179: Seminaire Bourbaki -- vol. 1968/69. Exposés 347-363. IV. 295 pages. 1971. DM 22,--

}{Vol. 179: Séminaire Bourbaki -- vol. 1968/69. Exposés 347-363. IV. 295 pages. 1971. DM 22,--

}{F7I-IX-213}Vol. 180: Séminaire Bourbaki -- vol. 1969/70. Exposés 364-381. IV, 310 pages. 1971. DM 22,--

Vol. 181: F. DeMeyer and E. Ingraham, Separable Algebras over Commutative Rings. V, 157 pages. 1971. DM 16.--

Vol. 182: L. D. Baumert. Cyclic Difference Sets. VI, 166 pages. 1971. DM 16,--

Vol. 183: Analytic Theory of Differential Equations. Edited by P. F. Hsieh and A. W. J. Stoddart. VI, 225 pages. 1971. DM 20,--

Vol. 184: Symposium on Several Complex Variables, Park City, Utah, 1970. Edited by R. M. Brooks. V, 234 pages. 1971. DM 20,--

Vol. 185: Several Complex Variables II, Maryland 1970. Edited by J. Horváth. III, 287 pages. 1971. DM 24,--

Vol. 186: Recent Trends in Graph Theory. Edited by M. Capobianco/J. B. Frechen/M. Krolik. VI, 219 pages. 1971. DM 18.--

Vol. 187: H. S. Shapiro, Topics in Approximation Theory. VIII, 275 pages. 1971. DM 22,--

Vol. 188: Symposium on Semantics of Algorithmic Languages. Edited by E. Engeler. VI, 372 pages. 1971. DM 26,--

Vol. 189: A. Weil, Dirichlet Series and Automorphic Forms. V, 164 pages. 1971. DM 16,--

Vol. 190: Martingales. A Report on a Meeting at Oberwolfach, May 17-23, 1970. Edited by H. Dinges. V, 75 pages. 1971. DM 16,--

Vol. 191: Séminaire de Probabilités V. Edited by P. A. Meyer. IV, 372 pages. 1971. DM 26,--

Vol. 192: Proceedings of Liverpool Singularities -- Symposium I. Edited by C. T. C. Wall. V, 319 pages. 1971. DM 24,--

Vol. 193: Symposium on the Theory of Numerical Analysis. Edited by J. Ll. Morris. VI, 152 pages. 1971. DM 16,--

Vol. 194: M. Berger, P. Gauduchon et E. Mazet. Le Spectre d'une Variété Riemannienne. VII, 251 pages. 1971. DM 22,--

Vol. 195: Reports of the Midwest Category Seminar V. Edited by J.W. Gray and S. Mac Lane.III, 255 pages. 1971. DM 22,--

Vol. 196: H-spaces -- Neuchâtel (Suisse)- Août 1970. Edited by F. Sigrist, V, 156 pages. 1971. DM 16,--

Vol. 197: Manifolds -- Amsterdam 1970. Edited by N. H. Kuiper. V, 231 pages. 1971. DM 20,--

Vol. 198: M. Hervé, Analytic and Plurisubharmonic Functions in Finite and Infinite Dimensional Spaces. VI, 90 pages. 1971. DM 16.--

Vol. 199: Ch. J. Mozzochi, On the Pointwise Convergence of Fourier Series. VII, 87 pages. 1971. DM 16,--

Vol. 2 00: U. Neri, Singular Integrals. VII, 272 pages. 1971. DM 22,--

Vol. 201 : J. H. van Lint, Coding Theory. VII, 136 pages. 1971. DM 16,--

Vol. 202: J. Benedetto, Harmonic Analysis on Totally Disconnected Sets. VIII, 261 pages. 1971. DM 22,--

Vol. 203: D. Knutson, Algebraic Spaces. VI, 261 pages. 1971. DM 22,--

Vol. 204: A. Zygmund, Intégrales Singulières. IV, 53 pages. 1971. DM 16,--

Vol. 205: Séminaire Pierre Lelong (Analyse) Année 1970. VI, 243 pages. 1971. DM 20,--

Vol. 206: Symposium on Differential Equations and Dynamical Systems. Edited by D. Chillingworth. XI, 173 pages. 1971. DM 16,--

Vol. 207: L. Bernstein, The Jacobi-Perron Algorithm -- Its Theory and Application. IV, 161 pages. 1971. DM 16,--

Vol. 208: A. Grothendieck and J. P. Murre, The Tame Fundamental Group of a Formal Neighbourhood of a Divisor with Normal Crossings on a Scheme. VIII, 133 pages. 1971. DM 16,--

Vol. 209: Proceedings of Liverpool Singularities \quad Symposium II. Edited by C. T. C. Wall. V, 280 pages. 1971. DM 22,--

Vol. 210: M. Eichler, Projective Varieties and Modular Forms. III, 118 pages. 1971. DM 16,--

Vol. 211: Théorie des Matroïdes. Edité par C. P. Bruter. III, 108 pages. 1971. DM 16,--

Vol. 212: B. Scarpellini, Proof Theory and Intuitionistic Systems. VII, 291 pages. 1971. DM 24,--

Vol. 213: H. Hogbe-Nlend, Théorie des Bornologies et Applications. V, 168 pages. 1971. DM 18,--

Vol. 214: M. Smorodinsky, Ergodic Theory, Entropy. V, 64 pages. 1971. DM 16,--

Vol. 215: P. Antonelli, D. Burghelea and P. J. Kahn, The Concordance-Homotopy Groups of Geometric Automorphism Groups. X, 140 pages. 1971. DM 16,--

Vol. 216: H. Maaß, Siegel's Modular Forms and Dirichlet Series. VII, 328 pages. 1971. DM 20,--

Vol. 217: T. J. Jech, Lectures in Set Theory with Particular Emphasis on the Method of Forcing. V, 137 pages. 1971. DM 16,--

Vol. 218: C. P. Schnorr, Zufälligkeit und Wahrscheinlichkeit. IV, 212.Seiten 1971. DM 20,--

Vol. 219: N. L. Alling and N. Greenleaf, Foundations of the Theory of Klein Surfaces. IX, 117 pages. 1971. DM 16,--

Vol. 220: W. A. Coppel, Disconjugacy. V, 148 pages. 1971. DM 16,--

Vol. 221: P. Gabriel und F. Ulmer, Lokal präsentierbare Kategorien. V, 200 Seiten. 1971. DM 18,--

Vol. 222: C. Meghea, Compactification des Espaces Harmoniques. III, 108 pages. 1971. DM 16,--

Vol. 223: U. Felgner, Models of ZF-Set Theory. VI, 173 pages. 1971. DM 16,--

\SGAvar{Vol. 224: Revètements Etales et Groupe Fondamental. (SGA 1). Dirigé par A. Grothendieck XXII, 447 pages. 1971. DM 30,--

}{Vol. 224: Revêtements Etales et Groupe Fondamental. (SGA 1). Dirigé par A. Grothendieck XXII, 447 pages. 1971. DM 30,--

}{F7I-IX-214}Vol. 225: Théorie des Intersections et Théorème de Riemann-Roch. (SGA 6). Dirigé par P. Berthelot, A. Grothendieck et L. Illusie. XII, 700 pages. 1971. DM 40,--

Vol. 226: Seminar on Potential Theory, II. Edited by H. Bauer. IV, 170 pages. 1971. DM 18,--

Vol. 227: H. L. Montgomery, Topics in Multiplicative Number Theory. IX, 178 pages. 1971. DM 18,--

Vol. 228: Conference on Applications of Numerical Analysis. Edited by J. Ll. Morris. X, 358 pages. 1971. DM 26,--

Vol. 229: J. Väisälä, Lectures on n-Dimensional Quasiconformal Mappings. XIV, 144 pages. 1971. DM 16,--

Vol. 230: L. Waelbroeck, Topological Vector Spaces and Algebras. VII, 158 pages. 1971. DM 16,--

Vol. 231: H. Reiter, L$^1$-Algebras and Segal Algebras. XI, 113 pages. 1971. DM 16,--

Vol. 232: T. H. Ganelius, Tauberian Remainder Theorems. VI, 75 pages. 1971. DM 16,--

Vol. 233: C. P. Tsokos and W. J. Padgett. Random Integral Equations with Applications to Stochastic Systems. VII, 174 pages. 1971. DM 18,--

Vol. 234: A. Andreotti and W. Stoll. Analytic and Algebraic Dependence of Meromorphic Functions. III, 390 pages. 1971. DM 26,--

Vol. 235: Global Differentiable Dynamics. Edited by O. Hájek, A. J. Lohwater, and R. McCann. X, 140 pages. 1971. DM 16,--

Vol. 236 : M. Barr, P. A. Grillet, and D. H. van Osdol. Exact Categories and Categories of Sheaves. VII, 239 pages. 1971, DM 20,--

Vol. 237: B. Stenström. Rings and Modules of Quotients. VII, 136 pages. 1971. DM 16,--

Vol. 238: Der kanonische Modul eines Cohen-Macaulay-Rings. Herausgegeben von Jürgen Herzog und Ernst Kunz. VI, 103 Seiten. 1971. DM 16,--

Vol. 239: L. Illusie, Complexe Cotangent et Déformations I. XV, 355 pages. 1971. DM 26,--

Vol. 240: A. Kerber, Representations of Permutation Groups I. VII, 192 pages. 1971. DM 18,--

Vol. 241: S. Kaneyuki, Homogeneous Bounded Domains and Siegel Domains. V, 89 pages. 1971. DM 16,--

Vol. 242: R. R. Coifman et G. Weiss, Analyse Harmonique Non-Commutative sur Certains Espaces. V, 160 pages. 1971. DM 16,--

Vol. 243: Japan-United States Seminar on Ordinary Differential and Functional Equations. Edited by M. Urabe. VIII, 332 pages. 1971. DM 26,--

Vol. 244: Séminaire Bourbaki -- vol. 1970/71. Exposés 382--399. IV, 356 pages. 1971. DM 26,--

Vol. 245: D. E. Cohen, Groups of Cohomological Dimension One. V, 99 pages. 1972. DM 16,--

\noindent\rule{\linewidth}{0.4pt}


\SGApage{I}{527}

\noindent\rule{\linewidth}{0.4pt}

Vol. 246: Lectures on Rings and Modules. Tulane University Ring and Operator Theory Year, 1970--1971. Volume I. X, 661 pages. 1972. DM 40,--

Vol. 247: Lectures on Operator Algebras. Tulane University Ring and Operator Theory Year, 1970--1971. Volume II. XI, 786 pages. 1972. DM 40,--

Vol. 248: Lectures on the Applications of Sheaves to Ring Theory. Tulane University Ring and Operator Theory Year, 1970--1971. Volume III. VIII, 315 pages. 1971. DM 26,--

Vol. 249: Symposium on Algebraic Topology. Edited by P. J. Hilton. VII, 111 pages. 1971. DM 16,--

Vol. 250: B. Jónsson, Topics in Universal Algebra. VI, 220 pages. 1972. DM 20,--

Vol. 251: The Theory of Arithmetic Functions. Edited by A. A. Gioia and D. L. Goldsmith VI, 287 pages. 1972. DM 24,--

Vol. 252: D. A. Stone, Stratified Polyhedra. IX, 193 pages. 1972. DM 18,--

Vol. 253: V. Komkov, Optimal Control Theory for the Damping of Vibrations of Simple Elastic Systems. V, 240 pages. 1972. DM 20,--

Vol. 254: C. U. Jensen, Les Foncteurs Dérivés de $\varprojlim$ et leurs Applications en Théorie des Modules. V, 103 pages. 1972. DM 16,--

Vol. 255: Conference in Mathematical Logic -- London '70. Edited by W. Hodges. VIII, 351 pages. 1972. DM 26,--

Vol. 256: C. A. Berenstein and M. A. Dostal, Analytically Uniform Spaces and their Applications to Convolution Equations. VII, 130 pages. 1972. DM 16,--

Vol. 257: R. B. Holmes, A Course on Optimization and Best Approximation. VIII, 233 pages. 1972. DM 20,--

Vol. 258: Séminaire de Probabilités VI. Edited by P. A. Meyer. VI, 253 pages. 1972. DM 22,--

Vol. 259: N. Moulis, Structures de Fredholm sur les Variétés Hilbertiennes. V, 123 pages. 1972. DM 16,--

Vol. 260: R. Godement and H. Jacquet, Zeta Functions of Simple Algebras. IX, 188 pages. 1972. DM 18,--

Vol. 261: A. Guichardet, Symmetric Hilbert Spaces and Related Topics. V, 197 pages. 1972. DM 18,--

Vol. 262: H. G. Zimmer, Computational Problems, Methods, and Results in Algebraic Number Theory. V, 103 pages. 1972. DM 16,--

Vol. 263: T. Parthasarathy, Selection Theorems and their Applications. VII, 101 pages. 1972. DM 16,--

Vol. 264: W. Messing, The Crystals Associated to Barsotti-Tate Groups: with Applications to Abelian Schemes. III, 190 pages. 1972. DM 18,--

Vol. 265: N. Saavedra Rivano, Catégories Tannakiennes. II, 418 pages. 1972. DM 26,--

Vol. 266: Conference on Harmonic Analysis. Edited by D. Gulick and R. L. Lipsman. VI, 323 pages. 1972. DM 24,--

Vol. 267: Numerische Lösung nichtlinearer partieller Differential- und Integro-Differentialgleichungen. Herausgegeben von R. Ansorge und W. Törnig, VI, 339 Seiten. 1972. DM 26,--

Vol. 268: C. G. Simader, On Dirichlet's Boundary Value Problem. IV, 238 pages. 1972. DM 20,--

Vol. 269: Théorie des Topos et Cohomologie Etale des Schémas. (SGA 4). Dirigé par M. Artin, A. Grothendieck et J. L. Verdier. XIX, 525 pages. 1972. DM 50,--

Vol. 270: Théorie des Topos et Cohomologie Etale des Schémas. Tome 2. (SGA 4). Dirigé par M. Artin, A. Grothendieck et J. L. Verdier. V, 314 pages. 1972. DM 50,--

Vol. 271: J. P. May, The Geometry of Iterated Loop Spaces. IX, 175 pages. 1972. DM 18,--

Vol. 272: K. R. Parthasarathy and K. Schmidt, Positive Definite Kernels, Continuous Tensor Products, and Central Limit Theorems of Probability Theory. VI, 107 pages. 1972. DM 16,--

Vol. 273: U. Seip, Kompakt erzeugte Vektorräume und Analysis. IX, 119 Seiten. 1972. DM 16,--

\SGAvar{Vol. 274: Toposes, Algebraic Geometry and Logic. Edited by. F. W. Lawvere. VI, 189 pages. 1972. DM 18,--

}{Vol. 274: Toposes, Algebraic Geometry and Logic. Edited by F. W. Lawvere. VI, 189 pages. 1972. DM 18,--

}{F7I-IX-215}Vol. 275: Séminaire Pierre Lelong (Analyse) Année 1970--1971. VI, 181 pages. 1972. DM 18,--

Vol. 276: A. Borel, Représentations de Groupes Localement Compacts. V, 98 pages. 1972. DM 16,--

Vol. 277: Séminaire Banach. Edité par C. Houzel. VII, 229 pages. 1972. DM 20,--

Vol. 278: H. Jacquet, Automorphic Forms on GL(2). Part II. XIII, 142 pages. 1972. DM 16,--

Vol. 279: R. Bott, S. Gitler and I. M. James, Lectures on Algebraic and Differential Topology. V, 174 pages. 1972. DM 18,--

Vol. 280: Conference on the Theory of Ordinary and Partial Differential Equations. Edited by W. N. Everitt and B. D. Sleeman. XV, 367 pages. 1972. DM 26,--

Vol. 281: Coherence in Categories. Edited by S. Mac Lane. VII, 235 pages. 1972. DM 20,--

Vol. 282: W. Klingenberg und P. Flaschel, Riemannsche Hilbertmannigfaltigkeiten. Periodische Geodätische. VII, 211 Seiten. 1972. DM 20,--

Vol. 283: L. Illusie, Complexe Cotangent et Déformations II. VII, 304 pages. 1972. DM 24,--

Vol. 284: P. A. Meyer, Martingales and Stochastic Integrals I. VI, 89 pages. 1972. DM 16,--

Vol. 285: P. de la Harpe, Classical Banach-Lie Algebras and Banach-Lie Groups of Operators in Hilbert Space. III, 160 pages. 1972. DM 16,--

Vol. 286: S. Murakami, On Automorphisms of Siegel Domains. V, 95 pages. 1972. DM 16,--

Vol. 288: Groupes de Monodromie en Géométrie Algébrique. (SGA 7 I). Dirigé par A. Grothendieck. IX, 523 pages. 1972. DM 50,--

\noindent\rule{\linewidth}{0.4pt}


\SGApage{I}{528}

ISBN 3-540-05987-3

ISBN 0-387-05987-3
\selectlanguage{french}

\ifnum\SGAform=1
\clearpage\section*{Apparat critique}\small
Les lectures ci-dessous sont données dans le TeX exact des témoins, afin de préserver les distinctions de notation. Les indices décimaux renvoient aux différences individuelles du passage marqué.

\par\medskip\noindent\hypertarget{app-F7I-IX-001.1}{}\textbf{F7I-IX-001.1}\quad p. 313; misprint\_pass.\par
\noindent Imprimé : {\ttfamily\footnotesize assymptotique\allowbreak{}}\par
\noindent Corrigé : {\ttfamily\footnotesize asymptotique\allowbreak{}}\par
\noindent Source : 7I:536. \nolinkurl{fixes/FIXLOG_7I_CURRENT.jsonl#822}\par

\par\medskip\noindent\hypertarget{app-F7I-IX-002.1}{}\textbf{F7I-IX-002.1}\quad p. 316; misprint\_pass.\par
\noindent Imprimé : {\ttfamily\footnotesize parait\allowbreak{}}\par
\noindent Corrigé : {\ttfamily\footnotesize paraît\allowbreak{}}\par
\noindent Source : 7I:538. \nolinkurl{fixes/FIXLOG_7I_CURRENT.jsonl#824}\par

\par\medskip\noindent\hypertarget{app-F7I-IX-002.2}{}\textbf{F7I-IX-002.2}\quad p. 316; misprint\_pass.\par
\noindent Imprimé : {\ttfamily\footnotesize eû\allowbreak{}}\par
\noindent Corrigé : {\ttfamily\footnotesize eu\allowbreak{}}\par
\noindent Source : 7I:537. \nolinkurl{fixes/FIXLOG_7I_CURRENT.jsonl#823}\par

\par\medskip\noindent\hypertarget{app-F7I-IX-003.1}{}\textbf{F7I-IX-003.1}\quad p. 317; misprint\_pass.\par
\noindent Imprimé : {\ttfamily\footnotesize toutesfois\allowbreak{}}\par
\noindent Corrigé : {\ttfamily\footnotesize toutefois\allowbreak{}}\par
\noindent Source : 7I:540. \nolinkurl{fixes/FIXLOG_7I_CURRENT.jsonl#826}\par

\par\medskip\noindent\hypertarget{app-F7I-IX-003.2}{}\textbf{F7I-IX-003.2}\quad p. 317; misprint\_pass.\par
\noindent Imprimé : {\ttfamily\footnotesize considéré\allowbreak{}}\par
\noindent Corrigé : {\ttfamily\footnotesize considérée\allowbreak{}}\par
\noindent Source : 7I:1115. \nolinkurl{fixes/FIXLOG_7I_CURRENT.jsonl#1033}\par

\par\medskip\noindent\hypertarget{app-F7I-IX-004.1}{}\textbf{F7I-IX-004.1}\quad p. 319; misprint\_pass.\par
\noindent Imprimé : {\ttfamily\footnotesize cancnique\allowbreak{}}\par
\noindent Corrigé : {\ttfamily\footnotesize canonique\allowbreak{}}\par
\noindent Source : 7I:541; 7I:1116. \nolinkurl{fixes/FIXLOG_7I_CURRENT.jsonl#827}; \nolinkurl{fixes/MISPRINTS_7I.json#1116}\par

\par\medskip\noindent\hypertarget{app-F7I-IX-004.2}{}\textbf{F7I-IX-004.2}\quad p. 319; misprint\_pass.\par
\noindent Imprimé : {\ttfamily\footnotesize X\allowbreak{}}\par
\noindent Corrigé : {\ttfamily\footnotesize K\allowbreak{}}\par
\noindent Source : 7I:1117. \nolinkurl{fixes/FIXLOG_7I_CURRENT.jsonl#1034}; \nolinkurl{fixes/MISPRINTS_7I.json#1117}\par

\par\medskip\noindent\hypertarget{app-F7I-IX-005.1}{}\textbf{F7I-IX-005.1}\quad p. 321; misprint\_pass.\par
\noindent Imprimé : {\ttfamily\footnotesize consacré\allowbreak{}}\par
\noindent Corrigé : {\ttfamily\footnotesize consacrée\allowbreak{}}\par
\noindent Source : 7I:542; 7I:541. \nolinkurl{fixes/FIXLOG_7I_CURRENT.jsonl#828}; \nolinkurl{fixes/MISPRINTS_7I.json#541}\par

\par\medskip\noindent\hypertarget{app-F7I-IX-006.1}{}\textbf{F7I-IX-006.1}\quad p. 322; printed misprint.\par
\noindent Imprimé : \emph{(vide)}\par
\noindent Corrigé : {\ttfamily\footnotesize '\allowbreak{}}\par
\noindent Source : 7I:542; 7I:1101. \nolinkurl{fixes/MISPRINTS_7I.json#542}; \nolinkurl{fixes/MISPRINTS_7I.json#1101}\par

\par\medskip\noindent\hypertarget{app-F7I-IX-007.1}{}\textbf{F7I-IX-007.1}\quad p. 323; misprint\_pass.\par
\noindent Imprimé : {\ttfamily\footnotesize formellements\allowbreak{}}\par
\noindent Corrigé : {\ttfamily\footnotesize formellement\allowbreak{}}\par
\noindent Source : 7I:544. \nolinkurl{fixes/FIXLOG_7I_CURRENT.jsonl#829}\par

\par\medskip\noindent\hypertarget{app-F7I-IX-008.1}{}\textbf{F7I-IX-008.1}\quad p. 324; misprint\_pass.\par
\noindent Imprimé : {\ttfamily\footnotesize donné\allowbreak{}}\par
\noindent Corrigé : {\ttfamily\footnotesize donnée\allowbreak{}}\par
\noindent Source : 7I:545; 7I:544; 7I:1103. \nolinkurl{fixes/FIXLOG_7I_CURRENT.jsonl#830}; \nolinkurl{fixes/MISPRINTS_7I.json#544}; \nolinkurl{fixes/MISPRINTS_7I.json#1103}\par

\par\medskip\noindent\hypertarget{app-F7I-IX-009.1}{}\textbf{F7I-IX-009.1}\quad p. 324; designated-witness attribution.\par
\noindent Imprimé : {\ttfamily\footnotesize  \textbackslash{}Phi'\allowbreak{}}\par
\noindent Corrigé : {\ttfamily\footnotesize  F\allowbreak{}}\par
\noindent Source : WITNESS-F7I-IX-009.1. \nolinkurl{fr/sga7i/diplomatic/expose_IX_body.tex}; \nolinkurl{fr/sga7i/current/expose_IX_body.tex}\par

\par\medskip\noindent\hypertarget{app-F7I-IX-009.2}{}\textbf{F7I-IX-009.2}\quad p. 324; designated-witness attribution.\par
\noindent Imprimé : {\ttfamily\footnotesize \textbackslash{}Psi\allowbreak{}}\par
\noindent Corrigé : {\ttfamily\footnotesize F\allowbreak{}}\par
\noindent Source : WITNESS-F7I-IX-009.2. \nolinkurl{fr/sga7i/diplomatic/expose_IX_body.tex}; \nolinkurl{fr/sga7i/current/expose_IX_body.tex}\par

\par\medskip\noindent\hypertarget{app-F7I-IX-009.3}{}\textbf{F7I-IX-009.3}\quad p. 324; designated-witness attribution.\par
\noindent Imprimé : {\ttfamily\footnotesize \textbackslash{}Psi\allowbreak{}}\par
\noindent Corrigé : {\ttfamily\footnotesize \textbackslash{}Phi'\allowbreak{}}\par
\noindent Source : WITNESS-F7I-IX-009.3. \nolinkurl{fr/sga7i/diplomatic/expose_IX_body.tex}; \nolinkurl{fr/sga7i/current/expose_IX_body.tex}\par

\par\medskip\noindent\hypertarget{app-F7I-IX-009.4}{}\textbf{F7I-IX-009.4}\quad p. 324; printed misprint.\par
\noindent Imprimé : {\ttfamily\footnotesize \textbackslash{}Psi\allowbreak{}}\par
\noindent Corrigé : {\ttfamily\footnotesize \textbackslash{}Phi'\allowbreak{}}\par
\noindent Source : 7I:1105. \nolinkurl{fixes/MISPRINTS_7I.json#1105}\par

\par\medskip\noindent\hypertarget{app-F7I-IX-009.5}{}\textbf{F7I-IX-009.5}\quad p. 324; printed misprint.\par
\noindent Imprimé : \emph{(vide)}\par
\noindent Corrigé : \emph{(1 espace(s), 0 retour(s) de ligne dans le témoin TeX)}\par
\noindent Source : 7I:1106. \nolinkurl{fixes/MISPRINTS_7I.json#1106}\par

\par\medskip\noindent\hypertarget{app-F7I-IX-010.1}{}\textbf{F7I-IX-010.1}\quad p. 325; misprint\_pass.\par
\noindent Imprimé : {\ttfamily\footnotesize isomorphe\allowbreak{}}\par
\noindent Corrigé : {\ttfamily\footnotesize isomorphisme\allowbreak{}}\par
\noindent Source : 7I:546; 7I:545; 7I:1107. \nolinkurl{fixes/FIXLOG_7I_CURRENT.jsonl#831}; \nolinkurl{fixes/MISPRINTS_7I.json#545}; \nolinkurl{fixes/MISPRINTS_7I.json#1107}\par

\par\medskip\noindent\hypertarget{app-F7I-IX-011.1}{}\textbf{F7I-IX-011.1}\quad p. 328; misprint\_pass.\par
\noindent Imprimé : {\ttfamily\footnotesize 5\allowbreak{}}\par
\noindent Corrigé : {\ttfamily\footnotesize 6\allowbreak{}}\par
\noindent Source : 7I:547. \nolinkurl{fixes/FIXLOG_7I_CURRENT.jsonl#832}\par

\par\medskip\noindent\hypertarget{app-F7I-IX-012.1}{}\textbf{F7I-IX-012.1}\quad p. 328; misprint\_pass.\par
\noindent Imprimé : {\ttfamily\footnotesize Il\allowbreak{}}\par
\noindent Corrigé : {\ttfamily\footnotesize Ils\allowbreak{}}\par
\noindent Source : 7I:548. \nolinkurl{fixes/FIXLOG_7I_CURRENT.jsonl#833}\par

\par\medskip\noindent\hypertarget{app-F7I-IX-013.1}{}\textbf{F7I-IX-013.1}\quad p. 329; misprint\_pass.\par
\noindent Imprimé : {\ttfamily\footnotesize -\allowbreak{}}\par
\noindent Corrigé : \emph{(1 espace(s), 0 retour(s) de ligne dans le témoin TeX)}\par
\noindent Source : 7I:549. \nolinkurl{fixes/FIXLOG_7I_CURRENT.jsonl#834}\par

\par\medskip\noindent\hypertarget{app-F7I-IX-014.1}{}\textbf{F7I-IX-014.1}\quad p. 329; printed misprint.\par
\noindent Imprimé : \emph{(vide)}\par
\noindent Corrigé : {\ttfamily\footnotesize 2.\allowbreak{}}\par
\noindent Source : 7I:549; 7I:1111. \nolinkurl{fixes/MISPRINTS_7I.json#549}; \nolinkurl{fixes/MISPRINTS_7I.json#1111}\par

\par\medskip\noindent\hypertarget{app-F7I-IX-014.2}{}\textbf{F7I-IX-014.2}\quad p. 329; designated-witness attribution.\par
\noindent Imprimé : {\ttfamily\footnotesize .5\allowbreak{}}\par
\noindent Corrigé : \emph{(vide)}\par
\noindent Source : WITNESS-F7I-IX-014.2. \nolinkurl{fr/sga7i/diplomatic/expose_IX_body.tex}; \nolinkurl{fr/sga7i/current/expose_IX_body.tex}\par

\par\medskip\noindent\hypertarget{app-F7I-IX-015.1}{}\textbf{F7I-IX-015.1}\quad p. 331; misprint\_pass.\par
\noindent Imprimé : {\ttfamily\footnotesize appelera\allowbreak{}}\par
\noindent Corrigé : {\ttfamily\footnotesize appellera\allowbreak{}}\par
\noindent Source : 7I:551. \nolinkurl{fixes/FIXLOG_7I_CURRENT.jsonl#835}\par

\par\medskip\noindent\hypertarget{app-F7I-IX-016.1}{}\textbf{F7I-IX-016.1}\quad p. 333; printed misprint.\par
\noindent Imprimé : {\ttfamily\footnotesize 6\allowbreak{}}\par
\noindent Corrigé : {\ttfamily\footnotesize 2\allowbreak{}}\par
\noindent Source : 7I:551; 7I:1119. \nolinkurl{fixes/MISPRINTS_7I.json#551}; \nolinkurl{fixes/MISPRINTS_7I.json#1119}\par

\par\medskip\noindent\hypertarget{app-F7I-IX-017.1}{}\textbf{F7I-IX-017.1}\quad p. 334; misprint\_pass.\par
\noindent Imprimé : {\ttfamily\footnotesize isogénes\allowbreak{}}\par
\noindent Corrigé : {\ttfamily\footnotesize isogènes\allowbreak{}}\par
\noindent Source : 7I:553. \nolinkurl{fixes/FIXLOG_7I_CURRENT.jsonl#836}\par

\par\medskip\noindent\hypertarget{app-F7I-IX-018.1}{}\textbf{F7I-IX-018.1}\quad p. 335; printed misprint.\par
\noindent Imprimé : \emph{(vide)}\par
\noindent Corrigé : {\ttfamily\footnotesize )\allowbreak{}}\par
\noindent Source : 7I:1121. \nolinkurl{fixes/MISPRINTS_7I.json#1121}\par

\par\medskip\noindent\hypertarget{app-F7I-IX-019.1}{}\textbf{F7I-IX-019.1}\quad p. 337; misprint\_pass.\par
\noindent Imprimé : {\ttfamily\footnotesize appelée\allowbreak{}}\par
\noindent Corrigé : {\ttfamily\footnotesize appelé\allowbreak{}}\par
\noindent Source : 7I:555; 7I:1131. \nolinkurl{fixes/FIXLOG_7I_CURRENT.jsonl#838}; \nolinkurl{fixes/MISPRINTS_7I.json#1131}\par

\par\medskip\noindent\hypertarget{app-F7I-IX-020.1}{}\textbf{F7I-IX-020.1}\quad p. 337; misprint\_pass.\par
\noindent Imprimé : {\ttfamily\footnotesize (\allowbreak{}}\par
\noindent Corrigé : \emph{(vide)}\par
\noindent Source : 7I:556. \nolinkurl{fixes/FIXLOG_7I_CURRENT.jsonl#839}\par

\par\medskip\noindent\hypertarget{app-F7I-IX-020.2}{}\textbf{F7I-IX-020.2}\quad p. 337; printed misprint.\par
\noindent Imprimé : {\ttfamily\footnotesize 1\allowbreak{}}\par
\noindent Corrigé : {\ttfamily\footnotesize 2\allowbreak{}}\par
\noindent Source : 7I:1133. \nolinkurl{fixes/MISPRINTS_7I.json#1133}\par

\par\medskip\noindent\hypertarget{app-F7I-IX-020.3}{}\textbf{F7I-IX-020.3}\quad p. 337; printed misprint.\par
\noindent Imprimé : {\ttfamily\footnotesize 6\allowbreak{}}\par
\noindent Corrigé : {\ttfamily\footnotesize 2\allowbreak{}}\par
\noindent Source : 7I:1133. \nolinkurl{fixes/MISPRINTS_7I.json#1133}\par

\par\medskip\noindent\hypertarget{app-F7I-IX-021.1}{}\textbf{F7I-IX-021.1}\quad p. 338; misprint\_pass.\par
\noindent Imprimé : {\ttfamily\footnotesize gra\textbackslash{}\textasciicircum{}ce\allowbreak{}}\par
\noindent Corrigé : {\ttfamily\footnotesize grâce\allowbreak{}}\par
\noindent Source : 7I:557. \nolinkurl{fixes/FIXLOG_7I_CURRENT.jsonl#840}\par

\par\medskip\noindent\hypertarget{app-F7I-IX-021.2}{}\textbf{F7I-IX-021.2}\quad p. 338; misprint\_pass.\par
\noindent Imprimé : {\ttfamily\footnotesize Thoérème\allowbreak{}}\par
\noindent Corrigé : {\ttfamily\footnotesize Théorème\allowbreak{}}\par
\noindent Source : 7I:558; 7I:557; 7I:1135. \nolinkurl{fixes/FIXLOG_7I_CURRENT.jsonl#841}; \nolinkurl{fixes/MISPRINTS_7I.json#557}; \nolinkurl{fixes/MISPRINTS_7I.json#1135}\par

\par\medskip\noindent\hypertarget{app-F7I-IX-021.3}{}\textbf{F7I-IX-021.3}\quad p. 338; printed misprint.\par
\noindent Imprimé : {\ttfamily\footnotesize ;\allowbreak{}}\par
\noindent Corrigé : {\ttfamily\footnotesize ,\allowbreak{}}\par
\noindent Source : 7I:1136. \nolinkurl{fixes/MISPRINTS_7I.json#1136}\par

\par\medskip\noindent\hypertarget{app-F7I-IX-021.4}{}\textbf{F7I-IX-021.4}\quad p. 338; misprint\_in\_print.\par
\noindent Imprimé : {\ttfamily\footnotesize w\allowbreak{}}\par
\noindent Corrigé : {\ttfamily\footnotesize W\allowbreak{}}\par
\noindent Source : FIX-7I-D-0559. \nolinkurl{fixes/FIXLOG_7I.jsonl#559}\par

\par\medskip\noindent\hypertarget{app-F7I-IX-022.1}{}\textbf{F7I-IX-022.1}\quad p. 339; designated-witness attribution.\par
\noindent Imprimé : \emph{(vide)}\par
\noindent Corrigé : {\ttfamily\footnotesize \textbackslash{}prime\allowbreak{}}\par
\noindent Source : WITNESS-F7I-IX-022.1. \nolinkurl{fr/sga7i/diplomatic/expose_IX_body.tex}; \nolinkurl{fr/sga7i/current/expose_IX_body.tex}\par

\par\medskip\noindent\hypertarget{app-F7I-IX-023.1}{}\textbf{F7I-IX-023.1}\quad p. 340; misprint\_pass.\par
\noindent Imprimé : {\ttfamily\footnotesize égale\allowbreak{}}\par
\noindent Corrigé : {\ttfamily\footnotesize égal\allowbreak{}}\par
\noindent Source : 7I:563. \nolinkurl{fixes/FIXLOG_7I_CURRENT.jsonl#843}\par

\par\medskip\noindent\hypertarget{app-F7I-IX-023.2}{}\textbf{F7I-IX-023.2}\quad p. 340; designated-witness attribution.\par
\noindent Imprimé : {\ttfamily\footnotesize 1\allowbreak{}}\par
\noindent Corrigé : {\ttfamily\footnotesize 2\allowbreak{}}\par
\noindent Source : WITNESS-F7I-IX-023.2. \nolinkurl{fr/sga7i/diplomatic/expose_IX_body.tex}; \nolinkurl{fr/sga7i/current/expose_IX_body.tex}\par

\par\medskip\noindent\hypertarget{app-F7I-IX-023.3}{}\textbf{F7I-IX-023.3}\quad p. 340; designated-witness attribution.\par
\noindent Imprimé : {\ttfamily\footnotesize 10\allowbreak{}}\par
\noindent Corrigé : {\ttfamily\footnotesize 4.7\allowbreak{}}\par
\noindent Source : WITNESS-F7I-IX-023.3. \nolinkurl{fr/sga7i/diplomatic/expose_IX_body.tex}; \nolinkurl{fr/sga7i/current/expose_IX_body.tex}\par

\par\medskip\noindent\hypertarget{app-F7I-IX-024.1}{}\textbf{F7I-IX-024.1}\quad p. 341; misprint\_pass.\par
\noindent Imprimé : {\ttfamily\footnotesize du\allowbreak{}}\par
\noindent Corrigé : {\ttfamily\footnotesize des\allowbreak{}}\par
\noindent Source : 7I:564. \nolinkurl{fixes/FIXLOG_7I_CURRENT.jsonl#844}\par

\par\medskip\noindent\hypertarget{app-F7I-IX-024.2}{}\textbf{F7I-IX-024.2}\quad p. 341; printed misprint.\par
\noindent Imprimé : \emph{(vide)}\par
\noindent Corrigé : {\ttfamily\footnotesize \textasciicircum{}\{\textbackslash{}perp\}\allowbreak{}}\par
\noindent Source : 7I:564; 7I:1123. \nolinkurl{fixes/MISPRINTS_7I.json#564}; \nolinkurl{fixes/MISPRINTS_7I.json#1123}\par

\par\medskip\noindent\hypertarget{app-F7I-IX-024.3}{}\textbf{F7I-IX-024.3}\quad p. 341; printed misprint.\par
\noindent Imprimé : \emph{(vide)}\par
\noindent Corrigé : {\ttfamily\footnotesize \textasciicircum{}\{\textbackslash{}perp\}\allowbreak{}}\par
\noindent Source : 7I:564; 7I:1123. \nolinkurl{fixes/MISPRINTS_7I.json#564}; \nolinkurl{fixes/MISPRINTS_7I.json#1123}\par

\par\medskip\noindent\hypertarget{app-F7I-IX-024.4}{}\textbf{F7I-IX-024.4}\quad p. 341; misprint\_pass.\par
\noindent Imprimé : {\ttfamily\footnotesize n\allowbreak{}}\par
\noindent Corrigé : {\ttfamily\footnotesize 2n\allowbreak{}}\par
\noindent Source : 7I:1124. \nolinkurl{fixes/FIXLOG_7I_CURRENT.jsonl#1035}\par

\par\medskip\noindent\hypertarget{app-F7I-IX-025.1}{}\textbf{F7I-IX-025.1}\quad p. 343; misprint\_pass.\par
\noindent Imprimé : {\ttfamily\footnotesize S\allowbreak{}}\par
\noindent Corrigé : {\ttfamily\footnotesize s\allowbreak{}}\par
\noindent Source : 7I:566; 7I:565; 7I:1125. \nolinkurl{fixes/FIXLOG_7I_CURRENT.jsonl#846}; \nolinkurl{fixes/MISPRINTS_7I.json#565}; \nolinkurl{fixes/MISPRINTS_7I.json#1125}\par

\par\medskip\noindent\hypertarget{app-F7I-IX-026.1}{}\textbf{F7I-IX-026.1}\quad p. 344; misprint\_pass.\par
\noindent Imprimé : {\ttfamily\footnotesize retrocompact\allowbreak{}}\par
\noindent Corrigé : {\ttfamily\footnotesize rétrocompact\allowbreak{}}\par
\noindent Source : 7I:567. \nolinkurl{fixes/FIXLOG_7I_CURRENT.jsonl#847}\par

\par\medskip\noindent\hypertarget{app-F7I-IX-027.1}{}\textbf{F7I-IX-027.1}\quad p. 345; misprint\_pass.\par
\noindent Imprimé : {\ttfamily\footnotesize élévant\allowbreak{}}\par
\noindent Corrigé : {\ttfamily\footnotesize élevant\allowbreak{}}\par
\noindent Source : 7I:568. \nolinkurl{fixes/FIXLOG_7I_CURRENT.jsonl#848}\par

\par\medskip\noindent\hypertarget{app-F7I-IX-027.2}{}\textbf{F7I-IX-027.2}\quad p. 345; misprint\_pass.\par
\noindent Imprimé : {\ttfamily\footnotesize s\allowbreak{}}\par
\noindent Corrigé : {\ttfamily\footnotesize \textbackslash{}mathrm\{red\}\allowbreak{}}\par
\noindent Source : 7I:569; 7I:568; 7I:1128. \nolinkurl{fixes/FIXLOG_7I_CURRENT.jsonl#849}; \nolinkurl{fixes/MISPRINTS_7I.json#568}; \nolinkurl{fixes/MISPRINTS_7I.json#1128}\par

\par\medskip\noindent\hypertarget{app-F7I-IX-028.1}{}\textbf{F7I-IX-028.1}\quad p. 345; misprint\_pass.\par
\noindent Imprimé : {\ttfamily\footnotesize premier\allowbreak{}}\par
\noindent Corrigé : {\ttfamily\footnotesize première\allowbreak{}}\par
\noindent Source : 7I:570; 7I:569; 7I:1129. \nolinkurl{fixes/FIXLOG_7I_CURRENT.jsonl#850}; \nolinkurl{fixes/MISPRINTS_7I.json#569}; \nolinkurl{fixes/MISPRINTS_7I.json#1129}\par

\par\medskip\noindent\hypertarget{app-F7I-IX-029.1}{}\textbf{F7I-IX-029.1}\quad p. 346; printed misprint.\par
\noindent Imprimé : {\ttfamily\footnotesize T\allowbreak{}}\par
\noindent Corrigé : {\ttfamily\footnotesize S\allowbreak{}}\par
\noindent Source : 7I:1140. \nolinkurl{fixes/MISPRINTS_7I.json#1140}\par

\par\medskip\noindent\hypertarget{app-F7I-IX-030.1}{}\textbf{F7I-IX-030.1}\quad p. 347; misprint\_pass.\par
\noindent Imprimé : \emph{(vide)}\par
\noindent Corrigé : {\ttfamily\footnotesize 3\allowbreak{} }\par
\noindent Source : 7I:572; 7I:571. \nolinkurl{fixes/FIXLOG_7I_CURRENT.jsonl#851}; \nolinkurl{fixes/MISPRINTS_7I.json#571}\par

\par\medskip\noindent\hypertarget{app-F7I-IX-030.2}{}\textbf{F7I-IX-030.2}\quad p. 347; printed misprint.\par
\noindent Imprimé : {\ttfamily\footnotesize VII\allowbreak{}}\par
\noindent Corrigé : {\ttfamily\footnotesize XVII\allowbreak{}}\par
\noindent Source : 7I:1142. \nolinkurl{fixes/MISPRINTS_7I.json#1142}\par

\par\medskip\noindent\hypertarget{app-F7I-IX-030.3}{}\textbf{F7I-IX-030.3}\quad p. 347; printed misprint.\par
\noindent Imprimé : \emph{(vide)}\par
\noindent Corrigé : {\ttfamily\footnotesize )\allowbreak{}}\par
\noindent Source : 7I:1142. \nolinkurl{fixes/MISPRINTS_7I.json#1142}\par

\par\medskip\noindent\hypertarget{app-F7I-IX-031.1}{}\textbf{F7I-IX-031.1}\quad p. 349; misprint\_pass.\par
\noindent Imprimé : {\ttfamily\footnotesize \$S\$\allowbreak{} }\par
\noindent Corrigé : \emph{(vide)}\par
\noindent Source : 7I:574; 7I:573; 7I:1143. \nolinkurl{fixes/FIXLOG_7I_CURRENT.jsonl#853}; \nolinkurl{fixes/MISPRINTS_7I.json#573}; \nolinkurl{fixes/MISPRINTS_7I.json#1143}\par

\par\medskip\noindent\hypertarget{app-F7I-IX-032.1}{}\textbf{F7I-IX-032.1}\quad p. 350; designated-witness attribution.\par
\noindent Imprimé : \emph{(vide)}\par
\noindent Corrigé : {\ttfamily\footnotesize \textasciicircum{}\{\textbackslash{}perp\}\allowbreak{}}\par
\noindent Source : WITNESS-F7I-IX-032.1. \nolinkurl{fr/sga7i/diplomatic/expose_IX_body.tex}; \nolinkurl{fr/sga7i/current/expose_IX_body.tex}\par

\par\medskip\noindent\hypertarget{app-F7I-IX-032.2}{}\textbf{F7I-IX-032.2}\quad p. 350; designated-witness attribution.\par
\noindent Imprimé : {\ttfamily\footnotesize \textasciicircum{}\{\textbackslash{}perp\}\allowbreak{}}\par
\noindent Corrigé : \emph{(vide)}\par
\noindent Source : WITNESS-F7I-IX-032.2. \nolinkurl{fr/sga7i/diplomatic/expose_IX_body.tex}; \nolinkurl{fr/sga7i/current/expose_IX_body.tex}\par

\par\medskip\noindent\hypertarget{app-F7I-IX-033.1}{}\textbf{F7I-IX-033.1}\quad p. 350; printed misprint.\par
\noindent Imprimé : \emph{(vide)}\par
\noindent Corrigé : {\ttfamily\footnotesize \textasciicircum{}\{\textbackslash{}perp\}\allowbreak{}}\par
\noindent Source : 7I:1145. \nolinkurl{fixes/MISPRINTS_7I.json#1145}\par

\par\medskip\noindent\hypertarget{app-F7I-IX-033.2}{}\textbf{F7I-IX-033.2}\quad p. 350; designated-witness attribution.\par
\noindent Imprimé : \emph{(vide)}\par
\noindent Corrigé : {\ttfamily\footnotesize /U\allowbreak{}}\par
\noindent Source : WITNESS-F7I-IX-033.2. \nolinkurl{fr/sga7i/diplomatic/expose_IX_body.tex}; \nolinkurl{fr/sga7i/current/expose_IX_body.tex}\par

\par\medskip\noindent\hypertarget{app-F7I-IX-033.3}{}\textbf{F7I-IX-033.3}\quad p. 350; designated-witness attribution.\par
\noindent Imprimé : \emph{(1 espace(s), 0 retour(s) de ligne dans le témoin TeX)}\par
\noindent Corrigé : {\ttfamily\footnotesize U/\allowbreak{}}\par
\noindent Source : WITNESS-F7I-IX-033.3. \nolinkurl{fr/sga7i/diplomatic/expose_IX_body.tex}; \nolinkurl{fr/sga7i/current/expose_IX_body.tex}\par

\par\medskip\noindent\hypertarget{app-F7I-IX-034.1}{}\textbf{F7I-IX-034.1}\quad p. 351; inherited corrected-layer difference.\par
\noindent Imprimé : {\ttfamily\footnotesize  \$(I\allowbreak{} \textbackslash{}qquad\allowbreak{} )\$\allowbreak{}}\par
\noindent Corrigé : \emph{(vide)}\par
\noindent Source : ADOPT-F7I-IX-034.1. \nolinkurl{fr/sga7i/diplomatic/expose_IX_body.tex}; \nolinkurl{fr/sga7i/current/expose_IX_body.tex}\par

\par\medskip\noindent\hypertarget{app-F7I-IX-035.1}{}\textbf{F7I-IX-035.1}\quad p. 352; misprint\_pass.\par
\noindent Imprimé : {\ttfamily\footnotesize point\allowbreak{}}\par
\noindent Corrigé : {\ttfamily\footnotesize points\allowbreak{}}\par
\noindent Source : 7I:578. \nolinkurl{fixes/FIXLOG_7I_CURRENT.jsonl#855}\par

\par\medskip\noindent\hypertarget{app-F7I-IX-035.2}{}\textbf{F7I-IX-035.2}\quad p. 352; misprint\_pass.\par
\noindent Imprimé : {\ttfamily\footnotesize de\allowbreak{}}\par
\noindent Corrigé : {\ttfamily\footnotesize le\allowbreak{}}\par
\noindent Source : 7I:579; 7I:578. \nolinkurl{fixes/FIXLOG_7I_CURRENT.jsonl#856}; \nolinkurl{fixes/MISPRINTS_7I.json#578}\par

\par\medskip\noindent\hypertarget{app-F7I-IX-035.3}{}\textbf{F7I-IX-035.3}\quad p. 352; misprint\_pass.\par
\noindent Imprimé : \emph{(vide)}\par
\noindent Corrigé : {\ttfamily\footnotesize  de\allowbreak{}}\par
\noindent Source : 7I:581. \nolinkurl{fixes/FIXLOG_7I_CURRENT.jsonl#857}\par

\par\medskip\noindent\hypertarget{app-F7I-IX-036.1}{}\textbf{F7I-IX-036.1}\quad p. 353; printed misprint.\par
\noindent Imprimé : {\ttfamily\footnotesize v\allowbreak{}}\par
\noindent Corrigé : {\ttfamily\footnotesize vi\allowbreak{}}\par
\noindent Source : 7I:581. \nolinkurl{fixes/MISPRINTS_7I.json#581}\par

\par\medskip\noindent\hypertarget{app-F7I-IX-036.2}{}\textbf{F7I-IX-036.2}\quad p. 353; designated-witness attribution.\par
\noindent Imprimé : {\ttfamily\footnotesize vi\allowbreak{}}\par
\noindent Corrigé : {\ttfamily\footnotesize iv\allowbreak{}}\par
\noindent Source : WITNESS-F7I-IX-036.2. \nolinkurl{fr/sga7i/diplomatic/expose_IX_body.tex}; \nolinkurl{fr/sga7i/current/expose_IX_body.tex}\par

\par\medskip\noindent\hypertarget{app-F7I-IX-037.1}{}\textbf{F7I-IX-037.1}\quad p. 354; printed misprint.\par
\noindent Imprimé : {\ttfamily\footnotesize \textbackslash{}neq\allowbreak{}}\par
\noindent Corrigé : {\ttfamily\footnotesize =\allowbreak{}}\par
\noindent Source : 7I:582. \nolinkurl{fixes/MISPRINTS_7I.json#582}\par

\par\medskip\noindent\hypertarget{app-F7I-IX-038.1}{}\textbf{F7I-IX-038.1}\quad p. 355; misprint\_pass.\par
\noindent Imprimé : {\ttfamily\footnotesize '\allowbreak{}}\par
\noindent Corrigé : \emph{(vide)}\par
\noindent Source : 7I:584; 7I:583. \nolinkurl{fixes/FIXLOG_7I_CURRENT.jsonl#858}; \nolinkurl{fixes/MISPRINTS_7I.json#583}\par

\par\medskip\noindent\hypertarget{app-F7I-IX-039.1}{}\textbf{F7I-IX-039.1}\quad p. 356; printed misprint.\par
\noindent Imprimé : {\ttfamily\footnotesize 4\allowbreak{}}\par
\noindent Corrigé : {\ttfamily\footnotesize 3\allowbreak{}}\par
\noindent Source : 7I:584. \nolinkurl{fixes/MISPRINTS_7I.json#584}\par

\par\medskip\noindent\hypertarget{app-F7I-IX-040.1}{}\textbf{F7I-IX-040.1}\quad p. 357; misprint\_pass.\par
\noindent Imprimé : {\ttfamily\footnotesize déduit\allowbreak{}}\par
\noindent Corrigé : {\ttfamily\footnotesize déduite\allowbreak{}}\par
\noindent Source : 7I:586. \nolinkurl{fixes/FIXLOG_7I_CURRENT.jsonl#859}\par

\par\medskip\noindent\hypertarget{app-F7I-IX-040.2}{}\textbf{F7I-IX-040.2}\quad p. 357; designated-witness attribution.\par
\noindent Imprimé : \emph{(0 espace(s), 1 retour(s) de ligne dans le témoin TeX)}\par
\noindent Corrigé : \emph{(0 espace(s), 2 retour(s) de ligne dans le témoin TeX)}\par
\noindent Source : WITNESS-F7I-IX-040.2. \nolinkurl{fr/sga7i/diplomatic/expose_IX_body.tex}; \nolinkurl{fr/sga7i/current/expose_IX_body.tex}\par

\par\medskip\noindent\hypertarget{app-F7I-IX-040.3}{}\textbf{F7I-IX-040.3}\quad p. 357; inherited corrected-layer difference.\par
\noindent Imprimé : \emph{(vide)}\par
\noindent Corrigé : {\ttfamily\footnotesize Exposé\allowbreak{} }\par
\noindent Source : ADOPT-F7I-IX-040.3. \nolinkurl{fr/sga7i/diplomatic/expose_IX_body.tex}; \nolinkurl{fr/sga7i/current/expose_IX_body.tex}\par

\par\medskip\noindent\hypertarget{app-F7I-IX-040.4}{}\textbf{F7I-IX-040.4}\quad p. 357; inherited corrected-layer difference.\par
\noindent Imprimé : {\ttfamily\footnotesize \textbackslash{}hspace\{2.5em\}\allowbreak{}}\par
\noindent Corrigé : {\ttfamily\footnotesize [référence\allowbreak{} interne\allowbreak{} omise\allowbreak{} dans\allowbreak{} la\allowbreak{} source]\allowbreak{}}\par
\noindent Source : ADOPT-F7I-IX-040.4. \nolinkurl{fr/sga7i/diplomatic/expose_IX_body.tex}; \nolinkurl{fr/sga7i/current/expose_IX_body.tex}\par

\par\medskip\noindent\hypertarget{app-F7I-IX-041.1}{}\textbf{F7I-IX-041.1}\quad p. 358; inherited corrected-layer difference.\par
\noindent Imprimé : {\ttfamily\footnotesize \textasciicircum{}\{o\}\allowbreak{}}\par
\noindent Corrigé : \emph{(vide)}\par
\noindent Source : ADOPT-F7I-IX-041.1. \nolinkurl{fr/sga7i/diplomatic/expose_IX_body.tex}; \nolinkurl{fr/sga7i/current/expose_IX_body.tex}\par

\par\medskip\noindent\hypertarget{app-F7I-IX-042.1}{}\textbf{F7I-IX-042.1}\quad p. 359; misprint\_pass.\par
\noindent Imprimé : {\ttfamily\footnotesize coefficients\allowbreak{}}\par
\noindent Corrigé : {\ttfamily\footnotesize coefficient\allowbreak{}}\par
\noindent Source : 7I:589. \nolinkurl{fixes/FIXLOG_7I_CURRENT.jsonl#860}\par

\par\medskip\noindent\hypertarget{app-F7I-IX-043.1}{}\textbf{F7I-IX-043.1}\quad p. 361; misprint\_pass.\par
\noindent Imprimé : {\ttfamily\footnotesize de\allowbreak{}}\par
\noindent Corrigé : {\ttfamily\footnotesize se\allowbreak{}}\par
\noindent Source : 7I:597; 7I:596. \nolinkurl{fixes/FIXLOG_7I_CURRENT.jsonl#861}; \nolinkurl{fixes/MISPRINTS_7I.json#596}\par

\par\medskip\noindent\hypertarget{app-F7I-IX-044.1}{}\textbf{F7I-IX-044.1}\quad p. 362; misprint\_pass.\par
\noindent Imprimé : {\ttfamily\footnotesize l'application\allowbreak{}}\par
\noindent Corrigé : {\ttfamily\footnotesize les\allowbreak{}}\par
\noindent Source : 7I:598; 7I:597. \nolinkurl{fixes/FIXLOG_7I_CURRENT.jsonl#862}; \nolinkurl{fixes/MISPRINTS_7I.json#597}\par

\par\medskip\noindent\hypertarget{app-F7I-IX-044.2}{}\textbf{F7I-IX-044.2}\quad p. 362; misprint\_pass.\par
\noindent Imprimé : {\ttfamily\footnotesize ensembliste\allowbreak{}}\par
\noindent Corrigé : {\ttfamily\footnotesize applications\allowbreak{} ensemblistes\allowbreak{}}\par
\noindent Source : 7I:598; 7I:597. \nolinkurl{fixes/FIXLOG_7I_CURRENT.jsonl#862}; \nolinkurl{fixes/MISPRINTS_7I.json#597}\par

\par\medskip\noindent\hypertarget{app-F7I-IX-044.3}{}\textbf{F7I-IX-044.3}\quad p. 362; misprint\_pass.\par
\noindent Imprimé : {\ttfamily\footnotesize jacente\allowbreak{}}\par
\noindent Corrigé : {\ttfamily\footnotesize jacentes\allowbreak{}}\par
\noindent Source : 7I:598; 7I:597. \nolinkurl{fixes/FIXLOG_7I_CURRENT.jsonl#862}; \nolinkurl{fixes/MISPRINTS_7I.json#597}\par

\par\medskip\noindent\hypertarget{app-F7I-IX-045.1}{}\textbf{F7I-IX-045.1}\quad p. 363; misprint\_pass.\par
\noindent Imprimé : {\ttfamily\footnotesize derniere\allowbreak{}}\par
\noindent Corrigé : {\ttfamily\footnotesize dernière\allowbreak{}}\par
\noindent Source : 7I:599. \nolinkurl{fixes/FIXLOG_7I_CURRENT.jsonl#863}\par

\par\medskip\noindent\hypertarget{app-F7I-IX-046.1}{}\textbf{F7I-IX-046.1}\quad p. 364; misprint\_pass.\par
\noindent Imprimé : {\ttfamily\footnotesize manisfestement\allowbreak{}}\par
\noindent Corrigé : {\ttfamily\footnotesize manifestement\allowbreak{}}\par
\noindent Source : 7I:601; 7I:600. \nolinkurl{fixes/FIXLOG_7I_CURRENT.jsonl#864}; \nolinkurl{fixes/MISPRINTS_7I.json#600}\par

\par\medskip\noindent\hypertarget{app-F7I-IX-047.1}{}\textbf{F7I-IX-047.1}\quad p. 366; misprint\_pass.\par
\noindent Imprimé : {\ttfamily\footnotesize en\allowbreak{}}\par
\noindent Corrigé : {\ttfamily\footnotesize un\allowbreak{}}\par
\noindent Source : 7I:590; 7I:589. \nolinkurl{fixes/FIXLOG_7I_CURRENT.jsonl#865}; \nolinkurl{fixes/MISPRINTS_7I.json#589}\par

\par\medskip\noindent\hypertarget{app-F7I-IX-048.1}{}\textbf{F7I-IX-048.1}\quad p. 367; inherited corrected-layer difference.\par
\noindent Imprimé : {\ttfamily\footnotesize T\allowbreak{}}\par
\noindent Corrigé : {\ttfamily\footnotesize (1-nT)\allowbreak{}}\par
\noindent Source : ADOPT-F7I-IX-048.1. \nolinkurl{fr/sga7i/diplomatic/expose_IX_body.tex}; \nolinkurl{fr/sga7i/current/expose_IX_body.tex}\par

\par\medskip\noindent\hypertarget{app-F7I-IX-049.1}{}\textbf{F7I-IX-049.1}\quad p. 368; printed misprint.\par
\noindent Imprimé : {\ttfamily\footnotesize n\allowbreak{}}\par
\noindent Corrigé : {\ttfamily\footnotesize m\allowbreak{}}\par
\noindent Source : 7I:591. \nolinkurl{fixes/MISPRINTS_7I.json#591}\par

\par\medskip\noindent\hypertarget{app-F7I-IX-050.1}{}\textbf{F7I-IX-050.1}\quad p. 368; misprint\_pass.\par
\noindent Imprimé : \emph{(vide)}\par
\noindent Corrigé : {\ttfamily\footnotesize  dans\allowbreak{}}\par
\noindent Source : 7I:593. \nolinkurl{fixes/FIXLOG_7I_CURRENT.jsonl#866}\par

\par\medskip\noindent\hypertarget{app-F7I-IX-050.2}{}\textbf{F7I-IX-050.2}\quad p. 368; misprint\_pass.\par
\noindent Imprimé : {\ttfamily\footnotesize montrent\allowbreak{}}\par
\noindent Corrigé : {\ttfamily\footnotesize montre\allowbreak{}}\par
\noindent Source : 7I:594. \nolinkurl{fixes/FIXLOG_7I_CURRENT.jsonl#867}\par

\par\medskip\noindent\hypertarget{app-F7I-IX-050.3}{}\textbf{F7I-IX-050.3}\quad p. 368; misprint\_pass.\par
\noindent Imprimé : {\ttfamily\footnotesize opèrant\allowbreak{}}\par
\noindent Corrigé : {\ttfamily\footnotesize opérant\allowbreak{}}\par
\noindent Source : 7I:595. \nolinkurl{fixes/FIXLOG_7I_CURRENT.jsonl#868}\par

\par\medskip\noindent\hypertarget{app-F7I-IX-051.1}{}\textbf{F7I-IX-051.1}\quad p. 369; misprint\_pass.\par
\noindent Imprimé : {\ttfamily\footnotesize il\allowbreak{}}\par
\noindent Corrigé : {\ttfamily\footnotesize ils\allowbreak{}}\par
\noindent Source : 7I:596. \nolinkurl{fixes/FIXLOG_7I_CURRENT.jsonl#869}\par

\par\medskip\noindent\hypertarget{app-F7I-IX-052.1}{}\textbf{F7I-IX-052.1}\quad p. 371; misprint\_in\_print.\par
\noindent Imprimé : \emph{(vide)}\par
\noindent Corrigé : {\ttfamily\footnotesize '\allowbreak{}}\par
\noindent Source : FIX-7I-D-0595. \nolinkurl{fixes/FIXLOG_7I.jsonl#595}\par

\par\medskip\noindent\hypertarget{app-F7I-IX-052.2}{}\textbf{F7I-IX-052.2}\quad p. 371; misprint\_in\_print.\par
\noindent Imprimé : {\ttfamily\footnotesize \{\allowbreak{}}\par
\noindent Corrigé : \emph{(vide)}\par
\noindent Source : FIX-7I-D-0595. \nolinkurl{fixes/FIXLOG_7I.jsonl#595}\par

\par\medskip\noindent\hypertarget{app-F7I-IX-052.3}{}\textbf{F7I-IX-052.3}\quad p. 371; misprint\_in\_print.\par
\noindent Imprimé : {\ttfamily\footnotesize '\}\allowbreak{}}\par
\noindent Corrigé : \emph{(vide)}\par
\noindent Source : FIX-7I-D-0595. \nolinkurl{fixes/FIXLOG_7I.jsonl#595}\par

\par\medskip\noindent\hypertarget{app-F7I-IX-052.4}{}\textbf{F7I-IX-052.4}\quad p. 371; misprint\_pass.\par
\noindent Imprimé : {\ttfamily\footnotesize obtenu\allowbreak{}}\par
\noindent Corrigé : {\ttfamily\footnotesize obtenue\allowbreak{}}\par
\noindent Source : 7I:603. \nolinkurl{fixes/FIXLOG_7I_CURRENT.jsonl#870}\par

\par\medskip\noindent\hypertarget{app-F7I-IX-053.1}{}\textbf{F7I-IX-053.1}\quad p. 372; misprint\_pass.\par
\noindent Imprimé : {\ttfamily\footnotesize là\allowbreak{}}\par
\noindent Corrigé : {\ttfamily\footnotesize la\allowbreak{}}\par
\noindent Source : 7I:604. \nolinkurl{fixes/FIXLOG_7I_CURRENT.jsonl#871}\par

\par\medskip\noindent\hypertarget{app-F7I-IX-054.1}{}\textbf{F7I-IX-054.1}\quad p. 373; misprint\_pass.\par
\noindent Imprimé : {\ttfamily\footnotesize hypohtèses\allowbreak{}}\par
\noindent Corrigé : {\ttfamily\footnotesize hypothèses\allowbreak{}}\par
\noindent Source : 7I:605; 7I:604. \nolinkurl{fixes/FIXLOG_7I_CURRENT.jsonl#872}; \nolinkurl{fixes/MISPRINTS_7I.json#604}\par

\par\medskip\noindent\hypertarget{app-F7I-IX-055.1}{}\textbf{F7I-IX-055.1}\quad p. 373; misprint\_pass.\par
\noindent Imprimé : {\ttfamily\footnotesize déduite\allowbreak{}}\par
\noindent Corrigé : {\ttfamily\footnotesize déduit\allowbreak{}}\par
\noindent Source : 7I:606. \nolinkurl{fixes/FIXLOG_7I_CURRENT.jsonl#873}\par

\par\medskip\noindent\hypertarget{app-F7I-IX-056.1}{}\textbf{F7I-IX-056.1}\quad p. 375; misprint\_pass.\par
\noindent Imprimé : {\ttfamily\footnotesize du\allowbreak{}}\par
\noindent Corrigé : {\ttfamily\footnotesize des\allowbreak{}}\par
\noindent Source : 7I:607. \nolinkurl{fixes/FIXLOG_7I_CURRENT.jsonl#874}\par

\par\medskip\noindent\hypertarget{app-F7I-IX-057.1}{}\textbf{F7I-IX-057.1}\quad p. 377; printed misprint.\par
\noindent Imprimé : {\ttfamily\footnotesize S\allowbreak{}}\par
\noindent Corrigé : {\ttfamily\footnotesize K\allowbreak{}}\par
\noindent Source : 7I:607. \nolinkurl{fixes/MISPRINTS_7I.json#607}\par

\par\medskip\noindent\hypertarget{app-F7I-IX-058.1}{}\textbf{F7I-IX-058.1}\quad p. 379; misprint\_pass.\par
\noindent Imprimé : \emph{(vide)}\par
\noindent Corrigé : {\ttfamily\footnotesize \_K\allowbreak{}}\par
\noindent Source : 7I:609; 7I:608. \nolinkurl{fixes/FIXLOG_7I_CURRENT.jsonl#875}; \nolinkurl{fixes/MISPRINTS_7I.json#608}\par

\par\medskip\noindent\hypertarget{app-F7I-IX-059.1}{}\textbf{F7I-IX-059.1}\quad p. 380; misprint\_pass.\par
\noindent Imprimé : {\ttfamily\footnotesize rappelera\allowbreak{}}\par
\noindent Corrigé : {\ttfamily\footnotesize rappellera\allowbreak{}}\par
\noindent Source : 7I:610; 7I:609. \nolinkurl{fixes/FIXLOG_7I_CURRENT.jsonl#876}; \nolinkurl{fixes/MISPRINTS_7I.json#609}\par

\par\medskip\noindent\hypertarget{app-F7I-IX-060.1}{}\textbf{F7I-IX-060.1}\quad p. 380; misprint\_pass.\par
\noindent Imprimé : {\ttfamily\footnotesize déploye\allowbreak{}}\par
\noindent Corrigé : {\ttfamily\footnotesize déploie\allowbreak{}}\par
\noindent Source : 7I:611; 7I:610. \nolinkurl{fixes/FIXLOG_7I_CURRENT.jsonl#877}; \nolinkurl{fixes/MISPRINTS_7I.json#610}\par

\par\medskip\noindent\hypertarget{app-F7I-IX-060.2}{}\textbf{F7I-IX-060.2}\quad p. 380; misprint\_pass.\par
\noindent Imprimé : \emph{(vide)}\par
\noindent Corrigé : {\ttfamily\footnotesize '\allowbreak{}}\par
\noindent Source : 7I:612; 7I:611. \nolinkurl{fixes/FIXLOG_7I_CURRENT.jsonl#878}; \nolinkurl{fixes/MISPRINTS_7I.json#611}\par

\par\medskip\noindent\hypertarget{app-F7I-IX-060.3}{}\textbf{F7I-IX-060.3}\quad p. 380; misprint\_pass.\par
\noindent Imprimé : {\ttfamily\footnotesize '\allowbreak{}}\par
\noindent Corrigé : \emph{(vide)}\par
\noindent Source : 7I:612; 7I:611. \nolinkurl{fixes/FIXLOG_7I_CURRENT.jsonl#878}; \nolinkurl{fixes/MISPRINTS_7I.json#611}\par

\par\medskip\noindent\hypertarget{app-F7I-IX-061.1}{}\textbf{F7I-IX-061.1}\quad p. 383; misprint\_pass.\par
\noindent Imprimé : {\ttfamily\footnotesize connexe\allowbreak{}}\par
\noindent Corrigé : {\ttfamily\footnotesize étale\allowbreak{}}\par
\noindent Source : 7I:613. \nolinkurl{fixes/FIXLOG_7I_CURRENT.jsonl#879}\par

\par\medskip\noindent\hypertarget{app-F7I-IX-062.1}{}\textbf{F7I-IX-062.1}\quad p. 385; misprint\_pass.\par
\noindent Imprimé : {\ttfamily\footnotesize Soir\allowbreak{}}\par
\noindent Corrigé : {\ttfamily\footnotesize Soit\allowbreak{}}\par
\noindent Source : 7I:614; 7I:613. \nolinkurl{fixes/FIXLOG_7I_CURRENT.jsonl#880}; \nolinkurl{fixes/MISPRINTS_7I.json#613}\par

\par\medskip\noindent\hypertarget{app-F7I-IX-063.1}{}\textbf{F7I-IX-063.1}\quad p. 386; misprint\_pass.\par
\noindent Imprimé : {\ttfamily\footnotesize X\allowbreak{}}\par
\noindent Corrigé : {\ttfamily\footnotesize G\allowbreak{}}\par
\noindent Source : 7I:615; 7I:614. \nolinkurl{fixes/FIXLOG_7I_CURRENT.jsonl#881}; \nolinkurl{fixes/MISPRINTS_7I.json#614}\par

\par\medskip\noindent\hypertarget{app-F7I-IX-064.1}{}\textbf{F7I-IX-064.1}\quad p. 388; inherited corrected-layer difference.\par
\noindent Imprimé : \emph{(vide)}\par
\noindent Corrigé : {\ttfamily\footnotesize \_K\allowbreak{}}\par
\noindent Source : ADOPT-F7I-IX-064.1. \nolinkurl{fr/sga7i/diplomatic/expose_IX_body.tex}; \nolinkurl{fr/sga7i/current/expose_IX_body.tex}\par

\par\medskip\noindent\hypertarget{app-F7I-IX-065.1}{}\textbf{F7I-IX-065.1}\quad p. 388; inherited corrected-layer difference.\par
\noindent Imprimé : \emph{(vide)}\par
\noindent Corrigé : {\ttfamily\footnotesize \_K\allowbreak{}}\par
\noindent Source : ADOPT-F7I-IX-065.1. \nolinkurl{fr/sga7i/diplomatic/expose_IX_body.tex}; \nolinkurl{fr/sga7i/current/expose_IX_body.tex}\par

\par\medskip\noindent\hypertarget{app-F7I-IX-066.1}{}\textbf{F7I-IX-066.1}\quad p. 389; printed misprint.\par
\noindent Imprimé : {\ttfamily\footnotesize \textbackslash{}overline\{\allowbreak{}}\par
\noindent Corrigé : \emph{(vide)}\par
\noindent Source : 7I:615. \nolinkurl{fixes/MISPRINTS_7I.json#615}\par

\par\medskip\noindent\hypertarget{app-F7I-IX-066.2}{}\textbf{F7I-IX-066.2}\quad p. 389; printed misprint.\par
\noindent Imprimé : {\ttfamily\footnotesize \}\allowbreak{}}\par
\noindent Corrigé : \emph{(vide)}\par
\noindent Source : 7I:615. \nolinkurl{fixes/MISPRINTS_7I.json#615}\par

\par\medskip\noindent\hypertarget{app-F7I-IX-067.1}{}\textbf{F7I-IX-067.1}\quad p. 390; designated-witness attribution.\par
\noindent Imprimé : \emph{(vide)}\par
\noindent Corrigé : {\ttfamily\footnotesize \textasciicircum{}\{\textbackslash{}times\}\allowbreak{}}\par
\noindent Source : WITNESS-F7I-IX-067.1. \nolinkurl{fr/sga7i/diplomatic/expose_IX_body.tex}; \nolinkurl{fr/sga7i/current/expose_IX_body.tex}\par

\par\medskip\noindent\hypertarget{app-F7I-IX-068.1}{}\textbf{F7I-IX-068.1}\quad p. 391; misprint\_pass.\par
\noindent Imprimé : \emph{(vide)}\par
\noindent Corrigé : {\ttfamily\footnotesize )\allowbreak{}}\par
\noindent Source : 7I:619; 7I:618. \nolinkurl{fixes/FIXLOG_7I_CURRENT.jsonl#882}; \nolinkurl{fixes/MISPRINTS_7I.json#618}\par

\par\medskip\noindent\hypertarget{app-F7I-IX-069.1}{}\textbf{F7I-IX-069.1}\quad p. 392; misprint\_pass.\par
\noindent Imprimé : {\ttfamily\footnotesize de\allowbreak{} }\par
\noindent Corrigé : \emph{(vide)}\par
\noindent Source : 7I:620; 7I:619. \nolinkurl{fixes/FIXLOG_7I_CURRENT.jsonl#883}; \nolinkurl{fixes/MISPRINTS_7I.json#619}\par

\par\medskip\noindent\hypertarget{app-F7I-IX-070.1}{}\textbf{F7I-IX-070.1}\quad p. 394; misprint\_in\_print.\par
\noindent Imprimé : {\ttfamily\footnotesize résoud\allowbreak{}}\par
\noindent Corrigé : {\ttfamily\footnotesize résout\allowbreak{}}\par
\noindent Source : FIX-7I-D-0607. \nolinkurl{fixes/FIXLOG_7I.jsonl#607}\par

\par\medskip\noindent\hypertarget{app-F7I-IX-071.1}{}\textbf{F7I-IX-071.1}\quad p. 395; printed misprint.\par
\noindent Imprimé : {\ttfamily\footnotesize '\allowbreak{}}\par
\noindent Corrigé : \emph{(vide)}\par
\noindent Source : 7I:621. \nolinkurl{fixes/MISPRINTS_7I.json#621}\par

\par\medskip\noindent\hypertarget{app-F7I-IX-071.2}{}\textbf{F7I-IX-071.2}\quad p. 395; designated-witness attribution.\par
\noindent Imprimé : {\ttfamily\footnotesize S\allowbreak{}}\par
\noindent Corrigé : {\ttfamily\footnotesize K\allowbreak{}}\par
\noindent Source : WITNESS-F7I-IX-071.2. \nolinkurl{fr/sga7i/diplomatic/expose_IX_body.tex}; \nolinkurl{fr/sga7i/current/expose_IX_body.tex}\par

\par\medskip\noindent\hypertarget{app-F7I-IX-072.1}{}\textbf{F7I-IX-072.1}\quad p. 396; printed misprint.\par
\noindent Imprimé : {\ttfamily\footnotesize Barsotit\allowbreak{}}\par
\noindent Corrigé : {\ttfamily\footnotesize Barsotti\allowbreak{}}\par
\noindent Source : 7I:623. \nolinkurl{fixes/MISPRINTS_7I.json#623}\par

\par\medskip\noindent\hypertarget{app-F7I-IX-072.2}{}\textbf{F7I-IX-072.2}\quad p. 396; misprint\_pass.\par
\noindent Imprimé : \emph{(vide)}\par
\noindent Corrigé : {\ttfamily\footnotesize \{\allowbreak{}}\par
\noindent Source : 7I:625; 7I:624. \nolinkurl{fixes/FIXLOG_7I_CURRENT.jsonl#885}; \nolinkurl{fixes/MISPRINTS_7I.json#624}\par

\par\medskip\noindent\hypertarget{app-F7I-IX-072.3}{}\textbf{F7I-IX-072.3}\quad p. 396; misprint\_pass.\par
\noindent Imprimé : \emph{(vide)}\par
\noindent Corrigé : {\ttfamily\footnotesize '\}\allowbreak{}}\par
\noindent Source : 7I:625; 7I:624. \nolinkurl{fixes/FIXLOG_7I_CURRENT.jsonl#885}; \nolinkurl{fixes/MISPRINTS_7I.json#624}\par

\par\medskip\noindent\hypertarget{app-F7I-IX-072.4}{}\textbf{F7I-IX-072.4}\quad p. 396; misprint\_pass.\par
\noindent Imprimé : \emph{(vide)}\par
\noindent Corrigé : {\ttfamily\footnotesize \{\allowbreak{}}\par
\noindent Source : 7I:625; 7I:624. \nolinkurl{fixes/FIXLOG_7I_CURRENT.jsonl#885}; \nolinkurl{fixes/MISPRINTS_7I.json#624}\par

\par\medskip\noindent\hypertarget{app-F7I-IX-072.5}{}\textbf{F7I-IX-072.5}\quad p. 396; misprint\_pass.\par
\noindent Imprimé : \emph{(vide)}\par
\noindent Corrigé : {\ttfamily\footnotesize '\}\allowbreak{}}\par
\noindent Source : 7I:625; 7I:624. \nolinkurl{fixes/FIXLOG_7I_CURRENT.jsonl#885}; \nolinkurl{fixes/MISPRINTS_7I.json#624}\par

\par\medskip\noindent\hypertarget{app-F7I-IX-073.1}{}\textbf{F7I-IX-073.1}\quad p. 397; designated-witness attribution.\par
\noindent Imprimé : {\ttfamily\footnotesize k\allowbreak{}}\par
\noindent Corrigé : {\ttfamily\footnotesize K\allowbreak{}}\par
\noindent Source : WITNESS-F7I-IX-073.1. \nolinkurl{fr/sga7i/diplomatic/expose_IX_body.tex}; \nolinkurl{fr/sga7i/current/expose_IX_body.tex}\par

\par\medskip\noindent\hypertarget{app-F7I-IX-073.2}{}\textbf{F7I-IX-073.2}\quad p. 397; designated-witness attribution.\par
\noindent Imprimé : {\ttfamily\footnotesize k\allowbreak{}}\par
\noindent Corrigé : {\ttfamily\footnotesize K\allowbreak{}}\par
\noindent Source : WITNESS-F7I-IX-073.2. \nolinkurl{fr/sga7i/diplomatic/expose_IX_body.tex}; \nolinkurl{fr/sga7i/current/expose_IX_body.tex}\par

\par\medskip\noindent\hypertarget{app-F7I-IX-073.3}{}\textbf{F7I-IX-073.3}\quad p. 397; misprint\_in\_print.\par
\noindent Imprimé : \emph{(vide)}\par
\noindent Corrigé : {\ttfamily\footnotesize '\allowbreak{}}\par
\noindent Source : FIX-7I-D-0609. \nolinkurl{fixes/FIXLOG_7I.jsonl#609}\par

\par\medskip\noindent\hypertarget{app-F7I-IX-073.4}{}\textbf{F7I-IX-073.4}\quad p. 397; misprint\_pass.\par
\noindent Imprimé : {\ttfamily\footnotesize 12\allowbreak{}}\par
\noindent Corrigé : {\ttfamily\footnotesize 13\allowbreak{}}\par
\noindent Source : 7I:628; 7I:627. \nolinkurl{fixes/FIXLOG_7I_CURRENT.jsonl#886}; \nolinkurl{fixes/MISPRINTS_7I.json#627}\par

\par\medskip\noindent\hypertarget{app-F7I-IX-073.5}{}\textbf{F7I-IX-073.5}\quad p. 397; misprint\_pass.\par
\noindent Imprimé : {\ttfamily\footnotesize a\allowbreak{}}\par
\noindent Corrigé : {\ttfamily\footnotesize b\allowbreak{}}\par
\noindent Source : 7I:628; 7I:627. \nolinkurl{fixes/FIXLOG_7I_CURRENT.jsonl#886}; \nolinkurl{fixes/MISPRINTS_7I.json#627}\par

\par\medskip\noindent\hypertarget{app-F7I-IX-073.6}{}\textbf{F7I-IX-073.6}\quad p. 397; misprint\_pass.\par
\noindent Imprimé : {\ttfamily\footnotesize 12\allowbreak{}}\par
\noindent Corrigé : {\ttfamily\footnotesize 13\allowbreak{}}\par
\noindent Source : 7I:629; 7I:628. \nolinkurl{fixes/FIXLOG_7I_CURRENT.jsonl#887}; \nolinkurl{fixes/MISPRINTS_7I.json#628}\par

\par\medskip\noindent\hypertarget{app-F7I-IX-073.7}{}\textbf{F7I-IX-073.7}\quad p. 397; misprint\_pass.\par
\noindent Imprimé : {\ttfamily\footnotesize a\allowbreak{}}\par
\noindent Corrigé : {\ttfamily\footnotesize b\allowbreak{}}\par
\noindent Source : 7I:629; 7I:628. \nolinkurl{fixes/FIXLOG_7I_CURRENT.jsonl#887}; \nolinkurl{fixes/MISPRINTS_7I.json#628}\par

\par\medskip\noindent\hypertarget{app-F7I-IX-073.8}{}\textbf{F7I-IX-073.8}\quad p. 397; misprint\_pass.\par
\noindent Imprimé : {\ttfamily\footnotesize égal\allowbreak{}}\par
\noindent Corrigé : {\ttfamily\footnotesize égale\allowbreak{}}\par
\noindent Source : 7I:630; 7I:629. \nolinkurl{fixes/FIXLOG_7I_CURRENT.jsonl#888}; \nolinkurl{fixes/MISPRINTS_7I.json#629}\par

\par\medskip\noindent\hypertarget{app-F7I-IX-074.1}{}\textbf{F7I-IX-074.1}\quad p. 398; printed misprint.\par
\noindent Imprimé : {\ttfamily\footnotesize Barsoti\allowbreak{}}\par
\noindent Corrigé : {\ttfamily\footnotesize Barsotti\allowbreak{}}\par
\noindent Source : 7I:631. \nolinkurl{fixes/MISPRINTS_7I.json#631}\par

\par\medskip\noindent\hypertarget{app-F7I-IX-074.2}{}\textbf{F7I-IX-074.2}\quad p. 398; designated-witness attribution.\par
\noindent Imprimé : {\ttfamily\footnotesize 6\allowbreak{}}\par
\noindent Corrigé : {\ttfamily\footnotesize 5\allowbreak{}}\par
\noindent Source : WITNESS-F7I-IX-074.2. \nolinkurl{fr/sga7i/diplomatic/expose_IX_body.tex}; \nolinkurl{fr/sga7i/current/expose_IX_body.tex}\par

\par\medskip\noindent\hypertarget{app-F7I-IX-074.3}{}\textbf{F7I-IX-074.3}\quad p. 398; designated-witness attribution.\par
\noindent Imprimé : {\ttfamily\footnotesize 5\allowbreak{}}\par
\noindent Corrigé : {\ttfamily\footnotesize 12\allowbreak{}}\par
\noindent Source : WITNESS-F7I-IX-074.3. \nolinkurl{fr/sga7i/diplomatic/expose_IX_body.tex}; \nolinkurl{fr/sga7i/current/expose_IX_body.tex}\par

\par\medskip\noindent\hypertarget{app-F7I-IX-075.1}{}\textbf{F7I-IX-075.1}\quad p. 399; misprint\_pass.\par
\noindent Imprimé : {\ttfamily\footnotesize da\allowbreak{}}\par
\noindent Corrigé : {\ttfamily\footnotesize de\allowbreak{}}\par
\noindent Source : 7I:634; 7I:633. \nolinkurl{fixes/FIXLOG_7I_CURRENT.jsonl#890}; \nolinkurl{fixes/MISPRINTS_7I.json#633}\par

\par\medskip\noindent\hypertarget{app-F7I-IX-076.1}{}\textbf{F7I-IX-076.1}\quad p. 401; printed misprint.\par
\noindent Imprimé : {\ttfamily\footnotesize M\allowbreak{}}\par
\noindent Corrigé : {\ttfamily\footnotesize L\allowbreak{}}\par
\noindent Source : 7I:824. \nolinkurl{fixes/MISPRINTS_7I.json#824}\par

\par\medskip\noindent\hypertarget{app-F7I-IX-076.2}{}\textbf{F7I-IX-076.2}\quad p. 401; designated-witness attribution.\par
\noindent Imprimé : {\ttfamily\footnotesize de\allowbreak{}}\par
\noindent Corrigé : {\ttfamily\footnotesize sur\allowbreak{}}\par
\noindent Source : WITNESS-F7I-IX-076.2. \nolinkurl{fr/sga7i/diplomatic/expose_IX_body.tex}; \nolinkurl{fr/sga7i/current/expose_IX_body.tex}\par

\par\medskip\noindent\hypertarget{app-F7I-IX-076.3}{}\textbf{F7I-IX-076.3}\quad p. 401; designated-witness attribution.\par
\noindent Imprimé : {\ttfamily\footnotesize G\allowbreak{}}\par
\noindent Corrigé : {\ttfamily\footnotesize A\allowbreak{}}\par
\noindent Source : WITNESS-F7I-IX-076.3. \nolinkurl{fr/sga7i/diplomatic/expose_IX_body.tex}; \nolinkurl{fr/sga7i/current/expose_IX_body.tex}\par

\par\medskip\noindent\hypertarget{app-F7I-IX-076.4}{}\textbf{F7I-IX-076.4}\quad p. 401; designated-witness attribution.\par
\noindent Imprimé : {\ttfamily\footnotesize G\allowbreak{}}\par
\noindent Corrigé : {\ttfamily\footnotesize A\allowbreak{}}\par
\noindent Source : WITNESS-F7I-IX-076.4. \nolinkurl{fr/sga7i/diplomatic/expose_IX_body.tex}; \nolinkurl{fr/sga7i/current/expose_IX_body.tex}\par

\par\medskip\noindent\hypertarget{app-F7I-IX-076.5}{}\textbf{F7I-IX-076.5}\quad p. 401; misprint\_pass.\par
\noindent Imprimé : \emph{(vide)}\par
\noindent Corrigé : {\ttfamily\footnotesize \textbackslash{}hat\{\allowbreak{}}\par
\noindent Source : 7I:827; 7I:826. \nolinkurl{fixes/FIXLOG_7I_CURRENT.jsonl#891}; \nolinkurl{fixes/MISPRINTS_7I.json#826}\par

\par\medskip\noindent\hypertarget{app-F7I-IX-076.6}{}\textbf{F7I-IX-076.6}\quad p. 401; misprint\_pass.\par
\noindent Imprimé : \emph{(vide)}\par
\noindent Corrigé : {\ttfamily\footnotesize \}\allowbreak{}}\par
\noindent Source : 7I:827; 7I:826. \nolinkurl{fixes/FIXLOG_7I_CURRENT.jsonl#891}; \nolinkurl{fixes/MISPRINTS_7I.json#826}\par

\par\medskip\noindent\hypertarget{app-F7I-IX-076.7}{}\textbf{F7I-IX-076.7}\quad p. 401; misprint\_pass.\par
\noindent Imprimé : \emph{(vide)}\par
\noindent Corrigé : {\ttfamily\footnotesize \textbackslash{}hat\{\allowbreak{}}\par
\noindent Source : 7I:828; 7I:827. \nolinkurl{fixes/FIXLOG_7I_CURRENT.jsonl#892}; \nolinkurl{fixes/MISPRINTS_7I.json#827}\par

\par\medskip\noindent\hypertarget{app-F7I-IX-076.8}{}\textbf{F7I-IX-076.8}\quad p. 401; misprint\_pass.\par
\noindent Imprimé : \emph{(vide)}\par
\noindent Corrigé : {\ttfamily\footnotesize \}\allowbreak{}}\par
\noindent Source : 7I:828; 7I:827. \nolinkurl{fixes/FIXLOG_7I_CURRENT.jsonl#892}; \nolinkurl{fixes/MISPRINTS_7I.json#827}\par

\par\medskip\noindent\hypertarget{app-F7I-IX-077.1}{}\textbf{F7I-IX-077.1}\quad p. 403; misprint\_pass.\par
\noindent Imprimé : {\ttfamily\footnotesize 3\allowbreak{}}\par
\noindent Corrigé : {\ttfamily\footnotesize 13\allowbreak{}}\par
\noindent Source : 7I:833; 7I:832. \nolinkurl{fixes/FIXLOG_7I_CURRENT.jsonl#894}; \nolinkurl{fixes/MISPRINTS_7I.json#832}\par

\par\medskip\noindent\hypertarget{app-F7I-IX-077.2}{}\textbf{F7I-IX-077.2}\quad p. 403; printed misprint.\par
\noindent Imprimé : {\ttfamily\footnotesize A\allowbreak{}}\par
\noindent Corrigé : {\ttfamily\footnotesize V\allowbreak{}}\par
\noindent Source : 7I:833. \nolinkurl{fixes/MISPRINTS_7I.json#833}\par

\par\medskip\noindent\hypertarget{app-F7I-IX-078.1}{}\textbf{F7I-IX-078.1}\quad p. 403; printed misprint.\par
\noindent Imprimé : {\ttfamily\footnotesize A\allowbreak{}}\par
\noindent Corrigé : {\ttfamily\footnotesize V\allowbreak{}}\par
\noindent Source : 7I:639; 7I:834. \nolinkurl{fixes/MISPRINTS_7I.json#639}; \nolinkurl{fixes/MISPRINTS_7I.json#834}\par

\par\medskip\noindent\hypertarget{app-F7I-IX-078.2}{}\textbf{F7I-IX-078.2}\quad p. 403; misprint\_pass.\par
\noindent Imprimé : \emph{(vide)}\par
\noindent Corrigé : {\ttfamily\footnotesize \textbackslash{}mathcal\{\allowbreak{}}\par
\noindent Source : 7I:642; 7I:641; 7I:836. \nolinkurl{fixes/FIXLOG_7I_CURRENT.jsonl#893}; \nolinkurl{fixes/MISPRINTS_7I.json#641}; \nolinkurl{fixes/MISPRINTS_7I.json#836}\par

\par\medskip\noindent\hypertarget{app-F7I-IX-078.3}{}\textbf{F7I-IX-078.3}\quad p. 403; misprint\_pass.\par
\noindent Imprimé : \emph{(vide)}\par
\noindent Corrigé : {\ttfamily\footnotesize \}\allowbreak{}}\par
\noindent Source : 7I:642; 7I:641; 7I:836. \nolinkurl{fixes/FIXLOG_7I_CURRENT.jsonl#893}; \nolinkurl{fixes/MISPRINTS_7I.json#641}; \nolinkurl{fixes/MISPRINTS_7I.json#836}\par

\par\medskip\noindent\hypertarget{app-F7I-IX-078.4}{}\textbf{F7I-IX-078.4}\quad p. 403; printed misprint.\par
\noindent Imprimé : {\ttfamily\footnotesize A\allowbreak{}}\par
\noindent Corrigé : {\ttfamily\footnotesize V\allowbreak{}}\par
\noindent Source : 7I:640; 7I:835. \nolinkurl{fixes/MISPRINTS_7I.json#640}; \nolinkurl{fixes/MISPRINTS_7I.json#835}\par

\par\medskip\noindent\hypertarget{app-F7I-IX-079.1}{}\textbf{F7I-IX-079.1}\quad p. 404; misprint\_pass.\par
\noindent Imprimé : \emph{(vide)}\par
\noindent Corrigé : {\ttfamily\footnotesize \textbackslash{}underline\{\allowbreak{}}\par
\noindent Source : 7I:643; 7I:838. \nolinkurl{fixes/FIXLOG_7I_CURRENT.jsonl#895}; \nolinkurl{fixes/MISPRINTS_7I.json#838}\par

\par\medskip\noindent\hypertarget{app-F7I-IX-079.2}{}\textbf{F7I-IX-079.2}\quad p. 404; misprint\_pass.\par
\noindent Imprimé : \emph{(vide)}\par
\noindent Corrigé : {\ttfamily\footnotesize \}\allowbreak{}}\par
\noindent Source : 7I:643; 7I:838. \nolinkurl{fixes/FIXLOG_7I_CURRENT.jsonl#895}; \nolinkurl{fixes/MISPRINTS_7I.json#838}\par

\par\medskip\noindent\hypertarget{app-F7I-IX-080.1}{}\textbf{F7I-IX-080.1}\quad p. 405; inherited corrected-layer difference.\par
\noindent Imprimé : \emph{(vide)}\par
\noindent Corrigé : {\ttfamily\footnotesize \_o\allowbreak{}}\par
\noindent Source : ADOPT-F7I-IX-080.1. \nolinkurl{fr/sga7i/diplomatic/expose_IX_body.tex}; \nolinkurl{fr/sga7i/current/expose_IX_body.tex}\par

\par\medskip\noindent\hypertarget{app-F7I-IX-080.2}{}\textbf{F7I-IX-080.2}\quad p. 405; inherited corrected-layer difference.\par
\noindent Imprimé : \emph{(vide)}\par
\noindent Corrigé : {\ttfamily\footnotesize \_n\allowbreak{}}\par
\noindent Source : ADOPT-F7I-IX-080.2. \nolinkurl{fr/sga7i/diplomatic/expose_IX_body.tex}; \nolinkurl{fr/sga7i/current/expose_IX_body.tex}\par

\par\medskip\noindent\hypertarget{app-F7I-IX-080.3}{}\textbf{F7I-IX-080.3}\quad p. 405; inherited corrected-layer difference.\par
\noindent Imprimé : \emph{(vide)}\par
\noindent Corrigé : {\ttfamily\footnotesize \_n\allowbreak{}}\par
\noindent Source : ADOPT-F7I-IX-080.3. \nolinkurl{fr/sga7i/diplomatic/expose_IX_body.tex}; \nolinkurl{fr/sga7i/current/expose_IX_body.tex}\par

\par\medskip\noindent\hypertarget{app-F7I-IX-081.1}{}\textbf{F7I-IX-081.1}\quad p. 406; misprint\_pass.\par
\noindent Imprimé : {\ttfamily\footnotesize fait\allowbreak{}}\par
\noindent Corrigé : {\ttfamily\footnotesize faut\allowbreak{}}\par
\noindent Source : 7I:644; 7I:643; 7I:818. \nolinkurl{fixes/FIXLOG_7I_CURRENT.jsonl#896}; \nolinkurl{fixes/MISPRINTS_7I.json#643}; \nolinkurl{fixes/MISPRINTS_7I.json#818}\par

\par\medskip\noindent\hypertarget{app-F7I-IX-081.2}{}\textbf{F7I-IX-081.2}\quad p. 406; misprint\_pass.\par
\noindent Imprimé : {\ttfamily\footnotesize scindé\allowbreak{}}\par
\noindent Corrigé : {\ttfamily\footnotesize scinde\allowbreak{}}\par
\noindent Source : 7I:645. \nolinkurl{fixes/FIXLOG_7I_CURRENT.jsonl#897}\par

\par\medskip\noindent\hypertarget{app-F7I-IX-082.1}{}\textbf{F7I-IX-082.1}\quad p. 406; misprint\_pass.\par
\noindent Imprimé : {\ttfamily\footnotesize satisfait\allowbreak{}}\par
\noindent Corrigé : {\ttfamily\footnotesize satisfaite\allowbreak{}}\par
\noindent Source : 7I:646; 7I:645; 7I:820. \nolinkurl{fixes/FIXLOG_7I_CURRENT.jsonl#898}; \nolinkurl{fixes/MISPRINTS_7I.json#645}; \nolinkurl{fixes/MISPRINTS_7I.json#820}\par

\par\medskip\noindent\hypertarget{app-F7I-IX-083.1}{}\textbf{F7I-IX-083.1}\quad p. 408; misprint\_pass.\par
\noindent Imprimé : {\ttfamily\footnotesize interprèter\allowbreak{}}\par
\noindent Corrigé : {\ttfamily\footnotesize interpréter\allowbreak{}}\par
\noindent Source : 7I:647. \nolinkurl{fixes/FIXLOG_7I_CURRENT.jsonl#899}\par

\par\medskip\noindent\hypertarget{app-F7I-IX-084.1}{}\textbf{F7I-IX-084.1}\quad p. 409; misprint\_pass.\par
\noindent Imprimé : {\ttfamily\footnotesize interprèter\allowbreak{}}\par
\noindent Corrigé : {\ttfamily\footnotesize interpréter\allowbreak{}}\par
\noindent Source : 7I:648. \nolinkurl{fixes/FIXLOG_7I_CURRENT.jsonl#900}\par

\par\medskip\noindent\hypertarget{app-F7I-IX-085.1}{}\textbf{F7I-IX-085.1}\quad p. 410; misprint\_pass.\par
\noindent Imprimé : {\ttfamily\footnotesize IV\allowbreak{}}\par
\noindent Corrigé : {\ttfamily\footnotesize III\allowbreak{}}\par
\noindent Source : 7I:649; 7I:648; 7I:823. \nolinkurl{fixes/FIXLOG_7I_CURRENT.jsonl#901}; \nolinkurl{fixes/MISPRINTS_7I.json#648}; \nolinkurl{fixes/MISPRINTS_7I.json#823}\par

\par\medskip\noindent\hypertarget{app-F7I-IX-086.1}{}\textbf{F7I-IX-086.1}\quad p. 411; misprint\_pass.\par
\noindent Imprimé : {\ttfamily\footnotesize 4\allowbreak{}}\par
\noindent Corrigé : {\ttfamily\footnotesize 5\allowbreak{}}\par
\noindent Source : 7I:840; 7I:839. \nolinkurl{fixes/FIXLOG_7I_CURRENT.jsonl#902}; \nolinkurl{fixes/MISPRINTS_7I.json#839}\par

\par\medskip\noindent\hypertarget{app-F7I-IX-086.2}{}\textbf{F7I-IX-086.2}\quad p. 411; misprint\_pass.\par
\noindent Imprimé : {\ttfamily\footnotesize 3\allowbreak{}}\par
\noindent Corrigé : {\ttfamily\footnotesize 4\allowbreak{}}\par
\noindent Source : 7I:841; 7I:840. \nolinkurl{fixes/FIXLOG_7I_CURRENT.jsonl#903}; \nolinkurl{fixes/MISPRINTS_7I.json#840}\par

\par\medskip\noindent\hypertarget{app-F7I-IX-086.3}{}\textbf{F7I-IX-086.3}\quad p. 411; misprint\_pass.\par
\noindent Imprimé : {\ttfamily\footnotesize 1\allowbreak{}}\par
\noindent Corrigé : {\ttfamily\footnotesize 2\allowbreak{}}\par
\noindent Source : 7I:841; 7I:840. \nolinkurl{fixes/FIXLOG_7I_CURRENT.jsonl#903}; \nolinkurl{fixes/MISPRINTS_7I.json#840}\par

\par\medskip\noindent\hypertarget{app-F7I-IX-087.1}{}\textbf{F7I-IX-087.1}\quad p. 413; misprint\_pass.\par
\noindent Imprimé : {\ttfamily\footnotesize aves\allowbreak{}}\par
\noindent Corrigé : {\ttfamily\footnotesize avec\allowbreak{}}\par
\noindent Source : 7I:650. \nolinkurl{fixes/FIXLOG_7I_CURRENT.jsonl#904}\par

\par\medskip\noindent\hypertarget{app-F7I-IX-088.1}{}\textbf{F7I-IX-088.1}\quad p. 414; printed misprint.\par
\noindent Imprimé : {\ttfamily\footnotesize 5\allowbreak{}}\par
\noindent Corrigé : {\ttfamily\footnotesize 3\allowbreak{}}\par
\noindent Source : 7I:842. \nolinkurl{fixes/MISPRINTS_7I.json#842}\par

\par\medskip\noindent\hypertarget{app-F7I-IX-089.1}{}\textbf{F7I-IX-089.1}\quad p. 415; printed misprint.\par
\noindent Imprimé : {\ttfamily\footnotesize 5\allowbreak{}}\par
\noindent Corrigé : {\ttfamily\footnotesize 3\allowbreak{}}\par
\noindent Source : 7I:844. \nolinkurl{fixes/MISPRINTS_7I.json#844}\par

\par\medskip\noindent\hypertarget{app-F7I-IX-090.1}{}\textbf{F7I-IX-090.1}\quad p. 416; printed misprint.\par
\noindent Imprimé : \emph{(vide)}\par
\noindent Corrigé : {\ttfamily\footnotesize '\allowbreak{}}\par
\noindent Source : 7I:845. \nolinkurl{fixes/MISPRINTS_7I.json#845}\par

\par\medskip\noindent\hypertarget{app-F7I-IX-090.2}{}\textbf{F7I-IX-090.2}\quad p. 416; printed misprint.\par
\noindent Imprimé : {\ttfamily\footnotesize \{\allowbreak{}}\par
\noindent Corrigé : \emph{(vide)}\par
\noindent Source : 7I:845. \nolinkurl{fixes/MISPRINTS_7I.json#845}\par

\par\medskip\noindent\hypertarget{app-F7I-IX-090.3}{}\textbf{F7I-IX-090.3}\quad p. 416; printed misprint.\par
\noindent Imprimé : {\ttfamily\footnotesize '\}\allowbreak{}}\par
\noindent Corrigé : \emph{(vide)}\par
\noindent Source : 7I:845. \nolinkurl{fixes/MISPRINTS_7I.json#845}\par

\par\medskip\noindent\hypertarget{app-F7I-IX-090.4}{}\textbf{F7I-IX-090.4}\quad p. 416; printed misprint.\par
\noindent Imprimé : {\ttfamily\footnotesize '\allowbreak{}}\par
\noindent Corrigé : \emph{(vide)}\par
\noindent Source : 7I:846. \nolinkurl{fixes/MISPRINTS_7I.json#846}\par

\par\medskip\noindent\hypertarget{app-F7I-IX-091.1}{}\textbf{F7I-IX-091.1}\quad p. 417; misprint\_pass.\par
\noindent Imprimé : \emph{(vide)}\par
\noindent Corrigé : {\ttfamily\footnotesize \textasciicircum{}\{\textbackslash{}natural\}\allowbreak{}}\par
\noindent Source : 7I:653; 7I:652; 7I:847. \nolinkurl{fixes/FIXLOG_7I_CURRENT.jsonl#905}; \nolinkurl{fixes/MISPRINTS_7I.json#652}; \nolinkurl{fixes/MISPRINTS_7I.json#847}\par

\par\medskip\noindent\hypertarget{app-F7I-IX-092.1}{}\textbf{F7I-IX-092.1}\quad p. 417; misprint\_pass.\par
\noindent Imprimé : {\ttfamily\footnotesize dans\allowbreak{}}\par
\noindent Corrigé : {\ttfamily\footnotesize sont\allowbreak{}}\par
\noindent Source : 7I:654; 7I:653; 7I:848. \nolinkurl{fixes/FIXLOG_7I_CURRENT.jsonl#906}; \nolinkurl{fixes/MISPRINTS_7I.json#653}; \nolinkurl{fixes/MISPRINTS_7I.json#848}\par

\par\medskip\noindent\hypertarget{app-F7I-IX-093.1}{}\textbf{F7I-IX-093.1}\quad p. 418; misprint\_pass.\par
\noindent Imprimé : \emph{(vide)}\par
\noindent Corrigé : {\ttfamily\footnotesize \textbackslash{}underline\{\allowbreak{}}\par
\noindent Source : 7I:655; 7I:654; 7I:849. \nolinkurl{fixes/FIXLOG_7I_CURRENT.jsonl#907}; \nolinkurl{fixes/MISPRINTS_7I.json#654}; \nolinkurl{fixes/MISPRINTS_7I.json#849}\par

\par\medskip\noindent\hypertarget{app-F7I-IX-093.2}{}\textbf{F7I-IX-093.2}\quad p. 418; misprint\_pass.\par
\noindent Imprimé : \emph{(vide)}\par
\noindent Corrigé : {\ttfamily\footnotesize \}\allowbreak{}}\par
\noindent Source : 7I:655; 7I:654; 7I:849. \nolinkurl{fixes/FIXLOG_7I_CURRENT.jsonl#907}; \nolinkurl{fixes/MISPRINTS_7I.json#654}; \nolinkurl{fixes/MISPRINTS_7I.json#849}\par

\par\medskip\noindent\hypertarget{app-F7I-IX-094.1}{}\textbf{F7I-IX-094.1}\quad p. 419; printed misprint.\par
\noindent Imprimé : {\ttfamily\footnotesize u\allowbreak{}}\par
\noindent Corrigé : {\ttfamily\footnotesize x\allowbreak{}}\par
\noindent Source : 7I:850. \nolinkurl{fixes/MISPRINTS_7I.json#850}\par

\par\medskip\noindent\hypertarget{app-F7I-IX-095.1}{}\textbf{F7I-IX-095.1}\quad p. 419; inherited corrected-layer difference.\par
\noindent Imprimé : {\ttfamily\footnotesize g\allowbreak{} \textbackslash{}longmapsto\allowbreak{} }\par
\noindent Corrigé : \emph{(vide)}\par
\noindent Source : ADOPT-F7I-IX-095.1. \nolinkurl{fr/sga7i/diplomatic/expose_IX_body.tex}; \nolinkurl{fr/sga7i/current/expose_IX_body.tex}\par

\par\medskip\noindent\hypertarget{app-F7I-IX-096.1}{}\textbf{F7I-IX-096.1}\quad p. 420; misprint\_pass.\par
\noindent Imprimé : {\ttfamily\footnotesize a\allowbreak{}}\par
\noindent Corrigé : {\ttfamily\footnotesize à\allowbreak{}}\par
\noindent Source : 7I:657; 7I:656; 7I:851. \nolinkurl{fixes/FIXLOG_7I_CURRENT.jsonl#908}; \nolinkurl{fixes/MISPRINTS_7I.json#656}; \nolinkurl{fixes/MISPRINTS_7I.json#851}\par

\par\medskip\noindent\hypertarget{app-F7I-IX-096.2}{}\textbf{F7I-IX-096.2}\quad p. 420; misprint\_in\_print.\par
\noindent Imprimé : {\ttfamily\footnotesize interprèter\allowbreak{}}\par
\noindent Corrigé : {\ttfamily\footnotesize interpréter\allowbreak{}}\par
\noindent Source : FIX-7I-D-0672. \nolinkurl{fixes/FIXLOG_7I.jsonl#672}\par

\par\medskip\noindent\hypertarget{app-F7I-IX-097.1}{}\textbf{F7I-IX-097.1}\quad p. 420; printed misprint.\par
\noindent Imprimé : {\ttfamily\footnotesize )\allowbreak{}}\par
\noindent Corrigé : \emph{(vide)}\par
\noindent Source : 7I:853. \nolinkurl{fixes/MISPRINTS_7I.json#853}\par

\par\medskip\noindent\hypertarget{app-F7I-IX-097.2}{}\textbf{F7I-IX-097.2}\quad p. 420; misprint\_in\_print.\par
\noindent Imprimé : \emph{(vide)}\par
\noindent Corrigé : {\ttfamily\footnotesize '\allowbreak{}}\par
\noindent Source : FIX-7I-D-0618; 7I:854. \nolinkurl{fixes/FIXLOG_7I.jsonl#618}; \nolinkurl{fixes/MISPRINTS_7I.json#854}\par

\par\medskip\noindent\hypertarget{app-F7I-IX-097.3}{}\textbf{F7I-IX-097.3}\quad p. 420; misprint\_pass.\par
\noindent Imprimé : \emph{(vide)}\par
\noindent Corrigé : {\ttfamily\footnotesize \textbackslash{}check\{\allowbreak{}}\par
\noindent Source : 7I:660; 7I:659. \nolinkurl{fixes/FIXLOG_7I_CURRENT.jsonl#909}; \nolinkurl{fixes/MISPRINTS_7I.json#659}\par

\par\medskip\noindent\hypertarget{app-F7I-IX-097.4}{}\textbf{F7I-IX-097.4}\quad p. 420; misprint\_pass.\par
\noindent Imprimé : \emph{(vide)}\par
\noindent Corrigé : {\ttfamily\footnotesize \}\allowbreak{}}\par
\noindent Source : 7I:660; 7I:659. \nolinkurl{fixes/FIXLOG_7I_CURRENT.jsonl#909}; \nolinkurl{fixes/MISPRINTS_7I.json#659}\par

\par\medskip\noindent\hypertarget{app-F7I-IX-098.1}{}\textbf{F7I-IX-098.1}\quad p. 421; misprint\_pass.\par
\noindent Imprimé : \emph{(vide)}\par
\noindent Corrigé : {\ttfamily\footnotesize (\allowbreak{}}\par
\noindent Source : 7I:661; 7I:660; 7I:856. \nolinkurl{fixes/FIXLOG_7I_CURRENT.jsonl#910}; \nolinkurl{fixes/MISPRINTS_7I.json#660}; \nolinkurl{fixes/MISPRINTS_7I.json#856}\par

\par\medskip\noindent\hypertarget{app-F7I-IX-098.2}{}\textbf{F7I-IX-098.2}\quad p. 421; designated-witness attribution.\par
\noindent Imprimé : {\ttfamily\footnotesize I\allowbreak{}}\par
\noindent Corrigé : {\ttfamily\footnotesize \{\textbackslash{}pi\_o\}\allowbreak{}}\par
\noindent Source : WITNESS-F7I-IX-098.2. \nolinkurl{fr/sga7i/diplomatic/expose_IX_body.tex}; \nolinkurl{fr/sga7i/current/expose_IX_body.tex}\par

\par\medskip\noindent\hypertarget{app-F7I-IX-099.1}{}\textbf{F7I-IX-099.1}\quad p. 421; misprint\_pass.\par
\noindent Imprimé : {\ttfamily\footnotesize munie\allowbreak{}}\par
\noindent Corrigé : {\ttfamily\footnotesize munis\allowbreak{}}\par
\noindent Source : 7I:663; 7I:662; 7I:857. \nolinkurl{fixes/FIXLOG_7I_CURRENT.jsonl#911}; \nolinkurl{fixes/MISPRINTS_7I.json#662}; \nolinkurl{fixes/MISPRINTS_7I.json#857}\par

\par\medskip\noindent\hypertarget{app-F7I-IX-100.1}{}\textbf{F7I-IX-100.1}\quad p. 421; printed misprint.\par
\noindent Imprimé : {\ttfamily\footnotesize P\allowbreak{}}\par
\noindent Corrigé : {\ttfamily\footnotesize Q\allowbreak{}}\par
\noindent Source : 7I:661; 7I:858. \nolinkurl{fixes/MISPRINTS_7I.json#661}; \nolinkurl{fixes/MISPRINTS_7I.json#858}\par

\par\medskip\noindent\hypertarget{app-F7I-IX-100.2}{}\textbf{F7I-IX-100.2}\quad p. 421; printed misprint.\par
\noindent Imprimé : {\ttfamily\footnotesize Q\allowbreak{}}\par
\noindent Corrigé : {\ttfamily\footnotesize P\allowbreak{}}\par
\noindent Source : 7I:661; 7I:858. \nolinkurl{fixes/MISPRINTS_7I.json#661}; \nolinkurl{fixes/MISPRINTS_7I.json#858}\par

\par\medskip\noindent\hypertarget{app-F7I-IX-101.1}{}\textbf{F7I-IX-101.1}\quad p. 422; misprint\_pass.\par
\noindent Imprimé : {\ttfamily\footnotesize appelerons\allowbreak{}}\par
\noindent Corrigé : {\ttfamily\footnotesize appellerons\allowbreak{}}\par
\noindent Source : 7I:664; 7I:663; 7I:859. \nolinkurl{fixes/FIXLOG_7I_CURRENT.jsonl#912}; \nolinkurl{fixes/MISPRINTS_7I.json#663}; \nolinkurl{fixes/MISPRINTS_7I.json#859}\par

\par\medskip\noindent\hypertarget{app-F7I-IX-102.1}{}\textbf{F7I-IX-102.1}\quad p. 423; printed misprint.\par
\noindent Imprimé : {\ttfamily\footnotesize F\allowbreak{}}\par
\noindent Corrigé : {\ttfamily\footnotesize E\allowbreak{}}\par
\noindent Source : 7I:665; 7I:861. \nolinkurl{fixes/MISPRINTS_7I.json#665}; \nolinkurl{fixes/MISPRINTS_7I.json#861}\par

\par\medskip\noindent\hypertarget{app-F7I-IX-102.2}{}\textbf{F7I-IX-102.2}\quad p. 423; printed misprint.\par
\noindent Imprimé : {\ttfamily\footnotesize F\allowbreak{}}\par
\noindent Corrigé : {\ttfamily\footnotesize E\allowbreak{}}\par
\noindent Source : 7I:666; 7I:862. \nolinkurl{fixes/MISPRINTS_7I.json#666}; \nolinkurl{fixes/MISPRINTS_7I.json#862}\par

\par\medskip\noindent\hypertarget{app-F7I-IX-103.1}{}\textbf{F7I-IX-103.1}\quad p. 425; printed misprint.\par
\noindent Imprimé : {\ttfamily\footnotesize d\allowbreak{}}\par
\noindent Corrigé : {\ttfamily\footnotesize c\allowbreak{}}\par
\noindent Source : 7I:667; 7I:863. \nolinkurl{fixes/MISPRINTS_7I.json#667}; \nolinkurl{fixes/MISPRINTS_7I.json#863}\par

\par\medskip\noindent\hypertarget{app-F7I-IX-103.2}{}\textbf{F7I-IX-103.2}\quad p. 425; misprint\_pass.\par
\noindent Imprimé : {\ttfamily\footnotesize décrit\allowbreak{}}\par
\noindent Corrigé : {\ttfamily\footnotesize décrite\allowbreak{}}\par
\noindent Source : 7I:669. \nolinkurl{fixes/FIXLOG_7I_CURRENT.jsonl#913}\par

\par\medskip\noindent\hypertarget{app-F7I-IX-103.3}{}\textbf{F7I-IX-103.3}\quad p. 425; misprint\_pass.\par
\noindent Imprimé : {\ttfamily\footnotesize donc\allowbreak{}}\par
\noindent Corrigé : {\ttfamily\footnotesize dont\allowbreak{}}\par
\noindent Source : 7I:670. \nolinkurl{fixes/FIXLOG_7I_CURRENT.jsonl#914}\par

\par\medskip\noindent\hypertarget{app-F7I-IX-103.4}{}\textbf{F7I-IX-103.4}\quad p. 425; printed misprint.\par
\noindent Imprimé : {\ttfamily\footnotesize R\allowbreak{}}\par
\noindent Corrigé : {\ttfamily\footnotesize Q\allowbreak{}}\par
\noindent Source : 7I:670; 7I:866. \nolinkurl{fixes/MISPRINTS_7I.json#670}; \nolinkurl{fixes/MISPRINTS_7I.json#866}\par

\par\medskip\noindent\hypertarget{app-F7I-IX-103.5}{}\textbf{F7I-IX-103.5}\quad p. 425; printed misprint.\par
\noindent Imprimé : {\ttfamily\footnotesize Q\allowbreak{}}\par
\noindent Corrigé : {\ttfamily\footnotesize P\allowbreak{}}\par
\noindent Source : 7I:670; 7I:866. \nolinkurl{fixes/MISPRINTS_7I.json#670}; \nolinkurl{fixes/MISPRINTS_7I.json#866}\par

\par\medskip\noindent\hypertarget{app-F7I-IX-104.1}{}\textbf{F7I-IX-104.1}\quad p. 425; printed misprint.\par
\noindent Imprimé : \emph{(vide)}\par
\noindent Corrigé : {\ttfamily\footnotesize '\allowbreak{}}\par
\noindent Source : 7I:867. \nolinkurl{fixes/MISPRINTS_7I.json#867}\par

\par\medskip\noindent\hypertarget{app-F7I-IX-105.1}{}\textbf{F7I-IX-105.1}\quad p. 426; designated-witness attribution.\par
\noindent Imprimé : \emph{(vide)}\par
\noindent Corrigé : {\ttfamily\footnotesize '\allowbreak{}}\par
\noindent Source : WITNESS-F7I-IX-105.1. \nolinkurl{fr/sga7i/diplomatic/expose_IX_body.tex}; \nolinkurl{fr/sga7i/current/expose_IX_body.tex}\par

\par\medskip\noindent\hypertarget{app-F7I-IX-106.1}{}\textbf{F7I-IX-106.1}\quad p. 428; misprint\_pass.\par
\noindent Imprimé : \emph{(vide)}\par
\noindent Corrigé : {\ttfamily\footnotesize )\allowbreak{}}\par
\noindent Source : 7I:674. \nolinkurl{fixes/FIXLOG_7I_CURRENT.jsonl#915}\par

\par\medskip\noindent\hypertarget{app-F7I-IX-106.2}{}\textbf{F7I-IX-106.2}\quad p. 428; misprint\_pass.\par
\noindent Imprimé : \emph{(vide)}\par
\noindent Corrigé : {\ttfamily\footnotesize \textbackslash{}check\{\allowbreak{}}\par
\noindent Source : 7I:675; 7I:674; 7I:871. \nolinkurl{fixes/FIXLOG_7I_CURRENT.jsonl#916}; \nolinkurl{fixes/MISPRINTS_7I.json#674}; \nolinkurl{fixes/MISPRINTS_7I.json#871}\par

\par\medskip\noindent\hypertarget{app-F7I-IX-106.3}{}\textbf{F7I-IX-106.3}\quad p. 428; misprint\_pass.\par
\noindent Imprimé : \emph{(vide)}\par
\noindent Corrigé : {\ttfamily\footnotesize \}\allowbreak{}}\par
\noindent Source : 7I:675; 7I:674; 7I:871. \nolinkurl{fixes/FIXLOG_7I_CURRENT.jsonl#916}; \nolinkurl{fixes/MISPRINTS_7I.json#674}; \nolinkurl{fixes/MISPRINTS_7I.json#871}\par

\par\medskip\noindent\hypertarget{app-F7I-IX-106.4}{}\textbf{F7I-IX-106.4}\quad p. 428; misprint\_pass.\par
\noindent Imprimé : {\ttfamily\footnotesize auite\allowbreak{}}\par
\noindent Corrigé : {\ttfamily\footnotesize suite\allowbreak{}}\par
\noindent Source : 7I:676; 7I:675; 7I:872. \nolinkurl{fixes/FIXLOG_7I_CURRENT.jsonl#917}; \nolinkurl{fixes/MISPRINTS_7I.json#675}; \nolinkurl{fixes/MISPRINTS_7I.json#872}\par

\par\medskip\noindent\hypertarget{app-F7I-IX-107.1}{}\textbf{F7I-IX-107.1}\quad p. 429; printed misprint.\par
\noindent Imprimé : {\ttfamily\footnotesize \textasciicircum{}\allowbreak{}}\par
\noindent Corrigé : {\ttfamily\footnotesize \_\allowbreak{}}\par
\noindent Source : 7I:676. \nolinkurl{fixes/MISPRINTS_7I.json#676}\par

\par\medskip\noindent\hypertarget{app-F7I-IX-107.2}{}\textbf{F7I-IX-107.2}\quad p. 429; misprint\_pass.\par
\noindent Imprimé : {\ttfamily\footnotesize oS\allowbreak{}}\par
\noindent Corrigé : {\ttfamily\footnotesize oU\allowbreak{}}\par
\noindent Source : 7I:678; 7I:677; 7I:874. \nolinkurl{fixes/FIXLOG_7I_CURRENT.jsonl#919}; \nolinkurl{fixes/MISPRINTS_7I.json#677}; \nolinkurl{fixes/MISPRINTS_7I.json#874}\par

\par\medskip\noindent\hypertarget{app-F7I-IX-108.1}{}\textbf{F7I-IX-108.1}\quad p. 429; misprint\_pass.\par
\noindent Imprimé : \emph{(vide)}\par
\noindent Corrigé : {\ttfamily\footnotesize )\allowbreak{}}\par
\noindent Source : 7I:679; 7I:678. \nolinkurl{fixes/FIXLOG_7I_CURRENT.jsonl#920}; \nolinkurl{fixes/MISPRINTS_7I.json#678}\par

\par\medskip\noindent\hypertarget{app-F7I-IX-109.1}{}\textbf{F7I-IX-109.1}\quad p. 431; printed misprint.\par
\noindent Imprimé : \emph{(vide)}\par
\noindent Corrigé : {\ttfamily\footnotesize Q,\allowbreak{}}\par
\noindent Source : 7I:881. \nolinkurl{fixes/MISPRINTS_7I.json#881}\par

\par\medskip\noindent\hypertarget{app-F7I-IX-109.2}{}\textbf{F7I-IX-109.2}\quad p. 431; printed misprint.\par
\noindent Imprimé : {\ttfamily\footnotesize ,Q\allowbreak{}}\par
\noindent Corrigé : \emph{(vide)}\par
\noindent Source : 7I:881. \nolinkurl{fixes/MISPRINTS_7I.json#881}\par

\par\medskip\noindent\hypertarget{app-F7I-IX-110.1}{}\textbf{F7I-IX-110.1}\quad p. 431; printed misprint.\par
\noindent Imprimé : \emph{(vide)}\par
\noindent Corrigé : {\ttfamily\footnotesize Q,\allowbreak{}}\par
\noindent Source : 7I:882. \nolinkurl{fixes/MISPRINTS_7I.json#882}\par

\par\medskip\noindent\hypertarget{app-F7I-IX-110.2}{}\textbf{F7I-IX-110.2}\quad p. 431; printed misprint.\par
\noindent Imprimé : {\ttfamily\footnotesize ,Q\allowbreak{}}\par
\noindent Corrigé : \emph{(vide)}\par
\noindent Source : 7I:882. \nolinkurl{fixes/MISPRINTS_7I.json#882}\par

\par\medskip\noindent\hypertarget{app-F7I-IX-110.3}{}\textbf{F7I-IX-110.3}\quad p. 431; printed misprint.\par
\noindent Imprimé : {\ttfamily\footnotesize P\allowbreak{}}\par
\noindent Corrigé : {\ttfamily\footnotesize Q\allowbreak{}}\par
\noindent Source : 7I:882. \nolinkurl{fixes/MISPRINTS_7I.json#882}\par

\par\medskip\noindent\hypertarget{app-F7I-IX-110.4}{}\textbf{F7I-IX-110.4}\quad p. 431; printed misprint.\par
\noindent Imprimé : {\ttfamily\footnotesize Q\allowbreak{}}\par
\noindent Corrigé : {\ttfamily\footnotesize P\allowbreak{}}\par
\noindent Source : 7I:882. \nolinkurl{fixes/MISPRINTS_7I.json#882}\par

\par\medskip\noindent\hypertarget{app-F7I-IX-111.1}{}\textbf{F7I-IX-111.1}\quad p. 432; printed misprint.\par
\noindent Imprimé : {\ttfamily\footnotesize 6\allowbreak{}}\par
\noindent Corrigé : {\ttfamily\footnotesize 7\allowbreak{}}\par
\noindent Source : 7I:883. \nolinkurl{fixes/MISPRINTS_7I.json#883}\par

\par\medskip\noindent\hypertarget{app-F7I-IX-112.1}{}\textbf{F7I-IX-112.1}\quad p. 433; printed misprint.\par
\noindent Imprimé : {\ttfamily\footnotesize donnée\allowbreak{}}\par
\noindent Corrigé : {\ttfamily\footnotesize donné\allowbreak{}}\par
\noindent Source : 7I:679. \nolinkurl{fixes/MISPRINTS_7I.json#679}\par

\par\medskip\noindent\hypertarget{app-F7I-IX-113.1}{}\textbf{F7I-IX-113.1}\quad p. 434; printed misprint.\par
\noindent Imprimé : {\ttfamily\footnotesize 6\allowbreak{}}\par
\noindent Corrigé : {\ttfamily\footnotesize 7\allowbreak{}}\par
\noindent Source : 7I:885. \nolinkurl{fixes/MISPRINTS_7I.json#885}\par

\par\medskip\noindent\hypertarget{app-F7I-IX-113.2}{}\textbf{F7I-IX-113.2}\quad p. 434; printed misprint.\par
\noindent Imprimé : {\ttfamily\footnotesize \_\textbackslash{}eta\allowbreak{}}\par
\noindent Corrigé : \emph{(vide)}\par
\noindent Source : 7I:680. \nolinkurl{fixes/MISPRINTS_7I.json#680}\par

\par\medskip\noindent\hypertarget{app-F7I-IX-114.1}{}\textbf{F7I-IX-114.1}\quad p. 435; printed misprint.\par
\noindent Imprimé : {\ttfamily\footnotesize n\allowbreak{}}\par
\noindent Corrigé : {\ttfamily\footnotesize S\allowbreak{}}\par
\noindent Source : 7I:681. \nolinkurl{fixes/MISPRINTS_7I.json#681}\par

\par\medskip\noindent\hypertarget{app-F7I-IX-114.2}{}\textbf{F7I-IX-114.2}\quad p. 435; punctuation.\par
\noindent Imprimé : {\ttfamily\footnotesize ,\allowbreak{}}\par
\noindent Corrigé : {\ttfamily\footnotesize .\allowbreak{}}\par
\noindent Source : FIX-7I-D-0626; FIX-7I-C-0541. \nolinkurl{fixes/FIXLOG_7I.jsonl#626}; \nolinkurl{fixes/FIXLOG_7I_CURRENT.jsonl#541}\par

\par\medskip\noindent\hypertarget{app-F7I-IX-115.1}{}\textbf{F7I-IX-115.1}\quad p. 435; misprint\_pass.\par
\noindent Imprimé : {\ttfamily\footnotesize 6\allowbreak{}}\par
\noindent Corrigé : {\ttfamily\footnotesize 5\allowbreak{}}\par
\noindent Source : 7I:890; 7I:889. \nolinkurl{fixes/FIXLOG_7I_CURRENT.jsonl#922}; \nolinkurl{fixes/MISPRINTS_7I.json#889}\par

\par\medskip\noindent\hypertarget{app-F7I-IX-116.1}{}\textbf{F7I-IX-116.1}\quad p. 436; designated-witness attribution.\par
\noindent Imprimé : {\ttfamily\footnotesize 6\allowbreak{}}\par
\noindent Corrigé : {\ttfamily\footnotesize 5\allowbreak{}}\par
\noindent Source : WITNESS-F7I-IX-116.1. \nolinkurl{fr/sga7i/diplomatic/expose_IX_body.tex}; \nolinkurl{fr/sga7i/current/expose_IX_body.tex}\par

\par\medskip\noindent\hypertarget{app-F7I-IX-117.1}{}\textbf{F7I-IX-117.1}\quad p. 437; printed misprint.\par
\noindent Imprimé : {\ttfamily\footnotesize F\allowbreak{}}\par
\noindent Corrigé : {\ttfamily\footnotesize E\allowbreak{}}\par
\noindent Source : 7I:683; 7I:877. \nolinkurl{fixes/MISPRINTS_7I.json#683}; \nolinkurl{fixes/MISPRINTS_7I.json#877}\par

\par\medskip\noindent\hypertarget{app-F7I-IX-117.2}{}\textbf{F7I-IX-117.2}\quad p. 437; inherited corrected-layer difference.\par
\noindent Imprimé : {\ttfamily\footnotesize 6\allowbreak{}}\par
\noindent Corrigé : {\ttfamily\footnotesize 5\allowbreak{}}\par
\noindent Source : ADOPT-F7I-IX-117.2. \nolinkurl{fr/sga7i/diplomatic/expose_IX_body.tex}; \nolinkurl{fr/sga7i/current/expose_IX_body.tex}\par

\par\medskip\noindent\hypertarget{app-F7I-IX-117.3}{}\textbf{F7I-IX-117.3}\quad p. 437; inherited corrected-layer difference.\par
\noindent Imprimé : {\ttfamily\footnotesize S\allowbreak{}}\par
\noindent Corrigé : {\ttfamily\footnotesize \textbackslash{}eta\allowbreak{}}\par
\noindent Source : ADOPT-F7I-IX-117.3. \nolinkurl{fr/sga7i/diplomatic/expose_IX_body.tex}; \nolinkurl{fr/sga7i/current/expose_IX_body.tex}\par

\par\medskip\noindent\hypertarget{app-F7I-IX-117.4}{}\textbf{F7I-IX-117.4}\quad p. 437; misprint\_pass.\par
\noindent Imprimé : \emph{(vide)}\par
\noindent Corrigé : {\ttfamily\footnotesize )\allowbreak{}}\par
\noindent Source : 7I:685; 7I:684; 7I:878. \nolinkurl{fixes/FIXLOG_7I_CURRENT.jsonl#923}; \nolinkurl{fixes/MISPRINTS_7I.json#684}; \nolinkurl{fixes/MISPRINTS_7I.json#878}\par

\par\medskip\noindent\hypertarget{app-F7I-IX-118.1}{}\textbf{F7I-IX-118.1}\quad p. 438; misprint\_pass.\par
\noindent Imprimé : {\ttfamily\footnotesize interprèter\allowbreak{}}\par
\noindent Corrigé : {\ttfamily\footnotesize interpréter\allowbreak{}}\par
\noindent Source : 7I:687; 7I:686; 7I:879. \nolinkurl{fixes/FIXLOG_7I_CURRENT.jsonl#924}; \nolinkurl{fixes/MISPRINTS_7I.json#686}; \nolinkurl{fixes/MISPRINTS_7I.json#879}\par

\par\medskip\noindent\hypertarget{app-F7I-IX-119.1}{}\textbf{F7I-IX-119.1}\quad p. 439; misprint\_pass.\par
\noindent Imprimé : \emph{(vide)}\par
\noindent Corrigé : {\ttfamily\footnotesize '\allowbreak{}}\par
\noindent Source : 7I:688; 7I:687; 7I:880. \nolinkurl{fixes/FIXLOG_7I_CURRENT.jsonl#925}; \nolinkurl{fixes/MISPRINTS_7I.json#687}; \nolinkurl{fixes/MISPRINTS_7I.json#880}\par

\par\medskip\noindent\hypertarget{app-F7I-IX-120.1}{}\textbf{F7I-IX-120.1}\quad p. 441; misprint\_pass.\par
\noindent Imprimé : {\ttfamily\footnotesize epéciales\allowbreak{}}\par
\noindent Corrigé : {\ttfamily\footnotesize spéciales\allowbreak{}}\par
\noindent Source : 7I:689; 7I:688; 7I:890. \nolinkurl{fixes/FIXLOG_7I_CURRENT.jsonl#926}; \nolinkurl{fixes/MISPRINTS_7I.json#688}; \nolinkurl{fixes/MISPRINTS_7I.json#890}\par

\par\medskip\noindent\hypertarget{app-F7I-IX-121.1}{}\textbf{F7I-IX-121.1}\quad p. 441; misprint\_pass.\par
\noindent Imprimé : {\ttfamily\footnotesize \textbackslash{}varphi\textasciicircum{}t\allowbreak{}}\par
\noindent Corrigé : {\ttfamily\footnotesize u\_\textbackslash{}ell\allowbreak{}}\par
\noindent Source : 7I:691; 7I:690; 7I:891. \nolinkurl{fixes/FIXLOG_7I_CURRENT.jsonl#928}; \nolinkurl{fixes/MISPRINTS_7I.json#690}; \nolinkurl{fixes/MISPRINTS_7I.json#891}\par

\par\medskip\noindent\hypertarget{app-F7I-IX-121.2}{}\textbf{F7I-IX-121.2}\quad p. 441; misprint\_pass.\par
\noindent Imprimé : {\ttfamily\footnotesize a\allowbreak{}}\par
\noindent Corrigé : {\ttfamily\footnotesize à\allowbreak{}}\par
\noindent Source : 7I:690; 7I:689; 7I:892. \nolinkurl{fixes/FIXLOG_7I_CURRENT.jsonl#927}; \nolinkurl{fixes/MISPRINTS_7I.json#689}; \nolinkurl{fixes/MISPRINTS_7I.json#892}\par

\par\medskip\noindent\hypertarget{app-F7I-IX-121.3}{}\textbf{F7I-IX-121.3}\quad p. 441; printed misprint.\par
\noindent Imprimé : \emph{(vide)}\par
\noindent Corrigé : {\ttfamily\footnotesize 1.\allowbreak{}}\par
\noindent Source : 7I:893. \nolinkurl{fixes/MISPRINTS_7I.json#893}\par

\par\medskip\noindent\hypertarget{app-F7I-IX-121.4}{}\textbf{F7I-IX-121.4}\quad p. 441; printed misprint.\par
\noindent Imprimé : {\ttfamily\footnotesize .1\allowbreak{}}\par
\noindent Corrigé : \emph{(vide)}\par
\noindent Source : 7I:893. \nolinkurl{fixes/MISPRINTS_7I.json#893}\par

\par\medskip\noindent\hypertarget{app-F7I-IX-122.1}{}\textbf{F7I-IX-122.1}\quad p. 442; misprint\_pass.\par
\noindent Imprimé : \emph{(vide)}\par
\noindent Corrigé : {\ttfamily\footnotesize \{\allowbreak{}}\par
\noindent Source : 7I:692; 7I:691. \nolinkurl{fixes/FIXLOG_7I_CURRENT.jsonl#929}; \nolinkurl{fixes/MISPRINTS_7I.json#691}\par

\par\medskip\noindent\hypertarget{app-F7I-IX-122.2}{}\textbf{F7I-IX-122.2}\quad p. 442; misprint\_pass.\par
\noindent Imprimé : \emph{(vide)}\par
\noindent Corrigé : {\ttfamily\footnotesize '\}\allowbreak{}}\par
\noindent Source : 7I:692; 7I:691. \nolinkurl{fixes/FIXLOG_7I_CURRENT.jsonl#929}; \nolinkurl{fixes/MISPRINTS_7I.json#691}\par

\par\medskip\noindent\hypertarget{app-F7I-IX-123.1}{}\textbf{F7I-IX-123.1}\quad p. 443; misprint\_pass.\par
\noindent Imprimé : {\ttfamily\footnotesize 2\allowbreak{}}\par
\noindent Corrigé : {\ttfamily\footnotesize 3\allowbreak{}}\par
\noindent Source : 7I:694. \nolinkurl{fixes/FIXLOG_7I_CURRENT.jsonl#931}\par

\par\medskip\noindent\hypertarget{app-F7I-IX-123.2}{}\textbf{F7I-IX-123.2}\quad p. 443; misprint\_pass.\par
\noindent Imprimé : {\ttfamily\footnotesize 3\allowbreak{}}\par
\noindent Corrigé : {\ttfamily\footnotesize 4\allowbreak{}}\par
\noindent Source : 7I:694. \nolinkurl{fixes/FIXLOG_7I_CURRENT.jsonl#931}\par

\par\medskip\noindent\hypertarget{app-F7I-IX-123.3}{}\textbf{F7I-IX-123.3}\quad p. 443; misprint\_pass.\par
\noindent Imprimé : {\ttfamily\footnotesize inmédiatement\allowbreak{}}\par
\noindent Corrigé : {\ttfamily\footnotesize immédiatement\allowbreak{}}\par
\noindent Source : 7I:693; 7I:692; 7I:896. \nolinkurl{fixes/FIXLOG_7I_CURRENT.jsonl#930}; \nolinkurl{fixes/MISPRINTS_7I.json#692}; \nolinkurl{fixes/MISPRINTS_7I.json#896}\par

\par\medskip\noindent\hypertarget{app-F7I-IX-124.1}{}\textbf{F7I-IX-124.1}\quad p. 444; misprint\_pass.\par
\noindent Imprimé : {\ttfamily\footnotesize Quand\allowbreak{}}\par
\noindent Corrigé : {\ttfamily\footnotesize Quant\allowbreak{}}\par
\noindent Source : 7I:695. \nolinkurl{fixes/FIXLOG_7I_CURRENT.jsonl#932}\par

\par\medskip\noindent\hypertarget{app-F7I-IX-124.2}{}\textbf{F7I-IX-124.2}\quad p. 444; misprint\_pass.\par
\noindent Imprimé : {\ttfamily\footnotesize appelera\allowbreak{}}\par
\noindent Corrigé : {\ttfamily\footnotesize appellera\allowbreak{}}\par
\noindent Source : 7I:696; 7I:695; 7I:898. \nolinkurl{fixes/FIXLOG_7I_CURRENT.jsonl#933}; \nolinkurl{fixes/MISPRINTS_7I.json#695}; \nolinkurl{fixes/MISPRINTS_7I.json#898}\par

\par\medskip\noindent\hypertarget{app-F7I-IX-125.1}{}\textbf{F7I-IX-125.1}\quad p. 445; misprint\_in\_print.\par
\noindent Imprimé : \emph{(vide)}\par
\noindent Corrigé : {\ttfamily\footnotesize \textbackslash{}mathbb\{\allowbreak{}}\par
\noindent Source : FIX-7I-D-0682. \nolinkurl{fixes/FIXLOG_7I.jsonl#682}\par

\par\medskip\noindent\hypertarget{app-F7I-IX-125.2}{}\textbf{F7I-IX-125.2}\quad p. 445; misprint\_in\_print.\par
\noindent Imprimé : \emph{(vide)}\par
\noindent Corrigé : {\ttfamily\footnotesize \}\allowbreak{}}\par
\noindent Source : FIX-7I-D-0682. \nolinkurl{fixes/FIXLOG_7I.jsonl#682}\par

\par\medskip\noindent\hypertarget{app-F7I-IX-126.1}{}\textbf{F7I-IX-126.1}\quad p. 445; printed misprint.\par
\noindent Imprimé : {\ttfamily\footnotesize k\allowbreak{}}\par
\noindent Corrigé : {\ttfamily\footnotesize K\allowbreak{}}\par
\noindent Source : 7I:696; 7I:900. \nolinkurl{fixes/MISPRINTS_7I.json#696}; \nolinkurl{fixes/MISPRINTS_7I.json#900}\par

\par\medskip\noindent\hypertarget{app-F7I-IX-127.1}{}\textbf{F7I-IX-127.1}\quad p. 446; misprint\_pass.\par
\noindent Imprimé : {\ttfamily\footnotesize géométrique\allowbreak{}}\par
\noindent Corrigé : {\ttfamily\footnotesize géométriquement\allowbreak{}}\par
\noindent Source : 7I:699. \nolinkurl{fixes/FIXLOG_7I_CURRENT.jsonl#934}\par

\par\medskip\noindent\hypertarget{app-F7I-IX-127.2}{}\textbf{F7I-IX-127.2}\quad p. 446; printed misprint.\par
\noindent Imprimé : {\ttfamily\footnotesize X\allowbreak{}}\par
\noindent Corrigé : {\ttfamily\footnotesize K\allowbreak{}}\par
\noindent Source : 7I:697; 7I:901. \nolinkurl{fixes/MISPRINTS_7I.json#697}; \nolinkurl{fixes/MISPRINTS_7I.json#901}\par

\par\medskip\noindent\hypertarget{app-F7I-IX-128.1}{}\textbf{F7I-IX-128.1}\quad p. 448; misprint\_pass.\par
\noindent Imprimé : {\ttfamily\footnotesize assymptotique\allowbreak{}}\par
\noindent Corrigé : {\ttfamily\footnotesize asymptotique\allowbreak{}}\par
\noindent Source : 7I:700. \nolinkurl{fixes/FIXLOG_7I_CURRENT.jsonl#935}\par

\par\medskip\noindent\hypertarget{app-F7I-IX-129.1}{}\textbf{F7I-IX-129.1}\quad p. 448; printed misprint.\par
\noindent Imprimé : {\ttfamily\footnotesize 1\allowbreak{}}\par
\noindent Corrigé : {\ttfamily\footnotesize \textbackslash{}ell\allowbreak{}}\par
\noindent Source : 7I:700; 7I:904. \nolinkurl{fixes/MISPRINTS_7I.json#700}; \nolinkurl{fixes/MISPRINTS_7I.json#904}\par

\par\medskip\noindent\hypertarget{app-F7I-IX-130.1}{}\textbf{F7I-IX-130.1}\quad p. 452; printed misprint.\par
\noindent Imprimé : {\ttfamily\footnotesize ,\allowbreak{}}\par
\noindent Corrigé : {\ttfamily\footnotesize .\allowbreak{}}\par
\noindent Source : 7I:905. \nolinkurl{fixes/MISPRINTS_7I.json#905}\par

\par\medskip\noindent\hypertarget{app-F7I-IX-131.1}{}\textbf{F7I-IX-131.1}\quad p. 452; misprint\_in\_print.\par
\noindent Imprimé : \emph{(1 espace(s), 0 retour(s) de ligne dans le témoin TeX)}\par
\noindent Corrigé : {\ttfamily\footnotesize \textbackslash{}mathbb\{\allowbreak{}}\par
\noindent Source : FIX-7I-D-0683. \nolinkurl{fixes/FIXLOG_7I.jsonl#683}\par

\par\medskip\noindent\hypertarget{app-F7I-IX-131.2}{}\textbf{F7I-IX-131.2}\quad p. 452; misprint\_in\_print.\par
\noindent Imprimé : \emph{(vide)}\par
\noindent Corrigé : {\ttfamily\footnotesize \}\allowbreak{}}\par
\noindent Source : FIX-7I-D-0683. \nolinkurl{fixes/FIXLOG_7I.jsonl#683}\par

\par\medskip\noindent\hypertarget{app-F7I-IX-132.1}{}\textbf{F7I-IX-132.1}\quad p. 454; printed misprint.\par
\noindent Imprimé : {\ttfamily\footnotesize 4\allowbreak{}}\par
\noindent Corrigé : {\ttfamily\footnotesize 3\allowbreak{}}\par
\noindent Source : 7I:907. \nolinkurl{fixes/MISPRINTS_7I.json#907}\par

\par\medskip\noindent\hypertarget{app-F7I-IX-133.1}{}\textbf{F7I-IX-133.1}\quad p. 455; misprint\_in\_print.\par
\noindent Imprimé : {\ttfamily\footnotesize ;\allowbreak{}}\par
\noindent Corrigé : {\ttfamily\footnotesize :\allowbreak{}}\par
\noindent Source : FIX-7I-D-0684. \nolinkurl{fixes/FIXLOG_7I.jsonl#684}\par

\par\medskip\noindent\hypertarget{app-F7I-IX-134.1}{}\textbf{F7I-IX-134.1}\quad p. 456; misprint\_pass.\par
\noindent Imprimé : {\ttfamily\footnotesize ,\allowbreak{}}\par
\noindent Corrigé : \emph{(vide)}\par
\noindent Source : 7I:704; 7I:703; 7I:909. \nolinkurl{fixes/FIXLOG_7I_CURRENT.jsonl#936}; \nolinkurl{fixes/MISPRINTS_7I.json#703}; \nolinkurl{fixes/MISPRINTS_7I.json#909}\par

\par\medskip\noindent\hypertarget{app-F7I-IX-135.1}{}\textbf{F7I-IX-135.1}\quad p. 456; printed misprint.\par
\noindent Imprimé : {\ttfamily\footnotesize \textbackslash{}Phi(\allowbreak{}}\par
\noindent Corrigé : {\ttfamily\footnotesize M\_\{\allowbreak{}}\par
\noindent Source : 7I:910. \nolinkurl{fixes/MISPRINTS_7I.json#910}\par

\par\medskip\noindent\hypertarget{app-F7I-IX-135.2}{}\textbf{F7I-IX-135.2}\quad p. 456; designated-witness attribution.\par
\noindent Imprimé : {\ttfamily\footnotesize )\allowbreak{}}\par
\noindent Corrigé : {\ttfamily\footnotesize \}\textbackslash{}otimes\allowbreak{}}\par
\noindent Source : WITNESS-F7I-IX-135.2. \nolinkurl{fr/sga7i/diplomatic/expose_IX_body.tex}; \nolinkurl{fr/sga7i/current/expose_IX_body.tex}\par

\par\medskip\noindent\hypertarget{app-F7I-IX-135.3}{}\textbf{F7I-IX-135.3}\quad p. 456; designated-witness attribution.\par
\noindent Imprimé : \emph{(vide)}\par
\noindent Corrigé : {\ttfamily\footnotesize D\_\{\textbackslash{}ell\}\allowbreak{}}\par
\noindent Source : WITNESS-F7I-IX-135.3. \nolinkurl{fr/sga7i/diplomatic/expose_IX_body.tex}; \nolinkurl{fr/sga7i/current/expose_IX_body.tex}\par

\par\medskip\noindent\hypertarget{app-F7I-IX-135.4}{}\textbf{F7I-IX-135.4}\quad p. 456; printed misprint.\par
\noindent Imprimé : {\ttfamily\footnotesize \textbackslash{}Phi\allowbreak{}}\par
\noindent Corrigé : {\ttfamily\footnotesize \textbackslash{}check\{M\}\allowbreak{}}\par
\noindent Source : 7I:910. \nolinkurl{fixes/MISPRINTS_7I.json#910}\par

\par\medskip\noindent\hypertarget{app-F7I-IX-135.5}{}\textbf{F7I-IX-135.5}\quad p. 456; printed misprint.\par
\noindent Imprimé : {\ttfamily\footnotesize (\allowbreak{}}\par
\noindent Corrigé : {\ttfamily\footnotesize \_\{\allowbreak{}}\par
\noindent Source : 7I:910. \nolinkurl{fixes/MISPRINTS_7I.json#910}\par

\par\medskip\noindent\hypertarget{app-F7I-IX-135.6}{}\textbf{F7I-IX-135.6}\quad p. 456; designated-witness attribution.\par
\noindent Imprimé : {\ttfamily\footnotesize )\allowbreak{}}\par
\noindent Corrigé : {\ttfamily\footnotesize \}\textbackslash{}otimes\allowbreak{}}\par
\noindent Source : WITNESS-F7I-IX-135.6. \nolinkurl{fr/sga7i/diplomatic/expose_IX_body.tex}; \nolinkurl{fr/sga7i/current/expose_IX_body.tex}\par

\par\medskip\noindent\hypertarget{app-F7I-IX-135.7}{}\textbf{F7I-IX-135.7}\quad p. 456; designated-witness attribution.\par
\noindent Imprimé : \emph{(vide)}\par
\noindent Corrigé : {\ttfamily\footnotesize D\_\{\textbackslash{}ell\}\allowbreak{}}\par
\noindent Source : WITNESS-F7I-IX-135.7. \nolinkurl{fr/sga7i/diplomatic/expose_IX_body.tex}; \nolinkurl{fr/sga7i/current/expose_IX_body.tex}\par

\par\medskip\noindent\hypertarget{app-F7I-IX-136.1}{}\textbf{F7I-IX-136.1}\quad p. 457; printed misprint.\par
\noindent Imprimé : {\ttfamily\footnotesize ,\allowbreak{}}\par
\noindent Corrigé : {\ttfamily\footnotesize .\allowbreak{}}\par
\noindent Source : 7I:704; 7I:911. \nolinkurl{fixes/MISPRINTS_7I.json#704}; \nolinkurl{fixes/MISPRINTS_7I.json#911}\par

\par\medskip\noindent\hypertarget{app-F7I-IX-136.2}{}\textbf{F7I-IX-136.2}\quad p. 457; printed misprint.\par
\noindent Imprimé : \emph{(vide)}\par
\noindent Corrigé : {\ttfamily\footnotesize '\allowbreak{}}\par
\noindent Source : 7I:705; 7I:912. \nolinkurl{fixes/MISPRINTS_7I.json#705}; \nolinkurl{fixes/MISPRINTS_7I.json#912}\par

\par\medskip\noindent\hypertarget{app-F7I-IX-136.3}{}\textbf{F7I-IX-136.3}\quad p. 457; misprint\_pass.\par
\noindent Imprimé : {\ttfamily\footnotesize où\allowbreak{}}\par
\noindent Corrigé : {\ttfamily\footnotesize ou\allowbreak{}}\par
\noindent Source : 7I:707. \nolinkurl{fixes/FIXLOG_7I_CURRENT.jsonl#937}\par

\par\medskip\noindent\hypertarget{app-F7I-IX-136.4}{}\textbf{F7I-IX-136.4}\quad p. 457; misprint\_pass.\par
\noindent Imprimé : {\ttfamily\footnotesize comptatibles\allowbreak{}}\par
\noindent Corrigé : {\ttfamily\footnotesize compatibles\allowbreak{}}\par
\noindent Source : 7I:708. \nolinkurl{fixes/FIXLOG_7I_CURRENT.jsonl#938}\par

\par\medskip\noindent\hypertarget{app-F7I-IX-137.1}{}\textbf{F7I-IX-137.1}\quad p. 460; printed misprint.\par
\noindent Imprimé : {\ttfamily\footnotesize 7\allowbreak{}}\par
\noindent Corrigé : {\ttfamily\footnotesize 6\allowbreak{}}\par
\noindent Source : 7I:709; 7I:916. \nolinkurl{fixes/MISPRINTS_7I.json#709}; \nolinkurl{fixes/MISPRINTS_7I.json#916}\par

\par\medskip\noindent\hypertarget{app-F7I-IX-137.2}{}\textbf{F7I-IX-137.2}\quad p. 460; misprint\_pass.\par
\noindent Imprimé : \emph{(vide)}\par
\noindent Corrigé : {\ttfamily\footnotesize '\allowbreak{}}\par
\noindent Source : 7I:711; 7I:710; 7I:917. \nolinkurl{fixes/FIXLOG_7I_CURRENT.jsonl#939}; \nolinkurl{fixes/MISPRINTS_7I.json#710}; \nolinkurl{fixes/MISPRINTS_7I.json#917}\par

\par\medskip\noindent\hypertarget{app-F7I-IX-137.3}{}\textbf{F7I-IX-137.3}\quad p. 460; inherited corrected-layer difference.\par
\noindent Imprimé : {\ttfamily\footnotesize 6\allowbreak{}}\par
\noindent Corrigé : {\ttfamily\footnotesize 7\allowbreak{}}\par
\noindent Source : ADOPT-F7I-IX-137.3. \nolinkurl{fr/sga7i/diplomatic/expose_IX_body.tex}; \nolinkurl{fr/sga7i/current/expose_IX_body.tex}\par

\par\medskip\noindent\hypertarget{app-F7I-IX-137.4}{}\textbf{F7I-IX-137.4}\quad p. 460; printed misprint.\par
\noindent Imprimé : \emph{(vide)}\par
\noindent Corrigé : {\ttfamily\footnotesize '\allowbreak{}}\par
\noindent Source : 7I:711; 7I:918. \nolinkurl{fixes/MISPRINTS_7I.json#711}; \nolinkurl{fixes/MISPRINTS_7I.json#918}\par

\par\medskip\noindent\hypertarget{app-F7I-IX-138.1}{}\textbf{F7I-IX-138.1}\quad p. 461; printed misprint.\par
\noindent Imprimé : {\ttfamily\footnotesize 7\allowbreak{}}\par
\noindent Corrigé : {\ttfamily\footnotesize 6\allowbreak{}}\par
\noindent Source : 7I:712; 7I:919. \nolinkurl{fixes/MISPRINTS_7I.json#712}; \nolinkurl{fixes/MISPRINTS_7I.json#919}\par

\par\medskip\noindent\hypertarget{app-F7I-IX-139.1}{}\textbf{F7I-IX-139.1}\quad p. 462; misprint\_pass.\par
\noindent Imprimé : {\ttfamily\footnotesize 8\allowbreak{}}\par
\noindent Corrigé : {\ttfamily\footnotesize 7\allowbreak{}}\par
\noindent Source : 7I:922; 7I:921. \nolinkurl{fixes/FIXLOG_7I_CURRENT.jsonl#940}; \nolinkurl{fixes/MISPRINTS_7I.json#921}\par

\par\medskip\noindent\hypertarget{app-F7I-IX-140.1}{}\textbf{F7I-IX-140.1}\quad p. 462; printed misprint.\par
\noindent Imprimé : {\ttfamily\footnotesize ,\allowbreak{} }\par
\noindent Corrigé : {\ttfamily\footnotesize /\allowbreak{}}\par
\noindent Source : 7I:920. \nolinkurl{fixes/MISPRINTS_7I.json#920}\par

\par\medskip\noindent\hypertarget{app-F7I-IX-141.1}{}\textbf{F7I-IX-141.1}\quad p. 463; printed misprint.\par
\noindent Imprimé : {\ttfamily\footnotesize \textbackslash{}Psi\allowbreak{}}\par
\noindent Corrigé : {\ttfamily\footnotesize \textbackslash{}Phi\allowbreak{}}\par
\noindent Source : 7I:922. \nolinkurl{fixes/MISPRINTS_7I.json#922}\par

\par\medskip\noindent\hypertarget{app-F7I-IX-142.1}{}\textbf{F7I-IX-142.1}\quad p. 464; inherited corrected-layer difference.\par
\noindent Imprimé : {\ttfamily\footnotesize  (cf.\allowbreak{} I\textbackslash{}hspace\{3em\})\allowbreak{}}\par
\noindent Corrigé : \emph{(vide)}\par
\noindent Source : ADOPT-F7I-IX-142.1. \nolinkurl{fr/sga7i/diplomatic/expose_IX_body.tex}; \nolinkurl{fr/sga7i/current/expose_IX_body.tex}\par

\par\medskip\noindent\hypertarget{app-F7I-IX-143.1}{}\textbf{F7I-IX-143.1}\quad p. 464; misprint\_pass.\par
\noindent Imprimé : {\ttfamily\footnotesize 2\allowbreak{}}\par
\noindent Corrigé : {\ttfamily\footnotesize 3\allowbreak{}}\par
\noindent Source : 7I:924. \nolinkurl{fixes/FIXLOG_7I_CURRENT.jsonl#941}\par

\par\medskip\noindent\hypertarget{app-F7I-IX-144.1}{}\textbf{F7I-IX-144.1}\quad p. 465; misprint\_pass.\par
\noindent Imprimé : \emph{(vide)}\par
\noindent Corrigé : {\ttfamily\footnotesize \_i\allowbreak{}}\par
\noindent Source : 7I:716. \nolinkurl{fixes/FIXLOG_7I_CURRENT.jsonl#942}\par

\par\medskip\noindent\hypertarget{app-F7I-IX-145.1}{}\textbf{F7I-IX-145.1}\quad p. 467; designated-witness attribution.\par
\noindent Imprimé : {\ttfamily\footnotesize S\allowbreak{}}\par
\noindent Corrigé : {\ttfamily\footnotesize X\allowbreak{}}\par
\noindent Source : WITNESS-F7I-IX-145.1. \nolinkurl{fr/sga7i/diplomatic/expose_IX_body.tex}; \nolinkurl{fr/sga7i/current/expose_IX_body.tex}\par

\par\medskip\noindent\hypertarget{app-F7I-IX-145.2}{}\textbf{F7I-IX-145.2}\quad p. 467; misprint\_pass.\par
\noindent Imprimé : \emph{(vide)}\par
\noindent Corrigé : {\ttfamily\footnotesize o\textasciicircum{}\allowbreak{}}\par
\noindent Source : 7I:717. \nolinkurl{fixes/FIXLOG_7I_CURRENT.jsonl#943}\par

\par\medskip\noindent\hypertarget{app-F7I-IX-145.3}{}\textbf{F7I-IX-145.3}\quad p. 467; misprint\_pass.\par
\noindent Imprimé : {\ttfamily\footnotesize \textasciicircum{}o\allowbreak{}}\par
\noindent Corrigé : \emph{(vide)}\par
\noindent Source : 7I:717. \nolinkurl{fixes/FIXLOG_7I_CURRENT.jsonl#943}\par

\par\medskip\noindent\hypertarget{app-F7I-IX-146.1}{}\textbf{F7I-IX-146.1}\quad p. 467; misprint\_pass.\par
\noindent Imprimé : {\ttfamily\footnotesize spéciale\allowbreak{}}\par
\noindent Corrigé : {\ttfamily\footnotesize générique\allowbreak{}}\par
\noindent Source : 7I:718; 7I:926. \nolinkurl{fixes/FIXLOG_7I_CURRENT.jsonl#944}; \nolinkurl{fixes/MISPRINTS_7I.json#926}\par

\par\medskip\noindent\hypertarget{app-F7I-IX-147.1}{}\textbf{F7I-IX-147.1}\quad p. 469; misprint\_pass.\par
\noindent Imprimé : {\ttfamily\footnotesize 12\allowbreak{}}\par
\noindent Corrigé : {\ttfamily\footnotesize 10\allowbreak{}}\par
\noindent Source : 7I:928. \nolinkurl{fixes/FIXLOG_7I_CURRENT.jsonl#945}\par

\par\medskip\noindent\hypertarget{app-F7I-IX-148.1}{}\textbf{F7I-IX-148.1}\quad p. 470; printed misprint.\par
\noindent Imprimé : {\ttfamily\footnotesize X\allowbreak{}}\par
\noindent Corrigé : {\ttfamily\footnotesize C\allowbreak{}}\par
\noindent Source : 7I:718; 7I:928. \nolinkurl{fixes/MISPRINTS_7I.json#718}; \nolinkurl{fixes/MISPRINTS_7I.json#928}\par

\par\medskip\noindent\hypertarget{app-F7I-IX-149.1}{}\textbf{F7I-IX-149.1}\quad p. 471; printed misprint.\par
\noindent Imprimé : {\ttfamily\footnotesize 2\allowbreak{}}\par
\noindent Corrigé : {\ttfamily\footnotesize 3\allowbreak{}}\par
\noindent Source : 7I:719; 7I:929. \nolinkurl{fixes/MISPRINTS_7I.json#719}; \nolinkurl{fixes/MISPRINTS_7I.json#929}\par

\par\medskip\noindent\hypertarget{app-F7I-IX-149.2}{}\textbf{F7I-IX-149.2}\quad p. 471; printed misprint.\par
\noindent Imprimé : {\ttfamily\footnotesize 2\allowbreak{}}\par
\noindent Corrigé : {\ttfamily\footnotesize 3\allowbreak{}}\par
\noindent Source : 7I:720; 7I:930. \nolinkurl{fixes/MISPRINTS_7I.json#720}; \nolinkurl{fixes/MISPRINTS_7I.json#930}\par

\par\medskip\noindent\hypertarget{app-F7I-IX-149.3}{}\textbf{F7I-IX-149.3}\quad p. 471; misprint\_pass.\par
\noindent Imprimé : {\ttfamily\footnotesize cellè\allowbreak{}}\par
\noindent Corrigé : {\ttfamily\footnotesize celle\allowbreak{}}\par
\noindent Source : 7I:722. \nolinkurl{fixes/FIXLOG_7I_CURRENT.jsonl#946}\par

\par\medskip\noindent\hypertarget{app-F7I-IX-150.1}{}\textbf{F7I-IX-150.1}\quad p. 472; misprint\_pass.\par
\noindent Imprimé : {\ttfamily\footnotesize arrêtes\allowbreak{}}\par
\noindent Corrigé : {\ttfamily\footnotesize arêtes\allowbreak{}}\par
\noindent Source : 7I:723. \nolinkurl{fixes/FIXLOG_7I_CURRENT.jsonl#947}\par

\par\medskip\noindent\hypertarget{app-F7I-IX-151.1}{}\textbf{F7I-IX-151.1}\quad p. 472; printed misprint.\par
\noindent Imprimé : {\ttfamily\footnotesize *\allowbreak{}}\par
\noindent Corrigé : {\ttfamily\footnotesize \}\textbackslash{}tag\{12.3.8\allowbreak{}}\par
\noindent Source : 7I:724. \nolinkurl{fixes/MISPRINTS_7I.json#724}\par

\par\medskip\noindent\hypertarget{app-F7I-IX-151.2}{}\textbf{F7I-IX-151.2}\quad p. 472; printed misprint.\par
\noindent Imprimé : {\ttfamily\footnotesize *\allowbreak{}}\par
\noindent Corrigé : \emph{(vide)}\par
\noindent Source : 7I:724. \nolinkurl{fixes/MISPRINTS_7I.json#724}\par

\par\medskip\noindent\hypertarget{app-F7I-IX-151.3}{}\textbf{F7I-IX-151.3}\quad p. 472; misprint\_pass.\par
\noindent Imprimé : {\ttfamily\footnotesize a\allowbreak{}}\par
\noindent Corrigé : {\ttfamily\footnotesize à\allowbreak{}}\par
\noindent Source : 7I:724; 7I:723; 7I:934. \nolinkurl{fixes/FIXLOG_7I_CURRENT.jsonl#948}; \nolinkurl{fixes/MISPRINTS_7I.json#723}; \nolinkurl{fixes/MISPRINTS_7I.json#934}\par

\par\medskip\noindent\hypertarget{app-F7I-IX-152.1}{}\textbf{F7I-IX-152.1}\quad p. 472; misprint\_pass.\par
\noindent Imprimé : {\ttfamily\footnotesize 3\allowbreak{}}\par
\noindent Corrigé : {\ttfamily\footnotesize 2\allowbreak{}}\par
\noindent Source : 7I:936. \nolinkurl{fixes/FIXLOG_7I_CURRENT.jsonl#949}\par

\par\medskip\noindent\hypertarget{app-F7I-IX-152.2}{}\textbf{F7I-IX-152.2}\quad p. 472; misprint\_in\_print.\par
\noindent Imprimé : {\ttfamily\footnotesize interprèter\allowbreak{}}\par
\noindent Corrigé : {\ttfamily\footnotesize interpréter\allowbreak{}}\par
\noindent Source : FIX-7I-D-0645; 7I:936. \nolinkurl{fixes/FIXLOG_7I.jsonl#645}; \nolinkurl{fixes/MISPRINTS_7I.json#936}\par

\par\medskip\noindent\hypertarget{app-F7I-IX-153.1}{}\textbf{F7I-IX-153.1}\quad p. 473; misprint\_pass.\par
\noindent Imprimé : {\ttfamily\footnotesize arrêtes\allowbreak{}}\par
\noindent Corrigé : {\ttfamily\footnotesize arêtes\allowbreak{}}\par
\noindent Source : 7I:727; 7I:726; 7I:937. \nolinkurl{fixes/FIXLOG_7I_CURRENT.jsonl#950}; \nolinkurl{fixes/MISPRINTS_7I.json#726}; \nolinkurl{fixes/MISPRINTS_7I.json#937}\par

\par\medskip\noindent\hypertarget{app-F7I-IX-153.2}{}\textbf{F7I-IX-153.2}\quad p. 473; misprint\_pass.\par
\noindent Imprimé : {\ttfamily\footnotesize arrêtes\allowbreak{}}\par
\noindent Corrigé : {\ttfamily\footnotesize arêtes\allowbreak{}}\par
\noindent Source : 7I:728. \nolinkurl{fixes/FIXLOG_7I_CURRENT.jsonl#951}\par

\par\medskip\noindent\hypertarget{app-F7I-IX-154.1}{}\textbf{F7I-IX-154.1}\quad p. 473; printed misprint.\par
\noindent Imprimé : \emph{(vide)}\par
\noindent Corrigé : {\ttfamily\footnotesize \_1\allowbreak{}}\par
\noindent Source : 7I:728; 7I:939. \nolinkurl{fixes/MISPRINTS_7I.json#728}; \nolinkurl{fixes/MISPRINTS_7I.json#939}\par

\par\medskip\noindent\hypertarget{app-F7I-IX-154.2}{}\textbf{F7I-IX-154.2}\quad p. 473; misprint\_pass.\par
\noindent Imprimé : \emph{(vide)}\par
\noindent Corrigé : {\ttfamily\footnotesize \_\textbackslash{}eta\allowbreak{}}\par
\noindent Source : 7I:941. \nolinkurl{fixes/FIXLOG_7I_CURRENT.jsonl#952}\par

\par\medskip\noindent\hypertarget{app-F7I-IX-155.1}{}\textbf{F7I-IX-155.1}\quad p. 474; misprint\_pass.\par
\noindent Imprimé : {\ttfamily\footnotesize billinéaire\allowbreak{}}\par
\noindent Corrigé : {\ttfamily\footnotesize bilinéaire\allowbreak{}}\par
\noindent Source : 7I:730. \nolinkurl{fixes/FIXLOG_7I_CURRENT.jsonl#953}\par

\par\medskip\noindent\hypertarget{app-F7I-IX-155.2}{}\textbf{F7I-IX-155.2}\quad p. 474; misprint\_pass.\par
\noindent Imprimé : {\ttfamily\footnotesize ambiguité\allowbreak{}}\par
\noindent Corrigé : {\ttfamily\footnotesize ambiguïté\allowbreak{}}\par
\noindent Source : 7I:943. \nolinkurl{fixes/FIXLOG_7I_CURRENT.jsonl#955}\par

\par\medskip\noindent\hypertarget{app-F7I-IX-155.3}{}\textbf{F7I-IX-155.3}\quad p. 474; misprint\_pass.\par
\noindent Imprimé : {\ttfamily\footnotesize standart\allowbreak{}}\par
\noindent Corrigé : {\ttfamily\footnotesize standard\allowbreak{}}\par
\noindent Source : 7I:731. \nolinkurl{fixes/FIXLOG_7I_CURRENT.jsonl#954}\par

\par\medskip\noindent\hypertarget{app-F7I-IX-156.1}{}\textbf{F7I-IX-156.1}\quad p. 475; misprint\_pass.\par
\noindent Imprimé : {\ttfamily\footnotesize ,\allowbreak{}}\par
\noindent Corrigé : {\ttfamily\footnotesize .\allowbreak{}}\par
\noindent Source : 7I:732. \nolinkurl{fixes/FIXLOG_7I_CURRENT.jsonl#956}\par

\par\medskip\noindent\hypertarget{app-F7I-IX-157.1}{}\textbf{F7I-IX-157.1}\quad p. 476; misprint\_pass.\par
\noindent Imprimé : {\ttfamily\footnotesize ,\allowbreak{}}\par
\noindent Corrigé : \emph{(vide)}\par
\noindent Source : 7I:733; 7I:732; 7I:956. \nolinkurl{fixes/FIXLOG_7I_CURRENT.jsonl#957}; \nolinkurl{fixes/MISPRINTS_7I.json#732}; \nolinkurl{fixes/MISPRINTS_7I.json#956}\par

\par\medskip\noindent\hypertarget{app-F7I-IX-157.2}{}\textbf{F7I-IX-157.2}\quad p. 476; misprint\_pass.\par
\noindent Imprimé : {\ttfamily\footnotesize sour\allowbreak{}}\par
\noindent Corrigé : {\ttfamily\footnotesize sous\allowbreak{}}\par
\noindent Source : 7I:734; 7I:733; 7I:957. \nolinkurl{fixes/FIXLOG_7I_CURRENT.jsonl#958}; \nolinkurl{fixes/MISPRINTS_7I.json#733}; \nolinkurl{fixes/MISPRINTS_7I.json#957}\par

\par\medskip\noindent\hypertarget{app-F7I-IX-158.1}{}\textbf{F7I-IX-158.1}\quad p. 477; misprint\_pass.\par
\noindent Imprimé : \emph{(vide)}\par
\noindent Corrigé : {\ttfamily\footnotesize )\allowbreak{}}\par
\noindent Source : 7I:735; 7I:734; 7I:958. \nolinkurl{fixes/FIXLOG_7I_CURRENT.jsonl#959}; \nolinkurl{fixes/MISPRINTS_7I.json#734}; \nolinkurl{fixes/MISPRINTS_7I.json#958}\par

\par\medskip\noindent\hypertarget{app-F7I-IX-159.1}{}\textbf{F7I-IX-159.1}\quad p. 478; printed misprint.\par
\noindent Imprimé : {\ttfamily\footnotesize F\allowbreak{}}\par
\noindent Corrigé : {\ttfamily\footnotesize f\allowbreak{}}\par
\noindent Source : 7I:959. \nolinkurl{fixes/MISPRINTS_7I.json#959}\par

\par\medskip\noindent\hypertarget{app-F7I-IX-159.2}{}\textbf{F7I-IX-159.2}\quad p. 478; misprint\_pass.\par
\noindent Imprimé : \emph{(vide)}\par
\noindent Corrigé : {\ttfamily\footnotesize \textbackslash{}underline\{\allowbreak{}}\par
\noindent Source : 7I:961. \nolinkurl{fixes/FIXLOG_7I_CURRENT.jsonl#960}\par

\par\medskip\noindent\hypertarget{app-F7I-IX-159.3}{}\textbf{F7I-IX-159.3}\quad p. 478; misprint\_pass.\par
\noindent Imprimé : \emph{(vide)}\par
\noindent Corrigé : {\ttfamily\footnotesize \}\allowbreak{}}\par
\noindent Source : 7I:961. \nolinkurl{fixes/FIXLOG_7I_CURRENT.jsonl#960}\par

\par\medskip\noindent\hypertarget{app-F7I-IX-159.4}{}\textbf{F7I-IX-159.4}\quad p. 478; misprint\_pass.\par
\noindent Imprimé : \emph{(vide)}\par
\noindent Corrigé : {\ttfamily\footnotesize \textbackslash{}underline\{\allowbreak{}}\par
\noindent Source : 7I:962; 7I:961. \nolinkurl{fixes/FIXLOG_7I_CURRENT.jsonl#961}; \nolinkurl{fixes/MISPRINTS_7I.json#961}\par

\par\medskip\noindent\hypertarget{app-F7I-IX-159.5}{}\textbf{F7I-IX-159.5}\quad p. 478; misprint\_pass.\par
\noindent Imprimé : \emph{(vide)}\par
\noindent Corrigé : {\ttfamily\footnotesize \}\allowbreak{}}\par
\noindent Source : 7I:962; 7I:961. \nolinkurl{fixes/FIXLOG_7I_CURRENT.jsonl#961}; \nolinkurl{fixes/MISPRINTS_7I.json#961}\par

\par\medskip\noindent\hypertarget{app-F7I-IX-160.1}{}\textbf{F7I-IX-160.1}\quad p. 480; printed misprint.\par
\noindent Imprimé : {\ttfamily\footnotesize B\allowbreak{}}\par
\noindent Corrigé : {\ttfamily\footnotesize P\allowbreak{}}\par
\noindent Source : 7I:962. \nolinkurl{fixes/MISPRINTS_7I.json#962}\par

\par\medskip\noindent\hypertarget{app-F7I-IX-160.2}{}\textbf{F7I-IX-160.2}\quad p. 480; misprint\_pass.\par
\noindent Imprimé : \emph{(vide)}\par
\noindent Corrigé : {\ttfamily\footnotesize \textbackslash{}underline\{\allowbreak{}}\par
\noindent Source : 7I:964; 7I:963. \nolinkurl{fixes/FIXLOG_7I_CURRENT.jsonl#962}; \nolinkurl{fixes/MISPRINTS_7I.json#963}\par

\par\medskip\noindent\hypertarget{app-F7I-IX-160.3}{}\textbf{F7I-IX-160.3}\quad p. 480; misprint\_pass.\par
\noindent Imprimé : \emph{(vide)}\par
\noindent Corrigé : {\ttfamily\footnotesize \}\allowbreak{}}\par
\noindent Source : 7I:964; 7I:963. \nolinkurl{fixes/FIXLOG_7I_CURRENT.jsonl#962}; \nolinkurl{fixes/MISPRINTS_7I.json#963}\par

\par\medskip\noindent\hypertarget{app-F7I-IX-161.1}{}\textbf{F7I-IX-161.1}\quad p. 480; misprint\_in\_print.\par
\noindent Imprimé : {\ttfamily\footnotesize \textbackslash{}underline\{\allowbreak{}}\par
\noindent Corrigé : \emph{(vide)}\par
\noindent Source : FIX-7I-D-0691. \nolinkurl{fixes/FIXLOG_7I.jsonl#691}\par

\par\medskip\noindent\hypertarget{app-F7I-IX-161.2}{}\textbf{F7I-IX-161.2}\quad p. 480; misprint\_in\_print.\par
\noindent Imprimé : {\ttfamily\footnotesize \}\allowbreak{}}\par
\noindent Corrigé : \emph{(vide)}\par
\noindent Source : FIX-7I-D-0691. \nolinkurl{fixes/FIXLOG_7I.jsonl#691}\par

\par\medskip\noindent\hypertarget{app-F7I-IX-161.3}{}\textbf{F7I-IX-161.3}\quad p. 480; designated-witness attribution.\par
\noindent Imprimé : \emph{(vide)}\par
\noindent Corrigé : {\ttfamily\footnotesize )\allowbreak{}}\par
\noindent Source : WITNESS-F7I-IX-161.3. \nolinkurl{fr/sga7i/diplomatic/expose_IX_body.tex}; \nolinkurl{fr/sga7i/current/expose_IX_body.tex}\par

\par\medskip\noindent\hypertarget{app-F7I-IX-162.1}{}\textbf{F7I-IX-162.1}\quad p. 481; misprint\_pass.\par
\noindent Imprimé : \emph{(vide)}\par
\noindent Corrigé : {\ttfamily\footnotesize )\allowbreak{}}\par
\noindent Source : 7I:740; 7I:739; 7I:945. \nolinkurl{fixes/FIXLOG_7I_CURRENT.jsonl#963}; \nolinkurl{fixes/MISPRINTS_7I.json#739}; \nolinkurl{fixes/MISPRINTS_7I.json#945}\par

\par\medskip\noindent\hypertarget{app-F7I-IX-163.1}{}\textbf{F7I-IX-163.1}\quad p. 481; misprint\_pass.\par
\noindent Imprimé : \emph{(vide)}\par
\noindent Corrigé : {\ttfamily\footnotesize \textbackslash{}underline\{\allowbreak{}}\par
\noindent Source : 7I:741; 7I:740; 7I:946. \nolinkurl{fixes/FIXLOG_7I_CURRENT.jsonl#964}; \nolinkurl{fixes/MISPRINTS_7I.json#740}; \nolinkurl{fixes/MISPRINTS_7I.json#946}\par

\par\medskip\noindent\hypertarget{app-F7I-IX-163.2}{}\textbf{F7I-IX-163.2}\quad p. 481; misprint\_pass.\par
\noindent Imprimé : \emph{(vide)}\par
\noindent Corrigé : {\ttfamily\footnotesize \}\allowbreak{}}\par
\noindent Source : 7I:741; 7I:740; 7I:946. \nolinkurl{fixes/FIXLOG_7I_CURRENT.jsonl#964}; \nolinkurl{fixes/MISPRINTS_7I.json#740}; \nolinkurl{fixes/MISPRINTS_7I.json#946}\par

\par\medskip\noindent\hypertarget{app-F7I-IX-164.1}{}\textbf{F7I-IX-164.1}\quad p. 482; printed misprint.\par
\noindent Imprimé : {\ttfamily\footnotesize x\allowbreak{}}\par
\noindent Corrigé : {\ttfamily\footnotesize X\allowbreak{}}\par
\noindent Source : 7I:742. \nolinkurl{fixes/MISPRINTS_7I.json#742}\par

\par\medskip\noindent\hypertarget{app-F7I-IX-165.1}{}\textbf{F7I-IX-165.1}\quad p. 482; printed misprint.\par
\noindent Imprimé : \emph{(vide)}\par
\noindent Corrigé : {\ttfamily\footnotesize \_\textbackslash{}lambda\allowbreak{}}\par
\noindent Source : 7I:743; 7I:949. \nolinkurl{fixes/MISPRINTS_7I.json#743}; \nolinkurl{fixes/MISPRINTS_7I.json#949}\par

\par\medskip\noindent\hypertarget{app-F7I-IX-165.2}{}\textbf{F7I-IX-165.2}\quad p. 482; designated-witness attribution.\par
\noindent Imprimé : \emph{(vide)}\par
\noindent Corrigé : {\ttfamily\footnotesize \textbackslash{}underline\{\allowbreak{}}\par
\noindent Source : WITNESS-F7I-IX-165.2. \nolinkurl{fr/sga7i/diplomatic/expose_IX_body.tex}; \nolinkurl{fr/sga7i/current/expose_IX_body.tex}\par

\par\medskip\noindent\hypertarget{app-F7I-IX-165.3}{}\textbf{F7I-IX-165.3}\quad p. 482; designated-witness attribution.\par
\noindent Imprimé : \emph{(vide)}\par
\noindent Corrigé : {\ttfamily\footnotesize \}\allowbreak{}}\par
\noindent Source : WITNESS-F7I-IX-165.3. \nolinkurl{fr/sga7i/diplomatic/expose_IX_body.tex}; \nolinkurl{fr/sga7i/current/expose_IX_body.tex}\par

\par\medskip\noindent\hypertarget{app-F7I-IX-165.4}{}\textbf{F7I-IX-165.4}\quad p. 482; designated-witness attribution.\par
\noindent Imprimé : \emph{(vide)}\par
\noindent Corrigé : {\ttfamily\footnotesize \textbackslash{}underline\{\allowbreak{}}\par
\noindent Source : WITNESS-F7I-IX-165.4. \nolinkurl{fr/sga7i/diplomatic/expose_IX_body.tex}; \nolinkurl{fr/sga7i/current/expose_IX_body.tex}\par

\par\medskip\noindent\hypertarget{app-F7I-IX-165.5}{}\textbf{F7I-IX-165.5}\quad p. 482; designated-witness attribution.\par
\noindent Imprimé : \emph{(vide)}\par
\noindent Corrigé : {\ttfamily\footnotesize \}\allowbreak{}}\par
\noindent Source : WITNESS-F7I-IX-165.5. \nolinkurl{fr/sga7i/diplomatic/expose_IX_body.tex}; \nolinkurl{fr/sga7i/current/expose_IX_body.tex}\par

\par\medskip\noindent\hypertarget{app-F7I-IX-165.6}{}\textbf{F7I-IX-165.6}\quad p. 482; misprint\_pass.\par
\noindent Imprimé : \emph{(vide)}\par
\noindent Corrigé : {\ttfamily\footnotesize \textbackslash{}underline\{\allowbreak{}}\par
\noindent Source : 7I:952. \nolinkurl{fixes/FIXLOG_7I_CURRENT.jsonl#965}\par

\par\medskip\noindent\hypertarget{app-F7I-IX-165.7}{}\textbf{F7I-IX-165.7}\quad p. 482; misprint\_pass.\par
\noindent Imprimé : \emph{(vide)}\par
\noindent Corrigé : {\ttfamily\footnotesize \}\allowbreak{}}\par
\noindent Source : 7I:952. \nolinkurl{fixes/FIXLOG_7I_CURRENT.jsonl#965}\par

\par\medskip\noindent\hypertarget{app-F7I-IX-166.1}{}\textbf{F7I-IX-166.1}\quad p. 483; printed misprint.\par
\noindent Imprimé : {\ttfamily\footnotesize \textbackslash{}lambda\allowbreak{}}\par
\noindent Corrigé : {\ttfamily\footnotesize y\allowbreak{}}\par
\noindent Source : 7I:952. \nolinkurl{fixes/MISPRINTS_7I.json#952}\par

\par\medskip\noindent\hypertarget{app-F7I-IX-167.1}{}\textbf{F7I-IX-167.1}\quad p. 484; designated-witness attribution.\par
\noindent Imprimé : \emph{(vide)}\par
\noindent Corrigé : {\ttfamily\footnotesize \textasciicircum{}\textbackslash{}lambda\allowbreak{}}\par
\noindent Source : WITNESS-F7I-IX-167.1. \nolinkurl{fr/sga7i/diplomatic/expose_IX_body.tex}; \nolinkurl{fr/sga7i/current/expose_IX_body.tex}\par

\par\medskip\noindent\hypertarget{app-F7I-IX-167.2}{}\textbf{F7I-IX-167.2}\quad p. 484; misprint\_pass.\par
\noindent Imprimé : {\ttfamily\footnotesize correspondants\allowbreak{}}\par
\noindent Corrigé : {\ttfamily\footnotesize correspondant\allowbreak{}}\par
\noindent Source : 7I:747; 7I:746; 7I:954. \nolinkurl{fixes/FIXLOG_7I_CURRENT.jsonl#967}; \nolinkurl{fixes/MISPRINTS_7I.json#746}; \nolinkurl{fixes/MISPRINTS_7I.json#954}\par

\par\medskip\noindent\hypertarget{app-F7I-IX-167.3}{}\textbf{F7I-IX-167.3}\quad p. 484; misprint\_pass.\par
\noindent Imprimé : {\ttfamily\footnotesize monodormie\allowbreak{}}\par
\noindent Corrigé : {\ttfamily\footnotesize monodromie\allowbreak{}}\par
\noindent Source : 7I:746; 7I:745; 7I:955. \nolinkurl{fixes/FIXLOG_7I_CURRENT.jsonl#966}; \nolinkurl{fixes/MISPRINTS_7I.json#745}; \nolinkurl{fixes/MISPRINTS_7I.json#955}\par

\par\medskip\noindent\hypertarget{app-F7I-IX-168.1}{}\textbf{F7I-IX-168.1}\quad p. 486; misprint\_pass.\par
\noindent Imprimé : {\ttfamily\footnotesize appelera\allowbreak{}}\par
\noindent Corrigé : {\ttfamily\footnotesize appellera\allowbreak{}}\par
\noindent Source : 7I:748; 7I:966. \nolinkurl{fixes/FIXLOG_7I_CURRENT.jsonl#968}; \nolinkurl{fixes/MISPRINTS_7I.json#966}\par

\par\medskip\noindent\hypertarget{app-F7I-IX-169.1}{}\textbf{F7I-IX-169.1}\quad p. 487; misprint\_pass.\par
\noindent Imprimé : {\ttfamily\footnotesize S\allowbreak{}}\par
\noindent Corrigé : {\ttfamily\footnotesize s\allowbreak{}}\par
\noindent Source : 7I:749. \nolinkurl{fixes/FIXLOG_7I_CURRENT.jsonl#969}\par

\par\medskip\noindent\hypertarget{app-F7I-IX-170.1}{}\textbf{F7I-IX-170.1}\quad p. 487; misprint\_pass.\par
\noindent Imprimé : {\ttfamily\footnotesize C\allowbreak{}}\par
\noindent Corrigé : {\ttfamily\footnotesize 0\allowbreak{}}\par
\noindent Source : 7I:750; 7I:749; 7I:968. \nolinkurl{fixes/FIXLOG_7I_CURRENT.jsonl#970}; \nolinkurl{fixes/MISPRINTS_7I.json#749}; \nolinkurl{fixes/MISPRINTS_7I.json#968}\par

\par\medskip\noindent\hypertarget{app-F7I-IX-171.1}{}\textbf{F7I-IX-171.1}\quad p. 488; misprint\_pass.\par
\noindent Imprimé : \emph{(vide)}\par
\noindent Corrigé : {\ttfamily\footnotesize \textbackslash{}underline\{\allowbreak{}}\par
\noindent Source : 7I:751. \nolinkurl{fixes/FIXLOG_7I_CURRENT.jsonl#971}\par

\par\medskip\noindent\hypertarget{app-F7I-IX-171.2}{}\textbf{F7I-IX-171.2}\quad p. 488; misprint\_pass.\par
\noindent Imprimé : \emph{(vide)}\par
\noindent Corrigé : {\ttfamily\footnotesize \}\allowbreak{}}\par
\noindent Source : 7I:751. \nolinkurl{fixes/FIXLOG_7I_CURRENT.jsonl#971}\par

\par\medskip\noindent\hypertarget{app-F7I-IX-171.3}{}\textbf{F7I-IX-171.3}\quad p. 488; misprint\_pass.\par
\noindent Imprimé : {\ttfamily\footnotesize s\allowbreak{}}\par
\noindent Corrigé : {\ttfamily\footnotesize S\allowbreak{}}\par
\noindent Source : 7I:752; 7I:970. \nolinkurl{fixes/FIXLOG_7I_CURRENT.jsonl#972}; \nolinkurl{fixes/MISPRINTS_7I.json#970}\par

\par\medskip\noindent\hypertarget{app-F7I-IX-171.4}{}\textbf{F7I-IX-171.4}\quad p. 488; misprint\_pass.\par
\noindent Imprimé : \emph{(vide)}\par
\noindent Corrigé : {\ttfamily\footnotesize \textbackslash{}underline\{\allowbreak{}}\par
\noindent Source : 7I:972; 7I:971. \nolinkurl{fixes/FIXLOG_7I_CURRENT.jsonl#973}; \nolinkurl{fixes/MISPRINTS_7I.json#971}\par

\par\medskip\noindent\hypertarget{app-F7I-IX-171.5}{}\textbf{F7I-IX-171.5}\quad p. 488; misprint\_pass.\par
\noindent Imprimé : \emph{(vide)}\par
\noindent Corrigé : {\ttfamily\footnotesize \}\allowbreak{}}\par
\noindent Source : 7I:972; 7I:971. \nolinkurl{fixes/FIXLOG_7I_CURRENT.jsonl#973}; \nolinkurl{fixes/MISPRINTS_7I.json#971}\par

\par\medskip\noindent\hypertarget{app-F7I-IX-172.1}{}\textbf{F7I-IX-172.1}\quad p. 489; printed misprint.\par
\noindent Imprimé : {\ttfamily\footnotesize 2\allowbreak{}}\par
\noindent Corrigé : {\ttfamily\footnotesize 1\allowbreak{}}\par
\noindent Source : 7I:972. \nolinkurl{fixes/MISPRINTS_7I.json#972}\par

\par\medskip\noindent\hypertarget{app-F7I-IX-173.1}{}\textbf{F7I-IX-173.1}\quad p. 489; designated-witness attribution.\par
\noindent Imprimé : {\ttfamily\footnotesize \_\allowbreak{}}\par
\noindent Corrigé : {\ttfamily\footnotesize \textasciicircum{}\allowbreak{}}\par
\noindent Source : WITNESS-F7I-IX-173.1. \nolinkurl{fr/sga7i/diplomatic/expose_IX_body.tex}; \nolinkurl{fr/sga7i/current/expose_IX_body.tex}\par

\par\medskip\noindent\hypertarget{app-F7I-IX-173.2}{}\textbf{F7I-IX-173.2}\quad p. 489; misprint\_pass.\par
\noindent Imprimé : {\ttfamily\footnotesize 3\allowbreak{}}\par
\noindent Corrigé : {\ttfamily\footnotesize s\allowbreak{}}\par
\noindent Source : 7I:755; 7I:754; 7I:974. \nolinkurl{fixes/FIXLOG_7I_CURRENT.jsonl#974}; \nolinkurl{fixes/MISPRINTS_7I.json#754}; \nolinkurl{fixes/MISPRINTS_7I.json#974}\par

\par\medskip\noindent\hypertarget{app-F7I-IX-173.3}{}\textbf{F7I-IX-173.3}\quad p. 489; misprint\_pass.\par
\noindent Imprimé : {\ttfamily\footnotesize nul\allowbreak{}}\par
\noindent Corrigé : {\ttfamily\footnotesize nulle\allowbreak{}}\par
\noindent Source : 7I:756; 7I:975. \nolinkurl{fixes/FIXLOG_7I_CURRENT.jsonl#975}; \nolinkurl{fixes/MISPRINTS_7I.json#975}\par

\par\medskip\noindent\hypertarget{app-F7I-IX-174.1}{}\textbf{F7I-IX-174.1}\quad p. 491; misprint\_pass.\par
\noindent Imprimé : {\ttfamily\footnotesize S\allowbreak{}}\par
\noindent Corrigé : {\ttfamily\footnotesize s\allowbreak{}}\par
\noindent Source : 7I:757. \nolinkurl{fixes/FIXLOG_7I_CURRENT.jsonl#976}\par

\par\medskip\noindent\hypertarget{app-F7I-IX-175.1}{}\textbf{F7I-IX-175.1}\quad p. 492; misprint\_pass.\par
\noindent Imprimé : {\ttfamily\footnotesize contruction\allowbreak{}}\par
\noindent Corrigé : {\ttfamily\footnotesize construction\allowbreak{}}\par
\noindent Source : 7I:758; 7I:978. \nolinkurl{fixes/FIXLOG_7I_CURRENT.jsonl#977}; \nolinkurl{fixes/MISPRINTS_7I.json#978}\par

\par\medskip\noindent\hypertarget{app-F7I-IX-176.1}{}\textbf{F7I-IX-176.1}\quad p. 493; misprint\_pass.\par
\noindent Imprimé : {\ttfamily\footnotesize \textbackslash{}mathbb\allowbreak{}}\par
\noindent Corrigé : {\ttfamily\footnotesize \textbackslash{}underline\allowbreak{}}\par
\noindent Source : 7I:759; 7I:979. \nolinkurl{fixes/FIXLOG_7I_CURRENT.jsonl#978}; \nolinkurl{fixes/MISPRINTS_7I.json#979}\par

\par\medskip\noindent\hypertarget{app-F7I-IX-177.1}{}\textbf{F7I-IX-177.1}\quad p. 494; inherited corrected-layer difference.\par
\noindent Imprimé : \emph{(vide)}\par
\noindent Corrigé : {\ttfamily\footnotesize -\allowbreak{}}\par
\noindent Source : ADOPT-F7I-IX-177.1. \nolinkurl{fr/sga7i/diplomatic/expose_IX_body.tex}; \nolinkurl{fr/sga7i/current/expose_IX_body.tex}\par

\par\medskip\noindent\hypertarget{app-F7I-IX-177.2}{}\textbf{F7I-IX-177.2}\quad p. 494; inherited corrected-layer difference.\par
\noindent Imprimé : {\ttfamily\footnotesize -\allowbreak{}}\par
\noindent Corrigé : \emph{(vide)}\par
\noindent Source : ADOPT-F7I-IX-177.2. \nolinkurl{fr/sga7i/diplomatic/expose_IX_body.tex}; \nolinkurl{fr/sga7i/current/expose_IX_body.tex}\par

\par\medskip\noindent\hypertarget{app-F7I-IX-178.1}{}\textbf{F7I-IX-178.1}\quad p. 495; misprint\_pass.\par
\noindent Imprimé : {\ttfamily\footnotesize \textbackslash{}Lambda\allowbreak{}}\par
\noindent Corrigé : {\ttfamily\footnotesize A\allowbreak{}}\par
\noindent Source : 7I:760; 7I:759; 7I:981. \nolinkurl{fixes/FIXLOG_7I_CURRENT.jsonl#979}; \nolinkurl{fixes/MISPRINTS_7I.json#759}; \nolinkurl{fixes/MISPRINTS_7I.json#981}\par

\par\medskip\noindent\hypertarget{app-F7I-IX-179.1}{}\textbf{F7I-IX-179.1}\quad p. 496; inherited corrected-layer difference.\par
\noindent Imprimé : \emph{(vide)}\par
\noindent Corrigé : {\ttfamily\footnotesize \textasciicircum{}\textbackslash{}natural\allowbreak{}}\par
\noindent Source : ADOPT-F7I-IX-179.1. \nolinkurl{fr/sga7i/diplomatic/expose_IX_body.tex}; \nolinkurl{fr/sga7i/current/expose_IX_body.tex}\par

\par\medskip\noindent\hypertarget{app-F7I-IX-179.2}{}\textbf{F7I-IX-179.2}\quad p. 496; misprint\_pass.\par
\noindent Imprimé : \emph{(vide)}\par
\noindent Corrigé : {\ttfamily\footnotesize \textbackslash{}mathcal\{E\}\allowbreak{}}\par
\noindent Source : 7I:761; 7I:982. \nolinkurl{fixes/FIXLOG_7I_CURRENT.jsonl#980}; \nolinkurl{fixes/MISPRINTS_7I.json#982}\par

\par\medskip\noindent\hypertarget{app-F7I-IX-180.1}{}\textbf{F7I-IX-180.1}\quad p. 497; misprint\_pass.\par
\noindent Imprimé : {\ttfamily\footnotesize défini\allowbreak{}}\par
\noindent Corrigé : {\ttfamily\footnotesize définie\allowbreak{}}\par
\noindent Source : 7I:984; 7I:983. \nolinkurl{fixes/FIXLOG_7I_CURRENT.jsonl#1007}; \nolinkurl{fixes/MISPRINTS_7I.json#983}\par

\par\medskip\noindent\hypertarget{app-F7I-IX-181.1}{}\textbf{F7I-IX-181.1}\quad p. 498; misprint\_pass.\par
\noindent Imprimé : {\ttfamily\footnotesize servie\allowbreak{}}\par
\noindent Corrigé : {\ttfamily\footnotesize servi\allowbreak{}}\par
\noindent Source : 7I:763. \nolinkurl{fixes/FIXLOG_7I_CURRENT.jsonl#982}\par

\par\medskip\noindent\hypertarget{app-F7I-IX-181.2}{}\textbf{F7I-IX-181.2}\quad p. 498; misprint\_pass.\par
\noindent Imprimé : {\ttfamily\footnotesize impcrtant\allowbreak{}}\par
\noindent Corrigé : {\ttfamily\footnotesize important\allowbreak{}}\par
\noindent Source : 7I:762. \nolinkurl{fixes/FIXLOG_7I_CURRENT.jsonl#981}\par

\par\medskip\noindent\hypertarget{app-F7I-IX-182.1}{}\textbf{F7I-IX-182.1}\quad p. 500; punctuation.\par
\noindent Imprimé : \emph{(vide)}\par
\noindent Corrigé : {\ttfamily\footnotesize ,\allowbreak{}}\par
\noindent Source : FIX-7I-D-0656; FIX-7I-C-0565. \nolinkurl{fixes/FIXLOG_7I.jsonl#656}; \nolinkurl{fixes/FIXLOG_7I_CURRENT.jsonl#565}\par

\par\medskip\noindent\hypertarget{app-F7I-IX-183.1}{}\textbf{F7I-IX-183.1}\quad p. 501; printed misprint.\par
\noindent Imprimé : {\ttfamily\footnotesize C\allowbreak{}}\par
\noindent Corrigé : {\ttfamily\footnotesize O\allowbreak{}}\par
\noindent Source : 7I:993. \nolinkurl{fixes/MISPRINTS_7I.json#993}\par

\par\medskip\noindent\hypertarget{app-F7I-IX-183.2}{}\textbf{F7I-IX-183.2}\quad p. 501; misprint\_pass.\par
\noindent Imprimé : {\ttfamily\footnotesize C\allowbreak{}}\par
\noindent Corrigé : {\ttfamily\footnotesize O\allowbreak{}}\par
\noindent Source : 7I:765; 7I:994. \nolinkurl{fixes/FIXLOG_7I_CURRENT.jsonl#984}; \nolinkurl{fixes/MISPRINTS_7I.json#994}\par

\par\medskip\noindent\hypertarget{app-F7I-IX-184.1}{}\textbf{F7I-IX-184.1}\quad p. 502; misprint\_pass.\par
\noindent Imprimé : {\ttfamily\footnotesize Sévéri\allowbreak{}}\par
\noindent Corrigé : {\ttfamily\footnotesize Severi\allowbreak{}}\par
\noindent Source : 7I:998; 7I:997. \nolinkurl{fixes/FIXLOG_7I_CURRENT.jsonl#1009}; \nolinkurl{fixes/MISPRINTS_7I.json#997}\par

\par\medskip\noindent\hypertarget{app-F7I-IX-184.2}{}\textbf{F7I-IX-184.2}\quad p. 502; misprint\_pass.\par
\noindent Imprimé : {\ttfamily\footnotesize standarts\allowbreak{}}\par
\noindent Corrigé : {\ttfamily\footnotesize standards\allowbreak{}}\par
\noindent Source : 7I:996; 7I:995. \nolinkurl{fixes/FIXLOG_7I_CURRENT.jsonl#1008}; \nolinkurl{fixes/MISPRINTS_7I.json#995}\par

\par\medskip\noindent\hypertarget{app-F7I-IX-184.3}{}\textbf{F7I-IX-184.3}\quad p. 502; printed misprint.\par
\noindent Imprimé : {\ttfamily\footnotesize S\allowbreak{}}\par
\noindent Corrigé : {\ttfamily\footnotesize \textbackslash{}hat\{\textbackslash{}underline\{O\}\}\allowbreak{}}\par
\noindent Source : 7I:996. \nolinkurl{fixes/MISPRINTS_7I.json#996}\par

\par\medskip\noindent\hypertarget{app-F7I-IX-185.1}{}\textbf{F7I-IX-185.1}\quad p. 503; printed misprint.\par
\noindent Imprimé : {\ttfamily\footnotesize entiers\allowbreak{}}\par
\noindent Corrigé : {\ttfamily\footnotesize entier\allowbreak{}}\par
\noindent Source : 7I:998. \nolinkurl{fixes/MISPRINTS_7I.json#998}\par

\par\medskip\noindent\hypertarget{app-F7I-IX-185.2}{}\textbf{F7I-IX-185.2}\quad p. 503; misprint\_in\_print.\par
\noindent Imprimé : {\ttfamily\footnotesize C\allowbreak{}}\par
\noindent Corrigé : {\ttfamily\footnotesize 0\allowbreak{}}\par
\noindent Source : FIX-7I-D-0697. \nolinkurl{fixes/FIXLOG_7I.jsonl#697}\par

\par\medskip\noindent\hypertarget{app-F7I-IX-185.3}{}\textbf{F7I-IX-185.3}\quad p. 503; misprint\_pass.\par
\noindent Imprimé : {\ttfamily\footnotesize standart\allowbreak{}}\par
\noindent Corrigé : {\ttfamily\footnotesize standard\allowbreak{}}\par
\noindent Source : 7I:1002; 7I:1001. \nolinkurl{fixes/FIXLOG_7I_CURRENT.jsonl#1011}; \nolinkurl{fixes/MISPRINTS_7I.json#1001}\par

\par\medskip\noindent\hypertarget{app-F7I-IX-186.1}{}\textbf{F7I-IX-186.1}\quad p. 504; misprint\_pass.\par
\noindent Imprimé : \emph{(vide)}\par
\noindent Corrigé : {\ttfamily\footnotesize \textbackslash{}underline\{\allowbreak{}}\par
\noindent Source : 7I:766; 7I:765. \nolinkurl{fixes/FIXLOG_7I_CURRENT.jsonl#985}; \nolinkurl{fixes/MISPRINTS_7I.json#765}\par

\par\medskip\noindent\hypertarget{app-F7I-IX-186.2}{}\textbf{F7I-IX-186.2}\quad p. 504; misprint\_pass.\par
\noindent Imprimé : \emph{(vide)}\par
\noindent Corrigé : {\ttfamily\footnotesize \}\allowbreak{}}\par
\noindent Source : 7I:766; 7I:765. \nolinkurl{fixes/FIXLOG_7I_CURRENT.jsonl#985}; \nolinkurl{fixes/MISPRINTS_7I.json#765}\par

\par\medskip\noindent\hypertarget{app-F7I-IX-187.1}{}\textbf{F7I-IX-187.1}\quad p. 505; printed misprint.\par
\noindent Imprimé : \emph{(vide)}\par
\noindent Corrigé : {\ttfamily\footnotesize \textbackslash{}underline\{\allowbreak{}}\par
\noindent Source : 7I:1002. \nolinkurl{fixes/MISPRINTS_7I.json#1002}\par

\par\medskip\noindent\hypertarget{app-F7I-IX-187.2}{}\textbf{F7I-IX-187.2}\quad p. 505; printed misprint.\par
\noindent Imprimé : \emph{(vide)}\par
\noindent Corrigé : {\ttfamily\footnotesize \}\allowbreak{}}\par
\noindent Source : 7I:1002. \nolinkurl{fixes/MISPRINTS_7I.json#1002}\par

\par\medskip\noindent\hypertarget{app-F7I-IX-188.1}{}\textbf{F7I-IX-188.1}\quad p. 506; misprint\_pass.\par
\noindent Imprimé : {\ttfamily\footnotesize connu\allowbreak{}}\par
\noindent Corrigé : {\ttfamily\footnotesize connus\allowbreak{}}\par
\noindent Source : 7I:988. \nolinkurl{fixes/FIXLOG_7I_CURRENT.jsonl#1012}\par

\par\medskip\noindent\hypertarget{app-F7I-IX-188.2}{}\textbf{F7I-IX-188.2}\quad p. 506; misprint\_pass.\par
\noindent Imprimé : \emph{(2 espace(s), 0 retour(s) de ligne dans le témoin TeX)}\par
\noindent Corrigé : {\ttfamily\footnotesize -\allowbreak{}}\par
\noindent Source : 7I:768; 7I:767; 7I:988. \nolinkurl{fixes/FIXLOG_7I_CURRENT.jsonl#987}; \nolinkurl{fixes/MISPRINTS_7I.json#767}; \nolinkurl{fixes/MISPRINTS_7I.json#988}\par

\par\medskip\noindent\hypertarget{app-F7I-IX-189.1}{}\textbf{F7I-IX-189.1}\quad p. 507; misprint\_pass.\par
\noindent Imprimé : {\ttfamily\footnotesize P\allowbreak{}}\par
\noindent Corrigé : {\ttfamily\footnotesize p\allowbreak{}}\par
\noindent Source : 7I:769. \nolinkurl{fixes/FIXLOG_7I_CURRENT.jsonl#988}\par

\par\medskip\noindent\hypertarget{app-F7I-IX-190.1}{}\textbf{F7I-IX-190.1}\quad p. 508; misprint\_pass.\par
\noindent Imprimé : {\ttfamily\footnotesize ridige\allowbreak{}}\par
\noindent Corrigé : {\ttfamily\footnotesize rigide\allowbreak{}}\par
\noindent Source : 7I:770. \nolinkurl{fixes/FIXLOG_7I_CURRENT.jsonl#989}\par

\par\medskip\noindent\hypertarget{app-F7I-IX-190.2}{}\textbf{F7I-IX-190.2}\quad p. 508; misprint\_pass.\par
\noindent Imprimé : {\ttfamily\footnotesize analytiquessur\allowbreak{}}\par
\noindent Corrigé : {\ttfamily\footnotesize analytiques\allowbreak{} sur\allowbreak{}}\par
\noindent Source : 7I:770. \nolinkurl{fixes/FIXLOG_7I_CURRENT.jsonl#989}\par

\par\medskip\noindent\hypertarget{app-F7I-IX-191.1}{}\textbf{F7I-IX-191.1}\quad p. 508; misprint\_pass.\par
\noindent Imprimé : {\ttfamily\footnotesize dé\allowbreak{}}\par
\noindent Corrigé : {\ttfamily\footnotesize de\allowbreak{}}\par
\noindent Source : 7I:772. \nolinkurl{fixes/FIXLOG_7I_CURRENT.jsonl#990}\par

\par\medskip\noindent\hypertarget{app-F7I-IX-192.1}{}\textbf{F7I-IX-192.1}\quad p. 511; printed misprint.\par
\noindent Imprimé : {\ttfamily\footnotesize stucture\allowbreak{}}\par
\noindent Corrigé : {\ttfamily\footnotesize structure\allowbreak{}}\par
\noindent Source : 7I:1003. \nolinkurl{fixes/MISPRINTS_7I.json#1003}\par

\par\medskip\noindent\hypertarget{app-F7I-IX-192.2}{}\textbf{F7I-IX-192.2}\quad p. 511; misprint\_pass.\par
\noindent Imprimé : {\ttfamily\footnotesize réconstituer\allowbreak{}}\par
\noindent Corrigé : {\ttfamily\footnotesize reconstituer\allowbreak{}}\par
\noindent Source : 7I:774. \nolinkurl{fixes/FIXLOG_7I_CURRENT.jsonl#992}\par

\par\medskip\noindent\hypertarget{app-F7I-IX-193.1}{}\textbf{F7I-IX-193.1}\quad p. 512; misprint\_pass.\par
\noindent Imprimé : {\ttfamily\footnotesize réconstituer\allowbreak{}}\par
\noindent Corrigé : {\ttfamily\footnotesize reconstituer\allowbreak{}}\par
\noindent Source : 7I:775; 7I:1005. \nolinkurl{fixes/FIXLOG_7I_CURRENT.jsonl#993}; \nolinkurl{fixes/MISPRINTS_7I.json#1005}\par

\par\medskip\noindent\hypertarget{app-F7I-IX-194.1}{}\textbf{F7I-IX-194.1}\quad p. 512; printed misprint.\par
\noindent Imprimé : \emph{(vide)}\par
\noindent Corrigé : {\ttfamily\footnotesize  p'\allowbreak{} \textbackslash{};:\textbackslash{};\allowbreak{}}\par
\noindent Source : 7I:777; 7I:1006. \nolinkurl{fixes/MISPRINTS_7I.json#777}; \nolinkurl{fixes/MISPRINTS_7I.json#1006}\par

\par\medskip\noindent\hypertarget{app-F7I-IX-194.2}{}\textbf{F7I-IX-194.2}\quad p. 512; printed misprint.\par
\noindent Imprimé : \emph{(vide)}\par
\noindent Corrigé : {\ttfamily\footnotesize \textbackslash{}underline\{\allowbreak{}}\par
\noindent Source : 7I:1007. \nolinkurl{fixes/MISPRINTS_7I.json#1007}\par

\par\medskip\noindent\hypertarget{app-F7I-IX-194.3}{}\textbf{F7I-IX-194.3}\quad p. 512; printed misprint.\par
\noindent Imprimé : \emph{(vide)}\par
\noindent Corrigé : {\ttfamily\footnotesize \}\allowbreak{}}\par
\noindent Source : 7I:1007. \nolinkurl{fixes/MISPRINTS_7I.json#1007}\par

\par\medskip\noindent\hypertarget{app-F7I-IX-194.4}{}\textbf{F7I-IX-194.4}\quad p. 512; designated-witness attribution.\par
\noindent Imprimé : {\ttfamily\footnotesize 4\allowbreak{}}\par
\noindent Corrigé : {\ttfamily\footnotesize 14\allowbreak{}}\par
\noindent Source : WITNESS-F7I-IX-194.4. \nolinkurl{fr/sga7i/diplomatic/expose_IX_body.tex}; \nolinkurl{fr/sga7i/current/expose_IX_body.tex}\par

\par\medskip\noindent\hypertarget{app-F7I-IX-194.5}{}\textbf{F7I-IX-194.5}\quad p. 512; misprint\_pass.\par
\noindent Imprimé : {\ttfamily\footnotesize supposant\allowbreak{}}\par
\noindent Corrigé : {\ttfamily\footnotesize supposons\allowbreak{}}\par
\noindent Source : 7I:776. \nolinkurl{fixes/FIXLOG_7I_CURRENT.jsonl#994}\par

\par\medskip\noindent\hypertarget{app-F7I-IX-195.1}{}\textbf{F7I-IX-195.1}\quad p. 514; printed misprint.\par
\noindent Imprimé : \emph{(vide)}\par
\noindent Corrigé : {\ttfamily\footnotesize )\allowbreak{}}\par
\noindent Source : 7I:778; 7I:1010. \nolinkurl{fixes/MISPRINTS_7I.json#778}; \nolinkurl{fixes/MISPRINTS_7I.json#1010}\par

\par\medskip\noindent\hypertarget{app-F7I-IX-196.1}{}\textbf{F7I-IX-196.1}\quad p. 515; misprint\_pass.\par
\noindent Imprimé : \emph{(vide)}\par
\noindent Corrigé : {\ttfamily\footnotesize de\allowbreak{} }\par
\noindent Source : 7I:780. \nolinkurl{fixes/FIXLOG_7I_CURRENT.jsonl#995}\par

\par\medskip\noindent\hypertarget{app-F7I-IX-197.1}{}\textbf{F7I-IX-197.1}\quad p. 515; misprint\_pass.\par
\noindent Imprimé : {\ttfamily\footnotesize dans\allowbreak{}}\par
\noindent Corrigé : {\ttfamily\footnotesize dont\allowbreak{}}\par
\noindent Source : 7I:781. \nolinkurl{fixes/FIXLOG_7I_CURRENT.jsonl#996}\par

\par\medskip\noindent\hypertarget{app-F7I-IX-197.2}{}\textbf{F7I-IX-197.2}\quad p. 515; printed misprint.\par
\noindent Imprimé : {\ttfamily\footnotesize \textbackslash{}mathbb\{\allowbreak{}}\par
\noindent Corrigé : \emph{(vide)}\par
\noindent Source : 7I:1013. \nolinkurl{fixes/MISPRINTS_7I.json#1013}\par

\par\medskip\noindent\hypertarget{app-F7I-IX-197.3}{}\textbf{F7I-IX-197.3}\quad p. 515; printed misprint.\par
\noindent Imprimé : {\ttfamily\footnotesize \}\allowbreak{}}\par
\noindent Corrigé : \emph{(vide)}\par
\noindent Source : 7I:1013. \nolinkurl{fixes/MISPRINTS_7I.json#1013}\par

\par\medskip\noindent\hypertarget{app-F7I-IX-198.1}{}\textbf{F7I-IX-198.1}\quad p. 516; misprint\_pass.\par
\noindent Imprimé : {\ttfamily\footnotesize infiniments\allowbreak{}}\par
\noindent Corrigé : {\ttfamily\footnotesize infiniment\allowbreak{}}\par
\noindent Source : 7I:783; 7I:782. \nolinkurl{fixes/FIXLOG_7I_CURRENT.jsonl#997}; \nolinkurl{fixes/MISPRINTS_7I.json#782}\par

\par\medskip\noindent\hypertarget{app-F7I-IX-198.2}{}\textbf{F7I-IX-198.2}\quad p. 516; designated-witness attribution.\par
\noindent Imprimé : {\ttfamily\footnotesize \textbackslash{}longrightarrow\allowbreak{}}\par
\noindent Corrigé : {\ttfamily\footnotesize =\allowbreak{}}\par
\noindent Source : WITNESS-F7I-IX-198.2. \nolinkurl{fr/sga7i/diplomatic/expose_IX_body.tex}; \nolinkurl{fr/sga7i/current/expose_IX_body.tex}\par

\par\medskip\noindent\hypertarget{app-F7I-IX-199.1}{}\textbf{F7I-IX-199.1}\quad p. 517; printed misprint.\par
\noindent Imprimé : {\ttfamily\footnotesize 22\allowbreak{}}\par
\noindent Corrigé : {\ttfamily\footnotesize 2.2\allowbreak{}}\par
\noindent Source : 7I:1015. \nolinkurl{fixes/MISPRINTS_7I.json#1015}\par

\par\medskip\noindent\hypertarget{app-F7I-IX-200.1}{}\textbf{F7I-IX-200.1}\quad p. 518; misprint\_pass.\par
\noindent Imprimé : {\ttfamily\footnotesize G\allowbreak{}}\par
\noindent Corrigé : {\ttfamily\footnotesize A\allowbreak{}}\par
\noindent Source : 7I:784. \nolinkurl{fixes/FIXLOG_7I_CURRENT.jsonl#998}\par

\par\medskip\noindent\hypertarget{app-F7I-IX-200.2}{}\textbf{F7I-IX-200.2}\quad p. 518; misprint\_pass.\par
\noindent Imprimé : \emph{(vide)}\par
\noindent Corrigé : {\ttfamily\footnotesize '\allowbreak{}}\par
\noindent Source : 7I:785; 7I:784. \nolinkurl{fixes/FIXLOG_7I_CURRENT.jsonl#999}; \nolinkurl{fixes/MISPRINTS_7I.json#784}\par

\par\medskip\noindent\hypertarget{app-F7I-IX-200.3}{}\textbf{F7I-IX-200.3}\quad p. 518; printed misprint.\par
\noindent Imprimé : {\ttfamily\footnotesize 8\allowbreak{}}\par
\noindent Corrigé : {\ttfamily\footnotesize 5.3\allowbreak{}}\par
\noindent Source : 7I:1018. \nolinkurl{fixes/MISPRINTS_7I.json#1018}\par

\par\medskip\noindent\hypertarget{app-F7I-IX-200.4}{}\textbf{F7I-IX-200.4}\quad p. 518; printed misprint.\par
\noindent Imprimé : {\ttfamily\footnotesize 10\allowbreak{}}\par
\noindent Corrigé : {\ttfamily\footnotesize 5.5\allowbreak{}}\par
\noindent Source : 7I:1018. \nolinkurl{fixes/MISPRINTS_7I.json#1018}\par

\par\medskip\noindent\hypertarget{app-F7I-IX-201.1}{}\textbf{F7I-IX-201.1}\quad p. 519; printed misprint.\par
\noindent Imprimé : {\ttfamily\footnotesize seul\allowbreak{}}\par
\noindent Corrigé : {\ttfamily\footnotesize seule\allowbreak{}}\par
\noindent Source : 7I:787; 7I:1019. \nolinkurl{fixes/MISPRINTS_7I.json#787}; \nolinkurl{fixes/MISPRINTS_7I.json#1019}\par

\par\medskip\noindent\hypertarget{app-F7I-IX-201.2}{}\textbf{F7I-IX-201.2}\quad p. 519; misprint\_pass.\par
\noindent Imprimé : \emph{(vide)}\par
\noindent Corrigé : {\ttfamily\footnotesize de\allowbreak{} }\par
\noindent Source : 7I:789; 7I:1020. \nolinkurl{fixes/FIXLOG_7I_CURRENT.jsonl#1000}; \nolinkurl{fixes/MISPRINTS_7I.json#1020}\par

\par\medskip\noindent\hypertarget{app-F7I-IX-202.1}{}\textbf{F7I-IX-202.1}\quad p. 520; printed misprint.\par
\noindent Imprimé : {\ttfamily\footnotesize 4\allowbreak{}}\par
\noindent Corrigé : {\ttfamily\footnotesize 14\allowbreak{}}\par
\noindent Source : 7I:789. \nolinkurl{fixes/MISPRINTS_7I.json#789}\par

\par\medskip\noindent\hypertarget{app-F7I-IX-202.2}{}\textbf{F7I-IX-202.2}\quad p. 520; designated-witness attribution.\par
\noindent Imprimé : \emph{(vide)}\par
\noindent Corrigé : {\ttfamily\footnotesize .4\allowbreak{}}\par
\noindent Source : WITNESS-F7I-IX-202.2. \nolinkurl{fr/sga7i/diplomatic/expose_IX_body.tex}; \nolinkurl{fr/sga7i/current/expose_IX_body.tex}\par

\par\medskip\noindent\hypertarget{app-F7I-IX-203.1}{}\textbf{F7I-IX-203.1}\quad p. 521; printed misprint.\par
\noindent Imprimé : {\ttfamily\footnotesize schémes\allowbreak{}}\par
\noindent Corrigé : {\ttfamily\footnotesize schemes\allowbreak{}}\par
\noindent Source : 7I:790; 7I:1022. \nolinkurl{fixes/MISPRINTS_7I.json#790}; \nolinkurl{fixes/MISPRINTS_7I.json#1022}\par

\par\medskip\noindent\hypertarget{app-F7I-IX-203.2}{}\textbf{F7I-IX-203.2}\quad p. 521; misprint\_in\_print.\par
\noindent Imprimé : {\ttfamily\footnotesize notesby\allowbreak{}}\par
\noindent Corrigé : {\ttfamily\footnotesize notes\allowbreak{} by\allowbreak{}}\par
\noindent Source : FIX-7I-D-0701. \nolinkurl{fixes/FIXLOG_7I.jsonl#701}\par

\par\medskip\noindent\hypertarget{app-F7I-IX-203.3}{}\textbf{F7I-IX-203.3}\quad p. 521; printed misprint.\par
\noindent Imprimé : {\ttfamily\footnotesize Noth\allowbreak{}}\par
\noindent Corrigé : {\ttfamily\footnotesize North\allowbreak{}}\par
\noindent Source : 7I:791; 7I:1024. \nolinkurl{fixes/MISPRINTS_7I.json#791}; \nolinkurl{fixes/MISPRINTS_7I.json#1024}\par

\par\medskip\noindent\hypertarget{app-F7I-IX-204.1}{}\textbf{F7I-IX-204.1}\quad p. 521; misprint\_in\_print.\par
\noindent Imprimé : {\ttfamily\footnotesize Endlichkeitssatzfür\allowbreak{}}\par
\noindent Corrigé : {\ttfamily\footnotesize Endlichkeitssatz\allowbreak{} für\allowbreak{}}\par
\noindent Source : FIX-7I-D-0702. \nolinkurl{fixes/FIXLOG_7I.jsonl#702}\par

\par\medskip\noindent\hypertarget{app-F7I-IX-204.2}{}\textbf{F7I-IX-204.2}\quad p. 521; printed misprint.\par
\noindent Imprimé : {\ttfamily\footnotesize Abbildnngen\allowbreak{}}\par
\noindent Corrigé : {\ttfamily\footnotesize Abbildungen\allowbreak{}}\par
\noindent Source : 7I:1026. \nolinkurl{fixes/MISPRINTS_7I.json#1026}\par

\par\medskip\noindent\hypertarget{app-F7I-IX-205.1}{}\textbf{F7I-IX-205.1}\quad p. 522; printed misprint.\par
\noindent Imprimé : {\ttfamily\footnotesize  et\allowbreak{}}\par
\noindent Corrigé : \emph{(vide)}\par
\noindent Source : 7I:793; 7I:1027. \nolinkurl{fixes/MISPRINTS_7I.json#793}; \nolinkurl{fixes/MISPRINTS_7I.json#1027}\par

\par\medskip\noindent\hypertarget{app-F7I-IX-206.1}{}\textbf{F7I-IX-206.1}\quad p. 523; printed misprint.\par
\noindent Imprimé : {\ttfamily\footnotesize Herman\allowbreak{}}\par
\noindent Corrigé : {\ttfamily\footnotesize Hermann\allowbreak{}}\par
\noindent Source : 7I:794; 7I:1028. \nolinkurl{fixes/MISPRINTS_7I.json#794}; \nolinkurl{fixes/MISPRINTS_7I.json#1028}\par

\par\medskip\noindent\hypertarget{app-F7I-IX-207.1}{}\textbf{F7I-IX-207.1}\quad p. 523; misprint\_in\_print.\par
\noindent Imprimé : {\ttfamily\footnotesize míméographíé\allowbreak{}}\par
\noindent Corrigé : {\ttfamily\footnotesize miméographié\allowbreak{}}\par
\noindent Source : FIX-7I-D-0703. \nolinkurl{fixes/FIXLOG_7I.jsonl#703}\par

\par\medskip\noindent\hypertarget{app-F7I-IX-208.1}{}\textbf{F7I-IX-208.1}\quad p. 525; misprint\_pass.\par
\noindent Imprimé : {\ttfamily\footnotesize Sèminaire\allowbreak{}}\par
\noindent Corrigé : {\ttfamily\footnotesize Séminaire\allowbreak{}}\par
\noindent Source : 7I:796. \nolinkurl{fixes/FIXLOG_7I_CURRENT.jsonl#1001}\par

\par\medskip\noindent\hypertarget{app-F7I-IX-209.1}{}\textbf{F7I-IX-209.1}\quad p. 525; misprint\_pass.\par
\noindent Imprimé : {\ttfamily\footnotesize Séminare\allowbreak{}}\par
\noindent Corrigé : {\ttfamily\footnotesize Séminaire\allowbreak{}}\par
\noindent Source : 7I:797. \nolinkurl{fixes/FIXLOG_7I_CURRENT.jsonl#1002}\par

\par\medskip\noindent\hypertarget{app-F7I-IX-210.1}{}\textbf{F7I-IX-210.1}\quad p. 525; misprint\_pass.\par
\noindent Imprimé : {\ttfamily\footnotesize Strasbuorg\allowbreak{}}\par
\noindent Corrigé : {\ttfamily\footnotesize Strasbourg\allowbreak{}}\par
\noindent Source : 7I:798. \nolinkurl{fixes/FIXLOG_7I_CURRENT.jsonl#1003}\par

\par\medskip\noindent\hypertarget{app-F7I-IX-211.1}{}\textbf{F7I-IX-211.1}\quad p. 525; misprint\_pass.\par
\noindent Imprimé : {\ttfamily\footnotesize Constributions\allowbreak{}}\par
\noindent Corrigé : {\ttfamily\footnotesize Contributions\allowbreak{}}\par
\noindent Source : 7I:799. \nolinkurl{fixes/FIXLOG_7I_CURRENT.jsonl#1004}\par

\par\medskip\noindent\hypertarget{app-F7I-IX-212.1}{}\textbf{F7I-IX-212.1}\quad p. 525; misprint\_pass.\par
\noindent Imprimé : {\ttfamily\footnotesize Reguliers\allowbreak{}}\par
\noindent Corrigé : {\ttfamily\footnotesize Réguliers\allowbreak{}}\par
\noindent Source : 7I:800. \nolinkurl{fixes/FIXLOG_7I_CURRENT.jsonl#1005}\par

\par\medskip\noindent\hypertarget{app-F7I-IX-212.2}{}\textbf{F7I-IX-212.2}\quad p. 525; misprint\_pass.\par
\noindent Imprimé : {\ttfamily\footnotesize Seminaire\allowbreak{}}\par
\noindent Corrigé : {\ttfamily\footnotesize Séminaire\allowbreak{}}\par
\noindent Source : 7I:801. \nolinkurl{fixes/FIXLOG_7I_CURRENT.jsonl#1006}\par

\par\medskip\noindent\hypertarget{app-F7I-IX-212.3}{}\textbf{F7I-IX-212.3}\quad p. 525; misprint\_pass.\par
\noindent Imprimé : {\ttfamily\footnotesize Algebres\allowbreak{}}\par
\noindent Corrigé : {\ttfamily\footnotesize Algèbres\allowbreak{}}\par
\noindent Source : 7I:801. \nolinkurl{fixes/FIXLOG_7I_CURRENT.jsonl#1006}\par

\par\medskip\noindent\hypertarget{app-F7I-IX-213.1}{}\textbf{F7I-IX-213.1}\quad p. 526; misprint\_pass.\par
\noindent Imprimé : {\ttfamily\footnotesize Seminaire\allowbreak{}}\par
\noindent Corrigé : {\ttfamily\footnotesize Séminaire\allowbreak{}}\par
\noindent Source : 7I:1040. \nolinkurl{fixes/FIXLOG_7I_CURRENT.jsonl#1013}\par

\par\medskip\noindent\hypertarget{app-F7I-IX-214.1}{}\textbf{F7I-IX-214.1}\quad p. 526; misprint\_pass.\par
\noindent Imprimé : {\ttfamily\footnotesize Revètements\allowbreak{}}\par
\noindent Corrigé : {\ttfamily\footnotesize Revêtements\allowbreak{}}\par
\noindent Source : 7I:1044. \nolinkurl{fixes/FIXLOG_7I_CURRENT.jsonl#1014}\par

\par\medskip\noindent\hypertarget{app-F7I-IX-215.1}{}\textbf{F7I-IX-215.1}\quad p. 527; misprint\_pass.\par
\noindent Imprimé : {\ttfamily\footnotesize .\allowbreak{}}\par
\noindent Corrigé : \emph{(vide)}\par
\noindent Source : 7I:1045. \nolinkurl{fixes/FIXLOG_7I_CURRENT.jsonl#1015}\par

\normalsize\fi
\part{SGA 7 II}
\clearpage\SGAtoc{Exposé X — INTERSECTIONS SUR LES SURFACES REGULIERES}
\selectlanguage{french}

\SGApage{II}{1}


\underline{\underline{INTERSECTIONS SUR LES SURFACES REGULIERES}}


\underline{par P. Deligne}

Dans la première partie de cet exposé, après quelques préliminaires de théorie des intersections (contenus dans $[5] \cup [7]$), on donne la démonstration d'après Artin d'un cas particulier d'un théorème de Raynaud [9] (1.13), et une démonstration a priori d'un théorème que Néron [8] vérifiait cas par cas (1.15). De (1.13) résultent assez aisément les théorèmes de [9] utilisés dans [2] et par là dans IX.

Dans la seconde partie, je poursuis sur ma lancée et raconte comment les ``italiens'' [1], [3] calculaient les nombres d'intersections, sur les surfaces. Les résultats principaux 2.9 et 2.11 se trouvent par ailleurs dans Northcott [12] et [13]. Pour terminer, je donne en 2.27, 2.30, 2.34 une relation conjecturale entre la dimension du schéma formel versel des modules d'une singularité d'une courbe réduite, l'invariant $\delta$ de la singularité et un invariant mesurant sa symétrie.

Cet exposé ne sera pas utilisé dans la suite du séminaire.


\SGAsectiontoc{1. Nombres d'intersection sur les surfaces arithmétiques}{1}
1. \underline{Nombres d'intersection sur les surfaces arithmétiques}

1.0. On suppose donné dans ce \S{} un schéma régulier $X$ de dimension 2, un diviseur réduit $Y$ de $X$, et un sous-corps $k$ de $H^0(Y,\underline{O}_Y)$. On suppose que $Y$ est une courbe (EGA II 7.4.1) propre sur $k$. Voici trois exemples de cette situation

\underline{Exemple} 1.1. Soit $(S,\eta,s)$ un trait et $X$ un $S$-schéma $f : X \longrightarrow S$, régulier, propre et plat sur $S$, à fibres purement de dimension un. On prend $Y = (f^{-1}(s))_{red}$ et $k = k(s)$. On se restreindra à ce cas à partir de 1.10.


\SGApage{II}{2}

\SGAvar{\underline{Exemple} 1.2. Soient $X_o$ le spectre d'un anneau local normal de dimension 2, $s$ le point fermé de $X_o$ et $f : X \longrightarrow X_o$ un $X_o$-schéma régulier et connexe, propre sur $X_o$, distinct de $X_o$ mais coîncidant avec $X_o$ au-dessus de $X_o \setminus \{s\}$. On prend $Y = (f^{-1}(s))_{red}$ et $k = k(s)$.

}{\underline{Exemple} 1.2. Soient $X_o$ le spectre d'un anneau local normal de dimension 2, $s$ le point fermé de $X_o$ et $f : X \longrightarrow X_o$ un $X_o$-schéma régulier et connexe, propre sur $X_o$, distinct de $X_o$ mais coïncidant avec $X_o$ au-dessus de $X_o \setminus \{s\}$. On prend $Y = (f^{-1}(s))_{red}$ et $k = k(s)$.

}{F7II-X-001}\underline{Exemple} 1.3. On prend pour $X$ une surface régulière complète sur un corps $k$ et pour $Y$ une courbe réduite tracée sur $X$.

Rappelons le lemme suivant (SGA 6 IV ex. 1.9).

\SGAvar{\underline{Lemme} 1.4. \underline{Soit} $Z_o$ \underline{un sous-schéma fermé d'un schéma noethérien} $Z$. \underline{Alors, l'inclusion de} $Z_o$ \underline{dans} $Z$ \underline{induit un isomorphisme du groupe de Grothendieck de la catégorie des faisceauxcchérents de} $O_{\mathbb{Z}_o}$ \underline{-modules dans le groupe de Grothendieck de la catégorie des faisceaux cohérents sur} $Z$ \underline{à support dans} $Z_o$ .

}{\underline{Lemme} 1.4. \underline{Soit} $Z_o$ \underline{un sous-schéma fermé d'un schéma noethérien} $Z$. \underline{Alors, l'inclusion de} $Z_o$ \underline{dans} $Z$ \underline{induit un isomorphisme du groupe de Grothendieck de la catégorie des faisceaux cohérents de} $\underline{O}_{Z_o}$ \underline{-modules dans le groupe de Grothendieck de la catégorie des faisceaux cohérents sur} $Z$ \underline{à support dans} $Z_o$ .

}{F7II-X-002}Ce lemme permet de définir la caractéristique d'Euler-Poincaré $\chi(\mathcal{F})$ d'un faisceau cohérent $\mathcal{F}$ sur $X$ à support dans $Y$, de façon à ce que,

(i) Si $\mathcal{F}$ est un faisceau cohérent de $\underline{O}_Y$-modules, alors

\begin{equation}\tag{1.4.1}
\chi(\mathcal{F}) \;=\; \dim_k H^0(\mathcal{F}) - \dim_k H^1(\mathcal{F})
\end{equation}
(ii) $\chi$ est additif relativement aux suites exactes courtes.

Si les faisceaux de cohomologie d'un complexe $K$ de $\underline{O}_X$-modules sont des faisceaux cohérents à support dans $Y$, on pose

\begin{equation}\tag{1.4.2}
\chi(K) \;=\; \Sigma \, (-1)^i \, \chi(\underline{H}^i(K)) \qquad .
\end{equation}
Si $\mathcal{F}$ est localisé en un point $y$ de $Y$, on a

\begin{equation}\tag{1.4.3}
\chi(\mathcal{F}) \;=\; \mathrm{lg}_{\underline{O}_y}(\mathcal{F}_y) \, . \, [k(y) : k] \qquad .
\end{equation}

\SGApage{II}{3}

Sous les hypothèses (1.1), (1.2) ou (1.3), on a

\begin{equation}\tag{1.4.4}
\chi(\mathcal{F}) \;=\; \mathrm{lg}\; f_{*}\mathcal{F} \;-\; \mathrm{lg}\; R^{1}f_{*}\mathcal{F} \quad .
\end{equation}
Soient $D$ un diviseur $\geq 0$ concentré sur $Y$, et $\mathcal{L}$ un faisceau inversible sur $D$. Il y a alors une et une seule façon de définir le degré $\deg_D(\mathcal{L})$ de sorte que

\begin{equation}\tag{1.4.5}
\left\{
\begin{array}{l}
\deg_D(\mathcal{L}' \otimes \mathcal{L}'') \;=\; \deg_D(\mathcal{L}') + \deg_D(\mathcal{L}'') \\[3ex]
\text{pour } \mathcal{L} \subset \underline{O} \text{ , on a } \deg_D(\mathcal{L}) \;=\; - \chi(\underline{O}/\mathcal{L}) \quad .
\end{array}
\right.
\end{equation}
On a alors la formule de Riemann-Roch sur $D$

\begin{equation}\tag{1.4.6}
\chi(\mathcal{L}) \;=\; \chi(\underline{O}_D) \;+\; \deg_D(\mathcal{L}) \quad .
\end{equation}
Soit $D = \sum n_i \, D_i$ la décomposition de $D$ en diviseurs irréductibles réduits ; on a

\begin{equation}\tag{1.4.7}
\deg_D(\mathcal{L}) \;=\; \sum n_i \, \deg_{D_i}(\mathcal{L}) \quad .
\end{equation}
\underline{Définition} 1.5. \underline{Soient} $D$ \underline{et} $E$ \underline{deux diviseurs} $\geq 0$ \underline{sur} $X$, \underline{et supposons} $D$ \underline{concentré sur} $Y$. \underline{On définit le nombre d'intersection de} $D$ \underline{et} $E$ \underline{(relativement au corps} $k$ (1.0)) \underline{par} :

\begin{equation*}
(D,E) \;=\; \chi(\underline{O}_D \overset{\mathbb{L}}{\otimes} \underline{O}_E) \;=\; \chi(\underline{O}_D \otimes \underline{O}_E) \;-\; \chi(\mathrm{Tor}^{1}(\underline{O}_D,\underline{O}_E)) \quad .
\end{equation*}
Sous les hypothèses faites, on a en effet $\mathrm{Tor}^{2}(\underline{O}_D,\underline{O}_E) = 0$ .

\underline{Proposition} 1.6. \underline{Sous les hypothèses} 1.5, \underline{on a} :

(i) \quad \underline{Le nombre d'intersection} $(D,E)$ \underline{est additif en les diviseurs} $D$ \underline{et} $E$ ; \underline{ceci permet}, \underline{par linéarité}, \underline{de définir le nombre d'intersection de deux diviseurs non nécessairement positifs.}


\SGApage{II}{4}

(ii) \quad $(D,E) \;=\; \deg_D(\underline{O}(E))$.

(iii) \quad \underline{Si} $D$ \underline{et} $E$ \underline{ne se coupent qu'en un nombre fini de points}, \underline{alors} $(D,E) = \chi(\underline{O}_D \otimes \underline{O}_E) \geq 0$, \underline{l'égalité ne pouvant avoir lieu que si} $D$ \underline{et} $E$ \underline{sont disjoints.}

(iv) \quad \underline{Si} $D$ \underline{et} $E$ \underline{sont concentrés sur} $Y$, \underline{alors} $(D,E) = (E,D)$.

(v) \quad \underline{Si} $N$ \underline{est le faisceau conormal à} $D$, \underline{alors}

\begin{equation*}
(D)^{2} \;=\; (D,D) \;=\; - \deg_D(N) \quad .
\end{equation*}
(vi) \quad \underline{Si} $E$, \underline{non nécessairement positif}, \underline{est le diviseur d'une fonction méromorphe sur} $X$, \underline{alors} $(D,E) = 0$.

\underline{Preuve}. Soit la suite exacte

\begin{equation*}
0 \longrightarrow \underline{O}(-E) \longrightarrow \underline{O}_X \longrightarrow \underline{O}_E \longrightarrow 0 \quad ,
\end{equation*}
elle fournit, par (1.4.6),

\begin{equation*}
\begin{array}{rl}
(D,E) \;\overset{\mathrm{dfn}}{=\!=}\; & \chi(\underline{O}_D \overset{\mathbb{L}}{\otimes} \underline{O}_E) \;=\; \chi(\underline{O}_D \overset{\mathbb{L}}{\otimes} \underline{O}_X) \;-\; \chi(\underline{O}_D \overset{\mathbb{L}}{\otimes} \underline{O}(-E)) \\[3ex]
= & \chi(\underline{O}_D) \;-\; \chi(\underline{O}(-E)\vert D) \;=\; - \deg_D (\underline{O}(-E)) \;=\; \deg_D (\underline{O}(E)) \quad .
\end{array}
\end{equation*}
Ceci prouve (ii), (vi) et la biadditivité en vertu de (1.4.5) et (1.4.7). Sous les hypothèses (iii), on a $\mathrm{Tor}^{1}(\underline{O}_D,\underline{O}_E) = 0$, d'où la conclusion. Sous les hypothèses (v), on a

\begin{equation*}
\underline{O}_D \otimes \underline{O}_D \;=\; \underline{O}_D \quad \text{et} \quad \mathrm{Tor}^{1}(\underline{O}_D,\underline{O}_D) \;\simeq\; N \quad (\text{SGA 6 VII 2.5}) \quad ,
\end{equation*}
d'où

\begin{equation*}
(D,D) \;=\; \chi(\underline{O}_D) \;-\; \chi(N) \;=\; - \deg_D(N) \quad .
\end{equation*}
Enfin, (iv) résulte de la symétrie du produit tensoriel.

On notera que (ii), (iv) et (vi) restent vrais sans supposer les diviseurs positifs.


\SGApage{II}{5}

\underline{Proposition} 1.7. \underline{Soient} $(Y_i)_{1 \leq i \leq k}$ \underline{les composantes irréductibles de} $Y$, $f$ \underline{une fonction méromorphe sur} $X$, \underline{et}

\begin{equation*}
\mathrm{div}(f) \;=\; \sum n_i\, Y_i \;+\; Z
\end{equation*}
\underline{son diviseur}, \underline{avec} $|Z| \cap Y$ \underline{fini}. \underline{Alors}, \underline{posant} $D_i = n_i\, Y_i$, \underline{on a}, \underline{pour} $(a_i) \in \mathbb{Q}^k$ :

\begin{equation*}
(\sum a_i\, D_i)^2 \;=\; -\sum_{i<j} (a_i - a_j)^2\, (D_i, D_j) \;-\; \sum a_i^2\, (D_i, Z) \qquad .
\end{equation*}
\underline{Preuve} (d'après [7]). En vertu de 1.6 (v) (vi), on a

\SGAvar{\begin{equation*}
\begin{aligned}
(\sum_i a_i\, D_i)^2 \;&=\; \sum_i a_i\, (D_i,\; \sum_j a_j\, D_j) \;=\; \sum_i a_i\, (D_i,\; \sum_j a_j\, D_j - a_i\ \mathrm{div}\ f) \\[2pt]
&=\; \sum a_i\, (D_i,\; \sum_{j \neq i} (a_j - a_i) D_j \;-\; a_i Z) \\[2pt]
&=\; \sum_{i \neq j} a_i\, (a_j - a_i)(D_i, D_j) \;-\; \sum_i a_i^2\, (D_i, Z) \\[2pt]
&=\; -\sum_{i<j} (a_i - a_j)^2\, (D_i\ D_j) \;-\; \sum_i a_i^2\, (D_i, Z)
\end{aligned}
\end{equation*}
\underline{Corollaire} 1.8. \underline{Sous les hypothèses de l'exemple} 1.1, \underline{prenons pour} $f$ \underline{une uniformisante} $\pi$ \underline{de} $S$. \underline{Ici}, $Z = 0$ \underline{et}

}{\begin{equation*}
\begin{aligned}
(\sum_i a_i\, D_i)^2 \;&=\; \sum_i a_i\, (D_i,\; \sum_j a_j\, D_j) \;=\; \sum_i a_i\, (D_i,\; \sum_j a_j\, D_j - a_i\ \mathrm{div}\ f) \\[2pt]
&=\; \sum a_i\, (D_i,\; \sum_{j \neq i} (a_j - a_i) D_j \;-\; a_i Z) \\[2pt]
&=\; \sum_{i \neq j} a_i\, (a_j - a_i)(D_i, D_j) \;-\; \sum_i a_i^2\, (D_i, Z) \\[2pt]
&=\; -\sum_{i<j} (a_i - a_j)^2\, (D_i, D_j) \;-\; \sum_i a_i^2\, (D_i, Z)
\end{aligned}
\end{equation*}
\underline{Corollaire} 1.8. \underline{Sous les hypothèses de l'exemple} 1.1, \underline{prenons pour} $f$ \underline{une uniformisante} $\pi$ \underline{de} $S$. \underline{Ici}, $Z = 0$ \underline{et}

}{F7II-X-003}\begin{equation*}
(\sum k_i\, Y_i)^2 \;=\; -\sum_{i<j} \frac{1}{n_i n_j}\; (k_i\, n_j - k_j\, n_i)^2\, (Y_i,\, Y_j) \qquad ;
\end{equation*}
\underline{donc la forme d'intersection est négative}, \underline{et si} $Y$ \underline{est connexe}, \underline{son noyau est réduit aux multiples de} $\mathrm{div}(\pi)$.

\underline{Corollaire} 1.9 (Mumford [7]). \underline{Sous les hypothèses de l'exemple} 1.2, \underline{prenons pour} $f$ \underline{un élément non nul de l'idéal maximal de l'anneau local} $\underline{O}_s$ . \underline{Ici}, $Z > 0$, \underline{de sorte que la forme d'intersection est définie négative}.


\SGApage{II}{6}

1.10. \underline{On se restreint}, \underline{pour la suite}, \underline{au cas considéré dans l'exemple} (1.1).

\SGAvar{Lichtenbaum [5] a prouvé que, dans ce cas, $X$ est automatiquement projectif sur $S$, de sorte qu'il n'est pas fatiguant de définir le complexe $Rf^{!}\underline{O}_S$ ([4] III 8.7 p. 190). Puisque $f$ est de dimension relative 1 et que $X$ et $S$ sont réguliers(Gorenstein suffirait) il existe ([4] V 8.3 et V 9.1) un faisceau inversible $\omega_{X/S}$ sur $S$, unique à isomorphisme unique près, tel que

}{Lichtenbaum [5] a prouvé que, dans ce cas, $X$ est automatiquement projectif sur $S$, de sorte qu'il n'est pas fatigant de définir le complexe $Rf^{!}\underline{O}_S$ ([4] III 8.7 p. 190). Puisque $f$ est de dimension relative 1 et que $X$ et $S$ sont réguliers (Gorenstein suffirait) il existe ([4] V 8.3 et V 9.1) un faisceau inversible $\omega_{X/S}$ sur $X$, unique à isomorphisme unique près, tel que

}{F7II-X-004}\begin{equation}\tag{1.10.1}
Rf^{!}\underline{O}_S \;\simeq\; \omega_{X/S}[1] \quad .
\end{equation}
On désigne par $K$ un quelconque diviseur (effectif ou non) tel que $\omega \simeq \underline{O}(K)$.

Soient $D$ un diviseur $\geq 0$ concentré sur $Y$ et $N$ son faisceau conormal :

\begin{center}
\begin{tikzcd}[column sep=large, row sep=large]
D \arrow[rr, "i"] \arrow[dr, swap, "g"] \& {} \& X \arrow[dl, "f"] \\
 {} \& S \&
{} \end{tikzcd}\quad .
\end{center}










On a

\begin{equation}\tag{1.10.2}
Rg^{!}\underline{O}_S = Ri^{!}Rf^{!}\underline{O}_S = Ri^{!}\omega_{X/S}[1] = \omega_{X/S}\big|\ D \otimes N^{\vee}[0]
\end{equation}
En vertu du théorème de dualité ([4] V 11.1 p.210),

\begin{equation}\tag{1.10.3}
Rg_{*}\, Rg^{!}\underline{O}_S \;\simeq\; \mathrm{RHom}(Rg_{*}\, \underline{O}_D,\underline{O}_S) \quad .
\end{equation}
De plus, pour tout $\underline{O}_S$-module $M$ concentré au point fermé, on a

\begin{equation}\tag{1.10.4}
\chi(\mathrm{RHom}(M,\underline{O}_S)) \;=\; -\,\chi(M)
\end{equation}
de sorte qu'ici, utilisant (1.4.4) et (1.10.3), on obtient

\begin{equation*}
\chi(Rg^{!}\underline{O}_S) \;=\; -\,\chi(\underline{O}_D) \quad .
\end{equation*}
Puisque (1.4.6)

\begin{equation*}
\chi(Rg^{!}\underline{O}_S) = \chi(\underline{O}_D) + \deg_D(Rg^{!}\underline{O}_S) \qquad , 
\end{equation*}

\SGApage{II}{7}

 on en conclut par $(1.10.2)$ que

\SGAvar{\begin{equation*}
- 2\,\chi(\underline{O}_D) \;=\; \deg_D\,(Rg^{!}\underline{O}_D) \;=\; \deg_D\,(\omega|D \otimes N^{\vee}) \;=\; \deg_D\,(\omega) - \deg_D(N) \qquad ,
\end{equation*}
ce qui, par 1.6 (ii) et(v) se réécrit

\underline{Proposition} 1.11. \underline{Sous les hypothèses} (1.1), \underline{quelque soit le diviseur} $D \geq 0$ \underline{concentré sur} $Y$, \underline{on a}

}{\begin{equation*}
- 2\,\chi(\underline{O}_D) \;=\; \deg_D\,(Rg^{!}\underline{O}_S) \;=\; \deg_D\,(\omega|D \otimes N^{\vee}) \;=\; \deg_D\,(\omega) - \deg_D(N) \qquad ,
\end{equation*}
ce qui, par 1.6 (ii) et (v) se réécrit

\underline{Proposition} 1.11. \underline{Sous les hypothèses} (1.1), \underline{quel que soit le diviseur} $D \geq 0$ \underline{concentré sur} $Y$, \underline{on a}

}{F7II-X-005}\begin{equation*}
\chi(\underline{O}_D) \;=\; - \frac{1}{2}\,(D,D+K) \qquad .
\end{equation*}
1.12. On désignera par $d$ le pgcd des multiplicités géométriques des composantes irréductibles de $X_s = f^{-1}(s)$, i.e. le pgcd des longueurs des anneaux locaux des points maximaux de la fibre spéciale géométrique de $f$. Si $X_{\eta}$ a un point rationnel, i.e. si $f : X \longrightarrow S$ a une section $x$, alors $d = 1$ ; plus précisément $f$ est lisse en $x(s)$. En effet, $x$ est une immersion fermée de schémas réguliers, donc est une immersion régulière et on applique EGA IV 17.12.1.

Voici la démonstration, d'après M. Artin, d'un cas particulier d'un théorème de M. Raynaud [9], qui suppose seulement $d$ premier à la caractéristique $p$ de $k$.

\underline{Théorème} 1.13. \underline{Avec les notations précédentes}, \underline{si} $Y$ \underline{est géométriquement connexe et si} $d = 1$, \underline{alors} $X$ \underline{est cohomologiquement plat sur} $S$ \underline{en dimension} 0 (EGA III 7.8.1), \underline{plus précisément} (EGA III 7.6.9 (ii)) \underline{on a}

\begin{equation*}
k \xrightarrow{\;\sim\;} H^0(X_s,\, \underline{O}_{X_s}).
\end{equation*}
On se ramène aisément au cas où $S$ est strictement local. Comme $Y$ est géométriquement connexe, $H^0(Y,\underline{O}_Y)$ est une extension radicielle de $k$, et puisque $(d,p)=1$, une au moins des composantes irréductibles de $Y$ est
\SGApage{II}{8}
géométriquement réduite, et on en déduit facilement que $H^0(Y,\underline{O}_Y) = k$.

Soit $D$ un diviseur sur $X$ tel que

\begin{equation*}
Y \leq D < X_s \qquad .
\end{equation*}
Puisque $d = 1$, $D$ n'est pas multiple rationnel de $X_s$, donc par 1.8

\begin{equation*}
(D,X_s-D) \;=\; - \,(D,D) \;>\; 0 \qquad .
\end{equation*}
Le diviseur $Y$ a donc une composante irréductible $E$ telle que

\begin{equation*}
0 < E \leq X_s-D \quad \text{et} \quad (D,E) > 0 \qquad .
\end{equation*}
Considérons le diagramme

\begin{center}
\begin{tikzcd}[column sep=large, row sep=large]
0 \arrow[r] \& \underline{O}(-D-E) \arrow[r] \arrow[d, "\varphi"] \& \underline{O} \arrow[r] \arrow[d, equal] \& \underline{O}_{D+E} \arrow[r] \arrow[d, "\psi"] \& 0 \\
0 \arrow[r] \& \underline{O}(-D) \arrow[r] \& \underline{O} \arrow[r] \& \underline{O}_D \arrow[r] \& 0
\end{tikzcd}\qquad ,
\end{center}
















le diagramme du serpent fournit un isomorphisme entre $\mathrm{Ker}(\psi)$ et $\mathrm{Coker}(\varphi)$. Ce conoyau n'est autre que la restriction de $\underline{O}(-D)$ à $E$, d'où une suite exacte

\SGAvar{\begin{equation*}
0 \longrightarrow \underline{O}(-D)|E \longrightarrow \underline{O}_{D+E} \longrightarrow \underline{O}_D \longrightarrow 0 \qquad .
\end{equation*}
Par construction $\deg_E(\underline{O}(-D) = - (D,E) < 0$, donc $H^0(\underline{O}(-D)|E) = 0$, et la flèche de restriction

}{\begin{equation*}
0 \longrightarrow \underline{O}(-D)|E \longrightarrow \underline{O}_{D+E} \longrightarrow \underline{O}_D \longrightarrow 0 \qquad .
\end{equation*}
Par construction $\deg_E(\underline{O}(-D)|E) = - (D,E) < 0$, donc $H^0(\underline{O}(-D)|E) = 0$, et la flèche de restriction

}{F7II-X-006}\begin{equation*}
0 \longrightarrow H^0(\underline{O}_{D+E}) \longrightarrow H^0(\underline{O}_D)
\end{equation*}
est injective. Ceci permet de construire par récurrence une suite de diviseurs

\begin{equation*}
Y = D_0 \leq D_1 \leq \ldots \leq D_n = X_s
\end{equation*}
telle que $H^0(\underline{O}_{D_i})$ s'injecte dans $H^0(\underline{O}_{D_{i-1}})$, donc $H^0(\underline{O}_{X_s}) = H^0(\underline{O}_Y) = k$, ce qui achève la démonstration de 1.13.


\SGApage{II}{9}

Le fait suivant est prouvé dans [5] ou [10].

\underline{Théorème} 1.14. (Castelnuovo, Enriques,...). \underline{Sous les hypothèses} (1.1), \underline{soit} $D$ \underline{un diviseur réduit irréductible contenu dans} $Y$. \underline{Supposons que} $H^1(D,\underline{O}_D) = 0$ \underline{et que,} \underline{posant} $k' = H^0(D,\underline{O}_D)$, \underline{on ait}

\SGAvar{\begin{equation*}
(D,D) \; = \; - [k' : k] \qquad .
\end{equation*}
\underline{Alors,} $D$ \underline{est une droite projective sur} $k'$ \underline{et il existe un} $S$-\underline{schéma régulier} $X_1$ , \underline{un point fermé} $x_1$ \underline{de} $X_1$ \underline{de corps résiduel} $k'$ \underline{et un isomorphisme entre} $X$ \underline{et le} $S$-\underline{schéma déduit de} $X_1$ \underline{par éclatement de} $x_1$, \underline{cet isomorphisme identifiant} $D$ \underline{à la fibre de} $x_1$ \underline{dans} $X_1$.

On appelle un diviseur $D$ possèdant ces propriétés une \underline{courbe exceptionnelle de première espèce}.

}{\begin{equation*}
(D,D) \; = \; - [k' : k] \qquad .
\end{equation*}
\underline{Alors,} $D$ \underline{est une droite projective sur} $k'$ \underline{et il existe un} $S$-\underline{schéma régulier} $X_1$ , \underline{un point fermé} $x_1$ \underline{de} $X_1$ \underline{de corps résiduel} $k'$ \underline{et un isomorphisme entre} $X$ \underline{et le} $S$-\underline{schéma déduit de} $X_1$ \underline{par éclatement de} $x_1$, \underline{cet isomorphisme identifiant} $D$ \underline{à la fibre au-dessus de} $x_1$ \underline{dans} $X$.

On appelle un diviseur $D$ possédant ces propriétés une \underline{courbe exceptionnelle de première espèce}.

}{F7II-X-007}Je ne résiste pas au plaisir de démontrer la proposition suivante de Néron [8] prop 1' p. 90 :

\underline{Proposition} 1.15. \underline{Supposons que} $X_\eta$ \underline{soit une courbe elliptique} \underline{(en particulier ait une section)} \underline{et que} $X$ \underline{ne contienne aucune courbe exceptionnelle de première espèce.} \underline{Alors,} \underline{avec la notation} (1.10.1), \underline{on a}

\begin{equation*}
f^*f_* \; \omega_{X/S} \qquad \stackrel{\sim}{\longrightarrow} \qquad \omega_{X/S}
\end{equation*}
\underline{universellement}.

Comme $d = 1$ (1.12), il suffit, d'après (1.13), de prouver que $\omega_{X/S}$ est isomorphe à $\underline{O}_X$ ; puisque $\omega_{X/S}$ est trivial sur la courbe elliptique $X_\eta$ , on sait a priori que

\begin{equation*}
\omega_{X/S} \qquad \simeq \qquad \underline{O}(\Sigma \, n_i \, Y_i) \qquad , 
\end{equation*}

\SGApage{II}{10}

 où les $Y_i$ sont les composantes irréductibles de $Y$, et puisque $\mathrm{div}(\pi) \sim 0$, on peut supposer les $n_i$ positifs. Si $\Sigma \, n_i \, Y_i$ est proportionnel au diviseur $\mathrm{div}(\pi)$, il est multiple entier de $\mathrm{div}(\pi)$ et on a gagné. Sinon, on a

\begin{equation*}
(\Sigma \, n_i \, Y_i)^2 < 0 \qquad ,
\end{equation*}
il existe $i$ tel que

\begin{equation*}
(Y_i, \; \Sigma \, n_j \, Y_j) < 0 \qquad ,
\end{equation*}
et on applique le lemme suivant :

\underline{Lemme} 1.16. \underline{Sous les hypothèses} 1.1, \underline{supposons que} $Y$ \underline{soit connexe et que} $Y$ \underline{soit réductible,} \underline{ou que} $\chi(\underline{O}_Y) \leq 0$. \underline{Alors,} \underline{pour qu'une composante irréductible} $Y_i$ \underline{de} $Y$ \underline{soit une courbe exceptionnelle de première espèce,} \underline{il faut et il suffit que} $(Y_i,K) < 0$ .

Soit $k' = H^0(Y_i,\underline{O}_{Y_i})$. La formule de Riemann-Roch 1.11 fournit

\begin{equation*}
\chi(\underline{O}_{Y_i}) = [k':k] \; (1-\dim_{k'} H^1(Y_i,\underline{O}_{Y_i})) = - \frac{1}{2} \; ((Y_i,Y_i) + (Y_i,K)) \qquad .
\end{equation*}
Si $Y_i$ est exceptionnelle de première espèce, on a donc $(Y_i,K) = -[k':k] < 0$ . Réciproquement, si $(Y_i,K) < 0$, alors $\chi(\underline{O}_{Y_i}) > 0$, ce qui implique que $\chi(\underline{O}_{Y_i}) = [k':k]$ et que $Y$ soit réductible, de sorte que par 1.8, $(Y_i,Y_i) < 0$. Les nombres négatifs $(Y_i,Y_i)$ et $(Y_i,K)$ sont divisibles par $[k':k]$, de sorte qu'on doit avoir

\begin{equation*}
(Y_i,Y_i) = (Y_i,K) = -[k':k] \qquad .
\end{equation*}
\underline{Exercice} (a). Sous les hypothèses (1.1), on suppose que $X_\eta$ est un espace principal homogène sous une courbe elliptique et que $X$ ne contienne aucune courbe exceptionnelle de première espèce. Montrer que sur $X$, la formule de Riemann-Roch (1.11) s'écrit

\begin{equation*}
\chi(\underline{O}_D) = - \frac{1}{2} \, (D,D)
\end{equation*}

\SGApage{II}{11}

(b) Supposons en outre le corps résiduel $k(s)$ algébriquement clos. Montrer que le diagramme des composantes irréductibles de $X_s$ est multiple de l'un des diagrammes considérés par Néron [8] p. 124-125, 4è colonne (cas a à c 8). Ce diagramme détermine les genres des composantes irréductibles réduites, leur multiplicité dans $X_s$, la nature et la disposition des singularités de $(X_s)_{\mathrm{red}}$ .

\SGAvar{On traitera d'abord le cas où $X_s$ est irréductible ; ce cas exclus, on prouvera que toute composante irréductible réduite est une droite projective de self-intersection $-$ 2. Si $X_s = \Sigma\ n_i\ D_i$ est la décomposition de $X_s$ en diviseurs irréductibles réduits, on réécrira cette condition

\begin{equation}\tag{*}
n_i \;=\; \frac{1}{2} \sum_{j \neq i} n_i \ (D_i\ D_j) \qquad\qquad .
\end{equation}
D'autre part, $X_s$ est connexe, de sorte que :

}{On traitera d'abord le cas où $X_s$ est irréductible ; ce cas exclu, on prouvera que toute composante irréductible réduite est une droite projective de self-intersection $-$ 2. Si $X_s = \Sigma\ n_i\ D_i$ est la décomposition de $X_s$ en diviseurs irréductibles réduits, on réécrira cette condition

\begin{equation}\tag{*}
n_i \;=\; \frac{1}{2} \sum_{j \neq i} n_j \ (D_i\ D_j) \qquad\qquad .
\end{equation}
D'autre part, $X_s$ est connexe, de sorte que :

}{F7II-X-008}(**) Quels que soient $i$ et $j$, il existe une chaîne $i = i_1, i_2 \ldots i_n = j$ avec $(D_{i_k}\ D_{i_{k+1}}) \neq 0$ .

\SGAvar{On cherchera alors toutes les solutions en nombre entiers de $(*) + (**)$. On montrera d'abord que, sauf dans les cas b2 et c2 de Néron, les nombres d'intersection $(D_i D_j)$ sont égaux à zéro ou un.


}{On cherchera alors toutes les solutions en nombres entiers de $(*) + (**)$. On montrera d'abord que, sauf dans les cas b2 et c2 de Néron, les nombres d'intersection $(D_i D_j)$ sont égaux à zéro ou un.


}{F7II-X-009}\SGAsectiontoc{2. Intersections italiennes}{2}
2. \underline{Intersections italiennes}

Soient $S'$ et $S$ deux surfaces projectives lisses sur un corps $k$ et $f : S' \longrightarrow S$ un morphisme birationnel. Si $D$ et $E$ sont deux diviseurs sur $S$, on a

\begin{equation*}
(D,E) = (f^* \, D, \ f^* \, E) \qquad\qquad , \qquad 
\end{equation*}

\SGApage{II}{12}

 comme on le voit, par exemple, en faisant bouger $D$ et $E$.

Il s'impose donc de donner une expression ``birationnelle'' pour le nombre d'intersection. On se propose d'expliquer comment procédaient les italiens pour ce faire. Ma principale source de renseignements est [1] où se trouve une longue bibliographie. Voir aussi Manin [6], et Northcott[12],[13].

(2.0) Soient $S$ le spectre d'un anneau local régulier de dimension 2, $s$ le point fermé de $S$ et $k$ un sous-corps de $k(s)$ dont $k(s)$ soit une extension finie. Les nombres d'intersections et les caractéristiques d'Euler-Poincaré seront calculés en terme de $k$.

Dans ce \S, un $S$-schéma $p : S' \longrightarrow S$ sera dit être un \underline{transformé birationnel de} $S$ s'il est régulier, connexe, et si $p$ est propre et induit un isomorphisme de $S' \setminus p^{-1}(s)$ sur $S \setminus \{s\}$. Tout ce qui sera dit de $S$ s'appliquera encore aux localisés de ses transformés birationnels en un point fermé.

(2.1) Soit $f : (S'',p'') \longrightarrow (S',p')$ un morphisme de $S$-schémas entre transformés birationnels de $S$. On sait alors [5] qu'il existe une et une seule factorisation de $f$ :

\begin{equation*}
f = f_n \ \ldots \ f_1 \qquad , \qquad f_i : S_{i+1} \longrightarrow S_i
\end{equation*}
telle que $S_{i+1}$ se déduise de $S_i$ par éclatement d'un ensemble fini $X_i$ de points fermés, et que $f_i(X_{i+1}) \subset X_i$ . L'ensemble $X_i$ est l'ensemble des points de $S_i$ en lesquels la projection de $S''$ sur $S_i$ n'est pas un isomorphisme.

En particulier, $f$ induit un isomorphisme au-dessus du complémentaire d'une partie finie de $S'$, qui contient les points génériques des composantes irréductibles de $p'^{-1}(s)$. Dès lors, $f^{-1}$ définit une injection
\SGApage{II}{13}
de l'ensemble $\mathrm{Irr}({p'}^{-1}(s))$ des composantes irréductibles de ${p'}^{-1}(s)$ dans $\mathrm{Irr}({p''}^{-1}(s))$.

Les transformés birationnels de $S$ n'ont pas de $S$-automorphisme . On désignera ici par $I$ l'ensemble des classes de $S$-isomorphie des transformés birationnels de $S$, ordonné par la relation de domination. L'ensemble $I$ est filtrant à droite, et de plus petit élément $S$.

\underline{Définition} 2.2. \underline{L'ensemble} $S^*$ \underline{des} points de $S$ infiniment voisins de $s$ \underline{est} \underline{l'ensemble limite inductive}

\begin{equation*}
S^* = \varinjlim_{I} \ \mathrm{Irr}\,(p^{-1}(s)) \qquad .
\end{equation*}
Si un élément $t \in S^*$ est représenté sur un transformé birationnel $S'$ de $S$ par une courbe complète réduite irréductible $T$, le corps $H^o(T,\underline{O}_T)$ ne dépend que de $t$, et non du choix de $S'$ ; on l'appelle le \underline{corps résiduel} $k(t)$ de $t$.

Soit $S_1$ l'éclaté de $S$ le long de $s$. La droite projective $s^*$ image réciproque de $\{s\}$ dans $S_1$ s'appelle la \underline{curvetta} de $s$ et définit un élément $s^*$ de $S^*$ de corps résiduel $k(s)$.

Soient $S'$ un transformé birationnel de $S$, $t$ un point fermé de $S'$ et $S'_t$ le spectre de l'anneau local de $S'$ en $t$. On définit de façon évidente une injection de ${S'}^*_t$ dans $S^*$ ; en particulier, $t$ définit un élément $t^*$ de $S^*$, de même corps résiduel. On désigne par ${S'}^*$ la réunion dans $S^*$ des ${S'}^*_t$ pour $t$ fermé dans $S'$.

La proposition suivante résulte de la description 2.1 des transformés birationnels de $S$ :


\SGApage{II}{14}

\underline{Proposition} 2.3. \underline{Quel que soit} $t \in S^*$, \underline{il existe une et une seule suite de transformés birationnels de} $S$ $(S_i)_{0 \leqslant i \leqslant n+1}$ \underline{et de points fermés} $t_i \in S_i$ $(0 \leqslant i \leqslant n)$ \underline{telle que}

(i)\quad $S_o = S$, \underline{et} $S_{i+1}$ \underline{est l'éclaté de} $S_i$ \underline{le long de} $\{t_i\}$

(ii)\quad $t_{i+1} \in t_i^*$\quad $(0 \leqslant i < n)$ \underline{et} $t = t_n^*$ .

\underline{De plus}, \underline{quel que soit le transformé birationnel} $(S',p')$ \underline{de} $S$, \underline{pour que} $t$ \underline{appartienne à l'image de} $\mathrm{Irr}({p'}^{-1}(s))$ \underline{dans} $S^*$, \underline{il faut et il suffit que} $S'$ \underline{domine} $S_{n+1}$.

On appelle $n$ l'ordre de $t$. On désignera ici par $S(t)$ (resp. par $S_+(t)$) le transformé birationnel $S_n$ (resp. $S_{n+1}$) de $S$. Par abus de notation, on désignera encore par $t$ le point fermé $t_n$ de $S_n$ et par $t^*$ son image réciproque $t_n^*$ dans $S_{n+1}$.

On munit $S^*$ de l'ordre suivant :

\SGAvar{\begin{equation*}
t_1 \prec t_2 \quad \Longleftrightarrow \quad S_+(t_2) \ \mbox{domine} \ S_+(t_1) \qquad .
\end{equation*}
Si $t_1 \prec t_2$, alors $k(t_2)$ est une extension de $k(t_1)$. L'ensemble ordonné $S^*$ a pour plus petit élément $s^*$, et toute partie majorée de $S^*$ est finie et totalement ordonnée ($S^*$ est un ``arbre de souche $s^*$''). De plus, l'ordre $n$ de $t \in S^*$ est la distance, dans l'arbre $S^*$, de $s^*$ à $t$ ($n = \#[s^*,t]$).

}{\begin{equation*}
t_1 \prec t_2 \quad \Longleftrightarrow \quad S_+(t_2) \ \mbox{domine} \ S_+(t_1) \qquad .
\end{equation*}
Si $t_1 \prec t_2$, alors $k(t_2)$ est une extension de $k(t_1)$. L'ensemble ordonné $S^*$ a pour plus petit élément $s^*$, et toute partie majorée de $S^*$ est finie et totalement ordonnée ($S^*$ est un ``arbre de souche $s^*$''). De plus, l'ordre $n$ de $t \in S^*$ est la distance, dans l'arbre $S^*$, de $s^*$ à $t$ ($n = \#]s^*,t]$).

}{F7II-X-010}La proposition suivante explicite la description 2.1 des transformés birationnels de $S$.

\underline{Proposition} 2.4. \underline{La fonction qui à un transformé birationnel} $p : S' \to S$ \underline{de} $S$ \underline{associe la partie} $\mathrm{Irr}(p^{-1}(s))$ \underline{de} $S^*$ \underline{induit une bijection entre l'ensemble} $I$ (2.1) \underline{et l'ensemble des parties finies de} $S^*$ \underline{qui}, \underline{avec tout élément},
\SGApage{II}{15}
\underline{contiennent ses prédécesseurs}.

\SGAvar{Soient $D$ est un diviseur $\geq 0$ sur $S$ et $t \in S^*$. On dit que $t$ \underline{appartient} à $D$ si le point fermé $t$ de $S(t)$ (2.3) appartient au transformé pur (= adhérence schématique de $D - \{s\}$) de $D$ sur $S(t)$.

}{Soit $D$ un diviseur $\geq 0$ sur $S$ et $t \in S^*$. On dit que $t$ \underline{appartient} à $D$ si le point fermé $t$ de $S(t)$ (2.3) appartient au transformé pur (= adhérence schématique de $D - \{s\}$) de $D$ sur $S(t)$.

}{F7II-X-011}\underline{Définition} 2.5. \underline{Soient} $t$, $t' \in S^*$. \underline{On dit que} $t'$ \underline{est} proche \underline{de} $t$ \underline{si} $t' \in S_+(t)^*$ (2.2) \underline{et appartient à la curvetta} $t^*$ \underline{de} $t$ (2.3).

\SGAvar{Si $t'$ est un successeur immédiat de $t$, alors $t'$ est proche de $t$ ; si dans une suite $(t_i)_{0 \leq i \leq n}$ chaque $t_{i+1}$ est successeur immédiat de $t_i$, et si $t_n$ est proche de $t_o$, alors $t_i$ est proche de $t_o$ pour $0 < i \leq n$ et $t_n$ n'est pas proche de $t_i$ pour $0 < i \leq n-1$.

}{Si $t'$ est un successeur immédiat de $t$, alors $t'$ est proche de $t$ ; si dans une suite $(t_i)_{0 \leq i \leq n}$ chaque $t_{i+1}$ est successeur immédiat de $t_i$, et si $t_n$ est proche de $t_0$, alors $t_i$ est proche de $t_0$ pour $0 < i \leq n$ et $t_n$ n'est pas proche de $t_i$ pour $0 < i < n-1$.

}{F7II-X-012}(2.6)  Soient $A$ un anneau local noethérien de dimension $n$, $T$ son spectre, $m$ son idéal maximal, $k = A/m$ son corps résiduel et $P = \sum_{0 \leq i < n} a_i X^i$ son polynôme de Hilbert-Samuel :

\SGAvar{\begin{equation*}
P(n) \;=\; \dim_k(m^n/m^{n+1}) \qquad \text{pour } n \text{ grand.}
\end{equation*}
Rappelons que la multiplicité $\mu_s$ dans $T$ du point fermé $s$ vérifie par définition

}{\begin{equation*}
P(r) \;=\; \dim_k(m^r/m^{r+1}) \qquad \text{pour } r \text{ grand.}
\end{equation*}
Rappelons que la multiplicité $\mu_s$ dans $T$ du point fermé $s$ vérifie par définition

}{F7II-X-013}\begin{equation*}
a_{n-1} \;=\; \mu_s/(n-1)\, ! \qquad .
\end{equation*}
Lorsque $n = 1$, on a donc simplement

\SGAvar{\begin{equation*}
\mu_s \;=\; \dim_k m^n/m^{n+1} \qquad \text{pour } n \text{ grand.}
\end{equation*}
En termes géométriques, $\mu_s$ peut se définir comme le degré, dans l'espace projectif sur $k$ $P_k(m/m^2)$, de la fibre spéciale de l'éclaté de $T$ le long du sous-schéma $\{s\}$.


}{\begin{equation*}
\mu_s \;=\; \dim_k m^r/m^{r+1} \qquad \text{pour } r \text{ grand.}
\end{equation*}
En termes géométriques, $\mu_s$ peut se définir comme le degré, dans l'espace projectif sur $k$ $P_k(m/m^2)$, de la fibre spéciale de l'éclaté de $T$ le long du sous-schéma $\{s\}$.


}{F7II-X-014}\SGApage{II}{16}

Rappelons que si $Y$ et $Z$ sont deux sous-schémas fermés d'un schéma $X$, alors, le transformé pur de $Y$ dans l'éclaté $\tilde{X}$ de $X$ le long de $Z$, adhérence schématique de $Y \setminus Z$ dans $\tilde{X}$, n'est autre que l'éclaté de $Y$ le long de $Y \cap Z$.

Supposons que $T$ soit un sous-schéma fermé d'un schéma régulier $U$ . Soient $p : U_1 \longrightarrow U$ l'éclaté de $U$ le long de $\{s\}$, $s^*$ l'image réciproque de $s$ dans $U_1$ et $T'$ le transformé pur de $T$ dans $U_1$. Ce qui précède montre que la multiplicité de $T$ en $s$ est le degré dans l'espace projectif $s^*$ du sous-schéma $s^* \cap T'$ .

Supposons en outre que $T$ soit un diviseur de $U$, soit $T^*$ son transformé total dans $U_1$ (le diviseur image réciproque) et définissons $\nu$ par la formule $T^* = \nu\, s^* + T'$. Si $X$ est une droite de $s^*$, et si on calcule les caractéristiques d'Euler-Poincaré comme sommes alternées de dimension sur $k$, on sait que

\begin{equation}\tag{2.6.1}
(X,s^*) = \chi(\underline{O}_X \overset{\mathbb{L}}{\otimes} \underline{O}_{s^*}) \;=\; \deg_X \underline{O}(s^*) = -1
\end{equation}
et on a vérifié que

\begin{equation}\tag{2.6.2}
(X,T') = \chi(\underline{O}_X \overset{\mathbb{L}}{\otimes} \underline{O}_{T'}) \;=\; \deg_{s^*} (s^* \cap T') = \mu_s \qquad .
\end{equation}
La formule de projection

\begin{equation*}
Rp_*(\underline{O}_X \overset{\mathbb{L}}{\otimes} \underline{O}_{T^*}) \;=\; Rp_*(\underline{O}_X \overset{\mathbb{L}}{\otimes} Lp^* \underline{O}_T) = Rp_* \underline{O}_X \overset{\mathbb{L}}{\otimes} \underline{O}_T
\end{equation*}
et son corollaire

\begin{equation*}
(X,T^*) \;=\; \nu(X,s^*) + (X,T') = 0
\end{equation*}
fournit, par (2.6.1), (2.6.2), l'identité

\begin{equation}\tag{2.6.3}
\mu_s \;=\; \nu \qquad .
\end{equation}

\SGApage{II}{17}

\underline{Proposition} 2.7. \underline{Soient} $T$ \underline{un diviseur} $\geq 0$ \underline{dans un schéma régulier} $U$, $s$ \underline{un point fermé de} $U$, $U_1$ \underline{l'éclaté de} $U$ \underline{le long de} $\{s\}$, $s^*$ \underline{l'image réciproque de} $s$ \underline{dans} $U_1$, $T'$ \underline{le transformé pur de} $T$ \underline{sur} $U_1$ \underline{et} $T^*$ \underline{son transformé total}.

\underline{Il y a identité entre}

(i) \underline{La multiplicité}, \underline{au sens de Samuel}, \underline{du point} $s$ \underline{de} $T$.

(ii) \underline{Le degré}, \underline{dans l'espace projectif} $s^*$, \underline{de} $s^* \cap T'$.

(iii) \underline{La multiplicité de} $s^*$ \underline{dans le diviseur} $T^*$.

\SGAvar{2.8. Reprenons les hypothèses de 2.0. Soit $D$ un diviseur $\geq 0$ sur $S$ et $t \in S^*$. La \underline{multiplicité} $\mu_t(D)$ \underline{de} $D$ \underline{en} $t$ est, par définition, la multiplicité en $t$ du transformé pur de $D$ sur $S(t)$ (2.3). La fonction de $t$ $\mu_t(D)$ a pour support l'ensemble des points infiniment voisins de $s$ qui appartiennent à $D$ (2.4). Lorsque $D$ est régulier, les nombres $\mu_t(D).[k(t) : k(s)]$ sont égaux à zéro ou un : en effet, les transformés pur $D'$ de $D$ sont isomorphes à $D$, et on applique la définition (2.6) des multiplicités.

}{2.8. Reprenons les hypothèses de 2.0. Soit $D$ un diviseur $\geq 0$ sur $S$ et $t \in S^*$. La \underline{multiplicité} $\mu_t(D)$ \underline{de} $D$ \underline{en} $t$ est, par définition, la multiplicité en $t$ du transformé pur de $D$ sur $S(t)$ (2.3). La fonction de $t$ $\mu_t(D)$ a pour support l'ensemble des points infiniment voisins de $s$ qui appartiennent à $D$ (2.4). Lorsque $D$ est régulier, les nombres $\mu_t(D).[k(t) : k(s)]$ sont égaux à zéro ou un : en effet, les transformés purs $D'$ de $D$ sont isomorphes à $D$, et on applique la définition (2.6) des multiplicités.

}{F7II-X-015}Les nombres d'intersections se calculent en termes de multiplicités :

\underline{Théorème} 2.9. \underline{Si les diviseurs} $D$ \underline{et} $E$ \underline{ne se coupent qu'en} $s$, \underline{on a}

\begin{equation*}
(D,E) \;=\; \sum_{t \in S^*} \mu_t(D)\mu_t(E)[k(t):k]
\end{equation*}
\underline{la somme étant étendue aux points infiniment voisins de} $s$ \underline{communs à} $D$ \underline{et} $E$.

\underline{Preuve}. Soit $p : S' \longrightarrow S$ un transformé birationnel de $S$. On désigne par $D^*$ et $E^*$ (resp. par $D'$ et $E'$) les transformés totaux (resp. purs) de $D$ et $E$.


\SGApage{II}{18}

\underline{Lemme} 2.9.1. (i) \underline{Le transformé total} $D^*$ \underline{est caractérisé par} :

\SGAvar{(a) $D^* - D'$ \underline{est concentré sur} $f^{-1}(s)$

(b) \underline{si} $X$ \underline{est concentré sur} $f^{-1}(s)$, \underline{alors} $(D^*,X) = 0$

}{(a) $D^* - D'$ \underline{est concentré sur} $p^{-1}(s)$

(b) \underline{si} $X$ \underline{est concentré sur} $p^{-1}(s)$, \underline{alors} $(D^*,X) = 0$

}{F7II-X-016}(ii) $(D,E) = (D^*,E^*) = (D^*,E') = (D',E^*)$.

\SGAvar{L'assertion (i)(a) est triviale, et (i)(b) résulte de 1.6 (vi). Que (a) et (b) déterminent $D^*$ résulte de 1.9. Par construction, le conoyau de la flèche de $\underline{O}_E$ dans $f_* \underline{O}_E$ est fini, d'où

\begin{equation*}
(D^*,E') = \chi(f^*\underline{O}_D \otimes \underline{O}_{E'}) = \chi(\underline{O}_D \otimes f_* \underline{O}_{E'}) = \chi(\underline{O}_D \otimes \underline{O}_E) = (D,E)
\end{equation*}
ce qui, par (i), entraîne (ii).

}{L'assertion (i)(a) est triviale, et (i)(b) résulte de 1.6 (vi). Que (a) et (b) déterminent $D^*$ résulte de 1.9. Par construction, le conoyau de la flèche de $\underline{O}_E$ dans $p_* \underline{O}_{E'}$ est fini, d'où

\begin{equation*}
(D^*,E') = \chi(p^*\underline{O}_D \otimes \underline{O}_{E'}) = \chi(\underline{O}_D \otimes p_* \underline{O}_{E'}) = \chi(\underline{O}_D \otimes \underline{O}_E) = (D,E)
\end{equation*}
ce qui, par (i), entraîne (ii).

}{F7II-X-017}Prenons pour $S'$ l'éclaté de $S$ en $s$. D'après (2.9.1), (2.6.3) et 2.7, on a

\SGAvar{\begin{equation}\tag{2.9.2}
\begin{aligned} (D,E) &= (D^*,E') = (D' + \mu_s(S)\, s^*,E')\\ &= (D',E') + \mu_s(D)(s^*,E') = (D',E') + \mu_s(D)\mu_s(E)[k(s):k] \;. \end{aligned}
\end{equation}
Le théorème 2.9. se prouve maintenant par récurrence sur $(D,E)$, par une application itérée de (2.9.2).

}{\begin{equation}\tag{2.9.2}
\begin{aligned} (D,E) &= (D^*,E') = (D' + \mu_s(D)\, s^*,E')\\ &= (D',E') + \mu_s(D)(s^*,E') = (D',E') + \mu_s(D)\mu_s(E)[k(s):k] \;. \end{aligned}
\end{equation}
Le théorème 2.9. se prouve maintenant par récurrence sur $(D,E)$, par une application itérée de (2.9.2).

}{F7II-X-018}(2.10) Soient $A$ un anneau local noethérien réduit de dimension un, $s$ le point fermé de $\mathrm{Spec}(A)$, $\tilde{A}$ le normalisé de $A$ et $k$ un corps dont $k(s)$ soit extension finie. Supposons que $\tilde{A}$ soit fini sur $A$ (i.e. que le complété $A^{\wedge}$ soit réduit). On appellera \underline{invariant} $\delta$ \underline{de la singularité de} $\mathrm{Spec}(A)$ \underline{en} $s$ la longueur, calculée sur $k$, du $A$-module $\tilde{A}/A$.


\SGApage{II}{19}

\underline{Théorème} 2.11. \underline{Soit} $D$ \underline{une courbe réduite tracée sur} $S$. \underline{Alors} :

(i) \underline{Quel que soit} $t \in S^*$, \underline{on a}, \underline{avec les notations} 2.8

\begin{equation*}
\mu_t(D) \;=\; \sum_{t' \text{ proche de } t} \mu_{t'}(D)[k(t'):k(t)] \qquad .
\end{equation*}
(ii) \underline{Supposons le normalisé} $\widetilde{D}$ \underline{de} $D$ \underline{fini sur} $D$, \underline{et soit} $\delta = \chi(\underline{O}_{\widetilde{D}}/\underline{O}_D)$ \underline{l'invariant de la singularité de} $D$ \underline{en} $s$. \underline{Pour tous les} $t \in S^*$, \underline{à l'exception d'un nombre fini}, $\mu_t(D)$ \underline{est égal à zéro ou un}, \underline{et on a}

\begin{equation*}
\delta \;=\; \sum_{t \in S^*} \frac{1}{2}\,\mu_t(D)\,(\mu_t(D)-1)\,[k(t):k] \qquad .
\end{equation*}
Soient $p : S_1 \longrightarrow S$ l'éclaté de $S$ le long de $\{s\}$, $D^*$ le transformé total de $D$ et $D'$ son transformé pur. En vertu de (2.7), on a

\begin{equation}\tag{2.11.1}
D^* = D' + \mu_s(D)\, s^* \qquad \text{et}
\end{equation}
\begin{equation}\tag{2.11.2}
\mu_s(D) = (s^*,D')\ .\ [k(s):k]^{-1} \qquad .
\end{equation}
La curvetta $s^*$ est non singulière, de sorte que, par 2.8 et la définition 2.5, la formule 2.9, appliquée aux diviseurs $s^*$ et $D'$ de $S_1$, s'écrit

\SGAvar{\begin{equation}\tag{2.11.3}
(s^*,D) \;=\; \sum_{t' \text{ proche de } s} \mu_{t'}(D)\ [k(t'):k] \qquad .
\end{equation}
L'assertion (i) pour $s = t$ résulte de (2.11.2) et (2.11.3), et le cas général s'obtient en appliquant ce résultat au localisé en $t$ de $S(t)$ (2.3).

}{\begin{equation}\tag{2.11.3}
(s^*,D') \;=\; \sum_{t' \text{ proche de } s} \mu_{t'}(D)\ [k(t'):k] \qquad .
\end{equation}
L'assertion (i) pour $s = t$ résulte de (2.11.2) et (2.11.3), et le cas général s'obtient en appliquant ce résultat au localisé en $t$ de $S(t)$ (2.3).

}{F7II-X-019}Désignons par $m$ l'idéal maximal qui définit $\{s\}$, et soit $n \geq 0$. On sait (SGA VII 3.5) que

\begin{equation}\tag{2.11.4}
p_*\ \underline{O}_{S_1}(-n\ s^*) = m^n \quad \text{et} \quad R^i p_*\ \underline{O}_{S_1}(-n\ s^*) = 0 \quad (i > 0) \qquad .
\end{equation}

\SGApage{II}{20}

On déduit de là, et de la suite exacte

\begin{equation*}
0 \longrightarrow \underline{O}_{S_1}(-n\ s^*) \longrightarrow \underline{O}_{S_1} \longrightarrow \underline{O}_{n\,s^*} \longrightarrow 0
\end{equation*}
que

\begin{equation}\tag{2.11.5}
p_*\ \underline{O}_{n\,s^*} \;=\; O_S/m^n \quad \text{et} \quad R^i p_*\ \underline{O}_{n\,s^*} \;=\; 0 \quad (i > 0)
\end{equation}
et que

\begin{equation}\tag{2.11.6}
\chi(\underline{O}_{n\,s^*}) \;=\; \chi(\underline{O}_S/m^n) \;=\; \frac{n(n+1)}{2} \ [k(s):k] \ .
\end{equation}
Faisons $n = 0$ dans (2.11.4). Quel que soit le complexe $K$ de faisceaux cohérents sur $S$, on a par la ``formule de projection''

\SGAvar{\begin{equation*}
K \xrightarrow{\ \sim\ } Rp_*\ Lp^*K \qquad (\simeq K \overset{\mathbb{L}}{\otimes} Rp_*\ \underline{O}_{S_1}) \qquad .
\end{equation*}
En particulier, puisque $Lf^*\underline{O}_D = f^*\ \underline{O}_D = \underline{O}_{D^*}$, on a

}{\begin{equation*}
K \xrightarrow{\ \sim\ } Rp_*\ Lp^*K \qquad (\simeq K \overset{\mathbb{L}}{\otimes} Rp_*\ \underline{O}_{S_1}) \qquad .
\end{equation*}
En particulier, puisque $Lp^*\underline{O}_D = p^*\ \underline{O}_D = \underline{O}_{D^*}$, on a

}{F7II-X-020}\begin{equation}\tag{2.11.7}
p_*\ \underline{O}_{D^*} = \underline{O}_D \quad \text{et} \quad R^i p_*\underline{O}_{D^*} = 0 \quad (i>0) \qquad .
\end{equation}
Posons $E = \mu_s(D)\ s^*$ . D'après (2.11.1), (2.9.1) (i) b et (2.6.1) on a

\begin{equation}\tag{2.11.8}
\chi(\underline{O}_{D'\cap E}) = (D',E) = (D^* - E,E) = -(E,E) = -\mu_s(D)^2(s^*,s^*) = \mu_s(D)^2\ [k(s) : k] \qquad .
\end{equation}
La suite exacte

\begin{equation*}
0 \longrightarrow O_{D^*} \longrightarrow O_{D'} \oplus \underline{O}_E \longrightarrow \underline{O}_{D'\cap E} \longrightarrow 0
\end{equation*}
fournit par image directe, grâce à (2.11.7) des suites exactes

\SGAvar{\begin{equation*}
0 \longrightarrow O_D \longrightarrow f_*\underline{O}_{D'} \oplus f_*\underline{O}_E \longrightarrow f_*\underline{O}_{D'\cap E} \longrightarrow 0
\end{equation*}
et

\begin{equation*}
0 \longrightarrow f_*\underline{O}_{D'}/\underline{O}_D \longrightarrow f_*\underline{O}_E \longrightarrow f_*\ \underline{O}_{D'\cap E} \longrightarrow 0 \qquad .
\end{equation*}

}{\begin{equation*}
0 \longrightarrow O_D \longrightarrow p_*\underline{O}_{D'} \oplus p_*\underline{O}_E \longrightarrow p_*\underline{O}_{D'\cap E} \longrightarrow 0
\end{equation*}
et

\begin{equation*}
0 \longrightarrow p_*\underline{O}_E \longrightarrow p_*\ \underline{O}_{D'\cap E} \longrightarrow p_*\underline{O}_{D'}/\underline{O}_D \longrightarrow 0 \qquad .
\end{equation*}

}{F7II-X-021}\SGApage{II}{21}

D'après (2.11.6) et (2.11.8), on a donc

\SGAvar{\begin{equation}\tag{2.11.9}
\begin{aligned}
\chi(f_*\underline{O}_{D'}/\underline{O}_D) &= \frac{1}{2}\,\mu_s(D)(\mu_s(D)+1)[k(s): k] - \mu_s(D)^2\,[k(s) : k] \\
&= \frac{1}{2}\,\mu_s(D)(\mu_s(D)-1)[k(s) : k] \quad .
\end{aligned}
\end{equation}
Les schémas $D$ et $D'$ sont réduits, ne diffèrent qu'en $s$, et $D'$ est fini sur $D$. Ces schémas ont donc même normalisé, et l'invariant $\delta$ de la singularité de $D$ en $s$ est somme de $\chi(f_*\underline{O}_{D'}/\underline{O}_D)$ et de la somme des invariants $S$ des singularités de $D'$ en les points de $D' \cap s^*$. Prouvons 2.11 (ii) par récurrence sur $\delta$. Si $\mu_s(D) > 1$, l'hypothèse de récurrence s'applique à $D'$, et 2.11 (ii) résulte de (2.11.9). Si $\mu_s(D) = 1$, alors $(D',s^*) = 1$ et $D' \cap s^*$ est un diviseur régulier de $D'$ ; par (EGA IV 19.1.1), $D'$ est régulier ; par (2.11.9), $D = D'$ est lui-même régulier et 2.11 (ii) résulte de 2.8.

}{\begin{equation}\tag{2.11.9}
\begin{aligned}
\chi(p_*\underline{O}_{D'}/\underline{O}_D) &= \mu_s(D)^2\,[k(s) : k] - \frac{1}{2}\,\mu_s(D)(\mu_s(D)+1)[k(s): k] \\
&= \frac{1}{2}\,\mu_s(D)(\mu_s(D)-1)[k(s) : k] \quad .
\end{aligned}
\end{equation}
Les schémas $D$ et $D'$ sont réduits, ne diffèrent qu'en $s$, et $D'$ est fini sur $D$. Ces schémas ont donc même normalisé, et l'invariant $\delta$ de la singularité de $D$ en $s$ est somme de $\chi(p_*\underline{O}_{D'}/\underline{O}_D)$ et de la somme des invariants $\delta$ des singularités de $D'$ en les points de $D' \cap s^*$. Prouvons 2.11 (ii) par récurrence sur $\delta$. Si $\mu_s(D) > 1$, l'hypothèse de récurrence s'applique à $D'$, et 2.11 (ii) résulte de (2.11.9). Si $\mu_s(D) = 1$, alors $(D',s^*) = 1$ et $D' \cap s^*$ est un diviseur régulier de $D'$ ; par (EGA IV 19.1.1), $D'$ est régulier ; par (2.11.9), $D = D'$ est lui-même régulier et 2.11 (ii) résulte de 2.8.

}{F7II-X-022}2.12. Soit $t \in S^*$ et $S'$ un transformé birationnel de $S$ qui domine $S_+(t)$ (2.3). On appellera \underline{transformé total} de $t$ sur $S'$, et on désignera encore par $t^*$, le transformé total sur $S'$ du diviseur $t^*$ de $S_+(t)$.

Soit $\nu$ une fonction à support fini de $S^*$ dans $\mathbf{N}$ et supposons que $\nu$ vérifie, quel que soit $t \in S^*$

\begin{equation}\tag{2.12.1}
\nu(t) \;\geq\; \sum_{t' \ \mathrm{proche\ de}\ t} \nu(t')\,[k(t') : k(t)]
\end{equation}
(inégalité des points proches). Soient $f : S' \longrightarrow S$ un transformé birationnel de $S$ qui domine les $S_+(t)$ pour $t$ dans le support de $\nu$, et $\nu^*$ le diviseur $\sum_{t \in S^*} \nu(t)\ t^*$ sur $S'$.


\SGApage{II}{22}

On désignera par $I(\nu)$ l'idéal de $\underline{O}_S$ image directe par $f$ de l'idéal de $\underline{O}_{S'}$ qui définit le diviseur $\nu^*$ :

\begin{equation*}
I(\nu) \;=\; f_*\underline{O}(-\nu^*) \quad .
\end{equation*}
D'après (2.11.4), cet idéal ne dépend pas du choix de $S'$. Les idéaux de $S$ obtenus par ce procédé sont ceux que Zariski [11] appelle ``complets''.

\underline{Théorème} 2.13. \underline{Sous les hypothèses précédentes}, \underline{on a}

\begin{equation*}
\chi(\underline{O}_S/I(\nu)) \;=\; \sum_{t \in S^*} \frac{\nu(t)(\nu(t)+1)}{2}\,[k(t) : k]
\end{equation*}
\underline{et}

\begin{equation*}
R^i f_*\,\underline{O}_{S'}(-\nu^*) \;=\; 0 \quad \underline{\mathrm{pour}}\ i > 0 \quad .
\end{equation*}
En vieux français, on exprime ceci, brièvement, en disant que : ``exiger qu'une courbe plane passe $n$ fois par un point $t$, projectif ou infiniment voisin, impose $\frac{n(n+1)}{2}$ conditions, en général indépendantes, aux coefficients de son équation''.

Prouvons tout d'abord le lemme

\SGAvar{\underline{Lemme} 2.14. \underline{Soient} $p : S_1 \longrightarrow S$ \underline{l'éclaté de} $S$ \underline{le long de} $\{s\}$, $T$ \underline{un sous-schéma fermé de} $S_1$ , \underline{d'idéal} $q$, $D$ \underline{un diviseur}, $n$ \underline{un entier}, $J = O(-s^*)$ \underline{l'idéal de} $S_1$ \underline{qui définit le diviseur} $s^*$, $Y = \mathrm{Spec}(\underline{O}_{S_1}/J^n q)$, $D$ \underline{un diviseur régulier sur} $S$ \underline{et} $D'$ \underline{son transformé pur sur} $S_1$. \underline{On suppose} $T \cap s^*$ \underline{et} $T \cap D'$ \underline{finis}.

}{\underline{Lemme} 2.14. \underline{Soient} $p : S_1 \longrightarrow S$ \underline{l'éclaté de} $S$ \underline{le long de} $\{s\}$, $T$ \underline{un sous-schéma fermé de} $S_1$ , \underline{d'idéal} $q$, $n$ \underline{un entier}, $J = O(-s^*)$ \underline{l'idéal de} $S_1$ \underline{qui définit le diviseur} $s^*$, $Y = \mathrm{Spec}(\underline{O}_{S_1}/J^n q)$, $D$ \underline{un diviseur régulier sur} $S$ \underline{et} $D'$ \underline{son transformé pur sur} $S_1$. \underline{On suppose} $T \cap s^*$ \underline{et} $T \cap D'$ \underline{finis}.

}{F7II-X-023}(i) \underline{Si} $\mathrm{lg}_{k(s)}\,(T \cap s^*) \leq n +1$, \underline{on a}

\begin{equation*}
R^i p_*\,(q.J^n) \;=\; 0 \quad \underline{\mathrm{pour}}\ i > 0 \ \underline{\mathrm{et}},\ \underline{\mathrm{pour}}\ T\ \underline{\mathrm{artinien}},
\end{equation*}
\begin{equation*}
\mathrm{lg}_{k(s)}(\underline{O}_S/p_*(q\,J^n)) \;=\; \frac{n(n+1)}{2} \;+\; \mathrm{lg}_{k(s)}\,(T)
\end{equation*}

\SGApage{II}{23}

(ii)  \underline{Si} $\mathrm{lg}_{k(s)} (T \cap s^*) \leq n$, \underline{on a}

\begin{equation*}
\mathrm{lg}_{k(s)} (D \cap \mathrm{Spec}(\underline{O}_S/p_*(q\,J^n))) \; = \; n + \mathrm{lg}_{k(s)} (D' \cap T) \quad .
\end{equation*}
Soit la suite exacte sur $S_1$

\SGAvar{\begin{equation}\tag{2.14.1}
0 \longrightarrow J^n. q \longrightarrow J^n \stackrel{u}{\longrightarrow} J^n \otimes \underline{O}_T \longrightarrow 0 \quad .
\end{equation}
Pour prouver que $R^ip_* J^n.q = 0$  $(i > 0)$, il suffit par (2.11.4), de montrer que la flèche $p_*(u)$ est surjective. D'après le théorème des fonctions holomorphes, il suffit de prouver que les flèches

}{\begin{equation}\tag{2.14.1}
0 \longrightarrow J^n. q \longrightarrow J^n \stackrel{u}{\longrightarrow} J^n \otimes \underline{O}_T \longrightarrow 0 \quad .
\end{equation}
Pour prouver que $R^ip_* J^n.q = 0$  $(i > 0)$, il suffit par (2.11.4), de montrer que la flèche $p_*(u)$ est surjective. D'après le théorème des fonctions formelles, il suffit de prouver que les flèches

}{F7II-X-024}\begin{equation*}
u_k \; : \; p_*(J^n \otimes \underline{O}/J^k) \longrightarrow p_*(J^n \otimes \underline{O}_T \otimes \underline{O}/J^k)
\end{equation*}
sont surjectives.

\SGAvar{On désignera par $\mathrm{Gr}_J$ un gradué associé pour la filtration $J$-adique. Le module $J^m/J^{m+1}$ est isomorphe à $\underline{O}_{s^*}(m)$ ; on a donc $R^1p_*J^m/J^{m+1} = 0$ et la flèche de $p_*(J^n/J^{n+k+1})$ dans $p_*(J^n/J^{n+k})$ est surjective. Pour prouver, par récurrence sur $k$, que les flèches $u_k$ sont surjectives, il suffit donc de prouver, la surjectivité des $v_k$ :

}{On désignera par $\mathrm{Gr}_J$ un gradué associé pour la filtration $J$-adique. Le module $J^m/J^{m+1}$ est isomorphe à $\underline{O}_{s^*}(m)$ ; on a donc $R^1p_*J^m/J^{m+1} = 0$ et la flèche de $p_*(J^n/J^{n+k+1})$ dans $p_*(J^n/J^{n+k})$ est surjective. Pour prouver, par récurrence sur $k$, que les flèches $u_k$ sont surjectives, il suffit donc de prouver la surjectivité des $v_k$ :

}{F7II-X-025}\begin{equation*}
v_k \; : \; p_*(J^n \otimes \mathrm{Gr}_J^k(\underline{O})) \longrightarrow p_*(J^n \otimes \mathrm{Gr}_J^k(\underline{O}_T)) \quad (k \geq 0) \quad .
\end{equation*}
Le module $\mathrm{Gr}_J^k(\underline{O}_T)$ est un quotient de $J^k/J^{k+1} \otimes \underline{O}_T$ ; il est donc de longueur sur $k(s)$ au plus $n+1$. Le noyau $N$ de la flèche surjective

\begin{equation*}
v_k' \; : \; J^n \otimes \mathrm{Gr}_J^k(\underline{O}) \sim \underline{O}_{s^*}(\,n+k) \longrightarrow J^n \otimes \mathrm{Gr}_J^k(\underline{O}_T)
\end{equation*}
est donc un faisceau inversible sur la droite projective $s^*$, de degré $\geq -1$. On a $R^1p_* N = 0$ et $v_k = p_*(v_k')$ est surjectif.

Si $T$ est artinien, de la suite exacte (2.14.1), on déduit maintenant que

\begin{equation*}
\mathrm{lg}_{k(s)} (p_*J^n/p_*(J^nq)) = \mathrm{lg}_{k(s)} (\underline{O}_T) \quad , 
\end{equation*}

\SGApage{II}{24}

 et 2.14 (i) résulte de (2.11.4) et (2.11.6).

\SGAvar{Soit $I_{D'\cap T}$ l'idéal (monogène de $\underline{O}_{D'}$ qui définit $D'\cap T$. Sous les hypothèses (ii), la suite exacte

\begin{equation*}
0 \longrightarrow I_{D'\cap T} \longrightarrow \underline{O}_{D'\cap T} \longrightarrow O_T \longrightarrow 0
\end{equation*}
montre que $D'\cap T$ vérifie l'hypothèse (i).

La projection $p$ établit un isomorphisme entre $D$ et $D'$ ; par cet isomorphisme, à l'idéal de $\underline{O}_D$ qui définit le sous-schéma $D \cap \mathrm{Spec}(O_S/p_*(q\,J^n))$ de $D$ correspond l'idéal de $D'$ engendré par les sections globales de $q\,J^n$ sur $S_1$. Cet idéal est contenu dans l'idéal de $D'$ qui définit $D'\cap Y$. Il lui est en fait égal, car toute section de $\underline{O}_{D'\cup Y}$ se prolonge en une section de $\underline{O}_{S_1}$ : si $I_{D'\cup Y}$ est l'idéal qui définit $D'\cup Y$, l'obstruction ou prolongement se trouve dans $R^1p_*(I_{D'\cup Y})$. On a

}{Soit $I_{D'\cap T}$ l'idéal (monogène) de $\underline{O}_{D'}$ qui définit $D'\cap T$. Sous les hypothèses (ii), la suite exacte

\begin{equation*}
0 \longrightarrow I_{D'\cap T} \longrightarrow \underline{O}_{D'\cup T} \longrightarrow O_T \longrightarrow 0
\end{equation*}
montre que $D'\cup T$ vérifie l'hypothèse (i).

La projection $p$ établit un isomorphisme entre $D$ et $D'$ ; par cet isomorphisme, à l'idéal de $\underline{O}_D$ qui définit le sous-schéma $D \cap \mathrm{Spec}(O_S/p_*(q\,J^n))$ de $D$ correspond l'idéal de $D'$ engendré par les sections globales de $q\,J^n$ sur $S_1$. Cet idéal est contenu dans l'idéal de $D'$ qui définit $D'\cap Y$. Il lui est en fait égal, car toute section de $\underline{O}_{D'\cup Y}$ se prolonge en une section de $\underline{O}_{S_1}$ : si $I_{D'\cup Y}$ est l'idéal qui définit $D'\cup Y$, l'obstruction au prolongement se trouve dans $R^1p_*(I_{D'\cup Y})$. On a

}{F7II-X-026}\begin{equation*}
\begin{array}{rl} I_{D'\cup Y} & = \underline{O}(-D')\cap q\,J^n = (\underline{O}(-D')\cap J^n) \cap q\,J^n = \underline{O}(-D')J^n \cap q\,J^n \\[2mm] & = (\underline{O}(-D')\cap q). J^n \end{array}
\end{equation*}
et, d'après (i) appliqué à $D'\cup T$, il n'y a pas d'obstruction.

On a donc

\begin{equation}\tag{2.14.2}
\mathrm{lg}_{k(s)} (D\cap \mathrm{Spec}(\underline{O}_S/p_*(q\,J^n))) = \mathrm{lg}_{k(s)} (D'\cap Y) \quad .
\end{equation}
De l'isomorphie

\begin{equation*}
J^n \otimes \underline{O}_{D'}/q.\underline{O}_{D'} \stackrel{\sim}{\longrightarrow} \underline{O}_{D'}.J^n/\underline{O}_{D'}.q.J^n
\end{equation*}
on déduit que la suite suivante est exacte

\SGAvar{\begin{equation*}
0 \longrightarrow J^n \otimes_{D'\cap T} \longrightarrow \underline{O}_{D'\cap Y} \longrightarrow \underline{O}_{D'\cap\ n\ s^*} \longrightarrow 0 \quad ,
\end{equation*}
de sorte que 


}{\begin{equation*}
0 \longrightarrow J^n \otimes \underline{O}_{D'\cap T} \longrightarrow \underline{O}_{D'\cap Y} \longrightarrow \underline{O}_{D'\cap\ n\ s^*} \longrightarrow 0 \quad ,
\end{equation*}
de sorte que 


}{F7II-X-027}\SGApage{II}{25}

\begin{equation}\tag{2.14.3}
\mathrm{lg}_{k(s)}\ (D' \cap Y) = n\ (D',s^*) + \mathrm{lg}_{k(s)}(D' \cap T) \quad .
\end{equation}
Puisque $(D',s^*) = 1$, 2.14 (ii) résulte de (2.14.2) et (2.14.3).

Prouvons maintenant, par récurrence sur le nombre d'éléments dans le support de $\nu$, les assertions de 2.13 et l'assertion suivante :

Si $D$ est un diviseur régulier sur $S$, alors,

\SGAvar{\begin{equation}\tag{2.14.4}
\chi(D \cap \mathrm{Spec}(\underline{O}_S/I(\nu)) = \sum_{t\ \mathrm{sur}\ D} \nu(t)\ [k(t) : k] \quad .
\end{equation}
Si $\nu = 0$, alors $\nu^* = 0$, $I(\nu) = \underline{O}_S$ et 2.13 et (2.14.4) sont triviaux (la seconde assertion de (2.13) résulte d'une application itérée de (2.11.4) pour $n = 0$). Sinon, on a $\nu(s) > 0$, et il existe un $S$-morphisme $g : S' \longrightarrow S_1$. Pour chaque point fermé $a$ de $S_1$, soient $S_{1a}$ le localisé de $S_1$ en $a$, $\nu_a$ la restriction de $\nu$ à $S_{1a}^*$ , $S'_a$ le localisé $S'_a = S' \times_{S_1} S_{1a}$, $I(\nu_a) = g_* \underline{O}_{S'}(-\nu_a^*)$, $T_a = \mathrm{Spec}(\underline{O}_{S_1}/I(\nu_a))$, $T = \cup\ T_a$ (réunion finie disjointe) et $q$ l'idéal de $T$. On a, avec les notations de 2.14

}{\begin{equation}\tag{2.14.4}
\chi(D \cap \mathrm{Spec}(\underline{O}_S/I(\nu))) = \sum_{t\ \mathrm{sur}\ D} \nu(t)\ [k(t) : k] \quad .
\end{equation}
Si $\nu = 0$, alors $\nu^* = 0$, $I(\nu) = \underline{O}_S$ et 2.13 et (2.14.4) sont triviaux (la seconde assertion de (2.13) résulte d'une application itérée de (2.11.4) pour $n = 0$). Sinon, on a $\nu(s) > 0$, et il existe un $S$-morphisme $g : S' \longrightarrow S_1$. Pour chaque point fermé $a$ de $S_1$, soient $S_{1a}$ le localisé de $S_1$ en $a$, $\nu_a$ la restriction de $\nu$ à $S_{1a}^*$ , $S'_a$ le localisé $S'_a = S' \times_{S_1} S_{1a}$, $I(\nu_a) = g_* \underline{O}_{S'}(-\nu_a^*)$, $T_a = \mathrm{Spec}(\underline{O}_{S_1}/I(\nu_a))$, $T = \cup\ T_a$ (réunion finie disjointe) et $q$ l'idéal de $T$. On a, avec les notations de 2.14

}{F7II-X-028}\begin{equation}\tag{2.14.5}
g_* \underline{O}(-\nu^*) = q.J^{\nu(s)} \quad .
\end{equation}
L'hypothèse de récurrence s'applique aux $(S_{1a}, S'_a, \nu_a)$. On a donc

\begin{equation}\tag{2.14.6}
R^i g_* \underline{O}_{S'}(-\nu^*) = 0 \ \mathrm{pour}\ i > 0
\end{equation}
\begin{equation}\tag{2.14.7}
\chi(T) = \sum_{t \in S^* \setminus \{s\}} \frac{\nu(t)\ (\nu(t)+1)}{2} [k(t) : k]
\end{equation}
et, d'après (2.14.4) appliqué au diviseur $s^*$ de $S_1$ et l'hypothèse (2.12.1)

\begin{equation*}
\mathrm{lg}_{k(s)}\ (T \cap s^*) = \sum_{t\ \mathrm{proche\ de}\ s} \nu(t)\ [k(t) : k(s)] \leq \nu(s) \quad .
\end{equation*}

\SGApage{II}{26}

Cette inégalité nous permet d'appliquer 2.14 ; de (2.14.5), on tire

\begin{equation*}
R^i p_* \ g_* \ \underline{O}(-\nu^*) = 0 \ \mathrm{pour}\ i > 0
\end{equation*}
\begin{equation*}
\mathrm{lg}_{k(s)}(\underline{O}_S/p_* \ g_* \ \underline{O}(-\nu^*)) = \frac{\nu(s)(\nu(s)+1)}{2} + \mathrm{lg}_{k(s)}\ (T)
\end{equation*}
de sorte que les assertions (2.13) résultent de (2.14.6) et (2.14.7). Enfin, 2.14 (ii) fournit

\SGAvar{\begin{equation*}
\mathrm{lg}_{k(s)}\ (D \cap \mathrm{Spec}\ (\underline{O}_S/I(\nu)) = \nu(s) + \mathrm{lg}_{k(s)}\ (D' \cap T) \ ,
\end{equation*}
soit par (2.14.4) appliqué à $D'$,

\begin{equation*}
\begin{aligned}
\chi(D \cap \mathrm{Spec}(\underline{O}_S/I(\nu)) &= \nu(s)\ [k(s) : k] + \sum_{t\ \mathrm{sur}\ D} \nu(t)\ [k(t) : k] \\
&= \sum_{t\ \mathrm{sur}\ D} \nu(t)\ [k(t) : k] \quad .
\end{aligned}
\end{equation*}
Ceci achève la démonstration de 2.13.

(2.15) Des bribes de ce qui précèdent se généralisent en dimension supérieure. Soit $S$ le spectre d'un anneau local régulier de dimension $n$ de point fermé $s$, $k(s)$ étant une extension finie de $k$. On peut encore considérer l'ensemble $I$ des schémas $f : S' \longrightarrow S$ déduit de $S$ par éclatements successifs de points fermés et poser comme en 2.1

\begin{equation*}
S^* = \varinjlim_{s' \in I} \ \mathrm{Irr}(f^{-1}(s)
\end{equation*}
$S^*$ s'appellera encore l'ensemble des \underline{points de} $S$ \underline{infiniment voisins de} $s$.

}{\begin{equation*}
\mathrm{lg}_{k(s)}\ (D \cap \mathrm{Spec}\ (\underline{O}_S/I(\nu))) = \nu(s) + \mathrm{lg}_{k(s)}\ (D' \cap T) \ ,
\end{equation*}
soit par (2.14.4) appliqué à $D'$,

\begin{equation*}
\begin{aligned}
\chi(D \cap \mathrm{Spec}(\underline{O}_S/I(\nu))) &= \nu(s)\ [k(s) : k] + \sum_{\substack{t\ \mathrm{sur}\ D\\ t\ne s}} \nu(t)\ [k(t) : k] \\
&= \sum_{t\ \mathrm{sur}\ D} \nu(t)\ [k(t) : k] \quad .
\end{aligned}
\end{equation*}
Ceci achève la démonstration de 2.13.

(2.15) Des bribes de ce qui précède se généralisent en dimension supérieure. Soit $S$ le spectre d'un anneau local régulier de dimension $n$ de point fermé $s$, $k(s)$ étant une extension finie de $k$. On peut encore considérer l'ensemble $I$ des schémas $f : S' \longrightarrow S$ déduit de $S$ par éclatements successifs de points fermés et poser comme en 2.1

\begin{equation*}
S^* = \varinjlim_{S' \in I} \ \mathrm{Irr}(f^{-1}(s))
\end{equation*}
$S^*$ s'appellera encore l'ensemble des \underline{points de} $S$ \underline{infiniment voisins de} $s$.

}{F7II-X-029}Si $X$ est un sous-schéma fermé de $S$, et si $t \in S^*$, on définit comme en 2.8 la multiplicité de $X$ en $t$.

Soient $f : S' \longrightarrow S$ un composé de morphismes d'éclatement de points fermés, $D$ un diviseur sur $S$, $D^* = f^* D$ son transformé total sur $S'$,
\SGApage{II}{27}
$E$ une courbe sur $S$, $E'$ son transformé pur sur $S'$ et supposons que $D \cap E$ soit réduit à $s$. La formule de projection donne encore

\begin{equation}\tag{2.16}
(D,E) \;=\; (D^*,E')
\end{equation}
Procédant comme en 2.9, on prouve à l'aide de (2.16) que

\begin{equation}\tag{2.17}
(D,E) \;=\; \sum_{t \in S^*} \mu_t(D) \, \mu_t(E) \, [k(t) : k] \qquad .
\end{equation}
Définissons les points proches comme en 2.5. Procédant comme en 2.11, on vérifie que

\SGAvar{\begin{equation}\tag{2.18}
\mu_t(E) \;=\; \sum_{t' \text{ proche de } t} \mu_{t'}(E) \, [k(t') : k] \qquad .
\end{equation}
Par contre, la formule 2.11 (ii) devient fausse en général (voir les exemples 2.25).

}{\begin{equation}\tag{2.18}
\mu_t(E) \;=\; \sum_{t' \text{ proche de } t} \mu_{t'}(E) \, [k(t') : k(t)] \qquad .
\end{equation}
Par contre, la formule 2.11 (ii) devient fausse en général (voir les exemples 2.25).

}{F7II-X-030}2.19. Soit $S$ une surface projective et régulière sur un corps $k$. On désigne par $S^*$ la somme disjointe, étendue aux points fermés de $S$

\begin{equation*}
S^* \;=\; \coprod_{s \in S} \operatorname{Spec} \, (\underline{O}_{s,S})^* \qquad ,
\end{equation*}
et on appelle $S^*$ l'ensemble des \underline{points infiniment voisins des} points de $S$.

\SGAvar{Soit $I$ l'ensemble des surfaces déduites de $S$ par éclatements successifs de points fermés. L'ensemble $I$ est cofinal dans l'ensemble des modèles réguliers du corps des fonctions de $S$, ordonné par domination. Pour $i \in I$, on désigne par $S_i$ la surface correspondantes, munie de $p_i : S_i \to S$. On désignera par $\mathrm{Div}^{\pm}(S_i)$ le groupe des diviseurs, positifs ou non, sur $S_i$.

}{Soit $I$ l'ensemble des surfaces déduites de $S$ par éclatements successifs de points fermés. L'ensemble $I$ est cofinal dans l'ensemble des modèles réguliers du corps des fonctions de $S$, ordonné par domination. Pour $i \in I$, on désigne par $S_i$ la surface correspondante, munie de $p_i : S_i \to S$. On désignera par $\mathrm{Div}^{\pm}(S_i)$ le groupe des diviseurs, positifs ou non, sur $S_i$.

}{F7II-X-031}\underline{Définition} 2.20. (i) \underline{Le groupe des} ``courbes'' \underline{sur} $S$ \underline{est l'ensemble limite inductive}

\begin{equation*}
\varinjlim_{i \in I} \mathrm{Div}^{\pm}(S_i)
\end{equation*}
(ii) \underline{La catégorie des} ``faisceaux inversibles'' \underline{sur} $S$ \underline{est la limite inductive}
\SGApage{II}{28}
(SGA 4 VI 6) \underline{des catégories de faisceaux inversibles sur les} $(S_i)_{i \in I}$ .

Ces objets sont des invariants birationnels de $S$.

\underline{Exemple} 2.21. (i) Chaque ``courbe'' $C$ sur $S$ définit un ``faisceau inversible'' $\underline{O}(C)$ sur $S$.

(ii) Chaque diviseur $C$ sur $S$ définit une ``courbe'' sur $S$, qu'on désignera par $C^*$.

(iii) Si $t \in S^*$, on définit comme en 2.3 la surface $S_+(t)$ et le diviseur $t^*$ sur $S_+(t)$ ; on désigne encore par $t^*$ la ``courbe'' définie par $t^*$.

(2.22) On vérifie facilement qu'une ``courbe'' $C$ (resp. un ``faisceau inversible'' $\mathcal{L}$) s'écrit d'une et d'une seule façon sous la forme

\begin{equation*}
C \;=\; C_o^* \;-\; \sum_{t \in S^*} i_t \, t^* \qquad \text{resp.} \qquad \mathcal{L} \;=\; \mathcal{L}_o \otimes \underline{O} \left( -\sum_{t \in S^*} i_t \, t^* \right)
\end{equation*}
où $C_o$ (resp. $\mathcal{L}_o$) est un diviseur (resp. un faisceau inversible) sur $S$ et où les $i_t$ sont des nombres entiers presque tous nuls. On dit que $C$ (resp. $\mathcal{L}$) est la courbe $C_o$, (resp. le faisceau inversible $\mathcal{L}_o$) affectée de multiplicités virtuelles $i_t$.

Le complété $\mathrm{Div}^{\pm} (S) \times \mathbb{Z}^{S^*}$ du groupe $\mathrm{Div}^{\pm} (S) \times \mathbb{Z}^{(S^*)}$ des ``courbes'' sur $S$ est encore un invariant birationnel de $S$ ; on l'appelle le groupe des ``courbes généralisées'' sur $S$. Sur $S$, une ``courbe généralisée'' $C$ se représente par un diviseur $C_o$, affecté d'un nombre pouvant être infini de multiplicités virtuelles.

On définit de même les ``faisceaux inversibles généralisés'' sur $S$.


\SGApage{II}{29}

\SGAvar{\underline{Exemple} 2.23. Supposons $k$ parfait. Alors, le faisceau inversible $\Omega^2_{S/k}$ , affecté d'une multiplicité virtuelle $-1$ en chaque point $t \in S^*$, et un invariant birationnel de $S$, soit $\Omega^2_{S/k}{}^*$ .

\underline{Exemple} 2.24.  Soit $f : S_1 \longrightarrow S_2$ un morphisme dominant de surfaces projectives régulières sur $k$. On définit sans difficulté l'image inverse d'une "courbe" ou d'un "faisceau inversible" sur $S_2$, en passant à un modèle de $S_2$ où cette "courbe" ou "faisceau inversible" est diviseur ou faisceau inversible usuel. Par passage à la limite (je n'ai pas vérifié si ce passage à la limite était possible), on définit de même l'image réciproque d'une "courbe généralisée", ou d'un"faisceau inversible généralisé" . Supposons $k$ parfait et $f$ génériquement étale. La \underline{différente} $\mathfrak{d}(f)$ de $f$ est alors le diviseur positif sur $S_1$ qui vérifie

\begin{equation*}
f^* \Omega^2_{S_2/k} (\mathfrak{d}(f)) \overset{\sim}{\longrightarrow} \Omega^2_{S_1/k} \qquad .
\end{equation*}
Cette différente n'est pas un invariant birationnel de $f$ ; pour obtenir un invariant birationnel, il faut considérer la "courbe généralisée" de coîncidences

}{\underline{Exemple} 2.23. Supposons $k$ parfait. Alors, le faisceau inversible $\Omega^2_{S/k}$ , affecté d'une multiplicité virtuelle $-1$ en chaque point $t \in S^*$, est un invariant birationnel de $S$, soit $\Omega^2_{S/k}{}^*$ .

\underline{Exemple} 2.24.  Soit $f : S_1 \longrightarrow S_2$ un morphisme dominant de surfaces projectives régulières sur $k$. On définit sans difficulté l'image inverse d'une "courbe" ou d'un "faisceau inversible" sur $S_2$, en passant à un modèle de $S_2$ où cette "courbe" ou "faisceau inversible" est diviseur ou faisceau inversible usuel. Par passage à la limite (je n'ai pas vérifié si ce passage à la limite était possible), on définit de même l'image réciproque d'une "courbe généralisée", ou d'un "faisceau inversible généralisé" . Supposons $k$ parfait et $f$ génériquement étale. La \underline{différente} $\mathfrak{d}(f)$ de $f$ est alors le diviseur positif sur $S_1$ qui vérifie

\begin{equation*}
f^* \Omega^2_{S_2/k} (\mathfrak{d}(f)) \overset{\sim}{\longrightarrow} \Omega^2_{S_1/k} \qquad .
\end{equation*}
Cette différente n'est pas un invariant birationnel de $f$ ; pour obtenir un invariant birationnel, il faut considérer la "courbe généralisée" de coïncidences

}{F7II-X-032}\begin{equation*}
\mathfrak{d}^*(f) = \mathfrak{d}(f) + \sum_{t \in S_1^*} t^* - f^* ( \sum_{t \in S_2^*} t^*)
\end{equation*}
soit, brièvement

\SGAvar{\begin{equation*}
\mathfrak{d}^*(f) = \Omega^2_{S_1/k}{}^* - f^* \Omega^2_{S_2/k}{}^* \qquad .
\end{equation*}
Cette "courbe généralisée" à un support contenu dans le fermé de $S_2$ où $f$ n'est pas étale.


}{\begin{equation*}
\mathfrak{d}^*(f) = \Omega^2_{S_1/k}{}^* - f^* \Omega^2_{S_2/k}{}^* \qquad .
\end{equation*}
Cette "courbe généralisée" a un support contenu dans le fermé de $S_1$ où $f$ n'est pas étale.


}{F7II-X-033}\SGApage{II}{30}

(2.25) \underline{Exemple}.

\SGAvar{Soient $k$ un corps algébriquement clos et $E_n$ l'espace affine à $n$ dimension sur $k$. Soit $C$ une courbe réduite dans $E_n$, présentant une singularité à l'origine. A cette singularité on associe son "\underline{diagramme d'Enriques}" (cf. [3]),qui est le graphe de sommets les points infiniment voisins de $O$ \underline{sur} $C$ (2.4 et 2.15), chaque sommet étant joint par un trait à ceux qui lui sont proches (2.5 et 2.15). Ce diagramme détermine les multiplicités de ces points sur $C$ (2.11 et 2.18) et, pour $n = 2$, l'invariant $\delta$ (2.10 et 2.11) de la singularité de $C$ à l'origine.

}{Soient $k$ un corps algébriquement clos et $E_n$ l'espace affine à $n$ dimensions sur $k$. Soit $C$ une courbe réduite dans $E_n$, présentant une singularité à l'origine. À cette singularité on associe son "\underline{diagramme d'Enriques}" (cf. [3]), qui est le graphe de sommets les points infiniment voisins de $O$ \underline{sur} $C$ (2.4 et 2.15), chaque sommet étant joint par un trait à ceux qui lui sont proches (2.5 et 2.15). Ce diagramme détermine les multiplicités de ces points sur $C$ (2.11 et 2.18) et, pour $n = 2$, l'invariant $\delta$ (2.10 et 2.11) de la singularité de $C$ à l'origine.

}{F7II-X-034}Pour les singularités planes les plus simples, on trouve :

\begin{center}
\small
\setlength{\tabcolsep}{5pt}
\renewcommand{\arraystretch}{1.35}
\begin{tabular}{c c c c}
Dessin & équation & $\delta$ & diagramme (avec multiplicités) \\

\begin{tikzpicture}[baseline=-.5ex,x=.55cm,y=.55cm]
  \draw[line width=.45pt] (-.8,0)--(.8,0);
\end{tikzpicture}
& $y=0$ & $0$ &
\begin{tikzpicture}[baseline=-.5ex,x=.56cm,y=.56cm]
  \foreach \x in {0,...,6}{
    \fill (\x,0) circle (1.05pt);
    \node[above=1pt,scale=.72] at (\x,0) {$1$};
  }
  \draw (0,0)--(6,0);
  \node[right=2pt,scale=.85] at (6,0) {$\cdots$};
\end{tikzpicture}
\\








\begin{tikzpicture}[baseline=-.5ex,x=.55cm,y=.55cm]
  \draw[line width=.45pt] (-.75,-.42)--(.75,.42);
  \draw[line width=.45pt] (-.75,.42)--(.75,-.42);
\end{tikzpicture}
& $xy=0$ & $1$ &
\begin{tikzpicture}[baseline=-.5ex,x=.72cm,y=.62cm]
  \fill (0,0) circle (1.05pt);
  \node[left=1pt,scale=.72] at (0,0) {$2$};
  \foreach \x in {1,...,5}{
    \fill (\x,.48) circle (1.05pt);
    \fill (\x,-.48) circle (1.05pt);
    \node[above=1pt,scale=.72] at (\x,.48) {$1$};
    \node[above=1pt,scale=.72] at (\x,-.48) {$1$};
  }
  \draw (0,0)--(1,.48)--(5,.48);
  \draw (0,0)--(1,-.48)--(5,-.48);
  \node[right=2pt,scale=.85] at (5,.48) {$\cdots$};
  \node[right=2pt,scale=.85] at (5,-.48) {$\cdots$};
\end{tikzpicture}
\\











\begin{tikzpicture}[baseline=-.5ex,x=.55cm,y=.55cm]
  \draw[line width=.45pt] (-.85,0)--(.85,0);
  \draw[line width=.45pt,domain=-.72:.72,samples=40]
    plot (\x,{.72*\x*\x});
\end{tikzpicture}
& $y(y-x^2)=0$ & $2$ &
\begin{tikzpicture}[baseline=-.5ex,x=.68cm,y=.62cm]
  \fill (0,0) circle (1.05pt);
  \fill (.9,0) circle (1.05pt);
  \node[above=1pt,scale=.72] at (0,0) {$2$};
  \node[above=1pt,scale=.72] at (.9,0) {$2$};
  \foreach \x/\lab in {1.8/1,2.7/1,3.6/1,4.5/0}{
    \fill (\x,.48) circle (1.05pt);
    \fill (\x,-.48) circle (1.05pt);
    \ifnum\lab=1
      \node[above=1pt,scale=.72] at (\x,.48) {$1$};
      \node[below=1pt,scale=.72] at (\x,-.48) {$1$};
    \fi
  }
  \draw (0,0)--(.9,0)--(1.8,.48)--(4.5,.48);
  \draw (.9,0)--(1.8,-.48)--(4.5,-.48);
  \node[right=2pt,scale=.72] at (4.5,-.48) {$\cdot\ \cdot$};
\end{tikzpicture}
\\












\begin{tikzpicture}[baseline=-.5ex,x=.55cm,y=.55cm]
  \draw[line width=.45pt]
    (0,0) .. controls (.28,.05) and (.55,.35) .. (.72,.78);
  \draw[line width=.45pt]
    (0,0) .. controls (.28,-.05) and (.55,-.35) .. (.72,-.78);
\end{tikzpicture}
& $y^2=x^3$ & $1$ &
\begin{tikzpicture}[baseline=-.5ex,x=.72cm,y=.65cm]
  \coordinate (A) at (0,.42);
  \coordinate (B) at (.9,.42);
  \coordinate (C) at (1.7,-.25);
  \coordinate (D) at (2.6,-.25);
  \coordinate (E) at (3.5,-.25);
  \foreach \p in {A,B,C,D,E}{\fill (\p) circle (1.05pt);}
  \node[above=1pt,scale=.72] at (A) {$2$};
  \foreach \p in {B,C,D,E}{\node[above=1pt,scale=.72] at (\p) {$1$};}
  \draw (A)--(B)--(C)--(D)--(E);
  \draw (A)--(C);
  \node[right=2pt,scale=.85] at (E) {$\cdots$};
\end{tikzpicture}
\end{tabular}
\end{center}











Si $n > 2$, l'invariant $\delta$ n'est plus déterminé par ce diagramme. Ainsi, la réunion de 3 droites concourantes coplanaires (resp. non coplanaires) a un invariant $\delta = 3$ (resp. $\delta = 2$), alors que le diagramme d'Enriques est le même. J'ignore si ce phénomène se présente aussi pour des courbes unibranches.

Voici deux exemples de diagrammes, relatifs au cas $n = 3$, qui ne peuvent se présenter pour des courbes planes (cf. 2.5).

On donne successivement les équations de la courbe, les équations
\SGApage{II}{31}
de la transformée pure après éclatement de l'origine, des équations paramétriques de la courbe, l'invariant $\delta$ et le diagramme d'Enriques.

\SGAvar{\begin{equation}\tag{2.25.1}
\left\{
\begin{array}{l}
\left\{\begin{array}{l} y^{2} = xz^{4} \\ z^{3} = x^{2} \end{array}\right. \;\;,\;\;
\left\{\begin{array}{l} (y/z)^{2} = (x/z)^{3} \\ z \;\;\;\;\;\; = (x/z)^{2} \end{array}\right. \;\;,\;\;
\left\{\begin{array}{l} x = u^{6} \\ y = u^{7} \\ z = u^{4} \end{array}\right. \;\;,\;\; \delta = 5 \;\;,
\end{array}
\right.
\end{equation}
\begin{center}
\begin{tikzcd}[column sep=small, row sep=small]
4 \arrow[r, dash] \arrow[rrd, dash] \arrow[rrdd, dash] \& 2 \arrow[rd, dash] \arrow[rdd, dash] \& {} \& {} \& {} \& {} \\
{} \& {} \& 1 \arrow[d, dash] \& {} \& {} \& {} \\
{} \& {} \& 1 \arrow[rd, dash] \& {} \& {} \& {} \\
{} \& {} \& {} \& 1 \arrow[rd, dash] \& {} \& {} \\
{} \& {} \& {} \& {} \& 1 \& {} \\
{} \& {} \& {} \& {} \& {} \& \cdots
\end{tikzcd}
\end{center}













}{\begin{equation}\tag{2.25.1}
\left\{
\begin{array}{l}
\left\{\begin{array}{l} y^{2} = xz^{2} \\ z^{3} = x^{2} \end{array}\right. \;\;,\;\;
\left\{\begin{array}{l} (y/z)^{2} = (x/z)^{3} \\ z \;\;\;\;\;\; = (x/z)^{2} \end{array}\right. \;\;,\;\;
\left\{\begin{array}{l} x = u^{6} \\ y = u^{7} \\ z = u^{4} \end{array}\right. \;\;,\;\; \delta = 5 \;\;,
\end{array}
\right.
\end{equation}
\begin{center}
\begin{tikzcd}[column sep=small, row sep=small]
4 \arrow[r, dash] \arrow[rrd, dash] \arrow[rrdd, dash] \& 2 \arrow[rd, dash] \arrow[rdd, dash] \& {} \& {} \& {} \& {} \\
{} \& {} \& 1 \arrow[d, dash] \& {} \& {} \& {} \\
{} \& {} \& 1 \arrow[rd, dash] \& {} \& {} \& {} \\
{} \& {} \& {} \& 1 \arrow[rd, dash] \& {} \& {} \\
{} \& {} \& {} \& {} \& 1 \& {} \\
{} \& {} \& {} \& {} \& {} \& \cdots
\end{tikzcd}
\end{center}













}{F7II-X-035}\begin{equation}\tag{2.25.2}
\left\{
\begin{array}{l}
\left\{\begin{array}{l} y^{2}z = x^{3} \\ z^{3} \;\; = y^{2} \end{array}\right. \;\;,\;\;
\left\{\begin{array}{l} (y/z)^{2} = (x/z)^{3} \\ z \;\;\;\;\;\; = (y/z)^{2} \end{array}\right. \;\;,\;\;
\left\{\begin{array}{l} x = u^{8} \\ y = u^{9} \\ z = u^{6} \end{array}\right. \;\;,\;\; \delta = 10 \;\;,
\end{array}
\right.
\end{equation}
\begin{center}
\begin{tikzcd}[column sep=small, row sep=small]
6 \arrow[r, dash] \arrow[rrd, dash] \arrow[rrdd, dash] \arrow[rrddd, dash] \arrow[rrdddd, dash] \& 2 \arrow[rd, dash] \arrow[rdd, dash] \& {} \& {} \& {} \& {} \\
{} \& {} \& 1 \arrow[d, dash] \& {} \& {} \& {} \\
{} \& {} \& 1 \arrow[d, dash] \& {} \& {} \& {} \\
{} \& {} \& 1 \arrow[d, dash] \& {} \& {} \& {} \\
{} \& {} \& 1 \arrow[rd, dash] \& {} \& {} \& {} \\
{} \& {} \& {} \& 1 \arrow[rd, dash] \& {} \& {} \\
{} \& {} \& {} \& {} \& 1 \& {} \\
{} \& {} \& {} \& {} \& {} \& \cdots
\end{tikzcd}
\end{center}














2.26. Le théorème 2.27 ci-dessous était énoncé sous la forme d'une conjecture fausse dans la première version (diffusée par l'IHES) de cet exposé. Il doit beaucoup à des conversations avec D.S. Rim, dont il englobe certains résultats (D.S. Rim, torsion differentials and deformations, theorem 2.7 (i)).


\SGApage{II}{32}

Soit  $C$  une courbe projective réduite sur un corps algébriquement clos  $k$ , et  $s \in C(k)$. On désignera par

$A_{s}$ : l'anneau local de  $C$  en  $s$  (ou son complété, cela reviendrait au même pour ce qui suit) ;

$\widetilde{A}_{s}$ : le normalisé de  $A_{s}$ ;

$\delta(s):\dim_{k}(\widetilde{A}_{s}/A_{s})$  ;

$m_{o}(s)$ : la dimension du schéma versel formel  $T$  des déformations de $\mathrm{Spec}(A_{s})$.

Pour  $E$  une composante irréductible de  $T$ , on posera encore $m_{o}^{E}(s) = \dim(E)$.

\SGAvar{Les modules  $\widetilde{\tau}$  des dérivations de  $\widetilde{A}_{s}$  et  $\tau$  des dérivations de  $A_{s}$  sont sans torsion ; ils s'injectent dans le module des dérivations du corps des fractions de  $A_{s}$ , pour  $A$  intègre, et dans le module des dérivations de l'anneau total des fractions de  $A_{s}$  (un produit de corps) dans le cas général. De plus, via ces inclusions,  $\tau$  et  $\widetilde{\tau}$  sont commensurables : $\tau/\tau\cap\widetilde{\tau}$  et  $\widetilde{\tau}/\tau\cap\widetilde{\tau}$  sont de dimension finie sur  $k$ .

}{Les modules  $\widetilde{\tau}$  des dérivations de  $\widetilde{A}_{s}$  et  $\tau$  des dérivations de  $A_{s}$  sont sans torsion ; ils s'injectent dans le module des dérivations du corps des fractions de  $A_{s}$ , pour  $A_s$  intègre, et dans le module des dérivations de l'anneau total des fractions de  $A_{s}$  (un produit de corps) dans le cas général. De plus, via ces inclusions,  $\tau$  et  $\widetilde{\tau}$  sont commensurables : $\tau/\tau\cap\widetilde{\tau}$  et  $\widetilde{\tau}/\tau\cap\widetilde{\tau}$  sont de dimension finie sur  $k$ .

}{F7II-X-036}On pose

\begin{equation*}
m_{1}(s) = [\widetilde{\tau} : \tau] \mathop{=}\limits_{\mathrm{dfn}} \dim_{k}(\widetilde{\tau}/\tau \cap \widetilde{\tau}) - \dim_{k}(\tau/\tau\cap\widetilde{\tau}) \quad .
\end{equation*}
\underline{Théorème} 2.27.  \underline{Soit}  $E$  \underline{une composante irréductible du schéma formel versel}  $T$  \underline{des déformations de}  $A_{s}$ . \underline{Si}  "\underline{la déformation de}  $A_{s}$  \underline{au-dessus du point générique}  $e$  \underline{de}  $E$  \underline{est lisse}", \underline{i.e.} \underline{si l'image du sous-schéma formel de non lissité} (\underline{défini par un idéal jacobien} : VI 5), \underline{fini sur}  $T$ , \underline{ne contient pas}  $e$ , \underline{alors}

\begin{equation*}
3\ \delta(s) = m_{o}^{E}(s) + m_{1}(s) \quad .
\end{equation*}

\SGApage{II}{33}

La validité de 2.27 ne dépend que du complété de $A_s$ . Remplaçant $(C,s)$ par $(C',s)$ formellement isomorphe à $(C,s)$ en $s$ , on peut supposer que les conditions suivantes sont vérifiées

(a)\quad $s$ est le seul point singulier de $C$ ;

(b)\quad $C$ n'a pas d'automorphismes infinitésimaux.

\SGAvar{La condition (b) nous garantit que la foncteur des déformations de $C$ est proreprésentable. Il est même effectivement pro-représentable, car il n'y a pas d'obstruction à relever un faisceau inversible (ample) donné sur $C$ . Soit $f : X_o \longrightarrow T_o$ la déformation universelle de $C$ et $t_o$ le point fermé de $T_o$ : $C \stackrel{\sim}{\relbar\joinrel\relbar} X_{ot_o}$.

\underline{Lemme} 2.28. \quad \underline{Il existe un schéma de type fini sur} $k$ $T_1$, $t_1 \in T_1(k)$ \underline{et} $f : X_1 \longrightarrow T_1$, \underline{propre et plat,} \underline{de fibre} $C$ \underline{en} $t_1$ \underline{et tel que pour tout} $u \in T_1(k)$, \underline{le complété de} $T_1$ \underline{en} $u$ \underline{soit un schéma formel de modules pour la fibre} $X_{1u}$ .

Voici une démonstration directe de ce corollaire de résultats généraux d'Artin. Soit $\mathcal{L}$ un faisceau inversible très ample sur $C$ , tel que $H^1(C,\mathcal{L}) = 0$, et $\mathbb{P} = \mathbb{P}(H^o(C,\mathcal{L}))$. On a $v : C \hookrightarrow \mathbb{P}$ . Soit $T_2$ le schéma de Hilbert, classifiant les sous-variétés de $\mathbb{P}$ de même polynôme de Hilbert que $v(C) \subset \mathbb{P}$ . Soit $t_2 \in T_2$ correspondant à $v(C)$. Le complété $\hat{T}_2$ de $T_2$ en $t_2$ est lisse sur $S_o$ , le fibré tangent vertical de $\hat{T}_2/T_o$ se déduisant du sous-fibré $L$ du fibré tangent de $T_2$ image de $\mathrm{Lie}(\underline{\mathrm{Aut}}(\mathbb{P}))$ pour l'action de $\underline{\mathrm{Aut}}(\mathbb{P})$ sur $T_2$. Au voisinage de $t_2$, soit $f_1 \ldots f_n$ $n = \dim \underline{\mathrm{Aut}}(\mathbb{P})$ équations définissant un sous-schéma $T'_1$ de $T_2$ passant par $t_2$ et de complété en $t_2$ étale sur $T_o$ . On prend pour $T_1$ un voisinage de $t$ dans $T'_1$ où le fibré tangent soit transverse à $L$ .


}{La condition (b) nous garantit que le foncteur des déformations de $C$ est proreprésentable. Il est même effectivement pro-représentable, car il n'y a pas d'obstruction à relever un faisceau inversible (ample) donné sur $C$ . Soit $f : X_o \longrightarrow T_o$ la déformation universelle de $C$ et $t_o$ le point fermé de $T_o$ : $C \stackrel{\sim}{\relbar\joinrel\relbar} X_{ot_o}$.

\underline{Lemme} 2.28. \quad \underline{Il existe un schéma} $T_1$ \underline{de type fini sur} $k$, \underline{un point} $t_1 \in T_1(k)$ \underline{et un morphisme} $f : X_1 \longrightarrow T_1$, \underline{propre et plat,} \underline{de fibre} $C$ \underline{en} $t_1$ \underline{et tel que pour tout} $u \in T_1(k)$, \underline{le complété de} $T_1$ \underline{en} $u$ \underline{soit un schéma formel de modules pour la fibre} $X_{1u}$ .

Voici une démonstration directe de ce corollaire de résultats généraux d'Artin. Soit $\mathcal{L}$ un faisceau inversible très ample sur $C$ , tel que $H^1(C,\mathcal{L}) = 0$, et $\mathbb{P} = \mathbb{P}(H^0(C,\mathcal{L}))$. On a $v : C \hookrightarrow \mathbb{P}$ . Soit $T_2$ le schéma de Hilbert, classifiant les sous-variétés de $\mathbb{P}$ de même polynôme de Hilbert que $v(C) \subset \mathbb{P}$ . Soit $t_2 \in T_2$ correspondant à $v(C)$. Le complété $\hat{T}_2$ de $T_2$ en $t_2$ est lisse sur $T_o$ , le fibré tangent vertical de $\hat{T}_2/T_o$ se déduisant du sous-fibré $L$ du fibré tangent de $T_2$ image de $\mathrm{Lie}(\underline{\mathrm{Aut}}(\mathbb{P}))$ pour l'action de $\underline{\mathrm{Aut}}(\mathbb{P})$ sur $T_2$. Au voisinage de $t_2$, soient $f_1, \ldots, f_n$ des équations, avec $n = \dim \underline{\mathrm{Aut}}(\mathbb{P})$, définissant un sous-schéma $T'_1$ de $T_2$ passant par $t_2$ et de complété en $t_2$ étale sur $T_o$ . On prend pour $T_1$ un voisinage de $t_2$ dans $T'_1$ où le fibré tangent soit transverse à $L$ .


}{F7II-X-037}\SGApage{II}{34}

\SGAvar{2.29. Au morphisme de foncteur

}{2.29. Au morphisme de foncteurs

}{F7II-X-038}\begin{equation*}
(\text{déformations de } C) \longrightarrow (\text{déformations de } A_s)
\end{equation*}
correspond un morphisme $p : T_o \longrightarrow T$ . D'après VI 2.10 \qquad , $p$ est lisse. Soit $t$ le point fermé de $T$ . L'espace tangent de Zariski de $p^{-1}(t)$ en $t_o$ classifie les déformations localement triviales de $C$ sur $k[\varepsilon]$ ($\varepsilon^2 = 0$). Posant $E_o = p^{-1}(E)$, on a donc

\begin{equation}\tag{2.29.1}
\dim(E_o) = \dim(E) + \dim H^1(C, T^1_{C/k}) \qquad .
\end{equation}
Puisque $X_o/T_o$ est plat, les fibres de $X_o/T_o$ ont toutes le même genre arithmétique $p$ ($1-p = \chi(\mathcal{O})$). Par hypothèse, la fibre au point générique de $E_o$ est lisse ; le schéma formel des modules d'une courbe propre et lisse étant lisse de dimension $3p-3$, on déduit de 2.28 que

\SGAvar{\begin{equation}\tag{2.29.2}
\dim(E) \;=\; 3p-3 \qquad .
\end{equation}
Soit $\widetilde{C}$ la normalisée de $C$, de genre arithmétique $\widetilde{p} = p - \delta(s)$. On a

\begin{equation}\tag{2.29.3}
\begin{aligned}\dim\, H^1(C,T^1_{C/k}) \;&=\; -\,\chi(C,T^1_{C/k}) \;=\; \\ &=\; -\,\chi(\widetilde{C},T^1_{\widetilde{C}/k}) + m_1(s) \;=\; -[3(p-\delta(s))-3] + m_1(s) \qquad .\end{aligned}
\end{equation}
Les formules (2.29.1) (2.29.2) (2.29.3) prouvent 2.27.

}{\begin{equation}\tag{2.29.2}
\dim(E_o) \;=\; 3p-3 \qquad .
\end{equation}
Soit $\widetilde{C}$ la normalisée de $C$, de genre arithmétique $\widetilde{p} = p - \delta(s)$. On a

\begin{equation}\tag{2.29.3}
\begin{aligned}\dim\, H^1(C,T^1_{C/k}) \;&=\; -\,\chi(C,T^1_{C/k}) \;=\; \\ &=\; -\,\chi(\widetilde{C},T^1_{\widetilde{C}/k}) + m_1(s) \;=\; [3(p-\delta(s))-3] + m_1(s) \qquad .\end{aligned}
\end{equation}
Les formules (2.29.1) (2.29.2) (2.29.3) prouvent 2.27.

}{F7II-X-039}\underline{Conjecture} 2.30. \quad \underline{Toute courbe propre réduite a une déformation lisse.}

Cette conjecture implique, par 2.27 et 2.28, que le schéma formel versel des déformations de $A_s$ est toujours équidimensionnel, de dimension $3\delta(s) - m_1(s)$.


\SGApage{II}{35}

2.31. Soit $C$ une courbe propre réduite sur $k$, de genre arithmétique $p$ . On désignera par

$M_o$ : la dimension du schéma formel versel $T$ des déformations de $C$ ;

$M_1$ : la dimension de l'algèbre de Lie de $\underline{\mathrm{Aut}}(C)$.

Pour $E$ une composante irréductible de $T$, on notera $M_o^E$ la dimension de $E$ .

L'analogue global de 2.27 est le suivant.

\underline{Proposition} 2.32. \underline{Si la déformation de} $C$ \underline{au-dessus du point générique de} $E$ \underline{est lisse, on a}

\begin{equation*}
3p-3 = M_o^E - M_1 \quad .
\end{equation*}
Soit $X \subset C(k)$ un ensemble fini de points lisses de $C$ , tel que tout champ de vecteur sur $C$, s'annulant en tout point $x \in X$, soit nul. Le foncteur des déformations de $(C,X)$ est proreprésentable, et vérifie une variante de 2.28 . Le schéma formel des modules de $(C,X)$ est lisse sur $T$ , de dimension relative $\# X - M_1$ , et de l'analogue de 2.28 on déduit que

\SGAvar{\begin{equation*}
M_o + ( \# X - M_1) \;=\; 3p - 3 + \# X \; ,
\end{equation*}
d'où 2.32.

}{\begin{equation*}
M_o^E + ( \# X - M_1) \;=\; 3p - 3 + \# X \; ,
\end{equation*}
d'où 2.32.

}{F7II-X-040}Supposons $k$ de \underline{caractéristique} $0$ . On a alors $M_1 = \dim \underline{\mathrm{Aut}}(C)$. Voici une interprétation analogue de $m_1(s)$.


\SGApage{II}{36}

\underline{Lemme} 2.33. \underline{Si} $k$ \underline{est de caractéristique} $0$, \underline{toute dérivation de} $A_s$ \underline{se prolonge en une dérivation de} $A_s^{\sim}$ .

\SGAvar{Si $D$ est une dérivation de $A_s$ , $\exp(Dt) = \sum D^n t^n/n!$ est un automorphisme du complété en $(s,0)$ de $A_s \otimes_k k[t]$. Par transport de structure, cet automorphisme agit sur le normalisé $(A_s^{\sim} \otimes k[t])^{\wedge}$ de cet anneau. Réduisant mod $t^2$ , on trouve un automorphisme $1 + Dt$ de $A_s^{\sim\wedge} \otimes_k k[t]/(t^2)$, et $D$ se trouve ainsi prolongé en une dérivation de $A_s^{\sim\wedge}$ . Puisqu'il se prolonge aussi au corps des fractions de $A_s$ (anneau total des fractions si $A_s$ n'est pas intègre), il se prolonge à $A_s^{\sim}$ .

Supposons $s$ singulier. Un argument analogue montre alors que le champ de vecteur sur $\mathrm{Spec}(A_s^{\sim})$ défini par $D$ s'annulle en tout point au-dessus de $s$ .

}{Si $D$ est une dérivation de $A_s$ , $\exp(Dt) = \sum_{n\geq0} D^n t^n/n!$ est un automorphisme du complété en $(s,0)$ de $A_s \otimes_k k[t]$. Par transport de structure, cet automorphisme agit sur le normalisé $(A_s^{\sim} \otimes k[t])^{\wedge}$ de cet anneau. Réduisant mod $t^2$ , on trouve un automorphisme $1 + Dt$ de $A_s^{\sim\wedge} \otimes_k k[t]/(t^2)$, et $D$ se trouve ainsi prolongé en une dérivation de $A_s^{\sim\wedge}$ . Puisqu'il se prolonge aussi au corps des fractions de $A_s$ (anneau total des fractions si $A_s$ n'est pas intègre), il se prolonge à $A_s^{\sim}$ .

Supposons $s$ singulier. Un argument analogue montre alors que le champ de vecteur sur $\mathrm{Spec}(A_s^{\sim})$ défini par $D$ s'annule en tout point au-dessus de $s$ .

}{F7II-X-041}Soit $c$ un idéal de $A_s^{\sim}$ contenu dans le conducteur de $A_s^{\sim}/A_s$ . Soit $G_c^{\sim} = \underline{\mathrm{Aut}}(A_s^{\sim}/c)$ et soit $G_c$ le sous-groupe de $G_c^{\sim}$ qui stabilise la sous-algèbre $A_s/c$ . Pour $d \subset c$ , $G_d^{\sim}$ stabilise $c/d$ et s'envoie sur $G_c^{\sim}$ (remarquer que $g \in G_c^{\sim}$ est déterminé par $g(t)$, pour $t$ une uniformisante, et que $g(t)$ est n'importe quelle uniformisante). L'espace homogène pointé $G_c^{\sim}/G_c$ est donc indépendant de $c$ . On note $G^{\sim}/G$.

\SGAvar{Soit $\tau_o^{\sim}$ le $A_s^{\sim}$-module des champs de vecteurs sur $\mathrm{Spec}(A_s^{\sim})$ qui s'annullent en les points au-dessus de $s$ . Pour $s$ singulier, on

}{Soit $\tau_o^{\sim}$ le $A_s^{\sim}$-module des champs de vecteurs sur $\mathrm{Spec}(A_s^{\sim})$ qui s'annulent en les points au-dessus de $s$ . Pour $s$ singulier, on a

}{F7II-X-042}\begin{equation*}
\dim_k(\tau^{\sim}/\tau_o^{\sim}) \;=\; \mbox{nombre de branches} \;;
\end{equation*}
\begin{equation*}
\tau_o^{\sim}/\tau \;\simeq\; \mbox{espace tangent à l'origine de } G^{\sim}/G \;.
\end{equation*}

\SGApage{II}{37}

\underline{Proposition} 2.34. \underline{En caractéristique} $0$ \underline{et pour} $s$ \underline{singulier}, \underline{l'entier} $m_1(s)$ \underline{est somme du nombre de branches et de la dimension de} $\widetilde{G}/G$ .


\underline{BIBLIOGRAPHIE}

[1] \quad P. Defrise. Etude locale des correspondances rationnelles entre surfaces algébriques - Thèse d'agrégation - Bruxelles 1949.

[2] \quad P. Deligne et D. Mumford. The irreducibility of the space of curves of a given genus - Publ. Math. IHES \underline{36} (1969) p. 75-110.

\SGAvar{[3] \quad Enriques - Chisini. Lezioni sulla teoria geometrica delle equazioni et delle funzioni algebriche - ed Zanichelli - Bologne 1924.

[4] \quad R. Hartshorne. Residues and duality - Lectures Notes in Math. \underline{20} - Springer 1966.

[5] \quad S. Lichtenbaum. Curves over discrete valuation rings - Am. J. of Math. 15 n° 2 April 1968. p. 380-405.

}{[3] \quad Enriques - Chisini. Lezioni sulla teoria geometrica delle equazioni e delle funzioni algebriche - ed Zanichelli - Bologne 1924.

[4] \quad R. Hartshorne. Residues and duality - Lecture Notes in Math. \underline{20} - Springer 1966.

[5] \quad S. Lichtenbaum. Curves over discrete valuation rings - Am. J. of Math. 90 n° 2 April 1968. p. 380-405.

}{F7II-X-043}[6] \quad I. Manin. Rational surfaces over perfect fields - Publ. Math. IHES n° 30. p. 55-114.

\SGAvar{[7] \quad D. Mumford. The topology of normal singularities of an algebraic surfaces and a criterion for simplicity - Publ. Math. IHES n° 9 p. 5-22.

}{[7] \quad D. Mumford. The topology of normal singularities of an algebraic surface and a criterion for simplicity - Publ. Math. IHES n° 9 p. 5-22.

}{F7II-X-044}[8] \quad A. Néron. Modèles minimaux des variétés abéliennes sur les corps locaux et globaux - Publ. Math. n° 21.


\SGApage{II}{38}

[9] \quad M. Raynaud. Spécialisation du foncteur de Picard. Critères numériques de représentabilité - CR Acad. Sci. Paris t. 264 p. 1001 - 1004 - Juin 1967.

[10] \quad Saint-Donat. Sur un théorème de G. Castelnuovo et F. Enriques - Thèse de 3° cycle - Lyon.

\SGAvar{[11] \quad O. Zariski. On polynomial ideals defined by infinitely near base points Am. J. Math. vol. 60-1938 p. 151.

}{[11] \quad O. Zariski. On polynomial ideals defined by infinitely near base points, Am. J. Math. vol. 60-1938 p. 151.

}{F7II-X-045}[12] \quad D.G. Northcott. On the genus formula for plane curves. J. London Math. Soc. 30 1955 - p. 376-382.

[13] \quad D.G. Northcott. Abstract dilatations and infinitely near points. Proc. Cambridge Philos. Soc. 52 1956 - p. 178-197.

La notion de ``point infiniment voisin'' est discutée dans l'appendice par Lipman au chapitre II de

\SGAvar{[14] \quad O. Zariski. Algebraic surfaces - Ergebnisse \underline{61}, Springer Verlag.
}{[14] \quad O. Zariski. Algebraic surfaces - Ergebnisse \underline{61}, Springer-Verlag.
}{F7II-X-046}\selectlanguage{french}

\ifnum\SGAform=1
\clearpage\section*{Apparat critique}\small
Les lectures ci-dessous sont données dans le TeX exact des témoins, afin de préserver les distinctions de notation. Les indices décimaux renvoient aux différences individuelles du passage marqué.

\par\medskip\noindent\hypertarget{app-F7II-X-001.1}{}\textbf{F7II-X-001.1}\quad p. 2; misprint\_in\_print.\par
\noindent Imprimé : {\ttfamily\footnotesize coîncidant\allowbreak{}}\par
\noindent Corrigé : {\ttfamily\footnotesize coïncidant\allowbreak{}}\par
\noindent Source : FIX-7II-D-0002. \nolinkurl{fixes/FIXLOG_7II.jsonl#2}\par

\par\medskip\noindent\hypertarget{app-F7II-X-002.1}{}\textbf{F7II-X-002.1}\quad p. 2; printed misprint.\par
\noindent Imprimé : {\ttfamily\footnotesize faisceauxcchérents\allowbreak{}}\par
\noindent Corrigé : {\ttfamily\footnotesize faisceaux\allowbreak{} cohérents\allowbreak{}}\par
\noindent Source : 7II:1. \nolinkurl{fixes/MISPRINTS_7II.json#1}\par

\par\medskip\noindent\hypertarget{app-F7II-X-002.2}{}\textbf{F7II-X-002.2}\quad p. 2; misprint\_in\_print.\par
\noindent Imprimé : \emph{(vide)}\par
\noindent Corrigé : {\ttfamily\footnotesize \textbackslash{}underline\{\allowbreak{}}\par
\noindent Source : FIX-7II-D-0003. \nolinkurl{fixes/FIXLOG_7II.jsonl#3}\par

\par\medskip\noindent\hypertarget{app-F7II-X-002.3}{}\textbf{F7II-X-002.3}\quad p. 2; misprint\_in\_print.\par
\noindent Imprimé : \emph{(vide)}\par
\noindent Corrigé : {\ttfamily\footnotesize \}\allowbreak{}}\par
\noindent Source : FIX-7II-D-0003. \nolinkurl{fixes/FIXLOG_7II.jsonl#3}\par

\par\medskip\noindent\hypertarget{app-F7II-X-002.4}{}\textbf{F7II-X-002.4}\quad p. 2; misprint\_in\_print.\par
\noindent Imprimé : {\ttfamily\footnotesize \textbackslash{}mathbb\{\allowbreak{}}\par
\noindent Corrigé : \emph{(vide)}\par
\noindent Source : FIX-7II-D-0003. \nolinkurl{fixes/FIXLOG_7II.jsonl#3}\par

\par\medskip\noindent\hypertarget{app-F7II-X-002.5}{}\textbf{F7II-X-002.5}\quad p. 2; misprint\_in\_print.\par
\noindent Imprimé : {\ttfamily\footnotesize \}\allowbreak{}}\par
\noindent Corrigé : \emph{(vide)}\par
\noindent Source : FIX-7II-D-0003. \nolinkurl{fixes/FIXLOG_7II.jsonl#3}\par

\par\medskip\noindent\hypertarget{app-F7II-X-003.1}{}\textbf{F7II-X-003.1}\quad p. 5; misprint\_in\_print.\par
\noindent Imprimé : {\ttfamily\footnotesize \textbackslash{}\allowbreak{} }\par
\noindent Corrigé : {\ttfamily\footnotesize ,\allowbreak{} }\par
\noindent Source : FIX-7II-D-0004. \nolinkurl{fixes/FIXLOG_7II.jsonl#4}\par

\par\medskip\noindent\hypertarget{app-F7II-X-004.1}{}\textbf{F7II-X-004.1}\quad p. 6; misprint\_pass.\par
\noindent Imprimé : {\ttfamily\footnotesize fatiguant\allowbreak{}}\par
\noindent Corrigé : {\ttfamily\footnotesize fatigant\allowbreak{}}\par
\noindent Source : 7II:5. \nolinkurl{fixes/FIXLOG_7II_CURRENT.jsonl#508}\par

\par\medskip\noindent\hypertarget{app-F7II-X-004.2}{}\textbf{F7II-X-004.2}\quad p. 6; designated-witness attribution.\par
\noindent Imprimé : \emph{(vide)}\par
\noindent Corrigé : \emph{(1 espace(s), 0 retour(s) de ligne dans le témoin TeX)}\par
\noindent Source : WITNESS-F7II-X-004.2. \nolinkurl{fr/sga7ii/diplomatic/expose_X_body.tex}; \nolinkurl{fr/sga7ii/current/expose_X_body.tex}\par

\par\medskip\noindent\hypertarget{app-F7II-X-004.3}{}\textbf{F7II-X-004.3}\quad p. 6; printed misprint.\par
\noindent Imprimé : {\ttfamily\footnotesize S\allowbreak{}}\par
\noindent Corrigé : {\ttfamily\footnotesize X\allowbreak{}}\par
\noindent Source : 7II:4. \nolinkurl{fixes/MISPRINTS_7II.json#4}\par

\par\medskip\noindent\hypertarget{app-F7II-X-005.1}{}\textbf{F7II-X-005.1}\quad p. 7; printed misprint.\par
\noindent Imprimé : {\ttfamily\footnotesize D\allowbreak{}}\par
\noindent Corrigé : {\ttfamily\footnotesize S\allowbreak{}}\par
\noindent Source : 7II:6. \nolinkurl{fixes/MISPRINTS_7II.json#6}\par

\par\medskip\noindent\hypertarget{app-F7II-X-005.2}{}\textbf{F7II-X-005.2}\quad p. 7; inherited corrected-layer difference.\par
\noindent Imprimé : \emph{(vide)}\par
\noindent Corrigé : \emph{(1 espace(s), 0 retour(s) de ligne dans le témoin TeX)}\par
\noindent Source : ADOPT-F7II-X-005.2. \nolinkurl{fr/sga7ii/diplomatic/expose_X_body.tex}; \nolinkurl{fr/sga7ii/current/expose_X_body.tex}\par

\par\medskip\noindent\hypertarget{app-F7II-X-005.3}{}\textbf{F7II-X-005.3}\quad p. 7; printed misprint.\par
\noindent Imprimé : {\ttfamily\footnotesize quelque\allowbreak{}}\par
\noindent Corrigé : {\ttfamily\footnotesize quel\allowbreak{} que\allowbreak{}}\par
\noindent Source : 7II:7. \nolinkurl{fixes/MISPRINTS_7II.json#7}\par

\par\medskip\noindent\hypertarget{app-F7II-X-006.1}{}\textbf{F7II-X-006.1}\quad p. 8; designated-witness attribution.\par
\noindent Imprimé : \emph{(vide)}\par
\noindent Corrigé : {\ttfamily\footnotesize |E)\allowbreak{}}\par
\noindent Source : WITNESS-F7II-X-006.1. \nolinkurl{fr/sga7ii/diplomatic/expose_X_body.tex}; \nolinkurl{fr/sga7ii/current/expose_X_body.tex}\par

\par\medskip\noindent\hypertarget{app-F7II-X-007.1}{}\textbf{F7II-X-007.1}\quad p. 9; designated-witness attribution.\par
\noindent Imprimé : \emph{(vide)}\par
\noindent Corrigé : {\ttfamily\footnotesize au-dessus\allowbreak{} }\par
\noindent Source : WITNESS-F7II-X-007.1. \nolinkurl{fr/sga7ii/diplomatic/expose_X_body.tex}; \nolinkurl{fr/sga7ii/current/expose_X_body.tex}\par

\par\medskip\noindent\hypertarget{app-F7II-X-007.2}{}\textbf{F7II-X-007.2}\quad p. 9; designated-witness attribution.\par
\noindent Imprimé : {\ttfamily\footnotesize \_1\allowbreak{}}\par
\noindent Corrigé : \emph{(vide)}\par
\noindent Source : WITNESS-F7II-X-007.2. \nolinkurl{fr/sga7ii/diplomatic/expose_X_body.tex}; \nolinkurl{fr/sga7ii/current/expose_X_body.tex}\par

\par\medskip\noindent\hypertarget{app-F7II-X-007.3}{}\textbf{F7II-X-007.3}\quad p. 9; misprint\_in\_print.\par
\noindent Imprimé : {\ttfamily\footnotesize possèdant\allowbreak{}}\par
\noindent Corrigé : {\ttfamily\footnotesize possédant\allowbreak{}}\par
\noindent Source : FIX-7II-D-0008. \nolinkurl{fixes/FIXLOG_7II.jsonl#8}\par

\par\medskip\noindent\hypertarget{app-F7II-X-008.1}{}\textbf{F7II-X-008.1}\quad p. 11; inherited corrected-layer difference.\par
\noindent Imprimé : {\ttfamily\footnotesize exclus\allowbreak{}}\par
\noindent Corrigé : {\ttfamily\footnotesize exclu\allowbreak{}}\par
\noindent Source : ADOPT-F7II-X-008.1. \nolinkurl{fr/sga7ii/diplomatic/expose_X_body.tex}; \nolinkurl{fr/sga7ii/current/expose_X_body.tex}\par

\par\medskip\noindent\hypertarget{app-F7II-X-008.2}{}\textbf{F7II-X-008.2}\quad p. 11; printed misprint.\par
\noindent Imprimé : {\ttfamily\footnotesize i\allowbreak{}}\par
\noindent Corrigé : {\ttfamily\footnotesize j\allowbreak{}}\par
\noindent Source : 7II:10. \nolinkurl{fixes/MISPRINTS_7II.json#10}\par

\par\medskip\noindent\hypertarget{app-F7II-X-009.1}{}\textbf{F7II-X-009.1}\quad p. 11; printed misprint.\par
\noindent Imprimé : {\ttfamily\footnotesize nombre\allowbreak{}}\par
\noindent Corrigé : {\ttfamily\footnotesize nombres\allowbreak{}}\par
\noindent Source : 7II:11. \nolinkurl{fixes/MISPRINTS_7II.json#11}\par

\par\medskip\noindent\hypertarget{app-F7II-X-010.1}{}\textbf{F7II-X-010.1}\quad p. 14; misprint\_pass.\par
\noindent Imprimé : {\ttfamily\footnotesize [\allowbreak{}}\par
\noindent Corrigé : {\ttfamily\footnotesize ]\allowbreak{}}\par
\noindent Source : 7II:12. \nolinkurl{fixes/FIXLOG_7II_CURRENT.jsonl#509}\par

\par\medskip\noindent\hypertarget{app-F7II-X-011.1}{}\textbf{F7II-X-011.1}\quad p. 15; designated-witness attribution.\par
\noindent Imprimé : {\ttfamily\footnotesize Soient\allowbreak{}}\par
\noindent Corrigé : {\ttfamily\footnotesize Soit\allowbreak{}}\par
\noindent Source : WITNESS-F7II-X-011.1. \nolinkurl{fr/sga7ii/diplomatic/expose_X_body.tex}; \nolinkurl{fr/sga7ii/current/expose_X_body.tex}\par

\par\medskip\noindent\hypertarget{app-F7II-X-011.2}{}\textbf{F7II-X-011.2}\quad p. 15; designated-witness attribution.\par
\noindent Imprimé : {\ttfamily\footnotesize  est\allowbreak{}}\par
\noindent Corrigé : \emph{(vide)}\par
\noindent Source : WITNESS-F7II-X-011.2. \nolinkurl{fr/sga7ii/diplomatic/expose_X_body.tex}; \nolinkurl{fr/sga7ii/current/expose_X_body.tex}\par

\par\medskip\noindent\hypertarget{app-F7II-X-012.1}{}\textbf{F7II-X-012.1}\quad p. 15; designated-witness attribution.\par
\noindent Imprimé : {\ttfamily\footnotesize o\allowbreak{}}\par
\noindent Corrigé : {\ttfamily\footnotesize 0\allowbreak{}}\par
\noindent Source : WITNESS-F7II-X-012.1. \nolinkurl{fr/sga7ii/diplomatic/expose_X_body.tex}; \nolinkurl{fr/sga7ii/current/expose_X_body.tex}\par

\par\medskip\noindent\hypertarget{app-F7II-X-012.2}{}\textbf{F7II-X-012.2}\quad p. 15; designated-witness attribution.\par
\noindent Imprimé : {\ttfamily\footnotesize o\allowbreak{}}\par
\noindent Corrigé : {\ttfamily\footnotesize 0\allowbreak{}}\par
\noindent Source : WITNESS-F7II-X-012.2. \nolinkurl{fr/sga7ii/diplomatic/expose_X_body.tex}; \nolinkurl{fr/sga7ii/current/expose_X_body.tex}\par

\par\medskip\noindent\hypertarget{app-F7II-X-012.3}{}\textbf{F7II-X-012.3}\quad p. 15; misprint\_pass.\par
\noindent Imprimé : {\ttfamily\footnotesize \textbackslash{}leq\allowbreak{}}\par
\noindent Corrigé : {\ttfamily\footnotesize <\allowbreak{}}\par
\noindent Source : 7II:14. \nolinkurl{fixes/FIXLOG_7II_CURRENT.jsonl#510}\par

\par\medskip\noindent\hypertarget{app-F7II-X-013.1}{}\textbf{F7II-X-013.1}\quad p. 15; designated-witness attribution.\par
\noindent Imprimé : {\ttfamily\footnotesize n\allowbreak{}}\par
\noindent Corrigé : {\ttfamily\footnotesize r\allowbreak{}}\par
\noindent Source : WITNESS-F7II-X-013.1. \nolinkurl{fr/sga7ii/diplomatic/expose_X_body.tex}; \nolinkurl{fr/sga7ii/current/expose_X_body.tex}\par

\par\medskip\noindent\hypertarget{app-F7II-X-013.2}{}\textbf{F7II-X-013.2}\quad p. 15; designated-witness attribution.\par
\noindent Imprimé : {\ttfamily\footnotesize n\allowbreak{}}\par
\noindent Corrigé : {\ttfamily\footnotesize r\allowbreak{}}\par
\noindent Source : WITNESS-F7II-X-013.2. \nolinkurl{fr/sga7ii/diplomatic/expose_X_body.tex}; \nolinkurl{fr/sga7ii/current/expose_X_body.tex}\par

\par\medskip\noindent\hypertarget{app-F7II-X-013.3}{}\textbf{F7II-X-013.3}\quad p. 15; designated-witness attribution.\par
\noindent Imprimé : {\ttfamily\footnotesize n\allowbreak{}}\par
\noindent Corrigé : {\ttfamily\footnotesize r\allowbreak{}}\par
\noindent Source : WITNESS-F7II-X-013.3. \nolinkurl{fr/sga7ii/diplomatic/expose_X_body.tex}; \nolinkurl{fr/sga7ii/current/expose_X_body.tex}\par

\par\medskip\noindent\hypertarget{app-F7II-X-013.4}{}\textbf{F7II-X-013.4}\quad p. 15; designated-witness attribution.\par
\noindent Imprimé : {\ttfamily\footnotesize n\allowbreak{}}\par
\noindent Corrigé : {\ttfamily\footnotesize r\allowbreak{}}\par
\noindent Source : WITNESS-F7II-X-013.4. \nolinkurl{fr/sga7ii/diplomatic/expose_X_body.tex}; \nolinkurl{fr/sga7ii/current/expose_X_body.tex}\par

\par\medskip\noindent\hypertarget{app-F7II-X-014.1}{}\textbf{F7II-X-014.1}\quad p. 15; designated-witness attribution.\par
\noindent Imprimé : {\ttfamily\footnotesize n\allowbreak{}}\par
\noindent Corrigé : {\ttfamily\footnotesize r\allowbreak{}}\par
\noindent Source : WITNESS-F7II-X-014.1. \nolinkurl{fr/sga7ii/diplomatic/expose_X_body.tex}; \nolinkurl{fr/sga7ii/current/expose_X_body.tex}\par

\par\medskip\noindent\hypertarget{app-F7II-X-014.2}{}\textbf{F7II-X-014.2}\quad p. 15; designated-witness attribution.\par
\noindent Imprimé : {\ttfamily\footnotesize n\allowbreak{}}\par
\noindent Corrigé : {\ttfamily\footnotesize r\allowbreak{}}\par
\noindent Source : WITNESS-F7II-X-014.2. \nolinkurl{fr/sga7ii/diplomatic/expose_X_body.tex}; \nolinkurl{fr/sga7ii/current/expose_X_body.tex}\par

\par\medskip\noindent\hypertarget{app-F7II-X-014.3}{}\textbf{F7II-X-014.3}\quad p. 15; designated-witness attribution.\par
\noindent Imprimé : {\ttfamily\footnotesize n\allowbreak{}}\par
\noindent Corrigé : {\ttfamily\footnotesize r\allowbreak{}}\par
\noindent Source : WITNESS-F7II-X-014.3. \nolinkurl{fr/sga7ii/diplomatic/expose_X_body.tex}; \nolinkurl{fr/sga7ii/current/expose_X_body.tex}\par

\par\medskip\noindent\hypertarget{app-F7II-X-015.1}{}\textbf{F7II-X-015.1}\quad p. 17; misprint\_in\_print.\par
\noindent Imprimé : {\ttfamily\footnotesize pur\allowbreak{}}\par
\noindent Corrigé : {\ttfamily\footnotesize purs\allowbreak{}}\par
\noindent Source : FIX-7II-D-0013. \nolinkurl{fixes/FIXLOG_7II.jsonl#13}\par

\par\medskip\noindent\hypertarget{app-F7II-X-016.1}{}\textbf{F7II-X-016.1}\quad p. 18; inherited corrected-layer difference.\par
\noindent Imprimé : {\ttfamily\footnotesize f\allowbreak{}}\par
\noindent Corrigé : {\ttfamily\footnotesize p\allowbreak{}}\par
\noindent Source : ADOPT-F7II-X-016.1. \nolinkurl{fr/sga7ii/diplomatic/expose_X_body.tex}; \nolinkurl{fr/sga7ii/current/expose_X_body.tex}\par

\par\medskip\noindent\hypertarget{app-F7II-X-016.2}{}\textbf{F7II-X-016.2}\quad p. 18; inherited corrected-layer difference.\par
\noindent Imprimé : {\ttfamily\footnotesize f\allowbreak{}}\par
\noindent Corrigé : {\ttfamily\footnotesize p\allowbreak{}}\par
\noindent Source : ADOPT-F7II-X-016.2. \nolinkurl{fr/sga7ii/diplomatic/expose_X_body.tex}; \nolinkurl{fr/sga7ii/current/expose_X_body.tex}\par

\par\medskip\noindent\hypertarget{app-F7II-X-017.1}{}\textbf{F7II-X-017.1}\quad p. 18; designated-witness attribution.\par
\noindent Imprimé : {\ttfamily\footnotesize f\allowbreak{}}\par
\noindent Corrigé : {\ttfamily\footnotesize p\allowbreak{}}\par
\noindent Source : WITNESS-F7II-X-017.1. \nolinkurl{fr/sga7ii/diplomatic/expose_X_body.tex}; \nolinkurl{fr/sga7ii/current/expose_X_body.tex}\par

\par\medskip\noindent\hypertarget{app-F7II-X-017.2}{}\textbf{F7II-X-017.2}\quad p. 18; printed misprint.\par
\noindent Imprimé : \emph{(vide)}\par
\noindent Corrigé : {\ttfamily\footnotesize \{\allowbreak{}}\par
\noindent Source : 7II:16. \nolinkurl{fixes/MISPRINTS_7II.json#16}\par

\par\medskip\noindent\hypertarget{app-F7II-X-017.3}{}\textbf{F7II-X-017.3}\quad p. 18; printed misprint.\par
\noindent Imprimé : \emph{(vide)}\par
\noindent Corrigé : {\ttfamily\footnotesize '\}\allowbreak{}}\par
\noindent Source : 7II:16. \nolinkurl{fixes/MISPRINTS_7II.json#16}\par

\par\medskip\noindent\hypertarget{app-F7II-X-017.4}{}\textbf{F7II-X-017.4}\quad p. 18; designated-witness attribution.\par
\noindent Imprimé : {\ttfamily\footnotesize f\allowbreak{}}\par
\noindent Corrigé : {\ttfamily\footnotesize p\allowbreak{}}\par
\noindent Source : WITNESS-F7II-X-017.4. \nolinkurl{fr/sga7ii/diplomatic/expose_X_body.tex}; \nolinkurl{fr/sga7ii/current/expose_X_body.tex}\par

\par\medskip\noindent\hypertarget{app-F7II-X-017.5}{}\textbf{F7II-X-017.5}\quad p. 18; inherited corrected-layer difference.\par
\noindent Imprimé : {\ttfamily\footnotesize f\allowbreak{}}\par
\noindent Corrigé : {\ttfamily\footnotesize p\allowbreak{}}\par
\noindent Source : ADOPT-F7II-X-017.5. \nolinkurl{fr/sga7ii/diplomatic/expose_X_body.tex}; \nolinkurl{fr/sga7ii/current/expose_X_body.tex}\par

\par\medskip\noindent\hypertarget{app-F7II-X-018.1}{}\textbf{F7II-X-018.1}\quad p. 18; printed misprint.\par
\noindent Imprimé : {\ttfamily\footnotesize S\allowbreak{}}\par
\noindent Corrigé : {\ttfamily\footnotesize D\allowbreak{}}\par
\noindent Source : 7II:17. \nolinkurl{fixes/MISPRINTS_7II.json#17}\par

\par\medskip\noindent\hypertarget{app-F7II-X-019.1}{}\textbf{F7II-X-019.1}\quad p. 19; printed misprint.\par
\noindent Imprimé : \emph{(vide)}\par
\noindent Corrigé : {\ttfamily\footnotesize '\allowbreak{}}\par
\noindent Source : 7II:18. \nolinkurl{fixes/MISPRINTS_7II.json#18}\par

\par\medskip\noindent\hypertarget{app-F7II-X-020.1}{}\textbf{F7II-X-020.1}\quad p. 20; inherited corrected-layer difference.\par
\noindent Imprimé : {\ttfamily\footnotesize Lf\allowbreak{}}\par
\noindent Corrigé : {\ttfamily\footnotesize Lp\allowbreak{}}\par
\noindent Source : ADOPT-F7II-X-020.1. \nolinkurl{fr/sga7ii/diplomatic/expose_X_body.tex}; \nolinkurl{fr/sga7ii/current/expose_X_body.tex}\par

\par\medskip\noindent\hypertarget{app-F7II-X-020.2}{}\textbf{F7II-X-020.2}\quad p. 20; designated-witness attribution.\par
\noindent Imprimé : {\ttfamily\footnotesize f\allowbreak{}}\par
\noindent Corrigé : {\ttfamily\footnotesize p\allowbreak{}}\par
\noindent Source : WITNESS-F7II-X-020.2. \nolinkurl{fr/sga7ii/diplomatic/expose_X_body.tex}; \nolinkurl{fr/sga7ii/current/expose_X_body.tex}\par

\par\medskip\noindent\hypertarget{app-F7II-X-021.1}{}\textbf{F7II-X-021.1}\quad p. 20; designated-witness attribution.\par
\noindent Imprimé : {\ttfamily\footnotesize f\allowbreak{}}\par
\noindent Corrigé : {\ttfamily\footnotesize p\allowbreak{}}\par
\noindent Source : WITNESS-F7II-X-021.1. \nolinkurl{fr/sga7ii/diplomatic/expose_X_body.tex}; \nolinkurl{fr/sga7ii/current/expose_X_body.tex}\par

\par\medskip\noindent\hypertarget{app-F7II-X-021.2}{}\textbf{F7II-X-021.2}\quad p. 20; designated-witness attribution.\par
\noindent Imprimé : {\ttfamily\footnotesize f\allowbreak{}}\par
\noindent Corrigé : {\ttfamily\footnotesize p\allowbreak{}}\par
\noindent Source : WITNESS-F7II-X-021.2. \nolinkurl{fr/sga7ii/diplomatic/expose_X_body.tex}; \nolinkurl{fr/sga7ii/current/expose_X_body.tex}\par

\par\medskip\noindent\hypertarget{app-F7II-X-021.3}{}\textbf{F7II-X-021.3}\quad p. 20; designated-witness attribution.\par
\noindent Imprimé : {\ttfamily\footnotesize f\allowbreak{}}\par
\noindent Corrigé : {\ttfamily\footnotesize p\allowbreak{}}\par
\noindent Source : WITNESS-F7II-X-021.3. \nolinkurl{fr/sga7ii/diplomatic/expose_X_body.tex}; \nolinkurl{fr/sga7ii/current/expose_X_body.tex}\par

\par\medskip\noindent\hypertarget{app-F7II-X-021.4}{}\textbf{F7II-X-021.4}\quad p. 20; inherited corrected-layer difference.\par
\noindent Imprimé : {\ttfamily\footnotesize f\allowbreak{}}\par
\noindent Corrigé : {\ttfamily\footnotesize p\_*\textbackslash{}underline\{O\}\_E\allowbreak{} \textbackslash{}longrightarrow\allowbreak{} p\_*\textbackslash{}\allowbreak{} \textbackslash{}underline\{O\}\_\{D'\textbackslash{}ca\allowbreak{}p\allowbreak{} E\}\allowbreak{} \textbackslash{}longrightarrow\allowbreak{} p\allowbreak{}}\par
\noindent Source : ADOPT-F7II-X-021.4. \nolinkurl{fr/sga7ii/diplomatic/expose_X_body.tex}; \nolinkurl{fr/sga7ii/current/expose_X_body.tex}\par

\par\medskip\noindent\hypertarget{app-F7II-X-021.5}{}\textbf{F7II-X-021.5}\quad p. 20; designated-witness attribution.\par
\noindent Imprimé : {\ttfamily\footnotesize f\_*\textbackslash{}underline\{O\}\_E\allowbreak{} \textbackslash{}longrightarrow\allowbreak{} f\_*\textbackslash{}\allowbreak{} \textbackslash{}underline\{O\}\_\{D'\textbackslash{}ca\allowbreak{}p\allowbreak{} E\}\allowbreak{} \textbackslash{}longrightarrow\allowbreak{} }\par
\noindent Corrigé : \emph{(vide)}\par
\noindent Source : WITNESS-F7II-X-021.5. \nolinkurl{fr/sga7ii/diplomatic/expose_X_body.tex}; \nolinkurl{fr/sga7ii/current/expose_X_body.tex}\par

\par\medskip\noindent\hypertarget{app-F7II-X-022.1}{}\textbf{F7II-X-022.1}\quad p. 21; inherited corrected-layer difference.\par
\noindent Imprimé : {\ttfamily\footnotesize f\allowbreak{}}\par
\noindent Corrigé : {\ttfamily\footnotesize p\allowbreak{}}\par
\noindent Source : ADOPT-F7II-X-022.1. \nolinkurl{fr/sga7ii/diplomatic/expose_X_body.tex}; \nolinkurl{fr/sga7ii/current/expose_X_body.tex}\par

\par\medskip\noindent\hypertarget{app-F7II-X-022.2}{}\textbf{F7II-X-022.2}\quad p. 21; printed misprint.\par
\noindent Imprimé : \emph{(vide)}\par
\noindent Corrigé : {\ttfamily\footnotesize  \textbackslash{}mu\_s(D)\textasciicircum{}2\textbackslash{},[k(s)\allowbreak{} :\allowbreak{} k]\allowbreak{} -\allowbreak{}}\par
\noindent Source : 7II:21. \nolinkurl{fixes/MISPRINTS_7II.json#21}\par

\par\medskip\noindent\hypertarget{app-F7II-X-022.3}{}\textbf{F7II-X-022.3}\quad p. 21; designated-witness attribution.\par
\noindent Imprimé : {\ttfamily\footnotesize :\allowbreak{} k]\allowbreak{} -\allowbreak{} \textbackslash{}mu\_s(D)\textasciicircum{}2\textbackslash{},[k(s)\allowbreak{} }\par
\noindent Corrigé : \emph{(vide)}\par
\noindent Source : WITNESS-F7II-X-022.3. \nolinkurl{fr/sga7ii/diplomatic/expose_X_body.tex}; \nolinkurl{fr/sga7ii/current/expose_X_body.tex}\par

\par\medskip\noindent\hypertarget{app-F7II-X-022.4}{}\textbf{F7II-X-022.4}\quad p. 21; designated-witness attribution.\par
\noindent Imprimé : {\ttfamily\footnotesize f\allowbreak{}}\par
\noindent Corrigé : {\ttfamily\footnotesize p\allowbreak{}}\par
\noindent Source : WITNESS-F7II-X-022.4. \nolinkurl{fr/sga7ii/diplomatic/expose_X_body.tex}; \nolinkurl{fr/sga7ii/current/expose_X_body.tex}\par

\par\medskip\noindent\hypertarget{app-F7II-X-022.5}{}\textbf{F7II-X-022.5}\quad p. 21; printed misprint.\par
\noindent Imprimé : {\ttfamily\footnotesize S\allowbreak{}}\par
\noindent Corrigé : {\ttfamily\footnotesize \textbackslash{}delta\allowbreak{}}\par
\noindent Source : 7II:20. \nolinkurl{fixes/MISPRINTS_7II.json#20}\par

\par\medskip\noindent\hypertarget{app-F7II-X-023.1}{}\textbf{F7II-X-023.1}\quad p. 22; inherited corrected-layer difference.\par
\noindent Imprimé : {\ttfamily\footnotesize ,\allowbreak{} \$D\$\allowbreak{} \textbackslash{}underline\{un\allowbreak{} diviseur\}\allowbreak{}}\par
\noindent Corrigé : \emph{(vide)}\par
\noindent Source : ADOPT-F7II-X-023.1. \nolinkurl{fr/sga7ii/diplomatic/expose_X_body.tex}; \nolinkurl{fr/sga7ii/current/expose_X_body.tex}\par

\par\medskip\noindent\hypertarget{app-F7II-X-024.1}{}\textbf{F7II-X-024.1}\quad p. 23; inherited corrected-layer difference.\par
\noindent Imprimé : {\ttfamily\footnotesize holomorphes\allowbreak{}}\par
\noindent Corrigé : {\ttfamily\footnotesize formelles\allowbreak{}}\par
\noindent Source : ADOPT-F7II-X-024.1. \nolinkurl{fr/sga7ii/diplomatic/expose_X_body.tex}; \nolinkurl{fr/sga7ii/current/expose_X_body.tex}\par

\par\medskip\noindent\hypertarget{app-F7II-X-025.1}{}\textbf{F7II-X-025.1}\quad p. 23; inherited corrected-layer difference.\par
\noindent Imprimé : {\ttfamily\footnotesize ,\allowbreak{}}\par
\noindent Corrigé : \emph{(vide)}\par
\noindent Source : ADOPT-F7II-X-025.1. \nolinkurl{fr/sga7ii/diplomatic/expose_X_body.tex}; \nolinkurl{fr/sga7ii/current/expose_X_body.tex}\par

\par\medskip\noindent\hypertarget{app-F7II-X-026.1}{}\textbf{F7II-X-026.1}\quad p. 24; printed misprint.\par
\noindent Imprimé : \emph{(vide)}\par
\noindent Corrigé : {\ttfamily\footnotesize )\allowbreak{}}\par
\noindent Source : 7II:22. \nolinkurl{fixes/MISPRINTS_7II.json#22}\par

\par\medskip\noindent\hypertarget{app-F7II-X-026.2}{}\textbf{F7II-X-026.2}\quad p. 24; printed misprint.\par
\noindent Imprimé : {\ttfamily\footnotesize \textbackslash{}cap\allowbreak{}}\par
\noindent Corrigé : {\ttfamily\footnotesize \textbackslash{}cup\allowbreak{}}\par
\noindent Source : 7II:23. \nolinkurl{fixes/MISPRINTS_7II.json#23}\par

\par\medskip\noindent\hypertarget{app-F7II-X-026.3}{}\textbf{F7II-X-026.3}\quad p. 24; printed misprint.\par
\noindent Imprimé : {\ttfamily\footnotesize \textbackslash{}cap\allowbreak{}}\par
\noindent Corrigé : {\ttfamily\footnotesize \textbackslash{}cup\allowbreak{}}\par
\noindent Source : 7II:24. \nolinkurl{fixes/MISPRINTS_7II.json#24}\par

\par\medskip\noindent\hypertarget{app-F7II-X-026.4}{}\textbf{F7II-X-026.4}\quad p. 24; printed misprint.\par
\noindent Imprimé : {\ttfamily\footnotesize ou\allowbreak{}}\par
\noindent Corrigé : {\ttfamily\footnotesize au\allowbreak{}}\par
\noindent Source : 7II:25. \nolinkurl{fixes/MISPRINTS_7II.json#25}\par

\par\medskip\noindent\hypertarget{app-F7II-X-027.1}{}\textbf{F7II-X-027.1}\quad p. 24; printed misprint.\par
\noindent Imprimé : \emph{(vide)}\par
\noindent Corrigé : {\ttfamily\footnotesize  \textbackslash{}underline\{O\}\allowbreak{}}\par
\noindent Source : 7II:26. \nolinkurl{fixes/MISPRINTS_7II.json#26}\par

\par\medskip\noindent\hypertarget{app-F7II-X-028.1}{}\textbf{F7II-X-028.1}\quad p. 25; printed misprint.\par
\noindent Imprimé : \emph{(vide)}\par
\noindent Corrigé : {\ttfamily\footnotesize )\allowbreak{}}\par
\noindent Source : 7II:27. \nolinkurl{fixes/MISPRINTS_7II.json#27}\par

\par\medskip\noindent\hypertarget{app-F7II-X-029.1}{}\textbf{F7II-X-029.1}\quad p. 26; printed misprint.\par
\noindent Imprimé : \emph{(vide)}\par
\noindent Corrigé : {\ttfamily\footnotesize )\allowbreak{}}\par
\noindent Source : 7II:28. \nolinkurl{fixes/MISPRINTS_7II.json#28}\par

\par\medskip\noindent\hypertarget{app-F7II-X-029.2}{}\textbf{F7II-X-029.2}\quad p. 26; printed misprint.\par
\noindent Imprimé : \emph{(vide)}\par
\noindent Corrigé : {\ttfamily\footnotesize )\allowbreak{}}\par
\noindent Source : 7II:29. \nolinkurl{fixes/MISPRINTS_7II.json#29}\par

\par\medskip\noindent\hypertarget{app-F7II-X-029.3}{}\textbf{F7II-X-029.3}\quad p. 26; designated-witness attribution.\par
\noindent Imprimé : \emph{(vide)}\par
\noindent Corrigé : {\ttfamily\footnotesize \textbackslash{}substack\{\allowbreak{}}\par
\noindent Source : WITNESS-F7II-X-029.3. \nolinkurl{fr/sga7ii/diplomatic/expose_X_body.tex}; \nolinkurl{fr/sga7ii/current/expose_X_body.tex}\par

\par\medskip\noindent\hypertarget{app-F7II-X-029.4}{}\textbf{F7II-X-029.4}\quad p. 26; designated-witness attribution.\par
\noindent Imprimé : \emph{(vide)}\par
\noindent Corrigé : {\ttfamily\footnotesize \textbackslash{}\textbackslash{}\allowbreak{} t\textbackslash{}ne\allowbreak{} s\}\allowbreak{}}\par
\noindent Source : WITNESS-F7II-X-029.4. \nolinkurl{fr/sga7ii/diplomatic/expose_X_body.tex}; \nolinkurl{fr/sga7ii/current/expose_X_body.tex}\par

\par\medskip\noindent\hypertarget{app-F7II-X-029.5}{}\textbf{F7II-X-029.5}\quad p. 26; printed misprint.\par
\noindent Imprimé : {\ttfamily\footnotesize précèdent\allowbreak{}}\par
\noindent Corrigé : {\ttfamily\footnotesize précède\allowbreak{}}\par
\noindent Source : 7II:31. \nolinkurl{fixes/MISPRINTS_7II.json#31}\par

\par\medskip\noindent\hypertarget{app-F7II-X-029.6}{}\textbf{F7II-X-029.6}\quad p. 26; misprint\_in\_print.\par
\noindent Imprimé : {\ttfamily\footnotesize s\allowbreak{}}\par
\noindent Corrigé : {\ttfamily\footnotesize S\allowbreak{}}\par
\noindent Source : FIX-7II-D-0022. \nolinkurl{fixes/FIXLOG_7II.jsonl#22}\par

\par\medskip\noindent\hypertarget{app-F7II-X-029.7}{}\textbf{F7II-X-029.7}\quad p. 26; printed misprint.\par
\noindent Imprimé : \emph{(vide)}\par
\noindent Corrigé : {\ttfamily\footnotesize )\allowbreak{}}\par
\noindent Source : 7II:33. \nolinkurl{fixes/MISPRINTS_7II.json#33}\par

\par\medskip\noindent\hypertarget{app-F7II-X-030.1}{}\textbf{F7II-X-030.1}\quad p. 27; printed misprint.\par
\noindent Imprimé : \emph{(vide)}\par
\noindent Corrigé : {\ttfamily\footnotesize (t)\allowbreak{}}\par
\noindent Source : 7II:34. \nolinkurl{fixes/MISPRINTS_7II.json#34}\par

\par\medskip\noindent\hypertarget{app-F7II-X-031.1}{}\textbf{F7II-X-031.1}\quad p. 27; printed misprint.\par
\noindent Imprimé : {\ttfamily\footnotesize correspondantes\allowbreak{}}\par
\noindent Corrigé : {\ttfamily\footnotesize correspondante\allowbreak{}}\par
\noindent Source : 7II:35. \nolinkurl{fixes/MISPRINTS_7II.json#35}\par

\par\medskip\noindent\hypertarget{app-F7II-X-032.1}{}\textbf{F7II-X-032.1}\quad p. 29; printed misprint.\par
\noindent Imprimé : {\ttfamily\footnotesize et\allowbreak{}}\par
\noindent Corrigé : {\ttfamily\footnotesize est\allowbreak{}}\par
\noindent Source : 7II:36. \nolinkurl{fixes/MISPRINTS_7II.json#36}\par

\par\medskip\noindent\hypertarget{app-F7II-X-032.2}{}\textbf{F7II-X-032.2}\quad p. 29; inherited corrected-layer difference.\par
\noindent Imprimé : \emph{(vide)}\par
\noindent Corrigé : \emph{(1 espace(s), 0 retour(s) de ligne dans le témoin TeX)}\par
\noindent Source : ADOPT-F7II-X-032.2. \nolinkurl{fr/sga7ii/diplomatic/expose_X_body.tex}; \nolinkurl{fr/sga7ii/current/expose_X_body.tex}\par

\par\medskip\noindent\hypertarget{app-F7II-X-032.3}{}\textbf{F7II-X-032.3}\quad p. 29; misprint\_in\_print.\par
\noindent Imprimé : {\ttfamily\footnotesize coîncidences\allowbreak{}}\par
\noindent Corrigé : {\ttfamily\footnotesize coïncidences\allowbreak{}}\par
\noindent Source : FIX-7II-D-0023. \nolinkurl{fixes/FIXLOG_7II.jsonl#23}\par

\par\medskip\noindent\hypertarget{app-F7II-X-033.1}{}\textbf{F7II-X-033.1}\quad p. 29; printed misprint.\par
\noindent Imprimé : {\ttfamily\footnotesize à\allowbreak{}}\par
\noindent Corrigé : {\ttfamily\footnotesize a\allowbreak{}}\par
\noindent Source : 7II:38. \nolinkurl{fixes/MISPRINTS_7II.json#38}\par

\par\medskip\noindent\hypertarget{app-F7II-X-033.2}{}\textbf{F7II-X-033.2}\quad p. 29; printed misprint.\par
\noindent Imprimé : {\ttfamily\footnotesize 2\allowbreak{}}\par
\noindent Corrigé : {\ttfamily\footnotesize 1\allowbreak{}}\par
\noindent Source : 7II:39. \nolinkurl{fixes/MISPRINTS_7II.json#39}\par

\par\medskip\noindent\hypertarget{app-F7II-X-034.1}{}\textbf{F7II-X-034.1}\quad p. 30; printed misprint.\par
\noindent Imprimé : {\ttfamily\footnotesize dimension\allowbreak{}}\par
\noindent Corrigé : {\ttfamily\footnotesize dimensions\allowbreak{}}\par
\noindent Source : 7II:40. \nolinkurl{fixes/MISPRINTS_7II.json#40}\par

\par\medskip\noindent\hypertarget{app-F7II-X-034.2}{}\textbf{F7II-X-034.2}\quad p. 30; inherited corrected-layer difference.\par
\noindent Imprimé : {\ttfamily\footnotesize A\allowbreak{}}\par
\noindent Corrigé : {\ttfamily\footnotesize À\allowbreak{}}\par
\noindent Source : ADOPT-F7II-X-034.2. \nolinkurl{fr/sga7ii/diplomatic/expose_X_body.tex}; \nolinkurl{fr/sga7ii/current/expose_X_body.tex}\par

\par\medskip\noindent\hypertarget{app-F7II-X-034.3}{}\textbf{F7II-X-034.3}\quad p. 30; inherited corrected-layer difference.\par
\noindent Imprimé : \emph{(vide)}\par
\noindent Corrigé : \emph{(1 espace(s), 0 retour(s) de ligne dans le témoin TeX)}\par
\noindent Source : ADOPT-F7II-X-034.3. \nolinkurl{fr/sga7ii/diplomatic/expose_X_body.tex}; \nolinkurl{fr/sga7ii/current/expose_X_body.tex}\par

\par\medskip\noindent\hypertarget{app-F7II-X-035.1}{}\textbf{F7II-X-035.1}\quad p. 31; misprint\_pass.\par
\noindent Imprimé : {\ttfamily\footnotesize 4\allowbreak{}}\par
\noindent Corrigé : {\ttfamily\footnotesize 2\allowbreak{}}\par
\noindent Source : 7II:41. \nolinkurl{fixes/FIXLOG_7II_CURRENT.jsonl#511}; \nolinkurl{fixes/MISPRINTS_7II.json#41}\par

\par\medskip\noindent\hypertarget{app-F7II-X-036.1}{}\textbf{F7II-X-036.1}\quad p. 32; designated-witness attribution.\par
\noindent Imprimé : \emph{(vide)}\par
\noindent Corrigé : {\ttfamily\footnotesize \_s\allowbreak{}}\par
\noindent Source : WITNESS-F7II-X-036.1. \nolinkurl{fr/sga7ii/diplomatic/expose_X_body.tex}; \nolinkurl{fr/sga7ii/current/expose_X_body.tex}\par

\par\medskip\noindent\hypertarget{app-F7II-X-037.1}{}\textbf{F7II-X-037.1}\quad p. 33; printed misprint.\par
\noindent Imprimé : {\ttfamily\footnotesize la\allowbreak{}}\par
\noindent Corrigé : {\ttfamily\footnotesize le\allowbreak{}}\par
\noindent Source : 7II:43. \nolinkurl{fixes/MISPRINTS_7II.json#43}\par

\par\medskip\noindent\hypertarget{app-F7II-X-037.2}{}\textbf{F7II-X-037.2}\quad p. 33; inherited corrected-layer difference.\par
\noindent Imprimé : \emph{(vide)}\par
\noindent Corrigé : {\ttfamily\footnotesize \}\allowbreak{}}\par
\noindent Source : ADOPT-F7II-X-037.2. \nolinkurl{fr/sga7ii/diplomatic/expose_X_body.tex}; \nolinkurl{fr/sga7ii/current/expose_X_body.tex}\par

\par\medskip\noindent\hypertarget{app-F7II-X-037.3}{}\textbf{F7II-X-037.3}\quad p. 33; inherited corrected-layer difference.\par
\noindent Imprimé : \emph{(vide)}\par
\noindent Corrigé : {\ttfamily\footnotesize \$T\_1\$\allowbreak{} \textbackslash{}underline\{\allowbreak{}}\par
\noindent Source : ADOPT-F7II-X-037.3. \nolinkurl{fr/sga7ii/diplomatic/expose_X_body.tex}; \nolinkurl{fr/sga7ii/current/expose_X_body.tex}\par

\par\medskip\noindent\hypertarget{app-F7II-X-037.4}{}\textbf{F7II-X-037.4}\quad p. 33; inherited corrected-layer difference.\par
\noindent Imprimé : \emph{(vide)}\par
\noindent Corrigé : {\ttfamily\footnotesize ,\allowbreak{}}\par
\noindent Source : ADOPT-F7II-X-037.4. \nolinkurl{fr/sga7ii/diplomatic/expose_X_body.tex}; \nolinkurl{fr/sga7ii/current/expose_X_body.tex}\par

\par\medskip\noindent\hypertarget{app-F7II-X-037.5}{}\textbf{F7II-X-037.5}\quad p. 33; inherited corrected-layer difference.\par
\noindent Imprimé : {\ttfamily\footnotesize \$T\_1\$,\allowbreak{}}\par
\noindent Corrigé : {\ttfamily\footnotesize \textbackslash{}underline\{un\allowbreak{} point\}\allowbreak{}}\par
\noindent Source : ADOPT-F7II-X-037.5. \nolinkurl{fr/sga7ii/diplomatic/expose_X_body.tex}; \nolinkurl{fr/sga7ii/current/expose_X_body.tex}\par

\par\medskip\noindent\hypertarget{app-F7II-X-037.6}{}\textbf{F7II-X-037.6}\quad p. 33; inherited corrected-layer difference.\par
\noindent Imprimé : \emph{(vide)}\par
\noindent Corrigé : {\ttfamily\footnotesize  un\allowbreak{} morphisme\allowbreak{}}\par
\noindent Source : ADOPT-F7II-X-037.6. \nolinkurl{fr/sga7ii/diplomatic/expose_X_body.tex}; \nolinkurl{fr/sga7ii/current/expose_X_body.tex}\par

\par\medskip\noindent\hypertarget{app-F7II-X-037.7}{}\textbf{F7II-X-037.7}\quad p. 33; inherited corrected-layer difference.\par
\noindent Imprimé : {\ttfamily\footnotesize o\allowbreak{}}\par
\noindent Corrigé : {\ttfamily\footnotesize 0\allowbreak{}}\par
\noindent Source : ADOPT-F7II-X-037.7. \nolinkurl{fr/sga7ii/diplomatic/expose_X_body.tex}; \nolinkurl{fr/sga7ii/current/expose_X_body.tex}\par

\par\medskip\noindent\hypertarget{app-F7II-X-037.8}{}\textbf{F7II-X-037.8}\quad p. 33; printed misprint.\par
\noindent Imprimé : {\ttfamily\footnotesize S\allowbreak{}}\par
\noindent Corrigé : {\ttfamily\footnotesize T\allowbreak{}}\par
\noindent Source : 7II:44. \nolinkurl{fixes/MISPRINTS_7II.json#44}\par

\par\medskip\noindent\hypertarget{app-F7II-X-037.9}{}\textbf{F7II-X-037.9}\quad p. 33; inherited corrected-layer difference.\par
\noindent Imprimé : {\ttfamily\footnotesize soit\allowbreak{}}\par
\noindent Corrigé : {\ttfamily\footnotesize soient\allowbreak{}}\par
\noindent Source : ADOPT-F7II-X-037.9. \nolinkurl{fr/sga7ii/diplomatic/expose_X_body.tex}; \nolinkurl{fr/sga7ii/current/expose_X_body.tex}\par

\par\medskip\noindent\hypertarget{app-F7II-X-037.10}{}\textbf{F7II-X-037.10}\quad p. 33; inherited corrected-layer difference.\par
\noindent Imprimé : \emph{(vide)}\par
\noindent Corrigé : {\ttfamily\footnotesize ,\allowbreak{}}\par
\noindent Source : ADOPT-F7II-X-037.10. \nolinkurl{fr/sga7ii/diplomatic/expose_X_body.tex}; \nolinkurl{fr/sga7ii/current/expose_X_body.tex}\par

\par\medskip\noindent\hypertarget{app-F7II-X-037.11}{}\textbf{F7II-X-037.11}\quad p. 33; inherited corrected-layer difference.\par
\noindent Imprimé : \emph{(vide)}\par
\noindent Corrigé : {\ttfamily\footnotesize ,\allowbreak{}}\par
\noindent Source : ADOPT-F7II-X-037.11. \nolinkurl{fr/sga7ii/diplomatic/expose_X_body.tex}; \nolinkurl{fr/sga7ii/current/expose_X_body.tex}\par

\par\medskip\noindent\hypertarget{app-F7II-X-037.12}{}\textbf{F7II-X-037.12}\quad p. 33; inherited corrected-layer difference.\par
\noindent Imprimé : \emph{(vide)}\par
\noindent Corrigé : {\ttfamily\footnotesize  des\allowbreak{} équations,\allowbreak{} avec\allowbreak{}}\par
\noindent Source : ADOPT-F7II-X-037.12. \nolinkurl{fr/sga7ii/diplomatic/expose_X_body.tex}; \nolinkurl{fr/sga7ii/current/expose_X_body.tex}\par

\par\medskip\noindent\hypertarget{app-F7II-X-037.13}{}\textbf{F7II-X-037.13}\quad p. 33; inherited corrected-layer difference.\par
\noindent Imprimé : {\ttfamily\footnotesize  équations\allowbreak{}}\par
\noindent Corrigé : {\ttfamily\footnotesize ,\allowbreak{}}\par
\noindent Source : ADOPT-F7II-X-037.13. \nolinkurl{fr/sga7ii/diplomatic/expose_X_body.tex}; \nolinkurl{fr/sga7ii/current/expose_X_body.tex}\par

\par\medskip\noindent\hypertarget{app-F7II-X-037.14}{}\textbf{F7II-X-037.14}\quad p. 33; printed misprint.\par
\noindent Imprimé : \emph{(vide)}\par
\noindent Corrigé : {\ttfamily\footnotesize \_2\allowbreak{}}\par
\noindent Source : 7II:45. \nolinkurl{fixes/MISPRINTS_7II.json#45}\par

\par\medskip\noindent\hypertarget{app-F7II-X-038.1}{}\textbf{F7II-X-038.1}\quad p. 34; printed misprint.\par
\noindent Imprimé : {\ttfamily\footnotesize foncteur\allowbreak{}}\par
\noindent Corrigé : {\ttfamily\footnotesize foncteurs\allowbreak{}}\par
\noindent Source : 7II:46. \nolinkurl{fixes/MISPRINTS_7II.json#46}\par

\par\medskip\noindent\hypertarget{app-F7II-X-039.1}{}\textbf{F7II-X-039.1}\quad p. 34; printed misprint.\par
\noindent Imprimé : \emph{(vide)}\par
\noindent Corrigé : {\ttfamily\footnotesize \_o\allowbreak{}}\par
\noindent Source : 7II:47. \nolinkurl{fixes/MISPRINTS_7II.json#47}\par

\par\medskip\noindent\hypertarget{app-F7II-X-039.2}{}\textbf{F7II-X-039.2}\quad p. 34; printed misprint.\par
\noindent Imprimé : {\ttfamily\footnotesize -\allowbreak{}}\par
\noindent Corrigé : \emph{(vide)}\par
\noindent Source : 7II:48. \nolinkurl{fixes/MISPRINTS_7II.json#48}\par

\par\medskip\noindent\hypertarget{app-F7II-X-040.1}{}\textbf{F7II-X-040.1}\quad p. 35; printed misprint.\par
\noindent Imprimé : \emph{(vide)}\par
\noindent Corrigé : {\ttfamily\footnotesize \textasciicircum{}E\allowbreak{}}\par
\noindent Source : 7II:49. \nolinkurl{fixes/MISPRINTS_7II.json#49}\par

\par\medskip\noindent\hypertarget{app-F7II-X-041.1}{}\textbf{F7II-X-041.1}\quad p. 36; inherited corrected-layer difference.\par
\noindent Imprimé : \emph{(vide)}\par
\noindent Corrigé : {\ttfamily\footnotesize \_\{n\textbackslash{}geq0\}\allowbreak{}}\par
\noindent Source : ADOPT-F7II-X-041.1. \nolinkurl{fr/sga7ii/diplomatic/expose_X_body.tex}; \nolinkurl{fr/sga7ii/current/expose_X_body.tex}\par

\par\medskip\noindent\hypertarget{app-F7II-X-041.2}{}\textbf{F7II-X-041.2}\quad p. 36; misprint\_in\_print.\par
\noindent Imprimé : {\ttfamily\footnotesize annulle\allowbreak{}}\par
\noindent Corrigé : {\ttfamily\footnotesize annule\allowbreak{}}\par
\noindent Source : FIX-7II-D-0026. \nolinkurl{fixes/FIXLOG_7II.jsonl#26}\par

\par\medskip\noindent\hypertarget{app-F7II-X-042.1}{}\textbf{F7II-X-042.1}\quad p. 36; misprint\_in\_print.\par
\noindent Imprimé : {\ttfamily\footnotesize annullent\allowbreak{}}\par
\noindent Corrigé : {\ttfamily\footnotesize annulent\allowbreak{}}\par
\noindent Source : FIX-7II-D-0027. \nolinkurl{fixes/FIXLOG_7II.jsonl#27}\par

\par\medskip\noindent\hypertarget{app-F7II-X-042.2}{}\textbf{F7II-X-042.2}\quad p. 36; printed misprint.\par
\noindent Imprimé : \emph{(vide)}\par
\noindent Corrigé : {\ttfamily\footnotesize  a\allowbreak{}}\par
\noindent Source : 7II:52. \nolinkurl{fixes/MISPRINTS_7II.json#52}\par

\par\medskip\noindent\hypertarget{app-F7II-X-043.1}{}\textbf{F7II-X-043.1}\quad p. 37; printed misprint.\par
\noindent Imprimé : {\ttfamily\footnotesize et\allowbreak{}}\par
\noindent Corrigé : {\ttfamily\footnotesize e\allowbreak{}}\par
\noindent Source : 7II:56. \nolinkurl{fixes/MISPRINTS_7II.json#56}\par

\par\medskip\noindent\hypertarget{app-F7II-X-043.2}{}\textbf{F7II-X-043.2}\quad p. 37; printed misprint.\par
\noindent Imprimé : {\ttfamily\footnotesize Lectures\allowbreak{}}\par
\noindent Corrigé : {\ttfamily\footnotesize Lecture\allowbreak{}}\par
\noindent Source : 7II:53. \nolinkurl{fixes/MISPRINTS_7II.json#53}\par

\par\medskip\noindent\hypertarget{app-F7II-X-043.3}{}\textbf{F7II-X-043.3}\quad p. 37; printed misprint.\par
\noindent Imprimé : {\ttfamily\footnotesize 15\allowbreak{}}\par
\noindent Corrigé : {\ttfamily\footnotesize 90\allowbreak{}}\par
\noindent Source : 7II:54. \nolinkurl{fixes/MISPRINTS_7II.json#54}\par

\par\medskip\noindent\hypertarget{app-F7II-X-044.1}{}\textbf{F7II-X-044.1}\quad p. 37; printed misprint.\par
\noindent Imprimé : {\ttfamily\footnotesize surfaces\allowbreak{}}\par
\noindent Corrigé : {\ttfamily\footnotesize surface\allowbreak{}}\par
\noindent Source : 7II:55. \nolinkurl{fixes/MISPRINTS_7II.json#55}\par

\par\medskip\noindent\hypertarget{app-F7II-X-045.1}{}\textbf{F7II-X-045.1}\quad p. 38; inherited corrected-layer difference.\par
\noindent Imprimé : \emph{(vide)}\par
\noindent Corrigé : {\ttfamily\footnotesize ,\allowbreak{}}\par
\noindent Source : ADOPT-F7II-X-045.1. \nolinkurl{fr/sga7ii/diplomatic/expose_X_body.tex}; \nolinkurl{fr/sga7ii/current/expose_X_body.tex}\par

\par\medskip\noindent\hypertarget{app-F7II-X-046.1}{}\textbf{F7II-X-046.1}\quad p. 38; inherited corrected-layer difference.\par
\noindent Imprimé : \emph{(1 espace(s), 0 retour(s) de ligne dans le témoin TeX)}\par
\noindent Corrigé : {\ttfamily\footnotesize -\allowbreak{}}\par
\noindent Source : ADOPT-F7II-X-046.1. \nolinkurl{fr/sga7ii/diplomatic/expose_X_body.tex}; \nolinkurl{fr/sga7ii/current/expose_X_body.tex}\par

\normalsize\fi
\clearpage\SGAtoc{Exposé XI — COHOMOLOGIE DES INTERSECTIONS COMPLETES}
\selectlanguage{french}

\SGApage{II}{39}


S G A 7


\underline{E X P O S E  XI}


\underline{\underline{COHOMOLOGIE DES INTERSECTIONS COMPLETES}}


\underline{par P. DELIGNE}


\SGAsectiontoc{1. Cohomologie de l'espace projectif et des intersections compl\`etes}{1}
1. \underline{Cohomologie de l'espace projectif et des intersections compl\`etes}

Soit $X$ une vari\'et\'e projective non singuli\`ere complexe, purement de dimension $n$, et soit $\underline{O}(1)$ un faisceau inversible ample sur $X$. D'apr\`es Akizuki et Nakano [1], il r\'esulte du vanishing theorem de Kodaira [4] que

\begin{equation}\tag{*}
H^i(X,\Omega^j_X (-1)) = 0 \quad \mathrm{pour}\ i+j < n \ .
\end{equation}
D'apr\`es [1] toujours, ce r\'esultat implique le ``th\'eor\`eme de Lefschetz faible'', en cohomologie complexe. On se propose de montrer, par force brutale, que si $X$ est un espace projectif sur un corps de caract\'eristique quelconque, alors (*) reste vrai. Ceci permettra, \`a l'imitation de [1], de comparer les cohomologies de Hodge et de De Rham des intersections compl\`etes, sur un corps de caract\'eristique quelconque, \`a celles de l'espace projectif.

On explicite ensuite quelques cons\'equences des r\'esultats obtenus. Au $n^{\circ}$ 2, on expose des r\'esultats qualitatifs d\'eduits de la formule num\'erique de Hirzebruch (2.2). Je tiens \`a remercier N. Katz de l'aide qu'il m'a apport\'ee dans la conception de ce num\'ero.

1.0. Soit $\underline{E}$ un module localement libre de rang $r+1$ sur un sch\'ema $S$. \underline{On suppose que} $r > 0$. Rappelons que $\mathbb{P}(\underline{E})$ d\'esigne l'espace projectif relatif sur $S$ des hyperplans de $\underline{E}$ . On d\'esignera par $p$ la projection de $\mathbb{P}(\underline{E})$ sur $S$ et par $\eta \in H^0(S,R^1 p_*(\Omega^1_{\mathbb{P}(\underline{E})/S}))$ la premi\`ere classe de Chern, style Hodge, du faisceau
\SGApage{II}{40}
inversible $\underline{O}(1)$. Rappelons que $\eta$ est l'image par la diff\'erentielle logarithmique $df/f : \underline{O}^* \longrightarrow \Omega^1_{\mathbb{P}(\underline{E})/S}$ de la classe de $\underline{O}(1)$ dans $H^0(S,R^1 p_* \underline{O}^*)$. Pour tout faisceau $\underline{F}$ de $\underline{O}$-modules sur $\mathbb{P}(\underline{E})$, on pose

\begin{equation*}
\underline{F}(n) = \underline{F} \otimes \underline{O}(1)^{\otimes n} \qquad (n \in \mathbb{Z}) \qquad .
\end{equation*}
\underline{Th\'eor\`eme} 1.1. \quad (i) \underline{Les faisceaux} $R^i p_*(\Omega^j_{\mathbb{P}(\underline{E})/S}(n))$ \underline{sont localement libres, et leur formation est compatible \`a tout changement de base.}

(ii) \underline{Pour} $0 \leq i \leq r$, \underline{la section} $\eta^i$ \underline{de} $R^i p_*(\Omega^i_{\mathbb{P}(\underline{E})/S})$ \underline{d\'efinit un isomorphisme entre ce faisceau et} $\underline{O}_S$ . \underline{De plus,}

\begin{equation*}
R^i p_* \ \Omega^j_{\mathbb{P}(\underline{E})/S} \ = \ 0 \qquad \underline{\mathrm{pour}}\ i \neq j \ \underline{\mathrm{ou}}\ i > r \qquad .
\end{equation*}
(iii) \underline{Si} $n \neq 0$, \underline{alors,} \underline{en dehors des deux cas suivants,} $R^i p_*(\Omega^j_{\mathbb{P}(\underline{E})/S}(n))$ \underline{est nul} :

\begin{equation}\tag{iii$_{a}$}
i = 0 \quad \underline{\mathrm{et}}\ n > j
\end{equation}
\begin{equation}\tag{iii$_{b}$}
i = r \quad \underline{\mathrm{et}}\ n < j-r \ .
\end{equation}
Soient $\mathbb{V}(\underline{E}) = \mathrm{Spec}\ (\mathrm{Sym}^*\ \underline{E})$ l'espace affine d\'efini par $\underline{E}$, $e$ sa section nulle, $\mathbb{V}(\underline{E})^*$ le compl\'ementaire $\mathbb{V}(\underline{E})^* = \mathbb{V}(\underline{E}) \setminus e(S)$, $a$ et $a'$ les projections de $\mathbb{V}(\underline{E})$ et $\mathbb{V}(\underline{E})^*$ sur $S$ et $q$ la projection de $\mathbb{V}(\underline{E})^*$ sur $\mathbb{P}(\underline{E})$ :

\begin{center}
\begin{tikzcd}[row sep=large, column sep=large]
\mathbb{V}(\underline{E}) \arrow[dr, "a"] \& \mathbb{V}(\underline{E})^* \arrow[l, hook'] \arrow[r, "q"] \arrow[d, "a'"] \& \mathbb{P}(\underline{E}) \arrow[dl, "p"] \\
{} \& S \arrow[ul, bend left, "e"] \&
{} \end{tikzcd}\quad .
\end{center}













L'espace affine \'epoint\'e $\mathbb{V}(\underline{O}(1))^*$ sur $\mathbb{P}(\underline{E})$ est canoniquement isomorphe \`a $\mathbb{V}(\underline{E})^*$, de sorte que, pour tout faisceau quasi-coh\'erent de $\underline{O}$-modules $\underline{F}$ sur $\mathbb{P}(\underline{E})$, on a canoniquement 


\SGApage{II}{41}

\begin{equation}\tag{1.1.1}
q_* \, q^* \, \underline{F} \;\simeq\; \sum_{n \in \mathbb{Z}} \underline{F}(n) \;\; .
\end{equation}
Cette graduation de $q_* \, q^* \, \underline{F}$ peut se déduire de l'action, sur ce faisceau, du groupe $\mathbb{G}_m$ des homothéties de $\mathbb{V}(\underline{E})$.

\SGAvar{Les champs de vecteurs relatifs sur $\mathbb{V}(\underline{E})$ s'identifient aux $S$-morphismes de $\mathbb{V}(\underline{E})$ dans $\mathbb{V}(\underline{E})$ (identifié à l'espace tangent à l'origine). Le champ de vecteur $X$ défini par l'application identique de $\mathbb{V}(\underline{E})$ est invariant par homothéties, Ce champ de vecteur est tangent aux fibres de $q$ et définit une trivialisation

}{Les champs de vecteurs relatifs sur $\mathbb{V}(\underline{E})$ s'identifient aux $S$-morphismes de $\mathbb{V}(\underline{E})$ dans $\mathbb{V}(\underline{E})$ (identifié à l'espace tangent à l'origine). Le champ de vecteur $X$ défini par l'application identique de $\mathbb{V}(\underline{E})$ est invariant par homothéties. Ce champ de vecteur est tangent aux fibres de $q$ et définit une trivialisation

}{F7II-XI-001}\begin{equation}\tag{1.1.2}
X \;:\; \Omega_{\mathbb{V}(\underline{E})^*/\,\mathbb{P}(\underline{E})} \longrightarrow \underline{O}_{\mathbb{V}(\underline{E})^*} \;\; .
\end{equation}
Si les $x_i$ forment une base de $\underline{E}$, on a

\begin{equation*}
X \;=\; \sum x_i \, dx^i \;\; .
\end{equation*}
Puisque $\Omega^1_{\mathbb{V}(\underline{E})^*/\mathbb{P}(\underline{E})}$ est de rang un, la suite exacte

\begin{equation*}
0 \longrightarrow q^*\Omega^1_{\mathbb{P}(\underline{E})/S} \longrightarrow \Omega^1_{\mathbb{V}(\underline{E})^*/S} \longrightarrow \Omega^1_{\mathbb{V}(\underline{E})^*/\mathbb{P}(\underline{E})} \longrightarrow 0
\end{equation*}
induit des suites exactes

\begin{equation*}
0 \longrightarrow q^*\Omega^i_{\mathbb{P}(\underline{E})/S} \longrightarrow \Omega^i_{\mathbb{V}(\underline{E})/S} \longrightarrow \Omega^1_{\mathbb{V}(\underline{E})^*/\mathbb{P}(\underline{E})} \otimes q^*\Omega^{i-1}_{\mathbb{P}(\underline{E})/S} \longrightarrow 0 \;\; .
\end{equation*}
Celles-ci se réécrivent (1.1.2)

\begin{equation}\tag{1.1.3}
0 \longrightarrow q^*\Omega^i_{\mathbb{P}(\underline{E})/S} \longrightarrow \Omega^i_{\mathbb{V}(\underline{E})^*/S} \xrightarrow{\;X\lfloor\;} q^*\Omega^{i-1}_{\mathbb{P}(\underline{E})/S} \longrightarrow 0 \;\; .
\end{equation}
On a canoniquement

\begin{equation}\tag{1.1.4}
\Omega^i_{\mathbb{V}(\underline{E})/S} \;=\; a^* \overset{i}{\wedge} \underline{E} \;\; .
\end{equation}
Les suites exactes courtes (1.1.3) s'obtiennent donc en fractionnant une suite exacte longue sur $\mathbb{V}(\underline{E})^*$

\begin{equation}\tag{1.1.5}
\ldots \longrightarrow {a'}^{*} \overset{i}{\wedge} \underline{E} \xrightarrow{\;X\lfloor\;} {a'}^{*} \overset{i-1}{\wedge} \underline{E} \xrightarrow{\;X\lfloor\;} {a'}^{*} \overset{i-2}{\wedge} \underline{E} \longrightarrow \ldots \;\; .
\end{equation}

\SGApage{II}{42}

\SGAvar{Cette suite exacte longue n'est autre que la restriction à $\mathbb{V}(\underline{E})^*$ de la résolution de Kozul de $\underline{O}_{e(S)}$ sur $\mathbb{V}(\underline{E})$.

}{Cette suite exacte longue n'est autre que la restriction à $\mathbb{V}(\underline{E})^*$ de la résolution de Koszul de $\underline{O}_{e(S)}$ sur $\mathbb{V}(\underline{E})$.

}{F7II-XI-002}Projetons la suite exacte longue (1.1.5) sur $\mathbb{P}(\underline{E})$ par le morphisme affine $q$. En vertu de (1.1.1) et de la formule $q_*{a'}^{*} = q_*q^*p^*$, on obtient une suite exacte longue

\SGAvar{\begin{equation}\tag{1.1.6}
\ldots \longrightarrow \sum_{n \in \mathbb{Z}} p^* \overset{i}{\wedge} \underline{E}(n) \xrightarrow{\;X\lfloor\;} \sum_{n \in \mathbb{Z}} p^* \overset{i-1}{\wedge} \underline{E}(n) \longrightarrow \ldots \;\; .
\end{equation}
Celle-ci se fractionne en suite exactes courtes, images de (1.1.3),

}{\begin{equation}\tag{1.1.6}
\ldots \longrightarrow \sum_{n \in \mathbb{Z}} p^* \overset{i}{\wedge} \underline{E}(n) \xrightarrow{\;X\lfloor\;} \sum_{n \in \mathbb{Z}} p^* \overset{i-1}{\wedge} \underline{E}(n) \longrightarrow \ldots \;\; .
\end{equation}
Celle-ci se fractionne en suites exactes courtes, images de (1.1.3),

}{F7II-XI-003}\begin{equation}\tag{1.1.7}
0 \longrightarrow \sum_{n \in \mathbb{Z}} \Omega^i_{\mathbb{P}(\underline{E})/S}(n) \longrightarrow \sum_{n \in \mathbb{Z}} p^* \overset{i}{\wedge} \underline{E}(n) \xrightarrow{\;X\lfloor\;} \sum_{n \in \mathbb{Z}} \Omega^{i-1}_{\mathbb{P}(\underline{E})/S}(n) \to 0 \;\; .
\end{equation}
Les suites exactes (1.1.7) sont sommes de suites exactes

\begin{equation}\tag{1.1.8}
0 \longrightarrow \Omega^i_{\mathbb{P}(\underline{E})/S}(n) \longrightarrow p^* \overset{i}{\wedge} \underline{E}(n-i) \xrightarrow{\;X\lfloor\;} \Omega^{i-1}_{\mathbb{P}(\underline{E})/S}(n) \longrightarrow 0 \;\; .
\end{equation}
Rappelons (EGA III 2.1.15 ou FAC) que les $R^i p_* \underline{O}(k)$ sont localement libres de formation compatible à tout changement de base, et que

\begin{equation}\tag{1.1.9}
R^i p_* \, \underline{O}(k) = 0 \ \text{si } i \neq 0,\ r, \ \text{ou si } i = 0,\ k < 0, \ \text{ou si } i = r,\ k > -r-1,
\end{equation}
et que $R^i p_* p^* \overset{j}{\wedge} \underline{E}(k)$ est donc nul sous ces mêmes conditions.

Les suites exactes (1.1.8) fournissent une résolution

\begin{equation}\tag{1.1.10}
0 \longrightarrow \Omega^i_{\mathbb{P}(\underline{E})/S}(n) \xrightarrow{\;E\;} p^* \overset{i}{\wedge} \underline{E}(n-i) \xrightarrow{\;X\lfloor\;} \ldots \longrightarrow p^* \overset{0}{\wedge} \underline{E}(n-0) \longrightarrow 0
\end{equation}
de $\Omega^i_{\mathbb{P}(\underline{E})/S}(n)$. La suite spectrale définie par cette résolution dégénère $(E_1 = E_\infty)$ pour des raisons de degré (1.1.9) ; pour $0 < j < r$, elle fournit

\begin{equation}\tag{1.1.11}
R^j p_* \Omega^i_{\mathbb{P}(\underline{E})/S}(n) = \mathcal{H} \, ( \ldots \xrightarrow{\;X\lfloor\;} p_* \overset{i-j}{\wedge} \underline{E}(n-i+j) \xrightarrow{\;X\lfloor\;} \ldots ) \;\; .
\end{equation}
Puisque $r \geq 1$, un argument de profondeur montre que

\begin{equation}\tag{1.1.12}
\sum_{n \in \mathbb{Z}} p_*p^* \overset{i}{\wedge} \underline{E}(n) = p_*q_*q^*p^* \overset{i}{\wedge} \underline{E} = {a'}_{*}{a'}^{*} \overset{i}{\wedge} \underline{E} = a_*a^* \overset{i}{\wedge} \underline{E} \;\; .
\end{equation}

\SGApage{II}{43}

 de sorte que $p_*p^* \wedge^i \underline{E}(n)$ est la partie homogène de poids $n + i$ de $a_*a^* \wedge^i \underline{E}$, soit $\mathrm{Sym}^n(\underline{E}) \otimes \wedge^i \underline{E}$ . Si on se rappelle que (1.1.5) provient de la résolution de Koszul, on trouve par (1.1.11) que, pour $0 < j < r$, $R^j p_* \Omega^i_{\mathbb{P}(\underline{E})/S}(n)$ n'est non nul que si $i-j = 0$ et $n-i+j = 0$, i.e. si $n = 0$ et $i = j$. Dans ce cas, de plus, $R^i p_* \Omega^i_{\mathbb{P}(\underline{E})/S}$ est isomorphe à $\underline{O}_S$ .

Pour $j = 0$, on déduit encore de (1.1.10) que

\SGAvar{\begin{equation}\tag{1.1.13}
p_*\Omega^i_{\mathbb{P}(\underline{E})/S}(n) \;=\; \mathrm{Ker}(X \llcorner \;:\; p_* \wedge^i \underline{E}(n-i) \longrightarrow p_* \wedge^{i-1} \underline{E}(n-i+1) \qquad ,
\end{equation}
et le même argument montre que $X \llcorner$ est injectif lorsque $n-i \leq 0$.

}{\begin{equation}\tag{1.1.13}
p_*\Omega^i_{\mathbb{P}(\underline{E})/S}(n) \;=\; \mathrm{Ker}(X \llcorner \;:\; p_* \wedge^i \underline{E}(n-i) \longrightarrow p_* \wedge^{i-1} \underline{E}(n-i+1)) \qquad ,
\end{equation}
et le même argument montre que $X \llcorner$ est injectif lorsque $n-i \leq 0$.

}{F7II-XI-004}En vertu de(1.1.8) pour $n = 0$, $i = r+1$, on a $\Omega^r_{\mathbb{P}(\underline{E})/S} \simeq p^* \wedge^{r+1} \underline{E}(-r-1)$. Lorsque $S$ est le spectre d'un corps, les assertions de nullité de groupes de cohomologie de 1.1 (ii) et (iii) résultent de ce qui précède pour $j \neq r$, et s'en déduisent par dualité ([2] V 11.2 p. 211) pour $j = r$. En particulier, quels que soient $j$ et $n$, il existe au plus un $i$ tel que $R^i p_* \Omega^j_{\mathbb{P}(\underline{E})/S}(n) \neq 0$, ce qui implique (i) puisque $\Omega^j_{\mathbb{P}(\underline{E})/S}$ est plat sur $S$ (le plus dur est de trouver la référence dans EGA III ; je suggère : 6.10.5 + 7.4.1 + 7.5.5 + un passage à la limite).

\SGAvar{Sachant que les $R^i p_* \Omega^j_{\mathbb{P}(\underline{E})/S}$ sont des faisceaux inversibles pour $0 \leq i \leq r$, il reste seulement à prouver que, pour $S$ le spectre d'un corps, on a $\eta^i \neq 0$. En effet $\eta^r$ est la classe de cohomologie d'un point rationnel (intersection de $r$ hyperplans), donc une base de $H^r(\mathbb{P}^r,\Omega^r)$ (FGA exp. 149 n$^{\circ}$4).

}{Sachant que les $R^i p_* \Omega^i_{\mathbb{P}(\underline{E})/S}$ sont des faisceaux inversibles pour $0 \leq i \leq r$, il reste seulement à prouver que, pour $S$ le spectre d'un corps, on a $\eta^i \neq 0$. En effet $\eta^r$ est la classe de cohomologie d'un point rationnel (intersection de $r$ hyperplans), donc une base de $H^r(\mathbb{P}^r,\Omega^r)$ (FGA exp. 149 n$^{\circ}$4).

}{F7II-XI-005}1.2. Il est possible de calculer, directement à partir de (1.1.10) ou (1.1.13), les caractéristiques d'Euler-Poincaré des $\Omega^i(n)$. Le résultat, peu appétissant, est

\SGAvar{\begin{equation*}
\chi(\mathbb{P}^r_k,\Omega^j_{\mathbb{P}^r_k}(n)) \;=\; \sum_{a \leq i \leq j} \; (-1)^{i-j} \binom{r+1}{i}\binom{r+n-i}{r} \; .
\end{equation*}

}{\begin{equation*}
\chi(\mathbb{P}^r_k,\Omega^j_{\mathbb{P}^r_k}(n)) \;=\; \sum_{0 \leq i \leq j} \; (-1)^{i-j} \binom{r+1}{i}\binom{r+n-i}{r} \; .
\end{equation*}

}{F7II-XI-006}\SGApage{II}{44}

\underline{Proposition} 1.3. \underline{Soient} $f : X \longrightarrow S$ \underline{un morphisme propre et lisse,} $\underline{W}$ \underline{un faisceau localement libre de rang} $c$ \underline{sur} $X$, $\underline{F}$ \underline{un faisceau cohérent sur} $X$, $d$ \underline{un entier} $> 0$ \underline{et} $s : \underline{W} \longrightarrow \underline{O}_X$ \underline{une section de} $\underline{W}^{\vee}$. \underline{On suppose que} $S$ \underline{est noethérien et que}

(a) \underline{Le sous-schéma} $H$ \underline{de} $X$ \underline{défini par l'annulation de} $s$ \underline{est lisse sur} $S$.

\SGAvar{(b) \underline{Localement sur} $X$, \underline{si les} $s_i$ \underline{sont les coordonnées de} $S$ \underline{dans une base de} $\underline{W}^{\vee}$ \underline{les} $s_i$ \underline{forment une suite} $\underline{O}_X$\underline{-régulier et} $\underline{F}$\underline{-régulière.}

}{(b) \underline{Localement sur} $X$, \underline{si les} $s_i$ \underline{sont les coordonnées de} $s$ \underline{dans une base de} $\underline{W}^{\vee}$ \underline{les} $s_i$ \underline{forment une suite} $\underline{O}_X$\underline{-régulière et} $\underline{F}$\underline{-régulière.}

}{F7II-XI-007}(c) \underline{Quelle que soit la suite d'entiers} $k_i$, \underline{non tous nuls, on a}

\begin{equation*}
R^i f_* \; (\underset{i}{\otimes} \; \wedge^{k_i} \; \underline{W} \otimes \Omega^j_{X/S} \otimes \underline{F}) \;=\; 0 \;\; \underline{\text{pour}} \;\; i+j < d \qquad .
\end{equation*}
\underline{Alors, on a} :

\SGAvar{\begin{equation*}\tag{i}
R^i f_* \; (\Omega^j_{H/S} \otimes \underline{F}) \; \overset{\sim}{\longrightarrow} \; R^i f_*(\Omega^j_{X/S} \otimes \underline{F}) \;\; \underline{\text{pour}} \;\; i+j < d-c,
\end{equation*}
\begin{equation}\tag{ii}
R^i f_* \; (\Omega^j_{H/S} \otimes \underline{F}) \; \hookrightarrow \; R^i f_*(\Omega^j_{X/S} \otimes \underline{F}) \;\; \underline{\text{pour}} \;\; i+j = d-c \qquad .
\end{equation}
Le faisceau conormal $N_{H/X}$ de $H$ dans $X$ est isomorphe à la restriction de $\underline{W}$ à $H$, de sorte que la suite exacte

\begin{equation}\tag{1.3.1}
0 \longrightarrow N_{H/X} \longrightarrow \Omega^1_{X/S} \otimes \underline{O}_Y \longrightarrow \Omega^1_{H/S} \longrightarrow 0
\end{equation}
définit une filtration décroissante de $\Omega^j_{X/S} \otimes \underline{O}_Y$ dont les quotients successifs sont donnés par

\begin{equation}\tag{1.3.2}
\mathrm{Gr}^k(\Omega^j_{X/S} \otimes \underline{O}_Y) \qquad \simeq \qquad \wedge^k \; \underline{W} \otimes \Omega^{j-k}_{H/S} \qquad .
\end{equation}
Il résulte de (b) que l'on a encore

\begin{equation}\tag{1.3.3}
\mathrm{Gr}^k(\Omega^j_{X/S} \otimes \underline{O}_Y \otimes \underline{F}) \;=\; \wedge^k \; \underline{W} \otimes \Omega^{j-k}_{H/S} \otimes \underline{F} \qquad .
\end{equation}
On prouvera (i) et (ii) par récurrence sur $j$, simultanément pour tous les $\underline{F}$ vérifiant (b) et (c). Ces assertions sont vides pour $j < 0$ ;
\SGApage{II}{45}
supposons donc (i) et (ii) prouvés pour $j' < j$. L'hypothèse de récurrence s'applique aux faisceaux $\stackrel{k}{\wedge} \underline{W} \otimes \underline{F}$ ; elle fournit, par (c) :

\begin{equation}\tag{1.3.4}
R^i f_* ( \stackrel{k}{\wedge}\underline{W} \otimes \Omega^{j-k}_{H/S} \otimes \underline{F}) = 0 \quad \mbox{pour } k > 0 \mbox{ et } i+j < d-c+k \quad .
\end{equation}
La filtration considérée plus haut du $\Omega^j_{X/S} \otimes \underline{O}_Y \otimes \underline{F}$ définit une suite spectrale

\begin{equation*}
E_1^{p,i-p} \; = \; R^i f_*( \stackrel{p}{\wedge} \underline{W} \otimes \Omega^{j-p}_{H/S} \otimes \underline{F}) \Longrightarrow R^i f_*(\Omega^j_{X/S} \otimes \underline{O}_Y \otimes \underline{F}) \quad .
\end{equation*}
Par (1.3.4), on a $E_1^{p,i-p} = 0$ si $p > 0$ et $i+j \leq d-c$ d'où l'on tire que

\begin{equation}\tag{1.3.5}
\left\{ \begin{array}{l} R^i f_*(\Omega^j_{X/S} \otimes \underline{O}_Y \otimes \underline{F}) \; \stackrel{\sim}{\longrightarrow} \; R^i f_*(\Omega^j_{H/S} \otimes \underline{F}) \mbox{ pour } i+j < d-c \\[4ex] R^i f_*(\Omega^j_{X/S} \otimes \underline{O}_Y \otimes \underline{F}) \; \lhook\joinrel\longrightarrow \; R^i f_*(\Omega^j_{H/S} \otimes \underline{F}) \mbox{ pour } i+j = d-c \quad . \end{array} \right.
\end{equation}
D'après l'hypothèse (b), pour une différentielle convenable, les $\stackrel{k}{\wedge} \underline{W} \otimes \Omega^j_{X/S} \otimes \underline{F}$ forment une résolution à gauche (de Koszul) de $\Omega^j_{X/S} \otimes \underline{O}_Y \otimes \underline{F}$. Cette filtration définit une suite spectrale

\begin{equation}\tag{1.3.6}
E_1^{-p,q} \; = \; R^q f_*( \stackrel{p}{\wedge}\underline{W} \otimes \Omega^j_{X/S} \otimes \underline{F}) \Longrightarrow R^{q-p} f_*(\Omega^j_{X/S} \otimes \underline{O}_Y \otimes \underline{F}) \quad .
\end{equation}
D'après l'hypothèse (c), on a $E_1^{-p,q} = 0$ pour $p > 0$ et $q+j < d$, ou, trivialement, pour $p \notin [0,c]$. Dès lors,

}{\begin{equation*}\tag{i}
R^i f_* \; (\Omega^j_{X/S} \otimes \underline{F}) \; \overset{\sim}{\longrightarrow} \; R^i f_*(\Omega^j_{H/S} \otimes \underline{F}) \;\; \underline{\text{pour}} \;\; i+j < d-c,
\end{equation*}
\begin{equation}\tag{ii}
R^i f_* \; (\Omega^j_{X/S} \otimes \underline{F}) \; \hookrightarrow \; R^i f_*(\Omega^j_{H/S} \otimes \underline{F}) \;\; \underline{\text{pour}} \;\; i+j = d-c \qquad .
\end{equation}
Le faisceau conormal $N_{H/X}$ de $H$ dans $X$ est isomorphe à la restriction de $\underline{W}$ à $H$, de sorte que la suite exacte

\begin{equation}\tag{1.3.1}
0 \longrightarrow N_{H/X} \longrightarrow \Omega^1_{X/S} \otimes \underline{O}_H \longrightarrow \Omega^1_{H/S} \longrightarrow 0
\end{equation}
définit une filtration décroissante de $\Omega^j_{X/S} \otimes \underline{O}_H$ dont les quotients successifs sont donnés par

\begin{equation}\tag{1.3.2}
\mathrm{Gr}^k(\Omega^j_{X/S} \otimes \underline{O}_H) \qquad \simeq \qquad \wedge^k \; \underline{W} \otimes \Omega^{j-k}_{H/S} \qquad .
\end{equation}
Il résulte de (b) que l'on a encore

\begin{equation}\tag{1.3.3}
\mathrm{Gr}^k(\Omega^j_{X/S} \otimes \underline{O}_H \otimes \underline{F}) \;=\; \wedge^k \; \underline{W} \otimes \Omega^{j-k}_{H/S} \otimes \underline{F} \qquad .
\end{equation}
On prouvera (i) et (ii) par récurrence sur $j$, simultanément pour tous les $\underline{F}$ vérifiant (b) et (c). Ces assertions sont vides pour $j < 0$ ;
\SGApage{II}{45}
supposons donc (i) et (ii) prouvés pour $j' < j$. L'hypothèse de récurrence s'applique aux faisceaux $\stackrel{k}{\wedge} \underline{W} \otimes \underline{F}$ ; elle fournit, par (c) :

\begin{equation}\tag{1.3.4}
R^i f_* ( \stackrel{k}{\wedge}\underline{W} \otimes \Omega^{j-k}_{H/S} \otimes \underline{F}) = 0 \quad \mbox{pour } k > 0 \mbox{ et } i+j < d-c+k \quad .
\end{equation}
La filtration considérée plus haut du $\Omega^j_{X/S} \otimes \underline{O}_H \otimes \underline{F}$ définit une suite spectrale

\begin{equation*}
E_1^{p,i-p} \; = \; R^i f_*( \stackrel{p}{\wedge} \underline{W} \otimes \Omega^{j-p}_{H/S} \otimes \underline{F}) \Longrightarrow R^i f_*(\Omega^j_{X/S} \otimes \underline{O}_H \otimes \underline{F}) \quad .
\end{equation*}
Par (1.3.4), on a $E_1^{p,i-p} = 0$ si $p > 0$ et $i+j \leq d-c$ d'où l'on tire que

\begin{equation}\tag{1.3.5}
\left\{ \begin{array}{l} R^i f_*(\Omega^j_{X/S} \otimes \underline{O}_H \otimes \underline{F}) \; \stackrel{\sim}{\longrightarrow} \; R^i f_*(\Omega^j_{H/S} \otimes \underline{F}) \mbox{ pour } i+j < d-c \\[4ex] R^i f_*(\Omega^j_{X/S} \otimes \underline{O}_H \otimes \underline{F}) \; \lhook\joinrel\longrightarrow \; R^i f_*(\Omega^j_{H/S} \otimes \underline{F}) \mbox{ pour } i+j = d-c \quad . \end{array} \right.
\end{equation}
D'après l'hypothèse (b), pour une différentielle convenable, les $\stackrel{k}{\wedge} \underline{W} \otimes \Omega^j_{X/S} \otimes \underline{F}$ forment une résolution à gauche (de Koszul) de $\Omega^j_{X/S} \otimes \underline{O}_H \otimes \underline{F}$. Cette filtration définit une suite spectrale

\begin{equation}\tag{1.3.6}
E_1^{-p,q} \; = \; R^q f_*( \stackrel{p}{\wedge}\underline{W} \otimes \Omega^j_{X/S} \otimes \underline{F}) \Longrightarrow R^{q-p} f_*(\Omega^j_{X/S} \otimes \underline{O}_H \otimes \underline{F}) \quad .
\end{equation}
D'après l'hypothèse (c), on a $E_1^{-p,q} = 0$ pour $p > 0$ et $q+j < d$, ou, trivialement, pour $p \notin [0,c]$. Dès lors,

}{F7II-XI-008}\begin{equation}\tag{1.3.7}
\left\{ \begin{array}{l} R^i f_*(\Omega^j_{X/S} \otimes \underline{F}) \; \stackrel{\sim}{\longrightarrow} \; R^i f_*(\Omega^j_{H/S} \otimes \underline{F}) \mbox{ pour } i+j < d-c \\[4ex] R^i f_*(\Omega^j_{X/S} \otimes \underline{F}) \; \lhook\joinrel\longrightarrow \; R^i f_*(\Omega^j_{H/S} \otimes \underline{F}) \mbox{ pour } i+j = d-c \quad . \end{array} \right.
\end{equation}
Les assertions (i) et (ii) résultent maintenant de (1.3.5) et (1.3.7).


\SGApage{II}{46}

1.4. Soient $\underline{E}$ un faisceau de modules localement libres de rang $r+1$ sur un schéma $S$, $\mathbb{P}(\underline{E})$ le fibré projectif correspondant, $d$ un entier $\leq r$ et $\underline{a} = (a_1 \ldots a_d)$ une suite d'entiers. Une \underline{intersection complète} relative dans $\mathbb{P}(\underline{E})$, de multidegré $\underline{a}$, est un sous-schéma de $\mathbb{P}(\underline{E})$, fibre par fibre de codimension $d$, qui, localement sur $S$ pour la topologie étale, est intersection de $d$ hypersurfaces de degrés respectifs $a_1 \ldots a_d$.

En vertu de 1.1, les hypothèses (b) et (c) de 1.3 sont vérifiées pour $X/S = \mathbb{P}(\underline{E})/S$, $\underline{F} = \underline{O}_X$ , $\underline{W} = \sum\limits_i \underline{O}(-a_i)$ et pour $s$ une section de $\underline{W}^{\vee}$ définissant une intersection complète relative $Y \subset \mathbb{P}(\underline{E})$ de multidegré $\underline{a}$.

\underline{Théorème} 1.5. \underline{Sous les hypothèses} 1.4, \underline{soit} $n = r-d$ \underline{la dimension relative de} $Y$ \underline{et} $\eta \in H^o(S,R^1f_*\Omega^1_{Y/S})$ \underline{la première classe de Chern de} $\underline{O}(1)$

\begin{center}
\begin{tikzcd}[column sep=large, row sep=large]
Y \arrow[rr, hook] \arrow[dr, dash, "f"'] \& {} \& \mathbb{P}(\underline{E}) \arrow[dl, "g"] \\
 {} \& S \&
{} \end{tikzcd}\quad .
\end{center}









\underline{Si} $Y$ \underline{est lisse sur} $S$, \underline{alors} :

(i) \underline{Les} $R^i f_* \Omega^j_{Y/S}$ \underline{sont localement libres}, \underline{et leur formation est compatible à tout changement de base}.

\SGAvar{(ii) \underline{La suite spectrale} $E_1^{pq} = R^q f_* \Omega^p_{Y/S} \Longrightarrow R^{p+q} f_*(\Omega^*_{Y/S})$ \underline{reliant les cohomologie de Hodge et de De Rham dégénère} $(E_1 = E_\infty)$, \underline{son aboutissement est localement libre et sa formation est compatible à tout changement de base}.

}{(ii) \underline{La suite spectrale} $E_1^{pq} = R^q f_* \Omega^p_{Y/S} \Longrightarrow R^{p+q} f_*(\Omega^*_{Y/S})$ \underline{reliant les cohomologies de Hodge et de De Rham dégénère} $(E_1 = E_\infty)$, \underline{son aboutissement est localement libre et sa formation est compatible à tout changement de base}.

}{F7II-XI-009}(iii)$_a$ \underline{On a}

\SGAvar{\begin{equation*}
R^i f_* \Omega^j_{X/S} = 0 \quad \underline{\mbox{si}} \; i \neq j \; \underline{\mbox{et}} \; i+j \neq n \; .
\end{equation*}

}{\begin{equation*}
R^i f_* \Omega^j_{Y/S} = 0 \quad \underline{\mbox{si}} \; i \neq j \; \underline{\mbox{et}} \; i+j \neq n \; .
\end{equation*}

}{F7II-XI-010}\SGApage{II}{47}

(iii)$_b$ $\eta^i$ \underline{est une base de} $R^i f_* \Omega^i_{Y/S}$ \underline{si} $2i < n$

(iii)$_c$ \underline{si} $n<2i\leq 2n$ \underline{et si} $\alpha_i$ \underline{est la base de} $R^i f_* \Omega^i_{Y/S}$ \underline{duale de la base} $\eta^{n-i}$ de $R^{n-i} f_* \Omega^{n-i}_{Y/S}$, \underline{alors}

\begin{equation*}
\eta^i \;=\; \left( \prod_1^d \; a_i \right) \alpha_i \qquad .
\end{equation*}
(iii)$_d$ \underline{Si} $n = 2m$, \underline{alors} $\eta^m \neq 0$ (\underline{si} $S \neq \emptyset$) \underline{et le quotient} $R^m f_* \Omega^m_{Y/S}/\underline{O}_S . \eta^m$ \underline{est localement libre}.

Il suffit de prouver 1.5 dans un cas ``universel'', ce qui permet de supposer $S$ lisse sur $\mathrm{Spec}(\mathbb{Z})$. Dans ce cas, l'assertion (ii) résulte de (i), car elle est vraie sur $\mathbb{C}$, donc sur tout corps de caractéristique $0$, d'après la théorie de Hodge. De plus, $S$ étant réduit, (i) résulte de (iii), supposé démontré pour $S$ le spectre d'un corps : en effet, d'après (iii), pour chaque $j$ les fonctions

\SGAvar{\begin{equation*}
s \longrightarrow \dim_{k(s)} \; (H^i(Y_s,\Omega^j_{Ys})
\end{equation*}
sont constantes sur $S$, sauf peut-être pour $i = j$, donc sont toutes constantes puisque leur somme alternée l'est (EGA III 7.9.2) et on conclut par EGA III 6.10.5 + 7.6.9 + 7.7.5.

}{\begin{equation*}
s \longrightarrow \dim_{k(s)} \; (H^i(Y_s,\Omega^j_{Ys}))
\end{equation*}
sont constantes sur $S$, sauf peut-être pour $i + j = n$, donc sont toutes constantes puisque leur somme alternée l'est (EGA III 7.9.2) et on conclut par EGA III 6.10.5 + 7.6.9 + 7.7.5.

}{F7II-XI-011}En vertu de (1.4) et (1.3), on a

\SGAvar{\begin{equation}\tag{1.5.1}
\left\{
\begin{array}{ccl}
R^i f_* \Omega^j_{Y/S} & \stackrel{\sim}{\longrightarrow} & R^i g_* \Omega^j_{\mathbb{P}(\underline{E})/S} \qquad \mbox{pour } i+j < n \\[3ex]
R^i f_* \Omega^j_{Y/S} & \lhook\joinrel\longrightarrow & R^i g_* \Omega^j_{\mathbb{P}(\underline{E})/S} \qquad \mbox{pour } i+j = n \qquad .
\end{array}
\right.
\end{equation}
Si $S$ est le spectre d'un corps, les assertions (iii)$_a$ (iii)$_b$ (iii)$_d$ résultent de (1.5.1) et de la dualité de Serre ([2] V 11.2 p.211) La classe de
\SGApage{II}{48}
cohomologie $\eta^n$ est la classe de cohomologie d'un $0$-cycle de degré $\prod_1^d a_i$ (le degré de l'intersection complète), de sorte que

}{\begin{equation}\tag{1.5.1}
\left\{
\begin{array}{ccl}
R^i f_* \Omega^j_{Y/S} & \stackrel{\sim}{\longrightarrow} & R^i g_* \Omega^j_{\mathbb{P}(\underline{E})/S} \qquad \mbox{pour } i+j < n \\[3ex]
R^i f_* \Omega^j_{Y/S} & \lhook\joinrel\longrightarrow & R^i g_* \Omega^j_{\mathbb{P}(\underline{E})/S} \qquad \mbox{pour } i+j = n \qquad .
\end{array}
\right.
\end{equation}
Si $S$ est le spectre d'un corps, les assertions (iii)$_a$ (iii)$_b$ (iii)$_d$ résultent de (1.5.1) et de la dualité de Serre ([2] V 11.2 p.211). La classe de
\SGApage{II}{48}
cohomologie $\eta^n$ est la classe de cohomologie d'un $0$-cycle de degré $\prod_1^d a_i$ (le degré de l'intersection complète), de sorte que

}{F7II-XI-012}\begin{equation*}
\eta^n \;=\; \left( \prod_1^d \; a_i \right) \alpha_n \qquad .
\end{equation*}
Par définition, on a $\eta^{n-i} \alpha_i = \alpha_n$, et(iii)$_c$ en résulte.

Les assertions (i) et (ii) sont donc vraies, et les assertions (iii), dans le cas universel, résultent de (i) et de (iii) supposé prouvé pour $S$ le spectre d'un corps.

\underline{Théorème} 1.6. \underline{Gardons les hypothèses et notations de} 1.4, 1.5, \underline{supposons} $Y$ \underline{lisse,} \underline{soit} $\ell$ \underline{un nombre premier premier aux caractéristiques résiduelles de} $S$, \underline{et désignons par} $\eta_\ell \in H^0(S,R^2 f_*( \mathbb{Z}_\ell(1)))$ \underline{la première classe de Chern} $\ell$\underline{-adique de} $\underline{O}(1)$ (SGA 5 VII). \underline{Alors}

(i) $R^i f_* \mathbb{Z}_\ell = 0$ \underline{si} $i$ \underline{est impair et distinct de} $n$

(ii) \underline{Si} $2i < n$, \underline{alors} $R^{2i} f_* \mathbb{Z}_\ell(i)$ \underline{est isomorphe à} $\mathbb{Z}_\ell$, \underline{et admet} $\eta_\ell^i$ \underline{pour base}

(iii) \underline{Si} $n < 2i \leq 2n$, \underline{alors} $R^{2i} f_* \mathbb{Z}_\ell(i)$ \underline{est isomorphe à} $\mathbb{Z}_\ell$, \underline{et admet} $\eta_\ell^i / \prod_1^d a_i$ \underline{pour base.}

(iv) \underline{Le} $\mathbb{Z}_\ell$\underline{-faisceau} $R^n f_* \mathbb{Z}_\ell$ \underline{est constant tordu constructible sans torsion et,} \underline{si} $n$ \underline{est pair,} \underline{alors} $\eta_\ell^{n/2} \neq 0$ (\underline{si} $S \neq \emptyset$) \underline{et} $R^n f_* \mathbb{Z}_\ell(n/2)/ \mathbb{Z}_\ell . \eta_\ell^{n/2}$ \underline{est sans torsion.}

Les $\mathbb{Z}_\ell$-faisceaux $R^i f_* \mathbb{Z}_\ell$ sont constants tordus constructibles, et leur formation est compatible à tout changement de base, d'après SGA 4 XVI 2.2. Il suffit donc de prouver 1.6 pour $S$ le spectre d'un corps algébriquement clos. Une application itérée de SGA 4 XIV 3.3 montre alors, pour tout $m$, que la flèche canonique

\SGAvar{\begin{equation*}
H^q_Y \; (\mathbb{P}(\underline{E}), \; \mathbb{Z}/\ell^m) \longrightarrow H^q(Y, \; \mathbb{Z}/\ell^m)
\end{equation*}

}{\begin{equation*}
H^q_Y \; (\mathbb{P}(\underline{E}), \; \mathbb{Z}/\ell^m) \longrightarrow H^q(\mathbb{P}(\underline{E}), \; \mathbb{Z}/\ell^m)
\end{equation*}

}{F7II-XI-013}\SGApage{II}{49}
est bijective pour $q > r + d$ et surjective pour $q = r + d$. De là, et de la dualité de Poincaré (SGA 4 XVIII) on déduit (SGA 5 VII) que la flèche de restriction

\begin{equation*}
H^q(\ \mathbb{P}(\underline{E}), \mathbb{Z}_\ell) \longrightarrow H^q(Y, \mathbb{Z}_\ell)
\end{equation*}
est bijective pour $q < n$, injective et de conoyau sans torsion pour $q = n$. Ceci prouve (i) pour $i < n$, (ii) et (iv). On vérifie comme en 1.5 que $\eta_\ell^n$ est $\prod_1^d a_i$ fois la classe canonique, et (i) et (iii) s'obtiennent dès lors par dualité de Poincaré (SGA 4 XVIII).

1.7. Si $Y$ est une intersection complète de dimension au moins 3 dans un espace projectif $\mathbb{P}^r(k)$ sur un corps $k$, on sait que $\mathrm{Pic}(Y)$ est isomorphe à $\mathbb{Z}$, et est engendré par $\underline{O}(1)$ (SGA 2 XII 3.7). En dimension 2, et si $Y$ est lisse, on sait encore prouver :

\underline{Théorème} 1.8. \underline{Soit} $Y$ \underline{une surface lisse \ intersection complète dans l'espace projectif} $\mathbb{P}^r(k)$ \underline{sur un corps algébriquement clos} $k$. \underline{Alors,} $\underline{\mathrm{Pic}}^\tau(Y) = 0$, \underline{et la classe de} $\underline{O}(1)$ \underline{n'est pas divisible dans le groupe de Néron-Severi} $\mathrm{NS}(Y)$.

Puisque $H^1(Y,\underline{O}) = 0$ (1.5 (iii)$_a$), on sait que $\underline{\mathrm{Pic}}^0(Y) = 0$. La nullité de $\underline{\mathrm{Pic}}^\tau(Y)$ résulte alors du lemme bien connu suivant, du fait que $H^2(Y, \mathbb{Z}_\ell(1))$ est sans torsion (1.6 (iv)) et du fait que $H^0(Y,\Omega^1_Y) = 0$ (1.5 (iii)$_a$).

\underline{Lemme} 1.9. \underline{Soit} $Y$ \underline{une variété algébrique complète sur un corps algébriquement clos} $k$.

(i) \underline{Si} $\ell$ \underline{est premier à la caractéristique de} $k$, \underline{alors la} $\ell$-\underline{torsion de} $\mathrm{NS}(Y)$ \underline{s'identifie à la} $\ell$-\underline{torsion de} $H^2(Y, \mathbb{Z}_\ell(1))$.
\SGApage{II}{50}
(ii) \underline{Si} $k$ \underline{est de caractéristique non nulle} $p$ \underline{et si} $Y$ \underline{est lisse, alors le noyau de la multiplication par} $p$ \underline{dans} $\mathrm{Pic}(Y)$ \underline{s'identifie au sous-groupe de} $H^0(Y,\Omega^1_Y)$ \underline{formé des différentielles localement logarithmiques.}

Les suites exactes de Kummer

\begin{equation*}
0 \longrightarrow \mu_{\ell^n} \longrightarrow G_m \xrightarrow{\ x \mapsto x^{\ell^n}\ } G_m \longrightarrow 0
\end{equation*}
définissent des suites exactes

\begin{equation*}
0 \longrightarrow \mathrm{Pic}(Y)/\ell^n\ \mathrm{Pic}(Y) \longrightarrow H^2(Y,\mu_{\ell^n}) \longrightarrow (H^2(G_m))_{\ell^n} \longrightarrow 0 \quad ,
\end{equation*}
où $\mathrm{Pic}(Y)/\ell^n\ \mathrm{Pic}(Y) = \mathrm{NS}(Y) \otimes \mathbb{Z}/\ell^n$. Par passage à la limite projective, on trouve

\begin{equation*}
0 \longrightarrow \mathrm{NS}(Y) \otimes \mathbb{Z}_\ell \longrightarrow H^2(Y, \mathbb{Z}_\ell(1)) \longrightarrow \mathrm{Hom}(\mathbb{Q}_\ell / \mathbb{Z}_\ell,\ H^2(G_m)) \longrightarrow 0 \quad ,
\end{equation*}
et le conoyau est sans torsion puisque $\mathbb{Q}_\ell / \mathbb{Z}_\ell$ est divisible. L'assertion (i) en résulte.

Pour prouver (ii), il suffit de noter que la suite exacte

\begin{equation*}
0 \longrightarrow \underline{O}^*_Y \xrightarrow{\ x \longmapsto x^p\ } \underline{O}^*_Y \xrightarrow{\ df/f\ } \Omega^1_Y
\end{equation*}
définit une suite exacte courte

\SGAvar{\begin{equation*}
0 \longrightarrow \underline{O}^*_Y \longrightarrow \underline{O}^*_Y \longrightarrow \Omega^1_{Y\,\mathrm{loc}\,\mathrm{log}} \longrightarrow 0
\end{equation*}
et de prendre la suite exacte longue de cohomologie correspondante. Ce raisonnement est dû à Cartier, qui prouve de plus que, si l'on désigne par $C$ l'opération de Cartier sur $H^0(Y,\Omega^1_Y)$, alors les différentielles localement logarithmique forment le noyau de $C - \mathrm{Id}$.

}{\begin{equation*}
0 \longrightarrow \underline{O}^*_Y \longrightarrow \underline{O}^*_Y \longrightarrow \Omega^1_{Y\,\mathrm{loc}\,\mathrm{log}} \longrightarrow 0
\end{equation*}
et de prendre la suite exacte longue de cohomologie correspondante. Ce raisonnement est dû à Cartier, qui prouve de plus que, si l'on désigne par $C$ l'opération de Cartier sur $H^0(Y,\Omega^1_Y)$, alors les différentielles localement logarithmiques forment le noyau de $C - \mathrm{Id}$.

}{F7II-XI-014}Supposons par l'absurde que la classe de $\underline{O}(1)$ dans $\mathrm{NS}(Y)$ soit de la forme $nu$, avec $n > 1$. Si $\ell$ divise $n$ et est premier à la caractéristique,
\SGApage{II}{51}
on a encore, en cohomologie $\ell$-adique

\begin{equation*}
\eta_\ell = n\ c_1(u)
\end{equation*}
ce qui contredit 1.6 (iv).

Enfin, si la caractéristique $p$ de $k$ est différente de $0$, et si $p$ divise $n$, on a encore, en cohomologie de Hodge

\begin{equation*}
\eta = n\ c_1(u) = 0 \qquad ,
\end{equation*}
ce qui contredit 1.5 (iii)$_d$ et achève la démonstration.


\SGAsectiontoc{2. Résultats numériques}{2}
2. \underline{Résultats numériques}

\SGAvar{2.1. Soient $\underline{a}=(a_i)_{1\leq i\leq d}$ une suite d'entiers $\geq 1$, et désignons par $V_n(\underline{a})$ une intersection complète lisse de dimension $n$ et de multidegré $\underline{a}$ dans un espace projectif $\mathbb{P}^{n+d}$ sur un corps $k$. On posera

}{2.1. Soit $\underline{a}=(a_i)_{1\leq i\leq d}$ une suite d'entiers $\geq 1$, et désignons par $V_n(\underline{a})$ une intersection complète lisse de dimension $n$ et de multidegré $\underline{a}$ dans un espace projectif $\mathbb{P}^{n+d}$ sur un corps $k$. On posera

}{F7II-XI-015}\begin{equation*}
h^{pq}\ (\underline{a}) = \dim_k H^q(V_{p+q}(\underline{a}),\ \Omega^p)
\end{equation*}
et

\begin{equation*}
h_o^{pq}(\underline{a}) = h^{pq}(\underline{a}) - \delta_{p,q} \qquad .
\end{equation*}
Les $h_o^{pq}(\underline{a})$ sont donc les nombres de Hodge de la cohomologie primitive de dimension moitié d'une intersection complète de multidegré $\underline{a}$. Ils ne dépendent que de $\underline{a}$, $p$ et $q$ (1.5). On posera

\begin{equation}\tag{2.1.1}
H(\underline{a}) = \sum_{p,q\geq 0} h_o^{p,q}\ y^p z^q \quad \in \mathbb{Z}[[y,z]] \qquad .
\end{equation}
Pour $d = 1$ et $\underline{a} = (a)$, on pose $h^{pq}(a) = h^{pq}(\underline{a})$ et $H(a) = H(\underline{a})$. Pour $d = 0$, on pose $\sup(a_i) = 1$.

La formule suivante est le cas particulier $k = 0$ de la formule (2) p.160 de Hirzebruch [3] 


\SGApage{II}{52}

\SGAvar{\begin{equation}\tag{2.2}
\sum_{n=0}^{\infty}\ \sum\ \chi(V_n(\underline{a}),\ \Omega^p)\ y^p\ z^{n+d}\ =\ \frac{1}{(1+yz)(1-z)}\ \prod_{1}^{d}\ \frac{(1+yz)^{a_i}-(1-z)^{a_i}}{(i+yz)^{a_i}+y(1-z)^{a_i}}
\end{equation}
Nous allons transformer cette formule en la suivante :

}{\begin{equation}\tag{2.2}
\sum_{n=0}^{\infty}\ \sum\ \chi(V_n(\underline{a}),\ \Omega^p)\ y^p\ z^{n+d}\ =\ \frac{1}{(1+yz)(1-z)}\ \prod_{1}^{d}\ \frac{(1+yz)^{a_i}-(1-z)^{a_i}}{(1+yz)^{a_i}+y(1-z)^{a_i}}
\end{equation}
Nous allons transformer cette formule en la suivante :

}{F7II-XI-016}\underline{Théorème} 2.3. (Hirzebruch). \underline{Avec les notations} 2.1, \underline{on a}

\begin{equation*}
H(\underline{a}) = \frac{1}{(1+y)(1+z)}\left[\ \prod_{1\leq i\leq d}\ \frac{(1+y)^{a_i}\ -\ (1+z)^{a_i}}{-(1+y)^{a_i}z\ +\ (1+z)^{a_i}y}\ \ -1\ \right] \qquad .
\end{equation*}
En vertu de 1.5, on a

\begin{equation*}
\chi(V_n(\underline{a}),\Omega^p) = (-1)^p + (-1)^{n-p}\ h_o^{p,n-p}(\underline{a})
\end{equation*}
de sorte que le premier membre $I$ de (2.2) se réécrit

\SGAvar{\begin{equation*}
\begin{aligned}
I &= \sum_{n=0}^{\infty}\ \sum_{p+q=n} ((-1)^p+(-1)^q\ h_o^{p,q}(\underline{a}))\ y^p z^{p+q+d}\\
&= \sum_{p,q\geq 0} (-1)^p(yz)^p z.^q z^d\ +\ \sum_{p,q\geq 0} (-1)^q\ h_o^{p,q}(a)\ (yz)^p z.^q\ z^d\\
&= z^d\left[\frac{1}{(1+yz)(1-z)}\ +\ H(\underline{a})(yz,-z)\right] \qquad .
\end{aligned}
\end{equation*}
Le deuxième membre $II$ de (2.2) se réécrit

}{\begin{equation*}
\begin{aligned}
I &= \sum_{n=0}^{\infty}\ \sum_{p+q=n} ((-1)^p+(-1)^q\ h_o^{p,q}(\underline{a}))\ y^p z^{p+q+d}\\
&= \sum_{p,q\geq 0} (-1)^p(yz)^p z^q z^d\ +\ \sum_{p,q\geq 0} (-1)^q\ h_o^{p,q}(a)\ (yz)^p z^q\ z^d\\
&= z^d\left[\frac{1}{(1+yz)(1-z)}\ +\ H(\underline{a})(yz,-z)\right] \qquad .
\end{aligned}
\end{equation*}
Le deuxième membre $II$ de (2.2) se réécrit

}{F7II-XI-017}\begin{equation*}
II = \frac{z^d}{(1+yz)(1-z)}\ \prod_{1}^{d}\ \frac{(1+yz)^{a_i}\ -\ (1-z)^{a_i}}{(1+yz)^{a_i}z+\ yz(1-z)^{a_i}}
\end{equation*}
de sorte que, simplifiant par $z^d$ dans $I = II$, on trouve que

\begin{equation*}
H(\underline{a})(yz,-z) = \frac{1}{(1+yz)(1-z)}\left[\ \prod_{1}^{d}\ \frac{(1+yz)^{a_i}\ -\ (1-z)^{a_i}}{(1+yz)^{a_i}z\ +\ (1-z)^{a_i}yz}\ \ -1\right]
\end{equation*}
et 2.3 se déduit de cette formule par changement de variable.


\SGApage{II}{53}

\underline{Corollaire} 2.4. \underline{Avec les notations} 2.1,

(i) \quad \underline{pour} $d = 1$, \underline{on a}

\SGAvar{\begin{equation*}
\begin{aligned}
H(a) \;&=\; -\ \frac{(1+y)^{a-1} - (1+z)^{a-1}}{(1+y)^{a}z \,-\, (1+z)^{a}y}\\[2ex]
&=\; \frac{\displaystyle\sum_{i,j\geq 0}\binom{a-1}{i+j+1}\ y^{i}z^{j}}{1 \,-\, \displaystyle\sum_{i,j\geq 1}\binom{a}{i+j}y^{i}\ z^{j}}\ .
\end{aligned}
\end{equation*}
(ii) \quad \underline{on a}, $|P|$ \underline{désignent le nombre d'élément d'une partie} $P$ \underline{de} $[1,d]$,

}{\begin{equation*}
\begin{aligned}
H(a) \;&=\; -\ \frac{(1+y)^{a-1} - (1+z)^{a-1}}{(1+y)^{a}z \,-\, (1+z)^{a}y}\\[2ex]
&=\; \frac{\displaystyle\sum_{i,j\geq 0}\binom{a-1}{i+j+1}\ y^{i}z^{j}}{1 \,-\, \displaystyle\sum_{i,j\geq 1}\binom{a}{i+j}y^{i}\ z^{j}}\ .
\end{aligned}
\end{equation*}
(ii) \quad \underline{on a}, $|P|$ \underline{désignant le nombre d'éléments d'une partie} $P$ \underline{de} $[1,d]$,

}{F7II-XI-018}\begin{multline*}
H(\underline{a}) \;=\; \sum_{\substack{P \subset [1,d] \\ P \neq \emptyset}} (1+y)^{|P|-1}\ (1+z)^{|P|-1}\ \prod_{i \in P} H(a_{i})\ .
\end{multline*}
D'après (2.3), on a

\SGAvar{\begin{equation*}
\begin{aligned}
H(a) \;&=\; \frac{-1}{(1+y)(1+z)}\left[\ \frac{(1+y)^{a} \,-\, (1+z)^{a}}{(1+z)^{a}\cdot z \,-\, (1+z)^{a}y} \;+\; 1\ \right]\\[2ex]
&=\; \frac{-1}{(1+y)(1+z)}\left[\ \frac{(1+y)^{a}(1+z)-(1+z)^{a}(1+y)}{(1+y)^{a}z \,-\, (1+z)^{a}y}\ \right]\\[2ex]
&=\; -\ \frac{(1+y)^{a-1} \,-\, (1+z)^{a-1}}{(1+y)^{a}z \,-\, (1+z)^{a}y}\ .
\end{aligned}
\end{equation*}
Soient $N$ et $D$ les numérateurs et dénominateurs de cette formule.

}{\begin{equation*}
\begin{aligned}
H(a) \;&=\; \frac{-1}{(1+y)(1+z)}\left[\ \frac{(1+y)^{a} \,-\, (1+z)^{a}}{(1+y)^{a}\cdot z \,-\, (1+z)^{a}y} \;+\; 1\ \right]\\[2ex]
&=\; \frac{-1}{(1+y)(1+z)}\left[\ \frac{(1+y)^{a}(1+z)-(1+z)^{a}(1+y)}{(1+y)^{a}z \,-\, (1+z)^{a}y}\ \right]\\[2ex]
&=\; -\ \frac{(1+y)^{a-1} \,-\, (1+z)^{a-1}}{(1+y)^{a}z \,-\, (1+z)^{a}y}\ .
\end{aligned}
\end{equation*}
Soient $N$ et $D$ les numérateurs et dénominateurs de cette formule.

}{F7II-XI-019}On a

\begin{equation*}
H(a) \;=\; -\,\frac{N}{D}\ , 
\end{equation*}

\SGApage{II}{54}

\SGAvar{\begin{equation*}
 N \;=\; \sum \binom{a-1}{k}\ (y^{k} \,-\, z^{k}) \;=\; (y-z) \sum_{i,j\geq 0}\binom{a-1}{i+j+1}\ y^{i}\ z^{j}
\end{equation*}
\begin{equation*}
\begin{aligned}
D \;&=\; \sum \binom{a}{k}\ (y^{k}z \,-\, yz^{k}) \;=\; (z-y) \;+\; yz \sum_{k\geq 1}\binom{a}{k}\!\left(y^{k-1} \,-\, z^{k-1}\right)\\[2ex]
&=\; (y-z)\ (-1 \;+\; yz \sum_{i,j\geq 0}\binom{a}{i+j+2}\ y^{i}z^{j}\\[2ex]
&=\; (y-z)\ (-1 \;+\; \sum_{i,j\geq 1}\binom{a}{i+j}y^{i}\ z^{j})
\end{aligned}
\end{equation*}
et l'assertion (i) en résulte.

}{\begin{equation*}
 N \;=\; \sum \binom{a-1}{k}\ (y^{k} \,-\, z^{k}) \;=\; (y-z) \sum_{i,j\geq 0}\binom{a-1}{i+j+1}\ y^{i}\ z^{j}
\end{equation*}
\begin{equation*}
\begin{aligned}
D \;&=\; \sum \binom{a}{k}\ (y^{k}z \,-\, yz^{k}) \;=\; (z-y) \;+\; yz \sum_{k\geq 1}\binom{a}{k}\!\left(y^{k-1} \,-\, z^{k-1}\right)\\[2ex]
&=\; (y-z)\ (-1 \;+\; yz \sum_{i,j\geq 0}\binom{a}{i+j+2}\ y^{i}z^{j})\\[2ex]
&=\; (y-z)\ (-1 \;+\; \sum_{i,j\geq 1}\binom{a}{i+j}y^{i}\ z^{j})
\end{aligned}
\end{equation*}
et l'assertion (i) en résulte.

}{F7II-XI-020}Inversons la formule de 2.3, pour $d = 1$. On trouve

\SGAvar{\begin{equation*}
-\ \frac{(1+y)^{a} \,-\, (1+z)^{a}}{(1+y)^{a}z \,-\, (1+z)^{a}y} \;=\; 1 \;+\; (1+x)(1+y)\ H(a)\ .
\end{equation*}
Si l'on substitue cette formule dans 2.3, on obtient

\begin{equation*}
\begin{aligned}
H(\underline{a}) \;&=\; \frac{1}{(1+y)(1+z)}\left[\ \prod_{1}^{d}\ (1+(1+x)(1+y)\ H(a_{i}))\qquad -\ 1\ \right]\\[2ex]
&=\;\frac{1}{(1+y)(1+z)}\left[\ \sum_{\substack{P \subset [i,d] \\ P \neq \emptyset}}\ (1+x)^{|P|}(1+y)^{|P|}\ \prod_{i \in P} H(a_{i})\ \right]
\end{aligned}
\end{equation*}
d'où résulte (ii).

}{\begin{equation*}
-\ \frac{(1+y)^{a} \,-\, (1+z)^{a}}{(1+y)^{a}z \,-\, (1+z)^{a}y} \;=\; 1 \;+\; (1+y)(1+z)\ H(a)\ .
\end{equation*}
Si l'on substitue cette formule dans 2.3, on obtient

\begin{equation*}
\begin{aligned}
H(\underline{a}) \;&=\; \frac{1}{(1+y)(1+z)}\left[\ \prod_{1}^{d}\ (1+(1+y)(1+z)\ H(a_{i}))\qquad -\ 1\ \right]\\[2ex]
&=\;\frac{1}{(1+y)(1+z)}\left[\ \sum_{\substack{P \subset [1,d] \\ P \neq \emptyset}}\ (1+y)^{|P|}(1+z)^{|P|}\ \prod_{i \in P} H(a_{i})\ \right]
\end{aligned}
\end{equation*}
d'où résulte (ii).

}{F7II-XI-021}\underline{Théorème} 2.5. \underline{Avec les notations} 2.1

(i) \quad \underline{Si} $p+q = p'+q'$ \underline{et si} $p \leq p' \leq q' \leq q$, \underline{alors}

\begin{equation*}
h_{o}^{pq}\ (\underline{a}) \;\leq\; h_{o}^{p'q'}(\underline{a})
\end{equation*}
(ii) \quad \underline{Si} $p \leq q$, \underline{pour que} $h_{o}^{pq}\ (\underline{a}) \neq 0$, \underline{il faut et il suffit que} 


\SGApage{II}{55}

\begin{equation*}
 p \;>\; \left[\frac{p+q+d \;-\; \sum a_i}{\sup(a_i)}\right] \qquad .
\end{equation*}
\underline{où} [ ] \underline{désigne la partie entière du quotient}.

Soit (*) la propriété suivante d'une série formelle $P = \sum a_{pq} \, y^p \, z^q$ :

(*)  les $a_{pq}$ sont positifs, $a_{pq} = a_{qp}$, et si $p+q = p'+q'$, et $p \leqslant p' \leqslant q' \leqslant q$, alors $a_{pq} \leqslant a_{p'q'}$ .

L'assertion (i) résulte aussitôt des formules (2.4), de l'identité $(1-x)^{-1} = \sum_{n \geqslant 0} x^n$ , et du lemme suivant, dont la démonstration est laissée au lecteur :

\underline{Lemme} 2.5.1. \underline{L'ensemble des séries formelles vérifiant} (*) \underline{est stable par addition et multiplication}.

On prouvera (ii) tout d'abord dans le cas particulier des hypersurfaces ($d = 1$). La formule à démontrer se réécrit alors, pour $p \leqslant q$,

\begin{equation}\tag{2.5.2}
h_o^{pq}(a) \neq 0 \quad \Longleftrightarrow \quad p \geqslant \left[\frac{p+q+1}{a}\right]
\end{equation}
\begin{equation*}
(\; \Longleftrightarrow \quad (p+1) \; a \geqslant p+q+2 \quad \mbox{trivialement}).
\end{equation*}
De la deuxième formule 2.4 (i), on déduit que $h_o^{pq}$ est non nul si et seulement si il existe des couples $(e,f)$ et $(c_i,d_i)$ tels que

\begin{equation}\tag{2.5.3}
\left\{
\begin{array}{l}
(e,f) \geqslant (0,0) \quad \mbox{et } e+f \leqslant a-2 \\[3mm]
(c_i,d_i) \geqslant (1,1) \mbox{ et } c_i+d_i \leqslant a \\[3mm]
(p,q) = (e,f) + \sum (c_i,d_i) \qquad .
\end{array}
\right.
\end{equation}

\SGApage{II}{56}

Pour $p+q$ donné, le plus petit $p$ possible s'obtient pour $e = 0$ et $c_i = 1$, et est le nombre d'indices $i$. Ce $p$ est donc la borne inférieure des $p$ pour lesquels on a

\begin{equation*}
p \;\geqslant\; \frac{(p+q)-(a-2)}{a} \qquad , \; \mbox{i.e.} \quad (p+1)a \geqslant p+q+2 \qquad .
\end{equation*}
Passons au cas général. Si $\underline{a}'$ se déduit de $\underline{a}$ en supprimant les $a_i$ égaux à 1, alors $H(\underline{a}) = H(\underline{a}')$. Vu la forme de la formule 2.5 (ii), on peut supposer que pour tout $i$, $a_i > 1$. Si $d = 0$, alors $h_o^{pq} = 0$ et l'assertion est correcte; on ne restreint pas la généralité en supposant que $d > 0$ et que $a_d = \sup(a_i)$.

\SGAvar{En vertu de 2.4 (ii), et (2.5.2), pour que $h_o^{pq}(\underline{a}) \neq 0$, il faut et il suffit qu'il existe une partie non vide $P$ de $[1,d]$, un couple $(e,f)$ et des couples $(c_i,d_i)_{i \leqslant i \leqslant d}$ tels que

\begin{equation}\tag{2.5.4}
\left\{
\begin{array}{l}
(|P|-1, \; |P|-1) \geqslant (e,f) \geqslant (0,0) \; , \\[3mm]
c_i,d_i \geqslant \left[{\displaystyle \frac{c_i+d_i+1}{a_i}}\right] \quad ; \quad (c_i,d_i) \geqslant (0,0) \quad , \\[3mm]
(p,q) = (e,f) + \sum (c_i,d_i) \qquad .
\end{array}
\right.
\end{equation}
Pour $n = p+q$ fixé, cherchons la plus petite valeur de $p$ possible. Si $n_i = c_i + d_i$, pour un tel $p$, les conditions (2.5.4) se réécrivent, en terme des $c_i$, $n_i$, $e$ et $f$ 


}{En vertu de 2.4 (ii), et (2.5.2), pour que $h_o^{pq}(\underline{a}) \neq 0$, il faut et il suffit qu'il existe une partie non vide $P$ de $[1,d]$, un couple $(e,f)$ et des couples $(c_i,d_i)_{i\in P}$ tels que

\begin{equation}\tag{2.5.4}
\left\{
\begin{array}{l}
(|P|-1, \; |P|-1) \geqslant (e,f) \geqslant (0,0) \; , \\[3mm]
c_i,d_i \geqslant \left[{\displaystyle \frac{c_i+d_i+1}{a_i}}\right] \quad ; \quad (c_i,d_i) \geqslant (0,0) \quad , \\[3mm]
(p,q) = (e,f) + \sum_{i\in P} (c_i,d_i) \qquad .
\end{array}
\right.
\end{equation}
Pour $n = p+q$ fixé, cherchons la plus petite valeur de $p$ possible. Si $n_i = c_i + d_i$, pour un tel $p$, les conditions (2.5.4) se réécrivent, en termes des $c_i$, $n_i$, $e$ et $f$ 


}{F7II-XI-022}\SGApage{II}{57}

\SGAvar{\begin{equation}\tag{2.5.5}
\left\{
\begin{array}{l}
a_i = 2 \Longrightarrow n_i \mbox{ pair} \qquad \mbox{(sinon (2.5.4) est impossible)} \\[3mm]
c_i \geq \left[ \displaystyle\frac{n_i+1}{a_i} \right] \quad ; \quad c_i \leq n_i \quad , \\[3mm]
(|P|-1, \; |P|-1) \geq (e,f) \geq (0,0) \quad , \\[3mm]
n = e+f+ \sum n_i \quad , \\[3mm]
p = e + \sum c_i \quad .
\end{array}
\right.
\end{equation}
Cette description montre tout d'abord que le plus petit $p$ est obtenu pour $P = [1,d]$, cas auquel on se limitera. Si on a

\begin{equation}\tag{2.5.6}
\sum (a_i-2) + (d-1) \geq n,
\end{equation}
alors il existe une solution avec $p = c_i = e = 0$. Sinon, si $a_d \neq 2$, le mieuxqu'on puisse faire est de prendre $n_i = a_i - 2$ pour $i \neq d$, de sorte que $c_i = 0$ pour $i \neq d$, $e = 0$, $f = d-1$ et $n_d = n - (d-1) - \sum\limits_{i \neq d} (a_i-2)$. On a alors

}{\begin{equation}\tag{2.5.5}
\left\{
\begin{array}{l}
a_i = 2 \Longrightarrow n_i \mbox{ pair} \qquad \mbox{(sinon (2.5.4) est impossible)} \\[3mm]
c_i \geq \left[ \displaystyle\frac{n_i+1}{a_i} \right] \quad ; \quad c_i \leq n_i \quad , \\[3mm]
(|P|-1, \; |P|-1) \geq (e,f) \geq (0,0) \quad , \\[3mm]
n = e+f+ \sum_{i\in P} n_i \quad , \\[3mm]
p = e + \sum_{i\in P} c_i \quad .
\end{array}
\right.
\end{equation}
Cette description montre tout d'abord que le plus petit $p$ est obtenu pour $P = [1,d]$, cas auquel on se limitera. Si on a

\begin{equation}\tag{2.5.6}
\sum (a_i-2) + (d-1) \geq n,
\end{equation}
alors il existe une solution avec $p = c_i = e = 0$. Sinon, si $a_d \neq 2$, le mieux qu'on puisse faire est de prendre $n_i = a_i - 2$ pour $i \neq d$, de sorte que $c_i = 0$ pour $i \neq d$, $e = 0$, $f = d-1$ et $n_d = n - (d-1) - \sum\limits_{i \neq d} (a_i-2)$. On a alors

}{F7II-XI-023}\begin{equation}\tag{2.5.7}
p = \left[ \displaystyle\frac{n-(d-1) - \sum\limits_{i \neq d} (a_i-2) + 1}{a_d} \right] \qquad ,
\end{equation}
et donc, pour $p \leq q$, on a $h^{pq}_{o}(\underline{a}) \neq 0$ si et seulement si

\begin{equation}\tag{2.5.8}
p \geq \left[ \displaystyle\frac{n- \sum\limits_{i \neq d} (a_i-1) + 1}{a_d} \right] \qquad .
\end{equation}
La formule (2.5.8) se réduit à (2.5.2) si $d = 1$ et est trivialement vraie si (2.5.6) est vérifié. Si $d \geq 2$, si (2.5.6) n'est pas vérifié, et si tous les $a_i$ sont égaux à $2$, le mieux qu'on puisse faire est encore de
\SGApage{II}{58}
prendre $n_i = 0$ $(i \neq d)$, $e = 0$, $f = d-1$ ou $d-2$, congru à $n$ modulo $2$, et $n_d = n-f$. Si $n \equiv d-1$ $(2)$, on prouve (2.5.7) comme plus haut. Sinon, on note que

\begin{equation*}
\left[ \displaystyle\frac{n - (d-2) + 1}{2} \right] = \left[ \displaystyle\frac{n - (d-1) + 1}{2} \right]
\end{equation*}
de sorte que (2.5.7), et par là (2.5.8), est valable même si $a_d = 2$.

Reste à vérifier que (2.5.8) équivaut à la formule 2.5 (ii). En effet, on a

\begin{equation*}
\left[ \displaystyle\frac{n - \sum\limits_{i \neq d} (a_i-1) + 1}{a_d} \right] = \left[ \displaystyle\frac{n + d - \sum\limits_{i \neq d} a_i}{a_d} \right] = \left[ \displaystyle\frac{p + q + d - \sum a_i}{a_d} \right] + 1 \quad .
\end{equation*}
2.6. \quad Si $X$ est une intersection complète de dimension $n$ dans $\mathbb{P}^r$, il résulte aussitôt de (2.5.8) que la partie primitive de $H^n(X)$ n'est nulle que si $X$ est une quadrique de dimension impaire.

\underline{Définition} 2.7. \underline{Si} $X$ \underline{est une variété algébrique propre et lisse sur} $\mathbb{C}$, \underline{le niveau de Hodge de} $H^n(X)$ \underline{est la borne supérieure des nombres} $q-p$ \underline{pour} $p+q = n$ \underline{et} $h^{pq}(X) \neq 0$.

\SGAvar{Le niveau de Hodge $\ell$ est donc congru à $n$ modulo $2$, et est au plus égal à $\inf(n,\dim(X) - n)$. Il est $\geq 0$, sauf si $H^n(X) = 0$, auquel cas il est égal à $- \infty$.

}{Le niveau de Hodge $\ell$ est donc congru à $n$ modulo $2$, et est au plus égal à $\inf(n,2\dim(X) - n)$. Il est $\geq 0$, sauf si $H^n(X) = 0$, auquel cas il est égal à $- \infty$.

}{F7II-XI-024}\underline{Corollaire} 2.8. \underline{Reprenons les notations} 2.1, \underline{en faisant} $k = \mathbb{C}$, \underline{et en supposant que} $a_d = \sup (a_i)$. \underline{Soit} $\ell$ \underline{un entier} $0 \leq \ell < n$, \underline{congru à} $n$ \underline{modulo} $2$. \underline{Pour que le niveau de Hodge de} $H^n(V_n(\underline{a}))$ \underline{soit} $\leq \ell$ , \underline{il faut et il suffit que} 


\SGApage{II}{59}

\begin{equation*}
 \sum_{i \neq d} \, (a_i - 1) \;\leq\; (\ell + 1) \, - \, (a_d - 2) \, \frac{n - \ell}{2} \qquad .
\end{equation*}
En vertu de (2.5.8), si le niveau de Hodge est $\geq 0$, il est la borne supérieure de 0 et de

\begin{equation*}
n \, - \, 2 \left[ \frac{n \, - \, \sum_{i \neq d} \, (a_i - 1) \, + \, 1}{a_d} \right]
\end{equation*}
et l'inéquation

\begin{equation*}
n \, - \, 2 \left[ \frac{n \, - \, \sum_{i \neq d} \, (a_i - 1) \, + \, 1}{a_d} \right] \;\leq\; \ell
\end{equation*}
peut se réécrire

\begin{equation*}
\frac{n \, - \, \ell}{2} \;\leq\; \left[ \frac{n - \, \sum_{i \neq d} \, (a_i - 1) \; + \; 1}{a_d} \right] \qquad ,
\end{equation*}
ou encore

\SGAvar{\begin{equation*}
\sum_{i \neq d} \, (a_i - 1) \;\leq\; 1 \, + \, n \, - \, a_d \; . \; \frac{n \, - \, \ell}{2} \;=\; \ell \, + \, 1 \, - \, (a_d - 2) \; \frac{n \, - \, \ell}{2} \qquad .
\end{equation*}
2.9. \quad A l'aide de la formule 2.8, il est facile d'énumérer, pour $\ell$ petit, les multidegrés des intersections complètes dans $\mathbb{P}^r$, de dimension $n > 1$, dont le $H^n$ est de niveau $\ell$. La table donnée dans la première version de cet exposé était néanmoins incorrecte, et celle donnée ci-dessous est extraite de celle plus complète donnée par M. Rapoport dans son article : Complément à l'article de P. Deligne "La conjecture de Weil pour les surfaces K 3", Inventiones Math. \underline{15} 227-236 (1972). Pour $\ell \leq 1$, les intersections complètes $V_n(\underline{a})$ de dimension $n$ et de multidegré $\underline{a} = (a_1,\ldots,a_d)$, avec $n > \ell$ , $d \geq 1$ et $a_i \geq 2$, de niveau de Hodge $\ell$ , sont les suivantes :


}{\begin{equation*}
\sum_{i \neq d} \, (a_i - 1) \;\leq\; 1 \, + \, n \, - \, a_d \; . \; \frac{n \, - \, \ell}{2} \;=\; \ell \, + \, 1 \, - \, (a_d - 2) \; \frac{n \, - \, \ell}{2} \qquad .
\end{equation*}
2.9. \quad À l'aide de la formule 2.8, il est facile d'énumérer, pour $\ell$ petit, les multidegrés des intersections complètes dans $\mathbb{P}^r$, de dimension $n > 1$, dont le $H^n$ est de niveau $\ell$. La table donnée dans la première version de cet exposé était néanmoins incorrecte, et celle donnée ci-dessous est extraite de celle plus complète donnée par M. Rapoport dans son article : Complément à l'article de P. Deligne "La conjecture de Weil pour les surfaces K 3", Inventiones Math. \underline{15} 227-236 (1972). Pour $\ell \leq 1$, les intersections complètes $V_n(\underline{a})$ de dimension $n$ et de multidegré $\underline{a} = (a_1,\ldots,a_d)$, avec $n > \ell$ , $d \geq 1$ et $a_i \geq 2$, de niveau de Hodge $\ell$ , sont les suivantes :


}{F7II-XI-025}\SGApage{II}{60}

\begin{equation*}
\begin{array}{ll}
\ell \;=\; -\infty & V_{2k+1}(2) \qquad\qquad\qquad (k \geq 0) \\[2ex]
\ell \;=\; 0 & V_{2k}(2) \; (k \geq 1), \; V_{2k}(2,2) \; (k \geq 1), \; V_2(3) \\[2ex]
\ell \;=\; 1 & V_{2k+1}(2,2) \; (k \geq 1), \; V_{2k+1}(2,2,2) \; (k \geq 1) \\
& V_3(3), \; V_5(3), \; V_3(3,2), \; V_3(4) \; .
\end{array}
\end{equation*}
2.10. \quad Pour $\underline{a}$ petit, les fonctions génératrices des $h_o^{pq}(\underline{a})$ s'écrivent, d'après 2.4

\begin{equation*}
\begin{array}{rcl}
H(2) & = & \dfrac{1}{1-yz} \\[3ex]
H(3) & = & \dfrac{2 \, + \, y \, + \, z}{1 \, - \, 3yz \, - \, y^2 z \, - \, yz^2} \\[3ex]
H(2,2) & = & \dfrac{y \, + \, z}{(1-yz)^2} \; + \; \dfrac{2}{(1-yz)^2} + \; \dfrac{1}{1-yz}
\end{array}
\end{equation*}
de sorte que

\begin{equation*}
\begin{array}{rcll}
h_o^{pp}(2) & = & 1 & , \\[2ex]
h_o^{11}(3) & = & 6 & , \\[2ex]
h_o^{pp}(2,2) & = & 2p \, + \, 3 & .
\end{array}
\end{equation*}

\SGApage{II}{61}


\underline{BIBLIOGRAPHIE}

[1] Y. AKIZUKI and S. NAKANO. Note on Kodaira - Spencer's proof of Lefschetz theorems - Proc. Jap. Acad. 30, 266-272, (1954).

[2] R. HARTSHORNE. Residues and duality. Lecture notes 20 Springer 66.

[3] F. HIRZEBRUCH. Topological methods in algebraic geometry 3é ed. Springer-Verlag 1966.

[4] K. KODAIRA. On a differential geometric method in the theory of analytic stacks. Proc. Nat. Acad. Sci. USA 39, 1268-1273, (1953).
\selectlanguage{french}

\ifnum\SGAform=1
\clearpage\section*{Apparat critique}\small
Les lectures ci-dessous sont données dans le TeX exact des témoins, afin de préserver les distinctions de notation. Les indices décimaux renvoient aux différences individuelles du passage marqué.

\par\medskip\noindent\hypertarget{app-F7II-XI-001.1}{}\textbf{F7II-XI-001.1}\quad p. 41; misprint\_pass.\par
\noindent Imprimé : {\ttfamily\footnotesize ,\allowbreak{}}\par
\noindent Corrigé : {\ttfamily\footnotesize .\allowbreak{}}\par
\noindent Source : 7II:57. \nolinkurl{fixes/FIXLOG_7II_CURRENT.jsonl#512}; \nolinkurl{fixes/MISPRINTS_7II.json#57}\par

\par\medskip\noindent\hypertarget{app-F7II-XI-002.1}{}\textbf{F7II-XI-002.1}\quad p. 42; misprint\_pass.\par
\noindent Imprimé : {\ttfamily\footnotesize Kozul\allowbreak{}}\par
\noindent Corrigé : {\ttfamily\footnotesize Koszul\allowbreak{}}\par
\noindent Source : 7II:59. \nolinkurl{fixes/FIXLOG_7II_CURRENT.jsonl#513}; \nolinkurl{fixes/MISPRINTS_7II.json#59}\par

\par\medskip\noindent\hypertarget{app-F7II-XI-003.1}{}\textbf{F7II-XI-003.1}\quad p. 42; misprint\_pass.\par
\noindent Imprimé : {\ttfamily\footnotesize suite\allowbreak{}}\par
\noindent Corrigé : {\ttfamily\footnotesize suites\allowbreak{}}\par
\noindent Source : 7II:60. \nolinkurl{fixes/FIXLOG_7II_CURRENT.jsonl#514}\par

\par\medskip\noindent\hypertarget{app-F7II-XI-004.1}{}\textbf{F7II-XI-004.1}\quad p. 43; inherited corrected-layer difference.\par
\noindent Imprimé : \emph{(vide)}\par
\noindent Corrigé : {\ttfamily\footnotesize )\allowbreak{}}\par
\noindent Source : ADOPT-F7II-XI-004.1. \nolinkurl{fr/sga7ii/diplomatic/expose_XI_body.tex}; \nolinkurl{fr/sga7ii/current/expose_XI_body.tex}\par

\par\medskip\noindent\hypertarget{app-F7II-XI-005.1}{}\textbf{F7II-XI-005.1}\quad p. 43; misprint\_pass.\par
\noindent Imprimé : {\ttfamily\footnotesize j\allowbreak{}}\par
\noindent Corrigé : {\ttfamily\footnotesize i\allowbreak{}}\par
\noindent Source : 7II:62. \nolinkurl{fixes/FIXLOG_7II_CURRENT.jsonl#515}; \nolinkurl{fixes/MISPRINTS_7II.json#62}\par

\par\medskip\noindent\hypertarget{app-F7II-XI-006.1}{}\textbf{F7II-XI-006.1}\quad p. 43; misprint\_pass.\par
\noindent Imprimé : {\ttfamily\footnotesize a\allowbreak{}}\par
\noindent Corrigé : {\ttfamily\footnotesize 0\allowbreak{}}\par
\noindent Source : 7II:63. \nolinkurl{fixes/FIXLOG_7II_CURRENT.jsonl#516}; \nolinkurl{fixes/MISPRINTS_7II.json#63}\par

\par\medskip\noindent\hypertarget{app-F7II-XI-007.1}{}\textbf{F7II-XI-007.1}\quad p. 44; misprint\_pass.\par
\noindent Imprimé : {\ttfamily\footnotesize S\allowbreak{}}\par
\noindent Corrigé : {\ttfamily\footnotesize s\allowbreak{}}\par
\noindent Source : 7II:64. \nolinkurl{fixes/FIXLOG_7II_CURRENT.jsonl#517}; \nolinkurl{fixes/MISPRINTS_7II.json#64}\par

\par\medskip\noindent\hypertarget{app-F7II-XI-007.2}{}\textbf{F7II-XI-007.2}\quad p. 44; misprint\_pass.\par
\noindent Imprimé : {\ttfamily\footnotesize régulier\allowbreak{}}\par
\noindent Corrigé : {\ttfamily\footnotesize régulière\allowbreak{}}\par
\noindent Source : 7II:65. \nolinkurl{fixes/FIXLOG_7II_CURRENT.jsonl#518}; \nolinkurl{fixes/MISPRINTS_7II.json#65}\par

\par\medskip\noindent\hypertarget{app-F7II-XI-008.1}{}\textbf{F7II-XI-008.1}\quad p. 44; printed misprint.\par
\noindent Imprimé : {\ttfamily\footnotesize H\allowbreak{}}\par
\noindent Corrigé : {\ttfamily\footnotesize X\allowbreak{}}\par
\noindent Source : 7II:66. \nolinkurl{fixes/MISPRINTS_7II.json#66}\par

\par\medskip\noindent\hypertarget{app-F7II-XI-008.2}{}\textbf{F7II-XI-008.2}\quad p. 44; printed misprint.\par
\noindent Imprimé : {\ttfamily\footnotesize X\allowbreak{}}\par
\noindent Corrigé : {\ttfamily\footnotesize H\allowbreak{}}\par
\noindent Source : 7II:66. \nolinkurl{fixes/MISPRINTS_7II.json#66}\par

\par\medskip\noindent\hypertarget{app-F7II-XI-008.3}{}\textbf{F7II-XI-008.3}\quad p. 44; printed misprint.\par
\noindent Imprimé : {\ttfamily\footnotesize H\allowbreak{}}\par
\noindent Corrigé : {\ttfamily\footnotesize X\allowbreak{}}\par
\noindent Source : 7II:67. \nolinkurl{fixes/MISPRINTS_7II.json#67}\par

\par\medskip\noindent\hypertarget{app-F7II-XI-008.4}{}\textbf{F7II-XI-008.4}\quad p. 44; printed misprint.\par
\noindent Imprimé : {\ttfamily\footnotesize X\allowbreak{}}\par
\noindent Corrigé : {\ttfamily\footnotesize H\allowbreak{}}\par
\noindent Source : 7II:67. \nolinkurl{fixes/MISPRINTS_7II.json#67}\par

\par\medskip\noindent\hypertarget{app-F7II-XI-008.5}{}\textbf{F7II-XI-008.5}\quad p. 44; misprint\_pass.\par
\noindent Imprimé : {\ttfamily\footnotesize Y\allowbreak{}}\par
\noindent Corrigé : {\ttfamily\footnotesize H\allowbreak{}}\par
\noindent Source : 7II:68. \nolinkurl{fixes/FIXLOG_7II_CURRENT.jsonl#519}\par

\par\medskip\noindent\hypertarget{app-F7II-XI-008.6}{}\textbf{F7II-XI-008.6}\quad p. 44; misprint\_pass.\par
\noindent Imprimé : {\ttfamily\footnotesize Y\allowbreak{}}\par
\noindent Corrigé : {\ttfamily\footnotesize H\allowbreak{}}\par
\noindent Source : 7II:68. \nolinkurl{fixes/FIXLOG_7II_CURRENT.jsonl#519}\par

\par\medskip\noindent\hypertarget{app-F7II-XI-008.7}{}\textbf{F7II-XI-008.7}\quad p. 44; misprint\_pass.\par
\noindent Imprimé : {\ttfamily\footnotesize Y\allowbreak{}}\par
\noindent Corrigé : {\ttfamily\footnotesize H\allowbreak{}}\par
\noindent Source : 7II:68. \nolinkurl{fixes/FIXLOG_7II_CURRENT.jsonl#519}\par

\par\medskip\noindent\hypertarget{app-F7II-XI-008.8}{}\textbf{F7II-XI-008.8}\quad p. 44; misprint\_pass.\par
\noindent Imprimé : {\ttfamily\footnotesize Y\allowbreak{}}\par
\noindent Corrigé : {\ttfamily\footnotesize H\allowbreak{}}\par
\noindent Source : 7II:68. \nolinkurl{fixes/FIXLOG_7II_CURRENT.jsonl#519}\par

\par\medskip\noindent\hypertarget{app-F7II-XI-008.9}{}\textbf{F7II-XI-008.9}\quad p. 45; misprint\_pass.\par
\noindent Imprimé : {\ttfamily\footnotesize Y\allowbreak{}}\par
\noindent Corrigé : {\ttfamily\footnotesize H\allowbreak{}}\par
\noindent Source : 7II:69. \nolinkurl{fixes/FIXLOG_7II_CURRENT.jsonl#520}\par

\par\medskip\noindent\hypertarget{app-F7II-XI-008.10}{}\textbf{F7II-XI-008.10}\quad p. 45; misprint\_pass.\par
\noindent Imprimé : {\ttfamily\footnotesize Y\allowbreak{}}\par
\noindent Corrigé : {\ttfamily\footnotesize H\allowbreak{}}\par
\noindent Source : 7II:69. \nolinkurl{fixes/FIXLOG_7II_CURRENT.jsonl#520}\par

\par\medskip\noindent\hypertarget{app-F7II-XI-008.11}{}\textbf{F7II-XI-008.11}\quad p. 45; misprint\_pass.\par
\noindent Imprimé : {\ttfamily\footnotesize Y\allowbreak{}}\par
\noindent Corrigé : {\ttfamily\footnotesize H\allowbreak{}}\par
\noindent Source : 7II:69. \nolinkurl{fixes/FIXLOG_7II_CURRENT.jsonl#520}\par

\par\medskip\noindent\hypertarget{app-F7II-XI-008.12}{}\textbf{F7II-XI-008.12}\quad p. 45; misprint\_pass.\par
\noindent Imprimé : {\ttfamily\footnotesize Y\allowbreak{}}\par
\noindent Corrigé : {\ttfamily\footnotesize H\allowbreak{}}\par
\noindent Source : 7II:69. \nolinkurl{fixes/FIXLOG_7II_CURRENT.jsonl#520}\par

\par\medskip\noindent\hypertarget{app-F7II-XI-008.13}{}\textbf{F7II-XI-008.13}\quad p. 45; misprint\_pass.\par
\noindent Imprimé : {\ttfamily\footnotesize Y\allowbreak{}}\par
\noindent Corrigé : {\ttfamily\footnotesize H\allowbreak{}}\par
\noindent Source : 7II:69. \nolinkurl{fixes/FIXLOG_7II_CURRENT.jsonl#520}\par

\par\medskip\noindent\hypertarget{app-F7II-XI-008.14}{}\textbf{F7II-XI-008.14}\quad p. 45; misprint\_pass.\par
\noindent Imprimé : {\ttfamily\footnotesize Y\allowbreak{}}\par
\noindent Corrigé : {\ttfamily\footnotesize H\allowbreak{}}\par
\noindent Source : 7II:69. \nolinkurl{fixes/FIXLOG_7II_CURRENT.jsonl#520}\par

\par\medskip\noindent\hypertarget{app-F7II-XI-009.1}{}\textbf{F7II-XI-009.1}\quad p. 46; misprint\_pass.\par
\noindent Imprimé : {\ttfamily\footnotesize cohomologie\allowbreak{}}\par
\noindent Corrigé : {\ttfamily\footnotesize cohomologies\allowbreak{}}\par
\noindent Source : 7II:70. \nolinkurl{fixes/FIXLOG_7II_CURRENT.jsonl#521}; \nolinkurl{fixes/MISPRINTS_7II.json#70}\par

\par\medskip\noindent\hypertarget{app-F7II-XI-010.1}{}\textbf{F7II-XI-010.1}\quad p. 46; misprint\_pass.\par
\noindent Imprimé : {\ttfamily\footnotesize X\allowbreak{}}\par
\noindent Corrigé : {\ttfamily\footnotesize Y\allowbreak{}}\par
\noindent Source : 7II:71. \nolinkurl{fixes/FIXLOG_7II_CURRENT.jsonl#522}; \nolinkurl{fixes/MISPRINTS_7II.json#71}\par

\par\medskip\noindent\hypertarget{app-F7II-XI-011.1}{}\textbf{F7II-XI-011.1}\quad p. 47; misprint\_pass.\par
\noindent Imprimé : \emph{(vide)}\par
\noindent Corrigé : {\ttfamily\footnotesize )\allowbreak{}}\par
\noindent Source : 7II:72. \nolinkurl{fixes/FIXLOG_7II_CURRENT.jsonl#523}; \nolinkurl{fixes/MISPRINTS_7II.json#72}\par

\par\medskip\noindent\hypertarget{app-F7II-XI-011.2}{}\textbf{F7II-XI-011.2}\quad p. 47; misprint\_pass.\par
\noindent Imprimé : \emph{(vide)}\par
\noindent Corrigé : {\ttfamily\footnotesize +\allowbreak{} j\allowbreak{} }\par
\noindent Source : 7II:73. \nolinkurl{fixes/FIXLOG_7II_CURRENT.jsonl#524}; \nolinkurl{fixes/MISPRINTS_7II.json#73}\par

\par\medskip\noindent\hypertarget{app-F7II-XI-011.3}{}\textbf{F7II-XI-011.3}\quad p. 47; misprint\_pass.\par
\noindent Imprimé : {\ttfamily\footnotesize j\allowbreak{}}\par
\noindent Corrigé : {\ttfamily\footnotesize n\allowbreak{}}\par
\noindent Source : 7II:73. \nolinkurl{fixes/FIXLOG_7II_CURRENT.jsonl#524}; \nolinkurl{fixes/MISPRINTS_7II.json#73}\par

\par\medskip\noindent\hypertarget{app-F7II-XI-012.1}{}\textbf{F7II-XI-012.1}\quad p. 47; misprint\_pass.\par
\noindent Imprimé : \emph{(vide)}\par
\noindent Corrigé : {\ttfamily\footnotesize .\allowbreak{}}\par
\noindent Source : 7II:74. \nolinkurl{fixes/FIXLOG_7II_CURRENT.jsonl#525}\par

\par\medskip\noindent\hypertarget{app-F7II-XI-013.1}{}\textbf{F7II-XI-013.1}\quad p. 48; misprint\_pass.\par
\noindent Imprimé : {\ttfamily\footnotesize Y\allowbreak{}}\par
\noindent Corrigé : {\ttfamily\footnotesize \textbackslash{}mathbb\{P\}(\textbackslash{}underlin\allowbreak{}e\{E\})\allowbreak{}}\par
\noindent Source : 7II:75. \nolinkurl{fixes/FIXLOG_7II_CURRENT.jsonl#526}; \nolinkurl{fixes/MISPRINTS_7II.json#75}\par

\par\medskip\noindent\hypertarget{app-F7II-XI-014.1}{}\textbf{F7II-XI-014.1}\quad p. 50; printed misprint.\par
\noindent Imprimé : {\ttfamily\footnotesize logarithmique\allowbreak{}}\par
\noindent Corrigé : {\ttfamily\footnotesize logarithmiques\allowbreak{}}\par
\noindent Source : 7II:76. \nolinkurl{fixes/MISPRINTS_7II.json#76}\par

\par\medskip\noindent\hypertarget{app-F7II-XI-015.1}{}\textbf{F7II-XI-015.1}\quad p. 51; inherited corrected-layer difference.\par
\noindent Imprimé : {\ttfamily\footnotesize Soient\allowbreak{}}\par
\noindent Corrigé : {\ttfamily\footnotesize Soit\allowbreak{}}\par
\noindent Source : ADOPT-F7II-XI-015.1. \nolinkurl{fr/sga7ii/diplomatic/expose_XI_body.tex}; \nolinkurl{fr/sga7ii/current/expose_XI_body.tex}\par

\par\medskip\noindent\hypertarget{app-F7II-XI-016.1}{}\textbf{F7II-XI-016.1}\quad p. 52; misprint\_pass.\par
\noindent Imprimé : {\ttfamily\footnotesize i\allowbreak{}}\par
\noindent Corrigé : {\ttfamily\footnotesize 1\allowbreak{}}\par
\noindent Source : 7II:77. \nolinkurl{fixes/FIXLOG_7II_CURRENT.jsonl#527}\par

\par\medskip\noindent\hypertarget{app-F7II-XI-017.1}{}\textbf{F7II-XI-017.1}\quad p. 52; misprint\_in\_print.\par
\noindent Imprimé : {\ttfamily\footnotesize .\allowbreak{}}\par
\noindent Corrigé : \emph{(vide)}\par
\noindent Source : FIX-7II-D-0040. \nolinkurl{fixes/FIXLOG_7II.jsonl#40}\par

\par\medskip\noindent\hypertarget{app-F7II-XI-017.2}{}\textbf{F7II-XI-017.2}\quad p. 52; misprint\_in\_print.\par
\noindent Imprimé : {\ttfamily\footnotesize .\allowbreak{}}\par
\noindent Corrigé : \emph{(vide)}\par
\noindent Source : FIX-7II-D-0041. \nolinkurl{fixes/FIXLOG_7II.jsonl#41}\par

\par\medskip\noindent\hypertarget{app-F7II-XI-018.1}{}\textbf{F7II-XI-018.1}\quad p. 53; misprint\_pass.\par
\noindent Imprimé : {\ttfamily\footnotesize désignent\allowbreak{}}\par
\noindent Corrigé : {\ttfamily\footnotesize désignant\allowbreak{}}\par
\noindent Source : 7II:81. \nolinkurl{fixes/FIXLOG_7II_CURRENT.jsonl#528}\par

\par\medskip\noindent\hypertarget{app-F7II-XI-018.2}{}\textbf{F7II-XI-018.2}\quad p. 53; misprint\_pass.\par
\noindent Imprimé : {\ttfamily\footnotesize élément\allowbreak{}}\par
\noindent Corrigé : {\ttfamily\footnotesize éléments\allowbreak{}}\par
\noindent Source : 7II:81. \nolinkurl{fixes/FIXLOG_7II_CURRENT.jsonl#528}\par

\par\medskip\noindent\hypertarget{app-F7II-XI-019.1}{}\textbf{F7II-XI-019.1}\quad p. 53; misprint\_pass.\par
\noindent Imprimé : {\ttfamily\footnotesize z\allowbreak{}}\par
\noindent Corrigé : {\ttfamily\footnotesize y\allowbreak{}}\par
\noindent Source : 7II:83. \nolinkurl{fixes/FIXLOG_7II_CURRENT.jsonl#529}; \nolinkurl{fixes/MISPRINTS_7II.json#83}\par

\par\medskip\noindent\hypertarget{app-F7II-XI-020.1}{}\textbf{F7II-XI-020.1}\quad p. 54; misprint\_pass.\par
\noindent Imprimé : \emph{(vide)}\par
\noindent Corrigé : {\ttfamily\footnotesize )\allowbreak{}}\par
\noindent Source : 7II:84. \nolinkurl{fixes/FIXLOG_7II_CURRENT.jsonl#530}\par

\par\medskip\noindent\hypertarget{app-F7II-XI-021.1}{}\textbf{F7II-XI-021.1}\quad p. 54; misprint\_pass.\par
\noindent Imprimé : {\ttfamily\footnotesize x\allowbreak{}}\par
\noindent Corrigé : {\ttfamily\footnotesize y\allowbreak{}}\par
\noindent Source : 7II:85. \nolinkurl{fixes/FIXLOG_7II_CURRENT.jsonl#531}; \nolinkurl{fixes/MISPRINTS_7II.json#85}\par

\par\medskip\noindent\hypertarget{app-F7II-XI-021.2}{}\textbf{F7II-XI-021.2}\quad p. 54; misprint\_pass.\par
\noindent Imprimé : {\ttfamily\footnotesize y\allowbreak{}}\par
\noindent Corrigé : {\ttfamily\footnotesize z\allowbreak{}}\par
\noindent Source : 7II:85. \nolinkurl{fixes/FIXLOG_7II_CURRENT.jsonl#531}; \nolinkurl{fixes/MISPRINTS_7II.json#85}\par

\par\medskip\noindent\hypertarget{app-F7II-XI-021.3}{}\textbf{F7II-XI-021.3}\quad p. 54; misprint\_pass.\par
\noindent Imprimé : {\ttfamily\footnotesize x\allowbreak{}}\par
\noindent Corrigé : {\ttfamily\footnotesize y\allowbreak{}}\par
\noindent Source : 7II:86. \nolinkurl{fixes/FIXLOG_7II_CURRENT.jsonl#532}; \nolinkurl{fixes/MISPRINTS_7II.json#86}\par

\par\medskip\noindent\hypertarget{app-F7II-XI-021.4}{}\textbf{F7II-XI-021.4}\quad p. 54; misprint\_pass.\par
\noindent Imprimé : {\ttfamily\footnotesize y\allowbreak{}}\par
\noindent Corrigé : {\ttfamily\footnotesize z\allowbreak{}}\par
\noindent Source : 7II:86. \nolinkurl{fixes/FIXLOG_7II_CURRENT.jsonl#532}; \nolinkurl{fixes/MISPRINTS_7II.json#86}\par

\par\medskip\noindent\hypertarget{app-F7II-XI-021.5}{}\textbf{F7II-XI-021.5}\quad p. 54; misprint\_pass.\par
\noindent Imprimé : {\ttfamily\footnotesize i\allowbreak{}}\par
\noindent Corrigé : {\ttfamily\footnotesize 1\allowbreak{}}\par
\noindent Source : 7II:88. \nolinkurl{fixes/FIXLOG_7II_CURRENT.jsonl#534}; \nolinkurl{fixes/MISPRINTS_7II.json#88}\par

\par\medskip\noindent\hypertarget{app-F7II-XI-021.6}{}\textbf{F7II-XI-021.6}\quad p. 54; designated-witness attribution.\par
\noindent Imprimé : {\ttfamily\footnotesize x\allowbreak{}}\par
\noindent Corrigé : {\ttfamily\footnotesize y\allowbreak{}}\par
\noindent Source : WITNESS-F7II-XI-021.6. \nolinkurl{fr/sga7ii/diplomatic/expose_XI_body.tex}; \nolinkurl{fr/sga7ii/current/expose_XI_body.tex}\par

\par\medskip\noindent\hypertarget{app-F7II-XI-021.7}{}\textbf{F7II-XI-021.7}\quad p. 54; designated-witness attribution.\par
\noindent Imprimé : {\ttfamily\footnotesize y\allowbreak{}}\par
\noindent Corrigé : {\ttfamily\footnotesize z\allowbreak{}}\par
\noindent Source : WITNESS-F7II-XI-021.7. \nolinkurl{fr/sga7ii/diplomatic/expose_XI_body.tex}; \nolinkurl{fr/sga7ii/current/expose_XI_body.tex}\par

\par\medskip\noindent\hypertarget{app-F7II-XI-022.1}{}\textbf{F7II-XI-022.1}\quad p. 56; designated-witness attribution.\par
\noindent Imprimé : \emph{(vide)}\par
\noindent Corrigé : {\ttfamily\footnotesize \textbackslash{}in\allowbreak{}}\par
\noindent Source : WITNESS-F7II-XI-022.1. \nolinkurl{fr/sga7ii/diplomatic/expose_XI_body.tex}; \nolinkurl{fr/sga7ii/current/expose_XI_body.tex}\par

\par\medskip\noindent\hypertarget{app-F7II-XI-022.2}{}\textbf{F7II-XI-022.2}\quad p. 56; designated-witness attribution.\par
\noindent Imprimé : {\ttfamily\footnotesize \textbackslash{}leqslant\allowbreak{} i\allowbreak{} \textbackslash{}leqslant\allowbreak{} d\allowbreak{}}\par
\noindent Corrigé : {\ttfamily\footnotesize P\allowbreak{}}\par
\noindent Source : WITNESS-F7II-XI-022.2. \nolinkurl{fr/sga7ii/diplomatic/expose_XI_body.tex}; \nolinkurl{fr/sga7ii/current/expose_XI_body.tex}\par

\par\medskip\noindent\hypertarget{app-F7II-XI-022.3}{}\textbf{F7II-XI-022.3}\quad p. 56; inherited corrected-layer difference.\par
\noindent Imprimé : \emph{(vide)}\par
\noindent Corrigé : {\ttfamily\footnotesize \_\{i\textbackslash{}in\allowbreak{} P\}\allowbreak{}}\par
\noindent Source : ADOPT-F7II-XI-022.3. \nolinkurl{fr/sga7ii/diplomatic/expose_XI_body.tex}; \nolinkurl{fr/sga7ii/current/expose_XI_body.tex}\par

\par\medskip\noindent\hypertarget{app-F7II-XI-022.4}{}\textbf{F7II-XI-022.4}\quad p. 56; printed misprint.\par
\noindent Imprimé : {\ttfamily\footnotesize terme\allowbreak{}}\par
\noindent Corrigé : {\ttfamily\footnotesize termes\allowbreak{}}\par
\noindent Source : 7II:90. \nolinkurl{fixes/MISPRINTS_7II.json#90}\par

\par\medskip\noindent\hypertarget{app-F7II-XI-023.1}{}\textbf{F7II-XI-023.1}\quad p. 57; inherited corrected-layer difference.\par
\noindent Imprimé : \emph{(vide)}\par
\noindent Corrigé : {\ttfamily\footnotesize \_\{i\textbackslash{}in\allowbreak{} P\}\allowbreak{}}\par
\noindent Source : ADOPT-F7II-XI-023.1. \nolinkurl{fr/sga7ii/diplomatic/expose_XI_body.tex}; \nolinkurl{fr/sga7ii/current/expose_XI_body.tex}\par

\par\medskip\noindent\hypertarget{app-F7II-XI-023.2}{}\textbf{F7II-XI-023.2}\quad p. 57; inherited corrected-layer difference.\par
\noindent Imprimé : \emph{(vide)}\par
\noindent Corrigé : {\ttfamily\footnotesize \_\{i\textbackslash{}in\allowbreak{} P\}\allowbreak{}}\par
\noindent Source : ADOPT-F7II-XI-023.2. \nolinkurl{fr/sga7ii/diplomatic/expose_XI_body.tex}; \nolinkurl{fr/sga7ii/current/expose_XI_body.tex}\par

\par\medskip\noindent\hypertarget{app-F7II-XI-023.3}{}\textbf{F7II-XI-023.3}\quad p. 57; printed misprint.\par
\noindent Imprimé : {\ttfamily\footnotesize mieuxqu\allowbreak{}}\par
\noindent Corrigé : {\ttfamily\footnotesize mieux\allowbreak{} qu\allowbreak{}}\par
\noindent Source : 7II:91. \nolinkurl{fixes/MISPRINTS_7II.json#91}\par

\par\medskip\noindent\hypertarget{app-F7II-XI-024.1}{}\textbf{F7II-XI-024.1}\quad p. 58; printed misprint.\par
\noindent Imprimé : \emph{(vide)}\par
\noindent Corrigé : {\ttfamily\footnotesize 2\allowbreak{}}\par
\noindent Source : 7II:92. \nolinkurl{fixes/MISPRINTS_7II.json#92}\par

\par\medskip\noindent\hypertarget{app-F7II-XI-025.1}{}\textbf{F7II-XI-025.1}\quad p. 59; inherited corrected-layer difference.\par
\noindent Imprimé : {\ttfamily\footnotesize A\allowbreak{}}\par
\noindent Corrigé : {\ttfamily\footnotesize À\allowbreak{}}\par
\noindent Source : ADOPT-F7II-XI-025.1. \nolinkurl{fr/sga7ii/diplomatic/expose_XI_body.tex}; \nolinkurl{fr/sga7ii/current/expose_XI_body.tex}\par

\normalsize\fi
\clearpage\SGAtoc{Exposé XII — QUADRIQUES}
\selectlanguage{french}

\SGApage{II}{62}


S G A 7


\underline{E X P O S E XII}


\underline{\underline{QUADRIQUES}}


\underline{par P. DELIGNE}

0.0. Le présent exposé rassemble quelques résultats bien connus sur les quadriques, dont nous aurons besoin pour démontrer la formule de Picard-Lefschetz.

0.1. Soit $m$ une section d'un $\mathcal{O}_S$-Module $M$ sur un schéma $S$. On dira que $m$ est \underline{partout non nulle} sur $S$ si, pour tout $s \in S$, l'image de $m$ dans le $k(s)$-vectoriel $M_s$ déduit de $M$ par le changement de base $\{s\} \longrightarrow S$ est non nulle. Même terminologie pour $S = \mathrm{Spec}(A)$ et $m$ un élément d'un $A$-module $M$.


\SGAsectiontoc{1. Formes quadratiques}{1}
1. \underline{Formes quadratiques}

1.1. Soient $S = \mathrm{Spec}(A)$ un schéma affine, $V$ un $A$-module localement libre de rang $n$ et $Q$ une forme quadratique sur $V$ (Bourb. Alg 9 \S{} 3 n$^\circ$4). Rappelons que $Q$ est dite \underline{non dégénérée} si la forme bilinéaire associée

\begin{equation*}
\Phi(x,y) = Q(x+y) - Q(x) - Q(y)
\end{equation*}
est non dégénérée, i.e. établit un isomorphisme entre $V$ et son dual $V^*$. Si 2 n'est pas inversible dans $A$, ceci ne peut se produire que pour $n$ pair (loc. cit.).


\SGApage{II}{63}

Les formes quadratiques sur $V$ forment un $A$-module projectif ; une forme quadratique $Q$ est ``partout non nulle'' sur $S$ (0.1) si et seulement si les $Q(v)$ engendrent l'idéal $A$ de $A$.

La notation $\mathbb{P}(V)$ sera prise au sens des EGA ; rappelons que si $A$ est un corps, alors $\mathbb{P}(V)$ est l'espace projectif dont l'ensemble des points rationnels est le quotient $(V^* - \{0\})/A^*$ .

La forme quadratique $Q$ s'identifie à une section du faisceau inversible $\mathcal{O}(2)$ sur l'espace projectif $\mathbb{P}(V^*)$ sur $S$. Pour $Q$ partout non nulle, l'équation (``homogène'') $Q = 0$ définit un sous-schéma de $\mathbb{P}(V^*)$, plat et purement de dimension relative $n-2$ sur $S$. Ce schéma s'appelle la \underline{quadrique définie par} $Q$. La forme $Q$ est dite \underline{ordinaire} si elle n'est nulle en aucun point de $S$ et que la quadrique qu'elle définit est lisse sur $S$.

Pour que $Q$ soit ordinaire, il faut et il suffit que les formes qui s'en déduisent après extension des scalaires de $A$ à un corps le soient. Si $A$ est un corps, alors :

a) pour $n$ pair ou $\mathrm{car}(A) \neq 2$ : $Q$ ordinaire $\Longleftrightarrow$ $Q$ non dégénéré ;

\SGAvar{b) pour $n$ impair et $\mathrm{card}(A) = 2$ : $Q$ est ordinaire si et seulement si le noyau $N$ de la forme bilinéaire (alternée) associée $\Phi$ est de dimension un, et que $Q$ n'est pas nul sur $N$.

}{b) pour $n$ impair et $\mathrm{car}(A) = 2$ : $Q$ est ordinaire si et seulement si le noyau $N$ de la forme bilinéaire (alternée) associée $\Phi$ est de dimension un, et que $Q$ n'est pas nul sur $N$.

}{F7II-XII-001}\underline{Proposition} 1.2. \underline{Si} $Q$ \underline{est ordinaire et si} $n = 2m$ (resp. $n = 2m +1$), \underline{alors}, \underline{localement pour la topologie étale sur} $S$, $V$ \underline{admet une base} \underline{$e$} \underline{telle que} 


\SGApage{II}{64}

\SGAvar{\begin{equation*}
 Q\Bigl( \sum_{1}^{n} x_i e_i \Bigr) \;=\; \sum_{1}^{m-1} x_i x_{i+m}
\end{equation*}
\begin{equation*}
\text{(resp.} \qquad Q\Bigl( \sum_{1}^{n} x_i e_i \Bigr) \;=\; \sum_{1}^{m-1} x_i x_{i+m} + \lambda x_{2m+1}^{2} \qquad (\lambda \text{ inversible}) \;.
\end{equation*}
Prouvons 1.2 par récurrence sur $m$ . Si $m = 0$, l'assertion est évidente. Si $m > 0$, la quadrique d'équation $Q = 0$ est lisse sur $S$ et à fibres géométriques non vides ; elle admet donc des sections localement pour la topologie étale. Localement, une telle section est définie par un élément $e \in V$, nul en aucun point de $\mathrm{Spec}(A)$, et tel que $Q(e) = 0$. En aucun point de $S$, $e$ n'est dans le noyau de $\Phi$ ; localement, on peut donc trouver $f'$ tel que $\Phi(e,f') = 1$. Si $f = -Q(f').e+f$, on a

}{\begin{equation*}
 Q\Bigl( \sum_{1}^{n} x_i e_i \Bigr) \;=\; \sum_{1}^{m} x_i x_{i+m}
\end{equation*}
\begin{equation*}
\text{(resp.} \qquad Q\Bigl( \sum_{1}^{n} x_i e_i \Bigr) \;=\; \sum_{1}^{m} x_i x_{i+m} + \lambda x_{2m+1}^{2} \qquad (\lambda \text{ inversible})) \;.
\end{equation*}
Prouvons 1.2 par récurrence sur $m$ . Si $m = 0$, l'assertion est évidente. Si $m > 0$, la quadrique d'équation $Q = 0$ est lisse sur $S$ et à fibres géométriques non vides ; elle admet donc des sections localement pour la topologie étale. Localement, une telle section est définie par un élément $e \in V$, nul en aucun point de $\mathrm{Spec}(A)$, et tel que $Q(e) = 0$. En aucun point de $S$, $e$ n'est dans le noyau de $\Phi$ ; localement, on peut donc trouver $f'$ tel que $\Phi(e,f') = 1$. Si $f = -Q(f').e+f'$, on a

}{F7II-XII-002}\begin{equation*}
Q(e) = Q(f) = 0 \;,\; \Phi(e,f) = 1 \quad .
\end{equation*}
Si $V_2$ est l'orthogonal du sous-module $V_1$ de $V$ engendré par $e$ et $f$, on a $V = V_1 \oplus V_2$, et on conclut en appliquant l'hypothèse de récurrence à $V_2$ .

1.3. Pour la définition et les principales propriétés de l'algèbre de Clifford de $Q$, on renvoie à (Bourb. Alg 9 \S{} 9). On désigne par $\rho$ l'application canonique de $V$ dans $C(Q)$.

\underline{Lemme} 1.3.1. \underline{Sous les hypothèses de} 1.1, \underline{il existe pour} $\lambda \in A$ \underline{un et un seul homomorphisme}

\begin{equation*}
[\lambda] \;:\; C^{+}(\lambda\, Q) \longrightarrow C^{+}(Q)
\end{equation*}
\underline{tel que}, \underline{pour} $x, y \in V$, \underline{on ait}

\begin{equation*}
[\lambda]\,(\rho(x).\rho(y)) \;=\; \lambda.\rho(x).\rho(y) \quad .
\end{equation*}

\SGApage{II}{65}

\SGAvar{Cet homomorphisme se déduit par passage au quotient de l'endomorphisme $[\lambda]$ de la partie $T^{+}(E)$ de l'algèbre tensorielle de $E$, qui vaut $\lambda^{n}$ sur $T^{2n}(E)$. En effet,

\begin{equation*}
[\lambda]\,(x \otimes x - \lambda\, Q(x)) \;=\; \lambda(x \otimes x - Q(x)) \quad .
\end{equation*}
Supposons $Q$ partout non nulle sur $S$. Alors, le lemme 1.2.1 définit un système transitif d'isomorphismes entre les algèbres $C^{+}(\lambda\, Q)$ pour $\lambda \in A^{*}$.

}{Cet homomorphisme se déduit par passage au quotient de l'endomorphisme $[\lambda]$ de la partie $T^{+}(V)$ de l'algèbre tensorielle de $V$, qui vaut $\lambda^{n}$ sur $T^{2n}(V)$. En effet,

\begin{equation*}
[\lambda]\,(x \otimes x - \lambda\, Q(x)) \;=\; \lambda(x \otimes x - Q(x)) \quad .
\end{equation*}
Supposons $Q$ partout non nulle sur $S$. Alors, le lemme 1.3.1 définit un système transitif d'isomorphismes entre les algèbres $C^{+}(\lambda\, Q)$ pour $\lambda \in A^{*}$.

}{F7II-XII-003}Dès lors, l'algèbre $C^{+}(Q)$ ne dépend (à isomorphisme unique près) que de $V$ et de la donnée de $Q$ à un facteur près ; en d'autres termes, $C^{+}(Q)$ ne dépend que de $V$ et du sous-schéma d'équation $Q = 0$ de $\mathbb{P}(V^{*})$. Chaque homothétie de rapport inversible est un automorphisme du couple $(V, X \subset \mathbb{P}(V^{*}))$ ; on vérifie que ces automorphismes agissent trivialement sur $C^{+}(V, X)$, et il en résulte que l'algèbre $C^{+}(V, X)$ ne dépend que du couple de schémas $X \subset \mathbb{P}(V^{*})$ sur $A$. On la notera $C^{+}(X, \mathbb{P}(V^{*}))$.

1.4. Soit $V$ un $A$-module localement libre de rang $2m > 0$ et $Q$ une forme quadratique non dégénérée sur $V$ . Soit donné une décomposition $V = W_1 \oplus W_2$ de $V$ en deux sous-espaces totalement singuliers. La forme $Q(w_1 + w_2)$ met alors $W_1$ et $W_2$ en dualité parfaite. Rappelons que pour $f \in W_2$ (identifié à $W_1^{*}$), le produit contracté par $f$ est l'endomorphisme de degré $-1$ de $\wedge W_1$ défini par

\SGAvar{\begin{equation*}
f \llcorner (e_1 \wedge \ldots \wedge e_p) \;=\; \sum (-1)^{i} < e_i , f > e_o \wedge \ldots \; \hat{e}_i \ldots \wedge e_p \quad .
\end{equation*}
Munissons le $A$-module $\wedge W_1$ de la $\mathbb{Z}/2$-graduation pour laquelle
\SGApage{II}{66}
$(\wedge W_1)^+ = \underset{i}{\oplus} \overset{2i}{\wedge} W_1$ et $(\wedge W_1)^- = \underset{i}{\oplus} \overset{2i+1}{\wedge} W_1$ . L'algèbre $\mathrm{End}(\wedge W_1)$ est alors $\mathbb{Z}/2$-gradué : les endomorphismes pairs (resp. impairs) sont ceux qui respectent la graduation (resp. sont de degré 1 (mod 2)).

On renvoie à (Bourb Alg 9 § 9 n°4) pour la démonstration de l'existence d'un unique isomorphisme d'algèbre $\mathbb{Z}/2$-graduées $s$ entre $C(Q)$ et $\mathrm{End}(\wedge W_1)$, tel que

\begin{equation*}
s(\rho((e,0)) = e \wedge \qquad \mbox{et} \qquad s(\rho((0,f)) = f \llcorner \quad .
\end{equation*}
Désignons provisoirement par $e(W_1,W_2)$ l'élément de $C^+(Q)$ tel que $s(e(W_1,W_2))$ soit l'identité sur $(\wedge W_1)^+$ et $0$ sur $(\wedge W_1)^-$ . Le centre $Z(Q)$ de $C^+(Q)$ est isomorphe à $A \oplus A$, cette décomposition étant définie par les idempotents $e(W_1,W_2)$ et $1 - e(W_2,W_1)$. De plus $C^+(Q)$ est une algèbre de matrice sur $Z(Q)$.

}{\begin{equation*}
f \llcorner (e_1 \wedge \ldots \wedge e_p) \;=\; \sum (-1)^{i-1} < e_i , f > e_1 \wedge \ldots \; \hat{e}_i \ldots \wedge e_p \quad .
\end{equation*}
Munissons le $A$-module $\wedge W_1$ de la $\mathbb{Z}/2$-graduation pour laquelle
\SGApage{II}{66}
$(\wedge W_1)^+ = \underset{i}{\oplus} \overset{2i}{\wedge} W_1$ et $(\wedge W_1)^- = \underset{i}{\oplus} \overset{2i+1}{\wedge} W_1$ . L'algèbre $\mathrm{End}(\wedge W_1)$ est alors $\mathbb{Z}/2$-graduée : les endomorphismes pairs (resp. impairs) sont ceux qui respectent la graduation (resp. sont de degré 1 (mod 2)).

On renvoie à (Bourb Alg 9 § 9 n°4) pour la démonstration de l'existence d'un unique isomorphisme d'algèbres $\mathbb{Z}/2$-graduées $s$ entre $C(Q)$ et $\mathrm{End}(\wedge W_1)$, tel que

\begin{equation*}
s(\rho((e,0))) = e \wedge \qquad \mbox{et} \qquad s(\rho((0,f))) = f \llcorner \quad .
\end{equation*}
Désignons provisoirement par $e(W_1,W_2)$ l'élément de $C^+(Q)$ tel que $s(e(W_1,W_2))$ soit l'identité sur $(\wedge W_1)^+$ et $0$ sur $(\wedge W_1)^-$ . Le centre $Z(Q)$ de $C^+(Q)$ est isomorphe à $A \oplus A$, cette décomposition étant définie par les idempotents $e(W_1,W_2)$ et $1 - e(W_1,W_2)$. De plus $C^+(Q)$ est une algèbre de matrice sur $Z(Q)$.

}{F7II-XII-004}De ceci et de 1.2, on déduit comme en loc.cit.

\underline{Proposition} 1.5. \underline{Sous les hypothèses de} 1.1, \underline{si} $n = 2m > 0$ \underline{et que} $Q$ \underline{est non dégénérée, alors le centre} $Z(Q)$ \underline{de} $C^+(Q)$ \underline{est une algèbre étale localement libre de rang} 2 \underline{sur} $A$ \underline{et} $C^+(Q)$ \underline{est une algèbre d'Azumaya sur} $Z(Q)$.

1.6. Soient $V$ et $Q$ comme en 1.4, et $W$ un facteur direct totalement singulier localement libre de rang $m$ de $V$. Soit $W'_1$ un supplémentaire de $W : V = W \oplus W'_1$ . Tout supplémentaire $W'_2$ de $W$ s'identifie au graphe d'une application $u$ de $W'_1$ dans $W$. La forme $\Phi$ met en dualité $W$ et $W'_1$ ; l'application $u$ est donc uniquement déterminée par la forme bilinéaire sur $W'_1$ 


\SGApage{II}{67}

\begin{equation*}
 B(x,y) = \Phi(x,u(y)) \quad ,
\end{equation*}
et $W'_2$ est totalement singulier si et seulement si, sur $W'_1$, on a

\SGAvar{\begin{equation*}
B(x,x) = - Q(x) \quad .
\end{equation*}
Puisqu'il existe toujours de telles formes $B$, $W$ admet des supplémentaires totalement singulier. L'un d'eux, soit $W'$, étant choisi, l'ensemble des supplémentaire totalement singulier $W'_o$ de $W$ s'identifie à l'ensemble des formes alternées sur $W'$, soit encore à l'ensemble des sections de l'espace affine $\mathbb{E}(\overset{2}{\wedge} W')$ sur $S$. Pour chaque section $x$ de $\mathbb{E}(\overset{2}{\wedge} W')$, soit $W'_o(x)$ le supplémentaire correspondant de $W$ et $e(x)$ l'idempotent $e(W, W'_o(x))$. La fonction $x \longmapsto e(x)$ est compatible à toute extension des scalaires, dont définit un morphisme de $S$-schémas $e$ de $\mathbb{E}(\overset{2}{\wedge} W')$ dans $\mathrm{Spec}(Z(Q))$. Le $S$-schéma $\mathbb{E}(\overset{2}{\wedge} W')$ est plat à fibres géométriques connexes, alors que $\mathrm{Spec}(Z(Q))$ est étale sur $S$. Le morphisme $e$ se factorise donc par une section de $\mathbb{E}(\overset{2}{\wedge} W')$ sur $S$, i.e. $e(x)$ est indépendant de $x$. Ceci justifie la

}{\begin{equation*}
B(x,x) = - Q(x) \quad .
\end{equation*}
Puisqu'il existe toujours de telles formes $B$, $W$ admet des supplémentaires totalement singuliers. L'un d'eux, soit $W'$, étant choisi, l'ensemble des supplémentaires totalement singuliers $W'_o$ de $W$ s'identifie à l'ensemble des formes alternées sur $W'$, soit encore à l'ensemble des sections de l'espace affine $\mathbb{E}(\overset{2}{\wedge} W')$ sur $S$. Pour chaque section $x$ de $\mathbb{E}(\overset{2}{\wedge} W')$, soit $W'_o(x)$ le supplémentaire correspondant de $W$ et $e(x)$ l'idempotent $e(W, W'_o(x))$. La fonction $x \longmapsto e(x)$ est compatible à toute extension des scalaires, donc définit un morphisme de $S$-schémas $e$ de $\mathbb{E}(\overset{2}{\wedge} W')$ dans $\mathrm{Spec}(Z(Q))$. Le $S$-schéma $\mathbb{E}(\overset{2}{\wedge} W')$ est plat à fibres géométriques connexes, alors que $\mathrm{Spec}(Z(Q))$ est étale sur $S$. Le morphisme $e$ se factorise donc par $S$, et par suite par une section de $\mathrm{Spec}(Z(Q))$ sur $S$, i.e. $e(x)$ est indépendant de $x$. Ceci justifie la

}{F7II-XII-005}\underline{Notation} 1.7. \underline{Pour} $W$ \underline{comme plus haut et} $W'$ \underline{un quelconque supplémentaire totalement isotrope de} $W$, \underline{on désigne par} $e(W)$ \underline{l'idempotent} $e(W,W') \in Z(Q)$.

\SGAvar{1.8. On a vu en 1.6 que, pour $\lambda$ inversible, $C^+(Q)$ s'identifie à $C^+(\lambda Q)$. Un argument de continuité analogue au précédent montre que cette identification transforme

}{1.8. On a vu en 1.3 que, pour $\lambda$ inversible, $C^+(Q)$ s'identifie à $C^+(\lambda Q)$. Un argument de continuité analogue au précédent montre que cette identification transforme

}{F7II-XII-006}\begin{equation*}
e(W) \in Z(Q) \quad \mbox{en} \quad e(W) \in Z(\lambda Q) \quad .
\end{equation*}

\SGApage{II}{68}

\SGAvar{1.9. Rappelons que si $B$ et $C$ sont deux algèbres $\mathbb{Z}/2$-gradués, alors, dans le produit tensoriel gradué $B \otimes' C$ de $B$ et $C$, la multiplication est définie par la règle de Koszul

}{1.9. Rappelons que si $B$ et $C$ sont deux algèbres $\mathbb{Z}/2$-graduées, alors, dans le produit tensoriel gradué $B \otimes' C$ de $B$ et $C$, la multiplication est définie par la règle de Koszul

}{F7II-XII-007}\begin{equation*}
(b_1 \otimes c_1)\ (b_2 \otimes c_2) = (-1)^{\deg(c_1)\deg(b_2)}\ b_1 b_2 \otimes c_1 c_2 \quad ,
\end{equation*}
lorsque les $b_i$ et $c_i$ sont homogènes.

Si $(V_1, Q_1)$ et $(V_2, Q_2)$ sont deux modules quadratiques, alors $C(Q_1 \oplus Q_2) = C(Q_1) \otimes' C(Q_2)$ (Bourb. loc. cit.).

\SGAvar{Si $P$ et $Q$ sont deux $A$-modules localement libres $\mathbb{Z}/2$-gradués de type fini sur $A$, on identifie les algèbres $\mathbb{Z}/2$-gradués $\mathrm{End}(P) \otimes' \mathrm{End}(Q)$ et $\mathrm{End}(P \otimes Q)$, en posant, pour $u$, $v$, $p$ et $q$ homogènes

}{Si $P$ et $Q$ sont deux $A$-modules localement libres $\mathbb{Z}/2$-gradués de type fini sur $A$, on identifie les algèbres $\mathbb{Z}/2$-graduées $\mathrm{End}(P) \otimes' \mathrm{End}(Q)$ et $\mathrm{End}(P \otimes Q)$, en posant, pour $u$, $v$, $p$ et $q$ homogènes

}{F7II-XII-008}\begin{equation*}
u \otimes v\ (p \otimes q) = (-1)^{\deg(v)\deg(p)} u(p) \otimes v(q) \quad .
\end{equation*}
\underline{Lemme} 1.10. \underline{Soit} $(V,Q)$ \underline{un module quadratique somme de deux modules quadratiques} $(V_1, Q_1)$ \underline{et} $(V_2, Q_2)$, \underline{avec} $V_i$ \underline{localement libre de rang} $2\ m_i$ \underline{et} $Q_i$ \underline{non dégénérée.} \underline{Soit} $V_i = W_i \oplus W'_i$ \underline{une décomposition de} $V_i$ \underline{en deux sous-espaces totalement singuliers} $(i = 1, 2)$. \underline{Alors le diagramme}

\begin{center}
\begin{tikzcd}[column sep=huge, row sep=large]
C(Q) \arrow[rr, dash, "\sim"] \arrow[d, dash, "\wr"] \& {} \& C(Q_1) \otimes' C(Q_2) \arrow[d, equal] \\
\mathrm{End}(\wedge(W_1 \oplus W_2)) \arrow[r, phantom, "\simeq"] \& \mathrm{End}(\wedge W_1 \otimes \wedge W_2) \arrow[r, phantom, "\simeq"] \& \mathrm{End}(\wedge W_1) \otimes' \mathrm{End}(\wedge W_2)
\end{tikzcd}
\end{center}













\underline{est commutatif.}


\SGApage{II}{69}

La vérification est laissée au lecteur. On déduit aussitôt de ce lemme que

\SGAvar{\begin{equation}\tag{1.10.1}
e(W_1 \oplus W_2) = e(W_1)e(W_2) + (1 - e(W_1))\ (1 - e(W_2))
\end{equation}
1.11. La construction suivante est empruntée à Micali et Villamayor. Si $p : Z \longrightarrow S$ est un revêtement étale double de $S$, on peut ragarder $Z$ comme un $\mathbb{Z}/2$-torseurs sur $S$. Désignons par $*$ l'addition des torseurs, et pour $s$ une section de $Z$, soit $e(s)$ l'idempotent de $\rho_* \mathcal{O}_Z$ qui vaut $1$ sur $s$, et $0$ ailleurs. Avec les notations de 1.10, il existe un unique isomorphisme

}{\begin{equation}\tag{1.10.1}
e(W_1 \oplus W_2) = e(W_1)e(W_2) + (1 - e(W_1))\ (1 - e(W_2))
\end{equation}
1.11. La construction suivante est empruntée à Micali et Villamayor. Si $p : Z \longrightarrow S$ est un revêtement étale double de $S$, on peut regarder $Z$ comme un $\mathbb{Z}/2$-torseur sur $S$. Désignons par $*$ l'addition des torseurs, et pour $s$ une section de $Z$, soit $e(s)$ l'idempotent de $p_* \mathcal{O}_Z$ qui vaut $1$ sur $s$, et $0$ ailleurs. Avec les notations de 1.10, il existe un unique isomorphisme

}{F7II-XII-009}\begin{equation}\tag{1.11.1}
\mathrm{Spec}(Z(Q_1)) * \mathrm{Spec}(Z(Q_2)) \longrightarrow \mathrm{Spec}(Z(Q))
\end{equation}
\begin{equation*}
+ : \mathrm{Spec}(Z(Q_1))\ \mathbf{x}_S\ \mathrm{Spec}(Z(Q_2)) \longrightarrow \mathrm{Spec}(Z(Q))
\end{equation*}
pour lequel, après tout changement de base, on ait

\begin{equation*}
e(s + t) = e(s)\ e(t) + (1 - e(s))\ (1 - e(t)) \quad .
\end{equation*}
Si $s(W)$ est la section de $\mathrm{Spec}(Z(Q))$ telle que $e(s(W)) = e(W)$, la formule 1.10.1 se récrit

\SGAvar{\begin{equation*}
s(W_1 \oplus W_2)\ =\ s(W_1) + s(W_2) \quad .
\end{equation*}
\underline{Proposition} 1.12. \underline{Soient} $A$ \underline{un corps,} $V$ \underline{un vectoriel de dimension} $2m > 0$ \underline{sur} $V$, $Q$ \underline{une forme quadratique non dégénérée sur} $V$ \underline{et} $W_1$, $W_2$ \underline{deux sous-espaces totalement isotropes de dimension} $m$ \underline{de} $V$. \underline{Alors}, $e(W_1) = e(W_2)$ \underline{si et seulement si} $\dim(W_1/W_1 \cap W_2)$ \underline{est pair.}


}{\begin{equation*}
s(W_1 \oplus W_2)\ =\ s(W_1) + s(W_2) \quad .
\end{equation*}
\underline{Proposition} 1.12. \underline{Soient} $A$ \underline{un corps,} $V$ \underline{un vectoriel de dimension} $2m > 0$ \underline{sur} $A$, $Q$ \underline{une forme quadratique non dégénérée sur} $V$ \underline{et} $W_1$, $W_2$ \underline{deux sous-espaces totalement isotropes de dimension} $m$ \underline{de} $V$. \underline{Alors}, $e(W_1) = e(W_2)$ \underline{si et seulement si} $\dim(W_1/W_1 \cap W_2)$ \underline{est pair.}


}{F7II-XII-010}\SGApage{II}{70}

(Cf. Bourbaki Alg 9 \S 6 ex. 28). Soit $e_1 \dots e_m\ f_1 \dots f_m$ une base de $V$ telle que

a) \qquad $Q( \Sigma x_i e_i + \Sigma y_i f_i) = \Sigma x_i y_i$ ;

b) \qquad $W_1$ est engendré par les $e_i$ ;

c) \qquad et $W_2$ est engendré par les $(e_i)_{1 \leqslant i \leqslant s}$ et les $(f_i)_{s < i \leqslant m}$ .

Alors $V$ est somme des $V_i = A e_i + A f_i$, et

\begin{equation*}
W_\alpha \;=\; \oplus \; W_\alpha \cap V_i \qquad (\alpha = 1,2) \quad .
\end{equation*}
La formule d'addition 1.11 permet dès lors de se limiter au cas où $\dim V = 2$. Si $W_1 = W_2$, l'assertion est évidente. Si au contraire $W_1 = Ae$, et $W_2 = Af$ avec $\Phi(e,f) = 1$, on a $e(W_1) = fe$ et $e(W_2) = ef = 1 - fe$.


\SGAvar{\SGAsectiontoc{2. Quadriques'}{2}
\SGAsectiontoc{2. Quadriques}{2}
2. \underline{Quadriques}'

}{\SGAsectiontoc{2. Quadriques}{2}
2. \underline{Quadriques}

}{F7II-XII-011}2.1. Soient $k$ un corps algébriquement clos et $V$ un vectoriel de dimension $n+2$ sur $k$. Une quadrique lisse de dimension $n$ sur $k$ est un schéma $X$ sur $k$ isomorphe au sous-schéma de $\mathbb{P}(V^*)$ défini par l'équation $Q=0$, pour $Q$ une forme quadratique ordinaire sur $k$.

\underline{Remarque} 2.2. Pour $n = 0,1,2$ $X$ est une quadrique lisse de dimension $n$ sur $k$ si et seulement si :

$n = 0$ \qquad : \qquad $X \approx \mathrm{Spec}(k) \amalg \mathrm{Spec}(k)$

$n = 1$ \qquad : \qquad $X \approx \mathbb{P}^1_k$

$n = 2$ \qquad : \qquad $X \approx \mathbb{P}^1_k \times \mathbb{P}^1_k$ .


\SGApage{II}{71}

\underline{Lemme} 2.3. \underline{Avec les notations de} 2.1, \underline{soit} $X$ \underline{défini par} $Q = 0$ .

(i) \qquad $\Omega^n_X \approx \mathcal{O}_X(-n)$

(ii) \quad $\mathrm{Pic}(X) \approx \left\{ \begin{array}{l} \mbox{$0$ \qquad \underline{pour} $n = 0$ .} \\ \mbox{$\mathbb{Z}$ \qquad \underline{pour} $n = 1$ ; $\mathcal{O}_X(1)$ \underline{est le carré d'un générateur.}} \\ \mbox{$\mathbb{Z} \oplus \mathbb{Z}$ \underline{pour} $n =2$.} \\ \mbox{$\mathbb{Z}$ \qquad \underline{pour} $n \geqslant 3$ ; $\mathcal{O}_X(1)$ \underline{est un générateur.}} \end{array} \right.$

(iii) \quad $H^i(X,\mathcal{O}_X) = 0$ \underline{pour} $i > 0$, \qquad $H^1(X,\mathcal{O}_X(1)) = 0$

(iv) \quad $H^0(X,\mathcal{O}_X(1)) \stackrel{\sim}{\longleftarrow} H^0( \mathbb{P}(V^*),\mathcal{O}_X(1)) = V^*$, \underline{et pour} $n > 0$, $H^0(X,\mathcal{O}_X) = k$ .

L'assertion (i) résulte de ce que le faisceau conormal de $X$ dans $\mathbb{P}(V^*)$ est $\mathcal{O}_X(-2)$ et que le faisceau canonique de $\mathbb{P}^r$ est $\mathcal{O}(-r-1)$. L'assertion (ii) résulte de 2.2 pour $n \leqslant 2$, de SGA 2 XII 3.7 pour $n \geqslant 3$. Enfin (iii) et (iv) résultent de FAC n$^\circ$ 78.

\underline{Définition} 2.4. \underline{Une quadrique lisse de dimension} $n$ \underline{sur un schéma} $S$ \underline{est un} $S$-\underline{schéma} $f : X \longrightarrow S$, \underline{propre et lisse sur} $S$, \underline{dont les fibres géométriques sont des quadriques lisses.}

Pour $n = 0$, une quadrique lisse de dimension $0$ sur $S$ n'est autre qu'un revêtement étale double de $S$. Pour $n = 1$, ce n'est autre qu'une variété de Severi-Brauer de dimension $1$ sur $S$.

2.5. Si $\rho : X \longrightarrow S$ est une quadrique lisse de dimension $n$ sur $S$, alors, d'après 2.3 (i), l'inverse du faisceau inversible $\Omega^n_{X/S}$ est ample.


\SGApage{II}{72}

\SGAvar{D'après 2.3 (ii), le schéma de Picard $\mathrm{Pic}_{X/S}$ est lisse sur $S$ (car $\mathrm{H}^2(X,\mathcal{O}_X) = 0$), non ramifié sur $S$ (car $\mathrm{H}^1(X,\mathcal{O}_X) = 0$) et essentiellement propre sur $S$ (car $X/S$ est lisse). C'est donc, localement sur $S$ pour la topologie étale, un groupe discret constant ayant la structure indiquée en 2.3 (iii).

En particulier, si $X$ a une section, il existe sur $X$ un faisceau inversible $L$ tel que $L^{\otimes(-n)}$ ait même image dans $\mathrm{Pic}_{X/S}$ que $\Omega^n_{X/S}$ , et deux tels faisceaux sont isomorphes localement sur $S$. D'après 2.3 (iii) et [1] 0.5 p. 19 $p_*L$ est localement libre de formation compatible à tout changement de base. Dès lors, $p_*L$ est localement libre de rang $n + 2$ (2.3 (iv)), l'application de $p^*p_*L$ dans $L$ est surjective, l'application canonique de $X$ dans $\mathbb{P}(p_*L)$ est une immersion fermée, elle identifie $X$ à un diviseur relatif de degré 2 de $\mathbb{P}(p_*L)$, et, localement sur $S$, $X$ est la quadrique de $\mathbb{P}(p_*L)$ définie par une forme quadratique ordinaire sur $(p_*L)^*$, unique à un facteur inversible près.

2.6. Supposons $n > 0$. Si $L_1$ et $L_2$ sont deux faisceaux inversibles du type considéré plus haut, alors, localement sur $S$, il existe un isomorphisme $u : L_1 \stackrel{\sim}{\longrightarrow} L_2$. Cet isomorphisme est unique à un facteur inversible près, car $p_*\mathcal{O}_X = \mathcal{O}_S$ (d'après 2.3 (iv)). Dès lors, les divers isomorphismes $u : L_1 \stackrel{\sim}{\longrightarrow} L_2$ induisant le \underline{même} isomorphisme $\mathbb{P}(u): \mathbb{P}(p_*L_1) \stackrel{\sim}{\longrightarrow} \mathbb{P}(p_*L_2)$, de sorte que $\mathbb{P}(p_*L_1)$ est canoniquement isomorphe à $\mathbb{P}(p_*L_2)$, et ne dépend que de $X/S$, non du choix d'un $L$. On le notera $P(X)$.


}{D'après 2.3 (iii), le schéma de Picard $\mathrm{Pic}_{X/S}$ est lisse sur $S$ (car $\mathrm{H}^2(X,\mathcal{O}_X) = 0$), non ramifié sur $S$ (car $\mathrm{H}^1(X,\mathcal{O}_X) = 0$) et essentiellement propre sur $S$ (car $X/S$ est lisse). C'est donc, localement sur $S$ pour la topologie étale, un groupe discret constant ayant la structure indiquée en 2.3 (ii).

En particulier, si $X$ a une section, il existe sur $X$ un faisceau inversible $L$ tel que $L^{\otimes(-n)}$ ait même image dans $\mathrm{Pic}_{X/S}$ que $\Omega^n_{X/S}$ , et deux tels faisceaux sont isomorphes localement sur $S$. D'après 2.3 (iii) et [1] 0.5 p. 19 $\rho_*L$ est localement libre de formation compatible à tout changement de base. Dès lors, $p_*L$ est localement libre de rang $n + 2$ (2.3 (iv)), l'application de $p^*p_*L$ dans $L$ est surjective, l'application canonique de $X$ dans $\mathbb{P}(p_*L)$ est une immersion fermée, elle identifie $X$ à un diviseur relatif de degré 2 de $\mathbb{P}(p_*L)$, et, localement sur $S$, $X$ est la quadrique de $\mathbb{P}(p_*L)$ définie par une forme quadratique ordinaire sur $(p_*L)^*$, unique à un facteur inversible près.

2.6. Supposons $n > 0$. Si $L_1$ et $L_2$ sont deux faisceaux inversibles du type considéré plus haut, alors, localement sur $S$, il existe un isomorphisme $u : L_1 \stackrel{\sim}{\longrightarrow} L_2$. Cet isomorphisme est unique à un facteur inversible près, car $p_*\mathcal{O}_X = \mathcal{O}_S$ (d'après 2.3 (iv)). Dès lors, les divers isomorphismes $u : L_1 \stackrel{\sim}{\longrightarrow} L_2$ induisent le \underline{même} isomorphisme $\mathbb{P}(u): \mathbb{P}(p_*L_1) \stackrel{\sim}{\longrightarrow} \mathbb{P}(p_*L_2)$, de sorte que $\mathbb{P}(p_*L_1)$ est canoniquement isomorphe à $\mathbb{P}(p_*L_2)$, et ne dépend que de $X/S$, non du choix d'un $L$. On le notera $P(X)$.


}{F7II-XII-012}\SGApage{II}{73}

L'espace projectif $P(X)$ est défini localement sur $S$ pour la topologie étale. Par descente, il définit une variété de Severi-Brauer sur $S$, encore notée $P(X)$, donc $X$ est un diviseur relatif de degré 2.

Localement sur $S$, le couple $X \hookrightarrow P(X)$ définit une algèbre de Clifford $C^+(X, P(X))$ (1.3). Par descente, on obtient une algèbre sur $S$, qu'on notera simplement $C^+(X)$.

2.7. Supposons que $n = 2m > 0$ . Le centre $Z$ de $C^+(X)$ est alors une $\mathcal{O}_S$-algèbre définissant un revêtement étale double $Z(X)$ de $S$.

Une \underline{génératrice} de $X$ est un sous-schéma $D$ de $X$ qui soit un sous-espace linéaire de dimension $m$ de $P(X)$.

\SGAvar{Localement sur $S$, pour $P(X) \approx \mathbb{P}(V^*)$ et $X$ défini par l'équation $Q = 0$, une génératrice $D$ de $X$ s'identifie par définition à un sous-Module localement facteur direct $W(D)$ de $V$, totalement singulier et de rang $m + 1$. D'après 1.8, l'idempotent $e(W(D))$ du centre de $C^+(X)$ ne dépend que de $D \subset X$. On désigne par $e(D)$ la section correspondante de $Z(X)$ (celle sur lequelle $e(W(D))$ vaut 1).

}{Localement sur $S$, pour $P(X) \approx \mathbb{P}(V^*)$ et $X$ défini par l'équation $Q = 0$, une génératrice $D$ de $X$ s'identifie par définition à un sous-Module localement facteur direct $W(D)$ de $V$, totalement singulier et de rang $m + 1$. D'après 1.8, l'idempotent $e(W(D))$ du centre de $C^+(X)$ ne dépend que de $D \subset X$. On désigne par $e(D)$ la section correspondante de $Z(X)$ (celle sur laquelle $e(W(D))$ vaut 1).

}{F7II-XII-013}Les génératrices de $X/S$ sont les $S$-points d'un schéma $\mbox{Gén}(X)$ projectif sur $S$, et $e$ définit un morphisme de $S$-schémas

\begin{equation}\tag{2.7.1}
e : \mbox{Gén}(X) \longrightarrow Z(X) \quad .
\end{equation}
Lorsque $n = 0$, on pose $Z(X) = X$, on appelle génératrice de $X$ une section de $X$ et on définit $e$ comme étant l'identité.


\SGApage{II}{74}

\underline{Proposition} 2.8. \underline{Sous les hypothèses de} 2.7, $p : \mathrm{G\acute{e}n}(X) \longrightarrow S$ \underline{est projectif et lisse, et} (2.7.1) \underline{est sa factorisation de Stein.}

La question est locale pour la topologie étale sur $S$. Localement, si $D$ est une génératrice, les génératrices $D'$ disjointes de $D$ s'identifient, d'après 1.6, aux $S$-points d'un espace affine sur $S$ convenable, et on recouvre ainsi $\mathrm{G\acute{e}n}(X)$ par des ouverts lisses sur $S$.

Pour prouver que les fibres géométriques de $e$ sont géométriquement connexes, il suffit de considérer le cas où $S$ est le spectre d'un corps algébriquement clos $k$, $X$ étant défini par l'équation $Q = 0$ dans $\mathbb{P}(V^*)$.

\SGAvar{Soient donc $W_1$ et $W^*_2$ deux sous-espaces totalement isotropes de rang $m + 1$ de $V$ tels que $e(W) = e(W')$, i.e. tels que $W/W \cap W'$ soit de dimension paire (1.12). Procédant comme en 1.12, on trouve une décomposition orthogonale

}{Soient donc $W_1$ et $W_2$ deux sous-espaces totalement isotropes de rang $m + 1$ de $V$ tels que $e(W_1) = e(W_2)$, i.e. tels que $W_1/W_1 \cap W_2$ soit de dimension paire (1.12). Procédant comme en 1.12, on trouve une décomposition orthogonale

}{F7II-XII-014}\begin{equation*}
V = V_o \oplus \bigoplus V_i
\end{equation*}
de $V$ telle que

a) \quad $W_\alpha = W_\alpha \cap V_o \oplus \bigoplus W_\alpha \cap V_i$ \qquad $(\alpha = 1,2)$

b) \quad $W_1 \cap V_o = W_2 \cap V_o$

c) \quad $\dim V_i = 4$, et $V_i = (W_1 \cap V_i) \oplus (W_2 \cap V_i)$

et il suffit de prouver que $W_1 \cap V_i$ et $W_2 \cap V_i$ appartiennent à une même famille connexe de sous-espaces totalement singuliers maximaux de $V_i$ . En d'autres termes, on se ramène au cas des quadriques de dimension 2. Une
\SGApage{II}{75}
telle quadrique est isomorphe à $\mathbb{P}^1 \times \mathbb{P}^1$, les génératrices étant les sous-schémas $\{t\} \times \mathbb{P}^1$ et $\mathbb{P}^1 \times \{t\}$ ; la proposition est dans ce cas évidente.


\SGAsectiontoc{3. Cohomologie des quadriques}{3}
3. \underline{Cohomologie des quadriques}

3.1. Soit $p : X \longrightarrow S$ une quadrique lisse de dimension $n$ sur $S$, et $\ell$ un nombre premier inversible sur $S$. On désignera par

\SGAvar{\begin{equation*}
\eta \in H^0(S, R^1 p_* \mathbb{Z}_\ell(1))
\end{equation*}
la ``classe de cohomologie d'une section hyperplane''. Localement sur $S$ pour la topologie étale, on a $P(X) \approx \mathbb{P}(V^*)$, et $\eta$ est alors par définition l'image dans $H^0(S, R^2 p_* \mathbb{Z}_\ell(1))$ de la première classe de Chern $c_1(\mathcal{O}(1)) \in H^2(X, \mathbb{Z}_\ell(1))$.

}{\begin{equation*}
\eta \in H^0(S, R^2 p_* \mathbb{Z}_\ell(1))
\end{equation*}
la ``classe de cohomologie d'une section hyperplane''. Localement sur $S$ pour la topologie étale, on a $P(X) \approx \mathbb{P}(V^*)$, et $\eta$ est alors par définition l'image dans $H^0(S, R^2 p_* \mathbb{Z}_\ell(1))$ de la première classe de Chern $c_1(\mathcal{O}(1)) \in H^2(X, \mathbb{Z}_\ell(1))$.

}{F7II-XII-015}3.2. Supposons que $n = 2m$ . Une génératrice $D$ de $X$ définit alors une classe de cohomologie $c\ell'(D)$ dans $H^n(X, \mathbb{Z}_\ell(m))$, d'image $c\ell(D)$ dans $H^0(S,R^n p_* \mathbb{Z}_\ell(m))$. Il résulte de 2.8 que cette dernière ne dépend que de la section $e(D)$ de $Z(X)$, d'où un morphisme de faisceaux étales sur $S$

\begin{equation}\tag{3.2.1}
c\ell : Z(X) \longrightarrow R^n p_* \mathbb{Z}_\ell(m) \quad .
\end{equation}
\underline{Théorème} 3.3. \underline{Soit} $p : X \longrightarrow S$ \underline{une quadrique lisse de dimension} $n$ \underline{sur} $S$ \underline{et} $\ell$ \underline{un nombre premier inversible sur} $S$

(i) \quad $R^{2i+1} p_* \mathbb{Z}_\ell = 0$ 


\SGApage{II}{76}

(ii) \underline{Si} $0 \leqslant 2i < n$, (resp. si $n< 2i \leqslant 2n$) \underline{alors} $R^{2i}p_{*} \mathbb{Z}_{\ell}(i)$ \underline{est canoniquement isomorphe au faisceau constant} $\mathbb{Z}_{\ell}$, \underline{et est engendré par} $\eta^{i}$ (resp. \underline{par} $\eta^{i}/2$)

(iii) \underline{Si} $n = 2m$, \underline{l'application déduite de} (3.2.1)

\SGAvar{\begin{equation}\tag{3.2.1}
\mathbb{Z}_{\ell}^{Z(X)} \longrightarrow R^{n}p_{*} \mathbb{Z}_{\ell}(m)
\end{equation}
\underline{est un isomorphisme}, \underline{et si} $\alpha$ \underline{et} $\beta$ \underline{sont deux sections disjointes de} $Z(X)$ (\underline{supposé isomorphe à} $S \amalg S$ ; c' \underline{est localement le cas}), \underline{on a}

}{\begin{equation}\tag{3.3.1}
\mathbb{Z}_{\ell}^{Z(X)} \longrightarrow R^{n}p_{*} \mathbb{Z}_{\ell}(m)
\end{equation}
\underline{est un isomorphisme}, \underline{et si} $\alpha$ \underline{et} $\beta$ \underline{sont deux sections disjointes de} $Z(X)$ (\underline{supposé isomorphe à} $S \amalg S$ ; c' \underline{est localement le cas}), \underline{on a}

}{F7II-XII-016}\begin{equation*}\tag{a}
\eta^{m} = c\ell(\alpha) + c\ell(\beta)
\end{equation*}
\begin{equation*}\tag{b}
\left\{ \begin{array}{l} \underline{\mathrm{pour}} \quad m \quad \underline{\mathrm{pair}} \; : \; \mathrm{Tr}(c\ell(\alpha)^{2}) = \mathrm{Tr}(c\ell(\beta)^{2}) = 1, \; c\ell(\alpha) \; c\ell(\beta) = 0 \\[2ex] \underline{\mathrm{pour}} \quad m \; \underline{\mathrm{impair}} \; : \; c\ell(\alpha)^{2} = c\ell(\beta)^{2} = 0 \; , \; \mathrm{Tr}(c\ell(\alpha).c\ell(\beta)) = 1. \end{array} \right.
\end{equation*}
\begin{equation*}\tag{c}
\eta.(c\ell(\alpha) - c\ell(\beta)) = 0.
\end{equation*}
Les assertions (i) et (ii) ne sont mises que pour mémoire (XI 1.6, complété par XI 2.6 lorsque $2i + 1 = n$ dans (i)).

Il suffit de vérifier (iii) lorsque $S$ est le spectre d'un corps algébriquement clos $k$ (il suffirait même de prendre $S = \mathrm{Spec}(\mathbb{C})$).

\SGAvar{\underline{Formule} (a). On peut supposer que $X$ est la quadrique d'équation $\sum\limits_{0}^{m} x_{i}x_{i+m+1} = 0$ dans $\mathbb{P}^{m+1}(k)$. La classe $\eta^{m}$ est la classe de cohomologie de la section plane d'équation

}{\underline{Formule} (a). On peut supposer que $X$ est la quadrique d'équation $\sum\limits_{0}^{m} x_{i}x_{i+m+1} = 0$ dans $\mathbb{P}^{2m+1}(k)$. La classe $\eta^{m}$ est la classe de cohomologie de la section plane d'équation

}{F7II-XII-017}\begin{equation*}
x_{i} = 0 \qquad (0 \leqslant i < m) \qquad ,
\end{equation*}
et ce cycle est somme des génératrices $D_{1}$ et $D_{2}$ d'équation
\SGApage{II}{77}
$D_{1}$ : $x_{i} = 0$ $(0 \leqslant i \leqslant m)$ \quad et

$D_{2}$ : $x_{i} = 0$ $(0 \leqslant i < m)$ et $x_{2m+1} = 0$ .

D'après 1.12, $e(D_{1}) \neq e(D_{2})$, de sorte que

\begin{equation*}
\eta^{m} = c\ell(D_{1}) + c\ell(D_{2}) = c\ell(\alpha) + c\ell(\beta) \qquad .
\end{equation*}
\underline{Formule} (b). Si des génératrices $D_{1}$ et $D_{2}$ sont disjointes (resp. se coupent transversalement en un point) on a $c\ell(D_{1}).c\ell(D_{2}) = 0$ (resp. $\mathrm{Tr}(c\ell(D_{1}).c\ell(D_{2})) = 1$). D'après 1.12, si $m$ est pair, on a $e(D_{1}) \neq e(D_{2})$ (resp. $e(D_{1}) = e(D_{2})$) ; si $m$ est impair, on a $e(D_{1}) = e(D_{2})$ (resp. $e(D_{1}) \neq e(D_{2})$). Les formules (b) en résultent.

\SGAvar{\underline{Formule} (c). Evident si $n = 0$. Sinon, il suffit d'après (ii) de prouver que $\eta^{m}.(c\ell(\alpha)-c\ell(\beta) = 0$. En effet, $\mathrm{Tr}(\eta^{m}(c\ell(\alpha)-c\ell(\beta)) = \mathrm{Tr}(c\ell(\alpha)^{2}-c\ell(\beta)^{2})=0$.

}{\underline{Formule} (c). Evident si $n = 0$. Sinon, il suffit d'après (ii) de prouver que $\eta^{m}.(c\ell(\alpha)-c\ell(\beta)) = 0$. En effet, $\mathrm{Tr}(\eta^{m}(c\ell(\alpha)-c\ell(\beta))) = \mathrm{Tr}(c\ell(\alpha)^{2}-c\ell(\beta)^{2})=0$.

}{F7II-XII-018}\underline{Vérification}.

\SGAvar{$m$ pair : $2 = \mathrm{tr}(\eta^{2}) = \mathrm{tr}(c\ell(\alpha^{2})+ 2c\ell(\alpha)c\ell(\beta) + c\ell(\beta)^{2}) = 1 + 1$

$m$ impair : $2 = \mathrm{tr}(\eta^{2}) = \mathrm{Tr}(c\ell(\alpha)^{2} + 2c\ell(\alpha)c\ell(\beta) + c\ell(\beta)^{2}) = 2$ .

}{$m$ pair : $2 = \mathrm{tr}(\eta^{2m}) = \mathrm{tr}(c\ell(\alpha)^{2}+ 2c\ell(\alpha)c\ell(\beta) + c\ell(\beta)^{2}) = 1 + 1$

$m$ impair : $2 = \mathrm{tr}(\eta^{2m}) = \mathrm{Tr}(c\ell(\alpha)^{2} + 2c\ell(\alpha)c\ell(\beta) + c\ell(\beta)^{2}) = 2$ .

}{F7II-XII-019}Que 3.3.1 soit un isomorphisme résulte maintenant de ce que les formes bilinéaires à deux variables

\begin{equation*}
B(X,Y) \; = \; x_{1}y_{2} + x_{2}y_{1} \qquad\qquad \mathrm{et}
\end{equation*}
\begin{equation*}
B(x,y) \; = \; x_{1}y_{1} + x_{2}y_{2}
\end{equation*}
sont de discriminant $\pm 1$ .


\SGApage{II}{78}

\underline{Vérification} 3.4. Soit $X$ une quadrique lisse de dimension $n$ sur un corps fini $\mathbb{F}_q$. Compte tenu de 3.3, la formule des traces de Lefschetz donne la formule classique

\begin{equation*}
\# X(\mathbb{F}_q) \;=\; \left\{ \begin{array}{ll} \displaystyle\sum_{0}^{\,n} q^{i} & \mbox{pour } n \mbox{ impair} \\[3ex] \displaystyle\sum_{0}^{\,n} q^{i} + \varepsilon\, q^{m} & \mbox{pour } n = 2\,m \quad (\varepsilon = \pm\, 1) \end{array} \right. \qquad .
\end{equation*}
Pour $n$ pair, on a $\varepsilon = 1$ si $X$ est déployée, i.e. \quad a une génératrice rationnelle, et $\varepsilon = -1$ dans le cas contraire.

3.5. Si $X$ est une quadrique lisse sur $S$ de dimension $n = 2m$, on appellera \underline{partie primitive} de $R^{n}f_{*}\mathbb{Z}_{\ell}(m)$ l'orthogonal de $\eta^{m}$ (avec les notations de 3.3, cet orthogonal est engendré par $c\ell(\alpha)-c\ell(\beta)$). On appellera \underline{quotient primitif} de $R^{n}f_{*}\mathbb{Z}_{\ell}(m)$ le quotient de $R^{n}f_{*}\mathbb{Z}_{\ell}(m)$ par $\mathbb{Z}_{\ell}\eta^{m}$. On définit partie et quotient primitifs de $R^{n}f_{*}\mathbb{Z}_{\ell}$ en tensorisant avec $\mathbb{Z}_{\ell}(-m)$.

3.6. Soient $V$ un fibré vectoriel localement libre de rang $n + 2$ sur $S$ et $X$ la quadrique lisse sur $S$ définie par une forme quadratique ordinaire $Q$ sur $V$. Soit de plus $H$ un hyperplan de $\mathbb{P}(V^{*})$ qui coupe transversalement $X$, et $X^{o} = X - H$. Les schémas $X$ et $Y = X \cap H$ sont des quadriques lisses sur $S$, de dimension $n$ et $n - 1$. Ceci permet d'utiliser 3.3 pour calculer la cohomologie de $X^{o} = X - Y$.


\SGApage{II}{79}

\begin{equation}\tag{3.6.1}
\begin{tikzcd}[column sep=huge, row sep=huge]
X^{o} \arrow[r, hook] \arrow[rd, "f"'] \& X \arrow[d, "p"] \& Y \arrow[l, hook] \arrow[ld, "q"] \\
 {} \& S \&
{} \end{tikzcd}\quad .
\end{equation}









Si $n = 0$, alors $Y = \emptyset$ et $X = X^{o}$. Supposons $n > 0$. On dispose de deux suites exactes longues duales l'une de l'autre

\SGAvar{\begin{equation}\tag{3.6.2}
\ldots \ \stackrel{\partial}{\longrightarrow}\ R^{i}f_{!}\mathbb{Z}_{\ell} \longrightarrow R^{i}p_{*}\mathbb{Z}_{\ell} \ \stackrel{r_{i}}{\longrightarrow}\ R^{i}q_{*}\mathbb{Z}_{\ell} \longrightarrow \ldots
\end{equation}
\begin{equation}\tag{3.6.3}
\ldots \ \stackrel{\partial}{\longrightarrow}\ R^{i-2}q_{!}\mathbb{Z}_{\ell}(-1) \longrightarrow R^{i}p_{*}\mathbb{Z}_{\ell} \longrightarrow R^{i}f_{*}\mathbb{Z}_{\ell} \longrightarrow \ldots \qquad .
\end{equation}
Pour $n$ pair, l'homomorphisme de restriction $r_{i}$ est un isomorphisme pour $i \neq n,2n$ ; pour $n$ impair, c'est un isomorphisme pour $i \neq n-1$. Pour $n = 2m$, et $\alpha$ la classe de cohomologie d'une génératrice, $r_{n}(\alpha) = \frac{1}{2}\,\eta^{m}$, de sorte que $r_{n}$ est surjectif de noyau la partie primitive de $R^{n}p_{*}\mathbb{Z}_{\ell}$. Pour $r = 2m + 1$, $r_{2m}(\eta^{m}) = \eta^{m}$, de sorte que $r_{2m}$ est injectif de conoyau le quotient primitif de $R^{2m}q_{*}\mathbb{Z}_{\ell}$. On en déduit que $R^{i}f_{!}\mathbb{Z}_{\ell} = 0$ pour $i \neq n,2n$, et que $R^{i}f_{!}\mathbb{Z}_{\ell}$ est un $\mathbb{Z}_{\ell}$-faisceau constant tordu libre de rang 1, pour $i = n,2n$. Par dualité, ou à l'aide de 3.5.3, on montre de même que $R^{i}f_{*}\mathbb{Z}_{\ell} = 0$ pour $i \neq 0,n$, que $f_{*}\mathbb{Z}_{\ell} = \mathbb{Z}_{\ell}$, et que $R^{n}f_{*}\mathbb{Z}_{\ell}$ est de rang un.

}{\begin{equation}\tag{3.6.2}
\ldots \ \stackrel{\partial}{\longrightarrow}\ R^{i}f_{!}\mathbb{Z}_{\ell} \longrightarrow R^{i}p_{*}\mathbb{Z}_{\ell} \ \stackrel{r_{i}}{\longrightarrow}\ R^{i}q_{*}\mathbb{Z}_{\ell} \longrightarrow \ldots
\end{equation}
\begin{equation}\tag{3.6.3}
\ldots \ \stackrel{\partial}{\longrightarrow}\ R^{i-2}q_{!}\mathbb{Z}_{\ell}(-1) \longrightarrow R^{i}p_{*}\mathbb{Z}_{\ell} \longrightarrow R^{i}f_{*}\mathbb{Z}_{\ell} \longrightarrow \ldots \qquad .
\end{equation}
Pour $n$ pair, l'homomorphisme de restriction $r_{i}$ est un isomorphisme pour $i \neq n,2n$ ; pour $n$ impair, c'est un isomorphisme pour $i \neq n-1,2n$. Pour $n = 2m$, et $\alpha$ la classe de cohomologie d'une génératrice, $r_{n}(\alpha) = \frac{1}{2}\,\eta^{m}$, de sorte que $r_{n}$ est surjectif de noyau la partie primitive de $R^{n}p_{*}\mathbb{Z}_{\ell}$. Pour $n = 2m + 1$, $r_{2m}(\eta^{m}) = \eta^{m}$, de sorte que $r_{2m}$ est injectif de conoyau le quotient primitif de $R^{2m}q_{*}\mathbb{Z}_{\ell}$. On en déduit que $R^{i}f_{!}\mathbb{Z}_{\ell} = 0$ pour $i \neq n,2n$, et que $R^{i}f_{!}\mathbb{Z}_{\ell}$ est un $\mathbb{Z}_{\ell}$-faisceau constant tordu libre de rang 1, pour $i = n,2n$. Par dualité, ou à l'aide de 3.6.3, on montre de même que $R^{i}f_{*}\mathbb{Z}_{\ell} = 0$ pour $i \neq 0,n$, que $f_{*}\mathbb{Z}_{\ell} = \mathbb{Z}_{\ell}$, et que $R^{n}f_{*}\mathbb{Z}_{\ell}$ est de rang un.

}{F7II-XII-020}Supposons que $n = 2m$ et reprenons les notations de 3.3.


\SGApage{II}{80}

\begin{center}
\begin{tikzcd}[column sep=large, row sep=large]
0 \arrow[r] \& R^n f_{!} \mathbb{Z}_\ell(m) \arrow[r] \arrow[d, "\varphi"] \& R^n p_{*} \mathbb{Z}_\ell(m) \arrow[r] \arrow[d, equal] \& R^n q_{*} \mathbb{Z}_\ell(m) \arrow[r] \& 0 \\
0 \& R^n f_{*} \mathbb{Z}_\ell(m) \arrow[l] \& R^n p_{*} \mathbb{Z}_\ell(m) \arrow[l] \& R^{n-2} q_{*} \mathbb{Z}_\ell(m-1) \arrow[l] \& 0 \arrow[l] \;.
\end{tikzcd}
\end{center}










Le faisceau $R^n f_{!} \mathbb{Z}_\ell(m)$ a alors deux générateurs opposés naturels $\delta$, d'image $\pm (c\ell(\alpha) - c\ell(\beta))$ dans $R^n p_{*} \mathbb{Z}_\ell(m)$. On a (3.3 (iii) b)

\begin{equation*}
\mathrm{Tr}(\delta^2) = (c\ell(\alpha) - c\ell(\beta))^2 = (-1)^m . 2
\end{equation*}
et $\pm \varphi(\delta)$ est le double des générateurs naturels duaux $\pm \delta'$ de $R^n f_{*} \mathbb{Z}_\ell(m)$, à savoir les images (opposées) de $c\ell(\alpha)$ et $c\ell(\beta)$.

Supposons que $n = 2m + 1$ et appliquons les notations de 3.3 à $Y$.

\SGAvar{\begin{center}
\begin{tikzcd}[column sep=large, row sep=large]
0 \arrow[r] \& R^{2m} p_{*} \mathbb{Z}_\ell(m) \arrow[r] \& R^{2m} q_{*} \mathbb{Z}_\ell(m) \arrow[r, "\partial"] \& R^n f_{!} \mathbb{Z}_\ell(m) \arrow[r] \arrow[d, "\varphi"] \& 0 \\
{} \& R^{2m+2} p_{*} \mathbb{Z}_\ell(m) \& R^{2m} q_{*} \mathbb{Z}_\ell(m-1) \arrow[l] \& R^n f_{*} \mathbb{Z}_\ell(m) \arrow[l] \& 0 \arrow[l] \;.
\end{tikzcd}
\end{center}










}{\begin{center}
\begin{tikzcd}[column sep=large, row sep=large]
0 \arrow[r] \& R^{2m} p_{*} \mathbb{Z}_\ell(m) \arrow[r] \& R^{2m} q_{*} \mathbb{Z}_\ell(m) \arrow[r, "\partial"] \& R^n f_{!} \mathbb{Z}_\ell(m) \arrow[r] \arrow[d, "\varphi"] \& 0 \\
{} \& R^{2m+2} p_{*} \mathbb{Z}_\ell(m+1) \& R^{2m} q_{*} \mathbb{Z}_\ell(m-1) \arrow[l] \& R^n f_{*} \mathbb{Z}_\ell(m) \arrow[l] \& 0 \arrow[l] \;.
\end{tikzcd}
\end{center}










}{F7II-XII-021}Le faisceau $R^n f_{!} \mathbb{Z}_\ell(m)$ a alors deux générateurs naturels opposés $\pm \delta$ : $\partial c\ell(\alpha)$ et $\partial c\ell(\beta)$. On a

\begin{equation*}
\delta . \varphi(\delta) = \delta^2 = (\partial c\ell(\alpha))^2 = \partial(c\ell(\alpha) . \partial c\ell(\alpha)) = 0
\end{equation*}
de sorte que $\varphi(\delta) = 0$ . D'autre part, $R^n f_{*} \mathbb{Z}_\ell(m+1)$ a deux générateurs naturels opposés $\pm \delta'$, d'image $\pm(c\ell(\alpha) - c\ell(\beta))$ dans $R^{2m} q_{*} \mathbb{Z}_\ell(m)$, et $\mathrm{Tr}(\delta . \delta') = \pm 1$.


\SGApage{II}{81}

Ces résultats sont rassemblés dans la table suivante

\underline{Table} 3.7.  \underline{Cohomologie des quadriques affines de dimension} $n$, $X^o = X - Y$.

A.  La cohomologie est sans torsion. Les nombres de Betti non nuls sont :

cohomologie sans support  :  $b_o = b_n = 1$

cohomologie à support propre :$b_{2n} = b_n = 1$

\underline{sauf} pour $n = 0$  :  $b_o = 2$.

B.  $n = 2m > 0$.  $R^n f_{!} \mathbb{Z}_\ell(m)$ = partie primitive de $R^{2m} p_{*} \mathbb{Z}_\ell(m)$

$R^n f_{*} \mathbb{Z}_\ell(m)$ = quotient primitif de $R^{2m} p_{*} \mathbb{Z}_\ell(m)$

$n = 2m + 1$   $R^n f_{!} \mathbb{Z}_\ell(m)$ = quotient primitif de $R^{2m} q_{*} \mathbb{Z}_\ell(m)$

$R^n f_{*} \mathbb{Z}_\ell(m+1)$ = partie primitive de $R^{2m} q_{*} \mathbb{Z}_\ell(m)$  .

Ces groupes ont, localement, des générateurs naturels définis au signe près, notés $\delta$ pour les groupes de cohomologie à support propre et $\delta'$ pour les autres.

\SGAvar{C.  $\varphi : R^n f_{!} \mathbb{Z}_\ell \longrightarrow R^n f_{*} \mathbb{Z}_\ell$ est nul pour $n$ impair, envoyè $\pm \delta$ sur $\pm 2 \delta'$ pour $n$ pair $> 0$. On a

}{C.  $\varphi : R^n f_{!} \mathbb{Z}_\ell \longrightarrow R^n f_{*} \mathbb{Z}_\ell$ est nul pour $n$ impair, envoie $\pm \delta$ sur $\pm 2 \delta'$ pour $n$ pair $> 0$. On a

}{F7II-XII-022}\begin{equation*}
\mathrm{Tr}(\delta\delta') = \pm 1
\end{equation*}
\begin{equation*}
\mathrm{Tr}(\delta^2) = \begin{cases} 0 & \text{si} \quad n = 2m + 1 \\ (-1)^m . 2 & \text{si} \ n = 2m \ . \end{cases}
\end{equation*}
\underline{Vérification} 3.8.  La quadrique affine complexe $S$ dans $\mathbb{C}^n$ d'équation $\Sigma \; z_i^2 = 1$ est difféomorphe au fibré tangent à la sphère réelle $S \cap \mathbb{R}^n$ .  Ceci concorde avec A et C.

\underline{BIBLIOGRAPHIE}

{[}1]  D. Mumford.  Geometric invariant theory. Ergebnisse \underline{34}, Springer Verlag 1965.


\selectlanguage{french}

\ifnum\SGAform=1
\clearpage\section*{Apparat critique}\small
Les lectures ci-dessous sont données dans le TeX exact des témoins, afin de préserver les distinctions de notation. Les indices décimaux renvoient aux différences individuelles du passage marqué.

\par\medskip\noindent\hypertarget{app-F7II-XII-001.1}{}\textbf{F7II-XII-001.1}\quad p. 63; misprint\_pass.\par
\noindent Imprimé : {\ttfamily\footnotesize card\allowbreak{}}\par
\noindent Corrigé : {\ttfamily\footnotesize car\allowbreak{}}\par
\noindent Source : 7II:93. \nolinkurl{fixes/FIXLOG_7II_CURRENT.jsonl#535}; \nolinkurl{fixes/MISPRINTS_7II.json#93}\par

\par\medskip\noindent\hypertarget{app-F7II-XII-002.1}{}\textbf{F7II-XII-002.1}\quad p. 64; designated-witness attribution.\par
\noindent Imprimé : {\ttfamily\footnotesize -1\allowbreak{}}\par
\noindent Corrigé : \emph{(vide)}\par
\noindent Source : WITNESS-F7II-XII-002.1. \nolinkurl{fr/sga7ii/diplomatic/expose_XII_body.tex}; \nolinkurl{fr/sga7ii/current/expose_XII_body.tex}\par

\par\medskip\noindent\hypertarget{app-F7II-XII-002.2}{}\textbf{F7II-XII-002.2}\quad p. 64; designated-witness attribution.\par
\noindent Imprimé : {\ttfamily\footnotesize -1\allowbreak{}}\par
\noindent Corrigé : \emph{(vide)}\par
\noindent Source : WITNESS-F7II-XII-002.2. \nolinkurl{fr/sga7ii/diplomatic/expose_XII_body.tex}; \nolinkurl{fr/sga7ii/current/expose_XII_body.tex}\par

\par\medskip\noindent\hypertarget{app-F7II-XII-002.3}{}\textbf{F7II-XII-002.3}\quad p. 64; misprint\_pass.\par
\noindent Imprimé : \emph{(vide)}\par
\noindent Corrigé : {\ttfamily\footnotesize )\allowbreak{}}\par
\noindent Source : 7II:95. \nolinkurl{fixes/FIXLOG_7II_CURRENT.jsonl#536}; \nolinkurl{fixes/MISPRINTS_7II.json#95}\par

\par\medskip\noindent\hypertarget{app-F7II-XII-002.4}{}\textbf{F7II-XII-002.4}\quad p. 64; printed misprint.\par
\noindent Imprimé : \emph{(vide)}\par
\noindent Corrigé : {\ttfamily\footnotesize '\allowbreak{}}\par
\noindent Source : 7II:96. \nolinkurl{fixes/MISPRINTS_7II.json#96}\par

\par\medskip\noindent\hypertarget{app-F7II-XII-003.1}{}\textbf{F7II-XII-003.1}\quad p. 65; misprint\_pass.\par
\noindent Imprimé : {\ttfamily\footnotesize E\allowbreak{}}\par
\noindent Corrigé : {\ttfamily\footnotesize V\allowbreak{}}\par
\noindent Source : 7II:99. \nolinkurl{fixes/FIXLOG_7II_CURRENT.jsonl#538}; \nolinkurl{fixes/MISPRINTS_7II.json#99}\par

\par\medskip\noindent\hypertarget{app-F7II-XII-003.2}{}\textbf{F7II-XII-003.2}\quad p. 65; misprint\_pass.\par
\noindent Imprimé : {\ttfamily\footnotesize E\allowbreak{}}\par
\noindent Corrigé : {\ttfamily\footnotesize V\allowbreak{}}\par
\noindent Source : 7II:99. \nolinkurl{fixes/FIXLOG_7II_CURRENT.jsonl#538}; \nolinkurl{fixes/MISPRINTS_7II.json#99}\par

\par\medskip\noindent\hypertarget{app-F7II-XII-003.3}{}\textbf{F7II-XII-003.3}\quad p. 65; misprint\_pass.\par
\noindent Imprimé : {\ttfamily\footnotesize E\allowbreak{}}\par
\noindent Corrigé : {\ttfamily\footnotesize V\allowbreak{}}\par
\noindent Source : 7II:99. \nolinkurl{fixes/FIXLOG_7II_CURRENT.jsonl#538}; \nolinkurl{fixes/MISPRINTS_7II.json#99}\par

\par\medskip\noindent\hypertarget{app-F7II-XII-003.4}{}\textbf{F7II-XII-003.4}\quad p. 65; misprint\_pass.\par
\noindent Imprimé : {\ttfamily\footnotesize 2\allowbreak{}}\par
\noindent Corrigé : {\ttfamily\footnotesize 3\allowbreak{}}\par
\noindent Source : 7II:97. \nolinkurl{fixes/FIXLOG_7II_CURRENT.jsonl#537}\par

\par\medskip\noindent\hypertarget{app-F7II-XII-004.1}{}\textbf{F7II-XII-004.1}\quad p. 65; printed misprint.\par
\noindent Imprimé : \emph{(vide)}\par
\noindent Corrigé : {\ttfamily\footnotesize -1\allowbreak{}}\par
\noindent Source : 7II:98. \nolinkurl{fixes/MISPRINTS_7II.json#98}\par

\par\medskip\noindent\hypertarget{app-F7II-XII-004.2}{}\textbf{F7II-XII-004.2}\quad p. 65; printed misprint.\par
\noindent Imprimé : {\ttfamily\footnotesize o\allowbreak{}}\par
\noindent Corrigé : {\ttfamily\footnotesize 1\allowbreak{}}\par
\noindent Source : 7II:98. \nolinkurl{fixes/MISPRINTS_7II.json#98}\par

\par\medskip\noindent\hypertarget{app-F7II-XII-004.3}{}\textbf{F7II-XII-004.3}\quad p. 66; misprint\_pass.\par
\noindent Imprimé : {\ttfamily\footnotesize gradué\allowbreak{}}\par
\noindent Corrigé : {\ttfamily\footnotesize graduée\allowbreak{}}\par
\noindent Source : 7II:100. \nolinkurl{fixes/FIXLOG_7II_CURRENT.jsonl#539}; \nolinkurl{fixes/MISPRINTS_7II.json#100}\par

\par\medskip\noindent\hypertarget{app-F7II-XII-004.4}{}\textbf{F7II-XII-004.4}\quad p. 66; misprint\_pass.\par
\noindent Imprimé : {\ttfamily\footnotesize algèbre\allowbreak{}}\par
\noindent Corrigé : {\ttfamily\footnotesize algèbres\allowbreak{}}\par
\noindent Source : 7II:101. \nolinkurl{fixes/FIXLOG_7II_CURRENT.jsonl#540}; \nolinkurl{fixes/MISPRINTS_7II.json#101}\par

\par\medskip\noindent\hypertarget{app-F7II-XII-004.5}{}\textbf{F7II-XII-004.5}\quad p. 66; misprint\_in\_print.\par
\noindent Imprimé : \emph{(vide)}\par
\noindent Corrigé : {\ttfamily\footnotesize )\allowbreak{}}\par
\noindent Source : FIX-7II-D-0046. \nolinkurl{fixes/FIXLOG_7II.jsonl#46}\par

\par\medskip\noindent\hypertarget{app-F7II-XII-004.6}{}\textbf{F7II-XII-004.6}\quad p. 66; misprint\_in\_print.\par
\noindent Imprimé : \emph{(vide)}\par
\noindent Corrigé : {\ttfamily\footnotesize )\allowbreak{}}\par
\noindent Source : FIX-7II-D-0047. \nolinkurl{fixes/FIXLOG_7II.jsonl#47}\par

\par\medskip\noindent\hypertarget{app-F7II-XII-004.7}{}\textbf{F7II-XII-004.7}\quad p. 66; misprint\_pass.\par
\noindent Imprimé : {\ttfamily\footnotesize 2\allowbreak{}}\par
\noindent Corrigé : {\ttfamily\footnotesize 1\allowbreak{}}\par
\noindent Source : 7II:104. \nolinkurl{fixes/FIXLOG_7II_CURRENT.jsonl#541}; \nolinkurl{fixes/MISPRINTS_7II.json#104}\par

\par\medskip\noindent\hypertarget{app-F7II-XII-004.8}{}\textbf{F7II-XII-004.8}\quad p. 66; misprint\_pass.\par
\noindent Imprimé : {\ttfamily\footnotesize 1\allowbreak{}}\par
\noindent Corrigé : {\ttfamily\footnotesize 2\allowbreak{}}\par
\noindent Source : 7II:104. \nolinkurl{fixes/FIXLOG_7II_CURRENT.jsonl#541}; \nolinkurl{fixes/MISPRINTS_7II.json#104}\par

\par\medskip\noindent\hypertarget{app-F7II-XII-005.1}{}\textbf{F7II-XII-005.1}\quad p. 67; misprint\_pass.\par
\noindent Imprimé : {\ttfamily\footnotesize singulier\allowbreak{}}\par
\noindent Corrigé : {\ttfamily\footnotesize singuliers\allowbreak{}}\par
\noindent Source : 7II:105. \nolinkurl{fixes/FIXLOG_7II_CURRENT.jsonl#542}; \nolinkurl{fixes/MISPRINTS_7II.json#105}\par

\par\medskip\noindent\hypertarget{app-F7II-XII-005.2}{}\textbf{F7II-XII-005.2}\quad p. 67; misprint\_pass.\par
\noindent Imprimé : {\ttfamily\footnotesize supplémentaire\allowbreak{}}\par
\noindent Corrigé : {\ttfamily\footnotesize supplémentaires\allowbreak{}}\par
\noindent Source : 7II:106. \nolinkurl{fixes/FIXLOG_7II_CURRENT.jsonl#543}; \nolinkurl{fixes/MISPRINTS_7II.json#106}\par

\par\medskip\noindent\hypertarget{app-F7II-XII-005.3}{}\textbf{F7II-XII-005.3}\quad p. 67; misprint\_pass.\par
\noindent Imprimé : {\ttfamily\footnotesize singulier\allowbreak{}}\par
\noindent Corrigé : {\ttfamily\footnotesize singuliers\allowbreak{}}\par
\noindent Source : 7II:106. \nolinkurl{fixes/FIXLOG_7II_CURRENT.jsonl#543}; \nolinkurl{fixes/MISPRINTS_7II.json#106}\par

\par\medskip\noindent\hypertarget{app-F7II-XII-005.4}{}\textbf{F7II-XII-005.4}\quad p. 67; printed misprint.\par
\noindent Imprimé : {\ttfamily\footnotesize dont\allowbreak{}}\par
\noindent Corrigé : {\ttfamily\footnotesize donc\allowbreak{}}\par
\noindent Source : 7II:107. \nolinkurl{fixes/MISPRINTS_7II.json#107}\par

\par\medskip\noindent\hypertarget{app-F7II-XII-005.5}{}\textbf{F7II-XII-005.5}\quad p. 67; designated-witness attribution.\par
\noindent Imprimé : \emph{(vide)}\par
\noindent Corrigé : {\ttfamily\footnotesize \$S\$,\allowbreak{} et\allowbreak{} par\allowbreak{} suite\allowbreak{} par\allowbreak{} }\par
\noindent Source : WITNESS-F7II-XII-005.5. \nolinkurl{fr/sga7ii/diplomatic/expose_XII_body.tex}; \nolinkurl{fr/sga7ii/current/expose_XII_body.tex}\par

\par\medskip\noindent\hypertarget{app-F7II-XII-005.6}{}\textbf{F7II-XII-005.6}\quad p. 67; printed misprint.\par
\noindent Imprimé : {\ttfamily\footnotesize \textbackslash{}mathbb\allowbreak{}}\par
\noindent Corrigé : {\ttfamily\footnotesize \textbackslash{}mathrm\allowbreak{}}\par
\noindent Source : 7II:108. \nolinkurl{fixes/MISPRINTS_7II.json#108}\par

\par\medskip\noindent\hypertarget{app-F7II-XII-005.7}{}\textbf{F7II-XII-005.7}\quad p. 67; printed misprint.\par
\noindent Imprimé : {\ttfamily\footnotesize E\allowbreak{}}\par
\noindent Corrigé : {\ttfamily\footnotesize Spec\allowbreak{}}\par
\noindent Source : 7II:108. \nolinkurl{fixes/MISPRINTS_7II.json#108}\par

\par\medskip\noindent\hypertarget{app-F7II-XII-005.8}{}\textbf{F7II-XII-005.8}\quad p. 67; printed misprint.\par
\noindent Imprimé : {\ttfamily\footnotesize \textbackslash{}overset\{2\}\{\textbackslash{}wedge\}\allowbreak{} W'\allowbreak{}}\par
\noindent Corrigé : {\ttfamily\footnotesize Z(Q)\allowbreak{}}\par
\noindent Source : 7II:108. \nolinkurl{fixes/MISPRINTS_7II.json#108}\par

\par\medskip\noindent\hypertarget{app-F7II-XII-006.1}{}\textbf{F7II-XII-006.1}\quad p. 67; printed misprint.\par
\noindent Imprimé : {\ttfamily\footnotesize 6\allowbreak{}}\par
\noindent Corrigé : {\ttfamily\footnotesize 3\allowbreak{}}\par
\noindent Source : 7II:109. \nolinkurl{fixes/MISPRINTS_7II.json#109}\par

\par\medskip\noindent\hypertarget{app-F7II-XII-007.1}{}\textbf{F7II-XII-007.1}\quad p. 68; misprint\_pass.\par
\noindent Imprimé : {\ttfamily\footnotesize gradués\allowbreak{}}\par
\noindent Corrigé : {\ttfamily\footnotesize graduées\allowbreak{}}\par
\noindent Source : 7II:110. \nolinkurl{fixes/FIXLOG_7II_CURRENT.jsonl#545}; \nolinkurl{fixes/MISPRINTS_7II.json#110}\par

\par\medskip\noindent\hypertarget{app-F7II-XII-008.1}{}\textbf{F7II-XII-008.1}\quad p. 68; printed misprint.\par
\noindent Imprimé : {\ttfamily\footnotesize gradués\allowbreak{}}\par
\noindent Corrigé : {\ttfamily\footnotesize graduées\allowbreak{}}\par
\noindent Source : 7II:111. \nolinkurl{fixes/MISPRINTS_7II.json#111}\par

\par\medskip\noindent\hypertarget{app-F7II-XII-009.1}{}\textbf{F7II-XII-009.1}\quad p. 69; misprint\_pass.\par
\noindent Imprimé : {\ttfamily\footnotesize ragarder\allowbreak{}}\par
\noindent Corrigé : {\ttfamily\footnotesize regarder\allowbreak{}}\par
\noindent Source : 7II:112. \nolinkurl{fixes/FIXLOG_7II_CURRENT.jsonl#547}; \nolinkurl{fixes/MISPRINTS_7II.json#112}\par

\par\medskip\noindent\hypertarget{app-F7II-XII-009.2}{}\textbf{F7II-XII-009.2}\quad p. 69; printed misprint.\par
\noindent Imprimé : {\ttfamily\footnotesize torseurs\allowbreak{}}\par
\noindent Corrigé : {\ttfamily\footnotesize torseur\allowbreak{}}\par
\noindent Source : 7II:113. \nolinkurl{fixes/MISPRINTS_7II.json#113}\par

\par\medskip\noindent\hypertarget{app-F7II-XII-009.3}{}\textbf{F7II-XII-009.3}\quad p. 69; misprint\_pass.\par
\noindent Imprimé : {\ttfamily\footnotesize \textbackslash{}rho\allowbreak{}}\par
\noindent Corrigé : {\ttfamily\footnotesize p\allowbreak{}}\par
\noindent Source : 7II:114. \nolinkurl{fixes/FIXLOG_7II_CURRENT.jsonl#549}; \nolinkurl{fixes/MISPRINTS_7II.json#114}\par

\par\medskip\noindent\hypertarget{app-F7II-XII-010.1}{}\textbf{F7II-XII-010.1}\quad p. 69; printed misprint.\par
\noindent Imprimé : {\ttfamily\footnotesize V\allowbreak{}}\par
\noindent Corrigé : {\ttfamily\footnotesize A\allowbreak{}}\par
\noindent Source : 7II:115. \nolinkurl{fixes/MISPRINTS_7II.json#115}\par

\par\medskip\noindent\hypertarget{app-F7II-XII-011.1}{}\textbf{F7II-XII-011.1}\quad p. 70; misprint\_pass.\par
\noindent Imprimé : {\ttfamily\footnotesize '\allowbreak{}}\par
\noindent Corrigé : \emph{(vide)}\par
\noindent Source : 7II:116. \nolinkurl{fixes/FIXLOG_7II_CURRENT.jsonl#550}; \nolinkurl{fixes/MISPRINTS_7II.json#116}\par

\par\medskip\noindent\hypertarget{app-F7II-XII-012.1}{}\textbf{F7II-XII-012.1}\quad p. 72; printed misprint.\par
\noindent Imprimé : {\ttfamily\footnotesize ii\allowbreak{}}\par
\noindent Corrigé : {\ttfamily\footnotesize iii\allowbreak{}}\par
\noindent Source : 7II:119. \nolinkurl{fixes/MISPRINTS_7II.json#119}\par

\par\medskip\noindent\hypertarget{app-F7II-XII-012.2}{}\textbf{F7II-XII-012.2}\quad p. 72; printed misprint.\par
\noindent Imprimé : {\ttfamily\footnotesize iii\allowbreak{}}\par
\noindent Corrigé : {\ttfamily\footnotesize ii\allowbreak{}}\par
\noindent Source : 7II:120. \nolinkurl{fixes/MISPRINTS_7II.json#120}\par

\par\medskip\noindent\hypertarget{app-F7II-XII-012.3}{}\textbf{F7II-XII-012.3}\quad p. 72; inherited corrected-layer difference.\par
\noindent Imprimé : {\ttfamily\footnotesize p\allowbreak{}}\par
\noindent Corrigé : {\ttfamily\footnotesize \textbackslash{}rho\allowbreak{}}\par
\noindent Source : ADOPT-F7II-XII-012.3. \nolinkurl{fr/sga7ii/diplomatic/expose_XII_body.tex}; \nolinkurl{fr/sga7ii/current/expose_XII_body.tex}\par

\par\medskip\noindent\hypertarget{app-F7II-XII-012.4}{}\textbf{F7II-XII-012.4}\quad p. 72; misprint\_pass.\par
\noindent Imprimé : {\ttfamily\footnotesize induisant\allowbreak{}}\par
\noindent Corrigé : {\ttfamily\footnotesize induisent\allowbreak{}}\par
\noindent Source : 7II:121. \nolinkurl{fixes/FIXLOG_7II_CURRENT.jsonl#551}; \nolinkurl{fixes/MISPRINTS_7II.json#121}\par

\par\medskip\noindent\hypertarget{app-F7II-XII-013.1}{}\textbf{F7II-XII-013.1}\quad p. 73; misprint\_in\_print.\par
\noindent Imprimé : {\ttfamily\footnotesize lequelle\allowbreak{}}\par
\noindent Corrigé : {\ttfamily\footnotesize laquelle\allowbreak{}}\par
\noindent Source : FIX-7II-D-0053. \nolinkurl{fixes/FIXLOG_7II.jsonl#53}\par

\par\medskip\noindent\hypertarget{app-F7II-XII-014.1}{}\textbf{F7II-XII-014.1}\quad p. 74; printed misprint.\par
\noindent Imprimé : {\ttfamily\footnotesize \textasciicircum{}*\allowbreak{}}\par
\noindent Corrigé : \emph{(vide)}\par
\noindent Source : 7II:124. \nolinkurl{fixes/MISPRINTS_7II.json#124}\par

\par\medskip\noindent\hypertarget{app-F7II-XII-014.2}{}\textbf{F7II-XII-014.2}\quad p. 74; printed misprint.\par
\noindent Imprimé : \emph{(vide)}\par
\noindent Corrigé : {\ttfamily\footnotesize \_1\allowbreak{}}\par
\noindent Source : 7II:125. \nolinkurl{fixes/MISPRINTS_7II.json#125}\par

\par\medskip\noindent\hypertarget{app-F7II-XII-014.3}{}\textbf{F7II-XII-014.3}\quad p. 74; printed misprint.\par
\noindent Imprimé : {\ttfamily\footnotesize '\allowbreak{}}\par
\noindent Corrigé : {\ttfamily\footnotesize \_2\allowbreak{}}\par
\noindent Source : 7II:125. \nolinkurl{fixes/MISPRINTS_7II.json#125}\par

\par\medskip\noindent\hypertarget{app-F7II-XII-014.4}{}\textbf{F7II-XII-014.4}\quad p. 74; printed misprint.\par
\noindent Imprimé : \emph{(vide)}\par
\noindent Corrigé : {\ttfamily\footnotesize \_1\allowbreak{}}\par
\noindent Source : 7II:125. \nolinkurl{fixes/MISPRINTS_7II.json#125}\par

\par\medskip\noindent\hypertarget{app-F7II-XII-014.5}{}\textbf{F7II-XII-014.5}\quad p. 74; printed misprint.\par
\noindent Imprimé : \emph{(vide)}\par
\noindent Corrigé : {\ttfamily\footnotesize \_1\allowbreak{}}\par
\noindent Source : 7II:125. \nolinkurl{fixes/MISPRINTS_7II.json#125}\par

\par\medskip\noindent\hypertarget{app-F7II-XII-014.6}{}\textbf{F7II-XII-014.6}\quad p. 74; printed misprint.\par
\noindent Imprimé : {\ttfamily\footnotesize '\allowbreak{}}\par
\noindent Corrigé : {\ttfamily\footnotesize \_2\allowbreak{}}\par
\noindent Source : 7II:125. \nolinkurl{fixes/MISPRINTS_7II.json#125}\par

\par\medskip\noindent\hypertarget{app-F7II-XII-015.1}{}\textbf{F7II-XII-015.1}\quad p. 75; printed misprint.\par
\noindent Imprimé : {\ttfamily\footnotesize 1\allowbreak{}}\par
\noindent Corrigé : {\ttfamily\footnotesize 2\allowbreak{}}\par
\noindent Source : 7II:126. \nolinkurl{fixes/MISPRINTS_7II.json#126}\par

\par\medskip\noindent\hypertarget{app-F7II-XII-016.1}{}\textbf{F7II-XII-016.1}\quad p. 76; printed misprint.\par
\noindent Imprimé : {\ttfamily\footnotesize 2\allowbreak{}}\par
\noindent Corrigé : {\ttfamily\footnotesize 3\allowbreak{}}\par
\noindent Source : 7II:127. \nolinkurl{fixes/MISPRINTS_7II.json#127}\par

\par\medskip\noindent\hypertarget{app-F7II-XII-017.1}{}\textbf{F7II-XII-017.1}\quad p. 76; printed misprint.\par
\noindent Imprimé : {\ttfamily\footnotesize m\allowbreak{}}\par
\noindent Corrigé : {\ttfamily\footnotesize 2m\allowbreak{}}\par
\noindent Source : 7II:128. \nolinkurl{fixes/MISPRINTS_7II.json#128}\par

\par\medskip\noindent\hypertarget{app-F7II-XII-018.1}{}\textbf{F7II-XII-018.1}\quad p. 77; misprint\_pass.\par
\noindent Imprimé : \emph{(vide)}\par
\noindent Corrigé : {\ttfamily\footnotesize )\allowbreak{}}\par
\noindent Source : 7II:129. \nolinkurl{fixes/FIXLOG_7II_CURRENT.jsonl#553}; \nolinkurl{fixes/MISPRINTS_7II.json#129}\par

\par\medskip\noindent\hypertarget{app-F7II-XII-018.2}{}\textbf{F7II-XII-018.2}\quad p. 77; misprint\_pass.\par
\noindent Imprimé : \emph{(vide)}\par
\noindent Corrigé : {\ttfamily\footnotesize )\allowbreak{}}\par
\noindent Source : 7II:130. \nolinkurl{fixes/FIXLOG_7II_CURRENT.jsonl#554}\par

\par\medskip\noindent\hypertarget{app-F7II-XII-019.1}{}\textbf{F7II-XII-019.1}\quad p. 77; printed misprint.\par
\noindent Imprimé : {\ttfamily\footnotesize 2\allowbreak{}}\par
\noindent Corrigé : {\ttfamily\footnotesize 2m\allowbreak{}}\par
\noindent Source : 7II:132. \nolinkurl{fixes/MISPRINTS_7II.json#132}\par

\par\medskip\noindent\hypertarget{app-F7II-XII-019.2}{}\textbf{F7II-XII-019.2}\quad p. 77; printed misprint.\par
\noindent Imprimé : \emph{(vide)}\par
\noindent Corrigé : {\ttfamily\footnotesize )\allowbreak{}}\par
\noindent Source : 7II:131. \nolinkurl{fixes/MISPRINTS_7II.json#131}\par

\par\medskip\noindent\hypertarget{app-F7II-XII-019.3}{}\textbf{F7II-XII-019.3}\quad p. 77; printed misprint.\par
\noindent Imprimé : {\ttfamily\footnotesize )\allowbreak{}}\par
\noindent Corrigé : \emph{(vide)}\par
\noindent Source : 7II:131. \nolinkurl{fixes/MISPRINTS_7II.json#131}\par

\par\medskip\noindent\hypertarget{app-F7II-XII-019.4}{}\textbf{F7II-XII-019.4}\quad p. 77; printed misprint.\par
\noindent Imprimé : {\ttfamily\footnotesize 2\allowbreak{}}\par
\noindent Corrigé : {\ttfamily\footnotesize 2m\allowbreak{}}\par
\noindent Source : 7II:133. \nolinkurl{fixes/MISPRINTS_7II.json#133}\par

\par\medskip\noindent\hypertarget{app-F7II-XII-020.1}{}\textbf{F7II-XII-020.1}\quad p. 79; misprint\_pass.\par
\noindent Imprimé : \emph{(vide)}\par
\noindent Corrigé : {\ttfamily\footnotesize ,2n\allowbreak{}}\par
\noindent Source : 7II:136. \nolinkurl{fixes/FIXLOG_7II_CURRENT.jsonl#555}\par

\par\medskip\noindent\hypertarget{app-F7II-XII-020.2}{}\textbf{F7II-XII-020.2}\quad p. 79; printed misprint.\par
\noindent Imprimé : {\ttfamily\footnotesize r\allowbreak{}}\par
\noindent Corrigé : {\ttfamily\footnotesize n\allowbreak{}}\par
\noindent Source : 7II:134. \nolinkurl{fixes/MISPRINTS_7II.json#134}\par

\par\medskip\noindent\hypertarget{app-F7II-XII-020.3}{}\textbf{F7II-XII-020.3}\quad p. 79; printed misprint.\par
\noindent Imprimé : {\ttfamily\footnotesize 5\allowbreak{}}\par
\noindent Corrigé : {\ttfamily\footnotesize 6\allowbreak{}}\par
\noindent Source : 7II:135. \nolinkurl{fixes/MISPRINTS_7II.json#135}\par

\par\medskip\noindent\hypertarget{app-F7II-XII-021.1}{}\textbf{F7II-XII-021.1}\quad p. 80; inherited corrected-layer difference.\par
\noindent Imprimé : \emph{(vide)}\par
\noindent Corrigé : {\ttfamily\footnotesize +1\allowbreak{}}\par
\noindent Source : ADOPT-F7II-XII-021.1. \nolinkurl{fr/sga7ii/diplomatic/expose_XII_body.tex}; \nolinkurl{fr/sga7ii/current/expose_XII_body.tex}\par

\par\medskip\noindent\hypertarget{app-F7II-XII-022.1}{}\textbf{F7II-XII-022.1}\quad p. 81; designated-witness attribution.\par
\noindent Imprimé : {\ttfamily\footnotesize envoyè\allowbreak{}}\par
\noindent Corrigé : {\ttfamily\footnotesize envoie\allowbreak{}}\par
\noindent Source : WITNESS-F7II-XII-022.1. \nolinkurl{fr/sga7ii/diplomatic/expose_XII_body.tex}; \nolinkurl{fr/sga7ii/current/expose_XII_body.tex}\par

\normalsize\fi
\clearpage\SGAtoc{Exposé XIII — Le formalisme des cycles évanescents}
\selectlanguage{french}
\SGApage{II}{82}


S G A  7


\underline{Exposé XIII}


\underline{\underline{Le formalisme des cycles évanescents}}


\underline{par  P. Deligne}

Introduction \hfill 2

\S{} 1.  Faisceaux d'ensembles \hfill 6

1.1  Faisceaux galoisiens \hfill 6

1.2  Le topos $Y \times_s S$ \hfill 8

1.3  Le foncteur $\Psi$ \hfill 13

1.4  La variation \hfill 16

\S{} 2.  Cycles évanescents \hfill 17

2.1  Le foncteur $R\Psi$ \hfill 17

2.2  Une compatibilité (morphismes trace) \hfill 26

2.3  Théorème de finitude (en égale caractéristique zéro) \hfill 29

2.4  Singularités isolées \hfill 30

\SGApage{II}{83}


\underline{Introduction}

Dans cet exposé, nous reprenons le formalisme introduit dans le n$^\circ$ 2 de l'exposé I de Grothendieck. Le seul point nouveau est l'introduction de l'analogue algébrique de la \underline{variation} de la théorie classique.

L'image topologique que je suis est la suivante. Soient $X$ un espace analytique, $D \subset \mathbb{C}$ un disque de centre $0$ , $D^* = D - \{0\}$ et $f : X \to D$ un morphisme propre (le cas non propre sera considéré ensuite). Remplaçons $D$ par un disque suffisamment petit centré en $0$ , $X$ par l'image réciproque de ce petit disque et notons encore $f : X \to D$ le résultat. D'après la variante analytique complexe du théorème d'isotopie de Thom, $f^{-1}(D^*)$ est un espace fibré topologique sur $D^*$ . J'ignore si on peut toujours construire un couple $(\Gamma, r)$ du type décrit ci-dessous, mais il semble plausible que oui.

\SGAvar{(a) $\Gamma$ est un système localement cohérent de trivialisations locales du fibré $f^{-1}(D^*)$ sur $D^*$ . En d'autre termes, $\Gamma$ définit une trivialisation $\pi_1(D^*)$-équivariante de l'image réciproque de $X$ sur un revêtement universel $\widetilde{D}^*$ de $D^*$ .

}{(a) $\Gamma$ est un système localement cohérent de trivialisations locales du fibré $f^{-1}(D^*)$ sur $D^*$ . En d'autres termes, $\Gamma$ définit une trivialisation $\pi_1(D^*)$-équivariante de l'image réciproque de $X$ sur un revêtement universel $\widetilde{D}^*$ de $D^*$ .

}{F7II-XIII-001}(b) $r$ est une rétraction de $X$ sur la fibre spéciale $X_o = f^{-1}(0)$ , compatible à $\Gamma$ : si, sur un ouvert connexe $U$ de $D^*$ , $\Gamma$ identifie $X_U = f^{-1}(U)$ à $X_1 \times U$ , $r = r(x_1, u)$ est indépendant de $u$ .

Soient donnés $(X, f, \Gamma , r)$ . Soient $t \in D$ et $X_t = f^{-1}(t)$ la ``fibre générale'' . Le fibré $\varphi^* X$ , image réciproque de $X/D$ par $\varphi : [0, 1] \to D^*$ : $x \longmapsto \exp(2 \pi i x).t$ , est trivialisé par $\Gamma$ . La \underline{monodromie} $T : X_t \longrightarrow X_t$ est l'homéomorphisme

\begin{equation*}
X_t = (\varphi^* X)_o \xrightarrow[(\Gamma)]{\sim} (\varphi^* X)_1 = X_t
\end{equation*}
Soit $r_t : X_t \longrightarrow X_o$ la restriction de $r$ à $X_t$ . On a $r_t T = r_t$ . On reconstruit comme suit $(X, f)$ à partir de $(X_t, X_o, T, r_t)$ :

($\alpha$) on recolle, dans $X_t \times [0, 1]$ , les extrémités $X_t \times \{0\}$ et $X_t \times \{1\}$ à l'aide de $T$ , pour obtenir $f' : X' \longrightarrow S^1$

($\beta$) $r_t$ définit $r' : X' \longrightarrow X_o$ ; $(X, f)$ s'identifie à
\SGApage{II}{84}
cône $(f')$ : (cône d'application de $r' : X' \longrightarrow X_o$) $\longrightarrow D =$ cône d'application de $(S^1 \to P^t)$

Soit $\Psi = (r, f) : X \longrightarrow X_o \times D$ . Les faisceaux images directes supérieures $R^i \Psi_* \mathbb{Z}$ de Leray vérifient la condition suivante.

(*)  $D^*$ est recouvert par des ouverts $U$ tels que la restriction du faisceau $F$ à $X_o \times U$ est image réciproque d'un faisceau sur $X_o$ .

\underline{Construction} 0.1, \underline{Les faisceaux} $F$ \underline{sur} $X_o \times D$ \underline{vérifiant} (*) \underline{s'identifient aux systèmes} $(F_o, F_t, T, \alpha)$ \underline{où}

$\alpha$)  $F_o$  \underline{est un faisceau sur} $X_o$  ;

$\beta$)  $F_t$  \underline{est un faisceau sur} $X_o$  , \underline{muni d'une action du générateur} $T$ \underline{de} $\pi_1(D^*,t)$ ;

$\gamma$)  $\alpha$  \underline{est un morphisme de} $F_o$ \underline{dans le faisceau des} $T$-\underline{invariants de} $F_t$ .

Si un faisceau $G$ sur $X_o \times D^*$ vérifie (*), et que $\widetilde{D}^*$ est le revêtement universel de $D^*$ (muni du point base $t$ ), l'image réciproque de $G$ sur $X_o \times \widetilde{D}^*$ est d'une et d'une seule façon image réciproque d'un faisceau $G_t$ sur $X_o$  ; $G_t$ est muni d'une action de $\pi_1(D^*, t)$ et $G$ se déduit de $\mathrm{pr}_1^*(G_t)$ sur $X_o \times \widetilde{D}^*$ par passage au quotient par $\pi_1(D^*, t)$ .

Un faisceau $F$ sur $X_o \times D$ est déterminé par

1) sa restriction à $X_o$ , soit $F_o$ ,

2) sa restriction à $X_o \times D^*$ , soit $G$ ,

3) un morphisme $\alpha : F_o \longrightarrow i^*j_*G$ (pour $i : X_o \hookrightarrow X_o \times D$ , $j : X_o \times D^* \hookrightarrow X_o \times D$) .

Si $F$ vérifie (*), $G$ est décrit par $(G_t, T)$ sur $X_o$ , et on trouve la description annoncée de $F$ .

\SGAvar{Pour les faisceau $R^i \Psi_* \mathbb{Z}$ , on a

}{Pour les faisceaux $R^i \Psi_* \mathbb{Z}$ , on a

}{F7II-XIII-002}\begin{equation*}
(R^i \Psi_* \mathbb{Z})_o \;=\; \begin{cases} \mathbb{Z} & \text{pour } i = 0 \\ 0 & \text{pour } i > 0 \ , \end{cases}
\end{equation*}
tandis que les $(R^i \Psi_* \mathbb{Z} )_t$ méritent (pour $i > 0$ ) le nom de \underline{faisceaux de cycles évanescents} (cf. la discussion au début du paragraphe 2 de I).


\SGApage{II}{85}

Cette discussion n'est utile que si on la raffine en passant à la catégorie dérivée: les $R^i \Psi_* \mathbb{Z}$ sont déduits d'un objet de la catégorie dérivée $D^+(X_o \times D)$ , et même d'un objet $R\Psi_* \mathbb{Z}$ de la catégorie dérivée de la catégorie des faisceaux abéliens sur $X_o \times D$ vérifiant (*). Cet objet détermine non seulement les faisceaux de cycles évanescents et leur monodromie, mais encore les morphismes ``variation'' classiques (cf. 1.4 et XIV 3.1).

C'est son analogue que nous définirons en géométrie algébrique. Si $S$ est un trait strictement hensélien et que $f : X \to S$ est un morphisme de schémas, nous définirons un objet $R\Psi(\mathbb{Z}/\ell^n)$ de la catégorie dérivée de la catégorie des faisceaux abéliens sur le topos produit $(X_o)_{\mathrm{et}} \times S_{\mathrm{et}}$ . Un faisceau sur ce topos s'interprète comme un système $(F_o, F_t, T,\alpha)$ où

$\alpha$)  $F_o$ est un faisceau sur $X_o$

\SGAvar{$\beta$)  $F_t$ est un faisceau sur $X_o$ , muni d'une action continue $T$ du groupe d'intertie $I$

}{$\beta$)  $F_t$ est un faisceau sur $X_o$ , muni d'une action continue $T$ du groupe d'inertie $I$

}{F7II-XIII-003}$\gamma$)  $\alpha$ est un morphisme de $F_o$ dans le faisceau des $I$-invariants de $F_t$ .

Nous aurons à travailler dans un cadre un peu plus général que celui indiqué ci-dessus:

a)  La construction de $R \Psi( \mathbb{Z} /\ell^n)$ ne suppose pas que $f$ soit propre. Pour l'interpréter topologiquement, il y aurait lieu de reprendre la construction qui précède, en remplaçant $X$ par un voisinage de la fibre spéciale $X_o$ .

b)  Pour construire $R \Psi(\mathbb{Z}/\ell^n)$ et démontrer ses propriétés, il faudra considérer non seulement le faisceau constant $\mathbb{Z}/\ell^n$ , mais encore des faisceaux généraux, voire des complexes de faisceaux.

c)  Lorsqu'on ne suppose pas que le corps résiduel $k$ du trait hensélien $S$ soit séparablement clos, il y a lieu de remplacer le topos produit $(X_o)_{\mathrm{et}} \times S_{\mathrm{et}}$ par le topos produit fibré $(X_o)_{\mathrm{et}} \times_{\mathrm{Spec}(k)_{\mathrm{et}}} S_{\mathrm{et}}$ . Le lecteur répugnant aux 2-produits fibrés de topos pourra prendre la description galoisienne 1.2.4 de ce topos comme en étant la définition.


\SGApage{II}{86}


0.2. \underline{Notations et terminologie} .

(0.2.1) \underline{trait} : spectre d'un anneau de valuation discrète .

(0.2.2) \underline{hensélien} : (pour un schéma) spectre d'un anneau local hensélien .

(0.2.3) \underline{strictement hensélien} : pour un anneau: hensélien à corps résiduel séparablement clos. Pour un schéma: spectre d'un tel anneau .

(0.2.4) \underline{point géométrique} de $S$ : morphisme $\bar{x} : \mathrm{Spec}(\bar{k}) \to S$ d'image $x \in S$ , $k(\bar{x}) =_{\mathrm{dfn}} \bar{k}$ étant une clôture séparable de $\kappa(x)$ . On dit que $\bar{x}$ est localisé en $x$ . Par abus de langage, nous appellerons parfois encore point géométrique un morphisme $\bar{x} : \mathrm{Spec}(\bar{k}) \to S$ tel que $\bar{k}$ soit une extension séparablement close de $\kappa(x)$ ; un tel morphisme se factorise de façon unique par un point géométrique (au sens propre) de $S$ .

(0.2.5) $s$, $\eta$, $\bar{s}$, $\bar{\eta}$, $\bar{S}$ : Si $S = \mathrm{Spec}(V)$ est un trait hensélien, nous noterons en principe $s$ le point fermé de $S$ , $\eta$ son point générique, $\bar{\eta}$ un point géométrique générique (= localisé en $\eta$ ) de $S$ et $\bar{s}$ le point géométrique localisé en $s$ correspondant (I 0.0.3). On note $\bar{S}$ le spectre du normalisé de $V$ dans $k(\bar{\eta})$ , de corps résiduel une extension inséparable de $k(\bar{s})$ . On note $\mathrm{Gal}(\bar{\eta}/\eta )$ le groupe de Galois de $k(\bar{\eta})$ sur $k(\eta)$ . De même pour $s$ . On a une suite exacte

\begin{equation*}
0 \longrightarrow I \longrightarrow \mathrm{Gal}(\bar{\eta}/\eta) \longrightarrow \mathrm{Gal}(\bar{s}/s) \longrightarrow 0
\end{equation*}
( $I$ groupe d'inertie).

(0.2.6) \underline{faisceau} (sur un schéma $X$ ) signifie toujours faisceau sur le site étale $X_{et}$ de $X$ .

\SGAvar{(0.2.7) \underline{premier à}: un faisceau de torsion $F$ sur un schéma $S$ (ou un groupe abélien de torsion $F$ ) est premier aux caractéristique résiduelles de $S$ si pour tout $s \in S$ , la multiplication par l'exposant caractéristique de $k(s)$ est un automorphisme de $F$ .


}{(0.2.7) \underline{premier à}: un faisceau de torsion $F$ sur un schéma $S$ (ou un groupe abélien de torsion $F$ ) est premier aux caractéristiques résiduelles de $S$ si pour tout $s \in S$ , la multiplication par l'exposant caractéristique de $k(s)$ est un automorphisme de $F$ .


}{F7II-XIII-004}\SGApage{II}{87}


\SGAsectiontoc{1. Faisceaux d'ensembles}{1}
1. \underline{Faisceaux d'ensembles}

Les faisceaux de cycles évanescents seront définis comme valeurs de foncteurs dérivés du foncteur exact à gauche $\Psi$ défini au n$^{\circ}$ 3 pour tout faisceau d'ensembles.


1.1. \underline{Faisceaux galoisiens.}

1.1.1. Soient $Y$ un schéma sur un corps $k$ de clôture séparable $\bar{k}$ et $\bar{Y} = Y \otimes_{k} \bar{k}$ . Soit $G$ un groupe profini et

\begin{equation*}
u : G \longrightarrow \mathrm{Gal}(\bar{k}/k)
\end{equation*}
un homomorphisme continu de $G$ dans le groupe de Galois de $\bar{k}$ sur $k$ . Le groupe de Galois $\mathrm{Gal}(\bar{k}/k)$ agit (à gauche) sur $\bar{Y}$ par transport de structure. Le groupe $G$ agit sur $\bar{Y}$ via $u$ .

\SGAvar{Soit $\mathfrak{F}$ un faisceau d'ensembles sur (le site étale de) $\bar{Y}$ . Une action de $G$ sur $\mathfrak{F}$ , \underline{compatible} à l'action de $G$ sur $\bar{Y}$ , est une action par automorphismesde $G$ sur $(\bar{Y}, \mathfrak{F} )$ , induisant l'action donnée sur $\bar{Y}$ . En d'autre termes, c'est un système d'isomorphismes

}{Soit $\mathfrak{F}$ un faisceau d'ensembles sur (le site étale de) $\bar{Y}$ . Une action de $G$ sur $\mathfrak{F}$ , \underline{compatible} à l'action de $G$ sur $\bar{Y}$ , est une action par automorphismes de $G$ sur $(\bar{Y}, \mathfrak{F} )$ , induisant l'action donnée sur $\bar{Y}$ . En d'autres termes, c'est un système d'isomorphismes

}{F7II-XIII-005}\begin{equation*}
\sigma(g) : u(g)_{*} \mathfrak{F} \longrightarrow \mathfrak{F} \;,
\end{equation*}
vérifiant $\sigma(gh) = \sigma(g)\, \sigma(h)$ . Si $G$ agit sur $\mathfrak{F}$ , alors, quel que soit $U$ étale sur $Y$ , définissant $\bar{U} = U \otimes_{k} \bar{k}$ étale sur $\bar{Y}$ , le groupe $G$ agit sur $\mathfrak{F}(\bar{U})$ .

\underline{Définition} 1.1.2. \underline{On dit que} $G$ \underline{agit continûment sur} $\mathfrak{F}$ \underline{si, pour tout} $U$ \underline{quasi-compact et étale sur} $Y$ , $G$ \underline{agit continûment sur l'ensemble discret} $\mathfrak{F}(\bar{U})$ .

Il reviendrait au même de ne considérer que les $U$ affines.

Si $\mathfrak{F}$ est un faisceau d'ensembles sur $Y$ , d'image réciproque $\bar{\mathfrak{F}}$ sur $\bar{Y}$ , alors le groupe $\mathrm{Gal}(\bar{k}/k)$ agit sur $\bar{\mathfrak{F}}$ , par transport de structure, de façon compatible à son action sur $\bar{Y}$ .

\underline{Rappel} 1.1.3. (i) \underline{L'action de} $\mathrm{Gal}(\bar{k}/k)$ \underline{sur} $\bar{\mathfrak{F}}$ \underline{est continue.}


\SGApage{II}{88}

(ii) \underline{Le foncteur}

$\mathfrak{F} \longmapsto \vec{\mathfrak{F}}$ \underline{muni de l'action de} $\mathrm{Gal}(\overline{k}/k)$

\underline{est une équivalence de la catégorie des faisceaux d'ensembles sur} $Y$ \underline{avec la catégorie des faisceaux d'ensembles sur} $\overline{Y}$ \underline{munis d'une action continue de} $\mathrm{Gal}(\overline{k}/k)$ \underline{compatible à l'action de} $\mathrm{Gal}(\overline{k}/k)$ \underline{sur} $\overline{Y}$ .

Exprimons $\overline{k}$ comme limite inductive d'extensions finies galoisiennes $k_i$ . Soient $\varphi_i : Y_i = Y \otimes_k k_i \longrightarrow Y$ et $\overline{\varphi} : \overline{Y} \longrightarrow Y$ les projections. Pour $U$ étale sur $Y$ , supposé affine, soit $U_i = U \otimes_k k_i$ . On a

\begin{equation*}
\vec{\mathfrak{F}}(\overline{U}) = \varinjlim \ \varphi_i^* \mathfrak{F}(U_i) \quad .
\end{equation*}
Sur $\varphi_i^* \mathfrak{F}(U_i)$ , $\mathrm{Gal}(\overline{k}/k)$ agit via $\mathrm{Gal}(k_i/k)$ , ce qui prouve (i).

Le foncteur (ii) a pour adjoint à droite le foncteur $\mathfrak{G} \longmapsto (\overline{\varphi}_* \mathfrak{G})^{\mathrm{Gal}(\overline{k}/k)}$ . Pour tout $\mathfrak{F}$ sur $Y$ , on a

\begin{equation*}
\overline{\varphi}_* \ \overline{\varphi}^* \ \mathfrak{F} = \varinjlim \ \varphi_{i*} \ \varphi_i^* \ \mathfrak{F}
\end{equation*}
Pour prouver que $\mathfrak{F} \overset{\sim}{\longrightarrow} (\overline{\varphi}_* \ \overline{\varphi}^* \ \mathfrak{F})^{\mathrm{Gal}(\overline{k}/k)}$ , il suffit donc de prouver que pour tout $i$ , $\mathfrak{F} \overset{\sim}{\longrightarrow} (\varphi_{i*} \ \varphi_i^* \mathfrak{F})^{\mathrm{Gal}(k_i/k)}$ . Il suffit même de prouver que cette flèche devient un isomorphisme après extension des scalaires de $k$ à $k_i$ ; $\mathrm{Spec}(k_i)$ devient alors somme de copies de $\mathrm{Spec}(k)$ , permutée de façon simplement transitive par $\mathrm{Gal}(k_i/k)$ , et l'assertion devient triviale.

Le foncteur exact (ii) est donc pleinement fidèle. Pour conclure, il reste à vérifier que la flèche

\begin{equation*}
\overline{\varphi}^*(\overline{\varphi}_* \ \mathfrak{G})^{\mathrm{Gal}(\overline{k}/k)} \longrightarrow \mathfrak{G}
\end{equation*}
est un épimorphisme. Soit $\overline{\varphi}_i : \overline{Y} \to Y_i$ . L'hypothèse donne que

\begin{equation*}
\varinjlim \ \overline{\varphi}_i^*(\overline{\varphi}_{i*} \ \mathfrak{G})^{\mathrm{Gal}(\overline{k}/k_i)} \longrightarrow \mathfrak{G} 
\end{equation*}

\SGApage{II}{89}

 est épimorphique. On a

\SGAvar{\begin{equation*}
\overline{\varphi}^*(\varphi_* \ \mathfrak{G})^{\mathrm{Gal}(\overline{k}/k)} \ = \ \overline{\varphi}_i^* \ (\varphi_i^* \ \varphi_{i*} \ ((\overline{\varphi}_{i*} \ \mathfrak{G})^{\mathrm{Gal}(\overline{k}/k_i)})^{\mathrm{Gal}(k_i/k)}) \quad ,
\end{equation*}
et on conclut en notant, par le même argument de descente galoisienne que plus haut, que pour $\mathfrak{G}'$ sur $Y_i$ muni d'une action de $\mathrm{Gal}(k_i/k)$ compatible à son action sur $Y_i$ , on a

}{\begin{equation*}
\overline{\varphi}^*(\overline{\varphi}_* \ \mathfrak{G})^{\mathrm{Gal}(\overline{k}/k)} \ = \ \overline{\varphi}_i^* \ (\varphi_i^* \ \varphi_{i*} \ ((\overline{\varphi}_{i*} \ \mathfrak{G})^{\mathrm{Gal}(\overline{k}/k_i)})^{\mathrm{Gal}(k_i/k)}) \quad ,
\end{equation*}
et on conclut en notant, par le même argument de descente galoisienne que plus haut, que pour $\mathfrak{G}'$ sur $Y_i$ muni d'une action de $\mathrm{Gal}(k_i/k)$ compatible à son action sur $Y_i$ , on a

}{F7II-XIII-006}\begin{equation*}
\varphi_i^*(\varphi_{i*} \ \mathfrak{G}')^{\mathrm{Gal}(k_i/k)} \overset{\sim}{\longrightarrow} \ \mathfrak{G}' \quad .
\end{equation*}

1.2. \underline{Le topos} $Y \times_s S$ .

\SGAvar{1.2.1. Soient $S$ un schéma local hensélien de point fermé $i : s \hookrightarrow S$ . Le foncteur $S' \longmapsto S' \times_S s = i^* S'$ est une équivalence de la catégorie des $S$-schémas finis étales $S'$ avec la catégorie des $s$-schémas finis étales $s'$ (spectres de $k(s)$-algèbres étales). Le foncteur inverse $\mathrm{sp}^*$ , du site étale de $s$ dans le site étale de $S$ , est un morphisme de sites

}{1.2.1. Soit $S$ un schéma local hensélien de point fermé $i : s \hookrightarrow S$ . Le foncteur $S' \longmapsto S' \times_S s = i^* S'$ est une équivalence de la catégorie des $S$-schémas finis étales $S'$ avec la catégorie des $s$-schémas finis étales $s'$ (spectres de $k(s)$-algèbres étales). Le foncteur inverse $\mathrm{sp}^*$ , du site étale de $s$ dans le site étale de $S$ , est un morphisme de sites

}{F7II-XIII-007}\begin{equation*}
\mathrm{sp} : S \longrightarrow s \quad ,
\end{equation*}
le \underline{morphisme de spécialisation} (cf. SGA4 VIII 7). Pour $s'$ étale sur $s$ , on dit que $\mathrm{sp}^*(s')$ se déduit de $S$ par \underline{extension du corps résiduel} . Pour $F$ un faisceau sur $S$ , on a un isomorphisme fonctoriel

\begin{equation}\tag{1.2.1.1}
\mathrm{sp}_* \ F = i^* \ F \quad .
\end{equation}
Pour tout ouvert de $S$ : $j : U \hookrightarrow S$ , ou pour tout point de $S$ : $j : \eta \hookrightarrow S$ , on notera encore $\mathrm{sp}$ le morphisme composé $\mathrm{sp} \circ j$ . Pour $F$ un faisceau sur $U$ (resp. sur $\eta$ ), on a

\begin{equation}\tag{1.2.1.2}
\mathrm{sp}_* \ F = i^* \ j_* \ F \quad .
\end{equation}
1.2.2. Soient $S$ un trait hensélien, et $s, \eta, \overline{s}, \overline{\eta}, I$ comme en 0.2.5.

\begin{center}
\begin{tikzcd}
s \arrow[r, hook, "i"] \& S \& \eta \arrow[l, hook, "j"']
\end{tikzcd}\quad .
\end{center}








\SGApage{II}{90}

Voici comment les faisceaux sur $s,\eta$, $S$ et le morphisme de spécialisation s'interprètent en termes galoisiens.

(a) Les faisceaux $F$ sur $s$ (resp. $\eta$) s'identifient, par $F \longrightarrow F_{\bar{s}}$ (resp. $F_{\bar{\eta}}$ ) aux ensembles munis d'une action continue de $\mathrm{Gal}(\bar{s}/s)$ (resp. $\mathrm{Gal}(\bar{\eta}/\eta)$) (SGA4 VIII 2.1, ou 1.1.3).

(b) Un faisceau $F$ sur $S$ définit des faisceaux $F_s = i^* F$ et $F_\eta = j^* F$ sur $s$ et $\eta$ , et une flèche d'adjonction $F \longrightarrow j_* j^* F$ induisant

\begin{equation*}
\varphi : F_s = i^* F \longrightarrow i^* j_* j^* F = \mathrm{sp}_* F_\eta \quad .
\end{equation*}
On sait (SGA4 IV 9.5.4) que le foncteur $F \longmapsto (F_s, F_\eta, \varphi)$ est une équivalence de la catégorie des faisceaux sur $S$ avec celle des triples $(F_s, F_\eta, \varphi)$

$(F_s$ faisceau sur $s$ , $F_\eta$ faisceau sur $\eta$, $\varphi : F_s \longrightarrow i^* j_* F_\eta )$ .

Via l'équivalence (a), le foncteur $i^* j_*$ s'identifie au foncteur ``invariants sous $I$'' . Les faisceaux sur $S$ s'identifient donc aux triples formés d'un ensemble $F_{\bar{s}}$ muni d'une action de $\mathrm{Gal}(\bar{s}/s)$ , d'un ensemble $F_{\bar{\eta}}$ muni d'une action de $\mathrm{Gal}(\bar{\eta}/\eta)$ et d'un morphisme équivariant $\varphi : F_{\bar{s}} \longrightarrow F_{\bar{\eta}}$ .

(c) Les morphismes $\mathrm{sp}$, $j$, $i$, $\mathrm{sp} = \mathrm{sp} \circ j$ s'expriment comme suit, via (a) (b)

\begin{equation*}
\begin{array}{llcl}
\mathrm{sp} : S \longrightarrow s \ \ : & \mathrm{sp}^*(F_{\bar{s}}) & = & (F_{\bar{s}}, F_{\bar{s}}, \mathrm{Id}) \\[2pt]
& \mathrm{sp}_*(F_{\bar{s}}, F_{\bar{\eta}}, \varphi) & = & F_{\bar{s}} \\[2pt]
j : \eta \longrightarrow S \ \ : & j^*(F_{\bar{s}}, F_{\bar{\eta}}, \varphi) & = & F_{\bar{\eta}} \\[2pt]
& j_*(F_{\bar{\eta}}) & = & (F^{I}_{\bar{\eta}}, F_{\bar{\eta}}, \mathrm{inclusion}) \\[2pt]
i : s \longrightarrow S \ \ : & i^*(F_{\bar{s}}, F_{\bar{\eta}}; \varphi) & = & F_{\bar{s}} \\[2pt]
& i_*(F_{\bar{s}}) & = & (F_{\bar{s}}, \{0\}, \mathrm{projection}) \\[2pt]
\mathrm{sp} = \mathrm{sp}\, j : \eta \longrightarrow s \ : & \mathrm{sp}^*(F_{\bar{s}}) & = & F_{\bar{s}} \\[2pt]
& \mathrm{sp}_*(F_{\bar{\eta}}) & = & F^{I}_{\bar{\eta}}
\end{array}
\end{equation*}
1.2.3. Soient $S$ un trait hensélien comme en 1.2.2 et $Y$ un schéma sur $s$ . Nous aurons à travailler avec le 2-produit fibré $Y \times_s S$ des topos étales de $Y$ et $S$ sur le topos étale de $s$ . D'après Giraud [1] , ce 2-produit fibré existe, et on voit facilement qu'il admet la description galoisienne donnée ci-dessous. Ceux qui
\SGApage{II}{91}
répugnent aux 2-produits fibrés pourront prendre cette description pour définition.

\underline{Construction} 1.2.4. \underline{Les faisceaux} $F$ \underline{sur} $Y \times_s S$ \underline{s'identifient aux triples} $(F_s, F_\eta, \varphi )$ \underline{où}

(a) $F_s$ \underline{est un faisceau sur} $Y$ , \underline{soit encore un faisceau} $F_{\bar{s}}$ \underline{sur} $\overline{Y} = Y \otimes_s \bar{s}$ , \underline{muni d'une action continue de} $\mathrm{Gal}(\bar{s}/s)$ \underline{compatible à l'action de} $\mathrm{Gal}(\bar{s}/s)$ \underline{sur} $\overline{Y}$ .

(b) $F_\eta$ \underline{est un faisceau} $F_{\bar{\eta}}$ \underline{sur} $\overline{Y}$ , \underline{muni d'une action continue de} $\mathrm{Gal}(\bar{\eta}/\eta)$ , \underline{compatible à l'action} (\underline{via} 1.1.1. ) \underline{de} $\mathrm{Gal}(\bar{\eta}/\eta)$ \underline{sur} $\overline{Y}$ .

(c) $\varphi$ \underline{est un morphisme équivariant de} $F_{\bar{s}}$ \underline{dans} $F_{\bar{\eta}}$ .

De même, un \underline{faisceau} $F_\eta$ \underline{sur} $Y \times_s \eta$ est un objet 1.2.4(b). On dispose d'un diagramme de topos

\begin{equation}\tag{1.2.4.1}
\begin{tikzcd}[column sep=large, row sep=small]
Y \arrow[dd] \& {} \& Y \arrow[r, hook, "i"] \arrow[dd] \& Y \times_s S \arrow[dd] \& Y \times_s \eta \arrow[l, hook', "j"'] \arrow[dd] \& {} \\
{} \& {} \arrow[l, "\mathrm{sp}"'] \& {} \& {} \& {} \& {} \\
s \& {} \& s \arrow[r, hook, "i"] \& S \& \eta \arrow[l, hook', "j"'] \& {.}
\end{tikzcd}
\end{equation}




























$Y \times_s S$ est réunion de l'ouvert $Y \times_s \eta$ et du fermé complémentaire $Y$ . Les formules 1.2.2.(c) restent valables pour $\mathrm{sp} : Y \times_s S \longrightarrow Y$ (encore noté $\mathrm{pr}_1$ ) $j : Y \times_s \eta \hookrightarrow Y \times_s S$, $i : Y \hookrightarrow Y \times_s S$ et $\mathrm{sp} : Y \times_s \eta \longrightarrow Y$ (encore noté $\mathrm{pr}_1$ ).

1.2.4.2. Si on prend 1.2.4 pour définition de $Y \times_s S$ , il y a lieu de montrer, par une construction directe, que $Y \times_s S$ ne dépend pas à isomorphisme unique près, du choix de la clôture séparable $k(\bar{\eta})$ de $k(\eta)$ . On le fait en interprétant $F_\eta$ de 1.2.4(b) comme un foncteur qui à chaque clôture séparable $k(\bar{\eta})$ de $k(\eta)$ , définissant une clôture séparable $k(\bar{s})$ de $k(s)$ , associe un faisceau sur $Y \otimes_{k(s)} k(\bar{s})$ : le $F_\eta$ de 1.2.4(b) est la restriction de ce foncteur à $\{\bar{\eta}\}$ .


\SGApage{II}{92}

1.2.5. Chaque point géométrique $\bar{x}$ de $Y_{\bar{s}}$ définit un point géométrique $(\bar{x},\bar{\eta})$ de $Y \times_s \eta$ , de foncteur fibre correspondant

\begin{equation*}
F_\eta \longmapsto (F_{\bar{\eta}})_{\bar{x}} \quad .
\end{equation*}
Il définit des points géométriques $(\bar{x},\bar{\eta})$ et $(\bar{x}, \bar{s})$ de $Y \times_s S$ par

\begin{equation*}
F_{(\bar{x}\,\bar{\eta})} = (F_{\bar{\eta}})_{\bar{x}} \quad \mbox{et} \quad F_{(\bar{x}\,\bar{s})} = (F_{\bar{s}})_{\bar{x}} \quad .
\end{equation*}
On vérifie qu'on obtient ainsi tous les points des topos $Y \times_s \eta$ et $Y \times_s S$ . Nous ne l'utiliserons pas. On vérifie aussitôt que ces ensembles de points sont conservatifs.

$^{*}$\underline{Remarque} 1.2.6. La construction 1.2.5 définit une équivalence de la catégorie $\mathrm{Point}(Y \times_s \eta)$ (resp. $\mathrm{Point}(Y \times_s S)$ ) des points de $Y \times_s \eta$ (resp.....) avec le 2-produit fibré $\mathrm{Point}(Y) \times_{\mathrm{Point}(s)} \mathrm{Point}(\eta)$ .

Cette équivalence résulte aussi de ce que $Y \times_s \eta$ (resp. $Y \times_s S$) est un 2-produit fibré de topos.$^{*}$

1.2.7. Par la fonctorialité des 2-produits fibrés, le diagramme (1.2.4.1) est fonctoriel en $Y$ et $S$ . En particulier

a) Tout morphisme $f : Y \longrightarrow Y'$ de schémas sur $s$ induit un morphisme de topos de $Y \times_s S$ dans $Y' \times_s S$ . On a

\begin{equation}\tag{1.2.7.1}
f^{*}(F_{\bar{s}}, F_{\bar{\eta}}, \varphi ) = (f^{*} F_{\bar{s}}, f^{*} F_{\bar{\eta}} , f^{*}(\varphi))
\end{equation}
et pour $f$ quasi-compact

\begin{equation}\tag{1.2.7.2}
f_{*}(F_{\bar{s}}, F_{\bar{\eta}} , \varphi) = (f_{*} F_{\bar{s}}, f_{*} F_{\bar{\eta}} , f_{*}(\varphi)) \quad .
\end{equation}
Des formules analogues valent pour $f : Y \times_s \eta \longrightarrow Y' \times_s \eta$ .

b) Tout morphisme surjectif de traits henséliens $f : S' \longrightarrow S$ induit un morphisme de topos de $Y_{s'} \times_{s'} S'$ dans $Y \times_s S$ . Notant encore $f$ le morphisme naturel de
\SGApage{II}{93}
$Y_{s'} = Y \otimes_{k(s)} k(s')$ dans $Y$ , on a

\begin{equation}\tag{1.2.7.3}
f^{*}(F_{\bar{s}}, F_{\bar{\eta}}, \varphi) = (f^{*}F_{\bar{s}}, f^{*} F_{\bar{\eta}} , f^{*}(\varphi)) \quad .
\end{equation}
Lorsque $f$ est un morphisme fini, $f_{*}$ s'exprime comme représentation induite.

1.2.8. Soit $f : Y \longrightarrow Y'$ un morphisme localement de type fini et séparé de $s$-schémas. On définit un foncteur ``image directe à support propre''

\begin{equation*}
\begin{array}{l} f_{!} : (\mbox{faisceaux d'ensembles pointés sur } Y \times_s S) \longrightarrow (\mbox{faisceaux d'ensembles} \\ \qquad\qquad\qquad\qquad\qquad\qquad\qquad \mbox{pointés sur } Y' \times_s S) \end{array}
\end{equation*}
par la formule

\begin{equation}\tag{1.2.8.1}
f_{!} (F_{\bar{s}}, F_{\bar{\eta}} , \varphi) = (f_{!} F_{\bar{s}} , f_{!} F_{\bar{\eta}} , f_{!}(\varphi)) \quad .
\end{equation}
Pour $f$ un plongement localement fermé, $f_{!}$ s'appelle le foncteur de prolongement par zéro.

Un foncteur analogue est défini pour $Y \times_s \eta$ .

1.2.9. Soit $f : Y \hookrightarrow Y'$ un plongement localement fermé de $s$-schémas. Le foncteur $f_{!}$ a un adjoint à droite $f^{!}$ , ``sections à support dans $Y$ '', vérifiant

\begin{equation}\tag{1.2.9.1}
f^{!} (F_{\bar{s}}, F_{\bar{\eta}} , \varphi) = (f^{!} F_{\bar{s}}, f^{!} F_{\bar{\eta}} , f^{!} (\varphi)) \quad .
\end{equation}
Supposons seulement $f$ quasi-fini . Le foncteur $f_{!}$ des faisceaux abéliens sur $Y \times_s S$ dans les faisceaux abéliens sur $Y' \times_s S$ admet encore un adjoint à droite $f^{!}$ , donné par 1.2.9.1 (cf. SGA4 XVIII 3.1.8). Pour $f$ étale, $f^{!} = f^{*}$.

Des foncteurs analogues sont définis pour $Y \times_s \eta$ .


\SGApage{II}{94}


1.3. \underline{Le foncteur} $\Psi$ .

1.3.1. Soient $S$ un trait hensélien comme en 0.2.5 et $p : X \longrightarrow S$ . On se propose de définir un foncteur exact à gauche $\Psi$ des faisceaux sur $X$ dans les faisceaux sur le topos $X_s \times_s S$ construit en 1.2. Ce foncteur n'est pas en général le foncteur image directe d'un morphisme de topos.

On définira aussi $\Psi_\eta$ : (faisceaux sur $X_\eta$ ) $\longrightarrow$ (faisceaux sur $X_s \times_s \eta$).

1.3.2. Soit $\bar X = X \times_s \bar S$ . On dispose d'un diagramme

\begin{equation}\tag{1.3.2.1}
\begin{tikzcd}
X_{\bar s} \arrow[r, hook, "\bar i"] \arrow[d] \& \bar X \arrow[d] \& X_{\bar \eta} \arrow[l, hook', "\bar j"'] \arrow[d] \\
X_s \arrow[r, hook] \& X \& X_\eta \arrow[l, hook']
\end{tikzcd}
\end{equation}











cartésien au dessus du diagramme

\begin{center}
\begin{tikzcd}
\bar s \arrow[r, hook] \arrow[d] \& \bar S \arrow[d] \& \bar \eta \arrow[l, hook'] \arrow[d] \\
s \arrow[r, hook] \& S \& \eta \arrow[l, hook']
\end{tikzcd}\quad .
\end{center}









Soit $F$ un faisceau sur $X_\eta$ , d'image réciproque $F_{\bar\eta}$ sur $X_{\bar\eta}$ . On pose

\begin{equation}\tag{1.3.2.2}
\Psi_\eta (F) \; = \; \bar i^* \, \bar j_* \, F_{\bar\eta} \qquad .
\end{equation}
Par transport de structure, ce faisceau est muni d'une action continue de $\mathrm{Gal}(\bar\eta / \eta)$ , compatible à l'action de $\mathrm{Gal}(\bar\eta/\eta)$ sur $X_{\bar s}$ . Le foncteur

\begin{equation}\tag{1.3.2.3}
\Psi_\eta : (X_{\eta \ \mathrm{et}})^{\sim} \longrightarrow (X_s \times_s \eta)^{\sim}
\end{equation}
est exact à gauche.

1.3.3. Soient $F$ un faisceau sur $X$ , $F_\eta$ sa restriction à $X_\eta$ et $F_s$ sa restriction à $X_s$ . On pose 


\SGApage{II}{95}

\begin{equation}\tag{1.3.3.1}
 (\Psi (F))_\eta \; = \; \Psi_\eta(F_\eta)
\end{equation}
\begin{equation}\tag{1.3.3.2}
(\Psi (F))_s \; = \; F_s \quad .
\end{equation}
Soient $\bar F$ l'image réciproque de $F$ sur $\bar X$ et $\varphi'$ le morphisme d'adjonction $\varphi' : \bar F \longrightarrow \bar j_* \, \bar j^* \, \bar F$ . Ce morphisme induit

\begin{equation*}
\varphi : F_{\bar s} = \bar i^* \, \bar F \longrightarrow \bar i^* \, \bar j_* \, \bar j^* \, \bar F = (\Psi_\eta(F_\eta))_{\bar\eta} \qquad .
\end{equation*}
Le triple $\Psi(F) = (\Psi(F)_s \, , \Psi(F)_\eta, \varphi)$ est un faisceau sur $X_s \times_s S$ . Le foncteur

\begin{equation}\tag{1.3.3.3}
\Psi : (X_{\mathrm{et}})^{\sim} \longrightarrow (X_s \times_s S)^{\sim}
\end{equation}
est exact à gauche.

1.3.4. Lorsqu'il n'y aura pas de risque de confusion, on posera

\begin{equation*}
\Psi_\eta(F) = (\, \Psi(F))_\eta \quad , \quad \Psi_{\bar\eta}(F) = (\, \Psi(F))_{\bar\eta} \quad \mathrm{et} \quad \Psi_{\bar\eta}(F) = (\Psi_\eta \, (F))_{\bar\eta} \qquad .
\end{equation*}
1.3.5. Dans la fin de ce numéro, nous dressons une liste de propriétés de fonctorialité de $\Psi$ (cf. 1.2.7 à 1.2.9). Elles seront utiles surtout sous leur forme dérivée (2.1).

Soit $f : X \longrightarrow X'$ un morphisme de $S$-schémas, d'où un diagramme

\SGAvar{\begin{equation}\tag{1.3.5.1}
\begin{tikzcd}
X_{\bar s} \arrow[r, hook, "\bar i"] \arrow[d, "f"] \& X \arrow[d, "f"] \& X_{\bar \eta} \arrow[l, hook', "\bar j"'] \arrow[d, "f"] \\
X'_{\bar s} \arrow[r, hook, "\bar i"] \& \bar X' \& X'_{\bar \eta} \arrow[l, hook', "\bar j"']
\end{tikzcd}\quad .
\end{equation}














}{\begin{equation}\tag{1.3.5.1}
\begin{tikzcd}
X_{\bar s} \arrow[r, hook, "\bar i"] \arrow[d, "f"] \& \bar X \arrow[d, "f"] \& X_{\bar \eta} \arrow[l, hook', "\bar j"'] \arrow[d, "f"] \\
X'_{\bar s} \arrow[r, hook, "\bar i"] \& \bar X' \& X'_{\bar \eta} \arrow[l, hook', "\bar j"']
\end{tikzcd}\quad .
\end{equation}














}{F7II-XIII-008}L'ambiguité dans la notation des flèches simplifiera les formules, et ne prête pas à confusion: une seule interprétation donne un sens aux flèches composées ci-dessous. De même, nous noterons $\Psi$ le foncteur (1.3.3.3) pour $X$ tant que pour $X'$ .

Les morphismes de foncteurs décrits ci-dessous pour $\Psi$ ont une variante pour $\Psi_\eta$ , que nous omettons en général d'expliciter.


\SGApage{II}{96}

1.3.6. Le morphisme de changement de base $\bar{i}^* f_* \longrightarrow f_* \bar{i}^*$ induit un morphisme de foncteurs

\begin{equation}\tag{1.3.6.1}
\Psi\, f_* \longrightarrow f_*\, \Psi
\end{equation}
qui est un isomorphisme pour $f$ propre (SGA4 XII 5.1(i)).

Le cas essentiel est celui où $X' = S$ . On trouve un morphisme

\begin{equation}\tag{1.3.6.2}
f_* \longrightarrow f_*\, \Psi
\end{equation}
de composante essentielle

\begin{equation}\tag{1.3.6.3}
\Gamma(X_{\bar{\eta}},\, F) \longrightarrow \Gamma(X_{\bar{s}},\, \Psi_{\bar{\eta}}(F)) \quad ;
\end{equation}
ce morphisme est un isomorphisme pour $f$ propre.

1.3.7. Le morphisme de changement de base $f^*\,\bar{j}_* \longrightarrow \bar{j}_*\, f^*$ induit un morphisme de foncteurs

\begin{equation}\tag{1.3.7.1}
f^*\Psi \longrightarrow \Psi\, f^* \quad .
\end{equation}
1.3.8. Supposons $f$ quasi-fini. Les foncteurs $f_!$ commutent aux changements de base; le morphisme $f_!\,\bar{j}_* \longrightarrow \bar{j}_*\, f_!$ induit un morphisme de foncteurs

\begin{equation}\tag{1.3.8.1}
f_!\,\Psi \longrightarrow \Psi f_! \quad .
\end{equation}
Lorsque $f$ est fini, c'est l'inverse de (1.3.6.1).

1.3.9. La situation est duale pour $f^!$ (on suppose que $f$ est quasi-fini pour travailler avec les faisceaux abéliens, ou est un plongement pour les faisceaux d'ensembles pointés). Ce foncteur commute aux images directes, et le morphisme $\bar{i}^*\, f^! \longrightarrow f^!\,\bar{i}^*$ induit un morphisme de foncteurs

\begin{equation}\tag{1.3.9.1}
\Psi\, f^! \longrightarrow f^!\Psi \quad .
\end{equation}
Pour $f$ étale, c'est l'inverse de (1.3.7.1).

1.3.10. Enfin, pour un changement de base $S' \stackrel{f}{\longrightarrow} S$ , on trouve un morphisme de
\SGApage{II}{97}
changement de base

\begin{equation}\tag{1.3.10.1}
f^*\Psi \longrightarrow \Psi\, f^* \quad .
\end{equation}

1.4. \underline{La variation.}

1.4.1. Soient $S$ un trait hensélien comme en 0.2.5 et $Y$ un schéma sur $s$ . Soit $\Lambda$ un anneau. On pourrait plus généralement prendre pour $\Lambda$ un faisceau d'anneaux sur $Y$ . Pour tout objet $K$ de la catégorie dérivée $D(Y \times_s S, \Lambda)$ de la catégorie des faisceaux de $\Lambda$-modules sur $Y \times_s S$ , nous nous proposons de définir un morphisme \underline{variation} .

\SGAvar{1.4.2. Un complexe de faisceau de $\Lambda$-modules $K$ sur $X \times_s S$ s'identifie à un triple $(K_s, K_\eta , \varphi)$ où

a) $K_s$ (resp. $K_\eta$ ) est un complexe de faisceaux $K_{\bar{s}}$ (resp. $K_{\bar{\eta}}$ ) sur $\bar{Y}$ , muni d'une action continue de $\mathrm{Gal}(\bar{s}/\bar{s})$ (resp. $\mathrm{Gal}(\bar{\eta}/\eta)$);

}{1.4.2. Un complexe de faisceaux de $\Lambda$-modules $K$ sur $Y \times_s S$ s'identifie à un triple $(K_s, K_\eta , \varphi)$ où

a) $K_s$ (resp. $K_\eta$ ) est un complexe de faisceaux $K_{\bar{s}}$ (resp. $K_{\bar{\eta}}$ ) sur $\bar{Y}$ , muni d'une action continue de $\mathrm{Gal}(\bar{s}/s)$ (resp. $\mathrm{Gal}(\bar{\eta}/\eta)$);

}{F7II-XIII-009}b) $\varphi$ est un morphisme équivariant $\varphi\colon K_{\bar{s}} \to K_{\bar{\eta}}$ .

Un tel complexe $K$ est toujours homotope à un complexe $K' = (K'_s, K'_\eta , \varphi')$ tel que $\varphi'$ soit injectif et que la suite exacte

\begin{equation}\tag{1.4.2.1}
0 \longrightarrow K'_{\bar{s}} \longrightarrow K'_{\bar{\eta}} \longrightarrow \mathrm{coker}(\varphi'\,) \longrightarrow 0
\end{equation}
soit scindée degré par degré. Il suffit de prendre $K_s = K'_s$ , et de prendre pour $K'_\eta$ la somme de $K_\eta$ et du cône de $\mathrm{sp}^* K_s$ . Posons $\Phi(K) = \mathrm{coker}(\varphi')$ . La suite exacte (1.4.2.1) définit un triangle distingué

\begin{equation}\tag{1.4.2.2}
\begin{tikzcd}[column sep=large, row sep=large]
{} \& \Phi(K) \arrow[dl] \& {} \\
\mathrm{sp}^*\, K_s \arrow[rr] \& {} \& K_\eta \arrow[ul]
\end{tikzcd}
\end{equation}







dans la catégorie $K(Y \times_s \eta , \Lambda)$ des complexes de faisceaux de $\Lambda$-modules à homotopie près sur $Y \times_s \eta$ . Ce triangle ne dépend que de $K \in \mathrm{Ob}\, K(Y \times_s S, \Lambda)$ .


\SGApage{II}{98}

\SGAvar{Cette construction passe à la catégorie dérivée, et associe à tout objet $K$ de $D(X \times_s S, \Lambda)$ un triangle distingué de $D(X \times_s \eta, \Lambda)$ .

}{Cette construction passe à la catégorie dérivée, et associe à tout objet $K$ de $D(Y \times_s S, \Lambda)$ un triangle distingué de $D(Y \times_s \eta, \Lambda)$ .

}{F7II-XIII-010}Le triangle (1.4.2.2) est utile surtout via la suite exacte longue d'hypercohomologie qui s'en déduit

\begin{equation*}
\dots \longrightarrow H^i(\overline{Y}, K_{\bar{s}}) \longrightarrow H^i(\overline{Y}, K_{\bar{\eta}}) \longrightarrow H^i(\overline{Y}, \Phi(K)_{\bar{\eta}}) \longrightarrow \dots \qquad .
\end{equation*}
1.4.3. Le groupe d'inertie $I$ agit trivialement sur $K'_{\bar{s}}$ . Pour tout $\sigma \in I$ , l'endomorphisme $\sigma - 1$ de $K'_{\bar{\eta}}$ (notation (1.4.2.1)) se factorise donc par un morphisme

\begin{equation*}
\mathrm{Var}(\sigma) : \mathrm{coker}(\varphi') \longrightarrow K'_{\bar{\eta}} \qquad .
\end{equation*}
Cette construction passe à la catégorie dérivée: pour $K \in \mathrm{Ob}\, D(Y \times_s S, \Lambda)$ et $\sigma \in I$ , elle définit la \underline{variation}

\begin{equation}\tag{1.4.3.1}
\mathrm{Var}(\sigma) : \Phi(K)_{\bar{\eta}} \longrightarrow K_{\bar{\eta}} \qquad .
\end{equation}
Notant $q$ la flèche de $K_{\eta}$ dans $\Phi(K)$ , on a

\begin{equation}\tag{1.4.3.2}
\sigma \;=\; 1 + \mathrm{Var}(\sigma)q \qquad\qquad\qquad (\text{sur } K_{\bar{\eta}} )
\end{equation}
\begin{equation}\tag{1.4.3.3}
\sigma \;=\; 1 + q\, \mathrm{Var}(\sigma) \qquad\qquad\qquad (\text{sur } \Phi(K)_{\bar{\eta}} )
\end{equation}
\begin{equation}\tag{1.4.3.4}
\mathrm{Var}(\sigma\, \tau) = \mathrm{Var}(\sigma)\, q\, \mathrm{Var}(\tau) + \mathrm{Var}(\sigma) + \mathrm{Var}(\tau) \qquad .
\end{equation}

\SGAsectiontoc{2. Cycles évanescents}{2}
2. \underline{Cycles évanescents}


2.1. \underline{Le foncteur} $R\,\Psi$ .

2.1.1. Soient $S$ un trait hensélien comme en (0.2.5) et $\Lambda$ un anneau de torsion premier à la caractéristique résiduelle $p$ de $S$ (0.2.7). Plus généralement, on pourrait prendre pour $\Lambda$ un faisceau de torsion premier aux caractéristiques résiduelles; le cas le plus important est celui où $\Lambda = \mathbb{Z}/\ell^{n}$ ( $\ell$ premier inversible dans $\mathcal{O}_{S}$ ).


\SGApage{II}{99}

Soit $X$ un schéma sur $S$ . Le foncteur $\Psi$ (resp. $\Psi_{\eta}$ ) transforme faisceaux de $\Lambda$-modules sur $X$ (resp. $X_{\eta}$ ) en faisceaux de $\Lambda$-modules sur $X_s \times_s S$ (resp. $X_s \times_s \eta$ ) . Soient $R\Psi$ et $R\Psi_{\eta}$ les foncteurs dérivés

\begin{equation*}
R\Psi \;:\; D^{+}(X, \Lambda) \longrightarrow D^{+}(X_s \times_s S, \Lambda) \qquad ,
\end{equation*}
\begin{equation*}
R\Psi_{\eta}\colon D^{+}(X_{\eta}, \Lambda) \longrightarrow D^{+} (X_s \times_s \eta \,, \Lambda ) \qquad ,
\end{equation*}
et $R\Phi$ le foncteur composé $\Phi \circ R\Psi$

\begin{equation*}
R\Phi \;:\; D^{+}(X, \Lambda) \longrightarrow D^{+}(X_s \times_s \eta \,, \Lambda) \qquad .
\end{equation*}
2.1.2. Soit le diagramme commutatif (1.3.2.1)

\begin{equation}\notag
\begin{tikzcd}
X_{\bar{s}} \arrow[r, hook, "\bar{i}"] \arrow[d] \& \bar{X} \arrow[d] \& X_{\bar{\eta}} \arrow[l, hook', "\bar{j}"'] \arrow[d] \\
X_{s} \arrow[r, hook, "i"] \& X \& X_{\eta} \arrow[l, hook', "j"']
\end{tikzcd}\quad .
\end{equation}

















\SGAvar{Si $F$ est un faisceau de modules injectif sur $X$ , sa restriction à l'ouvert $X_{\eta}$ et encore injective. Pour $F_{\eta}$ injectif sur $X_{\eta}$ , l'image réciproque $F_{\bar{\eta}}$ est acyclique. Les formules (1.3.3.1) et (1.3.3.2) se dérivent donc en

}{Si $F$ est un faisceau de modules injectif sur $X$ , sa restriction à l'ouvert $X_{\eta}$ est encore injective. Pour $F_{\eta}$ injectif sur $X_{\eta}$ , l'image réciproque $F_{\bar{\eta}}$ est acyclique. Les formules (1.3.3.1) et (1.3.3.2) se dérivent donc en

}{F7II-XIII-011}\begin{equation}\tag{2.1.2.1}
R\Psi(K)_{s} \;=\; i^{*} K \qquad\qquad\qquad (K \in \mathrm{Ob}\, D^{+}(X, \Lambda))
\end{equation}
\begin{equation}\tag{2.1.2.2}
R\Psi(K)_{\eta} \;=\; R\Psi_{\eta}(K_{\eta}) \qquad\qquad\qquad (K \in \mathrm{Ob}\, D^{+}(X, \Lambda))
\end{equation}
\begin{equation}\tag{2.1.2.3}
R\Psi_{\eta}(K)_{\bar{\eta}} \;=\; \bar{i}^{*}\, R\bar{j}_{*}\, K_{\bar{\eta}} \qquad\qquad\qquad (K \in \mathrm{Ob}\, D^{+}(X_{\eta}, \Lambda))
\end{equation}
Ceci permet de poser sans conflit

\begin{equation*}
R\Psi_{\eta}(K) = (R\Psi(K))_{\eta} \;,\; R\Psi_{\bar{\eta}}(K) = (R\Psi(K))_{\bar{\eta}} \quad \text{et} \quad R\Psi_{\bar{\eta}}(K) = (R\Psi_{\eta}(K))_{\bar{\eta}} \qquad .
\end{equation*}
Avec ces notations, le triangle distingué (1.4.2.2) s'écrit 


\SGApage{II}{100}

\SGAvar{\begin{equation}\tag{2.1.2.4}
\begin{tikzcd}[row sep=large, column sep=huge]
{} \& R\Phi(K) \arrow[dl] \& {} \\
\mathrm{sp}^*\, i^*\, K \arrow[rr] \& {} \& R\Psi_\eta(K_s) \arrow[ul] \quad .
\end{tikzcd}
\end{equation}







}{\begin{equation}\tag{2.1.2.4}
\begin{tikzcd}[row sep=large, column sep=huge]
{} \& R\Phi(K) \arrow[dl] \& {} \\
\mathrm{sp}^*\, i^*\, K \arrow[rr] \& {} \& R\Psi_\eta(K_\eta) \arrow[ul] \quad .
\end{tikzcd}
\end{equation}







}{F7II-XIII-012}Pour $\sigma \in I$ , on dispose d'un morphisme \underline{variation}

\begin{equation}\tag{2.1.2.5}
\mathrm{Var}(\sigma) : R\Phi(K)_{\bar\eta} \longrightarrow R\Psi_{\bar\eta}(K) \quad .
\end{equation}
Les faisceaux $R^i\Phi(\Lambda)$ (ou plus généralement $R^i\Phi(K)$ ) sont les \underline{faisceaux de cycles évanescents} .

\SGAvar{2.1.3. Soit $(\bar x, \bar\eta)$ le point géométrique de $X_s \times_s \eta$ défini par un point géométrique $\bar x$ de $X_{\bar s}$ (1.2.5). Calculons la fibre de $R\Psi_\eta(K_\eta)$ en $(\bar x, \bar\eta)$ à l'aide de (2.1.2.3) (cf. I 2.3). L'hensélisé strict $X_{(\bar x)}$ de $X$ en $\bar x$ est un schéma sur l'hensélisé strict $S_{nr}$ de $S$ , de point géométrique $\eta_{nr}$ .

\underline{Proposition} 2.1.4.. \underline{On a} $(R\Psi(K_\eta))_{(\bar x,\bar\eta)} = R\Gamma(X_{(\bar x)} \times_{\eta_{xr}} \bar\eta , K)$ ; \underline{en particulier,}

}{2.1.3. Soit $(\bar x, \bar\eta)$ le point géométrique de $X_s \times_s \eta$ défini par un point géométrique $\bar x$ de $X_{\bar s}$ (1.2.5). Calculons la fibre de $R\Psi_\eta(K_\eta)$ en $(\bar x, \bar\eta)$ à l'aide de (2.1.2.3) (cf. I 2.3). L'hensélisé strict $X_{(\bar x)}$ de $X$ en $\bar x$ est un schéma sur l'hensélisé strict $S_{nr}$ de $S$ , de point générique $\eta_{nr}$ .

\underline{Proposition} 2.1.4.. \underline{On a} $(R\Psi(K_\eta))_{(\bar x,\bar\eta)} = R\Gamma(X_{(\bar x)} \times_{\eta_{nr}} \bar\eta , K)$ ; \underline{en particulier,}

}{F7II-XIII-013}\begin{equation*}
R^i\Psi(K_\eta)_{(\bar x,\bar\eta)} = \mathbb{H}^i(X_{(\bar x)} \times_{\eta_{nr}} \bar\eta, K) \quad .
\end{equation*}
Appliquons SGA4 XV 2.1 (théorème d'acyclicité locale des morphismes lisses). On trouve la

\underline{Reformulation} 2.1.5. \underline{Si} $f$ \underline{est lisse et} $F$ \underline{un faisceau localement constant, on a} $R\Phi(F) = 0$ .

Plus généralement, si $K \in \mathrm{Ob}\ D^+(X, \Lambda)$ , alors $R\Phi(K) = 0$ là où $f$ est lisse et que les $\underline{H}^i(K)$ sont localement constants. Voir aussi 2.1.11.

2.1.6. Soit $f : Y \longrightarrow Y'$ un morphisme de schémas sur $s$ . Les foncteurs $f^*\ f_*\ f^!\ f_!$ (1.2.7 à 1.2.9) ont un analogue naturel dans la catégorie dérivée.

a) Le foncteur $f^*$ est exact, donc se dérive trivialement.


\SGApage{II}{101}

b) Le foncteur $f_*$ a un dérivé droit $Rf_*$ ; pour $f$ quasi-compact, on a

\begin{equation}\tag{2.1.6.1}
(Rf_*\, K)_{\bar s} = Rf_{\bar s\,*}(K_{\bar s}) \quad ;
\end{equation}
\begin{equation}\tag{2.1.6.2}
(Rf_*\, K)_{\bar\eta} = Rf_{\bar s\,*}(K_{\bar\eta}) \quad .
\end{equation}
c) Pour $f$ quasi-fini, le foncteur $f_!$ est exact, donc se dérive trivialement. Dans le cas général, le foncteur dérivé du foncteur $f_!$ est pathologique, et on procède comme dans SGA4 XVII. On suppose $f$ séparé de type fini, et $Y'$ noethérien. Il existe alors [2] une décomposition $f = \widetilde{f}\, j$ , où $j$ est une immersion ouverte et où $\widetilde{f}$ est propre, et on définit

\begin{equation*}
Rf_! = R\widetilde{f}_* \circ j_! \quad .
\end{equation*}
Les arguments de SGA4 XVII montrent que le foncteur $Rf_!$ est indépendant du choix de la factorisation $\widetilde{f}\, j$ . L'analogue de (2.1.6.1) (2.1.6.2) est valable.

d) Nous n'utiliserons $Rf^!$ que pour un morphisme quasi-fini, auquel cas c'est le foncteur dérivé de $f^!$ . On a encore l'analogue de (2.1.6.1) (2.1.6.2).

2.1.7. Passons en revue les propriétés de fonctorialité de $R\Psi$ (i.e. dérivons 1.3.5 à 1.3.10). Soit $f : X \longrightarrow X'$ un morphisme de $S$-schémas, et reprenons les notations de 1.3.5.

Le morphisme (1.3.6.1) définit

\begin{equation}\tag{2.1.7.1}
R\Psi\, Rf_* \longrightarrow Rf_*\, R\Psi \quad .
\end{equation}
D'après le théorème de changement de base pour un morphisme propre (SGA4 XII 5.1), \underline{ce morphisme est un isomorphisme si} $f$ \underline{est propre.}

Le morphisme (1.3.7.1) définit

\begin{equation}\tag{2.1.7.2}
f^*\, R\Psi \longrightarrow R\Psi\, f^* \quad .
\end{equation}
D'après le théorème de changement de base pour un morphisme lisse (SGA4 XVI 1.2), \underline{ce morphisme est un isomorphisme si} $f$ \underline{est lisse.}

L'isomorphisme (2.1.7.1) pour $f$ propre et le morphisme (1.3.8.1)
\SGApage{II}{102}
définissent un morphisme

\SGAvar{\begin{equation}\tag{2.1.7.3}
Rf_! \, R\Psi \longrightarrow R\Psi \, Rf_!
\end{equation}
qui, pour $f$ propre, est l'inverse de (2.1.8.1).

}{\begin{equation}\tag{2.1.7.3}
Rf_! \, R\Psi \longrightarrow R\Psi \, Rf_!
\end{equation}
qui, pour $f$ propre, est l'inverse de (2.1.7.1).

}{F7II-XIII-014}Pour $f$ quasi-fini, (1.3.9.1) se dérive en

\SGAvar{\begin{equation}\tag{2.1.7.4}
R\Psi \, Rf^! \longrightarrow Rf^! \, R\Psi
\end{equation}
qui, pour $f$ étale, est l'inverse de (2.1.7.1).

}{\begin{equation}\tag{2.1.7.4}
R\Psi \, Rf^! \longrightarrow Rf^! \, R\Psi
\end{equation}
qui, pour $f$ étale, est l'inverse de (2.1.7.2).

}{F7II-XIII-015}Soit un changement de trait $f : S' \longrightarrow S$ et $X' = X \times_S S'$

\begin{center}
\begin{tikzcd}
X' \arrow[r, "f"] \arrow[d, "p"] \& X \arrow[d, "p"] \\
S' \arrow[r, "f"] \& S
\end{tikzcd}
\end{center}








\SGAvar{(pour l'ambiguité dans les notations, cf. 1.3.5). On suppose $f$ surjectif. La flèche de changement de base (1.3.10.1) se dérive en

}{(pour l'ambiguïté dans les notations, cf. 1.3.5). On suppose $f$ surjectif. La flèche de changement de base (1.3.10.1) se dérive en

}{F7II-XIII-016}\begin{equation}\tag{2.1.7.5}
f^* \, R\Psi \longrightarrow R\Psi \, f^* \quad .
\end{equation}
Les mêmes énoncés valent pour $R\Psi_\eta$ .

2.1.8. Le cas essentiel de 2.1.7.1 et 2.1.7.3 est celui où $X' = S$ . On trouve alors pour $K \in \mathrm{Ob}\, D^+(X_{\bar{\eta}} , \Lambda)$ des morphismes

\begin{equation}\tag{2.1.8.1}
Rf_* \longrightarrow Rf_* \, R\Psi
\end{equation}
\begin{equation}\tag{2.1.8.2}
Rf_! \, R\Psi \longrightarrow Rf_!
\end{equation}
d'où

\begin{equation}\tag{2.1.8.3}
\mathbb{H}^i(X_{\bar{\eta}}, K) \longrightarrow \mathbb{H}^i(X_{\bar{s}}, R\Psi_{\bar{\eta}}(K))
\end{equation}
\begin{equation}\tag{2.1.8.4}
\mathbb{H}^i_c(X_{\bar{s}}, R\Psi_{\bar{\eta}}(K)) \longrightarrow \mathbb{H}^i_c(X_{\bar{\eta}}, K)
\end{equation}
et la commutativité des diagrammes 


\SGApage{II}{103}

\begin{equation}\tag{2.1.8.5}
\begin{tikzcd}
(R^i f_* K)_{\bar{s}} \arrow[r, "\mathrm{sp}"] \arrow[d] \& \mathbb{H}^i(X_{\bar{\eta}}, K) \arrow[d] \\
\mathbb{H}^i(X_{\bar{s}}, K) \arrow[r] \& \mathbb{H}^i(X_{\bar{s}}, R\Psi_{\bar{\eta}}(K))
\end{tikzcd}
\end{equation}













\begin{equation}\tag{2.1.8.6}
\begin{tikzcd}[column sep=large]
\mathbb{H}^i_c(X_{\bar{s}},K) \arrow[rr, "\mathrm{sp}"] \arrow[dr] \& {} \& \mathbb{H}^i_c(X_{\bar{\eta}}, K) \\
{} \& \mathbb{H}^i_c(X_{\bar{s}}, R\Psi_{\bar{\eta}}(K)) \arrow[ur] \&
{} \end{tikzcd}\quad .
\end{equation}













Pour $K \in \mathrm{Ob}\, D^+(X_\eta, \Lambda)$ et $\sigma \in I$ , les diagrammes suivants sont commutatifs

\SGAvar{\begin{equation}\tag{2.1.8.7}
\begin{tikzcd}
\mathbb{H}^i(X_{\bar{\eta}}, K) \arrow[rr, dash, "\sigma-1"] \arrow[d] \& {} \& \mathbb{H}^i(X_{\bar{\eta}}, K) \arrow[d] \\
\mathbb{H}^i(X_{\bar{s}}, R\Psi_{\bar{\eta}}(K)) \arrow[r] \& \mathbb{H}^i(X_{\bar{s}}, R\Phi_{\bar{\eta}}(K)) \arrow[r, "\mathrm{Var}(\sigma)"] \& \mathbb{H}^i(X_{\bar{s}}, R\Psi_{\bar{\eta}}(K))
\end{tikzcd}
\end{equation}














\begin{equation}\tag{2.1.8.8}
\begin{tikzcd}
\mathbb{H}^i_c(X_{\bar{\eta}}, K) \arrow[rr, "\sigma-1"] \& {} \& \mathbb{H}^i_c(X_{\bar{\eta}}, K) \\
\mathbb{H}^i_c(X_s, R\Psi_{\bar{\eta}}(K)) \arrow[r] \arrow[u] \& \mathbb{H}^i_c(X_{\bar{s}}, R\Phi_{\bar{\eta}}(K) \arrow[r, "\mathrm{Var}(\sigma)"] \& \mathbb{H}^i_c(X_{\bar{s}}, R\Psi_{\bar{\eta}}(K)) \arrow[u]
\end{tikzcd}\quad .
\end{equation}

















}{\begin{equation}\tag{2.1.8.7}
\begin{tikzcd}
\mathbb{H}^i(X_{\bar{\eta}}, K) \arrow[rr, "\sigma-1"] \arrow[d] \& {} \& \mathbb{H}^i(X_{\bar{\eta}}, K) \arrow[d] \\
\mathbb{H}^i(X_{\bar{s}}, R\Psi_{\bar{\eta}}(K)) \arrow[r] \& \mathbb{H}^i(X_{\bar{s}}, R\Phi_{\bar{\eta}}(K)) \arrow[r, "\mathrm{Var}(\sigma)"] \& \mathbb{H}^i(X_{\bar{s}}, R\Psi_{\bar{\eta}}(K))
\end{tikzcd}
\end{equation}














\begin{equation}\tag{2.1.8.8}
\begin{tikzcd}
\mathbb{H}^i_c(X_{\bar{\eta}}, K) \arrow[rr, "\sigma-1"] \& {} \& \mathbb{H}^i_c(X_{\bar{\eta}}, K) \\
\mathbb{H}^i_c(X_{\bar{s}}, R\Psi_{\bar{\eta}}(K)) \arrow[r] \arrow[u] \& \mathbb{H}^i_c(X_{\bar{s}}, R\Phi_{\bar{\eta}}(K)) \arrow[r, "\mathrm{Var}(\sigma)"] \& \mathbb{H}^i_c(X_{\bar{s}}, R\Psi_{\bar{\eta}}(K)) \arrow[u]
\end{tikzcd}\quad .
\end{equation}

















}{F7II-XIII-017}Pour $f$ propre, (2.1.8.1) et (2.1.8.2), ainsi que (2.1.8.3) et (2.1.8.4) sont des isomorphismes inverses l'un de l'autre. Le triangle distingué (2.1.2.4) fournit une suite exacte longue

\begin{equation}\tag{2.1.8.9}
\ldots \longrightarrow \mathbb{H}^i(X_{\bar{s}}, K) \longrightarrow \mathbb{H}^i(X_{\bar{\eta}}, K) \longrightarrow \mathbb{H}^i(X_{\bar{s}}, R\Phi_{\bar{\eta}}(K)) \longrightarrow \ldots
\end{equation}
Le complexe évanescent $R\Phi(K)$ exprime donc la différence entre les cohomologies des fibres géométriques $X_{\bar{s}}$ et $X_{\bar{\eta}}$ (à valeurs dans l'image réciproque de $K$ ). D'après (2.1.8.7), la variation $\mathrm{Var}(\sigma)$ détermine l'action du groupe d'inertie $I$
\SGApage{II}{104}
sur $\mathbb{H}^i(X_{\bar\eta}, K)$ . Ces faits sont à la base de toutes les applications de la méthode des cycles évanescents.

L'hypothèse de propreté faite ci-dessus sur $f$ peut se remplacer par des conditions de régularité à l'infini sur $f$ et $K$ . On a par exemple le résultat suivant (pour $K \in \mathrm{Ob}\ D(X_{\bar\eta}, \Lambda)$ ).

\underline{Proposition} 2.1.9. \underline{Supposons que} $f$ \underline{admette une factorisation} $f = f_1k$

\begin{center}
\begin{tikzcd}[column sep=huge]
X \arrow[r, hook, "k"] \& X_1 \arrow[r, "f_1"] \& S \;,
\end{tikzcd}
\end{center}






\underline{avec} $f_1$ \underline{propre, et} $k$ \underline{l'inclusion dans} $X_1$ \underline{du complément d'un diviseur à croisements normaux relatif} $D$ (\underline{par hypothèse,} $f_1$ \underline{est donc lisse dans un voisinage de} $D$ ). \underline{Si} , \underline{dans un voisinage} $U$ \underline{de} $D$ , \underline{les} $\underline{H}^i(K)$ \underline{sont localement constants sur} $U-(D \cup X_o)$ , \underline{et modérément ramifiés le long du diviseur à croisements normaux} $D \cup X_o$ , \underline{alors} (2.1.8.1) \underline{et} (2.1.8.2) \underline{sont des isomorphismes.}

\SGAvar{La suite spectrale d'hypercohomologie de $K$ montre qu'il suffit de traiter le cas où $K$ est réduit à un faisceau $F_\eta$ en degré $0$ . Remplaçant $S$ par son normalisé dans une extension finie (ce qui est loisible, comme on le voit sur les définitions), et appliquant le lemme d'Abhyankhar, on peut supposer que $F_\eta$ est restriction à $X_\eta$ d'un faisceau $F$ sur $X$ , localement constant sur $U \cap X$ et modérément ramifié le long de $D$ . Avec les notations de 2.1.2, soit $\overline{F}$ l'image réciproque de $F$ sur $\overline{X}$ et $\bar f : \overline{X} \to \overline{S}$ .

}{La suite spectrale d'hypercohomologie de $K$ montre qu'il suffit de traiter le cas où $K$ est réduit à un faisceau $F_\eta$ en degré $0$ . Remplaçant $S$ par son normalisé dans une extension finie (ce qui est loisible, comme on le voit sur les définitions), et appliquant le lemme d'Abhyankar, on peut supposer que $F_\eta$ est restriction à $X_\eta$ d'un faisceau $F$ sur $X$ , localement constant sur $U \cap X$ et modérément ramifié le long de $D$ . Avec les notations de 2.1.2, soit $\overline{F}$ l'image réciproque de $F$ sur $\overline{X}$ et $\bar f : \overline{X} \to \overline{S}$ .

}{F7II-XIII-018}Pour prouver que (2.1.8.1) est un isomorphisme, il suffit de prouver que $R\bar f_* (R\bar j_* \bar j^* \overline{F})$ commute au changement de base $\bar s \to \overline{S}$ . Dans le triangle

\begin{equation*}
\longrightarrow \overline{F} \longrightarrow R\bar j_* \bar j^* \overline{F} \longrightarrow \Delta \longrightarrow
\end{equation*}
les $\underline{H}^i(\Delta)$ sont à support dans $X_{\bar s}$ ; pour $\Delta$ , le résultat de changement de base considéré est trivial. Il suffit donc de prouver que $R\bar f_*(\overline{F})$ commute au changement de base $\bar s \to \overline{S}$ soit, vu le théorème de changement de base pour $f_1$ , que
\SGApage{II}{105}
$Rk_*(\overline{F})$ commute au changement de base $X_{1\bar s} \to X_1$ (2.1.10 ci-dessous).

Pour prouver que (2.1.8.2) est un isomorphisme, il suffit de prouver que, pour le faisceau $F$ , le prolongement par $0$ $k_!$ commute au foncteur $R\bar j_*$ , i.e. que, dans un voisinage de $D$ ,

$R\Phi(k_! F) = 0$ sur $X_{1\bar s}$ ( 2.1.11 ci-dessous).

\SGAvar{\underline{Lemme} 2.1.10. \underline{Soit} $f : X \to S$ \underline{un morphisme lisse,} $D \subset X$ \underline{un diviseur à croisements normaux relatif,} $U = X - D \stackrel{j}{\hookrightarrow} X$ \underline{et} $F$ \underline{un faisceau de} $\Lambda$\underline{-module sur} $U$ , \underline{localement constant et modérément ramifié le long de} $D$ . \underline{La formation de} $Rj_* F$ \underline{commute à tout changement de base} $T \to S$ .

La question est locale (pour la topologie étale) sur $X$ . On peut donc supposer qu'il existe $g : X' \to X$ , normalisé de $X$ dans un revêtement fini étale surjectif $g' : U' \to U$ modérément ramifié le long de $D$ , tel que $g'^* F$ soit constant sur $U'$ . Utilisant le lemme d'Abhyankhar, on peut supposer que de plus $X'$ est lisse sur $S$ et que $U'$ est le complément d'un diviseur à croisements normaux relatif $D'$ dans $X'$ .

}{\underline{Lemme} 2.1.10. \underline{Soit} $f : X \to S$ \underline{un morphisme lisse,} $D \subset X$ \underline{un diviseur à croisements normaux relatif,} $U = X - D \stackrel{j}{\hookrightarrow} X$ \underline{et} $F$ \underline{un faisceau de} $\Lambda$\underline{-modules sur} $U$ , \underline{localement constant et modérément ramifié le long de} $D$ . \underline{La formation de} $Rj_* F$ \underline{commute à tout changement de base} $T \to S$ .

La question est locale (pour la topologie étale) sur $X$ . On peut donc supposer qu'il existe $g : X' \to X$ , normalisé de $X$ dans un revêtement fini étale surjectif $g' : U' \to U$ modérément ramifié le long de $D$ , tel que $g'^* F$ soit constant sur $U'$ . Utilisant le lemme d'Abhyankar, on peut supposer que de plus $X'$ est lisse sur $S$ et que $U'$ est le complément d'un diviseur à croisements normaux relatif $D'$ dans $X'$ .

}{F7II-XIII-019}$F$ s'injecte dans $g'_* g'^* F$ et le conoyau vérifie encore les hypothèses de 2.1.10. Itérant la construction (le même $g$ continue d'ailleurs à marcher), on obtient une résolution de $F$ par des faisceaux de la forme $g'_* G$ avec $G$ constant et $g'$ comme plus haut, il suffit donc de démontrer 2.1.10 pour un tel faisceau, i.e. de démontrer 2.1.10 pour $(X'/S, D', G$ constant sur $U'$ ) : c'est là un cas particulier de SGA5 II 4.4.

\underline{Lemme} 2.1.11. \underline{Soient} $S$ \underline{un trait hensélien,} $f : X \to S$ \underline{un morphisme lisse,} $D \subset X$ \underline{un diviseur à croisements normaux relatif,} $U = X - D \stackrel{j}{\hookrightarrow} X$ \underline{et} $F$ \underline{un faisceau de} $\Lambda$\underline{-modules sur} $U$ , \underline{localement constant et modérément ramifié le long de} $D$ . \underline{On a}

\begin{equation*}
R\Phi(j_! F) = 0 \;.
\end{equation*}

\SGApage{II}{106}

\SGAvar{Le dévissage utilisé en 2.1.11 nous ramène au cas où $F$ est constant, de valeur $G$ . On peut supposer $D$ réunion de diverseurs lisses $(D_i)_{i \in I}$ . Pour $J \subset I$ , soit $D_J$ l'intersection des $D_i$ pour $i \in J$ . Le faisceau $F$ admet une résolution $F^{\cdot}$ , de composantes des sommes de faisceaux $G_J$ , constants sur $D_J$ et prolongés par $0$ sur $X$ . Il suffit donc de prouver 2.1.11 pour un tel faisceau, ce qui résulte de (2.1.5) appliqué aux $D_J/S$ .

}{Le dévissage utilisé en 2.1.10 nous ramène au cas où $F$ est constant, de valeur $G$ . On peut supposer $D$ réunion de diviseurs lisses $(D_i)_{i \in I}$ . Pour $J \subset I$ , soit $D_J$ l'intersection des $D_i$ pour $i \in J$ . Le faisceau $F$ admet une résolution $F^{\cdot}$ , de composantes des sommes de faisceaux $G_J$ , constants sur $D_J$ et prolongés par $0$ sur $X$ . Il suffit donc de prouver 2.1.11 pour un tel faisceau, ce qui résulte de (2.1.5) appliqué aux $D_J/S$ .

}{F7II-XIII-020}Je conjecture que 2.1.7.5 est un isomorphisme si $S$ et $S'$ sont excellents. Voici un résultat partiel (voir aussi 2.4.2).

\underline{Proposition} 2.1.12. \underline{Le morphisme} 2.1.7.5 \underline{est un isomorphisme dans les cas suivants}

(a) $S$ \underline{est excellent, de complété} $S'$ ;

(b) $S$ \underline{et} $S'$ \underline{sont d'égale caractéristique} $0$ .

Dans le cas (b), on se ramène à supposer $S$ et $S'$ strictement locaux, et qu'une uniformisante pour $S$ est encore une uniformisante pour $S'$ (étendre les scalaires sans changer $\bar{S}' \to \bar{S}$ ). Dans les cas (a) et (b), notant par $\bar{\phantom{x}}$ le normalisé dans une clôture séparable du corps des fractions, on a alors $\bar{S}' \simeq S' \times_s \bar{S}$ , et, revenant aux définitions, on trouve qu'il suffit de prouver que $S' \to S$ est universellement localement acyclique (appliquer SGA4 XVI 1.1). Ceci résulte de I 0.5.1 (basé sur la désingularisation de Néron à la Artin) et de SGA4 XV 2.1 (théorème d'acyclicité locale des morphismes lisses).

2.1.13. Supposons $f : X \to S$ de type fini. La dimension cohomologique (SGA4 XVII 1.2.9) de $R\Psi$ est alors au plus égale à $\dim(X_\eta)$ , ainsi qu'il résulte de 2.1.4 et de SGA4 XIV 3.2 (cf. I 4.2.(i)). Ceci permet d'étendre $R\Psi$ à la catégorie dérivée entière

\begin{equation*}
R\Psi : D(X, \Lambda) \longrightarrow D(X_s \times_s S, \Lambda) \quad ,
\end{equation*}
et $R\Psi$ transforme $D^-$ en $D^-$ . Soit $L \in \mathrm{Ob}\ D^-(\Lambda)$ un complexe borné supérieurement de $\Lambda$-modules à droite, et $K \in \mathrm{Ob}(D^-(X, \Lambda))$ . On dispose alors de 


\SGApage{II}{107}

\SGAvar{\begin{equation}\tag{2.1.13.1}
 L \overset{\mathbf{L}}{\otimes}_{\Lambda} R\Psi(K) \xrightarrow{\ \sim\ } R\Psi(L \overset{\mathbf{L}}{\otimes}_{\Lambda} K) \quad .
\end{equation}
Pour vérifier que 2.1.13.1 est un isomorphisme, on remplace tout d'abord $L$ par un complexe de $\Lambda$-modules libres, et on note qu'il suffit, car $R\Psi$ est triangulé de dimension cohomologique finie, de le vérifier après avoir remplacé $L$ par une des ses composantes. Par passage à la limite inductive, on se ramène au cas où $L$ est un module libre de type fini, puis au cas trivial où $L = \Lambda$ .

}{\begin{equation}\tag{2.1.13.1}
 L \overset{\mathbf{L}}{\otimes}_{\Lambda} R\Psi(K) \xrightarrow{\ \sim\ } R\Psi(L \overset{\mathbf{L}}{\otimes}_{\Lambda} K) \quad .
\end{equation}
Pour vérifier que 2.1.13.1 est un isomorphisme, on remplace tout d'abord $L$ par un complexe de $\Lambda$-modules libres, et on note qu'il suffit, car $R\Psi$ est triangulé de dimension cohomologique finie, de le vérifier après avoir remplacé $L$ par une de ses composantes. Par passage à la limite inductive, on se ramène au cas où $L$ est un module libre de type fini, puis au cas trivial où $L = \Lambda$ .

}{F7II-XIII-021}De (2.1.13.1) on déduit par un argument familier (cf. SGA4 XVII 5.2.11) que si $K \in \mathrm{Ob}\ D (X_\eta, \Lambda)$ est de Tor dimension finie $\leq d$ ($d \in \mathbb{Z}$) (SGA4 XVII 4.1.9), alors $R\Psi_\eta(K)$ est de même de Tor-dimension $\leq d$ .


2.2. \underline{Une compatibilité (morphismes traces).}

\SGAvar{2.2.1. Soient $S$ un trait hensélien et $u : X_s \to Y_s$ un morphisme quasi-fini plat séparé de présentation finie de $s$-schémas. Pour tout faisceau abélien $F$ sur $Y_s$ , un morphisme trace $\mathrm{Tr}_u : u_! u F \to F$ est défini (SGA4 XVII 6.2.3). Pour $F$ un faisceau abélien sur $Y_s \times_s S$ ou $Y_s \times_s \eta$ , on en déduit un morphisme analogue $\mathrm{Tr}_u : u_! u^* F \to F$ .

}{2.2.1. Soient $S$ un trait hensélien et $u : X_s \to Y_s$ un morphisme quasi-fini plat séparé de présentation finie de $s$-schémas. Pour tout faisceau abélien $F$ sur $Y_s$ , un morphisme trace $\mathrm{Tr}_u : u_! u^* F \to F$ est défini (SGA4 XVII 6.2.3). Pour $F$ un faisceau abélien sur $Y_s \times_s S$ ou $Y_s \times_s \eta$ , on en déduit un morphisme analogue $\mathrm{Tr}_u : u_! u^* F \to F$ .

}{F7II-XIII-022}Soit $u : X \to Y$ un morphisme quasi-fini plat séparé de présentation finie de $S$-schémas. Nous noterons encore $u$ le morphisme induit $X_s \times_s S \to Y_s \times_s S$ et nous noterons $\Psi$ les foncteurs 1.3.3 pour $X$ et pour $Y$ .

\underline{Proposition} 2.2.2. \underline{Pour tout faisceau abélien} $F$ \underline{sur} $Y$ , \underline{le diagramme suivant est commutatif}

\begin{center}
\begin{tikzcd}[column sep=large, row sep=large]
u_! u^* \Psi F \arrow[r, "{(1.3.7.1)}"] \arrow[d, "\mathrm{Tr}_u"] \& u_! \Psi u^* F \arrow[r, "{(1.3.8.1)}"] \& \Psi u_! u^* F \arrow[d, "\mathrm{Tr}_u"] \\
\Psi F \arrow[rr, equal] \& {} \& \Psi F
\end{tikzcd}\quad .
\end{center}










\SGApage{II}{108}

La vérification est laissée au lecteur. On a une variante pour $\Psi_\eta$ et $F$ sur $Y_\eta$ .

\SGAvar{2.2.3. Soient $\Lambda$ comme en 2.2.1 et $u : X \to Y$ un morphisme quasi-fini et plat de $S$-schémas séparés de type fini et plats

}{2.2.3. Soient $\Lambda$ comme en 2.1.1 et $u : X \to Y$ un morphisme quasi-fini et plat de $S$-schémas séparés de type fini et plats

}{F7II-XIII-023}\begin{center}
\begin{tikzcd}[column sep=large, row sep=large]
X \arrow[rr, "u"] \arrow[dr, "f"'] \& {} \& Y \arrow[dl, "g"] \\
 {} \& S \&
{} \end{tikzcd}\quad .
\end{center}






On note encore $f$ et $g$ les morphismes $X_s \times_s S \to S$ et $Y_s \times_s S \to S$ . On suppose que $\dim Y_\eta \leq n$ .

Pour tout $K \in \mathrm{Ob}\, D(S, \Lambda)$ , on a défini un morphisme trace (SGA4 XVIII 2.9)

\begin{equation*}
\mathrm{Tr}_f : Rf_! \, f^* K(n) [2n] \longrightarrow K \; .
\end{equation*}

\underline{Corollaire} 2.2.4. \underline{Le diagramme suivant est commutatif}

\begin{center}
\begin{tikzcd}[column sep=huge, row sep=huge]
Rg_! u_! u^* R\Psi\, g^* K(n)[2n] \arrow[r, "(1.3.7.1)"] \arrow[d, "\mathrm{Tr}_u"] \& Rf_! \, R\Psi\, f^* K(n)[2n] \arrow[r, "(2.1.10.1)"] \& Rf_! \, f^* K(n)[2n] \arrow[d, "\mathrm{Tr}_f"] \\
Rg_! \, R\Psi\, g^* K(n)[2n] \arrow[r, "(2.1.10.1)"] \& Rg_! \, g^* K(n) [2n] \arrow[r, "\mathrm{Tr}_g"] \& K
\end{tikzcd}
\end{center}











\underline{Si} $u$ \underline{est étale,} (1.3.7.1) \underline{est un isomorphisme.}

\SGAvar{Le diagramme (2.2.4) se déploye en le suivant, où les arguments sont omis

}{Le diagramme (2.2.4) se déploie en le suivant, où les arguments sont omis

}{F7II-XIII-024}\begin{center}
\begin{tikzcd}[column sep=large, row sep=huge]
Rg_! u_! u^* R\Psi\, g^* \arrow[r] \arrow[d, "\mathrm{Tr}_u"'] \arrow[rd, phantom, "(1)" description] \& Rg_! \, R\Psi\, u_! u^* g^* \arrow[r] \arrow[d, "\mathrm{Tr}_u"'] \arrow[rd, phantom, "(2)" description] \& Rg_! u_! u^* g^* \arrow[r, "\sim"] \arrow[d, "\mathrm{Tr}_u"'] \arrow[rd, phantom, "(3)" description] \& Rf_! f^* \arrow[d, "\mathrm{Tr}_f"] \\
Rg_! \, R\Psi\, g^* \arrow[r, equal] \& Rg_! \, R\Psi\, g^* \arrow[r] \& Rg_! \, g^* \arrow[r, "\mathrm{Tr}_g"] \& K
\end{tikzcd}\quad .
\end{center}
















Que (1) soit commutatif résulte de 2.2.2. Que (2) le soit est la
\SGApage{II}{109}
fonctorialité de (2.1.10.2). Que (3) le soit est la transitivité de la trace SGA4 XVIII 2.9. Var.3.

2.2.4.1.On a un résultat analogue pour $\Psi_\eta$ et $K \in \mathrm{Ob}\, D(\eta,\Lambda)$ .

2.2.5. Soit $\Phi$ une partie de $Y_s$ propre sur $s$ , et notons $Rg_\Phi$ le foncteur ``image directe à support dans $\Phi$ ''. Notons $\mathrm{Tr}_g$ le morphisme composé

\SGAvar{\begin{equation}\tag{2.2.5.1}
\mathrm{Tr}_g : Rf_\Phi \, R\Psi g^* K(n)[2n] \to Rg_! R\Psi g^* K(n)[2n] \to Rg_! g^* K(n)[2n] \to K \; ,
\end{equation}
et son analogue pour $\Psi_\eta$ et $K \in D(\eta, \Lambda)$ . Si $\Phi$ est un ensemble fini de points fermés, on a

}{\begin{equation}\tag{2.2.5.1}
\mathrm{Tr}_g : Rg_\Phi \, R\Psi g^* K(n)[2n] \to Rg_! R\Psi g^* K(n)[2n] \to Rg_! g^* K(n)[2n] \to K \; ,
\end{equation}
et son analogue pour $\Psi_\eta$ et $K \in D(\eta, \Lambda)$ . Si $\Phi$ est un ensemble fini de points fermés, on a

}{F7II-XIII-025}\begin{equation*}
Rg_\Phi \circ u_! \; = \; Rf_{u^{-1}\Phi} \; ,
\end{equation*}
d'où un morphisme trace $\mathrm{Tr}_u : Rf_{u^{-1}\Phi} u^* = Rg_\Phi \, u_! \, u^* \to Rg_\Phi$ .

D'après 2.2.4, le diagramme suivant est commutatif.

\begin{center}
\begin{tikzcd}[column sep=large, row sep=huge]
Rf_{u^{-1}\Phi} \, u^* R\Psi\, g^* K(n)[2n] \arrow[r] \arrow[d, "\mathrm{Tr}_u"] \& Rf_{u^{-1}\Phi} \, R\Psi\, f^* K(n) [2n] \arrow[r, "\mathrm{Tr}_f"] \& K \arrow[d, equal] \\
Rg_\Phi \, R\Psi\, g^* K(n) [2n] \arrow[rr, "\mathrm{Tr}_g"] \& {} \& K
\end{tikzcd}
\end{center}










Supposons $u$ étale, soit $x \in X_s$ et supposons que $x$ et $y = u(x)$ aient même corps résiduel. On a alors

\SGAvar{\begin{equation*}
Rf_{\{x\}} \; \overset{\sim}{\longrightarrow} \; Rg_{\{x\}} \; ;
\end{equation*}
cet isomorphisme, donné par $\mathrm{Tr}_u$ , exprime la nature locale pour la cohomologie étale de la cohomologie à support dans $x$ . Le diagramme 

}{\begin{equation*}
Rf_{\{x\}} \; \overset{\sim}{\longrightarrow} \; Rg_{\{y\}} \; ;
\end{equation*}
cet isomorphisme, donné par $\mathrm{Tr}_u$ , exprime la nature locale pour la cohomologie étale de la cohomologie à support dans $x$ . Le diagramme 

}{F7II-XIII-026}\begin{equation}\tag{2.2.5.2}
\begin{tikzcd}[column sep=huge, row sep=huge]
Rf_{\{x\}} \, R\Psi\, f^* K(n)[2n] \arrow[r, "\mathrm{Tr}_f"] \arrow[d, equal] \& K \arrow[d, equal] \\
Rg_{\{y\}} \, R\Psi\, g^* K(n)[2n] \arrow[r, "\mathrm{Tr}_g"] \& K
\end{tikzcd}
\end{equation}










\SGApage{II}{110}

 est commutatif.

\SGAvar{\underline{Scholie} 2.2.6.  La compatibilité 2.3.5.2 exprime que le morphisme trace (2.3.5.1) pour $\Phi = \{x\}$ , est un invariant local: il ne dépend que du $S$-schéma hensélisé de $X$ en $x$ . De même pour $R\Psi_{\eta}$ .


}{\underline{Scholie} 2.2.6.  La compatibilité 2.2.5.2 exprime que le morphisme trace (2.2.5.1) pour $\Phi = \{y\}$ , est un invariant local: il ne dépend que du $S$-schéma hensélisé de $X$ en $x$ . De même pour $R\Psi_{\eta}$ .


}{F7II-XIII-027}2.3.  \underline{Théorème de finitude (en égale caractéristique zéro).}

\underline{Théorème} 2.3.1.  \underline{Soient} $S$ \underline{un trait hensélien comme en} 0.2.5, $\Lambda$ \underline{un anneau de torsion noethérien et} $f : X \to S$ \underline{un morphisme de type fini. On suppose que} $S$ \underline{est} \underline{d}'égale caractéristique zéro . \underline{Alors, pour tout faisceau constructible de} $\Lambda$ \underline{-modules} $F$ \underline{sur} $X_{\eta}$ , \underline{les faisceaux} $R^{i}\Psi_{\bar{\eta}}(F)$ \underline{sur} $X_{\bar{s}}$ \underline{sont constructibles.}

On commence par recopier le début de la démonstration de SGA4 XIV 1.10 : appliquons SGA4 IX 2.14(ii). Il existe une résolution $F^{\bullet}$ de $F$

\begin{center}
\begin{tikzcd}
 {} \& F \arrow[d] \& {} \& {} \\
0 \arrow[r] \& F^{0} \arrow[r] \& F^{1} \arrow[r] \& \ldots
\end{tikzcd}
\end{center}






par des faisceaux de la forme $\underset{i}{\oplus} p'_{i*} G_{i}$ , où $p'_{i} : Y_{i\eta} \to X_{\eta}$ est fini et où $G_{i}$ est un faisceau constant constructible. On peut supposer $Y_{i\eta}$ réduit; soit $p_{i} : Y_{i} \to X$ le normalisé de $X$ dans $Y_{i\eta}$ . Utilisons la suite spectrale de la résolution $F^{\bullet}$ et appliquons 2.1.7.1 dans le cas trivial du morphisme fini $p_{i}$ ; on trouve qu'il suffit de prouver 2.3.1 pour $(Y_{i}, G_{i})$ .

Prouvons 2.3.1 pour $F$ constant défini par un $\Lambda$ -module $M$ . Rappelons que $\Lambda$ est une $\mathbb{Z}/n$-algèbre pour $n$ convenable. Le module $M$ se dévisse donc en modules $M_{i}$ sur $\Lambda/\ell_{i}\Lambda$ , avec $\ell_{i}|n$ premier convenable. Ceci nous ramène au cas où $\Lambda$ est une $\mathbb{Z}/\ell$-algèbre, $\ell$ premier. On a alors (2.1.13.1)

\begin{equation*}
M \otimes R^{i}\Psi_{\bar{\eta}}(\mathbb{Z}/\ell) \overset{\sim}{\longrightarrow} R^{i}\Psi_{\bar{\eta}}(M) \quad ; 
\end{equation*}

\SGApage{II}{111}

 il reste à traiter le cas où $F$ est le faisceau constant $\mathbb{Z}/\ell$ .

Soit $X_{\eta} = X'_{0} \supset \ldots \supset X'_{n}$ une suite de parties fermées de $X_{\eta}$ telle que $(X'_{i} - X'_{i+1})_{red}$ soit lisse sur $\eta$ . Soit $X_{i} = \bar{X}'_{i}$ . Le faisceau $\mathbb{Z}/\ell$ se dévisse en les faisceaux $\mathbb{Z}/\ell$ sur $(X'_{i} - X'_{i+1})$ , prolongés par $0$ . On en tire la

\underline{Réduction 1}  Il suffit de prouver 2.3.1 pour un faisceau $j_{!} \mathbb{Z}/\ell$ , pour $j : U \hookrightarrow X_{\eta}$ l'inclusion d'un ouvert lisse.

D'après la résolution des singularités, il existe $g : X_{1} \to X$ tel que

(a)  $g$ est propre, et induit un isomorphisme $g^{-1}(U) \overset{\sim}{\longrightarrow} U$ ;

(b)  $X_{1}$ est régulier et $X_{1} - g^{-1}(U)$ un diviseur à croisements normaux dans $X_{1}$ , réunion de diviseurs réguliers $D_{\alpha}$ .

Soit $j_{1} : g^{-1}(U) \hookrightarrow X_{1}$ . On a $Rg_{*} j_{1!} \mathbb{Z}/\ell = j_{!} \mathbb{Z}/\ell$ . Appliquons 2.1.7.1 et le théorème de finitude SGA4 XIV 1.1.; on trouve qu'il suffit de prouver (2.3.1) pour $(X_{1}, j_{1!} \mathbb{Z}/\ell)$ .

On a une résolution

\begin{equation*}
0 \longrightarrow j_{1!}\mathbb{Z}/\ell \longrightarrow \mathbb{Z}/\ell \longrightarrow \underset{\alpha}{\oplus} (\mathbb{Z}/\ell)_{D_{\alpha}} \longrightarrow \underset{\alpha,\beta}{\oplus} (\mathbb{Z}/\ell)_{D_{\alpha} \cap D_{\beta}} \longrightarrow \ldots \; ,
\end{equation*}
d'où la

\underline{Réduction 2}  Il suffit de prouver 2.3.1 pour le faisceau constant $\mathbb{Z}/\ell$ et un $S$-schéma régulier $X$ tel que $X_{s}$ soit un diviseur à croisements normaux dans $X$ .

Dans le cas auquel on s'est réduit, on utilise le calcul explicite I 3.3 (qui utilise le théorème de pureté SGA4 XIX 3.2 et la caractéristique 0).


2.4.  \underline{Singularités isolées.}

\SGAvar{2.4.1.  Dans ce numéro, nous supposons pour simplifier que $S$ est un trait strictement hensélien. Soient $\Lambda$ un anneau noethérien de torsion premier à la caractéristique résiduelle de $S$ , $f : X \to S$ un morphisme de type fini et $F$ un faisceau constructible de $\Lambda$-modules sur $X$ . On dit que $x \in X$ est un point de lissité de
\SGApage{II}{112}
$(X, F)$ si $f$ est lisse en $X$ et que $F$ est localement constant au voisinage de $x$ . On note $\Sigma$ l'ensemble des points de non lissité de $(X, F)$ , et $\Sigma_s$ la fibre spéciale de $\Sigma$ . On sait que $R\Phi(F) = 0$ en dehors de $\Sigma_s$ .

}{2.4.1.  Dans ce numéro, nous supposons pour simplifier que $S$ est un trait strictement hensélien. Soient $\Lambda$ un anneau noethérien de torsion premier à la caractéristique résiduelle de $S$ , $f : X \to S$ un morphisme de type fini et $F$ un faisceau constructible de $\Lambda$-modules sur $X$ . On dit que $x \in X$ est un point de lissité de
\SGApage{II}{112}
$(X, F)$ si $f$ est lisse en $x$ et que $F$ est localement constant au voisinage de $x$ . On note $\Sigma$ l'ensemble des points de non lissité de $(X, F)$ , et $\Sigma_s$ la fibre spéciale de $\Sigma$ . On sait que $R\Phi(F) = 0$ en dehors de $\Sigma_s$ .

}{F7II-XIII-028}\underline{Théorème} 2.4.2. \underline{En tout point isolé} $x$ \underline{de} $\Sigma_s$ , \underline{les} $R^i\Phi_{\bar\eta}(F)_x$ \underline{sont des $\Lambda$-modules de type fini, invariants par changement de trait} (\underline{via} 2.1.7.1).

En égale caractéristique $0$ , ces assertions ont été démontrées en 2.1.12 et 2.3.1, sans supposer que $x$ soit isolé.

Soit $x$ un point isolé de $\Sigma_s$ . Le théorème est de nature locale au voisinage de $x$ ; on peut donc supposer que $f$ soit propre. On dispose alors de la suite exacte longue (2.1.8.9)

\begin{equation*}
\dots \longrightarrow H^i(X_s, F) \longrightarrow H^i(X_{\bar\eta}, F) \longrightarrow \mathbb{H}^i(X_s, R\Phi_{\bar\eta}(F)) \longrightarrow \dots
\end{equation*}
Dans cette suite, $H^i(X_s, F)$ et $H^i(X_{\bar\eta}, F)$ sont invariants par changement de trait et des $\Lambda$-modules de type fini (SGA4 XIV 1.1); les $\mathbb{H}^i(X_s, R\Phi_{\bar\eta}(F))$ partagent donc ces propriétés. Rappelons le

\underline{Lemme} 2.4.3. \underline{Soient} $(T, \Lambda)$ \underline{un topos annelé et} $F$ \underline{un fermé de} $T$ . \underline{Soit} $i$ \underline{l'inclusion de} $F$ . \underline{Alors,} $i_* D(F, \Lambda) \longrightarrow D(T, \Lambda)$ \underline{est une équivalence de catégorie de} $D(F, \Lambda)$ \underline{avec la sous-catégorie de} $D(T, \Lambda)$ \underline{formée des} $K$ \underline{tels que} $\underline{H}^i(K)$ \underline{soit à support dans} $F$ .

En effet, pour un tel complexe $K$ , $K \longrightarrow i_*i^*K$ est un quasi-isomorphisme.

On peut donc regarder $R\Phi_{\bar\eta}(F)$ comme un objet de $D^b(\Sigma_s, \Lambda) = D^b(\Sigma_s - \{x\}, \Lambda) + D^b(\{x\}, \Lambda)$ . Par passage à un facteur direct, on trouve donc que $\mathbb{H}^i(\{x\}, R\Phi_{\bar\eta}(F)) = R^i\Phi_{\bar\eta}(F)_x$ vérifie 2.4.2, ce qui achève la démonstration.

\underline{Remarque} 2.4.4. Avec la même démonstration, on voit que, pour $K$ constructible dans $D^+(X, \Lambda)$ et $x$ isolé dans le support de $R\Phi_{\bar\eta}(K)$ , le $\Lambda$-module $R\Phi_{\bar\eta}(K)_x$ est de type fini; si, après un changement de trait, $x$ reste un point isolé du support, ce module est invariant par ce changement de trait.


\SGApage{II}{113}

\SGAvar{2.4.5. Soit $x \in X_s$ un point de non lissité isolé de $(X, F)$ . Plus généralement, soit $K \in \mathrm{Ob}\ D^b(X, \Lambda)$ et soit $x$ un point isolé du support de $R\Phi(K)$ . Soit $i$ l'inclusion de $x$ dans $X_s$ . Le complexe $i_*i^* R\Phi(K)$ , isomorphe à $R\Phi(K)$ au voisinage de $x$ , est facteur direct de $R\Phi(F)$ (2.4.3). Pour $\sigma \in I$ , le morphisme variation induit donc

}{2.4.5. Soit $x \in X_s$ un point de non lissité isolé de $(X, F)$ . Plus généralement, soit $K \in \mathrm{Ob}\ D^b(X, \Lambda)$ et soit $x$ un point isolé du support de $R\Phi(K)$ . Soit $i$ l'inclusion de $x$ dans $X_s$ . Le complexe $i_*i^* R\Phi(K)$ , isomorphe à $R\Phi(K)$ au voisinage de $x$ , est facteur direct de $R\Phi(K)$ (2.4.3). Pour $\sigma \in I$ , le morphisme variation induit donc

}{F7II-XIII-029}\begin{equation*}
\mathrm{Var}(\sigma) : i_*i^* R\Phi_{\bar\eta}(K) \longrightarrow R\Psi_{\bar\eta}(K) \quad , \ \text{soit}
\end{equation*}
\begin{equation}\tag{2.4.5.1}
\mathrm{Var}(\sigma) : (R\Phi_{\bar\eta}(K))_x \longrightarrow R\Gamma_{\{x\}}(R\Psi_{\bar\eta}(K)) \quad .
\end{equation}
2.4.6. Supposons que le support $\Sigma$ de $R\Phi(K)$ soit formé d'un nombre fini de points isolés. Pour tout $\sigma \in I$ , l'endomorphisme $\sigma - 1$ de $R\Psi_{\bar\eta}(K)$ est alors le composé

\SGAvar{\begin{equation}\tag{2.4.6.1}
\begin{tikzcd}[column sep=huge, row sep=huge]
R\Psi_{\bar\eta}(K) \arrow[r] \arrow[d, "\sigma-1"] \& R\Phi_{\bar\eta}(K) \ =\ \underset{x\in\ \Sigma}{\oplus}\ i_{x*}(R\Phi_{\bar\eta}(K))_x \arrow[d, "{\oplus\, i_{x*}\ \mathrm{Var}(\sigma)}"] \\
R\Psi_{\bar\eta}(K) \& \underset{x\in\ \Sigma}{\oplus}\ i_x\, R\Gamma_{\{x\}}\, R\Psi_{\bar\eta}(K)\ \ . \arrow[l]
\end{tikzcd}
\end{equation}












}{\begin{equation}\tag{2.4.6.1}
\begin{tikzcd}[column sep=huge, row sep=huge]
R\Psi_{\bar\eta}(K) \arrow[r] \arrow[d, "\sigma-1"] \& R\Phi_{\bar\eta}(K) \ =\ \underset{x\in\ \Sigma}{\oplus}\ i_{x*}(R\Phi_{\bar\eta}(K))_x \arrow[d, "{\oplus\, i_{x*}\ \mathrm{Var}(\sigma)}"] \\
R\Psi_{\bar\eta}(K) \& \underset{x\in\ \Sigma}{\oplus}\ i_{x*}\, R\Gamma_{\{x\}}\, R\Psi_{\bar\eta}(K)\ \ . \arrow[l]
\end{tikzcd}
\end{equation}












}{F7II-XIII-030}Si $f$ est propre, appliquant à (2.4.6.1) le foncteur $H^i$ , on trouve que l'endomorphisme $\sigma-1$ de $\mathbb{H}^i(X_{\bar\eta}, K)$ est le composé

\begin{equation}\tag{2.4.6.2}
\begin{tikzcd}[column sep=large, row sep=huge]
\mathbb{H}^i(X_{\bar\eta}, K) \arrow[r, "\sim"] \arrow[d, "\sigma-1"] \& \mathbb{H}^i(X_s, R\Psi_{\bar\eta}(K)) \arrow[r] \& \mathbb{H}^i(X_s, R\Phi_{\bar\eta}(K)) \ =\ \underset{x\in\ \Sigma}{\oplus}\ R^i\Phi_{\bar\eta}(K)_x \arrow[d, "{\oplus\, \mathrm{Var}(\sigma)}"] \\
\mathbb{H}^i(X_{\bar\eta}, K) \& \mathbb{H}^i(X_s, R\Psi_{\bar\eta}(K)) \arrow[l, "\sim"'] \& \underset{x\, \in\, \Sigma}{\oplus}\ \mathbb{H}^i_{\{x\}}(X_s, R\Psi_{\bar\eta}(K)) \arrow[l]
\end{tikzcd}
\end{equation}


















\SGAvar{Ceci fournit une description de l'action de $I$ sur $\mathbb{H}^i(X_{\bar\eta}, K)$ d'ingrédient essentiel les morphismes variation (2.4.4.1), de nature locale sur $X$ . On a par ailleurs la suite exacte longue
\SGApage{II}{114}

}{Ceci fournit une description de l'action de $I$ sur $\mathbb{H}^i(X_{\bar\eta}, K)$ d'ingrédient essentiel les morphismes variation (2.4.5.1), de nature locale sur $X$ . On a par ailleurs la suite exacte longue
\SGApage{II}{114}

}{F7II-XIII-031}\begin{equation}\tag{2.4.6.3}
\longrightarrow \mathbb{H}^i(X_s, K) \longrightarrow \mathbb{H}^i(X_{\overline{\eta}}, K) \longrightarrow \underset{x \in \Sigma}{\oplus} R^i\Phi_{\overline{\eta}}(K)_x \longrightarrow \cdots
\end{equation}
Le même argument vaut pour $f$ non nécessairement propre si $(X, K)$ vérifie 2.1.9. On définit par le diagramme (2.4.6.5) ci-dessous un morphisme

\begin{equation}\tag{2.4.6.4}
\mathrm{Var}(\sigma) \ : \ R\Gamma(X_{\overline{\eta}}, K) \longrightarrow R\Gamma_c(X_{\overline{\eta}}, K) \ ,
\end{equation}
tel que l'endomorphisme $\sigma-1$ de $R\Gamma(X_{\overline{\eta}}, K)$ (resp. $R\Gamma_c(X_{\overline{\eta}}, K)$) soit le composé $\alpha \, \mathrm{Var}(\sigma)$ (resp. $\mathrm{Var}(\sigma)\alpha$) de $\mathrm{Var}(\sigma)$ et de $\alpha : R\Gamma_c(X_{\overline{\eta}}, K) \longrightarrow R\Gamma(X_{\overline{\eta}}, K)$ .

\begin{equation}\tag{2.4.6.5}
\begin{tikzcd}[column sep=huge]
R\Gamma (X_{\overline{\eta}} , K) \arrow[r, "\sim"] \arrow[d, "\mathrm{Var}(\sigma)"] \& R\Gamma(X_s, R\Psi_{\overline{\eta}}(K)) \arrow[r] \& \underset{x \in \Sigma}{\oplus} (R\Phi_{\overline{\eta}}(K))_x \arrow[d, "\oplus"] \& {} \\
R\Gamma_c(X_{\overline{\eta}}, K) \arrow[r, "\sim"] \& R\Gamma_c(X_s, R\Psi_{\overline{\eta}}(K)) \& \underset{x \in \Sigma}{\oplus} R\Gamma_{\{x\}}(X_s,R\Psi_{\overline{\eta}}(K)) \arrow[l] \& .
\end{tikzcd}
\end{equation}














On dispose aussi, dans ce cas, des suites exactes longues

\begin{equation}\tag{2.4.6.6}
\begin{tikzcd}
\cdots \arrow[r] \& \mathbb{H}^i_c(X_s, K) \arrow[r] \arrow[d] \& \mathbb{H}^i_c(X_{\overline{\eta}}, K) \arrow[r] \arrow[d] \& \underset{x \in \Sigma}{\oplus} (R^i\Phi_{\overline{\eta}}(K))_x \arrow[r] \arrow[d, equal] \& \cdots \\
\cdots \arrow[r] \& \mathbb{H}^i(X_s, K) \arrow[r] \& \mathbb{H}^i(X_{\overline{\eta}}, K) \arrow[r] \& \underset{x \in \Sigma}{\oplus} R^i\Phi_{\overline{\eta}}(K)_x \arrow[r] \& \cdots
\end{tikzcd}
\end{equation}











\SGAvar{2.4.6.   Soit $X_{s(x)}$ l'hensélisé de $X_s$ en $x$ et $X^{*}_{s(x)} = X_{s(x)} - \{x\}$ . La suite exacte longue d'hypercohomologie de $(X_{s(x)}, \{x\})$ à coefficient dans $K_s$ , $R\Psi_{\overline{\eta}}(K)$ et $R\Phi_{\overline{\eta}}(K)$ fournit des suites exactes longues
\SGApage{II}{115}

\begin{equation}\tag{2.4.6.1}
\begin{tikzcd}[column sep=large]
{} \& 0 \arrow[r] \arrow[d] \& R^{n-1}\Phi_{\overline{\eta}}(K)_x \arrow[r] \arrow[d] \& R^{n-1}\Phi_{\overline{\eta}}(K)_x \arrow[r] \arrow[d] \& 0 \arrow[d] \& {} \\
{} \arrow[r] \& \mathbb{H}^{n-1}(X^{*}_{s(x)},K) \arrow[r, "j"] \arrow[d, equal] \& \mathbb{H}^{n}_{\{x\}}(X_s, K) \arrow[r] \arrow[d] \& \underline{H}^{n}(K)_x \arrow[r] \arrow[d] \& \mathbb{H}^{n}(X^{*}_{s(x)}, K) \arrow[r] \arrow[d, equal] \& {} \\
{} \arrow[r] \& \mathbb{H}^{n-1}(X^{*}_{s(x)}, K) \arrow[r, "j", "\mbox{\textcircled{1}}"'] \arrow[d] \& \mathbb{H}^{n}_{\{x\}}(X_s, R\Psi_{\overline{\eta}}(K) \arrow[r] \arrow[d] \& R^{n}\Psi_{\overline{\eta}}(K)_x \arrow[r] \arrow[d] \& \mathbb{H}^{n}(X_{s(x)},K) \arrow[r] \arrow[d] \& {} \\
{} \& 0 \arrow[r] \& R^{n}\Phi_{\overline{\eta}}(K)_x \arrow[r, equal] \& R^{n}\Phi_{\overline{\eta}}(K)_x \arrow[r] \& 0 \&
{} \end{tikzcd}
\end{equation}



























}{2.4.7.   Soit $X_{s(x)}$ l'hensélisé de $X_s$ en $x$ et $X^{*}_{s(x)} = X_{s(x)} - \{x\}$ . La suite exacte longue d'hypercohomologie de $(X_{s(x)}, \{x\})$ à coefficient dans $K_s$ , $R\Psi_{\overline{\eta}}(K)$ et $R\Phi_{\overline{\eta}}(K)$ fournit des suites exactes longues
\SGApage{II}{115}

\begin{equation}\tag{2.4.7.1}
\begin{tikzcd}[column sep=large]
{} \& 0 \arrow[r] \arrow[d] \& R^{n-1}\Phi_{\overline{\eta}}(K)_x \arrow[r] \arrow[d] \& R^{n-1}\Phi_{\overline{\eta}}(K)_x \arrow[r] \arrow[d] \& 0 \arrow[d] \& {} \\
{} \arrow[r] \& \mathbb{H}^{n-1}(X^{*}_{s(x)},K) \arrow[r, "j"] \arrow[d, equal] \& \mathbb{H}^{n}_{\{x\}}(X_s, K) \arrow[r] \arrow[d] \& \underline{H}^{n}(K)_x \arrow[r] \arrow[d] \& \mathbb{H}^{n}(X^{*}_{s(x)}, K) \arrow[r] \arrow[d, equal] \& {} \\
{} \arrow[r] \& \mathbb{H}^{n-1}(X^{*}_{s(x)}, K) \arrow[r, "j", "\mbox{\textcircled{1}}"'] \arrow[d] \& \mathbb{H}^{n}_{\{x\}}(X_s, R\Psi_{\overline{\eta}}(K)) \arrow[r] \arrow[d] \& R^{n}\Psi_{\overline{\eta}}(K)_x \arrow[r] \arrow[d] \& \mathbb{H}^{n}(X^{*}_{s(x)},K) \arrow[r] \arrow[d] \& {} \\
{} \& 0 \arrow[r] \& R^{n}\Phi_{\overline{\eta}}(K)_x \arrow[r, equal] \& R^{n}\Phi_{\overline{\eta}}(K)_x \arrow[r] \& 0 \&
{} \end{tikzcd}
\end{equation}



























}{F7II-XIII-032}L'image de la flèche \textcircled{1} est contenue dans les invariants sous l'inertie.


\underline{Bibliographie}

[1]    J. GIRAUD  - Classifying topos, p. 43-56 in Toposes, algebraic geometry and logic, Lecture Notes in Mathematics 274, Springer Verlag 1971.

[2]    M. NAGATA - A generalization of the imbedding problem. J. Math. Kyoto \underline{3} 1 (1963) p. 89-102.
\selectlanguage{french}

\ifnum\SGAform=1
\clearpage\section*{Apparat critique}\small
Les lectures ci-dessous sont données dans le TeX exact des témoins, afin de préserver les distinctions de notation. Les indices décimaux renvoient aux différences individuelles du passage marqué.

\par\medskip\noindent\hypertarget{app-F7II-XIII-001.1}{}\textbf{F7II-XIII-001.1}\quad p. 83; misprint\_pass.\par
\noindent Imprimé : {\ttfamily\footnotesize autre\allowbreak{}}\par
\noindent Corrigé : {\ttfamily\footnotesize autres\allowbreak{}}\par
\noindent Source : 7II:138. \nolinkurl{fixes/FIXLOG_7II_CURRENT.jsonl#557}; \nolinkurl{fixes/MISPRINTS_7II.json#138}\par

\par\medskip\noindent\hypertarget{app-F7II-XIII-002.1}{}\textbf{F7II-XIII-002.1}\quad p. 84; misprint\_pass.\par
\noindent Imprimé : {\ttfamily\footnotesize faisceau\allowbreak{}}\par
\noindent Corrigé : {\ttfamily\footnotesize faisceaux\allowbreak{}}\par
\noindent Source : 7II:139. \nolinkurl{fixes/FIXLOG_7II_CURRENT.jsonl#558}\par

\par\medskip\noindent\hypertarget{app-F7II-XIII-003.1}{}\textbf{F7II-XIII-003.1}\quad p. 85; printed misprint.\par
\noindent Imprimé : {\ttfamily\footnotesize intertie\allowbreak{}}\par
\noindent Corrigé : {\ttfamily\footnotesize inertie\allowbreak{}}\par
\noindent Source : 7II:140. \nolinkurl{fixes/MISPRINTS_7II.json#140}\par

\par\medskip\noindent\hypertarget{app-F7II-XIII-004.1}{}\textbf{F7II-XIII-004.1}\quad p. 86; misprint\_pass.\par
\noindent Imprimé : {\ttfamily\footnotesize caractéristique\allowbreak{}}\par
\noindent Corrigé : {\ttfamily\footnotesize caractéristiques\allowbreak{}}\par
\noindent Source : 7II:141. \nolinkurl{fixes/FIXLOG_7II_CURRENT.jsonl#559}\par

\par\medskip\noindent\hypertarget{app-F7II-XIII-005.1}{}\textbf{F7II-XIII-005.1}\quad p. 87; misprint\_pass.\par
\noindent Imprimé : {\ttfamily\footnotesize automorphismesde\allowbreak{}}\par
\noindent Corrigé : {\ttfamily\footnotesize automorphismes\allowbreak{} de\allowbreak{}}\par
\noindent Source : 7II:142. \nolinkurl{fixes/FIXLOG_7II_CURRENT.jsonl#560}\par

\par\medskip\noindent\hypertarget{app-F7II-XIII-005.2}{}\textbf{F7II-XIII-005.2}\quad p. 87; misprint\_pass.\par
\noindent Imprimé : {\ttfamily\footnotesize autre\allowbreak{}}\par
\noindent Corrigé : {\ttfamily\footnotesize autres\allowbreak{}}\par
\noindent Source : 7II:143. \nolinkurl{fixes/FIXLOG_7II_CURRENT.jsonl#561}; \nolinkurl{fixes/MISPRINTS_7II.json#143}\par

\par\medskip\noindent\hypertarget{app-F7II-XIII-006.1}{}\textbf{F7II-XIII-006.1}\quad p. 89; printed misprint.\par
\noindent Imprimé : \emph{(vide)}\par
\noindent Corrigé : {\ttfamily\footnotesize \textbackslash{}overline\{\allowbreak{}}\par
\noindent Source : 7II:144. \nolinkurl{fixes/MISPRINTS_7II.json#144}\par

\par\medskip\noindent\hypertarget{app-F7II-XIII-006.2}{}\textbf{F7II-XIII-006.2}\quad p. 89; printed misprint.\par
\noindent Imprimé : \emph{(vide)}\par
\noindent Corrigé : {\ttfamily\footnotesize \}\allowbreak{}}\par
\noindent Source : 7II:144. \nolinkurl{fixes/MISPRINTS_7II.json#144}\par

\par\medskip\noindent\hypertarget{app-F7II-XIII-007.1}{}\textbf{F7II-XIII-007.1}\quad p. 89; inherited corrected-layer difference.\par
\noindent Imprimé : {\ttfamily\footnotesize Soient\allowbreak{}}\par
\noindent Corrigé : {\ttfamily\footnotesize Soit\allowbreak{}}\par
\noindent Source : ADOPT-F7II-XIII-007.1. \nolinkurl{fr/sga7ii/diplomatic/expose_XIII_body.tex}; \nolinkurl{fr/sga7ii/current/expose_XIII_body.tex}\par

\par\medskip\noindent\hypertarget{app-F7II-XIII-008.1}{}\textbf{F7II-XIII-008.1}\quad p. 95; inherited corrected-layer difference.\par
\noindent Imprimé : \emph{(vide)}\par
\noindent Corrigé : {\ttfamily\footnotesize  \textbackslash{}bar\allowbreak{}}\par
\noindent Source : ADOPT-F7II-XIII-008.1. \nolinkurl{fr/sga7ii/diplomatic/expose_XIII_body.tex}; \nolinkurl{fr/sga7ii/current/expose_XIII_body.tex}\par

\par\medskip\noindent\hypertarget{app-F7II-XIII-009.1}{}\textbf{F7II-XIII-009.1}\quad p. 97; misprint\_pass.\par
\noindent Imprimé : {\ttfamily\footnotesize faisceau\allowbreak{}}\par
\noindent Corrigé : {\ttfamily\footnotesize faisceaux\allowbreak{}}\par
\noindent Source : 7II:146. \nolinkurl{fixes/FIXLOG_7II_CURRENT.jsonl#562}\par

\par\medskip\noindent\hypertarget{app-F7II-XIII-009.2}{}\textbf{F7II-XIII-009.2}\quad p. 97; printed misprint.\par
\noindent Imprimé : {\ttfamily\footnotesize X\allowbreak{}}\par
\noindent Corrigé : {\ttfamily\footnotesize Y\allowbreak{}}\par
\noindent Source : 7II:147. \nolinkurl{fixes/MISPRINTS_7II.json#147}\par

\par\medskip\noindent\hypertarget{app-F7II-XIII-009.3}{}\textbf{F7II-XIII-009.3}\quad p. 97; misprint\_pass.\par
\noindent Imprimé : {\ttfamily\footnotesize \textbackslash{}bar\{\allowbreak{}}\par
\noindent Corrigé : \emph{(vide)}\par
\noindent Source : 7II:148. \nolinkurl{fixes/FIXLOG_7II_CURRENT.jsonl#563}; \nolinkurl{fixes/MISPRINTS_7II.json#148}\par

\par\medskip\noindent\hypertarget{app-F7II-XIII-009.4}{}\textbf{F7II-XIII-009.4}\quad p. 97; misprint\_pass.\par
\noindent Imprimé : {\ttfamily\footnotesize \}\allowbreak{}}\par
\noindent Corrigé : \emph{(vide)}\par
\noindent Source : 7II:148. \nolinkurl{fixes/FIXLOG_7II_CURRENT.jsonl#563}; \nolinkurl{fixes/MISPRINTS_7II.json#148}\par

\par\medskip\noindent\hypertarget{app-F7II-XIII-010.1}{}\textbf{F7II-XIII-010.1}\quad p. 98; misprint\_pass.\par
\noindent Imprimé : {\ttfamily\footnotesize X\allowbreak{}}\par
\noindent Corrigé : {\ttfamily\footnotesize Y\allowbreak{}}\par
\noindent Source : 7II:149. \nolinkurl{fixes/FIXLOG_7II_CURRENT.jsonl#564}; \nolinkurl{fixes/MISPRINTS_7II.json#149}\par

\par\medskip\noindent\hypertarget{app-F7II-XIII-010.2}{}\textbf{F7II-XIII-010.2}\quad p. 98; misprint\_pass.\par
\noindent Imprimé : {\ttfamily\footnotesize X\allowbreak{}}\par
\noindent Corrigé : {\ttfamily\footnotesize Y\allowbreak{}}\par
\noindent Source : 7II:150. \nolinkurl{fixes/FIXLOG_7II_CURRENT.jsonl#565}; \nolinkurl{fixes/MISPRINTS_7II.json#150}\par

\par\medskip\noindent\hypertarget{app-F7II-XIII-011.1}{}\textbf{F7II-XIII-011.1}\quad p. 99; misprint\_pass.\par
\noindent Imprimé : {\ttfamily\footnotesize et\allowbreak{}}\par
\noindent Corrigé : {\ttfamily\footnotesize est\allowbreak{}}\par
\noindent Source : 7II:151. \nolinkurl{fixes/FIXLOG_7II_CURRENT.jsonl#566}\par

\par\medskip\noindent\hypertarget{app-F7II-XIII-012.1}{}\textbf{F7II-XIII-012.1}\quad p. 100; printed misprint.\par
\noindent Imprimé : {\ttfamily\footnotesize s\allowbreak{}}\par
\noindent Corrigé : {\ttfamily\footnotesize \textbackslash{}eta\allowbreak{}}\par
\noindent Source : 7II:152. \nolinkurl{fixes/MISPRINTS_7II.json#152}\par

\par\medskip\noindent\hypertarget{app-F7II-XIII-013.1}{}\textbf{F7II-XIII-013.1}\quad p. 100; misprint\_pass.\par
\noindent Imprimé : {\ttfamily\footnotesize géométrique\allowbreak{}}\par
\noindent Corrigé : {\ttfamily\footnotesize générique\allowbreak{}}\par
\noindent Source : 7II:153. \nolinkurl{fixes/FIXLOG_7II_CURRENT.jsonl#568}; \nolinkurl{fixes/MISPRINTS_7II.json#153}\par

\par\medskip\noindent\hypertarget{app-F7II-XIII-013.2}{}\textbf{F7II-XIII-013.2}\quad p. 100; misprint\_pass.\par
\noindent Imprimé : {\ttfamily\footnotesize xr\allowbreak{}}\par
\noindent Corrigé : {\ttfamily\footnotesize nr\allowbreak{}}\par
\noindent Source : 7II:155. \nolinkurl{fixes/FIXLOG_7II_CURRENT.jsonl#569}; \nolinkurl{fixes/MISPRINTS_7II.json#155}\par

\par\medskip\noindent\hypertarget{app-F7II-XIII-014.1}{}\textbf{F7II-XIII-014.1}\quad p. 102; printed misprint.\par
\noindent Imprimé : {\ttfamily\footnotesize 8\allowbreak{}}\par
\noindent Corrigé : {\ttfamily\footnotesize 7\allowbreak{}}\par
\noindent Source : 7II:156. \nolinkurl{fixes/MISPRINTS_7II.json#156}\par

\par\medskip\noindent\hypertarget{app-F7II-XIII-015.1}{}\textbf{F7II-XIII-015.1}\quad p. 102; printed misprint.\par
\noindent Imprimé : {\ttfamily\footnotesize 1\allowbreak{}}\par
\noindent Corrigé : {\ttfamily\footnotesize 2\allowbreak{}}\par
\noindent Source : 7II:157. \nolinkurl{fixes/MISPRINTS_7II.json#157}\par

\par\medskip\noindent\hypertarget{app-F7II-XIII-016.1}{}\textbf{F7II-XIII-016.1}\quad p. 102; misprint\_in\_print.\par
\noindent Imprimé : {\ttfamily\footnotesize ambiguité\allowbreak{}}\par
\noindent Corrigé : {\ttfamily\footnotesize ambiguïté\allowbreak{}}\par
\noindent Source : FIX-7II-D-0074. \nolinkurl{fixes/FIXLOG_7II.jsonl#74}\par

\par\medskip\noindent\hypertarget{app-F7II-XIII-017.1}{}\textbf{F7II-XIII-017.1}\quad p. 103; misprint\_pass.\par
\noindent Imprimé : {\ttfamily\footnotesize ,\allowbreak{} dash\allowbreak{}}\par
\noindent Corrigé : \emph{(vide)}\par
\noindent Source : 7II:159. \nolinkurl{fixes/FIXLOG_7II_CURRENT.jsonl#571}; \nolinkurl{fixes/MISPRINTS_7II.json#159}\par

\par\medskip\noindent\hypertarget{app-F7II-XIII-017.2}{}\textbf{F7II-XIII-017.2}\quad p. 103; printed misprint.\par
\noindent Imprimé : \emph{(vide)}\par
\noindent Corrigé : {\ttfamily\footnotesize \{\textbackslash{}bar\{\allowbreak{}}\par
\noindent Source : 7II:160. \nolinkurl{fixes/MISPRINTS_7II.json#160}\par

\par\medskip\noindent\hypertarget{app-F7II-XIII-017.3}{}\textbf{F7II-XIII-017.3}\quad p. 103; printed misprint.\par
\noindent Imprimé : \emph{(vide)}\par
\noindent Corrigé : {\ttfamily\footnotesize \}\}\allowbreak{}}\par
\noindent Source : 7II:160. \nolinkurl{fixes/MISPRINTS_7II.json#160}\par

\par\medskip\noindent\hypertarget{app-F7II-XIII-017.4}{}\textbf{F7II-XIII-017.4}\quad p. 103; designated-witness attribution.\par
\noindent Imprimé : \emph{(vide)}\par
\noindent Corrigé : {\ttfamily\footnotesize )\allowbreak{}}\par
\noindent Source : WITNESS-F7II-XIII-017.4. \nolinkurl{fr/sga7ii/diplomatic/expose_XIII_body.tex}; \nolinkurl{fr/sga7ii/current/expose_XIII_body.tex}\par

\par\medskip\noindent\hypertarget{app-F7II-XIII-018.1}{}\textbf{F7II-XIII-018.1}\quad p. 104; misprint\_pass.\par
\noindent Imprimé : {\ttfamily\footnotesize Abhyankhar\allowbreak{}}\par
\noindent Corrigé : {\ttfamily\footnotesize Abhyankar\allowbreak{}}\par
\noindent Source : 7II:162. \nolinkurl{fixes/FIXLOG_7II_CURRENT.jsonl#572}\par

\par\medskip\noindent\hypertarget{app-F7II-XIII-019.1}{}\textbf{F7II-XIII-019.1}\quad p. 105; misprint\_pass.\par
\noindent Imprimé : {\ttfamily\footnotesize module\allowbreak{}}\par
\noindent Corrigé : {\ttfamily\footnotesize modules\allowbreak{}}\par
\noindent Source : 7II:164. \nolinkurl{fixes/FIXLOG_7II_CURRENT.jsonl#574}; \nolinkurl{fixes/MISPRINTS_7II.json#164}\par

\par\medskip\noindent\hypertarget{app-F7II-XIII-019.2}{}\textbf{F7II-XIII-019.2}\quad p. 105; misprint\_pass.\par
\noindent Imprimé : {\ttfamily\footnotesize Abhyankhar\allowbreak{}}\par
\noindent Corrigé : {\ttfamily\footnotesize Abhyankar\allowbreak{}}\par
\noindent Source : 7II:163. \nolinkurl{fixes/FIXLOG_7II_CURRENT.jsonl#573}\par

\par\medskip\noindent\hypertarget{app-F7II-XIII-020.1}{}\textbf{F7II-XIII-020.1}\quad p. 106; misprint\_pass.\par
\noindent Imprimé : {\ttfamily\footnotesize 11\allowbreak{}}\par
\noindent Corrigé : {\ttfamily\footnotesize 10\allowbreak{}}\par
\noindent Source : 7II:165. \nolinkurl{fixes/FIXLOG_7II_CURRENT.jsonl#575}\par

\par\medskip\noindent\hypertarget{app-F7II-XIII-020.2}{}\textbf{F7II-XIII-020.2}\quad p. 106; misprint\_pass.\par
\noindent Imprimé : {\ttfamily\footnotesize diverseurs\allowbreak{}}\par
\noindent Corrigé : {\ttfamily\footnotesize diviseurs\allowbreak{}}\par
\noindent Source : 7II:166. \nolinkurl{fixes/FIXLOG_7II_CURRENT.jsonl#576}\par

\par\medskip\noindent\hypertarget{app-F7II-XIII-021.1}{}\textbf{F7II-XIII-021.1}\quad p. 107; misprint\_pass.\par
\noindent Imprimé : {\ttfamily\footnotesize des\allowbreak{}}\par
\noindent Corrigé : {\ttfamily\footnotesize de\allowbreak{}}\par
\noindent Source : 7II:167. \nolinkurl{fixes/FIXLOG_7II_CURRENT.jsonl#577}\par

\par\medskip\noindent\hypertarget{app-F7II-XIII-022.1}{}\textbf{F7II-XIII-022.1}\quad p. 107; printed misprint.\par
\noindent Imprimé : \emph{(vide)}\par
\noindent Corrigé : {\ttfamily\footnotesize \textasciicircum{}*\allowbreak{}}\par
\noindent Source : 7II:168. \nolinkurl{fixes/MISPRINTS_7II.json#168}\par

\par\medskip\noindent\hypertarget{app-F7II-XIII-023.1}{}\textbf{F7II-XIII-023.1}\quad p. 108; inherited corrected-layer difference.\par
\noindent Imprimé : {\ttfamily\footnotesize 2\allowbreak{}}\par
\noindent Corrigé : {\ttfamily\footnotesize 1\allowbreak{}}\par
\noindent Source : ADOPT-F7II-XIII-023.1. \nolinkurl{fr/sga7ii/diplomatic/expose_XIII_body.tex}; \nolinkurl{fr/sga7ii/current/expose_XIII_body.tex}\par

\par\medskip\noindent\hypertarget{app-F7II-XIII-024.1}{}\textbf{F7II-XIII-024.1}\quad p. 108; misprint\_pass.\par
\noindent Imprimé : {\ttfamily\footnotesize déploye\allowbreak{}}\par
\noindent Corrigé : {\ttfamily\footnotesize déploie\allowbreak{}}\par
\noindent Source : 7II:169. \nolinkurl{fixes/FIXLOG_7II_CURRENT.jsonl#579}\par

\par\medskip\noindent\hypertarget{app-F7II-XIII-025.1}{}\textbf{F7II-XIII-025.1}\quad p. 109; printed misprint.\par
\noindent Imprimé : {\ttfamily\footnotesize Rf\allowbreak{}}\par
\noindent Corrigé : {\ttfamily\footnotesize Rg\allowbreak{}}\par
\noindent Source : 7II:170. \nolinkurl{fixes/MISPRINTS_7II.json#170}\par

\par\medskip\noindent\hypertarget{app-F7II-XIII-026.1}{}\textbf{F7II-XIII-026.1}\quad p. 109; printed misprint.\par
\noindent Imprimé : {\ttfamily\footnotesize x\allowbreak{}}\par
\noindent Corrigé : {\ttfamily\footnotesize y\allowbreak{}}\par
\noindent Source : 7II:171. \nolinkurl{fixes/MISPRINTS_7II.json#171}\par

\par\medskip\noindent\hypertarget{app-F7II-XIII-027.1}{}\textbf{F7II-XIII-027.1}\quad p. 110; printed misprint.\par
\noindent Imprimé : {\ttfamily\footnotesize 3\allowbreak{}}\par
\noindent Corrigé : {\ttfamily\footnotesize 2\allowbreak{}}\par
\noindent Source : 7II:172. \nolinkurl{fixes/MISPRINTS_7II.json#172}\par

\par\medskip\noindent\hypertarget{app-F7II-XIII-027.2}{}\textbf{F7II-XIII-027.2}\quad p. 110; printed misprint.\par
\noindent Imprimé : {\ttfamily\footnotesize 3\allowbreak{}}\par
\noindent Corrigé : {\ttfamily\footnotesize 2\allowbreak{}}\par
\noindent Source : 7II:173. \nolinkurl{fixes/MISPRINTS_7II.json#173}\par

\par\medskip\noindent\hypertarget{app-F7II-XIII-027.3}{}\textbf{F7II-XIII-027.3}\quad p. 110; designated-witness attribution.\par
\noindent Imprimé : {\ttfamily\footnotesize x\allowbreak{}}\par
\noindent Corrigé : {\ttfamily\footnotesize y\allowbreak{}}\par
\noindent Source : WITNESS-F7II-XIII-027.3. \nolinkurl{fr/sga7ii/diplomatic/expose_XIII_body.tex}; \nolinkurl{fr/sga7ii/current/expose_XIII_body.tex}\par

\par\medskip\noindent\hypertarget{app-F7II-XIII-028.1}{}\textbf{F7II-XIII-028.1}\quad p. 112; misprint\_pass.\par
\noindent Imprimé : {\ttfamily\footnotesize X\allowbreak{}}\par
\noindent Corrigé : {\ttfamily\footnotesize x\allowbreak{}}\par
\noindent Source : 7II:174. \nolinkurl{fixes/FIXLOG_7II_CURRENT.jsonl#580}\par

\par\medskip\noindent\hypertarget{app-F7II-XIII-029.1}{}\textbf{F7II-XIII-029.1}\quad p. 113; printed misprint.\par
\noindent Imprimé : {\ttfamily\footnotesize F\allowbreak{}}\par
\noindent Corrigé : {\ttfamily\footnotesize K\allowbreak{}}\par
\noindent Source : 7II:175. \nolinkurl{fixes/MISPRINTS_7II.json#175}\par

\par\medskip\noindent\hypertarget{app-F7II-XIII-030.1}{}\textbf{F7II-XIII-030.1}\quad p. 113; misprint\_pass.\par
\noindent Imprimé : \emph{(vide)}\par
\noindent Corrigé : {\ttfamily\footnotesize \{\allowbreak{}}\par
\noindent Source : 7II:176. \nolinkurl{fixes/FIXLOG_7II_CURRENT.jsonl#581}\par

\par\medskip\noindent\hypertarget{app-F7II-XIII-030.2}{}\textbf{F7II-XIII-030.2}\quad p. 113; misprint\_pass.\par
\noindent Imprimé : \emph{(vide)}\par
\noindent Corrigé : {\ttfamily\footnotesize *\}\allowbreak{}}\par
\noindent Source : 7II:176. \nolinkurl{fixes/FIXLOG_7II_CURRENT.jsonl#581}\par

\par\medskip\noindent\hypertarget{app-F7II-XIII-031.1}{}\textbf{F7II-XIII-031.1}\quad p. 113; printed misprint.\par
\noindent Imprimé : {\ttfamily\footnotesize 4\allowbreak{}}\par
\noindent Corrigé : {\ttfamily\footnotesize 5\allowbreak{}}\par
\noindent Source : 7II:177. \nolinkurl{fixes/MISPRINTS_7II.json#177}\par

\par\medskip\noindent\hypertarget{app-F7II-XIII-032.1}{}\textbf{F7II-XIII-032.1}\quad p. 114; misprint\_pass.\par
\noindent Imprimé : {\ttfamily\footnotesize 6\allowbreak{}}\par
\noindent Corrigé : {\ttfamily\footnotesize 7\allowbreak{}}\par
\noindent Source : 7II:178. \nolinkurl{fixes/FIXLOG_7II_CURRENT.jsonl#582}\par

\par\medskip\noindent\hypertarget{app-F7II-XIII-032.2}{}\textbf{F7II-XIII-032.2}\quad p. 115; misprint\_pass.\par
\noindent Imprimé : {\ttfamily\footnotesize 6\allowbreak{}}\par
\noindent Corrigé : {\ttfamily\footnotesize 7\allowbreak{}}\par
\noindent Source : 7II:179. \nolinkurl{fixes/FIXLOG_7II_CURRENT.jsonl#583}\par

\par\medskip\noindent\hypertarget{app-F7II-XIII-032.3}{}\textbf{F7II-XIII-032.3}\quad p. 115; misprint\_pass.\par
\noindent Imprimé : \emph{(vide)}\par
\noindent Corrigé : {\ttfamily\footnotesize )\allowbreak{}}\par
\noindent Source : 7II:180. \nolinkurl{fixes/FIXLOG_7II_CURRENT.jsonl#584}\par

\par\medskip\noindent\hypertarget{app-F7II-XIII-032.4}{}\textbf{F7II-XIII-032.4}\quad p. 115; misprint\_pass.\par
\noindent Imprimé : \emph{(vide)}\par
\noindent Corrigé : {\ttfamily\footnotesize \textasciicircum{}\{*\}\allowbreak{}}\par
\noindent Source : 7II:181. \nolinkurl{fixes/FIXLOG_7II_CURRENT.jsonl#585}\par

\normalsize\fi
\clearpage\SGAtoc{Exposé XIV — Comparaison avec la théorie transcendante}
\selectlanguage{french}

\SGApage{II}{116}


S G A  7


\underline{Exposé XIV}


\underline{\underline{Comparaison avec la théorie transcendante}}


\underline{par P. Deligne}


\underline{Sommaire}

\quad Introduction \hfill 1

\S 1. \quad Formalisme transcendant des cycles évanescents \hfill 1

\quad 1.1 Le groupe fondamental du disque épointé: conventions de signes \hfill 2

\quad 1.2 Faisceaux d'ensembles \hfill 5

\quad 1.3 Le foncteur $R\Psi$ \hfill 10

\quad 1.4 Un cas particulier \hfill 12

\S 2 \quad Le théorème de comparaison \hfill 16

\S 3 \quad Singularités isolées \hfill 20

\quad 3.1 La variété des cycles évanescents \hfill 20

\quad 3.2 La formule de Picard-Lefschetz (cas transcendant) \hfill 31

\S 4 \quad Cohomologie de De Rham \hfill 38


\underline{Introduction}

Pour comparer la théorie des cycles évanescents de l'exposé XIII avec son analogue transcendant, sur $\mathbb{C}$ , nous procèderons en deux étapes .

a) Nous donnerons un formalisme transcendant calqué sur le formalisme de l'exposé XIII, et démontrerons un théorème de comparaison par les méthodes habituelles (déjà utilisées pour prouver le théorème de finitude en XIII 2.3)

b) Dans le cas de singularités isolées, nous expliciterons en termes plus concrets ce que signifie le formalisme transcendant introduit.

Outre le théorème de comparaison, cet exposé contient une démonstration de la formule de Picard-Lefschetz (cas transcendant).


\SGApage{II}{117}

Au \S 4, nous donnons un résultat préliminaire à la théorie de Hodge des cycles évanescents. Ce paragraphe ne servira pas dans la suite du séminaire.


\SGAsectiontoc{\S 1 Formalisme transcendant des cycles évanescents}{1}
\S 1 \underline{Formalisme transcendant des cycles évanescents}


1.1. \underline{Le groupe fondamental du disque épointé: conventions de signes.}

\SGAvar{Nous aurons à considérer des disques épointés $D^*$ , et leur groupe fondamental $\mathbb{Z}$ . Ce groupe étant abélien, des ambiguités de signes risquent d'apparaître. On les élimine en oubliant qu'il l'est.

}{Nous aurons à considérer des disques épointés $D^*$ , et leur groupe fondamental $\mathbb{Z}$ . Ce groupe étant abélien, des ambiguïtés de signes risquent d'apparaître. On les élimine en oubliant qu'il l'est.

}{F7II-XIV-001}1.1.1. Pour $X$ un bon espace topologique connexe et $x \in X$ , le groupe fondamental a pour éléments les classes d'homotopie de lacets issus de $x$ ; si $\alpha, \beta \in \pi_1(X, x)$ sont représentés par des lacets $a, b$ alors $\alpha \beta$ est représenté par le lacet obtenu en faisant suivre $b$ par $a$ .

1.1.2. Soit $F$ un faisceau localement constant sur $X$ . Pour tout chemin $a : [0, 1] \to X$ , l'image réciproque $a^*F$ de $F$ sur $[0, 1]$ est un faisceau localement constant, donc constant, d'où une bijection

\begin{equation*}
a : F_{a(0)} \xleftarrow{\ \sim\ } H^0([0, 1], a^*F) \xrightarrow{\ \sim\ } F_{a(1)} \ ,
\end{equation*}
qui ne dépend que de la classe d'homotopie de $a$ . Par cette construction, $\pi_1(X, x)$ agit (à gauche) sur $F_x$ . Le foncteur $F \to F_x$ est une équivalence de la catégorie des faisceaux localement constants sur $X$ avec celle des ensembles munis d'une action de $\pi_1(X, x)$ .

\SGAvar{1.1.3. Les $\pi_1(X, x)$ pour $x$ variable sont les fibre d'un faisceau en groupes localement constant $\Pi_1(X)$ qui agit sur tout faisceau localement constant $F$ . Si $\pi_1(X, x)$ est abélien, $\Pi_1(X)$ est constant, de valeur notée $\pi_1(X)$ .

}{1.1.3. Les $\pi_1(X, x)$ pour $x$ variable sont les fibres d'un faisceau en groupes localement constant $\Pi_1(X)$ qui agit sur tout faisceau localement constant $F$ . Si $\pi_1(X, x)$ est abélien, $\Pi_1(X)$ est constant, de valeur notée $\pi_1(X)$ .

}{F7II-XIV-002}1.1.4. Le groupe $\pi_1(X, x)$ agit sur le revêtement universel $p : (\tilde{X}, \tilde{x}) \to (X, x)$ de $(X, x)$ : si on regarde ce revêtement comme un espace étalé, donc un faisceau, sur $X$ (le faisceau de ses sections locales), il correspond par 1.1.2 à $\pi_1(X, x)$ ,
\SGApage{II}{118}
muni de $e$ , sur lequel $\pi_1(X, x)$ agit par translation à gauche; le groupe $\pi_1(X, x)$ agit sur ce $\pi_1$-ensemble (par automorphismes de $\pi_1$-ensembles): $\sigma \in \pi_1(X, x)$ agit par $x \longmapsto x\,\sigma^{-1}$ .

Cette construction identifie $\pi_1(X, x)$ à $\mathrm{Aut}_X(\tilde{X})$ .

Si $F$ est localement constant sur $X$ , la bijection

\begin{equation}\tag{1.1.4.1}
H^{0}(\tilde{X}, p^{*} F) \longrightarrow (p^{*}F)_{\tilde{x}} = F_{x}
\end{equation}
commute à l'action de $\pi_1(X, x)$ , agissant sur le membre de gauche par transport de structure, via son action sur $\tilde{X}$ , et sur le membre de droite par 1.1.2. L'équivalence 1.1.2 peut se reformuler en disant que $F \longmapsto (\tilde{X} \longmapsto H^{0}(\tilde{X}, p^{*}F))$ est une équivalence de la catégorie des faisceaux localement constants sur $X$ avec celle des foncteurs

\begin{equation*}
(\text{revêtements universels de } X\ ) \longrightarrow (\mathrm{Ens})\ .
\end{equation*}
1.1.5. Si $\pi_1(X, x)$ est abélien, donc indépendant de $x$ , l'action 1.1.4 de $\pi_1(X)$ sur $\tilde{X}$ ne dépend pas du point base: $\pi_1(X)$ agit sur tout revêtement universel $\tilde{X}$ de $X$ .

On prendra garde que cette action (qui sera seule considérée) correspond à l'inverse de l'action 1.1.3.

1.1.6. Soit $D$ un disque, i.e. un espace topologique homéomorphe au disque standard $\{z \in \mathbb{C}|\ |z| < 1\}$ . Soient $O \in D$ et $D^{*} = D -\{O\}$ . Supposons $D$ muni d'une structure de variété analytique, donc orienté. Le groupe fondamental $\pi_1(D^{*})$ est alors canoniquement isomorphe à $\mathbb{Z}$ : pour le disque standard, $O = 0$ et $t \in D^{*}$ , on prend pour générateur positif de $\pi_1(D^{*}, t)$ le lacet $x \longmapsto \exp(2\pi\, i\, x)t$ .

\SGAvar{1.1.7. Soit $z : D \to \mathbb{C}$ une uniformisante $((dz)_{o} \neq 0)$ . Soit $\tilde{D}^{*}$ un revêtement universel de $D$ . On note $\log z$ une quelconque fonction sur $\tilde{D}^{*}$ telle que
\SGApage{II}{119}
$\exp(\log z) = z$ . Le couple $(\tilde{D}^{*}, \log z)$ est unique à isomorphisme unique près.

}{1.1.7. Soit $z : D \to \mathbb{C}$ une uniformisante $((dz)_{o} \neq 0)$ . Soit $\tilde{D}^{*}$ un revêtement universel de $D^{*}$ . On note $\log z$ une quelconque fonction sur $\tilde{D}^{*}$ telle que
\SGApage{II}{119}
$\exp(\log z) = z$ . Le couple $(\tilde{D}^{*}, \log z)$ est unique à isomorphisme unique près.

}{F7II-XIV-003}On pose, pour $\alpha \in \mathbb{C}$ , $z^{\alpha} = \exp(\alpha \log z)$ . Si $T$ est le générateur positif de $\pi_1(D^{*})$ , agissant sur $\tilde{D}^{*}$ par 1.1.5 et sur les fonctions sur $\tilde{D}^{*}$ par transport de structure, on a

\begin{equation}\tag{1.1.7.1}
T \log(z) = \log(z) + 2\pi\, i
\end{equation}
\begin{equation}\tag{1.1.7.2}
T\ z^{\alpha} \quad = \exp(2\pi\, i\alpha)z^{\alpha} \quad .
\end{equation}
1.1.8. Soit $D$ un disque, $O \in D$ et $D^{*} = D-\{O\}$ . Pour $D' \subset D$ un disque plus petit, avec $O \in D'$ , la restriction à ${D'}^{*}$ est une équivalence de la catégorie des faisceaux localement constants sur $D^{*}$ avec celle des faisceaux localement constants sur ${D'}^{*}$ . Pour $t \in {D'}^{*}$ , on a en effet $\pi_1({D'}^{*}, t) \overset{\sim}{\longrightarrow} \pi_1(D^{*}, t)$ .

Supposons que $D$ soit muni d'une structure de variété différentiable, soient $T$ l'espace tangent à $D$ en $O$ , et $T^{*} = T-\{0\}$ . Pour $D' \subset D$ comme plus haut, et $f : D' \hookrightarrow T$ avec $f(O) = 0$ et $(df)_{o} = \mathrm{Id}$ , le composé

\SGAvar{\begin{equation*}
(\text{restriction à } D')^{-1} \ \circ \ f^{*}
\end{equation*}
est une équivalence $[f]$ de la catégorie des faisceaux localement constant sur $T^{*}$ avec celle des faisceaux localement constant sur $D^{*}$ . A isomorphisme canonique près, elle ne dépend que du germe de $f$ en $O$ , ce germe vérifiant $(df)_{o} = \mathrm{Id}$ . Un chemin d'un tel germe à un autre définit un isomorphisme d'équivalences et deux chemins homotopes le même isomorphisme; l'ensemble des germes considéré étant contractile (rétraction sur $f_{o} : (1-\lambda)f + \lambda f_{o}$ ) , on obtient une équivalence, bien définie à isomorphisme unique près, de la catégorie des faisceaux localement constants sur $T^{*}$ avec celle des faisceaux localement constants sur $D^{*}$ .

}{\begin{equation*}
(\text{restriction à } D')^{-1} \ \circ \ f^{*}
\end{equation*}
est une équivalence $[f]$ de la catégorie des faisceaux localement constants sur $T^{*}$ avec celle des faisceaux localement constants sur $D^{*}$ . A isomorphisme canonique près, elle ne dépend que du germe de $f$ en $O$ , ce germe vérifiant $(df)_{o} = \mathrm{Id}$ . Un chemin d'un tel germe à un autre définit un isomorphisme d'équivalences et deux chemins homotopes le même isomorphisme; l'ensemble des germes considéré étant contractile (rétraction sur $f_{o} : (1-\lambda)f + \lambda f_{o}$ ) , on obtient une équivalence, bien définie à isomorphisme unique près, de la catégorie des faisceaux localement constants sur $T^{*}$ avec celle des faisceaux localement constants sur $D^{*}$ .

}{F7II-XIV-004}1.1.9. Pour $D$ un disque, $O \in D$ , $D^{*} = D-\{O\}$ et $p : \tilde{D}^{*} \to D^{*}$ un revêtement universel de $D^{*}$ , nous noterons $p : \tilde{D} \to D$ l'espace topologique sur $D$ déduit de $\tilde{D}^{*}$ par adjonction d'un point $O$ , de sorte que
\SGApage{II}{120}
a) $\widetilde{D}^*$ est ouvert dans $\widetilde{D}$ ;

b) $p : \widetilde{D} \to D$ prolonge $p : \widetilde{D}^* \to D^*$ et $p(O) = O$ ;

\SGAvar{c) les $p^{-1}(U)$ , pour $U$ voisinage de $O$ dans $D$ , formant un système fondamental de voisinages de $O$ dans $\widetilde{D}$ .

}{c) les $p^{-1}(U)$ , pour $U$ voisinage de $O$ dans $D$ , forment un système fondamental de voisinages de $O$ dans $\widetilde{D}$ .

}{F7II-XIV-005}1.1.10 Soient $X$ une courbe algébrique complexe non singulière, $O \in X$ et $D$ un disque voisinage de $O$ dans $X$ . L'anneau des germes en $O$ des fonctions holomorphes sur $D$ , algébriques sur le corps $K(X)$ des fonctions rationnelles sur $X$ , est l'hensélisé $\mathcal{O}_o^h$ de l'anneau local $\mathcal{O}_o$ de $X$ en $O$ . Soit $\widetilde{D}^*$ un revêtement universel de $D^* = D-\{O\}$ . Le corps des germes en $O \in \widetilde{D}$ de fonctions méromorphes sur $\widetilde{D}^*$ , algébriques sur $K(X)$ , est une clôture algébrique $\overline{K}$ de $K(X)$ (et donc du corps des fractions $K_o^h$ de $\mathcal{O}_o^h$ ).

Le groupe $\pi_1(D^*)$ des $D$-automorphismes de $\widetilde{D}^*$ (1.1.5) agit par transport de structure sur $\overline{K}$ , d'où

\begin{equation}\tag{1.1.10.1}
\begin{tikzcd}
\pi_1(D^*) \arrow[r] \arrow[d, equal] \& \mathrm{Gal}(\overline{K}/K_o^h) \arrow[d, equal] \& {} \\
\mathbb{Z} \& \widehat{\mathbb{Z}}(1) \& .
\end{tikzcd}
\end{equation}












Il résulte de (1.1.7.2) que ce diagramme commute avec l'inclusion naturelle (sur $\mathbb{C}$ ) de $\mathbb{Z}$ dans $\widehat{\mathbb{Z}}(1)$ , composée de l'inclusion de $\mathbb{Z} \lhook\joinrel\longrightarrow \widehat{\mathbb{Z}}$ et de l'isomorphisme $\widehat{\mathbb{Z}} \overset{\sim}{\longrightarrow} \widehat{\mathbb{Z}}(1)$ , limite projective des isomorphismes $\exp(\frac{2\pi i}{n} k )$ $: \mathbb{Z}/n \to \mu_n$ .

1.1.11. Chaque faisceau $F$ sur $\mathrm{Spec}(K_o^h)$ définit un faisceau localement constant sur $D^*$ : on utilise l'équivalence 1.1.4 et, à chaque revêtement universel $\widetilde{D}^*$ de $D^*$ , on associe $F_{\overline{K}}$ pour le $\overline{K}$ construit en 1.1.10.


1.2. \underline{Faisceaux d'ensembles}.

1.2.1 Soient $D$ un disque (1.1.6), $O \in D$ et $D^* = D-\{O\}$ . Nous noterons $\mathcal{D}^*$
\SGApage{II}{121}
le topos des faisceaux d'ensembles localement constants sur $D^*$ . D'après 1.1.8, ce topos ne ``change'' pas si on remplace $D$ par un disque plus petit $D'$ , contenant $O$ , ni, pour $D$ muni d'une structure de variété différentiable, si on remplace $D$ par son espace tangent en $O$ .

Pour $t \in D^*$ , le foncteur $F \longmapsto F_t$ identifie ce topos au topos des ensembles munis d'une action de $\pi_1(D^*)$ (1.1.2). D'après (1.1.4), le foncteur $F \longmapsto (\widetilde{D}^* \to H^o(\widetilde{D}^*, p^* F))$ l'identifie au topos des foncteurs

\begin{equation*}
(\text{revêtements universels de } D^*) \longrightarrow (\text{Ens}) \;,
\end{equation*}
et ces identifications sont reliées par 1.1.4.1.

1.2.2. Soit $F$ un faisceau sur $D$ , dont la restriction à $D^*$ soit localement constante. Il existe un et un seul morphisme de faisceaux $\underline{F_o} \xrightarrow{\ \alpha\ } F$ qui soit l'identité en $O$ . Comme expliqué en XIII 0.1 (pour $X_o = (P^t)$ ), le foncteur $F \longmapsto (F_o,$ $F$ restreint à $D^*$ , $\alpha)$ identifie ces faisceaux aux triples formés d'un ensemble $F_o$ , d'un faisceau localement constant $F_\eta$ sur $D^*$ et de $\alpha$ : (faisceau constant de valeur $F_o$ sur $D^*$ ) $\longrightarrow$ (faisceau $F_\eta$ ).

1.2.3 Nous noterons $\mathcal{D}$ le topos des faisceaux d'ensembles sur $D$ dont la restriction à $D^*$ est localement constante. L'inclusion de $\mathcal{D}$ dans les faisceaux sur $D$ est le morphisme image inverse d'un morphisme de topos $D \to \mathcal{D}$ qui s'insère dans un diagramme commutatif

\begin{center}
\begin{tikzcd}
\{O\} \arrow[r, hook] \arrow[d, equal] \& D \arrow[d, "\beta"] \& D^* \arrow[l, hook'] \arrow[d, "\beta"] \& {} \\
\{O\} \arrow[r, hook] \& \mathcal{D} \& \mathcal{D}^* \arrow[l, hook'] \& .
\end{tikzcd}
\end{center}



















\SGAvar{Comme en 1.1.1, on vérifie que ce topos ne change pas quand on remplace $D$ par un disque plus petit ou par l'espace tangent à $D$ en $O$ .


}{Comme en 1.2.1, on vérifie que ce topos ne change pas quand on remplace $D$ par un disque plus petit ou par l'espace tangent à $D$ en $O$ .


}{F7II-XIV-006}\SGApage{II}{122}

1.2.4. Soit $Y$ un espace topologique. Le topos produit $Y \times \mathcal{D}$ admet les trois descriptions suivantes (pour leur équivalence, voir XIII 0.1) .

a) $Y \times \mathcal{D}$ est le topos des faisceaux $F$ sur $Y \times D$ vérifiant

(*) $D^*$ est recouvert par des ouverts $U$ tels que $F | Y \times U$ soit image réciproque d'un faisceau sur $Y$ .

b) Soit $t \in D^*$ . Le topos $Y \times \mathcal{D}$ est le topos des triples $(F_o, F_t, \alpha)$ où

$\alpha$) $F_o$ est un faisceau sur $Y$ ;

$\beta$) $F_t$ est un faisceau sur $Y$ , muni d'une action de $\pi_1(D^*)$ ;

$\gamma$) $\alpha$ est un morphisme de $F_o$ dans $F_t$ , d'image contenue dans les invariants.

c) Le topos $Y \times \mathcal{D}$ est le topos des triples $(F_o, F_\eta, \alpha)$ où

$\alpha'$) $F_o$ est un faisceau sur $Y$ ;

$\beta'$) $F_\eta$ est un foncteur, de valeurs notées $F_{\bar\eta}$ ,

\SGAvar{\begin{equation*}
(\text{revêtements universel de } D^*) \longrightarrow (\text{faisceaux sur } Y ) ;
\end{equation*}
$\gamma'$) $\alpha$ est un morphisme de $F_o$ (le foncteur constant de valeur $F_o$ des revêtements universels de $D^*$ dans les faisceaux sur $Y$ ) dans $F_\eta$ .

}{\begin{equation*}
(\text{revêtements universels de } D^*) \longrightarrow (\text{faisceaux sur } Y ) ;
\end{equation*}
$\gamma'$) $\alpha$ est un morphisme de $F_o$ (le foncteur constant de valeur $F_o$ des revêtements universels de $D^*$ dans les faisceaux sur $Y$ ) dans $F_\eta$ .

}{F7II-XIV-007}Peu nous chaut, d'ailleurs, que ce topos soit vraiment le topos produit de $Y$ et $\mathcal{D}$ . On définit de même $Y \times \mathcal{D}^*$ comme étant au choix

a*) le topos des faisceaux sur $Y \times D^*$ vérifiant (*) ;

b*) (pour $t \in D^*$) le topos des faisceaux $F_t$ sur $Y$ , munis d'une action de $\pi_1(D^*)$ ;

c*) le topos des foncteurs

\begin{equation*}
(\text{revêtements universels de } D^*) \longrightarrow (\text{faisceaux sur } Y ).
\end{equation*}
1.2.5 Le topos $Y \times \mathcal{D}$ a les mêmes propriétés de fonctorialité en $Y$ que les topos $Y \times_s S$ de l'exposé XIII, et nous nous dispenserons d'en faire la liste. Citons simplement, pour fixer les notations, le diagramme 


\SGApage{II}{123}

\begin{center}
\begin{tikzcd}[row sep=huge, column sep=huge]
Y \arrow[d, ""{name=A}] \& Y \arrow[d, ""{name=B}] \arrow[r, hook, "i"] \& Y \times \mathcal{D} \arrow[d] \& Y \times \mathcal{D}^* \arrow[d] \arrow[l, hook', "j"'] \\
(P^t) \& \{0\} \arrow[r, hook, "i"] \& \mathcal{D} \& \mathcal{D}^* \arrow[l, hook', "j"'] \arrow[from=B, to=A, "\mathrm{sp}"']
\end{tikzcd}
\end{center}



















1.2.6. Soit $K$ un complexe de faisceaux abéliens sur $Y \times \mathcal{D}$ . On procède comme en XIII 1.4.2 pour déduire de $K$ un triangle distingué dans $K(Y \times \mathcal{D}^*)$

\begin{center}
\begin{tikzcd}[row sep=large, column sep=large]
{} \& \Phi(K) \arrow[dl] \& {} \\
\mathrm{sp}^* K_o \arrow[rr] \& {} \& K_\eta \arrow[ul]
\end{tikzcd}
\end{center}






Cette construction passe à la catégorie dérivée et définit un foncteur

\begin{equation*}
D(Y \times \mathcal{D}) \longrightarrow (\text{triangles distingués de } D(Y \times \mathcal{D}^*)) .
\end{equation*}
Supposons $D$ orienté (par exemple analytique complexe) et soit $T$ le générateur positif de $\pi_1(D^*)$ (1.1.6). Soit $K \in \mathrm{Ob}\ D(Y \times \mathcal{D})$ . On définit comme en XIII 1.4.3 la \underline{variation} $V = \mathrm{Var}(T)$

\begin{equation*}
V : \Phi(K) \longrightarrow K_\eta .
\end{equation*}
L'endomorphisme T-Id de $K_\eta$ est le composé

\begin{equation*}
\text{T-Id} : K_\eta \longrightarrow \Phi(K) \longrightarrow K_\eta .
\end{equation*}
1.2.7. Soient $D$ un disque, $0 \in D$ , $D^* = D-\{0\}$ , $f : X \to D$ une application continue et $X_o = f^{-1}(0)$ . Nous nous proposons de définir un foncteur

\begin{equation*}
\Psi : (\text{faisceaux sur } X ) \longrightarrow (\text{faisceaux sur } X_o \times \mathcal{D}) .
\end{equation*}
Soient $\widetilde{D}^*$ un revêtement universel de $D^*$ , $\widetilde{D}$ comme en 1.1.9, 


\SGApage{II}{124}

\begin{equation*}
 \overline{X} = X \times_D \widetilde{D} \; , \quad X^* = X-X_0 = X \times_D D^* \quad \mbox{et} \quad \overline{X}^* = \overline{X} - X_0 = X \times_D \widetilde{D}^* \; .
\end{equation*}
On a un diagramme commutatif

\SGAvar{\begin{equation}\tag{1.2.7.1}
\begin{tikzcd}
X_0 \arrow[r, hook, "\overline{i}"] \arrow[d, equal] \& \overline{X} \arrow[d, "p"] \& X^* \arrow[l, hook, "\overline{j}"'] \arrow[d, "p'"] \\
X_0 \arrow[r, hook, "i"] \& X \& X^* \arrow[l, hook, "j"']
\end{tikzcd}
\end{equation}















}{\begin{equation}\tag{1.2.7.1}
\begin{tikzcd}
X_0 \arrow[r, hook, "\overline{i}"] \arrow[d, equal] \& \overline{X} \arrow[d, "p"] \& \overline{X}^* \arrow[l, hook, "\overline{j}"'] \arrow[d, "p'"] \\
X_0 \arrow[r, hook, "i"] \& X \& X^* \arrow[l, hook, "j"']
\end{tikzcd}
\end{equation}















}{F7II-XIV-008}cartésien au-dessus du diagramme

\begin{equation}\tag{1.2.7.2}
\begin{tikzcd}
\{0\} \arrow[r, hook, "\overline{i}"] \arrow[d, equal] \& \widetilde{D} \arrow[d, "p"] \& \widetilde{D}^* \arrow[l, hook, "\overline{j}"'] \arrow[d, "p'"] \\
\{0\} \arrow[r, hook, "i"] \& D \& D^* \arrow[l, hook, "j"']
\end{tikzcd}
\end{equation}















Pour $F$ un faisceau sur $X$ , on pose (avec les notations de 1.2.4 c))

\begin{equation*}
\Psi(F) = (i^*F, \; \overline{i}^*\overline{j}_*((jp')^* F), \; \alpha)
\end{equation*}
où $\alpha : i^*F = \overline{i}^* p^* F \longrightarrow \overline{i}^* \overline{j}_* ((jp')^* F) = \overline{i}^* \overline{j}_* \overline{j}^* p^* F$ est déduit, par application de $\overline{i}^*$ , de la flèche d'adjonction $\mathrm{Id} \to \overline{j}_* \overline{j}^*$ .

On définit aussi un foncteur

\begin{equation*}
\Psi_\eta : (\mbox{faisceaux sur } X^* ) \longrightarrow (\mbox{faisceaux sur } X_0 \times \mathcal{D}^*)
\end{equation*}
par $\Psi_\eta(F) = \overline{i}^* \overline{j}_*({p'}^* F)$ .

1.2.8. Les propriétés de fonctorialité de $\Psi$ sont les mêmes que celles de son analogue étudié en XIII 1.3.5 à 1.3.9.


\SGApage{II}{125}


1.3. \underline{Le foncteur} $R\Psi$ .

1.3.1. Soient $D$ un disque, $0 \in D$ , $D^* = D-\{0\}$, $f: X \to D$ une application continue et $X_0 = f^{-1}(0)$ . Les foncteurs $R\Psi$ et $R\Psi_\eta$ sont les foncteurs dérivés des foncteurs $\Psi$ et $\Psi_\eta$

\begin{equation*}
R\Psi : D^+(X) \longrightarrow D^+(X_0 \times \mathcal{D})
\end{equation*}
\begin{equation*}
R\Psi_\eta : D^+(X^*) \longrightarrow D^+(X_0 \times \mathcal{D}^*) \; .
\end{equation*}
On note $R\Phi$ le foncteur composé $\Phi \circ R\Psi$ .

\begin{equation*}
R\Phi : D^+(X) \longrightarrow D^+(X_0 \times \mathcal{D}^*) \; .
\end{equation*}
Les propriétés de ces foncteurs sont toutes pareilles à celles des foncteurs de même nom de XIII 2.1. Leur démonstration est bien plus élémentaire car, dans le cadre topologique, les théorèmes de changement de base pour un morphisme propre ou par un morphisme lisse (= localement isomorphe à $\mathbb{R}^n \times S \to S$) sont faciles à prouver. Voici les menus changements à apporter.

1.3.2. Le foncteur $Rf_!$ n'est défini que pour un morphisme séparé entre espaces localement compacts, et est le foncteur dérivé du foncteur image directe à supports propres $f_!$ .

1.3.3. L'isomorphisme fondamental (2.1.8.3), pour un morphisme propre (et séparé) $f : X \to D$ est à récrire comme suit. Soit $p : \widetilde{D}^* \to D^*$ un revêtement universel de $D^*$ . Pour tout voisinage $U$ de $0$ dans $D$ , soit $\widetilde{U}^* = p^{-1}(U)$ , et $X_{\widetilde{U}^*} = X \times_D \widetilde{U}^*$ . Pour $K \in \mathrm{Ob}\; D^+(f^{-1}(D^*))$ , on a

\begin{equation}\tag{1.3.3.1}
\varinjlim_U \; \mathbb{H}^i(X_{\widetilde{U}^*} , K) \xrightarrow{\;\sim\;} \mathbb{H}^i(X_0, R\Psi_{\overline{\eta}} (K)) \; .
\end{equation}
Cet isomorphisme est intéressant surtout si $X^* = f^{-1}(D^*)$ est un fibré topologique sur $D^*$ . Pour $U \subset D$ un disque tendant vers $0$ , la suite spectrale
\SGApage{II}{126}
de Leray de $X_{\widetilde{U}^*} \to \widetilde{U}^*$ montre alors que pour tout faisceau localement constant $\Lambda$ sur $X^*$ et pour tout $t \in \widetilde{U}^*$ , on a $H^i(X_{\widetilde{U}}, \Lambda) \xrightarrow{\sim} H^i(X_t, \Lambda)$ . Pour $t \in D^*$ , on a donc

\begin{equation}\tag{1.3.3.2}
H^i(X_t, \Lambda) \xrightarrow{\sim} H^i(X_o, R\Psi_t(\Lambda))
\end{equation}
($X/D$ propre, $X^*$ fibré, $\Lambda$ localement constant).

Pour vérifier que $X^*$ est un fibré sur $D^*$ , on utilise en pratique le théorème d'isotopie de Thom dont une forme est énoncée ci-dessous. (Pour la démonstration de ce théorème, dont aucune conséquence ne servira dans les exposés suivants, je renvoie à Whitney et Mather).

1.3.4. \underline{Préliminaire}: Soit $f : X \to Y$ un morphisme lisse ($=_{\mathrm{dfn}}$ sans point critique ) de variétés différentiables. Une \underline{connexion} $\nabla$ sur $X/Y$ est un sous-fibré vectoriel du fibré tangent de $X$ , supplémentaire de $\mathrm{Ker}(df)$ . Pour $u$ un champ de vecteurs sur $Y$ , il existe un et un seul champ de vecteurs $\nabla(u)$ sur $X$ , relevant $u$ , dans ce sous-fibré. On dit que $\nabla$ est intégrable si $\nabla([u, v]) = [\nabla(u), \nabla(v)]$ ; pour $f$ propre, une connexion intégrable définit des trivialisations locales sur $Y$ de $X/Y$ : $\alpha : f^{-1}(U) \xrightarrow{\sim} X_o \times U$ , compatible à la projection sur $U$ , avec $d\alpha (\nabla(u)) = (0, u)$ . Pour $f$ non propre, le même résultat vaut sous des conditions convenables à l'infini, par exemple si les $\nabla(u)g$ sont nuls pour $g$ une fonction d'exhaustion convenable.


1.3.5. \underline{Théorème d'isotopie de Thom.}

\underline{Soient} $D = \{z \in \mathbb{C} \mid |z| < 1\}$ , $D^* = D-\{0\}$ \underline{et} $f : X \to D$ \underline{un morphisme propre d'espaces analytiques. Pour} $D$ \underline{remplacé par un disque plus petit et} $X$ \underline{par son image réciproque, il existe une stratification} $\Sigma$ \underline{de} $X^* = f^{-1}(D^*)$ \underline{ayant les propriétés suivantes.}

(a) $\Sigma$ \underline{est une filtration} $X^* = F_o \supset F_1 \supset \ldots \supset F_n \supset F_{n+1} = \emptyset$ \underline{de} $X^*$ \underline{par des parties fermées dont l'adhérence} $\overline{F}_i$ \underline{dans} $X$ \underline{est un sous-espace analytique de} $X$ .


\SGApage{II}{127}

(b) \underline{Les strates ouvertes} $F_i - F_{i+1}$ \underline{sont lisses sur} $D^*$ .

(c) $\Sigma$ \underline{vérifie les conditions} (a) \underline{et} (b) \underline{de Whitney}.

\underline{Pour une telle stratification} $\Sigma$ , \underline{il existe sur les} $F_i - F_{i+1}$ \underline{des connexions intégrables} $\nabla_i$ \underline{telles que}

(i) \underline{localement sur} $D^*$ , $\nabla_i$ \underline{définit une trivialisation de} $F_i - F_{i+1}$ ;

(ii) \underline{ces trivialisations se recollent pour fournir des trivialisations locales de} $X^*/D^*$ .

\SGAvar{\underline{Enfin, si} $(Y_\alpha)$ \underline{est une famille finie de sous-espaces analytiques fermés de} $X$ , \underline{et quitte à retrécir} $D$ \underline{davantage, on peut trouver} $\Sigma$ \underline{tel que chaque} $Y_\alpha$ \underline{soit réunion de composantes connexes de strates ouvertes} $(F_i - F_{i+1})$ . \underline{Une connexion} $\nabla$ \underline{comme plus haut trivialise alors localement} $(X, (Y_\alpha))$ .

}{\underline{Enfin, si} $(Y_\alpha)$ \underline{est une famille finie de sous-espaces analytiques fermés de} $X$ , \underline{et quitte à rétrécir} $D$ \underline{davantage, on peut trouver} $\Sigma$ \underline{tel que chaque} $Y_\alpha$ \underline{soit réunion de composantes connexes de strates ouvertes} $(F_i - F_{i+1})$ . \underline{Une connexion} $\nabla$ \underline{comme plus haut trivialise alors localement} $(X, (Y_\alpha))$ .

}{F7II-XIV-009}1.3.6. On dispose encore de (1.3.3.2) pour l'hypercohomologie à valeurs dans un complexe de faisceaux $K$ lorsque $X$ est propre sur $D$ et que la condition suivante est vérifiée

(*) Pour tout $i \in \mathbb{Z}$ , $D^*$ est recouvert par des ouverts $U$ tels qu'il existe un isomorphisme compatible à la projection sur $U$ $\varphi : f^{-1}(U) \xrightarrow{\sim} Y \times U$ et un faisceau $G$ sur $Y$ tel que $\underline{H}^i(K)$ soit isomorphe, sur $f^{-1}(U)$ , à $\varphi^* \mathrm{pr}_1^* G$ .

\SGAvar{Sous cette hypothèse, le triangle 1.2.1 fournit une suite exacte longue

\begin{equation}\tag{1.3.6.1}
\ldots \longrightarrow \mathbb{H}^i(X_o, K) \longrightarrow \mathbb{H}^i(X_t, K) \longrightarrow \mathbb{H}^i(X_o, R\Phi_t(K)) \longrightarrow \ldots
\end{equation}
D'après le théorème de Thom, la condition (*) est vérifiée (quitte à retrécir $D$ ) lorsqu'il existe une stratification analytique complexe de $X$ telle que les $\underline{H}^i(K)$ soient localement constants sur chaque strate ouverte.


}{Sous cette hypothèse, le triangle 1.2.6 fournit une suite exacte longue

\begin{equation}\tag{1.3.6.1}
\ldots \longrightarrow \mathbb{H}^i(X_o, K) \longrightarrow \mathbb{H}^i(X_t, K) \longrightarrow \mathbb{H}^i(X_o, R\Phi_t(K)) \longrightarrow \ldots
\end{equation}
D'après le théorème de Thom, la condition (*) est vérifiée (quitte à rétrécir $D$ ) lorsqu'il existe une stratification analytique complexe de $X$ telle que les $\underline{H}^i(K)$ soient localement constants sur chaque strate ouverte.


}{F7II-XIV-010}1.4. \underline{Un cas particulier.}

Ce numéro ne sera pas utilisé dans la suite du séminaire. On y formalise une
\SGApage{II}{128}
construction esquissée dans l'introduction de XIII.

1.4.1. Soient $D$ un disque orienté, $O \in D$ et $D^* = D-\{O\}$ . Un \underline{espace topologique} $\mathfrak{X}^*$ sur $\mathfrak{D}^*$ est un foncteur

\SGAvar{\begin{equation*}
(\mbox{revêtements universel de } D^*\ ) \longrightarrow (\mbox{espaces topologiques}) \quad .
\end{equation*}
Une fois choisi un revêtement universel $\widetilde{D}^*$ , c'est un espace topologique $X_\eta$ muni d'une action de $\pi_1(D^*)$ ; pour $\widetilde{D}^*$ le revêtement universel de $(D^*, t)$ , on le note $\mathfrak{X}^*_t$ . \underline{Un faisceau} sur $\mathfrak{X}^*$ est un foncteur

}{\begin{equation*}
(\mbox{revêtements universels de } D^*\ ) \longrightarrow (\mbox{espaces topologiques}) \quad .
\end{equation*}
Une fois choisi un revêtement universel $\widetilde{D}^*$ , c'est un espace topologique $X_\eta$ muni d'une action de $\pi_1(D^*)$ ; pour $\widetilde{D}^*$ le revêtement universel de $(D^*, t)$ , on le note $\mathfrak{X}^*_t$ . \underline{Un faisceau} sur $\mathfrak{X}^*$ est un foncteur

}{F7II-XIV-011}\begin{equation*}
(\mbox{revêtements universels de } D^*\ ) \longrightarrow (\mbox{espaces topologiques munis d'un faisceau}),
\end{equation*}
soit encore un faisceau $\pi_1(D^*)$-équivariant sur $X_\eta$ .

A $\mathfrak{X}^*$ , on associe l'espace topologique $\beta^* \mathfrak{X}^* = X_\eta \times \widetilde{D}^*/\pi_1(D^*)$ sur $D^*$ . Tout faisceau $F$ sur $\mathfrak{X}^*$ définit par passage au quotient un faisceau $\beta^*F$ sur $\beta^* \mathfrak{X}^*$ , d'où un diagramme commutatif de topos

\begin{center}
\begin{tikzcd}
\beta^* \mathfrak{X}^* \arrow[r] \arrow[d] \& \mathfrak{X}^* \arrow[d] \\
D^* \arrow[r, "\beta"] \& \mathfrak{D}^*
\end{tikzcd}\quad .
\end{center}








1.4.2. Un \underline{espace topologique} $\mathfrak{X}$ sur $\mathfrak{D}$ consiste en

(a) un espace topologique $\mathfrak{X}^*$ sur $\mathfrak{D}^*$ ;

(b) un espace topologique $X_o$ ;

(c) un morphisme $\pi_1(D^*)$-invariant $\mathrm{sp} : X_\eta \to X_o$ ($\pi_1(D^*)$ agissant trivialement sur $X_o$) .

Un \underline{faisceau} sur $\mathfrak{X}$ consiste en un faisceau $F_o$ sur $X_o$ , un faisceau sur $\mathfrak{X}^*$ , donné par $F_\eta$ , et un morphisme équivariant

$\mathrm{sp} : \mathrm{sp}^* F_o \longrightarrow F_\eta$ (ou $F_o \longrightarrow \mathrm{sp}_* F_\eta$) . On dispose d'un diagramme commutatif 


\SGApage{II}{129}

\SGAvar{\begin{center}
\begin{tikzcd}
X_o \arrow[r, hook] \arrow[d] \& \mathfrak{X} \arrow[d, dash] \& \mathfrak{X}^* \arrow[l, hook] \arrow[d] \\
\{O\} \arrow[r, hook] \& \mathfrak{D} \& \mathfrak{D}^* \arrow[l, hook]
\end{tikzcd}\quad .
\end{center}












}{\begin{center}
\begin{tikzcd}
X_o \arrow[r, hook] \arrow[d] \& \mathfrak{X} \arrow[d] \& \mathfrak{X}^* \arrow[l, hook] \arrow[d] \\
\{O\} \arrow[r, hook] \& \mathfrak{D} \& \mathfrak{D}^* \arrow[l, hook]
\end{tikzcd}\quad .
\end{center}












}{F7II-XIV-012}On dispose de plus d'un morphisme de topos $\Psi : \mathfrak{X} \to X_o \times \mathfrak{D}$ , d'après la fonctorialité de la construction précédente appliquée à $(X_\eta , X_o, \mathrm{sp}) \longrightarrow (X_o, X_o, \mathrm{Id})$ . On a

\begin{equation*}
\Psi(F_o, F_\eta , \mathrm{sp}) = (F_o, \mathrm{sp}_* F_\eta , \mathrm{sp}) \quad .
\end{equation*}
Ce morphisme induit $\Psi_\eta : \mathfrak{X}^* \longrightarrow X_o \times \mathfrak{D}^*$ .

\SGAvar{1.4.3. A un espace topologique $\mathfrak{X}$ sur $\mathfrak{D}$ , on associe comme suit un espace topologique $X = \beta^* \mathfrak{X}$ sur $D$ . On adjoint à $X^* = \beta^* \mathfrak{X}^*$ la fibre $X_o$ en $O$ , en prenant pour système fondamental de voisinage de $x_o \in X_o$ les $V \cup (\mathrm{sp}^{-1}(V) \times \widetilde{U}^*/\pi_1(D^*))$ pour $V$ un voisinage de $x_o$ dans $X_o$ , $U$ un voisinage de $O$ dans $D$ et $\widetilde{U}^* = U \times_D \widetilde{D}^*$ . Notons $j$ (resp. $i$) l'inclusion de $X^*$ (resp. $X_o$) dans $X$ . Pour $F$ un faisceau sur $\mathfrak{X}^*$ , on a

}{1.4.3. A un espace topologique $\mathfrak{X}$ sur $\mathfrak{D}$ , on associe comme suit un espace topologique $X = \beta^* \mathfrak{X}$ sur $D$ . On adjoint à $X^* = \beta^* \mathfrak{X}^*$ la fibre $X_o$ en $O$ , en prenant pour système fondamental de voisinages de $x_o \in X_o$ les $V \cup (\mathrm{sp}^{-1}(V) \times \widetilde{U}^*/\pi_1(D^*))$ pour $V$ un voisinage de $x_o$ dans $X_o$ , $U$ un voisinage de $O$ dans $D$ et $\widetilde{U}^* = U \times_D \widetilde{D}^*$ . Notons $j$ (resp. $i$) l'inclusion de $X^*$ (resp. $X_o$) dans $X$ . Pour $F$ un faisceau sur $\mathfrak{X}^*$ , on a

}{F7II-XIV-013}\begin{equation*}
i^* j_* \beta^* F = (\mathrm{sp}_* F_\eta)^{\pi_1(D^*)} \quad .
\end{equation*}
On en tire un diagramme commutatif de topos

\begin{center}
\begin{tikzcd}
X \arrow[r, "\beta"] \arrow[d] \& \mathfrak{X} \arrow[d] \\
D \arrow[r, "\beta"] \& \mathfrak{D}
\end{tikzcd}
\end{center}






tel que, pour tout faisceau $F$ sur $\mathfrak{X}$ , on ait

\begin{equation}\tag{1.4.3.1}
\Psi(F) \stackrel{\sim}{\longrightarrow} \Psi(\beta^* F) \quad .
\end{equation}
\underline{Proposition} 1.4.4. \underline{La flèche} (1.4.3.1)\underline{se dérive en un} isomorphisme

\begin{equation*}
R\Psi \stackrel{\sim}{\longrightarrow} R\Psi \circ \beta^*
\end{equation*}

\SGApage{II}{130}

 \underline{entre foncteurs de} $D^+(\mathfrak{X})$ \underline{dans} $D^+(X_o \times \mathfrak{D})$ .

Soit $\bar{j}$ l'inclusion de $\bar{X}^* = X \times_D \widetilde{D}^* \simeq X_\eta \times \widetilde{D}^*$ dans $\bar{X} = X \times_D \bar{D}$ . La proposition se ramène à vérifier que pour tout faisceau abélien $F$ flasque sur $X_\eta$ , $\mathrm{pr}_1^* F$ est acyclique pour le foncteur $\bar{j}_*$ . Cette assertion, à son tour, se ramène à la suivante (pour $B = \widetilde{D}^*$ ), bien connue, qui se démontre par un argument d'homotopie.

Soit $F$ un faisceau abélien sur un espace $Y$ , et $B$ un disque. On a

\begin{equation*}
H^i(Y, F) \overset{\sim}{\longrightarrow} H^i(Y \times B, \mathrm{pr}_1^* F) \quad .
\end{equation*}
1.4.5. La proposition 1.4.4 fournit un procédé de calcul de $R\Psi$ pour les faisceaux sur $X$ de la forme $\beta^* F$ , ramenant pour l'essentiel le calcul de $R\Psi$ à celui de $R\, \mathrm{sp}_*$ .

1.4.6. Nous aurons plus tard à considérer le cas où $X_\eta$ est localement compact et où $X_o$ se déduit de $X_\eta$ en contractant en un point $O$ une partie compacte $V \subset X_\eta$ , $\mathrm{sp}$ étant la flèche de contraction et l'action de $\pi_1(D^*)$ étant triviale en dehors de $V$ . Le théorème de changement de base pour le morphisme propre $\mathrm{sp}$ fournit alors, pour $K \in \mathrm{Ob}\ D^+(\mathfrak{X}^*)$

\begin{equation}\tag{1.4.6.1}
R\Psi_\eta(\beta^* K)_o \;\; = R\Gamma (V, K)
\end{equation}
\begin{equation}\tag{1.4.6.2}
R\Gamma_{\{O\}} \, R\Psi_\eta(\beta^* K) \;\; = R\Gamma_V(X_\eta , K)
\end{equation}
(cohomologie à support) .

Pour $K$ sur $\mathfrak{X}$ , avec $\mathrm{sp} : \mathrm{sp}^* K_o \longrightarrow K_\eta$ un quasi-isomorphisme en dehors de $V$ , on a $R\Phi(K) = 0$ en dehors de $O$ . En particulier, pour un faisceau constant $\Lambda$ sur $X$ , on a

\begin{equation*}
R^i \Phi(\Lambda) \;\; = \widetilde{H}^i(V, \Lambda) \quad \text{concentré en } O \quad 
\end{equation*}

\SGApage{II}{131}

\SGAvar{ (groupes de cohomologie réduit).

}{ (groupes de cohomologie réduits).

}{F7II-XIV-014}1.4.7. Le cas qui nous intéresse est celui où $X_\eta$ est une variété différentiable et où $V \subset X_\eta$ est une sous-variété à bord, d'intérieur $V^o$ ; $X-\{O\}$ est alors une variété différentiable, obtenue en recollant $\beta^* X_\eta$ et le fibré constant sur $D$ de fibre $X_\eta - V$ . L'application

\begin{equation}\tag{1.4.7.1}
H^i_c(V^o, \Lambda) \overset{\sim}{\longrightarrow} H^i_V(X_\eta , \Lambda)
\end{equation}
est un isomorphisme. Ces groupes sont aussi notés $H^i(V \bmod \partial V, \Lambda)$ . L'isomorphisme (1.4.7.1) définit un isomorphisme

\begin{equation}\tag{1.4.7.2}
R\Gamma_c(V^o, \Lambda ) \overset{\sim}{\longrightarrow} R\Gamma_{\{O\}} \; R\Psi_\eta(\Lambda ) \quad ,
\end{equation}
et la \underline{variation} est définie (par le procédé de XIII 2.4.5 et via (1.4.7.2)) par un morphisme de source la cohomologie réduite de $V$

\SGAvar{\begin{equation}\tag{1.4.7.3}
\mathrm{Var} \; : \; R\widetilde{\Gamma}(V, \Lambda) \longrightarrow R\Gamma_c(V^o, \Lambda) \quad .
\end{equation}
$T$ induit un automorphisme, encore noté $T$ , de la paire $(V, \partial V)$ et la restriction $\partial T$ de $T$ à $\partial V$ est l'identité. Il doit être vrai que la variation 1.4.7.3 coincide avec le morphisme variation classique (3.1.9) attaché à un tel automorphisme $T$ . Un énoncé de ce type sera démontré en 3.1.10.


}{\begin{equation}\tag{1.4.7.3}
\mathrm{Var} \; : \; R\widetilde{\Gamma}(V, \Lambda) \longrightarrow R\Gamma_c(V^o, \Lambda) \quad .
\end{equation}
$T$ induit un automorphisme, encore noté $T$ , de la paire $(V, \partial V)$ et la restriction $\partial T$ de $T$ à $\partial V$ est l'identité. Il doit être vrai que la variation 1.4.7.3 coïncide avec le morphisme variation classique (3.1.9) attaché à un tel automorphisme $T$ . Un énoncé de ce type sera démontré en 3.1.10.


}{F7II-XIV-015}\SGAsectiontoc{§ 2. Le théorème de comparaison}{2}
§ 2. \underline{Le théorème de comparaison}

2.1. Pour tout schéma $X$ de type fini sur $\mathbb{C}$ , nous noterons $X_{cl}$ l'espace $X(\mathbb{C})$ , muni de sa topologie usuelle.

Soit $S$ une courbe lisse sur $\mathbb{C}$ , i.e. un schéma lisse séparé de type fini sur $\mathbb{C}$ , purement de dimension 1. Soient $O \in S(\mathbb{C})$ , $f : X \longrightarrow S$ un $S$-schéma séparé de type fini et $\Lambda$ un anneau noethérien de torsion (par exemple $\mathbb{Z}/n\,\mathbb{Z}$). On pose $X_o = f^{-1}(O)$ .


\SGApage{II}{132}

\SGAvar{2.2. Nous noterons $S^h$ le trait hensélien henselisé de $S$ en $O$ , et $\eta$ le point générique de $S^h$ . On dispose d'un foncteur

}{2.2. Nous noterons $S^h$ le trait hensélien hensélisé de $S$ en $O$ , et $\eta$ le point générique de $S^h$ . On dispose d'un foncteur

}{F7II-XIV-016}\begin{equation*}
\begin{aligned}
\Psi : (\text{faisceaux sur } X) &\xrightarrow{\ \text{restriction}\ } (\text{faisceaux sur } X \times_S S^h) \longrightarrow \\
&\xrightarrow{\ \Psi\ } (\text{faisceaux sur } X_o \times S^h) \ .
\end{aligned}
\end{equation*}
Le foncteur dérivé de $\Psi$ est le composé

\begin{equation*}
R\Psi : D(X, \Lambda) \longrightarrow D(X \times_S S^h, \Lambda) \xrightarrow{\ R\Psi\ } D(X_o \times S^h, \Lambda) \ .
\end{equation*}
2.3. Soit $D \subset S_{cl}$ un disque contenant $O$ . Le foncteur composé $\Psi_{cl}$ (ou simplement $\Psi$ )

\begin{equation*}
\begin{aligned}
\Psi_{cl} : (\text{faisceaux sur } X_{cl}) &\xrightarrow{\ \text{restriction}\ } (\text{faisceaux sur } X_{cl} \times_{S_{cl}} D) \longrightarrow \\
&\xrightarrow{\ \Psi\ } (\text{faisceaux sur } (X_o)_{cl} \times \mathcal{D})
\end{aligned}
\end{equation*}
est indépendant (à isomorphisme unique près) du choix de $D$ (rappelons que $\mathcal{D}$ ne change pas quand on remplace $D$ par un disque plus petit). Son dérivé est le composé

\begin{equation*}
R\Psi_{cl} : D(X_{cl}, \Lambda) \longrightarrow D(f_{cl}^{-1}(D), \Lambda) \xrightarrow{\ R\Psi\ } D(X_o \times \mathcal{D}, \Lambda) \ .
\end{equation*}
2.4. On définit de même

\begin{equation*}
\begin{aligned}
R\Psi_{\eta} &: D(X-X_o, \Lambda) \longrightarrow D(X_o \times \eta, \Lambda) \\
R\Psi_{\eta\ cl} &: D((X-X_o)_{cl}, \Lambda) \longrightarrow D(X_o \times \mathcal{D}^*, \Lambda) \ .
\end{aligned}
\end{equation*}
2.5. La construction 1.1.11 définit un morphisme de topos de $\mathcal{D}^*$ dans $\eta$ qui se prolonge en

\begin{equation*}
\varepsilon : \mathcal{D} \longrightarrow S^h \ .
\end{equation*}
Ce morphisme, et $\varepsilon : X_{o\ cl} \longrightarrow X_o$ (SGA4 XI 4), définissent

\begin{equation*}
\varepsilon : X_{o\ cl} \times \mathcal{D} \longrightarrow X_o \times S^h \ .
\end{equation*}

\SGApage{II}{133}

Si $\widetilde{D}^*$ est un revêtement universel de $D^*$ et que $k(\bar{\eta})$ est la clôture algébrique correspondante (1.1.10) de $k(\eta)$ , on a

\begin{equation*}
\varepsilon^*(F_o, F_\eta, \alpha) \ = \ (\varepsilon^* F_o, \ \varepsilon^* F_{\bar{\eta}}, \ \varepsilon^* \alpha) \ .
\end{equation*}
2.6. Soit $\widetilde{D}^*$ un revêtement universel de $D^*$ , $k(\bar{\eta})$ la clôture algébrique correspondante de $k(\eta)$ et $\bar{S}$ le normalisé de $S^h$ dans $k(\bar{\eta})$ . Ecrivons $\bar{S}$ comme limite projective de $S$-schémas affines $S_i$ ; pour chaque $i$ , et après avoir remplacé $D$ par un disque plus petit $D_i$ , on dispose d'un morphisme naturel de $S_{cl}$-espaces $\widetilde{D}_i \xrightarrow{\ a_i\ } S_{i_{cl}}$ . Soit $j_i : U_i \hookrightarrow S_i$ l'ouvert de $S_i$ complément de l'image $O$ du point fermé de $\bar{S}$ , et $i : O \hookrightarrow S_i$ .

Soit $F$ un faisceau sur $X$ . Soit $F_{cl} = \varepsilon^* F$ . Notons encore $j_i$ l'inclusion de $X \times_S U_i$ dans $X \times_S S_i$ , $i$ celle de $X_o$ dans $X \times_S S_i$ , et $F$ l'image réciproque de $F$ sur $X \times_S S_i$ . On a

\begin{equation*}
\Psi(F) \ = \ (i^* F, \ \varinjlim \, i^* j_{i*} \, j_i^* F, \ \varinjlim \, (\varphi_i)) \ .
\end{equation*}
Les foncteurs $\varepsilon^*$ et image réciproque commutent. Le morphisme $\varepsilon^* j_{i*} \longrightarrow (j_{i_{cl}})_* \, \varepsilon^*$ fournit donc

\begin{equation*}
c_i \ : \ (i^* \, j_{i*} \, j_i^* \, F)_{cl} \longrightarrow i^* \, j_{i*} \, j_i^* \, (F_{cl}) \ .
\end{equation*}
Soit $\bar{j}$ (resp. $\bar{j}_i$ ) l'inclusion de $\widetilde{D}^*$ (resp. $\widetilde{D}^*_i$ ) dans $\widetilde{D}$ (resp. $\widetilde{D}_i$), et les inclusions $X \times_S \widetilde{D}^* \hookrightarrow X \times_S \widetilde{D}$ , $X \times_S \widetilde{D}^*_i \hookrightarrow X \times_S \widetilde{D}_i$ qui s'en déduisent. La flèche d'adjonction $\mathrm{Id} \longrightarrow a_{i*} \, a_i^*$ définit

\begin{equation*}
i^* \, j_{i*} \, j_i^* \, F_{cl} \ \longrightarrow \ i^* \, \bar{j}_{i*} \, \bar{j}_i^* \, \bar{F}_{cl} \ = \ i^* \, \bar{j}_* \, \bar{j}^* \, \bar{F}_{cl} \ ,
\end{equation*}
d'où par passage à la limite

\begin{equation*}
c_\eta \ : \ \Psi_{\bar{\eta}} \, (F)_{cl} \longrightarrow \Psi_{cl}(F_{cl}) \ .
\end{equation*}
Ce morphisme définit le \underline{morphisme de comparaison} 


\SGApage{II}{134}

\begin{equation}\tag{2.6.1}
c \ : \ \varepsilon^* \Psi(F) \longrightarrow \Psi_{\mathrm{cl}}(\varepsilon^* F) \ .
\end{equation}
On définit de même

\begin{equation}\tag{2.6.2}
c_\eta \ : \ \varepsilon^* \Psi_\eta(F) \longrightarrow \Psi_\eta(\varepsilon^* F) \ .
\end{equation}
2.7. Les morphismes 2.6.1 et 2.6.2 se dérivent en

\begin{equation*}
c \ : \ \varepsilon^* R\Psi \longrightarrow R\Psi_{\mathrm{cl}} \circ \varepsilon^*
\end{equation*}
\begin{equation*}
c_\eta \ : \ \varepsilon^* R\Psi_\eta \longrightarrow R\Psi_{\eta\,\mathrm{cl}} \circ \varepsilon^* \ .
\end{equation*}
\underline{Théorème} 2.8. \underline{Soit} $K \in \mathrm{Ob}\ D(X, \Lambda)$ (\underline{resp}. $\mathrm{Ob}\ D(X-X_o, \Lambda)$) \underline{un complexe de faisceaux de} $\Lambda$ \underline{-modules dont les faisceaux de cohomologie} $\underline{H}^i(K)$ \underline{sont constructibles. Alors, le morphisme de comparaison}

\begin{equation*}
c \ : \ \varepsilon^* R\Psi(K) \longrightarrow R\Psi_{\mathrm{cl}}(\varepsilon^* K)
\end{equation*}
\begin{equation*}
(\mathrm{resp.} \qquad c \ : \ \varepsilon^* R\Psi_\eta(K) \longrightarrow R\Psi_{\eta\,\mathrm{cl}}(\varepsilon^* K) \ \ )
\end{equation*}
\underline{est un isomorphisme.}

\SGAvar{Il suffit de prouver l'assertion pour $R\Psi_\eta$ ; de plus, les deux membres de (2.7) étant des foncteurs triangulés de dimension cohomologique finie, il suffit de traiter le cas où $K$ est réduit à un faisceau $F$ (placé en degré $0$). La démonstration est toute pareille à celle du théorème de finitude XIII 2.3.1 : $R\Psi$ et $R\Psi_{\mathrm{cl}}$ ayant les mêmes propriétés de variance, on se ramène par les dévissages de loc. cit. (où on se permet de retrécir $S$ ) au cas où $X$ est régulier, $f^{-1}(O)$ un diviseur à croisements normaux dans $X$ et où $F$ est $\mathbb{Z}/\ell$ . On compare alors I 3.3 à son analogue transcendant.


}{Il suffit de prouver l'assertion pour $R\Psi_\eta$ ; de plus, les deux membres de (2.7) étant des foncteurs triangulés de dimension cohomologique finie, il suffit de traiter le cas où $K$ est réduit à un faisceau $F$ (placé en degré $0$). La démonstration est toute pareille à celle du théorème de finitude XIII 2.3.1 : $R\Psi$ et $R\Psi_{\mathrm{cl}}$ ayant les mêmes propriétés de variance, on se ramène par les dévissages de loc. cit. (où on se permet de rétrécir $S$ ) au cas où $X$ est régulier, $f^{-1}(O)$ un diviseur à croisements normaux dans $X$ et où $F$ est $\mathbb{Z}/\ell$ . On compare alors I 3.3 à son analogue transcendant.


}{F7II-XIV-017}\SGApage{II}{135}


\SGAsectiontoc{§ 3. Singularités isolées}{3}
§ 3. \underline{Singularités isolées}


3.1. \underline{La variété des cycles évanescents}.

3.1.1. Soient $D = \{z \in \mathbb{C} \mid |z| < 1\}$ , $X$ un espace analytique, $0 \in X$ et $f : X \to D$ un morphisme tel que $f(0) = 0$ . On suppose que $f$ est lisse en dehors de $0$ (en particulier, $X$ est lisse en dehors de $0$ ). On se propose d'étudier les cycles évanescents au voisinage de $0$ ; on ne restreint donc pas la généralité en supposant $X$ plongé dans $\mathbb{C}^n \times D \subset \mathbb{C}^{n+1}$ , de façon compatible à la projection sur $D$ , et de sorte que $0$ s'envoie sur $0$ .

3.1.2. Soit $\delta(x)$ une fonction analytique réelle sur $X$ (restriction à $X$ d'une fonction analytique réelle sur $\mathbb{C}^{n+1}$ ), telle que, pour $x \in X_o = f^{-1}(0)$ , on ait

(a) $\delta(x) \geq 0$ ;

(b) $\delta(x) = 0 \Leftrightarrow x = 0$ .

On prendra par exemple pour $\delta$ le carré de la distance euclidienne, dans $\mathbb{C}^{n+1}$ , de $x$ à $0$ . Le lieu des points critiques de la restriction de $\delta$ à $X_o$ est analytique réel donc, dans un voisinage convenable de $0$ , n'a qu'un nombre fini de composantes connexes. Remplaçant $X$ par un voisinage ouvert de $0$ et appliquant Sard, on peut donc obtenir la situation suivante

(c) pour $r > 0$ convenable, $X_o \cap \delta^{-1}([0, r])$ est compact,

(d) sur $X_o$ , $\delta$ n'a pas de valeur critique $\leq r$ .

Quitte à rétrécir $X$ davantage, on a alors

(e) il existe une fonction $C^\infty$ $\tau(s)$ $(0 \leq s \leq r)$ , avec $\tau(0) = 0$ et $\tau'(s) > 0$ pour $s > 0$ , telle que, pour $0 < s \leq r$ et $0 \leq |t|^2 \leq \tau(s)$ , $X_t = f^{-1}(t)$ soit transverse à l'hypersurface réelle $\delta(x) = s$ de $X$ et que la boîte

\begin{equation*}
B_{t,d} = \{x \in X \mid \delta(x) \leq d \ \text{ et } \ |f(x)|^2 < t\}
\end{equation*}
soit propre sur le disque ouvert $D_t$ de rayon $t^{1/2}$ .


\SGApage{II}{136}

Soit $L = X_o \cap \delta^{-1}(r)$ . D'après (c) et (d), il existe un homéomorphisme de $X_o \cap \delta^{-1}([0, r])$ avec le cône ($L \times [0, r]/$ contraction de $L \times \{0\}$ en un point), transformant $\delta$ en la $2^e$ coordonnée, et $C^\infty$ en dehors de l'origine.

\SGAvar{3.1.3. Soient $d \leq r$ et $t \leq r(d)$ . En dehors de $0$ , $B_{t,d}$ est une variété à coins et la projection $f$ de $B_{t,d}$ sur $D_t$ une fibration. En particulier

}{3.1.3. Soient $d \leq r$ et $t \leq \tau(d)$ . En dehors de $0$ , $B_{t,d}$ est une variété à coins et la projection $f$ de $B_{t,d}$ sur $D_t$ une fibration. En particulier

}{F7II-XIV-018}(a) $B_{t,d}$ induit sur $D_t^* = D_t - \{0\}$ un fibré $C^\infty$ en variétés compactes à bord;

(b) le bord $\delta(x) = d$ de ce fibré est prolongé en un fibré $C^\infty$ $\partial B_{t,d}$ en variétés compactes sur $D_t$ .

\SGAvar{Soit $y \in D_t^*$ . Nous appellerons la variété à bord $V = f^{-1}(y)$ la \underline{variété des cycles évanescents} . Le disque étant contractile, le fibré $\partial B_{t,d}$ sur $D_t$ est isomorphe au fibré constant de fibre $\partial V$ , et l'ensemble des isomorphismes de fibré $\partial B_{t,d} \xrightarrow[\alpha]{\sim} \partial V \times D_t$ , qui soit l'identité sur la fibre en $y$ , est contractile. Choisissons en un, noté $\alpha$ . Relevons le champ de vecteurs $\frac{\partial}{\partial\theta}$ sur $S_{|y|}$ (circonférence de rayon $|y|$ ) en un champ de vecteur $u$ sur $f^{-1}(S_{|y|})$ qui, au bord, soit compatible à $\alpha$ $(d\alpha(u) = (\frac{\partial}{\partial\theta} , 0))$ . L'ensemble de ces relèvements est contractile. Par intégration (tourner une fois dans le sens positif), $u$ définit un difféomorphisme $T$ de $V$ sur elle-même, qui est l'identité sur $\partial V$ ; c'est la \underline{transformation de monodromie} . Elle est bien définie modulo isotopies mod $\partial V$ .

En un sens qu'on laisse au lecteur le soins de préciser, $(V, T)$ est indépendant des choix de $t$ et $s$ . J'ignore s'il est indépendant du choix de $\delta$ . On vérifie toutefois qu'il en est ainsi lorsqu'on se restreint à prendre pour $\delta$ la restriction à $X$ d'une fonction $f$ sur $\mathbb{C}^{n+1}$ telle que, sur l'espace tangent de Zariski de $X_o$ en $0$ , on ait $df = 0$ et $d^2f > 0$ (définie positif). Le point est que, pour $\delta$ et $\delta'$ de ce type, on peut choisir $r$ comme en 3.1.2 (c) (d) uniformément pour tous les $\lambda\delta + (1-\lambda)\delta'$ $(0 \leq \lambda \leq 1)$ .

}{Soit $y \in D_t^*$ . Nous appellerons la variété à bord $V = f^{-1}(y)$ la \underline{variété des cycles évanescents} . Le disque étant contractile, le fibré $\partial B_{t,d}$ sur $D_t$ est isomorphe au fibré constant de fibre $\partial V$ , et l'ensemble des isomorphismes de fibré $\partial B_{t,d} \xrightarrow[\alpha]{\sim} \partial V \times D_t$ , qui soient l'identité sur la fibre en $y$ , est contractile. Choisissons en un, noté $\alpha$ . Relevons le champ de vecteurs $\frac{\partial}{\partial\theta}$ sur $S_{|y|}$ (circonférence de rayon $|y|$ ) en un champ de vecteur $u$ sur $f^{-1}(S_{|y|})$ qui, au bord, soit compatible à $\alpha$ $(d\alpha(u) = (\frac{\partial}{\partial\theta} , 0))$ . L'ensemble de ces relèvements est contractile. Par intégration (tourner une fois dans le sens positif), $u$ définit un difféomorphisme $T$ de $V$ sur elle-même, qui est l'identité sur $\partial V$ ; c'est la \underline{transformation de monodromie} . Elle est bien définie modulo isotopies mod $\partial V$ .

En un sens qu'on laisse au lecteur le soin de préciser, $(V, T)$ est indépendant des choix de $t$ et $d$ . J'ignore s'il est indépendant du choix de $\delta$ . On vérifie toutefois qu'il en est ainsi lorsqu'on se restreint à prendre pour $\delta$ la restriction à $X$ d'une fonction $f$ sur $\mathbb{C}^{n+1}$ telle que, sur l'espace tangent de Zariski de $X_o$ en $0$ , on ait $df = 0$ et $d^2f > 0$ (définie positive). Le point est que, pour $\delta$ et $\delta'$ de ce type, on peut choisir $r$ comme en 3.1.2 (c) (d) uniformément pour tous les $\lambda\delta + (1-\lambda)\delta'$ $(0 \leq \lambda \leq 1)$ .

}{F7II-XIV-019}3.1.4. Soient $(V, T)$ et $(V_1, T_1)$ déduit de $(V, T)$ par adjonction d'un collier:
\SGApage{II}{137}
$V_1 - V^o \simeq \partial V \times [0, 1]$ et $T_1$ est l'identité sur $V_1 - V^o$ .

Soit $f' : W \longrightarrow D_t$ l'espace sur $D_t$ suivant (cf. 1.4.3)

(a) Soient $\widetilde{D}_t^*$ un revêtement universel de $D_t^*$ de point de base $y$ et $W^*$ le fibré sur $D_t^*$ quotient de $V_1 \times \widetilde{D}_t^*$ par la relation engendrée par $(v, y) \simeq (T_1 v, e^{2\pi i} y)$ .

(b) Puisque $T_1 | V_1 - V^o$ est trivial, on peut recoller $W^*$ et le fibré trivial sur $D_t$ de fibre $V_1 - V$ en un espace $W_1$ sur $D_t$ .

(c) $W = W_1 \cup \{0\}$ , avec $f'(0) = 0$ . $W$ est muni de l'unique topologie induisant celle de $W_1$ , pour laquelle $f'$ soit continue et que ${f'}^{-1}(\bar{D}_{t'})$ $(t' < t)$ soit la compactification par un point de $W_1 \cap {f'}^{-1}(\bar{D}_{t'})$ .

On munit $W - \{0\}$ de la connexion intégrable (1.3.4) pour laquelle la trivialisation fournie par (a) de l'image réciproque de $W$ sur $\widetilde{D}^*$ soit horizontale.

\SGAvar{\underline{Proposition} 3.1.5. \underline{Il existe un homéomorphisme} $\Psi$ , \underline{compatible à la projection sur} $D_t$ \underline{et envoyant} $0$ \underline{sur} $0$ , \underline{de} $B_{s,t}$ \underline{avec} $W$ ; \underline{cet homéomorphisme peut être choisi un difféomorphisme en dehors de} $0$ .

}{\underline{Proposition} 3.1.5. \underline{Il existe un homéomorphisme} $\Psi$ , \underline{compatible à la projection sur} $D_t$ \underline{et envoyant} $0$ \underline{sur} $0$ , \underline{de} $B_{t,d}$ \underline{avec} $W$ ; \underline{cet homéomorphisme peut être choisi un difféomorphisme en dehors de} $0$ .

}{F7II-XIV-020}Ce résultat ne sera pas utilisé dans la suite du séminaire.

\SGAvar{Soient $d_o < d_1 \leq r$ , $t_o \leq \tau(d_o)$ et $y \in D_{t_o}^*$ . On a $B_{t_o,d_o} \subset B_{t_o, d_1}$ et $B_{t_o,d_1} - B_{t_o,d_o}^o$ est un fibré sur $D_{|t|_o}$ . Trivialisons-le et au-dessus de la circonférence $S_{|y|}$ , relevons le champs de vecteur $\frac{\partial}{\partial \theta}$ de façon compatible à cette trivialisation. Soit $T_1$ l'endomorphisme correspondant de $V_1 = f^{-1}(y) \cap B_{t_o,d_1}$ , et $V_o = f^{-1}(y) \cap B_{t_o,d_o}$ . Alors, $(V_o \subset V_1, T_1)$ est isomorphe à $(V \subset V_1, T_1)$ de 3.1.4. Soit $W$ construit à l'aide de $(V_o \subset V_1, T_1)$ . On a ${f'}^{-1}(y) = V_1$ , canoniquement. La démonstration montrera qu'on peut prendre $\Psi$ qui induise l'identité $f^{-1}(y) \xrightarrow{\sim} {f'}^{-1}(y)$ .

}{Soient $d_o < d_1 \leq r$ , $t_o \leq \tau(d_o)$ et $y \in D_{t_o}^*$ . On a $B_{t_o,d_o} \subset B_{t_o, d_1}$ et $B_{t_o,d_1} - B_{t_o,d_o}^o$ est un fibré sur $D_{t_o}$ . Trivialisons-le et au-dessus de la circonférence $S_{|y|}$ , relevons le champ de vecteurs $\frac{\partial}{\partial \theta}$ de façon compatible à cette trivialisation. Soit $T_1$ l'endomorphisme correspondant de $V_1 = f^{-1}(y) \cap B_{t_o,d_1}$ , et $V_o = f^{-1}(y) \cap B_{t_o,d_o}$ . Alors, $(V_o \subset V_1, T_1)$ est isomorphe à $(V \subset V_1, T_1)$ de 3.1.4. Soit $W$ construit à l'aide de $(V_o \subset V_1, T_1)$ . On a ${f'}^{-1}(y) = V_1$ , canoniquement. La démonstration montrera qu'on peut prendre $\Psi$ qui induise l'identité $f^{-1}(y) \xrightarrow{\sim} {f'}^{-1}(y)$ .

}{F7II-XIV-021}Les dessins suivants ($f$, et sa restriction à $f^{-1}(\mathbb{R})$ ) résument la démonstration


\SGApage{II}{138}

\begin{center}\emph{[Quatre dessins a main levee occupant le haut de la page, disposes en deux colonnes. En haut a gauche : un cylindre vertical vu en perspective, l'ellipse de bord superieure et l'ellipse de bord inferieure tracees en trait plein sur leur moitie avant et en pointille sur leur moitie arriere ; un segment vertical joint les deux ellipses en passant par l'axe et porte en son milieu un point epais, etiquete $0$ a sa gauche. En bas a gauche : un cercle trace a main levee, un petit point en son centre, etiquete $D_t$ en bas a droite du cercle. En haut a droite : un carre ; a l'interieur, des droites horizontales allant d'un bord a l'autre dans la partie superieure et dans la partie inferieure ; sous les droites du haut, une courbe d'un bord a l'autre qui s'abaisse vers le milieu, et au-dessus des droites du bas, une courbe qui se releve vers le milieu ; au centre un point epais ou viennent se rejoindre deux branches venues de la gauche et deux branches venues de la droite (image du pincement nodal), avec en outre six traits, trois de chaque cote (un oblique en haut, un horizontal, un oblique en bas), partant pres des bords gauche et droit et diriges vers ce point sans l'atteindre. En bas a droite : un segment horizontal portant un point epais en son milieu, a l'aplomb du point central du carre.]}\end{center}

\SGAvar{Soit $B^+ \subset B_{t_o,d_1}$ l'ensemble des $x$ tels que $|f(x)|^2 \leq \tau(\delta(x))$ . Soit $\varphi_n$ une suite de fonction $C^\infty$ de $[0, t_o]$ dans $(0, d_1]$ telle que, pour une suite $\tau_o > r_1 > \ldots$ tendant vers $0$ , on ait

}{Soit $B^+ \subset B_{t_o,d_1}$ l'ensemble des $x$ tels que $|f(x)|^2 \leq \tau(\delta(x))$ . Soit $\varphi_n$ une suite de fonctions $C^\infty$ de $[0, t_o]$ dans $(0, d_1]$ telle que, pour une suite $\tau_o > r_1 > \ldots$ tendant vers $0$ , on ait

}{F7II-XIV-022}(a) pour $|t|^2 \geq r_i$ , $\varphi_i$ est la réciproque de la fonction $\tau$ .

(b) pour $|t|^2 < r_i$ et $j > i$ , $\varphi_j(t) < \varphi_i(t)$

\begin{center}\emph{[Dessin a main levee. Axe vertical etiquete $d$ a son sommet, axe horizontal etiquete $t$ a droite. Sur l'axe vertical le niveau $d_1$ est marque ; sur l'axe horizontal l'abscisse $t_o$ est marquee. Un rectangle a pour sommets l'origine, le point de hauteur $d_1$ sur l'axe vertical et le point d'abscisse $t_o$ sur l'axe horizontal. Une courbe croissante et concave part de l'origine et monte jusqu'au bord droit du rectangle, a l'abscisse $t_o$ ; elle est etiquetee $t=\tau(d)$ a droite de son extremite. Deux autres courbes partent de l'axe vertical, a deux hauteurs inferieures a $d_1$, descendent legerement puis se raccordent a la courbe precedente ; la plus haute des deux est etiquetee $\varphi_o$.]}\end{center}

.

\SGAvar{L'ensemble $B^+ - \{0\}$ est réunion des fibrés sur $D_t$ $B_i^+$ : $\delta(x) \geq \varphi_i(|f(x)|^2)$ . Un argument standard, utilisant que $\delta$ n'a pas de points critiques sur les fibres de $B^+ - \{0\} \longrightarrow D_t$ , montre qu'il existe un difféomorphisme $\alpha_i$ de $B_i^+$ avec $B_{i+1}^+$ , compatible à la projection sur $D_t$ , et l'identité en dehors d'un petit voisinage du sous-espace
\SGApage{II}{139}
$(\delta(x) = \varphi_i(|f(x)|^2)$ , $|f(x)|^2 \leq r_i)$ de $B_i^+$ . Composant les $\alpha_i$ et passant à la limite, on obtient un difféomorphisme $\alpha$ de $B_o^+ - \{x | (x) = 0$ et $\delta(x) = \varphi_o(0)\}$ avec $B^+ - \{0\}$ . Transportons par $\alpha$ une trivialisation du fibré $B_o^+$ sur $D_t$ pour obtenir une connexion intégrable $\nabla'$ sur $B^+ - \{0\}$ . On peut alors

}{L'ensemble $B^+ - \{0\}$ est réunion des fibrés sur $D_t$ $B_i^+$ : $\delta(x) \geq \varphi_i(|f(x)|^2)$ . Un argument standard, utilisant que $\delta$ n'a pas de points critiques sur les fibres de $B^+ - \{0\} \longrightarrow D_t$ , montre qu'il existe un difféomorphisme $\alpha_i$ de $B_i^+$ avec $B_{i+1}^+$ , compatible à la projection sur $D_t$ , et l'identité en dehors d'un petit voisinage du sous-espace
\SGApage{II}{139}
$(\delta(x) = \varphi_i(|f(x)|^2)$ , $|f(x)|^2 \leq r_i)$ de $B_i^+$ . Composant les $\alpha_i$ et passant à la limite, on obtient un difféomorphisme $\alpha$ de $B_o^+ - \{x | f(x) = 0$ et $\delta(x) = \varphi_o(0)\}$ avec $B^+ - \{0\}$ . Transportons par $\alpha$ une trivialisation du fibré $B_o^+$ sur $D_t$ pour obtenir une connexion intégrable $\nabla'$ sur $B^+ - \{0\}$ . On peut alors

}{F7II-XIV-023}(a) prolonger la restriction de cette connexion à $f^{-1}(S_{|y|}) \cap B^+$ en une connexion sur $f^{-1}(S_{|y|}) / S_{|y|}$ (ce prolongement donnant $T_1$ ) .

(b) prolonger à $B_{t_o,d_1}$ le relèvement sur $B^+ - \{0\}$ fourni par $\nabla'$ du champ de vecteurs radial sur $D_t$ .

Il existe alors une et une seule connexion intégrable $\nabla$ sur $B_{t_o,d_1}$ qui prolonge $\nabla'$ , (a) et (b), et un et un seul difféomorphisme horizontal de $B_{t_o,d_1} - \{0\}$ avec $W - \{0\}$ qui soit l'identité sur $f^{-1}(y)$ . Il répond au problème .

3.1.6. \underline{Une généralisation}. Ce qui précède se généralise au cas où $f : X \to D$ est un morphisme tel que $0 \in X$ soit un point de non lissité isolé \underline{dans} $X_o$ . Le lieu de non lissité de $f$ est alors fini sur $D$ (quitte à rapetisser $X$ et $D$ ). La variété des cycles évanescents $V$ a alors un ensemble fini $\Sigma$ de points critiques; $V - \Sigma$ est une variété différentiable à bord (de bord compact, mais non compacte). La transformation de monodromie $T$ est l'identité sur le bord et laisse $\Sigma$ stable. Enfin, 3.1.5 reste valable, l'homéomorphisme fibré $\Psi$ étant un difféomorphisme en dehors du lieu critique, qu'il envoie sur ce qui dans $W$ provient de $\Sigma$ .

3.1.7. La proposition 3.1.5 permet d'utiliser les calculs de 1.4. On trouve ainsi que

(A) Les groupes de cycles évanescents $(R^i \Phi(\Lambda))_o$ sont les groupes de cohomologie réduits, à valeurs dans $\Lambda$ , de la variété des cycles évanescents $V$ .

(B) Les groupes $H^i_{\{0\}}(X_o, R\Psi(\Lambda))$ sont les groupes de cohomologie à support propre de $V^o$ , à valeurs dans $\Lambda$ .

(C) La variation 


\SGApage{II}{140}

\begin{equation*}
 \mathrm{Var} : \widetilde{H}^i(V, \Lambda) \longrightarrow H^i_c(V^o, \Lambda) = H^i(V \bmod \partial V, \Lambda)
\end{equation*}
est la variation classique, définie par l'automorphisme $T$ de $V$ (qui est l'identité sur $\partial V$ ).

3.1.8. Les assertions (A) (B) (C) sont en fait plus élémentaires que 3.1.5; en voici une démonstration-construction directe, qui nous dispensera de vérifier que les isomorphismes (A) (B) sont indépendants des choix faits dans la construction 3.1.5.

Soient $d \leq r$, $t \leq \tau(d)$ et $B$ la boîte ouverte

\SGAvar{\begin{equation*}
B = \{x \mid \delta(n) < d \ \ \text{et} \ \ |f(n)|^2 < t\} \ \ .
\end{equation*}
Soit $\widetilde{D}^*$ un revêtement universel de $D^*$ , et

}{\begin{equation*}
B = \{x \mid \delta(x) < d \ \ \text{et} \ \ |f(x)|^2 < t\} \ \ .
\end{equation*}
Soit $\widetilde{D}^*$ un revêtement universel de $D^*$ , et

}{F7II-XIV-024}\begin{equation}\tag{3.1.8.1}
\begin{tikzcd}[column sep=large, row sep=large]
\{0\} \arrow[r, hook, "\overline{i}"] \arrow[d, equal] \& \widetilde{D}_t \arrow[d, "p"] \& \widetilde{D}^*_t \arrow[l, hook', "\overline{j}"'] \arrow[d, "p'"] \\
\{0\} \arrow[r, hook, "i"] \& D_t \& D^*_t \arrow[l, hook', "j"']
\end{tikzcd}
\end{equation}












déduit de 1.2.7.2 par restriction à $D_t$ . Soit

\begin{equation}\tag{3.1.8.2}
\begin{tikzcd}[column sep=large, row sep=large]
B_o \arrow[r, hook, "\overline{i}"] \arrow[d, equal] \& \overline{B} \arrow[d, "p"] \& \overline{B}^* \arrow[l, hook', "\overline{j}"'] \arrow[d, "p'"] \\
B_o \arrow[r, hook, "i"] \& B \& B^* \arrow[l, hook', "j"']
\end{tikzcd}
\end{equation}












image réciproque de $f : B \longrightarrow D$ sur (3.1.8.1). On note encore $f$ la projection de (3.1.8.2) sur (3.1.8.1).

\SGAvar{(A) Par définition, $R^i\Psi_\eta(\Lambda)_o$ est la limite inductive, pour $U$ parcourant les voisinages de $0$ dans $X$ , des groupes de cohomologie $H^i(U \times_D \widetilde{D}^*, \Lambda)$ . Faisons parcourir à $U$ le système fondamental de $0$ formé des boîtes $B$ comme plus haut. Le produit fibré $U \times_D \widetilde{D}^* = \overline{B}^*$ est un fibré de fibre $V^o$ sur $D^*_t$ (contractile). Le système inductif des $H^i(\overline{B}^*, \Lambda)$ est donc constant, de valeur $H^i(V^o, \Lambda)=H^i(V, \Lambda)$ .


}{(A) Par définition, $R^i\Psi_\eta(\Lambda)_o$ est la limite inductive, pour $U$ parcourant les voisinages de $0$ dans $X$ , des groupes de cohomologie $H^i(U \times_D \widetilde{D}^*, \Lambda)$ . Faisons parcourir à $U$ le système fondamental de $0$ formé des boîtes $B$ comme plus haut. Le produit fibré $U \times_D \widetilde{D}^* = \overline{B}^*$ est un fibré de fibre $V^o$ sur $\widetilde{D}^*_t$ (contractile). Le système inductif des $H^i(\overline{B}^*, \Lambda)$ est donc constant, de valeur $H^i(V^o, \Lambda)=H^i(V, \Lambda)$ .


}{F7II-XIV-025}\SGApage{II}{141}

Les isomorphismes obtenus

\begin{equation}\tag{3.1.8.3}
R^i\Psi_\eta(\Lambda)_o = H^i(V, \Lambda)
\end{equation}
peuvent se préciser en un isomorphisme composé (cf. (B) ci-dessous).

\begin{equation}\tag{3.1.8.4}
\begin{tikzcd}[column sep=large, row sep=large]
R\Psi_\eta(\Lambda)_o \& R\Gamma(B_o, R\Psi_\eta(\Lambda)) \arrow[l, "\sim"'] \& (Rf_* R\overline{j}_* \Lambda)_o \arrow[l, "\sim"'] \& R\Gamma(\widetilde{D}_t, Rf_* R\overline{j}_* \Lambda) \arrow[l, "\sim"'] \arrow[d, equal] \\
{} \& R\Gamma(V, \Lambda) \arrow[r, "\sim"] \& R\Gamma(V^o, \Lambda) \& R\Gamma(\overline{B}^*, \Lambda) \arrow[l, "\sim"']
\end{tikzcd}\quad .
\end{equation}













Le complexe $R\Phi(\Lambda)_o$ est défini par le triangle distingué

\begin{equation*}
\Lambda \longrightarrow (R\Psi_\eta \, \Lambda)_o \longrightarrow R\Phi(\Lambda)_o \longrightarrow \quad .
\end{equation*}
Les $R^i\Phi(\Lambda)_o$ sont donc les groupes de cohomologie réduits de $V$ .

(B) Soit $\Phi$ l'ensemble des fermés de $\overline{B}^*$ dont l'adhérence dans $\overline{B}$ est propre sur $\widetilde{D}$ . Les flèches suivantes sont des isomorphismes

\begin{equation}\tag{3.1.8.5}
\begin{tikzcd}[column sep=huge, row sep=huge]
R\Gamma_{\{0\}}(R\Psi_\eta(\Lambda)) \arrow[r, "\sim", "\text{\textcircled{1}}"'] \& R\Gamma_c(B_o, R\Psi_\eta(\Lambda)) \& (Rf_! R\overline{j}_* \Lambda)_o \arrow[l, "\sim"', "\text{\textcircled{2}}"] \\
R\Gamma_c(V^o, \Lambda) \& R\Gamma_\Phi(\overline{B}^*, \Lambda) \arrow[l, "\sim"', "\text{\textcircled{5}}"] \arrow[r, equal, "\text{\textcircled{4}}"'] \& R\Gamma(\widetilde{D}_t, Rf_! R\overline{j}_* \Lambda) \arrow[u, "\text{\textcircled{3}}", "\wr"']
\end{tikzcd}\quad .
\end{equation}












( \textcircled{1} parce que $B_o$ est un cône (3.1.2) et que $R^i\Psi_\eta(\Lambda)$ est constant en dehors de $0$ ;

\textcircled{2} par le théorème de changement de base pour les images directes à support propre;

\textcircled{3} parce que $Rf_! \, R\overline{j}_* \, \Lambda$ est constant en dehors de $0$ ,

\textcircled{4} est la version ``catégorie dérivée'' de la suite spectrale d'un foncteur composé;

\textcircled{5} parce que $\widetilde{D}^*$ est contractile).

L'isomorphisme (3.1.8.5) fournit les isomorphismes voulus 


\SGApage{II}{142}

\begin{equation}\tag{3.1.8.6}
 R^i\Gamma_{\{0\}}(R\Psi_\eta(\Lambda)) = H^i_c(V^o, \Lambda) .
\end{equation}
Nous vérifierons (C) en 3.1.10.

3.1.9. (cf. [3]) Soient $X$ un espace topologique, $i : Y \hookrightarrow X$ une partie fermée de $X$, $T$ un automorphisme de $X$ , induisant l'identité sur $Y$ , et $\Lambda$ un groupe abélien. Pour définir la variation $\mathrm{Var}$ , on choisit une résolution $T$-équivariante $\Lambda^*_X$ de $\Lambda$ sur $X$ , une résolution $\Lambda^*_Y$ du faisceau $\Lambda$ sur $Y$ , sur laquelle $T$ agit trivialement par convention, et un morphisme équivariant

\SGAvar{\begin{equation*}
\varphi : \Lambda^*_X \longrightarrow i_*\Lambda^*_Y ,
\end{equation*}
qui résoud le morphisme évident $\varphi_o : \Lambda_X \longrightarrow i_* \Lambda_Y$ . On suppose les faisceaux $\Lambda^i_X$ et $\Lambda^i_Y$ acycliques pour les foncteurs qui seront considérés. Passant aux sections globales, on trouve

}{\begin{equation*}
\varphi : \Lambda^*_X \longrightarrow i_*\Lambda^*_Y ,
\end{equation*}
qui résout le morphisme évident $\varphi_o : \Lambda_X \longrightarrow i_* \Lambda_Y$ . On suppose les faisceaux $\Lambda^i_X$ et $\Lambda^i_Y$ acycliques pour les foncteurs qui seront considérés. Passant aux sections globales, on trouve

}{F7II-XIV-026}\begin{equation*}
\Gamma(\varphi) : \Gamma(\Lambda^*_X) \longrightarrow \Gamma(\Lambda^*_Y) .
\end{equation*}
Supposons pour simplifier les $\Gamma(\varphi)^i$ surjectifs (sinon, on les rend surjectifs par addition à $\Gamma(\Lambda^*_X)$ du cône de $\Gamma(\Lambda^*_Y)$ ). On pose

\SGAvar{\begin{equation*}
R\Gamma (X \bmod Y, \Lambda) = \mathrm{Ker}(\Gamma (\varphi)) .
\end{equation*}
Sous des hypothèses raisonables, si $\Phi$ est l'ensemble des fermés de $X$ disjoints de $Y$ , on a

}{\begin{equation*}
R\Gamma (X \bmod Y, \Lambda) = \mathrm{Ker}(\Gamma (\varphi)) .
\end{equation*}
Sous des hypothèses raisonnables, si $\Phi$ est l'ensemble des fermés de $X$ disjoints de $Y$ , on a

}{F7II-XIV-027}\begin{equation*}
\Gamma_\Phi(X-Y, \Lambda^*_X) \overset{\approx}{\longrightarrow} R\Gamma (X \bmod Y, \Lambda)
\end{equation*}
(quasi-isomorphisme). L'endomorphisme $T-1$ de $\Gamma(\Lambda^*_X)$ ( $T$ agit par transport de structure) est nul sur $\Lambda$ (les constantes) et se factorise par $\mathrm{Ker}(\Gamma (\varphi))$ , d'où dans la catégorie dérivée et sous les mêmes hypothèses, un morphisme variation

\begin{equation*}
\mathrm{Var} : R\widetilde{\Gamma} (X, \Lambda) \longrightarrow R\Gamma_\Phi (X-Y, \Lambda) 
\end{equation*}

\SGApage{II}{143}

 (sans hypothèse, on trouve $\mathrm{Var} : R\widetilde{\Gamma} (X, \Lambda) \longrightarrow R\Gamma (X \bmod Y, \Lambda)$ ).

\SGAvar{Cette description se simplifie si $T$ est l'identité dans un voisinage de $Y$ et, dans ce coisinage, est l'identité sur $\Lambda^*_X$ . L'endomorphisme $T-1$ de $\Gamma (\Lambda^*_X)$ se factorise alors par $\Gamma_\Phi(\Lambda^*_X)$ , et fournit la variation.

}{Cette description se simplifie si $T$ est l'identité dans un voisinage de $Y$ et, dans ce voisinage, est l'identité sur $\Lambda^*_X$ . L'endomorphisme $T-1$ de $\Gamma (\Lambda^*_X)$ se factorise alors par $\Gamma_\Phi(\Lambda^*_X)$ , et fournit la variation.

}{F7II-XIV-028}\underline{Proposition} 3.1.10. \underline{Via les isomorphismes} 3.1.8 (A) (B), \underline{la variation} $(R\Phi_\eta(\Lambda))_o \to R\Gamma_{\{0\}}(R\Psi_\eta(\Lambda))$ \underline{est égale à la variation} 3.1.9 \underline{définie par l'automorphisme} $T$ \underline{de} $V$ , \underline{identique sur} $\partial V$ .

Le lecteur est invité à faire un petit dessin plutôt que de lire la démonstration. Soient ${}_fB$ la boîte fermée de 3.1.2 et ${}_fC \subset {}_fB$ un collier le long du bord ( ${}_fC \simeq \partial_fB \times [0, 1]$ comme fibré sur $D_t$ via un isomorphisme qui envoie $\partial_fB \times \{1\}$ identiquement sur $\partial_fB$ ). Soit

\SGAvar{\begin{equation*}
{}_f\alpha : {}_fC \overset{\sim}{\longrightarrow} (\partial V \times [0,1]) \times D_t
\end{equation*}
une trivialisation du fibré ${}_fC$ sur $D_t$ , dont la fibre en $y \in D_t$ induite l'identité $\partial V = \partial(f^{-1}(y)) \overset{\sim}{\longrightarrow} \partial V \times \{1\}$ . Soit $C = {}_fC \cap B$ et

}{\begin{equation*}
{}_f\alpha : {}_fC \overset{\sim}{\longrightarrow} (\partial V \times [0,1]) \times D_t
\end{equation*}
une trivialisation du fibré ${}_fC$ sur $D_t$ , dont la fibre en $y \in D_t$ induit l'identité $\partial V = \partial(f^{-1}(y)) \overset{\sim}{\longrightarrow} \partial V \times \{1\}$ . Soit $C = {}_fC \cap B$ et

}{F7II-XIV-029}\begin{equation*}
\alpha : C \overset{\sim}{\longrightarrow} (\partial V \times [0,1[) \times D_t
\end{equation*}
induit par ${}_f\alpha$ . On garde les notations de 3.1.8.2.

Soient $R$ la relation d'équivalence sur $\bar{B}$

$x = y$ ou $p(x)$, $p(y) \in C$ et $\alpha(p(x))$, $\alpha(p(y))$ ont même image dans $\partial V \times [0,1[$ .

Soient ${}_1\bar{B} = \bar{B}/R$ , ${}_1\bar{B}^* = \bar{B}^*/R$ et ${}_1\bar{C}^* = p^{-1}(C)/R$ . On dispose d'un diagramme commutatif 


\SGApage{II}{144}

\SGAvar{\begin{center}
\begin{tikzcd}
B_o \arrow[r, hook, "\bar{\imath}"] \arrow[d, equal] \& \bar{B} \arrow[d, "q"] \& {}_1\bar{B}^* \arrow[l, hook, "\bar{\jmath}"'] \arrow[d, "q'"] \\
B_o \arrow[r, hook, "\bar{\imath}"] \& {}_1\bar{B} \& {}_1\bar{B}^* \arrow[l, hook, "\bar{\jmath}"']
\end{tikzcd}\quad .
\end{center}














Soit $\Lambda^*$ une résolution $\pi_1(D^*)$-équivariante de $\Lambda$ par des faisceaux acycliques sur $\bar{B}^*$ . Par définition, la variation se calcule à partir du morphisme de complexe

}{\begin{center}
\begin{tikzcd}
B_o \arrow[r, hook, "\bar{\imath}"] \arrow[d, equal] \& \bar{B} \arrow[d, "q"] \& \bar{B}^* \arrow[l, hook, "\bar{\jmath}"'] \arrow[d, "q'"] \\
B_o \arrow[r, hook, "\bar{\imath}"] \& {}_1\bar{B} \& {}_1\bar{B}^* \arrow[l, hook, "\bar{\jmath}"']
\end{tikzcd}\quad .
\end{center}














Soit $\Lambda^*$ une résolution $\pi_1(D^*)$-équivariante de $\Lambda$ par des faisceaux acycliques sur $\bar{B}^*$ . Par définition, la variation se calcule à partir du morphisme de complexes

}{F7II-XIV-030}\begin{equation*}
\Lambda\,[o] \longrightarrow \bar{\imath}^*\, \bar{\jmath}_*\,(\Lambda^*) \qquad .
\end{equation*}
Puisque $\Lambda \stackrel{\approx}{\longrightarrow} Rq_*\Lambda$ et $\Lambda \stackrel{\approx}{\longrightarrow} Rq'_*\Lambda$ , on peut aussi la calculer  à l'aide d'une résolution équivariante de $\Lambda$ sur ${}_1\bar{B}^*$ et du morphisme

\SGAvar{\begin{equation*}
\Lambda[o] \longrightarrow \bar{\imath}^*\, {}_1\bar{\jmath}_*(\Lambda) \qquad .
\end{equation*}
Via (A) (B), on identifie alors la variation au morphisme composé

}{\begin{equation*}
\Lambda[o] \longrightarrow \bar{\imath}^*\, {}_1\bar{\jmath}_*(\Lambda^*) \qquad .
\end{equation*}
Via (A) (B), on identifie alors la variation au morphisme composé

}{F7II-XIV-031}\begin{equation}\tag{3.1.10.1}
\begin{tikzcd}[column sep=large]
R\Gamma\,(V,\Lambda) \arrow[r, "\sim"] \& R\Gamma\,(V^o,\Lambda) \& R\Gamma\,({}_1\bar{B}^*,\ \Lambda) \arrow[l, "\sim"', "\text{\textcircled{1}}"] \arrow[d, "\mathrm{Var}"] \\
{} \& R\Gamma_c(V^o,\Lambda) \& R\Gamma\,({}_1\bar{B}^*\ \mathrm{mod}\,{}_1\bar{C}^*,\Lambda) \arrow[l, "\sim"', "\text{\textcircled{2}}"]
\end{tikzcd}\quad ,
\end{equation}












\SGAvar{où \textcircled{1} et \textcircled{2} sont déduits de l'inclusion de $V^o$ dans ${}_1\bar{B}^*$ , et où Var est la variation 3.1.9 pour ${}_1\bar{B}^*$ et l'action du générateur positif de $\pi_1(D^*)$ (ce groupe agit sur ${}_1\bar{B}^*$ via son action 1.15; il agit trivialement sur ${}_1\bar{C}^*$ ).

}{où \textcircled{1} et \textcircled{2} sont déduits de l'inclusion de $V^o$ dans ${}_1\bar{B}^*$ , et où Var est la variation 3.1.9 pour ${}_1\bar{B}^*$ et l'action du générateur positif de $\pi_1(D^*)$ (ce groupe agit sur ${}_1\bar{B}^*$ via son action 1.1.5; il agit trivialement sur ${}_1\bar{C}^*$ ).

}{F7II-XIV-032}Choisissons une connexion $\nabla$ sur $B^*/D^*$ , compatible à la trivialisation de $C$ . Cette connexion trivialise $\bar{B}^*$ et définit par passage au quotient $u : {}_1\bar{B}^* \longrightarrow V$ . Pour $T$ défini comme en 3.1.3, le diagramme 


\SGApage{II}{145}

\begin{center}
\begin{tikzcd}
{}_1\bar{B}^* \arrow[r, "u"] \arrow[d, "T"'] \& V \arrow[d, "T"'] \\
{}_1\bar{B}^* \arrow[r, "u"] \& V
\end{tikzcd}
\end{center}








 est commutatif. L'application $u$ induit des isomorphismes en cohomologie, de sorte que (3.1.10) résulte de 3.1.10.1 et de la fonctorialité de la variation.

La proposition 3.1.10 sera censée justifier un calcul explicite que nous ferons plus tard et où, comme ici, $T$ sera l'identité dans un voisinage de $\partial V$ . On prendra une classe de cohomologie $c$ , représentée par un courant $C$ , et on utilisera que $\mathrm{Var}(c)$ est représenté par le courant $TC - C$ , à support compact dans $V^o$ .

3.1.11. Pour calculer en pratique $T$ , on procède souvent en deux étapes

(a) On fixe $t$ et $d$ avec $t \leq \tau(d)$ , on choisit $y \in D_t$ , et on choisit une connexion $\nabla$ sur $f^{-1}(S_{|y|}) \cap B_{t,d}$ . Cette connexion $\nabla$ définit un difféomorphisme $T_o$ de $V_o = f^{-1}(y) \cap B_{t,d}$ sur lui-même, qui n'est pas nécessairement l'identité sur le bord.

(b) On choisit une trivialisation de $\partial B_{t,d}$ sur $D_t$ . L'intégration de $\nabla$ fournit un difféomorphisme $\varphi(\theta)$ de $\partial V$ avec $f^{-1}(e^{2\pi i\theta}y) \cap \partial B_{t,d}$ , identifié à $\partial V$ par la trivialisation. Pour $0 \leq \theta \leq 1$ , $\varphi(\theta)$ est un chemin dans $\mathrm{Diff}(\partial V)$ , qui relie l'identité à $\partial T_o$ .

Le couple $(T_o, \varphi)$ obtenu dépend des choix arbitraires faits. L'ensemble de ces choix étant contractile, il n'en dépend pas à homotopie près. Pour $T$ du type 3.1.3, $(T$, chemin trivial$)$ est un $(T_o, \varphi)$ du type considéré ici.

3.1.12. On sait que $\mathrm{Diff}(V) \longrightarrow \mathrm{Diff}(\partial V)$ est une fibration, de sorte que $T \longrightarrow (T$ , trivial$)$ induit une bijection entre

\SGAvar{a) l'ensemble des difféomorphismes $T : V \longrightarrow V$ avec $\partial T = \mathrm{Id}$ , à isotopie près,avec
\SGApage{II}{146}
b) l'ensemble des couples $(T_o, \varphi)$ , $T_o : V \stackrel{\sim}{\longrightarrow} V$ , $\varphi$ classe d'homotopie de chemins de $\mathrm{Id}$ à $\partial T$ , à isotopie près.

}{a) l'ensemble des difféomorphismes $T : V \longrightarrow V$ avec $\partial T = \mathrm{Id}$ , à isotopie près,avec
\SGApage{II}{146}
b) l'ensemble des couples $(T_o, \varphi)$ , $T_o : V \stackrel{\sim}{\longrightarrow} V$ , $\varphi$ classe d'homotopie de chemins de $\mathrm{Id}$ à $\partial T_o$ , à isotopie près.

}{F7II-XIV-033}Un couple $(T_o, \varphi)$ calculé par (a) (b) ci-dessus détermine donc $T$ . Concrètement, on calcule $T$ à l'aide de $(T_o, \varphi)$ comme suit:

($\alpha$) On prend un collier $[0, 1] \times \partial V \stackrel{\alpha}{\hookrightarrow} V$ de $\partial V$ dans $V$ (avec $\alpha(1, x) = x$ ), et un représentant de $\varphi$ tel que $\varphi : [0, 1] \longrightarrow \mathrm{Diff}(\partial V)$ , prolongé pour $x < 0$ par $\mathrm{Id}$ , soit $C^{\infty}$ . Soit, dans le collier, $\beta(x) = \mathrm{pr}_1 \alpha^{-1}(x)$ ( $\beta$ est sans point critique, à valeurs dans $[0, 1]$ , et vaut 1 sur $\partial V$ ) .

($\beta$) On prend $T = T_o$ en dehors du collier et, dans le collier,

\SGAvar{\begin{equation*}
T(x) = T_o(\varphi(\beta(x))^{-1}(x))
\end{equation*}
3.1.13. Voici une variante de 3.1.11(b) pour calculer la classe d'homotopie de chemins $\varphi$ . Sur chaque rayon $\arg(z) = \theta$ , on choisit une trivialisation du fibré $\partial B_{t,d}$ , variant continument avec $\theta$ . On identifie $(\partial V)_o = f^{-1}(0) \cap \partial B_{t,d}$ à $\partial V$ grâce à la trivialisation au-dessus du rayon passant par $y$ . Soit $\Psi_{\theta} : V \longrightarrow f^{-1}(e^{i\theta} y) \cap B_{t,d}$ déduit de l'intégration de $\nabla$ .Transportant $\partial \Psi_{\theta}$ le long des rayons, on obtient un difféomorphisme de $(\partial V)_o$ , donc de $\partial V$ , sur lui-même. C'est $\varphi(\theta)$ .


}{\begin{equation*}
T(x) = T_o(\varphi(\beta(x))^{-1}(x))
\end{equation*}
3.1.13. Voici une variante de 3.1.11(b) pour calculer la classe d'homotopie de chemins $\varphi$ . Sur chaque rayon $\arg(z) = \theta$ , on choisit une trivialisation du fibré $\partial B_{t,d}$ , variant continûment avec $\theta$ . On identifie $(\partial V)_o = f^{-1}(0) \cap \partial B_{t,d}$ à $\partial V$ grâce à la trivialisation au-dessus du rayon passant par $y$ . Soit $\Psi_{\theta} : V \longrightarrow f^{-1}(e^{i\theta} y) \cap B_{t,d}$ déduit de l'intégration de $\nabla$ .Transportant $\partial \Psi_{\theta}$ le long des rayons, on obtient un difféomorphisme de $(\partial V)_o$ , donc de $\partial V$ , sur lui-même. C'est $\varphi(\theta)$ .


}{F7II-XIV-034}3.2. \underline{La formule de Picard-Lefschetz (cas transcendant).}

3.2.1. Nous nous proposons d'expliciter les résultats de 3.1 pour l'application $f = \sum\limits_{1}^{n+1} z_i^2$ de $\mathbb{C}^{n+1}$ dans $\mathbb{C}$ . On prend pour $\delta$ la fonction $\sum |z_i|^2$ , et une boîte $B = B_{a,b}$

\begin{equation*}
| \sum z_i^{\,2} | \leq a^2
\end{equation*}
\begin{equation*}
\sum |z_i|^2 \leq b^2
\end{equation*}
avec $0 < a < b$ (ce n'est que pour simplifier les notations qu'on prend une boîte
\SGApage{II}{147}
fermée plutôt qu'une boîte 3.1.2).

3.2.2. La variété des cycles évanescents $V$ admet l'équation

\begin{equation}\tag{3.2.2.1}
\begin{array}{rcl}
\Sigma z_i^2 & = & a^2 \\
\Sigma |z_i|^2 & \leq & b^2
\end{array}
\end{equation}
Décomposons $z = (z_i)$ en parties réelle et imaginaire: $z = x + iy$ . Les équations précédentes se récrivent

\SGAvar{\begin{equation}\tag{3.2.2.2}
\begin{array}{rcl}
x^2 - y^2 & = & a^2 \\
x.y & = & 0 \qquad\qquad (\mathrm{i.e.}\ \ x \perp y\ \ \mathrm{dans}\ \ \mathbb{R}^n\ ) \\
x^2 + y^2 & \leq & b^2
\end{array}
\end{equation}
Prenons sur $V$ le nouveau système de coordonnées $X = x/|x|$ et $Y = y(\frac{b^2-a^2}{2})^{-1/2}$ . On a $x = X. (a^2 + y^2)^{1/2}$ et (3.2.2.2) se récrit

}{\begin{equation}\tag{3.2.2.2}
\begin{array}{rcl}
x^2 - y^2 & = & a^2 \\
x.y & = & 0 \qquad\qquad (\mathrm{i.e.}\ \ x \perp y\ \ \mathrm{dans}\ \ \mathbb{R}^{n+1}\ ) \\
x^2 + y^2 & \leq & b^2
\end{array}
\end{equation}
Prenons sur $V$ le nouveau système de coordonnées $X = x/|x|$ et $Y = y(\frac{b^2-a^2}{2})^{-1/2}$ . On a $x = X. (a^2 + y^2)^{1/2}$ et (3.2.2.2) se récrit

}{F7II-XIV-035}\begin{equation}\tag{3.2.2.3}
\begin{array}{rcl}
X^2 & = & 1 \\
XY & = & 0 \\
Y^2 & \leq & 1
\end{array}
\end{equation}
Ceci identifie $V$ au fibré en boules tangent à la sphère $S^n$ .

3.2.3. L'application $\theta \longmapsto \exp(2\pi i\theta)a$ identifie $\mathbb{R}$ au revêtement universel de la circonférence $S_a$ de rayon $a$ et centre $0$ de $\mathbb{C}$ . Le fibré sur $\mathbb{R}$ image réciproque de $B$ est isomorphe à $R \times V$ . Nous le trivialiserons à l'aide de l'isomorphisme $v \longmapsto \exp (\pi i\theta)v$ de $V$ avec la fibre de $B$ en $\exp(2\pi i\theta)a$ . Soit $T_o$ de cette trivialisation comme en 3.1.11. En coordonnées (3.2.2.3), on a $T_o(X,Y) = (-X, -Y)$

3.2.4. Le fibré sur le disque $D_a$ de rayon $a$ , bord du ``fibré'' $B$ , a pour équation

\begin{equation*}
\Sigma |z_i|^2 = b^2 \qquad , \qquad f(z) \in D_a \quad , 
\end{equation*}

\SGApage{II}{148}

 la projection sur le disque étant $f$ . Nous nous proposons de décrire une trivialisation de l'image réciproque de ce fibré au-dessus de chacun des rayons $e^{2\pi i \theta} \mathbb{R}^+ \cap D_a$ .

\SGAvar{Au dessus du rayon $\mathbb{R}^+ \cap D_a$ , on prend pour trivialisation l'isomorphisme avec $(\mathbb{R}^+ \cap D_a) \times \partial V$ donné en coordonnées (3.2.2.3) sur $\partial V$ par

}{Au-dessus du rayon $\mathbb{R}^+ \cap D_a$ , on prend pour trivialisation l'isomorphisme avec $(\mathbb{R}^+ \cap D_a) \times \partial V$ donné en coordonnées (3.2.2.3) sur $\partial V$ par

}{F7II-XIV-036}\begin{equation*}
z = x+iy \longmapsto (f(z), X, Y) = (f(z), x/|x|, y/|y|) .
\end{equation*}
Au-dessus de $(e^{2\pi i \theta} \mathbb{R}^+) \cap D_a$ , on le trivialise par la trivialisation image de la précédente par multiplication par $e^{\pi i \theta}$ . Appliquons 3.1.13 et prenons sur $\partial V$ les coordonnées $(X, Y)$

\begin{equation}\tag{3.2.4.1}
X^2 = Y^2 = 1, XY = 0 .
\end{equation}
On trouve que, si $\varphi(\theta)$ transforme $(X, Y)$ en $(X_\theta, Y_\theta)$ , on a

\SGAvar{\begin{equation}\tag{3.2.4.2}
X_\theta + i Y_\theta = e^{\pi i \theta}(X + i Y) \qquad\qquad (0 \leq \theta \leq 1)
\end{equation}
3.2.5. En résumé, la variété des cycles évanescents $V$ est le fibré en boules tangent à $S^n$ . La transformation de monodromie est représentée par $(T_o, \varphi)$ , où $T_o$ est la différentielle de l'application $x \longmapsto$ point antipodique, et où $\varphi$ est le chemin suivant de $\mathrm{Id}$ à $\partial T_o$ : par $\varphi(\theta)$ , on déplace par transport parallèle chaque vecteur tangent unitaire le long de la geodésique qu'il engendre d'une longueur $+ \pi \theta$ (le $+$ signifie: ``dans le sens vers lequel pointe le vecteur''). Pour $\theta = 1$ , on arrive au vecteur antipodique.

Par 3.1.13, on obtient la description suivante de $T$ : soit $\varphi$ une application $C^\infty$ $[0,1] \to [0,1]$ , valant $0$ près de $0$ et $1$ près de $1$ . $T$ consiste à déplacer $u$ , le long de la geodésique qu'il engendre, d'une longueur $- \pi(1-\varphi(|u|)$ (près de $|u| = 1$ , c'est l'identité, près de $|u| = 0$ , c'est l'application antipodique).

3.2.7. Soient $W$ une variété différentiable de dimension $x$ et $p : T(W) \to W$ son fibré tangent (vu comme une variété).


}{\begin{equation}\tag{3.2.4.2}
X_\theta + i Y_\theta = e^{\pi i \theta}(X + i Y) \qquad\qquad (0 \leq \theta \leq 1)
\end{equation}
3.2.5. En résumé, la variété des cycles évanescents $V$ est le fibré en boules tangent à $S^n$ . La transformation de monodromie est représentée par $(T_o, \varphi)$ , où $T_o$ est la différentielle de l'application $x \longmapsto$ point antipodique, et où $\varphi$ est le chemin suivant de $\mathrm{Id}$ à $\partial T_o$ : par $\varphi(\theta)$ , on déplace par transport parallèle chaque vecteur tangent unitaire le long de la géodésique qu'il engendre d'une longueur $+ \pi \theta$ (le $+$ signifie: ``dans le sens vers lequel pointe le vecteur''). Pour $\theta = 1$ , on arrive au vecteur antipodique.

Par 3.1.12, on obtient la description suivante de $T$ : soit $\varphi$ une application $C^\infty$ $[0,1] \to [0,1]$ , valant $0$ près de $0$ et $1$ près de $1$ . $T$ consiste à déplacer $u$ , le long de la géodésique qu'il engendre, d'une longueur $- \pi(1-\varphi(|u|))$ (près de $|u| = 1$ , c'est l'identité, près de $|u| = 0$ , c'est l'application antipodique).

3.2.7. Soient $W$ une variété différentiable de dimension $d$ et $p : T(W) \to W$ son fibré tangent (vu comme une variété).


}{F7II-XIV-037}\SGApage{II}{149}

\SGAvar{Si $(x_1,\ldots,x_n)$ est un système de coordonnées locales sur $W$ , $(x_1,\ldots,x_n, dx_1,\ldots,dx_n)$ est un système de coordonnées locales sur $T(W)$ . Soit $\varepsilon'$ l'orientation de $T(W)$ définie par ce système de coordonnées, elle ne dépend pas du choix des $x_i$ , donc est globalement définie. Supposons $W$ orientable. La sous-variété $W \subset T(W)$ (section 0), munie d'une orientation (resp. d'une orientation de son fibré normal), définit un élément $w$ de $H_d(T(W))$ (resp. $H^d_c(T(W))$ ). On sait que la self-intersection de $W \subset T(W)$ est la caractéristique d'Euler-Poincaré

}{Si $(x_1,\ldots,x_d)$ est un système de coordonnées locales sur $W$ , $(x_1,\ldots,x_d, dx_1,\ldots,dx_d)$ est un système de coordonnées locales sur $T(W)$ . Soit $\varepsilon'$ l'orientation de $T(W)$ définie par ce système de coordonnées, elle ne dépend pas du choix des $x_i$ , donc est globalement définie. Supposons $W$ orientable. La sous-variété $W \subset T(W)$ (section 0), munie d'une orientation (resp. d'une orientation de son fibré normal), définit un élément $w$ de $H_d(T(W))$ (resp. $H^d_c(T(W))$ ). On sait que la self-intersection de $W \subset T(W)$ est la caractéristique d'Euler-Poincaré

}{F7II-XIV-038}\begin{equation*}
\int w \wedge w = \chi(W)
\end{equation*}
(pour l'orientation $\varepsilon'$ ).

Dans le cas de la variété des cycles évanescents $V$ , la structure complexe de $V$ définit une orientation $\varepsilon$ . Elle est liée à l'orientation $\varepsilon'$ du fibré tangent à $S^n$ par

\begin{equation*}
\varepsilon = (-1)^{\frac{n(n-1)}{2}} \varepsilon' .
\end{equation*}
3.2.8. Le \underline{cycle évanescent} $\delta$ est bien défini au signe près. Pour $n \neq 0$ , c'est l'élément de $H^n_c(V^o)$ défini par la sous-variété $S^n \subset V^o$ . Pour $n = 0$ , $V$ est réduit à deux points et c'est la différence de ces deux points. Pour l'orientation complexe de $V$ , on a par 3.2.7

\begin{equation}\tag{3.2.8.1}
(\delta, \delta) = (-1)^{\frac{n(n-1)}{2}} (1 + (-1)^n) .
\end{equation}
3.2.9. Posons $\widetilde{H}^i_c(V^o) = \mathrm{Ker}\, (H^i_c(V^o) \xrightarrow{\ \mathrm{Tr}\ } \mathbb{Z})$ (c'est $H^i_c$ pour $i \neq 2n$) . On a alors

\begin{equation*}
\begin{aligned} \widetilde{H}^i_c(V^o) &= 0 \qquad \text{pour } i \neq n \\ \widetilde{H}^n_c(V^o) &= \mathbb{Z} \ , \quad \text{engendré par } \delta \ . \end{aligned}
\end{equation*}
Dualement, on a

\begin{equation*}
\begin{aligned} \widetilde{H}^i(V) &= 0 \qquad \text{pour } i \neq n \\ \widetilde{H}^n(V) &\xrightarrow[\sim]{(\ , \delta)} \mathbb{Z} \ . \end{aligned}
\end{equation*}

\SGApage{II}{150}

Pour $n \neq 0$ , l'image de $\delta \in H^n_c(V^o)$ dans $H^n(V)$ est nulle pour $n$ impair, et divisible par deux (et pas plus) pour $n$ pair.


3.2.10. Calculons la variation

\SGAvar{\begin{equation*}
\mathrm{Var} : \widetilde{H}^n(V) \longrightarrow H^n_c(V^o) \quad .
\end{equation*}
Nous utiliserons d'abord l'orientation $\epsilon'$ , et supposerons que $n \neq 0$ . Soit $\alpha$ une orientation du fibré tangent à $S^n$ , et soit $\delta \in H^n_c(V^o)$ défini par $S^n \subset V^o$ et l'orientation $\alpha$ de son fibré normal. Le cocycle $\delta' \in H^n(V)$ tel que $(\delta' \; \delta) = 1$ est représenté par une fibre de la fibration de $V^o$ sur $S^n$ ; le fibré normal à cette fibré est le fibré constant, de valeur l'espace tangent en un point de $S^n$ , et est muni de l'orientation $\alpha$ .

}{\begin{equation*}
\mathrm{Var} : \widetilde{H}^n(V) \longrightarrow H^n_c(V^o) \quad .
\end{equation*}
Nous utiliserons d'abord l'orientation $\epsilon'$ , et supposerons que $n \neq 0$ . Soit $\alpha$ une orientation du fibré tangent à $S^n$ , et soit $\delta \in H^n_c(V^o)$ défini par $S^n \subset V^o$ et l'orientation $\alpha$ de son fibré normal. Le cocycle $\delta' \in H^n(V)$ tel que $(\delta' \; \delta) = 1$ est représenté par une fibre de la fibration de $V^o$ sur $S^n$ ; le fibré normal à cette fibre est le fibré constant, de valeur l'espace tangent en un point de $S^n$ , et est muni de l'orientation $\alpha$ .

}{F7II-XIV-039}On notera que cette assertion ne dépend pas de la convention de signe choisie pour définir l'application

(sous-variété + orientation du fibré normal) $\longrightarrow$ cohomologie

(changer de convention change $\delta$ en $-\delta$ et $\delta'$ en $-\delta'$ ). Elle est valable si Tr : $H^{2n}_c(V, \mathbb{R}) \longrightarrow \mathbb{R}$ est défini comme suit:

\SGAvar{($\alpha$) $\Omega^*_V$ étant une résolution molle de $\underline{\mathbb{R}}$ , $\Gamma_c(\Omega^*_V)$ représente $R\Gamma(V, \mathbb{R})$ , d'où un isomorphisme $H^i(\Gamma_c(V, \; \Omega^*_V)) = H^i_c(V, \mathbb{R})$ (on n'utilise pas le cobord itéré);

}{($\alpha$) $\Omega^*_V$ étant une résolution molle de $\underline{\mathbb{R}}$ , $\Gamma_c(\Omega^*_V)$ représente $R\Gamma_c(V, \mathbb{R})$ , d'où un isomorphisme $H^i(\Gamma_c(V, \; \Omega^*_V)) = H^i_c(V, \mathbb{R})$ (on n'utilise pas le cobord itéré);

}{F7II-XIV-040}($\beta$) Tr est l'intégration.

\SGAvar{$\mathrm{T}(\delta')$ coincide avec $\delta'$ dans un voisinage de $\partial V$ , de sorte que $\mathrm{Var}(\delta') = \mathrm{T}(\delta') - \delta'$ est à support compact dans $V^o$ . On a

}{$\mathrm{T}(\delta')$ coïncide avec $\delta'$ dans un voisinage de $\partial V$ , de sorte que $\mathrm{Var}(\delta') = \mathrm{T}(\delta') - \delta'$ est à support compact dans $V^o$ . On a

}{F7II-XIV-041}\begin{equation*}
\mathrm{T}(\delta') - \delta' \sim -\delta \quad , \qquad \mbox{d'o\`u}
\end{equation*}
\begin{equation*}
\mathrm{Var}(\delta') = -(\delta' \delta) . \delta \quad .
\end{equation*}
Voici le dessin pour $n = 1$ . On dessine le fibré tangent à $S^1$ en faisant pointer vers l'extérieur les vecteurs pointant dans le sens trigonométrique choisi pour orienter $S^1$


\SGApage{II}{151}

\begin{center}
\begin{tikzcd}[row sep=large, column sep=large]
{} \& {} \arrow[l] \& {} \& {} \& \epsilon \& {} \\
{} \& {} \arrow[l] \& \delta' \& {} \& {} \& {} \arrow[u, "2"'] \arrow[l, "1"] \\
{} \& {} \& {} \& {} \& {} \& {} \\
{} \& {} \arrow[ul] \& {} \& {} \& {} \& {} \\
{} \& \delta \& {} \& {} \& {} \& {} \\
{} \& {} \arrow[dl] \& {} \& {} \& {} \arrow[dr] \& {} \\
{} \& {} \& {} \& {} \& {} \&
{} \end{tikzcd}
\end{center}























$\delta$, $\delta'$ et $\epsilon$

\begin{center}
\begin{tikzcd}[row sep=large, column sep=large]
{} \& {} \arrow[dl] \\
{} \&
{} \end{tikzcd}
\end{center}














$\mathrm{Var}(\delta')$

Pour l'orientation $\epsilon$ , on a donc

\begin{equation*}
\mathrm{Var}(x) = -\ \epsilon/\epsilon'\ (x\delta)\delta \quad .
\end{equation*}
\underline{Théorème} 3.2.11. \underline{Le groupe} $\widetilde{H}^n_c(V^o)$ \underline{est isomorphe à} $\mathbb{Z}$ , \underline{engendré par le cycle évanescent} $\delta$ . \underline{Le groupe} $\widetilde{H}^n(V)$ \underline{est aussi isomorphe à} $\mathbb{Z}$ , \underline{et dual du groupe précédent. Les autres} $\widetilde{H}^i_c(V^o)$ \underline{et} $\widetilde{H}^i(V)$ \underline{sont nuls.}

\underline{La variation est donnée par}

\begin{equation*}
\mathrm{Var}(x) = -\ (-1)^{\frac{n(n-1)}{2}}\ (x\delta)\delta \quad , \quad \underline{\mbox{et}}
\end{equation*}
\begin{equation*}
(\delta, \; \delta) = (-1)^{\frac{n(n-1)}{2}}\ (1+(-1)^n)
\end{equation*}
Selon la congruence de $n \bmod 4$ , on a la table suivante

\begin{equation*}
\begin{array}{lcccc}
n \;\bmod\; 4 & 0 & 1 & 2 & 3 \\[2pt]
(\delta, \; \delta) & 2 & 0 & -2 & 0 \\[2pt]
\mathrm{Var}(x) = \pm(x\delta)\delta & - & - & + & + \\[2pt]
\mathrm{T}(\delta) = (-1)^{n+1}\delta & -\delta & \delta & -\delta & \delta
\end{array}
\end{equation*}
A titre de contrôle, on note qu'on a bien

\begin{equation*}
\mathrm{T}(\delta) = \delta + \mathrm{Var}(\delta) \quad .
\end{equation*}

\SGApage{II}{152}


\SGAsectiontoc{4. Cohomologie de De Rham}{4}
4. \underline{Cohomologie de De Rham}

4.1.  Soient $D$ un disque analytique complexe, $0 \in D$ , $D^* = D - \{0\}$ et $\tilde{D}^*$ un revêtement universel de $D^*$ . Soient $X$ un espace analytique, $f : X \to D$ un morphisme et le diagramme (1.2.7.1)

\begin{equation}\tag{4.1.1}
\begin{tikzcd}[column sep=large, row sep=large]
X_o \arrow[r, hook, "\bar{i}"] \arrow[d, equal] \& \bar{X} \arrow[d, "p"] \& \bar{X}^* \arrow[l, hook, "\bar{j}"'] \arrow[d, "p'"] \\
X_o \arrow[r, hook, "i"] \& X \& X^* \arrow[l, hook, "j"']
\end{tikzcd}
\end{equation}













\SGAvar{On suppose que $X^* = X - X_o$ est non singulier (= lisse). Son revêtement $\bar{X}^*$ est donc lisse, on note $\Omega^*_{\bar{X}^*}$ le complexe de De Rham des faisceaux de formes différentielles holomorphes sur $\Omega^*_{\bar{X}^*}$ . C'est une résolution du faisceau constant $\mathbb{C}$ (lemme de Poincaré holomorphe), d'où une flèche

}{On suppose que $X^* = X - X_o$ est non singulier (= lisse). Son revêtement $\bar{X}^*$ est donc lisse, on note $\Omega^*_{\bar{X}^*}$ le complexe de De Rham des faisceaux de formes différentielles holomorphes sur $\bar{X}^*$ . C'est une résolution du faisceau constant $\mathbb{C}$ (lemme de Poincaré holomorphe), d'où une flèche

}{F7II-XIV-042}\begin{equation}\tag{4.1.2}
\bar{j}_* \, \Omega^*_{\bar{X}^*} \longrightarrow R\bar{j}_* \, \mathbb{C}
\end{equation}
dans $D^+(\bar{X}, \mathbb{C})$ .

\underline{Proposition} 4.2.  \underline{Le morphisme} (4.1.2) \underline{est un isomorphisme.}

\SGAvar{Il est clair que 4.1.1 est un quasi-isomorphisme en tout point de $\bar{X}^* \subset \bar{X}$ . Il suffit donc de montrer que tout point $x \in X_o$ a un système fondamental de voisinages $U$ dans $\bar{X}$ tels que

\begin{equation*}
H^q(\bar{j}^{-1}(U), \Omega^p) \;=\; 0 \quad \text{pour} \quad q > 0 \ .
\end{equation*}
Pour $V$ parcourant un système rondamental de $x$ dans $X$ , les $p^{-1}(V)$ forment un système fondamental de voisinages de $x$ dans $\bar{X}$ . En particulier, les $p^{-1}(V)$ pour $V$ un voisinage de Stein de $x$ dans $X$ forment un système fondamental de voisinages de $X$ . Pour $V$ de Stein, $V-X_o$ est encore de Stein (car $X_o \cap V$ est défini par une équation dans $V$ ), ainsi que son revêtement $p'^{-1}(V-X_o) = \bar{j}^{-1}p^{-1}(V)$,
\SGApage{II}{153}
Pour tout faisceau cohérent $\mathfrak{F}$ sur $\bar{X}^*$ , on a donc

}{Il est clair que 4.1.2 est un quasi-isomorphisme en tout point de $\bar{X}^* \subset \bar{X}$ . Il suffit donc de montrer que tout point $x \in X_o$ a un système fondamental de voisinages $U$ dans $\bar{X}$ tels que

\begin{equation*}
H^q(\bar{j}^{-1}(U), \Omega^p) \;=\; 0 \quad \text{pour} \quad q > 0 \ .
\end{equation*}
Pour $V$ parcourant un système fondamental de $x$ dans $X$ , les $p^{-1}(V)$ forment un système fondamental de voisinages de $x$ dans $\bar{X}$ . En particulier, les $p^{-1}(V)$ pour $V$ un voisinage de Stein de $x$ dans $X$ forment un système fondamental de voisinages de $x$ . Pour $V$ de Stein, $V-X_o$ est encore de Stein (car $X_o \cap V$ est défini par une équation dans $V$ ), ainsi que son revêtement $p'^{-1}(V-X_o) = \bar{j}^{-1}p^{-1}(V)$,
\SGApage{II}{153}
Pour tout faisceau cohérent $\mathfrak{F}$ sur $\bar{X}^*$ , on a donc

}{F7II-XIV-043}\begin{equation*}
H^q(\bar{j}^{-1} \, p^{-1}(V), \mathfrak{F}) \;=\; 0 \quad \text{pour} \quad q > 0 \ ,
\end{equation*}
ce qui achève la démonstration.

\underline{Corollaire} 4.3.  \underline{L'isomorphisme} (4.1.2) \underline{induit un isomorphisme}

\begin{equation}\tag{4.3.1}
\bar{i}^* \, \bar{j}_* \, \Omega^*_{\bar{X}^*} \xrightarrow{\;\sim\;} R\Psi_\eta(\mathbb{C}) \ .
\end{equation}
\underline{Variante} 4.4.  \underline{Soit} $V$ \underline{un faisceau localement constant de vectoriels complexes de dimension finie sur} $X^*$ (\underline{un système local complexe}). \underline{Le quasi-isomorphisme} $V \longrightarrow \Omega^*_{X^*}(V)$ (\underline{lemme de Poincaré) induit un quasi-isomorphisme}

\begin{equation}\tag{4.4.1}
\bar{i}^* \, \bar{j}_* \, \Omega^*_{\bar{X}^*} \, (p'^* V) \xrightarrow{\;\sim\;} R\Psi_\eta(V) \ .
\end{equation}
Nous nous proposons de montrer que 4.3.1 reste vrai lorsqu'on remplace $\bar{i}^* \, \bar{j}_* \, \Omega^*_{\bar{X}^*}$ par un sous-complexe défini par des conditions de croissance et de monodromie.

\SGAvar{4.5.  Soient $f : X^* \to D^*$ un espace topologique sur $D^*$ , $F$ un faisceau de $f^{\circ} \mathcal{O}_{D^*}$-modules sur $X^*$ et $\bar{F}$ son image réciproque sur $\bar{X}^* = X^* \times_{D^*} D^*$ .

}{4.5.  Soient $f : X^* \to D^*$ un espace topologique sur $D^*$ , $F$ un faisceau de $f^{\circ} \mathcal{O}_{D^*}$-modules sur $X^*$ et $\bar{F}$ son image réciproque sur $\bar{X}^* = X^* \times_{D^*} \tilde{D}^*$ .

}{F7II-XIV-044}Soit $z : D \hookrightarrow \mathbb{C}$ un plongement holomorphe de $D$ dans $\mathbb{C}$ , avec $z(0) = 0$ . Par la suite, nous aurons toujours le droit de rétrécir $D$ ; pour un disque assez petit, l'existence de $z$ est une trivialité. On suppose choisie sur $\tilde{D}^*$ une fonction $\log z$ tel que $\exp(\log z) = z$ , et on pose

\SGAvar{\begin{equation*}
z^\alpha \;=\; \exp(\alpha \log z) \qquad\qquad (\alpha \in \mathbb{C}) \ .
\end{equation*}
Le groupe $\pi_1(D^*)$ agit (par transport de structure) sur $H^o(\bar{X}^*, \bar{F})$ . Un élément $f \in H^o(\bar{X}^*, \bar{F})$ est de \underline{détermination finie} si l'ensemble de ses transformés sous $\pi_1(D^*)$ engendre un espace vectoriel $V(F)$ de dimension finie sur $\mathbb{C}$ . Soit $T$ le générateur positif de $\pi_1(D^*)$ . Pour que $f$ soit de détermination finie,
\SGApage{II}{154}
il faut et il suffit qu'il existe un polynôme $P \in \mathbb{C}[X]$ , non nul, tel que $P(T)(f) = 0$ . Le polynôme $P$ de plus petit degré tel que $P(T)(f) = 0$ est alors le polynôme minimal, égal au polynôme caractéristique, de $T$ agissant sur $V(f)$ . On dit que $f$ est de \underline{détermination finie unipotente} s'il existe $A$ tel que

}{\begin{equation*}
z^\alpha \;=\; \exp(\alpha \log z) \qquad\qquad (\alpha \in \mathbb{C}) \ .
\end{equation*}
Le groupe $\pi_1(D^*)$ agit (par transport de structure) sur $H^o(\bar{X}^*, \bar{F})$ . Un élément $f \in H^o(\bar{X}^*, \bar{F})$ est de \underline{détermination finie} si l'ensemble de ses transformés sous $\pi_1(D^*)$ engendre un espace vectoriel $V(f)$ de dimension finie sur $\mathbb{C}$ . Soit $T$ le générateur positif de $\pi_1(D^*)$ . Pour que $f$ soit de détermination finie,
\SGApage{II}{154}
il faut et il suffit qu'il existe un polynôme $P \in \mathbb{C}[X]$ , non nul, tel que $P(T)(f) = 0$ . Le polynôme $P$ de plus petit degré tel que $P(T)(f) = 0$ est alors le polynôme minimal, égal au polynôme caractéristique, de $T$ agissant sur $V(f)$ . On dit que $f$ est de \underline{détermination finie unipotente} s'il existe $A$ tel que

}{F7II-XIV-045}\begin{equation*}
(T-1)^A \; f = 0 \; .
\end{equation*}
Sur l'espace de ces sections, on définit $N = \frac{\log T}{2\pi\ i}$ par la somme finie

\begin{equation*}
N(f) = \frac{-1}{2\pi i} \sum_{n > 0} \frac{(1-T)^n}{n} \; f \; .
\end{equation*}
Pour $f$ (resp. $g$ ) une section de détermination finie unipotente de $H^0(D^*, \mathcal{O})$ (resp. $H^0(\bar{X}^*, \bar{F})$ ), on a

\SGAvar{\begin{equation*}
N(f\ g) = N(f).g + f.\ N(g)
\end{equation*}
(le logarithme d'un autmorphisme est une dérivation).

}{\begin{equation*}
N(f\ g) = N(f).g + f.\ N(g)
\end{equation*}
(le logarithme d'un automorphisme est une dérivation).

}{F7II-XIV-046}On dit que $f$ est de \underline{détermination finie quasi-unipotente} s'il existe $A$ et $B > 0$ avec

\SGAvar{\begin{equation*}
(T^B - 1)^A \; f = 0 \; .
\end{equation*}
\underline{Examples} 4.6. (pour $X = D^*$ , $F = \mathcal{O}$)

\begin{equation*}
(T - \exp(2\pi i\alpha))^{k+1} \; (z^{\alpha}(\log z)^k) = 0
\end{equation*}
\begin{equation*}
N(\log z) = 1
\end{equation*}
\begin{equation*}
N\frac{(\log z)^k}{k!}) = \frac{(\log z)^{k-1}}{(k-1)!} \qquad (0 \ \mbox{pour} \ k = 0) \; .
\end{equation*}
\underline{Proposition} 4.7. \underline{Pour qu'une section} $f \in H^0(\bar{X}^*, \bar{F})$ \underline{soit de détermination finie, il faut et il suffit qu'elle puisse s'écrire}

}{\begin{equation*}
(T^B - 1)^A \; f = 0 \; .
\end{equation*}
\underline{Exemples} 4.6. (pour $X = D^*$ , $F = \mathcal{O}$)

\begin{equation*}
(T - \exp(2\pi i\alpha))^{k+1} \; (z^{\alpha}(\log z)^k) = 0
\end{equation*}
\begin{equation*}
N(\log z) = 1
\end{equation*}
\begin{equation*}
N\left(\frac{(\log z)^k}{k!}\right) = \frac{(\log z)^{k-1}}{(k-1)!} \qquad (0 \ \mbox{pour} \ k = 0) \; .
\end{equation*}
\underline{Proposition} 4.7. \underline{Pour qu'une section} $f \in H^0(\bar{X}^*, \bar{F})$ \underline{soit de détermination finie, il faut et il suffit qu'elle puisse s'écrire}

}{F7II-XIV-047}\begin{equation}\tag{4.7.1}
f = \sum_{\alpha, k} \; f_{\alpha, k} \; z^{\alpha}\frac{(\log z)^k}{k!} \qquad ,
\end{equation}
\underline{avec} $0 \leq \Re(\alpha) < 1$ , $k \in \mathbb{N}$ , \underline{et les} $f_{\alpha, k} \in H^0(X^*, F)$ \underline{presque tous nuls. Pour} $f$ \underline{de détermination finie, la décomposition} (4.7.1) \underline{est unique.}


\SGApage{II}{155}

D'après 4.6, toute section (4.7.1) est de détermination finie. Soit réciproquement $f$ de détermination finie: $P(T) (f) = 0$ . Posons

\begin{equation*}
P(X) = \prod_{\beta \in \mathbb{C}} \; (X-\beta)^{n'(\beta)} \; = \prod_{0 \leq \Re(\alpha) < 1} \; (X-\exp(2\pi i\alpha))^{n(\alpha)}
\end{equation*}
avec $n(\alpha) = n'(\exp(2\pi i\alpha))$ presque toujours nul. Décomposant le $\mathbb{C}[X]$-module $V(f)$ en composante isotypique, on trouve que $f$ admet une et une seule décomposition $f = \sum f'_{\alpha}$ , avec

\begin{equation*}
(T-\exp(2\pi i\alpha))^{n(\alpha)} \; f'_{\alpha} = 0 \; .
\end{equation*}
Pour $Q_{\alpha} \in \mathbb{C}[X]$ un polynôme congru à 1 mod$(X- \exp(2\pi i\alpha))^{n(\alpha)}$ et à 0 mod $(X - \exp(2\pi i\alpha'))^{n(\alpha')}$ pour $\alpha' \neq \alpha$ , on a $f'_{\alpha} = Q_{\alpha}(T) (f)$ . Posons

\begin{equation}\tag{4.7.2}
f_{\alpha} = z^{-\alpha} \; f'_{\alpha} \; = \; z^{-\alpha} \; Q_{\alpha}(T) (f) \; .
\end{equation}
Puisque $T(z^{-\alpha} f'_{\alpha}) = \exp(-2\pi i\alpha) \; z^{-\alpha} \; T(f'_{\alpha})$ , on a

\begin{equation*}
(T-1)^{n(\alpha)} \; f_{\alpha} \; = 0 \; .
\end{equation*}
Si $f_{\alpha} = \sum f_{\alpha, k} \; \frac{(\log z)^k}{k!}$ , avec $f_{\alpha, k} \in H^0(X^*, F)$ , on a alors par 4.6 et l'identité $\sum (-1)^j {k \choose j} = \delta_{0, k}$

\begin{equation}\tag{4.7.3}
f_{\alpha, k} \; = \sum_{j \geq 0} \; (-1)^j \; \frac{(\log z)^j}{j!} \; N^{j+k} \; f_{\alpha} \quad ,
\end{equation}
d'où l'unicité. Si réciproquement $f_{\alpha}$ est une section de détermination finie unipotente, $f_{\alpha, k}$ défini par (4.7.3) est nul pour $k$ grand (car $N^k f_{\alpha} = 0$ pour $k$ grand) et est invariant sous $T$ , donc s'identifie à une section de $F$ sur $X^*$ . Montrons, ce qui revient au même, que $N f_{\alpha, k} = 0$ : on applique à $N^k f_{\alpha}$ l'identité

\begin{equation*}
N \sum_{j \geq 0} \; (-1)^j \; \frac{(\log z)^j}{j!} \; N^j \; = \; \sum_{j > 0} \; (-1)^j \; \frac{(\log z)^{j-1}}{(j-1)!} \; N^j \; + \sum_{j \geq 0} \; (-1)^j \; \frac{(\log z)^j}{j!} \; N^{j+1} = 0 \; .
\end{equation*}

\SGApage{II}{156}

On a enfin

\begin{equation*}
\sum_{k \geq 0} f_{\alpha,k} \; \frac{(\log z)^k}{k!} \; = \sum_{j,k \geq 0} \; (-1)^j \; \frac{(\log z)^{j+k}}{j! \; k!} \; N^{j+k} \; f_\alpha
\end{equation*}
\begin{equation*}
= \sum_{n \geq 0} \; (\sum \, (-1)^j \binom{n}{j}) \; \frac{(\log z)^n}{n!} \; N^n \; f_\alpha \; = \; f_\alpha \qquad ,
\end{equation*}
ce qui prouve 4.7.

\underline{Remarque} 4.8.1. La condition $\Re(\alpha) \in [0, 1[$ de 4.7 peut se remplacer par $\alpha \in C$ , pour $C$ n'importe quel domaine fondamental pour $\mathbb{C} \to \mathbb{C}/\mathbb{Z}$ .

\underline{Remarque} 4.8.2. Soit $f$ de détermination finie et sa décomposition (4.7.1). Pour $0 \leq \Re(\alpha) < 1$ , et $\alpha$ tel qu'un $f_{\alpha,k} \neq 0$ , soit $m(\alpha) = 1 + \sup \{k \mid f_{\alpha,k} \neq 0\}$ . Si tous les $f_{\alpha,k}$ sont nuls, posons $m(\alpha) = 0$ . Le polynôme caractéristique de $T$ agissant sur $V(f)$ est

\begin{equation*}
Pc(T, \; V(f)) = \prod \, (X - \exp(2\pi i \alpha))^{m(\alpha)}
\end{equation*}
En particulier, $f$ est de détermination finie quasi-unipotente si et seulement si dans (4.7.1), on peut restreindre $\alpha$ à être dans $\mathbb{Q}$ .

4.9. Revenons aux hypothèses plus restrictives de 4.1 et soit $F$ un faisceau cohérent sur $X$ dont la restriction à $X^*$ soit localement libre.

\SGAvar{\underline{Proposition} 4.10. \underline{Soit} $f$ \underline{une section de détermination finie de l'image réciproque} $\overline{F}$ \underline{de} $F$ \underline{sur} $\overline{X}^*$ . \underline{Pour que} $f$ \underline{soit à croissance modérée le long de} $X_o$ , \underline{il faut et il suffit que les} $f_{\alpha,k}$ \underline{de la décomposition} (4.7.1) \underline{soient méromorphe le long de} $X_o$ .

}{\underline{Proposition} 4.10. \underline{Soit} $f$ \underline{une section de détermination finie de l'image réciproque} $\overline{F}$ \underline{de} $F$ \underline{sur} $\overline{X}^*$ . \underline{Pour que} $f$ \underline{soit à croissance modérée le long de} $X_o$ , \underline{il faut et il suffit que les} $f_{\alpha,k}$ \underline{de la décomposition} (4.7.1) \underline{soient méromorphes le long de} $X_o$ .

}{F7II-XIV-048}Pour la définition de "à croissance modérée", je renvoie à [ 1 ] 2.10. Il est clair sur (4.7.1), (4.7.2), (4.7.3) que $f$ est à croissance modérée si et seulement si les $f_{\alpha,k}$ le sont; pour des sections de $F$ sur $X^*$ , on sait
\SGApage{II}{157}
que "à croissance modérée le long de $X_o$" équivaut à "méromorphe le long de $X_o$ ".

Pour $X = D$ et $F = \mathcal{O}$ , 4.10 a le corollaire élémentaire suivant.

\underline{Corollaire} 4.11. \underline{Soient} $K$ \underline{le corps de germes en} $0$ \underline{de fonctions méromorphes sur} $D$ , \underline{et} $K^m$ (\underline{resp.} $K^{mu}$ , $K^{mqu}$ ) \underline{la} $K$-\underline{algèbre des germes en} $0$ \underline{de fonctions holomorphes sur} $\widetilde{D}^*$ , \underline{à croissance modérée en} $0$ \underline{et de détermination finie (resp. de détermination finie unipotente, resp. quasi-unipotente}). \underline{Alors}, $K^m$ , $K^{mu}$ , $K^{mqu}$ \underline{admettent sur} $K$ \underline{les bases suivantes}

$K^m$ \qquad \underline{les} \quad $z^\alpha (\log z)^k$ \qquad\qquad\qquad ($0 \leq \Re(\alpha) < 1$ , $k \in \mathbb{N}$)

$K^{mu}$ \qquad \underline{les} \quad $(\log z)^k$

\SGAvar{$K^{muq}$ \qquad \underline{les} \quad $z^\alpha (\log z)^k$ \qquad\qquad\qquad ($\alpha \in \mathbb{Q}$, $0 \leq \alpha < 1$, $k \in \mathbb{N}$)

Sous les hypothèses de 4.9, notons $\Psi^m_\eta(F)$ (resp. $\Psi^{mu}_\eta(F)$ , resp. $\Psi^{muq}_\eta(F)$ ) le sous-faisceau de $\overline{i}^* \overline{j}_*(\overline{F})$ formé des sections de $\overline{F}$ au voisinage des points de $X_o$ qui soient à croissance modérée le long de $X_o$ et de détermination finie (resp. finie unipotente, resp. finie quasi-unipotente). Soit $F^m \subset j_* \, j^* F$ le faisceau sur $X$ des sections de $j_* j^* F$ méromorphes le long de $X_o$ . La proposition 4.10 a le corollaire suivant, indépendant des choix arbitraires de $z$ et $\log z$

}{$K^{mqu}$ \qquad \underline{les} \quad $z^\alpha (\log z)^k$ \qquad\qquad\qquad ($\alpha \in \mathbb{Q}$, $0 \leq \alpha < 1$, $k \in \mathbb{N}$)

Sous les hypothèses de 4.9, notons $\Psi^m_\eta(F)$ (resp. $\Psi^{mu}_\eta(F)$ , resp. $\Psi^{mqu}_\eta(F)$ ) le sous-faisceau de $\overline{i}^* \overline{j}_*(\overline{F})$ formé des sections de $\overline{F}$ au voisinage des points de $X_o$ qui soient à croissance modérée le long de $X_o$ et de détermination finie (resp. finie unipotente, resp. finie quasi-unipotente). Soit $F^m \subset j_* \, j^* F$ le faisceau sur $X$ des sections de $j_* j^* F$ méromorphes le long de $X_o$ . La proposition 4.10 a le corollaire suivant, indépendant des choix arbitraires de $z$ et $\log z$

}{F7II-XIV-049}\underline{Corollaire} 4.12. \underline{Les morphismes naturels}

\SGAvar{\begin{equation*}
\mathrm{(i)} \qquad K^m \otimes_K i^* F^m \longrightarrow \Psi^m_\eta(F) \quad ,
\end{equation*}
\begin{equation*}
\mathrm{(ii)} \qquad K^{mu} \otimes_K i^* F^m \longrightarrow \Psi^{mu}_\eta(F) \quad ,
\end{equation*}
\begin{equation*}
\mathrm{(iii)} \qquad K^{muq} \otimes_K i^* F^m \longrightarrow \Psi^{muq}_\eta(F)
\end{equation*}
\underline{sont des isomorphismes.}

}{\begin{equation*}
\mathrm{(i)} \qquad K^m \otimes_K i^* F^m \longrightarrow \Psi^m_\eta(F) \quad ,
\end{equation*}
\begin{equation*}
\mathrm{(ii)} \qquad K^{mu} \otimes_K i^* F^m \longrightarrow \Psi^{mu}_\eta(F) \quad ,
\end{equation*}
\begin{equation*}
\mathrm{(iii)} \qquad K^{mqu} \otimes_K i^* F^m \longrightarrow \Psi^{mqu}_\eta(F)
\end{equation*}
\underline{sont des isomorphismes.}

}{F7II-XIV-050}Le résultat principal de ce paragraphe est le suivant.

\underline{Théorème} 4.13. \underline{Le morphisme naturel, induit par} 4.3.1,

\begin{equation}\tag{4.13.1}
\Psi^{mqu}_\eta(\Omega^*_X) \; \longrightarrow \; R\Psi_\eta(\mathbb{C}) 
\end{equation}

\SGApage{II}{158}

 \underline{est un isomorphisme dans} $D^+(X_o, \mathbb{C})$ .

Nous prouverons en fait des résultats plus généraux (4.15  ) (4.17.  ) dont l'énoncé requiert quelques préliminaires que le lecteur intéressé seulement par 4.13 pourra se dispenser de lire.

\SGAvar{4.14. Soit $V$ un faisceau localement constant de $\mathbb{C}$-vectoriels de dimension finie sur $X^*$ (un système local). Soit $\mathcal{V} = V \otimes_{\mathbb{C}} \mathcal{O}$ le $\mathcal{O}$-module localement libre correspondant et $\nabla$ la connexion intégrable sur $\mathcal{V}$ pour laquelle les sections locales de $V$ soit horizontale. On sait que $\mathcal{V}$ admet un prolongement $\mathcal{V}^\sim$ en un faisceau cohérent sur $X$ , tel que la connexion $\nabla$ soit \underline{régulière} le long de $X_o$ relativement à ce prolongement. On sait aussi que deux tels prolongements sont méromorphiquement équivalents ([1] II 5.7); le faisceau $\Psi^m_\eta(\mathcal{V}^\sim)$ , (et ses variantes $\Psi^{mu}$ , $\Psi^{mqu}$) , ne dépend donc pas du $\mathcal{V}$ choisi; on le note $\Psi^m_\eta(\mathcal{V})$ . On définit de même $\Psi^m_\eta(\Omega^p_{X^*}(V)) = \Psi^m_\eta(\Omega^p_X \otimes \mathcal{V})$ , et ses variantes $\Psi^{mu}$ , $\Psi^{mqu}$ .

}{4.14. Soit $V$ un faisceau localement constant de $\mathbb{C}$-vectoriels de dimension finie sur $X^*$ (un système local). Soit $\mathcal{V} = V \otimes_{\mathbb{C}} \mathcal{O}$ le $\mathcal{O}$-module localement libre correspondant et $\nabla$ la connexion intégrable sur $\mathcal{V}$ pour laquelle les sections locales de $V$ soient horizontales. On sait que $\mathcal{V}$ admet un prolongement $\mathcal{V}^\sim$ en un faisceau cohérent sur $X$ , tel que la connexion $\nabla$ soit \underline{régulière} le long de $X_o$ relativement à ce prolongement. On sait aussi que deux tels prolongements sont méromorphiquement équivalents ([1] II 5.7); le faisceau $\Psi^m_\eta(\mathcal{V}^\sim)$ , (et ses variantes $\Psi^{mu}$ , $\Psi^{mqu}$) , ne dépend donc pas du $\mathcal{V}^\sim$ choisi; on le note $\Psi^m_\eta(\mathcal{V})$ . On définit de même $\Psi^m_\eta(\Omega^p_{X^*}(V)) = \Psi^m_\eta(\Omega^p_X \otimes \mathcal{V})$ , et ses variantes $\Psi^{mu}$ , $\Psi^{mqu}$ .

}{F7II-XIV-051}\underline{Proposition} 4.15. \underline{Le morphisme naturel, induit par} (4.4.1)

\begin{equation}\tag{4.15.1}
\Psi^m_\eta(\Omega^*_{X^*}(V)) \longrightarrow R\Psi_\eta(V)
\end{equation}
\underline{est un isomorphisme dans} $D^+(X_o, \mathbb{C})$ .

4.16. On dit que $V$ comme en 4.14 est \underline{unipotent} (resp. \underline{quasi-unipotent} ) le long de $X_o$ si, pour tout morphisme $u$ du disque standard $D_o$ dans $X$ , avec $u^{-1}(X_o) = \{0\}$ , la transformation de monodromie du système local $u^*V$ sur $D_o^*$ est unipotente (resp. quasi-unipotente).

\underline{Proposition} 4.17. \underline{Si} $V$ \underline{est} quasi-unipotent, \underline{le morphisme naturel induit par} (4.4.1)

\begin{equation}\tag{4.17.1}
\Psi^{mqu}_\eta(\Omega^*_{X^*}(V)) \longrightarrow R\Psi_\eta(V)
\end{equation}
\underline{est un isomorphisme dans} $D^+(X_o, \mathbb{C})$ .


\SGApage{II}{159}

Le théorème 4.13 est le cas particulier de 4.17 où $V = \underline{\mathbb{C}}$ .

4.18. Nous traiterons d'abord le cas où $X$ est lisse et où $(X_o)_{red}$ est un diviseur à croisements normaux dans $X$ . Les théorèmes étant locaux, on peut supposer que $f$ est l'application $f(\underline{z}) = \prod_1^d z_i^{\,n_i}$ , $(0 < d \leqslant e,\ n_i > 0\ )$ de $X = D^e$ dans $D$ , où $D$ désigne cette fois le disque standard $D = \{z \in \mathbb{C} \mid |z| < 1\}$ , et il suffit de prouver que la fibre en $0 \in D^e$ de (4.15.1) (resp. 4.17.1) est un quasi-isomorphisme.

Soit $U$ un voisinage de $0$ dans $D^e$ de la forme $|z_i| < a$ $(0 < a \leqslant 1)$. On définit $U_o$ , $\overline{U}$ , $U^*$, $\overline{U}^*$ en prenant l'image réciproque de $U$ dans 4.1.1. La cohomologie de $U^*$ ou $\overline{U}^*$ est indépendante de $a$ , de sorte que

\begin{equation*}
H^*(\overline{U}^*, V) \xrightarrow{\ \sim\ } (R^* \Psi_\eta(V))_o \qquad .
\end{equation*}
Le groupe fondamental $\mathbb{Z}^d$ de $X^* = (D^*)^d \times D^{e-d}$ (égal à celui de $U^*$ ) étant abélien, $V$ admet une filtration finie de quotients successifs des systèmes locaux de rang un. Par dévissage, nous pouvons donc supposer $V$ de rang un.

Pour $\mathcal{V}^\sim$ comme en 4.14, nous noterons $j_*^m\, \Omega^*_{X^*}(V)$ le sous-complexe de $j_*\, \Omega^*_{X^*}(V)$ de composantes les $(\Omega^p_X \otimes \mathcal{V}^\sim)^m$ . D'après [1] II 6.4.1 et 6.8, l'application naturelle déduite du quasi-isomorphisme $V \xrightarrow{\ \approx\ } \Omega^*_{X^*}(V)$ (lemme de Poincaré holomorphe),

\begin{equation}\tag{4.18.1}
j_*^m\, \Omega^*_{X^*}(V) \longrightarrow Rj_*\, V
\end{equation}
est un quasi-isomorphisme. On a

\SGAvar{\begin{equation}\tag{4.18.2}
H^*((j_*^m\, \Omega^*_{X^*}(V))_o) \xrightarrow{\ \sim\ } (R^*j_*\, V)_o \xleftarrow{\ \sim\ } H^*(U, V) \qquad .
\end{equation}
Ceci sera le seul ingrédient analytique de cette partie 4.18 de la démonstration.

}{\begin{equation}\tag{4.18.2}
H^*((j_*^m\, \Omega^*_{X^*}(V))_o) \xrightarrow{\ \sim\ } (R^*j_*\, V)_o \xleftarrow{\ \sim\ } H^*(U^*, V) \qquad .
\end{equation}
Ceci sera le seul ingrédient analytique de cette partie 4.18 de la démonstration.

}{F7II-XIV-052}\underline{Lemme} 4.18.3. \underline{Pour} $V$ \underline{de rang un, on a}

(i) \qquad \underline{pour} $V$ \underline{non trivial} : $H^*(U^*, V) = 0$


\SGApage{II}{160}

 (ii) \underline{pour} $V = \underline{\mathbb{C}}$ : $H^*(U^*, V) = \stackrel{*}{\Lambda} \mathbb{C}^d$ .

En effet, $U^*$ a le type d'homotopie d'un tore.

Soit $\eta$ le générateur de $H^1(D^*, \mathbb{Z})$ , et son image réciproque dans $H^1(U^*, \mathbb{Z})$ .

\underline{Lemme} 4.18.4. \underline{Pour} $V$ \underline{de rang un, l'application naturelle}

\SGAvar{\begin{equation*}
H^*(U^*, \mathbb{C}) \xrightarrow[(4.18.1)]{\sim} H^*((j^m_* \Omega^*_{X^*}(V))_o) \longrightarrow H^*(\Psi^{mu}_{\eta}(\Omega^*_{X^*}(V))_o)
\end{equation*}
\underline{identifie} $H^*(\Psi^{mu}_{\eta}(\Omega^*_{X^*}(V))_o)$ \underline{à} $H^*(U^*, V)/\eta H^*(U^*, V)$ .

}{\begin{equation*}
H^*(U^*, V) \xrightarrow[(4.18.1)]{\sim} H^*((j^m_* \Omega^*_{X^*}(V))_o) \longrightarrow H^*(\Psi^{mu}_{\eta}(\Omega^*_{X^*}(V))_o)
\end{equation*}
\underline{identifie} $H^*(\Psi^{mu}_{\eta}(\Omega^*_{X^*}(V))_o)$ \underline{à} $H^*(U^*, V)/\eta H^*(U^*, V)$ .

}{F7II-XIV-053}Utilisons 4.12 (ii), et l'identité $d \log f = \Sigma n_i \frac{dz_i}{z_i}$ . Le complexe $\Psi^{mu}_{\eta}(\Omega^*_{X^*}(V))_o$ s'identifie au complexe simple $\underline{\underline{s}} K$ associé (en prenant des sommes) au complexe double $K$ de composantes

\SGAvar{\begin{equation*}
K^{-p,n+p} = \frac{(\log f)^p}{p!} \; . \; (j^m_* \Omega^n_{X^*}(V))_o \quad .
\end{equation*}
Ce complexe est dans un mauvais quadrant $(p \leqslant 0, q \geqslant 0, p+q \geqslant 0)$; toutefois, $\underline{\underline{s}} K$ est limite inductive des sous-complexe $\underline{\underline{s}} \sigma_{p \geqslant p_o}(K)$ de composantes

}{\begin{equation*}
K^{-p,n+p} = \frac{(\log f)^p}{p!} \; . \; (j^m_* \Omega^n_{X^*}(V))_o \quad .
\end{equation*}
Ce complexe est dans un mauvais quadrant $(p \leqslant 0, q \geqslant 0, p+q \geqslant 0)$; toutefois, $\underline{\underline{s}} K$ est limite inductive des sous-complexes $\underline{\underline{s}} \sigma_{p \geqslant p_o}(K)$ de composantes

}{F7II-XIV-054}\begin{multline*}
(\underline{\underline{s}} \sigma_{p \geqslant p_o} K)^n = \bigoplus_{\substack{p+q=n \\ p \geqslant p_o}} K^{pq} \quad ,
\end{multline*}
et cela suffit à justifier l'argument ci-dessous.

La suite spectrale de la filtration par le premier degré s'écrit

\begin{equation*}
E_1^{-p,n+p} = H^n(U^*, V) \qquad (p \geqslant 0) ,
\end{equation*}
avec $d_1 = 2\pi i \eta \wedge : H^n(U^*, V) \longrightarrow H^{n+1}(U^*, V)$ .

Si $V$ est non trivial, $E_1 = 0$ par 4.18.3 et 4.18.4 se réduit à $0 = 0$ . Supposons que $V = \underline{\mathbb{C}}$ . Les $(E_1, d_1)$ sont alors des complexes de Koszul, éventuellement tronqués, de sorte que (acyclicité du complexe de Koszul)
\SGApage{II}{161}

\begin{equation*}
E_2^{pq} = 0 \qquad \text{pour } p \neq 0
\end{equation*}
\begin{equation*}
E_2^{0n} = H^n(U^*, V)/\eta H^{n-1}(U^*, V)
\end{equation*}
ce qui vérifie 4.18.4.

Pour $\beta \in \mathbb{C}^*$ , notons $V_\beta$ le système local de rang un sur $D^*$ , d'image réciproque sur $\widetilde{D}^*$ trivialisée, et de monodromie $\beta$ . Notons encore $V_\beta$ son image réciproque sur $U^*$ . Sur $\overline{U}^*$ , l'image réciproque de $V \otimes V_\beta$ est canoniquement isomorphe à $V$ .

\underline{Lemme} 4.18.5. \underline{Pour} $V$ \underline{de rang un, l'application}

\begin{equation*}
\bigoplus_{\beta \in \mathbb{C}^*} H^*(U^*, V \otimes V_\beta)/\eta H^*(U^*, V \otimes V_\beta) \longrightarrow H^*(\overline{U}^*, V)
\end{equation*}
\underline{est bijective.}

\SGAvar{Soit $k$ le pgcd des $n_i$ . L'espace $\overline{U}^*$ à $k$ composantes connexes, paramétrée par les racines $k^{\text{ièmes}}$ de l'unité: sur celle, $\overline{U}^*_\xi$ , d'indice $\xi$ , l'image réciproque de la fonction $z^{1/k}$ sur $\widetilde{D}^*$ coïncide avec l'image réciproque de la fonction $\xi : \prod_1^d z_i^{n_i/k}$ sur $U^*$ . La composante $\overline{U}^*_\xi$ a le type d'homotopie d'un tore, et la suite

\begin{equation*}
0 \longrightarrow \pi_1(\overline{U}^*_\xi) \longrightarrow \pi_1(U^*) \longrightarrow \pi_1(D^*)
\end{equation*}
est exacte. Notons $\varphi : \mathbb{Z}^d \to \mathbb{Z}$ l'application $\pi_1(U^*) \to \pi_1(D^*)$ . Si $\mathrm{Ker}(\varphi)$ agit non trivialement sur $V$ , les deux membres de 4.18.4 sont nuls. Sinon, remplaçant $V$ par un $V \otimes V_\beta$ , on peut supposer que $V$ est trivial, soit $V = \mathbb{C}$ . On a alors en cohomologie complexe

}{Soit $k$ le pgcd des $n_i$ . L'espace $\overline{U}^*$ a $k$ composantes connexes, paramétrées par les racines $k^{\text{ièmes}}$ de l'unité: sur celle, $\overline{U}^*_\xi$ , d'indice $\xi$ , l'image réciproque de la fonction $z^{1/k}$ sur $\widetilde{D}^*$ coïncide avec l'image réciproque de la fonction $\xi \cdot \prod_1^d z_i^{n_i/k}$ sur $U^*$ . La composante $\overline{U}^*_\xi$ a le type d'homotopie d'un tore, et la suite

\begin{equation*}
0 \longrightarrow \pi_1(\overline{U}^*_\xi) \longrightarrow \pi_1(U^*) \longrightarrow \pi_1(D^*)
\end{equation*}
est exacte. Notons $\varphi : \mathbb{Z}^d \to \mathbb{Z}$ l'application $\pi_1(U^*) \to \pi_1(D^*)$ . Si $\mathrm{Ker}(\varphi)$ agit non trivialement sur $V$ , les deux membres de 4.18.5 sont nuls. Sinon, remplaçant $V$ par un $V \otimes V_\beta$ , on peut supposer que $V$ est trivial, soit $V = \mathbb{C}$ . On a alors en cohomologie complexe

}{F7II-XIV-055}\begin{equation*}
H^0(\overline{U}^*) \otimes H^*(U^*)/\eta H^*(U^*) \xrightarrow{\;\sim\;} H^*(\overline{U}^*)
\end{equation*}
et on conclut en notant que $V_\beta$ est trivial sur $U^*$ si $\beta^k = 1$ , que

\begin{equation*}
\bigoplus_{\beta^k=1} H^0(U^*, V_\beta) = \bigoplus_{\beta^k=1} \mathbb{C} \xrightarrow{\;\sim\;} H^0(\overline{U}^*, \mathbb{C}) \quad ,
\end{equation*}

\SGApage{II}{162}
et que, pour $\beta^k \neq 1$ , $H^*(U^*, V_\beta) = 0$ .

\underline{Remarque} 4.18.6. Si $V$ est quasi-unipotent, on peut dans 4.18.5 remplacer la sommation sur $\mathbb{C}^*$ par une sommation sur les racines de l'unité: la preuve est la même.

Appliquant 4.12, on trouve un diagramme commutatif

\begin{center}
\begin{tikzcd}[column sep=large]
H^*\Psi^m(\Omega^*_{X^*}(V))_o \arrow[d] \& \underset{\beta}{\oplus} H^*\Psi^{mu}(\Omega^*_{X^*}(V\otimes V_\beta)) \arrow[l, "\sim"', "4.12"] \& \underset{\beta}{\oplus} H^*(U^*,V\otimes V_\beta)/\eta\, H^*(U^*, V\otimes V_\beta)) \arrow[l, "\sim"', "4.18.4"] \arrow[d, "\wr"', "4.18.5"] \\
H^* R\Psi(V)_o \& {} \& H^*(\bar{U}^*, V) \quad , \arrow[ll, "\sim"']
\end{tikzcd}
\end{center}



















dont 4.15 résulte. On obtient 4.17 en y remplaçant $\Psi^m$ par $\Psi^{mqu}$ , la sommation sur $\beta$ par une sommation sur les racines de l'unité et 4.18.5 par 4.18.6.

4.19. Déduisons 4.15 et 4.17 du cas particulier 4.18. Soit $u : Y \longrightarrow X$ une résolution de $(X, X_o)$ , c'est-à-dire un morphisme propre, induisant un isomorphisme $Y^* \stackrel{\sim}{\longrightarrow} X^*$ , tel que $Y$ soit lisse et que $(Y_o)_{\mathrm{red}}$ soit un diviseur à croisements normaux de $Y$ . Si $V$ est quasi-unipotent le long de $X_o$ , son image réciproque $u^*V$ est quasi-unipotente le long de $Y_o$ .

\SGAvar{Pour tout faisceau cohérent $F$ sur $X$ , notons $F^m$ ou $j^m_* j^* F$ le sous-faisceau de $j_* j^* F$ formé des sections méromorphe, si $f$ est l'équation de $X_o$ , $j^m_* j^* F$ est la limite inductive du système inductif

}{Pour tout faisceau cohérent $F$ sur $X$ , notons $F^m$ ou $j^m_* j^* F$ le sous-faisceau de $j_* j^* F$ formé des sections méromorphes, si $f$ est l'équation de $X_o$ , $j^m_* j^* F$ est la limite inductive du système inductif

}{F7II-XIV-056}\begin{equation}\tag{4.19.1}
F \stackrel{f}{\longrightarrow} F \stackrel{f}{\longrightarrow} F \longrightarrow \quad \ldots\ldots
\end{equation}
Le lemme suivant est prouvé dans [2].

\underline{Lemme} 4.19.2. \underline{Pour tout faisceau cohérent} $F$ \underline{sur} $X$ , \underline{on a}

\begin{equation*}
F^m \stackrel{\sim}{\longrightarrow} u_*(u^* F)^m
\end{equation*}
\underline{et} $R^i u_* (u^* F)^m = 0$ \underline{pour} $i > 0$ .


\SGApage{II}{163}

\SGAvar{D'après le théorème de Grauert, les $R^i u_* u^* F$ sont des faisceaux cohérents. Les noyaux et conoyaux de $F \longrightarrow u_* u^* F$ , et les $R^i u_* F$ $(i > 0)$ sont nuls en dehors de $X_o$ , donc localement annulés par une puissance de $f$ . Puisque $u$ est propre, $R^i u_*$ commute aux limites inductives filtrantes, en particulier à la limite inductive (4.19.1), et (4.19.2) en résulte.

}{D'après le théorème de Grauert, les $R^i u_* u^* F$ sont des faisceaux cohérents. Les noyaux et conoyaux de $F \longrightarrow u_* u^* F$ , et les $R^i u_* u^* F$ $(i > 0)$ sont nuls en dehors de $X_o$ , donc localement annulés par une puissance de $f$ . Puisque $u$ est propre, $R^i u_*$ commute aux limites inductives filtrantes, en particulier à la limite inductive (4.19.1), et (4.19.2) en résulte.

}{F7II-XIV-057}De 4.19.2 et 4.12 on déduit que, sous les hypothèses de loc. cit. et désignant encore par $u$ l'application de $Y_o$ dans $X_o$ induite par $u$ , on a:

\underline{Lemme} 4.19.3. \underline{Les morphismes}

\begin{equation*}
\Psi^m_\eta(F) \longrightarrow u_* \Psi^m_\eta(u^* F) \longrightarrow Ru_* \Psi^m_\eta(u^* F)
\end{equation*}
\underline{sont des isomorphismes.}

Que la seconde flèche soit un isomorphisme signifie que $R^i u_* \Psi^m_\eta(u^* F) = 0$ pour $i > 0$ . Ceci résulte de 4.12 et de $R^i u_*(i^{\cdot}u^* F^m) = 0$ pour $i > 0$ , qui résulte de $R^i u_*(u^* F^m) = 0$ pour $i > 0$ et du théorème de changement de base pour le morphisme propre $u$ .

Pour prouver 4.15, il reste à contempler le diagramme commutatif

\begin{center}
\begin{tikzcd}[column sep=large]
\Psi^m(\Omega^*_{X^*}(V)) \arrow[d] \arrow[r, equal, "(4.19.3)"'] \& u_* \Psi^m(\Omega^*_{Y^*}(u^{\cdot}V)) \arrow[r, "\sim", "(4.19.3)"'] \& Ru_* \Psi^m \Omega^*_{Y^*}(u^{\cdot}V) \arrow[d, "\wr"', "(4.18)"] \\
R\Psi(V) \arrow[rr, "(1.3.1; XIII 2.1.7.1)"] \& {} \& Ru_* R\Psi(u^* V) \quad .
\end{tikzcd}
\end{center}















On procède de même pour 4.17, en remplaçant partout $\Psi^m$ par $\Psi^{mqu}$ .

\underline{Remarque} 4.20. Le théorème 4.13 fournit une incarnation du théorème de monodromie au niveau des complexes.


\SGApage{II}{164}


\underline{Bibliographie}

[1] \hspace{2em} P. Deligne - Equations différentielles à points singuliers réguliers. Lecture Notes in Math. 163, Springer Verlag.

[2] \hspace{2em} A. Grothendieck - On the De Rham cohomology of algebraic varieties. Publ. Math. IHES \underline{29}, 95-103 (1966).

[3] \hspace{2em} K. Lamotke - Isolated critical points and monodromy.
\selectlanguage{french}

\ifnum\SGAform=1
\clearpage\section*{Apparat critique}\small
Les lectures ci-dessous sont données dans le TeX exact des témoins, afin de préserver les distinctions de notation. Les indices décimaux renvoient aux différences individuelles du passage marqué.

\par\medskip\noindent\hypertarget{app-F7II-XIV-001.1}{}\textbf{F7II-XIV-001.1}\quad p. 117; misprint\_pass.\par
\noindent Imprimé : {\ttfamily\footnotesize ambiguités\allowbreak{}}\par
\noindent Corrigé : {\ttfamily\footnotesize ambiguïtés\allowbreak{}}\par
\noindent Source : 7II:183. \nolinkurl{fixes/FIXLOG_7II_CURRENT.jsonl#587}\par

\par\medskip\noindent\hypertarget{app-F7II-XIV-002.1}{}\textbf{F7II-XIV-002.1}\quad p. 117; misprint\_pass.\par
\noindent Imprimé : {\ttfamily\footnotesize fibre\allowbreak{}}\par
\noindent Corrigé : {\ttfamily\footnotesize fibres\allowbreak{}}\par
\noindent Source : 7II:182. \nolinkurl{fixes/FIXLOG_7II_CURRENT.jsonl#586}; \nolinkurl{fixes/MISPRINTS_7II.json#182}\par

\par\medskip\noindent\hypertarget{app-F7II-XIV-003.1}{}\textbf{F7II-XIV-003.1}\quad p. 118; misprint\_pass.\par
\noindent Imprimé : \emph{(vide)}\par
\noindent Corrigé : {\ttfamily\footnotesize \textasciicircum{}\{*\}\allowbreak{}}\par
\noindent Source : 7II:184. \nolinkurl{fixes/FIXLOG_7II_CURRENT.jsonl#588}; \nolinkurl{fixes/MISPRINTS_7II.json#184}\par

\par\medskip\noindent\hypertarget{app-F7II-XIV-004.1}{}\textbf{F7II-XIV-004.1}\quad p. 119; misprint\_pass.\par
\noindent Imprimé : {\ttfamily\footnotesize constant\allowbreak{}}\par
\noindent Corrigé : {\ttfamily\footnotesize constants\allowbreak{}}\par
\noindent Source : 7II:185. \nolinkurl{fixes/FIXLOG_7II_CURRENT.jsonl#589}; \nolinkurl{fixes/MISPRINTS_7II.json#185}\par

\par\medskip\noindent\hypertarget{app-F7II-XIV-004.2}{}\textbf{F7II-XIV-004.2}\quad p. 119; misprint\_pass.\par
\noindent Imprimé : {\ttfamily\footnotesize constant\allowbreak{}}\par
\noindent Corrigé : {\ttfamily\footnotesize constants\allowbreak{}}\par
\noindent Source : 7II:186. \nolinkurl{fixes/FIXLOG_7II_CURRENT.jsonl#590}; \nolinkurl{fixes/MISPRINTS_7II.json#186}\par

\par\medskip\noindent\hypertarget{app-F7II-XIV-005.1}{}\textbf{F7II-XIV-005.1}\quad p. 120; misprint\_pass.\par
\noindent Imprimé : {\ttfamily\footnotesize formant\allowbreak{}}\par
\noindent Corrigé : {\ttfamily\footnotesize forment\allowbreak{}}\par
\noindent Source : 7II:187. \nolinkurl{fixes/FIXLOG_7II_CURRENT.jsonl#591}; \nolinkurl{fixes/MISPRINTS_7II.json#187}\par

\par\medskip\noindent\hypertarget{app-F7II-XIV-006.1}{}\textbf{F7II-XIV-006.1}\quad p. 121; misprint\_pass.\par
\noindent Imprimé : {\ttfamily\footnotesize 1\allowbreak{}}\par
\noindent Corrigé : {\ttfamily\footnotesize 2\allowbreak{}}\par
\noindent Source : 7II:188. \nolinkurl{fixes/FIXLOG_7II_CURRENT.jsonl#592}; \nolinkurl{fixes/MISPRINTS_7II.json#188}\par

\par\medskip\noindent\hypertarget{app-F7II-XIV-007.1}{}\textbf{F7II-XIV-007.1}\quad p. 122; printed misprint.\par
\noindent Imprimé : {\ttfamily\footnotesize universel\allowbreak{}}\par
\noindent Corrigé : {\ttfamily\footnotesize universels\allowbreak{}}\par
\noindent Source : 7II:189. \nolinkurl{fixes/MISPRINTS_7II.json#189}\par

\par\medskip\noindent\hypertarget{app-F7II-XIV-008.1}{}\textbf{F7II-XIV-008.1}\quad p. 124; printed misprint.\par
\noindent Imprimé : \emph{(vide)}\par
\noindent Corrigé : {\ttfamily\footnotesize \textbackslash{}overline\{\allowbreak{}}\par
\noindent Source : 7II:190. \nolinkurl{fixes/MISPRINTS_7II.json#190}\par

\par\medskip\noindent\hypertarget{app-F7II-XIV-008.2}{}\textbf{F7II-XIV-008.2}\quad p. 124; printed misprint.\par
\noindent Imprimé : \emph{(vide)}\par
\noindent Corrigé : {\ttfamily\footnotesize \}\allowbreak{}}\par
\noindent Source : 7II:190. \nolinkurl{fixes/MISPRINTS_7II.json#190}\par

\par\medskip\noindent\hypertarget{app-F7II-XIV-009.1}{}\textbf{F7II-XIV-009.1}\quad p. 127; misprint\_pass.\par
\noindent Imprimé : {\ttfamily\footnotesize retrécir\allowbreak{}}\par
\noindent Corrigé : {\ttfamily\footnotesize rétrécir\allowbreak{}}\par
\noindent Source : 7II:191. \nolinkurl{fixes/FIXLOG_7II_CURRENT.jsonl#594}; \nolinkurl{fixes/MISPRINTS_7II.json#191}\par

\par\medskip\noindent\hypertarget{app-F7II-XIV-010.1}{}\textbf{F7II-XIV-010.1}\quad p. 127; inherited corrected-layer difference.\par
\noindent Imprimé : {\ttfamily\footnotesize 1\allowbreak{}}\par
\noindent Corrigé : {\ttfamily\footnotesize 6\allowbreak{}}\par
\noindent Source : ADOPT-F7II-XIV-010.1. \nolinkurl{fr/sga7ii/diplomatic/expose_XIV_body.tex}; \nolinkurl{fr/sga7ii/current/expose_XIV_body.tex}\par

\par\medskip\noindent\hypertarget{app-F7II-XIV-010.2}{}\textbf{F7II-XIV-010.2}\quad p. 127; misprint\_pass.\par
\noindent Imprimé : {\ttfamily\footnotesize retrécir\allowbreak{}}\par
\noindent Corrigé : {\ttfamily\footnotesize rétrécir\allowbreak{}}\par
\noindent Source : 7II:192. \nolinkurl{fixes/FIXLOG_7II_CURRENT.jsonl#595}; \nolinkurl{fixes/MISPRINTS_7II.json#192}\par

\par\medskip\noindent\hypertarget{app-F7II-XIV-011.1}{}\textbf{F7II-XIV-011.1}\quad p. 128; printed misprint.\par
\noindent Imprimé : {\ttfamily\footnotesize universel\allowbreak{}}\par
\noindent Corrigé : {\ttfamily\footnotesize universels\allowbreak{}}\par
\noindent Source : 7II:193. \nolinkurl{fixes/MISPRINTS_7II.json#193}\par

\par\medskip\noindent\hypertarget{app-F7II-XIV-012.1}{}\textbf{F7II-XIV-012.1}\quad p. 129; misprint\_in\_print.\par
\noindent Imprimé : {\ttfamily\footnotesize ,\allowbreak{} dash\allowbreak{}}\par
\noindent Corrigé : \emph{(vide)}\par
\noindent Source : FIX-7II-D-0110. \nolinkurl{fixes/FIXLOG_7II.jsonl#110}\par

\par\medskip\noindent\hypertarget{app-F7II-XIV-013.1}{}\textbf{F7II-XIV-013.1}\quad p. 129; misprint\_pass.\par
\noindent Imprimé : {\ttfamily\footnotesize voisinage\allowbreak{}}\par
\noindent Corrigé : {\ttfamily\footnotesize voisinages\allowbreak{}}\par
\noindent Source : 7II:195. \nolinkurl{fixes/FIXLOG_7II_CURRENT.jsonl#597}; \nolinkurl{fixes/MISPRINTS_7II.json#195}\par

\par\medskip\noindent\hypertarget{app-F7II-XIV-014.1}{}\textbf{F7II-XIV-014.1}\quad p. 131; misprint\_pass.\par
\noindent Imprimé : {\ttfamily\footnotesize réduit\allowbreak{}}\par
\noindent Corrigé : {\ttfamily\footnotesize réduits\allowbreak{}}\par
\noindent Source : 7II:196. \nolinkurl{fixes/FIXLOG_7II_CURRENT.jsonl#598}; \nolinkurl{fixes/MISPRINTS_7II.json#196}\par

\par\medskip\noindent\hypertarget{app-F7II-XIV-015.1}{}\textbf{F7II-XIV-015.1}\quad p. 131; misprint\_pass.\par
\noindent Imprimé : {\ttfamily\footnotesize coincide\allowbreak{}}\par
\noindent Corrigé : {\ttfamily\footnotesize coïncide\allowbreak{}}\par
\noindent Source : 7II:197. \nolinkurl{fixes/FIXLOG_7II_CURRENT.jsonl#599}\par

\par\medskip\noindent\hypertarget{app-F7II-XIV-016.1}{}\textbf{F7II-XIV-016.1}\quad p. 132; misprint\_pass.\par
\noindent Imprimé : {\ttfamily\footnotesize henselisé\allowbreak{}}\par
\noindent Corrigé : {\ttfamily\footnotesize hensélisé\allowbreak{}}\par
\noindent Source : 7II:198. \nolinkurl{fixes/FIXLOG_7II_CURRENT.jsonl#600}\par

\par\medskip\noindent\hypertarget{app-F7II-XIV-017.1}{}\textbf{F7II-XIV-017.1}\quad p. 134; misprint\_pass.\par
\noindent Imprimé : {\ttfamily\footnotesize retrécir\allowbreak{}}\par
\noindent Corrigé : {\ttfamily\footnotesize rétrécir\allowbreak{}}\par
\noindent Source : 7II:199. \nolinkurl{fixes/FIXLOG_7II_CURRENT.jsonl#601}\par

\par\medskip\noindent\hypertarget{app-F7II-XIV-018.1}{}\textbf{F7II-XIV-018.1}\quad p. 136; misprint\_pass.\par
\noindent Imprimé : {\ttfamily\footnotesize r\allowbreak{}}\par
\noindent Corrigé : {\ttfamily\footnotesize \textbackslash{}tau\allowbreak{}}\par
\noindent Source : 7II:200. \nolinkurl{fixes/FIXLOG_7II_CURRENT.jsonl#602}; \nolinkurl{fixes/MISPRINTS_7II.json#200}\par

\par\medskip\noindent\hypertarget{app-F7II-XIV-019.1}{}\textbf{F7II-XIV-019.1}\quad p. 136; misprint\_pass.\par
\noindent Imprimé : {\ttfamily\footnotesize soit\allowbreak{}}\par
\noindent Corrigé : {\ttfamily\footnotesize soient\allowbreak{}}\par
\noindent Source : 7II:204. \nolinkurl{fixes/FIXLOG_7II_CURRENT.jsonl#606}\par

\par\medskip\noindent\hypertarget{app-F7II-XIV-019.2}{}\textbf{F7II-XIV-019.2}\quad p. 136; misprint\_pass.\par
\noindent Imprimé : {\ttfamily\footnotesize soins\allowbreak{}}\par
\noindent Corrigé : {\ttfamily\footnotesize soin\allowbreak{}}\par
\noindent Source : 7II:201. \nolinkurl{fixes/FIXLOG_7II_CURRENT.jsonl#603}; \nolinkurl{fixes/MISPRINTS_7II.json#201}\par

\par\medskip\noindent\hypertarget{app-F7II-XIV-019.3}{}\textbf{F7II-XIV-019.3}\quad p. 136; misprint\_pass.\par
\noindent Imprimé : {\ttfamily\footnotesize s\allowbreak{}}\par
\noindent Corrigé : {\ttfamily\footnotesize d\allowbreak{}}\par
\noindent Source : 7II:202. \nolinkurl{fixes/FIXLOG_7II_CURRENT.jsonl#604}; \nolinkurl{fixes/MISPRINTS_7II.json#202}\par

\par\medskip\noindent\hypertarget{app-F7II-XIV-019.4}{}\textbf{F7II-XIV-019.4}\quad p. 136; misprint\_pass.\par
\noindent Imprimé : {\ttfamily\footnotesize positif\allowbreak{}}\par
\noindent Corrigé : {\ttfamily\footnotesize positive\allowbreak{}}\par
\noindent Source : 7II:203. \nolinkurl{fixes/FIXLOG_7II_CURRENT.jsonl#605}; \nolinkurl{fixes/MISPRINTS_7II.json#203}\par

\par\medskip\noindent\hypertarget{app-F7II-XIV-020.1}{}\textbf{F7II-XIV-020.1}\quad p. 137; misprint\_pass.\par
\noindent Imprimé : {\ttfamily\footnotesize s\allowbreak{}}\par
\noindent Corrigé : {\ttfamily\footnotesize t\allowbreak{}}\par
\noindent Source : 7II:206. \nolinkurl{fixes/FIXLOG_7II_CURRENT.jsonl#608}; \nolinkurl{fixes/MISPRINTS_7II.json#206}\par

\par\medskip\noindent\hypertarget{app-F7II-XIV-020.2}{}\textbf{F7II-XIV-020.2}\quad p. 137; misprint\_pass.\par
\noindent Imprimé : {\ttfamily\footnotesize t\allowbreak{}}\par
\noindent Corrigé : {\ttfamily\footnotesize d\allowbreak{}}\par
\noindent Source : 7II:206. \nolinkurl{fixes/FIXLOG_7II_CURRENT.jsonl#608}; \nolinkurl{fixes/MISPRINTS_7II.json#206}\par

\par\medskip\noindent\hypertarget{app-F7II-XIV-021.1}{}\textbf{F7II-XIV-021.1}\quad p. 137; misprint\_pass.\par
\noindent Imprimé : {\ttfamily\footnotesize |\allowbreak{}}\par
\noindent Corrigé : \emph{(vide)}\par
\noindent Source : 7II:205. \nolinkurl{fixes/FIXLOG_7II_CURRENT.jsonl#607}; \nolinkurl{fixes/MISPRINTS_7II.json#205}\par

\par\medskip\noindent\hypertarget{app-F7II-XIV-021.2}{}\textbf{F7II-XIV-021.2}\quad p. 137; misprint\_pass.\par
\noindent Imprimé : {\ttfamily\footnotesize |\allowbreak{}}\par
\noindent Corrigé : \emph{(vide)}\par
\noindent Source : 7II:205. \nolinkurl{fixes/FIXLOG_7II_CURRENT.jsonl#607}; \nolinkurl{fixes/MISPRINTS_7II.json#205}\par

\par\medskip\noindent\hypertarget{app-F7II-XIV-021.3}{}\textbf{F7II-XIV-021.3}\quad p. 137; misprint\_pass.\par
\noindent Imprimé : {\ttfamily\footnotesize champs\allowbreak{}}\par
\noindent Corrigé : {\ttfamily\footnotesize champ\allowbreak{}}\par
\noindent Source : 7II:207. \nolinkurl{fixes/FIXLOG_7II_CURRENT.jsonl#609}; \nolinkurl{fixes/MISPRINTS_7II.json#207}\par

\par\medskip\noindent\hypertarget{app-F7II-XIV-021.4}{}\textbf{F7II-XIV-021.4}\quad p. 137; misprint\_pass.\par
\noindent Imprimé : {\ttfamily\footnotesize vecteur\allowbreak{}}\par
\noindent Corrigé : {\ttfamily\footnotesize vecteurs\allowbreak{}}\par
\noindent Source : 7II:207. \nolinkurl{fixes/FIXLOG_7II_CURRENT.jsonl#609}; \nolinkurl{fixes/MISPRINTS_7II.json#207}\par

\par\medskip\noindent\hypertarget{app-F7II-XIV-022.1}{}\textbf{F7II-XIV-022.1}\quad p. 138; misprint\_pass.\par
\noindent Imprimé : {\ttfamily\footnotesize fonction\allowbreak{}}\par
\noindent Corrigé : {\ttfamily\footnotesize fonctions\allowbreak{}}\par
\noindent Source : 7II:209. \nolinkurl{fixes/FIXLOG_7II_CURRENT.jsonl#610}; \nolinkurl{fixes/MISPRINTS_7II.json#209}\par

\par\medskip\noindent\hypertarget{app-F7II-XIV-023.1}{}\textbf{F7II-XIV-023.1}\quad p. 139; misprint\_pass.\par
\noindent Imprimé : \emph{(vide)}\par
\noindent Corrigé : {\ttfamily\footnotesize f\allowbreak{}}\par
\noindent Source : 7II:210. \nolinkurl{fixes/FIXLOG_7II_CURRENT.jsonl#611}; \nolinkurl{fixes/MISPRINTS_7II.json#210}\par

\par\medskip\noindent\hypertarget{app-F7II-XIV-024.1}{}\textbf{F7II-XIV-024.1}\quad p. 140; misprint\_pass.\par
\noindent Imprimé : {\ttfamily\footnotesize n\allowbreak{}}\par
\noindent Corrigé : {\ttfamily\footnotesize x\allowbreak{}}\par
\noindent Source : 7II:211. \nolinkurl{fixes/FIXLOG_7II_CURRENT.jsonl#612}; \nolinkurl{fixes/MISPRINTS_7II.json#211}\par

\par\medskip\noindent\hypertarget{app-F7II-XIV-024.2}{}\textbf{F7II-XIV-024.2}\quad p. 140; misprint\_pass.\par
\noindent Imprimé : {\ttfamily\footnotesize n\allowbreak{}}\par
\noindent Corrigé : {\ttfamily\footnotesize x\allowbreak{}}\par
\noindent Source : 7II:211. \nolinkurl{fixes/FIXLOG_7II_CURRENT.jsonl#612}; \nolinkurl{fixes/MISPRINTS_7II.json#211}\par

\par\medskip\noindent\hypertarget{app-F7II-XIV-025.1}{}\textbf{F7II-XIV-025.1}\quad p. 140; misprint\_pass.\par
\noindent Imprimé : \emph{(vide)}\par
\noindent Corrigé : {\ttfamily\footnotesize \textbackslash{}widetilde\{\allowbreak{}}\par
\noindent Source : 7II:212. \nolinkurl{fixes/FIXLOG_7II_CURRENT.jsonl#613}; \nolinkurl{fixes/MISPRINTS_7II.json#212}\par

\par\medskip\noindent\hypertarget{app-F7II-XIV-025.2}{}\textbf{F7II-XIV-025.2}\quad p. 140; misprint\_pass.\par
\noindent Imprimé : \emph{(vide)}\par
\noindent Corrigé : {\ttfamily\footnotesize \}\allowbreak{}}\par
\noindent Source : 7II:212. \nolinkurl{fixes/FIXLOG_7II_CURRENT.jsonl#613}; \nolinkurl{fixes/MISPRINTS_7II.json#212}\par

\par\medskip\noindent\hypertarget{app-F7II-XIV-026.1}{}\textbf{F7II-XIV-026.1}\quad p. 142; misprint\_pass.\par
\noindent Imprimé : {\ttfamily\footnotesize résoud\allowbreak{}}\par
\noindent Corrigé : {\ttfamily\footnotesize résout\allowbreak{}}\par
\noindent Source : 7II:213. \nolinkurl{fixes/FIXLOG_7II_CURRENT.jsonl#614}; \nolinkurl{fixes/MISPRINTS_7II.json#213}\par

\par\medskip\noindent\hypertarget{app-F7II-XIV-027.1}{}\textbf{F7II-XIV-027.1}\quad p. 142; misprint\_pass.\par
\noindent Imprimé : {\ttfamily\footnotesize raisonables\allowbreak{}}\par
\noindent Corrigé : {\ttfamily\footnotesize raisonnables\allowbreak{}}\par
\noindent Source : 7II:214. \nolinkurl{fixes/FIXLOG_7II_CURRENT.jsonl#615}; \nolinkurl{fixes/MISPRINTS_7II.json#214}\par

\par\medskip\noindent\hypertarget{app-F7II-XIV-028.1}{}\textbf{F7II-XIV-028.1}\quad p. 143; misprint\_pass.\par
\noindent Imprimé : {\ttfamily\footnotesize coisinage\allowbreak{}}\par
\noindent Corrigé : {\ttfamily\footnotesize voisinage\allowbreak{}}\par
\noindent Source : 7II:215. \nolinkurl{fixes/FIXLOG_7II_CURRENT.jsonl#616}; \nolinkurl{fixes/MISPRINTS_7II.json#215}\par

\par\medskip\noindent\hypertarget{app-F7II-XIV-029.1}{}\textbf{F7II-XIV-029.1}\quad p. 143; misprint\_pass.\par
\noindent Imprimé : {\ttfamily\footnotesize induite\allowbreak{}}\par
\noindent Corrigé : {\ttfamily\footnotesize induit\allowbreak{}}\par
\noindent Source : 7II:216. \nolinkurl{fixes/FIXLOG_7II_CURRENT.jsonl#617}\par

\par\medskip\noindent\hypertarget{app-F7II-XIV-030.1}{}\textbf{F7II-XIV-030.1}\quad p. 144; misprint\_pass.\par
\noindent Imprimé : {\ttfamily\footnotesize \{\}\_1\allowbreak{}}\par
\noindent Corrigé : \emph{(vide)}\par
\noindent Source : 7II:217. \nolinkurl{fixes/FIXLOG_7II_CURRENT.jsonl#618}; \nolinkurl{fixes/MISPRINTS_7II.json#217}\par

\par\medskip\noindent\hypertarget{app-F7II-XIV-030.2}{}\textbf{F7II-XIV-030.2}\quad p. 144; misprint\_pass.\par
\noindent Imprimé : {\ttfamily\footnotesize complexe\allowbreak{}}\par
\noindent Corrigé : {\ttfamily\footnotesize complexes\allowbreak{}}\par
\noindent Source : 7II:218. \nolinkurl{fixes/FIXLOG_7II_CURRENT.jsonl#619}\par

\par\medskip\noindent\hypertarget{app-F7II-XIV-031.1}{}\textbf{F7II-XIV-031.1}\quad p. 144; misprint\_pass.\par
\noindent Imprimé : \emph{(vide)}\par
\noindent Corrigé : {\ttfamily\footnotesize \textasciicircum{}*\allowbreak{}}\par
\noindent Source : 7II:219. \nolinkurl{fixes/FIXLOG_7II_CURRENT.jsonl#620}; \nolinkurl{fixes/MISPRINTS_7II.json#219}\par

\par\medskip\noindent\hypertarget{app-F7II-XIV-032.1}{}\textbf{F7II-XIV-032.1}\quad p. 144; misprint\_pass.\par
\noindent Imprimé : {\ttfamily\footnotesize 15\allowbreak{}}\par
\noindent Corrigé : {\ttfamily\footnotesize 1.5\allowbreak{}}\par
\noindent Source : 7II:220. \nolinkurl{fixes/FIXLOG_7II_CURRENT.jsonl#621}; \nolinkurl{fixes/MISPRINTS_7II.json#220}\par

\par\medskip\noindent\hypertarget{app-F7II-XIV-033.1}{}\textbf{F7II-XIV-033.1}\quad p. 146; misprint\_pass.\par
\noindent Imprimé : \emph{(vide)}\par
\noindent Corrigé : {\ttfamily\footnotesize \_o\allowbreak{}}\par
\noindent Source : 7II:221. \nolinkurl{fixes/FIXLOG_7II_CURRENT.jsonl#622}; \nolinkurl{fixes/MISPRINTS_7II.json#221}\par

\par\medskip\noindent\hypertarget{app-F7II-XIV-034.1}{}\textbf{F7II-XIV-034.1}\quad p. 146; misprint\_pass.\par
\noindent Imprimé : {\ttfamily\footnotesize continument\allowbreak{}}\par
\noindent Corrigé : {\ttfamily\footnotesize continûment\allowbreak{}}\par
\noindent Source : 7II:222. \nolinkurl{fixes/FIXLOG_7II_CURRENT.jsonl#623}; \nolinkurl{fixes/MISPRINTS_7II.json#222}\par

\par\medskip\noindent\hypertarget{app-F7II-XIV-035.1}{}\textbf{F7II-XIV-035.1}\quad p. 147; printed misprint.\par
\noindent Imprimé : \emph{(vide)}\par
\noindent Corrigé : {\ttfamily\footnotesize \{\allowbreak{}}\par
\noindent Source : 7II:223. \nolinkurl{fixes/MISPRINTS_7II.json#223}\par

\par\medskip\noindent\hypertarget{app-F7II-XIV-035.2}{}\textbf{F7II-XIV-035.2}\quad p. 147; printed misprint.\par
\noindent Imprimé : \emph{(vide)}\par
\noindent Corrigé : {\ttfamily\footnotesize +1\}\allowbreak{}}\par
\noindent Source : 7II:223. \nolinkurl{fixes/MISPRINTS_7II.json#223}\par

\par\medskip\noindent\hypertarget{app-F7II-XIV-036.1}{}\textbf{F7II-XIV-036.1}\quad p. 148; misprint\_pass.\par
\noindent Imprimé : \emph{(1 espace(s), 0 retour(s) de ligne dans le témoin TeX)}\par
\noindent Corrigé : {\ttfamily\footnotesize -\allowbreak{}}\par
\noindent Source : 7II:225. \nolinkurl{fixes/FIXLOG_7II_CURRENT.jsonl#624}; \nolinkurl{fixes/MISPRINTS_7II.json#225}\par

\par\medskip\noindent\hypertarget{app-F7II-XIV-037.1}{}\textbf{F7II-XIV-037.1}\quad p. 148; misprint\_pass.\par
\noindent Imprimé : {\ttfamily\footnotesize geodésique\allowbreak{}}\par
\noindent Corrigé : {\ttfamily\footnotesize géodésique\allowbreak{}}\par
\noindent Source : 7II:226. \nolinkurl{fixes/FIXLOG_7II_CURRENT.jsonl#625}; \nolinkurl{fixes/MISPRINTS_7II.json#226}\par

\par\medskip\noindent\hypertarget{app-F7II-XIV-037.2}{}\textbf{F7II-XIV-037.2}\quad p. 148; misprint\_pass.\par
\noindent Imprimé : {\ttfamily\footnotesize 13\allowbreak{}}\par
\noindent Corrigé : {\ttfamily\footnotesize 12\allowbreak{}}\par
\noindent Source : 7II:227. \nolinkurl{fixes/FIXLOG_7II_CURRENT.jsonl#626}; \nolinkurl{fixes/MISPRINTS_7II.json#227}\par

\par\medskip\noindent\hypertarget{app-F7II-XIV-037.3}{}\textbf{F7II-XIV-037.3}\quad p. 148; misprint\_pass.\par
\noindent Imprimé : {\ttfamily\footnotesize geodésique\allowbreak{}}\par
\noindent Corrigé : {\ttfamily\footnotesize géodésique\allowbreak{}}\par
\noindent Source : 7II:228. \nolinkurl{fixes/FIXLOG_7II_CURRENT.jsonl#627}; \nolinkurl{fixes/MISPRINTS_7II.json#228}\par

\par\medskip\noindent\hypertarget{app-F7II-XIV-037.4}{}\textbf{F7II-XIV-037.4}\quad p. 148; misprint\_pass.\par
\noindent Imprimé : \emph{(vide)}\par
\noindent Corrigé : {\ttfamily\footnotesize )\allowbreak{}}\par
\noindent Source : 7II:229. \nolinkurl{fixes/FIXLOG_7II_CURRENT.jsonl#628}; \nolinkurl{fixes/MISPRINTS_7II.json#229}\par

\par\medskip\noindent\hypertarget{app-F7II-XIV-037.5}{}\textbf{F7II-XIV-037.5}\quad p. 148; printed misprint.\par
\noindent Imprimé : {\ttfamily\footnotesize x\allowbreak{}}\par
\noindent Corrigé : {\ttfamily\footnotesize d\allowbreak{}}\par
\noindent Source : 7II:230. \nolinkurl{fixes/MISPRINTS_7II.json#230}\par

\par\medskip\noindent\hypertarget{app-F7II-XIV-038.1}{}\textbf{F7II-XIV-038.1}\quad p. 149; inherited corrected-layer difference.\par
\noindent Imprimé : {\ttfamily\footnotesize n\allowbreak{}}\par
\noindent Corrigé : {\ttfamily\footnotesize d\allowbreak{}}\par
\noindent Source : ADOPT-F7II-XIV-038.1. \nolinkurl{fr/sga7ii/diplomatic/expose_XIV_body.tex}; \nolinkurl{fr/sga7ii/current/expose_XIV_body.tex}\par

\par\medskip\noindent\hypertarget{app-F7II-XIV-038.2}{}\textbf{F7II-XIV-038.2}\quad p. 149; inherited corrected-layer difference.\par
\noindent Imprimé : {\ttfamily\footnotesize n\allowbreak{}}\par
\noindent Corrigé : {\ttfamily\footnotesize d\allowbreak{}}\par
\noindent Source : ADOPT-F7II-XIV-038.2. \nolinkurl{fr/sga7ii/diplomatic/expose_XIV_body.tex}; \nolinkurl{fr/sga7ii/current/expose_XIV_body.tex}\par

\par\medskip\noindent\hypertarget{app-F7II-XIV-038.3}{}\textbf{F7II-XIV-038.3}\quad p. 149; inherited corrected-layer difference.\par
\noindent Imprimé : {\ttfamily\footnotesize n\allowbreak{}}\par
\noindent Corrigé : {\ttfamily\footnotesize d\allowbreak{}}\par
\noindent Source : ADOPT-F7II-XIV-038.3. \nolinkurl{fr/sga7ii/diplomatic/expose_XIV_body.tex}; \nolinkurl{fr/sga7ii/current/expose_XIV_body.tex}\par

\par\medskip\noindent\hypertarget{app-F7II-XIV-039.1}{}\textbf{F7II-XIV-039.1}\quad p. 150; misprint\_pass.\par
\noindent Imprimé : {\ttfamily\footnotesize fibré\allowbreak{}}\par
\noindent Corrigé : {\ttfamily\footnotesize fibre\allowbreak{}}\par
\noindent Source : 7II:231. \nolinkurl{fixes/FIXLOG_7II_CURRENT.jsonl#629}; \nolinkurl{fixes/MISPRINTS_7II.json#231}\par

\par\medskip\noindent\hypertarget{app-F7II-XIV-040.1}{}\textbf{F7II-XIV-040.1}\quad p. 150; misprint\_pass.\par
\noindent Imprimé : \emph{(vide)}\par
\noindent Corrigé : {\ttfamily\footnotesize \_c\allowbreak{}}\par
\noindent Source : 7II:232. \nolinkurl{fixes/FIXLOG_7II_CURRENT.jsonl#630}; \nolinkurl{fixes/MISPRINTS_7II.json#232}\par

\par\medskip\noindent\hypertarget{app-F7II-XIV-041.1}{}\textbf{F7II-XIV-041.1}\quad p. 150; misprint\_pass.\par
\noindent Imprimé : {\ttfamily\footnotesize coincide\allowbreak{}}\par
\noindent Corrigé : {\ttfamily\footnotesize coïncide\allowbreak{}}\par
\noindent Source : 7II:233. \nolinkurl{fixes/FIXLOG_7II_CURRENT.jsonl#631}\par

\par\medskip\noindent\hypertarget{app-F7II-XIV-042.1}{}\textbf{F7II-XIV-042.1}\quad p. 152; misprint\_pass.\par
\noindent Imprimé : {\ttfamily\footnotesize \textbackslash{}Omega\textasciicircum{}*\_\{\allowbreak{}}\par
\noindent Corrigé : \emph{(vide)}\par
\noindent Source : 7II:243. \nolinkurl{fixes/FIXLOG_7II_CURRENT.jsonl#632}; \nolinkurl{fixes/MISPRINTS_7II.json#243}\par

\par\medskip\noindent\hypertarget{app-F7II-XIV-042.2}{}\textbf{F7II-XIV-042.2}\quad p. 152; misprint\_pass.\par
\noindent Imprimé : {\ttfamily\footnotesize \}\allowbreak{}}\par
\noindent Corrigé : \emph{(vide)}\par
\noindent Source : 7II:243. \nolinkurl{fixes/FIXLOG_7II_CURRENT.jsonl#632}; \nolinkurl{fixes/MISPRINTS_7II.json#243}\par

\par\medskip\noindent\hypertarget{app-F7II-XIV-043.1}{}\textbf{F7II-XIV-043.1}\quad p. 152; printed misprint.\par
\noindent Imprimé : {\ttfamily\footnotesize 1\allowbreak{}}\par
\noindent Corrigé : {\ttfamily\footnotesize 2\allowbreak{}}\par
\noindent Source : 7II:244. \nolinkurl{fixes/MISPRINTS_7II.json#244}\par

\par\medskip\noindent\hypertarget{app-F7II-XIV-043.2}{}\textbf{F7II-XIV-043.2}\quad p. 152; misprint\_pass.\par
\noindent Imprimé : {\ttfamily\footnotesize rondamental\allowbreak{}}\par
\noindent Corrigé : {\ttfamily\footnotesize fondamental\allowbreak{}}\par
\noindent Source : 7II:245. \nolinkurl{fixes/FIXLOG_7II_CURRENT.jsonl#633}; \nolinkurl{fixes/MISPRINTS_7II.json#245}\par

\par\medskip\noindent\hypertarget{app-F7II-XIV-043.3}{}\textbf{F7II-XIV-043.3}\quad p. 152; misprint\_pass.\par
\noindent Imprimé : {\ttfamily\footnotesize X\allowbreak{}}\par
\noindent Corrigé : {\ttfamily\footnotesize x\allowbreak{}}\par
\noindent Source : 7II:246. \nolinkurl{fixes/FIXLOG_7II_CURRENT.jsonl#634}\par

\par\medskip\noindent\hypertarget{app-F7II-XIV-044.1}{}\textbf{F7II-XIV-044.1}\quad p. 153; misprint\_pass.\par
\noindent Imprimé : \emph{(vide)}\par
\noindent Corrigé : {\ttfamily\footnotesize \textbackslash{}tilde\{\allowbreak{}}\par
\noindent Source : 7II:247. \nolinkurl{fixes/FIXLOG_7II_CURRENT.jsonl#635}; \nolinkurl{fixes/MISPRINTS_7II.json#247}\par

\par\medskip\noindent\hypertarget{app-F7II-XIV-044.2}{}\textbf{F7II-XIV-044.2}\quad p. 153; misprint\_pass.\par
\noindent Imprimé : \emph{(vide)}\par
\noindent Corrigé : {\ttfamily\footnotesize \}\allowbreak{}}\par
\noindent Source : 7II:247. \nolinkurl{fixes/FIXLOG_7II_CURRENT.jsonl#635}; \nolinkurl{fixes/MISPRINTS_7II.json#247}\par

\par\medskip\noindent\hypertarget{app-F7II-XIV-045.1}{}\textbf{F7II-XIV-045.1}\quad p. 153; printed misprint.\par
\noindent Imprimé : {\ttfamily\footnotesize F\allowbreak{}}\par
\noindent Corrigé : {\ttfamily\footnotesize f\allowbreak{}}\par
\noindent Source : 7II:248. \nolinkurl{fixes/MISPRINTS_7II.json#248}\par

\par\medskip\noindent\hypertarget{app-F7II-XIV-046.1}{}\textbf{F7II-XIV-046.1}\quad p. 154; misprint\_pass.\par
\noindent Imprimé : {\ttfamily\footnotesize autmorphisme\allowbreak{}}\par
\noindent Corrigé : {\ttfamily\footnotesize automorphisme\allowbreak{}}\par
\noindent Source : 7II:249. \nolinkurl{fixes/FIXLOG_7II_CURRENT.jsonl#636}; \nolinkurl{fixes/MISPRINTS_7II.json#249}\par

\par\medskip\noindent\hypertarget{app-F7II-XIV-047.1}{}\textbf{F7II-XIV-047.1}\quad p. 154; misprint\_pass.\par
\noindent Imprimé : {\ttfamily\footnotesize Examples\allowbreak{}}\par
\noindent Corrigé : {\ttfamily\footnotesize Exemples\allowbreak{}}\par
\noindent Source : 7II:250. \nolinkurl{fixes/FIXLOG_7II_CURRENT.jsonl#637}\par

\par\medskip\noindent\hypertarget{app-F7II-XIV-047.2}{}\textbf{F7II-XIV-047.2}\quad p. 154; misprint\_in\_print.\par
\noindent Imprimé : \emph{(vide)}\par
\noindent Corrigé : {\ttfamily\footnotesize \textbackslash{}left(\allowbreak{}}\par
\noindent Source : FIX-7II-D-0159. \nolinkurl{fixes/FIXLOG_7II.jsonl#159}\par

\par\medskip\noindent\hypertarget{app-F7II-XIV-047.3}{}\textbf{F7II-XIV-047.3}\quad p. 154; misprint\_in\_print.\par
\noindent Imprimé : \emph{(vide)}\par
\noindent Corrigé : {\ttfamily\footnotesize \textbackslash{}right\allowbreak{}}\par
\noindent Source : FIX-7II-D-0159. \nolinkurl{fixes/FIXLOG_7II.jsonl#159}\par

\par\medskip\noindent\hypertarget{app-F7II-XIV-048.1}{}\textbf{F7II-XIV-048.1}\quad p. 156; misprint\_pass.\par
\noindent Imprimé : {\ttfamily\footnotesize méromorphe\allowbreak{}}\par
\noindent Corrigé : {\ttfamily\footnotesize méromorphes\allowbreak{}}\par
\noindent Source : 7II:234. \nolinkurl{fixes/FIXLOG_7II_CURRENT.jsonl#638}\par

\par\medskip\noindent\hypertarget{app-F7II-XIV-049.1}{}\textbf{F7II-XIV-049.1}\quad p. 157; misprint\_pass.\par
\noindent Imprimé : {\ttfamily\footnotesize muq\allowbreak{}}\par
\noindent Corrigé : {\ttfamily\footnotesize mqu\allowbreak{}}\par
\noindent Source : 7II:235. \nolinkurl{fixes/FIXLOG_7II_CURRENT.jsonl#639}; \nolinkurl{fixes/MISPRINTS_7II.json#235}\par

\par\medskip\noindent\hypertarget{app-F7II-XIV-049.2}{}\textbf{F7II-XIV-049.2}\quad p. 157; misprint\_pass.\par
\noindent Imprimé : {\ttfamily\footnotesize muq\allowbreak{}}\par
\noindent Corrigé : {\ttfamily\footnotesize mqu\allowbreak{}}\par
\noindent Source : 7II:236. \nolinkurl{fixes/FIXLOG_7II_CURRENT.jsonl#640}\par

\par\medskip\noindent\hypertarget{app-F7II-XIV-050.1}{}\textbf{F7II-XIV-050.1}\quad p. 157; misprint\_pass.\par
\noindent Imprimé : {\ttfamily\footnotesize muq\allowbreak{}}\par
\noindent Corrigé : {\ttfamily\footnotesize mqu\allowbreak{}}\par
\noindent Source : 7II:237. \nolinkurl{fixes/FIXLOG_7II_CURRENT.jsonl#641}; \nolinkurl{fixes/MISPRINTS_7II.json#237}\par

\par\medskip\noindent\hypertarget{app-F7II-XIV-050.2}{}\textbf{F7II-XIV-050.2}\quad p. 157; misprint\_pass.\par
\noindent Imprimé : {\ttfamily\footnotesize muq\allowbreak{}}\par
\noindent Corrigé : {\ttfamily\footnotesize mqu\allowbreak{}}\par
\noindent Source : 7II:237. \nolinkurl{fixes/FIXLOG_7II_CURRENT.jsonl#641}; \nolinkurl{fixes/MISPRINTS_7II.json#237}\par

\par\medskip\noindent\hypertarget{app-F7II-XIV-051.1}{}\textbf{F7II-XIV-051.1}\quad p. 158; misprint\_pass.\par
\noindent Imprimé : {\ttfamily\footnotesize soit\allowbreak{}}\par
\noindent Corrigé : {\ttfamily\footnotesize soient\allowbreak{}}\par
\noindent Source : 7II:238. \nolinkurl{fixes/FIXLOG_7II_CURRENT.jsonl#642}; \nolinkurl{fixes/MISPRINTS_7II.json#238}\par

\par\medskip\noindent\hypertarget{app-F7II-XIV-051.2}{}\textbf{F7II-XIV-051.2}\quad p. 158; misprint\_pass.\par
\noindent Imprimé : {\ttfamily\footnotesize horizontale\allowbreak{}}\par
\noindent Corrigé : {\ttfamily\footnotesize horizontales\allowbreak{}}\par
\noindent Source : 7II:238. \nolinkurl{fixes/FIXLOG_7II_CURRENT.jsonl#642}; \nolinkurl{fixes/MISPRINTS_7II.json#238}\par

\par\medskip\noindent\hypertarget{app-F7II-XIV-051.3}{}\textbf{F7II-XIV-051.3}\quad p. 158; misprint\_pass.\par
\noindent Imprimé : \emph{(vide)}\par
\noindent Corrigé : {\ttfamily\footnotesize \textasciicircum{}\textbackslash{}sim\allowbreak{}}\par
\noindent Source : 7II:239. \nolinkurl{fixes/FIXLOG_7II_CURRENT.jsonl#643}; \nolinkurl{fixes/MISPRINTS_7II.json#239}\par

\par\medskip\noindent\hypertarget{app-F7II-XIV-052.1}{}\textbf{F7II-XIV-052.1}\quad p. 159; printed misprint.\par
\noindent Imprimé : \emph{(vide)}\par
\noindent Corrigé : {\ttfamily\footnotesize \textasciicircum{}*\allowbreak{}}\par
\noindent Source : 7II:240. \nolinkurl{fixes/MISPRINTS_7II.json#240}\par

\par\medskip\noindent\hypertarget{app-F7II-XIV-053.1}{}\textbf{F7II-XIV-053.1}\quad p. 160; printed misprint.\par
\noindent Imprimé : {\ttfamily\footnotesize \textbackslash{}mathbb\{C\}\allowbreak{}}\par
\noindent Corrigé : {\ttfamily\footnotesize V\allowbreak{}}\par
\noindent Source : 7II:241. \nolinkurl{fixes/MISPRINTS_7II.json#241}\par

\par\medskip\noindent\hypertarget{app-F7II-XIV-054.1}{}\textbf{F7II-XIV-054.1}\quad p. 160; misprint\_pass.\par
\noindent Imprimé : {\ttfamily\footnotesize complexe\allowbreak{}}\par
\noindent Corrigé : {\ttfamily\footnotesize complexes\allowbreak{}}\par
\noindent Source : 7II:242. \nolinkurl{fixes/FIXLOG_7II_CURRENT.jsonl#644}\par

\par\medskip\noindent\hypertarget{app-F7II-XIV-055.1}{}\textbf{F7II-XIV-055.1}\quad p. 161; printed misprint.\par
\noindent Imprimé : {\ttfamily\footnotesize à\allowbreak{}}\par
\noindent Corrigé : {\ttfamily\footnotesize a\allowbreak{}}\par
\noindent Source : 7II:253. \nolinkurl{fixes/MISPRINTS_7II.json#253}\par

\par\medskip\noindent\hypertarget{app-F7II-XIV-055.2}{}\textbf{F7II-XIV-055.2}\quad p. 161; misprint\_pass.\par
\noindent Imprimé : {\ttfamily\footnotesize paramétrée\allowbreak{}}\par
\noindent Corrigé : {\ttfamily\footnotesize paramétrées\allowbreak{}}\par
\noindent Source : 7II:254. \nolinkurl{fixes/FIXLOG_7II_CURRENT.jsonl#646}; \nolinkurl{fixes/MISPRINTS_7II.json#254}\par

\par\medskip\noindent\hypertarget{app-F7II-XIV-055.3}{}\textbf{F7II-XIV-055.3}\quad p. 161; misprint\_pass.\par
\noindent Imprimé : {\ttfamily\footnotesize :\allowbreak{}}\par
\noindent Corrigé : {\ttfamily\footnotesize \textbackslash{}cdot\allowbreak{}}\par
\noindent Source : 7II:255. \nolinkurl{fixes/FIXLOG_7II_CURRENT.jsonl#647}; \nolinkurl{fixes/MISPRINTS_7II.json#255}\par

\par\medskip\noindent\hypertarget{app-F7II-XIV-055.4}{}\textbf{F7II-XIV-055.4}\quad p. 161; printed misprint.\par
\noindent Imprimé : {\ttfamily\footnotesize 4\allowbreak{}}\par
\noindent Corrigé : {\ttfamily\footnotesize 5\allowbreak{}}\par
\noindent Source : 7II:256. \nolinkurl{fixes/MISPRINTS_7II.json#256}\par

\par\medskip\noindent\hypertarget{app-F7II-XIV-056.1}{}\textbf{F7II-XIV-056.1}\quad p. 162; misprint\_pass.\par
\noindent Imprimé : {\ttfamily\footnotesize méromorphe\allowbreak{}}\par
\noindent Corrigé : {\ttfamily\footnotesize méromorphes\allowbreak{}}\par
\noindent Source : 7II:257. \nolinkurl{fixes/FIXLOG_7II_CURRENT.jsonl#648}; \nolinkurl{fixes/MISPRINTS_7II.json#257}\par

\par\medskip\noindent\hypertarget{app-F7II-XIV-057.1}{}\textbf{F7II-XIV-057.1}\quad p. 163; misprint\_pass.\par
\noindent Imprimé : \emph{(vide)}\par
\noindent Corrigé : {\ttfamily\footnotesize *\allowbreak{} u\textasciicircum{}\allowbreak{}}\par
\noindent Source : 7II:258. \nolinkurl{fixes/FIXLOG_7II_CURRENT.jsonl#649}; \nolinkurl{fixes/MISPRINTS_7II.json#258}\par

\normalsize\fi
\clearpage\SGAtoc{Exposé XV — La formule de Picard-Lefschetz}
\selectlanguage{french}

\SGApage{II}{165}


SGA 7


\underline{Exposé XV}


\underline{\underline{La formule de Picard-Lefschetz}}


\underline{par P. DELIGNE}

\S1. \hspace{1em} Singularités quadratiques ordinaires (formes canoniques) \hfill 2

\hspace*{2em}1.1 \hspace{0.5em} Rappels \hfill 2

\hspace*{2em}1.2 \hspace{0.5em} Singularités quadratiques ordinaires \hfill 4

\hspace*{2em}1.3 \hspace{0.5em} Modules de singularités quadratiques \hfill 10

\S2 \hspace{1em} Calcul des cycles évanescents dans un cas quadratique ordinaire standard \hfill 13

\hspace*{2em}2.1 \hspace{0.5em} Préliminaire : la cohomologie d'un cône \hfill 13

\hspace*{2em}2.2 \hspace{0.5em} Cycles évanescents \hfill 17

\S3 \hspace{1em} La formule de Picard-Lefschetz \hfill 23

\hspace*{2em}3.1 \hspace{0.5em} Résultats préliminaires \hfill 23

\hspace*{2em}3.2 \hspace{0.5em} Dimension relative paire \hfill 24

\hspace*{2em}3.3 \hspace{0.5em} Dimension relative impaire \hfill 26

\hspace*{2em}3.4 \hspace{0.5em} Résumé \hfill 31


\SGApage{II}{166}


\SGAsectiontoc{\S1  Singularités quadratiques ordinaires (formes canoniques)}{1}
\S1  \underline{Singularités quadratiques ordinaires (formes canoniques)}

\SGAvar{Dans ce \S, on reprend les résultats de VI 6, en adoptant un point de vue "hensélien" plutôt que "formel". On donne des singularités quadratiques ordinaires et de leurs déformations des formes canoniques, sous lcsquelles elles se mettent localement pour la topologie étale.

}{Dans ce \S, on reprend les résultats de VI 6, en adoptant un point de vue "hensélien" plutôt que "formel". On donne des singularités quadratiques ordinaires et de leurs déformations des formes canoniques, sous lesquelles elles se mettent localement pour la topologie étale.

}{F7II-XV-001}Les résultats sont conséquence facile de VI 6, du théorème des fonctions implicites (sous la forme de Tougeron) et du théorème de R. Elkik sur les modules henséliens versels de singularités isolées. Lors du séminaire oral, on ne disposait pas encore de ce dernier; un théorème plus faible, ne couvrant que le cas (qui nous suffisait) des intersections complètes, avait été démontré et ne sera pas publié.


1.1.  \underline{Rappels}

\SGAvar{1.1.1  Soit $f : X \to S$ un morphisme, $X$ pouvant localement être décrit par $p$ équations dans l'espace affine relatif $A_S^{n+p}$ de dimension $n+p$ sur $S$ . L'idéal jacobien $J^n(X/S)$ a été défini en VI 5. Si $X$ est le sous-schéma de $A_S^{n+p}$ défini par les équations $f_i = 0$ $(1 < i \leq p)$ , c'est le faisceau d'idéaux engendré par les mineurs $p \times p$ de la matrice jacobienne $(\partial f_i/\partial x_j)$ . Le théorème suivant est bien connu. Un cas particulier de son analogue formel est démontré dans Bourbaki Alg. Comm. III \S4, n$^{\circ}$ 6, cor. 3 au Th. 2; l'analogue $C^{\infty}$ est démontré dans Tougeron (idéaux de fonctions differentiables - Thèse - mai 1967 - Rennes); le cas considéré ici est traité dans Artin, Algebraic approximation of structures over complete local rings, Publ. Math. IHES \underline{36} (Zariski's volume) lemme 5.10.

}{1.1.1  Soit $f : X \to S$ un morphisme, $X$ pouvant localement être décrit par $p$ équations dans l'espace affine relatif $A_S^{n+p}$ de dimension $n+p$ sur $S$ . L'idéal jacobien $J^n(X/S)$ a été défini en VI 5. Si $X$ est le sous-schéma de $A_S^{n+p}$ défini par les équations $f_i = 0$ $(1 \leq i \leq p)$ , c'est le faisceau d'idéaux engendré par les mineurs $p \times p$ de la matrice jacobienne $(\partial f_i/\partial x_j)$ . Le théorème suivant est bien connu. Un cas particulier de son analogue formel est démontré dans Bourbaki Alg. Comm. III \S4, n$^{\circ}$ 6, cor. 3 au Th. 2; l'analogue $C^{\infty}$ est démontré dans Tougeron (idéaux de fonctions différentiables - Thèse - mai 1967 - Rennes); le cas considéré ici est traité dans Artin, Algebraic approximation of structures over complete local rings, Publ. Math. IHES \underline{36} (Zariski's volume) lemme 5.10.

}{F7II-XV-002}1.1.2  \underline{Théorème} \underline{(Tougeron-Artin)}. \underline{Soient} $S$ \underline{le spectre d'un anneau local hensélien} $A$ , $f : X \to S$ \underline{un morphisme comme en 1.1.1.,} $a$ \underline{et} $I$ \underline{deux idéaux de} $A$ \underline{distincts de} $A$ , \underline{et} $\delta'$ \underline{un faisceau d'idéaux de} $X$ \underline{contenu dans} $J^n(X/S)$ . \underline{On suppose} $I$ \underline{de type fini. Soit} $s$ \underline{une section de} $f$ \underline{au-dessus de} $S/I$ : $s \in X(A/I)$ . \underline{Notons} $s^{*}\delta'^{2}$ \underline{l'idéal de} $A$ \underline{image réciproque dans} $A$ \underline{de l'idéal} $s^{*}\delta'^{2}$ de $A/I$ .

\underline{Si} $(s^{*}\delta'^{2}).a \supset I$ , \underline{alors, il existe une section} $s_1$ \underline{de} $f : s_1 \in X(A)$
\SGApage{II}{167}
\underline{congrue à} $s$ \underline{modulo} $(s^{*}\delta').a$ .

1.1.3.  Soit $S$ le spectre d'un anneau local hensélien, de point fermé $s$ . Soient $X'_o$ un schéma de type fini sur $k(s)$ , lisse en dehors d'un point fermé $x_o$ , et $X_o$ l'hensélisé de $X'_o$ en $x_o$ . Nous nous intéressons aux systèmes

\begin{center}
\begin{tikzcd}
X_o \arrow[r, hook] \arrow[d] \& X \arrow[d] \& {} \\
s \arrow[r, hook] \& T \arrow[r] \& S
\end{tikzcd}
\end{center}







du type suivant

(a)  $T$ est local hensélien essentiellement de présentation finie sur $S$(I 0.0.5); le morphisme $T \to S$ est local et induit un isomorphisme sur les corps résiduels;

(b)  $X$ est local hensélien plat et essentiellement de présentation finie sur $T$ ;

(c)  la fibre spéciale de $X/T$ est munie d'un $k(s)$ - isomorphisme avec $X_o$ .

Ce sont les $S$-\underline{déformations henséliennes} de $X_o$ . Un \underline{morphisme} de $(X,T)$ dans $(X',T')$ est un diagramme commutatif de $S$-schémas

\begin{center}
\begin{tikzcd}
X \arrow[r] \arrow[d] \& X' \arrow[d] \\
T \arrow[r] \& T'
\end{tikzcd}
\end{center}





induisant l'identité de $X_o$ par passage aux fibres spéciales.Une déformation $X/T$ est dite \underline{verselle} (minimally versal) si

(a)  Pour tout morphisme $X''/T'' \longrightarrow X'/T'$, tel que $T'' \to T'$ soit un plongement fermé, l'application

\begin{equation*}
\mathrm{Hom}(X'/T' , X/T) \longrightarrow \mathrm{Hom}(X''/T'' , X/T)
\end{equation*}
est surjective. En particulier (faire $X''/T'' = X_o/s$) , $\mathrm{Hom}(X'/T', X/T) \neq \emptyset$ .

(b)  Pour $T' = \mathrm{Spec}(k(s)[\varepsilon]/(\varepsilon^2))$ , l'application "image réciproque" : $\mathrm{Hom}(T',T) \to$ (classes d'isomorphie de déformations de $X_o$ sur $T'$) est bijective.


\SGApage{II}{168}

\SGAvar{1.1.4. \underline{Théorème (R. Elkik).} \underline{Il existe une} $S$-\underline{déformation hensélienne verselle de} $X_o$ \underline{(et une seule à isomorphisme non unique près). Si} $S$ \underline{est noethérien, pour que} $X/T$ \underline{en soit une, il faut et il suffit} $X$ \underline{induise sur le complété} $\widehat{T}$ \underline{de} $T$ \underline{une} $S$-\underline{déformation formelle verselle de} $X_o$ \underline{(minimally formal versal deformation, VI 1.9 et VI 4).}

}{1.1.4. \underline{Théorème (R. Elkik).} \underline{Il existe une} $S$-\underline{déformation hensélienne verselle de} $X_o$ \underline{(et une seule à isomorphisme non unique près). Si} $S$ \underline{est noethérien, pour que} $X/T$ \underline{en soit une, il faut et il suffit que} $X$ \underline{induise sur le complété} $\widehat{T}$ \underline{de} $T$ \underline{une} $S$-\underline{déformation formelle verselle de} $X_o$ \underline{(minimally formal versal deformation, VI 1.9 et VI 4).}

}{F7II-XV-003}Pour la démonstration de ce théorème délicat, on renvoie à R. Elkik (séminaire de l'ENS, 1971/72).


1.2. \underline{Singularités quadratiques ordinaires}

Soit $y$ un point fermé d'un schéma $Y$ de type fini sur un corps $k$ de caractéristique $p$ , et soit $n$ la dimension de $Y$ en $y$ .


\underline{Définition} 1.2.1 (= VI 6.6)

\SGAvar{(1) \underline{Pour} $k$ \underline{algébriquement clos, on dit que} $y$ \underline{est un point} quadratique ordinaire \underline{de} $Y$ \underline{si le complété} $\widehat{\mathcal{O}}_{Y,y}$ \underline{de l'anneau local de} $Y$ \underline{en} $y$ \underline{est isomorphe au quotient de} $k[[x_1,\ldots x_{n+1}]]$ \underline{par l'idéal engendré par une seule série formelle} $f$ , \underline{qui s'écrive}

\begin{equation}\tag{4.2.1}
f(\underline{X}) = Q(\underline{X}) = + \mbox{ termes d'ordre } > 2
\end{equation}
\underline{avec} $Q(\underline{X})$ \underline{une forme quadratique ordinaire.}

(ii) \underline{Si} $k$ \underline{est une clôture algébrique de} $k$ , \underline{on dit que} $y$ \underline{est un} point quadratique ordinaire \underline{de} $Y$ \underline{si les points de} $\overline{Y} = Y \otimes_k \overline{k}$ \underline{au-dessus de} $y$ \underline{sont des points quadratiques ordinaires de} $\overline{Y}$ .

}{(i) \underline{Pour} $k$ \underline{algébriquement clos, on dit que} $y$ \underline{est un point} quadratique ordinaire \underline{de} $Y$ \underline{si le complété} $\widehat{\mathcal{O}}_{Y,y}$ \underline{de l'anneau local de} $Y$ \underline{en} $y$ \underline{est isomorphe au quotient de} $k[[x_1,\ldots x_{n+1}]]$ \underline{par l'idéal engendré par une seule série formelle} $f$ , \underline{qui s'écrive}

\begin{equation}\tag{1.2.1.1}
f(\underline{X}) = Q(\underline{X}) + \mbox{ termes d'ordre } > 2
\end{equation}
\underline{avec} $Q(\underline{X})$ \underline{une forme quadratique ordinaire.}

(ii) \underline{Si} $\overline{k}$ \underline{est une clôture algébrique de} $k$ , \underline{on dit que} $y$ \underline{est un} point quadratique ordinaire \underline{de} $Y$ \underline{si les points de} $\overline{Y} = Y \otimes_k \overline{k}$ \underline{au-dessus de} $y$ \underline{sont des points quadratiques ordinaires de} $\overline{Y}$ .

}{F7II-XV-004}1.2.2. Si dans cette définition on remplace ``forme quadratique ordinaire'' par ``forme quadratique non dégénérée'', on obtient la notion de \underline{point quadratique non dégénéré} (VI 6.0). Pour que $y$ soit un point quadratique non dégénéré de $Y$ , il faut et il suffit qu'il soit un point quadratique ordinaire et que soit $p \neq 2$ ,
\SGApage{II}{169}
soit $n$ est impair.

\underline{Exemple} 1.2.3. Soit $Q(\underline{X})$ une forme quadratique non dégénérée à $n+1$ variables sur $k$ . Soit $Y_o$ le cône quadratique dans $\mathbb{E}^{n+1}_k$ d'équation $Q = 0$ . Alors, l'origine, qu'on notera $y_o$ , est un point quadratique non dégénéré de $Y_o$ .

\underline{Exemple} 1.2.4. Supposons que $p = 2$ et que $n = 2m$ . Soit $a \in k$ , et soit $Q$ un polynôme

\begin{equation*}
Q(\underline{X}) = (X_o^2 - a) + \sum_{0<i\leq j\leq 2m} a_{ij}x_ix_j \quad .
\end{equation*}
On suppose que la forme quadrique à $2m$ variables $\sum_{0<i\leq j} a_{ij}x_ix_j$ est non dégénérée, et on désigne par $y_o$ le point fermé de $\mathbb{E}^{n+1}_k$ de coordonnées $(\sqrt{a},0,\ldots,0)$ . Alors, $y_o$ est un point quadratique ordinaire du sous-schéma $Y_o$ de $\mathbb{E}^{n+1}_k$ d'équation $Q = 0$.

\underline{Exemple} 1.2.5. Soit $k'$ une extension finie séparable de $k$ et soit $Y_o$ un $k'$-schéma du type considéré en 1.2.3 (resp. 1.2.4). Si l'on regarde le schéma $Y_o$ comme un $k$-schéma, alors $y_o$ est un point quadratique non dégénéré (resp. ordinaire) de $Y_o$ . L'extension $k'$ de $k$ s'identifie à la plus grande sous-extension séparable de $k(y_o)$ .

\underline{Théorème} 1.2.6. \underline{Supposons que} $y$ \underline{soit un point quadratique ordinaire du} $k$-\underline{schéma} $Y$ , \underline{et soit} $k'$ \underline{la plus grande sous-extension séparable de} $k(y)$ . \underline{Il existe alors un} $k'$-\underline{schéma} $Y_o \subset \mathbb{E}^{n+1}_{k'}$ \underline{du type 1.2.3. ou 1.2.4., et un} $k$-\underline{isomorphisme} $\varphi$ \underline{entre} \underline{l'hensélisé} $Y_{(y)}$ \underline{de} $Y$ \underline{en} $y$ \underline{et l'hensélisé} $Y_{o(y_o)}$ \underline{de} $Y_o$ \underline{en} $y_o$ .

Notons tout d'abord qu'il existe un voisinage étale $Y'$ de $y$ dans $Y$

\begin{center}
\begin{tikzcd}[row sep=large, column sep=large]
{} \& Y' \arrow[d] \\
y \arrow[ur, hook] \arrow[r, hook] \& Y
\end{tikzcd}
\end{center}






qui soit un $k'$-schéma. Le schéma $Y_{(y)}$ est en effet canoniquement un $k'$-schéma.


\SGApage{II}{170}

Remplaçant $(k,Y,y)$ par $(k',Y',y)$ , ceci nous ramène au cas où $k = k'$ , ie. au cas où $k(y)$ est radiciel sur $k$ .

Le $k(y)$-vectoriel $(\Omega^1_{Y/k})_y$ est de dimension $n+1$ : il suffit de le vérifier après extension des scalaires de $k(y)$ à sa clôture algébrique, et après complétion en $y$ ; c'est alors évident. Un voisinage de $y$ dans $Y$ est donc un sous-schéma d'un $k$-schéma lisse $Z$ de dimension $n+1$ . De plus, quitte à remplacer $Y$ et $Z$ par des voisinages convenables de $y$ , on a que l'idéal de $\mathcal{O}_Z$ qui définit $Y$ est principal (même chanson que plus haut), défini par une équation $f$ .

\underline{Lemme} 1.2.7. \underline{Supposons que} $y$ \underline{soit un point quadratique non dégénéré de} $Y$ \underline{et que} $k(y)$ \underline{soit radiciel sur} $k$ . \underline{Alors, au voisinage de} $y$ , $J^n_{Y/k}$ (1.1.1) est \underline{monogène, concentré en} $y$ \underline{et de rang} 1 \underline{sur} $k$ . \underline{En particulier} $k(y) = k$ .

Il suffit de prouver la première assertion après passage à la clôture algébrique de $k$ , et après complétion en $y$ . Elle est alors évidente par calcul direct: avec les notations de 1.2.1, il résulte de Nakayama que les $\partial f/\partial x_i$ (a fortiori $f$ et les $\partial f/\partial x_i$) engendrent l'idéal maximal.

\SGAvar{Prouvons 1.2.6 pour $y$ non dégénéré. Nous pouvons supposer et supposerons que $k(y) = k$ , que $Y$ est défini par une équation $f$ dans un schéma lisse $Z$ , et qu'il existe un morphisme étale $x : Z \longrightarrow \mathbb{E}^{n+1}_k$ de $Z$ dans $\mathbb{E}^{n+1}_k$ (un ``système de coordonnées locales $x_i$'') qui envoye $y$ sur $0$ . Il existe alors une forme quadratique non dégénérée à $n+1$ variables $Q$ telle que, en $y$ ,

\begin{equation*}
f = Q(x_i) + \text{termes d'ordre} > 2 \;,
\end{equation*}
Soit $Y_o$ défini par $Q$ comme en 4.3.

}{Prouvons 1.2.6 pour $y$ non dégénéré. Nous pouvons supposer et supposerons que $k(y) = k$ , que $Y$ est défini par une équation $f$ dans un schéma lisse $Z$ , et qu'il existe un morphisme étale $x : Z \longrightarrow \mathbb{E}^{n+1}_k$ de $Z$ dans $\mathbb{E}^{n+1}_k$ (un ``système de coordonnées locales $x_i$'') qui envoie $y$ sur $0$ . Il existe alors une forme quadratique non dégénérée à $n+1$ variables $Q$ telle que, en $y$ ,

\begin{equation*}
f = Q(x_i) + \text{termes d'ordre} > 2 \;,
\end{equation*}
Soit $Y_o$ défini par $Q$ comme en 1.2.3.

}{F7II-XV-005}\underline{Sur} $Y$ , le système des $x_i$ est une ``solution approchée'' de l'équation $Q(x_i) = 0$ : si $m$ est l'idéal qui définit $y$ , on a

\begin{equation*}
Q(x_i) \in m^3 \qquad (\text{sur } Y) \;.
\end{equation*}

\SGApage{II}{171}

Puisque $Q$ est non dégénérée, l'idéal de $\mathbb{E}^{n+1}_k$ engendré par les $\frac{\partial Q}{\partial x_i}$ est l'idéal maximal en $0$ , et l'image réciproque de cet idéal par $x = (x_i) : Y \longrightarrow \mathbb{E}^{n+1}_k$ est $m$ . D'après le théorème des fonctions implicites 1.1.2. (appliqué avec $a = m$ et $\delta' = J$) il existe sur l'hensélisé $Y_{(y)}$ des fonctions $x'_i$ , congruents aux $x_i$ mod $m^2$ (le $s^*\delta'_S.a$ de loc. cit.), qui vérifient $Q(x'_i) = 0$ . Elles fournissent l'isomorphisme cherché entre $Y_{(y)}$ et $Y_{o(y_o)}$ .

\underline{Lemme} 1.2.8. \underline{Supposons que} $y$ \underline{soit un point quadratique ordinaire dégénéré} ($p = 2$, $n$ \underline{pair}), \underline{et que} $k(y)$ \underline{soit radiciel sur} $k$ . \underline{Alors, au voisinage de} $y$ , $T^1_{Y/k}$ \underline{est monogène, concentré en} $y$ \underline{et de rang} 2 \underline{sur} $k$ . \underline{En particulier, soit} $k(y) = k$ , \underline{soit} $k(y)$ \underline{est isomorphe à} $k(\sqrt{a})$ \underline{avec} $a \in k$ , $a \notin k^2$ .

Procédons comme en 1.2.7 : avec les notations de 1.2.1, il suffit de vérifier que

\SGAvar{\begin{equation*}
\dim_k(k[[X_1,\ldots,X_{n+1}]] \;/\; (f,\frac{\partial f}{\partial x_i}))) = 2 \;.
\end{equation*}
Dans un système de coordonnées convenable $x_o$ , $x_1,\ldots,x_{2m}$ , la série formelle $f$ s'écrit

}{\begin{equation*}
\dim_k(k[[X_1,\ldots,X_{n+1}]] \;/\; (f,\frac{\partial f}{\partial x_i})) = 2 \;.
\end{equation*}
Dans un système de coordonnées convenable $x_o$ , $x_1,\ldots,x_{2m}$ , la série formelle $f$ s'écrit

}{F7II-XV-006}\begin{equation*}
f = x_o^2 + \sum_1^m x_i\, x_{i+m} + R \qquad\qquad (\text{ordre } R \geq 3)
\end{equation*}
Soient alors $n'$ l'idéal engendré par $f$ et les $\partial f/\partial x_i$ , $n$ celui engendré par $x_o^2$ et les $x_i(i \neq 0)$ et $q$ l'idéal maximal. Prouvons que $n = n'$ . On a clairement $n' \subset n$ , et il suffit de vérifier que $n' + nq = n$ (Nakayama). Modulo $nq$ , on a $f \equiv x_o^2$ ; modulo $nq + (x_o^2)$ , on a $\partial f/\partial x_i \equiv x_{i+m}$ et $\partial f/\partial x_{i+m} \equiv x_i$ ($1 \leq i \leq m$), d'où l'assertion.

Sous les hypothèses de 1.2.8, si $k = k(y)$ , on désignera par $a$ un quelconque carré dans $k$ .


\SGApage{II}{172}

\underline{Lemme} 1.2.9. \underline{Soient} $V$ \underline{un fibré vectoriel sur un schéma} $S$ , \underline{et} $T$ \underline{un sous-schéma de} $\mathbb{E}_S(V)$ \underline{localement libre de rang} 2 \underline{sur} $S$ . \underline{Alors,} $T$ \underline{est contenu dans une et une seule droite de} $\mathbb{E}_S(V)$ .

Le problème est local (fppf) sur $S$ ; on peut donc supposer que $T/S$ admet une section. Par translation, on peut supposer que $T$ contient la section $0$ . Soit $n$ l'idéal de $T$ qui définit $0$ : $n$ est localement libre de rang un sur $S$ . Chaque forme linéaire homogène $v \in V$ sur $\mathbb{E}_S(V)$ définit un élément de $n$ , d'où un épimorphisme de $V$ dans $n$ . Son noyau définit la droite cherchée.

1.2.10. En particulier, si $A$ est un sous-schéma radiciel de rang 2 de l'espace affine $\mathbb{E}_k^{n+1}$ sur un corps $k$ de caractéristique 2, il existe $a \in k$ tel que, dans un système de coordonnées affines convenables, $A$ ait pour équations $x_o^2 - a = 0$ , $x_i = 0$ $(1 \leq i \leq n)$ . Il suffit de prendre les $x_i$ s'annulant sur la droite définie par $A$ .

\SGAvar{1.2.11. Prouvons 1.2.6 pour $y$ dégénéré ($p = 2$ et $n = 2m$) . nous pouvons supposer et supposerons que $k(y) = k(\sqrt{a})$ , que $Y$ est défini par une équation $f$ dans un schéma lisse $Z$ , et qu'il existe un système de coordonnées locales en $y$ $x = (x_i)_{0 \leq i \leq 2m} : Z \longrightarrow \mathbb{E}_k^{2m+1}$ qui envoye $y$ en le point fermé $(\sqrt{a},0,\ldots,0)$ , et tel que $J^n_{Y/k}$ soit isomorphe, au voisinage de $y$ , à $\mathcal{O}_Y /(x_o^2 - a , x_1,\ldots,x_{2m})$ .

}{1.2.11. Prouvons 1.2.6 pour $y$ dégénéré ($p = 2$ et $n = 2m$) . Nous pouvons supposer et supposerons que $k(y) = k(\sqrt{a})$ , que $Y$ est défini par une équation $f$ dans un schéma lisse $Z$ , et qu'il existe un système de coordonnées locales en $y$ $x = (x_i)_{0 \leq i \leq 2m} : Z \longrightarrow \mathbb{E}_k^{2m+1}$ qui envoie $y$ en le point fermé $(\sqrt{a},0,\ldots,0)$ , et tel que $J^n_{Y/k}$ soit isomorphe, au voisinage de $y$ , à $\mathcal{O}_Y /(x_o^2 - a , x_1,\ldots,x_{2m})$ .

}{F7II-XV-007}On a alors, quitte à remplacer $f$ par un multiple

\begin{equation*}
f = (x_o^2 - a) + \sum_{1 \leq i \leq j \leq 2m} a_{ij} \, x_i x_j + R \ , \qquad \mbox{où}
\end{equation*}
(a) la forme à $2m$ variables $\sum a_{ij} \, x_i x_j$ est non dégénérée;

(b) après extension des scalaires à $\overline{k}$ , $R$ est d'ordre $\geq 3$ en $x_o - \sqrt{a}$ et les $x_i$ .


\SGApage{II}{173}

Soit $Y_o$ défini par

\begin{equation*}
Q = x_o^2 - a + \sum a_{ij} \, x_i x_j
\end{equation*}
comme en 1.2.4.

\underline{Sur} $Y$ , on peut encore regarder le système des $x_i$ comme une solution approchée de l'équation $Q(X_i) = 0$ . Cette fois, l'idéal $\delta$ image réciproque par $\chi$ de l'idéal engendré par les $\partial Q/\partial X_i$ vérifie $\mathcal{O}_Y/\delta \simeq J^n_{Y/k}$ , ie est

\begin{equation*}
\delta = (x_o^2 - a \ , \ (x_i)_{1 \, \leq \, i}) \qquad \mbox{en } \ y \ .
\end{equation*}
Soit $q$ l'idéal maximal qui définit $y$ . Pour conclure comme dans le cas non dégénéré, par une application du théorème des fonctions implicites, il faut vérifier que

\begin{equation}\tag{1.2.11.1}
Q(x_i) \equiv 0 \qquad (\mbox{mod } \delta^2 q) \qquad \underline{\mbox{sur}} \ \ Y \ .
\end{equation}
Soit $\overline{q}$ l'idéal qui définit $y$ , après extension des scalaires à la clôture algébrique de $k$ . On vérifie que

\SGAvar{\begin{equation*}
\delta^2 \overline{q} \cap \mathcal{O}_{Y,y} = \delta^2 q \qquad ,
\end{equation*}
de sorte qu'il suffit de vérifier 4.11.1 après extension des scalaires à $\overline{k}$ : on peut supposer que $a = 0$ . En $y$ , on a alors

}{\begin{equation*}
\delta^2 \overline{q} \cap \mathcal{O}_{Y,y} = \delta^2 q \qquad ,
\end{equation*}
de sorte qu'il suffit de vérifier 1.2.11.1 après extension des scalaires à $\overline{k}$ : on peut supposer que $a = 0$ . En $y$ , on a alors

}{F7II-XV-008}\begin{equation*}
\delta^2 q = (x_o^5 \ , (x_o^3 \, x_i)_{1 \leq i} \ , \ (x_o x_i x_j)_{1 \leq i,j} \ , \ (x_i x_j x_k)_{1 \leq i,j,k}) \ .
\end{equation*}
De là, on tire que $\delta^2 q \supset q^5$ et que

\begin{equation*}
x_o^2 f \equiv x_o^4 \qquad (\mbox{mod } \delta^2 q) \ , \quad \mbox{d'où} \quad \delta^2 q + (x_o^2 f) \supset q^4 \ ,
\end{equation*}
\begin{equation*}
x_o f \equiv x_o^3 \qquad (\mbox{mod } \delta^2 q + q^4) \ , \quad \mbox{d'où} \quad \delta^2 q + (x_o f) \supset q^3 \ .
\end{equation*}

\SGApage{II}{174}

On a enfin

\begin{equation*}
Q(\chi_i) \equiv 0 \qquad (\delta^2 q + (f)) \qquad \underline{\mathrm{sur}} \quad Z \quad ,
\end{equation*}
ie 1.2.11.1.

\underline{Remarque} 1.2.12.  Soit $Q$ une forme quadratique non nécessairement homogène à $n+1$ variables, non identiquement nulle, et supposons que le sous-schéma $Y$ de $\mathbb{E}_k^{n+1}$ d'équation $Q$ présente en un point $y$ une singularité quadratique ordinaire. Par un changement linéaire (non homogène) de variables, on peut alors mettre $Y$ sous la forme 1.2.3 ou 1.2.4. Dans le cas non dégénéré, $k(y) = k$ , et il suffit de mettre $y$ à l'origine des coordonnées. Dans le cas dégénéré, il suffit que le support de $T^1_{Y/k}$ soit sur la droite $x_i = 0$  ($1 \leq i \leq 2m$)  (cf. 1.2.9.).


1.3  \underline{Modules de singularités quadratiques}

La proposition suivante résulte du théorème de R. Elkik 1.1.4. et de calculs explicites laissés au lecteur (essentiellement contenus dans VI 6).

\underline{Proposition} 1.3.1.  \underline{Soit} $S = \mathrm{Spec}(A)$ \underline{un schéma hensélien local de point fermé} $s$ , \underline{et soit} $Y_o \subset \mathbb{E}^{n+1}_{k(s)}$ \underline{un schéma du type 1.2.3. ou 1.2.4.}

(i) \underline{Prenons} $Y_o$ \underline{de type 1.2.3, défini par l'équation}

\begin{equation*}
Q(\underline{X}) = \sum_{i \leq j} \overline{a}_{ij} x_i x_j = 0 \qquad (\overline{a}_{ij} \in k(s))
\end{equation*}
\underline{et relevons les} $\overline{a}_{ij}$ \underline{en des éléments} $a_{ij}$ \underline{de} $A$ . \underline{Alors, le schéma hensélien versel sur} $S$ \underline{des déformations de la singularité de} $Y_o$ \underline{en} $y_o$ \underline{est}

\begin{equation*}
T = \mathrm{Spec}(A\{b\}) \quad ,
\end{equation*}
\underline{la déformation verselle étant l'hensélisé en} $y_o$ \underline{du} $T$-\underline{schéma} $Y \subset \mathbb{E}_T^{n+1}$ \underline{d'équation} 


\SGApage{II}{175}

\begin{equation*}
 \sum a_{ij} x_i x_j - b = 0 \quad .
\end{equation*}
(ii) \underline{Prenons} $Y_o$ \underline{de type 1.2.4, défini par l'équation}

\begin{equation*}
Q(\underline{x}) = (x_o^2 - \overline{a}) + \sum \overline{a}_{ij} x_i x_j = 0 \qquad (\overline{a},\overline{a}_{ij} \in k(s))
\end{equation*}
\underline{et relevons les} $\overline{a}_{ij}$ \underline{et} $\overline{a}$ \underline{en des éléments} $a_{ij}$ \underline{et} $a$ \underline{de} $A$ . \underline{Alors, le schéma hensélien versel sur} $S$ \underline{des déformations de la singularité de} $Y_o$ \underline{en} $y_o$ \underline{est}

\begin{equation*}
T = \mathrm{Spec}(A\{b,c\}) \quad ,
\end{equation*}
\underline{la déformation verselle étant l'hensélisé en} $y_o$ \underline{du} $T$- \underline{schéma} $Y \subset \mathbb{E}_T^{n+1}$ \underline{d'équation}

\begin{equation*}
(x_o^2 - a) + \sum a_{ij} x_i x_j + b x_o + c = 0
\end{equation*}
La proposition suivante résulte aussitôt de 1.2.6 et 1.3.1.

\underline{Corollaire} 1.3.2.  \underline{Soit} $S = \mathrm{Spec}(A)$ \underline{un schéma hensélien local de point fermé} $s$ \underline{et soit} $f : X \to S$ \underline{un morphisme plat de présentation finie. On suppose qu'en le point fermé} $x$ \underline{de} $X_s$ , $X_s$ \underline{présente une singularité quadratique ordinaire, que} $k(x)$ \underline{est radiciel sur} $k(s)$ , \underline{et que} $X_s$ \underline{est de dimension} $n$ .

(i) Pour $x$ non dégénéré : \underline{il existe une forme quadrique non dégénérée à} $n+1$ \underline{variables} $Q$ , \underline{à coefficients dans} $A$ , \underline{et} $b$ \underline{dans l'idéal maximal de} $A$ , \underline{tel que le} $S$-\underline{schéma hensélisé de} $X$ \underline{en} $x$ \underline{soit isomorphe à l'hensélisé en l'origine du sous-schéma de} $\mathbb{E}_S^{n+1}$ \underline{d'équation (1.3.1(i))} $Q - b = 0$ .

(ii) Pour $x$ dégénéré ($n = 2m$, $\mathrm{car}(k(s)) = 2$) : \underline{il existe une forme quadratique non homogène à} $n+1$ \underline{variables, à coefficients dans} $A$

\SGAvar{\begin{equation}\tag{4.14.2}
Q(\underline{X}) = x_o + b x_o + c + \sum_{1 \leq i \leq j \leq 2m} a_{ij} x_i x_j \quad , 
\end{equation}

}{\begin{equation}\tag{1.3.2.1}
Q(\underline{X}) = x_o^2 + b x_o + c + \sum_{1 \leq i \leq j \leq 2m} a_{ij} x_i x_j \quad , 
\end{equation}

}{F7II-XV-009}\SGApage{II}{176}

 \underline{telle que} $b$ \underline{soit dans l'idéal maximal, que la forme à} $2m$ \underline{variables} $\sum a_{ij} x_i x_j$ \underline{soit non dégénérée et que le} $S$-\underline{schéma hensélisé de} $X$ \underline{en} $x$ \underline{soit isomorphe à l'hensélisé en} $(\sqrt{c},0,\ldots,0)$ \underline{du sous-schéma de} $\mathbb{E}^{n+1}_S$ \underline{d'équation} $Q = 0$ .

\SGAvar{\underline{Remarque} 1.3.3. Dans le cas (i), si l'hensélisé $x$ de $X_s$ est isomorphe par un isomorphisme $\varphi$ à l'hensélisé en $O$ de $X_o \subset \mathbb{E}^{n+1}_k$ , d'équation $\sum \bar{a}_{ij} x_i x_j = 0$ , et si les $a_{ij}$ relèvent les $\bar{a}_{ij}$ , on peut prendre $Q = \sum a_{ij} x_i x_j$ .

}{\underline{Remarque} 1.3.3. Dans le cas (i), si l'hensélisé en $x$ de $X_s$ est isomorphe par un isomorphisme $\varphi$ à l'hensélisé en $O$ de $X_o \subset \mathbb{E}^{n+1}_k$ , d'équation $\sum \bar{a}_{ij} x_i x_j = 0$ , et si les $a_{ij}$ relèvent les $\bar{a}_{ij}$ , on peut prendre $Q = \sum a_{ij} x_i x_j$ .

}{F7II-XV-010}Dans le cas (ii), si l'hensélisé en $x$ de $X_s$ est isomorphe par un isomorphisme $\varphi$ à l'hensélisé en $(\sqrt{a},0,\ldots,0)$ de $X_o \subset \mathbb{E}^{n+1}_k$ d'équation $(x_o^2 + \bar{a}) + \sum \bar{a}_{ij} x_i x_j$ , et si les $a_{ij}$ relèvent les $\bar{a}_{ij}$ , on peut prendre $Q = x_o^2 + bx_o + c + \sum a_{ij} x_i x_j$ , avec $c \equiv \bar{a}$ mod l'idéal maximal.

Dans les deux cas, l'isomorphisme promis en 1.3.2 peut être pris prolongeant $\varphi$ .

\underline{Corollaire} 1.3.4. \underline{Soit} $f : X \to S$ \underline{comme en 1.3.2. Il existe un voisinage} $U$ \underline{de} $x$ \underline{dans} $X$ \underline{tel que les points de non lissité de} $f$ \underline{contenus dans} $U$ \underline{soient des points quadratiques ordinaires de leur fibre.}


\SGApage{II}{177}


\SGAsectiontoc{§2. Calcul des cycles évanescents dans un cas quadratique ordinaire standard}{2}
§2. \underline{Calcul des cycles évanescents dans un cas quadratique ordinaire standard}


2.1. \underline{Préliminaire : la cohomologie d'un cône.}

2.1.1. Soit $Y \subset \mathbb{P}^r$ une variété projective sur un corps algébriquement clos $k$ . Posons les notations suivantes.

$X \subset \mathbb{A}^{r+1}$ est le cône projetant affine, d'origine $O$ ;

$X_{(o)}$ est l'hensélisé de $X$ en $O$ ;

$X^* = X - \{O\}$ et $X^*_{(o)} = X_{(o)} - \{O\}$ ;

\SGAvar{$X_1 \subset \mathbb{P}^{r+1}$ est le cône projectant projectif de $Y$ ; on a $X = X_1 - Y$ ;

}{$X_1 \subset \mathbb{P}^{r+1}$ est le cône projetant projectif de $Y$ ; on a $X = X_1 - Y$ ;

}{F7II-XV-011}$\widetilde{X}$ , $\widetilde{X}_{(o)}$ , $\widetilde{X}_1$ sont les éclatés de $X$ , $X_{(o)}$ , $X_1$ en $O$ ;

$Y_o$ est l'image réciproque de $O$ dans $\widetilde{X}$ , $\widetilde{X}_{(o)}$ ou $\widetilde{X}_1$ ;

$h : \widetilde{X}_1 \to Y$ est la projection naturelle ;

$y_o , y_\infty : Y \to \widetilde{X}_1$ sont les sections de $h$ d'image $Y_o$ et $Y$ ;

$F$ est un groupe abélien de torsion premier à la caractéristique de $k$ ; tous les groupes de cohomologie considérés seront à coefficients dans $F$ , qui ne sera pas mentionné.

Notre but dans ce n$^{\circ}$ est de relier les cohomologies de $X^*_{(o)}$ et $Y$ .

\underline{Proposition} 2.1.2. \underline{On a}

\begin{equation*}
\mathrm{(i)} \quad H^i(X) \stackrel{\sim}{\longrightarrow} H^i(\{O\}) = \left\{ \begin{array}{ll} F & \underline{\mathrm{si}} \; i = 0 \\[1.5ex] 0 & \underline{\mathrm{si}} \; i > 0 \quad ; \end{array} \right.
\end{equation*}
\begin{equation*}
\mathrm{(ii)} \quad H^i_{\{O\}}(X) \stackrel{\sim}{\longrightarrow} H^i_c(X) \quad .
\end{equation*}
Cette proposition se démontre par un argument d'homotopie, qu'il nous va falloir formaliser.

\SGAvar{Soit $k$ un corps algébriquement clos, $U$ et $V$ deux schémas sur $k$, $\Lambda$ un anneau de torsion premier à la caractéristique de $k$ , $K \in \mathrm{Ob}\ D^+(U,\Lambda)$ et $L \in \mathrm{Ob}\ D^+(V,X)$ .


}{Soit $k$ un corps algébriquement clos, $U$ et $V$ deux schémas sur $k$, $\Lambda$ un anneau de torsion premier à la caractéristique de $k$ , $K \in \mathrm{Ob}\ D^+(U,\Lambda)$ et $L \in \mathrm{Ob}\ D^+(V,\Lambda)$ .


}{F7II-XV-012}\SGApage{II}{178}

Appelons ``morphisme'' $f$ de $(U,K)$ dans $(V,L)$ un couple formé de $f : U \to V$ et de $\varphi : f^*L \to K$ . Un tel morphisme induit sur la cohomologie

\begin{equation*}
f^* : H^*(V,L) \longrightarrow H^*(U,K) \quad .
\end{equation*}
Disons que deux ``morphismes'' $f_o$ et $f_1$ sont homotopes s'il existe un schéma connexe de type fini $T$ sur $k$ , deux points $0$ et $1$ de $T$ et un ``morphisme'' $H$ de $(U \times T , \mathrm{pr}_1^* K)$ dans $(V,L)$ qui induise $f_o$ et $f_1$ en prenant les fibres en $0$ et $1$ .

\underline{Lemme} 2.1.3. \underline{Si} $f_o$ \underline{est homotope à} $f_1$ , \underline{alors} $f_o^*$ \underline{coïncide avec} $f_1^*$ .

On sait qu'il existe une suite de points $0 = x_o,x_1,\ldots,x_n = 1$ de $T$ , des courbes lisses et connexes $\Gamma_i$ et des morphismes $a_i : \Gamma_i \to T$ tels que $x_i$ et $x_{i+1}$ soient dans l'image de $a_i$ . Pour construire les $\Gamma_i$ , on normalise un sous-schéma connexe de dimension 1 de $T$ qui contienne $0$ et $1$ . Ceci nous ramène au cas où $T$ est une courbe lisse. Soient

\begin{center}
\begin{tikzcd}[row sep=large, column sep=large]
U \times T \arrow[r, "\mathrm{pr}_1"] \arrow[d, "\mathrm{pr}_2"'] \& U \arrow[d, "f"] \\
T \arrow[r, "t"] \& k
\end{tikzcd}\quad .
\end{center}










En vertu du théorème de changement de base par un morphisme lisse, appliqué à $t$ , on a

\begin{equation*}
t^*Rf_* K \xrightarrow{\ \sim\ } R\mathrm{pr}_{2*} (\mathrm{pr}_1^* K) \quad .
\end{equation*}
L'application $f_i^*$ $(i = 0,1)$ se factorise alors en

\begin{center}
\begin{tikzcd}[row sep=large, column sep=large]
H^n(V,L) \arrow[r, "H^*"] \& H^n(U \times T,\mathrm{pr}_1^* K) \arrow[r] \& H^o(T,R^n\mathrm{pr}_{2*} (\mathrm{pr}_1^* K)) \arrow[dd, "\text{fibre en } i"] \\
 {} \& H^o(T,t^* H^n(U,K)) \arrow[ur, "\wr"] \& {} \\
 {} \& {} \& H^n(U,K) \arrow[ul, "\wr"]
\end{tikzcd}
\end{center}















\SGApage{II}{179}

 et est donc indépendante de $i$ .

2.1.3. Prouvons 2.1.2. L'application identique de $X$ est homotope (grâce aux homothéties) à l'application de $X$ dans $X$ constante de valeur $0$ ; (i) en résulte.

Utilisant des homothéties de rapport tendant vers l'infini, on voit que $Y$ est un rétracte par déformation de $X_1^*$ et (ii) se déduit du lemme des 5 et du diagramme

\begin{center}
\begin{tikzcd}[row sep=large, column sep=large]
{} \arrow[r, "\partial"] \& H^i_{\{0\}}(X) \arrow[r] \arrow[d] \& H^i(X_1) \arrow[r] \arrow[d, equal] \& H^i(X_1^*) \arrow[r, "\partial"] \arrow[d, "\wr"] \& {} \\
{} \arrow[r, "\partial"] \& H^i_c(X) \arrow[r] \& H^i(X_1) \arrow[r] \& H^i(Y) \arrow[r, "\partial"] \& {}
\end{tikzcd}\quad .
\end{center}














\underline{Corollaire} 2.1.4. \qquad $H^i(X^*) \xrightarrow{\ \sim\ } H^i(X^*_{(0)})$ .

On applique le lemme des 5 au diagramme

\begin{center}
\begin{tikzcd}[row sep=large, column sep=large]
{} \arrow[r] \& H^i_{\{0\}}(X) \arrow[r] \arrow[d, "\wr"'] \& H^i(X) \arrow[r] \arrow[d, "\wr", "(2.1.2(\mathrm{i}))"'] \& H^i(X^*) \arrow[r] \arrow[d] \& {} \\
{} \arrow[r] \& H^i_{\{0\}}(X_{(0)}) \arrow[r] \& H^i(X_{(0)}) \arrow[r] \& H^i(X^*_{(0)}) \arrow[r] \& {}
\end{tikzcd}\quad .
\end{center}













2.1.5. Utilisant la suite spectrale de Leray de $h|\widetilde{X}_1-Y_o$ et $h|\widetilde{X}_1-Y$ , on trouve que les morphismes de restriction

\begin{equation*}
H^i(\widetilde{X}_1-Y_o) \xrightarrow{\ \sim\ } H^i(Y)
\end{equation*}
\begin{equation*}
H^i(\widetilde{X}) \qquad \xrightarrow{\ \sim\ } H^i(Y_o) = H^i(Y)
\end{equation*}
sont des isomorphismes. Via ces isomorphismes, les suites exactes longues de cohomologie des couples $(\widetilde{X}_1-Y_o,Y)$ , $(\widetilde{X}_1-Y,Y_o)$ sont les lignes d'un diagramme 


\SGApage{II}{180}

\begin{equation}\tag{2.1.5.1}
\begin{tikzcd}[column sep=large]
H^{i-1}(X^*) \arrow[r] \arrow[d, equal] \& H^{i-2}(Y)(-1) \arrow[r, "\eta"'] \arrow[d, equal] \& H^i(Y) \arrow[r] \arrow[d, dash, "-1"] \& H^i(X^*) \arrow[r] \arrow[d, equal] \& {} \& {} \\
H^{i-1}(X^*) \arrow[r] \& H^{i-2}(Y)(-1) \arrow[r, "-\eta"'] \& H^i(Y) \arrow[r] \& H^i(X^*) \arrow[r] \& {} \& \mbox{.}
\end{tikzcd}
\end{equation}
















\underline{Lemme} 2.1.6.  \underline{Les flèches notées} $\eta$ \underline{et} $-\eta$ \underline{de 2.1.5.1 sont les cup-produits avec la classe de cohomologie d'une section hyperplane. Le diagramme 2.1.5.1 est commutatif.}

La première assertion résulte de ce que $\eta$ (resp. $-\eta$ ) est la restriction à $Y$ (resp. $Y_o \sim Y$) de la classe de cohomologie du faisceau inversible $\mathcal{O}(Y)$ (resp. $\mathcal{O}(Y_o)$) sur $\widetilde{X}_1 - Y_o$ (resp. $\widetilde{X}_1 - Y$) .

Utilisons la suite spectrale de Leray de $h$ : les deux lignes de (2.1.5.1) s'obtiennent en appliquant le foncteur $H(Y,\ \ )$ à des triangles distingués

\begin{center}
\begin{tikzcd}[column sep=large]
{} \& y_{\infty *}Ry^!_{\infty}F \arrow[r] \& R(h|\widetilde{X}_1-Y_o)_* F \arrow[r, "\mbox{\textcircled{\scriptsize 1}}"] \& R(h|X^*)_* F \arrow[r, "\mbox{\textcircled{\scriptsize 2}}"] \& {} \\
\mbox{et} \& y_{o*}Ry^!_o F \arrow[r] \& R(h|\widetilde{X}_1-Y)_* F \arrow[r] \& R(h|X^*)_* F \arrow[r] \& {}
\end{tikzcd}
\end{center}












les faisceaux de cohomologie non nuls des complexes sommets de ces triangles sont ceux d'indice (2), (0), (0 et 1) , de sorte que les flèches \textcircled{\scriptsize 1} et \textcircled{\scriptsize 2} sont déterminées par leur effet sur les faisceaux de cohomologie. Ceci nous ramène au cas où $Y$ est un point. Ce cas est laissé au lecteur.

2.1.7.  La suite exacte longue du couple $(\widetilde{X}_{(0)},Y_o)$ est un analogue local de la $2^e$ ligne de (2.1.5.1). On a

\begin{equation*}
H^i(\widetilde{X}_{(0)}) \stackrel{\sim}{\longrightarrow} H^i(Y_o)
\end{equation*}
(théorème de changement de base pour le morphisme propre $\widetilde{X}_{(0)} \to X_{(0)}$ ; la suite exacte longue en question s'écrit donc 


\SGApage{II}{181}

\SGAvar{\begin{equation}\tag{2.1 7.1}
\begin{tikzcd}[column sep=large]
H^{i-1}(X^*_{(0)}) \arrow[r] \& H^{i-2}(Y_o)(-1) \arrow[r, "-\eta"] \& H^i(Y_o) \arrow[r] \& H^i(X^*_{(0)}) \& {} \\
H^{i-1}(X^*) \arrow[r] \arrow[u, "\wr"] \& H^{i-2}(Y)(-1) \arrow[r, "-\eta"] \arrow[u, equal] \& H^i(Y) \arrow[r] \arrow[u, equal] \& H^i(X^*) \arrow[u, dash, "\wr"] \& \mbox{,}
\end{tikzcd}
\end{equation}
















}{\begin{equation}\tag{2.1.7.1}
\begin{tikzcd}[column sep=large]
H^{i-1}(X^*_{(0)}) \arrow[r] \& H^{i-2}(Y_o)(-1) \arrow[r, "-\eta"] \& H^i(Y_o) \arrow[r] \& H^i(X^*_{(0)}) \& {} \\
H^{i-1}(X^*) \arrow[r] \arrow[u, "\wr"] \& H^{i-2}(Y)(-1) \arrow[r, "-\eta"] \arrow[u, equal] \& H^i(Y) \arrow[r] \arrow[u, equal] \& H^i(X^*) \arrow[u, dash, "\wr"] \& \mbox{,}
\end{tikzcd}
\end{equation}
















}{F7II-XV-013}la seconde ligne de ce diagramme commutatif étant celle de (2.1.5.1).

Le lemme technique suivant nous servira en 2.2.7.

\SGAvar{\underline{Lemme} 2.7.8.

}{\underline{Lemme} 2.1.8.

}{F7II-XV-014}\begin{center}
\begin{tikzcd}[column sep=large]
H^{n-1}(Y_o) \arrow[d, equal] \& H^{n-1}(\widetilde{X}) \arrow[l, "\sim"'] \arrow[r] \& H^{n-1}(X^*) \arrow[r] \& H^n_{\{x\}}(X) \arrow[r] \& H^n_c(X_s) \arrow[d, equal] \\
H^{n-1}(Y) \arrow[rrrr, "\partial"] \& {} \& {} \& {} \& H^n_c(X_s)
\end{tikzcd}
\end{center}













\underline{est anticommutatif.}

D'après 2.1.7, il revient au même de prouver commutatif le diagramme

\begin{center}
\begin{tikzcd}[column sep=large]
H^{n-1}(X_1-\{o\}) \arrow[r, "\partial"] \arrow[d, "\wr"] \& H^n_{\{o\}}(X_1) = H^n_{\{o\}}(X) \arrow[d] \\
H^{n-1}(Y) \arrow[r, "\partial"] \& H^n_c(X)
\end{tikzcd}
\end{center}










C'est là un general non-sense.

2.2.  \underline{Cycles évanescents} .

2.2.1.  Soient $S$ un trait hensélien et $s,\eta,\bar{s},\bar{\eta}$ comme en XIII 0.2.5 , $\Lambda = \mathbb{Z}/k$ , avec $k$ premier à la caractéristique résiduelle de $S$ , $\mathbb{E}^{n+1}_S$ l'espace affine type de dimension $n+1$ sur $S$ et $X$ le sous-schéma fermé de $\mathbb{E}^{n+1}_S$ défini par une équation quadratique $Q = 0$ , non nulle modulo l'uniformisante 


\SGApage{II}{182}

\begin{equation*}
 Q(\underline{\mathbf{x}}) = \sum_{i \leq j} a_{ij}\, x_i x_j + \sum b_i x_i + c \quad .
\end{equation*}
Si $\mathbb{E}^{n+1}_S$ est le complément dans $\mathbb{P}^{n+1}_S$ de l'hyperplan à l'infini $H$ , alors $X$ se déduit de la quadratique $X_1$ de $\mathbb{P}^{n+1}_S$ d'équation homogène

\begin{equation}\tag{2.2.1.1}
\sum_{i \leqslant j} a_{ij}\, x_i x_j + \sum b_i x_i z + cz^2 = 0
\end{equation}
par soustraction de $Y = X_1 \cap H$ .

Faisons les hypothèses suivantes.

(a) $Y$ est une quadrique lisse sur $S$ , ie. la forme quadratique $\sum_{i \leq j} a_{ij}\, x_i x_j$ est ordinaire;

\SGAvar{(b) la quadrique $X'_s$ est singulière, ie. $X_{\bar{s}}$ est un cône quadratique.

}{(b) la quadrique $X_s$ est singulière, ie. $X_{\bar{s}}$ est un cône quadratique.

}{F7II-XV-015}On désigne par $x_o$ le point singulier de $X_s$ , et par $f$ la projection de $X$ sur $S$ .

\underline{Remarque} 2.2.2. Si $x_o$ est un point rationnel, on peut par translation le mettre à l'origine des coordonnées. Si $x_o = 0$ , les $b_i$ et $c$ sont divisibles par l'uniformisante.

Le sous-schéma $A$ de $X_s$ d'équations $\frac{\partial Q}{\partial x_i}$ est concentré en $x_o$ . Sauf dans le cas exceptionnel où $k(s)$ est de caractéristique 2 et où $n+1$ est pair, ce schéma est de degré un sur $k(s)$ , donc $x_o$ est rationnel. Dans le cas exceptionnel, il est ramifié sur $k(s)$ et de rang 2, de sorte que $k(x_o)$ est soit $k(s)$ , soit une extension quadratique inséparable de $k(s)$ . Les assertions faites sur $A$ se vérifient le plus aisément après passage à la clôture algébrique $k(\bar{s})$ de $k(s)$ .

\underline{Proposition} 2.2.3. \underline{Sous les hypothèses de 2.2.1, les morphismes} 


\SGApage{II}{183}

\begin{equation*}
 H^i(X_{\bar{\eta}} , \Lambda) \;\stackrel{\sim}{\longrightarrow}\; \mathbb{H}^i(X_{\bar{s}}, R\psi_{\bar{\eta}} (\Lambda)) \;\stackrel{\sim}{\longrightarrow}\; R^i\psi_{\bar{\eta}}(\Lambda)_{x_o}
\end{equation*}
\begin{equation*}
H_c^i(X_{\bar{\eta}} , \Lambda) \;\stackrel{\sim}{\longleftarrow}\; \mathbb{H}_c^i(X_{\bar{s}}, R\psi_{\bar{\eta}} (\Lambda)) \;\stackrel{\sim}{\longleftarrow}\; \mathbb{H}^i_{\{x_o\}}(X_{\bar{s}}, R\psi_{\bar{\eta}}(\Lambda))
\end{equation*}
\underline{sont des isomorphismes.}

Que les flèches de gauche soient des isomorphismes résulte de 2.1.8.6 et 2.1.10.5; que celles de droite en soient est un corollaire de 2.1.2; on dispose en effet d'un triangle distingué

\begin{equation*}
(\Lambda \mbox{ sur } X_{\bar{s}})[0] \longrightarrow R\psi_{\bar{\eta}}(\Lambda) \longrightarrow R\Phi_{\bar{\eta}}(\Lambda)
\end{equation*}
et 2.1.2 s'applique à $(\Lambda$ sur $X_{\bar{s}})$ , tandis que $R\Phi_{\bar{\eta}}(\Lambda)$ est à support dans $x_o$ .

\underline{Corollaire} 2.2.4. \underline{Si la fibre générale géométrique} $X_{\bar{\eta}}$ \underline{est singulière, ie. est encore un cône quadratique, les faisceaux} $R^i\Phi(\Lambda)$ \underline{sont tous nuls.}

\SGAvar{D'après 2.1.2, $\Lambda = H^o(X_{\bar{s}},\Lambda) \stackrel{\sim}{\longrightarrow} H^o(X_{\bar{\eta}},\Lambda)$ , et $0 = H^i(X_{\bar{s}},\Lambda) = H^i(X_{\bar{\eta}},\Lambda)$ . On applique alors la suite exacte longue XIII, 2.1.8.9 (valide par XIII, 2.1.9).

}{D'après 2.1.2, $\Lambda = H^0(X_{\bar{s}},\Lambda) \stackrel{\sim}{\longrightarrow} H^0(X_{\bar{\eta}},\Lambda)$ , et $0 = H^i(X_{\bar{s}},\Lambda) = H^i(X_{\bar{\eta}},\Lambda)$ . On applique alors la suite exacte longue XIII, 2.1.8.9 (valide par XIII, 2.1.9).

}{F7II-XV-016}2.2.5. Supposons maintenant que

($*$) la fibre générale $X_\eta$ est lisse.

Pour simplifier  nous supposerons aussi $S$ strictement hensélien. Il est facile de calculer les faisceaux de cycles évanescents à l'aide de 2.2.3 et XII 3.7. La cohomologie à coefficients dans $\Lambda$ , considérée ici, se déduit de la cohomologie $\ell$ - adique de loc. cit. par la formule des coefficients universels (cf. 2.1.13) . Voici les résultats.

\underline{A. Cohomologie sans support,}

\SGAvar{On a $R^i\psi_\eta(\Lambda) = 0$ si $i \neq 0,n$ .


}{On a $R^i\psi_{\bar\eta}(\Lambda) = 0$ si $i \neq 0,n$ .


}{F7II-XV-017}\SGApage{II}{184}

\SGAvar{1) \underline{si $n \neq 0$} : on a $\psi_{\overline{\eta}}(\Lambda) = \Lambda$    De plus, $R^n\psi_{\overline{\eta}}(\Lambda)$ est (non canoniquement) isomorphe au faisceau $\Lambda$ sur $x_o$ prolongé par $0$ sur $X_s$ .

}{1) \underline{si $n \neq 0$} : on a $\psi_{\overline{\eta}}(\Lambda) = \Lambda$ .    De plus, $R^n\psi_{\overline{\eta}}(\Lambda)$ est (non canoniquement) isomorphe au faisceau $\Lambda$ sur $x_o$ prolongé par $0$ sur $X_s$ .

}{F7II-XV-018}2) \underline{$n$ quelconque} : on vérifie que, même pour $n = 0$ , $\Phi^i_{\overline{\eta}}(\Lambda) = 0$ pour $i \neq n$ , et que $\Phi^n_{\overline{\eta}}(\Lambda)$ est (non canoniquement) isomorphe au faisceau $\Lambda$ sur $x_o$ prolongé par $0$ sur $X_s$ .


\underline{B\quad Cohomologie à support dans} $x_o$ .

On a $H^i_{\{x_o\}}(X_s,R\psi_{\overline{\eta}}(\Lambda)) = 0$ pour $i \neq n,2n$ . On dispose de plus d'un morphisme trace

\begin{equation}\tag{2.2.5.1}
\mathrm{tr} \; : \; H^{2n}_{\{x_o\}}(X_s,R\psi_{\overline{\eta}}(\Lambda(n))) \longrightarrow \Lambda \;.
\end{equation}
\underline{si $n \neq 0$} : le morphisme (2.2.5.1) est un isomorphisme. De plus, $H^n_{\{x_o\}}(X_s,R\psi_{\overline{\eta}}(\Lambda(n)))$ est (non canoniquement) isomorphe à $\Lambda$ .

\underline{C. Dualité.} Pour $a \in R^n\psi_{\overline{\eta}}(\Lambda)_{x_o}$ et $b \in H^n_{\{x_o\}}(X_s,R\psi_{\overline{\eta}}(\Lambda(n)))$, posons $(a,b) = \mathrm{Tr}(a \wedge b)$ , où $\mathrm{Tr}$ est le morphisme XIII 2.2.6

\SGAvar{\begin{equation*}
\mathrm{Tr} \; : \; H^{2n}_{\{x_o\}}(X_s,R\psi_{\eta}(\Lambda(n))) \longrightarrow \Lambda \;.
\end{equation*}
Il résulte de la dualité de Poincaré sur $X_{\overline{\eta}}$ et des isomorphismes 2.2.3 que $(a,b)$ met en dualité parfaite les $\Lambda$ - modules libres $R^n\psi_{\overline{\eta}}(\Lambda)_{x_o}$ et $H^n_{\{x_o\}}(X_s,R\psi_{\overline{\eta}}(\Lambda(n))$ .


}{\begin{equation*}
\mathrm{Tr} \; : \; H^{2n}_{\{x_o\}}(X_s,R\psi_{\bar\eta}(\Lambda(n))) \longrightarrow \Lambda \;.
\end{equation*}
Il résulte de la dualité de Poincaré sur $X_{\overline{\eta}}$ et des isomorphismes 2.2.3 que $(a,b)$ met en dualité parfaite les $\Lambda$ - modules libres $R^n\psi_{\overline{\eta}}(\Lambda)_{x_o}$ et $H^n_{\{x_o\}}(X_s,R\psi_{\overline{\eta}}(\Lambda(n)))$ .


}{F7II-XV-019}\underline{D. \quad $n = 2m > 0$}.

D'après XII 3.7, les groupes $H^n_{\{x_o\}}(X_s,R\psi_{\overline{\eta}}(\Lambda(m)))$ et $R^n\psi_{\overline{\eta}}(\Lambda(m))_{x_o}$ ont des générateurs naturels définis au signe près, soit $\delta$ et $\delta'$ . Soit $\varphi$ l'application naturelle du premier de ces groupes dans le second. On peut normaliser les signes de $\delta$ et $\delta'$ de sorte que

\begin{equation}\tag{2.2.5.2}
(\delta',\delta) = 1
\end{equation}
on a alors
\SGApage{II}{185}

\begin{equation}\tag{2.2.5.3}
\varphi(\delta) = (-1)^m . 2 . \delta' \qquad .
\end{equation}
Soit $Z$ l'extension quadratique séparable de $k(\eta)$ centre de la partie paire de l'algèbre de Clifford de la forme quadratique (2.2.1.1). Soit $\epsilon$ le caractère correspondant d'ordre 1 ou 2 du groupe d'inertie $I$. Pour $\sigma \in I$ , on a

\begin{equation}\tag{2.2.5.4}
\sigma\delta = \epsilon(\sigma)\delta
\end{equation}
\begin{equation}\tag{2.2.5.5}
\sigma\delta' = \epsilon(\sigma)\delta' \;.
\end{equation}
Les formules (2.2.5.2) (2.2.5.3) (2.2.5.4) en cohomologie à coefficients dans $\mathbb{Z}/2k$ impliquent que la variation est

\SGAvar{\begin{equation}\tag{2.2.5.6}
\mathrm{Var}(\sigma) \; (a) = \frac{\epsilon(\sigma)-1}{2} \; (-1)^m \; (a\delta).\delta \;.
\end{equation}

}{\begin{equation}\tag{2.2.5.6}
\mathrm{Var}(\sigma) \; (a) = \frac{\epsilon(\sigma)-1}{2} \; (-1)^m \; (a,\delta).\delta \;.
\end{equation}

}{F7II-XV-020}\underline{E. \quad $n = 2m + 1$ \quad - Cycles évanescents.}

D'après XII 3.7, les groupes $H^n_{\{x_o\}}(X_s,R\psi_{\overline{\eta}} \, \Lambda(m))$ et $R^n\psi_{\overline{\eta}}(\Lambda(m+1))_{x_o}$ ont des générateurs naturels définis au signe près, soit $\delta$ et $\delta'$ . Avec les notations de D , on peut normaliser les signes de $\delta$ et $\delta'$ de sorte que

\begin{equation}\tag{2.2.5.7}
(\delta',\delta \;) = 1
\end{equation}
; on a

\begin{equation}\tag{2.2.5.8}
\varphi(\delta) = 0 \;.
\end{equation}

\underline{F. \quad $n = 2m + 1$ \quad - La variation} .

Le groupe d'inertie $I$ agit trivialement sur la cohomologie de $X_{\eta}$ . Ceci résulte de XII 3.7 et du fait que $I$ agit trivialement sur la cohomologie de $Y_{\eta}$ , $Y$ étant propre et lisse sur $S$ .

Pour $\sigma$ dans le groupe d'inertie, on a nécessairement

\begin{equation}\tag{2.2.5.9}
D(\sigma)(\delta') = \lambda(\sigma).\delta \;\; , \;\; \mathrm{avec} \;\; \lambda(\sigma) \in \Lambda(1) \;.
\end{equation}
La formule d'additivité (XIII, 1.4.3.4) s'écrit ici

\begin{equation*}
\lambda(\sigma\tau) = \lambda(\sigma) + \lambda(\tau) \;.
\end{equation*}
Soit $\epsilon : I \longrightarrow \Lambda(1) = \mu_k$ l'homomorphisme tel que
\SGApage{II}{186}

\begin{equation*}
\sigma(\sqrt[k]{t}) \;=\; \epsilon(\sigma) \;\; \sqrt[k]{t}
\end{equation*}
pour  $t$  une uniformisante.  $\lambda$  est nécessairement un multiple de  $\epsilon$ , vu la structure du groupe d'inertie modéré. La formule (2.2.5.9) se récrit

\SGAvar{\begin{equation}\tag{2.2.5.10}
\mathrm{Var}(\sigma)\,(a) \;=\; \lambda_X . \epsilon(\sigma)(a\delta)\delta
\end{equation}
pour  $\lambda_X \in \Lambda$  convenable (dépendant de  $X/S$) . Nous déterminerons  $\lambda_X$  plus tard.

}{\begin{equation}\tag{2.2.5.10}
\mathrm{Var}(\sigma)\,(a) \;=\; \lambda_X . \epsilon(\sigma)(a,\delta)\delta
\end{equation}
pour  $\lambda_X \in \Lambda$  convenable (dépendant de  $X/S$) . Nous déterminerons  $\lambda_X$  plus tard.

}{F7II-XV-021}2.2.6. Nous allons donner des cycles  $\delta$  et  $\delta'$  de  (D) et (E) une description de nature locale pour la topologie étale.

Le cas où  $n$  est pair est facile : si  $2 \nmid k$  ($\Lambda = \mathbb{Z}/k$)  ,  $\pm \delta$  est caractérisé par (2.2.5.3) (2.2.5.4) , ie. par

\begin{equation}\tag{2.2.6.1}
(\delta,\delta) \;=\; (-1)^{m}.2 \qquad .
\end{equation}
En général,  $\pm \delta$  est l'image d'un quelconque cycle  $\widetilde{\delta}$ , en cohomologie à valeurs dans  $\mathbb{Z}/2^{a}k$  ($a$  grand) , qui vérifie (2.2.6.1).

Supposons  $n$  impair; $x_o$  est alors un point rationnel (2.2.2).

\SGAvar{Soient  $X_{s(o)}$  l'hensélisé en  $x_o$  de  $X_s$ ,  $\widetilde{X}_{s(o)}$  l'éclaté de  $X_{s(o)}$  en  $x_o$ ,  $Y_o$  l'image réciproque de  $x_o$  dans  $X_{s(o)}$  (une quadrique) et  $X^{*}_{s(o)} = X_{s(o)} - \{x_o\} = \widetilde{X}_{s(o)} - Y_o$ . Considérons la flèche composée  (cf. XIII, 2.4.6.1)

\begin{equation}\tag{2.2.6.2}
\begin{tikzcd}[column sep=huge, row sep=huge]
H^{n-1}(Y_o,\Lambda(m)) \& H^{n-1}(\widetilde{X}_{s(o)},\Lambda(m)) \arrow[l,"\sim"'] \arrow[r] \& H^{n-1}(X^{*}_{s(o)},\Lambda(m)) \arrow[d] \\
{} \& H^{n}_{\{x\}}(X_s,R\psi_{\bar{\eta}}(\Lambda(m))) \& H^{n}_{\{x\}}(X_s,\Lambda(m)) \arrow[l]
\end{tikzcd}\quad .
\end{equation}












\underline{Lemme} 2.2.7  \underline{La flèche 2.2.6.2 identifie}  $\mathbb{H}^{n}_{\{x\}}(X_s,R\psi_{\bar{\eta}}(\Lambda(m)))$  \underline{au quotient primitif de la cohomologie de dimension}  $2m$  \underline{de la quadrique}  $Y_o$  \underline{de dimension}  $2m$  ;  $\pm \delta$  \underline{est}
\SGApage{II}{187}
\underline{est l'image des générateurs naturels de ce quotient primitif.}

Les cycles  $\pm \delta$  sont caractérisés par leur image dans  $H^{i}_{c}(X_{\bar{\eta}},\Lambda(m))$ ; utilisant (2.1.7.1), ceci nous ramène à montrer que la flèche composée

}{Soient  $X_{s(o)}$  l'hensélisé en  $x_o$  de  $X_s$ ,  $\widetilde{X}_{s(o)}$  l'éclaté de  $X_{s(o)}$  en  $x_o$ ,  $Y_o$  l'image réciproque de  $x_o$  dans  $\widetilde{X}_{s(o)}$  (une quadrique) et  $X^{*}_{s(o)} = X_{s(o)} - \{x_o\} = \widetilde{X}_{s(o)} - Y_o$ . Considérons la flèche composée  (cf. XIII, 2.4.6.1)

\begin{equation}\tag{2.2.6.2}
\begin{tikzcd}[column sep=huge, row sep=huge]
H^{n-1}(Y_o,\Lambda(m)) \& H^{n-1}(\widetilde{X}_{s(o)},\Lambda(m)) \arrow[l,"\sim"'] \arrow[r] \& H^{n-1}(X^{*}_{s(o)},\Lambda(m)) \arrow[d] \\
{} \& H^{n}_{\{x_o\}}(X_s,R\psi_{\bar{\eta}}(\Lambda(m))) \& H^{n}_{\{x_o\}}(X_s,\Lambda(m)) \arrow[l]
\end{tikzcd}\quad .
\end{equation}












\underline{Lemme} 2.2.7  \underline{La flèche 2.2.6.2 identifie}  $\mathbb{H}^{n}_{\{x\}}(X_s,R\psi_{\bar{\eta}}(\Lambda(m)))$  \underline{au quotient primitif de la cohomologie de dimension}  $2m$  \underline{de la quadrique}  $Y_o$  \underline{de dimension}  $2m$  ;  $\pm \delta$  \underline{est}
\SGApage{II}{187}
\underline{l'image des générateurs naturels de ce quotient primitif.}

Les cycles  $\pm \delta$  sont caractérisés par leur image dans  $H^{n}_{c}(X_{\bar{\eta}},\Lambda(m))$ ; utilisant (2.1.7.1), ceci nous ramène à montrer que la flèche composée

}{F7II-XV-022}\begin{center}
\begin{tikzcd}[column sep=huge, row sep=huge]
H^{n-1}(Y_o,\Lambda(m)) \& H^{n-1}(\widetilde{X}_{s},\Lambda(m)) \arrow[l,"\sim"'] \arrow[r] \& H^{n-1}(X^{*}_{s},\Lambda(m)) \arrow[d] \\
H^{n}_{c}(X_{\bar{\eta}},\Lambda(m)) \& H^{n}_{c}(X_s,\Lambda(m)) \arrow[l,"\mathrm{sp}"'] \& H^{n}_{\{x\}}(X_s,\Lambda(m)) \arrow[l]
\end{tikzcd}
\end{center}











envoie les générateurs distingués du quotient primitif du premier groupe sur les générateurs distingués du dernier groupe.

Appliquons 2.1.8. La flèche se récrit (au signe près)

\begin{center}
\begin{tikzcd}[column sep=huge, row sep=huge]
H^{n-1}(Y_s,\Lambda(m)) \arrow[r,"\partial"] \arrow[d,equal] \& H^{n}_{c}(X_s,\Lambda(m)) \arrow[d] \\
H^{n-1}(Y_{\bar{\eta}},\Lambda(m)) \arrow[r] \& H^{n}_{c}(X_{\bar{\eta}},\Lambda(m))
\end{tikzcd}
\end{center}








et 2.2.7 en résulte.


\SGAsectiontoc{\S3.  La formule de Picard Lefschetz}{3}
\underline{\S3.  La formule de Picard Lefschetz}


\underline{3.1  Résultats préliminaires}

3.1.1.  Soient  $S$  un trait hensélien et  $f : X \to S$  un morphisme plat de type fini, purement de dimension relative  $n$  . On suppose que la fibre spéciale  $X_s$  est lisse, sauf en un ensemble fini de points  $\Sigma'$  , où  $X_s$  présente une singularité quadratique ordinaire.


\SGApage{II}{188}

Soit $\Sigma$ l'ensemble des points $x \in \Sigma'$ tel que la fibre générale $X_{\eta}$ soit lisse dans un voisinage de $x$ . Soit enfin $\Lambda$ un anneau de torsion, premier à la caractéristique résiduelle de $S$ .

\underline{Proposition} 3.1.2.  (i) \underline{Les faisceaux} $R^{i}\Phi(\Lambda)$ \underline{sont nuls pour} $i \neq n$.

\SGAvar{(ii)  \underline{Le faisceau} $R^{n}\Phi(\Lambda)$ \underline{est seul en dehors de} $\Sigma$ . \underline{Sa restriction à} $\Sigma$ \underline{est un faisceau de} $\Lambda$ - \underline{modules de rang 1}.

(iii)  \underline{Pour tout point géométrique} $\overline{x}$ \underline{de} $\Sigma$ , \underline{la forme} (a,b) \underline{(2.2.5 (C))} \underline{met en dualité les} $\Lambda$ - \underline{modules} $(R^{n}\psi(\Lambda))_{\overline{x}}$ \underline{et} $H^{n}_{\{\overline{x}\}}(R^{n}\psi_{\eta}(\Lambda))$ . \underline{Ceux-ci sont libres de rang un pour} $n \neq 0$ , \underline{de rang} 2 \underline{pour} $n = 0$ .

Les assertions de 3.1.2 sont de nature locale pour la topologie étale. Ceci permet de ne traiter que le cas où $S$ est strictement hensélien, et où $X$ est l'hypersurface de l'espace affine relatif $A^{n+1}_{S}$ défini par une équation quadratique vérifiant 2.2.1 (a) (b) (appliquer 1.3.2).

}{(ii)  \underline{Le faisceau} $R^{n}\Phi(\Lambda)$ \underline{est nul en dehors de} $\Sigma$ . \underline{Sa restriction à} $\Sigma$ \underline{est un faisceau de} $\Lambda$ - \underline{modules de rang 1}.

(iii)  \underline{Pour tout point géométrique} $\overline{x}$ \underline{de} $\Sigma$ , \underline{la forme} (a,b) \underline{(2.2.5 (C))} \underline{met en dualité les} $\Lambda$ - \underline{modules} $(R^{n}\psi(\Lambda))_{\overline{x}}$ \underline{et} $H^{n}_{\{\overline{x}\}}(X_{\bar s},R\psi_{\bar\eta}(\Lambda))$ . \underline{Ceux-ci sont libres de rang un pour} $n \neq 0$ , \underline{de rang} 2 \underline{pour} $n = 0$ .

Les assertions de 3.1.2 sont de nature locale pour la topologie étale. Ceci permet de ne traiter que le cas où $S$ est strictement hensélien, et où $X$ est l'hypersurface de l'espace affine relatif $A^{n+1}_{S}$ définie par une équation quadratique vérifiant 2.2.1 (a) (b) (appliquer 1.3.2).

}{F7II-XV-023}$f$ est lisse en tout point de $X_{s} - \Sigma'$ , de sorte que $R\Phi = 0$ en dehors de $\Sigma'$ . D'après 2.2.4, $R\Phi$ est même nul hors de $\Sigma$ . On peut maintenant supposer 2.2.5 (*). Les assertions (ii) et (iii) résultent alors de 2.2.5 (A) (B) (C) .


\underline{3.2. Dimension relative paire}.

On conserve les notations de 3.1, en prenant pour $n$ un entier pair $2m$ . On suppose pour simplifier $S$ strictement hensélien (le cas général s'obtiendrait par descente).

\SGAvar{\underline{Proposition} 3.2.1.  (i) \underline{Pour} $x \in \Sigma$ , \underline{le groupe} $H^{n}_{\{x\}}(R^{n}\psi_{\eta}(\Lambda(m))$ \underline{a un générateur naturel} $\delta$ , \underline{bien défini au signe près, caractérisé par les conditions d'être naturel en} $\Lambda$ \underline{et de vérifier} 


}{\underline{Proposition} 3.2.1.  (i) \underline{Pour} $x \in \Sigma$ , \underline{le groupe} $H^{n}_{\{x\}}(X_s,R\psi_{\bar\eta}(\Lambda(m)))$ \underline{a un générateur naturel} $\delta$ , \underline{bien défini au signe près, caractérisé par les conditions d'être naturel en} $\Lambda$ \underline{et de vérifier} 


}{F7II-XV-024}\SGApage{II}{189}

\begin{equation*}
 (\delta,\delta) \;=\; (-1)^{m}\, 2
\end{equation*}
(\underline{pour} $n = 0$ , \underline{ajouter la condition} $\mathrm{Tr}(\delta) = 0$) .

(ii)  \underline{Pour un caractère} $\varepsilon_{x} : I \longrightarrow \{\pm 1\}$ \underline{convenable (indépendant de} $\Lambda$ ) \underline{du groupe d'inertie} $I$ , \underline{la variation s'écrit}

\SGAvar{\begin{equation*}
\mathrm{Var}(\sigma)\ (a) \;=\; (-1)^{m}\Bigl(\frac{\varepsilon_{x}(\sigma)-1}{2}\Bigr)\ (a\delta)\delta
\end{equation*}
(\underline{pour} $a \in R^{n}\Phi_{\overline{\eta}}(\Lambda(m))$) . \underline{Il résulte que}

}{\begin{equation*}
\mathrm{Var}(\sigma)\ (a) \;=\; (-1)^{m}\Bigl(\frac{\varepsilon_{x}(\sigma)-1}{2}\Bigr)\ (a,\delta)\delta
\end{equation*}
(\underline{pour} $a \in R^{n}\Phi_{\overline{\eta}}(\Lambda(m))$) . \underline{Il résulte que}

}{F7II-XV-025}\begin{equation*}
\sigma(\delta) \;=\; \varepsilon_{x}(\sigma).\delta \;.
\end{equation*}
Passons à la limite pour ne traiter que le cas universel où $\Lambda$ serait $\mathbb{Z}_{\ell}$ . Au signe près,un seul $\delta$ peut alors vérifier les conditions de (i). Comme en 3.1.2, on se ramène alors au cas 2.2.5 (*), et on applique 2.2.5 (D).

\underline{Complément} 3.2.2.  Si, en $x \in \Sigma$ , l'hensélisé $X_{(x)}$ de $X$ en $x$ est isomorphe à l'hensélisé en $x_{o}$ du schéma projectif $Q$ d'équation,

\begin{equation*}
\sum_{i \leq j}\, a_{ij}\, x_{i}x_{j} \;=\; 0 \quad ,
\end{equation*}
alors, le caractère $\varepsilon_{x}$ est défini par l'extension quadratique séparable $Z(C^{+}(Q))$ du corps des fractions (ceci résulte de 2.2.5 (D)).

3.2.3.  En caractéristique résiduelle $\neq 2$ , on peut rendre 3.2.2 plus explicite. Dans ce cas, $I$ admet un seul caractère non trivial $\varepsilon$ d'ordre 2. Pour $t$ une uniformisante, on a

\begin{equation*}
\sigma(\sqrt{t}) \;=\; \varepsilon(\sigma)\sqrt{t} \quad .
\end{equation*}
D'autre part, $X_{(x)}$ est alors isomorphe à l'hensélisé en $0$ d'un schéma affine 


\SGApage{II}{190}

\begin{equation*}
\sum a_{ij}x_i x_j = b \quad ,
\end{equation*}
où $b$ est dans l'idéal maximal et où $Q(x)=\sum a_{ij}x_i x_j$ est non dégénérée. L'idéal $(b)$ définit l'image du sous-schéma de $X_{(x)}$ support de $J^n_{X_{(x)}/S}$. Le centre de l'algèbre de Clifford est isomorphe à

\begin{equation*}
k(\eta)\left(\left((-1)^{m+1}.2b\det(a_{ij})\right)^{1/2}\right)
\end{equation*}
(Bourbaki, Alg. Ch 9, §9, n$^{\circ}$ 4 Rmq finale), d'où

\begin{equation*}
\varepsilon_x = \varepsilon^{v(b)} \quad .
\end{equation*}
La variation est nulle si $b$ est de valuation paire, et sinon $\varepsilon_x=\varepsilon$.

\underline{3.3. Dimension relative impaire}

Pour démontrer la formule de Picard-Lefschetz en dimension relative impaire, nous devrons nous ramener au cas transcendant, traité en XIV 3.2.1.1. L'argument de spécialisation requis nous force à travailler sur une base qui ne soit plus de dimension un. Nous le ferons par des arguments ad hoc, adaptés au cas précis dont nous avons besoin : je ne dispose pas d'une théorie serviable sur une base de dimension $>1$, et j'ai préféré ne pas recouvrir mon ignorance de généralités.

\SGAvar{3.3.1. Soient $S=\mathrm{Spec}(A)$ un schéma strictement local noethérien et $f:X\to S$ un morphisme plat de type fini, purement de dimension relative $n$, dont la fibre spéciale $X_s$ n'a pour seuls points de non lissité que des points quadratiques non dégénérés. Soient $s$ le point fermé de $S$, $X_s$ la fibre spéciale et $x_o\in X_s$ un point singulier. D'après 1.3.2 (i), au voisinage de $x_o$, $X/S$ est localement isomorphe à un $S$-schéma défini par une équation $Q-b=0$ dans $A^{n+1}_S$, avec $b$ dans l'idéal maximal et $Q$ une forme quadratique non dégénérée. L'idéal $(b)$ de $A$ ne dépend que de $(X/S,x_o)$ : dans un voisinage de $x$, il existe une section $s\in X(A/(b))$ de $X$ sur $\mathrm{Spec}(A/(b))$, tel que $J^n_{X/S}$ soit isomorphe à $\mathcal{O}_{s(\mathrm{Spec}(A/b))}$. On note $s_x$ cette section et $b(x)$ un quelconque générateur de $(b)$.


}{3.3.1. Soient $S=\mathrm{Spec}(A)$ un schéma strictement local noethérien et $f:X\to S$ un morphisme plat de type fini, purement de dimension relative $n$, dont la fibre spéciale $X_s$ n'a pour seuls points de non lissité que des points quadratiques non dégénérés. Soient $s$ le point fermé de $S$, $X_s$ la fibre spéciale et $x_o\in X_s$ un point singulier. D'après 1.3.2 (i), au voisinage de $x_o$, $X/S$ est localement isomorphe à un $S$-schéma défini par une équation $Q-b=0$ dans $A^{n+1}_S$, avec $b$ dans l'idéal maximal et $Q$ une forme quadratique non dégénérée. L'idéal $(b)$ de $A$ ne dépend que de $(X/S,x_o)$ : dans un voisinage de $x$, il existe une section $s\in X(A/(b))$ de $X$ sur $\mathrm{Spec}(A/(b))$, telle que $J^n_{X/S}$ soit isomorphe à $\mathcal{O}_{s(\mathrm{Spec}(A/b))}$. On note $s_x$ cette section et $b(x)$ un quelconque générateur de $(b)$.


}{F7II-XV-026}\SGApage{II}{191}















3.3.2. Supposons $S$ normal, et soit $b\in A$. Pour $n$ invertible dans $A$ et $b\neq 0$, nous noterons $\varepsilon_b^{(n)}$ l'homomorphisme suivant du groupe de Galois du corps des fractions de $A$ dans $\mathbb{Z}/(n)(1)$

\SGAvar{\begin{equation*}
\varepsilon_b^{(n)}(\sigma) = (\sigma \sqrt[n]{b})/ \sqrt[n]{b} \quad .
\end{equation*}
On a $(\varepsilon_b^{(nm)})^m=\varepsilon_b^{(n)}$ . Si $\Lambda$ est un $\mathbb{Z}/n$ -algèbre, on note encore $\varepsilon_b$ le composé $\mathrm{Gal}\longrightarrow \mathbb{Z}/(n)(1)\longrightarrow \Lambda(1)$ . Il est indépendant de $n$ .

}{\begin{equation*}
\varepsilon_b^{(n)}(\sigma) = (\sigma \sqrt[n]{b})/ \sqrt[n]{b} \quad .
\end{equation*}
On a $(\varepsilon_b^{(nm)})^m=\varepsilon_b^{(n)}$ . Si $\Lambda$ est une $\mathbb{Z}/n$ -algèbre, on note encore $\varepsilon_b$ le composé $\mathrm{Gal}\longrightarrow \mathbb{Z}/(n)(1)\longrightarrow \Lambda(1)$ . Il est indépendant de $n$ .

}{F7II-XV-027}Le caractère $\varepsilon_b$ se factorise par $\pi_1(S-(\text{sous-schéma } b=0))$.

3.3.3. Supposons que $S$ soit un trait strictement hensélien, et que $n$ soit impair : $n=2m+1$. Soient $x$ un point singulier de $X_s$ dans un voisinage duquel $X_\eta$ soit lisse, et $X_1/S$ un sous-schéma d'équation $Q-b=0$ de $A^{n+1}_S$ comme en 1.3.2 (i). L'hypothèse sur $x$ signifie que $b$, nul modulo l'idéal maximal, n'est pas 0.

\SGAvar{Soient $X_{s(x)}$ l'hensélisé de $X_s$ en $x$, $\widetilde{X}_{s(x)}$ l'éclaté de $\widetilde{X}_s$ en $x$, $Y$ la quadrique image réciproque de $x$ et $\widetilde{X}^{*}_{s(x)}=X_{s(x)}-\{x\}=\widetilde{X}_{s(x)}-Y$. Le morphisme composé

}{Soient $X_{s(x)}$ l'hensélisé de $X_s$ en $x$, $\widetilde{X}_{s(x)}$ l'éclaté de $X_s$ en $x$, $Y$ la quadrique image réciproque de $x$ et $X^{*}_{s(x)}=X_{s(x)}-\{x\}=\widetilde{X}_{s(x)}-Y$. Le morphisme composé

}{F7II-XV-028}\begin{center}
\begin{tikzcd}
H^{n-1}(Y,\Lambda(m)) \& H^{n-1}(\widetilde{X}_{s(x)},\Lambda(m)) \arrow[l,"\sim"'] \arrow[r] \& H^{n-1}(X^*_{s(x)},\Lambda(m)) \arrow[d] \\
{} \& H^n_{\{x\}}(X_s,R\psi_{\bar\eta}(\Lambda(m))) \& H^n_{\{x\}}(X_s,\Lambda(m)) \arrow[l]
\end{tikzcd}
\end{center}
identifie ce dernier groupe au quotient primitif (XII 3.5), de la cohomologie de la quadrique $Y$. D'après (2.2.7), c'est en effet le cas dans une situation localement isomorphe. Les images des deux générateurs naturels de ce quotient (classes de cohomologie des génératrices de $Y$) sont deux vecteurs de base $\pm\delta$ de $H^n_{\{x\}}(X_s,R\psi_{\bar\eta}(\Lambda(m)))$.


\SGApage{II}{192}

3.3.4.  Revenons aux hypothèses plus générales de 3.3.1, avec $n = 2m + 1$ . Pour tout point singulier $x$ de $X_s$ , nous noterons encore $\pm \delta_x$ les sections de $R^n f_! \Lambda(m)$ sur $S$ image des générateurs du quotient primitif de $H^{n-1}(Y,\Lambda(m))$ (même notation que plus haut) par le morphisme composé

\begin{center}
\begin{tikzcd}
H^{n-1}(Y,\Lambda(m)) \& H^{n-1}(\widetilde{X}_{s(x)},\Lambda(m)) \arrow[l, "\sim"'] \arrow[r] \& H^{n-1}(X^{*}_{s(x)},\Lambda(m)) \arrow[d] \\
H^{0}(S,R^{n}f_{!}\Lambda(m)) \& H^{n}_{c}(X_{s},\Lambda(m)) \arrow[l, "\sim"'] \& H^{n}_{\{x\}}(X_{s},\Lambda(m)) \arrow[l]
\end{tikzcd}\quad .
\end{center}














$\pm \delta_x$ est le \underline{cycle évanescent} en $x$ .

Lorsque $S$ est un trait, $\pm \delta \in H^n_c(X_{\overline{\eta}},\Lambda(m))$ est l'image par XIII 2.1.8.4 de la classe (3.3.3) $\pm \delta \in H^n_{\{x\}}(X_s,R\psi_{\overline{\eta}}(\Lambda(m)))$ .

Nous prouverons simultanément les deux résultats suivants.

\underline{Proposition} 3.3.5. \underline{Pour} $S$ \underline{un trait hensélien et} $X/S$ \underline{comme en 3.3.3, la variation est donnée par}

\begin{equation*}
\mathrm{Var}(\sigma)(a) = (-1)^{m+1} \varepsilon_{b(x)}(\sigma)(a\delta)\delta \;.
\end{equation*}
\underline{Proposition} 3.3.6. \underline{Soit} $S$ \underline{strictement local régulier et} $f : X \to S$ \underline{un morphisme propre et plat de dimension relative} $n = 2m + 1$ \underline{tel que} $X_s$ \underline{présente un et un seul point de non lissité} $x$ - \underline{supposé quadratique non dégénéré. On suppose aussi que} $b(x)$ \underline{est un paramètre. Alors, Galois agit sur la cohomologie de la fibre générale géométrique par}

\begin{equation*}
\sigma(a) = a+(-1)^{m+1} \varepsilon_{b(x)}(\sigma)(a\delta_x)\delta_x \;.
\end{equation*}
(A) Soit $k$ un corps algébriquement clos de caractéristique $0$ , $S$ l'hensélisé
\SGApage{II}{193}
de $\mathrm{Spec}(k[T])$ en $0$ . Prouvons 3.3.5 lorsque $b(x)$ est une uniformisante.

\SGAvar{Par le principe de Lefschetz, on se ramène à supposer que $k = \mathbb{C}$ . La question étant locale, on peut supposer que $f : X \to S$ se déduit par changement de base de l'application $\sum z_i^2$ de $\mathbb{C}^{n+1}$ dans $\mathbb{C}$ . La construction donnée de $\pm \delta$ a un analogue transcendant, en cohomologique entière. Via XIV 2.1, $\pm \delta$ correspond donc à un cycle évanescent transcendant, image de

}{Par le principe de Lefschetz, on se ramène à supposer que $k = \mathbb{C}$ . La question étant locale, on peut supposer que $f : X \to S$ se déduit par changement de base de l'application $\sum z_i^2$ de $\mathbb{C}^{n+1}$ dans $\mathbb{C}$ . La construction donnée de $\pm \delta$ a un analogue transcendant, en cohomologie entière. Via XIV 2.1, $\pm \delta$ correspond donc à un cycle évanescent transcendant, image de

}{F7II-XV-029}\begin{equation*}
\pm \delta_{c\ell} \in H^{n}_{\{0\}}(X_o,R\psi_{c\ell\eta}(\mathbb{Z})) \;.
\end{equation*}
D'après XIV 2.1 et la propriété correspondante des cycles $\pm \delta$ algébriques, l'image de $\pm \delta_{c\ell}$ dans $H^n_{\{0\}}(X_o,R\psi_{c\ell\eta}(\mathbb{Z})) \otimes \mathbb{Z}_\ell$ est un générateur de ce groupe. Dès lors, $\pm \delta_{c\ell}$ sont deux générateurs de $H^n_{\{0\}}(X_o,R\psi_{c\ell\eta}(\mathbb{Z}))$ , et coïncident avec les cycles évanescents de XIV 3.2. La formule 3.3.5 résulte donc de XIV 3.2.11 (en toute rigueur, il faudrait avoir vérifié une compatibilité entre l'isomorphisme XIV 2.1 et les cup-produits et traces).

(B) Prouvons 3.3.6 pour $X/S$ comme en (A) (et de plus propre).

Cela résulte de XIII 2.4.6.2.

(C) Soient $p$ un nombre premier et $S_o$ l'hensélisé strict de $\mathrm{Spec}(\mathbb{Z}[T])$ en $(p,T)$ . Prouvons 3.3.6 pour $X/S_o$ tel que $b(x) = T$ .

\SGAvar{La restriction du faisceau $R^i f_* \Lambda$ à l'hensélisé strict de $S$ en $(T)$ est justifiable de (B). Ce faisceau étant localement constant sur le complément du diviseur lisse $T = 0$ , qui est d'inégale caractéristique, (C) résulte de (B) et du lemme d'Abhyankar.

}{La restriction du faisceau $R^i f_* \Lambda$ à l'hensélisé strict de $S_o$ en $(T)$ est justifiable de (B). Ce faisceau étant localement constant sur le complément du diviseur lisse $T = 0$ , qui est d'inégale caractéristique, (C) résulte de (B) et du lemme d'Abhyankar.

}{F7II-XV-030}(D) Dans au moins un cas où (C) s'applique, $\delta_x$ n'est pas un élément de torsion en cohomologie $\ell$-adique.


\SGApage{II}{194}

Soient $S_1$ l'hensélisé strict de $\mathrm{Spec}(\mathbb{Z})$ en $(p)$ et $V \subset \mathbb{P}^r/S_1$ une variété projective lisse sur $S_1$ , purement de dimension relative $n+1$ . On suppose qu'il existe un pinceau de Lefschetz de sections hyperplanes pour la fibre spéciale $V_p$ de $V/S_1$ (XVII, 2.5); ce pinceau est défini par une droite $d'$ dans la fibre spéciale de l'espace projectif $\check{\mathbb{P}}^r/S_1$ dual de $\mathbb{P}^r/S_1$ . Relevons $d'$ en une droite $d$ de $\check{\mathbb{P}}^r/S_1$ et soit $X \to d$ le pinceau des sections hyperplanes de $V$ paramétrées par $d$ (XVIII 3.1). Le lieu exceptionnel $\Delta$ de $d$ est étale sur $S_1$ (XVII 6) . Par localisation étale en un point de $\Delta_p \subset d$ , on obtient donc une situation (C).

\SGAvar{Pour vérifier que $\delta_x$ n'est pas de torsion, il suffit de la faire après extension des scalaires de $S_1$ à $\mathbb{C}$ . Prenons $V = \mathbb{P}^{n+1}$ , et le plongement projectif défini par $\mathcal{O}(N)$ . Pour $N \geq 2$ , on peut appliquer XVII 2.5 (valable pour tout $N$ dans ce cas particulier). On applique alors le lemme suivant.

}{Pour vérifier que $\delta_x$ n'est pas de torsion, il suffit de le faire après extension des scalaires de $S_1$ à $\mathbb{C}$ . Prenons $V = \mathbb{P}^{n+1}$ , et le plongement projectif défini par $\mathcal{O}(N)$ . Pour $N \geq 2$ , on peut appliquer XVII 2.5 (valable pour tout $N$ dans ce cas particulier). On applique alors le lemme suivant.

}{F7II-XV-031}\underline{Lemme} 3.3.7. \underline{Pour un pinceau de Lefschetz d'hypersurfaces de degré} $N$ \underline{grand dans} $\mathbb{P}^{n+1}(\mathbb{C})$ ($N \geq 3$ \underline{suffit), un cycle évanescent (en} $\mathbb{Q}_{\ell}$ - \underline{cohomologie) est non nul.}

Ce lemme résulte de la théorie des pinceaux de Lefschetz, par exemple de XVIII 6.6.1. Il n'y a pas de cercle vicieux à utiliser ici ce résultat de XVIII, basé sur la ``formule de Picard-Lefschetz'', car nous disposons déjà de (A). On peut aussi utiliser la variante plus directe XIX n$^{\circ}$ 4, purement transcendante et qui ne dépend que de XIV 3.2.

(E) Preuve de 3.3.5 en caractéristique résiduelle $p$ pour $(X/S, x \in X_s)$ .

Soit $X_o/S_o$ comme en (D). D'après 1.3.2 (i), il existe $g : S \to S_o$ tel que l'hensélisé en $x_o$ de $g^{*}X_o = X_o \times_{S_o} S$ soit $S$-isomorphe à l'hensélisé de $X$ en $x$ . Il suffit donc de prouver 3.3.5 pour $g^{*}X_o$ . On déduit par changement de base de (C) pour $X_o/S_o$ que

\begin{equation*}
\sigma(a) = a - (-1)^m \epsilon_{b(x)}(\sigma)(a\delta_x)\delta_x 
\end{equation*}

\SGApage{II}{195}

 dans $H^n(X_{o\bar{\eta}})$ , pour $\sigma$ dans le groupe d'inertie. Soit $\pi$ une uniformisante, comparons la formule ci-dessus avec (2.2.5.9), où $\epsilon = \epsilon_{\pi}$ . On trouve que

\begin{equation*}
\lambda_X \epsilon_{\pi}(\sigma)(a\delta_x)\delta_x = - (-1)^m \epsilon_{b(x)}(\sigma)(a\delta_x)\delta_x
\end{equation*}
dans $H^n(X_{o\bar{\eta}})$ . Puisque $\delta_x$ n'est pas de torsion, nous pouvons simplifier cette formule par $(a\delta_x)\delta_x$ ; la formule 2.2.5.9, pour la valeur obtenue de $\lambda_X$ , n'est autre que 3.3.5.

(F) Preuve de 3.3.6 (cas général)

\SGAvar{3.3.6 se réduit du cas général de 3.3.5 comme (B) et (C) de (A).


}{3.3.6 se déduit du cas général de 3.3.5 comme (B) et (C) de (A).


}{F7II-XV-032}\underline{3.4. \quad Résumé}.

Les propositions 3.1.2, 3.2.1, 3.2.2, 3.2.3 et 3.3.5 sont la``formule de Picard-Lefschetz''. Cette formule est l'analogue algébrique de XIV 3.2.11, mais prend une forme très différente selon la parité de la dimension relative.

Dans le cas le plus intéressant où $X/S$ est propre, les résultats peuvent se résumer comme suit.

\underline{Théorème} 3.4 \underline{(Formule de Picard-Lefschetz). Soient} $S$ \underline{un trait strictement hensélien et} $f : X \to S$ \underline{un morphisme propre et plat purement de dimension relative} $n = 2m$ \underline{ou} $2m + 1$ . \underline{On suppose que} $X_{\eta}$ \underline{est lisse et que} $X_s$ \underline{ne présente que des singularités quadratiques ordinaires. Soit} $\Sigma$ \underline{le lieu singulier de} $X_s$

\SGAvar{(i) \underline{à chaque} $x \in \Sigma$ \underline{est attaché un cycle évanescent} $\pm \delta_x \in H_n(X_{\bar{\eta}})(m)$ . \underline{Ces} $\delta_x$ \underline{sont deux à deux orthogonaux. Pour} $n$ \underline{pair, on a} $(\delta_x,\delta_x) = (-1)^m.2$ .

}{(i) \underline{à chaque} $x \in \Sigma$ \underline{est attaché un cycle évanescent} $\pm \delta_x \in H^n(X_{\bar{\eta}})(m)$ . \underline{Ces} $\delta_x$ \underline{sont deux à deux orthogonaux. Pour} $n$ \underline{pair, on a} $(\delta_x,\delta_x) = (-1)^m.2$ .

}{F7II-XV-033}(ii) \underline{On a} \qquad $H^i(X_s) \xrightarrow{\ \sim\ } H^i(X_{\bar{\eta}})$

\underline{pour} $i \neq n$ , $n+1$ \underline{et une suite exacte longue} 


\SGApage{II}{196}

\SGAvar{\begin{equation*}
 0 \longrightarrow H^n(X_s) \longrightarrow H^n(X_{\overline{\eta}}) \overset{(\ , \delta_x)}{\longrightarrow} \bigoplus_{x \in \Sigma} \Lambda(n-m) \longrightarrow H^{n+1}(X_s) \longrightarrow H^{n+1}(X_{\eta}) \longrightarrow 0
\end{equation*}
(iii) \underline{Des caractères locaux du groupe d'inertie sont définis} :

}{\begin{equation*}
 0 \longrightarrow H^n(X_s) \longrightarrow H^n(X_{\overline{\eta}}) \overset{(\ , \delta_x)}{\longrightarrow} \bigoplus_{x \in \Sigma} \Lambda(n-m) \longrightarrow H^{n+1}(X_s) \longrightarrow H^{n+1}(X_{\overline{\eta}}) \longrightarrow 0
\end{equation*}
(iii) \underline{Des caractères locaux du groupe d'inertie sont définis} :

}{F7II-XV-034}$n$ \quad \underline{pair} \qquad $\varepsilon_x : I \longrightarrow \{\pm 1\}$

$n$ \quad \underline{impair} \qquad $\varepsilon_x : I \longrightarrow \Lambda(1)$ .

\underline{Pour} $X$ \underline{régulier en} $x$ , \underline{ils sont surjectifs; pour} $x$ \underline{non dégénéré, ce sont l'unique caractère d'ordre 2 et le caractère canonique.}

\underline{L'action du groupe d'inertie} $I$ \underline{est fournie par}

\SGAvar{\begin{equation*}
n = 2m \qquad\qquad \sigma(a) = a + (-1)^m \sum_{x \in \Sigma} \frac{\varepsilon_x(\sigma)-1}{2}\, (a\delta_x)\delta_x
\end{equation*}
\begin{equation*}
n = 2m + 1 \qquad\qquad \sigma(a) = a - (-1)^m \sum_{x \in \Sigma} \varepsilon_x(\sigma)(a\delta_x)\delta_x \quad .
\end{equation*}
\underline{En particulier}

}{\begin{equation*}
n = 2m \qquad\qquad \sigma(a) = a + (-1)^m \sum_{x \in \Sigma} \frac{\varepsilon_x(\sigma)-1}{2}\, (a,\delta_x)\delta_x
\end{equation*}
\begin{equation*}
n = 2m + 1 \qquad\qquad \sigma(a) = a - (-1)^m \sum_{x \in \Sigma} \varepsilon_x(\sigma)(a,\delta_x)\delta_x \quad .
\end{equation*}
\underline{En particulier}

}{F7II-XV-035}$n$ \quad \underline{pair} \qquad $\sigma(\delta_x) = \varepsilon_x(\sigma)\delta_x$

$n$ \quad \underline{impair} \qquad $\sigma(\delta_x) = \delta_x$ .
\selectlanguage{french}

\ifnum\SGAform=1
\clearpage\section*{Apparat critique}\small
Les lectures ci-dessous sont données dans le TeX exact des témoins, afin de préserver les distinctions de notation. Les indices décimaux renvoient aux différences individuelles du passage marqué.

\par\medskip\noindent\hypertarget{app-F7II-XV-001.1}{}\textbf{F7II-XV-001.1}\quad p. 166; misprint\_pass.\par
\noindent Imprimé : {\ttfamily\footnotesize lcsquelles\allowbreak{}}\par
\noindent Corrigé : {\ttfamily\footnotesize lesquelles\allowbreak{}}\par
\noindent Source : 7II:259. \nolinkurl{fixes/FIXLOG_7II_CURRENT.jsonl#650}; \nolinkurl{fixes/MISPRINTS_7II.json#259}\par

\par\medskip\noindent\hypertarget{app-F7II-XV-002.1}{}\textbf{F7II-XV-002.1}\quad p. 166; misprint\_pass.\par
\noindent Imprimé : {\ttfamily\footnotesize <\allowbreak{}}\par
\noindent Corrigé : {\ttfamily\footnotesize \textbackslash{}leq\allowbreak{}}\par
\noindent Source : 7II:260. \nolinkurl{fixes/FIXLOG_7II_CURRENT.jsonl#651}; \nolinkurl{fixes/MISPRINTS_7II.json#260}\par

\par\medskip\noindent\hypertarget{app-F7II-XV-002.2}{}\textbf{F7II-XV-002.2}\quad p. 166; misprint\_pass.\par
\noindent Imprimé : {\ttfamily\footnotesize differentiables\allowbreak{}}\par
\noindent Corrigé : {\ttfamily\footnotesize différentiables\allowbreak{}}\par
\noindent Source : 7II:261. \nolinkurl{fixes/FIXLOG_7II_CURRENT.jsonl#652}; \nolinkurl{fixes/MISPRINTS_7II.json#261}\par

\par\medskip\noindent\hypertarget{app-F7II-XV-003.1}{}\textbf{F7II-XV-003.1}\quad p. 168; printed misprint.\par
\noindent Imprimé : \emph{(vide)}\par
\noindent Corrigé : {\ttfamily\footnotesize  que\allowbreak{}}\par
\noindent Source : 7II:262. \nolinkurl{fixes/MISPRINTS_7II.json#262}\par

\par\medskip\noindent\hypertarget{app-F7II-XV-004.1}{}\textbf{F7II-XV-004.1}\quad p. 168; misprint\_pass.\par
\noindent Imprimé : {\ttfamily\footnotesize 1\allowbreak{}}\par
\noindent Corrigé : {\ttfamily\footnotesize i\allowbreak{}}\par
\noindent Source : 7II:263. \nolinkurl{fixes/FIXLOG_7II_CURRENT.jsonl#654}; \nolinkurl{fixes/MISPRINTS_7II.json#263}\par

\par\medskip\noindent\hypertarget{app-F7II-XV-004.2}{}\textbf{F7II-XV-004.2}\quad p. 168; designated-witness attribution.\par
\noindent Imprimé : {\ttfamily\footnotesize 4\allowbreak{}}\par
\noindent Corrigé : {\ttfamily\footnotesize 1\allowbreak{}}\par
\noindent Source : WITNESS-F7II-XV-004.2. \nolinkurl{fr/sga7ii/diplomatic/expose_XV_body.tex}; \nolinkurl{fr/sga7ii/current/expose_XV_body.tex}\par

\par\medskip\noindent\hypertarget{app-F7II-XV-004.3}{}\textbf{F7II-XV-004.3}\quad p. 168; designated-witness attribution.\par
\noindent Imprimé : \emph{(vide)}\par
\noindent Corrigé : {\ttfamily\footnotesize .1\allowbreak{}}\par
\noindent Source : WITNESS-F7II-XV-004.3. \nolinkurl{fr/sga7ii/diplomatic/expose_XV_body.tex}; \nolinkurl{fr/sga7ii/current/expose_XV_body.tex}\par

\par\medskip\noindent\hypertarget{app-F7II-XV-004.4}{}\textbf{F7II-XV-004.4}\quad p. 168; printed misprint.\par
\noindent Imprimé : {\ttfamily\footnotesize  =\allowbreak{}}\par
\noindent Corrigé : \emph{(vide)}\par
\noindent Source : 7II:264. \nolinkurl{fixes/MISPRINTS_7II.json#264}\par

\par\medskip\noindent\hypertarget{app-F7II-XV-004.5}{}\textbf{F7II-XV-004.5}\quad p. 168; misprint\_pass.\par
\noindent Imprimé : \emph{(vide)}\par
\noindent Corrigé : {\ttfamily\footnotesize \textbackslash{}overline\{\allowbreak{}}\par
\noindent Source : 7II:265. \nolinkurl{fixes/FIXLOG_7II_CURRENT.jsonl#655}; \nolinkurl{fixes/MISPRINTS_7II.json#265}\par

\par\medskip\noindent\hypertarget{app-F7II-XV-004.6}{}\textbf{F7II-XV-004.6}\quad p. 168; misprint\_pass.\par
\noindent Imprimé : \emph{(vide)}\par
\noindent Corrigé : {\ttfamily\footnotesize \}\allowbreak{}}\par
\noindent Source : 7II:265. \nolinkurl{fixes/FIXLOG_7II_CURRENT.jsonl#655}; \nolinkurl{fixes/MISPRINTS_7II.json#265}\par

\par\medskip\noindent\hypertarget{app-F7II-XV-005.1}{}\textbf{F7II-XV-005.1}\quad p. 170; misprint\_pass.\par
\noindent Imprimé : {\ttfamily\footnotesize envoye\allowbreak{}}\par
\noindent Corrigé : {\ttfamily\footnotesize envoie\allowbreak{}}\par
\noindent Source : 7II:267. \nolinkurl{fixes/FIXLOG_7II_CURRENT.jsonl#656}; \nolinkurl{fixes/MISPRINTS_7II.json#267}\par

\par\medskip\noindent\hypertarget{app-F7II-XV-005.2}{}\textbf{F7II-XV-005.2}\quad p. 170; printed misprint.\par
\noindent Imprimé : {\ttfamily\footnotesize 4\allowbreak{}}\par
\noindent Corrigé : {\ttfamily\footnotesize 1.2\allowbreak{}}\par
\noindent Source : 7II:268. \nolinkurl{fixes/MISPRINTS_7II.json#268}\par

\par\medskip\noindent\hypertarget{app-F7II-XV-006.1}{}\textbf{F7II-XV-006.1}\quad p. 171; misprint\_pass.\par
\noindent Imprimé : {\ttfamily\footnotesize )\allowbreak{}}\par
\noindent Corrigé : \emph{(vide)}\par
\noindent Source : 7II:269. \nolinkurl{fixes/FIXLOG_7II_CURRENT.jsonl#657}; \nolinkurl{fixes/MISPRINTS_7II.json#269}\par

\par\medskip\noindent\hypertarget{app-F7II-XV-007.1}{}\textbf{F7II-XV-007.1}\quad p. 172; misprint\_pass.\par
\noindent Imprimé : {\ttfamily\footnotesize nous\allowbreak{}}\par
\noindent Corrigé : {\ttfamily\footnotesize Nous\allowbreak{}}\par
\noindent Source : 7II:270. \nolinkurl{fixes/FIXLOG_7II_CURRENT.jsonl#658}; \nolinkurl{fixes/MISPRINTS_7II.json#270}\par

\par\medskip\noindent\hypertarget{app-F7II-XV-007.2}{}\textbf{F7II-XV-007.2}\quad p. 172; misprint\_pass.\par
\noindent Imprimé : {\ttfamily\footnotesize envoye\allowbreak{}}\par
\noindent Corrigé : {\ttfamily\footnotesize envoie\allowbreak{}}\par
\noindent Source : 7II:271. \nolinkurl{fixes/FIXLOG_7II_CURRENT.jsonl#659}\par

\par\medskip\noindent\hypertarget{app-F7II-XV-008.1}{}\textbf{F7II-XV-008.1}\quad p. 173; printed misprint.\par
\noindent Imprimé : {\ttfamily\footnotesize 4\allowbreak{}}\par
\noindent Corrigé : {\ttfamily\footnotesize 1.2\allowbreak{}}\par
\noindent Source : 7II:272. \nolinkurl{fixes/MISPRINTS_7II.json#272}\par

\par\medskip\noindent\hypertarget{app-F7II-XV-009.1}{}\textbf{F7II-XV-009.1}\quad p. 175; printed misprint.\par
\noindent Imprimé : {\ttfamily\footnotesize 4\allowbreak{}}\par
\noindent Corrigé : {\ttfamily\footnotesize 1\allowbreak{}}\par
\noindent Source : 7II:275. \nolinkurl{fixes/MISPRINTS_7II.json#275}\par

\par\medskip\noindent\hypertarget{app-F7II-XV-009.2}{}\textbf{F7II-XV-009.2}\quad p. 175; printed misprint.\par
\noindent Imprimé : {\ttfamily\footnotesize 14\allowbreak{}}\par
\noindent Corrigé : {\ttfamily\footnotesize 3\allowbreak{}}\par
\noindent Source : 7II:275. \nolinkurl{fixes/MISPRINTS_7II.json#275}\par

\par\medskip\noindent\hypertarget{app-F7II-XV-009.3}{}\textbf{F7II-XV-009.3}\quad p. 175; printed misprint.\par
\noindent Imprimé : \emph{(vide)}\par
\noindent Corrigé : {\ttfamily\footnotesize .1\allowbreak{}}\par
\noindent Source : 7II:275. \nolinkurl{fixes/MISPRINTS_7II.json#275}\par

\par\medskip\noindent\hypertarget{app-F7II-XV-009.4}{}\textbf{F7II-XV-009.4}\quad p. 175; printed misprint.\par
\noindent Imprimé : \emph{(vide)}\par
\noindent Corrigé : {\ttfamily\footnotesize \textasciicircum{}2\allowbreak{}}\par
\noindent Source : 7II:274. \nolinkurl{fixes/MISPRINTS_7II.json#274}\par

\par\medskip\noindent\hypertarget{app-F7II-XV-010.1}{}\textbf{F7II-XV-010.1}\quad p. 176; misprint\_pass.\par
\noindent Imprimé : \emph{(vide)}\par
\noindent Corrigé : {\ttfamily\footnotesize  en\allowbreak{}}\par
\noindent Source : 7II:276. \nolinkurl{fixes/FIXLOG_7II_CURRENT.jsonl#660}; \nolinkurl{fixes/MISPRINTS_7II.json#276}\par

\par\medskip\noindent\hypertarget{app-F7II-XV-011.1}{}\textbf{F7II-XV-011.1}\quad p. 177; misprint\_in\_print.\par
\noindent Imprimé : {\ttfamily\footnotesize projectant\allowbreak{}}\par
\noindent Corrigé : {\ttfamily\footnotesize projetant\allowbreak{}}\par
\noindent Source : FIX-7II-D-0174. \nolinkurl{fixes/FIXLOG_7II.jsonl#174}\par

\par\medskip\noindent\hypertarget{app-F7II-XV-012.1}{}\textbf{F7II-XV-012.1}\quad p. 177; misprint\_pass.\par
\noindent Imprimé : {\ttfamily\footnotesize X\allowbreak{}}\par
\noindent Corrigé : {\ttfamily\footnotesize \textbackslash{}Lambda\allowbreak{}}\par
\noindent Source : 7II:278. \nolinkurl{fixes/FIXLOG_7II_CURRENT.jsonl#661}; \nolinkurl{fixes/MISPRINTS_7II.json#278}\par

\par\medskip\noindent\hypertarget{app-F7II-XV-013.1}{}\textbf{F7II-XV-013.1}\quad p. 181; printed misprint.\par
\noindent Imprimé : \emph{(1 espace(s), 0 retour(s) de ligne dans le témoin TeX)}\par
\noindent Corrigé : {\ttfamily\footnotesize .\allowbreak{}}\par
\noindent Source : 7II:279. \nolinkurl{fixes/MISPRINTS_7II.json#279}\par

\par\medskip\noindent\hypertarget{app-F7II-XV-014.1}{}\textbf{F7II-XV-014.1}\quad p. 181; misprint\_pass.\par
\noindent Imprimé : {\ttfamily\footnotesize 7\allowbreak{}}\par
\noindent Corrigé : {\ttfamily\footnotesize 1\allowbreak{}}\par
\noindent Source : 7II:280. \nolinkurl{fixes/FIXLOG_7II_CURRENT.jsonl#662}; \nolinkurl{fixes/MISPRINTS_7II.json#280}\par

\par\medskip\noindent\hypertarget{app-F7II-XV-015.1}{}\textbf{F7II-XV-015.1}\quad p. 182; printed misprint.\par
\noindent Imprimé : {\ttfamily\footnotesize '\allowbreak{}}\par
\noindent Corrigé : \emph{(vide)}\par
\noindent Source : 7II:282. \nolinkurl{fixes/MISPRINTS_7II.json#282}\par

\par\medskip\noindent\hypertarget{app-F7II-XV-016.1}{}\textbf{F7II-XV-016.1}\quad p. 183; inherited corrected-layer difference.\par
\noindent Imprimé : {\ttfamily\footnotesize o\allowbreak{}}\par
\noindent Corrigé : {\ttfamily\footnotesize 0\allowbreak{}}\par
\noindent Source : ADOPT-F7II-XV-016.1. \nolinkurl{fr/sga7ii/diplomatic/expose_XV_body.tex}; \nolinkurl{fr/sga7ii/current/expose_XV_body.tex}\par

\par\medskip\noindent\hypertarget{app-F7II-XV-016.2}{}\textbf{F7II-XV-016.2}\quad p. 183; inherited corrected-layer difference.\par
\noindent Imprimé : {\ttfamily\footnotesize o\allowbreak{}}\par
\noindent Corrigé : {\ttfamily\footnotesize 0\allowbreak{}}\par
\noindent Source : ADOPT-F7II-XV-016.2. \nolinkurl{fr/sga7ii/diplomatic/expose_XV_body.tex}; \nolinkurl{fr/sga7ii/current/expose_XV_body.tex}\par

\par\medskip\noindent\hypertarget{app-F7II-XV-017.1}{}\textbf{F7II-XV-017.1}\quad p. 183; inherited corrected-layer difference.\par
\noindent Imprimé : \emph{(vide)}\par
\noindent Corrigé : {\ttfamily\footnotesize \{\textbackslash{}bar\allowbreak{}}\par
\noindent Source : ADOPT-F7II-XV-017.1. \nolinkurl{fr/sga7ii/diplomatic/expose_XV_body.tex}; \nolinkurl{fr/sga7ii/current/expose_XV_body.tex}\par

\par\medskip\noindent\hypertarget{app-F7II-XV-017.2}{}\textbf{F7II-XV-017.2}\quad p. 183; inherited corrected-layer difference.\par
\noindent Imprimé : \emph{(vide)}\par
\noindent Corrigé : {\ttfamily\footnotesize \}\allowbreak{}}\par
\noindent Source : ADOPT-F7II-XV-017.2. \nolinkurl{fr/sga7ii/diplomatic/expose_XV_body.tex}; \nolinkurl{fr/sga7ii/current/expose_XV_body.tex}\par

\par\medskip\noindent\hypertarget{app-F7II-XV-018.1}{}\textbf{F7II-XV-018.1}\quad p. 184; misprint\_pass.\par
\noindent Imprimé : \emph{(vide)}\par
\noindent Corrigé : {\ttfamily\footnotesize  .\allowbreak{}}\par
\noindent Source : 7II:283. \nolinkurl{fixes/FIXLOG_7II_CURRENT.jsonl#663}; \nolinkurl{fixes/MISPRINTS_7II.json#283}\par

\par\medskip\noindent\hypertarget{app-F7II-XV-019.1}{}\textbf{F7II-XV-019.1}\quad p. 184; inherited corrected-layer difference.\par
\noindent Imprimé : \emph{(vide)}\par
\noindent Corrigé : {\ttfamily\footnotesize \textbackslash{}bar\allowbreak{}}\par
\noindent Source : ADOPT-F7II-XV-019.1. \nolinkurl{fr/sga7ii/diplomatic/expose_XV_body.tex}; \nolinkurl{fr/sga7ii/current/expose_XV_body.tex}\par

\par\medskip\noindent\hypertarget{app-F7II-XV-019.2}{}\textbf{F7II-XV-019.2}\quad p. 184; misprint\_pass.\par
\noindent Imprimé : \emph{(vide)}\par
\noindent Corrigé : {\ttfamily\footnotesize )\allowbreak{}}\par
\noindent Source : 7II:285. \nolinkurl{fixes/FIXLOG_7II_CURRENT.jsonl#664}\par

\par\medskip\noindent\hypertarget{app-F7II-XV-020.1}{}\textbf{F7II-XV-020.1}\quad p. 185; inherited corrected-layer difference.\par
\noindent Imprimé : \emph{(vide)}\par
\noindent Corrigé : {\ttfamily\footnotesize ,\allowbreak{}}\par
\noindent Source : ADOPT-F7II-XV-020.1. \nolinkurl{fr/sga7ii/diplomatic/expose_XV_body.tex}; \nolinkurl{fr/sga7ii/current/expose_XV_body.tex}\par

\par\medskip\noindent\hypertarget{app-F7II-XV-021.1}{}\textbf{F7II-XV-021.1}\quad p. 186; inherited corrected-layer difference.\par
\noindent Imprimé : \emph{(vide)}\par
\noindent Corrigé : {\ttfamily\footnotesize ,\allowbreak{}}\par
\noindent Source : ADOPT-F7II-XV-021.1. \nolinkurl{fr/sga7ii/diplomatic/expose_XV_body.tex}; \nolinkurl{fr/sga7ii/current/expose_XV_body.tex}\par

\par\medskip\noindent\hypertarget{app-F7II-XV-022.1}{}\textbf{F7II-XV-022.1}\quad p. 186; misprint\_pass.\par
\noindent Imprimé : \emph{(vide)}\par
\noindent Corrigé : {\ttfamily\footnotesize \textbackslash{}widetilde\{\allowbreak{}}\par
\noindent Source : 7II:286. \nolinkurl{fixes/FIXLOG_7II_CURRENT.jsonl#665}; \nolinkurl{fixes/MISPRINTS_7II.json#286}\par

\par\medskip\noindent\hypertarget{app-F7II-XV-022.2}{}\textbf{F7II-XV-022.2}\quad p. 186; misprint\_pass.\par
\noindent Imprimé : \emph{(vide)}\par
\noindent Corrigé : {\ttfamily\footnotesize \}\allowbreak{}}\par
\noindent Source : 7II:286. \nolinkurl{fixes/FIXLOG_7II_CURRENT.jsonl#665}; \nolinkurl{fixes/MISPRINTS_7II.json#286}\par

\par\medskip\noindent\hypertarget{app-F7II-XV-022.3}{}\textbf{F7II-XV-022.3}\quad p. 186; designated-witness attribution.\par
\noindent Imprimé : \emph{(vide)}\par
\noindent Corrigé : {\ttfamily\footnotesize \_o\allowbreak{}}\par
\noindent Source : WITNESS-F7II-XV-022.3. \nolinkurl{fr/sga7ii/diplomatic/expose_XV_body.tex}; \nolinkurl{fr/sga7ii/current/expose_XV_body.tex}\par

\par\medskip\noindent\hypertarget{app-F7II-XV-022.4}{}\textbf{F7II-XV-022.4}\quad p. 186; designated-witness attribution.\par
\noindent Imprimé : \emph{(vide)}\par
\noindent Corrigé : {\ttfamily\footnotesize \_o\allowbreak{}}\par
\noindent Source : WITNESS-F7II-XV-022.4. \nolinkurl{fr/sga7ii/diplomatic/expose_XV_body.tex}; \nolinkurl{fr/sga7ii/current/expose_XV_body.tex}\par

\par\medskip\noindent\hypertarget{app-F7II-XV-022.5}{}\textbf{F7II-XV-022.5}\quad p. 187; misprint\_pass.\par
\noindent Imprimé : {\ttfamily\footnotesize est\allowbreak{} }\par
\noindent Corrigé : \emph{(vide)}\par
\noindent Source : 7II:288. \nolinkurl{fixes/FIXLOG_7II_CURRENT.jsonl#666}; \nolinkurl{fixes/MISPRINTS_7II.json#288}\par

\par\medskip\noindent\hypertarget{app-F7II-XV-022.6}{}\textbf{F7II-XV-022.6}\quad p. 187; printed misprint.\par
\noindent Imprimé : {\ttfamily\footnotesize i\allowbreak{}}\par
\noindent Corrigé : {\ttfamily\footnotesize n\allowbreak{}}\par
\noindent Source : 7II:289. \nolinkurl{fixes/MISPRINTS_7II.json#289}\par

\par\medskip\noindent\hypertarget{app-F7II-XV-023.1}{}\textbf{F7II-XV-023.1}\quad p. 188; printed misprint.\par
\noindent Imprimé : {\ttfamily\footnotesize seul\allowbreak{}}\par
\noindent Corrigé : {\ttfamily\footnotesize nul\allowbreak{}}\par
\noindent Source : 7II:290. \nolinkurl{fixes/MISPRINTS_7II.json#290}\par

\par\medskip\noindent\hypertarget{app-F7II-XV-023.2}{}\textbf{F7II-XV-023.2}\quad p. 188; inherited corrected-layer difference.\par
\noindent Imprimé : \emph{(vide)}\par
\noindent Corrigé : {\ttfamily\footnotesize X\_\{\textbackslash{}bar\allowbreak{} s\},\allowbreak{}}\par
\noindent Source : ADOPT-F7II-XV-023.2. \nolinkurl{fr/sga7ii/diplomatic/expose_XV_body.tex}; \nolinkurl{fr/sga7ii/current/expose_XV_body.tex}\par

\par\medskip\noindent\hypertarget{app-F7II-XV-023.3}{}\textbf{F7II-XV-023.3}\quad p. 188; inherited corrected-layer difference.\par
\noindent Imprimé : {\ttfamily\footnotesize \textasciicircum{}\{n\}\allowbreak{}}\par
\noindent Corrigé : \emph{(vide)}\par
\noindent Source : ADOPT-F7II-XV-023.3. \nolinkurl{fr/sga7ii/diplomatic/expose_XV_body.tex}; \nolinkurl{fr/sga7ii/current/expose_XV_body.tex}\par

\par\medskip\noindent\hypertarget{app-F7II-XV-023.4}{}\textbf{F7II-XV-023.4}\quad p. 188; inherited corrected-layer difference.\par
\noindent Imprimé : \emph{(vide)}\par
\noindent Corrigé : {\ttfamily\footnotesize \textbackslash{}bar\allowbreak{}}\par
\noindent Source : ADOPT-F7II-XV-023.4. \nolinkurl{fr/sga7ii/diplomatic/expose_XV_body.tex}; \nolinkurl{fr/sga7ii/current/expose_XV_body.tex}\par

\par\medskip\noindent\hypertarget{app-F7II-XV-023.5}{}\textbf{F7II-XV-023.5}\quad p. 188; misprint\_pass.\par
\noindent Imprimé : {\ttfamily\footnotesize défini\allowbreak{}}\par
\noindent Corrigé : {\ttfamily\footnotesize définie\allowbreak{}}\par
\noindent Source : 7II:291. \nolinkurl{fixes/FIXLOG_7II_CURRENT.jsonl#667}; \nolinkurl{fixes/MISPRINTS_7II.json#291}\par

\par\medskip\noindent\hypertarget{app-F7II-XV-024.1}{}\textbf{F7II-XV-024.1}\quad p. 188; inherited corrected-layer difference.\par
\noindent Imprimé : \emph{(vide)}\par
\noindent Corrigé : {\ttfamily\footnotesize X\_s,\allowbreak{}}\par
\noindent Source : ADOPT-F7II-XV-024.1. \nolinkurl{fr/sga7ii/diplomatic/expose_XV_body.tex}; \nolinkurl{fr/sga7ii/current/expose_XV_body.tex}\par

\par\medskip\noindent\hypertarget{app-F7II-XV-024.2}{}\textbf{F7II-XV-024.2}\quad p. 188; inherited corrected-layer difference.\par
\noindent Imprimé : {\ttfamily\footnotesize \textasciicircum{}\{n\}\allowbreak{}}\par
\noindent Corrigé : \emph{(vide)}\par
\noindent Source : ADOPT-F7II-XV-024.2. \nolinkurl{fr/sga7ii/diplomatic/expose_XV_body.tex}; \nolinkurl{fr/sga7ii/current/expose_XV_body.tex}\par

\par\medskip\noindent\hypertarget{app-F7II-XV-024.3}{}\textbf{F7II-XV-024.3}\quad p. 188; inherited corrected-layer difference.\par
\noindent Imprimé : \emph{(vide)}\par
\noindent Corrigé : {\ttfamily\footnotesize \textbackslash{}bar\allowbreak{}}\par
\noindent Source : ADOPT-F7II-XV-024.3. \nolinkurl{fr/sga7ii/diplomatic/expose_XV_body.tex}; \nolinkurl{fr/sga7ii/current/expose_XV_body.tex}\par

\par\medskip\noindent\hypertarget{app-F7II-XV-024.4}{}\textbf{F7II-XV-024.4}\quad p. 188; inherited corrected-layer difference.\par
\noindent Imprimé : \emph{(vide)}\par
\noindent Corrigé : {\ttfamily\footnotesize )\allowbreak{}}\par
\noindent Source : ADOPT-F7II-XV-024.4. \nolinkurl{fr/sga7ii/diplomatic/expose_XV_body.tex}; \nolinkurl{fr/sga7ii/current/expose_XV_body.tex}\par

\par\medskip\noindent\hypertarget{app-F7II-XV-025.1}{}\textbf{F7II-XV-025.1}\quad p. 189; inherited corrected-layer difference.\par
\noindent Imprimé : \emph{(vide)}\par
\noindent Corrigé : {\ttfamily\footnotesize ,\allowbreak{}}\par
\noindent Source : ADOPT-F7II-XV-025.1. \nolinkurl{fr/sga7ii/diplomatic/expose_XV_body.tex}; \nolinkurl{fr/sga7ii/current/expose_XV_body.tex}\par

\par\medskip\noindent\hypertarget{app-F7II-XV-026.1}{}\textbf{F7II-XV-026.1}\quad p. 190; misprint\_pass.\par
\noindent Imprimé : {\ttfamily\footnotesize tel\allowbreak{}}\par
\noindent Corrigé : {\ttfamily\footnotesize telle\allowbreak{}}\par
\noindent Source : 7II:293. \nolinkurl{fixes/FIXLOG_7II_CURRENT.jsonl#668}; \nolinkurl{fixes/MISPRINTS_7II.json#293}\par

\par\medskip\noindent\hypertarget{app-F7II-XV-027.1}{}\textbf{F7II-XV-027.1}\quad p. 191; misprint\_pass.\par
\noindent Imprimé : {\ttfamily\footnotesize un\allowbreak{}}\par
\noindent Corrigé : {\ttfamily\footnotesize une\allowbreak{}}\par
\noindent Source : 7II:294. \nolinkurl{fixes/FIXLOG_7II_CURRENT.jsonl#669}; \nolinkurl{fixes/MISPRINTS_7II.json#294}\par

\par\medskip\noindent\hypertarget{app-F7II-XV-028.1}{}\textbf{F7II-XV-028.1}\quad p. 191; misprint\_in\_print.\par
\noindent Imprimé : {\ttfamily\footnotesize \textbackslash{}widetilde\{\allowbreak{}}\par
\noindent Corrigé : \emph{(vide)}\par
\noindent Source : FIX-7II-D-0196. \nolinkurl{fixes/FIXLOG_7II.jsonl#196}\par

\par\medskip\noindent\hypertarget{app-F7II-XV-028.2}{}\textbf{F7II-XV-028.2}\quad p. 191; misprint\_in\_print.\par
\noindent Imprimé : {\ttfamily\footnotesize \}\allowbreak{}}\par
\noindent Corrigé : \emph{(vide)}\par
\noindent Source : FIX-7II-D-0196. \nolinkurl{fixes/FIXLOG_7II.jsonl#196}\par

\par\medskip\noindent\hypertarget{app-F7II-XV-028.3}{}\textbf{F7II-XV-028.3}\quad p. 191; designated-witness attribution.\par
\noindent Imprimé : {\ttfamily\footnotesize \textbackslash{}widetilde\{\allowbreak{}}\par
\noindent Corrigé : \emph{(vide)}\par
\noindent Source : WITNESS-F7II-XV-028.3. \nolinkurl{fr/sga7ii/diplomatic/expose_XV_body.tex}; \nolinkurl{fr/sga7ii/current/expose_XV_body.tex}\par

\par\medskip\noindent\hypertarget{app-F7II-XV-028.4}{}\textbf{F7II-XV-028.4}\quad p. 191; designated-witness attribution.\par
\noindent Imprimé : {\ttfamily\footnotesize \}\allowbreak{}}\par
\noindent Corrigé : \emph{(vide)}\par
\noindent Source : WITNESS-F7II-XV-028.4. \nolinkurl{fr/sga7ii/diplomatic/expose_XV_body.tex}; \nolinkurl{fr/sga7ii/current/expose_XV_body.tex}\par

\par\medskip\noindent\hypertarget{app-F7II-XV-029.1}{}\textbf{F7II-XV-029.1}\quad p. 193; misprint\_pass.\par
\noindent Imprimé : {\ttfamily\footnotesize cohomologique\allowbreak{}}\par
\noindent Corrigé : {\ttfamily\footnotesize cohomologie\allowbreak{}}\par
\noindent Source : 7II:296. \nolinkurl{fixes/FIXLOG_7II_CURRENT.jsonl#670}; \nolinkurl{fixes/MISPRINTS_7II.json#296}\par

\par\medskip\noindent\hypertarget{app-F7II-XV-030.1}{}\textbf{F7II-XV-030.1}\quad p. 193; misprint\_pass.\par
\noindent Imprimé : \emph{(vide)}\par
\noindent Corrigé : {\ttfamily\footnotesize \_o\allowbreak{}}\par
\noindent Source : 7II:298. \nolinkurl{fixes/FIXLOG_7II_CURRENT.jsonl#671}; \nolinkurl{fixes/MISPRINTS_7II.json#298}\par

\par\medskip\noindent\hypertarget{app-F7II-XV-031.1}{}\textbf{F7II-XV-031.1}\quad p. 194; misprint\_pass.\par
\noindent Imprimé : {\ttfamily\footnotesize la\allowbreak{}}\par
\noindent Corrigé : {\ttfamily\footnotesize le\allowbreak{}}\par
\noindent Source : 7II:299. \nolinkurl{fixes/FIXLOG_7II_CURRENT.jsonl#672}\par

\par\medskip\noindent\hypertarget{app-F7II-XV-032.1}{}\textbf{F7II-XV-032.1}\quad p. 195; misprint\_pass.\par
\noindent Imprimé : {\ttfamily\footnotesize réduit\allowbreak{}}\par
\noindent Corrigé : {\ttfamily\footnotesize déduit\allowbreak{}}\par
\noindent Source : 7II:300. \nolinkurl{fixes/FIXLOG_7II_CURRENT.jsonl#673}\par

\par\medskip\noindent\hypertarget{app-F7II-XV-033.1}{}\textbf{F7II-XV-033.1}\quad p. 195; misprint\_pass.\par
\noindent Imprimé : {\ttfamily\footnotesize \_\allowbreak{}}\par
\noindent Corrigé : {\ttfamily\footnotesize \textasciicircum{}\allowbreak{}}\par
\noindent Source : 7II:301. \nolinkurl{fixes/FIXLOG_7II_CURRENT.jsonl#674}; \nolinkurl{fixes/MISPRINTS_7II.json#301}\par

\par\medskip\noindent\hypertarget{app-F7II-XV-034.1}{}\textbf{F7II-XV-034.1}\quad p. 196; misprint\_in\_print.\par
\noindent Imprimé : \emph{(vide)}\par
\noindent Corrigé : {\ttfamily\footnotesize \textbackslash{}overline\{\allowbreak{}}\par
\noindent Source : FIX-7II-D-0206. \nolinkurl{fixes/FIXLOG_7II.jsonl#206}\par

\par\medskip\noindent\hypertarget{app-F7II-XV-034.2}{}\textbf{F7II-XV-034.2}\quad p. 196; designated-witness attribution.\par
\noindent Imprimé : \emph{(vide)}\par
\noindent Corrigé : {\ttfamily\footnotesize \}\allowbreak{}}\par
\noindent Source : WITNESS-F7II-XV-034.2. \nolinkurl{fr/sga7ii/diplomatic/expose_XV_body.tex}; \nolinkurl{fr/sga7ii/current/expose_XV_body.tex}\par

\par\medskip\noindent\hypertarget{app-F7II-XV-035.1}{}\textbf{F7II-XV-035.1}\quad p. 196; inherited corrected-layer difference.\par
\noindent Imprimé : \emph{(vide)}\par
\noindent Corrigé : {\ttfamily\footnotesize ,\allowbreak{}}\par
\noindent Source : ADOPT-F7II-XV-035.1. \nolinkurl{fr/sga7ii/diplomatic/expose_XV_body.tex}; \nolinkurl{fr/sga7ii/current/expose_XV_body.tex}\par

\par\medskip\noindent\hypertarget{app-F7II-XV-035.2}{}\textbf{F7II-XV-035.2}\quad p. 196; inherited corrected-layer difference.\par
\noindent Imprimé : \emph{(vide)}\par
\noindent Corrigé : {\ttfamily\footnotesize ,\allowbreak{}}\par
\noindent Source : ADOPT-F7II-XV-035.2. \nolinkurl{fr/sga7ii/diplomatic/expose_XV_body.tex}; \nolinkurl{fr/sga7ii/current/expose_XV_body.tex}\par

\normalsize\fi
\clearpage\SGAtoc{Exposé XVI — La formule de Milnor}
\selectlanguage{french}

\SGApage{II}{197}


SGA 7


\underline{Exposé XVI}


\underline{\underline{La formule de Milnor}}


\underline{par P. Deligne}

Dans cet exposé, on transpose en caractéristique $p$ la formule donnée par Milnor [3] (voir aussi [2]), qui permet dans certains cas de calculer des dimensions d'espaces de cycles évanescents. La démonstration est de nature globale, et ne permet de traiter que le cas ``géométrique'' d'une base un trait hensélien d'égale caractéristique à corps résiduel parfait.


\SGAsectiontoc{1. Enoncé du problème}{1}
1. \underline{Enoncé du problème}

1.1. Soient $S$ un schéma régulier purement de dimension un, $s$ un point fermé de $S$ \underline{à corps résiduel parfait}, $f : X \to S$ un morphisme plat de type fini, $x_o$ un point fermé de la fibre $X_s$ et $\overline{x}_o$ un point géométrique de $X$ localisé en $x_o$ . On suppose que

(a) $x_o$ est un \underline{point de non lissité isolé} de $f$ : il existe $n$ et un voisinage $U$ de $x_o$ , purement de dimension relative $n$ , tel que $U - \{x_o\}$ soit lisse sur $S$ ;

(b) $X$ est \underline{régulier} en $x_o$ .

Sous ces hypothèses, $X$ est une intersection complète relative sur $S$ en $x_o$ . On renvoie à l'exposé VI. 3.10 pour la définition de $D^1_{X/S}$ . Dans le voisinage $U$ de $x_o$ , on a

\begin{equation*}
D^1_{X/S} = \underline{H}^1(R\,\mathrm{Hom}(L\Omega_{X/S},\mathcal{O}_X)) = \mathrm{Ext}^1(\Omega^1_{X/S},\mathcal{O}_X) \quad ;
\end{equation*}
dans ce voisinage, $D^1_{X/S}$ est de longueur finie et de support contenu dans $\{x_o\}$ .


\SGApage{II}{198}

\underline{Définition} 1.2. Le nombre de Milnor $\mu = \mu(X/S,\bar{x}_o)$ \underline{de} $X/S$ \underline{en} $\bar{x}_o$ \underline{est la longueur du} $\mathcal{O}_{X,x_o}$\underline{-module défini par} $D^1_{X/S}$ .

\underline{Exemple} 1.3. Prenons pour $S$ une courbe lisse sur un corps algébriquement clos $k$ . Remplaçant $X$ par un voisinage de $x_o$ , on peut supposer $X$ lisse sur $k$ . Le complexe cotangent se réduit alors à la résolution suivante de $\Omega^1_{X/S}$

\begin{equation*}
L\Omega_{X/S} = [f^*\Omega^1_{S/k} \to \Omega^1_{X/k}] \qquad .
\end{equation*}
Choisissons une coordonnée locale $t : S \to \mathbb{E}^1_k$ sur $S$ en $s$ ($t$ étale en $s$ ), désignons encore par $f$ le composé $t \circ f$ , et soit $(x_i)_{1 \leq i \leq n+1}$ un système de coordonnées locales sur $X$ en $x_o$ . A ce moment, en $x_o$ , $f^*\Omega^1_{S/k}$ a pour base $df$ , et $\Omega^1_{X/k}$ a pour base les $dx_i$ . Dès lors, $D^1_{X/S}$ s'identifie au quotient de $\mathcal{O}_X$ par l'idéal engendré par les $\frac{\partial f}{\partial x_i}$

\begin{equation}\tag{1.3.1}
\mu = \dim_k(\mathcal{O}_{X,x_o}/(\frac{\partial f}{\partial x_1},\ldots,\frac{\partial f}{\partial x_{n+1}}))
\end{equation}
C'est là la définition qui figure dans Milnor.

\underline{Exemple} 1.4. Supposons que $X$ soit le sous-schéma de $\mathbb{E}^{n+1}_S$ défini par une équation $g$ . Soit $i : X \hookrightarrow \mathbb{E}^{n+1}_S$ le morphisme d'inclusion. Le complexe cotangent se réduit alors à

\begin{equation*}
[(g)/(g)^2 \longrightarrow i^*\Omega^1_{\mathbb{E}^{n+1}/S}]
\end{equation*}
et $D^1_{X/S}$ est isomorphe au quotient de $\mathcal{O}_X$ par l'idéal engendré par les $(\partial g/\partial x_i)_{1 \leq i \leq n+1}$ , soit encore au quotient de $\mathcal{O}_{\mathbb{E}^{n+1}}$ par l'idéal engendré par $g$ et les $\partial g/\partial x_i$ .

\underline{Exemple} 1.5. Prenons pour $S$ le spectre d'un anneau de valuation discrète $V$ à corps résiduel parfait et de corps des fractions $K$ . Soit $X$ le normalisé de $S$
\SGApage{II}{199}
dans une extension séparable $K'$ de $K$ . D'après 1.4 et [5] III 7.14 pg 68, $D^1_{X/S}$ est alors isomorphe au quotient de $\mathcal{O}_X$ par la différente $\mathfrak{d}_{X/S}$ . Si $\bar{x}_o$ est un point géométrique de $X_s$ , localisé en $x_o \in X_s$ , le nombre de Milnor en $\bar{x}_o$ est donc

\SGAvar{\begin{equation}\tag{1.5.1}
\mu = \ell g(\mathcal{O}_{X,x_o}/\mathfrak{d}_{X/S}) \quad .
\end{equation}
Soient $\bar{s}$ un point géométrique de $S$ localisé au point fermé $s$ de $S$ , et soit $d_{K'/K}$ un élément de $V$ qui engendre l'idéal déscriminant, de valuation $v(d_{K'/K})$ . Il résulte de 1.5.1 que

}{\begin{equation}\tag{1.5.1}
\mu = \ell g(\mathcal{O}_{X,x_o}/\mathfrak{d}_{X/S}) \quad .
\end{equation}
Soient $\bar{s}$ un point géométrique de $S$ localisé au point fermé $s$ de $S$ , et soit $d_{K'/K}$ un élément de $V$ qui engendre l'idéal discriminant, de valuation $v(d_{K'/K})$ . Il résulte de 1.5.1 que

}{F7II-XVI-001}\begin{equation}\tag{1.5.2}
v(d_{K'/K}) = \sum_{x \in X(\bar{s})} \mu(X/S,x).
\end{equation}
1.6. Gardons les hypothèses de 1.1 , désignons par $\bar{s}$ le point géométrique de $S$ image de $\bar{x}_o$ , soit $S_{(\bar{s})}$ le spectre de l'hensélisé strict de $\mathcal{O}_{S,s}$ en $\bar{s}$ , et soit $X_{(\bar{s})}$ l'image réciproque de $X$ sur $S_{(\bar{s})}$ . Soient $\bar{\eta}$ le spectre d'une clôture algébrique du corps des fractions $k(\eta)$ de $S_{(\bar{s})}$ , et $I = \mathrm{Gal}(k(\bar{\eta})/k(\eta))$ le groupe d'inertie.

Soit $\ell$ un nombre premier distinct de la caractéristique résiduelle de $S$ . Dans un voisinage du point $\bar{x}_o$ de $X_{\bar{s}}$ , les faisceaux de cycles évanescents $\Phi^i_{\bar{\eta}}(\mathbb{Z}/\ell) = R^i\Phi_{\bar{\eta}}(\mathbb{Z}/\ell)$ sur $X_{\bar{s}}$ , relatifs au faisceau constant $\mathbb{Z}/\ell$ sur $X_{(\bar{s})}$ sont concentrés en $\bar{x}_o$ ; de plus, le groupe d'inertie $I$ agit sur les $\mathbb{Z}/\ell$ - vectoriels $\Phi^i_{\bar{\eta}}(\mathbb{Z}/\ell)_{\bar{x}_o}$ .

Ceux-ci sont de dimension finie (XIII 2.4.2).

1.7. Soit $I_1$ un quotient fini du groupe d'inertie. Rappelons que le \underline{caractère de Swan} $Sw$ de $I_1$ est la différence $a - u$ du caractère d'Artin $a$ et du caractère de la représentation d'augmentation $u$ . Rappelons aussi qu'il existe, à isomorphisme près, un seul $\mathbb{Z}_\ell[I_1]$ - module projectif $S(I_1)$ de caractère $Sw$ (voir [6]).


\SGApage{II}{200}

Soit $V$ un vectoriel de dimension finie sur $\mathbb{Z}/\ell$ , muni d'une action de $I$ : $\rho : I \to \mathrm{GL}(V)$ . Supposons que $\rho$ se factorise par $I_1$ . La \underline{dimension sauvage} de $V$ est alors

\SGAvar{\begin{equation}\tag{1.7.1}
\dim \mathrm{Sw}(V) = \dim_{\mathbb{Z}/\ell} \mathrm{Hom}_{I_1}(S(I_1) \otimes_{\mathbb{Z}/\ell} \mathbb{Z}/\ell \, , \, V) \quad .
\end{equation}
Cet entier est indépendant du choix de $I_1$ , car pour $H$ distingué dans $I_1$ , on a $S(I_1/H) \sim S(I_1)_H$ .

}{\begin{equation}\tag{1.7.1}
\dim \mathrm{Sw}(V) = \dim_{\mathbb{Z}/\ell} \mathrm{Hom}_{I_1}(S(I_1) \otimes_{\mathbb{Z}_\ell} \mathbb{Z}/\ell \, , \, V) \quad .
\end{equation}
Cet entier est indépendant du choix de $I_1$ , car pour $H$ distingué dans $I_1$ , on a $S(I_1/H) \sim S(I_1)_H$ .

}{F7II-XVI-002}On pose de plus

\begin{equation*}
\dim \mathrm{tot}(V) = \dim \mathrm{Sw}(V) + \dim_{\mathbb{Z}/\ell}(V)
\end{equation*}
et on appelle ce nombre \underline{la dimension totale} de $V$ .

\underline{Définition} 1.8. \underline{Avec la terminologie de 1.6 et 1.7, on appelle} nombre de cycles évanescents \underline{de} $X/S$ \underline{en} $\overline{x}_o$ (\underline{resp.} nombre sauvage de cycles évanescents, \underline{resp}. nombre total de cycles évanescents) \underline{l'entier}

\begin{equation*}
(-1)^n \dim \Phi_{\overline{x}_o} \underset{\mathrm{dfn}}{=} (-1)^n \sum_i (-1)^i \dim \Phi^i_{\overline{\eta}} \, (\mathbb{Z}/\ell)_{\overline{x}_o}
\end{equation*}
\begin{equation*}
(\mathrm{resp} \ (-1)^n \dim \mathrm{Sw} \ \Phi_{x_o} \underset{\mathrm{dfn}}{=} (-1)^n \sum_i (-1)^i \dim \mathrm{Sw} \ \Phi^i_{\overline{\eta}} \, (\mathbb{Z}/\ell)_{\overline{x}_o}
\end{equation*}
\begin{equation*}
\mathrm{resp} \ (-1)^n \dim \mathrm{tot} \ \Phi_{x_o} \underset{\mathrm{dfn}}{=} (-1)^n \sum_i (-1)^i \dim \mathrm{tot} \ \Phi^i_{\overline{\eta}} \, (\mathbb{Z}/\ell)_{\overline{x}_o}
\end{equation*}
\begin{equation*}
= (-1)^n(\dim \Phi_{x_o} + \dim \mathrm{Sw} \ \Phi_{x_o} ) \quad )
\end{equation*}
\underline{Conjecture} 1.9. \underline{Sous les hypothèses de 1.1, le nombre de Milnor} $\mu$ \underline{de} $X/S$ \underline{en} $\overline{x}_o$ \underline{est égal au nombre} total \underline{de cycles évanescents de} $X/S$ \underline{en} $\overline{x}_o$ .

\underline{Remarque} 1.10. Soit $S$ un trait hensélien de corps résiduel imparfait et $f : X \to S$
\SGApage{II}{201}
un morphisme lisse sauf en $x_o \in X_s$ . On suppose $X$ régulier en $x_o$ . La définition 1.2 du nombre de Milnor $\mu$ garde un sens sous ces hypothèses, mais je n'ai plus d'interprétation géométrique de $\mu$ à proposer.

1.11. Dans la suite de ce n$^{\circ}$, on suppose que $S = \mathrm{Spec}(V)$ est un trait hensélien à corps résiduel \underline{algébriquement clos}, de point fermé $s$ , que $f : X \to S$ est un morphisme plat de type fini purement de dimension relative $n$ , que $X$ est régulier, que $x_o$ est un point fermé de $X_s$ et que $X - \{x_o\}$ est lisse sur $S$ .

Ceci ne restreint pas la généralité.

\underline{Proposition} 1.12. \underline{Si} $n = 0$ (\underline{ie.} $X/S$ \underline{fini), la conjecture 1.9 est vraie pour} $(X/S,x_o)$ .

On peut supposer que $X$ est le normalisé de $S$ dans une extension finie $K'$ du corps des fractions $K$ . D'après 1.5, on a $\mu = v(d_{K'/K})$ .

\SGAvar{L'extension $K'$ est définie par un sous-groupe ouvert $H$ du groupe d'inertie $I$ . Soit $H' \subset H$ un sous-groupe distingué ouvert de $I$ et $I' = I/H'$ . Soit $s_{I/H}$ le caractère de la représentation de permutation de $I'$ sur $I/H$ , $a_I$ , le caractère de la représentation d'Artin, $\mathrm{Sw}_I$ , celui de la représentation de Swan et $u_{I'}$ , celui de la représentation d'augmentation. D'après la ``Führerdiskriminantenproduktformel'' d'Artin, l'idéal discriminant est le conducteur d'Artin de $s_{I/H}$ ([5] VI 3 cor 1 pg 111) :

}{L'extension $K'$ est définie par un sous-groupe ouvert $H$ du groupe d'inertie $I$ . Soit $H' \subset H$ un sous-groupe distingué ouvert de $I$ et $I' = I/H'$ . Soit $s_{I/H}$ le caractère de la représentation de permutation de $I'$ sur $I/H$ , $a_{I'}$ , le caractère de la représentation d'Artin, $\mathrm{Sw}_{I'}$ , celui de la représentation de Swan et $u_{I'}$ , celui de la représentation d'augmentation. D'après la ``Führerdiskriminantenproduktformel'' d'Artin, l'idéal discriminant est le conducteur d'Artin de $s_{I/H}$ ([5] VI 3 cor 1 pg 111) :

}{F7II-XVI-003}\begin{equation*}
\mu = (a_{I'},s_{I/H}) = |I'|^{-1} \sum_{i \in I'} a_{I'}(i)s_{I/H}(i) \quad .
\end{equation*}

\SGApage{II}{202}

On a

\begin{equation*}
a_{I'} = u_{I'} + \mathrm{Sw}_{I'}
\end{equation*}
\begin{equation*}
(u_{I'},s_{I/H}) = \# I/H - 1
\end{equation*}
\begin{equation*}
(\mathrm{Sw}_{I'},s_{I/H}) = \dim_{\mathbb{Z}/\ell} \mathrm{Hom}_I(S(I') \otimes_{\mathbb{Z}_\ell} \mathbb{Z}/\ell, \ \mathbb{Z}/\ell^{I/H})
\end{equation*}
\begin{equation*}
= \dim_{\mathbb{Z}/\ell} \mathrm{Hom}_I(S(I') \otimes_{\mathbb{Z}_\ell} \mathbb{Z}/\ell, \mathrm{Coker}( \mathbb{Z}/\ell \to \mathbb{Z}/\ell^{I/H}))
\end{equation*}
et donc

\begin{equation*}
\mu = (\# I/H - 1) + \dim \mathrm{Sw} (\mathrm{Coker}( \mathbb{Z}/\ell \to \mathbb{Z}/\ell^{I/H})) \quad .
\end{equation*}
On vérifie aussitôt que $\Phi^i = 0$ pour $i \neq 0$ et que

\begin{equation*}
\Phi^o_{x_o} = \mathrm{Coker}( \mathbb{Z}/\ell \to \mathbb{Z}/\ell^{I/H}) \quad ;
\end{equation*}
la formule en résulte.

\underline{Proposition} 1.13. \underline{Si} $X_s$ \underline{présente en} $x_o$ \underline{une singularité quadratique ordinaire, la conjecture 1.9 est vraie pour} $(X/S,x_o)$ .

\underline{Cas 1}. \quad \underline{$n = \dim(X_s) = 2m - 1$} .

\SGAvar{Le problème est local pour la topologie étale sur $X$ . On peut donc supposer que $X_s$ est le sous-schéma de $E_S^{n+1}$ défini par une équation $f$ :

}{Le problème est local pour la topologie étale sur $X$ . On peut donc supposer que $X$ est le sous-schéma de $\mathbb{E}_S^{n+1}$ défini par une équation $f$ :

}{F7II-XVI-004}\begin{equation*}
\sum_{i=1}^{m} x_i \ x_{i+m} + a = 0
\end{equation*}
$a$ étant dans l'idéal maximal de $V$ , et que $x_o$ est l'origine des coordonnées. De plus, pour que $X$ soit régulier en $x_o$ , il faut et il suffit que $a$ soit une uniformisante. D'après 1.4, on a

\begin{equation*}
\mu = \ell g(\mathcal{O}_{\mathbb{E}_S^{2m}} / (f, \frac{\partial f}{\partial x_i})) = 1 \quad .
\end{equation*}

\SGApage{II}{203}

D'après \hspace{2.5cm} , $\Phi^i_{x_o} = 0$ pour $i \neq n$ , et $\Phi^n_{x_o}$ est de dimension un, avec action triviale de l'inertie. Le nombre sauvage de cycles évanescents est donc nul, et le nombre total de cycles évanescents est 1, ce qui vérifie 1.9.

\underline{Cas 2}. \quad \underline{$n = \dim(X_s) = 2m \ , \ \mathrm{car}(k(s)) \neq 2$} .

Sous ces hypothèses, on peut prendre $X \subset \mathbb{E}_S^{n+1}$ défini par une équation $f$

\SGAvar{\begin{equation*}
\sum_{i=1}^{m} x_ix_{i+m} \ + \ x_{2m+1}^2 + a = 0 \quad ,
\end{equation*}
$a$ étant une uniformisante. Ici encore, $\mu = 1$ , $\Phi^i_{x_o} = 0$ pour $i \neq n$ et $\Phi^i_{x_o}$ est de dimension un. Cette fois le groupe d'inertie agit par $\pm 1$ sur $\Phi^n_{x_o}$ , donc de façon modérée. Le nombre sauvage de cycles évanescents est encore nul, et 1.9 est vérifié.

}{\begin{equation*}
\sum_{i=1}^{m} x_ix_{i+m} \ + \ x_{2m+1}^2 + a = 0 \quad ,
\end{equation*}
$a$ étant une uniformisante. Ici encore, $\mu = 1$ , $\Phi^i_{x_o} = 0$ pour $i \neq n$ et $\Phi^n_{x_o}$ est de dimension un. Cette fois le groupe d'inertie agit par $\pm 1$ sur $\Phi^n_{x_o}$ , donc de façon modérée. Le nombre sauvage de cycles évanescents est encore nul, et 1.9 est vérifié.

}{F7II-XVI-005}\underline{Cas 3}. \quad \underline{$n = \dim(X_s) = 2m \ , \ \mathrm{car}(k(s)) = 2$} .

Cette fois, on peut prendre $X \subset \mathbb{E}_S^{n+1}$ défini par l'équation $f$ :

\begin{equation*}
\sum_{i=1}^{m} x_ix_{i+m} + (x_{2m+1}^2 + ax_{2m+1} + b) = 0 \quad ,
\end{equation*}
avec $a$ et $b$ dans l'idéal maximal. Pour que $X$ soit régulier, il faut et suffit que $b$ soit une uniformisante; pour que la fibre générique de $X$ soit lisse, il faut et il suffit que $a^2 - 4b \neq 0$ .

Le nombre de Milnor $\mu$ de $X/S$ coïncide avec celui de $X_o \subset \mathbb{E}_S^1$ défini par l'équation $f_o$

\begin{equation}\tag{1}
x^2 + ax + b = 0 \quad .
\end{equation}
De même, les $\Phi^i_{x_o}$ sont nuls pour $i \neq n$ , $\Phi^n_{x_o}$ est de dimension un, $I$ agit sur $\Phi^n_{x_o}$ par $\pm 1$ et le sous-groupe d'indice deux de $I$ qui agit trivialement correspond à l'extension de $K$ d'équation (1). Le nombre total de cycles évanescents de 


\SGApage{II}{204}

$X/S$ en $x_o$ coïncide donc avec celui relatif à $X_o/S$, et la validité de 1.9 dans ce cas résulte de 1.12. Le nombre de Milnor est donné par (1.5)

\begin{equation}\tag{1.13.1}
\mu=v(a^2-4b) \quad .
\end{equation}

\SGAsectiontoc{2. Le cas ``géométrique''}{2}
2. \underline{Le cas ``géométrique''}

On se propose de montrer que la conjecture 1.9 est vraie en égale caractéristique. Le résultat géométrique clef est le suivant. La possibilité de l'utiliser m'a été suggérée par des conversations avec B. Saint-Donat.

\underline{Proposition} 2.1. \underline{Soient} $S$ \underline{une courbe projective lisse et connexe sur un corps algébriquement clos} $k$, $X$ \underline{une variété propre et lisse sur} $k$, \underline{purement de dimension} $n+1$, \underline{et} $f:X\to S$ \underline{un morphisme lisse en dehors d'un ensemble fini} $F$ \underline{de points de} $X$. \underline{Les caractéristiques d'Euler-Poincaré (en cohomologie étale) de} $S$, $X$ \underline{et de la fibre générique géométrique} $X_{\bar\eta}$ \underline{de} $f$ \underline{vérifient alors}

\begin{equation}\tag{2.1.1}
\chi(X)=\chi(S)\,\chi(X_{\bar\eta})-(-1)^n\sum_{x\in F}\mu(X/S,x) \quad .
\end{equation}

On sait que la caractéristique d'Euler-Poincaré de $X$ coïncide avec la trace (ou intégrale sur $X$) de la dernière classe de Chern du fibré tangent

\begin{equation*}
\chi(X)=\int c_{n+1}(\Theta_X)=(-1)^{n+1}\int c_{n+1}(\Omega_X^1) \quad .
\end{equation*}
Posons

\begin{equation*}
\Omega_X^1=\left(\Omega_X^1\otimes(f^*\Omega_S^1)^{\otimes-1}\right)\otimes f^*\Omega_S^1 \quad .
\end{equation*}
La classe totale de Chern du faisceau inversible $f^*\Omega_S^1$ est l'image réciproque de la classe totale de Chern de $\Omega_S^1$ sur $S$


\SGApage{II}{205}

\begin{equation*}
 c^{*}(f^{*} \Omega^{1}_{S}) = 1 + f^{*}c^{1}(S,\Omega^{1}_{S})
\end{equation*}
et les cup-puissances de $c^{1}$ sont nulles :

\begin{equation*}
(f^{*} c^{1}(S,\Omega^{1}_{S}))^{n} = f^{*}(c^{1}(S,\Omega^{1}_{S})^{n}) = 0 \qquad (n > 1) \quad .
\end{equation*}
On en déduit que

\begin{equation*}
c^{n+1}(\Omega^{1}_{X}) = c^{n}(\Omega^{1}_{X} \otimes (f^{*}\Omega^{1}_{S})^{\otimes -1}) . \; f^{*}c_{1}(S,\Omega^{1}_{S}) + c^{n+1}(\Omega^{1}_{X} \otimes (f^{*} \Omega^{1}_{S})^{\otimes -1}) .
\end{equation*}
Si $S$ est de genre $g$ , alors $c_{1}(S,\Omega^{1}_{S})$ est égal à $(2g - 2)$ fois la classe de cohomologie d'un quelconque point fermé $s$ de $S$ , d'où, pour $X_{s}$ non singulier,

\begin{equation*}
I = \int_{X} c^{n}(\Omega^{1}_{X} \otimes (f^{*} \Omega^{1}_{S})^{\otimes -1}) . \; f^{*}c_{1}(S,\Omega^{1}_{S})
\end{equation*}
\begin{equation*}
= (2g-2) . \int_{X} c^{n}(\Omega^{1}_{X} \otimes (f^{*} \Omega^{1}_{S})^{\otimes -1}) . \; f^{*}c\ell(s)
\end{equation*}
\begin{equation*}
= (2g-2) . \int_{X_{s}} c^{n}(X_{s},\Omega^{1}_{X} \otimes (f^{*} \Omega^{1}_{S})^{\otimes -1} \otimes \mathcal{O}_{X_{s}})
\end{equation*}
De plus, sur $X_{s}$ , $f^{*}\Omega^{1}_{S}$ est isomorphe à $\mathcal{O}_{X_{s}}$ et $\Omega^{1}_{X} \otimes_{\mathcal{O}_{X}} \mathcal{O}_{X_{s}}$ figure dans une suite exacte

\begin{equation*}
0 \to f^{*}\Omega^{1}_{S} \otimes \mathcal{O}_{X_{s}} \to \Omega^{1}_{X} \otimes \mathcal{O}_{X_{s}} \to \Omega^{1}_{X_{s}} \to 0 \quad .
\end{equation*}
De là, on tire que

\begin{equation*}
c^{n}(X_{s},\Omega^{1}_{X} \otimes (f^{*} \Omega^{1}_{S})^{\otimes -1} \otimes \mathcal{O}_{X_{s}}) = c^{n}(X_{s},\Omega^{1}_{X_{s}}) \qquad \mathrm{et}
\end{equation*}
\begin{equation*}
I = (2g-2) . \int_{X_{s}} c^{n}(X_{s},\Omega^{1}_{X_{s}}) = (-1)^{n+1} \chi(S) \chi(X_{s}) \quad .
\end{equation*}

\SGApage{II}{206}

L'application canonique

\begin{equation*}
f^* \; : \; f^*\Omega^1_S \; \longrightarrow \; \Omega^1_X
\end{equation*}
s'identifie à une section $a$ de $\Omega^1_X \otimes (f^*\Omega^1_S)^{\otimes -1}$ . Une fois choisis des systèmes de coordonnées locales sur $S$ et $X$ , ce qui identifie $\Omega^1_X \otimes (f^*\Omega^1_S)^{\otimes -1}$ à $\mathcal{O}_X^{n+1}$ , les coordonnées de $a$ sont les dérivées partielles $\frac{\partial f}{\partial x_i}$ . La section $a$ ne s'annule donc que sur $F$ , et d'après 1.3, la multiplicité du zéro de $a$ en $x \in F$ est égale au nombre de Milnor $\mu(X/S,x)$.

Rappelons que si une section $a$ d'un fibré vectoriel $G$ de rang $d$ sur un $k$-schéma propre et lisse $X$ purement de dimension $d$ ne s'annule qu'en des points isolés, alors la trace $\int_X c^d(G)$ est égale au nombre de zéros de $a$ , comptés avec leurs multiplicités. On a ici

\begin{equation*}
\mathrm{II} \; = \; \int_X c^{n+1}(\Omega^1_X \otimes (f^*\Omega^1_S)^{\otimes -1}) \; = \; \sum_{x \in F} \mu(X/S,x) \quad .
\end{equation*}
Regroupant les termes, on trouve

\begin{equation*}
\chi(X) \; = \; (-1)^{n+1}(\mathrm{I} + \mathrm{II}) \; = \; \chi(S).\chi(X_s) \; - \; (-1)^n \sum_{x \in F} \mu(X/S,x) \quad ,
\end{equation*}
d'où 2.1 puisque $\chi(X_s) = \chi(X_{\bar{\eta}})$ .

\underline{Corollaire} 2.2. \quad \underline{Sous les hypothèses précédentes, si} $n(X/S,x)$ \underline{est le nombre} total \underline{de cycles évanescents en} $x$ \underline{de} $X/S$ (1.8), \underline{on a}

\begin{equation*}
\sum_{x \in F} \mu(X/S,x) \; = \; \sum_{x \in F} n(X/S,x) \quad .
\end{equation*}
La suite spectrale de Leray pour $f$ fournit

\begin{equation*}
\chi(X) \; = \; \sum (-1)^{p+q} \dim_{\mathbb{Z}/\ell} H^p(S,R^q f_* \mathbb{Z}/\ell) \; = \; \sum_q (-1)^q \, \chi(S,R^q f_* \mathbb{Z}/\ell)
\end{equation*}

\SGApage{II}{207}

Pour tout faisceau constructible $G$ de $\mathbb{Z}/\ell$ - modules sur $S$ , on a ([4] ou SGA 5 X 7.2)

\begin{equation*}
\chi(S,G) \; = \; \chi(S) \, . \, \dim_{\mathbb{Z}/\ell}(G_{\bar{\eta}}) \; - \; \sum_{s \in S} \varepsilon_s(G) \quad ,
\end{equation*}
où,pour chaque $s \in S$ , si on identifie $\bar{\eta}$ à la clôture algébrique du corps des fractions de l'hensélisé de $S$ en $s$ , on a

\SGAvar{\begin{equation*}
\varepsilon_s(G) \; = \; \dim(G_{\bar{\eta}}) \; - \; \dim(G_s) \; - \; \dim \mathrm{Sw}(G_{\bar{\eta}}) \quad .
\end{equation*}
Pour $s \in S$ , on déduit des suites exactes XIII 2.4.6.3

}{\begin{equation*}
\varepsilon_s(G) \; = \; \dim(G_{\bar{\eta}}) \; - \; \dim(G_s) \; + \; \dim \mathrm{Sw}(G_{\bar{\eta}}) \quad .
\end{equation*}
Pour $s \in S$ , on déduit des suites exactes XIII 2.4.6.3

}{F7II-XVI-006}\begin{equation*}
\cdots \quad (R^q f_* \mathbb{Z}/\ell)_s \; \longrightarrow \; (R^q f_* \mathbb{Z}/\ell)_{\bar{\eta}} \; \longrightarrow \; \bigoplus_{x \in F(s)} \Phi^q_{\bar{\eta}} \, (\mathbb{Z}/\ell)_x \; \longrightarrow \; \cdots
\end{equation*}
que

\begin{equation*}\tag{a}
\begin{aligned} &\sum (-1)^q \dim(R^q f_* \mathbb{Z}/\ell)_{\bar{\eta}} \; - \; \dim(R^q f_* \mathbb{Z}/\ell)_s \\[1ex] &\qquad = \; \sum_{x \in F(s)} \; \sum (-1)^q \dim \Phi^q_{\bar{\eta}} \, (\mathbb{Z}/\ell)_x \end{aligned}
\end{equation*}
\begin{equation*}\tag{b}
\sum (-1)^q \dim \mathrm{Sw}(R^q f_* \mathbb{Z}/\ell)_{\bar{\eta}} \; = \; \sum_{x \in F(s)} \; \sum (-1)^q \dim \mathrm{Sw} \; \Phi^q_{\bar{\eta}} \, (\mathbb{Z}/\ell)_x
\end{equation*}
\begin{equation}\tag{c}
\sum (-1)^q \, \varepsilon_s(R^q f_* \mathbb{Z}/\ell) \; = \; \sum_{x \in F(s)} \dim \mathrm{tot} \; \Phi_x \quad .
\end{equation}
La formule d'Euler-Poincaré se réduit donc à

\begin{equation*}
\begin{aligned} \chi(X) \; &= \; \sum (-1)^q \, \chi(S,R^q f_* \mathbb{Z}/\ell) \\[1ex] &= \; \sum (-1)^q \, (\chi(S).\dim(R^q f_* \mathbb{Z}/\ell)_{\bar{\eta}} \; - \; \sum_s \varepsilon_s(R^q f_* \mathbb{Z}/\ell)) \\[1ex] &= \; \chi(S). \; \sum (-1)^q \dim(R^q f_* \mathbb{Z}/\ell)_{\bar{\eta}} \; - \; \sum_s \; \sum (-1)^q \, \varepsilon_s(R^q f_* \mathbb{Z}/\ell) \\[1ex] &= \; \chi(S). \; \chi(X_{\bar{\eta}}) \; - \; \sum_s \; \sum_{x \in F(s)} \dim \mathrm{tot} \; \Phi_x \\[1ex] &= \; \chi(S). \; \chi(X_{\bar{\eta}}) \; - \; (-1)^n \sum_{x \in F} n(X/S,x) \end{aligned}
\end{equation*}

\SGApage{II}{208}

Comparant cette formule avec (2.1.1), on trouve 2.2.

\underline{Corollaire} 2.3. \underline{Sous les hypothèses et avec les notations précédentes,} \underline{si} $x_o \in F$ \underline{et si pour tout} $x \in F$ \underline{distinct de} $x_o$ , \underline{la fibre de} $f$ \underline{passant par} $x$ \underline{présente en} $x$ \underline{un point singulier quadratique ordinaire, alors}

\begin{equation*}
\mu(X/S,x_o) = n(X/S,x_o) \quad .
\end{equation*}
Pour $x \neq x_o$ , on a d'après 1.13

\begin{equation*}
\mu(X/S,x) = n(X/S,x) \quad ,
\end{equation*}
et il reste à appliquer 2.2.

\underline{Théorème} 2.4. \underline{Soient} $S$ \underline{un trait hensélien d'égale caractéristique à corps résiduel parfait,} $f : X \to S$ \underline{un morphisme plat de type fini et} $\overline{x}$ \underline{un point géométrique de} $X$ , \underline{d'image dans} $X$ \underline{un point fermé} $x$ \underline{de} $X_s$ \underline{qui soit un point de non lissité isolé de} $X/S$ . \underline{On suppose que} $X$ \underline{est régulier en} $x$ . \underline{Alors, le nombre} total \underline{(1.8) de cycles évanescents de} $X/S$ \underline{en} $\overline{x}$ \underline{coïncide avec le nombre de Milnor} $\mu(X/S,x)$ .

\SGAvar{Soit $S'$ le trait hensélien déduit de $S$ par extension du corps résiduel de $k(s)$ à sa clôture séparable, et soit $X' = X \times_s S'$ . L'assertion 2.4 pour $X/S$ équivaut alors à 2.4 pour $X'/S'$ . Puisque $k(s)$ est déjà supposé parfait, on peut donc supposer le corps résiduel de $S$ algébriquement clos.

Pour un tel $S$ , soit $S'$ le complété de $S$ , et $X' = X \times_s S'$ . Le nombre de Milnor de $X/S$ en $x$ est le même pour $X/S$ et $X'/S'$ . D'après XIII 2.4.2 il en est de même du nombre total de cycles évanescents. Ceci permet de supposer que $S$ est de plus complet, donc isomorphe à $k[[t]]$ (avec $k$ algébriquement clos). On a alors, pour $X_s$ de dimension $n$ en $x$ :

}{Soit $S'$ le trait hensélien déduit de $S$ par extension du corps résiduel de $k(s)$ à sa clôture séparable, et soit $X' = X \times_S S'$ . L'assertion 2.4 pour $X/S$ équivaut alors à 2.4 pour $X'/S'$ . Puisque $k(s)$ est déjà supposé parfait, on peut donc supposer le corps résiduel de $S$ algébriquement clos.

Pour un tel $S$ , soit $S'$ le complété de $S$ , et $X' = X \times_S S'$ . Le nombre de Milnor de $X/S$ en $x$ est le même pour $X/S$ et $X'/S'$ . D'après XIII 2.4.2 il en est de même du nombre total de cycles évanescents. Ceci permet de supposer que $S$ est de plus complet, donc isomorphe à $k[[t]]$ (avec $k$ algébriquement clos). On a alors, pour $X_s$ de dimension $n$ en $x$ :

}{F7II-XVI-007}\underline{Proposition} 2.5. \underline{Il existe une variété projective et lisse} $Y$ \underline{sur} $k$ , \underline{purement de dimension} $n + 1$ , \underline{un morphisme} $f : Y \to \mathbb{P}^1_k$ , \underline{lisse en dehors d'un ensemble fini} $F$
\SGApage{II}{209}
\underline{de points de} $Y$ \underline{et} $y_o \in F$ \underline{tels que}

(a) \underline{si} $y \in F$ , $y \neq y_o$ , \underline{alors} $y$ \underline{est un point singulier quadratique ordinaire de la fibre de} $f$ \underline{passant par} $y$ ;

(b) \underline{on a} $f(y_o) = 0$ , \underline{si} $a : S \to \mathbb{P}^1_k$ \underline{est l'isomorphisme naturel entre} $S$ \underline{et le complété de} $\mathbb{P}^1_k$ \underline{en} $0$ , \underline{il existe un} $S$-\underline{isomorphisme entre l'hensélisé de} $X$ \underline{en} $x$ \underline{et l'hensélisé de} $a^{*}Y$ \underline{en} $y_o$ .

\SGAvar{Montrons que $2.5 \Rightarrow 2.4$. D'après 2.3 et 2.5 (a), le nombre de Milnor de $Y/T$ en $y_o$ coïncide avec le nombre total de cycles évanescents. Comme plus haut, d'après XIII 2.4.2 , on a le même résultat pour $a^{*}Y/S$ , donc aussi pour $(X/S,x)$, qui par 2.5 (b) est localement isomorphe à $(a^{*} Y/S,y_o)$ , pour la topologie étale.

}{Montrons que $2.5 \Rightarrow 2.4$. D'après 2.3 et 2.5 (a), le nombre de Milnor de $Y/\mathbb{P}^1_k$ en $y_o$ coïncide avec le nombre total de cycles évanescents. Comme plus haut, d'après XIII 2.4.2 , on a le même résultat pour $a^{*}Y/S$ , donc aussi pour $(X/S,x)$, qui par 2.5 (b) est localement isomorphe à $(a^{*} Y/S,y_o)$ , pour la topologie étale.

}{F7II-XVI-008}Prouvons 2.5. Le complété de l'anneau local $\mathcal{O}_{X,x}$ est isomorphe, en tant que $k[[t]]$ - algèbre, à un quotient

\begin{equation*}
\mathcal{O}_{X,x}^{\wedge} = k[[t,x_1 \ldots x_{n+1}]]/(f) \quad .
\end{equation*}
On déduit alors du théorème des fonctions implicites XV 1.1.2 ($=$ [1] lemme 5.10) qu'il existe un entier $N$ tel que la condition 2.5 (b) soit vérifiée dès que $\mathcal{O}_{Y,y_o}^{\wedge}$ est $k[[t]]$ - isomorphe à $k[[t,x_1 \cdots x_{n+1}]]/(f_1)$ , avec $f_1 - f$ dans la puissance $N$\textsuperscript{ième} de l'idéal maximal $m$ de $k[[t,x_1 \cdots x_{n+1}]]$ .

Si $(t_o,t_1)$ et $(x_o,\ldots,x_{n+1})$ sont les systèmes de coordonnées homogènes usuels sur $\mathbb{P}^1$ et $\mathbb{P}^{n+1}$ , on prendra pour $Y/\mathbb{P}^1$ la sous-variété de $\mathbb{P}^1 \times \mathbb{P}^{n+1}$ , définie par une équation

\begin{equation}\tag{2.5.1}
g(t_o,t_1 \; ; \; x_o,\ldots,x_{n+1}) = 0 \quad ,
\end{equation}
homogène de degré $d'$ en $(t_o,t_1)$ et homogène de degré $d''$ en les $x_i$ . En termes plus intrinsèques, $g$ sera une section de $\mathcal{O}(d',d'') = \mathrm{pr}_1^{*} \mathcal{O}(d') \otimes \mathrm{pr}_2^{*} \mathcal{O}(d'')$ .


\SGApage{II}{210}

La condition (b) sera alors vérifiée si

\SGAvar{\begin{equation}\tag{2.5.2}
g(1,t ; 1,x_1,\ldots,x_{n+1}) \equiv f \qquad \mathrm{mod}\ m^N .
\end{equation}
Fixons $d'$ et $d''$ . Les $g$ du type (2.5.1) sont alors les points (rationnels) d'un espace affine $A$ sur $k$ . Soit $u \in \mathbb{P}^1 \times \mathbb{P}^{n+1}$ . On désignera par $B(u,k)$ l'espace affine dont les points sont les $k$-jets de sections de $\mathcal{O}(d',d'')$ en $u$ . Si $u_o = (1,0;1,0,\ldots,0)$ et si $u \neq u_o$ , alors, pour $d'$ et $d''$ assez grand, l'application de restriction

}{\begin{equation}\tag{2.5.2}
g(1,t ; 1,x_1,\ldots,x_{n+1}) \equiv f \qquad \mathrm{mod}\ m^N .
\end{equation}
Fixons $d'$ et $d''$ . Les $g$ du type (2.5.1) sont alors les points (rationnels) d'un espace affine $A$ sur $k$ . Soit $u \in \mathbb{P}^1 \times \mathbb{P}^{n+1}$ . On désignera par $B(u,r)$ l'espace affine dont les points sont les $r$-jets de sections de $\mathcal{O}(d',d'')$ en $u$ . Si $u_o = (1,0;1,0,\ldots,0)$ et si $u \neq u_o$ , alors, pour $d'$ et $d''$ assez grands, l'application de restriction

}{F7II-XVI-009}\begin{equation}\tag{2.5.3}
A \to B(u_o,N) \times B(u,2) \quad ,
\end{equation}
toujours linéaire, est surjective. Soit $C(u) \subset B(u,2)$ l'ensemble des 2-jets $h \in B(u,2)$ qui soient le 2-jet d'un germe $\tilde{h}$ nul ou tel que

(a) $u$ se trouve sur la variété $V$ d'équation $\tilde{h} = 0$ ;

(b) $u$ n'est ni un point lisse, ni un point quadratique ordinaire de la fibre passant par $u$ de $V/\mathbb{P}^1$ .

On vérifie facilement que $C(u)$ est fermé de codimension $n + 3$ dans $B(u,2)$ .

\SGAvar{Soit $D$ le sous-espace affine de $A$ formé des $g$ vérifiant (2.5.2). Pour $u \neq u_o$ , soit $D(u)$ le sous-espace de $D$ formé des $g$ nuls ou tels que $g(u) = 0$ et que $u$ ne soit ni un point lisse, ni un point quadratique ordinaire, de la fibre passant par $u$ de la variété $V$ d'équation $g = 0$ . Alors, si $d'$ et $d''$ sont assez grands pour que (2.5.3) soit surjectif, $D(u)$ est de codimension $n+3$ dans $D$ . La réunion des $D(u)$ pour $u \neq u_o$ est donc de codimension $\geq 1 = (n+3) - (n+2)$ , et un élément général $g$ de $D$ n'est dans aucun des $D(u)$ . Il résoud le problème posé.


}{Soit $D$ le sous-espace affine de $A$ formé des $g$ vérifiant (2.5.2). Pour $u \neq u_o$ , soit $D(u)$ le sous-espace de $D$ formé des $g$ nuls ou tels que $g(u) = 0$ et que $u$ ne soit ni un point lisse, ni un point quadratique ordinaire, de la fibre passant par $u$ de la variété $V$ d'équation $g = 0$ . Alors, si $d'$ et $d''$ sont assez grands pour que (2.5.3) soit surjectif, $D(u)$ est de codimension $n+3$ dans $D$ . La réunion des $D(u)$ pour $u \neq u_o$ est donc de codimension $\geq 1 = (n+3) - (n+2)$ , et un élément général $g$ de $D$ n'est dans aucun des $D(u)$ . Il résout le problème posé.


}{F7II-XVI-010}\SGApage{II}{211}


\underline{Bibliographie}

[1] M. ARTIN : Algebraic approximation of structures over complete local rings. Publ. Math. IHES \underline{36} (1969) (Zariski's volume).

[2] LE DUNG TRANG : Singularités isolées des hypersurfaces complexes. Pub. Centre Math. Ecole Polytechnique, Mai 69.

[3] J. MILNOR : Singular points of complex hypersurfaces. Princeton univ. press, 1968.

[4] M. RAYNAUD : Caractéristique d'Euler-Poincaré d'un faisceau et cohomologie des variétés abéliennes. Séminaire Bourbaki 286, février 1965.

[5] J.P. SERRE : Corps locaux. Publ. Inst. Math. un. Nancago VIII. Hermann 1962.

[6] J.P. SERRE : SGA 5, IX. Introduction à la théorie de Brauer.Reproduit dans : Représentations linéaires des groupes (3\textsuperscript{e} partie). Hermann. Paris 1967.
\selectlanguage{french}

\ifnum\SGAform=1
\clearpage\section*{Apparat critique}\small
Les lectures ci-dessous sont données dans le TeX exact des témoins, afin de préserver les distinctions de notation. Les indices décimaux renvoient aux différences individuelles du passage marqué.

\par\medskip\noindent\hypertarget{app-F7II-XVI-001.1}{}\textbf{F7II-XVI-001.1}\quad p. 199; misprint\_pass.\par
\noindent Imprimé : {\ttfamily\footnotesize déscriminant\allowbreak{}}\par
\noindent Corrigé : {\ttfamily\footnotesize discriminant\allowbreak{}}\par
\noindent Source : 7II:303. \nolinkurl{fixes/FIXLOG_7II_CURRENT.jsonl#675}; \nolinkurl{fixes/MISPRINTS_7II.json#303}\par

\par\medskip\noindent\hypertarget{app-F7II-XVI-002.1}{}\textbf{F7II-XVI-002.1}\quad p. 200; misprint\_pass.\par
\noindent Imprimé : {\ttfamily\footnotesize /\allowbreak{}}\par
\noindent Corrigé : {\ttfamily\footnotesize \_\allowbreak{}}\par
\noindent Source : 7II:304. \nolinkurl{fixes/FIXLOG_7II_CURRENT.jsonl#676}; \nolinkurl{fixes/MISPRINTS_7II.json#304}\par

\par\medskip\noindent\hypertarget{app-F7II-XVI-003.1}{}\textbf{F7II-XVI-003.1}\quad p. 201; printed misprint.\par
\noindent Imprimé : \emph{(vide)}\par
\noindent Corrigé : {\ttfamily\footnotesize \{\allowbreak{}}\par
\noindent Source : 7II:305. \nolinkurl{fixes/MISPRINTS_7II.json#305}\par

\par\medskip\noindent\hypertarget{app-F7II-XVI-003.2}{}\textbf{F7II-XVI-003.2}\quad p. 201; printed misprint.\par
\noindent Imprimé : \emph{(vide)}\par
\noindent Corrigé : {\ttfamily\footnotesize '\}\allowbreak{}}\par
\noindent Source : 7II:305. \nolinkurl{fixes/MISPRINTS_7II.json#305}\par

\par\medskip\noindent\hypertarget{app-F7II-XVI-003.3}{}\textbf{F7II-XVI-003.3}\quad p. 201; misprint\_pass.\par
\noindent Imprimé : \emph{(vide)}\par
\noindent Corrigé : {\ttfamily\footnotesize \{\allowbreak{}}\par
\noindent Source : 7II:305. \nolinkurl{fixes/FIXLOG_7II_CURRENT.jsonl#677}\par

\par\medskip\noindent\hypertarget{app-F7II-XVI-003.4}{}\textbf{F7II-XVI-003.4}\quad p. 201; misprint\_pass.\par
\noindent Imprimé : \emph{(vide)}\par
\noindent Corrigé : {\ttfamily\footnotesize '\}\allowbreak{}}\par
\noindent Source : 7II:305. \nolinkurl{fixes/FIXLOG_7II_CURRENT.jsonl#677}\par

\par\medskip\noindent\hypertarget{app-F7II-XVI-004.1}{}\textbf{F7II-XVI-004.1}\quad p. 202; printed misprint.\par
\noindent Imprimé : {\ttfamily\footnotesize \_s\allowbreak{}}\par
\noindent Corrigé : \emph{(vide)}\par
\noindent Source : 7II:307. \nolinkurl{fixes/MISPRINTS_7II.json#307}\par

\par\medskip\noindent\hypertarget{app-F7II-XVI-004.2}{}\textbf{F7II-XVI-004.2}\quad p. 202; misprint\_in\_print.\par
\noindent Imprimé : \emph{(vide)}\par
\noindent Corrigé : {\ttfamily\footnotesize \textbackslash{}mathbb\{\allowbreak{}}\par
\noindent Source : FIX-7II-D-0212. \nolinkurl{fixes/FIXLOG_7II.jsonl#212}\par

\par\medskip\noindent\hypertarget{app-F7II-XVI-004.3}{}\textbf{F7II-XVI-004.3}\quad p. 202; misprint\_in\_print.\par
\noindent Imprimé : \emph{(vide)}\par
\noindent Corrigé : {\ttfamily\footnotesize \}\allowbreak{}}\par
\noindent Source : FIX-7II-D-0212. \nolinkurl{fixes/FIXLOG_7II.jsonl#212}\par

\par\medskip\noindent\hypertarget{app-F7II-XVI-005.1}{}\textbf{F7II-XVI-005.1}\quad p. 203; printed misprint.\par
\noindent Imprimé : {\ttfamily\footnotesize i\allowbreak{}}\par
\noindent Corrigé : {\ttfamily\footnotesize n\allowbreak{}}\par
\noindent Source : 7II:309. \nolinkurl{fixes/MISPRINTS_7II.json#309}\par

\par\medskip\noindent\hypertarget{app-F7II-XVI-006.1}{}\textbf{F7II-XVI-006.1}\quad p. 207; printed misprint.\par
\noindent Imprimé : {\ttfamily\footnotesize -\allowbreak{}}\par
\noindent Corrigé : {\ttfamily\footnotesize +\allowbreak{}}\par
\noindent Source : 7II:310. \nolinkurl{fixes/MISPRINTS_7II.json#310}\par

\par\medskip\noindent\hypertarget{app-F7II-XVI-007.1}{}\textbf{F7II-XVI-007.1}\quad p. 208; printed misprint.\par
\noindent Imprimé : {\ttfamily\footnotesize s\allowbreak{}}\par
\noindent Corrigé : {\ttfamily\footnotesize S\allowbreak{}}\par
\noindent Source : 7II:311. \nolinkurl{fixes/MISPRINTS_7II.json#311}\par

\par\medskip\noindent\hypertarget{app-F7II-XVI-007.2}{}\textbf{F7II-XVI-007.2}\quad p. 208; printed misprint.\par
\noindent Imprimé : {\ttfamily\footnotesize s\allowbreak{}}\par
\noindent Corrigé : {\ttfamily\footnotesize S\allowbreak{}}\par
\noindent Source : 7II:312. \nolinkurl{fixes/MISPRINTS_7II.json#312}\par

\par\medskip\noindent\hypertarget{app-F7II-XVI-008.1}{}\textbf{F7II-XVI-008.1}\quad p. 209; designated-witness attribution.\par
\noindent Imprimé : {\ttfamily\footnotesize T\allowbreak{}}\par
\noindent Corrigé : {\ttfamily\footnotesize \textbackslash{}mathbb\{P\}\textasciicircum{}1\_k\allowbreak{}}\par
\noindent Source : WITNESS-F7II-XVI-008.1. \nolinkurl{fr/sga7ii/diplomatic/expose_XVI_body.tex}; \nolinkurl{fr/sga7ii/current/expose_XVI_body.tex}\par

\par\medskip\noindent\hypertarget{app-F7II-XVI-009.1}{}\textbf{F7II-XVI-009.1}\quad p. 210; inherited corrected-layer difference.\par
\noindent Imprimé : {\ttfamily\footnotesize k\allowbreak{}}\par
\noindent Corrigé : {\ttfamily\footnotesize r\allowbreak{}}\par
\noindent Source : ADOPT-F7II-XVI-009.1. \nolinkurl{fr/sga7ii/diplomatic/expose_XVI_body.tex}; \nolinkurl{fr/sga7ii/current/expose_XVI_body.tex}\par

\par\medskip\noindent\hypertarget{app-F7II-XVI-009.2}{}\textbf{F7II-XVI-009.2}\quad p. 210; inherited corrected-layer difference.\par
\noindent Imprimé : {\ttfamily\footnotesize k\allowbreak{}}\par
\noindent Corrigé : {\ttfamily\footnotesize r\allowbreak{}}\par
\noindent Source : ADOPT-F7II-XVI-009.2. \nolinkurl{fr/sga7ii/diplomatic/expose_XVI_body.tex}; \nolinkurl{fr/sga7ii/current/expose_XVI_body.tex}\par

\par\medskip\noindent\hypertarget{app-F7II-XVI-009.3}{}\textbf{F7II-XVI-009.3}\quad p. 210; misprint\_pass.\par
\noindent Imprimé : {\ttfamily\footnotesize grand\allowbreak{}}\par
\noindent Corrigé : {\ttfamily\footnotesize grands\allowbreak{}}\par
\noindent Source : 7II:314. \nolinkurl{fixes/FIXLOG_7II_CURRENT.jsonl#678}; \nolinkurl{fixes/MISPRINTS_7II.json#314}\par

\par\medskip\noindent\hypertarget{app-F7II-XVI-010.1}{}\textbf{F7II-XVI-010.1}\quad p. 210; misprint\_pass.\par
\noindent Imprimé : {\ttfamily\footnotesize résoud\allowbreak{}}\par
\noindent Corrigé : {\ttfamily\footnotesize résout\allowbreak{}}\par
\noindent Source : 7II:315. \nolinkurl{fixes/FIXLOG_7II_CURRENT.jsonl#679}; \nolinkurl{fixes/MISPRINTS_7II.json#315}\par

\normalsize\fi
\clearpage\SGAtoc{Exposé XVII — PINCEAUX DE LEFSCHETZ: THEOREME D'EXISTENCE}
\selectlanguage{french}

\SGApage{II}{212}


S G A 7


\underline{EXPOSE XVII}


\underline{\underline{PINCEAUX DE LEFSCHETZ: THEOREME D'EXISTENCE}}


\underline{par N. M. KATZ}


\SGAsectiontoc{0. Introduction}{0}
0. \underline{Introduction}

Dans le présent exposé (qui, dans le séminaire oral, se plaçait avant les exposés XIV, XV consacrés à la formule de Picard-Lefschetz), nous introduisons les ``pinceaux de Lefschetz'' de sections hyperplanes d'un schéma projectif et lisse, nous en donnons des conditions d'existence; dans l'exposé suivant, nous en faisons une étude cohomologique élémentaire, basée sur la structure connue de la cohomologie des variétés éclatées. Le fait que certaines hypothèses sur les pinceaux de Lefschetz que nous sommes amenés à introduire sont souvent satisfaites résultera de la formule de Picard-Lefschetz de Exp. XV, de nature moins élémentaire que les résultats du présent exposé et du suivant.

Certains développements du présent exposé et du suivant se trouveront exposés ailleurs dans un cadre plus général (schémas de base généraux, hypothèses de quasi-projectivité au lieu de projectivité, etc.), notamment l'étude pré-cohomologique des pinceaux de sections hyperplanes (et des pinceaux de Lefschetz plus particulièrement) dans EGA V, et l'étude cohomologique des schémas éclatés dans
\SGApage{II}{213}
l'exposé de JOUANOLOU SGA 5 VII; le lecteur trouvera également dans EGA V des démonstrations où la chasse aux diagrammes intrinsèque remplace l'usage des coordonnées, qu'affectionne le rédacteur du présent exposé.

Je tiens à remercier P. DELIGNE et P. GRIFFITHS pour l'aide qu'ils m'ont apportée dans la conception de cet exposé. Dans tout cet exposé, $k$ désigne un corps algébriquement clos, sauf mention expresse du contraire.


\SGAvar{\SGAsectiontoc{1. Un rappel sur les singularités quadratiques (cf VI 6).}{1}
1. \underline{Un rappel sur les singularités quadratiques} (cf VI 6).

1.1. \quad Soit $Y$ un $k$-schéma purement de dimension $n$ . Un point fermé $y \in Y$ s'appelle une \underline{singularité quadratique ordinaire} de $Y$ si :

a) le complété $\hat{\underline{O}}_{Y,y}$ de l'anneau local de $y$ dans $Y$ est isomorphe au quotient $k[[T_1,\ldots,T_n]]/(Q(T))$ des series formelles en $n=1+\dim_k Y$ variables par l'idéal engendré par un élément $Q(T)$ .

}{\SGAsectiontoc{1. Un rappel sur les singularités quadratiques (cf VI 6).}{1}
1. \underline{Un rappel sur les singularités quadratiques} (cf VI 6).



1.1. \quad Soit $Y$ un $k$-schéma purement de dimension $n-1$ . Un point fermé $y \in Y$ s'appelle une \underline{singularité quadratique ordinaire} de $Y$ si :

a) le complété $\hat{\underline{O}}_{Y,y}$ de l'anneau local de $y$ dans $Y$ est isomorphe au quotient $k[[T_1,\ldots,T_n]]/(Q(T))$ des séries formelles en $n=1+\dim_k Y$ variables par l'idéal engendré par un élément $Q(T)$ .

}{F7II-XVII-001}b) $Q$ ne commence qu'en degré deux, et la sous-variété de $\mathbb{P}^{\,n-1}_k$ définie par l'annulation de la partie homogène de degré deux de $Q(T)$ est lisse sur $k$ si $n \geq 2$ , et la partie homogène de $Q(T)$ de degré deux est non-nulle si $n = 1$ .

Si la caractéristique de $k$ est différente de deux, \underline{ou} si $n$ est pair, la condition b) équivaut à la condition
\SGApage{II}{214}
b') $Q$ ne commence qu'en degré deux, et la matrice hessienne de $Q$

\begin{equation*}
\left(\ \frac{\partial^2 Q}{\partial T_i \partial T_j}\,(0)\ \right)_{1 \leq i,j \leq n}
\end{equation*}
est inversible,

\SGAvar{Comme on voit en s'appuyant sur le critère jacobienne de lissité.

}{comme on voit en s'appuyant sur le critère jacobien de lissité.

}{F7II-XVII-002}La  "forme standard "  pour la partie homogène de degré deux $Q_2$ d'une série $Q$ vérifiant b), forme qui peut être obtenue après un changement de variable linéaire, est

\SGAvar{\begin{equation}\tag{1.1.1}
\left\{
\begin{array}{l}
Q_2(T) \;=\; \displaystyle\sum_{i=1}^{[n/2]} T_{2i-1}\, T_{2i} \qquad \mathrm{si}\ \ n\ \ \mathrm{est\ pair} \\[4ex]
Q_2(T) \;=\; T_n^2 \;+\; \displaystyle\sum_{i=1}^{\left[\frac{n}{2}\right]} T_{2i-1}\, T_{2i} \quad \mathrm{si}\ \ n\ \ \mathrm{est\ impair,}
\end{array}
\right.
\end{equation}
comme on voit par la théorie des formes quadratique (Bourbaki, Alg.X).

1.2.    Un point fermé $y \in Y$ qui vérifie a), b) s'appelle \underline{singularité quadratique non-dégénérée}. Comme le déterminant de la matrice hessienne est identiquement nul en caractéristique deux si $n$ est impair, les singularités quadratiques non-dégénérées en dimension $n$ n'existent que si, ou bien $n$ est pair, ou bien si $k$ est de caractéristique différente de deux cas dans lesquels ce sont exactement les singularités quadratiques ordinaires.


}{\begin{equation}\tag{1.1.1}
\left\{
\begin{array}{l}
Q_2(T) \;=\; \displaystyle\sum_{i=1}^{[n/2]} T_{2i-1}\, T_{2i} \qquad \mathrm{si}\ \ n\ \ \mathrm{est\ pair} \\[4ex]
Q_2(T) \;=\; T_n^2 \;+\; \displaystyle\sum_{i=1}^{\left[\frac{n}{2}\right]} T_{2i-1}\, T_{2i} \quad \mathrm{si}\ \ n\ \ \mathrm{est\ impair,}
\end{array}
\right.
\end{equation}
comme on voit par la théorie des formes quadratiques (Bourbaki, Alg.IX).

1.2.    Un point fermé $y \in Y$ qui vérifie a), b') s'appelle \underline{singularité quadratique non-dégénérée}. Comme le déterminant de la matrice hessienne est identiquement nul en caractéristique deux si $n$ est impair, les singularités quadratiques non-dégénérées en dimension $n-1$ n'existent que si, ou bien $n$ est pair, ou bien si $k$ est de caractéristique différente de deux, cas dans lesquels ce sont exactement les singularités quadratiques ordinaires.


}{F7II-XVII-003}\SGApage{II}{215}


\SGAsectiontoc{2. Les pinceaux de Lefschetz : énoncé des résultats}{2}
2. \underline{Les pinceaux de Lefschetz : énoncé des résultats}

2.1.    Désignons par $P^r$ l'espace projectif à $r$ dimensions sur $k$ , par $\check{\mathbb{P}}^r$ l'espace projectif \underline{dual} des hyperplans dans $\mathbb{P}^r$ , et par $\mathrm{Gr}(1,\ \check{\mathbb{P}}^r)$ la variété des droites dans $\check{\mathbb{P}}^r$ . Un point $D \in \mathrm{Gr}(1,\ \check{\mathbb{P}}^r)$ s'appelle un \underline{pinceau d'hyperplans} dans $\mathbb{P}^r$ . Considérant $D$ comme une droite dans $\check{\mathbb{P}}^r$ , on écrit, par abus de notation, $D = \{H_t\}_{t \in D}$ , les $H_t$ étant les hyperplans de $\mathbb{P}^r$ qui sont les points de la droite $D$ . On appelle \underline{axe de} $D$ la sous-variété linéaire de $\mathbb{P}^r$ de codimension deux qui est l'intersection $H_o \cdot H_\infty$ de deux éléments distincts quelconques du pinceau. Considérant l'axe de $D$ comme un point de $\mathrm{Gr}(r-2,\ \mathbb{P}^r)$ , la variété des sous-variétés linéaires de dimension $r-2$ dans $\mathbb{P}^r$ , la flèche $D \longrightarrow$ l'axe de $D$ n'est autre que l'isomorphisme bien connu d'orthogonalité

\SGAvar{\begin{equation*}
\mathrm{Gr}(1,\ \check{\mathbb{P}}^r) \ \ \stackrel{\textstyle\sim}{\longrightarrow} \ \ \mathrm{Gr}(r-2,\ \mathbb{P}^r) \qquad .
\end{equation*}
2.2.    Soit $X$ un $k$-schéma, propre, lisse et irréductible, de dimension $n \geq 1$ , muni d'un plongement $X \hookrightarrow \mathbb{P}^r$ . Un pinceau $D = \{H_t\}_{t \in \mathbb{P}}$ d'hyperplans dans $\mathbb{P}^r$ s'appelle un \underline{pinceau de Lefschetz}   (par rapport au plongement  $X \hookrightarrow \mathbf{P}^r$) si

}{\begin{equation*}
\mathrm{Gr}(1,\ \check{\mathbb{P}}^r) \ \ \stackrel{\textstyle\sim}{\longrightarrow} \ \ \mathrm{Gr}(r-2,\ \mathbb{P}^r) \qquad .
\end{equation*}
2.2.    Soit $X$ un $k$-schéma, propre, lisse et irréductible, de dimension $n \geq 1$ , muni d'un plongement $X \hookrightarrow \mathbb{P}^r$ . Un pinceau $D = \{H_t\}_{t \in D}$ d'hyperplans dans $\mathbb{P}^r$ s'appelle un \underline{pinceau de Lefschetz}   (par rapport au plongement  $X \hookrightarrow \mathbf{P}^r$) si

}{F7II-XVII-004}a) l'axe de $D$ coupe transversalement $X$  (EGA IV 17.13.7);

b) il existe un ouvert dense $U$ de $D$ tel que pour tout $t \in U$ , l'hyperplan $H_t$ coupe transversalement $X$ ;

c) pour $t_o \in D-U$ , $H_{t_o}$ coupe transversalement $X$ sauf en
\SGApage{II}{216}
un point, qui est un point singulier quadratique ordinaire de  $X . H_{t_o}$  .

\SGAvar{Nous allons voir plus bas (3.2.1) que les pinceaux de Lefschetz forment un  \underline{ouvert}  de  $\mathrm{Gr}(1, \mathbb{P}^r)$  .

2.3.    Le plongement  $X \hookrightarrow \mathbb{P}^r$  s'appelle un  \underline{plongement de Lefschetz} si les pinceaux de Lefschetz forment un ouvert dense dans $\mathrm{Gr}(1, \check{\mathbb{P}}^r)$ $^{*)}$.

2.4.    Rappellons maintenant la notion de  \underline{multiple d'un plongement donné}  $X \hookrightarrow \mathbb{P}^r$  . C'est le composé du plongement donné  $X \hookrightarrow \mathbb{P}^r$ et du $d$-ième plongement de Segre,  $S_d : \mathbb{P}^r \hookrightarrow \mathbb{P}^{\binom{r+d}{d}-1}$  qui s'écrit en termes des coordonnés homogènes  $X_o,\ldots,X_r$  dans  $\mathbb{P}^r$

\begin{equation*}
(X_o,\ldots,X_r) \xrightarrow{\ \ S_d\ \ } (,\ldots,X^W,\ldots)
\end{equation*}
où  $X^W = X_o^{W_o} \ldots X_r^{W_o}$  parcourt les monômes de degré  $d$  en $X_o,\ldots,X_r$  . Remarquons qu'une "section de  $X$  par une hypersurface de degré  $d$ " $^{**)}$ pour un plongement donné n'est autre qu'une 

}{Nous allons voir plus bas (3.2.1) que les pinceaux de Lefschetz forment un  \underline{ouvert}  de  $\mathrm{Gr}(1, \check{\mathbb{P}}^r)$  .

2.3.    Le plongement  $X \hookrightarrow \mathbb{P}^r$  s'appelle un  \underline{plongement de Lefschetz} si les pinceaux de Lefschetz forment un ouvert dense dans $\mathrm{Gr}(1, \check{\mathbb{P}}^r)$ $^{*)}$.



2.4.    Rappelons maintenant la notion de  \underline{multiple d'un plongement donné}  $X \hookrightarrow \mathbb{P}^r$  . C'est le composé du plongement donné  $X \hookrightarrow \mathbb{P}^r$ et du $d$-ième plongement de Veronese,  $S_d : \mathbb{P}^r \hookrightarrow \mathbb{P}^{\binom{r+d}{d}-1}$  qui s'écrit en termes des coordonnées homogènes  $X_o,\ldots,X_r$  dans  $\mathbb{P}^r$

\begin{equation*}
(X_o,\ldots,X_r) \xrightarrow{\ \ S_d\ \ } (\ldots,X^W,\ldots)
\end{equation*}
où  $X^W = X_o^{W_o} \ldots X_r^{W_r}$  parcourt les monômes de degré  $d$  en $X_o,\ldots,X_r$  . Remarquons qu'une "section de  $X$  par une hypersurface de degré  $d$ " $^{**)}$ pour un plongement donné n'est autre qu'une 

}{F7II-XVII-005}$^{*)}$ou plus correctement, dans l'ensemble des points fermés de $\mathrm{Gr}(1,\check{\mathbb{P}}^r)$. Nous nous permettrons par la suite de tels abus de langage.

$^{**)}$En Anglais, on  dit plus commodément: hypersurface section of degree $d$ .


\SGApage{II}{217}

 "section hyperplane" pour  le $d$-ième multiple de ce plongement.

\underline{Théorème 2.5.}    \underline{Soit}  $X$  \underline{irréductible, propre et lisse sur}  $k$  \underline{de dimension} $n \geq 1$  ,  \underline{muni d'un plongement}  $i : X \hookrightarrow P^r$  . \underline{Alors}

2.5.1. \underline{Tout multiple $d$-ième de}  $i$  ,  \underline{avec}  $d \geq 2$  ,  \underline{est un plongement de Lefschetz.}

\SGAvar{2.5.2. \underline{Si}  $k$  \underline{est de caractéristique nulle}; \underline{le plongement donné}  $i$ \underline{est un plongement de Lefschetz.}

}{2.5.2. \underline{Si}  $k$  \underline{est de caractéristique nulle}, \underline{le plongement donné}  $i$ \underline{est un plongement de Lefschetz.}

}{F7II-XVII-006}La démonstration se coupe en plusieurs morceaux, qui suivent.


\SGAsectiontoc{3. La variété duale}{3}
3. \underline{La variété duale}

3.1.    Soit  $I$  l'Idéal dans  $\mathbb{P}^r$  qui définit  $X$  . On désigne par $N$  le faisceau conormal  $I/I^2$  , faisceau localement libre de rang $r-n$  sur  $X$  , et par  $\mathbb{P}(N)$  le fibré projectif au-dessus de  $X$ formé des droites dans  $\check{N}$  . (Pour  $X = \mathbb{P}^r$  ,  $\mathbb{P}(N)$  est vide.) Le fibré  $\mathbb{P}(N)$  s'identifie à la sous-variété de  $X \times \check{\mathbb{P}}^r$  formée des couples  $(x, H)$  tels que  $H$  soit tangent à  $X$  en  $x$  , par la flèche  $\mathbb{P}(N) \xrightarrow{\ \nu\ } X \times \check{\mathbb{P}}^r$

\begin{equation}\tag{3.1.1}
\mathbb{P}(I/I^2)_x \ni F \xrightarrow{\ \nu\ } (x, \text{l'hyperplan d'équation } \sum \frac{\partial F}{\partial X_i}(x)T_i\ ),
\end{equation}
$F$  désignant une section locale de  $I/I^2$  . On définit

\SGAvar{\begin{equation}\tag{3.1.2}
\varphi : \mathbb{P}(N) \longrightarrow \check{\mathbb{P}}^r
\end{equation}
comme étant le composé  $\mathbb{P}(N) \xrightarrow{\ \nu\ } X \times \check{\mathbb{P}}^r \stackrel{\mathrm{pr}_2}{\relbar\joinrel\relbar\joinrel\relbar} \check{\mathbb{P}}^r$ .


}{\begin{equation}\tag{3.1.2}
\varphi : \mathbb{P}(N) \longrightarrow \check{\mathbb{P}}^r
\end{equation}
comme étant le composé  $\mathbb{P}(N) \xrightarrow{\ \nu\ } X \times \check{\mathbb{P}}^r \xrightarrow{\ \mathrm{pr}_2\ } \check{\mathbb{P}}^r$ .


}{F7II-XVII-007}\SGApage{II}{218}

3.1.3. On désigne par $\check{X}$ l'image de $\varphi$ ; c'est l'ensemble des hyperplans qui ne coupent pas transversalement $X$ , et s'appelle la \underline{variété duale} de $X$ (plongé dans $\mathbb{P}^r$ par $i$) .

\underline{Proposition} 3.1.4. \underline{La variété duale} $\check{X}$ \underline{est irréductible et propre de dimension} $\leq r-1$ .

\underline{Démonstration.} C'est l'image par $\varphi$ de $\mathbb{P}(N)$ , lui-même irréductible de dimension $r-1$ si $X \neq \mathbb{P}^r$ , et vide si $X = \mathbb{P}^r$ .

\SGAvar{\underline{Remarque} 3.1.5. Il est possible de préciser l'interprétation de $\mathbb{P}(N)$ de la façon suivante. Soit $\underline{Y} \subset X \times \check{\mathbb{P}}^r$ le sous-schéma de $X \times \check{\mathbb{P}}^r$ ``d'incidence'' , dont la fibre en tout point $H$ de $\check{\mathbb{P}}^r(k)$ est $X.H$. Il est immédiat que $\underline{Y}$ est lisse de dimension $n+r-1$ , comme fibré projectif de dimension relative $r-1$ sur $X$ , donc $Y \longrightarrow \check{\mathbb{P}}^r$ fait de $\underline{Y}$ une intersection complète relative de dimension relative virtuelle $(n+r-1)-r = n-1$ . Ceci posé, le \underline{sous-schéma jacobien} $\underline{J}_{n-1}(\underline{Y}/\check{P}^r)$ (VI 5.2 ) \underline{n'est autre que} $\mathbb{P}(N)$ . Grâce à la compatibilité de la formation du sous-schéma jacobien d'une intersection complète relative avec tout changement de base, on en déduit que pour toute droite $D$ dans $\check{\mathbb{P}}^r$ , le sous-schéma jacobien $\underline{J}_{n-1}(\tilde{X}/D)$ , où $\tilde{X} = \underline{Y} \times_{\check{\mathbb{P}}^r} D$ , n'est autre que $\mathbb{P}(N) \times_{\check{\mathbb{P}}^r} D$ .

\underline{Proposition} 3.2. \underline{Le sous-ensemble} $B$ \underline{de} $\check{X}$ \underline{formé des hyperplans} $H$ \underline{tel que} $X \cdot H$ \underline{ait un et un seul point singulier, et que ce point}
\SGApage{II}{219}
\underline{soit quadratique ordinaire, est ouvert} .

C'est un cas particulier de XV 1.3.2, , s'appuyant la théorie des ``modules henséliens'' de R. ELKIK (XV 1.1.4). Le cas où on se restreint aux points quadratiques \underline{non dégénérés} (i.e. où $\mathrm{car}(k) \neq 2$ \underline{ou} $n=\dim X$ est paire)est d'ailleurs évident grâce à 3.3 ci-dessous (dont la démonstration n'utilise pas 3.2).

}{\underline{Remarque} 3.1.5. Il est possible de préciser l'interprétation de $\mathbb{P}(N)$ de la façon suivante. Soit $\underline{Y} \subset X \times \check{\mathbb{P}}^r$ le sous-schéma de $X \times \check{\mathbb{P}}^r$ ``d'incidence'' , dont la fibre en tout point $H$ de $\check{\mathbb{P}}^r(k)$ est $X.H$. Il est immédiat que $\underline{Y}$ est lisse de dimension $n+r-1$ , comme fibré projectif de dimension relative $r-1$ sur $X$ , donc $\underline{Y} \longrightarrow \check{\mathbb{P}}^r$ fait de $\underline{Y}$ une intersection complète relative de dimension relative virtuelle $(n+r-1)-r = n-1$ . Ceci posé, le \underline{sous-schéma jacobien} $\underline{J}_{n-1}(\underline{Y}/\check{P}^r)$ (VI 5.2 ) \underline{n'est autre que} $\mathbb{P}(N)$ . Grâce à la compatibilité de la formation du sous-schéma jacobien d'une intersection complète relative avec tout changement de base, on en déduit que pour toute droite $D$ dans $\check{\mathbb{P}}^r$ , le sous-schéma jacobien $\underline{J}_{n-1}(\tilde{X}/D)$ , où $\tilde{X} = \underline{Y} \times_{\check{\mathbb{P}}^r} D$ , n'est autre que $\mathbb{P}(N) \times_{\check{\mathbb{P}}^r} D$ .

\underline{Proposition} 3.2. \underline{Le sous-ensemble} $B$ \underline{de} $\check{X}$ \underline{formé des hyperplans} $H$ \underline{tels que} $X \cdot H$ \underline{ait un et un seul point singulier, et que ce point}
\SGApage{II}{219}
\underline{soit quadratique ordinaire, est ouvert} .

C'est un cas particulier de XV 1.3.2, s'appuyant sur la théorie des ``modules henséliens'' de R. ELKIK (XV 1.1.4). Le cas où on se restreint aux points quadratiques \underline{non dégénérés} (i.e. où $\mathrm{car}(k) \neq 2$ \underline{ou} $n=\dim X$ est paire)est d'ailleurs évident grâce à 3.3 ci-dessous (dont la démonstration n'utilise pas 3.2).

}{F7II-XVII-008}\underline{Corollaire} 3.2.1. \underline{Sous les hypothèses de} 3.2, \underline{désignons par} $F$ \underline{le fermé} $\check{X}-B$ \underline{de} $\check{X}$ . \underline{Alors les pinceaux de Lefschetz forment un ouvert dans} $\mathrm{Gr}(1, \check{\mathbb{P}}^r)$ , \underline{et les conditions suivantes sont équivalentes:}

3.2.2. \underline{Le plongement donné} $i : X \hookrightarrow \mathbb{P}^r$ \underline{est un plongement de Lefschetz;}

3.2.3. \underline{il existe un pinceau de Lefschetz d'hyperplans (pour le plongement donné).}

3.2.4. \underline{La codimension de} $F$ \underline{dans} $\check{\mathbb{P}}^r$ \underline{est} $\geq 2$ .

\underline{Démonstration.} Une droite $D$ dans $\check{\mathbb{P}}^r$ est un pinceau de Lefschetz si et seulement si elle vérifie:

(3.2.5) L'axe de $D$ coupe transversalement $X$ ,

\SGAvar{(3.2.6) $D$ n'est pas contenue dans $X$ ,

}{(3.2.6) $D$ n'est pas contenue dans $\check{X}$ ,

}{F7II-XVII-009}(3.2.7) $D$ ne rencontre pas $F$ .


\SGApage{II}{220}

\SGAvar{Les conditions (3.2.5) et (3.2.6) définissent des ouverts de $\mathrm{Gr}(1, \check{\mathbb{P}}^{r})$ , et, $F$ étant fermé , il en est de même pour (3.2.7), donc les pinceaux de Lefschetz forment un ouvert dans $\mathrm{Gr}(1, \check{\mathbb{P}}^{r})$ . Parce que $\dim X \leq r-1$ , l'ouvert (3.2.6) est non-vide, et donc dense. L'ouvert (3.2.5) est également non-vide, donc dense. En effet, soit $H$ un hyperplan qui coupe transversalement $X$ ; un tel $H$ existe car $\dim \check{X} \leq r-1$ . Parce que $\dim \check{(X \cdot G)} \leq r-1$ , il existe un hyperplan $G$ qui coupe transversalement $X \cdot H$ , et par construction la droite $D$ passant par $H$ et $G$ a une axe qui coupe transversalement $X$ . Donc le plongement est de Lefschetz si et seulement si l'ouvert (3.2.7) est non-vide (et donc dense). Or, $F$ étant fermé, (3.2.7) est non-vide si et seulement si $F$ est de codimension $\geq 2$ dans $\check{\mathbb{P}}^{r}$ (cf. EGA V, ou Mumford, Introduction to Algebraic Geometry, p. 88, Cor 3).

}{Les conditions (3.2.5) et (3.2.6) définissent des ouverts de $\mathrm{Gr}(1, \check{\mathbb{P}}^{r})$ , et, $F$ étant fermé , il en est de même pour (3.2.7), donc les pinceaux de Lefschetz forment un ouvert dans $\mathrm{Gr}(1, \check{\mathbb{P}}^{r})$ . Parce que $\dim \check{X} \leq r-1$ , l'ouvert (3.2.6) est non-vide, et donc dense. L'ouvert (3.2.5) est également non-vide, donc dense. En effet, soit $H$ un hyperplan qui coupe transversalement $X$ ; un tel $H$ existe car $\dim \check{X} \leq r-1$ . Parce que $\dim \check{(X \cdot H)} \leq r-1$ , il existe un hyperplan $G$ qui coupe transversalement $X \cdot H$ , et par construction la droite $D$ passant par $H$ et $G$ a un axe qui coupe transversalement $X$ . Donc le plongement est de Lefschetz si et seulement si l'ouvert (3.2.7) est non-vide (et donc dense). Or, $F$ étant fermé, (3.2.7) est non-vide si et seulement si $F$ est de codimension $\geq 2$ dans $\check{\mathbb{P}}^{r}$ (cf. EGA V, ou Mumford, Introduction to Algebraic Geometry, p. 88, Cor 3).

}{F7II-XVII-010}\underline{Corollaire} 3.2.8. \underline{Sous les hypothèses de} 3.2., \underline{supposons que} $i : X \hookrightarrow \mathbb{P}^{r}$ \underline{soit un plongement de Lefschetz. Soit} $H_{o}$ \underline{un hyperplan qui coupe transversalement} $X$ , \underline{et désignons par} $\mathrm{Gr}(1, \check{\mathbb{P}}^{r})_{H_{o}}$ \underline{la sous-variété de} $\mathrm{Gr}(1, \check{\mathbb{P}}^{r})$ \underline{formée des droites passant par} $H_{o}$ . \underline{Alors les pinceaux de Lefschetz passant par} $H_{o}$ \underline{forment un ouvert dense dans} $\mathrm{Gr}(1, \check{\mathbb{P}}^{r})_{H_{o}}$ .

\SGAvar{\underline{Démonstration.} D'après (3.2.1), les-dits pinceaux forment un ouvert dans $\mathrm{Gr}(1, \check{\mathbb{P}}^{r})_{H_{o}}$ . L'argument qu'on vient de donner pour
\SGApage{II}{221}
établir que les ouverts (3.2.5) et (3.2.6) sont non-vides démontre que leurs intersections avec $\mathrm{Gr}(1, \check{\mathbb{P}}^{r})_{H_{o}}$ sont également non-vides. Il reste à prouver qu'il y a des droites passant par $H_{o}$ qui ne rencontrent pas $F$ . Pour ceci, soit $T$ un hyperplan dans $\check{\mathbb{P}}^{r}$ qui ne contient pas $H_{o}$ . On a un isomorphisme

}{\underline{Démonstration.} D'après (3.2.1), lesdits pinceaux forment un ouvert dans $\mathrm{Gr}(1, \check{\mathbb{P}}^{r})_{H_{o}}$ . L'argument qu'on vient de donner pour
\SGApage{II}{221}
établir que les ouverts (3.2.5) et (3.2.6) sont non-vides démontre que leurs intersections avec $\mathrm{Gr}(1, \check{\mathbb{P}}^{r})_{H_{o}}$ sont également non-vides. Il reste à prouver qu'il y a des droites passant par $H_{o}$ qui ne rencontrent pas $F$ . Pour ceci, soit $T$ un hyperplan dans $\check{\mathbb{P}}^{r}$ qui ne contient pas $H_{o}$ . On a un isomorphisme

}{F7II-XVII-011}\begin{center}
\begin{tikzcd}[column sep=huge]
\mathrm{Gr}(1, \check{\mathbb{P}}^{r})_{H_{o}} \& T \arrow[l, "\sim"'] \\
\text{la droite passant par } H_{o} \text{ et } \tau \& \tau \arrow[l, mapsto]
\end{tikzcd}
\end{center}









\SGAvar{donc $\mathrm{Gr}(1, \check{\mathbb{P}}^{r})_{H_{o}}$ ``est'' un espace projectif de dimension $r-1$. Les droites la-dedans passant par un point de $F$ y forment un fermé de dimension $\leq r-2$ , étant l'image du morphisme

}{donc $\mathrm{Gr}(1, \check{\mathbb{P}}^{r})_{H_{o}}$ ``est'' un espace projectif de dimension $r-1$. Les droites là-dedans passant par un point de $F$ y forment un fermé de dimension $\leq r-2$ , étant l'image du morphisme

}{F7II-XVII-012}\begin{center}
\begin{tikzcd}[column sep=huge]
F \arrow[r] \& \mathrm{Gr}(1, \check{\mathbb{P}}^{r})_{H_{o}} \\
f \arrow[r, mapsto] \& \text{la droite passant par } H_{o} \text{ et par } f \; .
\end{tikzcd}
\end{center}







\SGAvar{\underline{Proposition} 3.3. \underline{Le morphisme} $\varphi : \mathbb{P}(N) \longrightarrow \check{\mathbb{P}}^{r}$ \underline{est net (= non ramifié) en un point fermé} $(x_{o}, H)$ \underline{si et seulement si le point} $x_{o}$ \underline{est une singularité quadratique non-dégénérée} (1.2) \underline{de} $X \cdot H$ . \underline{En particulier,} \underline{le sous-ensemble de} $\mathbb{P}(N)$ \underline{formé des points} $(x, H)$ \underline{tel que} $x$ \underline{soit une singularité quadratique non-dégénérée de} $X \cdot H$ \underline{est ouvert.}


}{\underline{Proposition} 3.3. \underline{Le morphisme} $\varphi : \mathbb{P}(N) \longrightarrow \check{\mathbb{P}}^{r}$ \underline{est net (= non ramifié) en un point fermé} $(x_{o}, H)$ \underline{si et seulement si le point} $x_{o}$ \underline{est une singularité quadratique non-dégénérée} (1.2) \underline{de} $X \cdot H$ . \underline{En particulier,} \underline{le sous-ensemble de} $\mathbb{P}(N)$ \underline{formé des points} $(x, H)$ \underline{tels que} $x$ \underline{soit une singularité quadratique non-dégénérée de} $X \cdot H$ \underline{est ouvert.}


}{F7II-XVII-013}\SGApage{II}{222}

\SGAvar{\underline{Démonstration}. Nous pouvons nous borner au cas $X \neq \mathbb{P}^r$ . Les propriétés prétendues équivalentes étant indépendantes du choix particulier des coordonnés, on peut les expliciter à l'aide de  coordonnés choisies d'une façon très particulière. Or, $X$ étant une sous-variété lisse de $\mathbb{P}^r$ de dimension $n \geq 1$ , on peut choisir des coordonnés homogènes $X_o,\dots,X_r$ dans $\mathbb{P}^r$ de telle façon que

}{\underline{Démonstration}. Nous pouvons nous borner au cas $X \neq \mathbb{P}^r$ . Les propriétés prétendues équivalentes étant indépendantes du choix particulier des coordonnées, on peut les expliciter à l'aide de  coordonnées choisies d'une façon très particulière. Or, $X$ étant une sous-variété lisse de $\mathbb{P}^r$ de dimension $n \geq 1$ , on peut choisir des coordonnées homogènes $X_o,\dots,X_r$ dans $\mathbb{P}^r$ de telle façon que

}{F7II-XVII-014}3.3.1.  le point $x_o \in \mathbb{P}^r$ soit $(1, 0,\dots,0)$ ;

3.3.2.  posant $x_i = X_i/X_o$ pour $i = 1,\dots,r$ , de sorte que le complété $\hat{\underline{O}}_{\mathbb{P}^r,x_o}$ s'identifie à $k[[x_1,\dots,x_r]]$ , l'idéal définissant la trace de $X$ dans $\mathrm{Spec}\ \hat{\underline{O}}_{\mathbb{P}^r,x_o}$ est engendré par des éléments $g_i$ $(i=1,\dots,r-n)$ dans $\underline{O}_{\mathbb{P}^r,x}$ qui, \underline{dans} $\hat{\underline{O}}_{\mathbb{P}^r,\, x_o}$ s'écrivent

\begin{equation*}
g_i = x_{n+i} - h_i(x_1,\dots,x_n) \quad , \quad i = 1,\dots,r-n
\end{equation*}
où les $h_i(x_1,\dots,x_n)$ sont des séries en $x_1,\dots,x_n$ seulement qui se trouvent dans le carré de l'idéal maximal de $k[[x_1,\dots,x_n]]$ .

\SGAvar{3.3.3.  l'hyperplan $H$ provient, via (3.1.1) de $g_{n-r}$ .

}{3.3.3.  l'hyperplan $H$ provient, via (3.1.1) de $g_{r-n}$ .

}{F7II-XVII-015}Il résulte de (3.3.2) que, posant $V = \mathrm{Spec}(\hat{\underline{O}}_{X,x_o})$, il y a une trivialisation

\SGAvar{\begin{equation*}
(I/I^2)|V \;\simeq\; \underline{O}_V^{\,r-n}
\end{equation*}
\begin{equation}\tag{3.3.4}
(\lambda_1,\dots,\lambda_{r-n}) \longmapsto \sum_1^{n-r} \lambda_i g_{n+i}
\end{equation}
Il résulte de ceci et de (3.3.3) qu'il existe un voisinage $W$ de $(x, H)$
\SGApage{II}{223}
dans $\mathbb{P}(N)|V$ tel qu'on ait $V \times \mathbb{A}^{r-n-1} \simeq W$ , l'isomorphisme étant donné par

\begin{equation}\tag{3.3.5}
(v,(\lambda_1,\dots,\lambda_{r-n-1})) \in V\times\mathbb{A}^{r-n-1} \longmapsto (v,(g_{r-n}+ \sum_1^{n-r-1} \lambda_i g_i) \in \mathbb{P}(I/I^2)|V
\end{equation}
(où l'on pose $\lambda_i = \Lambda_i / \Lambda_{n-r}$ ).

En particulier, le complété $\hat{\underline{O}}_{\mathbb{P}(\check{N}),(x,H)}$ de l'anneau local de $\mathbb{P}(N)$ en $(x, H)$ s'identifie à $k[[x_1,\dots,x_n,\lambda_1,\dots,\lambda_{r-n-1}]]$ . Exprimons maintenant $\varphi$ en termes de ces coordonnés, au niveau des complétés:

\begin{equation}\tag{3.3.6}
\varphi(x,\lambda) = \varphi\, (x_1,\dots,x_n,\lambda_1,\dots,\lambda_{r-n-1})
\end{equation}
\begin{equation*}
= \mbox{l'hyperplan d'équation} \quad \sum_{i=0}^{n+r} \varphi_i(x,\lambda)T_i = 0 \quad .
\end{equation*}
Comme l'hyperplan image $\sum \varphi_i\, T_i$ contient le point $(x)$ , point qui s'écrit en coordonnés homogènes

\begin{equation}\tag{3.3.7}
(x) = (1,x_1,\dots,x_n,\, h_1(x),\dots,h_{r-n}(x) \quad ,
\end{equation}
on a

\begin{equation}\tag{3.3.8}
\varphi_o(x,\lambda)= - \sum_{i=1}^{n} x_i\varphi_i(x,\lambda)- \sum_{i=1}^{r-n} h_i(x)\varphi_{n+i}(x,\lambda)
\end{equation}
dans $\hat{\underline{O}}_{\mathbb{P}(N),(x_o,H)}$ et, d'après (3.1.1), pour $i \geq 1$ , posant $D_i = \frac{\partial}{\partial x_i} \in \mathrm{Der}_k(\hat{\underline{O}}_{\mathbb{P}^r,x_o},\, \hat{\underline{O}}_{\mathbb{P}^r,x_o})$ :

\begin{equation}\tag{3.3.9}
\varphi_i(x,\lambda) = D_i(g_{r-n}+ \sum_{i=1}^{r-n-1} \lambda_i\, g_i)\, (x) \quad \mbox{dans} \quad \hat{\underline{O}}_{\mathbb{P}(N),(x_o H)}
\end{equation}

}{\begin{equation*}
(I/I^2)|V \;\simeq\; \underline{O}_V^{\,r-n}
\end{equation*}
\begin{equation}\tag{3.3.4}
(\Lambda_1,\dots,\Lambda_{r-n}) \longmapsto
\sum_{i=1}^{r-n}\Lambda_i g_i
\end{equation}
Il résulte de ceci et de (3.3.3) qu'il existe un voisinage $W$ de $(x, H)$
\SGApage{II}{223}
dans $\mathbb{P}(N)|V$ tel qu'on ait $V \times \mathbb{A}^{r-n-1} \simeq W$ , l'isomorphisme étant donné par

\begin{equation}\tag{3.3.5}
(v,(\lambda_1,\dots,\lambda_{r-n-1})) \in
V\times\mathbb{A}^{r-n-1} \longmapsto
\left(v,\left[g_{r-n}+\sum_{i=1}^{r-n-1}\lambda_i g_i\right]\right)
\in \mathbb{P}(I/I^2)|V
\end{equation}
(où l'on pose $\lambda_i = \Lambda_i/\Lambda_{r-n}$ ).

En particulier, le complété $\hat{\underline{O}}_{\mathbb{P}(N),(x,H)}$ de l'anneau local de $\mathbb{P}(N)$ en $(x, H)$ s'identifie à $k[[x_1,\dots,x_n,\lambda_1,\dots,\lambda_{r-n-1}]]$ . Exprimons maintenant $\varphi$ en termes de ces coordonnées, au niveau des complétés:

\begin{equation}\tag{3.3.6}
\varphi(x,\lambda)
= \mbox{l'hyperplan d'équation} \quad
\sum_{i=0}^{r}\varphi_i(x,\lambda)T_i=0 \quad .
\end{equation}
Comme l'hyperplan image $\sum \varphi_i\, T_i$ contient le point $(x)$ , point qui s'écrit en coordonnées homogènes

\begin{equation}\tag{3.3.7}
(x) = (1,x_1,\dots,x_n,\, h_1(x),\dots,h_{r-n}(x)) \quad ,
\end{equation}
on a

\begin{equation}\tag{3.3.8}
\varphi_o(x,\lambda)= - \sum_{i=1}^{n} x_i\varphi_i(x,\lambda)- \sum_{j=1}^{r-n} h_j(x)\varphi_{n+j}(x,\lambda)
\end{equation}
dans $\hat{\underline{O}}_{\mathbb{P}(N),(x_o,H)}$ et, d'après (3.1.1), pour $i \geq 1$ , posant $D_i = \frac{\partial}{\partial x_i} \in \mathrm{Der}_k(\hat{\underline{O}}_{\mathbb{P}^r,x_o},\, \hat{\underline{O}}_{\mathbb{P}^r,x_o})$ :

\begin{equation}\tag{3.3.9}
\varphi_i(x,\lambda) = D_i(g_{r-n}+ \sum_{j=1}^{r-n-1} \lambda_j\, g_j)\, (x) \quad \mbox{dans} \quad \hat{\underline{O}}_{\mathbb{P}(N),(x_o,H)}
\end{equation}

}{F7II-XVII-016}\SGApage{II}{224}

\SGAvar{Rappellant  $g_i = x_{n+i} - h_i(x_1,\dots,x_n)$  , on a

\begin{equation}\tag{3.3.10}
\varphi_i(x,\lambda)= -\frac{\partial}{\partial x_i}\left(h_{r-n} + \sum_{i=1}^{r-n-1} \lambda_i h_i\right) \qquad \mbox{pour} \quad 1 \leq i \leq n \;\; ,
\end{equation}
\begin{equation}\tag{3.3.11}
\varphi_i(x,\lambda) = \lambda_i \qquad \mbox{pour} \quad 1 \leq i \leq n \;\; ,
\end{equation}
et

}{Rappelant  $g_i = x_{n+i} - h_i(x_1,\dots,x_n)$  , on a

\begin{equation}\tag{3.3.10}
\varphi_i(x,\lambda)= -\frac{\partial}{\partial x_i}\left(h_{r-n} + \sum_{j=1}^{r-n-1} \lambda_j h_j\right) \qquad \mbox{pour} \quad 1 \leq i \leq n \;\; ,
\end{equation}
\begin{equation}\tag{3.3.11}
\varphi_{n+j}(x,\lambda) = \lambda_j \qquad \mbox{pour} \quad 1 \leq j \leq r-n-1 \;\; ,
\end{equation}
et

}{F7II-XVII-017}\begin{equation}\tag{3.3.12}
\varphi_r(x,\lambda) = 1
\end{equation}
Désignons par  $m$  l'idéal maximal de  $\hat{\underline{O}}_{\mathbb{P}(N),\,(x_o,H)}$  . Compte tenu de ce que  $h_i \equiv 0 \bmod m^2$  pour  $i = 1,\dots,r-n$  , on trouve

\begin{equation}\tag{3.3.13}
\varphi_o(x,\lambda) \equiv 0 \bmod m^2
\end{equation}
\begin{equation}\tag{3.3.14}
\varphi_i(x,\lambda) \equiv \; -\frac{\partial h_{r-n}}{\partial x_i} \; \bmod m^2 \qquad \mbox{pour} \quad i=1,\dots,n \;\; .
\end{equation}
Mettant ensemble les formules (3.3.11) - (3.3.14), on trouve que la matrice jacobienne de  $\varphi$ en le point  $(x_o,H)$  est 


\SGApage{II}{225}

\begin{equation}\tag{3.3.15}
 \begin{array}{r@{\;}l}
 & \hspace{1.2cm} n \hspace{3.6cm} r-n-1 \\[0.5ex]
\begin{array}{r}
1\,\{ \\[1ex]
n\,\Bigg\{\vphantom{\displaystyle\frac{\partial^2 h_{r-n}}{\partial x_1 \partial x_1}(0)} \\[1.5ex]
\vphantom{\displaystyle\frac{\partial^2 h_{r-n}}{\partial x_1 \partial x_n}(0)} \\[1.5ex]
\\
\\
r-n-1\,\Bigg\{ \\
\\
\\
\\
1\,\{
\end{array}
&
\left(
\begin{array}{ccccccccccc}
0 & & 0 & & 0 & & . & . & . & . & 0 \\[1ex]
-\displaystyle\frac{\partial^2 h_{r-n}}{\partial x_1 \partial x_1}(0) & .\;.\;.\;.\;. & -\displaystyle\frac{\partial^2 h_{r-n}}{\partial x_n \partial x_1}(0) & .\;.\;. & 0 & . & . & . & . & . & 0 \\[2ex]
-\displaystyle\frac{\partial^2 h_{r-n}}{\partial x_1 \partial x_n}(0) & .\;.\;.\;. & -\displaystyle\frac{\partial^2 h_{r-n}}{\partial x_n \partial x_n}(0) & & 0 & & . & . & . & . & 0 \\[2ex]
0 & & 0 & & 1 & 0 & & & & & 0 \\
0 & & . & & 0 & 1 & & 0 & . & . & 0 \\
 & & & & & 0 & & . & & & \\
 & & & & & & & & . & & . \\
 & & & & & & & & & . & 0 \\
0 & & 0 & & 0 & 0 & & . & . & 0 & 1 \\
0 & .\;.\;.\;.\;.\;.\;. & . & & . & . & & . & . & & 0
\end{array}
\right)
\end{array}
\end{equation}
Donc  $\varphi$ est net en  $(x_o,H)$  si et seulement si la matrice hessienne

\begin{equation*}
\left( \displaystyle\frac{\partial^2 h_{r-n}}{\partial x_i \; \partial x_j}(0) \right) \qquad 1 \leq i, \; j \leq n
\end{equation*}
est inversible. D'autre part, (3.3.2) et (3.3.3) nous assurent que le complété de l'anneau local de  $x_o$  en  $X \cdot H$  est

\begin{equation}\tag{3.3.16}
k[[x_1,\dots,x_n]] \,/\, (h_{r-n}(x_1,\dots,x_n)) 
\end{equation}

\SGApage{II}{226}

 et on conclut grâce à la définition (1.2).

\SGAvar{\underline{Exemple} 3.4. Supposons que $X$ soit une hypersurface dans $\mathbb{P}^r$ , définie par une équation homogène $F(X_o,\dots,X_r) = 0$ . Alors $\mathbb{P}(N) = X$ , et $\varphi : X \longrightarrow \check{\mathbb{P}}^r$ s'exprime, en coordonnés homogènes, par

}{\underline{Exemple} 3.4. Supposons que $X$ soit une hypersurface dans $\mathbb{P}^r$ , définie par une équation homogène $F(X_o,\dots,X_r) = 0$ . Alors $\mathbb{P}(N) = X$ , et $\varphi : X \longrightarrow \check{\mathbb{P}}^r$ s'exprime, en coordonnées homogènes, par

}{F7II-XVII-018}\begin{equation}\tag{3.4.1}
\varphi(X_o,\dots,X_r) \;=\; \left( \frac{\partial F}{\partial X_o},\dots, \frac{\partial F}{\partial X_r} \right) \;.
\end{equation}
C'est le \underline{``morphisme de GAUSS''} .

Donnons un exemple où $\varphi$ est partout ramifié (i.e. non net). Pour ceci, soit $p$ la caractéristique de $k$ , $r$ un entier $\geqslant 1$ , et prenons l'hypersurface $X$ d'équation

\begin{equation}\tag{3.4.2}
\sum_{i=0}^{r} X_i^{p^r+1} \;=\; 0 \;,
\end{equation}
pour laquelle on a

\SGAvar{\begin{equation}\tag{3.4.3}
\varphi(X_o,\dots,X_r) = (X_o^{p^r},\dots,X_r^{p^r})
\end{equation}
de sorte que $\check{X}$ est reduit à un point si $p = 0$ , et $\check{X} \simeq X$ , avec $\varphi$ le morphisme de Frobenius si $p > 0$ . On observera qu'il résulte de 3.3 (ou de 3.5.0 plus bas) que si $p > 2$ , le plongement envisagé n'est pas un plongement de Lefschetz.

}{\begin{equation}\tag{3.4.3}
\varphi(X_o,\dots,X_r) = (X_o^{p^r},\dots,X_r^{p^r})
\end{equation}
de sorte que $\check{X}$ est réduit à un point si $p = 0$ , et $\check{X} \simeq X$ , avec $\varphi$ le morphisme de Frobenius si $p > 0$ . On observera qu'il résulte de 3.3 (ou de 3.5.0 plus bas) que si $p > 2$ , le plongement envisagé n'est pas un plongement de Lefschetz.

}{F7II-XVII-019}Reprenons le cas général.


\SGApage{II}{227}

\SGAvar{\underline{Proposition} 3.5. \underline{Supposons que le morphisme} $\varphi : \mathbb{P}(N) \longrightarrow \check{\mathbb{P}}^r$ \underline{soit} \underline{génériquement net} (i.e. \underline{ne soit pas partout ramifié}). \underline{Alors} $\varphi$ \underline{induit} \underline{un morphisme birationnel de} $\mathbb{P}(N)$ \underline{sur} $\check{X}$ . \underline{Plus précisement, les trois} \underline{ouverts suivants de} $X$ \underline{sont identiques}:

}{\underline{Proposition} 3.5. \underline{Supposons que le morphisme} $\varphi : \mathbb{P}(N) \longrightarrow \check{\mathbb{P}}^r$ \underline{soit} \underline{génériquement net} (i.e. \underline{ne soit pas partout ramifié}). \underline{Alors} $\varphi$ \underline{induit} \underline{un morphisme birationnel de} $\mathbb{P}(N)$ \underline{sur} $\check{X}$ . \underline{Plus précisément, les trois} \underline{ouverts suivants de} $\check{X}$ \underline{sont identiques}:

}{F7II-XVII-020}a) \underline{le plus grand ouvert de lissité de} $\check{X}$ ,

b) \underline{le plus grand ouvert} $V$ \underline{de} $\check{X}$ \underline{tel que} $\varphi : \mathbb{P}(N) \longrightarrow \check{X}$ \underline{induise} \underline{un isomorphisme} $\varphi^{-1}(V) \stackrel{\sim}{\longrightarrow} V$ ,

c) \underline{l'ouvert} (cf. 3.2 ) \underline{de} $\check{X}$ \underline{formé des hyperplans} $H$ \underline{tels que} $X \cdot H$ \underline{ait un et un seul point singulier,} \underline{ce dernier étant quadratique non-dégénéré.}

\underline{Corollaire} 3.5.0 \underline{Supposons} car.$k \neq 2$ \underline{ou} $\dim X$ \underline{paire.} \underline{Le plongement} \underline{donné est un plongement de Lefschetz si et seulement si l'une des deux} \underline{conditions est satisfaite}: a) $\varphi$ \underline{est génériquement net,} \underline{ou} b) $\dim \check{X} \leqslant r-2$ .

En effet, la suffisance résulte du critère (3.2.4), car, par 3.5, on a ou $F =$ lieu sing$(\check{X})$ , donc $\dim F < \dim(\check{X}) = r-1$ , ou $\dim(\check{X}) \leqslant r-2$ . Pour la nécessité, on observe que si $D$ est un pinceau de Lefschetz et si $\dim \check{X} = r-1$ , il existe $x \in D \cap \check{X}$ , et par 3.3 et la définition (2.2 c)) $\varphi$ est net en les points de $\varphi^{-1}(x)$ , donc génériquement net, cqfd.

\underline{Démonstration de} 3.5 L'hypothèse implique $\dim \check{X} = \dim \mathbb{P}(N)(=r-1)$ , d'où résulte (3.2) que $\check{X}$ est une hypersurface irréductible dans $\check{\mathbb{P}}^r$ .


\SGApage{II}{228}

Soit $R$ le lieu de ramification de $\varphi$ , qui est par hypothèse un fermé de $\mathbb{P}(N)$ de dimension au plus $r-2$ , de sorte que $\varphi(R)$ est un fermé de $\check{X}$ de dimension au plus $r-2$ . Désignons par $V'$ l'ouvert dense de $\check{X}$

\begin{equation}\tag{3.5.1}
V' = (\check{X} - \varphi(R)) \cap (\check{X} - \text{lieu sing}(\check{X})) ,
\end{equation}
et par $\underline{F}$ son complémentaire dans $\check{X}$

\SGAvar{\begin{equation}\tag{3.5.1.1}
\underline{F} = \varphi(R) \cup \text{lieu sing}(\check{X}) .
\end{equation}
Nous allons prouver d'abord que $\varphi$ induit un isomorphisme de $\varphi^{-1}(V')$ avec $V'$ par la méthode naive de construire un inverse $\psi : V' \longrightarrow \varphi^{-1}(V')$. En effet, considérant $\varphi^{-1}(V')$ comme plongé dans $X \times V'$ via (3.1.1) , de sorte que $\varphi$ s'identifie à $\mathrm{pr}_2$ , on voit que la construction de $\psi$ équivaut à la construction d'un morphisme $\psi_1 : V' \longrightarrow \mathbb{P}^r$ tel que le diagramme

}{\begin{equation}\tag{3.5.1.1}
\underline{F} = \varphi(R) \cup \text{lieu sing}(\check{X}) .
\end{equation}
Nous allons prouver d'abord que $\varphi$ induit un isomorphisme de $\varphi^{-1}(V')$ avec $V'$ par la méthode naïve de construire un inverse $\psi : V' \longrightarrow \varphi^{-1}(V')$. En effet, considérant $\varphi^{-1}(V')$ comme plongé dans $X \times V'$ via (3.1.1) , de sorte que $\varphi$ s'identifie à $\mathrm{pr}_2$ , on voit que la construction de $\psi$ équivaut à la construction d'un morphisme $\psi_1 : V' \longrightarrow \mathbb{P}^r$ tel que le diagramme

}{F7II-XVII-021}\begin{equation}\tag{3.5.2}
\begin{tikzcd}
\mathbb{P}(N) \supset \varphi^{-1}(V') \arrow[r, "\varphi"] \arrow[d] \& V' \arrow[d, "\psi_1"] \\
X \arrow[r, hook] \& \mathbb{P}^r
\end{tikzcd}
\end{equation}












soit commutatif (on pose $\psi(v) = (\psi_1(v), v) \in \mathbb{P}^r \times V'$) .

\SGAvar{Pour ceci, remarquons que pour une hypersurface quelconque dans un espace projectif, le morphisme de GAUSS (3.4.1) est ``défini'' en tout point lisse. Appliquons ceci a l'hypersurface $\check{X}$ dans $\check{\mathbb{P}}^r$ ; comme $V'$ est
\SGApage{II}{229}
contenu dans $V$ , l'ouvert de lissité de $\check{X}$ , le morphisme de GAUSS nous donne un morphisme, noté $\psi_1$ :

\begin{equation}\tag{3.5.3}
\psi_1 : V' \longrightarrow \check{\check{\mathbb{P}}}^r \simeq \mathbb{P}^r .
\end{equation}
Il reste à vérifier la commutativité du diagramme (3.5.2). Soit alors $(x_o, H_o)$ un point fermé $\varphi^{-1}(V')$ ; il faut vérifier que $x_o$ , considéré comme étant un hyperplan dans $\check{\mathbb{P}}^r$ , est tangent à $\check{X}$ en $H_o$ , i.e. on va démontrer que pour $(x_o, H_o) \in \varphi^{-1}(V')$ , $H_o$ tangent à $X$ en $x_o \Rightarrow x_o$ tangent à $\check{X}$ en $H_o$ . Comme $\varphi$ est net en $(x_o, H_o)$ , et $\check{X}$ est lisse en $H_o = \varphi(x_o, H_o)$ , $\varphi$ induit un isomorphisme des espaces tangents

}{Pour ceci, remarquons que pour une hypersurface quelconque dans un espace projectif, le morphisme de GAUSS (3.4.1) est ``défini'' en tout point lisse. Appliquons ceci à l'hypersurface $\check{X}$ dans $\check{\mathbb{P}}^r$ ; comme $V'$ est
\SGApage{II}{229}
contenu dans $V$ , l'ouvert de lissité de $\check{X}$ , le morphisme de GAUSS nous donne un morphisme, noté $\psi_1$ :

\begin{equation}\tag{3.5.3}
\psi_1 : V' \longrightarrow \check{\check{\mathbb{P}}}^r \simeq \mathbb{P}^r .
\end{equation}
Il reste à vérifier la commutativité du diagramme (3.5.2). Soit alors $(x_o, H_o)$ un point fermé de $\varphi^{-1}(V')$ ; il faut vérifier que $x_o$ , considéré comme étant un hyperplan dans $\check{\mathbb{P}}^r$ , est tangent à $\check{X}$ en $H_o$ , i.e. on va démontrer que pour $(x_o, H_o) \in \varphi^{-1}(V')$ , $H_o$ tangent à $X$ en $x_o \Rightarrow x_o$ tangent à $\check{X}$ en $H_o$ . Comme $\varphi$ est net en $(x_o, H_o)$ , et $\check{X}$ est lisse en $H_o = \varphi(x_o, H_o)$ , $\varphi$ induit un isomorphisme des espaces tangents

}{F7II-XVII-022}\begin{equation}\tag{3.5.4}
T_{\mathbb{P}(N)}((x_o,H_o)) \overset{\varphi}{\overset{\sim}{\longrightarrow}} T_{\check{X}} (H_o) ,
\end{equation}
de sorte que la première partie de la démonstration s'achève par le lemme suivant.

\underline{Lemme} 3.6 (EULER). \underline{Soit} $(x_o, H_o)$ \underline{un point fermé de} $\mathbb{P}(N)$ , \underline{et} \underline{désignons par} $T_L$ \underline{le fibré tangent de l'hyperplan} $L$ \underline{dans} $\check{\mathbb{P}}^r$ \underline{défini par} $x_o$ . \underline{Alors dans} $T_{\check{\mathbb{P}}^r}(H_o)$ \underline{on a}

\SGAvar{\begin{equation*}
\varphi (T_{\mathbb{P}(N)}((x_o, H_o))) \subset T_L(H_o) .
\end{equation*}
\underline{Démonstration.} Choisissons des coordonnés homogènes $X_o,\ldots,X_r$ dans $\mathbb{P}^r$ , en termes desquelles $x_o$ devient le point $\underline{A} = (A_o,\ldots,A_r)$ .


}{\begin{equation*}
\varphi (T_{\mathbb{P}(N)}((x_o, H_o))) \subset T_L(H_o) .
\end{equation*}
\underline{Démonstration.} Choisissons des coordonnées homogènes $X_o,\ldots,X_r$ dans $\mathbb{P}^r$ , en termes desquelles $x_o$ devient le point $\underline{A} = (A_o,\ldots,A_r)$ .


}{F7II-XVII-023}\SGApage{II}{230}

Un point de $T_X(x_o)$ est un point de $\mathbb{P}^r$ à valeurs dans $k[\epsilon]/(\epsilon^2)$,

\begin{equation}\tag{3.6.1}
\underline{A} + \epsilon \, \underline{B} \quad ,
\end{equation}
tel que pour toute section locale $H$ de $I$ , on ait

\begin{equation}\tag{3.6.2}
\sum_{i=0}^{r} \; B_i \, \frac{\partial H}{\partial X_i}(\underline{A}) \;=\; 0 \quad .
\end{equation}
Choisissons de plus une section locale $H$ de $I$ qui donne lieu à $H_o$ par (3.1.1). Par la trivialité locale de $\mathbb{P}(N)$ au-dessus de $X$ , un point de $T_{\mathbb{P}(N)}((x_o,H_o))$ se représente par un point

\begin{equation}\tag{3.6.3}
(\underline{A} + \epsilon \, \underline{B} \, , \; H + \epsilon \, F) \quad ,
\end{equation}
où $F$ est une section locale de $I$ , et où $\underline{A}$ et $\underline{B}$ sont comme ci-dessus. L'image par $\varphi$ de (3.6.3) est l'hyperplan à coefficients dans $k[\epsilon]/(\epsilon^2)$ d'équation (c.f. 3.1.1)

\begin{equation}\tag{3.6.4}
\sum_{i=0}^{r} \; \frac{\partial}{\partial X_i} \, (H+\epsilon F) \; (\underline{A}+\epsilon\underline{B})T_i \quad ,
\end{equation}
de sorte que le lemme équivaut à la formule

\begin{equation}\tag{3.6.5}
\sum_{i=0}^{r} \; A_i \, \frac{\partial}{\partial X_i} \, (H+\epsilon F) \; (\underline{A}+\epsilon\underline{B}) = 0 \quad .
\end{equation}
Compte tenu de ce que $\epsilon^2 = 0$ , et la formule de Taylor, on a

\begin{equation}\tag{3.6.6}
\frac{\partial}{\partial X_i}(H+\epsilon F)(\underline{A}+\epsilon\underline{B}) = \frac{\partial H}{\partial X_i}(\underline{A})+\epsilon \, \frac{\partial F}{\partial X_i}(\underline{A})+\epsilon \sum_{j=0}^{r} B_j \, \frac{\partial^2 H}{\partial X_j \partial X_i} \, (\underline{A}) \quad , \quad 
\end{equation}

\SGApage{II}{231}

 donc il suffit de vérifier les trois formules suivantes:

\begin{equation}\tag{3.6.7}
\sum_{i=0}^{r} \; A_i \, \frac{\partial H}{\partial X_i}(\underline{A}) \;=\; 0 \quad ,
\end{equation}
\begin{equation}\tag{3.6.8}
\sum_{i=0}^{r} \; A_i \, \frac{\partial F}{\partial X_i}(\underline{A}) \;=\; 0 \quad ,
\end{equation}
\begin{equation}\tag{3.6.9}
\sum_{i,j=0}^{r} \; A_i B_j \, \frac{\partial^2 H}{\partial X_i \partial X_j} \; (\underline{A}) \;=\; 0 \quad .
\end{equation}
Or, $H$ et $F$ étant homogènes de degré zéro, la formule d'Euler donne

\begin{equation}\tag{3.6.10}
\sum X_i \, \frac{\partial H}{\partial X_i} = 0 \quad , \quad \sum \; X_i \, \frac{\partial F}{\partial X_i} \;=\; 0
\end{equation}
et, $\frac{\partial H}{\partial X_j}$ étant homogène de degré $-1$ , cette formule donne

\begin{equation*}
\sum_i \; X_i \, \frac{\partial^2 H}{\partial X_i \partial X_j} \;=\; - \, \frac{\partial H}{\partial X_j} \quad ,
\end{equation*}
de sorte que (3.6.9) se récrit

\begin{equation}\tag{3.6.11}
- \sum_j B_j \, \frac{\partial H}{\partial X_j}(\underline{A}) = 0 \quad ,
\end{equation}
qui n'est autre que (3.6.2).

Ceci démontre que $\varphi$ induit un isomorphisme

\begin{equation}\tag{3.6.12}
\varphi^{-1}(V') \; \stackrel{\sim}{\longrightarrow} \; V' \quad 
\end{equation}

\SGApage{II}{232}

 dont l'inverse est $\psi$ . Remarquons maintenant que le morphisme

\begin{equation}\tag{3.6.13}
\left\{ \begin{array}{l} \psi \ : \ V' \longrightarrow \ \mathbb{P}^r \times V' \subset \mathbb{P}^r \times \check{X} \\[2ex] \psi(v) \ = \ (\psi_1(v),v) \end{array} \right.
\end{equation}
s'étend en un morphisme

\begin{equation}\tag{3.6.14}
\left\{ \begin{array}{l} \widetilde{\psi} \ : \ V \longrightarrow \ \mathbb{P}^r \times V \\[2ex] \widetilde{\psi}(v) \ = \ (\widetilde{\psi}_1(v),v) \ \ , \end{array} \right.
\end{equation}
où l'on désigne par

\begin{equation}\tag{3.6.15}
\widetilde{\psi}_1 \ : \ V \longrightarrow \ \mathbb{P}^r
\end{equation}
le morphisme de GAUSS, vu comme défini sur $V$ tout entier.

On vient de démontrer l'inclusion

\begin{equation}\tag{3.6.16}
\psi(V') \ \subset \ \mathbb{P}(N)
\end{equation}
ce qui implique, $V'$ étant dense dans $V$ , et $\mathbb{P}(N)$ étant fermé dans $\mathbb{P}^r \times \check{X}$ , l'inclusion

\begin{equation}\tag{3.6.17}
\widetilde{\psi}(V) \subset \mathbb{P}(N)
\end{equation}
Il est évident qu'on a même

\begin{equation}\tag{3.6.18}
\widetilde{\psi}(V) \ \subset \ \varphi^{-1}(V) \ .
\end{equation}
Encore par "continuité" , les relations 


\SGApage{II}{233}

\begin{equation*}
 \left\{ \begin{array}{l} \psi \ \circ \ \varphi \ = \ \mathrm{id}_{\varphi^{-1}(V')} \\[3ex] \varphi \ \circ \ \psi \ = \ \mathrm{id}_{V'} \end{array} \right.
\end{equation*}
impliquent les relations

\begin{equation*}
\left\{ \begin{array}{l} \widetilde{\psi} \ \circ \ \varphi = \mathrm{id}_{\varphi^{-1}(V)} \\[3ex] \varphi \ \circ \ \widetilde{\psi} = \mathrm{id}_{V} \qquad , \end{array} \right.
\end{equation*}
ce qui démontre que les ouverts 3.5 a) et 3.5 b) sont identiques.

Démontrons que les ouverts 3.5 a) et 3.5 c) sont identiques. On remarque d'abord que, d'après 3.3, l'ouvert 3.5 c) est le sous-ensemble $\underline{U}$ de $\check{X}$ défini par

\begin{equation}\tag{3.6.19}
\underline{U} = \left\{\, x \in \check{X} \ \middle|\ \begin{array}{l} \varphi^{-1}(x) \ \mbox{contient exactement un point,} \\ \mbox{et en ce point } \varphi \mbox{ est net.} \end{array} \right\} \ .
\end{equation}
Or $\varphi^{-1}(\underline{U})$ est normal (étant lisse) et le morphisme induit

\SGAvar{\begin{equation*}
\varphi \ : \ \varphi^{-1}(\underline{U}) \longrightarrow \underline{U}
\end{equation*}
est birationnel (on l'a déja démontré), radiciel (par définition de $\underline{U}$) et propre. Donc

}{\begin{equation*}
\varphi \ : \ \varphi^{-1}(\underline{U}) \longrightarrow \underline{U}
\end{equation*}
est birationnel (on l'a déjà démontré), radiciel (par définition de $\underline{U}$) et propre. Donc

}{F7II-XVII-024}\begin{equation*}
\varphi \ : \ \varphi^{-1}(\underline{U}) \longrightarrow \underline{U}
\end{equation*}
est le \underline{normalisé} de $\underline{U}$ . En particulier, $\underline{U}$ est \underline{géométriquement unibranche} . Comme $\varphi : \varphi^{-1}(\underline{U}) \longrightarrow \underline{U}$ est \underline{net} , il s'en déduit (EGA IV 18.10.1) qu'il est \underline{étale} . Par EGA IV 18.10.18, on conclut alors que c'est une immersion ouverte, donc un isomorphisme, donc que
\SGApage{II}{234}
$\underline{U}$ est lisse, i.e. qu'on a l'inclusion 3.5 c) $\subset$ 3.5 a). D'autre part, (3.6.19) rend évidente l'inclusion 3.5 b) $\subset$ 3.5 c), et on conclut, compte tenu de ce que 3.5 a)$=$ 3.5 b).

Nous pouvons maintenant démontrer:

3.7.  Le théorème 2.5 est vrai dans le cas $\mathrm{car}.k \neq 2$ , ainsi que dans le cas $X$ de dimension paire.

\underline{Démonstration}.  Compte tenu de 3.5.0 et de 3.3, il suffit de vérifier :

(3.7.1) Pour tout point fermé $x_o \in X$, et tout entier $d \geqslant 2$, il existe une hypersurface $H$ de degré $d$ tangente à $X$ en $x_o$ , telle que $x_o$ soit une singularité quadratique ordinaire de $X.H$.

\SGAvar{Or, prenant convenablement des coordonnés homogènes $X_o,\ldots,X_r$ , on peut supposer que $x_o$ est le point $(1,0,\ldots,0)$ , et que les fonctions $X_i/X_o$ , $i = 1,\ldots,n$ , donnent des coordonnés locales sur $X$ dans un voisinage de $x_o$ . On prend pour $H$ l'hypersurface d'équation

}{Or, prenant convenablement des coordonnées homogènes $X_o,\ldots,X_r$ , on peut supposer que $x_o$ est le point $(1,0,\ldots,0)$ , et que les fonctions $X_i/X_o$ , $i = 1,\ldots,n$ , donnent des coordonnées locales sur $X$ dans un voisinage de $x_o$ . On prend pour $H$ l'hypersurface d'équation

}{F7II-XVII-025}\begin{equation}\tag{3.7.1.1}
H = X_o^{d-2} \sum_{i=1}^{n} X_i^2 \qquad \text{si} \quad \mathrm{car}.(k) \neq 2
\end{equation}
\begin{equation}\tag{3.7.1.2}
H = X_o^{d-2} \sum_{i=1}^{n/2} X_i \, X_{n/2+i} \qquad \text{si} \quad n \ \text{est pair}
\end{equation}
\begin{equation}\tag{3.7.1.3}
H = X_o^{d-2} \, (X_n^2 + \sum_{i=1}^{[n/2]} X_i \, X_{[n/2]+i}) \ \text{si} \quad n \ \text{est impair} \ .
\end{equation}

\SGAvar{\SGApage{II}{235}

Pour l'énoncé 2.5.2, remarquons qu'ou bien $\varphi$ est génériquement net, et 3.5.0 s'applique, ou bien $\varphi$ est partout ramifié, ce qui implique $\dim X \leqslant r-2$ , $k$ étant de caractéristique nulle. Dans ce dernier cas, le critère (3.2.4) est trivialement vérifié.


}{\SGApage{II}{235}



Pour l'énoncé 2.5.2, remarquons qu'ou bien $\varphi$ est génériquement net, et 3.5.0 s'applique, ou bien $\varphi$ est partout ramifié, ce qui implique $\dim \check{X} \leqslant r-2$ , $k$ étant de caractéristique nulle. Dans ce dernier cas, le critère (3.2.4) est trivialement vérifié.


}{F7II-XVII-026}\SGAsectiontoc{4.  Le cas général}{4}
4.  \underline{Le cas général}

Nous allons esquisser une autre méthode de démonstration de 2.5 valable en toute caractéristique, en admettant l'énoncé suivant, voisin de 3.1, et cas particulier de XV 1.1.4.

(4.1) \underline{Le sous-ensemble de} $\mathbb{P}(N)$ \underline{formé des points} $(x, H)$ \underline{tels que} $x$ \underline{soit une singularité quadratique ordinaire de} $X \cdot H$ \underline{est ouvert.}

Remarquons que 4.1 résulte de 3.3, sauf si $\mathrm{car}(k) = 2$ \underline{et} $\dim X$ impaire.

4.1.1. Désignons par $R_1$ le fermé de $\mathbb{P}(N)$ complément de l'ouvert 4.1, et par $F_1$ son image $\varphi(R_1) \subset \check{X} \subset \check{\mathbb{P}}^r$ .

\underline{Lemme} 4.1.2. \underline{Si l'ouvert} 4.1 \underline{est non-vide,} \underline{alors} $\operatorname{codim}_{\check{\mathbb{P}}^r}(F_1) \geqslant 2$

\underline{Démonstration}. Par hypothèse, $\dim(R_1) < \dim(\mathbb{P}(N)) = r-1$ , et $F_1 = \varphi(R_1)$ .

4.1.3. Désignons par $F_2$ la partie de $\check{X}$ formée des hyperplans qui sont tangents à $X$ en au moins deux points. On remarque que $F_2$ n'est
\SGApage{II}{236}
pas fermé, mais qu'on a

\SGAvar{4.1.4.\quad l'adhérence de $F_2 \subset F_1 \cup F_2$ ,

ceci à cause du fait que $F_1 \cup F_2 = F$ (cf. 3.2) est fermé. (L'inclusion 4.1.4. ``dit'' qu'une \underline{confluence} de deux singularités n'est pas quadratique ordinaire.) En tout cas la construction de 4.2.4. montrera que $F_2$ est constructible. Pour démontrer 2.5.1 nous allons appliquer le critère 3.2.4. Pour ceci, il faudra savoir que $\mathrm{codim}_{\check{\mathbb{P}}^r}(F_1) \geq 2$ et que $\mathrm{codim}_{\check{\mathbb{P}}^r}(F_2) \geq 2$. D'autre part, 3.8.1 et 4.1.2 nous assurent que, remplaçant le plongement donné par son $d$-ième multiple, $d \geq 2$ , on aura $\mathrm{codim}_{\check{\mathbb{P}}^r}(F_2) \geq 2$ . On conclut que 2.5.1 sera prouvé, une fois connue la

\underline{Proposition} 4.2. \underline{Soit} $X$ \underline{une sous-variété lisse et irréductible de} $\mathbb{P}^r$ , \underline{de dimension} $n \geq 1$ . \underline{Pour tout entier} $d \geq 2$ , \underline{la partie} $F_2$ \underline{de} $\check{\mathbb{P}}^{{n+d \choose d}-1} = \check{\mathbb{P}}^N$ \underline{formée des hypersurfaces de degré} $d$ \underline{dans} $\mathbb{P}^r$ \underline{qui sont tangentes à} $X$ \underline{en au moins deux points est de codimension} $\geq 2$ \underline{dans} $\check{\mathbb{P}}^N$ .

4.2.1. Commencons par le cas spécial $d = 2$ et $X$ une sous-variété linéaire de $\mathbb{P}^r$ . Or, il suffira de trouver une droite dans $\check{\mathbb{P}}^N$ qui ne rencontre pas $F$ (cf. 3.2.4). Choisissons des coordonnés homogènes $X_o,\ldots,X_r$ dans $\mathbb{P}^r$ de telle facon que $X$ soit défini par
\SGApage{II}{237}
$X_{n+1} = \ldots = X_r = 0$ , et des coordonnés homogènes $(\lambda, \mu)$ dans $\mathbb{P}^1$ .

}{\begin{equation}\tag{4.1.4}
\overline{F_2}\subset F_1\cup F_2 \quad ,
\end{equation}




ceci à cause du fait que $F_1 \cup F_2 = F$ (cf. 3.2) est fermé. (L'inclusion 4.1.4. ``dit'' qu'une \underline{confluence} de deux singularités n'est pas quadratique ordinaire.) En tout cas la construction de 4.2.4. montrera que $F_2$ est constructible. Pour démontrer 2.5.1 nous allons appliquer le critère 3.2.4. Pour ceci, il faudra savoir que $\mathrm{codim}_{\check{\mathbb{P}}^r}(F_1) \geq 2$ et que $\mathrm{codim}_{\check{\mathbb{P}}^r}(F_2) \geq 2$. D'autre part, 3.7.1 et 4.1.2 nous assurent que, remplaçant le plongement donné par son $d$-ième multiple, $d \geq 2$ , on aura $\mathrm{codim}_{\check{\mathbb{P}}^r}(F_1) \geq 2$ . On conclut que 2.5.1 sera prouvé, une fois connue la



\underline{Proposition} 4.2. \underline{Soit} $X$ \underline{une sous-variété lisse et irréductible de} $\mathbb{P}^r$ , \underline{de dimension} $n \geq 1$ . \underline{Pour tout entier} $d \geq 2$ , \underline{la partie} $F_2$ \underline{de} $\check{\mathbb{P}}^{{r+d \choose d}-1} = \check{\mathbb{P}}^N$ \underline{formée des hypersurfaces de degré} $d$ \underline{dans} $\mathbb{P}^r$ \underline{qui sont tangentes à} $X$ \underline{en au moins deux points est de codimension} $\geq 2$ \underline{dans} $\check{\mathbb{P}}^N$ .

4.2.1. Commençons par le cas spécial $d = 2$ et $X$ une sous-variété linéaire de $\mathbb{P}^r$ . Or, il suffira de trouver une droite dans $\check{\mathbb{P}}^N$ qui ne rencontre pas $F$ (cf. 3.2.4). Choisissons des coordonnées homogènes $X_o,\ldots,X_r$ dans $\mathbb{P}^r$ de telle façon que $X$ soit défini par
\SGApage{II}{237}
$X_{n+1} = \ldots = X_r = 0$ , et des coordonnées homogènes $(\lambda, \mu)$ dans $\mathbb{P}^1$ .

}{F7II-XVII-027}Pour $n$ pair, $n = 2k$ , on prend le pinceau de quadriques

\begin{equation}\tag{4.2.2}
H_{(\lambda,\mu)} = \lambda\ X_o^2 + (\lambda X_k + \mu X_o) X_{2k} + \sum_{i=1}^{k-1}\ (\lambda X_i + \mu X_{i+1}) X_{k+i} \quad ;
\end{equation}
pour $n$ impair, $n = 2k+1$ , on prend le pinceau

\begin{equation}\tag{4.2.3}
H_{(\lambda,\mu)} = (\lambda X_1 + \mu X_o) X_o + \sum_{i=1}^{k}\ (\lambda X_{2i} + \mu X_i) X_{2i+1} \quad .
\end{equation}
Dans les deux cas, on constate que $X \cdot H_{(\lambda,\mu)}$ est lisse pour $(\lambda,\mu) \neq (0,1)$ , et que $X \cdot H_{(0,1)}$ n'a qu'un point singulier, qui est même quadratique ordinaire, et donc que le pinceau ne coupe pas $F$ .

4.2.4. Désignons par $\mathrm{Btg}^d(X)$ la sous-variété de $\check{\mathbb{P}}^N \times (X \times X - \Delta(X \times X))$ formée des triples $(H,\ x_1,\ x_2)$ tels que $H$ est une hypersurface de degré $d$ dans $\mathbb{P}^r$ qui est tangente à $X$ en $x_1$ et en $x_2$ , et $x_1 \neq x_2$ . L'image de $\mathrm{Btg}^d(X)$ par $\mathrm{pr}_1$ est l'ensemble $F_2$ de (4.2), de sorte qu'il suffira de démontrer

\SGAvar{(4.2.5)\quad $\dim \mathrm{Btg}^d(X) \leq N-2$ si $d \geq 3$ , ou si $d = 2$ et $X$ n'est pas linéaire.

}{\begin{equation}\tag{4.2.5}
\dim \mathrm{Btg}^d(X) \leq N-2
\quad\mbox{si } d \geq 3,\ \mbox{ou si } d = 2
\mbox{ et } X \mbox{ n'est pas linéaire}.
\end{equation}

}{F7II-XVII-028}Pour ceci, désignons par $\mathrm{Lin}(X \times X)$ la sous-variété de $X \times X - \Delta(X \times X)$

\SGAvar{(4.2.6)\quad $\mathrm{Lin}(X \times X) = \{(x_1,x_2) \in X \times X \mid x_1 \neq x_2$ , et la droite engendrée par $x_1,x_2$ est tangente à $X$ en $x_1$ $\}$ \\ $\{(x_1,x_2) \in X \times X \mid x_1 \neq x_2$ et $x_2 \in X \cap \mathrm{Tg}\ (x_1)\}$ .


}{\begin{equation}\tag{4.2.6}
\begin{aligned}
\mathrm{Lin}(X \times X)
&= \{(x_1,x_2) \in X \times X \mid x_1 \neq x_2,\ 
\text{et la droite engendrée par }x_1,x_2
\text{ est tangente à }X\text{ en }x_1\}\\
&= \{(x_1,x_2) \in X \times X \mid x_1 \neq x_2,\ 
x_2 \in X \cap \mathrm{Tg}(x_1)\}.
\end{aligned}
\end{equation}


}{F7II-XVII-029}\SGApage{II}{238}

On a évidemment :

4.2.7.\quad \underline{Si} $X$ \underline{n'est pas une sous-variété linéaire de} $\mathbb{P}^r$ , \underline{alors}

\begin{equation*}
\dim \mathrm{Lin}(X \times X) < 2\dim X \quad .
\end{equation*}
Donc pour (4.2.5) il suffira de démontrer que les fibres du morphisme

\SGAvar{\begin{equation}\tag{4.2.8}
\mathrm{Btg}^d(X) \stackrel{\mathrm{pr}_{2,3}}{\longrightarrow} X \times X - \Delta(X \times X) \quad ,
\end{equation}
considérés comme des sous-variétés de $\check{\mathbb{P}}^N$ , sont

\begin{equation}\tag{4.2.9}
\left\{
\begin{array}{l}
\mbox{de codimension} \geq 2\dim(X)+1 \quad \mbox{au-dessus de} \quad \mathrm{Lin}(X{\times}X) \quad , \\[2ex]
\mbox{de codimension} \geq 2\dim(X)+2 \quad \mbox{au-dessus de} \quad X{\times}X-\Delta(X{\times}X)-\mathrm{Lin}(X{\times}X)
\end{array}
\right\} \quad .
\end{equation}
L'énoncé (4.2.9) est une conséquence du lemme suivant, en prenant pour $L_i$ l'espace tangent à $X$ en $x_i$ , $i=1,2$ .

\underline{Proposition} 4.3.\quad \underline{Soient} $x_1$ \underline{et} $x_2$ \underline{deux points distincts fermés de} $\mathbb{P}^r$ , \underline{et} $L_1$ \underline{et} $L_2$ \underline{deux sous-espaces linéaires de} $\mathbb{P}^r$ \underline{tels que} $x_1 \in L_1$ \underline{et} $x_2 \in L_2$ . \underline{Dans l'espace} $\check{\mathbb{P}}^N = \check{\mathbb{P}}^{\binom{r+d}{d}-d}$ \underline{des hypersurfaces de degré} $d$ , \underline{le sous-espace linéaire formé des hypersurfaces tangentes à} $L_1$ \underline{en} $x_1$ \underline{et à} $L_2$ \underline{en} $x_2$ \underline{est de codimension}

}{\begin{equation}\tag{4.2.8}
\mathrm{Btg}^d(X) \stackrel{\mathrm{pr}_{2,3}}{\longrightarrow} X \times X - \Delta(X \times X) \quad ,
\end{equation}
considérées comme des sous-variétés de $\check{\mathbb{P}}^N$ , sont

\begin{equation}\tag{4.2.9}
\left\{
\begin{array}{l}
\mbox{de codimension} \geq 2\dim(X)+1 \quad \mbox{au-dessus de} \quad \mathrm{Lin}(X{\times}X) \quad , \\[2ex]
\mbox{de codimension} \geq 2\dim(X)+2 \quad \mbox{au-dessus de} \quad X{\times}X-\Delta(X{\times}X)-\mathrm{Lin}(X{\times}X)
\end{array}
\right\} \quad .
\end{equation}
L'énoncé (4.2.9) est une conséquence du lemme suivant, en prenant pour $L_i$ l'espace tangent à $X$ en $x_i$ , $i=1,2$ .



\underline{Proposition} 4.3.\quad \underline{Soient} $x_1$ \underline{et} $x_2$ \underline{deux points distincts fermés de} $\mathbb{P}^r$ , \underline{et} $L_1$ \underline{et} $L_2$ \underline{deux sous-espaces linéaires de} $\mathbb{P}^r$ \underline{tels que} $x_1 \in L_1$ \underline{et} $x_2 \in L_2$ . \underline{Dans l'espace} $\check{\mathbb{P}}^N = \check{\mathbb{P}}^{\binom{r+d}{d}-1}$ \underline{des hypersurfaces de degré} $d$ , \underline{le sous-espace linéaire formé des hypersurfaces tangentes à} $L_1$ \underline{en} $x_1$ \underline{et à} $L_2$ \underline{en} $x_2$ \underline{est de codimension}

}{F7II-XVII-030}\begin{equation*}
\left\{
\begin{array}{l}
= \dim L_1+\dim L_2+2 \quad \mbox{si } d\geq 3, \mbox{ ou si } d=2 \mbox{ et } L_1 \cap L_2 \mbox{ ne contient pas} \\[1ex]
\qquad\qquad\qquad\qquad \mbox{la droite } (x_1,x_2) \mbox{ contenant } x_1,x_2 \\[2ex]
= \dim L_1+\dim L_2+1 \quad \mbox{si } d=2 \mbox{ et } L_1 \cap L_2 \mbox{ contient la droite } (x_1,x_2).
\end{array}
\right.
\end{equation*}

\SGApage{II}{239}

\SGAvar{\underline{Démonstration.}\quad Pour $d \geq 3$ , il faut prouver que $\dim L_1 + \dim L_2+2$ conditions linéaires sont indépendantes , de sorte qu'il suffit de traiter le cas $L_1 = L_2 = \mathbb{P}^r$ , i.e. de trouver ``combien'' d'hypersurfaces sont singulières en $x_1$ et en $x_2$ . Or, en termes de coordonnés homogènes choisies de telle façon que $x_1 = (1,0,...0)$ et $x_2 = (0,1,0...,0)$ , l'hypersurface d'équation $\Sigma A_W X^W$ est singulière en $(1,0,..,0)$ si et seulement si les coéfficients des monômes $X_o^{d-1}X_i$ , $i=0,\ldots,r$, sont tous nuls, et elle est singulière en $(0,1,0,..0)$ si et seulement si les coefficients des monômes $X_1^{d-1}X_i$ , $i=0,\ldots,r$ , sont tous nuls. Pour $d \geq 3$ , les monômes $X_o^{d-1}X_i$, $i=0,\ldots,r$ et $X_1^{d-1}X_i$, $i=0,\ldots,r$ sont tous distincts, ce qui donne l'énoncé voulu pour $d \geq 3$ .

}{\underline{Démonstration.}\quad Pour $d \geq 3$ , il faut prouver que $\dim L_1 + \dim L_2+2$ conditions linéaires sont indépendantes , de sorte qu'il suffit de traiter le cas $L_1 = L_2 = \mathbb{P}^r$ , i.e. de trouver ``combien'' d'hypersurfaces sont singulières en $x_1$ et en $x_2$ . Or, en termes de coordonnées homogènes choisies de telle façon que $x_1 = (1,0,...0)$ et $x_2 = (0,1,0...,0)$ , l'hypersurface d'équation $\Sigma A_W X^W$ est singulière en $(1,0,..,0)$ si et seulement si les coefficients des monômes $X_o^{d-1}X_i$ , $i=0,\ldots,r$, sont tous nuls, et elle est singulière en $(0,1,0,..0)$ si et seulement si les coefficients des monômes $X_1^{d-1}X_i$ , $i=0,\ldots,r$ , sont tous nuls. Pour $d \geq 3$ , les monômes $X_o^{d-1}X_i$, $i=0,\ldots,r$ et $X_1^{d-1}X_i$, $i=0,\ldots,r$ sont tous distincts, ce qui donne l'énoncé voulu pour $d \geq 3$ .

}{F7II-XVII-031}Pour $d = 2$ , le même calcul donne l'inégalité

\SGAvar{\begin{equation*}
\mathrm{codim} \geq \dim L_1 + \dim L_2 + 1 \quad ,
\end{equation*}
(car le monôme $X_o \, X_1$ est compté deux fois).

Vérifions que si $L_1 \cap L_2$ ne contient pas la droite $(x_1,x_2)$, alors les conditions de tangentialité sont indépendentes. On se ramène au cas $L_1 = \mathbb{P}^r$ , $L_2$ un hyperplan dans $\mathbb{P}^r$ qui ne contient pas la droite $(x_1,x_2)$ . On choisit des coordonnés homogènes $X_o,\ldots,X_r$ de telle façon que $x_1 = (1,0,\ldots,0)$ , $x_1 = (0,1,0...0)$ et $L_2$ soit l'hyperplan $X_o = 0$ . Or, une quadrique $\sum\limits_{i \leq j} A_{ij} \, X_i \, X_j$ est singulière en $(1,0,..,0)$ si et seulement si $A_{oj} = 0$ pour $j = 0,\ldots,r$ , et tangente à
\SGApage{II}{240}
$X_o = 0$ en $(0,1,\ldots,0)$ si et seulement si $A_{ij} = 0$ pour $j = 1,\ldots,r$ , ce qui fait ensemble $2r+1 = \dim L_1 + \dim L_2 + 2$ conditions indépendantes.

Pour le cas où $L_1 \cap L_2$ contient la droite $(x_1,x_2)$, on se ramène au cas $L_1 = L_2 =$ la droite $(x_1,x_2)$ . En termes de coordonnés homogènes dans lesquelles $x_1 = (1,0,..,0)$ , $x_2 = (0,1,0,...0)$ , la quadrique $\sum_{i \leq j} A_{ij} X_i X_j$ passe par $(1,0,..0)$ et y est tangente à $L_1$ si et seulement si $A_{oo} = A_{o1} = 0$ ; elle passe par $(0,1,0...0)$ et y est tangente à $L_2$ si et seulement si $A_{11} = A_{o,1} = 0$ , ce qui ne fait ensemble que $3 = \dim L_1 + \dim L_2 + 1$ conditions.


}{\begin{equation*}
\mathrm{codim} \geq \dim L_1 + \dim L_2 + 1 \quad ,
\end{equation*}
(car le monôme $X_o \, X_1$ est compté deux fois).


Vérifions que si $L_1 \cap L_2$ ne contient pas la droite $(x_1,x_2)$, alors les conditions de tangentialité sont indépendantes. On se ramène au cas $L_1 = \mathbb{P}^r$ , $L_2$ un hyperplan dans $\mathbb{P}^r$ qui ne contient pas la droite $(x_1,x_2)$ . On choisit des coordonnées homogènes $X_o,\ldots,X_r$ de telle façon que $x_1 = (1,0,\ldots,0)$ , $x_2 = (0,1,0...0)$ et $L_2$ soit l'hyperplan $X_o = 0$ . Or, une quadrique $\sum\limits_{i \leq j} A_{ij} \, X_i \, X_j$ est singulière en $(1,0,..,0)$ si et seulement si $A_{oj} = 0$ pour $j = 0,\ldots,r$ , et tangente à
\SGApage{II}{240}

$X_o = 0$ en $(0,1,\ldots,0)$ si et seulement si $A_{1j} = 0$ pour $j = 1,\ldots,r$ , ce qui fait ensemble $2r+1 = \dim L_1 + \dim L_2 + 2$ conditions indépendantes.

Pour le cas où $L_1 \cap L_2$ contient la droite $(x_1,x_2)$, on se ramène au cas $L_1 = L_2 =$ la droite $(x_1,x_2)$ . En termes de coordonnées homogènes dans lesquelles $x_1 = (1,0,..,0)$ , $x_2 = (0,1,0,...0)$ , la quadrique $\sum_{i \leq j} A_{ij} X_i X_j$ passe par $(1,0,..0)$ et y est tangente à $L_1$ si et seulement si $A_{oo} = A_{o1} = 0$ ; elle passe par $(0,1,0...0)$ et y est tangente à $L_2$ si et seulement si $A_{11} = A_{o,1} = 0$ , ce qui ne fait ensemble que $3 = \dim L_1 + \dim L_2 + 1$ conditions.


}{F7II-XVII-032}\SGAsectiontoc{5.\quad Le degré de la variété duale par voie "élémentaire"}{5}
5.\quad \underline{Le degré de la variété duale par voie "élémentaire"}

(Cf.XVIII 3.2. pour une autre méthode dans le cas où $\dim X$ est pair \underline{ou} $\mathrm{car}(k) \neq 2$ .)

Nous employons toujours les notations de 2.2, donc $X \longrightarrow \mathbb{P}^r$ est une sous-variété lisse irréductible de dimension $n$ .

5.1. Commençons par considérer le morphisme (3.1.1)

\begin{equation}\tag{5.1.0}
\nu : \mathbb{P}(N) \longrightarrow X \times \check{\mathbb{P}}^r
\end{equation}
d'un point de vue plus "fonctoriel" . (Je tiens à remercier S. KLEIMAN pour m'avoir signalé la formulation qui suit.) On a une suite exacte
\SGApage{II}{241}
de faisceaux cohérents sur $X$

\SGAvar{\begin{equation}\tag{5.1.1}
0 \longrightarrow \check{N} \longrightarrow \Omega^1_{\mathbb{P}/k} \,|\, X \longrightarrow \Omega^1_X \longrightarrow 0 \quad .
\end{equation}
D'autre part, "écrivant les formes différentielles sur $\mathbb{P}^r$ en coordonnés homogènes" , on trouve une suite exacte de faisceaux sur $\mathbb{P}^r$

}{\begin{equation}\tag{5.1.1}
0 \longrightarrow \check{N} \longrightarrow \Omega^1_{\mathbb{P}/k} \,|\, X \longrightarrow \Omega^1_X \longrightarrow 0 \quad .
\end{equation}
D'autre part, "écrivant les formes différentielles sur $\mathbb{P}^r$ en coordonnées homogènes" , on trouve une suite exacte de faisceaux sur $\mathbb{P}^r$

}{F7II-XVII-033}\begin{equation}\tag{5.1.2}
0 \longrightarrow \Omega^1_{\mathbb{P}^r/k} \longrightarrow (\underline{O}_{\mathbb{P}}(-1))^{r+1} \longrightarrow \underline{O}_{\mathbb{P}} \longrightarrow 0
\end{equation}
qui, restreinte à $X$ , donne alors une suite exacte

\SGAvar{\begin{equation}\tag{5.1.3}
0 \longrightarrow \Omega^1_{\mathbb{P}^r/k} \,|\, X \longrightarrow (\underline{O}_X(-1))^{r+1} \longrightarrow \underline{O}_{\mathbb{P}} \longrightarrow 0 \quad .
\end{equation}
Mettant ensemble 5.1.1 et 5.1.3, on trouve une inclusion

}{\begin{equation}\tag{5.1.3}
0 \longrightarrow \Omega^1_{\mathbb{P}^r/k} \,|\, X \longrightarrow (\underline{O}_X(-1))^{r+1} \longrightarrow \underline{O}_X \longrightarrow 0 \quad .
\end{equation}
Mettant ensemble 5.1.1 et 5.1.3, on trouve une inclusion

}{F7II-XVII-034}\begin{equation}\tag{5.1.4}
\check{N} \lhook\joinrel\longrightarrow (\underline{O}_X(-1))^{r+1} \quad ,
\end{equation}
qui donne, par transposition, une surjection,

\begin{equation}\tag{5.1.5}
\underline{O}_X(1)^{r+1} \longrightarrow N
\end{equation}
d'où une immersion fermée

\begin{equation}\tag{5.1.6}
\alpha : \mathbb{P}(N) \longrightarrow \mathbb{P}(\underline{O}_X(1)^{r+1}) \;\simeq\; X \times \check{\mathbb{P}}^r
\end{equation}
qui n'est autre que (5.1.0).


\SGApage{II}{242}


5.2.  \underline{L'idée du calcul}.

Nous supposons à partir de maintenant que $\check{X} = \varphi(\mathbb{P}(N))$ est de dimension $r-1$ , donc que c'est une hypersurface dans $\check{\mathbb{P}}^r$ . Le calcul de son degré se fait naturellement en introduisant les classes de Chern [1] .

Pour toute variété projective et lisse $T$ , nous notons par $\mathrm{CH}(T)$ son anneau de Chow , qui est un anneau gradué (par la codimension d'un cycle). Etant donné un élément $z \in \mathrm{CH}(T)$ , nous notons par

\SGAvar{\begin{equation*}
\int_T z
\end{equation*}
le degré de sa composante de degré zéro (qui est un $0$-cycle). La classe dans $\mathrm{CH}^1(T)$ d'un faisceau inversible $\underline{L}$ sur $T$ sera noté également $\underline{L}$ quand aucun confusion n'est à craindre (par exemple, quand elle se trouve sous le signe d'intégration $\int_T$ ) .

}{\begin{equation*}
\int_T z
\end{equation*}
le degré de sa composante de degré zéro (qui est un $0$-cycle). La classe dans $\mathrm{CH}^1(T)$ d'un faisceau inversible $\underline{L}$ sur $T$ sera notée également $\underline{L}$ quand aucune confusion n'est à craindre (par exemple, quand elle se trouve sous le signe d'intégration $\int_T$ ) .

}{F7II-XVII-035}Ceci posé, notant par $L$ la classe dans $\mathrm{CH}(\check{\mathbb{P}}^r)$ d'une droite, on a

\begin{equation}\tag{5.2.1}
\deg(\check{X}) = \int_{\check{\mathbb{P}}^r} L\cdot\check{X} \quad .
\end{equation}
Or, $\check{X}$ est ``ensemblistement'' l'image de $\mathbb{P}(N)$ par $\varphi$ (3.1.2). Mais, du point de vue de l'homomorphisme ``image directe''

\begin{equation*}
\varphi_* : \mathrm{CH}(\mathbb{P}(N)) \longrightarrow \mathrm{CH}(\check{\mathbb{P}}^r) \quad 
\end{equation*}

\SGApage{II}{243}

 on a

\begin{equation*}
\varphi_*( \mathbb{P}(N)) = \deg(\varphi)\cdot \check{X} \quad ,
\end{equation*}
de sorte que 5.2.1 se récrit

\SGAvar{\begin{equation}\tag{5.2.2}
\deg(\check{X}) = \frac{1}{\deg(\varphi)} \int_{\check{\mathbb{P}}^r} L\cdot\varphi_*( \mathbb{P}(N)) \quad ;
\end{equation}
par abus de notations, dans $\deg(\varphi)$ , $\varphi$ désigne évidemment le morphisme $\mathbb{P}(N) \longrightarrow \check{X}$ induit par (3.1.2), qui est un entier $\geqslant 0$ , nul si et seulement si $\dim X < r-1$ ($=\dim \mathbb{P}(N)$). Par la formule de projection pour $\varphi$ , ceci se récrit

}{\begin{equation}\tag{5.2.2}
\deg(\check{X}) = \frac{1}{\deg(\varphi)} \int_{\check{\mathbb{P}}^r} L\cdot\varphi_*( \mathbb{P}(N)) \quad ;
\end{equation}


par abus de notations, dans $\deg(\varphi)$ , $\varphi$ désigne évidemment le morphisme $\mathbb{P}(N) \longrightarrow \check{X}$ induit par (3.1.2), qui est un entier $\geqslant 0$ , nul si et seulement si $\dim \check{X} < r-1$ ($=\dim \mathbb{P}(N)$). Par la formule de projection pour $\varphi$ , ceci se récrit

}{F7II-XVII-036}\begin{equation}\tag{5.2.3}
\deg(\check{X}) = \frac{1}{\deg(\varphi)} \int_{\mathbb{P}(N)} \varphi^*(L) \quad .
\end{equation}
Désignons par $\eta \in \mathrm{CH}( \check{\mathbb{P}}^r)$ la classe d'un hyperplan, de sorte que $L = \eta^{r-1}$ , et (5.2.3) se récrit

\begin{equation}\tag{5.2.4}
\deg(\check{X}) = \frac{1}{\deg(\varphi)} \int_{\mathbb{P}(N)} (\varphi^*(\eta))^{r-1} \quad .
\end{equation}
Calculons maintenant la classe $\varphi^*(\eta)$ en termes ``concrets'' .

\SGAvar{Notons par $\xi$ (resp. $\tau$) le faisceau inversible canonique $\underline{O}_{\mathbb{P}(N)}(1)$ sur $\mathbb{P}(N)$ (resp. $\underline{O}_{\mathbb{P}(\underline{O}_X(1)^{r+1})}$ sur $\mathbb{P}(\underline{O}_X(1)^{r+1})$ ) . Par fonctorialité, nous avons (cf. (5.1.6))

}{Notons par $\xi$ (resp. $\tau$) le faisceau inversible canonique $\underline{O}_{\mathbb{P}(N)}(1)$ sur $\mathbb{P}(N)$ (resp. $\underline{O}_{\mathbb{P}(\underline{O}_X(1)^{r+1})}(1)$ sur $\mathbb{P}(\underline{O}_X(1)^{r+1})$ ) . Par fonctorialité, nous avons (cf. (5.1.6))

}{F7II-XVII-037}\begin{equation}\tag{5.2.5}
\alpha^*(\tau) = \xi \quad .
\end{equation}

\SGApage{II}{244}

\SGAvar{D'autre part, moyenant l'identification $\mathbb{P}(\underline{O}_X(1)^{r+1}) \simeq X \times \check{\mathbb{P}}^r$ , on a évidemment, posant $H = \underline{O}_X(1)$ ,

}{D'autre part, moyennant l'identification $\mathbb{P}(\underline{O}_X(1)^{r+1}) \simeq X \times \check{\mathbb{P}}^r$ , on a évidemment, posant $H = \underline{O}_X(1)$ ,

}{F7II-XVII-038}\begin{equation}\tag{5.2.6}
\tau = \mathrm{pr}_1^*(H) \otimes \mathrm{pr}_2^*(\eta) \quad .
\end{equation}
Désignons par $\pi : \mathbb{P}(N) \longrightarrow X$ le morphisme structural, on a

\SGAvar{\begin{equation*}
\begin{cases} \pi = \mathrm{pr}_1 \circ \alpha \\ \varphi = \mathrm{pr}_2 \circ \alpha \quad , \end{cases}
\end{equation*}
d'où, appliquant $\alpha^*$ aux deux membres de (5.2.6), on trouve, par (5.2.5)

\begin{equation}\tag{5.2.6}
\xi = \alpha^*(\tau) = \pi^*(H) \otimes \varphi^*(\eta) \quad .
\end{equation}
Dans $\mathrm{CH}(\ \mathbb{P}(N))$ , on a donc

}{\begin{equation*}
\begin{cases} \pi = \mathrm{pr}_1 \circ \alpha \\ \varphi = \mathrm{pr}_2 \circ \alpha \quad , \end{cases}
\end{equation*}
d'où, appliquant $\alpha^*$ aux deux membres de (5.2.6), on trouve, par (5.2.5)


\begin{equation}\tag{5.2.6$'$}
\xi = \alpha^*(\tau) = \pi^*(H) \otimes \varphi^*(\eta) \quad .
\end{equation}
Dans $\mathrm{CH}(\ \mathbb{P}(N))$ , on a donc

}{F7II-XVII-039}\begin{equation}\tag{5.2.7}
\varphi^*(\eta) = \xi - \pi^*(H) \quad ,
\end{equation}
et (5.2.3) se récrit

\SGAvar{\begin{equation}\tag{5.2.8}
\deg(\check{X}) = \frac{1}{\deg(\varphi)} \int_{\mathbb{P}(N)} (\xi - \pi^*(H))^{r-1} \quad .
\end{equation}
Le deuxieme membre de 5.2.8 ``se calcule'' grâce au lemme suivant (cf. 5.3.3) :

}{\begin{equation}\tag{5.2.8}
\deg(\check{X}) = \frac{1}{\deg(\varphi)} \int_{\mathbb{P}(N)} (\xi - \pi^*(H))^{r-1} \quad .
\end{equation}
Le deuxième membre de 5.2.8 ``se calcule'' grâce au lemme suivant (cf. 5.3.3) :

}{F7II-XVII-040}\underline{Lemme} 5.3. \underline{Soient} $X$ \underline{une variété quasi-projective}, \underline{lisse et connexe de dimension} $n$ , \underline{et} $M$ \underline{un faisceau  localement libre de rang fini} $a$ \underline{sur} $X$ . \underline{Notons par} $\xi$ \underline{le faisceau inversible} $\underline{O}_{\mathbb{P}(M)}(1)$ \underline{sur} $\mathbb{P}(M)$ ,
\SGApage{II}{245}
\underline{et par} $\pi : \mathbb{P}(M) \longrightarrow X$ \underline{le morphisme structural.} \underline{On sait que} $\mathrm{CH}(\ \mathbb{P}(M))$ \underline{est un} $\mathrm{CH}(X)$ \underline{module} (\underline{via} $\pi^*$ ) \underline{qui est libre,} \underline{de base} $1$ , $\xi, \ldots, \xi^{a-1}$ , \underline{où} $\xi$ \underline{satisfait à l'équation}

\begin{equation}\tag{5.3.1}
\xi^a - C_1\xi^{a-1} + C_2\xi^{a-2} - C_3\xi^{a-3} \ldots = 0 \quad ,
\end{equation}
\underline{les} $C_i$ \underline{dénotant les classes de Chern dans} $\mathrm{CH}(X)$ \underline{de} $M$ . \underline{On pose} $C(M) = 1+C_1+C_2+\ldots$, \underline{la classe de Chern totale de} $M$ ; \underline{c'est un élément inversible de} $\mathrm{CH}(X)$ .

\underline{Alors:}

(i) \underline{On a la formule}

\begin{equation}\tag{5.3.2}
\pi_*\!\left(\frac{\xi^{a-1}}{1 + \xi}\right) = \frac{1}{C(M)} \quad .
\end{equation}
(ii) \underline{Soit} $F(T) \in \mathrm{CH}(X)\ [T]$ \underline{un polynôme à coefficients dans} $\mathrm{CH}(X)$ , \underline{tel que le coefficient de} $T^i$ \underline{soit un élément homogène de degré} $n+a-1-i$ . \underline{Donc} $F(\xi)$ \underline{est un élément de degré} \underline{``maximum''} $n+a-1=\dim \mathbb{P}(M)$ \underline{dans} $\mathrm{CH}(\ \mathbb{P}(M))$. \underline{Ceci posé,} \underline{on a}

\begin{equation}\tag{5.3.3}
\int_{\mathbb{P}(M)} F(\xi) = (-1)^{a-1} \int_X \frac{F(-1)}{C(M)} \quad .
\end{equation}
\underline{Démonstration.} Commençons par déduire (5.3.3) de (5.3.2).

Ecrivons

\begin{equation}\tag{5.3.4}
\begin{gathered} \frac{1}{C(M)} = \sum_{j \geq 0} b_j \quad , \quad b_o = 1 \ , \ b_j \ \text{de degré } j \ \text{dans } \mathrm{CH}(X) \\ b_j = 0 \ \text{pour } j < 0 \quad , \end{gathered}
\end{equation}

\SGApage{II}{246}
de sorte que (5.3.2) se récrit

\begin{equation}\tag{5.3.5}
\pi_*(\xi^{a-1+j}) \;=\; (-1)^j b_j \qquad \mbox{pour} \quad a-1+j \geq 0 \ ,
\end{equation}
(compte tenu que $\pi_*(\xi^j) = 0$ si $j < a-1$ pour raison de degrés).

Déduisons-en (5.3.3). Par linéarité, il suffit de traiter le cas $F(T) = uT^{a-1+j}$ , $u$ homogène de degré $n-j$ dans $\mathrm{CH}(X)$ . Par (5.3.5) et la formule de projection pour $\pi$ , on a

\begin{equation}\tag{5.3.6}
\int_{\mathbb{P}(M)} F(\xi) \;=\; \int_{\mathbb{P}(M)} \pi^*(u) \cdot \xi^{a-1+j} \;=\; \int_X u \cdot \pi_*(\xi^{a-1+j}) \;=\; (-1)^j \int_X u \cdot b_j \ .
\end{equation}
D'autre part,

\begin{equation}\tag*{5.3.7.}
(-1)^{a-1} \int_X \frac{F(-1)}{C(M)} \;=\; (-1)^{a-1} \int_X \frac{(-1)^{a-1+j}u}{C(M)} \;=\; (-1)^j \int_X \sum_i \, u \, b_i \ ,
\end{equation}
et compte tenu de ce que $\deg(u) = n-j$ , et $\deg(b_i) = i$ , on a

\begin{equation*}
\int_X u \cdot b_i = 0 \qquad \mbox{si} \quad i \neq j \ ,
\end{equation*}
de sorte que (5.3.7) se récrit

\begin{equation}\tag{5.3.8}
(-1)^{a-1} \int_X \frac{F(-1)}{C(M)} \;=\; (-1)^j \int_X u \cdot b_j \ ,
\end{equation}
d'où l'implication annoncée.

\SGAvar{Démontrons maintenant (5.3.2), ou se qui revient au même, (5.3.5).


}{Démontrons maintenant (5.3.2), ou ce qui revient au même, (5.3.5).


}{F7II-XVII-041}\SGApage{II}{247}

On sait que

\begin{equation*}
\pi_*(\xi^j) = 0 \qquad \mbox{si} \quad j < a-1
\end{equation*}
pour une raison de degré, et que

\begin{equation*}
\pi_*(\xi^{a-1}) \;=\; 1_X
\end{equation*}
pour une raison géométrique ( $\xi^{a-1}$ rencontre chaque fibre "en un point"), de sorte que, par la formule de projection, on a

\SGAvar{\begin{equation}\tag{5.3.9}
\pi_*\left( \sum_{j=0}^{a-1} \pi^*(U_i) \, \xi^i \right) \;=\; U_{j-i} \ .
\end{equation}
Par (5.3.9), l'énoncé 5.3.5 est conséquence du lemme suivant  dans lequelle $\xi$ devient $T$ , et $C_i$ devient $S_i$, compte tenu de l'équation 5.3.1.

\underline{Lemme} 5.4. \underline{Soit} $a$ \underline{un entier} $> 0$ , \underline{et désignons par} $A$ \underline{l'anneau} $Z[[X_1,\ldots,X_a]]$ \underline{des séries formelles en} $a$ \underline{indéterminées.} \underline{Désignons par} $B$ \underline{la} $A$-\underline{algèbre}

}{\begin{equation}\tag{5.3.9}
\pi_*\left( \sum_{i=0}^{a-1} \pi^*(U_i) \, \xi^i \right)
\;=\; U_{a-1} \ .
\end{equation}
Par (5.3.9), l'énoncé 5.3.5 est conséquence du lemme suivant  dans lequel $\xi$ devient $T$ , et $C_i$ devient $S_i$, compte tenu de l'équation 5.3.1.

\underline{Lemme} 5.4. \underline{Soit} $a$ \underline{un entier} $> 0$ , \underline{et désignons par} $A$ \underline{l'anneau} $\mathbb{Z}[[X_1,\ldots,X_a]]$ \underline{des séries formelles en} $a$ \underline{indéterminées.} \underline{Désignons par} $B$ \underline{la} $A$-\underline{algèbre}

}{F7II-XVII-042}\begin{equation*}
B = A[[T]] \, / \, \left( \prod_{i=1}^{a} (T-X_i) \right) \ ,
\end{equation*}
\underline{de sorte que} $B$ \underline{est un} $A$-\underline{module libre} , \underline{à base} $1,T,\ldots,T^{a-1}$ ,

\begin{equation}\tag{5.4.1}
B = A \, T^{a-1} \oplus A \, T^{a-2} \oplus \ldots \oplus A \ .
\end{equation}

\SGApage{II}{248}

\underline{Désignons par} $S_i \in A$ \underline{la} $i$-\underline{ième fonction symétrique}

\begin{equation*}
S_i = \sum_{1 \leq \alpha_1 < \ldots < \alpha_i \leq a} X_{\alpha_1} \ldots X_{\alpha_i} \ ,
\end{equation*}
\underline{de sorte que}

\begin{equation}\tag{5.4.2}
\sum_{i=0}^{a} (-1)^i \, S_i \, T^{a-i} = 0 \quad \mbox{dans } B \ .
\end{equation}
\underline{L'élément} $1 + S_1 + S_2 + \ldots + S_a$ \underline{est inversible dans} $A$ , \underline{nous écrivons son inverse}

\begin{equation}\tag{5.4.3}
\frac{1}{\sum\limits_{i=0}^{a} S_i} = \sum_{j \geq 0} b_j \ , \qquad b_j \ \mbox{homogène de degré } j \mbox{ dans } A \ , \quad b_o = 1 \ .
\end{equation}
\underline{D'autre part, pour tout entier} $j \geq 0$ \underline{on définit un élément} $a_j \in A$ \underline{homogène de degré} $j$ \underline{en écrivant, par} 5.4.1

\begin{equation}\tag{5.4.4}
T^{a-1+j} = a_j \, T^{a-1} + ? \, T^{a-2} + ? \, T^{a-3} + \ldots + ? \ , \quad \mbox{les } ? \in A \ .
\end{equation}
\underline{Ceci posé, on a pour tout} $j \geq 0$

\begin{equation}\tag{5.4.5}
a_j = (-1)^j \, b_j \ .
\end{equation}
\underline{Démonstration}. Compte tenu de la définition (5.4.3) des $b_j$ , tout revient à démontrer qu'on a 


\SGApage{II}{249}

\begin{equation}\tag{5.4.6}
 \begin{array}{l} a_o = 1 \\[4pt] \sum\limits_{i=0}^{a} (-1)^i \, S_i \, a_{j-i} = 0 \qquad \mbox{si } j > 0 \ . \end{array}
\end{equation}
Or, on a évidemment $a_o = 1$ , et la deuxième relation (pour $j > 0$ donné) n'est autre que celle déduite ,en prenant le coefficient de $T^{a-1}$ (via 5.4.1), de la relation

\SGAvar{\begin{equation}\tag{5.4.7}
\sum_{i=0} (-1)^i \, S_i \, T^{a-1+j-i} = 0 \quad \mbox{dans } B \ ,
\end{equation}
relation obtenue en multipliant 5.4.2 par $T^{j-1}$ , cqfd.

}{\begin{equation}\tag{5.4.7}
\sum_{i=0}^{a} (-1)^i \, S_i \, T^{a-1+j-i} = 0 \quad \mbox{dans } B \ ,
\end{equation}
relation obtenue en multipliant 5.4.2 par $T^{j-1}$ , cqfd.

}{F7II-XVII-043}\underline{Corollaire} 5.5. \underline{On a la formule}

\begin{equation}\tag{5.5.1}
\deg(\check{X}) = \frac{(-1)^n}{\deg(\varphi)} \int_X \frac{(1+H)^{r-1}}{C(N)} \ .
\end{equation}
\underline{Démonstration}. On applique (5.3.3) (pour $M = N$ , le faisceau normal de $X \hookrightarrow \mathbb{P}^r$) à (5.2.8).

Ecrivons, pour une variété quasi-projective et lisse $S$ , sa classe de Chern totale (cf. 5.3)

\begin{equation}\tag{5.5.2}
C(S) \stackrel{\mathrm{dfn}}{=} C(T_S)
\end{equation}
où $T_S$ désigne le fibré \underline{tangent} de $S$ .


\SGApage{II}{250}

La duale de la suite exacte (5.1.2) montre qu'on a

\begin{equation}\tag{5.5.3}
C(\mathbb{P}^r) \;=\; (1 + H)^{r+1} \qquad \text{dans } \mathrm{CH}(\mathbb{P}^r) \quad .
\end{equation}
La duale de la suite exacte(5.1.1)montre qu'on a

\begin{equation}\tag{5.5.4}
C(N)\, C(X) = C(\mathbb{P}^r)|X \qquad \text{dans } \mathrm{CH}(X) \quad .
\end{equation}
Mettant ensemble (5.5.3) et (5.5.4),(5.5.1)se récrit

\underline{Corollaire} 5.6. \qquad \underline{On a la formule}

\SGAvar{\begin{equation}\tag{5.6.1}
\deg(\check{X}) \;=\; \frac{(-1)^r}{\deg(\varphi)} \; \int_X \frac{C(X)}{(1+H)^2} \quad .
\end{equation}
(5.7) Pour une variété $S$ projective et lisse, on définit sa \underline{caractéristique d'Euler-Poincaré par}

\begin{equation}\tag{5.7.1}
\chi(S) \;=\; \int_S C(S) \quad .
\end{equation}
\underline{Proposition} 5.7.2. \qquad \underline{Désignons par} $Y$ \underline{une section hyperplanelisse de} $X$ , \underline{et par} $\Delta$ \underline{une section lisse de} $X$ \underline{par un sous-espace linéaire de} $\mathbb{P}^r$ \underline{de codimension deux.} \underline{On a}

}{\begin{equation}\tag{5.6.1}
\deg(\check{X}) \;=\; \frac{(-1)^n}{\deg(\varphi)} \; \int_X \frac{C(X)}{(1+H)^2} \quad .
\end{equation}
(5.7) Pour une variété $S$ projective et lisse, on définit sa \underline{caractéristique d'Euler-Poincaré par}

\begin{equation}\tag{5.7.1}
\chi(S) \;=\; \int_S C(S) \quad .
\end{equation}
\underline{Proposition} 5.7.2. \qquad \underline{Désignons par} $Y$ \underline{une section hyperplane lisse de} $X$ , \underline{et par} $\Delta$ \underline{une section lisse de} $X$ \underline{par un sous-espace linéaire de} $\mathbb{P}^r$ \underline{de codimension deux.} \underline{On a}

}{F7II-XVII-044}\begin{equation}\tag{5.7.3}
\deg(\check{X}) = \frac{(-1)^n}{\deg(\varphi)} \quad \left[ \chi(X) + \chi(\Delta) - 2\chi(Y) \right] \quad .
\end{equation}

\SGApage{II}{251}

\underline{Démonstration.} \quad Notons $i : Y \hookrightarrow X$ l'inclusion. On a

\begin{equation*}
N_{Y \longrightarrow X} \;=\; H|Y \;=\; i^{*}H \quad ,
\end{equation*}
de sorte que la suite exacte

\begin{equation*}
0 \longrightarrow T_Y \longrightarrow i^{*}T_X \longrightarrow i^{*}H \longrightarrow 0
\end{equation*}
donne

\begin{equation}\tag{5.7.4}
C(Y) \;=\; i^{*} \left( \frac{C(X)}{1+H} \right) \quad .
\end{equation}
On a donc

\begin{equation*}
\chi(Y) \;=\; \int_Y i^{*} \left( \frac{C(X)}{1+H} \right)
\end{equation*}
et, en appliquant la formule de projection

\begin{equation}\tag{5.7.5}
\chi(Y) \;=\; \int_X i_{*}i^{*} \left( \frac{C(X)}{1+H} \right) \;=\; \int_X \frac{H \cdot C(X)}{1 + H} \quad .
\end{equation}
Itérant ce calcul, on trouve pour $\Delta$

\begin{equation}\tag{5.7.6}
\chi(\Delta) \;=\; \int_Y \frac{H\ C(Y)}{1 + H} \;=\; \int_Y i^{*} \left( \frac{H}{(1+H)} \cdot \frac{C(X)}{(1+H)} \right) \;=\; \int_X \frac{H^2\ C(X)}{(1+H)^2} \quad ,
\end{equation}
de sorte que

\begin{equation*}
\begin{aligned} \chi(X) - 2\chi(Y)+\chi(\Delta) &= \int_X C(X) \left[ 1 - \frac{2\ H}{1+H} + \frac{H^2}{(1+H)^2} \right] \\ &= \int_X \frac{C(X)}{(1 + H)^2} \quad , \end{aligned}
\end{equation*}
ce qui achève la démonstration, compte tenu de 5.6.1.


\SGApage{II}{252}


\SGAsectiontoc{6. Le cas d'une base générale}{6}
6. \underline{Le cas d'une base générale}

6.1. Soient $S$ un schéma et $\bar{s}$ un point géométrique de $S$, défini par une clôture algébrique du corps résiduel $k(s)$ pour $s \in S$. Désignons par $\mathbb{P}^r$ l'espace projectif à $r$ dimensions sur $S$ , par $\check{\mathbb{P}}^r$ l'espace projectif dual et par $I \subset \mathbb{P}^r \times_S \check{\mathbb{P}}^r$ la variété d'incidence. Soit $D \subset \check{\mathbb{P}}^r$ une droite, définie par un point à valeurs dans $S$ d'une grassmannienne convenable, et $\Delta$ l'axe du pinceau d'hyperplans correspondant (cf. 2.1).

\SGAvar{Soit $X$ un $S$-schéma propre et lisse, purement de dimension $n \geq 1$, muni d'un plongement projectif $X \hookrightarrow \mathbb{P}^r$ . Soient $\widetilde{X} = (X \times_S D) \cap I$ le schéma sur $D$ de fibres les sections hyperplanes de $X$ paramètrées par $D$, et $u$ la projection de $\widetilde{X}$ sur $D$.

}{Soit $X$ un $S$-schéma propre et lisse, purement de dimension $n \geq 1$, muni d'un plongement projectif $X \hookrightarrow \mathbb{P}^r$ . Soient $\widetilde{X} = (X \times_S D) \cap I$ le schéma sur $D$ de fibres les sections hyperplanes de $X$ paramétrées par $D$, et $u$ la projection de $\widetilde{X}$ sur $D$.

}{F7II-XVII-045}Notons par un indice $\bar{s}$ le changement de base par $\bar{s} \longrightarrow S$.

Supposons que $\Delta_{\bar{s}}$ soit transverse à $X_{\bar{s}}$ . Alors, au-dessus d'un voisinage $U$ de $s$ dans $S$, $\Delta$ est transverse à $X$ . On suppose dans ce qui suit que $U = S$ . Le schéma $\widetilde{X}$ est alors lisse sur $S$ , et un voisinage de l'image réciproque de $\Delta$ dans $\widetilde{X}$ est lisse sur $D$ .

Supposons de plus qu'une des sections hyperplanes de $X_{\bar{s}}$ appartenant au pinceau $D_{\bar{s}}$ soit lisse. Alors, après que l'on ait remplacé $S$ par un voisinage convenable de $s$, le sous-schéma $F$ de $\widetilde{X}$ défini par l'idéal jacobien $J^{n-1}(\widetilde{X}/D)$ est fini et plat sur $S$ .


\SGApage{II}{253}

6.2. \underline{Le lieu exceptionnel} du pinceau $D$ est le diviseur de Cartier relatif $E$ de $D/S$ défini par le $\mathcal{O}_D$-Module $u_*\mathcal{O}_F$, fini et plat sur $S$ . Localement sur $D$, si $\mathcal{O}_D^n \xrightarrow{\ v\ } \mathcal{O}_D^n \longrightarrow u_*\mathcal{O}_F \longrightarrow 0$ est une présentation de $u_*\mathcal{O}_F$ , $E$ a pour équation $\det(v) = 0$. La formation de $E$ est compatible à tout changement de base.

\SGAvar{Le diviseur $E_{\bar{s}}$ de $D_{\bar{s}}$ s'écrit $\Sigma\, n(P).P$, pour $n(P)$ la somme des colongeurs de l'idéal jacobien $J^{n-1}(\widetilde{X}/D)$ en les points critiques de $u$ d'image $P$ .

}{Le diviseur $E_{\bar{s}}$ de $D_{\bar{s}}$ s'écrit $\Sigma\, n(P).P$, pour $n(P)$ la somme des colongueurs de l'idéal jacobien $J^{n-1}(\widetilde{X}/D)$ en les points critiques de $u$ d'image $P$ .

}{F7II-XVII-046}\underline{Proposition} 6.3. \underline{Si} $D_{\bar{s}}$ \underline{est un pinceau de Lefschetz de sections hyperplanes de} $X_{\bar{s}}$ , \underline{et si} car $k(s) \neq 2$ \underline{ou que} $n$ \underline{est pair}, \underline{il existe un voisinage} $U$ \underline{de} $s$ \underline{dans} $S$ \underline{tel que}

(a) \underline{Pour tout point géométrique} $\bar{t}$ \underline{de} $U$, $D_{\bar{t}}$ \underline{est un pinceau de Lefschetz de sections hyperplanes de} $X_{\bar{t}}$

(b) \underline{Au-dessus de} $U$, \underline{le lieu exceptionnel} $E$ \underline{est étale sur} $S$.

Pour prouver (a) et (b), il suffit de noter que, sous les hypothèses faites, $D_{\bar{s}}$ vérifie la condition c) de 2.2 si et seulement si le diviseur $E_{\bar{s}}$ est réduit. On pourrait, en invoquant 3.3, n'utiliser que la moitié la plus facile $(\Longrightarrow)$ de cette équivalence.

\SGAvar{6.3. La proposition devient fausse dans le cas d'exception cité. Toutefois, si $S$ est \underline{de caractéristique} 2 et que $n$ est impair, il existe un voisinage $U$ , de $S$ au-dessus duquel (a) est vrai, et on peut définir un diviseur $E'$ étale sur $U$ tel que $E = 2 E'$. Nous ne ferons pas usage de ce résultat.


}{6.4. La proposition devient fausse dans le cas d'exception cité. Toutefois, si $S$ est \underline{de caractéristique} 2 et que $n$ est impair, il existe un voisinage $U$ , de $s$ au-dessus duquel (a) est vrai, et on peut définir un diviseur $E'$ étale sur $U$ tel que $E = 2 E'$. Nous ne ferons pas usage de ce résultat.


}{F7II-XVII-047}\underline{Références}

[1] GROTHENDIECK, A. \quad La théorie des Classes de Chern, \\ \hspace*{5.5em} Bull. Soc. Math. France, 86, 1958, 137-154.
\selectlanguage{french}

\ifnum\SGAform=1
\clearpage\section*{Apparat critique}\small
Les lectures ci-dessous sont données dans le TeX exact des témoins, afin de préserver les distinctions de notation. Les indices décimaux renvoient aux différences individuelles du passage marqué.

\par\medskip\noindent\hypertarget{app-F7II-XVII-001.1}{}\textbf{F7II-XVII-001.1}\quad p. 213; printed misprint.\par
\noindent Imprimé : \emph{(vide)}\par
\noindent Corrigé : {\ttfamily\footnotesize -1\allowbreak{}}\par
\noindent Source : 7II:316. \nolinkurl{fixes/MISPRINTS_7II.json#316}\par

\par\medskip\noindent\hypertarget{app-F7II-XVII-001.2}{}\textbf{F7II-XVII-001.2}\quad p. 213; misprint\_pass.\par
\noindent Imprimé : {\ttfamily\footnotesize series\allowbreak{}}\par
\noindent Corrigé : {\ttfamily\footnotesize séries\allowbreak{}}\par
\noindent Source : 7II:317. \nolinkurl{fixes/FIXLOG_7II_CURRENT.jsonl#680}; \nolinkurl{fixes/MISPRINTS_7II.json#317}\par

\par\medskip\noindent\hypertarget{app-F7II-XVII-002.1}{}\textbf{F7II-XVII-002.1}\quad p. 214; misprint\_pass.\par
\noindent Imprimé : {\ttfamily\footnotesize Comme\allowbreak{}}\par
\noindent Corrigé : {\ttfamily\footnotesize comme\allowbreak{}}\par
\noindent Source : 7II:318. \nolinkurl{fixes/FIXLOG_7II_CURRENT.jsonl#681}; \nolinkurl{fixes/MISPRINTS_7II.json#318}\par

\par\medskip\noindent\hypertarget{app-F7II-XVII-002.2}{}\textbf{F7II-XVII-002.2}\quad p. 214; misprint\_pass.\par
\noindent Imprimé : {\ttfamily\footnotesize jacobienne\allowbreak{}}\par
\noindent Corrigé : {\ttfamily\footnotesize jacobien\allowbreak{}}\par
\noindent Source : 7II:319. \nolinkurl{fixes/FIXLOG_7II_CURRENT.jsonl#682}; \nolinkurl{fixes/MISPRINTS_7II.json#319}\par

\par\medskip\noindent\hypertarget{app-F7II-XVII-003.1}{}\textbf{F7II-XVII-003.1}\quad p. 214; misprint\_pass.\par
\noindent Imprimé : {\ttfamily\footnotesize quadratique\allowbreak{}}\par
\noindent Corrigé : {\ttfamily\footnotesize quadratiques\allowbreak{}}\par
\noindent Source : 7II:320. \nolinkurl{fixes/FIXLOG_7II_CURRENT.jsonl#683}; \nolinkurl{fixes/MISPRINTS_7II.json#320}\par

\par\medskip\noindent\hypertarget{app-F7II-XVII-003.2}{}\textbf{F7II-XVII-003.2}\quad p. 214; misprint\_pass.\par
\noindent Imprimé : {\ttfamily\footnotesize X\allowbreak{}}\par
\noindent Corrigé : {\ttfamily\footnotesize IX\allowbreak{}}\par
\noindent Source : 7II:321. \nolinkurl{fixes/FIXLOG_7II_CURRENT.jsonl#684}; \nolinkurl{fixes/MISPRINTS_7II.json#321}\par

\par\medskip\noindent\hypertarget{app-F7II-XVII-003.3}{}\textbf{F7II-XVII-003.3}\quad p. 214; misprint\_pass.\par
\noindent Imprimé : \emph{(vide)}\par
\noindent Corrigé : {\ttfamily\footnotesize '\allowbreak{}}\par
\noindent Source : 7II:322. \nolinkurl{fixes/FIXLOG_7II_CURRENT.jsonl#685}; \nolinkurl{fixes/MISPRINTS_7II.json#322}\par

\par\medskip\noindent\hypertarget{app-F7II-XVII-003.4}{}\textbf{F7II-XVII-003.4}\quad p. 214; misprint\_pass.\par
\noindent Imprimé : \emph{(vide)}\par
\noindent Corrigé : {\ttfamily\footnotesize -1\allowbreak{}}\par
\noindent Source : 7II:323. \nolinkurl{fixes/FIXLOG_7II_CURRENT.jsonl#686}; \nolinkurl{fixes/MISPRINTS_7II.json#323}\par

\par\medskip\noindent\hypertarget{app-F7II-XVII-003.5}{}\textbf{F7II-XVII-003.5}\quad p. 214; misprint\_pass.\par
\noindent Imprimé : \emph{(vide)}\par
\noindent Corrigé : {\ttfamily\footnotesize ,\allowbreak{}}\par
\noindent Source : 7II:324. \nolinkurl{fixes/FIXLOG_7II_CURRENT.jsonl#687}; \nolinkurl{fixes/MISPRINTS_7II.json#324}\par

\par\medskip\noindent\hypertarget{app-F7II-XVII-004.1}{}\textbf{F7II-XVII-004.1}\quad p. 215; printed misprint.\par
\noindent Imprimé : {\ttfamily\footnotesize \textbackslash{}mathbb\{P\}\allowbreak{}}\par
\noindent Corrigé : {\ttfamily\footnotesize D\allowbreak{}}\par
\noindent Source : 7II:326. \nolinkurl{fixes/MISPRINTS_7II.json#326}\par

\par\medskip\noindent\hypertarget{app-F7II-XVII-005.1}{}\textbf{F7II-XVII-005.1}\quad p. 216; misprint\_pass.\par
\noindent Imprimé : \emph{(vide)}\par
\noindent Corrigé : {\ttfamily\footnotesize \textbackslash{}check\{\allowbreak{}}\par
\noindent Source : 7II:327. \nolinkurl{fixes/FIXLOG_7II_CURRENT.jsonl#688}; \nolinkurl{fixes/MISPRINTS_7II.json#327}\par

\par\medskip\noindent\hypertarget{app-F7II-XVII-005.2}{}\textbf{F7II-XVII-005.2}\quad p. 216; misprint\_pass.\par
\noindent Imprimé : \emph{(vide)}\par
\noindent Corrigé : {\ttfamily\footnotesize \}\allowbreak{}}\par
\noindent Source : 7II:327. \nolinkurl{fixes/FIXLOG_7II_CURRENT.jsonl#688}; \nolinkurl{fixes/MISPRINTS_7II.json#327}\par

\par\medskip\noindent\hypertarget{app-F7II-XVII-005.3}{}\textbf{F7II-XVII-005.3}\quad p. 216; misprint\_pass.\par
\noindent Imprimé : {\ttfamily\footnotesize Rappellons\allowbreak{}}\par
\noindent Corrigé : {\ttfamily\footnotesize Rappelons\allowbreak{}}\par
\noindent Source : 7II:328. \nolinkurl{fixes/FIXLOG_7II_CURRENT.jsonl#689}\par

\par\medskip\noindent\hypertarget{app-F7II-XVII-005.4}{}\textbf{F7II-XVII-005.4}\quad p. 216; inherited corrected-layer difference.\par
\noindent Imprimé : {\ttfamily\footnotesize Segre\allowbreak{}}\par
\noindent Corrigé : {\ttfamily\footnotesize Veronese\allowbreak{}}\par
\noindent Source : ADOPT-F7II-XVII-005.4. \nolinkurl{fr/sga7ii/diplomatic/expose_XVII_body.tex}; \nolinkurl{fr/sga7ii/current/expose_XVII_body.tex}\par

\par\medskip\noindent\hypertarget{app-F7II-XVII-005.5}{}\textbf{F7II-XVII-005.5}\quad p. 216; misprint\_pass.\par
\noindent Imprimé : {\ttfamily\footnotesize coordonnés\allowbreak{}}\par
\noindent Corrigé : {\ttfamily\footnotesize coordonnées\allowbreak{}}\par
\noindent Source : 7II:329. \nolinkurl{fixes/FIXLOG_7II_CURRENT.jsonl#690}\par

\par\medskip\noindent\hypertarget{app-F7II-XVII-005.6}{}\textbf{F7II-XVII-005.6}\quad p. 216; misprint\_pass.\par
\noindent Imprimé : {\ttfamily\footnotesize ,\allowbreak{}}\par
\noindent Corrigé : \emph{(vide)}\par
\noindent Source : 7II:330. \nolinkurl{fixes/FIXLOG_7II_CURRENT.jsonl#691}; \nolinkurl{fixes/MISPRINTS_7II.json#330}\par

\par\medskip\noindent\hypertarget{app-F7II-XVII-005.7}{}\textbf{F7II-XVII-005.7}\quad p. 216; printed misprint.\par
\noindent Imprimé : {\ttfamily\footnotesize o\allowbreak{}}\par
\noindent Corrigé : {\ttfamily\footnotesize r\allowbreak{}}\par
\noindent Source : 7II:331. \nolinkurl{fixes/MISPRINTS_7II.json#331}\par

\par\medskip\noindent\hypertarget{app-F7II-XVII-006.1}{}\textbf{F7II-XVII-006.1}\quad p. 217; misprint\_pass.\par
\noindent Imprimé : {\ttfamily\footnotesize ;\allowbreak{}}\par
\noindent Corrigé : {\ttfamily\footnotesize ,\allowbreak{}}\par
\noindent Source : 7II:333. \nolinkurl{fixes/FIXLOG_7II_CURRENT.jsonl#692}; \nolinkurl{fixes/MISPRINTS_7II.json#333}\par

\par\medskip\noindent\hypertarget{app-F7II-XVII-007.1}{}\textbf{F7II-XVII-007.1}\quad p. 217; misprint\_in\_print.\par
\noindent Imprimé : {\ttfamily\footnotesize \textbackslash{}stackrel\allowbreak{}}\par
\noindent Corrigé : {\ttfamily\footnotesize \textbackslash{}xrightarrow\allowbreak{}}\par
\noindent Source : FIX-7II-D-0225. \nolinkurl{fixes/FIXLOG_7II.jsonl#225}\par

\par\medskip\noindent\hypertarget{app-F7II-XVII-007.2}{}\textbf{F7II-XVII-007.2}\quad p. 217; misprint\_in\_print.\par
\noindent Imprimé : \emph{(vide)}\par
\noindent Corrigé : {\ttfamily\footnotesize \textbackslash{}\allowbreak{} }\par
\noindent Source : FIX-7II-D-0225. \nolinkurl{fixes/FIXLOG_7II.jsonl#225}\par

\par\medskip\noindent\hypertarget{app-F7II-XVII-007.3}{}\textbf{F7II-XVII-007.3}\quad p. 217; designated-witness attribution.\par
\noindent Imprimé : {\ttfamily\footnotesize \}\{\textbackslash{}relbar\textbackslash{}joinrel\textbackslash{}re\allowbreak{}lbar\textbackslash{}joinrel\textbackslash{}relbar\allowbreak{}}\par
\noindent Corrigé : {\ttfamily\footnotesize \textbackslash{}\allowbreak{} }\par
\noindent Source : WITNESS-F7II-XVII-007.3. \nolinkurl{fr/sga7ii/diplomatic/expose_XVII_body.tex}; \nolinkurl{fr/sga7ii/current/expose_XVII_body.tex}\par

\par\medskip\noindent\hypertarget{app-F7II-XVII-008.1}{}\textbf{F7II-XVII-008.1}\quad p. 218; misprint\_pass.\par
\noindent Imprimé : \emph{(vide)}\par
\noindent Corrigé : {\ttfamily\footnotesize \textbackslash{}underline\{\allowbreak{}}\par
\noindent Source : 7II:335. \nolinkurl{fixes/FIXLOG_7II_CURRENT.jsonl#693}\par

\par\medskip\noindent\hypertarget{app-F7II-XVII-008.2}{}\textbf{F7II-XVII-008.2}\quad p. 218; misprint\_pass.\par
\noindent Imprimé : \emph{(vide)}\par
\noindent Corrigé : {\ttfamily\footnotesize \}\allowbreak{}}\par
\noindent Source : 7II:335. \nolinkurl{fixes/FIXLOG_7II_CURRENT.jsonl#693}\par

\par\medskip\noindent\hypertarget{app-F7II-XVII-008.3}{}\textbf{F7II-XVII-008.3}\quad p. 218; misprint\_pass.\par
\noindent Imprimé : {\ttfamily\footnotesize tel\allowbreak{}}\par
\noindent Corrigé : {\ttfamily\footnotesize tels\allowbreak{}}\par
\noindent Source : 7II:337. \nolinkurl{fixes/FIXLOG_7II_CURRENT.jsonl#694}\par

\par\medskip\noindent\hypertarget{app-F7II-XVII-008.4}{}\textbf{F7II-XVII-008.4}\quad p. 219; designated-witness attribution.\par
\noindent Imprimé : {\ttfamily\footnotesize ,\allowbreak{} }\par
\noindent Corrigé : \emph{(vide)}\par
\noindent Source : WITNESS-F7II-XVII-008.4. \nolinkurl{fr/sga7ii/diplomatic/expose_XVII_body.tex}; \nolinkurl{fr/sga7ii/current/expose_XVII_body.tex}\par

\par\medskip\noindent\hypertarget{app-F7II-XVII-008.5}{}\textbf{F7II-XVII-008.5}\quad p. 219; misprint\_pass.\par
\noindent Imprimé : \emph{(vide)}\par
\noindent Corrigé : {\ttfamily\footnotesize  sur\allowbreak{}}\par
\noindent Source : 7II:339. \nolinkurl{fixes/FIXLOG_7II_CURRENT.jsonl#696}; \nolinkurl{fixes/MISPRINTS_7II.json#339}\par

\par\medskip\noindent\hypertarget{app-F7II-XVII-009.1}{}\textbf{F7II-XVII-009.1}\quad p. 219; printed misprint.\par
\noindent Imprimé : \emph{(vide)}\par
\noindent Corrigé : {\ttfamily\footnotesize \textbackslash{}check\{\allowbreak{}}\par
\noindent Source : 7II:340. \nolinkurl{fixes/MISPRINTS_7II.json#340}\par

\par\medskip\noindent\hypertarget{app-F7II-XVII-009.2}{}\textbf{F7II-XVII-009.2}\quad p. 219; printed misprint.\par
\noindent Imprimé : \emph{(vide)}\par
\noindent Corrigé : {\ttfamily\footnotesize \}\allowbreak{}}\par
\noindent Source : 7II:340. \nolinkurl{fixes/MISPRINTS_7II.json#340}\par

\par\medskip\noindent\hypertarget{app-F7II-XVII-010.1}{}\textbf{F7II-XVII-010.1}\quad p. 220; misprint\_pass.\par
\noindent Imprimé : \emph{(vide)}\par
\noindent Corrigé : {\ttfamily\footnotesize \textbackslash{}check\{\allowbreak{}}\par
\noindent Source : 7II:341. \nolinkurl{fixes/FIXLOG_7II_CURRENT.jsonl#697}; \nolinkurl{fixes/MISPRINTS_7II.json#341}\par

\par\medskip\noindent\hypertarget{app-F7II-XVII-010.2}{}\textbf{F7II-XVII-010.2}\quad p. 220; misprint\_pass.\par
\noindent Imprimé : \emph{(vide)}\par
\noindent Corrigé : {\ttfamily\footnotesize \}\allowbreak{}}\par
\noindent Source : 7II:341. \nolinkurl{fixes/FIXLOG_7II_CURRENT.jsonl#697}; \nolinkurl{fixes/MISPRINTS_7II.json#341}\par

\par\medskip\noindent\hypertarget{app-F7II-XVII-010.3}{}\textbf{F7II-XVII-010.3}\quad p. 220; misprint\_pass.\par
\noindent Imprimé : {\ttfamily\footnotesize G\allowbreak{}}\par
\noindent Corrigé : {\ttfamily\footnotesize H\allowbreak{}}\par
\noindent Source : 7II:342. \nolinkurl{fixes/FIXLOG_7II_CURRENT.jsonl#698}\par

\par\medskip\noindent\hypertarget{app-F7II-XVII-010.4}{}\textbf{F7II-XVII-010.4}\quad p. 220; misprint\_pass.\par
\noindent Imprimé : {\ttfamily\footnotesize une\allowbreak{}}\par
\noindent Corrigé : {\ttfamily\footnotesize un\allowbreak{}}\par
\noindent Source : 7II:343. \nolinkurl{fixes/FIXLOG_7II_CURRENT.jsonl#699}; \nolinkurl{fixes/MISPRINTS_7II.json#343}\par

\par\medskip\noindent\hypertarget{app-F7II-XVII-011.1}{}\textbf{F7II-XVII-011.1}\quad p. 220; misprint\_pass.\par
\noindent Imprimé : {\ttfamily\footnotesize les-dits\allowbreak{}}\par
\noindent Corrigé : {\ttfamily\footnotesize lesdits\allowbreak{}}\par
\noindent Source : 7II:344. \nolinkurl{fixes/FIXLOG_7II_CURRENT.jsonl#700}\par

\par\medskip\noindent\hypertarget{app-F7II-XVII-012.1}{}\textbf{F7II-XVII-012.1}\quad p. 221; misprint\_pass.\par
\noindent Imprimé : {\ttfamily\footnotesize la\allowbreak{}}\par
\noindent Corrigé : {\ttfamily\footnotesize là\allowbreak{}}\par
\noindent Source : 7II:345. \nolinkurl{fixes/FIXLOG_7II_CURRENT.jsonl#701}\par

\par\medskip\noindent\hypertarget{app-F7II-XVII-013.1}{}\textbf{F7II-XVII-013.1}\quad p. 221; misprint\_pass.\par
\noindent Imprimé : {\ttfamily\footnotesize tel\allowbreak{}}\par
\noindent Corrigé : {\ttfamily\footnotesize tels\allowbreak{}}\par
\noindent Source : 7II:346. \nolinkurl{fixes/FIXLOG_7II_CURRENT.jsonl#702}\par

\par\medskip\noindent\hypertarget{app-F7II-XVII-014.1}{}\textbf{F7II-XVII-014.1}\quad p. 222; misprint\_pass.\par
\noindent Imprimé : {\ttfamily\footnotesize coordonnés\allowbreak{}}\par
\noindent Corrigé : {\ttfamily\footnotesize coordonnées\allowbreak{}}\par
\noindent Source : 7II:347. \nolinkurl{fixes/FIXLOG_7II_CURRENT.jsonl#703}\par

\par\medskip\noindent\hypertarget{app-F7II-XVII-014.2}{}\textbf{F7II-XVII-014.2}\quad p. 222; misprint\_pass.\par
\noindent Imprimé : {\ttfamily\footnotesize coordonnés\allowbreak{}}\par
\noindent Corrigé : {\ttfamily\footnotesize coordonnées\allowbreak{}}\par
\noindent Source : 7II:348. \nolinkurl{fixes/FIXLOG_7II_CURRENT.jsonl#704}; \nolinkurl{fixes/MISPRINTS_7II.json#348}\par

\par\medskip\noindent\hypertarget{app-F7II-XVII-014.3}{}\textbf{F7II-XVII-014.3}\quad p. 222; misprint\_pass.\par
\noindent Imprimé : {\ttfamily\footnotesize coordonnés\allowbreak{}}\par
\noindent Corrigé : {\ttfamily\footnotesize coordonnées\allowbreak{}}\par
\noindent Source : 7II:349. \nolinkurl{fixes/FIXLOG_7II_CURRENT.jsonl#705}\par

\par\medskip\noindent\hypertarget{app-F7II-XVII-015.1}{}\textbf{F7II-XVII-015.1}\quad p. 222; printed misprint.\par
\noindent Imprimé : \emph{(vide)}\par
\noindent Corrigé : {\ttfamily\footnotesize r-\allowbreak{}}\par
\noindent Source : 7II:350. \nolinkurl{fixes/MISPRINTS_7II.json#350}\par

\par\medskip\noindent\hypertarget{app-F7II-XVII-015.2}{}\textbf{F7II-XVII-015.2}\quad p. 222; printed misprint.\par
\noindent Imprimé : {\ttfamily\footnotesize -r\allowbreak{}}\par
\noindent Corrigé : \emph{(vide)}\par
\noindent Source : 7II:350. \nolinkurl{fixes/MISPRINTS_7II.json#350}\par

\par\medskip\noindent\hypertarget{app-F7II-XVII-016.1}{}\textbf{F7II-XVII-016.1}\quad p. 222; designated-witness attribution.\par
\noindent Imprimé : {\ttfamily\footnotesize \textbackslash{}lambda\allowbreak{}}\par
\noindent Corrigé : {\ttfamily\footnotesize \textbackslash{}Lambda\allowbreak{}}\par
\noindent Source : WITNESS-F7II-XVII-016.1. \nolinkurl{fr/sga7ii/diplomatic/expose_XVII_body.tex}; \nolinkurl{fr/sga7ii/current/expose_XVII_body.tex}\par

\par\medskip\noindent\hypertarget{app-F7II-XVII-016.2}{}\textbf{F7II-XVII-016.2}\quad p. 222; designated-witness attribution.\par
\noindent Imprimé : {\ttfamily\footnotesize \textbackslash{}lambda\allowbreak{}}\par
\noindent Corrigé : {\ttfamily\footnotesize \textbackslash{}Lambda\allowbreak{}}\par
\noindent Source : WITNESS-F7II-XVII-016.2. \nolinkurl{fr/sga7ii/diplomatic/expose_XVII_body.tex}; \nolinkurl{fr/sga7ii/current/expose_XVII_body.tex}\par

\par\medskip\noindent\hypertarget{app-F7II-XVII-016.3}{}\textbf{F7II-XVII-016.3}\quad p. 222; designated-witness attribution.\par
\noindent Imprimé : \emph{(vide)}\par
\noindent Corrigé : {\ttfamily\footnotesize \{i=\allowbreak{}}\par
\noindent Source : WITNESS-F7II-XVII-016.3. \nolinkurl{fr/sga7ii/diplomatic/expose_XVII_body.tex}; \nolinkurl{fr/sga7ii/current/expose_XVII_body.tex}\par

\par\medskip\noindent\hypertarget{app-F7II-XVII-016.4}{}\textbf{F7II-XVII-016.4}\quad p. 222; designated-witness attribution.\par
\noindent Imprimé : \emph{(vide)}\par
\noindent Corrigé : {\ttfamily\footnotesize \}\allowbreak{}}\par
\noindent Source : WITNESS-F7II-XVII-016.4. \nolinkurl{fr/sga7ii/diplomatic/expose_XVII_body.tex}; \nolinkurl{fr/sga7ii/current/expose_XVII_body.tex}\par

\par\medskip\noindent\hypertarget{app-F7II-XVII-016.5}{}\textbf{F7II-XVII-016.5}\quad p. 222; designated-witness attribution.\par
\noindent Imprimé : \emph{(vide)}\par
\noindent Corrigé : {\ttfamily\footnotesize r-\allowbreak{}}\par
\noindent Source : WITNESS-F7II-XVII-016.5. \nolinkurl{fr/sga7ii/diplomatic/expose_XVII_body.tex}; \nolinkurl{fr/sga7ii/current/expose_XVII_body.tex}\par

\par\medskip\noindent\hypertarget{app-F7II-XVII-016.6}{}\textbf{F7II-XVII-016.6}\quad p. 222; designated-witness attribution.\par
\noindent Imprimé : {\ttfamily\footnotesize -r\allowbreak{}}\par
\noindent Corrigé : \emph{(vide)}\par
\noindent Source : WITNESS-F7II-XVII-016.6. \nolinkurl{fr/sga7ii/diplomatic/expose_XVII_body.tex}; \nolinkurl{fr/sga7ii/current/expose_XVII_body.tex}\par

\par\medskip\noindent\hypertarget{app-F7II-XVII-016.7}{}\textbf{F7II-XVII-016.7}\quad p. 222; designated-witness attribution.\par
\noindent Imprimé : {\ttfamily\footnotesize  \textbackslash{}lambda\allowbreak{}}\par
\noindent Corrigé : {\ttfamily\footnotesize \textbackslash{}Lambda\allowbreak{}}\par
\noindent Source : WITNESS-F7II-XVII-016.7. \nolinkurl{fr/sga7ii/diplomatic/expose_XVII_body.tex}; \nolinkurl{fr/sga7ii/current/expose_XVII_body.tex}\par

\par\medskip\noindent\hypertarget{app-F7II-XVII-016.8}{}\textbf{F7II-XVII-016.8}\quad p. 222; printed misprint.\par
\noindent Imprimé : {\ttfamily\footnotesize \{n+\allowbreak{}}\par
\noindent Corrigé : \emph{(vide)}\par
\noindent Source : 7II:351. \nolinkurl{fixes/MISPRINTS_7II.json#351}\par

\par\medskip\noindent\hypertarget{app-F7II-XVII-016.9}{}\textbf{F7II-XVII-016.9}\quad p. 222; printed misprint.\par
\noindent Imprimé : {\ttfamily\footnotesize \}\allowbreak{}}\par
\noindent Corrigé : \emph{(vide)}\par
\noindent Source : 7II:351. \nolinkurl{fixes/MISPRINTS_7II.json#351}\par

\par\medskip\noindent\hypertarget{app-F7II-XVII-016.10}{}\textbf{F7II-XVII-016.10}\quad p. 223; designated-witness attribution.\par
\noindent Imprimé : \emph{(1 espace(s), 0 retour(s) de ligne dans le témoin TeX)}\par
\noindent Corrigé : {\ttfamily\footnotesize  \textbackslash{}left\allowbreak{}}\par
\noindent Source : WITNESS-F7II-XVII-016.10. \nolinkurl{fr/sga7ii/diplomatic/expose_XVII_body.tex}; \nolinkurl{fr/sga7ii/current/expose_XVII_body.tex}\par

\par\medskip\noindent\hypertarget{app-F7II-XVII-016.11}{}\textbf{F7II-XVII-016.11}\quad p. 223; designated-witness attribution.\par
\noindent Imprimé : {\ttfamily\footnotesize (\allowbreak{}}\par
\noindent Corrigé : {\ttfamily\footnotesize \textbackslash{}left[\allowbreak{}}\par
\noindent Source : WITNESS-F7II-XVII-016.11. \nolinkurl{fr/sga7ii/diplomatic/expose_XVII_body.tex}; \nolinkurl{fr/sga7ii/current/expose_XVII_body.tex}\par

\par\medskip\noindent\hypertarget{app-F7II-XVII-016.12}{}\textbf{F7II-XVII-016.12}\quad p. 223; designated-witness attribution.\par
\noindent Imprimé : \emph{(vide)}\par
\noindent Corrigé : {\ttfamily\footnotesize \{i=\allowbreak{}}\par
\noindent Source : WITNESS-F7II-XVII-016.12. \nolinkurl{fr/sga7ii/diplomatic/expose_XVII_body.tex}; \nolinkurl{fr/sga7ii/current/expose_XVII_body.tex}\par

\par\medskip\noindent\hypertarget{app-F7II-XVII-016.13}{}\textbf{F7II-XVII-016.13}\quad p. 223; designated-witness attribution.\par
\noindent Imprimé : \emph{(vide)}\par
\noindent Corrigé : {\ttfamily\footnotesize \}\allowbreak{}}\par
\noindent Source : WITNESS-F7II-XVII-016.13. \nolinkurl{fr/sga7ii/diplomatic/expose_XVII_body.tex}; \nolinkurl{fr/sga7ii/current/expose_XVII_body.tex}\par

\par\medskip\noindent\hypertarget{app-F7II-XVII-016.14}{}\textbf{F7II-XVII-016.14}\quad p. 223; designated-witness attribution.\par
\noindent Imprimé : \emph{(vide)}\par
\noindent Corrigé : {\ttfamily\footnotesize r-\allowbreak{}}\par
\noindent Source : WITNESS-F7II-XVII-016.14. \nolinkurl{fr/sga7ii/diplomatic/expose_XVII_body.tex}; \nolinkurl{fr/sga7ii/current/expose_XVII_body.tex}\par

\par\medskip\noindent\hypertarget{app-F7II-XVII-016.15}{}\textbf{F7II-XVII-016.15}\quad p. 223; designated-witness attribution.\par
\noindent Imprimé : {\ttfamily\footnotesize -r\allowbreak{}}\par
\noindent Corrigé : \emph{(vide)}\par
\noindent Source : WITNESS-F7II-XVII-016.15. \nolinkurl{fr/sga7ii/diplomatic/expose_XVII_body.tex}; \nolinkurl{fr/sga7ii/current/expose_XVII_body.tex}\par

\par\medskip\noindent\hypertarget{app-F7II-XVII-016.16}{}\textbf{F7II-XVII-016.16}\quad p. 223; designated-witness attribution.\par
\noindent Imprimé : \emph{(vide)}\par
\noindent Corrigé : {\ttfamily\footnotesize \textbackslash{}right]\textbackslash{}right\allowbreak{}}\par
\noindent Source : WITNESS-F7II-XVII-016.16. \nolinkurl{fr/sga7ii/diplomatic/expose_XVII_body.tex}; \nolinkurl{fr/sga7ii/current/expose_XVII_body.tex}\par

\par\medskip\noindent\hypertarget{app-F7II-XVII-016.17}{}\textbf{F7II-XVII-016.17}\quad p. 223; designated-witness attribution.\par
\noindent Imprimé : \emph{(vide)}\par
\noindent Corrigé : {\ttfamily\footnotesize r-\allowbreak{}}\par
\noindent Source : WITNESS-F7II-XVII-016.17. \nolinkurl{fr/sga7ii/diplomatic/expose_XVII_body.tex}; \nolinkurl{fr/sga7ii/current/expose_XVII_body.tex}\par

\par\medskip\noindent\hypertarget{app-F7II-XVII-016.18}{}\textbf{F7II-XVII-016.18}\quad p. 223; designated-witness attribution.\par
\noindent Imprimé : {\ttfamily\footnotesize -r\allowbreak{}}\par
\noindent Corrigé : \emph{(vide)}\par
\noindent Source : WITNESS-F7II-XVII-016.18. \nolinkurl{fr/sga7ii/diplomatic/expose_XVII_body.tex}; \nolinkurl{fr/sga7ii/current/expose_XVII_body.tex}\par

\par\medskip\noindent\hypertarget{app-F7II-XVII-016.19}{}\textbf{F7II-XVII-016.19}\quad p. 223; printed misprint.\par
\noindent Imprimé : {\ttfamily\footnotesize \textbackslash{}check\{\allowbreak{}}\par
\noindent Corrigé : \emph{(vide)}\par
\noindent Source : 7II:355. \nolinkurl{fixes/MISPRINTS_7II.json#355}\par

\par\medskip\noindent\hypertarget{app-F7II-XVII-016.20}{}\textbf{F7II-XVII-016.20}\quad p. 223; printed misprint.\par
\noindent Imprimé : {\ttfamily\footnotesize \}\allowbreak{}}\par
\noindent Corrigé : \emph{(vide)}\par
\noindent Source : 7II:355. \nolinkurl{fixes/MISPRINTS_7II.json#355}\par

\par\medskip\noindent\hypertarget{app-F7II-XVII-016.21}{}\textbf{F7II-XVII-016.21}\quad p. 223; misprint\_pass.\par
\noindent Imprimé : {\ttfamily\footnotesize coordonnés\allowbreak{}}\par
\noindent Corrigé : {\ttfamily\footnotesize coordonnées\allowbreak{}}\par
\noindent Source : 7II:356. \nolinkurl{fixes/FIXLOG_7II_CURRENT.jsonl#706}; \nolinkurl{fixes/MISPRINTS_7II.json#356}\par

\par\medskip\noindent\hypertarget{app-F7II-XVII-016.22}{}\textbf{F7II-XVII-016.22}\quad p. 223; designated-witness attribution.\par
\noindent Imprimé : {\ttfamily\footnotesize  =\allowbreak{} \textbackslash{}varphi\textbackslash{},\allowbreak{} (x\_1,\textbackslash{}dots,x\_n,\textbackslash{}lamb\allowbreak{}da\_1,\textbackslash{}dots,\textbackslash{}lambda\_\{\allowbreak{}r-n-1\})\allowbreak{} \textbackslash{}end\{equation\}\allowbreak{} \textbackslash{}begin\{equation*\}\allowbreak{}}\par
\noindent Corrigé : \emph{(vide)}\par
\noindent Source : WITNESS-F7II-XVII-016.22. \nolinkurl{fr/sga7ii/diplomatic/expose_XVII_body.tex}; \nolinkurl{fr/sga7ii/current/expose_XVII_body.tex}\par

\par\medskip\noindent\hypertarget{app-F7II-XVII-016.23}{}\textbf{F7II-XVII-016.23}\quad p. 223; printed misprint.\par
\noindent Imprimé : {\ttfamily\footnotesize n+\allowbreak{}}\par
\noindent Corrigé : \emph{(vide)}\par
\noindent Source : 7II:357. \nolinkurl{fixes/MISPRINTS_7II.json#357}\par

\par\medskip\noindent\hypertarget{app-F7II-XVII-016.24}{}\textbf{F7II-XVII-016.24}\quad p. 223; designated-witness attribution.\par
\noindent Imprimé : {\ttfamily\footnotesize *\allowbreak{}}\par
\noindent Corrigé : \emph{(vide)}\par
\noindent Source : WITNESS-F7II-XVII-016.24. \nolinkurl{fr/sga7ii/diplomatic/expose_XVII_body.tex}; \nolinkurl{fr/sga7ii/current/expose_XVII_body.tex}\par

\par\medskip\noindent\hypertarget{app-F7II-XVII-016.25}{}\textbf{F7II-XVII-016.25}\quad p. 223; misprint\_pass.\par
\noindent Imprimé : {\ttfamily\footnotesize coordonnés\allowbreak{}}\par
\noindent Corrigé : {\ttfamily\footnotesize coordonnées\allowbreak{}}\par
\noindent Source : 7II:358. \nolinkurl{fixes/FIXLOG_7II_CURRENT.jsonl#707}\par

\par\medskip\noindent\hypertarget{app-F7II-XVII-016.26}{}\textbf{F7II-XVII-016.26}\quad p. 223; printed misprint.\par
\noindent Imprimé : \emph{(vide)}\par
\noindent Corrigé : {\ttfamily\footnotesize )\allowbreak{}}\par
\noindent Source : 7II:359. \nolinkurl{fixes/MISPRINTS_7II.json#359}\par

\par\medskip\noindent\hypertarget{app-F7II-XVII-016.27}{}\textbf{F7II-XVII-016.27}\quad p. 223; designated-witness attribution.\par
\noindent Imprimé : {\ttfamily\footnotesize i\allowbreak{}}\par
\noindent Corrigé : {\ttfamily\footnotesize j\allowbreak{}}\par
\noindent Source : WITNESS-F7II-XVII-016.27. \nolinkurl{fr/sga7ii/diplomatic/expose_XVII_body.tex}; \nolinkurl{fr/sga7ii/current/expose_XVII_body.tex}\par

\par\medskip\noindent\hypertarget{app-F7II-XVII-016.28}{}\textbf{F7II-XVII-016.28}\quad p. 223; designated-witness attribution.\par
\noindent Imprimé : {\ttfamily\footnotesize i\allowbreak{}}\par
\noindent Corrigé : {\ttfamily\footnotesize j\allowbreak{}}\par
\noindent Source : WITNESS-F7II-XVII-016.28. \nolinkurl{fr/sga7ii/diplomatic/expose_XVII_body.tex}; \nolinkurl{fr/sga7ii/current/expose_XVII_body.tex}\par

\par\medskip\noindent\hypertarget{app-F7II-XVII-016.29}{}\textbf{F7II-XVII-016.29}\quad p. 223; designated-witness attribution.\par
\noindent Imprimé : {\ttfamily\footnotesize i\allowbreak{}}\par
\noindent Corrigé : {\ttfamily\footnotesize j\allowbreak{}}\par
\noindent Source : WITNESS-F7II-XVII-016.29. \nolinkurl{fr/sga7ii/diplomatic/expose_XVII_body.tex}; \nolinkurl{fr/sga7ii/current/expose_XVII_body.tex}\par

\par\medskip\noindent\hypertarget{app-F7II-XVII-016.30}{}\textbf{F7II-XVII-016.30}\quad p. 223; designated-witness attribution.\par
\noindent Imprimé : {\ttfamily\footnotesize i\allowbreak{}}\par
\noindent Corrigé : {\ttfamily\footnotesize j\allowbreak{}}\par
\noindent Source : WITNESS-F7II-XVII-016.30. \nolinkurl{fr/sga7ii/diplomatic/expose_XVII_body.tex}; \nolinkurl{fr/sga7ii/current/expose_XVII_body.tex}\par

\par\medskip\noindent\hypertarget{app-F7II-XVII-016.31}{}\textbf{F7II-XVII-016.31}\quad p. 223; designated-witness attribution.\par
\noindent Imprimé : {\ttfamily\footnotesize i\allowbreak{}}\par
\noindent Corrigé : {\ttfamily\footnotesize j\allowbreak{}}\par
\noindent Source : WITNESS-F7II-XVII-016.31. \nolinkurl{fr/sga7ii/diplomatic/expose_XVII_body.tex}; \nolinkurl{fr/sga7ii/current/expose_XVII_body.tex}\par

\par\medskip\noindent\hypertarget{app-F7II-XVII-016.32}{}\textbf{F7II-XVII-016.32}\quad p. 223; designated-witness attribution.\par
\noindent Imprimé : {\ttfamily\footnotesize i\allowbreak{}}\par
\noindent Corrigé : {\ttfamily\footnotesize j\allowbreak{}}\par
\noindent Source : WITNESS-F7II-XVII-016.32. \nolinkurl{fr/sga7ii/diplomatic/expose_XVII_body.tex}; \nolinkurl{fr/sga7ii/current/expose_XVII_body.tex}\par

\par\medskip\noindent\hypertarget{app-F7II-XVII-016.33}{}\textbf{F7II-XVII-016.33}\quad p. 223; printed misprint.\par
\noindent Imprimé : \emph{(1 espace(s), 0 retour(s) de ligne dans le témoin TeX)}\par
\noindent Corrigé : {\ttfamily\footnotesize ,\allowbreak{}}\par
\noindent Source : 7II:360. \nolinkurl{fixes/MISPRINTS_7II.json#360}\par

\par\medskip\noindent\hypertarget{app-F7II-XVII-017.1}{}\textbf{F7II-XVII-017.1}\quad p. 224; misprint\_pass.\par
\noindent Imprimé : {\ttfamily\footnotesize Rappellant\allowbreak{}}\par
\noindent Corrigé : {\ttfamily\footnotesize Rappelant\allowbreak{}}\par
\noindent Source : 7II:361. \nolinkurl{fixes/FIXLOG_7II_CURRENT.jsonl#708}; \nolinkurl{fixes/MISPRINTS_7II.json#361}\par

\par\medskip\noindent\hypertarget{app-F7II-XVII-017.2}{}\textbf{F7II-XVII-017.2}\quad p. 224; inherited corrected-layer difference.\par
\noindent Imprimé : {\ttfamily\footnotesize i\allowbreak{}}\par
\noindent Corrigé : {\ttfamily\footnotesize j\allowbreak{}}\par
\noindent Source : ADOPT-F7II-XVII-017.2. \nolinkurl{fr/sga7ii/diplomatic/expose_XVII_body.tex}; \nolinkurl{fr/sga7ii/current/expose_XVII_body.tex}\par

\par\medskip\noindent\hypertarget{app-F7II-XVII-017.3}{}\textbf{F7II-XVII-017.3}\quad p. 224; inherited corrected-layer difference.\par
\noindent Imprimé : {\ttfamily\footnotesize i\allowbreak{}}\par
\noindent Corrigé : {\ttfamily\footnotesize j\allowbreak{}}\par
\noindent Source : ADOPT-F7II-XVII-017.3. \nolinkurl{fr/sga7ii/diplomatic/expose_XVII_body.tex}; \nolinkurl{fr/sga7ii/current/expose_XVII_body.tex}\par

\par\medskip\noindent\hypertarget{app-F7II-XVII-017.4}{}\textbf{F7II-XVII-017.4}\quad p. 224; inherited corrected-layer difference.\par
\noindent Imprimé : {\ttfamily\footnotesize i\allowbreak{}}\par
\noindent Corrigé : {\ttfamily\footnotesize j\allowbreak{}}\par
\noindent Source : ADOPT-F7II-XVII-017.4. \nolinkurl{fr/sga7ii/diplomatic/expose_XVII_body.tex}; \nolinkurl{fr/sga7ii/current/expose_XVII_body.tex}\par

\par\medskip\noindent\hypertarget{app-F7II-XVII-017.5}{}\textbf{F7II-XVII-017.5}\quad p. 224; designated-witness attribution.\par
\noindent Imprimé : {\ttfamily\footnotesize i\allowbreak{}}\par
\noindent Corrigé : {\ttfamily\footnotesize \{n+j\}\allowbreak{}}\par
\noindent Source : WITNESS-F7II-XVII-017.5. \nolinkurl{fr/sga7ii/diplomatic/expose_XVII_body.tex}; \nolinkurl{fr/sga7ii/current/expose_XVII_body.tex}\par

\par\medskip\noindent\hypertarget{app-F7II-XVII-017.6}{}\textbf{F7II-XVII-017.6}\quad p. 224; designated-witness attribution.\par
\noindent Imprimé : {\ttfamily\footnotesize i\allowbreak{}}\par
\noindent Corrigé : {\ttfamily\footnotesize j\allowbreak{}}\par
\noindent Source : WITNESS-F7II-XVII-017.6. \nolinkurl{fr/sga7ii/diplomatic/expose_XVII_body.tex}; \nolinkurl{fr/sga7ii/current/expose_XVII_body.tex}\par

\par\medskip\noindent\hypertarget{app-F7II-XVII-017.7}{}\textbf{F7II-XVII-017.7}\quad p. 224; designated-witness attribution.\par
\noindent Imprimé : {\ttfamily\footnotesize i\allowbreak{}}\par
\noindent Corrigé : {\ttfamily\footnotesize j\allowbreak{}}\par
\noindent Source : WITNESS-F7II-XVII-017.7. \nolinkurl{fr/sga7ii/diplomatic/expose_XVII_body.tex}; \nolinkurl{fr/sga7ii/current/expose_XVII_body.tex}\par

\par\medskip\noindent\hypertarget{app-F7II-XVII-017.8}{}\textbf{F7II-XVII-017.8}\quad p. 224; printed misprint.\par
\noindent Imprimé : \emph{(vide)}\par
\noindent Corrigé : {\ttfamily\footnotesize r-\allowbreak{}}\par
\noindent Source : 7II:362. \nolinkurl{fixes/MISPRINTS_7II.json#362}\par

\par\medskip\noindent\hypertarget{app-F7II-XVII-017.9}{}\textbf{F7II-XVII-017.9}\quad p. 224; printed misprint.\par
\noindent Imprimé : \emph{(vide)}\par
\noindent Corrigé : {\ttfamily\footnotesize -1\allowbreak{}}\par
\noindent Source : 7II:362. \nolinkurl{fixes/MISPRINTS_7II.json#362}\par

\par\medskip\noindent\hypertarget{app-F7II-XVII-018.1}{}\textbf{F7II-XVII-018.1}\quad p. 226; misprint\_pass.\par
\noindent Imprimé : {\ttfamily\footnotesize coordonnés\allowbreak{}}\par
\noindent Corrigé : {\ttfamily\footnotesize coordonnées\allowbreak{}}\par
\noindent Source : 7II:363. \nolinkurl{fixes/FIXLOG_7II_CURRENT.jsonl#709}; \nolinkurl{fixes/MISPRINTS_7II.json#363}\par

\par\medskip\noindent\hypertarget{app-F7II-XVII-019.1}{}\textbf{F7II-XVII-019.1}\quad p. 226; misprint\_pass.\par
\noindent Imprimé : {\ttfamily\footnotesize reduit\allowbreak{}}\par
\noindent Corrigé : {\ttfamily\footnotesize réduit\allowbreak{}}\par
\noindent Source : 7II:364. \nolinkurl{fixes/FIXLOG_7II_CURRENT.jsonl#710}; \nolinkurl{fixes/MISPRINTS_7II.json#364}\par

\par\medskip\noindent\hypertarget{app-F7II-XVII-020.1}{}\textbf{F7II-XVII-020.1}\quad p. 227; misprint\_pass.\par
\noindent Imprimé : {\ttfamily\footnotesize précisement\allowbreak{}}\par
\noindent Corrigé : {\ttfamily\footnotesize précisément\allowbreak{}}\par
\noindent Source : 7II:365. \nolinkurl{fixes/FIXLOG_7II_CURRENT.jsonl#711}; \nolinkurl{fixes/MISPRINTS_7II.json#365}\par

\par\medskip\noindent\hypertarget{app-F7II-XVII-020.2}{}\textbf{F7II-XVII-020.2}\quad p. 227; printed misprint.\par
\noindent Imprimé : \emph{(vide)}\par
\noindent Corrigé : {\ttfamily\footnotesize \textbackslash{}check\{\allowbreak{}}\par
\noindent Source : 7II:366. \nolinkurl{fixes/MISPRINTS_7II.json#366}\par

\par\medskip\noindent\hypertarget{app-F7II-XVII-020.3}{}\textbf{F7II-XVII-020.3}\quad p. 227; printed misprint.\par
\noindent Imprimé : \emph{(vide)}\par
\noindent Corrigé : {\ttfamily\footnotesize \}\allowbreak{}}\par
\noindent Source : 7II:366. \nolinkurl{fixes/MISPRINTS_7II.json#366}\par

\par\medskip\noindent\hypertarget{app-F7II-XVII-021.1}{}\textbf{F7II-XVII-021.1}\quad p. 228; misprint\_pass.\par
\noindent Imprimé : {\ttfamily\footnotesize naive\allowbreak{}}\par
\noindent Corrigé : {\ttfamily\footnotesize naïve\allowbreak{}}\par
\noindent Source : 7II:367. \nolinkurl{fixes/FIXLOG_7II_CURRENT.jsonl#712}; \nolinkurl{fixes/MISPRINTS_7II.json#367}\par

\par\medskip\noindent\hypertarget{app-F7II-XVII-022.1}{}\textbf{F7II-XVII-022.1}\quad p. 228; misprint\_pass.\par
\noindent Imprimé : {\ttfamily\footnotesize a\allowbreak{}}\par
\noindent Corrigé : {\ttfamily\footnotesize à\allowbreak{}}\par
\noindent Source : 7II:368. \nolinkurl{fixes/FIXLOG_7II_CURRENT.jsonl#713}; \nolinkurl{fixes/MISPRINTS_7II.json#368}\par

\par\medskip\noindent\hypertarget{app-F7II-XVII-022.2}{}\textbf{F7II-XVII-022.2}\quad p. 229; misprint\_pass.\par
\noindent Imprimé : \emph{(vide)}\par
\noindent Corrigé : {\ttfamily\footnotesize  de\allowbreak{}}\par
\noindent Source : 7II:369. \nolinkurl{fixes/FIXLOG_7II_CURRENT.jsonl#714}; \nolinkurl{fixes/MISPRINTS_7II.json#369}\par

\par\medskip\noindent\hypertarget{app-F7II-XVII-023.1}{}\textbf{F7II-XVII-023.1}\quad p. 229; misprint\_pass.\par
\noindent Imprimé : {\ttfamily\footnotesize coordonnés\allowbreak{}}\par
\noindent Corrigé : {\ttfamily\footnotesize coordonnées\allowbreak{}}\par
\noindent Source : 7II:370. \nolinkurl{fixes/FIXLOG_7II_CURRENT.jsonl#715}; \nolinkurl{fixes/MISPRINTS_7II.json#370}\par

\par\medskip\noindent\hypertarget{app-F7II-XVII-024.1}{}\textbf{F7II-XVII-024.1}\quad p. 233; misprint\_pass.\par
\noindent Imprimé : {\ttfamily\footnotesize déja\allowbreak{}}\par
\noindent Corrigé : {\ttfamily\footnotesize déjà\allowbreak{}}\par
\noindent Source : 7II:371. \nolinkurl{fixes/FIXLOG_7II_CURRENT.jsonl#716}\par

\par\medskip\noindent\hypertarget{app-F7II-XVII-025.1}{}\textbf{F7II-XVII-025.1}\quad p. 234; misprint\_pass.\par
\noindent Imprimé : {\ttfamily\footnotesize coordonnés\allowbreak{}}\par
\noindent Corrigé : {\ttfamily\footnotesize coordonnées\allowbreak{}}\par
\noindent Source : 7II:372. \nolinkurl{fixes/FIXLOG_7II_CURRENT.jsonl#717}\par

\par\medskip\noindent\hypertarget{app-F7II-XVII-025.2}{}\textbf{F7II-XVII-025.2}\quad p. 234; misprint\_pass.\par
\noindent Imprimé : {\ttfamily\footnotesize coordonnés\allowbreak{}}\par
\noindent Corrigé : {\ttfamily\footnotesize coordonnées\allowbreak{}}\par
\noindent Source : 7II:373. \nolinkurl{fixes/FIXLOG_7II_CURRENT.jsonl#718}\par

\par\medskip\noindent\hypertarget{app-F7II-XVII-026.1}{}\textbf{F7II-XVII-026.1}\quad p. 235; printed misprint.\par
\noindent Imprimé : \emph{(vide)}\par
\noindent Corrigé : {\ttfamily\footnotesize \textbackslash{}check\{\allowbreak{}}\par
\noindent Source : 7II:374. \nolinkurl{fixes/MISPRINTS_7II.json#374}\par

\par\medskip\noindent\hypertarget{app-F7II-XVII-026.2}{}\textbf{F7II-XVII-026.2}\quad p. 235; printed misprint.\par
\noindent Imprimé : \emph{(vide)}\par
\noindent Corrigé : {\ttfamily\footnotesize \}\allowbreak{}}\par
\noindent Source : 7II:374. \nolinkurl{fixes/MISPRINTS_7II.json#374}\par

\par\medskip\noindent\hypertarget{app-F7II-XVII-027.1}{}\textbf{F7II-XVII-027.1}\quad p. 236; designated-witness attribution.\par
\noindent Imprimé : \emph{(vide)}\par
\noindent Corrigé : {\ttfamily\footnotesize \textbackslash{}begin\{equation\}\textbackslash{}tag\allowbreak{}\{\allowbreak{}}\par
\noindent Source : WITNESS-F7II-XVII-027.1. \nolinkurl{fr/sga7ii/diplomatic/expose_XVII_body.tex}; \nolinkurl{fr/sga7ii/current/expose_XVII_body.tex}\par

\par\medskip\noindent\hypertarget{app-F7II-XVII-027.2}{}\textbf{F7II-XVII-027.2}\quad p. 236; designated-witness attribution.\par
\noindent Imprimé : {\ttfamily\footnotesize .\textbackslash{}quad\allowbreak{} l'adhérence\allowbreak{} de\allowbreak{} \$\allowbreak{}}\par
\noindent Corrigé : {\ttfamily\footnotesize \}\allowbreak{} \textbackslash{}overline\{\allowbreak{}}\par
\noindent Source : WITNESS-F7II-XVII-027.2. \nolinkurl{fr/sga7ii/diplomatic/expose_XVII_body.tex}; \nolinkurl{fr/sga7ii/current/expose_XVII_body.tex}\par

\par\medskip\noindent\hypertarget{app-F7II-XVII-027.3}{}\textbf{F7II-XVII-027.3}\quad p. 236; designated-witness attribution.\par
\noindent Imprimé : \emph{(1 espace(s), 0 retour(s) de ligne dans le témoin TeX)}\par
\noindent Corrigé : {\ttfamily\footnotesize \}\allowbreak{}}\par
\noindent Source : WITNESS-F7II-XVII-027.3. \nolinkurl{fr/sga7ii/diplomatic/expose_XVII_body.tex}; \nolinkurl{fr/sga7ii/current/expose_XVII_body.tex}\par

\par\medskip\noindent\hypertarget{app-F7II-XVII-027.4}{}\textbf{F7II-XVII-027.4}\quad p. 236; designated-witness attribution.\par
\noindent Imprimé : {\ttfamily\footnotesize \$\allowbreak{}}\par
\noindent Corrigé : {\ttfamily\footnotesize  \textbackslash{}quad\allowbreak{}}\par
\noindent Source : WITNESS-F7II-XVII-027.4. \nolinkurl{fr/sga7ii/diplomatic/expose_XVII_body.tex}; \nolinkurl{fr/sga7ii/current/expose_XVII_body.tex}\par

\par\medskip\noindent\hypertarget{app-F7II-XVII-027.5}{}\textbf{F7II-XVII-027.5}\quad p. 236; designated-witness attribution.\par
\noindent Imprimé : \emph{(0 espace(s), 2 retour(s) de ligne dans le témoin TeX)}\par
\noindent Corrigé : {\ttfamily\footnotesize  \textbackslash{}end\{equation\}\allowbreak{} }\par
\noindent Source : WITNESS-F7II-XVII-027.5. \nolinkurl{fr/sga7ii/diplomatic/expose_XVII_body.tex}; \nolinkurl{fr/sga7ii/current/expose_XVII_body.tex}\par

\par\medskip\noindent\hypertarget{app-F7II-XVII-027.6}{}\textbf{F7II-XVII-027.6}\quad p. 236; printed misprint.\par
\noindent Imprimé : {\ttfamily\footnotesize 8\allowbreak{}}\par
\noindent Corrigé : {\ttfamily\footnotesize 7\allowbreak{}}\par
\noindent Source : 7II:375. \nolinkurl{fixes/MISPRINTS_7II.json#375}\par

\par\medskip\noindent\hypertarget{app-F7II-XVII-027.7}{}\textbf{F7II-XVII-027.7}\quad p. 236; printed misprint.\par
\noindent Imprimé : {\ttfamily\footnotesize 2\allowbreak{}}\par
\noindent Corrigé : {\ttfamily\footnotesize 1\allowbreak{}}\par
\noindent Source : 7II:376. \nolinkurl{fixes/MISPRINTS_7II.json#376}\par

\par\medskip\noindent\hypertarget{app-F7II-XVII-027.8}{}\textbf{F7II-XVII-027.8}\quad p. 236; printed misprint.\par
\noindent Imprimé : {\ttfamily\footnotesize n\allowbreak{}}\par
\noindent Corrigé : {\ttfamily\footnotesize r\allowbreak{}}\par
\noindent Source : 7II:377. \nolinkurl{fixes/MISPRINTS_7II.json#377}\par

\par\medskip\noindent\hypertarget{app-F7II-XVII-027.9}{}\textbf{F7II-XVII-027.9}\quad p. 236; misprint\_pass.\par
\noindent Imprimé : {\ttfamily\footnotesize Commencons\allowbreak{}}\par
\noindent Corrigé : {\ttfamily\footnotesize Commençons\allowbreak{}}\par
\noindent Source : 7II:378. \nolinkurl{fixes/FIXLOG_7II_CURRENT.jsonl#719}; \nolinkurl{fixes/MISPRINTS_7II.json#378}\par

\par\medskip\noindent\hypertarget{app-F7II-XVII-027.10}{}\textbf{F7II-XVII-027.10}\quad p. 236; misprint\_pass.\par
\noindent Imprimé : {\ttfamily\footnotesize coordonnés\allowbreak{}}\par
\noindent Corrigé : {\ttfamily\footnotesize coordonnées\allowbreak{}}\par
\noindent Source : 7II:379. \nolinkurl{fixes/FIXLOG_7II_CURRENT.jsonl#720}; \nolinkurl{fixes/MISPRINTS_7II.json#379}\par

\par\medskip\noindent\hypertarget{app-F7II-XVII-027.11}{}\textbf{F7II-XVII-027.11}\quad p. 236; misprint\_pass.\par
\noindent Imprimé : {\ttfamily\footnotesize facon\allowbreak{}}\par
\noindent Corrigé : {\ttfamily\footnotesize façon\allowbreak{}}\par
\noindent Source : 7II:380. \nolinkurl{fixes/FIXLOG_7II_CURRENT.jsonl#721}; \nolinkurl{fixes/MISPRINTS_7II.json#380}\par

\par\medskip\noindent\hypertarget{app-F7II-XVII-027.12}{}\textbf{F7II-XVII-027.12}\quad p. 237; printed misprint.\par
\noindent Imprimé : {\ttfamily\footnotesize coordonnés\allowbreak{}}\par
\noindent Corrigé : {\ttfamily\footnotesize coordonnées\allowbreak{}}\par
\noindent Source : 7II:381. \nolinkurl{fixes/MISPRINTS_7II.json#381}\par

\par\medskip\noindent\hypertarget{app-F7II-XVII-028.1}{}\textbf{F7II-XVII-028.1}\quad p. 237; designated-witness attribution.\par
\noindent Imprimé : {\ttfamily\footnotesize (\allowbreak{}}\par
\noindent Corrigé : {\ttfamily\footnotesize \textbackslash{}begin\{equation\}\textbackslash{}tag\allowbreak{}\{\allowbreak{}}\par
\noindent Source : WITNESS-F7II-XVII-028.1. \nolinkurl{fr/sga7ii/diplomatic/expose_XVII_body.tex}; \nolinkurl{fr/sga7ii/current/expose_XVII_body.tex}\par

\par\medskip\noindent\hypertarget{app-F7II-XVII-028.2}{}\textbf{F7II-XVII-028.2}\quad p. 237; designated-witness attribution.\par
\noindent Imprimé : {\ttfamily\footnotesize )\textbackslash{}quad\allowbreak{} \$\allowbreak{}}\par
\noindent Corrigé : {\ttfamily\footnotesize \}\allowbreak{} }\par
\noindent Source : WITNESS-F7II-XVII-028.2. \nolinkurl{fr/sga7ii/diplomatic/expose_XVII_body.tex}; \nolinkurl{fr/sga7ii/current/expose_XVII_body.tex}\par

\par\medskip\noindent\hypertarget{app-F7II-XVII-028.3}{}\textbf{F7II-XVII-028.3}\quad p. 237; designated-witness attribution.\par
\noindent Imprimé : {\ttfamily\footnotesize \$\allowbreak{} }\par
\noindent Corrigé : {\ttfamily\footnotesize  \textbackslash{}quad\textbackslash{}mbox\{\allowbreak{}}\par
\noindent Source : WITNESS-F7II-XVII-028.3. \nolinkurl{fr/sga7ii/diplomatic/expose_XVII_body.tex}; \nolinkurl{fr/sga7ii/current/expose_XVII_body.tex}\par

\par\medskip\noindent\hypertarget{app-F7II-XVII-028.4}{}\textbf{F7II-XVII-028.4}\quad p. 237; designated-witness attribution.\par
\noindent Imprimé : {\ttfamily\footnotesize \$\allowbreak{}}\par
\noindent Corrigé : {\ttfamily\footnotesize \}\allowbreak{} }\par
\noindent Source : WITNESS-F7II-XVII-028.4. \nolinkurl{fr/sga7ii/diplomatic/expose_XVII_body.tex}; \nolinkurl{fr/sga7ii/current/expose_XVII_body.tex}\par

\par\medskip\noindent\hypertarget{app-F7II-XVII-028.5}{}\textbf{F7II-XVII-028.5}\quad p. 237; designated-witness attribution.\par
\noindent Imprimé : {\ttfamily\footnotesize \$\allowbreak{} }\par
\noindent Corrigé : \emph{(vide)}\par
\noindent Source : WITNESS-F7II-XVII-028.5. \nolinkurl{fr/sga7ii/diplomatic/expose_XVII_body.tex}; \nolinkurl{fr/sga7ii/current/expose_XVII_body.tex}\par

\par\medskip\noindent\hypertarget{app-F7II-XVII-028.6}{}\textbf{F7II-XVII-028.6}\quad p. 237; designated-witness attribution.\par
\noindent Imprimé : \emph{(1 espace(s), 0 retour(s) de ligne dans le témoin TeX)}\par
\noindent Corrigé : {\ttfamily\footnotesize \textbackslash{}\allowbreak{} \textbackslash{}mbox\{\allowbreak{}}\par
\noindent Source : WITNESS-F7II-XVII-028.6. \nolinkurl{fr/sga7ii/diplomatic/expose_XVII_body.tex}; \nolinkurl{fr/sga7ii/current/expose_XVII_body.tex}\par

\par\medskip\noindent\hypertarget{app-F7II-XVII-028.7}{}\textbf{F7II-XVII-028.7}\quad p. 237; designated-witness attribution.\par
\noindent Imprimé : {\ttfamily\footnotesize \$\allowbreak{}}\par
\noindent Corrigé : {\ttfamily\footnotesize \}\allowbreak{} }\par
\noindent Source : WITNESS-F7II-XVII-028.7. \nolinkurl{fr/sga7ii/diplomatic/expose_XVII_body.tex}; \nolinkurl{fr/sga7ii/current/expose_XVII_body.tex}\par

\par\medskip\noindent\hypertarget{app-F7II-XVII-028.8}{}\textbf{F7II-XVII-028.8}\quad p. 237; designated-witness attribution.\par
\noindent Imprimé : {\ttfamily\footnotesize \$\allowbreak{}}\par
\noindent Corrigé : {\ttfamily\footnotesize  \textbackslash{}mbox\{\allowbreak{}}\par
\noindent Source : WITNESS-F7II-XVII-028.8. \nolinkurl{fr/sga7ii/diplomatic/expose_XVII_body.tex}; \nolinkurl{fr/sga7ii/current/expose_XVII_body.tex}\par

\par\medskip\noindent\hypertarget{app-F7II-XVII-028.9}{}\textbf{F7II-XVII-028.9}\quad p. 237; designated-witness attribution.\par
\noindent Imprimé : {\ttfamily\footnotesize \$\allowbreak{}}\par
\noindent Corrigé : {\ttfamily\footnotesize \}\allowbreak{} }\par
\noindent Source : WITNESS-F7II-XVII-028.9. \nolinkurl{fr/sga7ii/diplomatic/expose_XVII_body.tex}; \nolinkurl{fr/sga7ii/current/expose_XVII_body.tex}\par

\par\medskip\noindent\hypertarget{app-F7II-XVII-028.10}{}\textbf{F7II-XVII-028.10}\quad p. 237; designated-witness attribution.\par
\noindent Imprimé : {\ttfamily\footnotesize \$\allowbreak{}}\par
\noindent Corrigé : {\ttfamily\footnotesize  \textbackslash{}mbox\{\allowbreak{}}\par
\noindent Source : WITNESS-F7II-XVII-028.10. \nolinkurl{fr/sga7ii/diplomatic/expose_XVII_body.tex}; \nolinkurl{fr/sga7ii/current/expose_XVII_body.tex}\par

\par\medskip\noindent\hypertarget{app-F7II-XVII-028.11}{}\textbf{F7II-XVII-028.11}\quad p. 237; designated-witness attribution.\par
\noindent Imprimé : \emph{(vide)}\par
\noindent Corrigé : {\ttfamily\footnotesize \}\allowbreak{}}\par
\noindent Source : WITNESS-F7II-XVII-028.11. \nolinkurl{fr/sga7ii/diplomatic/expose_XVII_body.tex}; \nolinkurl{fr/sga7ii/current/expose_XVII_body.tex}\par

\par\medskip\noindent\hypertarget{app-F7II-XVII-028.12}{}\textbf{F7II-XVII-028.12}\quad p. 237; designated-witness attribution.\par
\noindent Imprimé : \emph{(vide)}\par
\noindent Corrigé : {\ttfamily\footnotesize  \textbackslash{}end\{equation\}\allowbreak{}}\par
\noindent Source : WITNESS-F7II-XVII-028.12. \nolinkurl{fr/sga7ii/diplomatic/expose_XVII_body.tex}; \nolinkurl{fr/sga7ii/current/expose_XVII_body.tex}\par

\par\medskip\noindent\hypertarget{app-F7II-XVII-029.1}{}\textbf{F7II-XVII-029.1}\quad p. 237; designated-witness attribution.\par
\noindent Imprimé : {\ttfamily\footnotesize (\allowbreak{}}\par
\noindent Corrigé : {\ttfamily\footnotesize \textbackslash{}begin\{equation\}\textbackslash{}tag\allowbreak{}\{\allowbreak{}}\par
\noindent Source : WITNESS-F7II-XVII-029.1. \nolinkurl{fr/sga7ii/diplomatic/expose_XVII_body.tex}; \nolinkurl{fr/sga7ii/current/expose_XVII_body.tex}\par

\par\medskip\noindent\hypertarget{app-F7II-XVII-029.2}{}\textbf{F7II-XVII-029.2}\quad p. 237; designated-witness attribution.\par
\noindent Imprimé : {\ttfamily\footnotesize )\textbackslash{}quad\allowbreak{} \$\allowbreak{}}\par
\noindent Corrigé : {\ttfamily\footnotesize \}\allowbreak{} \textbackslash{}begin\{aligned\}\allowbreak{} }\par
\noindent Source : WITNESS-F7II-XVII-029.2. \nolinkurl{fr/sga7ii/diplomatic/expose_XVII_body.tex}; \nolinkurl{fr/sga7ii/current/expose_XVII_body.tex}\par

\par\medskip\noindent\hypertarget{app-F7II-XVII-029.3}{}\textbf{F7II-XVII-029.3}\quad p. 237; designated-witness attribution.\par
\noindent Imprimé : \emph{(1 espace(s), 0 retour(s) de ligne dans le témoin TeX)}\par
\noindent Corrigé : {\ttfamily\footnotesize  \&\allowbreak{}}\par
\noindent Source : WITNESS-F7II-XVII-029.3. \nolinkurl{fr/sga7ii/diplomatic/expose_XVII_body.tex}; \nolinkurl{fr/sga7ii/current/expose_XVII_body.tex}\par

\par\medskip\noindent\hypertarget{app-F7II-XVII-029.4}{}\textbf{F7II-XVII-029.4}\quad p. 237; designated-witness attribution.\par
\noindent Imprimé : {\ttfamily\footnotesize \$\allowbreak{} }\par
\noindent Corrigé : \emph{(vide)}\par
\noindent Source : WITNESS-F7II-XVII-029.4. \nolinkurl{fr/sga7ii/diplomatic/expose_XVII_body.tex}; \nolinkurl{fr/sga7ii/current/expose_XVII_body.tex}\par

\par\medskip\noindent\hypertarget{app-F7II-XVII-029.5}{}\textbf{F7II-XVII-029.5}\quad p. 237; designated-witness attribution.\par
\noindent Imprimé : \emph{(1 espace(s), 0 retour(s) de ligne dans le témoin TeX)}\par
\noindent Corrigé : {\ttfamily\footnotesize \textbackslash{}\allowbreak{} \textbackslash{}text\{\allowbreak{}}\par
\noindent Source : WITNESS-F7II-XVII-029.5. \nolinkurl{fr/sga7ii/diplomatic/expose_XVII_body.tex}; \nolinkurl{fr/sga7ii/current/expose_XVII_body.tex}\par

\par\medskip\noindent\hypertarget{app-F7II-XVII-029.6}{}\textbf{F7II-XVII-029.6}\quad p. 237; designated-witness attribution.\par
\noindent Imprimé : {\ttfamily\footnotesize \$\allowbreak{}}\par
\noindent Corrigé : {\ttfamily\footnotesize \}\allowbreak{}}\par
\noindent Source : WITNESS-F7II-XVII-029.6. \nolinkurl{fr/sga7ii/diplomatic/expose_XVII_body.tex}; \nolinkurl{fr/sga7ii/current/expose_XVII_body.tex}\par

\par\medskip\noindent\hypertarget{app-F7II-XVII-029.7}{}\textbf{F7II-XVII-029.7}\quad p. 237; designated-witness attribution.\par
\noindent Imprimé : {\ttfamily\footnotesize \$\allowbreak{}}\par
\noindent Corrigé : {\ttfamily\footnotesize  \textbackslash{}text\{\allowbreak{}}\par
\noindent Source : WITNESS-F7II-XVII-029.7. \nolinkurl{fr/sga7ii/diplomatic/expose_XVII_body.tex}; \nolinkurl{fr/sga7ii/current/expose_XVII_body.tex}\par

\par\medskip\noindent\hypertarget{app-F7II-XVII-029.8}{}\textbf{F7II-XVII-029.8}\quad p. 237; designated-witness attribution.\par
\noindent Imprimé : {\ttfamily\footnotesize \$\allowbreak{}}\par
\noindent Corrigé : {\ttfamily\footnotesize \}\allowbreak{}}\par
\noindent Source : WITNESS-F7II-XVII-029.8. \nolinkurl{fr/sga7ii/diplomatic/expose_XVII_body.tex}; \nolinkurl{fr/sga7ii/current/expose_XVII_body.tex}\par

\par\medskip\noindent\hypertarget{app-F7II-XVII-029.9}{}\textbf{F7II-XVII-029.9}\quad p. 237; designated-witness attribution.\par
\noindent Imprimé : {\ttfamily\footnotesize \$\allowbreak{}}\par
\noindent Corrigé : {\ttfamily\footnotesize \textbackslash{}text\{\allowbreak{}}\par
\noindent Source : WITNESS-F7II-XVII-029.9. \nolinkurl{fr/sga7ii/diplomatic/expose_XVII_body.tex}; \nolinkurl{fr/sga7ii/current/expose_XVII_body.tex}\par

\par\medskip\noindent\hypertarget{app-F7II-XVII-029.10}{}\textbf{F7II-XVII-029.10}\quad p. 237; designated-witness attribution.\par
\noindent Imprimé : {\ttfamily\footnotesize \$\allowbreak{}}\par
\noindent Corrigé : {\ttfamily\footnotesize \}\allowbreak{}}\par
\noindent Source : WITNESS-F7II-XVII-029.10. \nolinkurl{fr/sga7ii/diplomatic/expose_XVII_body.tex}; \nolinkurl{fr/sga7ii/current/expose_XVII_body.tex}\par

\par\medskip\noindent\hypertarget{app-F7II-XVII-029.11}{}\textbf{F7II-XVII-029.11}\quad p. 237; designated-witness attribution.\par
\noindent Imprimé : {\ttfamily\footnotesize \$\allowbreak{}}\par
\noindent Corrigé : {\ttfamily\footnotesize \textbackslash{}\}\textbackslash{}\textbackslash{}\allowbreak{} \&=\allowbreak{}}\par
\noindent Source : WITNESS-F7II-XVII-029.11. \nolinkurl{fr/sga7ii/diplomatic/expose_XVII_body.tex}; \nolinkurl{fr/sga7ii/current/expose_XVII_body.tex}\par

\par\medskip\noindent\hypertarget{app-F7II-XVII-029.12}{}\textbf{F7II-XVII-029.12}\quad p. 237; designated-witness attribution.\par
\noindent Imprimé : {\ttfamily\footnotesize \$\textbackslash{}\}\$\allowbreak{} \textbackslash{}\textbackslash{}\allowbreak{} \$\allowbreak{}}\par
\noindent Corrigé : \emph{(vide)}\par
\noindent Source : WITNESS-F7II-XVII-029.12. \nolinkurl{fr/sga7ii/diplomatic/expose_XVII_body.tex}; \nolinkurl{fr/sga7ii/current/expose_XVII_body.tex}\par

\par\medskip\noindent\hypertarget{app-F7II-XVII-029.13}{}\textbf{F7II-XVII-029.13}\quad p. 237; designated-witness attribution.\par
\noindent Imprimé : {\ttfamily\footnotesize \$\allowbreak{} et\allowbreak{} \$\allowbreak{}}\par
\noindent Corrigé : {\ttfamily\footnotesize ,\textbackslash{}\allowbreak{} }\par
\noindent Source : WITNESS-F7II-XVII-029.13. \nolinkurl{fr/sga7ii/diplomatic/expose_XVII_body.tex}; \nolinkurl{fr/sga7ii/current/expose_XVII_body.tex}\par

\par\medskip\noindent\hypertarget{app-F7II-XVII-029.14}{}\textbf{F7II-XVII-029.14}\quad p. 237; designated-witness attribution.\par
\noindent Imprimé : {\ttfamily\footnotesize \textbackslash{}\allowbreak{} }\par
\noindent Corrigé : \emph{(vide)}\par
\noindent Source : WITNESS-F7II-XVII-029.14. \nolinkurl{fr/sga7ii/diplomatic/expose_XVII_body.tex}; \nolinkurl{fr/sga7ii/current/expose_XVII_body.tex}\par

\par\medskip\noindent\hypertarget{app-F7II-XVII-029.15}{}\textbf{F7II-XVII-029.15}\quad p. 237; designated-witness attribution.\par
\noindent Imprimé : {\ttfamily\footnotesize \$\allowbreak{} }\par
\noindent Corrigé : \emph{(vide)}\par
\noindent Source : WITNESS-F7II-XVII-029.15. \nolinkurl{fr/sga7ii/diplomatic/expose_XVII_body.tex}; \nolinkurl{fr/sga7ii/current/expose_XVII_body.tex}\par

\par\medskip\noindent\hypertarget{app-F7II-XVII-029.16}{}\textbf{F7II-XVII-029.16}\quad p. 237; designated-witness attribution.\par
\noindent Imprimé : \emph{(vide)}\par
\noindent Corrigé : {\ttfamily\footnotesize  \textbackslash{}end\{aligned\}\allowbreak{} \textbackslash{}end\{equation\}\allowbreak{}}\par
\noindent Source : WITNESS-F7II-XVII-029.16. \nolinkurl{fr/sga7ii/diplomatic/expose_XVII_body.tex}; \nolinkurl{fr/sga7ii/current/expose_XVII_body.tex}\par

\par\medskip\noindent\hypertarget{app-F7II-XVII-030.1}{}\textbf{F7II-XVII-030.1}\quad p. 238; misprint\_pass.\par
\noindent Imprimé : {\ttfamily\footnotesize considérés\allowbreak{}}\par
\noindent Corrigé : {\ttfamily\footnotesize considérées\allowbreak{}}\par
\noindent Source : 7II:384. \nolinkurl{fixes/FIXLOG_7II_CURRENT.jsonl#723}; \nolinkurl{fixes/MISPRINTS_7II.json#384}\par

\par\medskip\noindent\hypertarget{app-F7II-XVII-030.2}{}\textbf{F7II-XVII-030.2}\quad p. 238; printed misprint.\par
\noindent Imprimé : {\ttfamily\footnotesize d\allowbreak{}}\par
\noindent Corrigé : {\ttfamily\footnotesize 1\allowbreak{}}\par
\noindent Source : 7II:385. \nolinkurl{fixes/MISPRINTS_7II.json#385}\par

\par\medskip\noindent\hypertarget{app-F7II-XVII-031.1}{}\textbf{F7II-XVII-031.1}\quad p. 239; printed misprint.\par
\noindent Imprimé : {\ttfamily\footnotesize coordonnés\allowbreak{}}\par
\noindent Corrigé : {\ttfamily\footnotesize coordonnées\allowbreak{}}\par
\noindent Source : 7II:386. \nolinkurl{fixes/MISPRINTS_7II.json#386}\par

\par\medskip\noindent\hypertarget{app-F7II-XVII-031.2}{}\textbf{F7II-XVII-031.2}\quad p. 239; misprint\_pass.\par
\noindent Imprimé : {\ttfamily\footnotesize coéfficients\allowbreak{}}\par
\noindent Corrigé : {\ttfamily\footnotesize coefficients\allowbreak{}}\par
\noindent Source : 7II:387. \nolinkurl{fixes/FIXLOG_7II_CURRENT.jsonl#725}; \nolinkurl{fixes/MISPRINTS_7II.json#387}\par

\par\medskip\noindent\hypertarget{app-F7II-XVII-032.1}{}\textbf{F7II-XVII-032.1}\quad p. 239; misprint\_in\_print.\par
\noindent Imprimé : {\ttfamily\footnotesize indépendentes\allowbreak{}}\par
\noindent Corrigé : {\ttfamily\footnotesize indépendantes\allowbreak{}}\par
\noindent Source : FIX-7II-D-0237. \nolinkurl{fixes/FIXLOG_7II.jsonl#237}\par

\par\medskip\noindent\hypertarget{app-F7II-XVII-032.2}{}\textbf{F7II-XVII-032.2}\quad p. 239; misprint\_pass.\par
\noindent Imprimé : {\ttfamily\footnotesize coordonnés\allowbreak{}}\par
\noindent Corrigé : {\ttfamily\footnotesize coordonnées\allowbreak{}}\par
\noindent Source : 7II:389. \nolinkurl{fixes/FIXLOG_7II_CURRENT.jsonl#726}; \nolinkurl{fixes/MISPRINTS_7II.json#389}\par

\par\medskip\noindent\hypertarget{app-F7II-XVII-032.3}{}\textbf{F7II-XVII-032.3}\quad p. 239; printed misprint.\par
\noindent Imprimé : {\ttfamily\footnotesize 1\allowbreak{}}\par
\noindent Corrigé : {\ttfamily\footnotesize 2\allowbreak{}}\par
\noindent Source : 7II:390. \nolinkurl{fixes/MISPRINTS_7II.json#390}\par

\par\medskip\noindent\hypertarget{app-F7II-XVII-032.4}{}\textbf{F7II-XVII-032.4}\quad p. 240; designated-witness attribution.\par
\noindent Imprimé : \emph{(0 espace(s), 1 retour(s) de ligne dans le témoin TeX)}\par
\noindent Corrigé : \emph{(0 espace(s), 2 retour(s) de ligne dans le témoin TeX)}\par
\noindent Source : WITNESS-F7II-XVII-032.4. \nolinkurl{fr/sga7ii/diplomatic/expose_XVII_body.tex}; \nolinkurl{fr/sga7ii/current/expose_XVII_body.tex}\par

\par\medskip\noindent\hypertarget{app-F7II-XVII-032.5}{}\textbf{F7II-XVII-032.5}\quad p. 240; printed misprint.\par
\noindent Imprimé : {\ttfamily\footnotesize ij\allowbreak{}}\par
\noindent Corrigé : {\ttfamily\footnotesize 1j\allowbreak{}}\par
\noindent Source : 7II:391. \nolinkurl{fixes/MISPRINTS_7II.json#391}\par

\par\medskip\noindent\hypertarget{app-F7II-XVII-032.6}{}\textbf{F7II-XVII-032.6}\quad p. 240; printed misprint.\par
\noindent Imprimé : {\ttfamily\footnotesize coordonnés\allowbreak{}}\par
\noindent Corrigé : {\ttfamily\footnotesize coordonnées\allowbreak{}}\par
\noindent Source : 7II:392. \nolinkurl{fixes/MISPRINTS_7II.json#392}\par

\par\medskip\noindent\hypertarget{app-F7II-XVII-033.1}{}\textbf{F7II-XVII-033.1}\quad p. 241; misprint\_pass.\par
\noindent Imprimé : {\ttfamily\footnotesize coordonnés\allowbreak{}}\par
\noindent Corrigé : {\ttfamily\footnotesize coordonnées\allowbreak{}}\par
\noindent Source : 7II:393. \nolinkurl{fixes/FIXLOG_7II_CURRENT.jsonl#728}\par

\par\medskip\noindent\hypertarget{app-F7II-XVII-034.1}{}\textbf{F7II-XVII-034.1}\quad p. 241; printed misprint.\par
\noindent Imprimé : {\ttfamily\footnotesize \{\textbackslash{}mathbb\{P\}\}\allowbreak{}}\par
\noindent Corrigé : {\ttfamily\footnotesize X\allowbreak{}}\par
\noindent Source : 7II:394. \nolinkurl{fixes/MISPRINTS_7II.json#394}\par

\par\medskip\noindent\hypertarget{app-F7II-XVII-035.1}{}\textbf{F7II-XVII-035.1}\quad p. 242; misprint\_pass.\par
\noindent Imprimé : {\ttfamily\footnotesize noté\allowbreak{}}\par
\noindent Corrigé : {\ttfamily\footnotesize notée\allowbreak{}}\par
\noindent Source : 7II:395. \nolinkurl{fixes/FIXLOG_7II_CURRENT.jsonl#729}; \nolinkurl{fixes/MISPRINTS_7II.json#395}\par

\par\medskip\noindent\hypertarget{app-F7II-XVII-035.2}{}\textbf{F7II-XVII-035.2}\quad p. 242; misprint\_pass.\par
\noindent Imprimé : {\ttfamily\footnotesize aucun\allowbreak{}}\par
\noindent Corrigé : {\ttfamily\footnotesize aucune\allowbreak{}}\par
\noindent Source : 7II:396. \nolinkurl{fixes/FIXLOG_7II_CURRENT.jsonl#730}; \nolinkurl{fixes/MISPRINTS_7II.json#396}\par

\par\medskip\noindent\hypertarget{app-F7II-XVII-036.1}{}\textbf{F7II-XVII-036.1}\quad p. 243; inherited corrected-layer difference.\par
\noindent Imprimé : \emph{(0 espace(s), 1 retour(s) de ligne dans le témoin TeX)}\par
\noindent Corrigé : \emph{(0 espace(s), 3 retour(s) de ligne dans le témoin TeX)}\par
\noindent Source : ADOPT-F7II-XVII-036.1. \nolinkurl{fr/sga7ii/diplomatic/expose_XVII_body.tex}; \nolinkurl{fr/sga7ii/current/expose_XVII_body.tex}\par

\par\medskip\noindent\hypertarget{app-F7II-XVII-036.2}{}\textbf{F7II-XVII-036.2}\quad p. 243; printed misprint.\par
\noindent Imprimé : \emph{(vide)}\par
\noindent Corrigé : {\ttfamily\footnotesize \textbackslash{}check\{\allowbreak{}}\par
\noindent Source : 7II:397. \nolinkurl{fixes/MISPRINTS_7II.json#397}\par

\par\medskip\noindent\hypertarget{app-F7II-XVII-036.3}{}\textbf{F7II-XVII-036.3}\quad p. 243; printed misprint.\par
\noindent Imprimé : \emph{(vide)}\par
\noindent Corrigé : {\ttfamily\footnotesize \}\allowbreak{}}\par
\noindent Source : 7II:397. \nolinkurl{fixes/MISPRINTS_7II.json#397}\par

\par\medskip\noindent\hypertarget{app-F7II-XVII-037.1}{}\textbf{F7II-XVII-037.1}\quad p. 243; misprint\_pass.\par
\noindent Imprimé : \emph{(vide)}\par
\noindent Corrigé : {\ttfamily\footnotesize (1)\allowbreak{}}\par
\noindent Source : 7II:398. \nolinkurl{fixes/FIXLOG_7II_CURRENT.jsonl#731}; \nolinkurl{fixes/MISPRINTS_7II.json#398}\par

\par\medskip\noindent\hypertarget{app-F7II-XVII-038.1}{}\textbf{F7II-XVII-038.1}\quad p. 244; misprint\_pass.\par
\noindent Imprimé : {\ttfamily\footnotesize moyenant\allowbreak{}}\par
\noindent Corrigé : {\ttfamily\footnotesize moyennant\allowbreak{}}\par
\noindent Source : 7II:399. \nolinkurl{fixes/FIXLOG_7II_CURRENT.jsonl#732}\par

\par\medskip\noindent\hypertarget{app-F7II-XVII-039.1}{}\textbf{F7II-XVII-039.1}\quad p. 244; designated-witness attribution.\par
\noindent Imprimé : \emph{(vide)}\par
\noindent Corrigé : {\ttfamily\footnotesize \$'\$\allowbreak{}}\par
\noindent Source : WITNESS-F7II-XVII-039.1. \nolinkurl{fr/sga7ii/diplomatic/expose_XVII_body.tex}; \nolinkurl{fr/sga7ii/current/expose_XVII_body.tex}\par

\par\medskip\noindent\hypertarget{app-F7II-XVII-040.1}{}\textbf{F7II-XVII-040.1}\quad p. 244; misprint\_pass.\par
\noindent Imprimé : {\ttfamily\footnotesize deuxieme\allowbreak{}}\par
\noindent Corrigé : {\ttfamily\footnotesize deuxième\allowbreak{}}\par
\noindent Source : 7II:401. \nolinkurl{fixes/FIXLOG_7II_CURRENT.jsonl#733}\par

\par\medskip\noindent\hypertarget{app-F7II-XVII-041.1}{}\textbf{F7II-XVII-041.1}\quad p. 246; misprint\_pass.\par
\noindent Imprimé : {\ttfamily\footnotesize se\allowbreak{}}\par
\noindent Corrigé : {\ttfamily\footnotesize ce\allowbreak{}}\par
\noindent Source : 7II:402. \nolinkurl{fixes/FIXLOG_7II_CURRENT.jsonl#734}; \nolinkurl{fixes/MISPRINTS_7II.json#402}\par

\par\medskip\noindent\hypertarget{app-F7II-XVII-042.1}{}\textbf{F7II-XVII-042.1}\quad p. 247; printed misprint.\par
\noindent Imprimé : {\ttfamily\footnotesize j\allowbreak{}}\par
\noindent Corrigé : {\ttfamily\footnotesize i\allowbreak{}}\par
\noindent Source : 7II:403. \nolinkurl{fixes/MISPRINTS_7II.json#403}\par

\par\medskip\noindent\hypertarget{app-F7II-XVII-042.2}{}\textbf{F7II-XVII-042.2}\quad p. 247; printed misprint.\par
\noindent Imprimé : {\ttfamily\footnotesize j\allowbreak{}}\par
\noindent Corrigé : {\ttfamily\footnotesize a\allowbreak{}}\par
\noindent Source : 7II:403. \nolinkurl{fixes/MISPRINTS_7II.json#403}\par

\par\medskip\noindent\hypertarget{app-F7II-XVII-042.3}{}\textbf{F7II-XVII-042.3}\quad p. 247; printed misprint.\par
\noindent Imprimé : {\ttfamily\footnotesize i\allowbreak{}}\par
\noindent Corrigé : {\ttfamily\footnotesize 1\allowbreak{}}\par
\noindent Source : 7II:403. \nolinkurl{fixes/MISPRINTS_7II.json#403}\par

\par\medskip\noindent\hypertarget{app-F7II-XVII-042.4}{}\textbf{F7II-XVII-042.4}\quad p. 247; misprint\_pass.\par
\noindent Imprimé : {\ttfamily\footnotesize lequelle\allowbreak{}}\par
\noindent Corrigé : {\ttfamily\footnotesize lequel\allowbreak{}}\par
\noindent Source : 7II:404. \nolinkurl{fixes/FIXLOG_7II_CURRENT.jsonl#735}; \nolinkurl{fixes/MISPRINTS_7II.json#404}\par

\par\medskip\noindent\hypertarget{app-F7II-XVII-042.5}{}\textbf{F7II-XVII-042.5}\quad p. 247; designated-witness attribution.\par
\noindent Imprimé : \emph{(vide)}\par
\noindent Corrigé : {\ttfamily\footnotesize \textbackslash{}mathbb\{\allowbreak{}}\par
\noindent Source : WITNESS-F7II-XVII-042.5. \nolinkurl{fr/sga7ii/diplomatic/expose_XVII_body.tex}; \nolinkurl{fr/sga7ii/current/expose_XVII_body.tex}\par

\par\medskip\noindent\hypertarget{app-F7II-XVII-042.6}{}\textbf{F7II-XVII-042.6}\quad p. 247; designated-witness attribution.\par
\noindent Imprimé : \emph{(vide)}\par
\noindent Corrigé : {\ttfamily\footnotesize \}\allowbreak{}}\par
\noindent Source : WITNESS-F7II-XVII-042.6. \nolinkurl{fr/sga7ii/diplomatic/expose_XVII_body.tex}; \nolinkurl{fr/sga7ii/current/expose_XVII_body.tex}\par

\par\medskip\noindent\hypertarget{app-F7II-XVII-043.1}{}\textbf{F7II-XVII-043.1}\quad p. 249; printed misprint.\par
\noindent Imprimé : \emph{(vide)}\par
\noindent Corrigé : {\ttfamily\footnotesize \}\textasciicircum{}\{a\allowbreak{}}\par
\noindent Source : 7II:405. \nolinkurl{fixes/MISPRINTS_7II.json#405}\par

\par\medskip\noindent\hypertarget{app-F7II-XVII-044.1}{}\textbf{F7II-XVII-044.1}\quad p. 250; printed misprint.\par
\noindent Imprimé : {\ttfamily\footnotesize r\allowbreak{}}\par
\noindent Corrigé : {\ttfamily\footnotesize n\allowbreak{}}\par
\noindent Source : 7II:406. \nolinkurl{fixes/MISPRINTS_7II.json#406}\par

\par\medskip\noindent\hypertarget{app-F7II-XVII-044.2}{}\textbf{F7II-XVII-044.2}\quad p. 250; misprint\_pass.\par
\noindent Imprimé : {\ttfamily\footnotesize hyperplanelisse\allowbreak{}}\par
\noindent Corrigé : {\ttfamily\footnotesize hyperplane\allowbreak{} lisse\allowbreak{}}\par
\noindent Source : 7II:407. \nolinkurl{fixes/FIXLOG_7II_CURRENT.jsonl#736}; \nolinkurl{fixes/MISPRINTS_7II.json#407}\par

\par\medskip\noindent\hypertarget{app-F7II-XVII-045.1}{}\textbf{F7II-XVII-045.1}\quad p. 252; misprint\_pass.\par
\noindent Imprimé : {\ttfamily\footnotesize paramètrées\allowbreak{}}\par
\noindent Corrigé : {\ttfamily\footnotesize paramétrées\allowbreak{}}\par
\noindent Source : 7II:408. \nolinkurl{fixes/FIXLOG_7II_CURRENT.jsonl#737}; \nolinkurl{fixes/MISPRINTS_7II.json#408}\par

\par\medskip\noindent\hypertarget{app-F7II-XVII-046.1}{}\textbf{F7II-XVII-046.1}\quad p. 253; misprint\_pass.\par
\noindent Imprimé : {\ttfamily\footnotesize colongeurs\allowbreak{}}\par
\noindent Corrigé : {\ttfamily\footnotesize colongueurs\allowbreak{}}\par
\noindent Source : 7II:409. \nolinkurl{fixes/FIXLOG_7II_CURRENT.jsonl#738}; \nolinkurl{fixes/MISPRINTS_7II.json#409}\par

\par\medskip\noindent\hypertarget{app-F7II-XVII-047.1}{}\textbf{F7II-XVII-047.1}\quad p. 253; misprint\_pass.\par
\noindent Imprimé : {\ttfamily\footnotesize 3\allowbreak{}}\par
\noindent Corrigé : {\ttfamily\footnotesize 4\allowbreak{}}\par
\noindent Source : 7II:410. \nolinkurl{fixes/FIXLOG_7II_CURRENT.jsonl#739}; \nolinkurl{fixes/MISPRINTS_7II.json#410}\par

\par\medskip\noindent\hypertarget{app-F7II-XVII-047.2}{}\textbf{F7II-XVII-047.2}\quad p. 253; misprint\_pass.\par
\noindent Imprimé : {\ttfamily\footnotesize S\allowbreak{}}\par
\noindent Corrigé : {\ttfamily\footnotesize s\allowbreak{}}\par
\noindent Source : 7II:411. \nolinkurl{fixes/FIXLOG_7II_CURRENT.jsonl#740}\par

\normalsize\fi
\clearpage\SGAtoc{Exposé XVIII — ETUDE COHOMOLOGIQUE DES PINCEAUX DE LEFSCHETZ}
\selectlanguage{french}

\SGApage{II}{254}


S G A 7


\underline{E X P O S E XVIII}


\underline{\underline{ETUDE COHOMOLOGIQUE DES PINCEAUX DE LEFSCHETZ}}


\underline{par N. Katz (*)}


\underline{SOMMAIRE}

1. La cohomologie d'un fibré projectif. \hfill 2

2. La cohomologie des variétés éclatées. \hfill 4

3. L'éclatement associé à un pinceau. \hfill 13

4. La cohomologie de l'éclatement associé à un pinceau. \hfill 17

5. Etude cohomologique du morphisme $\rho : \widetilde{X} \longrightarrow \mathbb{P}^{1}$ associé à un pinceau. \hfill 24

\quad 5.1. L'homomorphisme de restriction $H^{*}(\widetilde{X}) \longrightarrow H^{*}(Y)$. \hfill 24

\quad 5.2. Le ``théorème de Lefschetz vache'' (LV). \hfill 27

\quad 5.3. La condition (A). \hfill 33

\SGAvar{\quad 5.8. L'homomorphisme de Griffiths. \hfill 47

\quad 5.9. Un exemple où l'hypothèse (A) n'est pas vérifié. \hfill 57

6. Applications de la théorie locale des cycles évanescents (Picard-Lefschetz). \hfill 58

(*) D'après des notes (succintes) de GROTHENDIECK.


}{\quad 5.8. L'homomorphisme de Griffiths. \hfill 47


\quad 5.9. Un exemple où l'hypothèse (A) n'est pas vérifiée. \hfill 57

6. Applications de la théorie locale des cycles évanescents (Picard-Lefschetz). \hfill 58


(*) D'après des notes (succinctes) de GROTHENDIECK.


}{F7II-XVIII-001}\SGApage{II}{255}


\underline{INTRODUCTION}

Dans le présent exposé nous donnons la structure cohomologique élémentaire des pinceaux de Lefschetz. Les sections 1-4 donnent les faits ``bien connus'' sur la structure cohomologique des fibrés projectifs et des éclatements. Dans 5 nous faisons l'étude qui figure dans le titre, en supposant que $X$ satisfait au ``théorème de Lefschetz vache'' (5.2), et que le pinceau $D$ en question satisfait à une certaine condition (A) de nature locale sur $D$ (5.3). Les résultats principaux sont 5.6, 5.7 et 5.8.7, ce  dernier intervenant de façon essentielle dans la démonstration du théorème de GRIFFITHS dans XX. Dans la dernière section, de nature nettement moins élémentaire, nous nous appuyons sur l'exposé SGA 1 XIII de Mme. RAYNAUD sur le \underline{groupe fondamental} ``\underline{modéré}'', et sur la \underline{formule de Picard-Lefschetz} (XVI) pour démontrer :

(1) La condition (A) est souvent satisfaite (6.3, 6.4) ;

(2) Sauf dans le cas $\dim(X)$ impair et caractéristique $= 2$ , le groupe de monodromie d'un pinceau de Lefschetz agit irréductiblement sur la ``cohomologie évanescente'', et transitivement (modulo signe) sur l'ensemble des
\SGApage{II}{256}
"cycles évanescents locaux" associés (via la formule de Picard-Lefschetz) aux divers "points critiques" du pinceau (6.6 et 6.7).


\SGAsectiontoc{1. La cohomologie d'un fibré projectif}{1}
1. \underline{La cohomologie d'un fibré projectif}

\SGAvar{\underline{Théorème} 1.1. (DELIGNE). \underline{Soient} $f : X \longrightarrow S$ \underline{un morphisme propre}, $\ell$ \underline{un nombre premier inversible sur} $S$, $u \in H^2(X,\mathbb{Z}_\ell(1))$ (resp. $u \in H^2(X,\mathbb{Q}_\ell(1))$ \underline{une classe de cohomologie sur} $X$, \underline{et} $\mathfrak{F}$ \underline{un} $\mathbb{Z}_\ell$-\underline{faisceau} (resp. $\mathbb{Q}_\ell$-\underline{faisceau}) \underline{sur} $X$ . \underline{On suppose qu'il existe un entier} $n$ \underline{tel que pour tout entier} $j$ \underline{et tout entier} $i > 0$, \underline{l'application}

}{\underline{Théorème} 1.1. (DELIGNE). \underline{Soient} $f : X \longrightarrow S$ \underline{un morphisme propre}, $\ell$ \underline{un nombre premier inversible sur} $S$, $u \in H^2(X,\mathbb{Z}_\ell(1))$ (resp. $u \in H^2(X,\mathbb{Q}_\ell(1))$) \underline{une classe de cohomologie sur} $X$, \underline{et} $\mathfrak{F}$ \underline{un} $\mathbb{Z}_\ell$-\underline{faisceau} (resp. $\mathbb{Q}_\ell$-\underline{faisceau}) \underline{sur} $X$ . \underline{On suppose qu'il existe un entier} $n$ \underline{tel que pour tout entier} $j$ \underline{et tout entier} $i > 0$, \underline{l'application}

}{F7II-XVIII-002}\begin{equation}\tag{1.1.1}
R^{n-i} f_*(\mathfrak{F}(j)) \ \xrightarrow{\ \wedge u^i\ }\ R^{n+i} f_*(\mathfrak{F}(j+i)) \quad ,
\end{equation}
\underline{soit un isomorphisme de} $\mathbb{Z}_\ell$-\underline{faisceaux}(resp. \underline{de} $\mathbb{Q}_\ell$-\underline{faisceaux}) \underline{sur} $S$ . \underline{Alors pour tout morphisme} $g : S \longrightarrow T$, \underline{la suite spectrale de Leray}

\begin{equation*}
E_2^{p,q} = R^p g_* R^q f_*(\mathfrak{F}) \qquad \Longrightarrow\ R^{p+q}(gf)_*(\mathfrak{F})
\end{equation*}
\underline{est dégénérée.}

La démonstration se trouve dans [1, Prop. 2.4]. On remarquera que, grâce au théorème de changement de base pour un morphisme propre (SGA 4 XII  5.2), l'hypothèse se vérifie fibre par fibre.


\SGApage{II}{257}

\underline{Théorème} 1.2. (Cf. SGA 5 VII 2.2.1 ). \underline{Soit} $\mathcal{E} \longrightarrow X$ \underline{un fibré vectoriel}, \underline{localement libre de rang} $1 + r$. \underline{Désignons par} $f : \mathbb{P}(\mathcal{E}) \longrightarrow X$ \underline{le fibré projectif défini par} $\mathcal{E}$, \underline{et par} $u \in H^2(\mathbb{P}(\mathcal{E}),\mathbb{Z}_\ell(1))$ \underline{la première classe de Chern du faisceau inversible canonique} $\underline{O}_{\mathbb{P}(\mathcal{E})}(1)$. \underline{Alors, pour} $\ell$ \underline{un nombre premier inversible sur} $X$, \underline{on a}

a) \underline{la suite spectrale}

\begin{equation}\tag{1.2.1}
E_2^{p,q} = H^p(X,R^q f_* \mathbb{Z}_\ell) \ \Longrightarrow\ H^{p+q}(\mathbb{P}(\mathcal{E}), \mathbb{Z}_\ell)
\end{equation}
\underline{dégénère.}

\SGAvar{b) \underline{la flèche}

\begin{equation}\tag{1.2.2}
\underset{0 \leq j \leq n}{\oplus} \ H^{q-2-j}(X,\mathbb{Z}_\ell(-j) \ \longrightarrow\ H^q(\mathbb{P}(\mathcal{E}),\mathbb{Z}_\ell)
\end{equation}
\underline{défini par} $\oplus\ \tau_{q-2j} \ \longrightarrow\ \Sigma\ f^*(\tau_{q-2j}) \wedge u^j$

}{b) \underline{la flèche}



\begin{equation}\tag{1.2.2}
\underset{0 \leq j \leq r}{\oplus} \ H^{q-2j}(X,\mathbb{Z}_\ell(-j)) \ \longrightarrow\ H^q(\mathbb{P}(\mathcal{E}),\mathbb{Z}_\ell)
\end{equation}
\underline{définie par} $\oplus\ \tau_{q-2j} \ \longrightarrow\ \Sigma\ f^*(\tau_{q-2j}) \wedge u^j$

}{F7II-XVIII-003}\underline{est un isomorphisme.}

\underline{Démonstration.} Les fibres de $f$ étant des espaces projectifs de dimension $r$ on sait que

\begin{equation}\tag{1.2.3}
R^i f_* \mathbb{Z}_\ell \ \xleftarrow[\ \sim\ ]{\ \wedge^{i/2} u\ } \ \begin{cases} \mathbb{Z}_\ell \left( \dfrac{-i}{2} \right) & \text{pour } i \text{ pair, } 0 \leq i \leq 2r \quad , \\[2ex] 0 & \text{sinon} \quad . \end{cases}
\end{equation}
Le théorème 1.1 s'applique à la situation $(f, \ell, u, \mathbb{Z}_\ell, r)$, d'où a). D'après 1.2.3  on a 


\SGApage{II}{258}

\begin{equation}\tag{1.2.4}
 E_2^{a,b} \;=\; \begin{cases} H^a(X,\mathbb{Z}_{\ell}(\frac{-b}{2})) & \text{pour } b \text{ pair, } 0 \leq b \leq 2r \\[2ex] 0 & \text{sinon,} \end{cases}
\end{equation}
donc le gradué associé de $H^q(\mathbb{P}(\mathcal{E}), \mathbb{Z}_{\ell})$ pour la filtration provenant de la suite spectrale est $\bigoplus\limits_{0 \leq j \leq r} H^{q-2j}(X,\mathbb{Z}_{\ell}(-j))$. On voit tout de suite que, munissant $\bigoplus\limits_{0 \leq j \leq r} H^{q-2j}(X,\mathbb{Z}_{\ell}(-j))$ de la filtration

\SGAvar{\begin{multline*}
F^i \;=\; \bigoplus_{\substack{0 \leq j \leq r \\ q-2j \leq i}} H^{q-2j}(X,\mathbb{Z}_{\ell}(-j)),
\end{multline*}
la flèche 1.2.2 est compatible aux filtrations, et qu'elle induit l'identité sur les gradués associés, d'où b).


}{\begin{multline*}
F^i \;=\; \bigoplus_{\substack{0 \leq j \leq r \\ q-2j \geq i}} H^{q-2j}(X,\mathbb{Z}_{\ell}(-j)),
\end{multline*}
la flèche 1.2.2 est compatible aux filtrations, et qu'elle induit l'identité sur les gradués associés, d'où b).


}{F7II-XVIII-004}\SGAsectiontoc{2. La cohomologie des variétés éclatées (cf. SGA 5 VII  8)}{2}
2. \underline{La cohomologie des variétés éclatées} (cf. SGA 5 VII  8)

2.1. Soient $X$ un schéma lisse sur un corps $k$, et $\Delta \subset X$ un schéma lisse partout de codimension $1 + r$. On sait (EGA IV 19.4.9) qu'il existe un schéma \underline{lisse} $\widetilde{X}$ , muni d'un morphisme projectif

\begin{equation}\tag{2.1.1}
f : \widetilde{X} \longrightarrow X
\end{equation}
tel que, posant

\begin{equation}\tag{2.1.2}
\widetilde{\Delta} \;=\; f^{-1}(\Delta) \;=\; \Delta \times_X \widetilde{X} \quad ,
\end{equation}
$f$ induise un isomorphisme

\begin{equation}\tag{2.1.3}
f \,/\, \widetilde{X} - \widetilde{\Delta} \;:\; \widetilde{X} - \widetilde{\Delta} \xrightarrow{\ \sim\ } X - \Delta \quad ,
\end{equation}
et qu'il existe un $\Delta$-isomorphisme entre $\widetilde{\Delta}$ et $\check{\mathbb{P}}(\underline{N}_{\Delta/X})$, le fibré projectif
\SGApage{II}{259}
des droites dans le fibré normal $\underline{N}_{\Delta/X}$ de $\Delta$ dans $X$. On a donc un diagramme commutatif :

\SGAvar{\begin{equation}\tag{2.1.4}
\begin{tikzcd}[row sep=huge, column sep=small]
\widetilde{\Delta} \arrow[rd, "f|\widetilde{\Delta}"'] \& \simeq \& \check{\mathbb{P}}\,(\underline{N}_{\Delta|X}) \arrow[ld, "\text{proj. canonique}"] \\
{} \& \Delta \& .
\end{tikzcd}
\end{equation}















2.1.5. Dans le cas spécial $X = \mathrm{Spec}(A)$, $\Delta = \mathrm{Spec}(A/I)$ où $I = (g_o,\ldots,g_r)$, $\widetilde{X}$ s'obtient par le procédé de Hopf (cf. EGA IV 19.4.11) :

}{\begin{equation}\tag{2.1.4}
\begin{tikzcd}[row sep=huge, column sep=small]
\widetilde{\Delta} \arrow[rd, "f|\widetilde{\Delta}"'] \& \simeq \& \check{\mathbb{P}}\,(\underline{N}_{\Delta|X}) \arrow[ld, "\text{proj. canonique}"] \\
{} \& \Delta \& .
\end{tikzcd}
\end{equation}

















2.1.5. Dans le cas spécial $X = \mathrm{Spec}(A)$, $\Delta = \mathrm{Spec}(A/I)$ où $I = (g_0,\ldots,g_r)$, $\widetilde{X}$ s'obtient par le procédé de Hopf (cf. EGA IV 19.4.11) :

}{F7II-XVIII-005}On définit un morphisme de $X - \Delta$ dans $\mathbb{P}^r$ par

\SGAvar{\begin{equation*}
x \in X - \Delta \longrightarrow (g_o(x),\ldots,g_r(x) \in \mathbb{P}^r \quad ,
\end{equation*}
et on considère son graphe $\Gamma$ comme plongé dans $X \times \mathbb{P}^r$. Alors on définit $\widetilde{X} = $ l'adhérence schématique de $\Gamma$ dans $X \times \mathbb{P}^r$, le morphisme $f : \widetilde{X} \longrightarrow X$ s'obtient comme le composé

}{\begin{equation*}
x \in X - \Delta \longrightarrow (g_0(x),\ldots,g_r(x)) \in \mathbb{P}^r \quad ,
\end{equation*}
et on considère son graphe $\Gamma$ comme plongé dans $X \times \mathbb{P}^r$. Alors on définit $\widetilde{X} = $ l'adhérence schématique de $\Gamma$ dans $X \times \mathbb{P}^r$, le morphisme $f : \widetilde{X} \longrightarrow X$ s'obtient comme le composé

}{F7II-XVIII-006}\begin{equation*}
\widetilde{X} \lhook\joinrel\longrightarrow X \times \mathbb{P}^r \xrightarrow{\ \mathrm{pr}_1\ } X \quad .
\end{equation*}
2.1.6. De façon analogue (cf. encore EGA IV 19.4.11), si $X$ est projectif, et si $\Delta$ est l'intersection de $X$ avec $1 + r$ hypersurfaces,disons d'équations $G_i = 0$, $i = 0,\ldots,r$, on considère encore le morphisme

\SGAvar{\begin{equation*}
X - \Delta \longrightarrow \mathbb{P}^r
\end{equation*}
\begin{equation*}
x \in X- \Delta \longrightarrow (G_o(x),\ldots,G_r(x)),
\end{equation*}
et on obtient $\widetilde{X}$ comme l'adhérence de son graphe dans $X \times \mathbb{P}^r$. De façon
\SGApage{II}{260}
concrète, en termes des coordonnées homogènes $\lambda_o,\ldots,\lambda_r$ dans $\mathbb{P}^r$, $\widetilde{X}$ s'obtient comme le schéma de $X \times \mathbb{P}^r$ défini par les équations

}{\begin{equation*}
X - \Delta \longrightarrow \mathbb{P}^r
\end{equation*}
\begin{equation*}
x \in X- \Delta \longrightarrow (G_0(x),\ldots,G_r(x)),
\end{equation*}
et on obtient $\widetilde{X}$ comme l'adhérence de son graphe dans $X \times \mathbb{P}^r$. De façon
\SGApage{II}{260}
concrète, en termes des coordonnées homogènes $\lambda_0,\ldots,\lambda_r$ dans $\mathbb{P}^r$, $\widetilde{X}$ s'obtient comme le schéma de $X \times \mathbb{P}^r$ défini par les équations

}{F7II-XVIII-007}\begin{equation}\tag{2.1.6.1}
\lambda_i \; G_j(X) - \lambda_j G_i(X) \qquad 0 \leq i < j \leq r \qquad .
\end{equation}
On voit par ces équations que le point $(x,\lambda)$ se trouve dans $\widetilde{\Delta}$ si et seulement si $x \in \Delta$ ; autrement dit, $\check{\mathbb{P}} (\underline{N}_{\Delta/X})$ est le fibré projectif trivial

\SGAvar{\begin{equation}\tag{2.1.6.2}
\begin{tikzcd}[column sep=large, row sep=large]
\check{\mathbb{P}}(\underline{N}_{\Delta/X}) \arrow[rr, phantom, "\simeq"] \arrow[dr, "\text{proj.canonique}"'] \& {} \& \Delta \times \mathbb{P}^r \arrow[dl, "\mathrm{pr}_i"] \\
{} \& \Delta \& .
\end{tikzcd}
\end{equation}
















}{\begin{equation}\tag{2.1.6.2}
\begin{tikzcd}[column sep=large, row sep=large]
\check{\mathbb{P}}(\underline{N}_{\Delta/X}) \arrow[rr, phantom, "\simeq"] \arrow[dr, "\text{proj.canonique}"'] \& {} \& \Delta \times \mathbb{P}^r \arrow[dl, "\mathrm{pr}_1"] \\
{} \& \Delta \& .
\end{tikzcd}
\end{equation}
















}{F7II-XVIII-008}On a le diagramme commutatif suivant

\begin{equation}\tag{2.1.6.3}
\begin{tikzcd}[column sep=large, row sep=large]
\Delta \times \mathbb{P}^r \arrow[r, "\sim"] \arrow[dr, "\mathrm{pr}_1"] \& \widetilde{\Delta} \arrow[r, hook] \arrow[d, "f|\widetilde{\Delta}"] \& \widetilde{X} \arrow[r, hook] \arrow[d, "f"] \& X \times \mathbb{P}^r \arrow[dl, "\mathrm{pr}_1"] \\
{} \& \Delta \arrow[r, hook] \& X \& ,
\end{tikzcd}
\end{equation}




















et l'inclusion composée $\Delta \times \mathbb{P}^r \hookrightarrow X \times \mathbb{P}^r$ provient de $\Delta \hookrightarrow X$ par changement de base.

\underline{Théorème} 2.2.  \underline{Soient} $X$ \underline{lisse sur un corps} $k$, $\Delta$ \underline{un sous-schéma lisse partout de codimension} $1 + r$, \underline{et} $f : \widetilde{X} \longrightarrow X$ \underline{l'éclatement de} $X$ \underline{en} $\Delta$. \underline{Alors} :

(2.2.1) \underline{La suite spectrale de Leray}

\SGAvar{\begin{equation*}
E_2^{p,q} = H^p(X, R^q f_* \mathbb{Z}_\ell \quad \Longrightarrow \quad H^{p+q}(\widetilde{X}_o \;\; \mathbb{Z}_\ell)
\end{equation*}
\underline{dégénère}.


}{\begin{equation*}
E_2^{p,q} = H^p(X, R^q f_* \mathbb{Z}_\ell) \quad \Longrightarrow \quad H^{p+q}(\widetilde{X},\mathbb{Z}_\ell)
\end{equation*}
\underline{dégénère}.


}{F7II-XVIII-009}\SGApage{II}{261}

2.2.2.  \underline{Il existe un isomorphisme additif}

\begin{equation*}
\begin{gathered}
H^q(\widetilde{X}, \mathbb{Z}_\ell) \;\overset{\sim}{\longleftarrow}\; H^q(X, \mathbb{Z}_\ell) \oplus H^q(\widetilde{\Delta}, \mathbb{Z}_\ell)/\mathrm{Im}\; H^q(\Delta, \mathbb{Z}_\ell) \quad \simeq \\
\bigoplus_{j=1}^{r} H^{q-2j}(\Delta, \mathbb{Z}_\ell(-j)) \;\oplus\; H^q(X, \mathbb{Z}_\ell) \qquad .
\end{gathered}
\end{equation*}
\underline{Démonstration.}  Désignons par

\begin{equation}\tag{2.2.3}
g : \widetilde{\Delta} \longrightarrow \Delta
\end{equation}
la restriction de $f$ à $\widetilde{\Delta}$ , et par

\begin{equation}\tag{2.2.4}
\left\{\;
\begin{array}{l}
j : \Delta \hookrightarrow X \\[6ex]
i : \widetilde{\Delta} \hookrightarrow \widetilde{X}
\end{array}
\right.
\end{equation}
les inclusions, de façon à obtenir le diagramme commutatif

\begin{equation}\tag{2.2.5}
\begin{tikzcd}[column sep=huge, row sep=large]
\widetilde{X} \arrow[d, "f"'] \& \widetilde{\Delta} \arrow[l, "i"'] \arrow[d, "g"] \\
X \& \Delta \arrow[l, "j"'] \& .
\end{tikzcd}
\end{equation}













Celui-ci induit un homomorphisme des suites spectrales de Leray pour les morphismes $f$ et $g$, et le faisceau $\mathbb{Z}_\ell$ sur $\widetilde{X}$ et $\widetilde{\Delta}$ ,

\begin{equation}\tag{2.2.6}
E_r^{p,q} \, (f) \;\xrightarrow{\;\;(i^*,j^*)\;\;}\; E_r^{p,q}(g)
\end{equation}
D'autre part, $f$ induisant un isomorphisme de $\widetilde{X} - \widetilde{\Delta}$ avec $X - \Delta$ , nous avons, par le théorème de changement de base pour un morphisme propre, 


\SGApage{II}{262}

\begin{equation}\tag{2.2.7}
 R^q f_* \mathbb{Z}_\ell \quad \simeq \quad \begin{cases} \mathbb{Z}_\ell & \text{si } q = 0 \\[2ex] j_* R^q g_* \mathbb{Z}_\ell & \text{si } q > 0 \quad , \end{cases}
\end{equation}
d'où

\begin{multline}\tag{2.2.8}
E_2^{p,q}(f) \;\xrightarrow{\;\substack{(i^*,j^*)\\ \sim}\;}\; E_2^{p,q}(g) \qquad \text{pour } q \neq 0 \ .
\end{multline}
Nous allons démontrer par récurrence sur $r$ que

\begin{multline}\tag{2.2.9}
E_r^{p,q}(f) \;\xrightarrow{\;\substack{(i^*,j^*)\\ \sim}\;}\; E_r^{p,q}(g) \quad \text{pour } q \neq 0, \ r \geq 2
\end{multline}
et

\begin{equation}\tag{2.2.10}
d_r^{p,q}(f) = 0 \quad \text{pour } r \geq 2 \quad .
\end{equation}
Admettons pour l'instant la formule

\begin{equation}\tag{2.2.11}
d_r^{p,r-1}(f) = 0 \quad \text{pour } r \geq 2, \ \text{tout } p \ .
\end{equation}
Commençons par $r = 2$ . Nous avons un diagramme commutatif

\begin{equation}\tag{2.2.12}
\begin{tikzcd}[column sep=huge, row sep=large]
E_2^{p,q}(f) \arrow[r, "{(i^*,j^*)}"] \arrow[d, "{d_2^{p,q}(f)}"] \& E_2^{p,q}(g) \arrow[d, "{0 = d_2^{p,q}(g)}"] \\
E_2^{p+2,q-1}(f) \arrow[r, "{(i^*,j^*)}"] \& E_2^{p+2,q-1}(g) \& .
\end{tikzcd}
\end{equation}










Pour $q \leq 0$, on a $E_2^{p+2,q-1}(f) = 0$, donc $d_2^{p,q}(f) = 0$ . Pour $q = 1$, on a $d_2^{p,1}(f) = 0$ par (2.2.11), et pour $q > 1$, les flèches horizontales dans (2.2.12) sont des isomorphismes, et on sait que $d_2^{p,q}(g) = 0$ par 1.2 et 2.1.4. Ceci donne (2.2.10) pour $r = 2$, et (2.2.9) pour $r = 2$ n'est autre que (2.2.8).


\SGAvar{\SGApage{II}{263}

Supposons connues (2.2.9) et (2.2.10) pour une valeur de $r \geq 2$ . L'hypothèse (2.2.10) implique $E_r^{p,q}(f) \simeq E_{r+1}^{p,q}(f)$, et pat 1.2 on a $E_r^{p,q}(g) \simeq E_{r+1}^{p,q}(g)$, de sorte qu'on a le diagramme commutatif

\begin{equation}\tag{2.2.13}
\begin{tikzcd}[column sep=huge, row sep=large]
E_r^{p,q}(f) \arrow[r, "{(i^*,j^*)}"] \arrow[d, "\wr"'] \& E_r^{p,q}(g) \arrow[d, "\wr"] \\
R_{r+1}^{p,q}(f) \arrow[r, "{(i^*,j^*)}"] \& E_{r+1}^{p,q}(g) \& ,
\end{tikzcd}
\end{equation}












}{\SGApage{II}{263}


Supposons connues (2.2.9) et (2.2.10) pour une valeur de $r \geq 2$ . L'hypothèse (2.2.10) implique $E_r^{p,q}(f) \simeq E_{r+1}^{p,q}(f)$, et par 1.2 on a $E_r^{p,q}(g) \simeq E_{r+1}^{p,q}(g)$, de sorte qu'on a le diagramme commutatif

\begin{equation}\tag{2.2.13}
\begin{tikzcd}[column sep=huge, row sep=large]
E_r^{p,q}(f) \arrow[r, "{(i^*,j^*)}"] \arrow[d, "\wr"'] \& E_r^{p,q}(g) \arrow[d, "\wr"] \\
E_{r+1}^{p,q}(f) \arrow[r, "{(i^*,j^*)}"] \& E_{r+1}^{p,q}(g) \& ,
\end{tikzcd}
\end{equation}












}{F7II-XVIII-010}ce qui donne (2.2.9) pour $r+1$, compte tenu de (2.2.9) pour $r$. Pour déduire (2.2.10) pour $r+1$, on regarde le diagramme commutatif

\SGAvar{\begin{equation}\tag{2.2.14}
\begin{tikzcd}[column sep=huge, row sep=large]
E_{r+1}^{p,q}(f) \arrow[r, "{(i^*,j^*)}"] \arrow[d, "{d_{r+1}^{p,q}(f)}"] \& E_{r+1}^{p,q}(g) \arrow[d, "{0 = d_{r+1}^{p,q}(g)}"] \\
E_{r+1}^{p+r+1,q-r}(f) \arrow[r, "{(i^*,j^*)}"] \& E_{r+1}^{p+r+1,q-r}(g) \& .
\end{tikzcd}
\end{equation}












Pour $q \leq r-1$, on a $E_{r+1}^{p+r+1,q-r}(f) = 0$, donc $d_r^{p,q}(f) = 0$ . Pour $q = r$, on a (2.2.11), et pour $q > r$ les flèches horizontales dans (2.2.14) sont des isomorphismes, et $d_{r+1}^{p,q}(g)= 0$ par 1.2, d'où la conclusion voulue.

}{\begin{equation}\tag{2.2.14}
\begin{tikzcd}[column sep=huge, row sep=large]
E_{r+1}^{p,q}(f) \arrow[r, "{(i^*,j^*)}"] \arrow[d, "{d_{r+1}^{p,q}(f)}"] \& E_{r+1}^{p,q}(g) \arrow[d, "{0 = d_{r+1}^{p,q}(g)}"] \\
E_{r+1}^{p+r+1,q-r}(f) \arrow[r, "{(i^*,j^*)}"] \& E_{r+1}^{p+r+1,q-r}(g) \& .
\end{tikzcd}
\end{equation}













Pour $q \leq r-1$, on a $E_{r+1}^{p+r+1,q-r}(f) = 0$, donc $d_{r+1}^{p,q}(f) = 0$ . Pour $q = r$, on a (2.2.11), et pour $q > r$ les flèches horizontales dans (2.2.14) sont des isomorphismes, et $d_{r+1}^{p,q}(g)= 0$ par 1.2, d'où la conclusion voulue.

}{F7II-XVIII-011}Il reste à vérifier (2.2.11), qui signifie que le "edge homomorphisme" composé

\begin{equation}\tag{2.2.15}
E_2^{p,0}(f) \;\longrightarrow\; E_3^{p,0}(f) \longrightarrow E_4^{p,0}(f) \longrightarrow \ldots E_\infty^{p,0}(f)
\end{equation}
est un isomorphisme, i.e. que

\begin{equation}\tag{2.2.16}
f^* : H^p(X,\mathbb{Z}_\ell) \;\longrightarrow\; H^p(\widetilde{X},\ \mathbb{Z}_\ell) 
\end{equation}

\SGApage{II}{264}

 est injectif. Or, il existe un inverse à gauche, savoir \underline{l'homomorphisme de Gysin} \quad (SGA 5 IV)

\begin{equation}\tag{2.2.17}
f_{*} \;:\; H^{p}(\widetilde{X},\,\mathbb{Z}_{\ell}) \longrightarrow H^{p}(X,\mathbb{Z}_{\ell}) \qquad ,
\end{equation}
qui vérifie en effet,

\begin{equation}\tag{2.2.18}
f_{*}f^{*} = \mathrm{id}.
\end{equation}
comme il résulte de la formule de projection, et du fait que $f_{*}(1) = 1$, $f$ étant birationnel donc de degré un.

Ceci démontre (2.2.1). Quant à (2.2.2), remarquons que le diagramme (2.2.5) donne, par passage à la cohomologie, un diagramme commutatif

\begin{equation}\tag{2.2.19}
\begin{tikzcd}[row sep=huge, column sep=huge]
H^{p}(\widetilde{X},\mathbb{Z}_{\ell}) \arrow[r, "i^{*}"] \& H^{p}(\widetilde{\Delta},\mathbb{Z}_{\ell}) \\
H^{p}(X,\mathbb{Z}_{\ell}) \arrow[r, "j^{*}"] \arrow[u, "f^{*}"'] \& H^{p}(\Delta,\mathbb{Z}_{\ell}) \;\;, \arrow[u, "g^{*}"']
\end{tikzcd}
\end{equation}













d'où un morphisme

\begin{equation}\tag{2.2.20}
\frac{H^{p}(\widetilde{X},\mathbb{Z}_{\ell})}{f^{*}H^{p}(X,\mathbb{Z}_{\ell})} \;\stackrel{(i^{*},j^{*})}{\longrightarrow}\; \frac{H^{p}(\widetilde{\Delta}\,,\mathbb{Z}_{\ell})}{g^{*}H^{p}(\Delta,\mathbb{Z}_{\ell})} \qquad\qquad .
\end{equation}
Désignons par $F^{i}$ les filtrations de $H^{\cdot}(\widetilde{X},\mathbb{Z}_{\ell})$ et de $H^{\cdot}(\widetilde{\Delta},\mathbb{Z}_{\ell})$ provenant des suites spectrales de Leray, de sorte que (2.2.20) se récrit 


\SGApage{II}{265}

\SGAvar{\begin{equation}\tag{2.2.21}
\frac{H^{p}(\widetilde{X},\mathbb{Z}_{\ell})}{F^{p}H^{p}(\widetilde{X},\mathbb{Z}_{\ell})} \;\stackrel{(i^{*},j^{*})}{\longrightarrow}\; \frac{H^{p}(\widetilde{\Delta},\mathbb{Z}_{\ell})}{H^{p}F^{p}(\widetilde{\Delta},\mathbb{Z}_{\ell})} \qquad\qquad .
\end{equation}
Or, d'après (2.2.9), le gradué associé à la flèche (2.2.21) est un isomorphisme, donc (2.2.20) est un isomorphisme. D'autre part, (2.2.18) donne une décomposition en somme direct

}{\begin{equation}\tag{2.2.21}
\frac{H^{p}(\widetilde{X},\mathbb{Z}_{\ell})}{F^{p}H^{p}(\widetilde{X},\mathbb{Z}_{\ell})} \;\stackrel{(i^{*},j^{*})}{\longrightarrow}\; \frac{H^{p}(\widetilde{\Delta},\mathbb{Z}_{\ell})}{F^{p}H^{p}(\widetilde{\Delta},\mathbb{Z}_{\ell})} \qquad\qquad .
\end{equation}
Or, d'après (2.2.9), le gradué associé à la flèche (2.2.21) est un isomorphisme, donc (2.2.20) est un isomorphisme. D'autre part, (2.2.18) donne une décomposition en somme directe

}{F7II-XVIII-012}\begin{equation}\tag{2.2.22}
\begin{tikzcd}[column sep=7em]
H^{\cdot}(\widetilde{X},\mathbb{Z}_{\ell}) \arrow[r, shift right=1.2ex, "f_{*}\oplus(1-f^{*}f_{*})"'] \& H^{\cdot}(X,\mathbb{Z}_{\ell}) \oplus \mathrm{Ker}\, f_{*} \;\;, \arrow[l, shift right=1.2ex, "f^{*}+\mathrm{inclusion}"']
\end{tikzcd}
\end{equation}









et (1.2.2) nous donne une décomposition en somme directe

\SGAvar{\begin{equation}\tag{2.2.23}
\frac{H^{\cdot}(\widetilde{\Delta},\mathbb{Z}_{\ell})}{g^{*}H^{\cdot}(\Delta,\mathbb{Z}_{\ell})} \;\;\stackrel{\sim}{\longleftarrow}\;\; \bigoplus_{1\leq j\leq r} H^{\cdot-2j}(\Delta,\mathbb{Z}_{\ell}(-j)) \;\;;
\end{equation}
mettant ensemble l'isomorphisme (2.2.20) et les décompositions (2.2.22) et (2.2.23), on trouve (2.2.2), cqfd.

\underline{Remarque} 2.2.24. \quad Esquissons une deuxième démonstration de (2.2.2) nettement plus simple, valable dans le cas où $X$ est \underline{propre} sur $k$. Plus généralement, on trouve (sans hypothèse sur $X$) le résultat correspondant (2.2.26) en cohomologie à \ support propre. Nous avons les suites exactes de"cohomologie à supports propres", notée $H^{\cdot}_{c}$ , pour les couples $(\widetilde{X},\Delta)$ et $(X,\Delta)$ qui forment un diagramme commutatif 


}{\begin{equation}\tag{2.2.23}
\frac{H^{\cdot}(\widetilde{\Delta},\mathbb{Z}_{\ell})}{g^{*}H^{\cdot}(\Delta,\mathbb{Z}_{\ell})} \;\;\stackrel{\sim}{\longleftarrow}\;\; \bigoplus_{1\leq j\leq r} H^{\cdot-2j}(\Delta,\mathbb{Z}_{\ell}(-j)) \;\;;
\end{equation}
mettant ensemble l'isomorphisme (2.2.20) et les décompositions (2.2.22) et (2.2.23), on trouve (2.2.2), cqfd.


\underline{Remarque} 2.2.24. \quad Esquissons une deuxième démonstration de (2.2.2) nettement plus simple, valable dans le cas où $X$ est \underline{propre} sur $k$. Plus généralement, on trouve (sans hypothèse sur $X$) le résultat correspondant (2.2.26) en cohomologie à support propre. Nous avons les suites exactes de « cohomologie à supports propres », notée $H^{\cdot}_{c}$, pour les couples $(\widetilde{X},\widetilde{\Delta})$ et $(X,\Delta)$ qui forment un diagramme commutatif


}{F7II-XVIII-013}\SGApage{II}{266}

\begin{tikzcd}[column sep=small, row sep=large]
(2.2.25) \& {} \arrow[r] \& H^i_c(\widetilde{X}-\widetilde{\Delta},\mathbb{Z}_\ell) \arrow[r] \& H^i_c(\widetilde{X},\mathbb{Z}_\ell) \arrow[r] \& H^i_c(\widetilde{\Delta},\mathbb{Z}_\ell) \arrow[r] \& H^{i+1}_c(\widetilde{X}-\widetilde{\Delta},\mathbb{Z}_\ell) \\
(2.2.26) \& {} \arrow[r] \& H^i_c(X-\Delta,\mathbb{Z}_\ell) \arrow[r] \arrow[u, "\wr" left] \& H^i_c(X,\mathbb{Z}_\ell) \arrow[r] \arrow[u, hook, "f^*" right] \& H^i_c(\Delta,\mathbb{Z}_\ell) \arrow[r] \arrow[u, hook, "g^*" right] \& H^{i+1}_c(X-\Delta,\mathbb{Z}_\ell)\;.\arrow[u, "\wr" left]
\end{tikzcd}





















Or, le quotient du complexe acyclique (2.2.25) par le sous-complexe acyclique (2.2.26) nous donne le complexe \underline{acyclique}

\SGAvar{\begin{equation}\tag{2.2.27}
\longrightarrow 0 \longrightarrow \frac{H^i_c(\widetilde{X},\mathbb{Z}_\ell)}{f^*H^i_c(X,\mathbb{Z}_\ell)} \longrightarrow \frac{H^i_C(\widetilde{\Delta},\mathbb{Z}_\ell)}{g^*H^i_c(\Delta,\mathbb{Z}_\ell)} \longrightarrow 0 \longrightarrow
\end{equation}
i.e, on trouve, sans considérer des suites spectrales, que (2.2.20) est un isomorphisme. Dès lors, les raisonnements qu'on vient de donner s'appliquent pour donner (2.2.2).

}{\begin{equation}\tag{2.2.27}
\longrightarrow 0 \longrightarrow \frac{H^i_c(\widetilde{X},\mathbb{Z}_\ell)}{f^*H^i_c(X,\mathbb{Z}_\ell)} \longrightarrow \frac{H^i_c(\widetilde{\Delta},\mathbb{Z}_\ell)}{g^*H^i_c(\Delta,\mathbb{Z}_\ell)} \longrightarrow 0 \longrightarrow
\end{equation}
i.e., on trouve, sans considérer des suites spectrales, que (2.2.20) est un isomorphisme. Dès lors, les raisonnements qu'on vient de donner s'appliquent pour donner (2.2.2).

}{F7II-XVIII-014}\underline{Corollaire} 2.3.  \underline{Posons}

\SGAvar{\begin{equation*}
b_q(X) = \dim_{\mathbb{Q}_\ell} H^q(\chi,\mathbb{Q}_\ell), \quad \chi(X) = \Sigma(-1)^q b_q(X) \quad .
\end{equation*}
\underline{On a}

\begin{equation}\tag{2.3.1}
b_q(\widetilde{X}) = b_q(X) + \sum_{j=1}^{r} b_{q-j}(\Delta) \quad ,
\end{equation}
\begin{equation}\tag{2.3.2}
\chi(\widetilde{X}) = \chi(X) + r\,\chi(\Delta) \quad .
\end{equation}

}{\begin{equation*}
b_q(X) = \dim_{\mathbb{Q}_\ell} H^q(X,\mathbb{Q}_\ell), \quad \chi(X) = \Sigma(-1)^q b_q(X) \quad .
\end{equation*}
\underline{On a}

\begin{equation}\tag{2.3.1}
b_q(\widetilde{X}) = b_q(X) + \sum_{j=1}^{r} b_{q-2j}(\Delta) \quad ,
\end{equation}
\begin{equation}\tag{2.3.2}
\chi(\widetilde{X}) = \chi(X) + r\,\chi(\Delta) \quad .
\end{equation}

}{F7II-XVIII-015}\SGApage{II}{267}


\SGAvar{\SGAsectiontoc{3. L'éclatement associé à un pinceau}{3}
3. \underline{L'éclatement associé à un pinceau}

3.1.  Soient $X$ projectif, lisse et irréductible sur un corps $k$ algébriquement clos, de dimension $u$, et $D = \{H_\tau\}_{\tau \in \mathbb{P}^1}$ un pinceau d'hypersurfaces dont l'axe coupe transversalement $X$, en $\Delta$ . Considérons alors l'éclatement de $X$ en $\Delta$, que nous allons décrire explicitement.

En termes d'équations, désignant par $(\lambda,\mu)$ des coordonnés homogènes de $\mathbb{P}^1$, et par $F$ et $G$ deux formes linéaires tels que $F = G = 0$ définisse l'axe du pinceau, le pinceau s'écrit

}{\SGAsectiontoc{3. L'éclatement associé à un pinceau}{3}
3. \underline{L'éclatement associé à un pinceau}



3.1.  Soient $X$ projectif, lisse et irréductible sur un corps $k$ algébriquement clos, de dimension $n$, et $D = \{H_\tau\}_{\tau \in \mathbb{P}^1}$ un pinceau d'hypersurfaces dont l'axe coupe transversalement $X$, en $\Delta$ . Considérons alors l'éclatement de $X$ en $\Delta$, que nous allons décrire explicitement.

En termes d'équations, désignant par $(\lambda,\mu)$ des coordonnées homogènes de $\mathbb{P}^1$, et par $F$ et $G$ deux formes linéaires tels que $F = G = 0$ définisse l'axe du pinceau, le pinceau s'écrit

}{F7II-XVIII-016}\begin{equation}\tag{3.1.1}
\lambda F - \mu G \quad .
\end{equation}
Le sous-schéma $\Delta$ est par définition l'intersection de $X$ avec la sous-variété linéaire $G = F = 0$ . D'après 2.1.6,

\SGAvar{\begin{equation}\tag{3.1.2}
\widetilde{X} = \{(x,(\lambda,\mu) \in X \times \mathbb{P}^1 \mid \lambda F - \mu G = 0\} \quad ;
\end{equation}
c'est l'adhérence schématique, dans $X \times \mathbb{P}^1$, du graphe de l'application

}{\begin{equation}\tag{3.1.2}
\widetilde{X} = \{(x,(\lambda,\mu)) \in X \times \mathbb{P}^1 \mid \lambda F - \mu G = 0\} \quad ;
\end{equation}
c'est l'adhérence schématique, dans $X \times \mathbb{P}^1$, du graphe de l'application

}{F7II-XVIII-017}\begin{equation}\tag{3.1.3}
\begin{tikzcd}[row sep=large, column sep=huge]
X - \Delta \arrow[r] \& \mathbb{P}^1 \\
x \arrow[r] \& (G(x),F(x))\;.
\end{tikzcd}
\end{equation}









La grande vertu de $\widetilde{X}$ est d'être une variété projective (car plongée dans $X \times \mathbb{P}^1$) et lisse(car $\Delta$ l'est) sur $k$, qui admet un morphisme $\rho$ vers $\mathbb{P}^1$, à savoir

\begin{equation}\tag{3.1.4}
\rho : \widetilde{X} \longrightarrow X \times \mathbb{P}^1 \stackrel{\mathrm{pr}_2}{\longrightarrow} \mathbb{P}^1 \quad ,
\end{equation}
dont les fibres sont les sections hyperplanes de $X$ par les membres de notre pinceau

\begin{equation*}
\rho^{-1}(t) = X_t \quad \stackrel{\text{déf}}{=} \quad X.H_t \quad .
\end{equation*}

\SGApage{II}{268}


3.2. \underline{Le degré de la variété duale} : \underline{méthode topologique.}

Supposons en plus que le pinceau de 3.1 soit tel que

\SGAvar{\begin{equation}\tag{3.2.1}
X_t \ \text{est lisse pour tout } t \text{ contenu dans un ouvert non vide de } \mathbb{P}^1 ,
\end{equation}
\begin{equation}\tag{3.2.2}
\text{Si } X_t \text{ est singulier, elle n'a que des singularités quadratiques non dégénérées (XVII 1.1) ;}
\end{equation}
ces conditions sont vérifiées en particulier pour un \underline{pinceau de Lefschetz} (XVII 2.21) si car $k \neq 2$ ou si $X$ est de dimension paire.

Soient $\eta$ le point générique de $\mathbb{P}^1$, $\bar{\eta}$ un point géométrique au-dessus, $n$ la dimension de $X$, et $\chi$ la caractéristique d'Euler-Pincaré à valeurs dans le faisceau $\mathbb{Q}_\ell$ . Alors la formule d'Euler-Poincaré XVI 2.1 appliquée au morphisme $S$, nous donne

\begin{multline}\tag{3.2.3}
\chi(\widetilde{X}) \;=\; \chi(\mathbb{P}^1)\, \chi\,(X_{\bar{\eta}}) + (-1)^{\eta} \sum_{\substack{\text{pts sing.}\\ \text{des fibres}}} 1 \qquad ,
\end{multline}
formule qui, compte tenu de (2.2.2), se récrit

\begin{equation}\tag{3.2.4}
(-1)^{\eta} \left[\begin{array}{c} \text{nombre total des} \\[4pt] \text{points singuliers} \end{array}\right] = \chi(X) + \chi(\Delta) - 2\,\chi(X_{\bar{\eta}})
\end{equation}
On remarquera que la quantité $\chi(X) + \chi(\Delta) - 2\chi\,(X_{\bar{\eta}})$ est indépendante du choix particulier du pinceau, et ne dépend que du plongement de $X$ dans un $\mathbb{P}^r$, et même seulement de la classe de Chern de $\underline{O}_X(1)$, cf. [XVII, 5.7.5 et 5.7.6].


}{\begin{equation}\tag{3.2.1}
X_t \ \text{est lisse pour tout } t \text{ contenu dans un ouvert non vide de } \mathbb{P}^1 ,
\end{equation}
\begin{equation}\tag{3.2.2}
\text{Si } X_t \text{ est singulier, elle n'a que des singularités quadratiques non dégénérées (XVII 1.1) ;}
\end{equation}
ces conditions sont vérifiées en particulier pour un \underline{pinceau de Lefschetz} (XVII 2.21) si car $k \neq 2$ ou si $X$ est de dimension paire.



Soient $\eta$ le point générique de $\mathbb{P}^1$, $\bar{\eta}$ un point géométrique au-dessus, $n$ la dimension de $X$, et $\chi$ la caractéristique d'Euler-Poincaré à valeurs dans le faisceau $\mathbb{Q}_\ell$ . Alors la formule d'Euler-Poincaré XVI 2.1 appliquée au morphisme $\rho$, nous donne

\begin{multline}\tag{3.2.3}
\chi(\widetilde{X}) \;=\; \chi(\mathbb{P}^1)\, \chi\,(X_{\bar{\eta}}) + (-1)^{n} \sum_{\substack{\text{pts sing.}\\ \text{des fibres}}} 1 \qquad ,
\end{multline}
formule qui, compte tenu de (2.2.2), se récrit

\begin{equation}\tag{3.2.4}
(-1)^{n} \left[\begin{array}{c} \text{nombre total des} \\[4pt] \text{points singuliers} \end{array}\right] = \chi(X) + \chi(\Delta) - 2\,\chi(X_{\bar{\eta}})
\end{equation}
On remarquera que la quantité $\chi(X) + \chi(\Delta) - 2\chi\,(X_{\bar{\eta}})$ est indépendante du choix particulier du pinceau, et ne dépend que du plongement de $X$ dans un $\mathbb{P}^r$, et même seulement de la classe de Chern de $\underline{O}_X(1)$, cf. [XVII, 5.7.5 et 5.7.6].


}{F7II-XVIII-018}\SGApage{II}{269}

3.2.5. Explicitons ceci dans un cas concret. Prenons $X = \mathbb{P}^2$ , $\{H_\tau\}$ un pinceau de courbes de degré $d$ dans $\mathbb{P}^2$ qui vérifie (3.1), (3.2.1) et (3.2.2). On a

\begin{equation}\tag{3.2.6}
\begin{array}{l@{\qquad\qquad}l} \chi(X) = 3 & (\text{en vertu de 1.2}) \\[6pt] \chi(X_{\bar{\eta}}) = -d(d-3) & (\text{cf. SGA 5 VII 7.3}) \\[6pt] \chi(\Delta) = d^2 \qquad , & \end{array}
\end{equation}
d'où

\begin{equation}\tag{3.2.7}
\text{le nombre total de points singuliers est } \; 3\,(d-1)^2
\end{equation}
Prenons le cas particulier $d = 3$, et considérons la famille des courbes elliptiques (en caractéristique $\neq 3$)

\SGAvar{\begin{equation}\tag{3.2.8}
X^3 + Y^3 + Z^3 + 3\mu\, XYZ \qquad .
\end{equation}
Les quatre fibres singulières sont au-dessus de $\mu =$ racine cubique de l'unité, et $\mu = \infty$, et chaque fibre singulière est l'union de trois droites comme ceci

}{\begin{equation}\tag{3.2.8}
X^3 + Y^3 + Z^3 + 3\mu\, XYZ \qquad .
\end{equation}


Les quatre fibres singulières sont au-dessus de $\mu^3=-1$, et $\mu = \infty$, et chaque fibre singulière est l'union de trois droites comme ceci

}{F7II-XVIII-019}\begin{equation}\tag{3.2.9}
\begin{tikzcd}[row sep=4.5em, column sep=4.5em]
	{} \& {} \arrow[dl, dash, shorten <=-9mm, shorten >=-9mm] \arrow[dr, dash, shorten <=-9mm, shorten >=-9mm] \& {} \& {} \\
	{} \arrow[rr, dash, shorten <=-11mm, shorten >=-11mm] \& {} \& {} \& {,}
\end{tikzcd}
\end{equation}















et a exactement trois points singuliers, tous quadratiques non-dégénérés. Cela fait ensemble douze points singuliers, et vérifie (3.2.7) dans cet exemple.


\SGApage{II}{270}

\underline{Proposition} 3.2.10. \underline{Soit} $X \hookrightarrow \mathbb{P}^{r}$ \underline{un plongement de Lefschetz} (XVII 2.3), \underline{et supposons que} $n = \dim X$ \underline{soit pair,} ou \underline{que} car. $k \neq 2$ . \underline{Alors tout pinceau de Lefschetz,} \underline{considéré comme droite dans} $\check{\mathbb{P}}^{r}$, \underline{coupe} $\check{X}$ \underline{en}

\SGAvar{\begin{equation}\tag{3.2.11}
(-1)^{\eta} [ \chi(X) + \chi(\Delta) - 2 \chi (X_{\overline{\eta}})]
\end{equation}
\underline{points, et le coupe transversalement} (EGA IV 17.13.7).

\underline{Démonstration.} Par hypothèse, toute singularité d'une fibre du pinceau de Lefschetz est quadratique non-dégénérée, de sorte que (3.2.4) s'applique pour donner que le nombre des fibres singulières (chacune contenant exactement un point singulier) dans le pinceau $D$ est le nombre (3.2.11). Or, la fibre $X_{\tau}$ est singulière si et seulement si $\tau \in \check{X} \cap D$ . Pour voir la transversalité de l'intersection, on remarque que le nombre (3.2.11) est \underline{indépendant} du choix particulier du pinceau de Lefschetz, et que, le pinceau générique étant un pinceau de Lefschetz, ce nombre (3.2.11) est le degré de $\check{X}$ (i.e. son degré comme hypersurface dans $\check{\mathbb{P}}^{r}$, qui est (par convention) $0$ si $\dim \check{X} \leq r-2$). Or, si une droite $D$ coupe une hypersurface de degré $d > 0$ en $d$ points distincts $x_{i}$ ($i \leq i \leq d$), l'intersection est forcément transversale. En effet, la multiplicité d'intersection $\mu_{i} = \mathrm{rg}(D \cap \check{X}|x_{i})$ de $D$ et $\check{X}$ en un $x_{i}$ est $\geq 1$, donc comme $\sum\limits^{d} \mu_{i} = d$ , on a $\mu_{i} = 1$, pour tout $i$, i.e. $D \cap \check{X}$ est lisse, ce qui, équivaut à la transversalité, comme on voit aussitôt.


}{\begin{equation}\tag{3.2.11}
(-1)^{n} [ \chi(X) + \chi(\Delta) - 2 \chi (X_{\overline{\eta}})]
\end{equation}
\underline{points, et le coupe transversalement} (EGA IV 17.13.7).





\underline{Démonstration.} Par hypothèse, toute singularité d'une fibre du pinceau de Lefschetz est quadratique non-dégénérée, de sorte que (3.2.4) s'applique pour donner que le nombre des fibres singulières (chacune contenant exactement un point singulier) dans le pinceau $D$ est le nombre (3.2.11). Or, la fibre $X_{\tau}$ est singulière si et seulement si $\tau \in \check{X} \cap D$ . Pour voir la transversalité de l'intersection, on remarque que le nombre (3.2.11) est \underline{indépendant} du choix particulier du pinceau de Lefschetz, et que, le pinceau générique étant un pinceau de Lefschetz, ce nombre (3.2.11) est le degré de $\check{X}$ (i.e. son degré comme hypersurface dans $\check{\mathbb{P}}^{r}$, qui est (par convention) $0$ si $\dim \check{X} \leq r-2$). Or, si une droite $D$ coupe une hypersurface de degré $d > 0$ en $d$ points distincts $x_{i}$ ($1 \leq i \leq d$), l'intersection est forcément transversale. En effet, la multiplicité d'intersection $\mu_{i} = \operatorname{lg}(D \cap \check{X}|x_{i})$ de $D$ et $\check{X}$ en un $x_{i}$ est $\geq 1$, donc comme $\sum\limits_{i=1}^{d} \mu_{i} = d$ , on a $\mu_{i} = 1$, pour tout $i$, i.e. $D \cap \check{X}$ est lisse, ce qui équivaut à la transversalité, comme on voit aussitôt.


}{F7II-XVIII-020}\SGApage{II}{271}

\SGAvar{\underline{Remarque} 3.2.12. Pour un pinceau de Lefschetz $D$ dont on suppose déjà qu'il coupe $\check{X}$ transversalement (condition vérifiéeen tous cas pour $D$ assez générale), le résultat 3.2.10 peut être considéré aussi comme la conjonction de XVII 5.7.2 et du fait que $\deg \varphi = 1$, contenu dans XVII 3.5. Bien entendu, le fait que \underline{tout} pinceau de Lefschetz $D$ coupe $\check{X}$ transversalement, établi ici via la formule d'Euler-Poincaré XVI 2.1 , peut aussi se vérifier directement de façon élémentaire, en utilisant le fait que l'axe de $D$ coupe $X$ transversalement.


}{\underline{Remarque} 3.2.12. Pour un pinceau de Lefschetz $D$ dont on suppose déjà qu'il coupe $\check{X}$ transversalement (condition vérifiée en tous cas pour $D$ assez générale), le résultat 3.2.10 peut être considéré aussi comme la conjonction de XVII 5.7.2 et du fait que $\deg \varphi = 1$, contenu dans XVII 3.5. Bien entendu, le fait que \underline{tout} pinceau de Lefschetz $D$ coupe $\check{X}$ transversalement, établi ici via la formule d'Euler-Poincaré XVI 2.1 , peut aussi se vérifier directement de façon élémentaire, en utilisant le fait que l'axe de $D$ coupe $X$ transversalement.


}{F7II-XVIII-021}\SGAsectiontoc{4. La cohomologie de l'éclatement associé à un pinceau}{4}
4. \underline{La cohomologie de l'éclatement associé à un pinceau}

4.1. Nous reprenons la situation 3.1, résumée par le diagramme commutatif (cf. 2.1.6.3) :

\SGAvar{\begin{equation}\tag{4.1.1}
\begin{tikzcd}[row sep=huge, column sep=huge]
X \times \mathbb{P}^{1} \& \widetilde{X} \arrow[l, hook', "f \times \rho"] \arrow[d, "f"] \arrow[dl, "\rho"'] \& \widetilde{\Delta} \arrow[l, hook', "i"] \arrow[d, "g"] \arrow[r, "{(g,\rho_{i})}", "\sim"'] \& \Delta \times \mathbb{P}^{1} \\
\mathbb{P}^{1} \& X \& \Delta \arrow[l, hook', "j"]
\end{tikzcd}\;.
\end{equation}





















}{\begin{equation}\tag{4.1.1}
\begin{tikzcd}[row sep=huge, column sep=huge]
X \times \mathbb{P}^{1} \& \widetilde{X} \arrow[l, hook', "f \times \rho"] \arrow[d, "f"] \arrow[dl, "\rho"'] \& \widetilde{\Delta} \arrow[l, hook', "i"] \arrow[d, "g"] \arrow[r, "{(g,\rho_{i})}", "\sim"'] \& \Delta \times \mathbb{P}^{1} \\
\mathbb{P}^{1} \& X \& \Delta \arrow[l, hook', "j"]
\end{tikzcd}\;.
\end{equation}





















}{}Commençons par préciser (2.2.2) dans ce cas.

\underline{Proposition} 4.2. \underline{Sous les hypothèses de} 3.1:

(4.2.1) \underline{Les homomorphismes de} $\mathbb{Z}_{\ell}$\underline{-modules}

\begin{equation*}
H^{\bullet}(\widetilde{X},\mathbb{Z}_{\ell}) \; \begin{array}{c} \xrightarrow{\;f_{*}\oplus(1-f^{*}f_{*})\;} \\ \xleftarrow{\;f^{*}+\mathrm{ind}\;} \end{array} \; H^{\bullet}(X,\mathbb{Z}_{\ell}) \oplus \mathrm{Ker}(f_{*}:H^{\bullet}(\widetilde{X},\mathbb{Z}_{\ell}) \longrightarrow H^{\bullet}(X,\mathbb{Z}_{\ell})) \quad 
\end{equation*}





\SGApage{II}{272}

 \underline{sont des isomorphismes, inverses l'un de l'autre.}

(4.2.2) \underline{Les homomorphismes de} $\mathbb{Z}_{\ell}$\underline{-modules}

\SGAvar{\begin{center}
\begin{tikzcd}[column sep=huge]
H^{\bullet}(\widetilde{X},\mathbb{Z}_{\ell}) \arrow[r, shift left=0.7ex, "f_{*}\oplus(g_{*}i^{*})"] \& H^{\bullet}(X,\mathbb{Z}_{\ell}) \oplus H^{\bullet-2}(\Delta,\mathbb{Z}_{\ell}(-1)) \arrow[l, shift left=0.7ex, "f^{*}+i_{*}g^{*}"]
\end{tikzcd}
\end{center}






}{\begin{center}
\begin{tikzcd}[column sep=huge]
H^{\bullet}(\widetilde{X},\mathbb{Z}_{\ell}) \arrow[r, shift left=0.7ex, "f_{*}\oplus(-g_{*}i^{*})"] \& H^{\bullet}(X,\mathbb{Z}_{\ell}) \oplus H^{\bullet-2}(\Delta,\mathbb{Z}_{\ell}(-1)) \arrow[l, shift left=0.7ex, "f^{*}+i_{*}g^{*}"]
\end{tikzcd}
\end{center}






}{F7II-XVIII-023}\underline{sont des isomorphismes, inverses l'un de l'autre.}

(4.2.3) \underline{La structure d'algèbre sur} $H^{\bullet}(X,\mathbb{Z}_{\ell}) \oplus H^{\bullet-2}(\Delta,\mathbb{Z}_{\ell}(-1))$ \underline{déduite par} ``\underline{transport de structure}'' \underline{par} (4.2.2) \underline{du cup-produit dans} $H^{\bullet}(\widetilde{X},\mathbb{Z}_{\ell})$, \underline{s'exprime, en termes des éléments}

\SGAvar{\begin{equation*}
a,b \in H^{\bullet}(X,\mathbb{Z}_{\ell}) \;,\; x,y \in H^{\bullet-2}(\Delta,\mathbb{Z}_{\ell}(-1) \quad ,
\end{equation*}
\underline{par les formules} :

\begin{equation}\tag{4.2.3.1}
(0 \oplus x) \wedge (0 \oplus y) \;=\; -\, j_{*}(xy) \oplus 2xy \; c_{1}(\underline{O}_{\Delta}(1)) \quad ,
\end{equation}
\begin{equation}\tag{4.2.3.2}
(a \oplus 0) \wedge (b \oplus 0 \;=\; ab \oplus 0 \quad ,
\end{equation}
\begin{equation}\tag{4.2.3.3}
(a \oplus 0) \wedge (0 \oplus y) \;=\; 0 \oplus j^{*}(a)y \quad ,
\end{equation}
\begin{equation}\tag{4.2.3.4}
(0 \oplus x) \wedge (b \oplus 0) \;=\; 0 \oplus xj^{*}(b) \quad .
\end{equation}
\underline{En particulier, pour}

\begin{equation*}
a \oplus x \in H^{\bullet}(X,\mathbb{Z}_{\ell}) \;\; \oplus H^{\bullet-2}(\Delta,\mathbb{Z}_{\ell}(-1)
\end{equation*}
\begin{equation*}
b \oplus y \in H^{2n-\bullet}(X,\mathbb{Z}_{\ell}) \oplus H^{2n-2-\bullet}(\Delta,\mathbb{Z}_{\ell}(-1)) \quad , 
\end{equation*}

}{\begin{equation*}
a,b \in H^{\bullet}(X,\mathbb{Z}_{\ell}) \;,\; x,y \in H^{\bullet-2}(\Delta,\mathbb{Z}_{\ell}(-1)) \quad ,
\end{equation*}
\underline{par les formules} :

\begin{equation}\tag{4.2.3.1}
(0 \oplus x) \wedge (0 \oplus y) \;=\; -\, j_{*}(xy) \oplus 2xy \; c_{1}(\underline{O}_{\Delta}(1)) \quad ,
\end{equation}
\begin{equation}\tag{4.2.3.2}
(a \oplus 0) \wedge (b \oplus 0) \;=\; ab \oplus 0 \quad ,
\end{equation}
\begin{equation}\tag{4.2.3.3}
(a \oplus 0) \wedge (0 \oplus y) \;=\; 0 \oplus j^{*}(a)y \quad ,
\end{equation}
\begin{equation}\tag{4.2.3.4}
(0 \oplus x) \wedge (b \oplus 0) \;=\; 0 \oplus xj^{*}(b) \quad .
\end{equation}
\underline{En particulier, pour}

\begin{equation*}
a \oplus x \in H^{\bullet}(X,\mathbb{Z}_{\ell}) \;\; \oplus H^{\bullet-2}(\Delta,\mathbb{Z}_{\ell}(-1))
\end{equation*}
\begin{equation*}
b \oplus y \in H^{2n-\bullet}(X,\mathbb{Z}_{\ell}) \oplus H^{2n-2-\bullet}(\Delta,\mathbb{Z}_{\ell}(-1)) \quad , 
\end{equation*}

}{F7II-XVIII-024}\SGApage{II}{273}

 \underline{on a}

\SGAvar{\begin{equation}\tag{4.2.3.5}
(a \oplus x) \wedge (b \oplus y) \;=\; ab - j_{*}(xy) \quad .
\end{equation}
Autrement dit, la matrice d'intersection totalede $H^{\bullet}(\widetilde{X},\ \mathbb{Z}_{\ell})$ s'écrit :

\begin{equation*}
\left(
\begin{array}{c|c}
\left(\begin{array}{l}\mbox{la matrice totale}\\ \mbox{des intersections sur } X\end{array}\right) & 0 \\ \hline
0 & -\left(\begin{array}{l}\mbox{la matrice totale}\\ \mbox{des intersections}\\ \mbox{sur } \Delta\end{array}\right)
\end{array}
\right) .
\end{equation*}
\underline{Démonstration.}  On a déjà démontré (4.2.1) ;(cf. 2.2.22)). Pour (4.2.2), remarquons d'abord que, $g : \widetilde{\Delta} \longrightarrow \Delta$ étant un fibré projectif de rang un, la suite

}{\begin{equation}\tag{4.2.3.5}
(a \oplus x) \wedge (b \oplus y) \;=\; ab - j_{*}(xy) \quad .
\end{equation}
Autrement dit, la matrice d'intersection totale de $H^{\bullet}(\widetilde{X},\ \mathbb{Z}_{\ell})$ s'écrit :

\begin{equation*}
\left(
\begin{array}{c|c}
\left(\begin{array}{l}\mbox{la matrice totale}\\ \mbox{des intersections sur } X\end{array}\right) & 0 \\ \hline
0 & -\left(\begin{array}{l}\mbox{la matrice totale}\\ \mbox{des intersections}\\ \mbox{sur } \Delta\end{array}\right)
\end{array}
\right) .
\end{equation*}
\underline{Démonstration.}  On a déjà démontré (4.2.1) ;(cf. (2.2.22)). Pour (4.2.2), remarquons d'abord que, $g : \widetilde{\Delta} \longrightarrow \Delta$ étant un fibré projectif de rang un, la suite

}{F7II-XVIII-025}\begin{equation}\tag{4.2.4}
0 \longrightarrow H^{\bullet}(\Delta,\mathbb{Z}_{\ell}) \stackrel{g^{*}}{\longrightarrow} H^{\bullet}(\widetilde{\Delta},\ \mathbb{Z}_{\ell}) \stackrel{g_{*}}{\longrightarrow} H^{\bullet-2}(\Delta,\mathbb{Z}_{\ell}(-1)) \longrightarrow 0
\end{equation}
est exacte (cf. (1.2.2)). Compte tenu de (2.2.26), on obtient le diagramme commutatif  de suites exactes

\SGAvar{\begin{equation}\tag{4.2.5}
\begin{tikzcd}[column sep=large, row sep=huge]
0 \arrow[r] \& H^{\bullet}(X,\mathbb{Z}_{\ell}) \arrow[r, "f^{*}"] \arrow[d, leftrightarrow, "j^{*}"] \& H^{\bullet}(\widetilde{X},\mathbb{Z}_{\ell}) \arrow[r, "g_{*}i^{*}"] \arrow[d, "i^{*}"] \& H^{\bullet-2}(\Delta,\mathbb{Z}_{\ell}(-1)) \arrow[r] \arrow[d, Rightarrow, no head] \& 0 \\
0 \arrow[r] \& H^{\bullet}(\Delta,\mathbb{Z}_{\ell}) \arrow[r, "g^{*}"] \& H^{\bullet}(\widetilde{\Delta},\mathbb{Z}_{\ell}) \arrow[r, "g_{*}"] \& H^{\bullet-2}(\Delta,\mathbb{Z}_{\ell}(-1)) \arrow[r] \& 0\;.
\end{tikzcd}
\end{equation}
















}{\begin{equation}\tag{4.2.5}
\begin{tikzcd}[column sep=large, row sep=huge]
0 \arrow[r] \& H^{\bullet}(X,\mathbb{Z}_{\ell}) \arrow[r, "f^{*}"] \arrow[d, "j^{*}"] \& H^{\bullet}(\widetilde{X},\mathbb{Z}_{\ell}) \arrow[r, "g_{*}i^{*}"] \arrow[d, "i^{*}"] \& H^{\bullet-2}(\Delta,\mathbb{Z}_{\ell}(-1)) \arrow[r] \arrow[d, Rightarrow, no head] \& 0 \\
0 \arrow[r] \& H^{\bullet}(\Delta,\mathbb{Z}_{\ell}) \arrow[r, "g^{*}"] \& H^{\bullet}(\widetilde{\Delta},\mathbb{Z}_{\ell}) \arrow[r, "g_{*}"] \& H^{\bullet-2}(\Delta,\mathbb{Z}_{\ell}(-1)) \arrow[r] \& 0\;.
\end{tikzcd}
\end{equation}
















}{F7II-XVIII-026}\SGApage{II}{274}

Compte tenu du scindage (4.2.1), ceci donne un isomorphisme

\SGAvar{\begin{equation}\tag{4.2.6}
f_* \oplus g_* \ i^* : H^{\cdot}(\widetilde{X},\mathbb{Z}_\ell) \;\stackrel{\sim}{\longrightarrow}\; H^{\cdot}(X,\mathbb{Z}_\ell) \oplus H^{\cdot-2}(\Delta,\mathbb{Z}_\ell(-1)) \ .
\end{equation}
Le transposé de (4.2.5) donne un diagramme commutatif de suites exactes (après le décalage $. \stackrel{\sim}{\longrightarrow} 2n-.$)

\begin{equation}\tag{4.2.7}
\begin{tikzcd}[column sep=large, row sep=large]
0 \& H^{\cdot}(X,\mathbb{Z}_\ell) \arrow[l] \& H^{\cdot}(\widetilde{X},\mathbb{Z}_\ell) \arrow[l, "f_*"'] \& H^{\cdot-2}(\Delta,\mathbb{Z}_\ell(-1)) \arrow[l, "i_*g^*"'] \& 0 \arrow[l] \\
0 \& H^{\cdot-2}(\Delta,\mathbb{Z}_\ell(-2)) \arrow[l] \arrow[u, "j_*"'] \& H^{\cdot-2}(\widetilde{\Delta},\mathbb{Z}_\ell(-1)) \arrow[l, "g_*"'] \arrow[u, "i_*"'] \& H^{\cdot-2}(\Delta,\mathbb{Z}_\ell(-1)) \arrow[l, "g_*"'] \arrow[u, equal] \& 0 \arrow[l]
\end{tikzcd}
\end{equation}

















}{\begin{equation}\tag{4.2.6}
f_* \oplus (-g_* \ i^*) : H^{\cdot}(\widetilde{X},\mathbb{Z}_\ell) \;\stackrel{\sim}{\longrightarrow}\; H^{\cdot}(X,\mathbb{Z}_\ell) \oplus H^{\cdot-2}(\Delta,\mathbb{Z}_\ell(-1)) \ .
\end{equation}
Le transposé de (4.2.5) donne un diagramme commutatif de suites exactes (après le décalage $. \stackrel{\sim}{\longrightarrow} 2n-.$)

\begin{equation}\tag{4.2.7}
\begin{tikzcd}[column sep=large, row sep=large]
0 \& H^{\cdot}(X,\mathbb{Z}_\ell) \arrow[l] \& H^{\cdot}(\widetilde{X},\mathbb{Z}_\ell) \arrow[l, "f_*"'] \& H^{\cdot-2}(\Delta,\mathbb{Z}_\ell(-1)) \arrow[l, "i_*g^*"'] \& 0 \arrow[l] \\
0 \& H^{\cdot-4}(\Delta,\mathbb{Z}_\ell(-2)) \arrow[l] \arrow[u, "j_*"'] \& H^{\cdot-2}(\widetilde{\Delta},\mathbb{Z}_\ell(-1)) \arrow[l, "g_*"'] \arrow[u, "i_*"'] \& H^{\cdot-2}(\Delta,\mathbb{Z}_\ell(-1)) \arrow[l, "g^*"'] \arrow[u, equal] \& 0 \arrow[l]
\end{tikzcd}
\end{equation}

















}{F7II-XVIII-027}et on obtient un isomorphisme

\SGAvar{\begin{equation}\tag{4.2.8}
H^{\cdot}(X,\mathbb{Z}_\ell) \oplus H^{\cdot-2}(\Delta,\mathbb{Z}_\ell(-1)) \;\stackrel{f^*+i_*g^*}{\longrightarrow}\; H^{\cdot}(\widetilde{X},\mathbb{Z}_\ell) \qquad .
\end{equation}
Il reste à démontrer que le composé

\begin{equation}\tag{4.2.9}
H^{\cdot}(X,\mathbb{Z}_\ell) \oplus H^{\cdot-2}(\Delta,\mathbb{Z}_\ell(-1)) \;\stackrel{f^*+i_*g^*}{\longrightarrow}\; H^{\cdot}(\widetilde{X},\mathbb{Z}_\ell) \;\stackrel{f_*\oplus(g_*i^*)}{\longrightarrow}\; H^{\cdot}(X,\mathbb{Z}_\ell)\oplus H^{\cdot-2}(\Delta,\mathbb{Z}_\ell(-1))
\end{equation}
est l'identité. Or, on a

}{\begin{equation}\tag{4.2.8}
H^{\cdot}(X,\mathbb{Z}_\ell) \oplus H^{\cdot-2}(\Delta,\mathbb{Z}_\ell(-1)) \;\stackrel{f^*+i_*g^*}{\longrightarrow}_{\sim}\; H^{\cdot}(\widetilde{X},\mathbb{Z}_\ell) \qquad .
\end{equation}
Il reste à démontrer que le composé

\begin{equation}\tag{4.2.9}
H^{\cdot}(X,\mathbb{Z}_\ell) \oplus H^{\cdot-2}(\Delta,\mathbb{Z}_\ell(-1)) \;\stackrel{f^*+i_*g^*}{\longrightarrow}\; H^{\cdot}(\widetilde{X},\mathbb{Z}_\ell) \;\stackrel{f_*\oplus(-g_*i^*)}{\longrightarrow}\; H^{\cdot}(X,\mathbb{Z}_\ell)\oplus H^{\cdot-2}(\Delta,\mathbb{Z}_\ell(-1))
\end{equation}
est l'identité. Or, on a

}{F7II-XVIII-028}\begin{equation*}
f_*f^* = \text{id sur } H^{\cdot}(X,\mathbb{Z}_\ell) \ , \ \text{par } (2.2.18)
\end{equation*}
\begin{equation*}
g_*i^*f^* = 0 \qquad\qquad\qquad \text{par } (4.2.5)
\end{equation*}
\begin{equation*}
f_*i_*g^* = 0 \qquad\qquad\qquad \text{par } (4.2.7) \qquad ,
\end{equation*}
de sorte qu'il suffit de démontrer

\SGAvar{\begin{equation}\tag{4.2.10}
g_*i^*i_*g^*(X) = X \ \text{pour} \ X \in H^{\cdot}(\Delta,\mathbb{Z}_\ell(-1)) \qquad .
\end{equation}

}{\begin{equation}\tag{4.2.10}
-g_*i^*i_*g^*(x) = x \ \text{pour} \ x \in H^{\cdot}(\Delta,\mathbb{Z}_\ell(-1)) \qquad .
\end{equation}

}{F7II-XVIII-029}\SGApage{II}{275}

Pour vérifier (4.2.10), remarquons d'abord que $i : \widetilde{\Delta} \hookrightarrow \widetilde{X}$ étant l'inclusion d'une sous-variété lisse de codimension un, on a [SGA 5 VII]

\SGAvar{\begin{equation}\tag{4.2.11}
i^*i_*(x) \;=\; c_1(\underline{N}_{\widetilde{\Delta}/\widetilde{X}}) \ x
\end{equation}
\begin{equation}\tag{4.2.12}
i_*(x)i_*(y) \;=\; i_*(xy \ c_1(\underline{N}_{\widetilde{\Delta}/\widetilde{X}})) \qquad .
\end{equation}
Ecrivons $\xi_X = \underline{O}_X(1)$, $\xi_\Delta = \underline{O}_\Delta(1)$ , $\xi_{\mathbb{P}^1} = \underline{O}_{\mathbb{P}^1}(1)$, etc . En vertu de EGA II 8.1.7, on a, si $\widetilde{\mathfrak{J}}$ est l'Idéal $\widetilde{\Delta}$ dans $\widetilde{X}$

}{\begin{equation}\tag{4.2.11}
i^*i_*(x) \;=\; c_1(\underline{N}_{\widetilde{\Delta}/\widetilde{X}}) \ x
\end{equation}
\begin{equation}\tag{4.2.12}
i_*(x)i_*(y) \;=\; i_*(xy \ c_1(\underline{N}_{\widetilde{\Delta}/\widetilde{X}})) \qquad .
\end{equation}
Ecrivons $\xi_X = \underline{O}_X(1)$, $\xi_\Delta = \underline{O}_\Delta(1)$ , $\xi_{\mathbb{P}^1} = \underline{O}_{\mathbb{P}^1}(1)$, etc . En vertu de EGA II 8.1.7, on a, si $\widetilde{\mathfrak{J}}$ est l'Idéal de $\widetilde{\Delta}$ dans $\widetilde{X}$

}{F7II-XVIII-030}\begin{equation}\tag{4.2.13}
\underline{\check{N}}_{\widetilde{\Delta}/\widetilde{X}} \;\simeq\; \widetilde{\mathfrak{J}}/\widetilde{\mathfrak{J}}^2 \;\simeq\; \widetilde{\mathfrak{J}} \otimes_{\underline{O}_{\widetilde{X}}} O_{\widetilde{\Delta}} \;\simeq\; \underline{O}_{\widetilde{X}}(1) \otimes_{\underline{O}_{\widetilde{X}}} \underline{O}_{\widetilde{\Delta}} \;=\; \underline{O}_{\widetilde{\Delta}}(1) \qquad ,
\end{equation}
où $\widetilde{\Delta}$ est interprété comme $\mathbb{P}(\mathfrak{J}/\mathfrak{J}^2)$, $\mathfrak{J}$ étant l'Idéal de $\Delta$ dans $X$ . Or on a un isomorphisme canonique ( $\Delta$ étant représenté comme intersection de $X$ avec deux hyperplans  donnés dans $\check{\mathbb{P}}$ ) ,

\begin{equation*}
\mathfrak{J}/\mathfrak{J}^2 \;\simeq\; \underline{O}^2_\Delta(-1) \qquad ,
\end{equation*}
qui donne naissance à l'isomorphisme

\begin{equation*}
\widetilde{\Delta} = \mathbb{P}(\mathfrak{J}/\mathfrak{J}^2) \;\simeq\; \mathbb{P} \ (\underline{O}^2_\Delta(-1)) \;\stackrel{\alpha}{\simeq}\; \mathbb{P}(\underline{O}^2_\Delta) \;\simeq\; \Delta \times \mathbb{P}^1
\end{equation*}
où $\alpha$ est l'isomorphisme canonique EGA II 4.1.4. Par loc. cit. cet isomorphisme transforme

\SGAvar{\begin{equation*}
\underline{O}_{\widetilde{\Delta}}(1) = \underline{O}_{\mathbb{P}(\mathfrak{J}/\mathfrak{J}^2)}(1) = \underline{O}_{\mathbb{P}(\underline{O}^2_\Delta(-1))}(1) \ \text{en} \ \underline{O}_{\mathbb{P}(\underline{O}^2_\Delta)}(1) \otimes_{\underline{O}^2_\Delta} \underline{O}_\Delta(-1) = \underline{O}_{\mathbb{P}^1}(1) \otimes_K \underline{O}_\Delta(-1),
\end{equation*}
de sorte que (4.2.13) s'écrit

\begin{equation}\tag{4.2.14}
\check{N}_{\widetilde{\Delta}/\widetilde{X}} \;\simeq\; (\rho i)^*(\underline{O}_{\mathbb{P}}(1)) \otimes g^*(\underline{O}_\Delta(-1)) \qquad ,
\end{equation}
d'où en prenant le $c_1$ des deux membres, le 


}{\begin{equation*}
\underline{O}_{\widetilde{\Delta}}(1) = \underline{O}_{\mathbb{P}(\mathfrak{J}/\mathfrak{J}^2)}(1) = \underline{O}_{\mathbb{P}(\underline{O}^2_\Delta(-1))}(1) \ \text{en} \ \underline{O}_{\mathbb{P}(\underline{O}^2_\Delta)}(1) \otimes_{\underline{O}_\Delta} \underline{O}_\Delta(-1) = \underline{O}_{\mathbb{P}^1}(1) \otimes_k \underline{O}_\Delta(-1),
\end{equation*}
de sorte que (4.2.13) s'écrit

\begin{equation}\tag{4.2.14}
\check{N}_{\widetilde{\Delta}/\widetilde{X}} \;\simeq\; (\rho i)^*(\underline{O}_{\mathbb{P}^1}(1)) \otimes g^*(\underline{O}_\Delta(-1)) \qquad ,
\end{equation}
d'où en prenant le $c_1$ des deux membres, le 


}{F7II-XVIII-031}\SGApage{II}{276}

\underline{Lemme} 4.3. \underline{On a}

\begin{equation*}
c_1(\underline{N}_{\widetilde{\Delta}/\widetilde{X}}) = g^*(\xi_\Delta) - (\rho i)^*(\xi_{\mathbb{P}^1}) \qquad .
\end{equation*}
Reprenons la démonstration de (4.2.10). D'après (4.2.11) et (4.3), il suffit de démontrer

\SGAvar{\begin{equation*}
g_*(g^*(x)[(\rho i)^* (c_1(\xi_{\mathbb{P}^1})) - g^*(c_1(\xi_\Delta)]) = x \qquad .
\end{equation*}
Or, d'après (4.2.4), $g_*g^* = 0$, et, compte tenu de ce que

}{\begin{equation*}
g_*\bigl(g^*(x)\bigl[(\rho i)^* (c_1(\xi_{\mathbb{P}^1})) - g^*(c_1(\xi_\Delta))\bigr]\bigr) = x \qquad .
\end{equation*}
Or, d'après (4.2.4), $g_*g^* = 0$, et, compte tenu de ce que

}{F7II-XVIII-032}\begin{center}
\begin{tikzcd}
\widetilde{\Delta} \arrow[r, "g \times \rho_i"{yshift=7pt}, "\sim"] \& \Delta \times \mathbb{P}^1 \quad ,
\end{tikzcd}
\end{center}








on a

\SGAvar{\begin{equation}\tag{4.3.1}
g_*(g^*(x) (\rho i)^*(c_1(\xi_{\mathbb{P}^1})) = x
\end{equation}
\begin{equation}\tag{4.3.2}
g_*(g^*(x) c_1(\underline{N}_{\widetilde{\Delta}/\widetilde{X}})) = - x \qquad .
\end{equation}
Calculons maintenant la structure d'algèbre sur $H^{\bullet}(X,\mathbb{Z}_{\ell}) \oplus H^{\bullet-2}(\Delta, \mathbb{Z}_{\ell}(-1))$, déduite du cup-produit par transport de structure via (4.2.2). Soient $a$, $b$ des éléments de $H^{\bullet}(X,\mathbb{Z}_{\ell})$, $x,y$ des éléments de $H^{\bullet-2}(\Delta,\mathbb{Z}_{\ell}(-1))$. Commençons par (4.2.3.1).

\begin{equation*}
(0 \oplus x) \Lambda (0 \oplus y) = ? \oplus ?'
\end{equation*}
\begin{equation*}
\begin{aligned}
? &= f_*((i_*g^*(x))(i_*g^*(y)) \\
&= f_*(i_*(g^*(xy)\, c_1\, (N_{\widetilde{\Delta}/\widetilde{X}}))) \qquad \text{par (4.2.12)} \\
&= j_*g_*(g^*(xy)\, c_1\, (N_{\widetilde{\Delta}/\widetilde{X}})) \\
&= j_*(-xy) \qquad\qquad \text{par (4.3.2)} \;.
\end{aligned}
\end{equation*}

}{\begin{equation}\tag{4.3.1}
g_*\bigl(g^*(x) (\rho i)^*(c_1(\xi_{\mathbb{P}^1}))\bigr) = x
\end{equation}
\begin{equation}\tag{4.3.2}
g_*(g^*(x) c_1(\underline{N}_{\widetilde{\Delta}/\widetilde{X}})) = - x \qquad .
\end{equation}
Calculons maintenant la structure d'algèbre sur $H^{\bullet}(X,\mathbb{Z}_{\ell}) \oplus H^{\bullet-2}(\Delta, \mathbb{Z}_{\ell}(-1))$, déduite du cup-produit par transport de structure via (4.2.2). Soient $a$, $b$ des éléments de $H^{\bullet}(X,\mathbb{Z}_{\ell})$, $x,y$ des éléments de $H^{\bullet-2}(\Delta,\mathbb{Z}_{\ell}(-1))$. Commençons par (4.2.3.1).

\begin{equation*}
(0 \oplus x) \wedge (0 \oplus y) = ? \oplus ?'
\end{equation*}
\begin{equation*}
\begin{aligned}
? &= f_*\bigl((i_*g^*(x))(i_*g^*(y))\bigr) \\
&= f_*(i_*(g^*(xy)\, c_1\, (N_{\widetilde{\Delta}/\widetilde{X}}))) \qquad \text{par (4.2.12)} \\
&= j_*g_*(g^*(xy)\, c_1\, (N_{\widetilde{\Delta}/\widetilde{X}})) \\
&= j_*(-xy) \qquad\qquad \text{par (4.3.2)} \;.
\end{aligned}
\end{equation*}

}{F7II-XVIII-033}\SGApage{II}{277}

\SGAvar{\begin{equation*}
\begin{aligned}
?' &= -g_*i^*(i_*g^*(x) \wedge i_*g^*(y)) \\
&= -g_*i^*(i_*(g^*(xy)\, c_1(\underline{N}_{\widetilde{\Delta}/\widetilde{X}}))) \qquad \text{par (4.3.2)} \\
&= -g_*(g^*(xy)(c_1(N_{\widetilde{\Delta}/\widetilde{X}}))^2) \qquad \text{par (4.3.2)} \quad ,
\end{aligned}
\end{equation*}
or,

\begin{equation*}
\begin{aligned}
(c_1(\underline{N}_{\widetilde{\Delta}/\widetilde{X}}))^2 &= (g^*c_1(\underline{O}_{\Delta(1)}) - (\rho i)^* (c_1(\xi_{\mathbb{P}^1})))^2 \qquad \text{par (4.3)} \\
&= g^*((c_1(\xi_\Delta))^2) - 2g^*(c_1(\xi_\Delta))(\rho i)^*(c_1(\xi_{\mathbb{P}^1})) \\
&\quad ( + (\rho i)^* ((c_1( \xi_{\mathbb{P}^1}))^2) = 0 ) \quad ,
\end{aligned}
\end{equation*}
\begin{equation*}
\begin{aligned}
?' &= - g_*(g^*(xy)g^*(c_1(\xi_\Delta)^2)) + g_*(2g^*(xy)g^*(c_1(\xi_\Delta)(\rho i)^*(c_1(\xi_{\mathbb{P}^1})) \\
? &= 0 + 2xy\, c_1(\xi_\Delta) \qquad \text{par (4.2.4) et (4.3.1)} \quad ,
\end{aligned}
\end{equation*}
ce qui établit (4.2.3.1). Quant à (4.2.3.2) : $(a \oplus 0) \Lambda (b \oplus 0) = ab \oplus 0$, ceci exprime que $f^*$ est un homomorphisme d'algèbres.

Prouvons (4.2.3.3). Or, la formule de projection implique $\mathrm{Ker}\, f_* \cap \mathrm{image}\, f^* \subset \mathrm{Ker}\, f_*$ , de sorte que

\begin{equation*}
(a \oplus 0) \Lambda (0 \oplus y) = 0 \oplus ? \qquad ,
\end{equation*}
où

}{\begin{equation*}
\begin{aligned}
?' &= -g_*i^*(i_*g^*(x) \wedge i_*g^*(y)) \\
&= -g_*i^*(i_*(g^*(xy)\, c_1(\underline{N}_{\widetilde{\Delta}/\widetilde{X}}))) \qquad \text{par (4.2.12)} \\
&= -g_*(g^*(xy)(c_1(N_{\widetilde{\Delta}/\widetilde{X}}))^2) \qquad \text{par (4.2.11)} \quad ,
\end{aligned}
\end{equation*}
or,

\begin{equation*}
\begin{aligned}
(c_1(\underline{N}_{\widetilde{\Delta}/\widetilde{X}}))^2 &= (g^*c_1(\underline{O}_{\Delta}(1)) - (\rho i)^* (c_1(\xi_{\mathbb{P}^1})))^2 \qquad \text{par (4.3)} \\
&= g^*((c_1(\xi_\Delta))^2) - 2g^*(c_1(\xi_\Delta))(\rho i)^*(c_1(\xi_{\mathbb{P}^1})) \\
&\quad ( + (\rho i)^* ((c_1( \xi_{\mathbb{P}^1}))^2) = 0 ) \quad ,
\end{aligned}
\end{equation*}
\begin{equation*}
\begin{aligned}
?' &= - g_*\bigl(g^*(xy)g^*(c_1(\xi_\Delta)^2)\bigr)
      + g_*\bigl(2g^*(xy)g^*(c_1(\xi_\Delta))(\rho i)^*(c_1(\xi_{\mathbb{P}^1}))\bigr) \\
?' &= 0 + 2xy\, c_1(\xi_\Delta) \qquad \text{par (4.2.4) et (4.3.1)} \quad ,
\end{aligned}
\end{equation*}
ce qui établit (4.2.3.1). Quant à (4.2.3.2) : $(a \oplus 0) \wedge (b \oplus 0) = ab \oplus 0$, ceci exprime que $f^*$ est un homomorphisme d'algèbres.



Prouvons (4.2.3.3). Or, la formule de projection implique $\mathrm{Ker}\, f_* \cdot \mathrm{image}\, f^* \subset \mathrm{Ker}\, f_*$ , de sorte que

\begin{equation*}
(a \oplus 0) \wedge (0 \oplus y) = 0 \oplus ? \qquad ,
\end{equation*}
où

}{F7II-XVIII-034}\begin{equation}\tag{4.3.3}
? = - g_*i^*(f^*(a)\, i_*g^*(y)) 
\end{equation}

\SGApage{II}{278}

\SGAvar{\begin{equation*}
 = - g_*(i^*f^*(a) \; g^*(y) \; c_1(\underline{N}_{\widetilde{\Delta}/\widetilde{X}})) \qquad \mathrm{par} \; (4.3.2)
\end{equation*}
\begin{equation*}
= - g_*(g^*(j^*(a)) \; g^*(y) \; c_1(\underline{N}_{\widetilde{\Delta}/\widetilde{X}}))
\end{equation*}
\begin{equation*}
= \quad j^*(a).y
\end{equation*}
ce qui établit (4.2.3.3). On remarque enfin que (4.2.3.4) se déduit de (4.2.3.3) par l'anticommutativité du cup-produit.


}{\begin{equation*}
 = - g_*\bigl(i^*f^*(a) \; g^*(y) \; c_1(\underline{N}_{\widetilde{\Delta}/\widetilde{X}})\bigr) \qquad \mathrm{par} \; (4.2.11)
\end{equation*}
\begin{equation*}
= - g_*(g^*(j^*(a)) \; g^*(y) \; c_1(\underline{N}_{\widetilde{\Delta}/\widetilde{X}}))
\end{equation*}
\begin{equation*}
= \quad j^*(a).y
\end{equation*}
ce qui établit (4.2.3.3). On remarque enfin que (4.2.3.4) se déduit de (4.2.3.3) par l'anticommutativité du cup-produit.


}{F7II-XVIII-035}\SGAsectiontoc{5. Etude cohomologique du morphisme $\rho : \widetilde{X} \longrightarrow \mathbb{P}^1$ associé à un pinceau}{5}
5. \underline{Etude cohomologique du morphisme} $\rho : \widetilde{X} \longrightarrow \mathbb{P}^1$ \underline{associé à un pinceau}


5.1. \underline{L'homomorphisme de restriction} $H^*(\widetilde{X}) \longrightarrow H^*(Y)$.

Désignons par $Y$ une des fibres lisses de $\rho : \widetilde{X} \longrightarrow \mathbb{P}^1$, disons $Y = \rho^{-1}(t_o) = X.H_{t_o}$ . On a le diagramme commutatif

\begin{center}
\begin{tikzcd}[column sep=large, row sep=huge]
\mathbb{P}^1 \& {} \& \widetilde{X} \arrow[ll, "\rho"'] \arrow[d, "f"] \& {} \& \widetilde{\Delta} \arrow[ll, hook, "i"'] \arrow[d, "g"] \arrow[r, phantom, "\simeq"] \& \Delta \times \mathbb{P}^1 \\
\mathrm{Sp}(k)=\{t_o\} \arrow[u, hook] \& {} \& X \& {} \& \Delta \arrow[ll, hook, "j"'] \arrow[ur, hook, "\ell = \mathrm{id} \times \{t_o\}"'] \arrow[dlll, hook, "h"] \& {} \\
{} \& Y \arrow[ul] \arrow[uur, "k"] \arrow[ur, "m"'] \& {} \& {} \& {} \&
{} \quad . \end{tikzcd}
\end{center}


























\SGApage{II}{279}

\underline{Proposition} 5.1.1. \underline{Sous les hypothèses de} (3.1), $Y$ \underline{désignant une fibre lisse} :

(5.1.2) \quad \underline{L'homomorphisme restriction}

\begin{equation*}
k^* : H^\cdot(\widetilde{X},\mathbb{Z}_\ell) \longrightarrow H^\cdot(Y,\mathbb{Z}_\ell)
\end{equation*}
\underline{s'exprime} (\underline{via} (4.2.2)) \underline{par}

\SGAvar{\begin{equation*}
H^\cdot(X,\mathbb{Z}_\ell) \oplus H^{\cdot-2}(\Delta,\mathbb{Z}_\ell(-1)) \xrightarrow{\;m^*+h^*\;} H^\cdot(Y,\mathbb{Z}_\ell) \qquad .
\end{equation*}
(5.1.3) \quad \underline{La flèche de Gysin}

}{\begin{equation*}
H^\cdot(X,\mathbb{Z}_\ell) \oplus H^{\cdot-2}(\Delta,\mathbb{Z}_\ell(-1)) \xrightarrow{\;m^*+h_*\;} H^\cdot(Y,\mathbb{Z}_\ell) \qquad .
\end{equation*}
(5.1.3) \quad \underline{La flèche de Gysin}

}{F7II-XVIII-036}\begin{equation*}
k_* : H^{\cdot-2}(Y,\mathbb{Z}_\ell(-1)) \longrightarrow H^\cdot(\widetilde{X},\mathbb{Z}_\ell)
\end{equation*}
\underline{s'exprime par}

\begin{equation*}
H^{\cdot-2}(Y,\mathbb{Z}_\ell(-1)) \xrightarrow{\;m_* \oplus (-h^*)\;} H^\cdot(X,\mathbb{Z}_\ell) \oplus H^{\cdot-2}(\Delta,\mathbb{Z}_\ell(-1)).
\end{equation*}
\underline{Démonstration.} Les sous-schémas lisses $\widetilde{\Delta}$ et $Y$ de $\widetilde{X}$ se coupent transversalement en $\Delta$ , autrement dit, le carré de schémas lisses

\begin{equation}\tag{5.1.4}
\begin{tikzcd}[column sep=large, row sep=large]
\Delta \arrow[r, "\ell"] \arrow[d, hook, "h"] \& \widetilde{\Delta} \arrow[d, hook, "i"] \\
Y \arrow[r, "k"'] \& \widetilde{X}
\end{tikzcd}
\end{equation}











est cartésien, ce qui implique la formule [SGA 5, IV \hspace{2.2cm} ]

\begin{equation}\tag{5.1.5}
h_*\ell^* = k^*i_* \quad , \quad \ell_*h^* = i^*k_* \qquad .
\end{equation}

\SGApage{II}{280}

Le composé

\SGAvar{\begin{equation*}
H^{\bullet}(X,\mathbb{Z}_{\ell}) \oplus H^{\bullet-2}(\Delta,\mathbb{Z}_{\ell}(-1)) \stackrel{f^*+i_*g^*}{\longrightarrow} H^{\bullet}(\widetilde{X},\mathbb{Z}_{\ell}) \stackrel{k^*}{\longrightarrow} H^{\bullet}(Y,\mathbb{Z}_{\ell})
\end{equation*}
est la somme de $k^*f^* = m^*$ et de $k^*i_*g^*$ $= h_*\ell^*g^*$ par (5.1.5)), et $g.\ell = \mathrm{id}_{\Delta}$ , ce qui donne (5.1.2).

}{\begin{equation*}
H^{\bullet}(X,\mathbb{Z}_{\ell}) \oplus H^{\bullet-2}(\Delta,\mathbb{Z}_{\ell}(-1)) \stackrel{f^*+i_*g^*}{\longrightarrow} H^{\bullet}(\widetilde{X},\mathbb{Z}_{\ell}) \stackrel{k^*}{\longrightarrow} H^{\bullet}(Y,\mathbb{Z}_{\ell})
\end{equation*}
est la somme de $k^*f^* = m^*$ et de $k^*i_*g^*$ $= h_*\ell^*g^*$ par (5.1.5), et $g.\ell = \mathrm{id}_{\Delta}$ , ce qui donne (5.1.2).

}{F7II-XVIII-037}Le composé

\SGAvar{\begin{equation*}
H^{\bullet-2}(Y,\mathbb{Z}_{\ell(-1)}) \stackrel{k_*}{\longrightarrow} H^{\bullet}(\widetilde{X},\ \mathbb{Z}_{\ell}) \stackrel{f_*\oplus(g_*i^*)}{\longrightarrow} H^{\bullet}(X,\mathbb{Z}_{\ell}) \oplus H^{\bullet-2}(\Delta,\mathbb{Z}_{\ell}(-1))
\end{equation*}
est la somme de $f_*k_* = m_*$ et de $-g_*i^*k_* = g_*\ell_*h^* = -h^*$ ,
ce qui donne (5.1.3).

}{\begin{equation*}
H^{\bullet-2}(Y,\mathbb{Z}_{\ell}(-1)) \stackrel{k_*}{\longrightarrow} H^{\bullet}(\widetilde{X},\ \mathbb{Z}_{\ell}) \stackrel{f_*\oplus(-g_*i^*)}{\longrightarrow} H^{\bullet}(X,\mathbb{Z}_{\ell}) \oplus H^{\bullet-2}(\Delta,\mathbb{Z}_{\ell}(-1))
\end{equation*}
est la somme de $f_*k_* = m_*$ et de $-g_*i^*k_* = -g_*\ell_*h^* = -h^*$ ,
ce qui donne (5.1.3).

}{F7II-XVIII-038}\underline{Corollaire} 5.1.6. \underline{Sous les hypothèses de} 5.1.1, \underline{pour} $i \leq n-1$, \underline{les images de}

\begin{equation*}
m^* : H^i(X,\mathbb{Z}_{\ell}) \longrightarrow H^i(Y,\mathbb{Z}_{\ell})
\end{equation*}
\underline{et de}

\begin{equation*}
k^* : H^i(\widetilde{X},\mathbb{Z}_{\ell}) \longrightarrow H^i(Y,\mathbb{Z}_{\ell})
\end{equation*}
\underline{sont égales}.

\underline{Démonstration.} D'après le ``théorème de Lefschetz faible'' sur les sections hyperplanes (SGA 5 VII 7.1  ), pour $i \leq n-1$ la flèche

\SGAvar{\begin{equation}\tag{5.1.7}
H^{i-2}(X,\mathbb{Z}_{\ell}(-1)) \stackrel{j^*}{\longrightarrow} H^{i-2}(\Delta,\mathbb{Z}_{\ell}(-1))
\end{equation}
est un isomorphisme, de sorte que pour $i \leq n-1$, l'image de $h^*$ est égale, grâce à (5.1.2), à la somme de celle de $m^*$ et de celle du composé

\begin{equation*}
H^{i-2}(X,\mathbb{Z}_{\ell}(-1)) \stackrel{j^*}{\longrightarrow} H^{i-2}(\Delta,\mathbb{Z}_{\ell}(-1)) \stackrel{h^*}{\longrightarrow} H^i(Y,\mathbb{Z}_{\ell}) \qquad .
\end{equation*}

}{\begin{equation}\tag{5.1.7}
H^{i-2}(X,\mathbb{Z}_{\ell}(-1)) \stackrel{j^*}{\longrightarrow} H^{i-2}(\Delta,\mathbb{Z}_{\ell}(-1))
\end{equation}
est un isomorphisme, de sorte que pour $i \leq n-1$, l'image de $k^*$ est égale, grâce à (5.1.2), à la somme de celle de $m^*$ et de celle du composé

\begin{equation*}
H^{i-2}(X,\mathbb{Z}_{\ell}(-1)) \stackrel{j^*}{\longrightarrow} H^{i-2}(\Delta,\mathbb{Z}_{\ell}(-1)) \stackrel{h_*}{\longrightarrow} H^i(Y,\mathbb{Z}_{\ell}) \qquad .
\end{equation*}

}{F7II-XVIII-039}\SGApage{II}{281}

Or, par la formule de projection

\SGAvar{\begin{equation*}
h_*j^*(x) = h_*(h^*(m^*(x)) = m^*(x)h_*(1) = \xi_Y\ m^*(x)
\end{equation*}
\begin{equation*}
= m^*(x\ \xi_X) \quad \mbox{puisque} \quad m^*(\xi_X) = \xi_Y \qquad ,
\end{equation*}
d'où l'inclusion $\mathrm{Im}\ k^i \subset \mathrm{Im}\ m^i$ , cqfd.


}{\begin{equation*}
h_*j^*(x) = h_*(h^*(m^*(x))) = m^*(x)h_*(1) = \xi_Y\ m^*(x)
\end{equation*}
\begin{equation*}
= m^*(x\ \xi_X) \quad \mbox{puisque} \quad m^*(\xi_X) = \xi_Y \qquad ,
\end{equation*}
d'où l'inclusion $\mathrm{Im}\ k^i \subset \mathrm{Im}\ m^i$ , cqfd.


}{F7II-XVIII-040}5.2.  \underline{Le ``théorème de Lefschetz vache}'' (LV)

5.2.1.  Soit $S$ projectif de dimension $n$, sur le corps $k$ (alg. clos) muni d'un plongement projectif $S \hookrightarrow \mathbb{P}$ . Désignons par

\SGAvar{\begin{equation*}
u \in H^2(S,\mathbb{Z}_{\ell}(1))
\end{equation*}
la ``classe de cohomologie d'une section hyperplane '' (de sorte que $u$ s'obtient comme l'image réciproque par le plongement $S \hookrightarrow \mathbb{P}$ de la première classe de Chern de $\underline{O}_{\mathbb{P}}$ (1), classe qui se trouve dans $H^2(\mathbb{P}\ \ , \mathbb{Z}_{\ell}(1))$, et désignons par

\begin{equation*}
L_S : H^{\bullet}(S,\mathbb{Q}_{\ell}(j) \longrightarrow H^{\bullet+2}(S,\mathbb{Q}_{\ell}(j+1))
\end{equation*}
l'opération ``multiplication par u''.

}{\begin{equation*}
u \in H^2(S,\mathbb{Z}_{\ell}(1))
\end{equation*}
la ``classe de cohomologie d'une section hyperplane '' (de sorte que $u$ s'obtient comme l'image réciproque par le plongement $S \hookrightarrow \mathbb{P}$ de la première classe de Chern de $\underline{O}_{\mathbb{P}}$ (1), classe qui se trouve dans $H^2(\mathbb{P}\ \ , \mathbb{Z}_{\ell}(1))$), et désignons par

\begin{equation*}
L_S : H^{\bullet}(S,\mathbb{Q}_{\ell}(j)) \longrightarrow H^{\bullet+2}(S,\mathbb{Q}_{\ell}(j+1))
\end{equation*}
l'opération ``multiplication par u''.

}{F7II-XVIII-041}5.2.2. \underline{On dit  que} $S$ (\underline{muni de} $u$) \underline{vérifie} (LV), (le ``théorème de Lefschetz vache'') \underline{si pour tout} (\underline{ou un})\underline{entier} $j$ \underline{et tout entier} $i > 0$, \underline{les applications}

\begin{equation*}
L_S^i : H^{n-i}(S,\mathbb{Q}_{\ell}(j)) \longrightarrow H^{n+i}(S,\mathbb{Q}_{\ell}(i+j))
\end{equation*}
\underline{sont des isomorphismes}.


\SGApage{II}{282}

5.2.2.1.  On conjecture que cette condition est vérifiée si $S$ est lisse. C'est vrai si $k$ est de caractéristique nulle [6], d'où on conclut, grâce au ``théorème de spécialisation'', SGA 4 XVI 2.2 qu'il en est de même si $S$, avec son plongement projectif, se remonte en caractéristique nulle. Donc on trouve que c'est vérifié pour $S$ une intersection complète lisse. On sait également que c'est vrai pour $S$ lisse de dimension $\leq 2$ [2], et pour $S$ une variété abélienne (soit ``directement'' (LIEBERMAN [4] et [2]), soit par le théorème de MUMFORD [5] qui dit qu'une variété abélienne polarisée se remonte, avec sa polarisation, en caractéristique nulle).

Rappelons que pour une variété $S$ projective, on définit pour $i \leq n$ la ``\underline{partie primitive}'' de la cohomologie :

\SGAvar{\begin{equation}\tag{5.2.2.2}
\mathrm{Prim}^i(S,\mathbb{Q}_\ell(j)) = \mathrm{Ker}\ (L^{n-i+1} \ :\ H^i(S,\mathbb{Q}_\ell(j)) \longrightarrow H^{2n-i+2}(S,\mathbb{Q}_\ell(j+2n-2i+2)).
\end{equation}
\underline{Si} $S$ vérifie (LV), on a pour $i \leq n$ la décomposition :

}{\begin{equation}\tag{5.2.2.2}
\mathrm{Prim}^i(S,\mathbb{Q}_\ell(j)) = \mathrm{Ker}\ (L^{n-i+1} \ :\ H^i(S,\mathbb{Q}_\ell(j)) \longrightarrow H^{2n-i+2}(S,\mathbb{Q}_\ell(j+n-i+1))).
\end{equation}
\underline{Si} $S$ vérifie (LV), on a pour $i \leq n$ la décomposition :

}{F7II-XVIII-042}\begin{equation}\tag{5.2.2.3}
H^i(S,\mathbb{Q}_\ell(j)) \ \xleftarrow[\sim]{\ 1+L\ }\ \mathrm{Prim}^i(S,\mathbb{Q}_\ell(j)) \oplus H^{i-2}(S,\mathbb{Q}_\ell(j-1)) \qquad .
\end{equation}
\underline{Proposition} 5.2.3.  \underline{Soit} $X$ \underline{projectif et lisse,} \underline{vérifiant} (LV). \underline{Désignons} \underline{par} $\tau : Y \hookrightarrow X$ \underline{l'inclusion d'une section hyperplane lisse.} \underline{Alors}

\SGAvar{\begin{equation}\tag{5.2.3.1}
\mathrm{Prim}^n(X,\mathbb{Q}_\ell(j)) = \mathrm{Ker}(\tau^* \ :\ H^n(X,\mathbb{Q}_\ell(j)) \longrightarrow H^n(Y,\mathbb{Q}_\ell(j))) \qquad .
\end{equation}
\underline{Démonstration.}  On a, grâce à la formule de projection pour $m : Y \longrightarrow X$, le diagramme commutatif 


}{\begin{equation}\tag{5.2.3.1}
\mathrm{Prim}^n(X,\mathbb{Q}_\ell(j)) = \mathrm{Ker}(\tau^* \ :\ H^n(X,\mathbb{Q}_\ell(j)) \longrightarrow H^n(Y,\mathbb{Q}_\ell(j))) \qquad .
\end{equation}
\underline{Démonstration.}  On a, grâce à la formule de projection pour $\tau : Y \longrightarrow X$, le diagramme commutatif 


}{F7II-XVIII-043}\SGApage{II}{283}

\SGAvar{\begin{equation}\tag{5.2.3.2}
\begin{tikzcd}[column sep=huge, row sep=huge]
H^n(X,\mathbb{Q}_\ell(j)) \arrow[rr, "L"] \arrow[dr, "\tau^*"'] \& {} \& H^{n+2}(X,\mathbb{Q}_\ell(j+1)) \arrow[dl, "\tau_*"] \\
{} \& H^n(Y,\mathbb{Q}_\ell(j)) \& {} \& {.}
\end{tikzcd}
\end{equation}













}{\begin{equation}\tag{5.2.3.2}
\begin{tikzcd}[column sep=huge, row sep=huge]
H^n(X,\mathbb{Q}_\ell(j)) \arrow[rr, "L"] \arrow[dr, "\tau^*"'] \& {} \& H^{n+2}(X,\mathbb{Q}_\ell(j+1)) \\
{} \& H^n(Y,\mathbb{Q}_\ell(j)) \arrow[ur, "\tau_*"'] \& {} \& {.}
\end{tikzcd}
\end{equation}













}{F7II-XVIII-044}Par dualité et Lefschetz faible,

\SGAvar{\begin{equation*}
\tau_* \ :\ H^n(Y,\mathbb{Q}_\ell(j)) \longrightarrow H^{n+2}(X,\mathbb{Q}_\ell(j+1))
\end{equation*}
est la transposée d'une isomorphisme, donc un isomorphisme, d'où la conclusion.

}{\begin{equation*}
\tau_* \ :\ H^n(Y,\mathbb{Q}_\ell(j)) \longrightarrow H^{n+2}(X,\mathbb{Q}_\ell(j+1))
\end{equation*}
est la transposée d'un isomorphisme, donc un isomorphisme, d'où la conclusion.

}{F7II-XVIII-045}\underline{Lemme} 5.2.4.  \underline{Soient} $X$ \underline{projectif et lisse,} \underline{et} $\tau : Y \hookrightarrow X$ \underline{l'inclusion d'une section hyperplane lisse.} \underline{Alors} $X$ \underline{vérifie} (LV) \underline{si et seulement si les deux conditions suivantes sont vérifiées.}

5.2.4.1.  $Y$ \underline{vérifie} (LV) .

5.2.4.2.  \underline{La restriction du cup-produit} $H^{n-1}(Y,Q_\ell)$ \underline{au sous-espace} $\mathrm{Im}\ H^{n-1}(X,Q_\ell)$ \underline{est une forme non-dégénérée.}

\underline{Démonstration.}  On a

\begin{equation}\tag{5.2.4.3}
\left\{
\begin{aligned}
L_X &\ =\ \tau_*\tau^* \\[4ex]
L_Y &\ =\ \tau^*\tau_* \qquad \text{(par 4.2.11),}
\end{aligned}
\right. 
\end{equation}

\SGApage{II}{284}

 de sorte que

\begin{equation}\tag{5.2.4.4}
L_X^{i+1} \;=\; \tau_* L_Y^i \, \tau^* \qquad .
\end{equation}
Dans le diagramme commutatif

\begin{equation}\tag{5.2.4.5}
\begin{tikzcd}[column sep=huge, row sep=huge]
H^{n-i-1}(X,\mathbb{Q}_\ell(j)) \arrow[r, "L_X^{i+1}"] \arrow[d, "\tau^*"] \& H^{n+i+1}(X,\mathbb{Q}_\ell(j+i+1)) \\
H^{n-1-i}(Y,\mathbb{Q}_\ell(j)) \arrow[r, "L_Y^i"] \& H^{n-1+i}(Y,\mathbb{Q}_\ell(j+i)) \arrow[u, "\tau_*"']
\end{tikzcd}
\end{equation}












on sait que $\tau^*$ est un isomorphisme pour $i \geq 1$ par ``Lefschetz faible'', et que $\tau_*$ l'est par dualité. Donc $X$ vérifie (LV) si et seulement si $Y$ vérifie (LV), et la flèche

\begin{equation}\tag{5.2.4.6}
L_X : H^{n-1}(X,\mathbb{Q}_\ell(j)) \longrightarrow H^{n+1}(X,\mathbb{Q}_\ell(j+1))
\end{equation}
est un isomorphisme. Pour achever la démonstration, il suffit de voir l'équivalence de 5.2.4.2 et 5.2.4.6. Pour ceci, soit $a$, $b \in H^{n-1}(X,\mathbb{Q}_\ell(j))$. Par la dualité de Poincaré sur $X$, (5.2.4.6) équivaut à :

(5.2.4.7) La forme bilinéaire

\SGAvar{\begin{equation*}
H^{n-1}(X,\mathbb{Q}_\ell(j))\times H^{n-1}(X,\mathbb{Q}_\ell(j)) \longrightarrow H^{2n}(X,\mathbb{Q}_\ell(2j+2)
\end{equation*}
\begin{equation*}
(a,b) \longmapsto a.L_X(b)
\end{equation*}
est non-dégénérée.


}{\begin{equation}\tag{5.2.4.7}
H^{n-1}(X,\mathbb{Q}_\ell(j))\times H^{n-1}(X,\mathbb{Q}_\ell(j)) \longrightarrow H^{2n}(X,\mathbb{Q}_\ell(2j+1))
\end{equation}
\begin{equation*}
(a,b) \longmapsto a.L_X(b)
\end{equation*}
est non-dégénérée.


}{F7II-XVIII-046}\SGApage{II}{285}

\SGAvar{Mais par (5.2.5) et la formule de projection, on a

}{Mais par (5.2.4.3) et la formule de projection, on a

}{F7II-XVIII-047}\begin{equation*}
a.L_X(b) = a.\tau_*(\tau^*(b)) = \tau_*(\tau^*(a).\tau^*(b)),
\end{equation*}
d'où l'équivalence de (5.2.4.7) et (5.2.4.2).

\SGAvar{\underline{Corollaire} 5.2.5. \underline{Soient} $X$ \underline{projectif et lisse}, \underline{de dimension} $n$, \underline{et} $\tau : Y \hookrightarrow X$ \underline{l'inclusion d'une section hyperplan lisse}. \underline{Supposons que} $X$ \underline{satisfait à la condition} (LV). \underline{Alors pour} $i \geq n$, \underline{la flèche}

}{\underline{Corollaire} 5.2.5. \underline{Soient} $X$ \underline{projectif et lisse}, \underline{de dimension} $n$, \underline{et} $\tau : Y \hookrightarrow X$ \underline{l'inclusion d'une section hyperplane lisse}. \underline{Supposons que} $X$ \underline{satisfait à la condition} (LV). \underline{Alors pour} $i \geq n$, \underline{la flèche}

}{F7II-XVIII-048}\begin{equation*}
\tau^* : H^i(X,\mathbb{Q}_\ell(j)) \longrightarrow H^i(Y,\mathbb{Q}_\ell(j))
\end{equation*}
\underline{est surjective.}

\underline{Démonstration}. Par (5.2.4.3), on a, pour $i \geq 0$

\begin{equation*}
L_Y^{i+1} \;=\; \tau^* \, L_X^i \, \tau_* \quad ,
\end{equation*}
d'où le diagramme commutatif

\SGAvar{\begin{center}
\begin{tikzcd}[column sep=huge, row sep=huge]
H^{n-i}(X,\mathbb{Q}_\ell(j+1)) \arrow[r, "L_X^i"] \& H^{n-i}(X,\mathbb{Q}_\ell(j+i+1)) \arrow[d, "\tau^*"] \\
H^{n-i-2}(Y,\mathbb{Q}_\ell(j)) \arrow[r, "L_Y^{i+1}"{yshift=7pt}, "\sim"] \arrow[u, "\tau_*"'] \& H^{n+i}(Y,\mathbb{Q}_\ell(j+i+1))
\end{tikzcd}
\end{center}












}{\begin{center}
\begin{tikzcd}[column sep=huge, row sep=huge]
H^{n-i}(X,\mathbb{Q}_\ell(j+1)) \arrow[r, "L_X^i"] \& H^{n+i}(X,\mathbb{Q}_\ell(j+i+1)) \arrow[d, "\tau^*"] \\
H^{n-i-2}(Y,\mathbb{Q}_\ell(j)) \arrow[r, "L_Y^{i+1}"{yshift=7pt}, "\sim"] \arrow[u, "\tau_*"'] \& H^{n+i}(Y,\mathbb{Q}_\ell(j+i+1))
\end{tikzcd}
\end{center}












}{F7II-XVIII-049}ce qui donne le résultat voulu, cqfd.


\SGApage{II}{286}

\underline{Lemme} 5.2.6. \underline{Soit} $f : F \longrightarrow B$ \underline{un morphisme propre et lisse, avec} $B$ \underline{connexe. Désignons par} $\underline{O}_F(1)$ \underline{un module relativement ample sur} $F$, \underline{et munissons chaque fibre géométrique} $F_{\bar{b}}$ \underline{du module ample induit par} $\underline{O}_F(1)$. \underline{Alors, si une des fibres géométriques satisfait la condition} (LV), \underline{il en est de même pour toutes les autres.}

\SGAvar{\underline{Démonstration}. Notons par $\xi_F \in H^2(F,\mathbb{Z}_\ell(1))$ la première classe de Chern de $\underline{O}_F(1)$. Cette classe, et ses puissances $\xi_F^i \in H^{ii}(F,\mathbb{Z}_\ell(i))$ induisent des endomorphismes par cup-produit

\begin{equation}\tag{5.2.6.1}
\xi_F^i \; : \; R^{n-i}_{f_*}\mathbb{Q}_\ell(j) \longrightarrow R^{n+i}_{f_*}\mathbb{Q}_\ell(j+i)
\end{equation}
où $n$ désigne la dimension relative des fibres. Par le théorème de changement de base pour un morphisme propre [SGA 4 XII 5.2], toutes les fibres géométriques satisfont à la condition (LV) si et seulement si toutes les flèches (5.2.6.1) sont des isomorphismes. Or, par le théorème de spécialisation [SGA 4 XVI 2.2], les faisceaux $R^i f_*\mathbb{Q}_\ell(j)$ sont constants tordus, de sorte que les flèches (5.2.6.1) sont des isomorphismes si et seulement ils le sont au-dessus d'un point géométrique donné, cqfd.

\underline{Corollaire} 5.2.7. \underline{Soit} $X$ \underline{propective et lisse. Si une section hyperplane lisse vérifie} (LV), \underline{il en est de même pour} toute \underline{section hyperplane lisse.}

\underline{Démonstration}. En effet, la famille universelle des sections hyperplanes lisses de $X$ vérifie les hypothèses de 5.2.5.


}{\underline{Démonstration}. Notons par $\xi_F \in H^2(F,\mathbb{Z}_\ell(1))$ la première classe de Chern de $\underline{O}_F(1)$. Cette classe, et ses puissances $\xi_F^i \in H^{2i}(F,\mathbb{Z}_\ell(i))$ induisent des endomorphismes par cup-produit

\begin{equation}\tag{5.2.6.1}
\xi_F^i \; : \; R^{n-i}_{f_*}\mathbb{Q}_\ell(j) \longrightarrow R^{n+i}_{f_*}\mathbb{Q}_\ell(j+i)
\end{equation}
où $n$ désigne la dimension relative des fibres. Par le théorème de changement de base pour un morphisme propre [SGA 4 XII 5.2], toutes les fibres géométriques satisfont à la condition (LV) si et seulement si toutes les flèches (5.2.6.1) sont des isomorphismes. Or, par le théorème de spécialisation [SGA 4 XVI 2.2], les faisceaux $R^i f_*\mathbb{Q}_\ell(j)$ sont constants tordus, de sorte que les flèches (5.2.6.1) sont des isomorphismes si et seulement s'ils le sont au-dessus d'un point géométrique donné, cqfd.

\underline{Corollaire} 5.2.7. \underline{Soit} $X$ \underline{projective et lisse. Si une section hyperplane lisse vérifie} (LV), \underline{il en est de même pour} toute \underline{section hyperplane lisse.}

\underline{Démonstration}. En effet, la famille universelle des sections hyperplanes lisses de $X$ vérifie les hypothèses de 5.2.6.


}{F7II-XVIII-050}\SGApage{II}{287}

5.3. \underline{La condition} (A).

5.3.1. Soit

\begin{equation*}
\rho : \widetilde{X} \longrightarrow \mathbb{P}^1
\end{equation*}
la ``fibration'' associée à un pinceau dont l'axe coupe transversalement $X$, et dont la fibre générique est supposée lisse.

Désignons par $U$ un ouvert non-vide de $\mathbb{P}^1$ au-dessus duquel $\rho$ soit lisse, et par

\begin{equation}\tag{5.3.1.1}
\nu : U \hookrightarrow \mathbb{P}^1
\end{equation}
l'inclusion. Compte tenu de ce que les faisceaux $R^i_{\rho*}(\mathbb{Q}_\ell(j))$ sont \underline{constants tordus au-dessus de} $U$, on voit que les faisceaux

\begin{equation*}
\nu_*\nu^* \; R^i_{\rho*}(\mathbb{Q}_\ell(j))
\end{equation*}
sont essentiellement \underline{indépendants} du choix de $U$, et que les homomorphismes canoniques ``d'adjonction''

\begin{equation}\tag{5.3.2}
R^i_{\rho*}\mathbb{Q}_\ell(j) \longrightarrow \nu_*\nu^* \; R^i_{\rho*}(\mathbb{Q}_\ell(j))
\end{equation}
sont aussi indépendants du choix de $U$ .

5.3.3. Remarquons que pour un pinceau \underline{quelconque} de sections hyperplanes, le théorème de Lefschetz faible implique que pour $i \leq n-2$, l'immersion $\widetilde{X} \overset{(f\times\rho)}{\hookrightarrow} X\times\mathbb{P}^1$ donne lieu à un isomorphisme de faisceaux sur $\mathbb{P}^1$

\begin{equation*}
R^i_{\rho*}\mathbb{Q}_\ell(j) \; \xleftarrow{\ \sim\ } \; H^i(X,\mathbb{Q}_\ell(j))_{\mathbb{P}^1} \;\; , \;\; i \leq n-2
\end{equation*}

\SGApage{II}{288}

 i.e. $R^i_{\rho *}\mathbb{Q}_\ell(j)$ est un faisceau constant pour $i \leq n-2$ et, a fortiori, les flèches (5.3.2) sont des isomorphismes pour $i \leq n-2$ .

5.3.4. On sait, par la théorie locale des cycles évanescents et notamment la \underline{formule de Picard-Lefschetz} (XV 3.4 ), que, \underline{pour un pinceau de Lefschetz, les homomorphismes} 5.3.2 \underline{sont des isomorphismes pour} $i \neq n$ , \underline{et sont surjectifs pour} $i = n$.

5.3.5. On dit que $\rho$ , ou le pinceau dont il provient, \underline{vérifie la condition} (A), si \underline{tous} les homomorphismes (5.3.2) sont des isomorphismes (ou, ce qui revient au même, s'ils le sont pour tout $i$, et $j = 0$) .

On verra plus bas (6.3, 6.4) que cette condition (A) a une nette tendance à être vérifiée.

\underline{Lemme} 5.4. \underline{Si une des sections hyperplanes lisses de} $X$ (\underline{donc toute, par} 5.2.7) \underline{vérifie} (LV) (\underline{p.ex.} \underline{si} $X$ \underline{vérifie} (LV), cf. (5.2.4)), \underline{et si} $\rho$ \underline{vérifie l'hypothèse} (A), \underline{alors toute fibre} (\underline{y compris toute fibre singulière}) \underline{vérifie} (LV), \underline{et les faisceaux} $R^i\rho_*(Q_\ell(j))$ \underline{sont constants pour} $i \neq n-1$ .

\underline{Démonstration.} Désignons par

\SGAvar{\begin{equation*}
L : R^i\rho_*(\mathbb{Q}_\ell(j)) \longrightarrow R^{i+2}\rho_*(\mathbb{Q}_\ell(j+1))
\end{equation*}
le morphisme multiplication par $f^*c_1(\mathcal{O}_X(1) \in H^2(\widetilde{X},\mathbb{Z}_\ell(1))$; fibre par fibre, c'est l'opération envisagée dans (5.2.1). Par hypothèse, on a
\SGApage{II}{289}
un isomorphisme au-dessus de $U$

}{\begin{equation*}
L : R^i\rho_*(\mathbb{Q}_\ell(j)) \longrightarrow R^{i+2}\rho_*(\mathbb{Q}_\ell(j+1))
\end{equation*}
le morphisme multiplication par $f^*c_1(\mathcal{O}_X(1)) \in H^2(\widetilde{X},\mathbb{Z}_\ell(1))$; fibre par fibre, c'est l'opération envisagée dans (5.2.1). Par hypothèse, on a
\SGApage{II}{289}
un isomorphisme au-dessus de $U$

}{F7II-XVIII-051}\begin{equation}\tag{5.4.1}
\nu^*(L^i) : \nu^*R^{n-1-i}\rho_*(\mathbb{Q}_\ell(j)) \overset{\sim}{\longrightarrow} \nu^*R^{n-1+i}\rho_*(\mathbb{Q}_\ell(j+i))
\end{equation}
pour $i > 0$, tout $j$ . Appliquant le foncteur $\nu_*$ , on obtient des \underline{isomorphismes}

\begin{equation}\tag{5.4.2}
(\nu_*\nu^*) L^i : \nu_*\nu^* R^{n-1-i}\rho_*(\mathbb{Q}_\ell(j)) \overset{\sim}{\longrightarrow} \nu_*\nu^*R^{n-1+i}\rho_*(\mathbb{Q}_\ell(j+i)),
\end{equation}
qui, compte tenu de (A) s'identifient aux morphismes

\SGAvar{\begin{equation}\tag{5.4.3}
L^i : R^{n-1-i}\rho^*\mathbb{Q}_\ell(j) \overset{\sim}{\longrightarrow} R^{n-1+i}\rho_*\mathbb{Q}_\ell(j+i) .
\end{equation}
On conclut par 5.3.3.

}{\begin{equation}\tag{5.4.3}
L^i : R^{n-1-i}\rho_*\mathbb{Q}_\ell(j) \overset{\sim}{\longrightarrow} R^{n-1+i}\rho_*\mathbb{Q}_\ell(j+i) .
\end{equation}
On conclut par 5.3.3.

}{F7II-XVIII-052}5.4.4. Désignons par $Y$ une fibre lisse de $\rho : \widetilde{X} \longrightarrow \mathbb{P}^1$ . Par ``Lefschetz faible'', la flèche de restriction

\SGAvar{\begin{equation*}
H^{n-1}(X,\mathbb{Q}_\ell(j)) \overset{m^*}{\hookrightarrow} H^{n-1}(Y,\mathbb{Q}_\ell(j))
\end{equation*}
est injectif, et on a vu (5.2.4) que, si $X$ vérifie (LV), le cup-produit sur $H^{n-1}(Y,\mathbb{Q}_\ell(j))$, restreint à $m^* H^{n-1}(X,\mathbb{Q}_\ell(j))$, reste non dégénéré. L'orthogonal de $H^{n-1}(X,\mathbb{Q}_\ell(j))$ dans $H^{n-1}(Y,\mathbb{Q}_\ell(j))$ est donc un sous-espace complémentaire de $H^{n-1}(X,\mathbb{Q}_\ell(j))$ dans $H^{n-1}(Y,\mathbb{Q}_\ell(j))$, qu'on appelle, suivant la terminologie classique, l'\underline{espace de cohomologie évanescente de} $Y$ . On définit un sous-faisceau correspondant de $\nu_*(R^{n-1}\rho_*\mathbb{Q}_\ell(j))$, où $\nu : U \longrightarrow \mathbb{P}^1$ est un ouvert de lissité $\neq \emptyset$ de $\rho$ , le ``\underline{faisceau de cohomologie évanescente}'' $E^{n-1}\rho_*\mathbb{Q}_\ell(j)$ sur $\mathbb{P}^1$, par 


}{\begin{equation*}
H^{n-1}(X,\mathbb{Q}_\ell(j)) \overset{m^*}{\hookrightarrow} H^{n-1}(Y,\mathbb{Q}_\ell(j))
\end{equation*}
est injective, et on a vu (5.2.4) que, si $X$ vérifie (LV), le cup-produit sur $H^{n-1}(Y,\mathbb{Q}_\ell(j))$, restreint à $m^* H^{n-1}(X,\mathbb{Q}_\ell(j))$, reste non dégénéré. L'orthogonal de $H^{n-1}(X,\mathbb{Q}_\ell(j))$ dans $H^{n-1}(Y,\mathbb{Q}_\ell(j))$ est donc un sous-espace complémentaire de $H^{n-1}(X,\mathbb{Q}_\ell(j))$ dans $H^{n-1}(Y,\mathbb{Q}_\ell(j))$, qu'on appelle, suivant la terminologie classique, l'\underline{espace de cohomologie évanescente de} $Y$ . On définit un sous-faisceau correspondant de $\nu_*(R^{n-1}\rho_*\mathbb{Q}_\ell(j))$, où $\nu : U \longrightarrow \mathbb{P}^1$ est un ouvert de lissité $\neq \emptyset$ de $\rho$ , le ``\underline{faisceau de cohomologie évanescente}'' $E^{n-1}\rho_*\mathbb{Q}_\ell(j)$ sur $\mathbb{P}^1$, par 


}{F7II-XVIII-053}\SGApage{II}{290}

\begin{equation}\tag{5.4.4.1}
 \;\; E^{n-1}\rho_*\mathbb{Q}_\ell(j) \;=\; \nu_*\left\{\begin{array}{l}\text{l'orthogonal dans } \nu^*R^{n-1}\rho_*\mathbb{Q}_\ell(j) \text{ du}\\[2ex]\text{sous-faisceau constant } H^{n-1}(X,\mathbb{Q}_\ell(j))_U\end{array}\right\} \;\;.
\end{equation}
Il est clair que ce faisceau est indépendant du choix de $U$ .

\underline{Lemme} 5.5. \underline{Si} $X$ \underline{vérifie} (LV) \underline{et si les flèches} (5.3.2) \underline{sont des isomorphismes pour} $i = n - 1$, \underline{par exemple si} $\rho$ \underline{provient d'un pinceau de Lefschetz} (cf. (5.3.4)), \underline{on a une décomposition en somme directe}

\begin{equation}\tag{5.5.1}
R^{n-1}\rho_*\mathbb{Q}_\ell(j) \quad \simeq \quad H^{n-1}(X,\mathbb{Q}_\ell(j))_{\mathbb{P}^1} \oplus E^{n-1}\rho_*\mathbb{Q}_\ell(j) \qquad .
\end{equation}
\underline{Démonstration.} Par (LV), on a une telle décomposition au-dessus de $U$ :

\begin{equation}\tag{5.5.2}
\nu^*R^{n-1}\rho_*\mathbb{Q}_\ell(j) \quad \simeq \quad H^{n-1}(X,\mathbb{Q}_\ell(j))_U \oplus \nu^*\, E^{n-1}\rho_*\mathbb{Q}_\ell(j),
\end{equation}
à laquelle on applique $\nu_*$ .

\SGAvar{\underline{Théorème} 5.6. \underline{Soit} $X$ \underline{une variété projective, lisse et irréductible qui vérifie} (LV) (5.2.2). \underline{Soit} $\{\, H_\tau \,\}$ \underline{un pinceau d'hyperplans dont l'élément générique et l'axe coupent transversalement} $X$, \underline{en} $X_\eta$ \underline{et en} $\Delta$ \underline{respectivement. Supposons que la projection} $\rho : \widetilde{X} \longrightarrow \mathbb{P}^1$ \underline{vérifie la condition} (5.3.2). \underline{Alors} :

}{\underline{Théorème} 5.6. \underline{Soit} $X$ \underline{une variété projective, lisse et irréductible qui vérifie} (LV) (5.2.2). \underline{Soit} $\{\, H_\tau \,\}$ \underline{un pinceau d'hyperplans dont l'élément générique et l'axe coupent transversalement} $X$, \underline{en} $X_\eta$ \underline{et en} $\Delta$ \underline{respectivement. Supposons que la projection} $\rho : \widetilde{X} \longrightarrow \mathbb{P}^1$ \underline{vérifie la condition} (A) \underline{de} 5.3.5. \underline{Alors} :

}{F7II-XVIII-054}(5.6.1) \quad \underline{la suite spectrale de Leray}

\SGAvar{\begin{equation*}
E_2^{p,q} = H^p(\mathbb{P}^1,\; R^q\rho_*\mathbb{Q}_\ell) \;\Longrightarrow\; H^{m+q}(\widetilde{X},\; \mathbb{Q}_\ell)
\end{equation*}
\underline{dégénère.} (NB. \underline{Ceci reste valable si, au lieu de supposer que} $X$ \underline{satisfait la condition} (LV), \underline{on suppose seulement qu'une section hyperplane lisse} $Y$ \underline{y satisfait.})


}{\begin{equation*}
E_2^{p,q} = H^p(\mathbb{P}^1,\; R^q\rho_*\mathbb{Q}_\ell) \;\Longrightarrow\; H^{p+q}(\widetilde{X},\; \mathbb{Q}_\ell)
\end{equation*}
\underline{dégénère.} (NB. \underline{Ceci reste valable si, au lieu de supposer que} $X$ \underline{satisfait la condition} (LV), \underline{on suppose seulement qu'une section hyperplane lisse} $Y$ \underline{y satisfait.})


}{F7II-XVIII-055}\SGApage{II}{291}

5.6.2. \underline{Désignons par} $\eta$ \underline{le point générique de} $\mathbb{P}^1$, \underline{par} $\overline{\eta}$ \underline{le point générique géométrique,} \underline{par} $X_{\overline{\eta}}$ \underline{la fibre générique géométrique de} $\rho$, \underline{et par} $\pi$ \underline{le groupe de Galois de} $k(\overline{\eta})/k(\eta)$.

\underline{Pour} $q \neq n-1$, \underline{on a des isomorphismes canoniques} :

\SGAvar{\begin{align}
&E_2^{o,q} \; (=H^{o}(\mathbb{P}^1,R^q\rho_*\mathbb{Q}_\ell)) \quad \simeq \quad H^q(X_{\overline{\eta}},\mathbb{Q}_\ell) \tag{5.6.2.1}\\
\smash{\left\{\rule[-5ex]{0pt}{10ex}\right.}\;&E_2^{2,q}(=H^2(\mathbb{P}^1,R^q\rho_*\mathbb{Q}_\ell)) \quad \simeq \quad H^q(X_{\overline{\eta}},\mathbb{Q}_\ell(-1)) \tag{5.6.2.2}\\
&E_2^{p,q}(=H^p(\mathbb{P}^1,R^q\rho_*\mathbb{Q}_\ell)) \;=\; 0,\; p \neq 0,2 \qquad . \tag{5.6.2.3}
\end{align}
\underline{Pour} $q = n-1$, \underline{on a}

\begin{equation}\tag{5.6.2.4}
\begin{array}{r@{\quad \simeq \quad}l} E_2^{0,n-1}(=H^{o}(\mathbb{P}^1,R^{n-1}\rho_*\mathbb{Q}_\ell)) & H^{n-1}(X_{\overline{\eta}},\mathbb{Q}_\ell)^{\pi}\\[4ex]  & H^{n-1}(X,\mathbb{Q}_\ell) \;\;, \end{array}
\end{equation}
\begin{equation}\tag{5.6.2.5}
\begin{array}{r@{\quad \simeq \quad}l} E_2^{2,n-1}(=H^2(\mathbb{P}^1,R^{n-1}\rho_*\mathbb{Q}_\ell) & H^{n-1}(X_{\overline{\eta}},\mathbb{Q}_\ell(-1))^{\pi}\\[4ex]  & H^{n-1}(X,\mathbb{Q}_\ell(-1)) \qquad . \end{array}
\end{equation}
(\underline{On trouvera la valeur manquante} $E_2^{1,n-1}$ \underline{dans} 5.7 \underline{ci-dessous}).

\underline{Démonstration.} Le théorème 1.1 s'applique, par (5.4.3), pour donner (5.6.1). Quant aux trois premières formules (5.6.2), elles résultent du fait que les $R^q\rho_*\mathbb{Q}_\ell$ pour $q \neq n-1$ sont constants (5.4) ; pour (5.6.2.3) il faut utiliser bien sûr que $\mathbb{P}^1$ est simplement connexe. La première
\SGApage{II}{292}
égalité de (5.6.2.4) résulte de ce que $R^{n-1}\rho_*\mathbb{Q}_{\ell} \xrightarrow{\ \sim\ } \nu_*\nu^*R^{n-1}\rho_*\mathbb{Q}_{\ell}$ , de sorte qu'on a $H^0(\mathbb{P}^1,R^{n-1}\rho_*\mathbb{Q}_{\ell}) \simeq H^0(U,\nu^*R^{n-1}f_*\mathbb{Q}_{\ell})$, et de ce que $\nu^*R^{n-1}f_*\mathbb{Q}_{\ell}$ est un faisceau ``constant tordu'' sur $U$, donc équivalent au $\pi$-module $H^{n-1}(X_{\overline{\eta}},\mathbb{Q}_{\ell})$. La deuxième égalité de (5.6.2.4) résulte de la dégénerescence, compte tenu de 5.1.6.

}{\begin{align}
&E_2^{0,q} \; (=H^{0}(\mathbb{P}^1,R^q\rho_*\mathbb{Q}_\ell)) \quad \simeq \quad H^q(X_{\overline{\eta}},\mathbb{Q}_\ell) \tag{5.6.2.1}\\
\smash{\left\{\rule[-5ex]{0pt}{10ex}\right.}\;&E_2^{2,q}(=H^2(\mathbb{P}^1,R^q\rho_*\mathbb{Q}_\ell)) \quad \simeq \quad H^q(X_{\overline{\eta}},\mathbb{Q}_\ell(-1)) \tag{5.6.2.2}\\
&E_2^{p,q}(=H^p(\mathbb{P}^1,R^q\rho_*\mathbb{Q}_\ell)) \;=\; 0,\; p \neq 0,2 \qquad . \tag{5.6.2.3}
\end{align}
\underline{Pour} $q = n-1$, \underline{on a}

\begin{equation}\tag{5.6.2.4}
\begin{array}{r@{\quad \simeq \quad}l} E_2^{0,n-1}(=H^{0}(\mathbb{P}^1,R^{n-1}\rho_*\mathbb{Q}_\ell)) & H^{n-1}(X_{\overline{\eta}},\mathbb{Q}_\ell)^{\pi}\\[4ex]  & H^{n-1}(X,\mathbb{Q}_\ell) \;\;, \end{array}
\end{equation}
\begin{equation}\tag{5.6.2.5}
\begin{array}{r@{\quad \simeq \quad}l} E_2^{2,n-1}(=H^2(\mathbb{P}^1,R^{n-1}\rho_*\mathbb{Q}_\ell)) & H^{n-1}(X_{\overline{\eta}},\mathbb{Q}_\ell(-1))^{\pi}\\[4ex]  & H^{n-1}(X,\mathbb{Q}_\ell(-1)) \qquad . \end{array}
\end{equation}
(\underline{On trouvera la valeur manquante} $E_2^{1,n-1}$ \underline{dans} 5.7 \underline{ci-dessous}).

\underline{Démonstration.} Le théorème 1.1 s'applique, par (5.4.3), pour donner (5.6.1). Quant aux trois premières formules (5.6.2), elles résultent du fait que les $R^q\rho_*\mathbb{Q}_\ell$ pour $q \neq n-1$ sont constants (5.4) ; pour (5.6.2.3) il faut utiliser bien sûr que $\mathbb{P}^1$ est simplement connexe. La première
\SGApage{II}{292}
égalité de (5.6.2.4) résulte de ce que $R^{n-1}\rho_*\mathbb{Q}_{\ell} \xrightarrow{\ \sim\ } \nu_*\nu^*R^{n-1}\rho_*\mathbb{Q}_{\ell}$ , de sorte qu'on a $H^0(\mathbb{P}^1,R^{n-1}\rho_*\mathbb{Q}_{\ell}) \simeq H^0(U,\nu^*R^{n-1}\rho_*\mathbb{Q}_{\ell})$, et de ce que $\nu^*R^{n-1}\rho_*\mathbb{Q}_{\ell}$ est un faisceau ``constant tordu'' sur $U$, donc équivalent au $\pi$-module $H^{n-1}(X_{\overline{\eta}},\mathbb{Q}_{\ell})$. La deuxième égalité de (5.6.2.4) résulte de la dégénérescence, compte tenu de 5.1.6.

}{F7II-XVIII-056}Passons à (5.6.2.5). Désignons par $F^{\bullet}$ la filtration définie sur $H^{\bullet}(\widetilde{X},\mathbb{Q}_{\ell})$ par la suite spectrale. Par dégénérescence, $F^2\,H^{n+1}(\widetilde{X},\mathbb{Q}_{\ell}) = E_2^{2,n-1}$, et $\mathrm{gr}^1H^{n+1}(\widetilde{X},\mathbb{Q}_{\ell}) = E_2^{1,n} = 0$ (5.6.2.3), de sorte que (cf. 5.1.0)

\SGAvar{\begin{equation}\tag{5.6.3}
E_2^{2,n-1} = F^1H^{n+1}(\widetilde{X},\mathbb{Q}_{\ell}) = \mathrm{Ker}(H^{n+1}(\widetilde{X},\mathbb{Q}_{\ell}) \longrightarrow H^0(\mathbb{P}^1,R^{n+1}\rho_*\mathbb{Q}_{\ell})
\end{equation}
\begin{equation*}
(\text{par }(A)) = \mathrm{Ker}(k^* : H^{n+1}(\widetilde{X},\mathbb{Q}_{\ell}) \longrightarrow H^{n+1}(X_{\overline{\eta}},\mathbb{Q}_{\ell})) \quad .
\end{equation*}
Par (5.1.2), $k^*$ se récrit

}{\begin{equation}\tag{5.6.3}
E_2^{2,n-1} = F^1H^{n+1}(\widetilde{X},\mathbb{Q}_{\ell}) = \mathrm{Ker}(H^{n+1}(\widetilde{X},\mathbb{Q}_{\ell}) \longrightarrow H^0(\mathbb{P}^1,R^{n+1}\rho_*\mathbb{Q}_{\ell}))
\end{equation}
\begin{equation*}
(\text{par }(A)) = \mathrm{Ker}(k^* : H^{n+1}(\widetilde{X},\mathbb{Q}_{\ell}) \longrightarrow H^{n+1}(X_{\overline{\eta}},\mathbb{Q}_{\ell})) \quad .
\end{equation*}
Par (5.1.2), $k^*$ se récrit

}{F7II-XVIII-057}\begin{equation}\tag{5.6.4}
H^{n+1}(X,\mathbb{Q}_{\ell}) \oplus H^{n-1}(\Delta,\mathbb{Q}_{\ell}(-1)) \xrightarrow{\ m^*+h_*\ } H^{n+1}(X_{\overline{\eta}},\mathbb{Q}_{\ell}) \quad ,
\end{equation}
et, par (5.2.2.3) se récrit encore

\SGAvar{\begin{equation}\tag{5.6.5}
\mathrm{Prim}^{n-1}(X,\mathbb{Q}_{\ell}(-1) \oplus H^{n-3}(X,\mathbb{Q}_{\ell}(-2)) \oplus H^{n-1}(\Delta,\mathbb{Q}_{\ell}(-1)) \xrightarrow{\ m^*L+m^*L^2h_*\ } H^{n+1}(X_{\overline{\eta}},\mathbb{Q}_{\ell}
\end{equation}
Or, la flèche

\begin{equation*}
H^{n-3}(X,\mathbb{Q}_{\ell}(-2) \xrightarrow{\ m^*L^2\ } H^{n+1}(X_{\overline{\eta}},\mathbb{Q}_{\ell})
\end{equation*}
est un isomorphisme, par (LV), car $m^*L_X^2 = L^2_{X_{\overline{\eta}}} \cdot m^*$ (cf.(5.2.4) et (5.2.5)). Donc le noyau de (5.6.8) s'identifie (par $\mathrm{pr}_1 \oplus \mathrm{pr}_3$) à
\SGApage{II}{293}
$\mathrm{Prim}^{n-1}(X,\mathbb{Q}_{\ell}(-1)) \oplus H^{n-1}(\Delta,\mathbb{Q}_{\ell}(-1))$. Par (LV), $H^{n-1}(\Delta,\mathbb{Q}_{\ell}(-1)) \xleftarrow[\sim]{L} H^{n-3}(\Delta,\mathbb{Q}_{\ell}(-2)) \xleftarrow[\sim]{\;} H^{n-3}(X,\mathbb{Q}_{\ell}(-2))$, la dernière par Lefschetz faible. Par (5.2.2.3), on a donc

}{\begin{equation}\tag{5.6.5}
\mathrm{Prim}^{n-1}(X,\mathbb{Q}_{\ell}(-1)) \oplus H^{n-3}(X,\mathbb{Q}_{\ell}(-2)) \oplus H^{n-1}(\Delta,\mathbb{Q}_{\ell}(-1)) \xrightarrow{\ m^*L+m^*L^2+h_*\ } H^{n+1}(X_{\overline{\eta}},\mathbb{Q}_{\ell})
\end{equation}
Or, la flèche

\begin{equation*}
H^{n-3}(X,\mathbb{Q}_{\ell}(-2)) \xrightarrow{\ m^*L^2\ } H^{n+1}(X_{\overline{\eta}},\mathbb{Q}_{\ell})
\end{equation*}
est un isomorphisme, par (LV), car $m^*L_X^2 = L^2_{X_{\overline{\eta}}} \cdot m^*$ (cf.(5.2.4) et (5.2.5)). Donc le noyau de (5.6.5) s'identifie (par $\mathrm{pr}_1 \oplus \mathrm{pr}_3$) à
\SGApage{II}{293}
$\mathrm{Prim}^{n-1}(X,\mathbb{Q}_{\ell}(-1)) \oplus H^{n-1}(\Delta,\mathbb{Q}_{\ell}(-1))$. Par (LV), $H^{n-1}(\Delta,\mathbb{Q}_{\ell}(-1)) \xleftarrow[\sim]{L} H^{n-3}(\Delta,\mathbb{Q}_{\ell}(-2)) \xleftarrow[\sim]{\;} H^{n-3}(X,\mathbb{Q}_{\ell}(-2))$, la dernière par Lefschetz faible. Par (5.2.2.3), on a donc

}{F7II-XVIII-058}\begin{equation}\tag{5.6.6}
E_2^{2,n-1} \simeq H^{n-1}(X,\mathbb{Q}_{\ell}(-1)) \quad ,
\end{equation}
et l'autre égalité de (5.6.2.5) se déduit en ``tordant'' (5.6.2.4).

\underline{Remarque} 5.6.7. On peut démontrer (5.6.1) sous des hypothèses nettement plus faibles, comme DELIGNE m'a expliqué;

\underline{Théorème} 5.6.8. \underline{Soient} $X$ \underline{lisse de dimension $n$ sur un corps} $k$ \underline{algébriquement clos, et}

\begin{equation*}
\rho : X \longrightarrow \mathbb{P}^1_k
\end{equation*}
\underline{un morphisme projectif, dont la fibre générique est lisse. Supposons que} :

(5.6.8.1) \underline{la fibre générique géométrique, munie de la restriction du module relativement ample} $\underline{O}_X(1)$, \underline{satisfait à la condition} (LV);

(5.6.8.2) \underline{les flèches} (5.3.2) \underline{sont des isomorphismes pour} $i \neq n$, \underline{et sont surjectives pour} $i = n$ (\underline{condition satisfaite en particulier si} $\rho$ \underline{provient d'un pinceau de Lefschetz} (cf. 5.3.4)). \underline{Alors, la suite spectrale de Leray}

\begin{equation*}
E_2^{p,q} : H^p(\ \mathbb{P}^1,R^q\rho_*\mathbb{Q}_{\ell}) \Longrightarrow H^{p+q}(X,\mathbb{Q}_{\ell})
\end{equation*}
\underline{dégénère.}

\underline{Remarque} : Le même théorème vaut, avec la même démonstration, pour $\mathbb{P}^1_k$ remplacé par une courbe projective lisse irréductible quelconque.


\SGApage{II}{294}

\underline{Démonstration}. Nous allons identifier ,pour la démonstration, le faisceau $\mathbb{Q}_\ell(1)$ à $\mathbb{Q}_\ell$ (possible car on a pris $k$ algébriquement clos). On peut évidemment supposer $X$ connexe, et $\rho$ \underline{surjectif}, donc plat, à fibres de dimension $n-1$. Désignons par

\SGAvar{\begin{equation*}
\xi_X \; \in \; H^2(X,\mathbb{Q}_\ell)
\end{equation*}
la première classe de Chern de $\underline{O}_{X(1)}$. Soit

\begin{equation*}
\nu \; : \; U \longrightarrow \mathbb{P}^1
\end{equation*}
l'inclusion d'un ouvert de lissité de $\rho$ . Par 5.6.8.1 et.5.2.6, les flèches de cup-produit

}{\begin{equation*}
\xi_X \; \in \; H^2(X,\mathbb{Q}_\ell)
\end{equation*}
la première classe de Chern de $\underline{O}_X(1)$. Soit

\begin{equation*}
\nu \; : \; U \longrightarrow \mathbb{P}^1
\end{equation*}
l'inclusion d'un ouvert de lissité de $\rho$ . Par 5.6.8.1 et 5.2.6, les flèches de cup-produit

}{F7II-XVIII-059}\begin{equation*}
\xi_X^i \; : \; \nu^* R^{n-1-i}\rho_* \mathbb{Q}_\ell \longrightarrow \nu^* R^{n-1+i}\rho_* \mathbb{Q}_\ell
\end{equation*}
sont des isomorphismes, pour $i \geq 1$, de sorte que, en appliquant le foncteur $\nu_*$ , on trouve des isomorphismes (compte tenu de 5.6.8.2) pour $i \geq 2$

\begin{equation}\tag{5.6.8.3}
\xi_X^i \; : \; R^{n-1-i}\rho_* \mathbb{Q}_\ell \; \xrightarrow{\;\sim\;} \; R^{n-1+i}\rho_* \mathbb{Q}_\ell \; ,
\end{equation}
et, pour $i = 1$, un isomorphisme

\begin{equation}\tag{5.6.8.4}
\xi_X \; : \; R^{n-2}\rho_* \mathbb{Q}_\ell \; \xrightarrow{\;\sim\;} \; \nu_* \nu^* R^n \rho_* \mathbb{Q}_\ell \quad .
\end{equation}
Cette dernière flèche se factorise en

\begin{equation}\tag{5.6.8.5}
R^{n-2}\rho_* \mathbb{Q}_\ell \; \xrightarrow{\;\xi_X\;} \; R^n \rho_* \mathbb{Q}_\ell \; \xrightarrow{\;\text{flèche d'adjonction}\;} \; \nu_* \nu^* R^n \rho_* \mathbb{Q}_\ell \quad , 
\end{equation}

\SGApage{II}{295}

 de sorte qu'on a une décomposition canonique

\begin{multline}\tag{5.6.8.6}
R^n \rho_* \mathbb{Q}_\ell \; \xleftarrow{\substack{\xi_X + \mathrm{ind} \\ \sim}} \; R^{n-2}\rho_* \mathbb{Q}_\ell \oplus \mathrm{Local}^n \mathbb{Q}_\ell
\end{multline}
où l'on a posé

\begin{equation}\tag{5.6.8.7}
\mathrm{Local}^n \mathbb{Q}_\ell \; = \; \mathrm{Ker}(R^n \rho_* \mathbb{Q}_\ell \longrightarrow \nu_* \nu^* R^n \rho_* \mathbb{Q}_\ell) \quad .
\end{equation}
On remarque que le faisceau $\mathrm{Local}^n \mathbb{Q}_\ell$ est concentré aux points de $\mathbb{P}^1$ où la fibre de $\rho$ n'est pas lisse, de sorte que, compte tenu encore de (5.6.8.3), on trouve

\SGAvar{\begin{multline}\tag{5.6.8.8}
H^j(\mathbb{P}^1, \; R^{n-2-i}\rho_* \mathbb{Q}_\ell) \; \xrightarrow{\substack{\xi_X \\ \sim}} \; H^j(\mathbb{P}^1,R^{n-1+i}\mathbb{Q}_\ell) \quad \text{si } i,j \geq 1 \; .
\end{multline}
Démontrons maintenant la dégénérescence. Compte tenu de ce que $E_2^{p,q} = 0$ sauf pour $p = 0$, $1,2$, il suffit de prouver que les opérations

\begin{equation*}
d_2 \; : \; H^0(\, \mathbb{P}^1,R^i \rho_* \mathbb{Q}_\ell) \longrightarrow H^2(\, \mathbb{P}^1,R^{i-1}, \rho_* \mathbb{Q}_\ell)
\end{equation*}
sont nulles. Pour $q \leq n - 1$, posons

\begin{equation}\tag{5.6.8.9}
\mathrm{Prim}(E_2^{p,q}) \; = \; \mathrm{Ker}(\xi_X^{n-q} \; : \; E_2^{p,q} \longrightarrow E_2^{p,2n-q}) \quad .
\end{equation}
Grâce à (5.6.8.3) et (5.6.8.6), impliquant l'\underline{injectivité} de $\xi_X \; : \; H^p(\, \mathbb{P}^1,R^{n-2}\rho_* \mathbb{Q}_\ell) \longrightarrow H^p(\, \mathbb{P}^1, \mathbb{R}^1 \rho_* \mathbb{Q}_\ell)$, pour $q \leq n - 1$, on a

\begin{multline}\tag{5.6.8.10}
E_2^{p,q} \; \xleftarrow{\substack{1+\xi_X+\xi_X^2+\dots \\ \sim}} \; \mathrm{Prim}\,(E_2^{p,q}) \oplus \mathrm{Prim}(E_2^{p,q-2}) \oplus \mathrm{Prim}(E_2^{p,q-4}) \; .
\end{multline}

}{\begin{multline}\tag{5.6.8.8}
H^j(\mathbb{P}^1, \; R^{n-2-i}\rho_* \mathbb{Q}_\ell) \; \xrightarrow{\substack{\xi_X^{i+1} \\ \sim}} \; H^j(\mathbb{P}^1,R^{n+i}\rho_*\mathbb{Q}_\ell) \quad \text{si } i \geq 0,\; j \geq 1 \; .
\end{multline}
Démontrons maintenant la dégénérescence. Compte tenu de ce que $E_2^{p,q} = 0$ sauf pour $p = 0$, $1,2$, il suffit de prouver que les opérations

\begin{equation*}
d_2 \; : \; H^0(\, \mathbb{P}^1,R^i \rho_* \mathbb{Q}_\ell) \longrightarrow H^2(\, \mathbb{P}^1,R^{i-1}\rho_* \mathbb{Q}_\ell)
\end{equation*}
sont nulles. Pour $q \leq n - 1$, posons

\begin{equation}\tag{5.6.8.9}
\mathrm{Prim}(E_2^{p,q}) \; = \; \mathrm{Ker}(\xi_X^{n-q} \; : \; E_2^{p,q} \longrightarrow E_2^{p,2n-q}) \quad .
\end{equation}
Grâce à (5.6.8.3) et (5.6.8.6), impliquant l'\underline{injectivité} de $\xi_X \; : \; H^p(\, \mathbb{P}^1,R^{n-2}\rho_* \mathbb{Q}_\ell) \longrightarrow H^p(\, \mathbb{P}^1, R^n \rho_* \mathbb{Q}_\ell)$, pour $q \leq n - 1$, on a

\begin{multline}\tag{5.6.8.10}
E_2^{p,q} \; \xleftarrow{\substack{1+\xi_X+\xi_X^2+\dots \\ \sim}} \; \mathrm{Prim}\,(E_2^{p,q}) \oplus \mathrm{Prim}(E_2^{p,q-2}) \oplus \mathrm{Prim}(E_2^{p,q-4}) \oplus \cdots \; .
\end{multline}

}{F7II-XVIII-060}\SGApage{II}{296}

Pour voir qu'on a

\begin{equation*}
d_2^{p,q} \;=\; 0 \quad \text{pour } q \leq n - 1 \quad ,
\end{equation*}
remarquons simplement que le diagramme commutatif

\begin{center}
\begin{tikzcd}[column sep=huge, row sep=huge]
\mathrm{Prim}(E_2^{p,q}) \arrow[r, "\xi_X^{n-q}=0"] \arrow[d, "{d_2^{p,q}}"] \& E_2^{p,2n-q} \arrow[d, "{d_2^{p,2n-q}}"] \\
E_2^{p+2,q-1} \arrow[r, "\xi_X^{n-1-(q-1)}", "\sim"'] \& E_2^{p+2,2n-q-1}
\end{tikzcd}
\end{center}












\SGAvar{(la flèche du bas étant un isomorphisme grâce à (5.6.8.8)), entraîne que $d_2^{p,q}$ tue $\mathrm{Prim}(E_2^{p,q})$ pour $q \leq n - 1$, et donc, par (5.6.8.10), que $d_2$ tue $E_2^{p,q}$ pour $q \leq n - 1$. Par (5.6.8.3), ceci implique aussi que $d_2$ tue $E_2^{p,q}$ pour $q \geq n + 1$. Donc ilreste à prouver que l'homomorphisme

}{(la flèche du bas étant un isomorphisme grâce à (5.6.8.8)), entraîne que $d_2^{p,q}$ tue $\mathrm{Prim}(E_2^{p,q})$ pour $q \leq n - 1$, et donc, par (5.6.8.10), que $d_2$ tue $E_2^{p,q}$ pour $q \leq n - 1$. Par (5.6.8.3), ceci implique aussi que $d_2$ tue $E_2^{p,q}$ pour $q \geq n + 1$. Donc il reste à prouver que l'homomorphisme

}{F7II-XVIII-061}\begin{equation*}
d_2^{0,n} \;:\; E_2^{0,n} \;\longrightarrow\; E_2^{2,n-1}
\end{equation*}
est nul, ou, ce qui revient au même, que

\begin{equation}\tag{5.6.8.12}
\dim E_3^{2,n-1} \;=\; \dim E_2^{2,n-1} \quad .
\end{equation}
Nous avons déjà démontré qu'on a

\begin{equation}\tag{5.6.8.13}
\dim E_2^{p,q} \;=\; \dim E_3^{p,q} \quad \text{pour } (p,q) \neq (0,n),\; (2,n-1)
\end{equation}
et, pour une raison de degrés,

\begin{equation}\tag{5.6.8.14}
E_3^{p,q} \;=\; E_\infty^{p,q} \quad \text{pour } \underline{\text{tout}} \text{ couple } (p,q) \quad .
\end{equation}

\SGApage{II}{297}

Par(5.6.8.14), on a

\begin{equation*}
b_i(X) \;=\; \dim H^i(X,\mathbb{Q}_\ell) \;=\; \sum_{p+q=i} \dim E_3^{p,q} \quad ,
\end{equation*}
et, compte tenu de (5.6.8.13), on a

\SGAvar{\begin{equation}\tag{5.6.8.15}
b_{n-1}(X) \;=\; \dim E_2^{c,n-1} \,+\, \dim E_2^{1,n-2} \,+\, \dim_2^{2,n-3}
\end{equation}
et

}{\begin{equation}\tag{5.6.8.15}
b_{n-1}(X) \;=\; \dim E_2^{0,n-1} \,+\, \dim E_2^{1,n-2} \,+\, \dim E_2^{2,n-3}
\end{equation}
et

}{F7II-XVIII-062}\begin{equation}\tag{5.6.8.16}
b_{n+1}(X) \;=\; \dim E_2^{0,n+1} \,+\, \dim E_2^{1,n} \,+\, \dim E_3^{2,n-1} \quad .
\end{equation}
Par la dualité de Poincaré, on a

\begin{equation}\tag{5.6.8.17}
b_{n-1}(X) \;=\; b_{n+1}(X) \quad .
\end{equation}
Remarquons d'abord que, par (5.6.8.8), on a

\begin{equation*}
\dim E_2^{1,n} \;=\; \dim E_2^{1,n-2} \quad ,
\end{equation*}
de sorte que (par 5.6.8.17) la démonstration de (5.6.8.12) est terminée une fois connue

\begin{equation}\tag{5.6.8.18}
\dim E_2^{2,n-3} \;=\; \dim E_2^{0,n+1} \quad ,
\end{equation}
\begin{equation}\tag{5.6.8.19}
\dim E_2^{0,n-1} \;=\; \dim E_2^{2,n-1} \quad .
\end{equation}
Or, ces deux formules résultent de (5.6.8.2) et la dualité de Poincaré pour les fibres lisses de $\rho$, compte tenu du lemme suivant, appliqué à $\mathbb{P}^1$ et aux faisceaux $R^{n-3}\rho_*\mathbb{Q}_\ell$ et $R^{n-1}\rho_*\mathbb{Q}_\ell$ .


\SGApage{II}{298}

\underline{Lemme} 5.6.9. (cf. SGA 5 I \S 5 et SGA 4 XVIII\quad ). \underline{Soit} $D$ \underline{une courbe propre,} \underline{lisse et connexe sur un corps} $k$ \underline{algébriquement clos.} \underline{Désignons par} $i : \eta \longrightarrow D$ \underline{l'inclusion du point générique,} \underline{et par} $F$ \underline{un} $\mathbb{Q}_\ell$\underline{-faisceau constructible sur} $D$ \underline{tel que} $F \overset{\sim}{\longrightarrow} i_*i^*F$ . \underline{Alors}

\begin{equation*}
H^2(D,F) \quad \simeq \quad H^0(D,\check{F}(-1))^{\vee} \quad ,
\end{equation*}
\underline{où} $\check{F} \simeq i_*((i^*F)^{\vee})$.

\underline{Corollaire} 5.6.10. (Utilise Picard-Lefschetz). \underline{Soit} $\rho : \tilde{X} \longrightarrow \mathbb{P}^1$ \underline{la} ``\underline{fibration}'' \underline{provenant d'un pinceau de Lefschetz,} \underline{et supposons que les} \underline{fibres lisses satisfont à la condition} (LV). \underline{Alors}

\begin{equation}\tag{5.6.10.1}
H^i(X_{\bar{\eta}},\mathbb{Q}_\ell)^{\pi} \;=\; \mathrm{Im}\,(H^i(X,\mathbb{Q}_\ell) \longrightarrow H^i(X_{\bar{\eta}},\mathbb{Q}_\ell)) \quad .
\end{equation}
\underline{Démonstration.} Par 5.3.4. (qui utilise la formule de Picard-Lefschetz), 5.6.8 s'applique, et donne

\begin{equation*}
H^i(X_{\bar{\eta}},\mathbb{Q}_\ell)^{\pi} \;=\; \mathrm{Image}\; H^i(\tilde{X},\mathbb{Q}_\ell) \quad .
\end{equation*}
Il reste à voir que $H^i(X,\mathbb{Q}_\ell)$ et $H^i(\tilde{X},\mathbb{Q}_\ell)$ ont la même image dans $H^i(X_{\bar{\eta}},\mathbb{Q}_\ell)$. Or, ceci est vrai pour $i \leq n-1$ par 5.1.6, et pour $i \geq n$ par 5.2.5.

Passons maintenant au calcul le plus important :

\underline{Théorème} 5.7. \underline{Sous les hypothèses de} 5.6, \underline{posant}

\begin{equation*}
E^{n-2}(\Delta,\mathbb{Q}_\ell(j)) = \underline{\mbox{orthogonal dans}}\; H^{n-2}(\Delta,\mathbb{Q}_\ell(j)) \; \underline{\mbox{de l'image de}} \; H^{n-2}(X,\mathbb{Q}_\ell(j))
\end{equation*}
(\underline{partie} ``\underline{évanescente}'' \underline{de la cohomologie de} $\Delta$, cf. 5.4.4), \underline{on a une décomposition en somme directe} (cf. 4.2)

\begin{equation}\tag{5.7.1}
E_2^{1,n-1} = H^1(\mathbb{P}^1, R^{n-1}\rho_*\mathbb{Q}_\ell) \;\xleftarrow[f^*+i_*g^*]{\sim}\; \mathrm{Prim}^n(X,\mathbb{Q}_\ell) \oplus E^{n-2}(\Delta,\mathbb{Q}_\ell(-1)) \quad .
\end{equation}

\SGApage{II}{299}

\underline{Démonstration.} Par dégénérescence on a

\begin{equation*}
E_2^{1,n-1} \quad \simeq \quad F^1 H^n(\tilde{X},\mathbb{Q}_\ell) \Big/ F^2 H^n(\tilde{X},\mathbb{Q}_\ell) \quad .
\end{equation*}
On a par (A) :

\SGAvar{\begin{equation}\tag{5.7.2}
F^1 H^n(\tilde{X},\mathbb{Q}_\ell) = \mathrm{Ker}\,(k^* : H^n(\tilde{X},\mathbb{Q}_\ell) \longrightarrow H^n(X_{\bar{\eta}},\mathbb{Q}_\ell)) \quad ,
\end{equation}
(par 5.1.2) $\simeq \mathrm{Ker}: (H^n(X,\mathbb{Q}_\ell) \oplus H^{n-2}(\Delta,\mathbb{Q}_\ell(-1)) \xrightarrow{m^*+h_*} H^n(X_{\bar{\eta}},\mathbb{Q}_\ell))$ .

}{\begin{equation}\tag{5.7.2}
F^1 H^n(\tilde{X},\mathbb{Q}_\ell) = \mathrm{Ker}\,(k^* : H^n(\tilde{X},\mathbb{Q}_\ell) \longrightarrow H^n(X_{\bar{\eta}},\mathbb{Q}_\ell)) \quad ,
\end{equation}
(par 5.1.2) $\simeq \mathrm{Ker}\,(H^n(X,\mathbb{Q}_\ell) \oplus H^{n-2}(\Delta,\mathbb{Q}_\ell(-1)) \xrightarrow{m^*+h_*} H^n(X_{\bar{\eta}},\mathbb{Q}_\ell))$ .

}{F7II-XVIII-063}Par (LV), $k^*$ se récrit encore

\begin{equation}\tag{5.7.3}
\mathrm{Prim}^n(X,\mathbb{Q}_\ell) \oplus H^{n-2}(X,\mathbb{Q}_\ell(-1)) \oplus H^{n-2}(\Delta,\mathbb{Q}_\ell(-1)) \xrightarrow{\;m^*+m^*L+h_*\;} H^n(X_{\bar{\eta}},\mathbb{Q}_\ell),
\end{equation}
et la flèche du milieu

\begin{equation}\tag{5.7.4}
H^{n-2}(X,\mathbb{Q}_\ell(-1)) \xrightarrow{\;m^*L\;} H^n(X_{\bar{\eta}},\mathbb{Q}_\ell)
\end{equation}
est un isomorphisme (cf. 5.2.3 et (5.2.2.3)), de sorte que, désignant par

\begin{equation}\tag{5.7.5}
\wp : H^n(X,\mathbb{Q}_\ell) \longrightarrow \mathrm{Prim}^n(X,\mathbb{Q}_\ell)
\end{equation}
la projection sur la partie primitive, on a

\SGAvar{\begin{equation}\tag{5.7.6}
F^1 H^n(\tilde{X},\mathbb{Q}_\ell) \;\xrightarrow[\;\wp f_*\oplus(-g_*2^*)\;]{\sim}\; \mathrm{Prim}^n(X,\mathbb{Q}_\ell) \oplus H^{n-2}(\Delta,\mathbb{Q}_\ell(-1)) \quad .
\end{equation}
Par (5.6.2.2) et dégénérescence, on a

}{\begin{equation}\tag{5.7.6}
F^1 H^n(\tilde{X},\mathbb{Q}_\ell) \;\xrightarrow[\;\wp f_*\oplus(-g_*i^*)\;]{\sim}\; \mathrm{Prim}^n(X,\mathbb{Q}_\ell) \oplus H^{n-2}(\Delta,\mathbb{Q}_\ell(-1)) \quad .
\end{equation}
Par (5.6.2.2) et dégénérescence, on a

}{F7II-XVIII-064}\begin{equation}\tag{5.7.7}
F^2 H^n(\tilde{X},\mathbb{Q}_\ell) \quad \simeq \quad E_2^{2,n-2} \quad \simeq \quad H^{n-2}(X_{\bar{\eta}},\mathbb{Q}_\ell(-1)) \quad .
\end{equation}

\SGApage{II}{300}

\SGAvar{D'autre part, l'inclusion de $F^2H^N(\widetilde{X},\mathbb{Q}_\ell) = H^{n-2}(X_{\overline{\eta}},\mathbb{Q}_\ell(-1))$ dans $H^n(\widetilde{X},\mathbb{Q}_\ell)$ est par la flèche $k_*$ de \underline{Gysin}. ( En toute riqueur, la flèche de Gysin n'est définie que comme une flèche

\begin{equation*}
H^{i-2}(X_{\overline{\eta}},\mathbb{Q}_\ell)) \longrightarrow H^i(X \times_k k(\overline{\eta}),\mathbb{Q}_\ell) \quad ,
\end{equation*}
mais compte tenu de ce qu'on a un isomorphisme canonique (SGA 4 XII 5.4) pour tout corps $K/k$ algébriquement clos

}{D'autre part, l'inclusion de $F^2H^n(\widetilde{X},\mathbb{Q}_\ell) = H^{n-2}(X_{\overline{\eta}},\mathbb{Q}_\ell(-1))$ dans $H^n(\widetilde{X},\mathbb{Q}_\ell)$ est par la flèche $k_*$ de \underline{Gysin}. (En toute rigueur, la flèche de Gysin n'est définie que comme une flèche

\begin{equation*}
H^{i-2}(X_{\overline{\eta}},\mathbb{Q}_\ell(-1)) \longrightarrow H^i(X \times_k k(\overline{\eta}),\mathbb{Q}_\ell) \quad ,
\end{equation*}
mais compte tenu de ce qu'on a un isomorphisme canonique (SGA 4 XII 5.4) pour tout corps $K/k$ algébriquement clos

}{F7II-XVIII-065}\begin{equation*}
H^i(X,\mathbb{Q}_\ell) \stackrel{\sim}{\longrightarrow} H^i(X \otimes_k K,\mathbb{Q}_\ell) \quad ,
\end{equation*}
nous espérons que le lecteur nous permettra de tels abus.) D'après (5.1.3), la flèche de Gysin s'exprime par

\begin{equation}\tag{5.7.8}
H^{n-2}(X_{\overline{\eta}},\mathbb{Q}_\ell(-1)) \stackrel{m_* \ \oplus(-h^*)}{\longrightarrow} H^n(X,\mathbb{Q}_\ell) \oplus H^{n-2}(\Delta,\mathbb{Q}_\ell(-1)) \quad .
\end{equation}
Par Lefschetz faible,

\begin{equation*}
m_* \ H^{n-2}(X_{\overline{\eta}},\mathbb{Q}_\ell(-1)) = m_*m^* \ H^{n-2}(X,\mathbb{Q}_\ell(-1)) = L_X H^{n-2}(X,\mathbb{Q}_\ell(-1))
\end{equation*}
est l'orthogonal de la partie primitive, de sorte qu'on a le diagramme commutatif

\SGAvar{\begin{equation}\tag{5.7.9}
\begin{tikzcd}[column sep=huge, row sep=huge]
F^2H^n(\widetilde{X},\mathbb{Q}_\ell) \arrow[r, hook] \& F^1H^n(\widetilde{X},\mathbb{Q}_\ell) \arrow[d, "\wr"', "\wp f_*\oplus -g_*i^*"] \\
H^{n-2}(X_{\overline{\eta}},\mathbb{Q}_\ell(-1)) \arrow[u, "\wr", "k_*"'] \arrow[r, "0 \oplus h^*"] \& \mathrm{Prim}(X,\mathbb{Q}_\ell)\oplus H^{n-2}(\Delta,\mathbb{Q}_\ell(-1))
\end{tikzcd}
\end{equation}











}{\begin{equation}\tag{5.7.9}
\begin{tikzcd}[column sep=huge, row sep=huge]
F^2H^n(\widetilde{X},\mathbb{Q}_\ell) \arrow[r, hook] \& F^1H^n(\widetilde{X},\mathbb{Q}_\ell) \arrow[d, "\wr"', "\wp f_*\oplus -g_*i^*"] \\
H^{n-2}(X_{\overline{\eta}},\mathbb{Q}_\ell(-1)) \arrow[u, "\wr", "k_*"'] \arrow[r, "0 \oplus h^*"] \& \mathrm{Prim}^n(X,\mathbb{Q}_\ell)\oplus H^{n-2}(\Delta,\mathbb{Q}_\ell(-1))
\end{tikzcd}
\end{equation}











}{F7II-XVIII-066}ce qui donne bien (5.7.1), cqfd.


\SGApage{II}{301}

\SGAvar{\underline{Remarque} 5.7.10. Dans 5.7, l'hypothèse (A) peut-être remplacée par celle, plus faible, que pour $q = n$, les flèches

}{\underline{Remarque} 5.7.10. Dans 5.7, l'hypothèse (A) peut être remplacée par celle, plus faible, que pour $q = n$, les flèches

}{F7II-XVIII-067}\begin{equation}\tag{5.7.11}
H^q(\widetilde{X},\mathbb{Q}_\ell) \longrightarrow H^0(\mathbb{P}^1,R^q\rho_*\mathbb{Q}_\ell)
\end{equation}
et

\begin{equation}\tag{5.7.12}
H^q(\widetilde{X},\mathbb{Q}_\ell) \longrightarrow H^q(X_{\overline{\eta}},\mathbb{Q}_\ell)
\end{equation}
aient même noyau, comme le montre notre démonstration. En l'absence de cette condition, la conclusion peut tomber en défaut, cf. 5.9.


5.8. \underline{L'homomorphisme de Griffiths.}

5.8.1. Soit $f: T \longrightarrow S$ un morphisme de schémas. La suite spectrale de Leray définit une filtration $(F^i)$ sur les $H^{\cdot}(T,\mathbb{Q}_\ell(j))$. On définit

\SGAvar{\begin{equation}\tag{5.8.1.1}
\mathrm{Prim}^{\cdot}(T/S,\mathbb{Q}_\ell(j)) \ = \ F^1H^{\cdot}(T,\mathbb{Q}_\ell(j)).
\end{equation}
Grâce au fait que la suite spectrale de Leray est contenu dans le ``premier quadrant'', on a

}{\begin{equation}\tag{5.8.1.1}
\mathrm{Prim}^{\cdot}(T/S,\mathbb{Q}_\ell(j)) \ = \ F^1H^{\cdot}(T,\mathbb{Q}_\ell(j)).
\end{equation}
Grâce au fait que la suite spectrale de Leray est contenue dans le ``premier quadrant'', on a

}{F7II-XVIII-068}\begin{equation}\tag{5.8.1.2}
E_\infty^{0,\cdot} \ \lhook\joinrel\longrightarrow \ E_2^{0,\cdot}
\end{equation}
de sorte que

\SGAvar{\begin{equation}\tag{5.8.1.3}
\mathrm{Prim}^{\cdot}(T_{/S},\mathbb{Q}_\ell(j) \ = \ \mathrm{Ker} \ (H^{\cdot}(T,\mathbb{Q}_\ell(j)) \longrightarrow H^0(S,R^{\cdot}f_*\mathbb{Q}_\ell(j)) \quad .
\end{equation}
Pour la même raison, on a

}{\begin{equation}\tag{5.8.1.3}
\mathrm{Prim}^{\cdot}(T/S,\mathbb{Q}_\ell(j)) \ = \ \mathrm{Ker} \bigl(H^{\cdot}(T,\mathbb{Q}_\ell(j)) \longrightarrow H^0(S,R^{\cdot}f_*\mathbb{Q}_\ell(j))\bigr) \quad .
\end{equation}
Pour la même raison, on a

}{F7II-XVIII-069}\begin{equation}\tag{5.8.1.4}
E_\infty^{1,\cdot} \ \lhook\joinrel\longrightarrow \ E_2^{1,\cdot} \quad .
\end{equation}

\SGAvar{\SGApage{II}{302}














}{\SGApage{II}{302}









}{}Le composé

\SGAvar{\begin{equation}\tag{5.8.1.5}
\begin{tikzcd}[column sep=large, row sep=large]
\mathrm{Prim}^i(\mathrm{T} /\mathrm{S},\mathbb{Q}_\ell(j)) = \mathrm{F}^1\mathrm{H}^i(\mathrm{T},\mathbb{Q}_\ell(j)) \arrow[r] \& \mathrm{gr}^1\mathrm{H}^i(\mathrm{T},\mathbb{Q}_\ell(j)) = \mathrm{E}_\infty^{1,i-1} \arrow[d, hook] \\
{} \& \mathrm{H}^1(\mathrm{S},\mathrm{R}^{i-1}_{f_*}\mathbb{Q}_\ell(j)) = \mathrm{E}_2^{1,i-1}
\end{tikzcd}
\end{equation}

}{\begin{equation}\tag{5.8.1.5}
\begin{tikzcd}[column sep=large, row sep=large]
\mathrm{Prim}^i(T/S,\mathbb{Q}_\ell(j)) = F^1H^i(T,\mathbb{Q}_\ell(j)) \arrow[r] \& \mathrm{gr}^1H^i(T,\mathbb{Q}_\ell(j)) = E_\infty^{1,i-1} \arrow[d, hook] \\
{} \& H^1(S,R^{i-1}f_*\mathbb{Q}_\ell(j)) = E_2^{1,i-1}
\end{tikzcd}
\end{equation}

}{F7II-XVIII-071}sera dénoté

\SGAvar{\begin{equation}\tag{5.8.1.6}
\mathrm{U}^i : \mathrm{Prim}(\mathrm{T}_{/\mathrm{S}},\mathbb{Q}_\ell(j)) \longrightarrow \mathrm{H}^1(\mathrm{S},\mathrm{R}^{\cdot-1}f_*\mathbb{Q}_\ell(j)) \quad .
\end{equation}

5.8.2. Soit maintenant X projective, lisse, et irréductible, D $= \{\mathrm{H}_\tau\}$ un pinceau d'hyperplansdont l'élément générique et l'axe coupent transversalement X. Désignons par U $\overset{\nu}{\longrightarrow} \mathbb{P}^1$ un ouvert de lissité non vide de la projection $\rho : \widetilde{\mathrm{X}} \longrightarrow \mathbb{P}^1$ et soit $\widetilde{\mathrm{X}}|\mathrm{U} = \widetilde{\mathrm{X}} \times_{\mathbb{P}^1} \mathrm{U}$ . On a le diagramme commutatif

}{\begin{equation}\tag{5.8.1.6}
U^i : \mathrm{Prim}^i(T/S,\mathbb{Q}_\ell(j)) \longrightarrow H^1(S,R^{i-1}f_*\mathbb{Q}_\ell(j)) \quad .
\end{equation}

5.8.2. Soit maintenant X projective, lisse, et irréductible, D $= \{\mathrm{H}_\tau\}$ un pinceau d'hyperplans dont l'élément générique et l'axe coupent transversalement X. Désignons par U $\overset{\nu}{\longrightarrow} \mathbb{P}^1$ un ouvert de lissité non vide de la projection $\rho : \widetilde{\mathrm{X}} \longrightarrow \mathbb{P}^1$ et soit $\widetilde{\mathrm{X}}|\mathrm{U} = \widetilde{\mathrm{X}} \times_{\mathbb{P}^1} \mathrm{U}$ . On a le diagramme commutatif

}{F7II-XVIII-072}\begin{equation}\tag{5.8.2.1}
\begin{tikzcd}[column sep=large, row sep=large]
{} \& \widetilde{\mathrm{X}} \arrow[dl, "f"] \arrow[d, "\rho"] \& \widetilde{\mathrm{X}}|\mathrm{U} \arrow[l, "\alpha"'] \arrow[d, "\rho|\mathrm{U}"] \\
\mathrm{X} \& \mathbb{P}^1 \& \mathrm{U} \arrow[l, "\nu"] \& .
\end{tikzcd}
\end{equation}


\SGAvar{\SGApage{II}{303}






























}{\SGApage{II}{303}





















}{}Le morphisme $\rho|\mathrm{U} : \widetilde{\mathrm{X}}|\mathrm{U} \longrightarrow \mathrm{U}$ étant propre et lisse, on a, par le théorème de la spécialisation de la cohomologie (SGA 4 XVI 2.2)

\SGAvar{\begin{equation}\tag{5.8.2.2}
\mathrm{Prim}^{\cdot}(\widetilde{\mathrm{X}}|\mathrm{U}/\mathrm{U},\mathbb{Q}_\ell(j) = \mathrm{Ker}(\mathrm{H}^{\cdot}(\widetilde{\mathrm{X}}|\mathrm{U},\mathbb{Q}_\ell(j)) \longrightarrow \mathrm{H}^{\cdot}(\mathrm{X}_{\overline{\eta}},\mathbb{Q}_\ell(j))) \quad ,
\end{equation}

}{\begin{equation}\tag{5.8.2.2}
\mathrm{Prim}^{\cdot}(\widetilde{\mathrm{X}}|\mathrm{U}/\mathrm{U},\mathbb{Q}_\ell(j)) = \mathrm{Ker}(\mathrm{H}^{\cdot}(\widetilde{\mathrm{X}}|\mathrm{U},\mathbb{Q}_\ell(j)) \longrightarrow \mathrm{H}^{\cdot}(\mathrm{X}_{\overline{\eta}},\mathbb{Q}_\ell(j))) \quad ,
\end{equation}

}{F7II-XVIII-074}donc, d'après (5.2.3.1), on a

\underline{Proposition} 5.8.2.3. \underline{Sous les hypothèses} (5.8.2), \underline{pour} $x \in \mathrm{H}^n(\mathrm{X},\mathbb{Q}_\ell(j))$, \underline{on a l'équivalence}

\begin{equation}\tag{5.8.2.4}
x \in \mathrm{Prim}^n(\mathrm{X},\mathbb{Q}_\ell(j)) \Longleftrightarrow \alpha^*f^*(x) \in \mathrm{Prim}^n(\widetilde{\mathrm{X}}|\mathrm{U}/\mathrm{U},\mathbb{Q}_\ell(j)) \quad .
\end{equation}

\underline{Si de plus} $\rho$ \underline{vérifie l'hypothèse} (A) \underline{de} 5.3, \underline{on a}

\begin{equation}\tag{5.8.2.5}
x \in \mathrm{Prim}^n(\mathrm{X},\mathbb{Q}_\ell(j)) \Longleftrightarrow f^*(x) \in \mathrm{Prim}^n(\widetilde{\mathrm{X}}/ \mathbb{P}^1,\mathbb{Q}_\ell(j)) \quad .
\end{equation}

\SGAvar{\underline{Démonstration}.(5.8.2.4) est prouvé. Quant à (5.8.2.5), l'hypothèse (A) implique que $\mathrm{H}^o(\mathbb{P}^1,\mathrm{R}^n\rho_*\mathbb{Q}_\ell(j)) \overset{\nu^*}{\underset{\sim}{\longrightarrow}} \mathrm{H}^o(\mathrm{U},\nu^*\mathrm{R}^n\rho_*\mathbb{Q}_\ell(j))$, donc[$f^*(x)$ est primitif rel. à $\mathbb{P}^1\alpha$] $\Longrightarrow \alpha^*f^*(\alpha)$ l'est rel. à U .

}{\underline{Démonstration}. (5.8.2.4) est prouvé. Quant à (5.8.2.5), l'hypothèse (A) implique que $H^0(\mathbb{P}^1,R^n\rho_*\mathbb{Q}_\ell(j)) \xrightarrow[\nu^*]{\sim} H^0(U,\nu^*R^n\rho_*\mathbb{Q}_\ell(j))$, donc $f^*(x)$ est primitif relativement à $\mathbb{P}^1$ si et seulement si $\alpha^*f^*(x)$ l'est relativement à $U$.

}{F7II-XVIII-075}\underline{Définition} 5.8.3. \underline{Dans la situation} 5.8.2, \underline{l'homomorphisme composé}

\SGAvar{\begin{equation}\tag{5.8.3.1}
\begin{tikzcd}[column sep=huge, row sep=huge]
\mathrm{Prim}^n(\mathrm{X},\mathbb{Q}_\ell(j)) \arrow[r, "\alpha^*f^*"] \arrow[dr, "\mathrm{Grif}"] \& \mathrm{Prim}^n(\widetilde{\mathrm{X}}|\mathrm{U}/\mathrm{U},\mathbb{Q}_\ell(j)) \arrow[d, "{[\mathrm{U}^n\,?]}"] \\
{} \& \mathrm{H}^1(\mathrm{U},\nu^*\mathrm{R}^{n-1}\rho_*\mathbb{Q}_\ell(j))
\end{tikzcd}
\end{equation}

}{\begin{equation}\tag{5.8.3.1}
\begin{tikzcd}[column sep=huge, row sep=huge]
\mathrm{Prim}^n(\mathrm{X},\mathbb{Q}_\ell(j)) \arrow[r, "\alpha^*f^*"] \arrow[dr, "\mathrm{Grif}"] \& \mathrm{Prim}^n(\widetilde{\mathrm{X}}|\mathrm{U}/\mathrm{U},\mathbb{Q}_\ell(j)) \arrow[d, "{[U^n]}"] \\
{} \& \mathrm{H}^1(\mathrm{U},\nu^*\mathrm{R}^{n-1}\rho_*\mathbb{Q}_\ell(j))
\end{tikzcd}
\end{equation}

}{F7II-XVIII-076}\underline{s'appelle} l'homomorphisme de Griffiths (\underline{pour} X, \underline{relatif au pinceau donné} D).


\SGApage{II}{304}

Ajoutons au diagramme (5.8.2.1) la famille constante au-dessus de $U$, de fibre $X$

\begin{equation}\tag{5.8.3.2}
\begin{tikzcd}[row sep=large, column sep=large]
 {} \& \widetilde{X} \arrow[dl, "f"'] \arrow[dd] \& \widetilde{X}|U \arrow[l, "\alpha"'] \arrow[r, "f\times\rho|U"] \arrow[dd, "\rho|U"] \& X\times U \arrow[ddl, "\mathrm{pr}_2"] \\
X \& {} \& {} \& {} \\
 {} \& \mathbb{P}^1 \& U \arrow[l, "\nu"'] \& .
\end{tikzcd}
\end{equation}


















\underline{Théorème} 5.8.4. \underline{Supposons que} $X$ \underline{vérifie} (LV) (5.2.2), \underline{et que} $\rho$ \underline{vérifie} \underline{l'hypothèse} (A) (5.3). \underline{Alors pour tout} $U \xrightarrow{\ \nu\ } \mathbb{P}^1$, \underline{ouvert de lissité} $\neq \emptyset$ \underline{de} $\rho$, \underline{l'homomorphisme de Griffiths} ``\underline{modulo partie fixe}'', \underline{qui est}, \underline{par définition}, \underline{le composé}

\SGAvar{\begin{equation}\tag{5.8.4.1}
\begin{tikzcd}[row sep=huge, column sep=huge]
\mathrm{Prim}^n(X,\mathbb{Q}_\ell(j)) \arrow[r, "\mathrm{Grif}"] \& H^1(U,\nu^*R^{n-1}\rho_*\mathbb{Q}_\ell(j)) \arrow[d] \\
 {} \& H^1(U,E^{n-1}\rho_*(\mathbb{Q}_\ell(j)) \;\simeq\; \dfrac{H^1(U,\nu^*R^{n-1}\rho_*\mathbb{Q}_\ell(j))}{(f\times\rho|U)^*H^1(U,H^{n-1}(X,\mathbb{Q}_\ell(j))_U)}
\end{tikzcd}
\end{equation}












}{\begin{equation}\tag{5.8.4.1}
\begin{tikzcd}[row sep=huge, column sep=huge]
\mathrm{Prim}^n(X,\mathbb{Q}_\ell(j)) \arrow[r, "\mathrm{Grif}"] \& H^1(U,\nu^*R^{n-1}\rho_*\mathbb{Q}_\ell(j)) \arrow[d] \\
 {} \& H^1(U,\nu^*E^{n-1}\rho_*\mathbb{Q}_\ell(j)) \;\simeq\; \dfrac{H^1(U,\nu^*R^{n-1}\rho_*\mathbb{Q}_\ell(j))}{(f\times\rho|U)^*H^1(U,H^{n-1}(X,\mathbb{Q}_\ell(j))_U)}
\end{tikzcd}
\end{equation}












}{F7II-XVIII-077}\underline{est injectif}.


\SGApage{II}{305}

\underline{Démonstration.} Par (5.5.1), on a, au-dessus de $\mathbb{P}^1$ ,

\begin{equation}\tag{5.8.4.2}
R^{n-1}\rho_*\mathbb{Q}_\ell(j) \;\simeq\; H^{n-1}(X,\mathbb{Q}_\ell(j))_{\mathbb{P}^1} \oplus E^{n-1}\rho_*\mathbb{Q}_\ell(j) \quad ,
\end{equation}
et au-dessus de $U$, on a

\SGAvar{\begin{equation}\tag{5.8.4.3}
\nu^*R^{n-1}\rho_*\mathbb{Q}_\ell(j)) \;\simeq\; H^{n-1}(X,\mathbb{Q}_\ell(j))_U \oplus \nu^*E^{n-1}\rho_*\mathbb{Q}_\ell(j) \quad ,
\end{equation}
ce qui permet de récrire (5.8.4.1) comme indiqué :

}{\begin{equation}\tag{5.8.4.3}
\nu^*R^{n-1}\rho_*\mathbb{Q}_\ell(j) \;\simeq\; H^{n-1}(X,\mathbb{Q}_\ell(j))_U \oplus \nu^*E^{n-1}\rho_*\mathbb{Q}_\ell(j) \quad ,
\end{equation}
ce qui permet de récrire (5.8.4.1) comme indiqué :

}{F7II-XVIII-078}\begin{equation}\tag{5.8.4.4}
\begin{tikzcd}[row sep=huge, column sep=large]
\mathrm{Prim}^n(X,\mathbb{Q}_\ell(j)) \arrow[r, "\mathrm{Grif}"] \& H^1(U,\nu^*R^{n-1}\rho_*\mathbb{Q}_\ell(j)) \arrow[d] \\
 {} \& H^1(U,\nu^*E^{n-1}\rho_*\mathbb{Q}_\ell(j)) \quad .
\end{tikzcd}
\end{equation}










\SGAvar{D'après (5.7.1), la flèche (qui n'est défini que grâce à l'hypothèse (A))

\begin{equation}\tag{5.8.4.5}
u \;:\; \mathrm{Prim}^n(X,\mathbb{Q}_\ell(j)) \longrightarrow H^1(\mathbb{P}^1,R^{n-1}\rho_*\mathbb{Q}_\ell)
\end{equation}
est injective. La décomposition (5.8.4.2), compte tenu de ce que $\mathbb{P}^1$ est simplement connexe, donne

\begin{equation}\tag{5.8.4.6}
H^1(\mathbb{P}^1,R^{n-1}\rho_*\mathbb{Q}_\ell) \;\simeq\; H^1(\mathbb{P}^1,E^{n-1}\rho_*\mathbb{Q}_\ell(j)) \quad ,
\end{equation}
de sorte qu'on peut compléter (5.8.4.4) en un diagramme commutatif 


}{D'après (5.7.1), la flèche (qui n'est définie que grâce à l'hypothèse (A))

\begin{equation}\tag{5.8.4.5}
u \;:\; \mathrm{Prim}^n(X,\mathbb{Q}_\ell(j)) \longrightarrow H^1(\mathbb{P}^1,R^{n-1}\rho_*\mathbb{Q}_\ell(j))
\end{equation}
est injective. La décomposition (5.8.4.2), compte tenu de ce que $\mathbb{P}^1$ est simplement connexe, donne

\begin{equation}\tag{5.8.4.6}
H^1(\mathbb{P}^1,R^{n-1}\rho_*\mathbb{Q}_\ell(j)) \;\simeq\; H^1(\mathbb{P}^1,E^{n-1}\rho_*\mathbb{Q}_\ell(j)) \quad ,
\end{equation}
de sorte qu'on peut compléter (5.8.4.4) en un diagramme commutatif 


}{F7II-XVIII-079}\SGApage{II}{306}

\begin{equation}\tag{5.8.4.7}
\begin{tikzcd}[column sep=huge, row sep=huge]
\mathrm{Prim}^n(X,\mathbb{Q}_\ell(j)) \arrow[r, "\mathrm{Grif}"] \arrow[d, hook, "u"] \& H^1(U,\nu^*R^{n-1}\rho_*\mathbb{Q}_\ell(j)) \arrow[d] \\
H^1(\mathbb{P}^1,E^{n-1}\rho_*\mathbb{Q}_\ell(j)) \arrow[r, "\nu^*"] \& H^1(U,\nu^*E^{n-1}\rho_*\mathbb{Q}_\ell(j)) \quad .
\end{tikzcd}
\end{equation}











\SGAvar{Il suffit maintenant de voir que la flèche en bas est injective,ce qui résulte du fait

}{Il suffit maintenant de voir que la flèche en bas est injective, ce qui résulte du fait

}{F7II-XVIII-080}\begin{equation}\tag{5.8.4.8}
E^{n-1}\rho_*\mathbb{Q}_\ell(j) \;\simeq\; \nu_*\nu^*E^{n-1}\rho_*\mathbb{Q}_\ell(j)
\end{equation}
par le lemme (bien connu !) suivant :

\underline{Lemme} 5.8.5. \underline{Soient} $\nu : A \longrightarrow B$ \underline{un morphisme de topos}, \underline{et soit} $F$ \underline{un faisceau sur} $A$. \underline{Alors la flèche canonique}

\begin{equation}\tag{5.8.5.1}
H^1(B,\nu_*F) \longrightarrow H^1(A,F)
\end{equation}
\underline{est injective}.

\underline{Démonstration.} Dans une suite spectrale du premier quadrant, on a

\begin{equation}\tag{5.8.5.2}
E_2^{1,0} \;\simeq\; E_\infty^{1,0} \qquad .
\end{equation}
Donc pour la suite spectrale de Leray 

\begin{equation*}
E_2^{p,q} = H^p(B,R^q\nu_*F) \;\Longrightarrow\; H^{p+q}(A,F)
\end{equation*}

\SGApage{II}{307}

 on a

\begin{equation}\tag{5.8.5.3}
E_2^{1,0} \;=\; H^1(B,\nu_*F) \;\simeq\; E_\infty^{1,0} \;=\; F^1H^1(A,F) \qquad .
\end{equation}
Prenant $F = \nu^*E$, $E$ un faisceau sur $B$, on trouve

\begin{equation}\tag{5.8.5.4}
H^1(B,\ \nu_*\nu^*E) \lhook\joinrel\longrightarrow H^1(A,\nu^*E) \qquad .
\end{equation}
5.8.6. Soit $L/k(\eta)$ une extension finie du corps des fonctions rationnelles sur $\mathbb{P}^1$, et posons

\begin{equation}\tag{5.8.6.1}
S = \text{le normalisé de } \mathbb{P}^1 \text{ dans } L \qquad .
\end{equation}
Faisons le changement de base $S \longrightarrow \mathbb{P}^1$ dans la situation (5.8.3.2)

\begin{equation}\tag{5.8.6.1 bis}
\begin{tikzcd}[column sep=huge, row sep=large]
 {} \& \widetilde{X} \arrow[dl, "f"] \arrow[dd, "\rho"] \& \widetilde{X}_S \arrow[l, "\alpha_S"'] \arrow[dd, "\rho_S"] \\
X \& {} \& {} \\
 {} \& \mathbb{P}^1 \& S \arrow[l, "\pi"'] \quad .
\end{tikzcd}
\end{equation}











Pour tout ouvert $V \overset{\mu}{\hookrightarrow} S$ au-dessus duquel $\rho_S$ soit lisse, on obtient après le changement de base $V \overset{\pi\,\mu}{\longrightarrow} \mathbb{P}^1$ dans (5.8.3.2), 

\begin{equation}\tag{5.8.6.2}
\begin{tikzcd}[column sep=huge, row sep=large]
 {} \& \widetilde{X} \arrow[dl, "f"'] \arrow[dd, "\rho"] \& \widetilde{X}_V \arrow[l, "\alpha_V"'] \arrow[dd, "\rho_V"] \arrow[dr, "f\times\rho_V"] \& {} \\
X \& {} \& {} \& X\times V \arrow[dl, "\mathrm{pr}_2"] \\
 {} \& \mathbb{P}^1 \& V \arrow[l, "{\pi\,\mu}"'] \quad , \&
{} \end{tikzcd}
\end{equation}
















\SGAvar{\SGApage{II}{308}
























}{\SGApage{II}{308}























}{}et on a encore (cf. (5.8.2.3))

\SGAvar{\begin{equation}\tag{5.8.6.3}
(\alpha_{\mathrm{V}}f)^* \mathrm{Prim}^n(\mathrm{X},\mathbb{Q}_\ell(j)) \subset \mathrm{Prim}^n(\widetilde{\mathrm{X}}_{\mathrm{V}}/_{\mathrm{V}} , \mathbb{Q}_\ell(j)) \quad ,
\end{equation}

}{\begin{equation}\tag{5.8.6.3}
\alpha_V^*f^* \mathrm{Prim}^n(X,\mathbb{Q}_\ell(j)) \subset \mathrm{Prim}^n(\widetilde X_V/V,\mathbb{Q}_\ell(j)) \quad ,
\end{equation}

}{F7II-XVIII-082}de sorte qu'on peut parler encore de l'homomorphisme de Griffiths

\SGAvar{\begin{equation}\tag{5.8.6.4}
\begin{tikzcd}[column sep=huge]
\mathrm{Prim}^n(\mathrm{X},\mathbb{Q}_\ell(j)) \arrow[r, "\mathrm{Grif}"] \& \mathrm{H}^1(\mathrm{V},\mathrm{R}^{n-1}\rho_{\mathrm{V}_*}\mathbb{Q}_\ell(j))
\end{tikzcd}
\quad .
\end{equation}

}{\begin{equation}\tag{5.8.6.4}
\begin{tikzcd}[column sep=huge]
\mathrm{Prim}^n(X,\mathbb{Q}_\ell(j)) \arrow[r, "\mathrm{Grif}"] \& H^1(V,R^{n-1}\rho_{V*}\mathbb{Q}_\ell(j))
\end{tikzcd}
\quad .
\end{equation}

}{F7II-XVIII-083}\underline{Corollaire} 5.8.7. \underline{Avec les notations précédentes}, \underline{et sous les hypothèses de} 5.8.4, \underline{l'homomorphisme de Griffiths ``modulo partie fixe'', i.e. le composé}

\SGAvar{\begin{equation}\tag{5.8.7.1}
\begin{tikzcd}[column sep=huge, row sep=huge]
\mathrm{Prim}^n(\mathrm{X},\mathbb{Q}_\ell(j)) \arrow[r, "\mathrm{Grif}"] \& \mathrm{H}^1(\mathrm{V},\mathrm{R}^{n-1}\rho_{\mathrm{V}_*}\mathbb{Q}_\ell(j)) \arrow[d] \\
{} \& \mathrm{H}^1(\mathrm{V},(\pi\ \mu)^*\mathrm{E}^{n-1}\rho_*(\mathbb{Q}_\ell(j))) \ \simeq \ \dfrac{\mathrm{H}^1(\mathrm{V},\mathrm{R}^{n-1}\rho_{\mathrm{V}_*}\mathbb{Q}_\ell(j))}{(f\times\rho_{\mathrm{V}})^* \mathrm{H}^1(\mathrm{V},\mathrm{H}^{n-1}(\mathrm{X},\mathbb{Q}_\ell(j))_{\mathrm{V}})}
\end{tikzcd}
\end{equation}

}{\begin{equation}\tag{5.8.7.1}
\begin{tikzcd}[column sep=huge, row sep=huge]
\mathrm{Prim}^n(X,\mathbb{Q}_\ell(j)) \arrow[r, "\mathrm{Grif}"] \& H^1(V,R^{n-1}\rho_{V*}\mathbb{Q}_\ell(j)) \arrow[d] \\
{} \& H^1(V,(\pi\mu)^*E^{n-1}\rho_*\mathbb{Q}_\ell(j)) \ \simeq \ \dfrac{H^1(V,R^{n-1}\rho_{V*}\mathbb{Q}_\ell(j))}{(f\times\rho_V)^* H^1(V,H^{n-1}(X,\mathbb{Q}_\ell(j))_V)}
\end{tikzcd}
\end{equation}

}{F7II-XVIII-084}\underline{est injectif}.

\underline{Démonstration}. Par le théorème de la spécialisation pour la cohomologie, on a

\SGAvar{\begin{equation*}
\mathrm{R}^{n-1}\rho_{\mathrm{V}_*}\mathbb{Q}_\ell(j)) = (\pi\ \mu)^* \ \mathrm{R}^{n-1}\rho_*\mathbb{Q}_\ell(j)
\end{equation*}


}{\begin{equation*}
R^{n-1}\rho_{V*}\mathbb{Q}_\ell(j) = (\pi\mu)^*R^{n-1}\rho_*\mathbb{Q}_\ell(j)
\end{equation*}


}{F7II-XVIII-085}\SGApage{II}{309}













de sorte que par la décomposition (5.8.4.2), (5.8.7.1) se récrit bien comme indiqué

\begin{equation*}
\begin{tikzcd}[column sep=huge, row sep=huge]
\mathrm{Prim}^n(\mathrm{X},\mathbb{Q}_\ell(j)) \arrow[r, "\mathrm{Grif}"] \& \mathrm{H}^1(\mathrm{V},(\pi\ \mu)^*\mathrm{R}^{n-1}\rho_*\mathbb{Q}_\ell(j)) \arrow[d] \\
{} \& \mathrm{H}^1(\mathrm{V},(\pi\ \mu)^* \ \mathrm{E}^{n-1}\rho_*\mathbb{Q}_\ell(j)) \quad .
\end{tikzcd}
\end{equation*}

Il suffit donc de démontrer l'injectivité de

\SGAvar{\begin{equation*}
\begin{tikzcd}[column sep=huge]
\mathrm{Prim}^n(\mathrm{X},\mathbb{Q}_\ell(j)) \arrow[r] \& \mathrm{H}^1(\mathrm{V},\mu^* \ \pi^* \ \mathrm{E}^{n-1}\rho_*\mathbb{Q}_\ell(j))
\end{tikzcd}
\end{equation*}

}{\begin{equation}\tag{5.8.7.3}
\begin{tikzcd}[column sep=huge]
\mathrm{Prim}^n(\mathrm{X},\mathbb{Q}_\ell(j)) \arrow[r] \& \mathrm{H}^1(\mathrm{V},\mu^* \ \pi^* \ \mathrm{E}^{n-1}\rho_*\mathbb{Q}_\ell(j))
\end{tikzcd}
\end{equation}

}{F7II-XVIII-086}pour $\mathrm{V} \overset{\mu}{\hookrightarrow} \mathrm{S}$ arbitrairement petit. On est ramené au cas

\begin{equation*}
\left\{
\begin{aligned}
\mathrm{U} \ &\overset{\nu}{\longrightarrow} \ \mathbb{P}^1 \quad \text{un ouvert de lissité de } \rho \\
\mathrm{V} \ &\overset{\pi_{\mathrm{U}}}{\longrightarrow} \ \mathrm{U} \quad \text{le normalisé de U dans L} \quad .
\end{aligned}
\right.
\end{equation*}

Du diagramme commutatif cartésien

\SGAvar{\begin{equation*}
\begin{tikzcd}[column sep=huge, row sep=huge]
\mathrm{V} \arrow[r, hook, "\mu"] \arrow[d, "\pi"] \& \mathrm{S} \arrow[d, "\pi"] \\
\mathrm{U} \arrow[r, hook, "\nu"] \& \mathbb{P}^1
\end{tikzcd}
\end{equation*}


}{\begin{equation*}
\begin{tikzcd}[column sep=huge, row sep=huge]
\mathrm{V} \arrow[r, hook, "\mu"] \arrow[d, "\pi_U"] \& \mathrm{S} \arrow[d, "\pi"] \\
\mathrm{U} \arrow[r, hook, "\nu"] \& \mathbb{P}^1
\end{tikzcd}
\end{equation*}


}{F7II-XVIII-087}\SGApage{II}{310}

 on conclut que

\SGAvar{\begin{equation*}
\mu^*\pi^*\, E^{n-1}\rho_*\mathbb{Q}_\ell(j) \;=\; \pi_U^*\nu^*\, E^{n-1}\rho_*\mathbb{Q}_\ell
\end{equation*}
de sorte que (5.8.7.3) se factorise

}{\begin{equation*}
\mu^*\pi^*\, E^{n-1}\rho_*\mathbb{Q}_\ell(j) \;=\; \pi_U^*\nu^*\, E^{n-1}\rho_*\mathbb{Q}_\ell(j)
\end{equation*}
de sorte que (5.8.7.3) se factorise

}{F7II-XVIII-088}\begin{center}
\begin{tikzcd}[column sep=huge, row sep=huge]
\mathrm{Prim}^n(X,\mathbb{Q}_\ell(j)) \arrow[r] \& H^1(U,\nu^*\, E^{n-1}\rho_*\mathbb{Q}_\ell(j)) \arrow[d, "\pi_U^*"] \\
 {} \& H^1(V,\pi_U^*\nu^*\, E^{n-1}\rho_*\mathbb{Q}_\ell(j)) \quad .
\end{tikzcd}
\end{center}









Or, il suffit, grâce à (5.8.4), de vérifier que

\begin{equation*}
\pi_U^* \;:\; H^1(U,\nu^*\, E^{n-1}\rho_*\mathbb{Q}_\ell(j)) \longrightarrow H^1(V,\pi_U^*\nu^*\, E^{n-1}\rho_*\mathbb{Q}_\ell(j))
\end{equation*}
est injectif. Comme $\pi_U : V \longrightarrow U$ est fini et plat, la théorie de la trace [SGA 4 XVII 6.2] nous fournit, pour tout $\mathbb{Q}_\ell$-faisceau $F$ sur $U$ (p.ex. $F = \nu^*\, E^{n-1}\rho_*\mathbb{Q}_\ell(j)$) une application

\SGAvar{\begin{equation*}
\mathrm{tr}_{\pi_U} : H^1(\ V,\pi_U^*\ ) \longrightarrow H^1(U,F(j))
\end{equation*}
qui vérifie

\begin{equation*}
\mathrm{tr}_{\pi} \;\circ\; \pi^* \;=\; \deg(\pi_U).\mathrm{id} \quad ,
\end{equation*}
ce qui implique l'injectivité de $H^*(U,F) \longrightarrow H^*(V,\pi_U^*F)$, et achève la démonstration.


\SGApage{II}{311}


5.9. \underline{Un exemple où l'hypothèse} (A) \underline{n'est pas vérifiéé}.

Soit $X$ une surface quadratique lisse dans $\mathbb{P}^3$. Alors $X$ vérifie (LV) (5.2.4). Prenons un pinceau (de Lefschetz, si l'on veut) d'\underline{hyperplans} dont l'élément générique et l'axe coupent transversalement $X$ . Alors (comme m'a fait observer Deligne) \underline{la projection} $\rho : \widetilde{X} \longrightarrow \mathbb{P}^1$ \underline{ne vérifie pas} (A). En effet, $R^2\rho_*\mathbb{Q}_\ell$ n'est pas constant, car les fibres singulières sont des courbes coniques singulières, donc sont réductibles, et pour un pinceau de Lefschetz, sont chacun la réunion de deux droites qui se coupent en un point (c'est la définition d'un pinceau de Lefschetz dans ce cas simple). Donc

}{\begin{equation*}
\mathrm{tr}_{\pi_U} : H^1(V,\pi_U^*F) \longrightarrow H^1(U,F)
\end{equation*}
qui vérifie

\begin{equation*}
\mathrm{tr}_{\pi_U} \;\circ\; \pi_U^* \;=\; \deg(\pi_U)\,\mathrm{id} \quad ,
\end{equation*}
ce qui implique l'injectivité de $H^*(U,F) \longrightarrow H^*(V,\pi_U^*F)$, et achève la démonstration.


\SGApage{II}{311}



5.9. \underline{Un exemple où l'hypothèse} (A) \underline{n'est pas vérifiée}.

Soit $X$ une surface quadratique lisse dans $\mathbb{P}^3$. Alors $X$ vérifie (LV) (5.2.4). Prenons un pinceau (de Lefschetz, si l'on veut) d'\underline{hyperplans} dont l'élément générique et l'axe coupent transversalement $X$ . Alors (comme m'a fait observer Deligne) \underline{la projection} $\rho : \widetilde{X} \longrightarrow \mathbb{P}^1$ \underline{ne vérifie pas} (A). En effet, $R^2\rho_*\mathbb{Q}_\ell$ n'est pas constant, car les fibres singulières sont des courbes coniques singulières, donc sont réductibles, et pour un pinceau de Lefschetz, sont chacune la réunion de deux droites qui se coupent en un point (c'est la définition d'un pinceau de Lefschetz dans ce cas simple). Donc

}{F7II-XVIII-089}\begin{equation*}
\dim_{\mathbb{Q}_\ell}\left(\, H^2\,(\text{fibre sing.},\ \mathbb{Q}_\ell)\right) = 2
\end{equation*}
bien que

\begin{equation*}
\dim_{\mathbb{Q}_\ell}\left(\, H^2\,(\text{fibre lisse},\ \mathbb{Q}_\ell)\right) = 1 \quad .
\end{equation*}
De plus, la conclusion du théorème 5.8.4 est ici fausse. Désignons par $U \stackrel{\nu}{\hookrightarrow} \mathbb{P}^1$ un ouvert de lissité $\neq \emptyset$ de $\rho$ . Comme les fibres de $\rho$ au-dessus de $U$ sont des courbes lisses de genre zéro, on a

\begin{equation*}
\nu^*R^1\rho_*\mathbb{Q}_\ell \;=\; 0
\end{equation*}
et donc

\begin{equation*}
H^1(U,\nu^*R^1\rho_*\mathbb{Q}_\ell) \;=\; 0 \quad , 
\end{equation*}

\SGApage{II}{312}













de sorte que l'homomorphisme de Griffiths

\begin{equation*}
\begin{tikzcd}[column sep=huge]
\mathrm{Prim}^2(\mathrm{X},\mathbb{Q}_\ell) \arrow[r, "\mathrm{Grif}"] \& \mathrm{H}^1(\mathrm{U},\nu^*\mathrm{R}^1\rho_*\mathbb{Q}_\ell)
\end{tikzcd}
\end{equation*}

est nul. Mais sur une surface quadratique, on a

\begin{equation*}
\dim_{\mathbb{Q}_\ell} \left( \mathrm{Prim}^2(\mathrm{X},\mathbb{Q}_\ell) \right) = 1 \quad .
\end{equation*}

On voit de même que la conclusion de 5.7 est ici en défaut.

\SGAvar{\SGAsectiontoc{6. Application de la théorie locale des cycles evanescents (Picard-Lefschetz).}{6}
6. \underline{Application de la théorie locale des cycles evanescents} (\underline{Picard}-\underline{Lefschetz}).

}{\SGAsectiontoc{6. Application de la théorie locale des cycles évanescents (Picard-Lefschetz).}{6}
6. \underline{Application de la théorie locale des cycles évanescents} (\underline{Picard}-\underline{Lefschetz}).

}{F7II-XVIII-090}\underline{Proposition} 6.1. \underline{Désignons par} Z \underline{une hypersurface dans l'espace projectif} $\mathbb{P}^r_k$ , $r \geq 2$, k \underline{étant un corps algébriquement clos}. \underline{Soit}

\SGAvar{\begin{equation*}
i_o : \mathrm{D}_o \lhook\joinrel\longrightarrow \mathbb{P}^r
\end{equation*}

\underline{l'inclusion d'une droite qui ne coupe} Z \underline{qu'en son ouvert de lissité, et le coupe transversalement}. \underline{Choissisons un point géométrique} $\overline{a}$ \underline{de} $\mathrm{D}_o - \mathrm{D}_o \cap \mathrm{Z}$, \underline{et désignons par}

}{\begin{equation*}
i_o : \mathrm{D}_o \lhook\joinrel\longrightarrow \mathbb{P}^r
\end{equation*}


\underline{l'inclusion d'une droite qui ne coupe} Z \underline{qu'en son ouvert de lissité, et le coupe transversalement}. \underline{Choisissons un point géométrique} $\overline{a}$ \underline{de} $\mathrm{D}_o - \mathrm{D}_o \cap \mathrm{Z}$, \underline{et désignons par}

}{F7II-XVIII-091}\begin{equation*}
\pi_1^{\mathrm{mod}}(\mathrm{D}_o - \mathrm{D}_o \cap \mathrm{Z}, \ \overline{a})
\end{equation*}

\underline{le quotient de} $\pi_1(\mathrm{D}_o - \mathrm{D}_o \cap \mathrm{Z}, \overline{a})$ \underline{qui classifie les revêtements étales de} $\mathrm{D}_o - \mathrm{D}_o \cap \mathrm{Z}$ \underline{qui sont modérément ramifiés en tout point de} $\mathrm{D}_o \cap \mathrm{Z}$ .

\underline{Alors}


\SGApage{II}{313}











6.1.1. \underline{L'homomorphisme canonique}

\begin{equation}\tag{6.1.1.1}
i_{o_*} : \pi_1^{\mathrm{mod}}(\mathrm{D}_o - \mathrm{D}_o \cap \mathrm{Z}, \overline{a}) \longrightarrow \pi_1^{\mathrm{mod}}(\mathbb{P}^r - \mathrm{Z}, \overline{a})
\end{equation}

\underline{est surjectif}.

6.1.2. \underline{Chaque choix d'un point géométrique} $\overline{x}_o \in \mathrm{D}_o \cap \mathrm{Z}$ \underline{et d'un} ``\underline{chemin}'' \underline{dans} $\mathrm{D}_o$ \underline{entre} $\overline{x}_o$ \underline{et} $\overline{a}$ \underline{détermine un homomorphisme}

\begin{equation}\tag{6.1.2.1}
\prod\limits_{\ell \neq \mathrm{car}(k)} \mathbb{Z}_\ell(1) \longrightarrow \pi_1^{\mathrm{mod}}(\mathrm{D}_o - \mathrm{D}_o \cap \mathrm{Z}, \overline{a})
\end{equation}

\underline{dont l'image est} ``\underline{le}'' \underline{groupe d'inertie en} $\overline{x}_o$ . \underline{Prenant le composé de ce morphisme et de} $i_{o_*}$ , \underline{on trouve un homomorphisme}

\begin{equation}\tag{6.1.2.2}
\prod\limits_{\ell \neq \mathrm{car}(k)} \mathbb{Z}_\ell(1) \longrightarrow \pi_1^{\mathrm{mod}}(\mathbb{P}^r - \mathrm{Z}) \quad .
\end{equation}

\SGAvar{\underline{Si} Z \underline{est irréductible, tous ces homomorphismes sont conjugués entre eux}.

\underline{Démonstration}. Désignons par T le sous-schéma de la grassmanniene formée des droites dans $\mathbb{P}^r$ qui

}{\underline{Si} Z \underline{est irréductible, tous ces homomorphismes sont conjugués entre eux}.


\underline{Démonstration}. Désignons par T le sous-schéma de la grassmannienne formé des droites dans $\mathbb{P}^r$ qui

}{F7II-XVIII-092}(i) passent par $\overline{a}$

\SGAvar{(ii) ne coupent Z qu'en son ouvert de lissité, et l'y coupent transversalement, par $\eta$ son point générique, et par $i_\eta : \mathrm{D}_\eta \longrightarrow \mathbb{P}^r \otimes_k k(\eta)$ la droite correspondante (qui est donc une droite ``joingnant $\overline{a}$ au point générique de $\mathbb{P}^r_k$ ''). Choisissons aussi un point géométrique $\overline{\eta}$ au-dessus de $\eta$, et notons par $\mathrm{D}_{\overline{\eta}}$ la droite ``géométrique'' correspondante. Soit o le point T correspondant à la droite donnée $\mathrm{D}_o$ .


}{(ii) ne coupent Z qu'en son ouvert de lissité, et l'y coupent transversalement, par $\eta$ son point générique, et par $i_\eta : \mathrm{D}_\eta \longrightarrow \mathbb{P}^r \otimes_k k(\eta)$ la droite correspondante (qui est donc une droite ``joignant $\overline{a}$ au point générique de $\mathbb{P}^r_k$ ''). Choisissons aussi un point géométrique $\overline{\eta}$ au-dessus de $\eta$, et notons par $\mathrm{D}_{\overline{\eta}}$ la droite ``géométrique'' correspondante. Soit o le point de T correspondant à la droite donnée $\mathrm{D}_o$ .


}{F7II-XVIII-093}\SGApage{II}{314}

\SGAvar{D'après (SGA 1, XIII 2.8), il existe un homomorphisme de spécialisation (dépendant d'un chemin entre $\bar{\eta}$ et $O$ dans $T$, donc défini à conjugaison près)

\begin{equation}\tag{6.1.3}
\pi_1^{\mathrm{mod}}(D_{\bar{\eta}} - D_{\bar{\eta}} \cap Z), \bar{a}) \xrightarrow{\ \bar{\eta} \longrightarrow o\ } \pi_1^{\mathrm{mod}}(D_o - D_o \cap Z, \bar{a})
\end{equation}
tel que l'on ait :

6.1.4. Les homomorphisme (6.1.2.1) de monodromie locale

\begin{equation*}
\prod_{\ell \neq \mathrm{car}(k)} \mathbb{Z}_{\ell}(1) \longrightarrow \pi_1(D_o - D_o \cap Z, \bar{a})
\end{equation*}
s'obtiennent en prenant le composé de la flèche de spécialisation (6.1.3) avec les homomorphismes de monodromie locale pour $D_{\bar{\eta}}$ ,

}{D'après (SGA 1, XIII 2.8), il existe un homomorphisme de spécialisation (dépendant d'un chemin entre $\bar{\eta}$ et $o$ dans $T$, donc défini à conjugaison près)

\begin{equation}\tag{6.1.3}
\pi_1^{\mathrm{mod}}(D_{\bar{\eta}} - D_{\bar{\eta}} \cap Z, \bar{a}) \xrightarrow{\ \bar{\eta} \longrightarrow o\ } \pi_1^{\mathrm{mod}}(D_o - D_o \cap Z, \bar{a})
\end{equation}
tel que l'on ait :


6.1.4. Les homomorphismes (6.1.2.1) de monodromie locale

\begin{equation*}
\prod_{\ell \neq \mathrm{car}(k)} \mathbb{Z}_{\ell}(1) \longrightarrow \pi_1^{\mathrm{mod}}(D_o - D_o \cap Z, \bar{a})
\end{equation*}
s'obtiennent en prenant le composé de la flèche de spécialisation (6.1.3) avec les homomorphismes de monodromie locale pour $D_{\bar{\eta}}$ ,

}{F7II-XVIII-094}\begin{equation}\tag{6.1.4.1}
\prod_{\ell \neq \mathrm{car}(k)} \mathbb{Z}_{\ell}(1) \longrightarrow \pi_1^{\mathrm{mod}}(D_{\bar{\eta}} - D_{\bar{\eta}} \cap Z, \bar{a}) \xrightarrow{\ \bar{\eta} \longrightarrow o\ } \pi_1^{\mathrm{mod}}(D_o - D_o \cap Z, \bar{a}) \ ;
\end{equation}
(6.1.5) le diagramme suivant est commutatif

\SGAvar{\begin{equation}\notag
\begin{tikzcd}[column sep=huge, row sep=huge]
 {} \& \pi_1^{\mathrm{mod}}(\mathbb{P}^r-Z,\bar{a}) \& {} \\
\pi_1^{\mathrm{mod}}(D_{\eta} - D_{\eta} \cap Z,\bar{a}) \arrow[ur, "i_{\eta_*}"] \& {} \& \pi_1^{\mathrm{mod}}(D_o - D_o \cap Z,\bar{a}) \arrow[ul, "i_{o_*}"'] \\
 {} \& \pi_1^{\mathrm{mod}}(D_{\bar{\eta}} - D_{\bar{\eta}} \cap Z,\bar{a}) \arrow[uu, "i_{\bar{\eta}_*}"'] \arrow[ul] \arrow[ur, "\bar{\eta} \longrightarrow o"'] \&
{} \end{tikzcd} \quad .
\end{equation}















}{\begin{equation}\notag
\begin{tikzcd}[column sep=huge, row sep=huge]
 {} \& \pi_1^{\mathrm{mod}}(\mathbb{P}^r-Z,\bar{a}) \& {} \\
\pi_1^{\mathrm{mod}}(D_{\eta} - D_{\eta} \cap Z,\bar{a}) \arrow[ur, "i_{\eta_*}"] \& {} \& \pi_1^{\mathrm{mod}}(D_o - D_o \cap Z,\bar{a}) \arrow[ul, "i_{o_*}"'] \\
 {} \& \pi_1^{\mathrm{mod}}(D_{\bar{\eta}} - D_{\bar{\eta}} \cap Z,\bar{a}) \arrow[uu, "i_{\bar{\eta}*}"'] \arrow[ul] \arrow[ur, "\bar{\eta} \longrightarrow o"'] \&
{} \end{tikzcd} \quad .
\end{equation}















}{F7II-XVIII-095}\SGApage{II}{315}

Admettons ceci, et prouvons d'abord 6.1.1. Or, par la commutativité de 6.1.5, il suffit de démontrer que l'homomorphisme

\begin{equation*}
i_{\bar{\eta}_*} \ : \ \pi_1^{\mathrm{mod}}(D_{\bar{\eta}} - D_{\bar{\eta}} \cap Z, \bar{a}) \longrightarrow \pi_1^{\mathrm{mod}}(\mathbb{P}^r - Z, \bar{a})
\end{equation*}
est surjectif, et pour ceci il suffit de démontrer que l'homomorphisme analogue des $\pi_1$ \underline{immodérés}

\SGAvar{\begin{equation}\tag{6.1.6}
i_{\bar{\eta}} \ : \ \pi_1(D_{\bar{\eta}} - D_{\bar{\eta}} \cap Z, \bar{a}) \longrightarrow \pi_1(\mathbb{P}^r - Z, \bar{a})
\end{equation}
est surjectif, i.e. que pour tout revêtement étale $E \longrightarrow \mathbb{P}^r - Z$, avec $E$ connexe (et donc irréductible, $\mathbb{P}^r - Z$ étant lisse, donc unibranche (EGA IV 18.10.3)), son image réciproque $E_{\bar{\eta}} = E \times_{\mathbb{P}^r - Z} (D_{\bar{\eta}} - D_{\bar{\eta}} \cap Z)$ au-dessus de $D_{\bar{\eta}} - D_{\bar{\eta}} \cap Z$ reste géométriquement irréductible. Pour ceci, il suffit d'utiliser la variante suivante du théorème de Bertini (EGA V ....) (pour la flèche composé $E \longrightarrow \mathbb{P}^r - Z \hookrightarrow \mathbb{P}^r$, et $s = 1$):

}{\begin{equation}\tag{6.1.6}
i_{\bar{\eta}*} \ : \ \pi_1(D_{\bar{\eta}} - D_{\bar{\eta}} \cap Z, \bar{a}) \longrightarrow \pi_1(\mathbb{P}^r - Z, \bar{a})
\end{equation}
est surjectif, i.e. que pour tout revêtement étale $E \longrightarrow \mathbb{P}^r - Z$, avec $E$ connexe (et donc irréductible, $\mathbb{P}^r - Z$ étant lisse, donc unibranche (EGA IV 18.10.3)), son image réciproque $E_{\bar{\eta}} = E \times_{\mathbb{P}^r - Z} (D_{\bar{\eta}} - D_{\bar{\eta}} \cap Z)$ au-dessus de $D_{\bar{\eta}} - D_{\bar{\eta}} \cap Z$ reste géométriquement irréductible. Pour ceci, il suffit d'utiliser la variante suivante du théorème de Bertini (EGA V ....) (pour la flèche composée $E \longrightarrow \mathbb{P}^r - Z \hookrightarrow \mathbb{P}^r$, et $s = 1$):

}{F7II-XVIII-096}6.1.6.1. \underline{Théorème de Bertini} : Soit $k$ un corps, $\mathbb{P}^r = \mathbb{P}^r_k$, et $f : X \longrightarrow \mathbb{P}^r$ un morphisme non constant, où $X$ est un schéma algébrique géométriquement irréductible. Pour chaque entier $s < \dim \overline{f(X)}$, désignons par $L_t$ le sous-espace linéaire \underline{générique} de codimension $s$ de $\mathbb{P}^r$ . Alors $X_t = f^{-1}(L_t)$ est géométriquement irréductible.

Etablissons maintenant 6.1.2. D'après 6.1.4, et la commutativité de (6.1.5), il suffit de démontrer que tous les homomorphismes de monodromie locale (6.1.3) \underline{pour} $D_{\bar{\eta}}$ 


\SGApage{II}{316}

\begin{equation}\tag{6.1.7.1}
 \prod_{\ell \neq \mathrm{car}(k)} \mathbb{Z}_{\ell}(1) \longrightarrow \pi_1^{\mathrm{mod}}(D_{\overline{\eta}} - D_{\overline{\eta}} \cap Z, \overline{a})
\end{equation}
deviennent conjugués après la composition avec l'homomorphisme

\SGAvar{\begin{equation}\tag{6.1.7.2}
\pi_1^{\mathrm{mod}}(D_{\overline{\eta}} - D_{\overline{\eta}} \cap Z, \overline{a}) \longrightarrow \pi_1^{\mathrm{mod}}(D_{\eta} - D_{\eta} \cap Z, \overline{a}) \; .
\end{equation}
Pour prouver ceci, remarquons que pour tout $K$ fini et galoisien sur $k(\eta)$, $k(\eta) \subset K \subset K(\overline{\eta})$, le schéma

\begin{equation*}
D_K \stackrel{\mathrm{def}}{=} D_{\eta} \times_{k(\eta)} K
\end{equation*}
est un revêtement étale de $D_{\eta}$ , de sorte que $D_K - D_K \cap Z$ est un revêtement modéré de $D_{\eta} - D_{\eta} \cap Z$, d'où on déduit une action de $\pi_1^{\mathrm{mod}}(D_{\overline{\eta}} - D_{\eta} \cap Z, \overline{a})$ sur le schéma $D_K \cap Z$ . Nous prenons $K$ suffisament grand pour que tous les points de $D_K \cap Z$ soient $K$-rationnels, ceci étant possible car par hypothèse le schéma $D_{\eta} \cap Z$ est \underline{étale} sur $k(\eta)$.

}{\begin{equation}\tag{6.1.7.2}
\pi_1^{\mathrm{mod}}(D_{\overline{\eta}} - D_{\overline{\eta}} \cap Z, \overline{a}) \longrightarrow \pi_1^{\mathrm{mod}}(D_{\eta} - D_{\eta} \cap Z, \overline{a}) \; .
\end{equation}
Pour prouver ceci, remarquons que pour tout $K$ fini et galoisien sur $k(\eta)$, $k(\eta) \subset K \subset k(\overline{\eta})$, le schéma

\begin{equation*}
D_K \stackrel{\mathrm{def}}{=} D_{\eta} \times_{k(\eta)} K
\end{equation*}
est un revêtement étale de $D_{\eta}$, de sorte que $D_K - D_K \cap Z$ est un revêtement modéré de $D_{\eta} - D_{\eta} \cap Z$, d'où on déduit une action de $\pi_1^{\mathrm{mod}}(D_{\eta} - D_{\eta} \cap Z, \overline{a})$ sur le schéma $D_K \cap Z$. Nous prenons $K$ suffisamment grand pour que tous les points de $D_K \cap Z$ soient $K$-rationnels, ceci étant possible car par hypothèse le schéma $D_{\eta} \cap Z$ est \underline{étale} sur $k(\eta)$.

}{F7II-XVIII-097}Il reste à démontrer que le groupe $\pi_1^{\mathrm{mod}}(D_{\eta}-D_{\eta} \cap Z, \overline{a})$ agit de façon \underline{transitive} sur les points de $D_K \cap Z$, pour établir la conjugaison des homomorphismes composés d'inertie locale

\begin{equation*}
\prod_{\ell \neq \mathrm{car}(k)} \mathbb{Z}_{\ell}(1) \longrightarrow \pi_1^{\mathrm{mod}}(D_{\eta} - D_{\eta} \cap Z, \overline{a})
\end{equation*}
pour les divers points de $D_K \cap Z$. Or, l'action sur $D_K \cap Z$ de $\pi_1^{\mathrm{mod}}(D_{\eta} - D_{\eta} \cap Z, \overline{a})$ se factorise par le quotient canonique

\SGAvar{\begin{equation*}
\pi_1^{\mathrm{mod}}(D_{\eta} - D_{\eta} \cap Z, \overline{a}) \longrightarrow \pi_1(\eta, \overline{\eta}) = \mathrm{Gal}(/k(\eta) \quad , 
\end{equation*}

}{\begin{equation*}
\pi_1^{\mathrm{mod}}(D_{\eta} - D_{\eta} \cap Z, \overline{a}) \longrightarrow \pi_1(\eta, \overline{\eta}) = \mathrm{Gal}(k(\overline{\eta})/k(\eta)) \quad ,
\end{equation*}

}{F7II-XVIII-098}\SGApage{II}{317}

\SGAvar{ et l'action de $\mathrm{Gal}(/k(\eta))$ sur $K_K \cap Z$ est celle qui est compatible avec l'inclusion

\begin{equation*}
D_K \cap Z \hookrightarrow \mathbb{P}^r_k \times_k K \quad ,
\end{equation*}
où $\mathrm{Gal}(/k(\eta))$ n'agit sur $\mathbb{P}^r_k \times_k K$ que via le deuxième facteur. Ceci posé, on voit que l'énoncé voulu équivaut à ce que le schéma $D_{\eta} \cap Z$ n'a qu'un point (i.e. que c'est le spectre d'un \underline{corps}).Or, en termes terre à terre, ceci n'est autre que le fait suivant (du style \underline{Bertini}) (où l'on prend pour $Z$ une équation affine $F = 0$, et on suppose que le point $\overline{a}$ est l'origine) :

}{et l'action de $\mathrm{Gal}(k(\overline{\eta})/k(\eta))$ sur $D_K \cap Z$ est celle qui est compatible avec l'inclusion

\begin{equation*}
D_K \cap Z \hookrightarrow \mathbb{P}^r_k \times_k K \quad ,
\end{equation*}
où $\mathrm{Gal}(k(\overline{\eta})/k(\eta))$ n'agit sur $\mathbb{P}^r_k \times_k K$ que via le deuxième facteur. Ceci posé, on voit que l'énoncé voulu équivaut à ce que le schéma $D_{\eta} \cap Z$ n'ait qu'un point (i.e. que ce soit le spectre d'un \underline{corps}). Or, en termes terre à terre, ceci n'est autre que le fait suivant (du style \underline{Bertini}) (où l'on prend pour $Z$ une équation affine $F = 0$, et on suppose que le point $\overline{a}$ est l'origine) :

}{F7II-XVIII-099}\underline{Fait} 6.1.8. \underline{Soit} $k$ \underline{un corps algébriquement clos}, $r \geq 2$, \underline{et} $F(X_1,\ldots,X_r) \in k[X_1,\ldots,X_r]$ \underline{un polynôme irréductible.} \underline{Alors le polynôme}

\SGAvar{\begin{equation*}
G(T) = F (X_1 T,\ldots,X_r T) \in k (X_1,\ldots,X_r)[T]
\end{equation*}
\underline{est irréductible} \underline{(en tant que polynôme en} $T$ \underline{à coefficients dans} $k(X_1,\ldots,X_r)$.

}{\begin{equation*}
G(T) = F (X_1 T,\ldots,X_r T) \in k (X_1,\ldots,X_r)[T]
\end{equation*}
\underline{est irréductible} \underline{(en tant que polynôme en} $T$ \underline{à coefficients dans} $k(X_1,\ldots,X_r)$).

}{F7II-XVIII-100}Ceci termine la démonstration de 6.1.2, cqfd.

\underline{Remarque} 6.1.9. Une autre méthode de démonstration de 6.1.2, valable sous des conditions nettement plus générales, consisterait à s'assurer simplement d'une compatibilité générale pour les ``homomorphismes de monodromie'', disant que les homomorphismes composés de 6.1.2 sont conjugués de l'homomorphisme de monodromie relatif au point générique de $Z$ (donc sont conjugués entre eux !).


\SGApage{II}{318}

\SGAvar{\underline{Corollaire} 6.1.10. \underline{Soit} $X$ \underline{une variété projective, lisse et connexe de dimension} $n$ \underline{sur un corps} $k$ \underline{algébriquement clos} $k$. \underline{Supposons que} $n = \dim X$ \underline{soit pair, ou que} $\mathrm{card}(k) \neq 2$. \underline{Soit} $D$ \underline{un pinceau de Lefschetz pour} $X$, $\eta$ \underline{son point générique}, $\bar{\eta}$ \underline{un point géométrique au-dessus de} $\eta$ . \underline{Alors la représentation naturelle de} $\pi_1(D - D \cap \check{X}, \bar{\eta})$ \underline{sur} $H^{n-1}(X_{\bar{\eta}},\mathbb{Q}_\ell)$ \underline{se factorise par} $\pi_1^{\mathrm{mod}}(D - D \cap \check{X}, \bar{\eta})$. \underline{Si cette représentation est non triviale, alors} $D \cap \check{X} \neq \emptyset$ \underline{et le groupe d'inertie en tout point de} $D \cap \check{X}$ \underline{agit non trivialement.}

\underline{Démonstration}. Si $\dim (\check{X}) \leq r-2$, alor $D \cap \check{X}$ est vide, et il n'y a rien à prouver, compte tenu que $\mathbb{P}^1$ est simplement connexe. Sinon, $\check{X}$ est une hypersurface dans $\mathbb{P}^r$ (XVII, 3.1.4), et, par 3.2.10 et XVII 3.5, 6.1.1 s'applique (prenant $\check{X}$ pour $Z$, $D$ pour $D_o$). Le fait que l'action de $\pi_1$ sur $H^{n-1}(X_{\bar{\eta}} , \mathbb{Q}_\ell)$ soit modérée a été vue dans la théorie locale, avec la \underline{formule de Picard-Lefschetz} (XV 3.4) , rappelée plus bas. Compte tenu du fait (non trivial!) que $\pi_1^{\mathrm{mod}}(D - D \cap \check{X}, \bar{\eta})$ est engendré par les groupes d'inertie en les différents points de $D \cap \check{X}$ (SGA 1 XIII 2.10), on a la conclusion voulue. Nous allons utiliser le résultat précédent pour étudier maintenant la condition (A) de (5.3.5). Rappelons d'abord le résultat suivant (XV 3.4) , conséquence facile de la formule de \underline{Picard-Lefschetz} :


}{\underline{Corollaire} 6.1.10. \underline{Soit} $X$ \underline{une variété projective, lisse et connexe de dimension} $n$ \underline{sur un corps} $k$ \underline{algébriquement clos}. \underline{Supposons que} $n = \dim X$ \underline{soit pair, ou que} $\mathrm{car}(k) \neq 2$. \underline{Soit} $D$ \underline{un pinceau de Lefschetz pour} $X$, $\eta$ \underline{son point générique}, $\bar{\eta}$ \underline{un point géométrique au-dessus de} $\eta$ . \underline{Alors la représentation naturelle de} $\pi_1(D - D \cap \check{X}, \bar{\eta})$ \underline{sur} $H^{n-1}(X_{\bar{\eta}},\mathbb{Q}_\ell)$ \underline{se factorise par} $\pi_1^{\mathrm{mod}}(D - D \cap \check{X}, \bar{\eta})$. \underline{Si cette représentation est non triviale, alors} $D \cap \check{X} \neq \emptyset$ \underline{et le groupe d'inertie en tout point de} $D \cap \check{X}$ \underline{agit non trivialement.}

\underline{Démonstration}. Si $\dim (\check{X}) \leq r-2$, alors $D \cap \check{X}$ est vide, et il n'y a rien à prouver, compte tenu que $\mathbb{P}^1$ est simplement connexe. Sinon, $\check{X}$ est une hypersurface dans $\mathbb{P}^r$ (XVII, 3.1.4), et, par 3.2.10 et XVII 3.5, 6.1.1 s'applique (prenant $\check{X}$ pour $Z$, $D$ pour $D_o$). Le fait que l'action de $\pi_1$ sur $H^{n-1}(X_{\bar{\eta}} , \mathbb{Q}_\ell)$ soit modérée a été vu dans la théorie locale, avec la \underline{formule de Picard-Lefschetz} (XV 3.4) , rappelée plus bas. Compte tenu du fait (non trivial!) que $\pi_1^{\mathrm{mod}}(D - D \cap \check{X}, \bar{\eta})$ est engendré par les groupes d'inertie en les différents points de $D \cap \check{X}$ (SGA 1 XIII 2.10), on a la conclusion voulue. Nous allons utiliser le résultat précédent pour étudier maintenant la condition (A) de (5.3.5). Rappelons d'abord le résultat suivant (XV 3.4) , conséquence facile de la formule de \underline{Picard-Lefschetz} :


}{F7II-XVIII-101}\SGApage{II}{319}

\underline{Proposition} 6.2. \underline{Soit} $X$ \underline{une variété projective lisse et connexe de dimension} $n$ \underline{sur le corps algébriquement clos} $k$. \underline{Désignons par} $D$ \underline{un pinceau de Lefschetz}, $\eta$ \underline{son point générique}, $\bar{\eta}$ \underline{un point géométrique au-dessus de} $\eta$ . \underline{Alors le pinceau} $D$ \underline{satisfait à la condition} (A) \underline{si et seulement si le groupe d'inertie en tout point de} $D \cap \check{X}$ \underline{agit non trivialement sur} $H^{n-1}(X_{\bar{\eta}},\mathbb{Q}_\ell)$.

\underline{Théorème} 6.3. \underline{Sous les hypothèses de} 6.2, \underline{on a} :

(6.3.1) \underline{Si} $n = \dim X$ \underline{est impair, la condition} (A) \underline{est vérifiée pour tout pinceau de Lefschetz pour} $X$.

(6.3.2) \underline{Supposons que} $n = \dim X$ \underline{soit pair, ou que} $\mathrm{car}(k) \neq 2$. \underline{Alors les conditions suivantes sont équivalentes} :

(i) $D$ \underline{ne satisfait pas à la condition} (A) (5.3.5)

(ii) $\pi_1(D-D \cap \check{X})$ \underline{agit trivialement sur} $H^{n-1}(X_{\bar{\eta}},\mathbb{Q}_\ell)$, \underline{et} $D \cap \check{X} \neq \emptyset$.

(iii) \underline{Il existe un point de} $D \cap \check{X}$ \underline{dont le groupe d'inertie agit trivialement sur} $H^{n-1}(X_{\bar{\eta}},\mathbb{Q}_\ell)$.

(iv) \underline{Le faisceau} $R^{n-1}\rho_*\mathbb{Q}_\ell$ \underline{est un faisceau constant sur} $D$ \underline{et} $D \cap \check{X}$ \underline{est non vide}.

\SGAvar{(v) \underline{Le faisceau} $R^{n-1}\rho_*\mathbb{Q}_\ell|U$ \underline{est un faisceau constant sur} $U$ (\underline{où} $U$ \underline{est un ouvert non vide} $D$ \underline{au-dessus duquel} $\rho$ \underline{soit lisse}), \underline{et} $\rho$ \underline{n'est pas lisse}.


}{(v) \underline{Le faisceau} $R^{n-1}\rho_*\mathbb{Q}_\ell|U$ \underline{est un faisceau constant sur} $U$ (\underline{où} $U$ \underline{est un ouvert non vide de} $D$ \underline{au-dessus duquel} $\rho$ \underline{soit lisse}), \underline{et} $\rho$ \underline{n'est pas lisse}.


}{F7II-XVIII-102}\SGApage{II}{320}

(6.3.3)  \underline{Si de plus} $X$ \underline{satisfait à la condition} (LV) (5.2.2), \underline{ces conditions} (i) -(v) \underline{sont aussi équivalentes à la condition} :

(vi)  \underline{la flèche de restriction}

\begin{equation*}
H^{n-1}(X,\mathbb{Q}_\ell) \longrightarrow H^{n-1}(X_{\bar{\eta}},\mathbb{Q}_\ell)
\end{equation*}
\underline{est surjective} \underline{(donc un isomorphisme par} "\underline{le théorème de Lefschetz faible}").

\underline{Démonstration.}  L'assertion 6.3.1  résulte de XV 3.4.

\SGAvar{Passons aux équivalences de 6.3.2 ; (i) $\Longleftrightarrow$ (iii) par 6.2, (ii) $\Longleftrightarrow$ (iii) par 6.1.6, (ii) $\Longleftrightarrow$ (v) par tautologie, et (iv) $\Longleftrightarrow$ (v) par 5.3.4. Enfin, si $X$ satisfait à la condition (LV), 5.6.10 donne l'équivalence (ii) $\Longleftrightarrow$ (vi), cqfd.

\underline{Remarque} 6.3.4. On voit que la condition (A) ne \underline{dépend pas du choix du pinceau de Lefschetz} $D$, mais seulement de l'immersion (de Lefschetz) $i : X \longrightarrow \mathbb{P}^r$ envisagée. C'est en effet clair par (6.3.1) si $\dim X$ est impair, et si $\dim X$ est pair, la condition (A) signifie aussi que le système local sur $\mathbb{P}^r - X$ formé des $H^{n-1}(X \cap H_s, Q_\ell)$ est constant, comme il résulte du critère (ii), et du fait que le système local envisagé est modérément ramifié en le point générique de $\check{X}$ et de la surjectivité (6.1.2).

\underline{Corollaire} 6.4.  \underline{Sous les hypothèses de} 6.2, \underline{supposons que} $n = \dim X$ \underline{est paire}, \underline{ou que} $\mathrm{car}(k) \neq 2$ \underline{et que} $X$ \underline{satisfait à la condition} (LV). \underline{Alors il existe un entier} $d_o$ \underline{tel que pour} $d \geq d_o$, \underline{tout pinceau de}
\SGApage{II}{321}
\underline{Lefschetz de sectionsde} $X$ \underline{par des hypersurfaces de degré} $d$ \underline{satisfait à la condition} (A).

}{Passons aux équivalences de 6.3.2 ; (i) $\Longleftrightarrow$ (iii) par 6.2, (ii) $\Longleftrightarrow$ (iii) par 6.1.10, (ii) $\Longleftrightarrow$ (v) par tautologie, et (iv) $\Longleftrightarrow$ (v) par 5.3.4. Enfin, si $X$ satisfait à la condition (LV), 5.6.10 donne l'équivalence (ii) $\Longleftrightarrow$ (vi), cqfd.

\underline{Remarque} 6.3.4. On voit que la condition (A) ne \underline{dépend pas du choix du pinceau de Lefschetz} $D$, mais seulement de l'immersion (de Lefschetz) $i : X \longrightarrow \mathbb{P}^r$ envisagée. C'est en effet clair par (6.3.1) si $\dim X$ est impair, et si $\dim X$ est pair, la condition (A) signifie aussi que le système local sur $\mathbb{P}^r - \check{X}$ formé des $H^{n-1}(X \cap H_s, \mathbb{Q}_\ell)$ est constant, comme il résulte du critère (ii), et du fait que le système local envisagé est modérément ramifié en le point générique de $\check{X}$ et de la surjectivité (6.1.1).

\underline{Corollaire} 6.4.  \underline{Sous les hypothèses de} 6.2, \underline{supposons que} ($n = \dim X$ \underline{soit pair}, \underline{ou que} $\mathrm{car}(k) \neq 2$) \underline{et que} $X$ \underline{satisfasse à la condition} (LV). \underline{Alors il existe un entier} $d_o$ \underline{tel que pour} $d \geq d_o$, \underline{tout pinceau de}
\SGApage{II}{321}
\underline{Lefschetz de sections de} $X$ \underline{par des hypersurfaces de degré} $d$ \underline{satisfait à la condition} (A).

}{F7II-XVIII-103}\underline{Démonstration.}  D'après l'équivalence (i) $\Longleftrightarrow$ (vi) de (6.3.2), il suffit de voir que, désignant par $Y(d)$ une section hypersurface lisse de $X$ de degré $d$, on a

\begin{equation}\tag{6.4.1}
\dim H^{n-1}(Y(d),\mathbb{Q}_\ell) \longrightarrow \infty \ \mbox{pour} \ d \longrightarrow \infty \qquad .
\end{equation}
En effet, on a le résultat précis suivant :

\underline{Lemme} 6.4.2. \underline{Soit} $X$ \underline{une variété projective}, \underline{lisse}, \underline{et connexe de dimension} $n \geq 1$ \underline{sur un corps algébriquement clos} $k$, \underline{et désignons par} $Y(d)$ \underline{une section hypersurface de degré} $d$. \underline{Alors il existe un polynôme} $f(T) \in \mathbb{Z}[T]$ \underline{tel que}

(6.4.2.1)  $f(T) = \deg(X)\ T^n +$ \underline{un polynôme de degré} $\leq n-1$

(6.4.2.2)  \underline{pour} $d \geq 1$, $f(d) = \dim H^{n-1}(Y(d), \mathbb{Q}_\ell)$ \qquad .

\underline{Démonstration.}  D'après (XVII 5.7.5), désignant par $C(X)$ la classe de Chern totale de $X$, et par $H$ la première classe de Chern d'un hyperplan, on a la formule suivante pour la caractéristique d'Euler-Poincaré de $Y(d)$

\SGAvar{\begin{equation}\tag{6.4.2.3}
\chi(Y(d) \;=\; \int\limits_X \frac{d\ H.C(X)}{1 + dH} \;=\; \sum_{i=1}^{n} \ (-1)^{i-1} d^i \ \int\limits_X H^i . C_{n-i} \qquad .
\end{equation}

}{\begin{equation}\tag{6.4.2.3}
\chi(Y(d)) \;=\; \int\limits_X \frac{d\ H.C(X)}{1 + dH} \;=\; \sum_{i=1}^{n} \ (-1)^{i-1} d^i \ \int\limits_X H^i . C_{n-i} \qquad .
\end{equation}

}{F7II-XVIII-104}\SGApage{II}{322}

On a, d'autre part, par dualité de Poincaré et Lefschetz faible

\begin{equation}\tag{6.4.2.4}
\chi(Y(d)) = (-1)^{n-1}\dim H^{n-1}(Y(d),\mathbb{Q}_\ell) + 2 \sum_{i \geq 1} (-1)^{n-1-i}\dim H^{n-1-i}(X,\mathbb{Q}_\ell).
\end{equation}
Mettant ensemble (6.4.2.3) et (6.4.2.4) donne le résultat voulu.


6.5. \underline{Un rappel sur la formule de Picard-Lefschetz} (XV 3.4).

\SGAvar{Soient $S$ le spectre d'un anneau de valuation hensélien strict $\eta$ son point générique, $\overline{\eta}$ un point géométrique au-dessus de $\eta$, et $s$ son point fermé. Soit

}{Soient $S$ le spectre d'un anneau de valuation hensélien strict, $\eta$ son point générique, $\overline{\eta}$ un point géométrique au-dessus de $\eta$, et $s$ son point fermé. Soit

}{F7II-XVIII-105}\begin{equation*}
\rho : X \longrightarrow S
\end{equation*}
un morphisme propre, avec $X$ un schéma régulier connexe de dimension $n$, tel que $\rho$ est \underline{lisse} en dehors d'un point $x$, point qui se trouve dans la fibre spéciale, et qui est une singularité quadratique ordinaire de la fibre spéciale.

Désignons par $I$ le groupe d'inertie de $S$ en $s$. On sait (XV 3.4) que $I$ agit trivialement sur $H^i(X_{\overline{\eta}},\mathbb{Q}_\ell)$ pour $i \neq n-1$, et que son action sur $H^{n-1}(X_{\overline{\eta}},\mathbb{Q}_\ell)$ est donnée par les formules suivantes :

6.5.1. \underline{Si} $n = 2m+1$, alors $I$ agit à travers un quotient isomorphe à $\mathbb{Z}/2\mathbb{Z}$ (NB. si car.$k \neq 2$, il y a \underline{un seul} tel quotient). Il existe un élément $\delta \in H^{n-1}(X_{\overline{\eta}},\mathbb{Q}_\ell(m))$, uniquement déterminé au signe près, qui s'appelle le \underline{cycle évanescent en le point} $s$, tel qu'on ait 


\SGApage{II}{323}

\SGAvar{\begin{equation}\tag{6.5.1.1}
 \delta.\delta = (-1)^m 2 \quad \text{dans} \quad H^{2n-2}(X_{\overline{\eta}},\mathbb{Q}_\ell(n-1)) \simeq \mathbb{Q}_\ell \, ,
\end{equation}
(6.5.2.2)\quad l'élément non trivial $\sigma$ de $\mathbb{Z}/2\mathbb{Z}$ agit sur $H^i(X_{\overline{\eta}},\mathbb{Q}_\ell)$ par la \underline{symétrie}

\begin{equation*}
\sigma x = x - (-1)^m (x\delta)\delta \, .
\end{equation*}
[L'élément $x.\delta$ est dans $H^{2n-2}(X_{\overline{\eta}},\mathbb{Q}_\ell(m) \simeq \mathbb{Q}_\ell(-m)$, donc le produit $(x.\delta)\delta$ définit un élément de $H^{n-1}(X_{\overline{\eta}},\mathbb{Q}_\ell)$]

6.5.2. \underline{Si} $n = 2m + 2$ , alors $I$ agit à travers son quotient canonique $\mathbb{Z}_\ell(1)$. Il existe un élément $\delta \in H^{n-1}(X_{\overline{\eta}},\mathbb{Q}_\ell(m)$ , élément canoniquement déterminé au signe près et qui s'appelle le \underline{cycle évanescent en le point} $s$, tel que $\sigma \in \mathbb{Z}_\ell(1)$ agit sur $\chi \in H^{n-1}(X_{\overline{\eta}},\mathbb{Q}_\ell)$ par la \underline{transvection}

\begin{equation}\tag{6.5.2.1}
\sigma(x) = x - (-1)^m \sigma.(x.\delta)\delta \, .
\end{equation}
[Le produit $x.\delta$ est dans $H^{2n-2}(X_{\overline{\eta}},\mathbb{Q}_\ell(m) \simeq \mathbb{Q}_\ell(-m-1)$, donc son produit par $\sigma \in \mathbb{Z}_\ell(1)$ définit un élément $\sigma.(x.\delta) \in \mathbb{Q}_\ell(-m)$ , et cet élément, multiplié par $\delta \in H^{n-1}(X_{\overline{\eta}},\mathbb{Q}_\ell(m)$ donne l'élément $\sigma.(x.\delta)\delta \in H^{n-1}(X_{\overline{\eta}},\mathbb{Q}_\ell).$]

}{\begin{equation}\tag{6.5.1.1}
 \delta.\delta = (-1)^m 2 \quad \text{dans} \quad H^{2n-2}(X_{\overline{\eta}},\mathbb{Q}_\ell(n-1)) \simeq \mathbb{Q}_\ell \, ,
\end{equation}
(6.5.1.2)\quad l'élément non trivial $\sigma$ de $\mathbb{Z}/2\mathbb{Z}$ agit sur $H^{n-1}(X_{\overline{\eta}},\mathbb{Q}_\ell)$ par la \underline{symétrie}

\begin{equation*}
\sigma x = x - (-1)^m (x\delta)\delta \, .
\end{equation*}
[L'élément $x.\delta$ est dans $H^{2n-2}(X_{\overline{\eta}},\mathbb{Q}_\ell(m)) \simeq \mathbb{Q}_\ell(-m)$, donc le produit $(x.\delta)\delta$ définit un élément de $H^{n-1}(X_{\overline{\eta}},\mathbb{Q}_\ell)$.]

6.5.2. \underline{Si} $n = 2m + 2$, alors $I$ agit à travers son quotient canonique $\mathbb{Z}_\ell(1)$. Il existe un élément $\delta \in H^{n-1}(X_{\overline{\eta}},\mathbb{Q}_\ell(m))$, élément canoniquement déterminé au signe près et qui s'appelle le \underline{cycle évanescent en le point} $s$, tel que $\sigma \in \mathbb{Z}_\ell(1)$ agit sur $x \in H^{n-1}(X_{\overline{\eta}},\mathbb{Q}_\ell)$ par la \underline{transvection}

\begin{equation}\tag{6.5.2.1}
\sigma(x) = x - (-1)^m \sigma.(x.\delta)\delta \, .
\end{equation}
[Le produit $x.\delta$ est dans $H^{2n-2}(X_{\overline{\eta}},\mathbb{Q}_\ell(m)) \simeq \mathbb{Q}_\ell(-m-1)$, donc son produit par $\sigma \in \mathbb{Z}_\ell(1)$ définit un élément $\sigma.(x.\delta) \in \mathbb{Q}_\ell(-m)$, et cet élément, multiplié par $\delta \in H^{n-1}(X_{\overline{\eta}},\mathbb{Q}_\ell(m))$, donne l'élément $\sigma.(x.\delta)\delta \in H^{n-1}(X_{\overline{\eta}},\mathbb{Q}_\ell)$.]

}{F7II-XVIII-106}Reprenons le cas global :

\SGAvar{\underline{Théorème} 6.6. \underline{Soit} $X$ \underline{une variété projective lisse et connexe de dimension} $n$ \underline{sur un corps algébriquement clos} $k$. \underline{On suppose} $n$ \underline{pair ou} car $k \neq 2$. \underline{Soient} $D$ \underline{un pinceau de Lefschetz pour} $X$, $\eta$ \underline{son point générique}, $\overline{\eta}$ \underline{un}
\SGApage{II}{324}
\underline{point géométrique au-dessus de} $\eta$, \underline{et} $s_{1},\ldots,s_{N}$ \underline{les points de} $D \cap \check{X}$. \underline{Choisissons par chaque} $i$ \underline{un isomorphisme} $\ell_{i}$ \underline{entre} $\bar{\eta}$ \underline{et le point générique géométrique de} $S_{i}$ = \underline{le hensélisé strict de} $D$ \underline{en} $s_{i}$ , \underline{de sorte que le groupe} $I_{i}$ \underline{d'inertie de} $\delta_{i}$ \underline{en} $S_{i}$ \underline{agit sur} $H^{n-1}(X_{\bar{\eta}},\mathbb{Q}_{\ell}(j))$ \underline{et définit le cycle évanescent} (6.5) \underline{en} $s_{i}$

\begin{equation*}
\delta_{i} \in H^{n-1}(X,\mathbb{Q}_{\ell}(m) \ , \ m \ = \ \left\{ \begin{array}{ll} \frac{n-1}{2} & \underline{\mbox{si}}\ n\ \underline{\mbox{impair}} \\[2ex] \frac{n}{2} - 1 & \underline{\mbox{si}}\ n\ \underline{\mbox{pair}} \ . \end{array} \right.
\end{equation*}
\underline{On suppose les} $\ell_{i}$ \underline{choisis de telle façon que les} $I_{i} \subset \pi = \pi_{1}^{\mathrm{mod}}(D-D \cap \check{X},\bar{\eta})$ \underline{engendrent topologiquement le groupe} $\pi_{1}$ \underline{(ce qui est possible, en vertu de} SGA 1 XIII 2.10). \underline{Ceci posé}, \underline{on a} :

6.6.1. \underline{Si} $X$ \underline{satisfait à la condition} (LV), \underline{l'espace de cohomologie évanescente} $E^{n-1}(X_{\bar{\eta}},\mathbb{Q}_{\ell}(m)$ ( $\stackrel{\mathrm{dfn}}{=}$ \underline{l'orthogonal de} $H^{n-1}(X,\mathbb{Q}_{\ell}(m))$ \underline{dans} $H^{n-1}(X_{\bar{\eta}},\mathbb{Q}_{\ell}(m)$, cf. 5.4.4) \underline{est engendré par les cycles évanescents} $\delta_{i}$, $i=1,\ldots,N$ .

}{\underline{Théorème} 6.6. \underline{Soit} $X$ \underline{une variété projective lisse et connexe de dimension} $n$ \underline{sur un corps algébriquement clos} $k$. \underline{On suppose} $n$ \underline{pair ou} car $k \neq 2$. \underline{Soient} $D$ \underline{un pinceau de Lefschetz pour} $X$, $\eta$ \underline{son point générique}, $\overline{\eta}$ \underline{un}
\SGApage{II}{324}
\underline{point géométrique au-dessus de} $\eta$, \underline{et} $s_{1},\ldots,s_{N}$ \underline{les points de} $D \cap \check{X}$. \underline{Choisissons pour chaque} $i$ \underline{un isomorphisme} $\ell_{i}$ \underline{entre} $\bar{\eta}$ \underline{et le point générique géométrique de} $S_{i}$ = \underline{le hensélisé strict de} $D$ \underline{en} $s_{i}$, \underline{de sorte que le groupe} $I_{i}$ \underline{d'inertie en} $s_i$, \underline{identifié au moyen de} $\ell_i$, \underline{agisse sur} $H^{n-1}(X_{\bar{\eta}},\mathbb{Q}_{\ell}(j))$ \underline{et définisse le cycle évanescent} (6.5) \underline{en} $s_{i}$

\begin{equation*}
\delta_{i} \in H^{n-1}(X_{\bar\eta},\mathbb{Q}_{\ell}(m)), \qquad m \ = \ \left\{ \begin{array}{ll} \frac{n-1}{2} & \underline{\mbox{si}}\ n\ \underline{\mbox{impair}}, \\[2ex] \frac{n}{2} - 1 & \underline{\mbox{si}}\ n\ \underline{\mbox{pair}}. \end{array} \right.
\end{equation*}
\underline{On suppose les} $\ell_{i}$ \underline{choisis de telle façon que les} $I_{i} \subset \pi = \pi_{1}^{\mathrm{mod}}(D-D \cap \check{X},\bar{\eta})$ \underline{engendrent topologiquement le groupe} $\pi$ \underline{(ce qui est possible, en vertu de} SGA 1 XIII 2.10). \underline{Ceci posé}, \underline{on a} :

6.6.1. \underline{Si} $X$ \underline{satisfait à la condition} (LV), \underline{l'espace de cohomologie évanescente} $E^{n-1}(X_{\bar{\eta}},\mathbb{Q}_{\ell}(m))$ ( $\stackrel{\mathrm{dfn}}{=}$ \underline{l'orthogonal de} $H^{n-1}(X,\mathbb{Q}_{\ell}(m))$ \underline{dans} $H^{n-1}(X_{\bar{\eta}},\mathbb{Q}_{\ell}(m))$, cf. 5.4.4) \underline{est engendré par les cycles évanescents} $\delta_{i}$, $i=1,\ldots,N$.

}{F7II-XVIII-107}6.6.2. \underline{Les cycles évanescents} $\delta_{i}$ \underline{sont conjugués, au signe près, sous le groupe de monodromie globale} $\pi_{1}(D - D \cap \check{X})$.

\underline{Démonstration}. Pour prouver 6.6.1, notons que, par la formule de Picard-Lefschetz 6.5, on a, posant $\pi = \pi_{1}(D - D \cap \check{X})$, et grâce au fait que les $I_{i}$ engendrent $\pi$ :


\SGApage{II}{325}

\SGAvar{\begin{equation}\tag{6.6.3}
\begin{array}{l} H^{n-1}(X_{\bar{\eta}},\mathbb{Q}_{\ell}(m)^{\pi} \ = \ \mbox{l'orthogonal dans } H^{n-1}(X_{\bar{\eta}},\mathbb{Q}_{\ell}(m)) \\[1.5ex] \hspace{7em} \mbox{du sous-espace engendré par les } \delta_{i} \ . \end{array}
\end{equation}
Par 5.6.10, on a

\begin{equation}\tag{6.6.4}
H^{n-1}(X_{\bar{\eta}},\mathbb{Q}_{\ell}(m))^{\pi} \ = \ \mathrm{Im}(H^{n-1}(X,\mathbb{Q}_{\ell}(i))) \qquad ,
\end{equation}
de sorte que prenant les orthogonaux des deux membres de (6.6.3), on trouve 6.6.1.

Pour prouver 6.6.2, on remarque simplement que le cycle évanescent $\delta_{i}$ est déterminé au signe près, par la représentation du groupe d'inertie en $s_{i}$ . Donc on n'a qu'a appliquer 6.1.2, qui donne la conjugaison dans $\pi_{1}^{\mathrm{mod}}(\mathbb{P}^{r} - X,\bar{\eta})$ des homomorphismes d'inertie composés

}{\begin{equation}\tag{6.6.3}
\begin{array}{l} H^{n-1}(X_{\bar{\eta}},\mathbb{Q}_{\ell}(m))^{\pi} \ = \ \mbox{l'orthogonal dans } H^{n-1}(X_{\bar{\eta}},\mathbb{Q}_{\ell}(m)) \\[1.5ex] \hspace{7em} \mbox{du sous-espace engendré par les } \delta_{i} \ . \end{array}
\end{equation}
Par 5.6.10, on a

\begin{equation}\tag{6.6.4}
H^{n-1}(X_{\bar{\eta}},\mathbb{Q}_{\ell}(m))^{\pi} \ = \ \mathrm{Im}\bigl(H^{n-1}(X,\mathbb{Q}_{\ell}(m)) \longrightarrow H^{n-1}(X_{\bar\eta},\mathbb{Q}_{\ell}(m))\bigr) \qquad ,
\end{equation}
de sorte que prenant les orthogonaux des deux membres de (6.6.3), on trouve 6.6.1.

Pour prouver 6.6.2, on remarque simplement que le cycle évanescent $\delta_{i}$ est déterminé au signe près par la représentation du groupe d'inertie en $s_{i}$. Donc on n'a qu'à appliquer 6.1.2, qui donne la conjugaison dans $\pi_{1}^{\mathrm{mod}}(\check{\mathbb{P}}^{r} - \check X,\bar{\eta})$ des homomorphismes d'inertie composés

}{F7II-XVIII-108}\begin{center}
\begin{tikzcd}[column sep=huge, row sep=huge]
\mathbb{Z}_{\ell}(1) \arrow[r] \& \pi_{1}^{\mathrm{mod}}(D - D \cap \check{X},\ \bar{\eta}) \arrow[d] \& {} \\
{} \& \pi_{1}^{\mathrm{mod}}(\check{\mathbb{P}}^{r} - \check{X},\bar{\eta}) \& .
\end{tikzcd}
\end{center}














Ceci donne la \underline{conjugaison} (à signe près) des cycles évanescents qu'ils déterminent, par des éléments de $\pi_{1}^{\mathrm{mod}}(\check{\mathbb{P}}^{r} - \check{X},\ \bar{\eta})$, et la démonstration s'achève par 6.1.1, qui donne la surjectivité de

\begin{equation*}
\pi_{1}^{\mathrm{mod}}(D - D \cap \check{X}) \longrightarrow \pi_{1}^{\mathrm{mod}}(\check{\mathbb{P}}^{r} - \check{X}) \qquad .
\end{equation*}

\SGApage{II}{326}

\underline{Corollaire} 6.7. \underline{Sous les hypothèses de} 6.6, \underline{supposons que} $X$ \underline{satisfait à la condition} (LV). \underline{Alors} :

6.7.1. \underline{La représentation de} $\pi_1(D - D \cap \check{X})$ \underline{sur} $E^{n-1}(X_{\bar{\eta}},\mathbb{Q}_\ell)$ \underline{est absolument irréductible.}

6.7.2. \underline{Si} $n$ \underline{est pair,} \underline{tout sous-groupe d'indice fini dans} $\pi_1(D - D \cap \check{X})$ \underline{agit de façon absolument irréductible sur} $E^{n-1}(X_{\bar{\eta}},\mathbb{Q}_\ell)$.

\SGAvar{\underline{Démonstration.} On peut évidemment remplacer le module envisagé par le tordu $E = E^{n-1}(X_{\bar{\eta}},\mathbb{Q}_\ell(m)$. Soit $\overline{E} = E \otimes_{\mathbb{Q}_\ell} \overline{\mathbb{Q}_\ell}$ et soit $\chi$ un élément $\neq 0$ de $\overline{E}$. Pour démontrer 6.7.1, il faut prouver que le sous $\overline{\mathbb{Q}_\ell}$-vectoriel de $\overline{E}$ engendré par $\chi$ et ses conjugués sous $\pi$ est $\overline{E}$ tout entier. D'après 6.6.1 et 6.6.2, il suffit de prouver que ce sous-module contient \underline{un} des cycles évanescents $\delta_i$ . Mais comme $\chi$ satisfait (LV), la restriction de la forme de cup-produit à $H^{n-1}(X,\mathbb{Q}_\ell(m)) \lhook\joinrel\relbar\joinrel\relbar H^{n-1}(X_{\bar{\eta}},\mathbb{Q}_\ell(m))$ est non dégénérée, donc aussi la restriction à l'orthogonal $E^{n-1}(X_{\bar{\eta}},\mathbb{Q}_\ell(m))$, donc en vertu de (6.6.1) il existe au moins une valeur de $i$ telle que l'on ait :

\begin{equation*}
x.\delta_i \neq 0 \quad .
\end{equation*}
Pour cet $i$, il existe par 6.5 un élément $\sigma_i$ du groupe d'inertie en $s_i$ tel que

\begin{equation}\tag{6.7.3}
\sigma_i(x) - x = \text{un multiple non nul de } \delta_i \ ,
\end{equation}
ce qui achève la démonstration de 6.7.1.


}{\underline{Démonstration.} On peut évidemment remplacer le module envisagé par le tordu $E = E^{n-1}(X_{\bar{\eta}},\mathbb{Q}_\ell(m))$. Soit $\overline{E} = E \otimes_{\mathbb{Q}_\ell} \overline{\mathbb{Q}}_\ell$ et soit $\chi$ un élément $\neq 0$ de $\overline{E}$. Pour démontrer 6.7.1, il faut prouver que le sous-espace $\overline{\mathbb{Q}}_\ell$-vectoriel de $\overline{E}$ engendré par $\chi$ et ses conjugués sous $\pi$ est $\overline{E}$ tout entier. D'après 6.6.1 et 6.6.2, il suffit de prouver que ce sous-espace contient \underline{un} des cycles évanescents $\delta_i$. Mais comme $X$ satisfait (LV), la restriction de la forme de cup-produit à $H^{n-1}(X,\mathbb{Q}_\ell(m)) \lhook\joinrel\relbar\joinrel\relbar H^{n-1}(X_{\bar{\eta}},\mathbb{Q}_\ell(m))$ est non dégénérée, donc aussi la restriction à l'orthogonal $E^{n-1}(X_{\bar{\eta}},\mathbb{Q}_\ell(m))$ ; par 6.6.1 il existe donc au moins une valeur de $i$ telle que l'on ait :

\begin{equation*}
\chi.\delta_i \neq 0 \quad .
\end{equation*}
Pour cet $i$, il existe par 6.5 un élément $\sigma_i$ du groupe d'inertie en $s_i$ tel que

\begin{equation}\tag{6.7.3}
\sigma_i(\chi) - \chi = \text{un multiple non nul de } \delta_i \ ,
\end{equation}
ce qui achève la démonstration de 6.7.1.


}{F7II-XVIII-109}\SGApage{II}{327}

\SGAvar{Quant à 6.7.2, on a $n-1$ impair, donc $\delta.\delta = 0$, et pour toute puissance $k$.ième de $\sigma_i$ on a, en itérant 6.5.2,

\begin{equation}\tag{6.7.4}
\sigma_i^k(x) = \ \ X + k\,\sigma_i(x.\delta_i)\,\delta_i \quad .
\end{equation}
Donc un sous $Q_\ell$-module $M$ de $E$ est stable sous l'action de $\pi$ si et seulement il l'est sous un sous-groupe d'indice fini, car les deux conditions équivalent à la condition

\begin{equation*}
x \in M \implies \forall_i \text{ tel que le cup-produit } x.\delta_i \neq 0, \text{ alors } \delta_i \in M.
\end{equation*}
Donc on est ramené à l'énoncé 6.7.1 déjà prouvé.


}{Quant à 6.7.2, on a $n-1$ impair, donc $\delta_i.\delta_i = 0$, et, pour toute puissance $k$-ième de $\sigma_i$, on a, en itérant 6.5.2,

\begin{equation}\tag{6.7.4}
\sigma_i^k(\chi) = \chi + k\,\sigma_i(\chi.\delta_i)\,\delta_i \quad .
\end{equation}
Donc un sous-espace $\mathbb{Q}_\ell$-vectoriel $M$ de $E$ est stable sous l'action de $\pi$ si et seulement s'il l'est sous un sous-groupe d'indice fini, car les deux conditions équivalent à la condition

\begin{equation*}
\chi \in M \implies \forall_i \text{ tel que le cup-produit } \chi.\delta_i \neq 0, \text{ alors } \delta_i \in M.
\end{equation*}
Donc on est ramené à l'énoncé 6.7.1 déjà prouvé.


}{F7II-XVIII-110}\underline{BIBLIOGRAPHIE}

[1] P. DELIGNE. Théorème de Lefschetz et critères de dégénérescence de suites spectrales, Pub. Math. IHES 35, 1968, 259-278.

[2] S. KLEIMAN. Algebraic cycles and the Weil conjectures, in Dix Exposés sur la cohomologie des schémas, North Holland, 1968.

\SGAvar{[3] S. LEFSCHETZ. L'analyse situs et la Géométrie Algébrique, Gauthiers-Villars, Paris 1924.

[4] D. LIEBERMAN. Higher Picard Varieties, Amer. J. Math, \underline{15} (4) 1968 p.1165-1199.

}{[3] S. LEFSCHETZ. L'analyse situs et la Géométrie Algébrique, Gauthier-Villars, Paris 1924.

[4] D. LIEBERMAN. Higher Picard Varieties, Amer. J. Math, \underline{90} (4) 1968 p.1165-1199.

}{F7II-XVIII-111}[5] D. MUMFORD. non publié.

[6] A. WEIL. Variétés Kählériennes, Hermann, Paris 1958.
\selectlanguage{french}

\ifnum\SGAform=1
\clearpage\section*{Apparat critique}\small
Les lectures ci-dessous sont données dans le TeX exact des témoins, afin de préserver les distinctions de notation. Les indices décimaux renvoient aux différences individuelles du passage marqué.

\par\medskip\noindent\hypertarget{app-F7II-XVIII-001.1}{}\textbf{F7II-XVIII-001.1}\quad p. 254; printed misprint.\par
\noindent Imprimé : {\ttfamily\footnotesize vérifié\allowbreak{}}\par
\noindent Corrigé : {\ttfamily\footnotesize vérifiée\allowbreak{}}\par
\noindent Source : 7II:412. \nolinkurl{fixes/MISPRINTS_7II.json#412}\par

\par\medskip\noindent\hypertarget{app-F7II-XVIII-001.2}{}\textbf{F7II-XVIII-001.2}\quad p. 254; printed misprint.\par
\noindent Imprimé : {\ttfamily\footnotesize succintes\allowbreak{}}\par
\noindent Corrigé : {\ttfamily\footnotesize succinctes\allowbreak{}}\par
\noindent Source : 7II:413. \nolinkurl{fixes/MISPRINTS_7II.json#413}\par

\par\medskip\noindent\hypertarget{app-F7II-XVIII-002.1}{}\textbf{F7II-XVIII-002.1}\quad p. 256; misprint\_pass.\par
\noindent Imprimé : \emph{(vide)}\par
\noindent Corrigé : {\ttfamily\footnotesize )\allowbreak{}}\par
\noindent Source : 7II:414. \nolinkurl{fixes/FIXLOG_7II_CURRENT.jsonl#741}; \nolinkurl{fixes/MISPRINTS_7II.json#414}\par

\par\medskip\noindent\hypertarget{app-F7II-XVIII-003.1}{}\textbf{F7II-XVIII-003.1}\quad p. 257; printed misprint.\par
\noindent Imprimé : {\ttfamily\footnotesize n\allowbreak{}}\par
\noindent Corrigé : {\ttfamily\footnotesize r\allowbreak{}}\par
\noindent Source : 7II:415. \nolinkurl{fixes/MISPRINTS_7II.json#415}\par

\par\medskip\noindent\hypertarget{app-F7II-XVIII-003.2}{}\textbf{F7II-XVIII-003.2}\quad p. 257; printed misprint.\par
\noindent Imprimé : {\ttfamily\footnotesize 2-j\allowbreak{}}\par
\noindent Corrigé : {\ttfamily\footnotesize 2j\allowbreak{}}\par
\noindent Source : 7II:416. \nolinkurl{fixes/MISPRINTS_7II.json#416}\par

\par\medskip\noindent\hypertarget{app-F7II-XVIII-003.3}{}\textbf{F7II-XVIII-003.3}\quad p. 257; printed misprint.\par
\noindent Imprimé : \emph{(vide)}\par
\noindent Corrigé : {\ttfamily\footnotesize )\allowbreak{}}\par
\noindent Source : 7II:417. \nolinkurl{fixes/MISPRINTS_7II.json#417}\par

\par\medskip\noindent\hypertarget{app-F7II-XVIII-003.4}{}\textbf{F7II-XVIII-003.4}\quad p. 257; misprint\_pass.\par
\noindent Imprimé : {\ttfamily\footnotesize défini\allowbreak{}}\par
\noindent Corrigé : {\ttfamily\footnotesize définie\allowbreak{}}\par
\noindent Source : 7II:418. \nolinkurl{fixes/FIXLOG_7II_CURRENT.jsonl#742}; \nolinkurl{fixes/MISPRINTS_7II.json#418}\par

\par\medskip\noindent\hypertarget{app-F7II-XVIII-004.1}{}\textbf{F7II-XVIII-004.1}\quad p. 258; misprint\_pass.\par
\noindent Imprimé : {\ttfamily\footnotesize \textbackslash{}leq\allowbreak{}}\par
\noindent Corrigé : {\ttfamily\footnotesize \textbackslash{}geq\allowbreak{}}\par
\noindent Source : 7II:419. \nolinkurl{fixes/FIXLOG_7II_CURRENT.jsonl#743}; \nolinkurl{fixes/MISPRINTS_7II.json#419}\par

\par\medskip\noindent\hypertarget{app-F7II-XVIII-005.1}{}\textbf{F7II-XVIII-005.1}\quad p. 259; inherited corrected-layer difference.\par
\noindent Imprimé : {\ttfamily\footnotesize o\allowbreak{}}\par
\noindent Corrigé : {\ttfamily\footnotesize 0\allowbreak{}}\par
\noindent Source : ADOPT-F7II-XVIII-005.1. \nolinkurl{fr/sga7ii/diplomatic/expose_XVIII_body.tex}; \nolinkurl{fr/sga7ii/current/expose_XVIII_body.tex}\par

\par\medskip\noindent\hypertarget{app-F7II-XVIII-006.1}{}\textbf{F7II-XVIII-006.1}\quad p. 259; designated-witness attribution.\par
\noindent Imprimé : {\ttfamily\footnotesize o\allowbreak{}}\par
\noindent Corrigé : {\ttfamily\footnotesize 0\allowbreak{}}\par
\noindent Source : WITNESS-F7II-XVIII-006.1. \nolinkurl{fr/sga7ii/diplomatic/expose_XVIII_body.tex}; \nolinkurl{fr/sga7ii/current/expose_XVIII_body.tex}\par

\par\medskip\noindent\hypertarget{app-F7II-XVIII-006.2}{}\textbf{F7II-XVIII-006.2}\quad p. 259; printed misprint.\par
\noindent Imprimé : \emph{(vide)}\par
\noindent Corrigé : {\ttfamily\footnotesize )\allowbreak{}}\par
\noindent Source : 7II:420. \nolinkurl{fixes/MISPRINTS_7II.json#420}\par

\par\medskip\noindent\hypertarget{app-F7II-XVIII-007.1}{}\textbf{F7II-XVIII-007.1}\quad p. 259; inherited corrected-layer difference.\par
\noindent Imprimé : {\ttfamily\footnotesize o\allowbreak{}}\par
\noindent Corrigé : {\ttfamily\footnotesize 0\allowbreak{}}\par
\noindent Source : ADOPT-F7II-XVIII-007.1. \nolinkurl{fr/sga7ii/diplomatic/expose_XVIII_body.tex}; \nolinkurl{fr/sga7ii/current/expose_XVIII_body.tex}\par

\par\medskip\noindent\hypertarget{app-F7II-XVIII-007.2}{}\textbf{F7II-XVIII-007.2}\quad p. 260; inherited corrected-layer difference.\par
\noindent Imprimé : {\ttfamily\footnotesize o\allowbreak{}}\par
\noindent Corrigé : {\ttfamily\footnotesize 0\allowbreak{}}\par
\noindent Source : ADOPT-F7II-XVIII-007.2. \nolinkurl{fr/sga7ii/diplomatic/expose_XVIII_body.tex}; \nolinkurl{fr/sga7ii/current/expose_XVIII_body.tex}\par

\par\medskip\noindent\hypertarget{app-F7II-XVIII-008.1}{}\textbf{F7II-XVIII-008.1}\quad p. 260; printed misprint.\par
\noindent Imprimé : {\ttfamily\footnotesize i\allowbreak{}}\par
\noindent Corrigé : {\ttfamily\footnotesize 1\allowbreak{}}\par
\noindent Source : 7II:422. \nolinkurl{fixes/MISPRINTS_7II.json#422}\par

\par\medskip\noindent\hypertarget{app-F7II-XVIII-009.1}{}\textbf{F7II-XVIII-009.1}\quad p. 260; printed misprint.\par
\noindent Imprimé : \emph{(vide)}\par
\noindent Corrigé : {\ttfamily\footnotesize )\allowbreak{}}\par
\noindent Source : 7II:423. \nolinkurl{fixes/MISPRINTS_7II.json#423}\par

\par\medskip\noindent\hypertarget{app-F7II-XVIII-009.2}{}\textbf{F7II-XVIII-009.2}\quad p. 260; designated-witness attribution.\par
\noindent Imprimé : {\ttfamily\footnotesize \_o\allowbreak{} \textbackslash{};\textbackslash{};\allowbreak{} }\par
\noindent Corrigé : {\ttfamily\footnotesize ,\allowbreak{}}\par
\noindent Source : WITNESS-F7II-XVIII-009.2. \nolinkurl{fr/sga7ii/diplomatic/expose_XVIII_body.tex}; \nolinkurl{fr/sga7ii/current/expose_XVIII_body.tex}\par

\par\medskip\noindent\hypertarget{app-F7II-XVIII-010.1}{}\textbf{F7II-XVIII-010.1}\quad p. 263; printed misprint.\par
\noindent Imprimé : {\ttfamily\footnotesize pat\allowbreak{}}\par
\noindent Corrigé : {\ttfamily\footnotesize par\allowbreak{}}\par
\noindent Source : 7II:425. \nolinkurl{fixes/MISPRINTS_7II.json#425}\par

\par\medskip\noindent\hypertarget{app-F7II-XVIII-010.2}{}\textbf{F7II-XVIII-010.2}\quad p. 263; designated-witness attribution.\par
\noindent Imprimé : {\ttfamily\footnotesize  R\allowbreak{}}\par
\noindent Corrigé : {\ttfamily\footnotesize  E\allowbreak{}}\par
\noindent Source : WITNESS-F7II-XVIII-010.2. \nolinkurl{fr/sga7ii/diplomatic/expose_XVIII_body.tex}; \nolinkurl{fr/sga7ii/current/expose_XVIII_body.tex}\par

\par\medskip\noindent\hypertarget{app-F7II-XVIII-011.1}{}\textbf{F7II-XVIII-011.1}\quad p. 263; printed misprint.\par
\noindent Imprimé : \emph{(vide)}\par
\noindent Corrigé : {\ttfamily\footnotesize \{\allowbreak{}}\par
\noindent Source : 7II:427. \nolinkurl{fixes/MISPRINTS_7II.json#427}\par

\par\medskip\noindent\hypertarget{app-F7II-XVIII-011.2}{}\textbf{F7II-XVIII-011.2}\quad p. 263; printed misprint.\par
\noindent Imprimé : \emph{(vide)}\par
\noindent Corrigé : {\ttfamily\footnotesize +1\}\allowbreak{}}\par
\noindent Source : 7II:427. \nolinkurl{fixes/MISPRINTS_7II.json#427}\par

\par\medskip\noindent\hypertarget{app-F7II-XVIII-012.1}{}\textbf{F7II-XVIII-012.1}\quad p. 265; printed misprint.\par
\noindent Imprimé : {\ttfamily\footnotesize H\allowbreak{}}\par
\noindent Corrigé : {\ttfamily\footnotesize F\allowbreak{}}\par
\noindent Source : 7II:428. \nolinkurl{fixes/MISPRINTS_7II.json#428}\par

\par\medskip\noindent\hypertarget{app-F7II-XVIII-012.2}{}\textbf{F7II-XVIII-012.2}\quad p. 265; printed misprint.\par
\noindent Imprimé : {\ttfamily\footnotesize F\allowbreak{}}\par
\noindent Corrigé : {\ttfamily\footnotesize H\allowbreak{}}\par
\noindent Source : 7II:428. \nolinkurl{fixes/MISPRINTS_7II.json#428}\par

\par\medskip\noindent\hypertarget{app-F7II-XVIII-012.3}{}\textbf{F7II-XVIII-012.3}\quad p. 265; printed misprint.\par
\noindent Imprimé : {\ttfamily\footnotesize direct\allowbreak{}}\par
\noindent Corrigé : {\ttfamily\footnotesize directe\allowbreak{}}\par
\noindent Source : 7II:429. \nolinkurl{fixes/MISPRINTS_7II.json#429}\par

\par\medskip\noindent\hypertarget{app-F7II-XVIII-013.1}{}\textbf{F7II-XVIII-013.1}\quad p. 265; inherited corrected-layer difference.\par
\noindent Imprimé : {\ttfamily\footnotesize \textbackslash{}\allowbreak{} }\par
\noindent Corrigé : \emph{(vide)}\par
\noindent Source : ADOPT-F7II-XVIII-013.1. \nolinkurl{fr/sga7ii/diplomatic/expose_XVIII_body.tex}; \nolinkurl{fr/sga7ii/current/expose_XVIII_body.tex}\par

\par\medskip\noindent\hypertarget{app-F7II-XVIII-013.2}{}\textbf{F7II-XVIII-013.2}\quad p. 265; designated-witness attribution.\par
\noindent Imprimé : {\ttfamily\footnotesize "\allowbreak{}}\par
\noindent Corrigé : {\ttfamily\footnotesize  «\allowbreak{} }\par
\noindent Source : WITNESS-F7II-XVIII-013.2. \nolinkurl{fr/sga7ii/diplomatic/expose_XVIII_body.tex}; \nolinkurl{fr/sga7ii/current/expose_XVIII_body.tex}\par

\par\medskip\noindent\hypertarget{app-F7II-XVIII-013.3}{}\textbf{F7II-XVIII-013.3}\quad p. 265; designated-witness attribution.\par
\noindent Imprimé : {\ttfamily\footnotesize "\allowbreak{}}\par
\noindent Corrigé : {\ttfamily\footnotesize  »\allowbreak{}}\par
\noindent Source : WITNESS-F7II-XVIII-013.3. \nolinkurl{fr/sga7ii/diplomatic/expose_XVIII_body.tex}; \nolinkurl{fr/sga7ii/current/expose_XVIII_body.tex}\par

\par\medskip\noindent\hypertarget{app-F7II-XVIII-013.4}{}\textbf{F7II-XVIII-013.4}\quad p. 265; designated-witness attribution.\par
\noindent Imprimé : \emph{(1 espace(s), 0 retour(s) de ligne dans le témoin TeX)}\par
\noindent Corrigé : \emph{(vide)}\par
\noindent Source : WITNESS-F7II-XVIII-013.4. \nolinkurl{fr/sga7ii/diplomatic/expose_XVIII_body.tex}; \nolinkurl{fr/sga7ii/current/expose_XVIII_body.tex}\par

\par\medskip\noindent\hypertarget{app-F7II-XVIII-013.5}{}\textbf{F7II-XVIII-013.5}\quad p. 265; printed misprint.\par
\noindent Imprimé : \emph{(vide)}\par
\noindent Corrigé : {\ttfamily\footnotesize \textbackslash{}widetilde\{\allowbreak{}}\par
\noindent Source : 7II:430. \nolinkurl{fixes/MISPRINTS_7II.json#430}\par

\par\medskip\noindent\hypertarget{app-F7II-XVIII-013.6}{}\textbf{F7II-XVIII-013.6}\quad p. 265; printed misprint.\par
\noindent Imprimé : \emph{(vide)}\par
\noindent Corrigé : {\ttfamily\footnotesize \}\allowbreak{}}\par
\noindent Source : 7II:430. \nolinkurl{fixes/MISPRINTS_7II.json#430}\par

\par\medskip\noindent\hypertarget{app-F7II-XVIII-014.1}{}\textbf{F7II-XVIII-014.1}\quad p. 266; printed misprint.\par
\noindent Imprimé : {\ttfamily\footnotesize C\allowbreak{}}\par
\noindent Corrigé : {\ttfamily\footnotesize c\allowbreak{}}\par
\noindent Source : 7II:766. \nolinkurl{fixes/MISPRINTS_7II.json#766}\par

\par\medskip\noindent\hypertarget{app-F7II-XVIII-014.2}{}\textbf{F7II-XVIII-014.2}\quad p. 266; misprint\_pass.\par
\noindent Imprimé : \emph{(vide)}\par
\noindent Corrigé : {\ttfamily\footnotesize .\allowbreak{}}\par
\noindent Source : 7II:768; 7II:767. \nolinkurl{fixes/FIXLOG_7II_CURRENT.jsonl#821}; \nolinkurl{fixes/MISPRINTS_7II.json#767}\par

\par\medskip\noindent\hypertarget{app-F7II-XVIII-015.1}{}\textbf{F7II-XVIII-015.1}\quad p. 266; printed misprint.\par
\noindent Imprimé : {\ttfamily\footnotesize \textbackslash{}chi\allowbreak{}}\par
\noindent Corrigé : {\ttfamily\footnotesize X\allowbreak{}}\par
\noindent Source : 7II:433; 7II:768. \nolinkurl{fixes/MISPRINTS_7II.json#433}; \nolinkurl{fixes/MISPRINTS_7II.json#768}\par

\par\medskip\noindent\hypertarget{app-F7II-XVIII-015.2}{}\textbf{F7II-XVIII-015.2}\quad p. 266; printed misprint.\par
\noindent Imprimé : {\ttfamily\footnotesize j\allowbreak{}}\par
\noindent Corrigé : {\ttfamily\footnotesize 2j\allowbreak{}}\par
\noindent Source : 7II:434; 7II:769. \nolinkurl{fixes/MISPRINTS_7II.json#434}; \nolinkurl{fixes/MISPRINTS_7II.json#769}\par

\par\medskip\noindent\hypertarget{app-F7II-XVIII-016.1}{}\textbf{F7II-XVIII-016.1}\quad p. 267; printed misprint.\par
\noindent Imprimé : {\ttfamily\footnotesize u\allowbreak{}}\par
\noindent Corrigé : {\ttfamily\footnotesize n\allowbreak{}}\par
\noindent Source : 7II:770. \nolinkurl{fixes/MISPRINTS_7II.json#770}\par

\par\medskip\noindent\hypertarget{app-F7II-XVIII-016.2}{}\textbf{F7II-XVIII-016.2}\quad p. 267; misprint\_pass.\par
\noindent Imprimé : {\ttfamily\footnotesize coordonnés\allowbreak{}}\par
\noindent Corrigé : {\ttfamily\footnotesize coordonnées\allowbreak{}}\par
\noindent Source : 7II:436; 7II:771. \nolinkurl{fixes/FIXLOG_7II_CURRENT.jsonl#744}; \nolinkurl{fixes/MISPRINTS_7II.json#771}\par

\par\medskip\noindent\hypertarget{app-F7II-XVIII-017.1}{}\textbf{F7II-XVIII-017.1}\quad p. 267; printed misprint.\par
\noindent Imprimé : \emph{(vide)}\par
\noindent Corrigé : {\ttfamily\footnotesize )\allowbreak{}}\par
\noindent Source : 7II:437; 7II:772. \nolinkurl{fixes/MISPRINTS_7II.json#437}; \nolinkurl{fixes/MISPRINTS_7II.json#772}\par

\par\medskip\noindent\hypertarget{app-F7II-XVIII-018.1}{}\textbf{F7II-XVIII-018.1}\quad p. 268; printed misprint.\par
\noindent Imprimé : {\ttfamily\footnotesize Pincaré\allowbreak{}}\par
\noindent Corrigé : {\ttfamily\footnotesize Poincaré\allowbreak{}}\par
\noindent Source : 7II:438; 7II:773. \nolinkurl{fixes/MISPRINTS_7II.json#438}; \nolinkurl{fixes/MISPRINTS_7II.json#773}\par

\par\medskip\noindent\hypertarget{app-F7II-XVIII-018.2}{}\textbf{F7II-XVIII-018.2}\quad p. 268; printed misprint.\par
\noindent Imprimé : {\ttfamily\footnotesize S\allowbreak{}}\par
\noindent Corrigé : {\ttfamily\footnotesize \textbackslash{}rho\allowbreak{}}\par
\noindent Source : 7II:774. \nolinkurl{fixes/MISPRINTS_7II.json#774}\par

\par\medskip\noindent\hypertarget{app-F7II-XVIII-018.3}{}\textbf{F7II-XVIII-018.3}\quad p. 268; designated-witness attribution.\par
\noindent Imprimé : \emph{(0 espace(s), 1 retour(s) de ligne dans le témoin TeX)}\par
\noindent Corrigé : \emph{(0 espace(s), 2 retour(s) de ligne dans le témoin TeX)}\par
\noindent Source : WITNESS-F7II-XVIII-018.3. \nolinkurl{fr/sga7ii/diplomatic/expose_XVIII_body.tex}; \nolinkurl{fr/sga7ii/current/expose_XVIII_body.tex}\par

\par\medskip\noindent\hypertarget{app-F7II-XVIII-018.4}{}\textbf{F7II-XVIII-018.4}\quad p. 268; printed misprint.\par
\noindent Imprimé : {\ttfamily\footnotesize \textbackslash{}eta\allowbreak{}}\par
\noindent Corrigé : {\ttfamily\footnotesize n\allowbreak{}}\par
\noindent Source : 7II:775. \nolinkurl{fixes/MISPRINTS_7II.json#775}\par

\par\medskip\noindent\hypertarget{app-F7II-XVIII-018.5}{}\textbf{F7II-XVIII-018.5}\quad p. 268; printed misprint.\par
\noindent Imprimé : {\ttfamily\footnotesize \textbackslash{}eta\allowbreak{}}\par
\noindent Corrigé : {\ttfamily\footnotesize n\allowbreak{}}\par
\noindent Source : 7II:776. \nolinkurl{fixes/MISPRINTS_7II.json#776}\par

\par\medskip\noindent\hypertarget{app-F7II-XVIII-019.1}{}\textbf{F7II-XVIII-019.1}\quad p. 269; inherited corrected-layer difference.\par
\noindent Imprimé : \emph{(0 espace(s), 1 retour(s) de ligne dans le témoin TeX)}\par
\noindent Corrigé : \emph{(0 espace(s), 3 retour(s) de ligne dans le témoin TeX)}\par
\noindent Source : ADOPT-F7II-XVIII-019.1. \nolinkurl{fr/sga7ii/diplomatic/expose_XVIII_body.tex}; \nolinkurl{fr/sga7ii/current/expose_XVIII_body.tex}\par

\par\medskip\noindent\hypertarget{app-F7II-XVIII-019.2}{}\textbf{F7II-XVIII-019.2}\quad p. 269; inherited corrected-layer difference.\par
\noindent Imprimé : \emph{(1 espace(s), 0 retour(s) de ligne dans le témoin TeX)}\par
\noindent Corrigé : {\ttfamily\footnotesize \textasciicircum{}3\allowbreak{}}\par
\noindent Source : ADOPT-F7II-XVIII-019.2. \nolinkurl{fr/sga7ii/diplomatic/expose_XVIII_body.tex}; \nolinkurl{fr/sga7ii/current/expose_XVIII_body.tex}\par

\par\medskip\noindent\hypertarget{app-F7II-XVIII-019.3}{}\textbf{F7II-XVIII-019.3}\quad p. 269; inherited corrected-layer difference.\par
\noindent Imprimé : \emph{(vide)}\par
\noindent Corrigé : {\ttfamily\footnotesize -1\allowbreak{}}\par
\noindent Source : ADOPT-F7II-XVIII-019.3. \nolinkurl{fr/sga7ii/diplomatic/expose_XVIII_body.tex}; \nolinkurl{fr/sga7ii/current/expose_XVIII_body.tex}\par

\par\medskip\noindent\hypertarget{app-F7II-XVIII-019.4}{}\textbf{F7II-XVIII-019.4}\quad p. 269; inherited corrected-layer difference.\par
\noindent Imprimé : {\ttfamily\footnotesize  racine\allowbreak{} cubique\allowbreak{} de\allowbreak{} l'unité\allowbreak{}}\par
\noindent Corrigé : \emph{(vide)}\par
\noindent Source : ADOPT-F7II-XVIII-019.4. \nolinkurl{fr/sga7ii/diplomatic/expose_XVIII_body.tex}; \nolinkurl{fr/sga7ii/current/expose_XVIII_body.tex}\par

\par\medskip\noindent\hypertarget{app-F7II-XVIII-020.1}{}\textbf{F7II-XVIII-020.1}\quad p. 270; printed misprint.\par
\noindent Imprimé : {\ttfamily\footnotesize \textbackslash{}eta\allowbreak{}}\par
\noindent Corrigé : {\ttfamily\footnotesize n\allowbreak{}}\par
\noindent Source : 7II:777. \nolinkurl{fixes/MISPRINTS_7II.json#777}\par

\par\medskip\noindent\hypertarget{app-F7II-XVIII-020.2}{}\textbf{F7II-XVIII-020.2}\quad p. 270; printed misprint.\par
\noindent Imprimé : {\ttfamily\footnotesize i\allowbreak{}}\par
\noindent Corrigé : {\ttfamily\footnotesize 1\allowbreak{}}\par
\noindent Source : 7II:778. \nolinkurl{fixes/MISPRINTS_7II.json#778}\par

\par\medskip\noindent\hypertarget{app-F7II-XVIII-020.3}{}\textbf{F7II-XVIII-020.3}\quad p. 270; inherited corrected-layer difference.\par
\noindent Imprimé : {\ttfamily\footnotesize \textbackslash{}mathrm\allowbreak{}}\par
\noindent Corrigé : {\ttfamily\footnotesize \textbackslash{}operatorname\allowbreak{}}\par
\noindent Source : ADOPT-F7II-XVIII-020.3. \nolinkurl{fr/sga7ii/diplomatic/expose_XVIII_body.tex}; \nolinkurl{fr/sga7ii/current/expose_XVIII_body.tex}\par

\par\medskip\noindent\hypertarget{app-F7II-XVIII-020.4}{}\textbf{F7II-XVIII-020.4}\quad p. 270; inherited corrected-layer difference.\par
\noindent Imprimé : {\ttfamily\footnotesize rg\allowbreak{}}\par
\noindent Corrigé : {\ttfamily\footnotesize lg\allowbreak{}}\par
\noindent Source : ADOPT-F7II-XVIII-020.4. \nolinkurl{fr/sga7ii/diplomatic/expose_XVIII_body.tex}; \nolinkurl{fr/sga7ii/current/expose_XVIII_body.tex}\par

\par\medskip\noindent\hypertarget{app-F7II-XVIII-020.5}{}\textbf{F7II-XVIII-020.5}\quad p. 270; inherited corrected-layer difference.\par
\noindent Imprimé : \emph{(vide)}\par
\noindent Corrigé : {\ttfamily\footnotesize \_\{i=1\}\allowbreak{}}\par
\noindent Source : ADOPT-F7II-XVIII-020.5. \nolinkurl{fr/sga7ii/diplomatic/expose_XVIII_body.tex}; \nolinkurl{fr/sga7ii/current/expose_XVIII_body.tex}\par

\par\medskip\noindent\hypertarget{app-F7II-XVIII-020.6}{}\textbf{F7II-XVIII-020.6}\quad p. 270; misprint\_pass.\par
\noindent Imprimé : {\ttfamily\footnotesize ,\allowbreak{}}\par
\noindent Corrigé : \emph{(vide)}\par
\noindent Source : 7II:780. \nolinkurl{fixes/FIXLOG_7II_CURRENT.jsonl#822}\par

\par\medskip\noindent\hypertarget{app-F7II-XVIII-021.1}{}\textbf{F7II-XVIII-021.1}\quad p. 271; misprint\_in\_print.\par
\noindent Imprimé : {\ttfamily\footnotesize vérifiéeen\allowbreak{}}\par
\noindent Corrigé : {\ttfamily\footnotesize vérifiée\allowbreak{} en\allowbreak{}}\par
\noindent Source : FIX-7II-D-0629. \nolinkurl{fixes/FIXLOG_7II.jsonl#629}\par

\par\medskip\noindent\hypertarget{app-F7II-XVIII-023.1}{}\textbf{F7II-XVIII-023.1}\quad p. 272; designated-witness attribution.\par
\noindent Imprimé : \emph{(vide)}\par
\noindent Corrigé : {\ttfamily\footnotesize -\allowbreak{}}\par
\noindent Source : WITNESS-F7II-XVIII-023.1. \nolinkurl{fr/sga7ii/diplomatic/expose_XVIII_body.tex}; \nolinkurl{fr/sga7ii/current/expose_XVIII_body.tex}\par

\par\medskip\noindent\hypertarget{app-F7II-XVIII-024.1}{}\textbf{F7II-XVIII-024.1}\quad p. 272; printed misprint.\par
\noindent Imprimé : \emph{(vide)}\par
\noindent Corrigé : {\ttfamily\footnotesize )\allowbreak{}}\par
\noindent Source : 7II:794. \nolinkurl{fixes/MISPRINTS_7II.json#794}\par

\par\medskip\noindent\hypertarget{app-F7II-XVIII-024.2}{}\textbf{F7II-XVIII-024.2}\quad p. 272; printed misprint.\par
\noindent Imprimé : \emph{(vide)}\par
\noindent Corrigé : {\ttfamily\footnotesize )\allowbreak{}}\par
\noindent Source : 7II:795. \nolinkurl{fixes/MISPRINTS_7II.json#795}\par

\par\medskip\noindent\hypertarget{app-F7II-XVIII-024.3}{}\textbf{F7II-XVIII-024.3}\quad p. 272; printed misprint.\par
\noindent Imprimé : \emph{(vide)}\par
\noindent Corrigé : {\ttfamily\footnotesize )\allowbreak{}}\par
\noindent Source : 7II:796. \nolinkurl{fixes/MISPRINTS_7II.json#796}\par

\par\medskip\noindent\hypertarget{app-F7II-XVIII-025.1}{}\textbf{F7II-XVIII-025.1}\quad p. 273; misprint\_pass.\par
\noindent Imprimé : {\ttfamily\footnotesize totalede\allowbreak{}}\par
\noindent Corrigé : {\ttfamily\footnotesize totale\allowbreak{} de\allowbreak{}}\par
\noindent Source : 7II:447; 7II:797. \nolinkurl{fixes/FIXLOG_7II_CURRENT.jsonl#745}; \nolinkurl{fixes/MISPRINTS_7II.json#447}; \nolinkurl{fixes/MISPRINTS_7II.json#797}\par

\par\medskip\noindent\hypertarget{app-F7II-XVIII-025.2}{}\textbf{F7II-XVIII-025.2}\quad p. 273; misprint\_pass.\par
\noindent Imprimé : \emph{(vide)}\par
\noindent Corrigé : {\ttfamily\footnotesize (\allowbreak{}}\par
\noindent Source : 7II:448; 7II:798. \nolinkurl{fixes/FIXLOG_7II_CURRENT.jsonl#746}; \nolinkurl{fixes/MISPRINTS_7II.json#798}\par

\par\medskip\noindent\hypertarget{app-F7II-XVIII-026.1}{}\textbf{F7II-XVIII-026.1}\quad p. 273; printed misprint.\par
\noindent Imprimé : {\ttfamily\footnotesize ,\allowbreak{} leftrightarrow\allowbreak{}}\par
\noindent Corrigé : \emph{(vide)}\par
\noindent Source : 7II:449; 7II:799. \nolinkurl{fixes/MISPRINTS_7II.json#449}; \nolinkurl{fixes/MISPRINTS_7II.json#799}\par

\par\medskip\noindent\hypertarget{app-F7II-XVIII-027.1}{}\textbf{F7II-XVIII-027.1}\quad p. 274; inherited corrected-layer difference.\par
\noindent Imprimé : \emph{(vide)}\par
\noindent Corrigé : {\ttfamily\footnotesize (-\allowbreak{}}\par
\noindent Source : ADOPT-F7II-XVIII-027.1. \nolinkurl{fr/sga7ii/diplomatic/expose_XVIII_body.tex}; \nolinkurl{fr/sga7ii/current/expose_XVIII_body.tex}\par

\par\medskip\noindent\hypertarget{app-F7II-XVIII-027.2}{}\textbf{F7II-XVIII-027.2}\quad p. 274; inherited corrected-layer difference.\par
\noindent Imprimé : \emph{(vide)}\par
\noindent Corrigé : {\ttfamily\footnotesize )\allowbreak{}}\par
\noindent Source : ADOPT-F7II-XVIII-027.2. \nolinkurl{fr/sga7ii/diplomatic/expose_XVIII_body.tex}; \nolinkurl{fr/sga7ii/current/expose_XVIII_body.tex}\par

\par\medskip\noindent\hypertarget{app-F7II-XVIII-027.3}{}\textbf{F7II-XVIII-027.3}\quad p. 274; misprint\_pass.\par
\noindent Imprimé : {\ttfamily\footnotesize 2\allowbreak{}}\par
\noindent Corrigé : {\ttfamily\footnotesize 4\allowbreak{}}\par
\noindent Source : 7II:450. \nolinkurl{fixes/FIXLOG_7II_CURRENT.jsonl#747}\par

\par\medskip\noindent\hypertarget{app-F7II-XVIII-027.4}{}\textbf{F7II-XVIII-027.4}\quad p. 274; misprint\_pass.\par
\noindent Imprimé : {\ttfamily\footnotesize \_\allowbreak{}}\par
\noindent Corrigé : {\ttfamily\footnotesize \textasciicircum{}\allowbreak{}}\par
\noindent Source : 7II:451; 7II:801. \nolinkurl{fixes/FIXLOG_7II_CURRENT.jsonl#748}; \nolinkurl{fixes/MISPRINTS_7II.json#451}; \nolinkurl{fixes/MISPRINTS_7II.json#801}\par

\par\medskip\noindent\hypertarget{app-F7II-XVIII-028.1}{}\textbf{F7II-XVIII-028.1}\quad p. 274; designated-witness attribution.\par
\noindent Imprimé : \emph{(vide)}\par
\noindent Corrigé : {\ttfamily\footnotesize \_\{\textbackslash{}sim\}\allowbreak{}}\par
\noindent Source : WITNESS-F7II-XVIII-028.1. \nolinkurl{fr/sga7ii/diplomatic/expose_XVIII_body.tex}; \nolinkurl{fr/sga7ii/current/expose_XVIII_body.tex}\par

\par\medskip\noindent\hypertarget{app-F7II-XVIII-028.2}{}\textbf{F7II-XVIII-028.2}\quad p. 274; designated-witness attribution.\par
\noindent Imprimé : \emph{(vide)}\par
\noindent Corrigé : {\ttfamily\footnotesize -\allowbreak{}}\par
\noindent Source : WITNESS-F7II-XVIII-028.2. \nolinkurl{fr/sga7ii/diplomatic/expose_XVIII_body.tex}; \nolinkurl{fr/sga7ii/current/expose_XVIII_body.tex}\par

\par\medskip\noindent\hypertarget{app-F7II-XVIII-029.1}{}\textbf{F7II-XVIII-029.1}\quad p. 274; designated-witness attribution.\par
\noindent Imprimé : \emph{(0 espace(s), 1 retour(s) de ligne dans le témoin TeX)}\par
\noindent Corrigé : {\ttfamily\footnotesize  -\allowbreak{}}\par
\noindent Source : WITNESS-F7II-XVIII-029.1. \nolinkurl{fr/sga7ii/diplomatic/expose_XVIII_body.tex}; \nolinkurl{fr/sga7ii/current/expose_XVIII_body.tex}\par

\par\medskip\noindent\hypertarget{app-F7II-XVIII-029.2}{}\textbf{F7II-XVIII-029.2}\quad p. 274; printed misprint.\par
\noindent Imprimé : {\ttfamily\footnotesize X\allowbreak{}}\par
\noindent Corrigé : {\ttfamily\footnotesize x\allowbreak{}}\par
\noindent Source : 7II:802. \nolinkurl{fixes/MISPRINTS_7II.json#802}\par

\par\medskip\noindent\hypertarget{app-F7II-XVIII-029.3}{}\textbf{F7II-XVIII-029.3}\quad p. 274; printed misprint.\par
\noindent Imprimé : {\ttfamily\footnotesize X\allowbreak{}}\par
\noindent Corrigé : {\ttfamily\footnotesize x\allowbreak{}}\par
\noindent Source : 7II:802. \nolinkurl{fixes/MISPRINTS_7II.json#802}\par

\par\medskip\noindent\hypertarget{app-F7II-XVIII-029.4}{}\textbf{F7II-XVIII-029.4}\quad p. 274; printed misprint.\par
\noindent Imprimé : {\ttfamily\footnotesize X\allowbreak{}}\par
\noindent Corrigé : {\ttfamily\footnotesize x\allowbreak{}}\par
\noindent Source : 7II:802. \nolinkurl{fixes/MISPRINTS_7II.json#802}\par

\par\medskip\noindent\hypertarget{app-F7II-XVIII-030.1}{}\textbf{F7II-XVIII-030.1}\quad p. 275; printed misprint.\par
\noindent Imprimé : \emph{(vide)}\par
\noindent Corrigé : {\ttfamily\footnotesize de\allowbreak{} }\par
\noindent Source : 7II:803. \nolinkurl{fixes/MISPRINTS_7II.json#803}\par

\par\medskip\noindent\hypertarget{app-F7II-XVIII-031.1}{}\textbf{F7II-XVIII-031.1}\quad p. 275; misprint\_pass.\par
\noindent Imprimé : {\ttfamily\footnotesize \textasciicircum{}2\allowbreak{}}\par
\noindent Corrigé : \emph{(vide)}\par
\noindent Source : 7II:453; 7II:804. \nolinkurl{fixes/FIXLOG_7II_CURRENT.jsonl#750}; \nolinkurl{fixes/MISPRINTS_7II.json#453}; \nolinkurl{fixes/MISPRINTS_7II.json#804}\par

\par\medskip\noindent\hypertarget{app-F7II-XVIII-031.2}{}\textbf{F7II-XVIII-031.2}\quad p. 275; misprint\_pass.\par
\noindent Imprimé : {\ttfamily\footnotesize K\allowbreak{}}\par
\noindent Corrigé : {\ttfamily\footnotesize k\allowbreak{}}\par
\noindent Source : 7II:806. \nolinkurl{fixes/FIXLOG_7II_CURRENT.jsonl#823}\par

\par\medskip\noindent\hypertarget{app-F7II-XVIII-031.3}{}\textbf{F7II-XVIII-031.3}\quad p. 275; inherited corrected-layer difference.\par
\noindent Imprimé : \emph{(vide)}\par
\noindent Corrigé : {\ttfamily\footnotesize \textasciicircum{}1\allowbreak{}}\par
\noindent Source : ADOPT-F7II-XVIII-031.3. \nolinkurl{fr/sga7ii/diplomatic/expose_XVIII_body.tex}; \nolinkurl{fr/sga7ii/current/expose_XVIII_body.tex}\par

\par\medskip\noindent\hypertarget{app-F7II-XVIII-032.1}{}\textbf{F7II-XVIII-032.1}\quad p. 276; designated-witness attribution.\par
\noindent Imprimé : \emph{(vide)}\par
\noindent Corrigé : {\ttfamily\footnotesize \textbackslash{}bigl\allowbreak{}}\par
\noindent Source : WITNESS-F7II-XVIII-032.1. \nolinkurl{fr/sga7ii/diplomatic/expose_XVIII_body.tex}; \nolinkurl{fr/sga7ii/current/expose_XVIII_body.tex}\par

\par\medskip\noindent\hypertarget{app-F7II-XVIII-032.2}{}\textbf{F7II-XVIII-032.2}\quad p. 276; designated-witness attribution.\par
\noindent Imprimé : \emph{(vide)}\par
\noindent Corrigé : {\ttfamily\footnotesize \textbackslash{}bigl\allowbreak{}}\par
\noindent Source : WITNESS-F7II-XVIII-032.2. \nolinkurl{fr/sga7ii/diplomatic/expose_XVIII_body.tex}; \nolinkurl{fr/sga7ii/current/expose_XVIII_body.tex}\par

\par\medskip\noindent\hypertarget{app-F7II-XVIII-032.3}{}\textbf{F7II-XVIII-032.3}\quad p. 276; designated-witness attribution.\par
\noindent Imprimé : \emph{(vide)}\par
\noindent Corrigé : {\ttfamily\footnotesize )\textbackslash{}bigr\allowbreak{}}\par
\noindent Source : WITNESS-F7II-XVIII-032.3. \nolinkurl{fr/sga7ii/diplomatic/expose_XVIII_body.tex}; \nolinkurl{fr/sga7ii/current/expose_XVIII_body.tex}\par

\par\medskip\noindent\hypertarget{app-F7II-XVIII-032.4}{}\textbf{F7II-XVIII-032.4}\quad p. 276; designated-witness attribution.\par
\noindent Imprimé : \emph{(vide)}\par
\noindent Corrigé : {\ttfamily\footnotesize \textbackslash{}bigr\allowbreak{}}\par
\noindent Source : WITNESS-F7II-XVIII-032.4. \nolinkurl{fr/sga7ii/diplomatic/expose_XVIII_body.tex}; \nolinkurl{fr/sga7ii/current/expose_XVIII_body.tex}\par

\par\medskip\noindent\hypertarget{app-F7II-XVIII-033.1}{}\textbf{F7II-XVIII-033.1}\quad p. 276; designated-witness attribution.\par
\noindent Imprimé : \emph{(vide)}\par
\noindent Corrigé : {\ttfamily\footnotesize \textbackslash{}bigl\allowbreak{}}\par
\noindent Source : WITNESS-F7II-XVIII-033.1. \nolinkurl{fr/sga7ii/diplomatic/expose_XVIII_body.tex}; \nolinkurl{fr/sga7ii/current/expose_XVIII_body.tex}\par

\par\medskip\noindent\hypertarget{app-F7II-XVIII-033.2}{}\textbf{F7II-XVIII-033.2}\quad p. 276; designated-witness attribution.\par
\noindent Imprimé : \emph{(vide)}\par
\noindent Corrigé : {\ttfamily\footnotesize )\textbackslash{}bigr\allowbreak{}}\par
\noindent Source : WITNESS-F7II-XVIII-033.2. \nolinkurl{fr/sga7ii/diplomatic/expose_XVIII_body.tex}; \nolinkurl{fr/sga7ii/current/expose_XVIII_body.tex}\par

\par\medskip\noindent\hypertarget{app-F7II-XVIII-033.3}{}\textbf{F7II-XVIII-033.3}\quad p. 276; misprint\_in\_print.\par
\noindent Imprimé : {\ttfamily\footnotesize \textbackslash{}Lambda\allowbreak{}}\par
\noindent Corrigé : {\ttfamily\footnotesize \textbackslash{}wedge\allowbreak{}}\par
\noindent Source : FIX-7II-D-0631. \nolinkurl{fixes/FIXLOG_7II.jsonl#631}\par

\par\medskip\noindent\hypertarget{app-F7II-XVIII-033.4}{}\textbf{F7II-XVIII-033.4}\quad p. 276; designated-witness attribution.\par
\noindent Imprimé : \emph{(vide)}\par
\noindent Corrigé : {\ttfamily\footnotesize \textbackslash{}bigl\allowbreak{}}\par
\noindent Source : WITNESS-F7II-XVIII-033.4. \nolinkurl{fr/sga7ii/diplomatic/expose_XVIII_body.tex}; \nolinkurl{fr/sga7ii/current/expose_XVIII_body.tex}\par

\par\medskip\noindent\hypertarget{app-F7II-XVIII-033.5}{}\textbf{F7II-XVIII-033.5}\quad p. 276; designated-witness attribution.\par
\noindent Imprimé : \emph{(vide)}\par
\noindent Corrigé : {\ttfamily\footnotesize )\textbackslash{}bigr\allowbreak{}}\par
\noindent Source : WITNESS-F7II-XVIII-033.5. \nolinkurl{fr/sga7ii/diplomatic/expose_XVIII_body.tex}; \nolinkurl{fr/sga7ii/current/expose_XVIII_body.tex}\par

\par\medskip\noindent\hypertarget{app-F7II-XVIII-034.1}{}\textbf{F7II-XVIII-034.1}\quad p. 277; printed misprint.\par
\noindent Imprimé : {\ttfamily\footnotesize 3\allowbreak{}}\par
\noindent Corrigé : {\ttfamily\footnotesize 2\allowbreak{}}\par
\noindent Source : 7II:813. \nolinkurl{fixes/MISPRINTS_7II.json#813}\par

\par\medskip\noindent\hypertarget{app-F7II-XVIII-034.2}{}\textbf{F7II-XVIII-034.2}\quad p. 277; printed misprint.\par
\noindent Imprimé : {\ttfamily\footnotesize 2\allowbreak{}}\par
\noindent Corrigé : {\ttfamily\footnotesize 12\allowbreak{}}\par
\noindent Source : 7II:813. \nolinkurl{fixes/MISPRINTS_7II.json#813}\par

\par\medskip\noindent\hypertarget{app-F7II-XVIII-034.3}{}\textbf{F7II-XVIII-034.3}\quad p. 277; misprint\_pass.\par
\noindent Imprimé : {\ttfamily\footnotesize 3\allowbreak{}}\par
\noindent Corrigé : {\ttfamily\footnotesize 2\allowbreak{}}\par
\noindent Source : 7II:815. \nolinkurl{fixes/FIXLOG_7II_CURRENT.jsonl#824}\par

\par\medskip\noindent\hypertarget{app-F7II-XVIII-034.4}{}\textbf{F7II-XVIII-034.4}\quad p. 277; misprint\_pass.\par
\noindent Imprimé : {\ttfamily\footnotesize 2\allowbreak{}}\par
\noindent Corrigé : {\ttfamily\footnotesize 11\allowbreak{}}\par
\noindent Source : 7II:815. \nolinkurl{fixes/FIXLOG_7II_CURRENT.jsonl#824}\par

\par\medskip\noindent\hypertarget{app-F7II-XVIII-034.5}{}\textbf{F7II-XVIII-034.5}\quad p. 277; printed misprint.\par
\noindent Imprimé : \emph{(vide)}\par
\noindent Corrigé : {\ttfamily\footnotesize \}\allowbreak{}}\par
\noindent Source : 7II:815. \nolinkurl{fixes/MISPRINTS_7II.json#815}\par

\par\medskip\noindent\hypertarget{app-F7II-XVIII-034.6}{}\textbf{F7II-XVIII-034.6}\quad p. 277; printed misprint.\par
\noindent Imprimé : {\ttfamily\footnotesize \}\allowbreak{}}\par
\noindent Corrigé : \emph{(vide)}\par
\noindent Source : 7II:815. \nolinkurl{fixes/MISPRINTS_7II.json#815}\par

\par\medskip\noindent\hypertarget{app-F7II-XVIII-034.7}{}\textbf{F7II-XVIII-034.7}\quad p. 277; designated-witness attribution.\par
\noindent Imprimé : \emph{(vide)}\par
\noindent Corrigé : {\ttfamily\footnotesize \textbackslash{}bigl\allowbreak{}}\par
\noindent Source : WITNESS-F7II-XVIII-034.7. \nolinkurl{fr/sga7ii/diplomatic/expose_XVIII_body.tex}; \nolinkurl{fr/sga7ii/current/expose_XVIII_body.tex}\par

\par\medskip\noindent\hypertarget{app-F7II-XVIII-034.8}{}\textbf{F7II-XVIII-034.8}\quad p. 277; designated-witness attribution.\par
\noindent Imprimé : \emph{(vide)}\par
\noindent Corrigé : {\ttfamily\footnotesize \textbackslash{}bigr\allowbreak{}}\par
\noindent Source : WITNESS-F7II-XVIII-034.8. \nolinkurl{fr/sga7ii/diplomatic/expose_XVIII_body.tex}; \nolinkurl{fr/sga7ii/current/expose_XVIII_body.tex}\par

\par\medskip\noindent\hypertarget{app-F7II-XVIII-034.9}{}\textbf{F7II-XVIII-034.9}\quad p. 277; designated-witness attribution.\par
\noindent Imprimé : \emph{(vide)}\par
\noindent Corrigé : {\ttfamily\footnotesize \textbackslash{}bigl\allowbreak{}}\par
\noindent Source : WITNESS-F7II-XVIII-034.9. \nolinkurl{fr/sga7ii/diplomatic/expose_XVIII_body.tex}; \nolinkurl{fr/sga7ii/current/expose_XVIII_body.tex}\par

\par\medskip\noindent\hypertarget{app-F7II-XVIII-034.10}{}\textbf{F7II-XVIII-034.10}\quad p. 277; printed misprint.\par
\noindent Imprimé : \emph{(vide)}\par
\noindent Corrigé : {\ttfamily\footnotesize )\allowbreak{}}\par
\noindent Source : 7II:457; 7II:816. \nolinkurl{fixes/MISPRINTS_7II.json#457}; \nolinkurl{fixes/MISPRINTS_7II.json#816}\par

\par\medskip\noindent\hypertarget{app-F7II-XVIII-034.11}{}\textbf{F7II-XVIII-034.11}\quad p. 277; designated-witness attribution.\par
\noindent Imprimé : \emph{(vide)}\par
\noindent Corrigé : {\ttfamily\footnotesize \textbackslash{}bigr)\allowbreak{}}\par
\noindent Source : WITNESS-F7II-XVIII-034.11. \nolinkurl{fr/sga7ii/diplomatic/expose_XVIII_body.tex}; \nolinkurl{fr/sga7ii/current/expose_XVIII_body.tex}\par

\par\medskip\noindent\hypertarget{app-F7II-XVIII-034.12}{}\textbf{F7II-XVIII-034.12}\quad p. 277; printed misprint.\par
\noindent Imprimé : \emph{(vide)}\par
\noindent Corrigé : {\ttfamily\footnotesize '\allowbreak{}}\par
\noindent Source : 7II:458; 7II:817. \nolinkurl{fixes/MISPRINTS_7II.json#458}; \nolinkurl{fixes/MISPRINTS_7II.json#817}\par

\par\medskip\noindent\hypertarget{app-F7II-XVIII-034.13}{}\textbf{F7II-XVIII-034.13}\quad p. 277; misprint\_in\_print.\par
\noindent Imprimé : {\ttfamily\footnotesize \textbackslash{}Lambda\allowbreak{}}\par
\noindent Corrigé : {\ttfamily\footnotesize \textbackslash{}wedge\allowbreak{}}\par
\noindent Source : FIX-7II-D-0632. \nolinkurl{fixes/FIXLOG_7II.jsonl#632}\par

\par\medskip\noindent\hypertarget{app-F7II-XVIII-034.14}{}\textbf{F7II-XVIII-034.14}\quad p. 277; designated-witness attribution.\par
\noindent Imprimé : {\ttfamily\footnotesize \textbackslash{}cap\allowbreak{}}\par
\noindent Corrigé : {\ttfamily\footnotesize \textbackslash{}cdot\allowbreak{}}\par
\noindent Source : WITNESS-F7II-XVIII-034.14. \nolinkurl{fr/sga7ii/diplomatic/expose_XVIII_body.tex}; \nolinkurl{fr/sga7ii/current/expose_XVIII_body.tex}\par

\par\medskip\noindent\hypertarget{app-F7II-XVIII-034.15}{}\textbf{F7II-XVIII-034.15}\quad p. 277; misprint\_in\_print.\par
\noindent Imprimé : {\ttfamily\footnotesize \textbackslash{}Lambda\allowbreak{}}\par
\noindent Corrigé : {\ttfamily\footnotesize \textbackslash{}wedge\allowbreak{}}\par
\noindent Source : FIX-7II-D-0633. \nolinkurl{fixes/FIXLOG_7II.jsonl#633}\par

\par\medskip\noindent\hypertarget{app-F7II-XVIII-035.1}{}\textbf{F7II-XVIII-035.1}\quad p. 278; inherited corrected-layer difference.\par
\noindent Imprimé : \emph{(vide)}\par
\noindent Corrigé : {\ttfamily\footnotesize \textbackslash{}bigl\allowbreak{}}\par
\noindent Source : ADOPT-F7II-XVIII-035.1. \nolinkurl{fr/sga7ii/diplomatic/expose_XVIII_body.tex}; \nolinkurl{fr/sga7ii/current/expose_XVIII_body.tex}\par

\par\medskip\noindent\hypertarget{app-F7II-XVIII-035.2}{}\textbf{F7II-XVIII-035.2}\quad p. 278; designated-witness attribution.\par
\noindent Imprimé : \emph{(vide)}\par
\noindent Corrigé : {\ttfamily\footnotesize \textbackslash{}bigr\allowbreak{}}\par
\noindent Source : WITNESS-F7II-XVIII-035.2. \nolinkurl{fr/sga7ii/diplomatic/expose_XVIII_body.tex}; \nolinkurl{fr/sga7ii/current/expose_XVIII_body.tex}\par

\par\medskip\noindent\hypertarget{app-F7II-XVIII-035.3}{}\textbf{F7II-XVIII-035.3}\quad p. 278; misprint\_pass.\par
\noindent Imprimé : {\ttfamily\footnotesize 3\allowbreak{}}\par
\noindent Corrigé : {\ttfamily\footnotesize 2\allowbreak{}}\par
\noindent Source : 7II:820. \nolinkurl{fixes/FIXLOG_7II_CURRENT.jsonl#825}\par

\par\medskip\noindent\hypertarget{app-F7II-XVIII-035.4}{}\textbf{F7II-XVIII-035.4}\quad p. 278; misprint\_pass.\par
\noindent Imprimé : {\ttfamily\footnotesize 2\allowbreak{}}\par
\noindent Corrigé : {\ttfamily\footnotesize 11\allowbreak{}}\par
\noindent Source : 7II:820. \nolinkurl{fixes/FIXLOG_7II_CURRENT.jsonl#825}\par

\par\medskip\noindent\hypertarget{app-F7II-XVIII-036.1}{}\textbf{F7II-XVIII-036.1}\quad p. 279; printed misprint.\par
\noindent Imprimé : {\ttfamily\footnotesize \textasciicircum{}\allowbreak{}}\par
\noindent Corrigé : {\ttfamily\footnotesize \_\allowbreak{}}\par
\noindent Source : 7II:461; 7II:820. \nolinkurl{fixes/MISPRINTS_7II.json#461}; \nolinkurl{fixes/MISPRINTS_7II.json#820}\par

\par\medskip\noindent\hypertarget{app-F7II-XVIII-037.1}{}\textbf{F7II-XVIII-037.1}\quad p. 280; designated-witness attribution.\par
\noindent Imprimé : {\ttfamily\footnotesize )\allowbreak{}}\par
\noindent Corrigé : \emph{(vide)}\par
\noindent Source : WITNESS-F7II-XVIII-037.1. \nolinkurl{fr/sga7ii/diplomatic/expose_XVIII_body.tex}; \nolinkurl{fr/sga7ii/current/expose_XVIII_body.tex}\par

\par\medskip\noindent\hypertarget{app-F7II-XVIII-038.1}{}\textbf{F7II-XVIII-038.1}\quad p. 280; printed misprint.\par
\noindent Imprimé : \emph{(vide)}\par
\noindent Corrigé : {\ttfamily\footnotesize \}\allowbreak{}}\par
\noindent Source : 7II:822. \nolinkurl{fixes/MISPRINTS_7II.json#822}\par

\par\medskip\noindent\hypertarget{app-F7II-XVIII-038.2}{}\textbf{F7II-XVIII-038.2}\quad p. 280; printed misprint.\par
\noindent Imprimé : {\ttfamily\footnotesize \}\allowbreak{}}\par
\noindent Corrigé : \emph{(vide)}\par
\noindent Source : 7II:822. \nolinkurl{fixes/MISPRINTS_7II.json#822}\par

\par\medskip\noindent\hypertarget{app-F7II-XVIII-038.3}{}\textbf{F7II-XVIII-038.3}\quad p. 280; printed misprint.\par
\noindent Imprimé : \emph{(vide)}\par
\noindent Corrigé : {\ttfamily\footnotesize -\allowbreak{}}\par
\noindent Source : 7II:823. \nolinkurl{fixes/MISPRINTS_7II.json#823}\par

\par\medskip\noindent\hypertarget{app-F7II-XVIII-038.4}{}\textbf{F7II-XVIII-038.4}\quad p. 280; printed misprint.\par
\noindent Imprimé : \emph{(vide)}\par
\noindent Corrigé : {\ttfamily\footnotesize -\allowbreak{}}\par
\noindent Source : 7II:824. \nolinkurl{fixes/MISPRINTS_7II.json#824}\par

\par\medskip\noindent\hypertarget{app-F7II-XVIII-039.1}{}\textbf{F7II-XVIII-039.1}\quad p. 280; printed misprint.\par
\noindent Imprimé : {\ttfamily\footnotesize h\allowbreak{}}\par
\noindent Corrigé : {\ttfamily\footnotesize k\allowbreak{}}\par
\noindent Source : 7II:825. \nolinkurl{fixes/MISPRINTS_7II.json#825}\par

\par\medskip\noindent\hypertarget{app-F7II-XVIII-039.2}{}\textbf{F7II-XVIII-039.2}\quad p. 280; printed misprint.\par
\noindent Imprimé : {\ttfamily\footnotesize \textasciicircum{}\allowbreak{}}\par
\noindent Corrigé : {\ttfamily\footnotesize \_\allowbreak{}}\par
\noindent Source : 7II:826. \nolinkurl{fixes/MISPRINTS_7II.json#826}\par

\par\medskip\noindent\hypertarget{app-F7II-XVIII-040.1}{}\textbf{F7II-XVIII-040.1}\quad p. 281; printed misprint.\par
\noindent Imprimé : \emph{(vide)}\par
\noindent Corrigé : {\ttfamily\footnotesize )\allowbreak{}}\par
\noindent Source : 7II:780. \nolinkurl{fixes/MISPRINTS_7II.json#780}\par

\par\medskip\noindent\hypertarget{app-F7II-XVIII-041.1}{}\textbf{F7II-XVIII-041.1}\quad p. 281; misprint\_pass.\par
\noindent Imprimé : \emph{(vide)}\par
\noindent Corrigé : {\ttfamily\footnotesize )\allowbreak{}}\par
\noindent Source : 7II:469. \nolinkurl{fixes/FIXLOG_7II_CURRENT.jsonl#752}; \nolinkurl{fixes/MISPRINTS_7II.json#469}\par

\par\medskip\noindent\hypertarget{app-F7II-XVIII-041.2}{}\textbf{F7II-XVIII-041.2}\quad p. 281; printed misprint.\par
\noindent Imprimé : \emph{(vide)}\par
\noindent Corrigé : {\ttfamily\footnotesize )\allowbreak{}}\par
\noindent Source : 7II:782. \nolinkurl{fixes/MISPRINTS_7II.json#782}\par

\par\medskip\noindent\hypertarget{app-F7II-XVIII-042.1}{}\textbf{F7II-XVIII-042.1}\quad p. 282; printed misprint.\par
\noindent Imprimé : {\ttfamily\footnotesize 2n\allowbreak{}}\par
\noindent Corrigé : {\ttfamily\footnotesize n\allowbreak{}}\par
\noindent Source : 7II:471. \nolinkurl{fixes/MISPRINTS_7II.json#471}\par

\par\medskip\noindent\hypertarget{app-F7II-XVIII-042.2}{}\textbf{F7II-XVIII-042.2}\quad p. 282; printed misprint.\par
\noindent Imprimé : {\ttfamily\footnotesize 2i\allowbreak{}}\par
\noindent Corrigé : {\ttfamily\footnotesize i\allowbreak{}}\par
\noindent Source : 7II:471. \nolinkurl{fixes/MISPRINTS_7II.json#471}\par

\par\medskip\noindent\hypertarget{app-F7II-XVIII-042.3}{}\textbf{F7II-XVIII-042.3}\quad p. 282; printed misprint.\par
\noindent Imprimé : {\ttfamily\footnotesize 2\allowbreak{}}\par
\noindent Corrigé : {\ttfamily\footnotesize 1)\allowbreak{}}\par
\noindent Source : 7II:471. \nolinkurl{fixes/MISPRINTS_7II.json#471}\par

\par\medskip\noindent\hypertarget{app-F7II-XVIII-043.1}{}\textbf{F7II-XVIII-043.1}\quad p. 282; printed misprint.\par
\noindent Imprimé : {\ttfamily\footnotesize m\allowbreak{}}\par
\noindent Corrigé : {\ttfamily\footnotesize \textbackslash{}tau\allowbreak{}}\par
\noindent Source : 7II:472; 7II:785. \nolinkurl{fixes/MISPRINTS_7II.json#472}; \nolinkurl{fixes/MISPRINTS_7II.json#785}\par

\par\medskip\noindent\hypertarget{app-F7II-XVIII-044.1}{}\textbf{F7II-XVIII-044.1}\quad p. 283; designated-witness attribution.\par
\noindent Imprimé : {\ttfamily\footnotesize \textbackslash{}arrow[dl,\allowbreak{} "\textbackslash{}tau\_*"]\allowbreak{} }\par
\noindent Corrigé : \emph{(vide)}\par
\noindent Source : WITNESS-F7II-XVIII-044.1. \nolinkurl{fr/sga7ii/diplomatic/expose_XVIII_body.tex}; \nolinkurl{fr/sga7ii/current/expose_XVIII_body.tex}\par

\par\medskip\noindent\hypertarget{app-F7II-XVIII-044.2}{}\textbf{F7II-XVIII-044.2}\quad p. 283; designated-witness attribution.\par
\noindent Imprimé : \emph{(vide)}\par
\noindent Corrigé : {\ttfamily\footnotesize \textbackslash{}arrow[ur,\allowbreak{} "\textbackslash{}tau\_*"']\allowbreak{} }\par
\noindent Source : WITNESS-F7II-XVIII-044.2. \nolinkurl{fr/sga7ii/diplomatic/expose_XVIII_body.tex}; \nolinkurl{fr/sga7ii/current/expose_XVIII_body.tex}\par

\par\medskip\noindent\hypertarget{app-F7II-XVIII-045.1}{}\textbf{F7II-XVIII-045.1}\quad p. 283; printed misprint.\par
\noindent Imprimé : {\ttfamily\footnotesize une\allowbreak{}}\par
\noindent Corrigé : {\ttfamily\footnotesize un\allowbreak{}}\par
\noindent Source : 7II:473; 7II:787. \nolinkurl{fixes/MISPRINTS_7II.json#473}; \nolinkurl{fixes/MISPRINTS_7II.json#787}\par

\par\medskip\noindent\hypertarget{app-F7II-XVIII-046.1}{}\textbf{F7II-XVIII-046.1}\quad p. 284; inherited corrected-layer difference.\par
\noindent Imprimé : {\ttfamily\footnotesize *\allowbreak{}}\par
\noindent Corrigé : {\ttfamily\footnotesize \}\textbackslash{}tag\{5.2.4.7\allowbreak{}}\par
\noindent Source : ADOPT-F7II-XVIII-046.1. \nolinkurl{fr/sga7ii/diplomatic/expose_XVIII_body.tex}; \nolinkurl{fr/sga7ii/current/expose_XVIII_body.tex}\par

\par\medskip\noindent\hypertarget{app-F7II-XVIII-046.2}{}\textbf{F7II-XVIII-046.2}\quad p. 284; printed misprint.\par
\noindent Imprimé : {\ttfamily\footnotesize 2\allowbreak{}}\par
\noindent Corrigé : {\ttfamily\footnotesize 1)\allowbreak{}}\par
\noindent Source : 7II:474. \nolinkurl{fixes/MISPRINTS_7II.json#474}\par

\par\medskip\noindent\hypertarget{app-F7II-XVIII-046.3}{}\textbf{F7II-XVIII-046.3}\quad p. 284; designated-witness attribution.\par
\noindent Imprimé : {\ttfamily\footnotesize *\allowbreak{}}\par
\noindent Corrigé : \emph{(vide)}\par
\noindent Source : WITNESS-F7II-XVIII-046.3. \nolinkurl{fr/sga7ii/diplomatic/expose_XVIII_body.tex}; \nolinkurl{fr/sga7ii/current/expose_XVIII_body.tex}\par

\par\medskip\noindent\hypertarget{app-F7II-XVIII-047.1}{}\textbf{F7II-XVIII-047.1}\quad p. 285; printed misprint.\par
\noindent Imprimé : {\ttfamily\footnotesize 5\allowbreak{}}\par
\noindent Corrigé : {\ttfamily\footnotesize 4.3\allowbreak{}}\par
\noindent Source : 7II:790. \nolinkurl{fixes/MISPRINTS_7II.json#790}\par

\par\medskip\noindent\hypertarget{app-F7II-XVIII-048.1}{}\textbf{F7II-XVIII-048.1}\quad p. 285; misprint\_pass.\par
\noindent Imprimé : {\ttfamily\footnotesize hyperplan\allowbreak{}}\par
\noindent Corrigé : {\ttfamily\footnotesize hyperplane\allowbreak{}}\par
\noindent Source : 7II:476. \nolinkurl{fixes/FIXLOG_7II_CURRENT.jsonl#754}\par

\par\medskip\noindent\hypertarget{app-F7II-XVIII-049.1}{}\textbf{F7II-XVIII-049.1}\quad p. 285; printed misprint.\par
\noindent Imprimé : {\ttfamily\footnotesize -\allowbreak{}}\par
\noindent Corrigé : {\ttfamily\footnotesize +\allowbreak{}}\par
\noindent Source : 7II:792. \nolinkurl{fixes/MISPRINTS_7II.json#792}\par

\par\medskip\noindent\hypertarget{app-F7II-XVIII-050.1}{}\textbf{F7II-XVIII-050.1}\quad p. 286; printed misprint.\par
\noindent Imprimé : {\ttfamily\footnotesize ii\allowbreak{}}\par
\noindent Corrigé : {\ttfamily\footnotesize 2i\allowbreak{}}\par
\noindent Source : 7II:478. \nolinkurl{fixes/MISPRINTS_7II.json#478}\par

\par\medskip\noindent\hypertarget{app-F7II-XVIII-050.2}{}\textbf{F7II-XVIII-050.2}\quad p. 286; misprint\_pass.\par
\noindent Imprimé : \emph{(vide)}\par
\noindent Corrigé : {\ttfamily\footnotesize s'\allowbreak{}}\par
\noindent Source : 7II:479. \nolinkurl{fixes/FIXLOG_7II_CURRENT.jsonl#755}; \nolinkurl{fixes/MISPRINTS_7II.json#479}\par

\par\medskip\noindent\hypertarget{app-F7II-XVIII-050.3}{}\textbf{F7II-XVIII-050.3}\quad p. 286; printed misprint.\par
\noindent Imprimé : {\ttfamily\footnotesize propective\allowbreak{}}\par
\noindent Corrigé : {\ttfamily\footnotesize projective\allowbreak{}}\par
\noindent Source : 7II:480. \nolinkurl{fixes/MISPRINTS_7II.json#480}\par

\par\medskip\noindent\hypertarget{app-F7II-XVIII-050.4}{}\textbf{F7II-XVIII-050.4}\quad p. 286; printed misprint.\par
\noindent Imprimé : {\ttfamily\footnotesize 5\allowbreak{}}\par
\noindent Corrigé : {\ttfamily\footnotesize 6\allowbreak{}}\par
\noindent Source : 7II:481. \nolinkurl{fixes/MISPRINTS_7II.json#481}\par

\par\medskip\noindent\hypertarget{app-F7II-XVIII-051.1}{}\textbf{F7II-XVIII-051.1}\quad p. 288; printed misprint.\par
\noindent Imprimé : \emph{(vide)}\par
\noindent Corrigé : {\ttfamily\footnotesize )\allowbreak{}}\par
\noindent Source : 7II:483. \nolinkurl{fixes/MISPRINTS_7II.json#483}\par

\par\medskip\noindent\hypertarget{app-F7II-XVIII-052.1}{}\textbf{F7II-XVIII-052.1}\quad p. 289; printed misprint.\par
\noindent Imprimé : {\ttfamily\footnotesize \textasciicircum{}\allowbreak{}}\par
\noindent Corrigé : {\ttfamily\footnotesize \_\allowbreak{}}\par
\noindent Source : 7II:484. \nolinkurl{fixes/MISPRINTS_7II.json#484}\par

\par\medskip\noindent\hypertarget{app-F7II-XVIII-053.1}{}\textbf{F7II-XVIII-053.1}\quad p. 289; printed misprint.\par
\noindent Imprimé : {\ttfamily\footnotesize injectif\allowbreak{}}\par
\noindent Corrigé : {\ttfamily\footnotesize injective\allowbreak{}}\par
\noindent Source : 7II:485. \nolinkurl{fixes/MISPRINTS_7II.json#485}\par

\par\medskip\noindent\hypertarget{app-F7II-XVIII-054.1}{}\textbf{F7II-XVIII-054.1}\quad p. 290; designated-witness attribution.\par
\noindent Imprimé : \emph{(vide)}\par
\noindent Corrigé : {\ttfamily\footnotesize A)\allowbreak{} \textbackslash{}underline\{de\}\allowbreak{} }\par
\noindent Source : WITNESS-F7II-XVIII-054.1. \nolinkurl{fr/sga7ii/diplomatic/expose_XVIII_body.tex}; \nolinkurl{fr/sga7ii/current/expose_XVIII_body.tex}\par

\par\medskip\noindent\hypertarget{app-F7II-XVIII-054.2}{}\textbf{F7II-XVIII-054.2}\quad p. 290; designated-witness attribution.\par
\noindent Imprimé : {\ttfamily\footnotesize 2)\allowbreak{}}\par
\noindent Corrigé : {\ttfamily\footnotesize 5\allowbreak{}}\par
\noindent Source : WITNESS-F7II-XVIII-054.2. \nolinkurl{fr/sga7ii/diplomatic/expose_XVIII_body.tex}; \nolinkurl{fr/sga7ii/current/expose_XVIII_body.tex}\par

\par\medskip\noindent\hypertarget{app-F7II-XVIII-055.1}{}\textbf{F7II-XVIII-055.1}\quad p. 290; printed misprint.\par
\noindent Imprimé : {\ttfamily\footnotesize m\allowbreak{}}\par
\noindent Corrigé : {\ttfamily\footnotesize p\allowbreak{}}\par
\noindent Source : 7II:488. \nolinkurl{fixes/MISPRINTS_7II.json#488}\par

\par\medskip\noindent\hypertarget{app-F7II-XVIII-056.1}{}\textbf{F7II-XVIII-056.1}\quad p. 291; designated-witness attribution.\par
\noindent Imprimé : {\ttfamily\footnotesize o\allowbreak{}}\par
\noindent Corrigé : {\ttfamily\footnotesize 0\allowbreak{}}\par
\noindent Source : WITNESS-F7II-XVIII-056.1. \nolinkurl{fr/sga7ii/diplomatic/expose_XVIII_body.tex}; \nolinkurl{fr/sga7ii/current/expose_XVIII_body.tex}\par

\par\medskip\noindent\hypertarget{app-F7II-XVIII-056.2}{}\textbf{F7II-XVIII-056.2}\quad p. 291; designated-witness attribution.\par
\noindent Imprimé : {\ttfamily\footnotesize o\allowbreak{}}\par
\noindent Corrigé : {\ttfamily\footnotesize 0\allowbreak{}}\par
\noindent Source : WITNESS-F7II-XVIII-056.2. \nolinkurl{fr/sga7ii/diplomatic/expose_XVIII_body.tex}; \nolinkurl{fr/sga7ii/current/expose_XVIII_body.tex}\par

\par\medskip\noindent\hypertarget{app-F7II-XVIII-056.3}{}\textbf{F7II-XVIII-056.3}\quad p. 291; designated-witness attribution.\par
\noindent Imprimé : {\ttfamily\footnotesize o\allowbreak{}}\par
\noindent Corrigé : {\ttfamily\footnotesize 0\allowbreak{}}\par
\noindent Source : WITNESS-F7II-XVIII-056.3. \nolinkurl{fr/sga7ii/diplomatic/expose_XVIII_body.tex}; \nolinkurl{fr/sga7ii/current/expose_XVIII_body.tex}\par

\par\medskip\noindent\hypertarget{app-F7II-XVIII-056.4}{}\textbf{F7II-XVIII-056.4}\quad p. 291; printed misprint.\par
\noindent Imprimé : \emph{(vide)}\par
\noindent Corrigé : {\ttfamily\footnotesize )\allowbreak{}}\par
\noindent Source : 7II:489; 7II:898. \nolinkurl{fixes/MISPRINTS_7II.json#489}; \nolinkurl{fixes/MISPRINTS_7II.json#898}\par

\par\medskip\noindent\hypertarget{app-F7II-XVIII-056.5}{}\textbf{F7II-XVIII-056.5}\quad p. 292; inherited corrected-layer difference.\par
\noindent Imprimé : {\ttfamily\footnotesize f\allowbreak{}}\par
\noindent Corrigé : {\ttfamily\footnotesize \textbackslash{}rho\allowbreak{}}\par
\noindent Source : ADOPT-F7II-XVIII-056.5. \nolinkurl{fr/sga7ii/diplomatic/expose_XVIII_body.tex}; \nolinkurl{fr/sga7ii/current/expose_XVIII_body.tex}\par

\par\medskip\noindent\hypertarget{app-F7II-XVIII-056.6}{}\textbf{F7II-XVIII-056.6}\quad p. 292; inherited corrected-layer difference.\par
\noindent Imprimé : {\ttfamily\footnotesize f\allowbreak{}}\par
\noindent Corrigé : {\ttfamily\footnotesize \textbackslash{}rho\allowbreak{}}\par
\noindent Source : ADOPT-F7II-XVIII-056.6. \nolinkurl{fr/sga7ii/diplomatic/expose_XVIII_body.tex}; \nolinkurl{fr/sga7ii/current/expose_XVIII_body.tex}\par

\par\medskip\noindent\hypertarget{app-F7II-XVIII-056.7}{}\textbf{F7II-XVIII-056.7}\quad p. 292; printed misprint.\par
\noindent Imprimé : {\ttfamily\footnotesize dégénerescence\allowbreak{}}\par
\noindent Corrigé : {\ttfamily\footnotesize dégénérescence\allowbreak{}}\par
\noindent Source : 7II:495. \nolinkurl{fixes/MISPRINTS_7II.json#495}\par

\par\medskip\noindent\hypertarget{app-F7II-XVIII-057.1}{}\textbf{F7II-XVIII-057.1}\quad p. 292; misprint\_pass.\par
\noindent Imprimé : \emph{(vide)}\par
\noindent Corrigé : {\ttfamily\footnotesize )\allowbreak{}}\par
\noindent Source : 7II:490. \nolinkurl{fixes/FIXLOG_7II_CURRENT.jsonl#756}\par

\par\medskip\noindent\hypertarget{app-F7II-XVIII-058.1}{}\textbf{F7II-XVIII-058.1}\quad p. 292; printed misprint.\par
\noindent Imprimé : \emph{(vide)}\par
\noindent Corrigé : {\ttfamily\footnotesize )\allowbreak{}}\par
\noindent Source : 7II:491; 7II:901. \nolinkurl{fixes/MISPRINTS_7II.json#491}; \nolinkurl{fixes/MISPRINTS_7II.json#901}\par

\par\medskip\noindent\hypertarget{app-F7II-XVIII-058.2}{}\textbf{F7II-XVIII-058.2}\quad p. 292; printed misprint.\par
\noindent Imprimé : {\ttfamily\footnotesize 2h\allowbreak{}}\par
\noindent Corrigé : {\ttfamily\footnotesize 2+h\allowbreak{}}\par
\noindent Source : 7II:492. \nolinkurl{fixes/MISPRINTS_7II.json#492}\par

\par\medskip\noindent\hypertarget{app-F7II-XVIII-058.3}{}\textbf{F7II-XVIII-058.3}\quad p. 292; printed misprint.\par
\noindent Imprimé : \emph{(vide)}\par
\noindent Corrigé : {\ttfamily\footnotesize )\allowbreak{}}\par
\noindent Source : 7II:903. \nolinkurl{fixes/MISPRINTS_7II.json#903}\par

\par\medskip\noindent\hypertarget{app-F7II-XVIII-058.4}{}\textbf{F7II-XVIII-058.4}\quad p. 292; printed misprint.\par
\noindent Imprimé : \emph{(vide)}\par
\noindent Corrigé : {\ttfamily\footnotesize )\allowbreak{}}\par
\noindent Source : 7II:493. \nolinkurl{fixes/MISPRINTS_7II.json#493}\par

\par\medskip\noindent\hypertarget{app-F7II-XVIII-058.5}{}\textbf{F7II-XVIII-058.5}\quad p. 292; printed misprint.\par
\noindent Imprimé : {\ttfamily\footnotesize 8\allowbreak{}}\par
\noindent Corrigé : {\ttfamily\footnotesize 5\allowbreak{}}\par
\noindent Source : 7II:494. \nolinkurl{fixes/MISPRINTS_7II.json#494}\par

\par\medskip\noindent\hypertarget{app-F7II-XVIII-059.1}{}\textbf{F7II-XVIII-059.1}\quad p. 294; printed misprint.\par
\noindent Imprimé : {\ttfamily\footnotesize \{\allowbreak{}}\par
\noindent Corrigé : \emph{(vide)}\par
\noindent Source : 7II:496; 7II:906. \nolinkurl{fixes/MISPRINTS_7II.json#496}; \nolinkurl{fixes/MISPRINTS_7II.json#906}\par

\par\medskip\noindent\hypertarget{app-F7II-XVIII-059.2}{}\textbf{F7II-XVIII-059.2}\quad p. 294; printed misprint.\par
\noindent Imprimé : {\ttfamily\footnotesize \}\allowbreak{}}\par
\noindent Corrigé : \emph{(vide)}\par
\noindent Source : 7II:496; 7II:906. \nolinkurl{fixes/MISPRINTS_7II.json#496}; \nolinkurl{fixes/MISPRINTS_7II.json#906}\par

\par\medskip\noindent\hypertarget{app-F7II-XVIII-059.3}{}\textbf{F7II-XVIII-059.3}\quad p. 294; printed misprint.\par
\noindent Imprimé : {\ttfamily\footnotesize .\allowbreak{}}\par
\noindent Corrigé : \emph{(1 espace(s), 0 retour(s) de ligne dans le témoin TeX)}\par
\noindent Source : 7II:497. \nolinkurl{fixes/MISPRINTS_7II.json#497}\par

\par\medskip\noindent\hypertarget{app-F7II-XVIII-060.1}{}\textbf{F7II-XVIII-060.1}\quad p. 295; designated-witness attribution.\par
\noindent Imprimé : \emph{(vide)}\par
\noindent Corrigé : {\ttfamily\footnotesize \textasciicircum{}\{i+1\}\allowbreak{}}\par
\noindent Source : WITNESS-F7II-XVIII-060.1. \nolinkurl{fr/sga7ii/diplomatic/expose_XVIII_body.tex}; \nolinkurl{fr/sga7ii/current/expose_XVIII_body.tex}\par

\par\medskip\noindent\hypertarget{app-F7II-XVIII-060.2}{}\textbf{F7II-XVIII-060.2}\quad p. 295; designated-witness attribution.\par
\noindent Imprimé : {\ttfamily\footnotesize -1\allowbreak{}}\par
\noindent Corrigé : \emph{(vide)}\par
\noindent Source : WITNESS-F7II-XVIII-060.2. \nolinkurl{fr/sga7ii/diplomatic/expose_XVIII_body.tex}; \nolinkurl{fr/sga7ii/current/expose_XVIII_body.tex}\par

\par\medskip\noindent\hypertarget{app-F7II-XVIII-060.3}{}\textbf{F7II-XVIII-060.3}\quad p. 295; printed misprint.\par
\noindent Imprimé : \emph{(vide)}\par
\noindent Corrigé : {\ttfamily\footnotesize \textbackslash{}rho\_*\allowbreak{}}\par
\noindent Source : 7II:910. \nolinkurl{fixes/MISPRINTS_7II.json#910}\par

\par\medskip\noindent\hypertarget{app-F7II-XVIII-060.4}{}\textbf{F7II-XVIII-060.4}\quad p. 295; designated-witness attribution.\par
\noindent Imprimé : \emph{(vide)}\par
\noindent Corrigé : {\ttfamily\footnotesize  \textbackslash{}geq\allowbreak{} 0\allowbreak{}}\par
\noindent Source : WITNESS-F7II-XVIII-060.4. \nolinkurl{fr/sga7ii/diplomatic/expose_XVIII_body.tex}; \nolinkurl{fr/sga7ii/current/expose_XVIII_body.tex}\par

\par\medskip\noindent\hypertarget{app-F7II-XVIII-060.5}{}\textbf{F7II-XVIII-060.5}\quad p. 295; designated-witness attribution.\par
\noindent Imprimé : \emph{(vide)}\par
\noindent Corrigé : {\ttfamily\footnotesize \textbackslash{};\allowbreak{} }\par
\noindent Source : WITNESS-F7II-XVIII-060.5. \nolinkurl{fr/sga7ii/diplomatic/expose_XVIII_body.tex}; \nolinkurl{fr/sga7ii/current/expose_XVIII_body.tex}\par

\par\medskip\noindent\hypertarget{app-F7II-XVIII-060.6}{}\textbf{F7II-XVIII-060.6}\quad p. 295; printed misprint.\par
\noindent Imprimé : {\ttfamily\footnotesize ,\allowbreak{} }\par
\noindent Corrigé : \emph{(vide)}\par
\noindent Source : 7II:501. \nolinkurl{fixes/MISPRINTS_7II.json#501}\par

\par\medskip\noindent\hypertarget{app-F7II-XVIII-060.7}{}\textbf{F7II-XVIII-060.7}\quad p. 295; printed misprint.\par
\noindent Imprimé : {\ttfamily\footnotesize \textbackslash{}mathbb\{\allowbreak{}}\par
\noindent Corrigé : \emph{(vide)}\par
\noindent Source : 7II:502. \nolinkurl{fixes/MISPRINTS_7II.json#502}\par

\par\medskip\noindent\hypertarget{app-F7II-XVIII-060.8}{}\textbf{F7II-XVIII-060.8}\quad p. 295; printed misprint.\par
\noindent Imprimé : {\ttfamily\footnotesize \}\allowbreak{}}\par
\noindent Corrigé : \emph{(vide)}\par
\noindent Source : 7II:502. \nolinkurl{fixes/MISPRINTS_7II.json#502}\par

\par\medskip\noindent\hypertarget{app-F7II-XVIII-060.9}{}\textbf{F7II-XVIII-060.9}\quad p. 295; printed misprint.\par
\noindent Imprimé : {\ttfamily\footnotesize 1\allowbreak{}}\par
\noindent Corrigé : {\ttfamily\footnotesize n\allowbreak{}}\par
\noindent Source : 7II:502. \nolinkurl{fixes/MISPRINTS_7II.json#502}\par

\par\medskip\noindent\hypertarget{app-F7II-XVIII-060.10}{}\textbf{F7II-XVIII-060.10}\quad p. 295; inherited corrected-layer difference.\par
\noindent Imprimé : \emph{(vide)}\par
\noindent Corrigé : {\ttfamily\footnotesize \textbackslash{}oplus\allowbreak{} \textbackslash{}cdots\allowbreak{} }\par
\noindent Source : ADOPT-F7II-XVIII-060.10. \nolinkurl{fr/sga7ii/diplomatic/expose_XVIII_body.tex}; \nolinkurl{fr/sga7ii/current/expose_XVIII_body.tex}\par

\par\medskip\noindent\hypertarget{app-F7II-XVIII-061.1}{}\textbf{F7II-XVIII-061.1}\quad p. 296; printed misprint.\par
\noindent Imprimé : {\ttfamily\footnotesize ilreste\allowbreak{}}\par
\noindent Corrigé : {\ttfamily\footnotesize il\allowbreak{} reste\allowbreak{}}\par
\noindent Source : 7II:913. \nolinkurl{fixes/MISPRINTS_7II.json#913}\par

\par\medskip\noindent\hypertarget{app-F7II-XVIII-062.1}{}\textbf{F7II-XVIII-062.1}\quad p. 297; printed misprint.\par
\noindent Imprimé : {\ttfamily\footnotesize c\allowbreak{}}\par
\noindent Corrigé : {\ttfamily\footnotesize 0\allowbreak{}}\par
\noindent Source : 7II:504; 7II:914. \nolinkurl{fixes/MISPRINTS_7II.json#504}; \nolinkurl{fixes/MISPRINTS_7II.json#914}\par

\par\medskip\noindent\hypertarget{app-F7II-XVIII-062.2}{}\textbf{F7II-XVIII-062.2}\quad p. 297; printed misprint.\par
\noindent Imprimé : \emph{(vide)}\par
\noindent Corrigé : {\ttfamily\footnotesize  E\allowbreak{}}\par
\noindent Source : 7II:505; 7II:915. \nolinkurl{fixes/MISPRINTS_7II.json#505}; \nolinkurl{fixes/MISPRINTS_7II.json#915}\par

\par\medskip\noindent\hypertarget{app-F7II-XVIII-063.1}{}\textbf{F7II-XVIII-063.1}\quad p. 299; designated-witness attribution.\par
\noindent Imprimé : {\ttfamily\footnotesize :\allowbreak{} }\par
\noindent Corrigé : {\ttfamily\footnotesize \textbackslash{},\allowbreak{}}\par
\noindent Source : WITNESS-F7II-XVIII-063.1. \nolinkurl{fr/sga7ii/diplomatic/expose_XVIII_body.tex}; \nolinkurl{fr/sga7ii/current/expose_XVIII_body.tex}\par

\par\medskip\noindent\hypertarget{app-F7II-XVIII-064.1}{}\textbf{F7II-XVIII-064.1}\quad p. 299; printed misprint.\par
\noindent Imprimé : {\ttfamily\footnotesize 2\allowbreak{}}\par
\noindent Corrigé : {\ttfamily\footnotesize i\allowbreak{}}\par
\noindent Source : 7II:506. \nolinkurl{fixes/MISPRINTS_7II.json#506}\par

\par\medskip\noindent\hypertarget{app-F7II-XVIII-065.1}{}\textbf{F7II-XVIII-065.1}\quad p. 300; printed misprint.\par
\noindent Imprimé : {\ttfamily\footnotesize N\allowbreak{}}\par
\noindent Corrigé : {\ttfamily\footnotesize n\allowbreak{}}\par
\noindent Source : 7II:507. \nolinkurl{fixes/MISPRINTS_7II.json#507}\par

\par\medskip\noindent\hypertarget{app-F7II-XVIII-065.2}{}\textbf{F7II-XVIII-065.2}\quad p. 300; designated-witness attribution.\par
\noindent Imprimé : \emph{(1 espace(s), 0 retour(s) de ligne dans le témoin TeX)}\par
\noindent Corrigé : \emph{(vide)}\par
\noindent Source : WITNESS-F7II-XVIII-065.2. \nolinkurl{fr/sga7ii/diplomatic/expose_XVIII_body.tex}; \nolinkurl{fr/sga7ii/current/expose_XVIII_body.tex}\par

\par\medskip\noindent\hypertarget{app-F7II-XVIII-065.3}{}\textbf{F7II-XVIII-065.3}\quad p. 300; printed misprint.\par
\noindent Imprimé : {\ttfamily\footnotesize riqueur\allowbreak{}}\par
\noindent Corrigé : {\ttfamily\footnotesize rigueur\allowbreak{}}\par
\noindent Source : 7II:508; 7II:918. \nolinkurl{fixes/MISPRINTS_7II.json#508}; \nolinkurl{fixes/MISPRINTS_7II.json#918}\par

\par\medskip\noindent\hypertarget{app-F7II-XVIII-065.4}{}\textbf{F7II-XVIII-065.4}\quad p. 300; printed misprint.\par
\noindent Imprimé : \emph{(vide)}\par
\noindent Corrigé : {\ttfamily\footnotesize (-1\allowbreak{}}\par
\noindent Source : 7II:509; 7II:919. \nolinkurl{fixes/MISPRINTS_7II.json#509}; \nolinkurl{fixes/MISPRINTS_7II.json#919}\par

\par\medskip\noindent\hypertarget{app-F7II-XVIII-066.1}{}\textbf{F7II-XVIII-066.1}\quad p. 300; inherited corrected-layer difference.\par
\noindent Imprimé : \emph{(vide)}\par
\noindent Corrigé : {\ttfamily\footnotesize \textasciicircum{}n\allowbreak{}}\par
\noindent Source : ADOPT-F7II-XVIII-066.1. \nolinkurl{fr/sga7ii/diplomatic/expose_XVIII_body.tex}; \nolinkurl{fr/sga7ii/current/expose_XVIII_body.tex}\par

\par\medskip\noindent\hypertarget{app-F7II-XVIII-067.1}{}\textbf{F7II-XVIII-067.1}\quad p. 301; misprint\_pass.\par
\noindent Imprimé : {\ttfamily\footnotesize -\allowbreak{}}\par
\noindent Corrigé : \emph{(1 espace(s), 0 retour(s) de ligne dans le témoin TeX)}\par
\noindent Source : 7II:510; 7II:954. \nolinkurl{fixes/FIXLOG_7II_CURRENT.jsonl#757}; \nolinkurl{fixes/MISPRINTS_7II.json#510}; \nolinkurl{fixes/MISPRINTS_7II.json#954}\par

\par\medskip\noindent\hypertarget{app-F7II-XVIII-068.1}{}\textbf{F7II-XVIII-068.1}\quad p. 301; printed misprint.\par
\noindent Imprimé : {\ttfamily\footnotesize contenu\allowbreak{}}\par
\noindent Corrigé : {\ttfamily\footnotesize contenue\allowbreak{}}\par
\noindent Source : 7II:511; 7II:955. \nolinkurl{fixes/MISPRINTS_7II.json#511}; \nolinkurl{fixes/MISPRINTS_7II.json#955}\par

\par\medskip\noindent\hypertarget{app-F7II-XVIII-069.1}{}\textbf{F7II-XVIII-069.1}\quad p. 301; designated-witness attribution.\par
\noindent Imprimé : {\ttfamily\footnotesize \_\{\allowbreak{}}\par
\noindent Corrigé : \emph{(vide)}\par
\noindent Source : WITNESS-F7II-XVIII-069.1. \nolinkurl{fr/sga7ii/diplomatic/expose_XVIII_body.tex}; \nolinkurl{fr/sga7ii/current/expose_XVIII_body.tex}\par

\par\medskip\noindent\hypertarget{app-F7II-XVIII-069.2}{}\textbf{F7II-XVIII-069.2}\quad p. 301; designated-witness attribution.\par
\noindent Imprimé : {\ttfamily\footnotesize \}\allowbreak{}}\par
\noindent Corrigé : \emph{(vide)}\par
\noindent Source : WITNESS-F7II-XVIII-069.2. \nolinkurl{fr/sga7ii/diplomatic/expose_XVIII_body.tex}; \nolinkurl{fr/sga7ii/current/expose_XVIII_body.tex}\par

\par\medskip\noindent\hypertarget{app-F7II-XVIII-069.3}{}\textbf{F7II-XVIII-069.3}\quad p. 301; printed misprint.\par
\noindent Imprimé : \emph{(vide)}\par
\noindent Corrigé : {\ttfamily\footnotesize )\allowbreak{}}\par
\noindent Source : 7II:956. \nolinkurl{fixes/MISPRINTS_7II.json#956}\par

\par\medskip\noindent\hypertarget{app-F7II-XVIII-069.4}{}\textbf{F7II-XVIII-069.4}\quad p. 301; designated-witness attribution.\par
\noindent Imprimé : {\ttfamily\footnotesize \textbackslash{}\allowbreak{} }\par
\noindent Corrigé : {\ttfamily\footnotesize \textbackslash{}bigl\allowbreak{}}\par
\noindent Source : WITNESS-F7II-XVIII-069.4. \nolinkurl{fr/sga7ii/diplomatic/expose_XVIII_body.tex}; \nolinkurl{fr/sga7ii/current/expose_XVIII_body.tex}\par

\par\medskip\noindent\hypertarget{app-F7II-XVIII-069.5}{}\textbf{F7II-XVIII-069.5}\quad p. 301; designated-witness attribution.\par
\noindent Imprimé : \emph{(vide)}\par
\noindent Corrigé : {\ttfamily\footnotesize \textbackslash{}bigr)\allowbreak{}}\par
\noindent Source : WITNESS-F7II-XVIII-069.5. \nolinkurl{fr/sga7ii/diplomatic/expose_XVIII_body.tex}; \nolinkurl{fr/sga7ii/current/expose_XVIII_body.tex}\par

\par\medskip\noindent\hypertarget{app-F7II-XVIII-071.1}{}\textbf{F7II-XVIII-071.1}\quad p. 302; inherited corrected-layer difference.\par
\noindent Imprimé : {\ttfamily\footnotesize \textbackslash{}mathrm\{\allowbreak{}}\par
\noindent Corrigé : \emph{(vide)}\par
\noindent Source : ADOPT-F7II-XVIII-071.1. \nolinkurl{fr/sga7ii/diplomatic/expose_XVIII_body.tex}; \nolinkurl{fr/sga7ii/current/expose_XVIII_body.tex}\par

\par\medskip\noindent\hypertarget{app-F7II-XVIII-071.2}{}\textbf{F7II-XVIII-071.2}\quad p. 302; inherited corrected-layer difference.\par
\noindent Imprimé : {\ttfamily\footnotesize \}\allowbreak{} }\par
\noindent Corrigé : \emph{(vide)}\par
\noindent Source : ADOPT-F7II-XVIII-071.2. \nolinkurl{fr/sga7ii/diplomatic/expose_XVIII_body.tex}; \nolinkurl{fr/sga7ii/current/expose_XVIII_body.tex}\par

\par\medskip\noindent\hypertarget{app-F7II-XVIII-071.3}{}\textbf{F7II-XVIII-071.3}\quad p. 302; inherited corrected-layer difference.\par
\noindent Imprimé : {\ttfamily\footnotesize \textbackslash{}mathrm\{\allowbreak{}}\par
\noindent Corrigé : \emph{(vide)}\par
\noindent Source : ADOPT-F7II-XVIII-071.3. \nolinkurl{fr/sga7ii/diplomatic/expose_XVIII_body.tex}; \nolinkurl{fr/sga7ii/current/expose_XVIII_body.tex}\par

\par\medskip\noindent\hypertarget{app-F7II-XVIII-071.4}{}\textbf{F7II-XVIII-071.4}\quad p. 302; inherited corrected-layer difference.\par
\noindent Imprimé : {\ttfamily\footnotesize \}\allowbreak{}}\par
\noindent Corrigé : \emph{(vide)}\par
\noindent Source : ADOPT-F7II-XVIII-071.4. \nolinkurl{fr/sga7ii/diplomatic/expose_XVIII_body.tex}; \nolinkurl{fr/sga7ii/current/expose_XVIII_body.tex}\par

\par\medskip\noindent\hypertarget{app-F7II-XVIII-071.5}{}\textbf{F7II-XVIII-071.5}\quad p. 302; inherited corrected-layer difference.\par
\noindent Imprimé : {\ttfamily\footnotesize \textbackslash{}mathrm\{\allowbreak{}}\par
\noindent Corrigé : \emph{(vide)}\par
\noindent Source : ADOPT-F7II-XVIII-071.5. \nolinkurl{fr/sga7ii/diplomatic/expose_XVIII_body.tex}; \nolinkurl{fr/sga7ii/current/expose_XVIII_body.tex}\par

\par\medskip\noindent\hypertarget{app-F7II-XVIII-071.6}{}\textbf{F7II-XVIII-071.6}\quad p. 302; inherited corrected-layer difference.\par
\noindent Imprimé : {\ttfamily\footnotesize \}\allowbreak{}}\par
\noindent Corrigé : \emph{(vide)}\par
\noindent Source : ADOPT-F7II-XVIII-071.6. \nolinkurl{fr/sga7ii/diplomatic/expose_XVIII_body.tex}; \nolinkurl{fr/sga7ii/current/expose_XVIII_body.tex}\par

\par\medskip\noindent\hypertarget{app-F7II-XVIII-071.7}{}\textbf{F7II-XVIII-071.7}\quad p. 302; inherited corrected-layer difference.\par
\noindent Imprimé : {\ttfamily\footnotesize 1\textbackslash{}mathrm\{H\}\allowbreak{}}\par
\noindent Corrigé : {\ttfamily\footnotesize 1H\allowbreak{}}\par
\noindent Source : ADOPT-F7II-XVIII-071.7. \nolinkurl{fr/sga7ii/diplomatic/expose_XVIII_body.tex}; \nolinkurl{fr/sga7ii/current/expose_XVIII_body.tex}\par

\par\medskip\noindent\hypertarget{app-F7II-XVIII-071.8}{}\textbf{F7II-XVIII-071.8}\quad p. 302; inherited corrected-layer difference.\par
\noindent Imprimé : {\ttfamily\footnotesize \textbackslash{}mathrm\{\allowbreak{}}\par
\noindent Corrigé : \emph{(vide)}\par
\noindent Source : ADOPT-F7II-XVIII-071.8. \nolinkurl{fr/sga7ii/diplomatic/expose_XVIII_body.tex}; \nolinkurl{fr/sga7ii/current/expose_XVIII_body.tex}\par

\par\medskip\noindent\hypertarget{app-F7II-XVIII-071.9}{}\textbf{F7II-XVIII-071.9}\quad p. 302; inherited corrected-layer difference.\par
\noindent Imprimé : {\ttfamily\footnotesize \}\allowbreak{}}\par
\noindent Corrigé : \emph{(vide)}\par
\noindent Source : ADOPT-F7II-XVIII-071.9. \nolinkurl{fr/sga7ii/diplomatic/expose_XVIII_body.tex}; \nolinkurl{fr/sga7ii/current/expose_XVIII_body.tex}\par

\par\medskip\noindent\hypertarget{app-F7II-XVIII-071.10}{}\textbf{F7II-XVIII-071.10}\quad p. 302; inherited corrected-layer difference.\par
\noindent Imprimé : {\ttfamily\footnotesize 1\textbackslash{}mathrm\{H\}\allowbreak{}}\par
\noindent Corrigé : {\ttfamily\footnotesize 1H\allowbreak{}}\par
\noindent Source : ADOPT-F7II-XVIII-071.10. \nolinkurl{fr/sga7ii/diplomatic/expose_XVIII_body.tex}; \nolinkurl{fr/sga7ii/current/expose_XVIII_body.tex}\par

\par\medskip\noindent\hypertarget{app-F7II-XVIII-071.11}{}\textbf{F7II-XVIII-071.11}\quad p. 302; inherited corrected-layer difference.\par
\noindent Imprimé : {\ttfamily\footnotesize \textbackslash{}mathrm\{\allowbreak{}}\par
\noindent Corrigé : \emph{(vide)}\par
\noindent Source : ADOPT-F7II-XVIII-071.11. \nolinkurl{fr/sga7ii/diplomatic/expose_XVIII_body.tex}; \nolinkurl{fr/sga7ii/current/expose_XVIII_body.tex}\par

\par\medskip\noindent\hypertarget{app-F7II-XVIII-071.12}{}\textbf{F7II-XVIII-071.12}\quad p. 302; inherited corrected-layer difference.\par
\noindent Imprimé : {\ttfamily\footnotesize \}\allowbreak{}}\par
\noindent Corrigé : \emph{(vide)}\par
\noindent Source : ADOPT-F7II-XVIII-071.12. \nolinkurl{fr/sga7ii/diplomatic/expose_XVIII_body.tex}; \nolinkurl{fr/sga7ii/current/expose_XVIII_body.tex}\par

\par\medskip\noindent\hypertarget{app-F7II-XVIII-071.13}{}\textbf{F7II-XVIII-071.13}\quad p. 302; designated-witness attribution.\par
\noindent Imprimé : {\ttfamily\footnotesize \textbackslash{}mathrm\{\allowbreak{}}\par
\noindent Corrigé : \emph{(vide)}\par
\noindent Source : WITNESS-F7II-XVIII-071.13. \nolinkurl{fr/sga7ii/diplomatic/expose_XVIII_body.tex}; \nolinkurl{fr/sga7ii/current/expose_XVIII_body.tex}\par

\par\medskip\noindent\hypertarget{app-F7II-XVIII-071.14}{}\textbf{F7II-XVIII-071.14}\quad p. 302; designated-witness attribution.\par
\noindent Imprimé : {\ttfamily\footnotesize \}\allowbreak{}}\par
\noindent Corrigé : \emph{(vide)}\par
\noindent Source : WITNESS-F7II-XVIII-071.14. \nolinkurl{fr/sga7ii/diplomatic/expose_XVIII_body.tex}; \nolinkurl{fr/sga7ii/current/expose_XVIII_body.tex}\par

\par\medskip\noindent\hypertarget{app-F7II-XVIII-071.15}{}\textbf{F7II-XVIII-071.15}\quad p. 302; designated-witness attribution.\par
\noindent Imprimé : {\ttfamily\footnotesize \textbackslash{}mathrm\{\allowbreak{}}\par
\noindent Corrigé : \emph{(vide)}\par
\noindent Source : WITNESS-F7II-XVIII-071.15. \nolinkurl{fr/sga7ii/diplomatic/expose_XVIII_body.tex}; \nolinkurl{fr/sga7ii/current/expose_XVIII_body.tex}\par

\par\medskip\noindent\hypertarget{app-F7II-XVIII-071.16}{}\textbf{F7II-XVIII-071.16}\quad p. 302; designated-witness attribution.\par
\noindent Imprimé : {\ttfamily\footnotesize \}\allowbreak{}}\par
\noindent Corrigé : \emph{(vide)}\par
\noindent Source : WITNESS-F7II-XVIII-071.16. \nolinkurl{fr/sga7ii/diplomatic/expose_XVIII_body.tex}; \nolinkurl{fr/sga7ii/current/expose_XVIII_body.tex}\par

\par\medskip\noindent\hypertarget{app-F7II-XVIII-071.17}{}\textbf{F7II-XVIII-071.17}\quad p. 302; designated-witness attribution.\par
\noindent Imprimé : {\ttfamily\footnotesize \textbackslash{}mathrm\{\allowbreak{}}\par
\noindent Corrigé : \emph{(vide)}\par
\noindent Source : WITNESS-F7II-XVIII-071.17. \nolinkurl{fr/sga7ii/diplomatic/expose_XVIII_body.tex}; \nolinkurl{fr/sga7ii/current/expose_XVIII_body.tex}\par

\par\medskip\noindent\hypertarget{app-F7II-XVIII-071.18}{}\textbf{F7II-XVIII-071.18}\quad p. 302; designated-witness attribution.\par
\noindent Imprimé : {\ttfamily\footnotesize \}\allowbreak{}}\par
\noindent Corrigé : \emph{(vide)}\par
\noindent Source : WITNESS-F7II-XVIII-071.18. \nolinkurl{fr/sga7ii/diplomatic/expose_XVIII_body.tex}; \nolinkurl{fr/sga7ii/current/expose_XVIII_body.tex}\par

\par\medskip\noindent\hypertarget{app-F7II-XVIII-071.19}{}\textbf{F7II-XVIII-071.19}\quad p. 302; designated-witness attribution.\par
\noindent Imprimé : {\ttfamily\footnotesize \textbackslash{}mathrm\{\allowbreak{}}\par
\noindent Corrigé : \emph{(vide)}\par
\noindent Source : WITNESS-F7II-XVIII-071.19. \nolinkurl{fr/sga7ii/diplomatic/expose_XVIII_body.tex}; \nolinkurl{fr/sga7ii/current/expose_XVIII_body.tex}\par

\par\medskip\noindent\hypertarget{app-F7II-XVIII-071.20}{}\textbf{F7II-XVIII-071.20}\quad p. 302; designated-witness attribution.\par
\noindent Imprimé : {\ttfamily\footnotesize \}\allowbreak{}}\par
\noindent Corrigé : \emph{(vide)}\par
\noindent Source : WITNESS-F7II-XVIII-071.20. \nolinkurl{fr/sga7ii/diplomatic/expose_XVIII_body.tex}; \nolinkurl{fr/sga7ii/current/expose_XVIII_body.tex}\par

\par\medskip\noindent\hypertarget{app-F7II-XVIII-071.21}{}\textbf{F7II-XVIII-071.21}\quad p. 302; printed misprint.\par
\noindent Imprimé : {\ttfamily\footnotesize \_\{\allowbreak{}}\par
\noindent Corrigé : \emph{(vide)}\par
\noindent Source : 7II:961. \nolinkurl{fixes/MISPRINTS_7II.json#961}\par

\par\medskip\noindent\hypertarget{app-F7II-XVIII-071.22}{}\textbf{F7II-XVIII-071.22}\quad p. 302; designated-witness attribution.\par
\noindent Imprimé : {\ttfamily\footnotesize \}\allowbreak{}}\par
\noindent Corrigé : \emph{(vide)}\par
\noindent Source : WITNESS-F7II-XVIII-071.22. \nolinkurl{fr/sga7ii/diplomatic/expose_XVIII_body.tex}; \nolinkurl{fr/sga7ii/current/expose_XVIII_body.tex}\par

\par\medskip\noindent\hypertarget{app-F7II-XVIII-071.23}{}\textbf{F7II-XVIII-071.23}\quad p. 302; designated-witness attribution.\par
\noindent Imprimé : {\ttfamily\footnotesize \textbackslash{}mathrm\{\allowbreak{}}\par
\noindent Corrigé : \emph{(vide)}\par
\noindent Source : WITNESS-F7II-XVIII-071.23. \nolinkurl{fr/sga7ii/diplomatic/expose_XVIII_body.tex}; \nolinkurl{fr/sga7ii/current/expose_XVIII_body.tex}\par

\par\medskip\noindent\hypertarget{app-F7II-XVIII-071.24}{}\textbf{F7II-XVIII-071.24}\quad p. 302; designated-witness attribution.\par
\noindent Imprimé : {\ttfamily\footnotesize \}\allowbreak{}}\par
\noindent Corrigé : \emph{(vide)}\par
\noindent Source : WITNESS-F7II-XVIII-071.24. \nolinkurl{fr/sga7ii/diplomatic/expose_XVIII_body.tex}; \nolinkurl{fr/sga7ii/current/expose_XVIII_body.tex}\par

\par\medskip\noindent\hypertarget{app-F7II-XVIII-072.1}{}\textbf{F7II-XVIII-072.1}\quad p. 302; designated-witness attribution.\par
\noindent Imprimé : {\ttfamily\footnotesize \textbackslash{}mathrm\{\allowbreak{}}\par
\noindent Corrigé : \emph{(vide)}\par
\noindent Source : WITNESS-F7II-XVIII-072.1. \nolinkurl{fr/sga7ii/diplomatic/expose_XVIII_body.tex}; \nolinkurl{fr/sga7ii/current/expose_XVIII_body.tex}\par

\par\medskip\noindent\hypertarget{app-F7II-XVIII-072.2}{}\textbf{F7II-XVIII-072.2}\quad p. 302; designated-witness attribution.\par
\noindent Imprimé : {\ttfamily\footnotesize \}\allowbreak{}}\par
\noindent Corrigé : \emph{(vide)}\par
\noindent Source : WITNESS-F7II-XVIII-072.2. \nolinkurl{fr/sga7ii/diplomatic/expose_XVIII_body.tex}; \nolinkurl{fr/sga7ii/current/expose_XVIII_body.tex}\par

\par\medskip\noindent\hypertarget{app-F7II-XVIII-072.3}{}\textbf{F7II-XVIII-072.3}\quad p. 302; printed misprint.\par
\noindent Imprimé : \emph{(vide)}\par
\noindent Corrigé : {\ttfamily\footnotesize \textasciicircum{}i\allowbreak{}}\par
\noindent Source : 7II:516; 7II:959. \nolinkurl{fixes/MISPRINTS_7II.json#516}; \nolinkurl{fixes/MISPRINTS_7II.json#959}\par

\par\medskip\noindent\hypertarget{app-F7II-XVIII-072.4}{}\textbf{F7II-XVIII-072.4}\quad p. 302; designated-witness attribution.\par
\noindent Imprimé : {\ttfamily\footnotesize \textbackslash{}mathrm\{\allowbreak{}}\par
\noindent Corrigé : \emph{(vide)}\par
\noindent Source : WITNESS-F7II-XVIII-072.4. \nolinkurl{fr/sga7ii/diplomatic/expose_XVIII_body.tex}; \nolinkurl{fr/sga7ii/current/expose_XVIII_body.tex}\par

\par\medskip\noindent\hypertarget{app-F7II-XVIII-072.5}{}\textbf{F7II-XVIII-072.5}\quad p. 302; designated-witness attribution.\par
\noindent Imprimé : {\ttfamily\footnotesize \}\_\{\allowbreak{}}\par
\noindent Corrigé : \emph{(vide)}\par
\noindent Source : WITNESS-F7II-XVIII-072.5. \nolinkurl{fr/sga7ii/diplomatic/expose_XVIII_body.tex}; \nolinkurl{fr/sga7ii/current/expose_XVIII_body.tex}\par

\par\medskip\noindent\hypertarget{app-F7II-XVIII-072.6}{}\textbf{F7II-XVIII-072.6}\quad p. 302; designated-witness attribution.\par
\noindent Imprimé : {\ttfamily\footnotesize \textbackslash{}mathrm\{\allowbreak{}}\par
\noindent Corrigé : \emph{(vide)}\par
\noindent Source : WITNESS-F7II-XVIII-072.6. \nolinkurl{fr/sga7ii/diplomatic/expose_XVIII_body.tex}; \nolinkurl{fr/sga7ii/current/expose_XVIII_body.tex}\par

\par\medskip\noindent\hypertarget{app-F7II-XVIII-072.7}{}\textbf{F7II-XVIII-072.7}\quad p. 302; designated-witness attribution.\par
\noindent Imprimé : {\ttfamily\footnotesize \}\}\allowbreak{}}\par
\noindent Corrigé : \emph{(vide)}\par
\noindent Source : WITNESS-F7II-XVIII-072.7. \nolinkurl{fr/sga7ii/diplomatic/expose_XVIII_body.tex}; \nolinkurl{fr/sga7ii/current/expose_XVIII_body.tex}\par

\par\medskip\noindent\hypertarget{app-F7II-XVIII-072.8}{}\textbf{F7II-XVIII-072.8}\quad p. 302; designated-witness attribution.\par
\noindent Imprimé : {\ttfamily\footnotesize \textbackslash{}mathrm\{\allowbreak{}}\par
\noindent Corrigé : \emph{(vide)}\par
\noindent Source : WITNESS-F7II-XVIII-072.8. \nolinkurl{fr/sga7ii/diplomatic/expose_XVIII_body.tex}; \nolinkurl{fr/sga7ii/current/expose_XVIII_body.tex}\par

\par\medskip\noindent\hypertarget{app-F7II-XVIII-072.9}{}\textbf{F7II-XVIII-072.9}\quad p. 302; designated-witness attribution.\par
\noindent Imprimé : {\ttfamily\footnotesize \}\allowbreak{}}\par
\noindent Corrigé : \emph{(vide)}\par
\noindent Source : WITNESS-F7II-XVIII-072.9. \nolinkurl{fr/sga7ii/diplomatic/expose_XVIII_body.tex}; \nolinkurl{fr/sga7ii/current/expose_XVIII_body.tex}\par

\par\medskip\noindent\hypertarget{app-F7II-XVIII-072.10}{}\textbf{F7II-XVIII-072.10}\quad p. 302; designated-witness attribution.\par
\noindent Imprimé : {\ttfamily\footnotesize \textbackslash{}mathrm\{\allowbreak{}}\par
\noindent Corrigé : \emph{(vide)}\par
\noindent Source : WITNESS-F7II-XVIII-072.10. \nolinkurl{fr/sga7ii/diplomatic/expose_XVIII_body.tex}; \nolinkurl{fr/sga7ii/current/expose_XVIII_body.tex}\par

\par\medskip\noindent\hypertarget{app-F7II-XVIII-072.11}{}\textbf{F7II-XVIII-072.11}\quad p. 302; designated-witness attribution.\par
\noindent Imprimé : {\ttfamily\footnotesize \}\allowbreak{}}\par
\noindent Corrigé : \emph{(vide)}\par
\noindent Source : WITNESS-F7II-XVIII-072.11. \nolinkurl{fr/sga7ii/diplomatic/expose_XVIII_body.tex}; \nolinkurl{fr/sga7ii/current/expose_XVIII_body.tex}\par

\par\medskip\noindent\hypertarget{app-F7II-XVIII-072.12}{}\textbf{F7II-XVIII-072.12}\quad p. 302; designated-witness attribution.\par
\noindent Imprimé : {\ttfamily\footnotesize \textbackslash{}mathrm\{\allowbreak{}}\par
\noindent Corrigé : \emph{(vide)}\par
\noindent Source : WITNESS-F7II-XVIII-072.12. \nolinkurl{fr/sga7ii/diplomatic/expose_XVIII_body.tex}; \nolinkurl{fr/sga7ii/current/expose_XVIII_body.tex}\par

\par\medskip\noindent\hypertarget{app-F7II-XVIII-072.13}{}\textbf{F7II-XVIII-072.13}\quad p. 302; designated-witness attribution.\par
\noindent Imprimé : {\ttfamily\footnotesize \}\allowbreak{}}\par
\noindent Corrigé : \emph{(vide)}\par
\noindent Source : WITNESS-F7II-XVIII-072.13. \nolinkurl{fr/sga7ii/diplomatic/expose_XVIII_body.tex}; \nolinkurl{fr/sga7ii/current/expose_XVIII_body.tex}\par

\par\medskip\noindent\hypertarget{app-F7II-XVIII-072.14}{}\textbf{F7II-XVIII-072.14}\quad p. 302; printed misprint.\par
\noindent Imprimé : {\ttfamily\footnotesize \textbackslash{}cdot\allowbreak{}}\par
\noindent Corrigé : {\ttfamily\footnotesize i\allowbreak{}}\par
\noindent Source : 7II:960. \nolinkurl{fixes/MISPRINTS_7II.json#960}\par

\par\medskip\noindent\hypertarget{app-F7II-XVIII-072.15}{}\textbf{F7II-XVIII-072.15}\quad p. 302; printed misprint.\par
\noindent Imprimé : {\ttfamily\footnotesize hyperplansdont\allowbreak{}}\par
\noindent Corrigé : {\ttfamily\footnotesize hyperplans\allowbreak{} dont\allowbreak{}}\par
\noindent Source : 7II:514; 7II:958. \nolinkurl{fixes/MISPRINTS_7II.json#514}; \nolinkurl{fixes/MISPRINTS_7II.json#958}\par

\par\medskip\noindent\hypertarget{app-F7II-XVIII-074.1}{}\textbf{F7II-XVIII-074.1}\quad p. 303; printed misprint.\par
\noindent Imprimé : \emph{(vide)}\par
\noindent Corrigé : {\ttfamily\footnotesize )\allowbreak{}}\par
\noindent Source : 7II:962. \nolinkurl{fixes/MISPRINTS_7II.json#962}\par

\par\medskip\noindent\hypertarget{app-F7II-XVIII-075.1}{}\textbf{F7II-XVIII-075.1}\quad p. 303; inherited corrected-layer difference.\par
\noindent Imprimé : \emph{(vide)}\par
\noindent Corrigé : \emph{(1 espace(s), 0 retour(s) de ligne dans le témoin TeX)}\par
\noindent Source : ADOPT-F7II-XVIII-075.1. \nolinkurl{fr/sga7ii/diplomatic/expose_XVIII_body.tex}; \nolinkurl{fr/sga7ii/current/expose_XVIII_body.tex}\par

\par\medskip\noindent\hypertarget{app-F7II-XVIII-075.2}{}\textbf{F7II-XVIII-075.2}\quad p. 303; inherited corrected-layer difference.\par
\noindent Imprimé : {\ttfamily\footnotesize \textbackslash{}mathrm\{\allowbreak{}}\par
\noindent Corrigé : \emph{(vide)}\par
\noindent Source : ADOPT-F7II-XVIII-075.2. \nolinkurl{fr/sga7ii/diplomatic/expose_XVIII_body.tex}; \nolinkurl{fr/sga7ii/current/expose_XVIII_body.tex}\par

\par\medskip\noindent\hypertarget{app-F7II-XVIII-075.3}{}\textbf{F7II-XVIII-075.3}\quad p. 303; inherited corrected-layer difference.\par
\noindent Imprimé : {\ttfamily\footnotesize \}\allowbreak{}}\par
\noindent Corrigé : \emph{(vide)}\par
\noindent Source : ADOPT-F7II-XVIII-075.3. \nolinkurl{fr/sga7ii/diplomatic/expose_XVIII_body.tex}; \nolinkurl{fr/sga7ii/current/expose_XVIII_body.tex}\par

\par\medskip\noindent\hypertarget{app-F7II-XVIII-075.4}{}\textbf{F7II-XVIII-075.4}\quad p. 303; inherited corrected-layer difference.\par
\noindent Imprimé : {\ttfamily\footnotesize o\allowbreak{}}\par
\noindent Corrigé : {\ttfamily\footnotesize 0\allowbreak{}}\par
\noindent Source : ADOPT-F7II-XVIII-075.4. \nolinkurl{fr/sga7ii/diplomatic/expose_XVIII_body.tex}; \nolinkurl{fr/sga7ii/current/expose_XVIII_body.tex}\par

\par\medskip\noindent\hypertarget{app-F7II-XVIII-075.5}{}\textbf{F7II-XVIII-075.5}\quad p. 303; inherited corrected-layer difference.\par
\noindent Imprimé : {\ttfamily\footnotesize \textbackslash{}mathrm\{\allowbreak{}}\par
\noindent Corrigé : \emph{(vide)}\par
\noindent Source : ADOPT-F7II-XVIII-075.5. \nolinkurl{fr/sga7ii/diplomatic/expose_XVIII_body.tex}; \nolinkurl{fr/sga7ii/current/expose_XVIII_body.tex}\par

\par\medskip\noindent\hypertarget{app-F7II-XVIII-075.6}{}\textbf{F7II-XVIII-075.6}\quad p. 303; inherited corrected-layer difference.\par
\noindent Imprimé : {\ttfamily\footnotesize \}\allowbreak{}}\par
\noindent Corrigé : \emph{(vide)}\par
\noindent Source : ADOPT-F7II-XVIII-075.6. \nolinkurl{fr/sga7ii/diplomatic/expose_XVIII_body.tex}; \nolinkurl{fr/sga7ii/current/expose_XVIII_body.tex}\par

\par\medskip\noindent\hypertarget{app-F7II-XVIII-075.7}{}\textbf{F7II-XVIII-075.7}\quad p. 303; inherited corrected-layer difference.\par
\noindent Imprimé : {\ttfamily\footnotesize \textbackslash{}overset\{\allowbreak{}}\par
\noindent Corrigé : {\ttfamily\footnotesize \textbackslash{}xrightarrow[\allowbreak{}}\par
\noindent Source : ADOPT-F7II-XVIII-075.7. \nolinkurl{fr/sga7ii/diplomatic/expose_XVIII_body.tex}; \nolinkurl{fr/sga7ii/current/expose_XVIII_body.tex}\par

\par\medskip\noindent\hypertarget{app-F7II-XVIII-075.8}{}\textbf{F7II-XVIII-075.8}\quad p. 303; inherited corrected-layer difference.\par
\noindent Imprimé : {\ttfamily\footnotesize \}\{\textbackslash{}underset\allowbreak{}}\par
\noindent Corrigé : {\ttfamily\footnotesize ]\allowbreak{}}\par
\noindent Source : ADOPT-F7II-XVIII-075.8. \nolinkurl{fr/sga7ii/diplomatic/expose_XVIII_body.tex}; \nolinkurl{fr/sga7ii/current/expose_XVIII_body.tex}\par

\par\medskip\noindent\hypertarget{app-F7II-XVIII-075.9}{}\textbf{F7II-XVIII-075.9}\quad p. 303; inherited corrected-layer difference.\par
\noindent Imprimé : {\ttfamily\footnotesize \{\textbackslash{}longrightarrow\}\}\allowbreak{}}\par
\noindent Corrigé : \emph{(vide)}\par
\noindent Source : ADOPT-F7II-XVIII-075.9. \nolinkurl{fr/sga7ii/diplomatic/expose_XVIII_body.tex}; \nolinkurl{fr/sga7ii/current/expose_XVIII_body.tex}\par

\par\medskip\noindent\hypertarget{app-F7II-XVIII-075.10}{}\textbf{F7II-XVIII-075.10}\quad p. 303; inherited corrected-layer difference.\par
\noindent Imprimé : {\ttfamily\footnotesize \textbackslash{}mathrm\{\allowbreak{}}\par
\noindent Corrigé : \emph{(vide)}\par
\noindent Source : ADOPT-F7II-XVIII-075.10. \nolinkurl{fr/sga7ii/diplomatic/expose_XVIII_body.tex}; \nolinkurl{fr/sga7ii/current/expose_XVIII_body.tex}\par

\par\medskip\noindent\hypertarget{app-F7II-XVIII-075.11}{}\textbf{F7II-XVIII-075.11}\quad p. 303; inherited corrected-layer difference.\par
\noindent Imprimé : {\ttfamily\footnotesize \}\allowbreak{}}\par
\noindent Corrigé : \emph{(vide)}\par
\noindent Source : ADOPT-F7II-XVIII-075.11. \nolinkurl{fr/sga7ii/diplomatic/expose_XVIII_body.tex}; \nolinkurl{fr/sga7ii/current/expose_XVIII_body.tex}\par

\par\medskip\noindent\hypertarget{app-F7II-XVIII-075.12}{}\textbf{F7II-XVIII-075.12}\quad p. 303; inherited corrected-layer difference.\par
\noindent Imprimé : {\ttfamily\footnotesize o\allowbreak{}}\par
\noindent Corrigé : {\ttfamily\footnotesize 0\allowbreak{}}\par
\noindent Source : ADOPT-F7II-XVIII-075.12. \nolinkurl{fr/sga7ii/diplomatic/expose_XVIII_body.tex}; \nolinkurl{fr/sga7ii/current/expose_XVIII_body.tex}\par

\par\medskip\noindent\hypertarget{app-F7II-XVIII-075.13}{}\textbf{F7II-XVIII-075.13}\quad p. 303; inherited corrected-layer difference.\par
\noindent Imprimé : {\ttfamily\footnotesize \textbackslash{}mathrm\{\allowbreak{}}\par
\noindent Corrigé : \emph{(vide)}\par
\noindent Source : ADOPT-F7II-XVIII-075.13. \nolinkurl{fr/sga7ii/diplomatic/expose_XVIII_body.tex}; \nolinkurl{fr/sga7ii/current/expose_XVIII_body.tex}\par

\par\medskip\noindent\hypertarget{app-F7II-XVIII-075.14}{}\textbf{F7II-XVIII-075.14}\quad p. 303; inherited corrected-layer difference.\par
\noindent Imprimé : {\ttfamily\footnotesize \}\allowbreak{}}\par
\noindent Corrigé : \emph{(vide)}\par
\noindent Source : ADOPT-F7II-XVIII-075.14. \nolinkurl{fr/sga7ii/diplomatic/expose_XVIII_body.tex}; \nolinkurl{fr/sga7ii/current/expose_XVIII_body.tex}\par

\par\medskip\noindent\hypertarget{app-F7II-XVIII-075.15}{}\textbf{F7II-XVIII-075.15}\quad p. 303; designated-witness attribution.\par
\noindent Imprimé : {\ttfamily\footnotesize \textbackslash{}mathrm\{\allowbreak{}}\par
\noindent Corrigé : \emph{(vide)}\par
\noindent Source : WITNESS-F7II-XVIII-075.15. \nolinkurl{fr/sga7ii/diplomatic/expose_XVIII_body.tex}; \nolinkurl{fr/sga7ii/current/expose_XVIII_body.tex}\par

\par\medskip\noindent\hypertarget{app-F7II-XVIII-075.16}{}\textbf{F7II-XVIII-075.16}\quad p. 303; designated-witness attribution.\par
\noindent Imprimé : {\ttfamily\footnotesize \}\allowbreak{}}\par
\noindent Corrigé : \emph{(vide)}\par
\noindent Source : WITNESS-F7II-XVIII-075.16. \nolinkurl{fr/sga7ii/diplomatic/expose_XVIII_body.tex}; \nolinkurl{fr/sga7ii/current/expose_XVIII_body.tex}\par

\par\medskip\noindent\hypertarget{app-F7II-XVIII-075.17}{}\textbf{F7II-XVIII-075.17}\quad p. 303; designated-witness attribution.\par
\noindent Imprimé : {\ttfamily\footnotesize [\allowbreak{}}\par
\noindent Corrigé : \emph{(1 espace(s), 0 retour(s) de ligne dans le témoin TeX)}\par
\noindent Source : WITNESS-F7II-XVIII-075.17. \nolinkurl{fr/sga7ii/diplomatic/expose_XVIII_body.tex}; \nolinkurl{fr/sga7ii/current/expose_XVIII_body.tex}\par

\par\medskip\noindent\hypertarget{app-F7II-XVIII-075.18}{}\textbf{F7II-XVIII-075.18}\quad p. 303; designated-witness attribution.\par
\noindent Imprimé : {\ttfamily\footnotesize rel.\allowbreak{}}\par
\noindent Corrigé : {\ttfamily\footnotesize relativement\allowbreak{}}\par
\noindent Source : WITNESS-F7II-XVIII-075.18. \nolinkurl{fr/sga7ii/diplomatic/expose_XVIII_body.tex}; \nolinkurl{fr/sga7ii/current/expose_XVIII_body.tex}\par

\par\medskip\noindent\hypertarget{app-F7II-XVIII-075.19}{}\textbf{F7II-XVIII-075.19}\quad p. 303; printed misprint.\par
\noindent Imprimé : {\ttfamily\footnotesize \textbackslash{}alpha\allowbreak{}}\par
\noindent Corrigé : \emph{(vide)}\par
\noindent Source : 7II:518; 7II:963. \nolinkurl{fixes/MISPRINTS_7II.json#518}; \nolinkurl{fixes/MISPRINTS_7II.json#963}\par

\par\medskip\noindent\hypertarget{app-F7II-XVIII-075.20}{}\textbf{F7II-XVIII-075.20}\quad p. 303; designated-witness attribution.\par
\noindent Imprimé : {\ttfamily\footnotesize ]\allowbreak{}}\par
\noindent Corrigé : {\ttfamily\footnotesize  si\allowbreak{} et\allowbreak{} seulement\allowbreak{} si\allowbreak{}}\par
\noindent Source : WITNESS-F7II-XVIII-075.20. \nolinkurl{fr/sga7ii/diplomatic/expose_XVIII_body.tex}; \nolinkurl{fr/sga7ii/current/expose_XVIII_body.tex}\par

\par\medskip\noindent\hypertarget{app-F7II-XVIII-075.21}{}\textbf{F7II-XVIII-075.21}\quad p. 303; designated-witness attribution.\par
\noindent Imprimé : {\ttfamily\footnotesize \textbackslash{}Longrightarrow\allowbreak{} }\par
\noindent Corrigé : \emph{(vide)}\par
\noindent Source : WITNESS-F7II-XVIII-075.21. \nolinkurl{fr/sga7ii/diplomatic/expose_XVIII_body.tex}; \nolinkurl{fr/sga7ii/current/expose_XVIII_body.tex}\par

\par\medskip\noindent\hypertarget{app-F7II-XVIII-075.22}{}\textbf{F7II-XVIII-075.22}\quad p. 303; printed misprint.\par
\noindent Imprimé : {\ttfamily\footnotesize \textbackslash{}alpha\allowbreak{}}\par
\noindent Corrigé : {\ttfamily\footnotesize x\allowbreak{}}\par
\noindent Source : 7II:519; 7II:964. \nolinkurl{fixes/MISPRINTS_7II.json#519}; \nolinkurl{fixes/MISPRINTS_7II.json#964}\par

\par\medskip\noindent\hypertarget{app-F7II-XVIII-075.23}{}\textbf{F7II-XVIII-075.23}\quad p. 303; designated-witness attribution.\par
\noindent Imprimé : {\ttfamily\footnotesize rel.\allowbreak{}}\par
\noindent Corrigé : {\ttfamily\footnotesize relativement\allowbreak{}}\par
\noindent Source : WITNESS-F7II-XVIII-075.23. \nolinkurl{fr/sga7ii/diplomatic/expose_XVIII_body.tex}; \nolinkurl{fr/sga7ii/current/expose_XVIII_body.tex}\par

\par\medskip\noindent\hypertarget{app-F7II-XVIII-075.24}{}\textbf{F7II-XVIII-075.24}\quad p. 303; designated-witness attribution.\par
\noindent Imprimé : \emph{(vide)}\par
\noindent Corrigé : {\ttfamily\footnotesize \$\allowbreak{}}\par
\noindent Source : WITNESS-F7II-XVIII-075.24. \nolinkurl{fr/sga7ii/diplomatic/expose_XVIII_body.tex}; \nolinkurl{fr/sga7ii/current/expose_XVIII_body.tex}\par

\par\medskip\noindent\hypertarget{app-F7II-XVIII-075.25}{}\textbf{F7II-XVIII-075.25}\quad p. 303; designated-witness attribution.\par
\noindent Imprimé : \emph{(1 espace(s), 0 retour(s) de ligne dans le témoin TeX)}\par
\noindent Corrigé : {\ttfamily\footnotesize \$\allowbreak{}}\par
\noindent Source : WITNESS-F7II-XVIII-075.25. \nolinkurl{fr/sga7ii/diplomatic/expose_XVIII_body.tex}; \nolinkurl{fr/sga7ii/current/expose_XVIII_body.tex}\par

\par\medskip\noindent\hypertarget{app-F7II-XVIII-076.1}{}\textbf{F7II-XVIII-076.1}\quad p. 303; inherited corrected-layer difference.\par
\noindent Imprimé : {\ttfamily\footnotesize \textbackslash{}mathrm\{\allowbreak{}}\par
\noindent Corrigé : \emph{(vide)}\par
\noindent Source : ADOPT-F7II-XVIII-076.1. \nolinkurl{fr/sga7ii/diplomatic/expose_XVIII_body.tex}; \nolinkurl{fr/sga7ii/current/expose_XVIII_body.tex}\par

\par\medskip\noindent\hypertarget{app-F7II-XVIII-076.2}{}\textbf{F7II-XVIII-076.2}\quad p. 303; inherited corrected-layer difference.\par
\noindent Imprimé : {\ttfamily\footnotesize \}\allowbreak{}}\par
\noindent Corrigé : \emph{(vide)}\par
\noindent Source : ADOPT-F7II-XVIII-076.2. \nolinkurl{fr/sga7ii/diplomatic/expose_XVIII_body.tex}; \nolinkurl{fr/sga7ii/current/expose_XVIII_body.tex}\par

\par\medskip\noindent\hypertarget{app-F7II-XVIII-076.3}{}\textbf{F7II-XVIII-076.3}\quad p. 303; inherited corrected-layer difference.\par
\noindent Imprimé : {\ttfamily\footnotesize \textbackslash{},?\allowbreak{}}\par
\noindent Corrigé : \emph{(vide)}\par
\noindent Source : ADOPT-F7II-XVIII-076.3. \nolinkurl{fr/sga7ii/diplomatic/expose_XVIII_body.tex}; \nolinkurl{fr/sga7ii/current/expose_XVIII_body.tex}\par

\par\medskip\noindent\hypertarget{app-F7II-XVIII-077.1}{}\textbf{F7II-XVIII-077.1}\quad p. 304; designated-witness attribution.\par
\noindent Imprimé : \emph{(vide)}\par
\noindent Corrigé : {\ttfamily\footnotesize \textbackslash{}nu\textasciicircum{}*\allowbreak{}}\par
\noindent Source : WITNESS-F7II-XVIII-077.1. \nolinkurl{fr/sga7ii/diplomatic/expose_XVIII_body.tex}; \nolinkurl{fr/sga7ii/current/expose_XVIII_body.tex}\par

\par\medskip\noindent\hypertarget{app-F7II-XVIII-077.2}{}\textbf{F7II-XVIII-077.2}\quad p. 304; printed misprint.\par
\noindent Imprimé : {\ttfamily\footnotesize (\allowbreak{}}\par
\noindent Corrigé : \emph{(vide)}\par
\noindent Source : 7II:520; 7II:965. \nolinkurl{fixes/MISPRINTS_7II.json#520}; \nolinkurl{fixes/MISPRINTS_7II.json#965}\par

\par\medskip\noindent\hypertarget{app-F7II-XVIII-078.1}{}\textbf{F7II-XVIII-078.1}\quad p. 305; printed misprint.\par
\noindent Imprimé : {\ttfamily\footnotesize )\allowbreak{}}\par
\noindent Corrigé : \emph{(vide)}\par
\noindent Source : 7II:521; 7II:966. \nolinkurl{fixes/MISPRINTS_7II.json#521}; \nolinkurl{fixes/MISPRINTS_7II.json#966}\par

\par\medskip\noindent\hypertarget{app-F7II-XVIII-079.1}{}\textbf{F7II-XVIII-079.1}\quad p. 305; printed misprint.\par
\noindent Imprimé : {\ttfamily\footnotesize défini\allowbreak{}}\par
\noindent Corrigé : {\ttfamily\footnotesize définie\allowbreak{}}\par
\noindent Source : 7II:522; 7II:967. \nolinkurl{fixes/MISPRINTS_7II.json#522}; \nolinkurl{fixes/MISPRINTS_7II.json#967}\par

\par\medskip\noindent\hypertarget{app-F7II-XVIII-079.2}{}\textbf{F7II-XVIII-079.2}\quad p. 305; inherited corrected-layer difference.\par
\noindent Imprimé : \emph{(vide)}\par
\noindent Corrigé : {\ttfamily\footnotesize (j)\allowbreak{}}\par
\noindent Source : ADOPT-F7II-XVIII-079.2. \nolinkurl{fr/sga7ii/diplomatic/expose_XVIII_body.tex}; \nolinkurl{fr/sga7ii/current/expose_XVIII_body.tex}\par

\par\medskip\noindent\hypertarget{app-F7II-XVIII-079.3}{}\textbf{F7II-XVIII-079.3}\quad p. 305; inherited corrected-layer difference.\par
\noindent Imprimé : \emph{(vide)}\par
\noindent Corrigé : {\ttfamily\footnotesize (j)\allowbreak{}}\par
\noindent Source : ADOPT-F7II-XVIII-079.3. \nolinkurl{fr/sga7ii/diplomatic/expose_XVIII_body.tex}; \nolinkurl{fr/sga7ii/current/expose_XVIII_body.tex}\par

\par\medskip\noindent\hypertarget{app-F7II-XVIII-080.1}{}\textbf{F7II-XVIII-080.1}\quad p. 306; printed misprint.\par
\noindent Imprimé : \emph{(vide)}\par
\noindent Corrigé : \emph{(1 espace(s), 0 retour(s) de ligne dans le témoin TeX)}\par
\noindent Source : 7II:949. \nolinkurl{fixes/MISPRINTS_7II.json#949}\par

\par\medskip\noindent\hypertarget{app-F7II-XVIII-082.1}{}\textbf{F7II-XVIII-082.1}\quad p. 308; inherited corrected-layer difference.\par
\noindent Imprimé : {\ttfamily\footnotesize (\allowbreak{}}\par
\noindent Corrigé : \emph{(vide)}\par
\noindent Source : ADOPT-F7II-XVIII-082.1. \nolinkurl{fr/sga7ii/diplomatic/expose_XVIII_body.tex}; \nolinkurl{fr/sga7ii/current/expose_XVIII_body.tex}\par

\par\medskip\noindent\hypertarget{app-F7II-XVIII-082.2}{}\textbf{F7II-XVIII-082.2}\quad p. 308; inherited corrected-layer difference.\par
\noindent Imprimé : {\ttfamily\footnotesize \{\textbackslash{}mathrm\{\allowbreak{}}\par
\noindent Corrigé : \emph{(vide)}\par
\noindent Source : ADOPT-F7II-XVIII-082.2. \nolinkurl{fr/sga7ii/diplomatic/expose_XVIII_body.tex}; \nolinkurl{fr/sga7ii/current/expose_XVIII_body.tex}\par

\par\medskip\noindent\hypertarget{app-F7II-XVIII-082.3}{}\textbf{F7II-XVIII-082.3}\quad p. 308; inherited corrected-layer difference.\par
\noindent Imprimé : {\ttfamily\footnotesize \}\}\allowbreak{}}\par
\noindent Corrigé : {\ttfamily\footnotesize \textasciicircum{}*\allowbreak{}}\par
\noindent Source : ADOPT-F7II-XVIII-082.3. \nolinkurl{fr/sga7ii/diplomatic/expose_XVIII_body.tex}; \nolinkurl{fr/sga7ii/current/expose_XVIII_body.tex}\par

\par\medskip\noindent\hypertarget{app-F7II-XVIII-082.4}{}\textbf{F7II-XVIII-082.4}\quad p. 308; inherited corrected-layer difference.\par
\noindent Imprimé : {\ttfamily\footnotesize )\allowbreak{}}\par
\noindent Corrigé : \emph{(vide)}\par
\noindent Source : ADOPT-F7II-XVIII-082.4. \nolinkurl{fr/sga7ii/diplomatic/expose_XVIII_body.tex}; \nolinkurl{fr/sga7ii/current/expose_XVIII_body.tex}\par

\par\medskip\noindent\hypertarget{app-F7II-XVIII-082.5}{}\textbf{F7II-XVIII-082.5}\quad p. 308; inherited corrected-layer difference.\par
\noindent Imprimé : {\ttfamily\footnotesize \textbackslash{}mathrm\{\allowbreak{}}\par
\noindent Corrigé : \emph{(vide)}\par
\noindent Source : ADOPT-F7II-XVIII-082.5. \nolinkurl{fr/sga7ii/diplomatic/expose_XVIII_body.tex}; \nolinkurl{fr/sga7ii/current/expose_XVIII_body.tex}\par

\par\medskip\noindent\hypertarget{app-F7II-XVIII-082.6}{}\textbf{F7II-XVIII-082.6}\quad p. 308; inherited corrected-layer difference.\par
\noindent Imprimé : {\ttfamily\footnotesize \}\allowbreak{}}\par
\noindent Corrigé : \emph{(vide)}\par
\noindent Source : ADOPT-F7II-XVIII-082.6. \nolinkurl{fr/sga7ii/diplomatic/expose_XVIII_body.tex}; \nolinkurl{fr/sga7ii/current/expose_XVIII_body.tex}\par

\par\medskip\noindent\hypertarget{app-F7II-XVIII-082.7}{}\textbf{F7II-XVIII-082.7}\quad p. 308; inherited corrected-layer difference.\par
\noindent Imprimé : {\ttfamily\footnotesize \{\textbackslash{}mathrm\{\allowbreak{}}\par
\noindent Corrigé : \emph{(1 espace(s), 0 retour(s) de ligne dans le témoin TeX)}\par
\noindent Source : ADOPT-F7II-XVIII-082.7. \nolinkurl{fr/sga7ii/diplomatic/expose_XVIII_body.tex}; \nolinkurl{fr/sga7ii/current/expose_XVIII_body.tex}\par

\par\medskip\noindent\hypertarget{app-F7II-XVIII-082.8}{}\textbf{F7II-XVIII-082.8}\quad p. 308; inherited corrected-layer difference.\par
\noindent Imprimé : {\ttfamily\footnotesize \}\}\allowbreak{}}\par
\noindent Corrigé : \emph{(vide)}\par
\noindent Source : ADOPT-F7II-XVIII-082.8. \nolinkurl{fr/sga7ii/diplomatic/expose_XVIII_body.tex}; \nolinkurl{fr/sga7ii/current/expose_XVIII_body.tex}\par

\par\medskip\noindent\hypertarget{app-F7II-XVIII-082.9}{}\textbf{F7II-XVIII-082.9}\quad p. 308; inherited corrected-layer difference.\par
\noindent Imprimé : {\ttfamily\footnotesize \{\textbackslash{}mathrm\{\allowbreak{}}\par
\noindent Corrigé : \emph{(vide)}\par
\noindent Source : ADOPT-F7II-XVIII-082.9. \nolinkurl{fr/sga7ii/diplomatic/expose_XVIII_body.tex}; \nolinkurl{fr/sga7ii/current/expose_XVIII_body.tex}\par

\par\medskip\noindent\hypertarget{app-F7II-XVIII-082.10}{}\textbf{F7II-XVIII-082.10}\quad p. 308; inherited corrected-layer difference.\par
\noindent Imprimé : {\ttfamily\footnotesize \}\}\allowbreak{}}\par
\noindent Corrigé : \emph{(vide)}\par
\noindent Source : ADOPT-F7II-XVIII-082.10. \nolinkurl{fr/sga7ii/diplomatic/expose_XVIII_body.tex}; \nolinkurl{fr/sga7ii/current/expose_XVIII_body.tex}\par

\par\medskip\noindent\hypertarget{app-F7II-XVIII-082.11}{}\textbf{F7II-XVIII-082.11}\quad p. 308; inherited corrected-layer difference.\par
\noindent Imprimé : {\ttfamily\footnotesize \_\{\textbackslash{}mathrm\{\allowbreak{}}\par
\noindent Corrigé : \emph{(vide)}\par
\noindent Source : ADOPT-F7II-XVIII-082.11. \nolinkurl{fr/sga7ii/diplomatic/expose_XVIII_body.tex}; \nolinkurl{fr/sga7ii/current/expose_XVIII_body.tex}\par

\par\medskip\noindent\hypertarget{app-F7II-XVIII-082.12}{}\textbf{F7II-XVIII-082.12}\quad p. 308; inherited corrected-layer difference.\par
\noindent Imprimé : {\ttfamily\footnotesize \}\}\allowbreak{} }\par
\noindent Corrigé : \emph{(vide)}\par
\noindent Source : ADOPT-F7II-XVIII-082.12. \nolinkurl{fr/sga7ii/diplomatic/expose_XVIII_body.tex}; \nolinkurl{fr/sga7ii/current/expose_XVIII_body.tex}\par

\par\medskip\noindent\hypertarget{app-F7II-XVIII-083.1}{}\textbf{F7II-XVIII-083.1}\quad p. 308; inherited corrected-layer difference.\par
\noindent Imprimé : {\ttfamily\footnotesize \textbackslash{}mathrm\{\allowbreak{}}\par
\noindent Corrigé : \emph{(vide)}\par
\noindent Source : ADOPT-F7II-XVIII-083.1. \nolinkurl{fr/sga7ii/diplomatic/expose_XVIII_body.tex}; \nolinkurl{fr/sga7ii/current/expose_XVIII_body.tex}\par

\par\medskip\noindent\hypertarget{app-F7II-XVIII-083.2}{}\textbf{F7II-XVIII-083.2}\quad p. 308; inherited corrected-layer difference.\par
\noindent Imprimé : {\ttfamily\footnotesize \}\allowbreak{}}\par
\noindent Corrigé : \emph{(vide)}\par
\noindent Source : ADOPT-F7II-XVIII-083.2. \nolinkurl{fr/sga7ii/diplomatic/expose_XVIII_body.tex}; \nolinkurl{fr/sga7ii/current/expose_XVIII_body.tex}\par

\par\medskip\noindent\hypertarget{app-F7II-XVIII-083.3}{}\textbf{F7II-XVIII-083.3}\quad p. 308; inherited corrected-layer difference.\par
\noindent Imprimé : {\ttfamily\footnotesize \textbackslash{}mathrm\{\allowbreak{}}\par
\noindent Corrigé : \emph{(vide)}\par
\noindent Source : ADOPT-F7II-XVIII-083.3. \nolinkurl{fr/sga7ii/diplomatic/expose_XVIII_body.tex}; \nolinkurl{fr/sga7ii/current/expose_XVIII_body.tex}\par

\par\medskip\noindent\hypertarget{app-F7II-XVIII-083.4}{}\textbf{F7II-XVIII-083.4}\quad p. 308; inherited corrected-layer difference.\par
\noindent Imprimé : {\ttfamily\footnotesize \}\allowbreak{}}\par
\noindent Corrigé : \emph{(vide)}\par
\noindent Source : ADOPT-F7II-XVIII-083.4. \nolinkurl{fr/sga7ii/diplomatic/expose_XVIII_body.tex}; \nolinkurl{fr/sga7ii/current/expose_XVIII_body.tex}\par

\par\medskip\noindent\hypertarget{app-F7II-XVIII-083.5}{}\textbf{F7II-XVIII-083.5}\quad p. 308; inherited corrected-layer difference.\par
\noindent Imprimé : {\ttfamily\footnotesize \textbackslash{}mathrm\{\allowbreak{}}\par
\noindent Corrigé : \emph{(vide)}\par
\noindent Source : ADOPT-F7II-XVIII-083.5. \nolinkurl{fr/sga7ii/diplomatic/expose_XVIII_body.tex}; \nolinkurl{fr/sga7ii/current/expose_XVIII_body.tex}\par

\par\medskip\noindent\hypertarget{app-F7II-XVIII-083.6}{}\textbf{F7II-XVIII-083.6}\quad p. 308; inherited corrected-layer difference.\par
\noindent Imprimé : {\ttfamily\footnotesize \}\allowbreak{}}\par
\noindent Corrigé : \emph{(vide)}\par
\noindent Source : ADOPT-F7II-XVIII-083.6. \nolinkurl{fr/sga7ii/diplomatic/expose_XVIII_body.tex}; \nolinkurl{fr/sga7ii/current/expose_XVIII_body.tex}\par

\par\medskip\noindent\hypertarget{app-F7II-XVIII-083.7}{}\textbf{F7II-XVIII-083.7}\quad p. 308; inherited corrected-layer difference.\par
\noindent Imprimé : {\ttfamily\footnotesize \textbackslash{}mathrm\{\allowbreak{}}\par
\noindent Corrigé : \emph{(vide)}\par
\noindent Source : ADOPT-F7II-XVIII-083.7. \nolinkurl{fr/sga7ii/diplomatic/expose_XVIII_body.tex}; \nolinkurl{fr/sga7ii/current/expose_XVIII_body.tex}\par

\par\medskip\noindent\hypertarget{app-F7II-XVIII-083.8}{}\textbf{F7II-XVIII-083.8}\quad p. 308; inherited corrected-layer difference.\par
\noindent Imprimé : {\ttfamily\footnotesize \}\allowbreak{}}\par
\noindent Corrigé : \emph{(vide)}\par
\noindent Source : ADOPT-F7II-XVIII-083.8. \nolinkurl{fr/sga7ii/diplomatic/expose_XVIII_body.tex}; \nolinkurl{fr/sga7ii/current/expose_XVIII_body.tex}\par

\par\medskip\noindent\hypertarget{app-F7II-XVIII-083.9}{}\textbf{F7II-XVIII-083.9}\quad p. 308; inherited corrected-layer difference.\par
\noindent Imprimé : {\ttfamily\footnotesize \textbackslash{}mathrm\{\allowbreak{}}\par
\noindent Corrigé : \emph{(vide)}\par
\noindent Source : ADOPT-F7II-XVIII-083.9. \nolinkurl{fr/sga7ii/diplomatic/expose_XVIII_body.tex}; \nolinkurl{fr/sga7ii/current/expose_XVIII_body.tex}\par

\par\medskip\noindent\hypertarget{app-F7II-XVIII-083.10}{}\textbf{F7II-XVIII-083.10}\quad p. 308; inherited corrected-layer difference.\par
\noindent Imprimé : {\ttfamily\footnotesize \}\_\allowbreak{}}\par
\noindent Corrigé : \emph{(vide)}\par
\noindent Source : ADOPT-F7II-XVIII-083.10. \nolinkurl{fr/sga7ii/diplomatic/expose_XVIII_body.tex}; \nolinkurl{fr/sga7ii/current/expose_XVIII_body.tex}\par

\par\medskip\noindent\hypertarget{app-F7II-XVIII-084.1}{}\textbf{F7II-XVIII-084.1}\quad p. 308; inherited corrected-layer difference.\par
\noindent Imprimé : {\ttfamily\footnotesize \textbackslash{}mathrm\{\allowbreak{}}\par
\noindent Corrigé : \emph{(vide)}\par
\noindent Source : ADOPT-F7II-XVIII-084.1. \nolinkurl{fr/sga7ii/diplomatic/expose_XVIII_body.tex}; \nolinkurl{fr/sga7ii/current/expose_XVIII_body.tex}\par

\par\medskip\noindent\hypertarget{app-F7II-XVIII-084.2}{}\textbf{F7II-XVIII-084.2}\quad p. 308; inherited corrected-layer difference.\par
\noindent Imprimé : {\ttfamily\footnotesize \}\allowbreak{}}\par
\noindent Corrigé : \emph{(vide)}\par
\noindent Source : ADOPT-F7II-XVIII-084.2. \nolinkurl{fr/sga7ii/diplomatic/expose_XVIII_body.tex}; \nolinkurl{fr/sga7ii/current/expose_XVIII_body.tex}\par

\par\medskip\noindent\hypertarget{app-F7II-XVIII-084.3}{}\textbf{F7II-XVIII-084.3}\quad p. 308; inherited corrected-layer difference.\par
\noindent Imprimé : {\ttfamily\footnotesize \textbackslash{}mathrm\{\allowbreak{}}\par
\noindent Corrigé : \emph{(vide)}\par
\noindent Source : ADOPT-F7II-XVIII-084.3. \nolinkurl{fr/sga7ii/diplomatic/expose_XVIII_body.tex}; \nolinkurl{fr/sga7ii/current/expose_XVIII_body.tex}\par

\par\medskip\noindent\hypertarget{app-F7II-XVIII-084.4}{}\textbf{F7II-XVIII-084.4}\quad p. 308; inherited corrected-layer difference.\par
\noindent Imprimé : {\ttfamily\footnotesize \}\allowbreak{}}\par
\noindent Corrigé : \emph{(vide)}\par
\noindent Source : ADOPT-F7II-XVIII-084.4. \nolinkurl{fr/sga7ii/diplomatic/expose_XVIII_body.tex}; \nolinkurl{fr/sga7ii/current/expose_XVIII_body.tex}\par

\par\medskip\noindent\hypertarget{app-F7II-XVIII-084.5}{}\textbf{F7II-XVIII-084.5}\quad p. 308; inherited corrected-layer difference.\par
\noindent Imprimé : {\ttfamily\footnotesize \textbackslash{}mathrm\{\allowbreak{}}\par
\noindent Corrigé : \emph{(vide)}\par
\noindent Source : ADOPT-F7II-XVIII-084.5. \nolinkurl{fr/sga7ii/diplomatic/expose_XVIII_body.tex}; \nolinkurl{fr/sga7ii/current/expose_XVIII_body.tex}\par

\par\medskip\noindent\hypertarget{app-F7II-XVIII-084.6}{}\textbf{F7II-XVIII-084.6}\quad p. 308; inherited corrected-layer difference.\par
\noindent Imprimé : {\ttfamily\footnotesize \}\allowbreak{}}\par
\noindent Corrigé : \emph{(vide)}\par
\noindent Source : ADOPT-F7II-XVIII-084.6. \nolinkurl{fr/sga7ii/diplomatic/expose_XVIII_body.tex}; \nolinkurl{fr/sga7ii/current/expose_XVIII_body.tex}\par

\par\medskip\noindent\hypertarget{app-F7II-XVIII-084.7}{}\textbf{F7II-XVIII-084.7}\quad p. 308; inherited corrected-layer difference.\par
\noindent Imprimé : {\ttfamily\footnotesize \textbackslash{}mathrm\{\allowbreak{}}\par
\noindent Corrigé : \emph{(vide)}\par
\noindent Source : ADOPT-F7II-XVIII-084.7. \nolinkurl{fr/sga7ii/diplomatic/expose_XVIII_body.tex}; \nolinkurl{fr/sga7ii/current/expose_XVIII_body.tex}\par

\par\medskip\noindent\hypertarget{app-F7II-XVIII-084.8}{}\textbf{F7II-XVIII-084.8}\quad p. 308; inherited corrected-layer difference.\par
\noindent Imprimé : {\ttfamily\footnotesize \}\allowbreak{}}\par
\noindent Corrigé : \emph{(vide)}\par
\noindent Source : ADOPT-F7II-XVIII-084.8. \nolinkurl{fr/sga7ii/diplomatic/expose_XVIII_body.tex}; \nolinkurl{fr/sga7ii/current/expose_XVIII_body.tex}\par

\par\medskip\noindent\hypertarget{app-F7II-XVIII-084.9}{}\textbf{F7II-XVIII-084.9}\quad p. 308; inherited corrected-layer difference.\par
\noindent Imprimé : {\ttfamily\footnotesize \textbackslash{}mathrm\{\allowbreak{}}\par
\noindent Corrigé : \emph{(vide)}\par
\noindent Source : ADOPT-F7II-XVIII-084.9. \nolinkurl{fr/sga7ii/diplomatic/expose_XVIII_body.tex}; \nolinkurl{fr/sga7ii/current/expose_XVIII_body.tex}\par

\par\medskip\noindent\hypertarget{app-F7II-XVIII-084.10}{}\textbf{F7II-XVIII-084.10}\quad p. 308; inherited corrected-layer difference.\par
\noindent Imprimé : {\ttfamily\footnotesize \}\_\allowbreak{}}\par
\noindent Corrigé : \emph{(vide)}\par
\noindent Source : ADOPT-F7II-XVIII-084.10. \nolinkurl{fr/sga7ii/diplomatic/expose_XVIII_body.tex}; \nolinkurl{fr/sga7ii/current/expose_XVIII_body.tex}\par

\par\medskip\noindent\hypertarget{app-F7II-XVIII-084.11}{}\textbf{F7II-XVIII-084.11}\quad p. 308; inherited corrected-layer difference.\par
\noindent Imprimé : {\ttfamily\footnotesize \textbackslash{}mathrm\{\allowbreak{}}\par
\noindent Corrigé : \emph{(vide)}\par
\noindent Source : ADOPT-F7II-XVIII-084.11. \nolinkurl{fr/sga7ii/diplomatic/expose_XVIII_body.tex}; \nolinkurl{fr/sga7ii/current/expose_XVIII_body.tex}\par

\par\medskip\noindent\hypertarget{app-F7II-XVIII-084.12}{}\textbf{F7II-XVIII-084.12}\quad p. 308; inherited corrected-layer difference.\par
\noindent Imprimé : {\ttfamily\footnotesize \}\allowbreak{}}\par
\noindent Corrigé : \emph{(vide)}\par
\noindent Source : ADOPT-F7II-XVIII-084.12. \nolinkurl{fr/sga7ii/diplomatic/expose_XVIII_body.tex}; \nolinkurl{fr/sga7ii/current/expose_XVIII_body.tex}\par

\par\medskip\noindent\hypertarget{app-F7II-XVIII-084.13}{}\textbf{F7II-XVIII-084.13}\quad p. 308; inherited corrected-layer difference.\par
\noindent Imprimé : {\ttfamily\footnotesize \textbackslash{}mathrm\{\allowbreak{}}\par
\noindent Corrigé : \emph{(vide)}\par
\noindent Source : ADOPT-F7II-XVIII-084.13. \nolinkurl{fr/sga7ii/diplomatic/expose_XVIII_body.tex}; \nolinkurl{fr/sga7ii/current/expose_XVIII_body.tex}\par

\par\medskip\noindent\hypertarget{app-F7II-XVIII-084.14}{}\textbf{F7II-XVIII-084.14}\quad p. 308; inherited corrected-layer difference.\par
\noindent Imprimé : {\ttfamily\footnotesize \}\allowbreak{}}\par
\noindent Corrigé : \emph{(vide)}\par
\noindent Source : ADOPT-F7II-XVIII-084.14. \nolinkurl{fr/sga7ii/diplomatic/expose_XVIII_body.tex}; \nolinkurl{fr/sga7ii/current/expose_XVIII_body.tex}\par

\par\medskip\noindent\hypertarget{app-F7II-XVIII-084.15}{}\textbf{F7II-XVIII-084.15}\quad p. 308; inherited corrected-layer difference.\par
\noindent Imprimé : {\ttfamily\footnotesize \textbackslash{}\allowbreak{} }\par
\noindent Corrigé : \emph{(vide)}\par
\noindent Source : ADOPT-F7II-XVIII-084.15. \nolinkurl{fr/sga7ii/diplomatic/expose_XVIII_body.tex}; \nolinkurl{fr/sga7ii/current/expose_XVIII_body.tex}\par

\par\medskip\noindent\hypertarget{app-F7II-XVIII-084.16}{}\textbf{F7II-XVIII-084.16}\quad p. 308; inherited corrected-layer difference.\par
\noindent Imprimé : {\ttfamily\footnotesize \textbackslash{}mathrm\{\allowbreak{}}\par
\noindent Corrigé : \emph{(vide)}\par
\noindent Source : ADOPT-F7II-XVIII-084.16. \nolinkurl{fr/sga7ii/diplomatic/expose_XVIII_body.tex}; \nolinkurl{fr/sga7ii/current/expose_XVIII_body.tex}\par

\par\medskip\noindent\hypertarget{app-F7II-XVIII-084.17}{}\textbf{F7II-XVIII-084.17}\quad p. 308; inherited corrected-layer difference.\par
\noindent Imprimé : {\ttfamily\footnotesize \}\allowbreak{}}\par
\noindent Corrigé : \emph{(vide)}\par
\noindent Source : ADOPT-F7II-XVIII-084.17. \nolinkurl{fr/sga7ii/diplomatic/expose_XVIII_body.tex}; \nolinkurl{fr/sga7ii/current/expose_XVIII_body.tex}\par

\par\medskip\noindent\hypertarget{app-F7II-XVIII-084.18}{}\textbf{F7II-XVIII-084.18}\quad p. 308; inherited corrected-layer difference.\par
\noindent Imprimé : {\ttfamily\footnotesize (\allowbreak{}}\par
\noindent Corrigé : \emph{(vide)}\par
\noindent Source : ADOPT-F7II-XVIII-084.18. \nolinkurl{fr/sga7ii/diplomatic/expose_XVIII_body.tex}; \nolinkurl{fr/sga7ii/current/expose_XVIII_body.tex}\par

\par\medskip\noindent\hypertarget{app-F7II-XVIII-084.19}{}\textbf{F7II-XVIII-084.19}\quad p. 308; inherited corrected-layer difference.\par
\noindent Imprimé : {\ttfamily\footnotesize )\allowbreak{}}\par
\noindent Corrigé : \emph{(vide)}\par
\noindent Source : ADOPT-F7II-XVIII-084.19. \nolinkurl{fr/sga7ii/diplomatic/expose_XVIII_body.tex}; \nolinkurl{fr/sga7ii/current/expose_XVIII_body.tex}\par

\par\medskip\noindent\hypertarget{app-F7II-XVIII-084.20}{}\textbf{F7II-XVIII-084.20}\quad p. 308; inherited corrected-layer difference.\par
\noindent Imprimé : {\ttfamily\footnotesize \textbackslash{}mathrm\{\allowbreak{}}\par
\noindent Corrigé : \emph{(vide)}\par
\noindent Source : ADOPT-F7II-XVIII-084.20. \nolinkurl{fr/sga7ii/diplomatic/expose_XVIII_body.tex}; \nolinkurl{fr/sga7ii/current/expose_XVIII_body.tex}\par

\par\medskip\noindent\hypertarget{app-F7II-XVIII-084.21}{}\textbf{F7II-XVIII-084.21}\quad p. 308; inherited corrected-layer difference.\par
\noindent Imprimé : {\ttfamily\footnotesize \}\allowbreak{}}\par
\noindent Corrigé : \emph{(vide)}\par
\noindent Source : ADOPT-F7II-XVIII-084.21. \nolinkurl{fr/sga7ii/diplomatic/expose_XVIII_body.tex}; \nolinkurl{fr/sga7ii/current/expose_XVIII_body.tex}\par

\par\medskip\noindent\hypertarget{app-F7II-XVIII-084.22}{}\textbf{F7II-XVIII-084.22}\quad p. 308; inherited corrected-layer difference.\par
\noindent Imprimé : {\ttfamily\footnotesize \textbackslash{}mathrm\{\allowbreak{}}\par
\noindent Corrigé : \emph{(vide)}\par
\noindent Source : ADOPT-F7II-XVIII-084.22. \nolinkurl{fr/sga7ii/diplomatic/expose_XVIII_body.tex}; \nolinkurl{fr/sga7ii/current/expose_XVIII_body.tex}\par

\par\medskip\noindent\hypertarget{app-F7II-XVIII-084.23}{}\textbf{F7II-XVIII-084.23}\quad p. 308; inherited corrected-layer difference.\par
\noindent Imprimé : {\ttfamily\footnotesize \}\allowbreak{}}\par
\noindent Corrigé : \emph{(vide)}\par
\noindent Source : ADOPT-F7II-XVIII-084.23. \nolinkurl{fr/sga7ii/diplomatic/expose_XVIII_body.tex}; \nolinkurl{fr/sga7ii/current/expose_XVIII_body.tex}\par

\par\medskip\noindent\hypertarget{app-F7II-XVIII-084.24}{}\textbf{F7II-XVIII-084.24}\quad p. 308; inherited corrected-layer difference.\par
\noindent Imprimé : {\ttfamily\footnotesize \textbackslash{}mathrm\{\allowbreak{}}\par
\noindent Corrigé : \emph{(vide)}\par
\noindent Source : ADOPT-F7II-XVIII-084.24. \nolinkurl{fr/sga7ii/diplomatic/expose_XVIII_body.tex}; \nolinkurl{fr/sga7ii/current/expose_XVIII_body.tex}\par

\par\medskip\noindent\hypertarget{app-F7II-XVIII-084.25}{}\textbf{F7II-XVIII-084.25}\quad p. 308; inherited corrected-layer difference.\par
\noindent Imprimé : {\ttfamily\footnotesize \}\allowbreak{}}\par
\noindent Corrigé : \emph{(vide)}\par
\noindent Source : ADOPT-F7II-XVIII-084.25. \nolinkurl{fr/sga7ii/diplomatic/expose_XVIII_body.tex}; \nolinkurl{fr/sga7ii/current/expose_XVIII_body.tex}\par

\par\medskip\noindent\hypertarget{app-F7II-XVIII-084.26}{}\textbf{F7II-XVIII-084.26}\quad p. 308; inherited corrected-layer difference.\par
\noindent Imprimé : {\ttfamily\footnotesize \textbackslash{}mathrm\{\allowbreak{}}\par
\noindent Corrigé : \emph{(vide)}\par
\noindent Source : ADOPT-F7II-XVIII-084.26. \nolinkurl{fr/sga7ii/diplomatic/expose_XVIII_body.tex}; \nolinkurl{fr/sga7ii/current/expose_XVIII_body.tex}\par

\par\medskip\noindent\hypertarget{app-F7II-XVIII-084.27}{}\textbf{F7II-XVIII-084.27}\quad p. 308; inherited corrected-layer difference.\par
\noindent Imprimé : {\ttfamily\footnotesize \}\_\allowbreak{}}\par
\noindent Corrigé : \emph{(vide)}\par
\noindent Source : ADOPT-F7II-XVIII-084.27. \nolinkurl{fr/sga7ii/diplomatic/expose_XVIII_body.tex}; \nolinkurl{fr/sga7ii/current/expose_XVIII_body.tex}\par

\par\medskip\noindent\hypertarget{app-F7II-XVIII-084.28}{}\textbf{F7II-XVIII-084.28}\quad p. 308; inherited corrected-layer difference.\par
\noindent Imprimé : {\ttfamily\footnotesize \{\textbackslash{}mathrm\{\allowbreak{}}\par
\noindent Corrigé : \emph{(vide)}\par
\noindent Source : ADOPT-F7II-XVIII-084.28. \nolinkurl{fr/sga7ii/diplomatic/expose_XVIII_body.tex}; \nolinkurl{fr/sga7ii/current/expose_XVIII_body.tex}\par

\par\medskip\noindent\hypertarget{app-F7II-XVIII-084.29}{}\textbf{F7II-XVIII-084.29}\quad p. 308; inherited corrected-layer difference.\par
\noindent Imprimé : {\ttfamily\footnotesize \}\}\allowbreak{}}\par
\noindent Corrigé : \emph{(vide)}\par
\noindent Source : ADOPT-F7II-XVIII-084.29. \nolinkurl{fr/sga7ii/diplomatic/expose_XVIII_body.tex}; \nolinkurl{fr/sga7ii/current/expose_XVIII_body.tex}\par

\par\medskip\noindent\hypertarget{app-F7II-XVIII-084.30}{}\textbf{F7II-XVIII-084.30}\quad p. 308; inherited corrected-layer difference.\par
\noindent Imprimé : {\ttfamily\footnotesize \textbackslash{}mathrm\{\allowbreak{}}\par
\noindent Corrigé : \emph{(vide)}\par
\noindent Source : ADOPT-F7II-XVIII-084.30. \nolinkurl{fr/sga7ii/diplomatic/expose_XVIII_body.tex}; \nolinkurl{fr/sga7ii/current/expose_XVIII_body.tex}\par

\par\medskip\noindent\hypertarget{app-F7II-XVIII-084.31}{}\textbf{F7II-XVIII-084.31}\quad p. 308; inherited corrected-layer difference.\par
\noindent Imprimé : {\ttfamily\footnotesize \}\allowbreak{}}\par
\noindent Corrigé : \emph{(vide)}\par
\noindent Source : ADOPT-F7II-XVIII-084.31. \nolinkurl{fr/sga7ii/diplomatic/expose_XVIII_body.tex}; \nolinkurl{fr/sga7ii/current/expose_XVIII_body.tex}\par

\par\medskip\noindent\hypertarget{app-F7II-XVIII-084.32}{}\textbf{F7II-XVIII-084.32}\quad p. 308; inherited corrected-layer difference.\par
\noindent Imprimé : {\ttfamily\footnotesize \textbackslash{}mathrm\{\allowbreak{}}\par
\noindent Corrigé : \emph{(vide)}\par
\noindent Source : ADOPT-F7II-XVIII-084.32. \nolinkurl{fr/sga7ii/diplomatic/expose_XVIII_body.tex}; \nolinkurl{fr/sga7ii/current/expose_XVIII_body.tex}\par

\par\medskip\noindent\hypertarget{app-F7II-XVIII-084.33}{}\textbf{F7II-XVIII-084.33}\quad p. 308; inherited corrected-layer difference.\par
\noindent Imprimé : {\ttfamily\footnotesize \}\allowbreak{}}\par
\noindent Corrigé : \emph{(vide)}\par
\noindent Source : ADOPT-F7II-XVIII-084.33. \nolinkurl{fr/sga7ii/diplomatic/expose_XVIII_body.tex}; \nolinkurl{fr/sga7ii/current/expose_XVIII_body.tex}\par

\par\medskip\noindent\hypertarget{app-F7II-XVIII-084.34}{}\textbf{F7II-XVIII-084.34}\quad p. 308; inherited corrected-layer difference.\par
\noindent Imprimé : {\ttfamily\footnotesize \textbackslash{}mathrm\{\allowbreak{}}\par
\noindent Corrigé : \emph{(vide)}\par
\noindent Source : ADOPT-F7II-XVIII-084.34. \nolinkurl{fr/sga7ii/diplomatic/expose_XVIII_body.tex}; \nolinkurl{fr/sga7ii/current/expose_XVIII_body.tex}\par

\par\medskip\noindent\hypertarget{app-F7II-XVIII-084.35}{}\textbf{F7II-XVIII-084.35}\quad p. 308; inherited corrected-layer difference.\par
\noindent Imprimé : {\ttfamily\footnotesize \}\allowbreak{}}\par
\noindent Corrigé : \emph{(vide)}\par
\noindent Source : ADOPT-F7II-XVIII-084.35. \nolinkurl{fr/sga7ii/diplomatic/expose_XVIII_body.tex}; \nolinkurl{fr/sga7ii/current/expose_XVIII_body.tex}\par

\par\medskip\noindent\hypertarget{app-F7II-XVIII-084.36}{}\textbf{F7II-XVIII-084.36}\quad p. 308; inherited corrected-layer difference.\par
\noindent Imprimé : {\ttfamily\footnotesize \textbackslash{}mathrm\{\allowbreak{}}\par
\noindent Corrigé : \emph{(vide)}\par
\noindent Source : ADOPT-F7II-XVIII-084.36. \nolinkurl{fr/sga7ii/diplomatic/expose_XVIII_body.tex}; \nolinkurl{fr/sga7ii/current/expose_XVIII_body.tex}\par

\par\medskip\noindent\hypertarget{app-F7II-XVIII-084.37}{}\textbf{F7II-XVIII-084.37}\quad p. 308; inherited corrected-layer difference.\par
\noindent Imprimé : {\ttfamily\footnotesize \}\allowbreak{}}\par
\noindent Corrigé : \emph{(vide)}\par
\noindent Source : ADOPT-F7II-XVIII-084.37. \nolinkurl{fr/sga7ii/diplomatic/expose_XVIII_body.tex}; \nolinkurl{fr/sga7ii/current/expose_XVIII_body.tex}\par

\par\medskip\noindent\hypertarget{app-F7II-XVIII-084.38}{}\textbf{F7II-XVIII-084.38}\quad p. 308; inherited corrected-layer difference.\par
\noindent Imprimé : {\ttfamily\footnotesize \{\textbackslash{}mathrm\{\allowbreak{}}\par
\noindent Corrigé : \emph{(vide)}\par
\noindent Source : ADOPT-F7II-XVIII-084.38. \nolinkurl{fr/sga7ii/diplomatic/expose_XVIII_body.tex}; \nolinkurl{fr/sga7ii/current/expose_XVIII_body.tex}\par

\par\medskip\noindent\hypertarget{app-F7II-XVIII-084.39}{}\textbf{F7II-XVIII-084.39}\quad p. 308; inherited corrected-layer difference.\par
\noindent Imprimé : {\ttfamily\footnotesize \}\}\allowbreak{}}\par
\noindent Corrigé : \emph{(vide)}\par
\noindent Source : ADOPT-F7II-XVIII-084.39. \nolinkurl{fr/sga7ii/diplomatic/expose_XVIII_body.tex}; \nolinkurl{fr/sga7ii/current/expose_XVIII_body.tex}\par

\par\medskip\noindent\hypertarget{app-F7II-XVIII-085.1}{}\textbf{F7II-XVIII-085.1}\quad p. 308; designated-witness attribution.\par
\noindent Imprimé : {\ttfamily\footnotesize \textbackslash{}mathrm\{\allowbreak{}}\par
\noindent Corrigé : \emph{(vide)}\par
\noindent Source : WITNESS-F7II-XVIII-085.1. \nolinkurl{fr/sga7ii/diplomatic/expose_XVIII_body.tex}; \nolinkurl{fr/sga7ii/current/expose_XVIII_body.tex}\par

\par\medskip\noindent\hypertarget{app-F7II-XVIII-085.2}{}\textbf{F7II-XVIII-085.2}\quad p. 308; designated-witness attribution.\par
\noindent Imprimé : {\ttfamily\footnotesize \}\allowbreak{}}\par
\noindent Corrigé : \emph{(vide)}\par
\noindent Source : WITNESS-F7II-XVIII-085.2. \nolinkurl{fr/sga7ii/diplomatic/expose_XVIII_body.tex}; \nolinkurl{fr/sga7ii/current/expose_XVIII_body.tex}\par

\par\medskip\noindent\hypertarget{app-F7II-XVIII-085.3}{}\textbf{F7II-XVIII-085.3}\quad p. 308; designated-witness attribution.\par
\noindent Imprimé : {\ttfamily\footnotesize \textbackslash{}mathrm\{\allowbreak{}}\par
\noindent Corrigé : \emph{(vide)}\par
\noindent Source : WITNESS-F7II-XVIII-085.3. \nolinkurl{fr/sga7ii/diplomatic/expose_XVIII_body.tex}; \nolinkurl{fr/sga7ii/current/expose_XVIII_body.tex}\par

\par\medskip\noindent\hypertarget{app-F7II-XVIII-085.4}{}\textbf{F7II-XVIII-085.4}\quad p. 308; designated-witness attribution.\par
\noindent Imprimé : {\ttfamily\footnotesize \}\_\allowbreak{}}\par
\noindent Corrigé : \emph{(vide)}\par
\noindent Source : WITNESS-F7II-XVIII-085.4. \nolinkurl{fr/sga7ii/diplomatic/expose_XVIII_body.tex}; \nolinkurl{fr/sga7ii/current/expose_XVIII_body.tex}\par

\par\medskip\noindent\hypertarget{app-F7II-XVIII-085.5}{}\textbf{F7II-XVIII-085.5}\quad p. 308; designated-witness attribution.\par
\noindent Imprimé : {\ttfamily\footnotesize )\allowbreak{}}\par
\noindent Corrigé : \emph{(vide)}\par
\noindent Source : WITNESS-F7II-XVIII-085.5. \nolinkurl{fr/sga7ii/diplomatic/expose_XVIII_body.tex}; \nolinkurl{fr/sga7ii/current/expose_XVIII_body.tex}\par

\par\medskip\noindent\hypertarget{app-F7II-XVIII-085.6}{}\textbf{F7II-XVIII-085.6}\quad p. 308; designated-witness attribution.\par
\noindent Imprimé : {\ttfamily\footnotesize \textbackslash{}\allowbreak{} }\par
\noindent Corrigé : \emph{(vide)}\par
\noindent Source : WITNESS-F7II-XVIII-085.6. \nolinkurl{fr/sga7ii/diplomatic/expose_XVIII_body.tex}; \nolinkurl{fr/sga7ii/current/expose_XVIII_body.tex}\par

\par\medskip\noindent\hypertarget{app-F7II-XVIII-085.7}{}\textbf{F7II-XVIII-085.7}\quad p. 308; designated-witness attribution.\par
\noindent Imprimé : {\ttfamily\footnotesize  \textbackslash{}\allowbreak{} \textbackslash{}mathrm\{\allowbreak{}}\par
\noindent Corrigé : \emph{(vide)}\par
\noindent Source : WITNESS-F7II-XVIII-085.7. \nolinkurl{fr/sga7ii/diplomatic/expose_XVIII_body.tex}; \nolinkurl{fr/sga7ii/current/expose_XVIII_body.tex}\par

\par\medskip\noindent\hypertarget{app-F7II-XVIII-085.8}{}\textbf{F7II-XVIII-085.8}\quad p. 308; designated-witness attribution.\par
\noindent Imprimé : {\ttfamily\footnotesize \}\allowbreak{}}\par
\noindent Corrigé : \emph{(vide)}\par
\noindent Source : WITNESS-F7II-XVIII-085.8. \nolinkurl{fr/sga7ii/diplomatic/expose_XVIII_body.tex}; \nolinkurl{fr/sga7ii/current/expose_XVIII_body.tex}\par

\par\medskip\noindent\hypertarget{app-F7II-XVIII-086.1}{}\textbf{F7II-XVIII-086.1}\quad p. 309; designated-witness attribution.\par
\noindent Imprimé : {\ttfamily\footnotesize *\allowbreak{}}\par
\noindent Corrigé : {\ttfamily\footnotesize \}\textbackslash{}tag\{5.8.7.3\allowbreak{}}\par
\noindent Source : WITNESS-F7II-XVIII-086.1. \nolinkurl{fr/sga7ii/diplomatic/expose_XVIII_body.tex}; \nolinkurl{fr/sga7ii/current/expose_XVIII_body.tex}\par

\par\medskip\noindent\hypertarget{app-F7II-XVIII-086.2}{}\textbf{F7II-XVIII-086.2}\quad p. 309; inherited corrected-layer difference.\par
\noindent Imprimé : {\ttfamily\footnotesize *\allowbreak{}}\par
\noindent Corrigé : \emph{(vide)}\par
\noindent Source : ADOPT-F7II-XVIII-086.2. \nolinkurl{fr/sga7ii/diplomatic/expose_XVIII_body.tex}; \nolinkurl{fr/sga7ii/current/expose_XVIII_body.tex}\par

\par\medskip\noindent\hypertarget{app-F7II-XVIII-087.1}{}\textbf{F7II-XVIII-087.1}\quad p. 309; designated-witness attribution.\par
\noindent Imprimé : \emph{(vide)}\par
\noindent Corrigé : {\ttfamily\footnotesize \_U\allowbreak{}}\par
\noindent Source : WITNESS-F7II-XVIII-087.1. \nolinkurl{fr/sga7ii/diplomatic/expose_XVIII_body.tex}; \nolinkurl{fr/sga7ii/current/expose_XVIII_body.tex}\par

\par\medskip\noindent\hypertarget{app-F7II-XVIII-088.1}{}\textbf{F7II-XVIII-088.1}\quad p. 310; printed misprint.\par
\noindent Imprimé : \emph{(vide)}\par
\noindent Corrigé : {\ttfamily\footnotesize (j)\allowbreak{}}\par
\noindent Source : 7II:952. \nolinkurl{fixes/MISPRINTS_7II.json#952}\par

\par\medskip\noindent\hypertarget{app-F7II-XVIII-089.1}{}\textbf{F7II-XVIII-089.1}\quad p. 310; designated-witness attribution.\par
\noindent Imprimé : {\ttfamily\footnotesize \textbackslash{}\allowbreak{} }\par
\noindent Corrigé : \emph{(vide)}\par
\noindent Source : WITNESS-F7II-XVIII-089.1. \nolinkurl{fr/sga7ii/diplomatic/expose_XVIII_body.tex}; \nolinkurl{fr/sga7ii/current/expose_XVIII_body.tex}\par

\par\medskip\noindent\hypertarget{app-F7II-XVIII-089.2}{}\textbf{F7II-XVIII-089.2}\quad p. 310; designated-witness attribution.\par
\noindent Imprimé : {\ttfamily\footnotesize \textbackslash{}\allowbreak{} }\par
\noindent Corrigé : {\ttfamily\footnotesize F\allowbreak{}}\par
\noindent Source : WITNESS-F7II-XVIII-089.2. \nolinkurl{fr/sga7ii/diplomatic/expose_XVIII_body.tex}; \nolinkurl{fr/sga7ii/current/expose_XVIII_body.tex}\par

\par\medskip\noindent\hypertarget{app-F7II-XVIII-089.3}{}\textbf{F7II-XVIII-089.3}\quad p. 310; designated-witness attribution.\par
\noindent Imprimé : {\ttfamily\footnotesize (j)\allowbreak{}}\par
\noindent Corrigé : \emph{(vide)}\par
\noindent Source : WITNESS-F7II-XVIII-089.3. \nolinkurl{fr/sga7ii/diplomatic/expose_XVIII_body.tex}; \nolinkurl{fr/sga7ii/current/expose_XVIII_body.tex}\par

\par\medskip\noindent\hypertarget{app-F7II-XVIII-089.4}{}\textbf{F7II-XVIII-089.4}\quad p. 310; inherited corrected-layer difference.\par
\noindent Imprimé : \emph{(vide)}\par
\noindent Corrigé : {\ttfamily\footnotesize \_U\allowbreak{}}\par
\noindent Source : ADOPT-F7II-XVIII-089.4. \nolinkurl{fr/sga7ii/diplomatic/expose_XVIII_body.tex}; \nolinkurl{fr/sga7ii/current/expose_XVIII_body.tex}\par

\par\medskip\noindent\hypertarget{app-F7II-XVIII-089.5}{}\textbf{F7II-XVIII-089.5}\quad p. 310; inherited corrected-layer difference.\par
\noindent Imprimé : \emph{(vide)}\par
\noindent Corrigé : {\ttfamily\footnotesize \_U\allowbreak{}}\par
\noindent Source : ADOPT-F7II-XVIII-089.5. \nolinkurl{fr/sga7ii/diplomatic/expose_XVIII_body.tex}; \nolinkurl{fr/sga7ii/current/expose_XVIII_body.tex}\par

\par\medskip\noindent\hypertarget{app-F7II-XVIII-089.6}{}\textbf{F7II-XVIII-089.6}\quad p. 310; inherited corrected-layer difference.\par
\noindent Imprimé : {\ttfamily\footnotesize .\allowbreak{}}\par
\noindent Corrigé : {\ttfamily\footnotesize \textbackslash{},\allowbreak{}}\par
\noindent Source : ADOPT-F7II-XVIII-089.6. \nolinkurl{fr/sga7ii/diplomatic/expose_XVIII_body.tex}; \nolinkurl{fr/sga7ii/current/expose_XVIII_body.tex}\par

\par\medskip\noindent\hypertarget{app-F7II-XVIII-089.7}{}\textbf{F7II-XVIII-089.7}\quad p. 311; printed misprint.\par
\noindent Imprimé : {\ttfamily\footnotesize vérifiéé\allowbreak{}}\par
\noindent Corrigé : {\ttfamily\footnotesize vérifiée\allowbreak{}}\par
\noindent Source : 7II:526; 7II:920. \nolinkurl{fixes/MISPRINTS_7II.json#526}; \nolinkurl{fixes/MISPRINTS_7II.json#920}\par

\par\medskip\noindent\hypertarget{app-F7II-XVIII-089.8}{}\textbf{F7II-XVIII-089.8}\quad p. 311; misprint\_pass.\par
\noindent Imprimé : {\ttfamily\footnotesize chacun\allowbreak{}}\par
\noindent Corrigé : {\ttfamily\footnotesize chacune\allowbreak{}}\par
\noindent Source : 7II:527; 7II:921. \nolinkurl{fixes/FIXLOG_7II_CURRENT.jsonl#758}; \nolinkurl{fixes/MISPRINTS_7II.json#921}\par

\par\medskip\noindent\hypertarget{app-F7II-XVIII-090.1}{}\textbf{F7II-XVIII-090.1}\quad p. 312; misprint\_pass.\par
\noindent Imprimé : {\ttfamily\footnotesize evanescents\allowbreak{}}\par
\noindent Corrigé : {\ttfamily\footnotesize évanescents\allowbreak{}}\par
\noindent Source : 7II:528. \nolinkurl{fixes/FIXLOG_7II_CURRENT.jsonl#759}\par

\par\medskip\noindent\hypertarget{app-F7II-XVIII-091.1}{}\textbf{F7II-XVIII-091.1}\quad p. 312; printed misprint.\par
\noindent Imprimé : {\ttfamily\footnotesize Choissisons\allowbreak{}}\par
\noindent Corrigé : {\ttfamily\footnotesize Choisissons\allowbreak{}}\par
\noindent Source : 7II:529; 7II:923. \nolinkurl{fixes/MISPRINTS_7II.json#529}; \nolinkurl{fixes/MISPRINTS_7II.json#923}\par

\par\medskip\noindent\hypertarget{app-F7II-XVIII-092.1}{}\textbf{F7II-XVIII-092.1}\quad p. 313; printed misprint.\par
\noindent Imprimé : {\ttfamily\footnotesize grassmanniene\allowbreak{}}\par
\noindent Corrigé : {\ttfamily\footnotesize grassmannienne\allowbreak{}}\par
\noindent Source : 7II:530; 7II:924. \nolinkurl{fixes/MISPRINTS_7II.json#530}; \nolinkurl{fixes/MISPRINTS_7II.json#924}\par

\par\medskip\noindent\hypertarget{app-F7II-XVIII-092.2}{}\textbf{F7II-XVIII-092.2}\quad p. 313; designated-witness attribution.\par
\noindent Imprimé : {\ttfamily\footnotesize formée\allowbreak{}}\par
\noindent Corrigé : {\ttfamily\footnotesize formé\allowbreak{}}\par
\noindent Source : WITNESS-F7II-XVIII-092.2. \nolinkurl{fr/sga7ii/diplomatic/expose_XVIII_body.tex}; \nolinkurl{fr/sga7ii/current/expose_XVIII_body.tex}\par

\par\medskip\noindent\hypertarget{app-F7II-XVIII-093.1}{}\textbf{F7II-XVIII-093.1}\quad p. 313; printed misprint.\par
\noindent Imprimé : {\ttfamily\footnotesize joingnant\allowbreak{}}\par
\noindent Corrigé : {\ttfamily\footnotesize joignant\allowbreak{}}\par
\noindent Source : 7II:531; 7II:925. \nolinkurl{fixes/MISPRINTS_7II.json#531}; \nolinkurl{fixes/MISPRINTS_7II.json#925}\par

\par\medskip\noindent\hypertarget{app-F7II-XVIII-093.2}{}\textbf{F7II-XVIII-093.2}\quad p. 313; printed misprint.\par
\noindent Imprimé : \emph{(vide)}\par
\noindent Corrigé : {\ttfamily\footnotesize de\allowbreak{} }\par
\noindent Source : 7II:532; 7II:926. \nolinkurl{fixes/MISPRINTS_7II.json#532}; \nolinkurl{fixes/MISPRINTS_7II.json#926}\par

\par\medskip\noindent\hypertarget{app-F7II-XVIII-094.1}{}\textbf{F7II-XVIII-094.1}\quad p. 314; misprint\_pass.\par
\noindent Imprimé : {\ttfamily\footnotesize O\allowbreak{}}\par
\noindent Corrigé : {\ttfamily\footnotesize o\allowbreak{}}\par
\noindent Source : 7II:533. \nolinkurl{fixes/FIXLOG_7II_CURRENT.jsonl#760}\par

\par\medskip\noindent\hypertarget{app-F7II-XVIII-094.2}{}\textbf{F7II-XVIII-094.2}\quad p. 314; printed misprint.\par
\noindent Imprimé : {\ttfamily\footnotesize )\allowbreak{}}\par
\noindent Corrigé : \emph{(vide)}\par
\noindent Source : 7II:534; 7II:928. \nolinkurl{fixes/MISPRINTS_7II.json#534}; \nolinkurl{fixes/MISPRINTS_7II.json#928}\par

\par\medskip\noindent\hypertarget{app-F7II-XVIII-094.3}{}\textbf{F7II-XVIII-094.3}\quad p. 314; printed misprint.\par
\noindent Imprimé : {\ttfamily\footnotesize homomorphisme\allowbreak{}}\par
\noindent Corrigé : {\ttfamily\footnotesize homomorphismes\allowbreak{}}\par
\noindent Source : 7II:535; 7II:929. \nolinkurl{fixes/MISPRINTS_7II.json#535}; \nolinkurl{fixes/MISPRINTS_7II.json#929}\par

\par\medskip\noindent\hypertarget{app-F7II-XVIII-094.4}{}\textbf{F7II-XVIII-094.4}\quad p. 314; printed misprint.\par
\noindent Imprimé : \emph{(vide)}\par
\noindent Corrigé : {\ttfamily\footnotesize \textasciicircum{}\{\textbackslash{}mathrm\{mod\}\}\allowbreak{}}\par
\noindent Source : 7II:930. \nolinkurl{fixes/MISPRINTS_7II.json#930}\par

\par\medskip\noindent\hypertarget{app-F7II-XVIII-095.1}{}\textbf{F7II-XVIII-095.1}\quad p. 314; designated-witness attribution.\par
\noindent Imprimé : {\ttfamily\footnotesize \_\allowbreak{}}\par
\noindent Corrigé : \emph{(vide)}\par
\noindent Source : WITNESS-F7II-XVIII-095.1. \nolinkurl{fr/sga7ii/diplomatic/expose_XVIII_body.tex}; \nolinkurl{fr/sga7ii/current/expose_XVIII_body.tex}\par

\par\medskip\noindent\hypertarget{app-F7II-XVIII-096.1}{}\textbf{F7II-XVIII-096.1}\quad p. 315; inherited corrected-layer difference.\par
\noindent Imprimé : \emph{(vide)}\par
\noindent Corrigé : {\ttfamily\footnotesize *\allowbreak{}}\par
\noindent Source : ADOPT-F7II-XVIII-096.1. \nolinkurl{fr/sga7ii/diplomatic/expose_XVIII_body.tex}; \nolinkurl{fr/sga7ii/current/expose_XVIII_body.tex}\par

\par\medskip\noindent\hypertarget{app-F7II-XVIII-096.2}{}\textbf{F7II-XVIII-096.2}\quad p. 315; printed misprint.\par
\noindent Imprimé : {\ttfamily\footnotesize composé\allowbreak{}}\par
\noindent Corrigé : {\ttfamily\footnotesize composée\allowbreak{}}\par
\noindent Source : 7II:536; 7II:931. \nolinkurl{fixes/MISPRINTS_7II.json#536}; \nolinkurl{fixes/MISPRINTS_7II.json#931}\par

\par\medskip\noindent\hypertarget{app-F7II-XVIII-097.1}{}\textbf{F7II-XVIII-097.1}\quad p. 316; printed misprint.\par
\noindent Imprimé : {\ttfamily\footnotesize K\allowbreak{}}\par
\noindent Corrigé : {\ttfamily\footnotesize k\allowbreak{}}\par
\noindent Source : 7II:932. \nolinkurl{fixes/MISPRINTS_7II.json#932}\par

\par\medskip\noindent\hypertarget{app-F7II-XVIII-097.2}{}\textbf{F7II-XVIII-097.2}\quad p. 316; inherited corrected-layer difference.\par
\noindent Imprimé : \emph{(1 espace(s), 0 retour(s) de ligne dans le témoin TeX)}\par
\noindent Corrigé : \emph{(vide)}\par
\noindent Source : ADOPT-F7II-XVIII-097.2. \nolinkurl{fr/sga7ii/diplomatic/expose_XVIII_body.tex}; \nolinkurl{fr/sga7ii/current/expose_XVIII_body.tex}\par

\par\medskip\noindent\hypertarget{app-F7II-XVIII-097.3}{}\textbf{F7II-XVIII-097.3}\quad p. 316; printed misprint.\par
\noindent Imprimé : {\ttfamily\footnotesize \textbackslash{}overline\{\allowbreak{}}\par
\noindent Corrigé : \emph{(vide)}\par
\noindent Source : 7II:538. \nolinkurl{fixes/MISPRINTS_7II.json#538}\par

\par\medskip\noindent\hypertarget{app-F7II-XVIII-097.4}{}\textbf{F7II-XVIII-097.4}\quad p. 316; printed misprint.\par
\noindent Imprimé : {\ttfamily\footnotesize \}\allowbreak{}}\par
\noindent Corrigé : \emph{(vide)}\par
\noindent Source : 7II:538. \nolinkurl{fixes/MISPRINTS_7II.json#538}\par

\par\medskip\noindent\hypertarget{app-F7II-XVIII-097.5}{}\textbf{F7II-XVIII-097.5}\quad p. 316; designated-witness attribution.\par
\noindent Imprimé : \emph{(1 espace(s), 0 retour(s) de ligne dans le témoin TeX)}\par
\noindent Corrigé : \emph{(vide)}\par
\noindent Source : WITNESS-F7II-XVIII-097.5. \nolinkurl{fr/sga7ii/diplomatic/expose_XVIII_body.tex}; \nolinkurl{fr/sga7ii/current/expose_XVIII_body.tex}\par

\par\medskip\noindent\hypertarget{app-F7II-XVIII-097.6}{}\textbf{F7II-XVIII-097.6}\quad p. 316; printed misprint.\par
\noindent Imprimé : {\ttfamily\footnotesize suffisament\allowbreak{}}\par
\noindent Corrigé : {\ttfamily\footnotesize suffisamment\allowbreak{}}\par
\noindent Source : 7II:539; 7II:934. \nolinkurl{fixes/MISPRINTS_7II.json#539}; \nolinkurl{fixes/MISPRINTS_7II.json#934}\par

\par\medskip\noindent\hypertarget{app-F7II-XVIII-098.1}{}\textbf{F7II-XVIII-098.1}\quad p. 316; printed misprint.\par
\noindent Imprimé : \emph{(vide)}\par
\noindent Corrigé : {\ttfamily\footnotesize k(\textbackslash{}overline\{\textbackslash{}eta\})\allowbreak{}}\par
\noindent Source : 7II:935. \nolinkurl{fixes/MISPRINTS_7II.json#935}\par

\par\medskip\noindent\hypertarget{app-F7II-XVIII-098.2}{}\textbf{F7II-XVIII-098.2}\quad p. 316; printed misprint.\par
\noindent Imprimé : \emph{(vide)}\par
\noindent Corrigé : {\ttfamily\footnotesize )\allowbreak{}}\par
\noindent Source : 7II:935. \nolinkurl{fixes/MISPRINTS_7II.json#935}\par

\par\medskip\noindent\hypertarget{app-F7II-XVIII-099.1}{}\textbf{F7II-XVIII-099.1}\quad p. 317; designated-witness attribution.\par
\noindent Imprimé : \emph{(1 espace(s), 0 retour(s) de ligne dans le témoin TeX)}\par
\noindent Corrigé : \emph{(vide)}\par
\noindent Source : WITNESS-F7II-XVIII-099.1. \nolinkurl{fr/sga7ii/diplomatic/expose_XVIII_body.tex}; \nolinkurl{fr/sga7ii/current/expose_XVIII_body.tex}\par

\par\medskip\noindent\hypertarget{app-F7II-XVIII-099.2}{}\textbf{F7II-XVIII-099.2}\quad p. 317; printed misprint.\par
\noindent Imprimé : \emph{(vide)}\par
\noindent Corrigé : {\ttfamily\footnotesize k(\textbackslash{}overline\{\textbackslash{}eta\})\allowbreak{}}\par
\noindent Source : 7II:936. \nolinkurl{fixes/MISPRINTS_7II.json#936}\par

\par\medskip\noindent\hypertarget{app-F7II-XVIII-099.3}{}\textbf{F7II-XVIII-099.3}\quad p. 317; printed misprint.\par
\noindent Imprimé : {\ttfamily\footnotesize K\allowbreak{}}\par
\noindent Corrigé : {\ttfamily\footnotesize D\allowbreak{}}\par
\noindent Source : 7II:937. \nolinkurl{fixes/MISPRINTS_7II.json#937}\par

\par\medskip\noindent\hypertarget{app-F7II-XVIII-099.4}{}\textbf{F7II-XVIII-099.4}\quad p. 317; printed misprint.\par
\noindent Imprimé : \emph{(vide)}\par
\noindent Corrigé : {\ttfamily\footnotesize k(\textbackslash{}overline\{\textbackslash{}eta\})\allowbreak{}}\par
\noindent Source : 7II:938. \nolinkurl{fixes/MISPRINTS_7II.json#938}\par

\par\medskip\noindent\hypertarget{app-F7II-XVIII-099.5}{}\textbf{F7II-XVIII-099.5}\quad p. 317; inherited corrected-layer difference.\par
\noindent Imprimé : {\ttfamily\footnotesize a\allowbreak{}}\par
\noindent Corrigé : {\ttfamily\footnotesize ait\allowbreak{}}\par
\noindent Source : ADOPT-F7II-XVIII-099.5. \nolinkurl{fr/sga7ii/diplomatic/expose_XVIII_body.tex}; \nolinkurl{fr/sga7ii/current/expose_XVIII_body.tex}\par

\par\medskip\noindent\hypertarget{app-F7II-XVIII-099.6}{}\textbf{F7II-XVIII-099.6}\quad p. 317; inherited corrected-layer difference.\par
\noindent Imprimé : {\ttfamily\footnotesize c'est\allowbreak{}}\par
\noindent Corrigé : {\ttfamily\footnotesize ce\allowbreak{} soit\allowbreak{}}\par
\noindent Source : ADOPT-F7II-XVIII-099.6. \nolinkurl{fr/sga7ii/diplomatic/expose_XVIII_body.tex}; \nolinkurl{fr/sga7ii/current/expose_XVIII_body.tex}\par

\par\medskip\noindent\hypertarget{app-F7II-XVIII-099.7}{}\textbf{F7II-XVIII-099.7}\quad p. 317; inherited corrected-layer difference.\par
\noindent Imprimé : \emph{(vide)}\par
\noindent Corrigé : \emph{(1 espace(s), 0 retour(s) de ligne dans le témoin TeX)}\par
\noindent Source : ADOPT-F7II-XVIII-099.7. \nolinkurl{fr/sga7ii/diplomatic/expose_XVIII_body.tex}; \nolinkurl{fr/sga7ii/current/expose_XVIII_body.tex}\par

\par\medskip\noindent\hypertarget{app-F7II-XVIII-100.1}{}\textbf{F7II-XVIII-100.1}\quad p. 317; printed misprint.\par
\noindent Imprimé : \emph{(vide)}\par
\noindent Corrigé : {\ttfamily\footnotesize )\allowbreak{}}\par
\noindent Source : 7II:542; 7II:939. \nolinkurl{fixes/MISPRINTS_7II.json#542}; \nolinkurl{fixes/MISPRINTS_7II.json#939}\par

\par\medskip\noindent\hypertarget{app-F7II-XVIII-101.1}{}\textbf{F7II-XVIII-101.1}\quad p. 318; printed misprint.\par
\noindent Imprimé : {\ttfamily\footnotesize  \$k\$\allowbreak{}}\par
\noindent Corrigé : \emph{(vide)}\par
\noindent Source : 7II:940. \nolinkurl{fixes/MISPRINTS_7II.json#940}\par

\par\medskip\noindent\hypertarget{app-F7II-XVIII-101.2}{}\textbf{F7II-XVIII-101.2}\quad p. 318; printed misprint.\par
\noindent Imprimé : {\ttfamily\footnotesize card\allowbreak{}}\par
\noindent Corrigé : {\ttfamily\footnotesize car\allowbreak{}}\par
\noindent Source : 7II:543; 7II:941. \nolinkurl{fixes/MISPRINTS_7II.json#543}; \nolinkurl{fixes/MISPRINTS_7II.json#941}\par

\par\medskip\noindent\hypertarget{app-F7II-XVIII-101.3}{}\textbf{F7II-XVIII-101.3}\quad p. 318; printed misprint.\par
\noindent Imprimé : {\ttfamily\footnotesize alor\allowbreak{}}\par
\noindent Corrigé : {\ttfamily\footnotesize alors\allowbreak{}}\par
\noindent Source : 7II:942. \nolinkurl{fixes/MISPRINTS_7II.json#942}\par

\par\medskip\noindent\hypertarget{app-F7II-XVIII-101.4}{}\textbf{F7II-XVIII-101.4}\quad p. 318; misprint\_pass.\par
\noindent Imprimé : {\ttfamily\footnotesize vue\allowbreak{}}\par
\noindent Corrigé : {\ttfamily\footnotesize vu\allowbreak{}}\par
\noindent Source : 7II:943. \nolinkurl{fixes/FIXLOG_7II_CURRENT.jsonl#903}\par

\par\medskip\noindent\hypertarget{app-F7II-XVIII-102.1}{}\textbf{F7II-XVIII-102.1}\quad p. 319; printed misprint.\par
\noindent Imprimé : \emph{(vide)}\par
\noindent Corrigé : {\ttfamily\footnotesize  de\allowbreak{}}\par
\noindent Source : 7II:546; 7II:944. \nolinkurl{fixes/MISPRINTS_7II.json#546}; \nolinkurl{fixes/MISPRINTS_7II.json#944}\par

\par\medskip\noindent\hypertarget{app-F7II-XVIII-103.1}{}\textbf{F7II-XVIII-103.1}\quad p. 320; misprint\_pass.\par
\noindent Imprimé : {\ttfamily\footnotesize 6\allowbreak{}}\par
\noindent Corrigé : {\ttfamily\footnotesize 10\allowbreak{}}\par
\noindent Source : 7II:945. \nolinkurl{fixes/FIXLOG_7II_CURRENT.jsonl#904}; \nolinkurl{fixes/MISPRINTS_7II.json#945}\par

\par\medskip\noindent\hypertarget{app-F7II-XVIII-103.2}{}\textbf{F7II-XVIII-103.2}\quad p. 320; printed misprint.\par
\noindent Imprimé : \emph{(vide)}\par
\noindent Corrigé : {\ttfamily\footnotesize \textbackslash{}check\{\allowbreak{}}\par
\noindent Source : 7II:547; 7II:946. \nolinkurl{fixes/MISPRINTS_7II.json#547}; \nolinkurl{fixes/MISPRINTS_7II.json#946}\par

\par\medskip\noindent\hypertarget{app-F7II-XVIII-103.3}{}\textbf{F7II-XVIII-103.3}\quad p. 320; printed misprint.\par
\noindent Imprimé : \emph{(vide)}\par
\noindent Corrigé : {\ttfamily\footnotesize \}\allowbreak{}}\par
\noindent Source : 7II:547; 7II:946. \nolinkurl{fixes/MISPRINTS_7II.json#547}; \nolinkurl{fixes/MISPRINTS_7II.json#946}\par

\par\medskip\noindent\hypertarget{app-F7II-XVIII-103.4}{}\textbf{F7II-XVIII-103.4}\quad p. 320; printed misprint.\par
\noindent Imprimé : \emph{(vide)}\par
\noindent Corrigé : {\ttfamily\footnotesize \textbackslash{}mathbb\{\allowbreak{}}\par
\noindent Source : 7II:947. \nolinkurl{fixes/MISPRINTS_7II.json#947}\par

\par\medskip\noindent\hypertarget{app-F7II-XVIII-103.5}{}\textbf{F7II-XVIII-103.5}\quad p. 320; printed misprint.\par
\noindent Imprimé : \emph{(vide)}\par
\noindent Corrigé : {\ttfamily\footnotesize \}\allowbreak{}}\par
\noindent Source : 7II:947. \nolinkurl{fixes/MISPRINTS_7II.json#947}\par

\par\medskip\noindent\hypertarget{app-F7II-XVIII-103.6}{}\textbf{F7II-XVIII-103.6}\quad p. 320; misprint\_pass.\par
\noindent Imprimé : {\ttfamily\footnotesize 2\allowbreak{}}\par
\noindent Corrigé : {\ttfamily\footnotesize 1\allowbreak{}}\par
\noindent Source : 7II:948. \nolinkurl{fixes/FIXLOG_7II_CURRENT.jsonl#905}; \nolinkurl{fixes/MISPRINTS_7II.json#948}\par

\par\medskip\noindent\hypertarget{app-F7II-XVIII-103.7}{}\textbf{F7II-XVIII-103.7}\quad p. 320; inherited corrected-layer difference.\par
\noindent Imprimé : \emph{(vide)}\par
\noindent Corrigé : {\ttfamily\footnotesize (\allowbreak{}}\par
\noindent Source : ADOPT-F7II-XVIII-103.7. \nolinkurl{fr/sga7ii/diplomatic/expose_XVIII_body.tex}; \nolinkurl{fr/sga7ii/current/expose_XVIII_body.tex}\par

\par\medskip\noindent\hypertarget{app-F7II-XVIII-103.8}{}\textbf{F7II-XVIII-103.8}\quad p. 320; inherited corrected-layer difference.\par
\noindent Imprimé : {\ttfamily\footnotesize est\allowbreak{}}\par
\noindent Corrigé : {\ttfamily\footnotesize soit\allowbreak{}}\par
\noindent Source : ADOPT-F7II-XVIII-103.8. \nolinkurl{fr/sga7ii/diplomatic/expose_XVIII_body.tex}; \nolinkurl{fr/sga7ii/current/expose_XVIII_body.tex}\par

\par\medskip\noindent\hypertarget{app-F7II-XVIII-103.9}{}\textbf{F7II-XVIII-103.9}\quad p. 320; inherited corrected-layer difference.\par
\noindent Imprimé : {\ttfamily\footnotesize paire\allowbreak{}}\par
\noindent Corrigé : {\ttfamily\footnotesize pair\allowbreak{}}\par
\noindent Source : ADOPT-F7II-XVIII-103.9. \nolinkurl{fr/sga7ii/diplomatic/expose_XVIII_body.tex}; \nolinkurl{fr/sga7ii/current/expose_XVIII_body.tex}\par

\par\medskip\noindent\hypertarget{app-F7II-XVIII-103.10}{}\textbf{F7II-XVIII-103.10}\quad p. 320; inherited corrected-layer difference.\par
\noindent Imprimé : \emph{(vide)}\par
\noindent Corrigé : {\ttfamily\footnotesize )\allowbreak{}}\par
\noindent Source : ADOPT-F7II-XVIII-103.10. \nolinkurl{fr/sga7ii/diplomatic/expose_XVIII_body.tex}; \nolinkurl{fr/sga7ii/current/expose_XVIII_body.tex}\par

\par\medskip\noindent\hypertarget{app-F7II-XVIII-103.11}{}\textbf{F7II-XVIII-103.11}\quad p. 320; inherited corrected-layer difference.\par
\noindent Imprimé : {\ttfamily\footnotesize satisfait\allowbreak{}}\par
\noindent Corrigé : {\ttfamily\footnotesize satisfasse\allowbreak{}}\par
\noindent Source : ADOPT-F7II-XVIII-103.11. \nolinkurl{fr/sga7ii/diplomatic/expose_XVIII_body.tex}; \nolinkurl{fr/sga7ii/current/expose_XVIII_body.tex}\par

\par\medskip\noindent\hypertarget{app-F7II-XVIII-103.12}{}\textbf{F7II-XVIII-103.12}\quad p. 321; printed misprint.\par
\noindent Imprimé : {\ttfamily\footnotesize sectionsde\allowbreak{}}\par
\noindent Corrigé : {\ttfamily\footnotesize sections\allowbreak{} de\allowbreak{}}\par
\noindent Source : 7II:548. \nolinkurl{fixes/MISPRINTS_7II.json#548}\par

\par\medskip\noindent\hypertarget{app-F7II-XVIII-104.1}{}\textbf{F7II-XVIII-104.1}\quad p. 321; printed misprint.\par
\noindent Imprimé : \emph{(vide)}\par
\noindent Corrigé : {\ttfamily\footnotesize )\allowbreak{}}\par
\noindent Source : 7II:549. \nolinkurl{fixes/MISPRINTS_7II.json#549}\par

\par\medskip\noindent\hypertarget{app-F7II-XVIII-105.1}{}\textbf{F7II-XVIII-105.1}\quad p. 322; printed misprint.\par
\noindent Imprimé : \emph{(vide)}\par
\noindent Corrigé : {\ttfamily\footnotesize ,\allowbreak{}}\par
\noindent Source : 7II:550. \nolinkurl{fixes/MISPRINTS_7II.json#550}\par

\par\medskip\noindent\hypertarget{app-F7II-XVIII-106.1}{}\textbf{F7II-XVIII-106.1}\quad p. 323; printed misprint.\par
\noindent Imprimé : {\ttfamily\footnotesize 2\allowbreak{}}\par
\noindent Corrigé : {\ttfamily\footnotesize 1\allowbreak{}}\par
\noindent Source : 7II:551. \nolinkurl{fixes/MISPRINTS_7II.json#551}\par

\par\medskip\noindent\hypertarget{app-F7II-XVIII-106.2}{}\textbf{F7II-XVIII-106.2}\quad p. 323; printed misprint.\par
\noindent Imprimé : {\ttfamily\footnotesize i\allowbreak{}}\par
\noindent Corrigé : {\ttfamily\footnotesize \{n-1\}\allowbreak{}}\par
\noindent Source : 7II:552. \nolinkurl{fixes/MISPRINTS_7II.json#552}\par

\par\medskip\noindent\hypertarget{app-F7II-XVIII-106.3}{}\textbf{F7II-XVIII-106.3}\quad p. 323; printed misprint.\par
\noindent Imprimé : \emph{(vide)}\par
\noindent Corrigé : {\ttfamily\footnotesize )\allowbreak{}}\par
\noindent Source : 7II:553. \nolinkurl{fixes/MISPRINTS_7II.json#553}\par

\par\medskip\noindent\hypertarget{app-F7II-XVIII-106.4}{}\textbf{F7II-XVIII-106.4}\quad p. 323; designated-witness attribution.\par
\noindent Imprimé : \emph{(vide)}\par
\noindent Corrigé : {\ttfamily\footnotesize .\allowbreak{}}\par
\noindent Source : WITNESS-F7II-XVIII-106.4. \nolinkurl{fr/sga7ii/diplomatic/expose_XVIII_body.tex}; \nolinkurl{fr/sga7ii/current/expose_XVIII_body.tex}\par

\par\medskip\noindent\hypertarget{app-F7II-XVIII-106.5}{}\textbf{F7II-XVIII-106.5}\quad p. 323; designated-witness attribution.\par
\noindent Imprimé : \emph{(1 espace(s), 0 retour(s) de ligne dans le témoin TeX)}\par
\noindent Corrigé : \emph{(vide)}\par
\noindent Source : WITNESS-F7II-XVIII-106.5. \nolinkurl{fr/sga7ii/diplomatic/expose_XVIII_body.tex}; \nolinkurl{fr/sga7ii/current/expose_XVIII_body.tex}\par

\par\medskip\noindent\hypertarget{app-F7II-XVIII-106.6}{}\textbf{F7II-XVIII-106.6}\quad p. 323; designated-witness attribution.\par
\noindent Imprimé : \emph{(vide)}\par
\noindent Corrigé : {\ttfamily\footnotesize )\allowbreak{}}\par
\noindent Source : WITNESS-F7II-XVIII-106.6. \nolinkurl{fr/sga7ii/diplomatic/expose_XVIII_body.tex}; \nolinkurl{fr/sga7ii/current/expose_XVIII_body.tex}\par

\par\medskip\noindent\hypertarget{app-F7II-XVIII-106.7}{}\textbf{F7II-XVIII-106.7}\quad p. 323; designated-witness attribution.\par
\noindent Imprimé : \emph{(1 espace(s), 0 retour(s) de ligne dans le témoin TeX)}\par
\noindent Corrigé : \emph{(vide)}\par
\noindent Source : WITNESS-F7II-XVIII-106.7. \nolinkurl{fr/sga7ii/diplomatic/expose_XVIII_body.tex}; \nolinkurl{fr/sga7ii/current/expose_XVIII_body.tex}\par

\par\medskip\noindent\hypertarget{app-F7II-XVIII-106.8}{}\textbf{F7II-XVIII-106.8}\quad p. 323; printed misprint.\par
\noindent Imprimé : {\ttfamily\footnotesize \textbackslash{}chi\allowbreak{}}\par
\noindent Corrigé : {\ttfamily\footnotesize x\allowbreak{}}\par
\noindent Source : 7II:555. \nolinkurl{fixes/MISPRINTS_7II.json#555}\par

\par\medskip\noindent\hypertarget{app-F7II-XVIII-106.9}{}\textbf{F7II-XVIII-106.9}\quad p. 323; printed misprint.\par
\noindent Imprimé : \emph{(vide)}\par
\noindent Corrigé : {\ttfamily\footnotesize )\allowbreak{}}\par
\noindent Source : 7II:556. \nolinkurl{fixes/MISPRINTS_7II.json#556}\par

\par\medskip\noindent\hypertarget{app-F7II-XVIII-106.10}{}\textbf{F7II-XVIII-106.10}\quad p. 323; designated-witness attribution.\par
\noindent Imprimé : \emph{(1 espace(s), 0 retour(s) de ligne dans le témoin TeX)}\par
\noindent Corrigé : \emph{(vide)}\par
\noindent Source : WITNESS-F7II-XVIII-106.10. \nolinkurl{fr/sga7ii/diplomatic/expose_XVIII_body.tex}; \nolinkurl{fr/sga7ii/current/expose_XVIII_body.tex}\par

\par\medskip\noindent\hypertarget{app-F7II-XVIII-106.11}{}\textbf{F7II-XVIII-106.11}\quad p. 323; designated-witness attribution.\par
\noindent Imprimé : \emph{(vide)}\par
\noindent Corrigé : {\ttfamily\footnotesize )\allowbreak{}}\par
\noindent Source : WITNESS-F7II-XVIII-106.11. \nolinkurl{fr/sga7ii/diplomatic/expose_XVIII_body.tex}; \nolinkurl{fr/sga7ii/current/expose_XVIII_body.tex}\par

\par\medskip\noindent\hypertarget{app-F7II-XVIII-106.12}{}\textbf{F7II-XVIII-106.12}\quad p. 323; designated-witness attribution.\par
\noindent Imprimé : \emph{(vide)}\par
\noindent Corrigé : {\ttfamily\footnotesize ,\allowbreak{}}\par
\noindent Source : WITNESS-F7II-XVIII-106.12. \nolinkurl{fr/sga7ii/diplomatic/expose_XVIII_body.tex}; \nolinkurl{fr/sga7ii/current/expose_XVIII_body.tex}\par

\par\medskip\noindent\hypertarget{app-F7II-XVIII-106.13}{}\textbf{F7II-XVIII-106.13}\quad p. 323; designated-witness attribution.\par
\noindent Imprimé : \emph{(vide)}\par
\noindent Corrigé : {\ttfamily\footnotesize \$\allowbreak{}}\par
\noindent Source : WITNESS-F7II-XVIII-106.13. \nolinkurl{fr/sga7ii/diplomatic/expose_XVIII_body.tex}; \nolinkurl{fr/sga7ii/current/expose_XVIII_body.tex}\par

\par\medskip\noindent\hypertarget{app-F7II-XVIII-106.14}{}\textbf{F7II-XVIII-106.14}\quad p. 323; designated-witness attribution.\par
\noindent Imprimé : {\ttfamily\footnotesize \$\allowbreak{}}\par
\noindent Corrigé : \emph{(vide)}\par
\noindent Source : WITNESS-F7II-XVIII-106.14. \nolinkurl{fr/sga7ii/diplomatic/expose_XVIII_body.tex}; \nolinkurl{fr/sga7ii/current/expose_XVIII_body.tex}\par

\par\medskip\noindent\hypertarget{app-F7II-XVIII-107.1}{}\textbf{F7II-XVIII-107.1}\quad p. 324; inherited corrected-layer difference.\par
\noindent Imprimé : {\ttfamily\footnotesize par\allowbreak{}}\par
\noindent Corrigé : {\ttfamily\footnotesize pour\allowbreak{}}\par
\noindent Source : ADOPT-F7II-XVIII-107.1. \nolinkurl{fr/sga7ii/diplomatic/expose_XVIII_body.tex}; \nolinkurl{fr/sga7ii/current/expose_XVIII_body.tex}\par

\par\medskip\noindent\hypertarget{app-F7II-XVIII-107.2}{}\textbf{F7II-XVIII-107.2}\quad p. 324; inherited corrected-layer difference.\par
\noindent Imprimé : \emph{(1 espace(s), 0 retour(s) de ligne dans le témoin TeX)}\par
\noindent Corrigé : \emph{(vide)}\par
\noindent Source : ADOPT-F7II-XVIII-107.2. \nolinkurl{fr/sga7ii/diplomatic/expose_XVIII_body.tex}; \nolinkurl{fr/sga7ii/current/expose_XVIII_body.tex}\par

\par\medskip\noindent\hypertarget{app-F7II-XVIII-107.3}{}\textbf{F7II-XVIII-107.3}\quad p. 324; designated-witness attribution.\par
\noindent Imprimé : \emph{(vide)}\par
\noindent Corrigé : {\ttfamily\footnotesize en\}\allowbreak{} \$s\_i\$,\allowbreak{} \textbackslash{}underline\{identifié\allowbreak{} au\allowbreak{} moyen\allowbreak{} }\par
\noindent Source : WITNESS-F7II-XVIII-107.3. \nolinkurl{fr/sga7ii/diplomatic/expose_XVIII_body.tex}; \nolinkurl{fr/sga7ii/current/expose_XVIII_body.tex}\par

\par\medskip\noindent\hypertarget{app-F7II-XVIII-107.4}{}\textbf{F7II-XVIII-107.4}\quad p. 324; designated-witness attribution.\par
\noindent Imprimé : {\ttfamily\footnotesize \textbackslash{}delta\allowbreak{}}\par
\noindent Corrigé : {\ttfamily\footnotesize \textbackslash{}ell\allowbreak{}}\par
\noindent Source : WITNESS-F7II-XVIII-107.4. \nolinkurl{fr/sga7ii/diplomatic/expose_XVIII_body.tex}; \nolinkurl{fr/sga7ii/current/expose_XVIII_body.tex}\par

\par\medskip\noindent\hypertarget{app-F7II-XVIII-107.5}{}\textbf{F7II-XVIII-107.5}\quad p. 324; designated-witness attribution.\par
\noindent Imprimé : {\ttfamily\footnotesize \{\allowbreak{}}\par
\noindent Corrigé : \emph{(vide)}\par
\noindent Source : WITNESS-F7II-XVIII-107.5. \nolinkurl{fr/sga7ii/diplomatic/expose_XVIII_body.tex}; \nolinkurl{fr/sga7ii/current/expose_XVIII_body.tex}\par

\par\medskip\noindent\hypertarget{app-F7II-XVIII-107.6}{}\textbf{F7II-XVIII-107.6}\quad p. 324; designated-witness attribution.\par
\noindent Imprimé : {\ttfamily\footnotesize \}\allowbreak{}}\par
\noindent Corrigé : \emph{(vide)}\par
\noindent Source : WITNESS-F7II-XVIII-107.6. \nolinkurl{fr/sga7ii/diplomatic/expose_XVIII_body.tex}; \nolinkurl{fr/sga7ii/current/expose_XVIII_body.tex}\par

\par\medskip\noindent\hypertarget{app-F7II-XVIII-107.7}{}\textbf{F7II-XVIII-107.7}\quad p. 324; designated-witness attribution.\par
\noindent Imprimé : \emph{(vide)}\par
\noindent Corrigé : {\ttfamily\footnotesize ,\allowbreak{}}\par
\noindent Source : WITNESS-F7II-XVIII-107.7. \nolinkurl{fr/sga7ii/diplomatic/expose_XVIII_body.tex}; \nolinkurl{fr/sga7ii/current/expose_XVIII_body.tex}\par

\par\medskip\noindent\hypertarget{app-F7II-XVIII-107.8}{}\textbf{F7II-XVIII-107.8}\quad p. 324; designated-witness attribution.\par
\noindent Imprimé : {\ttfamily\footnotesize en\}\allowbreak{} \$S\_\{i\}\$\allowbreak{} \textbackslash{}underline\{agit\allowbreak{}}\par
\noindent Corrigé : {\ttfamily\footnotesize agisse\allowbreak{}}\par
\noindent Source : WITNESS-F7II-XVIII-107.8. \nolinkurl{fr/sga7ii/diplomatic/expose_XVIII_body.tex}; \nolinkurl{fr/sga7ii/current/expose_XVIII_body.tex}\par

\par\medskip\noindent\hypertarget{app-F7II-XVIII-107.9}{}\textbf{F7II-XVIII-107.9}\quad p. 324; designated-witness attribution.\par
\noindent Imprimé : {\ttfamily\footnotesize définit\allowbreak{}}\par
\noindent Corrigé : {\ttfamily\footnotesize définisse\allowbreak{}}\par
\noindent Source : WITNESS-F7II-XVIII-107.9. \nolinkurl{fr/sga7ii/diplomatic/expose_XVIII_body.tex}; \nolinkurl{fr/sga7ii/current/expose_XVIII_body.tex}\par

\par\medskip\noindent\hypertarget{app-F7II-XVIII-107.10}{}\textbf{F7II-XVIII-107.10}\quad p. 324; designated-witness attribution.\par
\noindent Imprimé : \emph{(vide)}\par
\noindent Corrigé : {\ttfamily\footnotesize \_\{\textbackslash{}bar\textbackslash{}eta\}\allowbreak{}}\par
\noindent Source : WITNESS-F7II-XVIII-107.10. \nolinkurl{fr/sga7ii/diplomatic/expose_XVIII_body.tex}; \nolinkurl{fr/sga7ii/current/expose_XVIII_body.tex}\par

\par\medskip\noindent\hypertarget{app-F7II-XVIII-107.11}{}\textbf{F7II-XVIII-107.11}\quad p. 324; designated-witness attribution.\par
\noindent Imprimé : {\ttfamily\footnotesize  \textbackslash{}\allowbreak{} }\par
\noindent Corrigé : {\ttfamily\footnotesize )\allowbreak{}}\par
\noindent Source : WITNESS-F7II-XVIII-107.11. \nolinkurl{fr/sga7ii/diplomatic/expose_XVIII_body.tex}; \nolinkurl{fr/sga7ii/current/expose_XVIII_body.tex}\par

\par\medskip\noindent\hypertarget{app-F7II-XVIII-107.12}{}\textbf{F7II-XVIII-107.12}\quad p. 324; designated-witness attribution.\par
\noindent Imprimé : {\ttfamily\footnotesize \textbackslash{}\allowbreak{} }\par
\noindent Corrigé : {\ttfamily\footnotesize \textbackslash{}qquad\allowbreak{} }\par
\noindent Source : WITNESS-F7II-XVIII-107.12. \nolinkurl{fr/sga7ii/diplomatic/expose_XVIII_body.tex}; \nolinkurl{fr/sga7ii/current/expose_XVIII_body.tex}\par

\par\medskip\noindent\hypertarget{app-F7II-XVIII-107.13}{}\textbf{F7II-XVIII-107.13}\quad p. 324; inherited corrected-layer difference.\par
\noindent Imprimé : \emph{(vide)}\par
\noindent Corrigé : {\ttfamily\footnotesize ,\allowbreak{}}\par
\noindent Source : ADOPT-F7II-XVIII-107.13. \nolinkurl{fr/sga7ii/diplomatic/expose_XVIII_body.tex}; \nolinkurl{fr/sga7ii/current/expose_XVIII_body.tex}\par

\par\medskip\noindent\hypertarget{app-F7II-XVIII-107.14}{}\textbf{F7II-XVIII-107.14}\quad p. 324; inherited corrected-layer difference.\par
\noindent Imprimé : {\ttfamily\footnotesize  \textbackslash{}\allowbreak{} }\par
\noindent Corrigé : \emph{(vide)}\par
\noindent Source : ADOPT-F7II-XVIII-107.14. \nolinkurl{fr/sga7ii/diplomatic/expose_XVIII_body.tex}; \nolinkurl{fr/sga7ii/current/expose_XVIII_body.tex}\par

\par\medskip\noindent\hypertarget{app-F7II-XVIII-107.15}{}\textbf{F7II-XVIII-107.15}\quad p. 324; inherited corrected-layer difference.\par
\noindent Imprimé : {\ttfamily\footnotesize \_\{1\}\allowbreak{}}\par
\noindent Corrigé : \emph{(vide)}\par
\noindent Source : ADOPT-F7II-XVIII-107.15. \nolinkurl{fr/sga7ii/diplomatic/expose_XVIII_body.tex}; \nolinkurl{fr/sga7ii/current/expose_XVIII_body.tex}\par

\par\medskip\noindent\hypertarget{app-F7II-XVIII-107.16}{}\textbf{F7II-XVIII-107.16}\quad p. 324; printed misprint.\par
\noindent Imprimé : \emph{(vide)}\par
\noindent Corrigé : {\ttfamily\footnotesize )\allowbreak{}}\par
\noindent Source : 7II:561. \nolinkurl{fixes/MISPRINTS_7II.json#561}\par

\par\medskip\noindent\hypertarget{app-F7II-XVIII-107.17}{}\textbf{F7II-XVIII-107.17}\quad p. 324; printed misprint.\par
\noindent Imprimé : \emph{(vide)}\par
\noindent Corrigé : {\ttfamily\footnotesize )\allowbreak{}}\par
\noindent Source : 7II:562. \nolinkurl{fixes/MISPRINTS_7II.json#562}\par

\par\medskip\noindent\hypertarget{app-F7II-XVIII-107.18}{}\textbf{F7II-XVIII-107.18}\quad p. 324; inherited corrected-layer difference.\par
\noindent Imprimé : \emph{(1 espace(s), 0 retour(s) de ligne dans le témoin TeX)}\par
\noindent Corrigé : \emph{(vide)}\par
\noindent Source : ADOPT-F7II-XVIII-107.18. \nolinkurl{fr/sga7ii/diplomatic/expose_XVIII_body.tex}; \nolinkurl{fr/sga7ii/current/expose_XVIII_body.tex}\par

\par\medskip\noindent\hypertarget{app-F7II-XVIII-108.1}{}\textbf{F7II-XVIII-108.1}\quad p. 325; printed misprint.\par
\noindent Imprimé : \emph{(vide)}\par
\noindent Corrigé : {\ttfamily\footnotesize )\allowbreak{}}\par
\noindent Source : 7II:563. \nolinkurl{fixes/MISPRINTS_7II.json#563}\par

\par\medskip\noindent\hypertarget{app-F7II-XVIII-108.2}{}\textbf{F7II-XVIII-108.2}\quad p. 325; designated-witness attribution.\par
\noindent Imprimé : \emph{(vide)}\par
\noindent Corrigé : {\ttfamily\footnotesize \textbackslash{}bigl\allowbreak{}}\par
\noindent Source : WITNESS-F7II-XVIII-108.2. \nolinkurl{fr/sga7ii/diplomatic/expose_XVIII_body.tex}; \nolinkurl{fr/sga7ii/current/expose_XVIII_body.tex}\par

\par\medskip\noindent\hypertarget{app-F7II-XVIII-108.3}{}\textbf{F7II-XVIII-108.3}\quad p. 325; printed misprint.\par
\noindent Imprimé : {\ttfamily\footnotesize i\allowbreak{}}\par
\noindent Corrigé : {\ttfamily\footnotesize m\allowbreak{}}\par
\noindent Source : 7II:564. \nolinkurl{fixes/MISPRINTS_7II.json#564}\par

\par\medskip\noindent\hypertarget{app-F7II-XVIII-108.4}{}\textbf{F7II-XVIII-108.4}\quad p. 325; designated-witness attribution.\par
\noindent Imprimé : \emph{(vide)}\par
\noindent Corrigé : {\ttfamily\footnotesize  \textbackslash{}longrightarrow\allowbreak{} H\textasciicircum{}\{n-1\}(X\_\{\textbackslash{}bar\textbackslash{}eta\}\allowbreak{},\textbackslash{}mathbb\{Q\}\_\{\textbackslash{}ell\}(m\allowbreak{}))\textbackslash{}bigr\allowbreak{}}\par
\noindent Source : WITNESS-F7II-XVIII-108.4. \nolinkurl{fr/sga7ii/diplomatic/expose_XVIII_body.tex}; \nolinkurl{fr/sga7ii/current/expose_XVIII_body.tex}\par

\par\medskip\noindent\hypertarget{app-F7II-XVIII-108.5}{}\textbf{F7II-XVIII-108.5}\quad p. 325; designated-witness attribution.\par
\noindent Imprimé : {\ttfamily\footnotesize ,\allowbreak{}}\par
\noindent Corrigé : \emph{(vide)}\par
\noindent Source : WITNESS-F7II-XVIII-108.5. \nolinkurl{fr/sga7ii/diplomatic/expose_XVIII_body.tex}; \nolinkurl{fr/sga7ii/current/expose_XVIII_body.tex}\par

\par\medskip\noindent\hypertarget{app-F7II-XVIII-108.6}{}\textbf{F7II-XVIII-108.6}\quad p. 325; designated-witness attribution.\par
\noindent Imprimé : \emph{(1 espace(s), 0 retour(s) de ligne dans le témoin TeX)}\par
\noindent Corrigé : \emph{(vide)}\par
\noindent Source : WITNESS-F7II-XVIII-108.6. \nolinkurl{fr/sga7ii/diplomatic/expose_XVIII_body.tex}; \nolinkurl{fr/sga7ii/current/expose_XVIII_body.tex}\par

\par\medskip\noindent\hypertarget{app-F7II-XVIII-108.7}{}\textbf{F7II-XVIII-108.7}\quad p. 325; printed misprint.\par
\noindent Imprimé : {\ttfamily\footnotesize a\allowbreak{}}\par
\noindent Corrigé : {\ttfamily\footnotesize à\allowbreak{}}\par
\noindent Source : 7II:565. \nolinkurl{fixes/MISPRINTS_7II.json#565}\par

\par\medskip\noindent\hypertarget{app-F7II-XVIII-108.8}{}\textbf{F7II-XVIII-108.8}\quad p. 325; designated-witness attribution.\par
\noindent Imprimé : \emph{(vide)}\par
\noindent Corrigé : {\ttfamily\footnotesize \textbackslash{}check\{\allowbreak{}}\par
\noindent Source : WITNESS-F7II-XVIII-108.8. \nolinkurl{fr/sga7ii/diplomatic/expose_XVIII_body.tex}; \nolinkurl{fr/sga7ii/current/expose_XVIII_body.tex}\par

\par\medskip\noindent\hypertarget{app-F7II-XVIII-108.9}{}\textbf{F7II-XVIII-108.9}\quad p. 325; designated-witness attribution.\par
\noindent Imprimé : \emph{(vide)}\par
\noindent Corrigé : {\ttfamily\footnotesize \}\allowbreak{}}\par
\noindent Source : WITNESS-F7II-XVIII-108.9. \nolinkurl{fr/sga7ii/diplomatic/expose_XVIII_body.tex}; \nolinkurl{fr/sga7ii/current/expose_XVIII_body.tex}\par

\par\medskip\noindent\hypertarget{app-F7II-XVIII-108.10}{}\textbf{F7II-XVIII-108.10}\quad p. 325; designated-witness attribution.\par
\noindent Imprimé : \emph{(vide)}\par
\noindent Corrigé : {\ttfamily\footnotesize  \textbackslash{}check\allowbreak{}}\par
\noindent Source : WITNESS-F7II-XVIII-108.10. \nolinkurl{fr/sga7ii/diplomatic/expose_XVIII_body.tex}; \nolinkurl{fr/sga7ii/current/expose_XVIII_body.tex}\par

\par\medskip\noindent\hypertarget{app-F7II-XVIII-109.1}{}\textbf{F7II-XVIII-109.1}\quad p. 326; printed misprint.\par
\noindent Imprimé : \emph{(vide)}\par
\noindent Corrigé : {\ttfamily\footnotesize )\allowbreak{}}\par
\noindent Source : 7II:567; 7II:968. \nolinkurl{fixes/MISPRINTS_7II.json#567}; \nolinkurl{fixes/MISPRINTS_7II.json#968}\par

\par\medskip\noindent\hypertarget{app-F7II-XVIII-109.2}{}\textbf{F7II-XVIII-109.2}\quad p. 326; designated-witness attribution.\par
\noindent Imprimé : \emph{(vide)}\par
\noindent Corrigé : {\ttfamily\footnotesize \}\allowbreak{}}\par
\noindent Source : WITNESS-F7II-XVIII-109.2. \nolinkurl{fr/sga7ii/diplomatic/expose_XVIII_body.tex}; \nolinkurl{fr/sga7ii/current/expose_XVIII_body.tex}\par

\par\medskip\noindent\hypertarget{app-F7II-XVIII-109.3}{}\textbf{F7II-XVIII-109.3}\quad p. 326; designated-witness attribution.\par
\noindent Imprimé : {\ttfamily\footnotesize \}\allowbreak{}}\par
\noindent Corrigé : \emph{(vide)}\par
\noindent Source : WITNESS-F7II-XVIII-109.3. \nolinkurl{fr/sga7ii/diplomatic/expose_XVIII_body.tex}; \nolinkurl{fr/sga7ii/current/expose_XVIII_body.tex}\par

\par\medskip\noindent\hypertarget{app-F7II-XVIII-109.4}{}\textbf{F7II-XVIII-109.4}\quad p. 326; designated-witness attribution.\par
\noindent Imprimé : \emph{(vide)}\par
\noindent Corrigé : {\ttfamily\footnotesize -espace\allowbreak{}}\par
\noindent Source : WITNESS-F7II-XVIII-109.4. \nolinkurl{fr/sga7ii/diplomatic/expose_XVIII_body.tex}; \nolinkurl{fr/sga7ii/current/expose_XVIII_body.tex}\par

\par\medskip\noindent\hypertarget{app-F7II-XVIII-109.5}{}\textbf{F7II-XVIII-109.5}\quad p. 326; designated-witness attribution.\par
\noindent Imprimé : \emph{(vide)}\par
\noindent Corrigé : {\ttfamily\footnotesize \}\allowbreak{}}\par
\noindent Source : WITNESS-F7II-XVIII-109.5. \nolinkurl{fr/sga7ii/diplomatic/expose_XVIII_body.tex}; \nolinkurl{fr/sga7ii/current/expose_XVIII_body.tex}\par

\par\medskip\noindent\hypertarget{app-F7II-XVIII-109.6}{}\textbf{F7II-XVIII-109.6}\quad p. 326; designated-witness attribution.\par
\noindent Imprimé : {\ttfamily\footnotesize \}\allowbreak{}}\par
\noindent Corrigé : \emph{(vide)}\par
\noindent Source : WITNESS-F7II-XVIII-109.6. \nolinkurl{fr/sga7ii/diplomatic/expose_XVIII_body.tex}; \nolinkurl{fr/sga7ii/current/expose_XVIII_body.tex}\par

\par\medskip\noindent\hypertarget{app-F7II-XVIII-109.7}{}\textbf{F7II-XVIII-109.7}\quad p. 326; designated-witness attribution.\par
\noindent Imprimé : {\ttfamily\footnotesize module\allowbreak{}}\par
\noindent Corrigé : {\ttfamily\footnotesize espace\allowbreak{}}\par
\noindent Source : WITNESS-F7II-XVIII-109.7. \nolinkurl{fr/sga7ii/diplomatic/expose_XVIII_body.tex}; \nolinkurl{fr/sga7ii/current/expose_XVIII_body.tex}\par

\par\medskip\noindent\hypertarget{app-F7II-XVIII-109.8}{}\textbf{F7II-XVIII-109.8}\quad p. 326; designated-witness attribution.\par
\noindent Imprimé : \emph{(1 espace(s), 0 retour(s) de ligne dans le témoin TeX)}\par
\noindent Corrigé : \emph{(vide)}\par
\noindent Source : WITNESS-F7II-XVIII-109.8. \nolinkurl{fr/sga7ii/diplomatic/expose_XVIII_body.tex}; \nolinkurl{fr/sga7ii/current/expose_XVIII_body.tex}\par

\par\medskip\noindent\hypertarget{app-F7II-XVIII-109.9}{}\textbf{F7II-XVIII-109.9}\quad p. 326; printed misprint.\par
\noindent Imprimé : {\ttfamily\footnotesize \textbackslash{}chi\allowbreak{}}\par
\noindent Corrigé : {\ttfamily\footnotesize X\allowbreak{}}\par
\noindent Source : 7II:568; 7II:971. \nolinkurl{fixes/MISPRINTS_7II.json#568}; \nolinkurl{fixes/MISPRINTS_7II.json#971}\par

\par\medskip\noindent\hypertarget{app-F7II-XVIII-109.10}{}\textbf{F7II-XVIII-109.10}\quad p. 326; inherited corrected-layer difference.\par
\noindent Imprimé : {\ttfamily\footnotesize ,\allowbreak{}}\par
\noindent Corrigé : \emph{(vide)}\par
\noindent Source : ADOPT-F7II-XVIII-109.10. \nolinkurl{fr/sga7ii/diplomatic/expose_XVIII_body.tex}; \nolinkurl{fr/sga7ii/current/expose_XVIII_body.tex}\par

\par\medskip\noindent\hypertarget{app-F7II-XVIII-109.11}{}\textbf{F7II-XVIII-109.11}\quad p. 326; inherited corrected-layer difference.\par
\noindent Imprimé : {\ttfamily\footnotesize donc\allowbreak{}}\par
\noindent Corrigé : {\ttfamily\footnotesize ;\allowbreak{}}\par
\noindent Source : ADOPT-F7II-XVIII-109.11. \nolinkurl{fr/sga7ii/diplomatic/expose_XVIII_body.tex}; \nolinkurl{fr/sga7ii/current/expose_XVIII_body.tex}\par

\par\medskip\noindent\hypertarget{app-F7II-XVIII-109.12}{}\textbf{F7II-XVIII-109.12}\quad p. 326; designated-witness attribution.\par
\noindent Imprimé : {\ttfamily\footnotesize en\allowbreak{}}\par
\noindent Corrigé : {\ttfamily\footnotesize par\allowbreak{}}\par
\noindent Source : WITNESS-F7II-XVIII-109.12. \nolinkurl{fr/sga7ii/diplomatic/expose_XVIII_body.tex}; \nolinkurl{fr/sga7ii/current/expose_XVIII_body.tex}\par

\par\medskip\noindent\hypertarget{app-F7II-XVIII-109.13}{}\textbf{F7II-XVIII-109.13}\quad p. 326; inherited corrected-layer difference.\par
\noindent Imprimé : {\ttfamily\footnotesize vertu\allowbreak{} de\allowbreak{} (\allowbreak{}}\par
\noindent Corrigé : \emph{(vide)}\par
\noindent Source : ADOPT-F7II-XVIII-109.13. \nolinkurl{fr/sga7ii/diplomatic/expose_XVIII_body.tex}; \nolinkurl{fr/sga7ii/current/expose_XVIII_body.tex}\par

\par\medskip\noindent\hypertarget{app-F7II-XVIII-109.14}{}\textbf{F7II-XVIII-109.14}\quad p. 326; inherited corrected-layer difference.\par
\noindent Imprimé : {\ttfamily\footnotesize )\allowbreak{}}\par
\noindent Corrigé : \emph{(vide)}\par
\noindent Source : ADOPT-F7II-XVIII-109.14. \nolinkurl{fr/sga7ii/diplomatic/expose_XVIII_body.tex}; \nolinkurl{fr/sga7ii/current/expose_XVIII_body.tex}\par

\par\medskip\noindent\hypertarget{app-F7II-XVIII-109.15}{}\textbf{F7II-XVIII-109.15}\quad p. 326; inherited corrected-layer difference.\par
\noindent Imprimé : \emph{(vide)}\par
\noindent Corrigé : {\ttfamily\footnotesize  donc\allowbreak{}}\par
\noindent Source : ADOPT-F7II-XVIII-109.15. \nolinkurl{fr/sga7ii/diplomatic/expose_XVIII_body.tex}; \nolinkurl{fr/sga7ii/current/expose_XVIII_body.tex}\par

\par\medskip\noindent\hypertarget{app-F7II-XVIII-109.16}{}\textbf{F7II-XVIII-109.16}\quad p. 326; inherited corrected-layer difference.\par
\noindent Imprimé : {\ttfamily\footnotesize x\allowbreak{}}\par
\noindent Corrigé : {\ttfamily\footnotesize \textbackslash{}chi\allowbreak{}}\par
\noindent Source : ADOPT-F7II-XVIII-109.16. \nolinkurl{fr/sga7ii/diplomatic/expose_XVIII_body.tex}; \nolinkurl{fr/sga7ii/current/expose_XVIII_body.tex}\par

\par\medskip\noindent\hypertarget{app-F7II-XVIII-109.17}{}\textbf{F7II-XVIII-109.17}\quad p. 326; inherited corrected-layer difference.\par
\noindent Imprimé : {\ttfamily\footnotesize x\allowbreak{}}\par
\noindent Corrigé : {\ttfamily\footnotesize \textbackslash{}chi\allowbreak{}}\par
\noindent Source : ADOPT-F7II-XVIII-109.17. \nolinkurl{fr/sga7ii/diplomatic/expose_XVIII_body.tex}; \nolinkurl{fr/sga7ii/current/expose_XVIII_body.tex}\par

\par\medskip\noindent\hypertarget{app-F7II-XVIII-109.18}{}\textbf{F7II-XVIII-109.18}\quad p. 326; inherited corrected-layer difference.\par
\noindent Imprimé : {\ttfamily\footnotesize x\allowbreak{}}\par
\noindent Corrigé : {\ttfamily\footnotesize \textbackslash{}chi\allowbreak{}}\par
\noindent Source : ADOPT-F7II-XVIII-109.18. \nolinkurl{fr/sga7ii/diplomatic/expose_XVIII_body.tex}; \nolinkurl{fr/sga7ii/current/expose_XVIII_body.tex}\par

\par\medskip\noindent\hypertarget{app-F7II-XVIII-110.1}{}\textbf{F7II-XVIII-110.1}\quad p. 327; inherited corrected-layer difference.\par
\noindent Imprimé : \emph{(vide)}\par
\noindent Corrigé : {\ttfamily\footnotesize \_i\allowbreak{}}\par
\noindent Source : ADOPT-F7II-XVIII-110.1. \nolinkurl{fr/sga7ii/diplomatic/expose_XVIII_body.tex}; \nolinkurl{fr/sga7ii/current/expose_XVIII_body.tex}\par

\par\medskip\noindent\hypertarget{app-F7II-XVIII-110.2}{}\textbf{F7II-XVIII-110.2}\quad p. 327; inherited corrected-layer difference.\par
\noindent Imprimé : \emph{(vide)}\par
\noindent Corrigé : {\ttfamily\footnotesize \_i\allowbreak{}}\par
\noindent Source : ADOPT-F7II-XVIII-110.2. \nolinkurl{fr/sga7ii/diplomatic/expose_XVIII_body.tex}; \nolinkurl{fr/sga7ii/current/expose_XVIII_body.tex}\par

\par\medskip\noindent\hypertarget{app-F7II-XVIII-110.3}{}\textbf{F7II-XVIII-110.3}\quad p. 327; inherited corrected-layer difference.\par
\noindent Imprimé : \emph{(vide)}\par
\noindent Corrigé : {\ttfamily\footnotesize ,\allowbreak{}}\par
\noindent Source : ADOPT-F7II-XVIII-110.3. \nolinkurl{fr/sga7ii/diplomatic/expose_XVIII_body.tex}; \nolinkurl{fr/sga7ii/current/expose_XVIII_body.tex}\par

\par\medskip\noindent\hypertarget{app-F7II-XVIII-110.4}{}\textbf{F7II-XVIII-110.4}\quad p. 327; designated-witness attribution.\par
\noindent Imprimé : {\ttfamily\footnotesize .\allowbreak{}}\par
\noindent Corrigé : {\ttfamily\footnotesize -\allowbreak{}}\par
\noindent Source : WITNESS-F7II-XVIII-110.4. \nolinkurl{fr/sga7ii/diplomatic/expose_XVIII_body.tex}; \nolinkurl{fr/sga7ii/current/expose_XVIII_body.tex}\par

\par\medskip\noindent\hypertarget{app-F7II-XVIII-110.5}{}\textbf{F7II-XVIII-110.5}\quad p. 327; designated-witness attribution.\par
\noindent Imprimé : \emph{(vide)}\par
\noindent Corrigé : {\ttfamily\footnotesize ,\allowbreak{}}\par
\noindent Source : WITNESS-F7II-XVIII-110.5. \nolinkurl{fr/sga7ii/diplomatic/expose_XVIII_body.tex}; \nolinkurl{fr/sga7ii/current/expose_XVIII_body.tex}\par

\par\medskip\noindent\hypertarget{app-F7II-XVIII-110.6}{}\textbf{F7II-XVIII-110.6}\quad p. 327; designated-witness attribution.\par
\noindent Imprimé : {\ttfamily\footnotesize x\allowbreak{}}\par
\noindent Corrigé : {\ttfamily\footnotesize \textbackslash{}chi\allowbreak{}}\par
\noindent Source : WITNESS-F7II-XVIII-110.6. \nolinkurl{fr/sga7ii/diplomatic/expose_XVIII_body.tex}; \nolinkurl{fr/sga7ii/current/expose_XVIII_body.tex}\par

\par\medskip\noindent\hypertarget{app-F7II-XVIII-110.7}{}\textbf{F7II-XVIII-110.7}\quad p. 327; designated-witness attribution.\par
\noindent Imprimé : {\ttfamily\footnotesize \textbackslash{}\allowbreak{} \textbackslash{}\allowbreak{} X\allowbreak{}}\par
\noindent Corrigé : {\ttfamily\footnotesize \textbackslash{}chi\allowbreak{}}\par
\noindent Source : WITNESS-F7II-XVIII-110.7. \nolinkurl{fr/sga7ii/diplomatic/expose_XVIII_body.tex}; \nolinkurl{fr/sga7ii/current/expose_XVIII_body.tex}\par

\par\medskip\noindent\hypertarget{app-F7II-XVIII-110.8}{}\textbf{F7II-XVIII-110.8}\quad p. 327; designated-witness attribution.\par
\noindent Imprimé : {\ttfamily\footnotesize x\allowbreak{}}\par
\noindent Corrigé : {\ttfamily\footnotesize \textbackslash{}chi\allowbreak{}}\par
\noindent Source : WITNESS-F7II-XVIII-110.8. \nolinkurl{fr/sga7ii/diplomatic/expose_XVIII_body.tex}; \nolinkurl{fr/sga7ii/current/expose_XVIII_body.tex}\par

\par\medskip\noindent\hypertarget{app-F7II-XVIII-110.9}{}\textbf{F7II-XVIII-110.9}\quad p. 327; designated-witness attribution.\par
\noindent Imprimé : \emph{(vide)}\par
\noindent Corrigé : {\ttfamily\footnotesize -espace\allowbreak{}}\par
\noindent Source : WITNESS-F7II-XVIII-110.9. \nolinkurl{fr/sga7ii/diplomatic/expose_XVIII_body.tex}; \nolinkurl{fr/sga7ii/current/expose_XVIII_body.tex}\par

\par\medskip\noindent\hypertarget{app-F7II-XVIII-110.10}{}\textbf{F7II-XVIII-110.10}\quad p. 327; printed misprint.\par
\noindent Imprimé : \emph{(vide)}\par
\noindent Corrigé : {\ttfamily\footnotesize \textbackslash{}mathbb\{\allowbreak{}}\par
\noindent Source : 7II:974. \nolinkurl{fixes/MISPRINTS_7II.json#974}\par

\par\medskip\noindent\hypertarget{app-F7II-XVIII-110.11}{}\textbf{F7II-XVIII-110.11}\quad p. 327; printed misprint.\par
\noindent Imprimé : \emph{(vide)}\par
\noindent Corrigé : {\ttfamily\footnotesize \}\allowbreak{}}\par
\noindent Source : 7II:974. \nolinkurl{fixes/MISPRINTS_7II.json#974}\par

\par\medskip\noindent\hypertarget{app-F7II-XVIII-110.12}{}\textbf{F7II-XVIII-110.12}\quad p. 327; designated-witness attribution.\par
\noindent Imprimé : {\ttfamily\footnotesize module\allowbreak{}}\par
\noindent Corrigé : {\ttfamily\footnotesize vectoriel\allowbreak{}}\par
\noindent Source : WITNESS-F7II-XVIII-110.12. \nolinkurl{fr/sga7ii/diplomatic/expose_XVIII_body.tex}; \nolinkurl{fr/sga7ii/current/expose_XVIII_body.tex}\par

\par\medskip\noindent\hypertarget{app-F7II-XVIII-110.13}{}\textbf{F7II-XVIII-110.13}\quad p. 327; misprint\_pass.\par
\noindent Imprimé : \emph{(vide)}\par
\noindent Corrigé : {\ttfamily\footnotesize s'\allowbreak{}}\par
\noindent Source : 7II:572. \nolinkurl{fixes/FIXLOG_7II_CURRENT.jsonl#761}\par

\par\medskip\noindent\hypertarget{app-F7II-XVIII-110.14}{}\textbf{F7II-XVIII-110.14}\quad p. 327; inherited corrected-layer difference.\par
\noindent Imprimé : {\ttfamily\footnotesize x\allowbreak{}}\par
\noindent Corrigé : {\ttfamily\footnotesize \textbackslash{}chi\allowbreak{}}\par
\noindent Source : ADOPT-F7II-XVIII-110.14. \nolinkurl{fr/sga7ii/diplomatic/expose_XVIII_body.tex}; \nolinkurl{fr/sga7ii/current/expose_XVIII_body.tex}\par

\par\medskip\noindent\hypertarget{app-F7II-XVIII-110.15}{}\textbf{F7II-XVIII-110.15}\quad p. 327; inherited corrected-layer difference.\par
\noindent Imprimé : {\ttfamily\footnotesize x\allowbreak{}}\par
\noindent Corrigé : {\ttfamily\footnotesize \textbackslash{}chi\allowbreak{}}\par
\noindent Source : ADOPT-F7II-XVIII-110.15. \nolinkurl{fr/sga7ii/diplomatic/expose_XVIII_body.tex}; \nolinkurl{fr/sga7ii/current/expose_XVIII_body.tex}\par

\par\medskip\noindent\hypertarget{app-F7II-XVIII-111.1}{}\textbf{F7II-XVIII-111.1}\quad p. 327; printed misprint.\par
\noindent Imprimé : {\ttfamily\footnotesize Gauthiers\allowbreak{}}\par
\noindent Corrigé : {\ttfamily\footnotesize Gauthier\allowbreak{}}\par
\noindent Source : 7II:573; 7II:975. \nolinkurl{fixes/MISPRINTS_7II.json#573}; \nolinkurl{fixes/MISPRINTS_7II.json#975}\par

\par\medskip\noindent\hypertarget{app-F7II-XVIII-111.2}{}\textbf{F7II-XVIII-111.2}\quad p. 327; printed misprint.\par
\noindent Imprimé : {\ttfamily\footnotesize 15\allowbreak{}}\par
\noindent Corrigé : {\ttfamily\footnotesize 90\allowbreak{}}\par
\noindent Source : 7II:574; 7II:977. \nolinkurl{fixes/MISPRINTS_7II.json#574}; \nolinkurl{fixes/MISPRINTS_7II.json#977}\par

\normalsize\fi
\clearpage\SGAtoc{Exposé XIX — Le théorème de Noether}
\selectlanguage{french}

\SGApage{II}{328}


S G A \quad 7


\underline{Exposé XIX}


\underline{\underline{Le théorème de Noether}}


\underline{par \quad P. Deligne}


\underline{Sommaire}

1. Enoncé des théorèmes \hfill 1

2. Complément à XVIII \hfill 4

3. Preuve de 1.3 \hfill 5

4. Deuxième complément à XVIII \hfill 8

5. Les cas d'exception au théorème de Noether. \hfill 11


\SGAsectiontoc{1. Enoncé des théorèmes}{1}
\underline{1. Enoncé des théorèmes}

Le présent exposé est une mise au goût du jour de la démonstration de Lefschetz [3] page 108 du théorème de Noether (1.2) ci-dessous.

1.1 \qquad Soit $k$ un corps algébriquement clos de caractéristique $p$ . Soient $d \geq 1$ et $\underline{a} = (a_i)_{1 \leq i \leq d}$ une suite d'entiers $\geq 1$ . Une intersection complète $V$ de dimension $n$ et de multidegré $\underline{a}$ dans l'espace projectif $\mathbb{P}^{n+d}$ sur $k$ sera dite \underline{générique} si, sur une extension de $k$ , elle est projectivement isomorphe à une intersection complète définie par des équations dont les coefficients sont algébriquement indépendants sur le corps premier. Dans ce qui suit et sauf mention expresse du contraire, \underline{on suppose que} $d \geq 1$ et que $a_i \geq 2$ .

\underline{Théorème} 1.2 \qquad \underline{Soit} $V \subset \mathbb{P}^{2+d}$ \underline{une surface intersection complète générique de multidegré} $\underline{a}$ . \underline{On suppose que} $\underline{a} \neq (2), (3), (2, 2)$ . \underline{Alors, toute courbe (réduite) tracée sur} $V$ \underline{est l'intersection (schématique) de} $V$ \underline{et d'une hypersurface de} $\mathbb{P}^{2+d}$ .

Montrons que le théorème résulte de son corollaire 


\SGApage{II}{329}

\begin{equation}\tag{1.2.1}
 \mathrm{Pic}\ (\mathbb{P}^{2+d}) = \mathbb{Z} \overset{\sim}{\longrightarrow} \mathrm{Pic}(V)\ .
\end{equation}
Si $C$ est une courbe tracée sur $V$ , (1.2.1) implique que $\mathcal{O}(C) \simeq \mathcal{O}(n)$ pour $n$ convenable; si $H$ est une section hyperplane, $n = (\Pi\ a_i)^{-1}.(C, H) > 0$ . Il reste à prouver que le système linéaire induit sur $V$ par $|n\ H|$ est complet. Ecrivons $V$ comme intersection d'hypersurfaces $H_i$ de degré $a_i$ et soit $V_i = H_1 \cap \ldots \cap H_i$ . On vérifie par récurrence sur $i$ que le système linéaire induit par $|n\ H|$ sur $V_i$ est complet: dans la suite exacte

\begin{equation*}
H^0(V_i,\ \mathcal{O}(n)) \longrightarrow H^0(V_{i+1},\ \mathcal{O}(n)) \longrightarrow H^1(V_i,\ \mathcal{O}(n - a_{i+1}))\ ,
\end{equation*}
le dernier groupe est nul (FAC) .

\underline{Théorème} 1.3 \qquad \underline{Soit} $V \subset \mathbb{P}^{2n+d}$ \underline{une intersection complète de dimension paire} $2n \neq 0$ \underline{et de multidegré} $\underline{a}$ \underline{générique.} \underline{On suppose vérifiée l'une des conditions suivantes}

(a) \quad $k$ \underline{n'est pas de caractéristique} $2$ , \underline{et} $(2n; \underline{a})$ \underline{est distinct de} $(2n; 2)$, $(2n; 2, 2)$ , $(2; 3)$ .

(b) \quad \underline{le nombre de Hodge} $h^{2n,0}$ \underline{est non nul,} \underline{i.e.} $\Sigma\ (a_i-1) > 2n$ .

\underline{Alors, toute classe de cohomologie algébrique dans} $H^{2n}(V,\ \mathbb{Q}_\ell(n))$ \underline{est un multiple rationnel de la classe} $\eta^n$ \underline{des sections linéaires.}

\SGAvar{Déduisons 1.2.1 de 1.3. Les cas d'exceptions dans 1.2 sont exactement ceux pour lesquels $\Sigma(d_i -1) \leq 2$ ; on peut donc appliquer 1.3 (b) à la surface $V$ de 1.2. On sait que $\mathrm{Pic}^{\tau}(V) = 0$ (XI 1.8); donc que $\mathrm{Pic}(V)$ s'injecte dans $H^2(V,\ \mathbb{Q}_\ell(1))$ . D'après 1.3, tout élément de $\mathrm{Pic}(V)$ est un multiple rationnel de la classe $\eta$ de $\mathcal{O}(1)$ , et on conclut par (XI 1.8).

}{Déduisons 1.2.1 de 1.3. Les cas d'exceptions dans 1.2 sont exactement ceux pour lesquels $\Sigma(a_i -1) \leq 2$ ; on peut donc appliquer 1.3 (b) à la surface $V$ de 1.2. On sait que $\mathrm{Pic}^{\tau}(V) = 0$ (XI 1.8); donc que $\mathrm{Pic}(V)$ s'injecte dans $H^2(V,\ \mathbb{Q}_\ell(1))$ . D'après 1.3, tout élément de $\mathrm{Pic}(V)$ est un multiple rationnel de la classe $\eta$ de $\mathcal{O}(1)$ , et on conclut par (XI 1.8).

}{F7II-XIX-001}Il est clair que la validité de 1.2 et 1.3 ne dépend pas de l'intersection complète générique choisie, mais seulement de $p, n$ et $\underline{a}$ .


\SGApage{II}{330}

\underline{Théorème} 1.4. \quad \underline{Soit} $X \subset \mathbb{P}$ \underline{une variété projective lisse purement de dimension} $2n + 1$ \underline{sur} $k$ . \underline{Soit} $\ell \subset \check{\mathbb{P}}$ \underline{un pinceau de Lefschetz de sections hyperplanes de} $X$ . \underline{Pour} $k$ \underline{de caractéristique} 2, \underline{on suppose que} $\ell$ \underline{est même un pinceau de Lefschetz générique.}

\underline{On suppose que} $X$ \underline{vérifie} (LV)  (XVIII 5.2.2).

\underline{Soit} $V$ \underline{la section hyperplane générique dans le pinceau} $\ell$ . \underline{L'une des deux conditions suivantes est vérifiée.}

(a) \underline{Aucune classe de cohomologie} $\neq 0$ \underline{dans la partie évanescente} $\mathrm{Ev}(V, \mathbb{Q}_\ell(n))$ \underline{du groupe de cohomologie} $H^{2n}(V, \mathbb{Q}_\ell(n))$ \underline{n'est algébrique.}

(b) \underline{Les classes de cohomologie algébrique engendrent} $\mathrm{Ev}(V, \mathbb{Q}_\ell(n))$ .

Précisons le sens du mot générique dans 1.4.

\SGAvar{(a) Soit $G$ la grassmanienne des droites de l'espace projectif $\check{\mathbb{P}}$ dual de $\mathbb{P}$ . Appelons ``point géométrique'' de $G$ un morphisme du spectre d'un corps algébriquement clos dans $G$ .

}{(a) Soit $G$ la grassmannienne des droites de l'espace projectif $\check{\mathbb{P}}$ dual de $\mathbb{P}$ . Appelons ``point géométrique'' de $G$ un morphisme du spectre d'un corps algébriquement clos dans $G$ .

}{F7II-XIX-002}Tout point géométrique $\bar{s} \to G$ définit un pinceau de sections hyperplanes $\ell$ de la variété déduite de $X$ par extension des scalaires de $k$ à $k(\bar{s})$ . On dit que $\ell$ est générique si $\bar{s}$ est localisé au point générique de $G$ .

(b) Soit $\ell$ un pinceau de sections hyperplanes (une droite de $\check{\mathbb{P}}$ ). Tout point géométrique $\bar{s} \to \ell$ définit une section hyperplane de la variété déduite de $X$ par extension des scalaires de $k$ à $k(\bar{s})$ . On dit qu'elle est générique si $\bar{s}$ est localisé au point générique de $\ell$ .

Soient $\bar{s}$ comme ci-dessus, d'image dans $\ell$ le point générique $s$ et $S \subset \ell$ l'ensemble fini des points de $\ell$ donnant lieu à une section hyperplane singulière. Le groupe de monodromie $\pi_1(\ell - S, \bar{s})$ , quotient de $\mathrm{Gal}(k(\bar{s})/k(s))$ , agit sur $\mathrm{Ev}(V, \mathbb{Q}_\ell(n))$ . Si $k$ n'est pas de caractéristique 2, cette action est irréductible (XVIII 6.7); au n$^{\circ}$ 2 ci-dessous, on montrera que ce résultat reste valable sans hypothèse de caractéristique pour un pinceau générique. Le sous-espace de $\mathrm{Ev}(V, \mathbb{Q}_\ell(n))$ engendré par les classes de cohomologie algébrique est stable sous le groupe de monodromie (i.e. sous Galois); c'est donc tout ou rien, ce qui prouve 1.4.


\SGApage{II}{331}

\underline{Corollaire} 1.5 \quad \underline{Soit} $V \subset \mathbb{P}^{2n+d}$ \underline{une intersection complète de dimension paire} $2n$ \underline{et de multidegré} $\underline{a}$ \underline{générique. On a soit}

(a) \underline{Toute classe de cohomologie algébrique dans} $H^{2n}(V, \mathbb{Q}_\ell(n))$ \underline{est multiple rationnel de} $\eta^n$ .

(b) $H^{2n}(V, \mathbb{Q}_\ell(n))$ \underline{est engendré par des classes de cohomologie algébrique.}

\underline{Si} (b) \underline{est vérifié, pour} toute \underline{intersection complète de dimension} $2n$ \underline{et multidegré} $\underline{a}$ \underline{sur un corps algébriquement clos de caractéristique} $p$ , $H^{2n}(V, \mathbb{Q}_\ell(n))$ \underline{est engendré par les classes algébriques.}

On applique 1.4 pour $X$ une intersection complète de dimension $2n+1$ et de multidegré $(a_1 \ldots a_{d-1})$ générique et pour $\ell$ un pinceau générique de sections de $X$ par des hypersurfaces de degré $a_d$ (cf. XVII 2.4). En fait, on pourrait aussi prendre $X$ de multidegré $(a_1 \ldots a_d)$ et $\ell$ un pinceau générique de sections hyperplanes. La remarque finale résulte de ce qu'une classe de cohomologie algébrique reste algébrique par spécialisation.


\SGAsectiontoc{2. Complément à XVIII}{2}
\underline{2. Complément à XVIII}

\SGAvar{\underline{Proposition} 2.1 \quad \underline{Soient} $X \subset \mathbb{P}$ \underline{une variété projective lisse purement de dimension paire} $2n$ \underline{sur un corps algébriquement clos} $k$ , \underline{et} $\ell \subset \check{\mathbb{P}}$ \underline{un pinceau de Lefschetz générique de sections hyperplanes de} $X$ . \underline{Soient} $s$ \underline{un point géométrique de} $\ell$ , $H_s$ \underline{la section hyperplane correspondante,} $S \subset \ell$ \underline{l'ensemble des valeurs exceptionnelles du pinceau et} $\pi_1(\ell - S, s)$ \underline{le groupe de monodromie.}

}{\underline{Proposition} 2.1 \quad \underline{Soient} $X \subset \mathbb{P}$ \underline{une variété projective lisse purement de dimension impaire} $2n+1$ \underline{sur un corps algébriquement clos} $k$ , \underline{et} $\ell \subset \check{\mathbb{P}}$ \underline{un pinceau de Lefschetz générique de sections hyperplanes de} $X$ . \underline{Soient} $s$ \underline{un point géométrique de} $\ell$ , $H_s$ \underline{la section hyperplane correspondante,} $S \subset \ell$ \underline{l'ensemble des valeurs exceptionnelles du pinceau et} $\pi_1(\ell - S, s)$ \underline{le groupe de monodromie.}

}{F7II-XIX-003}\underline{On suppose que} $X$ \underline{vérifie} (LV) .

\underline{Alors, la représentation de} $\pi_1(\ell - S, s)$ \underline{sur la cohomologie évanescente} $\mathrm{Ev}(H_s, \mathbb{Q}_\ell) \subset H^{2n}(H_s, \mathbb{Q}_\ell)$ \underline{est absolument irréductible.}

\SGAvar{Rappelons que $\mathrm{Ev}(H_s, \mathbb{Q}_\ell)$ est l'orthogonal de l'image de $H^{2n}(X, \mathbb{Q}_\ell)$ dans $H^{2n}(H_{s,} \mathbb{Q}_\ell)$ , image égale au sous-groupe des invariants sous le groupe de monodromie (XVIII 5.6.10).


}{Rappelons que $\mathrm{Ev}(H_s, \mathbb{Q}_\ell)$ est l'orthogonal de l'image de $H^{2n}(X, \mathbb{Q}_\ell)$ dans $H^{2n}(H_s, \mathbb{Q}_\ell)$ , image égale au sous-groupe des invariants sous le groupe de monodromie (XVIII 5.6.10).


}{F7II-XIX-004}\SGApage{II}{332}

\SGAvar{La droite projective $\mathbb{P}^1$ est simplement connexe. Le groupe $\pi_1(\ell - S, s)$ est donc (topologiquement) engendré par les sous-groupes d'inertie en les points de $S$ - et leurs conjugués. Son image dans $\mathrm{GL}(H^n(H_s, \mathbb{Q}_\ell))$ est donc (topologiquement) engendrée par les $x \longrightarrow x - (-1)^n(x\ \delta)\delta$ , pour $\delta$ conjugué à un cycle évanescent. L'image de $H^{2n}(X, \mathbb{Q}_\ell)$ dans $H^{2n}(H_s, \mathbb{Q}_\ell)$ est donc l'orthogonal du sous-espace engendré par les cycles évanescents et leur conjugués: $\mathrm{Ev}(H_s, \mathbb{Q}_\ell)$ est engendré par les cycles évanescents et leur conjugués.

}{La droite projective $\mathbb{P}^1$ est simplement connexe. Le groupe $\pi_1(\ell - S, s)$ est donc (topologiquement) engendré par les sous-groupes d'inertie en les points de $S$ - et leurs conjugués. Son image dans $\mathrm{GL}(H^{2n}(H_s, \mathbb{Q}_\ell))$ est donc (topologiquement) engendrée par les $x \longrightarrow x - (-1)^n(x\ \delta)\delta$ , pour $\delta$ conjugué à un cycle évanescent. L'image de $H^{2n}(X, \mathbb{Q}_\ell)$ dans $H^{2n}(H_s, \mathbb{Q}_\ell)$ est donc l'orthogonal du sous-espace engendré par les cycles évanescents et leurs conjugués: $\mathrm{Ev}(H_s, \mathbb{Q}_\ell)$ est engendré par les cycles évanescents et leurs conjugués.

}{F7II-XIX-005}Pour un pinceau de Lefschetz \underline{générique} , les arguments du n° 6 de XVIII montrent encore que les cycles évanescents $\delta$ relatifs aux différents points de $S$ sont conjugués entre eux.

\SGAvar{Soit $V \subset \mathrm{Ev}(H_s, \mathbb{Q}_\ell)$ un sous-espace invariant $\neq 0$ . Il existe un conjugué $\delta$ d'un cycle évanescent et $v \in V$ tels que $(v, \delta) \neq 0$ (c'est ici qu'on utilise (LV)). Puisque $v - (-1)^m(v\ \delta)\delta \in V$ , on a $\delta \in V$ . Le sous-espace $V$ , contenant un cycle évanescent (et tous ses conjugués) coïncide avec $\mathrm{Ev}(H_s, \mathbb{Q}_\ell)$ .

}{Soit $V \subset \mathrm{Ev}(H_s, \mathbb{Q}_\ell)$ un sous-espace invariant $\neq 0$ . Il existe un conjugué $\delta$ d'un cycle évanescent et $v \in V$ tels que $(v, \delta) \neq 0$ (c'est ici qu'on utilise (LV)). Puisque $v - (-1)^n(v\ \delta)\delta \in V$ , on a $\delta \in V$ . Le sous-espace $V$ , contenant un cycle évanescent (et tous ses conjugués) coïncide avec $\mathrm{Ev}(H_s, \mathbb{Q}_\ell)$ .

}{F7II-XIX-006}Cet argument reste valable après extension du corps de coefficients $\mathbb{Q}_\ell$ , d'où l'absolue irréductibilité.

\underline{Remarque} 2.2 \quad Lorsqu'on omet l'hypothèse (LV) , les arguments qui précèdent prouvent encore que tout sous-espace invariant de $H^{2n}(H_s, \mathbb{Q}_\ell)$ contient $\mathrm{Ev}(H_s, \mathbb{Q}_\ell)$ ou est contenu dans l'image de $H^{2n}(X, \mathbb{Q}_\ell)$ . Le même résultat vaut pour $X$ de dimension paire (et tout pinceau de Lefschetz).


\SGAsectiontoc{3. Preuve de 1.3}{3}
\underline{3. Preuve de 1.3}

\SGAvar{3.1 \quad Il s'agit de prouver que, sous les hypothèses du théorème, 1.5 (b) n'est pas vérifié. Le cas où $h^{0,2n} \neq 0$ sera traité par Katz en XXI 4.1 (la condition $h^{0,2n} \neq 0$ équivaut à $\Sigma\, d_i - (2n + r + 1) \geq 0$ , ce nombre étant l'entier $m$ tel que $\Omega^{2n}_V \simeq \mathcal{O}(m)$ ). Les résultats du présent exposé ne servirons plus dans la suite de séminaire, de sorte qu'il n'y a pas de cercle vicieux.


}{3.1 \quad Il s'agit de prouver que, sous les hypothèses du théorème, 1.5 (b) n'est pas vérifié. Le cas où $h^{0,2n} \neq 0$ sera traité par Katz en XXI 4.1 (la condition $h^{0,2n} \neq 0$ équivaut à $\Sigma\, a_i - (2n + r + 1) \geq 0$ , ce nombre étant l'entier $m$ tel que $\Omega^{2n}_V \simeq \mathcal{O}(m)$ ). Les résultats du présent exposé ne serviront plus dans la suite du séminaire, de sorte qu'il n'y a pas de cercle vicieux.


}{F7II-XIX-007}\SGApage{II}{333}

\SGAvar{3.2 \quad Lorsque $k = \mathbb{C}$ , on démontre 1.3 en notant que si 1.5 (b) est vérifié, la cohomologie de $V$ est purement de type $(n,n)$ , la cohomologie algébrique étant de type $(n,n)$ . Ceci n'est le cas que dans les cas exceptionnels exclus de 1.3 (XI 2.9). Par le principe de Lefschetz, 1.3 est donc vrai en caractérisitque $0$ .

}{3.2 \quad Lorsque $k = \mathbb{C}$ , on démontre 1.3 en notant que si 1.5 (b) est vérifié, la cohomologie de $V$ est purement de type $(n,n)$ , la cohomologie algébrique étant de type $(n,n)$ . Ceci n'est le cas que dans les cas exceptionnels exclus de 1.3 (XI 2.9). Par le principe de Lefschetz, 1.3 est donc vrai en caractéristique $0$ .

}{F7II-XIX-008}\underline{Proposition} 3.3 \quad \underline{Plaçons-nous sous les hypothèses de} 1.4, \underline{et supposons} (b) \underline{vérifié.} \underline{Soit} $\mathrm{Ev}(V, \mathbb{Q})$ \underline{le} $\mathbb{Q}$\underline{-sous-espace vectoriel de} $\mathrm{Ev}(V, \mathbb{Q}_\ell(n))$ \underline{engendré par les classes de cohomologie algébrique.}

(a) \underline{On a} $\mathrm{Ev}(V, \mathbb{Q}) \otimes \mathbb{Q}_\ell \stackrel{\sim}{\longrightarrow} \mathrm{Ev}(V, \mathbb{Q}_\ell(n))$ . \underline{Sur} $\mathrm{Ev}(V, \mathbb{Q})$ , \underline{la forme d'intersection est rationnelle définie, de signe} $(-1)^n$ .

\SGAvar{(b) \underline{Les cycles évanescents et leur conjugués sont dans} $\mathrm{Ev}(V, \mathbb{Q})$ ; \underline{ils y forment un système de racine irréductible de type} A, D \underline{ou} E .

(c) \underline{Le groupe de monodromie (image de} $\pi_1(\ell - S, \bar{s})$ \underline{dans} $\mathrm{GL}(\mathrm{Ev}(V, \mathbb{Q}))$ ) \underline{est fini; c'est le groupe de Weil du système de racines} (b) .

}{(b) \underline{Les cycles évanescents et leurs conjugués sont dans} $\mathrm{Ev}(V, \mathbb{Q})$ ; \underline{ils y forment un système de racines irréductible de type} A, D \underline{ou} E .

(c) \underline{Le groupe de monodromie (image de} $\pi_1(\ell - S, \bar{s})$ \underline{dans} $\mathrm{GL}(\mathrm{Ev}(V, \mathbb{Q}))$ ) \underline{est fini; c'est le groupe de Weyl du système de racines} (b) .

}{F7II-XIX-009}Sauf le caractère défini de la forme d'intersection, (a) résulte de ce que la forme d'intersection est non dégénérée sur $\mathrm{Ev}(V, \mathbb{Q}_\ell(n))$ (par (LV)) et rationnelle sur $\mathrm{Ev}(V, \mathbb{Q})$ .

La condition 1.4 (b) équivaut à dire qu'il existe sur $V$ des cycles algébriques modulo équivalence numérique non invariants par la monodromie. Cette condition est donc indépendante de $\ell$ . Soit $x$ un tel cycle algébrique, et $\delta$ le cycle évanescent relatif à $s_o \in S$ et à un chemin ch de $\bar{s}$ (définissant $V$ ) à $s_o$ . Pour ch convenable, $(x\ \delta) \neq 0$ . Le transformé galoisien $\sigma(x) = x-(-1)^n(x\delta)\delta$ de $x$ est encore algébrique; $(x\delta)\delta$ est donc algébrique, image d'un cycle $y$ indépendant de $\ell$ . Le nombre rationnel $\frac{(-1)^n}{2}(y, y) = (x\delta)^2$ est donc indépendant de $\ell$ ; c'est un carré dans tous les corps $\mathbb{Q}_\ell$ , donc un carré, et $\pm (x, \delta)$ est dans $\mathbb{Q}$ et indépendant de $\ell$ . Ceci montre non seulement que $\delta \in \mathrm{Ev}(V, \mathbb{Q})$ , mais encore que le cycle algébrique modulo équivalence numérique $\pm\, \delta = (\sigma(x) -x)/(x, \delta)$ , est indépendant de $\ell$ .


\SGApage{II}{334}

\SGAvar{Un cycle algébrique sur une variété $V$ est toujours défini sur une extension finie de tout corps de définition de $V$ . Il en résulte que $\pi_1(\ell - S, s)$ agit sur $\mathrm{Ev}(V, \mathbb{Q})$ via un groupe fini. Ce groupe est engendré par les réflexions associées aux cycles évanescents $\delta$ et leur conjugués (XVIII 6.6.1 et n$^\circ$ 2). Les $\delta$ et leur conjugués forment donc un système de racines de groupe de Weil le groupe de monodromie. De plus, toutes les racines de ce système de racines sont conjuguées sous le groupe de Weil: il est de type $A$, $D$ ou $E$ .

La forme d'intersection est, à un facteur près, l'unique forme invariante sous le groupe de Weil. Elle est donc définie. Son signe se détermine en calculant

}{Un cycle algébrique sur une variété $V$ est toujours défini sur une extension finie de tout corps de définition de $V$ . Il en résulte que $\pi_1(\ell - S, s)$ agit sur $\mathrm{Ev}(V, \mathbb{Q})$ via un groupe fini. Ce groupe est engendré par les réflexions associées aux cycles évanescents $\delta$ et leurs conjugués (XVIII 6.6.1 et n$^\circ$ 2). Les $\delta$ et leurs conjugués forment donc un système de racines ayant pour groupe de Weyl le groupe de monodromie. De plus, toutes les racines de ce système de racines sont conjuguées sous le groupe de Weyl : il est de type $A$, $D$ ou $E$ .

La forme d'intersection est, à un facteur près, l'unique forme invariante sous le groupe de Weyl. Elle est donc définie. Son signe se détermine en calculant

}{F7II-XIX-010}\begin{equation*}
(\delta, \delta) = (-1)^n . 2 \quad .
\end{equation*}
\underline{Proposition} 3.4 \quad \underline{Si} $k = \mathbb{C}$ , \underline{et sous les hypothèses de} 1.4, \underline{les conditions suivantes sont équivalentes}

(i) \underline{le groupe de monodromie est fini};

(ii) \underline{la cohomologie évanescente} $\mathrm{Ev}(V)$ \underline{est purement de type} $(n, n)$ .

Soit $\widetilde{X} \longrightarrow \ell$ le fibré de fibres les sections hyperplanes appartenant au pinceau $\ell$ , et $V_s$ la fibre en $s \in \ell(\mathbb{C})$ . Les $\mathrm{Ev}(V_s)$ forment une variation algébrique de structures de Hodge sur $\ell$ . D'après le théorème de la partie fixe [2] 4.1.2 appliqué à un revêtement fini de $\ell$ (cf. [2] 4.1.3.3), (i) implique que la structure de Hodge des $\mathrm{Ev}(V_s)$ est localement constante. Le groupe de monodromie agit donc par automorphismes de la structure de Hodge. Pour $T$ une transformation de Picard-Lefschetz, $\mathrm{Im}(T-1)$ est une sous-structure de Hodge de rang un, donc de type $(n, n)$ , engendrée par un cycle évanescent $\delta$ . Les cycles évanescents étant de type $(n, n)$ , $\mathrm{Ev}(V_s)$ tout entier l'est.

Réciproquement, si $\mathrm{Ev}(V)$ est purement de type $(n, n)$ , la forme d'intersection sur $\mathrm{Ev}(V)$ est définie et le groupe de monodromie fini (cf. [2] 4.1.3.3).


\SGApage{II}{335}

\SGAvar{\underline{Remarque} 3.5 \quad Si les conditions équivalentes de 3.3 sont vérifiées, on voit comme en 3.3 que les cycles évanescents (et leurs conjugués) forment un système de racines de type $A$, $D$ ou $E$ dont le groupe de monodromie est le groupe de Weil.

}{\underline{Remarque} 3.5 \quad Si les conditions équivalentes de 3.4 sont vérifiées, on voit comme en 3.3 que les cycles évanescents (et leurs conjugués) forment un système de racines de type $A$, $D$ ou $E$ dont le groupe de monodromie est le groupe de Weyl.

}{F7II-XIX-011}3.6. \quad Prouvons 1.3 pour $k$ de caractéristique finie $\neq 2$ . Soient $S = \mathrm{Spec}(V)$ le spectre d'un anneau de valuation discrète d'inégale caractéristique de corps résiduel $k$ , et $X \subset \mathbb{P}^{2n+d}_S$ une intersection complète de dimension relative $2n+1$ et de multidegré $(a_1 \ldots a_{d-1})$ sur $S$ . Soit $\ell$ un pinceau de Lefschetz de sections hypersurfaces de degré $a_d$ de $X$ (en fibre spéciale, $\ell_s$ est un pinceau de Lefschetz pour $X_s$ ). Soit $f : \widetilde{X} \longrightarrow \ell \simeq \mathbb{P}^1_S$ le fibré de fibres les sections du pinceau. Le lieu exceptionnel $T \subset \ell$ est étale sur $S$ (XVII 6; on utilise ici que $p \neq 2$ ). Le faisceau $R^{2n} f_* \mathbb{Q}_\ell$ est localement constant sur $\ell - T$ ; il est donc modérément ramifié et sa monodromie est la même en caractéristique $p$ et en caractéristique $0$. Si 1.5 (b) est vérifié, on déduit donc de 3.3 (c) et 3.4 que le groupe de cohomologie $H^{2n}(V)$ des intersections complètes \underline{complexes} de dimension $2n$ et multidegré \underline{$a$} est purement de type $(n, n)$ , ce qui n'arrive que dans les cas exclus par 1.3 (XI 2.9).


\SGAsectiontoc{4. Deuxième complément à XVIII}{4}
\underline{4. Deuxième complément à XVIII}

\SGAvar{Dans ce numéro, nous travaillons sur $\mathbb{C}$ et en cohomologie entière (transcendante). Nous reprenons, avec des méthodes plus proches de celles de Lefschetz, certains des résultats de XVIII. Dans XVIII, ils n'étaient prouvés qu'après tensorisation avec $\mathbb{Q}$ (on plutôt $\mathbb{Q}_\ell$ , ce qui revient au même).

}{Dans ce numéro, nous travaillons sur $\mathbb{C}$ et en cohomologie entière (transcendante). Nous reprenons, avec des méthodes plus proches de celles de Lefschetz, certains des résultats de XVIII. Dans XVIII, ils n'étaient prouvés qu'après tensorisation avec $\mathbb{Q}$ (ou plutôt $\mathbb{Q}_\ell$ , ce qui revient au même).

}{F7II-XIX-012}4.1 \quad Soient $X \subset \mathbb{P}$ une variété projective complexe lisse purement de dimension $n + 1$ et $\ell$ un pinceau de Lefschetz de sections hyperplanes de $X$ , d'axe $A$ . On désigne aussi par $\ell$ l'espace des paramètres du pinceau, isomorphe à $\mathbb{P}^1$ . Soit $S \subset \ell$ l'ensemble des valeurs exceptionnelles. Pour $s \in \ell$ , on notera $H_s$ la section hyperplane correspondante de $X$ (de dimension $n$).


\SGApage{II}{336}

Soit $T(A)$ un voisinage tubulaire fermé de $A \cap X$ dans $X$ , de bord transverse à tous les $H_s$ .

\SGAvar{Soient $0$, $\infty$ deux points dans $\ell - S$ , et soit $\alpha$ un difféomorphisme orienté de $\ell - S$ vers $\mathbb{P}^1(\mathbb{C})$ envoyant $0$ sur $0$ , $\infty$ sur $\infty$ et $S$ dans la circonférence de rayon un centrée en $0$ . Le composé $\alpha f$ de $\alpha$ et de la projection sur $\ell$ fait de $X - T(A)^o$ un fibré singulier en variétés à bord sur $\mathbb{P}^1(\mathbb{C})$ (une singularité au-dessus de chaque point de $\alpha(S)$ ). Nous noterons $c_s$ le segment de droite joignant $0$ à $\alpha(s)$ .

}{Soient $0$, $\infty$ deux points dans $\ell - S$ , et soit $\alpha$ un difféomorphisme orienté de $\ell$ vers $\mathbb{P}^1(\mathbb{C})$ envoyant $0$ sur $0$ , $\infty$ sur $\infty$ et $S$ dans la circonférence de rayon un centrée en $0$ . Le composé $\alpha f$ de $\alpha$ et de la projection sur $\ell$ fait de $X - T(A)^o$ un fibré singulier en variétés à bord sur $\mathbb{P}^1(\mathbb{C})$ (une singularité au-dessus de chaque point de $\alpha(S)$ ). Nous noterons $c_s$ le segment de droite joignant $0$ à $\alpha(s)$ .

}{F7II-XIX-013}Omettant la mention de $\alpha$ , on a la figure

\SGAvar{\begin{center}
\begin{tikzcd}[row sep=huge, column sep=huge]
{} \& s_2 \& s_1 \\
\cdots \& 0 \arrow[u, dash, "{c_2}"'] \arrow[ur, dash, "{c_1}"'] \arrow[ul, dash]
\end{tikzcd}\ .
\end{center}












Au voisinage de $s \in S$ et du point singulier $x_s$ de $H_s \subset X$ , on se trouve (dans des coordonnées locales convenables) dans la situation étudié en XIV 3.2. Si $\alpha$ est choisi holomorphe près de $s$ , on peut même trouver des coordonnées locales sur $X$ et $\ell$ telles que $f$ s'écrive $\Sigma z_i^2$ et que $\alpha^{-1}(c_s)$ soit " $z$ réel positif" . Au-dessus de l'intersection d'un voisinage de $s$ et de $\alpha^{-1}(c_s) - s$ , on peut alors considérer le fibré en sphères réelles $\Delta'_s \subset X$ , défini par " $z_i$ réel" . Pour $t \to s$ , la fibre $\Delta'_{st}$ de $\Delta'_s$ en $t$ tend vers $x_s$ . On notera $\Delta_s$ la réunion de $\{x_s\}$ et d'un fibré en sphères au-dessus de $\alpha^{-1}(c_s) - \{s\}$ , contenu dans $X - T(A)$ et qui prolonge $\Delta'_s$ . $\Delta_s$ est un disque plongé dans $X$ , de bord $H_o \cap \Delta_s$ une sphère qui représente le cycle évanescent $\delta_s$ en $s$ , transporté dans $H_o$ le long du chemin $\alpha^{-1}(c_s)$ .

}{\begin{center}
\begin{tikzpicture}[x=1.2cm,y=1.2cm]
  \coordinate (O) at (0,0);
  \coordinate (L) at (-1.8,1.8);
  \coordinate (S2) at (0,2.8);
  \coordinate (S1) at (2.2,2.0);
  \draw (O)--(L);
  \draw (O)--(S2);
  \draw (O)--(S1);
  \node[right] at (O) {$0$};
  \node[above right] at (S2) {$s_2$};
  \node[above right] at (S1) {$s_1$};
  \node[right] at (0,1.65) {$c_2$};
  \node[below right] at (1.25,1.14) {$c_1$};
  \fill (-0.73,0.76) circle (0.035);
  \fill (-0.55,0.91) circle (0.035);
  \fill (-0.37,1.03) circle (0.035);
\end{tikzpicture}\ .
\end{center}












Au voisinage de $s \in S$ et du point singulier $x_s$ de $H_s \subset X$ , on se trouve (dans des coordonnées locales convenables) dans la situation étudiée en XIV 3.2. Si $\alpha$ est choisi holomorphe près de $s$ , on peut même trouver des coordonnées locales sur $X$ et $\ell$ telles que $f$ s'écrive $\Sigma z_i^2$ et que $\alpha^{-1}(c_s)$ soit " $z$ réel positif" . Au-dessus de l'intersection d'un voisinage de $s$ et de $\alpha^{-1}(c_s) - s$ , on peut alors considérer le fibré en sphères réelles $\Delta'_s \subset X$ , défini par " $z_i$ réel" . Pour $t \to s$ , la fibre $\Delta'_{st}$ de $\Delta'_s$ en $t$ tend vers $x_s$ . On notera $\Delta_s$ la réunion de $\{x_s\}$ et d'un fibré en sphères au-dessus de $\alpha^{-1}(c_s) - \{s\}$ , contenu dans $X - T(A)$ et qui prolonge $\Delta'_s$ . $\Delta_s$ est un disque plongé dans $X$ , de bord $H_o \cap \Delta_s$ une sphère qui représente le cycle évanescent $\delta_s$ en $s$ , transporté dans $H_o$ le long du chemin $\alpha^{-1}(c_s)$ .

}{F7II-XIX-014}Appliquons la théorie de Morse à la fonction $|\alpha f|^2$ , pour étudier comment change l'homotopie de

\begin{equation*}
\{x \in X| \ x \in T(A) \quad \mathrm{ou} \quad |\alpha f(x)|^2 \leq r \} \quad 
\end{equation*}

\SGApage{II}{337}

 lorsque $r$ croît. Soit

\begin{equation*}
X^- = \{x \in X \mid x \in T(A) \quad \mathrm{ou} \quad \alpha f(x) \neq \infty \} \ .
\end{equation*}
On trouve que l'inclusion

\begin{equation*}
H_o \cup \bigcup_{s \in S} \Delta_s \hookrightarrow T(A) \cup H_o \cup \bigcup_{s \in S} \Delta_s \hookrightarrow X^-
\end{equation*}
est une équivalence d'homotopie. On sait, encore par la théorie de Morse, que $X - X^-$ peut se rétracter par déformation sur un sous-complexe de dimension réelle $\leq n$ . Au total, on a donc le résultat suivant

\underline{Théorème} 4.2. \quad \underline{Il existe un sous-complexe} $\Sigma$ \underline{compact et de dimension (réelle)} $\leq n$ \underline{de} $H_\infty^{\mathrm{aff}} = H_\infty - A$ \underline{tel que l'inclusion}

\begin{equation*}
H_o \cup \bigcup_{s \in S} \Delta_s \longrightarrow X - \Sigma
\end{equation*}
\underline{soit une équivalence d'homotopie.}

\SGAvar{4.3 \quad La soustraction de $\Sigma$ est sans effet sur les groupes de cohomologie, d'homologie ou d'homotopie d'indice $i \leq n$ de $X$ . Puisque $H_o \cup \bigcup_{s \in S} \Delta_s$ se déduit de $H_o$ en attachant des $n+1$-disque le long de sphères représentant les cycles évanescents, on trouve en homologie

\begin{equation*}
\mathbb{Z}^S \longrightarrow H_n(V_o) \longrightarrow H_n(X) \longrightarrow 0 \ .
\end{equation*}
La suite exacte analogue en cohomologie, et celle qui s'en déduit par dualité de Poincaré, forment un diagramme commutatif

\begin{equation}\tag{4.3.1}
\begin{tikzcd}[row sep=large, column sep=large]
{} \& \mathbb{Z}^S \& {} \& {} \\
\mathbb{Z}^S \arrow[r, "{\Sigma n_s \delta_s}"] \& H^n(V_o) \arrow[u, "{( \ , \delta_s)}"'] \arrow[r] \& H^{n+2}(X) \arrow[r] \& 0 \\
{} \& H^n(X) \arrow[u] \arrow[ur, "{\eta\Lambda}"'] \& {} \& {} \\
{} \& 0 \arrow[u] \& {} \&
{} \end{tikzcd}
\end{equation}
















}{4.3 \quad La soustraction de $\Sigma$ est sans effet sur les groupes de cohomologie, d'homologie ou d'homotopie d'indice $i \leq n$ de $X$ . Puisque $H_o \cup \bigcup_{s \in S} \Delta_s$ se déduit de $H_o$ en attachant des $n+1$-disques le long de sphères représentant les cycles évanescents, on trouve en homologie

\begin{equation*}
\mathbb{Z}^S \longrightarrow H_n(H_o) \longrightarrow H_n(X) \longrightarrow 0 \ .
\end{equation*}
La suite exacte analogue en cohomologie, et celle qui s'en déduit par dualité de Poincaré, forment un diagramme commutatif

\begin{equation}\tag{4.3.1}
\begin{tikzcd}[row sep=large, column sep=large]
{} \& \mathbb{Z}^S \& {} \& {} \\
\mathbb{Z}^S \arrow[r, "{\Sigma n_s \delta_s}"] \& H^n(H_o) \arrow[u, "{( \ , \delta_s)}"'] \arrow[r] \& H^{n+2}(X) \arrow[r] \& 0 \\
{} \& H^n(X) \arrow[u] \arrow[ur, "{\eta\wedge}"'] \& {} \& {} \\
{} \& 0 \arrow[u] \& {} \&
{} \end{tikzcd}
\end{equation}
















}{F7II-XIX-015}\SGApage{II}{338}


\SGAsectiontoc{5. Les cas d'exception au théorème de Noether}{5}
5. \underline{Les cas d'exception au théorème de Noether}

Dans ce numéro, nous persistons à travailler sur $\mathbb{C}$ , en cohomologie entière.

5.1 \quad Soient $2n$, $d$ et $\underline{a} = (a_1,\ldots,a_d)$ $(d \geq 1, a_i \geq 2)$ tels que la cohomologie de dimension $2n$ des intersections complètes lisses $V$ de dimension $2n$ et de multidegré $\underline{a}$ soit purement de type $(n, n)$ . En d'autres termes, $n = 0$ ou $(2n; \underline{a})$ est égal à $(2n; 2)$, $(2n; 2, 2)$ ou $(2; 3)$ (XI 2.9 ).

Supposons que $V$ se déduise d'une intersection complète $X \subset \mathbb{P}^{2n+d}$ (resp. $X \subset \mathbb{P}^{2n+d+1}$) de dimension $2n + 1$ et de multidegré $(a_1,\ldots,a_{d-1})$ (resp. $(a_1,\ldots,a_d)$ ) en prenant une section par une hypersurface de degré $a_d$ (resp. une section hyperplane) appartenant à un pinceau de Lefschetz $\ell$ .

La \underline{partie évanescente} $\mathrm{Ev}(V)$ de $H^{2n}(V)$ , égale à la \underline{partie primitive} de cette cohomologie est, d'après 4.3.1, engendrée \underline{sur} $\mathbb{Z}$ par les cycles évanescents. C'est l'orthogonal de la puissance $\eta^n$ de la classe de cohomologie $\eta$ d'une section hyperplane.

\SGAvar{Puisque sur $H^{2n}(V)$ la forme d'intersection est unimodulaire, et que $\eta^{2n}$ n'est pas le multiple $by$ $(b > 1)$ d'une autre classe de cohomologie (XI 1.6 (iv)), le discriminant de la forme d'intersection sur $\mathrm{Ev}(V)$ est égal à

}{Puisque sur $H^{2n}(V)$ la forme d'intersection est unimodulaire, et que $\eta^n$ n'est pas le multiple $b y$ $(b > 1)$ d'une autre classe de cohomologie (XI 1.6 (iv)), le discriminant de la forme d'intersection sur $\mathrm{Ev}(V)$ est égal à

}{F7II-XIX-016}\begin{equation*}
(\eta^n, \eta^n) = \Pi\, a_i \quad .
\end{equation*}
Soit $R$ le système de racines dans $H^{2n}(V)$ formé des cycles évanescents et de leurs conjugués. On a sur $R$ les informations suivantes

(a) \quad $R$ est vide ou de type $A$, $D$ ou $E$ ;

(b) \quad le rang de $R$ est donné en XI 1.10 ;

(c) \quad le discriminant est $\Pi\, a_i$ . Dans les tables, ce discriminant est l'ordre du centre du groupe simple simplement connexe complexe correspondant. L'argument par lequel nous avons calculé le discriminant montre que ce centre, dual du quotient du réseau des poids par celui des racines, est cyclique.


\SGApage{II}{339}

Dans tous les cas de figure, ces informations déterminent $R$ complètement si $n \neq 0$ . On trouve:

\SGAvar{\underline{Proposition} 5.2 \quad (i) \underline{Pour} $V$ \underline{uhe quadrique de dimension paire}, $R$ \underline{est de type} $A_1$

}{\underline{Proposition} 5.2 \quad (i) \underline{Pour} $V$ \underline{une quadrique de dimension paire}, $R$ \underline{est de type} $A_1$

}{F7II-XIX-017}(ii) \underline{Pour} $V$ \underline{une surface cubique}, $R$ \underline{est de type} $E_6$ .

(iii) \underline{Pour} $V$ \underline{de dimension} $2n$ , \underline{et l'intersection de deux quadriques}, $R$ \underline{est de type} $D_{2n+3}$ .

\SGAvar{5.3. \quad Supposons que $n \neq 0$ . L'ensemble $R^1$ des vecteurs de longueur $\pm 2$ dans $\mathrm{Ev}(V)$ est aussi un système de racine, ayant les mêmes caractéristiques (a) (b) (c) que $R$ . On a donc $R = R^1$ . Soit $\pi_1$ le groupe fondamental de l'espace des paramètres de toutes les intersections complètes $V \subset \mathbb{P}^{2n+d}$ du type considéré. L'image de $\pi_1$ dans $\mathrm{GL}(\mathrm{Ev}(V))$ est le groupe de Weil de $R$ . Dans tous les cas de figure, c'est le groupe d'automorphismes du triple

\begin{equation*}
(H^{2n}(V),\ \eta^n\ ,\ \mbox{forme d'intersection}).
\end{equation*}
5.4. \quad Avec $R$ remplacé par $R^1$ de 5.3, le résultat 4.2 (ii) est bien connu [1] et pour les petites valeurs de $n$ , 4.2 (iii) était connu, au moins par Manin ($n = 0$ est trivial; une surface intersection de deux quadriques est une surface de del Pezzo). Une étude détaillée de l'intersection de deux quadriques a été faite par M. Reid [4]. Il déduit 4.2 (iii) d'une étude de la géométrie des sous-espaces linéaires de dimension $n$ sur une intersection de deux quadriques $V \subset \mathbb{P}^{2n+2}$ .


}{5.3. \quad Supposons que $n \neq 0$ . L'ensemble $R^1$ des vecteurs de longueur $\pm 2$ dans $\mathrm{Ev}(V)$ est aussi un système de racines, ayant les mêmes caractéristiques (a) (b) (c) que $R$ . On a donc $R = R^1$ . Soit $\pi_1$ le groupe fondamental de l'espace des paramètres de toutes les intersections complètes $V \subset \mathbb{P}^{2n+d}$ du type considéré. L'image de $\pi_1$ dans $\mathrm{GL}(\mathrm{Ev}(V))$ est le groupe de Weyl de $R$ . Dans tous les cas de figure, c'est le groupe d'automorphismes du triple

\begin{equation*}
(H^{2n}(V),\ \eta^n\ ,\ \mbox{forme d'intersection}).
\end{equation*}
5.4. \quad Avec $R$ remplacé par $R^1$ de 5.3, le résultat 5.2 (ii) est bien connu [1] et pour les petites valeurs de $n$ , 5.2 (iii) était connu, au moins par Manin ($n = 0$ est trivial; une surface intersection de deux quadriques est une surface de del Pezzo). Une étude détaillée de l'intersection de deux quadriques a été faite par M. Reid [4]. Il déduit 5.2 (iii) d'une étude de la géométrie des sous-espaces linéaires de dimension $n$ sur une intersection de deux quadriques $V \subset \mathbb{P}^{2n+2}$ .


}{F7II-XIX-018}\underline{Bibliographie}

[1] \quad E. CARTAN. Quelques remarques sur les 28 bitangentes d'une quartique plane et les 27 droites d'une surface cubique.\\ Bull. Sc. Math. \underline{70} 1946 p. 42-45, Oeuvres, partie I, tome 2 p. 1353.

[2] \quad P. DELIGNE. Théorie de Hodge II. Publ. Math. IHES \underline{40}, 1971 p. 5-58.


\SGApage{II}{340}

[3]\quad S. LEFSCHETZ. L'analysis situs et la géométrie algébrique. Gauthier-Villars 1924 (réédité dans Selected papers - Chelsea).

[4]\quad M. REID. Thèse, Cambridge 1972.

[5]\quad J.P. SERRE. Faisceaux algébriques cohérents. Ann. of Math. LXI (1955) p. 197-278 (cité FAC)
\selectlanguage{french}

\ifnum\SGAform=1
\clearpage\section*{Apparat critique}\small
Les lectures ci-dessous sont données dans le TeX exact des témoins, afin de préserver les distinctions de notation. Les indices décimaux renvoient aux différences individuelles du passage marqué.

\par\medskip\noindent\hypertarget{app-F7II-XIX-001.1}{}\textbf{F7II-XIX-001.1}\quad p. 329; printed misprint.\par
\noindent Imprimé : {\ttfamily\footnotesize d\allowbreak{}}\par
\noindent Corrigé : {\ttfamily\footnotesize a\allowbreak{}}\par
\noindent Source : 7II:575; 7II:978. \nolinkurl{fixes/MISPRINTS_7II.json#575}; \nolinkurl{fixes/MISPRINTS_7II.json#978}\par

\par\medskip\noindent\hypertarget{app-F7II-XIX-002.1}{}\textbf{F7II-XIX-002.1}\quad p. 330; misprint\_pass.\par
\noindent Imprimé : {\ttfamily\footnotesize grassmanienne\allowbreak{}}\par
\noindent Corrigé : {\ttfamily\footnotesize grassmannienne\allowbreak{}}\par
\noindent Source : 7II:979. \nolinkurl{fixes/FIXLOG_7II_CURRENT.jsonl#906}\par

\par\medskip\noindent\hypertarget{app-F7II-XIX-003.1}{}\textbf{F7II-XIX-003.1}\quad p. 331; printed misprint.\par
\noindent Imprimé : {\ttfamily\footnotesize paire\allowbreak{}}\par
\noindent Corrigé : {\ttfamily\footnotesize impaire\allowbreak{}}\par
\noindent Source : 7II:576; 7II:980. \nolinkurl{fixes/MISPRINTS_7II.json#576}; \nolinkurl{fixes/MISPRINTS_7II.json#980}\par

\par\medskip\noindent\hypertarget{app-F7II-XIX-003.2}{}\textbf{F7II-XIX-003.2}\quad p. 331; printed misprint.\par
\noindent Imprimé : \emph{(vide)}\par
\noindent Corrigé : {\ttfamily\footnotesize +1\allowbreak{}}\par
\noindent Source : 7II:576; 7II:980. \nolinkurl{fixes/MISPRINTS_7II.json#576}; \nolinkurl{fixes/MISPRINTS_7II.json#980}\par

\par\medskip\noindent\hypertarget{app-F7II-XIX-004.1}{}\textbf{F7II-XIX-004.1}\quad p. 331; inherited corrected-layer difference.\par
\noindent Imprimé : {\ttfamily\footnotesize \{\allowbreak{}}\par
\noindent Corrigé : \emph{(vide)}\par
\noindent Source : ADOPT-F7II-XIX-004.1. \nolinkurl{fr/sga7ii/diplomatic/expose_XIX_body.tex}; \nolinkurl{fr/sga7ii/current/expose_XIX_body.tex}\par

\par\medskip\noindent\hypertarget{app-F7II-XIX-004.2}{}\textbf{F7II-XIX-004.2}\quad p. 331; inherited corrected-layer difference.\par
\noindent Imprimé : {\ttfamily\footnotesize \}\allowbreak{}}\par
\noindent Corrigé : \emph{(vide)}\par
\noindent Source : ADOPT-F7II-XIX-004.2. \nolinkurl{fr/sga7ii/diplomatic/expose_XIX_body.tex}; \nolinkurl{fr/sga7ii/current/expose_XIX_body.tex}\par

\par\medskip\noindent\hypertarget{app-F7II-XIX-005.1}{}\textbf{F7II-XIX-005.1}\quad p. 332; printed misprint.\par
\noindent Imprimé : {\ttfamily\footnotesize n\allowbreak{}}\par
\noindent Corrigé : {\ttfamily\footnotesize \{2n\}\allowbreak{}}\par
\noindent Source : 7II:577; 7II:981. \nolinkurl{fixes/MISPRINTS_7II.json#577}; \nolinkurl{fixes/MISPRINTS_7II.json#981}\par

\par\medskip\noindent\hypertarget{app-F7II-XIX-005.2}{}\textbf{F7II-XIX-005.2}\quad p. 332; misprint\_pass.\par
\noindent Imprimé : {\ttfamily\footnotesize leur\allowbreak{}}\par
\noindent Corrigé : {\ttfamily\footnotesize leurs\allowbreak{}}\par
\noindent Source : 7II:579; 7II:982. \nolinkurl{fixes/FIXLOG_7II_CURRENT.jsonl#762}; \nolinkurl{fixes/MISPRINTS_7II.json#579}; \nolinkurl{fixes/MISPRINTS_7II.json#982}\par

\par\medskip\noindent\hypertarget{app-F7II-XIX-005.3}{}\textbf{F7II-XIX-005.3}\quad p. 332; misprint\_pass.\par
\noindent Imprimé : {\ttfamily\footnotesize leur\allowbreak{}}\par
\noindent Corrigé : {\ttfamily\footnotesize leurs\allowbreak{}}\par
\noindent Source : 7II:580; 7II:983. \nolinkurl{fixes/FIXLOG_7II_CURRENT.jsonl#763}; \nolinkurl{fixes/MISPRINTS_7II.json#580}; \nolinkurl{fixes/MISPRINTS_7II.json#983}\par

\par\medskip\noindent\hypertarget{app-F7II-XIX-006.1}{}\textbf{F7II-XIX-006.1}\quad p. 332; printed misprint.\par
\noindent Imprimé : {\ttfamily\footnotesize m\allowbreak{}}\par
\noindent Corrigé : {\ttfamily\footnotesize n\allowbreak{}}\par
\noindent Source : 7II:578. \nolinkurl{fixes/MISPRINTS_7II.json#578}\par

\par\medskip\noindent\hypertarget{app-F7II-XIX-007.1}{}\textbf{F7II-XIX-007.1}\quad p. 332; designated-witness attribution.\par
\noindent Imprimé : {\ttfamily\footnotesize d\allowbreak{}}\par
\noindent Corrigé : {\ttfamily\footnotesize a\allowbreak{}}\par
\noindent Source : WITNESS-F7II-XIX-007.1. \nolinkurl{fr/sga7ii/diplomatic/expose_XIX_body.tex}; \nolinkurl{fr/sga7ii/current/expose_XIX_body.tex}\par

\par\medskip\noindent\hypertarget{app-F7II-XIX-007.2}{}\textbf{F7II-XIX-007.2}\quad p. 332; printed misprint.\par
\noindent Imprimé : {\ttfamily\footnotesize servirons\allowbreak{}}\par
\noindent Corrigé : {\ttfamily\footnotesize serviront\allowbreak{}}\par
\noindent Source : 7II:581; 7II:985. \nolinkurl{fixes/MISPRINTS_7II.json#581}; \nolinkurl{fixes/MISPRINTS_7II.json#985}\par

\par\medskip\noindent\hypertarget{app-F7II-XIX-007.3}{}\textbf{F7II-XIX-007.3}\quad p. 332; printed misprint.\par
\noindent Imprimé : {\ttfamily\footnotesize de\allowbreak{}}\par
\noindent Corrigé : {\ttfamily\footnotesize du\allowbreak{}}\par
\noindent Source : 7II:582; 7II:986. \nolinkurl{fixes/MISPRINTS_7II.json#582}; \nolinkurl{fixes/MISPRINTS_7II.json#986}\par

\par\medskip\noindent\hypertarget{app-F7II-XIX-008.1}{}\textbf{F7II-XIX-008.1}\quad p. 333; printed misprint.\par
\noindent Imprimé : {\ttfamily\footnotesize caractérisitque\allowbreak{}}\par
\noindent Corrigé : {\ttfamily\footnotesize caractéristique\allowbreak{}}\par
\noindent Source : 7II:583; 7II:987. \nolinkurl{fixes/MISPRINTS_7II.json#583}; \nolinkurl{fixes/MISPRINTS_7II.json#987}\par

\par\medskip\noindent\hypertarget{app-F7II-XIX-009.1}{}\textbf{F7II-XIX-009.1}\quad p. 333; printed misprint.\par
\noindent Imprimé : {\ttfamily\footnotesize leur\allowbreak{}}\par
\noindent Corrigé : {\ttfamily\footnotesize leurs\allowbreak{}}\par
\noindent Source : 7II:988. \nolinkurl{fixes/MISPRINTS_7II.json#988}\par

\par\medskip\noindent\hypertarget{app-F7II-XIX-009.2}{}\textbf{F7II-XIX-009.2}\quad p. 333; printed misprint.\par
\noindent Imprimé : {\ttfamily\footnotesize racine\allowbreak{}}\par
\noindent Corrigé : {\ttfamily\footnotesize racines\allowbreak{}}\par
\noindent Source : 7II:585; 7II:989. \nolinkurl{fixes/MISPRINTS_7II.json#585}; \nolinkurl{fixes/MISPRINTS_7II.json#989}\par

\par\medskip\noindent\hypertarget{app-F7II-XIX-009.3}{}\textbf{F7II-XIX-009.3}\quad p. 333; inherited corrected-layer difference.\par
\noindent Imprimé : {\ttfamily\footnotesize Weil\allowbreak{}}\par
\noindent Corrigé : {\ttfamily\footnotesize Weyl\allowbreak{}}\par
\noindent Source : ADOPT-F7II-XIX-009.3. \nolinkurl{fr/sga7ii/diplomatic/expose_XIX_body.tex}; \nolinkurl{fr/sga7ii/current/expose_XIX_body.tex}\par

\par\medskip\noindent\hypertarget{app-F7II-XIX-010.1}{}\textbf{F7II-XIX-010.1}\quad p. 334; printed misprint.\par
\noindent Imprimé : {\ttfamily\footnotesize leur\allowbreak{}}\par
\noindent Corrigé : {\ttfamily\footnotesize leurs\allowbreak{}}\par
\noindent Source : 7II:990. \nolinkurl{fixes/MISPRINTS_7II.json#990}\par

\par\medskip\noindent\hypertarget{app-F7II-XIX-010.2}{}\textbf{F7II-XIX-010.2}\quad p. 334; printed misprint.\par
\noindent Imprimé : {\ttfamily\footnotesize leur\allowbreak{}}\par
\noindent Corrigé : {\ttfamily\footnotesize leurs\allowbreak{}}\par
\noindent Source : 7II:991. \nolinkurl{fixes/MISPRINTS_7II.json#991}\par

\par\medskip\noindent\hypertarget{app-F7II-XIX-010.3}{}\textbf{F7II-XIX-010.3}\quad p. 334; designated-witness attribution.\par
\noindent Imprimé : {\ttfamily\footnotesize de\allowbreak{}}\par
\noindent Corrigé : {\ttfamily\footnotesize ayant\allowbreak{} pour\allowbreak{}}\par
\noindent Source : WITNESS-F7II-XIX-010.3. \nolinkurl{fr/sga7ii/diplomatic/expose_XIX_body.tex}; \nolinkurl{fr/sga7ii/current/expose_XIX_body.tex}\par

\par\medskip\noindent\hypertarget{app-F7II-XIX-010.4}{}\textbf{F7II-XIX-010.4}\quad p. 334; designated-witness attribution.\par
\noindent Imprimé : {\ttfamily\footnotesize Weil\allowbreak{}}\par
\noindent Corrigé : {\ttfamily\footnotesize Weyl\allowbreak{}}\par
\noindent Source : WITNESS-F7II-XIX-010.4. \nolinkurl{fr/sga7ii/diplomatic/expose_XIX_body.tex}; \nolinkurl{fr/sga7ii/current/expose_XIX_body.tex}\par

\par\medskip\noindent\hypertarget{app-F7II-XIX-010.5}{}\textbf{F7II-XIX-010.5}\quad p. 334; inherited corrected-layer difference.\par
\noindent Imprimé : {\ttfamily\footnotesize Weil\allowbreak{}}\par
\noindent Corrigé : {\ttfamily\footnotesize Weyl\allowbreak{} }\par
\noindent Source : ADOPT-F7II-XIX-010.5. \nolinkurl{fr/sga7ii/diplomatic/expose_XIX_body.tex}; \nolinkurl{fr/sga7ii/current/expose_XIX_body.tex}\par

\par\medskip\noindent\hypertarget{app-F7II-XIX-010.6}{}\textbf{F7II-XIX-010.6}\quad p. 334; inherited corrected-layer difference.\par
\noindent Imprimé : {\ttfamily\footnotesize Weil\allowbreak{}}\par
\noindent Corrigé : {\ttfamily\footnotesize Weyl\allowbreak{}}\par
\noindent Source : ADOPT-F7II-XIX-010.6. \nolinkurl{fr/sga7ii/diplomatic/expose_XIX_body.tex}; \nolinkurl{fr/sga7ii/current/expose_XIX_body.tex}\par

\par\medskip\noindent\hypertarget{app-F7II-XIX-011.1}{}\textbf{F7II-XIX-011.1}\quad p. 335; misprint\_pass.\par
\noindent Imprimé : {\ttfamily\footnotesize 3\allowbreak{}}\par
\noindent Corrigé : {\ttfamily\footnotesize 4\allowbreak{}}\par
\noindent Source : 7II:588. \nolinkurl{fixes/FIXLOG_7II_CURRENT.jsonl#764}\par

\par\medskip\noindent\hypertarget{app-F7II-XIX-011.2}{}\textbf{F7II-XIX-011.2}\quad p. 335; inherited corrected-layer difference.\par
\noindent Imprimé : {\ttfamily\footnotesize Weil\allowbreak{}}\par
\noindent Corrigé : {\ttfamily\footnotesize Weyl\allowbreak{}}\par
\noindent Source : ADOPT-F7II-XIX-011.2. \nolinkurl{fr/sga7ii/diplomatic/expose_XIX_body.tex}; \nolinkurl{fr/sga7ii/current/expose_XIX_body.tex}\par

\par\medskip\noindent\hypertarget{app-F7II-XIX-012.1}{}\textbf{F7II-XIX-012.1}\quad p. 335; printed misprint.\par
\noindent Imprimé : {\ttfamily\footnotesize on\allowbreak{}}\par
\noindent Corrigé : {\ttfamily\footnotesize ou\allowbreak{}}\par
\noindent Source : 7II:589; 7II:993. \nolinkurl{fixes/MISPRINTS_7II.json#589}; \nolinkurl{fixes/MISPRINTS_7II.json#993}\par

\par\medskip\noindent\hypertarget{app-F7II-XIX-013.1}{}\textbf{F7II-XIX-013.1}\quad p. 336; printed misprint.\par
\noindent Imprimé : {\ttfamily\footnotesize  -\allowbreak{} S\allowbreak{}}\par
\noindent Corrigé : \emph{(vide)}\par
\noindent Source : 7II:591; 7II:994. \nolinkurl{fixes/MISPRINTS_7II.json#591}; \nolinkurl{fixes/MISPRINTS_7II.json#994}\par

\par\medskip\noindent\hypertarget{app-F7II-XIX-014.1}{}\textbf{F7II-XIX-014.1}\quad p. 336; inherited corrected-layer difference.\par
\noindent Imprimé : {\ttfamily\footnotesize tikzcd\allowbreak{}}\par
\noindent Corrigé : {\ttfamily\footnotesize tikzpicture\allowbreak{}}\par
\noindent Source : ADOPT-F7II-XIX-014.1. \nolinkurl{fr/sga7ii/diplomatic/expose_XIX_body.tex}; \nolinkurl{fr/sga7ii/current/expose_XIX_body.tex}\par

\par\medskip\noindent\hypertarget{app-F7II-XIX-014.2}{}\textbf{F7II-XIX-014.2}\quad p. 336; inherited corrected-layer difference.\par
\noindent Imprimé : {\ttfamily\footnotesize row\allowbreak{}}\par
\noindent Corrigé : {\ttfamily\footnotesize x=1.2cm,y=1.2cm]\allowbreak{} \textbackslash{}coordinate\allowbreak{}}\par
\noindent Source : ADOPT-F7II-XIX-014.2. \nolinkurl{fr/sga7ii/diplomatic/expose_XIX_body.tex}; \nolinkurl{fr/sga7ii/current/expose_XIX_body.tex}\par

\par\medskip\noindent\hypertarget{app-F7II-XIX-014.3}{}\textbf{F7II-XIX-014.3}\quad p. 336; inherited corrected-layer difference.\par
\noindent Imprimé : {\ttfamily\footnotesize sep=huge\allowbreak{}}\par
\noindent Corrigé : {\ttfamily\footnotesize (O)\allowbreak{} at\allowbreak{} (0\allowbreak{}}\par
\noindent Source : ADOPT-F7II-XIX-014.3. \nolinkurl{fr/sga7ii/diplomatic/expose_XIX_body.tex}; \nolinkurl{fr/sga7ii/current/expose_XIX_body.tex}\par

\par\medskip\noindent\hypertarget{app-F7II-XIX-014.4}{}\textbf{F7II-XIX-014.4}\quad p. 336; inherited corrected-layer difference.\par
\noindent Imprimé : \emph{(vide)}\par
\noindent Corrigé : {\ttfamily\footnotesize 0);\allowbreak{} \textbackslash{}coordinate\allowbreak{}}\par
\noindent Source : ADOPT-F7II-XIX-014.4. \nolinkurl{fr/sga7ii/diplomatic/expose_XIX_body.tex}; \nolinkurl{fr/sga7ii/current/expose_XIX_body.tex}\par

\par\medskip\noindent\hypertarget{app-F7II-XIX-014.5}{}\textbf{F7II-XIX-014.5}\quad p. 336; inherited corrected-layer difference.\par
\noindent Imprimé : {\ttfamily\footnotesize column\allowbreak{}}\par
\noindent Corrigé : {\ttfamily\footnotesize (L)\allowbreak{}}\par
\noindent Source : ADOPT-F7II-XIX-014.5. \nolinkurl{fr/sga7ii/diplomatic/expose_XIX_body.tex}; \nolinkurl{fr/sga7ii/current/expose_XIX_body.tex}\par

\par\medskip\noindent\hypertarget{app-F7II-XIX-014.6}{}\textbf{F7II-XIX-014.6}\quad p. 336; inherited corrected-layer difference.\par
\noindent Imprimé : {\ttfamily\footnotesize sep=huge\allowbreak{}}\par
\noindent Corrigé : {\ttfamily\footnotesize at\allowbreak{} (-1.8,1.8);\allowbreak{} \textbackslash{}coordinate\allowbreak{} (S2)\allowbreak{} at\allowbreak{} (0,2.8);\allowbreak{} \textbackslash{}coordinate\allowbreak{} (S1)\allowbreak{} at\allowbreak{} (2.2,2.0);\allowbreak{} \textbackslash{}draw\allowbreak{} (O)--(L);\allowbreak{} \textbackslash{}draw\allowbreak{} (O)--(S2);\allowbreak{} \textbackslash{}draw\allowbreak{} (O)--(S1);\allowbreak{} \textbackslash{}node[right\allowbreak{}}\par
\noindent Source : ADOPT-F7II-XIX-014.6. \nolinkurl{fr/sga7ii/diplomatic/expose_XIX_body.tex}; \nolinkurl{fr/sga7ii/current/expose_XIX_body.tex}\par

\par\medskip\noindent\hypertarget{app-F7II-XIX-014.7}{}\textbf{F7II-XIX-014.7}\quad p. 336; inherited corrected-layer difference.\par
\noindent Imprimé : \emph{(0 espace(s), 1 retour(s) de ligne dans le témoin TeX)}\par
\noindent Corrigé : {\ttfamily\footnotesize  at\allowbreak{} (O)\allowbreak{} }\par
\noindent Source : ADOPT-F7II-XIX-014.7. \nolinkurl{fr/sga7ii/diplomatic/expose_XIX_body.tex}; \nolinkurl{fr/sga7ii/current/expose_XIX_body.tex}\par

\par\medskip\noindent\hypertarget{app-F7II-XIX-014.8}{}\textbf{F7II-XIX-014.8}\quad p. 336; inherited corrected-layer difference.\par
\noindent Imprimé : \emph{(vide)}\par
\noindent Corrigé : {\ttfamily\footnotesize \$0\$\allowbreak{}}\par
\noindent Source : ADOPT-F7II-XIX-014.8. \nolinkurl{fr/sga7ii/diplomatic/expose_XIX_body.tex}; \nolinkurl{fr/sga7ii/current/expose_XIX_body.tex}\par

\par\medskip\noindent\hypertarget{app-F7II-XIX-014.9}{}\textbf{F7II-XIX-014.9}\quad p. 336; inherited corrected-layer difference.\par
\noindent Imprimé : \emph{(vide)}\par
\noindent Corrigé : {\ttfamily\footnotesize ;\allowbreak{} \textbackslash{}node[above\allowbreak{}}\par
\noindent Source : ADOPT-F7II-XIX-014.9. \nolinkurl{fr/sga7ii/diplomatic/expose_XIX_body.tex}; \nolinkurl{fr/sga7ii/current/expose_XIX_body.tex}\par

\par\medskip\noindent\hypertarget{app-F7II-XIX-014.10}{}\textbf{F7II-XIX-014.10}\quad p. 336; inherited corrected-layer difference.\par
\noindent Imprimé : {\ttfamily\footnotesize \&\allowbreak{}}\par
\noindent Corrigé : {\ttfamily\footnotesize right]\allowbreak{}}\par
\noindent Source : ADOPT-F7II-XIX-014.10. \nolinkurl{fr/sga7ii/diplomatic/expose_XIX_body.tex}; \nolinkurl{fr/sga7ii/current/expose_XIX_body.tex}\par

\par\medskip\noindent\hypertarget{app-F7II-XIX-014.11}{}\textbf{F7II-XIX-014.11}\quad p. 336; inherited corrected-layer difference.\par
\noindent Imprimé : \emph{(vide)}\par
\noindent Corrigé : {\ttfamily\footnotesize at\allowbreak{} (S2)\allowbreak{} \{\$\allowbreak{}}\par
\noindent Source : ADOPT-F7II-XIX-014.11. \nolinkurl{fr/sga7ii/diplomatic/expose_XIX_body.tex}; \nolinkurl{fr/sga7ii/current/expose_XIX_body.tex}\par

\par\medskip\noindent\hypertarget{app-F7II-XIX-014.12}{}\textbf{F7II-XIX-014.12}\quad p. 336; inherited corrected-layer difference.\par
\noindent Imprimé : \emph{(vide)}\par
\noindent Corrigé : {\ttfamily\footnotesize \$\};\allowbreak{} \textbackslash{}node[above\allowbreak{}}\par
\noindent Source : ADOPT-F7II-XIX-014.12. \nolinkurl{fr/sga7ii/diplomatic/expose_XIX_body.tex}; \nolinkurl{fr/sga7ii/current/expose_XIX_body.tex}\par

\par\medskip\noindent\hypertarget{app-F7II-XIX-014.13}{}\textbf{F7II-XIX-014.13}\quad p. 336; inherited corrected-layer difference.\par
\noindent Imprimé : {\ttfamily\footnotesize \&\allowbreak{}}\par
\noindent Corrigé : {\ttfamily\footnotesize right]\allowbreak{}}\par
\noindent Source : ADOPT-F7II-XIX-014.13. \nolinkurl{fr/sga7ii/diplomatic/expose_XIX_body.tex}; \nolinkurl{fr/sga7ii/current/expose_XIX_body.tex}\par

\par\medskip\noindent\hypertarget{app-F7II-XIX-014.14}{}\textbf{F7II-XIX-014.14}\quad p. 336; inherited corrected-layer difference.\par
\noindent Imprimé : \emph{(vide)}\par
\noindent Corrigé : {\ttfamily\footnotesize at\allowbreak{} (S1)\allowbreak{} \{\$\allowbreak{}}\par
\noindent Source : ADOPT-F7II-XIX-014.14. \nolinkurl{fr/sga7ii/diplomatic/expose_XIX_body.tex}; \nolinkurl{fr/sga7ii/current/expose_XIX_body.tex}\par

\par\medskip\noindent\hypertarget{app-F7II-XIX-014.15}{}\textbf{F7II-XIX-014.15}\quad p. 336; inherited corrected-layer difference.\par
\noindent Imprimé : \emph{(vide)}\par
\noindent Corrigé : {\ttfamily\footnotesize \$\};\allowbreak{} \textbackslash{}node[right]\allowbreak{}}\par
\noindent Source : ADOPT-F7II-XIX-014.15. \nolinkurl{fr/sga7ii/diplomatic/expose_XIX_body.tex}; \nolinkurl{fr/sga7ii/current/expose_XIX_body.tex}\par

\par\medskip\noindent\hypertarget{app-F7II-XIX-014.16}{}\textbf{F7II-XIX-014.16}\quad p. 336; inherited corrected-layer difference.\par
\noindent Imprimé : {\ttfamily\footnotesize \textbackslash{}\textbackslash{}\allowbreak{} \textbackslash{}cdots\allowbreak{}}\par
\noindent Corrigé : {\ttfamily\footnotesize at\allowbreak{}}\par
\noindent Source : ADOPT-F7II-XIX-014.16. \nolinkurl{fr/sga7ii/diplomatic/expose_XIX_body.tex}; \nolinkurl{fr/sga7ii/current/expose_XIX_body.tex}\par

\par\medskip\noindent\hypertarget{app-F7II-XIX-014.17}{}\textbf{F7II-XIX-014.17}\quad p. 336; inherited corrected-layer difference.\par
\noindent Imprimé : {\ttfamily\footnotesize \&\allowbreak{}}\par
\noindent Corrigé : {\ttfamily\footnotesize (0,1.65)\allowbreak{}}\par
\noindent Source : ADOPT-F7II-XIX-014.17. \nolinkurl{fr/sga7ii/diplomatic/expose_XIX_body.tex}; \nolinkurl{fr/sga7ii/current/expose_XIX_body.tex}\par

\par\medskip\noindent\hypertarget{app-F7II-XIX-014.18}{}\textbf{F7II-XIX-014.18}\quad p. 336; inherited corrected-layer difference.\par
\noindent Imprimé : {\ttfamily\footnotesize 0\allowbreak{} \textbackslash{}arrow[u,\allowbreak{} dash,\allowbreak{} "\allowbreak{}}\par
\noindent Corrigé : \emph{(vide)}\par
\noindent Source : ADOPT-F7II-XIX-014.18. \nolinkurl{fr/sga7ii/diplomatic/expose_XIX_body.tex}; \nolinkurl{fr/sga7ii/current/expose_XIX_body.tex}\par

\par\medskip\noindent\hypertarget{app-F7II-XIX-014.19}{}\textbf{F7II-XIX-014.19}\quad p. 336; inherited corrected-layer difference.\par
\noindent Imprimé : \emph{(vide)}\par
\noindent Corrigé : {\ttfamily\footnotesize \$\allowbreak{}}\par
\noindent Source : ADOPT-F7II-XIX-014.19. \nolinkurl{fr/sga7ii/diplomatic/expose_XIX_body.tex}; \nolinkurl{fr/sga7ii/current/expose_XIX_body.tex}\par

\par\medskip\noindent\hypertarget{app-F7II-XIX-014.20}{}\textbf{F7II-XIX-014.20}\quad p. 336; inherited corrected-layer difference.\par
\noindent Imprimé : \emph{(vide)}\par
\noindent Corrigé : {\ttfamily\footnotesize \$\allowbreak{}}\par
\noindent Source : ADOPT-F7II-XIX-014.20. \nolinkurl{fr/sga7ii/diplomatic/expose_XIX_body.tex}; \nolinkurl{fr/sga7ii/current/expose_XIX_body.tex}\par

\par\medskip\noindent\hypertarget{app-F7II-XIX-014.21}{}\textbf{F7II-XIX-014.21}\quad p. 336; inherited corrected-layer difference.\par
\noindent Imprimé : {\ttfamily\footnotesize "'\allowbreak{}}\par
\noindent Corrigé : {\ttfamily\footnotesize ;\allowbreak{} \textbackslash{}node[below\allowbreak{} right\allowbreak{}}\par
\noindent Source : ADOPT-F7II-XIX-014.21. \nolinkurl{fr/sga7ii/diplomatic/expose_XIX_body.tex}; \nolinkurl{fr/sga7ii/current/expose_XIX_body.tex}\par

\par\medskip\noindent\hypertarget{app-F7II-XIX-014.22}{}\textbf{F7II-XIX-014.22}\quad p. 336; inherited corrected-layer difference.\par
\noindent Imprimé : {\ttfamily\footnotesize \textbackslash{}arrow[ur\allowbreak{}}\par
\noindent Corrigé : {\ttfamily\footnotesize at\allowbreak{} (1.25\allowbreak{}}\par
\noindent Source : ADOPT-F7II-XIX-014.22. \nolinkurl{fr/sga7ii/diplomatic/expose_XIX_body.tex}; \nolinkurl{fr/sga7ii/current/expose_XIX_body.tex}\par

\par\medskip\noindent\hypertarget{app-F7II-XIX-014.23}{}\textbf{F7II-XIX-014.23}\quad p. 336; inherited corrected-layer difference.\par
\noindent Imprimé : \emph{(vide)}\par
\noindent Corrigé : {\ttfamily\footnotesize 1.14)\allowbreak{}}\par
\noindent Source : ADOPT-F7II-XIX-014.23. \nolinkurl{fr/sga7ii/diplomatic/expose_XIX_body.tex}; \nolinkurl{fr/sga7ii/current/expose_XIX_body.tex}\par

\par\medskip\noindent\hypertarget{app-F7II-XIX-014.24}{}\textbf{F7II-XIX-014.24}\quad p. 336; inherited corrected-layer difference.\par
\noindent Imprimé : {\ttfamily\footnotesize dash,\allowbreak{} "\allowbreak{}}\par
\noindent Corrigé : \emph{(vide)}\par
\noindent Source : ADOPT-F7II-XIX-014.24. \nolinkurl{fr/sga7ii/diplomatic/expose_XIX_body.tex}; \nolinkurl{fr/sga7ii/current/expose_XIX_body.tex}\par

\par\medskip\noindent\hypertarget{app-F7II-XIX-014.25}{}\textbf{F7II-XIX-014.25}\quad p. 336; inherited corrected-layer difference.\par
\noindent Imprimé : \emph{(vide)}\par
\noindent Corrigé : {\ttfamily\footnotesize \$\allowbreak{}}\par
\noindent Source : ADOPT-F7II-XIX-014.25. \nolinkurl{fr/sga7ii/diplomatic/expose_XIX_body.tex}; \nolinkurl{fr/sga7ii/current/expose_XIX_body.tex}\par

\par\medskip\noindent\hypertarget{app-F7II-XIX-014.26}{}\textbf{F7II-XIX-014.26}\quad p. 336; inherited corrected-layer difference.\par
\noindent Imprimé : \emph{(vide)}\par
\noindent Corrigé : {\ttfamily\footnotesize \$\allowbreak{}}\par
\noindent Source : ADOPT-F7II-XIX-014.26. \nolinkurl{fr/sga7ii/diplomatic/expose_XIX_body.tex}; \nolinkurl{fr/sga7ii/current/expose_XIX_body.tex}\par

\par\medskip\noindent\hypertarget{app-F7II-XIX-014.27}{}\textbf{F7II-XIX-014.27}\quad p. 336; inherited corrected-layer difference.\par
\noindent Imprimé : {\ttfamily\footnotesize "']\allowbreak{}}\par
\noindent Corrigé : {\ttfamily\footnotesize ;\allowbreak{} \textbackslash{}fill\allowbreak{}}\par
\noindent Source : ADOPT-F7II-XIX-014.27. \nolinkurl{fr/sga7ii/diplomatic/expose_XIX_body.tex}; \nolinkurl{fr/sga7ii/current/expose_XIX_body.tex}\par

\par\medskip\noindent\hypertarget{app-F7II-XIX-014.28}{}\textbf{F7II-XIX-014.28}\quad p. 336; inherited corrected-layer difference.\par
\noindent Imprimé : {\ttfamily\footnotesize \textbackslash{}arrow[ul\allowbreak{}}\par
\noindent Corrigé : {\ttfamily\footnotesize (-0.73\allowbreak{}}\par
\noindent Source : ADOPT-F7II-XIX-014.28. \nolinkurl{fr/sga7ii/diplomatic/expose_XIX_body.tex}; \nolinkurl{fr/sga7ii/current/expose_XIX_body.tex}\par

\par\medskip\noindent\hypertarget{app-F7II-XIX-014.29}{}\textbf{F7II-XIX-014.29}\quad p. 336; inherited corrected-layer difference.\par
\noindent Imprimé : \emph{(vide)}\par
\noindent Corrigé : {\ttfamily\footnotesize 0.76)\allowbreak{}}\par
\noindent Source : ADOPT-F7II-XIX-014.29. \nolinkurl{fr/sga7ii/diplomatic/expose_XIX_body.tex}; \nolinkurl{fr/sga7ii/current/expose_XIX_body.tex}\par

\par\medskip\noindent\hypertarget{app-F7II-XIX-014.30}{}\textbf{F7II-XIX-014.30}\quad p. 336; designated-witness attribution.\par
\noindent Imprimé : {\ttfamily\footnotesize dash]\allowbreak{}}\par
\noindent Corrigé : {\ttfamily\footnotesize circle\allowbreak{} (0.035);\allowbreak{} \textbackslash{}fill\allowbreak{} (-0.55,0.91)\allowbreak{} circle\allowbreak{} (0.035);\allowbreak{} \textbackslash{}fill\allowbreak{} (-0.37,1.03)\allowbreak{} circle\allowbreak{} (0.035);\allowbreak{}}\par
\noindent Source : WITNESS-F7II-XIX-014.30. \nolinkurl{fr/sga7ii/diplomatic/expose_XIX_body.tex}; \nolinkurl{fr/sga7ii/current/expose_XIX_body.tex}\par

\par\medskip\noindent\hypertarget{app-F7II-XIX-014.31}{}\textbf{F7II-XIX-014.31}\quad p. 336; designated-witness attribution.\par
\noindent Imprimé : {\ttfamily\footnotesize tikzcd\allowbreak{}}\par
\noindent Corrigé : {\ttfamily\footnotesize tikzpicture\allowbreak{}}\par
\noindent Source : WITNESS-F7II-XIX-014.31. \nolinkurl{fr/sga7ii/diplomatic/expose_XIX_body.tex}; \nolinkurl{fr/sga7ii/current/expose_XIX_body.tex}\par

\par\medskip\noindent\hypertarget{app-F7II-XIX-014.32}{}\textbf{F7II-XIX-014.32}\quad p. 336; printed misprint.\par
\noindent Imprimé : {\ttfamily\footnotesize étudié\allowbreak{}}\par
\noindent Corrigé : {\ttfamily\footnotesize étudiée\allowbreak{}}\par
\noindent Source : 7II:995. \nolinkurl{fixes/MISPRINTS_7II.json#995}\par

\par\medskip\noindent\hypertarget{app-F7II-XIX-015.1}{}\textbf{F7II-XIX-015.1}\quad p. 337; misprint\_pass.\par
\noindent Imprimé : {\ttfamily\footnotesize disque\allowbreak{}}\par
\noindent Corrigé : {\ttfamily\footnotesize disques\allowbreak{}}\par
\noindent Source : 7II:592. \nolinkurl{fixes/FIXLOG_7II_CURRENT.jsonl#765}\par

\par\medskip\noindent\hypertarget{app-F7II-XIX-015.2}{}\textbf{F7II-XIX-015.2}\quad p. 337; inherited corrected-layer difference.\par
\noindent Imprimé : {\ttfamily\footnotesize V\allowbreak{}}\par
\noindent Corrigé : {\ttfamily\footnotesize H\allowbreak{}}\par
\noindent Source : ADOPT-F7II-XIX-015.2. \nolinkurl{fr/sga7ii/diplomatic/expose_XIX_body.tex}; \nolinkurl{fr/sga7ii/current/expose_XIX_body.tex}\par

\par\medskip\noindent\hypertarget{app-F7II-XIX-015.3}{}\textbf{F7II-XIX-015.3}\quad p. 337; inherited corrected-layer difference.\par
\noindent Imprimé : {\ttfamily\footnotesize V\allowbreak{}}\par
\noindent Corrigé : {\ttfamily\footnotesize H\allowbreak{}}\par
\noindent Source : ADOPT-F7II-XIX-015.3. \nolinkurl{fr/sga7ii/diplomatic/expose_XIX_body.tex}; \nolinkurl{fr/sga7ii/current/expose_XIX_body.tex}\par

\par\medskip\noindent\hypertarget{app-F7II-XIX-015.4}{}\textbf{F7II-XIX-015.4}\quad p. 337; misprint\_in\_print.\par
\noindent Imprimé : {\ttfamily\footnotesize \textbackslash{}Lambda\allowbreak{}}\par
\noindent Corrigé : {\ttfamily\footnotesize \textbackslash{}wedge\allowbreak{}}\par
\noindent Source : FIX-7II-D-0738. \nolinkurl{fixes/FIXLOG_7II.jsonl#738}\par

\par\medskip\noindent\hypertarget{app-F7II-XIX-016.1}{}\textbf{F7II-XIX-016.1}\quad p. 338; printed misprint.\par
\noindent Imprimé : {\ttfamily\footnotesize \{2n\}\allowbreak{}}\par
\noindent Corrigé : {\ttfamily\footnotesize n\allowbreak{}}\par
\noindent Source : 7II:998. \nolinkurl{fixes/MISPRINTS_7II.json#998}\par

\par\medskip\noindent\hypertarget{app-F7II-XIX-016.2}{}\textbf{F7II-XIX-016.2}\quad p. 338; designated-witness attribution.\par
\noindent Imprimé : {\ttfamily\footnotesize by\allowbreak{}}\par
\noindent Corrigé : {\ttfamily\footnotesize b\allowbreak{} y\allowbreak{}}\par
\noindent Source : WITNESS-F7II-XIX-016.2. \nolinkurl{fr/sga7ii/diplomatic/expose_XIX_body.tex}; \nolinkurl{fr/sga7ii/current/expose_XIX_body.tex}\par

\par\medskip\noindent\hypertarget{app-F7II-XIX-017.1}{}\textbf{F7II-XIX-017.1}\quad p. 339; printed misprint.\par
\noindent Imprimé : {\ttfamily\footnotesize uhe\allowbreak{}}\par
\noindent Corrigé : {\ttfamily\footnotesize une\allowbreak{}}\par
\noindent Source : 7II:999. \nolinkurl{fixes/MISPRINTS_7II.json#999}\par

\par\medskip\noindent\hypertarget{app-F7II-XIX-018.1}{}\textbf{F7II-XIX-018.1}\quad p. 339; printed misprint.\par
\noindent Imprimé : {\ttfamily\footnotesize racine\allowbreak{}}\par
\noindent Corrigé : {\ttfamily\footnotesize racines\allowbreak{}}\par
\noindent Source : 7II:595; 7II:1000. \nolinkurl{fixes/MISPRINTS_7II.json#595}; \nolinkurl{fixes/MISPRINTS_7II.json#1000}\par

\par\medskip\noindent\hypertarget{app-F7II-XIX-018.2}{}\textbf{F7II-XIX-018.2}\quad p. 339; printed misprint.\par
\noindent Imprimé : {\ttfamily\footnotesize Weil\allowbreak{}}\par
\noindent Corrigé : {\ttfamily\footnotesize Weyl\allowbreak{}}\par
\noindent Source : 7II:1001. \nolinkurl{fixes/MISPRINTS_7II.json#1001}\par

\par\medskip\noindent\hypertarget{app-F7II-XIX-018.3}{}\textbf{F7II-XIX-018.3}\quad p. 339; printed misprint.\par
\noindent Imprimé : {\ttfamily\footnotesize 4\allowbreak{}}\par
\noindent Corrigé : {\ttfamily\footnotesize 5\allowbreak{}}\par
\noindent Source : 7II:597; 7II:1002. \nolinkurl{fixes/MISPRINTS_7II.json#597}; \nolinkurl{fixes/MISPRINTS_7II.json#1002}\par

\par\medskip\noindent\hypertarget{app-F7II-XIX-018.4}{}\textbf{F7II-XIX-018.4}\quad p. 339; printed misprint.\par
\noindent Imprimé : {\ttfamily\footnotesize 4\allowbreak{}}\par
\noindent Corrigé : {\ttfamily\footnotesize 5\allowbreak{}}\par
\noindent Source : 7II:598; 7II:1003. \nolinkurl{fixes/MISPRINTS_7II.json#598}; \nolinkurl{fixes/MISPRINTS_7II.json#1003}\par

\par\medskip\noindent\hypertarget{app-F7II-XIX-018.5}{}\textbf{F7II-XIX-018.5}\quad p. 339; printed misprint.\par
\noindent Imprimé : {\ttfamily\footnotesize 4\allowbreak{}}\par
\noindent Corrigé : {\ttfamily\footnotesize 5\allowbreak{}}\par
\noindent Source : 7II:599; 7II:1004. \nolinkurl{fixes/MISPRINTS_7II.json#599}; \nolinkurl{fixes/MISPRINTS_7II.json#1004}\par

\normalsize\fi
\clearpage\SGAtoc{Exposé XX — LE THEOREME DE GRIFFITHS}
\selectlanguage{french}

\SGApage{II}{341}


S G A 7


\underline{E X P O S E \ XX}


\underline{\underline{LE THEOREME DE GRIFFITHS}}


\underline{par N. KATZ}


\SGAsectiontoc{0. Introduction}{0}
0. \underline{Introduction}

\SGAvar{On démontre le théorème de GRIFFITHS [1], qui donne une facon systématique (bien que non-constructive) de ``trouver'' des cycles algèbriques qui sont $\mathbb{Q}_\ell$-cohomologues à zéro, mais donc aucun multiple n'est algébriquement équivalent à zéro. La démonstration donnée ici (due à GROTHENDIECK) est la traduction en termes purement algébriques de la démonstration originelle, plus au moins transcendante, de GRIFFITHS. Dans toutes les deux, c'est le théorème d'irréductibilité XVIII 6.7 qui est un point-clef (cf.3.1.6 et 3.2).


}{On démontre le théorème de GRIFFITHS [1], qui donne une façon systématique (bien que non-constructive) de ``trouver'' des cycles algébriques qui sont $\mathbb{Q}_\ell$-cohomologues à zéro, mais dont aucun multiple n'est algébriquement équivalent à zéro. La démonstration donnée ici (due à GROTHENDIECK) est la traduction en termes purement algébriques de la démonstration originelle, plus ou moins transcendante, de GRIFFITHS. Dans les deux, c'est le théorème d'irréductibilité XVIII 6.7 qui est un point-clef (cf. 3.1.6 et 3.2).


}{F7II-XX-001}\SGAsectiontoc{1. Le formalisme des classes primitives}{1}
1. \underline{Le formalisme des classes primitives}

\SGAvar{Dans toute la suite, on suppose fixé un nombre premier $\ell$ . Tous les schémas envisagés sont supposés à caractéristiques $\neq \ell$ .

}{Dans toute la suite, on suppose fixé un nombre premier $\ell$. Tous les schémas envisagés sont supposés de caractéristique $\neq \ell$.

}{F7II-XX-002}Considérons la situation

\begin{equation}\tag{1.0.1}
\begin{tikzcd}[column sep=huge, row sep=huge]
X \arrow[rr, "\pi"] \arrow[dr, "f"'] \& {} \& Y \arrow[dl, "g"] \\
{} \& S \&
{} \end{tikzcd}
\end{equation}









dans laquelle on suppose 


\SGApage{II}{342}

\SGAvar{\begin{equation}\tag{1.0.3}
f,\ g \quad \text{lisses de dimensions relatives} \quad \ell \quad d_1 \ \text{et} \ d_2 \ ,
\end{equation}
\begin{equation}\tag{1.0.4}
\pi \quad \text{propre} \ .
\end{equation}
Compte tenu de ce que $X$ et $Y$ sont lisses sur $S$ , la propreté de $\pi$ nous donne un \underline{morphisme de Gysin} (SGA 5 IV \qquad )

}{\begin{equation}\tag{1.0.3}
f,\ g \quad \text{lisses de dimensions relatives} \quad d_1 \ \text{et} \ d_2 \ ,
\end{equation}
\begin{equation}\tag{1.0.4}
\pi \quad \text{propre} \ .
\end{equation}
Compte tenu de ce que $X$ et $Y$ sont lisses sur $S$ , la propreté de $\pi$ nous donne un \underline{morphisme de Gysin} (SGA 5 IV \qquad )

}{F7II-XX-003}\begin{equation}\tag{1.0.5}
\pi_* : H^q(X,\ \mathbb{Q}_\ell(i)) \longrightarrow H^{q-2(d_1-d_2)}(Y,\ \mathbb{Q}_\ell(i-(d_1-d_2))) \ .
\end{equation}
Localisant sur $S$ (au sens étale), on obtient également un morphisme de faisceaux sur $S$

\SGAvar{\begin{equation}\tag{1.0.6}
\pi_* : R^q f_*(\mathbb{Q}_\ell(i) \longrightarrow R^{q-2(d_1-d_2)} g_*(\mathbb{Q}_\ell(i-d_1-d_2)) \ ,
\end{equation}
et, appliquant le foncteur $H^p(S, \qquad )$ , un morphisme

\begin{equation}\tag{1.0.7}
\pi_*: H^p(S,R^q f_*\mathbb{Q}_\ell(i)) \longrightarrow H^p(S,R^{q-2(d_1-d_2)} g_*(\mathbb{Q}_\ell(i-d_1+d_2))) \ .
\end{equation}
\underline{Définition} 1.0.8. \underline{Soient} $(E_r^{p,q}\ ,\ d_r^{p,q})$ \underline{et} $({}'E_r^{p,q}\ ,\ {}'d_r^{p,q})$ \underline{deux suites spectrales (dans une catégorie abélienne donnée)}, \underline{et} $(a, b) \in \mathbb{Z} \times \mathbb{Z}$ \underline{Le décalage de} $({}'E_r^{p,q}\ ,\ {}'d_r^{p,q})$ \underline{par} $(a, b)$ \underline{est la suite spectrale} $({}''E_r^{p,q},\ {}''d^{p,q})$ , \underline{avec}

\begin{equation}\tag{1.0.9}
\left\{
\begin{aligned}
{}''E_r^{p,q} &= {}'E_r^{p+a,q+b} \\
{}''d_r^{p,q} &= d_r^{p+a,q+b}
\end{aligned}
\right.
\end{equation}

}{\begin{equation}\tag{1.0.6}
\pi_* : R^q f_*\mathbb{Q}_\ell(i) \longrightarrow R^{q-2(d_1-d_2)} g_*\mathbb{Q}_\ell(i-d_1+d_2) \ ,
\end{equation}
et, appliquant le foncteur $H^p(S, \qquad )$ , un morphisme

\begin{equation}\tag{1.0.7}
\pi_*: H^p(S,R^q f_*\mathbb{Q}_\ell(i)) \longrightarrow H^p(S,R^{q-2(d_1-d_2)} g_*(\mathbb{Q}_\ell(i-d_1+d_2))) \ .
\end{equation}
\underline{Définition} 1.0.8. \underline{Soient} $(E_r^{p,q}\ ,\ d_r^{p,q})$ \underline{et} $({}'E_r^{p,q}\ ,\ {}'d_r^{p,q})$ \underline{deux suites spectrales (dans une catégorie abélienne donnée)}, \underline{et} $(a, b) \in \mathbb{Z} \times \mathbb{Z}$. \underline{Le décalage de} $({}'E_r^{p,q}\ ,\ {}'d_r^{p,q})$ \underline{par} $(a, b)$ \underline{est la suite spectrale} $({}''E_r^{p,q},\ {}''d_r^{p,q})$ , \underline{avec}

\begin{equation}\tag{1.0.9}
\left\{
\begin{aligned}
{}''E_r^{p,q} &= {}'E_r^{p+a,q+b} \\
{}''d_r^{p,q} &= {}'d_r^{p+a,q+b}
\end{aligned}
\right.
\end{equation}

}{F7II-XX-004}\SGApage{II}{343}

1.0.10. \underline{Un morphisme de suites spectrales de} $(E_r^{p,q}\ ,\ d_r^{p,q})$ \underline{dans} $({}''E_r^{p,q}\ ,\ {}''d_r^{p,q})$ \underline{s'appelle un morphisme de bidegré} $(a, b)$ \underline{de} $(E_r^{p,q}\ ,\ d_r^{p,q})$ \underline{dans} $({}'E_r^{p,q}\ ,\ {}'d_r^{p,q})$ .

Le lemme suivant est une conséquence triviale de la théorie du \underline{morphisme de Gysin} (SGA 5, IV).

\SGAvar{\underline{Lemme} 1.1. \underline{Il existe un morphisme canonique} $\pi_*$ \underline{de bidegré} $(a,-2(d_1-d_2))$ \underline{des suites spectrales de Leray}

}{\underline{Lemme} 1.1. \underline{Il existe un morphisme canonique} $\pi_*$ \underline{de bidegré} $(0,-2(d_1-d_2))$ \underline{des suites spectrales de Leray}

}{F7II-XX-005}\begin{equation}\tag{1.1.0}
R^p \alpha_* R^q f_* \mathbb{Q}_\ell(i) \Longrightarrow R^{p+q}(\alpha f)_* \ \mathbb{Q}_\ell(i)
\end{equation}
\underline{et}

\begin{equation}\tag{1.1.1}
R^p \alpha_* R^q g_* \mathbb{Q}_\ell(i-d_1+d_2) \Longrightarrow R^{p+q}(\alpha g)_* \ \mathbb{Q}_\ell(i-d_1+d_2) \quad ,
\end{equation}
\underline{qui sur le terme} $E_2$ \underline{est donné par} (1.0.7), \underline{et sur l'aboutissement par} (1.0.5).

1.2. Considérons maintenant la situation

\begin{equation}\tag{1.2.1}
\begin{tikzcd}[column sep=large, row sep=large]
{} \& X \times_S Y \arrow[dl, "\mathrm{pr}_1"'] \arrow[dr, "\mathrm{pr}_2"] \& {} \\
X \arrow[dr, "f"'] \& {} \& Y \arrow[dl, "g"] \\
{} \& S \arrow[d, "\alpha"] \& {} \\
{} \& T \&
{} \end{tikzcd}
\end{equation}














\SGAvar{,
\SGApage{II}{344}
dans laquelle $S$, $X$, $f$, $Y$, $g$ sont comme dans (1.0.2) - (1.0.4), mais en supposant maintenant également $f$ propre.

}{,
\SGApage{II}{344}
dans laquelle $S$, $X$, $f$, $Y$, $g$ sont comme dans (1.0.1) - (1.0.4), mais en supposant maintenant également $f$ propre.

}{F7II-XX-006}Soit

\begin{equation}\tag{1.2.2}
z \in H^a(X \times_S Y, \mathbb{Q}_\ell(b)).
\end{equation}
Alors $z$ induit une application (encore notée $z$ )

\begin{equation}\tag{1.2.3}
z : H^q(X, \mathbb{Q}_\ell(i)) \longrightarrow H^{q+a-2d_1} (Y, \mathbb{Q}_\ell(i+b-d_1))
\end{equation}
par la formule

\SGAvar{\begin{equation}\tag{1.2.4}
z(x) = \mathrm{pr}_{2*}(z \wedge \mathrm{pr}_1^* (x)) \qquad ,
\end{equation}
où $\mathrm{pr}_{2*}$ est l'homomorphisme de Gysin  Par localisation étale sur $S$ , $z$ définit également

}{\begin{equation}\tag{1.2.4}
z(x) = \mathrm{pr}_{2*}(z \wedge \mathrm{pr}_1^* (x)) \qquad ,
\end{equation}
où $\mathrm{pr}_{2*}$ est l'homomorphisme de Gysin. Par localisation étale sur $S$ , $z$ définit également

}{F7II-XX-007}\begin{equation}\tag{1.2.5}
z : R^q f_* \mathbb{Q}_\ell(i) \longrightarrow R^{q+a-2d_1} g_* (\mathbb{Q}_\ell(i+b-d_1)) \quad ,
\end{equation}
et, appliquant $H^p(S, \text{---} )$ , des applications

\SGAvar{\begin{equation}\tag{1.2.6}
H^p(S, R^q f_* \mathbb{Q}_\ell(i)) \longrightarrow H^p(S, R^{q+a-2d_1} g_*(\mathbb{Q}_\ell(i+b-d_1)) \ .
\end{equation}
\underline{Lemme} 1.3.  \underline{La classe}  $z$  \underline{définit un morphisme de bidegré} $(0,\ b-2d_1)$ \underline{des suites spectrales de Leray}

}{\begin{equation}\tag{1.2.6}
H^p(S, R^q f_* \mathbb{Q}_\ell(i)) \longrightarrow H^p(S, R^{q+a-2d_1} g_*(\mathbb{Q}_\ell(i+b-d_1))) \ .
\end{equation}
\underline{Lemme} 1.3.  \underline{La classe}  $z$  \underline{définit un morphisme de bidegré} $(0,\ a-2d_1)$ \underline{des suites spectrales de Leray}

}{F7II-XX-008}\begin{equation}\tag{1.2.7}
H^p(S, R^q f_* \ \mathbb{Q}_\ell(i)) \quad \Rightarrow \quad H^{p+q}(X, \mathbb{Q}_\ell(i))
\end{equation}
et

\begin{equation}\tag{1.2.8}
H^p(S, R^q g_* \ \mathbb{Q}_\ell(i+b-d_1)) \Rightarrow H^{p+q}(Y, \mathbb{Q}_\ell(i+b-d_1)) \quad , 
\end{equation}

\SGApage{II}{345}

 \underline{qui sur les termes} $E_2$ \underline{est donné par} (1.2.6), \underline{et par} (1.2.3) \underline{sur les aboutissements}.

\underline{Démonstration}. Compte tenu de la formule (1.2.4) et de ce que la "multiplication par $z$ " définit un morphisme de bidegré $(0,\ a)$ de suites spectrales de Leray

\SGAvar{\begin{equation}\tag{1.2.9}
H^p(S,\ R^q(\mathrm{pr}_1\ f)_* \ \mathbb{Q}_\ell(i)) \quad \Rightarrow \quad H^{p+q}(X \times_S Y,\ \mathbb{Q}_\ell(i))
\end{equation}
\begin{center}
\begin{tikzcd}
{} \arrow[d, "z"'] \\
{}
\end{tikzcd}
\end{center}







\begin{equation}\tag{1.2.10}
H^p(S,\ R^{q+a}(\mathrm{pr}_1\ f)_* \ \mathbb{Q}_\ell(i+b)) \Rightarrow H^{p+q+a}(X \times_S Y,\ \mathbb{Q}_\ell(i+b)) \quad ,
\end{equation}
on n'a qu'à appliquer 1.1. à la situation

\begin{equation}\tag{1.2.11}
\begin{tikzcd}[column sep=large, row sep=large]
X \times_S Y \arrow[rr, "\mathrm{pr}_2"] \arrow[dr, "f\ \mathrm{pr}_1"] \& {} \& Y \arrow[dl, "g"] \\
{} \& S \&
{} \end{tikzcd}
\end{equation}










}{\begin{equation}\tag{1.2.9}
H^p(S,\ R^q(f\circ\mathrm{pr}_1)_* \ \mathbb{Q}_\ell(i)) \quad \Rightarrow \quad H^{p+q}(X \times_S Y,\ \mathbb{Q}_\ell(i))
\end{equation}
\begin{center}
\begin{tikzcd}
{} \arrow[d, "z"'] \\
{}
\end{tikzcd}
\end{center}







\begin{equation}\tag{1.2.10}
H^p(S,\ R^{q+a}(f\circ\mathrm{pr}_1)_* \ \mathbb{Q}_\ell(i+b)) \Rightarrow H^{p+q+a}(X \times_S Y,\ \mathbb{Q}_\ell(i+b)) \quad ,
\end{equation}
on n'a qu'à appliquer 1.1. à la situation

\begin{equation}\tag{1.2.11}
\begin{tikzcd}[column sep=large, row sep=large]
X \times_S Y \arrow[rr, "\mathrm{pr}_2"] \arrow[dr, "f\circ\mathrm{pr}_1"] \& {} \& Y \arrow[dl, "g"] \\
{} \& S \&
{} \end{tikzcd}
\end{equation}










}{F7II-XX-009}.

\underline{Corollaire} 1.4.  \underline{Dans la situation} (1.2.1) - (1.2.2), \underline{le morphisme} (1.2.3)

\SGAvar{\begin{equation}\tag{1.4.1}
z : H^q(X,\ \mathbb{Q}_\ell(i)) \longrightarrow H^{q+a-2d_1} (Y,\ \mathbb{Q}_\ell(i+b-d_1))
\end{equation}
\underline{transforme classe primitive} (XVIII. 5.8.1) \underline{en classe primitive}, \underline{et le diagramme suivant est commutatif}.


}{\begin{equation}\tag{1.4.1}
z : H^q(X,\ \mathbb{Q}_\ell(i)) \longrightarrow H^{q+a-2d_1} (Y,\ \mathbb{Q}_\ell(i+b-d_1))
\end{equation}
\underline{transforme toute classe primitive} (XVIII. 5.8.1) \underline{en classe primitive}, \underline{et le diagramme suivant est commutatif}.


}{F7II-XX-010}\SGApage{II}{346}

\begin{equation}\tag{1.4.2}
\begin{tikzcd}[column sep=large, row sep=large]
\mathrm{Prim}^{q}(X/S,\ \mathbb{Q}_\ell(i)) \arrow[r, "z"] \arrow[d, "u_f"] \& \mathrm{Prim}^{q+a-2d_1}(Y/S,\ \mathbb{Q}_\ell(i+b-d_1)) \arrow[d, "u_g"] \\
H^{1}(S,\ R^{q-1}f_*\mathbb{Q}_\ell(i)) \arrow[r, "z"] \& H^{1}(S,\ R^{q+a-2d_1-1}\,g_*\mathbb{Q}_\ell(i+b-d_1))
\end{tikzcd}
\end{equation}









\underline{(où} $u_f$, $u_g$ \underline{sont les homomorphismes canoniques de} XVIII 5.8.1).

1.5. Soit $\bar{\eta}$ un point géométrique de $S$ (au-dessus d'un point $\eta$ de $S$ ), et soit

\begin{equation}\tag{1.5.1}
z_{\bar{\eta}} \ \in \ H^{a}(X_{\bar{\eta}} \times_{k(\bar{\eta})} Y_{\bar{\eta}} \ , \ \mathbb{Q}_\ell(b))
\end{equation}
la classe induite par $z$ sur la fibre géométrique correspondante de

\SGAvar{\begin{equation}\tag{1.5.2}
X \times_S Y \ \longrightarrow \ S \quad .
\end{equation}
\underline{Proposition} 1.5.3. \underline{Supposons} $f$, $g$ \underline{propres,} $S$ \underline{annexe}, \underline{et soit} $q \in \mathbb{Z}$ . \underline{Si la classe} $z_{\bar{\eta}}$ \underline{induit l'homomorphisme nul}

\begin{equation}\tag{1.5.3}
0 = z_{\bar{\eta}} \ : \ H^{q-1}(X_{\bar{\eta}} \ , \ \mathbb{Q}_\ell(i)) \ \longrightarrow \ H^{q+a-2d_1-1}(Y_{\bar{\eta}}, \ \mathbb{Q}_\ell(i+b-d_1)) \ ,
\end{equation}
\underline{alors, pour tout} $X \in \mathrm{Prim}^{q}(X/S,\ \mathbb{Q}_\ell(i))$ , \underline{on a}

}{\begin{equation}\tag{1.5.2}
X \times_S Y \ \longrightarrow \ S \quad .
\end{equation}
\underline{Proposition} 1.5.3. \underline{Supposons} $f$, $g$ \underline{propres,} $S$ \underline{connexe}, \underline{et soit} $q \in \mathbb{Z}$ . \underline{Si la classe} $z_{\bar{\eta}}$ \underline{induit l'homomorphisme nul}

\begin{equation}\tag{1.5.3}
0 = z_{\bar{\eta}} \ : \ H^{q-1}(X_{\bar{\eta}} \ , \ \mathbb{Q}_\ell(i)) \ \longrightarrow \ H^{q+a-2d_1-1}(Y_{\bar{\eta}}, \ \mathbb{Q}_\ell(i+b-d_1)) \ ,
\end{equation}
\underline{alors, pour tout} $x \in \mathrm{Prim}^{q}(X/S,\ \mathbb{Q}_\ell(i))$ , \underline{on a}

}{F7II-XX-011}\begin{equation}\tag{1.5.4}
u_g(z(x)) = 0 \quad \mathrm{dans} \quad H^{1}(S, R^{q+a-2d_1-1}\,g_*\mathbb{Q}_\ell(i+b-d_1)) \quad .
\end{equation}

\SGApage{II}{347}

\underline{Démonstration} Par la commutativité de (1.4.2), il suffit de voir que $z$ induit l'homomorphisme nul

\SGAvar{\begin{equation}\tag{1.5.5}
z : R^{q-1}f_*\,\mathbb{Q}_\ell(i) \ \stackrel{0}{\longrightarrow} \ R^{q+a-2d_1-1}\,g_*\mathbb{Q}_\ell(i+b-d_1)\,.
\end{equation}
Or $f$ et $g$ étant propres et lisses, les faisceaux $R^{\bullet}f_*\mathbb{Q}_\ell(\cdot)$ et $R^{\bullet}g_*\mathbb{Q}_\ell(\cdot)$ sont constantstordus (SGA 4 XVI 2.2), donc équivalents à leurs valeurs en $\bar{\eta}$ , vus comme modules sous $\pi_1(S, \bar{\eta})$ .

}{\begin{equation}\tag{1.5.5}
z : R^{q-1}f_*\,\mathbb{Q}_\ell(i) \ \stackrel{0}{\longrightarrow} \ R^{q+a-2d_1-1}\,g_*\mathbb{Q}_\ell(i+b-d_1)\,.
\end{equation}
Or $f$ et $g$ étant propres et lisses, les faisceaux $R^{\bullet}f_*\mathbb{Q}_\ell(\cdot)$ et $R^{\bullet}g_*\mathbb{Q}_\ell(\cdot)$ sont constants tordus (SGA 4 XVI 2.2), donc équivalents à leurs valeurs en $\bar{\eta}$ , vus comme modules sous $\pi_1(S, \bar{\eta})$ .

}{F7II-XX-012}1.6. \underline{Une Construction de Manin}

Supposons données deux sections de $f : X \longrightarrow S$ :

\begin{equation}\tag{1.6.0}
\begin{tikzcd}[column sep=huge]
S \arrow[r, bend left=45, "t_1"] \arrow[r, bend right=45, "t_2"'] \& X \arrow[l, "f" description]
\end{tikzcd} \quad ,
\end{equation}









à chacune on associe sa classe de cohomologie

\SGAvar{\begin{equation}\tag{1.6.1}
\mathrm{cl}(t_i) \ = \ t_{i_*}(1) \ \in \ H^{2d_1} (X, \ \mathbb{Q}_\ell(d_1)) \quad .
\end{equation}
Si \underline{les fibres de} $f : X \longrightarrow S$ \underline{sont géométriquement connexes}, (plus généralement si pour $s \in S$ , $f(s)$ et $g(s)$ appartiennent à la même composante connexe de $X_s$ ) on voit bien que

}{\begin{equation}\tag{1.6.1}
\mathrm{cl}(t_i) \ = \ t_{i_*}(1) \ \in \ H^{2d_1} (X, \ \mathbb{Q}_\ell(d_1)) \quad .
\end{equation}
Si \underline{les fibres de} $f : X \longrightarrow S$ \underline{sont géométriquement connexes}, (plus généralement si pour $s \in S$ , $t_1(s)$ et $t_2(s)$ appartiennent à la même composante connexe de $X_s$ ) on voit bien que

}{F7II-XX-013}\begin{equation}\tag{1.6.2}
\mathrm{cl}(t_1) \ - \ \mathrm{cl}(t_2) \ \in \ \mathrm{Prim}^{2d_1} (X/S, \ \mathbb{Q}_\ell(d_1)) \quad .
\end{equation}

\SGApage{II}{348}

On désigne par

\begin{equation}\tag{1.6.3}
u_f(t_1-t_2) \in H^1(S, R^{2d_1-1}f_*\mathbb{Q}_\ell(d_1))
\end{equation}
son image par

\begin{equation}\tag{1.6.4}
u_f : \mathrm{Prim}^{2d_1}(X/S, \mathbb{Q}_\ell(d_1)) \longrightarrow H^1(S, R^{2d_1-1}f_*\mathbb{Q}_\ell(d_1)) \quad .
\end{equation}
(C'est un analogue $\mu(t_1 - t_2)$ de cet élément que Manin étudie dans [5]).

1.6.5. Soit maintenant $z$ un cycle algébrique sur $X \times_S Y$ de codimension $\geq b$ sur chaque fibre de $X \times_S Y$ $^{(*)}$, qui détermine une classe de cohomologie (SGA 5 IV\hspace{2em}):

\begin{equation}\tag{1.6.6}
\mathrm{cl}(z) \in H^{2b}(X \times_S Y, \mathbb{Q}_\ell(b)) \quad .
\end{equation}
Cette classe définit une application (cf (1.2.3))

\begin{equation}\tag{1.6.7}
\mathrm{cl}(z) : H^{2d_1}(X, \mathbb{Q}_\ell(d_1)) \longrightarrow H^{2b}(Y, \mathbb{Q}_\ell(b)) \quad ,
\end{equation}
et si les cycles $z(t_1)$ et $z(t_2)$ sont définis, on a (grâce à la compatibilité SGA 5 IV\hspace{3em} du cup-produit avec les intersections de cycles):

(*) Si $S$ est lisse sur un schéma $T$ (p. ex. $T=\mathrm{Spec}\ k$, $k$ un corps), il suffirait de supposer $z$ de codimension $\geq b$ sur chaque fibre au-dessus de $T$ .


\SGApage{II}{349}

\begin{equation}\tag{1.6.8}
\mathrm{cl}(z) (\mathrm{cl}(t_1) - \mathrm{cl}(t_2)) = \mathrm{cl}(z(t_1) - z(t_2)) \quad .
\end{equation}
La commutativité de (1.4.2) nous donne

\begin{equation}\tag{1.6.9}
\mathrm{cl}(z) (u_f(t_1-t_2)) = u_g(\mathrm{cl}(z(t_1) - z(t_2)))
\end{equation}
Comme corollaire de 1.5.3 , on a alors:

\underline{Proposition} 1.7. \underline{Soient} $t_1$ \underline{et} $t_2$ \underline{deux sections de} $f : X \longrightarrow S$ , $f$ \underline{étant supposé à fibres connexes}, \underline{et soit} $z$ \underline{un cycle algébrique} \underline{sur} $X \times_S Y$ \underline{de codimension} $b$ , \underline{tel que le cycle algébrique}

\SGAvar{\begin{equation}\tag{1.7.1}
z(t_1) - z(t_2)
\end{equation}
\underline{soit défini}. \underline{Alors sa classe de cohomologie est primitive} (XVIII 5.8.1). \underline{Si le cycle algébrique} $Z_{\overline{\eta}}$ \underline{sur} $X_{\overline{\eta}} \times Y_{\overline{\eta}}$ \underline{induit l'homomorphisme nul}

}{\begin{equation}\tag{1.7.1}
z(t_1) - z(t_2)
\end{equation}
\underline{soit défini}. \underline{Alors sa classe de cohomologie est primitive} (XVIII 5.8.1). \underline{Si le cycle algébrique} $z_{\overline{\eta}}$ \underline{sur} $X_{\overline{\eta}} \times Y_{\overline{\eta}}$ \underline{induit l'homomorphisme nul}

}{F7II-XX-014}\begin{equation}\tag{1.7.2}
z_{\overline{\eta}} : H^{2d_1-1} (X_{\overline{\eta}}, \mathbb{Q}_\ell(d_1)) \stackrel{0}{\longrightarrow} H^{2b-1}(Y_{\overline{\eta}}, \mathbb{Q}_\ell(b))
\end{equation}
\underline{alors}

\begin{equation}\tag{1.7.3}
u_g(z(t_1) - z(t_2)) = 0 \quad \underline{\mbox{dans}} \quad H^1(S, R^{2b-1}g_* \mathbb{Q}_\ell(b)) \quad ,
\end{equation}
(\underline{où} $u_g$ \underline{est l'homomorphisme de} XVIII 5.8.1 \underline{intervenant dans} 1.4).


\SGApage{II}{350}


\SGAvar{\SGAsectiontoc{2.  Un rappel sur le niveau (Cf. [2], 9.7, 10.1] )}{2}
2.  \underline{Un rappel sur le niveau} (Cf. [2], 9.7, 10.1] )

}{\SGAsectiontoc{2.  Un rappel sur le niveau (Cf. [2], 9.7, 10.1)}{2}
2.  \underline{Un rappel sur le niveau} (Cf. [2], 9.7, 10.1)

}{F7II-XX-015}2.0.1. Soit  $X$  propre et lisse sur un corps  $k$  algébriquement clos. On a vu (1.2.3) que, pour tout

\begin{equation}\tag{2.0.2}
T \text{ propre, lisse et connexe sur } k \text{ de dimension } d ,
\end{equation}
et tout

\begin{equation}\tag{2.0.3}
\text{cycle algébrique } z \text{ de codimension } b \text{ sur } T \times X,
\end{equation}
la classe de cohomologie

\begin{equation}\tag{2.0.4}
\mathrm{cl}(z) \in H^{2b}(T \times X , \mathbb{Q}_\ell(b))
\end{equation}
donne lieu aux applications

\begin{equation}\tag{2.0.5}
\mathrm{cl}(z) : H^q(T, Q_\ell(i)) \longrightarrow H^{q+2b-2d} (X, \mathbb{Q}_\ell(i+b-d)) .
\end{equation}
On définit \underline{la filtration par le co-niveau}  $(N^j)$  de  $H^q(X,\mathbb{Q}_\ell(i))$

\begin{equation}\tag{2.0.6}
N^j H^q(X,\mathbb{Q}_\ell(i)) = \sum_{(T,z)} \text{images} \quad \mathrm{cl}(z) : H^{q-2j}(T,\mathbb{Q}_\ell(i-j)) \longrightarrow H^q(X,\mathbb{Q}_\ell(i)),
\end{equation}
la somme étant étendue aux couples

\SGAvar{\begin{equation}\tag{2.0.7}
(T, z) \qquad \begin{array}{l} T/k \text{ propre lisse et connexe de dimension } d \\ z \text{ cycle algébrique sur } T \times X \text{ de codimension } j+d. \end{array}
\end{equation}
Il faut remarquer qu'on parle parfois de la filtration  $F^j$  \underline{croissante} ``par le niveau''  , défini par

\begin{equation}\tag{2.0.8}
F^j H^q(X, \mathbb{Q}_\ell(i)) = N^{\left[\frac{g-j}{2}\right]} H^q(X, \mathbb{Q}_\ell(i)) .
\end{equation}

}{\begin{equation}\tag{2.0.7}
(T, z) \qquad \begin{array}{l} T/k \text{ propre lisse et connexe de dimension } d \\ z \text{ cycle algébrique sur } T \times X \text{ de codimension } j+d. \end{array}
\end{equation}
Il faut remarquer qu'on parle parfois de la filtration  $F^j$  \underline{croissante} ``par le niveau''  , définie par

\begin{equation}\tag{2.0.8}
F^j H^q(X, \mathbb{Q}_\ell(i)) = N^{\left[\frac{q-j}{2}\right]} H^q(X, \mathbb{Q}_\ell(i)) .
\end{equation}

}{F7II-XX-016}\SGApage{II}{351}

Remplaçant  $(T, z)$  par  $(\mathbb{P}^1 \times T, \{pt\} \times z)$  , on voit que

\begin{equation}\tag{2.0.9}
N^j H^q(X, \mathbb{Q}_\ell(i)) \subset N^{j-1}H^q(X, \mathbb{Q}_\ell(i)) .
\end{equation}
Evidemment on a:

\SGAvar{\underline{Proposition} 2.1,   \underline{Soient}  $X$  \underline{propre,} \underline{lisse et géométriquement connexe sur un corps}  $k$ ,  $\bar{k}$  \underline{une clôture algébrique de}  $k$ ,  $X_{\bar{k}} = X \otimes_k \bar{k}$  .  \underline{Alors}  $\mathrm{gal}(\bar{k}/k)$  , \underline{agissant sur}  $H^q(X_{\bar{k}}, \mathbb{Q}_\ell(i))$, \underline{préserve la filtration par le co-niveau.}

}{\underline{Proposition} 2.1.   \underline{Soient}  $X$  \underline{propre,} \underline{lisse et géométriquement connexe sur un corps}  $k$ ,  $\bar{k}$  \underline{une clôture algébrique de}  $k$ ,  $X_{\bar{k}} = X \otimes_k \bar{k}$  .  \underline{Alors}  $\mathrm{gal}(\bar{k}/k)$  , \underline{agissant sur}  $H^q(X_{\bar{k}}, \mathbb{Q}_\ell(i))$, \underline{préserve la filtration par le co-niveau.}

}{F7II-XX-017}Des raisonnements standards (cf [2]) montrent les énoncés suivants,

\underline{Proposition} 2.2.  \underline{La filtration par le niveau est invariante par extension des corps algébriquement clos}  $k'/k$  .

\underline{Proposition} 2.3.   \underline{Soient}  $S$  \underline{un schéma réduit et irréductible, à point générique}  $\eta$  , \underline{et}

\begin{equation*}
\pi : X \longrightarrow S
\end{equation*}
\underline{un morphisme propre et lisse,} \underline{à fibres géométriquement connexes.} \underline{Si,} \underline{pour des entiers}  $j$  \underline{et}  $q$  \underline{donnés,} \underline{on a pour la fibre générique géométrique}  $X_{\bar{\eta}}$

\begin{equation}\tag{2.3.1}
H^q(X_{\bar{\eta}}, \mathbb{Q}_\ell) = N^j H^q(X_{\bar{\eta}}, \mathbb{Q}_\ell) ,
\end{equation}
\underline{alors il existe un ouvert non-vide}  $U \subset S$  \underline{tel que en tout point}  $u \in U$  , \underline{on ait pour la fibre géométrique}  $X_{\bar{u}}$

\begin{equation}\tag{2.3.2}
H^q(X_{\bar{u}}, \mathbb{Q}_\ell) = N^j H^q(X_{\bar{u}}, \mathbb{Q}_\ell)
\end{equation}

\SGApage{II}{352}


\SGAsectiontoc{3. Application aux pinceaux de Lefschetz}{3}
3. \underline{Application aux pinceaux de Lefschetz}

3.1. \quad A partir de maintenant, on fixe :

(3.1.0) $k$ , un corps algébriquement clos,

(3.1.1) $X \subset \mathbb{P}^{r}$ , un $k$-schéma projectif, lisse et connexe de dimension paire $n$ , qui vérifie (LV) (XVIII 5.2.2) .

Prenons

\SGAvar{(3.1.2) $D = \{H_{t}\}_{t \in \mathbb{P}^{1}}$ un pinceau de Lefschetz d'hypersurfacesde degré $d$ , à fibre générique

}{(3.1.2) $D = \{H_{t}\}_{t \in \mathbb{P}^{1}}$ un pinceau de Lefschetz d'hypersurfaces de degré $d$ , à fibre générique

}{F7II-XX-018}(3.1.3) $X_{\eta}$ , $\eta =$ le point générique de $\mathbb{P}^{1}$ .

Posons (cf. XVIII 5.4.4)

(3.1.4) $m : X_{\bar{\eta}} \hookrightarrow X \times_{k} k(\bar{\eta})$ l'inclusion ,

(3.1.5) $E^{n-1}(X_{\bar{\eta}}, \mathbb{Q}_{\ell}(i)) =$ l'orthogonal de $m^{*}H^{n-1}(X,\mathbb{Q}_{\ell}(i))$ dans $H^{n-1}(X_{\bar{\eta}},\mathbb{Q}_{\ell}(i))$ .

Nous avons vu (XVIII 6.7.2) que la \underline{formule de Picard-Lefschetz} entraîne

\SGAvar{(3.1.6) Si $n$ est pair, tout sous-groupe d'indice fini dans $\mathrm{gal}(k(\bar{\eta})/k(\eta))$ agit de façon irréductible sur $E^{2n-1}(X_{\bar{\eta}},\mathbb{Q}_{\ell}(i))$ .

}{(3.1.6) Si $n$ est pair, tout sous-groupe d'indice fini dans $\mathrm{gal}(k(\bar{\eta})/k(\eta))$ agit de façon irréductible sur $E^{n-1}(X_{\bar{\eta}},\mathbb{Q}_{\ell}(i))$ .

}{F7II-XX-019}\underline{Proposition} 3.2. \quad \underline{Soit} $D = \{H_{t}\}$ \underline{un pinceau de Lefschetz}, \underline{et supposons} $n = 2m$ \underline{pair.} \underline{Alors ou bien} 


\SGApage{II}{353}

\begin{equation}\tag{3.2.1}
 H^{n-1}(X_{\bar{\eta}},\mathbb{Q}_{\ell}(i)) \subset m^{*}H^{n-1}(X,\mathbb{Q}_{\ell}(i)) + N^{m-1}H^{n-1}(X_{\bar{\eta}},\mathbb{Q}_{\ell}(i))
\end{equation}
\underline{ou bien:}

(3.2.2) \underline{Pour toute courbe} $C$ \underline{propre et lisse sur} $k(\bar{\eta})$ , \underline{et tout} \underline{cycle algébrique} $z$ \underline{sur} $C \times X_{\bar{\eta}}$ \underline{de codimension} $m$ , \underline{l'homomorphisme induit} (1.2.3)

\begin{equation}\tag{3.2.3}
\mathrm{cl}(z) : H^{1}(C, \mathbb{Q}_{\ell}(1+i-m)) \longrightarrow H^{2m-1}(X_{\bar{\eta}},\mathbb{Q}_{\ell}(i))
\end{equation}
\underline{a son image dans} $m^{*}H^{2m-1}(X,\mathbb{Q}_{\ell}(i))$ .

\SGAvar{\underline{Démonstration.} Supposons que l'image d'un tel homomorphisme (3.2.3) n'est pas contenu dans $m^{*}H^{n-1}(X,\mathbb{Q}_{\ell}(i))$ , ou, ce qui revient au même, que son image dans $E^{n-1}(X_{\bar{\eta}},\mathbb{Q}_{\ell}(i)) = H^{n-1}(X_{\bar{\eta}},\mathbb{Q}_{\ell}(i)) \quad H^{n-1}(X,Q_{\ell}(i))$ n'est pas nulle. Comme cette image est évidemment stable sous les opérations d'un sous-groupe d'indice fini dans $\mathrm{gal}(k(\bar{\eta})/k(\eta))$ , l'irréductibilité (3.1.6) de $E^{n-1}(X_{\bar{\eta}},\mathbb{Q}_{\ell}(i)$ se récrit en disant

}{\underline{Démonstration.} Supposons que l'image d'un tel homomorphisme (3.2.3) n'est pas contenue dans $m^{*}H^{n-1}(X,\mathbb{Q}_{\ell}(i))$ , ou, ce qui revient au même, que son image dans $E^{n-1}(X_{\bar{\eta}},\mathbb{Q}_{\ell}(i)) = H^{n-1}(X_{\bar{\eta}},\mathbb{Q}_{\ell}(i)) / H^{n-1}(X,\mathbb{Q}_{\ell}(i))$ n'est pas nulle. Comme cette image est évidemment stable sous l'action d'un sous-groupe d'indice fini dans $\mathrm{gal}(k(\bar{\eta})/k(\eta))$ , l'irréductibilité (3.1.6) de $E^{n-1}(X_{\bar{\eta}},\mathbb{Q}_{\ell}(i))$ se récrit en disant

}{F7II-XX-020}\begin{equation*}
H^{n-1}(X_{\bar{\eta}},Q_{\ell}(i)) = m^{*}H^{n-1}(X,Q_{\ell}(i)) + \text{l'image de } (3.2.3)
\end{equation*}
et, par définition,

\begin{equation*}
\text{l'image de } (3.2.3) = \text{l'image par } \mathrm{cl}(z) \text{ de } H^{1}(C,\mathbb{Q}_{\ell}(1+i-m)) \subset N^{m-1}H^{n-1}(X_{\bar{\eta}},\mathbb{Q}_{\ell}(i)) ,
\end{equation*}
de sorte qu'on a bien trouvé

\begin{equation*}
H^{n-1}(X_{\bar{\eta}},\mathbb{Q}_{\ell}(i)) = m^{*}H^{n-1}(X,\mathbb{Q}_{\ell}(i)) + N^{m-1}H^{n-1}(X_{\bar{\eta}},\mathbb{Q}_{\ell}(i))
\end{equation*}

\SGApage{II}{354}


3.3. \underline{La condition (B)}.

3.3.1. On dira que l'inclusion

\begin{equation*}
m : X_{\overline{\eta}} \hookrightarrow X \times_{k} k(\overline{\eta})
\end{equation*}
vérifie la condition (B)  s'il existe

(3.3.2) $V_{\overline{\eta}}$ , un cycle algébrique de codimension $n-1$ sur $X_{\overline{\eta}} \times (X \times_{k} k(\overline{\eta}))$ tel que l'application qu'il définit

\SGAvar{\begin{equation}\tag{3.3.3}
V_{\overline{\eta}} ; H^{n-1}(X_{\overline{\eta}},\mathbb{Q}_{\ell}(i)) \longrightarrow H^{n-1}(X,\mathbb{Q}_{\ell}(i))
\end{equation}
soit inverse à gauche de

\begin{equation}\tag{3.3.4}
m^* : H^{n-1}(X,\mathbb{Q}_{\ell}(i)) \longrightarrow H^{n-1}(X_{\eta},\mathbb{Q}_{\ell}(i)) \quad ,
\end{equation}
i.e. tel que l'on ait

\begin{equation}\tag{3.3.5}
V_{\overline{\eta}} \; m^* = \mathrm{id} .
\end{equation}
(3.4). Cette condition (B) est vérifiée si l'opération $\Lambda$ [4, 1.4.2.1] est algébrique (cf. [4, 2.12 (ii)]. On conjecture que $\Lambda$ est toujours algébrique, donc que la condition (B) est toujours vérifiée, et on sait le vérifier dans de nombreux cas [4]. \underline{Elle est vérifiée trivialement pour} $X$ \underline{une intersection complète de dimension} paire car alors on sait (IX 1.6) que

}{\begin{equation}\tag{3.3.3}
V_{\overline{\eta}} : H^{n-1}(X_{\overline{\eta}},\mathbb{Q}_{\ell}(i)) \longrightarrow H^{n-1}(X,\mathbb{Q}_{\ell}(i))
\end{equation}
soit inverse à gauche de

\begin{equation}\tag{3.3.4}
m^* : H^{n-1}(X,\mathbb{Q}_{\ell}(i)) \longrightarrow H^{n-1}(X_{\overline{\eta}},\mathbb{Q}_{\ell}(i)) \quad ,
\end{equation}
i.e. tel que l'on ait

\begin{equation}\tag{3.3.5}
V_{\overline{\eta}} \; m^* = \mathrm{id} .
\end{equation}
(3.4). Cette condition (B) est vérifiée si l'opération $\Lambda$ [4, 1.4.2.1] est algébrique (cf. [4, 2.12 (ii)]). On conjecture que $\Lambda$ est toujours algébrique, donc que la condition (B) est toujours vérifiée, et on sait le vérifier dans de nombreux cas [4]. \underline{Elle est vérifiée trivialement pour} $X$ \underline{une intersection complète de dimension} paire car alors on sait (IX 1.6) que

}{F7II-XX-021}\begin{equation}\tag{3.4.1}
H^{n-1}(X,\mathbb{Q}_{\ell}(i)) \;=\; 0 .
\end{equation}

\SGApage{II}{355}


\SGAsectiontoc{4. Le théorème de GRIFFITHS}{4}
4. \underline{Le théorème de GRIFFITHS}

\SGAvar{\underline{Théorème} 4.1. \underline{Soit} $X$ \underline{projectif,} \underline{lisse et connexe sur} $k$ , \underline{de dimension paire} $n = 2m \geq 4$ \underline{Soit} $D$ \underline{un pinceau de Lefschetz de sections par des hypersurfaces de degré} $d$ . \underline{On suppose}

}{\underline{Théorème} 4.1. \underline{Soit} $X$ \underline{projectif,} \underline{lisse et connexe sur} $k$ , \underline{de dimension paire} $n = 2m \geq 4$. \underline{Soit} $D$ \underline{un pinceau de Lefschetz de sections par des hypersurfaces de degré} $d$ . \underline{On suppose}

}{F7II-XX-022}(4.1.0) $X$ \underline{vérifie} (L V) (XVIII 5.2.2),

(4.1.1) \underline{le pinceau} $D$ \underline{vérifie l'hypothèse} (A) (XVIII 5.3)

\SGAvar{(4.1.2) \underline{le pinceau} $D$ \underline{vérifie} (3.2.2), \underline{i e.} (3.2) \underline{la cohomologie évanescente} $E^{n-1}(X_{\overline{\eta}},\mathbb{Q}_{\ell})$ \underline{n'est pas  de niveau} $\leq 1$ ,

(4.1.3) \underline{l'inclusion} $m^* : X_{\overline{\eta}} \hookrightarrow X \times_{k} k(\overline{\eta})$ \underline{vérifie l'hypothèse} (B) (3.3).

\underline{Alors,} \underline{si un cycle algébrique} $z$ \underline{de codimension} $m$ \underline{sur} $X$ \underline{est tel que sa restriction(en tant que cycle algébrique) à} $X_{\overline{\eta}}$ \underline{satisfasse}

}{(4.1.2) \underline{le pinceau} $D$ \underline{vérifie} (3.2.2), \underline{i.e.} (3.2) \underline{la cohomologie évanescente} $E^{n-1}(X_{\overline{\eta}},\mathbb{Q}_{\ell})$ \underline{n'est pas  de niveau} $\leq 1$ ,

(4.1.3) \underline{l'inclusion} $m : X_{\overline{\eta}} \hookrightarrow X \times_{k} k(\overline{\eta})$ \underline{vérifie l'hypothèse} (B) (3.3).

\underline{Alors,} \underline{si un cycle algébrique} $z$ \underline{de codimension} $m$ \underline{sur} $X$ \underline{est tel que sa restriction (en tant que cycle algébrique) à} $X_{\overline{\eta}}$ \underline{satisfasse}

}{F7II-XX-023}(4.1.4) $z_{\overline{\eta}} = z|X_{\overline{\eta}}$ \underline{est algébriquement équivalent à zéro}

\underline{Alors} (4.1.5) $z$ \underline{est} \underline{$\mathbb{Q}_{\ell}$-cohomologue à zéro,} \underline{i.e.} $\mathrm{cl}(z) \in H^{2m}(X,\mathbb{Q}_{\ell}(m))$ \underline{est nul}.


4.2. \underline{Une remarque sur le cas où} $X$ \underline{est une intersection complète}.

\SGAvar{Comme on a déjà remarqué (cf. XVIII 5.2.5, et (3.4)),  les hypothèses (4.1.0) et (4.1.3) sont alors automatiquement vérifiées . Quant à l'hypothèse (4.1.1), on a vu (XVIII 6.3.4, 6.4) qu'elle est vérifiée pour un pinceau d'hypersurfaces de degré $d > 2n$  D'autre part, on va prouver dans le prochain exposé la condition (4.1 2) pour un pinceau ``général''d'hypersurfaces de degré $d > 2n$  (XXI 5.2 ). Donc un pinceau ``général''d'hypersurfaces de degré $d > 2n$  vérifie toutes les hypothèses du théorème 4.1.


}{Comme on l'a déjà remarqué (cf. XVIII 5.2.5 et (3.4)), les hypothèses (4.1.0) et (4.1.3) sont alors automatiquement vérifiées. Quant à l'hypothèse (4.1.1), on a vu (XVIII 6.3.4, 6.4) qu'elle est vérifiée pour un pinceau d'hypersurfaces de degré $d > 2n$. D'autre part, on va prouver dans le prochain exposé la condition (4.1.2) pour un pinceau ``général'' d'hypersurfaces de degré $d > 2n$ (XXI 5.2). Donc un pinceau ``général'' d'hypersurfaces de degré $d > 2n$ vérifie toutes les hypothèses du théorème 4.1.


}{F7II-XX-024}\SGApage{II}{356}

\SGAvar{4.3.   Revenant au cas où $X$ est quelconque, il semble très plausible que pour un degré $d$ assez grand, toutes les conditions de 4.1 sont vérifiées pour $D$ , du moins pour $D$ assez général  On a déjà vu que c'est le cas pour (4.1.1) (XVIII 6.4), et on a signalé qu'on conjecture que (4.1.0) et (4.1.3) sont vérifiés en tous cas (et c'est vrai en tous cas pour (4.1.0) si  car $k = 0$ ). Quant à la condition (4.1.2), elle est vérifiée (grâce à la théorie de Hodge) pour $d$ grand si  car $k = 0$, et des arguments heuristiques (que le lecteur trouvera en appendice à l'exposé suivant) rendent plausibles qu'il en est encore de même si car $k > 0$.    Plus généralement, ces arguments (théorie de Hodge, resp cohomologie cristalline) prouvent (resp  rendent plausible) que pour une variété projective lisse $X$ de dimension $n$, et pour tout pinceau de Lefschetz $D$ d'hypersurfaces de degré $d$ assez grand (resp. tout pinceau $D$ ``assez général''  d'hypersurfaces de degré $d$ ) la cohomologie évanescente $E^{n-1}(X_{\bar{\eta}}, \mathbb{Q}_{\ell})$ est de niveau égal à $n-1$  (i.e. n'est pas de niveau  $\leqslant n-3$, i.e. n'est pas de coniveau 1 ).


}{4.3. Revenant au cas où $X$ est quelconque, il semble très plausible que, pour un degré $d$ assez grand, toutes les conditions de 4.1 soient vérifiées pour $D$, du moins pour $D$ assez général. On a déjà vu que c'est le cas pour (4.1.1) (XVIII 6.4), et on a signalé qu'on conjecture que (4.1.0) et (4.1.3) sont vérifiées dans tous les cas (et c'est vrai dans tous les cas pour (4.1.0) si $\operatorname{car}(k)=0$). Quant à la condition (4.1.2), elle est vérifiée (grâce à la théorie de Hodge) pour $d$ grand si $\operatorname{car}(k)=0$, et des arguments heuristiques (que le lecteur trouvera en appendice à l'exposé suivant) rendent plausible qu'il en est encore de même si $\operatorname{car}(k)>0$. Plus généralement, ces arguments (théorie de Hodge, resp. cohomologie cristalline) prouvent (resp. rendent plausible) que, pour une variété projective lisse $X$ de dimension $n$, et pour tout pinceau de Lefschetz $D$ d'hypersurfaces de degré $d$ assez grand (resp. tout pinceau $D$ ``assez général'' d'hypersurfaces de degré $d$), la cohomologie évanescente $E^{n-1}(X_{\bar{\eta}}, \mathbb{Q}_{\ell})$ est de niveau égal à $n-1$ (i.e. n'est pas de niveau $\leqslant n-3$, i.e. n'est pas de coniveau 1).


}{F7II-XX-025}4.4.  \underline{Démonstration du théorème de Griffiths.}

\SGAvar{Soit $z$ un cycle algébrique de codimension $n$ sur $X$ , tel que

}{Soit $z$ un cycle algébrique de codimension $m$ sur $X$, tel que

}{F7II-XX-026}\begin{equation}\tag{4.4.1}
z_{\bar{\eta}} = z|X_{\bar{\eta}} \quad \text{est algébriquement équivalent à zéro} \quad .
\end{equation}
Par définition, il existe

\begin{equation}\tag{4.4.2}
\text{une courbe} \quad C_{\bar{\eta}} \quad \text{, propre, lisse et connexe sur} \quad k(\bar{\eta}) , 
\end{equation}

\SGApage{II}{357}

\SGAvar{\begin{equation}\tag{4.4.3}
 \text{deux sections} \; t_{i}(\bar{\eta}) : \mathrm{Sp}(k(\bar{\eta})) \longrightarrow C \quad , \; i=1,2 \quad ,
\end{equation}
\begin{equation}\tag{4.4.4}
\text{un cycle algébrique} \quad P_{\bar{\eta}} \quad \text{sur} \quad C_{\bar{\eta}} \times X_{\bar{\eta}} \quad \text{, de codimension} \quad n \; ,
\end{equation}
tels que

}{\begin{equation}\tag{4.4.3}
 \text{deux sections} \; t_{i}(\bar{\eta}) : \mathrm{Sp}(k(\bar{\eta})) \longrightarrow C_{\bar\eta} \quad , \; i=1,2 \quad ,
\end{equation}
\begin{equation}\tag{4.4.4}
\text{un cycle algébrique} \quad P_{\bar{\eta}} \quad \text{sur} \quad C_{\bar{\eta}} \times X_{\bar{\eta}} \quad \text{, de codimension} \quad m \; ,
\end{equation}
tels que

}{F7II-XX-027}\begin{equation}\tag{4.4.5}
z_{\bar{\eta}} = P_{\bar{\eta}}(t_{1}(\bar{\eta})) - P_{\bar{\eta}}(t_{2}(\bar{\eta})) \quad \text{sur} \quad X_{\bar{\eta}} \quad .
\end{equation}
Considérons le cycle $P_{\bar{\eta}}$ comme une \underline{correspondance algébrique} sur $(C_{\bar{\eta}}, X_{\bar{\eta}})$ , définissant une correspondance cohomologique, encore notée $P_{\bar{\eta}}$, ou encore un homomorphisme

\SGAvar{\begin{equation}\tag{4.4.6}
P_{\bar{\eta}} : H^{*}(C_{\bar{\eta}},\mathbb{Q}_{\ell}(i)) \longrightarrow H^{*+2(m-1)}(X_{\bar{\eta}}, \mathbb{Q}_{\ell}((i+m-1)) \quad .
\end{equation}
Considérons de même la correspondance algébrique $V_{\bar{\eta}}$ (3.3.2) sur $(X_{\bar{\eta}}, X \otimes_{k} k(\bar{\eta}))$ donnée par l'hypothèse (B) de (4.1.3), définissant

\begin{equation}\tag{4.4.7}
V_{\bar{\eta}} : H^{*}(X_{\bar{\eta}}, \mathbb{Q}_{\ell}(i+m-1)) \longrightarrow H^{*}(X\otimes_{k}k(\bar{\eta}), \mathbb{Q}_{\ell}(i+m-1)) \quad ,
\end{equation}
d'où une correspondance algébrique composée $Q_{\bar{\eta}} = V_{\bar{\eta}} \circ P$ sur $(C_{\bar{\eta}}, X \otimes_{k} k(\bar{\eta}))$ , définissant

}{\begin{equation}\tag{4.4.6}
P_{\bar{\eta}} : H^{*}(C_{\bar{\eta}},\mathbb{Q}_{\ell}(i)) \longrightarrow H^{*+2(m-1)}(X_{\bar{\eta}}, \mathbb{Q}_{\ell}(i+m-1)) \quad .
\end{equation}
Considérons de même la correspondance algébrique $V_{\bar{\eta}}$ (3.3.2) sur $(X_{\bar{\eta}}, X \otimes_{k} k(\bar{\eta}))$ donnée par l'hypothèse (B) de (4.1.3), définissant

\begin{equation}\tag{4.4.7}
V_{\bar{\eta}} : H^{*}(X_{\bar{\eta}}, \mathbb{Q}_{\ell}(i+m-1)) \longrightarrow H^{*}(X\otimes_{k}k(\bar{\eta}), \mathbb{Q}_{\ell}(i+m-1)) \quad ,
\end{equation}
d'où une correspondance algébrique composée $Q_{\bar{\eta}} = V_{\bar{\eta}} \circ P_{\bar{\eta}}$ sur $(C_{\bar{\eta}}, X \otimes_{k} k(\bar{\eta}))$ , définissant

}{F7II-XX-028}\begin{equation}\tag{4.4.8}
Q_{\bar{\eta}} = V_{\bar{\eta}} \circ P_{\bar{\eta}} : H^{*}(C_{\bar{\eta}},\mathbb{Q}_{\ell}(i)) \longrightarrow H^{*+2(m-1)}(X\otimes_{k}k(\bar{\eta}),\mathbb{Q}_{\ell}(i+m-1))
\end{equation}
Considérons l'inclusion

\SGAvar{\begin{equation*}
m : X_{\bar{\eta}} \longrightarrow X \otimes_{k} k(\bar{\eta}) \quad ,
\end{equation*}
donnant la correspondance algébrique $m^{*}$ sur $(X \otimes_{k} k(\bar{\eta}), X_{\bar{\eta}})$ , d'où la correspondance composée $m^{*} Q_{\bar{\eta}} \;\; {}_{\bar{\eta}}$ sur $(C_{\bar{\eta}}, X_{\bar{\eta}})$ , définissant 


}{\begin{equation*}
m : X_{\bar{\eta}} \longrightarrow X \otimes_{k} k(\bar{\eta}) \quad ,
\end{equation*}
donnant la correspondance algébrique $m^{*}$ sur $(X \otimes_{k} k(\bar{\eta}), X_{\bar{\eta}})$, d'où la correspondance composée $m^{*} Q_{\bar{\eta}}$ sur $(C_{\bar{\eta}}, X_{\bar{\eta}})$, définissant 


}{F7II-XX-029}\SGApage{II}{358}

\begin{equation}\tag{4.4.9}
 m^*Q_{\bar\eta} = m^*V_{\bar\eta} P_{\bar\eta} : H^*(C_{\bar\eta}, \mathbb{Q}_\ell(i)) \longrightarrow H^{*+2(m-1)}(X_{\bar\eta},\mathbb{Q}_\ell(i+m-1)) \quad .
\end{equation}
Considérons la correspondance algébrique

\SGAvar{\begin{equation}\tag{4.4.10}
R_{\bar\eta} = P_{\bar\eta} m^*Q_{\bar\eta} = (\mathrm{id}-m^*V_{\bar\eta})P_{\bar\eta} : H^*(C_{\bar\eta},\mathbb{Q}_\ell(i)) \longrightarrow H^{*+2(m-1)}(X_{\bar\eta},\mathbb{Q}_\ell(i+m-1)).
\end{equation}
Comme $V_{\bar\eta}$ est inverse à gauche de $m^*$ , $(\mathrm{id}-m^*V_{\bar\eta})$ est un projecteur dans $H^*(X_{\bar\eta}, -)$ , dont l'image est un supplémentaire de $\mathrm{Im}\ (m^* : H^*(X \otimes_k k(\bar\eta),-) \longrightarrow H^*(X_{\bar\eta},-))$ . Comme en vertu de l'hypothèse (4.1.2) l'homomorphisme

}{\begin{equation}\tag{4.4.10}
R_{\bar\eta} = P_{\bar\eta} - m^*Q_{\bar\eta} = (\mathrm{id}-m^*V_{\bar\eta})P_{\bar\eta} : H^*(C_{\bar\eta},\mathbb{Q}_\ell(i)) \longrightarrow H^{*+2(m-1)}(X_{\bar\eta},\mathbb{Q}_\ell(i+m-1)).
\end{equation}
Comme $V_{\bar\eta}$ est inverse à gauche de $m^*$ , $(\mathrm{id}-m^*V_{\bar\eta})$ est un projecteur dans $H^*(X_{\bar\eta}, -)$ , dont l'image est un supplémentaire de $\mathrm{Im}\ (m^* : H^*(X \otimes_k k(\bar\eta),-) \longrightarrow H^*(X_{\bar\eta},-))$ . Comme en vertu de l'hypothèse (4.1.2) l'homomorphisme

}{F7II-XX-030}\begin{equation}\tag{4.4.11}
R_{\bar\eta} : H^1(C_{\bar\eta},\mathbb{Q}_\ell(i)) \longrightarrow H^{2m-1}(X_{\bar\eta},\mathbb{Q}_\ell(i+m-1))
\end{equation}
a son image dans $\mathrm{Im}\ m^*$ , on trouve que

\SGAvar{(4.4.12) \quad l'homomorphisme (4.4.11) est nul

Choissisons maintenant une extension finie

}{(4.4.12) \quad l'homomorphisme (4.4.11) est nul.

Choisissons maintenant une extension finie

}{F7II-XX-031}\begin{equation*}
L/k(\eta)
\end{equation*}
au-dessus de laquelle les objets

\begin{equation*}
C_{\bar\eta}, \ t_1(\bar\eta), \ t_2(\bar\eta), \ Q_{\bar\eta} \ , \ R_{\bar\eta}
\end{equation*}
soient déjà définis, et

$U$  un ouvert $\neq \emptyset$  dans le normalisé de $\mathbb{P}^1$  dans  $L$ .

Prenant $U$ assez petit, il existe une courbe relative propre et lisse, à fibres géométriquement connexes 


\SGApage{II}{359}

\begin{equation*}
 C \longrightarrow U \ ,
\end{equation*}
dont la fibre générique géométrique soit $C_{\bar\eta}$ , et deux sections

\begin{equation*}
t_i : U \longrightarrow C \quad i = 1, 2
\end{equation*}
dont les valeurs en $\bar\eta$ sont $t_i(\bar\eta)$ .

Soient

\begin{equation*}
\rho : \widetilde{X} \longrightarrow \mathbb{P}^1
\end{equation*}
la projection (XVIII 3.1.4) dont les fibres sont les sections du pinceau,

\begin{equation*}
\widetilde{X}_U = \widetilde{X} \times_{\mathbb{P}^1} U \ ,
\end{equation*}
et

\begin{equation*}
M : \widetilde{X}_U \hookrightarrow X \times U
\end{equation*}
l'immersion canonique de $\widetilde{X}_U$ dans $X \times U$ (considéré comme la famille constante sur $U$ de valeur $X$ ) ,

\SGAvar{\begin{equation*}
z_U = M^* \ \mathrm{pr}_i^*(z) \ ,
\end{equation*}
le cycle algébrique sur $\widetilde{X}_U$ , induit par $z$ . Prenant $U$ assez petit, il existe un cycle algébrique $R$ sur $C \times_U \widetilde{X}_U$ dont la valeur en $C_{\bar\eta} \times X_{\bar\eta}$ est $R_{\bar\eta}$ , et un cycle algébrique $Q$ sur $C \times X \times U$ dont la valeur en $C_{\bar\eta} \times (X \otimes_k k(\bar\eta))$ est $Q_{\bar\eta}$ , tels que l'on ait

}{\begin{equation*}
z_U = M^* \ \mathrm{pr}_1^*(z) \ ,
\end{equation*}
le cycle algébrique sur $\widetilde{X}_U$ , induit par $z$ . Prenant $U$ assez petit, il existe un cycle algébrique $R$ sur $C \times_U \widetilde{X}_U$ dont la valeur en $C_{\bar\eta} \times X_{\bar\eta}$ est $R_{\bar\eta}$ , et un cycle algébrique $Q$ sur $C \times X \times U$ dont la valeur en $C_{\bar\eta} \times (X \otimes_k k(\bar\eta))$ est $Q_{\bar\eta}$ , tels que l'on ait

}{F7II-XX-032}\begin{equation}\tag{4.4.13}
z_U = R(t_1) - R(t_2) + (\mathrm{id}\times M)^*[Q(t_1) - Q(t_2)] \quad .
\end{equation}
Pour démontrer le théorème, il suffira, grâce à l'injectivité de
\SGApage{II}{360}
l'homomorphisme de Griffiths  ``modulo partie fixe'' (XVIII 5.8.7)

\SGAvar{\begin{equation*}
\mathrm{Prim}^{2m}(X,\mathbb{Q}_\ell(n)) \ \xrightarrow{\ \mathrm{grif}\ }\ \frac{H^1(U,R^{2m-1}\rho_{U*}\mathbb{Q}_\ell(m))}{M^*H^1(U,H^{2m-1}(X,\mathbb{Q}_\ell(m))} \qquad ,
\end{equation*}
de démontrer que l'image de  $z_U$  par

\begin{equation*}
u_0 \ : \ \mathrm{Prim}^{2m}(\widetilde{X}_{U/U},\mathbb{Q}_\ell(m)) \longrightarrow H^1(U,R^{2m-1}\rho_{U*}\mathbb{Q}_\ell(m))
\end{equation*}
se trouve dans  $M^*(H^1(U,H^{2m-1}(X,\mathbb{Q}_\ell(m))$  .

Or, compte tenu de la décomposition (4.4.13), et de la commutativité de (1.4.2)

\begin{center}
\begin{tikzcd}[column sep=huge, row sep=huge]
\mathrm{Prim}^{2}(C/U,\mathbb{Q}_\ell(1)) \arrow[r, "Q"] \& \mathrm{Prim}^{2m}(X\times U/U,\mathbb{Q}_\ell(n)) \arrow[r, "M^*"] \arrow[d, "u_{pr_2}"] \& \mathrm{Prim}^{2m}(\widetilde{X}_{U/U},\mathbb{Q}_\ell(m)) \arrow[d, "u_0"] \\
 {} \& H^1(U,H^{2m-1}(X,\mathbb{Q}_\ell(m)) \arrow[r, "M^*"] \& H^1(U,R^{2m-1}\rho_{U*}\mathbb{Q}_\ell(m))\ ,
\end{tikzcd}
\end{center}












}{\begin{equation*}
\mathrm{Prim}^{2m}(X,\mathbb{Q}_\ell(m)) \ \xrightarrow{\ \mathrm{grif}\ }\ \frac{H^1(U,R^{2m-1}\rho_{U*}\mathbb{Q}_\ell(m))}{M^*H^1(U,H^{2m-1}(X,\mathbb{Q}_\ell(m)))} \qquad ,
\end{equation*}
de démontrer que l'image de  $z_U$  par

\begin{equation*}
u_0 \ : \ \mathrm{Prim}^{2m}(\widetilde{X}_{U/U},\mathbb{Q}_\ell(m)) \longrightarrow H^1(U,R^{2m-1}\rho_{U*}\mathbb{Q}_\ell(m))
\end{equation*}
se trouve dans $M^*\bigl(H^1(U,H^{2m-1}(X,\mathbb{Q}_\ell(m)))\bigr)$.

Or, compte tenu de la décomposition (4.4.13), et de la commutativité de (1.4.2),

\begin{center}
\begin{tikzcd}[column sep=huge, row sep=huge]
\mathrm{Prim}^{2}(C/U,\mathbb{Q}_\ell(1)) \arrow[r, "Q"] \& \mathrm{Prim}^{2m}(X\times U/U,\mathbb{Q}_\ell(m)) \arrow[r, "M^*"] \arrow[d, "u_{pr_2}"] \& \mathrm{Prim}^{2m}(\widetilde{X}_{U/U},\mathbb{Q}_\ell(m)) \arrow[d, "u_0"] \\
 {} \& H^1(U,H^{2m-1}(X,\mathbb{Q}_\ell(m))) \arrow[r, "M^*"] \& H^1(U,R^{2m-1}\rho_{U*}\mathbb{Q}_\ell(m))\ ,
\end{tikzcd}
\end{center}












}{F7II-XX-033}il suffit de démontrer

\SGAvar{\begin{equation*}
u_0(R(t_1) - R(t_2)) = 0 \quad \text{dans} \quad H^1(U,R^{2n-1}\rho_{U*}\mathbb{Q}_\ell(n)) \quad ,
\end{equation*}
qui est conséquence de (4.3.12), grâce à 1.7.


\SGAsectiontoc{5.  Le Groupe de Griffiths}{5}
5.  \underline{Le Groupe de Griffiths}

}{\begin{equation*}
u_0(R(t_1) - R(t_2)) = 0 \quad \text{dans} \quad H^1(U,R^{2m-1}\rho_{U*}\mathbb{Q}_\ell(m)) \quad ,
\end{equation*}
qui est conséquence de (4.4.12), grâce à 1.7.


\SGAsectiontoc{5.  Le groupe de Griffiths}{5}
5.  \underline{Le groupe de Griffiths}

}{F7II-XX-034}5.0.  Soit  $S$  propre et lisse sur un corps  $k$  algébriquement clos. Désignons par
\SGApage{II}{361}
(5.0.1)   $Z^i(S)$  le groupe des cycles algébriques de codimension $i$  sur $S$

\SGAvar{(5.0.2)   $Z^i_{\mathrm{alg}}(S)$ le sous-groupe de  $Z^i(S)$  formé des cycles algébriquement équivalent à zéro,

et, pour tout nombre premier  $\ell \neq$ car $(k)$

}{(5.0.2)   $Z^i_{\mathrm{alg}}(S)$ le sous-groupe de  $Z^i(S)$  formé des cycles algébriquement équivalents à zéro,

et, pour tout nombre premier  $\ell \neq \operatorname{car}(k)$,

}{F7II-XX-035}(5.0.3)   $Z^i_{\ell-\mathrm{coh}}(S)$  le sous-groupe de  $Z^i(S)$  formé des cycles dont la classe de cohomologie dans  $H^{2i}(S,\mathbb{Q}_\ell(i))$  est nulle.

\SGAvar{(On conjecture [4] que ce groupe est indépendant de  $\ell$  , ce qui est vrai du moins si  car.$k=0$ . )

}{(On conjecture [4] que ce groupe est indépendant de $\ell$, ce qui est vrai du moins si $\operatorname{car}(k)=0$.)

}{F7II-XX-036}On pose

\SGAvar{\begin{equation}\tag{5.0.4}
\mathrm{Grif}^i_\ell(S) = Z^i_{\ell-\mathrm{coh}}(S)/\ Z^i_{\mathrm{alg}}(S) \quad .
\end{equation}
Si de plus  $S$  est projectif, et muni d'une immersion projective, on pose

}{\begin{equation}\tag{5.0.4}
\mathrm{Grif}^i_\ell(S) = Z^i_{\ell-\mathrm{coh}}(S)/\ Z^i_{\mathrm{alg}}(S) \quad .
\end{equation}
Si de plus $S$ est projectif et muni d'une immersion projective, on pose

}{F7II-XX-037}(5.0.5)   $Z^{i,\mathrm{prim}}(S) =$ le sous-groupe de  $Z^i(S)$  formé des cycles qui sont primitifs (i.e. dont la classe de cohomologie l'est),

(5.0.6)   $\mathrm{Prim}^{2i}_{\mathrm{alg}}(S,\mathbb{Q}_\ell(i)) =$ le sous-$\mathbb{Q}$-vectoriel de  $H^{2i}(S,\mathbb{Q}_\ell(i))$ engendré par des classes de cohomologie des cycles algébriques primitifs.

\underline{Corollaire} 5.1    \underline{Sous les hypothèses de}  4.1, \underline{on a}

\SGAvar{\begin{equation}\tag*{5.1.1,}
\dim_{\mathbb{Q}}(\mathrm{Grif}^n(X_{\overline{\eta}})\otimes_{\mathbb{Z}}\mathbb{Q}) \ \geq \ \dim_Q(\mathrm{Prim}^{2n}_{\mathrm{alg}}(X,\mathbb{Q}_\ell(n)))
\end{equation}
En effet, d'après 4.1 5, le noyau de la flèche composée

\begin{equation*}
\mathbb{Q}\otimes_{\mathbb{Z}}Z^{n,\mathrm{prim}}(X) \ \xrightarrow{\ \text{rest. à } X_{\overline{\eta}}\ }\ \mathbb{Q} \otimes_{\mathbb{Z}}Z^n_{\ell-\mathrm{coh}}(X_{\overline{\eta}}) \longrightarrow \mathrm{Grif}^n(X_{\overline{\eta}}) \quad 
\end{equation*}

}{\begin{equation}\tag{5.1.1}
\dim_{\mathbb{Q}}(\mathrm{Grif}^m_\ell(X_{\overline{\eta}})\otimes_{\mathbb{Z}}\mathbb{Q}) \ \geq \ \dim_{\mathbb Q}(\mathrm{Prim}^{2m}_{\mathrm{alg}}(X,\mathbb{Q}_\ell(m)))
\end{equation}
En effet, d'après (4.1.5), le noyau de la flèche composée

\begin{equation*}
\mathbb{Q}\otimes_{\mathbb{Z}}Z^{m,\mathrm{prim}}(X) \ \xrightarrow{\ \text{rest. à } X_{\overline{\eta}}\ }\ \mathbb{Q} \otimes_{\mathbb{Z}}Z^m_{\ell-\mathrm{coh}}(X_{\overline{\eta}}) \longrightarrow \mathrm{Grif}^m_\ell(X_{\overline{\eta}}) \quad 
\end{equation*}

}{F7II-XX-038}\SGApage{II}{362}

\SGAvar{ est \underline{contenu} dans $Z^n_{\ell-\mathrm{coh}}(X) \otimes_{\mathbb{Z}} \mathbb{Q}$ , de sorte que $\mathrm{Prim}^{2n}_{\mathrm{alg}}(X,\mathbb{Q}_{\ell}(n))$ est un quotient d'un sous-groupe de $\mathrm{Grif}^n(X_{\overline{\eta}}) \otimes_{\mathbb{Z}} \mathbb{Q}$ , ce qui achève la démonstration

}{ est \underline{contenu} dans $Z^m_{\ell-\mathrm{coh}}(X) \otimes_{\mathbb{Z}} \mathbb{Q}$, de sorte que $\mathrm{Prim}^{2m}_{\mathrm{alg}}(X,\mathbb{Q}_{\ell}(m))$ est un quotient d'un sous-groupe de $\mathrm{Grif}^m_\ell(X_{\overline{\eta}}) \otimes_{\mathbb{Z}} \mathbb{Q}$, ce qui achève la démonstration.

}{F7II-XX-039}\underline{Remarque} 5.2.  Pour $X$ une hypersurface quadrique dans $\mathbb{P}^{2m+1}$ on a

\SGAvar{\begin{equation*}
\mathrm{dim}_{\mathbb{Q}_{\ell}}(\mathrm{Prim}^{2m}_{\mathrm{alg}}(X,\mathbb{Q}_{\ell}(n)) = 1
\end{equation*}
\underline{En particulier} (4.2), \underline{pour} $X$ \underline{une quadrique dans} $\mathbb{P}^5$ , \underline{la section de} $X$ \underline{par une hypersurface générique de degré} $> 5$ \underline{contient une courbe,} $Q_{\ell}$\underline{-cohomologie à zéro}, \underline{dont aucun multiple n'est algébriquement équivalent à zéro.}


}{\begin{equation*}
\dim_{\mathbb{Q}_{\ell}}(\mathrm{Prim}^{2m}_{\mathrm{alg}}(X,\mathbb{Q}_{\ell}(m))) = 1
\end{equation*}
\underline{En particulier} ($m=2$), \underline{pour} $X$ \underline{une quadrique dans} $\mathbb{P}^5$, \underline{la section de} $X$ \underline{par une hypersurface générique de degré} $> 5$ \underline{contient une courbe,} $\mathbb{Q}_{\ell}$\underline{-cohomologue à zéro}, \underline{dont aucun multiple non nul n'est algébriquement équivalent à zéro.}


}{F7II-XX-040}\underline{REFERENCES}

[1]\quad Griffiths, P.  On the Periods of Certain Rational Integrals III, à paraître dans Annals of Mathematics.

[2]\quad Grothendieck, A.  Le Groupe de Brauer III (esp \S\S 9, 10) dans Dix Exposés sur la Cohomologie des Schémas, North Holland, 1968.

[3]\quad Jouanolou, J.P.  Cohomologie $\ell$-adique, Thèse, Paris 1969.

[4]\quad Kleiman, S.  Algebraic Cycles and the Weil Conjectures, in Dix Exposés sur la Cohomologie des Schémas, North Holland, 1968.

[5]\quad Manin, J.  Rational Points of Algebraic Curves over Function Fields, Translations Amer. Math. Soc.
\selectlanguage{french}

\ifnum\SGAform=1
\clearpage\section*{Apparat critique}\small
Les lectures ci-dessous sont données dans le TeX exact des témoins, afin de préserver les distinctions de notation. Les indices décimaux renvoient aux différences individuelles du passage marqué.

\par\medskip\noindent\hypertarget{app-F7II-XX-001.1}{}\textbf{F7II-XX-001.1}\quad p. 341; printed misprint.\par
\noindent Imprimé : {\ttfamily\footnotesize facon\allowbreak{}}\par
\noindent Corrigé : {\ttfamily\footnotesize façon\allowbreak{}}\par
\noindent Source : 7II:600; 7II:1032. \nolinkurl{fixes/MISPRINTS_7II.json#600}; \nolinkurl{fixes/MISPRINTS_7II.json#1032}\par

\par\medskip\noindent\hypertarget{app-F7II-XX-001.2}{}\textbf{F7II-XX-001.2}\quad p. 341; printed misprint.\par
\noindent Imprimé : {\ttfamily\footnotesize algèbriques\allowbreak{}}\par
\noindent Corrigé : {\ttfamily\footnotesize algébriques\allowbreak{}}\par
\noindent Source : 7II:1033. \nolinkurl{fixes/MISPRINTS_7II.json#1033}\par

\par\medskip\noindent\hypertarget{app-F7II-XX-001.3}{}\textbf{F7II-XX-001.3}\quad p. 341; printed misprint.\par
\noindent Imprimé : {\ttfamily\footnotesize donc\allowbreak{}}\par
\noindent Corrigé : {\ttfamily\footnotesize dont\allowbreak{}}\par
\noindent Source : 7II:602; 7II:1034. \nolinkurl{fixes/MISPRINTS_7II.json#602}; \nolinkurl{fixes/MISPRINTS_7II.json#1034}\par

\par\medskip\noindent\hypertarget{app-F7II-XX-001.4}{}\textbf{F7II-XX-001.4}\quad p. 341; printed misprint.\par
\noindent Imprimé : {\ttfamily\footnotesize au\allowbreak{}}\par
\noindent Corrigé : {\ttfamily\footnotesize ou\allowbreak{}}\par
\noindent Source : 7II:603; 7II:1035. \nolinkurl{fixes/MISPRINTS_7II.json#603}; \nolinkurl{fixes/MISPRINTS_7II.json#1035}\par

\par\medskip\noindent\hypertarget{app-F7II-XX-001.5}{}\textbf{F7II-XX-001.5}\quad p. 341; designated-witness attribution.\par
\noindent Imprimé : {\ttfamily\footnotesize  toutes\allowbreak{}}\par
\noindent Corrigé : \emph{(vide)}\par
\noindent Source : WITNESS-F7II-XX-001.5. \nolinkurl{fr/sga7ii/diplomatic/expose_XX_body.tex}; \nolinkurl{fr/sga7ii/current/expose_XX_body.tex}\par

\par\medskip\noindent\hypertarget{app-F7II-XX-001.6}{}\textbf{F7II-XX-001.6}\quad p. 341; inherited corrected-layer difference.\par
\noindent Imprimé : \emph{(vide)}\par
\noindent Corrigé : \emph{(1 espace(s), 0 retour(s) de ligne dans le témoin TeX)}\par
\noindent Source : ADOPT-F7II-XX-001.6. \nolinkurl{fr/sga7ii/diplomatic/expose_XX_body.tex}; \nolinkurl{fr/sga7ii/current/expose_XX_body.tex}\par

\par\medskip\noindent\hypertarget{app-F7II-XX-002.1}{}\textbf{F7II-XX-002.1}\quad p. 341; inherited corrected-layer difference.\par
\noindent Imprimé : \emph{(1 espace(s), 0 retour(s) de ligne dans le témoin TeX)}\par
\noindent Corrigé : \emph{(vide)}\par
\noindent Source : ADOPT-F7II-XX-002.1. \nolinkurl{fr/sga7ii/diplomatic/expose_XX_body.tex}; \nolinkurl{fr/sga7ii/current/expose_XX_body.tex}\par

\par\medskip\noindent\hypertarget{app-F7II-XX-002.2}{}\textbf{F7II-XX-002.2}\quad p. 341; inherited corrected-layer difference.\par
\noindent Imprimé : {\ttfamily\footnotesize à\allowbreak{}}\par
\noindent Corrigé : {\ttfamily\footnotesize de\allowbreak{}}\par
\noindent Source : ADOPT-F7II-XX-002.2. \nolinkurl{fr/sga7ii/diplomatic/expose_XX_body.tex}; \nolinkurl{fr/sga7ii/current/expose_XX_body.tex}\par

\par\medskip\noindent\hypertarget{app-F7II-XX-002.3}{}\textbf{F7II-XX-002.3}\quad p. 341; inherited corrected-layer difference.\par
\noindent Imprimé : {\ttfamily\footnotesize caractéristiques\allowbreak{}}\par
\noindent Corrigé : {\ttfamily\footnotesize caractéristique\allowbreak{}}\par
\noindent Source : ADOPT-F7II-XX-002.3. \nolinkurl{fr/sga7ii/diplomatic/expose_XX_body.tex}; \nolinkurl{fr/sga7ii/current/expose_XX_body.tex}\par

\par\medskip\noindent\hypertarget{app-F7II-XX-002.4}{}\textbf{F7II-XX-002.4}\quad p. 341; inherited corrected-layer difference.\par
\noindent Imprimé : \emph{(1 espace(s), 0 retour(s) de ligne dans le témoin TeX)}\par
\noindent Corrigé : \emph{(vide)}\par
\noindent Source : ADOPT-F7II-XX-002.4. \nolinkurl{fr/sga7ii/diplomatic/expose_XX_body.tex}; \nolinkurl{fr/sga7ii/current/expose_XX_body.tex}\par

\par\medskip\noindent\hypertarget{app-F7II-XX-003.1}{}\textbf{F7II-XX-003.1}\quad p. 342; printed misprint.\par
\noindent Imprimé : {\ttfamily\footnotesize  \textbackslash{}quad\allowbreak{} \textbackslash{}ell\allowbreak{}}\par
\noindent Corrigé : \emph{(vide)}\par
\noindent Source : 7II:604; 7II:1036. \nolinkurl{fixes/MISPRINTS_7II.json#604}; \nolinkurl{fixes/MISPRINTS_7II.json#1036}\par

\par\medskip\noindent\hypertarget{app-F7II-XX-004.1}{}\textbf{F7II-XX-004.1}\quad p. 342; designated-witness attribution.\par
\noindent Imprimé : {\ttfamily\footnotesize (\allowbreak{}}\par
\noindent Corrigé : \emph{(vide)}\par
\noindent Source : WITNESS-F7II-XX-004.1. \nolinkurl{fr/sga7ii/diplomatic/expose_XX_body.tex}; \nolinkurl{fr/sga7ii/current/expose_XX_body.tex}\par

\par\medskip\noindent\hypertarget{app-F7II-XX-004.2}{}\textbf{F7II-XX-004.2}\quad p. 342; designated-witness attribution.\par
\noindent Imprimé : {\ttfamily\footnotesize (\allowbreak{}}\par
\noindent Corrigé : \emph{(vide)}\par
\noindent Source : WITNESS-F7II-XX-004.2. \nolinkurl{fr/sga7ii/diplomatic/expose_XX_body.tex}; \nolinkurl{fr/sga7ii/current/expose_XX_body.tex}\par

\par\medskip\noindent\hypertarget{app-F7II-XX-004.3}{}\textbf{F7II-XX-004.3}\quad p. 342; designated-witness attribution.\par
\noindent Imprimé : {\ttfamily\footnotesize -\allowbreak{}}\par
\noindent Corrigé : {\ttfamily\footnotesize +\allowbreak{}}\par
\noindent Source : WITNESS-F7II-XX-004.3. \nolinkurl{fr/sga7ii/diplomatic/expose_XX_body.tex}; \nolinkurl{fr/sga7ii/current/expose_XX_body.tex}\par

\par\medskip\noindent\hypertarget{app-F7II-XX-004.4}{}\textbf{F7II-XX-004.4}\quad p. 342; designated-witness attribution.\par
\noindent Imprimé : {\ttfamily\footnotesize )\allowbreak{}}\par
\noindent Corrigé : \emph{(vide)}\par
\noindent Source : WITNESS-F7II-XX-004.4. \nolinkurl{fr/sga7ii/diplomatic/expose_XX_body.tex}; \nolinkurl{fr/sga7ii/current/expose_XX_body.tex}\par

\par\medskip\noindent\hypertarget{app-F7II-XX-004.5}{}\textbf{F7II-XX-004.5}\quad p. 342; printed misprint.\par
\noindent Imprimé : \emph{(vide)}\par
\noindent Corrigé : {\ttfamily\footnotesize .\allowbreak{}}\par
\noindent Source : 7II:1039. \nolinkurl{fixes/MISPRINTS_7II.json#1039}\par

\par\medskip\noindent\hypertarget{app-F7II-XX-004.6}{}\textbf{F7II-XX-004.6}\quad p. 342; printed misprint.\par
\noindent Imprimé : \emph{(vide)}\par
\noindent Corrigé : {\ttfamily\footnotesize \_r\allowbreak{}}\par
\noindent Source : 7II:607; 7II:1040. \nolinkurl{fixes/MISPRINTS_7II.json#607}; \nolinkurl{fixes/MISPRINTS_7II.json#1040}\par

\par\medskip\noindent\hypertarget{app-F7II-XX-004.7}{}\textbf{F7II-XX-004.7}\quad p. 342; printed misprint.\par
\noindent Imprimé : \emph{(vide)}\par
\noindent Corrigé : {\ttfamily\footnotesize \{\}'\allowbreak{}}\par
\noindent Source : 7II:608; 7II:1041. \nolinkurl{fixes/MISPRINTS_7II.json#608}; \nolinkurl{fixes/MISPRINTS_7II.json#1041}\par

\par\medskip\noindent\hypertarget{app-F7II-XX-005.1}{}\textbf{F7II-XX-005.1}\quad p. 343; printed misprint.\par
\noindent Imprimé : {\ttfamily\footnotesize a\allowbreak{}}\par
\noindent Corrigé : {\ttfamily\footnotesize 0\allowbreak{}}\par
\noindent Source : 7II:609; 7II:1042. \nolinkurl{fixes/MISPRINTS_7II.json#609}; \nolinkurl{fixes/MISPRINTS_7II.json#1042}\par

\par\medskip\noindent\hypertarget{app-F7II-XX-006.1}{}\textbf{F7II-XX-006.1}\quad p. 344; designated-witness attribution.\par
\noindent Imprimé : {\ttfamily\footnotesize 2\allowbreak{}}\par
\noindent Corrigé : {\ttfamily\footnotesize 1\allowbreak{}}\par
\noindent Source : WITNESS-F7II-XX-006.1. \nolinkurl{fr/sga7ii/diplomatic/expose_XX_body.tex}; \nolinkurl{fr/sga7ii/current/expose_XX_body.tex}\par

\par\medskip\noindent\hypertarget{app-F7II-XX-007.1}{}\textbf{F7II-XX-007.1}\quad p. 344; printed misprint.\par
\noindent Imprimé : \emph{(2 espace(s), 0 retour(s) de ligne dans le témoin TeX)}\par
\noindent Corrigé : {\ttfamily\footnotesize .\allowbreak{} }\par
\noindent Source : 7II:1043. \nolinkurl{fixes/MISPRINTS_7II.json#1043}\par

\par\medskip\noindent\hypertarget{app-F7II-XX-008.1}{}\textbf{F7II-XX-008.1}\quad p. 344; printed misprint.\par
\noindent Imprimé : \emph{(vide)}\par
\noindent Corrigé : {\ttfamily\footnotesize )\allowbreak{}}\par
\noindent Source : 7II:1044. \nolinkurl{fixes/MISPRINTS_7II.json#1044}\par

\par\medskip\noindent\hypertarget{app-F7II-XX-008.2}{}\textbf{F7II-XX-008.2}\quad p. 344; printed misprint.\par
\noindent Imprimé : {\ttfamily\footnotesize b\allowbreak{}}\par
\noindent Corrigé : {\ttfamily\footnotesize a\allowbreak{}}\par
\noindent Source : 7II:611; 7II:1045. \nolinkurl{fixes/MISPRINTS_7II.json#611}; \nolinkurl{fixes/MISPRINTS_7II.json#1045}\par

\par\medskip\noindent\hypertarget{app-F7II-XX-009.1}{}\textbf{F7II-XX-009.1}\quad p. 345; designated-witness attribution.\par
\noindent Imprimé : \emph{(vide)}\par
\noindent Corrigé : {\ttfamily\footnotesize f\textbackslash{}circ\allowbreak{}}\par
\noindent Source : WITNESS-F7II-XX-009.1. \nolinkurl{fr/sga7ii/diplomatic/expose_XX_body.tex}; \nolinkurl{fr/sga7ii/current/expose_XX_body.tex}\par

\par\medskip\noindent\hypertarget{app-F7II-XX-009.2}{}\textbf{F7II-XX-009.2}\quad p. 345; designated-witness attribution.\par
\noindent Imprimé : {\ttfamily\footnotesize \textbackslash{}\allowbreak{} f\allowbreak{}}\par
\noindent Corrigé : \emph{(vide)}\par
\noindent Source : WITNESS-F7II-XX-009.2. \nolinkurl{fr/sga7ii/diplomatic/expose_XX_body.tex}; \nolinkurl{fr/sga7ii/current/expose_XX_body.tex}\par

\par\medskip\noindent\hypertarget{app-F7II-XX-009.3}{}\textbf{F7II-XX-009.3}\quad p. 345; designated-witness attribution.\par
\noindent Imprimé : \emph{(vide)}\par
\noindent Corrigé : {\ttfamily\footnotesize f\textbackslash{}circ\allowbreak{}}\par
\noindent Source : WITNESS-F7II-XX-009.3. \nolinkurl{fr/sga7ii/diplomatic/expose_XX_body.tex}; \nolinkurl{fr/sga7ii/current/expose_XX_body.tex}\par

\par\medskip\noindent\hypertarget{app-F7II-XX-009.4}{}\textbf{F7II-XX-009.4}\quad p. 345; designated-witness attribution.\par
\noindent Imprimé : {\ttfamily\footnotesize \textbackslash{}\allowbreak{} f\allowbreak{}}\par
\noindent Corrigé : \emph{(vide)}\par
\noindent Source : WITNESS-F7II-XX-009.4. \nolinkurl{fr/sga7ii/diplomatic/expose_XX_body.tex}; \nolinkurl{fr/sga7ii/current/expose_XX_body.tex}\par

\par\medskip\noindent\hypertarget{app-F7II-XX-009.5}{}\textbf{F7II-XX-009.5}\quad p. 345; inherited corrected-layer difference.\par
\noindent Imprimé : {\ttfamily\footnotesize \textbackslash{}\allowbreak{} }\par
\noindent Corrigé : {\ttfamily\footnotesize \textbackslash{}circ\allowbreak{}}\par
\noindent Source : ADOPT-F7II-XX-009.5. \nolinkurl{fr/sga7ii/diplomatic/expose_XX_body.tex}; \nolinkurl{fr/sga7ii/current/expose_XX_body.tex}\par

\par\medskip\noindent\hypertarget{app-F7II-XX-010.1}{}\textbf{F7II-XX-010.1}\quad p. 345; inherited corrected-layer difference.\par
\noindent Imprimé : \emph{(vide)}\par
\noindent Corrigé : {\ttfamily\footnotesize toute\allowbreak{} }\par
\noindent Source : ADOPT-F7II-XX-010.1. \nolinkurl{fr/sga7ii/diplomatic/expose_XX_body.tex}; \nolinkurl{fr/sga7ii/current/expose_XX_body.tex}\par

\par\medskip\noindent\hypertarget{app-F7II-XX-011.1}{}\textbf{F7II-XX-011.1}\quad p. 346; printed misprint.\par
\noindent Imprimé : {\ttfamily\footnotesize annexe\allowbreak{}}\par
\noindent Corrigé : {\ttfamily\footnotesize connexe\allowbreak{}}\par
\noindent Source : 7II:613; 7II:1005. \nolinkurl{fixes/MISPRINTS_7II.json#613}; \nolinkurl{fixes/MISPRINTS_7II.json#1005}\par

\par\medskip\noindent\hypertarget{app-F7II-XX-011.2}{}\textbf{F7II-XX-011.2}\quad p. 346; printed misprint.\par
\noindent Imprimé : {\ttfamily\footnotesize X\allowbreak{}}\par
\noindent Corrigé : {\ttfamily\footnotesize x\allowbreak{}}\par
\noindent Source : 7II:614; 7II:1006. \nolinkurl{fixes/MISPRINTS_7II.json#614}; \nolinkurl{fixes/MISPRINTS_7II.json#1006}\par

\par\medskip\noindent\hypertarget{app-F7II-XX-012.1}{}\textbf{F7II-XX-012.1}\quad p. 347; printed misprint.\par
\noindent Imprimé : {\ttfamily\footnotesize constantstordus\allowbreak{}}\par
\noindent Corrigé : {\ttfamily\footnotesize constants\allowbreak{} tordus\allowbreak{}}\par
\noindent Source : 7II:615; 7II:1007. \nolinkurl{fixes/MISPRINTS_7II.json#615}; \nolinkurl{fixes/MISPRINTS_7II.json#1007}\par

\par\medskip\noindent\hypertarget{app-F7II-XX-013.1}{}\textbf{F7II-XX-013.1}\quad p. 347; printed misprint.\par
\noindent Imprimé : {\ttfamily\footnotesize f\allowbreak{}}\par
\noindent Corrigé : {\ttfamily\footnotesize t\_1\allowbreak{}}\par
\noindent Source : 7II:616; 7II:1008. \nolinkurl{fixes/MISPRINTS_7II.json#616}; \nolinkurl{fixes/MISPRINTS_7II.json#1008}\par

\par\medskip\noindent\hypertarget{app-F7II-XX-013.2}{}\textbf{F7II-XX-013.2}\quad p. 347; printed misprint.\par
\noindent Imprimé : {\ttfamily\footnotesize g\allowbreak{}}\par
\noindent Corrigé : {\ttfamily\footnotesize t\_2\allowbreak{}}\par
\noindent Source : 7II:616; 7II:1008. \nolinkurl{fixes/MISPRINTS_7II.json#616}; \nolinkurl{fixes/MISPRINTS_7II.json#1008}\par

\par\medskip\noindent\hypertarget{app-F7II-XX-014.1}{}\textbf{F7II-XX-014.1}\quad p. 349; printed misprint.\par
\noindent Imprimé : {\ttfamily\footnotesize Z\allowbreak{}}\par
\noindent Corrigé : {\ttfamily\footnotesize z\allowbreak{}}\par
\noindent Source : 7II:617; 7II:1009. \nolinkurl{fixes/MISPRINTS_7II.json#617}; \nolinkurl{fixes/MISPRINTS_7II.json#1009}\par

\par\medskip\noindent\hypertarget{app-F7II-XX-015.1}{}\textbf{F7II-XX-015.1}\quad p. 350; printed misprint.\par
\noindent Imprimé : {\ttfamily\footnotesize ]\allowbreak{} }\par
\noindent Corrigé : \emph{(vide)}\par
\noindent Source : 7II:618; 7II:1010. \nolinkurl{fixes/MISPRINTS_7II.json#618}; \nolinkurl{fixes/MISPRINTS_7II.json#1010}\par

\par\medskip\noindent\hypertarget{app-F7II-XX-016.1}{}\textbf{F7II-XX-016.1}\quad p. 350; printed misprint.\par
\noindent Imprimé : {\ttfamily\footnotesize défini\allowbreak{}}\par
\noindent Corrigé : {\ttfamily\footnotesize définie\allowbreak{}}\par
\noindent Source : 7II:620; 7II:1012. \nolinkurl{fixes/MISPRINTS_7II.json#620}; \nolinkurl{fixes/MISPRINTS_7II.json#1012}\par

\par\medskip\noindent\hypertarget{app-F7II-XX-016.2}{}\textbf{F7II-XX-016.2}\quad p. 350; printed misprint.\par
\noindent Imprimé : {\ttfamily\footnotesize g\allowbreak{}}\par
\noindent Corrigé : {\ttfamily\footnotesize q\allowbreak{}}\par
\noindent Source : 7II:1013. \nolinkurl{fixes/MISPRINTS_7II.json#1013}\par

\par\medskip\noindent\hypertarget{app-F7II-XX-017.1}{}\textbf{F7II-XX-017.1}\quad p. 351; printed misprint.\par
\noindent Imprimé : {\ttfamily\footnotesize ,\allowbreak{}}\par
\noindent Corrigé : {\ttfamily\footnotesize .\allowbreak{}}\par
\noindent Source : 7II:1014. \nolinkurl{fixes/MISPRINTS_7II.json#1014}\par

\par\medskip\noindent\hypertarget{app-F7II-XX-018.1}{}\textbf{F7II-XX-018.1}\quad p. 352; printed misprint.\par
\noindent Imprimé : {\ttfamily\footnotesize hypersurfacesde\allowbreak{}}\par
\noindent Corrigé : {\ttfamily\footnotesize hypersurfaces\allowbreak{} de\allowbreak{}}\par
\noindent Source : 7II:623; 7II:1015. \nolinkurl{fixes/MISPRINTS_7II.json#623}; \nolinkurl{fixes/MISPRINTS_7II.json#1015}\par

\par\medskip\noindent\hypertarget{app-F7II-XX-019.1}{}\textbf{F7II-XX-019.1}\quad p. 352; printed misprint.\par
\noindent Imprimé : {\ttfamily\footnotesize 2n\allowbreak{}}\par
\noindent Corrigé : {\ttfamily\footnotesize n\allowbreak{}}\par
\noindent Source : 7II:624; 7II:1016. \nolinkurl{fixes/MISPRINTS_7II.json#624}; \nolinkurl{fixes/MISPRINTS_7II.json#1016}\par

\par\medskip\noindent\hypertarget{app-F7II-XX-020.1}{}\textbf{F7II-XX-020.1}\quad p. 353; printed misprint.\par
\noindent Imprimé : {\ttfamily\footnotesize contenu\allowbreak{}}\par
\noindent Corrigé : {\ttfamily\footnotesize contenue\allowbreak{}}\par
\noindent Source : 7II:627; 7II:1017. \nolinkurl{fixes/MISPRINTS_7II.json#627}; \nolinkurl{fixes/MISPRINTS_7II.json#1017}\par

\par\medskip\noindent\hypertarget{app-F7II-XX-020.2}{}\textbf{F7II-XX-020.2}\quad p. 353; printed misprint.\par
\noindent Imprimé : {\ttfamily\footnotesize \textbackslash{}quad\allowbreak{}}\par
\noindent Corrigé : {\ttfamily\footnotesize /\allowbreak{}}\par
\noindent Source : 7II:1018. \nolinkurl{fixes/MISPRINTS_7II.json#1018}\par

\par\medskip\noindent\hypertarget{app-F7II-XX-020.3}{}\textbf{F7II-XX-020.3}\quad p. 353; printed misprint.\par
\noindent Imprimé : \emph{(vide)}\par
\noindent Corrigé : {\ttfamily\footnotesize \textbackslash{}mathbb\{\allowbreak{}}\par
\noindent Source : 7II:1018. \nolinkurl{fixes/MISPRINTS_7II.json#1018}\par

\par\medskip\noindent\hypertarget{app-F7II-XX-020.4}{}\textbf{F7II-XX-020.4}\quad p. 353; printed misprint.\par
\noindent Imprimé : \emph{(vide)}\par
\noindent Corrigé : {\ttfamily\footnotesize \}\allowbreak{}}\par
\noindent Source : 7II:1018. \nolinkurl{fixes/MISPRINTS_7II.json#1018}\par

\par\medskip\noindent\hypertarget{app-F7II-XX-020.5}{}\textbf{F7II-XX-020.5}\quad p. 353; inherited corrected-layer difference.\par
\noindent Imprimé : {\ttfamily\footnotesize les\allowbreak{} opérations\allowbreak{}}\par
\noindent Corrigé : {\ttfamily\footnotesize l'action\allowbreak{}}\par
\noindent Source : ADOPT-F7II-XX-020.5. \nolinkurl{fr/sga7ii/diplomatic/expose_XX_body.tex}; \nolinkurl{fr/sga7ii/current/expose_XX_body.tex}\par

\par\medskip\noindent\hypertarget{app-F7II-XX-020.6}{}\textbf{F7II-XX-020.6}\quad p. 353; printed misprint.\par
\noindent Imprimé : \emph{(vide)}\par
\noindent Corrigé : {\ttfamily\footnotesize )\allowbreak{}}\par
\noindent Source : 7II:1019. \nolinkurl{fixes/MISPRINTS_7II.json#1019}\par

\par\medskip\noindent\hypertarget{app-F7II-XX-021.1}{}\textbf{F7II-XX-021.1}\quad p. 354; misprint\_in\_print.\par
\noindent Imprimé : {\ttfamily\footnotesize ;\allowbreak{}}\par
\noindent Corrigé : {\ttfamily\footnotesize :\allowbreak{}}\par
\noindent Source : FIX-7II-D-0741. \nolinkurl{fixes/FIXLOG_7II.jsonl#741}\par

\par\medskip\noindent\hypertarget{app-F7II-XX-021.2}{}\textbf{F7II-XX-021.2}\quad p. 354; printed misprint.\par
\noindent Imprimé : \emph{(vide)}\par
\noindent Corrigé : {\ttfamily\footnotesize \textbackslash{}overline\{\allowbreak{}}\par
\noindent Source : 7II:1022. \nolinkurl{fixes/MISPRINTS_7II.json#1022}\par

\par\medskip\noindent\hypertarget{app-F7II-XX-021.3}{}\textbf{F7II-XX-021.3}\quad p. 354; printed misprint.\par
\noindent Imprimé : \emph{(vide)}\par
\noindent Corrigé : {\ttfamily\footnotesize \}\allowbreak{}}\par
\noindent Source : 7II:1022. \nolinkurl{fixes/MISPRINTS_7II.json#1022}\par

\par\medskip\noindent\hypertarget{app-F7II-XX-021.4}{}\textbf{F7II-XX-021.4}\quad p. 354; printed misprint.\par
\noindent Imprimé : \emph{(vide)}\par
\noindent Corrigé : {\ttfamily\footnotesize )\allowbreak{}}\par
\noindent Source : 7II:1023. \nolinkurl{fixes/MISPRINTS_7II.json#1023}\par

\par\medskip\noindent\hypertarget{app-F7II-XX-022.1}{}\textbf{F7II-XX-022.1}\quad p. 355; printed misprint.\par
\noindent Imprimé : \emph{(vide)}\par
\noindent Corrigé : {\ttfamily\footnotesize .\allowbreak{}}\par
\noindent Source : 7II:630; 7II:1024. \nolinkurl{fixes/MISPRINTS_7II.json#630}; \nolinkurl{fixes/MISPRINTS_7II.json#1024}\par

\par\medskip\noindent\hypertarget{app-F7II-XX-023.1}{}\textbf{F7II-XX-023.1}\quad p. 355; misprint\_pass.\par
\noindent Imprimé : \emph{(1 espace(s), 0 retour(s) de ligne dans le témoin TeX)}\par
\noindent Corrigé : {\ttfamily\footnotesize .\allowbreak{}}\par
\noindent Source : 7II:1025; 7II:631. \nolinkurl{fixes/FIXLOG_7II_CURRENT.jsonl#891}; \nolinkurl{fixes/MISPRINTS_7II.json#631}; \nolinkurl{fixes/MISPRINTS_7II.json#1025}\par

\par\medskip\noindent\hypertarget{app-F7II-XX-023.2}{}\textbf{F7II-XX-023.2}\quad p. 355; printed misprint.\par
\noindent Imprimé : {\ttfamily\footnotesize \textasciicircum{}*\allowbreak{}}\par
\noindent Corrigé : \emph{(vide)}\par
\noindent Source : 7II:632. \nolinkurl{fixes/MISPRINTS_7II.json#632}\par

\par\medskip\noindent\hypertarget{app-F7II-XX-023.3}{}\textbf{F7II-XX-023.3}\quad p. 355; printed misprint.\par
\noindent Imprimé : \emph{(vide)}\par
\noindent Corrigé : \emph{(1 espace(s), 0 retour(s) de ligne dans le témoin TeX)}\par
\noindent Source : 7II:1027. \nolinkurl{fixes/MISPRINTS_7II.json#1027}\par

\par\medskip\noindent\hypertarget{app-F7II-XX-024.1}{}\textbf{F7II-XX-024.1}\quad p. 355; inherited corrected-layer difference.\par
\noindent Imprimé : \emph{(vide)}\par
\noindent Corrigé : {\ttfamily\footnotesize l'\allowbreak{}}\par
\noindent Source : ADOPT-F7II-XX-024.1. \nolinkurl{fr/sga7ii/diplomatic/expose_XX_body.tex}; \nolinkurl{fr/sga7ii/current/expose_XX_body.tex}\par

\par\medskip\noindent\hypertarget{app-F7II-XX-024.2}{}\textbf{F7II-XX-024.2}\quad p. 355; inherited corrected-layer difference.\par
\noindent Imprimé : {\ttfamily\footnotesize ,\allowbreak{}}\par
\noindent Corrigé : \emph{(vide)}\par
\noindent Source : ADOPT-F7II-XX-024.2. \nolinkurl{fr/sga7ii/diplomatic/expose_XX_body.tex}; \nolinkurl{fr/sga7ii/current/expose_XX_body.tex}\par

\par\medskip\noindent\hypertarget{app-F7II-XX-024.3}{}\textbf{F7II-XX-024.3}\quad p. 355; inherited corrected-layer difference.\par
\noindent Imprimé : \emph{(1 espace(s), 0 retour(s) de ligne dans le témoin TeX)}\par
\noindent Corrigé : \emph{(vide)}\par
\noindent Source : ADOPT-F7II-XX-024.3. \nolinkurl{fr/sga7ii/diplomatic/expose_XX_body.tex}; \nolinkurl{fr/sga7ii/current/expose_XX_body.tex}\par

\par\medskip\noindent\hypertarget{app-F7II-XX-024.4}{}\textbf{F7II-XX-024.4}\quad p. 355; printed misprint.\par
\noindent Imprimé : \emph{(2 espace(s), 0 retour(s) de ligne dans le témoin TeX)}\par
\noindent Corrigé : {\ttfamily\footnotesize .\allowbreak{} }\par
\noindent Source : 7II:633; 7II:1028. \nolinkurl{fixes/MISPRINTS_7II.json#633}; \nolinkurl{fixes/MISPRINTS_7II.json#1028}\par

\par\medskip\noindent\hypertarget{app-F7II-XX-024.5}{}\textbf{F7II-XX-024.5}\quad p. 355; printed misprint.\par
\noindent Imprimé : \emph{(1 espace(s), 0 retour(s) de ligne dans le témoin TeX)}\par
\noindent Corrigé : {\ttfamily\footnotesize .\allowbreak{}}\par
\noindent Source : 7II:634; 7II:1029. \nolinkurl{fixes/MISPRINTS_7II.json#634}; \nolinkurl{fixes/MISPRINTS_7II.json#1029}\par

\par\medskip\noindent\hypertarget{app-F7II-XX-024.6}{}\textbf{F7II-XX-024.6}\quad p. 355; printed misprint.\par
\noindent Imprimé : \emph{(vide)}\par
\noindent Corrigé : \emph{(1 espace(s), 0 retour(s) de ligne dans le témoin TeX)}\par
\noindent Source : 7II:1030. \nolinkurl{fixes/MISPRINTS_7II.json#1030}\par

\par\medskip\noindent\hypertarget{app-F7II-XX-024.7}{}\textbf{F7II-XX-024.7}\quad p. 355; designated-witness attribution.\par
\noindent Imprimé : \emph{(1 espace(s), 0 retour(s) de ligne dans le témoin TeX)}\par
\noindent Corrigé : \emph{(vide)}\par
\noindent Source : WITNESS-F7II-XX-024.7. \nolinkurl{fr/sga7ii/diplomatic/expose_XX_body.tex}; \nolinkurl{fr/sga7ii/current/expose_XX_body.tex}\par

\par\medskip\noindent\hypertarget{app-F7II-XX-024.8}{}\textbf{F7II-XX-024.8}\quad p. 355; printed misprint.\par
\noindent Imprimé : \emph{(vide)}\par
\noindent Corrigé : \emph{(1 espace(s), 0 retour(s) de ligne dans le témoin TeX)}\par
\noindent Source : 7II:1031. \nolinkurl{fixes/MISPRINTS_7II.json#1031}\par

\par\medskip\noindent\hypertarget{app-F7II-XX-025.1}{}\textbf{F7II-XX-025.1}\quad p. 356; inherited corrected-layer difference.\par
\noindent Imprimé : \emph{(vide)}\par
\noindent Corrigé : {\ttfamily\footnotesize ,\allowbreak{}}\par
\noindent Source : ADOPT-F7II-XX-025.1. \nolinkurl{fr/sga7ii/diplomatic/expose_XX_body.tex}; \nolinkurl{fr/sga7ii/current/expose_XX_body.tex}\par

\par\medskip\noindent\hypertarget{app-F7II-XX-025.2}{}\textbf{F7II-XX-025.2}\quad p. 356; designated-witness attribution.\par
\noindent Imprimé : {\ttfamily\footnotesize sont\allowbreak{}}\par
\noindent Corrigé : {\ttfamily\footnotesize soient\allowbreak{}}\par
\noindent Source : WITNESS-F7II-XX-025.2. \nolinkurl{fr/sga7ii/diplomatic/expose_XX_body.tex}; \nolinkurl{fr/sga7ii/current/expose_XX_body.tex}\par

\par\medskip\noindent\hypertarget{app-F7II-XX-025.3}{}\textbf{F7II-XX-025.3}\quad p. 356; designated-witness attribution.\par
\noindent Imprimé : \emph{(1 espace(s), 0 retour(s) de ligne dans le témoin TeX)}\par
\noindent Corrigé : \emph{(vide)}\par
\noindent Source : WITNESS-F7II-XX-025.3. \nolinkurl{fr/sga7ii/diplomatic/expose_XX_body.tex}; \nolinkurl{fr/sga7ii/current/expose_XX_body.tex}\par

\par\medskip\noindent\hypertarget{app-F7II-XX-025.4}{}\textbf{F7II-XX-025.4}\quad p. 356; printed misprint.\par
\noindent Imprimé : \emph{(2 espace(s), 0 retour(s) de ligne dans le témoin TeX)}\par
\noindent Corrigé : {\ttfamily\footnotesize .\allowbreak{} }\par
\noindent Source : 7II:1060. \nolinkurl{fixes/MISPRINTS_7II.json#1060}\par

\par\medskip\noindent\hypertarget{app-F7II-XX-025.5}{}\textbf{F7II-XX-025.5}\quad p. 356; inherited corrected-layer difference.\par
\noindent Imprimé : {\ttfamily\footnotesize vérifiés\allowbreak{}}\par
\noindent Corrigé : {\ttfamily\footnotesize vérifiées\allowbreak{}}\par
\noindent Source : ADOPT-F7II-XX-025.5. \nolinkurl{fr/sga7ii/diplomatic/expose_XX_body.tex}; \nolinkurl{fr/sga7ii/current/expose_XX_body.tex}\par

\par\medskip\noindent\hypertarget{app-F7II-XX-025.6}{}\textbf{F7II-XX-025.6}\quad p. 356; inherited corrected-layer difference.\par
\noindent Imprimé : {\ttfamily\footnotesize en\allowbreak{}}\par
\noindent Corrigé : {\ttfamily\footnotesize dans\allowbreak{}}\par
\noindent Source : ADOPT-F7II-XX-025.6. \nolinkurl{fr/sga7ii/diplomatic/expose_XX_body.tex}; \nolinkurl{fr/sga7ii/current/expose_XX_body.tex}\par

\par\medskip\noindent\hypertarget{app-F7II-XX-025.7}{}\textbf{F7II-XX-025.7}\quad p. 356; inherited corrected-layer difference.\par
\noindent Imprimé : \emph{(vide)}\par
\noindent Corrigé : {\ttfamily\footnotesize  les\allowbreak{}}\par
\noindent Source : ADOPT-F7II-XX-025.7. \nolinkurl{fr/sga7ii/diplomatic/expose_XX_body.tex}; \nolinkurl{fr/sga7ii/current/expose_XX_body.tex}\par

\par\medskip\noindent\hypertarget{app-F7II-XX-025.8}{}\textbf{F7II-XX-025.8}\quad p. 356; inherited corrected-layer difference.\par
\noindent Imprimé : {\ttfamily\footnotesize en\allowbreak{}}\par
\noindent Corrigé : {\ttfamily\footnotesize dans\allowbreak{}}\par
\noindent Source : ADOPT-F7II-XX-025.8. \nolinkurl{fr/sga7ii/diplomatic/expose_XX_body.tex}; \nolinkurl{fr/sga7ii/current/expose_XX_body.tex}\par

\par\medskip\noindent\hypertarget{app-F7II-XX-025.9}{}\textbf{F7II-XX-025.9}\quad p. 356; inherited corrected-layer difference.\par
\noindent Imprimé : \emph{(vide)}\par
\noindent Corrigé : {\ttfamily\footnotesize  les\allowbreak{}}\par
\noindent Source : ADOPT-F7II-XX-025.9. \nolinkurl{fr/sga7ii/diplomatic/expose_XX_body.tex}; \nolinkurl{fr/sga7ii/current/expose_XX_body.tex}\par

\par\medskip\noindent\hypertarget{app-F7II-XX-025.10}{}\textbf{F7II-XX-025.10}\quad p. 356; inherited corrected-layer difference.\par
\noindent Imprimé : {\ttfamily\footnotesize  car\allowbreak{}}\par
\noindent Corrigé : \emph{(vide)}\par
\noindent Source : ADOPT-F7II-XX-025.10. \nolinkurl{fr/sga7ii/diplomatic/expose_XX_body.tex}; \nolinkurl{fr/sga7ii/current/expose_XX_body.tex}\par

\par\medskip\noindent\hypertarget{app-F7II-XX-025.11}{}\textbf{F7II-XX-025.11}\quad p. 356; inherited corrected-layer difference.\par
\noindent Imprimé : \emph{(vide)}\par
\noindent Corrigé : {\ttfamily\footnotesize \textbackslash{}operatorname\{car\}(\allowbreak{}}\par
\noindent Source : ADOPT-F7II-XX-025.11. \nolinkurl{fr/sga7ii/diplomatic/expose_XX_body.tex}; \nolinkurl{fr/sga7ii/current/expose_XX_body.tex}\par

\par\medskip\noindent\hypertarget{app-F7II-XX-025.12}{}\textbf{F7II-XX-025.12}\quad p. 356; inherited corrected-layer difference.\par
\noindent Imprimé : \emph{(1 espace(s), 0 retour(s) de ligne dans le témoin TeX)}\par
\noindent Corrigé : {\ttfamily\footnotesize )\allowbreak{}}\par
\noindent Source : ADOPT-F7II-XX-025.12. \nolinkurl{fr/sga7ii/diplomatic/expose_XX_body.tex}; \nolinkurl{fr/sga7ii/current/expose_XX_body.tex}\par

\par\medskip\noindent\hypertarget{app-F7II-XX-025.13}{}\textbf{F7II-XX-025.13}\quad p. 356; inherited corrected-layer difference.\par
\noindent Imprimé : \emph{(1 espace(s), 0 retour(s) de ligne dans le témoin TeX)}\par
\noindent Corrigé : \emph{(vide)}\par
\noindent Source : ADOPT-F7II-XX-025.13. \nolinkurl{fr/sga7ii/diplomatic/expose_XX_body.tex}; \nolinkurl{fr/sga7ii/current/expose_XX_body.tex}\par

\par\medskip\noindent\hypertarget{app-F7II-XX-025.14}{}\textbf{F7II-XX-025.14}\quad p. 356; inherited corrected-layer difference.\par
\noindent Imprimé : {\ttfamily\footnotesize  car\allowbreak{}}\par
\noindent Corrigé : \emph{(vide)}\par
\noindent Source : ADOPT-F7II-XX-025.14. \nolinkurl{fr/sga7ii/diplomatic/expose_XX_body.tex}; \nolinkurl{fr/sga7ii/current/expose_XX_body.tex}\par

\par\medskip\noindent\hypertarget{app-F7II-XX-025.15}{}\textbf{F7II-XX-025.15}\quad p. 356; inherited corrected-layer difference.\par
\noindent Imprimé : \emph{(vide)}\par
\noindent Corrigé : {\ttfamily\footnotesize \textbackslash{}operatorname\{car\}(\allowbreak{}}\par
\noindent Source : ADOPT-F7II-XX-025.15. \nolinkurl{fr/sga7ii/diplomatic/expose_XX_body.tex}; \nolinkurl{fr/sga7ii/current/expose_XX_body.tex}\par

\par\medskip\noindent\hypertarget{app-F7II-XX-025.16}{}\textbf{F7II-XX-025.16}\quad p. 356; inherited corrected-layer difference.\par
\noindent Imprimé : \emph{(1 espace(s), 0 retour(s) de ligne dans le témoin TeX)}\par
\noindent Corrigé : {\ttfamily\footnotesize )\allowbreak{}}\par
\noindent Source : ADOPT-F7II-XX-025.16. \nolinkurl{fr/sga7ii/diplomatic/expose_XX_body.tex}; \nolinkurl{fr/sga7ii/current/expose_XX_body.tex}\par

\par\medskip\noindent\hypertarget{app-F7II-XX-025.17}{}\textbf{F7II-XX-025.17}\quad p. 356; printed misprint.\par
\noindent Imprimé : {\ttfamily\footnotesize plausibles\allowbreak{}}\par
\noindent Corrigé : {\ttfamily\footnotesize plausible\allowbreak{}}\par
\noindent Source : 7II:635; 7II:1061. \nolinkurl{fixes/MISPRINTS_7II.json#635}; \nolinkurl{fixes/MISPRINTS_7II.json#1061}\par

\par\medskip\noindent\hypertarget{app-F7II-XX-025.18}{}\textbf{F7II-XX-025.18}\quad p. 356; designated-witness attribution.\par
\noindent Imprimé : \emph{(vide)}\par
\noindent Corrigé : {\ttfamily\footnotesize \$\textbackslash{}operatorname\{\allowbreak{}}\par
\noindent Source : WITNESS-F7II-XX-025.18. \nolinkurl{fr/sga7ii/diplomatic/expose_XX_body.tex}; \nolinkurl{fr/sga7ii/current/expose_XX_body.tex}\par

\par\medskip\noindent\hypertarget{app-F7II-XX-025.19}{}\textbf{F7II-XX-025.19}\quad p. 356; designated-witness attribution.\par
\noindent Imprimé : {\ttfamily\footnotesize  \$\allowbreak{}}\par
\noindent Corrigé : {\ttfamily\footnotesize \}(\allowbreak{}}\par
\noindent Source : WITNESS-F7II-XX-025.19. \nolinkurl{fr/sga7ii/diplomatic/expose_XX_body.tex}; \nolinkurl{fr/sga7ii/current/expose_XX_body.tex}\par

\par\medskip\noindent\hypertarget{app-F7II-XX-025.20}{}\textbf{F7II-XX-025.20}\quad p. 356; designated-witness attribution.\par
\noindent Imprimé : \emph{(1 espace(s), 0 retour(s) de ligne dans le témoin TeX)}\par
\noindent Corrigé : {\ttfamily\footnotesize )\allowbreak{}}\par
\noindent Source : WITNESS-F7II-XX-025.20. \nolinkurl{fr/sga7ii/diplomatic/expose_XX_body.tex}; \nolinkurl{fr/sga7ii/current/expose_XX_body.tex}\par

\par\medskip\noindent\hypertarget{app-F7II-XX-025.21}{}\textbf{F7II-XX-025.21}\quad p. 356; printed misprint.\par
\noindent Imprimé : \emph{(vide)}\par
\noindent Corrigé : {\ttfamily\footnotesize .\allowbreak{}}\par
\noindent Source : 7II:1062. \nolinkurl{fixes/MISPRINTS_7II.json#1062}\par

\par\medskip\noindent\hypertarget{app-F7II-XX-025.22}{}\textbf{F7II-XX-025.22}\quad p. 356; printed misprint.\par
\noindent Imprimé : \emph{(2 espace(s), 0 retour(s) de ligne dans le témoin TeX)}\par
\noindent Corrigé : {\ttfamily\footnotesize .\allowbreak{} }\par
\noindent Source : 7II:1063. \nolinkurl{fixes/MISPRINTS_7II.json#1063}\par

\par\medskip\noindent\hypertarget{app-F7II-XX-025.23}{}\textbf{F7II-XX-025.23}\quad p. 356; designated-witness attribution.\par
\noindent Imprimé : \emph{(vide)}\par
\noindent Corrigé : {\ttfamily\footnotesize ,\allowbreak{}}\par
\noindent Source : WITNESS-F7II-XX-025.23. \nolinkurl{fr/sga7ii/diplomatic/expose_XX_body.tex}; \nolinkurl{fr/sga7ii/current/expose_XX_body.tex}\par

\par\medskip\noindent\hypertarget{app-F7II-XX-025.24}{}\textbf{F7II-XX-025.24}\quad p. 356; inherited corrected-layer difference.\par
\noindent Imprimé : \emph{(1 espace(s), 0 retour(s) de ligne dans le témoin TeX)}\par
\noindent Corrigé : \emph{(vide)}\par
\noindent Source : ADOPT-F7II-XX-025.24. \nolinkurl{fr/sga7ii/diplomatic/expose_XX_body.tex}; \nolinkurl{fr/sga7ii/current/expose_XX_body.tex}\par

\par\medskip\noindent\hypertarget{app-F7II-XX-025.25}{}\textbf{F7II-XX-025.25}\quad p. 356; inherited corrected-layer difference.\par
\noindent Imprimé : \emph{(vide)}\par
\noindent Corrigé : {\ttfamily\footnotesize ,\allowbreak{}}\par
\noindent Source : ADOPT-F7II-XX-025.25. \nolinkurl{fr/sga7ii/diplomatic/expose_XX_body.tex}; \nolinkurl{fr/sga7ii/current/expose_XX_body.tex}\par

\par\medskip\noindent\hypertarget{app-F7II-XX-025.26}{}\textbf{F7II-XX-025.26}\quad p. 356; inherited corrected-layer difference.\par
\noindent Imprimé : \emph{(1 espace(s), 0 retour(s) de ligne dans le témoin TeX)}\par
\noindent Corrigé : \emph{(vide)}\par
\noindent Source : ADOPT-F7II-XX-025.26. \nolinkurl{fr/sga7ii/diplomatic/expose_XX_body.tex}; \nolinkurl{fr/sga7ii/current/expose_XX_body.tex}\par

\par\medskip\noindent\hypertarget{app-F7II-XX-026.1}{}\textbf{F7II-XX-026.1}\quad p. 356; printed misprint.\par
\noindent Imprimé : {\ttfamily\footnotesize n\allowbreak{}}\par
\noindent Corrigé : {\ttfamily\footnotesize m\allowbreak{}}\par
\noindent Source : 7II:1064. \nolinkurl{fixes/MISPRINTS_7II.json#1064}\par

\par\medskip\noindent\hypertarget{app-F7II-XX-026.2}{}\textbf{F7II-XX-026.2}\quad p. 356; designated-witness attribution.\par
\noindent Imprimé : \emph{(1 espace(s), 0 retour(s) de ligne dans le témoin TeX)}\par
\noindent Corrigé : \emph{(vide)}\par
\noindent Source : WITNESS-F7II-XX-026.2. \nolinkurl{fr/sga7ii/diplomatic/expose_XX_body.tex}; \nolinkurl{fr/sga7ii/current/expose_XX_body.tex}\par

\par\medskip\noindent\hypertarget{app-F7II-XX-027.1}{}\textbf{F7II-XX-027.1}\quad p. 357; designated-witness attribution.\par
\noindent Imprimé : \emph{(vide)}\par
\noindent Corrigé : {\ttfamily\footnotesize \_\{\textbackslash{}bar\textbackslash{}eta\}\allowbreak{}}\par
\noindent Source : WITNESS-F7II-XX-027.1. \nolinkurl{fr/sga7ii/diplomatic/expose_XX_body.tex}; \nolinkurl{fr/sga7ii/current/expose_XX_body.tex}\par

\par\medskip\noindent\hypertarget{app-F7II-XX-027.2}{}\textbf{F7II-XX-027.2}\quad p. 357; printed misprint.\par
\noindent Imprimé : {\ttfamily\footnotesize n\allowbreak{}}\par
\noindent Corrigé : {\ttfamily\footnotesize m\allowbreak{}}\par
\noindent Source : 7II:1066. \nolinkurl{fixes/MISPRINTS_7II.json#1066}\par

\par\medskip\noindent\hypertarget{app-F7II-XX-028.1}{}\textbf{F7II-XX-028.1}\quad p. 357; printed misprint.\par
\noindent Imprimé : {\ttfamily\footnotesize (\allowbreak{}}\par
\noindent Corrigé : \emph{(vide)}\par
\noindent Source : 7II:636; 7II:1067. \nolinkurl{fixes/MISPRINTS_7II.json#636}; \nolinkurl{fixes/MISPRINTS_7II.json#1067}\par

\par\medskip\noindent\hypertarget{app-F7II-XX-028.2}{}\textbf{F7II-XX-028.2}\quad p. 357; printed misprint.\par
\noindent Imprimé : \emph{(vide)}\par
\noindent Corrigé : {\ttfamily\footnotesize \_\{\textbackslash{}bar\{\textbackslash{}eta\}\}\allowbreak{}}\par
\noindent Source : 7II:1068. \nolinkurl{fixes/MISPRINTS_7II.json#1068}\par

\par\medskip\noindent\hypertarget{app-F7II-XX-029.1}{}\textbf{F7II-XX-029.1}\quad p. 357; designated-witness attribution.\par
\noindent Imprimé : \emph{(1 espace(s), 0 retour(s) de ligne dans le témoin TeX)}\par
\noindent Corrigé : \emph{(vide)}\par
\noindent Source : WITNESS-F7II-XX-029.1. \nolinkurl{fr/sga7ii/diplomatic/expose_XX_body.tex}; \nolinkurl{fr/sga7ii/current/expose_XX_body.tex}\par

\par\medskip\noindent\hypertarget{app-F7II-XX-029.2}{}\textbf{F7II-XX-029.2}\quad p. 357; printed misprint.\par
\noindent Imprimé : {\ttfamily\footnotesize \_\{\textbackslash{}bar\{\textbackslash{}eta\}\}\allowbreak{} \textbackslash{};\textbackslash{};\allowbreak{} \{\}\allowbreak{}}\par
\noindent Corrigé : \emph{(vide)}\par
\noindent Source : 7II:1069. \nolinkurl{fixes/MISPRINTS_7II.json#1069}\par

\par\medskip\noindent\hypertarget{app-F7II-XX-029.3}{}\textbf{F7II-XX-029.3}\quad p. 357; designated-witness attribution.\par
\noindent Imprimé : \emph{(1 espace(s), 0 retour(s) de ligne dans le témoin TeX)}\par
\noindent Corrigé : \emph{(vide)}\par
\noindent Source : WITNESS-F7II-XX-029.3. \nolinkurl{fr/sga7ii/diplomatic/expose_XX_body.tex}; \nolinkurl{fr/sga7ii/current/expose_XX_body.tex}\par

\par\medskip\noindent\hypertarget{app-F7II-XX-030.1}{}\textbf{F7II-XX-030.1}\quad p. 358; printed misprint.\par
\noindent Imprimé : \emph{(vide)}\par
\noindent Corrigé : {\ttfamily\footnotesize  -\allowbreak{}}\par
\noindent Source : 7II:1070. \nolinkurl{fixes/MISPRINTS_7II.json#1070}\par

\par\medskip\noindent\hypertarget{app-F7II-XX-031.1}{}\textbf{F7II-XX-031.1}\quad p. 358; designated-witness attribution.\par
\noindent Imprimé : \emph{(vide)}\par
\noindent Corrigé : {\ttfamily\footnotesize .\allowbreak{}}\par
\noindent Source : WITNESS-F7II-XX-031.1. \nolinkurl{fr/sga7ii/diplomatic/expose_XX_body.tex}; \nolinkurl{fr/sga7ii/current/expose_XX_body.tex}\par

\par\medskip\noindent\hypertarget{app-F7II-XX-031.2}{}\textbf{F7II-XX-031.2}\quad p. 358; printed misprint.\par
\noindent Imprimé : {\ttfamily\footnotesize Choissisons\allowbreak{}}\par
\noindent Corrigé : {\ttfamily\footnotesize Choisissons\allowbreak{}}\par
\noindent Source : 7II:640; 7II:1072. \nolinkurl{fixes/MISPRINTS_7II.json#640}; \nolinkurl{fixes/MISPRINTS_7II.json#1072}\par

\par\medskip\noindent\hypertarget{app-F7II-XX-032.1}{}\textbf{F7II-XX-032.1}\quad p. 359; printed misprint.\par
\noindent Imprimé : {\ttfamily\footnotesize i\allowbreak{}}\par
\noindent Corrigé : {\ttfamily\footnotesize 1\allowbreak{}}\par
\noindent Source : 7II:1073. \nolinkurl{fixes/MISPRINTS_7II.json#1073}\par

\par\medskip\noindent\hypertarget{app-F7II-XX-033.1}{}\textbf{F7II-XX-033.1}\quad p. 360; printed misprint.\par
\noindent Imprimé : {\ttfamily\footnotesize n\allowbreak{}}\par
\noindent Corrigé : {\ttfamily\footnotesize m\allowbreak{}}\par
\noindent Source : 7II:1074. \nolinkurl{fixes/MISPRINTS_7II.json#1074}\par

\par\medskip\noindent\hypertarget{app-F7II-XX-033.2}{}\textbf{F7II-XX-033.2}\quad p. 360; printed misprint.\par
\noindent Imprimé : \emph{(vide)}\par
\noindent Corrigé : {\ttfamily\footnotesize )\allowbreak{}}\par
\noindent Source : 7II:1075. \nolinkurl{fixes/MISPRINTS_7II.json#1075}\par

\par\medskip\noindent\hypertarget{app-F7II-XX-033.3}{}\textbf{F7II-XX-033.3}\quad p. 360; designated-witness attribution.\par
\noindent Imprimé : \emph{(vide)}\par
\noindent Corrigé : {\ttfamily\footnotesize \textbackslash{}bigl\allowbreak{}}\par
\noindent Source : WITNESS-F7II-XX-033.3. \nolinkurl{fr/sga7ii/diplomatic/expose_XX_body.tex}; \nolinkurl{fr/sga7ii/current/expose_XX_body.tex}\par

\par\medskip\noindent\hypertarget{app-F7II-XX-033.4}{}\textbf{F7II-XX-033.4}\quad p. 360; designated-witness attribution.\par
\noindent Imprimé : \emph{(vide)}\par
\noindent Corrigé : {\ttfamily\footnotesize )\textbackslash{}bigr)\allowbreak{}}\par
\noindent Source : WITNESS-F7II-XX-033.4. \nolinkurl{fr/sga7ii/diplomatic/expose_XX_body.tex}; \nolinkurl{fr/sga7ii/current/expose_XX_body.tex}\par

\par\medskip\noindent\hypertarget{app-F7II-XX-033.5}{}\textbf{F7II-XX-033.5}\quad p. 360; designated-witness attribution.\par
\noindent Imprimé : \emph{(2 espace(s), 0 retour(s) de ligne dans le témoin TeX)}\par
\noindent Corrigé : \emph{(vide)}\par
\noindent Source : WITNESS-F7II-XX-033.5. \nolinkurl{fr/sga7ii/diplomatic/expose_XX_body.tex}; \nolinkurl{fr/sga7ii/current/expose_XX_body.tex}\par

\par\medskip\noindent\hypertarget{app-F7II-XX-033.6}{}\textbf{F7II-XX-033.6}\quad p. 360; designated-witness attribution.\par
\noindent Imprimé : \emph{(vide)}\par
\noindent Corrigé : {\ttfamily\footnotesize ,\allowbreak{}}\par
\noindent Source : WITNESS-F7II-XX-033.6. \nolinkurl{fr/sga7ii/diplomatic/expose_XX_body.tex}; \nolinkurl{fr/sga7ii/current/expose_XX_body.tex}\par

\par\medskip\noindent\hypertarget{app-F7II-XX-033.7}{}\textbf{F7II-XX-033.7}\quad p. 360; printed misprint.\par
\noindent Imprimé : {\ttfamily\footnotesize n\allowbreak{}}\par
\noindent Corrigé : {\ttfamily\footnotesize m\allowbreak{}}\par
\noindent Source : 7II:1078. \nolinkurl{fixes/MISPRINTS_7II.json#1078}\par

\par\medskip\noindent\hypertarget{app-F7II-XX-033.8}{}\textbf{F7II-XX-033.8}\quad p. 360; printed misprint.\par
\noindent Imprimé : \emph{(vide)}\par
\noindent Corrigé : {\ttfamily\footnotesize )\allowbreak{}}\par
\noindent Source : 7II:1079. \nolinkurl{fixes/MISPRINTS_7II.json#1079}\par

\par\medskip\noindent\hypertarget{app-F7II-XX-034.1}{}\textbf{F7II-XX-034.1}\quad p. 360; printed misprint.\par
\noindent Imprimé : {\ttfamily\footnotesize 2n\allowbreak{}}\par
\noindent Corrigé : {\ttfamily\footnotesize 2m\allowbreak{}}\par
\noindent Source : 7II:646; 7II:1080. \nolinkurl{fixes/MISPRINTS_7II.json#646}; \nolinkurl{fixes/MISPRINTS_7II.json#1080}\par

\par\medskip\noindent\hypertarget{app-F7II-XX-034.2}{}\textbf{F7II-XX-034.2}\quad p. 360; printed misprint.\par
\noindent Imprimé : {\ttfamily\footnotesize n\allowbreak{}}\par
\noindent Corrigé : {\ttfamily\footnotesize m\allowbreak{}}\par
\noindent Source : 7II:646; 7II:1080. \nolinkurl{fixes/MISPRINTS_7II.json#646}; \nolinkurl{fixes/MISPRINTS_7II.json#1080}\par

\par\medskip\noindent\hypertarget{app-F7II-XX-034.3}{}\textbf{F7II-XX-034.3}\quad p. 360; printed misprint.\par
\noindent Imprimé : {\ttfamily\footnotesize 3\allowbreak{}}\par
\noindent Corrigé : {\ttfamily\footnotesize 4\allowbreak{}}\par
\noindent Source : 7II:645; 7II:1081. \nolinkurl{fixes/MISPRINTS_7II.json#645}; \nolinkurl{fixes/MISPRINTS_7II.json#1081}\par

\par\medskip\noindent\hypertarget{app-F7II-XX-034.4}{}\textbf{F7II-XX-034.4}\quad p. 360; designated-witness attribution.\par
\noindent Imprimé : {\ttfamily\footnotesize Groupe\allowbreak{}}\par
\noindent Corrigé : {\ttfamily\footnotesize groupe\allowbreak{}}\par
\noindent Source : WITNESS-F7II-XX-034.4. \nolinkurl{fr/sga7ii/diplomatic/expose_XX_body.tex}; \nolinkurl{fr/sga7ii/current/expose_XX_body.tex}\par

\par\medskip\noindent\hypertarget{app-F7II-XX-035.1}{}\textbf{F7II-XX-035.1}\quad p. 361; printed misprint.\par
\noindent Imprimé : {\ttfamily\footnotesize équivalent\allowbreak{}}\par
\noindent Corrigé : {\ttfamily\footnotesize équivalents\allowbreak{}}\par
\noindent Source : 7II:647. \nolinkurl{fixes/MISPRINTS_7II.json#647}\par

\par\medskip\noindent\hypertarget{app-F7II-XX-035.2}{}\textbf{F7II-XX-035.2}\quad p. 361; designated-witness attribution.\par
\noindent Imprimé : {\ttfamily\footnotesize \$\allowbreak{}}\par
\noindent Corrigé : \emph{(vide)}\par
\noindent Source : WITNESS-F7II-XX-035.2. \nolinkurl{fr/sga7ii/diplomatic/expose_XX_body.tex}; \nolinkurl{fr/sga7ii/current/expose_XX_body.tex}\par

\par\medskip\noindent\hypertarget{app-F7II-XX-035.3}{}\textbf{F7II-XX-035.3}\quad p. 361; designated-witness attribution.\par
\noindent Imprimé : \emph{(vide)}\par
\noindent Corrigé : {\ttfamily\footnotesize \textbackslash{}operatorname\{\allowbreak{}}\par
\noindent Source : WITNESS-F7II-XX-035.3. \nolinkurl{fr/sga7ii/diplomatic/expose_XX_body.tex}; \nolinkurl{fr/sga7ii/current/expose_XX_body.tex}\par

\par\medskip\noindent\hypertarget{app-F7II-XX-035.4}{}\textbf{F7II-XX-035.4}\quad p. 361; designated-witness attribution.\par
\noindent Imprimé : {\ttfamily\footnotesize  \$\allowbreak{}}\par
\noindent Corrigé : {\ttfamily\footnotesize \}\allowbreak{}}\par
\noindent Source : WITNESS-F7II-XX-035.4. \nolinkurl{fr/sga7ii/diplomatic/expose_XX_body.tex}; \nolinkurl{fr/sga7ii/current/expose_XX_body.tex}\par

\par\medskip\noindent\hypertarget{app-F7II-XX-035.5}{}\textbf{F7II-XX-035.5}\quad p. 361; designated-witness attribution.\par
\noindent Imprimé : \emph{(vide)}\par
\noindent Corrigé : {\ttfamily\footnotesize ,\allowbreak{}}\par
\noindent Source : WITNESS-F7II-XX-035.5. \nolinkurl{fr/sga7ii/diplomatic/expose_XX_body.tex}; \nolinkurl{fr/sga7ii/current/expose_XX_body.tex}\par

\par\medskip\noindent\hypertarget{app-F7II-XX-036.1}{}\textbf{F7II-XX-036.1}\quad p. 361; inherited corrected-layer difference.\par
\noindent Imprimé : \emph{(2 espace(s), 0 retour(s) de ligne dans le témoin TeX)}\par
\noindent Corrigé : \emph{(vide)}\par
\noindent Source : ADOPT-F7II-XX-036.1. \nolinkurl{fr/sga7ii/diplomatic/expose_XX_body.tex}; \nolinkurl{fr/sga7ii/current/expose_XX_body.tex}\par

\par\medskip\noindent\hypertarget{app-F7II-XX-036.2}{}\textbf{F7II-XX-036.2}\quad p. 361; inherited corrected-layer difference.\par
\noindent Imprimé : \emph{(2 espace(s), 0 retour(s) de ligne dans le témoin TeX)}\par
\noindent Corrigé : {\ttfamily\footnotesize  \$\textbackslash{}operatorname\{\allowbreak{}}\par
\noindent Source : ADOPT-F7II-XX-036.2. \nolinkurl{fr/sga7ii/diplomatic/expose_XX_body.tex}; \nolinkurl{fr/sga7ii/current/expose_XX_body.tex}\par

\par\medskip\noindent\hypertarget{app-F7II-XX-036.3}{}\textbf{F7II-XX-036.3}\quad p. 361; inherited corrected-layer difference.\par
\noindent Imprimé : {\ttfamily\footnotesize .\$\allowbreak{}}\par
\noindent Corrigé : {\ttfamily\footnotesize \}(\allowbreak{}}\par
\noindent Source : ADOPT-F7II-XX-036.3. \nolinkurl{fr/sga7ii/diplomatic/expose_XX_body.tex}; \nolinkurl{fr/sga7ii/current/expose_XX_body.tex}\par

\par\medskip\noindent\hypertarget{app-F7II-XX-036.4}{}\textbf{F7II-XX-036.4}\quad p. 361; inherited corrected-layer difference.\par
\noindent Imprimé : \emph{(vide)}\par
\noindent Corrigé : {\ttfamily\footnotesize )\allowbreak{}}\par
\noindent Source : ADOPT-F7II-XX-036.4. \nolinkurl{fr/sga7ii/diplomatic/expose_XX_body.tex}; \nolinkurl{fr/sga7ii/current/expose_XX_body.tex}\par

\par\medskip\noindent\hypertarget{app-F7II-XX-036.5}{}\textbf{F7II-XX-036.5}\quad p. 361; inherited corrected-layer difference.\par
\noindent Imprimé : \emph{(1 espace(s), 0 retour(s) de ligne dans le témoin TeX)}\par
\noindent Corrigé : \emph{(vide)}\par
\noindent Source : ADOPT-F7II-XX-036.5. \nolinkurl{fr/sga7ii/diplomatic/expose_XX_body.tex}; \nolinkurl{fr/sga7ii/current/expose_XX_body.tex}\par

\par\medskip\noindent\hypertarget{app-F7II-XX-036.6}{}\textbf{F7II-XX-036.6}\quad p. 361; inherited corrected-layer difference.\par
\noindent Imprimé : \emph{(1 espace(s), 0 retour(s) de ligne dans le témoin TeX)}\par
\noindent Corrigé : \emph{(vide)}\par
\noindent Source : ADOPT-F7II-XX-036.6. \nolinkurl{fr/sga7ii/diplomatic/expose_XX_body.tex}; \nolinkurl{fr/sga7ii/current/expose_XX_body.tex}\par

\par\medskip\noindent\hypertarget{app-F7II-XX-037.1}{}\textbf{F7II-XX-037.1}\quad p. 361; designated-witness attribution.\par
\noindent Imprimé : {\ttfamily\footnotesize ,\allowbreak{}}\par
\noindent Corrigé : \emph{(vide)}\par
\noindent Source : WITNESS-F7II-XX-037.1. \nolinkurl{fr/sga7ii/diplomatic/expose_XX_body.tex}; \nolinkurl{fr/sga7ii/current/expose_XX_body.tex}\par

\par\medskip\noindent\hypertarget{app-F7II-XX-038.1}{}\textbf{F7II-XX-038.1}\quad p. 361; inherited corrected-layer difference.\par
\noindent Imprimé : {\ttfamily\footnotesize *\allowbreak{}}\par
\noindent Corrigé : \emph{(vide)}\par
\noindent Source : ADOPT-F7II-XX-038.1. \nolinkurl{fr/sga7ii/diplomatic/expose_XX_body.tex}; \nolinkurl{fr/sga7ii/current/expose_XX_body.tex}\par

\par\medskip\noindent\hypertarget{app-F7II-XX-038.2}{}\textbf{F7II-XX-038.2}\quad p. 361; inherited corrected-layer difference.\par
\noindent Imprimé : {\ttfamily\footnotesize ,\allowbreak{}}\par
\noindent Corrigé : \emph{(vide)}\par
\noindent Source : ADOPT-F7II-XX-038.2. \nolinkurl{fr/sga7ii/diplomatic/expose_XX_body.tex}; \nolinkurl{fr/sga7ii/current/expose_XX_body.tex}\par

\par\medskip\noindent\hypertarget{app-F7II-XX-038.3}{}\textbf{F7II-XX-038.3}\quad p. 361; inherited corrected-layer difference.\par
\noindent Imprimé : {\ttfamily\footnotesize n\allowbreak{}}\par
\noindent Corrigé : {\ttfamily\footnotesize m\_\textbackslash{}ell\allowbreak{}}\par
\noindent Source : ADOPT-F7II-XX-038.3. \nolinkurl{fr/sga7ii/diplomatic/expose_XX_body.tex}; \nolinkurl{fr/sga7ii/current/expose_XX_body.tex}\par

\par\medskip\noindent\hypertarget{app-F7II-XX-038.4}{}\textbf{F7II-XX-038.4}\quad p. 361; designated-witness attribution.\par
\noindent Imprimé : \emph{(vide)}\par
\noindent Corrigé : {\ttfamily\footnotesize \{\textbackslash{}mathbb\allowbreak{} }\par
\noindent Source : WITNESS-F7II-XX-038.4. \nolinkurl{fr/sga7ii/diplomatic/expose_XX_body.tex}; \nolinkurl{fr/sga7ii/current/expose_XX_body.tex}\par

\par\medskip\noindent\hypertarget{app-F7II-XX-038.5}{}\textbf{F7II-XX-038.5}\quad p. 361; inherited corrected-layer difference.\par
\noindent Imprimé : \emph{(vide)}\par
\noindent Corrigé : {\ttfamily\footnotesize \}\allowbreak{}}\par
\noindent Source : ADOPT-F7II-XX-038.5. \nolinkurl{fr/sga7ii/diplomatic/expose_XX_body.tex}; \nolinkurl{fr/sga7ii/current/expose_XX_body.tex}\par

\par\medskip\noindent\hypertarget{app-F7II-XX-038.6}{}\textbf{F7II-XX-038.6}\quad p. 361; printed misprint.\par
\noindent Imprimé : {\ttfamily\footnotesize 2n\allowbreak{}}\par
\noindent Corrigé : {\ttfamily\footnotesize 2m\allowbreak{}}\par
\noindent Source : 7II:1050. \nolinkurl{fixes/MISPRINTS_7II.json#1050}\par

\par\medskip\noindent\hypertarget{app-F7II-XX-038.7}{}\textbf{F7II-XX-038.7}\quad p. 361; designated-witness attribution.\par
\noindent Imprimé : {\ttfamily\footnotesize n\allowbreak{}}\par
\noindent Corrigé : {\ttfamily\footnotesize m\allowbreak{}}\par
\noindent Source : WITNESS-F7II-XX-038.7. \nolinkurl{fr/sga7ii/diplomatic/expose_XX_body.tex}; \nolinkurl{fr/sga7ii/current/expose_XX_body.tex}\par

\par\medskip\noindent\hypertarget{app-F7II-XX-038.8}{}\textbf{F7II-XX-038.8}\quad p. 361; designated-witness attribution.\par
\noindent Imprimé : \emph{(vide)}\par
\noindent Corrigé : {\ttfamily\footnotesize (\allowbreak{}}\par
\noindent Source : WITNESS-F7II-XX-038.8. \nolinkurl{fr/sga7ii/diplomatic/expose_XX_body.tex}; \nolinkurl{fr/sga7ii/current/expose_XX_body.tex}\par

\par\medskip\noindent\hypertarget{app-F7II-XX-038.9}{}\textbf{F7II-XX-038.9}\quad p. 361; printed misprint.\par
\noindent Imprimé : \emph{(1 espace(s), 0 retour(s) de ligne dans le témoin TeX)}\par
\noindent Corrigé : {\ttfamily\footnotesize .\allowbreak{}}\par
\noindent Source : 7II:1049. \nolinkurl{fixes/MISPRINTS_7II.json#1049}\par

\par\medskip\noindent\hypertarget{app-F7II-XX-038.10}{}\textbf{F7II-XX-038.10}\quad p. 361; designated-witness attribution.\par
\noindent Imprimé : \emph{(vide)}\par
\noindent Corrigé : {\ttfamily\footnotesize )\allowbreak{}}\par
\noindent Source : WITNESS-F7II-XX-038.10. \nolinkurl{fr/sga7ii/diplomatic/expose_XX_body.tex}; \nolinkurl{fr/sga7ii/current/expose_XX_body.tex}\par

\par\medskip\noindent\hypertarget{app-F7II-XX-038.11}{}\textbf{F7II-XX-038.11}\quad p. 361; designated-witness attribution.\par
\noindent Imprimé : {\ttfamily\footnotesize n\allowbreak{}}\par
\noindent Corrigé : {\ttfamily\footnotesize m\allowbreak{}}\par
\noindent Source : WITNESS-F7II-XX-038.11. \nolinkurl{fr/sga7ii/diplomatic/expose_XX_body.tex}; \nolinkurl{fr/sga7ii/current/expose_XX_body.tex}\par

\par\medskip\noindent\hypertarget{app-F7II-XX-038.12}{}\textbf{F7II-XX-038.12}\quad p. 361; designated-witness attribution.\par
\noindent Imprimé : {\ttfamily\footnotesize n\allowbreak{}}\par
\noindent Corrigé : {\ttfamily\footnotesize m\allowbreak{}}\par
\noindent Source : WITNESS-F7II-XX-038.12. \nolinkurl{fr/sga7ii/diplomatic/expose_XX_body.tex}; \nolinkurl{fr/sga7ii/current/expose_XX_body.tex}\par

\par\medskip\noindent\hypertarget{app-F7II-XX-038.13}{}\textbf{F7II-XX-038.13}\quad p. 361; inherited corrected-layer difference.\par
\noindent Imprimé : {\ttfamily\footnotesize n\allowbreak{}}\par
\noindent Corrigé : {\ttfamily\footnotesize m\_\textbackslash{}ell\allowbreak{}}\par
\noindent Source : ADOPT-F7II-XX-038.13. \nolinkurl{fr/sga7ii/diplomatic/expose_XX_body.tex}; \nolinkurl{fr/sga7ii/current/expose_XX_body.tex}\par

\par\medskip\noindent\hypertarget{app-F7II-XX-039.1}{}\textbf{F7II-XX-039.1}\quad p. 362; designated-witness attribution.\par
\noindent Imprimé : {\ttfamily\footnotesize n\allowbreak{}}\par
\noindent Corrigé : {\ttfamily\footnotesize m\allowbreak{}}\par
\noindent Source : WITNESS-F7II-XX-039.1. \nolinkurl{fr/sga7ii/diplomatic/expose_XX_body.tex}; \nolinkurl{fr/sga7ii/current/expose_XX_body.tex}\par

\par\medskip\noindent\hypertarget{app-F7II-XX-039.2}{}\textbf{F7II-XX-039.2}\quad p. 362; designated-witness attribution.\par
\noindent Imprimé : \emph{(1 espace(s), 0 retour(s) de ligne dans le témoin TeX)}\par
\noindent Corrigé : \emph{(vide)}\par
\noindent Source : WITNESS-F7II-XX-039.2. \nolinkurl{fr/sga7ii/diplomatic/expose_XX_body.tex}; \nolinkurl{fr/sga7ii/current/expose_XX_body.tex}\par

\par\medskip\noindent\hypertarget{app-F7II-XX-039.3}{}\textbf{F7II-XX-039.3}\quad p. 362; designated-witness attribution.\par
\noindent Imprimé : {\ttfamily\footnotesize 2n\allowbreak{}}\par
\noindent Corrigé : {\ttfamily\footnotesize 2m\allowbreak{}}\par
\noindent Source : WITNESS-F7II-XX-039.3. \nolinkurl{fr/sga7ii/diplomatic/expose_XX_body.tex}; \nolinkurl{fr/sga7ii/current/expose_XX_body.tex}\par

\par\medskip\noindent\hypertarget{app-F7II-XX-039.4}{}\textbf{F7II-XX-039.4}\quad p. 362; designated-witness attribution.\par
\noindent Imprimé : {\ttfamily\footnotesize n\allowbreak{}}\par
\noindent Corrigé : {\ttfamily\footnotesize m\allowbreak{}}\par
\noindent Source : WITNESS-F7II-XX-039.4. \nolinkurl{fr/sga7ii/diplomatic/expose_XX_body.tex}; \nolinkurl{fr/sga7ii/current/expose_XX_body.tex}\par

\par\medskip\noindent\hypertarget{app-F7II-XX-039.5}{}\textbf{F7II-XX-039.5}\quad p. 362; designated-witness attribution.\par
\noindent Imprimé : {\ttfamily\footnotesize n\allowbreak{}}\par
\noindent Corrigé : {\ttfamily\footnotesize m\_\textbackslash{}ell\allowbreak{}}\par
\noindent Source : WITNESS-F7II-XX-039.5. \nolinkurl{fr/sga7ii/diplomatic/expose_XX_body.tex}; \nolinkurl{fr/sga7ii/current/expose_XX_body.tex}\par

\par\medskip\noindent\hypertarget{app-F7II-XX-039.6}{}\textbf{F7II-XX-039.6}\quad p. 362; designated-witness attribution.\par
\noindent Imprimé : \emph{(1 espace(s), 0 retour(s) de ligne dans le témoin TeX)}\par
\noindent Corrigé : \emph{(vide)}\par
\noindent Source : WITNESS-F7II-XX-039.6. \nolinkurl{fr/sga7ii/diplomatic/expose_XX_body.tex}; \nolinkurl{fr/sga7ii/current/expose_XX_body.tex}\par

\par\medskip\noindent\hypertarget{app-F7II-XX-039.7}{}\textbf{F7II-XX-039.7}\quad p. 362; printed misprint.\par
\noindent Imprimé : \emph{(vide)}\par
\noindent Corrigé : {\ttfamily\footnotesize .\allowbreak{}}\par
\noindent Source : 7II:650. \nolinkurl{fixes/MISPRINTS_7II.json#650}\par

\par\medskip\noindent\hypertarget{app-F7II-XX-040.1}{}\textbf{F7II-XX-040.1}\quad p. 362; designated-witness attribution.\par
\noindent Imprimé : {\ttfamily\footnotesize \textbackslash{}mathrm\{dim\}\allowbreak{}}\par
\noindent Corrigé : {\ttfamily\footnotesize \textbackslash{}dim\allowbreak{}}\par
\noindent Source : WITNESS-F7II-XX-040.1. \nolinkurl{fr/sga7ii/diplomatic/expose_XX_body.tex}; \nolinkurl{fr/sga7ii/current/expose_XX_body.tex}\par

\par\medskip\noindent\hypertarget{app-F7II-XX-040.2}{}\textbf{F7II-XX-040.2}\quad p. 362; printed misprint.\par
\noindent Imprimé : {\ttfamily\footnotesize n\allowbreak{}}\par
\noindent Corrigé : {\ttfamily\footnotesize m)\allowbreak{}}\par
\noindent Source : 7II:648. \nolinkurl{fixes/MISPRINTS_7II.json#648}\par

\par\medskip\noindent\hypertarget{app-F7II-XX-040.3}{}\textbf{F7II-XX-040.3}\quad p. 362; designated-witness attribution.\par
\noindent Imprimé : {\ttfamily\footnotesize 4.\allowbreak{}}\par
\noindent Corrigé : {\ttfamily\footnotesize \$m=\allowbreak{}}\par
\noindent Source : WITNESS-F7II-XX-040.3. \nolinkurl{fr/sga7ii/diplomatic/expose_XX_body.tex}; \nolinkurl{fr/sga7ii/current/expose_XX_body.tex}\par

\par\medskip\noindent\hypertarget{app-F7II-XX-040.4}{}\textbf{F7II-XX-040.4}\quad p. 362; designated-witness attribution.\par
\noindent Imprimé : \emph{(vide)}\par
\noindent Corrigé : {\ttfamily\footnotesize \$\allowbreak{}}\par
\noindent Source : WITNESS-F7II-XX-040.4. \nolinkurl{fr/sga7ii/diplomatic/expose_XX_body.tex}; \nolinkurl{fr/sga7ii/current/expose_XX_body.tex}\par

\par\medskip\noindent\hypertarget{app-F7II-XX-040.5}{}\textbf{F7II-XX-040.5}\quad p. 362; inherited corrected-layer difference.\par
\noindent Imprimé : \emph{(1 espace(s), 0 retour(s) de ligne dans le témoin TeX)}\par
\noindent Corrigé : \emph{(vide)}\par
\noindent Source : ADOPT-F7II-XX-040.5. \nolinkurl{fr/sga7ii/diplomatic/expose_XX_body.tex}; \nolinkurl{fr/sga7ii/current/expose_XX_body.tex}\par

\par\medskip\noindent\hypertarget{app-F7II-XX-040.6}{}\textbf{F7II-XX-040.6}\quad p. 362; font.\par
\noindent Imprimé : \emph{(vide)}\par
\noindent Corrigé : {\ttfamily\footnotesize \textbackslash{}mathbb\{\allowbreak{}}\par
\noindent Source : FIX-7II-D-0752. \nolinkurl{fixes/FIXLOG_7II.jsonl#752}\par

\par\medskip\noindent\hypertarget{app-F7II-XX-040.7}{}\textbf{F7II-XX-040.7}\quad p. 362; font.\par
\noindent Imprimé : \emph{(vide)}\par
\noindent Corrigé : {\ttfamily\footnotesize \}\allowbreak{}}\par
\noindent Source : FIX-7II-D-0752. \nolinkurl{fixes/FIXLOG_7II.jsonl#752}\par

\par\medskip\noindent\hypertarget{app-F7II-XX-040.8}{}\textbf{F7II-XX-040.8}\quad p. 362; printed misprint.\par
\noindent Imprimé : {\ttfamily\footnotesize cohomologie\allowbreak{}}\par
\noindent Corrigé : {\ttfamily\footnotesize cohomologue\allowbreak{}}\par
\noindent Source : 7II:649; 7II:1056. \nolinkurl{fixes/MISPRINTS_7II.json#649}; \nolinkurl{fixes/MISPRINTS_7II.json#1056}\par

\par\medskip\noindent\hypertarget{app-F7II-XX-040.9}{}\textbf{F7II-XX-040.9}\quad p. 362; designated-witness attribution.\par
\noindent Imprimé : \emph{(vide)}\par
\noindent Corrigé : {\ttfamily\footnotesize non\allowbreak{} nul\allowbreak{} }\par
\noindent Source : WITNESS-F7II-XX-040.9. \nolinkurl{fr/sga7ii/diplomatic/expose_XX_body.tex}; \nolinkurl{fr/sga7ii/current/expose_XX_body.tex}\par

\normalsize\fi
\clearpage\SGAtoc{Exposé XXI — LE NIVEAU DE LA COHOMOLOGIE DES INTERSECTIONS COMPLETES}
\selectlanguage{french}

\SGApage{II}{363}


S G A  7


\underline{E X P O S E  XXI}


\underline{\underline{LE NIVEAU DE LA COHOMOLOGIE DES INTERSECTIONS COMPLETES}}

\underline{par N. Katz}

avec un Appendice \quad par P. Deligne ($\S$ 5)


\underline{SOMMAIRE}

0.  Introduction. \hfill 1

1.  La matrice de Hasse-Witt d'une intersection complète. \hfill 2

2.  Le coniveau dans le cas d'un corps de définition fini. \hfill 14

3.  Application aux intersections complètes. \hfill 16

4.  Les conclusions. \hfill 19

5.  Théorèmes d'intégralité. \hfill 21


\SGApage{II}{364}


\SGAsectiontoc{0. Introduction}{0}
0. \underline{Introduction}

\SGAvar{Soient $X$ une variété projective et lisse de dimension $n \geq 3$, $Y(d)$ une section par une hypersurface de degré $d$, et $E^{n-1}(Y(d),\mathbb{Q}_\ell)$ la cohomologie ``évanescente'' de $Y(d)$ (XVIII 5.4.4). Le but du présent exposé et de démontrer que pour $d$ suffisamment grand, et $Y(d)$ suffisament générale, il n'est pas vrai que tout $E^{n-1}(Y(d),\mathbb{Q}_\ell)$ ``provienne'' du $H^i(i < n-1)$ de variétés algébriques propres et lisses, par des correspondances algébriques. En caractéristique nulle, c'est une conséquence triviale de la bigraduation de Hodge sur la cohomologie complexe, et de ce que, pour la cohomologie cohérente, on a

\begin{equation*}
\dim H^{n-1}(Y(d),\ \mathcal{O}_{Y(d)}) \longrightarrow \infty \ \mbox{pour} \ d \longrightarrow \infty \quad .
\end{equation*}
La méthode employée en caractéristique $p$ est encore de relier la cohomologie cohérente $H^{n-1}(Y(d),\ \mathcal{O}_{Y(d)})$ (un vectoriel en car. $p$) à la cohomologie étale $H^{n-1}_{\mathrm{et}}(Y(d),\mathbb{Q}_\ell)$. Pour ceci, on se place sur un corps
\SGApage{II}{365}
de base fini, où les deux sont liés par la fonction $\zeta$ de la variété $Y(d)$, fonction qui s'exprime en termes de la cohomologie $\ell$-adique, et dont la ``réduction modulo $p$'' s'exprime en termes de la cohomologie cohérente (cf. XXII 3.9). Cette méthode marche sans problème pour les intersections complètes, grâce à la connaissance détaillée des deux espèces de cohomologie (étale et cohérente) dans ce cas, et donne le résultat principal 4.2 comme conséquence d'un résultat (1.4) sur les opérations de Hasse-Witt (1.0).

Il existe des arguments heuristiques dûs à A. GROTHENDIECK et s'appuyant sur le yoga de la cohomologie cristalline, qui rendent plausible l'énoncé général pour toute $X$ projective et lisse, par essentiellement la même méthode.

Même dans le cas particulier des intersections complètes, la démonstration utilise de façon essentielle (2.2) le théorème d'intégralité des valeurs propres de Frobenius 2.1, démontré récemment par P. DELIGNE, et qui est prouvé dans l'Appendice I.


}{Soient $X$ une variété projective et lisse de dimension $n \geq 3$, $Y(d)$ une section par une hypersurface de degré $d$, et $E^{n-1}(Y(d),\mathbb{Q}_\ell)$ la cohomologie ``évanescente'' de $Y(d)$ (XVIII 5.4.4). Le but du présent exposé est de démontrer que pour $d$ suffisamment grand et $Y(d)$ suffisamment générale, il n'est pas vrai que toute la cohomologie $E^{n-1}(Y(d),\mathbb{Q}_\ell)$ ``provienne'' des $H^i$ ($i < n-1$) de variétés algébriques propres et lisses, par des correspondances algébriques. En caractéristique nulle, c'est une conséquence triviale de la bigraduation de Hodge sur la cohomologie complexe, et de ce que, pour la cohomologie cohérente, on a

\begin{equation*}
\dim H^{n-1}(Y(d),\ \mathcal{O}_{Y(d)}) \longrightarrow \infty \ \mbox{pour} \ d \longrightarrow \infty \quad .
\end{equation*}
La méthode employée en caractéristique $p$ est encore de relier la cohomologie cohérente $H^{n-1}(Y(d),\ \mathcal{O}_{Y(d)})$ (un espace vectoriel en caractéristique $p$) à la cohomologie étale $H^{n-1}_{\mathrm{et}}(Y(d),\mathbb{Q}_\ell)$. Pour ceci, on se place sur un corps
\SGApage{II}{365}
de base fini, où les deux sont liés par la fonction $\zeta$ de la variété $Y(d)$, fonction qui s'exprime en termes de la cohomologie $\ell$-adique, et dont la ``réduction modulo $p$'' s'exprime en termes de la cohomologie cohérente (cf. XXII 3.9). Cette méthode marche sans problème pour les intersections complètes, grâce à la connaissance détaillée des deux espèces de cohomologie (étale et cohérente) dans ce cas, et donne le résultat principal 4.2 comme conséquence d'un résultat (1.4) sur les opérations de Hasse-Witt (1.0).

Il existe des arguments heuristiques dus à A. GROTHENDIECK et s'appuyant sur le yoga de la cohomologie cristalline, qui rendent plausible l'énoncé général pour toute $X$ projective et lisse, par essentiellement la même méthode.

Même dans le cas particulier des intersections complètes, la démonstration utilise de façon essentielle (2.2) le théorème d'intégralité des valeurs propres de Frobenius 2.1, démontré récemment par P. DELIGNE et prouvé ici dans l'Appendice ($\S$ 5).


}{F7II-XXI-001}\SGAsectiontoc{1. La matrice de Hasse-Witt d'une intersection complète}{1}
1. \underline{La matrice de Hasse-Witt d'une intersection complète}

\SGAvar{\underline{Définition} 1.0. \underline{Un schéma propre sur un corps} $k$ \underline{d'exposant caractéristique} $p$ \underline{s'appelle} spécial \underline{si l'opération additive} $p$-\underline{linéaire}

}{\underline{Définition} 1.0. \underline{Un schéma propre} $X$ \underline{sur un corps} $k$ \underline{d'exposant caractéristique} $p$ \underline{s'appelle} spécial \underline{si l'opération additive} $p$-\underline{linéaire}

}{F7II-XXI-002}\begin{equation}\tag{1.0.1}
\mathfrak{F}_p : \mathcal{O}_X \longrightarrow \mathcal{O}_X \quad , \quad g \longmapsto g^p
\end{equation}
\underline{induit une opération}

\begin{equation}\tag{1.0.2}
\mathfrak{F}_p : H^{\dim X}(X,\mathcal{O}_X) \longrightarrow H^{\dim X}(X,\mathcal{O}_X) 
\end{equation}

\SGApage{II}{366}

 \underline{qui est nilpotente} (\underline{ou, ce qui revient au même, posant} $M = \dim_k(H^{\dim X}(X,\mathcal{O}_X))$, \underline{si l'opération} $p^M$-\underline{linéaire}

\SGAvar{\begin{equation}\tag{1.0.3}
\mathfrak{F}_p^M = \mathfrak{F}_{p^M} : \mathcal{O}_X \longrightarrow \mathcal{O}_X \quad , \quad g \longmapsto g^{p^M}
\end{equation}
\underline{tue} $H^{\dim X}(X,\mathcal{O}_X)$.

}{\begin{equation}\tag{1.0.3}
\mathfrak{F}_p^M = \mathfrak{F}_{p^M} : \mathcal{O}_X \longrightarrow \mathcal{O}_X \quad , \quad g \longmapsto g^{p^M}
\end{equation}
\underline{tue} $H^{\dim X}(X,\mathcal{O}_X)$).

}{F7II-XXI-003}Rappelons [2] que pour tout $i \in \mathbb{Z}$ , l'homomorphisme

\begin{equation*}
H^i(X,\mathcal{O}_X) \stackrel{\mathfrak{F}_p}{\longrightarrow} H^i(X,\mathcal{O}_X)
\end{equation*}
induit par (1.0.1) s'appelle classiquement \underline{l'homomorphisme de Hasse-Witt} dans $H^i(X,\mathcal{O}_X)$. On notera que si $p = 1$, dire que $X$ est non-spécial signifie simplement que $H^{\dim X}(X,\mathcal{O}_X) \neq 0$.

\underline{Théorème} 1.1. \underline{Fixons} \underline{un exposant caractéristique} $p$, \underline{une dimension} $n \geq 1$, \underline{un entier} $r \geq 1$ \underline{et un multidegré}$(d_1,\dots,d_r)$. \underline{Si on a}

\SGAvar{\begin{equation}\tag{1.1.0}
\sum_{i=1}^{r} d_i \geq (n+r)+1
\end{equation}
\underline{alors, parmi toute les intersections complètes} (\underline{de dimension} $n$) \underline{dans} $\mathbb{P}^{n+r}$ \underline{de multidegré} $(d_1,\dots,d_r)$ \underline{en caractéristique} $p$, \underline{un élément} \underline{``assez général'' est non-spécial} (1.0).

}{\begin{equation}\tag{1.1.0}
\sum_{i=1}^{r} d_i \geq (n+r)+1
\end{equation}
\underline{alors, parmi toutes les intersections complètes} (\underline{de dimension} $n$) \underline{dans} $\mathbb{P}^{n+r}$ \underline{de multidegré} $(d_1,\dots,d_r)$ \underline{en caractéristique} $p$, \underline{un élément} \underline{``assez général'' est non-spécial} (1.0).

}{F7II-XXI-004}\underline{Démonstration.} Désignons par

\SGAvar{\begin{equation}\tag{1.1.1}
f : X_U \longrightarrow U
\end{equation}
la``famille universelle'' des intersections complètes dans $\mathbb{P}^{n+r}$ de
\SGApage{II}{367}
multidegré $(d_1,\dots,d_r)$ en caractéristique $p$ (muni des équations de définition), de sorte que $U$ est un ouvert dans un produit de $r$ espaces projectifs sur $\mathbb{F}_p$ , d'une dimension impressionnante.

}{\begin{equation}\tag{1.1.1}
f : X_U \longrightarrow U
\end{equation}
la ``famille universelle'' des intersections complètes dans $\mathbb{P}^{n+r}$ de
\SGApage{II}{367}
multidegré $(d_1,\dots,d_r)$ en caractéristique $p$ (muni des équations de définition), de sorte que $U$ est un ouvert dans un produit de $r$ espaces projectifs sur $\mathbb{F}_p$ , d'une dimension impressionnante.

}{F7II-XXI-005}On sait que le $\mathcal{O}_U$-module

\SGAvar{\begin{equation}\tag{1.1.2}
R^n f_* \mathcal{O}_{X_U}
\end{equation}
est localement libre, et sa formation commute à tout changement de base (XI 1.5 (i)) et, par l'hypothèse (1.1.0), que son rang $M$ est non nul (XI 2.5). Cela donne le théroème si $p = 1$, pour toute intersection complète vérifiante (1.1.0).

}{\begin{equation}\tag{1.1.2}
R^n f_* \mathcal{O}_{X_U}
\end{equation}
est localement libre, et sa formation commute à tout changement de base (XI 1.5 (i)), et que, par l'hypothèse (1.1.0), son rang $M$ est non nul (XI 2.5). Cela donne le théorème si $p = 1$, pour toute intersection complète vérifiant (1.1.0).

}{F7II-XXI-006}Supposons $p \geq 2$. L'opération de Hasse-Witt itérée

\begin{equation}\tag{1.1.3}
\mathfrak{F}_p^M = \mathfrak{F}_{p^M} : R^n f_* \mathcal{O}_{X_U} \longrightarrow R^n f_* \mathcal{O}_{X_U}
\end{equation}
étant $p^M$-linéaire sur $\mathcal{O}_U$ , et $R^n f_* \mathcal{O}_{X_U}$ étant localement libre, l'action (1.1.3) s'exprime localement sur $U$ par une matrice $H$ de type $(M,M)$ à valeurs dans $\mathcal{O}_U$ , de sorte que

\begin{equation}\tag{1.1.4}
\{ u \in U \mid X_u \mbox{ est non-spéciale}\} \mbox{ est ouvert dans } U
\end{equation}
(car c'est l'ensemble des $u$ tels que $H(u) \neq 0$), et, $U$ étant irréductible, il suffit de trouver une valeur de $u$ telle que $X_u$ soit non spéciale, ce qu'on va faire sur un corps fini.

\SGAvar{\underline{Proposition} 1.1.5. \underline{Soit} $X$ \underline{une intersection complète sur} $\mathbb{F}_q$. \underline{Si} $X$ \underline{n'a pas des points} $\mathbb{F}_q$\underline{-rationnels,} \underline{alors} $X$ \underline{est non spéciale.}


}{\underline{Proposition} 1.1.5. \underline{Soit} $X$ \underline{une intersection complète sur} $\mathbb{F}_q$. \underline{Si} $X$ \underline{n'a pas de points} $\mathbb{F}_q$\underline{-rationnels,} \underline{alors} $X$ \underline{est non spéciale.}


}{F7II-XXI-007}\SGApage{II}{368}

\underline{Démonstration.} En effet, il suffit évidemment de prouver que pour $X$ une intersection complète, on a l'implication

\begin{equation}\tag{1.1.5.1}
X \text{ spéciale} \Longrightarrow \mathrm{card}(X(\mathbb{F}_q)) \equiv 1 \bmod p \quad .
\end{equation}
Or, d'après la formule de la trace (prouvée dans l'exposé suivant XXII, 3.2), on a

\begin{equation*}
\mathrm{card}(X(\mathbb{F}_q)) \equiv \sum_{i \geq 0} (-1)^i \, \mathrm{tr}(\mathfrak{F}_q \mid H^i(X,\mathcal{O}_X)) \quad .
\end{equation*}
Compte tenu de ce que $H^i(X,\mathcal{O}_X)=0$ si $i \neq 0, n=\dim(X)$ et $H^0(X,\mathcal{O}_X) \simeq \mathbb{F}_q$ (XI, 1.5 (iii)), ceci se récrit

\begin{equation*}
\mathrm{card}(X(\mathbb{F}_q)) \equiv 1 + (-1)^n \, \mathrm{tr}(\mathfrak{F}_q \mid H^n(X,\mathcal{O}_X)) \quad .
\end{equation*}
Si $X$ est spéciale, alors par définition l'opération de $\mathfrak{F}_q$ sur $H^n(X,\mathcal{O}_X)$ est nilpotente, donc a une trace nulle, d'où l'implication (1.1.5.1) et la proposition.

\underline{Proposition} 1.2. \underline{Quelle que soit la puissance} $q$ \underline{de} $p$, \underline{si}

\SGAvar{\begin{equation*}
\sum_1^r d_i \geq n+r+1 \quad ,
\end{equation*}
\underline{il existe une intersection complète dans} $\mathbb{P}^{n+r}_{\mathbb{F}_q}$ \underline{de multidegré} $(d_1,\dots,d_r)$ \underline{qui est sans point} $\mathbb{F}_q$\underline{-rationnel}.

}{\begin{equation*}
\sum_{i=1}^r d_i \geq n+r+1 \quad ,
\end{equation*}
\underline{il existe une intersection complète dans} $\mathbb{P}^{n+r}_{\mathbb{F}_q}$ \underline{de multidegré} $(d_1,\dots,d_r)$ \underline{qui est sans point} $\mathbb{F}_q$\underline{-rationnel}.

}{F7II-XXI-008}\underline{Démonstration.}

Commençons par le cas $r=1$. Pour tout $d \geq n+2$, il faut trouver une forme homogène de degré $d$ 


\SGApage{II}{369}

\SGAvar{\begin{equation}\tag{1.2.1}
 H(X_1),\dots,X_{n+2}) \in \mathbb{F}_q[X_1,\dots,X_{n+2}]
\end{equation}
telle que

\begin{equation}\tag{1.2.2}
(\alpha_1,\dots,\alpha_{n+2}) \in (\mathbb{F}_q)^{n+2}, \; H\,(\alpha_1,\dots,\alpha_{n+2}) = 0 \Longrightarrow \alpha_1 = \dots = \alpha_{n+2} = 0 \quad .
\end{equation}
Pour ceci, posons $\mathbb{E} = \mathbb{F}_{q^d}$ , et soient $\beta_1,\dots,\beta_{n+2}$ des éléments de $\mathbb{E}$ qui sont linéairement indépendant sur $\mathbb{F}_q$ . On prend

}{\begin{equation}\tag{1.2.1}
 H(X_1,\dots,X_{n+2}) \in \mathbb{F}_q[X_1,\dots,X_{n+2}]
\end{equation}
telle que

\begin{equation}\tag{1.2.2}
(\alpha_1,\dots,\alpha_{n+2}) \in (\mathbb{F}_q)^{n+2}, \; H\,(\alpha_1,\dots,\alpha_{n+2}) = 0 \Longrightarrow \alpha_1 = \dots = \alpha_{n+2} = 0 \quad .
\end{equation}
Pour ceci, posons $\mathbb{E} = \mathbb{F}_{q^d}$ , et soient $\beta_1,\dots,\beta_{n+2}$ des éléments de $\mathbb{E}$ qui sont linéairement indépendants sur $\mathbb{F}_q$ . On prend

}{F7II-XXI-009}\begin{equation}\tag{1.2.3}
H(X_1,\dots,X_{n+2}) \;=\; N_{\mathbb{E}/\mathbb{F}_q} \Bigl( \sum_{i=1}^{n+2} \beta_i X_i \Bigr) \quad .
\end{equation}
Dans le cas général, compte tenu de ce que

\SGAvar{\begin{equation}\tag{1.2.4}
\sum_i^r d_i \geq n+r+1 \quad ,
\end{equation}
on peut trouver $r$ entiers

}{\begin{equation}\tag{1.2.4}
\sum_{i=1}^r d_i \geq n+r+1 \quad ,
\end{equation}
en posant $a_0=0$, on peut trouver $r$ entiers

}{F7II-XXI-010}\begin{equation}\tag{1.2.5}
1 \leq a_1 < a_2 < \dots < a_r = n+r+1
\end{equation}
vérifiant

\SGAvar{\begin{equation}\tag{1.2.6}
\left\{ \begin{array}{lcl} d_1 & \geq & a_o \\[3ex] d_2 & \geq & a_2 - a_1 \\[3ex] d_r & \geq & a_r - a_{r-1} \end{array} \right. \quad .
\end{equation}

}{\begin{equation}\tag{1.2.6}
\left\{ \begin{array}{lcl} d_1 & \geq & a_1 \\[3ex] d_2 & \geq & a_2 - a_1 \\[3ex] d_r & \geq & a_r - a_{r-1} \end{array} \right. \quad .
\end{equation}

}{F7II-XXI-011}\SGApage{II}{370}

On prend, pour $i = 1,\dots,r$, une forme homogène de degré $d_i$

\SGAvar{\begin{equation*}
H_i \, (X_{(a_{i-1}+1)},\dots,X_{a_i}) \ \in \ \mathbb{F}_q[X_{a_{i-1}+1},\dots,X_{a_i}]
\end{equation*}
en $a_i - a_{i-1}$ variable  qui a la propriété (1.2.2). L'intersection complète voulue est celle définie par les formes homogènes

\begin{equation*}
H_1(X_1,\dots,X_{a_1}), \ \ \dots\ , \ \ H_r(X_{a_{r-1}},\dots,X_{a_r}) \qquad .
\end{equation*}
\underline{Théorème} 1.3.  \underline{Soit} $X$ \underline{une intersection complète}, \underline{irréductible}, \underline{de dimension} $n \geq 1$ \underline{sur un corps} $k$ \underline{d'exposant caractéristique} $p$. \underline{Pour tout entier} $d \geq 1 + n$, \underline{une section assez générale de} $X$ \underline{par une hypersurface de degré} $d$ \underline{est non-spéciale} (1.0).

\underline{Démonstration}. Les sections $X.H_u$ de $X$ (par des hypersurfaces) de degré $d$ sont (paramétérisées par) les droites dans $H^0(X,\mathcal{O}_X(d))$, i.e. par

\begin{equation}\tag{1.3.1}
\mathbb{P}(H^0(X,\mathcal{O}_X(d)) = U
\end{equation}
et au-dessus de $U$, nous avons la ``section universelle de $X$ par une hypersurface de degré $d$ ''

}{\begin{equation*}
H_i \, (X_{(a_{i-1}+1)},\dots,X_{a_i}) \ \in \ \mathbb{F}_q[X_{a_{i-1}+1},\dots,X_{a_i}]
\end{equation*}
en $a_i - a_{i-1}$ variables qui a la propriété (1.2.2). L'intersection complète voulue est celle définie par les formes homogènes

\begin{equation*}
H_1(X_1,\dots,X_{a_1}), \ \ \dots\ , \ \ H_r(X_{a_{r-1}+1},\dots,X_{a_r}) \qquad .
\end{equation*}
\underline{Théorème} 1.3.  \underline{Soit} $X$ \underline{une intersection complète}, \underline{irréductible}, \underline{de dimension} $n \geq 1$ \underline{sur un corps} $k$ \underline{d'exposant caractéristique} $p$. \underline{Pour tout entier} $d \geq 1 + n$, \underline{une section assez générale de} $X$ \underline{par une hypersurface de degré} $d$ \underline{est non-spéciale} (1.0).

\underline{Démonstration}. Les sections $X.H_u$ de $X$ (par des hypersurfaces) de degré $d$ sont (paramétrisées par) les droites dans $H^0(X,\mathcal{O}_X(d))$, i.e. par

\begin{equation}\tag{1.3.1}
\mathbb{P}(H^0(X,\mathcal{O}_X(d))) = U
\end{equation}
et au-dessus de $U$, nous avons la ``section universelle de $X$ par une hypersurface de degré $d$ ''

}{F7II-XXI-012}\begin{equation}\tag{1.3.2}
f : X.H_u \longrightarrow U \qquad ,
\end{equation}
soit $U'$ l'ouvert des $u \in U$ tels que $\dim(X_u) = n-1$. Nous savons par (XI 1.5 (i) et 2.5) que le $\mathcal{O}_U$-module

\begin{equation}\tag{1.3.3}
R^{n-1}f_*(\mathcal{O}_{X.H_u}) \ \mid \ U' 
\end{equation}

\SGApage{II}{371}

 est un faisceau localement libre dont la formation commute à tout changement de base, et est de rang $M \geq 1$. Cela donne le théorème si $p = 1$, pour toute section par une hypersurface de degré $\geq 1+n$ .

\SGAvar{On suppose donc que $p > 0$ ; l'opération $\mathfrak{F}_{p^n}$ s'exprime encore localement en termes d'une matrice $M \times M$ à valeurs dans $\mathcal{O}_u$, de sorte que

}{On suppose donc que $p > 1$ ; l'opération $\mathfrak{F}_{p^M}$ s'exprime encore localement en termes d'une matrice $M \times M$ à valeurs dans $\mathcal{O}_U$, de sorte que

}{F7II-XXI-013}\begin{equation}\tag{1.3.4}
\{ \, u \in U' | X.H_u \text{ est non-spéciale} \} \ \text{ est ouvert dans } U' \ .
\end{equation}
Il suffit donc de démontrer que l'ensemble (1.3.4) est \underline{non-vide}.

Traitons d'abord le cas

\begin{equation}\tag{1.3.5}
k = \mathbb{F}_q \qquad\qquad .
\end{equation}
Quitte à remplacer $\mathbb{F}_q$ par une extension finie, on peut supposer

\begin{equation}\tag{1.3.6}
\mathrm{card}(X(\mathbb{F}_q)) \ > \ 0
\end{equation}
\begin{equation}\tag{1.3.7}
\begin{aligned}
&\text{il existe des coordonnées homogènes } X_1,\dots,X_{n+r+1} \\
&\text{dans l'espace ambiant } \mathbb{P}^{n+r}_k, \text{ telles que } X \text{ ne rencontre pas} \\
&\text{la variété linéaire } X_1=\dots=X_{n+1}=0.
\end{aligned}
\end{equation}
On prend alors une forme homogène de degré $d$

\SGAvar{\begin{equation}\tag{1.3.8}
H(X_1,\dots,X_{n+1}) \ \in \ \mathbb{F}_q[X_1,\dots, X_{n+1}]
\end{equation}
qui a la propriété (1.2.2) (ceci est possible car $d \geq 1+n$). On prend pour $H$ l'hypersurface dans $\mathbb{P}^{n+r}$ d'équation $H(X_1,\dots,X_{n+1}) = 0$. Alors
\SGApage{II}{372}
l'intersection $X.H$ est sans point $\mathbb{F}_q$-rationnel par construction, ce qui montre à la fois, grâce à (1.3.6), que $X \not\subset H$, donc $X.H$ est de dimension $n-1$ ($X$ étant irréductible) i.e. $H \in U'$, et que $X.H$ est non-spécial (1.5).

}{\begin{equation}\tag{1.3.8}
H(X_1,\dots,X_{n+1}) \ \in \ \mathbb{F}_q[X_1,\dots, X_{n+1}]
\end{equation}
qui a la propriété (1.2.2) (ceci est possible car $d \geq 1+n$). On prend pour $H$ l'hypersurface dans $\mathbb{P}^{n+r}$ d'équation $H(X_1,\dots,X_{n+1}) = 0$. Alors
\SGApage{II}{372}
l'intersection $X.H$ est sans point $\mathbb{F}_q$-rationnel par construction, ce qui montre à la fois, grâce à (1.3.6), que $X \not\subset H$, donc $X.H$ est de dimension $n-1$ ($X$ étant irréductible) i.e. $H \in U'$, et que $X.H$ est non-spécial (1.1.5).

}{F7II-XXI-014}Passons au cas général. Regardant les équations qui définissent $X$, on voit qu'on peut supposer que $k$ est de type fini sur $\mathbb{F}_p$, et qu'il existe un sous-anneau $B$ de $k$, de type fini (comme anneau) sur $\mathbb{F}_p$ dont le corps des fractions est $k$, tel que (1.3.9) $X \subset \mathbb{P}^{n+r}_k$ est la fibre générique d'une intersection complète relative $\widetilde{X} \subset \mathbb{P}^{n+r}_B$ qui est, \underline{fibre par fibre}, de même multidegré que $X$, et géométriquement irréductible.

Désignons par

\begin{equation}\tag{1.3.10}
\rho : \widetilde{X} \longrightarrow S = \mathrm{Spec}(B)
\end{equation}
la projection, par

\begin{equation}\tag{1.3.11}
E, \ \mbox{le } \mathcal{O}_S\mbox{-module localement libre } \rho_*\mathcal{O}_{\widetilde{X}}(d) \ ,
\end{equation}
et par

\SGAvar{\begin{equation}\tag{1.3.12}
\pi : T = \check{\mathbb{P}}(E) \longrightarrow S
\end{equation}
le fibré associé des droites dans $\mathbb{E}$ . Notons que la fibre générique de (1.3.12) n'est autre que $U = \mathbb{P}(H^0(X,\mathcal{O}_X(d))$. Au-dessus de $T$ nous avons la section universelle de $\widetilde{X} \subset \mathbb{P}^{n+r}_B$ par des hypersurfaces de degré $d$, notée 


}{\begin{equation}\tag{1.3.12}
\pi : T = \check{\mathbb{P}}(E) \longrightarrow S
\end{equation}
le fibré associé des droites dans $\mathbb{E}$ . Notons que la fibre générique de (1.3.12) n'est autre que $U = \mathbb{P}(H^0(X,\mathcal{O}_X(d)))$. Au-dessus de $T$ nous avons la section universelle de $\widetilde{X} \subset \mathbb{P}^{n+r}_B$ par des hypersurfaces de degré $d$, notée 


}{F7II-XXI-015}\SGApage{II}{373}

\SGAvar{\begin{equation}\tag{1.3.13}
 \alpha : \widetilde{X}.H_T \longrightarrow T \ .
\end{equation}
Ces copstructions s'insèrent dans le diagramme à carrés cartésiens

\begin{equation}\tag{1.3.14}
\begin{tikzcd}
\widetilde{X}.H_T \arrow[d, "\alpha"] \& X.H_u = (\widetilde{X}.H_T) \times_U T \arrow[l] \arrow[d] \\
T \arrow[d, "\pi"] \& U = T \times_S \mathrm{Sp}(k) \arrow[l] \arrow[d] \\
S \& \mathrm{Sp}(k) \arrow[l] \& ,
\end{tikzcd}
\end{equation}


















}{\begin{equation}\tag{1.3.13}
 \alpha : \widetilde{X}.H_T \longrightarrow T \ .
\end{equation}
Ces constructions s'insèrent dans le diagramme à carrés cartésiens

\begin{equation}\tag{1.3.14}
\begin{tikzcd}
\widetilde{X}.H_T \arrow[d, "\alpha"] \& X.H_u = (\widetilde{X}.H_T) \times_T U \arrow[l] \arrow[d] \\
T \arrow[d, "\pi"] \& U = T \times_S \mathrm{Sp}(k) \arrow[l] \arrow[d] \\
S \& \mathrm{Sp}(k) \arrow[l] \& ,
\end{tikzcd}
\end{equation}


















}{F7II-XXI-016}et nous avons encore que le $\mathcal{O}_T$-module

\SGAvar{\begin{equation}\tag{1.3.15}
R^{n-1}\alpha_*(\mathcal{O}_{\widetilde{X}.H_T})
\end{equation}
au-dessus de l'ouvert $T' \subset T$ formé des $t$ tels que $\dim \alpha^{-1}(t) = n-1$, est localement libre, et se formation y commute à tout changement de base. Compte tenu de ce que les schémas $T$ et $U$ ont le \underline{même} point générique, pour démontrer que (1.3.4) est non vide, il suffit de vérifier que l'ouvert de $T'$

}{\begin{equation}\tag{1.3.15}
R^{n-1}\alpha_*(\mathcal{O}_{\widetilde{X}.H_T})
\end{equation}
au-dessus de l'ouvert $T' \subset T$ formé des $t$ tels que $\dim \alpha^{-1}(t) = n-1$, est localement libre, et sa formation y commute à tout changement de base. Compte tenu de ce que les schémas $T$ et $U$ ont le \underline{même} point générique, pour démontrer que (1.3.4) est non vide, il suffit de vérifier que l'ouvert de $T'$

}{F7II-XXI-017}\begin{equation}\tag{1.3.16}
\{\ t \in T' \mid \widetilde{X}_{\pi(t)}.H_t \ \mbox{est non-spéciale}\}
\end{equation}
est non vide. Or, le cas (1.3.5) déjà traité démontre que pour \underline{tout point}
\SGApage{II}{374}
\underline{fermé} $s \in S$, la fibre $\pi^{-1}(s)$ a une intersection non vide avec l'ouvert (1.3.16), Q.E.D.

Des raisonnements tout à fait semblables permettent de démontrer:

\SGAvar{\underline{Théorème} 1.4. \underline{Soit} $X \subset \mathbb{P}^N$ \underline{un schéma projectif de Cohen-Macaulay,} \underline{équidimensionnel de dimension} $n \geq 1$, \underline{sur un corps} $k$, \underline{alors il existe} \underline{un entier} $d_o$ \underline{tel que pour} $d \geq d_o$, \underline{désignant par} $X.H \stackrel{i}{\longrightarrow} X$ \underline{l'inclusion} \underline{d'une section par une hypersurface} générale \underline{de degré} $d$, \underline{l'action de}

}{\underline{Théorème} 1.4. \underline{Soit} $X \subset \mathbb{P}^N$ \underline{un schéma projectif de Cohen-Macaulay,} \underline{équidimensionnel de dimension} $n \geq 1$, \underline{sur un corps} $k$, \underline{alors il existe} \underline{un entier} $d_0$ \underline{tel que pour} $d \geq d_0$, \underline{désignant par} $X.H \stackrel{i}{\longrightarrow} X$ \underline{l'inclusion} \underline{d'une section par une hypersurface} générale \underline{de degré} $d$, \underline{l'action de}

}{F7II-XXI-018}\begin{equation}\tag{1.4.0}
\mathfrak{F}_p \ \mathrm{sur}\ H^{n-1}(X.H,\mathcal{O}_{X.H}) \ /\ i^{*}H^{n-1}(X,\mathcal{O}_X)
\end{equation}
\underline{soit} non nilpotente.

\underline{Démonstration.} On se ramène, tout comme ci-dessus, au cas $k = \mathbb{F}_q$, et au problème de trouver un \underline{seul} $H$ (pour un degré donné assez grand et tel que $\dim(H^{n-1}(X.H,\mathcal{O}_{X.H}))$ soit indépendant de $H$ ; cf. 1.4.4) qui rend vrai (1.4.0). Quitte à faire agrandir $q$ , on peut supposer que les composantes irréductibles $X_i$ de $X$ sont géométriquement irréductibles, satisfont

\begin{equation}\tag{1.4.1}
\mathrm{Card}\ X_i(\ \mathbb{F}_q) \geq 2 \quad ,
\end{equation}
et qu'on ait

\SGAvar{(1.4.2) il y a des coordonnés projectives $X_o,\ldots,X_N$ telles que $X$ ne coupe pas la variété linéaire $X_o = \ldots = X_n = 0$ .

}{(1.4.2) il y a des coordonnées projectives $X_0,\ldots,X_N$ telles que $X$ ne coupe pas la variété linéaire $X_0 = \ldots = X_n = 0$ .

}{F7II-XXI-019}Par un théorème d'évanescence de Serre, ([1] ou SGA 2 XII 1.4), il existe un entier $d_1$ tel que pour tout entier $d \geq d_1$ , on ait :


\SGApage{II}{375}

\begin{equation}\tag{1.4.3}
H^i(X,\mathcal{O}_X(-d)) = 0 \ \mathrm{si}\ i \neq n = \dim X \quad .
\end{equation}
Désignant par $H$ une hypersurface de degré $d \geq d_1$, on a (par 1.4.3) :

\SGAvar{\begin{equation}\tag{1.4.4}
\begin{aligned} H^i(X,\mathcal{O}_X) &\stackrel{\sim}{\longrightarrow} H^i(X.H,\mathcal{O}_{X.H}) \qquad \mathrm{si}\ \ i \leq n - 2 \\[1ex] H^{n-1}(X,\mathcal{O}_X) &\lhook\joinrel\longrightarrow H^{n-1}(X.H,\mathcal{O}_{X.H})\ \mathrm{injectif}\ . \end{aligned}
\end{equation}
Nous allons démontrer le théorème en prenant $d_o = \max\ (1+N,d_1)$. Pour tout entier $d \geq d_o$, il faut trouver une hypersurface $H$ de degré $d$, définie sur $\mathbb{F}_q$, qui ne contienne aucun $X_i$, et telle que

\begin{equation}\tag{1.4.5}
\mathrm{trace}(\mathfrak{F}_q\ \mathrm{sur}\ H^{n-1}(X.H,\mathcal{O}_{X.H})\ /\ H^{n-1}(X,\mathcal{O}_X)) \neq 0 \quad .
\end{equation}
D'après la formule de la trace (XX, 3.2.1), on a

\begin{equation}\tag{1.4.6}
\begin{aligned}
\mathrm{card}(X.H\ (\mathbb{F}_q))\mathrm{mod}\ p
&= \sum_{i=0}^{n-1}(-1)^i
\mathrm{tr}(\mathfrak{F}_q\ \mathrm{sur}\ H^i(X.H,\mathcal{O}_{X.H})) \quad . \\
&\stackrel{(1.4.4)}{=} \sum_{i=0}^{n-1}(-1)^i
\mathrm{tr}(\mathfrak{F}_q\ \mathrm{sur}\ H^i(X,\mathcal{O}_X)) \\
&\qquad +(1)^{n-1}\mathrm{tr}\!\left(\mathfrak{F}_q\ \mathrm{sur}\
\left(\frac{H^{n-1}(X.H,\mathcal{O}_{X.H})}{H^{n-1}(X,\mathcal{O}_X)}\right.\right.
\end{aligned}
\end{equation}
Désignons par $a \in \mathbb{Z}/p\ \mathbb{Z}$ , la quantité

}{\begin{equation}\tag{1.4.4}
\begin{aligned} H^i(X,\mathcal{O}_X) &\stackrel{\sim}{\longrightarrow} H^i(X.H,\mathcal{O}_{X.H}) \qquad \mathrm{si}\ \ i \leq n - 2 \\[1ex] H^{n-1}(X,\mathcal{O}_X) &\lhook\joinrel\longrightarrow H^{n-1}(X.H,\mathcal{O}_{X.H})\ \mathrm{injectif}\ . \end{aligned}
\end{equation}
Nous allons démontrer le théorème en prenant $d_0 = \max\ (1+N,d_1)$. Pour tout entier $d \geq d_0$, il faut trouver une hypersurface $H$ de degré $d$, définie sur $\mathbb{F}_q$, qui ne contienne aucun $X_i$, et telle que

\begin{equation}\tag{1.4.5}
\mathrm{trace}(\mathfrak{F}_q\ \mathrm{sur}\ H^{n-1}(X.H,\mathcal{O}_{X.H})\ /\ H^{n-1}(X,\mathcal{O}_X)) \neq 0 \quad .
\end{equation}
D'après la formule de la trace (XXII, 3.2.1), on a

\begin{equation}\tag{1.4.6}
\begin{aligned}
\mathrm{card}(X.H\ (\mathbb{F}_q))\mathrm{mod}\ p
&\equiv \sum_{i=0}^{n-1}(-1)^i
\mathrm{tr}(\mathfrak{F}_q\ \mathrm{sur}\ H^i(X.H,\mathcal{O}_{X.H})) \quad . \\
&\stackrel{(1.4.4)}{\equiv} \sum_{i=0}^{n-1}(-1)^i
\mathrm{tr}(\mathfrak{F}_q\ \mathrm{sur}\ H^i(X,\mathcal{O}_X)) \\
&\qquad +(-1)^{n-1}\mathrm{tr}\!\left(\mathfrak{F}_q\ \mathrm{sur}\
\left(\frac{H^{n-1}(X.H,\mathcal{O}_{X.H})}{H^{n-1}(X,\mathcal{O}_X)}\right)\right)
\end{aligned}
\end{equation}
Désignons par $a \in \mathbb{Z}/p\ \mathbb{Z}$ , la quantité

}{F7II-XXI-020}\begin{equation*}
\sum_{i=0}^{n-1} (-1)^i \mathrm{tr}(\mathfrak{F}_q\ \mathrm{sur}\ H^i(X,\mathcal{O}_X)) \quad .
\end{equation*}

\SGApage{II}{376}

\SGAvar{Il suffit alors, pour tout $d \geq d_o$, de trouver une $H$ défini sur $\mathbb{F}_q$, de degré $d$, qui ne contient pas $X$, et tel que

}{Il suffit alors, pour tout $d \geq d_0$, de trouver un $H$ défini sur $\mathbb{F}_q$, de degré $d$, qui ne contient pas $X$, et tel que

}{F7II-XXI-021}\begin{equation}\tag{1.4.7}
\mathrm{Card}(X.H\ (\mathbb{F}_q)) \not\equiv a \bmod p \ .
\end{equation}
Il y a lieu de distinguer deux cas. Si $a \neq 0$, il suffit de prendre un $H$ sur $\mathbb{F}_q$ tel que $(X.H)\ (\mathbb{F}_q)$ soit vide, tout comme dans (1.3.8). Un tel $H$ vérifie (1.4.7), et ne contient aucun $X_i$ (par hypothèse on sait que $\mathrm{card}(X_i(\mathbb{F}_q)) \geq 2$).

\SGAvar{Si $a = 0$, il suffira de trouver un $H$ tel que $\mathrm{Card}(X.H)(\mathbb{F}_q)=1$ (1.4.7 sera vérifié, et $X_i \not\subset H$, car $\mathrm{Card}(X_i)(\mathbb{F}_q) \geq 2$). Pour ceci, on choisit des coordonnées projectives $X_o,\dots,X_N$ telles que $(1,0,\dots,0)$ soit un point de $X$. On prend $H$ d'équation $H(X_1,\dots,X_N) = 0$, où $H$ est une forme homogène telle que (cf. (1.2.2))

\begin{equation}\tag{1.4.8}
\alpha_1,\dots,\alpha_N \in \mathbb{F}_q, \ H(\alpha_1,.,\alpha_N) = 0 \Longrightarrow \alpha_1=\dots=\alpha_N = 0 \ .
\end{equation}
Alors, l'hypersurface définie par $H$ dans $\mathbb{P}^N$ n'a que le point $(1,0,\dots,0)$ comme point à valeurs dans $\mathbb{F}_q$, donc a fortiori, $X.H$ n'a que ce point comme $\mathbb{F}_q$-point. cqfd.


}{Si $a = 0$, il suffira de trouver un $H$ tel que $\mathrm{Card}((X.H)(\mathbb{F}_q))=1$ (1.4.7 sera vérifié, et $X_i \not\subset H$, car $\mathrm{Card}(X_i(\mathbb{F}_q)) \geq 2$). Pour ceci, on choisit des coordonnées projectives $X_0,\dots,X_N$ telles que $(1,0,\dots,0)$ soit un point de $X$. On prend $H$ d'équation $H(X_1,\dots,X_N) = 0$, où $H$ est une forme homogène telle que (cf. (1.2.2))

\begin{equation}\tag{1.4.8}
\alpha_1,\dots,\alpha_N \in \mathbb{F}_q, \ H(\alpha_1,\dots,\alpha_N) = 0 \Longrightarrow \alpha_1=\dots=\alpha_N = 0 \ .
\end{equation}
Alors, l'hypersurface définie par $H$ dans $\mathbb{P}^N$ n'a que le point $(1,0,\dots,0)$ comme point à valeurs dans $\mathbb{F}_q$, donc a fortiori, $X.H$ n'a que ce point comme $\mathbb{F}_q$-point. cqfd.


}{F7II-XXI-022}\SGApage{II}{377}


\SGAsectiontoc{2. Le coniveau dans le cas d'un corps de définition fini}{2}
2. \underline{Le coniveau dans le cas d'un corps de définition fini}

2.0. Soit $X$ un schéma séparé de type fini sur $\mathbb{F}_q$. Pour tout nombre premier $\ell \neq \mathrm{car}(\mathbb{F}_q)$, on désigne par

\SGAvar{\begin{equation}\tag{2.0.1}
H^i_c(\overline{X},\mathbb{Q}_\ell)
\end{equation}
les groupes de cohomologie à supports propres de $\overline{X} = X \otimes_{\mathbb{F}_q} \overline{\mathbb{F}}_q$, qui sont des $\mathbb{Q}_\ell$-vectoriels de rang fini sur lesquels opère le groupe de Galois $\mathrm{Gal}(/\mathbb{F}_q)$ de $\mathbb{F}_q$ . Désignons par

\begin{equation}\tag{2.0.2}
\varphi_q \in \mathrm{Gal}(/\mathbb{F}_q) \quad \mbox{l'automorphisme de Frobenius} \ .
\end{equation}
On sait (SGA 5 XIII 3.2.1) que la fonction zêta de $X$ (XXII 3.0.1) est donné par la formule

}{\begin{equation}\tag{2.0.1}
H^i_c(\overline{X},\mathbb{Q}_\ell)
\end{equation}
les groupes de cohomologie à supports propres de $\overline{X} = X \otimes_{\mathbb{F}_q} \overline{\mathbb{F}}_q$, qui sont des espaces vectoriels de dimension finie sur $\mathbb{Q}_\ell$, sur lesquels opère le groupe de Galois $\mathrm{Gal}(\overline{\mathbb{F}}_q/\mathbb{F}_q)$ de $\mathbb{F}_q$ . Désignons par

\begin{equation}\tag{2.0.2}
\varphi_q \in \mathrm{Gal}(\overline{\mathbb{F}}_q/\mathbb{F}_q) \quad \mbox{l'automorphisme de Frobenius} \ .
\end{equation}
On sait (SGA 5 XIII 3.2.1) que la fonction zêta de $X$ (XXII 3.0.1) est donnée par la formule

}{F7II-XXI-023}\begin{equation}\tag{2.0.3}
Z(t,X/\mathbb{F}_q) = \prod_{i \geq 0} \det (1-t\ \varphi_q^{-1}|H^i_c(\overline{X},\mathbb{Q}_\ell))^{(-1)^{i+1}} \ .
\end{equation}
DELIGNE a démontré récemment, en généralisant un résultat dû à WASHNITZER dans le cas $X$ projectif et lisse, le théorème suivant, qui est prouvé dans l'appendice plus bas (5.5.3) :

\underline{Théorème} 2.1. \underline{Les valeurs propres de} $\varphi_q^{-1}$ \underline{dans} $H^i_c(\overline{X},\mathbb{Q}_\ell)$ \underline{sont des entiers algébriques}, \underline{et}, \underline{si} $i > \dim X$, \underline{elles sont divisibles} (\underline{en tant qu'entiers algébriques}) \underline{par} $q^{i-\dim X}$ .


\SGApage{II}{378}

En particulier, on obtient (tautologiquement),

(2.1 bis) \quad les valeurs propres de $\varphi_q^{-1}$ dans $H^i_c(\overline{X},\mathbb{Q}_\ell(-j))$ sont des entiers algébriques divisibles par $q^j$ .

\underline{Corollaire} 2.2. \underline{Soit} $X$ \underline{propre lisse et géométriquement connexe sur} $\mathbb{F}_q$ . \underline{Alors les valeurs propres de} $\varphi_q^{-1}$ \underline{dans} $N^jH^q(\overline{X},\mathbb{Q}_\ell)$ (\underline{notations de} XX 2.0.6) \underline{sont des entiers algébriques divisibles par} $q^j$ .

\SGAvar{\underline{Démonstration.} L'énoncé est invariant par extension finie $\mathbb{F}_{q^M}/\mathbb{F}_q$ (qui revient à prendre les puissances. $M$\textsuperscript{èmes} des valeurs propres envisagées). Quitte à faire un tel changement de base, on a, par définition un nombre fini de schémas $T_i$ projectifs, lisses, et géométriquement connexes sur $\mathbb{F}_q$, des dimensions $d_i$, et de cycles algébriques

}{\underline{Démonstration.} L'énoncé est invariant par extension finie $\mathbb{F}_{q^M}/\mathbb{F}_q$ (ce qui revient à élever les valeurs propres envisagées à la puissance $M$). Quitte à faire un tel changement de base, on a, par définition un nombre fini de schémas $T_i$ projectifs, lisses, et géométriquement connexes sur $\mathbb{F}_q$, des dimensions $d_i$, et de cycles algébriques

}{F7II-XXI-024}$Z_i$ sur $T_i \times_{\mathbb{F}_q} X$ de codimension $j + d_i$

\SGAvar{tels que \underline{l'homomorphisme de Gal}$(/\mathbb{F}_q)$-\underline{modules}

\begin{equation*}
H^{q-2j}(\overline{T}_i,\ \mathbb{Q}_\ell(-j) \ \stackrel{\oplus \mathrm{Cl}(Z_i)}{\longrightarrow} \ N^jH^q(\overline{X},\mathbb{Q}_\ell)
\end{equation*}
soit surjectif, et on conclut par 2.1 bis.


}{tels que \underline{l'homomorphisme de }$\mathrm{Gal}(\overline{\mathbb{F}}_q/\mathbb{F}_q)$-\underline{modules}

\begin{equation*}
\bigoplus_i H^{q-2j}(\overline{T}_i,\ \mathbb{Q}_\ell(-j)) \ \stackrel{\oplus_i \mathrm{Cl}(Z_i)}{\longrightarrow} \ N^jH^q(\overline{X},\mathbb{Q}_\ell)
\end{equation*}
soit surjectif, et on conclut par 2.1 bis.


}{F7II-XXI-025}\SGApage{II}{379}


\SGAsectiontoc{3. Application aux intersections complètes}{3}
3. \underline{Application aux intersections complètes}

\SGAvar{3.0. Soit $X$ une intersection complète dans $\mathbb{P}^N$ lisse de dimension $n \geq 1$ sur un corps $k$. Pour $\ell \neq \mathrm{car}(k)$, on sait par ``Lefschetz faible'' et la dualité que (XI 1.6)

}{3.0. Soit $X$ une intersection complète lisse de dimension $n \geq 1$ dans $\mathbb{P}^N$ sur un corps $k$. Pour $\ell \neq \mathrm{car}(k)$, on sait par ``Lefschetz faible'' et la dualité que (XI 1.6)

}{F7II-XXI-026}\begin{equation}\tag{3.0.1}
H^i(\overline{X},\mathbb{Q}_\ell) \;=\; \left\{ \begin{array}{l} 0 \ \mbox{si $i$ impair et $i \neq n$} \\[1.5ex] \mathbb{Q}_\ell\!\left(\frac{-i}{2}\right) \ \mbox{si $i$ pair $\leq 2n$ et $i \neq n$} \\[1.5ex] H^n(\overline{X},\mathbb{Q}_\ell) \ \mbox{si $i=n$} \quad . \end{array} \right.
\end{equation}
Pour $k = \mathbb{F}_q$, on obtient donc une formule pour la fonction zêta :

\SGAvar{\begin{equation}\tag{3.0.2}
Z(t,X/\mathbb{F}_q) \ \prod_{i=0}^{n} \ (1-q^i t) = \det(1-t\,\varphi_q^{-1} \mid \mathrm{Prim}^n(\overline{X},\mathbb{Q}_\ell))^{(-1)^{n+1}} \quad ,
\end{equation}
(on a employé da décomposition primitive, et (3.0.1).)

}{\begin{equation}\tag{3.0.2}
Z(t,X/\mathbb{F}_q) \ \prod_{i=0}^{n} \ (1-q^i t) = \det(1-t\,\varphi_q^{-1} \mid \mathrm{Prim}^n(\overline{X},\mathbb{Q}_\ell))^{(-1)^{n+1}} \quad ,
\end{equation}
(on a employé la décomposition primitive, et (3.0.1).)

}{F7II-XXI-027}\underline{Corollaire} 3.0.3. \underline{Sous les hypothèses} 3.0, \underline{avec} $k = \mathbb{F}_q$ , \underline{on a une congruence modulo} $p$

\begin{equation*}
\det(1-t\,\varphi_q^{-1} \mid H^n(\overline{X},\mathbb{Q}_\ell)) \;\equiv\; \det(1-t\,\mathfrak{F}_q \mid H^n(X,\mathcal{O}_X)) \quad .
\end{equation*}
\underline{Démonstration.} D'après la formule de congruence (XXII, 3.2.1), et compte tenu de ce que

\SGAvar{\begin{equation}\tag{3.0.4}
H^i(X,\mathcal{O}_X) \;\simeq\; \left\{ \begin{array}{ll} \mathbb{F}_q & i = 0 \\[1.5ex] 0 & 1 \leq \; \leq n - 1 \\[1.5ex] 0 & i \geq n+1 \end{array} \right.
\end{equation}

}{\begin{equation}\tag{3.0.4}
H^i(X,\mathcal{O}_X) \;\simeq\; \left\{ \begin{array}{ll} \mathbb{F}_q & i = 0 \\[1.5ex] 0 & 1 \leq i \leq n - 1 \\[1.5ex] 0 & i \geq n+1 \end{array} \right.
\end{equation}

}{F7II-XXI-028}\SGApage{II}{380}
on a

\begin{equation}\tag{3.0.5}
Z(t,X/\mathbb{F}_q) \;\equiv\; (1-t)^{-1}\ \det\,(1-t\mathfrak{F}_q|H^n(X,\mathcal{O}_X))^{(-1)^{n+1}} \quad \mathrm{mod}\ p
\end{equation}
de sorte que (3.0.5) et (3.0.2) donnent la congruence

\begin{equation}\tag{3.0.6}
\det(1-t\ \varphi_q^{-1}|\mathrm{Prim}^n(\overline{X},\mathbb{Q}_\ell)) \;\equiv\; \det(1-t\ \mathfrak{F}_q|H^n(X,\mathcal{O}_X)) \quad \mathrm{mod}\ p \quad .
\end{equation}
La conclusion voulue s'en déduit, compte tenu de ce que

\begin{equation}\tag{3.0.7}
H^n(\overline{X},\mathbb{Q}_\ell) \;\simeq\; \begin{cases} \mathrm{Prim}^n(\overline{X},\mathbb{Q}_\ell) & \text{si } n \text{ est impair} \\[2mm] \mathrm{Prim}^n(\overline{X},\mathbb{Q}_\ell) \oplus \mathbb{Q}_\ell\!\left(\frac{-n}{2}\right) & \text{si } n \text{ est pair} \end{cases} \quad .
\end{equation}
\underline{Corollaire} 3.1. \underline{Soit} $X$ \underline{une intersection complète}, \underline{lisse de dimension} $n \geq 1$ \underline{sur} $\mathbb{F}_q$ . \underline{Si}

\begin{equation}\tag{3.1.0}
H^n(\overline{X},\mathbb{Q}_\ell) \;=\; N^1H^n(\overline{X},\mathbb{Q}_\ell)
\end{equation}
\underline{alors}

\begin{equation}\tag{3.1.1}
\mathfrak{F}_p\ \underline{\text{est nilpotent sur}}\ H^n(X,\mathcal{O}_X) \quad .
\end{equation}
\underline{Démonstration}. D'après (2.2), (3.1.0) implique qu'on a

\begin{equation}\tag{3.1.2}
\det(1-t\ \varphi_q^{-1}\ |\ H^n(\overline{X},\mathbb{Q}_\ell)) \;\equiv\; 1 \quad \mathrm{mod}\ p\ (\text{et même mod } q) \quad ,
\end{equation}
et ceci, avec (3.0.3), implique qu'on a

\begin{equation}\tag{3.1.3}
\det(1-t\ \mathfrak{F}_q\ |\ H^n(X,\mathcal{O}_X)) \;=\; 1 \quad , 
\end{equation}

\SGApage{II}{381}

 ce qui implique que $\mathfrak{F}_q$ est nilpotent sur $H^n(X,\mathcal{O}_X)$, et donc que $\mathfrak{F}_p$ l'est, cqfd.

\underline{Proposition} 3.2. Soit $X$ une intersection complète, lisse de dimension $n \geq 1$ sur un corps $k$. Si

\begin{equation}\tag{3.2.1}
H^n(\overline{X},\mathbb{Q}_\ell) \;=\; N^1H^n(\overline{X},\mathbb{Q}_\ell)
\end{equation}
alors $X$ est spéciale, i.e.

\begin{equation}\tag{3.2.2}
\mathfrak{F}_p\ \text{est nilpotente sur}\ H^n(X,\mathcal{O}_X) \quad .
\end{equation}
\underline{Démonstration}. On peut supposer $k$ de type fini sur $\mathbb{F}_p$ (le cas $p = 0$ étant trivial), et qu'il existe une sous-$\mathbb{F}_p$-algèbre $B$ de type fini de $k$, telle que $X/k$ soit la fibre générique d'une intersection complète lisse relative $\widetilde{X}$ :

\SGAvar{\begin{equation}\tag{3.2.3}
\begin{tikzcd}
\widetilde{X} \arrow[d, "\pi"'] \& X = \widetilde{X}/_{\eta} \arrow[l] \arrow[d] \\
\mathrm{Spec}(B) = S \& \eta = \mathrm{Sp}(k) \;\; . \arrow[l]
\end{tikzcd}
\end{equation}












}{\begin{equation}\tag{3.2.3}
\begin{tikzcd}
\widetilde{X} \arrow[d, "\pi"'] \& X = \widetilde{X}_{\eta} \arrow[l] \arrow[d] \\
\mathrm{Spec}(B) = S \& \eta = \mathrm{Sp}(k) \;\; . \arrow[l]
\end{tikzcd}
\end{equation}












}{F7II-XXI-029}On sait que le $B$-module $R^n\pi_*(\mathcal{O}_{\widetilde{X}})$ est localement libre, et sa formation commute à tout changement de base, donc l'ensemble des points 


\SGApage{II}{382}

\begin{equation}\tag{3.2.4}
 \{\, s \in S \mid X_s \text{ est spéciale} \,\}
\end{equation}
est un fermé de $S$.

Par (XX, 2.3), quitte à rapetisser $S$, nous pouvons supposer que

\begin{equation}\tag{3.2.5}
\forall\, s \in S, \quad H^n(X_{\bar{s}}, \mathbb{Q}_\ell) \;=\; N^1 H^n(X_{\bar{s}}, \mathbb{Q}_\ell) \quad .
\end{equation}
Par (3.1), pour tout point \underline{fermé} $s$ de $S$, $X_s$ est spécial. Il s'ensuit que le fermé (3.2.4) de $S$ est $S$ tout entier, donc, en particulier, que $X = \tilde{X}_\eta$ est spécial.


\SGAsectiontoc{4. Les conclusions}{4}
4. \underline{Les conclusions}

Mettant ensemble 1.1 et 3.2, on trouve :

\underline{Théorème} 4.1. \underline{Fixons une caractéristique} $p$, \underline{une dimension} $n$, \underline{et un multidegré} $(d_1,\ldots,d_r)$. \underline{Désignons par}

\begin{equation}\tag{4.1.0}
X_V \longrightarrow V
\end{equation}
\underline{l'intersection complète lisse de dimension} $n$ \underline{et multidegré} $(d_1,\ldots,d_r)$ universelle \underline{en caractéristique} $p$.

\underline{Si}

\SGAvar{\begin{equation}\tag{4.1.1}
\sum_{1}^{r} d_i \geqslant n + r + 1 \quad ,
\end{equation}
\underline{alors il existe un ouvert dense} $U \subset V$ \underline{tel que}

}{\begin{equation}\tag{4.1.1}
\sum_{i=1}^{r} d_i \geqslant n + r + 1 \quad ,
\end{equation}
\underline{alors il existe un ouvert dense} $U \subset V$ \underline{tel que}

}{F7II-XXI-030}\begin{equation}\tag{4.1.2}
\underline{\text{pour}}\ u \in U, \ \underline{\text{on a}}\ H^n(X_{\bar{u}}, \mathbb{Q}_\ell) \neq N^1 H^n(X_{\bar{u}}, \mathbb{Q}_\ell) \quad .
\end{equation}

\SGApage{II}{383}

\SGAvar{Mettant ensemble 1.3 et 3.2, on trouve de meme .

}{Mettant ensemble 1.3 et 3.2, on trouve de même.

}{F7II-XXI-031}\underline{Théorème} 4.2. \underline{Soit} $X$ \underline{une intersection complète,} \underline{lisse de dimension} $n \geqslant 1$ \underline{sur un corps} $k$ . \underline{Fixons un entier}

\begin{equation}\tag{4.2.1}
d \geqslant 1 + n \quad .
\end{equation}
\underline{Désignons par}

\begin{equation}\tag{4.2.2}
X.H_u \longrightarrow U
\end{equation}
\underline{la famille} universelle \underline{des sections de} $X$ \underline{par des hypersurfaces de degré} $d$, \underline{qui sont lisses de dimension} $n - 1$. \underline{Il existe un ouvert dense} $T \subset U$ \underline{tel que}

\begin{equation}\tag{4.2.3}
\underline{\text{pour}}\ t \in T, \ \underline{\text{on a}}\ H^{n-1}(X.H_{\bar{t}}, \mathbb{Q}_\ell) \neq N^1 H^{n-1}(X.H_{\bar{t}}, \mathbb{Q}_\ell) \quad .
\end{equation}

\underline{BIBLIOGRAPHIE}

[1]  J.P.  SERRE.  Faisceaux algébriques cohérents, Annals of Math, 61, 1955, pp.197-278.

[2]  J.P.  SERRE.  Sur la topologie des variétés algébriques en caractéristique $p$, dans Symposium Internacional de topologia algebraica, Mexico 1958.


\SGApage{II}{384}


\SGAsectiontoc{§ 5. Appendice , Théorèmes d'intégralité. Par P. Deligne}{5}
§ 5. \underline{Appendice} , \underline{Théorèmes d'intégralité}. Par P. Deligne

\SGAvar{5.0. Soient $p$ un nombre premier, $q$ une puissance de $p$, $\mathbb{F}_q$ un corps à $q$ éléments, $E$ le complété d'un corps de nombresen une place divisant un nombre premier $\ell$ distinct de $p$, et $T$ un ensemble de nombres premiers (par exemple $T = \emptyset$).

}{5.0. Soient $p$ un nombre premier, $q$ une puissance de $p$, $\mathbb{F}_q$ un corps à $q$ éléments, $E$ le complété d'un corps de nombres en une place divisant un nombre premier $\ell$ distinct de $p$, et $T$ un ensemble de nombres premiers (par exemple $T = \emptyset$).

}{F7II-XXI-032}Un élément $\alpha$ d'une clôture algébrique de $E$ est dit \underline{T-entier} s'il est algébrique sur $\mathbb{Q}$ et entier sur $\mathbb{Z}\bigl[(1/r)_{r \in T}\bigr]$. Un endomorphisme $F$ d'un vectoriel de dimension finie sur $E$ est dit \underline{T-entier} si ses valeurs propres sont T-entières.

Soient $k$ un corps fini à $q'$ éléments, $\bar{k}$ une clôture algébrique de $k$ et $\mathfrak{G} = \operatorname{Gal}(\bar{k}/k)$ le groupe de Galois de $k$. On désigne par $f \in \mathfrak{G}$ la substitution de Frobenius, donnée par

\begin{equation*}
f(x) = x^{q'} \qquad ,
\end{equation*}
et par $F$ le Frobenius ``géométrique'', inverse de $f$ dans $\mathfrak{G}$. Une représentation $\rho$ de $\mathfrak{G}$ dans un vectoriel de dimension finie sur $E$ est dite T-entière si $\rho(F)$ est T-entier. Si $X$ est un schéma de type fini sur $\mathbb{F}_q$ et $x$ un point fermé de $X$, on désigne par $F_x$ l'élément de Frobenius géométrique du groupe de Galois de $k(x)$ ; si $\bar{x}$ est un point géométrique de $X$ localisé en $x$ et $\underline{G}$ un E-faisceau sur $X$, $F_x$ agit par transport de structure sur la fibre $G_{\bar{x}}$ .

\underline{Définition} 5.1. \underline{Un} E-\underline{faisceau constructible} $\underline{G}$ \underline{sur un schéma} $X$ \underline{de type fini sur} $\mathbb{F}_q$ \underline{est dit} T-\underline{entier si, pour tout point géométrique} $\bar{x}$ \underline{de} $X$ \underline{localisé en un point fermé} $x$ \underline{de} $X$, $F_x$ \underline{est un automorphisme} T-\underline{entier de} $\underline{G}_{\bar{x}}$ .


\SGApage{II}{385}

Si $a : X \longrightarrow \mathrm{Spec}(\mathbb{F}_q)$ est un schéma de type fini sur $\mathbb{F}_q$ et $\underline{G}$ un E-faisceau constructible sur $X$, on désignera par $\underline{H}^*(X,\underline{G})$ les E-faisceaux $R^*a_*(\underline{G})$ sur $\mathrm{Spec}(\mathbb{F}_q)_{\acute{e}t}$ ; si $\overline{\mathbb{F}}_q$ est une clôture algébrique de $\mathbb{F}_q$ et si $\overline{X}$ est le schéma déduit de $X$ par extension des scalaires de $\mathbb{F}_q$ à $\overline{\mathbb{F}}_q$, les $\underline{H}^*(X,\underline{G})$ s'identifient aux groupes de cohomologie $H^*(\overline{X},\underline{G})$, considérés comme modules galoisiens, i.e. munis d'une action du Frobenius géométrique $F$. De même en cohomologie à support propre.

\underline{Lemme} 5.2.1. \underline{Soit} $\underline{G}$ \underline{un} E-\underline{faisceau constructible sur un schéma} $X$ \underline{de type fini sur} $\mathbb{F}_q$ \underline{et de dimension} $\leq n$.

(i) \underline{Il existe une partie fermée} $Y$ \underline{de dimension} $0$ \underline{de} $X$ \underline{telle que la flèche évidente de} $\underline{H}^0(X,\underline{G})$ \underline{dans} $\underline{H}^0(Y,\underline{G})$ \underline{soit injective.}

(ii) \underline{Si} $X$ \underline{est séparé, et si} $U$ \underline{est un ouvert de} $X$ \underline{dont le complémentaire} $Z$ \underline{est de dimension} $< n$, \underline{il existe une partie fermée} $Y$ \underline{de dimension} $0$ \underline{de} $U$ \underline{et une flèche surjective de} $\underline{H}^0(Y,\underline{G})(-n)$ \underline{dans} $\underline{H}^{2n}_c(X,G)$.

L'assertion (i) est évidente. Prouvons (ii). Quitte à remplacer $X$ par $X_{\mathrm{red}}$ et à rétrécir $U$, on peut supposer que $U$ est lisse et que $G|U$ est constant tordu. On a, puisque $\dim(Z) < n$,

\begin{equation*}
0 = \underline{H}^{2n-1}_c(Z,\underline{G}) \longrightarrow \underline{H}^{2n}_c(U,\underline{G}) \stackrel{\sim}{\longrightarrow} \underline{H}^{2n}_c(X,\underline{G}) \longrightarrow \underline{H}^{2n}_c(Z,G) = 0
\end{equation*}
et, par dualité de Poincaré (SGA 4 XVIII et [2])

\begin{equation*}
\underline{H}^{2n}_c(U,\underline{G}) \quad \simeq \quad \underline{H}^0(U,\underline{G}^{\vee}(n))^{\vee} \qquad .
\end{equation*}
On conclut par (i).


\SGApage{II}{386}

\SGAvar{Théorème 5.2.2. Soient $X$ un schéma séparé de type fini sur $\mathbb{F}_q$ et de dimension $\leq n$, et $\underline{G}$ un E-faisceau constructible T-entier. Alors, les endomorphismes $F$ et $q^{n-i} F$ \underline{de} $H^i_c (X,\underline{G})$ sont T-entiers.

}{Théorème 5.2.2. Soient $X$ un schéma séparé de type fini sur $\mathbb{F}_q$ et de dimension $\leq n$, et $\underline{G}$ un E-faisceau constructible T-entier. Alors, les endomorphismes $F$ et $q^{n-i} F$ \underline{de} $\underline{H}^i_c (X,\underline{G})$ sont T-entiers.

}{F7II-XXI-033}Le cas où $\dim X \leq 0$ est trivial, et les cas $i = 0$ et $i = 2n$ du théorème résultent de 5.2.1. Soit $Z(X,\underline{G},t)$ la série formelle

\SGAvar{\begin{equation*}
Z(X,\underline{G},t) \;=\; \prod_{x \in X} \det (1-F_x \, t^{\deg(x)} \;;\; \underline{G}_x)^{-1} \in F[[t]] \;.
\end{equation*}
D'après (SGA 5 XV 3.2), on a

}{\begin{equation*}
Z(X,\underline{G},t) \;=\; \prod_{x \in X} \det (1-F_x \, t^{\deg(x)} \;;\; \underline{G}_x)^{-1} \in E[[t]] \;.
\end{equation*}
D'après (SGA 5 XV 3.2), on a

}{F7II-XXI-034}\begin{equation}\tag{5.2.3}
Z(X,\underline{G},t) \;=\; \prod_i \det (1-Ft \;;\; \underline{H}^i_c(X,\underline{G}))^{(-1)^{i+1}} \;.
\end{equation}
Prouvons 5.2.2 pour $n = 1$. On a alors $\underline{H}^i_c(X,\underline{G}) = 0$ pour $i > 2$ (SGA 4 X 4.3 ou XVII 5.2.8.1). Par hypothèse, la série formelle $Z(X,\underline{G},t)$ est à coefficients T-entiers. D'après les cas $i = 0$ et $i = 2n = 2$ de 5.2.2, le produit de $Z(X,\underline{G},t)$ par

\begin{equation*}
\det(1-Ft \;;\; \underline{H}^0_c(X,\underline{G})). \; \det(1-Ft \;;\; \underline{H}^2_c(X,\underline{G}))
\end{equation*}
est encore à coefficients T-entiers. Ce produit n'est autre que $\det(1-Ft \;;\; \underline{H}^1_c(X,\underline{G}))$, et le théorème en résulte pour $n = 1$.

Procédons maintenant par récurrence sur $n = \dim(X)$.

Soit $U$ un ouvert de Zariski de $X$, dont le complément $Z$ soit de dimension $< n$, et tel qu'il existe un morphisme $f : U \longrightarrow Y$, de but une courbe et à fibres de dimension $\leq n-1$.


\SGApage{II}{387}

\begin{center}
\begin{tikzcd}
U \arrow[r, hook, "j"] \arrow[d, "f"'] \& X \& Z \arrow[l, hook', "i"'] \\
Y \& .
\end{tikzcd}
\end{center}







La suite exacte

\begin{equation*}
0 \longrightarrow j_! j^* \underline{G} \longrightarrow \underline{G} \longrightarrow i_! i^* \underline{G} \longrightarrow 0
\end{equation*}
induit des suites exactes

\begin{equation*}
\underline{H}^i_c(U,\underline{G}) \longrightarrow \underline{H}^i_c(X,\underline{G}) \longrightarrow \underline{H}^i_c(Z,\underline{G})
\end{equation*}
de sorte que par récurrence, il suffit de vérifier le théorème pour le faisceau $\underline{G}$ sur $U$. Grâce à l'hypothèse de récurrence appliquée aux fibres de $f$, et au théorème de changement de base pour un morphisme propre, on sait que les faisceaux $R^i f_! \underline{G}$ sur $Y$ sont T-entiers, et que les faisceaux $R^i f_! \underline{G}(i-(n-1))$ sont T-entiers. La suite spectrale de Leray pour $f$ s'écrit

\begin{equation*}
E_2^{pq} = \underline{H}^p_c(Y,R^q f_! \underline{G}) \;\Longrightarrow\; \underline{H}^{p+q}_c(U,\underline{G}) \;.
\end{equation*}
En vertu de ce qui précède, et du théorème déjà connu pour $Y$, les $E_2^{pq}$ sont T-entiers, et les $E_2^{pq}(p+q-n)$ le sont aussi. Le théorème en résulte.

\underline{Corollaire} 5.3. \underline{Soient} $X$ \underline{un schéma de type fini sur} $\mathbb{F}_q$, \underline{de dimension} $\leq 1$, $U$ \underline{un ouvert de} $X$, $j : U \hookrightarrow X$ \underline{le morphisme d'inclusion et} $\underline{G}$ \underline{un} E-\underline{faisceau constructible sur} $U$. \underline{Alors} 


\SGApage{II}{388}

(a)\quad \underline{Si} $\underline{G}$ \underline{est} $T$-\underline{entier}, \underline{alors} $j_*\underline{G}$ \underline{est} $T$-\underline{entier}.

(b)\quad \underline{Si pour tout} $x \in U$, $\underline{G}^\vee_x$ \underline{est} $T$-\underline{entier}, \underline{alors}, \underline{pour tout} $x \in X$, $(j_*\underline{G})_x^\vee$ \underline{est} $T$-\underline{entier}.

On se ramène au cas $X$ affine, puis $X$ projectif.

Définissons $\underline{H}$ par la suite exacte

\begin{equation*}
0 \longrightarrow j_!\underline{G} \longrightarrow j_*\underline{G} \longrightarrow \underline{H} \longrightarrow 0 \quad .
\end{equation*}
La suite exacte de cohomologie correspondante fournit, si $Y$ est le support (fini) de $\underline{H}$, une suite exacte

\SGAvar{\begin{equation*}
\underline{H}^0(U,\underline{G}) \longrightarrow \underline{H}^0(Y,\underline{H}) \longrightarrow \underline{H}^1_c(U,\underline{G}) \quad .
\end{equation*}
On sait, par 5.2.1 et par le théorème 5.2 que $\underline{H}^0(U,\underline{G})$ et $\underline{H}^1_c(U,\underline{G})$ sont $T$-entiers ; il en résulte que $\underline{H}^0(Y,\underline{H})$ est $T$-entier, ce qui prouve (a).

Si $X'$ est le normalisé de la partie purement de dimension un de $X_{\mathrm{red}}$ et $U'$ un ouvert de $U$, dense dans le lieu lisse de dimension 1 de $U_{\mathrm{red}}$, et tel que $\underline{G}$ restreint à $U$ soit constant tordu :

}{\begin{equation*}
\underline{H}^0(U,\underline{G}) \longrightarrow \underline{H}^0(Y,\underline{H}) \longrightarrow \underline{H}^1_c(U,\underline{G}) \quad .
\end{equation*}
On sait, par 5.2.1 et par le théorème 5.2.2 que $\underline{H}^0(U,\underline{G})$ et $\underline{H}^1_c(U,\underline{G})$ sont $T$-entiers ; il en résulte que $\underline{H}^0(Y,\underline{H})$ est $T$-entier, ce qui prouve (a).

Si $X'$ est le normalisé de la partie purement de dimension un de $X_{\mathrm{red}}$ et $U'$ un ouvert de $U$, dense dans le lieu lisse de dimension 1 de $U_{\mathrm{red}}$, et tel que $\underline{G}$ restreint à $U'$ soit constant tordu :

}{F7II-XXI-035}\begin{center}
\begin{tikzcd}
U' \arrow[r, hook, "j'"'] \arrow[d, hook] \& X' \arrow[d, hook, "p"] \& {} \\
U \arrow[r, hook, "j"'] \& X \& Y \arrow[l, hook'] \& ,
\end{tikzcd}
\end{center}















\SGAvar{alors $p_*j'_*\,\underline{G} \mid Y \simeq j_*\underline{G}|Y$. Il suffit donc d'étudier le cas où $X$ est lisse de dimension un et où $\underline{G}|U$ est constant tordu. Le faisceau $\underline{G}^\vee$ est alors $T$-entier, de sorte qu'on sait déjà que pour tout $x \in X$ $(j_* \underline{G}^\vee)_x$
\SGApage{II}{389}
est $T$-entier. Désignnons par $\mathfrak{G}_x$ le groupe de Galois du corps des fractions de $\hat{\mathcal{O}}_{X,x}$ pour $x \in X$, et par $\mathfrak{I}_x$ le groupe d'inertie. Le faisceau $\underline{G}$ définit alors une représentation, encore notée $\underline{G}$, de $\mathfrak{G}_x$, et le $\mathfrak{G}_x/\mathfrak{I}_x = \operatorname{Gal}(\overline{k(x)}/k(x))$-module $(j_*\underline{G}^\vee)_x$ (resp $(j_*\underline{G})_x^\vee$) s'identifie alors au $\mathfrak{G}_x/\mathfrak{I}_x$-module des invariants (resp. des coinvariants) dans $\underline{G}^\vee$. Il reste à prouver :

\underline{Lemme} 5.3.1.\quad \underline{Soient} $K$ \underline{un corps local de corps résiduel} $\mathbb{F}_q$, $G$ \underline{son groupe de Galois}, $I$ \underline{le groupe d'inertie}, \underline{et} $\rho$ \underline{une représentation continue de} $G$ \underline{dans un vectoriel} $V$ \underline{de dimension} $n$ \underline{sur} $E$ (cf. 5.0), $V^I$ \underline{l'espace des invariants sous} $I$, \underline{et} $V_I$ \underline{l'espace des coinvariants sous} $I$, \underline{considérés comme} $G/I$-\underline{modules}. \underline{Soient} $\alpha_1 \ldots \alpha_n$ \underline{les valeurs propres de} $F$ \underline{agissant sur} $V^I$. \underline{Alors, il existe des entiers} $m_i \geq 0$ \underline{tels que les valeurs propres de} $F$ \underline{agissant sur} $V_I$ \underline{soient les} $\alpha_i q^{m_i}$.

}{alors $p_*j'_*\,\underline{G} \mid Y \simeq j_*\underline{G}|Y$. Il suffit donc d'étudier le cas où $X$ est lisse de dimension un et où $\underline{G}|U$ est constant tordu. Le faisceau $\underline{G}^\vee$ est alors $T$-entier, de sorte qu'on sait déjà que pour tout $x \in X$ $(j_* \underline{G}^\vee)_x$
\SGApage{II}{389}
est $T$-entier. Désignons par $\mathfrak{G}_x$ le groupe de Galois du corps des fractions de $\hat{\mathcal{O}}_{X,x}$ pour $x \in X$, et par $\mathfrak{I}_x$ le groupe d'inertie. Le faisceau $\underline{G}$ définit alors une représentation, encore notée $\underline{G}$, de $\mathfrak{G}_x$, et le $\mathfrak{G}_x/\mathfrak{I}_x = \operatorname{Gal}(\overline{k(x)}/k(x))$-module $(j_*\underline{G}^\vee)_x$ (resp. $(j_*\underline{G})_x^\vee$) s'identifie alors au $\mathfrak{G}_x/\mathfrak{I}_x$-module des invariants (resp. des coinvariants) dans $\underline{G}^\vee$. Il reste à prouver :

\underline{Lemme} 5.3.1.\quad \underline{Soient} $K$ \underline{un corps local de corps résiduel} $\mathbb{F}_q$, $G$ \underline{son groupe de Galois}, $I$ \underline{le groupe d'inertie}, \underline{et} $\rho$ \underline{une représentation continue de} $G$ \underline{dans un vectoriel} $V$ \underline{de dimension} $n$ \underline{sur} $E$ (cf. 5.0), $V^I$ \underline{l'espace des invariants sous} $I$, \underline{et} $V_I$ \underline{l'espace des coinvariants sous} $I$, \underline{considérés comme} $G/I$-\underline{modules}. \underline{Posons} $r=\dim_E V^I$ \underline{et soient} $\alpha_1,\ldots,\alpha_r$ \underline{les valeurs propres de} $F$ \underline{agissant sur} $V^I$. \underline{Alors, il existe des entiers} $m_i \geq 0$ $(1\leq i\leq r)$ \underline{tels que les valeurs propres de} $F$ \underline{agissant sur} $V_I$ \underline{soient les} $\alpha_i q^{m_i}$.

}{F7II-XXI-036}Soit $I'$ le plus grand sous-groupe de $I$ d'ordre premier à $\ell$. Le groupe $\rho(I')$ est fini, étant un sous-groupe compact d'ordre premier à $\ell$ de $\mathrm{GL}(V)$. On a donc

\begin{equation*}
V_I \quad \simeq \quad (V^{\rho(I')})_I
\end{equation*}
et, remplaçant $V$ par $V^{I'}$, on peut supposer que $I'$ agit trivialement sur $V$. On a $I/I' \simeq \mathbb{Z}_\ell(1)$ (en tant que $G/I$-modules) ; soit $T$ un générateur topologique de $I/I'$, et, pour chaque $\alpha$, soit $V^\alpha$ le sous-espace propre généralisé de $T$ dans $V$ de valeur propre $\alpha$ : on a $V = \oplus\, V^{(\alpha)}$ ; et la décomposition $V = V^{(1)} \oplus \sum_{\alpha \neq 1} V^{(\alpha)}$ est stable par $G$, et compatible
\SGApage{II}{390}
au passage au dual ; on peut donc remplacer $V$ par $V^{(1)}$, ce qui permet de supposer que $I$ agit de façon unipotente.

\SGAvar{\underline{Lemme} 5.3.2. \underline{Soient} $V$ \underline{un vectoriel de dimension finie sur un corps} $k$ \underline{de caractéristique} $0$, $I$ \underline{un groupe additif à un paramètre}, $\rho$ $I \longrightarrow \mathrm{Gl}(V)$ \underline{une représentation de} $I$ \underline{dans} $V$, $N$ \underline{un générateur infinitésimal de} $I$ \underline{et posons}

}{\underline{Lemme} 5.3.2. \underline{Soient} $V$ \underline{un vectoriel de dimension finie sur un corps} $k$ \underline{de caractéristique} $0$, $I$ \underline{un groupe additif à un paramètre}, $\rho : I \longrightarrow \mathrm{Gl}(V)$ \underline{une représentation de} $I$ \underline{dans} $V$, $N$ \underline{un générateur infinitésimal de} $I$ \underline{et posons}

}{F7II-XXI-037}\begin{equation*}
V^{I,k} = V^I \ \cap \ \mathrm{Im}(N^k)
\end{equation*}
\begin{equation*}
V_{I,k} = \mathrm{Im}(\mathrm{Ker}(N^{k+1}) \longrightarrow V_I) \quad .
\end{equation*}
\underline{Alors, la flèche} $N^k$ \underline{induit un isomorphisme}

\begin{equation*}
N^k : V_{I,k}/V_{I,k-1} \ \overset{\sim}{\longrightarrow} \ V^{I,k}/V^{I,k+1}
\end{equation*}
\underline{et définit des isomorphismes canoniques}

\SGAvar{\begin{equation*}
V_{I,k}/V_{I,k-1} \otimes I^{\otimes k} \ \simeq \ V^{I,k}/V^{I,k+1}
\end{equation*}
Les représentations $V$ de $I$ s'identifient, via $N$, aux $k[T]$-modules tués par une puissance de $T$. Il suffit de vérifier 3.2 lorsque $V$ correspond à un $k[T]$-module irréductible $k[T]/T^k$, ce qui est trivial.

}{\begin{equation*}
V_{I,k}/V_{I,k-1} \otimes I^{\otimes k} \ \simeq \ V^{I,k}/V^{I,k+1}
\end{equation*}
Les représentations $V$ de $I$ s'identifient, via $N$, aux $k[T]$-modules tués par une puissance de $T$. Il suffit de vérifier 5.3.2 lorsque $V$ correspond à un $k[T]$-module indécomposable $k[T]/T^k$, ce qui est trivial.

}{F7II-XXI-038}Achevons la démonstration de 5.3.1. Le sorite 5.3.2 fournit des filtrations croissantes et décroissantes $V_{I,k}$ et $V^{I,k}$ de $V_I$ et $V^I$, et des isomorphismes de $G/I$-modules

\begin{equation*}
V^{I,k}/V^{I,k+1} \ \simeq \ V_{I,k}/V_{I,k-1}(k) \quad ,
\end{equation*}
d'où les assertions.


\SGApage{II}{391}

\SGAvar{\underline{Théorème} 5.4. \underline{Soient} $X$ \underline{un schéma de type fini sur} $\mathbb{F}_q$, \underline{séparé et de dimension} $n$, \underline{et} $\underline{G}$ \underline{un} $E$-\underline{faisceau constructible sur} $X$. \underline{On suppose que pour tout} $x \in X$, \underline{l'endomorphisme} $F_x$ \underline{de} $\underline{G}_x^\vee$ \underline{est} $T$-\underline{entier}. \underline{Alors l'endomorphisme} $q^i F^\vee$ \underline{de} $\underline{H}_c^i(X,G)^\vee$ \underline{est} $T$-\underline{entier}, \underline{ainsi que l'endomorphisme} $q^n F^\vee$ \underline{de} $H_c^i(X,G)^\vee$.

}{\underline{Théorème} 5.4. \underline{Soient} $X$ \underline{un schéma de type fini sur} $\mathbb{F}_q$, \underline{séparé et de dimension} $n$, \underline{et} $\underline{G}$ \underline{un} $E$-\underline{faisceau constructible sur} $X$. \underline{On suppose que pour tout} $x \in X$, \underline{l'endomorphisme} $F_x$ \underline{de} $\underline{G}_x^\vee$ \underline{est} $T$-\underline{entier}. \underline{Alors l'endomorphisme} $q^i F^\vee$ \underline{de} $\underline{H}_c^i(X,G)^\vee$ \underline{est} $T$-\underline{entier}, \underline{ainsi que l'endomorphisme} $q^n F^\vee$ \underline{de} $\underline{H}_c^i(X,G)^\vee$.

}{F7II-XXI-039}On procède comme dans le théorème 5.2 pour se ramener au cas où $X$ est une courbe lisse et où $\underline{G}$ est constant tordu sur $X$. Soit $j$ l'inclusion de $X$ dans la courbe projective et lisse $X'$ qui compactifie $X$.

On sait (dualité de Poincaré) que $\underline{H}^i(X',j_*G)$ est dual de $\underline{H}^{2-i}(X',j_* G^\vee(1))$. En vertu du corollaire 5.3 (a), le théorème 5.2 s'applique à $j_*G^\vee$ :

\begin{align*}
\underline{H}^0(X',j_*G^\vee) &\ \simeq \ \underline{H}^2(X',j_*G)^\vee(-1) && \text{est \(T\)-entier,} \\
\underline{H}^1(X',j_*G^\vee) &\ \simeq \ \underline{H}^1(X',j_*G)^\vee(-1) && \text{est \(T\)-entier,} \\
\underline{H}^2(X',j_*G^\vee)(1) &\simeq \underline{H}^0(X',j_*G)^\vee && \text{est \(T\)-entier,}
\end{align*}
ce qui prouve le théorème pour $j_*G$. Définissons $H$ par la suite exacte

\begin{equation*}
0 \longrightarrow j_!G \longrightarrow j_*G \longrightarrow H \longrightarrow 0 \quad .
\end{equation*}
La suite exacte de cohomologie fournit, si $Y$ est le support de $H$,

\begin{center}
\begin{tikzcd}[row sep=large, column sep=large]
 {} \& \underline{H}^0_c(X,G) \arrow[r, hook] \& \underline{H}^0(X',j_*G) \& , \\
\underline{H}^0(Y,H) \arrow[r] \& \underline{H}^1_c(X,G) \arrow[r] \& \underline{H}^1(X',j_*G) \& , \\
 {} \& \underline{H}^2_c(X,G) \arrow[r, "\sim"] \& \underline{H}^2(X',j_*G) \& ,
\end{tikzcd}
\end{center}















\SGApage{II}{392}
de sorte que le théorème pour $(X,G)$ résulte du théorème pour $(X',j_*G)$.

5.5. La méthode précédente de fibration en courbes permet aussi d'obtenir des majorations à la Lang-Weil [3] sur les valeurs absolues complexes des valeurs propres de Frobenius.

Soient $X$ un schéma de type fini sur $\mathbb{F}_q$, $\underline{G}$ un E-faisceau constructible sur $X$, et $\lambda \in \mathbb{R}$. Si, pour tout point $x \in X$, avec $\#\, k(x) = q_x$, les valeurs propres de $F_x$ agissant sur $\underline{G}_x$ sont des nombres algébriques $\alpha$ dont toutes les valeurs absolues complexes vérifient

\begin{equation*}
|\alpha| \;\leqslant\; q_x^{\lambda} \qquad ,
\end{equation*}
on écrira

\begin{equation*}
\log |\underline{G}| \;\leqslant\; \lambda \qquad .
\end{equation*}
\underline{Lemme} 5.5.1. \underline{Soit} $X$ \underline{une courbe sur} $\mathbb{F}_q$, \underline{par quoi on entend ici un schéma} $X$ \underline{de type fini sur} $\mathbb{F}_q$, \underline{séparé et de dimension} $\leqslant 1$. \underline{Soit} $\underline{G}$ \underline{un} $E$-\underline{faisceau constructible sur} $X$ \underline{tel que} $\log |\underline{G}| \leqslant \lambda$ \underline{avec} $\lambda \in \mathbb{R}$. \underline{Alors les valeurs absolues complexes des valeurs propres} $\alpha$ \underline{de} $F$ \underline{agissant sur} $\underline{H}^1_c(X,G)$ \underline{vérifient} $|\alpha| \leqslant q^{\lambda+1}$.

Soit $\sigma$ un plongement complexe de E, ou simplement un plongement complexe de la fermeture algébrique de $\mathbb{Q}$ dans E.

Il suffit de vérifier que, quel que soit $\sigma$, les zéros $\beta$ de la fraction rationnelle $\sigma Z(X,\underline{G},t)$ vérifient $|\beta^{-1}| \leqslant q^{\lambda+1}$. On sait déjà en
\SGApage{II}{393}
effet par 5.2.1 si $\alpha$ est une valeur propre de Frobenius agissant sur $\underline{H}^0_c(X,\underline{G})$ ou sur $\underline{H}^2_c(X,\underline{G})$, on a $|\sigma \alpha| \leqslant q^{\lambda+1}$.

Vérifions que la série

\begin{equation*}
\log \sigma Z(X,\underline{G},t) \;=\; \sum_k \frac{t^k}{k} . \sum_{x \in X(\mathbb{F}_{q^k})} \sigma \; \mathrm{Tr}(F,x^*\underline{G})
\end{equation*}
est convergente pour $|t| < q^{-\lambda-1}$ . Représentant $X$ comme revêtement ramifié de $\mathbb{P}^1$, on obtient trivialement une majoration

\begin{equation*}
\# \; X(\mathbb{F}_{q^k}) \;\leqslant\; A.q^k \qquad .
\end{equation*}
Par hypothèse, on a d'autre part

\SGAvar{\begin{equation*}
\left| \; \sigma \; \mathrm{Tr}(F,x^*G) \right| \;\leqslant\; B.(q^k)^{\lambda} \qquad .
\end{equation*}
Le coefficient $c_k$ de $t^k$ dans $\sigma Z (X,\underline{G},t)$ vérifie donc

}{\begin{equation*}
\left| \; \sigma \; \mathrm{Tr}(F,x^*G) \right| \;\leqslant\; B.(q^k)^{\lambda} \qquad .
\end{equation*}
Le coefficient $c_k$ de $t^k$ dans $\log \sigma Z (X,\underline{G},t)$ vérifie donc

}{F7II-XXI-040}\begin{equation*}
\left| \, c_k \right| \;\leqslant\; \frac{1}{k} \;.\; C. \; (q^{1+\lambda})^k
\end{equation*}
et l'assertion en résulte.

\underline{Théorème} 5.5.2. \underline{Soient} $X$ \underline{un schéma de type fini sur} $\mathbb{F}_q$, \underline{séparé et de dimension} $\leqslant n$, \underline{et} $G$ \underline{un} $E$-\underline{faisceau constructible sur} $X$ \underline{tel que} $\log|G| \leqslant \lambda$ \underline{avec} $\lambda \in \mathbb{R}$ .

\SGAvar{\underline{Alors, les valeurs absolues complexes des valeurs propres} $\alpha$ \underline{de Frobenius agissant sur} $H^i_c(X,\underline{G})$ \underline{vérifient}

}{\underline{Alors, les valeurs absolues complexes des valeurs propres} $\alpha$ \underline{de Frobenius agissant sur} $\underline{H}^i_c(X,\underline{G})$ \underline{vérifient}

}{F7II-XXI-041}\begin{equation*}
|\alpha| \;\leqslant\; q^{\lambda+i} \qquad \underline{\mbox{et}}
\end{equation*}
\begin{equation*}
|\alpha| \;\leqslant\; q^{\lambda+n} \qquad .
\end{equation*}

\SGApage{II}{394}

\SGAvar{Si $n \leq 0$, 5.5.2 est trivial. Si $n = 1$, les cas $i = 0$ et $i = 2$ de 5.5.2 résultent de 5.2.1, tandis que le cas $i = 1$ résulte de 5.5.1. On procède alors par récurrence sur $n$, comme dans la démonstration de 5.5.2.

}{Si $n \leq 0$, 5.5.2 est trivial. Si $n = 1$, les cas $i = 0$ et $i = 2$ de 5.5.2 résultent de 5.2.1, tandis que le cas $i = 1$ résulte de 5.5.1. On procède alors par récurrence sur $n$, comme dans la démonstration de 5.2.2.

}{F7II-XXI-042}\underline{Corollaire} 5.5.3. \underline{Soit} $X$ \underline{un schéma séparé de type fini de dimension} $n$ \underline{sur} $\mathbb{F}_q$ . \underline{Alors, les valeurs propres} $\alpha$ \underline{de l'endomorphisme de Frobenius géométrique} $F$ \underline{de} $\underline{H}^i_c(X,\mathbb{Q}_\ell)$ \underline{sont des entiers algébriques}. \underline{De plus}

(i)\quad \underline{Pour} $\ell' \neq p$, $\alpha$ \underline{est une unité} $\ell'$\underline{-adique}.

(ii)\quad \underline{Pour} $i > n$, $\alpha$ \underline{est divisible par} $q^{i-n}$.

(iii)\quad \underline{Pour} $i \leq n$, (resp. \underline{pour} $i \geq n$), $q^i \alpha^{-1}$ (resp. $q^n \, \alpha^{-1}$) \underline{est un entier algébrique}.

(iv)\quad \underline{Pour} $0 < i$, \underline{les valeurs absolues complexes de} $\alpha$ \underline{vérifient}

\begin{equation*}
|\alpha| \leq q^{i-\frac{1}{2}} \qquad \mathrm{et} \qquad |\alpha| \leq q^n \quad .
\end{equation*}
D'après 5.2.2, $\alpha$ et $q^{n-i} \, \alpha$ sont des entiers algébriques. Les assertions (i) à (iii) résultent alors de 5.4. Pour $i > n$, (iv) résulte de 5.5.2. Pour obtenir (iv) quand $i \leq n$, il va falloir conjuguer 5.5.2 à l'hypothèse de Riemann pour les courbes. Procédons par récurrence sur $n$, le cas $n = 0$ étant trivial. Si $U$ est un ouvert de $X$, de complémentaire $Z$ de dimension $< n$, la suite exacte

\SGAvar{\begin{equation*}
\underline{H}^i_c(U,\mathbb{Q}_\ell) \longrightarrow \underline{H}^i_c(X,\mathbb{Q}_\ell) \longrightarrow \underline{H}^i_c(Z,\mathbb{Q}_\ell)
\end{equation*}
montre qu'il suffit de prouver 5.2.2 pour $U$. Quitte à remplacer au préalable $X$ par $X_{\mathrm{red}}$, ceci nous ramène au cas où $X$ est lisse affine,
\SGApage{II}{395}
et tel qu'il existe un morphisme lisse $f : X \longrightarrow S$ vérifiant les conditions suivantes :

}{\begin{equation*}
\underline{H}^i_c(U,\mathbb{Q}_\ell) \longrightarrow \underline{H}^i_c(X,\mathbb{Q}_\ell) \longrightarrow \underline{H}^i_c(Z,\mathbb{Q}_\ell)
\end{equation*}
montre qu'il suffit de prouver l'assertion (iv) pour $U$. Quitte à remplacer au préalable $X$ par $X_{\mathrm{red}}$, ceci nous ramène au cas où $X$ est lisse affine,
\SGApage{II}{395}
et tel qu'il existe un morphisme lisse $f : X \longrightarrow S$ vérifiant les conditions suivantes :

}{F7II-XXI-043}(a)\quad $\dim(S) \leq n-1$

(b)\quad les fibres de $f$ sont purement de dimension 1.

Puisque les fibres de $f$ sont des courbes affines, on a

\begin{equation*}
R^0 f_! \; \mathbb{Q}_\ell = 0 \quad .
\end{equation*}
D'après Weil [4], on a

\begin{equation}\tag{5.5.3.1}
\log \, \bigl| R^1 f_! \; \mathbb{Q}_\ell \bigr| \leq 1/2 \quad ,
\end{equation}
et d'après 5.2.1 appliqué aux fibres de $f$, on a

\begin{equation}\tag{5.5.3.2}
\log \, \bigl| R^2 f_! \; \mathbb{Q}_\ell \bigr| \leq 1 \quad .
\end{equation}
La suite spectrale de Leray pour $f$ se réduit ici à une suite exacte longue

\begin{equation*}
\underline{H}^{i-2}_c(S,R^2 f_! \mathbb{Q}_\ell) \longrightarrow \underline{H}^i_c(X,\mathbb{Q}_\ell) \longrightarrow \underline{H}^{i-1}_c(S,R^1 f_! \mathbb{Q}_\ell) \qquad .
\end{equation*}
D'après (5.5.3.1) (resp. 5.5.3.2) et 5.5.2, les valeurs absolues complexes des valeurs propres $\alpha$ de Frobenius agissant sur $\underline{H}^{i-1}_c(S,R^1 f_! \mathbb{Q}_\ell)$ (resp. sur $\underline{H}^{i-2}_c(S,R^2 f_! \mathbb{Q}_\ell)$) vérifient

\SGAvar{\begin{equation*}
|\alpha| \leq q^{i-1+\frac{1}{2}} \quad = q^{i-\frac{1}{2}}
\end{equation*}
\begin{equation*}
(\mathrm{resp.} \; |\alpha| \leq q^{i-2+1} \quad = q^{i-1} \leq q^{i-\frac{1}{2}} \qquad .
\end{equation*}
Le corollaire en résulte.


}{\begin{equation*}
|\alpha| \leq q^{i-1+\frac{1}{2}} \quad = q^{i-\frac{1}{2}}
\end{equation*}
\begin{equation*}
(\mathrm{resp.} \; |\alpha| \leq q^{i-2+1} \quad = q^{i-1} \leq q^{i-\frac{1}{2}}) \qquad .
\end{equation*}
Le corollaire en résulte.


}{F7II-XXI-044}\SGApage{II}{396}

Pour obtenir quelques résultats analogues aux précédents, en cohomologie sans supports, il semble nécessaire de disposer de la résolution des singularités. En dimension $\leq 1$, la solution du problème est fournie par le corollaire 5.3.

\underline{Théorème} 5.6. \underline{Soit} $N$ \underline{un entier, et supposons que l'on dispose de la résolution des singularités, sous la forme forte de Hironaka, pour les schémas de dimension} $\leq N$ \underline{de type fini sur} $\mathbb{F}_q$ (\underline{jusqu'ici, on peut donc prendre} $N = 2$ [1]). \underline{Soient} $f : X \longrightarrow Y$ \underline{un morphisme de schémas de type fini sur} $\mathbb{F}_q$ \underline{et} $\underline{G}$ \underline{un} $E$\underline{-faisceau constructible} $T$\underline{-entier sur} $X$. \underline{Alors, si} $\dim(X) \leq N$, \underline{les} $E$\underline{-faisceaux} $R^i f_*\underline{G}$ \underline{sur} $Y$ \underline{sont} $T$\underline{-entiers.}

\SGAvar{La démonstration procède par récurrence sur la dimension $n \leq N$ de $X$. Soit $j : U \longrightarrow X$ un ouvert de $X$, lisse, affine, tel que $\underline{G}$ restreint à $U$ soit constant tordu et tel que la dimension de son complément $Z$ soit $< n$. Utilisant l'hypothèse de résolution, on sait que les $R^i j_*(j^*\underline{G})$ sont constructibles (SGA 5 I \quad , [2]); pour $i > 0$, ils sont concentrés sur $Z$ et, pour $i = 0$, $j_*(j^*G)$ ne différe de $\underline{G}$ que sur $Z$. Il en résulte, d'après l'hypothèse de récurrence, qu'il suffit de vérifier le théorème pour le faisceau $j^*G$ et les morphisme $j$ et $f\,j$. Il suffit de vérifier le théorème sous l'hypothèse additionnelle que $X$ soit affine lisse et $\underline{G}$ constant tordu (i.e. fourni par une représentation continue du groupe fondamental de $X$ dans un E-vectoriel de dimension finie).


}{La démonstration procède par récurrence sur la dimension $n \leq N$ de $X$. Soit $j : U \longrightarrow X$ un ouvert de $X$, lisse, affine, tel que $\underline{G}$ restreint à $U$ soit constant tordu et tel que la dimension de son complément $Z$ soit $< n$. Utilisant l'hypothèse de résolution, on sait que les $R^i j_*(j^*\underline{G})$ sont constructibles (SGA 5 I \quad , [2]); pour $i > 0$, ils sont concentrés sur $Z$ et, pour $i = 0$, $j_*(j^*G)$ ne diffère de $\underline{G}$ que sur $Z$. Il en résulte, d'après l'hypothèse de récurrence, qu'il suffit de vérifier le théorème pour le faisceau $j^*G$ et les morphismes $j$ et $f\circ j$. Il suffit de vérifier le théorème sous l'hypothèse additionnelle que $X$ soit affine lisse et $\underline{G}$ constant tordu (i.e. fourni par une représentation continue du groupe fondamental de $X$ dans un E-vectoriel de dimension finie).


}{F7II-XXI-045}\SGApage{II}{397}

\SGAvar{Soient $R$ l'anneau des entiers de $E$, $\lambda$ une uniformisante, $e$ l'indice de ramification absolu et $b$ un entier $> \frac{e}{\ell-1}$ . Puisque $X$ est lisse, donc normal et à groupe fondamental profini, il existe un R-faisceau $\underline{G}^o$ constant tordu tel que $\underline{G} \simeq \underline{G}^o \otimes_R E$. Soit $p : X' \longrightarrow X$ le revêtement étale de $X$ qui trivialise le faisceau localement constant $\underline{G}^o \otimes_R R/\lambda^b$. Le faisceau $\underline{G}$ est facteur direct de $p_*p^*\underline{G}$, de sorte qu'il suffit de vérifier le théorème pour $p^*\underline{G}$ et $f \circ p$. Il reste donc à traiter le cas où les hypothèses suivantes sont vérifiées. :

}{Soient $R$ l'anneau des entiers de $E$, $\lambda$ une uniformisante, $e$ l'indice de ramification absolu et $b$ un entier $> \frac{e}{\ell-1}$ . Puisque $X$ est lisse, donc normal et à groupe fondamental profini, il existe un R-faisceau $\underline{G}^o$ constant tordu tel que $\underline{G} \simeq \underline{G}^o \otimes_R E$. Soit $p : X' \longrightarrow X$ le revêtement étale de $X$ qui trivialise le faisceau localement constant $\underline{G}^o \otimes_R R/\lambda^b$. Le faisceau $\underline{G}$ est facteur direct de $p_*p^*\underline{G}$, de sorte qu'il suffit de vérifier le théorème pour $p^*\underline{G}$ et $f \circ p$. Il reste donc à traiter le cas où les hypothèses suivantes sont vérifiées :

}{F7II-XXI-046}(*)  $X$ est affine et lisse, et il existe un R-faisceau $\underline{G}^o$, constant tordu, tel que $\underline{G} \simeq \underline{G}^o \otimes_R E$ et que le faisceau $\underline{G}^o \otimes_R R/\lambda^b$ soit constant .

Factorisons le morphisme $f$ en un morphisme propre et une immersion ouverte $f = \bar{f} j$ :

\begin{center}
\begin{tikzcd}
X \arrow[r, hook, "j", swap] \arrow[d, "f"'] \& X' \arrow[dl, "\bar{f}"] \\
Y \&
{} \end{tikzcd}
\end{center}










.

Le théorème a déjà été établi pour $\bar{f}$ (5.2), de sorte qu'il suffit de l'établir pour $j$. D'après la résolution des singularités, appliquée à $X'$, on pouvait supposer que $X'$ soit lisse et que $X$ soit le complément d'un diviseur à croisements normaux de $X'$, réunion de diviseurs lisses. L'hypothèse (*) implique que le faisceau $\underline{G}$ est modérément ramifié le long de ce diviseur. On a alors :


\SGApage{II}{398}

\underline{Lemme} 5.6.1. \underline{Soient} $X$ \underline{un schéma lisse de type fini sur} $\mathbb{F}_q$, $D$ \underline{un diviseur à croisements normaux dans} $X$, \underline{réunion de diviseurs lisses} $D_i$, $j : U \hookrightarrow X$ \underline{le complément de} $D$ \underline{dans} $X$ \underline{et} $\underline{G}$ \underline{un} $E$\underline{-faisceau constant tordu sur} $U$, \underline{modérément ramifié le long de} $D$. \underline{Alors, si} $\underline{G}$ \underline{est} $T$\underline{-entier, les faisceaux} $R^i j_*\underline{G}$ \underline{sont} $T$\underline{-entiers.}

On admettra (cf. SGA 1 XIII\hspace{2.5em}) :

\underline{Lemme} 5.6.2. (i) \underline{Sous les hypothèses} 5.6.1, \underline{soient} $D'$ \underline{une composante de} $D$, $D''$ \underline{la réunion des autres composantes,} $U' = X-D''$ \underline{et} $D^{\prime o} = D'\cap U'$ :

\begin{center}
\begin{tikzcd}
{} \& D^{\prime o} \arrow[r, hook, "j'_D", swap] \arrow[d, "i^{o}"] \& D' \arrow[d, "i"] \\
U \arrow[r, hook, "j_1", swap] \& U' \arrow[r, hook, "j'", swap] \& X
\end{tikzcd}
\end{center}










.

\underline{Alors, la flèche de changement de base}

\SGAvar{\begin{equation*}
i^* Rj'_* (Rj_{1*} \underline{G}) \longrightarrow Rj'_{D*}\, i^{o*} (Rj_{1*}\underline{G})
\end{equation*}
\underline{est un isomorphisme, et les faisceaux} $i^{o*}Rj_{1*}\underline{G}$ \underline{sont constant tordus, modérément ramifiés le long de} $D' - D^{\prime o}$.

}{\begin{equation*}
i^* Rj'_* (Rj_{1*} \underline{G}) \longrightarrow Rj'_{D*}\, i^{o*} (Rj_{1*}\underline{G})
\end{equation*}
\underline{est un isomorphisme, et les faisceaux} $i^{o*}Rj_{1*}\underline{G}$ \underline{sont constants tordus, modérément ramifiés le long de} $D' - D^{\prime o}$.

}{F7II-XXI-047}(ii) \underline{Si} $C$ \underline{est une courbe lisse sur} $X$, \underline{transversale à} $D$ \underline{en un point lisse} $x$ \underline{de} $D$,

\begin{center}
\begin{tikzcd}
C \cap U \arrow[r, "j_c", swap] \arrow[d, "k'"] \& C \arrow[d, "k"] \\
U \arrow[r, "j", swap] \& X
\end{tikzcd}
\end{center}









,
\SGApage{II}{399}
\underline{alors, la flèche de changement de base}

\begin{equation*}
k^* Rj_*\underline{G} \longrightarrow Rj_{c*}\, k^{\prime *}\underline{G}
\end{equation*}
\underline{est un isomorphisme en} $x$.

\SGAvar{Déduisons 5.6.1 de 5.6.2 en raisonnant par récurrence sur la dimension, et en admettant provisoirement 5.6.1 dans le cas des courbes. Il résulte de 5.6.2 (ii) et de cette hypothèse, que les faisceaux $R^i j_*\underline{G}$ sont $T$-entiers en tout point lisse de $D$. En particulier, avec les notations de 5.6.2 (i), les faisceaux $R^i j_{1*}\underline{G}$ sont $T$-entiers, de sorte que, d'après 6.2 (i) et l'hypothèse de récurrence appliquée à $i^{o*} Rj_{1*}\underline{G}$ et à $j'_D$ , les faisceaux $i^* R^i j_*\underline{G}$ sont $T$-entiers, quelle que soit la composante choisie $D'$ de $D$. Ceci prouve la conclusion de 5.6.1.

}{Déduisons 5.6.1 de 5.6.2 en raisonnant par récurrence sur la dimension, et en admettant provisoirement 5.6.1 dans le cas des courbes. Il résulte de 5.6.2 (ii) et de cette hypothèse, que les faisceaux $R^i j_*\underline{G}$ sont $T$-entiers en tout point lisse de $D$. En particulier, avec les notations de 5.6.2 (i), les faisceaux $R^i j_{1*}\underline{G}$ sont $T$-entiers, de sorte que, d'après 5.6.2 (i) et l'hypothèse de récurrence appliquée à $i^{o*} Rj_{1*}\underline{G}$ et à $j'_D$ , les faisceaux $i^* R^i j_*\underline{G}$ sont $T$-entiers, quelle que soit la composante choisie $D'$ de $D$. Ceci prouve la conclusion de 5.6.1.

}{F7II-XXI-048}Reste le cas des courbes. Pour une courbe, sous les hypothèses 5.6.1, si $x$ est un point de $X - U$, si $\mathfrak{G}$ est le groupe de Galois du corps des fractions de $\hat{\mathcal{O}}_{x,X}$ et $\mathfrak{I}$ le groupe d'inertie, si on désigne encore par $\underline{G}$ la représentation de $\mathfrak{G}$ définie par $\underline{G}$ et si $\bar{x}$ est un point géométrique de $X$ localisé en $x$, on a des isomorphismes de $\mathfrak{G}/\mathfrak{I}$-modules galoisiens :

\begin{equation*}
(j_*\underline{G})_{\bar{x}} \;\simeq\; \underline{G}^{\mathfrak{I}}
\end{equation*}
\begin{equation*}
(R^1 j_*\underline{G})_{\bar{x}} \;\simeq\; \underline{G}_{\mathfrak{I}}(-1)
\end{equation*}
\begin{equation*}
(R^i j_*\underline{G})_{\bar{x}} = 0 \ \text{si} \ i \geq 2
\end{equation*}
et l'assertion résulte de 5.3 (a) et 5.3.1.

\SGApage{II}{400}

\underline{BIBLIOGRAPHIE}

[1]  S.S. ABHYANKAR.  Resolution of singularities of embedded algebraic surfaces. Acad. Press. New York 1966.

[2]  J.P. JOUANOLOU.  Cohomologie $\ell$-adique des schémas. A paraître dans Lecture Notes of Math.

[3]  S. LANG and A. WEIL.  Number of points of varieties in finite fields. Am. J. Math. 76 (1954) 819-827.

[4]  A. WEIL.  Sur les courbes algébriques et les variétés qui s'en déduisent. Act. Sci. Ind 1041. Hermann 1948.
\selectlanguage{french}

\ifnum\SGAform=1
\clearpage\section*{Apparat critique}\small
Les lectures ci-dessous sont données dans le TeX exact des témoins, afin de préserver les distinctions de notation. Les indices décimaux renvoient aux différences individuelles du passage marqué.

\par\medskip\noindent\hypertarget{app-F7II-XXI-001.1}{}\textbf{F7II-XXI-001.1}\quad p. 364; printed misprint.\par
\noindent Imprimé : {\ttfamily\footnotesize et\allowbreak{}}\par
\noindent Corrigé : {\ttfamily\footnotesize est\allowbreak{}}\par
\noindent Source : 7II:651. \nolinkurl{fixes/MISPRINTS_7II.json#651}\par

\par\medskip\noindent\hypertarget{app-F7II-XXI-001.2}{}\textbf{F7II-XXI-001.2}\quad p. 364; designated-witness attribution.\par
\noindent Imprimé : {\ttfamily\footnotesize ,\allowbreak{}}\par
\noindent Corrigé : \emph{(vide)}\par
\noindent Source : WITNESS-F7II-XXI-001.2. \nolinkurl{fr/sga7ii/diplomatic/expose_XXI_body.tex}; \nolinkurl{fr/sga7ii/current/expose_XXI_body.tex}\par

\par\medskip\noindent\hypertarget{app-F7II-XXI-001.3}{}\textbf{F7II-XXI-001.3}\quad p. 364; printed misprint.\par
\noindent Imprimé : {\ttfamily\footnotesize suffisament\allowbreak{}}\par
\noindent Corrigé : {\ttfamily\footnotesize suffisamment\allowbreak{}}\par
\noindent Source : 7II:652. \nolinkurl{fixes/MISPRINTS_7II.json#652}\par

\par\medskip\noindent\hypertarget{app-F7II-XXI-001.4}{}\textbf{F7II-XXI-001.4}\quad p. 364; designated-witness attribution.\par
\noindent Imprimé : {\ttfamily\footnotesize tout\allowbreak{}}\par
\noindent Corrigé : {\ttfamily\footnotesize toute\allowbreak{} la\allowbreak{} cohomologie\allowbreak{}}\par
\noindent Source : WITNESS-F7II-XXI-001.4. \nolinkurl{fr/sga7ii/diplomatic/expose_XXI_body.tex}; \nolinkurl{fr/sga7ii/current/expose_XXI_body.tex}\par

\par\medskip\noindent\hypertarget{app-F7II-XXI-001.5}{}\textbf{F7II-XXI-001.5}\quad p. 364; designated-witness attribution.\par
\noindent Imprimé : {\ttfamily\footnotesize du\allowbreak{}}\par
\noindent Corrigé : {\ttfamily\footnotesize des\allowbreak{}}\par
\noindent Source : WITNESS-F7II-XXI-001.5. \nolinkurl{fr/sga7ii/diplomatic/expose_XXI_body.tex}; \nolinkurl{fr/sga7ii/current/expose_XXI_body.tex}\par

\par\medskip\noindent\hypertarget{app-F7II-XXI-001.6}{}\textbf{F7II-XXI-001.6}\quad p. 364; inherited corrected-layer difference.\par
\noindent Imprimé : \emph{(vide)}\par
\noindent Corrigé : {\ttfamily\footnotesize \$\allowbreak{} }\par
\noindent Source : ADOPT-F7II-XXI-001.6. \nolinkurl{fr/sga7ii/diplomatic/expose_XXI_body.tex}; \nolinkurl{fr/sga7ii/current/expose_XXI_body.tex}\par

\par\medskip\noindent\hypertarget{app-F7II-XXI-001.7}{}\textbf{F7II-XXI-001.7}\quad p. 364; inherited corrected-layer difference.\par
\noindent Imprimé : \emph{(vide)}\par
\noindent Corrigé : {\ttfamily\footnotesize \$\allowbreak{}}\par
\noindent Source : ADOPT-F7II-XXI-001.7. \nolinkurl{fr/sga7ii/diplomatic/expose_XXI_body.tex}; \nolinkurl{fr/sga7ii/current/expose_XXI_body.tex}\par

\par\medskip\noindent\hypertarget{app-F7II-XXI-001.8}{}\textbf{F7II-XXI-001.8}\quad p. 364; inherited corrected-layer difference.\par
\noindent Imprimé : \emph{(vide)}\par
\noindent Corrigé : {\ttfamily\footnotesize \$\allowbreak{}}\par
\noindent Source : ADOPT-F7II-XXI-001.8. \nolinkurl{fr/sga7ii/diplomatic/expose_XXI_body.tex}; \nolinkurl{fr/sga7ii/current/expose_XXI_body.tex}\par

\par\medskip\noindent\hypertarget{app-F7II-XXI-001.9}{}\textbf{F7II-XXI-001.9}\quad p. 364; inherited corrected-layer difference.\par
\noindent Imprimé : {\ttfamily\footnotesize \$\allowbreak{}}\par
\noindent Corrigé : \emph{(vide)}\par
\noindent Source : ADOPT-F7II-XXI-001.9. \nolinkurl{fr/sga7ii/diplomatic/expose_XXI_body.tex}; \nolinkurl{fr/sga7ii/current/expose_XXI_body.tex}\par

\par\medskip\noindent\hypertarget{app-F7II-XXI-001.10}{}\textbf{F7II-XXI-001.10}\quad p. 364; inherited corrected-layer difference.\par
\noindent Imprimé : \emph{(vide)}\par
\noindent Corrigé : {\ttfamily\footnotesize espace\allowbreak{} }\par
\noindent Source : ADOPT-F7II-XXI-001.10. \nolinkurl{fr/sga7ii/diplomatic/expose_XXI_body.tex}; \nolinkurl{fr/sga7ii/current/expose_XXI_body.tex}\par

\par\medskip\noindent\hypertarget{app-F7II-XXI-001.11}{}\textbf{F7II-XXI-001.11}\quad p. 364; inherited corrected-layer difference.\par
\noindent Imprimé : {\ttfamily\footnotesize car.\allowbreak{}}\par
\noindent Corrigé : {\ttfamily\footnotesize caractéristique\allowbreak{}}\par
\noindent Source : ADOPT-F7II-XXI-001.11. \nolinkurl{fr/sga7ii/diplomatic/expose_XXI_body.tex}; \nolinkurl{fr/sga7ii/current/expose_XXI_body.tex}\par

\par\medskip\noindent\hypertarget{app-F7II-XXI-001.12}{}\textbf{F7II-XXI-001.12}\quad p. 365; printed misprint.\par
\noindent Imprimé : {\ttfamily\footnotesize dûs\allowbreak{}}\par
\noindent Corrigé : {\ttfamily\footnotesize dus\allowbreak{}}\par
\noindent Source : 7II:1059. \nolinkurl{fixes/MISPRINTS_7II.json#1059}\par

\par\medskip\noindent\hypertarget{app-F7II-XXI-001.13}{}\textbf{F7II-XXI-001.13}\quad p. 365; inherited corrected-layer difference.\par
\noindent Imprimé : {\ttfamily\footnotesize ,\allowbreak{}}\par
\noindent Corrigé : \emph{(vide)}\par
\noindent Source : ADOPT-F7II-XXI-001.13. \nolinkurl{fr/sga7ii/diplomatic/expose_XXI_body.tex}; \nolinkurl{fr/sga7ii/current/expose_XXI_body.tex}\par

\par\medskip\noindent\hypertarget{app-F7II-XXI-001.14}{}\textbf{F7II-XXI-001.14}\quad p. 365; inherited corrected-layer difference.\par
\noindent Imprimé : {\ttfamily\footnotesize qui\allowbreak{}}\par
\noindent Corrigé : {\ttfamily\footnotesize prouvé\allowbreak{}}\par
\noindent Source : ADOPT-F7II-XXI-001.14. \nolinkurl{fr/sga7ii/diplomatic/expose_XXI_body.tex}; \nolinkurl{fr/sga7ii/current/expose_XXI_body.tex}\par

\par\medskip\noindent\hypertarget{app-F7II-XXI-001.15}{}\textbf{F7II-XXI-001.15}\quad p. 365; inherited corrected-layer difference.\par
\noindent Imprimé : {\ttfamily\footnotesize est\allowbreak{} prouvé\allowbreak{}}\par
\noindent Corrigé : {\ttfamily\footnotesize ici\allowbreak{}}\par
\noindent Source : ADOPT-F7II-XXI-001.15. \nolinkurl{fr/sga7ii/diplomatic/expose_XXI_body.tex}; \nolinkurl{fr/sga7ii/current/expose_XXI_body.tex}\par

\par\medskip\noindent\hypertarget{app-F7II-XXI-001.16}{}\textbf{F7II-XXI-001.16}\quad p. 365; inherited corrected-layer difference.\par
\noindent Imprimé : {\ttfamily\footnotesize I\allowbreak{}}\par
\noindent Corrigé : {\ttfamily\footnotesize (\$\textbackslash{}S\$\allowbreak{} 5)\allowbreak{}}\par
\noindent Source : ADOPT-F7II-XXI-001.16. \nolinkurl{fr/sga7ii/diplomatic/expose_XXI_body.tex}; \nolinkurl{fr/sga7ii/current/expose_XXI_body.tex}\par

\par\medskip\noindent\hypertarget{app-F7II-XXI-002.1}{}\textbf{F7II-XXI-002.1}\quad p. 365; misprint\_pass.\par
\noindent Imprimé : \emph{(vide)}\par
\noindent Corrigé : {\ttfamily\footnotesize \}\allowbreak{}}\par
\noindent Source : 7II:654; 7II:653. \nolinkurl{fixes/FIXLOG_7II_CURRENT.jsonl#766}; \nolinkurl{fixes/MISPRINTS_7II.json#653}\par

\par\medskip\noindent\hypertarget{app-F7II-XXI-002.2}{}\textbf{F7II-XXI-002.2}\quad p. 365; misprint\_pass.\par
\noindent Imprimé : \emph{(vide)}\par
\noindent Corrigé : {\ttfamily\footnotesize \$X\$\allowbreak{} \textbackslash{}underline\{\allowbreak{}}\par
\noindent Source : 7II:654; 7II:653. \nolinkurl{fixes/FIXLOG_7II_CURRENT.jsonl#766}; \nolinkurl{fixes/MISPRINTS_7II.json#653}\par

\par\medskip\noindent\hypertarget{app-F7II-XXI-003.1}{}\textbf{F7II-XXI-003.1}\quad p. 366; misprint\_pass.\par
\noindent Imprimé : \emph{(vide)}\par
\noindent Corrigé : {\ttfamily\footnotesize )\allowbreak{}}\par
\noindent Source : 7II:1082. \nolinkurl{fixes/FIXLOG_7II_CURRENT.jsonl#892}; \nolinkurl{fixes/MISPRINTS_7II.json#1082}\par

\par\medskip\noindent\hypertarget{app-F7II-XXI-004.1}{}\textbf{F7II-XXI-004.1}\quad p. 366; printed misprint.\par
\noindent Imprimé : {\ttfamily\footnotesize toute\allowbreak{}}\par
\noindent Corrigé : {\ttfamily\footnotesize toutes\allowbreak{}}\par
\noindent Source : 7II:1083. \nolinkurl{fixes/MISPRINTS_7II.json#1083}\par

\par\medskip\noindent\hypertarget{app-F7II-XXI-005.1}{}\textbf{F7II-XXI-005.1}\quad p. 366; inherited corrected-layer difference.\par
\noindent Imprimé : \emph{(vide)}\par
\noindent Corrigé : \emph{(1 espace(s), 0 retour(s) de ligne dans le témoin TeX)}\par
\noindent Source : ADOPT-F7II-XXI-005.1. \nolinkurl{fr/sga7ii/diplomatic/expose_XXI_body.tex}; \nolinkurl{fr/sga7ii/current/expose_XXI_body.tex}\par

\par\medskip\noindent\hypertarget{app-F7II-XXI-006.1}{}\textbf{F7II-XXI-006.1}\quad p. 367; designated-witness attribution.\par
\noindent Imprimé : \emph{(vide)}\par
\noindent Corrigé : {\ttfamily\footnotesize ,\allowbreak{}}\par
\noindent Source : WITNESS-F7II-XXI-006.1. \nolinkurl{fr/sga7ii/diplomatic/expose_XXI_body.tex}; \nolinkurl{fr/sga7ii/current/expose_XXI_body.tex}\par

\par\medskip\noindent\hypertarget{app-F7II-XXI-006.2}{}\textbf{F7II-XXI-006.2}\quad p. 367; designated-witness attribution.\par
\noindent Imprimé : \emph{(vide)}\par
\noindent Corrigé : {\ttfamily\footnotesize  que\allowbreak{}}\par
\noindent Source : WITNESS-F7II-XXI-006.2. \nolinkurl{fr/sga7ii/diplomatic/expose_XXI_body.tex}; \nolinkurl{fr/sga7ii/current/expose_XXI_body.tex}\par

\par\medskip\noindent\hypertarget{app-F7II-XXI-006.3}{}\textbf{F7II-XXI-006.3}\quad p. 367; designated-witness attribution.\par
\noindent Imprimé : {\ttfamily\footnotesize  que\allowbreak{}}\par
\noindent Corrigé : \emph{(vide)}\par
\noindent Source : WITNESS-F7II-XXI-006.3. \nolinkurl{fr/sga7ii/diplomatic/expose_XXI_body.tex}; \nolinkurl{fr/sga7ii/current/expose_XXI_body.tex}\par

\par\medskip\noindent\hypertarget{app-F7II-XXI-006.4}{}\textbf{F7II-XXI-006.4}\quad p. 367; printed misprint.\par
\noindent Imprimé : {\ttfamily\footnotesize théroème\allowbreak{}}\par
\noindent Corrigé : {\ttfamily\footnotesize théorème\allowbreak{}}\par
\noindent Source : 7II:655; 7II:1084. \nolinkurl{fixes/MISPRINTS_7II.json#655}; \nolinkurl{fixes/MISPRINTS_7II.json#1084}\par

\par\medskip\noindent\hypertarget{app-F7II-XXI-006.5}{}\textbf{F7II-XXI-006.5}\quad p. 367; printed misprint.\par
\noindent Imprimé : {\ttfamily\footnotesize vérifiante\allowbreak{}}\par
\noindent Corrigé : {\ttfamily\footnotesize vérifiant\allowbreak{}}\par
\noindent Source : 7II:1085. \nolinkurl{fixes/MISPRINTS_7II.json#1085}\par

\par\medskip\noindent\hypertarget{app-F7II-XXI-007.1}{}\textbf{F7II-XXI-007.1}\quad p. 367; printed misprint.\par
\noindent Imprimé : {\ttfamily\footnotesize des\allowbreak{}}\par
\noindent Corrigé : {\ttfamily\footnotesize de\allowbreak{}}\par
\noindent Source : 7II:1086. \nolinkurl{fixes/MISPRINTS_7II.json#1086}\par

\par\medskip\noindent\hypertarget{app-F7II-XXI-008.1}{}\textbf{F7II-XXI-008.1}\quad p. 368; inherited corrected-layer difference.\par
\noindent Imprimé : \emph{(vide)}\par
\noindent Corrigé : {\ttfamily\footnotesize \{i=\allowbreak{}}\par
\noindent Source : ADOPT-F7II-XXI-008.1. \nolinkurl{fr/sga7ii/diplomatic/expose_XXI_body.tex}; \nolinkurl{fr/sga7ii/current/expose_XXI_body.tex}\par

\par\medskip\noindent\hypertarget{app-F7II-XXI-008.2}{}\textbf{F7II-XXI-008.2}\quad p. 368; inherited corrected-layer difference.\par
\noindent Imprimé : \emph{(vide)}\par
\noindent Corrigé : {\ttfamily\footnotesize \}\allowbreak{}}\par
\noindent Source : ADOPT-F7II-XXI-008.2. \nolinkurl{fr/sga7ii/diplomatic/expose_XXI_body.tex}; \nolinkurl{fr/sga7ii/current/expose_XXI_body.tex}\par

\par\medskip\noindent\hypertarget{app-F7II-XXI-009.1}{}\textbf{F7II-XXI-009.1}\quad p. 369; printed misprint.\par
\noindent Imprimé : {\ttfamily\footnotesize )\allowbreak{}}\par
\noindent Corrigé : \emph{(vide)}\par
\noindent Source : 7II:657; 7II:1087. \nolinkurl{fixes/MISPRINTS_7II.json#657}; \nolinkurl{fixes/MISPRINTS_7II.json#1087}\par

\par\medskip\noindent\hypertarget{app-F7II-XXI-009.2}{}\textbf{F7II-XXI-009.2}\quad p. 369; printed misprint.\par
\noindent Imprimé : {\ttfamily\footnotesize indépendant\allowbreak{}}\par
\noindent Corrigé : {\ttfamily\footnotesize indépendants\allowbreak{}}\par
\noindent Source : 7II:1088. \nolinkurl{fixes/MISPRINTS_7II.json#1088}\par

\par\medskip\noindent\hypertarget{app-F7II-XXI-010.1}{}\textbf{F7II-XXI-010.1}\quad p. 369; printed misprint.\par
\noindent Imprimé : \emph{(vide)}\par
\noindent Corrigé : {\ttfamily\footnotesize \{\allowbreak{}}\par
\noindent Source : 7II:1089. \nolinkurl{fixes/MISPRINTS_7II.json#1089}\par

\par\medskip\noindent\hypertarget{app-F7II-XXI-010.2}{}\textbf{F7II-XXI-010.2}\quad p. 369; printed misprint.\par
\noindent Imprimé : \emph{(vide)}\par
\noindent Corrigé : {\ttfamily\footnotesize =1\}\allowbreak{}}\par
\noindent Source : 7II:1089. \nolinkurl{fixes/MISPRINTS_7II.json#1089}\par

\par\medskip\noindent\hypertarget{app-F7II-XXI-010.3}{}\textbf{F7II-XXI-010.3}\quad p. 369; designated-witness attribution.\par
\noindent Imprimé : \emph{(vide)}\par
\noindent Corrigé : {\ttfamily\footnotesize en\allowbreak{} posant\allowbreak{} \$a\_0=0\$,\allowbreak{} }\par
\noindent Source : WITNESS-F7II-XXI-010.3. \nolinkurl{fr/sga7ii/diplomatic/expose_XXI_body.tex}; \nolinkurl{fr/sga7ii/current/expose_XXI_body.tex}\par

\par\medskip\noindent\hypertarget{app-F7II-XXI-011.1}{}\textbf{F7II-XXI-011.1}\quad p. 369; printed misprint.\par
\noindent Imprimé : {\ttfamily\footnotesize o\allowbreak{}}\par
\noindent Corrigé : {\ttfamily\footnotesize 1\allowbreak{}}\par
\noindent Source : 7II:1090. \nolinkurl{fixes/MISPRINTS_7II.json#1090}\par

\par\medskip\noindent\hypertarget{app-F7II-XXI-012.1}{}\textbf{F7II-XXI-012.1}\quad p. 370; designated-witness attribution.\par
\noindent Imprimé : {\ttfamily\footnotesize variable\allowbreak{} }\par
\noindent Corrigé : {\ttfamily\footnotesize variables\allowbreak{} }\par
\noindent Source : WITNESS-F7II-XXI-012.1. \nolinkurl{fr/sga7ii/diplomatic/expose_XXI_body.tex}; \nolinkurl{fr/sga7ii/current/expose_XXI_body.tex}\par

\par\medskip\noindent\hypertarget{app-F7II-XXI-012.2}{}\textbf{F7II-XXI-012.2}\quad p. 370; printed misprint.\par
\noindent Imprimé : \emph{(vide)}\par
\noindent Corrigé : {\ttfamily\footnotesize +1\allowbreak{}}\par
\noindent Source : 7II:662; 7II:1092. \nolinkurl{fixes/MISPRINTS_7II.json#662}; \nolinkurl{fixes/MISPRINTS_7II.json#1092}\par

\par\medskip\noindent\hypertarget{app-F7II-XXI-012.3}{}\textbf{F7II-XXI-012.3}\quad p. 370; misprint\_pass.\par
\noindent Imprimé : {\ttfamily\footnotesize paramétérisées\allowbreak{}}\par
\noindent Corrigé : {\ttfamily\footnotesize paramétrisées\allowbreak{}}\par
\noindent Source : 7II:1093. \nolinkurl{fixes/FIXLOG_7II_CURRENT.jsonl#893}; \nolinkurl{fixes/MISPRINTS_7II.json#1093}\par

\par\medskip\noindent\hypertarget{app-F7II-XXI-012.4}{}\textbf{F7II-XXI-012.4}\quad p. 370; printed misprint.\par
\noindent Imprimé : \emph{(vide)}\par
\noindent Corrigé : {\ttfamily\footnotesize )\allowbreak{}}\par
\noindent Source : 7II:1094. \nolinkurl{fixes/MISPRINTS_7II.json#1094}\par

\par\medskip\noindent\hypertarget{app-F7II-XXI-013.1}{}\textbf{F7II-XXI-013.1}\quad p. 371; misprint\_pass.\par
\noindent Imprimé : {\ttfamily\footnotesize 0\allowbreak{}}\par
\noindent Corrigé : {\ttfamily\footnotesize 1\allowbreak{}}\par
\noindent Source : 7II:1105. \nolinkurl{fixes/FIXLOG_7II_CURRENT.jsonl#894}\par

\par\medskip\noindent\hypertarget{app-F7II-XXI-013.2}{}\textbf{F7II-XXI-013.2}\quad p. 371; designated-witness attribution.\par
\noindent Imprimé : {\ttfamily\footnotesize n\allowbreak{}}\par
\noindent Corrigé : {\ttfamily\footnotesize M\allowbreak{}}\par
\noindent Source : WITNESS-F7II-XXI-013.2. \nolinkurl{fr/sga7ii/diplomatic/expose_XXI_body.tex}; \nolinkurl{fr/sga7ii/current/expose_XXI_body.tex}\par

\par\medskip\noindent\hypertarget{app-F7II-XXI-013.3}{}\textbf{F7II-XXI-013.3}\quad p. 371; inherited corrected-layer difference.\par
\noindent Imprimé : {\ttfamily\footnotesize u\allowbreak{}}\par
\noindent Corrigé : {\ttfamily\footnotesize U\allowbreak{}}\par
\noindent Source : ADOPT-F7II-XXI-013.3. \nolinkurl{fr/sga7ii/diplomatic/expose_XXI_body.tex}; \nolinkurl{fr/sga7ii/current/expose_XXI_body.tex}\par

\par\medskip\noindent\hypertarget{app-F7II-XXI-014.1}{}\textbf{F7II-XXI-014.1}\quad p. 372; designated-witness attribution.\par
\noindent Imprimé : \emph{(vide)}\par
\noindent Corrigé : {\ttfamily\footnotesize 1.\allowbreak{}}\par
\noindent Source : WITNESS-F7II-XXI-014.1. \nolinkurl{fr/sga7ii/diplomatic/expose_XXI_body.tex}; \nolinkurl{fr/sga7ii/current/expose_XXI_body.tex}\par

\par\medskip\noindent\hypertarget{app-F7II-XXI-015.1}{}\textbf{F7II-XXI-015.1}\quad p. 372; misprint\_in\_print.\par
\noindent Imprimé : \emph{(vide)}\par
\noindent Corrigé : {\ttfamily\footnotesize )\allowbreak{}}\par
\noindent Source : FIX-7II-D-0757. \nolinkurl{fixes/FIXLOG_7II.jsonl#757}\par

\par\medskip\noindent\hypertarget{app-F7II-XXI-016.1}{}\textbf{F7II-XXI-016.1}\quad p. 373; printed misprint.\par
\noindent Imprimé : {\ttfamily\footnotesize copstructions\allowbreak{}}\par
\noindent Corrigé : {\ttfamily\footnotesize constructions\allowbreak{}}\par
\noindent Source : 7II:1108. \nolinkurl{fixes/MISPRINTS_7II.json#1108}\par

\par\medskip\noindent\hypertarget{app-F7II-XXI-016.2}{}\textbf{F7II-XXI-016.2}\quad p. 373; printed misprint.\par
\noindent Imprimé : \emph{(vide)}\par
\noindent Corrigé : {\ttfamily\footnotesize T\allowbreak{} }\par
\noindent Source : 7II:664; 7II:1109. \nolinkurl{fixes/MISPRINTS_7II.json#664}; \nolinkurl{fixes/MISPRINTS_7II.json#1109}\par

\par\medskip\noindent\hypertarget{app-F7II-XXI-016.3}{}\textbf{F7II-XXI-016.3}\quad p. 373; printed misprint.\par
\noindent Imprimé : {\ttfamily\footnotesize  T\allowbreak{}}\par
\noindent Corrigé : \emph{(vide)}\par
\noindent Source : 7II:664; 7II:1109. \nolinkurl{fixes/MISPRINTS_7II.json#664}; \nolinkurl{fixes/MISPRINTS_7II.json#1109}\par

\par\medskip\noindent\hypertarget{app-F7II-XXI-017.1}{}\textbf{F7II-XXI-017.1}\quad p. 373; printed misprint.\par
\noindent Imprimé : {\ttfamily\footnotesize se\allowbreak{}}\par
\noindent Corrigé : {\ttfamily\footnotesize sa\allowbreak{}}\par
\noindent Source : 7II:665; 7II:1110. \nolinkurl{fixes/MISPRINTS_7II.json#665}; \nolinkurl{fixes/MISPRINTS_7II.json#1110}\par

\par\medskip\noindent\hypertarget{app-F7II-XXI-018.1}{}\textbf{F7II-XXI-018.1}\quad p. 374; inherited corrected-layer difference.\par
\noindent Imprimé : {\ttfamily\footnotesize o\allowbreak{}}\par
\noindent Corrigé : {\ttfamily\footnotesize 0\allowbreak{}}\par
\noindent Source : ADOPT-F7II-XXI-018.1. \nolinkurl{fr/sga7ii/diplomatic/expose_XXI_body.tex}; \nolinkurl{fr/sga7ii/current/expose_XXI_body.tex}\par

\par\medskip\noindent\hypertarget{app-F7II-XXI-018.2}{}\textbf{F7II-XXI-018.2}\quad p. 374; inherited corrected-layer difference.\par
\noindent Imprimé : {\ttfamily\footnotesize o\allowbreak{}}\par
\noindent Corrigé : {\ttfamily\footnotesize 0\allowbreak{}}\par
\noindent Source : ADOPT-F7II-XXI-018.2. \nolinkurl{fr/sga7ii/diplomatic/expose_XXI_body.tex}; \nolinkurl{fr/sga7ii/current/expose_XXI_body.tex}\par

\par\medskip\noindent\hypertarget{app-F7II-XXI-019.1}{}\textbf{F7II-XXI-019.1}\quad p. 374; printed misprint.\par
\noindent Imprimé : {\ttfamily\footnotesize coordonnés\allowbreak{}}\par
\noindent Corrigé : {\ttfamily\footnotesize coordonnées\allowbreak{}}\par
\noindent Source : 7II:1111. \nolinkurl{fixes/MISPRINTS_7II.json#1111}\par

\par\medskip\noindent\hypertarget{app-F7II-XXI-019.2}{}\textbf{F7II-XXI-019.2}\quad p. 374; designated-witness attribution.\par
\noindent Imprimé : {\ttfamily\footnotesize o\allowbreak{}}\par
\noindent Corrigé : {\ttfamily\footnotesize 0\allowbreak{}}\par
\noindent Source : WITNESS-F7II-XXI-019.2. \nolinkurl{fr/sga7ii/diplomatic/expose_XXI_body.tex}; \nolinkurl{fr/sga7ii/current/expose_XXI_body.tex}\par

\par\medskip\noindent\hypertarget{app-F7II-XXI-019.3}{}\textbf{F7II-XXI-019.3}\quad p. 374; designated-witness attribution.\par
\noindent Imprimé : {\ttfamily\footnotesize o\allowbreak{}}\par
\noindent Corrigé : {\ttfamily\footnotesize 0\allowbreak{}}\par
\noindent Source : WITNESS-F7II-XXI-019.3. \nolinkurl{fr/sga7ii/diplomatic/expose_XXI_body.tex}; \nolinkurl{fr/sga7ii/current/expose_XXI_body.tex}\par

\par\medskip\noindent\hypertarget{app-F7II-XXI-020.1}{}\textbf{F7II-XXI-020.1}\quad p. 375; inherited corrected-layer difference.\par
\noindent Imprimé : {\ttfamily\footnotesize o\allowbreak{}}\par
\noindent Corrigé : {\ttfamily\footnotesize 0\allowbreak{}}\par
\noindent Source : ADOPT-F7II-XXI-020.1. \nolinkurl{fr/sga7ii/diplomatic/expose_XXI_body.tex}; \nolinkurl{fr/sga7ii/current/expose_XXI_body.tex}\par

\par\medskip\noindent\hypertarget{app-F7II-XXI-020.2}{}\textbf{F7II-XXI-020.2}\quad p. 375; inherited corrected-layer difference.\par
\noindent Imprimé : {\ttfamily\footnotesize o\allowbreak{}}\par
\noindent Corrigé : {\ttfamily\footnotesize 0\allowbreak{}}\par
\noindent Source : ADOPT-F7II-XXI-020.2. \nolinkurl{fr/sga7ii/diplomatic/expose_XXI_body.tex}; \nolinkurl{fr/sga7ii/current/expose_XXI_body.tex}\par

\par\medskip\noindent\hypertarget{app-F7II-XXI-020.3}{}\textbf{F7II-XXI-020.3}\quad p. 375; misprint\_pass.\par
\noindent Imprimé : {\ttfamily\footnotesize XX\allowbreak{}}\par
\noindent Corrigé : {\ttfamily\footnotesize XXII\allowbreak{}}\par
\noindent Source : 7II:1112. \nolinkurl{fixes/FIXLOG_7II_CURRENT.jsonl#895}\par

\par\medskip\noindent\hypertarget{app-F7II-XXI-020.4}{}\textbf{F7II-XXI-020.4}\quad p. 375; inherited corrected-layer difference.\par
\noindent Imprimé : {\ttfamily\footnotesize =\allowbreak{}}\par
\noindent Corrigé : {\ttfamily\footnotesize \textbackslash{}equiv\allowbreak{}}\par
\noindent Source : ADOPT-F7II-XXI-020.4. \nolinkurl{fr/sga7ii/diplomatic/expose_XXI_body.tex}; \nolinkurl{fr/sga7ii/current/expose_XXI_body.tex}\par

\par\medskip\noindent\hypertarget{app-F7II-XXI-020.5}{}\textbf{F7II-XXI-020.5}\quad p. 375; designated-witness attribution.\par
\noindent Imprimé : {\ttfamily\footnotesize =\allowbreak{}}\par
\noindent Corrigé : {\ttfamily\footnotesize \textbackslash{}equiv\allowbreak{}}\par
\noindent Source : WITNESS-F7II-XXI-020.5. \nolinkurl{fr/sga7ii/diplomatic/expose_XXI_body.tex}; \nolinkurl{fr/sga7ii/current/expose_XXI_body.tex}\par

\par\medskip\noindent\hypertarget{app-F7II-XXI-020.6}{}\textbf{F7II-XXI-020.6}\quad p. 375; printed misprint.\par
\noindent Imprimé : \emph{(vide)}\par
\noindent Corrigé : {\ttfamily\footnotesize -\allowbreak{}}\par
\noindent Source : 7II:1113. \nolinkurl{fixes/MISPRINTS_7II.json#1113}\par

\par\medskip\noindent\hypertarget{app-F7II-XXI-020.7}{}\textbf{F7II-XXI-020.7}\quad p. 375; misprint\_in\_print.\par
\noindent Imprimé : {\ttfamily\footnotesize .\allowbreak{}}\par
\noindent Corrigé : {\ttfamily\footnotesize )\allowbreak{}}\par
\noindent Source : FIX-7II-D-0762. \nolinkurl{fixes/FIXLOG_7II.jsonl#762}\par

\par\medskip\noindent\hypertarget{app-F7II-XXI-020.8}{}\textbf{F7II-XXI-020.8}\quad p. 375; misprint\_in\_print.\par
\noindent Imprimé : {\ttfamily\footnotesize .\allowbreak{}}\par
\noindent Corrigé : {\ttfamily\footnotesize )\allowbreak{}}\par
\noindent Source : FIX-7II-D-0762. \nolinkurl{fixes/FIXLOG_7II.jsonl#762}\par

\par\medskip\noindent\hypertarget{app-F7II-XXI-021.1}{}\textbf{F7II-XXI-021.1}\quad p. 376; designated-witness attribution.\par
\noindent Imprimé : {\ttfamily\footnotesize o\allowbreak{}}\par
\noindent Corrigé : {\ttfamily\footnotesize 0\allowbreak{}}\par
\noindent Source : WITNESS-F7II-XXI-021.1. \nolinkurl{fr/sga7ii/diplomatic/expose_XXI_body.tex}; \nolinkurl{fr/sga7ii/current/expose_XXI_body.tex}\par

\par\medskip\noindent\hypertarget{app-F7II-XXI-021.2}{}\textbf{F7II-XXI-021.2}\quad p. 376; printed misprint.\par
\noindent Imprimé : {\ttfamily\footnotesize une\allowbreak{}}\par
\noindent Corrigé : {\ttfamily\footnotesize un\allowbreak{}}\par
\noindent Source : 7II:669; 7II:1095. \nolinkurl{fixes/MISPRINTS_7II.json#669}; \nolinkurl{fixes/MISPRINTS_7II.json#1095}\par

\par\medskip\noindent\hypertarget{app-F7II-XXI-022.1}{}\textbf{F7II-XXI-022.1}\quad p. 376; inherited corrected-layer difference.\par
\noindent Imprimé : \emph{(vide)}\par
\noindent Corrigé : {\ttfamily\footnotesize (\allowbreak{}}\par
\noindent Source : ADOPT-F7II-XXI-022.1. \nolinkurl{fr/sga7ii/diplomatic/expose_XXI_body.tex}; \nolinkurl{fr/sga7ii/current/expose_XXI_body.tex}\par

\par\medskip\noindent\hypertarget{app-F7II-XXI-022.2}{}\textbf{F7II-XXI-022.2}\quad p. 376; inherited corrected-layer difference.\par
\noindent Imprimé : \emph{(vide)}\par
\noindent Corrigé : {\ttfamily\footnotesize )\allowbreak{}}\par
\noindent Source : ADOPT-F7II-XXI-022.2. \nolinkurl{fr/sga7ii/diplomatic/expose_XXI_body.tex}; \nolinkurl{fr/sga7ii/current/expose_XXI_body.tex}\par

\par\medskip\noindent\hypertarget{app-F7II-XXI-022.3}{}\textbf{F7II-XXI-022.3}\quad p. 376; inherited corrected-layer difference.\par
\noindent Imprimé : {\ttfamily\footnotesize )\allowbreak{}}\par
\noindent Corrigé : \emph{(vide)}\par
\noindent Source : ADOPT-F7II-XXI-022.3. \nolinkurl{fr/sga7ii/diplomatic/expose_XXI_body.tex}; \nolinkurl{fr/sga7ii/current/expose_XXI_body.tex}\par

\par\medskip\noindent\hypertarget{app-F7II-XXI-022.4}{}\textbf{F7II-XXI-022.4}\quad p. 376; inherited corrected-layer difference.\par
\noindent Imprimé : \emph{(vide)}\par
\noindent Corrigé : {\ttfamily\footnotesize )\allowbreak{}}\par
\noindent Source : ADOPT-F7II-XXI-022.4. \nolinkurl{fr/sga7ii/diplomatic/expose_XXI_body.tex}; \nolinkurl{fr/sga7ii/current/expose_XXI_body.tex}\par

\par\medskip\noindent\hypertarget{app-F7II-XXI-022.5}{}\textbf{F7II-XXI-022.5}\quad p. 376; designated-witness attribution.\par
\noindent Imprimé : {\ttfamily\footnotesize o\allowbreak{}}\par
\noindent Corrigé : {\ttfamily\footnotesize 0\allowbreak{}}\par
\noindent Source : WITNESS-F7II-XXI-022.5. \nolinkurl{fr/sga7ii/diplomatic/expose_XXI_body.tex}; \nolinkurl{fr/sga7ii/current/expose_XXI_body.tex}\par

\par\medskip\noindent\hypertarget{app-F7II-XXI-022.6}{}\textbf{F7II-XXI-022.6}\quad p. 376; misprint\_in\_print.\par
\noindent Imprimé : {\ttfamily\footnotesize .\allowbreak{}}\par
\noindent Corrigé : {\ttfamily\footnotesize \textbackslash{}dots\allowbreak{}}\par
\noindent Source : FIX-7II-D-0755. \nolinkurl{fixes/FIXLOG_7II.jsonl#755}\par

\par\medskip\noindent\hypertarget{app-F7II-XXI-023.1}{}\textbf{F7II-XXI-023.1}\quad p. 377; inherited corrected-layer difference.\par
\noindent Imprimé : \emph{(vide)}\par
\noindent Corrigé : {\ttfamily\footnotesize espaces\allowbreak{} vectoriels\allowbreak{} de\allowbreak{} dimension\allowbreak{} finie\allowbreak{} sur\allowbreak{} }\par
\noindent Source : ADOPT-F7II-XXI-023.1. \nolinkurl{fr/sga7ii/diplomatic/expose_XXI_body.tex}; \nolinkurl{fr/sga7ii/current/expose_XXI_body.tex}\par

\par\medskip\noindent\hypertarget{app-F7II-XXI-023.2}{}\textbf{F7II-XXI-023.2}\quad p. 377; designated-witness attribution.\par
\noindent Imprimé : {\ttfamily\footnotesize -vectoriels\allowbreak{} de\allowbreak{} rang\allowbreak{} fini\allowbreak{}}\par
\noindent Corrigé : {\ttfamily\footnotesize ,\allowbreak{}}\par
\noindent Source : WITNESS-F7II-XXI-023.2. \nolinkurl{fr/sga7ii/diplomatic/expose_XXI_body.tex}; \nolinkurl{fr/sga7ii/current/expose_XXI_body.tex}\par

\par\medskip\noindent\hypertarget{app-F7II-XXI-023.3}{}\textbf{F7II-XXI-023.3}\quad p. 377; printed misprint.\par
\noindent Imprimé : \emph{(vide)}\par
\noindent Corrigé : {\ttfamily\footnotesize \textbackslash{}overline\{\textbackslash{}mathbb\{F\}\allowbreak{}\}\_q\allowbreak{}}\par
\noindent Source : 7II:1097. \nolinkurl{fixes/MISPRINTS_7II.json#1097}\par

\par\medskip\noindent\hypertarget{app-F7II-XXI-023.4}{}\textbf{F7II-XXI-023.4}\quad p. 377; printed misprint.\par
\noindent Imprimé : \emph{(vide)}\par
\noindent Corrigé : {\ttfamily\footnotesize \textbackslash{}overline\{\textbackslash{}mathbb\{F\}\allowbreak{}\}\_q\allowbreak{}}\par
\noindent Source : 7II:1098. \nolinkurl{fixes/MISPRINTS_7II.json#1098}\par

\par\medskip\noindent\hypertarget{app-F7II-XXI-023.5}{}\textbf{F7II-XXI-023.5}\quad p. 377; printed misprint.\par
\noindent Imprimé : {\ttfamily\footnotesize donné\allowbreak{}}\par
\noindent Corrigé : {\ttfamily\footnotesize donnée\allowbreak{}}\par
\noindent Source : 7II:1099. \nolinkurl{fixes/MISPRINTS_7II.json#1099}\par

\par\medskip\noindent\hypertarget{app-F7II-XXI-024.1}{}\textbf{F7II-XXI-024.1}\quad p. 378; designated-witness attribution.\par
\noindent Imprimé : \emph{(vide)}\par
\noindent Corrigé : {\ttfamily\footnotesize ce\allowbreak{} }\par
\noindent Source : WITNESS-F7II-XXI-024.1. \nolinkurl{fr/sga7ii/diplomatic/expose_XXI_body.tex}; \nolinkurl{fr/sga7ii/current/expose_XXI_body.tex}\par

\par\medskip\noindent\hypertarget{app-F7II-XXI-024.2}{}\textbf{F7II-XXI-024.2}\quad p. 378; designated-witness attribution.\par
\noindent Imprimé : {\ttfamily\footnotesize prendre\allowbreak{}}\par
\noindent Corrigé : {\ttfamily\footnotesize élever\allowbreak{}}\par
\noindent Source : WITNESS-F7II-XXI-024.2. \nolinkurl{fr/sga7ii/diplomatic/expose_XXI_body.tex}; \nolinkurl{fr/sga7ii/current/expose_XXI_body.tex}\par

\par\medskip\noindent\hypertarget{app-F7II-XXI-024.3}{}\textbf{F7II-XXI-024.3}\quad p. 378; designated-witness attribution.\par
\noindent Imprimé : {\ttfamily\footnotesize  puissances.\allowbreak{} \$M\$\textbackslash{}textsuperscript\{\allowbreak{}èmes\}\allowbreak{} des\allowbreak{}}\par
\noindent Corrigé : \emph{(vide)}\par
\noindent Source : WITNESS-F7II-XXI-024.3. \nolinkurl{fr/sga7ii/diplomatic/expose_XXI_body.tex}; \nolinkurl{fr/sga7ii/current/expose_XXI_body.tex}\par

\par\medskip\noindent\hypertarget{app-F7II-XXI-024.4}{}\textbf{F7II-XXI-024.4}\quad p. 378; designated-witness attribution.\par
\noindent Imprimé : \emph{(vide)}\par
\noindent Corrigé : {\ttfamily\footnotesize  à\allowbreak{} la\allowbreak{} puissance\allowbreak{} \$M\$\allowbreak{}}\par
\noindent Source : WITNESS-F7II-XXI-024.4. \nolinkurl{fr/sga7ii/diplomatic/expose_XXI_body.tex}; \nolinkurl{fr/sga7ii/current/expose_XXI_body.tex}\par

\par\medskip\noindent\hypertarget{app-F7II-XXI-025.1}{}\textbf{F7II-XXI-025.1}\quad p. 378; designated-witness attribution.\par
\noindent Imprimé : \emph{(vide)}\par
\noindent Corrigé : {\ttfamily\footnotesize \}\$\textbackslash{}mathrm\{\allowbreak{}}\par
\noindent Source : WITNESS-F7II-XXI-025.1. \nolinkurl{fr/sga7ii/diplomatic/expose_XXI_body.tex}; \nolinkurl{fr/sga7ii/current/expose_XXI_body.tex}\par

\par\medskip\noindent\hypertarget{app-F7II-XXI-025.2}{}\textbf{F7II-XXI-025.2}\quad p. 378; designated-witness attribution.\par
\noindent Imprimé : {\ttfamily\footnotesize \$\allowbreak{}}\par
\noindent Corrigé : \emph{(vide)}\par
\noindent Source : WITNESS-F7II-XXI-025.2. \nolinkurl{fr/sga7ii/diplomatic/expose_XXI_body.tex}; \nolinkurl{fr/sga7ii/current/expose_XXI_body.tex}\par

\par\medskip\noindent\hypertarget{app-F7II-XXI-025.3}{}\textbf{F7II-XXI-025.3}\quad p. 378; printed misprint.\par
\noindent Imprimé : \emph{(vide)}\par
\noindent Corrigé : {\ttfamily\footnotesize \textbackslash{}overline\{\textbackslash{}mathbb\{F\}\allowbreak{}\}\_q\allowbreak{}}\par
\noindent Source : 7II:674; 7II:1101. \nolinkurl{fixes/MISPRINTS_7II.json#674}; \nolinkurl{fixes/MISPRINTS_7II.json#1101}\par

\par\medskip\noindent\hypertarget{app-F7II-XXI-025.4}{}\textbf{F7II-XXI-025.4}\quad p. 378; designated-witness attribution.\par
\noindent Imprimé : \emph{(vide)}\par
\noindent Corrigé : {\ttfamily\footnotesize \textbackslash{}bigoplus\_i\allowbreak{} }\par
\noindent Source : WITNESS-F7II-XXI-025.4. \nolinkurl{fr/sga7ii/diplomatic/expose_XXI_body.tex}; \nolinkurl{fr/sga7ii/current/expose_XXI_body.tex}\par

\par\medskip\noindent\hypertarget{app-F7II-XXI-025.5}{}\textbf{F7II-XXI-025.5}\quad p. 378; printed misprint.\par
\noindent Imprimé : \emph{(vide)}\par
\noindent Corrigé : {\ttfamily\footnotesize )\allowbreak{}}\par
\noindent Source : 7II:1102. \nolinkurl{fixes/MISPRINTS_7II.json#1102}\par

\par\medskip\noindent\hypertarget{app-F7II-XXI-025.6}{}\textbf{F7II-XXI-025.6}\quad p. 378; designated-witness attribution.\par
\noindent Imprimé : \emph{(vide)}\par
\noindent Corrigé : {\ttfamily\footnotesize \_i\allowbreak{}}\par
\noindent Source : WITNESS-F7II-XXI-025.6. \nolinkurl{fr/sga7ii/diplomatic/expose_XXI_body.tex}; \nolinkurl{fr/sga7ii/current/expose_XXI_body.tex}\par

\par\medskip\noindent\hypertarget{app-F7II-XXI-026.1}{}\textbf{F7II-XXI-026.1}\quad p. 379; inherited corrected-layer difference.\par
\noindent Imprimé : {\ttfamily\footnotesize dans\allowbreak{} \$\textbackslash{}mathbb\{P\}\textasciicircum{}N\$\allowbreak{} }\par
\noindent Corrigé : \emph{(vide)}\par
\noindent Source : ADOPT-F7II-XXI-026.1. \nolinkurl{fr/sga7ii/diplomatic/expose_XXI_body.tex}; \nolinkurl{fr/sga7ii/current/expose_XXI_body.tex}\par

\par\medskip\noindent\hypertarget{app-F7II-XXI-026.2}{}\textbf{F7II-XXI-026.2}\quad p. 379; inherited corrected-layer difference.\par
\noindent Imprimé : \emph{(vide)}\par
\noindent Corrigé : {\ttfamily\footnotesize \$\allowbreak{} dans\allowbreak{} \$\textbackslash{}mathbb\{P\}\textasciicircum{}N\allowbreak{}}\par
\noindent Source : ADOPT-F7II-XXI-026.2. \nolinkurl{fr/sga7ii/diplomatic/expose_XXI_body.tex}; \nolinkurl{fr/sga7ii/current/expose_XXI_body.tex}\par

\par\medskip\noindent\hypertarget{app-F7II-XXI-027.1}{}\textbf{F7II-XXI-027.1}\quad p. 379; printed misprint.\par
\noindent Imprimé : {\ttfamily\footnotesize da\allowbreak{}}\par
\noindent Corrigé : {\ttfamily\footnotesize la\allowbreak{}}\par
\noindent Source : 7II:1103. \nolinkurl{fixes/MISPRINTS_7II.json#1103}\par

\par\medskip\noindent\hypertarget{app-F7II-XXI-028.1}{}\textbf{F7II-XXI-028.1}\quad p. 379; printed misprint.\par
\noindent Imprimé : {\ttfamily\footnotesize \textbackslash{};\allowbreak{}}\par
\noindent Corrigé : {\ttfamily\footnotesize i\allowbreak{}}\par
\noindent Source : 7II:1104. \nolinkurl{fixes/MISPRINTS_7II.json#1104}\par

\par\medskip\noindent\hypertarget{app-F7II-XXI-029.1}{}\textbf{F7II-XXI-029.1}\quad p. 381; inherited corrected-layer difference.\par
\noindent Imprimé : {\ttfamily\footnotesize /\allowbreak{}}\par
\noindent Corrigé : \emph{(vide)}\par
\noindent Source : ADOPT-F7II-XXI-029.1. \nolinkurl{fr/sga7ii/diplomatic/expose_XXI_body.tex}; \nolinkurl{fr/sga7ii/current/expose_XXI_body.tex}\par

\par\medskip\noindent\hypertarget{app-F7II-XXI-030.1}{}\textbf{F7II-XXI-030.1}\quad p. 382; inherited corrected-layer difference.\par
\noindent Imprimé : \emph{(vide)}\par
\noindent Corrigé : {\ttfamily\footnotesize i=\allowbreak{}}\par
\noindent Source : ADOPT-F7II-XXI-030.1. \nolinkurl{fr/sga7ii/diplomatic/expose_XXI_body.tex}; \nolinkurl{fr/sga7ii/current/expose_XXI_body.tex}\par

\par\medskip\noindent\hypertarget{app-F7II-XXI-031.1}{}\textbf{F7II-XXI-031.1}\quad p. 383; printed misprint.\par
\noindent Imprimé : {\ttfamily\footnotesize meme\allowbreak{} }\par
\noindent Corrigé : {\ttfamily\footnotesize même\allowbreak{}}\par
\noindent Source : 7II:678; 7II:1124. \nolinkurl{fixes/MISPRINTS_7II.json#678}; \nolinkurl{fixes/MISPRINTS_7II.json#1124}\par

\par\medskip\noindent\hypertarget{app-F7II-XXI-032.1}{}\textbf{F7II-XXI-032.1}\quad p. 384; misprint\_in\_print.\par
\noindent Imprimé : {\ttfamily\footnotesize nombresen\allowbreak{}}\par
\noindent Corrigé : {\ttfamily\footnotesize nombres\allowbreak{} en\allowbreak{}}\par
\noindent Source : FIX-7II-D-0320. \nolinkurl{fixes/FIXLOG_7II.jsonl#320}\par

\par\medskip\noindent\hypertarget{app-F7II-XXI-033.1}{}\textbf{F7II-XXI-033.1}\quad p. 386; misprint\_corrected.\par
\noindent Imprimé : \emph{(vide)}\par
\noindent Corrigé : {\ttfamily\footnotesize \textbackslash{}underline\{\allowbreak{}}\par
\noindent Source : FIX-7II-C-0882. \nolinkurl{fixes/FIXLOG_7II_CURRENT.jsonl#882}\par

\par\medskip\noindent\hypertarget{app-F7II-XXI-033.2}{}\textbf{F7II-XXI-033.2}\quad p. 386; misprint\_corrected.\par
\noindent Imprimé : \emph{(vide)}\par
\noindent Corrigé : {\ttfamily\footnotesize \}\allowbreak{}}\par
\noindent Source : FIX-7II-C-0882. \nolinkurl{fixes/FIXLOG_7II_CURRENT.jsonl#882}\par

\par\medskip\noindent\hypertarget{app-F7II-XXI-034.1}{}\textbf{F7II-XXI-034.1}\quad p. 386; printed misprint.\par
\noindent Imprimé : {\ttfamily\footnotesize F\allowbreak{}}\par
\noindent Corrigé : {\ttfamily\footnotesize E\allowbreak{}}\par
\noindent Source : 7II:1115. \nolinkurl{fixes/MISPRINTS_7II.json#1115}\par

\par\medskip\noindent\hypertarget{app-F7II-XXI-035.1}{}\textbf{F7II-XXI-035.1}\quad p. 388; misprint\_pass.\par
\noindent Imprimé : \emph{(vide)}\par
\noindent Corrigé : {\ttfamily\footnotesize .2\allowbreak{}}\par
\noindent Source : 7II:1117. \nolinkurl{fixes/FIXLOG_7II_CURRENT.jsonl#896}\par

\par\medskip\noindent\hypertarget{app-F7II-XXI-035.2}{}\textbf{F7II-XXI-035.2}\quad p. 388; printed misprint.\par
\noindent Imprimé : \emph{(vide)}\par
\noindent Corrigé : {\ttfamily\footnotesize '\allowbreak{}}\par
\noindent Source : 7II:680. \nolinkurl{fixes/MISPRINTS_7II.json#680}\par

\par\medskip\noindent\hypertarget{app-F7II-XXI-036.1}{}\textbf{F7II-XXI-036.1}\quad p. 389; printed misprint.\par
\noindent Imprimé : {\ttfamily\footnotesize Désignnons\allowbreak{}}\par
\noindent Corrigé : {\ttfamily\footnotesize Désignons\allowbreak{}}\par
\noindent Source : 7II:1119. \nolinkurl{fixes/MISPRINTS_7II.json#1119}\par

\par\medskip\noindent\hypertarget{app-F7II-XXI-036.2}{}\textbf{F7II-XXI-036.2}\quad p. 389; misprint\_pass.\par
\noindent Imprimé : \emph{(vide)}\par
\noindent Corrigé : {\ttfamily\footnotesize .\allowbreak{}}\par
\noindent Source : 7II:1120. \nolinkurl{fixes/FIXLOG_7II_CURRENT.jsonl#897}\par

\par\medskip\noindent\hypertarget{app-F7II-XXI-036.3}{}\textbf{F7II-XXI-036.3}\quad p. 389; inherited corrected-layer difference.\par
\noindent Imprimé : {\ttfamily\footnotesize Soient\allowbreak{}}\par
\noindent Corrigé : {\ttfamily\footnotesize Posons\}\allowbreak{} \$r=\textbackslash{}dim\_E\allowbreak{} V\textasciicircum{}I\$\allowbreak{} \textbackslash{}underline\{et\allowbreak{} soient\allowbreak{}}\par
\noindent Source : ADOPT-F7II-XXI-036.3. \nolinkurl{fr/sga7ii/diplomatic/expose_XXI_body.tex}; \nolinkurl{fr/sga7ii/current/expose_XXI_body.tex}\par

\par\medskip\noindent\hypertarget{app-F7II-XXI-036.4}{}\textbf{F7II-XXI-036.4}\quad p. 389; inherited corrected-layer difference.\par
\noindent Imprimé : \emph{(1 espace(s), 0 retour(s) de ligne dans le témoin TeX)}\par
\noindent Corrigé : {\ttfamily\footnotesize ,\allowbreak{}}\par
\noindent Source : ADOPT-F7II-XXI-036.4. \nolinkurl{fr/sga7ii/diplomatic/expose_XXI_body.tex}; \nolinkurl{fr/sga7ii/current/expose_XXI_body.tex}\par

\par\medskip\noindent\hypertarget{app-F7II-XXI-036.5}{}\textbf{F7II-XXI-036.5}\quad p. 389; inherited corrected-layer difference.\par
\noindent Imprimé : \emph{(1 espace(s), 0 retour(s) de ligne dans le témoin TeX)}\par
\noindent Corrigé : {\ttfamily\footnotesize ,\allowbreak{}}\par
\noindent Source : ADOPT-F7II-XXI-036.5. \nolinkurl{fr/sga7ii/diplomatic/expose_XXI_body.tex}; \nolinkurl{fr/sga7ii/current/expose_XXI_body.tex}\par

\par\medskip\noindent\hypertarget{app-F7II-XXI-036.6}{}\textbf{F7II-XXI-036.6}\quad p. 389; designated-witness attribution.\par
\noindent Imprimé : {\ttfamily\footnotesize n\allowbreak{}}\par
\noindent Corrigé : {\ttfamily\footnotesize r\allowbreak{}}\par
\noindent Source : WITNESS-F7II-XXI-036.6. \nolinkurl{fr/sga7ii/diplomatic/expose_XXI_body.tex}; \nolinkurl{fr/sga7ii/current/expose_XXI_body.tex}\par

\par\medskip\noindent\hypertarget{app-F7II-XXI-036.7}{}\textbf{F7II-XXI-036.7}\quad p. 389; inherited corrected-layer difference.\par
\noindent Imprimé : \emph{(vide)}\par
\noindent Corrigé : {\ttfamily\footnotesize \$\allowbreak{} \$(1\textbackslash{}leq\allowbreak{} i\textbackslash{}leq\allowbreak{} r)\allowbreak{}}\par
\noindent Source : ADOPT-F7II-XXI-036.7. \nolinkurl{fr/sga7ii/diplomatic/expose_XXI_body.tex}; \nolinkurl{fr/sga7ii/current/expose_XXI_body.tex}\par

\par\medskip\noindent\hypertarget{app-F7II-XXI-037.1}{}\textbf{F7II-XXI-037.1}\quad p. 390; misprint\_pass.\par
\noindent Imprimé : {\ttfamily\footnotesize \$\allowbreak{}}\par
\noindent Corrigé : \emph{(vide)}\par
\noindent Source : 7II:1122; 7II:683. \nolinkurl{fixes/FIXLOG_7II_CURRENT.jsonl#898}; \nolinkurl{fixes/MISPRINTS_7II.json#683}\par

\par\medskip\noindent\hypertarget{app-F7II-XXI-037.2}{}\textbf{F7II-XXI-037.2}\quad p. 390; misprint\_pass.\par
\noindent Imprimé : {\ttfamily\footnotesize \$\allowbreak{}}\par
\noindent Corrigé : {\ttfamily\footnotesize :\allowbreak{} }\par
\noindent Source : 7II:1122; 7II:683. \nolinkurl{fixes/FIXLOG_7II_CURRENT.jsonl#898}; \nolinkurl{fixes/MISPRINTS_7II.json#683}\par

\par\medskip\noindent\hypertarget{app-F7II-XXI-038.1}{}\textbf{F7II-XXI-038.1}\quad p. 390; printed misprint.\par
\noindent Imprimé : \emph{(vide)}\par
\noindent Corrigé : {\ttfamily\footnotesize 5.\allowbreak{}}\par
\noindent Source : 7II:684; 7II:1123. \nolinkurl{fixes/MISPRINTS_7II.json#684}; \nolinkurl{fixes/MISPRINTS_7II.json#1123}\par

\par\medskip\noindent\hypertarget{app-F7II-XXI-038.2}{}\textbf{F7II-XXI-038.2}\quad p. 390; designated-witness attribution.\par
\noindent Imprimé : {\ttfamily\footnotesize irréductible\allowbreak{}}\par
\noindent Corrigé : {\ttfamily\footnotesize indécomposable\allowbreak{}}\par
\noindent Source : WITNESS-F7II-XXI-038.2. \nolinkurl{fr/sga7ii/diplomatic/expose_XXI_body.tex}; \nolinkurl{fr/sga7ii/current/expose_XXI_body.tex}\par

\par\medskip\noindent\hypertarget{app-F7II-XXI-039.1}{}\textbf{F7II-XXI-039.1}\quad p. 391; misprint\_in\_print.\par
\noindent Imprimé : \emph{(vide)}\par
\noindent Corrigé : {\ttfamily\footnotesize \textbackslash{}underline\{\allowbreak{}}\par
\noindent Source : FIX-7II-D-0768. \nolinkurl{fixes/FIXLOG_7II.jsonl#768}\par

\par\medskip\noindent\hypertarget{app-F7II-XXI-039.2}{}\textbf{F7II-XXI-039.2}\quad p. 391; misprint\_in\_print.\par
\noindent Imprimé : \emph{(vide)}\par
\noindent Corrigé : {\ttfamily\footnotesize \}\allowbreak{}}\par
\noindent Source : FIX-7II-D-0768. \nolinkurl{fixes/FIXLOG_7II.jsonl#768}\par

\par\medskip\noindent\hypertarget{app-F7II-XXI-040.1}{}\textbf{F7II-XXI-040.1}\quad p. 393; inherited corrected-layer difference.\par
\noindent Imprimé : \emph{(vide)}\par
\noindent Corrigé : {\ttfamily\footnotesize \textbackslash{}log\allowbreak{} }\par
\noindent Source : ADOPT-F7II-XXI-040.1. \nolinkurl{fr/sga7ii/diplomatic/expose_XXI_body.tex}; \nolinkurl{fr/sga7ii/current/expose_XXI_body.tex}\par

\par\medskip\noindent\hypertarget{app-F7II-XXI-041.1}{}\textbf{F7II-XXI-041.1}\quad p. 393; misprint\_in\_print.\par
\noindent Imprimé : \emph{(vide)}\par
\noindent Corrigé : {\ttfamily\footnotesize \textbackslash{}underline\{\allowbreak{}}\par
\noindent Source : FIX-7II-D-0772. \nolinkurl{fixes/FIXLOG_7II.jsonl#772}\par

\par\medskip\noindent\hypertarget{app-F7II-XXI-041.2}{}\textbf{F7II-XXI-041.2}\quad p. 393; misprint\_in\_print.\par
\noindent Imprimé : \emph{(vide)}\par
\noindent Corrigé : {\ttfamily\footnotesize \}\allowbreak{}}\par
\noindent Source : FIX-7II-D-0772. \nolinkurl{fixes/FIXLOG_7II.jsonl#772}\par

\par\medskip\noindent\hypertarget{app-F7II-XXI-042.1}{}\textbf{F7II-XXI-042.1}\quad p. 394; designated-witness attribution.\par
\noindent Imprimé : {\ttfamily\footnotesize 5\allowbreak{}}\par
\noindent Corrigé : {\ttfamily\footnotesize 2\allowbreak{}}\par
\noindent Source : WITNESS-F7II-XXI-042.1. \nolinkurl{fr/sga7ii/diplomatic/expose_XXI_body.tex}; \nolinkurl{fr/sga7ii/current/expose_XXI_body.tex}\par

\par\medskip\noindent\hypertarget{app-F7II-XXI-043.1}{}\textbf{F7II-XXI-043.1}\quad p. 394; designated-witness attribution.\par
\noindent Imprimé : {\ttfamily\footnotesize 5.2.2\allowbreak{}}\par
\noindent Corrigé : {\ttfamily\footnotesize l'assertion\allowbreak{} (iv)\allowbreak{}}\par
\noindent Source : WITNESS-F7II-XXI-043.1. \nolinkurl{fr/sga7ii/diplomatic/expose_XXI_body.tex}; \nolinkurl{fr/sga7ii/current/expose_XXI_body.tex}\par

\par\medskip\noindent\hypertarget{app-F7II-XXI-044.1}{}\textbf{F7II-XXI-044.1}\quad p. 395; misprint\_pass.\par
\noindent Imprimé : \emph{(vide)}\par
\noindent Corrigé : {\ttfamily\footnotesize )\allowbreak{}}\par
\noindent Source : 7II:688; 7II:687; 7II:1130. \nolinkurl{fixes/FIXLOG_7II_CURRENT.jsonl#767}; \nolinkurl{fixes/MISPRINTS_7II.json#687}; \nolinkurl{fixes/MISPRINTS_7II.json#1130}\par

\par\medskip\noindent\hypertarget{app-F7II-XXI-045.1}{}\textbf{F7II-XXI-045.1}\quad p. 396; printed misprint.\par
\noindent Imprimé : {\ttfamily\footnotesize différe\allowbreak{}}\par
\noindent Corrigé : {\ttfamily\footnotesize diffère\allowbreak{}}\par
\noindent Source : 7II:688; 7II:827. \nolinkurl{fixes/MISPRINTS_7II.json#688}; \nolinkurl{fixes/MISPRINTS_7II.json#827}\par

\par\medskip\noindent\hypertarget{app-F7II-XXI-045.2}{}\textbf{F7II-XXI-045.2}\quad p. 396; printed misprint.\par
\noindent Imprimé : {\ttfamily\footnotesize morphisme\allowbreak{}}\par
\noindent Corrigé : {\ttfamily\footnotesize morphismes\allowbreak{}}\par
\noindent Source : 7II:828. \nolinkurl{fixes/MISPRINTS_7II.json#828}\par

\par\medskip\noindent\hypertarget{app-F7II-XXI-045.3}{}\textbf{F7II-XXI-045.3}\quad p. 396; designated-witness attribution.\par
\noindent Imprimé : {\ttfamily\footnotesize \textbackslash{},\allowbreak{}}\par
\noindent Corrigé : {\ttfamily\footnotesize \textbackslash{}circ\allowbreak{} }\par
\noindent Source : WITNESS-F7II-XXI-045.3. \nolinkurl{fr/sga7ii/diplomatic/expose_XXI_body.tex}; \nolinkurl{fr/sga7ii/current/expose_XXI_body.tex}\par

\par\medskip\noindent\hypertarget{app-F7II-XXI-046.1}{}\textbf{F7II-XXI-046.1}\quad p. 397; printed misprint.\par
\noindent Imprimé : {\ttfamily\footnotesize .\allowbreak{}}\par
\noindent Corrigé : \emph{(vide)}\par
\noindent Source : 7II:690; 7II:829. \nolinkurl{fixes/MISPRINTS_7II.json#690}; \nolinkurl{fixes/MISPRINTS_7II.json#829}\par

\par\medskip\noindent\hypertarget{app-F7II-XXI-047.1}{}\textbf{F7II-XXI-047.1}\quad p. 398; misprint\_pass.\par
\noindent Imprimé : {\ttfamily\footnotesize constant\allowbreak{}}\par
\noindent Corrigé : {\ttfamily\footnotesize constants\allowbreak{}}\par
\noindent Source : 7II:692. \nolinkurl{fixes/FIXLOG_7II_CURRENT.jsonl#768}\par

\par\medskip\noindent\hypertarget{app-F7II-XXI-048.1}{}\textbf{F7II-XXI-048.1}\quad p. 399; printed misprint.\par
\noindent Imprimé : \emph{(vide)}\par
\noindent Corrigé : {\ttfamily\footnotesize 5.\allowbreak{}}\par
\noindent Source : 7II:692; 7II:831. \nolinkurl{fixes/MISPRINTS_7II.json#692}; \nolinkurl{fixes/MISPRINTS_7II.json#831}\par

\normalsize\fi
\clearpage\SGAtoc{Exposé XXII — UNE FORMULE DE CONGRUENCE POUR LA FONCTION \texorpdfstring{$\zeta$}{zeta}}
\selectlanguage{french}
\SGApage{II}{401}

S G A  7


\underline{E X P O S E  XXII}


\underline{\underline{UNE FORMULE DE CONGRUENCE POUR LA FONCTION}} $\zeta$

\underline{par N. KATZ}


\SGAsectiontoc{0. Introduction}{0}
0. \underline{Introduction}

\SGAvar{Soit $X$ un schéma propre sur $\mathbb F_q$. Sa fonction zêta, considérée
comme élément de $\mathbb Z[[T]]$, peut-être réduite modulo $p$ ; on veut la
calculer modulo $p$. On trouve une expression du type auquel on pouvait
s'attendre depuis les formules de points fixes du type ``Woods Hole'', comme
produit alterné de polynômes caractéristiques d'un ``Frobenius'' convenable
agissant sur l'un quelconque des trois espaces de cohomologie suivants :

}{Soit $X$ un schéma propre sur $\mathbb F_q$. Sa fonction zêta, considérée
comme élément de $\mathbb Z[[T]]$, peut être réduite modulo $p$ ; on veut la
calculer modulo $p$. On trouve une expression du type auquel on pouvait
s'attendre depuis les formules de points fixes du type ``Woods Hole'', comme
produit alterné de polynômes caractéristiques d'un ``Frobenius'' convenable
agissant sur l'un quelconque des trois espaces de cohomologie suivants :

}{F7II-XXII-001}1) la cohomologie cohérente $H^i(X,\mathcal O_X)$ (3.1.1)

2) la cohomologie de De Rham
$H^i_{\mathrm{DR}}(X)=\mathbb H^i(X,\Omega_X^\bullet)$ (3.1.1)

3) la cohomologie étale
$H^i_{\mathrm{et}}(X,\mathbb Z/p\mathbb Z)$ (3.1.2).

Cette expression doit être vue comme une amélioration d'un cas spécial de la
formule des traces de Woods Hole : celui du faisceau structurel et de
l'endomorphisme de Frobenius. Ce cas est, en effet, un corollaire de la
formule de congruence (cf. 3.2.1).

La méthode de cet exposé est complètement naïve ; elle consiste à se ramener
au cas d'une hypersurface, ensuite à ``réduire modulo $p$''

\SGApage{II}{402}

\SGAvar{le calcul de DWORK pour la vraie fonction zêta d'une hypersurface. En
``réduisant'' soigneusement, on trouve l'amélioration 5.0.6 due à AX [1]
(et également à DWORK dans le cas lisse [4, \S 7])du théorème bien connu de
WARNING-CHEVALLEY [8].

Une autre méthode, entièrement cohomologique et dans l'esprit de SGA 5,
s'appliquant également à des ``coefficient tordus'', vient d'être trouvée par
DELIGNE, et devrait figurer dans la réédition de SGA 5 ; elle repose sur la
formule des traces type ``Woods Hole'', et la``formule de Künneth symétrique''
de DELIGNE (SGA 4 XVII 5.5.21).

}{le calcul de DWORK pour la vraie fonction zêta d'une hypersurface. En
``réduisant'' soigneusement, on trouve l'amélioration 5.0.6 due à AX [1]
(et également à DWORK dans le cas lisse [4, \S 7]) du théorème bien connu de
WARNING-CHEVALLEY [8].

Une autre méthode, entièrement cohomologique et dans l'esprit de SGA 5,
s'appliquant également à des ``coefficients tordus'', vient d'être trouvée par
DELIGNE, et devrait figurer dans la réédition de SGA 5 ; elle repose sur la
formule des traces type ``Woods Hole'', et la ``formule de Künneth symétrique''
de DELIGNE (SGA 4 XVII 5.5.21).

}{F7II-XXII-002}\SGAsectiontoc{1. Un rappel sur les opérations ``$p$-linéaires''}{1}
1. \underline{Un rappel sur les opérations} ``$p$-\underline{linéaires}''

\SGAvar{1.0. Soit $V$ un espace vectoriel de dimension finie sur un corps $k$ de
caractéristique $p>0$. Désignons par $q$ une puissance de $p$. Una application

}{1.0. Soit $V$ un espace vectoriel de dimension finie sur un corps $k$ de
caractéristique $p>0$. Désignons par $q$ une puissance de $p$. Une application

}{F7II-XXII-003}\begin{equation}\tag{1.0.1}
\varphi:V\longrightarrow V
\end{equation}

\SGAvar{s'appelle $q$-\underline{linéaire} si elle additive et si elle vérifie

}{s'appelle $q$-\underline{linéaire} si elle est additive et si elle vérifie

}{F7II-XXII-004}\begin{equation}\tag{1.0.2}
\varphi(\lambda v)=\lambda^q\varphi(v)
\qquad \forall\lambda\in k,\ v\in V.
\end{equation}

1.0.3. On fera attention au fait que l'itéré $n$'ième $\varphi^n$ d'un
endomorphisme $q$-linéaire est $q^n$-linéaire, et que, si $k$ n'est pas
parfait, l'image $\varphi(V)$ d'un endomorphisme $q$-linéaire n'est pas
nécessairement un $k$-sous-espace vectoriel.

1.0.4. Pour chaque extension $K/k$, un endomorphisme $q$-linéaire $\varphi$
de $V$ donne lieu à un unique endomorphisme $q$-linéaire $\varphi_K$ de
$V_K=V\otimes_k K$, satisfaisant à la formule

\begin{equation}\tag{1.0.5}
\varphi_K(\mu\otimes v)=\mu^q\varphi(v)
\qquad\text{pour }\mu\in K,\ v\in V.
\end{equation}

\SGApage{II}{403}

1.0.6. On dit qu'un endomorphisme $q$-linéaire $\varphi$ de $V$ est
semi-simple si son image \underline{engendre} $V$ comme $k$-vectoriel, et
qu'il est nilpotent s'il admet un itéré nul.

\SGAvar{On définit des $k$-sous-espace stables par $\varphi$

}{On définit des $k$-sous-espaces stables par $\varphi$

}{F7II-XXII-005}\begin{equation}\tag{1.0.7}
V_{\mathrm{ss}}=
\bigcap_{n\geq1}
\bigl(\text{le $k$-sous-espace engendré par }\varphi^n(V)\bigr)
\end{equation}

\begin{equation}\tag{1.0.8}
V_{\mathrm{nilp}}=
\bigcup_{n\geq1}\bigl(\text{le noyau de }\varphi^n:V\longrightarrow V\bigr).
\end{equation}

Évidemment, $\varphi$ restreint à $V_{\mathrm{ss}}$ est semi-simple, et
$\varphi$ restreint à $V_{\mathrm{nilp}}$ est nilpotent. \underline{Si} $k$ \underline{est parfait},
on a

\begin{equation}\tag{1.0.9}
V=V_{\mathrm{ss}}\oplus V_{\mathrm{nilp}},
\end{equation}

mais, si $k$ n'est pas parfait, (1.0.9) devient inexact, car il y a des
endomorphismes $q$-linéaires qui sont injectifs sans être semi-simples.

Pour une extension $K/k$ quelconque, on a

\begin{equation}\tag{1.0.10}
(V_K)_{\mathrm{ss}}=V_{\mathrm{ss}}\otimes_k K.
\end{equation}

\underline{Proposition} 1.1. \underline{Soit $\varphi$ un endomorphisme
$q$-linéaire d'un espace vectoriel $V$ de dimension finie sur un corps
séparablement clos $k$. Si $\varphi$ est semi-simple alors, posant}

\SGAvar{\[
V^{1-\varphi}=\text{\underline{les $\mathbb F_q$-vectoriel formé des points fixes de}
$\varphi$}=\Ker(1-\varphi),
\]

}{\[
V^{1-\varphi}=\text{\underline{le $\mathbb F_q$-vectoriel formé des points fixes de}
$\varphi$}=\Ker(1-\varphi),
\]

}{F7II-XXII-006}\underline{on a}

\begin{equation}\tag{1.1.1}
V\xleftarrow{\ \sim\ }
(V^{1-\varphi})\otimes_{\mathbb F_q}k.
\end{equation}

\SGApage{II}{404}

\underline{Corollaire} 1.1.2.
\[
V^{1-\varphi^n}\xleftarrow{\ \sim\ }
(V^{1-\varphi})\otimes_{\mathbb F_q}\mathbb F_{q^n}.
\]

\underline{Démonstration.} (LANG, cf. [6], p. 119).

Soit $\underline{e}$ une $k$-base de $V$, et écrivons

\begin{equation}\tag{1.1.3}
\varphi(\underline{e})=A\underline{e},\qquad A\in\mathrm{GL}(n,k)
\end{equation}

Désignons par

\begin{equation}\tag{1.1.4}
F:\mathrm{GL}(n,k)\longrightarrow\mathrm{GL}(n,k)
\end{equation}

la flèche ``élévation des coordonnées à la $q$'ième puissance''. Alors pour
$g\in\mathrm{GL}(n,k)$, $g\underline{e}$ est une base de $V$ formée de points fixes de
$\varphi$ si et seulement si

\begin{equation}\tag{1.1.5}
F(g)Ag^{-1}=1.
\end{equation}

Désignons par $\bar k$ une clôture algébrique de $k$. Le
\underline{groupe} $\mathrm{GL}(n,\bar k)$ \underline{agit} sur la variété
$\mathrm{GL}(n,\bar k)$ par

\begin{equation}\tag{1.1.6}
\left\{
\begin{aligned}
\mathrm{GL}(n,\bar k)\times\mathrm{GL}(n,\bar k)
&\longrightarrow\mathrm{GL}(n,\bar k)\\
(g,A)&\longmapsto F(g)Ag^{-1}.
\end{aligned}
\right.
\end{equation}

Pour $A$ fixé, la flèche

\begin{equation}\tag{1.1.7}
g\longmapsto F(g)Ag^{-1}
\end{equation}

\SGAvar{est étale comme on voit en calculant son application tangente, donc les
orbites sont ouvertes.Comme $\mathrm{GL}(n,\bar k)$ est irréductible, il n'y
a qu'une orbite ; autrement dit, la flèche

\begin{equation}\tag{1.1.8}
\begin{aligned}
\mathrm{Gl}(n,\bar k)&\longrightarrow\mathrm{Gl}(n,\bar k)\\
g&\longrightarrow F(g^{-1})\cdot g
\end{aligned}
\end{equation}

}{est étale comme on voit en calculant son application tangente, donc les
orbites sont ouvertes. Comme $\mathrm{GL}(n,\bar k)$ est irréductible, il n'y
a qu'une orbite ; autrement dit, la flèche

\begin{equation}\tag{1.1.8}
\begin{aligned}
\mathrm{GL}(n,\bar k)&\longrightarrow\mathrm{GL}(n,\bar k)\\
g&\longrightarrow F(g^{-1})\cdot g
\end{aligned}
\end{equation}

}{F7II-XXII-007}\SGApage{II}{405}

\SGAvar{est surjective. Parce qu'elle est aussi étale, il s'ensuit, $k$ étant
séparablement clos que la flèche

}{est surjective. Parce qu'elle est aussi étale, il s'ensuit, $k$ étant
séparablement clos, que la flèche

}{F7II-XXII-008}\begin{equation}\tag{1.1.9}
\begin{aligned}
\mathrm{GL}(n,k)&\longrightarrow\mathrm{GL}(n,k)\\
g&\longrightarrow F(g^{-1})\cdot g
\end{aligned}
\end{equation}

est surjective, ce qui veut dire que $V$ admet une $k$-base formée des points
fixes de $\varphi$. Il est alors clair, $\varphi$ étant $q$-linéaire, que tout
point fixe de $\varphi$ se trouve dans le $\mathbb F_q$-vectoriel engendré par
une $k$-base formée des points fixes de $\varphi$, ce qui prouve 1.1.

\underline{Corollaire} 1.1.10. \underline{Soit $\varphi$ un endomorphisme
$q$-linéaire d'un espace vectoriel $V$ de dimension finie sur un corps
séparablement clos $k$. Alors on a un isomorphisme canonique}

\begin{equation}\tag{1.1.11}
V_{\mathrm{ss}}\xleftarrow{\ \sim\ }
V^{1-\varphi}\otimes_{\mathbb F_q}k.
\end{equation}

Il suffit en effet d'appliquer 1.1 à $V_{\mathrm{ss}}$.

\SGAvar{\underline{Proposition} 1.2. \underline{Soit $\varphi$ un endomorphisme
$q$-linéaire d'un espace vectoriel de dimension finie sur un corps
séparablement clos $k$. Alors l'application additive}

}{\underline{Proposition} 1.2. \underline{Soit $\varphi$ un endomorphisme
$q$-linéaire d'un espace vectoriel $V$ de dimension finie sur un corps
séparablement clos $k$. Alors l'application additive}

}{F7II-XXII-009}\[
1-\varphi:V\longrightarrow V
\]

\underline{est surjective.}

\underline{Démonstration.} On a une suite exacte de groupes abéliens

\[
0\longrightarrow V_{\mathrm{ss}}\longrightarrow V
\longrightarrow V/V_{\mathrm{ss}}\longrightarrow0,
\]

\SGApage{II}{406}

sur laquelle agit $1-\varphi$. Par le ``lemme du serpent'', il suffit de
traiter $V_{\mathrm{ss}}$, évident par 1.1 en employant une base formée des
points fixes, et $V/V_{\mathrm{ss}}$, sur lequel $\varphi$ est nilpotente,
donc sur lequel $1-\varphi$ est inversible.

\SGAsectiontoc{2. Cohomologie cohérente et cohomologie étale}{2}
2. \underline{Cohomologie cohérente et cohomologie étale}

\underline{Proposition} 2.0. \underline{Soit} $\bar X$ \underline{un schéma propre sur
un corps} $k$ \underline{séparablement clos, et désignons par}

\begin{equation}\tag{2.0.1}
\overline{\mathfrak F}_p:H^{\cdot}(\bar X,\mathcal O_{\bar X})
\longrightarrow H^{\cdot}(\bar X,\mathcal O_{\bar X})
\end{equation}

\underline{l'application} $p$-\underline{linéaire induite par}

\begin{equation}\tag{2.0.2}
\overline{\mathfrak F}_p:\mathcal O_{\bar X}\longrightarrow\mathcal O_{\bar X},
\qquad \overline{\mathfrak F}_p(f)=f^p.
\end{equation}

\underline{Désignons par}

\[
H^{\cdot}_{\mathrm{et}}(\bar X,\mathbb Z/p\mathbb Z)
\]

\underline{les groupes de cohomologie étales à coefficients dans le faisceau
constant} $\mathbb Z/p\mathbb Z$. \underline{Alors on a des isomorphismes canoniques :}

\begin{equation}\tag{2.0.3}
H^i_{\mathrm{et}}(\bar X,\mathbb Z/p\mathbb Z)
\simeq
(H^i(\bar X,\mathcal O_{\bar X}))^{1-\overline{\mathfrak F}_p}.
\end{equation}

\underline{Démonstration.} On a une suite exacte de faisceaux sur
$\bar X_{\mathrm{et}}$.

\begin{equation}\tag{2.0.4}
0\longrightarrow\mathbb Z/p\mathbb Z\longrightarrow\mathcal O_{\bar X}
\xrightarrow{\ 1-\mathfrak F_p\ }\mathcal O_{\bar X}\longrightarrow0
\end{equation}

qui donne une suite exacte de cohomologie

\begin{equation}\tag{2.0.5}
\cdots\longrightarrow H^i_{\mathrm{et}}(\bar X,\mathbb Z/p\mathbb Z)
\longrightarrow H^i_{\mathrm{et}}(\bar X,\mathcal O_{\bar X})
\xrightarrow{\ 1-\mathfrak F_p\ }
H^i_{\mathrm{et}}(\bar X,\mathcal O_{\bar X})\longrightarrow\cdots .
\end{equation}

\SGAvar{On a
$H^i_{\mathrm{et}}(\bar X,\mathcal O_{\bar X})\simeq
H^i(\bar X,\mathcal O_{\bar X})$ ; c'est un $k$-vectoriel \underline{de dimension finie}
$\bar X$ étant propre. En appliquant 1.2, on trouve des suites exactes

}{On a
$H^i_{\mathrm{et}}(\bar X,\mathcal O_{\bar X})\simeq
H^i(\bar X,\mathcal O_{\bar X})$ ; c'est un $k$-vectoriel \underline{de dimension finie},
$\bar X$ étant propre. En appliquant 1.2, on trouve des suites exactes

}{F7II-XXII-010}\SGApage{II}{407}

\begin{equation}\tag{2.0.6}
0\longrightarrow H^i_{\mathrm{et}}(\bar X,\mathbb Z/p\mathbb Z)
\longrightarrow H^i(\bar X,\mathcal O_{\bar X})
\xrightarrow{\ 1-\mathfrak F_p\ }
H^i(\bar X,\mathcal O_{\bar X})\longrightarrow0,
\end{equation}

cqfd. Utilisant 1.1.10, on conclut de (2.0) :

\underline{Corollaire} 2.1. \underline{Pour} $\bar X$ \underline{propre sur un corps
séparablement clos} $k$,

\begin{equation}\tag{2.1.1}
H^i(\bar X,\mathcal O_{\bar X})_{\mathrm{ss}}
\xleftarrow{\ \sim\ }
H^i_{\mathrm{et}}(\bar X,\mathbb Z/p\mathbb Z)
\otimes_{\mathbb F_p}k.
\end{equation}

2.2. Soit maintenant $f:X\longrightarrow\Spec(\mathbb F_q)$ un schéma
propre sur $\mathbb F_q$. Désignons par $k$ une clôture algébrique de
$\mathbb F_q$, et posons

\begin{equation}\tag{2.2.1}
\bar X=X\times_{\mathbb F_q}k.
\end{equation}

\SGAvar{Le faisceau étale $R^if_*(\mathbb Z/p\mathbb Z)$ sur
$\operatorname{Sp}(\mathbb F_q)$ s'identifie à
$H^i(\bar X,\mathbb Z/p\mathbb Z)$, vu comme
$\Gal(k/\mathbb F_q)$ module. Désignons par

}{Le faisceau étale $R^if_*(\mathbb Z/p\mathbb Z)$ sur
$\operatorname{Sp}(\mathbb F_q)$ s'identifie à
$H^i(\bar X,\mathbb Z/p\mathbb Z)$, vu comme
$\Gal(k/\mathbb F_q)$-module. Désignons par

}{F7II-XXII-011}\begin{equation}\tag{2.2.2}
\mathfrak F_q:H^i(X,\mathcal O_X)\longrightarrow H^i(X,\mathcal O_X)
\end{equation}

l'opération linéaire (elle est $q$-linéaire, mais on est sur $\mathbb F_q$!)
induite par

\begin{equation}\tag{2.2.3}
\mathfrak F_q:\mathcal O_X\longrightarrow\mathcal O_X,
\qquad \mathfrak F_q(f)=f^q.
\end{equation}

\underline{Proposition} 2.2.4. \underline{Désignons par}
$\varphi_q\in\Gal(k/\mathbb F_q)$ \underline{l'élément de} ``\underline{Frobenius}'',
$\varphi_q(\lambda)=\lambda^q$ \underline{et soit} $i\in\mathbb Z$. \underline{On a une identité de
polynômes caractéristiques}

\begin{equation}\tag{2.2.4.1}
\det_{\mathbb F_p}\!\left(
1-t\varphi_q^{-1}\mid H^i_{\mathrm{et}}(\bar X,\mathbb Z/p\mathbb Z)
\right)
=
\det_{\mathbb F_q}\!\left(
1-t\mathfrak F_q\mid H^i(X,\mathcal O_X)
\right).
\end{equation}

\underline{où le premier membre est le polynôme caractéristique de l'élément}
$\varphi_q^{-1}$ \underline{dans} $\pi_1(\operatorname{Sp}(\mathbb F_q))$ \underline{agissant sur le
module galoisien} $H^i(\bar X,\mathbb Z/p\mathbb Z)$.

En effet, 2.2.4 est un corollaire de la

\SGApage{II}{408}

\underline{Proposition} 2.2.5. \underline{Soit} $X$ \underline{propre sur}
$\Spec(\mathbb F_q)$. \underline{Il y a un isomorphisme canonique}

\begin{equation}\tag{2.2.6}
H^i(X,\mathcal O_X)_{\mathrm{ss}}\otimes_{\mathbb F_q}k
\simeq
H^i_{\mathrm{et}}(\bar X,\mathbb Z/p\mathbb Z)
\otimes_{\mathbb F_p}k
\end{equation}

\SGAvar{\underline{qui transforme l'automorphisme}
$\mathfrak F_q\otimes_{\mathbb F_q}k$ \underline{du premier membre en l'automorphisme}
$\varphi_q^{-1}\otimes_{\mathbb F_q}k$ \underline{du second} (\underline{qu'on peut aussi interpréter
comme l'endomorphisme} $\mathrm{fr}_q^{X*}\otimes_{\mathbb F_p}k$, \underline{où}
$\mathrm{fr}_q^X:\bar X\longrightarrow\bar X$ \underline{est le} $k$-\underline{endomorphisme de
Frobenius.} cf. SGA 5, XVIII, 2.3.4).

}{\underline{qui transforme l'automorphisme}
$\mathfrak F_q\otimes_{\mathbb F_q}k$ \underline{du premier membre en l'automorphisme}
$\varphi_q^{-1}\otimes_{\mathbb F_p}k$ \underline{du second} (\underline{qu'on peut aussi interpréter
comme l'endomorphisme} $\mathrm{fr}_q^{X*}\otimes_{\mathbb F_p}k$, \underline{où}
$\mathrm{fr}_q^X:\bar X\longrightarrow\bar X$ \underline{est le} $k$-\underline{endomorphisme de
Frobenius.} cf. SGA 5, XVIII, 2.3.4).

}{F7II-XXII-012}\underline{Démonstration.} Compte tenu de ce que

\begin{equation}\tag{2.2.7}
H^i(\bar X,\mathcal O_{\bar X})\simeq
H^i(X,\mathcal O_X)\otimes_{\mathbb F_q}k,
\end{equation}

on a, par 1.0.10,

\SGAvar{\begin{equation}\tag{2.2.8}
H^i(\bar X,\mathcal O_{\bar X})_{\mathrm{ss}}\simeq
H^i(X,\mathcal O_X)_{\mathrm{ss}}\otimes_{\mathbb F_q}k,
\end{equation}

}{\begin{equation}\tag{2.2.8}
H^i(\bar X,\mathcal O_{\bar X})_{\mathrm{ss}}
\simeq H^i(X,\mathcal O_X)_{\mathrm{ss}}\otimes_{\mathbb F_q}k,
\end{equation}

}{}de sorte qu'on obtient l'isomorphisme 2.2.6 en composant les isomorphismes
2.2.8 et 2.1.1.

\SGAvar{Compte tenu de 2.2.7 l'opération $q$-linéaire

}{Compte tenu de 2.2.7, l'opération $q$-linéaire

}{F7II-XXII-014}\begin{equation}\tag{2.2.9}
\mathfrak F_q:H^i(X,\mathcal O_X)\longrightarrow H^i(X,\mathcal O_X)
\end{equation}

\SGAvar{admet une extension $q$-linéaire (cf. 1.0.5))

}{admet une extension $q$-linéaire (cf. 1.0.5)

}{F7II-XXII-015}\begin{equation}\tag{2.2.10}
\overline{\mathfrak F}_q:H^i(\bar X,\mathcal O_{\bar X})
\longrightarrow H^i(\bar X,\mathcal O_{\bar X}),
\end{equation}

et admet aussi une extension $k$-linéaire, notée encore
$\mathfrak F_q$

\begin{equation}\tag{2.2.11}
\mathfrak F_q:H^i(\bar X,\mathcal O_{\bar X})
\longrightarrow H^i(\bar X,\mathcal O_{\bar X}).
\end{equation}

\SGApage{II}{409}

On désigne par

\begin{equation}\tag{2.2.12}
\varphi_q:H^i(\bar X,\mathcal O_{\bar X})
\longrightarrow H^i(\bar X,\mathcal O_{\bar X})
\end{equation}

\SGAvar{l'opération $q$-linéaire dont les pointes fixes sont $H^i(X,\mathcal O_X)$.
Les opérateurs $\mathfrak F_q$, $\varphi_q$ et
$\overline{\mathfrak F}_q$, moyennant (2.2.7), s'écrivent aussi
$\mathfrak F_q\otimes\mathrm{id}_k$, $\mathrm{id}\otimes\varphi_q$,
$\mathfrak F_k\otimes\varphi_k$, donc commutent entre eux, et on a

}{l'opération $q$-linéaire dont les points fixes sont $H^i(X,\mathcal O_X)$.
Les opérateurs $\mathfrak F_q$, $\varphi_q$ et
$\overline{\mathfrak F}_q$, moyennant (2.2.7), s'écrivent aussi
$\mathfrak F_q\otimes\mathrm{id}_k$, $\mathrm{id}\otimes\varphi_q$,
$\mathfrak F_q\otimes\varphi_q$, donc commutent entre eux, et on a

}{F7II-XXII-016}\begin{equation}\tag{2.2.13}
\overline{\mathfrak F}_q
=\varphi_q\mathfrak F_q
=\mathfrak F_q\circ\varphi_q,
\end{equation}

et ils commutent tous à $\overline{\mathfrak F}_p$, de sorte que $\mathfrak F_q$ et
$\varphi_q$ agissent sur
$H^i(\bar X,\mathcal O_{\bar X})^{1-\overline{\mathfrak F}_q}$ et sur
$H^i(\bar X,\mathcal O_{\bar X})^{1-\overline{\mathfrak F}_p}$. Or,
$\mathfrak F_q$ et $\varphi_q$ sont des inverses l'une à l'autre \underline{sur}
$H^i(\bar X,\mathcal O_{\bar X})^{1-\overline{\mathfrak F}_q}$, par (2.2.13).
Donc l'opération de $\mathfrak F_q$ sur

\SGAvar{\[
H^i(\bar X,\mathcal O_{\bar X})_{\mathrm{ss}}
\simeq
H^i(\bar X,\mathcal O_{\bar X})^{1-\overline{\mathfrak F}_q}
\otimes_{\mathbb F_q}k
\qquad\text{(cf. 1.1.11)}
\]

}{\[
H^i(\bar X,\mathcal O_{\bar X})_{\mathrm{ss}}
\simeq H^i(\bar X,\mathcal O_{\bar X})^{1-\overline{\mathfrak F}_q}
\otimes_{\mathbb F_q}k
\qquad\text{(cf. 1.1.11)}
\]

}{}est l'inverse de l'opération de $\varphi_q\otimes k$ sur

\[
H^i(\bar X,\mathcal O_{\bar X})^{1-\overline{\mathfrak F}_q}
\otimes_{\mathbb F_q}k.
\]

Écrivant

\SGAvar{\begin{equation}\tag{2.2.14}
H^i(\bar X,\mathcal O_{\bar X})^{1-\mathfrak F_q}
\simeq
H^i(\bar X,\mathcal O_{\bar X})^{1-\mathfrak F_p}
\otimes_{\mathbb F_p}\mathbb F_q,
\end{equation}

l'automorphisme $\varphi_q$ de
$H^i(\bar X,\mathcal O_{\bar X})^{1-\mathfrak F_q}$ devient, via
2.2.14, l'automorphisme $\varphi_q\otimes\mathbb F_q$ de
$H^i(\bar X,\mathcal O_{\bar X})^{1-\overline{\mathfrak F}_p}
\otimes_{\mathbb F_p}\mathbb F_q$, de sorte que l'automorphisme
$\varphi_q\otimes k$ de $H^i(\bar X,\mathcal O_{\bar X})_{\mathrm{ss}}$
devient, via 2.1.1, l'automorphisme $\varphi_q\otimes k$ de
$H^i_{\mathrm{et}}(\bar X,\mathbb Z/p\mathbb Z)\otimes_{\mathbb F_p}k$,
ce qui achève la démonstration (2.2.5).

}{\begin{equation}\tag{2.2.14}
H^i(\bar X,\mathcal O_{\bar X})^{1-\overline{\mathfrak F}_q}
\simeq
H^i(\bar X,\mathcal O_{\bar X})^{1-\overline{\mathfrak F}_p}
\otimes_{\mathbb F_p}\mathbb F_q,
\end{equation}

l'automorphisme $\varphi_q$ de
$H^i(\bar X,\mathcal O_{\bar X})^{1-\overline{\mathfrak F}_q}$ devient, via
2.2.14, l'automorphisme $\varphi_q\otimes\mathbb F_q$ de
$H^i(\bar X,\mathcal O_{\bar X})^{1-\overline{\mathfrak F}_p}
\otimes_{\mathbb F_p}\mathbb F_q$, de sorte que l'automorphisme
$\varphi_q\otimes k$ de $H^i(\bar X,\mathcal O_{\bar X})_{\mathrm{ss}}$
devient, via 2.1.1, l'automorphisme $\varphi_q^{-1}\otimes k$ de
$H^i_{\mathrm{et}}(\bar X,\mathbb Z/p\mathbb Z)\otimes_{\mathbb F_p}k$,
ce qui achève la démonstration (2.2.5).

}{F7II-XXII-018}\SGApage{II}{410}

\SGAsectiontoc{3. La formule de congruence : énoncés et équivalences élémentaires}{3}
3. \underline{La formule de congruence : énoncés et équivalences élémentaires}

\SGAvar{Je tiens à remercier Grothendieck et Deligne pour m'avoir expliqué les
variantes (6.1.2), (3.1.1.2) et (3.1.3) de (3.1.1),

}{Je tiens à remercier Grothendieck et Deligne pour m'avoir expliqué les
variantes (3.1.2), (3.2.1)--(3.2.2) et (3.3.1) de (3.1.1).

}{F7II-XXII-019}3.0. Soit $X$ un schéma de type fini sur $\mathbb F_q$. On définit le degré
d'un point fermé $y\in X$ comme le degré de son corps résiduel
$\mathbb F_q(y)/\mathbb F_q$. On définit la fonction zêta de
$X/\mathbb F_q$ comme élément de $\mathbb Z[[t]]$, par la formule
(cf. aussi SGA 5 XIII 3.1)

\begin{equation}\tag{3.0.1}
Z(t,X/\mathbb F_q)=\prod_y(1-t^{\deg(y)})^{-1}.
\end{equation}

Par la ``fonction zêta modulo $p$'' on veut dire toujours l'image de
$Z(t,X/\mathbb F_q)$ dans $\mathbb F_q[[t]]$ par la flèche de ``réduction
modulo $p\mathbb Z[[t]]$''

\[
\mathbb Z[[t]]\longrightarrow\mathbb F_q[[t]].
\]

\underline{Théorème} 3.1. \underline{Soit} $X$ \underline{un schéma propre sur}
$\mathbb F_q$. \underline{Alors}

\begin{equation}\tag{3.1.1}
Z(t,X/\mathbb F_q)
\equiv
\prod_{i\geq0}
\det\!\left(1-t\mathfrak F_q\mid H^i(X,\mathcal O_X)\right)^{(-1)^{i+1}}
\pmod p
\end{equation}

\SGAvar{\underline{ou, ce qui est équivalent par} (2.2.4).

}{\underline{ou, ce qui est équivalent par} (2.2.4),

}{F7II-XXII-020}\begin{equation}\tag{3.1.2}
Z(t,X/\mathbb F_q)
\equiv
\prod_{i\geq0}
\det\!\left(1-t\varphi_q^{-1}\mid
H^i_{\mathrm{et}}(\bar X,\mathbb Z/p\mathbb Z)\right)^{(-1)^{i+1}}
\pmod p
\end{equation}

\underline{Corollaire} 3.2. \underline{Soit} $X$ \underline{un schéma propre sur}
$\mathbb F_q$.

\underline{Alors}

\SGAvar{\begin{equation}\tag{3.2.1}
\sum_{i\geq0}(-1)^i
\operatorname{trace}\!\left(\mathfrak F_q\mid H^i(X,\mathcal O_X)\right.
\equiv \#X(\mathbb F_q)\pmod p
\end{equation}

}{\begin{equation}\tag{3.2.1}
\sum_{i\geq0}(-1)^i
\operatorname{trace}\!\left(\mathfrak F_q\mid H^i(X,\mathcal O_X)\right)
\equiv \#X(\mathbb F_q)\pmod p
\end{equation}

}{F7II-XXII-021}\underline{ou, ce qui est équivalent par} (2.2.4)

\SGApage{II}{411}

\begin{equation}\tag{3.2.2}
\sum_{i\geq0}(-1)^i
\operatorname{tr}\!\left(\varphi_q^{-1}\mid
H^i_{\mathrm{et}}(\bar X,\mathbb Z/p\mathbb Z)\right)
\equiv \#X(\mathbb F_q)\mod p.
\end{equation}

On prend les coefficients de $t$ dans les congruences (3.1.1) et (3.1.2).

\SGAvar{\underline{Théorème} 3.3. \underline{Soit} $X$ \underline{un schéma propre sur}
$\mathbb F_q$ \underline{on désigne par} $H^i_{\mathrm{DR}}(X/\mathbb F_q)$ \underline{les}
$\mathbb F_q$-\underline{vectoriels} $\mathbb H^i(X,\Omega^\bullet_{X/\mathbb F_q})$, \underline{sur lesquels}
$\mathfrak F_q$ \underline{agit}. \underline{On a}

}{\underline{Théorème} 3.3. \underline{Soit} $X$ \underline{un schéma propre sur}
$\mathbb F_q$ ; \underline{on désigne par} $H^i_{\mathrm{DR}}(X/\mathbb F_q)$ \underline{les}
$\mathbb F_q$-\underline{vectoriels} $\mathbb H^i(X,\Omega^\bullet_{X/\mathbb F_q})$, \underline{sur lesquels}
$\mathfrak F_q$ \underline{agit}. \underline{On a}

}{F7II-XXII-022}\begin{equation}\tag{3.3.1}
Z(t,X/\mathbb F_q)
\equiv
\prod_{i\geq0}
\det\!\left(1-t\mathfrak F_q\mid
H^i_{\mathrm{DR}}(X/\mathbb F_q)\right)^{(-1)^{i+1}}
\mod p.
\end{equation}

En effet, utilisant la suite spectrale

\begin{equation}\tag{3.3.2}
E_1^{p,q}=H^q(X,\Omega^p_{X/\mathbb F_q})
\Longrightarrow H^{p+q}_{\mathrm{DR}}(X/\mathbb F_q)
\end{equation}

\SGAvar{sur laquelle $\mathfrak F_q$ agit (via l'endomorphisme
$\mathfrak F_q=\mathrm{fr}_p^{X*}$ déduit de
$\mathrm{fr}_q^X:X\longrightarrow X$) en tuant $E_1^{p,q}$ pour $q\ne0$, on
voit que (3.3.1) équivaut à (3.1.1).

}{sur laquelle $\mathfrak F_q$ agit via l'endomorphisme
$\mathfrak F_q=(\mathrm{fr}_q^X)^*$ déduit de
$\mathrm{fr}_q^X:X\longrightarrow X$, en tuant $E_1^{p,q}$ pour $p\ne0$, on
voit que (3.3.1) équivaut à (3.1.1).

}{F7II-XXII-023}Donnons un petit complément.

La formation de la ``partie semi-simple'' (cf. 1.0.7) étant un foncteur
exact, on déduit de (3.3.2) une suite spectrale

\begin{equation}\tag{3.3.3}
(E_1^{p,q})_{\mathrm{ss}}
\Longrightarrow H^{p+q}_{\mathrm{DR}}(X/\mathbb F_q)_{\mathrm{ss}}.
\end{equation}

\SGAvar{Comme $(E_1^{p,q})_{\mathrm{ss}}=0$ pour $q\ne0$, on en déduit pour tout
$i$ un isomorphisme

}{Comme $(E_1^{p,q})_{\mathrm{ss}}=0$ pour $p\ne0$, on en déduit pour tout
$i$ un isomorphisme

}{F7II-XXII-024}\begin{equation}\tag{3.3.4}
H^i_{\mathrm{DR}}(X/\mathbb F_q)_{\mathrm{ss}}
\xrightarrow{\ \sim\ }
H^i(X,\mathcal O_X)_{\mathrm{ss}},
\end{equation}

commutant à $\mathfrak F_q$, d'où

\begin{equation}\tag{3.3.5}
\det\!\left(1-t\mathfrak F_q\mid H^i_{\mathrm{DR}}(X/\mathbb F_q)\right)
=
\det\!\left(1-t\mathfrak F_q\mid H^i(X,\mathcal O_X)\right).
\end{equation}

\SGApage{II}{412}

\SGAvar{\underline{Théorème} 3.4. \underline{Soit} $X$ \underline{un schéma séparé de
type fini sur} $\mathbb F_q$. \underline{Alors désignant par} $H_c^i$ \underline{la cohomologie
étale à supports propres} (SGA 4 XVII 5.1.9 (iii)), \underline{on a}

}{\underline{Théorème} 3.4. \underline{Soit} $X$ \underline{un schéma séparé de
type fini sur} $\mathbb F_q$. \underline{Alors, désignant par} $H_c^i$ \underline{la cohomologie
étale à supports propres} (SGA 4 XVII 5.1.9 (iii)), \underline{on a}

}{F7II-XXII-025}\begin{equation}\tag{3.4.1}
Z(t,X/\mathbb F_q)
\equiv
\prod_{i\geq0}
\det\!\left(1-t\varphi_q^{-1}\mid
H_c^i(\bar X,\mathbb Z/p\mathbb Z)\right)^{(-1)^{i+1}}
\mod p.
\end{equation}

Commençons par démontrer l'équivalence de (3.4.1) et (3.1.1).

Écrivons, pour le moment

\begin{equation}\tag{3.4.2}
Z(X)\quad\text{au lieu de }Z(t,X/\mathbb F_q)
\end{equation}

\begin{equation}\tag{3.4.3}
Z'(X)\quad\text{au lieu de }\quad
\prod\det\!\left(1-t\varphi_q^{-1}\mid
H_c^i(\bar X,\mathbb Z/p\mathbb Z)\right)^{(-1)^{i+1}}.
\end{equation}

\underline{Lemme} 3.5. \underline{Si} $\{U_i\}$ \underline{est un recouvrement fini de}
$X$ \underline{par des sous-schémas ouverts on a}

\begin{equation}\tag{3.5.1}
Z'(X)=
\prod_{j\geq0}
\prod_{1\leq i_0<\cdots<i_j\leq n}
Z'(U_{i_0}\cap\cdots\cap U_{i_j})^{(-1)^j}.
\end{equation}

\underline{Démonstration.} Résulte facilement de la suite spectrale de
localisation (SGA 4 XVII 6.2.10)

\SGAvar{\[
E_1^{-a,b}=\bigoplus_{i_0<\cdots<i_b}
H_c^a(\bar U_{i_0}\cap\cdots\cap\bar U_{i_b},\mathbb Z/p\mathbb Z)
\Longrightarrow H_c^{-a+b}(\bar X,\mathbb Z/p\mathbb Z).
\]

}{\[
E_1^{-a,b}=\bigoplus_{i_0<\cdots<i_a}
H_c^b(\bar U_{i_0}\cap\cdots\cap\bar U_{i_a},\mathbb Z/p\mathbb Z)
\Longrightarrow H_c^{-a+b}(\bar X,\mathbb Z/p\mathbb Z).
\]

}{F7II-XXII-026}On peut aussi l'obtenir par récurrence sur $n$, en appliquant 3.6 ci-dessous.

\underline{Lemme} 3.6. \underline{Si} $U$ \underline{est un ouvert de} $X$, $Y$ \underline{le fermé}
$X-U$,

\begin{equation}\tag{3.6.1}
Z'(X)=Z'(U)Z'(Y)
\end{equation}

\underline{Démonstration.} Résulte de la suite exacte de cohomologie

\begin{equation}\tag{3.6.2}
\longrightarrow H_c^i(\bar U,\mathbb Z/p\mathbb Z)
\longrightarrow H_c^i(\bar X,\mathbb Z/p\mathbb Z)
\longrightarrow H_c^i(\bar Y,\mathbb Z/p\mathbb Z)\longrightarrow.
\end{equation}

\SGApage{II}{413}

\SGAvar{\underline{Lemme} 3.7 (Moebius). \underline{Soit} $B$ \underline{un ensemble fini},
$\{A_i\}$ \underline{un recouvrement fini par des sous-ensembles}. \underline{Alors}

}{\underline{Lemme} 3.7 (Möbius). \underline{Soit} $B$ \underline{un ensemble fini},
$\{A_i\}$ \underline{un recouvrement fini par des sous-ensembles}. \underline{Alors}

}{F7II-XXII-027}\begin{equation}\tag{3.7.1}
\operatorname{card}(B)=
\sum_{j\geq0}(-1)^j
\sum_{i_0<\cdots<i_j}
\operatorname{card}(A_{i_0}\cap\cdots\cap A_{i_j}).
\end{equation}

\underline{Démonstration.} On a la suite spectrale de Leray pour la
cohomologie rationnelle de l'espace discret $B$ et le recouvrement ouvert
$\{A_i\}$

\SGAvar{\[
E_1^{p,q}=\bigoplus_q
H^p(A_{i_0}\cap\cdots\cap A_{i_q},\mathbb Q)
\Longrightarrow H^{p+q}(B,\mathbb Q)
\]

}{\[
E_1^{p,q}=\bigoplus_{i_0<\cdots<i_q}
H^p(A_{i_0}\cap\cdots\cap A_{i_q},\mathbb Q)
\Longrightarrow H^{p+q}(B,\mathbb Q)
\]

}{F7II-XXII-028}et on calcule la caractéristique d'Euler-Poincaré de $B$.

\underline{Lemme} 3.8. \underline{Si} $\{U_i\}$ \underline{est un recouvrement fini de}
$X$ \underline{par des sous-schémas, alors}

\begin{equation}\tag{3.8.1}
Z(X)=
\prod_{j\geq0}
\prod_{i_0<\cdots<i_j}
Z(U_{i_0}\cap\cdots\cap U_{i_j})^{(-1)^j}.
\end{equation}

\underline{Démonstration.} On a la formule

\begin{equation}\tag{3.8.2}
\operatorname{card}(X(\mathbb F_{q^n}))
=\sum_{r\mid n}r.\operatorname{card}\{y\in X\mid\deg(y)=r\}
\end{equation}

qui entraîne

\begin{equation}\tag{3.8.3}
Z(X)=\exp\!\left(
\sum_{n\geq1}\operatorname{card}(X(\mathbb F_{q^n}))\frac{t^n}{n}
\right),
\end{equation}

et on conclut par 3.7.1.

\SGAvar{\underline{Lemme} 3.9. \underline{Supposons que la formule} (3.4.1) (\underline{ou ce qui
revient au même}, (3.1.2) \underline{i}.\underline{e}. (3.1.1)) \underline{soit vraie pour} $X$ \underline{une hypersurface
projective}. \underline{Alors elle est vraie pour tout} $X$ \underline{séparé de type fini}.

}{\underline{Lemme} 3.9. \underline{Supposons que la formule} (3.4.1) (\underline{ou ce qui
revient au même}, (3.1.2), \underline{c'est-à-dire} (3.1.1)) \underline{soit vraie pour} $X$ \underline{une
hypersurface projective}. \underline{Alors elle est vraie pour tout} $X$ \underline{séparé de type
fini}.

}{F7II-XXII-029}\SGApage{II}{414}

\SGAvar{\underline{Démonstration.} Par 3.8 et 3.6, (3.4.1) est vraie pour le
complémentaire d'une hypersurface, $\mathbb P^n-H$. Compte tenu de ce que $\bigcap(\mathbb P^n-H_i)=\mathbb P^n-\bigcup_iH_i$ est encore le complément d'une hypersurface, le recouvrement ouvert
$\{\mathbb P^n-H_i\}$ de $\mathbb P^n-\bigcap H_i$ permet d'appliquer 3.8 et
3.5 pour déduire (3.1.3) pour $\mathbb P^n-\bigcap H_i$. Par 3.8 et 3.6,
(3.4.1) est vrai pour $\bigcap H_i$, i.e. pour tout $X$ projectif. Une
nouvelle application de 3.8 et 3.5 donne 3.1.3 pour tout $X$ quasi-projectif,
en particulier pour tout $X$ affine. Finalement, pour $X$ séparé de type
fini, on prend un recouvrement fini par des ouverts affines, et on applique
3.8 et 3.5.

}{\underline{Démonstration.} Par 3.8 et 3.6, (3.4.1) est vraie pour le
complémentaire d'une hypersurface, $\mathbb P^n-H$. Compte tenu de ce que $\bigcap(\mathbb P^n-H_i)=\mathbb P^n-\bigcup_iH_i$ est encore le complément d'une hypersurface, le recouvrement ouvert
$\{\mathbb P^n-H_i\}$ de $\mathbb P^n-\bigcap H_i$ permet d'appliquer 3.8 et
3.5 pour déduire (3.4.1) pour $\mathbb P^n-\bigcap H_i$. Par 3.8 et 3.6,
(3.4.1) est vraie pour $\bigcap H_i$, c'est-à-dire pour tout $X$ projectif.
Une nouvelle application de 3.8 et 3.5 donne (3.4.1) pour tout $X$
quasi-projectif, en particulier pour tout $X$ affine. Finalement, pour $X$
séparé de type fini, on prend un recouvrement fini par des ouverts affines,
et on applique 3.8 et 3.5.

}{F7II-XXII-030}\SGAsectiontoc{4. Le calcul de la fonction zêta d'une hypersurface à la Dwork}{4}
4. \underline{Le calcul de la fonction zêta d'une hypersurface à la Dwork}

4.0. Soit

\begin{equation}\tag{4.0.1}
f(X)=f(X_1,\ldots,X_{n+2})\in\mathbb F_q[X]
\quad\text{homogène de degré }d,
\end{equation}

dont l'annulation définit l'hypersurface $V\subset\mathbb P^{n+1}$, et
désignons par $V^*$ l'ouvert de $V$ où
$\prod_{i=1}^{n+2}X_i$ est inversible, $V^*_{\mathrm{aff}}$ le cône affine
de $V^*$.

Choisissons un caractère additif non-trivial

\begin{equation}\tag{4.0.2}
\chi:\mathbb F_{q^\nu}\longrightarrow\mu_p(C),
\end{equation}

$C$ étant un corps alg. clos fixé de car. nulle
(p.ex. $C=\mathbb C$), de sorte qu'on a

\SGAvar{\begin{equation}\tag{4.0.3}
\text{pour }x\in\mathbb F_{q^\nu}\qquad
\sum_{y\in\mathbb F_{q^\nu}}\chi(xy)=
\begin{cases}
q^\nu&\text{si }\chi=0\\
0&\text{si }\chi\ne0.
\end{cases}
\end{equation}

Pour $\chi=(\chi_1,\ldots,\chi_{n+2})\in(\mathbb F_{q^\nu}^*)^{n+2}$, on a donc

}{\begin{equation}\tag{4.0.3}
\text{pour }x\in\mathbb F_{q^\nu}\qquad
\sum_{y\in\mathbb F_{q^\nu}}\chi(xy)=
\begin{cases}
q^\nu&\text{si }x=0\\
0&\text{si }x\ne0.
\end{cases}
\end{equation}

Pour $x=(x_1,\ldots,x_{n+2})\in(\mathbb F_{q^\nu}^*)^{n+2}$, on a donc

}{F7II-XXII-031}\SGApage{II}{415}

\SGAvar{\begin{equation}\tag{4.0.4}
1+\sum_{X_0\in\mathbb F_{q^\nu}^*}\chi(x_0f(x))
=\sum_{x_0\in\mathbb F_{q^\nu}}\chi(x_0f(x))
=
\begin{cases}
q^\nu&\text{si }f(x)=0\\
0&\text{sinon}.
\end{cases}
\end{equation}

}{\begin{equation}\tag{4.0.4}
1+\sum_{x_0\in\mathbb F_{q^\nu}^*}\chi(x_0f(x))
=\sum_{x_0\in\mathbb F_{q^\nu}}\chi(x_0f(x))
=
\begin{cases}
q^\nu&\text{si }f(x)=0\\
0&\text{sinon}.
\end{cases}
\end{equation}

}{F7II-XXII-032}En sommant (4.0.4) sur $x\in(\mathbb F_{q^\nu}^*)^{n+2}$, on obtient

\begin{equation}\tag{4.0.5}
\operatorname{card}V^*_{\mathrm{aff}}(\mathbb F_{q^\nu})
=\frac1{q^\nu}
\sum_{x\in(\mathbb F_{q^\nu}^*)^{n+2}}
\left\{1+\sum_{x_0\in\mathbb F_{q^\nu}^*}\chi(x_0f(x))\right\}
\end{equation}

et, compte tenu de ce que

\begin{equation}\tag{4.0.6}
\operatorname{card}V^*_{\mathrm{aff}}(\mathbb F_{q^\nu})
=(q^\nu-1)\operatorname{card}V^*(\mathbb F_{q^\nu})
\end{equation}

(4.0.5) se récrit

\begin{equation}\tag{4.0.7}
(q^\nu-1)\operatorname{card}V^*(\mathbb F_{q^\nu})
=\frac{(q^\nu-1)^{n+2}}{q^\nu}
+\frac1{q^\nu}
\sum_{(x_0,x)\in(\mathbb F_{q^\nu}^*)^{n+3}}\chi(x_0f(x)).
\end{equation}

Écrivons explicitement la forme $f(X)$ comme somme de monômes

\begin{equation}\tag{4.0.8}
f(X)=\sum a_wX^w,
\end{equation}

où $w=(w_1,\ldots,w_{n+2})$ avec $\sum w_i=d$,
$X^w\overset{\mathrm{dfn}}{=}X_1^{w_1}\cdots X_{n+2}^{w_{n+2}}$, de sorte
que (4.0.7) se récrit

\begin{equation}\tag{4.0.9}
(q^\nu-1)\operatorname{card}V^*(\mathbb F_{q^\nu})
=\frac{(q^\nu-1)^{n+2}}{q^\nu}
+\frac1{q^\nu}
\sum_{(x_0,x)\in(\mathbb F_{q^\nu}^*)^{n+3}}
\prod_w\chi(a_wx_0x^w).
\end{equation}

4.1. Il faut maintenant expliciter notre ``caractère additif à valeurs en
caractéristique zéro''. On a envie de prendre simplement

\SGAvar{\begin{equation}\tag{4.1.0}
\chi:x\longmapsto
\exp\!\left(2\pi i\,\operatorname{tr}_{\mathbb F_{q^\nu}/\mathbb F_p}(x)\right.
=\xi_p^{\operatorname{tr}_{\mathbb F_{q^\nu}/\mathbb H_p}(x)}.
\end{equation}

}{\begin{equation}\tag{4.1.0}
\chi:x\longmapsto
\exp\!\left(2\pi i\,\operatorname{tr}_{\mathbb F_{q^\nu}/\mathbb F_p}(x)/p\right)
=\xi_p^{\operatorname{tr}_{\mathbb F_{q^\nu}/\mathbb F_p}(x)}.
\end{equation}

}{F7II-XXII-033}\SGApage{II}{416}

Désignons par

\begin{align}
\tag{4.1.1}
\Omega&\quad\text{la clôture algébrique de }\mathbb Q_p\\
\tag{4.1.2}
\operatorname{ord}:\Omega^*&\longrightarrow\mathbb Q
\quad\text{l'ordinal, normalisé par }\operatorname{ord}(p)=1\\
\tag{4.1.3}
\xi_p&\in\Omega
\quad\text{une racine $p$-ième primitive de l'unité.}
\end{align}

On constate (cf. [2]) qu'il existe

\begin{equation}\tag{4.1.4}
\pi\in\mathbb Q_p(\xi_p),\qquad
\operatorname{ord}(\pi)=\frac1{p-1}
\end{equation}

tel que

\begin{equation}\tag{4.1.5}
\sum_{n\geq0}\frac{\pi^{p^n}}{p^n}=0,
\end{equation}

\SGAvar{(NB. On constate aussitôt que le premier membre converge dans
$Q_p(\xi_p)$.)

}{(NB. On constate aussitôt que le premier membre converge dans
$\mathbb Q_p(\xi_p)$.)

}{F7II-XXII-034}On définit \underline{l'exponentielle d'Artin-Hasse}

\SGAvar{\begin{equation}\tag{4.1.6}
E(X)=\exp\!\left(\sum_{n\geq0}\frac{X^{p^n}}{p^n}\right)
\in Q_p[[X]]
\end{equation}

}{\begin{equation}\tag{4.1.6}
E(X)=\exp\!\left(\sum_{n\geq0}\frac{X^{p^n}}{p^n}\right)
\in\mathbb Q_p[[X]]
\end{equation}

}{F7II-XXII-035}qui, oh miracle, se trouve dans $\mathbb Z_p[[X]]$. On définit

\begin{equation}\tag{4.1.7}
\Theta(X)=E(\pi X),
\end{equation}

dont les propriétés fondamentales sont résumées dans le lemme suivant :

\underline{Lemme} 4.2.

\SGAvar{\begin{equation}\tag{4.2.1}
\Theta(X)=\sum B_nX^n,\qquad B_0=1\quad A_1=\pi,\quad
\operatorname{ord}(B_n)\geq\frac n{p-1}.
\end{equation}

\begin{equation}\tag{4.2.2}
\Theta(1)=\xi_p,\quad
\text{\uline{une racine $p$'ième primitive de l'unité.}}
\end{equation}

(4.2.3) \underline{Soit $t\in\Omega$, et supposons que, pour un entier
$a>0$ donné, $t=t^{p^a}$ (i.e. $t$ est un ``représentant de Teichmüller''
dans $\mathbb W(\mathbb F_{q^a})\subset\Omega$). Alors
$\sum_{j=0}^{a-1}t^{p^j}\in\mathbb Z_p$, de sorte qu'on peut définir la
série formelle}

}{\begin{equation}\tag{4.2.1}
\Theta(X)=\sum B_nX^n,\qquad B_0=1\quad B_1=\pi,\quad
\operatorname{ord}(B_n)\geq\frac n{p-1}.
\end{equation}

\begin{equation}\tag{4.2.2}
\Theta(1)=\xi_p,\quad
\text{\uline{une racine $p$-ième primitive de l'unité.}}
\end{equation}

(4.2.3) \underline{Soit $t\in\Omega$, et supposons que, pour un entier
$a>0$ donné, $t=t^{p^a}$ (i.e. $t$ est un ``représentant de Teichmüller''
dans $\mathbb W(\mathbb F_{p^a})\subset\Omega$). Alors
$\sum_{j=0}^{a-1}t^{p^j}\in\mathbb Z_p$, de sorte qu'on peut définir la
série formelle}

}{F7II-XXII-036}\SGApage{II}{417}

\[
(\Theta(X))^{\sum_{j=0}^{a-1}t^{p^j}},
\]

et on a une identité de séries formelles

\SGAvar{\begin{equation}\tag{4.2.3.1}
\prod_{j=0}^{a-1}\Theta(t^{p^j}X)
=(\Theta(X))^{\sum_{j=0}^{a-1}t^j}.
\end{equation}

\underline{Démonstration.} Compte tenu de ce qui
$E(X)\in\mathbb Z_p[[X]]$, (4.2.1) se déduit de la définition. Quant à
(4.2.2), on a

}{\begin{equation}\tag{4.2.3.1}
\prod_{j=0}^{a-1}\Theta(t^{p^j}X)
=(\Theta(X))^{\sum_{j=0}^{a-1}t^{p^j}}.
\end{equation}

\underline{Démonstration.} Compte tenu de ce que
$E(X)\in\mathbb Z_p[[X]]$, (4.2.1) se déduit de la définition. Quant à
(4.2.2), on a

}{F7II-XXII-037}\begin{equation}\tag{4.2.4}
\log(\Theta(1))=0
\qquad\text{par le choix de $\pi$ (4.1.5)}
\end{equation}

et

\begin{equation}\tag{4.2.5}
\operatorname{ord}(\Theta(1)-1)=\frac1{p-1}
\qquad\text{par (4.2.1).}
\end{equation}

Par (4.2.4), il existe un entier $n\geq0$ tel que
$\Theta(1)^{p^n}=1$, (c'est la caractérisation du noyau du logarithme) et
(4.2.5) entraîne que $n=1$.

\SGAvar{Quant à (4.2.3.1), en prenant les logarithmes, il résulte d'un calcul évident,
une fois prouvé que l'on a bien
$\sum_{j=0}^{a-1}t^{p^j}\in\mathbb Z_p$. Or on voit aussitôt que les
$t^{p^j}$ ($0\leq j\leq a-1$) sont les conjugués de
$t\in\mathbb W(\mathbb F_{q^a})$ sur $\mathbb Z_p$, donc

\[
\sum_{j=0}^{a-1}t^{p^j}
=\operatorname{Tr}_{\mathbb W(\mathbb F_{q^a})/\mathbb Z_p}(t)\in\mathbb Z_p.
\]

}{Quant à (4.2.3.1), en prenant les logarithmes, il résulte d'un calcul évident,
une fois prouvé que l'on a bien
$\sum_{j=0}^{a-1}t^{p^j}\in\mathbb Z_p$. Or on voit aussitôt que les
$t^{p^j}$ ($0\leq j\leq a-1$) sont les conjugués de
$t\in\mathbb W(\mathbb F_{p^a})$ sur $\mathbb Z_p$, donc

\[
\sum_{j=0}^{a-1}t^{p^j}
=\operatorname{Tr}_{\mathbb W(\mathbb F_{p^a})/\mathbb Z_p}(t)\in\mathbb Z_p.
\]

}{F7II-XXII-038}En prenant la valeur en $X=1$ des deux membres de (4.2.3.1) nous trouvons le

\underline{Corollaire} 4.3. \underline{Soit $t\in\Omega$ vérifiant $t^{p^a}=t$.
Désignons par $\bar t\in\mathbb F_{p^a}$ sa classe résiduelle. Alors on a}

\begin{equation}\tag{4.3.1}
\prod_{j=0}^{a-1}\Theta(t^{p^j})
=\xi_p^{\operatorname{tr}_{\mathbb F_{p^a}/\mathbb F_p}(\bar t)}.
\end{equation}

\SGApage{II}{418}

4.4. Rappelons que le corps résiduel de l'anneau des entiers de
$\mathbb Q_p(\xi_p)$ est $\mathbb F_p$. Soient

\begin{equation}\tag{4.4.0}
K=\text{l'extension non-ramifiée de $\mathbb Q_p(\xi_p)$ à corps résiduel
$\mathbb F_q$},
\end{equation}

\SGAvar{\begin{equation}\tag{4.4.1}
\tau\in\Gal(K/Q_p(\xi_p))
\end{equation}

}{\begin{equation}\tag{4.4.1}
\tau\in\Gal(K/\mathbb Q_p(\xi_p))
\end{equation}

}{F7II-XXII-039}le "Frobenius", i.e. l'automorphisme se réduisant suivant l'automorphisme
de Frobenius de $\mathbb F_q/\mathbb F_p$, $x\longmapsto x^p$,

\begin{equation}\tag{4.4.2}
\text{Soit }F(X)=\sum A_WX^W
\text{ le relèvement en $K[X_1,\ldots,X_{n+2}]$ de $f(X)$}
\end{equation}

\SGAvar{dont les coefficients vérifient $A_W^q=A_W$. Mettant ensemble (4.3.1) et
(4.0.9), on trouve (compte tenu de la définition (4.1.0) de $X$) :

}{dont les coefficients vérifient $A_W^q=A_W$. Mettant ensemble (4.3.1) et
(4.0.9), on trouve (compte tenu de la définition (4.1.0) de $\chi$) :

}{F7II-XXII-040}\begin{equation}\tag{4.4.3}
\left\{
\begin{aligned}
&(q^\nu-1)\operatorname{card}V^*(\mathbb F_{q^\nu})
-\frac{(q^\nu-1)^{n+2}}{q^\nu}\\
&\quad=\frac1{q^\nu}
\sum_{(X_0,X)\in(\mu_{q^\nu-1}(\Omega))^{n+3}}
\prod_W\prod_{j=0}^{\operatorname{ord}(q^\nu)-1}
\Theta\!\left(A_W^{p^j}X_0^{p^j}X^{p^jW}\right).
\end{aligned}
\right.
\end{equation}

Définissons

\[
H(X)=H(X_0,\ldots,X_{n+2})\in K[[X_0,\ldots,X_{n+2}]]
\]

par

\begin{equation}\tag{4.4.4}
H(X)=\prod_W
\Theta\!\left(A_WX_0X_1^{W_1}\cdots X_{n+2}^{W_{n+2}}\right).
\end{equation}

Nous faisons agir $\Gal(K/\mathbb Q_p(\xi_p))$ sur
$K[[X_0,\ldots,X_{n+2}]]$ à travers son action sur les coefficients ;
l'action composante de $\sigma\in\Gal(K/\mathbb Q_p(\xi_p))$ sera notée

\begin{equation}\tag{4.4.5}
G(X)\longmapsto G^\sigma(X).
\end{equation}

\SGApage{II}{419}

On a en particulier

\SGAvar{\[
H^{\tau^j}(X)=\prod_W
\Theta\!\left(A_W^{\tau^j}X_0X_1^{W}\cdots
X_{n+2}^{W_{n+2}}\right).
\]

}{\[
H^{\tau^j}(X)=\prod_W
\Theta\!\left(A_W^{\tau^j}X_0X_1^{W_1}\cdots
X_{n+2}^{W_{n+2}}\right).
\]

}{F7II-XXII-041}Compte tenu de ce que $A_W^{\tau^j}=A_W^{p^j}$, (4.4.3) se récrit alors

\begin{equation}\tag{4.4.6}
\begin{aligned}
(q^\nu-1)\operatorname{Card}V^*(\mathbb F_{q^\nu})
&=\frac{(q^\nu-1)^{n+2}}{q^\nu}\\
&\quad+\frac1{q^\nu}
\sum_{(X_0,X)\in(\mu_{q^\nu-1}(\Omega))^{n+3}}
\prod_{j=0}^{\operatorname{ord}(q^\nu)-1}H^{\tau^j}(X^{p^j})
\end{aligned}
\end{equation}

4.5. Nous allons maintenant exprimer la somme dans (4.4.6) comme une trace.

Pour chaque $b\in\mathbb Q$, on définit une $K$-algèbre de Banach $p$-adique

\begin{equation}\tag{4.5.1}
L(b)=\text{le sous-anneau de $K[[X_0,\ldots,X_{n+2}]]$ formé des séries}
\end{equation}

qui s'écrivent

\begin{equation}\tag{4.5.2}
\sum_{dU_0=\sum_{i\geq1}U_i}
B_UX_0^{U_0}X_1^{U_1}\cdots X_{n+2}^{U_{n+2}}
\qquad (U=(U_0,\ldots,U_{n+2}))
\end{equation}

et qui vérifient

\begin{equation}\tag{4.5.3}
\inf_U\bigl(\operatorname{ord}(B_U)-bU_0\bigr)>-\infty.
\end{equation}

C'est une algèbre de Banach $p$-adique pour la fonction d'ordinal
(l'opposé du logarithme de la norme)

\begin{equation}\tag{4.5.4}
\inf_U\bigl(\operatorname{ord}(B_U)-bU_0\bigr)
=\operatorname{ord}\!\left(\sum B_UX^U\right).
\end{equation}

On vérifie sur 4.4.4 et (4.2.1) que l'on a

\[
H(X)\in L\!\left(\frac1{p-1}\right).
\]

\SGApage{II}{420}

Introduisons des opérateurs.

Pour chaque entier $k>0$, on introduit un opérateur continu d'espaces de
Banach sur $K$

\begin{equation}\tag{4.5.5}
\left\{
\begin{aligned}
\psi_k:L(b)&\longrightarrow L(kb)\\
\psi_k\!\left(\sum B_UX^U\right)&=\sum B_{kU}X^U.
\end{aligned}
\right.
\end{equation}

Pour chaque couple $(b,b')$, $b'<b$, il y a un
\underline{opérateur de} "\underline{restriction}" qui est complètement continu
(i.e. une limite uniforme des opérateurs de rang fini)

\begin{equation}\tag{4.5.6}
\left\{
\begin{aligned}
L(b)&\xrightarrow{\ \text{"restriction"}\ }L(b')\\
\sum B_UX^U&\longmapsto\sum B_UX^U.
\end{aligned}
\right.
\end{equation}

Chaque élément $G$ de l'algèbre $L(b)$ donne l'endomorphisme continu de la
multiplication par $G$

\begin{equation}\tag{4.5.7}
\begin{aligned}
G:L(b)&\longrightarrow L(b)\\
\eta&\longmapsto\eta G.
\end{aligned}
\end{equation}

Ces opérateurs sont liés par le lemme suivant, dont la démonstration se
trouve dans [2] (cf. aussi [4]).

\SGApage{II}{421}

\underline{Lemme} 4.5.8 (\underline{ou lemme de la trace} (DWORK)).
\underline{Soient} $b\in\mathbb Q$, $b>0$, $k\in\mathbb Z$, $k\geq2$, \underline{et}
$G\in L(b/k)$. \underline{Alors le composé}

\[
L(b)\xrightarrow{\ \text{"restriction"}\ }L(b/k)
\xrightarrow{\ \times G\ }L(b/k)
\xrightarrow{\ \psi_k\ }L(b),
\]

\underline{noté} $\psi_k\circ G$, \underline{est complètement continu, et sa trace}
[cf. [4]] \underline{est donnée par la formule}

\begin{equation}\tag{4.5.9}
\operatorname{tr}(\psi_k\circ G)
=\frac1{(k-1)^{n+2}}
\sum_{X\in(\mu_{k-1}(\Omega))^{n+3}}G(X).
\end{equation}

Appliquant ce lemme au cas

\[
\left\{
\begin{aligned}
k&=q^\nu\\
G(X)&=\prod_{j=0}^{\operatorname{ord}(q^\nu)-1}
H^{\tau^j}(X^{p^j})
\in L\!\left(\frac{p}{(p-1)q^\nu}\right)\\
b&=\frac1{p-1}
\end{aligned}
\right.
\]

(4.4.6) se récrit

\SGAvar{\begin{equation}\tag{4.5.10}
\begin{aligned}
(q-1)\operatorname{card}V^*(\mathbb F_{q^\nu})
&=\frac{(q^\nu-1)^{n+2}}{q^\nu}\\
&\quad+\frac{(q^\nu-1)^{n+3}}{q^\nu}
\operatorname{tr}\!\left(
\left.\psi_q
\prod_{j=0}^{\operatorname{ord}(q^\nu)-1}H^{\tau^j}(X^{p^j})
\right|L\!\left(\frac1{p-1}\right)\right).
\end{aligned}
\end{equation}

Soit $L^0(b)$ l'idéal des éléments de $L(b)$ qui "s'annulent en $0$",
i.e. "sans terme constant". Il est stable sous $\psi\circ G$, et on a
évidemment

}{\begin{equation}\tag{4.5.10}
\begin{aligned}
(q^\nu-1)\operatorname{card}V^*(\mathbb F_{q^\nu})
&=\frac{(q^\nu-1)^{n+2}}{q^\nu}\\
&\quad+\frac{(q^\nu-1)^{n+3}}{q^\nu}
\operatorname{tr}\!\left(
\left.\psi_{q^\nu}
\prod_{j=0}^{\operatorname{ord}(q^\nu)-1}H^{\tau^j}(X^{p^j})
\right|L\!\left(\frac1{p-1}\right)\right).
\end{aligned}
\end{equation}

Soit $L^0(b)$ l'idéal des éléments de $L(b)$ qui "s'annulent en $0$",
i.e. "sans terme constant". Il est stable sous $\psi_k\circ G$, et on a
évidemment

}{F7II-XXII-042}\begin{equation}\tag{4.5.11}
\operatorname{tr}(\psi_k\circ G)
=G(0)+\operatorname{tr}\!\left(
\left.\psi_k\circ G\right|L^0\right).
\end{equation}

Compte tenu de ce que $H(0)=1$, (4.5.10) se récrit

\SGAvar{\begin{equation}\tag{4.5.12}
\begin{aligned}
\operatorname{card}V^*(\mathbb F_{q^\nu})
&=(q^\nu-1)^{n+1}\\
&\quad+\frac{(q^\nu-1)^{n+2}}{q}
\operatorname{tr}\!\left(
\left.\psi_q\circ
\prod_{j=0}^{\operatorname{ord}(q^\nu)-1}H^{\tau^j}(X^{p^j})
\right|L^0\!\left(\frac1{p-1}\right)\right).
\end{aligned}
\end{equation}

}{\begin{equation}\tag{4.5.12}
\begin{aligned}
\operatorname{card}V^*(\mathbb F_{q^\nu})
&=(q^\nu-1)^{n+1}\\
&\quad+\frac{(q^\nu-1)^{n+2}}{q^\nu}
\operatorname{tr}\!\left(
\left.\psi_{q^\nu}\circ
\prod_{j=0}^{\operatorname{ord}(q^\nu)-1}H^{\tau^j}(X^{p^j})
\right|L^0\!\left(\frac1{p-1}\right)\right).
\end{aligned}
\end{equation}

}{F7II-XXII-043}\SGApage{II}{422}

Considérons maintenant l'opérateur $\tau^{-1}$-linéaire

\begin{equation}\tag{4.5.13}
\tau^{-1}\circ\psi_p\circ H(X):L(b)\longrightarrow L(b).
\end{equation}

On constate que ses itérés sont donnés par

\SGAvar{\begin{equation}\tag{5.4.14}
(\tau^{-1}\circ\psi_p\circ H(X))^a
=\tau^{-a}\circ\psi_{p^a}\circ
H(X)H^\tau(X^p)\cdots H^{\tau^{a-1}}(X^{p^{a-1}}),
\end{equation}

}{\begin{equation}\tag{4.5.14}
(\tau^{-1}\circ\psi_p\circ H(X))^a
=\tau^{-a}\circ\psi_{p^a}\circ
H(X)H^\tau(X^p)\cdots H^{\tau^{a-1}}(X^{p^{a-1}}),
\end{equation}

}{F7II-XXII-044}et on note que, si $\operatorname{ord}(q)\mid a$, on a
$\tau^{-a}=\operatorname{id}$. Ceci posé, on définit une opération
$\tau^{-1}$-linéaire

\begin{equation}\tag{4.5.15}
\left\{
\begin{aligned}
\alpha:L^0\!\left(\frac1{p-1}\right)
&\longrightarrow L^0\!\left(\frac1{p-1}\right)\\
\alpha&=\frac1p\,\tau^{-1}\circ\psi_p\circ H(X),
\end{aligned}
\right.
\end{equation}

de sorte que (4.5.12) se récrit

\SGAvar{\begin{equation}\tag{4.5.16}
\operatorname{card}V^*(\mathbb F_{q^\nu})
=(q^\nu-1)^{n+1}
+(q^\nu-1)^{n+2}
\operatorname{tr}\!\left(
\left.(\alpha)^{\operatorname{ord}(q^\nu)}
\right|L^0\!\left(\frac1{p-1}\right)\right..
\end{equation}

}{\begin{equation}\tag{4.5.16}
\operatorname{card}V^*(\mathbb F_{q^\nu})
=(q^\nu-1)^{n+1}
+(q^\nu-1)^{n+2}
\operatorname{tr}\!\left(
\left.(\alpha)^{\operatorname{ord}(q^\nu)}
\right|L^0\!\left(\frac1{p-1}\right)\right).
\end{equation}

}{F7II-XXII-045}Désignons par

\begin{equation}\tag{4.5.17}
S=\text{l'ensemble }\{1,\ldots,n+2\}.
\end{equation}

Pour tout sous-ensemble $A\subset S$, $A\ne\varnothing$, on définit un
sous-espace de $L(b)$ :

\begin{equation}\tag{4.5.18}
\begin{aligned}
L^A(b)={}&\text{ensemble des séries $\sum B_UX^U\in L(b)$ telles que, pour chaque
monôme $X^U$ avec $B_U\ne0$,}\\
&\text{les $U_i>0$ sont exactement $U_0$ et les $U_j$, $j\in A$}.
\end{aligned}
\end{equation}

On a

\begin{equation}\tag{4.5.19}
L^0(b)=\bigoplus_{\substack{A\subset S\\A\ne\varnothing}}L^A(b)
\end{equation}

\SGApage{II}{423}

et, pour chaque sous-ensemble

\[
B\subset S,\qquad B\ne\varnothing,
\]

on a

\begin{equation}\tag{4.5.20}
\bigoplus_{A\supset B}L^A(b)
=L(b)\cap\text{ l'idéal }
\left(\prod_{i\in B}X_i\right)K[[X]].
\end{equation}

\SGAvar{Évidemment chaque sous-espace $\bigoplus_{A\subseteq B}L^A(b)$ est stable
par $\alpha$ ; il en résulte qu'il en est de même pour
$\bigoplus_{A\supsetneq B}L^A(b)$, qui est somme de sous-espaces du type
précédent.

}{Évidemment chaque sous-espace $\bigoplus_{A\supset B}L^A(b)$ est stable
par $\alpha$ ; il en résulte qu'il en est de même pour
$\bigoplus_{A\supsetneq B}L^A(b)$, qui est somme de sous-espaces du type
précédent.

}{F7II-XXII-046}On définit

\begin{equation}\tag{4.5.21}
\alpha_B=\text{l'endomorphisme $\tau^{-1}$-linéaire induit par $\alpha$ sur}
\quad
\frac{\displaystyle\bigoplus_{A\supset B}
L^A\!\left(\frac1{p-1}\right)}
{\displaystyle\bigoplus_{A\supsetneq B}
L^A\!\left(\frac1{p-1}\right)}
\simeq L^B\!\left(\frac1{p-1}\right).
\end{equation}

Explicitement, on obtient $\alpha_B$ par la formule

\begin{equation}\tag{4.5.22}
\alpha_B=\frac{\tau^{-1}}p\,\psi_p\,H_B(X),
\end{equation}

opération sous laquelle $L^B$ est stable, où $H_B(X)$ se déduit de $H(X)$
en y faisant la substitution $X_i=0$ pour $i\notin\{0\}\cup B$.

Récrivons encore (4.5.16)

\SGAvar{\begin{equation}\tag{4.5.2.3}
\operatorname{card}V^*(\mathbb F_{q^\nu})
=(q^\nu-1)^{n+1}
+(q^\nu-1)^{n+2}
\sum_{\substack{A\subset S\\A\ne\varnothing}}
\operatorname{tr}\!\left((\alpha_A)^{\operatorname{ord}(q^\nu)}\right).
\end{equation}

}{\begin{equation}\tag{4.5.23}
\operatorname{card}V^*(\mathbb F_{q^\nu})
=(q^\nu-1)^{n+1}
+(q^\nu-1)^{n+2}
\sum_{\substack{A\subset S\\A\ne\varnothing}}
\operatorname{tr}\!\left((\alpha_A)^{\operatorname{ord}(q^\nu)}\right).
\end{equation}

}{F7II-XXII-047}Définissons, pour $A\subset S$, $A\ne\varnothing$, l'hypersurface

\begin{equation}\tag{4.5.24}
V_A\subset\mathbb P^{\operatorname{card}(A)-1}
\simeq
\mathbb P^{n+1}\cap
\{X_i=0\ \forall i\notin A\}
\end{equation}

\SGAvar{définie par $f_A(X)$, le polynôme déduit de $f(X)$ y faisant la
substitution en $X_i=0$ pour $i\notin A$.

}{définie par $f_A(X)$, le polynôme déduit de $f(X)$ en y faisant la
substitution $X_i=0$ pour $i\notin A$.

}{F7II-XXII-048}\SGApage{II}{424}

Compte tenu de la compatibilité de la formation de la série $H$ à partir de
$f(X)$ à l'opération de substitution de 0 pour certaines des variables
(4.5.23) donne, pour chaque $A\subseteq S$, $A\ne\varnothing$

\SGAvar{\begin{equation}\tag{4.5.25}
\operatorname{card}(V_A^*(\mathbb F_{q^\nu}))
=(q^\nu-1)^{\operatorname{car}(A)-1}
+(q^\nu-1)^{\operatorname{card}(A)}
\sum_{B\subset A}
\operatorname{tr}\!\left((\alpha_B))^{\operatorname{ord}(q^\nu)}\right).
\end{equation}

}{\begin{equation}\tag{4.5.25}
\operatorname{card}(V_A^*(\mathbb F_{q^\nu}))
=(q^\nu-1)^{\operatorname{card}(A)-1}
+(q^\nu-1)^{\operatorname{card}(A)}
\sum_{B\subset A}
\operatorname{tr}\!\left((\alpha_B)^{\operatorname{ord}(q^\nu)}\right).
\end{equation}

}{F7II-XXII-049}D'autre part, on a

\begin{equation}\tag{4.5.26}
\operatorname{card}V(\mathbb F_{q^\nu})
=\sum_{\substack{A\subset S\\A\ne\varnothing}}
\operatorname{card}V_A^*(\mathbb F_{q^\nu}).
\end{equation}

Après une sommation facile, (4.5.25) donne alors

\begin{equation}\tag{4.5.27}
\begin{aligned}
\operatorname{card}V(\mathbb F_{q^\nu})
={}&1+q^\nu+q^{2\nu}+\cdots
+q^{\nu(\operatorname{card}(S)-1)}\\
&+\sum_{\substack{B\subseteq S\\B\ne\varnothing}}
(q^\nu-1)^{\operatorname{card}(B)}
q^{\nu\operatorname{card}(S-B)}
\operatorname{tr}\!\left((\alpha_B)^{\operatorname{ord}(q^\nu)}\right).
\end{aligned}
\end{equation}

\SGAvar{Notons que $\alpha$ donc les $\alpha_B$ étant $\tau^{-1}$ linéaire, les
$\alpha_B^{\operatorname{ord}(q)}$ sont $K$-linéaires. Définissons alors les
séries caractéristiques

}{Notons que $\alpha$ donc les $\alpha_B$ étant $\tau^{-1}$-linéaire, les
$\alpha_B^{\operatorname{ord}(q)}$ sont $K$-linéaires. Définissons alors les
séries caractéristiques

}{F7II-XXII-050}\[
\Delta_B(t)\in1+tK[[t]]
\]

\begin{equation}\tag{4.5.28}
\Delta_B(t)=\det\!\left(
1-t\alpha_B^{\operatorname{ord}(q)}
\mathrel{\Big|}L^B\!\left(\frac1{p-1}\right)\right).
\end{equation}

Introduisons des endomorphismes de $1+tK[[t]]$

\begin{equation}\tag{4.5.29}
\varphi:g(t)\longmapsto g^\varphi(t)=g(qt)
\end{equation}

et

\begin{equation}\tag{4.5.30}
\delta=1-\varphi:g(t)\longmapsto
g^\delta(t)=g(t)/g(qt)
\end{equation}

\SGApage{II}{425}

Compte tenu de ce que

\begin{equation}\tag{4.5.31}
Z(t,V/\mathbb F_q)
=\exp\!\left(
\sum_{\nu\geq1}\operatorname{card}(V(\mathbb F_{q^\nu}))
\frac{t^\nu}{\nu}\right)
\end{equation}

et

\begin{equation}\tag{4.5.32}
\det\!\left(1-t\alpha_B^{\operatorname{ord}(q)}\right)
=\exp\!\left(
-\sum_{\nu\geq1}
\operatorname{tr}\!\left((\alpha_B)^{\nu\operatorname{ord}q}\right)
\frac{t^\nu}{\nu}\right)
\end{equation}

(4.5.27) se récrit

\SGAvar{\begin{equation}\tag{4.5.23}
Z(t,V/\mathbb F_q)\prod_{i=0}^{n+1}(1-q^it)
=\prod_{j=1}^{n+2}
\left(
\prod_{\substack{B\subset S\\\operatorname{card}(B)=j}}
\Delta_B(t)
\right)^{-(-\delta)^j\varphi^{n+2-j}}.
\end{equation}

}{\begin{equation}\tag{4.5.33}
Z(t,V/\mathbb F_q)\prod_{i=0}^{n+1}(1-q^it)
=\prod_{j=1}^{n+2}
\left(
\prod_{\substack{B\subset S\\\operatorname{card}(B)=j}}
\Delta_B(t)
\right)^{-(-\delta)^j\varphi^{n+2-j}}.
\end{equation}

}{F7II-XXII-051}4.6. Nous allons étudier maintenant l'intégralité des séries $\Delta_B$.
Posons

\begin{equation}\tag{4.6.1}
\mathcal O=\text{l'anneau des entiers de $K$}
\end{equation}

et définissons une \underline{filtration} sur $1+t\mathcal O[[t]]$ par

\begin{equation}\tag{4.6.2}
F^i(1+t\mathcal O[[t]])
=1+q^it\mathcal O[[q^it]]
=\varphi^i(1+t\mathcal O[[t]]).
\end{equation}

On a

\begin{align}
\varphi(F^i)&\subset F^{i+1} \tag{4.6.3}\\
\delta(F^i)&\subset F^i \tag{4.6.4}\\
\delta&\text{ induit l'identité sur }F^i/F^{i+1}. \tag{4.6.5}
\end{align}

\underline{Lemme} 4.6.6.
$\Delta_B(t)\in1+t\mathcal O[[t]]$, \underline{et si un entier} $j$ \underline{vérifie}

\begin{equation}\tag{4.6.7}
jd<\operatorname{card}(B)
\end{equation}

\underline{alors on a}

\begin{equation}\tag{4.6.8}
\Delta_B(t)\in F^j(1+t\mathcal O[[t]]).
\end{equation}

\SGApage{II}{426}

\underline{Démonstration.}
Les espaces de Banach
$L^B\!\left(\frac1{p-1}\right)$ sont tous \og du type $B(I)$\fg{} au sens de
Serre [7], i.e. ils admettent des \og bases orthonormales\fg{}, à savoir, les
monômes

\begin{equation}\tag{4.6.9}
\pi^{U_0}X_0^{U_0}X_1^{U_1}\cdots X_{n+2}^{U_{n+2}}
\end{equation}

\SGAvar{\noindent
qui y appartiennent. Que c' une base \og orthonormale\fg{} veut dire par
définition qu'on a

}{\noindent
qui y appartiennent. Que c'est une base \og orthonormale\fg{} veut dire par
définition qu'on a

}{F7II-XXII-052}\begin{equation}\tag{4.6.10}
\operatorname{ord}\!\left(
\sum C_U\pi^{U_0}X_0^{U_0}\cdots X_{n+2}^{U_{n+2}}
\right)
=\min_U\operatorname{ord}(C_U).
\end{equation}

Soit

\SGAvar{\begin{equation}\tag{4.6.11}
\mathbf R_B=\text{le $\mathcal O$-sous-module de
$L^B\!\left(\frac1{p-1}\right)$ formé des éléments d'ordinal $\geq0$}.
\end{equation}
Pour démontrer le lemme, il suffit de démontrer

}{\begin{equation}\tag{4.6.11}
\mathbf R_B=\text{le $\mathcal O$-sous-module de
$L^B\!\left(\frac1{p-1}\right)$ formé des éléments d'ordre $\geq0$}.
\end{equation}
Pour démontrer le lemme, il suffit de démontrer

}{F7II-XXII-053}\begin{equation}\tag{4.6.12}
\text{si }jd<\operatorname{card}(B),\text{ alors }
(\alpha_B)^{\operatorname{ord}(q)}R_B\subset q^jR_B
\end{equation}

ou, a fortiori, de démontrer

\begin{equation}\tag{4.6.13}
\text{si }jd<\operatorname{card}(B),\text{ alors }
\alpha_B(R_B)\subset p^jR_B.
\end{equation}

Or, on a (4.5.22)

\begin{equation}\tag{4.6.14}
\alpha_B=\frac{\tau^{-1}}p\,\psi_p\circ H_B(X).
\end{equation}

Compte tenu de ce que

\begin{equation}\tag{4.6.15}
H(X)\in L\!\left(\frac1{p-1}\right),
\qquad \operatorname{ord}(H(X))\geq0
\end{equation}

on a

\begin{equation}\tag{4.6.16}
H_B(X)\cdot R_B\subset R_B
\end{equation}

et évidemment

\begin{equation}\tag{4.6.17}
\tau^{-1}(R_B)=R_B,
\end{equation}

\SGApage{II}{427}

de sorte que (4.6.13) sera conséquence de

\begin{equation}\tag{4.6.18}
\text{si }jd<\operatorname{card}(B),\text{ alors }
\psi_p(R_B)\subset p^{j+1}R_B.
\end{equation}

Vérifions (4.6.18) monôme par monôme. On a

\SGAvar{\begin{equation}\tag{4.6.19}
\begin{cases}
\psi_p\!\left(\pi^{U_0}X^{U_0}\cdots X_{n+2}^{U_{n+2}}\right)=0,
&\text{sauf si }\forall i\quad p\mid U_i\\
\text{et}\\[3pt]
\psi_p\!\left(
\pi^{pU_0}X_0^{pU_0}\cdots X_{n+2}^{pU_{n+2}}
\right)
=\pi^{(p-1)U_0}\cdot
\underbrace{\pi^{U_0}X_0^{U_0}\cdots X_{n+2}^{U_{n+2}}}_{\in R_B}.
\end{cases}
\end{equation}

}{\begin{equation}\tag{4.6.19}
\begin{cases}
\psi_p\!\left(\pi^{U_0}X_0^{U_0}\cdots X_{n+2}^{U_{n+2}}\right)=0,
&\text{sauf si }\forall i\quad p\mid U_i\\
\text{et}\\[3pt]
\psi_p\!\left(
\pi^{pU_0}X_0^{pU_0}\cdots X_{n+2}^{pU_{n+2}}
\right)
=\pi^{(p-1)U_0}\cdot
\underbrace{\pi^{U_0}X_0^{U_0}\cdots X_{n+2}^{U_{n+2}}}_{\in R_B}.
\end{cases}
\end{equation}

}{F7II-XXII-054}Or, les conditions pour que l'on ait
$X_0^{U_0}\cdots X_{n+2}^{U_{n+2}}
\in L^B\!\left(\frac1{p-1}\right)$ sont

\begin{equation}\tag{4.6.20}
\begin{cases}
dU_0=\displaystyle\sum_{i\geq1}U_i,\\
\text{pour }i\geq1,\text{ on a }U_i>0\Longleftrightarrow i\in B
\end{cases}
\end{equation}

d'où, en particulier,

\begin{equation}\tag{4.6.21}
dU_0\geq\operatorname{card}(B).
\end{equation}

On a donc

\begin{equation}\tag{4.6.22}
\operatorname{ord}\!\left(\pi^{(p-1)U_0}\right)
=U_0\geq\frac{\operatorname{card}(B)}d>j
\end{equation}

et, comme $U_0\in\mathbb Z$,

\begin{equation}\tag{4.6.23}
\operatorname{ord}\!\left(\pi^{(p-1)U_0}\right)
=U_0\geq j+1.
\end{equation}

Mettant ceci ensemble avec (4.6.19), on trouve (4.6.18).

\SGApage{II}{428}

\SGAsectiontoc{5. Un théorème d'AX [1], (cf. aussi [4], section 7).}{5}
\section*{5. \underline{Un théorème d'AX} [1], (cf. aussi [4], section 7).}

\underline{Théorème 5.0.}
\underline{Soient} $V$ \underline{une hypersurface dans} $\mathbb P^{n+1}$ \underline{de degré} $d>0$, \underline{définie
sur} $\mathbb F_q$, \underline{et} $j$ \underline{un entier} : \underline{si on a}

\begin{equation}\tag{5.0.1}
jd<n+2,
\end{equation}

\underline{alors}

\begin{equation}\tag{5.0.2}
Z(t,V/\mathbb F_q)\prod_{i=0}^{n+1}(1-q^it)
\in1+q^jt\mathbb Z[[q^jt]].
\end{equation}

et l'inverse de tout zéro ou pôle de

\begin{equation}\tag{5.0.3}
Z(t,V/\mathbb F_q)\prod_{i=0}^{n+1}(1-q^it)
\end{equation}

\noindent
\underline{est de la forme} $q^j\cdot(\text{\underline{un entier algébrique}})$.

\underline{Théorème 5.0.bis.}
\underline{Soit} $V_{\mathrm{aff}}$ \underline{une hypersurface dans} $\mathbb A^{n+2}$ \underline{définie
par une forme homogène de degré} $d$, \underline{définie sur} $\mathbb F_q$. \underline{Si on a}

\SGAvar{\begin{equation}\tag{5.0.7}
jd<n+2
\end{equation}

}{\begin{equation}\tag{5.0.4}
jd<n+2
\end{equation}

}{F7II-XXII-055}\underline{alors}

\begin{equation}\tag{5.0.5}
Z(t,V_{\mathrm{aff}}/\mathbb F_q)
\in1+q^jt\mathbb Z[[q^jt]].
\end{equation}

\underline{En particulier},

\begin{equation}\tag{5.0.6}
\operatorname{card}\bigl(V_{\mathrm{aff}}(\mathbb F_q)\bigr)
\equiv0\mod{q^j},
\end{equation}

\underline{et par suite}

\begin{equation}\tag{5.0.6 bis}
\operatorname{card}\bigl(V(\mathbb F_q)\bigr)
\geq\frac{q^j-1}{q-1}>0
\end{equation}

\noindent
(\underline{ce qui améliore le th. de CHEVALLEY--WARNING} [8]).

\SGApage{II}{429}

\SGAvar{\underline{Démonstration.}
L'implication (5.0.5) $\Rightarrow$ (5.0.6) est claire, en prenant le
coefficient de $t$ dans la fonction zêta. Les assertions (5.0.4) et (5.0.2)
sont équivalentes, compte tenu de ce que

}{\underline{Démonstration.}
L'implication (5.0.5) $\Rightarrow$ (5.0.6) est claire, en prenant le
coefficient de $t$ dans la fonction zêta. Les assertions (5.0.5) et (5.0.2)
sont équivalentes, compte tenu de ce que

}{F7II-XXII-056}\begin{equation}\tag{5.0.7}
\operatorname{card}V_{\mathrm{aff}}(\mathbb F_{q^\nu})
=1+(q^\nu-1)\operatorname{card}V(\mathbb F_{q^\nu}),
\end{equation}

d'où

\SGAvar{\begin{equation}\tag{5.0.8}
Z(t,V_{\mathrm{aff}}/\mathbb F_q)
=\frac{Z(t,V/\mathbb F_q)}
{Z(qt,V/\mathbb F_q)(1-t)}
\end{equation}

}{\begin{equation}\tag{5.0.8}
Z(t,V_{\mathrm{aff}}/\mathbb F_q)
=\frac{Z(qt,V/\mathbb F_q)}
{Z(t,V/\mathbb F_q)(1-t)}
\end{equation}

}{F7II-XXII-057}et

\SGAvar{\begin{equation}\tag{5.0.9}
Z(t,V/\mathbb F_q)
=\prod_{i\geq0}\left[(1-q^if)
Z(q^it,V_{\mathrm{aff}}/\mathbb F_q)\right].
\end{equation}

}{\begin{equation}\tag{5.0.9}
Z(t,V/\mathbb F_q)
=\prod_{i\geq0}\left[(1-q^it)
Z(q^it,V_{\mathrm{aff}}/\mathbb F_q)\right]^{-1}.
\end{equation}

}{F7II-XXII-058}Les énoncés (5.0.2) et (5.0.3) sont équivalents encore, par le théorème de
Fatou [3], grâce à la rationalité de la fonction zêta. Comme on a

\begin{equation}\tag{5.0.10}
Z(t,V/\mathbb F_q)\in\mathbb Z[[t]],
\end{equation}

il suffira de démontrer

\begin{equation}\tag{5.0.11}
Z(t,V/\mathbb F_q)\prod_{i=0}^{n+1}(1-q^it)
\in F^j(1+t\mathcal O[[t]]).
\end{equation}

Compte tenu de (4.5.33) et (4.6.4)--(4.6.5), il suffit de démontrer, pour
$B\subset\{1,\ldots,n+2\}$, $B\ne\varnothing$, qu'on a

\SGAvar{\begin{equation}\tag{5.0.12}
\Lambda_B(t)\in F^{j+\operatorname{card}(B)-(n+2)}.
\end{equation}

Nous pourrons terminer la démonstration en appliquant le lemme (4.4.6), une
fois démontrée l'assertion :

}{\begin{equation}\tag{5.0.12}
\Delta_B(t)\in F^{j+\operatorname{card}(B)-(n+2)}.
\end{equation}

Nous pourrons terminer la démonstration en appliquant le lemme (4.6.6), une
fois démontrée l'assertion :

}{F7II-XXII-059}\begin{equation}\tag{5.0.13}
\text{si }jd<n+2,\text{ alors }
\bigl(j+\operatorname{card}(B)-(n+2)\bigr)d
<\operatorname{card}(B),
\end{equation}

qui se récrit

\SGApage{II}{430}

\begin{equation}\tag{5.0.14}
\text{si }jd<n+2,\text{ alors }
(j-(n+2))d<\operatorname{card}(B)(1-d).
\end{equation}

Mais

\[
\inf_{\substack{B\subset\{1,\ldots,n+2\}\\B\ne\varnothing}}
\operatorname{card}(B)(1-d)
=(n+2)(1-d),
\]

de sorte qu'il suffit de démontrer

\begin{equation}\tag{5.0.15}
\text{si }jd<n+2,\text{ alors }
(j-(n+2))d<(n+2)(1-d),
\end{equation}

ce qui est vrai.

\underline{Remarque 5.1.}
La condition

\begin{equation}\tag{5.1.1}
jd<n+2
\end{equation}

admet une interprétation cohomologique. Soit $V$ une hypersurface \underline{lisse} de
degré $d$ et de dimension $n$. Posons, comme dans XI 2.1, pour $a+b=n$

\begin{equation}\tag{5.1.2}
h_0^{a,b}
=\dim H^b(V,\underline{\Omega}_V^a)-\delta_{a,b}
=\dim E^{a,b}(V)
\end{equation}

où $E^{a,b}(V)$ est la composante de bidegré $(a,b)$ de la \og cohomologie
évanescente de Hodge\fg{} (comparer XVII 5.4.4)

\SGAvar{\begin{equation}\tag{5.1.3}
E^{a,b}(V)
=\operatorname{Coker}\!\left(
H^b(\mathbb P^{n+1},\underline{\Omega}_{\mathbb P^{n+1}}^a)
\hookrightarrow H^b(V,\underline{\Omega}_V^a)\right..
\end{equation}

}{\begin{equation}\tag{5.1.3}
E^{a,b}(V)
=\operatorname{Coker}\!\left(
H^b(\mathbb P^{n+1},\underline{\Omega}_{\mathbb P^{n+1}}^a)
\hookrightarrow H^b(V,\underline{\Omega}_V^a)\right).
\end{equation}

}{F7II-XXII-060}Alors, par XI, 2.8,

\begin{equation}\tag{5.1.4}
jd<n+2
\Longleftrightarrow
\left[h_0^{a,n-a}=0\text{ pour }a<j\text{ et }a>n-j\right],
\end{equation}

i.e. la condition signifie que le \og niveau de Hodge\fg{} (XI 2.7) de la
\og cohomologie de Hodge évanescente\fg{} de $V$ :

\SGApage{II}{431}
\begin{equation*}
E_{\mathrm{DR}}^n(V)
=\operatorname{Coker}\!\left(
H_{\mathrm{Hdg}}^n(\mathbb P^{n+1})
\hookrightarrow H_{\mathrm{Hdg}}^n(V)\right),
\end{equation*}


\SGAvar{est $\leq n-2j$, i.e. son co-niveau est $\geq j$. D'autre part, la
conclusion signifie aussi que les valeurs propres du Frobenius
$\mathfrak F_q$ opérant sur la cohomologie évanescente $\ell$-adique
($\ell\ne p$) sont des entiers algébriques (XXI A.2) \underline{divisibles par} $q^j$.
Le résultat de AX 5.0. vient donc à l'appui de la conjecture suivante,
suggérée par le yoga du niveau des motifs, et qui met en relations les
classiques conjectures de HODGE [5] et de TATE [5] :

}{est $\leq n-2j$, i.e. son co-niveau est $\geq j$. D'autre part, la
conclusion signifie aussi que les valeurs propres du Frobenius
$\mathfrak F_q$ opérant sur la cohomologie évanescente $\ell$-adique
($\ell\ne p$) sont des entiers algébriques (XXI A.2) \underline{divisibles par} $q^j$.
Le résultat de AX 5.0 vient donc à l'appui de la conjecture suivante,
suggérée par le yoga du niveau des motifs, et qui met en relation les
classiques conjectures de HODGE [5] et de TATE [5] :

}{F7II-XXII-061}\underline{Conjecture} 5.2.
\underline{Soient} $f:X\to S$ \underline{un morphisme propre et lisse}, \underline{avec} $S$ \underline{intègre}, \underline{de type
fini sur} $\operatorname{Spec}(\mathbb Z)$ \underline{et dominant}
$\operatorname{Spec}(\mathbb Z)$, $i,j$ \underline{des entiers}. \underline{Supposons que le
co-niveau de Hodge du} $H^i$ \underline{de la fibre générique de} $f$ \underline{soit} $\geq j$.
\underline{Alors pour tout point} $s\in S$, \underline{le coniveau} (XX 2.0.6) \underline{de}
$H^i(X_s,\mathbb Q_\ell)$ \underline{est} $\geq j$, \underline{donc} (XX 2.2) \underline{si} $s$ \underline{est un point
fermé de} $S$, \underline{les valeurs propres de Frobenius opérant sur}
$H^i(X_s,\mathbb Q_\ell)$ \underline{sont des entiers algébriques divisibles par}
$q^j$, \underline{où} $q=\operatorname{card}(k(s))$.

\SGAsectiontoc{6. Fin de la démonstration de la formule de congruence.}{6}
\section*{6. Fin de la démonstration de la formule de congruence.}

\SGAvar{Compte tenu de (4.6.6) et (4.6.3)--(4.6.4), (4.5.24) donne une formule de
congruence

}{Compte tenu de (4.6.6) et (4.6.3)--(4.6.4), (4.5.33) donne une formule de
congruence

}{F7II-XXII-062}\begin{equation}\tag{6.0.1}
Z(t,V/\mathbb F_q)(1-t)
\equiv\Delta_S(t)^{(-1)^{n+1}}
\mod{1+qt\mathcal O[[t]]}.
\end{equation}

\SGAvar{Passons au calcul de $\Delta_S(t)$ modulo $\pi\mathcal O[[t]]$. Définissons
(cf. (4.6.11) :

\begin{equation}\tag{6.0.2}
W_S=R_B\otimes_{\mathcal O}\mathbb F_q,
\end{equation}

}{Passons au calcul de $\Delta_S(t)$ modulo $\pi\mathcal O[[t]]$. Définissons
(cf. (4.6.11)) :

\begin{equation}\tag{6.0.2}
W_S=R_S\otimes_{\mathcal O}\mathbb F_q,
\end{equation}

}{F7II-XXII-063}\SGApage{II}{432}

\begin{equation}\tag{6.0.3}
\bar\alpha_S:W_S\longrightarrow W_S
\end{equation}

\noindent
l'opération $\tau^{-1}$-linéaire induite par $\alpha_S$,

\SGAvar{\begin{equation}\tag{6.0.4}
R_S(t)\text{ le sous-module de $R_S$ forme des éléments de degré $i$ en $X_0$},
\end{equation}

}{\begin{equation}\tag{6.0.4}
R_S(i)\text{ le sous-module de $R_S$ formé des éléments de degré $i$ en $X_0$},
\end{equation}

}{F7II-XXII-064}\begin{equation}\tag{6.0.5}
W_S(i)=R_S(i)\otimes_{\mathcal O}\mathbb F_q.
\end{equation}

\underline{Proposition} 6.0.6.
\underline{On a l'inclusion}

\[
\alpha_S(R_S)\subset R_S(1)+pR_S.
\]

\underline{Démonstration.}
On a la factorisation (4.6.14)

\begin{equation}\tag{6.0.7}
\alpha_S=\tau^{-1}\circ\frac1p\,\psi_p\circ H(X).
\end{equation}

On sait (cf. (4.6.15))

\begin{equation}\tag{6.0.8}
H(X)\cdot R_S\subset R_S,
\end{equation}

et évidemment

\begin{equation}\tag{6.0.9}
\tau^{-1}\bigl(R_S(1)+pR_S\bigr)=R_S(1)+pR_S,
\end{equation}

de sorte qu'il suffit de vérifier

\begin{equation}\tag{6.0.10}
\frac1p\,\psi_pR_S\subset R_S(1)+pR_S,
\end{equation}

ce qui est clair par (4.6.19).

\underline{Corollaire} 6.0.11.
\underline{On a l'inclusion}

\[
\bar\alpha_S W_S\subset W_S(1).
\]

\SGApage{II}{433}

\underline{Corollaire} 6.0.12.
\underline{On a la congruence}

\[
\Delta_S(t)\equiv
\det\!\left(
1-t(\bar\alpha_S)^{\operatorname{ord}(q)}
\mathrel{\Big|}W_S(1)\right)
\pmod{\pi\mathcal O[[t]]}.
\]

\SGAvar{\underline{Proposition} 6.1.
\underline{L'opération} $\tau^{1}$-\underline{linéaire}
$\bar\alpha_S:W_S(1)\to W_S(1)$ \underline{est égale à l'opération}

}{\underline{Proposition} 6.1.
\underline{L'opération} $\tau^{-1}$-\underline{linéaire}
$\bar\alpha_S:W_S(1)\to W_S(1)$ \underline{est égale à l'opération}

}{F7II-XXII-065}\begin{equation}\tag{6.1.1}
\tau^{-1}\psi_p\circ(X_0f(X))^{p-1}.
\end{equation}

\underline{Démonstration.}
Écrivons

\begin{equation}\tag{6.1.2}
H(X)=1+\sum_{i\geq1}H_i(X)
\end{equation}

où

\begin{equation}\tag{6.1.3}
H_i(X)\text{ est de degré $i$ en $X_0$}.
\end{equation}

D'après (4.6.19) encore, on a

\begin{equation}\tag{6.1.4}
\frac1p\,\psi_pR_S(i)\subset pR_S
\qquad\text{si }i\ne p.
\end{equation}

Compte tenu de ce que

\begin{equation}\tag{6.1.5}
H_i(X)\cdot R_S(j)\subset R_S(i+j),
\end{equation}

on a

\begin{equation}\tag{6.1.6}
\left(
\alpha_S-\frac{\tau^{-1}}p\,\psi_pH_{p-1}(X)
\right)(R_S(1))
\subset pR_S.
\end{equation}

On a donc

\begin{equation}\tag{6.1.7}
\left[
(\alpha_S)^{\operatorname{ord}(q)}
-
\left(
\frac{\tau^{-1}}p\,\psi_p\cdot H_{p-1}(X)
\right)^{\operatorname{ord}(q)}
\right]R_S(1)
\subset pR_S.
\end{equation}

Calculons maintenant $H_{p-1}(X)$. On a (cf. (4.4.4))

\SGApage{II}{434}

\SGAvar{\begin{equation}\tag{6.1.8}
\begin{aligned}
H(X)
&=\prod_W\Theta(A_WX_0X^W)\\
&=\prod_W\exp\!\left(
\sum_{n\geq0}\frac{\pi A_WX_0X^W)^{p^n}}{p^n}
\right)\\
&=\prod_W\prod_{n\geq0}
\exp\!\left(\frac{(\pi A_WX_0X^W)^{p^n}}{p^n}\right)\\
&\equiv\prod_W\exp(\pi A_WX_0X^W)
\pmod{\text{degré }p}\\
&\equiv\exp\!\left(\pi\sum A_WX_0X^W\right)
\pmod{\text{degré }p}\\
&\equiv\exp(\pi X_0F(X))
\pmod{\text{degré }p},
\end{aligned}
\end{equation}

}{\begin{equation}\tag{6.1.8}
\begin{aligned}
H(X)
&=\prod_W\Theta(A_WX_0X^W)\\
&=\prod_W\exp\!\left(
\sum_{n\geq0}\frac{(\pi A_WX_0X^W)^{p^n}}{p^n}
\right)\\
&=\prod_W\prod_{n\geq0}
\exp\!\left(\frac{(\pi A_WX_0X^W)^{p^n}}{p^n}\right)\\
&\equiv\prod_W\exp(\pi A_WX_0X^W)
\pmod{\text{degré }p}\\
&\equiv\exp\!\left(\pi\sum A_WX_0X^W\right)
\pmod{\text{degré }p}\\
&\equiv\exp(\pi X_0F(X))
\pmod{\text{degré }p},
\end{aligned}
\end{equation}

}{F7II-XXII-066}d'où

\begin{equation}\tag{6.1.9}
H_{p-1}(X)
\equiv\frac{\pi^{p-1}X_0^{p-1}}{(p-1)!}\,F(X)^{p-1}.
\end{equation}

On constate facilement que

\begin{equation}\tag{6.1.10}
\frac{\pi^{p-1}}p\equiv-1\pmod p
\end{equation}

donc

\SGAvar{\begin{equation}\tag{6.1.11}
\frac1p\,H_{p-1}(X)
\equiv(X_0f(X)^{p-1}\pmod p,
\end{equation}

}{\begin{equation}\tag{6.1.11}
\frac1p\,H_{p-1}(X)
\equiv(X_0f(X))^{p-1}\pmod p,
\end{equation}

}{F7II-XXII-067}ce qui achève la démonstration de 6.1.

\underline{Corollaire} 6.1.12.
\underline{L'opérateur linéaire}

\[
(\bar\alpha_S)^{\operatorname{ord}(q)}:W_S(1)\longrightarrow W_S(1)
\]

\noindent
\underline{n'est autre que} $\psi_q\circ(X_0f(X))^{q-1}$.

\SGApage{II}{435}

\underline{Démonstration.}
D'après (6.1.1)

\SGAvar{\begin{equation*}
\begin{aligned}
(\bar\alpha_S)^{\operatorname{ord}(q)}
&=
\left(
\tau^{-1}\psi_p\circ(X_0f(X))^{p-1}
\right)^{\operatorname{ord}(q)}\\
&=
\psi_q
\prod_{j=0}^{\operatorname{ord}(a)-1}
\left(X_0^{p^j}f^{\tau^j}(X^{p^j})\right)^{p-1}\\
&=
\psi_q
\prod_{j=0}^{\operatorname{ord}(q)-1}
(X_0f(X))^{p^j(p-1)}\\
&=\psi_q\circ(X_0f(X))^{q-1}.
\end{aligned}
\end{equation*}

Mettant ensemble (6.1.12) (6.0.12), et (6.0.1), on trouve,

}{\begin{equation*}
\begin{aligned}
(\bar\alpha_S)^{\operatorname{ord}(q)}
&=
\left(
\tau^{-1}\psi_p\circ(X_0f(X))^{p-1}
\right)^{\operatorname{ord}(q)}\\
&=
\psi_q\circ
\prod_{j=0}^{\operatorname{ord}(q)-1}
\left(X_0^{p^j}f^{\tau^j}(X^{p^j})\right)^{p-1}\\
&=
\psi_q\circ
\prod_{j=0}^{\operatorname{ord}(q)-1}
(X_0f(X))^{p^j(p-1)}\\
&=\psi_q\circ(X_0f(X))^{q-1}.
\end{aligned}
\end{equation*}

Mettant ensemble (6.1.12), (6.0.12), et (6.0.1), on trouve

}{F7II-XXII-068}\underline{Corollaire} 6.1.13.
\underline{On a la congruence mod} $p$ :

\begin{equation*}
Z(t,V/\mathbb F_q)
\equiv
\frac{
\det\!\left(
1-t\,\psi_q(X_0f(X))^{q-1}
\mathrel{\Big|}W_S(1)
\right)^{(-1)^{n+1}}
}{1-t}.
\end{equation*}

\subsection*{6.2.}
Il reste à calculer les groupes de cohomologie de $\mathcal O_V$.

Posons

\begin{equation}\tag{6.2.1}
\mathfrak F_q
=\text{l'opération \og élévation à la puissance $q$'ième\fg{} dans $\mathcal O_V$}.
\end{equation}

Nous avons une suite exacte de faisceaux sur
$\mathbb P^{n+1}=\mathbb P$ :

\SGAvar{\begin{equation}\tag{6.2.2}
0\longrightarrow\mathcal O_{\mathbb P}(-d)
\xrightarrow{\,f\,}\mathcal O_{\mathbb P}
\longrightarrow\mathcal O_V\longrightarrow\mathcal O,
\end{equation}

}{\begin{equation}\tag{6.2.2}
0\longrightarrow\mathcal O_{\mathbb P}(-d)
\xrightarrow{\,f\,}\mathcal O_{\mathbb P}
\longrightarrow\mathcal O_V\longrightarrow0,
\end{equation}

}{F7II-XXII-069}et le diagramme suivant est commutatif

\SGAvar{\begin{equation}\tag{6.2.3}
\begin{tikzcd}[column sep=large,row sep=large]
0 \arrow[r] \&
\mathcal O_{\mathbb P}(-d) \arrow[r,"f"] \arrow[d,"f^{q-1}\mathfrak F_q"'] \&
\mathcal O_{\mathbb P} \arrow[r] \arrow[d,"\mathfrak F_q"'] \&
\mathcal O_V \arrow[r] \arrow[d,"\mathfrak F_q"'] \&
\mathcal O\\
0 \arrow[r] \&
\mathcal O_{\mathbb P}(-d) \arrow[r,"f"] \&
\mathcal O_{\mathbb P} \arrow[r] \&
\mathcal O_V \arrow[r] \&
\mathcal O .
\end{tikzcd}
\end{equation}

}{\begin{equation}\tag{6.2.3}
\begin{tikzcd}[column sep=large,row sep=large]
0 \arrow[r] \&
\mathcal O_{\mathbb P}(-d) \arrow[r,"f"] \arrow[d,"f^{q-1}\mathfrak F_q"'] \&
\mathcal O_{\mathbb P} \arrow[r] \arrow[d,"\mathfrak F_q"'] \&
\mathcal O_V \arrow[r] \arrow[d,"\mathfrak F_q"'] \&
0\\
0 \arrow[r] \&
\mathcal O_{\mathbb P}(-d) \arrow[r,"f"] \&
\mathcal O_{\mathbb P} \arrow[r] \&
\mathcal O_V \arrow[r] \&
0 .
\end{tikzcd}
\end{equation}

}{F7II-XXII-070}\SGApage{II}{436}

Compte tenu de ce que

\begin{equation}\tag{6.2.4}
\begin{cases}
H^0(\mathbb P^{n+1},\mathcal O_{\mathbb P})\simeq\mathbb F_q\\
H^i(\mathbb P^{n+1},\mathcal O_{\mathbb P})=0&i\ne0\\
H^i(\mathbb P^{n+1},\mathcal O_{\mathbb P}(-d))=0&i\ne0,n+1,
\end{cases}
\end{equation}

\SGAvar{la suite exacte longue de cohomologie associé à (6.2.2), donne

}{la suite exacte longue de cohomologie associée à (6.2.2), donne

}{F7II-XXII-071}\begin{equation}\tag{6.2.5}
\begin{cases}
H^0(V,\mathcal O_V)\simeq\mathbb F_q\\
H^i(V,\mathcal O_V)=0&\text{si }i\ne0,n\\
H^n(V,\mathcal O_V)
\simeq H^{n+1}(\mathbb P^{n+1},\mathcal O_{\mathbb P}(-d)).
\end{cases}
\end{equation}

Le diagramme commutatif (6.2.3) donne alors

\begin{equation}\tag{6.2.6}
\det\!\left(1-t\mathfrak F_q
\mathrel{\Big|}H^0(V,\mathcal O_V)\right)=1-t
\end{equation}

et

\SGAvar{\begin{equation}\tag{6.2.7}
\det\!\left(1-t\mathfrak F_q
\mathrel{\Big|}H^n(V,\mathcal O_V)\right)
=
\det\!\left(
1-t(f^{q-1}\mathfrak F_q)
\mathrel{\Big|}
H^{n+1}(\mathbb P^{n+1},\mathcal O_{\mathbb P}(-d))
\right..
\end{equation}

}{\begin{equation}\tag{6.2.7}
\det\!\left(1-t\mathfrak F_q
\mathrel{\Big|}H^n(V,\mathcal O_V)\right)
=
\det\!\left(
1-t(f^{q-1}\mathfrak F_q)
\mathrel{\Big|}
H^{n+1}(\mathbb P^{n+1},\mathcal O_{\mathbb P}(-d))
\right).
\end{equation}

}{F7II-XXII-072}Il ne reste qu'à établir

\SGAvar{\begin{equation}\tag{6.2.8}
\begin{aligned}
&\det\!\left(
1-t(f^{q-1}\mathfrak F_q)
\mathrel{\Big|}
H^{n+1}(\mathbb P^{n+1},\mathcal O_{\mathbb P}(-d))
\right.\\
&\qquad=
\det\!\left(
1-t\,(\psi_q\cdot(X_0f(X))^{q-1}
\mathrel{\Big|}W_S(1)
\right)
\end{aligned}
\end{equation}

}{\begin{equation}\tag{6.2.8}
\begin{aligned}
&\det\!\left(
1-t(f^{q-1}\mathfrak F_q)
\mathrel{\Big|}
H^{n+1}(\mathbb P^{n+1},\mathcal O_{\mathbb P}(-d))
\right)\\
&\qquad=
\det\!\left(
1-t\,\psi_q\cdot(X_0f(X))^{q-1}
\mathrel{\Big|}W_S(1)
\right)
\end{aligned}
\end{equation}

}{F7II-XXII-073}ce qui est une conséquence de la

\underline{Proposition} 6.3.
\underline{Il y a un isomorphisme}

\SGAvar{\begin{equation}\tag{6.3.1}
W_S(1)\simeq
\operatorname{Hom}_{\mathbb F_q}\!\left(
H^{n+1}(\mathbb P^{n+1}.\mathcal O_{\mathbb P}(-d),
\mathbb F_q\right)
\end{equation}

}{\begin{equation}\tag{6.3.1}
W_S(1)\simeq
\operatorname{Hom}_{\mathbb F_q}\!\left(
H^{n+1}(\mathbb P^{n+1},\mathcal O_{\mathbb P}(-d)),
\mathbb F_q\right)
\end{equation}

}{F7II-XXII-074}\underline{par rapport auquel l'opérateur}

\begin{equation}\tag{6.3.2}
\psi_q\circ(X_0f(X))^{q-1}:W_S(1)\longrightarrow W_S(1)
\end{equation}

\SGApage{II}{437}

\underline{est le transposé de}

\begin{equation}\tag{6.3.3}
f(X)^{q-1}\mathfrak F_q:
H^{n+1}(\mathbb P^{n+1},\mathcal O_{\mathbb P}(-d))
\longrightarrow
H^{n+1}(\mathbb P^{n+1},\mathcal O_{\mathbb P}(-d)).
\end{equation}

\underline{Démonstration.}
En calculant le groupe
$H^{n+1}(\mathbb P^{n+1},\mathcal O_{\mathbb P}(-d))$ à l'aide du
recouvrement standard de $\mathbb P^{n+1}$, on trouve

\SGAvar{\begin{equation}\tag{6.3.4}
\begin{aligned}
H^{i+1}(\mathbb P^{n+1},\mathcal O_{\mathbb P}(-d))
&=\text{le $\mathbb F_q$-vectoriel de base des monômes}\\
&\quad X_1^{W_1}\cdots X_{n+2}^{W_{n+2}}
\qquad\text{les }W_i\in\mathbb Z\text{ vérifiant }\sum W_i=-d\\
&=\text{le $\mathbb F_q$-sous-vectoriel de base des monômes}\\
&\quad X_1^{W_1}\cdots X_{n+2}^{W_{n+2}}
\text{ qui vérifient en plus : }\exists i\text{ tel que }W_i\geq0.
\end{aligned}
\end{equation}

}{\begin{equation}\tag{6.3.4}
\begin{aligned}
H^{n+1}(\mathbb P^{n+1},\mathcal O_{\mathbb P}(-d))
&=\text{le $\mathbb F_q$-vectoriel de base des monômes}\\
&\quad X_1^{W_1}\cdots X_{n+2}^{W_{n+2}}
\qquad\text{les }W_i\in\mathbb Z\text{ vérifiant }\sum W_i=-d\\
&\mathrel{\big/}\text{le $\mathbb F_q$-sous-vectoriel de base des monômes}\\
&\quad X_1^{W_1}\cdots X_{n+2}^{W_{n+2}}
\text{ qui vérifient en plus : }\exists i\text{ tel que }W_i\geq0.
\end{aligned}
\end{equation}

}{F7II-XXII-075}D'autre part, $W_S(1)=R_S(1)\otimes_{\mathcal O}\mathbb F_q$, et on a

\begin{equation}\tag{6.3.5}
\begin{aligned}
R_S(1)
&=\text{le $\mathcal O$-module libre de base les monômes}\\
&\quad
(\pi X_0)X_1^{U_1}\cdots X_{n+2}^{U_{n+2}}
\text{ vérifiant }
\sum U_i=d,\quad U_i>0\ \forall_i.
\end{aligned}
\end{equation}

L'accouplement

\begin{equation}\tag{6.3.6}
(\ ,\ ):
W_S(1)\times
H^{n+1}(\mathbb P^{n+1},\mathcal O_{\mathbb P}(-d))
\longrightarrow\mathbb F_q
\end{equation}

défini par

\SGAvar{\begin{equation}\tag{6.3.7}
\left(
\text{classe de }\pi X_0X_1^{U_1}\cdots X_{n+2}^{U_{n+2}},
X_1^{W_1}\cdots X_{n+2}^{W_{n+2}}
\right)
=
\begin{cases}
1&\text{si }U_1=-W_i\text{ pour tout }i\\
0&\text{sinon}
\end{cases}
\end{equation}

est évidemment une dualité parfaite de $\mathbb F_q$-vectoriels, et on
constate facilement que $\psi_q\cdot(X_0f(X))^{q-1}$ est la transposée de
$f(X)^{q-1}\cdot f_q$, ce qui prouve 6.3.

}{\begin{equation}\tag{6.3.7}
\left(
\text{classe de }\pi X_0X_1^{U_1}\cdots X_{n+2}^{U_{n+2}},
X_1^{W_1}\cdots X_{n+2}^{W_{n+2}}
\right)
=
\begin{cases}
1&\text{si }U_i=-W_i\text{ pour tout }i\\
0&\text{sinon}
\end{cases}
\end{equation}

est évidemment une dualité parfaite de $\mathbb F_q$-vectoriels, et on
constate facilement que $\psi_q\cdot(X_0f(X))^{q-1}$ est la transposée de
$f(X)^{q-1}\cdot \mathfrak F_q$, ce qui prouve 6.3.

}{F7II-XXII-076}\SGApage{II}{438}

\section*{Bibliographie}

\SGAvar{\begin{enumerate}[label={[\arabic*]},leftmargin=2.4em]
\item AX, J. Zeroes of Polynomials Over Finite Fields,
  Amer. J. Math. 86, 1964, 255--261.
\item DWORK, B. On the Rationality of the Zéta Function of an
  Algebraic Variety, Amer. J. Math. 82, 1960, 631--648.
\item DWORK, B. An Algebraic Variety, in Proc. of the Woods-Hole
  Summer Institute 1964 (miméo).
\item DWORK, B. On the Zeta Function of Hypersurface II,
  Ann. Math. (2), 80, 1964, 227--299.
\item GROTHENDIECK, A., Le Groupe de Brauer III, dans
  Dix Exposés sur la Cohomologie des Schémas,
  North Holland Pub. Co, 1968.
\item SERRE, J.P. Groupes Algébriques et Corps de Classes,
  Hermann, Paris, 1959.
\item SERRE, J.P. Endomorphismes Complétement, Continus des Espaces
  de Banach $p$-Adiques, Pub. Math. I.H.E.S. n\textsuperscript{o} 12, 1962.
\item CHEVALLEY, C. Démonstration d'une hypothèse de M. Artin,
  Abhandlungen Math. Sem. Hamburg, 11, 1936, 73--75.
\end{enumerate}
}{\begin{enumerate}[label={[\arabic*]},leftmargin=2.4em]
\item AX, J. Zeroes of Polynomials Over Finite Fields,
  Amer. J. Math. 86, 1964, 255--261.
\item DWORK, B. On the Rationality of the Zeta Function of an
  Algebraic Variety, Amer. J. Math. 82, 1960, 631--648.
\item DWORK, B. An Algebraic Variety, in Proc. of the Woods-Hole
  Summer Institute 1964 (miméo).
\item DWORK, B. On the Zeta Function of Hypersurface II,
  Ann. Math. (2), 80, 1964, 227--299.
\item GROTHENDIECK, A., Le Groupe de Brauer III, dans
  Dix Exposés sur la Cohomologie des Schémas,
  North Holland Pub. Co, 1968.
\item SERRE, J.P. Groupes Algébriques et Corps de Classes,
  Hermann, Paris, 1959.
\item SERRE, J.P. Endomorphismes Complètement Continus des Espaces
  de Banach $p$-Adiques, Pub. Math. I.H.E.S. n\textsuperscript{o} 12, 1962.
\item CHEVALLEY, C. Démonstration d'une hypothèse de M. Artin,
  Abhandlungen Math. Sem. Hamburg, 11, 1936, 73--75.
\end{enumerate}
}{F7II-XXII-077}\selectlanguage{french}

\ifnum\SGAform=1
\clearpage\section*{Apparat critique}\small
Les lectures ci-dessous sont données dans le TeX exact des témoins, afin de préserver les distinctions de notation. Les indices décimaux renvoient aux différences individuelles du passage marqué.

\par\medskip\noindent\hypertarget{app-F7II-XXII-001.1}{}\textbf{F7II-XXII-001.1}\quad p. 401; printed misprint.\par
\noindent Imprimé : {\ttfamily\footnotesize -\allowbreak{}}\par
\noindent Corrigé : \emph{(1 espace(s), 0 retour(s) de ligne dans le témoin TeX)}\par
\noindent Source : 7II:693; 7II:832. \nolinkurl{fixes/MISPRINTS_7II.json#693}; \nolinkurl{fixes/MISPRINTS_7II.json#832}\par

\par\medskip\noindent\hypertarget{app-F7II-XXII-002.1}{}\textbf{F7II-XXII-002.1}\quad p. 402; designated-witness attribution.\par
\noindent Imprimé : \emph{(vide)}\par
\noindent Corrigé : \emph{(1 espace(s), 0 retour(s) de ligne dans le témoin TeX)}\par
\noindent Source : WITNESS-F7II-XXII-002.1. \nolinkurl{fr/sga7ii/diplomatic/expose_XXII_body.tex}; \nolinkurl{fr/sga7ii/current/expose_XXII_body.tex}\par

\par\medskip\noindent\hypertarget{app-F7II-XXII-002.2}{}\textbf{F7II-XXII-002.2}\quad p. 402; printed misprint.\par
\noindent Imprimé : {\ttfamily\footnotesize coefficient\allowbreak{}}\par
\noindent Corrigé : {\ttfamily\footnotesize coefficients\allowbreak{}}\par
\noindent Source : 7II:694; 7II:833. \nolinkurl{fixes/MISPRINTS_7II.json#694}; \nolinkurl{fixes/MISPRINTS_7II.json#833}\par

\par\medskip\noindent\hypertarget{app-F7II-XXII-002.3}{}\textbf{F7II-XXII-002.3}\quad p. 402; inherited corrected-layer difference.\par
\noindent Imprimé : \emph{(vide)}\par
\noindent Corrigé : \emph{(1 espace(s), 0 retour(s) de ligne dans le témoin TeX)}\par
\noindent Source : ADOPT-F7II-XXII-002.3. \nolinkurl{fr/sga7ii/diplomatic/expose_XXII_body.tex}; \nolinkurl{fr/sga7ii/current/expose_XXII_body.tex}\par

\par\medskip\noindent\hypertarget{app-F7II-XXII-003.1}{}\textbf{F7II-XXII-003.1}\quad p. 402; printed misprint.\par
\noindent Imprimé : {\ttfamily\footnotesize Una\allowbreak{}}\par
\noindent Corrigé : {\ttfamily\footnotesize Une\allowbreak{}}\par
\noindent Source : 7II:695; 7II:834. \nolinkurl{fixes/MISPRINTS_7II.json#695}; \nolinkurl{fixes/MISPRINTS_7II.json#834}\par

\par\medskip\noindent\hypertarget{app-F7II-XXII-004.1}{}\textbf{F7II-XXII-004.1}\quad p. 402; printed misprint.\par
\noindent Imprimé : \emph{(vide)}\par
\noindent Corrigé : {\ttfamily\footnotesize est\allowbreak{} }\par
\noindent Source : 7II:835. \nolinkurl{fixes/MISPRINTS_7II.json#835}\par

\par\medskip\noindent\hypertarget{app-F7II-XXII-005.1}{}\textbf{F7II-XXII-005.1}\quad p. 403; printed misprint.\par
\noindent Imprimé : {\ttfamily\footnotesize espace\allowbreak{}}\par
\noindent Corrigé : {\ttfamily\footnotesize espaces\allowbreak{}}\par
\noindent Source : 7II:697; 7II:836. \nolinkurl{fixes/MISPRINTS_7II.json#697}; \nolinkurl{fixes/MISPRINTS_7II.json#836}\par

\par\medskip\noindent\hypertarget{app-F7II-XXII-006.1}{}\textbf{F7II-XXII-006.1}\quad p. 403; printed misprint.\par
\noindent Imprimé : {\ttfamily\footnotesize les\allowbreak{}}\par
\noindent Corrigé : {\ttfamily\footnotesize le\allowbreak{}}\par
\noindent Source : 7II:837. \nolinkurl{fixes/MISPRINTS_7II.json#837}\par

\par\medskip\noindent\hypertarget{app-F7II-XXII-007.1}{}\textbf{F7II-XXII-007.1}\quad p. 404; designated-witness attribution.\par
\noindent Imprimé : \emph{(vide)}\par
\noindent Corrigé : \emph{(1 espace(s), 0 retour(s) de ligne dans le témoin TeX)}\par
\noindent Source : WITNESS-F7II-XXII-007.1. \nolinkurl{fr/sga7ii/diplomatic/expose_XXII_body.tex}; \nolinkurl{fr/sga7ii/current/expose_XXII_body.tex}\par

\par\medskip\noindent\hypertarget{app-F7II-XXII-007.2}{}\textbf{F7II-XXII-007.2}\quad p. 404; printed misprint.\par
\noindent Imprimé : {\ttfamily\footnotesize Gl\allowbreak{}}\par
\noindent Corrigé : {\ttfamily\footnotesize GL\allowbreak{}}\par
\noindent Source : 7II:838. \nolinkurl{fixes/MISPRINTS_7II.json#838}\par

\par\medskip\noindent\hypertarget{app-F7II-XXII-007.3}{}\textbf{F7II-XXII-007.3}\quad p. 404; printed misprint.\par
\noindent Imprimé : {\ttfamily\footnotesize Gl\allowbreak{}}\par
\noindent Corrigé : {\ttfamily\footnotesize GL\allowbreak{}}\par
\noindent Source : 7II:838. \nolinkurl{fixes/MISPRINTS_7II.json#838}\par

\par\medskip\noindent\hypertarget{app-F7II-XXII-008.1}{}\textbf{F7II-XXII-008.1}\quad p. 405; inherited corrected-layer difference.\par
\noindent Imprimé : \emph{(vide)}\par
\noindent Corrigé : {\ttfamily\footnotesize ,\allowbreak{}}\par
\noindent Source : ADOPT-F7II-XXII-008.1. \nolinkurl{fr/sga7ii/diplomatic/expose_XXII_body.tex}; \nolinkurl{fr/sga7ii/current/expose_XXII_body.tex}\par

\par\medskip\noindent\hypertarget{app-F7II-XXII-009.1}{}\textbf{F7II-XXII-009.1}\quad p. 405; printed misprint.\par
\noindent Imprimé : \emph{(vide)}\par
\noindent Corrigé : {\ttfamily\footnotesize \$V\$\allowbreak{} }\par
\noindent Source : 7II:839. \nolinkurl{fixes/MISPRINTS_7II.json#839}\par

\par\medskip\noindent\hypertarget{app-F7II-XXII-010.1}{}\textbf{F7II-XXII-010.1}\quad p. 406; misprint\_in\_print.\par
\noindent Imprimé : \emph{(vide)}\par
\noindent Corrigé : {\ttfamily\footnotesize ,\allowbreak{}}\par
\noindent Source : FIX-7II-D-0358. \nolinkurl{fixes/FIXLOG_7II.jsonl#358}\par

\par\medskip\noindent\hypertarget{app-F7II-XXII-011.1}{}\textbf{F7II-XXII-011.1}\quad p. 407; inherited corrected-layer difference.\par
\noindent Imprimé : \emph{(1 espace(s), 0 retour(s) de ligne dans le témoin TeX)}\par
\noindent Corrigé : {\ttfamily\footnotesize -\allowbreak{}}\par
\noindent Source : ADOPT-F7II-XXII-011.1. \nolinkurl{fr/sga7ii/diplomatic/expose_XXII_body.tex}; \nolinkurl{fr/sga7ii/current/expose_XXII_body.tex}\par

\par\medskip\noindent\hypertarget{app-F7II-XXII-012.1}{}\textbf{F7II-XXII-012.1}\quad p. 408; printed misprint.\par
\noindent Imprimé : {\ttfamily\footnotesize q\allowbreak{}}\par
\noindent Corrigé : {\ttfamily\footnotesize p\allowbreak{}}\par
\noindent Source : 7II:841. \nolinkurl{fixes/MISPRINTS_7II.json#841}\par

\par\medskip\noindent\hypertarget{app-F7II-XXII-014.1}{}\textbf{F7II-XXII-014.1}\quad p. 408; inherited corrected-layer difference.\par
\noindent Imprimé : \emph{(vide)}\par
\noindent Corrigé : {\ttfamily\footnotesize ,\allowbreak{}}\par
\noindent Source : ADOPT-F7II-XXII-014.1. \nolinkurl{fr/sga7ii/diplomatic/expose_XXII_body.tex}; \nolinkurl{fr/sga7ii/current/expose_XXII_body.tex}\par

\par\medskip\noindent\hypertarget{app-F7II-XXII-015.1}{}\textbf{F7II-XXII-015.1}\quad p. 408; printed misprint.\par
\noindent Imprimé : {\ttfamily\footnotesize )\allowbreak{}}\par
\noindent Corrigé : \emph{(vide)}\par
\noindent Source : 7II:842. \nolinkurl{fixes/MISPRINTS_7II.json#842}\par

\par\medskip\noindent\hypertarget{app-F7II-XXII-016.1}{}\textbf{F7II-XXII-016.1}\quad p. 409; printed misprint.\par
\noindent Imprimé : {\ttfamily\footnotesize pointes\allowbreak{}}\par
\noindent Corrigé : {\ttfamily\footnotesize points\allowbreak{}}\par
\noindent Source : 7II:843. \nolinkurl{fixes/MISPRINTS_7II.json#843}\par

\par\medskip\noindent\hypertarget{app-F7II-XXII-016.2}{}\textbf{F7II-XXII-016.2}\quad p. 409; printed misprint.\par
\noindent Imprimé : {\ttfamily\footnotesize k\allowbreak{}}\par
\noindent Corrigé : {\ttfamily\footnotesize q\allowbreak{}}\par
\noindent Source : 7II:844. \nolinkurl{fixes/MISPRINTS_7II.json#844}\par

\par\medskip\noindent\hypertarget{app-F7II-XXII-016.3}{}\textbf{F7II-XXII-016.3}\quad p. 409; printed misprint.\par
\noindent Imprimé : {\ttfamily\footnotesize k\allowbreak{}}\par
\noindent Corrigé : {\ttfamily\footnotesize q\allowbreak{}}\par
\noindent Source : 7II:844. \nolinkurl{fixes/MISPRINTS_7II.json#844}\par

\par\medskip\noindent\hypertarget{app-F7II-XXII-018.1}{}\textbf{F7II-XXII-018.1}\quad p. 409; misprint\_in\_print.\par
\noindent Imprimé : \emph{(vide)}\par
\noindent Corrigé : {\ttfamily\footnotesize \textbackslash{}overline\{\allowbreak{}}\par
\noindent Source : FIX-7II-D-0658. \nolinkurl{fixes/FIXLOG_7II.jsonl#658}\par

\par\medskip\noindent\hypertarget{app-F7II-XXII-018.2}{}\textbf{F7II-XXII-018.2}\quad p. 409; misprint\_in\_print.\par
\noindent Imprimé : \emph{(vide)}\par
\noindent Corrigé : {\ttfamily\footnotesize \}\allowbreak{}}\par
\noindent Source : FIX-7II-D-0658. \nolinkurl{fixes/FIXLOG_7II.jsonl#658}\par

\par\medskip\noindent\hypertarget{app-F7II-XXII-018.3}{}\textbf{F7II-XXII-018.3}\quad p. 409; formula.\par
\noindent Imprimé : \emph{(vide)}\par
\noindent Corrigé : {\ttfamily\footnotesize \textbackslash{}overline\{\allowbreak{}}\par
\noindent Source : FIX-7II-D-0375; FIX-7II-C-0306. \nolinkurl{fixes/FIXLOG_7II.jsonl#375}; \nolinkurl{fixes/FIXLOG_7II_CURRENT.jsonl#306}\par

\par\medskip\noindent\hypertarget{app-F7II-XXII-018.4}{}\textbf{F7II-XXII-018.4}\quad p. 409; formula.\par
\noindent Imprimé : \emph{(vide)}\par
\noindent Corrigé : {\ttfamily\footnotesize \}\allowbreak{}}\par
\noindent Source : FIX-7II-D-0375; FIX-7II-C-0306. \nolinkurl{fixes/FIXLOG_7II.jsonl#375}; \nolinkurl{fixes/FIXLOG_7II_CURRENT.jsonl#306}\par

\par\medskip\noindent\hypertarget{app-F7II-XXII-018.5}{}\textbf{F7II-XXII-018.5}\quad p. 409; misprint\_in\_print.\par
\noindent Imprimé : \emph{(vide)}\par
\noindent Corrigé : {\ttfamily\footnotesize \textbackslash{}overline\{\allowbreak{}}\par
\noindent Source : FIX-7II-D-0660. \nolinkurl{fixes/FIXLOG_7II.jsonl#660}\par

\par\medskip\noindent\hypertarget{app-F7II-XXII-018.6}{}\textbf{F7II-XXII-018.6}\quad p. 409; misprint\_in\_print.\par
\noindent Imprimé : \emph{(vide)}\par
\noindent Corrigé : {\ttfamily\footnotesize \}\allowbreak{}}\par
\noindent Source : FIX-7II-D-0660. \nolinkurl{fixes/FIXLOG_7II.jsonl#660}\par

\par\medskip\noindent\hypertarget{app-F7II-XXII-018.7}{}\textbf{F7II-XXII-018.7}\quad p. 409; inherited corrected-layer difference.\par
\noindent Imprimé : \emph{(vide)}\par
\noindent Corrigé : {\ttfamily\footnotesize \textasciicircum{}\{-1\}\allowbreak{}}\par
\noindent Source : ADOPT-F7II-XXII-018.7. \nolinkurl{fr/sga7ii/diplomatic/expose_XXII_body.tex}; \nolinkurl{fr/sga7ii/current/expose_XXII_body.tex}\par

\par\medskip\noindent\hypertarget{app-F7II-XXII-019.1}{}\textbf{F7II-XXII-019.1}\quad p. 410; designated-witness attribution.\par
\noindent Imprimé : {\ttfamily\footnotesize 6\allowbreak{}}\par
\noindent Corrigé : {\ttfamily\footnotesize 3\allowbreak{}}\par
\noindent Source : WITNESS-F7II-XXII-019.1. \nolinkurl{fr/sga7ii/diplomatic/expose_XXII_body.tex}; \nolinkurl{fr/sga7ii/current/expose_XXII_body.tex}\par

\par\medskip\noindent\hypertarget{app-F7II-XXII-019.2}{}\textbf{F7II-XXII-019.2}\quad p. 410; designated-witness attribution.\par
\noindent Imprimé : {\ttfamily\footnotesize 1\allowbreak{}}\par
\noindent Corrigé : {\ttfamily\footnotesize 2\allowbreak{}}\par
\noindent Source : WITNESS-F7II-XXII-019.2. \nolinkurl{fr/sga7ii/diplomatic/expose_XXII_body.tex}; \nolinkurl{fr/sga7ii/current/expose_XXII_body.tex}\par

\par\medskip\noindent\hypertarget{app-F7II-XXII-019.3}{}\textbf{F7II-XXII-019.3}\quad p. 410; designated-witness attribution.\par
\noindent Imprimé : \emph{(vide)}\par
\noindent Corrigé : {\ttfamily\footnotesize )--(3.2\allowbreak{}}\par
\noindent Source : WITNESS-F7II-XXII-019.3. \nolinkurl{fr/sga7ii/diplomatic/expose_XXII_body.tex}; \nolinkurl{fr/sga7ii/current/expose_XXII_body.tex}\par

\par\medskip\noindent\hypertarget{app-F7II-XXII-019.4}{}\textbf{F7II-XXII-019.4}\quad p. 410; designated-witness attribution.\par
\noindent Imprimé : \emph{(vide)}\par
\noindent Corrigé : {\ttfamily\footnotesize 3.\allowbreak{}}\par
\noindent Source : WITNESS-F7II-XXII-019.4. \nolinkurl{fr/sga7ii/diplomatic/expose_XXII_body.tex}; \nolinkurl{fr/sga7ii/current/expose_XXII_body.tex}\par

\par\medskip\noindent\hypertarget{app-F7II-XXII-019.5}{}\textbf{F7II-XXII-019.5}\quad p. 410; designated-witness attribution.\par
\noindent Imprimé : {\ttfamily\footnotesize .3\allowbreak{}}\par
\noindent Corrigé : \emph{(vide)}\par
\noindent Source : WITNESS-F7II-XXII-019.5. \nolinkurl{fr/sga7ii/diplomatic/expose_XXII_body.tex}; \nolinkurl{fr/sga7ii/current/expose_XXII_body.tex}\par

\par\medskip\noindent\hypertarget{app-F7II-XXII-019.6}{}\textbf{F7II-XXII-019.6}\quad p. 410; printed misprint.\par
\noindent Imprimé : {\ttfamily\footnotesize ,\allowbreak{}}\par
\noindent Corrigé : {\ttfamily\footnotesize .\allowbreak{}}\par
\noindent Source : 7II:848. \nolinkurl{fixes/MISPRINTS_7II.json#848}\par

\par\medskip\noindent\hypertarget{app-F7II-XXII-020.1}{}\textbf{F7II-XXII-020.1}\quad p. 410; printed misprint.\par
\noindent Imprimé : {\ttfamily\footnotesize .\allowbreak{}}\par
\noindent Corrigé : {\ttfamily\footnotesize ,\allowbreak{}}\par
\noindent Source : 7II:849. \nolinkurl{fixes/MISPRINTS_7II.json#849}\par

\par\medskip\noindent\hypertarget{app-F7II-XXII-021.1}{}\textbf{F7II-XXII-021.1}\quad p. 410; printed misprint.\par
\noindent Imprimé : {\ttfamily\footnotesize .\allowbreak{}}\par
\noindent Corrigé : {\ttfamily\footnotesize )\allowbreak{}}\par
\noindent Source : 7II:850. \nolinkurl{fixes/MISPRINTS_7II.json#850}\par

\par\medskip\noindent\hypertarget{app-F7II-XXII-022.1}{}\textbf{F7II-XXII-022.1}\quad p. 411; misprint\_pass.\par
\noindent Imprimé : \emph{(vide)}\par
\noindent Corrigé : {\ttfamily\footnotesize  ;\allowbreak{}}\par
\noindent Source : 7II:887. \nolinkurl{fixes/FIXLOG_7II_CURRENT.jsonl#899}\par

\par\medskip\noindent\hypertarget{app-F7II-XXII-023.1}{}\textbf{F7II-XXII-023.1}\quad p. 411; designated-witness attribution.\par
\noindent Imprimé : {\ttfamily\footnotesize (\allowbreak{}}\par
\noindent Corrigé : \emph{(vide)}\par
\noindent Source : WITNESS-F7II-XXII-023.1. \nolinkurl{fr/sga7ii/diplomatic/expose_XXII_body.tex}; \nolinkurl{fr/sga7ii/current/expose_XXII_body.tex}\par

\par\medskip\noindent\hypertarget{app-F7II-XXII-023.2}{}\textbf{F7II-XXII-023.2}\quad p. 411; designated-witness attribution.\par
\noindent Imprimé : \emph{(vide)}\par
\noindent Corrigé : {\ttfamily\footnotesize (\allowbreak{}}\par
\noindent Source : WITNESS-F7II-XXII-023.2. \nolinkurl{fr/sga7ii/diplomatic/expose_XXII_body.tex}; \nolinkurl{fr/sga7ii/current/expose_XXII_body.tex}\par

\par\medskip\noindent\hypertarget{app-F7II-XXII-023.3}{}\textbf{F7II-XXII-023.3}\quad p. 411; printed misprint.\par
\noindent Imprimé : {\ttfamily\footnotesize p\allowbreak{}}\par
\noindent Corrigé : {\ttfamily\footnotesize q\allowbreak{}}\par
\noindent Source : 7II:708; 7II:886. \nolinkurl{fixes/MISPRINTS_7II.json#708}; \nolinkurl{fixes/MISPRINTS_7II.json#886}\par

\par\medskip\noindent\hypertarget{app-F7II-XXII-023.4}{}\textbf{F7II-XXII-023.4}\quad p. 411; designated-witness attribution.\par
\noindent Imprimé : {\ttfamily\footnotesize \{\allowbreak{}}\par
\noindent Corrigé : \emph{(vide)}\par
\noindent Source : WITNESS-F7II-XXII-023.4. \nolinkurl{fr/sga7ii/diplomatic/expose_XXII_body.tex}; \nolinkurl{fr/sga7ii/current/expose_XXII_body.tex}\par

\par\medskip\noindent\hypertarget{app-F7II-XXII-023.5}{}\textbf{F7II-XXII-023.5}\quad p. 411; designated-witness attribution.\par
\noindent Imprimé : \emph{(vide)}\par
\noindent Corrigé : {\ttfamily\footnotesize )\textasciicircum{}\allowbreak{}}\par
\noindent Source : WITNESS-F7II-XXII-023.5. \nolinkurl{fr/sga7ii/diplomatic/expose_XXII_body.tex}; \nolinkurl{fr/sga7ii/current/expose_XXII_body.tex}\par

\par\medskip\noindent\hypertarget{app-F7II-XXII-023.6}{}\textbf{F7II-XXII-023.6}\quad p. 411; designated-witness attribution.\par
\noindent Imprimé : {\ttfamily\footnotesize \}\allowbreak{}}\par
\noindent Corrigé : \emph{(vide)}\par
\noindent Source : WITNESS-F7II-XXII-023.6. \nolinkurl{fr/sga7ii/diplomatic/expose_XXII_body.tex}; \nolinkurl{fr/sga7ii/current/expose_XXII_body.tex}\par

\par\medskip\noindent\hypertarget{app-F7II-XXII-023.7}{}\textbf{F7II-XXII-023.7}\quad p. 411; designated-witness attribution.\par
\noindent Imprimé : {\ttfamily\footnotesize )\allowbreak{}}\par
\noindent Corrigé : {\ttfamily\footnotesize ,\allowbreak{}}\par
\noindent Source : WITNESS-F7II-XXII-023.7. \nolinkurl{fr/sga7ii/diplomatic/expose_XXII_body.tex}; \nolinkurl{fr/sga7ii/current/expose_XXII_body.tex}\par

\par\medskip\noindent\hypertarget{app-F7II-XXII-023.8}{}\textbf{F7II-XXII-023.8}\quad p. 411; designated-witness attribution.\par
\noindent Imprimé : {\ttfamily\footnotesize q\allowbreak{}}\par
\noindent Corrigé : {\ttfamily\footnotesize p\allowbreak{}}\par
\noindent Source : WITNESS-F7II-XXII-023.8. \nolinkurl{fr/sga7ii/diplomatic/expose_XXII_body.tex}; \nolinkurl{fr/sga7ii/current/expose_XXII_body.tex}\par

\par\medskip\noindent\hypertarget{app-F7II-XXII-024.1}{}\textbf{F7II-XXII-024.1}\quad p. 411; designated-witness attribution.\par
\noindent Imprimé : {\ttfamily\footnotesize q\allowbreak{}}\par
\noindent Corrigé : {\ttfamily\footnotesize p\allowbreak{}}\par
\noindent Source : WITNESS-F7II-XXII-024.1. \nolinkurl{fr/sga7ii/diplomatic/expose_XXII_body.tex}; \nolinkurl{fr/sga7ii/current/expose_XXII_body.tex}\par

\par\medskip\noindent\hypertarget{app-F7II-XXII-025.1}{}\textbf{F7II-XXII-025.1}\quad p. 412; designated-witness attribution.\par
\noindent Imprimé : \emph{(vide)}\par
\noindent Corrigé : {\ttfamily\footnotesize ,\allowbreak{}}\par
\noindent Source : WITNESS-F7II-XXII-025.1. \nolinkurl{fr/sga7ii/diplomatic/expose_XXII_body.tex}; \nolinkurl{fr/sga7ii/current/expose_XXII_body.tex}\par

\par\medskip\noindent\hypertarget{app-F7II-XXII-026.1}{}\textbf{F7II-XXII-026.1}\quad p. 412; misprint\_pass.\par
\noindent Imprimé : {\ttfamily\footnotesize b\allowbreak{}}\par
\noindent Corrigé : {\ttfamily\footnotesize a\allowbreak{}}\par
\noindent Source : 7II:888. \nolinkurl{fixes/FIXLOG_7II_CURRENT.jsonl#900}; \nolinkurl{fixes/MISPRINTS_7II.json#888}\par

\par\medskip\noindent\hypertarget{app-F7II-XXII-026.2}{}\textbf{F7II-XXII-026.2}\quad p. 412; misprint\_pass.\par
\noindent Imprimé : {\ttfamily\footnotesize a\allowbreak{}}\par
\noindent Corrigé : {\ttfamily\footnotesize b\allowbreak{}}\par
\noindent Source : 7II:888. \nolinkurl{fixes/FIXLOG_7II_CURRENT.jsonl#900}; \nolinkurl{fixes/MISPRINTS_7II.json#888}\par

\par\medskip\noindent\hypertarget{app-F7II-XXII-026.3}{}\textbf{F7II-XXII-026.3}\quad p. 412; misprint\_pass.\par
\noindent Imprimé : {\ttfamily\footnotesize b\allowbreak{}}\par
\noindent Corrigé : {\ttfamily\footnotesize a\allowbreak{}}\par
\noindent Source : 7II:888. \nolinkurl{fixes/FIXLOG_7II_CURRENT.jsonl#900}; \nolinkurl{fixes/MISPRINTS_7II.json#888}\par

\par\medskip\noindent\hypertarget{app-F7II-XXII-027.1}{}\textbf{F7II-XXII-027.1}\quad p. 413; designated-witness attribution.\par
\noindent Imprimé : {\ttfamily\footnotesize Moebius\allowbreak{}}\par
\noindent Corrigé : {\ttfamily\footnotesize Möbius\allowbreak{}}\par
\noindent Source : WITNESS-F7II-XXII-027.1. \nolinkurl{fr/sga7ii/diplomatic/expose_XXII_body.tex}; \nolinkurl{fr/sga7ii/current/expose_XXII_body.tex}\par

\par\medskip\noindent\hypertarget{app-F7II-XXII-028.1}{}\textbf{F7II-XXII-028.1}\quad p. 413; printed misprint.\par
\noindent Imprimé : \emph{(vide)}\par
\noindent Corrigé : {\ttfamily\footnotesize \{i\_0<\textbackslash{}cdots<i\_\allowbreak{}}\par
\noindent Source : 7II:889. \nolinkurl{fixes/MISPRINTS_7II.json#889}\par

\par\medskip\noindent\hypertarget{app-F7II-XXII-028.2}{}\textbf{F7II-XXII-028.2}\quad p. 413; printed misprint.\par
\noindent Imprimé : \emph{(vide)}\par
\noindent Corrigé : {\ttfamily\footnotesize \}\allowbreak{}}\par
\noindent Source : 7II:889. \nolinkurl{fixes/MISPRINTS_7II.json#889}\par

\par\medskip\noindent\hypertarget{app-F7II-XXII-029.1}{}\textbf{F7II-XXII-029.1}\quad p. 413; inherited corrected-layer difference.\par
\noindent Imprimé : \emph{(vide)}\par
\noindent Corrigé : {\ttfamily\footnotesize ,\allowbreak{}}\par
\noindent Source : ADOPT-F7II-XXII-029.1. \nolinkurl{fr/sga7ii/diplomatic/expose_XXII_body.tex}; \nolinkurl{fr/sga7ii/current/expose_XXII_body.tex}\par

\par\medskip\noindent\hypertarget{app-F7II-XXII-029.2}{}\textbf{F7II-XXII-029.2}\quad p. 413; designated-witness attribution.\par
\noindent Imprimé : {\ttfamily\footnotesize i\allowbreak{}}\par
\noindent Corrigé : {\ttfamily\footnotesize c'est-à-dire\allowbreak{}}\par
\noindent Source : WITNESS-F7II-XXII-029.2. \nolinkurl{fr/sga7ii/diplomatic/expose_XXII_body.tex}; \nolinkurl{fr/sga7ii/current/expose_XXII_body.tex}\par

\par\medskip\noindent\hypertarget{app-F7II-XXII-029.3}{}\textbf{F7II-XXII-029.3}\quad p. 413; inherited corrected-layer difference.\par
\noindent Imprimé : {\ttfamily\footnotesize .\textbackslash{}underline\{e\}.\allowbreak{}}\par
\noindent Corrigé : \emph{(vide)}\par
\noindent Source : ADOPT-F7II-XXII-029.3. \nolinkurl{fr/sga7ii/diplomatic/expose_XXII_body.tex}; \nolinkurl{fr/sga7ii/current/expose_XXII_body.tex}\par

\par\medskip\noindent\hypertarget{app-F7II-XXII-029.4}{}\textbf{F7II-XXII-029.4}\quad p. 413; inherited corrected-layer difference.\par
\noindent Imprimé : \emph{(vide)}\par
\noindent Corrigé : {\ttfamily\footnotesize  hypersurface\allowbreak{}}\par
\noindent Source : ADOPT-F7II-XXII-029.4. \nolinkurl{fr/sga7ii/diplomatic/expose_XXII_body.tex}; \nolinkurl{fr/sga7ii/current/expose_XXII_body.tex}\par

\par\medskip\noindent\hypertarget{app-F7II-XXII-029.5}{}\textbf{F7II-XXII-029.5}\quad p. 413; inherited corrected-layer difference.\par
\noindent Imprimé : {\ttfamily\footnotesize hypersurface\allowbreak{} }\par
\noindent Corrigé : \emph{(vide)}\par
\noindent Source : ADOPT-F7II-XXII-029.5. \nolinkurl{fr/sga7ii/diplomatic/expose_XXII_body.tex}; \nolinkurl{fr/sga7ii/current/expose_XXII_body.tex}\par

\par\medskip\noindent\hypertarget{app-F7II-XXII-030.1}{}\textbf{F7II-XXII-030.1}\quad p. 414; printed misprint.\par
\noindent Imprimé : \emph{(vide)}\par
\noindent Corrigé : {\ttfamily\footnotesize 4.\allowbreak{}}\par
\noindent Source : 7II:890. \nolinkurl{fixes/MISPRINTS_7II.json#890}\par

\par\medskip\noindent\hypertarget{app-F7II-XXII-030.2}{}\textbf{F7II-XXII-030.2}\quad p. 414; printed misprint.\par
\noindent Imprimé : {\ttfamily\footnotesize .3\allowbreak{}}\par
\noindent Corrigé : \emph{(vide)}\par
\noindent Source : 7II:890. \nolinkurl{fixes/MISPRINTS_7II.json#890}\par

\par\medskip\noindent\hypertarget{app-F7II-XXII-030.3}{}\textbf{F7II-XXII-030.3}\quad p. 414; designated-witness attribution.\par
\noindent Imprimé : {\ttfamily\footnotesize vrai\allowbreak{}}\par
\noindent Corrigé : {\ttfamily\footnotesize vraie\allowbreak{}}\par
\noindent Source : WITNESS-F7II-XXII-030.3. \nolinkurl{fr/sga7ii/diplomatic/expose_XXII_body.tex}; \nolinkurl{fr/sga7ii/current/expose_XXII_body.tex}\par

\par\medskip\noindent\hypertarget{app-F7II-XXII-030.4}{}\textbf{F7II-XXII-030.4}\quad p. 414; designated-witness attribution.\par
\noindent Imprimé : {\ttfamily\footnotesize i.e.\allowbreak{}}\par
\noindent Corrigé : {\ttfamily\footnotesize c'est-à-dire\allowbreak{}}\par
\noindent Source : WITNESS-F7II-XXII-030.4. \nolinkurl{fr/sga7ii/diplomatic/expose_XXII_body.tex}; \nolinkurl{fr/sga7ii/current/expose_XXII_body.tex}\par

\par\medskip\noindent\hypertarget{app-F7II-XXII-030.5}{}\textbf{F7II-XXII-030.5}\quad p. 414; designated-witness attribution.\par
\noindent Imprimé : \emph{(vide)}\par
\noindent Corrigé : {\ttfamily\footnotesize  Une\allowbreak{}}\par
\noindent Source : WITNESS-F7II-XXII-030.5. \nolinkurl{fr/sga7ii/diplomatic/expose_XXII_body.tex}; \nolinkurl{fr/sga7ii/current/expose_XXII_body.tex}\par

\par\medskip\noindent\hypertarget{app-F7II-XXII-030.6}{}\textbf{F7II-XXII-030.6}\quad p. 414; designated-witness attribution.\par
\noindent Imprimé : {\ttfamily\footnotesize Une\allowbreak{} }\par
\noindent Corrigé : \emph{(vide)}\par
\noindent Source : WITNESS-F7II-XXII-030.6. \nolinkurl{fr/sga7ii/diplomatic/expose_XXII_body.tex}; \nolinkurl{fr/sga7ii/current/expose_XXII_body.tex}\par

\par\medskip\noindent\hypertarget{app-F7II-XXII-030.7}{}\textbf{F7II-XXII-030.7}\quad p. 414; designated-witness attribution.\par
\noindent Imprimé : \emph{(vide)}\par
\noindent Corrigé : {\ttfamily\footnotesize (\allowbreak{}}\par
\noindent Source : WITNESS-F7II-XXII-030.7. \nolinkurl{fr/sga7ii/diplomatic/expose_XXII_body.tex}; \nolinkurl{fr/sga7ii/current/expose_XXII_body.tex}\par

\par\medskip\noindent\hypertarget{app-F7II-XXII-030.8}{}\textbf{F7II-XXII-030.8}\quad p. 414; designated-witness attribution.\par
\noindent Imprimé : \emph{(vide)}\par
\noindent Corrigé : {\ttfamily\footnotesize 4.\allowbreak{}}\par
\noindent Source : WITNESS-F7II-XXII-030.8. \nolinkurl{fr/sga7ii/diplomatic/expose_XXII_body.tex}; \nolinkurl{fr/sga7ii/current/expose_XXII_body.tex}\par

\par\medskip\noindent\hypertarget{app-F7II-XXII-030.9}{}\textbf{F7II-XXII-030.9}\quad p. 414; designated-witness attribution.\par
\noindent Imprimé : {\ttfamily\footnotesize .3\allowbreak{}}\par
\noindent Corrigé : {\ttfamily\footnotesize )\allowbreak{}}\par
\noindent Source : WITNESS-F7II-XXII-030.9. \nolinkurl{fr/sga7ii/diplomatic/expose_XXII_body.tex}; \nolinkurl{fr/sga7ii/current/expose_XXII_body.tex}\par

\par\medskip\noindent\hypertarget{app-F7II-XXII-031.1}{}\textbf{F7II-XXII-031.1}\quad p. 414; misprint\_in\_print.\par
\noindent Imprimé : {\ttfamily\footnotesize \textbackslash{}chi\allowbreak{}}\par
\noindent Corrigé : {\ttfamily\footnotesize x\allowbreak{}}\par
\noindent Source : FIX-7II-D-0405; 7II:892. \nolinkurl{fixes/FIXLOG_7II.jsonl#405}; \nolinkurl{fixes/MISPRINTS_7II.json#892}\par

\par\medskip\noindent\hypertarget{app-F7II-XXII-031.2}{}\textbf{F7II-XXII-031.2}\quad p. 414; misprint\_in\_print.\par
\noindent Imprimé : {\ttfamily\footnotesize \textbackslash{}chi\allowbreak{}}\par
\noindent Corrigé : {\ttfamily\footnotesize x\allowbreak{}}\par
\noindent Source : FIX-7II-D-0406; 7II:893. \nolinkurl{fixes/FIXLOG_7II.jsonl#406}; \nolinkurl{fixes/MISPRINTS_7II.json#893}\par

\par\medskip\noindent\hypertarget{app-F7II-XXII-031.3}{}\textbf{F7II-XXII-031.3}\quad p. 414; printed misprint.\par
\noindent Imprimé : {\ttfamily\footnotesize \textbackslash{}chi\allowbreak{}}\par
\noindent Corrigé : {\ttfamily\footnotesize x\allowbreak{}}\par
\noindent Source : 7II:894. \nolinkurl{fixes/MISPRINTS_7II.json#894}\par

\par\medskip\noindent\hypertarget{app-F7II-XXII-031.4}{}\textbf{F7II-XXII-031.4}\quad p. 414; printed misprint.\par
\noindent Imprimé : {\ttfamily\footnotesize \textbackslash{}chi\allowbreak{}}\par
\noindent Corrigé : {\ttfamily\footnotesize x\allowbreak{}}\par
\noindent Source : 7II:894. \nolinkurl{fixes/MISPRINTS_7II.json#894}\par

\par\medskip\noindent\hypertarget{app-F7II-XXII-031.5}{}\textbf{F7II-XXII-031.5}\quad p. 414; printed misprint.\par
\noindent Imprimé : {\ttfamily\footnotesize \textbackslash{}chi\allowbreak{}}\par
\noindent Corrigé : {\ttfamily\footnotesize x\allowbreak{}}\par
\noindent Source : 7II:894. \nolinkurl{fixes/MISPRINTS_7II.json#894}\par

\par\medskip\noindent\hypertarget{app-F7II-XXII-032.1}{}\textbf{F7II-XXII-032.1}\quad p. 415; printed misprint.\par
\noindent Imprimé : {\ttfamily\footnotesize X\allowbreak{}}\par
\noindent Corrigé : {\ttfamily\footnotesize x\allowbreak{}}\par
\noindent Source : 7II:895. \nolinkurl{fixes/MISPRINTS_7II.json#895}\par

\par\medskip\noindent\hypertarget{app-F7II-XXII-033.1}{}\textbf{F7II-XXII-033.1}\quad p. 415; printed misprint.\par
\noindent Imprimé : \emph{(vide)}\par
\noindent Corrigé : {\ttfamily\footnotesize /p\allowbreak{}}\par
\noindent Source : 7II:896. \nolinkurl{fixes/MISPRINTS_7II.json#896}\par

\par\medskip\noindent\hypertarget{app-F7II-XXII-033.2}{}\textbf{F7II-XXII-033.2}\quad p. 415; printed misprint.\par
\noindent Imprimé : {\ttfamily\footnotesize .\allowbreak{}}\par
\noindent Corrigé : {\ttfamily\footnotesize )\allowbreak{}}\par
\noindent Source : 7II:896. \nolinkurl{fixes/MISPRINTS_7II.json#896}\par

\par\medskip\noindent\hypertarget{app-F7II-XXII-033.3}{}\textbf{F7II-XXII-033.3}\quad p. 415; misprint\_in\_print.\par
\noindent Imprimé : {\ttfamily\footnotesize H\allowbreak{}}\par
\noindent Corrigé : {\ttfamily\footnotesize F\allowbreak{}}\par
\noindent Source : FIX-7II-D-0717. \nolinkurl{fixes/FIXLOG_7II.jsonl#717}\par

\par\medskip\noindent\hypertarget{app-F7II-XXII-034.1}{}\textbf{F7II-XXII-034.1}\quad p. 416; misprint\_in\_print.\par
\noindent Imprimé : \emph{(vide)}\par
\noindent Corrigé : {\ttfamily\footnotesize \textbackslash{}mathbb\allowbreak{} }\par
\noindent Source : FIX-7II-D-0667. \nolinkurl{fixes/FIXLOG_7II.jsonl#667}\par

\par\medskip\noindent\hypertarget{app-F7II-XXII-035.1}{}\textbf{F7II-XXII-035.1}\quad p. 416; misprint\_in\_print.\par
\noindent Imprimé : \emph{(vide)}\par
\noindent Corrigé : {\ttfamily\footnotesize \textbackslash{}mathbb\allowbreak{}}\par
\noindent Source : FIX-7II-D-0668. \nolinkurl{fixes/FIXLOG_7II.jsonl#668}\par

\par\medskip\noindent\hypertarget{app-F7II-XXII-036.1}{}\textbf{F7II-XXII-036.1}\quad p. 416; printed misprint.\par
\noindent Imprimé : {\ttfamily\footnotesize A\allowbreak{}}\par
\noindent Corrigé : {\ttfamily\footnotesize B\allowbreak{}}\par
\noindent Source : 7II:715. \nolinkurl{fixes/MISPRINTS_7II.json#715}\par

\par\medskip\noindent\hypertarget{app-F7II-XXII-036.2}{}\textbf{F7II-XXII-036.2}\quad p. 416; printed misprint.\par
\noindent Imprimé : {\ttfamily\footnotesize '\allowbreak{}}\par
\noindent Corrigé : {\ttfamily\footnotesize -\allowbreak{}}\par
\noindent Source : 7II:854. \nolinkurl{fixes/MISPRINTS_7II.json#854}\par

\par\medskip\noindent\hypertarget{app-F7II-XXII-036.3}{}\textbf{F7II-XXII-036.3}\quad p. 416; printed misprint.\par
\noindent Imprimé : {\ttfamily\footnotesize q\allowbreak{}}\par
\noindent Corrigé : {\ttfamily\footnotesize p\allowbreak{}}\par
\noindent Source : 7II:716. \nolinkurl{fixes/MISPRINTS_7II.json#716}\par

\par\medskip\noindent\hypertarget{app-F7II-XXII-037.1}{}\textbf{F7II-XXII-037.1}\quad p. 417; printed misprint.\par
\noindent Imprimé : \emph{(vide)}\par
\noindent Corrigé : {\ttfamily\footnotesize \{p\textasciicircum{}\allowbreak{}}\par
\noindent Source : 7II:717; 7II:856. \nolinkurl{fixes/MISPRINTS_7II.json#717}; \nolinkurl{fixes/MISPRINTS_7II.json#856}\par

\par\medskip\noindent\hypertarget{app-F7II-XXII-037.2}{}\textbf{F7II-XXII-037.2}\quad p. 417; designated-witness attribution.\par
\noindent Imprimé : \emph{(vide)}\par
\noindent Corrigé : {\ttfamily\footnotesize \}\allowbreak{}}\par
\noindent Source : WITNESS-F7II-XXII-037.2. \nolinkurl{fr/sga7ii/diplomatic/expose_XXII_body.tex}; \nolinkurl{fr/sga7ii/current/expose_XXII_body.tex}\par

\par\medskip\noindent\hypertarget{app-F7II-XXII-037.3}{}\textbf{F7II-XXII-037.3}\quad p. 417; misprint\_in\_print.\par
\noindent Imprimé : {\ttfamily\footnotesize qui\allowbreak{}}\par
\noindent Corrigé : {\ttfamily\footnotesize que\allowbreak{}}\par
\noindent Source : FIX-7II-D-0427; 7II:857. \nolinkurl{fixes/FIXLOG_7II.jsonl#427}; \nolinkurl{fixes/MISPRINTS_7II.json#857}\par

\par\medskip\noindent\hypertarget{app-F7II-XXII-038.1}{}\textbf{F7II-XXII-038.1}\quad p. 417; designated-witness attribution.\par
\noindent Imprimé : {\ttfamily\footnotesize q\allowbreak{}}\par
\noindent Corrigé : {\ttfamily\footnotesize p\allowbreak{}}\par
\noindent Source : WITNESS-F7II-XXII-038.1. \nolinkurl{fr/sga7ii/diplomatic/expose_XXII_body.tex}; \nolinkurl{fr/sga7ii/current/expose_XXII_body.tex}\par

\par\medskip\noindent\hypertarget{app-F7II-XXII-038.2}{}\textbf{F7II-XXII-038.2}\quad p. 417; designated-witness attribution.\par
\noindent Imprimé : {\ttfamily\footnotesize q\allowbreak{}}\par
\noindent Corrigé : {\ttfamily\footnotesize p\allowbreak{}}\par
\noindent Source : WITNESS-F7II-XXII-038.2. \nolinkurl{fr/sga7ii/diplomatic/expose_XXII_body.tex}; \nolinkurl{fr/sga7ii/current/expose_XXII_body.tex}\par

\par\medskip\noindent\hypertarget{app-F7II-XXII-039.1}{}\textbf{F7II-XXII-039.1}\quad p. 418; misprint\_in\_print.\par
\noindent Imprimé : \emph{(vide)}\par
\noindent Corrigé : {\ttfamily\footnotesize \textbackslash{}mathbb\allowbreak{} }\par
\noindent Source : FIX-7II-D-0674. \nolinkurl{fixes/FIXLOG_7II.jsonl#674}\par

\par\medskip\noindent\hypertarget{app-F7II-XXII-040.1}{}\textbf{F7II-XXII-040.1}\quad p. 418; printed misprint.\par
\noindent Imprimé : {\ttfamily\footnotesize X\allowbreak{}}\par
\noindent Corrigé : {\ttfamily\footnotesize \textbackslash{}chi\allowbreak{}}\par
\noindent Source : 7II:861. \nolinkurl{fixes/MISPRINTS_7II.json#861}\par

\par\medskip\noindent\hypertarget{app-F7II-XXII-041.1}{}\textbf{F7II-XXII-041.1}\quad p. 419; designated-witness attribution.\par
\noindent Imprimé : \emph{(vide)}\par
\noindent Corrigé : {\ttfamily\footnotesize \_1\allowbreak{}}\par
\noindent Source : WITNESS-F7II-XXII-041.1. \nolinkurl{fr/sga7ii/diplomatic/expose_XXII_body.tex}; \nolinkurl{fr/sga7ii/current/expose_XXII_body.tex}\par

\par\medskip\noindent\hypertarget{app-F7II-XXII-042.1}{}\textbf{F7II-XXII-042.1}\quad p. 421; printed misprint.\par
\noindent Imprimé : \emph{(vide)}\par
\noindent Corrigé : {\ttfamily\footnotesize \textasciicircum{}\textbackslash{}nu\allowbreak{}}\par
\noindent Source : 7II:723. \nolinkurl{fixes/MISPRINTS_7II.json#723}\par

\par\medskip\noindent\hypertarget{app-F7II-XXII-042.2}{}\textbf{F7II-XXII-042.2}\quad p. 421; printed misprint.\par
\noindent Imprimé : \emph{(vide)}\par
\noindent Corrigé : {\ttfamily\footnotesize \{\allowbreak{}}\par
\noindent Source : 7II:865. \nolinkurl{fixes/MISPRINTS_7II.json#865}\par

\par\medskip\noindent\hypertarget{app-F7II-XXII-042.3}{}\textbf{F7II-XXII-042.3}\quad p. 421; printed misprint.\par
\noindent Imprimé : \emph{(vide)}\par
\noindent Corrigé : {\ttfamily\footnotesize \textasciicircum{}\textbackslash{}nu\}\allowbreak{}}\par
\noindent Source : 7II:865. \nolinkurl{fixes/MISPRINTS_7II.json#865}\par

\par\medskip\noindent\hypertarget{app-F7II-XXII-042.4}{}\textbf{F7II-XXII-042.4}\quad p. 421; misprint\_in\_print.\par
\noindent Imprimé : \emph{(vide)}\par
\noindent Corrigé : {\ttfamily\footnotesize \_k\allowbreak{}}\par
\noindent Source : FIX-7II-D-0688. \nolinkurl{fixes/FIXLOG_7II.jsonl#688}\par

\par\medskip\noindent\hypertarget{app-F7II-XXII-043.1}{}\textbf{F7II-XXII-043.1}\quad p. 421; printed misprint.\par
\noindent Imprimé : \emph{(vide)}\par
\noindent Corrigé : {\ttfamily\footnotesize \textasciicircum{}\textbackslash{}nu\allowbreak{}}\par
\noindent Source : 7II:725; 7II:867. \nolinkurl{fixes/MISPRINTS_7II.json#725}; \nolinkurl{fixes/MISPRINTS_7II.json#867}\par

\par\medskip\noindent\hypertarget{app-F7II-XXII-043.2}{}\textbf{F7II-XXII-043.2}\quad p. 421; printed misprint.\par
\noindent Imprimé : \emph{(vide)}\par
\noindent Corrigé : {\ttfamily\footnotesize \{\allowbreak{}}\par
\noindent Source : 7II:868. \nolinkurl{fixes/MISPRINTS_7II.json#868}\par

\par\medskip\noindent\hypertarget{app-F7II-XXII-043.3}{}\textbf{F7II-XXII-043.3}\quad p. 421; printed misprint.\par
\noindent Imprimé : \emph{(vide)}\par
\noindent Corrigé : {\ttfamily\footnotesize \textasciicircum{}\textbackslash{}nu\}\allowbreak{}}\par
\noindent Source : 7II:868. \nolinkurl{fixes/MISPRINTS_7II.json#868}\par

\par\medskip\noindent\hypertarget{app-F7II-XXII-044.1}{}\textbf{F7II-XXII-044.1}\quad p. 422; printed misprint.\par
\noindent Imprimé : \emph{(vide)}\par
\noindent Corrigé : {\ttfamily\footnotesize 4.\allowbreak{}}\par
\noindent Source : 7II:727; 7II:869. \nolinkurl{fixes/MISPRINTS_7II.json#727}; \nolinkurl{fixes/MISPRINTS_7II.json#869}\par

\par\medskip\noindent\hypertarget{app-F7II-XXII-044.2}{}\textbf{F7II-XXII-044.2}\quad p. 422; printed misprint.\par
\noindent Imprimé : {\ttfamily\footnotesize .4\allowbreak{}}\par
\noindent Corrigé : \emph{(vide)}\par
\noindent Source : 7II:727; 7II:869. \nolinkurl{fixes/MISPRINTS_7II.json#727}; \nolinkurl{fixes/MISPRINTS_7II.json#869}\par

\par\medskip\noindent\hypertarget{app-F7II-XXII-045.1}{}\textbf{F7II-XXII-045.1}\quad p. 422; printed misprint.\par
\noindent Imprimé : {\ttfamily\footnotesize .\allowbreak{}}\par
\noindent Corrigé : {\ttfamily\footnotesize )\allowbreak{}}\par
\noindent Source : 7II:870. \nolinkurl{fixes/MISPRINTS_7II.json#870}\par

\par\medskip\noindent\hypertarget{app-F7II-XXII-046.1}{}\textbf{F7II-XXII-046.1}\quad p. 423; inherited corrected-layer difference.\par
\noindent Imprimé : {\ttfamily\footnotesize \textbackslash{}subseteq\allowbreak{}}\par
\noindent Corrigé : {\ttfamily\footnotesize \textbackslash{}supset\allowbreak{}}\par
\noindent Source : ADOPT-F7II-XXII-046.1. \nolinkurl{fr/sga7ii/diplomatic/expose_XXII_body.tex}; \nolinkurl{fr/sga7ii/current/expose_XXII_body.tex}\par

\par\medskip\noindent\hypertarget{app-F7II-XXII-047.1}{}\textbf{F7II-XXII-047.1}\quad p. 423; printed misprint.\par
\noindent Imprimé : {\ttfamily\footnotesize 2.3\allowbreak{}}\par
\noindent Corrigé : {\ttfamily\footnotesize 23\allowbreak{}}\par
\noindent Source : 7II:729; 7II:872. \nolinkurl{fixes/MISPRINTS_7II.json#729}; \nolinkurl{fixes/MISPRINTS_7II.json#872}\par

\par\medskip\noindent\hypertarget{app-F7II-XXII-048.1}{}\textbf{F7II-XXII-048.1}\quad p. 423; printed misprint.\par
\noindent Imprimé : \emph{(vide)}\par
\noindent Corrigé : {\ttfamily\footnotesize en\allowbreak{} }\par
\noindent Source : 7II:873. \nolinkurl{fixes/MISPRINTS_7II.json#873}\par

\par\medskip\noindent\hypertarget{app-F7II-XXII-048.2}{}\textbf{F7II-XXII-048.2}\quad p. 423; designated-witness attribution.\par
\noindent Imprimé : {\ttfamily\footnotesize  en\allowbreak{}}\par
\noindent Corrigé : \emph{(vide)}\par
\noindent Source : WITNESS-F7II-XXII-048.2. \nolinkurl{fr/sga7ii/diplomatic/expose_XXII_body.tex}; \nolinkurl{fr/sga7ii/current/expose_XXII_body.tex}\par

\par\medskip\noindent\hypertarget{app-F7II-XXII-049.1}{}\textbf{F7II-XXII-049.1}\quad p. 424; printed misprint.\par
\noindent Imprimé : {\ttfamily\footnotesize car\allowbreak{}}\par
\noindent Corrigé : {\ttfamily\footnotesize card\allowbreak{}}\par
\noindent Source : 7II:731. \nolinkurl{fixes/MISPRINTS_7II.json#731}\par

\par\medskip\noindent\hypertarget{app-F7II-XXII-049.2}{}\textbf{F7II-XXII-049.2}\quad p. 424; printed misprint.\par
\noindent Imprimé : {\ttfamily\footnotesize )\allowbreak{}}\par
\noindent Corrigé : \emph{(vide)}\par
\noindent Source : 7II:875. \nolinkurl{fixes/MISPRINTS_7II.json#875}\par

\par\medskip\noindent\hypertarget{app-F7II-XXII-050.1}{}\textbf{F7II-XXII-050.1}\quad p. 424; misprint\_in\_print.\par
\noindent Imprimé : \emph{(1 espace(s), 0 retour(s) de ligne dans le témoin TeX)}\par
\noindent Corrigé : {\ttfamily\footnotesize -\allowbreak{}}\par
\noindent Source : FIX-7II-D-0471. \nolinkurl{fixes/FIXLOG_7II.jsonl#471}\par

\par\medskip\noindent\hypertarget{app-F7II-XXII-051.1}{}\textbf{F7II-XXII-051.1}\quad p. 425; printed misprint.\par
\noindent Imprimé : {\ttfamily\footnotesize 23\allowbreak{}}\par
\noindent Corrigé : {\ttfamily\footnotesize 33\allowbreak{}}\par
\noindent Source : 7II:734; 7II:876. \nolinkurl{fixes/MISPRINTS_7II.json#734}; \nolinkurl{fixes/MISPRINTS_7II.json#876}\par

\par\medskip\noindent\hypertarget{app-F7II-XXII-052.1}{}\textbf{F7II-XXII-052.1}\quad p. 426; printed misprint.\par
\noindent Imprimé : \emph{(vide)}\par
\noindent Corrigé : {\ttfamily\footnotesize est\allowbreak{}}\par
\noindent Source : 7II:877. \nolinkurl{fixes/MISPRINTS_7II.json#877}\par

\par\medskip\noindent\hypertarget{app-F7II-XXII-053.1}{}\textbf{F7II-XXII-053.1}\quad p. 426; misprint\_in\_print.\par
\noindent Imprimé : {\ttfamily\footnotesize ordinal\allowbreak{}}\par
\noindent Corrigé : {\ttfamily\footnotesize ordre\allowbreak{}}\par
\noindent Source : FIX-7II-D-0483. \nolinkurl{fixes/FIXLOG_7II.jsonl#483}\par

\par\medskip\noindent\hypertarget{app-F7II-XXII-054.1}{}\textbf{F7II-XXII-054.1}\quad p. 427; misprint\_in\_print.\par
\noindent Imprimé : \emph{(vide)}\par
\noindent Corrigé : {\ttfamily\footnotesize \_0\allowbreak{}}\par
\noindent Source : FIX-7II-D-0489. \nolinkurl{fixes/FIXLOG_7II.jsonl#489}\par

\par\medskip\noindent\hypertarget{app-F7II-XXII-055.1}{}\textbf{F7II-XXII-055.1}\quad p. 428; printed misprint.\par
\noindent Imprimé : {\ttfamily\footnotesize 7\allowbreak{}}\par
\noindent Corrigé : {\ttfamily\footnotesize 4\allowbreak{}}\par
\noindent Source : 7II:738; 7II:879. \nolinkurl{fixes/MISPRINTS_7II.json#738}; \nolinkurl{fixes/MISPRINTS_7II.json#879}\par

\par\medskip\noindent\hypertarget{app-F7II-XXII-056.1}{}\textbf{F7II-XXII-056.1}\quad p. 429; misprint\_pass.\par
\noindent Imprimé : {\ttfamily\footnotesize 4\allowbreak{}}\par
\noindent Corrigé : {\ttfamily\footnotesize 5\allowbreak{}}\par
\noindent Source : 7II:740. \nolinkurl{fixes/FIXLOG_7II_CURRENT.jsonl#769}\par

\par\medskip\noindent\hypertarget{app-F7II-XXII-057.1}{}\textbf{F7II-XXII-057.1}\quad p. 429; misprint\_pass.\par
\noindent Imprimé : {\ttfamily\footnotesize t\allowbreak{}}\par
\noindent Corrigé : {\ttfamily\footnotesize qt\allowbreak{}}\par
\noindent Source : 7II:881. \nolinkurl{fixes/FIXLOG_7II_CURRENT.jsonl#902}\par

\par\medskip\noindent\hypertarget{app-F7II-XXII-057.2}{}\textbf{F7II-XXII-057.2}\quad p. 429; misprint\_pass.\par
\noindent Imprimé : {\ttfamily\footnotesize qt\allowbreak{}}\par
\noindent Corrigé : {\ttfamily\footnotesize t\allowbreak{}}\par
\noindent Source : 7II:881. \nolinkurl{fixes/FIXLOG_7II_CURRENT.jsonl#902}\par

\par\medskip\noindent\hypertarget{app-F7II-XXII-058.1}{}\textbf{F7II-XXII-058.1}\quad p. 429; printed misprint.\par
\noindent Imprimé : {\ttfamily\footnotesize if\allowbreak{}}\par
\noindent Corrigé : {\ttfamily\footnotesize it\allowbreak{}}\par
\noindent Source : 7II:740; 7II:882. \nolinkurl{fixes/MISPRINTS_7II.json#740}; \nolinkurl{fixes/MISPRINTS_7II.json#882}\par

\par\medskip\noindent\hypertarget{app-F7II-XXII-058.2}{}\textbf{F7II-XXII-058.2}\quad p. 429; misprint\_pass.\par
\noindent Imprimé : \emph{(vide)}\par
\noindent Corrigé : {\ttfamily\footnotesize \textasciicircum{}\{-1\}\allowbreak{}}\par
\noindent Source : 7II:881. \nolinkurl{fixes/FIXLOG_7II_CURRENT.jsonl#902}\par

\par\medskip\noindent\hypertarget{app-F7II-XXII-059.1}{}\textbf{F7II-XXII-059.1}\quad p. 429; misprint\_in\_print.\par
\noindent Imprimé : {\ttfamily\footnotesize \textbackslash{}Lambda\allowbreak{}}\par
\noindent Corrigé : {\ttfamily\footnotesize \textbackslash{}Delta\allowbreak{}}\par
\noindent Source : FIX-7II-D-0696. \nolinkurl{fixes/FIXLOG_7II.jsonl#696}\par

\par\medskip\noindent\hypertarget{app-F7II-XXII-059.2}{}\textbf{F7II-XXII-059.2}\quad p. 429; printed misprint.\par
\noindent Imprimé : {\ttfamily\footnotesize 4\allowbreak{}}\par
\noindent Corrigé : {\ttfamily\footnotesize 6\allowbreak{}}\par
\noindent Source : 7II:884. \nolinkurl{fixes/MISPRINTS_7II.json#884}\par

\par\medskip\noindent\hypertarget{app-F7II-XXII-060.1}{}\textbf{F7II-XXII-060.1}\quad p. 430; misprint\_in\_print.\par
\noindent Imprimé : {\ttfamily\footnotesize .\allowbreak{}}\par
\noindent Corrigé : {\ttfamily\footnotesize )\allowbreak{}}\par
\noindent Source : FIX-7II-D-0518. \nolinkurl{fixes/FIXLOG_7II.jsonl#518}\par

\par\medskip\noindent\hypertarget{app-F7II-XXII-061.1}{}\textbf{F7II-XXII-061.1}\quad p. 431; misprint\_in\_print.\par
\noindent Imprimé : {\ttfamily\footnotesize .\allowbreak{}}\par
\noindent Corrigé : \emph{(vide)}\par
\noindent Source : FIX-7II-D-0522. \nolinkurl{fixes/FIXLOG_7II.jsonl#522}\par

\par\medskip\noindent\hypertarget{app-F7II-XXII-061.2}{}\textbf{F7II-XXII-061.2}\quad p. 431; designated-witness attribution.\par
\noindent Imprimé : {\ttfamily\footnotesize relations\allowbreak{}}\par
\noindent Corrigé : {\ttfamily\footnotesize relation\allowbreak{}}\par
\noindent Source : WITNESS-F7II-XXII-061.2. \nolinkurl{fr/sga7ii/diplomatic/expose_XXII_body.tex}; \nolinkurl{fr/sga7ii/current/expose_XXII_body.tex}\par

\par\medskip\noindent\hypertarget{app-F7II-XXII-062.1}{}\textbf{F7II-XXII-062.1}\quad p. 431; inherited corrected-layer difference.\par
\noindent Imprimé : {\ttfamily\footnotesize 24\allowbreak{}}\par
\noindent Corrigé : {\ttfamily\footnotesize 33\allowbreak{}}\par
\noindent Source : ADOPT-F7II-XXII-062.1. \nolinkurl{fr/sga7ii/diplomatic/expose_XXII_body.tex}; \nolinkurl{fr/sga7ii/current/expose_XXII_body.tex}\par

\par\medskip\noindent\hypertarget{app-F7II-XXII-063.1}{}\textbf{F7II-XXII-063.1}\quad p. 431; misprint\_in\_print.\par
\noindent Imprimé : \emph{(vide)}\par
\noindent Corrigé : {\ttfamily\footnotesize )\allowbreak{}}\par
\noindent Source : FIX-7II-D-0530. \nolinkurl{fixes/FIXLOG_7II.jsonl#530}\par

\par\medskip\noindent\hypertarget{app-F7II-XXII-063.2}{}\textbf{F7II-XXII-063.2}\quad p. 431; printed misprint.\par
\noindent Imprimé : {\ttfamily\footnotesize B\allowbreak{}}\par
\noindent Corrigé : {\ttfamily\footnotesize S\allowbreak{}}\par
\noindent Source : 7II:744. \nolinkurl{fixes/MISPRINTS_7II.json#744}\par

\par\medskip\noindent\hypertarget{app-F7II-XXII-064.1}{}\textbf{F7II-XXII-064.1}\quad p. 432; misprint\_in\_print.\par
\noindent Imprimé : {\ttfamily\footnotesize t\allowbreak{}}\par
\noindent Corrigé : {\ttfamily\footnotesize i\allowbreak{}}\par
\noindent Source : FIX-7II-D-0533. \nolinkurl{fixes/FIXLOG_7II.jsonl#533}\par

\par\medskip\noindent\hypertarget{app-F7II-XXII-064.2}{}\textbf{F7II-XXII-064.2}\quad p. 432; misprint\_in\_print.\par
\noindent Imprimé : {\ttfamily\footnotesize forme\allowbreak{}}\par
\noindent Corrigé : {\ttfamily\footnotesize formé\allowbreak{}}\par
\noindent Source : FIX-7II-D-0534. \nolinkurl{fixes/FIXLOG_7II.jsonl#534}\par

\par\medskip\noindent\hypertarget{app-F7II-XXII-065.1}{}\textbf{F7II-XXII-065.1}\quad p. 433; misprint\_in\_print.\par
\noindent Imprimé : \emph{(vide)}\par
\noindent Corrigé : {\ttfamily\footnotesize -\allowbreak{}}\par
\noindent Source : FIX-7II-D-0544. \nolinkurl{fixes/FIXLOG_7II.jsonl#544}\par

\par\medskip\noindent\hypertarget{app-F7II-XXII-066.1}{}\textbf{F7II-XXII-066.1}\quad p. 434; misprint\_in\_print.\par
\noindent Imprimé : \emph{(vide)}\par
\noindent Corrigé : {\ttfamily\footnotesize (\allowbreak{}}\par
\noindent Source : FIX-7II-D-0552. \nolinkurl{fixes/FIXLOG_7II.jsonl#552}\par

\par\medskip\noindent\hypertarget{app-F7II-XXII-067.1}{}\textbf{F7II-XXII-067.1}\quad p. 434; misprint\_in\_print.\par
\noindent Imprimé : \emph{(vide)}\par
\noindent Corrigé : {\ttfamily\footnotesize )\allowbreak{}}\par
\noindent Source : FIX-7II-D-0559. \nolinkurl{fixes/FIXLOG_7II.jsonl#559}\par

\par\medskip\noindent\hypertarget{app-F7II-XXII-068.1}{}\textbf{F7II-XXII-068.1}\quad p. 435; designated-witness attribution.\par
\noindent Imprimé : \emph{(vide)}\par
\noindent Corrigé : {\ttfamily\footnotesize \textbackslash{}circ\allowbreak{}}\par
\noindent Source : WITNESS-F7II-XXII-068.1. \nolinkurl{fr/sga7ii/diplomatic/expose_XXII_body.tex}; \nolinkurl{fr/sga7ii/current/expose_XXII_body.tex}\par

\par\medskip\noindent\hypertarget{app-F7II-XXII-068.2}{}\textbf{F7II-XXII-068.2}\quad p. 435; printed misprint.\par
\noindent Imprimé : {\ttfamily\footnotesize a\allowbreak{}}\par
\noindent Corrigé : {\ttfamily\footnotesize q\allowbreak{}}\par
\noindent Source : 7II:750. \nolinkurl{fixes/MISPRINTS_7II.json#750}\par

\par\medskip\noindent\hypertarget{app-F7II-XXII-068.3}{}\textbf{F7II-XXII-068.3}\quad p. 435; extra.\par
\noindent Imprimé : \emph{(vide)}\par
\noindent Corrigé : {\ttfamily\footnotesize \textbackslash{}circ\allowbreak{}}\par
\noindent Source : FIX-7II-D-0565. \nolinkurl{fixes/FIXLOG_7II.jsonl#565}\par

\par\medskip\noindent\hypertarget{app-F7II-XXII-068.4}{}\textbf{F7II-XXII-068.4}\quad p. 435; inherited corrected-layer difference.\par
\noindent Imprimé : \emph{(vide)}\par
\noindent Corrigé : {\ttfamily\footnotesize ,\allowbreak{}}\par
\noindent Source : ADOPT-F7II-XXII-068.4. \nolinkurl{fr/sga7ii/diplomatic/expose_XXII_body.tex}; \nolinkurl{fr/sga7ii/current/expose_XXII_body.tex}\par

\par\medskip\noindent\hypertarget{app-F7II-XXII-068.5}{}\textbf{F7II-XXII-068.5}\quad p. 435; inherited corrected-layer difference.\par
\noindent Imprimé : {\ttfamily\footnotesize ,\allowbreak{}}\par
\noindent Corrigé : \emph{(vide)}\par
\noindent Source : ADOPT-F7II-XXII-068.5. \nolinkurl{fr/sga7ii/diplomatic/expose_XXII_body.tex}; \nolinkurl{fr/sga7ii/current/expose_XXII_body.tex}\par

\par\medskip\noindent\hypertarget{app-F7II-XXII-069.1}{}\textbf{F7II-XXII-069.1}\quad p. 435; misprint\_in\_print.\par
\noindent Imprimé : {\ttfamily\footnotesize \textbackslash{}mathcal\allowbreak{} O\allowbreak{}}\par
\noindent Corrigé : {\ttfamily\footnotesize 0\allowbreak{}}\par
\noindent Source : FIX-7II-D-0569. \nolinkurl{fixes/FIXLOG_7II.jsonl#569}\par

\par\medskip\noindent\hypertarget{app-F7II-XXII-070.1}{}\textbf{F7II-XXII-070.1}\quad p. 435; misprint\_in\_print.\par
\noindent Imprimé : {\ttfamily\footnotesize \textbackslash{}mathcal\allowbreak{} O\allowbreak{}}\par
\noindent Corrigé : {\ttfamily\footnotesize 0\allowbreak{}}\par
\noindent Source : FIX-7II-D-0570. \nolinkurl{fixes/FIXLOG_7II.jsonl#570}\par

\par\medskip\noindent\hypertarget{app-F7II-XXII-070.2}{}\textbf{F7II-XXII-070.2}\quad p. 435; misprint\_in\_print.\par
\noindent Imprimé : {\ttfamily\footnotesize \textbackslash{}mathcal\allowbreak{} O\allowbreak{}}\par
\noindent Corrigé : {\ttfamily\footnotesize 0\allowbreak{}}\par
\noindent Source : FIX-7II-D-0571. \nolinkurl{fixes/FIXLOG_7II.jsonl#571}\par

\par\medskip\noindent\hypertarget{app-F7II-XXII-071.1}{}\textbf{F7II-XXII-071.1}\quad p. 436; misprint\_in\_print.\par
\noindent Imprimé : {\ttfamily\footnotesize associé\allowbreak{}}\par
\noindent Corrigé : {\ttfamily\footnotesize associée\allowbreak{}}\par
\noindent Source : FIX-7II-D-0575. \nolinkurl{fixes/FIXLOG_7II.jsonl#575}\par

\par\medskip\noindent\hypertarget{app-F7II-XXII-072.1}{}\textbf{F7II-XXII-072.1}\quad p. 436; misprint\_in\_print.\par
\noindent Imprimé : {\ttfamily\footnotesize .\allowbreak{}}\par
\noindent Corrigé : {\ttfamily\footnotesize )\allowbreak{}}\par
\noindent Source : FIX-7II-D-0579. \nolinkurl{fixes/FIXLOG_7II.jsonl#579}\par

\par\medskip\noindent\hypertarget{app-F7II-XXII-073.1}{}\textbf{F7II-XXII-073.1}\quad p. 436; misprint\_in\_print.\par
\noindent Imprimé : {\ttfamily\footnotesize .\allowbreak{}}\par
\noindent Corrigé : {\ttfamily\footnotesize )\allowbreak{}}\par
\noindent Source : FIX-7II-D-0581. \nolinkurl{fixes/FIXLOG_7II.jsonl#581}\par

\par\medskip\noindent\hypertarget{app-F7II-XXII-073.2}{}\textbf{F7II-XXII-073.2}\quad p. 436; misprint\_in\_print.\par
\noindent Imprimé : {\ttfamily\footnotesize (\allowbreak{}}\par
\noindent Corrigé : \emph{(vide)}\par
\noindent Source : FIX-7II-D-0583. \nolinkurl{fixes/FIXLOG_7II.jsonl#583}\par

\par\medskip\noindent\hypertarget{app-F7II-XXII-074.1}{}\textbf{F7II-XXII-074.1}\quad p. 436; misprint\_in\_print.\par
\noindent Imprimé : {\ttfamily\footnotesize .\allowbreak{}}\par
\noindent Corrigé : {\ttfamily\footnotesize ,\allowbreak{}}\par
\noindent Source : FIX-7II-D-0587. \nolinkurl{fixes/FIXLOG_7II.jsonl#587}\par

\par\medskip\noindent\hypertarget{app-F7II-XXII-074.2}{}\textbf{F7II-XXII-074.2}\quad p. 436; misprint\_in\_print.\par
\noindent Imprimé : \emph{(vide)}\par
\noindent Corrigé : {\ttfamily\footnotesize )\allowbreak{}}\par
\noindent Source : FIX-7II-D-0588. \nolinkurl{fixes/FIXLOG_7II.jsonl#588}\par

\par\medskip\noindent\hypertarget{app-F7II-XXII-075.1}{}\textbf{F7II-XXII-075.1}\quad p. 437; printed misprint.\par
\noindent Imprimé : {\ttfamily\footnotesize i\allowbreak{}}\par
\noindent Corrigé : {\ttfamily\footnotesize n\allowbreak{}}\par
\noindent Source : 7II:760. \nolinkurl{fixes/MISPRINTS_7II.json#760}\par

\par\medskip\noindent\hypertarget{app-F7II-XXII-075.2}{}\textbf{F7II-XXII-075.2}\quad p. 437; designated-witness attribution.\par
\noindent Imprimé : {\ttfamily\footnotesize =\allowbreak{}}\par
\noindent Corrigé : {\ttfamily\footnotesize \textbackslash{}mathrel\{\textbackslash{}big/\}\allowbreak{}}\par
\noindent Source : WITNESS-F7II-XXII-075.2. \nolinkurl{fr/sga7ii/diplomatic/expose_XXII_body.tex}; \nolinkurl{fr/sga7ii/current/expose_XXII_body.tex}\par

\par\medskip\noindent\hypertarget{app-F7II-XXII-076.1}{}\textbf{F7II-XXII-076.1}\quad p. 437; misprint\_in\_print.\par
\noindent Imprimé : {\ttfamily\footnotesize 1\allowbreak{}}\par
\noindent Corrigé : {\ttfamily\footnotesize i\allowbreak{}}\par
\noindent Source : FIX-7II-D-0598. \nolinkurl{fixes/FIXLOG_7II.jsonl#598}\par

\par\medskip\noindent\hypertarget{app-F7II-XXII-076.2}{}\textbf{F7II-XXII-076.2}\quad p. 437; misprint\_in\_print.\par
\noindent Imprimé : {\ttfamily\footnotesize f\allowbreak{}}\par
\noindent Corrigé : {\ttfamily\footnotesize \textbackslash{}mathfrak\allowbreak{} F\allowbreak{}}\par
\noindent Source : FIX-7II-D-0602. \nolinkurl{fixes/FIXLOG_7II.jsonl#602}\par

\par\medskip\noindent\hypertarget{app-F7II-XXII-077.1}{}\textbf{F7II-XXII-077.1}\quad p. 438; misprint\_in\_print.\par
\noindent Imprimé : {\ttfamily\footnotesize Zéta\allowbreak{}}\par
\noindent Corrigé : {\ttfamily\footnotesize Zeta\allowbreak{}}\par
\noindent Source : FIX-7II-D-0607. \nolinkurl{fixes/FIXLOG_7II.jsonl#607}\par

\par\medskip\noindent\hypertarget{app-F7II-XXII-077.2}{}\textbf{F7II-XXII-077.2}\quad p. 438; designated-witness attribution.\par
\noindent Imprimé : {\ttfamily\footnotesize Complétement,\allowbreak{}}\par
\noindent Corrigé : {\ttfamily\footnotesize Complètement\allowbreak{}}\par
\noindent Source : WITNESS-F7II-XXII-077.2. \nolinkurl{fr/sga7ii/diplomatic/expose_XXII_body.tex}; \nolinkurl{fr/sga7ii/current/expose_XXII_body.tex}\par

\normalsize\fi

\ifnum\SGAform=1
\clearpage
\clearpage

\section*{Lectures conservées et questions éditoriales non tranchées}

Ces notes présentent des alternatives qui ne sont pas insérées comme corrections nouvelles. Elles sont distinctes de l’apparat numéroté : une question sans différence entre les formes diplomatique et corrigée ne crée aucune entrée de variante. « Needs the scan » signifie que la lecture matérielle exige le scan, absent de cette construction. La couche diplomatique fournie reste le témoin de la lecture imprimée.

\subsection*{ED-G27-x-hat — 7I:123}

Lecture de la source conservée : \textbackslash{}mathfrak X\_\{(\textbackslash{}hat x)\} in the supplied source layers

Alternative : earlier transcription \textbackslash{}mathcal X\_\{(\textbackslash{}mathfrak x)\}

Conserver le chapeau des deux couches fournies. Le commentaire antérieur hésitait sur la forme physique du caractère surfrappé ; seul le scan permettrait de vérifier cette question paléographique.

Statut : needs the scan. Source : \nolinkurl{conventions/GLOBAL_ISSUES_SGA.md:27}

\subsection*{ED-G27-missing-reference — 7I:128}

Lecture de la source conservée : is smooth by and therefore

Alternative : is smooth by 4.13(a) and therefore

Conserver la référence laissée vide dans le texte imprimé. Plusieurs résultats de lissité peuvent intervenir ; aucun témoin fourni ne donne le numéro manquant. La proposition 4.13(a) reste une suggestion éditoriale.

Statut : retained; authorial reference unresolved. Source : \nolinkurl{conventions/GLOBAL_ISSUES_SGA.md:27}

\subsection*{ED-G32-run-in-headings — 7I:163,165 (and analogous numbered headings)}

Lecture de la source conservée : Attached layers use separate paragraphs after numbered headings

Alternative : possible printed run-in headings with following prose on the same paragraph

Conserver les alinéas des couches fournies. Les corrections de mise en ligne ont été retenues à faible confiance et les autres occurrences n’ont pas été vérifiées. Le scan est nécessaire pour établir les retraits et les titres courants.

Statut : needs the scan. Source : \nolinkurl{conventions/GLOBAL_ISSUES_SGA.md:32}

\subsection*{ED-G36-E-or-X — 7I:180}

Lecture de la source conservée : d\_1\textasciicircum{}*d\_0\textasciicircum{}*(E)\textbackslash{}simeq(d\_0d\_1)\textasciicircum{}*(E)

Alternative : d\_1\textasciicircum{}*d\_0\textasciicircum{}*(X)\textbackslash{}simeq(d\_0d\_1)\textasciicircum{}*(X)

Conserver E. L’extension considérée juste avant est X, mais l’isomorphisme de transitivité est valable pour une extension quelconque E ; une notation générique de l’auteur reste possible.

Statut : retained; no determinate emendation. Source : \nolinkurl{conventions/GLOBAL_ISSUES_SGA.md:36}

\subsection*{ED-G36-psi-footnote — 7I:180}

Lecture de la source conservée : \textbackslash{}psi\_\{L.\}(A)

Alternative : \textbackslash{}underline\textbackslash{}Psi\_\{L.\}(A)

Conserver le psi minuscule de la couche diplomatique. Le texte principal emploie un Psi majuscule souligné, mais le dossier ne permet pas de distinguer une variante de glyphe d’une coquille. Le scan est nécessaire.

Statut : needs the scan. Source : \nolinkurl{conventions/GLOBAL_ISSUES_SGA.md:36}

\subsection*{ED-G37-duplicate-number — 7I:183}

Lecture de la source conservée : two displays tagged (3.2.3)

Alternative : second display (3.2.3\textbackslash{}prime) or another new tag

Conserver les deux numéros (3.2.3). Le numéro (3.2.4) est déjà employé ensuite ; aucun numéro voulu par l’auteur ne peut être rétabli avec certitude. Le diagramme et le morphisme c(X,alpha) sont distingués par leur contenu.

Statut : retained; intended numbering unresolved. Source : \nolinkurl{conventions/GLOBAL_ISSUES_SGA.md:37}

\subsection*{ED-G38-speck — 7II:202}

Lecture de la source conservée : (\textbackslash{}mathrm\{Sw\}\_\{I\textbackslash{}prime\},s\_\{I/H\}) without an extra lower dot

Alternative : a possible printed dot below the opening parenthesis

Ne pas ajouter de point sous la parenthèse. Aucune opération mathématique n’est définie par cette marque possible ; seul le scan peut distinguer ponctuation, encre et défaut matériel.

Statut : needs the scan. Source : \nolinkurl{conventions/GLOBAL_ISSUES_SGA.md:38}

\subsection*{ED-G41-exton-sion — 7I:195}

Lecture de la source conservée : extension (both attached layers)

Alternative : exton-sion (held low-confidence visual reading)

Conserver extension, lecture commune des deux couches fournies. La lecture exton-sion a été retenue à faible confiance ; le scan permettrait de distinguer un o d’un e endommagé et de vérifier le trait de coupure de fin de ligne.

Statut : needs the scan. Source : \nolinkurl{conventions/GLOBAL_ISSUES_SGA.md:41}

\subsection*{ED-G45-mathematical-repair — 7II:237}

Lecture de la source conservée : the printed pencil in (4.2.3)

Alternative : a proposed replacement pencil with an isolated ordinary quadratic singularity

Conserver le pinceau fourni dans (4.2.3). Le registre signale, pour n=2k+1 au moins 7, un plan singulier au lieu d’un point double ordinaire ; il distingue explicitement une réparation mathématique d’une lecture récupérée. La formule de remplacement détaillée n’est pas jointe et ne peut être inventée.

Statut : retained; proposed mathematical replacement unavailable. Source : \nolinkurl{conventions/GLOBAL_ISSUES_SGA.md:45}

\subsection*{ED-G48-inner-underlines — 7II:248 and other internal commas}

Lecture de la source conservée : source-layer continuous underlines through some inner commas

Alternative : stop and restart each underline around the comma

Conserver le soulignement fourni à chaque occurrence. Le registre précise que le dactylographe passait parfois sous les virgules et parfois non ; aucune substitution globale n’est justifiée sans le scan de chaque lieu.

Statut : needs the scan. Source : \nolinkurl{conventions/GLOBAL_ISSUES_SGA.md:48}

\subsection*{ED-G58-comma — 7II:332}

Lecture de la source conservée : H\_s,

Alternative : H\_\{s,\}

Conserver la virgule comme séparation des arguments, après H\_s. La position basse observée à faible confiance ne définit pas un nouvel indice. Le scan est nécessaire pour examiner la ligne de base physique.

Statut : needs the scan. Source : \nolinkurl{conventions/GLOBAL_ISSUES_SGA.md:58}

\subsection*{ED-G63-miracle — 7II:416}

Lecture de la source conservée : oh miracle (both attached layers)

Alternative : ô miracle (earlier held reading)

Conserver oh miracle, lecture commune des deux couches fournies. L’observation antérieure ô miracle n’avait pas été établie ; seule la consultation du scan permettrait de trancher indépendamment.

Statut : needs the scan. Source : \nolinkurl{conventions/GLOBAL_ISSUES_SGA.md:63}

\subsection*{ED-G64-quotient — 7II:423}

Lecture de la source conservée : the displayed fraction in (4.5.21)

Alternative : a slash presentation of the same quotient

Conserver la fraction de (4.5.21). Les deux couches donnent les mêmes numérateur et dénominateur ; la barre oblique proposée concerne uniquement la disposition graphique du quotient. Le scan est nécessaire pour établir cette disposition.

Statut : needs the scan. Source : \nolinkurl{conventions/GLOBAL_ISSUES_SGA.md:64}

\subsection*{ED-G64-subset-stroke — 7II:424}

Lecture de la source conservée : \textbackslash{}subseteq in the supplied source

Alternative : \textbackslash{}subset without the lower stroke

Conserver le signe subseteq de la couche fournie. Le trait inférieur peut relever d’une surfrappe ; la convention historique de subset peut aussi être non stricte. Le sens seul ne décide donc pas de la forme physique.

Statut : needs the scan. Source : \nolinkurl{conventions/GLOBAL_ISSUES_SGA.md:64}

\subsection*{ED-G64-Delta-Lambda — 7II:424–429}

Lecture de la source conservée : \textbackslash{}Delta\_B throughout the defining discussion; \textbackslash{}Lambda\_B in printed (5.0.12)

Alternative : change all seven \textbackslash{}Delta\_B to \textbackslash{}Lambda\_B, or retain current \textbackslash{}Delta\_B at (5.0.12)

Conserver les glyphes de la couche diplomatique et la correction locale Delta de la couche courante. Delta\_B est défini en (4.5.28) et la formule (5.0.12) renvoie à la même série caractéristique. Le remplacement global des sept occurrences demanderait leurs scans et n’est pas effectué.

Statut : needs the scan. Source : \nolinkurl{conventions/GLOBAL_ISSUES_SGA.md:64}

\subsection*{ED-G67-Dwork-psi — 7II:436–438}

Lecture de la source conservée : \textbackslash{}psi\_q in the attached layers

Alternative : \textbackslash{}Psi\_q suggested by the typewriter capital-form glyph

Conserver psi pour l’opérateur de Dwork déjà défini dans l’exposé. La hauteur de capitale du glyphe dactylographié ne suffit pas à établir un second opérateur ni une distinction de casse. La convention 78 exclut les simples variantes de forme du texte diplomatique.

Statut : needs the scan. Source : \nolinkurl{conventions/GLOBAL_ISSUES_SGA.md:67}

\subsection*{ED-G73-statement-underlines — 7II:403,405 and similar statements}

Lecture de la source conservée : source underlines sometimes include mathematical symbols and final punctuation

Alternative : split each underline around symbols and punctuation according to the print

Conserver le balisage fourni, sauf les réparations précises documentées séparément. Le registre renvoie ces soulignements à la vérification visuelle au cas par cas de la convention 48 ; aucune règle globale ne permet de retrouver le geste du dactylographe.

Statut : needs the scan. Source : \nolinkurl{conventions/GLOBAL_ISSUES_SGA.md:73}

\fi

\end{document}
